\IfFileExists{stacks-project-book.cls}{%
\documentclass{stacks-project-book}
}{%
\documentclass{amsbook}
}

\usepackage{verbatim}
\newenvironment{reference}{\comment}{\endcomment}
\newenvironment{slogan}{\comment}{\endcomment}
\newenvironment{history}{\comment}{\endcomment}

\usepackage[all]{xy}

\xyoption{2cell}
\UseAllTwocells

\usepackage{multicol}

\usepackage{lmodern}
\usepackage[T1]{fontenc}
\usepackage[french]{babel}
\AtBeginDocument{\shorthandoff{?}}


\usepackage{hyperref}


\theoremstyle{plain}
\newtheorem{theorem}[subsection]{Théorème}
\newtheorem{proposition}[subsection]{Proposition}
\newtheorem{lemma}[subsection]{Lemme}

\theoremstyle{definition}
\newtheorem{definition}[subsection]{Définition}
\newtheorem{example}[subsection]{Exemple}
\newtheorem{exercise}[subsection]{Exercice}
\newtheorem{situation}[subsection]{Situation}

\theoremstyle{remark}
\newtheorem{remark}[subsection]{Remarque}
\newtheorem{remarks}[subsection]{Remarques}

\renewcommand{\proofname}{Démonstration}

\numberwithin{equation}{subsection}

\def\lim{\mathop{\mathrm{lim}}\nolimits}
\def\colim{\mathop{\mathrm{colim}}\nolimits}
\def\Spec{\mathop{\mathrm{Spec}}}
\def\Hom{\mathop{\mathrm{Hom}}\nolimits}
\def\Ext{\mathop{\mathrm{Ext}}\nolimits}
\def\SheafHom{\mathop{\mathcal{H}\!\mathit{om}}\nolimits}
\def\SheafExt{\mathop{\mathcal{E}\!\mathit{xt}}\nolimits}
\def\Sch{\mathit{Sch}}
\def\Mor{\mathop{\mathrm{Mor}}\nolimits}
\def\Ob{\mathop{\mathrm{Ob}}\nolimits}
\def\Sh{\mathop{\mathit{Sh}}\nolimits}
\def\NL{\mathop{N\!L}\nolimits}
\def\CH{\mathop{\mathrm{CH}}\nolimits}
\def\proetale{{pro\text{-}\acute{e}tale}}
\def\etale{{\acute{e}tale}}
\def\QCoh{\mathit{QCoh}}
\def\Ker{\mathop{\mathrm{Ker}}}
\def\Im{\mathop{\mathrm{Im}}}
\def\Coker{\mathop{\mathrm{Coker}}}
\def\Coim{\mathop{\mathrm{Coim}}}

\DeclareMathSymbol{\boxtimes}{\mathbin}{AMSa}{"02}

\def\QCohstack{\mathcal{QC}\!\mathit{oh}}
\def\Cohstack{\mathcal{C}\!\mathit{oh}}
\def\Spacesstack{\mathcal{S}\!\mathit{paces}}
\def\Quotfunctor{\mathrm{Quot}}
\def\Hilbfunctor{\mathrm{Hilb}}
\def\Curvesstack{\mathcal{C}\!\mathit{urves}}
\def\Polarizedstack{\mathcal{P}\!\mathit{olarized}}
\def\Complexesstack{\mathcal{C}\!\mathit{omplexes}}
\def\Pic{\mathop{\mathrm{Pic}}\nolimits}
\def\Picardstack{\mathcal{P}\!\mathit{ic}}
\def\Picardfunctor{\mathrm{Pic}}
\def\Deformationcategory{\mathcal{D}\!\mathit{ef}}
\begin{document}
\title{Le projet Stacks --- édition française, partie 5}
\maketitle
\tableofcontents

% OK, start here.
%

\chapter{Topologies sur les schémas}


\phantomsection
\label{topologies-section-phantom}




\section{Introduction}
\label{topologies-section-introduction}

\noindent
Dans ce chapitre, nous décrivons les différentes topologies sur la
catégorie des schémas. On pourra consulter, entre autres, \cite{SGA1} et \cite{Ner}.
Avant de le faire, nous souhaitons signaler qu'il existe de nombreux
choix de sites (au sens de
Sites, définition \ref{sites-definition-site}) qui donnent lieu à
la même notion de faisceau sur la catégorie sous-jacente. Ainsi,
nos choix peuvent différer légèrement de ceux des références,
mais conduisent finalement aux mêmes groupes de cohomologie, etc.

\section{La méthode générale}
\label{topologies-section-procedure}

\noindent
Dans cette section, nous exposons une méthode générale de construction des
sites avec lesquels nous travaillerons. Supposons que l'on veuille étudier les faisceaux
sur les schémas relativement à une topologie $\tau$. Pour obtenir
un site, comme dans Sites, définition \ref{sites-definition-site},
formé de schémas munis de cette topologie, il faut prendre quelques précautions. En effet,
on ne peut pas simplement dire « considérons tous les schémas munis de la topologie de Zariski »,
car on obtiendrait ainsi une catégorie « grande ». Dans chaque section de
ce chapitre, nous procéderons plutôt comme suit :
\begin{enumerate}
\item On définit une classe $\text{Cov}_\tau$ de recouvrements de schémas
satisfaisant aux axiomes de Sites, définition \ref{sites-definition-site}.
Un recouvrement ouvert de Zariski d'un
schéma sera toujours un recouvrement pour $\tau$.
\item On distingue, dans la catégorie des schémas affines, une notion de
$\tau$-recouvrement standard.
\item On définit ce qu'est un gros $\tau$-site « absolu » $\Sch_\tau$.
Il s'agit des sites obtenus en choisissant convenablement un ensemble de schémas
et un ensemble de recouvrements.
\item Pour tout objet $S$ de $\Sch_\tau$,
on définit le gros $\tau$-site $(\Sch/S)_\tau$ et, lorsque
$\tau$ s'y prête, le petit\footnote{Les mots gros et
petit ne renvoient pas ici à la taille des catégories
correspondantes.} $\tau$-site $S_\tau$.
\item En outre, il existe un site $(\textit{Aff}/S)_\tau$ fondé sur la
notion de $\tau$-recouvrement standard d'affines\footnote{Dans le cas
de la topologie ph, nous nous écartons très légèrement de cette méthode ; voir la
définition \ref{topologies-definition-big-small-ph} et la discussion qui l'entoure.}
dont la catégorie des faisceaux
est équivalente à la catégorie des faisceaux sur $(\Sch/S)_\tau$.
\end{enumerate}
Cette construction est quelque peu malcommode, car elle ne fournit pas de choix canonique
du gros $\tau$-site d'un schéma, ni même du petit
$\tau$-site d'un schéma. Si l'on accepte d'ignorer les difficultés
ensemblistes, on peut travailler avec des classes et obtenir
des gros et petits sites canoniques…







\section{La topologie de Zariski}
\label{topologies-section-zariski}

\begin{definition}
\label{topologies-definition-zariski-covering}
Soit $T$ un schéma. Un {\it recouvrement de Zariski de $T$} est une famille
de morphismes de schémas $\{f_i : T_i \to T\}_{i \in I}$
telle que chaque $f_i$ soit une immersion ouverte et que
$T = \bigcup f_i(T_i)$.
\end{definition}

\noindent
Cela définit une classe (propre) de recouvrements.
Montrons à présent que cette notion satisfait aux conditions de
Sites, définition \ref{sites-definition-site}.

\begin{lemma}
\label{topologies-lemma-zariski}
Soit $T$ un schéma.
\begin{enumerate}
\item Si $T' \to T$ est un isomorphisme, alors $\{T' \to T\}$
est un recouvrement de Zariski de $T$.
\item Si $\{T_i \to T\}_{i\in I}$ est un recouvrement de Zariski et si, pour tout
$i$, on dispose d'un recouvrement de Zariski $\{T_{ij} \to T_i\}_{j\in J_i}$, alors
$\{T_{ij} \to T\}_{i \in I, j\in J_i}$ est un recouvrement de Zariski.
\item Si $\{T_i \to T\}_{i\in I}$ est un recouvrement de Zariski
et si $T' \to T$ est un morphisme de schémas, alors
$\{T' \times_T T_i \to T'\}_{i\in I}$ est un recouvrement de Zariski.
\end{enumerate}
\end{lemma}

\begin{proof}
Omis.
\end{proof}

\begin{lemma}
\label{topologies-lemma-zariski-affine}
Soit $T$ un schéma affine. Soit $\{T_i \to T\}_{i \in I}$ un
recouvrement de Zariski de $T$. Il existe alors un recouvrement de Zariski
$\{U_j \to T\}_{j = 1, \ldots, m}$ qui raffine
$\{T_i \to T\}_{i \in I}$ et tel que chaque $U_j$ soit un ouvert
principal de $T$ ; voir
Schémas, définition \ref{schemes-definition-standard-covering}.
On peut en outre choisir chaque $U_j$ comme ouvert de l'un des $T_i$.
\end{lemma}

\begin{proof}
Cela résulte de ce que $T$ est quasi-compact et que les ouverts principaux forment une base
de sa topologie. C'est également démontré dans
Schémas, lemme \ref{schemes-lemma-standard-open}.
\end{proof}

\noindent
On définit donc comme suit les recouvrements standards d'affines correspondants.

\begin{definition}
\label{topologies-definition-standard-Zariski}
Comparer à Schémas, définition \ref{schemes-definition-standard-covering}.
Soit $T$ un schéma affine. Un {\it recouvrement de Zariski standard}
de $T$ est un recouvrement de Zariski $\{U_j \to T\}_{j = 1, \ldots, m}$
tel que chaque $U_j \to T$ induise un isomorphisme sur un ouvert affine principal
de $T$.
\end{definition}

\begin{definition}
\label{topologies-definition-big-zariski-site}
Un {\it gros site de Zariski} est un site $\Sch_{Zar}$, au sens de
Sites, définition \ref{sites-definition-site}, construit comme suit :
\begin{enumerate}
\item Choisir un ensemble quelconque de schémas $S_0$, ainsi qu'un ensemble de recouvrements de Zariski
$\text{Cov}_0$ entre ces schémas.
\item Prendre pour catégorie sous-jacente à $\Sch_{Zar}$
une catégorie $\Sch_\alpha$ construite comme dans
Ensembles, lemme \ref{sets-lemma-construct-category}, à partir de l'ensemble $S_0$.
\item Prendre pour recouvrements de $\Sch_{Zar}$ un ensemble de recouvrements obtenu comme dans
Ensembles, lemme \ref{sets-lemma-coverings-site}, à partir de la
catégorie $\Sch_\alpha$, de la classe des recouvrements de Zariski
et de l'ensemble $\text{Cov}_0$ choisi ci-dessus.
\end{enumerate}
\end{definition}

\noindent
Sites, lemme \ref{sites-lemma-choice-set-coverings-immaterial}, montre
qu'une fois la catégorie $\Sch_\alpha$ choisie, la
catégorie des faisceaux sur $\Sch_\alpha$ ne dépend pas du
choix des recouvrements effectué au point (3) ci-dessus. Autrement dit, le topos
$\Sh(\Sch_{Zar})$ ne dépend que du choix de
la catégorie $\Sch_\alpha$. Il est démontré dans
Ensembles, lemme \ref{sets-lemma-what-is-in-it}, que ces catégories
sont stables par de nombreuses constructions de géométrie algébrique, par exemple
les produits fibrés et le passage aux sous-schémas ouverts ou fermés. On peut également montrer
que le choix précis de $\Sch_\alpha$ importe
assez peu ; voir la section \ref{topologies-section-change-alpha}.

\medskip\noindent
Une autre méthode consisterait à supposer l'existence d'un
cardinal fortement inaccessible et à définir $\Sch_{Zar}$
comme la catégorie des schémas contenus dans un univers choisi, les
recouvrements étant les recouvrements de Zariski contenus dans ce même
univers.

\medskip\noindent
Avant de poursuivre par l'introduction du gros site de Zariski d'un
schéma $S$, observons que la topologie d'un gros site de Zariski
$\Sch_{Zar}$ est, en un certain sens, induite par la topologie de Zariski
sur la catégorie de tous les schémas.

\begin{lemma}
\label{topologies-lemma-zariski-induced}
Soit $\Sch_{Zar}$ un gros site de Zariski comme dans la
définition \ref{topologies-definition-big-zariski-site}.
Soit $T \in \Ob(\Sch_{Zar})$.
Soit $\{T_i \to T\}_{i \in I}$ un recouvrement de Zariski quelconque de $T$.
Il existe un recouvrement $\{U_j \to T\}_{j \in J}$ de $T$ dans le site
$\Sch_{Zar}$ qui soit tautologiquement équivalent (voir
Sites, définition \ref{sites-definition-combinatorial-tautological})
à $\{T_i \to T\}_{i \in I}$.
\end{lemma}

\begin{proof}
Puisque chaque $T_i \to T$ est une immersion ouverte, il résulte de
Ensembles, lemme \ref{sets-lemma-what-is-in-it},
que chaque $T_i$ est isomorphe à un objet $V_i$ de $\Sch_{Zar}$.
Le recouvrement $\{V_i \to T\}_{i \in I}$ est tautologiquement équivalent
à $\{T_i \to T\}_{i \in I}$ (en utilisant dans les deux sens l'application identité de $I$).
En outre, $\{V_i \to T\}_{i \in I}$ est combinatoirement équivalent à un
recouvrement $\{U_j \to T\}_{j \in J}$ de $T$ dans le site $\Sch_{Zar}$ d'après
Ensembles, lemme \ref{sets-lemma-coverings-site}.
\end{proof}

\begin{definition}
\label{topologies-definition-big-small-Zariski}
Soit $S$ un schéma. Soit $\Sch_{Zar}$ un gros site de Zariski
contenant $S$.
\begin{enumerate}
\item Le {\it gros site de Zariski de $S$}, noté
$(\Sch/S)_{Zar}$, est le site $\Sch_{Zar}/S$
introduit dans Sites, section \ref{sites-section-localize}.
\item Le {\it petit site de Zariski de $S$}, que nous notons
$S_{Zar}$, est la sous-catégorie pleine de $(\Sch/S)_{Zar}$
dont les objets sont les $U/S$ tels que $U \to S$ soit une immersion ouverte.
Un recouvrement de $S_{Zar}$ est tout recouvrement $\{U_i \to U\}$ de
$(\Sch/S)_{Zar}$ avec $U \in \Ob(S_{Zar})$.
\item Le {\it gros site de Zariski affine de $S$}, noté
$(\textit{Aff}/S)_{Zar}$, est la sous-catégorie pleine de
$(\Sch/S)_{Zar}$ formée des objets $U/S$ tels que $U$ soit un
schéma affine. Un recouvrement de $(\textit{Aff}/S)_{Zar}$ est tout recouvrement
$\{U_i \to U\}$ de $(\Sch/S)_{Zar}$ avec $U \in \Ob((\textit{Aff}/S)_{Zar})$
qui soit un recouvrement de Zariski standard.
\item Le {\it petit site de Zariski affine de $S$}, noté
$S_{affine, Zar}$, est la sous-catégorie pleine de $S_{Zar}$
dont les objets sont les $U/S$ tels que $U$ soit un schéma affine.
Un recouvrement de $S_{affine, Zar}$ est tout recouvrement $\{U_i \to U\}$ de
$S_{Zar}$ avec $U \in \Ob(S_{affine, Zar})$
qui soit un recouvrement de Zariski standard.
\end{enumerate}
\end{definition}

\noindent
Il n'est pas tout à fait évident que le petit site de Zariski,
le gros site de Zariski affine et le petit site de Zariski affine
soient bien des sites. Vérifions-le.

\begin{lemma}
\label{topologies-lemma-verify-site-Zariski}
Soit $S$ un schéma. Soit $\Sch_{Zar}$ un gros site de Zariski
contenant $S$. Les structures $S_{Zar}$, $(\textit{Aff}/S)_{Zar}$
et $S_{affine, Zar}$ définies ci-dessus sont des sites.
\end{lemma}

\begin{proof}
Montrons que $S_{Zar}$ est un site. C'est une catégorie munie d'un
ensemble donné de familles de morphismes de but fixé. Il faut donc
établir les propriétés (1), (2) et (3) de
Sites, définition \ref{sites-definition-site}.
Puisque $(\Sch/S)_{Zar}$ est un site, il suffit de prouver
que, pour tout recouvrement $\{U_i \to U\}$ de $(\Sch/S)_{Zar}$
avec $U \in \Ob(S_{Zar})$, on a aussi $U_i \in \Ob(S_{Zar})$.
Cela résulte des définitions,
car la composée de deux immersions ouvertes est une immersion ouverte.

\medskip\noindent
Montrons que $(\textit{Aff}/S)_{Zar}$ est un site.
En raisonnant comme ci-dessus, il suffit de montrer que l'ensemble
des recouvrements de Zariski standards d'affines satisfait aux propriétés
(1), (2) et (3) de
Sites, définition \ref{sites-definition-site}.
Soit $R$ un anneau. Soient $f_1, \ldots, f_n \in R$ engendrant l'idéal unité.
Pour chaque $i \in \{1, \ldots, n\}$, soient $g_{i1}, \ldots, g_{in_i} \in R_{f_i}$
des éléments engendrant l'idéal unité de $R_{f_i}$. On peut écrire
$g_{ij} = f_{ij}/f_i^{e_{ij}}$. Après avoir remplacé, si nécessaire,
$f_{ij}$ par $f_i f_{ij}$, on a
$D(f_{ij}) \subset D(f_i) \cong \Spec(R_{f_i})$
égal à $D(g_{ij}) \subset \Spec(R_{f_i})$. On voit donc que
la famille de morphismes $\{D(g_{ij}) \to \Spec(R)\}$
est un recouvrement de Zariski standard. Ces considérations
montrent que (2) vaut pour les recouvrements de Zariski standards.
Nous omettons la vérification de (1) et de (3).

\medskip\noindent
Nous omettons la démonstration du fait que $S_{affine, Zar}$ est un site.
\end{proof}

\begin{lemma}
\label{topologies-lemma-fibre-products-Zariski}
Soit $S$ un schéma. Soit $\Sch_{Zar}$ un gros site de Zariski
contenant $S$. Les catégories sous-jacentes aux sites
$\Sch_{Zar}$, $(\Sch/S)_{Zar}$,
$S_{Zar}$, $(\textit{Aff}/S)_{Zar}$ et $S_{affine, Zar}$ possèdent des produits fibrés.
Dans chaque cas, le foncteur évident vers la catégorie $\Sch$ de
tous les schémas commute aux produits fibrés. Les catégories
$(\Sch/S)_{Zar}$ et $S_{Zar}$ possèdent toutes deux un objet final,
à savoir $S/S$.
\end{lemma}

\begin{proof}
Pour $\Sch_{Zar}$, c'est vrai par construction ; voir
Ensembles, lemme \ref{sets-lemma-what-is-in-it}.
Supposons donnés des morphismes de schémas $U \to S$, $V \to U$, $W \to U$
avec $U, V, W \in \Ob(\Sch_{Zar})$.
Le produit fibré $V \times_U W$ dans $\Sch_{Zar}$
est un produit fibré dans $\Sch$ ; c'est aussi
le produit fibré de $V/S$ et $W/S$ au-dessus de $U/S$ dans
la catégorie de tous les schémas au-dessus de $S$, et donc également un
produit fibré dans $(\Sch/S)_{Zar}$.
Cela démontre le résultat pour $(\Sch/S)_{Zar}$.
Si $U \to S$, $V \to U$ et $W \to U$ sont des immersions ouvertes, alors
$V \times_U W \to S$ l'est aussi, d'où le résultat pour $S_{Zar}$.
Si $U, V, W$ sont affines, $V \times_U W$ l'est également, d'où le
résultat pour $(\textit{Aff}/S)_{Zar}$ et $S_{affine, Zar}$.
\end{proof}

\noindent
Vérifions à présent que le gros, resp.\ petit, site affine définit le même
topos que le gros, resp.\ petit, site.

\begin{lemma}
\label{topologies-lemma-affine-big-site-Zariski}
Soit $S$ un schéma. Soit $\Sch_{Zar}$ un gros site de Zariski
contenant $S$.
Le foncteur $(\textit{Aff}/S)_{Zar} \to (\Sch/S)_{Zar}$
est un foncteur cocontinu spécial. Il induit donc une équivalence
de topos de $\Sh((\textit{Aff}/S)_{Zar})$ vers
$\Sh((\Sch/S)_{Zar})$.
\end{lemma}

\begin{proof}
La notion de foncteur cocontinu spécial est introduite dans
Sites, définition \ref{sites-definition-special-cocontinuous-functor}.
Il faut donc vérifier les hypothèses (1) -- (5) de
Sites, lemme \ref{sites-lemma-equivalence}.
Notons le foncteur d'inclusion
$u : (\textit{Aff}/S)_{Zar} \to (\Sch/S)_{Zar}$.
La cocontinuité signifie simplement que tout recouvrement de Zariski de
$T/S$, avec $T$ affine, peut être raffiné par un recouvrement de Zariski standard de $T$.
C'est le contenu du
lemme \ref{topologies-lemma-zariski-affine}.
Ainsi, (1) est satisfaite. Le foncteur $u$ est continu, simplement parce qu'un recouvrement
de Zariski standard est un recouvrement de Zariski. Ainsi, (2) est satisfaite.
Les assertions (3) et (4) résultent immédiatement de ce que $u$ est
pleinement fidèle. Enfin, la condition (5) résulte de ce que
tout schéma possède un recouvrement ouvert affine.
\end{proof}

\begin{lemma}
\label{topologies-lemma-alternative-zariski}
Soit $S$ un schéma. Soit $\Sch_{Zar}$ un gros site de Zariski
contenant $S$. Le foncteur $S_{affine, Zar} \to S_{Zar}$
est un foncteur cocontinu spécial. Il induit donc une équivalence
de topos de $\Sh(S_{affine, Zar})$ vers $\Sh(S_{Zar})$.
\end{lemma}

\begin{proof}
Omis. Indication : comparer à la démonstration du
lemme \ref{topologies-lemma-affine-big-site-Zariski}.
\end{proof}

\noindent
Vérifions que la notion de faisceau sur le petit site de Zariski
correspond à celle de faisceau sur $S$.

\begin{lemma}
\label{topologies-lemma-Zariski-usual}
La catégorie des faisceaux sur $S_{Zar}$ est équivalente à la
catégorie des faisceaux sur l'espace topologique sous-jacent à $S$.
\end{lemma}

\begin{proof}
Nous utiliserons à plusieurs reprises le fait que, pour tout objet
$U/S$ de $S_{Zar}$, le morphisme $U \to S$ est un isomorphisme
sur un sous-schéma ouvert.
Soit $\mathcal{F}$ un faisceau sur $S$. On définit alors un faisceau
sur $S_{Zar}$ par la règle $\mathcal{F}'(U/S) = \mathcal{F}(\Im(U \to S))$.
Réciproquement, choisissons pour tout sous-schéma ouvert $U \subset S$ un objet
$U'/S \in \Ob(S_{Zar})$ tel que $\Im(U' \to S) = U$
(il faut ici utiliser Ensembles, lemme \ref{sets-lemma-what-is-in-it}).
Étant donné un faisceau $\mathcal{G}$ sur $S_{Zar}$, on définit un faisceau sur $S$ en posant
$\mathcal{G}'(U) = \mathcal{G}(U'/S)$. Pour voir que $\mathcal{G}'$ est
un faisceau, on utilise que, pour tout recouvrement ouvert $U = \bigcup_{i \in I} U_i$,
le recouvrement $\{U_i \to U\}_{i \in I}$
est combinatoirement équivalent à un recouvrement $\{U_j' \to U'\}_{j \in J}$
dans $S_{Zar}$ d'après Ensembles, lemme \ref{sets-lemma-coverings-site},
et l'on applique Sites, lemme \ref{sites-lemma-tautological-same-sheaf}.
Détails omis.
\end{proof}

\noindent
Désormais, nous ne ferons plus de distinction entre un faisceau sur
$S_{Zar}$ et un faisceau sur $S$. Nous utiliserons toujours les procédés
de la démonstration du lemme pour passer d'une notion à l'autre.
Établissons maintenant quelques relations entre les topos
associés à ces sites.

\begin{lemma}
\label{topologies-lemma-put-in-T}
Soit $\Sch_{Zar}$ un gros site de Zariski.
Soit $f : T \to S$ un morphisme dans $\Sch_{Zar}$.
Le foncteur $T_{Zar} \to (\Sch/S)_{Zar}$
est cocontinu et induit un morphisme de topos
$$
i_f :
\Sh(T_{Zar})
\longrightarrow
\Sh((\Sch/S)_{Zar})
$$
Pour tout faisceau $\mathcal{G}$ sur $(\Sch/S)_{Zar}$,
on a la formule $(i_f^{-1}\mathcal{G})(U/T) = \mathcal{G}(U/S)$.
Le foncteur $i_f^{-1}$ possède aussi un adjoint à gauche $i_{f, !}$ qui commute
aux produits fibrés et aux égalisateurs.
\end{lemma}

\begin{proof}
Notons le foncteur $u : T_{Zar} \to (\Sch/S)_{Zar}$.
Autrement dit, étant donnée une immersion ouverte $j : U \to T$ correspondant
à un objet de $T_{Zar}$, on pose $u(U \to T) = (f \circ j : U \to S)$.
Ce foncteur commute aux produits fibrés ; voir le
lemme \ref{topologies-lemma-fibre-products-Zariski}.
De plus, $T_{Zar}$ possède des égalisateurs (car deux morphismes de même
source et de même but sont égaux) et $u$ commute à ceux-ci.
Il est manifestement cocontinu.
Il est aussi continu, puisque $u$ transforme les recouvrements en recouvrements et
commute aux produits fibrés. Le lemme résulte donc de
Sites, lemmes \ref{sites-lemma-when-shriek}
et \ref{sites-lemma-preserve-equalizers}.
\end{proof}

\begin{lemma}
\label{topologies-lemma-at-the-bottom}
Soit $S$ un schéma. Soit $\Sch_{Zar}$ un gros site de Zariski
contenant $S$.
Le foncteur d'inclusion $S_{Zar} \to (\Sch/S)_{Zar}$
satisfait aux hypothèses de Sites, lemme \ref{sites-lemma-bigger-site},
et induit donc un morphisme de sites
$$
\pi_S : (\Sch/S)_{Zar} \longrightarrow S_{Zar}
$$
et un morphisme de topos
$$
i_S : \Sh(S_{Zar}) \longrightarrow \Sh((\Sch/S)_{Zar})
$$
tels que $\pi_S \circ i_S = \text{id}$. De plus, $i_S = i_{\text{id}_S}$,
où $i_{\text{id}_S}$ est celui du lemme \ref{topologies-lemma-put-in-T}. En particulier, le
foncteur $i_S^{-1} = \pi_{S, *}$ est décrit par la règle
$i_S^{-1}(\mathcal{G})(U/S) = \mathcal{G}(U/S)$.
\end{lemma}

\begin{proof}
Dans ce cas, le foncteur $u : S_{Zar} \to (\Sch/S)_{Zar}$,
outre les propriétés relevées dans la démonstration du
lemme \ref{topologies-lemma-put-in-T} ci-dessus, est pleinement fidèle
et transforme l'objet final en l'objet final.
Le lemme en résulte.
\end{proof}

\begin{definition}
\label{topologies-definition-restriction-small-zariski}
Dans la situation du
lemme \ref{topologies-lemma-at-the-bottom},
le foncteur $i_S^{-1} = \pi_{S, *}$ est souvent
appelé la {\it restriction au petit site de Zariski} et, pour un faisceau
$\mathcal{F}$ sur le gros site de Zariski, on note $\mathcal{F}|_{S_{Zar}}$
cette restriction.
\end{definition}

\noindent
Avec cette notation, pour un faisceau $\mathcal{F}$ sur le
gros site et un faisceau $\mathcal{G}$ sur le petit site, on a
\begin{align*}
\Mor_{\Sh(S_{Zar})}(\mathcal{F}|_{S_{Zar}}, \mathcal{G})
& =
\Mor_{\Sh((\Sch/S)_{Zar})}(\mathcal{F},
i_{S, *}\mathcal{G}) \\
\Mor_{\Sh(S_{Zar})}(\mathcal{G}, \mathcal{F}|_{S_{Zar}})
& =
\Mor_{\Sh((\Sch/S)_{Zar})}(\pi_S^{-1}\mathcal{G},
\mathcal{F})
\end{align*}
De plus, on a $(i_{S, *}\mathcal{G})|_{S_{Zar}} = \mathcal{G}$
et $(\pi_S^{-1}\mathcal{G})|_{S_{Zar}} = \mathcal{G}$.

\begin{lemma}
\label{topologies-lemma-morphism-big}
Soit $\Sch_{Zar}$ un gros site de Zariski.
Soit $f : T \to S$ un morphisme dans $\Sch_{Zar}$.
Le foncteur
$$
u : (\Sch/T)_{Zar} \longrightarrow (\Sch/S)_{Zar},
\quad
V/T \longmapsto V/S
$$
est cocontinu et possède un adjoint à droite continu
$$
v : (\Sch/S)_{Zar} \longrightarrow (\Sch/T)_{Zar},
\quad
(U \to S) \longmapsto (U \times_S T \to T).
$$
Ils induisent le même morphisme de topos
$$
f_{big} :
\Sh((\Sch/T)_{Zar})
\longrightarrow
\Sh((\Sch/S)_{Zar})
$$
On a $f_{big}^{-1}(\mathcal{G})(U/T) = \mathcal{G}(U/S)$.
On a $f_{big, *}(\mathcal{F})(U/S) = \mathcal{F}(U \times_S T/T)$.
De plus, $f_{big}^{-1}$ possède un adjoint à gauche $f_{big!}$ qui commute aux
produits fibrés et aux égalisateurs.
\end{lemma}

\begin{proof}
Le foncteur $u$ est cocontinu, continu et commute aux produits fibrés
et aux égalisateurs (détails omis ; comparer à la démonstration du
lemme \ref{topologies-lemma-put-in-T}).
Par conséquent,
Sites, lemmes \ref{sites-lemma-when-shriek} et
\ref{sites-lemma-preserve-equalizers},
s'appliquent et donnent la formule
pour $f_{big}^{-1}$ ainsi que l'existence de $f_{big!}$. De plus,
le foncteur $v$ est adjoint à droite car, pour $U/T$ et $V/S$ donnés,
on a $\Mor_S(u(U), V) = \Mor_T(U, V \times_S T)$,
comme voulu. On peut donc appliquer
Sites, lemmes \ref{sites-lemma-have-functor-other-way} et
\ref{sites-lemma-have-functor-other-way-morphism},
pour obtenir la formule pour $f_{big, *}$.
\end{proof}

\begin{lemma}
\label{topologies-lemma-morphism-big-small}
Soit $\Sch_{Zar}$ un gros site de Zariski.
Soit $f : T \to S$ un morphisme dans $\Sch_{Zar}$.
\begin{enumerate}
\item On a $i_f = f_{big} \circ i_T$, où $i_f$ est celui du
lemme \ref{topologies-lemma-put-in-T} et $i_T$ celui du
lemme \ref{topologies-lemma-at-the-bottom}.
\item Le foncteur $S_{Zar} \to T_{Zar}$,
$(U \to S) \mapsto (U \times_S T \to T)$, est continu et induit
un morphisme de sites
$$
f_{small} :
T_{Zar}
\longrightarrow
S_{Zar}
$$
Les foncteurs $f_{small}^{-1}$ et $f_{small, *}$ coïncident avec
les notions usuelles $f^{-1}$ et $f_*$ lorsque l'on identifie les faisceaux
sur $T_{Zar}$, resp.\ $S_{Zar}$, aux faisceaux sur $T$, resp.\ $S$,
au moyen du lemme \ref{topologies-lemma-Zariski-usual}.
\item On a un diagramme commutatif de morphismes de sites
$$
\xymatrix{
T_{Zar} \ar[d]_{f_{small}} &
(\Sch/T)_{Zar} \ar[d]^{f_{big}} \ar[l]^{\pi_T} \\
S_{Zar} &
(\Sch/S)_{Zar} \ar[l]_{\pi_S}
}
$$
de sorte que $f_{small} \circ \pi_T = \pi_S \circ f_{big}$ comme morphismes de topos.
\item On a $f_{small} = \pi_S \circ f_{big} \circ i_T = \pi_S \circ i_f$.
\end{enumerate}
\end{lemma}

\begin{proof}
L'égalité $i_f = f_{big} \circ i_T$ résulte de
l'égalité $i_f^{-1} = i_T^{-1} \circ f_{big}^{-1}$, qui est
claire d'après les descriptions de ces foncteurs données ci-dessus.
Cela démontre (1).

\medskip\noindent
Assertion (2) : voir Sites, exemple \ref{sites-example-continuous-map}.

\medskip\noindent
L'assertion (3) résulte de ce que $\pi_S$ et $\pi_T$ sont donnés par
les foncteurs d'inclusion, tandis que $f_{small}$ et $f_{big}$ le sont par le
foncteur de changement de base $U \mapsto U \times_S T$.

\medskip\noindent
L'assertion (4) se déduit de (3) par précomposition avec $i_T$.
\end{proof}

\noindent
Dans la situation du lemme, avec la terminologie de la
définition \ref{topologies-definition-restriction-small-zariski},
on a, pour tout faisceau $\mathcal{F}$ sur le gros site de Zariski de $T$,
$$
(f_{big, *}\mathcal{F})|_{S_{Zar}} =
f_{small, *}(\mathcal{F}|_{T_{Zar}}),
$$
Cette égalité résulte immédiatement de la commutativité du diagramme de
sites du lemme, puisque la restriction au petit site de Zariski de
$T$, resp.\ $S$, est donnée par $\pi_{T, *}$, resp.\ $\pi_{S, *}$. La formule
analogue faisant intervenir images inverses et restrictions est fausse.

\begin{lemma}
\label{topologies-lemma-composition}
Étant donnés des schémas $X$, $Y$, $Z$ dans $(\Sch/S)_{Zar}$
et des morphismes $f : X \to Y$, $g : Y \to Z$, on a
$g_{big} \circ f_{big} = (g \circ f)_{big}$ et
$g_{small} \circ f_{small} = (g \circ f)_{small}$.
\end{lemma}

\begin{proof}
Cela résulte de la description simple de l'image directe
et de l'image inverse pour les foncteurs entre gros sites donnée par le
lemme \ref{topologies-lemma-morphism-big}. Pour les foncteurs
entre petits sites, c'est le contenu de
Faisceaux, lemme \ref{sheaves-lemma-pushforward-composition},
via l'identification du lemme \ref{topologies-lemma-Zariski-usual}.
\end{proof}

\begin{lemma}
\label{topologies-lemma-morphism-big-small-cartesian-diagram}
Soit $\Sch_{Zar}$ un gros site de Zariski. Considérons un diagramme cartésien
$$
\xymatrix{
T' \ar[r]_{g'} \ar[d]_{f'} & T \ar[d]^f \\
S' \ar[r]^g & S
}
$$
dans $\Sch_{Zar}$. Alors
$i_g^{-1} \circ f_{big, *} = f'_{small, *} \circ (i_{g'})^{-1}$
et $g_{big}^{-1} \circ f_{big, *} = f'_{big, *} \circ (g'_{big})^{-1}$.
\end{lemma}

\begin{proof}
Puisque le diagramme est cartésien, on a, pour $U'/S'$,
l'égalité $U' \times_{S'} T' = U' \times_S T$. Ainsi,
$i_g^{-1} \circ f_{big, *}$ et $f'_{small, *} \circ (i_{g'})^{-1}$
envoient tous deux un faisceau $\mathcal{F}$ sur $(\Sch/T)_{Zar}$ sur le faisceau
$U' \mapsto \mathcal{F}(U' \times_{S'} T')$ sur $S'_{Zar}$
(utiliser les lemmes \ref{topologies-lemma-put-in-T} et \ref{topologies-lemma-morphism-big-small}).
La seconde égalité se démontre de la même manière ou peut se
déduire du très général
Sites, lemme \ref{sites-lemma-localize-morphism}.
\end{proof}

\noindent
On peut considérer un faisceau sur le gros site de Zariski de $S$ comme une collection
de faisceaux « usuels » sur tous les schémas au-dessus de $S$.

\begin{lemma}
\label{topologies-lemma-characterize-sheaf-big}
Soit $S$ un schéma contenu dans un gros site de Zariski $\Sch_{Zar}$.
Un faisceau $\mathcal{F}$ sur le gros site de Zariski $(\Sch/S)_{Zar}$
est donné par les données suivantes :
\begin{enumerate}
\item pour tout $T/S \in \Ob((\Sch/S)_{Zar})$, un faisceau
$\mathcal{F}_T$ sur $T$ ;
\item pour tout $f : T' \to T$ dans
$(\Sch/S)_{Zar}$, une application
$c_f : f^{-1}\mathcal{F}_T \to \mathcal{F}_{T'}$.
\end{enumerate}
Ces données sont soumises aux conditions suivantes :
\begin{enumerate}
\item[(a)] pour tous $f : T' \to T$ et $g : T'' \to T'$ dans
$(\Sch/S)_{Zar}$, la composée $c_g \circ g^{-1}c_f$
est égale à $c_{f \circ g}$ ;
\item[(b)] si $f : T' \to T$ dans $(\Sch/S)_{Zar}$ est une
immersion ouverte, alors $c_f$ est un isomorphisme.
\end{enumerate}
\end{lemma}

\begin{proof}
Ce lemme résulte d'un énoncé purement faisceautique discuté dans
Sites, remarque \ref{sites-remark-variant-glue-sheaves-absolute}.
Nous en donnons aussi une démonstration directe dans ce cas.

\medskip\noindent
Étant donné un faisceau $\mathcal{F}$ sur $\Sh((\Sch/S)_{Zar})$,
on pose $\mathcal{F}_T = i_p^{-1}\mathcal{F}$, où $p : T \to S$
est le morphisme structural. Remarquons que
$\mathcal{F}_T(U) = \mathcal{F}(U'/S)$ pour tout ouvert $U \subset T$
et toute immersion ouverte $U' \to T$ dans $(\Sch/T)_{Zar}$
d'image $U$ ; voir les lemmes \ref{topologies-lemma-Zariski-usual} et \ref{topologies-lemma-put-in-T}.
Ainsi, étant donnés $f : T' \to T$ au-dessus de $S$ et $U, U' \to T$, on obtient une
application canonique $\mathcal{F}_T(U) = \mathcal{F}(U'/S) \to \mathcal{F}(U'\times_T T'/S)
= \mathcal{F}_{T'}(f^{-1}(U))$, celle du milieu étant l'application de restriction
de $\mathcal{F}$ relativement au morphisme
$U' \times_T T' \to U'$ au-dessus de $S$. La famille de ces applications est
compatible aux restrictions et définit donc une $f$-application $c_f$
de $\mathcal{F}_T$ vers $\mathcal{F}_{T'}$ ; voir
Faisceaux, définition \ref{sheaves-definition-f-map} et la discussion
qui l'entoure. Il est clair que $c_{f \circ g}$ est la composée de
$c_f$ et $c_g$, puisque composer les applications de restriction de $\mathcal{F}$
donne encore une application de restriction.

\medskip\noindent
Réciproquement, étant donné un système $(\mathcal{F}_T, c_f)$ comme dans le lemme,
on peut définir un préfaisceau $\mathcal{F}$ sur $\Sh((\Sch/S)_{Zar})$
en posant simplement $\mathcal{F}(T/S) = \mathcal{F}_T(T)$. Comme application de restriction,
pour $f : T' \to T$ donné, on définit, pour $s \in \mathcal{F}(T)$,
l'image inverse $f^*(s)$ comme étant $c_f(s)$ (où l'on considère de nouveau $c_f$ comme
une $f$-application). La condition imposée aux $c_f$ garantit que
les images inverses satisfont à la propriété de fonctorialité requise.
Nous omettons la vérification qu'il s'agit d'un faisceau.
Il est clair que les constructions ainsi définies sont mutuellement inverses.
\end{proof}























\section{La topologie \'etale}
\label{topologies-section-etale}

\noindent
Soit $S$ un schéma. Nous voulons définir la topologie \'etale sur
la catégorie des schémas au-dessus de $S$. Conformément à notre principe général,
nous introduisons d'abord la notion de recouvrement \'etale.

\begin{definition}
\label{topologies-definition-etale-covering}
Soit $T$ un schéma. Un {\it recouvrement \'etale de $T$} est une famille
de morphismes de schémas $\{f_i : T_i \to T\}_{i \in I}$
telle que chaque $f_i$ soit \'etale et que $T = \bigcup f_i(T_i)$.
\end{definition}

\begin{lemma}
\label{topologies-lemma-zariski-etale}
Tout recouvrement de Zariski est un recouvrement \'etale.
\end{lemma}

\begin{proof}
Cela résulte immédiatement des définitions et du fait qu'une immersion ouverte
est un morphisme \'etale ; voir
Morphismes, lemme \ref{morphisms-lemma-open-immersion-etale}.
\end{proof}

\noindent
Montrons maintenant que cette notion satisfait aux conditions de
Sites, définition \ref{sites-definition-site}.

\begin{lemma}
\label{topologies-lemma-etale}
Soit $T$ un schéma.
\begin{enumerate}
\item Si $T' \to T$ est un isomorphisme, alors $\{T' \to T\}$
est un recouvrement \'etale de $T$.
\item Si $\{T_i \to T\}_{i\in I}$ est un recouvrement \'etale et si, pour tout
$i$, on a un recouvrement \'etale $\{T_{ij} \to T_i\}_{j\in J_i}$, alors
$\{T_{ij} \to T\}_{i \in I, j\in J_i}$ est un recouvrement \'etale.
\item Si $\{T_i \to T\}_{i\in I}$ est un recouvrement \'etale
et si $T' \to T$ est un morphisme de schémas, alors
$\{T' \times_T T_i \to T'\}_{i\in I}$ est un recouvrement \'etale.
\end{enumerate}
\end{lemma}

\begin{proof}
Omis.
\end{proof}

\begin{lemma}
\label{topologies-lemma-etale-affine}
Soit $T$ un schéma affine.
Soit $\{T_i \to T\}_{i \in I}$ un recouvrement \'etale de $T$.
Il existe alors un recouvrement \'etale
$\{U_j \to T\}_{j = 1, \ldots, m}$ qui raffine
$\{T_i \to T\}_{i \in I}$ et tel que chaque $U_j$ soit un schéma
affine. On peut en outre choisir chaque $U_j$ comme ouvert affine
de l'un des $T_i$.
\end{lemma}

\begin{proof}
Omis.
\end{proof}

\noindent
On définit donc comme suit les recouvrements standards d'affines correspondants.

\begin{definition}
\label{topologies-definition-standard-etale}
Soit $T$ un schéma affine. Un {\it recouvrement \'etale standard}
de $T$ est une famille $\{f_j : U_j \to T\}_{j = 1, \ldots, m}$
telle que chaque $U_j$ soit affine et \'etale sur $T$, et que
$T = \bigcup f_j(U_j)$.
\end{definition}

\noindent
Dans la définition ci-dessus, nous ne supposons {\bf pas} que les morphismes $f_j$ soient
étales standards. En effet, sous cette hypothèse, les recouvrements étales
standards ne définiraient pas un site sur $\textit{Aff}/S$, notamment à cause de
Algèbre, lemme \ref{algebra-lemma-standard-etale}, partie (4).
En revanche, un morphisme étale d'affines est automatiquement
lisse standard ; voir
Algèbre, lemme \ref{algebra-lemma-etale-standard-smooth}.
Ainsi, un recouvrement étale standard est un recouvrement lisse
standard et un recouvrement syntomique standard.

\begin{definition}
\label{topologies-definition-big-etale-site}
Un {\it gros site étale} est un site $\Sch_\etale$, au sens de
Sites, définition \ref{sites-definition-site}, construit comme suit :
\begin{enumerate}
\item Choisir un ensemble quelconque de schémas $S_0$, ainsi qu'un ensemble de recouvrements étales
$\text{Cov}_0$ entre ces schémas.
\item Prendre comme catégorie sous-jacente une catégorie $\Sch_\alpha$
construite comme dans Ensembles, lemme \ref{sets-lemma-construct-category},
à partir de l'ensemble $S_0$.
\item Choisir un ensemble de recouvrements comme dans
Ensembles, lemme \ref{sets-lemma-coverings-site}, à partir de la
catégorie $\Sch_\alpha$, de la classe des recouvrements étales
et de l'ensemble $\text{Cov}_0$ choisi ci-dessus.
\end{enumerate}
\end{definition}

\noindent
Voir les remarques suivant la définition \ref{topologies-definition-big-zariski-site}
pour la motivation et les explications relatives à la définition des gros sites.

\medskip\noindent
Avant de poursuivre par l'introduction du gros site étale d'un
schéma $S$, observons que la topologie d'un gros site étale
$\Sch_\etale$ est, en un certain sens, induite par la topologie
étale sur la catégorie de tous les schémas.

\begin{lemma}
\label{topologies-lemma-etale-induced}
Soit $\Sch_\etale$ un gros site étale comme dans la
définition \ref{topologies-definition-big-etale-site}.
Soit $T \in \Ob(\Sch_\etale)$.
Soit $\{T_i \to T\}_{i \in I}$ un recouvrement étale quelconque de $T$.
\begin{enumerate}
\item Il existe un recouvrement $\{U_j \to T\}_{j \in J}$ de $T$ dans le site
$\Sch_\etale$ qui raffine $\{T_i \to T\}_{i \in I}$.
\item Si $\{T_i \to T\}_{i \in I}$ est un recouvrement étale standard, alors
il est tautologiquement équivalent à un recouvrement dans $\Sch_\etale$.
\item Si $\{T_i \to T\}_{i \in I}$ est un recouvrement de Zariski, alors
il est tautologiquement équivalent à un recouvrement dans $\Sch_\etale$.
\end{enumerate}
\end{lemma}

\begin{proof}
Pour chaque $i$, choisissons un recouvrement ouvert affine $T_i = \bigcup_{j \in J_i} T_{ij}$
tel que chaque $T_{ij}$ s'envoie dans un sous-schéma ouvert affine de $T$. D'après le
lemme \ref{topologies-lemma-etale},
le raffinement $\{T_{ij} \to T\}_{i \in I, j \in J_i}$ est lui aussi un recouvrement étale
de $T$. On peut donc supposer chaque $T_i$ affine et s'envoyant dans
un ouvert affine $W_i$ de $T$. En appliquant
Ensembles, lemme \ref{sets-lemma-what-is-in-it},
on voit que $W_i$ est isomorphe à un objet de $\Sch_\etale$.
Mais alors $T_i$, en tant que schéma de type fini au-dessus de $W_i$,
est isomorphe à un objet $V_i$ de $\Sch_\etale$ par une seconde
application de
Ensembles, lemme \ref{sets-lemma-what-is-in-it}.
Le recouvrement $\{V_i \to T\}_{i \in I}$ raffine $\{T_i \to T\}_{i \in I}$
(puisqu'ils sont isomorphes).
De plus, $\{V_i \to T\}_{i \in I}$ est combinatoirement équivalent à un
recouvrement $\{U_j \to T\}_{j \in J}$ de $T$ dans le site
$\Sch_\etale$ d'après
Ensembles, lemme \ref{sets-lemma-what-is-in-it}.
Le recouvrement $\{U_j \to T\}_{j \in J}$ est un raffinement comme en (1).
Dans la situation de (2), (3), chacun des
schémas $T_i$ est isomorphe à un objet de $\Sch_\etale$ d'après
Ensembles, lemme \ref{sets-lemma-what-is-in-it},
et une nouvelle application de
Ensembles, lemme \ref{sets-lemma-coverings-site},
donne le résultat voulu.
\end{proof}

\begin{definition}
\label{topologies-definition-big-small-etale}
Soit $S$ un schéma. Soit $\Sch_\etale$ un gros site
étale contenant $S$.
\begin{enumerate}
\item Le {\it gros site étale de $S$}, noté
$(\Sch/S)_\etale$, est le site
$\Sch_\etale/S$ introduit dans
Sites, section \ref{sites-section-localize}.
\item Le {\it petit site étale de $S$}, que nous notons
$S_\etale$, est la sous-catégorie pleine de
$(\Sch/S)_\etale$
dont les objets sont les $U/S$ tels que $U \to S$ soit étale.
Un recouvrement de $S_\etale$ est tout recouvrement $\{U_i \to U\}$ de
$(\Sch/S)_\etale$ avec $U \in \Ob(S_\etale)$.
\item Le {\it gros site étale affine de $S$}, noté
$(\textit{Aff}/S)_\etale$, est la sous-catégorie pleine de
$(\Sch/S)_\etale$ dont les objets sont les $U/S$
tels que $U$ soit un schéma affine.
Un recouvrement de $(\textit{Aff}/S)_\etale$ est tout recouvrement
$\{U_i \to U\}$ de $(\Sch/S)_\etale$ avec $U \in \Ob((\textit{Aff}/S)_\etale)$
qui soit un recouvrement étale standard.
\item Le {\it petit site étale affine de $S$}, noté
$S_{affine, \etale}$, est la sous-catégorie pleine de $S_\etale$
dont les objets sont les $U/S$ tels que $U$ soit un schéma affine.
Un recouvrement de $S_{affine, \etale}$ est tout recouvrement $\{U_i \to U\}$
de $S_\etale$ avec $U \in \Ob(S_{affine, \etale})$ qui soit un recouvrement
étale standard.
\end{enumerate}
\end{definition}

\noindent
Il n'est pas tout à fait évident que le gros site étale affine,
le petit site étale et le petit site étale affine soient des sites.
Vérifions-le.

\begin{lemma}
\label{topologies-lemma-verify-site-etale}
Soit $S$ un schéma. Soit $\Sch_\etale$ un gros site
étale contenant $S$.
Les structures $S_\etale$, $(\textit{Aff}/S)_\etale$ et $S_{affine, \etale}$
sont des sites.
\end{lemma}

\begin{proof}
Montrons que $S_\etale$ est un site. C'est une catégorie munie d'un
ensemble donné de familles de morphismes de but fixé. Il faut donc
établir les propriétés (1), (2) et (3) de
Sites, définition \ref{sites-definition-site}.
Puisque $(\Sch/S)_\etale$ est un site, il suffit de prouver
que, pour tout recouvrement $\{U_i \to U\}$ de $(\Sch/S)_\etale$
avec $U \in \Ob(S_\etale)$, on a aussi
$U_i \in \Ob(S_\etale)$.
Cela résulte des définitions, car la composée de morphismes étales
est un morphisme étale.

\medskip\noindent
Montrons que $(\textit{Aff}/S)_\etale$ est un site.
En raisonnant comme ci-dessus, il suffit de montrer que l'ensemble
des recouvrements étales standards d'affines satisfait aux propriétés
(1), (2) et (3) de
Sites, définition \ref{sites-definition-site}.
C'est clair : par exemple, étant donné un recouvrement étale
standard $\{T_i \to T\}_{i\in I}$ et, pour chaque
$i$, un recouvrement étale standard $\{T_{ij} \to T_i\}_{j\in J_i}$,
$\{T_{ij} \to T\}_{i \in I, j\in J_i}$ est un recouvrement étale standard
car $\bigcup_{i\in I} J_i$ est fini et chaque $T_{ij}$ est affine.

\medskip\noindent
Nous omettons la démonstration du fait que $S_{affine, \'etale}$ est un site.
\end{proof}

\begin{lemma}
\label{topologies-lemma-fibre-products-etale}
Soit $S$ un schéma. Soit $\Sch_\etale$ un gros site
étale contenant $S$. Les catégories sous-jacentes des sites
$\Sch_\etale$, $(\Sch/S)_\etale$, $S_\etale$, $(\textit{Aff}/S)_\etale$,
et $S_{affine, \etale}$ admettent des produits fibrés.
Dans chaque cas, le foncteur évident vers la catégorie $\Sch$ de
tous les schémas commute à la formation des produits fibrés. Les catégories
$(\Sch/S)_\etale$ et $S_\etale$ possèdent toutes deux un
objet final, à savoir $S/S$.
\end{lemma}

\begin{proof}
Pour $\Sch_\etale$, cela est vrai par construction ; voir
Ensembles, lemme \ref{sets-lemma-what-is-in-it}.
Supposons donnés des morphismes $U \to S$, $V \to U$, $W \to U$
de schémas, avec $U, V, W \in \Ob(\Sch_\etale)$.
Le produit fibré $V \times_U W$ dans $\Sch_\etale$
est un produit fibré dans $\Sch$ et
est le produit fibré de $V/S$ et $W/S$ au-dessus de $U/S$ dans
la catégorie de tous les schémas au-dessus de $S$, donc aussi un
produit fibré dans $(\Sch/S)_\etale$.
Cela prouve le résultat pour $(\Sch/S)_\etale$.
Si $U \to S$, $V \to U$ et $W \to U$ sont étales, alors
$V \times_U W \to S$ l'est aussi, d'où le résultat pour $S_\etale$.
Si $U, V, W$ sont affines, $V \times_U W$ l'est aussi, d'où le
résultat pour $(\textit{Aff}/S)_\etale$ et $S_{affine, \etale}$.
\end{proof}

\noindent
Vérifions ensuite que le gros, resp.\ petit, site affine définit le même
topos que le gros, resp.\ petit, site.

\begin{lemma}
\label{topologies-lemma-affine-big-site-etale}
Soit $S$ un schéma. Soit $\Sch_\etale$ un gros site
étale contenant $S$.
Le foncteur
$(\textit{Aff}/S)_\etale \to (\Sch/S)_\etale$
est cocontinu spécial et induit une équivalence de topos de
$\Sh((\textit{Aff}/S)_\etale)$ vers
$\Sh((\Sch/S)_\etale)$.
\end{lemma}

\begin{proof}
La notion de foncteur cocontinu spécial est introduite dans
Sites, définition \ref{sites-definition-special-cocontinuous-functor}.
Il faut donc vérifier les hypothèses (1) -- (5) de
Sites, lemme \ref{sites-lemma-equivalence}.
Notons le foncteur d'inclusion
$u : (\textit{Aff}/S)_\etale \to (\Sch/S)_\etale$.
La cocontinuité signifie simplement que tout recouvrement étale de
$T/S$, avec $T$ affine, admet un raffinement par un recouvrement étale standard de $T$.
C'est le contenu du
lemme \ref{topologies-lemma-etale-affine}.
Ainsi, (1) est vérifiée. Le foncteur $u$ est continu simplement parce qu'un recouvrement
étale standard est un recouvrement étale. Ainsi, (2) est vérifiée.
Les points (3) et (4) découlent immédiatement du fait que $u$ est
pleinement fidèle. Enfin, la condition (5) découle du
fait que tout schéma admet un recouvrement ouvert affine.
\end{proof}

\begin{lemma}
\label{topologies-lemma-alternative}
Soit $S$ un schéma. Soit $\Sch_\etale$ un gros site
étale contenant $S$. Le foncteur $S_{affine, \etale} \to S_\etale$
est cocontinu spécial et induit une équivalence de topos de
$\Sh(S_{affine, \etale})$ vers $\Sh(S_\etale)$.
\end{lemma}

\begin{proof}
Omis. Indication : comparer avec la démonstration du
lemme \ref{topologies-lemma-affine-big-site-etale}.
\end{proof}

\noindent
Établissons maintenant quelques relations entre les topos
associés à ces sites.

\begin{lemma}
\label{topologies-lemma-put-in-T-etale}
Soit $\Sch_\etale$ un gros site étale.
Soit $f : T \to S$ un morphisme dans $\Sch_\etale$.
Le foncteur $T_\etale \to (\Sch/S)_\etale$
est cocontinu et induit un morphisme de topos
$$
i_f :
\Sh(T_\etale)
\longrightarrow
\Sh((\Sch/S)_\etale)
$$
Pour un faisceau $\mathcal{G}$ sur $(\Sch/S)_\etale$,
on a la formule $(i_f^{-1}\mathcal{G})(U/T) = \mathcal{G}(U/S)$.
Le foncteur $i_f^{-1}$ admet aussi un adjoint à gauche $i_{f, !}$ qui commute
aux produits fibrés et aux égalisateurs.
\end{lemma}

\begin{proof}
Notons le foncteur $u : T_\etale \to (\Sch/S)_\etale$.
Autrement dit, à un morphisme étale $j : U \to T$ correspondant
à un objet de $T_\etale$, on associe $u(U \to T) = (f \circ j : U \to S)$.
Ce foncteur commute aux produits fibrés ; voir le
lemme \ref{topologies-lemma-fibre-products-etale}.
Soient $a, b : U \to V$ deux morphismes dans $T_\etale$.
Dans ce cas, l'égalisateur de $a$ et $b$ (dans la catégorie des schémas) est
$$
V \times_{\Delta_{V/T}, V \times_T V, (a, b)} U
$$
qui est un produit fibré de schémas étales au-dessus de $T$, donc étale
au-dessus de $T$. Ainsi, $T_\etale$ admet des égalisateurs et $u$ commute avec eux.
Il est clairement cocontinu.
Il est aussi continu, car $u$ envoie les recouvrements sur des recouvrements et
commute aux produits fibrés. Le lemme résulte donc de
Sites, lemmes \ref{sites-lemma-when-shriek}
et \ref{sites-lemma-preserve-equalizers}.
\end{proof}

\begin{lemma}
\label{topologies-lemma-at-the-bottom-etale}
Soit $S$ un schéma. Soit $\Sch_\etale$ un gros site
étale contenant $S$.
Le foncteur d'inclusion $S_\etale \to (\Sch/S)_\etale$
satisfait aux hypothèses de Sites, lemme \ref{sites-lemma-bigger-site},
et induit donc un morphisme de sites
$$
\pi_S : (\Sch/S)_\etale \longrightarrow S_\etale
$$
et un morphisme de topos
$$
i_S : \Sh(S_\etale) \longrightarrow \Sh((\Sch/S)_\etale)
$$
tel que $\pi_S \circ i_S = \text{id}$. De plus, $i_S = i_{\text{id}_S}$,
où $i_{\text{id}_S}$ est défini comme dans le lemme \ref{topologies-lemma-put-in-T-etale}.
En particulier, le foncteur $i_S^{-1} = \pi_{S, *}$ est décrit par la règle
$i_S^{-1}(\mathcal{G})(U/S) = \mathcal{G}(U/S)$.
\end{lemma}

\begin{proof}
Dans ce cas, le foncteur
$u : S_\etale \to (\Sch/S)_\etale$,
outre les propriétés établies dans la démonstration du
lemme \ref{topologies-lemma-put-in-T-etale} ci-dessus, est pleinement fidèle
et envoie l'objet final sur l'objet final.
Le lemme résulte de Sites, lemme \ref{sites-lemma-bigger-site}.
\end{proof}

\begin{definition}
\label{topologies-definition-restriction-small-etale}
Dans la situation du
lemme \ref{topologies-lemma-at-the-bottom-etale},
le foncteur $i_S^{-1} = \pi_{S, *}$ est souvent
appelé la {\it restriction au petit site étale} et, pour un faisceau
$\mathcal{F}$ sur le gros site étale, nous notons
$\mathcal{F}|_{S_\etale}$ cette restriction.
\end{definition}

\noindent
Avec cette notation, pour un faisceau $\mathcal{F}$ sur le
gros site et un faisceau $\mathcal{G}$ sur le petit site, on a
\begin{align*}
\Mor_{\Sh(S_\etale)}(
\mathcal{F}|_{S_\etale},
\mathcal{G})
& =
\Mor_{\Sh((\Sch/S)_\etale)}(
\mathcal{F},
i_{S, *}\mathcal{G}) \\
\Mor_{\Sh(S_\etale)}(
\mathcal{G},
\mathcal{F}|_{S_\etale})
& =
\Mor_{\Sh((\Sch/S)_\etale)}(
\pi_S^{-1}\mathcal{G},
\mathcal{F})
\end{align*}
De plus, on a $(i_{S, *}\mathcal{G})|_{S_\etale} = \mathcal{G}$
et $(\pi_S^{-1}\mathcal{G})|_{S_\etale} = \mathcal{G}$.

\begin{lemma}
\label{topologies-lemma-morphism-big-etale}
Soit $\Sch_\etale$ un gros site étale.
Soit $f : T \to S$ un morphisme dans $\Sch_\etale$.
Le foncteur
$$
u :
(\Sch/T)_\etale
\longrightarrow
(\Sch/S)_\etale,
\quad
V/T \longmapsto V/S
$$
est cocontinu et admet un adjoint à droite continu
$$
v :
(\Sch/S)_\etale
\longrightarrow
(\Sch/T)_\etale,
\quad
(U \to S) \longmapsto (U \times_S T \to T).
$$
Ils induisent le même morphisme de topos
$$
f_{big} :
\Sh((\Sch/T)_\etale)
\longrightarrow
\Sh((\Sch/S)_\etale)
$$
On a $f_{big}^{-1}(\mathcal{G})(U/T) = \mathcal{G}(U/S)$.
On a $f_{big, *}(\mathcal{F})(U/S) = \mathcal{F}(U \times_S T/T)$.
En outre, $f_{big}^{-1}$ admet un adjoint à gauche $f_{big!}$ qui commute aux
produits fibrés et aux égalisateurs.
\end{lemma}

\begin{proof}
Le foncteur $u$ est cocontinu, continu et commute aux produits fibrés
et aux égalisateurs (détails omis ; comparer à la démonstration du
lemme \ref{topologies-lemma-put-in-T-etale}).
Par conséquent,
Sites, lemmes \ref{sites-lemma-when-shriek} et
\ref{sites-lemma-preserve-equalizers}
s'appliquent et donnent la formule
pour $f_{big}^{-1}$ ainsi que l'existence de $f_{big!}$. De plus,
le foncteur $v$ est un adjoint à droite car, étant donnés $U/T$ et $V/S$,
on a $\Mor_S(u(U), V) = \Mor_T(U, V \times_S T)$,
comme voulu. On peut donc appliquer
Sites, lemmes \ref{sites-lemma-have-functor-other-way} et
\ref{sites-lemma-have-functor-other-way-morphism} pour obtenir la
formule pour $f_{big, *}$.
\end{proof}

\begin{lemma}
\label{topologies-lemma-morphism-big-small-etale}
Soit $\Sch_\etale$ un gros site étale.
Soit $f : T \to S$ un morphisme dans $\Sch_\etale$.
\begin{enumerate}
\item On a $i_f = f_{big} \circ i_T$, où $i_f$ est défini comme dans le
lemme \ref{topologies-lemma-put-in-T-etale} et $i_T$ comme dans le
lemme \ref{topologies-lemma-at-the-bottom-etale}.
\item Le foncteur $S_\etale \to T_\etale$,
$(U \to S) \mapsto (U \times_S T \to T)$, est continu et induit
un morphisme de sites
$$
f_{small} : T_\etale \longrightarrow S_\etale
$$
On a $f_{small, *}(\mathcal{F})(U/S) = \mathcal{F}(U \times_S T/T)$.
\item On a un diagramme commutatif de morphismes de sites
$$
\xymatrix{
T_\etale \ar[d]_{f_{small}} &
(\Sch/T)_\etale \ar[d]^{f_{big}} \ar[l]^{\pi_T}\\
S_\etale &
(\Sch/S)_\etale \ar[l]_{\pi_S}
}
$$
de sorte que $f_{small} \circ \pi_T = \pi_S \circ f_{big}$ comme morphismes de topos.
\item On a $f_{small} = \pi_S \circ f_{big} \circ i_T = \pi_S \circ i_f$.
\end{enumerate}
\end{lemma}

\begin{proof}
L'égalité $i_f = f_{big} \circ i_T$ résulte de
l'égalité $i_f^{-1} = i_T^{-1} \circ f_{big}^{-1}$, qui découle
des descriptions de ces foncteurs données ci-dessus.
Cela prouve (1).

\medskip\noindent
Le foncteur
$u :
S_\etale
\to
T_\etale$, $u(U \to S) = (U \times_S T \to T)$
envoie les recouvrements sur des recouvrements et commute aux produits fibrés ;
voir le lemme \ref{topologies-lemma-etale} (3) et \ref{topologies-lemma-fibre-products-etale}.
De plus, $S_\etale$ et $T_\etale$ admettent tous deux des objets finaux,
à savoir $S/S$ et $T/T$, et $u(S/S) = T/T$. Par conséquent, d'après
Sites, proposition \ref{sites-proposition-get-morphism},
le foncteur $u$ correspond à un morphisme de sites
$T_\etale \to S_\etale$. Celui-ci donne à son tour naissance au
morphisme de topos ; voir
Sites, lemme \ref{sites-lemma-morphism-sites-topoi}. La description
de l'image directe résulte immédiatement de ces références.

\medskip\noindent
Le point (3) résulte du fait que $\pi_S$ et $\pi_T$ sont donnés par les
foncteurs d'inclusion, et $f_{small}$ et $f_{big}$ par les
foncteurs de changement de base $U \mapsto U \times_S T$.

\medskip\noindent
L'assertion (4) résulte de (3) en précomposant avec $i_T$.
\end{proof}

\noindent
Dans la situation du lemme, avec la terminologie de la
définition \ref{topologies-definition-restriction-small-etale},
on a, pour $\mathcal{F}$ faisceau sur le gros site étale de $T$,
$$
(f_{big, *}\mathcal{F})|_{S_\etale} =
f_{small, *}(\mathcal{F}|_{T_\etale}),
$$
Cette égalité résulte de la commutativité du diagramme de
sites du lemme, puisque la restriction au petit site étale de
$T$, resp.\ $S$, est donnée par $\pi_{T, *}$, resp.\ $\pi_{S, *}$. Une formule
analogue faisant intervenir images inverses et restrictions est fausse.

\begin{lemma}
\label{topologies-lemma-composition-etale}
Étant donnés des schémas $X$, $Y$, $Y$ de $\Sch_\etale$
et des morphismes $f : X \to Y$, $g : Y \to Z$, on a
$g_{big} \circ f_{big} = (g \circ f)_{big}$ et
$g_{small} \circ f_{small} = (g \circ f)_{small}$.
\end{lemma}

\begin{proof}
Cela résulte de la description simple de l'image directe
et de l'image inverse pour les foncteurs sur les gros sites donnée par le
lemme \ref{topologies-lemma-morphism-big-etale}. Pour les foncteurs
sur les petits sites, cela résulte de la description des
foncteurs image directe dans le lemme \ref{topologies-lemma-morphism-big-small-etale}.
\end{proof}

\begin{lemma}
\label{topologies-lemma-morphism-big-small-cartesian-diagram-etale}
Soit $\Sch_\etale$ un gros site étale. Considérons un diagramme cartésien
$$
\xymatrix{
T' \ar[r]_{g'} \ar[d]_{f'} & T \ar[d]^f \\
S' \ar[r]^g & S
}
$$
dans $\Sch_\etale$. Alors
$i_g^{-1} \circ f_{big, *} = f'_{small, *} \circ (i_{g'})^{-1}$
et $g_{big}^{-1} \circ f_{big, *} = f'_{big, *} \circ (g'_{big})^{-1}$.
\end{lemma}

\begin{proof}
Le diagramme étant cartésien, on a, pour $U'/S'$,
$U' \times_{S'} T' = U' \times_S T$. Par conséquent, les deux foncteurs
$i_g^{-1} \circ f_{big, *}$ et $f'_{small, *} \circ (i_{g'})^{-1}$
envoient un faisceau $\mathcal{F}$ sur $(\Sch/T)_\etale$ sur le faisceau
$U' \mapsto \mathcal{F}(U' \times_{S'} T')$ sur $S'_\etale$
(utiliser les lemmes \ref{topologies-lemma-put-in-T-etale} et \ref{topologies-lemma-morphism-big-etale}).
La seconde égalité peut se démontrer de la même manière ou se
déduire du très général
Sites, lemme \ref{sites-lemma-localize-morphism}.
\end{proof}

\noindent
On peut considérer un faisceau sur le gros site étale de $S$ comme une collection
de faisceaux « usuels » sur tous les schémas au-dessus de $S$.

\begin{lemma}
\label{topologies-lemma-characterize-sheaf-big-etale}
Soit $S$ un schéma contenu dans un gros site étale
$\Sch_\etale$.
Un faisceau $\mathcal{F}$ sur le gros site étale
$(\Sch/S)_\etale$ est décrit par les données suivantes :
\begin{enumerate}
\item pour tout $T/S \in \Ob((\Sch/S)_\etale)$, un faisceau
$\mathcal{F}_T$ sur $T_\etale$,
\item pour tout $f : T' \to T$ dans
$(\Sch/S)_\etale$, une application
$c_f : f_{small}^{-1}\mathcal{F}_T \to \mathcal{F}_{T'}$.
\end{enumerate}
Ces données sont soumises aux conditions suivantes :
\begin{enumerate}
\item[(a)] pour tous $f : T' \to T$ et $g : T'' \to T'$ dans
$(\Sch/S)_\etale$, la composée
$c_g \circ g_{small}^{-1}c_f$ est égale à $c_{f \circ g}$, et
\item[(b)] si $f : T' \to T$ dans $(\Sch/S)_\etale$
est étale, alors $c_f$ est un isomorphisme.
\end{enumerate}
\end{lemma}

\begin{proof}
Ce lemme résulte d'un énoncé purement faisceautique discuté dans
Sites, remarque \ref{sites-remark-variant-glue-sheaves-absolute}.
Nous donnons aussi une démonstration directe dans ce cas.

\medskip\noindent
Étant donné un faisceau $\mathcal{F}$ sur $\Sh((\Sch/S)_\etale)$,
on pose $\mathcal{F}_T = i_p^{-1}\mathcal{F}$, où $p : T \to S$
est le morphisme structural. Remarquons que
$\mathcal{F}_T(U) = \mathcal{F}(U/S)$ pour tout $U \to T$
dans $T_\etale$ ; voir le lemme \ref{topologies-lemma-put-in-T-etale}.
Ainsi, étant donnés $f : T' \to T$ au-dessus de $S$ et $U \to T$, on obtient une
application canonique $\mathcal{F}_T(U) = \mathcal{F}(U/S) \to \mathcal{F}(U \times_T T'/S)
= \mathcal{F}_{T'}(U \times_T T')$, où l'application médiane est l'application de restriction
de $\mathcal{F}$ relativement au morphisme
$U \times_T T' \to U$ au-dessus de $S$. La famille de ces applications est
compatible aux restrictions et définit donc une application
$c'_f : \mathcal{F}_T \to f_{small, *}\mathcal{F}_{T'}$, où
$u : T_\etale \to T'_\etale$ est le foncteur de changement de base
associé à $f$. Par l'adjonction de $f_{small, *}$ (voir
Sites, section \ref{sites-section-continuous-functors}) avec
$f_{small}^{-1}$, cela revient à une application
$c_f : f_{small}^{-1}\mathcal{F}_T \to \mathcal{F}_{T'}$.
Il est clair que $c'_{f \circ g}$ est la composée de
$c'_f$ et $f_{small, *}c'_g$, puisque la composée d'applications de restriction
de $\mathcal{F}$ est encore une application de restriction ; on obtient ainsi la
relation voulue entre $c_f$, $c_g$ et $c_{f \circ g}$.

\medskip\noindent
Réciproquement, étant donné un système $(\mathcal{F}_T, c_f)$ comme dans le lemme,
on peut définir un préfaisceau $\mathcal{F}$ sur
$\Sh((\Sch/S)_\etale)$
en posant simplement $\mathcal{F}(T/S) = \mathcal{F}_T(T)$. Comme application de restriction,
étant donné $f : T' \to T$, on définit, pour $s \in \mathcal{F}(T)$,
l'image inverse $f^*(s)$ comme $c_f(s)$, où l'on considère de nouveau $c_f$ comme
une application $\mathcal{F}_T \to f_{small, *}\mathcal{F}_{T'}$.
La condition sur les $c_f$ garantit que
les images inverses satisfont à la fonctorialité requise.
Nous omettons la vérification qu'il s'agit d'un faisceau.
Il est clair que les constructions ainsi définies sont inverses l'une de l'autre.
\end{proof}























\section{La topologie lisse}
\label{topologies-section-smooth}

\noindent
Dans cette section, nous définissons la topologie lisse.
Cela est quelque peu gratuit, puisqu'on verra plus loin (voir
Compléments sur les morphismes, section \ref{more-morphisms-section-etale-over-smooth})
que cette topologie définit le même topos que la
topologie étale. Elle demeure néanmoins naturelle et intervient
à l'occasion.

\begin{definition}
\label{topologies-definition-smooth-covering}
Soit $T$ un schéma. Un {\it recouvrement lisse de $T$} est une famille
de morphismes de schémas $\{f_i : T_i \to T\}_{i \in I}$
telle que chaque $f_i$ soit lisse et
que $T = \bigcup f_i(T_i)$.
\end{definition}

\begin{lemma}
\label{topologies-lemma-zariski-etale-smooth}
Tout recouvrement étale est un recouvrement lisse et, à plus forte raison,
tout recouvrement de Zariski est un recouvrement lisse.
\end{lemma}

\begin{proof}
Cela résulte des définitions et du fait qu'un morphisme étale est
lisse ; voir
Morphismes, définition \ref{morphisms-definition-etale}
et le lemme \ref{topologies-lemma-zariski-etale}.
\end{proof}

\noindent
Montrons maintenant que cette notion satisfait aux conditions de
Sites, définition \ref{sites-definition-site}.

\begin{lemma}
\label{topologies-lemma-smooth}
Soit $T$ un schéma.
\begin{enumerate}
\item Si $T' \to T$ est un isomorphisme, alors $\{T' \to T\}$
est un recouvrement lisse de $T$.
\item Si $\{T_i \to T\}_{i\in I}$ est un recouvrement lisse et si, pour tout
$i$, on a un recouvrement lisse $\{T_{ij} \to T_i\}_{j\in J_i}$, alors
$\{T_{ij} \to T\}_{i \in I, j\in J_i}$ est un recouvrement lisse.
\item Si $\{T_i \to T\}_{i\in I}$ est un recouvrement lisse
et si $T' \to T$ est un morphisme de schémas, alors
$\{T' \times_T T_i \to T'\}_{i\in I}$ est un recouvrement lisse.
\end{enumerate}
\end{lemma}

\begin{proof}
Omis.
\end{proof}

\begin{lemma}
\label{topologies-lemma-smooth-affine}
Soit $T$ un schéma affine.
Soit $\{T_i \to T\}_{i \in I}$ un recouvrement lisse de $T$.
Il existe alors un recouvrement lisse
$\{U_j \to T\}_{j = 1, \ldots, m}$ qui raffine
$\{T_i \to T\}_{i \in I}$, tel que chaque $U_j$ soit un schéma
affine et que chaque morphisme $U_j \to T$ soit lisse
standard ; voir Morphismes, définition \ref{morphisms-definition-smooth}.
On peut en outre choisir chaque $U_j$ comme ouvert affine de l'un des $T_i$.
\end{lemma}

\begin{proof}
Omis ; voir toutefois Algèbre, lemme \ref{algebra-lemma-smooth-syntomic}.
\end{proof}

\noindent
On définit donc comme suit les recouvrements standards d'affines correspondants.

\begin{definition}
\label{topologies-definition-standard-smooth}
Soit $T$ un schéma affine. Un {\it recouvrement lisse standard}
de $T$ est une famille $\{f_j : U_j \to T\}_{j = 1, \ldots, m}$
telle que chaque $U_j$ soit affine, que $U_j \to T$ soit lisse standard
et que $T = \bigcup f_j(U_j)$.
\end{definition}

\begin{definition}
\label{topologies-definition-big-smooth-site}
Un {\it gros site lisse} est un site $\Sch_{smooth}$ au sens de
Sites, définition \ref{sites-definition-site}, construit comme suit :
\begin{enumerate}
\item Choisir un ensemble quelconque de schémas $S_0$ et un ensemble de recouvrements lisses
$\text{Cov}_0$ entre ces schémas.
\item Prendre comme catégorie sous-jacente une catégorie $\Sch_\alpha$
construite comme dans Ensembles, lemme \ref{sets-lemma-construct-category},
à partir de l'ensemble $S_0$.
\item Choisir un ensemble de recouvrements comme dans
Ensembles, lemme \ref{sets-lemma-coverings-site}, à partir de la
catégorie $\Sch_\alpha$, de la classe des recouvrements lisses
et de l'ensemble $\text{Cov}_0$ choisi ci-dessus.
\end{enumerate}
\end{definition}

\noindent
Voir les remarques suivant la définition \ref{topologies-definition-big-zariski-site}
pour la motivation et les explications relatives à la définition des gros sites.

\medskip\noindent
Avant de poursuivre par l'introduction du gros site lisse d'un
schéma $S$, observons que la topologie d'un gros site lisse
$\Sch_{smooth}$ est, en un certain sens, induite par la topologie lisse
sur la catégorie de tous les schémas.

\begin{lemma}
\label{topologies-lemma-smooth-induced}
Soit $\Sch_{smooth}$ un gros site lisse comme dans la
définition \ref{topologies-definition-big-smooth-site}.
Soit $T \in \Ob(\Sch_{smooth})$.
Soit $\{T_i \to T\}_{i \in I}$ un recouvrement lisse quelconque de $T$.
\begin{enumerate}
\item Il existe un recouvrement $\{U_j \to T\}_{j \in J}$ de $T$ dans le site
$\Sch_{smooth}$ qui raffine $\{T_i \to T\}_{i \in I}$.
\item Si $\{T_i \to T\}_{i \in I}$ est un recouvrement lisse standard, alors
il est tautologiquement équivalent à un recouvrement de $\Sch_{smooth}$.
\item Si $\{T_i \to T\}_{i \in I}$ est un recouvrement de Zariski, alors
il est tautologiquement équivalent à un recouvrement de $\Sch_{smooth}$.
\end{enumerate}
\end{lemma}

\begin{proof}
Pour chaque $i$, choisissons un recouvrement ouvert affine $T_i = \bigcup_{j \in J_i} T_{ij}$
tel que chaque $T_{ij}$ s'envoie dans un sous-schéma ouvert affine de $T$. D'après le
lemme \ref{topologies-lemma-smooth},
le raffinement $\{T_{ij} \to T\}_{i \in I, j \in J_i}$ est lui aussi un recouvrement lisse
de $T$. On peut donc supposer chaque $T_i$ affine et s'envoyant dans
un ouvert affine $W_i$ de $T$. En appliquant
Ensembles, lemme \ref{sets-lemma-what-is-in-it},
on voit que $W_i$ est isomorphe à un objet de $\Sch_{smooth}$.
Mais alors $T_i$, comme schéma de type fini au-dessus de $W_i$,
est isomorphe à un objet $V_i$ de $\Sch_{smooth}$ par une seconde
application de
Ensembles, lemme \ref{sets-lemma-what-is-in-it}.
Le recouvrement $\{V_i \to T\}_{i \in I}$ raffine $\{T_i \to T\}_{i \in I}$
(puisqu'ils sont isomorphes).
De plus, $\{V_i \to T\}_{i \in I}$ est combinatoirement équivalent à un
recouvrement $\{U_j \to T\}_{j \in J}$ de $T$ dans le site
$\Sch_{smooth}$ d'après
Ensembles, lemme \ref{sets-lemma-what-is-in-it}.
Le recouvrement $\{U_j \to T\}_{j \in J}$ est un raffinement comme en (1).
Dans la situation de (2), (3), chacun des
schémas $T_i$ est isomorphe à un objet de $\Sch_{smooth}$ d'après
Ensembles, lemme \ref{sets-lemma-what-is-in-it},
et une nouvelle application de
Ensembles, lemme \ref{sets-lemma-coverings-site},
donne le résultat voulu.
\end{proof}

\begin{definition}
\label{topologies-definition-big-small-smooth}
Soit $S$ un schéma. Soit $\Sch_{smooth}$ un gros site
lisse contenant $S$.
\begin{enumerate}
\item Le {\it gros site lisse de $S$}, noté
$(\Sch/S)_{smooth}$, est le site $\Sch_{smooth}/S$
introduit dans Sites, section \ref{sites-section-localize}.
\item Le {\it gros site lisse affine de $S$}, noté
$(\textit{Aff}/S)_{smooth}$, est la sous-catégorie pleine de
$(\Sch/S)_{smooth}$ dont les objets sont les $U/S$ affines.
Un recouvrement de $(\textit{Aff}/S)_{smooth}$ est tout recouvrement
$\{U_i \to U\}$ de $(\Sch/S)_{smooth}$ qui soit un
recouvrement lisse standard.
\end{enumerate}
\end{definition}

\noindent
Vérifions ensuite que le gros site affine définit le même
topos que le gros site.

\begin{lemma}
\label{topologies-lemma-affine-big-site-smooth}
Soit $S$ un schéma. Soit $\Sch_{smooth}$ un gros site
lisse contenant $S$.
Le foncteur
$(\textit{Aff}/S)_{smooth} \to (\Sch/S)_{smooth}$
est cocontinu spécial et induit une équivalence de topos de
$\Sh((\textit{Aff}/S)_{smooth})$ vers
$\Sh((\Sch/S)_{smooth})$.
\end{lemma}

\begin{proof}
La notion de foncteur cocontinu spécial est introduite dans
Sites, définition \ref{sites-definition-special-cocontinuous-functor}.
Il faut donc vérifier les hypothèses (1) -- (5) de
Sites, lemme \ref{sites-lemma-equivalence}.
Notons le foncteur d'inclusion
$u : (\textit{Aff}/S)_{smooth} \to (\Sch/S)_{smooth}$.
La cocontinuité signifie simplement que tout recouvrement lisse de
$T/S$, avec $T$ affine, admet un raffinement par un recouvrement lisse standard de $T$.
C'est le contenu du
lemme \ref{topologies-lemma-smooth-affine}.
Ainsi, (1) est vérifiée. Le foncteur $u$ est continu simplement parce qu'un recouvrement
lisse standard est un recouvrement lisse. Ainsi, (2) est vérifiée.
Les points (3) et (4) découlent immédiatement du fait que $u$ est
pleinement fidèle. Enfin, la condition (5) découle du
fait que tout schéma admet un recouvrement ouvert affine.
\end{proof}

\noindent
À suivre...

\begin{lemma}
\label{topologies-lemma-morphism-big-smooth}
Soit $\Sch_{smooth}$ un gros site lisse.
Soit $f : T \to S$ un morphisme dans $\Sch_{smooth}$.
Le foncteur
$$
u : (\Sch/T)_{smooth} \longrightarrow (\Sch/S)_{smooth},
\quad
V/T \longmapsto V/S
$$
est cocontinu et admet un adjoint à droite continu
$$
v : (\Sch/S)_{smooth} \longrightarrow (\Sch/T)_{smooth},
\quad
(U \to S) \longmapsto (U \times_S T \to T).
$$
Ils induisent le même morphisme de topos
$$
f_{big} :
\Sh((\Sch/T)_{smooth})
\longrightarrow
\Sh((\Sch/S)_{smooth})
$$
On a $f_{big}^{-1}(\mathcal{G})(U/T) = \mathcal{G}(U/S)$.
On a $f_{big, *}(\mathcal{F})(U/S) = \mathcal{F}(U \times_S T/T)$.
En outre, $f_{big}^{-1}$ admet un adjoint à gauche $f_{big!}$ qui commute aux
produits fibrés et aux égalisateurs.
\end{lemma}

\begin{proof}
Le foncteur $u$ est cocontinu, continu et commute aux produits fibrés
et aux égalisateurs. Par conséquent,
Sites, lemmes \ref{sites-lemma-when-shriek} et
\ref{sites-lemma-preserve-equalizers}
s'appliquent et donnent la formule
pour $f_{big}^{-1}$ ainsi que l'existence de $f_{big!}$. De plus,
le foncteur $v$ est un adjoint à droite car, étant donnés $U/T$ et $V/S$,
on a $\Mor_S(u(U), V) = \Mor_T(U, V \times_S T)$,
comme voulu. On peut donc appliquer
Sites, lemmes \ref{sites-lemma-have-functor-other-way} et
\ref{sites-lemma-have-functor-other-way-morphism} pour obtenir la
formule pour $f_{big, *}$.
\end{proof}











\section{La topologie syntomique}
\label{topologies-section-syntomic}

\noindent
Dans cette section, nous définissons la topologie syntomique.
Cette topologie est assez intéressante : elle possède souvent
les mêmes groupes de cohomologie que la topologie fppf,
tout en étant techniquement plus facile à manier.

\begin{definition}
\label{topologies-definition-syntomic-covering}
Soit $T$ un schéma. Un {\it recouvrement syntomique de $T$} est une famille
de morphismes de schémas $\{f_i : T_i \to T\}_{i \in I}$
telle que chaque $f_i$ soit syntomique et
que $T = \bigcup f_i(T_i)$.
\end{definition}

\begin{lemma}
\label{topologies-lemma-zariski-etale-smooth-syntomic}
Tout recouvrement lisse est un recouvrement syntomique et, à plus forte raison,
tout recouvrement étale ou de Zariski est un recouvrement syntomique.
\end{lemma}

\begin{proof}
Cela résulte des définitions et du fait qu'un
morphisme lisse est syntomique ; voir
Morphismes, lemme \ref{morphisms-lemma-smooth-syntomic}
et le lemme \ref{topologies-lemma-zariski-etale-smooth}.
\end{proof}

\noindent
Montrons maintenant que cette notion satisfait aux conditions de
Sites, définition \ref{sites-definition-site}.

\begin{lemma}
\label{topologies-lemma-syntomic}
Soit $T$ un schéma.
\begin{enumerate}
\item Si $T' \to T$ est un isomorphisme, alors $\{T' \to T\}$
est un recouvrement syntomique de $T$.
\item Si $\{T_i \to T\}_{i\in I}$ est un recouvrement syntomique et si, pour tout
$i$, on a un recouvrement syntomique $\{T_{ij} \to T_i\}_{j\in J_i}$, alors
$\{T_{ij} \to T\}_{i \in I, j\in J_i}$ est un recouvrement syntomique.
\item Si $\{T_i \to T\}_{i\in I}$ est un recouvrement syntomique
et si $T' \to T$ est un morphisme de schémas, alors
$\{T' \times_T T_i \to T'\}_{i\in I}$ est un recouvrement syntomique.
\end{enumerate}
\end{lemma}

\begin{proof}
Omis.
\end{proof}

\begin{lemma}
\label{topologies-lemma-syntomic-affine}
Soit $T$ un schéma affine.
Soit $\{T_i \to T\}_{i \in I}$ un recouvrement syntomique de $T$.
Il existe alors un recouvrement syntomique
$\{U_j \to T\}_{j = 1, \ldots, m}$ qui raffine
$\{T_i \to T\}_{i \in I}$, tel que chaque $U_j$ soit un schéma
affine et que chaque morphisme $U_j \to T$ soit syntomique
standard ; voir Morphismes, définition \ref{morphisms-definition-syntomic}.
On peut en outre choisir chaque $U_j$ comme ouvert affine de l'un des $T_i$.
\end{lemma}

\begin{proof}
Omis ; voir toutefois Algèbre, lemme \ref{algebra-lemma-syntomic}.
\end{proof}

\noindent
On définit donc comme suit les recouvrements standards d'affines correspondants.

\begin{definition}
\label{topologies-definition-standard-syntomic}
Soit $T$ un schéma affine. Un {\it recouvrement syntomique standard} de $T$ est
une famille $\{f_j : U_j \to T\}_{j = 1, \ldots, m}$ telle que chaque $U_j$ soit
affine, que $U_j \to T$ soit syntomique standard et que $T = \bigcup f_j(U_j)$.
\end{definition}

\begin{definition}
\label{topologies-definition-big-syntomic-site}
Un {\it gros site syntomique} est un site $\Sch_{syntomic}$ au sens de
Sites, définition \ref{sites-definition-site}, construit comme suit :
\begin{enumerate}
\item Choisir un ensemble quelconque de schémas $S_0$ et un ensemble de recouvrements syntomiques
$\text{Cov}_0$ entre ces schémas.
\item Prendre comme catégorie sous-jacente une catégorie $\Sch_\alpha$
construite comme dans Ensembles, lemme \ref{sets-lemma-construct-category},
à partir de l'ensemble $S_0$.
\item Choisir un ensemble de recouvrements comme dans
Ensembles, lemme \ref{sets-lemma-coverings-site}, à partir de la
catégorie $\Sch_\alpha$, de la classe des recouvrements syntomiques
et de l'ensemble $\text{Cov}_0$ choisi ci-dessus.
\end{enumerate}
\end{definition}

\noindent
Voir les remarques suivant la définition \ref{topologies-definition-big-zariski-site}
pour la motivation et les explications relatives à la définition des gros sites.

\medskip\noindent
Avant de poursuivre par l'introduction du gros site syntomique d'un
schéma $S$, observons que la topologie d'un gros site syntomique
$\Sch_{syntomic}$ est, en un certain sens, induite par la topologie syntomique
sur la catégorie de tous les schémas.

\begin{lemma}
\label{topologies-lemma-syntomic-induced}
Soit $\Sch_{syntomic}$ un gros site syntomique comme dans la
définition \ref{topologies-definition-big-syntomic-site}.
Soit $T \in \Ob(\Sch_{syntomic})$.
Soit $\{T_i \to T\}_{i \in I}$ un recouvrement syntomique quelconque de $T$.
\begin{enumerate}
\item Il existe un recouvrement $\{U_j \to T\}_{j \in J}$ de $T$ dans le site
$\Sch_{syntomic}$ qui raffine $\{T_i \to T\}_{i \in I}$.
\item Si $\{T_i \to T\}_{i \in I}$ est un recouvrement syntomique standard, alors
il est tautologiquement équivalent à un recouvrement dans $\Sch_{syntomic}$.
\item Si $\{T_i \to T\}_{i \in I}$ est un recouvrement de Zariski, alors
il est tautologiquement équivalent à un recouvrement dans $\Sch_{syntomic}$.
\end{enumerate}
\end{lemma}

\begin{proof}
Pour chaque $i$, choisissons un recouvrement ouvert affine $T_i = \bigcup_{j \in J_i} T_{ij}$
tel que chaque $T_{ij}$ s'envoie dans un sous-schéma ouvert affine de $T$. D'après le
lemme \ref{topologies-lemma-syntomic},
le raffinement $\{T_{ij} \to T\}_{i \in I, j \in J_i}$ est lui aussi un recouvrement syntomique
de $T$. On peut donc supposer chaque $T_i$ affine et s'envoyant dans
un ouvert affine $W_i$ de $T$. En appliquant
Ensembles, lemme \ref{sets-lemma-what-is-in-it},
on voit que $W_i$ est isomorphe à un objet de $\Sch_{syntomic}$.
Mais alors $T_i$, comme schéma de type fini au-dessus de $W_i$,
est isomorphe à un objet $V_i$ de $\Sch_{syntomic}$ par une seconde
application de
Ensembles, lemme \ref{sets-lemma-what-is-in-it}.
Le recouvrement $\{V_i \to T\}_{i \in I}$ raffine $\{T_i \to T\}_{i \in I}$
(puisqu'ils sont isomorphes).
De plus, $\{V_i \to T\}_{i \in I}$ est combinatoirement équivalent à un
recouvrement $\{U_j \to T\}_{j \in J}$ de $T$ dans le site
$\Sch_{syntomic}$ d'après
Ensembles, lemme \ref{sets-lemma-what-is-in-it}.
Le recouvrement $\{U_j \to T\}_{j \in J}$ est un recouvrement comme en (1).
Dans la situation de (2), (3), chacun des
schémas $T_i$ est isomorphe à un objet de $\Sch_{syntomic}$ d'après
Ensembles, lemme \ref{sets-lemma-what-is-in-it},
et une nouvelle application de
Ensembles, lemme \ref{sets-lemma-coverings-site},
donne le résultat voulu.
\end{proof}

\begin{definition}
\label{topologies-definition-big-small-syntomic}
Soit $S$ un schéma. Soit $\Sch_{syntomic}$ un gros site
syntomique contenant $S$.
\begin{enumerate}
\item Le {\it gros site syntomique de $S$}, noté
$(\Sch/S)_{syntomic}$, est le site $\Sch_{syntomic}/S$
introduit dans Sites, section \ref{sites-section-localize}.
\item Le {\it gros site syntomique affine de $S$}, noté
$(\textit{Aff}/S)_{syntomic}$, est la sous-catégorie pleine de
$(\Sch/S)_{syntomic}$ dont les objets sont les $U/S$ affines.
Un recouvrement de $(\textit{Aff}/S)_{syntomic}$ est tout recouvrement
$\{U_i \to U\}$ de $(\Sch/S)_{syntomic}$ qui soit un
recouvrement syntomique standard.
\end{enumerate}
\end{definition}

\noindent
Vérifions ensuite que le gros site affine définit le même
topos que le gros site.

\begin{lemma}
\label{topologies-lemma-affine-big-site-syntomic}
Soit $S$ un schéma. Soit $\Sch_{syntomic}$ un gros site
syntomique contenant $S$.
Le foncteur
$(\textit{Aff}/S)_{syntomic} \to (\Sch/S)_{syntomic}$
est cocontinu spécial et induit une équivalence de topos de
$\Sh((\textit{Aff}/S)_{syntomic})$ vers
$\Sh((\Sch/S)_{syntomic})$.
\end{lemma}

\begin{proof}
La notion de foncteur cocontinu spécial est introduite dans
Sites, définition \ref{sites-definition-special-cocontinuous-functor}.
Il faut donc vérifier les hypothèses (1) -- (5) de
Sites, lemme \ref{sites-lemma-equivalence}.
Notons le foncteur d'inclusion
$u : (\textit{Aff}/S)_{syntomic} \to (\Sch/S)_{syntomic}$.
La cocontinuité signifie simplement que tout recouvrement syntomique de
$T/S$, avec $T$ affine, admet un raffinement par un recouvrement syntomique standard de $T$.
C'est le contenu du
lemme \ref{topologies-lemma-syntomic-affine}.
Ainsi, (1) est vérifiée. Le foncteur $u$ est continu simplement parce qu'un recouvrement
syntomique standard est un recouvrement syntomique. Ainsi, (2) est vérifiée.
Les points (3) et (4) découlent immédiatement du fait que $u$ est
pleinement fidèle. Enfin, la condition (5) découle du
fait que tout schéma admet un recouvrement ouvert affine.
\end{proof}

\noindent
À suivre...

\begin{lemma}
\label{topologies-lemma-morphism-big-syntomic}
Soit $\Sch_{syntomic}$ un gros site syntomique.
Soit $f : T \to S$ un morphisme dans $\Sch_{syntomic}$.
Le foncteur
$$
u : (\Sch/T)_{syntomic} \longrightarrow (\Sch/S)_{syntomic},
\quad
V/T \longmapsto V/S
$$
est cocontinu et admet un adjoint à droite continu
$$
v : (\Sch/S)_{syntomic} \longrightarrow (\Sch/T)_{syntomic},
\quad
(U \to S) \longmapsto (U \times_S T \to T).
$$
Ils induisent le même morphisme de topos
$$
f_{big} :
\Sh((\Sch/T)_{syntomic})
\longrightarrow
\Sh((\Sch/S)_{syntomic})
$$
On a $f_{big}^{-1}(\mathcal{G})(U/T) = \mathcal{G}(U/S)$.
On a $f_{big, *}(\mathcal{F})(U/S) = \mathcal{F}(U \times_S T/T)$.
En outre, $f_{big}^{-1}$ admet un adjoint à gauche $f_{big!}$ qui commute aux
produits fibrés et aux égalisateurs.
\end{lemma}

\begin{proof}
Le foncteur $u$ est cocontinu, continu et commute aux produits fibrés
et aux égalisateurs. Par conséquent,
Sites, lemmes \ref{sites-lemma-when-shriek} et
\ref{sites-lemma-preserve-equalizers}
s'appliquent et donnent la formule
pour $f_{big}^{-1}$ ainsi que l'existence de $f_{big!}$. De plus,
le foncteur $v$ est un adjoint à droite car, étant donnés $U/T$ et $V/S$,
on a $\Mor_S(u(U), V) = \Mor_T(U, V \times_S T)$,
comme voulu. On peut donc appliquer
Sites, lemmes \ref{sites-lemma-have-functor-other-way} et
\ref{sites-lemma-have-functor-other-way-morphism} pour obtenir la
formule pour $f_{big, *}$.
\end{proof}













\section{La topologie fppf}
\label{topologies-section-fppf}

\noindent
Soit $S$ un schéma. Nous voulons définir la topologie fppf\footnote{
Les lettres fppf signifient « fid\`element plat de pr\'esentation finie ».} sur
la catégorie des schémas au-dessus de $S$. Conformément à notre principe général,
nous introduisons d'abord la notion de recouvrement fppf.

\begin{definition}
\label{topologies-definition-fppf-covering}
Soit $T$ un schéma. Un {\it recouvrement fppf de $T$} est une famille
de morphismes de schémas $\{f_i : T_i \to T\}_{i \in I}$
telle que chaque $f_i$ soit plat, localement de présentation finie, et
que $T = \bigcup f_i(T_i)$.
\end{definition}

\begin{lemma}
\label{topologies-lemma-zariski-etale-smooth-syntomic-fppf}
Tout recouvrement syntomique est un recouvrement fppf et, à plus forte raison,
tout recouvrement lisse, étale ou de Zariski est un recouvrement fppf.
\end{lemma}

\begin{proof}
Cela résulte des définitions et du fait qu'un morphisme syntomique
est plat et localement de présentation finie ; voir
Morphismes, lemmes
\ref{morphisms-lemma-syntomic-locally-finite-presentation} et
\ref{morphisms-lemma-syntomic-flat},
ainsi que
le lemme \ref{topologies-lemma-zariski-etale-smooth-syntomic}.
\end{proof}

\noindent
Montrons maintenant que cette notion satisfait aux conditions de
Sites, définition \ref{sites-definition-site}.

\begin{lemma}
\label{topologies-lemma-fppf}
Soit $T$ un schéma.
\begin{enumerate}
\item Si $T' \to T$ est un isomorphisme, alors $\{T' \to T\}$
est un recouvrement fppf de $T$.
\item Si $\{T_i \to T\}_{i\in I}$ est un recouvrement fppf et si, pour tout
$i$, on a un recouvrement fppf $\{T_{ij} \to T_i\}_{j\in J_i}$, alors
$\{T_{ij} \to T\}_{i \in I, j\in J_i}$ est un recouvrement fppf.
\item Si $\{T_i \to T\}_{i\in I}$ est un recouvrement fppf
et si $T' \to T$ est un morphisme de schémas, alors
$\{T' \times_T T_i \to T'\}_{i\in I}$ est un recouvrement fppf.
\end{enumerate}
\end{lemma}

\begin{proof}
La première assertion est claire.
La deuxième résulte de ce que la composée de morphismes plats est plate
(voir Morphismes, lemme \ref{morphisms-lemma-composition-flat})
et que la composée de morphismes de présentation finie est
de présentation finie
(voir Morphismes, lemme \ref{morphisms-lemma-composition-finite-presentation}).
La troisième résulte de ce que le changement de base d'un morphisme plat est plat
(voir Morphismes, lemme \ref{morphisms-lemma-base-change-flat})
et que le changement de base d'un morphisme de présentation finie est
de présentation finie
(voir Morphismes, lemme \ref{morphisms-lemma-base-change-finite-presentation}).
De plus, le changement de base d'une famille surjective de morphismes est surjectif
(démonstration omise).
\end{proof}

\begin{lemma}
\label{topologies-lemma-fppf-affine}
Soit $T$ un schéma affine.
Soit $\{T_i \to T\}_{i \in I}$ un recouvrement fppf de $T$.
Il existe alors un recouvrement fppf
$\{U_j \to T\}_{j = 1, \ldots, m}$ qui raffine
$\{T_i \to T\}_{i \in I}$ et tel que chaque $U_j$ soit un schéma
affine. On peut en outre choisir chaque $U_j$ comme ouvert affine
de l'un des $T_i$.
\end{lemma}

\begin{proof}
Cela résulte directement des définitions, puisqu'un
morphisme plat et localement de présentation finie est ouvert ;
voir Morphismes, lemme \ref{morphisms-lemma-fppf-open}.
\end{proof}

\noindent
On définit donc comme suit les recouvrements standards d'affines correspondants.

\begin{definition}
\label{topologies-definition-standard-fppf}
Soit $T$ un schéma affine. Un {\it recouvrement fppf standard}
de $T$ est une famille $\{f_j : U_j \to T\}_{j = 1, \ldots, m}$
telle que chaque $U_j$ soit affine, plat et de présentation finie sur $T$,
et que $T = \bigcup f_j(U_j)$.
\end{definition}

\begin{definition}
\label{topologies-definition-big-fppf-site}
Un {\it gros site fppf} est un site $\Sch_{fppf}$ au sens de
Sites, définition \ref{sites-definition-site}, construit comme suit :
\begin{enumerate}
\item Choisir un ensemble quelconque de schémas $S_0$ et un ensemble de recouvrements fppf
$\text{Cov}_0$ entre ces schémas.
\item Prendre comme catégorie sous-jacente une catégorie $\Sch_\alpha$
construite comme dans Ensembles, lemme \ref{sets-lemma-construct-category},
à partir de l'ensemble $S_0$.
\item Choisir un ensemble de recouvrements comme dans
Ensembles, lemme \ref{sets-lemma-coverings-site}, à partir de la
catégorie $\Sch_\alpha$, de la classe des recouvrements fppf
et de l'ensemble $\text{Cov}_0$ choisi ci-dessus.
\end{enumerate}
\end{definition}

\noindent
Voir les remarques suivant la définition \ref{topologies-definition-big-zariski-site}
pour la motivation et les explications relatives à la définition des gros sites.

\medskip\noindent
Avant de poursuivre par l'introduction du gros site fppf d'un
schéma $S$, observons que la topologie d'un gros site fppf
$\Sch_{fppf}$ est, en un certain sens, induite par la topologie fppf
sur la catégorie de tous les schémas.

\begin{lemma}
\label{topologies-lemma-fppf-induced}
Soit $\Sch_{fppf}$ un gros site fppf comme dans la
définition \ref{topologies-definition-big-fppf-site}.
Soit $T \in \Ob(\Sch_{fppf})$.
Soit $\{T_i \to T\}_{i \in I}$ un recouvrement fppf quelconque de $T$.
\begin{enumerate}
\item Il existe un recouvrement $\{U_j \to T\}_{j \in J}$ de $T$ dans le site
$\Sch_{fppf}$ qui raffine $\{T_i \to T\}_{i \in I}$.
\item Si $\{T_i \to T\}_{i \in I}$ est un recouvrement fppf standard, alors
il est tautologiquement équivalent à un recouvrement de $\Sch_{fppf}$.
\item Si $\{T_i \to T\}_{i \in I}$ est un recouvrement de Zariski, alors
il est tautologiquement équivalent à un recouvrement de $\Sch_{fppf}$.
\end{enumerate}
\end{lemma}

\begin{proof}
Pour chaque $i$, choisissons un recouvrement ouvert affine $T_i = \bigcup_{j \in J_i} T_{ij}$
tel que chaque $T_{ij}$ s'envoie dans un sous-schéma ouvert affine de $T$. D'après le
lemme \ref{topologies-lemma-fppf},
le raffinement $\{T_{ij} \to T\}_{i \in I, j \in J_i}$ est lui aussi un recouvrement fppf
de $T$. On peut donc supposer chaque $T_i$ affine et s'envoyant dans
un ouvert affine $W_i$ de $T$. En appliquant
Ensembles, lemme \ref{sets-lemma-what-is-in-it},
on voit que $W_i$ est isomorphe à un objet de $\Sch_{fppf}$.
Mais alors $T_i$, comme schéma de type fini au-dessus de $W_i$,
est isomorphe à un objet $V_i$ de $\Sch_{fppf}$ par une seconde
application de
Ensembles, lemme \ref{sets-lemma-what-is-in-it}.
Le recouvrement $\{V_i \to T\}_{i \in I}$ raffine $\{T_i \to T\}_{i \in I}$
(puisqu'ils sont isomorphes).
De plus, $\{V_i \to T\}_{i \in I}$ est combinatoirement équivalent à un
recouvrement $\{U_j \to T\}_{j \in J}$ de $T$ dans le site
$\Sch_{fppf}$ d'après
Ensembles, lemme \ref{sets-lemma-what-is-in-it}.
Le recouvrement $\{U_j \to T\}_{j \in J}$ est un raffinement comme en (1).
Dans la situation de (2), (3), chacun des
schémas $T_i$ est isomorphe à un objet de $\Sch_{fppf}$ d'après
Ensembles, lemme \ref{sets-lemma-what-is-in-it},
et une nouvelle application de
Ensembles, lemme \ref{sets-lemma-coverings-site},
donne le résultat voulu.
\end{proof}

\begin{definition}
\label{topologies-definition-big-small-fppf}
Soit $S$ un schéma. Soit $\Sch_{fppf}$ un gros site
fppf contenant $S$.
\begin{enumerate}
\item Le {\it gros site fppf de $S$}, noté
$(\Sch/S)_{fppf}$, est le site $\Sch_{fppf}/S$
introduit dans Sites, section \ref{sites-section-localize}.
\item Le {\it gros site fppf affine de $S$}, noté
$(\textit{Aff}/S)_{fppf}$, est la sous-catégorie pleine de
$(\Sch/S)_{fppf}$ dont les objets sont les $U/S$ affines.
Un recouvrement de $(\textit{Aff}/S)_{fppf}$ est tout recouvrement
$\{U_i \to U\}$ de $(\Sch/S)_{fppf}$ qui soit un
recouvrement fppf standard.
\end{enumerate}
\end{definition}

\noindent
Il n'est pas tout à fait évident que
le gros site fppf affine soit un site. Vérifions-le.

\begin{lemma}
\label{topologies-lemma-verify-site-fppf}
Soit $S$ un schéma. Soit $\Sch_{fppf}$ un gros site
fppf contenant $S$. Alors $(\textit{Aff}/S)_{fppf}$ est un site.
\end{lemma}

\begin{proof}
Montrons que $(\textit{Aff}/S)_{fppf}$ est un site.
En raisonnant comme dans la démonstration du lemme \ref{topologies-lemma-verify-site-etale},
il suffit de montrer que l'ensemble
des recouvrements fppf standards d'affines satisfait aux propriétés
(1), (2) et (3) de
Sites, définition \ref{sites-definition-site}.
C'est clair : par exemple, étant donné un recouvrement fppf
standard $\{T_i \to T\}_{i\in I}$ et, pour chaque
$i$, un recouvrement fppf standard $\{T_{ij} \to T_i\}_{j\in J_i}$,
$\{T_{ij} \to T\}_{i \in I, j\in J_i}$ est un recouvrement fppf standard
car $\bigcup_{i\in I} J_i$ est fini et chaque $T_{ij}$ est affine.
\end{proof}

\begin{lemma}
\label{topologies-lemma-fibre-products-fppf}
Soit $S$ un schéma. Soit $\Sch_{fppf}$ un gros site
fppf contenant $S$. Les catégories sous-jacentes des sites
$\Sch_{fppf}$, $(\Sch/S)_{fppf}$,
et $(\textit{Aff}/S)_{fppf}$ admettent des produits fibrés.
Dans chaque cas, le foncteur évident vers la catégorie $\Sch$ de
tous les schémas commute à la formation des produits fibrés. La catégorie
$(\Sch/S)_{fppf}$ possède un objet final, à savoir $S/S$.
\end{lemma}

\begin{proof}
Pour $\Sch_{fppf}$, cela est vrai par construction ; voir
Ensembles, lemme \ref{sets-lemma-what-is-in-it}.
Supposons donnés des morphismes $U \to S$, $V \to U$, $W \to U$
de schémas, avec $U, V, W \in \Ob(\Sch_{fppf})$.
Le produit fibré $V \times_U W$ dans $\Sch_{fppf}$
est un produit fibré dans $\Sch$ et
est le produit fibré de $V/S$ et $W/S$ au-dessus de $U/S$ dans
la catégorie de tous les schémas au-dessus de $S$, donc aussi un
produit fibré dans $(\Sch/S)_{fppf}$.
Cela prouve le résultat pour $(\Sch/S)_{fppf}$.
Si $U, V, W$ sont affines, $V \times_U W$ l'est aussi, d'où le
résultat pour $(\textit{Aff}/S)_{fppf}$.
\end{proof}

\noindent
Vérifions ensuite que le gros site affine définit le même
topos que le gros site.

\begin{lemma}
\label{topologies-lemma-affine-big-site-fppf}
Soit $S$ un schéma. Soit $\Sch_{fppf}$ un gros site
fppf contenant $S$.
Le foncteur $(\textit{Aff}/S)_{fppf} \to (\Sch/S)_{fppf}$
est cocontinu et induit une équivalence de topos de
$\Sh((\textit{Aff}/S)_{fppf})$ vers
$\Sh((\Sch/S)_{fppf})$.
\end{lemma}

\begin{proof}
La notion de foncteur cocontinu spécial est introduite dans
Sites, définition \ref{sites-definition-special-cocontinuous-functor}.
Il faut donc vérifier les hypothèses (1) -- (5) de
Sites, lemme \ref{sites-lemma-equivalence}.
Notons le foncteur d'inclusion
$u : (\textit{Aff}/S)_{fppf} \to (\Sch/S)_{fppf}$.
La cocontinuité signifie simplement que tout recouvrement fppf de
$T/S$, avec $T$ affine, admet un raffinement par un recouvrement fppf standard de $T$.
C'est le contenu du
lemme \ref{topologies-lemma-fppf-affine}.
Ainsi, (1) est vérifiée. Le foncteur $u$ est continu simplement parce qu'un recouvrement
fppf standard est un recouvrement fppf. Ainsi, (2) est vérifiée.
Les points (3) et (4) découlent immédiatement du fait que $u$ est
pleinement fidèle. Enfin, la condition (5) découle du
fait que tout schéma admet un recouvrement ouvert affine.
\end{proof}

\noindent
Établissons maintenant quelques relations entre les topos
associés à ces sites.

\begin{lemma}
\label{topologies-lemma-morphism-big-fppf}
Soit $\Sch_{fppf}$ un gros site fppf.
Soit $f : T \to S$ un morphisme dans $\Sch_{fppf}$.
Le foncteur
$$
u : (\Sch/T)_{fppf} \longrightarrow (\Sch/S)_{fppf},
\quad
V/T \longmapsto V/S
$$
est cocontinu et admet un adjoint à droite continu
$$
v : (\Sch/S)_{fppf} \longrightarrow (\Sch/T)_{fppf},
\quad
(U \to S) \longmapsto (U \times_S T \to T).
$$
Ils induisent le même morphisme de topos
$$
f_{big} :
\Sh((\Sch/T)_{fppf})
\longrightarrow
\Sh((\Sch/S)_{fppf})
$$
On a $f_{big}^{-1}(\mathcal{G})(U/T) = \mathcal{G}(U/S)$.
On a $f_{big, *}(\mathcal{F})(U/S) = \mathcal{F}(U \times_S T/T)$.
En outre, $f_{big}^{-1}$ admet un adjoint à gauche $f_{big!}$ qui commute aux
produits fibrés et aux égalisateurs.
\end{lemma}

\begin{proof}
Le foncteur $u$ est cocontinu, continu et commute aux produits fibrés
et aux égalisateurs. Par conséquent,
Sites, lemmes \ref{sites-lemma-when-shriek} et
\ref{sites-lemma-preserve-equalizers}
s'appliquent et donnent la formule
pour $f_{big}^{-1}$ ainsi que l'existence de $f_{big!}$. De plus,
le foncteur $v$ est un adjoint à droite car, étant donnés $U/T$ et $V/S$,
on a $\Mor_S(u(U), V) = \Mor_T(U, V \times_S T)$,
comme voulu. On peut donc appliquer
Sites, lemmes \ref{sites-lemma-have-functor-other-way} et
\ref{sites-lemma-have-functor-other-way-morphism} pour obtenir la
formule pour $f_{big, *}$.
\end{proof}

\begin{lemma}
\label{topologies-lemma-composition-fppf}
Étant donnés des schémas $X$, $Y$, $Z$ dans $(\Sch/S)_{fppf}$
et des morphismes $f : X \to Y$, $g : Y \to Z$, on a
$g_{big} \circ f_{big} = (g \circ f)_{big}$.
\end{lemma}

\begin{proof}
Cela résulte de la description simple de l'image directe
et de l'image inverse pour les foncteurs sur les gros sites donnée par le
lemme \ref{topologies-lemma-morphism-big-fppf}.
\end{proof}



















































\section{La topologie ph}
\label{topologies-section-ph}

\noindent
Dans cette section, nous définissons la topologie ph. C'est la topologie
engendrée par les recouvrements de Zariski et les morphismes propres surjectifs, voir
le lemme \ref{topologies-lemma-characterize-sheaf}.

\medskip\noindent
Nous reprenons nos notations et notre terminologie
de l'article \cite{ph} de Goodwillie et Lichtenbaum.
Ces auteurs montrent que, si l'on se restreint à la sous-catégorie des
schémas noethériens, la topologie ph coïncide avec la
« topologie h » telle que Voevodsky l'a définie à l'origine : c'est
la topologie engendrée par les recouvrements ouverts de Zariski et les
morphismes de type fini qui sont universellement submersifs.
Ils montrent aussi que les deux topologies ne coïncident pas sur les
schémas non noethériens, voir \cite[exemple 4.5]{ph}.
Nous revenons à (notre version de) la topologie h dans
Compléments sur la platitude, section \ref{flat-section-h}.

\medskip\noindent
Avant de pouvoir définir les recouvrements de notre topologie, il nous faut
quelques préliminaires.

\begin{definition}
\label{topologies-definition-standard-ph-covering}
Soit $T$ un schéma affine. Un {\it recouvrement ph standard} est une famille
$\{f_j : U_j \to T\}_{j = 1, \ldots, m}$ construite à partir d'un
morphisme propre surjectif $f : U \to T$ et d'un recouvrement ouvert affine
$U = \bigcup_{j = 1, \ldots, m} U_j$, en posant $f_j = f|_{U_j}$.
\end{definition}

\noindent
Il résulte immédiatement du lemme de Chow que l'on peut raffiner un
recouvrement ph standard par un recouvrement ph standard correspondant à
un morphisme projectif surjectif.

\begin{lemma}
\label{topologies-lemma-base-change-standard-ph}
Soit $\{f_j : U_j \to T\}_{j = 1, \ldots, m}$ un recouvrement ph standard.
Soit $T' \to T$ un morphisme de schémas affines.
Alors $\{U_j \times_T T' \to T'\}_{j = 1, \ldots, m}$ est un
recouvrement ph standard.
\end{lemma}

\begin{proof}
Soient $f : U \to T$ un morphisme propre surjectif et
$U = \bigcup_{j = 1, \ldots, m} U_j$ un recouvrement ouvert affine comme dans la
définition \ref{topologies-definition-standard-ph-covering}.
Alors $U \times_T T' \to T'$ est propre et surjectif
(Morphismes, lemmes \ref{morphisms-lemma-base-change-surjective} et
\ref{morphisms-lemma-base-change-proper}).
De plus, $U \times_T T' = \bigcup_{j = 1, \ldots, m} U_j \times_T T'$
est un recouvrement ouvert affine.
Ceci achève la démonstration.
\end{proof}

\begin{lemma}
\label{topologies-lemma-refine-by-standard-ph}
Soit $T$ un schéma affine. Chacun des types de familles suivants
de morphismes de but $T$ admet un raffinement qui est un recouvrement ph standard :
\begin{enumerate}
\item tout recouvrement ouvert de Zariski de $T$,
\item $\{W_{ji} \to T\}_{j = 1, \ldots, m, i = 1, \ldots n_j}$
où $\{W_{ji} \to U_j\}_{i = 1, \ldots, n_j}$
et $\{U_j \to T\}_{j = 1, \ldots, m}$ sont des recouvrements ph standards.
\end{enumerate}
\end{lemma}

\begin{proof}
L'assertion (1) résulte du fait que tout recouvrement ouvert de Zariski de $T$
peut être raffiné par un recouvrement ouvert affine fini.

\medskip\noindent
Preuve de (3). Choisissons $U \to T$ propre et surjectif et
$U = \bigcup_{j = 1, \ldots, m} U_j$ comme dans la
définition \ref{topologies-definition-standard-ph-covering}.
Choisissons $W_j \to U_j$ propre et surjectif et $W_j = \bigcup W_{ji}$
comme dans la définition \ref{topologies-definition-standard-ph-covering}.
Par le lemme de Chow (Limits, lemme \ref{limits-lemma-chow-finite-type}),
il existe $W'_j \to W_j$ propre et surjectif
et des immersions fermées $W'_j \to \mathbf{P}^{e_j}_{U_j}$.
Ainsi, après avoir remplacé $W_j$ par $W'_j$ et $W_j = \bigcup W_{ji}$
par un recouvrement ouvert affine convenable de $W'_j$, on peut supposer
qu'il existe une immersion fermée $W_j \subset \mathbf{P}^{e_j}_{U_j}$
pour tout $j = 1, \ldots, m$.

\medskip\noindent
Soit $\overline{W}_j \subset \mathbf{P}^{e_j}_U$ l'adhérence
schématique de $W_j$. Alors $W_j \subset \overline{W}_j$
est un sous-schéma ouvert ; en fait, $W_j$ est l'image inverse de
$U_j \subset U$ par le morphisme $\overline{W}_j \to U$.
(Pour le voir, utilisons le fait que $W_j \to \mathbf{P}^{e_j}_U$ est quasi-compact
et que, par suite, la formation de l'image schématique commute
à la restriction aux ouverts, voir
Morphismes, section \ref{morphisms-section-scheme-theoretic-image}.)
Munissons $Z_j = U \setminus U_j$
de la structure de sous-schéma fermé réduit induite.
Alors
$$
V_j = \overline{W}_j \amalg Z_j \to U
$$
est propre et surjectif, et le sous-schéma ouvert $W_j \subset V_j$
est l'image inverse de $U_j$. Ainsi, pour $v \in V_j$, $v \not \in W_j$,
on peut choisir un voisinage ouvert affine $v \in V_{j, v} \subset V_j$
dont l'image est contenue dans $U_{j'}$ pour un certain $1 \leq j' \leq m$.

\medskip\noindent
Pour achever la démonstration, considérons le morphisme propre surjectif
$$
V = V_1 \times_U V_2 \times_U \ldots \times_U V_m
\longrightarrow U \longrightarrow T
$$
et le recouvrement de $V$ par les ouverts affines
$$
V_{1, v_1} \times_U \ldots \times_U V_{j - 1, v_{j - 1}}
\times_U W_{j i} \times_U
V_{j + 1, v_{j + 1}} \times_U \ldots \times_U V_{m, v_m}
$$
Ceux-ci forment bien un recouvrement, car tout point de $U$ appartient
à un certain $U_j$ et l'image inverse de $U_j$ dans $V$ est égale à
$V_1 \times \ldots \times V_{j - 1} \times W_j \times V_{j + 1} \times
\ldots \times V_m$. Remarquons que le morphisme de
l'ouvert affine ci-dessus vers $T$ se factorise par $W_{ji}$ ;
on obtient donc un raffinement. Enfin, il suffit d'un nombre fini
de ces ouverts affines puisque $V$ est quasi-compact (comme schéma propre
sur le schéma affine $T$).
\end{proof}

\begin{definition}
\label{topologies-definition-ph-covering}
Soit $T$ un schéma. Un {\it recouvrement ph de $T$} est une famille
de morphismes de schémas $\{f_i : T_i \to T\}_{i \in I}$ telle
que $f_i$ soit localement de type fini et que, pour tout
ouvert affine $U \subset T$, il existe un recouvrement ph standard
$\{U_j \to U\}_{j = 1, \ldots, m}$ qui raffine la famille
$\{T_i \times_T U \to U\}_{i \in I}$.
\end{definition}

\noindent
Un recouvrement ph standard est un recouvrement ph d'après le
lemme \ref{topologies-lemma-base-change-standard-ph}.

\begin{lemma}
\label{topologies-lemma-zariski-ph}
Un recouvrement de Zariski est un recouvrement ph\footnote{Nous verrons
dans Compléments sur les morphismes, lemme \ref{more-morphisms-lemma-fppf-ph}, que les
recouvrements fppf (et donc les recouvrements syntomiques, lisses ou \'etales)
sont eux aussi des recouvrements ph.}.
\end{lemma}

\begin{proof}
Cela résulte du fait qu'un recouvrement de Zariski d'un schéma affine
peut être raffiné par un recouvrement ph standard d'après le
lemme \ref{topologies-lemma-refine-by-standard-ph}.
\end{proof}

\begin{lemma}
\label{topologies-lemma-surjective-proper-ph}
Soit $f : Y \to X$ un morphisme propre surjectif de schémas.
Alors $\{Y \to X\}$ est un recouvrement ph.
\end{lemma}

\begin{proof}
Omise.
\end{proof}

\begin{lemma}
\label{topologies-lemma-refine-by-ph}
Soit $T$ un schéma. Soit $\{f_i : T_i \to T\}_{i \in I}$ une famille
de morphismes telle que $f_i$ soit localement de type fini pour tout $i$.
Les conditions suivantes sont équivalentes :
\begin{enumerate}
\item $\{T_i \to T\}_{i \in I}$ est un recouvrement ph,
\item il existe un recouvrement ph qui raffine $\{T_i \to T\}_{i \in I}$, et
\item $\{\coprod_{i \in I} T_i \to T\}$ est un recouvrement ph.
\end{enumerate}
\end{lemma}

\begin{proof}
L'équivalence de (1) et (2) résulte immédiatement de la
définition \ref{topologies-definition-ph-covering}
et du fait qu'un raffinement d'un raffinement est un raffinement.
En vertu de l'équivalence de (1) et (2), et puisque
$\{T_i \to T\}_{i \in I}$ raffine $\{\coprod_{i \in I} T_i \to T\}$,
on voit que (1) implique (3). Supposons enfin que (3) soit vérifiée.
Soit $U \subset T$ un ouvert affine et soit
$\{U_j \to U\}_{j = 1, \ldots, m}$ un recouvrement ph standard
qui raffine $\{U \times_T \coprod_{i \in I} T_i \to U\}$.
Cela signifie que, pour tout $j$, on dispose d'un morphisme
$$
h_j :
U_j
\longrightarrow
U \times_T \coprod\nolimits_{i \in I} T_i =
\coprod\nolimits_{i \in I} U \times_T T_i
$$
au-dessus de $U$. Comme $U_j$ est quasi-compact, on obtient
des décompositions en réunion disjointe $U_j = \coprod_{i \in I} U_{j, i}$
en sous-schémas ouverts et fermés, presque tous vides,
telles que $h_j|_{U_{j, i}}$ envoie $U_{j, i}$ dans $U \times_T T_i$.
Il s'ensuit que
$$
\{U_{j, i} \to U\}_{j = 1, \ldots, m,\ i \in I,\ U_{j, i} \not = \emptyset}
$$
est un recouvrement ph standard (un petit détail est omis) qui raffine
$\{U \times_T T_i \to U\}_{i \in I}$. Ainsi, (1) est vérifiée.
\end{proof}

\noindent
Montrons ensuite que cette notion satisfait les conditions de
Sites, définition \ref{sites-definition-site}.

\begin{lemma}
\label{topologies-lemma-ph}
Soit $T$ un schéma.
\begin{enumerate}
\item Si $T' \to T$ est un isomorphisme, alors $\{T' \to T\}$
est un recouvrement ph de $T$.
\item Si $\{T_i \to T\}_{i\in I}$ est un recouvrement ph et si, pour tout
$i$, on a un recouvrement ph $\{T_{ij} \to T_i\}_{j\in J_i}$, alors
$\{T_{ij} \to T\}_{i \in I, j\in J_i}$ est un recouvrement ph.
\item Si $\{T_i \to T\}_{i\in I}$ est un recouvrement ph
et si $T' \to T$ est un morphisme de schémas, alors
$\{T' \times_T T_i \to T'\}_{i\in I}$ est un recouvrement ph.
\end{enumerate}
\end{lemma}

\begin{proof}
L'assertion (1) est claire.

\medskip\noindent
Preuve de (3). Le changement de base $T_i \times_T T' \to T'$ est
localement de type fini d'après
Morphismes, lemme \ref{morphisms-lemma-base-change-finite-type}.
Il suffit donc de vérifier la condition sur les ouverts affines.
Soit $U' \subset T'$ un sous-schéma ouvert affine.
Comme $U'$ est quasi-compact, on peut trouver un recouvrement ouvert
affine fini $U' = U'_1 \cup \ldots \cup U'$ tel que
l'image de $U'_j \to T$ soit contenue dans un ouvert affine $U_j \subset T$.
Choisissons un recouvrement ph standard $\{U_{jl} \to U_j\}_{l = 1, \ldots, n_j}$
qui raffine $\{T_i \times_T U_j \to U_j\}$.
D'après le lemme \ref{topologies-lemma-base-change-standard-ph},
le changement de base $\{U_{jl} \times_{U_j} U'_j \to U'_j\}$
est un recouvrement ph standard. Remarquons que $\{U'_j \to U'\}$
est lui aussi un recouvrement ph standard.
D'après le lemme \ref{topologies-lemma-refine-by-standard-ph},
la famille $\{U_{jl} \times_{U_j} U'_j \to U'\}$
peut être raffinée par un recouvrement ph standard. Comme
$\{U_{jl} \times_{U_j} U'_j \to U'\}$ raffine $\{T_i \times_T U' \to U'\}$,
on conclut.

\medskip\noindent
Preuve de (2). La composition préserve la propriété d'être localement de type fini,
voir Morphismes, lemme \ref{morphisms-lemma-composition-finite-type}.
Il suffit donc de vérifier la condition sur les ouverts affines.
Soit $U \subset T$ un ouvert affine. Choisissons d'abord un recouvrement ph standard
$\{U_k \to U\}_{k = 1, \ldots, m}$ qui raffine
$\{T_i \times_T U \to U\}$.
Supposons que le raffinement soit donné par des morphismes $U_k \to T_{i_k}$ sur $T$.
Alors
$$
\{T_{i_kj} \times_{T_{i_k}} U_k \to U_k\}_{j \in J_{i_k}}
$$
est un recouvrement ph d'après l'assertion (3). Comme $U_k$ est affine,
on peut trouver un recouvrement ph standard
$\{U_{ka} \to U_k\}_{a = 1, \ldots, b_k}$ qui raffine cette famille.
On applique alors le lemme \ref{topologies-lemma-refine-by-standard-ph}
pour voir que $\{U_{ka} \to U\}$ peut être raffiné par un
recouvrement ph standard. Comme $\{U_{ka} \to U\}$
raffine $\{T_{ij} \times_T U \to U\}$, ceci achève la démonstration.
\end{proof}

\begin{definition}
\label{topologies-definition-big-ph-site}
Un {\it gros site ph} est un site $\Sch_{ph}$ comme dans
Sites, définition \ref{sites-definition-site}, construit comme suit :
\begin{enumerate}
\item Choisir un ensemble quelconque de schémas $S_0$ et un ensemble quelconque
$\text{Cov}_0$ de recouvrements ph entre ces schémas.
\item Prendre pour catégorie sous-jacente une catégorie quelconque $\Sch_\alpha$
construite comme dans Ensembles, lemme \ref{sets-lemma-construct-category},
à partir de l'ensemble $S_0$.
\item Choisir un ensemble quelconque de recouvrements comme dans
Ensembles, lemme \ref{sets-lemma-coverings-site}, à partir de la
catégorie $\Sch_\alpha$, de la classe des recouvrements ph
et de l'ensemble $\text{Cov}_0$ choisi ci-dessus.
\end{enumerate}
\end{definition}

\noindent
Voir les remarques qui suivent la définition \ref{topologies-definition-big-zariski-site}
pour la motivation et les explications concernant la définition des gros sites.

\medskip\noindent
Avant de poursuivre l'introduction du gros site ph d'un
schéma $S$, signalons que la topologie d'un gros site ph
$\Sch_{ph}$ est, en un certain sens, induite par la topologie ph
sur la catégorie de tous les schémas.

\begin{lemma}
\label{topologies-lemma-ph-induced}
Soit $\Sch_{ph}$ un gros site ph comme dans la
définition \ref{topologies-definition-big-ph-site}.
Soit $T \in \Ob(\Sch_{ph})$.
Soit $\{T_i \to T\}_{i \in I}$ un recouvrement ph quelconque de $T$.
\begin{enumerate}
\item Il existe un recouvrement $\{U_j \to T\}_{j \in J}$ de $T$ dans le site
$\Sch_{ph}$ qui raffine $\{T_i \to T\}_{i \in I}$.
\item Si $\{T_i \to T\}_{i \in I}$ est un recouvrement ph standard, alors
il est tautologiquement équivalent à un recouvrement de $\Sch_{ph}$.
\item Si $\{T_i \to T\}_{i \in I}$ est un recouvrement de Zariski, alors
il est tautologiquement équivalent à un recouvrement de $\Sch_{ph}$.
\end{enumerate}
\end{lemma}

\begin{proof}
Pour tout $i$, choisissons un recouvrement ouvert affine $T_i = \bigcup_{j \in J_i} T_{ij}$
tel que chaque $T_{ij}$ s'envoie dans un sous-schéma ouvert affine de $T$. D'après les
lemmes \ref{topologies-lemma-zariski-ph} et \ref{topologies-lemma-ph},
le raffinement $\{T_{ij} \to T\}_{i \in I, j \in J_i}$ est lui aussi un recouvrement ph
de $T$. On peut donc supposer que chaque $T_i$ est affine et s'envoie dans
un ouvert affine $W_i$ de $T$. En appliquant
Ensembles, lemme \ref{sets-lemma-what-is-in-it}
on voit que $W_i$ est isomorphe à un objet de $\Sch_{ph}$.
Mais alors $T_i$, comme schéma de type fini sur $W_i$,
est isomorphe à un objet $V_i$ de $\Sch_{ph}$ par une seconde
application de
Ensembles, lemme \ref{sets-lemma-what-is-in-it}.
Le recouvrement $\{V_i \to T\}_{i \in I}$ raffine $\{T_i \to T\}_{i \in I}$
(car ils sont isomorphes).
De plus, $\{V_i \to T\}_{i \in I}$ est combinatoirement équivalent à un
recouvrement $\{U_j \to T\}_{j \in J}$ de $T$ dans le site
$\Sch_{ph}$ d'après
Ensembles, lemme \ref{sets-lemma-what-is-in-it}.
Le recouvrement $\{U_j \to T\}_{j \in J}$ est un raffinement comme en (1).
Dans les situations de (2) et (3), chacun des
schémas $T_i$ est isomorphe à un objet de $\Sch_{ph}$ d'après
Ensembles, lemme \ref{sets-lemma-what-is-in-it},
et une autre application de
Ensembles, lemme \ref{sets-lemma-coverings-site}
donne le résultat voulu.
\end{proof}

\begin{definition}
\label{topologies-definition-big-small-ph}
Soit $S$ un schéma. Soit $\Sch_{ph}$ un gros site ph contenant $S$.
\begin{enumerate}
\item Le {\it gros site ph de $S$}, noté
$(\Sch/S)_{ph}$, est le site $\Sch_{ph}/S$
introduit dans Sites, section \ref{sites-section-localize}.
\item Le {\it gros site ph affine de $S$}, noté
$(\textit{Aff}/S)_{ph}$, est la sous-catégorie pleine de
$(\Sch/S)_{ph}$ dont les objets sont les $U/S$ affines.
Un recouvrement de $(\textit{Aff}/S)_{ph}$ est tout recouvrement fini
$\{U_i \to U\}$ de $(\Sch/S)_{ph}$ tel que $U_i$ et $U$ soient affines.
\end{enumerate}
\end{definition}

\noindent
Remarquons que les recouvrements de $(\textit{Aff}/S)_{ph}$ ne sont {\it pas}
donnés par les recouvrements ph standards. La raison en est simplement que
le deuxième axiome de Sites, définition \ref{sites-definition-site}, échouerait.
Les recouvrements de $(\textit{Aff}/S)_{ph}$ sont plutôt les familles finies
$\{U_i \to U\}$ de morphismes de type fini entre objets affines
de $(\Sch/S)_{ph}$ qui peuvent être raffinées par un recouvrement ph standard.
Énonçons et démontrons explicitement que le gros site ph affine est un site.

\begin{lemma}
\label{topologies-lemma-verify-site-ph}
Soit $S$ un schéma. Soit $\Sch_{ph}$ un gros site ph
contenant $S$. Alors $(\textit{Aff}/S)_{ph}$ est un site.
\end{lemma}

\begin{proof}
En raisonnant comme dans la démonstration du lemme \ref{topologies-lemma-verify-site-etale},
il suffit de montrer que la famille des recouvrements ph finis
$\{U_i \to U\}$, où $U$ et les $U_i$ sont affines, satisfait les propriétés
(1), (2) et (3) de
Sites, définition \ref{sites-definition-site}.
C'est clair car, par exemple, étant donné un recouvrement ph fini
$\{T_i \to T\}_{i\in I}$ où $T_i, T$ sont affines, et, pour tout
$i$, un recouvrement ph fini $\{T_{ij} \to T_i\}_{j\in J_i}$ où $T_{ij}$ est
affine, $\{T_{ij} \to T\}_{i \in I, j\in J_i}$ est alors un recouvrement ph
(lemme \ref{topologies-lemma-ph}), $\bigcup_{i\in I} J_i$ est fini et
chaque $T_{ij}$ est affine.
\end{proof}

\begin{lemma}
\label{topologies-lemma-fibre-products-ph}
Soit $S$ un schéma. Soit $\Sch_{ph}$ un gros site ph
contenant $S$. Les catégories sous-jacentes aux sites
$\Sch_{ph}$, $(\Sch/S)_{ph}$ et $(\textit{Aff}/S)_{ph}$ ont des produits fibrés.
Dans chaque cas, le foncteur évident vers la catégorie $\Sch$ de
tous les schémas commute aux produits fibrés. La catégorie
$(\Sch/S)_{ph}$ a un objet final, à savoir $S/S$.
\end{lemma}

\begin{proof}
Pour $\Sch_{ph}$, c'est vrai par construction, voir
Ensembles, lemme \ref{sets-lemma-what-is-in-it}.
Supposons donnés des morphismes de schémas $U \to S$, $V \to U$, $W \to U$
avec $U, V, W \in \Ob(\Sch_{ph})$.
Le produit fibré $V \times_U W$ dans $\Sch_{ph}$
est un produit fibré dans $\Sch$ et
est le produit fibré de $V/S$ et $W/S$ au-dessus de $U/S$ dans
la catégorie de tous les schémas au-dessus de $S$, et donc aussi un
produit fibré dans $(\Sch/S)_{ph}$.
Ceci démontre le résultat pour $(\Sch/S)_{ph}$.
Si $U, V, W$ sont affines, $V \times_U W$ l'est aussi, d'où le
résultat pour $(\textit{Aff}/S)_{ph}$.
\end{proof}

\noindent
Vérifions ensuite que le gros site affine définit le même
topos que le gros site.

\begin{lemma}
\label{topologies-lemma-affine-big-site-ph}
Soit $S$ un schéma. Soit $\Sch_{ph}$ un gros site ph
contenant $S$.
Le foncteur $(\textit{Aff}/S)_{ph} \to (\Sch/S)_{ph}$
est cocontinu et induit une équivalence de topos de
$\Sh((\textit{Aff}/S)_{ph})$ vers
$\Sh((\Sch/S)_{ph})$.
\end{lemma}

\begin{proof}
La notion de foncteur cocontinu spécial est introduite dans
Sites, définition \ref{sites-definition-special-cocontinuous-functor}.
Il faut donc vérifier les hypothèses (1) -- (5) de
Sites, lemme \ref{sites-lemma-equivalence}.
Notons le foncteur d'inclusion
$u : (\textit{Aff}/S)_{ph} \to (\Sch/S)_{ph}$.
La cocontinuité résulte du fait que tout recouvrement ph de
$T/S$, avec $T$ affine, peut être raffiné par un recouvrement ph standard de $T$
par définition. Ainsi, (1) est vérifiée. Le foncteur $u$ est continu simplement
parce qu'un recouvrement ph fini d'un affine par des affines est un recouvrement ph.
Ainsi, (2) est vérifiée.
Les assertions (3) et (4) résultent immédiatement du fait que $u$ est
pleinement fidèle. Enfin, la condition (5) résulte du
fait que tout schéma possède un recouvrement ouvert affine (qui est
un recouvrement ph).
\end{proof}

\begin{lemma}
\label{topologies-lemma-characterize-sheaf}
Soit $\mathcal{F}$ un préfaisceau sur $(\Sch/S)_{ph}$.
Alors $\mathcal{F}$ est un faisceau si et seulement si
\begin{enumerate}
\item $\mathcal{F}$ satisfait la condition de faisceau pour les
recouvrements de Zariski, et
\item si $f : V \to U$ est propre et surjectif, alors
$\mathcal{F}(U)$ s'envoie bijectivement sur l'égalisateur
des deux applications $\mathcal{F}(V) \to \mathcal{F}(V \times_U V)$.
\end{enumerate}
De plus, en présence de (1), la propriété (2) équivaut à la propriété
\begin{enumerate}
\item[(2')] de faisceau pour $\{V \to U\}$ comme en (2), avec $U$ affine.
\end{enumerate}
\end{lemma}

\begin{proof}
Montrons que, si (1) et (2) sont vérifiées, alors $\mathcal{F}$ est un faisceau.
Soit $\{T_i \to T\}$ un recouvrement ph, c'est-à-dire un recouvrement dans $(\Sch/S)_{ph}$.
Vérifions la condition de faisceau pour ce recouvrement.
Soient $s_i \in \mathcal{F}(T_i)$ des sections dont les restrictions donnent la même
section sur $T_i \times_T T_{i'}$. Montrons qu'il existe une
unique section $s \in \mathcal{F}(T)$ dont la restriction est $s_i$ sur $T_i$.
Soit $T = \bigcup U_j$ un recouvrement ouvert affine.
D'après la propriété (1), il suffit de produire des sections $s_j \in \mathcal{F}(U_j)$
qui coïncident sur $U_j \cap U_{j'}$ pour construire $s$.
Considérons les recouvrements ph $\{T_i \times_T U_j \to U_j\}$.
Alors les $s_{ji} = s_i|_{T_i \times_T U_j}$ sont des sections qui coïncident
sur $(T_i \times_T U_j) \times_{U_j} (T_{i'} \times_T U_j)$.
Choisissons un morphisme propre surjectif $V_j \to U_j$ et un recouvrement ouvert
affine fini $V_j = \bigcup V_{jk}$ tel que le recouvrement ph standard
$\{V_{jk} \to U_j\}$ raffine $\{T_i \times_T U_j \to U_j\}$.
Si $s_{jk} \in \mathcal{F}(V_{jk})$
désigne l'image inverse de $s_{ji}$ sur $V_{jk}$ par les
morphismes sous-entendus, les $s_{jk}$ se recollent en une section
$s'_j \in \mathcal{F}(V_j)$. En utilisant encore une fois l'accord sur les intersections,
on voit que $s'_j$ appartient à l'égalisateur des deux
applications $\mathcal{F}(V_j) \to \mathcal{F}(V_j \times_{U_j} V_j)$.
Ainsi, d'après (2), $s'_j$ provient d'une unique section
$s_j \in \mathcal{F}(U_j)$. Nous omettons de vérifier que ces
sections $s_j$ possèdent toutes les propriétés voulues.

\medskip\noindent
Preuve de l'équivalence de (2) et (2') en présence de (1).
Supposons que $V \to U$ soit un morphisme de $(\Sch/S)_{ph}$ qui soit
propre et surjectif. Choisissons un
recouvrement ouvert affine $U = \bigcup U_i$ et posons $V_i = V \times_U U_i$.
On voit alors que $\mathcal{F}(U) \to \mathcal{F}(V)$
est injective, car $\mathcal{F}(U_i) \to \mathcal{F}(V_i)$
est injective d'après (2') et $\mathcal{F}(U) \to \prod \mathcal{F}(U_i)$
est injective d'après (1). Supposons enfin donné un élément
$t \in \mathcal{F}(V)$ de l'égalisateur des deux applications
$\mathcal{F}(V) \to \mathcal{F}(V \times_U V)$.
Alors $t|_{V_i}$ appartient à l'égalisateur des deux applications
$\mathcal{F}(V_i) \to \mathcal{F}(V_i \times_{U_i} V_i)$
pour tout $i$. On obtient donc une unique section $s_i \in \mathcal{F}(U_i)$
qui s'envoie sur $t|_{V_i}$ pour tout $i$ d'après (2').
Nous omettons de vérifier que $s_i|_{U_i \cap U_j} = s_j|_{U_i \cap U_j}$
pour tous $i, j$ ; cela utilise la propriété d'unicité que l'on vient d'établir.
La propriété de faisceau pour le recouvrement $U = \bigcup U_i$ fournit
une section $s \in \mathcal{F}(U)$. Nous omettons de démontrer que $s$
s'envoie sur $t$ dans $\mathcal{F}(V)$.
\end{proof}

\noindent
Établissons ensuite quelques relations entre les topos
associés à ces sites.

\begin{lemma}
\label{topologies-lemma-morphism-big-ph}
Soit $\Sch_{ph}$ un gros site ph.
Soit $f : T \to S$ un morphisme de $\Sch_{ph}$.
Le foncteur
$$
u : (\Sch/T)_{ph} \longrightarrow (\Sch/S)_{ph},
\quad
V/T \longmapsto V/S
$$
est cocontinu et admet un adjoint à droite continu
$$
v : (\Sch/S)_{ph} \longrightarrow (\Sch/T)_{ph},
\quad
(U \to S) \longmapsto (U \times_S T \to T).
$$
Ils induisent le même morphisme de topos
$$
f_{big} :
\Sh((\Sch/T)_{ph})
\longrightarrow
\Sh((\Sch/S)_{ph})
$$
On a $f_{big}^{-1}(\mathcal{G})(U/T) = \mathcal{G}(U/S)$.
On a $f_{big, *}(\mathcal{F})(U/S) = \mathcal{F}(U \times_S T/T)$.
De plus, $f_{big}^{-1}$ admet un adjoint à gauche $f_{big!}$ qui commute aux
produits fibrés et aux égalisateurs.
\end{lemma}

\begin{proof}
Le foncteur $u$ est cocontinu, continu et commute aux produits fibrés
et aux égalisateurs. Par conséquent,
Sites, lemmes \ref{sites-lemma-when-shriek} et
\ref{sites-lemma-preserve-equalizers}
s'appliquent et donnent la formule
pour $f_{big}^{-1}$ ainsi que l'existence de $f_{big!}$. De plus,
le foncteur $v$ est adjoint à droite car, étant donnés $U/T$ et $V/S$,
on a $\Mor_S(u(U), V) = \Mor_T(U, V \times_S T)$,
comme voulu. On peut donc appliquer
Sites, lemmes \ref{sites-lemma-have-functor-other-way} et
\ref{sites-lemma-have-functor-other-way-morphism} pour obtenir la
formule pour $f_{big, *}$.
\end{proof}

\begin{lemma}
\label{topologies-lemma-composition-ph}
Étant donnés des schémas $X$, $Y$, $Y$ dans $(\Sch/S)_{ph}$
et des morphismes $f : X \to Y$, $g : Y \to Z$, on a
$g_{big} \circ f_{big} = (g \circ f)_{big}$.
\end{lemma}

\begin{proof}
Cela résulte de la description simple de l'image directe
et de l'image inverse des foncteurs sur les gros sites donnée dans le
lemme \ref{topologies-lemma-morphism-big-ph}.
\end{proof}



































\section{La topologie fpqc}
\label{topologies-section-fpqc}

\begin{definition}
\label{topologies-definition-fpqc-covering}
Soit $T$ un schéma. Un {\it recouvrement fpqc de $T$} est une famille
de morphismes de schémas $\{f_i : T_i \to T\}_{i \in I}$
telle que chaque $f_i$ soit plat et que, pour tout ouvert affine
$U \subset T$, il existe $n \geq 0$, une application
$a : \{1, \ldots, n\} \to I$ et des ouverts affines
$V_j \subset T_{a(j)}$, $j = 1, \ldots, n$,
tels que $\bigcup_{j = 1}^n f_{a(j)}(V_j) = U$.
\end{definition}

\noindent
Cette condition implique bien que $T = \bigcup f_i(T_i)$.
Il est un peu plus difficile de reconnaître un recouvrement fpqc ; nous donnons donc
quelques lemmes à cet effet.

\begin{lemma}
\label{topologies-lemma-recognize-fpqc-covering}
Soit $T$ un schéma. Soit $\{f_i : T_i \to T\}_{i \in I}$ une famille de
morphismes de schémas de but $T$. Les conditions suivantes sont équivalentes :
\begin{enumerate}
\item $\{f_i : T_i \to T\}_{i \in I}$ est un recouvrement fpqc,
\item chaque $f_i$ est plat et, pour tout ouvert affine $U \subset T$,
il existe des ouverts quasi-compacts
$U_i \subset T_i$, presque tous vides,
tels que $U = \bigcup f_i(U_i)$,
\item chaque $f_i$ est plat et il existe un recouvrement ouvert affine
$T = \bigcup_{\alpha \in A} U_\alpha$ et, pour tout $\alpha \in A$,
il existe $i_{\alpha, 1}, \ldots, i_{\alpha, n(\alpha)} \in I$
et des ouverts quasi-compacts $U_{\alpha, j} \subset T_{i_{\alpha, j}}$ tels que
$U_\alpha =
\bigcup_{j = 1, \ldots, n(\alpha)} f_{i_{\alpha, j}}(U_{\alpha, j})$.
\end{enumerate}
Si $T$ est quasi-séparé, ces conditions équivalent aussi à
\begin{enumerate}
\item[(4)] chaque $f_i$ est plat et, pour tout $t \in T$, il existe
$i_1, \ldots, i_n \in I$ et des ouverts quasi-compacts $U_j \subset T_{i_j}$
tels que $\bigcup_{j = 1, \ldots, n} f_{i_j}(U_j)$ soit un
voisinage (non nécessairement ouvert) de $t$ dans $T$.
\end{enumerate}
\end{lemma}

\begin{proof}
Nous omettons la démonstration de l'équivalence de (1), (2) et (3).
Supposons désormais que $T$ soit quasi-séparé.
Montrons que (4) implique (2). Soit $U \subset T$ un ouvert affine.
Pour démontrer (2), il suffit de montrer que, pour tout $t \in U$, il existe
un nombre fini d'ouverts quasi-compacts $U_j \subset T_{i_j}$ tels que
$f_{i_j}(U_j) \subset U$ et que $\bigcup f_{i_j}(U_j)$
soit un voisinage de $t$ dans $U$. Par hypothèse, il existe bien
un nombre fini d'ouverts quasi-compacts $U'_j \subset T_{i_j}$ tels que
$\bigcup f_{i_j}(U'_j)$ soit un voisinage de $t$ dans $T$.
Puisque $T$ est quasi-séparé, on voit que
$U_j = U'_j \cap f_{i_j}^{-1}(U) = (f_{i_j}|_{U'_j})^{-1}(U)$
est quasi-compact d'après
Schémas, lemme \ref{schemes-lemma-quasi-compact-permanence}.
Ainsi, (2) est vérifiée. Comme il est clair que (2) implique
(4), la démonstration est achevée.
\end{proof}

\begin{example}
\label{topologies-example-not-fpqc}
Soit $X = X_1 \cup X_2$, où $X_1 = X_2 = \Spec(k[t_1,t_2,\dots])$
est le schéma non quasi-séparé de
Schémas, exemple \ref{schemes-example-not-quasi-separated}, et soit
$Y = X_1 \amalg \Spec(\mathcal{O}_{X_2, 0})$.
Alors le morphisme évident $f : Y \to X$
est plat et surjectif et, puisque $Y$ est quasi-compact,
satisfait trivialement à (4) du lemme \ref{topologies-lemma-recognize-fpqc-covering}.
Mais $\{f : Y \to X\}$ n'est pas un recouvrement fpqc. En effet, affirmons qu'il
n'existe aucun ouvert quasi-compact $W$ de $Y$ dont l'image par $f$ soit $X_2$.
En effet, $f^{-1}(X_2) = X_1 \setminus \{0\} \amalg \Spec(\mathcal{O}_{X_2, 0})$.
Donc, si $W$ existe, alors $W = U \amalg \Spec(\mathcal{O}_{X_2, 0})$
pour un ouvert quasi-compact $U$ de $X_1 \setminus \{0\} = X_2 \setminus \{0\}$.
Mais on peut alors écrire $U = D(f_1) \cup \ldots \cup D(f_n)$ pour
certains $f_i \in k[t_1, t_2, \ldots]$ qui s'annulent en $0$.
Supposons que $f_1, \ldots, f_n$ n'utilisent que les variables
$x_1, \ldots, x_m$. Alors $p = (0, \ldots, 0, 1, 0, \ldots)$, où $1$
occupe la $(m + 1)$-ième place, est un point fermé de $X_2$ qui n'appartient pas à l'image
de $W$.
\end{example}

\begin{lemma}
\label{topologies-lemma-disjoint-union-is-fpqc-covering}
Soit $T$ un schéma. Soit $\{f_i : T_i \to T\}_{i \in I}$ une famille de
morphismes de schémas de but $T$. Les conditions suivantes sont équivalentes :
\begin{enumerate}
\item $\{f_i : T_i \to T\}_{i \in I}$ est un recouvrement fpqc, et
\item en posant $T' = \coprod_{i \in I} T_i$ et $f = \coprod_{i \in I} f_i$,
la famille $\{f : T' \to T\}$ est un recouvrement fpqc.
\end{enumerate}
\end{lemma}

\begin{proof}
Supposons que $U \subset T$ soit un ouvert affine. Si (1) est vérifiée, il existe
$i_1, \ldots, i_n \in I$ et des ouverts affines $U_j \subset T_{i_j}$ tels que
$U = \bigcup_{j = 1, \ldots, n} f_{i_j}(U_j)$. Alors
$U_1 \amalg \ldots \amalg U_n \subset T'$ est un ouvert quasi-compact dont l'image
est $U$. Ainsi, $\{f : T' \to T\}$ est un recouvrement fpqc d'après le
lemme \ref{topologies-lemma-recognize-fpqc-covering}.
Réciproquement, si (2) est vérifiée, il existe un ouvert quasi-compact
$U' \subset T'$ tel que $U = f(U')$. Alors $U_j = U' \cap T_j$ est un ouvert quasi-compact
de $T_j$ et est vide pour presque tout $j$. D'après le
lemme \ref{topologies-lemma-recognize-fpqc-covering}, on voit que (1) est vérifiée.
\end{proof}

\begin{lemma}
\label{topologies-lemma-family-flat-dominated-covering}
Soit $T$ un schéma. Soit $\{f_i : T_i \to T\}_{i \in I}$ une famille de
morphismes de schémas de but $T$. Supposons que
\begin{enumerate}
\item chaque $f_i$ soit plat, et
\item la famille $\{f_i : T_i \to T\}_{i \in I}$ puisse être raffinée par un
recouvrement fpqc de $T$.
\end{enumerate}
Alors $\{f_i : T_i \to T\}_{i \in I}$ est un recouvrement fpqc de $T$.
\end{lemma}

\begin{proof}
Soit $\{g_j : X_j \to T\}_{j \in J}$ un recouvrement fpqc qui raffine
$\{f_i : T_i \to T\}$. Supposons que $U \subset T$ soit un ouvert affine.
Choisissons $j_1, \ldots, j_m \in J$ et des ouverts affines $V_k \subset X_{j_k}$
tel que $U = \bigcup g_{j_k}(V_k)$. Pour tout $j$, choisissons $i_j \in I$
et un morphisme $h_j : X_j \to T_{i_j}$ tel que $g_j = f_{i_j} \circ h_j$.
Comme $h_{j_k}(V_k)$ est quasi-compact, on peut trouver un ouvert quasi-compact
$h_{j_k}(V_k) \subset U_k \subset f_{i_{j_k}}^{-1}(U)$.
Alors $U = \bigcup f_{i_{j_k}}(U_k)$. On en conclut que
$\{f_i : T_i \to T\}_{i \in I}$ est un recouvrement fpqc d'après le
lemme \ref{topologies-lemma-recognize-fpqc-covering}.
\end{proof}

\begin{lemma}
\label{topologies-lemma-family-flat-fpqc-local-covering}
Soit $T$ un schéma. Soit $\{f_i : T_i \to T\}_{i \in I}$ une famille de
morphismes de schémas de but $T$. Supposons que
\begin{enumerate}
\item chaque $f_i$ soit plat, et
\item il existe un recouvrement fpqc
$\{g_j : S_j \to T\}_{j \in J}$ tel que chacun des
$\{S_j \times_T T_i \to S_j\}_{i \in I}$ soit un recouvrement fpqc.
\end{enumerate}
Alors $\{f_i : T_i \to T\}_{i \in I}$ est un recouvrement fpqc de $T$.
\end{lemma}

\begin{proof}
Nous utiliserons le lemme \ref{topologies-lemma-recognize-fpqc-covering} sans autre
mention. Soit $U \subset T$ un ouvert affine. D'après (2), il existe
des ouverts quasi-compacts $V_j \subset S_j$ pour $j \in J$, presque tous vides, tels que
$U = \bigcup g_j(V_j)$. Pour tout $j$, on peut alors choisir des ouverts quasi-compacts
$W_{ij} \subset S_j \times_T T_i$ pour $i \in I$, presque tous vides,
tels que $V_j = \bigcup_i \text{pr}_1(W_{ij})$. Ainsi,
$\{S_j \times_T T_i \to T\}$ est un recouvrement fpqc.
Comme ce recouvrement raffine $\{f_i : T_i \to T\}$, on conclut par le
lemme \ref{topologies-lemma-family-flat-dominated-covering}.
\end{proof}

\begin{lemma}
\label{topologies-lemma-zariski-etale-smooth-syntomic-fppf-fpqc}
Tout recouvrement fppf est un recouvrement fpqc et, a fortiori,
tout recouvrement syntomique, lisse, \'etale ou de Zariski est un recouvrement fpqc.
\end{lemma}

\begin{proof}
Montrons qu'un recouvrement fppf est un recouvrement fpqc ; le
reste résultera alors du
lemme \ref{topologies-lemma-zariski-etale-smooth-syntomic-fppf}.
Soit $\{f_i : U_i \to U\}_{i \in I}$ un recouvrement fppf.
Par définition, les $f_i$ sont plats, ce qui vérifie la première
condition de la définition \ref{topologies-definition-fpqc-covering}. Pour vérifier la
seconde, soit $V \subset U$ un sous-ensemble ouvert affine.
Écrivons $f_i^{-1}(V) = \bigcup_{j \in J_i} V_{ij}$
pour des ouverts affines $V_{ij} \subset U_i$. Comme chaque $f_i$ est ouvert
(Morphismes, lemme \ref{morphisms-lemma-fppf-open}), on voit que
$V = \bigcup_{i\in I} \bigcup_{j \in J_i} f_i(V_{ij})$
est un recouvrement ouvert de $V$.
Comme $V$ est quasi-compact, ce recouvrement admet un
raffinement fini. Ceci achève la démonstration.
\end{proof}

\noindent
La topologie fpqc\footnote{Les lettres fpqc signifient
« fid\`element plat quasi-compacte ».}
ne peut pas être traitée de la même manière que la topologie
fppf\footnote{Plus précisément, l'analogue du
lemme \ref{topologies-lemma-fppf-induced} pour la topologie fpqc n'est pas valable.}.
En effet, supposons que $R$ soit un anneau non nul. Nous verrons au
lemme \ref{topologies-lemma-no-set-of-fpqc-covers-is-initial}
qu'il n'existe pas d'ensemble $A$ de recouvrements fpqc de $\Spec(R)$
tel que tout recouvrement fpqc puisse être raffiné par un élément de $A$.
Si $R = k$ est un corps, la raison de cette absence de borne est
qu'il n'existe aucune extension de corps de $k$ qui contienne toute extension de corps
de $k$.

\medskip\noindent
Si l'on néglige les difficultés ensemblistes, on rencontre des préfaisceaux
qui n'admettent pas de faisceau associé, voir
\cite[théorème 5.5]{Waterhouse-fpqc-sheafification}.
Une option d'un certain intérêt consiste à ne considérer que les extensions d'anneaux
fidèlement plates $R \to R'$ où le cardinal de $R'$ est convenablement borné.
(Et si l'on considère tous les schémas d'un univers fixé, comme dans SGA4, on
borne le cardinal par un cardinal fortement inaccessible.)
Il n'est toutefois pas très clair de savoir ce qui se passe si l'on remplace ce cardinal
par un cardinal plus grand.

\medskip\noindent
Pour ces raisons, nous n'introduisons pas de sites fpqc et ne considérons pas
la cohomologie relativement à la topologie fpqc.

\medskip\noindent
En revanche, étant donné un foncteur contravariant
$F : \Sch^{opp} \to \textit{Ensembles}$
il est légitime de demander si $F$ satisfait la propriété de faisceau
pour la topologie fpqc, voir ci-dessous.
On peut en outre s'interroger sur la descente des objets
pour la topologie fpqc, etc. En termes simples, pour certains résultats, le bon
degré de généralité consiste à travailler avec des recouvrements fpqc.

\begin{lemma}
\label{topologies-lemma-fpqc}
Soit $T$ un schéma.
\begin{enumerate}
\item Si $T' \to T$ est un isomorphisme, alors $\{T' \to T\}$
est un recouvrement fpqc de $T$.
\item Si $\{T_i \to T\}_{i\in I}$ est un recouvrement fpqc et si, pour tout
$i$, on a un recouvrement fpqc $\{T_{ij} \to T_i\}_{j\in J_i}$, alors
$\{T_{ij} \to T\}_{i \in I, j\in J_i}$ est un recouvrement fpqc.
\item Si $\{T_i \to T\}_{i\in I}$ est un recouvrement fpqc
et si $T' \to T$ est un morphisme de schémas, alors
$\{T' \times_T T_i \to T'\}_{i\in I}$ est un recouvrement fpqc.
\end{enumerate}
\end{lemma}

\begin{proof}
L'assertion (1) est immédiate. Rappelons que le composé de morphismes plats
est plat et que le changement de base d'un morphisme plat est plat
(Morphismes, lemmes \ref{morphisms-lemma-base-change-flat} et
\ref{morphisms-lemma-composition-flat}).
On peut donc appliquer le lemme \ref{topologies-lemma-recognize-fpqc-covering}
dans chaque cas pour vérifier que nos familles de morphismes sont des recouvrements fpqc.

\medskip\noindent
Preuve de (2). Supposons que $\{T_i \to T\}_{i\in I}$ soit un recouvrement fpqc et que, pour tout
$i$, on ait un recouvrement fpqc $\{f_{ij} : T_{ij} \to T_i\}_{j\in J_i}$.
Soit $U \subset T$ un ouvert affine. On peut trouver
des ouverts quasi-compacts $U_i \subset T_i$ pour $i \in I$, presque tous vides,
tels que $U = \bigcup f_i(U_i)$. Pour tout $i$, on peut alors choisir
des ouverts quasi-compacts $U_{ij} \subset T_{ij}$ pour $j \in J_i$, presque tous vides,
tels que $U_i = \bigcup_j f_{ij}(U_{ij})$. Ainsi,
$\{T_{ij} \to T\}$ est un recouvrement fpqc.

\medskip\noindent
Preuve de (3). Supposons que $\{T_i \to T\}_{i\in I}$ soit un recouvrement fpqc
et que $T' \to T$ soit un morphisme de schémas. Soit $U' \subset T'$ un ouvert affine
dont l'image est contenue dans l'ouvert affine $U \subset T$. Choisissons
des ouverts quasi-compacts $U_i \subset T_i$, presque tous vides,
tels que $U = \bigcup f_i(U_i)$. Alors $U' \times_U U_i$ est
un ouvert quasi-compact de $T' \times_T T_i$ et
$U' = \bigcup \text{pr}_1(U' \times_U U_i)$. Puisque $T'$
peut être recouvert par de tels ouverts affines $U' \subset T'$, on voit
que $\{T' \times_T T_i \to T'\}_{i\in I}$ est un recouvrement fpqc d'après le
lemme \ref{topologies-lemma-recognize-fpqc-covering}.
\end{proof}

\begin{lemma}
\label{topologies-lemma-fpqc-affine}
Soit $T$ un schéma affine.
Soit $\{T_i \to T\}_{i \in I}$ un recouvrement fpqc de $T$.
Il existe alors un recouvrement fpqc
$\{U_j \to T\}_{j = 1, \ldots, n}$ qui est un raffinement
de $\{T_i \to T\}_{i \in I}$ tel que chaque $U_j$ soit un schéma
affine. De plus, on peut choisir chaque $U_j$ ouvert affine
dans l'un des $T_i$.
\end{lemma}

\begin{proof}
Cela résulte directement de la définition.
\end{proof}

\begin{definition}
\label{topologies-definition-standard-fpqc}
Soit $T$ un schéma affine. Un {\it recouvrement fpqc standard}
de $T$ est une famille $\{f_j : U_j \to T\}_{j = 1, \ldots, n}$
où chaque $U_j$ est affine, plat sur $T$, et $T = \bigcup f_j(U_j)$.
\end{definition}

\noindent
Puisque nous n'introduisons pas le site affine, il faut montrer directement
que la famille de tous les recouvrements fpqc standards satisfait les
axiomes.

\begin{lemma}
\label{topologies-lemma-fpqc-affine-axioms}
Soit $T$ un schéma affine.
\begin{enumerate}
\item Si $T' \to T$ est un isomorphisme, alors $\{T' \to T\}$
est un recouvrement fpqc standard de $T$.
\item Si $\{T_i \to T\}_{i\in I}$ est un recouvrement fpqc standard et si, pour tout
$i$, on a un recouvrement fpqc standard $\{T_{ij} \to T_i\}_{j\in J_i}$, alors
$\{T_{ij} \to T\}_{i \in I, j\in J_i}$ est un recouvrement fpqc standard.
\item Si $\{T_i \to T\}_{i\in I}$ est un recouvrement fpqc standard
et si $T' \to T$ est un morphisme de schémas affines, alors
$\{T' \times_T T_i \to T'\}_{i\in I}$ est un recouvrement fpqc standard.
\end{enumerate}
\end{lemma}

\begin{proof}
Cela résulte formellement du fait que les composés et les changements de base
de morphismes plats sont plats
(Morphismes, lemmes \ref{morphisms-lemma-base-change-flat} et
\ref{morphisms-lemma-composition-flat})
et que les produits fibrés de schémas affines sont affines
(Schémas, lemme \ref{schemes-lemma-fibre-product-affines}).
\end{proof}

\begin{lemma}
\label{topologies-lemma-fpqc-covering-affines-mapping-in}
Soit $T$ un schéma. Soit $\{f_i : T_i \to T\}_{i \in I}$ une famille de
morphismes de schémas de but $T$. Supposons que
\begin{enumerate}
\item chaque $f_i$ soit plat, et
\item pour tout schéma affine
$Z$ et tout morphisme $h : Z \to T$, il existe un recouvrement fpqc standard
$\{Z_j \to Z\}_{j = 1, \ldots, n}$ qui raffine la famille
$\{T_i \times_T Z \to Z\}_{i \in I}$.
\end{enumerate}
Alors $\{f_i : T_i \to T\}_{i \in I}$ est un recouvrement fpqc de $T$.
\end{lemma}

\begin{proof}
Soit $T = \bigcup U_\alpha$ un recouvrement ouvert affine.
Pour tout $\alpha$, la famille obtenue par image inverse $\{T_i \times_T U_\alpha \to U_\alpha\}$
peut être raffinée par un recouvrement fpqc standard ; c'est donc un
recouvrement fpqc d'après le lemme
\ref{topologies-lemma-family-flat-dominated-covering}.
Comme $\{U_\alpha \to T\}$ est un recouvrement fpqc, on en conclut que
$\{T_i \to T\}$ est un recouvrement fpqc d'après le
lemme \ref{topologies-lemma-family-flat-fpqc-local-covering}.
\end{proof}

\begin{definition}
\label{topologies-definition-sheaf-property-fpqc}
Soit $F$ un foncteur contravariant de la catégorie
des schémas vers celle des ensembles.
\begin{enumerate}
\item Soit $\{U_i \to T\}_{i \in I}$ une famille de morphismes
de schémas de but fixé.
On dit que $F$ {\it satisfait la propriété de faisceau pour la famille donnée}
si, pour toute famille d'éléments $\xi_i \in F(U_i)$ telle que
$\xi_i|_{U_i \times_T U_j} = \xi_j|_{U_i \times_T U_j}$
il existe un unique élément
$\xi \in F(T)$ tel que $\xi_i = \xi|_{U_i}$ dans $F(U_i)$.
\item On dit que $F$ {\it satisfait la propriété de faisceau pour la
topologie fpqc} s'il satisfait la propriété de faisceau pour tout
recouvrement fpqc.
\end{enumerate}
\end{definition}

\noindent
Nous évitons autant que possible de dire « $F$ est un faisceau » dans cette
situation, puisque nous ne définissons pas de catégorie de faisceaux fpqc,
comme nous l'avons expliqué ci-dessus.

\begin{lemma}
\label{topologies-lemma-sheaf-property-fpqc}
Soit $F$ un foncteur contravariant de la catégorie
des schémas vers celle des ensembles. Alors $F$ satisfait
la propriété de faisceau pour la topologie fpqc si et seulement
s'il satisfait
\begin{enumerate}
\item la propriété de faisceau pour tout recouvrement de Zariski, et
\item la propriété de faisceau pour tout recouvrement fpqc standard.
\end{enumerate}
De plus, en présence de (1), la propriété (2) équivaut à la
propriété
\begin{enumerate}
\item[(2')] de faisceau pour $\{V \to U\}$,
où $V$, $U$ sont affines et $V \to U$ est fidèlement plat.
\end{enumerate}
\end{lemma}

\begin{proof}
Supposons (1) et (2) vérifiées.
Soit $\{f_i : T_i \to T\}_{i \in I}$ un recouvrement fpqc. Soit $s_i \in F(T_i)$
une famille d'éléments telle que $s_i$ et $s_j$ aient la même image
dans $F(T_i \times_T T_j)$. Soit $W \subset T$ le plus grand ouvert
tel qu'il existe un unique $s \in F(W)$ satisfaisant
$s|_{f_i^{-1}(W)} = s_i|_{f_i^{-1}(W)}$ pour tout $i$.
Un tel plus grand ouvert existe puisque $F$ satisfait la
propriété de faisceau pour les recouvrements de Zariski ; en fait, $W$ est la
réunion de tous les ouverts qui ont cette propriété. Soit $t \in T$.
Montrons que $t \in W$. Pour cela, choisissons un ouvert affine
$t \in U \subset T$ et montrons qu'il existe un unique
$s \in F(U)$ tel que
$s|_{f_i^{-1}(U)} = s_i|_{f_i^{-1}(U)}$ pour tout $i$.

\medskip\noindent
D'après le lemme \ref{topologies-lemma-fpqc-affine}, on peut trouver un recouvrement fpqc standard
$\{U_j \to U\}_{j = 1, \ldots, n}$ qui raffine $\{U \times_T T_i \to U\}$,
disons par des morphismes $h_j : U_j \to T_{i_j}$. D'après (2), on obtient un unique élément
$s \in F(U)$ tel que $s|_{U_j} = F(h_j)(s_{i_j})$. Remarquons que, pour tout
schéma $V \to U$ sur $U$, il existe une unique section $s_V \in F(V)$
dont la restriction vaut $F(h_j \circ \text{pr}_2)(s_{i_j})$ sur
$V \times_U U_j$ pour $j = 1, \ldots, n$. En effet, c'est vrai si $V$
est affine d'après (2), puisque $\{V \times_U U_j \to V\}$ est un recouvrement fpqc standard,
et le cas général résulte de (1) et du cas affine en choisissant un
recouvrement ouvert affine de $V$. En particulier, $s_V = s|_V$.
En prenant maintenant $V = U \times_T T_i$ et en utilisant
$s_{i_j}|_{T_{i_j} \times_T T_i} = s_i|_{T_{i_j} \times_T T_i}$,
on en conclut que $s|_{U \times_T T_i} = s_V = s_i|_{U \times_T T_i}$,
ce qu'il fallait démontrer.

\medskip\noindent
Preuve de l'équivalence de (2) et (2') en présence de (1).
Si $\{T_i \to T\}$ est un recouvrement fpqc standard, alors
$\coprod T_i \to T$ est un morphisme fidèlement plat de schémas affines.
En présence de (1), on a $F(\coprod T_i) = \prod F(T_i)$
et, de même,
$F((\coprod T_i) \times_T (\coprod T_i)) = \prod F(T_i \times_T T_{i'})$.
Ainsi, la condition de faisceau pour $\{T_i \to T\}$ et pour $\{\coprod T_i \to T\}$
est la même.
\end{proof}

\noindent
Le lemme suivant sert seulement à signaler que les difficultés ensemblistes
se présentent bel et bien et que la plupart des lecteurs devraient les ignorer.

\begin{lemma}
\label{topologies-lemma-no-set-of-fpqc-covers-is-initial}
Soit $R$ un anneau non nul. Il n'existe aucun ensemble $A$ de
recouvrements fpqc de $\Spec(R)$ tel que tout recouvrement fpqc puisse
être raffiné par un élément de $A$.
\end{lemma}

\begin{proof}
Expliquons d'abord le cas où $R = k$ est un corps. Pour tout ensemble $I$, considérons
l'extension transcendante pure
$k_I = k(\{t_i\}_{i \in I})/k$. Comme $k \to k_I$ est fidèlement plat,
on voit que $\{\Spec(k_I) \to \Spec(k)\}$ est un recouvrement fpqc.
Soit $A$ un ensemble et, pour tout $\alpha \in A$, soit
$\mathcal{U}_\alpha = \{S_{\alpha, j} \to \Spec(k)\}_{j \in J_\alpha}$ un
recouvrement fpqc. Si $\mathcal{U}_\alpha$ raffine $\{\Spec(k_I) \to \Spec(k)\}$,
alors les morphismes $S_{\alpha, j} \to \Spec(k)$ se factorisent par
$\Spec(k_I)$. Comme $\mathcal{U}_\alpha$ est un recouvrement,
au moins un certain $S_{\alpha, j}$ est non vide. Choisissons un
point $s \in S_{\alpha, j}$. Puisque l'on a la factorisation
$S_{\alpha, j} \to \Spec(k_I) \to \Spec(k)$
on obtient un homomorphisme de corps $k_I \to \kappa(s)$.
En particulier, on voit que la cardinalité de $\kappa(s)$
est au moins celle de $I$. Ainsi, si l'on choisit pour $I$ un ensemble
de cardinalité supérieure à celles des corps résiduels
de tous les schémas $S_{\alpha, j}$, une telle factorisation
n'existe pas et le lemme est démontré pour $R = k$.

\medskip\noindent
Cas général. Comme $R$ est non nul, il possède un idéal premier maximal
$\mathfrak m$ de corps résiduel $\kappa$. Soit $I$ un ensemble et
considérons $R_I = S_I^{-1} R[\{t_i\}_{i \in I}]$,
où $S_I \subset R[\{t_i\}_{i \in I}]$ est la partie multiplicative
formée des $f \in R[\{t_i\}_{i \in I}]$ tels que l'image de $f$ soit
un élément non nul de $R/\mathfrak p[\{t_i\}_{i \in I}]$ pour
tout idéal premier $\mathfrak p$ de $R$. Alors $R_I$ est une $R$-algèbre
fidèlement plate et $\{\Spec(R_I) \to \Spec(R)\}$ est un
recouvrement fpqc. Nous laissons au lecteur le soin de montrer que
$R_I \otimes_R \kappa \cong \kappa(\{t_i\}_{i \in I}) = \kappa_I$
avec les notations ci-dessus (indication : utiliser la surjectivité de $R \to \kappa$
et le fait que tout $f \in R[\{t_i\}_{i \in I}]$ dont l'un des monômes apparaît
avec le coefficient $1$ appartient à $S_I$). Soit $A$ un ensemble et,
pour tout $\alpha \in A$, soit
$\mathcal{U}_\alpha = \{S_{\alpha, j} \to \Spec(R)\}_{j \in J_\alpha}$ un
recouvrement fpqc. Si $\mathcal{U}_\alpha$ raffine $\{\Spec(R_I) \to \Spec(R)\}$,
alors, par changement de base, on en conclut que
$\{S_{\alpha, j} \times_{\Spec(R)} \Spec(\kappa) \to \Spec(\kappa)\}$
raffine $\{\Spec(\kappa_I) \to \Spec(\kappa)\}$.
Par le résultat du paragraphe précédent, il existe donc un $I$
pour lequel ce n'est pas le cas, ce qui démontre le lemme.
\end{proof}












\section{La topologie V}
\label{topologies-section-V}

\noindent
La topologie V est plus fine que toutes les autres topologies de ce chapitre.
Grosso modo, elle est engendrée par les recouvrements de Zariski et
par les morphismes quasi-compacts satisfaisant une propriété de relèvement
des spécialisations (lemme \ref{topologies-lemma-refine-qcqs-V}).
Cependant, le procédé que nous emploierons pour définir les recouvrements V est un peu différent.
Nous définirons d'abord les recouvrements V standards des schémas affines, puis
nous les utiliserons pour définir les recouvrements V en général. Remarque typographique :
dans la littérature, on emploie parfois « $v$-recouvrement » au lieu de « recouvrement V ».

\begin{definition}
\label{topologies-definition-standard-V-covering}
Soit $T$ un schéma affine. Un {\it recouvrement V standard} est une famille finie
$\{T_j \to T\}_{j = 1, \ldots, m}$, où les $T_j$ sont affines,
telle que, pour tout morphisme $g : \Spec(V) \to T$, où $V$
est un anneau de valuation, il existe une extension $V \subset W$ d'anneaux de valuation
(Compléments d'algèbre, définition
\ref{more-algebra-definition-extension-valuation-rings}),
un indice $1 \leq j \leq m$ et un diagramme commutatif
$$
\xymatrix{
\Spec(W) \ar[r] \ar[d] & T_j \ar[d] \\
\Spec(V) \ar[r]^g & T
}
$$
\end{definition}

\noindent
Nous démontrons d'abord quelques lemmes élémentaires sur cette notion.

\begin{lemma}
\label{topologies-lemma-standard-fpqc-standard-V}
Tout recouvrement fpqc standard est un recouvrement V standard.
\end{lemma}

\begin{proof}
Soit $\{X_i \to X\}_{i = 1, \ldots, n}$ un recouvrement fpqc standard
(définition \ref{topologies-definition-standard-fpqc}). Soit $g : \Spec(V) \to X$
un morphisme, où $V$ est un anneau de valuation. Soit $x \in X$
l'image du point fermé de $\Spec(V)$. Choisissons un $i$ et un point
$x_i \in X_i$ d'image $x$. Alors $\Spec(V) \times_X X_i$
possède un point $x'_i$ dont l'image est le point fermé de $\Spec(V)$.
Comme $\Spec(V) \times_X X_i \to \Spec(V)$ est plat,
on peut trouver une spécialisation $x''_i \leadsto x'_i$ de
points de $\Spec(V) \times_X X_i$, où $x''_i$ s'envoie sur
le point générique de $\Spec(V)$ ; voir
Morphismes, lemme \ref{morphisms-lemma-generalizations-lift-flat}. D'après
Schémas, lemme \ref{schemes-lemma-points-specialize},
on peut choisir un anneau de valuation $W$ et un morphisme
$h : \Spec(W) \to \Spec(V) \times_X X_i$ tels que $h$
envoie le point générique de $\Spec(W)$ sur $x''_i$ et le
point fermé de $\Spec(W)$ sur $x'_i$. On obtient un
diagramme commutatif
$$
\xymatrix{
\Spec(W) \ar[r] \ar[d] & X_i \ar[d] \\
\Spec(V) \ar[r] & X
}
$$
où $V \to W$ est une extension d'anneaux de valuation.
Cela démontre le lemme.
\end{proof}

\begin{lemma}
\label{topologies-lemma-standard-ph-standard-V}
Tout recouvrement ph standard est un recouvrement V standard.
\end{lemma}

\begin{proof}
Soit $T$ un schéma affine. Soit $f : U \to T$ un morphisme propre
surjectif. Soit $U = \bigcup_{j = 1, \ldots, m} U_j$ un recouvrement ouvert
affine fini. Il faut montrer que $\{U_j \to T\}$ est un recouvrement V
standard ; voir la définition
\ref{topologies-definition-standard-ph-covering}.
Soit $g : \Spec(V) \to T$
un morphisme, où $V$ est un anneau de valuation de corps des fractions $K$.
Comme $U \to T$ est surjectif, on peut choisir une extension de corps
$L/K$ et un diagramme commutatif
$$
\xymatrix{
\Spec(L) \ar[rr] \ar[d] & & U \ar[d] \\
\Spec(K) \ar[r] & \Spec(V) \ar[r]^g & T
}
$$
D'après Algèbre, lemme \ref{algebra-lemma-dominate}, on peut choisir un anneau
de valuation $W \subset L$ dominant $V$. Le critère valuatif de
propreté (Morphismes, lemme \ref{morphisms-lemma-characterize-proper})
fournit alors le morphisme $h$ dans le diagramme
commutatif
$$
\xymatrix{
\Spec(L) \ar[r] \ar[d] & \Spec(W) \ar[r]_h \ar[d] & U \ar[d] \\
\Spec(K) \ar[r] & \Spec(V) \ar[r]^g & X
}
$$
Comme $\Spec(W)$ possède un unique point fermé, on voit que $\Im(h)$
est contenue dans $U_j$ pour un certain $j$. Ainsi, $h : \Spec(W) \to U_j$
est le relèvement cherché, et $\{U_j \to T\}$ est donc un recouvrement V standard.
\end{proof}

\begin{lemma}
\label{topologies-lemma-base-change-standard-V}
Soit $\{T_j \to T\}_{j = 1, \ldots, m}$ un recouvrement V standard.
Soit $T' \to T$ un morphisme de schémas affines.
Alors $\{T_j \times_T T' \to T'\}_{j = 1, \ldots, m}$ est un
recouvrement V standard.
\end{lemma}

\begin{proof}
Soit $\Spec(V) \to T'$ un morphisme, où $V$ est un anneau de valuation.
Par hypothèse, on peut trouver une extension d'anneaux de valuation $V \subset W$,
un $i$ et un diagramme commutatif
$$
\xymatrix{
\Spec(W) \ar[r] \ar[d] & T_i \ar[d] \\
\Spec(V) \ar[r] & T
}
$$
La propriété universelle des produits fibrés
donne le morphisme $\Spec(W) \to T' \times_T T_i$
voulu.
\end{proof}

\begin{lemma}
\label{topologies-lemma-composition-standard-V}
Soit $T$ un schéma affine. Soit $\{T_j \to T\}_{j = 1, \ldots, m}$
un recouvrement V standard. Soit $\{T_{ji} \to T_j\}_{i = 1, \ldots n_j}$
un recouvrement V standard. Alors $\{T_{ji} \to T\}_{i, j}$
est un recouvrement V standard.
\end{lemma}

\begin{proof}
Cela résulte formellement de l'observation suivante : si
$V \subset W$ et $W \subset \Omega$ sont des extensions
d'anneaux de valuation, alors $V \subset \Omega$ est une extension
d'anneaux de valuation.
\end{proof}

\begin{lemma}
\label{topologies-lemma-refine-standard-V}
Soit $T$ un schéma affine. Soit $\{T_j \to T\}_{j = 1, \ldots, m}$
une famille de morphismes telle que $T_j$ soit affine pour tout $j$.
Les conditions suivantes sont équivalentes :
\begin{enumerate}
\item $\{T_j \to T\}_{j = 1, \ldots, m}$ est un recouvrement V standard ;
\item il existe un recouvrement V standard qui raffine
$\{T_j \to T\}_{j = 1, \ldots, m}$ ;
\item $\{\coprod_{j = 1, \ldots, m} T_j \to T\}$ est un recouvrement
V standard.
\end{enumerate}
\end{lemma}

\begin{proof}
Omis. Indications : cela résulte presque immédiatement de la définition.
Le seul point qui mérite quelque attention est qu'un
morphisme du spectre d'un anneau local vers
$\coprod_{j = 1, \ldots, m} T_j$ se factorise nécessairement par l'un des $T_j$.
\end{proof}

\begin{definition}
\label{topologies-definition-V-covering}
Soit $T$ un schéma. Un {\it recouvrement V de $T$} est une famille
de morphismes de schémas $\{T_i \to T\}_{i \in I}$ telle que, pour tout
ouvert affine $U \subset T$, il existe un recouvrement V standard
$\{U_j \to U\}_{j = 1, \ldots, m}$ qui raffine la famille
$\{T_i \times_T U \to U\}_{i \in I}$.
\end{definition}

\noindent
La topologie V présente les mêmes problèmes ensemblistes que
la topologie fpqc. Nous nous abstiendrons donc de définir des sites V
et nous ne considérerons pas de cohomologie pour la topologie V.
En revanche, étant donné un foncteur
$F : \Sch^{opp} \to \textit{Ensembles}$
il est légitime de demander si $F$ satisfait la propriété de faisceau
pour la topologie V ; voir ci-dessous.
On peut en outre s'interroger sur la descente d'objets
pour la topologie V, etc.

\begin{lemma}
\label{topologies-lemma-refine-by-V}
Soit $T$ un schéma. Soit $\{f_i : T_i \to T\}_{i \in I}$ une famille
de morphismes. Les conditions suivantes sont équivalentes :
\begin{enumerate}
\item $\{T_i \to T\}_{i \in I}$ est un recouvrement V ;
\item il existe un recouvrement V qui raffine $\{T_i \to T\}_{i \in I}$ ;
\item $\{\coprod_{i \in I} T_i \to T\}$ est un recouvrement V.
\end{enumerate}
\end{lemma}

\begin{proof}
Omis. Indication : comparer avec la démonstration du
lemme \ref{topologies-lemma-refine-by-ph}.
\end{proof}

\begin{lemma}
\label{topologies-lemma-V}
Soit $T$ un schéma.
\begin{enumerate}
\item Si $T' \to T$ est un isomorphisme, alors $\{T' \to T\}$
est un recouvrement V de $T$.
\item Si $\{T_i \to T\}_{i\in I}$ est un recouvrement V et si, pour chaque
$i$, on a un recouvrement V $\{T_{ij} \to T_i\}_{j\in J_i}$, alors
$\{T_{ij} \to T\}_{i \in I, j\in J_i}$ est un recouvrement V.
\item Si $\{T_i \to T\}_{i\in I}$ est un recouvrement V
et si $T' \to T$ est un morphisme de schémas, alors
$\{T' \times_T T_i \to T'\}_{i\in I}$ est un recouvrement V.
\end{enumerate}
\end{lemma}

\begin{proof}
L'assertion (1) est claire.

\medskip\noindent
Démonstration de (3).
Soit $U' \subset T'$ un sous-schéma ouvert affine.
Comme $U'$ est quasi-compact, on peut trouver un recouvrement ouvert affine
fini $U' = U'_1 \cup \ldots \cup U'$ tel que
$U'_j \to T$ se factorise par un ouvert affine $U_j \subset T$.
Choisissons un recouvrement V standard $\{U_{jl} \to U_j\}_{l = 1, \ldots, n_j}$
qui raffine $\{T_i \times_T U_j \to U_j\}$.
D'après le lemme \ref{topologies-lemma-base-change-standard-V},
le changement de base $\{U_{jl} \times_{U_j} U'_j \to U'_j\}$
est un recouvrement V standard. Notons que $\{U'_j \to U'\}$
est un recouvrement V standard
(par exemple d'après le lemme \ref{topologies-lemma-standard-fpqc-standard-V}).
D'après le lemme \ref{topologies-lemma-composition-standard-V}, 
la famille $\{U_{jl} \times_{U_j} U'_j \to U'\}$
est un recouvrement V standard. Comme
$\{U_{jl} \times_{U_j} U'_j \to U'\}$ raffine $\{T_i \times_T U' \to U'\}$,
on conclut.

\medskip\noindent
Démonstration de (2).
Soit $U \subset T$ un ouvert affine. Choisissons d'abord un recouvrement V standard
$\{U_k \to U\}_{k = 1, \ldots, m}$ qui raffine
$\{T_i \times_T U \to U\}$.
Supposons que le raffinement soit donné par des morphismes $U_k \to T_{i_k}$ au-dessus de $T$.
Alors
$$
\{T_{i_kj} \times_{T_{i_k}} U_k \to U_k\}_{j \in J_{i_k}}
$$
est un recouvrement V d'après (3). Comme $U_k$ est affine,
on peut trouver un recouvrement V standard
$\{U_{ka} \to U_k\}_{a = 1, \ldots, b_k}$ qui raffine cette famille.
On applique alors le lemme \ref{topologies-lemma-composition-standard-V}
pour voir que $\{U_{ka} \to U\}$ est un recouvrement V standard qui
raffine $\{T_{ij} \times_T U \to U\}$. Cela achève la démonstration.
\end{proof}

\begin{lemma}
\label{topologies-lemma-zariski-etale-smooth-syntomic-fppf-fpqc-ph-V}
Tout recouvrement fpqc est un recouvrement V. A fortiori,
tout recouvrement fppf, syntomique, lisse, étale ou de Zariski est un recouvrement V.
De même, tout recouvrement ph est un recouvrement V.
\end{lemma}

\begin{proof}
Tout recouvrement fpqc peut, localement sur des ouverts affines, être raffiné par un
recouvrement fpqc standard ; voir le lemme \ref{topologies-lemma-fpqc-affine}.
Tout recouvrement fpqc standard est un recouvrement V standard ; voir le
lemme \ref{topologies-lemma-standard-fpqc-standard-V}. La première assertion résulte donc
de notre définition des recouvrements V au moyen des recouvrements V standards.
La conclusion pour les recouvrements fppf, syntomiques, lisses, étales ou de Zariski
résulte de ce que ceux-ci sont des recouvrements fpqc ; voir le
lemme \ref{topologies-lemma-zariski-etale-smooth-syntomic-fppf-fpqc}.

\medskip\noindent
L'assertion sur les recouvrements ph résulte de la même manière du
lemme \ref{topologies-lemma-standard-ph-standard-V}.
\end{proof}

\begin{definition}
\label{topologies-definition-sheaf-property-V}
Soit $F$ un foncteur contravariant de la catégorie
des schémas dans celle des ensembles. On dit que
$F$ {\it satisfait la propriété de faisceau pour la topologie V}
s'il satisfait la propriété de faisceau pour tout recouvrement V
(voir la définition \ref{topologies-definition-sheaf-property-fpqc}).
\end{definition}

\noindent
Nous évitons d'employer la terminologie « $F$ est un faisceau » dans cette
situation, puisque nous ne définissons pas de catégorie des faisceaux V,
comme nous l'avons expliqué plus haut.

\begin{lemma}
\label{topologies-lemma-sheaf-property-V}
Soit $F$ un foncteur contravariant de la catégorie
des schémas dans celle des ensembles. Alors $F$ satisfait
la propriété de faisceau pour la topologie V si et seulement
s'il satisfait aux deux conditions suivantes :
\begin{enumerate}
\item la propriété de faisceau pour tout recouvrement de Zariski ;
\item la propriété de faisceau pour tout recouvrement V standard.
\end{enumerate}
En outre, sous l'hypothèse (1), la propriété (2) équivaut à la
propriété suivante :
\begin{enumerate}
\item[(2')] la propriété de faisceau pour un recouvrement V standard
de la forme $\{V \to U\}$, c'est-à-dire constitué d'une seule flèche.
\end{enumerate}
\end{lemma}

\begin{proof}
Supposons (1) et (2) satisfaites.
Soit $\{f_i : T_i \to T\}_{i \in I}$ un recouvrement V. Soit $s_i \in F(T_i)$
une famille d'éléments telle que $s_i$ et $s_j$ s'envoient sur le même élément
de $F(T_i \times_T T_j)$. Soit $W \subset T$ le plus grand ouvert
tel qu'il existe un unique $s \in F(W)$ vérifiant
$s|_{f_i^{-1}(W)} = s_i|_{f_i^{-1}(W)}$ pour tout $i$.
Un tel plus grand ouvert existe parce que $F$ satisfait la
propriété de faisceau pour les recouvrements de Zariski ; en fait, $W$ est la
réunion de tous les ouverts possédant cette propriété. Soit $t \in T$.
Nous allons montrer que $t \in W$. Pour cela, choisissons un ouvert affine
$t \in U \subset T$ et montrons qu'il existe un unique
$s \in F(U)$ tel que
$s|_{f_i^{-1}(U)} = s_i|_{f_i^{-1}(U)}$ pour tout $i$.

\medskip\noindent
On peut trouver un recouvrement V standard
$\{U_j \to U\}_{j = 1, \ldots, n}$ qui raffine $\{U \times_T T_i \to U\}$,
disons au moyen de morphismes $h_j : U_j \to T_{i_j}$. D'après (2), on obtient un unique élément
$s \in F(U)$ tel que $s|_{U_j} = F(h_j)(s_{i_j})$. Notons que, pour tout
schéma $V \to U$ au-dessus de $U$, il existe une unique section $s_V \in F(V)$
qui se restreint en $F(h_j \circ \text{pr}_2)(s_{i_j})$ sur
$V \times_U U_j$ pour $j = 1, \ldots, n$. En effet, cela est vrai si $V$
est affine d'après (2), puisque $\{V \times_U U_j \to V\}$ est un recouvrement V standard
(lemme \ref{topologies-lemma-base-change-standard-V}) ;
dans le cas général, cela résulte de (1) et du cas affine en choisissant un
recouvrement ouvert affine de $V$. En particulier, $s_V = s|_V$.
En prenant maintenant $V = U \times_T T_i$ et en utilisant l'égalité
$s_{i_j}|_{T_{i_j} \times_T T_i} = s_i|_{T_{i_j} \times_T T_i}$
on conclut que $s|_{U \times_T T_i} = s_V = s_i|_{U \times_T T_i}$,
ce qu'il fallait démontrer.

\medskip\noindent
Démonstration de l'équivalence de (2) et (2') sous l'hypothèse (1).
Supposons que $\{T_i \to T\}_{i = 1, \ldots, n}$ soit un recouvrement V standard. Alors
$\coprod_{i = 1, \ldots, n} T_i \to T$ est un morphisme de schémas affines
qui est manifestement aussi un recouvrement V standard.
Sous l'hypothèse (1), on a $F(\coprod T_i) = \prod F(T_i)$
et, de même,
$F((\coprod T_i) \times_T (\coprod T_i)) = \prod F(T_i \times_T T_{i'})$.
Ainsi, la condition de faisceau pour $\{T_i \to T\}$ et pour $\{\coprod T_i \to T\}$
est la même.
\end{proof}

\noindent
Le lemme suivant montre que la propriété d'être un recouvrement V est
liée à la possibilité de relever les spécialisations.

\begin{lemma}
\label{topologies-lemma-refine-qcqs-V}
Soit $X \to Y$ un morphisme quasi-compact de schémas.
Les conditions suivantes sont équivalentes :
\begin{enumerate}
\item $\{X \to Y\}$ est un recouvrement V ;
\item pour tout anneau de valuation $V$ et tout morphisme $g : \Spec(V) \to Y$,
il existe une extension d'anneaux de valuation $V \subset W$
et un diagramme commutatif
$$
\xymatrix{
\Spec(W) \ar[r] \ar[d] & X \ar[d] \\
\Spec(V) \ar[r] & Y
}
$$
\item pour tout morphisme $Z \to Y$ et toute spécialisation $z' \leadsto z$
de points de $Z$, il existe une spécialisation $w' \leadsto w$
de points de $Z \times_Y X$ qui s'envoie sur $z' \leadsto z$.
\end{enumerate}
\end{lemma}

\begin{proof}
Supposons (1) et soit $g : \Spec(V) \to Y$ comme dans (2).
Comme $V$ est un anneau local, il existe un ouvert affine $U \subset Y$
tel que $g$ se factorise par $U$. D'après la définition \ref{topologies-definition-V-covering},
on peut trouver un recouvrement V standard $\{U_j \to U\}$ qui raffine
$\{X \times_Y U \to U\}$. D'après la définition \ref{topologies-definition-standard-V-covering},
on peut trouver un $j$, une extension d'anneaux de valuation $V \subset W$
et un diagramme commutatif
$$
\xymatrix{
\Spec(W) \ar[r] \ar[d] & U_j \ar[d] \ar@{..>}[r] & X \ar[ld] \\
\Spec(V) \ar[r] & Y
}
$$
La propriété de raffinement du recouvrement fournit la flèche pointillée
qui rend le diagramme commutatif, et l'on voit que (2) est satisfaite.

\medskip\noindent
Supposons (2) et soient $Z \to Y$ et $z' \leadsto z$ comme dans (3).
D'après Schémas, lemme \ref{schemes-lemma-points-specialize},
on peut trouver un anneau de valuation $V$ et un morphisme $\Spec(V) \to Z$
tels que le point fermé de $\Spec(V)$ s'envoie sur $z$ et que le
point générique de $\Spec(V)$ s'envoie sur $z'$. D'après (2), on peut trouver une
extension d'anneaux de valuation $V \subset W$ et un
diagramme commutatif
$$
\xymatrix{
\Spec(W) \ar[rr] \ar[d] & & X \ar[d] \\
\Spec(V) \ar[r] & Z \ar[r] & Y
}
$$
Les points générique et fermé de $\Spec(W)$ s'envoient sur des points
$w' \leadsto w$ de $Z \times_Y X$ par le
morphisme induit $\Spec(W) \to Z \times_Y X$.
Cela montre que (3) est satisfaite.

\medskip\noindent
Supposons (3) satisfaite et soit $U \subset Y$ un ouvert affine.
Choisissons un recouvrement ouvert affine fini
$U \times_Y X = \bigcup_{j = 1, \ldots, m} U_j$.
C'est possible puisque $X \to Y$ est quasi-compact.
Nous affirmons que $\{U_j \to U\}$ est un recouvrement V standard.
Cette affirmation entraîne (1) et achève la démonstration du lemme.
Pour la démontrer, soit $V$ un anneau de valuation
et soit $g : \Spec(V) \to U$ un morphisme.
D'après (3), on trouve une spécialisation $w' \leadsto w$ de
points de
$$
T = \Spec(V) \times_X Y = \Spec(V) \times_U (U \times_X Y)
$$
telle que $w'$ s'envoie sur le point générique de $\Spec(V)$
et que $w$ s'envoie sur le point fermé de $\Spec(V)$.
D'après Schémas, lemme \ref{schemes-lemma-points-specialize},
on peut trouver un anneau de valuation $W$ et un morphisme
$\Spec(W) \to T$ tels que le point générique de $\Spec(W)$
s'envoie sur $w'$ et que le point fermé de $\Spec(W)$ s'envoie sur $w$.
La composée $\Spec(W) \to T \to \Spec(V)$ correspond
à une inclusion $V \subset W$ qui fait de $W$ une
extension de l'anneau de valuation $V$.
Comme $T = \bigcup \Spec(V) \times_U U_j$ est un recouvrement
ouvert, on voit que $\Spec(W) \to T$ se factorise par
$\Spec(V) \times_U U_j$ pour un certain $j$. On obtient donc
un diagramme commutatif
$$
\xymatrix{
\Spec(W) \ar[d] \ar[r] & U_j \ar[d] \\
\Spec(V) \ar[r] & U
}
$$
et la démonstration de l'affirmation est achevée.
\end{proof}

\noindent
Tout recouvrement V fournit une famille de morphismes universellement submersive.
La réciproque de ce lemme est fausse ; voir
Exemples, section \ref{examples-section-universally-submersive-not-V}.

\begin{lemma}
\label{topologies-lemma-V-covering-universally-submersive}
Soit $\{f_i : X_i \to X\}_{i \in I}$ un recouvrement V.
Alors
$$
\coprod\nolimits_{i \in I} f_i :
\coprod\nolimits_{i \in I} X_i
\longrightarrow
X
$$
est un morphisme de schémas universellement submersif (Morphismes, définition
\ref{morphisms-definition-submersive}).
\end{lemma}

\begin{proof}
Nous utiliserons sans autre mention que le changement de base d'un recouvrement V
est un recouvrement V (lemme \ref{topologies-lemma-V}). Il suffit donc
de montrer que le morphisme est submersif.
Le caractère submersif est manifestement local pour la topologie de Zariski sur la base.
On peut donc supposer $X$ affine. Alors $\{X_i \to X\}$
peut être raffiné par un recouvrement V standard $\{Y_j \to X\}$.
S'il est possible de montrer que $\coprod Y_j \to X$ est submersif,
alors, puisqu'il existe une factorisation $\coprod Y_j \to \coprod X_i \to X$,
on conclut que $\coprod X_i \to X$ est submersif.
Posons $Y = \coprod Y_j$ et considérons le morphisme de schémas affines
$f : Y \to X$.
D'après le lemme \ref{topologies-lemma-refine-qcqs-V}, on sait que toute
spécialisation $x' \leadsto x$ dans $X$ se relève en une spécialisation
$y' \leadsto y$ dans $Y$.
Ainsi, si $T \subset X$ est une partie telle que $f^{-1}(T)$
soit fermée dans $Y$, alors $T \subset X$ est stable par spécialisation.
Comme $f^{-1}(T) \subset Y$, muni de la structure réduite
de sous-schéma fermé induite, est un schéma affine, on conclut que
$T \subset X$ est fermé d'après Algèbre, lemme
\ref{algebra-lemma-image-stable-specialization-closed}.
Ainsi, $f$ est submersif.
\end{proof}





















\section{Changement de topologies}
\label{topologies-section-change-topologies}

\noindent
Soit $f : X \to Y$ un morphisme de schémas au-dessus d'un schéma de base $S$.
On dispose alors des morphismes de sites suivants\footnote{Nous n'avons
pas inclus la comparaison entre la topologie ph et les autres ;
voir à ce sujet
Compléments sur les morphismes, remarque \ref{more-morphisms-remark-change-topologies}.}
(moyennant des choix convenables de sites comme dans la remarque
\ref{topologies-remark-choice-sites} ci-dessous) :
\begin{enumerate}
\item $(\Sch/X)_{fppf} \longrightarrow (\Sch/Y)_{fppf}$,
\item $(\Sch/X)_{fppf} \longrightarrow (\Sch/Y)_{syntomic}$,
\item $(\Sch/X)_{fppf} \longrightarrow (\Sch/Y)_{smooth}$,
\item $(\Sch/X)_{fppf} \longrightarrow
(\Sch/Y)_\etale$,
\item $(\Sch/X)_{fppf} \longrightarrow (\Sch/Y)_{Zar}$,
\item $(\Sch/X)_{syntomic} \longrightarrow (\Sch/Y)_{syntomic}$,
\item $(\Sch/X)_{syntomic} \longrightarrow (\Sch/Y)_{smooth}$,
\item $(\Sch/X)_{syntomic} \longrightarrow
(\Sch/Y)_\etale$,
\item $(\Sch/X)_{syntomic} \longrightarrow (\Sch/Y)_{Zar}$,
\item $(\Sch/X)_{smooth} \longrightarrow (\Sch/Y)_{smooth}$,
\item $(\Sch/X)_{smooth} \longrightarrow
(\Sch/Y)_\etale$,
\item $(\Sch/X)_{smooth} \longrightarrow (\Sch/Y)_{Zar}$,
\item $(\Sch/X)_\etale \longrightarrow
(\Sch/Y)_\etale$,
\item $(\Sch/X)_\etale \longrightarrow (\Sch/Y)_{Zar}$,
\item $(\Sch/X)_{Zar} \longrightarrow (\Sch/Y)_{Zar}$,
\item $(\Sch/X)_{fppf} \longrightarrow Y_\etale$,
\item $(\Sch/X)_{syntomic} \longrightarrow Y_\etale$,
\item $(\Sch/X)_{smooth} \longrightarrow Y_\etale$,
\item $(\Sch/X)_\etale \longrightarrow Y_\etale$,
\item $(\Sch/X)_{fppf} \longrightarrow Y_{Zar}$,
\item $(\Sch/X)_{syntomic} \longrightarrow Y_{Zar}$,
\item $(\Sch/X)_{smooth} \longrightarrow Y_{Zar}$,
\item $(\Sch/X)_\etale \longrightarrow Y_{Zar}$,
\item $(\Sch/X)_{Zar} \longrightarrow Y_{Zar}$,
\item $X_\etale \longrightarrow Y_\etale$,
\item $X_\etale \longrightarrow Y_{Zar}$,
\item $X_{Zar} \longrightarrow Y_{Zar}$,
\end{enumerate}
Dans chaque cas, le foncteur continu sous-jacent
$\Sch/Y \to \Sch/X$, ou
$Y_\tau \to \Sch/X$
est le foncteur $Y'/Y \mapsto X \times_Y Y'/X$. En effet, dans les sections
précédentes, nous avons construit les morphismes
$f_{big} : (\Sch/X)_\tau \to (\Sch/Y)_\tau$
et
$f_{small} : X_\tau \to Y_\tau$
pour $\tau$ comme ci-dessus.
Nous avons aussi construit les morphismes de sites
$\pi_Y : (\Sch/Y)_\tau \to Y_\tau$ pour
$\tau \in \{\etale, Zariski\}$.
D'autre part, il est clair que le foncteur identité
$(\Sch/X)_\tau \to (\Sch/X)_{\tau'}$ définit
un morphisme de sites lorsque $\tau$ est une topologie plus fine que
$\tau'$. En composant ces morphismes, on obtient donc la liste des morphismes possibles
ci-dessus.

\medskip\noindent
La description simple du foncteur sous-jacent montre clairement que, pour des
morphismes de schémas $X \to Y \to Z$, la composée de deux des
morphismes de sites ci-dessus, par exemple
$$
(\Sch/X)_{\tau_0} \longrightarrow
(\Sch/Y)_{\tau_1} \longrightarrow
(\Sch/Z)_{\tau_2}
$$
est le morphisme de sites correspondant associé au morphisme de
schémas $X \to Z$.

\begin{remark}
\label{topologies-remark-choice-sites}
Prenons une catégorie $\Sch_\alpha$ construite comme dans
Ensembles, lemme \ref{sets-lemma-construct-category},
en partant de l'ensemble de schémas $\{X, Y, S\}$. Choisissons un ensemble de
recouvrements $\text{Cov}_{fppf}$ sur $\Sch_\alpha$ comme dans
Ensembles, lemme \ref{sets-lemma-coverings-site},
en partant de la catégorie $\Sch_\alpha$ et de la classe des recouvrements
fppf. Notons $\Sch_{fppf}$ le gros site fppf ainsi
obtenu. Ensuite, pour $\tau \in \{Zariski, \etale, smooth, syntomic\}$,
munissons $\Sch_\tau$ de la même catégorie sous-jacente que
$\Sch_{fppf}$, avec pour recouvrements
$\text{Cov}_\tau \subset \text{Cov}_{fppf}$, simplement le sous-ensemble des
recouvrements $\tau$. On vérifie immédiatement que cela définit
un gros site $\Sch_\tau$.
\end{remark}









\section{Changement de gros sites}
\label{topologies-section-change-alpha}

\noindent
Dans cette section, nous expliquons ce qui se passe lorsqu'on change de gros
site de Zariski, fppf ou étale.

\medskip\noindent
Soient $\tau, \tau' \in \{Zariski, \etale, smooth, syntomic, fppf\}$.
Étant donnés deux gros sites $\Sch_\tau$ et $\Sch'_{\tau'}$,
on dit que
{\it $\Sch_\tau$ est contenu dans $\Sch'_{\tau'}$} si
$\Ob(\Sch_\tau) \subset \Ob(\Sch'_{\tau'})$
et
$\text{Cov}(\Sch_\tau) \subset \text{Cov}(\Sch'_{\tau'})$.
Dans ce cas, $\tau$ est plus fine que $\tau'$ ; par exemple, aucun site
fppf ne peut être contenu dans un site étale.

\begin{lemma}
\label{topologies-lemma-contained-in}
Tout ensemble de gros sites de Zariski est contenu dans un même gros site de Zariski.
Il en va de même, mutatis mutandis, pour les gros sites fppf et les gros sites étales.
\end{lemma}

\begin{proof}
Cela résulte de ce que la réunion d'un ensemble d'ensembles est un ensemble et que les
constructions de Ensembles, lemmes \ref{sets-lemma-construct-category} et
\ref{sets-lemma-coverings-site}
permettent de partir de n'importe quel ensemble initialement donné de schémas
et de recouvrements.
\end{proof}

\begin{lemma}
\label{topologies-lemma-change-alpha}
Soit $\tau \in \{Zariski, \etale, smooth, syntomic, fppf\}$.
Supposons donnés deux gros sites $\Sch_\tau$ et $\Sch'_\tau$.
Supposons que $\Sch_\tau$ soit contenu dans $\Sch'_\tau$.
Le foncteur d'inclusion $\Sch_\tau \to \Sch'_\tau$ satisfait
aux hypothèses de Sites, lemme \ref{sites-lemma-bigger-site}.
Il existe des morphismes de topos
\begin{eqnarray*}
g : \Sh(\Sch_\tau) &
\longrightarrow &
\Sh(\Sch'_\tau) \\
f : \Sh(\Sch'_\tau) &
\longrightarrow &
\Sh(\Sch_\tau)
\end{eqnarray*}
tels que $f \circ g \cong \text{id}$. Pour tout objet $S$
de $\Sch_\tau$, le foncteur d'inclusion
$(\Sch/S)_\tau \to (\Sch'/S)_\tau$ satisfait lui aussi
aux hypothèses de Sites, lemme \ref{sites-lemma-bigger-site}.
On obtient donc de même des morphismes
\begin{eqnarray*}
g : \Sh((\Sch/S)_\tau) &
\longrightarrow &
\Sh((\Sch'/S)_\tau) \\
f : \Sh((\Sch'/S)_\tau) &
\longrightarrow &
\Sh((\Sch/S)_\tau)
\end{eqnarray*}
avec $f \circ g \cong \text{id}$.
\end{lemma}

\begin{proof}
Les hypothèses (b), (c) et (e) de
Sites, lemme \ref{sites-lemma-bigger-site},
sont immédiates pour les foncteurs
$\Sch_\tau \to \Sch'_\tau$ et
$(\Sch/S)_\tau \to (\Sch'/S)_\tau$. La propriété (a) résulte du
lemme \ref{topologies-lemma-zariski-induced},
\ref{topologies-lemma-etale-induced},
\ref{topologies-lemma-smooth-induced},
\ref{topologies-lemma-syntomic-induced}, ou
\ref{topologies-lemma-fppf-induced}.
La propriété (d) est satisfaite parce que
les produits fibrés existent dans les catégories $\Sch_\tau$, $\Sch'_\tau$
et sont compatibles avec les produits fibrés dans la catégorie des schémas.
\end{proof}

\noindent
Discussion :
Le foncteur $g^{-1} = f_*$ est simplement le foncteur de restriction qui associe
à un faisceau $\mathcal{G}$ sur $\Sch'_\tau$ sa restriction
$\mathcal{G}|_{\Sch_\tau}$. Ainsi, ce lemme affirme simplement qu'étant donné
un faisceau d'ensembles $\mathcal{F}$ sur $\Sch_\tau$, il existe un
faisceau canonique $\mathcal{F}'$ sur $\Sch'_\tau$ tel que
$\mathcal{F}|_{\Sch'_\tau} = \mathcal{F}'$. En fait, le faisceau
$\mathcal{F}'$ admet la description suivante : c'est le faisceau associé
au préfaisceau
$$
\Sch'_\tau \longrightarrow \textit{Ensembles}, \quad
V \longmapsto \colim_{V \to U} \mathcal{F}(U)
$$
où $U$ est un objet de $\Sch_\tau$. Cela résulte de l'égalité
$\mathcal{F}' = f^{-1}\mathcal{F} = (u_p\mathcal{F})^\#$ d'après
Sites, lemmes \ref{sites-lemma-when-shriek} et \ref{sites-lemma-bigger-site}.






\section{Prolongement de foncteurs}
\label{topologies-section-extending-functors}

\noindent
Commençons par un exemple simple qui explique notre démarche.
Soit $R$ un anneau. Supposons que $F$ soit un foncteur défini sur la catégorie
$\mathcal{C}$ des $R$-algèbres de la forme
$$
A = R[x_1, \ldots, x_n]/(f_1, \ldots, f_m)
$$
où $n, m \geq 0$ sont des entiers et $f_1, \ldots, f_m \in R[x_1, \ldots, x_n]$
des éléments. Pour toute $R$-algèbre $B$, on peut alors définir
$$
F'(B) = \colim_{A \to B,\ A \in \mathcal{C}} F(A)
$$
Il se trouve que $F'$ est l'unique foncteur sur la catégorie
de toutes les $R$-algèbres qui prolonge $F$ et commute aux limites
inductives filtrantes. Le même procédé fonctionne dans la catégorie des schémas
si l'on impose que notre foncteur soit un faisceau de Zariski.

\begin{lemma}
\label{topologies-lemma-extend}
Soit $S$ un schéma. Soit $\mathcal{C}$ une sous-catégorie pleine
de la catégorie $\Sch/S$ de tous les schémas au-dessus de $S$. Supposons que
\begin{enumerate}
\item si $X \to S$ est un objet de $\mathcal{C}$ et si
$U \subset X$ est un ouvert affine, alors $U \to S$ est isomorphe
à un objet de $\mathcal{C}$ ;
\item si $V$ est un schéma affine au-dessus d'un ouvert affine $U \subset S$
tel que $V \to U$ soit de présentation finie, alors $V \to S$ est isomorphe
à un objet de $\mathcal{C}$.
\end{enumerate}
Soit $F : \mathcal{C}^{opp} \to \textit{Ensembles}$ un foncteur.
Supposons que
\begin{enumerate}
\item[(a)] pour tout recouvrement de Zariski $\{f_i : X_i \to X\}_{i \in I}$
où $X, X_i$ sont des objets de $\mathcal{C}$, la condition de
faisceau soit satisfaite par $F$ et cette famille\footnote{Comme nous
ne savons pas si $X_i \times_X X_j$ appartient à $\mathcal{C}$,
il faut l'interpréter comme suit : d'après la propriété (1),
il existe des recouvrements de Zariski $\{U_{ijk} \to X_i \times_X X_j\}_{k \in K_{ij}}$
où $U_{ijk}$ est un objet de $\mathcal{C}$. La condition de faisceau
signifie alors que $F(X)$ est l'égalisateur des deux applications
de $\prod F(X_i)$ vers $\prod F(U_{ijk})$.} ;
\item[(b)] si $X = \lim X_i$ est une limite projective filtrante de schémas affines
au-dessus de $S$, où $X, X_i$ sont des objets de $\mathcal{C}$, alors
$F(X) = \colim F(X_i)$.
\end{enumerate}
Il existe alors une manière unique de prolonger $F$ en un foncteur
$F' : (\Sch/S)^{opp} \to \textit{Ensembles}$ satisfaisant aux analogues de
(a) et (b), c'est-à-dire que $F'$ satisfait la condition de faisceau pour tout
recouvrement de Zariski et que $F'(X) = \colim F'(X_i)$ dès que $X = \lim X_i$
est une limite projective filtrante de schémas affines au-dessus de $S$.
\end{lemma}

\begin{proof}
L'idée consiste d'abord à prolonger $F$ à une collection suffisamment grande de
schémas affines au-dessus de $S$, puis à utiliser la propriété de faisceau de Zariski pour le prolonger
à tous les schémas.

\medskip\noindent
Supposons que $V$ soit un schéma affine au-dessus de $S$ dont le morphisme structural
$V \to S$ se factorise par un ouvert affine $U \subset S$. Dans ce cas,
on peut écrire
$$
V = \lim V_i
$$
comme une limite projective filtrante, où $V_i \to U$ est de présentation finie
et $V_i$ est affine. Voir Algèbre, lemme \ref{algebra-lemma-ring-colimit-fp}.
D'après les conditions (1) et (2),
on peut remplacer les $V_i$ par des objets de $\mathcal{C}$.
Observons que $V_i \to S$ est localement de présentation finie
(si $S$ est quasi-séparé, ces morphismes sont même
de présentation finie). Posons alors
$$
F'(V) = \colim F(V_i)
$$
On peut en fait donner l'expression plus canonique
$$
F'(V) = \colim_{V \to V'} F(V')
$$
où la limite inductive est prise sur la catégorie des morphismes $V \to V'$ au-dessus de $S$
pour lesquels $V'$ est un objet de $\mathcal{C}$ dont le morphisme
structural $V' \to S$ est localement de présentation finie.
Cette expression coïncide avec la première parce que, d'après
Limits, proposition
\ref{limits-proposition-characterize-locally-finite-presentation}
notre système projectif $V_i$ est cofinal dans cette catégorie !
Enfin, notons que si $V$ était un objet de $\mathcal{C}$,
alors $F'(V) = F(V)$ d'après l'hypothèse (b).

\medskip\noindent
La seconde formule fait de $F'$ un foncteur contravariant
sur la catégorie formée des schémas affines $V$ au-dessus de $S$ dont le
morphisme structural se factorise par un ouvert affine de $S$.
Soit $V$ un tel schéma affine au-dessus de $S$ et
supposons que $V = \bigcup_{k = 1, \ldots, n} V_k$ soit un recouvrement ouvert fini
par des schémas affines. Il est alors légitime de demander si la condition de faisceau
est satisfaite par $F'$ et ce recouvrement ouvert.
C'est vrai et facile à montrer : écrivons $V = \lim V_i$ comme au
paragraphe précédent. D'après Limits, lemme \ref{limits-lemma-descend-opens},
pour tout $i$ suffisamment grand, on peut trouver des ouverts affines
$V_{i, k} \subset V_i$ compatibles aux morphismes de transition
et dont l'image inverse est $V_k$ dans $V$. Ainsi,
$$
F'(V_k) = \colim F(V_{i, k})
\quad\text{et}\quad
F'(V_k \cap V_l) = \colim F(V_{i, k} \cap V_{i, l})
$$
À strictement parler, il faut, dans ces formules, remplacer $V_{i, k}$ et
$V_{i, k} \cap V_{i, l}$ par des objets affines isomorphes de $\mathcal{C}$
avant d'appliquer le foncteur $F$.
Comme $I$ est filtrant, les limites inductives commutent aux égalisateurs.
La condition de faisceau (b) pour $F$ et les recouvrements de Zariski
$\{V_{i, k} \to V_i\}$ entraîne donc la condition de
faisceau pour $F'$ et ce recouvrement.

\medskip\noindent
Soit $X$ un schéma quelconque au-dessus de $S$. Notons $\mathcal{B}_X$
la collection des ouverts affines de $X$ dont le morphisme structural
vers $S$ se factorise par un ouvert affine de $S$. Il est clair que
$\mathcal{B}_X$ est une base de la topologie de $X$.
D'après le résultat du paragraphe précédent et Faisceaux, lemme
\ref{sheaves-lemma-cofinal-systems-coverings-standard-case}
on voit que $F'$ est un faisceau sur $\mathcal{B}_X$.
La restriction de $F'$ à $\mathcal{B}_X$
se prolonge donc uniquement en un faisceau $F'_X$ sur $X$ ; voir
Faisceaux, lemme \ref{sheaves-lemma-extend-off-basis}.
Si $X$ est un objet de $\mathcal{C}$, on dispose d'une
identification canonique $F'_X(X) = F(X)$, car $F'$ et $F$ coïncident
sur les objets où ils sont tous deux définis et
parce que $F$ satisfait la condition de faisceau pour les
recouvrements de Zariski.

\medskip\noindent
Soit $f : X \to Y$ un morphisme de schémas au-dessus de $S$.
On obtient une unique $f$-application de $F'_Y$ vers $F'_X$ compatible
avec les applications $F'(V) \to F'(U)$ pour tous
$U \in \mathcal{B}_X$ et $V \in \mathcal{B}_Y$
tels que $f(U) \subset V$ ; voir
Faisceaux, lemme \ref{sheaves-lemma-f-map-basis-above-and-below-structures}.
Nous omettons la vérification que ces applications se composent
correctement pour des morphismes $X \to Y \to Z$ de schémas au-dessus de $S$.
Nous omettons aussi la vérification que, si $f$ est un morphisme
de $\mathcal{C}$, alors l'application induite $F'_Y(Y) \to F'_X(X)$
coïncide avec l'application $F(Y) \to F(X)$ via les identifications
$F'_X(X) = F(X)$ et $F'_Y(Y) = F(Y)$ ci-dessus.
On voit ainsi que le prolongement cherché de
$F$ est le foncteur qui envoie $X/S$ sur $F'_X(X)$.

\medskip\noindent
La propriété (a) pour le foncteur $X \mapsto F'_X(X)$ est presque immédiate
d'après la construction ; nous omettons les détails.
Supposons que $X = \lim_{i \in I} X_i$
soit une limite projective filtrante de schémas affines au-dessus de $S$. Il faut montrer que
$$
F'_X(X) = \colim_{i \in I} F'_{X_i}(X_i)
$$
Supposons d'abord qu'il existe un $i \in I$ tel que
$X_i \to S$ se factorise par un ouvert affine $U \subset S$.
Alors $F'$ est défini sur $X$ et sur $X_{i'}$ pour $i' \geq i$,
et l'on a $F'_{X_{i'}}(X_{i'}) = F'(X_{i'})$ pour
$i' \geq i$ ainsi que $F'_X(X) = F'(X)$. Dans ce cas, toute flèche
$X \to V$, où $V$ est localement de présentation finie
au-dessus de $S$, se factorise sous la forme $X \to X_{i'} \to V$ pour un certain
$i' \geq i$ ; voir Limits, proposition
\ref{limits-proposition-characterize-locally-finite-presentation}.
On a donc
\begin{align*}
F'_X(X)
& =
F'(X)  \\
& = \colim_{X \to V} F(V) \\
& =
\colim_{i' \geq i} \colim_{X_{i'} \to V} F(V) \\
& =
\colim_{i' \geq i} F'(X_{i'}) \\
& =
\colim_{i' \geq i} F'_{X_{i'}}(X_{i'}) \\
& =
\colim_{i' \in I} F'_{X_{i'}}(X_{i'})
\end{align*}
comme voulu. Enfin, dans le cas général, choisissons un $i \in I$ et un
recouvrement ouvert affine fini $V_i = V_{i, 1} \cup \ldots \cup V_{i, n}$
tel que $V_{i, k} \to S$ se factorise par un ouvert affine de $S$.
Soient $V_k \subset V$ et $V_{i', k}$, pour $i' \geq i$,
les images inverses de $V_{i, k}$.
Le cas précédent donne
$$
F'_{V_k}(V_k) = \colim_{i' \geq i} F'_{V_{i', k}}(V_{i', k})
$$
et
$$
F'_{V_k \cap V_l}(V_k \cap V_l) =
\colim_{i' \geq i}
F'_{V_{i', k} \cap V_{i', l}}(V_{i', k} \cap V_{i', l})
$$
La propriété de faisceau et l'exactitude des limites inductives filtrantes
donnent $F'_X(X) = \colim_{i \in I} F'_{X_i}(X_i)$
également dans ce cas. Cela achève la démonstration de la propriété (b)
et donc celle du lemme.
\end{proof}

\begin{lemma}
\label{topologies-lemma-limit-fppf-topology}
Soit $\tau \in \{Zariski, \etale, smooth, syntomic, fppf\}$.
Soit $T$ un schéma affine écrit comme limite
$T = \lim_{i \in I} T_i$ d'un système projectif filtrant de schémas affines.
\begin{enumerate}
\item Soit $\mathcal{V} = \{V_j \to T\}_{j = 1, \ldots, m}$ un
recouvrement $\tau$ standard de $T$; voir les définitions
\ref{topologies-definition-standard-Zariski},
\ref{topologies-definition-standard-etale},
\ref{topologies-definition-standard-smooth},
\ref{topologies-definition-standard-syntomic} et
\ref{topologies-definition-standard-fppf}.
Il existe alors un indice $i$ et un recouvrement $\tau$ standard
$\mathcal{V}_i = \{V_{i, j} \to T_i\}_{j = 1, \ldots, m}$
dont le changement de base $T \times_{T_i} \mathcal{V}_i$ à $T$
est isomorphe à $\mathcal{V}$.
\item Soient $\mathcal{V}_i$, $\mathcal{V}'_i$ deux recouvrements
$\tau$ standards de $T_i$. Si
$f : T \times_{T_i} \mathcal{V}_i \to T \times_{T_i} \mathcal{V}'_i$ est
un morphisme de recouvrements de $T$, il existe un indice
$i' \geq i$ et un morphisme
$f_{i'} : T_{i'} \times_{T_i} \mathcal{V} \to
T_{i'} \times_{T_i} \mathcal{V}'_i$
dont le changement de base à $T$ est $f$.
\item Si
$f, g : \mathcal{V} \to \mathcal{V}'_i$
sont des morphismes de recouvrements $\tau$ standards de $T_i$ dont
les changements de base $f_T, g_T$ à $T$ sont égaux, alors il existe un
indice $i' \geq i$ tel que $f_{T_{i'}} = g_{T_{i'}}$.
\end{enumerate}
Autrement dit, la catégorie des recouvrements $\tau$ standards de $T$ est
la limite inductive, sur $I$, des catégories des recouvrements $\tau$ standards de $T_i$.
\end{lemma}

\begin{proof}
Démontrons l'énoncé pour $\tau = fppf$.
D'après Limits, lemme \ref{limits-lemma-descend-finite-presentation},
la catégorie des schémas de présentation finie sur $T$ est la
limite inductive, sur $I$, des catégories des schémas de présentation finie sur $T_i$. D'après
Limits, lemmes \ref{limits-lemma-descend-affine-finite-presentation}
et \ref{limits-lemma-descend-flat-finite-presentation},
il en va de même pour la catégorie des schémas affines, plats et
de présentation finie sur $T$.
Pour achever la démonstration du lemme, il suffit de montrer que, si
$\{V_{j, i} \to T_i\}_{j = 1, \ldots, m}$ est une famille finie de
morphismes plats de présentation finie, avec $V_{j, i}$ affine, et si le
changement de base $\coprod_j T \times_{T_i} V_{j, i} \to T$ est surjectif,
alors, pour un certain $i' \geq i$, le morphisme
$\coprod T_{i'} \times_{T_i} V_{j, i} \to T_{i'}$ est surjectif.
Notons $W_{i'} \subset T_{i'}$, respectivement $W \subset T$, l'image.
Par hypothèse, on a bien sûr $W = T$.
Puisque les morphismes sont plats et de présentation finie,
$W_i$ est un ouvert quasi-compact de $T_i$; voir
Morphismes, lemme \ref{morphisms-lemma-fppf-open}.
En outre, $W = T \times_{T_i} W_i$ (la formation de l'image commute
au changement de base). Par conséquent, d'après
Limits, lemme \ref{limits-lemma-descend-opens},
on a $W_{i'} = T_{i'}$ pour un indice $i'$ assez grand,
ce qui conclut.

\medskip\noindent
Pour $\tau \in \{Zariski, \etale, smooth, syntomic\}$, tout recouvrement $\tau$ standard
est un recouvrement fppf standard. La pleine fidélité du foncteur
en résulte. Il reste seulement à montrer que, si un recouvrement fppf standard
$\mathcal{V}_i$ est donné pour un certain $i$ et si $\mathcal{V}_i \times_{T_i} T$
est un recouvrement $\tau$ standard, alors $\mathcal{V}_i \times_{T_i} T_{i'}$
est un recouvrement $\tau$ standard pour tout $i' \gg i$. Cela résulte immédiatement
de Limits, lemmes
\ref{limits-lemma-descend-open-immersion},
\ref{limits-lemma-descend-etale},
\ref{limits-lemma-descend-smooth} et
\ref{limits-lemma-descend-syntomic}.
\end{proof}

\begin{lemma}
\label{topologies-lemma-extend-sheaf-general}
Supposons que $S$, $\mathcal{C}$ et $F$ satisfassent aux conditions (1), (2), (a) et (b) du
lemme \ref{topologies-lemma-extend}, et notons $F' : (\Sch/S)^{opp} \to \textit{Ensembles}$
le prolongement unique construit dans ce lemme. Soit
$\tau \in \{Zariski, \etale, smooth, syntomic, fppf\}$. Supposons que
\begin{enumerate}
\item[(c)] pour tout recouvrement $\tau$ standard $\{V_i \to V\}_{i = 1, \ldots, n}$
par des schémas affines dans $\Sch/S$, tel que $V \to S$ se factorise par un ouvert affine
$U \subset S$ et que $V \to U$ soit de présentation finie, la condition de faisceau
vaut pour $F$ et $\{V_i \to V\}_{i = 1, \ldots, n}$\footnote{Cela a
un sens, puisque $V$, $V_i$ et $V_i \times_V V_j$ sont isomorphes à des objets
de $\mathcal{C}$ d'après (2).}.
\end{enumerate}
Alors $F'$ satisfait à la condition de faisceau pour tout recouvrement $\tau$.
\end{lemma}

\begin{proof}
Soit $X$ un schéma sur $S$, et soit $\{X_i \to X\}_{i \in I}$
un recouvrement $\tau$. Soient $s_i \in F'(X_i)$ des éléments tels que
$s_i$ et $s_j$ aient même image dans $F'(X_i \times_X X_j)$
pour tous $i, j \in I$. Il faut montrer qu'il existe un unique
élément $s \in F'(X)$ dont la restriction est $s_i \in F'(X_i)$ pour tout $i \in I$.

\medskip\noindent
Cas particulier : $X$ est affine, son morphisme structural
est à valeurs dans un ouvert affine $U$ de $S$, et le recouvrement
$\{X_i \to X\}_{i \in I}$ est un recouvrement $\tau$ standard.
Dans ce cas, on peut écrire
$$
X = \lim V_k
$$
comme limite projective filtrante, où $V_k \to U$ est de présentation finie
et $V_k$ est affine. Voir Algèbre, lemme \ref{algebra-lemma-ring-colimit-fp}.
D'après lemme \ref{topologies-lemma-limit-fppf-topology}, il existe un $k$ et un recouvrement
$\tau$ standard $\{V_{k, i} \to V_k\}_{i \in I}$ dont le changement de base
à $X$ est le recouvrement donné. Pour $k' \geq k$, notons
$\{V_{k', i} \to V_{k'}\}_{i \in I}$ le changement de base à $V_{k'}$
de ce recouvrement. On a alors
\begin{align*}
F'(X)
& =
\colim_{k' \geq k} F(V_k) \\
& =
\colim_{k' \geq k}
\text{Equalizer}(
\xymatrix{
\prod F(V_{k', i})
\ar@<1ex>[r] \ar@<-1ex>[r] &
\prod F(V_{k', i} \times_{V_{k'}} V_{k', j})
}
\\
& =
\text{Equalizer}(
\xymatrix{
\colim_{k' \geq k}
\prod F(V_{k', i})
\ar@<1ex>[r] \ar@<-1ex>[r] &
\colim_{k' \geq k}
\prod F(V_{k', i} \times_{V_{k'}} V_{k', j})
}
\\
& =
\text{Equalizer}(
\xymatrix{
\prod F'(X_i)
\ar@<1ex>[r] \ar@<-1ex>[r] &
\prod F'(X_i \times_X X_j)
}
\end{align*}
La première égalité résulte de la construction de $F'$. La deuxième résulte de
l'hypothèse (c). La troisième résulte de l'exactitude des limites inductives filtrantes.
La quatrième résulte encore de la construction de $F'$.
On obtient ainsi la propriété de faisceau pour $F'$
relativement à $\{X_i \to X\}_{i \in I}$.

\medskip\noindent
Cas général. Choisissons un recouvrement ouvert affine $X = \bigcup U_k$ tel que
chaque $U_k$ soit à valeurs dans un ouvert affine de $S$. Pour tout $k$, choisissons
un recouvrement $\tau$ standard $\{V_{k, j} \to U_k\}_{j = 1, \ldots, m_k}$
qui raffine $\{X_i \times_X U_k \to U_k\}_{i \in I}$. Pour chaque
$j \in \{1, \ldots, m_k\}$, choisissons un indice $i_{k, j} \in I$
et un morphisme $g_{k, j} : V_{k, j} \to X_{i_{k, j}}$ au-dessus de $X$.
Soit $s_{k, j}$ l'élément de $F'(V_{k, j})$ obtenu en restreignant
$s_{i_{k, j}}$ par $g_{k, j}$. Observons que $s_{k, j}$ et $s_{k', j'}$
ont même restriction dans $F'(V_{k, j} \times_X V_{k', j'})$
pour tous $k$ et $k'$, tout $j \in \{1, \ldots, m_k\}$
et tout $j' \in \{1, \ldots, m_{k'}\}$; nous omettons la vérification.
En particulier, le paragraphe précédent donne un unique
élément $s_k \in F'(U_k)$ dont la restriction est $s_{k, j}$ pour tout $j$.
Avec ces notations, nous pouvons achever la démonstration.

\medskip\noindent
Unicité de $s$ : elle résulte de ce que $F'$ satisfait à la
propriété de faisceau pour les recouvrements de Zariski et que $s|_{U_k}$ doit être égal
à $s_k$, puisque tous deux ont pour restriction $s_{k, j}$ pour tout $j$.
Cette unicité montre alors que $s_k$ et $s_{k'}$ doivent induire
la même section de $F'$ sur (le schéma non affine) $U_k \cap U_{k'}$,
car ces sections induisent la même section sur le
recouvrement $\tau$ $\{V_{k, j} \times_X V_{k', j'} \to U_k \cap U_{k'}\}$.
La propriété de faisceau pour les recouvrements de Zariski donne donc une unique
section $s$ de $F'$ sur $X$ dont la restriction à $U_k$ est $s_k$.
Nous omettons la vérification (analogue à la précédente) que $s$ se restreint
en $s_i$ sur $X_i$.
\end{proof}

\begin{lemma}
\label{topologies-lemma-extend-sheaf}
Soit $\tau \in \{Zariski, \etale, smooth, syntomic, fppf\}$.
Soit $S$ un schéma contenu dans un gros site $\Sch_\tau$.
Soit $F : (\Sch/S)_\tau^{opp} \to \textit{Ensembles}$ un $\tau$-faisceau
satisfaisant à la propriété (b) de lemme \ref{topologies-lemma-extend} pour
$\mathcal{C} = (\Sch/S)_\tau$. Alors le prolongement
$F'$ de $F$ à la catégorie de tous les schémas sur $S$
satisfait à la condition de faisceau pour tout recouvrement $\tau$.
\end{lemma}

\begin{proof}
Cela résulte de lemme \ref{topologies-lemma-extend-sheaf-general}
appliqué avec $\mathcal{C} = (\Sch/S)_\tau$.
Les conditions (1), (2), (a) et (b) de lemme \ref{topologies-lemma-extend}
sont satisfaites; nous omettons les détails. On obtient ainsi l'unique prolongement $F'$
à la catégorie de tous les schémas sur $S$. Observons enfin que
tout recouvrement $\tau$ standard est tautologiquement équivalent à un
recouvrement dans $(\Sch/S)_\tau$; voir Ensembles, lemme \ref{sets-lemma-what-is-in-it}
ainsi que lemmes \ref{topologies-lemma-zariski-induced},
\ref{topologies-lemma-etale-induced},
\ref{topologies-lemma-smooth-induced},
\ref{topologies-lemma-syntomic-induced} et
\ref{topologies-lemma-fppf-induced}. D'après
Sites, lemme \ref{sites-lemma-tautological-same-sheaf}
la propriété de faisceau se conserve par équivalence tautologique
des recouvrements. Le fait que $F$ soit un $\tau$-faisceau entraîne donc
la propriété (c) de lemme \ref{topologies-lemma-extend-sheaf-general},
ce qui conclut.
\end{proof}









% OK, start here.
%

\chapter{Descente}


\phantomsection
\label{descent-section-phantom}


\section{Introduction}
\label{descent-section-introduction}

\noindent
Dans le chapitre consacré aux topologies sur les schémas
(voir Topologies, section \ref{topologies-section-introduction}), nous avons
introduit les recouvrements de Zariski, étales, fppf, lisses, syntomiques et
fpqc de schémas. Dans le présent chapitre, nous examinons les structures sur
les schémas qui peuvent être descendues au moyen de tels recouvrements. Voir, par
exemple, \cite{Gr-I}, \cite{Gr-II}, \cite{Gr-III}, \cite{Gr-IV}, \cite{Gr-V}
et \cite{Gr-VI}. Il s'agit également d'introduire les notions de descente, de
donnée de descente et de donnée de descente effective dans le cadre moins
formel des questions de descente pour les faisceaux quasi-cohérents, les
schémas, etc. La notion formelle de champ sur un site est étudiée dans le
chapitre consacré aux champs (voir Champs, section
\ref{stacks-section-introduction}).

\section{Données de descente pour les faisceaux quasi-cohérents}
\label{descent-section-equivalence}

\noindent
Dans ce chapitre, nous adoptons la convention suivant laquelle les
projections $\text{pr}_i : X \times \ldots \times X \to X$ sont numérotées à
partir de $i = 0$. Ainsi, nous avons
$\text{pr}_0, \text{pr}_1 : X \times X  \to X$,
$\text{pr}_0, \text{pr}_1, \text{pr}_2 : X \times X \times X  \to X$, etc.


\begin{definition}
\label{descent-definition-descent-datum-quasi-coherent}
Soit $S$ un schéma. Soit $\{f_i : S_i \to S\}_{i \in I}$ une famille de
morphismes de but $S$.
\begin{enumerate}
\item Une {\it donnée de descente $(\mathcal{F}_i, \varphi_{ij})$ pour les
faisceaux quasi-cohérents} relativement à la famille donnée est constituée
d'un faisceau quasi-cohérent $\mathcal{F}_i$ sur $S_i$ pour chaque $i \in I$
et d'un isomorphisme de $\mathcal{O}_{S_i \times_S S_j}$-modules
quasi-cohérents
$\varphi_{ij} : \text{pr}_0^*\mathcal{F}_i \to \text{pr}_1^*\mathcal{F}_j$
pour chaque paire $(i, j) \in I^2$,
tel que, pour tout triplet d'indices $(i, j, k) \in I^3$, le diagramme
$$
\xymatrix{
\text{pr}_0^*\mathcal{F}_i \ar[rd]_{\text{pr}_{01}^*\varphi_{ij}}
\ar[rr]_{\text{pr}_{02}^*\varphi_{ik}} & &
\text{pr}_2^*\mathcal{F}_k \\
& \text{pr}_1^*\mathcal{F}_j \ar[ru]_{\text{pr}_{12}^*\varphi_{jk}} &
}
$$
de $\mathcal{O}_{S_i \times_S S_j \times_S S_k}$-modules soit commutatif.
C'est la {\it condition de cocycle}.
\item Un {\it morphisme
$\psi : (\mathcal{F}_i, \varphi_{ij}) \to
(\mathcal{F}'_i, \varphi'_{ij})$ de données de descente} est une famille
$\psi = (\psi_i)_{i\in I}$ de morphismes de $\mathcal{O}_{S_i}$-modules
$\psi_i : \mathcal{F}_i \to \mathcal{F}'_i$ telle que tous les diagrammes
$$
\xymatrix{
\text{pr}_0^*\mathcal{F}_i \ar[r]_{\varphi_{ij}} \ar[d]_{\text{pr}_0^*\psi_i}
& \text{pr}_1^*\mathcal{F}_j \ar[d]^{\text{pr}_1^*\psi_j} \\
\text{pr}_0^*\mathcal{F}'_i \ar[r]^{\varphi'_{ij}} &
\text{pr}_1^*\mathcal{F}'_j \\
}
$$
soient commutatifs.
\end{enumerate}
\end{definition}

\noindent
Un bon exemple à garder à l'esprit est le suivant. Supposons que
$S = \bigcup S_i$ soit un recouvrement ouvert. Nous avons alors rencontré
les données de descente pour les faisceaux d'ensembles dans Faisceaux,
section \ref{sheaves-section-glueing-sheaves}, où nous les avons appelées
``données de recollement pour les faisceaux d'ensembles relativement au
recouvrement donné''. Nous avons en outre démontré que la catégorie des
données de recollement est équivalente à la catégorie des faisceaux sur $S$.
Nous établirons l'analogue dans le cadre ci-dessus lorsque
$\{S_i \to S\}_{i\in I}$ est un recouvrement fpqc.

\medskip\noindent
Dans le cas extrême où le recouvrement $\{S \to S\}$ est donné par
$\text{id}_S$, une donnée de descente est nécessairement de la forme
$(\mathcal{F}, \text{id}_\mathcal{F})$. La condition de cocycle garantit que
l'identité de $\mathcal{F}$ est le seul morphisme admis dans ce cas. Le lemme
suivant montre notamment qu'à tout faisceau quasi-cohérent de
$\mathcal{O}_S$-modules est associée une unique donnée de descente relativement
à toute famille donnée.

\begin{lemma}
\label{descent-lemma-refine-descent-datum}
Soient $\mathcal{U} = \{U_i \to U\}_{i \in I}$ et
$\mathcal{V} = \{V_j \to V\}_{j \in J}$ des familles de morphismes de
schémas à but fixé. Soit
$(g, \alpha : I \to J, (g_i)) : \mathcal{U} \to \mathcal{V}$ un morphisme
de familles de morphismes à but fixé, voir Sites, définition
\ref{sites-definition-morphism-coverings}. Soit
$(\mathcal{F}_j, \varphi_{jj'})$ une donnée de descente pour les faisceaux
quasi-cohérents relativement à la famille $\{V_j \to V\}_{j \in J}$. Alors :
\begin{enumerate}
\item le système
$$
\left(g_i^*\mathcal{F}_{\alpha(i)},
(g_i \times g_{i'})^*\varphi_{\alpha(i)\alpha(i')}\right)
$$
est une donnée de descente relativement à la famille
$\{U_i \to U\}_{i \in I}$ ;
\item cette construction est
fonctorielle en la donnée de descente $(\mathcal{F}_j, \varphi_{jj'})$ ;
\item étant donné un second morphisme
$(g', \alpha' : I \to J, (g'_i))$ de familles de morphismes à but fixé tel
que $g = g'$, il existe un isomorphisme fonctoriel de données de descente
$$
(g_i^*\mathcal{F}_{\alpha(i)},
(g_i \times g_{i'})^*\varphi_{\alpha(i)\alpha(i')})
\cong
((g'_i)^*\mathcal{F}_{\alpha'(i)},
(g'_i \times g'_{i'})^*\varphi_{\alpha'(i)\alpha'(i')}).
$$
\end{enumerate}
\end{lemma}

\begin{proof}
Omis. Indication : les morphismes
$g_i^*\mathcal{F}_{\alpha(i)} \to (g'_i)^*\mathcal{F}_{\alpha'(i)}$ qui
donnent l'isomorphisme des données de descente en (3) sont les images
inverses des morphismes $\varphi_{\alpha(i)\alpha'(i)}$ par
$(g_i, g'_i) : U_i \to V_{\alpha(i)} \times_V V_{\alpha'(i)}$.
\end{proof}

\noindent
Toute famille $\mathcal{U} = \{S_i \to S\}_{i \in I}$ raffine d'une unique
manière le recouvrement trivial $\{S \to S\}$. Si $\mathcal{F}$ est un
faisceau quasi-cohérent sur $S$, nous notons simplement
$(\mathcal{F}|_{S_i}, can)$ la donnée de descente relativement à
$\mathcal{U}$ obtenue par le procédé ci-dessus.

\begin{definition}
\label{descent-definition-descent-datum-effective-quasi-coherent}
Soit $S$ un schéma. Soit $\{S_i \to S\}_{i \in I}$ une famille de morphismes
de but $S$.
\begin{enumerate}
\item Soit $\mathcal{F}$ un $\mathcal{O}_S$-module quasi-cohérent. Nous
appelons l'unique donnée de descente sur $\mathcal{F}$ relativement au
recouvrement $\{S \to S\}$ la {\it donnée de descente triviale}.
\item L'image inverse de la donnée de descente triviale sur
$\{S_i \to S\}$ est appelée la {\it donnée de descente canonique}. Notation :
$(\mathcal{F}|_{S_i}, can)$.
\item Une donnée de descente $(\mathcal{F}_i, \varphi_{ij})$ pour les
faisceaux quasi-cohérents relativement au recouvrement donné est dite
{\it effective} s'il existe un faisceau quasi-cohérent $\mathcal{F}$ sur $S$
tel que $(\mathcal{F}_i, \varphi_{ij})$ soit isomorphe à
$(\mathcal{F}|_{S_i}, can)$.
\end{enumerate}
\end{definition}

\begin{lemma}
\label{descent-lemma-zariski-descent-effective}
Soit $S$ un schéma et soit $S = \bigcup U_i$ un recouvrement ouvert. Toute
donnée de descente sur des faisceaux quasi-cohérents relativement à la famille
$\mathcal{U} = \{U_i \to S\}$ est effective. De plus, le foncteur de la
catégorie des $\mathcal{O}_S$-modules quasi-cohérents vers la catégorie des
données de descente relativement à $\mathcal{U}$ est pleinement fidèle.
\end{lemma}

\begin{proof}
Cela résulte immédiatement de Faisceaux, section
\ref{sheaves-section-glueing-sheaves}, et du fait que la quasi-cohérence est
une propriété locale, voir Modules, définition
\ref{modules-definition-quasi-coherent}.
\end{proof}

\noindent
Pour aller plus loin, nous devons d'abord étudier le cas des modules sur des
anneaux.










\section{Descente pour les modules}
\label{descent-section-descent-modules}

\noindent
Soit $R \to A$ un morphisme d'anneaux. D'après Simplicial, exemple
\ref{simplicial-example-push-outs-simplicial-object}, il donne naissance à une
$R$-algèbre cosimpliciale
$$
\xymatrix{
A
\ar@<1ex>[r]
\ar@<-1ex>[r]
&
A \otimes_R A
\ar@<0ex>[l]
\ar@<2ex>[r]
\ar@<0ex>[r]
\ar@<-2ex>[r]
&
A \otimes_R A \otimes_R A
\ar@<1ex>[l]
\ar@<-1ex>[l]
}
$$
Notons-la $(A/R)_\bullet$, de sorte que $(A/R)_n$ est le produit tensoriel de
$(n + 1)$ copies de $A$ sur $R$. Étant donné une application
$\varphi : [n] \to [m]$, le morphisme de $R$-algèbres
$(A/R)_\bullet(\varphi)$ est l'application
$$
a_0 \otimes \ldots \otimes a_n
\longmapsto
\prod\nolimits_{\varphi(i) = 0} a_i
\otimes
\prod\nolimits_{\varphi(i) = 1} a_i
\otimes \ldots \otimes
\prod\nolimits_{\varphi(i) = m} a_i
$$
où, par convention, le produit vide vaut $1$. Les premiers morphismes, avec
les notations de Simplicial, section
\ref{simplicial-section-cosimplicial-object}, sont donc
$$
\begin{matrix}
\delta^1_0 & : & a_0 & \mapsto & 1 \otimes a_0 \\
\delta^1_1 & : & a_0 & \mapsto & a_0 \otimes 1 \\
\sigma^0_0 & : & a_0 \otimes a_1 & \mapsto & a_0a_1 \\
\delta^2_0 & : & a_0 \otimes a_1 & \mapsto & 1 \otimes a_0 \otimes a_1 \\
\delta^2_1 & : & a_0 \otimes a_1 & \mapsto & a_0 \otimes 1 \otimes a_1 \\
\delta^2_2 & : & a_0 \otimes a_1 & \mapsto & a_0 \otimes a_1 \otimes 1 \\
\sigma^1_0 & : & a_0 \otimes a_1 \otimes a_2 & \mapsto & a_0a_1 \otimes a_2 \\
\sigma^1_1 & : & a_0 \otimes a_1 \otimes a_2 & \mapsto & a_0 \otimes a_1a_2
\end{matrix}
$$
et ainsi de suite.

\medskip\noindent
Un $R$-module $M$ donne naissance au $(A/R)_\bullet$-module cosimplicial
$(A/R)_\bullet \otimes_R M$. Autrement dit,
$M_n = (A/R)_n \otimes_R M$, et les morphismes de $R$-algèbres
$(A/R)_n \to (A/R)_m$ servent à définir les morphismes correspondants sur
$M \otimes_R (A/R)_\bullet$.

\medskip\noindent
L'analogue, pour les modules, d'une donnée de descente pour les faisceaux
quasi-cohérents est le suivant.

\begin{definition}
\label{descent-definition-descent-datum-modules}
Soit $R \to A$ un morphisme d'anneaux.
\begin{enumerate}
\item Une {\it donnée de descente $(N, \varphi)$ pour les modules
relativement à $R \to A$} est constituée d'un $A$-module $N$ et d'un
isomorphisme de $A \otimes_R A$-modules
$$
\varphi : N \otimes_R A \to A \otimes_R N
$$
tel que la {\it condition de cocycle} soit satisfaite : le diagramme de
morphismes de $A \otimes_R A \otimes_R A$-modules
$$
\xymatrix{
N \otimes_R A \otimes_R A \ar[rr]_{\varphi_{02}}
\ar[rd]_{\varphi_{01}}
& &
A \otimes_R A \otimes_R N \\
& A \otimes_R N \otimes_R A \ar[ru]_{\varphi_{12}} &
}
$$
est commutatif (voir ci-dessous pour les notations).
\item Un {\it morphisme $(N, \varphi) \to (N', \varphi')$ de données de
descente} est un morphisme de $A$-modules $\psi : N \to N'$ tel que le
diagramme
$$
\xymatrix{
N \otimes_R A \ar[r]_\varphi \ar[d]_{\psi \otimes \text{id}_A} &
A \otimes_R N \ar[d]^{\text{id}_A \otimes \psi} \\
N' \otimes_R A \ar[r]^{\varphi'} &
A \otimes_R N'
}
$$
soit commutatif.
\end{enumerate}
\end{definition}

\noindent
Dans cette définition, nous posons
$\varphi_{01} = \varphi \otimes \text{id}_A$,
$\varphi_{12} = \text{id}_A \otimes \varphi$, et
$\varphi_{02}(n \otimes 1 \otimes 1) = \sum a_i \otimes 1 \otimes n_i$
si $\varphi(n \otimes 1) = \sum a_i \otimes n_i$. Ce sont trois
homomorphismes de $A \otimes_R A \otimes_R A$-modules. De manière
équivalente, nous avons
$$
\varphi_{ij}
=
\varphi \otimes_{(A/R)_1, \ (A/R)_\bullet(\tau^2_{ij})} (A/R)_2
$$
où $\tau^2_{ij} : [1] \to [2]$ est l'application
$0 \mapsto i$, $1 \mapsto j$. En effet,
$(A/R)_{\bullet}(\tau^2_{02})(a_0 \otimes a_1) =
a_0 \otimes 1 \otimes a_1$, et de même pour les autres\footnote{Remarquons
que $\tau^2_{ij} = \delta^2_k$ si
$\{i, j, k\} = [2] = \{0, 1, 2\}$, voir Simplicial, définition
\ref{simplicial-definition-face-degeneracy}.}.

\medskip\noindent
Nous avons besoin de quelques notations supplémentaires pour énoncer le lemme
suivant. Soit $(N, \varphi)$ une donnée de descente relativement à un
morphisme d'anneaux $R \to A$. Pour $n \geq 0$ et $i \in [n]$, posons
$$
N_{n, i} =
A \otimes_R
\ldots
\otimes_R A \otimes_R N \otimes_R A \otimes_R
\ldots
\otimes_R A
$$
où le facteur $N$ occupe la $i$-ième place. C'est un $(A/R)_n$-module.
Introduisons les applications $\tau^n_i : [0] \to [n]$, $0 \mapsto i$.
Alors
$$
N_{n, i} = N \otimes_{(A/R)_0, \ (A/R)_\bullet(\tau^n_i)} (A/R)_n
$$
Pour $0 \leq i \leq j \leq n$, soit
$\tau^n_{ij} : [1] \to [n]$ l'application qui envoie $0$ sur $i$ et $1$ sur
$j$. Comme ci-dessus, l'homomorphisme $\varphi$ induit des isomorphismes
$$
\varphi^n_{ij}
=
\varphi \otimes_{(A/R)_1, \ (A/R)_\bullet(\tau^n_{ij})} (A/R)_n :
N_{n, i} \longrightarrow N_{n, j}
$$
de $(A/R)_n$-modules lorsque $i < j$. Si $i = j$, nous posons
$\varphi^n_{ij} = \text{id}$. Tous ces morphismes étant des isomorphismes, ils
permettent de déplacer le facteur $N$ vers n'importe quelle place. La
condition de cocycle signifie précisément que le résultat ne dépend pas de la
manière dont on effectue ce déplacement (par exemple en composant deux de ces
isomorphismes ou en une seule étape). Enfin, pour toute application
$\beta : [n] \to [m]$, nous définissons le morphisme
$$
N_{\beta, i} : N_{n, i} \to N_{m, \beta(i)}
$$
comme l'unique application semi-linéaire relativement à l'homomorphisme
$(A/R)_\bullet(\beta)$ telle que
$$
N_{\beta, i}(1 \otimes \ldots \otimes n \otimes \ldots \otimes 1)
=
1 \otimes \ldots \otimes n \otimes \ldots \otimes 1
$$
pour tout $n \in N$. Cela suggère le lemme suivant.

\begin{lemma}
\label{descent-lemma-descent-datum-cosimplicial}
Soit $R \to A$ un morphisme d'anneaux. À toute donnée de descente
$(N, \varphi)$, on peut associer un $(A/R)_\bullet$-module cosimplicial
$N_\bullet$\footnote{Il faudrait en réalité écrire
$(N, \varphi)_\bullet$.} en posant $N_n = N_{n, n}$ et, pour
$\beta : [n] \to [m]$, en définissant
$$
N_\bullet(\beta) = (\varphi^m_{\beta(n)m}) \circ N_{\beta, n} :
N_{n, n} \longrightarrow N_{m, m}.
$$
Cette construction est fonctorielle en la donnée de descente.
\end{lemma}

\begin{proof}
Voici les premières applications, où
$\varphi(n \otimes 1) = \sum \alpha_i \otimes x_i$ :
$$
\begin{matrix}
\delta^1_0 & : & N & \to & A \otimes N & n & \mapsto & 1 \otimes n \\
\delta^1_1 & : & N & \to & A \otimes N & n & \mapsto &
\sum \alpha_i \otimes x_i\\
\sigma^0_0 & : & A \otimes N & \to & N & a_0 \otimes n & \mapsto & a_0n \\
\delta^2_0 & : & A \otimes N & \to & A \otimes A \otimes N &
a_0 \otimes n & \mapsto & 1 \otimes a_0 \otimes n \\
\delta^2_1 & : & A \otimes N & \to & A \otimes A \otimes N &
a_0 \otimes n & \mapsto & a_0 \otimes 1 \otimes n \\
\delta^2_2 & : & A \otimes N & \to & A \otimes A \otimes N &
a_0 \otimes n & \mapsto & \sum a_0 \otimes \alpha_i \otimes x_i \\
\sigma^1_0 & : & A \otimes A \otimes N & \to & A \otimes N &
a_0 \otimes a_1 \otimes n & \mapsto & a_0a_1 \otimes n \\
\sigma^1_1 & : & A \otimes A \otimes N & \to & A \otimes N &
a_0 \otimes a_1 \otimes n & \mapsto & a_0 \otimes a_1n
\end{matrix}
$$
avec les notations de
Simplicial, section \ref{simplicial-section-cosimplicial-object}.
Vérifions d'abord les deux identités $\sigma^0_0 \circ \delta^1_0 = \text{id}$
et $\sigma^0_0 \circ \delta^1_1 = \text{id}$.
La première, $\sigma^0_0 \circ \delta^1_0 = \text{id}$, résulte immédiatement
de la description explicite des morphismes ci-dessus.
Pour démontrer la seconde, il faut utiliser la condition de cocycle
(car elle n'est pas vérifiée pour un isomorphisme arbitraire
$\varphi : N \otimes_R A \to A \otimes_R N$). Posons
$p = \sigma^0_0 \circ \delta^1_1 : N \to N$. D'après la description
des applications ci-dessus, $p$ s'identifie aussi à
$$
p = \varphi \otimes \text{id} :
N = (N \otimes_R A) \otimes_{(A \otimes_R A)} A
\longrightarrow
(A \otimes_R N) \otimes_{(A \otimes_R A)} A = N
$$
Puisque $\varphi$ est un isomorphisme, $p$ en est un également.
Écrivons $\varphi(n \otimes 1) = \sum \alpha_i \otimes x_i$, avec
$\alpha_i \in A$ et $x_i \in N$. Alors $p(n) = \sum \alpha_ix_i$.
Écrivons ensuite
$\varphi(x_i \otimes 1) = \sum \alpha_{ij} \otimes y_j$, avec
$\alpha_{ij} \in A$ et $y_j \in N$. La condition de cocycle donne
$$
\sum \alpha_i \otimes \alpha_{ij} \otimes y_j
=
\sum \alpha_i \otimes 1 \otimes x_i.
$$
Cela signifie que $p(n) = \sum \alpha_ix_i = \sum \alpha_i\alpha_{ij}y_j =
\sum \alpha_i p(x_i) = p(p(n))$. Ainsi, $p$ est un projecteur ; comme
c'est un isomorphisme, c'est l'identité.

\medskip\noindent
Pour achever de montrer que $N_\bullet$ est un module cosimplicial,
il faut vérifier les cinq types de relations de
Simplicial, remarque \ref{simplicial-remark-relations-cosimplicial}.
Les relations qui ne font intervenir que des composées de $\sigma$ sont
évidentes. Celles qui ne font intervenir que des composées de $\delta$ se
ramènent à la condition de cocycle pour $\varphi$.
On vérifie exactement comme ci-dessus les relations
$\sigma_j \circ \delta_j = \sigma_j \circ \delta_{j + 1} = \text{id}$.
Enfin, les autres relations entre composées de $\delta$ et de $\sigma$
sont vérifiées quel que soit $\varphi$.
\end{proof}

\noindent
À tout $R$-module $M$, on peut associer une donnée de descente canonique,
à savoir $(M \otimes_R A, can)$, où
$can : (M \otimes_R A) \otimes_R A \to A \otimes_R (M \otimes_R A)$
est l'application évidente :
$(m \otimes a) \otimes a' \mapsto a \otimes (m \otimes a')$.

\begin{lemma}
\label{descent-lemma-canonical-descent-datum-cosimplicial}
Soit $R \to A$ un homomorphisme d'anneaux.
Soit $M$ un $R$-module. Le
$(A/R)_\bullet$-module cosimplicial associé à la donnée de descente canonique
est isomorphe au module cosimplicial $(A/R)_\bullet \otimes_R M$.
\end{lemma}

\begin{proof}
Démonstration omise.
\end{proof}

\begin{definition}
\label{descent-definition-descent-datum-effective-module}
Soit $R \to A$ un homomorphisme d'anneaux.
Une donnée de descente $(N, \varphi)$ est dite {\it effective}
s'il existe un $R$-module $M$ et un isomorphisme de données de descente
de $(M \otimes_R A, can)$ vers
$(N, \varphi)$.
\end{definition}

\noindent
Soit $R \to A$ un homomorphisme d'anneaux.
Soit $(N, \varphi)$ une donnée de descente.
Considérons le complexe de cochaînes $s(N_\bullet)$ associé
à $N_\bullet$ (voir
Simplicial, section \ref{simplicial-section-dold-kan-cosimplicial}).
Il est de la forme
$$
N \to A \otimes_R N \to A \otimes_R A \otimes_R N \to \ldots
$$
On peut en décrire les différentielles.
La première est l'application
$$
n \longmapsto 1 \otimes n - \varphi(n \otimes 1).
$$
La seconde est donnée, sur les tenseurs purs, par
$$
a \otimes n \longmapsto 1 \otimes a \otimes n
- a \otimes 1 \otimes n + a \otimes \varphi(n \otimes 1).
$$
La règle se poursuit de manière évidente.

\medskip\noindent
Dans le cas particulier où
$N = A \otimes_R M$, pour tout $m \in M$,
l'élément $1 \otimes m$ appartient au noyau de la première différentielle
du complexe de cochaînes associé au module cosimplicial
$(A/R)_\bullet \otimes_R M$. On obtient donc un complexe de cochaînes augmenté
\begin{equation}
\label{descent-equation-extended-complex}
0 \to M \to A \otimes_R M \to A \otimes_R A \otimes_R M \to \ldots
\end{equation}
On place ici $0$ en degré $-2$,
le module $M$ en degré $-1$, le module $A \otimes_R M$
en degré $0$, et ainsi de suite. Ce complexe est de la forme
$$
0 \to R \to A \to A \otimes_R A \to A \otimes_R A \otimes_R A \to \ldots
$$
lorsque $M = R$.

\begin{lemma}
\label{descent-lemma-with-section-exact}
Supposons que $R \to A$ admette une section.
Alors, pour tout $R$-module $M$, le complexe de cochaînes augmenté
(\ref{descent-equation-extended-complex}) est exact.
\end{lemma}

\begin{proof}
D'après
Simplicial, lemme \ref{simplicial-lemma-push-outs-simplicial-object-w-section},
l'application $R \to (A/R)_\bullet$ est un homotopisme
de $R$-algèbres cosimpliciales
(où $R$ désigne la $R$-algèbre cosimpliciale constante).
Par suite, $M \to (A/R)_\bullet \otimes_R M$ est
un homotopisme dans la catégorie des
$R$-modules cosimpliciaux, car $\otimes_R M$ est
un foncteur de la catégorie des $R$-algèbres vers celle
des $R$-modules ; voir
Simplicial, lemme \ref{simplicial-lemma-functorial-homotopy}.
L'application induite entre les complexes associés est donc
un homotopisme ; voir
Simplicial, lemme \ref{simplicial-lemma-homotopy-s-Q}.
Puisque le complexe associé au $R$-module cosimplicial constant
$M$ est le complexe
$$
\xymatrix{
M \ar[r]^0 & M \ar[r]^1 & M \ar[r]^0 & M \ar[r]^1 & M \ldots
}
$$
le résultat s'ensuit (le complexe augmenté s'obtient simplement en ajoutant
un terme $M$ au début).
\end{proof}

\begin{lemma}
\label{descent-lemma-ff-exact}
Supposons que $R \to A$ soit fidèlement plat ; voir
Algèbre, définition \ref{algebra-definition-flat}.
Alors, pour tout $R$-module $M$, le complexe de cochaînes augmenté
(\ref{descent-equation-extended-complex}) est exact.
\end{lemma}

\begin{proof}
Supposons que l'on puisse trouver un homomorphisme d'anneaux fidèlement plat
$R \to R'$ tel que le résultat soit vrai pour l'homomorphisme
$R' \to A' = R' \otimes_R A$. Le résultat vaut alors aussi pour $R \to A$.
En effet, pour tout $R$-module $M$, le module cosimplicial
$(M \otimes_R R') \otimes_{R'} (A'/R')_\bullet$
s'identifie au module cosimplicial $R' \otimes_R (M \otimes_R (A/R)_\bullet)$.
L'annulation de la cohomologie du complexe associé à
$(M \otimes_R R') \otimes_{R'} (A'/R')_\bullet$ entraîne donc celle
de la cohomologie du complexe associé à
$M \otimes_R (A/R)_\bullet$, par fidèle platitude de $R \to R'$.
Il en va de même pour l'annulation des groupes de cohomologie en degrés
$-1$ et $0$ du complexe augmenté (démonstration omise).

\medskip\noindent
Or une telle extension fidèlement plate existe : il suffit de prendre $R' = A$,
car l'homomorphisme d'anneaux $R' = A \to A' = A \otimes_R A$ admet la section
$a \otimes a' \mapsto aa'$, et le
lemme \ref{descent-lemma-with-section-exact} s'applique.
\end{proof}

\noindent
Voici le lien entre ce complexe et la question de l'effectivité.

\begin{lemma}
\label{descent-lemma-recognize-effective}
Soit $R \to A$ un homomorphisme d'anneaux fidèlement plat.
Soit $(N, \varphi)$ une donnée de descente.
Alors $(N, \varphi)$ est effective si et seulement si l'application canonique
$$
A \otimes_R H^0(s(N_\bullet)) \longrightarrow N
$$
est un isomorphisme.
\end{lemma}

\begin{proof}
Si $(N, \varphi)$ est effective, on peut écrire $N = A \otimes_R M$,
avec $\varphi = can$. On a alors $H^0(s(N_\bullet)) = M$, d'après
les lemmes \ref{descent-lemma-canonical-descent-datum-cosimplicial}
et \ref{descent-lemma-ff-exact}. Réciproquement, supposons que l'application du lemme
soit un isomorphisme. Posons $M = H^0(s(N_\bullet))$.
C'est un $R$-sous-module de $N$,
à savoir $M = \{n \in N \mid 1 \otimes n = \varphi(n \otimes 1)\}$.
Il reste seulement à vérifier que, par l'isomorphisme
$A \otimes_R M \to N$,
la donnée de descente canonique correspond à $\varphi$.
Nous omettons cette vérification.
\end{proof}

\begin{lemma}
\label{descent-lemma-descent-descends}
Soit $R \to A$ un homomorphisme d'anneaux fidèlement plat, et soit $R \to R'$
un homomorphisme fidèlement plat. Posons $A' = R' \otimes_R A$.
Si toutes les données de descente relatives à $R' \to A'$ sont effectives,
il en est de même de toutes les données de descente relatives à $R \to A$.
\end{lemma}

\begin{proof}
Soit $(N, \varphi)$ une donnée de descente relative à $R \to A$.
Posons $N' = R' \otimes_R N = A' \otimes_A N$, et notons
$\varphi' = \text{id}_{R'} \otimes \varphi$ le changement de base
de la donnée de descente $\varphi$. Alors $(N', \varphi')$ est
une donnée de descente relative à $R' \to A'$ et
$H^0(s(N'_\bullet)) = R' \otimes_R H^0(s(N_\bullet))$.
De plus, l'application
$A' \otimes_{R'} H^0(s(N'_\bullet)) \to N'$
s'identifie au changement de base de l'homomorphisme de $A$-modules
$A \otimes_R H^0(s(N)) \to N$ par l'homomorphisme fidèlement plat
$A \to A'$. On conclut donc par le lemme \ref{descent-lemma-recognize-effective}.
\end{proof}

\noindent
Voici le résultat principal de cette section.
Sa démonstration peut paraître un peu laborieuse ; pour une approche
plus conceptuelle, voir la remarque
\ref{descent-remark-homotopy-equivalent-cosimplicial-algebras} ci-dessous.

\begin{proposition}
\label{descent-proposition-descent-module}
\begin{slogan}
Descente effective des modules par les homomorphismes d'anneaux fidèlement plats.
\end{slogan}
Soit $R \to A$ un homomorphisme d'anneaux fidèlement plat.
Alors
\begin{enumerate}
\item toute donnée de descente de modules relative à $R \to A$
est effective ;
\item le foncteur $M \mapsto (A \otimes_R M, can)$ de la catégorie des $R$-modules
vers celle des données de descente est une équivalence ;
\item le foncteur inverse est donné par $(N, \varphi) \mapsto H^0(s(N_\bullet))$.
\end{enumerate}
\end{proposition}

\begin{proof}
Nous démontrons seulement (1) et omettons les démonstrations de (2) et (3).
Puisque $R \to A$ est fidèlement plat, il existe un changement de base
fidèlement plat $R \to R'$ tel que $R' \to A' = R' \otimes_R A$
admette une section (il suffit de prendre $R' = A$, comme dans la preuve du
lemme \ref{descent-lemma-ff-exact}). Ainsi, par le
lemme \ref{descent-lemma-descent-descends}, on peut supposer que
$R \to A$ admet une section, notée $\sigma : A \to R$.
Soit $(N, \varphi)$ une donnée de descente relative à $R \to A$.
Posons
$$
M = H^0(s(N_\bullet)) = \{n \in N \mid 1 \otimes n = \varphi(n \otimes 1)\}
\subset
N
$$
D'après le lemme \ref{descent-lemma-recognize-effective}, il suffit de montrer que
$A \otimes_R M \to N$ est un isomorphisme.

\medskip\noindent
Soit $n \in N$. Écrivons
$\varphi(n \otimes 1) = \sum a_i \otimes x_i$, avec
$a_i \in A$ et $x_i \in N$. D'après le lemme \ref{descent-lemma-descent-datum-cosimplicial},
on a $n = \sum a_i x_i$ dans $N$ (puisque
$\sigma^0_0 \circ \delta^1_1 = \text{id}$ dans tout objet cosimplicial).
Écrivons ensuite $\varphi(x_i \otimes 1) = \sum a_{ij} \otimes y_j$, avec
$a_{ij} \in A$ et $y_j \in N$.
La condition de cocycle s'écrit
$$
\sum a_i \otimes a_{ij} \otimes y_j = \sum a_i \otimes 1 \otimes x_i
$$
dans $A \otimes_R A \otimes_R N$. On en déduit deux identités :
\begin{enumerate}
\item en appliquant $\sigma$ au premier facteur $A$, on obtient
$\sum \sigma(a_i) \varphi(x_i \otimes 1) = \sum \sigma(a_i) \otimes x_i$ ;
\item en appliquant $\sigma$ au facteur $A$ du milieu, on obtient
$\sum_i a_i \otimes \sum_j \sigma(a_{ij}) y_j = \sum a_i \otimes x_i$.
\end{enumerate}
Le point (1) montre que $\sum \sigma(a_i) x_i \in M$. En l'appliquant à
$x_i$, on voit que $\sum \sigma(a_{ij})y_i \in M$ pour tout $i$.
En effectuant les multiplications dans l'égalité (2), on obtient
$\sum_i a_i (\sum_j \sigma(a_{ij}) y_j) = \sum a_i x_i = n$.
Ainsi, $A \otimes_R M \to N$ est surjectif.
Enfin, supposons que $m_i \in M$ et $\sum a_i m_i = 0$.
En appliquant $\varphi$ à
$\sum a_im_i \otimes 1$, on obtient $\sum a_i \otimes m_i = 0$.
Autrement dit, $A \otimes_R M \to N$ est injectif, ce qui achève la preuve.
\end{proof}

\begin{remark}
\label{descent-remark-standard-covering}
Soit $R$ un anneau. Soient $f_1, \ldots, f_n\in R$ des éléments engendrant
l'idéal unité. L'anneau $A = \prod_i R_{f_i}$ est une
$R$-algèbre fidèlement plate. L'anneau cosimplicial $(A/R)_\bullet$
a pour composante de degré $n$ l'anneau
$$
\prod\nolimits_{i_0, \ldots, i_n} R_{f_{i_0}\ldots f_{i_n}}
$$
Les résultats ci-dessus redonnent donc
Algèbre, lemmes \ref{algebra-lemma-standard-covering},
\ref{algebra-lemma-cover-module} et \ref{algebra-lemma-glue-modules}.
Ils vont toutefois plus loin, grâce à l'exactitude en degrés supérieurs.
Celle-ci entraîne en effet la trivialité de la cohomologie de {\v C}ech
des faisceaux quasi-cohérents sur les schémas affines.
On obtient ainsi une seconde démonstration de
Cohomologie des schémas, lemme
\ref{coherent-lemma-cech-cohomology-quasi-coherent-trivial}.
\end{remark}

\begin{remark}
\label{descent-remark-homotopy-equivalent-cosimplicial-algebras}
Soit $R$ un anneau. Soit $A_\bullet$ une $R$-algèbre cosimpliciale.
Dans ce cadre, une donnée de descente correspond à un
$A_\bullet$-module cosimplicial $M_\bullet$ tel que, pour tous
$n, m \geq 0$ et tout $\varphi : [n] \to [m]$,
l'application $M(\varphi) : M_n \to M_m$ induise un isomorphisme
$$
M_n \otimes_{A_n, A(\varphi)} A_m \longrightarrow M_m.
$$
Appelons un tel module cosimplicial un {\it module cartésien}.
Dans ce cadre, la démonstration de la proposition \ref{descent-proposition-descent-module}
se décompose selon les étapes suivantes :
\begin{enumerate}
\item Si $R \to R'$ et $R \to A$ sont fidèlement plats,
les données de descente relatives à $A/R$ sont effectives dès que
celles relatives à $(R' \otimes_R A)/R'$ le sont.
\item Soit $A$ une $R$-algèbre. Les données de descente relatives à $A/R$
correspondent aux $(A/R)_\bullet$-modules cartésiens.
\item Si $R \to A$ admet une section, $(A/R)_\bullet$ est homotopiquement
équivalent à $R$, la $R$-algèbre cosimpliciale constante de valeur $R$.
\item Si $A_\bullet \to B_\bullet$ est une équivalence d'homotopie
de $R$-algèbres cosimpliciales, le foncteur
$M_\bullet \mapsto M_\bullet \otimes_{A_\bullet} B_\bullet$
induit une équivalence entre la catégorie des
$A_\bullet$-modules cartésiens et celle des $B_\bullet$-modules cartésiens.
\end{enumerate}
Pour (1), voir le lemme \ref{descent-lemma-descent-descends}.
Le point (2) utilise le lemme \ref{descent-lemma-descent-datum-cosimplicial}.
Le point (3) a été établi dans la preuve du lemme \ref{descent-lemma-with-section-exact}
(il repose sur Simplicial,
lemme \ref{simplicial-lemma-push-outs-simplicial-object-w-section}).
Enfin, le point (4) est immédiat si l'on adopte le bon point de vue !
\end{remark}








\section{Descente pour les morphismes universellement injectifs}
\label{descent-section-descent-universally-injective}

\noindent
De nombreuses constructions en géométrie algébrique font appel aux techniques
de {\it descente} : on construit des objets sur un espace donné en travaillant
d'abord sur un espace un peu plus grand qui se projette sur celui-ci,
ou l'on vérifie une propriété d'un espace ou d'un morphisme après
changement de base par un morphisme de recouvrement.
L'utilité de ces techniques dépend bien entendu de l'identification
d'une classe suffisamment large de {\it morphismes de descente effective}.
Dès les débuts de la géométrie algébrique moderne développée par Grothendieck,
on a montré que les morphismes {\it quasi-compacts} et {\it fidèlement plats}
permettent de faire descendre effectivement les objets, les morphismes
et nombre de leurs propriétés.

\medskip\noindent
Comme d'habitude, cet énoncé se ramène à une propriété des anneaux et des modules.
Pour qu'un homomorphisme $f: R \to S$ soit un morphisme de descente effective
pour les modules, Grothendieck a montré qu'il suffit que $f$ soit fidèlement plat.
Cette condition exclut toutefois de nombreux exemples naturels :
ainsi, tout homomorphisme d'anneaux scindé est un morphisme de descente effective.
Un exemple naturel intervient même dans la démonstration de la descente
fidèlement plate : pour tout homomorphisme d'anneaux $f: R \to S$,
l'homomorphisme $1_S \otimes f: S \to S \otimes_R S$
admet la multiplication pour rétraction, qu'il soit plat ou non.

\medskip\noindent
On peut donc chercher une condition naturelle sur un homomorphisme d'anneaux
qui entraîne la descente effective des modules et englobe à la fois
les homomorphismes fidèlement plats et les homomorphismes scindés.
Le lecteur sera peut-être surpris (l'auteur l'a du moins été)
d'apprendre qu'une réponse complète est connue depuis les environs de 1970 !
Il est facile de vérifier que, pour que $f: R \to S$ soit un morphisme
de descente effective pour les modules, il faut que $f$ soit
{\it universellement injectif} dans la catégorie des $R$-modules :
pour tout $R$-module $M$, l'application
$1_M  \otimes f: M \to M \otimes_R S$ doit être injective.
Or cette condition est aussi suffisante.
Par exemple, si $f$ est scindé dans la catégorie des $R$-modules
(sans l'être nécessairement dans celle des anneaux),
alors $f$ est un morphisme de descente effective pour les modules.

\medskip\noindent
L'histoire de ce résultat est assez complexe : il a d'abord été énoncé
par Olivier \cite{olivier}, qui appelait {\it purs} les morphismes
universellement injectifs, mais sans indication claire de démonstration.
On peut le déduire des travaux de Joyal et Tierney \cite{joyal-tierney} ;
cependant, à notre connaissance, la première démonstration autonome publiée
est celle de Mesablishvili \cite{mesablishvili1}.
Le premier objectif de cette section est d'exposer cette démonstration.
Hormis la correction de coquilles, nous modifions peu la présentation originale,
si ce n'est pour expliciter le lien entre la définition usuelle
d'une donnée de descente en géométrie algébrique et celle
employée dans \cite{mesablishvili1}.
La démonstration est entièrement catégorique et peut donc être formulée
dans le langage des monades, ce qui permet de l'appliquer
dans d'autres contextes ; voir \cite{janelidze-tholen}.

\medskip\noindent
Le second objectif est de rassembler des résultats sur les propriétés
des modules, des algèbres et des morphismes que l'on peut faire descendre le long
d'homomorphismes d'anneaux universellement injectifs.
Les cas des modules de type fini et des modules plats ont été traités
par Mesablishvili \cite{mesablishvili2}.


\subsection{Préliminaires de théorie des catégories}
\label{descent-subsection-category-prelims}

\noindent
Commençons par rappeler quelques notions élémentaires de théorie des catégories
qui simplifieront l'exposé. Dans cette sous-section, on fixe une catégorie ambiante.

\medskip\noindent
Étant donnés deux morphismes $g_1, g_2: B \to C$, rappelons qu'un
{\it égalisateur} de $g_1$ et $g_2$ est un morphisme $f: A \to B$ tel que
$g_1 \circ f = g_2 \circ f$ et universel pour cette propriété.
Cela signifie que tout diagramme commutatif
$$
\xymatrix{A' \ar[rd]^e \ar@/^1.5pc/[rrd] \ar@{-->}[d] & & \\
A \ar[r]^f & B \ar@<1ex>[r]^{g_1} \ar@<-1ex>[r]_{g_2} & 
C
}
$$
dont on a omis la flèche en pointillé peut être complété de façon unique.
On dit également que le diagramme
\begin{equation}
\label{descent-equation-equalizer}
\xymatrix{
A \ar[r]^f  & B \ar@<1ex>[r]^{g_1} \ar@<-1ex>[r]_{g_2} & C
}
\end{equation}
est un égalisateur. En inversant les flèches, on obtient la définition
d'un {\it coégalisateur}.
Voir Catégories, sections \ref{categories-section-equalizers} et
\ref{categories-section-coequalizers}.

\medskip\noindent
La propriété d'être un égalisateur étant définie par une propriété universelle,
elle n'est généralement pas préservée par l'application d'un foncteur covariant.
Comme pour les monomorphismes et les épimorphismes, on peut, dans certains cas,
remédier à cette difficulté en exhibant des scindages.

\begin{definition}
\label{descent-definition-split-equalizer}
Un {\it égalisateur scindé} est un diagramme (\ref{descent-equation-equalizer})
tel que $g_1 \circ f = g_2 \circ f$, pour lequel il existe des morphismes auxiliaires
$h : B \to A$ et $i : C \to B$ vérifiant
\begin{equation}
\label{descent-equation-split-equalizer-conditions}
h \circ f = 1_A, \quad f \circ h = i \circ g_1, \quad i \circ g_2 = 1_B.
\end{equation}
\end{definition}

\noindent
Les identités entre les flèches imposent alors que (\ref{descent-equation-equalizer})
soit un égalisateur : l'application $e$ se factorise de façon unique par $f$
en écrivant $e = f \circ (h \circ e)$.
Ainsi, l'image d'un égalisateur scindé par un foncteur covariant
est un égalisateur scindé ; son image par un foncteur contravariant
est un {\it coégalisateur scindé}, dont la définition est évidente.

\subsection{Morphismes universellement injectifs}
\label{descent-subsection-universally-injective}

\noindent
Rappelons que $\textit{Rings}$ désigne la catégorie des anneaux commutatifs
avec $1$. Pour un objet $R$ de $\textit{Rings}$, on note $\text{Mod}_R$
la catégorie des $R$-modules.

\begin{remark}
\label{descent-remark-reflects}
Tout foncteur $F : \mathcal{A} \to \mathcal{B}$ entre catégories abéliennes
qui est exact et envoie les objets non nuls sur des objets non nuls
reflète les injections et les surjections. En effet, l'exactitude implique
que $F$ préserve les noyaux et les conoyaux (comparer avec
Homologie, section \ref{homology-section-functors}).
Par exemple, si $f : R \to S$ est un homomorphisme d'anneaux fidèlement plat,
alors $\bullet \otimes_R S: \text{Mod}_R \to \text{Mod}_S$ possède ces propriétés.
\end{remark}

\noindent
Soit $R$ un anneau. Rappelons qu'un morphisme $f : M \to N$ de $\text{Mod}_R$
est {\it universellement injectif} si, pour tout $P \in \text{Mod}_R$,
le morphisme $f \otimes 1_P: M \otimes_R P \to N \otimes_R P$ est injectif.
Voir Algèbre, définition \ref{algebra-definition-universally-injective}.

\begin{definition}
\label{descent-definition-universally-injective}
Un homomorphisme d'anneaux $f: R \to S$ est {\it universellement injectif}
s'il est universellement injectif comme morphisme de $\text{Mod}_R$.
\end{definition}

\begin{example}
\label{descent-example-split-injection-universally-injective}
Toute injection scindée dans $\text{Mod}_R$ est universellement injective.
En particulier, toute injection scindée dans $\textit{Rings}$
est universellement injective.
\end{example}

\begin{example}
\label{descent-example-cover-universally-injective}
Pour un anneau $R$ et des éléments $f_1, \ldots, f_n \in R$ engendrant l'idéal
unité, le morphisme $R \to R_{f_1} \oplus \ldots \oplus R_{f_n}$
est universellement injectif. Bien que cela résulte immédiatement du
lemme \ref{descent-lemma-faithfully-flat-universally-injective},
il est instructif de le vérifier directement : on se ramène aussitôt
au cas où $R$ est local. L'un des $f_i$ est alors inversible,
et l'application $R \to R_{f_i}$ est donc un isomorphisme.
\end{example}

\begin{lemma}
\label{descent-lemma-faithfully-flat-universally-injective}
Tout homomorphisme d'anneaux fidèlement plat est universellement injectif.
\end{lemma}

\begin{proof}
C'est une reformulation de
Algèbre, lemme \ref{algebra-lemma-faithfully-flat-universally-injective}.
\end{proof}

\noindent
L'observation essentielle de \cite{mesablishvili1} est que l'injectivité
universelle se reformule utilement en termes de scindage, grâce à la
construction usuelle d'un cogénérateur injectif dans $\text{Mod}_R$.

\begin{definition}
\label{descent-definition-C}
Soit $R$ un anneau. On définit le foncteur contravariant
{\it $C$} $ : \text{Mod}_R \to \text{Mod}_R$ par
$$
C(M) = \Hom_{\textit{Ab}}(M, \mathbf{Q}/\mathbf{Z}),
$$
l'action de $R$ sur $C(M)$ étant donnée par $rf(s) = f(rs)$.
\end{definition}

\noindent
Ce foncteur était noté $M \mapsto M^\vee$ dans
Compléments d'algèbre, section \ref{more-algebra-section-injectives-modules}.

\begin{lemma}
\label{descent-lemma-C-is-faithful}
Pour tout anneau $R$, le foncteur $C : \text{Mod}_R \to \text{Mod}_R$
est exact et permet de reconnaître les injections et les surjections.
\end{lemma}

\begin{proof}
L'exactitude est établie dans
Compléments d'algèbre, lemme \ref{more-algebra-lemma-vee-exact}.
Les autres propriétés s'en déduisent ; voir la remarque \ref{descent-remark-reflects}.
\end{proof}

\begin{remark}
\label{descent-remark-adjunction}
Nous utiliserons fréquemment l'adjonction usuelle entre $\Hom$ et
le produit tensoriel, sous la forme de l'isomorphisme naturel de foncteurs
contravariants
\begin{equation}
\label{descent-equation-adjunction}
C(\bullet_1 \otimes_R \bullet_2) \cong \Hom_R(\bullet_1, C(\bullet_2)): 
\text{Mod}_R \times \text{Mod}_R \to \text{Mod}_R
\end{equation}
qui associe à $f: M_1 \otimes_R M_2 \to \mathbf{Q}/\mathbf{Z}$
l'application $m_1 \mapsto 
(m_2 \mapsto f(m_1 \otimes m_2))$. Voir
Algèbre, lemme \ref{algebra-lemma-hom-from-tensor-product-variant}.
Il en résulte que, si
$$
\xymatrix@C=9pc{
C(M) \ar@<1ex>[r] \ar@<-1ex>[r] & C(N) \ar[r] & C(P)
}
$$
est un diagramme de coégalisateur scindé dans $\text{Mod}_R$,
il en est de même de
$$
\xymatrix@C=9pc{
C(M \otimes_R Q) \ar@<1ex>[r] \ar@<-1ex>[r] & C(N \otimes_R Q) \ar[r] & C(P 
\otimes_R Q)
}
$$
pour tout $Q \in \text{Mod}_R$.
\end{remark}

\begin{lemma}
\label{descent-lemma-split-surjection}
Soit $R$ un anneau. Un morphisme $f: M \to N$ de $\text{Mod}_R$
est universellement injectif si et seulement si
$C(f): C(N) \to C(M)$ est une surjection scindée.
\end{lemma}

\begin{proof}
D'après (\ref{descent-equation-adjunction}), pour tout $P \in \text{Mod}_R$,
on a un diagramme commutatif
$$
\xymatrix@C=9pc{
\Hom_R( P, C(N)) \ar[r]_{\Hom_R(P,C(f))} \ar[d]^{\cong} &
\Hom_R(P,C(M)) \ar[d]^{\cong} \\
C(P \otimes_R N ) \ar[r]^{C(1_{P} \otimes f)} & C(P \otimes_R M ).
}
$$
Si $f$ est universellement injectif, alors $1_{C(M)} \otimes f: C(M) \otimes_R M \to 
C(M) \otimes_R N$ est injectif ; les deux flèches horizontales du diagramme
ci-dessus sont donc surjectives pour $P = C(M)$.
On peut ainsi relever $1_{C(M)} \in \Hom_R(C(M), C(M))$
en un élément $g \in \Hom_R(C(N), C(M))$ qui scinde $C(f)$.
Réciproquement, si $C(f)$ est une surjection scindée,
les deux flèches horizontales du diagramme ci-dessus sont surjectives.
Le lemme \ref{descent-lemma-C-is-faithful} entraîne alors que $1_{P} \otimes f$ est injectif.
\end{proof}

\begin{remark}
\label{descent-remark-functorial-splitting}
Soit $f: M \to N$ un morphisme universellement injectif de $\text{Mod}_R$.
En choisissant un scindage $g$ de $C(f)$, on peut construire
un scindage fonctoriel de $C(1_P \otimes f)$ pour tout $P \in \text{Mod}_R$.
En effet, d'après (\ref{descent-equation-adjunction}), il suffit de scinder
$\Hom_R(P, 
C(f))$ fonctoriellement en $P$,
ce que réalise l'application $g \circ \bullet$.
\end{remark}


\subsection{Descente pour les modules et leurs morphismes}
\label{descent-subsection-descent-modules-morphisms}

\noindent
Dans toute cette sous-section, on fixe un homomorphisme d'anneaux $f: R \to S$.
Comme on l'a vu dans la section \ref{descent-section-descent-modules},
les données de descente de modules peuvent se formuler dans le langage
des algèbres cosimpliciales. Nous préférons ici une présentation plus élémentaire.

\medskip\noindent
Pour $i = 1, 2, 3$, notons $S_i$ le produit tensoriel de $i$ copies
de $S$ sur $R$.
Définissons les homomorphismes d'anneaux $\delta_0^1, \delta_1^1: S_1 \to S_2$,
$\delta_{01}^1, \delta_{02}^1, \delta_{12}^1: S_1 \to S_3$ et
$\delta_0^2, \delta_1^2, \delta_2^2: S_2 \to S_3$ par les formules
\begin{align*}
\delta^1_0  (a_0) & =  1 \otimes a_0 \\
\delta^1_1  (a_0) & = a_0 \otimes 1 \\
\delta^2_0  (a_0 \otimes a_1) & =  1 \otimes a_0 \otimes a_1 \\
\delta^2_1  (a_0 \otimes a_1) & =  a_0 \otimes 1 \otimes a_1 \\
\delta^2_2  (a_0 \otimes a_1) & =  a_0 \otimes a_1 \otimes 1 \\
\delta_{01}^1(a_0) & = 1 \otimes 1 \otimes a_0 \\
\delta_{02}^1(a_0) & = 1 \otimes a_0 \otimes 1 \\
\delta_{12}^1(a_0) & = a_0 \otimes 1 \otimes 1.
\end{align*}
Autrement dit, l'indice supérieur indique l'anneau de départ,
et l'indice inférieur les positions où insérer des facteurs égaux à 1.
(Cette notation est compatible avec celle introduite dans la
section \ref{descent-section-descent-modules}.)

\medskip\noindent
Rappelons\footnote{Plus précisément, notre $\theta$ est ici l'inverse de
$\varphi$ dans la définition \ref{descent-definition-descent-datum-modules}.}
que, d'après la définition \ref{descent-definition-descent-datum-modules}, pour
$M \in \text{Mod}_S$, une {\it donnée de descente} sur $M$ relative à $f$
est un isomorphisme
$$
\theta :
M \otimes_{S,\delta^1_0} S_2
\longrightarrow
M \otimes_{S,\delta^1_1} S_2
$$
de $S_2$-modules satisfaisant la {\it condition de cocycle}
\begin{equation}
\label{descent-equation-cocycle-condition}
(\theta \otimes \delta_2^2) \circ (\theta \otimes \delta_2^0) = (\theta \otimes 
\delta_2^1):
M \otimes_{S, \delta^1_{01}} S_3 \to M \otimes_{S,\delta^1_{12}} S_3.
\end{equation}
Notons $DD_{S/R}$ la catégorie des $S$-modules munis de données de descente
relatives à $f$.

\medskip\noindent
Par exemple, pour $M_0 \in \text{Mod}_R$, le choix d'un isomorphisme
$M \cong M_0 \otimes_R S$ fournit une donnée de descente en identifiant
naturellement $M \otimes_{S,\delta^1_0} S_2$ et
$M \otimes_{S,\delta^1_1} S_2$ à $M_0 \otimes_R S_2$.
Cette construction définit en particulier un foncteur
$f^*: \text{Mod}_R \to DD_{S/R}$.

\begin{definition}
\label{descent-definition-effective-descent}
Le foncteur $f^*: \text{Mod}_R \to DD_{S/R}$
est appelé {\it extension des scalaires par $f$}.
On dit que $f$ est un {\it morphisme de descente pour les modules}
si $f^*$ est pleinement fidèle.
On dit que $f$ est un {\it morphisme de descente effective pour les modules}
si $f^*$ est une équivalence de catégories.
\end{definition}

\noindent
Notre but est de montrer que, lorsque $f$ est universellement injectif,
on peut utiliser $\theta$ pour retrouver $M_0$ à l'intérieur de $M$.
Cette construction repose de manière essentielle sur certains
diagrammes égalisateurs.

\begin{lemma}
\label{descent-lemma-equalizer-M}
Pour $(M,\theta) \in DD_{S/R}$, le diagramme
\begin{equation}
\label{descent-equation-equalizer-M}
\xymatrix@C=8pc{
M \ar[r]^{\theta \circ (1_M \otimes \delta_0^1)} &
M \otimes_{S, \delta_1^1} S_2
\ar@<1ex>[r]^{(\theta \otimes \delta_2^2) \circ (1_M \otimes \delta^2_0)}
\ar@<-1ex>[r]_{1_{M \otimes S_2} \otimes \delta^2_1} & 
M \otimes_{S, \delta_{12}^1} S_3
}
\end{equation}
est un égalisateur scindé.
\end{lemma}

\begin{proof}
Définissons les homomorphismes d'anneaux $\sigma^0_0: S_2 \to S_1$ et
$\sigma_0^1, 
\sigma_1^1: S_3 \to S_2$ par les formules
\begin{align*}
\sigma^0_0 (a_0 \otimes a_1) & = a_0a_1 \\
\sigma^1_0 (a_0 \otimes a_1 \otimes a_2) & = a_0a_1 \otimes a_2 \\
\sigma^1_1 (a_0 \otimes a_1 \otimes a_2) & = a_0 \otimes a_1a_2.
\end{align*}
On prend pour morphismes auxiliaires
$1_M \otimes \sigma_0^0: M \otimes_{S, \delta_1^1} S_2 \to M$
et $1_M \otimes \sigma_0^1: M \otimes_{S,\delta_{12}^1} S_3 \to M \otimes_{S, 
\delta_1^1} S_2$.
Parmi les compatibilités requises dans (\ref{descent-equation-split-equalizer-conditions}),
la première s'obtient en tensorisant la condition de cocycle
(\ref{descent-equation-cocycle-condition}) par $\sigma_1^1$ ;
les autres sont immédiates.
\end{proof}

\begin{lemma}
\label{descent-lemma-equalizer-CM}
Pour $(M, \theta) \in DD_{S/R}$, le diagramme
\begin{equation}
\label{descent-equation-coequalizer-CM}
\xymatrix@C=8pc{
C(M \otimes_{S, \delta_{12}^1} S_3)
\ar@<1ex>[r]^{C((\theta \otimes \delta_2^2) \circ (1_M \otimes \delta^2_0))}
\ar@<-1ex>[r]_{C(1_{M \otimes S_2} \otimes \delta^2_1)} &
C(M \otimes_{S, \delta_1^1} S_2 )
\ar[r]^{C(\theta \circ (1_M \otimes \delta_0^1))} & C(M).
}
\end{equation}
obtenu en appliquant $C$ à (\ref{descent-equation-equalizer-M})
est un coégalisateur scindé.
\end{lemma}

\begin{proof}
Démonstration omise.
\end{proof}

\begin{lemma}
\label{descent-lemma-equalizer-S}
Le diagramme
\begin{equation}
\label{descent-equation-equalizer-S}
\xymatrix@C=8pc{
S_1 \ar[r]^{\delta^1_1} &
S_2 \ar@<1ex>[r]^{\delta^2_2} \ar@<-1ex>[r]_{\delta^2_1} & 
S_3
}
\end{equation}
est un égalisateur scindé.
\end{lemma}

\begin{proof}
Dans le lemme \ref{descent-lemma-equalizer-M}, prendre $(M, \theta) = f^*(S)$.
\end{proof}

\noindent
Cela suggère la définition d'un candidat au rôle de quasi-inverse de $f^*$.

\begin{definition}
\label{descent-definition-pushforward}
On définit le foncteur {\it $f_*$} $: DD_{S/R} \to \text{Mod}_R$
en prenant pour $f_*(M, \theta)$ le $R$-sous-module de $M$
tel que le diagramme
\begin{equation}
\label{descent-equation-equalizer-f}
\xymatrix@C=8pc{f_*(M,\theta) \ar[r] & M \ar@<1ex>^{\theta \circ (1_M \otimes 
\delta_0^1)}[r] \ar@<-1ex>_{1_M \otimes \delta_1^1}[r] & 
M \otimes_{S, \delta_1^1} S_2 
}
\end{equation}
soit un égalisateur.
\end{definition}

\noindent
Le lemme \ref{descent-lemma-equalizer-M} et le fait que le foncteur de restriction
$\text{Mod}_S \to \text{Mod}_R$ est adjoint à droite au foncteur
d'extension des scalaires $\bullet \otimes_R S: \text{Mod}_R \to \text{Mod}_S$
montrent que $f_*$ est adjoint à droite à $f^*$.

\medskip\noindent
Nous pouvons maintenant démontrer le lemme clé. Dans le cas fidèlement plat,
il est immédiat (voir la remarque \ref{descent-remark-descent-lemma}),
mais le cas général demande une démonstration.

\begin{lemma}
\label{descent-lemma-descent-lemma}
Si $f$ est universellement injectif, le diagramme
\begin{equation}
\label{descent-equation-equalizer-f2}
\xymatrix@C=8pc{
f_*(M, \theta) \otimes_R S
\ar[r]^{\theta \circ (1_M \otimes \delta_0^1)} &
M \otimes_{S, \delta_1^1} S_2 
\ar@<1ex>[r]^{(\theta \otimes \delta_2^2) \circ (1_M \otimes \delta^2_0)}
\ar@<-1ex>[r]_{1_{M \otimes S_2} \otimes \delta^2_1} &
M \otimes_{S, \delta_{12}^1} S_3
}
\end{equation}
obtenu en tensorisant (\ref{descent-equation-equalizer-f}) sur $R$ par $S$
est un égalisateur.
\end{lemma}

\begin{proof}
D'après le lemme \ref{descent-lemma-split-surjection} et la remarque
\ref{descent-remark-functorial-splitting}, l'application
$C(1_N \otimes f): C(N \otimes_R S) \to C(N)$ admet une section
fonctorielle en $N$. On obtient ainsi les flèches verticales supérieures
du diagramme commutatif
$$
\xymatrix@C=8pc{
C(M \otimes_{S, \delta_1^1} S_2)
\ar@<1ex>^{C(\theta \circ (1_M \otimes \delta_0^1))}[r]
\ar@<-1ex>_{C(1_M \otimes \delta_1^1)}[r] \ar[d] &
C(M) \ar[r]\ar[d] & C(f_*(M,\theta)) \ar@{-->}[d] \\
C(M \otimes_{S,\delta_{12}^1} S_3)
\ar@<1ex>^{C((\theta \otimes \delta_2^2) \circ (1_M \otimes \delta^2_0))}[r]
\ar@<-1ex>_{C(1_{M \otimes S_2} \otimes \delta^2_1)}[r] \ar[d] &
C(M \otimes_{S, \delta_1^1} S_2 )
\ar[r]^{C(\theta \circ (1_M \otimes \delta_0^1))}
\ar[d]^{C(1_M \otimes \delta_1^1)} &
C(M) \ar[d] \ar@{=}[dl] \\
C(M \otimes_{S, \delta_1^1} S_2)
\ar@<1ex>[r]^{C(\theta \circ (1_M \otimes \delta_0^1))}
\ar@<-1ex>[r]_{C(1_M \otimes \delta_1^1)} &
C(M) \ar[r] &
C(f_*(M,\theta))
}
$$
dans lequel les composées le long des colonnes sont les morphismes
identités. La deuxième ligne est le diagramme coégalisateur
(\ref{descent-equation-coequalizer-CM}), d'où la flèche en pointillés.
Le carré supérieur droit fournit des morphismes auxiliaires
$C(f_*(M,\theta)) \to 
C(M)$ et $C(M) \to C(M\otimes_{S,\delta_1^1} S_2)$ qui montrent
que la première ligne est un diagramme coégalisateur scindé.
D'après la remarque \ref{descent-remark-adjunction}, on peut tensoriser par $S$
à l'intérieur de $C$ pour obtenir le diagramme coégalisateur scindé
$$
\xymatrix@C=8pc{
C(M \otimes_{S,\delta_2^2 \circ \delta_1^1} S_3)
\ar@<1ex>^{C((\theta \otimes \delta_2^2) \circ (1_M \otimes \delta^2_0))}[r] 
\ar@<-1ex>_{C(1_{M \otimes S_2} \otimes \delta^2_1)}[r] &
C(M \otimes_{S, \delta_1^1} S_2 )
\ar[r]^{C(\theta \circ (1_M \otimes \delta_0^1))} &
C(f_*(M,\theta) \otimes_R S).
}
$$
Le lemme \ref{descent-lemma-C-is-faithful} montre alors que
(\ref{descent-equation-equalizer-f2}) est lui aussi un égalisateur.
\end{proof}

\begin{remark}
\label{descent-remark-descent-lemma}
Si $f$ est une injection scindée dans $\text{Mod}_R$, on peut simplifier
la démonstration en scindant directement $f$, sans utiliser $C$.
Le cas où $f$ est fidèlement plat est encore plus simple :
la conclusion du lemme \ref{descent-lemma-descent-lemma} est immédiate,
car la tensorisation sur $R$ par $S$ préserve tous les égalisateurs.
\end{remark}

\begin{theorem}
\label{descent-theorem-descent}
Les conditions suivantes sont équivalentes.
\begin{enumerate}
\item[(a)] Le morphisme $f$ est un morphisme de descente pour les modules.
\item[(b)] Le morphisme $f$ est un morphisme de descente effective pour les modules.
\item[(c)] Le morphisme $f$ est universellement injectif.
\end{enumerate}
\end{theorem}

\begin{proof}
Il est clair que (b) implique (a). Montrons que (a) implique (c).
Si $f$ n'est pas universellement injectif, il existe
$M \in \text{Mod}_R$ tel que l'application
$1_M \otimes f: M \to M \otimes_R S$ ait un noyau non nul $N$.
La projection naturelle $M \to M/N$ n'est pas un isomorphisme,
mais son image dans $DD_{S/R}$ en est un.
Ainsi, $f^*$ n'est pas pleinement fidèle.

\medskip\noindent
Montrons enfin que (c) implique (b). D'après les lemmes
\ref{descent-lemma-equalizer-M} et \ref{descent-lemma-descent-lemma}, pour
$(M, \theta) \in DD_{S/R}$, l'application naturelle
$f^* f_*(M,\theta) \to M$ est un isomorphisme de $S$-modules.
Pour $M_0 \in \text{Mod}_R$, montrons que $M_0 \to f_*f^*M_0$
est un isomorphisme. Par adjonction, la composée
$f^*M_0 \to f^*f_*f^*M_0 \to f^*M_0$ est l'identité, et ce qui précède
montre que la deuxième flèche est un isomorphisme. Ainsi,
$f^*M_0 \to f^*f_*f^*M_0$ est un isomorphisme. Le lemme
\ref{algebra-lemma-base-change-along-universally-injective-ring-map}
du chapitre Algèbre montre que $M_0 \to f_* f^* M_0$ est un isomorphisme.
Par conséquent, $f_*$ et $f^*$ sont des foncteurs quasi-inverses,
ce qui achève la démonstration.
\end{proof}

\subsection{Descente des propriétés des modules}
\label{descent-subsection-descent-properties-modules}

\noindent
Dans toute cette sous-section, on fixe un morphisme d'anneaux
universellement injectif $f : R \to S$, un objet $M \in \text{Mod}_R$
et un morphisme d'anneaux $R \to A$. Nous cherchons les propriétés de
$M$ ou de $A$ qui peuvent se vérifier après extension des scalaires
par $f$. Commençons par quelques résultats de \cite{mesablishvili2}.

\begin{lemma}
\label{descent-lemma-flat-to-injective}
Si $M \in \text{Mod}_R$ est plat, alors $C(M)$ est un $R$-module injectif.
\end{lemma}

\begin{proof}
Soit $0 \to N \to P \to Q \to 0$ une suite exacte dans $\text{Mod}_R$.
Puisque $M$ est plat, la suite
$$
0 \to N \otimes_R M \to P \otimes_R M \to Q \otimes_R M \to 0
$$
est exacte. Le lemme \ref{descent-lemma-C-is-faithful} montre que la suite
$$
0 \to C(Q \otimes_R M) \to C(P \otimes_R M) \to C(N \otimes_R M) \to 0
$$
est exacte. En utilisant (\ref{descent-equation-adjunction}), celle-ci s'écrit
$$
0 \to \Hom_R(Q, C(M)) \to \Hom_R(P, C(M)) \to \Hom_R(N, C(M)) \to 0.
$$
Ainsi, $C(M)$ est un objet injectif de $\text{Mod}_R$.
\end{proof}

\begin{theorem}
\label{descent-theorem-descend-module-properties}
Si $M \otimes_R S$ possède l'une des propriétés suivantes comme $S$-module :
\begin{enumerate}
\item[(a)]
être de type fini ;
\item[(b)]
être de présentation finie ;
\item[(c)]
être plat ;
\item[(d)]
être fidèlement plat ;
\item[(e)]
être projectif de type fini ;
\end{enumerate}
alors $M$ possède la même propriété comme $R$-module, et réciproquement.
\end{theorem}

\begin{proof}
Pour démontrer (a), choisissons une famille finie $\{n_i\}$ de générateurs
de $M \otimes_R S$ dans $\text{Mod}_S$.
Écrivons chaque $n_i$ sous la forme $\sum_j m_{ij} \otimes s_{ij}$,
avec $m_{ij} \in M$ et $s_{ij} \in S$.
Soit $F$ le $R$-module libre de type fini de base $e_{ij}$, et soit
$F \to M$ le morphisme de $R$-modules qui envoie $e_{ij}$ sur $m_{ij}$.
Alors $F \otimes_R S\to M \otimes_R S$ est surjectif,
donc $\Coker(F \to M) \otimes_R S$ est nul, d'où $\Coker(F \to M)$ est nul.
Ceci démontre (a).

\medskip\noindent
Pour (b), supposons que $M \otimes_R S$ soit de présentation finie.
D'après (a), $M$ est de type fini. Choisissons une surjection
$R^{\oplus n} \to M$ de noyau $K$.
La suite $K \otimes_R S \to S^{\oplus r} \to M \otimes_R S \to 0$ est exacte.
D'après le lemme \ref{algebra-lemma-extension} du chapitre Algèbre,
le noyau de $S^{\oplus r} \to M \otimes_R S$ est un $S$-module de type fini.
Il existe donc un nombre fini d'éléments $k_1, \ldots, k_t \in K$
tels que les images de $k_i \otimes 1$ dans $S^{\oplus r}$
engendrent le noyau de $S^{\oplus r} \to M \otimes_R S$.
Soit $K' \subset K$ le sous-module engendré par $k_1, \ldots, k_t$.
Alors $M' = R^{\oplus r}/K'$ est un $R$-module de présentation finie,
muni d'un morphisme $M' \to M$ tel que
$M' \otimes_R S \to M \otimes_R S$ soit un isomorphisme.
Ainsi, $M' \cong M$, comme voulu.

\medskip\noindent
Pour démontrer (c), soit $0 \to M' \to M'' \to M \to 0$
une suite exacte courte dans $\text{Mod}_R$.
Le foncteur $\bullet \otimes_R S$ étant exact à droite,
l'application $M'' \otimes_R S \to M \otimes_R S$ est surjective.
D'après le lemme \ref{descent-lemma-C-is-faithful}, l'application
$C(M \otimes_R S) \to C(M'' \otimes_R S)$ est donc injective.
Si $M \otimes_R S$ est plat, le lemme \ref{descent-lemma-flat-to-injective}
montre que $C(M \otimes_R S)$ est un objet injectif de $\text{Mod}_S$.
L'injection $C(M \otimes_R S) \to C(M'' \otimes_R S)$ est donc scindée
dans $\text{Mod}_S$, et par suite dans $\text{Mod}_R$.
Comme $C(M \otimes_R S) \to C(M)$ est une surjection scindée
d'après le lemme \ref{descent-lemma-split-surjection}, il en résulte que
$C(M) \to C(M'')$ est une injection scindée dans $\text{Mod}_R$.
Autrement dit, la suite
$$
0 \to C(M) \to C(M'') \to C(M') \to 0
$$
est exacte scindée. Pour $N \in \text{Mod}_R$,
la relation (\ref{descent-equation-adjunction}) montre que la suite
$$
0 \to C(M \otimes_R N) \to C(M'' \otimes_R N) \to C(M' \otimes_R N) \to 0
$$
est exacte scindée. D'après le lemme \ref{descent-lemma-C-is-faithful}, la suite
$$
0 \to M' \otimes_R N \to M'' \otimes_R N \to M \otimes_R N \to 0
$$
est exacte. Il en résulte que $M$ est plat sur $R$.
En effet, en prenant pour $M'$ un module libre se surjectant sur $M$,
on obtient $\text{Tor}_1^R(M, N) = 0$ pour tout module $N$,
et on applique le lemme \ref{algebra-lemma-characterize-flat}
du chapitre Algèbre. Ceci démontre (c).

\medskip\noindent
Pour déduire (d) de (c), remarquons que si $N \in \text{Mod}_R$
et si $M \otimes_R N$ est nul, alors
$M \otimes_R S \otimes_S (N \otimes_R S) \cong (M \otimes_R N) \otimes_R 
S$ est nul. Ainsi, $N \otimes_R S$ est nul, donc $N$ est nul.

\medskip\noindent
Pour déduire (e), il suffit maintenant de rappeler que $M$
est projectif de type fini si et seulement s'il est de présentation finie
et plat. Voir le lemme \ref{algebra-lemma-finite-projective}
du chapitre Algèbre.
\end{proof}

\noindent
Il existe une variante pour les $R$-algèbres.

\begin{theorem}
\label{descent-theorem-descend-algebra-properties}
Si $A \otimes_R S$ possède l'une des propriétés suivantes comme $S$-algèbre :
\begin{enumerate}
\item[(a)]
être de type fini ;
\item[(b)]
être de présentation finie ;
\item[(c)]
être formellement non ramifiée ;
\item[(d)]
être non ramifiée ;
\item[(e)]
être \'etale ;
\end{enumerate}
alors $A$ possède la même propriété comme $R$-algèbre,
et la réciproque est bien entendu vraie.
\end{theorem}

\begin{proof}
Pour démontrer (a), choisissons une famille finie $\{x_i\}$ de générateurs
de $A \otimes_R S$ sur $S$.
Écrivons chaque $x_i$ sous la forme $\sum_j y_{ij} \otimes s_{ij}$,
avec $y_{ij} \in A$ et $s_{ij} \in S$.
Soit $F$ la $R$-algèbre de polynômes en les variables $e_{ij}$,
et soit $F \to M$ le morphisme de $R$-algèbres qui envoie
$e_{ij}$ sur $y_{ij}$. Alors $F \otimes_R S\to A \otimes_R S$ est surjectif,
donc $\Coker(F \to A) \otimes_R S$ est nul, d'où $\Coker(F \to A)$ est nul.
Ceci démontre (a).

\medskip\noindent
Pour (b), supposons que $A \otimes_R S$ soit une $S$-algèbre de présentation finie.
D'après (a), $A$ est de type fini sur $R$.
Choisissons une surjection $R[x_1, \ldots, x_n] \to A$ de noyau $I$.
La suite $I \otimes_R S \to S[x_1, \ldots, x_n] \to A \otimes_R S \to 0$
est exacte. D'après le lemme
\ref{algebra-lemma-finite-presentation-independent} du chapitre Algèbre,
le noyau de $S[x_1, \ldots, x_n] \to A \otimes_R S$
est un idéal de type fini.
Il existe donc un nombre fini d'éléments $y_1, \ldots, y_t \in I$
tels que les images de $y_i \otimes 1$ dans $S[x_1, \ldots, x_n]$
engendrent le noyau de $S[x_1, \ldots, x_n] \to A \otimes_R S$.
Soit $I' \subset I$ l'idéal engendré par $y_1, \ldots, y_t$.
Alors $A' = R[x_1, \ldots, x_n]/I'$ est une $R$-algèbre de présentation finie,
munie d'un morphisme $A' \to A$ tel que
$A' \otimes_R S \to A \otimes_R S$ soit un isomorphisme.
Ainsi, $A' \cong A$, comme voulu.

\medskip\noindent
Pour démontrer (c), rappelons que $A$ est formellement non ramifiée sur $R$
si et seulement si le module des différentielles relatives $\Omega_{A/R}$
est nul ; voir le lemme \ref{algebra-lemma-characterize-formally-unramified}
du chapitre Algèbre ou \cite[Proposition~17.2.1]{EGA4}.
Comme $\Omega_{(A \otimes_R S)/S} = \Omega_{A/R} \otimes_R S$,
la nullité descend en vertu du théorème \ref{descent-theorem-descent}.

\medskip\noindent
Pour déduire (d) des cas précédents, rappelons que $A$ est non ramifiée
sur $R$ si et seulement si $A$ est formellement non ramifiée
et de type fini sur $R$ ; voir le lemme
\ref{algebra-lemma-formally-unramified-unramified} du chapitre Algèbre.

\medskip\noindent
Pour démontrer (e), rappelons que, d'après le lemme
\ref{algebra-lemma-etale-flat-unramified-finite-presentation} du chapitre Algèbre
ou \cite[Th\'eor\`eme~17.6.1]{EGA4}, l'algèbre $A$ est \'etale sur $R$
si et seulement si $A$ est plate, non ramifiée et de présentation finie sur $R$.
\end{proof}

\begin{remark}
\label{descent-remark-when-locally-split}
La situation serait plus simple si l'on disposait d'un morphisme d'anneaux
fidèlement plat $g: R \to T$ tel que $T \to S \otimes_R T$
possède une structure supplémentaire.
Par exemple, si l'on pouvait faire en sorte que $T \to S \otimes_R T$
soit scindé dans $\textit{Rings}$, toute propriété d'un module ou d'une algèbre
qui est stable par extension des scalaires et que l'on peut faire descendre le
long de morphismes fidèlement plats pourrait aussi être descendue le long de
morphismes universellement injectifs.
On pourrait naturellement chercher $g$ tel que $T$ soit non seulement
fidèlement plat, mais encore injectif dans $\text{Mod}_R$.
Cependant, même pour $R = \mathbf{Z}$, un tel morphisme ne peut exister.
\end{remark}
















\section{Descente fpqc des faisceaux quasi-cohérents}
\label{descent-section-fpqc-descent-quasi-coherent}

\noindent
La principale application de la descente plate pour les modules
est l'énoncé correspondant de descente des faisceaux quasi-cohérents
relativement aux recouvrements fpqc.

\begin{lemma}
\label{descent-lemma-standard-fpqc-covering}
Soit $S$ un schéma affine.
Soit $\mathcal{U} = \{f_i : U_i \to S\}_{i = 1, \ldots, n}$
un recouvrement fpqc standard de $S$ ; voir la définition
\ref{topologies-definition-standard-fpqc} du chapitre Topologies.
Toute donnée de descente sur des faisceaux quasi-cohérents
relativement à $\mathcal{U} = \{U_i \to S\}$ est effective.
De plus, le foncteur de la catégorie des $\mathcal{O}_S$-modules
quasi-cohérents vers la catégorie des données de descente
relativement à $\mathcal{U}$ est pleinement fidèle.
\end{lemma}

\begin{proof}
Il s'agit de la proposition \ref{descent-proposition-descent-module},
reformulée en termes de schémas.
Remarquons d'abord qu'une donnée de descente $\xi$ sur des faisceaux
quasi-cohérents relativement à $\mathcal{U}$ revient exactement
à une donnée de descente $\xi'$ sur des faisceaux quasi-cohérents
relativement au recouvrement
$\mathcal{U}' = \{\coprod_{i = 1, \ldots, n} U_i \to S\}$.
De plus, l'effectivité de $\xi$ équivaut à celle de $\xi'$.
On peut donc supposer $n = 1$, c'est-à-dire $\mathcal{U} = \{U \to S\}$,
où $U$ et $S$ sont affines.
Dans ce cas, les données de descente correspondent aux données
de descente sur les modules relativement au morphisme d'anneaux
$$
\Gamma(S, \mathcal{O})
\longrightarrow
\Gamma(U, \mathcal{O}).
$$
Comme $U \to S$ est surjectif et plat, ce morphisme d'anneaux
est fidèlement plat. La proposition \ref{descent-proposition-descent-module}
s'applique donc et donne le résultat.
\end{proof}

\begin{proposition}
\label{descent-proposition-fpqc-descent-quasi-coherent}
Soit $S$ un schéma.
Soit $\mathcal{U} = \{\varphi_i : U_i \to S\}$ un recouvrement fpqc ;
voir la définition \ref{topologies-definition-fpqc-covering} du chapitre Topologies.
Toute donnée de descente sur des faisceaux quasi-cohérents
relativement à $\mathcal{U} = \{U_i \to S\}$ est effective.
De plus, le foncteur de la catégorie des $\mathcal{O}_S$-modules
quasi-cohérents vers la catégorie des données de descente
relativement à $\mathcal{U}$ est pleinement fidèle.
\end{proposition}

\begin{proof}
Soit $S = \bigcup_{j \in J} V_j$ un recouvrement par des ouverts affines.
Pour $j, j' \in J$, notons $V_{jj'} = V_j \cap V_{j'}$
l'intersection, qui n'est pas nécessairement affine.
Pour un ouvert $V \subset S$, notons
$\mathcal{U}_V = \{V \times_S U_i \to V\}_{i \in I}$ ;
c'est un recouvrement fpqc d'après le lemme
\ref{topologies-lemma-fpqc} du chapitre Topologies.
Par définition d'un recouvrement fpqc, il existe, pour chaque $j \in J$,
un ensemble fini $K_j$, une application $\underline{i} : K_j \to I$
et des ouverts affines $U_{\underline{i}(k), k} \subset U_{\underline{i}(k)}$,
$k \in K_j$, tels que
$\mathcal{V}_j = \{U_{\underline{i}(k), k} \to V_j\}_{k \in K_j}$
soit un recouvrement fpqc standard de $V_j$.
De plus, $\mathcal{V}_j$ est un raffinement de $\mathcal{U}_{V_j}$.
On a le diagramme
$$
\xymatrix{
\mathcal{V}_j \ar[r] \ar@{~>}[d] &
\mathcal{U}_{V_j} \ar[r] \ar@{~>}[d] &
\mathcal{U} \ar@{~>}[d] \\
V_j \ar@{=}[r] & V_j \ar[r] & S
}
$$
où les flèches horizontales supérieures sont des morphismes de familles
de morphismes de but fixé ; voir la définition
\ref{sites-definition-morphism-coverings} du chapitre Sites.

\medskip\noindent
Pour démontrer la proposition, nous établissons successivement
la fidélité, la plénitude et la surjectivité essentielle
du foncteur des faisceaux quasi-cohérents vers les données de descente.

\medskip\noindent
Fidélité. Soient $\mathcal{F}$ et $\mathcal{G}$ des faisceaux quasi-cohérents
sur $S$, et soient $a, b : \mathcal{F} \to \mathcal{G}$
des morphismes de $\mathcal{O}_S$-modules.
Supposons que $\varphi_i^*(a) = \varphi_i^*(b)$ pour tout $i$.
Soit $s \in S$. On a $s = \varphi_i(u)$ pour certains $i \in I$
et $u \in U_i$.
Puisque $\mathcal{O}_{S, s} \to \mathcal{O}_{U_i, u}$ est plat,
donc fidèlement plat par le lemme \ref{algebra-lemma-local-flat-ff}
du chapitre Algèbre, on a
$a_s = b_s : \mathcal{F}_s \to \mathcal{G}_s$. Ainsi, $a = b$.

\medskip\noindent
Pleine fidélité. Soient $\mathcal{F}$ et $\mathcal{G}$ des faisceaux
quasi-cohérents sur $S$, et soient
$a_i : \varphi_i^*\mathcal{F} \to \varphi_i^*\mathcal{G}$
des morphismes de $\mathcal{O}_{U_i}$-modules tels que
$\text{pr}_0^*a_i = \text{pr}_1^*a_j$ sur $U_i \times_U U_j$.
Par image inverse, on obtient des morphismes
$$
a_k :
\varphi_{i(k)}^*\mathcal{F}|_{U_{\underline{i}(k), k}}
\longrightarrow
\varphi_{i(k)}^*\mathcal{G}|_{U_{\underline{i}(k), k}}
$$
pour $k \in K_j$, avec les notations précédentes.
Le lemme \ref{descent-lemma-refine-descent-datum} montre qu'ils définissent
un morphisme entre les données de descente canoniques sur $\mathcal{V}_j$.
Le lemme \ref{descent-lemma-standard-fpqc-covering} fournit donc
d'uniques morphismes $a_j : \mathcal{F}|_{V_j} \to \mathcal{G}|_{V_j}$
correspondants. Pour montrer que $a_j|_{V_{jj'}} = a_{j'}|_{V_{jj'}}$,
on remarque que $a_j$ et $a_{j'}$ correspondent tous deux à l'image inverse
du morphisme $(a_i)_{i \in I}$ de données de descente canoniques sur
tout recouvrement raffinant à la fois $\mathcal{V}_{j, V_{jj'}}$
et $\mathcal{V}_{j', V_{jj'}}$, puis on utilise la fidélité déjà démontrée.
On peut prendre, par exemple, le recouvrement
$\mathcal{V}_{jj'} =
\{V_k \times_S V_{k'} \to V_{jj'}\}_{k \in K_j, k' \in K_{j'}}$.

\medskip\noindent
Surjectivité essentielle. Soit $\xi = (\mathcal{F}_i, \varphi_{ii'})$
une donnée de descente sur des faisceaux quasi-cohérents relativement
au recouvrement $\mathcal{U}$.
Par image inverse, elle donne des données de descente $\xi_j$
sur des faisceaux quasi-cohérents relativement aux recouvrements
$\mathcal{V}_j$ des $V_j$.
Le lemme \ref{descent-lemma-standard-fpqc-covering} fournit des faisceaux
quasi-cohérents $\mathcal{F}_j$ sur $V_j$ dont la donnée de descente
canonique associée est isomorphe à $\xi_j$.
La pleine fidélité démontrée ci-dessus fournit des isomorphismes
$$
\phi_{jj'} :
\mathcal{F}_j|_{V_{jj'}}
\longrightarrow
\mathcal{F}_{j'}|_{V_{jj'}}
$$
correspondant à l'isomorphisme entre les images inverses
de $\xi_j$ et de $\xi_{j'}$ sur $\mathcal{V}_{jj'}$.
Pour vérifier que les applications $\phi_{jj'}$ satisfont
la condition de cocycle, on utilise la fidélité déjà démontrée
sur les intersections triples $V_{jj'j''}$.
Le lemme \ref{descent-lemma-zariski-descent-effective} montre alors
que les faisceaux $\mathcal{F}_j$ se recollent en un faisceau
quasi-cohérent $\mathcal{F}$, comme voulu.
Il reste à vérifier que la donnée de descente canonique relative à
$\mathcal{U}$ associée à $\mathcal{F}$ est isomorphe à la donnée de descente
initiale. Cette vérification est omise.
\end{proof}









\section{Descente galoisienne des faisceaux quasi-cohérents}
\label{descent-section-galois-descent}

\noindent
La descente galoisienne des faisceaux quasi-cohérents est un cas particulier
de leur descente fpqc. Dans cette section, nous expliquons comment
reformuler une donnée de descente galoisienne en une donnée de descente fpqc,
puis nous appliquons les résultats précédents.

\medskip\noindent
Soit $k'/k$ une extension de corps. Alors $\{\Spec(k') \to \Spec(k)\}$
est un recouvrement fpqc. Soit $X$ un schéma sur $k$.
Pour une $k$-algèbre $A$, posons $X_A = X \times_{\Spec(k)} \Spec(A)$.
Le lemme \ref{topologies-lemma-fpqc} du chapitre Topologies montre que
$\{X_{k'} \to X\}$ est un recouvrement fpqc. On a
$$
X_{k'} \times_X X_{k'} = X_{k' \otimes_k k'}
\quad\text{et}\quad
X_{k'} \times_X X_{k'} \times_X X_{k'} = X_{k' \otimes_k k' \otimes_k k'}
$$
Une donnée de descente sur des faisceaux quasi-cohérents relativement à
$\{X_{k'} \to X\}$ consiste donc en un faisceau quasi-cohérent $\mathcal{F}$
sur $X_{k'}$ et en un isomorphisme
$\varphi : \text{pr}_0^*\mathcal{F} \to \text{pr}_1^*\mathcal{F}$
sur $X_{k' \otimes_k k'}$ satisfaisant la condition de cocycle évidente
sur $X_{k' \otimes_k k' \otimes_k k'}$.
Précisons ce que cela signifie dans le cas d'une extension galoisienne.

\medskip\noindent
Soit $k'/k$ une extension galoisienne finie de groupe de Galois
$G = \text{Gal}(k'/k)$. On a des isomorphismes de $k$-algèbres
$$
k' \otimes_k k' \longrightarrow \prod\nolimits_{\sigma \in G} k',\quad
a \otimes b \longrightarrow \prod a\sigma(b)
$$
et
$$
k' \otimes_k k' \otimes_k k' \longrightarrow
\prod\nolimits_{(\sigma, \tau) \in G \times G} k',\quad
a \otimes b \otimes c \longrightarrow \prod a\sigma(b)\sigma(\tau(c))
$$
Nous choisissons ici $a\sigma(b)\sigma(\tau(c))$, et non
$a\sigma(b)\tau(c)$, afin de simplifier les formules qui suivent,
bien que ce choix ne soit pas indispensable. Pour $\sigma \in G$, notons
$$
f_\sigma = \text{id}_X \times \Spec(\sigma) :
X_{k'} \longrightarrow X_{k'}
$$
Puisque $\Spec(-)$ est contravariant, on a
$f_{\sigma \tau} = f_\tau \circ f_\sigma$, et non la composée dans l'autre ordre.
Le premier isomorphisme ci-dessus fournit une identification
$$
X_{k' \otimes_k k'} = \coprod\nolimits_{\sigma \in G} X_{k'}
$$
par laquelle $\text{pr}_0$ correspond à l'application
$$
\coprod\nolimits_{\sigma \in G} X_{k'}
\xrightarrow{\coprod \text{id}}
X_{k'}
$$
et $\text{pr}_1$ à l'application
$$
\coprod\nolimits_{\sigma \in G} X_{k'}
\xrightarrow{\coprod f_\sigma}
X_{k'}
$$
Une donnée de descente $\varphi$ sur $\mathcal{F}$ au-dessus de $X_{k'}$
correspond donc à une famille d'isomorphismes
$\varphi_\sigma : \mathcal{F} \to f_\sigma^*\mathcal{F}$.
Pour expliciter la condition de cocycle, utilisons l'identification
fournie par le second isomorphisme d'algèbres ci-dessus :
$$
X_{k' \otimes_k k' \otimes_k k'} =
\coprod\nolimits_{(\sigma, \tau) \in G \times G} X_{k'}.
$$
Par cette identification, $\text{pr}_{01}$ correspond à l'application
$$
\coprod\nolimits_{(\sigma, \tau) \in G \times G} X_{k'}
\longrightarrow
\coprod\nolimits_{\sigma \in G} X_{k'}
$$
qui envoie la composante d'indice $(\sigma, \tau)$ sur celle d'indice
$\sigma$ par le morphisme identité.
L'application $\text{pr}_{12}$ correspond à l'application
$$
\coprod\nolimits_{(\sigma, \tau) \in G \times G} X_{k'}
\longrightarrow
\coprod\nolimits_{\sigma \in G} X_{k'}
$$
qui envoie la composante d'indice $(\sigma, \tau)$ sur celle d'indice $\tau$
par le morphisme $f_\sigma$. Enfin, $\text{pr}_{02}$ correspond à l'application
$$
\coprod\nolimits_{(\sigma, \tau) \in G \times G} X_{k'}
\longrightarrow
\coprod\nolimits_{\sigma \in G} X_{k'}
$$
qui envoie la composante d'indice $(\sigma, \tau)$ sur celle d'indice
$\sigma\tau$ par le morphisme identité. La condition de cocycle
$$
\text{pr}_{02}^*\varphi = \text{pr}_{12}^*\varphi \circ \text{pr}_{01}^*\varphi
$$
se traduit donc par une condition pour chaque couple $(\sigma, \tau)$, à savoir
$$
\varphi_{\sigma\tau} = f_\sigma^*\varphi_\tau \circ \varphi_\sigma
$$
comme applications $\mathcal{F} \to f_{\sigma\tau}^*\mathcal{F}$.
(Les sources et les buts coïncident : par exemple, le but de $\varphi_\sigma$
est $f_\sigma^*\mathcal{F}$, et la source de
$f_\sigma^*\varphi_\tau$ est également $f_\sigma^*\mathcal{F}$.)

\begin{lemma}
\label{descent-lemma-galois-descent}
Soit $k'/k$ une extension galoisienne (finie) de groupe de Galois $G$.
Soit $X$ un schéma sur $k$. La catégorie des $\mathcal{O}_X$-modules
quasi-cohérents est équivalente à la catégorie des systèmes
$(\mathcal{F}, (\varphi_\sigma)_{\sigma \in G})$ où
\begin{enumerate}
\item $\mathcal{F}$ est un module quasi-cohérent sur $X_{k'}$ ;
\item $\varphi_\sigma : \mathcal{F} \to f_\sigma^*\mathcal{F}$
est un isomorphisme de modules ;
\item $\varphi_{\sigma\tau} = f_\sigma^*\varphi_\tau \circ \varphi_\sigma$
pour tous $\sigma, \tau \in G$.
\end{enumerate}
Ici, $f_\sigma = \text{id}_X \times \Spec(\sigma) : X_{k'} \to X_{k'}$.
\end{lemma}

\begin{proof}
Comme nous l'avons vu, une donnée
$(\mathcal{F}, (\varphi_\sigma)_{\sigma \in G})$ comme dans l'énoncé
revient à une donnée de descente relativement au recouvrement fpqc
$\{X_{k'} \to X\}$. Le lemme résulte donc de la proposition
\ref{descent-proposition-fpqc-descent-quasi-coherent}.
\end{proof}

\noindent
Voici une légère généralisation. Supposons que l'on dispose d'un morphisme
fini \'etale surjectif $X \to Y$, d'un groupe fini $G$ et d'un morphisme de groupes
$G^{opp} \to \text{Aut}_Y(X), \sigma \mapsto f_\sigma$
tel que l'application
$$
G \times X \longrightarrow X \times_Y X,\quad
(\sigma, x) \longmapsto (x, f_\sigma(x))
$$
soit un isomorphisme. Le même résultat s'applique.

\begin{lemma}
\label{descent-lemma-galois-descent-more-general}
Soient $X \to Y$, $G$ et $f_\sigma : X \to X$ comme ci-dessus.
La catégorie des $\mathcal{O}_Y$-modules quasi-cohérents est équivalente
à la catégorie des systèmes $(\mathcal{F}, (\varphi_\sigma)_{\sigma \in G})$ où
\begin{enumerate}
\item $\mathcal{F}$ est un $\mathcal{O}_X$-module quasi-cohérent ;
\item $\varphi_\sigma : \mathcal{F} \to f_\sigma^*\mathcal{F}$
est un isomorphisme de modules ;
\item $\varphi_{\sigma\tau} = f_\sigma^*\varphi_\tau \circ \varphi_\sigma$
pour tous $\sigma, \tau \in G$.
\end{enumerate}
\end{lemma}

\begin{proof}
Puisque $X \to Y$ est fini \'etale surjectif, $\{X \to Y\}$
est un recouvrement fpqc. Comme
$G \times X \to X \times_Y X$, $(\sigma, x) \mapsto (x, f_\sigma(x))$
est un isomorphisme, l'application
$G \times G \times X \to X \times_Y X \times_Y X$,
$(\sigma, \tau, x) \mapsto (x, f_\sigma(x), f_{\sigma\tau}(x))$
en est un également. Ces identifications montrent que la catégorie
des données de l'énoncé s'identifie à celle des données de descente
sur les faisceaux quasi-cohérents relativement au recouvrement
$\{x \to Y\}$. Le lemme résulte donc de la proposition
\ref{descent-proposition-fpqc-descent-quasi-coherent}.
\end{proof}









\section{Descente des propriétés de finitude des modules}
\label{descent-section-descent-finiteness}

\noindent
Dans cette section, nous montrons que certaines conditions de finitude
sur un module quasi-cohérent peuvent se vérifier sur les membres d'un
recouvrement.

\begin{lemma}
\label{descent-lemma-finite-type-descends}
Soit $X$ un schéma.
Soit $\mathcal{F}$ un $\mathcal{O}_X$-module quasi-cohérent.
Soit $\{f_i : X_i \to X\}_{i \in I}$ un recouvrement fpqc tel que
chaque $f_i^*\mathcal{F}$ soit un $\mathcal{O}_{X_i}$-module de type fini.
Alors $\mathcal{F}$ est un $\mathcal{O}_X$-module de type fini.
\end{lemma}

\begin{proof}
Démonstration omise. Pour le cas affine, voir
Algèbre, lemme \ref{algebra-lemma-descend-properties-modules}.
\end{proof}

\begin{lemma}
\label{descent-lemma-finite-type-descends-fppf}
Soit $f : (X, \mathcal{O}_X) \to (Y, \mathcal{O}_Y)$ un morphisme
d'espaces localement annelés. Soit $\mathcal{F}$ un faisceau de $\mathcal{O}_Y$-modules.
Si
\begin{enumerate}
\item l'application continue sous-jacente à $f$ est ouverte,
\item $f$ est surjectif et plat, et
\item $f^*\mathcal{F}$ est de type fini,
\end{enumerate}
alors $\mathcal{F}$ est de type fini.
\end{lemma}

\begin{proof}
Soit $y \in Y$ un point. Choisissons un point $x \in X$ d'image $y$.
Choisissons un ouvert $x \in U \subset X$ et des éléments $s_1, \ldots, s_n$
de $f^*\mathcal{F}(U)$ qui engendrent $f^*\mathcal{F}$ sur $U$.
Puisque $f^*\mathcal{F} =
f^{-1}\mathcal{F} \otimes_{f^{-1}\mathcal{O}_Y} \mathcal{O}_X$,
nous pouvons, quitte à restreindre $U$, supposer que $s_i = \sum t_{ij} \otimes a_{ij}$
avec $t_{ij} \in f^{-1}\mathcal{F}(U)$ et $a_{ij} \in \mathcal{O}_X(U)$.
En restreignant encore $U$, nous pouvons supposer que $t_{ij}$ provient
d'une section $s_{ij} \in \mathcal{F}(V)$ pour un ouvert $V \subset Y$
tel que $f(U) \subset V$. Soit $N$ le nombre de sections $s_{ij}$ et
considérons l'application
$$
\sigma = (s_{ij}) : \mathcal{O}_V^{\oplus N} \to \mathcal{F}|_V
$$
Le choix des sections montre que $f^*\sigma|_U$ est surjective.
Ainsi, pour tout $u \in U$, l'application
$$
\sigma_{f(u)} \otimes_{\mathcal{O}_{Y, f(u)}} \mathcal{O}_{X, u} :
\mathcal{O}_{X, u}^{\oplus N}
\longrightarrow
\mathcal{F}_{f(u)} \otimes_{\mathcal{O}_{Y, f(u)}} \mathcal{O}_{X, u}
$$
est surjective. Comme $f$ est plat, l'homomorphisme local d'anneaux
$\mathcal{O}_{Y, f(u)} \to \mathcal{O}_{X, u}$ est plat, donc
fidèlement plat (Algèbre, lemme \ref{algebra-lemma-local-flat-ff}).
Par conséquent, $\sigma_{f(u)}$ est surjective. Puisque $f$ est ouvert, $f(U)$ est
un voisinage ouvert de $y$, ce qui achève la démonstration.
\end{proof}

\begin{lemma}
\label{descent-lemma-finite-presentation-descends}
Soit $X$ un schéma.
Soit $\mathcal{F}$ un $\mathcal{O}_X$-module quasi-cohérent.
Soit $\{f_i : X_i \to X\}_{i \in I}$ un recouvrement fpqc tel que
chaque $f_i^*\mathcal{F}$ soit un $\mathcal{O}_{X_i}$-module de présentation
finie. Alors $\mathcal{F}$ est un $\mathcal{O}_X$-module
de présentation finie.
\end{lemma}

\begin{proof}
Démonstration omise. Pour le cas affine, voir
Algèbre, lemme \ref{algebra-lemma-descend-properties-modules}.
\end{proof}

\begin{lemma}
\label{descent-lemma-locally-generated-by-r-sections-descends}
Soit $X$ un schéma.
Soit $\mathcal{F}$ un $\mathcal{O}_X$-module quasi-cohérent.
Soit $\{f_i : X_i \to X\}_{i \in I}$ un recouvrement fpqc tel que
chaque $f_i^*\mathcal{F}$ soit localement engendré par $r$ sections comme
$\mathcal{O}_{X_i}$-module. Alors $\mathcal{F}$ est localement engendré par
$r$ sections comme $\mathcal{O}_X$-module.
\end{lemma}

\begin{proof}
D'après le lemme \ref{descent-lemma-finite-type-descends}, le module $\mathcal{F}$
est de type fini. Le lemme de Nakayama
(Algèbre, lemme \ref{algebra-lemma-NAK}) implique donc que $\mathcal{F}$
est engendré par $r$ sections au voisinage d'un point $x \in X$
si et seulement si $\dim_{\kappa(x)} \mathcal{F}_x \otimes \kappa(x) \leq r$.
Choisissons un $i$ et un point $x_i \in X_i$ d'image $x$. Alors
$\dim_{\kappa(x)} \mathcal{F}_x \otimes \kappa(x) = 
\dim_{\kappa(x_i)} (f_i^*\mathcal{F})_{x_i} \otimes \kappa(x_i)$,
et cette dimension est $\leq r$, puisque $f_i^*\mathcal{F}$ est localement engendré par $r$
sections.
\end{proof}

\begin{lemma}
\label{descent-lemma-flat-descends}
Soit $X$ un schéma.
Soit $\mathcal{F}$ un $\mathcal{O}_X$-module quasi-cohérent.
Soit $\{f_i : X_i \to X\}_{i \in I}$ un recouvrement fpqc tel que
chaque $f_i^*\mathcal{F}$ soit un $\mathcal{O}_{X_i}$-module plat.
Alors $\mathcal{F}$ est un $\mathcal{O}_X$-module plat.
\end{lemma}

\begin{proof}
Démonstration omise. Pour le cas affine, voir
Algèbre, lemme \ref{algebra-lemma-descend-properties-modules}.
\end{proof}

\begin{lemma}
\label{descent-lemma-finite-locally-free-descends}
Soit $X$ un schéma.
Soit $\mathcal{F}$ un $\mathcal{O}_X$-module quasi-cohérent.
Soit $\{f_i : X_i \to X\}_{i \in I}$ un recouvrement fpqc tel que
chaque $f_i^*\mathcal{F}$ soit un $\mathcal{O}_{X_i}$-module localement libre de type fini.
Alors $\mathcal{F}$ est un $\mathcal{O}_X$-module localement libre de type fini.
\end{lemma}

\begin{proof}
Cela résulte du fait qu'un faisceau quasi-cohérent est localement libre de type fini
si et seulement s'il est de présentation finie et plat, voir
Algèbre, lemme \ref{algebra-lemma-finite-projective}.
En effet, si chaque $f_i^*\mathcal{F}$ est plat et de présentation finie,
il en est de même de $\mathcal{F}$ d'après les
lemmes \ref{descent-lemma-flat-descends} et
\ref{descent-lemma-finite-presentation-descends}.
\end{proof}

\noindent
La définition d'un faisceau quasi-cohérent localement projectif figure dans
Propriétés, section \ref{properties-section-locally-projective}.

\begin{lemma}
\label{descent-lemma-locally-projective-descends}
Soit $X$ un schéma.
Soit $\mathcal{F}$ un $\mathcal{O}_X$-module quasi-cohérent.
Soit $\{f_i : X_i \to X\}_{i \in I}$ un recouvrement fpqc tel que
chaque $f_i^*\mathcal{F}$ soit un $\mathcal{O}_{X_i}$-module localement projectif.
Alors $\mathcal{F}$ est un $\mathcal{O}_X$-module localement projectif.
\end{lemma}

\begin{proof}
Démonstration omise. Pour les recouvrements de Zariski, c'est
Propriétés, lemme \ref{properties-lemma-locally-projective}.
Pour le cas affine, c'est
Algèbre, théorème \ref{algebra-theorem-ffdescent-projectivity}.
\end{proof}

\begin{remark}
\label{descent-remark-locally-free-descends}
Être localement libre est une propriété des modules quasi-cohérents qui
ne descend pas pour la topologie fpqc. En effet, supposons que
$R$ soit un anneau et que $M$ soit un $R$-module projectif qui est
une somme directe dénombrable $M = \bigoplus L_n$ de modules localement
libres de rang 1, sans être lui-même localement libre, voir
Exemples, lemme \ref{examples-lemma-projective-not-locally-free}.
Alors $M$ devient libre après le changement de base fidèlement plat
$$
R \longrightarrow
\bigoplus\nolimits_{m \geq 1}
\bigoplus\nolimits_{(i_1, \ldots, i_m) \in \mathbf{Z}^{\oplus m}}
L_1^{\otimes i_1} \otimes_R \ldots \otimes_R L_m^{\otimes i_m}
$$
Mais nous ne savons pas ce qu'il en est pour les recouvrements fppf.
Autrement dit, nous ne connaissons pas la réponse à la question suivante :
supposons que $A \to B$ soit un homomorphisme d'anneaux fidèlement
plat et de présentation finie. Soit $M$ un $A$-module
tel que $M \otimes_A B$ soit libre. Est-ce que $M$ est un
$A$-module localement libre ? Lorsque $A$ est noethérien, la réponse
est affirmative. Cela résulte des résultats de \cite{Bass}. Mais, en général,
nous ne connaissons pas la réponse. Si vous la connaissez, ou si vous disposez
d'une référence, veuillez écrire à
\href{mailto:stacks.project@gmail.com}{stacks.project@gmail.com}.
\end{remark}

\noindent
Nous ajoutons ici deux résultats liés aux précédents, mais de nature
légèrement différente.

\begin{lemma}
\label{descent-lemma-finite-over-finite-module}
Soit $f : X \to Y$ un morphisme de schémas.
Soit $\mathcal{F}$ un $\mathcal{O}_X$-module quasi-cohérent.
Supposons que $f$ soit un morphisme fini.
Alors $\mathcal{F}$ est un $\mathcal{O}_X$-module de type fini
si et seulement si $f_*\mathcal{F}$ est un $\mathcal{O}_Y$-module de type
fini.
\end{lemma}

\begin{proof}
Comme $f$ est fini, il est affine. On se ramène donc au cas où
$f$ est le morphisme $\Spec(B) \to \Spec(A)$ défini
par un homomorphisme fini d'anneaux $A \to B$.
Alors $\mathcal{F} = \widetilde{M}$ est le faisceau de modules
associé au $B$-module $M$.
Notons que $M$ est de type fini comme $B$-module si et seulement si
$M$ est de type fini comme $A$-module, voir
Algèbre, lemme \ref{algebra-lemma-finite-module-over-finite-extension}.
En combinant cela avec
Propriétés, lemme \ref{properties-lemma-finite-type-module},
on obtient le résultat.
\end{proof}

\begin{lemma}
\label{descent-lemma-finite-finitely-presented-module}
Soit $f : X \to Y$ un morphisme de schémas.
Soit $\mathcal{F}$ un $\mathcal{O}_X$-module quasi-cohérent.
Supposons que $f$ soit fini et de présentation finie.
Alors $\mathcal{F}$ est un $\mathcal{O}_X$-module de présentation finie
si et seulement si $f_*\mathcal{F}$ est un $\mathcal{O}_Y$-module de présentation
finie.
\end{lemma}

\begin{proof}
Comme $f$ est fini, il est affine. On se ramène donc au cas où
$f$ est le morphisme $\Spec(B) \to \Spec(A)$ défini
par un homomorphisme d'anneaux fini et de présentation finie $A \to B$.
Alors $\mathcal{F} = \widetilde{M}$ est le faisceau de modules
associé au $B$-module $M$.
Notons que $M$ est de présentation finie comme $B$-module si et seulement si
$M$ est de présentation finie comme $A$-module, voir
Algèbre, lemme \ref{algebra-lemma-finite-finitely-presented-extension}.
En combinant cela avec
Propriétés, lemme \ref{properties-lemma-finite-presentation-module},
on obtient le résultat.
\end{proof}
















\section{Faisceaux quasi-cohérents et topologies, I}
\label{descent-section-quasi-coherent-sheaves}

\noindent
Les résultats de cette section donnent une équivalence naturelle entre
la catégorie des modules quasi-cohérents sur un schéma $S$ et la catégorie
des modules quasi-cohérents sur de nombreux sites associés à $S$
dans le chapitre consacré aux topologies.

\medskip\noindent
Soit $S$ un schéma.
Soit $\mathcal{F}$ un $\mathcal{O}_S$-module quasi-cohérent.
Considérons le foncteur
\begin{equation}
\label{descent-equation-quasi-coherent-presheaf}
(\Sch/S)^{opp} \longrightarrow \textit{Ab},
\quad
(f : T \to S) \longmapsto \Gamma(T, f^*\mathcal{F}).
\end{equation}

\begin{lemma}
\label{descent-lemma-sheaf-condition-holds}
Soit $S$ un schéma.
Soit $\mathcal{F}$ un $\mathcal{O}_S$-module quasi-cohérent.
Soit $\tau \in \{Zariski, \linebreak[0] \etale, \linebreak[0] smooth,
\linebreak[0] syntomic, \linebreak[0] fppf, \linebreak[0] fpqc\}$.
Le foncteur défini dans (\ref{descent-equation-quasi-coherent-presheaf})
satisfait la condition de faisceau relativement à tout recouvrement
pour la topologie $\tau$
$\{T_i \to T\}_{i \in I}$ d'un schéma $T$ sur $S$.
\end{lemma}

\begin{proof}
Pour $\tau \in \{Zariski, \linebreak[0] \etale, \linebreak[0] smooth,
\linebreak[0] syntomic, \linebreak[0] fppf\}$, un recouvrement pour la topologie $\tau$
est aussi un recouvrement fpqc, voir
Topologies, lemmes
\ref{topologies-lemma-zariski-etale},
\ref{topologies-lemma-zariski-etale-smooth},
\ref{topologies-lemma-zariski-etale-smooth-syntomic},
\ref{topologies-lemma-zariski-etale-smooth-syntomic-fppf} et
\ref{topologies-lemma-zariski-etale-smooth-syntomic-fppf-fpqc}.
Il suffit donc de démontrer le théorème
pour un recouvrement fpqc. Supposons que $\{f_i : T_i \to T\}_{i \in I}$
soit un recouvrement fpqc, avec $f : T \to S$ donné. Supposons
donnée une famille de sections $s_i \in \Gamma(T_i , f_i^*f^*\mathcal{F})$
telle que $s_i|_{T_i \times_T T_j} = s_j|_{T_i \times_T T_j}$.
Il faut trouver la section correspondante $s \in \Gamma(T, f^*\mathcal{F})$.
On peut interpréter les $s_i$ comme une famille de morphismes
$\varphi_i : f_i^*\mathcal{O}_T = \mathcal{O}_{T_i} \to f_i^*f^*\mathcal{F}$
compatibles avec les données de descente canoniques associées aux
faisceaux quasi-cohérents $\mathcal{O}_T$ et $f^*\mathcal{F}$ sur $T$.
D'après la proposition \ref{descent-proposition-fpqc-descent-quasi-coherent},
on peut donc les faire descendre, de manière unique,
en une application $\mathcal{O}_T \to f^*\mathcal{F}$, qui fournit
la section $s$ cherchée.
\end{proof}

\noindent
On peut en particulier donner la définition suivante.

\begin{definition}
\label{descent-definition-structure-sheaf}
Soit $\tau \in \{Zariski, \linebreak[0] \etale, \linebreak[0]
smooth, \linebreak[0] syntomic, \linebreak[0] fppf\}$.
Soit $S$ un schéma.
Soit $\Sch_\tau$ un grand site contenant $S$.
Soit $\mathcal{F}$ un $\mathcal{O}_S$-module quasi-cohérent.
\begin{enumerate}
\item Le {\it faisceau structural du grand site $(\Sch/S)_\tau$}
est le faisceau d'anneaux $T/S \mapsto \Gamma(T, \mathcal{O}_T)$, noté
$\mathcal{O}$ ou $\mathcal{O}_S$.
\item Si $\tau = Zariski$ ou $\tau = \etale$, le
{\it faisceau structural du petit site} $S_{Zar}$ ou $S_\etale$
est le faisceau d'anneaux $T/S \mapsto \Gamma(T, \mathcal{O}_T)$,
noté $\mathcal{O}$ ou $\mathcal{O}_S$.
\item Le {\it faisceau de $\mathcal{O}$-modules associé à
$\mathcal{F}$} sur le grand site $(\Sch/S)_\tau$
est le faisceau de $\mathcal{O}$-modules
$(f : T \to S) \mapsto \Gamma(T, f^*\mathcal{F})$,
noté $\mathcal{F}^a$ (et souvent simplement $\mathcal{F}$).
\item Si $\tau = Zariski$ ou $\tau = \etale$, le
{\it faisceau de $\mathcal{O}$-modules associé à $\mathcal{F}$}
sur le petit site $S_{Zar}$ ou $S_\etale$ est le faisceau de
$\mathcal{O}$-modules $(f : T \to S) \mapsto \Gamma(T, f^*\mathcal{F})$,
noté $\mathcal{F}^a$ (et souvent simplement $\mathcal{F}$).
\end{enumerate}
\end{definition}

\noindent
Nous employons donc la même notation $\mathcal{F}^a$ dans chaque cas.
Cela ne crée pas de confusion, puisque, par définition, la règle
qui définit le faisceau $\mathcal{F}^a$ est indépendante du site
considéré.

\begin{remark}
\label{descent-remark-Zariski-site-space}
Dans Topologies, lemme \ref{topologies-lemma-Zariski-usual},
nous avons vu que le petit site de Zariski d'un schéma $S$ est
équivalent à $S$ vu comme espace topologique, au sens où leurs
catégories de faisceaux sont naturellement équivalentes. Maintenant que $S_{Zar}$
est aussi muni d'un faisceau structural $\mathcal{O}$, les
faisceaux de modules sur le site annelé $(S_{Zar}, \mathcal{O})$
s'identifient aux faisceaux de modules sur l'espace annelé $(S, \mathcal{O}_S)$.
\end{remark}

\begin{remark}
\label{descent-remark-change-topologies-ringed}
Soit $f : T \to S$ un morphisme de schémas.
Chacun des morphismes de sites $f_{sites}$ énumérés dans
Topologies, section \ref{topologies-section-change-topologies},
devient un morphisme de sites annelés. En effet, chacun de ces morphismes de sites
$f_{sites} : (\Sch/T)_\tau \to (\Sch/S)_{\tau'}$, ou
$f_{sites} : (\Sch/S)_\tau \to S_{\tau'}$, est défini par le foncteur
continu $S'/S \mapsto T \times_S S'/S$. Ainsi, pour $S'/S$ donné, on prend pour
$$
f_{sites}^\sharp :
\mathcal{O}(S'/S)
\longrightarrow
f_{sites, *}\mathcal{O}(S'/S) =
\mathcal{O}(T \times_S S'/T)
$$
l'homomorphisme usuel
$\text{pr}_{S'}^\sharp : \mathcal{O}(S') \to \mathcal{O}(T \times_S S')$.
De même, le morphisme
$i_f : \Sh(T_\tau) \to \Sh((\Sch/S)_\tau)$,
pour $\tau \in \{Zar, \etale\}$, voir
Topologies, lemmes \ref{topologies-lemma-put-in-T} et
\ref{topologies-lemma-put-in-T-etale},
devient un morphisme de topos annelés, puisque $i_f^{-1}\mathcal{O} = \mathcal{O}$.
Voici quelques cas particuliers :
\begin{enumerate}
\item Le morphisme de grands sites
$f_{big} : (\Sch/X)_{fppf} \to (\Sch/Y)_{fppf}$
devient un morphisme de sites annelés
$$
(f_{big}, f_{big}^\sharp) :
((\Sch/X)_{fppf}, \mathcal{O}_X)
\longrightarrow
((\Sch/Y)_{fppf}, \mathcal{O}_Y)
$$
au sens de Modules sur les sites, définition \ref{sites-modules-definition-ringed-site}.
Il en va de même pour les grands sites syntomique, lisse, \'etale et de Zariski.
\item Le morphisme de petits sites
$f_{small} : X_\etale \to Y_\etale$
devient un morphisme de sites annelés
$$
(f_{small}, f_{small}^\sharp) :
(X_\etale, \mathcal{O}_X)
\longrightarrow
(Y_\etale, \mathcal{O}_Y)
$$
au sens de Modules sur les sites, définition \ref{sites-modules-definition-ringed-site}.
Il en va de même pour le petit site de Zariski.
\end{enumerate}
\end{remark}

\noindent
Soit $S$ un schéma. Étant donné un $\mathcal{O}$-module sur, par exemple,
$(\Sch/S)_{Zar}$, son image inverse sur, par exemple, $(\Sch/S)_{fppf}$
est simplement le faisceau associé pour la topologie fppf. Pour comparer
les grands et les petits sites, on dispose du résultat suivant.

\begin{lemma}
\label{descent-lemma-compare-sites}
Soit $S$ un schéma. Notons
$$
\begin{matrix}
\text{id}_{\tau, Zar} & : & (\Sch/S)_\tau \to S_{Zar}, &
\tau \in \{Zar, \etale, smooth, syntomic, fppf\} \\
\text{id}_{\tau, \etale} & : &
(\Sch/S)_\tau \to S_\etale, &
\tau \in \{\etale, smooth, syntomic, fppf\} \\
\text{id}_{small, \etale, Zar} & : & S_\etale \to S_{Zar},
\end{matrix}
$$
les morphismes de sites annelés de la
remarque \ref{descent-remark-change-topologies-ringed}.
Soit $\mathcal{F}$ un faisceau de $\mathcal{O}_S$-modules,
que l'on considère comme un faisceau de $\mathcal{O}$-modules sur $S_{Zar}$. Alors
\begin{enumerate}
\item $(\text{id}_{\tau, Zar})^*\mathcal{F}$ est le faisceau associé pour la topologie $\tau$
au faisceau de Zariski
$$
(f : T \to S) \longmapsto \Gamma(T, f^*\mathcal{F})
$$
sur $(\Sch/S)_\tau$, et
\item $(\text{id}_{small, \etale, Zar})^*\mathcal{F}$ est le
faisceau associé pour la topologie \'etale au faisceau de Zariski
$$
(f : T \to S) \longmapsto \Gamma(T, f^*\mathcal{F})
$$
sur $S_\etale$.
\end{enumerate}
Soit $\mathcal{G}$ un faisceau de $\mathcal{O}$-modules
sur $S_\etale$. Alors
\begin{enumerate}
\item[(3)] $(\text{id}_{\tau, \etale})^*\mathcal{G}$ est le
faisceau associé pour la topologie $\tau$ au faisceau pour la topologie \'etale
$$
(f : T \to S) \longmapsto \Gamma(T, f_{small}^*\mathcal{G})
$$
où $f_{small} : T_\etale \to S_\etale$
est le morphisme de petits sites annelés pour la topologie \'etale de la
remarque \ref{descent-remark-change-topologies-ringed}.
\end{enumerate}
\end{lemma}

\begin{proof}
Démonstration de (1). Remarquons d'abord que le résultat est vrai pour $\tau = Zar$,
car on dispose alors du morphisme de topos
$i_f : \Sh(T_{Zar}) \to \Sh((\Sch/S)_{Zar})$
tel que $\text{id}_{\tau, Zar} \circ i_f = f_{small}$ comme morphismes
$T_{Zar} \to S_{Zar}$, voir
Topologies, lemmes \ref{topologies-lemma-put-in-T} et
\ref{topologies-lemma-morphism-big-small}.
Par transitivité de l'image inverse (voir
Modules sur les sites,
lemme \ref{sites-modules-lemma-push-pull-composition-modules}),
on obtient
$i_f^*(\text{id}_{\tau, Zar})^*\mathcal{F} = f_{small}^*\mathcal{F}$,
comme souhaité. D'après la remarque précédant ce lemme,
$(\text{id}_{\tau, Zar})^*\mathcal{F}$ est donc le faisceau associé pour la topologie $\tau$
au préfaisceau $T \mapsto \Gamma(T, f^*\mathcal{F})$.

\medskip\noindent
La démonstration de (3) est identique à celle de (1), à ceci près qu'elle
utilise
Topologies, lemmes \ref{topologies-lemma-put-in-T-etale} et
\ref{topologies-lemma-morphism-big-small-etale}.
Nous omettons la démonstration de (2).
\end{proof}

\begin{remark}
\label{descent-remark-change-topologies-ringed-sites}
La remarque \ref{descent-remark-change-topologies-ringed}
et le
lemme \ref{descent-lemma-compare-sites}
ont les applications suivantes :
\begin{enumerate}
\item Soit $S$ un schéma.
La construction $\mathcal{F} \mapsto \mathcal{F}^a$
est l'image inverse par le morphisme de sites annelés
$\text{id}_{\tau, Zar} : ((\Sch/S)_\tau, \mathcal{O})
\to (S_{Zar}, \mathcal{O})$
ou par le morphisme
$\text{id}_{small, \etale, Zar} :
(S_\etale, \mathcal{O}) \to (S_{Zar}, \mathcal{O})$.
\item Soit $f : X \to Y$ un morphisme de schémas.
Pour chacun des morphismes $f_{sites}$ de sites annelés de la
remarque \ref{descent-remark-change-topologies-ringed},
on a
$$
(f^*\mathcal{F})^a = f_{sites}^*\mathcal{F}^a.
$$
Cela résulte de (1) et de la compatibilité de l'image inverse avec
la composition des morphismes de sites annelés, voir
Modules sur les sites,
lemme \ref{sites-modules-lemma-push-pull-composition-modules}.
\end{enumerate}
\end{remark}

\begin{lemma}
\label{descent-lemma-quasi-coherent-gives-quasi-coherent}
Soit $S$ un schéma.
Soit $\mathcal{F}$ un $\mathcal{O}_S$-module quasi-cohérent.
Soit $\tau \in \{Zariski, \linebreak[0] \etale, \linebreak[0]
smooth, \linebreak[0] syntomic, \linebreak[0] fppf\}$.
\begin{enumerate}
\item Le faisceau $\mathcal{F}^a$ est un $\mathcal{O}$-module
quasi-cohérent sur $(\Sch/S)_\tau$, au sens de
Modules sur les sites, définition \ref{sites-modules-definition-site-local}.
\item Si $\tau = Zariski$ ou $\tau = \etale$, alors le faisceau
$\mathcal{F}^a$ est un $\mathcal{O}$-module quasi-cohérent sur
$S_{Zar}$ ou $S_\etale$, au sens de
Modules sur les sites, définition \ref{sites-modules-definition-site-local}.
\end{enumerate}
\end{lemma}

\begin{proof}
Soit $\{S_i \to S\}$ un recouvrement de Zariski tel que l'on ait des
suites exactes
$$
\bigoplus\nolimits_{k \in K_i} \mathcal{O}_{S_i} \longrightarrow
\bigoplus\nolimits_{j \in J_i} \mathcal{O}_{S_i} \longrightarrow
\mathcal{F}|_{S_i} \longrightarrow 0
$$
pour certains ensembles d'indices $K_i$ et $J_i$. Cela est possible par
la définition d'un faisceau quasi-cohérent sur un espace annelé
(voir Modules, définition \ref{modules-definition-quasi-coherent}).

\medskip\noindent
Preuve de (1). Soit $\tau \in \{Zariski, \linebreak[0] fppf, \linebreak[0]
\etale, \linebreak[0] smooth, \linebreak[0] syntomic\}$.
Il est clair que $\mathcal{F}^a|_{(\Sch/S_i)_\tau}$ s'insère aussi dans
une suite exacte
$$
\bigoplus\nolimits_{k \in K_i} \mathcal{O}|_{(\Sch/S_i)_\tau}
\longrightarrow
\bigoplus\nolimits_{j \in J_i} \mathcal{O}|_{(\Sch/S_i)_\tau}
\longrightarrow
\mathcal{F}^a|_{(\Sch/S_i)_\tau} \longrightarrow 0
$$
Ainsi $\mathcal{F}^a$ est quasi-cohérent par Modules sur les sites,
lemme \ref{sites-modules-lemma-local-final-object}.

\medskip\noindent
Preuve de (2). Soit $\tau = \etale$. Il est clair que
$\mathcal{F}^a|_{(S_i)_\etale}$ s'insère aussi dans une suite exacte
$$
\bigoplus\nolimits_{k \in K_i} \mathcal{O}|_{(S_i)_\etale}
\longrightarrow
\bigoplus\nolimits_{j \in J_i} \mathcal{O}|_{(S_i)_\etale}
\longrightarrow
\mathcal{F}^a|_{(S_i)_\etale} \longrightarrow 0
$$
Ainsi $\mathcal{F}^a$ est quasi-cohérent par Modules sur les sites,
lemme \ref{sites-modules-lemma-local-final-object}.
Le cas $\tau = Zariski$ est analogue (il est même tautologique, puisque
les topos annelés correspondants coïncident).
\end{proof}

\begin{lemma}
\label{descent-lemma-fully-faithful-associated}
Soit $S$ un schéma.
Soit $\tau \in \{Zariski, \linebreak[0] \etale, \linebreak[0]
smooth, \linebreak[0] syntomic, \linebreak[0] fppf\}$.
Chacun des foncteurs $\mathcal{F} \mapsto \mathcal{F}^a$ de la
définition \ref{descent-definition-structure-sheaf}
$$
\QCoh(\mathcal{O}_S) \to \QCoh((\Sch/S)_\tau, \mathcal{O})
\quad\text{ou}\quad
\QCoh(\mathcal{O}_S) \to \QCoh(S_\tau, \mathcal{O})
$$
est pleinement fidèle.
\end{lemma}

\begin{proof}
Le lemme \ref{descent-lemma-quasi-coherent-gives-quasi-coherent} montre que l'on
obtient bien les foncteurs indiqués. On identifie les
$\mathcal{O}_S$-modules sur $S$ aux modules sur
$(S_{Zar}, \mathcal{O}_S)$. Le foncteur
$\mathcal{F} \mapsto \mathcal{F}^a$, défini sur les modules
quasi-cohérents $\mathcal{F}$, est l'image inverse par un morphisme $f$
de sites annelés; voir la remarque
\ref{descent-remark-change-topologies-ringed-sites}. Dans chaque cas, le foncteur
$f_*$ est la restriction suivant le foncteur d'inclusion
$S_{Zar} \to S_\tau$ ou $S_{Zar} \to (\Sch/S)_\tau$ (voir, dans
Topologies, section \ref{topologies-section-change-topologies}, la
construction de ces morphismes de sites). En combinant cela avec la
description de $f^*\mathcal{F} = \mathcal{F}^a$, on obtient
$f_*f^*\mathcal{F} = \mathcal{F}$ lorsque $\mathcal{F}$ est
quasi-cohérent. Par conséquent,
$$
\Hom_\mathcal{O}(\mathcal{F}^a, \mathcal{G}^a) =
\Hom_\mathcal{O}(f^*\mathcal{F}, f^*\mathcal{G}) =
\Hom_{\mathcal{O}_S}(\mathcal{F}, f_*f^*\mathcal{G}) =
\Hom_{\mathcal{O}_S}(\mathcal{F}, \mathcal{G})
$$
ce qu'il fallait démontrer.
\end{proof}

\begin{proposition}
\label{descent-proposition-equivalence-quasi-coherent}
Soit $S$ un schéma.
Soit $\tau \in \{Zariski, \linebreak[0] \etale, \linebreak[0]
smooth, \linebreak[0] syntomic, \linebreak[0] fppf\}$.
\begin{enumerate}
\item Le foncteur $\mathcal{F} \mapsto \mathcal{F}^a$ définit une
équivalence de catégories
$$
\QCoh(\mathcal{O}_S)
\longrightarrow
\QCoh((\Sch/S)_\tau, \mathcal{O})
$$
entre la catégorie des faisceaux quasi-cohérents sur $S$ et celle des
$\mathcal{O}$-modules quasi-cohérents sur le grand site de $S$ pour la
topologie $\tau$.
\item Si $\tau = Zariski$ ou $\tau = \etale$, le foncteur
$\mathcal{F} \mapsto \mathcal{F}^a$ définit une équivalence de catégories
$$
\QCoh(\mathcal{O}_S)
\longrightarrow
\QCoh(S_\tau, \mathcal{O})
$$
entre la catégorie des faisceaux quasi-cohérents sur $S$ et celle des
$\mathcal{O}$-modules quasi-cohérents sur le petit site de $S$ pour la
topologie $\tau$.
\end{enumerate}
\end{proposition}

\begin{proof}
Le lemme \ref{descent-lemma-quasi-coherent-gives-quasi-coherent} montre que le
foncteur est bien défini. Il est pleinement fidèle par le lemme
\ref{descent-lemma-fully-faithful-associated}. Pour achever la preuve, nous allons
montrer qu'un $\mathcal{O}$-module quasi-cohérent sur $(\Sch/S)_\tau$
fournit une donnée de descente pour des faisceaux quasi-cohérents
relativement à un recouvrement de $S$ pour la topologie $\tau$. Une fois
cette donnée
construite, la proposition \ref{descent-proposition-fpqc-descent-quasi-coherent}
donnera le faisceau quasi-cohérent correspondant sur $S$.

\medskip\noindent
Soit $\mathcal{G}$ un $\mathcal{O}$-module quasi-cohérent sur le grand
site de $S$ pour la topologie $\tau$. D'après Modules sur les sites,
définition \ref{sites-modules-definition-site-local}, il existe un
recouvrement $\{S_i \to S\}_{i \in I}$ de $S$ pour la topologie $\tau$,
tel que chacune des restrictions
$\mathcal{G}|_{(\Sch/S_i)_\tau}$ admette une présentation globale
$$
\bigoplus\nolimits_{k \in K_i} \mathcal{O}|_{(\Sch/S_i)_\tau}
\longrightarrow
\bigoplus\nolimits_{j \in J_i} \mathcal{O}|_{(\Sch/S_i)_\tau}
\longrightarrow
\mathcal{G}|_{(\Sch/S_i)_\tau} \longrightarrow 0
$$
pour certains ensembles d'indices $J_i$ et $K_i$. Nous affirmons qu'il
en résulte que $\mathcal{G}|_{(\Sch/S_i)_\tau}$ est de la forme
$\mathcal{F}_i^a$ pour un faisceau quasi-cohérent $\mathcal{F}_i$ sur
$S_i$. Cela est clair
pour les sommes directes
$\bigoplus\nolimits_{k \in K_i} \mathcal{O}|_{(\Sch/S_i)_\tau}$ et
$\bigoplus\nolimits_{j \in J_i} \mathcal{O}|_{(\Sch/S_i)_\tau}$.
Ainsi $\mathcal{G}|_{(\Sch/S_i)_\tau}$ est le conoyau d'un morphisme
$\varphi : \mathcal{K}_i^a \to \mathcal{L}_i^a$, où
$\mathcal{K}_i$ et $\mathcal{L}_i$ sont des faisceaux quasi-cohérents
sur $S_i$. Par la pleine fidélité de $(\ )^a$, on a
$\varphi = \phi^a$ pour un morphisme de faisceaux quasi-cohérents
$\phi : \mathcal{K}_i \to \mathcal{L}_i$ sur $S_i$. On a donc bien
$\mathcal{G}|_{(\Sch/S_i)_\tau} \cong \Coker(\phi)^a$.

\medskip\noindent
Comme $\mathcal{G}$ est défini sur toute la catégorie $(\Sch/S)_\tau$,
on a
$$
(\text{pr}_0^*\mathcal{F}_i)^a
\cong
\mathcal{G}|_{(\Sch/(S_i \times_S S_j))_\tau}
\cong
(\text{pr}_1^*\mathcal{F})^a
$$
comme $\mathcal{O}$-modules sur $(\Sch/(S_i \times_S S_j))_\tau$.
La pleine fidélité donne alors des isomorphismes canoniques
$$
\phi_{ij} :
\text{pr}_0^*\mathcal{F}_i
\longrightarrow
\text{pr}_1^*\mathcal{F}_j
$$
de modules quasi-cohérents sur $S_i \times_S S_j$. Nous omettons la
vérification de la condition de cocycle. Cette condition étant satisfaite,
l'effectivité de la descente des faisceaux quasi-cohérents pour le
recouvrement $\{S_i \to S\}$ (proposition
\ref{descent-proposition-fpqc-descent-quasi-coherent}) fournit un faisceau
quasi-cohérent $\mathcal{F}$ sur $S$, muni d'isomorphismes
$\mathcal{F}|_{S_i} \cong \mathcal{F}_i$ compatibles avec la donnée de
descente. Autrement dit, on dispose d'isomorphismes de
$\mathcal{O}$-modules
$$
\phi_i :
\mathcal{F}^a|_{(\Sch/S_i)_\tau}
\longrightarrow
\mathcal{G}|_{(\Sch/S_i)_\tau}
$$
qui coïncident sur $S_i \times_S S_j$. Comme
$\SheafHom_\mathcal{O}(\mathcal{F}^a, \mathcal{G})$ est un faisceau
(Modules sur les sites, lemme \ref{sites-modules-lemma-internal-hom}),
ces isomorphismes se recollent en un morphisme de $\mathcal{O}$-modules
$\mathcal{F}^a \to \mathcal{G}$ qui redonne les isomorphismes $\phi_i$
ci-dessus. Ce morphisme est un isomorphisme, ce qui
achève la preuve.

\medskip\noindent
Le cas des sites $S_\etale$ et $S_{Zar}$ se démontre exactement de la
même manière.
\end{proof}

\begin{lemma}
\label{descent-lemma-equivalence-quasi-coherent-properties}
Soit $S$ un schéma.
Soit $\tau \in \{Zariski, \linebreak[0] \etale, \linebreak[0]
smooth, \linebreak[0] syntomic, \linebreak[0] fppf\}$.
Soit $\mathcal{P}$ l'une des propriétés de modules\footnote{La liste est
la suivante : libre, libre de rang fini, engendré par ses sections globales,
engendré par $r$ sections globales, engendré par un nombre fini de sections
globales, admettant une présentation globale, admettant une présentation
globale finie, localement libre, localement libre de rang fini, localement
engendré par des sections, localement engendré par $r$ sections, de type
fini, de présentation finie, cohérent ou plat.} définies dans Modules sur
les sites, définitions \ref{sites-modules-definition-global},
\ref{sites-modules-definition-site-local} et
\ref{sites-modules-definition-flat}. Les équivalences de catégories
$$
\QCoh(\mathcal{O}_S)
\longrightarrow
\QCoh((\Sch/S)_\tau, \mathcal{O})
\quad\text{et}\quad
\QCoh(\mathcal{O}_S)
\longrightarrow
\QCoh(S_\tau, \mathcal{O})
$$
définies par $\mathcal{F} \mapsto \mathcal{F}^a$ dans la proposition
\ref{descent-proposition-equivalence-quasi-coherent} vérifient
$$
\mathcal{F}\text{ vérifie }\mathcal{P}
\Leftrightarrow
\mathcal{F}^a\text{ vérifie }\mathcal{P}\text{ comme }\mathcal{O}\text{-module}
$$
sauf, peut-être, lorsque $\mathcal{P}$ est ``localement libre'' ou
``cohérent''. Si $\mathcal{P}=$``cohérent'',
l'équivalence vaut pour
$\QCoh(\mathcal{O}_S) \to \QCoh(S_\tau, \mathcal{O})$ lorsque $S$ est
localement noethérien et que $\tau$ est Zariski ou \'etale.
\end{lemma}

\begin{proof}
C'est immédiat pour les propriétés globales, c'est-à-dire celles de
Modules sur les sites, définition \ref{sites-modules-definition-global}.
Pour les propriétés locales, le lemme
\ref{sites-modules-lemma-local-final-object} de Modules sur les sites
permet de traduire ``$\mathcal{F}^a$ vérifie $\mathcal{P}$'' en une
propriété portant sur les membres d'un recouvrement de $X$. Le résultat
découle alors des lemmes \ref{descent-lemma-finite-type-descends},
\ref{descent-lemma-finite-presentation-descends},
\ref{descent-lemma-locally-generated-by-r-sections-descends},
\ref{descent-lemma-flat-descends} et
\ref{descent-lemma-finite-locally-free-descends}. Pour un module quasi-cohérent,
être cohérent revient à être de type fini sur un schéma localement
noethérien (voir Cohomologie des schémas, lemme
\ref{coherent-lemma-coherent-Noetherian}); on se ramène donc au cas des
modules de type fini (détails omis).
\end{proof}










\section{Cohomologie des modules quasi-cohérents et topologies}
\label{descent-section-quasi-coherent-cohomology}

\noindent
Dans cette section, nous montrons que la cohomologie des modules
quasi-cohérents est indépendante du choix de la topologie.

\begin{lemma}
\label{descent-lemma-standard-covering-Cech}
Soit $S$ un schéma. Supposons que
\begin{enumerate}
\item[(a)] $\tau \in \{Zariski, \linebreak[0] fppf, \linebreak[0]
\etale, \linebreak[0] smooth, \linebreak[0] syntomic\}$ et
$\mathcal{C} = (\Sch/S)_\tau$, ou que
\item[(b)] $\tau = \etale$ et $\mathcal{C} = S_\etale$, ou que
\item[(c)] $\tau = Zariski$ et $\mathcal{C} = S_{Zar}$.
\end{enumerate}
Soit $\mathcal{F}$ un faisceau abélien sur $\mathcal{C}$.
Soit $U \in \Ob(\mathcal{C})$ affine.
Soit $\mathcal{U} = \{U_i \to U\}_{i = 1, \ldots, n}$ un recouvrement
affine standard pour la topologie $\tau$ dans $\mathcal{C}$. Alors
\begin{enumerate}
\item $\mathcal{V} = \{\coprod_{i = 1, \ldots, n} U_i \to U\}$ est un
recouvrement de $U$ pour la topologie $\tau$,
\item $\mathcal{U}$ est un raffinement de $\mathcal{V}$, et
\item l'application induite entre les complexes de {\v C}ech
(Cohomologie sur les sites,
équation (\ref{sites-cohomology-equation-map-cech-complexes}))
$$
\check{\mathcal{C}}^\bullet(\mathcal{V}, \mathcal{F})
\longrightarrow
\check{\mathcal{C}}^\bullet(\mathcal{U}, \mathcal{F})
$$
est un isomorphisme de complexes.
\end{enumerate}
\end{lemma}

\begin{proof}
Cela résulte de l'égalité
$$
(\coprod\nolimits_{i_0 = 1, \ldots, n} U_{i_0}) \times_U
\ldots \times_U
(\coprod\nolimits_{i_p = 1, \ldots, n} U_{i_p})
=
\coprod\nolimits_{i_0, \ldots, i_p \in \{1, \ldots, n\}}
U_{i_0} \times_U \ldots \times_U U_{i_p}
$$
et du fait que
$\mathcal{F}(\coprod_a V_a) = \prod_a \mathcal{F}(V_a)$, puisque les
réunions disjointes sont des recouvrements pour la topologie $\tau$.
\end{proof}

\begin{lemma}
\label{descent-lemma-standard-covering-Cech-quasi-coherent}
Soit $S$ un schéma. Soit $\mathcal{F}$ un faisceau quasi-cohérent sur
$S$. Soient $\tau$, $\mathcal{C}$, $U$ et $\mathcal{U}$ comme dans le
lemme \ref{descent-lemma-standard-covering-Cech}. Il existe alors un isomorphisme
de complexes
$$
\check{\mathcal{C}}^\bullet(\mathcal{U}, \mathcal{F}^a)
\cong
s((A/R)_\bullet \otimes_R M)
$$
(voir la section \ref{descent-section-descent-modules}), où
$R = \Gamma(U, \mathcal{O}_U)$, $M = \Gamma(U, \mathcal{F}^a)$ et
$R \to A$ est un homomorphisme d'anneaux fidèlement plat. En particulier,
$$
\check{H}^p(\mathcal{U}, \mathcal{F}^a) = 0
$$
pour tout $p \geq 1$.
\end{lemma}

\begin{proof}
Le lemme \ref{descent-lemma-standard-covering-Cech} montre que
$\check{\mathcal{C}}^\bullet(\mathcal{U}, \mathcal{F}^a)$ est isomorphe
à $\check{\mathcal{C}}^\bullet(\mathcal{V}, \mathcal{F}^a)$, où
$\mathcal{V} = \{V \to U\}$ et
$V = \coprod_{i = 1, \ldots n} U_i$ est encore affine. Posons
$A = \Gamma(V, \mathcal{O}_V)$. Puisque $\{V \to U\}$ est un
recouvrement pour la topologie $\tau$, l'homomorphisme $R \to A$ est fidèlement plat.
D'autre part, la définition de $\mathcal{F}^a$ donne pour terme de degré
$p$ du complexe $\check{\mathcal{C}}^p(\mathcal{V}, \mathcal{F}^a)$ :
$$
\Gamma(V \times_U \ldots \times_U V, \mathcal{F}^a)
=
\Gamma(\Spec(A \otimes_R \ldots \otimes_R A), \mathcal{F}^a)
=
A \otimes_R \ldots \otimes_R A \otimes_R M
$$
Nous omettons la vérification que les applications du complexe de
{\v C}ech coïncident avec celles du complexe
$s((A/R)_\bullet \otimes_R M)$. L'annulation de la cohomologie est le
lemme \ref{descent-lemma-ff-exact}.
\end{proof}

\begin{proposition}
\label{descent-proposition-same-cohomology-quasi-coherent}
\begin{slogan}
La cohomologie des faisceaux quasi-cohérents est la même quelle que soit
la topologie choisie.
\end{slogan}
Soit $S$ un schéma. Soit $\mathcal{F}$ un faisceau quasi-cohérent sur
$S$. Soit $\tau \in \{Zariski, \linebreak[0] \etale, \linebreak[0]
smooth, \linebreak[0] syntomic, \linebreak[0] fppf\}$.
\begin{enumerate}
\item Il existe un isomorphisme canonique
$$
H^q(S, \mathcal{F}) = H^q((\Sch/S)_\tau, \mathcal{F}^a).
$$
\item Il existe des isomorphismes canoniques
$$
H^q(S, \mathcal{F}) =
H^q(S_{Zar}, \mathcal{F}^a) =
H^q(S_\etale, \mathcal{F}^a).
$$
\end{enumerate}
\end{proposition}

\begin{proof}
Pour $q = 0$, le résultat est immédiat d'après la définition de
$\mathcal{F}^a$. Soit $\mathcal{C} = (\Sch/S)_\tau$, ou
$\mathcal{C} = S_\etale$, ou $\mathcal{C} = S_{Zar}$.

\medskip\noindent
Nous allons appliquer le lemme
\ref{sites-cohomology-lemma-cech-vanish-collection} de Cohomologie sur les
sites avec $\mathcal{F} = \mathcal{F}^a$, en prenant pour
$\mathcal{B} \subset \Ob(\mathcal{C})$ l'ensemble des schémas affines de
$\mathcal{C}$, et pour $\text{Cov} \subset \text{Cov}_\mathcal{C}$
l'ensemble des recouvrements affines standards pour la topologie $\tau$.
L'hypothèse (3) de ce lemme est satisfaite par le lemme
\ref{descent-lemma-standard-covering-Cech-quasi-coherent}. On en conclut que
$H^p(U, \mathcal{F}^a) = 0$ pour tout objet affine $U$ de $\mathcal{C}$.

\medskip\noindent
Soit maintenant $U \in \Ob(\mathcal{C})$ un objet séparé. Notons
$f : U \to S$ le morphisme structural. Soit $U = \bigcup U_i$ un
recouvrement ouvert affine. On peut aussi le considérer comme le
recouvrement $\mathcal{U} = \{U_i \to U\}$ de $U$ pour la topologie $\tau$ dans
$\mathcal{C}$. Remarquons que
$U_{i_0} \times_U \ldots \times_U U_{i_p} =
U_{i_0} \cap \ldots \cap U_{i_p}$ est affine, puisque $U$ est supposé
séparé. D'après Cohomologie sur les sites, lemme
\ref{sites-cohomology-lemma-cech-spectral-sequence-application}, et le
résultat précédent, on a
$$
H^p(U, \mathcal{F}^a) = \check{H}^p(\mathcal{U}, \mathcal{F}^a)
= H^p(U, f^*\mathcal{F})
$$
la dernière égalité provenant de Cohomologie des schémas, lemme
\ref{coherent-lemma-cech-cohomology-quasi-coherent}. En particulier, si
$S$ est séparé, on peut prendre $U = S$ et $f = \text{id}_S$, ce qui
démontre la proposition. Nous suggérons au lecteur de passer le reste de
la preuve (ou de le réécrire afin d'en rendre l'exposé plus clair).

\medskip\noindent
Choisissons une résolution injective
$\mathcal{F} \to \mathcal{I}^\bullet$ sur $S$. Choisissons une
résolution injective $\mathcal{F}^a \to \mathcal{J}^\bullet$
sur $\mathcal{C}$.
Notons $\mathcal{J}^n|_S$ la restriction de $\mathcal{J}^n$ aux ouverts
de $S$; c'est un faisceau sur l'espace topologique $S$, car les
recouvrements ouverts sont des recouvrements pour la topologie $\tau$. On obtient un
complexe
$$
0 \to \mathcal{F} \to \mathcal{J}^0|_S \to \mathcal{J}^1|_S \to \ldots
$$
qui est exact, puisque ses sections sur tout ouvert affine $U \subset S$
forment une suite exacte (par l'annulation de
$H^p(U, \mathcal{F}^a)$ pour $p > 0$ obtenue ci-dessus). Par Catégories
dérivées, lemme \ref{derived-lemma-morphisms-lift}, il existe donc un
morphisme de complexes
$\mathcal{J}^\bullet|_S \to \mathcal{I}^\bullet$ qui induit notamment
un morphisme
$$
R\Gamma(\mathcal{C}, \mathcal{F}^a)
=
\Gamma(S, \mathcal{J}^\bullet)
\longrightarrow
\Gamma(S, \mathcal{I}^\bullet)
=
R\Gamma(S, \mathcal{F}).
$$
En cohomologie, on obtient le morphisme
$H^n(\mathcal{C}, \mathcal{F}^a) \to H^n(S, \mathcal{F})$, dont il faut
montrer qu'elle est un isomorphisme. Soit
$\mathcal{U} : S = \bigcup U_i$ un recouvrement ouvert affine, que l'on
peut aussi voir comme un recouvrement pour la topologie $\tau$. Ce qui précède fournit un
morphisme de complexes doubles
$$
\check{\mathcal{C}}^\bullet(\mathcal{U}, \mathcal{J})
=
\check{\mathcal{C}}^\bullet(\mathcal{U}, \mathcal{J}|_S)
\longrightarrow
\check{\mathcal{C}}^\bullet(\mathcal{U}, \mathcal{I}).
$$
Il induit un morphisme de suites spectrales
$$
{}^\tau\! E_2^{p, q} = \check{H}^p(\mathcal{U}, \underline{H}^q(\mathcal{F}^a))
\longrightarrow
E_2^{p, q} = \check{H}^p(\mathcal{U}, \underline{H}^q(\mathcal{F}))
$$
La première suite spectrale converge vers
$H^{p + q}(\mathcal{C}, \mathcal{F})$, et la seconde vers
$H^{p + q}(S, \mathcal{F})$. D'autre part, nous avons vu que les
morphismes induits ${}^\tau\! E_2^{p, q} \to E_2^{p, q}$ sont
bijectives (toutes les intersections sont séparées, car ce sont des
ouverts d'affines). Les morphismes
$H^n(\mathcal{C}, \mathcal{F}^a) \to H^n(S, \mathcal{F})$ sont donc elles
aussi des isomorphismes, ce qui achève la preuve.
\end{proof}

\begin{proposition}
\label{descent-proposition-equivalence-quasi-coherent-functorial}
Soit $f : T \to S$ un morphisme de schémas.
\begin{enumerate}
\item Les équivalences de catégories de la
proposition \ref{descent-proposition-equivalence-quasi-coherent}
sont compatibles avec l'image inverse.
Plus précisément, on a $f^*(\mathcal{G}^a) = (f^*\mathcal{G})^a$
pour tout faisceau quasi-cohérent $\mathcal{G}$ sur $S$.
\item Les équivalences de catégories de la
proposition \ref{descent-proposition-equivalence-quasi-coherent}, partie (1),
ne sont en général {\bf pas} compatibles avec l'image directe.
\item Si $f$ est quasi-compact et quasi-séparé, et si
$\tau \in \{Zariski, \etale\}$, alors $f_*$ et $f_{small, *}$
préservent les faisceaux quasi-cohérents et le diagramme
$$
\xymatrix{
\QCoh(\mathcal{O}_T)
\ar[rr]_{f_*} \ar[d]_{\mathcal{F} \mapsto \mathcal{F}^a} & &
\QCoh(\mathcal{O}_S)
\ar[d]^{\mathcal{G} \mapsto \mathcal{G}^a} \\
\QCoh(T_\tau, \mathcal{O}) \ar[rr]^{f_{small, *}} & &
\QCoh(S_\tau, \mathcal{O})
}
$$
est commutatif, c'est-à-dire que
$f_{small, *}(\mathcal{F}^a) = (f_*\mathcal{F})^a$.
\end{enumerate}
\end{proposition}

\begin{proof}
La partie (1) résulte de la discussion de la
Remarque \ref{descent-remark-change-topologies-ringed-sites}.
La partie (2) constitue simplement un avertissement, que l'on peut expliquer
comme suit. Tout d'abord, l'énoncé ne peut pas être formulé précisément,
car $f_*$ ne transforme pas en général les faisceaux quasi-cohérents en
faisceaux quasi-cohérents. Même lorsque c'est le cas pour $f$ (et pour tout
changement de base de $f$), la compatibilité sur les grands sites signifierait
que la formation de $f_*\mathcal{F}$ commute à tout changement de base, ce
qui n'est pas vrai en général. Un exemple explicite est l'immersion ouverte
quasi-compacte
$j : X = \mathbf{A}^2_k \setminus \{0\} \to \mathbf{A}^2_k = Y$,
où $k$ est un corps. On a $j_*\mathcal{O}_X = \mathcal{O}_Y$,
mais, après le changement de base vers $\Spec(k)$ par le morphisme $0$,
on constate que l'image directe est nulle.

\medskip\noindent
Démontrons (3) dans le cas $\tau = \etale$. Le morphisme $f$, ainsi que tout
changement de base de $f$, envoie les faisceaux quasi-cohérents
en faisceaux quasi-cohérents ; voir Schémas, lemme
\ref{schemes-lemma-push-forward-quasi-coherent}.
L'égalité $f_{small, *}(\mathcal{F}^a) = (f_*\mathcal{F})^a$
signifie que, pour tout morphisme \'etale $g : U \to S$, on a
$\Gamma(U, g^*f_*\mathcal{F}) = \Gamma(U \times_S T, (g')^*\mathcal{F})$,
où $g' : U \times_S T \to T$ est la projection. Cela résulte du
Lemme \ref{coherent-lemma-flat-base-change-cohomology} de Cohomologie des schémas.
\end{proof}

\begin{lemma}
\label{descent-lemma-higher-direct-images-small-etale}
Soit $f : T \to S$ un morphisme quasi-compact et quasi-séparé de schémas.
Soit $\mathcal{F}$ un faisceau quasi-cohérent sur $T$. Pour la topologie
\'etale comme pour la topologie de Zariski, il existe des isomorphismes canoniques
$R^if_{small, *}(\mathcal{F}^a) = (R^if_*\mathcal{F})^a$.
\end{lemma}

\begin{proof}
Nous le démontrons pour la topologie \'etale ; nous omettons la démonstration
pour la topologie de Zariski. D'après Cohomologie des schémas, lemme
\ref{coherent-lemma-quasi-coherence-higher-direct-images},
les faisceaux $R^if_*\mathcal{F}$ sont quasi-cohérents, de sorte que l'énoncé
a un sens. Le faisceau $R^if_{small, *}\mathcal{F}^a$ est le faisceau associé
au préfaisceau
$$
U \longmapsto H^i(U \times_S T, \mathcal{F}^a)
$$
où $g : U \to S$ est un objet de $S_\etale$ ; voir Cohomologie sur les sites,
lemme \ref{sites-cohomology-lemma-higher-direct-images}.
Suivant nos conventions, le membre de droite est la cohomologie \'etale
de la restriction de $\mathcal{F}^a$ au site localisé
$T_\etale/U \times_S T$, lequel est égal à
$(U \times_S T)_\etale$. D'après la
proposition \ref{descent-proposition-same-cohomology-quasi-coherent},
ce préfaisceau est identique au préfaisceau
$$
U \longmapsto
H^i(U \times_S T, (g')^*\mathcal{F}),
$$
où $g' : U \times_S T \to T$ est la projection. Si $U$ est affine,
ce groupe est égal à $H^0(U, R^if'_*(g')^*\mathcal{F})$ ; voir
Cohomologie des schémas, lemme
\ref{coherent-lemma-quasi-coherence-higher-direct-images-application}.
D'après Cohomologie des schémas, lemme
\ref{coherent-lemma-flat-base-change-cohomology},
ce dernier groupe est égal à $H^0(U, g^*R^if_*\mathcal{F})$, qui est la valeur
de $(R^if_*\mathcal{F})^a$ en $U$.
Ainsi, les valeurs des faisceaux de modules
$R^if_{small, *}(\mathcal{F}^a)$ et $(R^if_*\mathcal{F})^a$
sur tout objet affine de $S_\etale$ sont canoniquement isomorphes ;
il s'ensuit que ces faisceaux sont canoniquement isomorphes.
\end{proof}




\section{Faisceaux quasi-cohérents et topologies, II}
\label{descent-section-quasi-coherent-sheaves-bis}

\noindent
Nous poursuivons la comparaison entre les modules quasi-cohérents sur un
schéma $S$ et les modules quasi-cohérents sur chacun des sites associés à
$S$ dans le chapitre sur les topologies.


\begin{lemma}
\label{descent-lemma-compare-etale-zariski-flat}
Dans le lemme \ref{descent-lemma-compare-sites}, le morphisme de sites annelés
$\text{id}_{small, \etale, Zar} : S_\etale \to S_{Zar}$ est plat.
\end{lemma}

\begin{proof}
Notons $\epsilon = \text{id}_{small, \etale, Zar}$, et notons
$\mathcal{O}_\etale$ et $\mathcal{O}_{Zar}$ les faisceaux structuraux sur
$S_\etale$ et $S_{Zar}$. Il faut montrer que $\mathcal{O}_\etale$ est un
$\epsilon^{-1}\mathcal{O}_{Zar}$-module plat. Rappelons qu'un morphisme
\'etale est ouvert ; voir Morphismes, lemme
\ref{morphisms-lemma-etale-open}. Il résulte de la construction de l'image
inverse des faisceaux que $\epsilon^{-1}\mathcal{O}_{Zar}$ est le faisceau
associé au préfaisceau $\mathcal{O}'$ sur $S_\etale$ qui associe à un
morphisme \'etale $f : V \to S$ l'anneau $\mathcal{O}_S(f(V))$. Si $V$ et
$U = f(V) \subset S$ sont tous deux affines, alors $V \to U$ est un
morphisme \'etale entre schémas affines et correspond donc à un homomorphisme
d'anneaux \'etale. Comme tout homomorphisme d'anneaux \'etale est plat,
l'homomorphisme
$\mathcal{O}_S(U) = \mathcal{O}'(V) \to
\mathcal{O}_\etale(V) = \mathcal{O}_V(V)$ est plat.
Enfin, pour tout morphisme \'etale $f : V \to S$, c'est-à-dire pour tout
objet de $S_\etale$, il existe un recouvrement ouvert affine
$V = \bigcup V_i$ tel que $f(V_i)$ soit un ouvert affine de $S$ pour tout
$i$\footnote{En effet, pour $y \in V$, choisissons un ouvert affine
$y \in V' \subset V$ tel que $f(V')$ soit contenu dans un ouvert affine
$U \subset S$. Choisissons ensuite un ouvert affine
$f(y) \in U' \subset f(V')$. Alors $V'' = f^{-1}(U') \subset V'$ est affine,
puisqu'il est égal à $U' \times_U V'$, et $f(V'') = U'$ est également affine.}.
Le résultat découle alors de Modules sur les sites, lemme
\ref{sites-modules-lemma-flatness-sheafification-refined}.
\end{proof}

\begin{lemma}
\label{descent-lemma-equivalence-quasi-coherent-limits}
Soit $S$ un schéma.
Soit $\tau \in \{Zariski, \linebreak[0] \etale,
\linebreak[0] smooth, \linebreak[0] syntomic, \linebreak[0] fppf\}$.
Les foncteurs
$$
\QCoh(\mathcal{O}_S)
\longrightarrow
\textit{Mod}((\Sch/S)_\tau, \mathcal{O})
\quad\text{et}\quad
\QCoh(\mathcal{O}_S)
\longrightarrow
\textit{Mod}(S_\tau, \mathcal{O})
$$
définis par la règle $\mathcal{F} \mapsto \mathcal{F}^a$ de la
proposition \ref{descent-proposition-equivalence-quasi-coherent}
ont les propriétés suivantes :
\begin{enumerate}
\item ils sont pleinement fidèles ;
\item ils commutent aux sommes directes ;
\item ils commutent aux colimites ;
\item ils sont exacts à droite ;
\item ils sont exacts comme foncteurs
$\QCoh(\mathcal{O}_S) \to \textit{Mod}(S_\etale, \mathcal{O})$ ;
\item ils ne sont en général {\bf pas} exacts comme foncteurs
$\QCoh(\mathcal{O}_S) \to
\textit{Mod}((\Sch/S)_\tau, \mathcal{O})$ ;
\item pour deux $\mathcal{O}_S$-modules quasi-cohérents
$\mathcal{F}$, $\mathcal{G}$, on a
$(\mathcal{F} \otimes_{\mathcal{O}_S} \mathcal{G})^a =
\mathcal{F}^a \otimes_\mathcal{O} \mathcal{G}^a$ ;
\item si $\tau = \etale$ ou $\tau = Zariski$, et si deux
$\mathcal{O}_S$-modules quasi-cohérents $\mathcal{F}$, $\mathcal{G}$ sont
donnés, tels que $\mathcal{F}$ soit de présentation finie, on a
$(\SheafHom_{\mathcal{O}_S}(\mathcal{F}, \mathcal{G}))^a =
\SheafHom_\mathcal{O}(\mathcal{F}^a, \mathcal{G}^a)$ dans
$\textit{Mod}(S_\tau, \mathcal{O})$ ;
\item pour deux $\mathcal{O}_S$-modules quasi-cohérents
$\mathcal{F}$, $\mathcal{G}$, on n'a en général {\bf pas}
$(\SheafHom_{\mathcal{O}_S}(\mathcal{F}, \mathcal{G}))^a =
\SheafHom_\mathcal{O}(\mathcal{F}^a, \mathcal{G}^a)$
dans $\textit{Mod}((\Sch/S)_\tau, \mathcal{O})$, même si
$\mathcal{F}$ est de présentation finie ;
\item pour toute suite exacte courte
$0 \to \mathcal{F}_1^a \to \mathcal{E} \to \mathcal{F}_2^a \to 0$
de $\mathcal{O}$-modules, le module $\mathcal{E}$ est
quasi-cohérent\footnote{Attention : cette formulation est trompeuse. Voir la
partie (6).}, c'est-à-dire que $\mathcal{E}$ appartient à l'image essentielle
du foncteur.
\end{enumerate}
\end{lemma}

\begin{proof}
La partie (1) a été établie dans la
proposition \ref{descent-proposition-equivalence-quasi-coherent}.

\medskip\noindent
Nous avons vu dans Schémas, section
\ref{schemes-section-quasi-coherent}, qu'une colimite de faisceaux
quasi-cohérents sur un schéma est encore un faisceau quasi-cohérent.
En outre, dans la remarque
\ref{descent-remark-change-topologies-ringed-sites}, nous avons vu que
$\mathcal{F} \mapsto \mathcal{F}^a$ est le foncteur d'image inverse d'un
morphisme de sites annelés ; il commute donc à toutes les colimites, par
Modules sur les sites, lemme
\ref{sites-modules-lemma-exactness-pushforward-pullback}.
Cela prouve (3), ainsi que son cas particulier (2).

\medskip\noindent
Cela montre aussi que le foncteur est exact à droite, c'est-à-dire qu'il
commute aux colimites finies, ce qui prouve (4).

\medskip\noindent
Le foncteur $\QCoh(\mathcal{O}_S) \to
\textit{Mod}(S_\etale, \mathcal{O})$,
$\mathcal{F} \mapsto \mathcal{F}^a$, est exact à gauche parce qu'un
morphisme \'etale est plat ; voir Morphismes, lemme
\ref{morphisms-lemma-etale-flat}. Cela prouve (5).

\medskip\noindent
Pour établir (6), supposons que $S = \Spec(\mathbf{Z})$.
Alors $2 : \mathcal{O}_S \to \mathcal{O}_S$ est injectif, mais le morphisme
associé de $\mathcal{O}$-modules sur $(\Sch/S)_\tau$ ne l'est pas, car
$2 : \mathbf{F}_2 \to \mathbf{F}_2$ n'est pas injectif et
$\Spec(\mathbf{F}_2)$ est un objet de $(\Sch/S)_\tau$.

\medskip\noindent
La partie (7) est vraie parce que, comme indiqué ci-dessus, le foncteur
$\mathcal{F} \mapsto \mathcal{F}^a$ est le foncteur d'image inverse d'un
morphisme de sites annelés, et qu'un tel foncteur commute aux produits
tensoriels par Modules sur les sites, lemme
\ref{sites-modules-lemma-tensor-product-pullback}.

\medskip\noindent
La partie (8) est évidente si $\tau = Zariski$, car la catégorie des
$\mathcal{O}$-modules sur $S_{Zar}$ est la même que celle des
$\mathcal{O}_S$-modules sur l'espace topologique $S$.
Si $\tau = \etale$, la partie (8) résulte du fait que, comme indiqué
ci-dessus, le foncteur $\mathcal{F} \mapsto \mathcal{F}^a$ est le foncteur
d'image inverse associé au morphisme plat de sites annelés
$(S_\etale, \mathcal{O}) \to (S_{Zar}, \mathcal{O}_S)$ ; voir le
lemme \ref{descent-lemma-compare-etale-zariski-flat}. L'image inverse par un
morphisme plat de sites annelés commute à la formation du Hom interne dont
le premier argument est un module de présentation finie, par Modules sur
les sites, lemme
\ref{sites-modules-lemma-pullback-internal-hom}.

\medskip\noindent
Pour établir (9), supposons que $S = \Spec(\mathbf{Z})$. Posons
$\mathcal{F} = \Coker(2 : \mathcal{O}_S \to \mathcal{O}_S)$
et $\mathcal{G} = \mathcal{O}_S$.
Alors $\mathcal{F}^a = \Coker(2 : \mathcal{O} \to \mathcal{O})$
et $\mathcal{G}^a = \mathcal{O}$. Par conséquent,
$\SheafHom_\mathcal{O}(\mathcal{F}^a, \mathcal{G}^a) = \mathcal{O}[2]$
est le sous-faisceau de $2$-torsion de $\mathcal{O}$, qui n'est pas nul ;
voir la preuve de (6). En revanche, le module
$\SheafHom_{\mathcal{O}_S}(\mathcal{F}, \mathcal{G})$ est nul.

\medskip\noindent
Preuve de (10). Soit
$0 \to \mathcal{F}_1^a \to \mathcal{E} \to \mathcal{F}_2^a \to 0$
une suite exacte courte de $\mathcal{O}$-modules, où $\mathcal{F}_1$ et
$\mathcal{F}_2$ sont quasi-cohérents sur $S$. Considérons sa restriction
$$
0 \to \mathcal{F}_1 \to \mathcal{E}|_{S_{Zar}} \to \mathcal{F}_2
$$
à $S_{Zar}$. D'après la
proposition \ref{descent-proposition-same-cohomology-quasi-coherent}, on a, pour
tout ouvert affine $U \subset S$,
$H^1(U, \mathcal{F}_1^a) = H^1(U, \mathcal{F}_1) = 0$.
La suite ci-dessus est donc également exacte à droite. D'après Schémas,
section \ref{schemes-section-quasi-coherent}, on en déduit que
$\mathcal{F} = \mathcal{E}|_{S_{Zar}}$ est quasi-cohérent.
On obtient ainsi un diagramme commutatif
$$
\xymatrix{
& \mathcal{F}_1^a \ar[r] \ar[d] &
\mathcal{F}^a \ar[r] \ar[d] &
\mathcal{F}_2^a \ar[r] \ar[d] & 0 \\
0 \ar[r] &
\mathcal{F}_1^a \ar[r] &
\mathcal{E} \ar[r] &
\mathcal{F}_2^a \ar[r] & 0
}
$$
Pour achever la preuve, il suffit de montrer que la ligne supérieure est
également exacte à droite. Notons de nouveau $U = \Spec(A) \subset S$ un
ouvert affine de $S$. Nous avons vu plus haut que
$0 \to \mathcal{F}_1(U) \to \mathcal{E}(U) \to \mathcal{F}_2(U) \to 0$
est exacte. Pour tout schéma affine $V/U$, avec $V = \Spec(B)$, le
morphisme $\mathcal{F}_1^a(V) \to \mathcal{E}(V)$ est injectif.
Par définition, on a
$\mathcal{F}_1^a(V) = \mathcal{F}_1(U) \otimes_A B$.
L'injection $\mathcal{F}_1^a(V) \to \mathcal{E}(V)$ se factorise en
$$
\mathcal{F}_1(U) \otimes_A B \to
\mathcal{E}(U) \otimes_A B \to \mathcal{E}(V)
$$
En considérant les $A$-algèbres $B$ de la forme $B = A \oplus M$, on voit
que $\mathcal{F}_1(U) \to \mathcal{E}(U)$ est universellement injectif
(voir Algèbre, Définition
\ref{algebra-definition-universally-injective}). Comme
$\mathcal{E}(U) = \mathcal{F}(U)$, on en déduit que
$\mathcal{F}_1 \to \mathcal{F}$ reste injectif après tout changement de
base, ou, de manière équivalente, que
$\mathcal{F}_1^a \to \mathcal{F}^a$ est injectif.
\end{proof}

\begin{lemma}
\label{descent-lemma-properties-quasi-coherent}
Soit $S$ un schéma. La catégorie $\QCoh(S_\etale, \mathcal{O})$ des modules
quasi-cohérents sur $S_\etale$ possède les propriétés suivantes :
\begin{enumerate}
\item Toute somme directe de faisceaux quasi-cohérents est quasi-cohérente.
\item Toute colimite de faisceaux quasi-cohérents est quasi-cohérente.
\item Le noyau et le conoyau d'un morphisme de faisceaux quasi-cohérents
sont quasi-cohérents.
\item Étant donnée une suite exacte courte de $\mathcal{O}$-modules
$0 \to \mathcal{F}_1 \to \mathcal{F}_2 \to \mathcal{F}_3 \to 0$,
si deux de ces trois modules sont quasi-cohérents, le troisième l'est aussi.
\item Le produit tensoriel de deux $\mathcal{O}$-modules quasi-cohérents
est quasi-cohérent.
\item Étant donnés deux $\mathcal{O}$-modules quasi-cohérents
$\mathcal{F}$, $\mathcal{G}$ tels que $\mathcal{F}$ soit de présentation
finie, le Hom interne
$\SheafHom_\mathcal{O}(\mathcal{F}, \mathcal{G})$
est quasi-cohérent.
\end{enumerate}
\end{lemma}

\begin{proof}
Les faits correspondants sont vrais pour les modules quasi-cohérents sur le
schéma $S$ ; voir Schémas, section
\ref{schemes-section-quasi-coherent}. La preuve consiste à utiliser le
lemme \ref{descent-lemma-equivalence-quasi-coherent-limits} pour les transporter à
$S_\etale$.

\medskip\noindent
Preuve de (1). Soient $\mathcal{F}_i$, $i \in I$, une famille d'objets de
$\QCoh(S_\etale, \mathcal{O})$. Écrivons
$\mathcal{F}_i = \mathcal{G}_i^a$, où les $\mathcal{G}_i$ sont des modules
quasi-cohérents sur $S$. Alors
$\bigoplus \mathcal{F}_i = (\bigoplus \mathcal{G}_i)^a$ par le lemme cité,
ce qui donne la conclusion.

\medskip\noindent
Preuve de (2). Soit
$\mathcal{I} \to \QCoh(S_\etale, \mathcal{O})$,
$i \mapsto \mathcal{F}_i$, un diagramme. Écrivons
$\mathcal{F}_i = \mathcal{G}_i^a$ ; on obtient ainsi un diagramme
$\mathcal{I} \to \QCoh(\mathcal{O}_S)$. Alors
$\colim \mathcal{F}_i = (\colim \mathcal{G}_i)^a$ par le lemme cité,
ce qui donne la conclusion.

\medskip\noindent
Preuve de (3). Soit $a : \mathcal{F} \to \mathcal{F}'$ une flèche de
$\QCoh(S_\etale, \mathcal{O})$. Écrivons $a = b^a$ pour un morphisme
$b : \mathcal{G} \to \mathcal{G}'$ de modules quasi-cohérents sur $S$.
Par le lemme cité, on a $\Ker(a) = \Ker(b)^a$ et
$\Coker(a) = \Coker(b)^a$, d'où le résultat.

\medskip\noindent
Preuve de (4). Cela résulte de (3), sauf dans le cas où l'on sait que
$\mathcal{F}_1$ et $\mathcal{F}_3$ sont quasi-cohérents. Dans ce cas,
écrivons $\mathcal{F}_1 = \mathcal{G}_1^a$ et
$\mathcal{F}_3 = \mathcal{G}_3^a$, avec $\mathcal{G}_i$ quasi-cohérent sur
$S$. On conclut par le lemme
\ref{descent-lemma-equivalence-quasi-coherent-limits}, partie (10).

\medskip\noindent
Preuve de (5). Soient $\mathcal{F}$ et $\mathcal{F}'$ dans
$\QCoh(S_\etale, \mathcal{O})$. Écrivons
$\mathcal{F} = \mathcal{G}^a$ et
$\mathcal{F}' = (\mathcal{G}')^a$, avec $\mathcal{G}$ et $\mathcal{G}'$
quasi-cohérents sur $S$. Par le lemme cité, on a
$\mathcal{F} \otimes_\mathcal{O} \mathcal{F}' =
(\mathcal{G} \otimes_{\mathcal{O}_S} \mathcal{G}')^a$,
d'où la conclusion.

\medskip\noindent
Preuve de (6). Soient $\mathcal{F}$ et $\mathcal{G}$ dans
$\QCoh(S_\etale, \mathcal{O})$, avec $\mathcal{F}$ de présentation finie.
Écrivons $\mathcal{F} = \mathcal{H}^a$ et
$\mathcal{G} = (\mathcal{I})^a$, avec $\mathcal{H}$ et $\mathcal{I}$
quasi-cohérents sur $S$. Par le
lemme \ref{descent-lemma-equivalence-quasi-coherent-properties}, le module
$\mathcal{H}$ est de présentation finie. Par le
lemme \ref{descent-lemma-equivalence-quasi-coherent-limits}, partie (8), on a
$\SheafHom_\mathcal{O}(\mathcal{F}, \mathcal{G}) =
(\SheafHom_{\mathcal{O}_S}(\mathcal{H}, \mathcal{I}))^a$,
d'où la conclusion.
\end{proof}

\begin{lemma}
\label{descent-lemma-properties-quasi-coherent-on-big}
Soit $S$ un schéma.
Soit $\tau \in \{Zariski, \linebreak[0] \etale, \linebreak[0]
smooth, \linebreak[0] syntomic, \linebreak[0] fppf\}$.
La catégorie $\QCoh((\Sch/S)_\tau, \mathcal{O})$ des modules
quasi-cohérents sur $(\Sch/S)_\tau$ possède les propriétés suivantes :
\begin{enumerate}
\item Toute somme directe de faisceaux quasi-cohérents est quasi-cohérente.
\item Toute colimite de faisceaux quasi-cohérents est quasi-cohérente.
\item Le conoyau d'un morphisme de faisceaux quasi-cohérents est
quasi-cohérent.
\item Étant donnée une suite exacte courte de $\mathcal{O}$-modules
$0 \to \mathcal{F}_1 \to \mathcal{F}_2 \to \mathcal{F}_3 \to 0$, si
$\mathcal{F}_1$ et $\mathcal{F}_3$ sont quasi-cohérents, alors
$\mathcal{F}_2$ l'est aussi.
\item Le produit tensoriel de deux $\mathcal{O}$-modules quasi-cohérents
est quasi-cohérent.
\item Étant donnés deux $\mathcal{O}$-modules quasi-cohérents
$\mathcal{F}$, $\mathcal{G}$ tels que $\mathcal{F}$ soit localement libre
de rang fini, le Hom interne
$\SheafHom_\mathcal{O}(\mathcal{F}, \mathcal{G})$
est quasi-cohérent.
\end{enumerate}
\end{lemma}

\begin{proof}
Les faits correspondants sont vrais pour les modules quasi-cohérents sur le
schéma $S$ ; voir Schémas, section
\ref{schemes-section-quasi-coherent}. La preuve consiste à utiliser le
lemme \ref{descent-lemma-equivalence-quasi-coherent-limits} pour les transporter à
$(\Sch/S)_\tau$.

\medskip\noindent
Preuve de (1). Soient $\mathcal{F}_i$, $i \in I$, une famille d'objets de
$\QCoh((\Sch/S)_\tau, \mathcal{O})$. Écrivons
$\mathcal{F}_i = \mathcal{G}_i^a$, où les $\mathcal{G}_i$ sont des modules
quasi-cohérents sur $S$. Alors
$\bigoplus \mathcal{F}_i = (\bigoplus \mathcal{G}_i)^a$ par le lemme cité,
ce qui donne la conclusion.

\medskip\noindent
Preuve de (2). Soit
$\mathcal{I} \to \QCoh((\Sch/S)_\tau, \mathcal{O})$,
$i \mapsto \mathcal{F}_i$, un diagramme. Écrivons
$\mathcal{F}_i = \mathcal{G}_i^a$ ; on obtient ainsi un diagramme
$\mathcal{I} \to \QCoh(\mathcal{O}_S)$. Alors
$\colim \mathcal{F}_i = (\colim \mathcal{G}_i)^a$ par le lemme cité,
ce qui donne la conclusion.

\medskip\noindent
Preuve de (3). Soit $a : \mathcal{F} \to \mathcal{F}'$ une flèche de
$\QCoh((\Sch/S)_\tau, \mathcal{O})$. Écrivons $a = b^a$ pour un morphisme
$b : \mathcal{G} \to \mathcal{G}'$ de modules quasi-cohérents sur $S$.
Par le lemme cité, on a $\Coker(a) = \Coker(b)^a$ (car un conoyau est une
colimite), d'où le résultat.

\medskip\noindent
Preuve de (4). Écrivons $\mathcal{F}_1 = \mathcal{G}_1^a$ et
$\mathcal{F}_3 = \mathcal{G}_3^a$, avec $\mathcal{G}_i$ quasi-cohérent sur
$S$. On conclut par le Lemme
\ref{descent-lemma-equivalence-quasi-coherent-limits}, partie (10).

\medskip\noindent
Preuve de (5). Soient $\mathcal{F}$ et $\mathcal{F}'$ dans
$\QCoh((\Sch/S)_\tau, \mathcal{O})$. Écrivons
$\mathcal{F} = \mathcal{G}^a$ et
$\mathcal{F}' = (\mathcal{G}')^a$, avec $\mathcal{G}$ et $\mathcal{G}'$
quasi-cohérents sur $S$. Par le lemme cité, on a
$\mathcal{F} \otimes_\mathcal{O} \mathcal{F}' =
(\mathcal{G} \otimes_{\mathcal{O}_S} \mathcal{G}')^a$,
d'où la conclusion.

\medskip\noindent
Preuve de (6). Écrivons $\mathcal{F} = \mathcal{H}^a$ pour un
$\mathcal{O}_S$-module quasi-cohérent. D'après le
lemme \ref{descent-lemma-equivalence-quasi-coherent-properties}, le module
$\mathcal{H}$ est localement libre de rang fini. La question est locale sur
$S$ pour la topologie de Zariski (détails omis), de sorte que l'on peut
supposer $\mathcal{H} = \mathcal{O}_S^{\oplus n}$ libre de rang fini.
Alors $\mathcal{F} = \mathcal{O}^{\oplus n}$ et
$\SheafHom_\mathcal{O}(\mathcal{F}, \mathcal{G}) = \mathcal{G}^{\oplus n}$
est quasi-cohérent.
\end{proof}

\begin{example}
\label{descent-example-internal-hom-not-qcoh}
Soit $S$ un schéma. Soient $\mathcal{F}$ et $\mathcal{G}$ des modules
quasi-cohérents sur $(\Sch/S)_\tau$, pour l'une des topologies $\tau$
considérées dans le lemme \ref{descent-lemma-properties-quasi-coherent-on-big}.
En général, $\SheafHom_\mathcal{O}(\mathcal{F}, \mathcal{G})$ n'est pas
quasi-cohérent, même si $\mathcal{F}$ est de présentation finie.
Prenons en effet $S = \Spec(\mathbf{Z})$,
$\mathcal{F} = \Coker(2 : \mathcal{O} \to \mathcal{O})$ et
$\mathcal{G} = \mathcal{O}$. Alors
$\SheafHom_\mathcal{O}(\mathcal{F}, \mathcal{G}) = \mathcal{O}[2]$
est le sous-faisceau de $2$-torsion de $\mathcal{O}$, qui n'est pas
quasi-cohérent.
\end{example}

\begin{lemma}
\label{descent-lemma-qc-colimits}
Soit $S$ un schéma.
\begin{enumerate}
\item La catégorie $\QCoh((\Sch/S)_{fppf}, \mathcal{O})$ admet des
colimites, et celles-ci coïncident avec les colimites dans les catégories
$\textit{Mod}((\Sch/S)_{Zar}, \mathcal{O})$,
$\textit{Mod}((\Sch/S)_\etale, \mathcal{O})$ et
$\textit{Mod}((\Sch/S)_{fppf}, \mathcal{O})$.
\item Étant donnés $\mathcal{F}, \mathcal{G}$ dans
$\QCoh((\Sch/S)_{fppf}, \mathcal{O})$, les produits tensoriels
$\mathcal{F} \otimes_\mathcal{O} \mathcal{G}$ calculés dans
$\textit{Mod}((\Sch/S)_{Zar}, \mathcal{O})$,
$\textit{Mod}((\Sch/S)_\etale, \mathcal{O})$ ou
$\textit{Mod}((\Sch/S)_{fppf}, \mathcal{O})$ coïncident, et leur valeur
commune est un objet de $\QCoh((\Sch/S)_{fppf}, \mathcal{O})$.
\item Étant donnés $\mathcal{F}, \mathcal{G}$ dans
$\QCoh((\Sch/S)_{fppf}, \mathcal{O})$, avec $\mathcal{F}$ localement libre
de rang fini (pour la topologie fppf, ou, de manière équivalente, pour la
topologie \'etale, ou encore pour la topologie de Zariski), les Hom internes
$\SheafHom_\mathcal{O}(\mathcal{F}, \mathcal{G})$ calculés dans
$\textit{Mod}((\Sch/S)_{Zar}, \mathcal{O})$,
$\textit{Mod}((\Sch/S)_\etale, \mathcal{O})$ ou
$\textit{Mod}((\Sch/S)_{fppf}, \mathcal{O})$ coïncident, et leur valeur
commune est un objet de $\QCoh((\Sch/S)_{fppf}, \mathcal{O})$.
\end{enumerate}
\end{lemma}

\begin{proof}
Ce lemme rassemble les résultats précédents sous une forme légèrement
différente. Tout d'abord, le
Lemme \ref{descent-lemma-properties-quasi-coherent-on-big} montre déjà que le
résultat de chacune des constructions (1), (2) et (3) appartient à
$\QCoh((\Sch/S)_\tau, \mathcal{O})$. Il reste, dans chaque cas, à montrer
que ce résultat ne dépend pas de la topologie utilisée. Le point essentiel
est que l'équivalence
$\QCoh(\mathcal{O}_S) \to \QCoh((\Sch/S)_\tau, \mathcal{O})$,
$\mathcal{F} \mapsto \mathcal{F}^a$, de la
proposition \ref{descent-proposition-equivalence-quasi-coherent}, est donnée par la
même formule, indépendamment du choix de la topologie
$\tau \in \{Zariski, \etale, fppf\}$.

\medskip\noindent
Preuve de (1). Soit
$\mathcal{I} \to \QCoh((\Sch/S)_{fppf}, \mathcal{O})$,
$i \mapsto \mathcal{F}_i$, un diagramme. Écrivons
$\mathcal{F}_i = \mathcal{G}_i^a$ ; on obtient ainsi un diagramme
$\mathcal{I} \to \QCoh(\mathcal{O}_S)$. Alors
$\colim \mathcal{F}_i = (\colim \mathcal{G}_i)^a$ dans
$\textit{Mod}((\Sch/S)_\tau, \mathcal{O})$ pour
$\tau \in \{Zariski, \etale, fppf\}$, par le
lemme \ref{descent-lemma-equivalence-quasi-coherent-limits}. Cela prouve (1).

\medskip\noindent
Preuve de (2). Écrivons $\mathcal{F} = \mathcal{H}^a$ et
$\mathcal{G} = (\mathcal{I})^a$, avec $\mathcal{H}$ et $\mathcal{I}$
quasi-cohérents sur $S$. Alors
$\mathcal{F} \otimes_\mathcal{O} \mathcal{G} =
(\mathcal{H} \otimes_\mathcal{O} \mathcal{I})^a$ dans
$\textit{Mod}((\Sch/S)_\tau, \mathcal{O})$ pour
$\tau \in \{Zariski, \etale, fppf\}$, par le
lemme \ref{descent-lemma-equivalence-quasi-coherent-limits}. Cela prouve (2).

\medskip\noindent
Preuve de (3). Soient $\mathcal{F}$ et $\mathcal{G}$ dans
$\QCoh((\Sch/S)_{fppf}, \mathcal{O})$. Écrivons
$\mathcal{F} = \mathcal{H}^a$, avec $\mathcal{H}$ quasi-cohérent sur $S$.
D'après le lemme \ref{descent-lemma-equivalence-quasi-coherent-properties}, on a
\begin{align*}
&\mathcal{F}\text{ est localement libre de rang fini} \\
&\qquad\text{pour la topologie fppf} \\
& \Leftrightarrow
\mathcal{H}\text{ est localement libre de rang fini sur }S \\
& \Leftrightarrow
\mathcal{F}\text{ est localement libre de rang fini} \\
&\qquad\text{pour la topologie \'etale} \\
& \Leftrightarrow
\mathcal{H}\text{ est localement libre de rang fini sur }S \\
& \Leftrightarrow
\mathcal{F}\text{ est localement libre de rang fini} \\
&\qquad\text{pour la topologie de Zariski}
\end{align*}
Cela explique la parenthèse de la partie (3). Si ces conditions équivalentes
sont satisfaites, alors $\mathcal{H}$ est localement libre de rang fini.
La construction de $\SheafHom_\mathcal{O}(\mathcal{F}, \mathcal{G})$ dans
Modules sur les sites, section
\ref{sites-modules-section-internal-hom}, ne dépend que de $\mathcal{F}$ et
$\mathcal{G}$ comme préfaisceaux de modules ; seul le fait que le résultat
$\SheafHom$ soit un faisceau dépend du fait que $\mathcal{F}$ et
$\mathcal{G}$ soient eux-mêmes des faisceaux.
\end{proof}










\section{Modules quasi-cohérents et schémas affines}
\label{descent-section-alternative-quasi-coherent}

\noindent
Soit $S$ un schéma\footnote{Dans cette section, comme dans Topologies,
section \ref{topologies-section-change-topologies}, nous choisissons les sites
$(\Sch/S)_\tau$ de façon qu'ils aient la même catégorie sous-jacente pour
$\tau \in \{Zariski, \etale, smooth, syntomic, fppf\}$. Les sites
$(\textit{Aff}/S)_\tau$ ont alors eux aussi la même catégorie sous-jacente.}.
Soit $\tau \in \{Zariski, \etale, smooth, syntomic, fppf\}$.
Rappelons que $(\textit{Aff}/S)_\tau$ est la sous-catégorie pleine de
$(\Sch/S)_\tau$ formée des objets affines, munie de la structure de site
pour laquelle les recouvrements sont les recouvrements $\tau$ standards.
D'après les lemmes
\ref{topologies-lemma-affine-big-site-Zariski},
\ref{topologies-lemma-affine-big-site-etale},
\ref{topologies-lemma-affine-big-site-smooth},
\ref{topologies-lemma-affine-big-site-syntomic} et
\ref{topologies-lemma-affine-big-site-fppf} de Topologies, il existe une équivalence de
topos
$g : \Sh((\textit{Aff}/S)_\tau) \to \Sh((\Sch/S)_\tau)$
dont le foncteur d'image inverse est donné par restriction.
Rappelons que $\mathcal{O}$ désigne le faisceau structural sur
$(\Sch/S)_\tau$ et notons, provisoirement et avec toute la précision voulue,
$\mathcal{O}_{\textit{Aff}}$ la restriction de $\mathcal{O}$ à
$(\textit{Aff}/S)_\tau$. On obtient alors une équivalence
\begin{equation}
\label{descent-equation-alternative-ringed}
(\Sh((\textit{Aff}/S)_\tau), \mathcal{O}_{\textit{Aff}})
\longrightarrow
(\Sh((\Sch/S)_\tau), \mathcal{O})
\end{equation}
de topos annelés. Cela permet également de comparer les modules
quasi-cohérents.

\begin{lemma}
\label{descent-lemma-quasi-coherent-alternative}
Soit $S$ un schéma. Soit $\mathcal{F}$ un préfaisceau de
$\mathcal{O}_{\textit{Aff}}$-modules sur $(\textit{Aff}/S)_{fppf}$.
Les conditions suivantes sont équivalentes :
\begin{enumerate}
\item pour tout morphisme $U \to U'$ de $(\textit{Aff}/S)_{fppf}$,
l'application
$\mathcal{F}(U') \otimes_{\mathcal{O}(U')} \mathcal{O}(U) \to \mathcal{F}(U)$
est un isomorphisme ;
\item $\mathcal{F}$ est un faisceau sur $(\textit{Aff}/S)_{Zar}$ et un
module quasi-cohérent sur le site annelé
$((\textit{Aff}/S)_{Zar}, \mathcal{O}_{\textit{Aff}})$ au sens de Modules
sur les sites, définition \ref{sites-modules-definition-site-local} ;
\item la même condition que (2) pour la topologie \'etale ;
\item la même condition que (2) pour la topologie lisse ;
\item la même condition que (2) pour la topologie syntomique ;
\item la même condition que (2) pour la topologie fppf ;
\item $\mathcal{F}$ correspond à un module quasi-cohérent sur
$(\Sch/S)_{Zar}$,
$(\Sch/S)_\etale$,
$(\Sch/S)_{smooth}$,
$(\Sch/S)_{syntomic}$ ou
$(\Sch/S)_{fppf}$ par l'équivalence
(\ref{descent-equation-alternative-ringed}) ;
\item $\mathcal{F}$ provient d'un unique $\mathcal{O}_S$-module
quasi-cohérent $\mathcal{G}$ par le procédé décrit à la section
\ref{descent-section-quasi-coherent-sheaves}.
\end{enumerate}
\end{lemma}

\begin{proof}
La notion de module quasi-cohérent étant intrinsèque (Modules sur les sites,
lemme \ref{sites-modules-lemma-special-locally-free}), l'équivalence
(\ref{descent-equation-alternative-ringed}) induit une équivalence entre les
catégories de modules quasi-cohérents. La
proposition \ref{descent-proposition-equivalence-quasi-coherent} affirme que la
catégorie des modules quasi-cohérents sur $\Sch/S$ est indépendante de la
topologie choisie, et que (8) équivaut à (7). Les conditions
(2)--(8) sont donc toutes équivalentes.

\medskip\noindent
Supposons satisfaites les conditions équivalentes (2)--(8), et soit
$\mathcal{G}$ comme en (8). Soit $h : U \to U' \to S$ un morphisme de
$\textit{Aff}/S$. Notons $f : U \to S$ et $f' : U' \to S$ les morphismes
structuraux, de sorte que $f = f' \circ h$. On a
$\mathcal{F}(U') = \Gamma(U', (f')^*\mathcal{G})$ et
$\mathcal{F}(U) = \Gamma(U, f^*\mathcal{G}) = \Gamma(U, h^*(f')^*\mathcal{G})$.
La condition (1) résulte donc du lemme
\ref{schemes-lemma-widetilde-pullback} de Schémas.

\medskip\noindent
Supposons (1). Pour achever la preuve, il suffit d'établir (2).
Soit $U$ un objet de $(\textit{Aff}/S)_{Zar}$. Écrivons
$U = \Spec(R)$. Un recouvrement ouvert standard
$U = U_1 \cup \ldots \cup U_n$ est donné par $U_i = D(f_i)$, où les
éléments $f_1, \ldots, f_n \in R$ engendrent l'idéal unité de $R$.
La propriété (1) donne
$$
\mathcal{F}(U_i) =
\mathcal{F}(U) \otimes_R R_{f_i} =
\mathcal{F}(U)_{f_i}
$$
et
$$
\mathcal{F}(U_i \cap U_j) =
\mathcal{F}(U) \otimes_R R_{f_if_j} =
\mathcal{F}(U)_{f_if_j}
$$
Le lemme \ref{algebra-lemma-cover-module} d'Algèbre montre alors que
$\mathcal{F}$ est un faisceau sur $(\textit{Aff}/S)_{Zar}$. Choisissons une
présentation
$$
\bigoplus\nolimits_{k \in K} R
\longrightarrow
\bigoplus\nolimits_{l \in L} R
\longrightarrow
\mathcal{F}(U)
\longrightarrow 0
$$
par des $R$-modules libres. La propriété (1) et l'exactitude à droite du
produit tensoriel montrent que, pour tout morphisme $U' \to U$ dans
$(\textit{Aff}/S)_{Zar}$, on obtient une présentation
$$
\bigoplus\nolimits_{k \in K} \mathcal{O}_{Aff}(U')
\longrightarrow
\bigoplus\nolimits_{l \in L} \mathcal{O}_{Aff}(U')
\longrightarrow
\mathcal{F}(U')
\longrightarrow 0
$$
Autrement dit, la restriction de $\mathcal{F}$ à la catégorie localisée
$(\textit{Aff}/S)_{Zar}/U$ admet une présentation
$$
\bigoplus\nolimits_{k \in K} \mathcal{O}_{Aff}|_{(\textit{Aff}/S)_{Zar}/U}
\longrightarrow
\bigoplus\nolimits_{l \in L} \mathcal{O}_{Aff}|_{(\textit{Aff}/S)_{Zar}/U}
\longrightarrow
\mathcal{F}|_{(\textit{Aff}/S)_{Zar}/U}
\longrightarrow 0
$$
Malgré cette notation peu élégante, cela achève la preuve.
\end{proof}

\noindent
Poursuivons la discussion commencée dans l'introduction de cette section.
Soit $\tau \in \{Zariski, \etale\}$. Rappelons que $S_{affine, \tau}$ est la
sous-catégorie pleine de $S_\tau$ formée des objets affines, munie de la
structure de site pour laquelle les recouvrements sont les recouvrements
$\tau$ standards. Voir les définitions
\ref{topologies-definition-big-small-Zariski} et
\ref{topologies-definition-big-small-etale} de Topologies. D'après les lemmes
\ref{topologies-lemma-alternative-zariski}, respectivement\  
\ref{topologies-lemma-alternative} de Topologies, il existe une équivalence de topos
$g : \Sh(S_{affine, \tau}) \to \Sh(S_\tau)$, dont le foncteur d'image inverse
est donné par restriction. Rappelons que $\mathcal{O}$ désigne le faisceau
structural sur $S_\tau$ et notons, provisoirement et avec toute la précision
voulue, $\mathcal{O}_{affine}$ la restriction de $\mathcal{O}$ à
$S_{affine, \tau}$. On obtient alors une équivalence
\begin{equation}
\label{descent-equation-alternative-small-ringed}
(\Sh(S_{affine, \tau}), \mathcal{O}_{affine})
\longrightarrow
(\Sh(S_\tau), \mathcal{O})
\end{equation}
de topos annelés. Cela permet également de comparer les modules
quasi-cohérents.

\begin{lemma}
\label{descent-lemma-quasi-coherent-alternative-small}
Soit $S$ un schéma. Soit $\tau \in \{Zariski, \etale\}$.
Soit $\mathcal{F}$ un préfaisceau de $\mathcal{O}_{affine}$-modules sur
$S_{affine, \tau}$. Les conditions suivantes sont équivalentes :
\begin{enumerate}
\item pour tout morphisme $U \to U'$ de $S_{affine, \tau}$, l'application
$\mathcal{F}(U') \otimes_{\mathcal{O}(U')} \mathcal{O}(U) \to \mathcal{F}(U)$
est un isomorphisme ;
\item $\mathcal{F}$ est un faisceau sur $S_{affine, \tau}$ et un module
quasi-cohérent sur le site annelé
$(S_{affine, \tau}, \mathcal{O}_{affine})$ au sens de Modules sur les sites,
définition \ref{sites-modules-definition-site-local} ;
\item $\mathcal{F}$ correspond à un module quasi-cohérent sur $S_\tau$ par
l'équivalence (\ref{descent-equation-alternative-small-ringed}) ;
\item $\mathcal{F}$ provient d'un unique $\mathcal{O}_S$-module
quasi-cohérent $\mathcal{G}$ par le procédé décrit à la section
\ref{descent-section-quasi-coherent-sheaves}.
\end{enumerate}
\end{lemma}

\begin{proof}
Démontrons l'énoncé dans le cas de la topologie \'etale.

\medskip\noindent
Supposons (1). Pour montrer que $\mathcal{F}$ est un faisceau, soit
$\mathcal{U} = \{U_i \to U\}_{i = 1, \ldots, n}$ un recouvrement de
$S_{affine, \etale}$. La condition de faisceau pour $\mathcal{F}$ et
$\mathcal{U}$, compte tenu de notre hypothèse sur $\mathcal{F}$, se ramène à
montrer que
$$
0 \to \mathcal{F}(U) \to
\prod \mathcal{F}(U) \otimes_{\mathcal{O}(U)} \mathcal{O}(U_i) \to
\prod \mathcal{F}(U) \otimes_{\mathcal{O}(U)} \mathcal{O}(U_i \times_U U_j)
$$
est exacte. En effet, $\mathcal{O}(U) \to \prod \mathcal{O}(U_i)$ est
fidèlement plat, par le lemme \ref{descent-lemma-standard-covering-Cech} et parce que
les recouvrements de $S_{affine, \etale}$ sont, par définition, les
recouvrements \'etales standards ; on peut donc appliquer le lemme
\ref{descent-lemma-ff-exact}.
Montrons ensuite que $\mathcal{F}$ est quasi-cohérent sur
$S_{affine, \etale}$. Pour $U$ dans $S_{affine, \etale}$, posons
$R = \mathcal{O}(U)$ et choisissons une présentation
$$
\bigoplus\nolimits_{k \in K} R
\longrightarrow
\bigoplus\nolimits_{l \in L} R
\longrightarrow
\mathcal{F}(U)
\longrightarrow 0
$$
par des $R$-modules libres. La propriété (1) et l'exactitude à droite du
produit tensoriel montrent que, pour tout morphisme $U' \to U$ dans
$S_{affine, \etale}$, on obtient une présentation
$$
\bigoplus\nolimits_{k \in K} \mathcal{O}(U')
\longrightarrow
\bigoplus\nolimits_{l \in L} \mathcal{O}(U')
\longrightarrow
\mathcal{F}(U')
\longrightarrow 0
$$
Autrement dit, la restriction de $\mathcal{F}$ à la catégorie localisée
$S_{affine, \etale}/U$ admet une présentation
$$
\bigoplus\nolimits_{k \in K} \mathcal{O}_{affine}|_{S_{affine, \etale}/U}
\longrightarrow
\bigoplus\nolimits_{l \in L} \mathcal{O}_{affine}|_{S_{affine, \etale}/U}
\longrightarrow
\mathcal{F}|_{S_{affine, \etale}/U}
\longrightarrow 0
$$
comme il le faut pour montrer que $\mathcal{F}$ est quasi-cohérent.
Malgré cette notation peu élégante, nous avons ainsi démontré que (1)
implique (2).

\medskip\noindent
La notion de module quasi-cohérent étant intrinsèque (Modules sur les sites,
lemme \ref{sites-modules-lemma-special-locally-free}), l'équivalence
(\ref{descent-equation-alternative-small-ringed}) induit une équivalence entre les
catégories de modules quasi-cohérents. Les conditions (2) et (3) sont donc
équivalentes.

\medskip\noindent
L'équivalence de (3) et (4) résulte de la
proposition \ref{descent-proposition-equivalence-quasi-coherent}.

\medskip\noindent
Supposons (4) et démontrons (1). Soit $\mathcal{G}$ comme en (4).
Soit $h : U \to U' \to S$ un morphisme de $S_{affine, \etale}$.
Notons $f : U \to S$ et $f' : U' \to S$ les morphismes structuraux, de sorte
que $f = f' \circ h$. On a
$\mathcal{F}(U') = \Gamma(U', (f')^*\mathcal{G})$ et
$\mathcal{F}(U) = \Gamma(U, f^*\mathcal{G}) = \Gamma(U, h^*(f')^*\mathcal{G})$.
La condition (1) résulte donc du lemme
\ref{schemes-lemma-widetilde-pullback} de Schémas.

\medskip\noindent
Nous omettons la démonstration dans le cas de la topologie de Zariski.
\end{proof}





\section{Modules parasites}
\label{descent-section-parasitic}

\noindent
Les modules parasites sont ceux qui deviennent nuls lorsqu'on les restreint
aux schémas plats sur le schéma de base. Voici la définition précise.

\begin{definition}
\label{descent-definition-parasitic}
Soit $S$ un schéma. Soit $\tau \in \{Zar, \etale,
smooth, syntomic, fppf\}$. Soit $\mathcal{F}$ un préfaisceau de
$\mathcal{O}$-modules sur $(\Sch/S)_\tau$.
\begin{enumerate}
\item On dit que $\mathcal{F}$ est
{\it parasite}\footnote{Cette terminologie n'est peut-être pas standard.}
si, pour tout morphisme plat $U \to S$, on a $\mathcal{F}(U) = 0$.
\item On dit que $\mathcal{F}$ est {\it parasite pour la topologie $\tau$}
si, pour tout recouvrement $\tau$,
$\{U_i \to S\}_{i \in I}$, on a $\mathcal{F}(U_i) = 0$ pour tout $i$.
\end{enumerate}
\end{definition}

\noindent
Si $\tau = fppf$, cela signifie que $\mathcal{F}|_{U_{Zar}} = 0$ chaque fois
que $U \to S$ est plat et localement de présentation finie ; les autres cas
sont analogues.

\begin{lemma}
\label{descent-lemma-cohomology-parasitic}
Soit $S$ un schéma. Soit $\tau \in \{Zar, \etale, smooth,
syntomic, fppf\}$. Soit $\mathcal{G}$ un préfaisceau de
$\mathcal{O}$-modules sur $(\Sch/S)_\tau$.
\begin{enumerate}
\item Si $\mathcal{G}$ est parasite pour la topologie $\tau$, alors
$H^p_\tau(U, \mathcal{G}) = 0$ pour tout $U$ ouvert dans $S$,
resp.\ \'etale sur $S$,
resp.\ lisse sur $S$,
resp.\ syntomique sur $S$,
resp.\ plat et localement de présentation finie sur $S$.
\item Si $\mathcal{G}$ est parasite, alors
$H^p_\tau(U, \mathcal{G}) = 0$ pour tout $U$ plat sur $S$.
\end{enumerate}
\end{lemma}

\begin{proof}
Démontrons-le lorsque $\tau = fppf$ ; les autres cas se démontrent exactement
de la même manière. L'hypothèse signifie que $\mathcal{G}(U) = 0$ pour tout
$U \to S$ plat et localement de présentation finie. Appliquons le lemme
\ref{sites-cohomology-lemma-cech-vanish-collection} de Cohomologie sur les
sites au sous-ensemble
$\mathcal{B} \subset \Ob((\Sch/S)_{fppf})$ formé des $U \to S$ plats et
localement de présentation finie, et à la collection $\text{Cov}$ de tous les
recouvrements fppf des éléments de $\mathcal{B}$.
\end{proof}

\begin{lemma}
\label{descent-lemma-direct-image-parasitic}
Soit $f : T \to S$ un morphisme de schémas. Pour tout
$\mathcal{O}$-module parasite sur $(\Sch/T)_\tau$, l'image directe
$f_*\mathcal{F}$ et les images directes supérieures $R^if_*\mathcal{F}$
sont des $\mathcal{O}$-modules parasites sur $(\Sch/S)_\tau$.
\end{lemma}

\begin{proof}
Rappelons que $R^if_*\mathcal{F}$ est le faisceau associé au préfaisceau
$$
U \mapsto H^i((\Sch/U \times_S T)_\tau, \mathcal{F})
$$
voir le lemme \ref{sites-cohomology-lemma-higher-direct-images} de
Cohomologie sur les sites.
Si $U \to S$ est plat, alors $U \times_S T \to T$ est plat par changement de
base. Le groupe affiché est donc nul par le
lemme \ref{descent-lemma-cohomology-parasitic}. Si $\{U_i \to U\}$ est un
recouvrement $\tau$, alors $U_i \times_S T \to T$ est également plat.
Il est par conséquent clair que le faisceau associé au préfaisceau affiché
est nul sur les schémas $U$ plats sur $S$.
\end{proof}

\begin{lemma}
\label{descent-lemma-quasi-coherent-and-flat-base-change}
Soit $S$ un schéma. Soit $\tau \in \{Zar, \etale\}$.
Soit $\mathcal{G}$ un faisceau de $\mathcal{O}$-modules sur
$(\Sch/S)_{fppf}$ tel que
\begin{enumerate}
\item $\mathcal{G}|_{S_\tau}$ est quasi-cohérent ;
\item pour tout morphisme plat et localement de présentation finie
$g : U \to S$, l'application canonique
$g_{\tau, small}^*(\mathcal{G}|_{S_\tau}) \to \mathcal{G}|_{U_\tau}$
est un isomorphisme.
\end{enumerate}
Alors $H^p(U, \mathcal{G}) = H^p(U, \mathcal{G}|_{U_\tau})$ pour tout $U$
plat et localement de présentation finie sur $S$.
\end{lemma}

\begin{proof}
Soit $\mathcal{F}$ l'image inverse de $\mathcal{G}|_{S_\tau}$ sur le grand
site fppf $(\Sch/S)_{fppf}$. Remarquons que $\mathcal{F}$ est quasi-cohérent.
Il existe une application de comparaison canonique
$\varphi : \mathcal{F} \to \mathcal{G}$ qui, en vertu des hypothèses (1) et
(2), induit un isomorphisme
$\mathcal{F}|_{U_\tau} \to \mathcal{G}|_{U_\tau}$ pour tout
$g : U \to S$ plat et localement de présentation finie.
Ainsi, dans les suites exactes courtes
$$
0 \to \Ker(\varphi) \to \mathcal{F} \to \Im(\varphi) \to 0
$$
et
$$
0 \to \Im(\varphi) \to \mathcal{G} \to \Coker(\varphi) \to 0
$$
les faisceaux $\Ker(\varphi)$ et $\Coker(\varphi)$ sont parasites pour la
topologie fppf. Le lemme \ref{descent-lemma-cohomology-parasitic} montre donc que
$H^p(U, \mathcal{F}) \to H^p(U, \mathcal{G})$ est un isomorphisme lorsque
$g : U \to S$ est plat et localement de présentation finie.
Le résultat étant vrai pour $\mathcal{F}$ d'après la
proposition \ref{descent-proposition-same-cohomology-quasi-coherent}, la preuve est
achevée.
\end{proof}













\section[Recouvrements fpqc et épimorphismes]{Les recouvrements fpqc sont des épimorphismes effectifs universels}
\label{descent-section-fpqc-universal-effective-epimorphisms}

\noindent
Nous appliquons les résultats précédents pour démontrer un résultat
intéressant, à savoir le lemme
\ref{descent-lemma-fpqc-universal-effective-epimorphisms}.
D'après Sites, section \ref{sites-section-representable-sheaves},
ce lemme implique que les préfaisceaux représentables sur chacun des sites
$(\Sch/S)_\tau$ sont des faisceaux pour
$\tau \in \{Zariski, fppf, \etale, smooth, syntomic\}$. Commençons
par démontrer un lemme auxiliaire.

\begin{lemma}
\label{descent-lemma-equiv-fibre-product}
Pour un schéma $X$, notons $|X|$ son ensemble sous-jacent.
Soit $f : X \to S$ un morphisme de schémas.
Alors
$$
|X \times_S X| \to |X| \times_{|S|} |X|
$$
est surjectif.
\end{lemma}

\begin{proof}
Cela résulte immédiatement de la description des points du produit
fibré donnée dans Schémas, lemme
\ref{schemes-lemma-points-fibre-product}.
\end{proof}


\begin{lemma}
\label{descent-lemma-universal-effective-epimorphism-affine}
Soit $\{f_i : X_i \to X\}_{i \in I}$ une famille de morphismes de schémas affines.
Les assertions suivantes sont équivalentes :
\begin{enumerate}
\item pour tout $\mathcal{O}_X$-module quasi-cohérent $\mathcal{F}$, on a
$$
\Gamma(X, \mathcal{F}) =
\text{Égalisateur}\left(
\xymatrix{
\prod\nolimits_{i \in I} \Gamma(X_i, f_i^*\mathcal{F})
\ar@<1ex>[r] \ar@<-1ex>[r] &
\prod\nolimits_{i, j \in I}
\Gamma(X_i \times_X X_j, (f_i \times f_j)^*\mathcal{F})
}
\right)
$$
\item $\{f_i : X_i \to X\}_{i \in I}$ est un épimorphisme effectif universel
(Sites, définition \ref{sites-definition-universal-effective-epimorphisms})
dans la catégorie des schémas affines.
\end{enumerate}
\end{lemma}

\begin{proof}
Supposons (2) et soit $\mathcal{F}$ un $\mathcal{O}_X$-module quasi-cohérent.
Considérons le schéma
(Constructions, section \ref{constructions-section-spec})
$$
X' = \underline{\Spec}_X(\mathcal{O}_X \oplus \mathcal{F})
$$
où $\mathcal{O}_X \oplus \mathcal{F}$ est la $\mathcal{O}_X$-algèbre dont
la multiplication est
$(f, s)(f', s') = (ff', fs' + f's)$.
Si $s_i \in \Gamma(X_i, f_i^*\mathcal{F})$ est une section,
alors $s_i$ détermine un unique élément de
$$
\Gamma(X' \times_X X_i, \mathcal{O}_{X' \times_X X_i}) =
\Gamma(X_i, \mathcal{O}_{X_i}) \oplus \Gamma(X_i, f_i^*\mathcal{F})
$$
Nous omettons la démonstration de cette égalité.
Si $(s_i)_{i \in I}$ appartient à l'égalisateur de (1), alors, en utilisant
l'égalité
$$
\Mor(T, \mathbf{A}^1_\mathbf{Z}) = \Gamma(T, \mathcal{O}_T)
$$
valable pour tout schéma $T$, on voit que ces sections définissent une famille
de morphismes $h_i : X' \times_X X_i \to \mathbf{A}^1_\mathbf{Z}$ telle que
$h_i \circ \text{pr}_1 = h_j \circ \text{pr}_2$ comme morphismes
$(X' \times_X X_i) \times_{X'} (X' \times_X X_j) \to \mathbf{A}^1_\mathbf{Z}$.
Puisque nous supposons (2), nous obtenons un morphisme
$h : X' \to \mathbf{A}^1_\mathbf{Z}$ compatible avec les morphismes $h_i$,
qui détermine à son tour un élément
$s \in \Gamma(X, \mathcal{F})$.
Nous omettons la vérification que $s$ s'envoie sur $s_i$ dans
$\Gamma(X_i, f_i^*\mathcal{F})$.

\medskip\noindent
Supposons (1). Soient $T$ un schéma affine et $h_i : X_i \to T$
une famille de morphismes telle que
$h_i \circ \text{pr}_1 = h_j \circ \text{pr}_2$ sur
$X_i \times_X X_j$ pour tous $i, j \in I$. Alors
$$
\prod h_i^\sharp :
\Gamma(T, \mathcal{O}_T)
\to
\prod \Gamma(X_i, \mathcal{O}_{X_i})
$$
prend ses valeurs dans l'égalisateur et l'on obtient un homomorphisme d'anneaux
$\Gamma(T, \mathcal{O}_T) \to \Gamma(X, \mathcal{O}_X)$ en appliquant
l'hypothèse du lemme à $\mathcal{F} = \mathcal{O}_X$.
Cet homomorphisme d'anneaux correspond à un morphisme $h : X \to T$
tel que $h_i = h \circ f_i$. Notre famille est donc un épimorphisme effectif.

\medskip\noindent
Soit $p : Y \to X$ un morphisme de schémas affines. Nous allons montrer que
les changements de base $g_i : Y_i \to Y$ des $f_i$ forment un épimorphisme
effectif, en appliquant le résultat du paragraphe précédent.
En effet, si $\mathcal{G}$ est un $\mathcal{O}_Y$-module quasi-cohérent, alors
$$
\Gamma(Y, \mathcal{G}) = \Gamma(X, p_*\mathcal{G}),\quad
\Gamma(Y_i, g_i^*\mathcal{G}) = \Gamma(X, f_i^*p_*\mathcal{G}),
$$
et
$$
\Gamma(Y_i \times_Y Y_j, (g_i \times g_j)^*\mathcal{G}) =
\Gamma(X, (f_i \times f_j)^*p_*\mathcal{G})
$$
par la formule triviale de changement de base
(Cohomologie des schémas, lemme \ref{coherent-lemma-affine-base-change}).
Ainsi, l'assertion (1) du lemme vaut pour la famille $g_i$.
\end{proof}

\begin{lemma}
\label{descent-lemma-universal-effective-epimorphism-surjective}
Soit $\{f_i : X_i \to X\}_{i \in I}$ une famille de morphismes de schémas.
\begin{enumerate}
\item Si la famille est un épimorphisme effectif universel dans la catégorie
des schémas, alors $\coprod f_i$ est surjectif.
\item Si $X$ et les $X_i$ sont affines et si la famille est un épimorphisme
effectif universel dans la catégorie des schémas affines, alors
$\coprod f_i$ est surjectif.
\end{enumerate}
\end{lemma}

\begin{proof}
Omis. Indication : effectuer le changement de base par
$\Spec(\kappa(x)) \to X$ pour voir que tout $x \in X$ doit appartenir à
l'image.
\end{proof}

\begin{lemma}
\label{descent-lemma-check-universal-effective-epimorphism-affine}
Soit $\{f_i : X_i \to X\}_{i \in I}$ une famille de morphismes de schémas.
Si, pour tout morphisme $Y \to X$ avec $Y$ affine, la famille des changements
de base $g_i : Y_i \to Y$ forme un épimorphisme effectif, alors la famille des
$f_i$ est un épimorphisme effectif universel dans la catégorie des schémas.
\end{lemma}

\begin{proof}
Soit $Y \to X$ un morphisme de schémas. Il faut montrer que les changements
de base $g_i : Y_i \to Y$ forment un épimorphisme effectif.
Pour cela, supposons donnés un schéma $T$ et des morphismes $h_i : Y_i \to T$
tels que $h_i \circ \text{pr}_1 = h_j \circ \text{pr}_2$ sur
$Y_i \times_Y Y_j$.
Choisissons un recouvrement ouvert affine $Y = \bigcup V_\alpha$.
Notons $V_{\alpha, i}$ l'image inverse de $V_\alpha$ dans $Y_i$.
Alors $V_{\alpha, i}  \to V_\alpha$ est le changement de base de $f_i$ par
$V_\alpha \to X$. Par hypothèse, la famille des restrictions
$h_i|_{V_{\alpha, i}}$ provient donc d'un morphisme de schémas
$h_\alpha : V_\alpha \to T$.
Nous laissons au lecteur le soin de montrer que ces morphismes coïncident sur
les intersections et définissent le morphisme voulu $Y \to T$.
Voir la discussion de Schémas, section
\ref{schemes-section-glueing-schemes}.
\end{proof}

\begin{lemma}
\label{descent-lemma-universal-effective-epimorphism}
Soit $\{f_i : X_i \to X\}_{i \in I}$ une famille de morphismes de schémas
affines. Supposons que l'hypothèse équivalente du lemme
\ref{descent-lemma-universal-effective-epimorphism-affine} soit satisfaite et que,
de plus, pour tout morphisme de schémas affines $Y \to X$, l'application
$$
\coprod X_i \times_X Y \longrightarrow Y
$$
soit une application submersive d'espaces topologiques
(Topologie, définition \ref{topology-definition-submersive}).
Alors notre famille de morphismes est un épimorphisme effectif universel dans
la catégorie des schémas.
\end{lemma}

\begin{proof}
D'après le lemme \ref{descent-lemma-check-universal-effective-epimorphism-affine},
il suffit d'effectuer sur notre famille le changement de base par
$Y \to X$, avec $Y$ affine. Posons $Y_i = X_i \times_X Y$.
Soient $T$ un schéma et $h_i : Y_i \to T$ une famille de morphismes tels que
$h_i \circ \text{pr}_1 = h_j \circ \text{pr}_2$ sur $Y_i \times_Y Y_j$.
Remarquons que l'ensemble $Y$ est le coégalisateur des deux applications de
$\coprod Y_i \times_Y Y_j$ vers $\coprod Y_i$.
La surjectivité résulte du cas affine du lemme
\ref{descent-lemma-universal-effective-epimorphism-surjective}, et l'injectivité du
lemme \ref{descent-lemma-equiv-fibre-product}.
Il existe donc une application entre ensembles sous-jacents $h : Y \to T$
compatible avec les applications $h_i$. La seconde condition du lemme montre
que $h$ est continue. Ainsi, si $y \in Y$ et si $U \subset T$ est un voisinage
ouvert affine de $h(y)$, on peut trouver un ouvert affine $V \subset Y$ tel que
$h(V) \subset U$.
En posant $V_i = Y_i \times_Y V = X_i \times_X V$, on peut appliquer le
résultat démontré au lemme
\ref{descent-lemma-universal-effective-epimorphism-affine} pour voir que
$h|_V : V \to U \subset T$ provient d'un unique morphisme de schémas affines
$h_V : V \to U$ qui coïncide avec $h_i|_{V_i}$ comme morphisme de schémas pour
tout $i$. En recollant ces $h_V$
(voir Schémas, section \ref{schemes-section-glueing-schemes}), on obtient le
morphisme voulu $Y \to T$.
\end{proof}

\begin{lemma}
\label{descent-lemma-open-fpqc-covering}
Soit $\{f_i : T_i \to T\}_{i \in I}$ un recouvrement fpqc.
Supposons que, pour chaque $i$, on dispose d'un ouvert $W_i \subset T_i$ tel
que, pour tous $i, j \in I$,
$\text{pr}_0^{-1}(W_i) = \text{pr}_1^{-1}(W_j)$ comme ouverts de
$T_i \times_T T_j$. Alors il existe un unique ouvert $W \subset T$ tel que
$W_i = f_i^{-1}(W)$ pour tout $i$.
\end{lemma}

\begin{proof}
Appliquons le lemme \ref{descent-lemma-equiv-fibre-product} à l'application
$\coprod_{i \in I} T_i \to T$.
Il en résulte qu'il existe une partie $W \subset T$ telle que
$W_i = f_i^{-1}(W)$ pour tout $i$, à savoir $W = \bigcup f_i(W_i)$.
Pour voir que $W$ est ouvert, on peut travailler localement pour la topologie
de Zariski sur $T$. On peut donc supposer que $T$ est affine.
D'après Topologies, définition
\ref{topologies-definition-fpqc-covering}, on peut choisir un recouvrement
fpqc standard $\{g_j : V_j \to T\}_{j \in J}$ qui raffine
$\{T_i \to T\}_{i \in I}$. Soient $\alpha : J \to I$ et
$h_j : V_j \to T_{\alpha(j)}$ comme dans Sites, définition
\ref{sites-definition-morphism-coverings}.
Alors $g_j^{-1}(W) = h_j^{-1}(W_{\alpha(j)})$.
On peut donc supposer que $\{f_i : T_i \to T\}$ est un recouvrement fpqc
standard. Dans ce cas, on applique Morphismes, lemme
\ref{morphisms-lemma-fpqc-quotient-topology}, au morphisme
$\coprod T_i \to T$ pour conclure que $W$ est ouvert.
\end{proof}

\begin{lemma}
\label{descent-lemma-fpqc-universal-effective-epimorphisms}
Soit $\{T_i \to T\}$ un recouvrement fpqc, voir Topologies, définition
\ref{topologies-definition-fpqc-covering}.
Alors $\{T_i \to T\}$ est un épimorphisme effectif universel dans la catégorie
des schémas, voir Sites, définition
\ref{sites-definition-universal-effective-epimorphisms}.
Autrement dit, tout foncteur représentable sur la catégorie des schémas
vérifie la condition de faisceau pour la topologie fpqc, voir Topologies,
définition \ref{topologies-definition-sheaf-property-fpqc}.
\end{lemma}

\begin{proof}
Soit $S$ un schéma. Il faut montrer l'assertion suivante : étant donnés des
morphismes $\varphi_i : T_i \to S$ tels que
$\varphi_i|_{T_i \times_T T_j} = \varphi_j|_{T_i \times_T T_j}$, il existe un
unique morphisme $T \to S$ qui se restreint à $\varphi_i$ sur chaque $T_i$.
Autrement dit, il faut montrer que le foncteur
$h_S = \Mor_{\Sch}( - , S)$ vérifie la condition de faisceau pour la topologie
fpqc.

\medskip\noindent
Si $\{T_i \to T\}$ est un recouvrement de Zariski, cela résulte de Schémas,
lemme \ref{schemes-lemma-glue}. Topologies, lemme
\ref{topologies-lemma-sheaf-property-fpqc}, nous ramène donc au cas d'un
recouvrement $\{X \to Y\}$ donné par un unique morphisme plat surjectif de
schémas affines.

\medskip\noindent
Première démonstration. Le lemme \ref{descent-lemma-sheaf-condition-holds} donne la
condition de faisceau pour les modules quasi-cohérents relativement à
$\{X \to Y\}$. D'après le lemme \ref{descent-lemma-open-fpqc-covering}, le morphisme
$X \to Y$ est universellement submersif. On peut donc appliquer le lemme
\ref{descent-lemma-universal-effective-epimorphism} pour voir que $\{X \to Y\}$ est un
épimorphisme effectif universel.

\medskip\noindent
Deuxième démonstration. Soit $R \to A$ l'homomorphisme d'anneaux fidèlement
plat correspondant à notre morphisme plat surjectif $\pi : X \to Y$.
Soit $f : X \to S$ un morphisme tel que
$f \circ \text{pr}_1 = f \circ \text{pr}_2$ comme morphismes
$X \times_Y X = \Spec(A \otimes_R A) \to S$.
D'après le lemme \ref{descent-lemma-equiv-fibre-product}, on voit que, comme
application entre ensembles sous-jacents, $f$ est de la forme
$f = g \circ \pi$ pour une application ensembliste
$g : \Spec(R) \to S$. D'après Morphismes, lemme
\ref{morphisms-lemma-fpqc-quotient-topology}, et puisque $f$ est continue,
l'application $g$ est continue.

\medskip\noindent
Fixons $y \in Y = \Spec(R)$. Choisissons un ouvert affine $U \subset S$
contenant $g(y)$, disons $U = \Spec(B)$.
D'après ce qui précède, on peut choisir $r \in R$ tel que
$y \in D(r) \subset g^{-1}(U)$.
La restriction de $f$ à $\pi^{-1}(D(r))$, à valeurs dans $U$, correspond à un
homomorphisme d'anneaux $B \to A_r$. Les deux homomorphismes d'anneaux induits
$B \to A_r \otimes_{R_r} A_r = (A \otimes_R A)_r$ sont égaux par hypothèse
sur $f$. Remarquons que $R_r \to A_r$ est fidèlement plat.
D'après le lemme \ref{descent-lemma-ff-exact}, l'égalisateur des deux flèches
$A_r \to A_r \otimes_{R_r} A_r$ est $R_r$.
On en conclut que $B \to A_r$ se factorise de manière unique par un
homomorphisme $B \to R_r$. Celui-ci donne à son tour un morphisme de schémas
$D(r) \to U \to S$, voir Schémas, lemme
\ref{schemes-lemma-morphism-into-affine}.

\medskip\noindent
Qu'avons-nous démontré jusqu'ici ? Nous avons montré que, pour tout idéal
premier $\mathfrak p \subset R$, il existe un ouvert affine standard
$D(r) \subset \Spec(R)$ tel que le morphisme
$f|_{\pi^{-1}(D(r))} : \pi^{-1}(D(r)) \to S$ se factorise de manière unique
par un certain morphisme de schémas $D(r) \to S$.
Nous omettons la vérification que ces morphismes se recollent pour donner le
morphisme voulu $\Spec(R) \to S$.
\end{proof}

\begin{lemma}
\label{descent-lemma-coequalizer-fpqc-local}
Considérons des schémas $X, Y, Z$, des morphismes $a, b : X \to Y$ et un
morphisme $c : Y \to Z$ tel que $c \circ a = c \circ b$. Posons
$d = c \circ a = c \circ b$. S'il existe un recouvrement fpqc
$\{Z_i \to Z\}$ tel que
\begin{enumerate}
\item pour tout $i$, le morphisme $Y \times_{c, Z} Z_i \to Z_i$ soit le
coégalisateur de $(a, 1) : X \times_{d, Z} Z_i \to Y \times_{c, Z} Z_i$ et de
$(b, 1) : X \times_{d, Z} Z_i \to Y \times_{c, Z} Z_i$, et
\item pour tous $i$ et $i'$, le morphisme
$Y \times_{c, Z} (Z_i \times_Z Z_{i'}) \to (Z_i \times_Z Z_{i'})$ soit le
coégalisateur de
$(a, 1) : X \times_{d, Z} (Z_i \times_Z Z_{i'}) \to
Y \times_{c, Z} (Z_i \times_Z Z_{i'})$ et de
$(b, 1) : X \times_{d, Z} (Z_i \times_Z Z_{i'}) \to
Y \times_{c, Z} (Z_i \times_Z Z_{i'})$,
\end{enumerate}
alors $c$ est le coégalisateur de $a$ et $b$.
\end{lemma}

\begin{proof}
En effet, pour un schéma $T$, la donnée d'un morphisme $Z \to T$ équivaut à
celle d'une famille de morphismes $Z_i \to T$ qui coïncident sur les
intersections, d'après le lemme
\ref{descent-lemma-fpqc-universal-effective-epimorphisms}.
\end{proof}















\section[Descente : finitude et lissité]{Descente des propriétés de finitude et de lissité des morphismes}
\label{descent-section-descent-finiteness-morphisms}

\noindent
Dans cette section, nous montrons que plusieurs propriétés des morphismes
(être lisse, être localement de présentation finie, etc.) se descendent bien
par des morphismes fidèlement plats. Nous commençons par une version algébrique.
(Le lecteur ``noethérien'' devrait consulter le lemme
\ref{descent-lemma-finite-type-local-source-fppf-algebra} plutôt que le lemme suivant.)

\begin{lemma}
\label{descent-lemma-flat-finitely-presented-permanence-algebra}
Soient $R \to A \to B$ des homomorphismes d'anneaux.
Supposons que $R \to B$ soit de présentation finie et que $A \to B$ soit
fidèlement plat et de présentation finie. Alors $R \to A$ est de présentation
finie.
\end{lemma}

\begin{proof}
Considérons l'algèbre $C = B \otimes_A B$ munie des deux homomorphismes
$p, q : B \to C$ donnés par $p(b) = b \otimes 1$ et $q(b) = 1 \otimes b$.
Bien entendu, les deux composés $A \to B \to C$ sont égaux.
Comme $p : B \to C$ est plat et de présentation finie (c'est le changement de
base de $A \to B$), l'homomorphisme d'anneaux $R \to C$ est de présentation
finie (comme composé de $R \to B \to C$).

\medskip\noindent
Nous allons utiliser le critère d'Algèbre, lemme
\ref{algebra-lemma-characterize-finite-presentation}, pour montrer que
$R \to A$ est de présentation finie.
Soit $S$ une $R$-algèbre quelconque et supposons que
$S = \colim_{\lambda \in \Lambda} S_\lambda$ soit écrite comme une colimite
filtrante de $R$-algèbres. Soit $A \to S$ un homomorphisme de $R$-algèbres.
Il faut montrer que $A \to S$ se factorise par l'une des $S_\lambda$.
Considérons les anneaux $B' = S \otimes_A B$ et
$C' = S \otimes_A C = B' \otimes_S B'$.
Puisque $B$ est fidèlement plat et de présentation finie sur $A$, l'anneau
$B'$ est lui aussi fidèlement plat et de présentation finie sur $S$.
D'après Algèbre, lemme
\ref{algebra-lemma-flat-finite-presentation-limit-flat}, partie (2), appliqué
au couple $(S \to B', B')$ et au système $(S_\lambda)$, il existe un
$\lambda_0 \in \Lambda$ et une $S_{\lambda_0}$-algèbre plate et de présentation
finie $B_{\lambda_0}$ telle que
$B' = S \otimes_{S_{\lambda_0}} B_{\lambda_0}$.
Pour $\lambda \geq \lambda_0$, posons
$B_\lambda = S_\lambda \otimes_{S_{\lambda_0}} B_{\lambda_0}$ et
$C_\lambda = B_\lambda \otimes_{S_\lambda} B_\lambda$.

\medskip\noindent
Nous interrompons le fil de l'argument pour montrer que
$S_\lambda \to B_\lambda$ est fidèlement plat pour $\lambda$ assez grand.
(Cela devrait en réalité faire l'objet d'un lemme distinct ailleurs, peut-être
dans le chapitre sur les limites.)
Comme $\Spec(B_{\lambda_0}) \to \Spec(S_{\lambda_0})$ est plat et de
présentation finie, il est ouvert (voir Morphismes, lemme
\ref{morphisms-lemma-fppf-open}).
Soit $I \subset S_{\lambda_0}$ un idéal tel que
$V(I) \subset \Spec(S_{\lambda_0})$ soit le complémentaire de l'image.
Remarquons que la formation de l'image commute au changement de base.
Puisque $\Spec(B') \to \Spec(S)$ est surjectif et que
$B' = B_{\lambda_0} \otimes_{S_{\lambda_0}} S$, on a $IS = S$.
Il existe donc $\lambda \geq \lambda_0$ tel que
$IS_{\lambda} = S_\lambda$. Pour cet indice et tous les indices
$\lambda$ plus grands, le morphisme
$\Spec(B_\lambda) \to \Spec(S_\lambda)$ est surjectif.

\medskip\noindent
Par analogie avec la notation du premier paragraphe de la démonstration,
notons $p_\lambda, q_\lambda : B_\lambda \to C_\lambda$ les deux
homomorphismes canoniques. Alors
$B' = \colim_{\lambda \geq \lambda_0} B_\lambda$ et
$C' = \colim_{\lambda \geq \lambda_0} C_\lambda$.
Comme $B$ et $C$ sont de présentation finie sur $R$, il existe
(en appliquant plusieurs fois Algèbre, lemme
\ref{algebra-lemma-characterize-finite-presentation}) un
$\lambda \geq \lambda_0$ et des homomorphismes de $R$-algèbres
$B \to B_\lambda$, $C \to C_\lambda$ tels que le diagramme
$$
\xymatrix{
C \ar[rr] & &
C_\lambda \\
B \ar[rr]
\ar@<1ex>[u]^-p
\ar@<-1ex>[u]_-q
& &
B_\lambda
\ar@<1ex>[u]^-{p_\lambda}
\ar@<-1ex>[u]_-{q_\lambda}
}
$$
soit commutatif. Cela signifie que $A \to B \to B_\lambda$ prend ses valeurs
dans l'égalisateur de $p_\lambda$ et $q_\lambda$.
D'après le lemme \ref{descent-lemma-ff-exact}, $S_\lambda$ est l'égalisateur de
$p_\lambda$ et $q_\lambda$. On obtient ainsi l'homomorphisme d'anneaux voulu
$A \to S_\lambda$, ce qui conclut.
\end{proof}

\noindent
Voici une version plus simple qui traite de la propriété d'être de type fini.

\begin{lemma}
\label{descent-lemma-finite-type-local-source-fppf-algebra}
Soient $R \to A \to B$ des homomorphismes d'anneaux.
Supposons que $R \to B$ soit de type fini et que $A \to B$ soit fidèlement plat
et de présentation finie. Alors $R \to A$ est de type fini.
\end{lemma}

\begin{proof}
D'après Algèbre, lemme
\ref{algebra-lemma-descend-faithfully-flat-finite-presentation}, il existe un
diagramme commutatif
$$
\xymatrix{
R \ar[r] \ar@{=}[d] &
A_0 \ar[d] \ar[r] &
B_0 \ar[d] \\
R \ar[r] & A \ar[r] & B
}
$$
où $R \to A_0$ est de présentation finie, $A_0 \to B_0$ est fidèlement plat
et de présentation finie, et $B = A \otimes_{A_0} B_0$.
Puisque $R \to B$ est de type fini par hypothèse, on peut adjoindre quelques
éléments à $A_0$ et supposer que l'homomorphisme $B_0 \to B$ est surjectif.
Dans ce cas, comme $A_0 \to B_0$ est fidèlement plat et que
$$
(A_0 \to A) \otimes_{A_0} B_0 \cong (B_0 \to B)
$$
est surjectif, $A_0 \to A$ est lui aussi surjectif, ce qui conclut.
\end{proof}

\begin{lemma}
\label{descent-lemma-flat-finitely-presented-permanence}
\begin{reference}
\cite[IV, 17.7.5 (i) and (ii)]{EGA}.
\end{reference}
Soit
$$
\xymatrix{
X \ar[rr]_f \ar[rd]_p & &
Y \ar[dl]^q \\
& S
}
$$
un diagramme commutatif de morphismes de schémas. Supposons que $f$ soit
surjectif, plat et localement de présentation finie, et que $p$ soit localement
de présentation finie (resp.\ localement de type fini).
Alors $q$ est localement de présentation finie (resp.\ localement de type fini).
\end{lemma}

\begin{proof}
La question est locale sur $S$ et sur $Y$. On peut donc supposer que $S$ et
$Y$ sont affines. Puisque $f$ est plat et localement de présentation finie,
$f$ est ouvert (Morphismes, lemme \ref{morphisms-lemma-fppf-open}).
Comme $Y$ est quasi-compact, il existe donc un nombre fini d'ouverts affines
$X_i \subset X$ tels que $Y = \bigcup f(X_i)$.
On peut manifestement remplacer $X$ par $\coprod X_i$ et supposer ainsi que
$X$ est lui aussi affine. Dans ce cas, le lemme équivaut au lemme
\ref{descent-lemma-flat-finitely-presented-permanence-algebra}
(resp. au lemme \ref{descent-lemma-finite-type-local-source-fppf-algebra}) ci-dessus.
\end{proof}

\noindent
Nous allons nous en servir pour améliorer certains des résultats antérieurs
sur les morphismes.

\begin{lemma}
\label{descent-lemma-syntomic-smooth-etale-permanence}
Soit
$$
\xymatrix{
X \ar[rr]_f \ar[rd]_p & &
Y \ar[dl]^q \\
& S
}
$$
un diagramme commutatif de morphismes de schémas. Supposons que
\begin{enumerate}
\item $f$ soit surjectif et syntomique (resp.\ lisse, resp.\ \'etale),
\item $p$ soit syntomique (resp.\ lisse, resp.\ \'etale).
\end{enumerate}
Alors $q$ est syntomique (resp.\ lisse, resp.\ \'etale).
\end{lemma}

\begin{proof}
Combiner Morphismes, lemmes
\ref{morphisms-lemma-syntomic-permanence},
\ref{morphisms-lemma-smooth-permanence} et
\ref{morphisms-lemma-etale-permanence-two}, avec le lemme
\ref{descent-lemma-flat-finitely-presented-permanence} ci-dessus.
\end{proof}

\noindent
En fait, on peut renforcer ce résultat comme suit.

\begin{lemma}
\label{descent-lemma-smooth-permanence}
Soit
$$
\xymatrix{
X \ar[rr]_f \ar[rd]_p & &
Y \ar[dl]^q \\
& S
}
$$
un diagramme commutatif de morphismes de schémas. Supposons que
\begin{enumerate}
\item $f$ soit surjectif, plat et localement de présentation finie,
\item $p$ soit lisse (resp.\ \'etale).
\end{enumerate}
Alors $q$ est lisse (resp.\ \'etale).
\end{lemma}

\begin{proof}
Supposons (1) et que $p$ soit lisse. Le lemme
\ref{descent-lemma-flat-finitely-presented-permanence} montre que $q$ est localement de
présentation finie. Morphismes, lemme
\ref{morphisms-lemma-flat-permanence}, montre que $q$ est plat.
Il suffit donc de montrer que les fibres de $q$ sont lisses, voir le lemme
\ref{morphisms-lemma-smooth-flat-smooth-fibres} de Morphismes.
Appliquons le lemme \ref{varieties-lemma-flat-under-smooth} de Variétés aux
morphismes plats surjectifs $X_s \to Y_s$, pour $s \in S$, et concluons.
Nous omettons la démonstration du cas \'etale.
\end{proof}

\begin{remark}
\label{descent-remark-smooth-permanence}
Sous les hypothèses (1), avec $p$ lisse, du lemme
\ref{descent-lemma-smooth-permanence}, le morphisme $X \to Y$ n'est pas nécessairement
lisse. Un contre-exemple est donné par $S = \Spec(k)$,
$X = \Spec(k[s])$, $Y = \Spec(k[t])$ et $f$ défini par
$t \mapsto s^2$. Voir cependant le lemme
\ref{descent-lemma-syntomic-permanence} pour des renseignements sur la structure de
$f$.
\end{remark}

\begin{lemma}
\label{descent-lemma-syntomic-permanence}
Soit
$$
\xymatrix{
X \ar[rr]_f \ar[rd]_p & &
Y \ar[dl]^q \\
& S
}
$$
un diagramme commutatif de morphismes de schémas. Supposons que
\begin{enumerate}
\item $f$ soit surjectif, plat et localement de présentation finie,
\item $p$ soit syntomique.
\end{enumerate}
Alors $q$ et $f$ sont tous deux syntomiques.
\end{lemma}

\begin{proof}
Le lemme \ref{descent-lemma-flat-finitely-presented-permanence} montre que $q$
est de présentation finie. Le lemme
\ref{morphisms-lemma-flat-permanence} de Morphismes
montre que $q$ est plat.
D'après le lemme
\ref{morphisms-lemma-syntomic-locally-standard-syntomic} de Morphismes,
il suffit donc de montrer que les anneaux locaux des fibres de
$Y \to S$ et de $X \to Y$ sont des anneaux locaux d'intersection
complète. Pour cela, nous pouvons prendre la fibre de
$X \to Y \to S$ en un point $s \in S$ ; autrement dit, nous pouvons
supposer que $S$ est le spectre d'un corps. Choisissons un point
$x \in X$, d'image $y \in Y$, et considérons l'homomorphisme d'anneaux
$$
\mathcal{O}_{Y, y} \longrightarrow \mathcal{O}_{X, x}
$$
C'est un homomorphisme local plat d'anneaux locaux noethériens.
L'anneau local $\mathcal{O}_{X, x}$ est d'intersection complète.
Nous pouvons donc appliquer le résultat d'Avramov, voir le lemme
\ref{dpa-lemma-avramov} d'Algèbre à puissances divisées, et conclure que
$\mathcal{O}_{Y, y}$ et
$\mathcal{O}_{X, x}/\mathfrak m_y\mathcal{O}_{X, x}$ sont tous deux
des anneaux d'intersection complète.
\end{proof}

\noindent
Le type de lemme suivant est parfois utile.

\begin{lemma}
\label{descent-lemma-curiosity}
Soient $X \to Y \to Z$ des morphismes de schémas.
Soit $P$ l'une des propriétés suivantes des morphismes de schémas :
être plat, être localement de type fini, être localement de présentation
finie. Supposons que $X \to Z$ possède la propriété $P$ et que
$\{X \to Y\}$ puisse être raffiné par un recouvrement fppf de $Y$.
Alors $Y \to Z$ possède la propriété $P$.
\end{lemma}

\begin{proof}
Soient $\Spec(C) \subset Z$ un ouvert affine et
$\Spec(B) \subset Y$ un ouvert affine dont l'image est contenue dans
$\Spec(C)$. L'hypothèse sur $X \to Y$ permet de trouver un recouvrement
fppf affine standard $\{\Spec(B_j) \to \Spec(B)\}$ et des relèvements
$x_j : \Spec(B_j) \to X$. Comme $\Spec(B_j)$ est quasi-compact, nous
pouvons trouver un nombre fini d'ouverts affines
$\Spec(A_i) \subset X$ au-dessus de $\Spec(B)$ tels que l'image de
chaque $x_j$ soit contenue dans la réunion $\bigcup \Spec(A_i)$.
Après avoir remplacé chaque $\Spec(B_j)$ par un recouvrement de Zariski
affine standard, nous pouvons donc supposer que nous disposons d'un
recouvrement fppf affine standard
$\{\Spec(B_i) \to \Spec(B)\}$ tel que chaque morphisme
$\Spec(B_i) \to Y$ se factorise par un ouvert affine
$\Spec(A_i) \subset X$ situé au-dessus de $\Spec(B)$. Autrement dit,
nous avons des homomorphismes d'anneaux $C \to B \to A_i \to B_i$
pour chaque $i$. Nous pouvons aussi considérer
$$
C \to B \to A = \prod A_i \to B' = \prod B_i
$$
et l'homomorphisme d'anneaux $B \to \prod B_i$ est fidèlement plat et
de présentation finie.

\medskip\noindent
Le cas $P = \text{plat}$. Dans ce cas, nous savons que $C \to A$ est plat et
nous devons montrer que $C \to B$ est plat. Soit
$N \to N' \to N''$ une suite exacte de $C$-modules. Nous voulons
montrer que la suite
$N \otimes_C B \to N' \otimes_C B \to N'' \otimes_C B$ est exacte.
Notons $H$ sa cohomologie et $H'$ celle de la suite
$N \otimes_C B' \to N' \otimes_C B' \to N'' \otimes_C B'$.
Comme $B \to B'$ est plat, nous avons $H' = H \otimes_B B'$.
D'autre part, la suite
$N \otimes_C A \to N' \otimes_C A \to N'' \otimes_C A$ est exacte,
donc d'homologie nulle. Le morphisme $H \to H'$ est par conséquent
nul, puisqu'il se factorise par le module nul. Ainsi $H' = 0$.
Puisque $B \to B'$ est fidèlement plat, nous concluons que $H = 0$,
comme souhaité.

\medskip\noindent
Le cas $P = \text{localement de type fini}$. Dans ce cas, nous savons que
$C \to A$ est de type fini et nous devons montrer que $C \to B$ est de
type fini. Puisque $B \to B'$ est de présentation finie (donc de type
fini), l'homomorphisme $A \to B'$ est de type fini ; voir le lemme
\ref{algebra-lemma-compose-finite-type} d'Algèbre. Par conséquent,
$C \to B'$ est de type fini, et nous concluons par le
Lemme \ref{descent-lemma-finite-type-local-source-fppf-algebra}.

\medskip\noindent
Le cas $P = \text{localement de présentation finie}$. Dans ce cas, nous savons
que $C \to A$ est de présentation finie et nous devons montrer que
$C \to B$ est de présentation finie. Puisque $B \to B'$ est de
présentation finie et que $B \to A$ est de type fini, l'homomorphisme
$A \to B'$ est de présentation finie ; voir le lemme
\ref{algebra-lemma-compose-finite-type} d'Algèbre. Par conséquent,
$C \to B'$ est de présentation finie, et nous concluons par le
Lemme \ref{descent-lemma-flat-finitely-presented-permanence-algebra}.
\end{proof}










\section{Propriétés locales des schémas}
\label{descent-section-descending-properties}

\noindent
Il arrive souvent que l'on puisse établir qu'une certaine propriété est
vérifiée par les membres d'un recouvrement d'un schéma. Dans bien des cas,
le schéma possède alors lui aussi cette propriété. Par exemple, si $S$ est
un schéma et si $f : S' \to S$ est un morphisme plat surjectif tel que $S'$
soit réduit, alors $S$ est réduit. On peut le démontrer en considérant les
anneaux locaux et en appliquant le lemme
\ref{algebra-lemma-descent-reduced} d'Algèbre. On dit que la propriété d'être
réduit {\it se descend par les morphismes plats surjectifs}. Des résultats de
ce type sont réunis dans Algèbre, section
\ref{algebra-section-descending-properties}, et, pour les schémas, dans la
section \ref{descent-section-variants}. Des résultats analogues sur la descente des
propriétés de morphismes figurent dans la
section \ref{descent-section-descent-finiteness-morphisms}.

\medskip\noindent
En revanche, il existe des morphismes plats surjectifs $f : S' \to S$ pour
lesquels $S$ est réduit mais $S'$ ne l'est pas, par exemple
$\Spec(k[x]/(x^2)) \to \Spec(k)$. La propriété d'être réduit n'est donc pas
{\it préservée par les morphismes plats}. La propriété d'avoir des corps
résiduels infinis est préservée par les morphismes plats (mais ne se descend
évidemment pas par les morphismes plats surjectifs). Des résultats de ce
type sont réunis dans Algèbre, section
\ref{algebra-section-ascending-properties}.

\medskip\noindent
Enfin, on dit qu'une propriété est {\it locale pour la topologie plate}
si elle est préservée par les morphismes plats et se descend par les
morphismes plats surjectifs. Un exemple assez élémentaire est la propriété d'avoir des corps
résiduels d'une caractéristique donnée. Afin de préciser cette notion et de
la relier aux diverses topologies sur les schémas, posons la définition
formelle suivante.

\begin{definition}
\label{descent-definition-property-local}
Soit $\mathcal{P}$ une propriété des schémas. Soit
$\tau \in \{fpqc, \linebreak[0] fppf, \linebreak[0] syntomic, \linebreak[0]
smooth, \linebreak[0] \etale, \linebreak[0] Zariski\}$. On dit que
$\mathcal{P}$ est {\it locale pour la topologie $\tau$} si, pour tout
recouvrement pour la topologie $\tau$, $\{S_i \to S\}_{i \in I}$ (voir
Topologies, section \ref{topologies-section-procedure}), on a
$$
S \text{ vérifie }\mathcal{P}
\Leftrightarrow
\text{chaque }S_i \text{ vérifie }\mathcal{P}.
$$
\end{definition}

\noindent
Bien entendu, puisque les isomorphismes sont toujours des recouvrements,
on constate (ou l'on exige) que la propriété $\mathcal{P}$ est vérifiée par
$S$ si et seulement si elle est vérifiée par tout schéma $S'$ isomorphe à
$S$. En fait, si
$\tau = fpqc, \linebreak[0] fppf, \linebreak[0] syntomic,
\linebreak[0] smooth, \linebreak[0] \etale$ ou $Zariski$, si $S$ vérifie
$\mathcal{P}$ et si $S' \to S$ est respectivement plat, plat et localement
de présentation finie, syntomique, lisse, \'etale, ou une immersion ouverte,
alors $S'$ vérifie $\mathcal{P}$. En effet, on peut toujours compléter
$\{S' \to S\}$ en un recouvrement pour la topologie $\tau$.

\medskip\noindent
Nous avons les implications suivantes :
$\mathcal{P}$ est locale pour la topologie fpqc
$\Rightarrow$
$\mathcal{P}$ est locale pour la topologie fppf
$\Rightarrow$
$\mathcal{P}$ est locale pour la topologie syntomique
$\Rightarrow$
$\mathcal{P}$ est locale pour la topologie lisse
$\Rightarrow$
$\mathcal{P}$ est locale pour la topologie \'etale
$\Rightarrow$
$\mathcal{P}$ est locale pour la topologie de Zariski.
Cela résulte des lemmes
\ref{topologies-lemma-zariski-etale},
\ref{topologies-lemma-zariski-etale-smooth},
\ref{topologies-lemma-zariski-etale-smooth-syntomic},
\ref{topologies-lemma-zariski-etale-smooth-syntomic-fppf} et
\ref{topologies-lemma-zariski-etale-smooth-syntomic-fppf-fpqc} de Topologies.

\begin{lemma}
\label{descent-lemma-descending-properties}
Soit $\mathcal{P}$ une propriété des schémas.
Soit $\tau \in \{fpqc, \linebreak[0] fppf, \linebreak[0]
\etale, \linebreak[0] smooth, \linebreak[0] syntomic\}$. Supposons que
\begin{enumerate}
\item la propriété soit locale pour la topologie de Zariski,
\item pour tout morphisme de schémas affines $S' \to S$ qui est plat,
plat de présentation finie, \'etale, lisse ou syntomique selon que $\tau$
vaut fpqc, fppf, \'etale, smooth ou syntomic, la propriété $\mathcal{P}$
soit vérifiée par $S'$ si la propriété $\mathcal{P}$ est vérifiée par $S$, et
\item pour tout morphisme surjectif de schémas affines $S' \to S$ qui est
plat, plat de présentation finie, \'etale, lisse ou syntomique selon que
$\tau$ vaut fpqc, fppf, \'etale, smooth ou syntomic, la propriété
$\mathcal{P}$ soit vérifiée par $S$ si la propriété $\mathcal{P}$ est
vérifiée par $S'$.
\end{enumerate}
Alors $\mathcal{P}$ est locale sur la base pour la topologie $\tau$.
\end{lemma}

\begin{proof}
Cela résulte presque immédiatement de la définition d'un recouvrement pour
la topologie $\tau$ ; voir les définitions
\ref{topologies-definition-fpqc-covering},
\ref{topologies-definition-fppf-covering},
\ref{topologies-definition-etale-covering},
\ref{topologies-definition-smooth-covering} ou
\ref{topologies-definition-syntomic-covering} de Topologies, et les lemmes
\ref{topologies-lemma-fpqc-affine},
\ref{topologies-lemma-fppf-affine},
\ref{topologies-lemma-etale-affine},
\ref{topologies-lemma-smooth-affine} ou
\ref{topologies-lemma-syntomic-affine} de Topologies. Nous omettons les détails.
\end{proof}

\begin{remark}
\label{descent-remark-descending-properties-standard}
Dans le lemme \ref{descent-lemma-descending-properties}, si $\tau = smooth$, on
peut supposer, dans la condition (3), que le morphisme est un morphisme
lisse standard (surjectif). Il en va de même lorsque $\tau = syntomic$ ou
$\tau = \etale$.
\end{remark}





\section{Propriétés des schémas locales pour la topologie fppf}
\label{descent-section-descending-properties-fppf}

\noindent
Dans cette section, nous donnons quelques propriétés de schémas qui sont
locales sur la base pour la topologie fppf.

\begin{lemma}
\label{descent-lemma-Noetherian-local-fppf}
La propriété $\mathcal{P}(S) =$``$S$ est localement noethérien'' est locale
pour la topologie fppf.
\end{lemma}

\begin{proof}
Nous allons appliquer le lemme \ref{descent-lemma-descending-properties}.
Tout d'abord, la propriété ``être localement noethérien'' est locale pour la
topologie de Zariski. Cela résulte immédiatement de la définition ; voir
Propriétés, définition \ref{properties-definition-noetherian}.
Ensuite, si $S' \to S$ est un morphisme plat de présentation finie entre
schémas affines et si $S$ est localement noethérien, alors $S'$ est
localement noethérien ; c'est le lemme
\ref{morphisms-lemma-finite-type-noetherian} de Morphismes. Enfin, si $S' \to S$ est un
morphisme plat, surjectif et de présentation finie entre schémas affines et
si $S'$ est localement noethérien, alors $S$ est localement noethérien ; cela
résulte du lemme \ref{algebra-lemma-descent-Noetherian} d'Algèbre.
Les conditions (1), (2) et (3) du lemme
\ref{descent-lemma-descending-properties} sont donc satisfaites, ce qui conclut.
\end{proof}

\begin{lemma}
\label{descent-lemma-Jacobson-local-fppf}
La propriété $\mathcal{P}(S) =$``$S$ est jacobsonien'' est locale pour la
topologie fppf.
\end{lemma}

\begin{proof}
Nous allons appliquer le Lemme \ref{descent-lemma-descending-properties}.
Tout d'abord, la propriété ``être jacobsonien'' est locale pour la topologie
de Zariski ; c'est Propriétés, Lemme
\ref{properties-lemma-locally-jacobson}. Ensuite, si $S' \to S$ est un
morphisme plat de présentation finie entre schémas affines et si $S$ est
jacobsonien, alors $S'$ est jacobsonien ; c'est Morphismes, Lemme
\ref{morphisms-lemma-Jacobson-universally-Jacobson}. Enfin, supposons que
$f : S' \to S$ soit un morphisme plat, surjectif et de présentation finie
entre schémas affines, et que $S'$ soit jacobsonien ; nous devons montrer
que $S$ est jacobsonien. Écrivons $S = \Spec(A)$ et $S' = \Spec(B)$, le
morphisme $S' \to S$ étant donné par $A \to B$. Alors $A \to B$ est de
présentation finie et fidèlement plat. De plus, l'anneau $B$ est jacobsonien ;
voir Propriétés, Lemme \ref{properties-lemma-locally-jacobson}.

\medskip\noindent
D'après Algèbre, Lemme \ref{algebra-lemma-fppf-fpqf}, il existe un diagramme
$$
\xymatrix{
B \ar[rr] & & B' \\
& A \ar[ru] \ar[lu] &
}
$$
tel que $A \to B'$ soit de présentation finie, fidèlement plat et quasi-fini.
En particulier, $B \to B'$ est de type fini et Algèbre, Proposition
\ref{algebra-proposition-Jacobson-permanence}, montre que $B'$ est
jacobsonien. Nous pouvons donc supposer que $A \to B$ est quasi-fini, en plus
d'être fidèlement plat et de présentation finie.

\medskip\noindent
Supposons que $A$ ne soit pas jacobsonien et cherchons une contradiction.
D'après Algèbre, Lemme \ref{algebra-lemma-characterize-jacobson}, il existe
un idéal premier non maximal $\mathfrak p \subset A$ et un élément $f \in A$,
$f \not \in \mathfrak p$, tels que
$V(\mathfrak p) \cap D(f) = \{\mathfrak p\}$.

\medskip\noindent
On obtient une contradiction comme suit. Choisissons d'abord un idéal maximal
$\mathfrak m$ de $A$ contenant $\mathfrak p$. Choisissons un idéal premier
$\mathfrak m' \subset B$ au-dessus de $\mathfrak m$ (il en existe un parce
que $A \to B$ est fidèlement plat ; voir Algèbre, Lemme
\ref{algebra-lemma-ff-rings}). Comme $A \to B$ est plat, le théorème de
descente, voir Algèbre, Lemme \ref{algebra-lemma-flat-going-down}, fournit
un idéal premier $\mathfrak q \subset \mathfrak m'$ au-dessus de
$\mathfrak p$. En particulier, $\mathfrak q$ n'est pas maximal. Une nouvelle
application d'Algèbre, Lemme \ref{algebra-lemma-characterize-jacobson},
montre donc que l'ensemble $V(\mathfrak q) \cap D(f)$ est infini (c'est ici
que nous utilisons enfin le fait que $B$ est jacobsonien). Tous les points de
$V(\mathfrak q) \cap D(f)$ ont pour image
$V(\mathfrak p) \cap D(f) = \{\mathfrak p\}$. La fibre au-dessus de
$\mathfrak p$ est donc infinie. Cela contredit le fait que $A \to B$ est
quasi-fini ; voir Algèbre, Lemme \ref{algebra-lemma-quasi-finite}, ou, plus
explicitement, Morphismes, Lemme \ref{morphisms-lemma-quasi-finite}.
Le lemme est ainsi démontré.
\end{proof}

\begin{lemma}
\label{descent-lemma-locally-finite-nr-irred-local-fppf}
La propriété $\mathcal{P}(S) =$``tout ouvert quasi-compact de $S$ possède
un nombre fini de composantes irréductibles'' est locale pour la topologie
fppf.
\end{lemma}

\begin{proof}
Nous allons appliquer le lemme \ref{descent-lemma-descending-properties}.
Tout d'abord, $\mathcal{P}$ est locale pour la topologie de Zariski.
Ensuite, supposons que $T \to S$ soit un morphisme plat de présentation
finie entre schémas affines et que $S$ possède un nombre fini de composantes
irréductibles ; il en va alors de même de $T$. En effet, comme $T \to S$ est
plat, les points génériques de $T$ s'envoient sur les points génériques de
$S$ ; voir Morphismes, Lemme
\ref{morphisms-lemma-generalizations-lift-flat}. Il suffit donc de montrer
que, pour $s \in S$, la fibre $T_s$ possède un nombre fini de points
génériques. Or $T_s$ est un schéma affine de type fini sur $\kappa(s)$ ;
voir Morphismes, Lemme \ref{morphisms-lemma-base-change-finite-type}.
Ainsi, $T_s$ est noethérien et possède un nombre fini de composantes
irréductibles (Morphismes, Lemme
\ref{morphisms-lemma-finite-type-noetherian}, et Propriétés, Lemme
\ref{properties-lemma-Noetherian-irreducible-components}). Enfin, si
$T \to S$ est un morphisme plat, surjectif et de présentation finie entre
schémas affines et si $T$ possède un nombre fini de composantes
irréductibles, il en va de même de $S$. Cela résulte de Topologie, Lemme
\ref{topology-lemma-surjective-continuous-irreducible-components}.
Les conditions (1), (2) et (3) du lemme
\ref{descent-lemma-descending-properties} sont donc satisfaites, ce qui conclut.
\end{proof}




\section{Propriétés des schémas locales pour la topologie syntomique}
\label{descent-section-descending-properties-syntomic}

\noindent
Dans cette section, nous donnons quelques propriétés de schémas qui sont
locales sur la base pour la topologie syntomique.

\begin{lemma}
\label{descent-lemma-Sk-local-syntomic}
La propriété $\mathcal{P}(S) =$``$S$ est localement noethérien et vérifie
$(S_k)$'' est locale pour la topologie syntomique.
\end{lemma}

\begin{proof}
Nous allons vérifier les conditions (1), (2) et (3) du lemme
\ref{descent-lemma-descending-properties}. Comme un morphisme syntomique est plat de
présentation finie (Morphismes, Lemmes
\ref{morphisms-lemma-syntomic-flat} et
\ref{morphisms-lemma-syntomic-locally-finite-presentation}), la propriété
``être localement noethérien'' a déjà été traitée dans la preuve du
lemme \ref{descent-lemma-Noetherian-local-fppf}. Nous l'utiliserons sans autre
mention. Tout d'abord, $\mathcal{P}$ est locale pour la topologie de Zariski.
Cela résulte immédiatement de la définition ; voir Cohomologie des schémas,
Définition \ref{coherent-definition-depth}. Ensuite, si $S' \to S$ est un
morphisme syntomique entre schémas affines, si $S$ vérifie $\mathcal{P}$,
alors $S'$ vérifie $\mathcal{P}$. C'est Algèbre, Lemme
\ref{algebra-lemma-Sk-goes-up} (appliquer Morphismes, Lemme
\ref{morphisms-lemma-syntomic-characterize}, ainsi qu'Algèbre, Définition
\ref{algebra-definition-lci} et Lemme \ref{algebra-lemma-lci-CM}).
Enfin, si $S' \to S$ est un morphisme syntomique surjectif entre schémas
affines et si $S'$ vérifie $\mathcal{P}$, alors $S$ vérifie $\mathcal{P}$ ;
c'est Algèbre, Lemme \ref{algebra-lemma-descent-Sk}.
Les conditions (1), (2) et (3) du lemme
\ref{descent-lemma-descending-properties} sont donc satisfaites, ce qui conclut.
\end{proof}

\begin{lemma}
\label{descent-lemma-CM-local-syntomic}
La propriété $\mathcal{P}(S) =$``$S$ est de Cohen-Macaulay'' est locale pour
la topologie syntomique.
\end{lemma}

\begin{proof}
Cela résulte du lemme \ref{descent-lemma-Sk-local-syntomic}, puisqu'un schéma est de
Cohen-Macaulay si et seulement s'il est localement noethérien et vérifie
$(S_k)$ pour tout $k \geq 0$ ; voir Propriétés, Lemme
\ref{properties-lemma-scheme-CM-iff-all-Sk}.
\end{proof}







\section{Propriétés des schémas locales pour la topologie lisse}
\label{descent-section-descending-properties-smooth}

\noindent
Dans cette section, nous donnons quelques propriétés de schémas qui sont
locales sur la base pour la topologie lisse.

\begin{lemma}
\label{descent-lemma-reduced-local-smooth}
La propriété $\mathcal{P}(S) =$``$S$ est réduit'' est locale pour la topologie
lisse.
\end{lemma}

\begin{proof}
Nous allons appliquer le lemme \ref{descent-lemma-descending-properties}.
Tout d'abord, la propriété ``être réduit'' est locale pour la topologie de
Zariski. Cela résulte immédiatement de la définition ; voir Schémas,
Définition \ref{schemes-definition-reduced}. Ensuite, si $S' \to S$ est un
morphisme lisse entre schémas affines et si $S$ est réduit, alors $S'$ est
réduit ; c'est Algèbre, Lemme \ref{algebra-lemma-reduced-goes-up}.
Enfin, si $S' \to S$ est un morphisme lisse surjectif entre schémas affines
et si $S'$ est réduit, alors $S$ est réduit ; c'est Algèbre, Lemme
\ref{algebra-lemma-descent-reduced}. Les conditions (1), (2) et (3) du
lemme \ref{descent-lemma-descending-properties} sont donc satisfaites, ce qui conclut.
\end{proof}

\begin{lemma}
\label{descent-lemma-normal-local-smooth}
\begin{slogan}
La normalité est locale pour la topologie lisse.
\end{slogan}
La propriété $\mathcal{P}(S) =$``$S$ est normal'' est locale pour la topologie
lisse.
\end{lemma}

\begin{proof}
Nous allons appliquer le lemme \ref{descent-lemma-descending-properties}.
Tout d'abord, la propriété ``être normal'' est locale pour la topologie de
Zariski. Cela résulte immédiatement de la définition ; voir Propriétés,
Définition \ref{properties-definition-normal}. Ensuite, si $S' \to S$ est
un morphisme lisse entre schémas affines et si $S$ est normal, alors $S'$ est
normal ; c'est Algèbre, Lemme \ref{algebra-lemma-normal-goes-up}. Enfin, si
$S' \to S$ est un morphisme lisse surjectif entre schémas affines et si $S'$
est normal, alors $S$ est normal ; c'est Algèbre, Lemme
\ref{algebra-lemma-descent-normal}. Les conditions (1), (2) et (3) du
lemme \ref{descent-lemma-descending-properties} sont donc satisfaites, ce qui conclut.
\end{proof}

\begin{lemma}
\label{descent-lemma-Rk-local-smooth}
La propriété $\mathcal{P}(S) =$``$S$ est localement noethérien et vérifie
$(R_k)$'' est locale pour la topologie lisse.
\end{lemma}

\begin{proof}
Nous allons vérifier les conditions (1), (2) et (3) du lemme
\ref{descent-lemma-descending-properties}. Comme un morphisme lisse est plat et de
présentation finie (voir les lemmes \ref{morphisms-lemma-smooth-flat} et
\ref{morphisms-lemma-smooth-locally-finite-presentation} de Morphismes), la propriété
``être localement noethérien'' a déjà été traitée dans la preuve du
lemme \ref{descent-lemma-Noetherian-local-fppf}. Nous l'utiliserons sans autre
mention. Tout d'abord, $\mathcal{P}$ est locale pour la topologie de Zariski.
Cela résulte immédiatement de la définition \ref{properties-definition-Rk}
de Propriétés. Ensuite, si $S' \to S$ est un morphisme
lisse entre schémas affines et si $S$ vérifie $\mathcal{P}$, alors $S'$
vérifie $\mathcal{P}$. Cela résulte du lemme
\ref{algebra-lemma-Rk-goes-up} d'Algèbre (où l'on utilise le lemme
\ref{morphisms-lemma-smooth-characterize} de Morphismes, ainsi que les lemmes
\ref{algebra-lemma-base-change-smooth} et
\ref{algebra-lemma-characterize-smooth-over-field} d'Algèbre). Enfin, si $S' \to S$
est un morphisme lisse surjectif entre schémas affines et si $S'$ vérifie
$\mathcal{P}$, alors $S$ vérifie $\mathcal{P}$ ; c'est Algèbre, Lemme
\ref{algebra-lemma-descent-Rk}. Les conditions (1), (2) et (3) du
lemme \ref{descent-lemma-descending-properties} sont donc satisfaites, ce qui conclut.
\end{proof}

\begin{lemma}
\label{descent-lemma-regular-local-smooth}
La propriété $\mathcal{P}(S) =$``$S$ est régulier'' est locale pour la
topologie lisse.
\end{lemma}

\begin{proof}
Cela résulte du lemme \ref{descent-lemma-Rk-local-smooth}, puisqu'un schéma localement
noethérien est régulier si et seulement s'il est localement noethérien et
vérifie $(R_k)$ pour tout $k \geq 0$.
\end{proof}

\begin{lemma}
\label{descent-lemma-Nagata-local-smooth}
La propriété $\mathcal{P}(S) =$``$S$ est un schéma de Nagata'' est locale pour
la topologie lisse.
\end{lemma}

\begin{proof}
Nous allons vérifier les conditions (1), (2) et (3) du lemme
\ref{descent-lemma-descending-properties}. Tout d'abord, la propriété d'être un schéma
de Nagata est locale pour la topologie de Zariski ; c'est le lemme
\ref{properties-lemma-locally-nagata} de Propriétés. Ensuite, si $S' \to S$ est un
morphisme lisse entre schémas affines et si $S$ est un schéma de Nagata,
alors $S'$ est un schéma de Nagata ; c'est le lemme
\ref{morphisms-lemma-finite-type-nagata} de Morphismes. Enfin, si $S' \to S$ est un
morphisme lisse surjectif entre schémas affines et si $S'$ est un schéma de
Nagata, alors $S$ est un schéma de Nagata ; c'est le lemme
\ref{algebra-lemma-descent-nagata} d'Algèbre. Les conditions (1), (2) et (3) du
lemme \ref{descent-lemma-descending-properties} sont donc satisfaites, ce qui conclut.
\end{proof}





\section{Variantes sur la descente des propriétés}
\label{descent-section-variants}

\noindent
Il est parfois possible de faire descendre des propriétés qui ne sont pas
locales. Nous réunissons dans cette section des résultats de ce type. Voir
aussi la section \ref{descent-section-descent-finiteness-morphisms}, consacrée à la
descente de propriétés de morphismes, comme la lissité.

\begin{lemma}
\label{descent-lemma-descend-reduced}
Si $f : X \to Y$ est un morphisme plat surjectif de schémas et si $X$ est
réduit, alors $Y$ est réduit.
\end{lemma}

\begin{proof}
Le résultat s'obtient en considérant les anneaux locaux (voir la définition
\ref{schemes-definition-reduced} de Schémas) et en appliquant le lemme
\ref{algebra-lemma-descent-reduced} d'Algèbre.
\end{proof}

\begin{lemma}
\label{descent-lemma-descend-regular}
Soit $f : X \to Y$ un morphisme d'espaces algébriques. Si $f$ est localement
de présentation finie, plat et surjectif, et si $X$ est régulier, alors $Y$
est régulier.
\end{lemma}

\begin{proof}
Ce lemme se ramène à l'énoncé algébrique suivant : si $A \to B$ est un
homomorphisme d'anneaux fidèlement plat et de présentation finie, et si $B$
est noethérien et régulier, alors $A$ est noethérien et régulier. Nous voyons
que $A$ est noethérien d'après le lemme
\ref{algebra-lemma-descent-Noetherian} d'Algèbre, et régulier d'après le lemme
\ref{algebra-lemma-flat-under-regular} d'Algèbre.
\end{proof}









\section{Germes de schémas}
\label{descent-section-germs}

\begin{definition}
\label{descent-definition-germs}
Germes de schémas.
\begin{enumerate}
\item Un couple $(X, x)$ formé d'un schéma $X$ et d'un point $x \in X$ est
appelé le {\it germe de $X$ en $x$}.
\item Un {\it morphisme de germes} $f : (X, x) \to (S, s)$ est une classe
d'équivalence de morphismes de schémas $f : U \to S$ tels que $f(x) = s$,
où $U \subset X$ est un voisinage ouvert de $x$. Deux tels morphismes
$f$, $f'$ sont dits équivalents si et seulement si $f$ et $f'$ coïncident
sur un voisinage ouvert de $x$.
\item On définit la {\it composition des morphismes de germes} en composant
des représentants (cette opération est bien définie).
\end{enumerate}
\end{definition}

\noindent
Avant de poursuivre, il nous faut encore une définition.

\begin{definition}
\label{descent-definition-etale-morphism-germs}
Soit $f : (X, x) \to (S, s)$ un morphisme de germes. On dit que $f$ est
{\it \'etale} (resp.\ {\it lisse}) s'il existe un représentant
$f : U \to S$ de $f$ qui soit un morphisme \'etale (resp.\ un morphisme
lisse) de schémas.
\end{definition}








\section{Propriétés locales des germes}
\label{descent-section-properties-germs-local}

\begin{definition}
\label{descent-definition-local-at-point}
Soit $\mathcal{P}$ une propriété des germes de schémas. On dit que
$\mathcal{P}$ est {\it locale pour la topologie \'etale}
(resp.\ {\it locale pour la topologie lisse}) si, pour tout morphisme de
germes \'etale (resp.\ lisse) $(U', u') \to (U, u)$, on a
$\mathcal{P}(U, u) \Leftrightarrow \mathcal{P}(U', u')$.
\end{definition}

\noindent
Soit $(X, x)$ un germe de schéma. La dimension de $X$ en $x$ est le minimum
des dimensions des voisinages ouverts de $x$ dans $X$, et tout voisinage
ouvert suffisamment petit possède cette dimension. C'est donc un invariant
de la classe d'isomorphisme du germe. On la note simplement $\dim_x(X)$.
Le lemme suivant affirme que la relation $\dim_x(X) = d$ est une propriété
locale des germes pour la topologie \'etale.

\begin{lemma}
\label{descent-lemma-dimension-at-point-local}
Soit $f : U \to V$ un morphisme \'etale de schémas. Soient $u \in U$ et
$v = f(u)$. Alors $\dim_u(U) = \dim_v(V)$.
\end{lemma}

\begin{proof}
Dans l'énoncé, $\dim_u(U)$ désigne la dimension de $U$ en $u$, définie dans
la définition \ref{topology-definition-Krull} de Topologie, comme le minimum des
dimensions de Krull des voisinages ouverts de $u$ dans $U$. Il en va de même
pour $\dim_v(V)$.

\medskip\noindent
Montrons que $\dim_v(V) \geq \dim_u(U)$. Soit $V'$ un voisinage ouvert de
$v$ dans $V$. Il existe un voisinage ouvert $U'$ de $u$ dans $U$, contenu
dans $f^{-1}(V')$, tel que $\dim_u(U) = \dim(U')$. Soit
$Z_0 \subset Z_1 \subset \ldots \subset Z_n$ une chaîne de sous-schémas
fermés irréductibles de $U'$. Si $\xi_i \in Z_i$ est le point générique,
nous avons les spécialisations
$\xi_n \leadsto \xi_{n - 1} \leadsto \ldots \leadsto \xi_0$. Elles donnent
les spécialisations
$f(\xi_n) \leadsto f(\xi_{n - 1}) \leadsto \ldots \leadsto f(\xi_0)$ dans
$V'$. De plus, $f(\xi_j) \not = f(\xi_i)$ si $i \not = j$, car les fibres de
$f$ sont discrètes (voir le lemme
\ref{morphisms-lemma-etale-over-field} de Morphismes). Ainsi, $\dim(V') \geq n$.
L'inégalité $\dim_v(V) \geq \dim_u(U)$ en résulte formellement.

\medskip\noindent
Montrons que $\dim_u(U) \geq \dim_v(V)$. Soit $U'$ un voisinage ouvert de
$u$ dans $U$. L'ensemble $V' = f(U')$ est un voisinage ouvert de $v$ d'après
le lemme \ref{morphisms-lemma-fppf-open} de Morphismes. Ainsi,
$\dim(V') \geq \dim_v(V)$. Choisissons une chaîne
$Z_0 \subset Z_1 \subset \ldots \subset Z_n$ de sous-schémas fermés
irréductibles de $V'$. Si $\xi_i \in Z_i$ est le point générique, nous avons
les spécialisations
$\xi_n \leadsto \xi_{n - 1} \leadsto \ldots \leadsto \xi_0$. Comme
$\xi_0 \in f(U')$, il existe un point $\eta_0 \in U'$ tel que
$f(\eta_0) = \xi_0$. Considérons l'homomorphisme d'anneaux locaux
$$
\mathcal{O}_{V', \xi_0} \longrightarrow \mathcal{O}_{U', \eta_0}
$$
qui est plat et local d'après le lemme
\ref{morphisms-lemma-etale-flat} de Morphismes. Les points $\xi_i$ correspondent aux idéaux
premiers de l'anneau de gauche d'après le lemme
\ref{schemes-lemma-specialize-points} de Schémas. Le théorème de descente (voir la
section \ref{algebra-section-going-up} d'Algèbre) appliqué à l'homomorphisme affiché
fournit une suite de spécialisations
$\eta_n \leadsto \eta_{n - 1} \leadsto \ldots \leadsto \eta_0$ dans $U'$,
qui s'envoie sur la suite
$\xi_n \leadsto \xi_{n - 1} \leadsto \ldots \leadsto \xi_0$ par $f$.
Il s'ensuit que $\dim_u(U) \geq \dim_v(V)$.
\end{proof}

\noindent
Soit $(X, x)$ un germe de schéma. La classe d'isomorphisme de l'anneau local
$\mathcal{O}_{X, x}$ est un invariant du germe. Le lemme suivant affirme que
la propriété $\dim(\mathcal{O}_{X, x}) = d$ est locale pour la topologie
\'etale sur les germes.

\begin{lemma}
\label{descent-lemma-dimension-local-ring-local}
Soit $f : U \to V$ un morphisme \'etale de schémas. Soient $u \in U$ et
$v = f(u)$. Alors
$\dim(\mathcal{O}_{U, u}) = \dim(\mathcal{O}_{V, v})$.
\end{lemma}

\begin{proof}
L'énoncé algébrique à démontrer est le suivant : si $A \to B$ est un
homomorphisme d'anneaux \'etale, si $\mathfrak q$ est un idéal premier de
$B$ situé au-dessus de $\mathfrak p \subset A$, alors
$\dim(A_{\mathfrak p}) = \dim(B_{\mathfrak q})$. C'est le lemme
\ref{more-algebra-lemma-dimension-etale-extension} de Compléments d'algèbre.
\end{proof}

\noindent
Soit $(X, x)$ un germe de schéma. La classe d'isomorphisme de l'anneau local
$\mathcal{O}_{X, x}$ est un invariant du germe. Le lemme suivant affirme que
la propriété ``$\mathcal{O}_{X, x}$ est régulier'' est locale pour la
topologie \'etale sur les germes.

\begin{lemma}
\label{descent-lemma-regular-local-ring-local}
Soit $f : U \to V$ un morphisme \'etale de schémas. Soient $u \in U$ et
$v = f(u)$. Alors $\mathcal{O}_{U, u}$ est un anneau local régulier si et
seulement si $\mathcal{O}_{V, v}$ est un anneau local régulier.
\end{lemma}

\begin{proof}
L'énoncé algébrique à démontrer est le suivant : si $A \to B$ est un
homomorphisme d'anneaux \'etale, si $\mathfrak q$ est un idéal premier de
$B$ situé au-dessus de $\mathfrak p \subset A$, alors $A_{\mathfrak p}$ est
régulier si et seulement si $B_{\mathfrak q}$ est régulier. C'est le lemme
\ref{more-algebra-lemma-regular-etale-extension} de Compléments d'algèbre.
\end{proof}







\section{Propriétés des morphismes locales sur la cible}
\label{descent-section-descending-properties-morphisms}

\noindent
Soit $f : X \to Y$ un morphisme de schémas. Soit $g : Y' \to Y$ un
morphisme de schémas. Notons $f' : X' \to Y'$ le changement de base de $f$
par $g$ :
$$
\xymatrix{
X' \ar[d]_{f'} \ar[r]_{g'} & X \ar[d]^f \\
Y' \ar[r]^g & Y
}
$$
Soit $\mathcal{P}$ une propriété des morphismes de schémas. On peut se
demander si (a) $\mathcal{P}(f) \Rightarrow \mathcal{P}(f')$ et si,
réciproquement, (b) $\mathcal{P}(f') \Rightarrow \mathcal{P}(f)$. Si (a) est
vérifiée chaque fois que $g$ est plat, on dit que $\mathcal{P}$ est stable
par changement de base plat. Si (b) est vérifiée chaque fois que $g$ est plat
et surjectif, on dit que $\mathcal{P}$ se descend par les changements de base
plats surjectifs. Si $\mathcal{P}$ est stable par changement de base plat et
se descend par les changements de base plats surjectifs, on dit que
$\mathcal{P}$ est locale sur la cible pour la topologie plate. Comparer avec
la discussion de la section \ref{descent-section-descending-properties}. Cette notion
est très importante ; nous la formalisons dans la définition suivante.

\begin{definition}
\label{descent-definition-property-morphisms-local}
Soit $\mathcal{P}$ une propriété des morphismes de schémas au-dessus d'une
base. Soit $\tau \in \{fpqc, fppf, syntomic, smooth, \etale, Zariski\}$.
On dit que $\mathcal{P}$ est {\it locale sur la base pour $\tau$}, ou
{\it locale sur la cible pour $\tau$}, ou encore {\it locale sur la base
pour la topologie $\tau$}, si, pour tout recouvrement pour la topologie
$\tau$, $\{Y_i \to Y\}_{i \in I}$ (voir la section
\ref{topologies-section-procedure} de Topologies), et pour tout morphisme de schémas
$f : X \to Y$ au-dessus de $S$, on a
$$
f \text{ vérifie }\mathcal{P}
\Leftrightarrow
\text{chaque }Y_i \times_Y X \to Y_i\text{ vérifie }\mathcal{P}.
$$
\end{definition}

\noindent
Bien entendu, puisque les isomorphismes sont toujours des recouvrements, on
constate (ou l'on exige) que la propriété $\mathcal{P}$ est vérifiée par
$X \to Y$ si et seulement si elle est vérifiée par toute flèche
$X' \to Y'$ isomorphe à $X \to Y$. Si une propriété est locale sur la cible
pour $\tau$, elle est stable par changement de base par les morphismes qui
interviennent dans les recouvrements pour $\tau$. Voici un énoncé formel.

\begin{lemma}
\label{descent-lemma-pullback-property-local-target}
Soit $\tau \in \{fpqc, fppf, syntomic, smooth, \etale, Zariski\}$.
Soit $\mathcal{P}$ une propriété des morphismes locale sur la cible pour
$\tau$. Soit $f : X \to Y$ un morphisme vérifiant $\mathcal{P}$. Pour tout
morphisme $Y' \to Y$ qui est plat, resp.\ plat et localement de présentation
finie, resp.\ syntomique, resp.\ \'etale, resp.\ une immersion ouverte, le
changement de base $f' : Y' \times_Y X \to Y'$ de $f$ vérifie
$\mathcal{P}$.
\end{lemma}

\begin{proof}
Cela tient au fait que l'on peut compléter $Y' \to Y$ en une famille de
morphismes qui forme un recouvrement pour la topologie $\tau$.
\end{proof}

\noindent
Une conséquence simple et souvent utilisée est la suivante. Si
$f : X \to Y$ vérifie une propriété $\mathcal{P}$ locale sur la cible pour
$\tau$, et si $f(X) \subset V$ pour un sous-schéma ouvert $V \subset Y$,
alors le morphisme induit $X \to V$ vérifie lui aussi $\mathcal{P}$.
En effet, le changement de base de $f$ par $V \to Y$ est $X \to V$.

\begin{lemma}
\label{descent-lemma-largest-open-of-the-base}
Soit $\tau \in \{fppf, syntomic, smooth, \etale\}$. Soit $\mathcal{P}$ une
propriété des morphismes locale sur la cible pour $\tau$. Pour tout morphisme
de schémas $f : X \to Y$, il existe un plus grand ouvert $W(f) \subset Y$
tel que la restriction $X_{W(f)} \to W(f)$ vérifie $\mathcal{P}$. En outre,
\begin{enumerate}
\item si $g : Y' \to Y$ est plat et localement de présentation finie,
syntomique, lisse ou \'etale, et si le changement de base
$f' : X_{Y'} \to Y'$ vérifie $\mathcal{P}$, alors $g(Y') \subset W(f)$,
\item si $g : Y' \to Y$ est plat et localement de présentation finie,
syntomique, lisse ou \'etale, alors $W(f') = g^{-1}(W(f))$, et
\item si $\{g_i : Y_i \to Y\}$ est un recouvrement pour la topologie $\tau$,
alors $g_i^{-1}(W(f)) = W(f_i)$, où $f_i$ est le changement de base de $f$
par $Y_i \to Y$.
\end{enumerate}
\end{lemma}

\begin{proof}
Considérons la réunion $W$ des images $g(Y') \subset Y$ des morphismes
$g : Y' \to Y$ ayant les propriétés suivantes :
\begin{enumerate}
\item $g$ est plat et localement de présentation finie, syntomique, lisse
ou \'etale, et
\item le changement de base $Y' \times_{g, Y} X \to Y'$ vérifie la propriété
$\mathcal{P}$.
\end{enumerate}
Un tel morphisme $g$ étant ouvert (voir le lemme
\ref{morphisms-lemma-fppf-open} de Morphismes), $W \subset Y$ est un ouvert de $Y$.
Puisque $\mathcal{P}$ est locale pour la topologie $\tau$, la restriction
$X_W \to W$ vérifie $\mathcal{P}$ : on dispose en effet d'un recouvrement
$\{Y' \to W\}$ de $W$ dont les changements de base vérifient
$\mathcal{P}$. Cela prouve l'existence et montre que $W(f)$ possède la
propriété (1). Pour (2), on a $W(f') \supset g^{-1}(W(f))$, car
$\mathcal{P}$ est stable par changement de base plat et localement de
présentation finie, syntomique, lisse ou \'etale ; voir le
lemme \ref{descent-lemma-pullback-property-local-target}. D'autre part, si
$Y'' \subset Y'$ est un ouvert tel que $X_{Y''} \to Y''$ vérifie
$\mathcal{P}$, alors $Y'' \to Y$ se factorise par $W$ par construction,
c'est-à-dire $Y'' \subset g^{-1}(W(f))$. Cela prouve (2). L'assertion (3)
résulte de (2), puisque chaque morphisme $Y_i \to Y$ est plat et localement
de présentation finie, syntomique, lisse ou \'etale, par définition d'un
recouvrement pour la topologie $\tau$.
\end{proof}

\begin{lemma}
\label{descent-lemma-descending-properties-morphisms}
Soit $\mathcal{P}$ une propriété des morphismes de schémas au-dessus d'une
base. Soit $\tau \in \{fpqc, fppf, \etale, lisse, syntomique\}$. Supposons que
\begin{enumerate}
\item la propriété soit stable par changement de base plat, plat et
localement de présentation finie, \'etale, lisse ou syntomique selon que
$\tau$ vaut fpqc, fppf, \'etale, smooth ou syntomic (comparer avec la
définition \ref{schemes-definition-preserved-by-base-change} de Schémas),
\item la propriété soit locale sur la base pour la topologie de Zariski,
\item pour tout morphisme surjectif de schémas affines $S' \to S$ qui est
plat, plat de présentation finie, \'etale, lisse ou syntomique selon que
$\tau$ vaut fpqc, fppf, \'etale, smooth ou syntomic, et pour tout morphisme de
schémas $f : X \to S$, la propriété $\mathcal{P}$ soit vérifiée par $f$ si
la propriété $\mathcal{P}$ est vérifiée par le changement de base
$f' : X' = S' \times_S X \to S'$.
\end{enumerate}
Alors $\mathcal{P}$ est locale sur la base pour la topologie $\tau$.
\end{lemma}

\begin{proof}
Cela résulte presque immédiatement de la définition d'un recouvrement pour
la topologie $\tau$ ; voir les définitions
\ref{topologies-definition-fpqc-covering},
\ref{topologies-definition-fppf-covering},
\ref{topologies-definition-etale-covering},
\ref{topologies-definition-smooth-covering} ou
\ref{topologies-definition-syntomic-covering} de Topologies, et les lemmes
\ref{topologies-lemma-fpqc-affine},
\ref{topologies-lemma-fppf-affine},
\ref{topologies-lemma-etale-affine},
\ref{topologies-lemma-smooth-affine} ou
\ref{topologies-lemma-syntomic-affine} de Topologies. Nous omettons les détails.
\end{proof}

\begin{remark}
\label{descent-remark-descending-properties-morphisms-standard}
(Il s'agit d'une répétition de la remarque
\ref{descent-remark-descending-properties-standard}.) Dans le
lemme \ref{descent-lemma-descending-properties-morphisms}, si $\tau = smooth$, on
peut supposer, dans la condition (3), que le morphisme est un morphisme lisse
standard (surjectif). Il en va de même lorsque $\tau = syntomic$ ou
$\tau = \etale$.
\end{remark}




\section[Localité sur la cible (fpqc)]{Propriétés des morphismes locales sur la cible pour la topologie fpqc}
\label{descent-section-descending-properties-morphisms-fpqc}

\noindent
Dans cette section, nous donnons un grand nombre de propriétés de morphismes
de schémas qui sont locales sur la base pour la topologie fpqc. En revanche,
nous montrerons dans la section
\ref{examples-section-non-descending-property-projective} d'Exemples que les propriétés
``projectif'' et ``quasi-projectif'' ne sont pas locales sur la base, même
pour la topologie de Zariski.

\begin{lemma}
\label{descent-lemma-descending-property-quasi-compact}
La propriété $\mathcal{P}(f) =$``$f$ est quasi-compact'' est locale sur la
base pour la topologie fpqc.
\end{lemma}

\begin{proof}
Tout changement de base d'un morphisme quasi-compact est quasi-compact ; voir
le lemme \ref{schemes-lemma-quasi-compact-preserved-base-change} de Schémas.
La quasi-compacité est locale sur la base pour la topologie de Zariski ; voir
le lemme \ref{schemes-lemma-quasi-compact-affine} de Schémas. Enfin, soit
$S' \to S$ un morphisme plat surjectif de schémas affines, et soit
$f : X \to S$ un morphisme. Supposons que le changement de base
$f' : X' \to S'$ soit quasi-compact. Alors $X'$ est quasi-compact et
$X' \to X$ est surjectif. Par conséquent, $X$ est quasi-compact, donc $f$
est quasi-compact. Le lemme \ref{descent-lemma-descending-properties-morphisms}
s'applique et conclut.
\end{proof}

\begin{lemma}
\label{descent-lemma-descending-property-quasi-separated}
La propriété $\mathcal{P}(f) =$``$f$ est quasi-séparé'' est locale sur la
base pour la topologie fpqc.
\end{lemma}

\begin{proof}
Tout changement de base d'un morphisme quasi-séparé est quasi-séparé ; voir
le lemme \ref{schemes-lemma-separated-permanence} de Schémas. La quasi-séparation
est locale sur la base pour la topologie de Zariski (cela résulte de la
définition ou du lemme
\ref{schemes-lemma-characterize-quasi-separated} de Schémas). Enfin, soit $S' \to S$
un morphisme plat surjectif de schémas affines, et soit $f : X \to S$ un
morphisme. Supposons que le changement de base $f' : X' \to S'$ soit
quasi-séparé. Cela signifie que
$\Delta' : X' \to X'\times_{S'} X'$ est quasi-compact. Or $\Delta'$ est le
changement de base de $\Delta : X \to X \times_S X$ par $S' \to S$.
Le lemme \ref{descent-lemma-descending-property-quasi-compact} montre donc que
$\Delta$ est quasi-compact et, par suite, que $f$ est quasi-séparé.
Le lemme \ref{descent-lemma-descending-properties-morphisms} s'applique et conclut.
\end{proof}

\begin{lemma}
\label{descent-lemma-descending-property-universally-closed}
La propriété $\mathcal{P}(f) =$``$f$ est universellement fermé'' est locale
sur la base pour la topologie fpqc.
\end{lemma}

\begin{proof}
Par définition, tout changement de base d'un morphisme universellement fermé
est universellement fermé. Cette propriété est locale sur la base pour la
topologie de Zariski (par définition ou d'après le lemme
\ref{morphisms-lemma-universally-closed-local-on-the-base} de Morphismes). Enfin, soit
$S' \to S$ un morphisme plat surjectif de schémas affines, et soit
$f : X \to S$ un morphisme. Supposons que le changement de base
$f' : X' \to S'$ soit universellement fermé. Soit $T \to S$ un morphisme
quelconque. Considérons le diagramme
$$
\xymatrix{
X' \ar[d] &
S' \times_S T \times_S X \ar[d] \ar[r] \ar[l] &
T \times_S X \ar[d] \\
S' &
S' \times_S T \ar[r] \ar[l] &
T
}
$$
dont les deux carrés sont cartésiens. L'hypothèse implique que la flèche
verticale médiane est fermée. Les flèches horizontales de droite sont plates,
quasi-compactes et surjectives, puisqu'elles s'obtiennent par changement de
base de $S' \to S$. Un sous-ensemble de $T$ est fermé si et seulement si son
image inverse dans $S' \times_S T$ est fermée ; voir le lemme
\ref{morphisms-lemma-fpqc-quotient-topology} de Morphismes. Une chasse au diagramme
immédiate montre que la flèche verticale de droite est également fermée ; on
en conclut que $X \to S$ est universellement fermé. Le lemme
\ref{descent-lemma-descending-properties-morphisms} s'applique et conclut.
\end{proof}

\begin{lemma}
\label{descent-lemma-descending-property-universally-open}
La propriété $\mathcal{P}(f) =$``$f$ est universellement ouvert'' est locale
sur la base pour la topologie fpqc.
\end{lemma}

\begin{proof}
La démonstration est identique à celle du
lemme \ref{descent-lemma-descending-property-universally-closed}.
\end{proof}

\begin{lemma}
\label{descent-lemma-descending-property-universally-submersive}
La propriété $\mathcal{P}(f) =$``$f$ est universellement submersif'' est
locale sur la base pour la topologie fpqc.
\end{lemma}

\begin{proof}
La démonstration est identique à celle du
lemme \ref{descent-lemma-descending-property-universally-closed}, en utilisant qu'un
morphisme plat, quasi-compact et surjectif est universellement submersif
d'après le lemme \ref{morphisms-lemma-fpqc-quotient-topology} de Morphismes.
\end{proof}

\begin{lemma}
\label{descent-lemma-descending-property-separated}
La propriété $\mathcal{P}(f) =$``$f$ est séparé'' est locale sur la base pour
la topologie fpqc.
\end{lemma}

\begin{proof}
Tout changement de base d'un morphisme séparé est séparé ; voir le lemme
\ref{schemes-lemma-separated-permanence} de Schémas. La séparation est locale sur
la base pour la topologie de Zariski (par définition ou d'après le lemme
\ref{schemes-lemma-characterize-separated} de Schémas). Enfin, soit $S' \to S$ un
morphisme plat surjectif de schémas affines, et soit $f : X \to S$ un
morphisme. Supposons que le changement de base $f' : X' \to S'$ soit séparé.
Cela signifie que $\Delta' : X' \to X'\times_{S'} X'$ est une immersion
fermée, donc universellement fermée. Or $\Delta'$ est le changement de base
de $\Delta : X \to X \times_S X$ par $S' \to S$. Le
lemme \ref{descent-lemma-descending-property-universally-closed} montre donc que
$\Delta$ est universellement fermé. Comme c'est une immersion (Schémas,
lemme \ref{schemes-lemma-diagonal-immersion} de Schémas), $\Delta$ est une immersion
fermée. Ainsi $f$ est séparé. Le
lemme \ref{descent-lemma-descending-properties-morphisms} s'applique et conclut.
\end{proof}

\begin{lemma}
\label{descent-lemma-descending-property-surjective}
La propriété $\mathcal{P}(f) =$``$f$ est surjectif'' est locale sur la base
pour la topologie fpqc.
\end{lemma}

\begin{proof}
C'est immédiat.
\end{proof}

\begin{lemma}
\label{descent-lemma-descending-property-dominant}
La propriété $\mathcal{P}(f) =$``$f$ est quasi-compact et dominant'' est
locale sur la base pour la topologie fpqc.
\end{lemma}

\begin{proof}
D'après le lemme \ref{morphisms-lemma-base-change-dominant} de Morphismes, les
morphismes quasi-compacts dominants sont stables par changement de base plat.
La réciproque est plus facile. En effet, la quasi-compacité est locale sur la
base pour la topologie fpqc d'après le
lemme \ref{descent-lemma-descending-property-quasi-compact}. La dominance est
manifestement locale sur la base pour la topologie de Zariski. Enfin, soit
$S' \to S$ un morphisme surjectif de schémas (non nécessairement plat), et
soit $f : X \to S$ un morphisme. Supposons que le changement de base
$f' : X' \to S'$ soit dominant. Alors $X' \to S'\to S$ est le composé de
deux morphismes dominants, donc est dominant. Comme il s'agit aussi du composé
$X' \to X \to S$, il s'ensuit que $X\to S$ est dominant.
\end{proof}

\noindent
Les morphismes dominants ne sont pas stables par changement de base plat en
toute généralité, mais le cas particulier suivant mérite d'être signalé.
Voir aussi le lemme \ref{descent-lemma-descending-fppf-property-dominant}.

\begin{lemma}
\label{descent-lemma-descending-property-dominant-over-field}
Soit $E/k$ une extension de corps. Un morphisme $X \to Y$ au-dessus de $k$
est alors dominant si et seulement si le changement de base $X_E \to Y_E$
est dominant.
\end{lemma}

\begin{proof}
D'après le lemme
\ref{morphisms-lemma-scheme-over-field-universally-open} de Morphismes, le morphisme
$\Spec(E) \to \Spec(k)$ est universellement ouvert. Ainsi $Y_E \to Y$ est
ouvert. Par conséquent, si $X \to Y$ est dominant, le lemme
\ref{morphisms-lemma-open-base-change-dominant} de Morphismes montre que $X_E \to Y_E$
est dominant. Réciproquement, supposons que $X_E \to Y_E$ soit dominant.
Le morphisme $Y_E \to Y$ est surjectif d'après le lemme
\ref{morphisms-lemma-base-change-surjective} de Morphismes ; le composé
$X_E \to Y_E \to Y$ est donc dominant. Mais ce composé est aussi
$X_E \to X \to Y$ ; il s'ensuit que $X \to Y$ est dominant.
\end{proof}

\begin{lemma}
\label{descent-lemma-descending-property-universally-injective}
La propriété $\mathcal{P}(f) =$``$f$ est universellement injectif'' est locale
sur la base pour la topologie fpqc.
\end{lemma}

\begin{proof}
Tout changement de base d'un morphisme universellement injectif est
universellement injectif (c'est formel). Cette propriété est locale sur la
base pour la topologie de Zariski, ce qui résulte immédiatement de la
définition. Enfin, soit $S' \to S$ un morphisme plat surjectif de schémas
affines, et soit $f : X \to S$ un morphisme. Supposons que le changement de
base $f' : X' \to S'$ soit universellement injectif. Soit $K$ un corps, et
soient $a, b : \Spec(K) \to X$ deux morphismes tels que
$f \circ a = f \circ b$. Puisque $S' \to S$ est surjectif, la discussion de
la section \ref{schemes-section-points} de Schémas fournit une extension de corps
$K'/K$ et un morphisme $\Spec(K')
\to S'$ tels que le diagramme suivant, formé de flèches pleines, soit commutatif :
$$
\xymatrix{
\Spec(K') \ar[rrd] \ar@{-->}[rd]_{a', b'} \ar[dd] \\
 &
X' \ar[r] \ar[d] &
S' \ar[d] \\
\Spec(K) \ar[r]^{a, b} &
X \ar[r] &
S
}
$$
Le carré étant cartésien, on obtient les deux flèches pointillées $a'$, $b'$
qui rendent le diagramme commutatif. Comme $X' \to S'$ est universellement
injectif, nous avons $a' = b'$ d'après le lemme
\ref{morphisms-lemma-universally-injective} de Morphismes. Cela impose évidemment $a = b$
(voir la discussion de la section \ref{schemes-section-points} de Schémas). Le
lemme \ref{descent-lemma-descending-properties-morphisms} s'applique et conclut.

\medskip\noindent
On peut aussi utiliser la caractérisation d'un morphisme universellement
injectif par la surjectivité de sa diagonale ; voir le lemme
\ref{morphisms-lemma-universally-injective} de Morphismes. Le résultat découle alors du fait
que la surjectivité est locale sur la base pour la topologie fpqc ; voir le
lemme \ref{descent-lemma-descending-property-surjective}. (Indication : utiliser que
le changement de base de la diagonale est la diagonale du changement de base.)
\end{proof}

\begin{lemma}
\label{descent-lemma-descending-property-universal-homeomorphism}
La propriété $\mathcal{P}(f) =$``$f$ est un homéomorphisme universel'' est
locale sur la base pour la topologie fpqc.
\end{lemma}

\begin{proof}
On peut procéder exactement comme dans la démonstration du
lemme \ref{descent-lemma-descending-property-universally-closed}. On peut aussi
observer qu'une application d'espaces topologiques est un homéomorphisme si
et seulement si elle est injective, surjective et ouverte. Un homéomorphisme
universel équivaut donc à un morphisme surjectif, universellement injectif et
universellement ouvert. Le résultat découle ainsi des lemmes
\ref{descent-lemma-descending-property-surjective},
\ref{descent-lemma-descending-property-universally-injective} et
\ref{descent-lemma-descending-property-universally-open}.
\end{proof}

\begin{lemma}
\label{descent-lemma-descending-property-locally-finite-type}
La propriété $\mathcal{P}(f) =$``$f$ est localement de type fini'' est locale
sur la base pour la topologie fpqc.
\end{lemma}

\begin{proof}
La propriété d'être localement de type fini est stable par changement de base ;
voir le lemme \ref{morphisms-lemma-base-change-finite-type} de Morphismes. Elle est
locale sur la base pour la topologie de Zariski ; voir le lemme
\ref{morphisms-lemma-locally-finite-type-characterize} de Morphismes. Enfin, soit
$S' \to S$ un morphisme plat surjectif de schémas affines, et soit
$f : X \to S$ un morphisme. Supposons que le changement de base
$f' : X' \to S'$ soit localement de type fini. Soit $U \subset X$ un ouvert
affine. Alors $U' = S' \times_S U$ est affine et de type fini sur $S'$.
Écrivons $S = \Spec(R)$, $S' = \Spec(R')$, $U = \Spec(A)$ et
$U' = \Spec(A')$. Nous savons que $R \to R'$ est fidèlement plat, que
$A' = R' \otimes_R A$ et que $R' \to A'$ est de type fini. Il faut montrer
que $R \to A$ est de type fini. C'est le lemme
\ref{algebra-lemma-finite-type-descends} d'Algèbre. Il s'ensuit que $f$ est localement
de type fini. Le lemme \ref{descent-lemma-descending-properties-morphisms}
s'applique et conclut.
\end{proof}

\begin{lemma}
\label{descent-lemma-descending-property-locally-finite-presentation}
La propriété $\mathcal{P}(f) =$``$f$ est localement de présentation finie''
est locale sur la base pour la topologie fpqc.
\end{lemma}

\begin{proof}
La propriété d'être localement de présentation finie est stable par
changement de base ; voir le lemme
\ref{morphisms-lemma-base-change-finite-presentation} de Morphismes. Elle est locale sur la
base pour la topologie de Zariski ; voir le lemme
\ref{morphisms-lemma-locally-finite-presentation-characterize} de Morphismes. Enfin, soit
$S' \to S$ un morphisme plat surjectif de schémas affines, et soit
$f : X \to S$ un morphisme. Supposons que le changement de base
$f' : X' \to S'$ soit localement de présentation finie. Soit $U \subset X$
un ouvert affine. Alors $U' = S' \times_S U$ est affine et de type fini sur
$S'$. Écrivons $S = \Spec(R)$, $S' = \Spec(R')$, $U = \Spec(A)$ et
$U' = \Spec(A')$. Nous savons que $R \to R'$ est fidèlement plat, que
$A' = R' \otimes_R A$ et que $R' \to A'$ est de présentation finie. Il faut
montrer que $R \to A$ est de présentation finie. C'est le lemme
\ref{algebra-lemma-finite-presentation-descends} d'Algèbre. Il s'ensuit que $f$ est
localement de présentation finie. Le
lemme \ref{descent-lemma-descending-properties-morphisms} s'applique et conclut.
\end{proof}

\begin{lemma}
\label{descent-lemma-descending-property-finite-type}
La propriété $\mathcal{P}(f) =$``$f$ est de type fini''
est locale sur la base pour la topologie fpqc.
\end{lemma}

\begin{proof}
Combiner les lemmes \ref{descent-lemma-descending-property-quasi-compact}
et \ref{descent-lemma-descending-property-locally-finite-type}.
\end{proof}

\begin{lemma}
\label{descent-lemma-descending-property-finite-presentation}
La propriété $\mathcal{P}(f) =$``$f$ est de présentation finie''
est locale sur la base pour la topologie fpqc.
\end{lemma}

\begin{proof}
Combiner les lemmes \ref{descent-lemma-descending-property-quasi-compact},
\ref{descent-lemma-descending-property-quasi-separated} et
\ref{descent-lemma-descending-property-locally-finite-presentation}.
\end{proof}

\begin{lemma}
\label{descent-lemma-descending-property-proper}
La propriété $\mathcal{P}(f) =$``$f$ est propre''
est locale sur la base pour la topologie fpqc.
\end{lemma}

\begin{proof}
Le lemme résulte de la combinaison des
lemmes \ref{descent-lemma-descending-property-universally-closed},
\ref{descent-lemma-descending-property-separated}
et \ref{descent-lemma-descending-property-finite-type}.
\end{proof}

\begin{lemma}
\label{descent-lemma-descending-property-flat}
La propriété $\mathcal{P}(f) =$``$f$ est plat''
est locale sur la base pour la topologie fpqc.
\end{lemma}

\begin{proof}
La platitude est préservée par tout changement de base, voir
le lemme \ref{morphisms-lemma-base-change-flat} de Morphismes.
Elle est locale pour la topologie de Zariski sur la base par définition.
Enfin, soit
$S' \to S$ un morphisme plat surjectif de schémas affines,
et soit $f : X \to S$ un morphisme. Supposons que le changement de base
$f' : X' \to S'$ soit plat.
Soit $U \subset X$ un ouvert affine. Alors $U' = S' \times_S U$
est affine. Écrivons
$S = \Spec(R)$,
$S' = \Spec(R')$,
$U = \Spec(A)$ et
$U' = \Spec(A')$.
Nous savons que $R \to R'$ est fidèlement plat,
que $A' = R' \otimes_R A$ et que $R' \to A'$ est plat.
Il s'agit de montrer que $R \to A$ est plat.
Cela résulte immédiatement du lemme
\ref{algebra-lemma-flatness-descends} d'Algèbre.
Ainsi $f$ est plat.
Le lemme \ref{descent-lemma-descending-properties-morphisms} s'applique donc et conclut.
\end{proof}

\begin{lemma}
\label{descent-lemma-descending-property-open-immersion}
La propriété $\mathcal{P}(f) =$``$f$ est une immersion ouverte''
est locale sur la base pour la topologie fpqc.
\end{lemma}

\begin{proof}
La propriété d'être une immersion ouverte est stable par changement de base,
voir le lemme \ref{schemes-lemma-base-change-immersion} de Schémas.
Elle est locale pour la topologie de Zariski sur la base
(cela est évident).

\medskip\noindent
Soit $S' \to S$ un morphisme plat surjectif de schémas affines,
et soit $f : X \to S$ un morphisme. Supposons que le changement de base
$f' : X' \to S'$ soit une immersion ouverte. Nous affirmons que $f$ est une
immersion ouverte.
Alors $f'$ est universellement ouvert et universellement injectif.
Nous en déduisons que $f$ est universellement ouvert par le
lemme \ref{descent-lemma-descending-property-universally-open}, et
universellement injectif par le
lemme \ref{descent-lemma-descending-property-universally-injective}.
En particulier, $f(X) \subset S$ est ouvert. S'il est possible de montrer,
pour tout ouvert affine $U \subset f(X)$, que $f^{-1}(U) \to U$
est un isomorphisme, alors $f$ est une immersion ouverte et la preuve est
achevée. Si $U' \subset S'$ désigne l'image inverse de $U$,
alors $U' \to U$ est un morphisme fidèlement plat entre schémas affines et
$(f')^{-1}(U') \to U'$ est un isomorphisme (car $f'(X')$ contient $U'$
par le choix de $U$). Nous sommes ainsi ramenés au cas traité dans le
paragraphe suivant.

\medskip\noindent
Soit $S' \to S$ un morphisme plat surjectif de schémas affines,
soit $f : X \to S$ un morphisme, et supposons que le changement de base
$f' : X' \to S'$ soit un isomorphisme. Il faut montrer que $f$ est lui aussi
un isomorphisme. Il est clair que $f$ est surjectif, universellement injectif
et universellement ouvert (voir les arguments ci-dessus pour les deux
dernières propriétés). Ainsi $f$ est bijectif, c'est-à-dire que $f$ est un
homéomorphisme. Le morphisme $f$ est donc affine par le lemme
\ref{morphisms-lemma-homeomorphism-affine} de Morphismes. Comme
$$
\mathcal{O}(S') \to
\mathcal{O}(X') =
\mathcal{O}(S') \otimes_{\mathcal{O}(S)} \mathcal{O}(X)
$$
est un isomorphisme et que $\mathcal{O}(S) \to \mathcal{O}(S')$
est fidèlement plat, il en résulte que
$\mathcal{O}(S) \to \mathcal{O}(X)$ est un isomorphisme.
Ainsi $f$ est un isomorphisme. Cela démontre l'affirmation ci-dessus.
Le lemme \ref{descent-lemma-descending-properties-morphisms} s'applique donc et conclut.
\end{proof}

\begin{lemma}
\label{descent-lemma-descending-property-isomorphism}
La propriété $\mathcal{P}(f) =$``$f$ est un isomorphisme''
est locale sur la base pour la topologie fpqc.
\end{lemma}

\begin{proof}
Combiner les lemmes \ref{descent-lemma-descending-property-surjective}
et \ref{descent-lemma-descending-property-open-immersion}.
\end{proof}

\begin{lemma}
\label{descent-lemma-descending-property-affine}
La propriété $\mathcal{P}(f) =$``$f$ est affine''
est locale sur la base pour la topologie fpqc.
\end{lemma}

\begin{proof}
Tout changement de base d'un morphisme affine est affine, voir
le lemme \ref{morphisms-lemma-base-change-affine} de Morphismes.
La propriété d'être affine est locale pour la topologie de Zariski sur la base,
voir le lemme \ref{morphisms-lemma-characterize-affine} de Morphismes.
Enfin, soit
$g : S' \to S$ un morphisme plat surjectif de schémas affines,
et soit $f : X \to S$ un morphisme. Supposons que le changement de base
$f' : X' \to S'$ soit affine. Autrement dit, $X'$ est affine, disons
$X' = \Spec(A')$. Écrivons aussi $S = \Spec(R)$
et $S' = \Spec(R')$. Il faut montrer que $X$ est affine.

\medskip\noindent
Les lemmes \ref{descent-lemma-descending-property-quasi-compact}
et \ref{descent-lemma-descending-property-separated} montrent que
$X \to S$ est séparé et quasi-compact. Par conséquent,
$f_*\mathcal{O}_X$ est un faisceau quasi-cohérent de $\mathcal{O}_S$-algèbres,
voir le lemme \ref{schemes-lemma-push-forward-quasi-coherent} de Schémas.
Ainsi $f_*\mathcal{O}_X = \widetilde{A}$ pour une certaine $R$-algèbre $A$.
Bien entendu, $A = \Gamma(X, \mathcal{O}_X)$.
De plus, par changement de base plat
(voir par exemple le lemme
\ref{coherent-lemma-flat-base-change-cohomology} de Cohomologie des schémas),
nous avons $g^*f_*\mathcal{O}_X = f'_*\mathcal{O}_{X'}$.
Autrement dit, $A' = R' \otimes_R A$.
Considérons le morphisme canonique
$$
X \longrightarrow \Spec(A)
$$
au-dessus de $S$ donné par le lemme
\ref{schemes-lemma-morphism-into-affine} de Schémas.
D'après ce qui précède, le changement de base de ce morphisme à $S'$ est un
isomorphisme. Il s'agit donc d'un isomorphisme par le
lemme \ref{descent-lemma-descending-property-isomorphism}.
Le lemme \ref{descent-lemma-descending-properties-morphisms} s'applique donc et conclut.
\end{proof}

\begin{lemma}
\label{descent-lemma-descending-property-closed-immersion}
La propriété $\mathcal{P}(f) =$``$f$ est une immersion fermée''
est locale sur la base pour la topologie fpqc.
\end{lemma}

\begin{proof}
Soit $f : X \to Y$ un morphisme de schémas.
Soit $\{Y_i \to Y\}$ un recouvrement fpqc.
Supposons que chaque $f_i : Y_i \times_Y X \to Y_i$
soit une immersion fermée.
Cela implique que chaque $f_i$ est affine, voir
le lemme \ref{morphisms-lemma-closed-immersion-affine} de Morphismes.
Le lemme \ref{descent-lemma-descending-property-affine}
permet d'en déduire que $f$ est affine. Il reste à montrer que
$\mathcal{O}_Y \to f_*\mathcal{O}_X$ est surjectif.
Pour tout $y \in Y$, il existe un $i$ et un point
$y_i \in Y_i$ d'image $y$.
Par le lemme \ref{coherent-lemma-flat-base-change-cohomology}
de Cohomologie des schémas,
le faisceau $f_{i, *}(\mathcal{O}_{Y_i \times_Y X})$
est l'image inverse de $f_*\mathcal{O}_X$.
Par hypothèse, c'est un quotient de $\mathcal{O}_{Y_i}$.
Par conséquent,
$$
\Big(
\mathcal{O}_{Y, y} \longrightarrow (f_*\mathcal{O}_X)_y
\Big)
\otimes_{\mathcal{O}_{Y, y}} \mathcal{O}_{Y_i, y_i}
$$
est surjectif. Puisque $\mathcal{O}_{Y_i, y_i}$ est fidèlement
plat sur $\mathcal{O}_{Y, y}$, il s'ensuit que
$\mathcal{O}_{Y, y} \longrightarrow (f_*\mathcal{O}_X)_y$ est surjectif,
comme voulu.
\end{proof}

\begin{lemma}
\label{descent-lemma-descending-property-quasi-affine}
La propriété $\mathcal{P}(f) =$``$f$ est quasi-affine''
est locale sur la base pour la topologie fpqc.
\end{lemma}

\begin{proof}
Soit $f : X \to Y$ un morphisme de schémas.
Soit $\{g_i : Y_i \to Y\}$ un recouvrement fpqc.
Supposons que chaque $f_i : Y_i \times_Y X \to Y_i$
soit quasi-affine.
Alors chaque $f_i$ est quasi-compact et séparé.
Les lemmes \ref{descent-lemma-descending-property-quasi-compact}
et \ref{descent-lemma-descending-property-separated}
impliquent donc que $f$ est quasi-compact et séparé.
Considérons le faisceau de $\mathcal{O}_Y$-algèbres
$\mathcal{A} = f_*\mathcal{O}_X$.
Par le lemme \ref{schemes-lemma-push-forward-quasi-coherent} de Schémas,
c'est une $\mathcal{O}_Y$-algèbre quasi-cohérente.
Considérons le morphisme canonique
$$
j : X \longrightarrow \underline{\Spec}_Y(\mathcal{A})
$$
voir le lemme \ref{constructions-lemma-canonical-morphism} de Constructions.
Par changement de base plat
(voir par exemple le lemme
\ref{coherent-lemma-flat-base-change-cohomology} de Cohomologie des schémas),
nous avons $g_i^*f_*\mathcal{O}_X = f_{i, *}\mathcal{O}_{X'}$,
où les $g_i : Y_i \to Y$ sont les morphismes plats donnés.
Ainsi, le changement de base $j_i$ de $j$ par $g_i$ est le morphisme canonique
du lemme \ref{constructions-lemma-canonical-morphism} de Constructions,
pour le morphisme $f_i$. Par hypothèse et par
le lemme \ref{morphisms-lemma-characterize-quasi-affine} de Morphismes,
tous ces morphismes $j_i$ sont des immersions ouvertes quasi-compactes.
Les lemmes \ref{descent-lemma-descending-property-quasi-compact} et
\ref{descent-lemma-descending-property-open-immersion} montrent donc que
$j$ est une immersion ouverte quasi-compacte.
Une nouvelle application du
lemme \ref{morphisms-lemma-characterize-quasi-affine} de Morphismes
montre alors que $f$ est quasi-affine.
\end{proof}

\begin{lemma}
\label{descent-lemma-descending-property-quasi-compact-immersion}
La propriété $\mathcal{P}(f) =$``$f$ est une immersion quasi-compacte''
est locale sur la base pour la topologie fpqc.
\end{lemma}

\begin{proof}
Soit $f : X \to Y$ un morphisme de schémas.
Soit $\{Y_i \to Y\}$ un recouvrement fpqc.
Notons $X_i = Y_i \times_Y X$ et $f_i : X_i \to Y_i$
le changement de base de $f$. Notons aussi
$q_i : Y_i \to Y$ les morphismes plats donnés.
Supposons que chaque $f_i$ soit une immersion quasi-compacte.
Par le lemme \ref{schemes-lemma-immersions-monomorphisms} de Schémas,
chaque $f_i$ est séparé.
Les lemmes \ref{descent-lemma-descending-property-quasi-compact} et
\ref{descent-lemma-descending-property-separated}
impliquent donc que $f$ est quasi-compact et séparé.
Soit $X \to Z \to Y$ la factorisation de $f$ par son
image schématique. Par le
lemme \ref{morphisms-lemma-quasi-compact-scheme-theoretic-image} de Morphismes,
le sous-schéma fermé $Z \subset Y$ est défini par le
faisceau quasi-cohérent d'idéaux
$\mathcal{I} = \Ker(\mathcal{O}_Y \to f_*\mathcal{O}_X)$,
car $f$ est quasi-compact. Par changement de base plat
(voir par exemple le lemme
\ref{coherent-lemma-flat-base-change-cohomology} de Cohomologie des schémas ;
nous utilisons ici le fait que $f$ est séparé),
nous voyons que $f_{i, *}\mathcal{O}_{X_i}$ est l'image inverse
$q_i^*f_*\mathcal{O}_X$.
Ainsi, $Y_i \times_Y Z$ est défini par le
faisceau quasi-cohérent d'idéaux $q_i^*\mathcal{I} =
\Ker(\mathcal{O}_{Y_i} \to f_{i, *}\mathcal{O}_{X_i})$.
Par le lemme \ref{morphisms-lemma-quasi-compact-immersion} de Morphismes,
les morphismes $X_i \to Y_i \times_Y Z$
sont des immersions ouvertes. Le
lemme \ref{descent-lemma-descending-property-open-immersion}
montre alors que $X \to Z$ est une immersion ouverte, et donc que
$f$ est une immersion, comme voulu
(nous avons déjà vu qu'il était quasi-compact).
\end{proof}

\begin{lemma}
\label{descent-lemma-descending-property-integral}
La propriété $\mathcal{P}(f) =$``$f$ est entier''
est locale sur la base pour la topologie fpqc.
\end{lemma}

\begin{proof}
Un morphisme entier est la même chose qu'un morphisme affine et
universellement fermé. Voir
le lemme \ref{morphisms-lemma-integral-universally-closed} de Morphismes.
Le lemme résulte donc de la combinaison des
lemmes \ref{descent-lemma-descending-property-universally-closed}
et \ref{descent-lemma-descending-property-affine}.
\end{proof}

\begin{lemma}
\label{descent-lemma-descending-property-finite}
La propriété $\mathcal{P}(f) =$``$f$ est fini''
est locale sur la base pour la topologie fpqc.
\end{lemma}

\begin{proof}
Un morphisme fini est la même chose qu'un morphisme entier
qui est localement de type fini. Voir
le lemme \ref{morphisms-lemma-finite-integral} de Morphismes.
Le lemme résulte donc de la combinaison des
lemmes \ref{descent-lemma-descending-property-locally-finite-type}
et \ref{descent-lemma-descending-property-integral}.
\end{proof}

\begin{lemma}
\label{descent-lemma-descending-property-quasi-finite}
Les propriétés
$\mathcal{P}(f) =$``$f$ est localement quasi-fini''
et
$\mathcal{P}(f) =$``$f$ est quasi-fini''
sont locales sur la base pour la topologie fpqc.
\end{lemma}

\begin{proof}
Soit $f : X \to S$ un morphisme de schémas, et soit $\{S_i \to S\}$
un recouvrement fpqc tel que chaque changement de base
$f_i : X_i \to S_i$ soit localement quasi-fini.
Nous avons déjà vu
(lemme \ref{descent-lemma-descending-property-locally-finite-type})
que la propriété ``être localement de type fini'' est locale sur la base pour la
topologie fpqc ; ainsi, $f$ est localement de type fini.
Il résulte alors du
lemme \ref{morphisms-lemma-base-change-quasi-finite} de Morphismes
que $f$ est localement quasi-fini. Le cas quasi-fini en découle,
puisque nous avons déjà vu que la propriété ``être quasi-compact'' est locale
sur la base pour la topologie fpqc
(lemme \ref{descent-lemma-descending-property-quasi-compact}).
\end{proof}

\begin{lemma}
\label{descent-lemma-descending-property-relative-dimension-d}
La propriété $\mathcal{P}(f) =$``$f$ est localement de type fini
de dimension relative $d$'' est locale sur la base pour la topologie fpqc.
\end{lemma}

\begin{proof}
Cela résulte immédiatement du fait que la propriété d'être localement de type
fini est locale sur la base pour la topologie fpqc, et du
lemme \ref{morphisms-lemma-dimension-fibre-after-base-change} de Morphismes.
\end{proof}

\begin{lemma}
\label{descent-lemma-descending-property-syntomic}
La propriété $\mathcal{P}(f) =$``$f$ est syntomique''
est locale sur la base pour la topologie fpqc.
\end{lemma}

\begin{proof}
Un morphisme est syntomique si et seulement s'il est localement de présentation
finie, plat, et si ses fibres sont localement d'intersection complète.
Nous avons déjà vu que la platitude et la propriété d'être localement de
présentation finie sont locales sur la base pour la topologie fpqc (lemmes
\ref{descent-lemma-descending-property-flat} et
\ref{descent-lemma-descending-property-locally-finite-presentation}).
Le résultat pour les morphismes syntomiques découle donc du
lemme \ref{morphisms-lemma-set-points-where-fibres-lci} de Morphismes.
\end{proof}

\begin{lemma}
\label{descent-lemma-descending-property-smooth}
La propriété $\mathcal{P}(f) =$``$f$ est lisse''
est locale sur la base pour la topologie fpqc.
\end{lemma}

\begin{proof}
Un morphisme est lisse si et seulement s'il est localement de présentation
finie, plat, et si ses fibres sont lisses. Nous avons déjà vu que la platitude
et la propriété d'être localement de présentation finie sont locales sur la
base pour la topologie fpqc (lemmes
\ref{descent-lemma-descending-property-flat} et
\ref{descent-lemma-descending-property-locally-finite-presentation}).
Le résultat pour les morphismes lisses découle donc du
lemme \ref{morphisms-lemma-set-points-where-fibres-smooth} de Morphismes.
\end{proof}

\begin{lemma}
\label{descent-lemma-descending-property-unramified}
La propriété $\mathcal{P}(f) =$``$f$ est non ramifié''
est locale sur la base pour la topologie fpqc.
La propriété $\mathcal{P}(f) =$``$f$ est G-non ramifié''
est locale sur la base pour la topologie fpqc.
\end{lemma}

\begin{proof}
Un morphisme est non ramifié (resp.\ G-non ramifié) si et seulement s'il est
localement de type fini (resp.\ de présentation finie)
et si son morphisme diagonal est une immersion ouverte (voir
le lemme \ref{morphisms-lemma-diagonal-unramified-morphism} de Morphismes).
Nous avons déjà vu que les propriétés d'être localement de type fini
(resp.\ localement de présentation finie) et d'être une immersion ouverte sont
locales sur la base pour la topologie fpqc (lemmes
\ref{descent-lemma-descending-property-locally-finite-presentation},
\ref{descent-lemma-descending-property-locally-finite-type} et
\ref{descent-lemma-descending-property-open-immersion}).
Le résultat en découle formellement.
\end{proof}

\begin{lemma}
\label{descent-lemma-descending-property-etale}
La propriété $\mathcal{P}(f) =$``$f$ est \'etale''
est locale sur la base pour la topologie fpqc.
\end{lemma}

\begin{proof}
Un morphisme est \'etale si et seulement s'il est plat et G-non ramifié.
Voir le lemme \ref{morphisms-lemma-flat-unramified-etale} de Morphismes.
Nous avons déjà vu que la platitude et la propriété d'être G-non ramifié
sont locales sur la base pour la topologie fpqc (lemmes
\ref{descent-lemma-descending-property-flat} et
\ref{descent-lemma-descending-property-unramified}).
Le résultat en découle.
\end{proof}

\begin{lemma}
\label{descent-lemma-descending-property-finite-locally-free}
La propriété $\mathcal{P}(f) =$``$f$ est fini localement libre''
est locale sur la base pour la topologie fpqc.
Soit $d \geq 0$.
La propriété $\mathcal{P}(f) =$``$f$ est fini localement libre de degré $d$''
est locale sur la base pour la topologie fpqc.
\end{lemma}

\begin{proof}
La propriété d'être fini localement libre équivaut à être
fini, plat et localement de présentation finie
(voir le lemme \ref{morphisms-lemma-finite-flat} de Morphismes).
Cela résulte donc des lemmes
\ref{descent-lemma-descending-property-finite},
\ref{descent-lemma-descending-property-flat} et
\ref{descent-lemma-descending-property-locally-finite-presentation}.
Si $f : Z \to U$ est fini localement libre, et si $\{U_i \to U\}$ est une
famille surjective de morphismes telle que chaque image inverse
$Z \times_U U_i \to U_i$ soit de degré $d$, alors $Z \to U$ est de degré $d$;
on peut par exemple lire le degré en un point $u \in U$ sur la fibre
$(f_*\mathcal{O}_Z)_u \otimes_{\mathcal{O}_{U, u}} \kappa(u)$.
\end{proof}

\begin{lemma}
\label{descent-lemma-descending-property-monomorphism}
La propriété $\mathcal{P}(f) =$``$f$ est un monomorphisme''
est locale sur la base pour la topologie fpqc.
\end{lemma}

\begin{proof}
Soit $f : X \to S$ un morphisme de schémas.
Soit $\{S_i \to S\}$ un recouvrement fpqc, et supposons que chacun des
changements de base $f_i : X_i \to S_i$ de $f$ soit un monomorphisme.
Soient $a, b : T \to X$ deux morphismes tels que
$f \circ a = f \circ b$. Il faut montrer que $a = b$.
Puisque $f_i$ est un monomorphisme, nous avons $a_i = b_i$, où
$a_i, b_i : S_i \times_S T \to X_i$ sont les changements de base.
En particulier, les composés
$S_i \times_S T \to T \to X$ sont égaux.
Comme $\coprod S_i \times_S T \to T$
est un épimorphisme (voir
par ex.\ le lemme \ref{descent-lemma-fpqc-universal-effective-epimorphisms}),
nous concluons que $a = b$.
\end{proof}

\begin{lemma}
\label{descent-lemma-descending-property-regular-immersion}
Les propriétés
\begin{enumerate}
\item[] $\mathcal{P}(f) =$``$f$ est une immersion Koszul-régulière'',
\item[] $\mathcal{P}(f) =$``$f$ est une immersion $H_1$-régulière'', et
\item[] $\mathcal{P}(f) =$``$f$ est une immersion quasi-régulière''
\end{enumerate}
sont locales sur la base pour la topologie fpqc.
\end{lemma}

\begin{proof}
Nous allons utiliser le critère du
lemme \ref{descent-lemma-descending-properties-morphisms}.
Par la définition \ref{divisors-definition-regular-immersion} de Diviseurs,
la propriété d'être une immersion Koszul-régulière
(resp.\ $H_1$-régulière, quasi-régulière) est locale pour la topologie de
Zariski sur la base. Par le
lemme \ref{divisors-lemma-flat-base-change-regular-immersion} de Diviseurs,
la propriété d'être une immersion Koszul-régulière
(resp.\ $H_1$-régulière, quasi-régulière)
est préservée par changement de base plat.
La dernière hypothèse (3) du
lemme \ref{descent-lemma-descending-properties-morphisms}
se traduit par l'énoncé algébrique suivant :
soit $A \to B$ un homomorphisme fidèlement plat d'anneaux.
Soit $I \subset A$ un idéal.
Si $IB$ est localement sur $\Spec(B)$ engendré par une suite Koszul-régulière
(resp.\ $H_1$-régulière, quasi-régulière) de $B$, alors $I \subset A$
est localement sur $\Spec(A)$ engendré par une suite Koszul-régulière
(resp.\ $H_1$-régulière, quasi-régulière) de $A$. C'est le
lemme \ref{more-algebra-lemma-flat-descent-regular-ideal}
de Compléments d'algèbre.
\end{proof}





\section{Propriétés des morphismes locales sur la cible pour la topologie fppf}
\label{descent-section-descending-properties-morphisms-fppf}

\noindent
Dans cette section, nous déterminons certaines propriétés des morphismes de
schémas dont nous n'avons pas (encore) pu montrer qu'elles sont locales sur la
base pour la topologie fpqc, mais qui le sont pour la topologie fppf.

\begin{lemma}
\label{descent-lemma-descending-fppf-property-immersion}
La propriété $\mathcal{P}(f) =$``$f$ est une immersion''
est locale sur la base pour la topologie fppf.
\end{lemma}

\begin{proof}
La propriété d'être une immersion est stable par changement de base,
voir le lemme \ref{schemes-lemma-base-change-immersion} de Schémas.
Elle est locale pour la topologie de Zariski sur la base.
Enfin, soit
$\pi : S' \to S$ un morphisme surjectif de schémas affines,
plat et localement de présentation finie.
Notons que $\pi : S' \to S$ est ouvert par
le lemme \ref{morphisms-lemma-fppf-open} de Morphismes.
Soit $f : X \to S$ un morphisme.
Supposons que le changement de base $f' : X' \to S'$ soit une immersion.
En particulier, $f'(X') = \pi^{-1}(f(X))$ est localement fermé.
Par le lemme \ref{topology-lemma-open-morphism-quotient-topology} de Topologie,
il s'ensuit que $f(X) \subset S$
est localement fermé. Soit $Z \subset S$
le fermé $Z = \overline{f(X)} \setminus f(X)$.
Une nouvelle application du lemme
\ref{topology-lemma-open-morphism-quotient-topology} de Topologie
montre que $f'(X')$ est fermé dans $S' \setminus Z'$.
Nous pouvons donc appliquer le
lemme \ref{descent-lemma-descending-property-closed-immersion}
au recouvrement fpqc $\{S' \setminus Z' \to S \setminus Z\}$
et conclure que $f : X \to S \setminus Z$ est une immersion fermée.
Autrement dit, $f$ est une immersion.
Le lemme \ref{descent-lemma-descending-properties-morphisms} s'applique donc et conclut.
\end{proof}

\begin{lemma}
\label{descent-lemma-descending-fppf-property-dominant}
La propriété $\mathcal{P}(f) =$``$f$ est dominant''
est locale sur la base pour la topologie fppf.
\end{lemma}

\begin{proof}
Par le lemme \ref{morphisms-lemma-open-base-change-dominant} de Morphismes,
les morphismes dominants sont préservés
par changement de base par des morphismes ouverts, donc
par des morphismes plats localement de présentation finie
(voir le lemme \ref{morphisms-lemma-fppf-open} de Morphismes).
La réciproque est plus facile.
En effet, la propriété d'être dominant est manifestement locale pour la
topologie de Zariski sur la base.
Soit ensuite $S' \to S$ un morphisme surjectif de schémas
(pas nécessairement plat ni localement de présentation finie),
et soit $f : X \to S$ un morphisme. Supposons que le changement de base
$f' : X' \to S'$ soit dominant. Alors $X' \to S'\to S$ est un composé
de deux morphismes dominants, donc est dominant. Comme il s'agit aussi du
composé $X' \to X \to S$, il s'ensuit que $X\to S$ est dominant.
\end{proof}






\section[Application de la descente fpqc]{Application de la descente fpqc des propriétés des morphismes}
\label{descent-section-application-descending-properties-morphisms}

\noindent
Le lemme suivant peut sembler quelque peu anodin, mais il constitue un outil
utile pour étudier les propriétés, pour un morphisme, d'être \'etale ou non
ramifié.

\begin{lemma}
\label{descent-lemma-flat-surjective-quasi-compact-monomorphism-isomorphism}
Soit $f : X \to Y$ un monomorphisme plat, quasi-compact et surjectif.
Alors ce morphisme est un isomorphisme.
\end{lemma}

\begin{proof}
Puisque $f$ est un morphisme plat, quasi-compact et surjectif,
$\{X \to Y\}$ est un recouvrement fpqc de $Y$.
La diagonale $\Delta : X \to X \times_Y X$ est un isomorphisme
(voir le lemme \ref{schemes-lemma-monomorphism} de Schémas).
Il s'ensuit que le changement de base de $f$ par $f$ est un
isomorphisme. Ainsi, $f$ est un isomorphisme par le
lemme \ref{descent-lemma-descending-property-isomorphism}.
\end{proof}

\noindent
Ce lemme permet de démontrer le résultat important suivant ; nous donnons aussi
une démonstration qui évite la descente fpqc.
Nous discuterons plus en détail de ce résultat et de résultats connexes dans
la section \ref{etale-section-topological-etale} du chapitre Morphismes étales.

\begin{lemma}
\label{descent-lemma-universally-injective-etale-open-immersion}
Tout morphisme universellement injectif et \'etale est une immersion ouverte.
\end{lemma}

\begin{proof}[Première démonstration]
Soit $f : X \to Y$ un morphisme \'etale universellement injectif.
Alors $f$ est ouvert
(voir le lemme \ref{morphisms-lemma-etale-open} de Morphismes) ;
nous pouvons donc remplacer $Y$ par $f(X)$ et supposer $f$ surjectif.
Le morphisme $f$ est alors bijectif et ouvert, donc est un homéomorphisme.
En particulier, $f$ est quasi-compact. Par le
lemme \ref{descent-lemma-flat-surjective-quasi-compact-monomorphism-isomorphism},
il suffit donc de montrer que $f$ est un monomorphisme.
Comme $X \to Y$ est \'etale, le morphisme
$\Delta_{X/Y} : X \to X \times_Y X$ est une immersion ouverte par
le lemme \ref{morphisms-lemma-diagonal-unramified-morphism} de Morphismes
(et le lemme \ref{morphisms-lemma-flat-unramified-etale} de Morphismes).
Comme $f$ est universellement injectif, $\Delta_{X/Y}$ est aussi surjectif,
voir le lemme \ref{morphisms-lemma-universally-injective} de Morphismes.
Ainsi, $\Delta_{X/Y}$ est un isomorphisme, c'est-à-dire que
$X \to Y$ est un monomorphisme.
\end{proof}

\begin{proof}[Seconde démonstration]
Soit $f : X \to Y$ un morphisme \'etale universellement injectif.
Alors $f$ est ouvert (voir le lemme \ref{morphisms-lemma-etale-open} de Morphismes) ;
nous pouvons donc remplacer $Y$ par $f(X)$ et supposer $f$ surjectif.
Les hypothèses restant satisfaites après tout changement de base, nous en
concluons que $f$ est un homéomorphisme universel. Par conséquent, $f$ est
entier, voir le
lemme \ref{morphisms-lemma-universal-homeomorphism} de Morphismes.
Il s'ensuit que $f$ est fini par
le lemme \ref{morphisms-lemma-finite-integral} de Morphismes.
Il s'ensuit que $f$ est fini localement libre par
le lemme \ref{morphisms-lemma-finite-flat} de Morphismes.
Pour achever la démonstration, il suffit de voir que $f$ est fini localement
libre de degré $1$ (un morphisme fini localement libre de degré $1$
est un isomorphisme).
Il existe une décomposition de $Y$ en sous-schémas ouverts et fermés
$V_d$ telle que $f^{-1}(V_d) \to V_d$ soit fini localement libre de
degré $d$, voir le lemme \ref{morphisms-lemma-finite-locally-free} de Morphismes.
Si $V_d$ est non vide, nous pouvons choisir un morphisme
$\Spec(k) \to V_d \subset Y$, où $k$ est un corps algébriquement clos
(il suffit de prendre la clôture algébrique du corps résiduel d'un point de
$V_d$).
Alors $\Spec(k) \times_Y X \to \Spec(k)$ est une réunion disjointe de
copies de $\Spec(k)$, par le
lemme \ref{morphisms-lemma-etale-over-field} de Morphismes,
et parce que $k$ est algébriquement clos.
Mais, comme $f$ est universellement injectif, il ne peut y avoir qu'une seule
copie; ainsi $d = 1$, comme voulu.
\end{proof}

\noindent
La caractérisation suivante des morphismes plats universellement injectifs
permet de reformuler légèrement les hypothèses du lemme précédent.

\begin{lemma}
\label{descent-lemma-flat-universally-injective}
Soit $f : X \to Y$ un morphisme de schémas. Notons $X^0$ l'ensemble
des points génériques des composantes irréductibles de $X$. Si
\begin{enumerate}
\item $f$ est plat et séparé,
\item pour $\xi \in X^0$, nous avons $\kappa(f(\xi)) = \kappa(\xi)$, et
\item si $\xi, \xi' \in X^0$, $\xi \not = \xi'$, alors
$f(\xi) \not = f(\xi')$,
\end{enumerate}
alors $f$ est universellement injectif.
\end{lemma}

\begin{proof}
Il faut montrer que $\Delta : X \to X \times_Y X$ est surjectif, voir le
lemme \ref{morphisms-lemma-universally-injective} de Morphismes.
Comme $X \to Y$ est séparé, l'image de $\Delta$ est fermée.
Ainsi, si $\Delta$ n'est pas surjectif, il existe un point générique
$\eta \in X \times_S X$ d'une composante irréductible de $X \times_S X$
qui n'appartient pas à l'image de $\Delta$. La projection
$\text{pr}_1 : X \times_Y X \to X$
est plate, car elle se déduit du morphisme plat $X \to Y$ par changement de
base, voir le lemme \ref{morphisms-lemma-base-change-flat} de Morphismes.
Les générisations se relèvent donc le long de $\text{pr}_1$, voir le
lemme \ref{morphisms-lemma-generalizations-lift-flat} de Morphismes.
Nous en concluons que $\xi = \text{pr}_1(\eta) \in X^0$.
Cependant, les hypothèses (2) et (3) garantissent que le schéma
$(X \times_Y X)_{f(\xi)}$ possède au plus un point pour tout $\xi \in X^0$.
Autrement dit, $\Delta(\xi) = \eta$, contradiction.
\end{proof}

\noindent
Nous pouvons donc reformuler le
lemme \ref{descent-lemma-universally-injective-etale-open-immersion} comme suit.

\begin{lemma}
\label{descent-lemma-characterize-open-immersion}
Soit $f : X \to Y$ un morphisme de schémas. Notons $X^0$ l'ensemble
des points génériques des composantes irréductibles de $X$. Si
\begin{enumerate}
\item $f$ est \'etale et séparé,
\item pour $\xi \in X^0$, nous avons $\kappa(f(\xi)) = \kappa(\xi)$, et
\item si $\xi, \xi' \in X^0$, $\xi \not = \xi'$, alors
$f(\xi) \not = f(\xi')$,
\end{enumerate}
alors $f$ est une immersion ouverte.
\end{lemma}

\begin{proof}
Cela résulte immédiatement des lemmes
\ref{descent-lemma-flat-universally-injective} et
\ref{descent-lemma-universally-injective-etale-open-immersion}.
\end{proof}

\begin{lemma}
\label{descent-lemma-descending-property-proper-over-base}
Soit $f : X \to Y$ un morphisme de schémas localement de type fini.
Soit $Z$ un fermé de $X$. S'il existe un recouvrement fpqc
$\{Y_i \to Y\}$ tel que l'image inverse $Z_i \subset Y_i \times_Y X$
soit propre sur $Y_i$
(voir la définition \ref{coherent-definition-proper-over-base}
de Cohomologie des schémas),
alors $Z$ est propre sur $Y$.
\end{lemma}

\begin{proof}
Munissons $Z$ de la structure réduite induite de sous-schéma fermé, voir la
définition \ref{schemes-definition-reduced-induced-scheme} de Schémas.
Pour tout $i$, le changement de base $Y_i \times_Y Z$ est un sous-schéma
fermé de $Y_i \times_Y X$ dont le fermé sous-jacent est $Z_i$.
Par définition (au moyen du lemme
\ref{coherent-lemma-closed-proper-over-base} de Cohomologie des schémas),
nous en concluons que les projections $Y_i \times_Y Z \to Y_i$ sont des
morphismes propres. Le morphisme $Z \to Y$ est donc propre par le
lemme \ref{descent-lemma-descending-property-proper}.
Ainsi, $Z$ est propre sur $Y$ par définition.
\end{proof}

\begin{lemma}
\label{descent-lemma-descending-property-ample}
Soit $f : X \to S$ un morphisme de schémas.
Soit $\mathcal{L}$ un $\mathcal{O}_X$-module inversible.
Soit $\{g_i : S_i \to S\}_{i \in I}$ un recouvrement fpqc.
Soit $f_i : X_i \to S_i$ le changement de base de $f$, et soit
$\mathcal{L}_i$ l'image inverse de $\mathcal{L}$ sur $X_i$.
Les assertions suivantes sont équivalentes :
\begin{enumerate}
\item $\mathcal{L}$ est ample sur $X/S$, et
\item $\mathcal{L}_i$ est ample sur $X_i/S_i$
pour tout $i \in I$.
\end{enumerate}
\end{lemma}

\begin{proof}
L'implication (1) $\Rightarrow$ (2) résulte du
lemme \ref{morphisms-lemma-ample-base-change} de Morphismes.
Supposons que $\mathcal{L}_i$ soit ample sur $X_i/S_i$ pour tout $i \in I$.
Par la définition \ref{morphisms-definition-relatively-ample} de Morphismes,
cela implique que $X_i \to S_i$ est quasi-compact et, par
le lemme \ref{morphisms-lemma-relatively-ample-separated} de Morphismes,
que $X_i \to S$ est séparé.
Ainsi, $f$ est quasi-compact et séparé par les
lemmes \ref{descent-lemma-descending-property-quasi-compact} et
\ref{descent-lemma-descending-property-separated}.

\medskip\noindent
Cela signifie que
$\mathcal{A} = \bigoplus_{d \geq 0} f_*\mathcal{L}^{\otimes d}$
est une $\mathcal{O}_S$-algèbre graduée quasi-cohérente
(voir le lemme \ref{schemes-lemma-push-forward-quasi-coherent} de Schémas).
De plus, la formation de $\mathcal{A}$ commute aux changements de base plats par
le lemme \ref{coherent-lemma-flat-base-change-cohomology}
de Cohomologie des schémas.
En particulier, si nous posons
$\mathcal{A}_i = \bigoplus_{d \geq 0} f_{i, *}\mathcal{L}_i^{\otimes d}$,
alors $\mathcal{A}_i = g_i^*\mathcal{A}$.
Il s'ensuit que les morphismes naturels
$\psi_d : f^*\mathcal{A}_d \to \mathcal{L}^{\otimes d}$
de $\mathcal{O}_X$-modules donnent par image inverse les morphismes naturels
$\psi_{i, d} : f_i^*(\mathcal{A}_i)_d \to \mathcal{L}_i^{\otimes d}$
de $\mathcal{O}_{X_i}$-modules. Comme $\mathcal{L}_i$ est ample sur $X_i/S_i$,
pour tout point $x_i \in X_i$, il existe un $d \geq 1$ tel que
$f_i^*(\mathcal{A}_i)_d \to \mathcal{L}_i^{\otimes d}$
soit surjectif sur les germes en $x_i$. Cela résulte soit directement
de la définition d'un module relativement ample, soit de
le lemme \ref{morphisms-lemma-characterize-relatively-ample} de Morphismes.
Si $x \in X$, nous pouvons choisir un $i$ et un $x_i \in X_i$
d'image $x$. Comme $\mathcal{O}_{X, x} \to \mathcal{O}_{X_i, x_i}$
est plat, donc fidèlement plat, nous en concluons que pour tout $x \in X$,
il existe un $d \geq 1$ tel que
$f^*\mathcal{A}_d \to \mathcal{L}^{\otimes d}$
soit surjectif sur les germes en $x$.
Cela implique que l'ouvert $U(\psi) \subset X$ du
lemme \ref{constructions-lemma-invertible-map-into-relative-proj} de Constructions,
associé au morphisme
$\psi : f^*\mathcal{A} \to \bigoplus_{d \geq 0} \mathcal{L}^{\otimes d}$
d'algèbres graduées sur $\mathcal{O}_X$,
est égal à $X$. Considérons le morphisme correspondant
$$
r_{\mathcal{L}, \psi} : X \longrightarrow \underline{\text{Proj}}_S(\mathcal{A})
$$
Il résulte clairement de ce qui précède que le changement de base de
$r_{\mathcal{L}, \psi}$ à $S_i$ est le morphisme
$r_{\mathcal{L}_i, \psi_i}$, qui est une immersion ouverte par le
lemme \ref{morphisms-lemma-characterize-relatively-ample} de Morphismes.
Ainsi, $r_{\mathcal{L}, \psi}$ est une immersion ouverte par le
lemme \ref{descent-lemma-descending-property-open-immersion},
et nous concluons que $\mathcal{L}$ est ample sur $X/S$ par
le lemme \ref{morphisms-lemma-characterize-relatively-ample} de Morphismes.
\end{proof}















\section{Propriétés des morphismes locales sur la source}
\label{descent-section-properties-morphisms-local-source}

\noindent
Il arrive souvent que l'on puisse établir qu'un morphisme possède une certaine
propriété après l'avoir précomposé par un autre morphisme. Dans de nombreux cas,
cela implique que le morphisme initial possède lui aussi cette propriété.
Nous formalisons cette situation dans la définition suivante.

\begin{definition}
\label{descent-definition-property-morphisms-local-source}
Soit $\mathcal{P}$ une propriété de morphismes de schémas.
Soit $\tau \in \{Zariski, \linebreak[0] fpqc, \linebreak[0] fppf, \linebreak[0]
\etale, \linebreak[0] smooth, \linebreak[0] syntomic\}$.
On dit que $\mathcal{P}$ est
{\it locale sur la source pour $\tau$}, ou
{\it locale sur la source pour la topologie $\tau$}, si pour tout
morphisme de schémas $f : X \to Y$ au-dessus de $S$, et pour tout
recouvrement pour la topologie $\tau$, $\{X_i \to X\}_{i \in I}$, on a
$$
f \text{ vérifie }\mathcal{P}
\Leftrightarrow
\text{chaque }X_i \to Y\text{ vérifie }\mathcal{P}.
$$
\end{definition}

\noindent
Bien entendu, puisque les isomorphismes sont toujours des recouvrements,
nous constatons (ou exigeons) que la propriété $\mathcal{P}$ est vérifiée par
$X \to Y$ si et seulement si elle l'est par toute flèche $X' \to Y'$
isomorphe à $X \to Y$. Si une propriété est locale sur la source pour $\tau$,
elle est préservée par précomposition avec les morphismes qui interviennent
dans les recouvrements pour la topologie $\tau$. Voici un énoncé précis.

\begin{lemma}
\label{descent-lemma-precompose-property-local-source}
Soit $\tau \in \{fpqc, fppf, syntomic, smooth, \etale, Zariski\}$.
Soit $\mathcal{P}$ une propriété de morphismes locale sur la source pour
$\tau$. Soit $f : X \to Y$ vérifiant la propriété $\mathcal{P}$.
Pour tout morphisme $a : X' \to X$ qui est
plat, resp.\ plat et localement de présentation finie, resp.\ syntomique,
resp.\ \'etale, resp.\ une immersion ouverte, le composé
$f \circ a : X' \to Y$ vérifie la propriété $\mathcal{P}$.
\end{lemma}

\begin{proof}
Cela résulte du fait que l'on peut compléter $X' \to X$ en une famille de
morphismes qui forme un $\tau$-recouvrement.
\end{proof}

\begin{lemma}
\label{descent-lemma-largest-open-of-the-source}
Soit $\tau \in \{fppf, syntomic, smooth, \etale\}$.
Soit $\mathcal{P}$ une propriété de morphismes locale sur la source pour
$\tau$. Pour tout morphisme de schémas $f : X \to Y$, il existe
un plus grand ouvert $W(f) \subset X$ tel que la restriction
$f|_{W(f)} : W(f) \to Y$ vérifie $\mathcal{P}$. De plus,
si $g : X' \to X$ est plat et localement de présentation finie,
syntomique, lisse ou \'etale, et si $f' = f \circ g : X' \to Y$, alors
$g^{-1}(W(f)) = W(f')$.
\end{lemma}

\begin{proof}
Considérons la réunion $W$ des images $g(X') \subset X$ des
morphismes $g : X' \to X$ qui vérifient les propriétés suivantes :
\begin{enumerate}
\item $g$ est plat et localement de présentation finie, syntomique,
lisse ou \'etale, et
\item le composé $X' \to X \to Y$ vérifie la propriété $\mathcal{P}$.
\end{enumerate}
Comme tout tel morphisme $g$ est ouvert (voir le
lemme \ref{morphisms-lemma-fppf-open} de Morphismes),
$W \subset X$ est un ouvert de $X$. Puisque $\mathcal{P}$ est locale pour la
topologie $\tau$, la restriction $f|_W : W \to Y$ vérifie la propriété
$\mathcal{P}$ : nous disposons en effet d'un recouvrement pour la topologie $\tau$,
$\{X' \to W\}$ de $W$ dont les images inverses vérifient $\mathcal{P}$.
Cela démontre l'existence de $W(f)$.
La compatibilité énoncée dans la dernière phrase résulte immédiatement de la
construction de $W(f)$.
\end{proof}

\begin{lemma}
\label{descent-lemma-properties-morphisms-local-source}
Soit $\mathcal{P}$ une propriété de morphismes de schémas.
Soit $\tau \in \{fpqc, \linebreak[0] fppf, \linebreak[0]
\etale, \linebreak[0] smooth, \linebreak[0] syntomic\}$.
Supposons que
\begin{enumerate}
\item la propriété soit préservée par précomposition avec un morphisme
plat, plat localement de présentation finie, \'etale, lisse ou syntomique
selon que $\tau$ est fpqc, fppf, \'etale, lisse ou syntomique,
\item la propriété soit locale sur la source pour la topologie de Zariski,
\item la propriété soit locale sur la cible pour la topologie de Zariski,
\item pour tout morphisme de schémas affines $f : X \to Y$, et pour tout
morphisme surjectif de schémas affines $X' \to X$ qui est plat, plat de
présentation finie, \'etale, lisse ou syntomique selon que $\tau$ est
fpqc, fppf, \'etale, lisse ou syntomique, la propriété $\mathcal{P}$ soit
vérifiée par $f$ si la propriété $\mathcal{P}$ est vérifiée par le composé
$f' : X' \to Y$.
\end{enumerate}
Alors $\mathcal{P}$ est locale sur la source pour la topologie $\tau$.
\end{lemma}

\begin{proof}
Cela résulte presque immédiatement de la définition d'un
recouvrement pour la topologie $\tau$ ; voir les définitions
\ref{topologies-definition-fpqc-covering}
\ref{topologies-definition-fppf-covering}
\ref{topologies-definition-etale-covering}
\ref{topologies-definition-smooth-covering}, ou
\ref{topologies-definition-syntomic-covering} de Topologies,
et les lemmes
\ref{topologies-lemma-fpqc-affine},
\ref{topologies-lemma-fppf-affine},
\ref{topologies-lemma-etale-affine},
\ref{topologies-lemma-smooth-affine}, ou
\ref{topologies-lemma-syntomic-affine} de Topologies.
Détails omis. (Indication : utiliser la localité sur la source et sur le but
pour ramener la vérification de la propriété $\mathcal{P}$ au cas d'un
morphisme entre schémas affines. Appliquer ensuite (1) et (4).)
\end{proof}

\begin{remark}
\label{descent-remark-properties-morphisms-local-source-standard}
(Cette remarque reprend les remarques
\ref{descent-remark-descending-properties-standard}
et \ref{descent-remark-descending-properties-morphisms-standard} ci-dessus.)
Dans le lemme \ref{descent-lemma-properties-morphisms-local-source}, si
$\tau = smooth$, on peut supposer, dans la condition (4), que le morphisme est
un morphisme lisse standard (surjectif).
De même lorsque $\tau = syntomic$ ou $\tau = \etale$.
\end{remark}



\section[Localité sur la source (fpqc)]{Propriétés des morphismes locales sur la source pour la topologie fpqc}
\label{descent-section-fpqc-local-source}

\noindent
Voici quelques propriétés de morphismes qui sont locales sur la source pour la
topologie fpqc.

\begin{lemma}
\label{descent-lemma-flat-fpqc-local-source}
La propriété $\mathcal{P}(f)=$``$f$ est plat'' est locale sur la source pour la
topologie fpqc.
\end{lemma}

\begin{proof}
Puisque la platitude se définit au moyen des homomorphismes d'anneaux locaux
(voir la définition \ref{morphisms-definition-flat} de Morphismes),
il faut établir le fait algébrique suivant : supposons que
$A \to B \to C$ soient des homomorphismes locaux d'anneaux locaux, et que
$B \to C$ soit plat. Alors $A \to B$ est plat si et seulement si
$A \to C$ est plat.
Si $A \to B$ est plat, alors $A \to C$ est plat par
le lemme \ref{algebra-lemma-composition-flat} d'Algèbre.
Réciproquement, supposons $A \to C$ plat.
Notons que $B \to C$ est fidèlement plat, voir
le lemme \ref{algebra-lemma-local-flat-ff} d'Algèbre.
Ainsi, $A \to B$ est plat par
le lemme \ref{algebra-lemma-flat-permanence} d'Algèbre.
(Voir aussi le lemme \ref{morphisms-lemma-flat-permanence} de Morphismes
pour une démonstration directe.)
\end{proof}

\begin{lemma}
\label{descent-lemma-injective-local-rings-fpqc-local-source}
La propriété
$\mathcal{P}(f : X \to Y)=$``pour tout $x \in X$, l'homomorphisme d'anneaux
locaux $\mathcal{O}_{Y, f(x)} \to \mathcal{O}_{X, x}$ est injectif''
est locale sur la source pour la topologie fpqc.
\end{lemma}

\begin{proof}
Omis. Il s'agit seulement d'une tentative (probablement malavisée) de badiner.
\end{proof}





\section{Propriétés des morphismes locales sur la source pour la topologie fppf}
\label{descent-section-fppf-local-source}

\noindent
Voici quelques propriétés de morphismes qui sont locales sur la source pour la
topologie fppf.

\begin{lemma}
\label{descent-lemma-locally-finite-presentation-fppf-local-source}
La propriété $\mathcal{P}(f)=$``$f$ est localement de présentation finie''
est locale sur la source pour la topologie fppf.
\end{lemma}

\begin{proof}
La propriété d'être localement de présentation finie est locale sur la source
et sur la cible pour la topologie de Zariski, voir le
lemme \ref{morphisms-lemma-locally-finite-presentation-characterize} de Morphismes.
Elle est préservée par composition, voir
le lemme \ref{morphisms-lemma-composition-finite-presentation} de Morphismes.
Cela démontre (1), (2) et (3) du
lemme \ref{descent-lemma-properties-morphisms-local-source}.
La dernière condition (4) est le
lemme \ref{descent-lemma-flat-finitely-presented-permanence-algebra}.
Le résultat en découle.
\end{proof}

\begin{lemma}
\label{descent-lemma-locally-finite-type-fppf-local-source}
La propriété $\mathcal{P}(f)=$``$f$ est localement de type fini''
est locale sur la source pour la topologie fppf.
\end{lemma}

\begin{proof}
La propriété d'être localement de type fini est locale sur la source et sur la
cible pour la topologie de Zariski, voir le
lemme \ref{morphisms-lemma-locally-finite-type-characterize} de Morphismes.
Elle est préservée par composition, voir
le lemme \ref{morphisms-lemma-composition-finite-type} de Morphismes,
et tout morphisme plat localement de présentation finie est localement de type
fini, voir
le lemme \ref{morphisms-lemma-finite-presentation-finite-type} de Morphismes.
Cela démontre (1), (2) et (3) du
lemme \ref{descent-lemma-properties-morphisms-local-source}.
La dernière condition (4) est le
lemme \ref{descent-lemma-finite-type-local-source-fppf-algebra}.
Le résultat en découle.
\end{proof}

\begin{lemma}
\label{descent-lemma-open-fppf-local-source}
La propriété $\mathcal{P}(f)=$``$f$ est ouvert''
est locale sur la source pour la topologie fppf.
\end{lemma}

\begin{proof}
La propriété d'être un morphisme ouvert est manifestement locale sur la source
et sur la cible pour la topologie de Zariski.
Elle est préservée par composition, voir
le lemme \ref{morphisms-lemma-composition-open} de Morphismes,
et tout morphisme plat de présentation finie est ouvert, voir
le lemme \ref{morphisms-lemma-fppf-open} de Morphismes.
Cela démontre (1), (2) et (3) du
lemme \ref{descent-lemma-properties-morphisms-local-source}.
La dernière condition (4) résulte de
lemme \ref{morphisms-lemma-fpqc-quotient-topology} de Morphismes.
Le résultat en découle.
\end{proof}

\begin{lemma}
\label{descent-lemma-universally-open-fppf-local-source}
La propriété $\mathcal{P}(f)=$``$f$ est universellement ouvert''
est locale sur la source pour la topologie fppf.
\end{lemma}

\begin{proof}
Soit $f : X \to Y$ un morphisme de schémas.
Soit $\{X_i \to X\}_{i \in I}$ un recouvrement fppf.
Notons $f_i : X_i \to X$ les composés. Il faut montrer que $f$ est
universellement ouvert si et seulement si chaque $f_i$ l'est.
Si $f$ est universellement ouvert, alors chaque $f_i$ l'est aussi, puisque les
morphismes $X_i \to X$ sont universellement ouverts et que les composés de
morphismes universellement ouverts sont universellement ouverts
(voir les lemmes \ref{morphisms-lemma-fppf-open}
et \ref{morphisms-lemma-composition-open} de Morphismes).
Réciproquement, supposons chaque $f_i$ universellement ouvert.
Soit $Y' \to Y$ un morphisme de schémas.
Notons $X' = Y' \times_Y X$ et $X'_i = Y' \times_Y X_i$.
Le système $\{X_i' \to X'\}_{i \in I}$ est encore un recouvrement fppf.
Les morphismes $f'_i : X_i' \to Y'$ sont ouverts par hypothèse.
Le lemme \ref{descent-lemma-open-fppf-local-source} montre donc que
$f' : X' \to Y'$ est ouvert, comme voulu.
\end{proof}



\section{Propriétés des morphismes locales sur la source pour la topologie syntomique}
\label{descent-section-syntomic-local-source}

\noindent
Voici quelques propriétés de morphismes qui sont locales sur la source pour la
topologie syntomique.

\begin{lemma}
\label{descent-lemma-syntomic-syntomic-local-source}
La propriété $\mathcal{P}(f)=$``$f$ est syntomique''
est locale sur la source pour la topologie syntomique.
\end{lemma}

\begin{proof}
Combiner le lemme \ref{descent-lemma-properties-morphisms-local-source} avec le
lemme \ref{morphisms-lemma-syntomic-characterize} de Morphismes
(localité sur la source et sur la cible pour la topologie de Zariski), le
lemme \ref{morphisms-lemma-composition-syntomic} de Morphismes
(précomposition), et le
lemme \ref{descent-lemma-syntomic-smooth-etale-permanence} (partie (4)).
\end{proof}




\section{Propriétés des morphismes locales sur la source pour la topologie lisse}
\label{descent-section-smooth-local-source}

\noindent
Voici quelques propriétés de morphismes qui sont locales sur la source pour la
topologie lisse. Notons aussi le résultat (à certains égards plus fort)
sur la descente de la lissité le long de morphismes plats,
lemme \ref{descent-lemma-smooth-permanence}.

\begin{lemma}
\label{descent-lemma-smooth-smooth-local-source}
La propriété $\mathcal{P}(f)=$``$f$ est lisse''
est locale sur la source pour la topologie lisse.
\end{lemma}

\begin{proof}
Combiner le lemme \ref{descent-lemma-properties-morphisms-local-source} avec le
lemme \ref{morphisms-lemma-smooth-characterize} de Morphismes
(localité sur la source et sur la cible pour la topologie de Zariski), le
lemme \ref{morphisms-lemma-composition-smooth} de Morphismes
(précomposition), et le
lemme \ref{descent-lemma-syntomic-smooth-etale-permanence} (partie (4)).
\end{proof}



\section[Localité sur la source (\'etale)]{Propriétés des morphismes locales sur la source pour la topologie \'etale}
\label{descent-section-etale-local-source}

\noindent
Voici quelques propriétés de morphismes qui sont locales sur la source pour la
topologie \'etale.

\begin{lemma}
\label{descent-lemma-etale-etale-local-source}
La propriété $\mathcal{P}(f)=$``$f$ est \'etale''
est locale sur la source pour la topologie \'etale.
\end{lemma}

\begin{proof}
Combiner le lemme \ref{descent-lemma-properties-morphisms-local-source} avec le
lemme \ref{morphisms-lemma-etale-characterize} de Morphismes
(localité sur la source et sur la cible pour la topologie de Zariski), le
lemme \ref{morphisms-lemma-composition-etale} de Morphismes
(précomposition), et le
lemme \ref{descent-lemma-syntomic-smooth-etale-permanence} (partie (4)).
\end{proof}

\begin{lemma}
\label{descent-lemma-locally-quasi-finite-etale-local-source}
La propriété $\mathcal{P}(f)=$``$f$ est localement quasi-fini''
est locale sur la source pour la topologie \'etale.
\end{lemma}

\begin{proof}
Nous allons utiliser le
lemme \ref{descent-lemma-properties-morphisms-local-source}.
Par le lemme \ref{morphisms-lemma-locally-quasi-finite-characterize} de
Morphismes, la propriété d'être localement quasi-fini est locale sur la source
et sur la cible pour la topologie de Zariski. Par les lemmes
\ref{morphisms-lemma-composition-quasi-finite} et
\ref{morphisms-lemma-etale-locally-quasi-finite} de Morphismes,
la précomposition d'un morphisme localement quasi-fini par un morphisme
\'etale est localement quasi-finie.
Enfin, supposons que $X \to Y$ soit un morphisme de schémas affines et que
$X' \to X$ soit un morphisme surjectif \'etale de schémas affines tel que
$X' \to Y$ soit localement quasi-fini. Alors $X' \to Y$ est de type fini, et
le lemme \ref{descent-lemma-finite-type-local-source-fppf-algebra}
montre que $X \to Y$ est lui aussi de type fini.
De plus, par hypothèse, $X' \to Y$ a des fibres finies ; il en va donc de même
de $X \to Y$. Nous concluons que $X \to Y$ est quasi-fini par
le lemme \ref{morphisms-lemma-quasi-finite} de Morphismes.
Cela démontre la dernière hypothèse du
lemme \ref{descent-lemma-properties-morphisms-local-source}
et achève la démonstration.
\end{proof}

\begin{lemma}
\label{descent-lemma-unramified-etale-local-source}
La propriété $\mathcal{P}(f)=$``$f$ est non ramifié''
est locale sur la source pour la topologie \'etale.
La propriété $\mathcal{P}(f)=$``$f$ est G-non ramifié''
est locale sur la source pour la topologie \'etale.
\end{lemma}

\begin{proof}
Nous allons utiliser le
lemme \ref{descent-lemma-properties-morphisms-local-source}.
Par le lemme \ref{morphisms-lemma-unramified-characterize} de Morphismes,
la propriété d'être non ramifié (resp.\ G-non ramifié)
est locale sur la source et sur la cible pour la topologie de Zariski. Par les
lemmes \ref{morphisms-lemma-composition-unramified} et
\ref{morphisms-lemma-etale-smooth-unramified} de Morphismes,
la précomposition d'un morphisme non ramifié (resp.\ G-non ramifié)
par un morphisme \'etale donne un morphisme non ramifié
(resp.\ G-non ramifié).
Enfin, supposons que $X \to Y$ soit un morphisme de schémas affines et que
$f : X' \to X$ soit un morphisme surjectif \'etale de schémas affines tel que
$X' \to Y$ soit non ramifié (resp.\ G-non ramifié).
Alors $X' \to Y$ est de type fini (resp.\ de présentation finie), et le
lemme \ref{descent-lemma-finite-type-local-source-fppf-algebra}
(resp.\ le lemme \ref{descent-lemma-flat-finitely-presented-permanence-algebra})
montre que $X \to Y$ est lui aussi de type fini
(resp.\ de présentation finie). Par
le lemme \ref{morphisms-lemma-triangle-differentials-smooth} de Morphismes,
nous avons une suite exacte courte
$$
0 \to f^*\Omega_{X/Y} \to \Omega_{X'/Y} \to \Omega_{X'/X} \to 0.
$$
Comme $X' \to Y$ est non ramifié, le terme du milieu est nul.
Puisque $f$ est fidèlement plat, il s'ensuit que $\Omega_{X/Y} = 0$.
Ainsi, $X \to Y$ est non ramifié (resp.\ G-non ramifié), voir le
lemme \ref{morphisms-lemma-unramified-omega-zero} de Morphismes.
Cela démontre la dernière hypothèse du
lemme \ref{descent-lemma-properties-morphisms-local-source}
et achève la démonstration.
\end{proof}




\section[Localité source-cible (\'etale)]{Propriétés des morphismes locales sur la source et la cible pour la topologie \'etale}
\label{descent-section-properties-etale-local-source-target}

\noindent
Soit $\mathcal{P}$ une propriété de morphismes de schémas. La phrase
``$\mathcal{P}$ est locale sur la source et la cible pour la topologie \'etale''
a un sens intuitif. Il s'avère toutefois que cette notion ne revient pas à
demander que $\mathcal{P}$ soit à la fois locale sur la source et locale sur
la cible pour la topologie \'etale.
Avant d'approfondir cette question, donnons deux exemples volontairement
élémentaires.

\begin{example}
\label{descent-example-silly-one}
Considérons la propriété $\mathcal{P}$ des morphismes de schémas définie par
la règle $\mathcal{P}(X \to Y) = $``$Y$ est localement noethérien''.
Le lecteur vérifiera qu'elle est locale sur la source et locale sur la cible
pour la topologie \'etale (démonstration omise, voir le
lemme \ref{descent-lemma-Noetherian-local-fppf}).
Mais il n'est {\bf pas} vrai que, si $f : X \to Y$ vérifie $\mathcal{P}$
et si $g : Y \to Z$ est \'etale, alors $g \circ f$ vérifie $\mathcal{P}$.
En effet, $f$ pourrait être l'identité de $Y$ et $g$ une immersion ouverte
d'un schéma localement noethérien $Y$ dans un schéma non localement
noethérien $Z$.
\end{example}

\noindent
L'exemple suivant est, en un sens, encore pire.

\begin{example}
\label{descent-example-silly-two}
Considérons la propriété $\mathcal{P}$ des morphismes de schémas définie par
la règle $\mathcal{P}(f : X \to Y) = $``pour tout $y \in Y$ qui est une
spécialisation d'un certain $f(x)$, $x \in X$, l'anneau local
$\mathcal{O}_{Y, y}$ est noethérien''.
Vérifions qu'elle est locale sur la source et sur la cible pour la topologie
\'etale. Nous utiliserons librement le
lemme \ref{schemes-lemma-specialize-points} de Schémas.

\medskip\noindent
Localité sur la cible : soit $\{g_i : Y_i \to Y\}$ un recouvrement pour la
topologie \'etale.
Soit $f_i : X_i \to Y_i$ le changement de base de $f$, et notons
$h_i : X_i \to X$ la projection. Supposons $\mathcal{P}(f)$.
Soit $f(x_i) \leadsto y_i$ une spécialisation.
Alors $f(h_i(x_i)) \leadsto g_i(y_i)$; ainsi, $\mathcal{P}(f)$ implique que
$\mathcal{O}_{Y, g_i(y_i)}$ est noethérien.
De plus, $\mathcal{O}_{Y, g_i(y_i)} \to \mathcal{O}_{Y_i, y_i}$ est une
localisation d'un homomorphisme \'etale d'anneaux.
Par conséquent, $\mathcal{O}_{Y_i, y_i}$ est noethérien par le
lemme \ref{algebra-lemma-Noetherian-permanence} d'Algèbre.
Réciproquement, supposons $\mathcal{P}(f_i)$ pour tout $i$.
Soit $f(x) \leadsto y$ une spécialisation. Choisissons un $i$ et un
$y_i \in Y_i$ d'image $y$.
Comme $x$ peut être considéré comme un point de
$\Spec(\mathcal{O}_{Y, y}) \times_Y X$ et que
$\mathcal{O}_{Y, y} \to \mathcal{O}_{Y_i, y_i}$ est fidèlement plat,
il existe un point
$x_i \in \Spec(\mathcal{O}_{Y_i, y_i}) \times_Y X$
d'image $x$. Alors $x_i \in X_i$, et $f_i(x_i)$ se spécialise en $y_i$.
Ainsi, nous voyons que $\mathcal{O}_{Y_i, y_i}$ est noethérien par
$\mathcal{P}(f_i)$, ce qui implique que $\mathcal{O}_{Y, y}$ est noethérien par
le lemme \ref{algebra-lemma-descent-Noetherian} d'Algèbre.

\medskip\noindent
Localité sur la source : soit $\{h_i : X_i \to X\}$ un recouvrement pour la
topologie \'etale.
Soit $f_i : X_i \to Y$ le composé $f \circ h_i$.
Supposons $\mathcal{P}(f)$. Soit $f(x_i) \leadsto y$ une spécialisation.
Alors $f(h_i(x_i)) \leadsto y$; ainsi, $\mathcal{P}(f)$ implique que
$\mathcal{O}_{Y, y}$ est noethérien. Par conséquent, $\mathcal{P}(f_i)$ est
vérifiée. Réciproquement, supposons $\mathcal{P}(f_i)$ pour tout $i$.
Soit $f(x) \leadsto y$ une spécialisation. Choisissons un $i$ et un
$x_i \in X_i$ d'image $x$.
Alors $y$ est une spécialisation de $f_i(x_i) = f(x)$.
Par conséquent, $\mathcal{P}(f_i)$ implique que $\mathcal{O}_{Y, y}$ est
noethérien, comme voulu.

\medskip\noindent
Nous affirmons qu'il existe un diagramme commutatif
$$
\xymatrix{
U \ar[d]_a \ar[r]_h & V \ar[d]^b \\
X \ar[r]^f & Y
}
$$
dont les flèches verticales sont surjectives, chacune étant \'etale, tel que $h$ vérifie
$\mathcal{P}$ et que $f$ ne vérifie pas $\mathcal{P}$. En effet, posons
$$
Y =
\Spec\Big(
\mathbf{C}[x_n; n \in \mathbf{Z}]/(x_n x_m; n \not = m)
\Big)
$$
et prenons pour $X \subset Y$ le sous-schéma ouvert complémentaire du point
dont toutes les coordonnées $x_n = 0$. Posons $U = X$ et
$V = X \amalg Y$; soient $a, b$ les morphismes évidents, et soit
$h : U \to V$ l'inclusion de $U = X$ dans la première composante de $V$.
L'affirmation précédente est vraie parce que $U$ est localement noethérien,
mais que $Y$ ne l'est pas.
\end{example}

\noindent
Quelle devrait être la bonne notion pour une propriété locale sur la source et
la cible pour la topologie \'etale ? Par analogie avec la
définition \ref{morphisms-definition-property-local} de Morphismes, nous pensons
que ce devrait être la suivante.

\begin{definition}
\label{descent-definition-local-source-target}
Soit $\mathcal{P}$ une propriété de morphismes de schémas.
Nous disons que $\mathcal{P}$ est {\it locale sur la source et la cible pour
la topologie \'etale} si
\begin{enumerate}
\item (stabilité par précomposition avec des morphismes étales)
si $f : X \to Y$ est \'etale et si $g : Y \to Z$ possède $\mathcal{P}$,
alors $g \circ f$ possède $\mathcal{P}$,
\item (stabilité par changement de base \'etale)
si $f : X \to Y$ possède $\mathcal{P}$ et si $Y' \to Y$ est \'etale, alors
le changement de base $f' : Y' \times_Y X \to Y'$ possède $\mathcal{P}$, et
\item (localité) pour un morphisme $f : X \to Y$, les assertions suivantes
sont équivalentes :
\begin{enumerate}
\item $f$ possède $\mathcal{P}$,
\item pour tout $x \in X$, il existe un diagramme commutatif
$$
\xymatrix{
U \ar[d]_a \ar[r]_h & V \ar[d]^b \\
X \ar[r]^f & Y
}
$$
dont les flèches verticales sont étales et un $u \in U$ tel que
$a(u) = x$, de sorte que $h$ possède $\mathcal{P}$.
\end{enumerate}
\end{enumerate}
\end{definition}

\noindent
Cette définition exclut en fait le comportement observé dans les
exemples \ref{descent-example-silly-one} et \ref{descent-example-silly-two}.
Nous la comparerons à celle de l'article \cite{DM} de Deligne et Mumford dans
la remarque \ref{descent-remark-compare-definitions}.
En outre, une propriété locale sur la source et la cible pour la topologie
\'etale est locale sur la source et locale sur la cible pour cette topologie.
Enfin, la réciproque est presque vraie, comme nous le verrons dans le
lemme \ref{descent-lemma-etale-local-source-target}.

\begin{lemma}
\label{descent-lemma-local-source-target-implies}
Soit $\mathcal{P}$ une propriété de morphismes de schémas qui est locale sur
la source et la cible pour la topologie \'etale. Alors
\begin{enumerate}
\item $\mathcal{P}$ est locale sur la source pour la topologie \'etale,
\item $\mathcal{P}$ est locale sur la cible pour la topologie \'etale,
\item $\mathcal{P}$ est stable par postcomposition avec des morphismes
\'etale : si $f : X \to Y$ possède $\mathcal{P}$ et si $g : Y \to Z$ est
\'etale, alors $g \circ f$ possède $\mathcal{P}$, et
\item $\mathcal{P}$ possède la propriété de permanence suivante : étant donnés
$f : X \to Y$ et $g : Y \to Z$ de type \'etale tels que $g \circ f$ possède
$\mathcal{P}$, le morphisme $f$ possède $\mathcal{P}$.
\end{enumerate}
\end{lemma}

\begin{proof}
Donnons tous les détails.

\medskip\noindent
Preuve de (1). Soit $f : X \to Y$ un morphisme de schémas.
Soit $\{X_i \to X\}_{i \in I}$ un recouvrement de $X$ pour la topologie
\'etale. Si chaque
composé $h_i : X_i \to Y$ possède $\mathcal{P}$, alors, pour tout $x \in X$,
on peut trouver un $i \in I$ et un point $x_i \in X_i$ d'image $x$. Le
morphisme de germes $(X_i, x_i) \to (X, x)$ est alors \'etale, le morphisme
$\text{id}_Y : Y \to Y$ est \'etale, et $h_i$ est comme dans la partie (3) de
la définition \ref{descent-definition-local-source-target}.
Nous voyons donc que $f$ possède $\mathcal{P}$.
Réciproquement, si $f$ possède $\mathcal{P}$, alors chaque $X_i \to Y$ possède
$\mathcal{P}$ d'après la partie (1) de la
définition \ref{descent-definition-local-source-target}.

\medskip\noindent
Preuve de (2). Soit $f : X \to Y$ un morphisme de schémas.
Soit $\{Y_i \to Y\}_{i \in I}$ un recouvrement de $Y$ pour la topologie
\'etale.
Écrivons $X_i = Y_i \times_Y X$ et notons $h_i : X_i \to Y_i$ le changement
de base de $f$. Si chaque $h_i : X_i \to Y_i$ possède $\mathcal{P}$, alors,
pour tout $x \in X$, choisissons un $i \in I$ et un point $x_i \in X_i$
d'image $x$. Le morphisme de germes $(X_i, x_i) \to (X, x)$ est alors
\'etale, le morphisme $Y_i \to Y$ est \'etale, et $h_i$ est comme dans la
partie (3) de la définition \ref{descent-definition-local-source-target}.
Nous voyons donc que $f$ possède $\mathcal{P}$.
Réciproquement, si $f$ possède $\mathcal{P}$, alors chaque $X_i \to Y_i$
possède $\mathcal{P}$ d'après la partie (2) de la
définition \ref{descent-definition-local-source-target}.

\medskip\noindent
Preuve de (3). Supposons que $f : X \to Y$ possède $\mathcal{P}$ et que
$g : Y \to Z$ soit \'etale. Pour tout $x \in X$, on peut considérer
$(X, x) \to (X, x)$ comme un morphisme de germes \'etale, tandis que
$Y \to Z$ est un morphisme \'etale et que $h = f$ est comme dans la partie (3)
de la définition \ref{descent-definition-local-source-target}.
Nous voyons donc que $g \circ f$ possède $\mathcal{P}$.

\medskip\noindent
Preuve de (4). Soient $f : X \to Y$ un morphisme et $g : Y \to Z$ un
morphisme \'etale tels que $g \circ f$ possède $\mathcal{P}$. La partie (2) de
la définition \ref{descent-definition-local-source-target} montre alors que
$\text{pr}_Y : Y \times_Z X \to Y$ possède $\mathcal{P}$. Or le morphisme
$(f, 1) : X \to Y \times_Z X$ est \'etale, car c'est une section de la
projection \'etale $\text{pr}_X : Y \times_Z X \to X$ ; voir le
lemme \ref{morphisms-lemma-etale-permanence} de Morphismes.
Par conséquent, $f = \text{pr}_Y \circ (f, 1)$ possède $\mathcal{P}$ d'après
la partie (1) de la définition \ref{descent-definition-local-source-target}.
\end{proof}

\noindent
Le lemme suivant est l'analogue du
lemme \ref{morphisms-lemma-locally-P-characterize} de Morphismes.

\begin{lemma}
\label{descent-lemma-local-source-target-characterize}
Soit $\mathcal{P}$ une propriété de morphismes de schémas locale sur la source
et la cible pour la topologie \'etale. Soit $f : X \to Y$ un morphisme de
schémas. Les assertions suivantes sont équivalentes :
\begin{enumerate}
\item[(a)] $f$ possède la propriété $\mathcal{P}$,
\item[(b)] pour tout $x \in X$, il existe un morphisme de germes \'etale
$a : (U, u) \to (X, x)$, un morphisme \'etale $b : V \to Y$ et un morphisme
$h : U \to V$ tels que $f \circ a = b \circ h$ et que $h$ possède
$\mathcal{P}$,
\item[(c)] pour tout diagramme commutatif
$$
\xymatrix{
U \ar[d]_a \ar[r]_h & V \ar[d]^b \\
X \ar[r]^f & Y
}
$$
avec $a$, $b$ étales, le morphisme $h$ possède $\mathcal{P}$,
\item[(d)] pour un diagramme comme en (c) dans lequel $a : U \to X$ est
surjectif, le morphisme $h$ possède $\mathcal{P}$,
\item[(e)] il existe un recouvrement pour la topologie \'etale
$\{Y_i \to Y\}_{i \in I}$ tel
que chaque changement de base $Y_i \times_Y X \to Y_i$ possède
$\mathcal{P}$,
\item[(f)] il existe un recouvrement pour la topologie \'etale
$\{X_i \to X\}_{i \in I}$ tel
que chaque composé $X_i \to Y$ possède $\mathcal{P}$,
\item[(g)] il existe un recouvrement pour la topologie \'etale
$\{Y_i \to Y\}_{i \in I}$ et, pour tout $i \in I$, un recouvrement pour la
topologie \'etale
$\{X_{ij} \to Y_i \times_Y X\}_{j \in J_i}$ tels que chaque morphisme
$X_{ij} \to Y_i$ possède $\mathcal{P}$.
\end{enumerate}
\end{lemma}

\begin{proof}
L'équivalence de (a) et (b) fait partie de la
définition \ref{descent-definition-local-source-target}.
L'équivalence de (a) et (e) est la partie (2) du
lemme \ref{descent-lemma-local-source-target-implies}.
L'équivalence de (a) et (f) est la partie (1) du
lemme \ref{descent-lemma-local-source-target-implies}.
Comme (a) est désormais équivalente à (e) et à (f), il s'ensuit que (a) est
équivalente à (g).

\medskip\noindent
Il est clair que (c) implique (a). Si (a) est vérifiée, alors, pour tout
diagramme comme en (c), le morphisme $f \circ a$ possède $\mathcal{P}$ d'après
la partie (1) de la définition \ref{descent-definition-local-source-target} ; le
morphisme $h$ possède alors $\mathcal{P}$ d'après la partie (4) du
lemme \ref{descent-lemma-local-source-target-implies}.
Ainsi, (a) et (c) sont équivalentes. Il est clair que (c) implique (d).
Pour voir que (d) implique (a), supposons donné un diagramme comme en (c),
avec $a : U \to X$ surjectif et $h$ possédant $\mathcal{P}$.
Alors $b \circ h$ possède $\mathcal{P}$ d'après la partie (3) du
lemme \ref{descent-lemma-local-source-target-implies}.
Puisque $\{a : U \to X\}$ est un recouvrement pour la topologie \'etale,
nous concluons que
$f$ possède $\mathcal{P}$ d'après la partie (1) du
lemme \ref{descent-lemma-local-source-target-implies}.
\end{proof}

\noindent
Il semble que le résultat du lemme suivant ne soit pas une pure formalité :
il utilise effectivement la géométrie des morphismes \'etale.

\begin{lemma}
\label{descent-lemma-etale-local-source-target}
Soit $\mathcal{P}$ une propriété de morphismes de schémas. Supposons que
\begin{enumerate}
\item $\mathcal{P}$ soit locale sur la source pour la topologie \'etale,
\item $\mathcal{P}$ soit locale sur la cible pour la topologie \'etale, et
\item $\mathcal{P}$ soit stable par postcomposition avec les immersions
ouvertes : si $f : X \to Y$ possède $\mathcal{P}$ et si $Y \subset Z$ est un
sous-schéma ouvert, alors $X \to Z$ possède $\mathcal{P}$.
\end{enumerate}
Alors $\mathcal{P}$ est locale sur la source et la cible pour la topologie
\'etale.
\end{lemma}

\begin{proof}
Soit $\mathcal{P}$ une propriété de morphismes de schémas satisfaisant les
conditions (1), (2) et (3) du lemme. Le
lemme \ref{descent-lemma-precompose-property-local-source} montre que $\mathcal{P}$ est
stable par précomposition avec les morphismes étales. Le
lemme \ref{descent-lemma-pullback-property-local-target} montre que $\mathcal{P}$ est
stable par changement de base \'etale. Il suffit donc de prouver la partie (3)
de la définition \ref{descent-definition-local-source-target}.

\medskip\noindent
Plus précisément, supposons que $f : X \to Y$ soit un morphisme de schémas
satisfaisant la partie (3)(b) de la
définition \ref{descent-definition-local-source-target}. Autrement dit, pour tout
$x \in X$, il existe un morphisme \'etale $a_x : U_x \to X$, un point
$u_x \in U_x$ d'image $x$, un morphisme \'etale $b_x : V_x \to Y$ et un
morphisme $h_x : U_x \to V_x$ tels que $f \circ a_x = b_x \circ h_x$ et que
$h_x$ possède $\mathcal{P}$. La preuve du lemme sera achevée dès que nous
aurons montré que $f$ possède $\mathcal{P}$. Posons $U = \coprod U_x$,
$a = \coprod a_x$, $V = \coprod V_x$, $b = \coprod b_x$ et
$h = \coprod h_x$. Nous obtenons un diagramme commutatif
$$
\xymatrix{
U \ar[d]_a \ar[r]_h & V \ar[d]^b \\
X \ar[r]^f & Y
}
$$
où $a$, $b$ sont étales et où $a$ est surjectif. Remarquons que $h$
possède $\mathcal{P}$, puisque chaque $h_x$ la possède et que $\mathcal{P}$ est
locale sur la cible pour la topologie \'etale. Comme $a$ est surjectif et que
$\mathcal{P}$ est locale sur la source pour la topologie \'etale, il suffit de
prouver que $b \circ h$ possède $\mathcal{P}$. Le lemme se ramène donc à
montrer que $\mathcal{P}$ est stable par postcomposition avec un morphisme
\'etale.

\medskip\noindent
Dans la suite de la preuve, soit $f : X \to Y$ un morphisme possédant la
propriété $\mathcal{P}$, et soit $g : Y \to Z$ un morphisme \'etale.
Considérons les assertions suivantes :
\begin{enumerate}
\item[(-)] Sans hypothèse supplémentaire, $g \circ f$ possède la propriété
$\mathcal{P}$.
\item[(A)] Lorsque $Z$ est affine, $g \circ f$ possède la propriété
$\mathcal{P}$.
\item[(AA)] Lorsque $X$ et $Z$ sont affines, $g \circ f$ possède la propriété
$\mathcal{P}$.
\item[(AAA)] Lorsque $X$, $Y$ et $Z$ sont affines, $g \circ f$ possède la
propriété $\mathcal{P}$.
\end{enumerate}
Une fois (-) prouvée, la démonstration du lemme sera achevée.

\medskip\noindent
Assertion 1 : (AAA) $\Rightarrow$ (AA).
En effet, soient $f : X \to Y$, $g : Y \to Z$ comme ci-dessus, avec $X$ et
$Z$ affines. Comme $X$ est affine, donc quasi-compact, il existe un nombre fini
d'ouverts affines $Y_i \subset Y$, $i = 1, \ldots, n$, tels que
$X = \bigcup_{i = 1, \ldots, n} f^{-1}(Y_i)$. Posons
$X_i = f^{-1}(Y_i)$. D'après le
lemme \ref{descent-lemma-pullback-property-local-target}, chacun des morphismes
$X_i \to Y_i$ possède $\mathcal{P}$. Par conséquent,
$\coprod_{i = 1, \ldots, n} X_i \to \coprod_{i = 1, \ldots, n} Y_i$ possède
$\mathcal{P}$, puisque $\mathcal{P}$ est locale sur la cible pour la topologie
\'etale. En appliquant (AAA) au morphisme
$\coprod_{i = 1, \ldots, n} X_i \to \coprod_{i = 1, \ldots, n} Y_i$ et au
morphisme \'etale $\coprod_{i = 1, \ldots, n} Y_i \to Z$, nous voyons que
$\coprod_{i = 1, \ldots, n} X_i \to Z$ possède $\mathcal{P}$. Or
$\{\coprod_{i = 1, \ldots, n} X_i \to X\}$ est un recouvrement pour la
topologie \'etale ; comme $\mathcal{P}$ est locale sur la source pour cette
topologie, nous concluons que $X \to Z$ possède $\mathcal{P}$, comme voulu.

\medskip\noindent
Assertion 2 : (AAA) $\Rightarrow$ (A).
En effet, soient $f : X \to Y$, $g : Y \to Z$ comme ci-dessus, avec $Z$
affine. Choisissons un recouvrement par des ouverts affines
$X = \bigcup X_i$. Puisque $\mathcal{P}$ est locale sur la source pour la
topologie \'etale, chaque $f|_{X_i} : X_i \to Y$ possède $\mathcal{P}$. D'après
(AA), qui découle de (AAA) par l'Assertion 1, chaque $X_i \to Z$ possède
$\mathcal{P}$, pour tout $i$. Puisque $\{X_i \to X\}$ est un recouvrement
pour la topologie \'etale et que
$\mathcal{P}$ est locale sur la source pour cette topologie, nous
concluons que $X \to Z$ possède $\mathcal{P}$.

\medskip\noindent
Assertion 3 : (AAA) $\Rightarrow$ (-).
En effet, soient $f : X \to Y$, $g : Y \to Z$ comme ci-dessus. Choisissons un
recouvrement par des ouverts affines $Z = \bigcup Z_i$. Posons
$Y_i = g^{-1}(Z_i)$ et $X_i = f^{-1}(Y_i)$. D'après le
lemme \ref{descent-lemma-pullback-property-local-target}, chacun des morphismes
$X_i \to Y_i$ possède $\mathcal{P}$. D'après (A), qui découle de (AAA) par
l'Assertion 2, chaque $X_i \to Z_i$ possède $\mathcal{P}$, pour tout $i$.
Comme $\mathcal{P}$ est locale sur la cible et que
$X_i = (g \circ f)^{-1}(Z_i)$, nous concluons que
$X \to Z$ possède $\mathcal{P}$.

\medskip\noindent
Ainsi, pour prouver le lemme, il suffit d'établir (AAA).
Soient $f : X \to Y$ et $g : Y \to Z$ comme ci-dessus, avec $X, Y, Z$
affines. Remarquons qu'un morphisme \'etale entre schémas affines a des fibres
universellement bornées ; voir les lemmes
\ref{morphisms-lemma-etale-locally-quasi-finite} et
\ref{morphisms-lemma-locally-quasi-finite-qc-source-universally-bounded} de
Morphismes.
Nous pouvons donc raisonner par récurrence sur l'entier $n$ qui borne le degré
des fibres de $Y \to Z$. Voir le
lemme \ref{morphisms-lemma-etale-universally-bounded} de Morphismes pour une description de
cet entier dans le cas d'un morphisme \'etale.
Si $n = 1$, alors $Y \to Z$ est une immersion ouverte ; voir le
lemme \ref{descent-lemma-universally-injective-etale-open-immersion}. Le résultat
découle alors de l'hypothèse (3) du lemme. Supposons $n > 1$.

\medskip\noindent
Considérons le diagramme commutatif suivant
$$
\xymatrix{
X \times_Z Y \ar[d] \ar[r]_{f_Y} &
Y \times_Z Y \ar[d] \ar[r]_-{\text{pr}} &
Y \ar[d] \\
X \ar[r]^f &
Y \ar[r]^g &
Z
}
$$
Remarquons que l'on a une décomposition en sous-schémas ouverts et fermés
$Y \times_Z Y = \Delta_{Y/Z}(Y) \amalg Y'$ ; voir le
lemme \ref{morphisms-lemma-diagonal-unramified-morphism} de Morphismes.
Puisqu'elle est obtenue par changement de base, la seconde projection
$\text{pr} : Y \times_Z Y \to Y$ a des fibres de degré borné par $n$ ; voir le
lemme \ref{morphisms-lemma-base-change-universally-bounded} de Morphismes.
D'autre part, $\text{pr}|_{\Delta(Y)} : \Delta(Y) \to Y$ est un isomorphisme,
et chacune de ses fibres possède exactement un point. En appliquant le
lemme \ref{morphisms-lemma-etale-universally-bounded} de Morphismes, nous
concluons que les degrés des fibres de la restriction
$\text{pr}|_{Y'} : Y' \to Y$ sont bornés par $n - 1$.
Posons $X' = f_Y^{-1}(Y')$. Considérons le diagramme
$$
\xymatrix{
X \amalg X' \ar@{=}[d] \ar[r]_-{f \amalg f'} &
\Delta(Y) \amalg Y' \ar@{=}[d] \ar[r] &
Y \ar@{=}[d] \\
X \times_Z Y \ar[r]^{f_Y} &
Y \times_Z Y \ar[r]^-{\text{pr}} &
Y
}
$$
Comme $\mathcal{P}$ est locale sur la cible pour la topologie \'etale et donc
stable par changement de base \'etale (voir le
lemme \ref{descent-lemma-pullback-property-local-target}), nous voyons que $f_Y$
possède $\mathcal{P}$. Puisque $\mathcal{P}$ est locale sur la source pour la
topologie \'etale, $f' = f_Y|_{X'}$ possède $\mathcal{P}$. Par l'hypothèse
de récurrence, $X' \to Y$ possède $\mathcal{P}$. Comme $\mathcal{P}$ est locale
sur la source et que
$\{X \to X \times_Z Y, X' \to X \times_Y Z\}$ est un recouvrement pour la
topologie \'etale,
nous concluons que $\text{pr} \circ f_Y$ possède $\mathcal{P}$.
Remarquons que $g \circ f$ peut être considéré comme un morphisme
$g \circ f : X \to g(Y)$. Or $\text{pr} \circ f_Y$ est le changement de base
de $g \circ f : X \to g(Y)$ par le recouvrement pour la topologie \'etale
$\{Y \to g(Y)\}$ ; puisque $\mathcal{P}$ est locale sur la cible pour cette
topologie, nous concluons que $g \circ f : X \to g(Y)$ possède la
propriété $\mathcal{P}$. Enfin, en appliquant encore une fois l'hypothèse (3)
du lemme, nous concluons que $g \circ f : X \to Z$ possède la propriété
$\mathcal{P}$.
\end{proof}

\begin{remark}
\label{descent-remark-list-local-source-target}
En utilisant le lemme \ref{descent-lemma-etale-local-source-target} et les résultats
des sections précédentes de ce chapitre, on obtient aisément une liste de
types de morphismes dont la propriété est locale sur la source et la cible pour
la topologie \'etale. Dans chaque cas, nous indiquons le lemme qui établit la
localité sur la source, puis celui qui établit la localité sur la cible pour
cette topologie. La troisième hypothèse du
lemme \ref{descent-lemma-etale-local-source-target} se vérifie immédiatement dans
chaque cas ; nous omettons cette vérification. Voici la liste :
\begin{enumerate}
\item plat, voir les lemmes \ref{descent-lemma-flat-fpqc-local-source} et
\ref{descent-lemma-descending-property-flat},
\item localement de présentation finie, voir les lemmes
\ref{descent-lemma-locally-finite-presentation-fppf-local-source} et
\ref{descent-lemma-descending-property-locally-finite-presentation},
\item localement de type fini, voir les lemmes
\ref{descent-lemma-locally-finite-type-fppf-local-source} et
\ref{descent-lemma-descending-property-locally-finite-type},
\item universellement ouvert, voir les lemmes
\ref{descent-lemma-universally-open-fppf-local-source} et
\ref{descent-lemma-descending-property-universally-open},
\item syntomique, voir les lemmes \ref{descent-lemma-syntomic-syntomic-local-source} et
\ref{descent-lemma-descending-property-syntomic},
\item lisse, voir les lemmes \ref{descent-lemma-smooth-smooth-local-source} et
\ref{descent-lemma-descending-property-smooth},
\item \'etale, voir les lemmes \ref{descent-lemma-etale-etale-local-source} et
\ref{descent-lemma-descending-property-etale},
\item localement quasi-fini, voir les lemmes
\ref{descent-lemma-locally-quasi-finite-etale-local-source} et
\ref{descent-lemma-descending-property-quasi-finite},
\item non ramifié, voir les lemmes \ref{descent-lemma-unramified-etale-local-source} et
\ref{descent-lemma-descending-property-unramified},
\item G-non ramifié, voir les lemmes
\ref{descent-lemma-unramified-etale-local-source} et
\ref{descent-lemma-descending-property-unramified}, et
\item tout autre type qu'il sera utile d'ajouter.
\end{enumerate}
\end{remark}

\begin{remark}
\label{descent-remark-compare-definitions}
À ce stade, trois définitions peuvent exprimer qu'une propriété
$\mathcal{P}$ de morphismes est ``locale sur la source et la cible pour la
topologie \'etale'' :
\begin{enumerate}
\item[(ST)] $\mathcal{P}$ est locale sur la source et locale sur la cible pour
la topologie \'etale,
\item[(DM)] (définition donnée par Deligne et Mumford dans l'article
\cite[p.~100]{DM}) pour tout diagramme
$$
\xymatrix{
U \ar[d]_a \ar[r]_h & V \ar[d]^b \\
X \ar[r]^f & Y
}
$$
dont les flèches verticales sont surjectives et étales, on a
$\mathcal{P}(h) \Leftrightarrow \mathcal{P}(f)$, et
\item[(SP)] $\mathcal{P}$ est locale sur la source et la cible pour la
topologie \'etale.
\end{enumerate}
Nous avons vu dans cette section que (SP) $\Rightarrow$ (DM) $\Rightarrow$
(ST). Les exemples \ref{descent-example-silly-one} et \ref{descent-example-silly-two} montrent
qu'aucune des deux implications réciproques n'est vraie. Enfin, le
lemme \ref{descent-lemma-etale-local-source-target} montre que la différence disparaît
pour les propriétés de morphismes stables par postcomposition avec les
immersions ouvertes, ce qui sera toujours le cas en pratique.
\end{remark}

\begin{lemma}
\label{descent-lemma-etale-etale-local-source-target}
Soit $\mathcal{P}$ une propriété de morphismes de schémas locale sur la source
et la cible pour la topologie \'etale. Étant donné un diagramme commutatif
de schémas
$$
\vcenter{
\xymatrix{
X' \ar[d]_{g'} \ar[r]_{f'} & Y' \ar[d]^g \\
X \ar[r]^f & Y
}
}
\quad\text{avec les points}\quad
\vcenter{
\xymatrix{
x' \ar[d] \ar[r] & y' \ar[d] \\
x \ar[r] & y
}
}
$$
tels que $g'$ soit \'etale en $x'$ et que $g$ soit \'etale en $y'$, on a
$x \in W(f) \Leftrightarrow x' \in W(f')$, où $W(-)$ est défini dans le
lemme \ref{descent-lemma-largest-open-of-the-source}.
\end{lemma}

\begin{proof}
Le lemme \ref{descent-lemma-largest-open-of-the-source} s'applique, puisque
$\mathcal{P}$ est locale sur la source pour la topologie \'etale d'après le
lemme \ref{descent-lemma-local-source-target-implies}.

\medskip\noindent
Supposons $x \in W(f)$. Soient $U' \subset X'$ et $V' \subset Y'$ des
voisinages ouverts de $x'$ et $y'$ tels que $f'(U') \subset V'$,
$g'(U') \subset W(f)$, et tels que $g'|_{U'}$ et $g|_{V'}$ soient tous deux
étales.
Alors $f \circ g'|_{U'} = g \circ f'|_{U'}$ possède $\mathcal{P}$ d'après la
propriété (1) de la définition \ref{descent-definition-local-source-target}.
Le morphisme $f'|_{U'} : U' \to V'$ possède donc la propriété $\mathcal{P}$
d'après la partie (4) du lemme \ref{descent-lemma-local-source-target-implies}.
Puis, par la partie (3) du lemme \ref{descent-lemma-local-source-target-implies}, nous
concluons que $f'_{U'} : U' \to Y'$ possède $\mathcal{P}$.
Ainsi, $U' \subset W(f')$ par définition, et donc $x' \in W(f')$.

\medskip\noindent
Supposons $x' \in W(f')$. Soient $U' \subset X'$ et $V' \subset Y'$ des
voisinages ouverts de $x'$ et $y'$ tels que $f'(U') \subset V'$,
$U' \subset W(f')$, et tels que $g'|_{U'}$ et $g|_{V'}$ soient tous deux
étales.
Alors $U' \to Y'$ possède $\mathcal{P}$ par définition de $W(f')$.
Le morphisme $U' \to V'$ possède donc $\mathcal{P}$ d'après la partie (4) du
lemme \ref{descent-lemma-local-source-target-implies}.
Ensuite, $U' \to Y$ possède $\mathcal{P}$ d'après la partie (3) du
lemme \ref{descent-lemma-local-source-target-implies}.
Soit $U \subset X$ l'image du morphisme \'etale, donc ouvert,
$g'|_U' : U' \to X$. Alors $\{U' \to U\}$ est un recouvrement pour la
topologie \'etale, et
nous concluons que $U \to Y$ possède $\mathcal{P}$ d'après la partie (1) du
lemme \ref{descent-lemma-local-source-target-implies}.
Ainsi, $U \subset W(f)$ par définition, et donc $x \in W(f)$.
\end{proof}

\begin{lemma}
\label{descent-lemma-orbits}
Soit $k$ un corps. Soit $n \geq 2$. Pour $1 \leq i, j \leq n$ avec
$i \not = j$ et $d \geq 0$, notons $T_{i, j, d}$ l'automorphisme de
$\mathbf{A}^n_k$ donné en coordonnées par
$$
(x_1, \ldots, x_n) \longmapsto
(x_1, \ldots, x_{i - 1}, x_i + x_j^d, x_{i + 1}, \ldots, x_n)
$$
Soit $W \subset \mathbf{A}^n_k$ un sous-schéma ouvert non vide tel que
$T_{i, j, d}(W) = W$ pour tous $i, j, d$ comme ci-dessus. Alors soit
$W = \mathbf{A}^n_k$, soit la caractéristique de $k$ est $p > 0$ et
$\mathbf{A}^n_k \setminus W$ est un ensemble fini de points fermés dont les
coordonnées sont algébriques sur $\mathbf{F}_p$.
\end{lemma}

\begin{proof}
Pour établir le résultat, nous pouvons remplacer $k$ par une extension
quelconque. Soit $Z$ une composante irréductible de
$\mathbf{A}^n_k \setminus W$. Supposons $\dim(Z) \geq 1$, et cherchons une
contradiction. Il existe alors une extension de corps $k'/k$ et un point à
valeurs dans $k'$, $\xi = (\xi_1, \ldots, \xi_n) \in (k')^n$, de
$Z_{k'} \subset \mathbf{A}^n_{k'}$, tel qu'au moins l'un des
$x_1, \ldots, x_n$ soit transcendant sur le corps premier. Affirmation :
l'orbite de $\xi$ sous le groupe engendré par les transformations
$T_{i, j, d}$ est dense pour la topologie de Zariski dans
$\mathbf{A}^n_{k'}$. Cette affirmation donnera la contradiction cherchée.

\medskip\noindent
Si la caractéristique de $k'$ est nulle, les opérateurs $T_{i, j, 0}$
suffisent déjà, puisqu'ils transforment $\xi$ en les points
$$
(\xi_1 + a_1, \ldots, \xi_n + a_n)
$$
pour $(a_1, \ldots, a_n) \in \mathbf{Z}_{\geq 0}^n$ arbitraire. Si la
caractéristique est $p > 0$, nous pouvons supposer, après renumérotation, que
$\xi_n$ est transcendant sur $\mathbf{F}_p$. En appliquant successivement les
opérateurs $T_{i, n, d}$ pour $i < n$, nous voyons que l'orbite de $\xi$
contient les éléments
$$
(\xi_1 + P_1(\xi_n), \ldots, \xi_{n - 1} + P_{n - 1}(\xi_n), \xi_n)
$$
pour $(P_1, \ldots, P_{n - 1}) \in \mathbf{F}_p[t]$ arbitraire. La fermeture
de l'orbite pour la topologie de Zariski contient donc l'hyperplan d'équation
$x_n = \xi_n$. En répétant l'argument avec une autre coordonnée, nous
concluons que cette fermeture contient $x_i = \xi_i + P(\xi_n)$ pour tout
$P \in \mathbf{F}_p[t]$ tel que $\xi_i + P(\xi_n)$ soit transcendant sur
$\mathbf{F}_p$. Puisqu'il existe une infinité de tels $P$, l'affirmation en
résulte.

\medskip\noindent
Bien entendu, l'argument du paragraphe précédent s'applique aussi lorsque
$Z = \{z\}$ est de dimension $0$ et que les coordonnées de $z$ dans
$\kappa(z)$ ne sont pas algébriques sur $\mathbf{F}_p$. Le lemme en résulte.
\end{proof}

\begin{lemma}
\label{descent-lemma-etale-tau-local-source-target}
Soit $\mathcal{P}$ une propriété de morphismes de schémas. Supposons que
\begin{enumerate}
\item $\mathcal{P}$ soit locale sur la source pour la topologie \'etale,
\item $\mathcal{P}$ soit locale sur la cible pour la topologie lisse,
\item $\mathcal{P}$ soit stable par postcomposition avec les immersions
ouvertes : si $f : X \to Y$ possède $\mathcal{P}$ et si $Y \subset Z$ est un
sous-schéma ouvert, alors $X \to Z$ possède $\mathcal{P}$.
\end{enumerate}
Étant donné un diagramme commutatif de schémas
$$
\vcenter{
\xymatrix{
X' \ar[d]_{g'} \ar[r]_{f'} & Y' \ar[d]^g \\
X \ar[r]^f & Y
}
}
\quad\text{avec les points}\quad
\vcenter{
\xymatrix{
x' \ar[d] \ar[r] & y' \ar[d] \\
x \ar[r] & y
}
}
$$
tels que $g$ soit lisse en $y'$ et que $X' \to X \times_Y Y'$ soit \'etale
en $x'$, on a $x \in W(f) \Leftrightarrow x' \in W(f')$, où $W(-)$ est défini
dans le lemme \ref{descent-lemma-largest-open-of-the-source}.
\end{lemma}

\begin{proof}
Comme $\mathcal{P}$ est locale sur la source pour la topologie \'etale, nous
voyons que $x \in W(f)$ si et seulement si l'image de $x$ dans
$X \times_Y Y'$ appartient à $W(X \times_Y Y' \to Y')$. Nous pouvons donc
supposer que le diagramme du lemme est cartésien.

\medskip\noindent
Supposons $x \in W(f)$. Comme $\mathcal{P}$ est locale sur la cible pour la
topologie lisse, nous voyons que
$(g')^{-1}W(f) = W(f) \times_Y Y' \to Y'$ possède
$\mathcal{P}$. Ainsi, $(g')^{-1}W(f) \subset W(f')$. Nous concluons que
$x' \in W(f')$.

\medskip\noindent
Supposons $x' \in W(f')$. Pour tout voisinage ouvert $V' \subset Y'$ de $y'$,
on peut remplacer $Y'$ par $V'$ et $X'$ par $U' = (f')^{-1}V'$, puisque
$V' \to Y'$ est lisse et que le changement de base
$W(f') \cap U' \to V'$ de $W(f') \to Y'$ possède la propriété $\mathcal{P}$.
Nous pouvons donc supposer qu'il existe un morphisme \'etale
$Y' \to \mathbf{A}^n_Y$ au-dessus de $Y$ ; voir le
lemme \ref{morphisms-lemma-smooth-etale-over-affine-space} de Morphismes.
Considérons le diagramme
$$
\xymatrix{
X' \ar[r] \ar[d] & Y' \ar[d] \\
\mathbf{A}^n_X \ar[r]_{f_n} \ar[d] & \mathbf{A}^n_Y \ar[d] \\
X \ar[r]^f & Y
}
$$
D'après le lemme \ref{descent-lemma-etale-local-source-target} (et puisque les
recouvrements pour la topologie \'etale sont des recouvrements pour la
topologie lisse), nous voyons que $\mathcal{P}$ est locale sur la source et la
cible pour la topologie \'etale. Le
lemme \ref{descent-lemma-etale-etale-local-source-target} montre que
$W(f')$ est l'image inverse de l'ouvert $W(f_n) \subset \mathbf{A}^n_X$. En
particulier, $W(f_n)$ contient un point au-dessus de $x$. Après avoir remplacé
$X$ par l'image de $W(f_n)$, qui est ouverte, nous pouvons supposer que
$W(f_n) \to X$ est surjectif. Affirmation :
$W(f_n) = \mathbf{A}^n_X$. Cette affirmation implique que $f$ possède
$\mathcal{P}$, puisque $\mathcal{P}$ est locale sur la cible pour la topologie
lisse et que $\{\mathbf{A}^n_Y \to Y\}$ est un recouvrement pour cette
topologie.

\medskip\noindent
Essentiellement, l'affirmation résulte du fait que
$W(f_n) \subset \mathbf{A}^n_X$ est un ouvert ``invariant par translation''
qui rencontre toute fibre de $\mathbf{A}^n_X \to X$. Toutefois, pour produire
un argument de cette forme, il faut effectuer une localisation \'etale sur
$Y$ afin de disposer d'assez de translations, ce qui devient quelque peu
fastidieux. Nous utilisons plutôt les automorphismes du
lemme \ref{descent-lemma-orbits} et des morphismes étales d'espaces affines.
Nous pouvons supposer $n \geq 2$. En effet, si $n = 0$, le résultat est acquis.
Si $n = 1$, considérons le diagramme
$$
\xymatrix{
\mathbf{A}^2_X \ar[r]_{f_2} \ar[d]_p & \mathbf{A}^2_Y \ar[d] \\
\mathbf{A}^1_X \ar[r]^{f_1} & \mathbf{A}^1_Y
}
$$
Nous avons $p^{-1}(W(f_1)) \subset W(f_2)$ ; voir le premier paragraphe de la
preuve. Ainsi, $W(f_2) \to X$ est encore surjectif, et nous pouvons travailler
avec $f_2$. Supposons $n \geq 2$.

\medskip\noindent
Pour tous $1 \leq i, j \leq n$ avec $i \not = j$ et $d \geq 0$, notons
$T_{i, j, d}$ l'automorphisme de $\mathbf{A}^n$ défini dans le
lemme \ref{descent-lemma-orbits}. Nous obtenons alors un diagramme commutatif
$$
\xymatrix{
\mathbf{A}^n_X \ar[r]_{f_n} \ar[d]_{T_{i, j, d}} &
\mathbf{A}^n_Y \ar[d]^{T_{i, j, d}} \\
\mathbf{A}^n_X \ar[r]^{f_n} & \mathbf{A}^n_Y
}
$$
dont les flèches verticales sont des isomorphismes. Nous concluons que
$T_{i, j, d}(W(f_n)) = W(f_n)$. En appliquant le lemme \ref{descent-lemma-orbits}, nous
concluons que, pour tout $x \in X$, la fibre
$W(f_n)_x \subset \mathbf{A}^n_x$ est soit $\mathbf{A}^n_x$, ce qui est le
résultat cherché, soit telle que $\kappa(x)$ est de caractéristique $p > 0$ et
que $W(f_n)_x$ est le complémentaire d'un ensemble fini
$Z_x \subset \mathbf{A}^n_x$ de points fermés. La seconde possibilité ne peut
pas se produire. Considérons en effet le morphisme
$T_p : \mathbf{A}^n \to \mathbf{A}^n$ donné par
$$
(x_1, \ldots, x_n) \mapsto (x_1 - x_1^p, \ldots, x_n - x_n^p)
$$
Comme ci-dessus, nous obtenons un diagramme commutatif
$$
\xymatrix{
\mathbf{A}^n_X \ar[r]_{f_n} \ar[d]_{T_p} &
\mathbf{A}^n_Y \ar[d]^{T_p} \\
\mathbf{A}^n_X \ar[r]^{f_n} & \mathbf{A}^n_Y
}
$$
Le morphisme $T_p : \mathbf{A}^n_X \to \mathbf{A}^n_X$ est \'etale en tout
point au-dessus de $x$, et le morphisme
$T_p : \mathbf{A}^n_Y \to \mathbf{A}^n_Y$ est \'etale en tout point au-dessus
de l'image de $x$ dans $Y$. (Détails omis ; indication : calculer les
dérivées.) Nous concluons que
$$
T_p^{-1}(W) \cap \mathbf{A}^n_x = W \cap \mathbf{A}^n_x
$$
d'après le lemme \ref{descent-lemma-etale-etale-local-source-target} (nous avons déjà
vu que $\mathcal{P}$ est locale sur la source et la cible pour la topologie
\'etale). Puisque $T_p : \mathbf{A}^n_x \to \mathbf{A}^n_x$ est fini \'etale de
degré $p^n > 1$, si $Z_x$ n'est pas vide, il contient $T_p^{-1}(Z_x)$, qui est
strictement plus grand. Cette contradiction achève la preuve.
\end{proof}





\section[Morphismes de germes : localité étale]{Propriétés des morphismes de germes locales sur la source et la cible
pour la topologie étale}
\label{descent-section-local-source-target-at-point}

\noindent
Dans cette section, nous examinons l'analogue des résultats de la
section \ref{descent-section-properties-etale-local-source-target}
pour les morphismes de germes de schémas.

\begin{definition}
\label{descent-definition-local-source-target-at-point}
Soit $\mathcal{Q}$ une propriété des morphismes de germes de schémas.
On dit que $\mathcal{Q}$ est {\it locale sur la source et la cible pour la
topologie \'etale} si, pour tout diagramme commutatif
$$
\xymatrix{
(U', u') \ar[d]_a \ar[r]_{h'} & (V', v') \ar[d]^b \\
(U, u) \ar[r]^h & (V, v)
}
$$
de germes dont les flèches verticales sont étales, on a
$\mathcal{Q}(h) \Leftrightarrow \mathcal{Q}(h')$.
\end{definition}

\begin{lemma}
\label{descent-lemma-local-source-target-global-implies-local}
Soit $\mathcal{P}$ une propriété des morphismes de schémas
qui est locale sur la source et la cible pour la topologie \'etale.
Considérons la propriété $\mathcal{Q}$ des morphismes de germes définie par
$$
\mathcal{Q}((X, x) \to (S, s))
\Leftrightarrow
\text{il existe un représentant }U \to S
\text{ qui vérifie }\mathcal{P}
$$
Alors $\mathcal{Q}$ est locale sur la source et la cible pour la topologie
\'etale au sens de la
définition \ref{descent-definition-local-source-target-at-point}.
\end{lemma}

\begin{proof}
Si un morphisme de germes $(X, x) \to (S, s)$ vérifie $\mathcal{Q}$,
alors il existe des voisinages ouverts arbitrairement petits
$U \subset X$ de $x$ et $V \subset S$ de $s$ tels qu'un représentant
$U \to V$ de $(X, x) \to (S, s)$ vérifie $\mathcal{P}$.
Cela résulte du lemme \ref{descent-lemma-local-source-target-implies}. Soit
$$
\xymatrix{
(U', u') \ar[r]_{h'} \ar[d]_a & (V', v') \ar[d]^b \\
(U, u) \ar[r]^h & (V, v)
}
$$
un diagramme comme dans la
définition \ref{descent-definition-local-source-target-at-point}.
Choisissons $U_1 \subset U$ et un représentant $h_1 : U_1 \to V$ de $h$.
Choisissons $V'_1 \subset V'$ et un représentant \'etale
$b_1 : V'_1 \to V$ de $b$
(définition \ref{descent-definition-etale-morphism-germs}).
Choisissons $U'_1 \subset U'$ et des représentants
$a_1 : U'_1 \to U_1$ et $h'_1 : U'_1 \to V'_1$ de $a$ et $h'$,
avec $a_1$ \'etale. Après avoir rétréci $U'_1$, nous pouvons supposer que
$h_1 \circ a_1 = b_1 \circ h'_1$. D'après la remarque liminaire de la
démonstration, il s'agit de montrer que
$u' \in W(h'_1) \Leftrightarrow u \in W(h_1)$, où $W(-)$ est défini comme
au lemme \ref{descent-lemma-largest-open-of-the-source}.
Le résultat découle donc du
lemme \ref{descent-lemma-etale-etale-local-source-target}.
\end{proof}

\begin{lemma}
\label{descent-lemma-local-source-target-local-implies-global}
Soit $\mathcal{P}$ une propriété des morphismes de schémas qui est
locale sur la source et la cible pour la topologie \'etale. Soit $Q$ la
propriété associée des morphismes de germes, voir le
lemme \ref{descent-lemma-local-source-target-global-implies-local}.
Soit $f : X \to Y$ un morphisme de schémas. Les conditions suivantes
sont équivalentes :
\begin{enumerate}
\item $f$ vérifie la propriété $\mathcal{P}$, et
\item pour tout $x \in X$, le morphisme de germes
$(X, x) \to (Y, f(x))$ vérifie la propriété $\mathcal{Q}$.
\end{enumerate}
\end{lemma}

\begin{proof}
L'implication (1) $\Rightarrow$ (2) résulte directement des définitions.
L'implication (2) $\Rightarrow$ (1) résulte aussi de la partie (3) de la
définition \ref{descent-definition-local-source-target}.
\end{proof}

\noindent
Un morphisme de germes $(X, x) \to (S, s)$ détermine un homomorphisme
bien défini d'anneaux locaux. Le lemme suivant a donc un sens.

\begin{lemma}
\label{descent-lemma-flat-at-point}
La propriété des morphismes de germes
$$
\mathcal{P}((X, x) \to (S, s)) =
\mathcal{O}_{S, s} \to \mathcal{O}_{X, x}\text{ est plat}
$$
est locale sur la source et la cible pour la topologie \'etale.
\end{lemma}

\begin{proof}
Étant donné un diagramme comme dans la
définition \ref{descent-definition-local-source-target-at-point}, nous obtenons le
diagramme suivant d'homomorphismes locaux d'anneaux locaux :
$$
\xymatrix{
\mathcal{O}_{U', u'} & \mathcal{O}_{V', v'} \ar[l] \\
\mathcal{O}_{U, u} \ar[u] & \mathcal{O}_{V, v} \ar[l] \ar[u]
}
$$
Les flèches verticales sont des localisations d'homomorphismes d'anneaux
étales ; en particulier, elles sont essentiellement de présentation finie,
plates et non ramifiées (voir
la section \ref{algebra-section-etale} d'Algèbre).
Ces homomorphismes verticaux sont donc fidèlement plats, voir le
lemme \ref{algebra-lemma-local-flat-ff} d'Algèbre.
Si la flèche horizontale supérieure est plate, la flèche horizontale
inférieure est plate par application du
lemme \ref{algebra-lemma-flat-permanence} d'Algèbre,
avec $R = \mathcal{O}_{V, v}$, $S = \mathcal{O}_{U, u}$ et
$M = \mathcal{O}_{U', u'}$. Si la flèche horizontale inférieure est
plate, alors l'homomorphisme d'anneaux
$$
\mathcal{O}_{V', v'} \otimes_{\mathcal{O}_{V, v}} \mathcal{O}_{U, u}
\longleftarrow
\mathcal{O}_{V', v'}
$$
est plat d'après le
lemme \ref{algebra-lemma-flat-base-change} d'Algèbre.
Enfin, l'homomorphisme d'anneaux
$$
\mathcal{O}_{U', u'}
\longleftarrow
\mathcal{O}_{V', v'} \otimes_{\mathcal{O}_{V, v}} \mathcal{O}_{U, u}
$$
est la localisation d'un homomorphisme entre des extensions de type \'etale de
$\mathcal{O}_{U, u}$ ; il est donc plat d'après le
lemme \ref{algebra-lemma-map-between-etale} d'Algèbre.
\end{proof}

\begin{lemma}
\label{descent-lemma-etale-on-fiber}
Considérons un diagramme commutatif de morphismes de schémas
$$
\xymatrix{
U' \ar[r] \ar[d] & V' \ar[d] \\
U \ar[r] & V
}
$$
dont les flèches verticales sont étales, et un point $v' \in V'$
dont l'image est $v \in V$. Alors le morphisme entre les fibres
$U'_{v'} \to U_v$ est \'etale.
\end{lemma}

\begin{proof}
Le morphisme $U'_v \to U_v$ est \'etale, puisqu'il s'obtient par
changement de base à partir du morphisme \'etale $U' \to U$.
Le schéma $U'_v$ est un schéma sur $V'_v$. D'après le
lemme \ref{morphisms-lemma-etale-over-field} de Morphismes, le schéma
$V'_v$ est une réunion disjointe de spectres d'extensions finies séparables
de $\kappa(v)$. L'une de ces composantes est
$v' = \Spec(\kappa(v'))$. Par conséquent, $U'_{v'}$ est un sous-schéma
ouvert et fermé de $U'_v$, et le composé
$U'_{v'} \to U'_v \to U_v$ est \'etale (comme composé d'une immersion
ouverte et d'un morphisme \'etale, voir la
section \ref{morphisms-section-etale} de Morphismes).
\end{proof}

\noindent
Étant donné un morphisme de germes de schémas $(X, x) \to (S, s)$,
on peut définir sa {\it fibre} comme la classe d'isomorphisme du germe
$(U_s, x)$, où $U \to S$ est un représentant quelconque. Par abus de
notation, nous écrirons souvent simplement $(X_s, x)$.

\begin{lemma}
\label{descent-lemma-dimension-local-ring-fibre}
Soit $d \in \{0, 1, 2, \ldots, \infty\}$.
La propriété des morphismes de germes
$$
\mathcal{P}_d((X, x) \to (S, s)) =
\text{l'anneau local }
\mathcal{O}_{X_s, x}
\text{ de la fibre est de dimension }d
$$
est locale sur la source et la cible pour la topologie \'etale.
\end{lemma}

\begin{proof}
Étant donné un diagramme comme dans la
définition \ref{descent-definition-local-source-target-at-point}, on obtient un
morphisme \'etale entre les fibres $U'_{v'} \to U_v$ envoyant $u'$ sur $u$,
voir le lemme \ref{descent-lemma-etale-on-fiber}. Le résultat découle donc du
lemme \ref{descent-lemma-dimension-local-ring-local}.
\end{proof}

\begin{lemma}
\label{descent-lemma-transcendence-degree-at-point}
Soit $r \in \{0, 1, 2, \ldots, \infty\}$.
La propriété des morphismes de germes
$$
\mathcal{P}_r((X, x) \to (S, s))
\Leftrightarrow
\text{trdeg}_{\kappa(s)} \kappa(x) = r
$$
est locale sur la source et la cible pour la topologie \'etale.
\end{lemma}

\begin{proof}
Étant donné un diagramme comme dans la
définition \ref{descent-definition-local-source-target-at-point}, on obtient le
diagramme suivant d'homomorphismes locaux d'anneaux locaux :
$$
\xymatrix{
\mathcal{O}_{U', u'} & \mathcal{O}_{V', v'} \ar[l] \\
\mathcal{O}_{U, u} \ar[u] & \mathcal{O}_{V, v} \ar[l] \ar[u]
}
$$
Les flèches verticales sont des localisations d'homomorphismes d'anneaux
étales ; elles sont en particulier non ramifiées (voir la
section \ref{algebra-section-etale} d'Algèbre).
Par conséquent, $\kappa(u')/\kappa(u)$ et $\kappa(v')/\kappa(v)$ sont des
extensions finies séparables. On a donc
$\text{trdeg}_{\kappa(v)} \kappa(u) = \text{trdeg}_{\kappa(v')} \kappa(u)$,
ce qui démontre le lemme.
\end{proof}

\noindent
Soit $(X, x)$ un germe de schéma. La dimension de $X$ en $x$ est le minimum
des dimensions des voisinages ouverts de $x$ dans $X$, et tout voisinage
ouvert suffisamment petit a cette dimension. C'est donc un invariant de la
classe d'isomorphisme du germe. On le note simplement $\dim_x(X)$.

\begin{lemma}
\label{descent-lemma-dimension-at-point}
Soit $d \in \{0, 1, 2, \ldots, \infty\}$.
La propriété des morphismes de germes
$$
\mathcal{P}_d((X, x) \to (S, s))
\Leftrightarrow
\dim_x (X_s) = d
$$
est locale sur la source et la cible pour la topologie \'etale.
\end{lemma}

\begin{proof}
Étant donné un diagramme comme dans la
définition \ref{descent-definition-local-source-target-at-point}, on obtient un
morphisme \'etale entre les fibres $U'_{v'} \to U_v$ envoyant $u'$ sur $u$,
voir le lemme \ref{descent-lemma-etale-on-fiber}. L'égalité
$\dim_u(U_v) = \dim_{u'}(U'_{v'})$ résulte alors du
lemme \ref{descent-lemma-dimension-at-point-local}.
\end{proof}







\section{Données de descente pour les schémas au-dessus de schémas}
\label{descent-section-descent-datum}

\noindent
La plupart des raisonnements de cette section sont formels et ne reposent que
sur la définition d'une donnée de descente. Dans
la section \ref{spaces-simplicial-section-simplicial-descent} d'Espaces
simpliciaux, nous examinerons le lien
avec les schémas simpliciaux, ce qui éclairera quelque peu la situation.

\begin{definition}
\label{descent-definition-descent-datum}
Soit $f : X \to S$ un morphisme de schémas.
\begin{enumerate}
\item Soit $V \to X$ un schéma au-dessus de $X$.
Une {\it donnée de descente pour $V/X/S$} est un isomorphisme
$\varphi : V \times_S X \to X \times_S V$ de schémas au-dessus de
$X \times_S X$ satisfaisant la {\it condition de cocycle}, c'est-à-dire tel
que le diagramme
$$
\xymatrix{
V \times_S X \times_S X \ar[rd]^{\varphi_{01}} \ar[rr]_{\varphi_{02}} &
&
X \times_S X \times_S V\\
&
X \times_S V \times_S X \ar[ru]^{\varphi_{12}}
}
$$
soit commutatif (avec les notations évidentes).
\item On dit aussi que le couple $(V/X, \varphi)$ est
une {\it donnée de descente relativement à $X \to S$}.
\item Un {\it morphisme
$g : (V/X, \varphi) \to (V'/X, \varphi')$ de données de descente
relativement à $X \to S$} est un morphisme $g : V \to V'$ de schémas
au-dessus de $X$ tel que le diagramme
$$
\xymatrix{
V \times_S X \ar[r]_{\varphi} \ar[d]_{g \times \text{id}_X} &
X \times_S V \ar[d]^{\text{id}_X \times g} \\
V' \times_S X \ar[r]^{\varphi'} & X \times_S V'
}
$$
soit commutatif.
\end{enumerate}
\end{definition}

\noindent
La définition précédente donne lieu à toutes sortes d'identités
``miraculeuses''. Par exemple, l'image inverse de $\varphi$ par le
morphisme diagonal $\Delta : X \to X \times_S X$ peut être considéré comme
un morphisme $\Delta^*\varphi : V \to V$.
En effet, $X \times_{\Delta, X \times_S X} (V \times_S X) = V$ et
$X \times_{\Delta, X \times_S X} (X \times_S V) = V$.
En fait, $\Delta^*\varphi$ est l'identité.
Voilà un bon exercice si cette construction ne vous est pas familière.

\begin{remark}
\label{descent-remark-easier}
Soit $X \to S$ un morphisme de schémas. Soit $(V/X, \varphi)$ une donnée
de descente relativement à $X \to S$. On peut voir l'isomorphisme
$\varphi$ comme un isomorphisme
$$
(X \times_S X) \times_{\text{pr}_0, X} V
\longrightarrow
(X \times_S X) \times_{\text{pr}_1, X} V
$$
de schémas au-dessus de $X \times_S X$. De façon informelle, on peut donc
voir $\varphi$ comme un morphisme
$\varphi : \text{pr}_0^*V \to \text{pr}_1^*V$\footnote{Malheureusement,
le sens choisi ici pour notre flèche est opposé à la convention attendue. Dans les
définitions \ref{descent-definition-descent-datum} et
\ref{descent-definition-descent-datum-for-family-of-morphisms}, nous aurions dû
prendre le sens opposé à celui de la
définition \ref{descent-definition-descent-datum-quasi-coherent}, conformément au
principe général selon lequel ``fonctions'' et ``espaces'' sont
duaux.}. La condition de cocycle affirme alors que
$\text{pr}_{02}^*\varphi =
\text{pr}_{12}^*\varphi \circ \text{pr}_{01}^*\varphi$.
De ce point de vue, la situation est très analogue à celle d'une donnée de
descente sur des faisceaux quasi-cohérents.
\end{remark}

\noindent
Voici la définition lorsqu'on considère une famille de morphismes de même but.

\begin{definition}
\label{descent-definition-descent-datum-for-family-of-morphisms}
Soit $S$ un schéma.
Soit $\{X_i \to S\}_{i \in I}$ une famille de morphismes de but $S$.
\begin{enumerate}
\item Une {\it donnée de descente $(V_i, \varphi_{ij})$ relativement à la
famille $\{X_i \to S\}$} consiste en un schéma $V_i$ au-dessus de $X_i$
pour chaque $i \in I$, et en un isomorphisme
$\varphi_{ij} : V_i \times_S X_j \to X_i \times_S V_j$
de schémas au-dessus de $X_i \times_S X_j$ pour chaque paire
$(i, j) \in I^2$, tels que, pour tout triplet d'indices
$(i, j, k) \in I^3$, le diagramme
$$
\xymatrix{
V_i \times_S X_j \times_S X_k
\ar[rd]^{\text{pr}_{01}^*\varphi_{ij}}
\ar[rr]_{\text{pr}_{02}^*\varphi_{ik}} &
&
X_i \times_S X_j \times_S V_k\\
&
X_i \times_S V_j \times_S X_k
\ar[ru]^{\text{pr}_{12}^*\varphi_{jk}}
}
$$
de schémas au-dessus de $X_i \times_S X_j \times_S X_k$ soit commutatif
(avec les notations évidentes).
\item Un {\it morphisme
$\psi : (V_i, \varphi_{ij}) \to (V'_i, \varphi'_{ij})$
de données de descente} consiste en une famille
$\psi = (\psi_i)_{i \in I}$ de morphismes de
$X_i$-schémas $\psi_i : V_i \to V'_i$ telle que tous les diagrammes
$$
\xymatrix{
V_i \times_S X_j \ar[r]_{\varphi_{ij}} \ar[d]_{\psi_i \times \text{id}} &
X_i \times_S V_j \ar[d]^{\text{id} \times \psi_j} \\
V'_i \times_S X_j \ar[r]^{\varphi'_{ij}} & X_i \times_S V'_j
}
$$
soient commutatifs.
\end{enumerate}
\end{definition}

\noindent
Cette notion apparaît naturellement, par exemple, lorsqu'on se demande si la
catégorie fibrée des courbes relatives est un champ pour la topologie fpqc
(ce n'est pas le cas -- du moins si l'on s'en tient aux schémas).

\begin{remark}
\label{descent-remark-easier-family}
Soit $S$ un schéma.
Soit $\{X_i \to S\}_{i \in I}$ une famille de morphismes de but $S$.
Soit $(V_i, \varphi_{ij})$ une donnée de descente relativement à
$\{X_i \to S\}$. On peut voir les isomorphismes $\varphi_{ij}$ comme des
isomorphismes
$$
(X_i \times_S X_j) \times_{\text{pr}_0, X_i} V_i
\longrightarrow
(X_i \times_S X_j) \times_{\text{pr}_1, X_j} V_j
$$
de schémas au-dessus de $X_i \times_S X_j$. De façon informelle, on peut
donc voir $\varphi_{ij}$ comme un isomorphisme
$\text{pr}_0^*V_i \to \text{pr}_1^*V_j$ au-dessus de
$X_i \times_S X_j$. La condition de cocycle affirme alors que
$\text{pr}_{02}^*\varphi_{ik} =
\text{pr}_{12}^*\varphi_{jk} \circ \text{pr}_{01}^*\varphi_{ij}$.
De ce point de vue, la situation est très analogue à celle d'une donnée de
descente sur des faisceaux quasi-cohérents.
\end{remark}

\noindent
Le lemme suivant explique pourquoi nous travaillerons généralement avec une
famille réduite à un seul morphisme.

\begin{lemma}
\label{descent-lemma-family-is-one}
Soit $S$ un schéma.
Soit $\{X_i \to S\}_{i \in I}$ une famille de morphismes de but $S$.
Posons $X = \coprod_{i \in I} X_i$ et considérons-le comme un $S$-schéma.
Il existe une équivalence canonique de catégories
$$
\begin{matrix}
\text{catégorie des données de descente } \\
\text{relativement à la famille } \{X_i \to S\}_{i \in I}
\end{matrix}
\longrightarrow
\begin{matrix}
\text{ catégorie des données de descente} \\
\text{ relativement à } X/S
\end{matrix}
$$
qui envoie $(V_i, \varphi_{ij})$ sur $(V, \varphi)$, où
$V = \coprod_{i\in I} V_i$ et $\varphi = \coprod \varphi_{ij}$.
\end{lemma}

\begin{proof}
On a $X \times_S X = \coprod_{ij} X_i \times_S X_j$, et de même pour les
produits fibrés d'ordre supérieur. Se donner un morphisme $V \to X$ équivaut
exactement à se donner une famille $V_i \to X_i$. De même, se donner une
donnée de descente $\varphi$ équivaut exactement à se donner une famille
$\varphi_{ij}$.
\end{proof}

\begin{lemma}
\label{descent-lemma-pullback}
Changement de base des données de descente pour les schémas au-dessus de
schémas.
\begin{enumerate}
\item Soit
$$
\xymatrix{
X' \ar[r]_f \ar[d]_{a'} & X \ar[d]^a \\
S' \ar[r]^h & S
}
$$
un diagramme commutatif de morphismes de schémas.
La construction
$$
(V \to X, \varphi) \longmapsto f^*(V \to X, \varphi) = (V' \to X', \varphi')
$$
où $V' = X' \times_X V$ et où $\varphi'$ est défini comme le composé
$$
\xymatrix{
V' \times_{S'} X' \ar@{=}[r] &
(X' \times_X V) \times_{S'} X' \ar@{=}[r] &
(X' \times_{S'} X') \times_{X \times_S X} (V \times_S X)
\ar[d]^{\text{id} \times \varphi} \\
X' \times_{S'} V' \ar@{=}[r] &
X' \times_{S'} (X' \times_X V) &
(X' \times_{S'} X') \times_{X \times_S X} (X \times_S V) \ar@{=}[l]
}
$$
définit un foncteur de la catégorie des données de descente relativement à
$X \to S$ vers la catégorie des données de descente relativement à
$X' \to S'$.
\item Étant donnés deux morphismes $f_i : X' \to X$, $i = 0, 1$, qui rendent
le diagramme commutatif, les foncteurs $f_0^*$ et $f_1^*$ sont canoniquement
isomorphes.
\end{enumerate}
\end{lemma}

\begin{proof}
Nous omettons la preuve de (1), mais remarquons que le morphisme $\varphi'$
est le morphisme $(f \times f)^*\varphi$ dans les notations introduites à la
remarque \ref{descent-remark-easier}. Pour (2), indiquons le morphisme
$f_0^*V \to f_1^*V$ qui donne l'isomorphisme de foncteurs. Comme $f_0$ et
$f_1$ s'insèrent tous deux dans le diagramme commutatif, il existe un unique
morphisme $r : X' \to X \times_S X$ tel que
$f_i = \text{pr}_i \circ r$. On prend alors
\begin{eqnarray*}
f_0^*V & = &
X' \times_{f_0, X} V \\
& = &
X' \times_{\text{pr}_0 \circ r, X} V \\
& = &
X' \times_{r, X \times_S X} (X \times_S X) \times_{\text{pr}_0, X} V \\
& \xrightarrow{\varphi} &
X' \times_{r, X \times_S X} (X \times_S X) \times_{\text{pr}_1, X} V \\
& = &
X' \times_{\text{pr}_1 \circ r, X} V \\
& = &
X' \times_{f_1, X} V \\
& = & f_1^*V
\end{eqnarray*}
Nous omettons la vérification.
\end{proof}

\begin{definition}
\label{descent-definition-pullback-functor}
Avec $S, S', X, X', f, a, a', h$ comme dans le
lemme \ref{descent-lemma-pullback}, le foncteur
$$
(V, \varphi) \longmapsto f^*(V, \varphi)
$$
construit dans ce lemme est appelé le {\it foncteur de changement de base}
sur les données de descente.
\end{definition}

\begin{lemma}[Changement de base des données de descente pour les schémas
au-dessus de familles]
\label{descent-lemma-pullback-family}
Soient $\mathcal{U} = \{U_i \to S'\}_{i \in I}$ et
$\mathcal{V} = \{V_j \to S\}_{j \in J}$ des familles de morphismes de même
but. Soient $\alpha : I \to J$, $h : S' \to S$ et
$g_i : U_i \to V_{\alpha(i)}$ un morphisme de familles de morphismes de même
but, voir la définition \ref{sites-definition-morphism-coverings} de Sites.
\begin{enumerate}
\item Soit $(Y_j, \varphi_{jj'})$ une donnée de descente relativement à la
famille $\{V_j \to S'\}$. Le système
$$
\left(
g_i^*Y_{\alpha(i)},
(g_i \times g_{i'})^*\varphi_{\alpha(i)\alpha(i')}
\right)
$$
(avec les notations de la remarque \ref{descent-remark-easier-family}) est une donnée
de descente relativement à $\mathcal{V}$.
\item Cette construction définit un foncteur entre les données de descente
relativement à $\mathcal{U}$ et les données de descente relativement à
$\mathcal{V}$.
\item Étant donnés un second $\alpha' : I \to J$, un morphisme
$h' : S' \to S$ et un morphisme de familles de morphismes de même but
$g'_i : U_i \to V_{\alpha'(i)}$, si $h = h'$, alors les deux foncteurs ainsi
obtenus entre les données de descente sont canoniquement isomorphes.
\item Ces foncteurs coïncident, via le
lemme \ref{descent-lemma-family-is-one}, avec les foncteurs de changement de base
construits au lemme \ref{descent-lemma-pullback}.
\end{enumerate}
\end{lemma}

\begin{proof}
Cela résulte du lemme \ref{descent-lemma-pullback} par la correspondance du
lemme \ref{descent-lemma-family-is-one}.
\end{proof}

\begin{definition}
\label{descent-definition-pullback-functor-family}
Avec $\mathcal{U} = \{U_i \to S'\}_{i \in I}$,
$\mathcal{V} = \{V_j \to S\}_{j \in J}$, $\alpha : I \to J$,
$h : S' \to S$ et $g_i : U_i \to V_{\alpha(i)}$ comme dans le
lemme \ref{descent-lemma-pullback-family}, le foncteur
$$
(Y_j, \varphi_{jj'}) \longmapsto
(g_i^*Y_{\alpha(i)}, (g_i \times g_{i'})^*\varphi_{\alpha(i)\alpha(i')})
$$
construit dans ce lemme est appelé le {\it foncteur de changement de base}
sur les données de descente.
\end{definition}

\noindent
Si $\mathcal{U}$ et $\mathcal{V}$ ont le même but $S$, et si
$\mathcal{U}$ raffine $\mathcal{V}$ (voir la
définition \ref{sites-definition-morphism-coverings} de Sites), mais qu'aucun
couple explicite $(\alpha, g_i)$ n'est donné, on peut néanmoins parler du
foncteur de changement de base : le
lemme \ref{descent-lemma-pullback-family} montre que le choix de ce couple est
indifférent (à isomorphisme canonique près).


\begin{definition}
\label{descent-definition-effective}
Soit $S$ un schéma.
Soit $f : X \to S$ un morphisme de schémas.
\begin{enumerate}
\item Étant donné un schéma $U$ au-dessus de $S$, on dispose de la
{\it donnée de descente triviale} de $U$ relativement à
$\text{id} : S \to S$, à savoir le morphisme identité de $U$.
\item Le lemme \ref{descent-lemma-pullback} fournit une
{\it donnée de descente canonique} sur $X \times_S U$ relativement à
$X \to S$, obtenue par changement de base de la donnée de descente triviale
le long de $f$. Nous noterons souvent cette donnée
$(X \times_S U, can)$.
\item Une donnée de descente $(V, \varphi)$ relativement à $X/S$ est dite
{\it effective} si $(V, \varphi)$ est isomorphe à la donnée de descente
canonique $(X \times_S U, can)$ pour un certain schéma $U$ au-dessus de $S$.
\end{enumerate}
\end{definition}

\noindent
Ainsi, être effective signifie qu'il existe un schéma $U$ au-dessus de $S$
et un isomorphisme $\psi : V \to X \times_S U$ de $X$-schémas tel que
$\varphi$ soit égal au composé
$$
V \times_S X \xrightarrow{\psi \times \text{id}_X}
X \times_S U \times_S X =
X \times_S X \times_S U
\xrightarrow{\text{id}_X \times \psi^{-1}}
X \times_S V
$$

\begin{definition}
\label{descent-definition-effective-family}
Soit $S$ un schéma.
Soit $\{X_i \to S\}$ une famille de morphismes de but $S$.
\begin{enumerate}
\item Étant donné un schéma $U$ au-dessus de $S$, on dispose d'une
{\it donnée de descente canonique} sur la famille de schémas
$X_i \times_S U$, obtenue par changement de base de la donnée de descente
triviale de $U$ relativement à $\{\text{id} : S \to S\}$.
Nous notons cette donnée de descente $(X_i \times_S U, can)$.
\item Une donnée de descente $(V_i, \varphi_{ij})$ relativement à
$\{X_i \to S\}$ est dite {\it effective} s'il existe un schéma $U$ au-dessus
de $S$ tel que $(V_i, \varphi_{ij})$ soit isomorphe à
$(X_i \times_S U, can)$.
\end{enumerate}
\end{definition}












\section{Pleine fidélité des foncteurs de changement de base}
\label{descent-section-fully-faithful}

\noindent
Le foncteur de changement de base entre données de descente pour des
recouvrements fpqc est en fait pleinement fidèle. Autrement dit, les
morphismes de schémas satisfont à la descente fpqc. Le but de cette section
est de le démontrer. Le lecteur est plutôt invité à en chercher lui-même une
preuve. L'idée essentielle consiste à utiliser le
lemme \ref{descent-lemma-fpqc-universal-effective-epimorphisms}.

\begin{lemma}
\label{descent-lemma-surjective-flat-epi}
Tout morphisme surjectif et plat est un épimorphisme dans la catégorie des
schémas.
\end{lemma}

\begin{proof}
Soit $h : X' \to X$ un morphisme surjectif et plat, et soient
$a, b : X \to Y$ des morphismes tels que $a \circ h = b \circ h$.
Comme $h$ est surjectif, les morphismes $a$ et $b$ coïncident sur les
espaces topologiques sous-jacents. Choisissons $x' \in X'$ et posons
$x = h(x')$ et $y = a(x) = b(x)$. Considérons les homomorphismes d'anneaux
locaux
$$
a^\sharp_x, b^\sharp_x : \mathcal{O}_{Y, y} \to \mathcal{O}_{X, x}
$$
Ils deviennent égaux après composition avec l'homomorphisme local plat
$h^\sharp_{x'} : \mathcal{O}_{X, x} \to \mathcal{O}_{X', x'}$.
Comme tout homomorphisme local plat est fidèlement plat
(lemme \ref{algebra-lemma-local-flat-ff} d'Algèbre), on en déduit que
$h^\sharp_{x'}$ est injectif. Ainsi $a^\sharp_x = b^\sharp_x$, ce qui
entraîne $a = b$, comme voulu.
\end{proof}

\begin{lemma}
\label{descent-lemma-ff-base-change-faithful}
Soit $h : S' \to S$ un morphisme de schémas surjectif et plat.
Le foncteur de changement de base
$$
\Sch/S \longrightarrow \Sch/S', \quad
X \longmapsto S' \times_S X
$$
est fidèle.
\end{lemma}

\begin{proof}
Soient $X_1$ et $X_2$ des schémas au-dessus de $S$.
Soient $\alpha, \beta : X_2 \to X_1$ des morphismes au-dessus de $S$.
Si les changements de base de $\alpha$ et $\beta$ sont égaux, on obtient le
diagramme commutatif
$$
\xymatrix{
X_2 \ar[d]^\alpha &
S' \times_S X_2 \ar[l] \ar[d] \ar[r] &
X_2 \ar[d]^\beta \\
X_1 &
S' \times_S X_1 \ar[l] \ar[r] &
X_1
}
$$
Il suffit donc de montrer que $S' \times_S X_2 \to X_2$ est un épimorphisme.
Puisqu'il est obtenu par changement de base à partir d'un morphisme surjectif
et plat, il est surjectif et plat (voir les lemmes
\ref{morphisms-lemma-base-change-surjective} et
\ref{morphisms-lemma-base-change-flat} de Morphismes). Le résultat découle
alors du lemme \ref{descent-lemma-surjective-flat-epi}.
\end{proof}

\begin{lemma}
\label{descent-lemma-faithful}
Dans la situation du lemme \ref{descent-lemma-pullback}, supposons que
$f : X' \to X$ soit surjectif et plat. Alors le foncteur de changement de
base est fidèle.
\end{lemma}

\begin{proof}
Soient $(V_i, \varphi_i)$, $i = 1, 2$, deux données de descente pour
$X \to S$. Soient $\alpha, \beta : V_1 \to V_2$ des morphismes de données
de descente. Supposons que $f^*\alpha = f^*\beta$. Il faut montrer que
$\alpha = \beta$. Or $\alpha$ et $\beta$ sont des morphismes de schémas
au-dessus de $X$, et $f^*\alpha$ et $f^*\beta$ sont simplement les
changements de base de $\alpha$ et $\beta$, considérés comme morphismes
au-dessus de $X'$.
Le résultat découle donc du
lemme \ref{descent-lemma-ff-base-change-faithful}.
\end{proof}

\noindent
Voici le lemme essentiel de cette section.

\begin{lemma}
\label{descent-lemma-fully-faithful}
Dans la situation du lemme \ref{descent-lemma-pullback}, supposons que
\begin{enumerate}
\item $\{f : X' \to X\}$ soit un recouvrement fpqc (par exemple si $f$ est
surjectif, plat et quasi-compact), et
\item $S = S'$.
\end{enumerate}
Alors le foncteur de changement de base est pleinement fidèle.
\end{lemma}

\begin{proof}
L'hypothèse (1) implique que $f$ est surjectif et plat. Le foncteur de
changement de base est donc fidèle d'après le lemme \ref{descent-lemma-faithful}.
Soient $(V, \varphi)$ et $(W, \psi)$ deux données de descente relativement à
$X \to S$. Posons $(V', \varphi') = f^*(V, \varphi)$ et
$(W', \psi') = f^*(W, \psi)$.
Soit $\alpha' : V' \to W'$ un morphisme de données de descente pour $X'$
au-dessus de $S$. Il faut montrer qu'il existe un morphisme
$\alpha : V \to W$ de données de descente pour $X$ au-dessus de $S$ dont le
changement de base est $\alpha'$.

\medskip\noindent
Rappelons que $V'$ est le changement de base de $V$ par $f$ et que
$\varphi'$ est le changement de base de $\varphi$ par $f \times f$
(voir la remarque \ref{descent-remark-easier}).
Par hypothèse, le diagramme
$$
\xymatrix{
V' \times_S X' \ar[r]_{\varphi'} \ar[d]_{\alpha' \times \text{id}} &
X' \times_S V' \ar[d]^{\text{id} \times \alpha'} \\
W' \times_S X' \ar[r]^{\psi'} &
X' \times_S W'
}
$$
est commutatif. Nous affirmons que les deux composés
$$
\xymatrix{
V' \times_V V' \ar[r]^-{\text{pr}_i} &
V' \ar[r]^{\alpha'} &
W' \ar[r] &
W
}
, \quad i = 0, 1
$$
sont égaux. Le lecteur gagnera à le démontrer lui-même plutôt qu'à lire la
suite de ce paragraphe. (Merci de nous écrire si vous trouvez un argument
simple et élégant.)
Soient $v_0, v_1$ des points de $V'$ ayant la même image $v \in V$.
Soit $x_i \in X'$ l'image de $v_i$, et soit $x$ le point de $X$ qui est
l'image de $v$ dans $X$. Autrement dit, $v_i = (x_i, v)$ dans
$V' = X' \times_X V$. Écrivons $\varphi(v, x) = (x, v')$ pour un certain
point $v'$ de $V$. Cela est possible parce que $\varphi$ est un morphisme
au-dessus de $X \times_S X$. Posons $v_i' = (x_i, v')$, qui est un point de
$V'$. Un calcul utilisant la définition de $\varphi'$ donne alors
$\varphi'(v_i, x_j) = (x_i, v'_j)$. Posons
$w_i = \alpha'(v_i)$ et $w'_i = \alpha'(v_i')$.
On peut écrire $w_i = (x_i, u_i)$ pour un certain point $u_i$ de $W$, et
$w_i' = (x_i, u'_i)$ pour un certain point $u_i'$ de $W$.
Notre assertion équivaut à $u_0 = u_1$. Un calcul formel à partir de la
définition de $\psi'$ (voir le lemme \ref{descent-lemma-pullback}) montre que la
commutativité du diagramme précédent signifie que
$$
((x_i, x_j), \psi(u_i, x)) = ((x_i, x_j), (x, u'_j))
$$
comme points de
$(X' \times_S X') \times_{X \times_S X} (X \times_S W)$ pour tous
$i, j \in \{0, 1\}$. Il en résulte que
$\psi(u_0, x) = \psi(u_1, x)$, puis que $u_0 = u_1$ en appliquant
$\psi^{-1}$. Cela démontre l'assertion, car l'argument précédent est formel
et l'on peut prendre des points à valeurs dans un schéma (autrement dit,
prendre $(v_0, v_1) = \text{id}_{V' \times_V V'}$).

\medskip\noindent
Nous pouvons maintenant appliquer le
lemme \ref{descent-lemma-fpqc-universal-effective-epimorphisms}.
En effet, $\{V' \to V\}$ est un recouvrement fpqc, puisqu'il est obtenu par
changement de base à partir du morphisme $f : X' \to X$. Par conséquent,
d'après le lemme \ref{descent-lemma-fpqc-universal-effective-epimorphisms}, le
morphisme $\alpha' : V' \to W' \to W$ se factorise par un unique morphisme
$\alpha : V \to W$, dont le changement de base est nécessairement
$\alpha'$. Enfin, le diagramme
$$
\xymatrix{
V \times_S X \ar[r]_{\varphi} \ar[d]_{\alpha \times \text{id}} &
X \times_S V \ar[d]^{\text{id} \times \alpha} \\
W \times_S X \ar[r]^{\psi} & X \times_S W
}
$$
est commutatif, car son changement de base à $X' \times_S X'$ l'est, et le
morphisme $X' \times_S X' \to X \times_S X$ est surjectif et plat (appliquer
le lemme \ref{descent-lemma-ff-base-change-faithful}). Ainsi, $\alpha$ est bien un
morphisme de données de descente
$(V, \varphi) \to (W, \psi)$, comme voulu.
\end{proof}

\noindent
Les deux lemmes qui suivent sont devenus superflus grâce à l'amélioration de
l'exposé précédent. Ils n'en restent pas moins vrais !

\begin{lemma}
\label{descent-lemma-pullback-selfmap}
Soit $X \to S$ un morphisme de schémas.
Soit $f : X \to X$ un endomorphisme de $X$ au-dessus de $S$.
Dans ce cas, le changement de base par $f$ est isomorphe au foncteur identité
de la catégorie des données de descente relativement à $X \to S$.
\end{lemma}

\begin{proof}
Cela résulte immédiatement du lemme \ref{descent-lemma-pullback}, puisque celui-ci
donne $f^* \cong \text{id}^*$.
\end{proof}

\begin{lemma}
\label{descent-lemma-morphism-with-section-equivalence}
Soit $f : X' \to X$ un morphisme de schémas au-dessus d'un schéma de base
$S$. Supposons qu'il existe un morphisme $g : X \to X'$ au-dessus de $S$,
par exemple si $f$ admet une section. Alors le foncteur de changement de base
du lemme \ref{descent-lemma-pullback} définit une équivalence entre les catégories de
données de descente relativement à $X/S$ et à $X'/S$.
\end{lemma}

\begin{proof}
Soit $g : X \to X'$ un morphisme au-dessus de $S$. Le
lemme \ref{descent-lemma-pullback-selfmap} montre que les foncteurs
$f^* \circ g^* = (g \circ f)^*$ et $g^* \circ f^* = (f \circ g)^*$ sont
isomorphes aux foncteurs identité correspondants, comme voulu.
\end{proof}

\begin{lemma}
\label{descent-lemma-morphism-source-faithfully-flat}
Soit $f : X \to X'$ un morphisme de schémas au-dessus d'un schéma de base
$S$. Supposons que $X \to S$ soit surjectif et plat. Alors le foncteur de
changement de base du lemme \ref{descent-lemma-pullback} est un foncteur fidèle de la
catégorie des données de descente relativement à $X'/S$ vers la catégorie des
données de descente relativement à $X/S$.
\end{lemma}

\begin{proof}
On peut factoriser $X \to X'$ en
$X \to X \times_S X' \to X'$. Le premier morphisme est une section de la
première projection ; il induit donc une équivalence de catégories de données
de descente d'après le
lemme \ref{descent-lemma-morphism-with-section-equivalence}. Le second morphisme est
surjectif et plat ; il induit donc un foncteur fidèle d'après le
lemme \ref{descent-lemma-faithful}.
\end{proof}

\begin{lemma}
\label{descent-lemma-morphism-source-fpqc-covering}
Soit $f : X \to X'$ un morphisme de schémas au-dessus d'un schéma de base
$S$. Supposons que $\{X \to S\}$ soit un recouvrement fpqc (par exemple si
$f$ est surjectif, plat et quasi-compact). Alors le foncteur de changement de
base du lemme \ref{descent-lemma-pullback} est un foncteur pleinement fidèle de la
catégorie des données de descente relativement à $X'/S$ vers la catégorie des
données de descente relativement à $X/S$.
\end{lemma}

\begin{proof}
On peut factoriser $X \to X'$ en
$X \to X \times_S X' \to X'$. Le premier morphisme est une section de la
première projection ; il induit donc une équivalence de catégories de données
de descente d'après le
lemme \ref{descent-lemma-morphism-with-section-equivalence}. Le second morphisme est un
recouvrement fpqc ; il induit donc un foncteur pleinement fidèle d'après le
lemme \ref{descent-lemma-fully-faithful}.
\end{proof}

\begin{lemma}
\label{descent-lemma-fpqc-refinement-coverings-fully-faithful}
Soit $S$ un schéma. Soient
$\mathcal{U} = \{U_i \to S\}_{i \in I}$ et
$\mathcal{V} = \{V_j \to S\}_{j \in J}$ des familles de morphismes de même
but $S$. Soient $\alpha : I \to J$, $\text{id} : S \to S$ et
$g_i : U_i \to V_{\alpha(i)}$ un morphisme de familles de morphismes de même
but ; voir la définition \ref{sites-definition-morphism-coverings} de Sites.
Supposons que, pour tout $j \in J$, la famille
$\{g_i : U_i \to V_j\}_{\alpha(i) = j}$ soit un recouvrement fpqc de $V_j$.
Alors le foncteur de changement de base
$$
\text{données de descente relativement à }
\mathcal{V}
\longrightarrow
\text{données de descente relativement à }
\mathcal{U}
$$
du lemme \ref{descent-lemma-pullback-family} est pleinement fidèle.
\end{lemma}

\begin{proof}
Considérons le morphisme de schémas
$$
g :
X = \coprod\nolimits_{i \in I} U_i
\longrightarrow
Y = \coprod\nolimits_{j \in J} V_j
$$
au-dessus de $S$ qui, sur la composante d'indice $i$, prend ses valeurs dans
la composante d'indice $\alpha(i)$ au moyen du morphisme
$g_{\alpha(i)}$. Nous affirmons que $\{g : X \to Y\}$ est un recouvrement
fpqc de schémas. En effet, d'après le
lemme \ref{topologies-lemma-disjoint-union-is-fpqc-covering} de Topologies,
pour tout $j$,
le morphisme $\{\coprod_{\alpha(i) = j} U_i \to V_j\}$ est un recouvrement
fpqc. Ainsi, pour tout ouvert affine $V \subset V_j$ (que l'on peut considérer
comme un ouvert affine de $Y$), il existe un nombre fini d'ouverts affines
$W_1, \ldots, W_n \subset \coprod_{\alpha(i) = j} U_i$ (que l'on peut
considérer comme des ouverts affines de $X$) tels que
$V = \bigcup_{i = 1, \ldots, n} g(W_i)$. On dispose donc de suffisamment
d'ouverts affines de $Y$ recouverts par un nombre fini d'ouverts affines de
$X$ pour appliquer la partie (3) du
lemme \ref{topologies-lemma-recognize-fpqc-covering} de Topologies, ce qui prouve
l'assertion. Notons $DD(X/S)$, resp.\ $DD(\mathcal{U})$, la catégorie des
données de descente relativement à $X/S$, resp.\ à $\mathcal{U}$, et de même
pour $Y/S$ et $\mathcal{V}$. Considérons le diagramme
$$
\xymatrix{
DD(Y/S) \ar[r] & DD(X/S) \\
DD(\mathcal{V}) \ar[u]^{\text{lemme }\ref{descent-lemma-family-is-one}} \ar[r] &
DD(\mathcal{U}) \ar[u]_{\text{lemme }\ref{descent-lemma-family-is-one}}
}
$$
Ce diagramme est commutatif, comme le montre la preuve du
lemme \ref{descent-lemma-pullback-family}. Les flèches verticales sont des
équivalences. Le résultat découle donc du lemme \ref{descent-lemma-fully-faithful},
qui montre que la flèche horizontale supérieure du diagramme est pleinement
fidèle.
\end{proof}

\noindent
Le lemme suivant montre que, pour vérifier qu'une donnée est effective, on
peut toujours raffiner pour la topologie de Zariski la famille donnée de
morphismes de but $S$.

\begin{lemma}
\label{descent-lemma-Zariski-refinement-coverings-equivalence}
Soit $S$ un schéma. Soient
$\mathcal{U} = \{U_i \to S\}_{i \in I}$ et
$\mathcal{V} = \{V_j \to S\}_{j \in J}$ des familles de morphismes de même
but $S$. Soient $\alpha : I \to J$, $\text{id} : S \to S$ et
$g_i : U_i \to V_{\alpha(i)}$ un morphisme de familles de morphismes de même
but ; voir la définition \ref{sites-definition-morphism-coverings} de Sites.
Supposons que, pour tout $j \in J$, la famille
$\{g_i : U_i \to V_j\}_{\alpha(i) = j}$ soit un recouvrement de Zariski (voir
la définition \ref{topologies-definition-zariski-covering} de Topologies) de
$V_j$. Alors le foncteur de
changement de base
$$
\text{données de descente relativement à }
\mathcal{V}
\longrightarrow
\text{données de descente relativement à }
\mathcal{U}
$$
du lemme \ref{descent-lemma-pullback-family} est une équivalence de catégories.
En particulier, la catégorie des schémas au-dessus de $S$ est équivalente à
la catégorie des données de descente relativement à n'importe quel
recouvrement de Zariski de $S$.
\end{lemma}

\begin{proof}
Le foncteur est pleinement fidèle d'après le
lemme \ref{descent-lemma-fpqc-refinement-coverings-fully-faithful}.
Indiquons comment montrer qu'il est essentiellement surjectif.
Soit $(X_i, \varphi_{ii'})$ une donnée de descente relativement à
$\mathcal{U}$. Fixons $j \in J$ et posons
$I_j = \{i \in I \mid \alpha(i) = j\}$. Pour $i, i' \in I_j$, il existe un
morphisme canonique
$$
c_{ii'} : U_i \times_{g_i, V_j, g_{i'}} U_{i'} \to U_i \times_S U_{i'}.
$$
On peut donc effectuer le changement de base de $\varphi_{ii'}$ par ce
morphisme et poser $\psi_{ii'} = c_{ii'}^*\varphi_{ii'}$ pour
$i, i' \in I_j$. On obtient ainsi une donnée de descente
$(X_i, \psi_{ii'})$ relativement au recouvrement de Zariski
$\{g_i : U_i \to V_j\}_{i \in I_j}$. Remarquons que $\psi_{ii'}$ est un
isomorphisme de l'ouvert $X_{i, U_i \times_{V_j} U_{i'}}$ de $X_i$ sur
l'ouvert correspondant de $X_{i'}$. Il résulte de la
section \ref{schemes-section-glueing-schemes} de Schémas que l'on peut recoller
$(X_i, \psi_{ii'})$ en un schéma $Y_j$ au-dessus de $V_j$. De plus, les
morphismes $\varphi_{ii'}$, pour $i \in I_j$ et $i' \in I_{j'}$, se
recollent en un morphisme
$\varphi_{jj'} : Y_j \times_S V_{j'} \to V_j \times_S Y_{j'}$ satisfaisant
la condition de cocycle (détails omis). On obtient ainsi la donnée de descente
recherchée $(Y_j, \varphi_{jj'})$ relativement à $\mathcal{V}$.
\end{proof}

\begin{lemma}
\label{descent-lemma-refine-coverings-fully-faithful}
Soit $S$ un schéma. Soient
$\mathcal{U} = \{U_i \to S\}_{i \in I}$ et
$\mathcal{V} = \{V_j \to S\}_{j \in J}$ des recouvrements fpqc de $S$.
Si $\mathcal{U}$ est un raffinement de $\mathcal{V}$, alors le foncteur de
changement de base
$$
\text{données de descente relativement à }
\mathcal{V}
\longrightarrow
\text{données de descente relativement à }
\mathcal{U}
$$
est pleinement fidèle. En particulier, la catégorie des schémas au-dessus de
$S$ s'identifie à une sous-catégorie pleine de la catégorie des données de
descente relativement à tout recouvrement fpqc de $S$.
\end{lemma}

\begin{proof}
Considérons le recouvrement fpqc
$\mathcal{W} = \{U_i \times_S V_j \to S\}_{(i, j) \in I \times J}$ de $S$.
Il raffine à la fois $\mathcal{U}$ et $\mathcal{V}$. On a donc un diagramme
$2$-commutatif de foncteurs et de catégories
$$
\xymatrix{
DD(\mathcal{V}) \ar[rd] \ar[rr] & & DD(\mathcal{U}) \ar[ld] \\
& DD(\mathcal{W}) &
}
$$
Les notations sont celles de la preuve du
lemme \ref{descent-lemma-fpqc-refinement-coverings-fully-faithful}, et la
commutativité résulte de la partie (3) du lemme \ref{descent-lemma-pullback-family}.
Il suffit donc manifestement de montrer que les foncteurs
$DD(\mathcal{V}) \to DD(\mathcal{W})$ et
$DD(\mathcal{U}) \to DD(\mathcal{W})$ sont pleinement fidèles. Cela découle
du lemme \ref{descent-lemma-fpqc-refinement-coverings-fully-faithful}, comme voulu.
\end{proof}

\begin{remark}
\label{descent-remark-morphisms-of-schemes-satisfy-fpqc-descent}
Le lemme \ref{descent-lemma-refine-coverings-fully-faithful} affirme que les
morphismes de schémas satisfont à la descente fpqc. Autrement dit, étant
donnés un schéma $S$ et des schémas $X$, $Y$ au-dessus de $S$, le foncteur
$$
(\Sch/S)^{opp} \longrightarrow \textit{Ensembles},
\quad
T \longmapsto \Mor_T(X_T, Y_T)
$$
satisfait à la condition de faisceau pour la topologie fpqc. Le cas le plus
simple est le suivant. Supposons que $T \to S$ soit un morphisme surjectif et
plat entre schémas affines. Soit $\psi_0 : X_T \to Y_T$ un morphisme de
schémas au-dessus de $T$, compatible avec les données de descente canoniques.
Il existe alors un unique morphisme $\psi : X \to Y$ dont le changement de
base à $T$ est $\psi_0$. Ce cas particulier découle en fait directement du
lemme \ref{descent-lemma-fully-faithful}. Ce lemme est lui-même une conséquence
formelle des deux faits suivants : (a) le foncteur de changement de base par
un morphisme fidèlement plat est fidèle, voir le
lemme \ref{descent-lemma-ff-base-change-faithful}, et (b) un schéma satisfait à la
condition de faisceau pour la topologie fpqc, voir le
lemme \ref{descent-lemma-fpqc-universal-effective-epimorphisms}.
\end{remark}

\begin{lemma}
\label{descent-lemma-effective-for-fpqc-is-local-upstairs}
Soit $X \to S$ un morphisme de schémas surjectif, quasi-compact et plat.
Soit $(V, \varphi)$ une donnée de descente relativement à $X/S$. Supposons que,
pour tout $v \in V$, il existe un sous-schéma ouvert $v \in W \subset V$ tel
que $\varphi(W \times_S X) \subset X \times_S W$ et que la donnée de descente
$(W, \varphi|_{W \times_S X})$ soit effective. Alors $(V, \varphi)$ est
effective.
\end{lemma}

\begin{proof}
Soit $V = \bigcup W_i$ un recouvrement ouvert tel que
$\varphi(W_i \times_S X) \subset X \times_S W_i$ et que la donnée de descente
$(W_i, \varphi|_{W_i \times_S X})$ soit effective. Soit $U_i \to S$ un
schéma, et soit
$\alpha_i : (X \times_S U_i, can) \to (W_i, \varphi|_{W_i \times_S X})$ un
isomorphisme de données de descente. Pour chaque paire d'indices $(i, j)$,
considérons l'ouvert
$\alpha_i^{-1}(W_i \cap W_j) \subset X \times_S U_i$. Comme toutes les
données sont compatibles avec les données de descente et que
$\{X \to S\}$ est un recouvrement fpqc, le
lemme \ref{descent-lemma-open-fpqc-covering} fournit un ouvert
$U_{ij} \subset U_i$ tel que
$\alpha_i^{-1}(W_i \cap W_j) = X \times_S U_{ij}$. Le morphisme identité de
$W_i \cap W_j$ est compatible avec les données de descente ; il provient donc
d'un unique morphisme $\varphi_{ij} : U_{ij} \to U_{ji}$ au-dessus de $S$
(voir la remarque \ref{descent-remark-morphisms-of-schemes-satisfy-fpqc-descent}).
Alors $(U_i, U_{ij}, \varphi_{ij})$ est une donnée de recollement au sens de
la section \ref{schemes-section-glueing-schemes} de Schémas (preuve omise). On peut
donc supposer qu'il existe un schéma $U$ au-dessus de $S$ tel que
$U_i \subset U$ soit ouvert, que $U_{ij} = U_i \cap U_j$ et que
$\varphi_{ij} = \text{id}_{U_i \cap U_j}$ ; voir le
lemme \ref{schemes-lemma-glue} de Schémas. Après changement de base à $X$, les
$\alpha_i$ fournissent l'isomorphisme recherché
$\alpha : X \times_S U \to V$.
\end{proof}







\section{Descente de types de morphismes}
\label{descent-section-descending-types-morphisms}

\noindent
Dans la suite, nous étudions si les données de descente de schémas
relativement à un recouvrement fpqc sont effectives. La première remarque est
que ce n'est pas toujours le cas. Nous le verrons dans
l'exemple \ref{spaces-example-non-representable-descent} d'Espaces algébriques. Même les morphismes
projectifs ne satisfont pas toujours à la descente pour les recouvrements
fpqc, d'après le
lemme \ref{examples-lemma-non-effective-descent-projective} d'Exemples.

\medskip\noindent
En revanche, lorsque les schémas que l'on cherche à descendre sont
particulièrement simples, il arrive que, pour des classes entières de schémas,
les données de descente soient effectives. Nous introduisons ici une
terminologie qui décrit ce phénomène abstraitement, bien qu'elle puisse prêter
à confusion si elle est mal employée par la suite.

\begin{definition}
\label{descent-definition-descending-types-morphisms}
Soit $\mathcal{P}$ une propriété de morphismes de schémas au-dessus d'une
base. Soit
$\tau \in \{Zariski, fpqc, fppf, \etale, lisse, syntomique\}$. On dit que
{\it les morphismes de type $\mathcal{P}$ satisfont à la descente pour les
recouvrements pour la topologie $\tau$} si, pour tout recouvrement
$\mathcal{U} : \{U_i \to S\}_{i \in I}$ pour la topologie $\tau$ (voir la
section \ref{topologies-section-procedure} de Topologies), toute donnée de descente
$(X_i, \varphi_{ij})$ relativement à $\mathcal{U}$ telle que chaque morphisme
$X_i \to U_i$ vérifie la propriété $\mathcal{P}$ est effective.
\end{definition}

\noindent
Remarquons que, dans chacun des cas considérés, nous avons déjà vu que le
foncteur des schémas au-dessus de $S$ vers les données de descente relativement
à $\mathcal{U}$ est pleinement fidèle : cela résulte du
lemme \ref{descent-lemma-refine-coverings-fully-faithful} et du fait, établi dans
Topologies, que tout recouvrement $\tau$ est aussi un recouvrement fpqc. Nous
avons également vu que les données de descente sont toujours effectives pour
les recouvrements de Zariski
(lemme \ref{descent-lemma-Zariski-refinement-coverings-equivalence}). Pour éviter les
raisonnements erronés, il peut être prudent de n'étudier la notion précédente
que lorsque $\mathcal{P}$ est préservée par tout changement de base, ou du
moins locale sur la base pour la topologie $\tau$ (voir la
définition \ref{descent-definition-property-morphisms-local}). En effet, on risquerait
sinon d'invoquer $\mathcal{P}$ au cours d'une descente encore inachevée.

\medskip\noindent
Voici l'inévitable lemme qui ramène la question au cas d'un recouvrement donné
par un unique morphisme entre schémas affines.

\begin{lemma}
\label{descent-lemma-descending-types-morphisms}
Soit $\mathcal{P}$ une propriété de morphismes de schémas au-dessus d'une
base. Soit $\tau \in \{fpqc, fppf, \etale, lisse, syntomique\}$. Supposons que
\begin{enumerate}
\item $\mathcal{P}$ soit préservée par tout changement de base (voir la
définition \ref{schemes-definition-preserved-by-base-change} de Schémas),
\item si $Y_j \to V_j$, $j = 1, \ldots, m$, vérifient $\mathcal{P}$, alors
$\coprod Y_j \to \coprod V_j$ la vérifie aussi, et
\item pour tout morphisme surjectif de schémas affines $X \to S$ qui est
plat, plat de présentation finie, \'etale, lisse ou syntomique suivant que
$\tau$ vaut fpqc, fppf, \'etale, lisse ou syntomique, toute donnée de descente
$(V, \varphi)$ relativement à $X$ au-dessus de $S$ telle que $\mathcal{P}$
soit vérifiée par $V \to X$ est effective.
\end{enumerate}
Alors les morphismes de type $\mathcal{P}$ satisfont à la descente pour les
recouvrements pour la topologie $\tau$.
\end{lemma}

\begin{proof}
Soit $S$ un schéma. Soit
$\mathcal{U} = \{\varphi_i : U_i \to S\}_{i \in I}$ un recouvrement de $S$
pour la topologie $\tau$. Soit $(X_i, \varphi_{ii'})$ une donnée de descente relativement
à $\mathcal{U}$, et supposons que chaque morphisme $X_i \to U_i$ vérifie la
propriété $\mathcal{P}$. Il faut montrer qu'il existe un schéma $X \to S$ tel
que $(X_i, \varphi_{ii'}) \cong (U_i \times_S X, can)$.

\medskip\noindent
Avant d'entamer la preuve proprement dite, remarquons que, pour toute famille
de morphismes $\mathcal{V} : \{V_j \to S\}$ et tout morphisme de familles
$\mathcal{V} \to \mathcal{U}$, si l'on effectue le changement de base de la
donnée de descente $(X_i, \varphi_{ii'})$ en une donnée de descente
$(Y_j, \varphi_{jj'})$ au-dessus de $\mathcal{V}$, chacun des morphismes
$Y_j \to V_j$ vérifie encore la propriété $\mathcal{P}$. Cela résulte de
l'hypothèse (1), selon laquelle $\mathcal{P}$ est préservée par tout changement
de base, et de la définition du changement de base (voir la
définition \ref{descent-definition-pullback-functor-family}). Nous utiliserons ce fait
sans le rappeler davantage.

\medskip\noindent
Commençons par démontrer le lemme lorsque $S$ est affine. D'après les lemmes
\ref{topologies-lemma-fpqc-affine},
\ref{topologies-lemma-fppf-affine},
\ref{topologies-lemma-etale-affine},
\ref{topologies-lemma-smooth-affine} ou
\ref{topologies-lemma-syntomic-affine} de Topologies, il existe un recouvrement
standard pour la topologie $\tau$
$\mathcal{V} : \{V_j \to S\}_{j = 1, \ldots, m}$ qui raffine
$\mathcal{U}$. Le foncteur de changement de base
$DD(\mathcal{U}) \to DD(\mathcal{V})$ entre catégories de données de descente
est pleinement fidèle d'après le
lemme \ref{descent-lemma-refine-coverings-fully-faithful}. Il suffit donc de montrer
que la donnée de descente au-dessus du recouvrement standard pour la topologie $\tau$
$\mathcal{V}$ est effective. L'hypothèse (2) montre que
$\coprod Y_j \to \coprod V_j$ vérifie la propriété $\mathcal{P}$. Le
lemme \ref{descent-lemma-family-is-one} ramène alors la question au recouvrement
$\{\coprod_{j = 1, \ldots, m} V_j \to S\}$, pour lequel le résultat a été
supposé dans l'hypothèse (3) du lemme. Le lemme est donc démontré lorsque $S$
est affine.

\medskip\noindent
Supposons maintenant $S$ quelconque. Soit $V \subset S$ un ouvert affine.
Les propriétés d'un site montrent que la famille
$\mathcal{U}_V = \{V \times_S U_i \to V\}_{i \in I}$ est un recouvrement de
$V$ pour la topologie $\tau$. Notons $(X_i, \varphi_{ii'})_V$ la restriction, ou le
changement de base, de la donnée de descente donnée à $\mathcal{U}_V$. D'après
ce qui précède, on obtient un schéma $X_V$ au-dessus de $V$ dont la donnée de
descente canonique relativement à $\mathcal{U}_V$ est isomorphe à
$(X_i, \varphi_{ii'})_V$. Supposons que $V' \subset V$ soit un ouvert affine
de $V$. Alors $X_{V'}$ et $V' \times_V X_V$ ont tous deux une donnée de
descente canonique isomorphe à $(X_i, \varphi_{ii'})_{V'}$. Le
lemme \ref{descent-lemma-refine-coverings-fully-faithful} fournit donc de nouveau un
morphisme canonique $\rho^V_{V'} : X_{V'} \to X_V$ au-dessus de $S$ qui
identifie $X_{V'}$ à l'image inverse de $V'$ dans $X_V$. Nous omettons de
vérifier que, pour des ouverts affines $V'' \subset V' \subset V$ de $S$, on a
$\rho^V_{V''} = \rho^V_{V'} \circ \rho^{V'}_{V''}$.

\medskip\noindent
D'après le lemme \ref{constructions-lemma-relative-glueing} de Constructions, les
données $(X_V, \rho^V_{V'})$ se recollent en un schéma $X \to S$. De plus, on
dispose d'isomorphismes $V \times_S X \to X_V$ qui redonnent les morphismes
$\rho^V_{V'}$. En déroulant la construction des schémas $X_V$, on obtient des
isomorphismes
$$
V \times_S U_i \times_S X
\longrightarrow
V \times_S X_i
$$
compatibles avec les morphismes $\varphi_{ii'}$ et avec la restriction aux
ouverts affines plus petits de $X$. Il en résulte que la donnée de descente
canonique sur $U_i \times_S X$ est isomorphe à la donnée de descente donnée,
ce qui conclut la preuve.
\end{proof}










\section{Descente des morphismes affines}
\label{descent-section-affine}

\noindent
Dans cette section, nous montrons que
``les morphismes affines satisfont à la descente pour les recouvrements
pour la topologie fpqc''. En voici l'énoncé formel.

\begin{lemma}
\label{descent-lemma-affine}
Soit $S$ un schéma. Soit $\{X_i \to S\}_{i\in I}$ un recouvrement fpqc ; voir
la définition \ref{topologies-definition-fpqc-covering} de Topologies. Soit
$(V_i/X_i, \varphi_{ij})$ une donnée de descente relativement à
$\{X_i \to S\}$. Si chaque morphisme $V_i \to X_i$ est affine, alors la
donnée de descente est effective.
\end{lemma}

\begin{proof}
La propriété d'être affine pour un morphisme de schémas est locale sur la base
et préservée par tout changement de base ; voir les lemmes
\ref{morphisms-lemma-characterize-affine} et
\ref{morphisms-lemma-base-change-affine} de Morphismes. Le
lemme \ref{descent-lemma-descending-types-morphisms} s'applique donc, et il suffit de
démontrer l'énoncé lorsque le recouvrement fpqc est donné par un unique
morphisme plat et surjectif de schémas affines $\{X \to S\}$. Écrivons
$X = \Spec(A)$ et $S = \Spec(R)$, de sorte que $R \to A$ soit un
homomorphisme d'anneaux fidèlement plat. Soit $(V, \varphi)$ une donnée de
descente relativement à $X$ au-dessus de $S$, et supposons que $V \to X$ soit
affine. Puisque $V \to X$ est affine, on a $V = \Spec(B)$ pour une
$A$-algèbre $B$ (voir la
définition \ref{morphisms-definition-affine} de Morphismes). L'isomorphisme $\varphi$
correspond à un isomorphisme d'anneaux
$$
\varphi^\sharp :
B \otimes_R A \longleftarrow A \otimes_R B
$$
d'$A \otimes_R A$-algèbres. La condition de cocycle sur $\varphi$ signifie que
le diagramme
$$
\xymatrix{
B \otimes_R A \otimes_R A & &
A \otimes_R A \otimes_R B \ar[ll] \ar[ld]\\
& A \otimes_R B \otimes_R A \ar[lu] &
}
$$
est commutatif. En inversant ces flèches, on obtient une donnée de descente
pour les modules relativement à $R \to A$, au sens de la
définition \ref{descent-definition-descent-datum-modules}. On peut donc appliquer la
proposition \ref{descent-proposition-descent-module} pour obtenir un $R$-module
$C = \Ker(B \to A \otimes_R B)$ et un isomorphisme
$A \otimes_R C \cong B$ compatible avec les données de descente. Pour toute
paire $c, c' \in C$, le produit $cc'$ dans $B$ appartient à $C$, puisque le
morphisme $\varphi$ est un homomorphisme d'algèbres. Ainsi, $C$ est une
$R$-algèbre dont le changement de base à $A$ est isomorphe à $B$ de manière
compatible avec les données de descente. En appliquant $\Spec$, on obtient un
schéma $U$ au-dessus de $S$ tel que
$(V, \varphi) \cong (X \times_S U, can)$, comme voulu.
\end{proof}

\begin{lemma}
\label{descent-lemma-closed-immersion}
Soit $S$ un schéma. Soit $\{X_i \to S\}_{i\in I}$ un recouvrement fpqc ; voir
la définition \ref{topologies-definition-fpqc-covering} de Topologies. Soit
$(V_i/X_i, \varphi_{ij})$ une donnée de descente relativement à
$\{X_i \to S\}$. Si chaque morphisme $V_i \to X_i$ est une immersion fermée,
alors la donnée de descente est effective.
\end{lemma}

\begin{proof}
Cela résulte du fait qu'une immersion fermée est un morphisme affine
(lemme \ref{morphisms-lemma-closed-immersion-affine} de Morphismes), si bien que
le lemme \ref{descent-lemma-affine} s'applique.
\end{proof}


\section{Descente des morphismes quasi-affines}
\label{descent-section-quasi-affine}

\noindent
Dans cette section, nous montrons que
``les morphismes quasi-affines satisfont à la descente pour les recouvrements
pour la topologie fpqc''. En voici l'énoncé formel.

\begin{lemma}
\label{descent-lemma-quasi-affine}
Soit $S$ un schéma. Soit $\{X_i \to S\}_{i\in I}$ un recouvrement fpqc ; voir
la définition \ref{topologies-definition-fpqc-covering} de Topologies. Soit
$(V_i/X_i, \varphi_{ij})$ une donnée de descente relativement à
$\{X_i \to S\}$. Si chaque morphisme $V_i \to X_i$ est quasi-affine, alors la
donnée de descente est effective.
\end{lemma}

\begin{proof}
La propriété d'être quasi-affine pour un morphisme de schémas est préservée
par tout changement de base ; voir les lemmes
\ref{morphisms-lemma-characterize-quasi-affine} et
\ref{morphisms-lemma-base-change-quasi-affine} de Morphismes. Le
lemme \ref{descent-lemma-descending-types-morphisms} s'applique donc, et il suffit de
démontrer l'énoncé lorsque le recouvrement fpqc est donné par un unique
morphisme plat et surjectif de schémas affines $\{X \to S\}$. Écrivons
$X = \Spec(A)$ et $S = \Spec(R)$, de sorte que $R \to A$ soit un
homomorphisme d'anneaux fidèlement plat. Soit $(V, \varphi)$ une donnée de
descente relativement à $X$ au-dessus de $S$, et supposons que
$\pi : V \to X$ soit quasi-affine.

\medskip\noindent
D'après le
lemme \ref{morphisms-lemma-characterize-quasi-affine} de Morphismes, cela signifie que
$$
V \longrightarrow \underline{\Spec}_X(\pi_*\mathcal{O}_V) = W
$$
est une immersion ouverte quasi-compacte de schémas au-dessus de $X$. Les
projections $\text{pr}_i : X \times_S X \to X$ sont plates, et l'on a donc
$$
\text{pr}_0^*\pi_*\mathcal{O}_V =
(\pi \times \text{id}_X)_*\mathcal{O}_{V \times_S X}, \quad
\text{pr}_1^*\pi_*\mathcal{O}_V =
(\text{id}_X \times \pi)_*\mathcal{O}_{X \times_S V}
$$
par changement de base plat (lemme
\ref{coherent-lemma-flat-base-change-cohomology} de Cohomologie des schémas). Ainsi,
l'isomorphisme $\varphi : V \times_S X \to X \times_S V$ (qui est un
isomorphisme au-dessus de $X \times_S X$) induit un isomorphisme de faisceaux
quasi-cohérents d'algèbres
$$
\varphi^\sharp :
\text{pr}_0^*\pi_*\mathcal{O}_V
\longrightarrow
\text{pr}_1^*\pi_*\mathcal{O}_V
$$
sur $X \times_S X$. La condition de cocycle pour $\varphi$ implique la
condition de cocycle pour $\varphi^\sharp$. Autrement dit, on obtient une
donnée de descente $\varphi'$ sur le schéma affine $W$ relativement à $X$
au-dessus de $S$, avec la propriété supplémentaire que le morphisme $V \to W$
est un morphisme de données de descente. Le lemme \ref{descent-lemma-affine} (ou
l'effectivité de la descente pour les algèbres quasi-cohérentes) fournit donc
un schéma $U' \to S$ et un isomorphisme
$(W, \varphi') \cong (X \times_S U', can)$ de données de descente. Notons au
passage que $U'$ est affine d'après le
lemme \ref{descent-lemma-descending-property-affine}.

\medskip\noindent
On peut maintenant considérer $V$ comme l'ouvert quasi-compact
$V \subset X \times_S U'$ qui est stable par la donnée de descente
$$
can : X \times_S U' \times_S X \to X \times_S X \times_S U',
(x_0, u', x_1) \mapsto (x_0, x_1, u').
$$
Autrement dit, $(x_0, u') \in V \Rightarrow (x_1, u') \in V$ pour tous
$x_0, x_1, u'$ ayant la même image dans $S$. Comme $X \to S$ est surjectif,
on voit immédiatement que $V$ est l'image inverse d'un sous-ensemble
$U \subset U'$ par le morphisme $X \times_S U' \to U'$. Comme $X \to S$ est
quasi-compact, plat et surjectif, le morphisme $X \times_S U' \to U'$ est lui
aussi quasi-compact, plat et surjectif. D'après le
lemme \ref{morphisms-lemma-fpqc-quotient-topology} de Morphismes, le sous-ensemble
$U \subset U'$ est donc ouvert, ce qui achève la preuve.
\end{proof}












\section{Données de descente en termes de faisceaux}
\label{descent-section-descent-data-sheaves}


\noindent
Voici une autre manière d'envisager les données de descente dans le cas d'un
recouvrement sur un site.

\begin{lemma}
\label{descent-lemma-descent-data-sheaves}
Soit $\tau \in \{Zariski, fppf, \etale, lisse, syntomique\}$\footnote{L'absence
de fpqc n'est pas une coquille. Voir la discussion à la
section \ref{topologies-section-fpqc} de Topologies.}. Soit $\Sch_\tau$ un grand
site pour la topologie $\tau$. Soit $S \in \Ob(\Sch_\tau)$. Soit
$\{S_i \to S\}_{i \in I}$ un recouvrement du site $(\Sch/S)_\tau$. Il existe
une équivalence de catégories
$$
\left\{
\begin{matrix}
\text{données de descente }(X_i, \varphi_{ii'})\text{ telles que}\\
\text{chaque }X_i \in \Ob((\Sch/S)_\tau)
\end{matrix}
\right\}
\leftrightarrow
\left\{
\begin{matrix}
\text{faisceaux }F\text{ sur }(\Sch/S)_\tau\text{ tels que}\\
\text{chaque }h_{S_i} \times F\text{ est représentable}
\end{matrix}
\right\}.
$$
De plus,
\begin{enumerate}
\item les objets représentant $h_{S_i} \times F$ dans le membre de droite
correspondent aux schémas $X_i$ dans le membre de gauche, et
\item le faisceau $F$ est représentable si et seulement si la donnée de
descente correspondante $(X_i, \varphi_{ii'})$ est effective.
\end{enumerate}
\end{lemma}

\begin{proof}
Nous avons vu à la section \ref{descent-section-fpqc-universal-effective-epimorphisms}
que les préfaisceaux représentables sont des faisceaux sur le site
$(\Sch/S)_\tau$. De plus, le lemme de Yoneda (lemme
\ref{categories-lemma-yoneda} de Catégories) garantit que les morphismes entre
faisceaux représentables correspondent bijectivement aux morphismes entre les
objets qui les représentent. Nous utiliserons ces remarques sans plus les
mentionner au cours de la démonstration.

\medskip\noindent
Construisons le foncteur de droite à gauche. Soit $F$ un faisceau sur
$(\Sch/S)_\tau$ tel que chaque $h_{S_i} \times F$ soit représentable. Dans ce
cas, prenons pour $X_i$ un objet représentant dans $(\Sch/S)_\tau$. Il est muni
d'un morphisme $X_i \to S_i$. Alors $X_i \times_S S_{i'}$ et
$S_i \times_S X_{i'}$ représentent tous deux le faisceau
$h_{S_i} \times F \times h_{S_{i'}}$ ; on obtient donc un isomorphisme
$$
\varphi_{ii'} : X_i \times_S S_{i'} \to S_i \times_S X_{i'}.
$$
On vérifie immédiatement que les morphismes $\varphi_{ii'}$ sont définis
au-dessus de $S_i \times_S S_{i'}$ et satisfont à la condition de cocycle. Le
foncteur de droite à gauche est donné par la construction
$F \mapsto (X_i, \varphi_{ii'})$.

\medskip\noindent
Construisons un foncteur de gauche à droite. Pour chaque $i$, notons $F_i$ le
faisceau $h_{X_i}$. Les isomorphismes $\varphi_{ii'}$ donnent des isomorphismes
$$
\varphi_{ii'} :
F_i \times h_{S_{i'}}
\longrightarrow
h_{S_i} \times F_{i'}
$$
au-dessus de $h_{S_i} \times h_{S_{i'}}$. Posons $F$ égal au coégalisateur du
diagramme suivant
$$
\xymatrix{
\coprod_{i, i'} F_i \times h_{S_{i'}}
\ar@<1ex>[rr]^-{\text{pr}_0}
\ar@<-1ex>[rr]_-{\text{pr}_1 \circ \varphi_{ii'}}
& &
\coprod_i F_i \ar[r]
&
F
}
$$
La condition de cocycle garantit que $h_{S_i} \times F$ est isomorphe à
$F_i$, donc représentable. Le foncteur de gauche à droite est donné par la
construction $(X_i, \varphi_{ii'}) \mapsto F$.

\medskip\noindent
Nous omettons de vérifier que ces constructions donnent des foncteurs
quasi-inverses l'un de l'autre. Les assertions finales (1) et (2) découlent
des constructions.
\end{proof}

\begin{remark}
\label{descent-remark-what-product-means}
Dans l'énoncé du lemme \ref{descent-lemma-descent-data-sheaves}, la condition que
$h_{S_i} \times F$ soit représentable équivaut à ce que la restriction de $F$
à $(\Sch/S_i)_\tau$ soit représentable.
\end{remark}





















% OK, start here.
%

\chapter{Catégories dérivées des schémas}



\phantomsection
\label{perfect-section-phantom}


\section{Introduction}
\label{perfect-section-introduction}

\noindent
Dans ce chapitre, nous étudions les catégories dérivées de modules sur les schémas.
La plupart des résultats présentés ici se trouvent dans
\cite{TT}, \cite{Bokstedt-Neeman}, \cite{BvdB} et \cite{LN}.
Il existe bien entendu de nombreuses autres références.


\section{Conventions}
\label{perfect-section-conventions}

\noindent
Si $\mathcal{A}$ est une catégorie abélienne et $M$ un objet
de $\mathcal{A}$, nous notons aussi $M$ l'objet de $K(\mathcal{A})$
et/ou de $D(\mathcal{A})$ correspondant au complexe qui vaut
$M$ en degré $0$ et qui est nul en tout autre degré.

\medskip\noindent
Pour un anneau $A$, on note $K(A)$ la catégorie homotopique
des complexes de $A$-modules et $D(A)$ la catégorie dérivée associée.
De même, pour un espace annelé $(X, \mathcal{O}_X)$, le symbole
$K(\mathcal{O}_X)$ désigne la catégorie homotopique des complexes de
$\mathcal{O}_X$-modules et $D(\mathcal{O}_X)$ la catégorie dérivée
associée.










\section{Catégorie dérivée des modules quasi-cohérents}
\label{perfect-section-derived-quasi-coherent}

\noindent
Dans cette section, nous étudions les rapports entre les modules quasi-cohérents
et les modules quelconques sur un schéma $X$. On pourra consulter
\cite[appendice B]{TT}. D'après la discussion de
Schémas, section \ref{schemes-section-quasi-coherent},
le plongement
$\QCoh(\mathcal{O}_X) \subset \textit{Mod}(\mathcal{O}_X)$
fait de $\QCoh(\mathcal{O}_X)$ une sous-catégorie de Serre faible de
la catégorie des $\mathcal{O}_X$-modules. Notons
$$
D_\QCoh(\mathcal{O}_X) \subset D(\mathcal{O}_X)
$$
la sous-catégorie des complexes dont les faisceaux de cohomologie sont quasi-cohérents ; voir
Catégories dérivées, section \ref{derived-section-triangulated-sub}.
Nous obtenons ainsi un foncteur canonique
\begin{equation}
\label{perfect-equation-compare}
D(\QCoh(\mathcal{O}_X))
\longrightarrow
D_\QCoh(\mathcal{O}_X)
\end{equation}
voir Catégories dérivées, équation (\ref{derived-equation-compare}).

\begin{lemma}
\label{perfect-lemma-quasi-coherence-direct-sums}
Soit $X$ un schéma. Alors $D_\QCoh(\mathcal{O}_X)$
admet des sommes directes.
\end{lemma}

\begin{proof}
D'après Injectifs, lemme \ref{injectives-lemma-derived-products},
la catégorie dérivée $D(\mathcal{O}_X)$ admet des sommes directes,
qui se calculent en prenant les sommes directes terme à terme de représentants quelconques.
Il est donc clair que le faisceau de cohomologie d'une somme directe est la
somme directe des faisceaux de cohomologie, puisque le foncteur somme directe
est exact (dans toute catégorie abélienne de Grothendieck). Le lemme
en résulte, car une somme directe de faisceaux quasi-cohérents est quasi-cohérente ; voir
Schémas, section \ref{schemes-section-quasi-coherent}.
\end{proof}

\noindent
Nous aurons besoin de quelques résultats sur les limites dérivées. Signalons
que, dans le lemme ci-dessous, la limite dérivée ne sera généralement pas
un objet de $D_\QCoh$.

\begin{lemma}
\label{perfect-lemma-Rlim-quasi-coherent}
Soit $X$ un schéma. Soit $(K_n)$ un système projectif dans
$D_\QCoh(\mathcal{O}_X)$, de limite dérivée
$K = R\lim K_n$ dans $D(\mathcal{O}_X)$. Supposons que $H^q(K_{n + 1}) \to H^q(K_n)$
soit surjectif pour tous $q \in \mathbf{Z}$ et $n \geq 1$.
Alors
\begin{enumerate}
\item $H^q(K) = \lim H^q(K_n)$,
\item $R\lim H^q(K_n) = \lim H^q(K_n)$, et
\item pour tout ouvert affine $U \subset X$, on a
$H^p(U, \lim H^q(K_n)) = 0$ pour $p > 0$.
\end{enumerate}
\end{lemma}

\begin{proof}
Soit $\mathcal{B}$ l'ensemble des ouverts affines de $X$.
Puisque $H^q(K_n)$ est quasi-cohérent, on a $H^p(U, H^q(K_n)) = 0$
pour $U \in \mathcal{B}$ d'après Cohomologie des schémas, lemme
\ref{coherent-lemma-quasi-coherent-affine-cohomology-zero}.
De plus, les applications $H^0(U, H^q(K_{n + 1})) \to H^0(U, H^q(K_n))$
sont surjectives pour $U \in \mathcal{B}$ d'après
Schémas, lemme \ref{schemes-lemma-equivalence-quasi-coherent}.
L'assertion (1) résulte de Cohomologie, lemme
\ref{cohomology-lemma-derived-limit-suitable-system},
dont nous venons de vérifier les hypothèses.
Les assertions (2) et (3) résultent de
Cohomologie, lemme \ref{cohomology-lemma-inverse-limit-is-derived-limit}.
\end{proof}

\noindent
Le lemme suivant nous aidera à « calculer » un foncteur dérivé à droite
sur un objet de $D_\QCoh(\mathcal{O}_X)$.

\begin{lemma}
\label{perfect-lemma-nice-K-injective}
Soit $X$ un schéma. Soit $E$ un objet de
$D_\QCoh(\mathcal{O}_X)$. Alors le morphisme $E \to R\lim \tau_{\geq -n}E$ de
Catégories dérivées, remarque
\ref{derived-remark-map-into-derived-limit-truncations},
est un isomorphisme\footnote{En particulier,
$E$ admet un représentant K-injectif d'après
Catégories dérivées, lemme \ref{derived-lemma-difficulty-K-injectives}.}.
\end{lemma}

\begin{proof}
Notons $\mathcal{H}^i = H^i(E)$ le $i$-ième faisceau de cohomologie de $E$.
Soit $\mathcal{B}$ l'ensemble des ouverts affines de $X$. Alors
$H^p(U, \mathcal{H}^i) = 0$ pour tous $p > 0$, tous $i \in \mathbf{Z}$
et tous $U \in \mathcal{B}$ ; voir
Cohomologie des schémas, lemme
\ref{coherent-lemma-quasi-coherent-affine-cohomology-zero}.
Le lemme résulte donc de
Cohomologie, lemme \ref{cohomology-lemma-is-limit-dimension}.
\end{proof}

\begin{lemma}
\label{perfect-lemma-application-nice-K-injective}
Soit $X$ un schéma. Soient $F : \textit{Mod}(\mathcal{O}_X) \to \textit{Ab}$
un foncteur additif et $N \geq 0$ un entier. Supposons que
\begin{enumerate}
\item $F$ commute aux produits directs dénombrables,
\item $R^pF(\mathcal{F}) = 0$ pour tout $p \geq N$ et tout $\mathcal{F}$
quasi-cohérent.
\end{enumerate}
Alors, pour $E \in D_\QCoh(\mathcal{O}_X)$,
\begin{enumerate}
\item $H^i(RF(\tau_{\leq a}E)) \to H^i(RF(E))$ est un isomorphisme
pour $i \leq a$,
\item $H^i(RF(E)) \to H^i(RF(\tau_{\geq b - N + 1}E))$ est un isomorphisme
pour $i \geq b$,
\item si $H^i(E) = 0$ pour $i \not \in [a, b]$, où
$-\infty \leq a \leq b \leq \infty$, alors $H^i(RF(E)) = 0$
pour $i \not \in [a, b + N - 1]$.
\end{enumerate}
\end{lemma}

\begin{proof}
L'assertion (1) est
Catégories dérivées, lemme \ref{derived-lemma-negative-vanishing}.

\medskip\noindent
Démontrons l'assertion (2). Posons $E_n = \tau_{\geq -n}E$. On a
$E = R\lim E_n$ ; voir le lemme \ref{perfect-lemma-nice-K-injective}. Par suite,
$RF(E) = R\lim RF(E_n)$ dans $D(\textit{Ab})$ d'après Injectifs, lemme
\ref{injectives-lemma-RF-commutes-with-Rlim}. Ainsi, pour tout $i \in \mathbf{Z}$,
on a une suite exacte courte
$$
0 \to R^1\lim H^{i - 1}(RF(E_n)) \to H^i(RF(E)) \to \lim H^i(RF(E_n)) \to 0
$$
voir Compléments sur l'algèbre, remarque
\ref{more-algebra-remark-compare-derived-limit}.
Pour démontrer (2), nous montrerons que le terme de gauche est nul
et que le terme de droite est égal à $H^i(RF(E_{-b + N - 1}))$
pour tout $b$ tel que $i \geq b$.

\medskip\noindent
Pour tout $n$, on a un triangle distingué
$$
H^{-n}(E)[n] \to E_n \to E_{n - 1} \to H^{-n}(E)[n + 1]
$$
(Catégories dérivées, remarque
\ref{derived-remark-truncation-distinguished-triangle})
dans $D(\mathcal{O}_X)$. Puisque $H^{-n}(E)$ est quasi-cohérent, on a
$$
H^i(RF(H^{-n}(E)[n])) = R^{i + n}F(H^{-n}(E)) = 0
$$
pour $i + n \geq N$ et
$$
H^i(RF(H^{-n}(E)[n + 1])) = R^{i + n + 1}F(H^{-n}(E)) = 0
$$
pour $i + n + 1 \geq N$. On en déduit que
$$
H^i(RF(E_n)) \to H^i(RF(E_{n - 1}))
$$
est un isomorphisme pour $n \geq N - i$. Les systèmes $H^i(RF(E_n))$ vérifient donc tous
la condition ML, et le terme en $R^1\lim$ de notre suite exacte
courte est nul (voir la discussion de
Compléments sur l'algèbre, section \ref{more-algebra-section-Rlim}).
De plus, le système $H^i(RF(E_n))$ est constant à partir
de $n = N - i - 1$, comme voulu.

\medskip\noindent
Démontrons (3). Sous l'hypothèse faite sur $E$, on a
$\tau_{\leq a - 1}E = 0$, et l'annulation
de $H^i(RF(E))$ pour $i \leq a - 1$ résulte de (1).
De même, on a $\tau_{\geq b + 1}E = 0$, d'où
l'annulation de $H^i(RF(E))$ pour $i \geq b + N$ d'après
l'assertion (2).
\end{proof}

\noindent
Le lemme suivant joue un rôle essentiel dans nombre des
résultats de ce chapitre.

\begin{lemma}
\label{perfect-lemma-affine-compare-bounded}
Soit $X = \Spec(A)$ un schéma affine. Tous les foncteurs du diagramme
$$
\xymatrix{
D(\QCoh(\mathcal{O}_X)) \ar[rr]_{(\ref{perfect-equation-compare})}
& &
D_\QCoh(\mathcal{O}_X) \ar[ld]^{R\Gamma(X, -)} \\
& D(A) \ar[lu]^{\widetilde{\ \ }}
}
$$
sont des équivalences de catégories triangulées. De plus, pour $E$ dans
$D_\QCoh(\mathcal{O}_X)$, on a $H^0(X, E) = H^0(X, H^0(E))$.
\end{lemma}

\begin{proof}
Le foncteur $R\Gamma(X, -)$ donne un foncteur
$D(\mathcal{O}_X) \to D(A)$, donc, par restriction, un foncteur
\begin{equation}
\label{perfect-equation-back}
R\Gamma(X, -) : D_\QCoh(\mathcal{O}_X) \longrightarrow D(A).
\end{equation}
Nous montrerons que ce foncteur est quasi-inverse de (\ref{perfect-equation-compare})
moyennant l'équivalence entre les modules quasi-cohérents sur $X$ et
la catégorie des $A$-modules.

\medskip\noindent
Précisons ce point. Notons $(Y, \mathcal{O}_Y)$ l'espace réduit à un point, muni du faisceau
d'anneaux défini par $A$. Notons
$\pi : (X, \mathcal{O}_X) \to (Y, \mathcal{O}_Y)$
le morphisme évident d'espaces annelés.
Alors $R\Gamma(X, -)$ s'identifie à $R\pi_*$ et le foncteur
(\ref{perfect-equation-compare}), moyennant l'équivalence
$\textit{Mod}(\mathcal{O}_Y) = \text{Mod}_A = \QCoh(\mathcal{O}_X)$,
s'identifie à $L\pi^* = \pi^* = \widetilde{\ }$ (voir
Modules, lemme \ref{modules-lemma-construct-quasi-coherent-sheaves} et
Schémas, lemmes \ref{schemes-lemma-compare-constructions} et
\ref{schemes-lemma-equivalence-quasi-coherent}). Ainsi les foncteurs
$$
\xymatrix{
D(A) \ar@<1ex>[r] & D(\mathcal{O}_X) \ar@<1ex>[l]
}
$$
sont adjoints (d'après Cohomologie, lemme \ref{cohomology-lemma-adjoint}). En
particulier, on obtient des morphismes d'adjonction canoniques
$$
a : \widetilde{R\Gamma(X, E)} \longrightarrow E
$$
pour $E$ dans $D(\mathcal{O}_X)$, et
$$
b : M^\bullet \longrightarrow R\Gamma(X, \widetilde{M^\bullet})
$$
pour $M^\bullet$ un complexe de $A$-modules.

\medskip\noindent
Soit $E$ un objet de $D_\QCoh(\mathcal{O}_X)$. On peut appliquer
le lemme \ref{perfect-lemma-application-nice-K-injective}
au foncteur $F(-) = \Gamma(X, -)$
avec $N = 1$, d'après Cohomologie des schémas, lemme
\ref{coherent-lemma-quasi-coherent-affine-cohomology-zero}.
Par suite,
$$
H^0(R\Gamma(X, E)) = H^0(R\Gamma(X, \tau_{\geq 0}E)) = \Gamma(X, H^0(E))
$$
(la dernière égalité résulte de la définition de la troncature canonique).
Nous en déduirons que les morphismes d'adjonction $a$ et $b$
induisent des isomorphismes $H^0(a)$ et $H^0(b)$. Ainsi $a$ et $b$
sont des quasi-isomorphismes (puisque l'énoncé est invariant par décalage),
ce qui démontre le lemme.

\medskip\noindent
Dans les deux cas, on utilise l'exactitude du foncteur $\widetilde{\ }$
(Schémas, lemme \ref{schemes-lemma-spec-sheaves}). Elle
entraîne en effet que
$$
H^0\left(\widetilde{R\Gamma(X, E)}\right) =
\widetilde{H^0(R\Gamma(X, E))} =
\widetilde{\Gamma(X, H^0(E))}
$$
qui est égal à $H^0(E)$, car $H^0(E)$ est quasi-cohérent. Ainsi
$H^0(a)$ est un isomorphisme. Dans l'autre sens, on a
$$
H^0(R\Gamma(X, \widetilde{M^\bullet})) =
\Gamma(X, H^0(\widetilde{M^\bullet})) =
\Gamma(X, \widetilde{H^0(M^\bullet)}) =
H^0(M^\bullet)
$$
ce qui prouve que $H^0(b)$ est un isomorphisme.
\end{proof}

\begin{lemma}
\label{perfect-lemma-affine-K-flat}
Soit $X = \Spec(A)$ un schéma affine. Si $K^\bullet$ est un complexe K-plat
de $A$-modules, alors $\widetilde{K^\bullet}$ est un complexe K-plat
de $\mathcal{O}_X$-modules.
\end{lemma}

\begin{proof}
D'après Compléments sur l'algèbre, lemme \ref{more-algebra-lemma-base-change-K-flat},
$K^\bullet \otimes_A A_\mathfrak p$ est un complexe K-plat
de $A_\mathfrak p$-modules pour tout $\mathfrak p \in \Spec(A)$.
On déduit donc de
Cohomologie, lemme \ref{cohomology-lemma-check-K-flat-stalks}
(et de
Schémas, lemme \ref{schemes-lemma-spec-sheaves})
que $\widetilde{K^\bullet}$ est K-plat.
\end{proof}

\begin{lemma}
\label{perfect-lemma-quasi-coherence-pushforward}
Si $f : X \to Y$ est un morphisme de schémas affines défini par le morphisme d'anneaux
$A \to B$, le diagramme
$$
\xymatrix{
D(B) \ar[d] \ar[r] & D_\QCoh(\mathcal{O}_X) \ar[d]^{Rf_*} \\
D(A) \ar[r] & D_\QCoh(\mathcal{O}_Y)
}
$$
est commutatif.
\end{lemma}

\begin{proof}
Cela résulte du lemme \ref{perfect-lemma-affine-compare-bounded},
compte tenu de l'égalité $R\Gamma(Y, Rf_*K) = R\Gamma(X, K)$ donnée par
Cohomologie, lemme \ref{cohomology-lemma-Leray-unbounded}.
\end{proof}

\begin{lemma}
\label{perfect-lemma-quasi-coherence-pullback}
Soit $f : Y \to X$ un morphisme de schémas.
\begin{enumerate}
\item Le foncteur $Lf^*$ envoie $D_\QCoh(\mathcal{O}_X)$
dans $D_\QCoh(\mathcal{O}_Y)$.
\item Si $X$ et $Y$ sont affines et si $f$ est défini par le morphisme d'anneaux
$A \to B$, le diagramme
$$
\xymatrix{
D(B) \ar[r] & D_\QCoh(\mathcal{O}_Y) \\
D(A) \ar[r] \ar[u]^{- \otimes_A^\mathbf{L} B} &
D_\QCoh(\mathcal{O}_X) \ar[u]_{Lf^*}
}
$$
est commutatif.
\end{enumerate}
\end{lemma}

\begin{proof}
Montrons d'abord que le diagramme
$$
\xymatrix{
D(B) \ar[r] & D(\mathcal{O}_Y) \\
D(A) \ar[r] \ar[u]^{- \otimes_A^\mathbf{L} B} &
D(\mathcal{O}_X) \ar[u]_{Lf^*}
}
$$
est commutatif. Cela résulte immédiatement du lemme \ref{perfect-lemma-affine-K-flat} et
de la construction des foncteurs en question. Pour démontrer (1), soit
$E$ un objet de $D_\QCoh(\mathcal{O}_X)$. Pour vérifier que
$Lf^*E$ a des faisceaux de cohomologie quasi-cohérents, on peut travailler localement sur $X$.
Notons que $Lf^*$ est compatible à la restriction aux sous-schémas ouverts.
On peut donc supposer que $f$ est un morphisme de schémas affines comme dans (2).
Le lemme \ref{perfect-lemma-affine-compare-bounded} montre alors que
$E$ provient d'un complexe de $A$-modules. La commutativité du premier
diagramme de la démonstration donne l'énoncé analogue pour $Lf^*E$, d'où (1).
\end{proof}

\begin{lemma}
\label{perfect-lemma-quasi-coherence-tensor-product}
Soit $X$ un schéma.
\begin{enumerate}
\item Pour des objets $K, L$ de $D_\QCoh(\mathcal{O}_X)$,
le produit tensoriel dérivé $K \otimes^\mathbf{L}_{\mathcal{O}_X} L$ appartient à
$D_\QCoh(\mathcal{O}_X)$.
\item Si $X = \Spec(A)$ est affine, on a
$$
\widetilde{M^\bullet} \otimes_{\mathcal{O}_X}^\mathbf{L} \widetilde{K^\bullet}
=
\widetilde{M^\bullet \otimes_A^\mathbf{L} K^\bullet}
$$
pour tout couple de complexes de $A$-modules $K^\bullet$, $M^\bullet$.
\end{enumerate}
\end{lemma}

\begin{proof}
L'égalité (2) résulte immédiatement du lemme \ref{perfect-lemma-affine-K-flat}
et de la construction du produit tensoriel dérivé.
Pour démontrer (1), soient $K, L$ des objets de $D_\QCoh(\mathcal{O}_X)$.
Pour vérifier que $K \otimes^\mathbf{L} L$ appartient à
$D_\QCoh(\mathcal{O}_X)$, on peut travailler localement sur $X$, donc
supposer $X = \Spec(A)$ affine. D'après
le lemme \ref{perfect-lemma-affine-compare-bounded}, on peut représenter
$K$ et $L$ par des complexes de $A$-modules. L'assertion (2) donne alors
le résultat.
\end{proof}



\section{Image directe totale}
\label{perfect-section-total-direct-image}

\noindent
Le lemme suivant est l'analogue de
Cohomologie des schémas, lemme
\ref{coherent-lemma-quasi-coherence-higher-direct-images}.

\begin{lemma}
\label{perfect-lemma-quasi-coherence-direct-image}
Soit $f : X \to S$ un morphisme de schémas.
Supposons que $f$ soit quasi-séparé et quasi-compact.
\begin{enumerate}
\item Le foncteur $Rf_*$ envoie $D_\QCoh(\mathcal{O}_X)$
dans $D_\QCoh(\mathcal{O}_S)$.
\item Si $S$ est quasi-compact, il existe un entier $N = N(X, S, f)$
tel que, pour tout objet $E$ de $D_\QCoh(\mathcal{O}_X)$
vérifiant $H^m(E) = 0$ pour $m > 0$, on ait
$H^m(Rf_*E) = 0$ pour $m \geq N$.
\item Plus précisément, si $S$ est quasi-compact, on peut trouver $N = N(X, S, f)$
tel que, pour tout morphisme de schémas $S' \to S$,
la même conclusion vaille pour le foncteur $R(f')_*$,
où $f' : X' \to S'$ est obtenu à partir de $f$ par changement de base.
\end{enumerate}
\end{lemma}

\begin{proof}
Soit $E$ un objet de $D_\QCoh(\mathcal{O}_X)$. Pour démontrer (1), il faut
montrer que $Rf_*E$ a des faisceaux de cohomologie quasi-cohérents. La question est locale
sur $S$ ; on peut donc supposer $S$ quasi-compact. Choisissons $N = N(X, S, f)$ comme dans
Cohomologie des schémas, lemme
\ref{coherent-lemma-quasi-coherence-higher-direct-images}.
Alors $R^pf_*\mathcal{F} = 0$ pour tout $\mathcal{O}_X$-module quasi-cohérent
$\mathcal{F}$ et tout $p \geq N$, et cela reste vrai après changement de base.

\medskip\noindent
Supposons d'abord $E$ borné inférieurement. Montrons que (1), (2) et (3) sont vérifiées
pour un tel $E$, avec le choix de $N$ fait ci-dessus. Dans ce cas, on peut par exemple utiliser la
suite spectrale
$$
R^pf_*H^q(E) \Rightarrow R^{p + q}f_*E
$$
(Catégories dérivées, lemme \ref{derived-lemma-two-ss-complex-functor}),
la quasi-cohérence de $R^pf_*H^q(E)$ et l'annulation de $R^pf_*H^q(E)$
pour $p \geq N$, pour obtenir (1), (2) et (3).

\medskip\noindent
Démontrons ensuite (2) et (3). Supposons $H^m(E) = 0$ pour $m > 0$.
Soit $U \subset S$ un ouvert affine. D'après Cohomologie des schémas, lemme
\ref{coherent-lemma-quasi-coherence-higher-direct-images-application},
et le choix de $N$,
on a $H^p(f^{-1}(U), \mathcal{F}) = 0$ pour $p \geq N$
et tout $\mathcal{O}_X$-module quasi-cohérent $\mathcal{F}$.
On peut donc appliquer le lemme \ref{perfect-lemma-application-nice-K-injective}
au foncteur $\Gamma(f^{-1}(U), -)$ pour voir que
$$
R\Gamma(U, Rf_*E) = R\Gamma(f^{-1}(U), E)
$$
a une cohomologie nulle en degrés $\geq N$. Puisque cela vaut pour
tout ouvert affine $U \subset S$, on en déduit que $H^m(Rf_*E) = 0$
pour $m \geq N$.

\medskip\noindent
Démontrons maintenant (1) dans le cas général. Rappelons qu'il existe un
triangle distingué
$$
\tau_{\leq -n - 1}E \to E \to \tau_{\geq -n}E \to
(\tau_{\leq -n - 1}E)[1]
$$
dans $D(\mathcal{O}_X)$ ; voir Catégories dérivées, remarque
\ref{derived-remark-truncation-distinguished-triangle}.
D'après (2), les faisceaux de cohomologie de $Rf_*\tau_{\leq -n - 1}E$
sont nuls en degrés $\geq -n + N$.
Ainsi, étant donné un entier $q$, on voit que $R^qf_*E$ est égal
à $R^qf_*\tau_{\geq -n}E$ pour un certain $n$, et le résultat
ci-dessus s'applique.
\end{proof}

\begin{lemma}
\label{perfect-lemma-acyclicity-lemma}
Soit $f : X \to S$ un morphisme quasi-séparé et quasi-compact
de schémas. Soit $\mathcal{F}^\bullet$ un complexe de
$\mathcal{O}_X$-modules quasi-cohérents dont chaque terme est acyclique à droite pour $f_*$.
Alors $f_*\mathcal{F}^\bullet$ représente $Rf_*\mathcal{F}^\bullet$
dans $D(\mathcal{O}_S)$.
\end{lemma}

\begin{proof}
On dispose toujours d'un morphisme canonique
$f_*\mathcal{F}^\bullet \to Rf_*\mathcal{F}^\bullet$.
Il s'agit de montrer qu'il induit un isomorphisme sur les faisceaux de cohomologie.
L'énoncé étant invariant par décalage, il suffit de montrer que
$H^0(f_*(\mathcal{F}^\bullet)) \to R^0f_*\mathcal{F}^\bullet$
est un isomorphisme. L'énoncé est local sur $S$ ; on peut donc
supposer $S$ affine. D'après
le lemme \ref{perfect-lemma-quasi-coherence-direct-image},
on a $R^0f_*\mathcal{F}^\bullet = R^0f_*\tau_{\geq -n}\mathcal{F}^\bullet$
pour tout $n$ assez grand. On peut donc supposer $\mathcal{F}^\bullet$
borné inférieurement. Puisque chaque $\mathcal{F}^n$ est acyclique à droite pour $f_*$ par
hypothèse, le morphisme $f_*\mathcal{F}^\bullet \to Rf_*\mathcal{F}^\bullet$
est un quasi-isomorphisme d'après le lemme d'acyclicité de Leray (Catégories dérivées, lemme
\ref{derived-lemma-leray-acyclicity}).
\end{proof}

\begin{lemma}
\label{perfect-lemma-acyclicity-lemma-global}
Soit $X$ un schéma quasi-séparé et quasi-compact.
Soit $\mathcal{F}^\bullet$ un complexe de
$\mathcal{O}_X$-modules quasi-cohérents dont chaque terme est acyclique à droite pour $\Gamma(X, -)$.
Alors $\Gamma(X, \mathcal{F}^\bullet)$ représente
$R\Gamma(X, \mathcal{F}^\bullet)$ dans $D(\Gamma(X, \mathcal{O}_X)$.
\end{lemma}

\begin{proof}
Appliquer le lemme \ref{perfect-lemma-acyclicity-lemma} au morphisme canonique
$X \to \Spec(\Gamma(X, \mathcal{O}_X))$. Nous omettons quelques détails.
\end{proof}

\begin{lemma}
\label{perfect-lemma-spectral-sequence}
Soit $X$ un schéma quasi-séparé et quasi-compact. Pour tout objet
$K$ de $D_\QCoh(\mathcal{O}_X)$, la suite spectrale
$$
E_2^{i, j} = H^i(X, H^j(K)) \Rightarrow H^{i + j}(X, K)
$$
de Cohomologie, exemple \ref{cohomology-example-spectral-sequence},
est bornée et converge.
\end{lemma}

\begin{proof}
La construction de la suite spectrale donnée dans
Cohomologie, lemme \ref{cohomology-lemma-spectral-sequence-filtered-object},
à l'aide de la filtration définie par $\tau_{\leq -p}K$, montre qu'il
suffit de vérifier que, pour tout $n \in \mathbf{Z}$, on a
$$
H^n(X, \tau_{\leq -p}K) = 0 \text{ pour } p \gg 0
$$
et
$$
H^n(X, K) = H^n(X, \tau_{\leq -p}K) \text{ pour } p \ll 0
$$
La première assertion résulte du lemme \ref{perfect-lemma-application-nice-K-injective}
appliqué à $F = \Gamma(X, -)$ et de la borne donnée dans
Cohomologie des schémas, lemme
\ref{coherent-lemma-quasi-coherence-higher-direct-images}.
La seconde est vérifiée dès que
$-p \leq n$, pour tout espace annelé $(X, \mathcal{O}_X)$ et tout
$K \in D(\mathcal{O}_X)$.
\end{proof}

\begin{lemma}
\label{perfect-lemma-quasi-coherence-pushforward-direct-sums}
Soit $f : X \to S$ un morphisme quasi-séparé et quasi-compact de
schémas. Alors
$Rf_* : D_\QCoh(\mathcal{O}_X) \to D_\QCoh(\mathcal{O}_S)$
commute aux sommes directes.
\end{lemma}

\begin{proof}
Soit $E_i$ une famille d'objets de $D_\QCoh(\mathcal{O}_X)$,
et posons $E = \bigoplus E_i$. Nous voulons montrer que le morphisme
$$
\bigoplus Rf_*E_i \longrightarrow Rf_*E
$$
est un isomorphisme. Nous montrerons qu'il induit un isomorphisme sur les
faisceaux de cohomologie en degré $0$, ce qui entraînera le lemme.
Pour cela, on peut travailler localement sur $S$, ce qui permet de
supposer $S$ quasi-compact.
Choisissons un entier $N$ comme dans le lemme \ref{perfect-lemma-quasi-coherence-direct-image}.
Alors $R^0f_*E = R^0f_*\tau_{\geq -N}E$ et
$R^0f_*E_i = R^0f_*\tau_{\geq -N}E_i$ d'après le lemme cité. Notons que
$\tau_{\geq -N}E = \bigoplus \tau_{\geq -N}E_i$.
On peut donc supposer que tous les $E_i$ ont des faisceaux de cohomologie
nuls en degrés $< -N$. Utilisons ensuite les suites spectrales
$$
R^pf_*H^q(E) \Rightarrow R^{p + q}f_*E
\quad\text{et}\quad
R^pf_*H^q(E_i) \Rightarrow R^{p + q}f_*E_i
$$
(Catégories dérivées, lemme \ref{derived-lemma-two-ss-complex-functor})
pour nous ramener au cas d'une somme directe de faisceaux quasi-cohérents.
Ce cas est traité dans
Cohomologie des schémas, lemme \ref{coherent-lemma-colimit-cohomology}.
\end{proof}








\section{Morphismes affines}
\label{perfect-section-affine-morphisms}

\noindent
Dans cette section, nous réunissons quelques résultats sur l'image directe
par un morphisme affine de schémas.

\begin{lemma}
\label{perfect-lemma-pushforward-affine-morphism}
Soit $f : X \to S$ un morphisme affine de schémas. Soit $\mathcal{F}^\bullet$
un complexe de $\mathcal{O}_X$-modules quasi-cohérents. Alors
$f_*\mathcal{F}^\bullet = Rf_*\mathcal{F}^\bullet$.
\end{lemma}

\begin{proof}
Combiner le lemme \ref{perfect-lemma-acyclicity-lemma} et
Cohomologie des schémas, lemme \ref{coherent-lemma-relative-affine-vanishing}.
On peut aussi travailler sur les ouverts affines de $S$
et utiliser le lemme \ref{perfect-lemma-quasi-coherence-pushforward}.
\end{proof}

\begin{lemma}
\label{perfect-lemma-affine-morphism}
Soit $f : X \to S$ un morphisme affine de schémas.
Alors
$Rf_* : D_\QCoh(\mathcal{O}_X) \to D_\QCoh(\mathcal{O}_S)$
est conservatif.
\end{lemma}

\begin{proof}
L'énoncé signifie qu'un morphisme $\alpha : E \to F$ de
$D_\QCoh(\mathcal{O}_X)$ est un isomorphisme si
$Rf_*\alpha$ est un isomorphisme. On peut le vérifier sur les faisceaux de cohomologie.
En particulier, la question est locale sur $S$. On peut donc supposer $S$
affine, et il en est alors de même de $X$. Dans ce cas, l'énoncé résulte immédiatement de
la description des catégories dérivées
$D_\QCoh(\mathcal{O}_X)$ et
$D_\QCoh(\mathcal{O}_S)$ donnée dans
le lemme \ref{perfect-lemma-affine-compare-bounded}.
Nous omettons quelques détails.
\end{proof}

\begin{lemma}
\label{perfect-lemma-affine-morphism-pull-push}
Soit $f : X \to S$ un morphisme affine de schémas.
Pour $E$ dans $D_\QCoh(\mathcal{O}_S)$, on a
$Rf_* Lf^* E = E \otimes^\mathbf{L}_{\mathcal{O}_S} f_*\mathcal{O}_X$.
\end{lemma}

\begin{proof}
Puisque $f$ est affine, le morphisme $f_*\mathcal{O}_X \to Rf_*\mathcal{O}_X$
est un isomorphisme
(Cohomologie des schémas, lemme \ref{coherent-lemma-relative-affine-vanishing}).
On dispose d'un morphisme canonique $E \otimes^\mathbf{L} f_*\mathcal{O}_X =
E \otimes^\mathbf{L} Rf_*\mathcal{O}_X \to Rf_* Lf^* E$
adjoint au morphisme
$$
Lf^*(E \otimes^\mathbf{L} Rf_*\mathcal{O}_X) =
Lf^*E \otimes^\mathbf{L} Lf^*Rf_*\mathcal{O}_X \longrightarrow
Lf^* E \otimes^\mathbf{L} \mathcal{O}_X = Lf^* E
$$
déduit de $1 : Lf^*E \to Lf^*E$ et du morphisme canonique
$Lf^*Rf_*\mathcal{O}_X \to \mathcal{O}_X$. Pour vérifier que le morphisme ainsi construit
est un isomorphisme, on peut travailler localement sur $S$. On peut donc supposer
$S$ affine, et il en est alors de même de $X$. Dans ce cas, l'énoncé résulte immédiatement de
la description des catégories dérivées
$D_\QCoh(\mathcal{O}_X)$ et
$D_\QCoh(\mathcal{O}_S)$ et du foncteur $Lf^*$ donnée dans
les lemmes \ref{perfect-lemma-affine-compare-bounded} et
\ref{perfect-lemma-quasi-coherence-pullback}.
Nous omettons quelques détails.
\end{proof}

\noindent
Soit $Y$ un schéma. Soit $\mathcal{A}$ un faisceau de $\mathcal{O}_Y$-algèbres.
Nous noterons $D_\QCoh(\mathcal{A})$ l'image réciproque de
$D_\QCoh(\mathcal{O}_X)$ par le foncteur de restriction
$D(\mathcal{A}) \to D(\mathcal{O}_X)$. Autrement dit, $K \in D(\mathcal{A})$
appartient à $D_\QCoh(\mathcal{A})$ si et seulement si ses faisceaux de cohomologie sont
quasi-cohérents en tant que $\mathcal{O}_X$-modules. Si $\mathcal{A}$ est lui-même quasi-cohérent,
cela revient à demander que les faisceaux de cohomologie soient quasi-cohérents
en tant que $\mathcal{A}$-modules ; voir
Morphismes, lemme \ref{morphisms-lemma-affine-equivalence-modules}.

\begin{lemma}
\label{perfect-lemma-affine-morphism-equivalence}
Soit $f : X \to Y$ un morphisme affine de schémas. Alors $f_*$ induit
une équivalence
$$
\Phi : D_\QCoh(\mathcal{O}_X) \longrightarrow D_\QCoh(f_*\mathcal{O}_X)
$$
dont le composé avec $D_\QCoh(f_*\mathcal{O}_X) \to D_\QCoh(\mathcal{O}_Y)$
est $Rf_* : D_\QCoh(\mathcal{O}_X) \to D_\QCoh(\mathcal{O}_Y)$.
\end{lemma}

\begin{proof}
Rappelons que, pour calculer $Rf_*$ sur un objet $K \in D_\QCoh(\mathcal{O}_X)$,
on choisit un complexe K-injectif $\mathcal{I}^\bullet$ de
$\mathcal{O}_X$-modules représentant $K$, et l'on prend $f_*\mathcal{I}^\bullet$.
On définit donc $\Phi(K)$ comme le complexe $f_*\mathcal{I}^\bullet$
considéré comme un complexe de $f_*\mathcal{O}_X$-modules.
Notons $g : (X, \mathcal{O}_X) \to (Y, f_*\mathcal{O}_X)$ le
morphisme évident d'espaces annelés. Alors $g$ est un morphisme plat
d'espaces annelés (voir ci-dessous la description des fibres), et
$\Phi$ est la restriction de $Rg_*$ à $D_\QCoh(\mathcal{O}_X)$.
Nous affirmons que $Lg^*$ en est un quasi-inverse. Observons d'abord que
$Lg^*$ envoie $D_\QCoh(f_*\mathcal{O}_X)$ dans $D_\QCoh(\mathcal{O}_X)$,
car $g^*$ transforme les modules quasi-cohérents en modules quasi-cohérents
(Modules, lemme \ref{modules-lemma-pullback-quasi-coherent}).
Pour achever la démonstration, il suffit de montrer que
les morphismes d'adjonction
$$
Lg^*\Phi(K) = Lg^*Rg_*K \to K
\quad\text{et}\quad
M \to Rg_*Lg^*M = \Phi(Lg^*M)
$$
sont des isomorphismes pour $K \in D_\QCoh(\mathcal{O}_X)$ et
$M \in D_\QCoh(f_*\mathcal{O}_X)$. La question étant locale,
on peut supposer $Y$ affine, et donc aussi $X$.

\medskip\noindent
Supposons $Y = \Spec(B)$ et $X = \Spec(A)$. Soit
$\mathfrak p = x \in \Spec(A) = X$ un point dont l'image est
$\mathfrak q = y \in \Spec(B) = Y$. Alors
$(f_*\mathcal{O}_X)_y = A_\mathfrak q$ et $\mathcal{O}_{X, x} = A_\mathfrak p$ ;
le morphisme $g$ est donc plat. Par suite, $g^*$ est exact et $H^i(Lg^*M) = g^*H^i(M)$
pour tout $M$ dans $D(f_*\mathcal{O}_X)$.
Pour $K \in D_\QCoh(\mathcal{O}_X)$, on a
$$
H^i(\Phi(K)) = H^i(Rf_*K) = f_*H^i(K)
$$
d'après l'annulation des images directes supérieures
(Cohomologie des schémas, lemme \ref{coherent-lemma-relative-affine-vanishing})
et le lemme \ref{perfect-lemma-application-nice-K-injective} (un petit détail est omis).
Il suffit donc de montrer que
$$
g^*g_*\mathcal{F} \to \mathcal{F}
\quad\text{et}\quad
\mathcal{G} \to g_*g^*\mathcal{G}
$$
sont des isomorphismes, où $\mathcal{F}$ est
un $\mathcal{O}_X$-module quasi-cohérent et $\mathcal{G}$
un $f_*\mathcal{O}_X$-module quasi-cohérent. Cela résulte de
Morphismes, lemme \ref{morphisms-lemma-affine-equivalence-modules}.
\end{proof}




\section{Cohomologie à support dans un sous-ensemble fermé}
\label{perfect-section-cohomology-support}

\noindent
Nous développons, pour les schémas et les modules quasi-cohérents,
les résultats de Cohomologie, sections
\ref{cohomology-section-cohomology-support} et
\ref{cohomology-section-cohomology-support-bis}.

\begin{definition}
\label{perfect-definition-supported-on}
Soit $X$ un schéma. Soit $E$ un objet de $D(\mathcal{O}_X)$.
Soit $T \subset X$ une partie fermée.
On dit que $E$ est {\it à support dans $T$} si les
faisceaux de cohomologie $H^i(E)$ sont à support dans $T$.
\end{definition}

\noindent
Reprenons une partie de la discussion de
Cohomologie, section \ref{cohomology-section-cohomology-support-bis},
dans la situation de la définition.
Soit $X$ un schéma. Soit $T \subset X$ une partie fermée.
La catégorie des $\mathcal{O}_X$-modules dont le
support est contenu dans $T$ est une sous-catégorie de Serre de la
catégorie de tous les $\mathcal{O}_X$-modules ; voir
Homologie, définition \ref{homology-definition-serre-subcategory}
et
Modules, lemme \ref{modules-lemma-support-section-closed}.
On peut donc appliquer la discussion de
Catégories dérivées, section \ref{derived-section-triangulated-sub},
et l'on obtient une sous-catégorie triangulée strictement pleine et saturée de
$D(\mathcal{O}_X)$, formée des objets à support dans $T$ ; nous
noterons cette catégorie $D_T(\mathcal{O}_X)$ dans la suite.

\medskip\noindent
Dans la situation de la définition \ref{perfect-definition-supported-on},
notons $i : T \to X$ l'inclusion. Rappelons que, d'après
Cohomologie, section \ref{cohomology-section-cohomology-support-bis},
on dispose dans cette situation d'un foncteur
$R\mathcal{H}_T : D(\mathcal{O}_X) \to D(i^{-1}\mathcal{O}_X)$
adjoint à droite de $i_* : D(i^{-1}\mathcal{O}_X) \to D(\mathcal{O}_X)$.

\begin{lemma}
\label{perfect-lemma-support-quasi-coherent}
Soit $X$ un schéma. Soit $T \subset X$ une partie fermée telle que
$X \setminus T$ soit un ouvert rétrocompact de $X$. Soit $i : T \to X$
l'inclusion.
\begin{enumerate}
\item Pour $E$ dans $D_\QCoh(\mathcal{O}_X)$, on a
$i_*R\mathcal{H}_T(E)$ dans $D_{\QCoh, T}(\mathcal{O}_X)$.
\item Le foncteur
$i_* \circ R\mathcal{H}_T : D_\QCoh(\mathcal{O}_X) \to
D_{\QCoh, T}(\mathcal{O}_X)$ est adjoint à droite du foncteur d'inclusion
$D_{\QCoh, T}(\mathcal{O}_X) \to D_\QCoh(\mathcal{O}_X)$.
\end{enumerate}
\end{lemma}

\begin{proof}
Posons $U = X \setminus T$ et notons $j : U \to X$ l'inclusion. D'après
Cohomologie, lemme \ref{cohomology-lemma-triangle-sections-with-support-sheaves},
il existe un triangle distingué
$$
i_*R\mathcal{H}_T(E) \to E \to Rj_*(E|_U) \to i_*R\mathcal{H}_Z(E)[1]
$$
dans $D(\mathcal{O}_X)$. D'après le lemme \ref{perfect-lemma-quasi-coherence-direct-image},
le complexe $Rj_*(E|_U)$ a des faisceaux de cohomologie quasi-cohérents
(c'est ici que l'on utilise la rétrocompacité de $U$ dans $X$).
On obtient ainsi (1). L'assertion (2) en résulte, compte tenu de
l'adjonction des foncteurs donnée dans
Cohomologie, lemme \ref{cohomology-lemma-complexes-with-support-on-closed}.
\end{proof}

\begin{lemma}
\label{perfect-lemma-support-direct-sums}
Soit $X$ un schéma. Soit $T \subset X$ une partie fermée telle que
$X \setminus T$ soit un ouvert rétrocompact de $X$. Alors, pour une famille
d'objets $E_i$, $i \in I$ de $D_\QCoh(\mathcal{O}_X)$, on a
$R\mathcal{H}_T(\bigoplus E_i) = \bigoplus R\mathcal{H}_T(E_i)$.
\end{lemma}

\begin{proof}
Posons $U = X \setminus T$ et notons $j : U \to X$ l'inclusion. D'après
Cohomologie, lemme \ref{cohomology-lemma-triangle-sections-with-support-sheaves},
il existe un triangle distingué
$$
i_*R\mathcal{H}_T(E) \to E \to Rj_*(E|_U) \to i_*R\mathcal{H}_Z(E)[1]
$$
dans $D(\mathcal{O}_X)$ pour tout $E$ dans $D(\mathcal{O}_X)$. Le foncteur
$E \mapsto Rj_*(E|_U)$ commute aux sommes directes sur $D_\QCoh(\mathcal{O}_X)$
d'après le lemme \ref{perfect-lemma-quasi-coherence-pushforward-direct-sums}.
Il en résulte qu'il en est de même du foncteur $i_* \circ R\mathcal{H}_T$
(nous omettons les détails). Puisque $i_* : D(i^{-1}\mathcal{O}_X) \to D_T(\mathcal{O}_X)$
est une équivalence
(Cohomologie, lemme \ref{cohomology-lemma-complexes-with-support-on-closed}),
on conclut.
\end{proof}

\begin{remark}
\label{perfect-remark-support-c-equations}
Soit $X$ un schéma. Soient $f_1, \ldots, f_c \in \Gamma(X, \mathcal{O}_X)$.
Notons $Z \subset X$ le sous-schéma fermé défini par $f_1, \ldots, f_c$.
Pour $0 \leq p < c$ et $1 \leq i_0 < \ldots < i_p \leq c$, on note
$U_{i_0 \ldots i_p} \subset X$ le sous-schéma ouvert sur lequel
$f_{i_0} \ldots f_{i_p}$ est inversible. Pour tout $\mathcal{O}_X$-module
$\mathcal{F}$, on pose
$$
\mathcal{F}_{i_0 \ldots i_p} =
(U_{i_0 \ldots i_p} \to X)_*(\mathcal{F}|_{U_{i_0 \ldots i_p}})
$$
Dans cette situation, le {\it complexe de {\v C}ech alterné augmenté}
est le complexe de $\mathcal{O}_X$-modules
\begin{equation}
\label{perfect-equation-extended-alternating}
0 \to \mathcal{F} \to
\bigoplus\nolimits_{i_0} \mathcal{F}_{i_0} \to
\ldots \to
\bigoplus\nolimits_{i_0 < \ldots < i_p} \mathcal{F}_{i_0 \ldots i_p} \to
\ldots \to \mathcal{F}_{1 \ldots c} \to 0
\end{equation}
où $\mathcal{F}$ est placé en degré $0$. Les morphismes se construisent
comme suit. Étant donnés
$1 \leq i_0 < \ldots < i_{p + 1} \leq c$ et $0 \leq j \leq p + 1$, on
dispose du morphisme canonique
$$
\mathcal{F}_{i_0 \ldots \hat i_j \ldots i_{p + 1}} \to
\mathcal{F}_{i_0 \ldots i_p}
$$
provenant de l'inclusion
$U_{i_0 \ldots i_p} \subset U_{i_0 \ldots \hat i_j \ldots i_{p + 1}}$.
Les différentielles du complexe alterné augmenté sont définies à l'aide de ces
morphismes canoniques, affectés du signe $(-1)^j$.
\end{remark}

\begin{lemma}
\label{perfect-lemma-extended-alternating-zero}
Avec $X$, $f_1, \ldots, f_c \in \Gamma(X, \mathcal{O}_X)$ et
$\mathcal{F}$ comme dans la remarque \ref{perfect-remark-support-c-equations},
la restriction du complexe (\ref{perfect-equation-extended-alternating}) à
$X \setminus Z$ est acyclique.
\end{lemma}

\noindent
Signalons que ce lemme vaut plus généralement pour tout
complexe de {\v C}ech alterné augmenté défini comme dans la
remarque \ref{perfect-remark-support-c-equations} à partir d'un recouvrement
ouvert fini $X \setminus Z = U_1 \cup \ldots \cup U_c$.

\begin{proof}
Soit $W \subset X \setminus Z$ un ouvert. En évaluant le complexe
de faisceaux (\ref{perfect-equation-extended-alternating}) sur $W$, on obtient le
complexe
$$
\mathcal{F}(W) \to \bigoplus\nolimits_{i_0} \mathcal{F}(U_{i_0} \cap W) \to
\bigoplus\nolimits_{i_0 < i_1} \mathcal{F}(U_{i_0i_1} \cap W) \to \ldots
$$
Autrement dit, on obtient le complexe de {\v C}ech ordonné augmenté
associé au recouvrement $W = \bigcup U_i \cap W$ et à l'ordre
usuel sur $\{1, \ldots, c\}$ ; voir
Cohomologie, section \ref{cohomology-section-alternating-cech}.
D'après Cohomologie, lemme \ref{cohomology-lemma-alternating-cech-trivial},
ce complexe est homotope à zéro dès que $W$ est contenu dans
$V(f_i)$ pour un certain $1 \leq i \leq c$. Cela achève la démonstration.
\end{proof}

\begin{remark}
\label{perfect-remark-extended-alternating-map-to-support}
Soient $X$, $f_1, \ldots, f_c \in \Gamma(X, \mathcal{O}_X)$ et
$\mathcal{F}$ comme dans la remarque \ref{perfect-remark-support-c-equations}.
Notons $\mathcal{F}^\bullet$ le complexe
(\ref{perfect-equation-extended-alternating}). D'après
le lemme \ref{perfect-lemma-extended-alternating-zero},
les faisceaux de cohomologie de $\mathcal{F}^\bullet$
sont à support dans $Z$ ; par suite, $\mathcal{F}^\bullet$ est un objet de
$D_Z(\mathcal{O}_X)$. D'autre part, l'égalité
$\mathcal{F}^0 = \mathcal{F}$ définit un morphisme canonique
$\mathcal{F}^\bullet \to \mathcal{F}$ dans $D(\mathcal{O}_X)$.
Puisque $i_* \circ R\mathcal{H}_Z$ est adjoint à droite du
foncteur d'inclusion $D_Z(\mathcal{O}_X) \to D(\mathcal{O}_X)$, voir
Cohomologie, lemme \ref{cohomology-lemma-complexes-with-support-on-closed},
on obtient un diagramme commutatif canonique
$$
\xymatrix{
\mathcal{F}^\bullet \ar[rd] \ar[rr] & & \mathcal{F} \\
& i_*R\mathcal{H}_Z(\mathcal{F}) \ar[ru]
}
$$
dans $D(\mathcal{O}_X)$, fonctoriel par rapport au $\mathcal{O}_X$-module $\mathcal{F}$.
\end{remark}

\begin{lemma}
\label{perfect-lemma-extended-alternating-represented}
Soient $X$, $f_1, \ldots, f_c \in \Gamma(X, \mathcal{O}_X)$ et
$\mathcal{F}$ comme dans la remarque \ref{perfect-remark-support-c-equations}.
Si $\mathcal{F}$ est quasi-cohérent, le complexe
(\ref{perfect-equation-extended-alternating}) représente
$i_* R\mathcal{H}_Z(\mathcal{F})$ dans $D_Z(\mathcal{O}_X)$.
\end{lemma}

\begin{proof}
Notons $\mathcal{F}^\bullet$ le complexe
(\ref{perfect-equation-extended-alternating}).
L'énoncé du lemme signifie que le morphisme
$\mathcal{F}^\bullet \to i_*R\mathcal{H}_Z(\mathcal{F})$
de la remarque \ref{perfect-remark-extended-alternating-map-to-support}
est un isomorphisme. Puisque $\mathcal{F}^\bullet$ appartient à
$D_Z(\mathcal{O}_X)$ (voir la remarque citée), on a
$i_*R\mathcal{H}_Z(\mathcal{F}^\bullet) = \mathcal{F}^\bullet$
d'après Cohomologie, lemme \ref{cohomology-lemma-complexes-with-support-on-closed}.
Le morphisme $U_{i_0 \ldots i_p} \to X$ est affine,
car il est donné, au-dessus des ouverts affines de $X$, par inversion de la fonction
$f_{i_0} \ldots f_{i_p}$. On a donc
$$
\mathcal{F}_{i_0 \ldots i_p} =
(U_{i_0 \ldots i_p} \to X)_*\mathcal{F}|_{U_{i_0 \ldots i_p}} =
R(U_{i_0 \ldots i_p} \to X)_*\mathcal{F}|_{U_{i_0 \ldots i_p}}
$$
d'après Cohomologie des schémas, lemme \ref{coherent-lemma-relative-affine-vanishing},
et l'hypothèse de quasi-cohérence de $\mathcal{F}$. On en déduit que
$R\mathcal{H}_Z(\mathcal{F}_{i_0 \ldots i_p}) = 0$ d'après Cohomologie, lemme
\ref{cohomology-lemma-sections-support-in-closed-disjoint-open}.
Ainsi $i_*R\mathcal{H}_Z(\mathcal{F}^p) = 0$ pour $p > 0$.
En réunissant ces résultats, on obtient
$$
\mathcal{F}^\bullet = i_*R\mathcal{H}_Z(\mathcal{F}^\bullet) =
i_*R\mathcal{H}_Z(\mathcal{F})
$$
comme voulu.
\end{proof}

\begin{lemma}
\label{perfect-lemma-supported-trivial-vanishing}
Soit $X$ un schéma. Soit $T \subset X$ une partie fermée qui peut
être définie localement par au plus $c$ éléments du faisceau structural.
Alors $\mathcal{H}^i_Z(\mathcal{F}) = 0$ pour $i > c$ et tout
$\mathcal{O}_X$-module quasi-cohérent $\mathcal{F}$.
\end{lemma}

\begin{proof}
Cela résulte immédiatement de la description locale de
$R\mathcal{H}_T(\mathcal{F})$ donnée dans
le lemme \ref{perfect-lemma-extended-alternating-represented}.
\end{proof}

\begin{lemma}
\label{perfect-lemma-supported-vanishing}
Soit $X$ un schéma. Soit $Z \subset X$ une partie fermée qui peut
être définie localement par une suite régulière au sens de Koszul ayant $c$ éléments.
Alors $\mathcal{H}^i_Z(\mathcal{F}) = 0$ pour $i \not = c$ et tout
$\mathcal{O}_X$-module plat et quasi-cohérent $\mathcal{F}$.
\end{lemma}

\begin{proof}
D'après la description de $R\mathcal{H}_Z(\mathcal{F})$ donnée dans
le lemme \ref{perfect-lemma-extended-alternating-represented}, on se ramène
à l'énoncé algébrique suivant : étant donnés un anneau $R$,
une suite régulière au sens de Koszul $f_1, \ldots, f_c \in R$ et un
$R$-module plat $M$, le complexe de {\v C}ech alterné augmenté
$M \to \bigoplus\nolimits_{i_0} M_{f_{i_0}} \to
\bigoplus\nolimits_{i_0 < i_1} M_{f_{i_0}f_{i_1}} \to
\ldots \to M_{f_1 \ldots f_c}$
de Compléments sur l'algèbre, section \ref{more-algebra-section-alternating-cech},
n'a de cohomologie qu'en degré $c$. D'après Compléments sur l'algèbre, lemme
\ref{more-algebra-lemma-vanishing-extended-alternating-koszul},
on obtient l'annulation voulue pour le complexe de
{\v C}ech alterné augmenté de $R$. Puisque le complexe associé à $M$ s'obtient
en tensorisant celui-ci par le $R$-module plat $M$
(Compléments sur l'algèbre, lemme
\ref{more-algebra-lemma-extended-alternating-form-module}),
on conclut.
\end{proof}

\begin{remark}
\label{perfect-remark-supported-map-c-equations}
Soient $X$, $f_1, \ldots, f_c \in \Gamma(X, \mathcal{O}_X)$ et
$\mathcal{F}$ comme dans la remarque \ref{perfect-remark-support-c-equations}.
Il existe un morphisme $\mathcal{O}_X|_Z$-linéaire canonique
$$
c_{f_1, \ldots, f_c} :
i^*\mathcal{F}
\longrightarrow
\mathcal{H}^c_Z(\mathcal{F})
$$
fonctoriel en $\mathcal{F}$. En effet, si l'on note $\mathcal{F}^\bullet$ le
complexe de {\v C}ech alterné augmenté (\ref{perfect-equation-extended-alternating}),
on dispose du morphisme canonique
$\mathcal{F}^\bullet \to i_*R\mathcal{H}_Z(\mathcal{F})$
de la remarque \ref{perfect-remark-extended-alternating-map-to-support}.
Il définit un morphisme canonique
$$
\Coker\left(\bigoplus \mathcal{F}_{1 \ldots \hat i \ldots c} \to
\mathcal{F}_{1 \ldots c}\right)
\longrightarrow
i_*\mathcal{H}^c_Z(\mathcal{F})
$$
sur les faisceaux de cohomologie en degré $c$.
Étant donnée une section locale $s$ de $\mathcal{F}$, on peut considérer la
section locale
$$
\frac{s}{f_1 \ldots f_c}
$$
de $\mathcal{F}_{1 \ldots c}$. La classe de cette section dans le conoyau
ci-dessus ne dépend que de $s$ modulo l'image de
$(f_1, \ldots, f_c) : \mathcal{F}^{\oplus c} \to \mathcal{F}$.
Puisque $i_*i^*\mathcal{F}$ est égal au conoyau de
$(f_1, \ldots, f_c) : \mathcal{F}^{\oplus c} \to \mathcal{F}$,
on obtient un morphisme de $\mathcal{O}_X$-modules
$i_*i^*\mathcal{F} \to i_*\mathcal{H}_Z^c(\mathcal{F})$.
La pleine fidélité de $i_*$ fournit alors le morphisme $c_{f_1, \ldots, f_c}$.
\end{remark}

\begin{example}
\label{perfect-example-affine-supported-map-c-equations}
Soient $X = \Spec(A)$ affine, $f_1, \ldots, f_c \in A$ et
$\mathcal{F} = \widetilde{M}$ pour un $A$-module $M$. Le morphisme
$c_{f_1, \ldots, f_c}$ de la remarque \ref{perfect-remark-supported-map-c-equations}
se décrit comme le morphisme
$$
M/(f_1, \ldots, f_c)M
\longrightarrow
\Coker\left(
\bigoplus M_{f_1 \ldots \hat f_i \ldots f_c} \to M_{f_1 \ldots f_c}
\right)
$$
qui envoie la classe de $s \in M$ sur la classe de $s/f_1 \ldots f_c$
dans le conoyau.
\end{example}

\begin{lemma}
\label{perfect-lemma-supported-map-determinant}
Soient $X$, $f_1, \ldots, f_c \in \Gamma(X, \mathcal{O}_X)$ et
$\mathcal{F}$ comme dans la remarque \ref{perfect-remark-support-c-equations}.
Soient $a_{ji} \in \Gamma(X, \mathcal{O}_X)$ pour $1 \leq i, j \leq c$,
et posons $g_j = \sum_{i = 1, \ldots, c} a_{ji}f_i$. Supposons que $g_1, \ldots, g_c$
définissent $Z$ comme sous-schéma. Si $\mathcal{F}$ est quasi-cohérent, on a
$$
c_{f_1, \ldots, f_c} = \det(a_{ji}) c_{g_1, \ldots, g_c}
$$
où $c_{f_1, \ldots, f_c}$ et $c_{g_1, \ldots, g_c}$ sont définis dans la
remarque \ref{perfect-remark-supported-map-c-equations}.
\end{lemma}

\begin{proof}
Nous montrerons que $c_{f_1, \ldots, f_c}(s) =
\det(a_{ij}) c_{g_1, \ldots, g_c}(s)$ en tant que sections globales de
$\mathcal{H}_Z(\mathcal{F})$, pour tout $s \in \mathcal{F}(X)$.
Cela suffit, car on obtient alors le même résultat pour les
sections sur tout sous-schéma ouvert de $X$.
Pour cela, pour $1 \leq i_0 < \ldots < i_p \leq c$
et $1 \leq j_0 < \ldots < j_q \leq c$, notons
$U_{i_0 \ldots i_p} \subset X$,
$V_{j_0 \ldots j_q} \subset X$ et
$W_{i_0 \ldots i_p, j_0 \ldots j_q} \subset X$
les sous-schémas ouverts où, respectivement, $f_{i_0} \ldots f_{i_p}$ est inversible,
$g_{j_0} \ldots g_{j_q}$ est inversible, et
$f_{i_0} \ldots f_{i_p}g_{j_0} \ldots g_{j_q}$ est inversible.
On note $\mathcal{F}_{i_0 \ldots i_p}$,
resp.\ $\mathcal{F}'_{j_0 \ldots j_q}$,
$\mathcal{F}''_{i_0 \ldots i_p, j_0 \ldots j_q}$
l'image directe sur $X$ de la restriction de $\mathcal{F}$ à
$U_{i_0 \ldots i_p}$, resp.\ $V_{j_0 \ldots j_q}$,
resp.\ $W_{i_0 \ldots i_p, j_0 \ldots j_q}$.
On obtient alors trois complexes de {\v C}ech alternés augmentés
$$
\mathcal{F}^\bullet :
\mathcal{F} \to \bigoplus\nolimits_{i_0} \mathcal{F}_{i_0} \to
\bigoplus\nolimits_{i_0 < i_1} \mathcal{F}_{i_0i_1} \to \ldots
$$
et
$$
(\mathcal{F}')^\bullet :
\mathcal{F} \to \bigoplus\nolimits_{j_0} \mathcal{F}'_{j_0} \to
\bigoplus\nolimits_{j_0 < j_1} \mathcal{F}'_{j_0j_1} \to \ldots
$$
et
$$
(\mathcal{F}'')^\bullet :
\mathcal{F} \to
\bigoplus\nolimits_{i_0} \mathcal{F}_{i_0} \oplus
\bigoplus\nolimits_{j_0} \mathcal{F}'_{j_0} \to
\bigoplus\nolimits_{i_0 < i_1} \mathcal{F}_{i_0i_1} \oplus
\bigoplus\nolimits_{i_0, j_0} \mathcal{F}''_{i_0, j_0} \oplus
\bigoplus\nolimits_{j_0 < j_1} \mathcal{F}'_{j_0j_1} \to
\ldots
$$
dont les différentielles sont celles qui interviennent dans la définition de
(\ref{perfect-equation-extended-alternating}).
On dispose de morphismes de complexes
$$
(\mathcal{F}'')^\bullet \to \mathcal{F}^\bullet
\quad\text{et}\quad
(\mathcal{F}'')^\bullet \to (\mathcal{F}')^\bullet
$$
donnés terme à terme par les projections (et induisant donc
l'identité en degré $0$). Notons que, d'après
le lemme \ref{perfect-lemma-extended-alternating-represented},
chacun de ces complexes représente
$i_*R\mathcal{H}_Z(\mathcal{F})$, et ces morphismes
représentent l'identité de cet objet. Il suffit donc
de trouver un élément
$$
\sigma \in H^c((\mathcal{F}'')^\bullet(X))
$$
dont les images par ces deux morphismes soient $c_{f_1, \ldots, f_c}(s)$ et $\det(a_{ji})c_{g_1, \ldots, g_c}(s)$.
On peut en fait donner explicitement un cocycle
représentant $\sigma$. On prend
$$
\sigma_{1 \ldots c} = \frac{s}{f_1 \ldots f_c} \in \mathcal{F}_{1 \ldots c}(X)
\quad\text{et}\quad
\sigma'_{1 \ldots c} = \frac{\det(a_{ji})s}{g_1 \ldots g_c} \in
\mathcal{F}'_{1 \ldots c}(X)
$$
et l'on prend
$$
\sigma_{i_0 \ldots i_p, j_0 \ldots j_{c - p - 2}} =
\frac{\lambda(i_0 \ldots i_p, j_0 \ldots j_{c - p - 2})s}%
{f_{i_0} \ldots f_{i_p}g_{j_0} \ldots g_{j_{c - p - 2}}}
\in \mathcal{F}''_{i_0 \ldots i_p, j_0 \ldots j_{c - p - 2}}(X)
$$
où $\lambda(i_0 \ldots i_p, j_0 \ldots j_{c - p - 2})$
est le coefficient de $e_1 \wedge \ldots \wedge e_c$ dans
l'expression formelle
$$
e_{i_0} \wedge \ldots \wedge e_{i_p} \wedge
(a_{j_01} e_1 + \ldots + a_{j_0c}e_c) \wedge \ldots \wedge
(a_{j_{c - p - 2}1} e_1 + \ldots + a_{j_{c - p - 2}c}e_c)
$$
Pour vérifier que $\sigma$ est un cocycle, il faut montrer que, pour
$1 \leq i_0 < \ldots < i_p \leq c$ et
$1 \leq j_0 < \ldots < j_{c - p - 1} \leq c$,
on a
\begin{align*}
0 & =
\sum\nolimits_{a = 0, \ldots, p} (-1)^a
f_{i_a} \lambda(i_0 \ldots \hat i_a \ldots i_p, j_0 \ldots j_{c - p - 1}) \\
& +
\sum\nolimits_{b = 0, \ldots, c - p - 1} (-1)^{p + b + 1}g_{j_b}
\lambda(i_0 \ldots i_p, j_0 \ldots \hat j_b \ldots j_{c - p - 1})
\end{align*}
Le plus simple est peut-être d'observer que l'expression formelle
$$
\xi = e_{i_0} \wedge \ldots \wedge e_{i_p} \wedge
(a_{j_01} e_1 + \ldots + a_{j_0c}e_c) \wedge \ldots \wedge
(a_{j_{c - p - 1}1} e_1 + \ldots + a_{j_{c - p - 1}c}e_c)
$$
est égale à $0$, puisqu'elle appartient à la $(c + 1)$-ième puissance extérieure du module libre
de base $e_1, \ldots, e_c$, et que l'expression ci-dessus est l'image de
$\xi$ par la différentielle de Koszul qui envoie $e_i \to f_i$. Nous omettons quelques
détails.
\end{proof}

\begin{lemma}
\label{perfect-lemma-supported-map-global}
Soit $X$ un schéma. Soit $Z \to X$ une immersion fermée de présentation finie
dont le faisceau conormal $\mathcal{C}_{Z/X}$ est localement libre de rang $c$.
Il existe alors un morphisme canonique
$$
c :
\wedge^c(\mathcal{C}_{Z/X})^\vee \otimes_{\mathcal{O}_Z} i^*\mathcal{F}
\longrightarrow
\mathcal{H}_Z^c(\mathcal{F})
$$
fonctoriel par rapport au module quasi-cohérent $\mathcal{F}$.
\end{lemma}

\begin{proof}
Cela résulte de la construction de la
remarque \ref{perfect-remark-supported-map-c-equations}
et de l'indépendance du choix
des générateurs du faisceau d'idéaux établie dans le
lemme \ref{perfect-lemma-supported-map-determinant}.
Nous omettons quelques détails.
\end{proof}

\begin{remark}
\label{perfect-remark-supported-functorial}
Soit $g : X' \to X$ un morphisme de schémas. Soient
$f_1, \ldots, f_c \in \Gamma(X, \mathcal{O}_X)$.
Posons $f'_i = g^\sharp(f_i) \in \Gamma(X', \mathcal{O}_{X'})$.
Notons $Z \subset X$, resp.\ $Z' \subset X'$, le sous-schéma
fermé défini par $f_1, \ldots, f_c$, resp.\ $f'_1, \ldots, f'_c$.
Alors $Z' = Z \times_X X'$. Notons $h : Z' \to Z$ le morphisme
de schémas induit.
Soit $\mathcal{F}$ un $\mathcal{O}_X$-module.
Posons $\mathcal{F}' = g^*\mathcal{F}$.
Dans cette situation, si $\mathcal{F}$ est quasi-cohérent, le diagramme
$$
\xymatrix{
(i')^{-1}\mathcal{O}_{X'} \otimes_{h^{-1}i^{-1}\mathcal{O}_X}
h^{-1}\mathcal{H}^c_Z(\mathcal{F}) \ar[r] &
\mathcal{H}_{Z'}^c(\mathcal{F}') \\
h^*i^*\mathcal{F} \ar[r] \ar[u]_-{c_{f_1, \ldots, f_c}} &
(i')^*\mathcal{F}' \ar[u]^-{c_{f'_1, \ldots, f'_c}}
}
$$
est commutatif ; la flèche horizontale supérieure est
le morphisme de Cohomologie, remarque \ref{cohomology-remark-support-functorial},
sur les faisceaux de cohomologie en degré $c$. En effet,
notons $\mathcal{F}^\bullet$, resp.\ $(\mathcal{F}')^\bullet$,
le complexe de {\v C}ech alterné augmenté construit dans la
remarque \ref{perfect-remark-support-c-equations}
à l'aide de $\mathcal{F}, f_1, \ldots, f_c$,
resp.\ $\mathcal{F}', f'_1, \ldots, f'_c$.
Notons que $(\mathcal{F}')^\bullet = g^*\mathcal{F}^\bullet$.
Alors, sans supposer $\mathcal{F}$ quasi-cohérent, le diagramme
$$
\xymatrix{
i'_* L(g|_{Z'})^* R\mathcal{H}_Z(\mathcal{F}) \ar[r] \ar@{=}[d] &
i'_*R\mathcal{H}_{Z'}(Lg^*\mathcal{F}) \ar[d] \\
Lg^*i_*R\mathcal{H}_Z(\mathcal{F}) &
i'_*R\mathcal{H}_{Z'}(\mathcal{F}') \\
Lg^*(\mathcal{F}^\bullet) \ar[u] \ar[r] &
(\mathcal{F}')^\bullet \ar[u]
}
$$
est commutatif, où $g|_{Z'} : (Z', (i')^{-1}\mathcal{O}_{X'}) \to
(Z, i^{-1}\mathcal{O}_X)$ est le morphisme induit d'espaces annelés.
La flèche horizontale supérieure est donnée dans
Cohomologie, remarque \ref{cohomology-remark-support-functorial},
qui explique aussi le signe d'égalité.
Les flèches ascendantes proviennent de la
remarque \ref{perfect-remark-extended-alternating-map-to-support}.
La flèche horizontale inférieure est le morphisme $Lg^*\mathcal{F}^\bullet
\to g^*\mathcal{F}^\bullet = (\mathcal{F}')^\bullet$, et la flèche
descendante est induite par
$Lg^*\mathcal{F} \to g^*\mathcal{F} = \mathcal{F}'$.
Le diagramme est commutatif, car les deux parcours donnent
deux flèches $Lg^*\mathcal{F}^\bullet \to
i'_*R\mathcal{H}_{Z'}(\mathcal{F}')$ dont le composé avec
$i'_*R\mathcal{H}_{Z'}(\mathcal{F}') \to \mathcal{F}'$
est le morphisme canonique $Lg^*\mathcal{F}^\bullet \to \mathcal{F}'$.
Nous omettons quelques détails. La commutativité du premier diagramme
s'obtient alors en passant aux faisceaux de cohomologie en degré
$c$ dans ce diagramme et en utilisant la compatibilité à l'image réciproque de la construction du morphisme
$i^*\mathcal{F} \to
\Coker(\bigoplus \mathcal{F}_{1 \ldots \hat i \ldots c} \to
\mathcal{F}_{1 \ldots c})$
employée dans la remarque \ref{perfect-remark-supported-map-c-equations}.
\end{remark}







\section{Le cohérateur}
\label{perfect-section-coherator}

\noindent
Soit $X$ un schéma. Le {\it cohérateur} est un foncteur
$$
Q_X :
\textit{Mod}(\mathcal{O}_X)
\longrightarrow
\QCoh(\mathcal{O}_X)
$$
adjoint à droite du foncteur d'inclusion
$\QCoh(\mathcal{O}_X) \to \textit{Mod}(\mathcal{O}_X)$.
Il existe pour tout schéma $X$ ; de plus, le morphisme d'adjonction
$Q_X(\mathcal{F}) \to \mathcal{F}$ est un isomorphisme pour tout
module quasi-cohérent $\mathcal{F}$ ; voir
Propriétés, proposition \ref{properties-proposition-coherator}.
Puisque $Q_X$ est exact à gauche (comme adjoint à droite), on peut considérer son
foncteur dérivé à droite
$$
RQ_X :
D(\mathcal{O}_X)
\longrightarrow
D(\QCoh(\mathcal{O}_X)).
$$
Puisque $Q_X$ est adjoint à droite du foncteur d'inclusion
$\QCoh(\mathcal{O}_X) \to \textit{Mod}(\mathcal{O}_X)$,
le foncteur $RQ_X$ est adjoint à droite du foncteur canonique
$D(\QCoh(\mathcal{O}_X)) \to D(\mathcal{O}_X)$ d'après
Catégories dérivées, lemme \ref{derived-lemma-derived-adjoint-functors}.

\medskip\noindent
Dans cette section, nous étudierons le foncteur $RQ_X$. Dans la
section \ref{perfect-section-better-coherator},
nous étudierons l'adjoint à droite, étroitement lié au précédent, du foncteur d'inclusion
$D_\QCoh(\mathcal{O}_X) \to D(\mathcal{O}_X)$ (lorsqu'il existe).

\begin{lemma}
\label{perfect-lemma-affine-pushforward}
Soit $f : X \to Y$ un morphisme affine de schémas.
Alors $f_*$ définit un foncteur dérivé
$f_* : D(\QCoh(\mathcal{O}_X)) \to D(\QCoh(\mathcal{O}_Y))$.
Ce foncteur rend le diagramme
$$
\xymatrix{
D(\QCoh(\mathcal{O}_X)) \ar[d]_{f_*} \ar[r] &
D_\QCoh(\mathcal{O}_X) \ar[d]^{Rf_*} \\
D(\QCoh(\mathcal{O}_Y)) \ar[r] &
D_\QCoh(\mathcal{O}_Y)
}
$$
commutatif.
\end{lemma}

\begin{proof}
Le foncteur
$f_* : \QCoh(\mathcal{O}_X) \to \QCoh(\mathcal{O}_Y)$
est exact ; voir
Cohomologie des schémas, lemme \ref{coherent-lemma-relative-affine-vanishing}.
Par suite, $f_*$ définit un foncteur dérivé
$f_* : D(\QCoh(\mathcal{O}_X)) \to D(\QCoh(\mathcal{O}_Y))$
en appliquant simplement $f_*$ à un complexe représentant quelconque ; voir
Catégories dérivées, lemme \ref{derived-lemma-right-derived-exact-functor}.
Le diagramme est commutatif d'après le lemme \ref{perfect-lemma-pushforward-affine-morphism}.
\end{proof}

\begin{lemma}
\label{perfect-lemma-flat-pushforward-coherator}
Soit $f : X \to Y$ un morphisme de schémas. Supposons $f$
quasi-compact, quasi-séparé et plat. Alors, en notant
$$
\Phi : D(\QCoh(\mathcal{O}_X)) \to D(\QCoh(\mathcal{O}_Y))
$$
le foncteur dérivé à droite de
$f_* : \QCoh(\mathcal{O}_X) \to \QCoh(\mathcal{O}_Y)$,
on a $RQ_Y \circ Rf_* = \Phi \circ RQ_X$.
\end{lemma}

\begin{proof}
Nous le démontrerons en montrant que $RQ_Y \circ Rf_*$ et $\Phi \circ RQ_X$
sont adjoints à droite du même foncteur
$D(\QCoh(\mathcal{O}_Y)) \to D(\mathcal{O}_X)$.

\medskip\noindent
Puisque $f$ est quasi-compact et quasi-séparé,
$f_*$ préserve la quasi-cohérence ; voir
Schémas, lemme \ref{schemes-lemma-push-forward-quasi-coherent}.
Rappelons que $\QCoh(\mathcal{O}_X)$ est une catégorie abélienne de Grothendieck
(Propriétés, proposition \ref{properties-proposition-coherator}).
Par suite, tout $K$ dans $D(\QCoh(\mathcal{O}_X))$
peut être représenté par un complexe K-injectif $\mathcal{I}^\bullet$
de $\QCoh(\mathcal{O}_X)$ ; voir
Injectifs, théorème
\ref{injectives-theorem-K-injective-embedding-grothendieck}.
On peut alors définir $\Phi(K) = f_*\mathcal{I}^\bullet$.

\medskip\noindent
Puisque $f$ est plat, le foncteur $f^*$ est exact. Par suite, $f^*$ définit
$f^* : D(\mathcal{O}_Y) \to D(\mathcal{O}_X)$ ainsi que
$f^* : D(\QCoh(\mathcal{O}_Y)) \to D(\QCoh(\mathcal{O}_X))$.
Le foncteur $f^* = Lf^* : D(\mathcal{O}_Y) \to D(\mathcal{O}_X)$
est adjoint à gauche de
$Rf_* : D(\mathcal{O}_X) \to D(\mathcal{O}_Y)$ ;
voir Cohomologie, lemme \ref{cohomology-lemma-adjoint}.
De même, le foncteur
$f^* : D(\QCoh(\mathcal{O}_Y)) \to D(\QCoh(\mathcal{O}_X))$
est adjoint à gauche de
$\Phi : D(\QCoh(\mathcal{O}_X)) \to D(\QCoh(\mathcal{O}_Y))$
d'après Catégories dérivées, lemme \ref{derived-lemma-derived-adjoint-functors}.

\medskip\noindent
Soient $A$ un objet de $D(\QCoh(\mathcal{O}_Y))$ et
$E$ un objet de $D(\mathcal{O}_X)$. Alors
\begin{align*}
\Hom_{D(\QCoh(\mathcal{O}_Y))}(A, RQ_Y(Rf_*E))
& =
\Hom_{D(\mathcal{O}_Y)}(A, Rf_*E) \\
& =
\Hom_{D(\mathcal{O}_X)}(f^*A, E) \\
& =
\Hom_{D(\QCoh(\mathcal{O}_X))}(f^*A, RQ_X(E)) \\
& =
\Hom_{D(\QCoh(\mathcal{O}_Y))}(A, \Phi(RQ_X(E)))
\end{align*}
On en déduit l'assertion voulue.
\end{proof}

\begin{lemma}
\label{perfect-lemma-affine-coherator}
Soit $X = \Spec(A)$ un schéma affine. Alors
\begin{enumerate}
\item $Q_X : \textit{Mod}(\mathcal{O}_X) \to \QCoh(\mathcal{O}_X)$
est le foncteur
qui envoie $\mathcal{F}$ sur le $\mathcal{O}_X$-module quasi-cohérent
associé au $A$-module $\Gamma(X, \mathcal{F})$,
\item $RQ_X : D(\mathcal{O}_X) \to D(\QCoh(\mathcal{O}_X))$
est le foncteur qui envoie $E$ sur le complexe de
$\mathcal{O}_X$-modules quasi-cohérents associé à l'objet $R\Gamma(X, E)$ de $D(A)$,
\item restreint à $D_\QCoh(\mathcal{O}_X)$, le foncteur
$RQ_X$ définit un quasi-inverse de (\ref{perfect-equation-compare}).
\end{enumerate}
\end{lemma}

\begin{proof}
Le foncteur $Q_X$ est le foncteur
$$
\mathcal{F} \mapsto \widetilde{\Gamma(X, \mathcal{F})}
$$
d'après Schémas, lemme \ref{schemes-lemma-compare-constructions}.
Cela entraîne immédiatement (1) et (2). La troisième assertion
résulte de la démonstration du
lemme \ref{perfect-lemma-affine-compare-bounded}.
\end{proof}

\noindent
Nous pouvons maintenant établir un critère
pour que le foncteur $D(\QCoh(\mathcal{O}_X)) \to D_\QCoh(\mathcal{O}_X)$
soit une équivalence.

\begin{lemma}
\label{perfect-lemma-argument-proves}
Soit $X$ un schéma quasi-compact et quasi-séparé. Supposons que,
pour tout ouvert affine $U \subset X$, le foncteur dérivé à droite
$$
\Phi : D(\QCoh(\mathcal{O}_U)) \to D(\QCoh(\mathcal{O}_X))
$$
du foncteur exact à gauche
$j_* : \QCoh(\mathcal{O}_U) \to \QCoh(\mathcal{O}_X)$
s'inscrive dans un diagramme commutatif
$$
\xymatrix{
D(\QCoh(\mathcal{O}_U)) \ar[d]_\Phi \ar[r]_{i_U} &
D_\QCoh(\mathcal{O}_U) \ar[d]^{Rj_*} \\
D(\QCoh(\mathcal{O}_X)) \ar[r]^{i_X} &
D_\QCoh(\mathcal{O}_X)
}
$$
Alors le foncteur (\ref{perfect-equation-compare})
$$
D(\QCoh(\mathcal{O}_X))
\longrightarrow
D_\QCoh(\mathcal{O}_X)
$$
est une équivalence dont un quasi-inverse est donné par $RQ_X$.
\end{lemma}

\begin{proof}
Soit $E$ un objet de $D_\QCoh(\mathcal{O}_X)$ et
soit $A$ un objet de $D(\QCoh(\mathcal{O}_X))$.
Il faut montrer que les morphismes d'adjonction
$$
A \to RQ_X(i_X(A))
\quad\text{et}\quad
i_X(RQ_X(E)) \to E
$$
sont des isomorphismes. Considérons l'hypothèse
$H_n$ : les morphismes d'adjonction ci-dessus sont des isomorphismes
dès que $E$ et $i_X(A)$ sont à support
(définition \ref{perfect-definition-supported-on})
dans une partie fermée de $X$
contenue dans la réunion de $n$ ouverts affines de $X$.
Nous démontrerons $H_n$ par récurrence sur $n$.

\medskip\noindent
Cas initial : $n = 0$. Dans ce cas, $E = 0$, donc le morphisme
$i_X(RQ_X(E)) \to E$ est un isomorphisme. De même, $i_X(A) = 0$.
Les faisceaux de cohomologie de $i_X(A)$ sont donc nuls. Puisque le foncteur d'inclusion
$\QCoh(\mathcal{O}_X) \to \textit{Mod}(\mathcal{O}_X)$
est pleinement fidèle et exact, on en déduit que les objets de cohomologie
de $A$ sont nuls, c'est-à-dire que $A = 0$ ; le morphisme
$A \to RQ_X(i_X(A))$ est donc lui aussi un isomorphisme.

\medskip\noindent
Étape de récurrence. Supposons que $E$ et $i_X(A)$ soient à support dans une
partie fermée $T$ de $X$ contenue dans $U_1 \cup \ldots \cup U_n$,
où les $U_i \subset X$ sont des ouverts affines. Posons $U = U_n$.
Considérons les triangles distingués
$$
A \to \Phi(A|_U) \to A' \to A[1]
\quad\text{et}\quad
E \to Rj_*(E|_U) \to E' \to E[1]
$$
où $\Phi$ est comme dans l'énoncé du lemme.
Notons que $E \to Rj_*(E|_U)$ est un quasi-isomorphisme sur $U = U_n$.
Puisque $i_X \circ \Phi = Rj_* \circ i_U$ par hypothèse
et que $i_X(A)|_U = i_U(A|_U)$,
le morphisme $i_X(A) \to i_X(\Phi(A|_U))$ est un quasi-isomorphisme sur $U$.
Ainsi $i_X(A')$ et $E'$ sont à support dans la partie
fermée $T \setminus U$ de $X$, contenue dans
$U_1 \cup \ldots \cup U_{n - 1}$.
Par hypothèse de récurrence, l'énoncé est vrai pour $A'$ et $E'$. D'après
Catégories dérivées, lemme \ref{derived-lemma-third-isomorphism-triangle},
il suffit de montrer que les morphismes
$$
\Phi(A|_U) \to RQ_X(i_X(\Phi(A|_U)))
\quad\text{et}\quad
i_X(RQ_X(Rj_*E|_U)) \to Rj_*(E|_U)
$$
sont des isomorphismes. Par hypothèse et d'après
le lemme \ref{perfect-lemma-flat-pushforward-coherator}
(le morphisme d'inclusion $j : U \to X$ est plat, quasi-compact et
quasi-séparé), on a
$$
RQ_X(i_X(\Phi(A|_U))) = RQ_X(Rj_*(i_U(A|_U))) = \Phi(RQ_U(i_U(A|_U)))
$$
et
$$
i_X(RQ_X(Rj_*(E|_U))) = i_X(\Phi(RQ_U(E|_U))) = Rj_*(i_U(RQ_U(E|_U)))
$$
Enfin, les morphismes
$$
A|_U \to RQ_U(i_U(A|_U))
\quad\text{et}\quad
i_U(RQ_U(E|_U)) \to E|_U
$$
sont des isomorphismes d'après le lemme \ref{perfect-lemma-affine-coherator}. Le résultat en découle.
\end{proof}

\begin{proposition}
\label{perfect-proposition-quasi-compact-affine-diagonal}
Soit $X$ un schéma quasi-compact à diagonale affine.
Alors le foncteur (\ref{perfect-equation-compare})
$$
D(\QCoh(\mathcal{O}_X))
\longrightarrow
D_\QCoh(\mathcal{O}_X)
$$
est une équivalence dont un quasi-inverse est donné par $RQ_X$.
\end{proposition}

\begin{proof}
Soit $U \subset X$ un ouvert affine. Alors le morphisme
$U \to X$ est affine d'après
Morphismes, lemme \ref{morphisms-lemma-affine-permanence}.
L'hypothèse du lemme \ref{perfect-lemma-argument-proves}
est donc vérifiée d'après le lemme \ref{perfect-lemma-affine-pushforward}, ce qui conclut.
\end{proof}

\begin{lemma}
\label{perfect-lemma-direct-image-coherator}
Soit $f : X \to Y$ un morphisme de schémas.
Supposons $X$ et $Y$ quasi-compacts et à diagonale affine.
Alors, en notant
$$
\Phi : D(\QCoh(\mathcal{O}_X)) \to D(\QCoh(\mathcal{O}_Y))
$$
le foncteur dérivé à droite de
$f_* : \QCoh(\mathcal{O}_X) \to \QCoh(\mathcal{O}_Y)$,
le diagramme
$$
\xymatrix{
D(\QCoh(\mathcal{O}_X)) \ar[d]_\Phi \ar[r] &
D_\QCoh(\mathcal{O}_X) \ar[d]^{Rf_*} \\
D(\QCoh(\mathcal{O}_Y)) \ar[r] &
D_\QCoh(\mathcal{O}_Y)
}
$$
est commutatif.
\end{lemma}

\begin{proof}
Observons que les flèches horizontales du diagramme sont
des équivalences de catégories d'après la
proposition \ref{perfect-proposition-quasi-compact-affine-diagonal}.
On peut donc identifier ces catégories (et de même pour
les autres schémas quasi-compacts à diagonale affine).
L'énoncé du lemme affirme que le morphisme canonique
$\Phi(K) \to Rf_*(K)$ est un isomorphisme pour tout $K$ dans
$D(\QCoh(\mathcal{O}_X))$. Notons que, si $K_1 \to K_2 \to K_3 \to K_1[1]$
est un triangle distingué dans $D(\QCoh(\mathcal{O}_X))$ et
si l'énoncé est vrai pour deux des trois objets, il l'est
aussi pour le troisième.

\medskip\noindent
Soit $U \subset X$ un ouvert affine. Puisque la diagonale de $X$ est affine,
le morphisme d'inclusion $j : U \to X$
est affine (Morphismes, lemme \ref{morphisms-lemma-affine-permanence}).
De même, le composé $g = f \circ j : U \to Y$ est affine.
Soit $\mathcal{I}^\bullet$ un complexe K-injectif dans $\QCoh(\mathcal{O}_U)$.
Puisque $j_* : \QCoh(\mathcal{O}_U) \to \QCoh(\mathcal{O}_X)$
admet un adjoint à gauche exact
$j^* : \QCoh(\mathcal{O}_X) \to \QCoh(\mathcal{O}_U)$,
le complexe $j_*\mathcal{I}^\bullet$ est K-injectif
dans $\QCoh(\mathcal{O}_X)$ ; voir
Catégories dérivées, lemme \ref{derived-lemma-adjoint-preserve-K-injectives}.
Il en résulte que
$$
\Phi(j_*\mathcal{I}^\bullet) =
f_*j_*\mathcal{I}^\bullet =
g_*\mathcal{I}^\bullet
$$
D'après le lemme \ref{perfect-lemma-affine-pushforward},
$j_*\mathcal{I}^\bullet$ représente $Rj_*\mathcal{I}^\bullet$ et
$g_*\mathcal{I}^\bullet$ représente $Rg_*\mathcal{I}^\bullet$.
D'autre part, on a $Rf_* \circ Rj_* = Rg_*$.
Ainsi $f_*j_*\mathcal{I}^\bullet$ représente $Rf_*(j_*\mathcal{I}^\bullet)$.
On en déduit que le lemme est vrai pour tout complexe
de la forme $j_*\mathcal{G}^\bullet$, où $\mathcal{G}^\bullet$
est un complexe de modules quasi-cohérents sur $U$. (Notons que, si
$\mathcal{G}^\bullet \to \mathcal{I}^\bullet$ est un quasi-isomorphisme,
alors $j_*\mathcal{G}^\bullet \to j_*\mathcal{I}^\bullet$ est aussi un
quasi-isomorphisme, puisque $j_*$ est un foncteur exact
sur les modules quasi-cohérents.)

\medskip\noindent
Soit $\mathcal{F}^\bullet$ un complexe de
$\mathcal{O}_X$-modules quasi-cohérents. Soit $T \subset X$ une partie fermée
telle que le support de $\mathcal{F}^p$ soit contenu dans $T$
pour tout $p$. Nous raisonnons par récurrence sur le nombre minimal $n$
d'ouverts affines $U_1, \ldots, U_n$ tels que
$T \subset U_1 \cup \ldots \cup U_n$. Le cas initial $n = 0$ est immédiat.
Si $n \geq 1$, posons $U = U_1$ et notons $j : U \to X$
l'immersion ouverte, comme ci-dessus. Considérons le morphisme de complexes
$c : \mathcal{F}^\bullet \to j_*j^*\mathcal{F}^\bullet$.
On obtient deux suites exactes courtes de complexes :
$$
0 \to \Ker(c) \to \mathcal{F}^\bullet \to \Im(c) \to 0
$$
et
$$
0 \to \Im(c) \to j_*j^*\mathcal{F}^\bullet \to \Coker(c) \to 0
$$
Les complexes $\Ker(c)$ et $\Coker(c)$ sont à support
dans $T \setminus U \subset U_2 \cup \ldots \cup U_n$, et le résultat
est vrai pour eux par récurrence. Il est vrai pour
$j_*j^*\mathcal{F}^\bullet$ d'après la discussion du
paragraphe précédent. On conclut en considérant les triangles distingués
associés aux suites exactes courtes et en utilisant l'observation
faite au début de la démonstration.
\end{proof}

\begin{remark}[Avertissement]
\label{perfect-remark-warning-coherator}
Soit $X$ un schéma quasi-compact à diagonale affine. Bien que l'on sache
que $D(\QCoh(\mathcal{O}_X)) = D_\QCoh(\mathcal{O}_X)$ d'après la
proposition \ref{perfect-proposition-quasi-compact-affine-diagonal},
des phénomènes inattendus peuvent
se produire, et il est facile de se tromper dans ces questions. On se gardera notamment
de déduire sans précaution de cette égalité que les foncteurs dérivés coïncident.
Par exemple, soit $U \subset X$ un ouvert quasi-compact. On peut alors
considérer les foncteurs dérivés à droite supérieurs
$$
R^i(\QCoh)\Gamma(U, -) : \QCoh(\mathcal{O}_X) \to \textit{Ab}
$$
du foncteur exact à gauche $\Gamma(U, -)$. Puisqu'ils forment un
$\delta$-foncteur universel, et puisque les foncteurs $H^i(U, -)$ (définis pour tous les
faisceaux abéliens sur $X$), restreints à $\QCoh(\mathcal{O}_X)$, forment
un $\delta$-foncteur, on obtient des transformations canoniques
$$
t^i : R^i(\QCoh)\Gamma(U, -) \to H^i(U, -).
$$
Ces transformations ne sont pas en général des isomorphismes, même si $X = \Spec(A)$
est affine ! En effet, on a $R^1(\QCoh)\Gamma(U, \widetilde{I}) = 0$
si $I$ est un $A$-module injectif, par construction des foncteurs dérivés à droite
et par l'équivalence entre $\QCoh(\mathcal{O}_X)$ et $\text{Mod}_A$.
Mais Exemples, lemme \ref{examples-lemma-nonvanishing},
montre qu'il existe $A$, $I$ et $U$ tels que
$H^1(U, \widetilde{I}) \not = 0$.
\end{remark}




\section{Le cohérateur pour les schémas noethériens}
\label{perfect-section-coherator-Noetherian}

\noindent
Dans le cas des schémas noethériens, on peut utiliser le lemme suivant.

\begin{lemma}
\label{perfect-lemma-injective-quasi-coherent-sheaf-Noetherian}
Soit $X$ un schéma noethérien. Soit $\mathcal{J}$ un objet injectif
de $\QCoh(\mathcal{O}_X)$. Alors $\mathcal{J}$
est un faisceau flasque de $\mathcal{O}_X$-modules.
\end{lemma}

\begin{proof}
Soient $U \subset X$ un ouvert et $s \in \mathcal{J}(U)$ une section.
Soit $\mathcal{I} \subset \mathcal{O}_X$ le faisceau quasi-cohérent
d'idéaux définissant la structure de schéma réduit induite sur $X \setminus U$
(voir Schémas, définition \ref{schemes-definition-reduced-induced-scheme}).
D'après Cohomologie des schémas, lemme \ref{coherent-lemma-homs-over-open},
la section $s$ correspond à un morphisme $\sigma : \mathcal{I}^n \to \mathcal{J}$
pour un certain $n$. Puisque $\mathcal{J}$ est un objet injectif de
$\QCoh(\mathcal{O}_X)$, on peut prolonger $\sigma$ en un morphisme
$\tilde s : \mathcal{O}_X \to \mathcal{J}$. Alors $\tilde s$ correspond
à une section globale de $\mathcal{J}$ dont la restriction est $s$.
\end{proof}

\begin{lemma}
\label{perfect-lemma-Noetherian-pushforward}
Soit $f : X \to Y$ un morphisme de schémas noethériens.
Alors le foncteur $f_*$ sur les faisceaux quasi-cohérents admet un foncteur
dérivé à droite
$\Phi : D(\QCoh(\mathcal{O}_X)) \to D(\QCoh(\mathcal{O}_Y))$
tel que le diagramme
$$
\xymatrix{
D(\QCoh(\mathcal{O}_X)) \ar[d]_{\Phi} \ar[r] &
D_\QCoh(\mathcal{O}_X) \ar[d]^{Rf_*} \\
D(\QCoh(\mathcal{O}_Y)) \ar[r] &
D_\QCoh(\mathcal{O}_Y)
}
$$
soit commutatif.
\end{lemma}

\begin{proof}
Puisque $X$ et $Y$ sont des schémas noethériens, le morphisme est quasi-compact
et quasi-séparé (voir
Propriétés, lemme \ref{properties-lemma-locally-Noetherian-quasi-separated}
et
Schémas, remarque \ref{schemes-remark-quasi-compact-and-quasi-separated}).
Ainsi $f_*$ préserve la quasi-cohérence ; voir
Schémas, lemme \ref{schemes-lemma-push-forward-quasi-coherent}.
Soit maintenant $K$ un objet de $D(\QCoh(\mathcal{O}_X))$.
Puisque $\QCoh(\mathcal{O}_X)$ est une catégorie abélienne de Grothendieck
(Propriétés, proposition \ref{properties-proposition-coherator}), on peut
représenter $K$ par un complexe K-injectif $\mathcal{I}^\bullet$
tel que chaque $\mathcal{I}^n$ soit un objet injectif de
$\QCoh(\mathcal{O}_X)$ ; voir
Injectifs, théorème
\ref{injectives-theorem-K-injective-embedding-grothendieck}.
Le foncteur $\Phi$ est donc défini en posant
$$
\Phi(K) = f_*\mathcal{I}^\bullet
$$
où le membre de droite est considéré comme un objet de
$D(\QCoh(\mathcal{O}_Y))$. Pour achever la démonstration du lemme,
il suffit de montrer que le morphisme canonique
$$
f_*\mathcal{I}^\bullet \longrightarrow Rf_*\mathcal{I}^\bullet
$$
est un isomorphisme dans $D(\mathcal{O}_Y)$. Pour cela, d'après
le lemme \ref{perfect-lemma-acyclicity-lemma},
il suffit de montrer que $\mathcal{I}^n$ est acyclique à droite pour
$f_*$, pour tout $n \in \mathbf{Z}$.
Or $\mathcal{I}^n$ est flasque d'après
le lemme \ref{perfect-lemma-injective-quasi-coherent-sheaf-Noetherian},
et les modules flasques sont acycliques à droite pour $f_*$ d'après
Cohomologie, lemme \ref{cohomology-lemma-flasque-acyclic-pushforward}.
\end{proof}

\begin{proposition}
\label{perfect-proposition-Noetherian}
Soit $X$ un schéma noethérien. Alors le foncteur (\ref{perfect-equation-compare})
$$
D(\QCoh(\mathcal{O}_X))
\longrightarrow
D_\QCoh(\mathcal{O}_X)
$$
est une équivalence dont un quasi-inverse est donné par $RQ_X$.
\end{proposition}

\begin{proof}
Cela résulte du lemme \ref{perfect-lemma-argument-proves} et du
lemme \ref{perfect-lemma-Noetherian-pushforward}.
\end{proof}






\section{Complexes de Koszul}
\label{perfect-section-koszul}

\noindent
Soit $A$ un anneau et soit $f_1, \ldots, f_r$ une suite d'éléments
de $A$. Nous avons défini le complexe de Koszul
$K_\bullet(f_1, \ldots, f_r)$ dans
Compléments sur l'algèbre, définition \ref{more-algebra-definition-koszul-complex}.
C'est un complexe de chaînes concentré en degrés $r, \ldots, 0$.
Nous en faisons un complexe de cochaînes $K^\bullet(f_1, \ldots, f_r)$
en posant $K^{-n}(f_1, \ldots, f_r) = K_n(f_1, \ldots, f_r)$
et en conservant les mêmes différentielles. Dans la suite de cette section, tous
les complexes seront des complexes de cochaînes.

\medskip\noindent
Nous définissons un complexe $I^\bullet(f_1, \ldots, f_r)$
de sorte que l'on ait un triangle distingué
$$
I^\bullet(f_1, \ldots, f_r) \to
A \to
K^\bullet(f_1, \ldots, f_r) \to
I^\bullet(f_1, \ldots, f_r)[1]
$$
dans $K(A)$.
Autrement dit, on pose
$$
I^i(f_1, \ldots, f_r) =
\left\{
\begin{matrix}
K^{i - 1}(f_1, \ldots, f_r) & \text{si } i \leq 0 \\
0 & \text{sinon}
\end{matrix}
\right.
$$
et l'on prend l'opposée de la différentielle de $K^\bullet(f_1, \ldots, f_r)$.
Les morphismes du triangle distingué sont les morphismes évidents. Notons que
$I^0(f_1, \ldots, f_r) = A^{\oplus r} \to A$ est donné par
la multiplication par $f_i$ sur le $i$-ième facteur.
Ainsi $I^\bullet(f_1, \ldots, f_r) \to A$ se factorise en
$$
I^\bullet(f_1, \ldots, f_r) \to I \to A
$$
où $I = (f_1, \ldots, f_r)$. On a en fait une suite exacte courte
$$
0 \to H^{-1}(K^\bullet(f_1, \ldots, f_s)) \to
H^0(I^\bullet(f_1, \ldots, f_s)) \to I \to 0
$$
et, pour tout $i < 0$, on a
$H^i(I^\bullet(f_1, \ldots, f_r)) = H^{i - 1}(K^\bullet(f_1, \ldots, f_r))$.
Notons que, pour une seconde suite $g_1, \ldots, g_r$ d'éléments de $A$,
on dispose de morphismes canoniques
$$
I^\bullet(f_1g_1, \ldots, f_rg_r) \to I^\bullet(f_1, \ldots, f_r)
\quad\text{et}\quad
K^\bullet(f_1g_1, \ldots, f_rg_r) \to K^\bullet(f_1, \ldots, f_r)
$$
compatibles aux morphismes décrits ci-dessus. Le premier de ces morphismes est
donné par la multiplication par $g_i$ sur le $i$-ième facteur de
$I^0(f_1g_1, \ldots, f_rg_r) = A^{\oplus r}$. En particulier, étant donnés
$f_1, \ldots, f_r$, on obtient un système projectif de complexes
\begin{equation}
\label{perfect-equation-system}
I^\bullet(f_1, \ldots, f_r) \leftarrow
I^\bullet(f_1^2, \ldots, f_r^2) \leftarrow
I^\bullet(f_1^3, \ldots, f_r^3) \leftarrow \ldots
\end{equation}
qui jouera un rôle important dans la suite.
Fixons quelques notations pour énoncer commodément les lemmes suivants.

\begin{situation}
\label{perfect-situation-complex}
Ici $A$ est un anneau et $f_1, \ldots, f_r$ une suite d'éléments de $A$.
On pose $X = \Spec(A)$ et $U = D(f_1) \cup \ldots \cup D(f_r) \subset X$.
On note $\mathcal{U} : U = \bigcup_{i = 1, \ldots, r} D(f_i)$ le
recouvrement ouvert ainsi défini de $U$.
\end{situation}

\noindent
Notre premier lemme montre que les complexes ci-dessus permettent de calculer
la cohomologie des faisceaux quasi-cohérents sur $U$.

\begin{lemma}
\label{perfect-lemma-alternating-cech-complex}
Dans la situation \ref{perfect-situation-complex}, soit $M$ un $A$-module, et
notons $\mathcal{F}$ le $\mathcal{O}_X$-module associé. Alors
il existe un isomorphisme canonique de complexes
$$
\Psi : \colim_e \Hom_A(I^\bullet(f_1^e, \ldots, f_r^e), M)
\longrightarrow
\check{\mathcal{C}}_{alt}^\bullet(\mathcal{U}, \mathcal{F})
$$
fonctoriel en $M$, où les différentielles du complexe $\Hom$
sont les transposées des différentielles de
$I^\bullet(f_1^e, \ldots, f_r^e)$.
\end{lemma}

\begin{proof}
Rappelons que le complexe de {\v C}ech alterné est le sous-complexe
du complexe de {\v C}ech usuel formé des cochaînes alternées ; voir
Cohomologie, section \ref{cohomology-section-alternating-cech}.
Comme d'habitude, on considère une $p$-cochaîne de
$\check{\mathcal{C}}_{alt}^\bullet(\mathcal{U}, \mathcal{F})$
comme une fonction alternée $s$ sur $\{1, \ldots, r\}^{p + 1}$
dont la valeur $s_{i_0\ldots i_p}$ en $(i_0, \ldots, i_p)$ appartient à
$M_{f_{i_0}\ldots f_{i_p}} = \mathcal{F}(U_{i_0\ldots i_p})$.
D'autre part, une $p$-cochaîne $t$ de
$\Hom^\bullet(I^\bullet(f_1^e, \ldots, f_r^e), M)$
est une application $t : \wedge^{p + 1}(A^{\oplus r}) \to M$.
Notons $[i] \in A^{\oplus r}$ le $i$-ième vecteur de base et
posons
$$
[i_0, \ldots, i_p] = [i_0] \wedge \ldots \wedge [i_p]
\in \wedge^{p + 1}(A^{\oplus r})
$$
Pour $t$ comme ci-dessus, on pose
$$
\Psi(t)_{i_0 \ldots i_p} =
(-1)^p \frac{t([i_0, \ldots, i_p])}{f_{i_0}^e\ldots f_{i_p}^e}
$$
Il est clair que $\Psi(t)$ est une cochaîne alternée.
Cette règle est compatible aux morphismes de transition
du système, car le morphisme de transition
$$
I^\bullet(f_1^e, \ldots, f_r^e) \leftarrow
I^\bullet(f_1^{e + 1}, \ldots, f_r^{e + 1}),
$$
de (\ref{perfect-equation-system}) envoie $[i_0, \ldots, i_p]$
sur $f_{i_0}\ldots f_{i_p}[i_0, \ldots, i_p]$.
La description des localisés
$M_{f_{i_0} \ldots f_{i_p}}$ donnée dans
Algèbre, lemme \ref{algebra-lemma-localization-colimit},
montre que la règle $\Psi$ définit, après passage à la limite inductive, un isomorphisme des modules de cochaînes
en degré $p$. Pour achever la démonstration, il faut montrer que ce morphisme
est compatible aux différentielles. Pour $t$ comme ci-dessus, on calcule
\begin{align*}
d(\Psi(t))_{i_0 \ldots i_{p + 1}}
& =
\sum\nolimits_{j = 0}^{p + 1} (-1)^j
\Psi(t)_{i_0\ldots \hat i_j \ldots i_{p + 1}} \\
& =
(-1)^p
\sum\nolimits_{j = 0}^{p + 1} (-1)^j t([i_0 \ldots \hat i_j \ldots i_{p + 1}])
(f_{i_0} \ldots \hat f_{i_j} \ldots f_{i_p})^{-e}
\end{align*}
Rappelons que les différentielles de $I^\bullet(f_1^e, \ldots, f_r^e)$
sont les opposées de celles de $K^\bullet(f_1, \ldots, f_r)$.
Ainsi
\begin{align*}
\Psi(d(t))_{i_0 \ldots i_{p + 1}} & =
(-1)^{p + 1}
d(t)([i_0, \ldots, i_{p + 1}]) (f_{i_0} \ldots f_{i_{p + 1}})^{-e} \\
& =
(-1)^{p + 1}
t(d([i_0, \ldots, i_{p + 1}])) (f_{i_0} \ldots f_{i_{p + 1}})^{-e} \\
& =
(-1)^{p + 1}
t(-\sum\nolimits_{j = 0}^{p + 1}
(-1)^j f_{i_j}^e [i_0, \ldots, \hat i_j, \ldots i_{p + 1}])
(f_{i_0} \ldots f_{i_{p + 1}})^{-e} \\
& =
-(-1)^{p + 1}
\sum\nolimits_{j = 0}^{p + 1}
(-1)^j t([i_0, \ldots, \hat i_j, \ldots i_{p + 1}])
(f_{i_0} \ldots \hat f_{i_j} \ldots f_{i_p})^{-e}
\end{align*}
Les deux formules coïncident, ce qui achève la démonstration.
\end{proof}

\noindent
Soient un complexe borné $I^\bullet$ de $A$-modules et un
complexe de $A$-modules $M^\bullet$. On dispose du
complexe Hom correspondant $\Hom^\bullet(I^\bullet, M^\bullet)$. Son terme
en degré $n$ est
$$
\bigoplus\nolimits_{p + q = n} \Hom_A(I^{-q}, M^p)
$$
et sa différentielle est décrite dans Compléments sur l'algèbre, section
\ref{more-algebra-section-hom-complexes}.
Puisque le complexe $I^\bullet$ n'a qu'un nombre fini de termes non nuls, la
somme directe ci-dessus est finie.
Les conventions pour le complexe total associé au
complexe de {\v C}ech d'un complexe sont celles de
Cohomologie, section \ref{cohomology-section-cech-cohomology-of-complexes}.

\begin{lemma}
\label{perfect-lemma-alternating-cech-complex-complex}
Dans la situation \ref{perfect-situation-complex}, soit $M^\bullet$ un
complexe de $A$-modules, et
notons $\mathcal{F}^\bullet$ le complexe associé de
$\mathcal{O}_X$-modules. Alors
il existe un isomorphisme canonique de complexes
$$
\colim_e \Hom^\bullet(I^\bullet(f_1^e, \ldots, f_r^e), M^\bullet)
\longrightarrow
\text{Tot}(\check{\mathcal{C}}_{alt}^\bullet(\mathcal{U}, \mathcal{F}^\bullet))
$$
fonctoriel en $M^\bullet$.
\end{lemma}

\begin{proof}
Considérons le complexe double $F^{\bullet, \bullet}$ de termes
$F^{p, q} = \mathcal{C}_{alt}^p(\mathcal{U}, \mathcal{F}^q)$
étudié dans
Cohomologie, section \ref{cohomology-section-cech-cohomology-of-complexes}.
Considérons le complexe double $G^{\bullet, \bullet}$ de termes
$G^{p, q} = \colim_e \Hom_A(I^{-p}(f_1^e, \ldots, f_r^e), M^q)$ et
dont les différentielles sont données par fonctorialité (sans intervention
de signes). Les morphismes $\psi^{p, q} : G^{p, q} \to F^{p, q}$
construits dans la démonstration du lemme \ref{perfect-lemma-alternating-cech-complex}
sont des isomorphismes compatibles aux différentielles $d_1$ (d'après le lemme)
et $d_2$ (ce qui est clair). Cependant, les différentielles $d$ des
complexes figurant à gauche et à droite de la flèche de l'énoncé
ont des signes différents. Pour $g \in G^{p, q}$, on a en effet
$$
d(g) =  d_2(g) - (-1)^{p + q} d_1(g) 
$$
(voir Compléments sur l'algèbre, section \ref{more-algebra-section-hom-complexes}),
tandis que la différentielle, pour $f \in F^{p, q}$, est donnée par
$$
d(f) = d_1(f) + (-1)^p d_2(f)
$$
On peut donc corriger les signes en multipliant $\psi^{p, q}$ par
$(-1)^{pq + p(p - 1)/2}$.
\end{proof}

\begin{lemma}
\label{perfect-lemma-alternating-cech-complex-complex-computes-cohomology}
Dans la situation \ref{perfect-situation-complex}, soit $\mathcal{F}^\bullet$
un complexe de $\mathcal{O}_X$-modules quasi-cohérents. Alors
il existe un isomorphisme canonique
$$
\text{Tot}(\check{\mathcal{C}}_{alt}^\bullet(\mathcal{U}, \mathcal{F}^\bullet))
\longrightarrow
R\Gamma(U, \mathcal{F}^\bullet)
$$
dans $D(A)$, fonctoriel en $\mathcal{F}^\bullet$.
\end{lemma}

\begin{proof}
Soit $\mathcal{B}$ l'ensemble des ouverts affines de $U$. Puisque les groupes de
cohomologie supérieure d'un module quasi-cohérent sur un schéma affine sont nuls
(Cohomologie des schémas, lemme
\ref{coherent-lemma-quasi-coherent-affine-cohomology-zero}),
il s'agit d'un cas particulier de
Cohomologie, lemme \ref{cohomology-lemma-alternating-cech-complex-complex-ss}.
\end{proof}

\noindent
Dans la situation \ref{perfect-situation-complex}, notons $I_e$ l'objet de
$D(\mathcal{O}_X)$ correspondant au complexe de $A$-modules
$I^\bullet(f_1^e, \ldots, f_r^e)$ par l'équivalence du
lemme \ref{perfect-lemma-affine-compare-bounded}. Les morphismes
(\ref{perfect-equation-system}) donnent un système
$$
I_1 \leftarrow
I_2 \leftarrow
I_3 \leftarrow \ldots
$$
De plus, on dispose d'un système compatible de morphismes $I_e \to \mathcal{O}_X$
qui deviennent des isomorphismes après restriction à $U$. Ainsi, pour
tout objet $E$ de $D(\mathcal{O}_X)$, il existe un morphisme canonique
\begin{equation}
\label{perfect-equation-comparison}
\colim_e \Hom_{D(\mathcal{O}_X)}(I_e, E) \longrightarrow H^0(U, E)
\end{equation}
obtenu en envoyant un morphisme $I_e \to E$ sur sa restriction à $U$
et en utilisant l'égalité
$\Hom_{D(\mathcal{O}_U)}(\mathcal{O}_U, E|_U) = H^0(U, E)$.

\begin{proposition}
\label{perfect-proposition-represent-cohomology-class-on-open}
Dans la situation \ref{perfect-situation-complex}, pour tout objet $E$
de $D_\QCoh(\mathcal{O}_X)$, le morphisme
(\ref{perfect-equation-comparison}) est un isomorphisme.
\end{proposition}

\begin{proof}
D'après le lemme \ref{perfect-lemma-affine-compare-bounded}, on peut supposer que $E$
est donné par un complexe de faisceaux quasi-cohérents $\mathcal{F}^\bullet$.
Soit $M^\bullet = \Gamma(X, \mathcal{F}^\bullet)$ le complexe correspondant
de $A$-modules. D'après
les lemmes \ref{perfect-lemma-alternating-cech-complex-complex} et
\ref{perfect-lemma-alternating-cech-complex-complex-computes-cohomology},
on a des quasi-isomorphismes
$$
\colim_e \Hom^\bullet(I^\bullet(f_1^e, \ldots, f_r^e), M^\bullet)
\longrightarrow
\text{Tot}(\check{\mathcal{C}}_{alt}^\bullet(\mathcal{U}, \mathcal{F}^\bullet))
\longrightarrow
R\Gamma(U, \mathcal{F}^\bullet)
$$
D'après Compléments sur l'algèbre, lemme \ref{more-algebra-lemma-RHom-out-of-projective} et
équation (\ref{more-algebra-equation-h0-RHom}),
le $H^0$ du complexe
$\Hom^\bullet(I^\bullet(f_1^e, \ldots, f_r^e), M^\bullet)$
calcule le $\Hom$ dans $D(A)$. En prenant $H^0$ des deux côtés, on obtient donc
$$
\colim_e \Hom_{D(A)}(I^\bullet(f_1^e, \ldots, f_r^e), M^\bullet)
=
H^0(U, E)
$$
Puisque $\Hom_{D(A)}(I^\bullet(f_1^e, \ldots, f_r^e), M^\bullet) =
\Hom_{D(\mathcal{O}_X)}(I_e, E)$ d'après
le lemme \ref{perfect-lemma-affine-compare-bounded}, le résultat en découle.
\end{proof}

\noindent
Dans la situation \ref{perfect-situation-complex}, notons $K_e$ l'objet de
$D(\mathcal{O}_X)$ correspondant au complexe de $A$-modules
$K^\bullet(f_1^e, \ldots, f_r^e)$ par l'équivalence du
lemme \ref{perfect-lemma-affine-compare-bounded}. On a ainsi des triangles
distingués
$$
I_e \to \mathcal{O}_X \to K_e \to I_e[1]
$$
et un système
$$
K_1 \leftarrow
K_2 \leftarrow
K_3 \leftarrow \ldots
$$
compatible au système $(I_e)$.
De plus, on dispose d'un système compatible de morphismes
$$
K_e \to H^0(K_e) = \mathcal{O}_X/(f_1^e, \ldots, f_r^e)
$$

\begin{lemma}
\label{perfect-lemma-represent-cohomology-class-on-closed}
Dans la situation \ref{perfect-situation-complex}, soit $E$ un objet de
$D_\QCoh(\mathcal{O}_X)$.
Supposons que $H^i(E)|_U = 0$ pour $i = - r + 1, \ldots, 0$.
Alors, étant donné $s \in H^0(X, E)$, il existe un $e \geq 0$ et
un morphisme $K_e \to E$ tels que $s$ appartienne à l'image de
$H^0(X, K_e) \to H^0(X, E)$.
\end{lemma}

\begin{proof}
Puisque $U$ est recouvert par $r$ ouverts affines, on a $H^j(U, \mathcal{F}) = 0$
pour $j \geq r$ et tout module quasi-cohérent
(Cohomologie des schémas, lemme \ref{coherent-lemma-vanishing-nr-affines}).
D'après le lemme \ref{perfect-lemma-application-nice-K-injective}, $H^0(U, E)$
est égal à $H^0(U, \tau_{\geq -r + 1}E)$. On dispose d'une
suite spectrale
$$
H^j(U, H^i(\tau_{\geq -r + 1}E)) \Rightarrow H^{i + j}(U, \tau_{\geq -N}E)
$$
voir Catégories dérivées, lemme \ref{derived-lemma-two-ss-complex-functor}.
Ainsi $H^0(U, E) = 0$ par l'hypothèse d'annulation des faisceaux de cohomologie de $E$.
On en déduit que $s|_U = 0$.
Considérons $s$ comme un morphisme $\mathcal{O}_X \to E$ dans $D(\mathcal{O}_X)$.
D'après la proposition \ref{perfect-proposition-represent-cohomology-class-on-open},
le composé $I_e \to \mathcal{O}_X \to E$ est nul pour un certain $e$.
Le triangle distingué $I_e \to \mathcal{O}_X \to K_e \to I_e[1]$
fournit alors un morphisme $K_e \to E$ tel que $s$ soit le composé
$\mathcal{O}_X \to K_e \to E$.
\end{proof}

\section{Complexes pseudo-cohérents et parfaits}
\label{perfect-section-spell-out}

\noindent
Dans cette section, nous établissons le lien entre les notions
générales définies dans
Cohomologie, sections \ref{cohomology-section-strictly-perfect},
\ref{cohomology-section-pseudo-coherent},
\ref{cohomology-section-tor} et
\ref{cohomology-section-perfect},
et les notions correspondantes pour les complexes de modules, dans
Compléments sur l'algèbre, sections
\ref{more-algebra-section-pseudo-coherent},
\ref{more-algebra-section-tor} et
\ref{more-algebra-section-perfect}.

\begin{lemma}
\label{perfect-lemma-pseudo-coherent}
Soit $X$ un schéma. Si $E$ est un objet $m$-pseudo-cohérent
de $D(\mathcal{O}_X)$, alors $H^i(E)$ est un
$\mathcal{O}_X$-module quasi-cohérent pour $i > m$, et $H^m(E)$ est un quotient
d'un $\mathcal{O}_X$-module quasi-cohérent.
Si $E$ est pseudo-cohérent, alors $E$ est un objet de
$D_\QCoh(\mathcal{O}_X)$.
\end{lemma}

\begin{proof}
Localement sur $X$, il existe un complexe strictement parfait $\mathcal{E}^\bullet$
tel que $H^i(E)$ soit isomorphe à $H^i(\mathcal{E}^\bullet)$ pour $i > m$
et que $H^m(E)$ soit un quotient de $H^m(\mathcal{E}^\bullet)$. Les faisceaux
$\mathcal{E}^i$ sont des facteurs directs de modules libres de rang fini,
donc sont quasi-cohérents. Le lemme en résulte.
\end{proof}

\begin{lemma}
\label{perfect-lemma-pseudo-coherent-affine}
Soit $X = \Spec(A)$ un schéma affine. Soit $M^\bullet$ un
complexe de $A$-modules et soit $E$ l'objet correspondant
de $D(\mathcal{O}_X)$. Alors $E$ est $m$-pseudo-cohérent
(resp.\ pseudo-cohérent) en tant qu'objet de $D(\mathcal{O}_X)$
si et seulement si $M^\bullet$ est $m$-pseudo-cohérent (resp.\ pseudo-cohérent)
en tant que complexe de $A$-modules.
\end{lemma}

\begin{proof}
Il résulte immédiatement des définitions que, si $M^\bullet$ est
$m$-pseudo-cohérent, il en est de même de $E$. Pour établir la réciproque, supposons
$E$ $m$-pseudo-cohérent. Puisque $X = \Spec(A)$ est quasi-compact et que
les ouverts principaux forment une base de sa topologie, on peut trouver un
recouvrement par des ouverts principaux $X = D(f_1) \cup \ldots \cup D(f_n)$, des complexes strictement
parfaits $\mathcal{E}_i^\bullet$ sur $D(f_i)$ et des
morphismes $\alpha_i : \mathcal{E}_i^\bullet \to E|_{U_i}$ induisant
des isomorphismes sur $H^j$ pour $j > m$ et des surjections sur $H^m$.
D'après Cohomologie, lemme \ref{cohomology-lemma-local-actual},
quitte à raffiner le recouvrement ouvert,
on peut supposer que $\alpha_i$ est donné par un morphisme de complexes
$\mathcal{E}_i^\bullet \to \widetilde{M^\bullet}|_{U_i}$
pour tout $i$. D'après Modules, lemme
\ref{modules-lemma-direct-summand-of-locally-free-is-locally-free},
les termes $\mathcal{E}_i^n$ sont des modules localement libres de rang fini.
Ainsi, quitte à raffiner le recouvrement ouvert, on peut supposer chaque
$\mathcal{E}_i^n$ libre de rang fini sur $\mathcal{O}_{U_i}$.
Il résulte de la définition que $M^\bullet_{f_i}$ est
un complexe $m$-pseudo-cohérent de $A_{f_i}$-modules.
On conclut en appliquant
Compléments sur l'algèbre, lemme \ref{more-algebra-lemma-glue-pseudo-coherent}.

\medskip\noindent
Le cas « pseudo-cohérent » résulte du fait que $E$ est
pseudo-cohérent si et seulement si $E$ est $m$-pseudo-cohérent pour
tout $m$ (par définition), et qu'il en est de même de $M^\bullet$
d'après Compléments sur l'algèbre, lemme \ref{more-algebra-lemma-pseudo-coherent}.
\end{proof}

\begin{lemma}
\label{perfect-lemma-identify-pseudo-coherent-noetherian}
Soit $X$ un schéma noethérien. Soit $E$ un objet de
$D_\QCoh(\mathcal{O}_X)$. Pour $m \in \mathbf{Z}$, les
conditions suivantes sont équivalentes :
\begin{enumerate}
\item $H^i(E)$ est cohérent pour $i \geq m$ et nul pour $i \gg 0$,
\item $E$ est $m$-pseudo-cohérent.
\end{enumerate}
En particulier, $E$ est pseudo-cohérent si et seulement si $E$ est un objet
de $D^-_{\textit{Coh}}(\mathcal{O}_X)$.
\end{lemma}

\begin{proof}
Puisque $X$ est quasi-compact, dans chacun des cas (1) et (2), l'objet $E$
est borné supérieurement. La question est donc locale sur $X$, et l'on peut supposer
$X$ affine. Écrivons $X = \Spec(A)$ pour un anneau noethérien $A$.
Dans ce cas, $E$ correspond à un complexe de $A$-modules $M^\bullet$
d'après le lemme \ref{perfect-lemma-affine-compare-bounded}. D'après
le lemme \ref{perfect-lemma-pseudo-coherent-affine},
$E$ est $m$-pseudo-cohérent si et seulement si $M^\bullet$
est $m$-pseudo-cohérent. D'autre part, $H^i(E)$ est cohérent
si et seulement si $H^i(M^\bullet)$ est un $A$-module de type fini
(Propriétés, lemme \ref{properties-lemma-finite-type-module}).
Le résultat découle donc de Compléments sur l'algèbre, lemme
\ref{more-algebra-lemma-Noetherian-pseudo-coherent}.
\end{proof}

\begin{lemma}
\label{perfect-lemma-tor-dimension-affine}
Soit $X = \Spec(A)$ un schéma affine. Soit $M^\bullet$ un
complexe de $A$-modules et soit $E$ l'objet correspondant
de $D(\mathcal{O}_X)$. Alors
\begin{enumerate}
\item $E$ a une amplitude de Tor contenue dans $[a, b]$ si et seulement si $M^\bullet$
a une amplitude de Tor contenue dans $[a, b]$.
\item $E$ est de dimension de Tor finie si et seulement si $M^\bullet$
est de dimension de Tor finie.
\end{enumerate}
\end{lemma}

\begin{proof}
L'assertion (2) résulte immédiatement de (1). Pour démontrer (1), nous
utiliserons l'équivalence $D(A) = D_\QCoh(X)$ du
lemme \ref{perfect-lemma-affine-compare-bounded}
sans autre mention.
Supposons que $M^\bullet$ ait une amplitude de Tor contenue dans $[a, b]$. Alors $K^\bullet$
est isomorphe dans $D(A)$ à un complexe $K^\bullet$ de $A$-modules plats
tel que $K^i = 0$ pour $i \not \in [a, b]$ ; voir
Compléments sur l'algèbre, lemme \ref{more-algebra-lemma-tor-amplitude}.
Alors $E$ est isomorphe à $\widetilde{K^\bullet}$. Puisque chaque
$\widetilde{K^i}$ est un $\mathcal{O}_X$-module plat,
$E$ a une amplitude de Tor contenue dans $[a, b]$ d'après
Cohomologie, lemme \ref{cohomology-lemma-tor-amplitude}.

\medskip\noindent
Supposons que $E$ ait une amplitude de Tor contenue dans $[a, b]$. Alors $E$ est borné,
donc $M^\bullet$ appartient à $K^-(A)$. On peut ainsi remplacer $M^\bullet$
par un complexe borné supérieurement de $A$-modules. On peut même choisir
une résolution projective et supposer que $M^\bullet$ est un complexe borné supérieurement
de $A$-modules libres. Pour tout $A$-module $N$, on a alors
$$
E \otimes_{\mathcal{O}_X}^\mathbf{L} \widetilde{N}
\cong
\widetilde{M^\bullet} \otimes_{\mathcal{O}_X}^\mathbf{L} \widetilde{N}
\cong
\widetilde{M^\bullet \otimes_A N}
$$
dans $D(\mathcal{O}_X)$. L'annulation des faisceaux de cohomologie du
membre de gauche entraîne donc que $M^\bullet$ a une amplitude de Tor contenue dans $[a, b]$.
\end{proof}

\begin{lemma}
\label{perfect-lemma-tor-dimension-rel-affine}
Soit $f : X \to S$ un morphisme de schémas affines correspondant
au morphisme d'anneaux $R \to A$. Soit $M^\bullet$ un
complexe de $A$-modules et soit $E$ l'objet correspondant
de $D(\mathcal{O}_X)$. Alors
\begin{enumerate}
\item $E$, en tant qu'objet de $D(f^{-1}\mathcal{O}_S)$, a une amplitude de Tor contenue dans
$[a, b]$ si et seulement si $M^\bullet$ a une amplitude de Tor contenue dans $[a, b]$
en tant qu'objet de $D(R)$.
\item $E$ est localement de dimension de Tor finie en tant qu'objet de
$D(f^{-1}\mathcal{O}_S)$ si et seulement si $M^\bullet$
est de dimension de Tor finie en tant qu'objet de $D(R)$.
\end{enumerate}
\end{lemma}

\begin{proof}
Considérons un idéal premier $\mathfrak q \subset A$ au-dessus de $\mathfrak p \subset R$.
Soient $x \in X$ et $s = f(x) \in S$ les points correspondants.
Alors $(f^{-1}\mathcal{O}_S)_x = \mathcal{O}_{S, s} = R_\mathfrak p$
et $E_x = M^\bullet_\mathfrak q$. Compte tenu de ces identifications, on établit
l'équivalence comme suit.

\medskip\noindent
Si $M^\bullet$ a une amplitude de Tor contenue dans $[a, b]$ en tant que complexe de $R$-modules,
il en est de même du localisé de $M^\bullet$ en tout idéal premier de $A$.
On déduit alors de
Cohomologie, lemme \ref{cohomology-lemma-tor-amplitude-stalk},
que $E$ a une amplitude de Tor contenue dans $[a, b]$ en tant que complexe de faisceaux
de $f^{-1}\mathcal{O}_S$-modules.
Réciproquement, supposons que $E$ ait une amplitude de Tor contenue dans $[a, b]$
en tant qu'objet de $D(f^{-1}\mathcal{O}_S)$.
On en déduit, à l'aide du lemme qui vient d'être cité, que
$M^\bullet_\mathfrak q$ a une amplitude de Tor contenue dans $[a, b]$
en tant que complexe de $R_\mathfrak p$-modules pour tout
idéal premier $\mathfrak q \subset A$ au-dessus de $\mathfrak p \subset R$.
D'après Compléments sur l'algèbre, lemme \ref{more-algebra-lemma-tor-amplitude-localization},
il en résulte que $M^\bullet$ a une amplitude de Tor contenue dans $[a, b]$
en tant que complexe de $R$-modules.
Cela achève la démonstration de (1).

\medskip\noindent
Puisque $X$ est quasi-compact, si $E$ est localement de dimension de Tor finie
en tant que complexe de $f^{-1}\mathcal{O}_S$-modules, alors $E$
a en fait une amplitude de Tor contenue dans $[a, b]$ pour certains $a, b$, en tant que complexe de
$f^{-1}\mathcal{O}_S$-modules. Ainsi (2) résulte de (1).
\end{proof}

\begin{lemma}
\label{perfect-lemma-tor-qc-qs}
Soit $X$ un schéma quasi-séparé. Soit $E$ un objet
de $D_\QCoh(\mathcal{O}_X)$. Soient $a \leq b$. Les
conditions suivantes sont équivalentes :
\begin{enumerate}
\item $E$ a une amplitude de Tor contenue dans $[a, b]$,
\item pour tout $\mathcal{F}$ dans $\QCoh(\mathcal{O}_X)$,
on a $H^i(E \otimes_{\mathcal{O}_X}^\mathbf{L} \mathcal{F}) = 0$
pour $i \not \in [a, b]$.
\end{enumerate}
\end{lemma}

\begin{proof}
Il est clair que (1) entraîne (2). Supposons (2). Soit $U \subset X$
un ouvert affine. Puisque $X$ est quasi-séparé, le morphisme $j : U \to X$
est quasi-compact et séparé ; ainsi $j_*$ transforme les modules quasi-cohérents
en modules quasi-cohérents
(Schémas, lemme \ref{schemes-lemma-push-forward-quasi-coherent}).
Le foncteur
$\QCoh(\mathcal{O}_X) \to \QCoh(\mathcal{O}_U)$
est donc essentiellement surjectif. Il en résulte que la condition (2)
entraîne l'annulation de
$H^i(E|_U \otimes_{\mathcal{O}_U}^\mathbf{L} \mathcal{G})$
pour $i \not \in [a, b]$, pour tout $\mathcal{O}_U$-module quasi-cohérent
$\mathcal{G}$. Écrivons $U = \Spec(A)$, et soit $M^\bullet$ le
complexe de $A$-modules correspondant à $E|_U$ d'après
le lemme \ref{perfect-lemma-affine-compare-bounded}.
Nous venons de montrer que $M^\bullet \otimes_A^\mathbf{L} N$
a des groupes de cohomologie nuls en dehors de l'intervalle $[a, b]$ ;
autrement dit, $M^\bullet$ a une amplitude de Tor contenue dans $[a, b]$.
D'après le lemme \ref{perfect-lemma-tor-dimension-affine},
on en déduit que $E|_U$ a une amplitude de Tor contenue dans $[a, b]$.
Cela démontre le lemme.
\end{proof}

\begin{lemma}
\label{perfect-lemma-perfect-affine}
Soit $X = \Spec(A)$ un schéma affine. Soit $M^\bullet$ un
complexe de $A$-modules et soit $E$ l'objet correspondant
de $D(\mathcal{O}_X)$. Alors $E$ est un objet parfait de $D(\mathcal{O}_X)$
si et seulement si $M^\bullet$ est parfait en tant qu'objet de $D(A)$.
\end{lemma}

\begin{proof}
C'est une conséquence formelle des
lemmes \ref{perfect-lemma-pseudo-coherent-affine} et
\ref{perfect-lemma-tor-dimension-affine},
de Cohomologie, lemme \ref{cohomology-lemma-perfect}, et de
Compléments sur l'algèbre, lemme \ref{more-algebra-lemma-perfect}.
\end{proof}

\noindent
Notre description des complexes pseudo-cohérents
sur les schémas permet de montrer que certains Hom internes
sont quasi-cohérents.

\begin{lemma}
\label{perfect-lemma-quasi-coherence-internal-hom}
Soit $X$ un schéma.
\begin{enumerate}
\item Si $L$ appartient à $D^+_\QCoh(\mathcal{O}_X)$ et si
$K$ dans $D(\mathcal{O}_X)$ est pseudo-cohérent, alors
$R\SheafHom(K, L)$ appartient à $D_\QCoh(\mathcal{O}_X)$
et est localement borné inférieurement.
\item Si $L$ appartient à $D_\QCoh(\mathcal{O}_X)$ et si
$K$ dans $D(\mathcal{O}_X)$ est parfait, alors
$R\SheafHom(K, L)$ appartient à $D_\QCoh(\mathcal{O}_X)$.
\item Si $X = \Spec(A)$ est affine et si $K, L \in D(A)$, alors
$$
R\SheafHom(\widetilde{K}, \widetilde{L}) = \widetilde{R\Hom_A(K, L)}
$$
dans les deux cas suivants :
\begin{enumerate}
\item $K$ est pseudo-cohérent et $L$ est borné inférieurement,
\item $K$ est parfait et $L$ est quelconque.
\end{enumerate}
\item Si $X = \Spec(A)$ et si $K, L$ appartiennent à $D(A)$, le $n$-ième
faisceau de cohomologie de $R\SheafHom(\widetilde{K}, \widetilde{L})$
est le faisceau associé au préfaisceau
$$
X \supset D(f) \longmapsto \Ext^n_{A_f}(K \otimes_A A_f, L \otimes_A A_f)
$$
pour $f \in A$.
\end{enumerate}
\end{lemma}

\begin{proof}
La construction du Hom interne dans la catégorie dérivée de
$\mathcal{O}_X$ commute à la localisation (voir
Cohomologie, section \ref{cohomology-section-internal-hom}).
Pour démontrer (1) et (2), on peut donc remplacer $X$ par un ouvert affine.
D'après les lemmes \ref{perfect-lemma-affine-compare-bounded},
\ref{perfect-lemma-pseudo-coherent-affine} et
\ref{perfect-lemma-perfect-affine},
il suffit alors de démontrer (3).

\medskip\noindent
L'assertion (3) résulte du calcul du
Hom interne de Cohomologie, lemme
\ref{cohomology-lemma-Rhom-complex-of-direct-summands-finite-free},
en représentant $K$ par un complexe borné supérieurement (resp.\ borné) de
$A$-modules projectifs de type fini et $L$ par un complexe borné inférieurement
(resp.\ quelconque) de $A$-modules.

\medskip\noindent
Pour démontrer (4), rappelons que, sur tout espace annelé, le $n$-ième faisceau de cohomologie de
$R\SheafHom(A, B)$ est le faisceau associé au préfaisceau
$$
U \mapsto \Hom_{D(U)}(A|_U, B|_U[n]) =
\Ext^n_{D(\mathcal{O}_U)}(A|_U, B|_U)
$$
Voir Cohomologie, section \ref{cohomology-section-internal-hom}.
D'autre part, la restriction de $\widetilde{K}$ à un ouvert principal
$D(f)$ est l'image de $K \otimes_A A_f$, et de même pour $L$.
L'assertion (4) résulte donc de l'équivalence de catégories du
lemme \ref{perfect-lemma-affine-compare-bounded}.
\end{proof}

\begin{lemma}
\label{perfect-lemma-internal-hom-evaluate-tensor-isomorphism}
Soit $X$ un schéma. Soient $K, L, M$ des objets de $D_\QCoh(\mathcal{O}_X)$.
Le morphisme
$$
K \otimes_{\mathcal{O}_X}^\mathbf{L} R\SheafHom(M, L)
\longrightarrow
R\SheafHom(M, K \otimes_{\mathcal{O}_X}^\mathbf{L} L)
$$
de Cohomologie, lemme \ref{cohomology-lemma-internal-hom-diagonal-better},
est un isomorphisme dans les cas suivants :
\begin{enumerate}
\item $M$ est parfait, ou
\item $K$ est parfait, ou
\item $M$ est pseudo-cohérent, $L \in D^+(\mathcal{O}_X)$, et $K$ est de dimension
de Tor finie.
\end{enumerate}
\end{lemma}

\begin{proof}
Le lemme \ref{perfect-lemma-quasi-coherence-internal-hom}
ramène les cas (1) et (3) au cas affine, traité dans
Compléments sur l'algèbre, lemme
\ref{more-algebra-lemma-internal-hom-evaluate-tensor-isomorphism}.
(Il faut aussi utiliser les lemmes \ref{perfect-lemma-pseudo-coherent-affine},
\ref{perfect-lemma-perfect-affine} et \ref{perfect-lemma-tor-dimension-affine}
pour passer à l'énoncé algébrique.)
Si $K$ est parfait, sans autre hypothèse, on
ne sait pas si les deux termes de la flèche appartiennent à $D_\QCoh(\mathcal{O}_X)$,
mais le résultat reste vrai, car on peut travailler localement et se ramener
au cas où $K$ est un complexe borné de modules libres de rang fini,
auquel cas il est clair.
\end{proof}





\section{Catégorie dérivée des modules cohérents}
\label{perfect-section-derived-coherent}

\noindent
Soit $X$ un schéma localement noethérien. Dans ce cas, la catégorie
$\textit{Coh}(\mathcal{O}_X) \subset \textit{Mod}(\mathcal{O}_X)$
des $\mathcal{O}_X$-modules cohérents est une sous-catégorie de Serre faible ; voir
Homologie, section \ref{homology-section-serre-subcategories}
et
Cohomologie des schémas, lemme \ref{coherent-lemma-coherent-abelian-Noetherian}.
Notons
$$
D_{\textit{Coh}}(\mathcal{O}_X) \subset D(\mathcal{O}_X)
$$
la sous-catégorie des complexes dont les faisceaux de cohomologie sont cohérents ; voir
Catégories dérivées, section \ref{derived-section-triangulated-sub}.
On obtient ainsi un foncteur canonique
\begin{equation}
\label{perfect-equation-compare-coherent}
D(\textit{Coh}(\mathcal{O}_X))
\longrightarrow
D_{\textit{Coh}}(\mathcal{O}_X)
\end{equation}
voir Catégories dérivées, équation (\ref{derived-equation-compare}).

\begin{lemma}
\label{perfect-lemma-coh-to-qcoh}
Soit $X$ un schéma noethérien. Alors le foncteur
$$
D^-(\textit{Coh}(\mathcal{O}_X))
\longrightarrow
D^-_{\textit{Coh}(\mathcal{O}_X)}(\QCoh(\mathcal{O}_X))
$$
est une équivalence.
\end{lemma}

\begin{proof}
Observons que $\textit{Coh}(\mathcal{O}_X) \subset \QCoh(\mathcal{O}_X)$
est une sous-catégorie de Serre ; voir
Homologie, définition \ref{homology-definition-serre-subcategory} et
lemme \ref{homology-lemma-characterize-serre-subcategory}, ainsi que
Cohomologie des schémas, lemmes
\ref{coherent-lemma-coherent-abelian-Noetherian} et
\ref{coherent-lemma-coherent-Noetherian-quasi-coherent-sub-quotient}.
D'autre part, si $\mathcal{G} \to \mathcal{F}$ est une surjection
d'un $\mathcal{O}_X$-module quasi-cohérent sur un
$\mathcal{O}_X$-module cohérent, il existe un sous-module cohérent
$\mathcal{G}' \subset \mathcal{G}$ qui s'envoie surjectivement sur $\mathcal{F}$.
En effet, on peut écrire $\mathcal{G}$ comme réunion filtrante de ses sous-modules
cohérents d'après
Propriétés, lemme \ref{properties-lemma-quasi-coherent-colimit-finite-type},
et l'un d'entre eux convient alors.
Le lemme résulte donc de
Catégories dérivées, lemme \ref{derived-lemma-fully-faithful-embedding}.
\end{proof}

\begin{proposition}
\label{perfect-proposition-DCoh}
Soit $X$ un schéma noethérien. Alors les foncteurs
$$
D^-(\textit{Coh}(\mathcal{O}_X))
\longrightarrow
D^-_{\textit{Coh}}(\mathcal{O}_X)
\quad\text{et}\quad
D^b(\textit{Coh}(\mathcal{O}_X))
\longrightarrow
D^b_{\textit{Coh}}(\mathcal{O}_X)
$$
sont des équivalences.
\end{proposition}

\begin{proof}
Considérons le diagramme commutatif
$$
\xymatrix{
D^-(\textit{Coh}(\mathcal{O}_X)) \ar[r] \ar[d] &
D^-_{\textit{Coh}}(\mathcal{O}_X) \ar[d] \\
D^-(\QCoh(\mathcal{O}_X)) \ar[r] &
D^-_\QCoh(\mathcal{O}_X)
}
$$
D'après le lemme \ref{perfect-lemma-coh-to-qcoh}, la flèche verticale de gauche est pleinement fidèle.
D'après la proposition \ref{perfect-proposition-Noetherian}, la flèche inférieure est une équivalence.
Par construction, la flèche verticale de droite est pleinement fidèle.
On en déduit que la flèche horizontale supérieure est pleinement fidèle.
Si $K$ est un objet de $D^-_{\textit{Coh}}(\mathcal{O}_X)$,
l'objet $K'$ de $D^-(\QCoh(\mathcal{O}_X))$ qui lui correspond
par la proposition \ref{perfect-proposition-Noetherian} a des
faisceaux de cohomologie cohérents. Ainsi $K'$ appartient à l'image essentielle
de la flèche verticale de gauche, d'après le lemme \ref{perfect-lemma-coh-to-qcoh},
et la flèche horizontale supérieure est donc essentiellement surjective.
Cela achève la démonstration dans le cas borné supérieurement. Le cas
borné en découle immédiatement.
\end{proof}

\begin{lemma}
\label{perfect-lemma-direct-image-coherent}
Soit $S$ un schéma noethérien. Soit $f : X \to S$ un morphisme de schémas
localement de type fini. Soit $E$ un objet de
$D^b_{\textit{Coh}}(\mathcal{O}_X)$ tel que le support de $H^i(E)$
soit propre sur $S$ pour tout $i$.
Alors $Rf_*E$ est un objet de $D^b_{\textit{Coh}}(\mathcal{O}_S)$.
\end{lemma}

\begin{proof}
Considérons la suite spectrale
$$
R^pf_*H^q(E) \Rightarrow R^{p + q}f_*E
$$
voir Catégories dérivées, lemme \ref{derived-lemma-two-ss-complex-functor}.
Par hypothèse et d'après
Cohomologie des schémas, lemme
\ref{coherent-lemma-support-proper-over-base-pushforward},
les faisceaux $R^pf_*H^q(E)$ sont cohérents. Ainsi
$R^{p + q}f_*E$ est cohérent, c'est-à-dire que $Rf_*E \in D_{\textit{Coh}}(\mathcal{O}_S)$.
La bornitude inférieure est immédiate. La bornitude supérieure
résulte de
Cohomologie des schémas, lemme
\ref{coherent-lemma-quasi-coherence-higher-direct-images},
ou du
lemme \ref{perfect-lemma-quasi-coherence-direct-image}.
\end{proof}

\begin{lemma}
\label{perfect-lemma-direct-image-coherent-bdd-below}
Soit $S$ un schéma noethérien. Soit $f : X \to S$ un morphisme de schémas
localement de type fini. Soit $E$ un objet de
$D^+_{\textit{Coh}}(\mathcal{O}_X)$ tel que le support de $H^i(E)$
soit propre sur $S$ pour tout $i$.
Alors $Rf_*E$ est un objet de $D^+_{\textit{Coh}}(\mathcal{O}_S)$.
\end{lemma}

\begin{proof}
La démonstration est la même que celle du
lemme \ref{perfect-lemma-direct-image-coherent}.
On peut aussi le déduire du
lemme \ref{perfect-lemma-direct-image-coherent}
en appliquant le foncteur exact $Rf_*$ aux
triangles distingués
$\tau_{\leq a}E \to E \to \tau_{\geq a + 1}E \to \tau_{\leq a}E[1]$.
\end{proof}

\begin{lemma}
\label{perfect-lemma-coherent-internal-hom}
Soit $X$ un schéma localement noethérien. Si $L$ appartient à
$D^+_{\textit{Coh}}(\mathcal{O}_X)$ et si $K$ appartient à
$D^-_{\textit{Coh}}(\mathcal{O}_X)$, alors
$R\SheafHom(K, L)$ appartient à $D^+_{\textit{Coh}}(\mathcal{O}_X)$.
\end{lemma}

\begin{proof}
Il suffit de le démontrer lorsque $X$ est le spectre
d'un anneau noethérien $A$.
D'après le lemme \ref{perfect-lemma-identify-pseudo-coherent-noetherian},
$K$ est pseudo-cohérent.
On peut alors utiliser le lemme \ref{perfect-lemma-quasi-coherence-internal-hom}
pour ramener la question à l'énoncé algébrique suivant :
pour $L \in D^+_{\textit{Coh}}(A)$ et $K$ dans $D^-_{\textit{Coh}}(A)$,
$R\Hom_A(K, L)$ appartient à $D^+_{\textit{Coh}}(A)$.
Puisque $L$ est borné inférieurement et $K$ borné supérieurement, on dispose d'une
suite spectrale convergente
$$
\Ext^p_A(K, H^q(L)) \Rightarrow \text{Ext}^{p + q}_A(K, L)
$$
ainsi que de suites spectrales convergentes
$$
\Ext^i_A(H^{-j}(K), H^q(L)) \Rightarrow \text{Ext}^{i + j}_A(K, H^q(L))
$$
Voir Injectifs, remarques \ref{injectives-remark-spectral-sequences-ext}
et \ref{injectives-remark-spectral-sequences-ext-variant}.
Cela achève la démonstration, car les modules $\Ext^p_A(M, N)$
sont de type fini pour des $A$-modules de type fini $M$, $N$, d'après
Algèbre, lemme \ref{algebra-lemma-ext-noetherian}.
\end{proof}

\begin{lemma}
\label{perfect-lemma-perfect-on-noetherian}
Soit $X$ un schéma noethérien. Soit $E$ dans $D(\mathcal{O}_X)$ un objet parfait.
Alors
\begin{enumerate}
\item $E$ appartient à $D^b_{\textit{Coh}}(\mathcal{O}_X)$,
\item si $L$ appartient à $D_{\textit{Coh}}(\mathcal{O}_X)$, alors
$E \otimes_{\mathcal{O}_X}^\mathbf{L} L$ et
$R\SheafHom_{\mathcal{O}_X}(E, L)$ appartiennent à
$D_{\textit{Coh}}(\mathcal{O}_X)$,
\item si $L$ appartient à $D^b_{\textit{Coh}}(\mathcal{O}_X)$, alors
$E \otimes_{\mathcal{O}_X}^\mathbf{L} L$ et
$R\SheafHom_{\mathcal{O}_X}(E, L)$ appartiennent à
$D^b_{\textit{Coh}}(\mathcal{O}_X)$,
\item si $L$ appartient à $D^+_{\textit{Coh}}(\mathcal{O}_X)$, alors
$E \otimes_{\mathcal{O}_X}^\mathbf{L} L$ et
$R\SheafHom_{\mathcal{O}_X}(E, L)$ appartiennent à
$D^+_{\textit{Coh}}(\mathcal{O}_X)$,
\item si $L$ appartient à $D^-_{\textit{Coh}}(\mathcal{O}_X)$, alors
$E \otimes_{\mathcal{O}_X}^\mathbf{L} L$ et
$R\SheafHom_{\mathcal{O}_X}(E, L)$ appartiennent à
$D^-_{\textit{Coh}}(\mathcal{O}_X)$.
\end{enumerate}
\end{lemma}

\begin{proof}
Puisque $X$ est quasi-compact, chacune de ces assertions se vérifie
sur les membres d'un recouvrement ouvert quelconque de $X$.
On peut donc supposer $E$ représenté par
un complexe borné $\mathcal{E}^\bullet$ de modules libres de rang fini ; voir
Cohomologie, lemme \ref{cohomology-lemma-perfect-on-locally-ringed}.
Dans ce cas, chacune des assertions est claire, car
$R\SheafHom_{\mathcal{O}_X}(E, L)$
et $E \otimes_{\mathcal{O}_X}^\mathbf{L} L$ se calculent au
niveau des complexes à l'aide de $\mathcal{E}^\bullet$ ; voir
Cohomologie, lemmes \ref{cohomology-lemma-Rhom-strictly-perfect} et
\ref{cohomology-lemma-bounded-flat-K-flat}. Nous omettons quelques détails.
\end{proof}

\begin{lemma}
\label{perfect-lemma-ext-finite}
Soit $A$ un anneau noethérien. Soit $X$ un schéma propre sur $A$.
Pour $L$ dans
$D^+_{\textit{Coh}}(\mathcal{O}_X)$ et $K$ dans
$D^-_{\textit{Coh}}(\mathcal{O}_X)$, les $A$-modules
$\Ext_{\mathcal{O}_X}^n(K, L)$ sont de type fini.
\end{lemma}

\begin{proof}
Rappelons que
$$
\Ext_{\mathcal{O}_X}^n(K, L) =
H^n(X, R\SheafHom_{\mathcal{O}_X}(K, L)) =
H^n(\Spec(A), Rf_*R\SheafHom_{\mathcal{O}_X}(K, L))
$$
voir Cohomologie, lemme \ref{cohomology-lemma-section-RHom-over-U}
et Cohomologie, section \ref{cohomology-section-Leray}.
Le résultat découle donc des
lemmes \ref{perfect-lemma-coherent-internal-hom} et
\ref{perfect-lemma-direct-image-coherent-bdd-below}.
\end{proof}

\begin{lemma}
\label{perfect-lemma-perfect-on-regular}
Soit $X$ un schéma régulier. Alors
tout objet de $D^b_{\textit{Coh}}(\mathcal{O}_X)$ est parfait.
Si $X$ est quasi-compact, c'est-à-dire noethérien et régulier,
alors, réciproquement, tout objet parfait de
$D(\mathcal{O}_X)$ appartient à $D^b_{\textit{Coh}}(\mathcal{O}_X)$.
\end{lemma}

\begin{proof}
Soit $K$ un objet de $D^b_{\textit{Coh}}(\mathcal{O}_X)$.
Pour vérifier que $K$ est parfait, on peut travailler sur les ouverts affines de $X$
(voir Cohomologie, section \ref{cohomology-section-perfect}).
Alors $K$ est parfait d'après le lemme \ref{perfect-lemma-perfect-affine} et
Compléments sur l'algèbre, lemme \ref{more-algebra-lemma-regular-perfect}.
La réciproque est le lemme \ref{perfect-lemma-perfect-on-noetherian}.
\end{proof}



\section{Descente des propriétés de finitude des complexes}
\label{perfect-section-descent-finiteness}

\noindent
Cette section est l'analogue de
Descente, section \ref{descent-section-descent-finiteness},
pour les objets de la catégorie dérivée d'un schéma.
Le résultat le plus simple de ce type est sans doute le suivant.

\begin{lemma}
\label{perfect-lemma-tor-amplitude-descends}
Soit $f : X \to Y$ un morphisme surjectif et plat de schémas
(ou, plus généralement, d'espaces localement annelés).
Soit $E \in D(\mathcal{O}_Y)$. Soient $a, b \in \mathbf{Z}$.
Alors $E$ a une amplitude de Tor contenue dans $[a, b]$ si et seulement si
$Lf^*E$ a une amplitude de Tor contenue dans $[a, b]$.
\end{lemma}

\begin{proof}
L'image réciproque préserve toujours l'amplitude de Tor ; voir
Cohomologie, lemme \ref{cohomology-lemma-tor-amplitude-pullback}.
On peut vérifier sur les fibres que l'amplitude de Tor est contenue dans $[a, b]$ ; voir
Cohomologie, lemme \ref{cohomology-lemma-tor-amplitude-stalk}.
Un homomorphisme local et plat d'anneaux est fidèlement plat d'après
Algèbre, lemme \ref{algebra-lemma-local-flat-ff}.
Le résultat découle donc de
Compléments sur l'algèbre, lemme
\ref{more-algebra-lemma-flat-descent-tor-amplitude}.
\end{proof}

\begin{lemma}
\label{perfect-lemma-pseudo-coherent-descends-fpqc}
Soit $\{f_i : X_i \to X\}$ un recouvrement fpqc de schémas. Soit
$E \in D_\QCoh(\mathcal{O}_X)$. Soit $m \in \mathbf{Z}$.
Alors $E$ est $m$-pseudo-cohérent si et seulement si chaque
$Lf_i^*E$ est $m$-pseudo-cohérent.
\end{lemma}

\begin{proof}
L'image réciproque préserve toujours la $m$-pseudo-cohérence ; voir
Cohomologie, lemme \ref{cohomology-lemma-pseudo-coherent-pullback}.
Réciproquement, supposons que $Lf_i^*E$ soit $m$-pseudo-cohérent pour tout $i$.
Soit $U \subset X$ un ouvert affine. Il suffit de démontrer que
$E|_U$ est $m$-pseudo-cohérent. Puisque $\{f_i : X_i \to X\}$ est un
recouvrement fpqc, on peut trouver un nombre fini d'ouverts affines $V_j \subset X_{a(j)}$
tels que $f_{a(j)}(V_j) \subset U$ et $U = \bigcup f_{a(j)}(V_j)$.
Posons $V = \coprod V_i$.
On peut donc remplacer $X$ par $U$ et $\{f_i : X_i \to X\}$ par
$\{V \to U\}$, et supposer que $X$ est affine et que le recouvrement
est donné par un unique morphisme surjectif et plat $\{f : Y \to X\}$
de schémas affines. Dans ce cas, le résultat découle de
Compléments sur l'algèbre, lemme \ref{more-algebra-lemma-flat-descent-pseudo-coherent},
à l'aide des lemmes \ref{perfect-lemma-affine-compare-bounded} et
\ref{perfect-lemma-pseudo-coherent-affine}.
\end{proof}

\begin{lemma}
\label{perfect-lemma-pseudo-coherent-descends-fppf}
Soit $\{f_i : X_i \to X\}$ un recouvrement fppf de schémas. Soit
$E \in D(\mathcal{O}_X)$. Soit $m \in \mathbf{Z}$.
Alors $E$ est $m$-pseudo-cohérent si et seulement si chaque
$Lf_i^*E$ est $m$-pseudo-cohérent.
\end{lemma}

\begin{proof}
L'image réciproque préserve toujours la $m$-pseudo-cohérence ; voir
Cohomologie, lemme \ref{cohomology-lemma-pseudo-coherent-pullback}.
Réciproquement, supposons que $Lf_i^*E$ soit $m$-pseudo-cohérent pour tout $i$.
Soit $U \subset X$ un ouvert affine. Il suffit de démontrer que
$E|_U$ est $m$-pseudo-cohérent. Puisque $\{f_i : X_i \to X\}$ est un
recouvrement fppf, on peut trouver un nombre fini d'ouverts affines $V_j \subset X_{a(j)}$
tels que $f_{a(j)}(V_j) \subset U$ et $U = \bigcup f_{a(j)}(V_j)$.
Posons $V = \coprod V_i$.
On peut donc remplacer $X$ par $U$ et $\{f_i : X_i \to X\}$ par
$\{V \to U\}$, et supposer que $X$ est affine et que le recouvrement
est donné par un unique morphisme surjectif et plat $\{f : Y \to X\}$
de présentation finie.

\medskip\noindent
Puisque $f$ est plat, le foncteur dérivé $Lf^*$ est donné par $f^*$, et $f^*$
est exact. Ainsi $H^i(Lf^*E) = f^*H^i(E)$. Puisque $Lf^*E$ est $m$-pseudo-cohérent,
on a $Lf^*E \in D^-(\mathcal{O}_Y)$. Puisque $f$ est surjectif et plat,
on a $E \in D^-(\mathcal{O}_X)$. Soit $i \in \mathbf{Z}$ le plus grand
entier tel que $H^i(E)$ soit non nul. Si $i < m$, la conclusion est acquise. Sinon,
$f^*H^i(E)$ est un $\mathcal{O}_Y$-module de type fini d'après
Cohomologie, lemme \ref{cohomology-lemma-finite-cohomology}.
D'après Descente, lemme \ref{descent-lemma-finite-type-descends-fppf},
le $\mathcal{O}_X$-module $H^i(E)$ est alors de type fini.
Ainsi, quitte à remplacer $X$ par les membres d'un recouvrement ouvert affine fini,
on peut supposer qu'il existe un morphisme
$$
\alpha : \mathcal{O}_X^{\oplus n}[-i] \longrightarrow E
$$
tel que $H^i(\alpha)$ soit surjectif. Soit $C$ le cône de $\alpha$
dans $D(\mathcal{O}_X)$. En prenant l'image réciproque sur $Y$ et en utilisant
Cohomologie, lemme \ref{cohomology-lemma-cone-pseudo-coherent},
on voit que $Lf^*C$ est $m$-pseudo-cohérent. De plus, $H^j(C) = 0$
pour $j \geq i$. Par récurrence sur $i$, on en déduit que $C$ est
$m$-pseudo-cohérent. En appliquant de nouveau
Cohomologie, lemme \ref{cohomology-lemma-cone-pseudo-coherent},
on conclut.
\end{proof}

\begin{lemma}
\label{perfect-lemma-perfect-descends-fpqc}
Soit $\{f_i : X_i \to X\}$ un recouvrement fpqc de schémas. Soit
$E \in D(\mathcal{O}_X)$. Alors $E$ est parfait
si et seulement si chaque $Lf_i^*E$ est parfait.
\end{lemma}

\begin{proof}
L'image réciproque préserve toujours les complexes parfaits ; voir
Cohomologie, lemme \ref{cohomology-lemma-perfect-pullback}.
Réciproquement, supposons que $Lf_i^*E$ soit parfait pour tout $i$.
Les faisceaux de cohomologie de chaque $Lf_i^*E$ sont alors quasi-cohérents ; voir
le lemme \ref{perfect-lemma-pseudo-coherent}
et
Cohomologie, lemme \ref{cohomology-lemma-perfect}.
Puisque les morphismes $f_i$ sont plats, on a $H^p(Lf_i^*E) = f_i^*H^p(E)$.
Les faisceaux de cohomologie de $E$ sont donc quasi-cohérents d'après
Descente, proposition \ref{descent-proposition-fpqc-descent-quasi-coherent}.
Cela étant, le lemme résulte formellement de
Cohomologie, lemme \ref{cohomology-lemma-perfect},
et des
lemmes \ref{perfect-lemma-tor-amplitude-descends} et
\ref{perfect-lemma-pseudo-coherent-descends-fpqc}.
\end{proof}

\begin{lemma}
\label{perfect-lemma-closed-push-pseudo-coherent}
Soit $i : Z \to X$ un morphisme d'espaces annelés tel que
$i$ soit une immersion fermée des espaces topologiques sous-jacents et que
$i_*\mathcal{O}_Z$ soit pseudo-cohérent en tant que $\mathcal{O}_X$-module.
Soit $E \in D(\mathcal{O}_Z)$. Alors $E$ est $m$-pseudo-cohérent
si et seulement si $Ri_*E$ est $m$-pseudo-cohérent.
\end{lemma}

\begin{proof}
Tout au long de cette démonstration, on utilisera l'exactitude du foncteur $i_*$ et,
par suite, l'égalité $Ri_* = i_*$ ; voir Modules, lemme \ref{modules-lemma-i-star-exact}.

\medskip\noindent
Supposons $E$ $m$-pseudo-cohérent. Soit $x \in X$. Nous allons trouver un voisinage
de $x$ sur lequel $i_*E$ soit $m$-pseudo-cohérent. Si $x \not \in Z$,
c'est clair. On peut donc supposer $x \in Z$. Nous utiliserons le fait
que les $U \cap Z$, pour $x \in U \subset X$ ouvert, forment un système fondamental de
voisinages de $x$ dans $Z$. Quitte à restreindre $X$, on peut supposer $E$
borné supérieurement. Nous raisonnons par récurrence sur
le plus grand entier $p$ tel que $H^p(E)$ soit non nul. Si $p < m$,
il n'y a rien à démontrer. Si $p \geq m$, alors $H^p(E)$ est un
$\mathcal{O}_Z$-module de type fini ; voir
Cohomologie, lemme \ref{cohomology-lemma-finite-cohomology}.
On peut donc choisir, quitte à restreindre $X$, un morphisme
$\mathcal{O}_Z^{\oplus n}[-p] \to E$ qui induit une surjection
$\mathcal{O}_Z^{\oplus n} \to H^p(E)$. Choisissons un triangle distingué
$$
\mathcal{O}_Z^{\oplus n}[-p] \to E \to C \to \mathcal{O}_Z^{\oplus n}[-p + 1]
$$
On voit que $H^j(C) = 0$ pour $j \geq p$ et que $C$ est $m$-pseudo-cohérent
d'après Cohomologie, lemme \ref{cohomology-lemma-cone-pseudo-coherent}.
Par récurrence, $i_*C$ est $m$-pseudo-cohérent sur $X$.
Puisque $i_*\mathcal{O}_Z$ est aussi $m$-pseudo-cohérent sur $X$, on déduit
du triangle distingué
$$
i_*\mathcal{O}_Z^{\oplus n}[-p] \to i_*E \to i_*C \to
i_*\mathcal{O}_Z^{\oplus n}[-p + 1]
$$
et de
Cohomologie, lemme \ref{cohomology-lemma-cone-pseudo-coherent},
que $i_*E$ est $m$-pseudo-cohérent.

\medskip\noindent
Supposons que $i_*E$ soit $m$-pseudo-cohérent. Soit $z \in Z$.
Nous allons trouver un voisinage de $z$ sur lequel $E$
soit $m$-pseudo-cohérent. Nous utiliserons le fait
que les $U \cap Z$, pour $z \in U \subset X$ ouvert, forment un système fondamental de
voisinages de $z$ dans $Z$. Quitte à restreindre $X$, on peut supposer $i_*E$
et donc $E$ bornés supérieurement. Nous raisonnons par récurrence sur
le plus grand entier $p$ tel que $H^p(E)$ soit non nul. Si $p < m$,
il n'y a rien à démontrer. Si $p \geq m$, alors $H^p(i_*E) = i_*H^p(E)$
est un $\mathcal{O}_X$-module de type fini ; voir
Cohomologie, lemme \ref{cohomology-lemma-finite-cohomology}.
Choisissons un complexe $\mathcal{E}^\bullet$ de $\mathcal{O}_Z$-modules
représentant $E$. On peut choisir, quitte à restreindre $X$,
un morphisme $\alpha : \mathcal{O}_X^{\oplus n}[-p] \to i_*\mathcal{E}^\bullet$
qui induit une surjection
$\mathcal{O}_X^{\oplus n} \to i_*H^p(\mathcal{E}^\bullet)$.
Par adjonction, on obtient un morphisme
$\alpha : \mathcal{O}_Z^{\oplus n}[-p] \to \mathcal{E}^\bullet$
qui induit une surjection
$\mathcal{O}_Z^{\oplus n} \to H^p(\mathcal{E}^\bullet)$.
Choisissons un triangle distingué
$$
\mathcal{O}_Z^{\oplus n}[-p] \to E \to C \to \mathcal{O}_Z^{\oplus n}[-p + 1]
$$
On voit que $H^j(C) = 0$ pour $j \geq p$. Du triangle distingué
$$
i_*\mathcal{O}_Z^{\oplus n}[-p] \to i_*E \to i_*C \to
i_*\mathcal{O}_Z^{\oplus n}[-p + 1]
$$
du fait que $i_*\mathcal{O}_Z$ est pseudo-cohérent
et de
Cohomologie, lemme \ref{cohomology-lemma-cone-pseudo-coherent},
on déduit que $i_*C$ est $m$-pseudo-cohérent.
Par récurrence, $C$ est $m$-pseudo-cohérent.
En appliquant de nouveau Cohomologie, lemme \ref{cohomology-lemma-cone-pseudo-coherent},
on en déduit que $E$ est $m$-pseudo-cohérent.
\end{proof}

\begin{lemma}
\label{perfect-lemma-finite-push-pseudo-coherent}
Soit $f : X \to Y$ un morphisme fini de schémas tel que
$f_*\mathcal{O}_X$ soit pseudo-cohérent en tant que
$\mathcal{O}_Y$-module\footnote{Cela signifie que $f$ est pseudo-cohérent ; voir
Compléments sur les morphismes, lemme
\ref{more-morphisms-lemma-finite-pseudo-coherent}.}.
Soit $E \in D_\QCoh(\mathcal{O}_X)$. Alors $E$ est $m$-pseudo-cohérent
si et seulement si $Rf_*E$ est $m$-pseudo-cohérent.
\end{lemma}

\begin{proof}
Il s'agit de la traduction de
Compléments sur l'algèbre, lemme \ref{more-algebra-lemma-finite-push-pseudo-coherent},
dans le langage des schémas. Pour effectuer ce passage, utiliser les
lemmes \ref{perfect-lemma-affine-compare-bounded} et
\ref{perfect-lemma-pseudo-coherent-affine}.
\end{proof}


\section{Relèvement des complexes}
\label{perfect-section-lift}

\noindent
Soit $U \subset X$ un sous-espace ouvert d'un espace annelé,
et notons $j : U \to X$ le morphisme d'inclusion. Le foncteur
$D(\mathcal{O}_X) \to D(\mathcal{O}_U)$ est essentiellement surjectif car
$Rj_*$ est un inverse à droite de la restriction.
Dans cette section, nous étendons ce résultat aux complexes dont les faisceaux
de cohomologie sont quasi-cohérents, etc.

\begin{lemma}
\label{perfect-lemma-lift-quasi-coherent}
Soient $X$ un schéma et $j : U \to X$ une immersion ouverte
quasi-compacte. Les foncteurs
$$
D_\QCoh(\mathcal{O}_X) \to D_\QCoh(\mathcal{O}_U)
\quad\text{et}\quad
D^+_\QCoh(\mathcal{O}_X) \to D^+_\QCoh(\mathcal{O}_U)
$$
sont essentiellement surjectifs. Si $X$ est quasi-compact, les foncteurs
$$
D^-_\QCoh(\mathcal{O}_X) \to D^-_\QCoh(\mathcal{O}_U)
\quad\text{et}\quad
D^b_\QCoh(\mathcal{O}_X) \to D^b_\QCoh(\mathcal{O}_U)
$$
sont essentiellement surjectifs.
\end{lemma}

\begin{proof}
L'argument précédant le lemme s'applique au premier cas, car $Rj_*$
envoie $D_\QCoh(\mathcal{O}_U)$ dans $D_\QCoh(\mathcal{O}_X)$
d'après le lemme \ref{perfect-lemma-quasi-coherence-direct-image}.
Il est clair que $Rj_*$ envoie
$D^+_\QCoh(\mathcal{O}_U)$ dans
$D^+_\QCoh(\mathcal{O}_X)$,
ce qui donne l'assertion relative aux complexes bornés inférieurement.
Enfin, le lemme \ref{perfect-lemma-quasi-coherence-direct-image}
assure que $Rj_*$ envoie
$D^-_\QCoh(\mathcal{O}_U)$ dans
$D^-_\QCoh(\mathcal{O}_X)$
si $X$ est quasi-compact. La dernière assertion résulte de ces deux faits.
\end{proof}

\begin{lemma}
\label{perfect-lemma-lift-coherent}
Soient $X$ un schéma noethérien et $j : U \to X$ une immersion ouverte.
Le foncteur
$D^b_{\textit{Coh}}(\mathcal{O}_X) \to D^b_{\textit{Coh}}(\mathcal{O}_U)$
est essentiellement surjectif.
\end{lemma}

\begin{proof}
Soit $K$ un objet de $D^b_{\textit{Coh}}(\mathcal{O}_U)$.
D'après la proposition \ref{perfect-proposition-DCoh}, on peut représenter $K$ par un complexe
borné $\mathcal{F}^\bullet$ de $\mathcal{O}_U$-modules cohérents.
Écrivons $\mathcal{F}^i = 0$ pour $i \not \in [a, b]$, avec $a \leq b$.
Puisque $j$ est quasi-compact et séparé, les termes du complexe borné
$j_*\mathcal{F}^\bullet$ sont des modules quasi-cohérents sur $X$, voir
Schémas, lemme \ref{schemes-lemma-push-forward-quasi-coherent}.
Choisissons par récurrence un sous-module cohérent
$\mathcal{G}^i \subset j_*\mathcal{F}^i$ de la façon suivante.
Pour $i = a$, choisissons un sous-module cohérent quelconque
$\mathcal{G}^a \subset j_*\mathcal{F}^a$ dont la restriction
à $U$ soit $\mathcal{F}^a$. C'est possible d'après
Propriétés, lemme \ref{properties-lemma-extend}.
Pour $i > a$, choisissons d'abord un sous-module cohérent quelconque
$\mathcal{H}^i \subset j_*\mathcal{F}^i$
dont la restriction à $U$ soit $\mathcal{F}^i$,
puis posons
$\mathcal{G}^i = \Im(\mathcal{H}^i \oplus \mathcal{G}^{i - 1}
\to j_*\mathcal{F}^i)$. Il est clair que
$\mathcal{G}^\bullet \subset j_*\mathcal{F}^\bullet$
est un complexe borné de $\mathcal{O}_X$-modules cohérents
dont la restriction à $U$ est $\mathcal{F}^\bullet$, comme voulu.
\end{proof}

\begin{lemma}
\label{perfect-lemma-lift-pseudo-coherent}
Soient $X$ un schéma affine et $U \subset X$ un sous-schéma ouvert
quasi-compact. Pour tout objet pseudo-cohérent $E$ de $D(\mathcal{O}_U)$,
il existe un complexe borné supérieurement de $\mathcal{O}_X$-modules libres de rang fini
dont la restriction à $U$ est isomorphe à $E$.
\end{lemma}

\begin{proof}
D'après le lemme \ref{perfect-lemma-pseudo-coherent}, $E$ est un objet de
$D_\QCoh(\mathcal{O}_U)$. D'après le
lemme \ref{perfect-lemma-lift-quasi-coherent},
on peut supposer que $E = E'|U$ pour un objet $E'$ de
$D_\QCoh(\mathcal{O}_X)$.
Écrivons $X = \Spec(A)$. D'après le lemme \ref{perfect-lemma-affine-compare-bounded},
on peut trouver un complexe $M^\bullet$ de $A$-modules dont le
complexe de $\mathcal{O}_X$-modules associé représente $E'$.

\medskip\noindent
Choisissons $f_1, \ldots, f_r \in A$ tels que $U = D(f_1) \cup \ldots \cup D(f_r)$.
D'après le lemme \ref{perfect-lemma-pseudo-coherent-affine}, les complexes
$M^\bullet_{f_j}$ sont des complexes pseudo-cohérents de $A_{f_j}$-modules.
Soit $n$ un entier. Supposons donné un morphisme de complexes
$\alpha : F^\bullet \to M^\bullet$, où $F^\bullet$ est
borné supérieurement, $F^i = 0$ pour $i < n$, et chaque $F^i$ est un
$R$-module libre de rang fini, tel que
$$
H^i(\alpha_{f_j}) : H^i(F^\bullet_{f_j}) \to H^i(M^\bullet_{f_j})
$$
soit un isomorphisme pour $i > n$ et soit surjectif pour $i = n$. On a le diagramme
$$
\xymatrix{
& F^n \ar[r] \ar[d]^\alpha & F^{n + 1} \ar[d]^\alpha \ar[r] & \ldots \\
M^{n-1} \ar[r] & M^n \ar[r] & M^{n + 1} \ar[r] & \ldots
}
$$
Puisque chaque $M^\bullet_{f_j}$ a sa cohomologie nulle
en degrés assez grands, on peut trouver un tel morphisme pour $n \gg 0$.
Par récurrence sur $n$, nous allons le prolonger en un morphisme
de complexes $F^\bullet \to M^\bullet$
tel que $H^i(\alpha_{f_j})$ soit un isomorphisme
pour tout $i$. Le lemme en résultera en prenant $\widetilde{F^\bullet}$.

\medskip\noindent
L'étape de récurrence consiste à prolonger le diagramme
ci-dessus en ajoutant $F^{n - 1}$. Soit $C^\bullet$ le cône de $\alpha$
(Catégories dérivées, définition \ref{derived-definition-cone}).
La suite exacte longue de cohomologie montre que
$H^i(C^\bullet_{f_j}) = 0$ pour $i \geq n$. D'après
Compléments sur l'algèbre, lemme \ref{more-algebra-lemma-cone-pseudo-coherent},
$C^\bullet_{f_j}$ est $(n - 1)$-pseudo-cohérent. D'après
Compléments sur l'algèbre, lemme \ref{more-algebra-lemma-finite-cohomology},
$H^{n - 1}(C^\bullet_{f_j})$ est un $A_{f_j}$-module de type fini.
Choisissons un $A$-module libre de rang fini $F^{n - 1}$ et un morphisme de $A$-modules
$\beta : F^{n - 1} \to C^{n - 1}$ tels que le composé
$F^{n - 1} \to C^{n - 1} \to C^n$ soit nul et que
$F^{n - 1}_{f_j}$ se surjecte sur $H^{n - 1}(C^\bullet_{f_j})$.
(Certains détails sont omis ; indication : chasser les dénominateurs.)
Puisque $C^{n - 1} = M^{n - 1} \oplus F^n$,
on peut écrire $\beta = (\alpha^{n - 1}, -d^{n - 1})$. La nullité du
composé $F^{n - 1} \to C^{n - 1} \to C^n$ implique
que ces applications s'insèrent dans un morphisme de complexes
$$
\xymatrix{
& F^{n - 1} \ar[d]^{\alpha^{n - 1}} \ar[r]_{d^{n - 1}} &
F^n \ar[r] \ar[d]^\alpha &
F^{n + 1} \ar[d]^\alpha \ar[r] & \ldots \\
\ldots \ar[r] &
M^{n - 1} \ar[r] & M^n \ar[r] & M^{n + 1} \ar[r] & \ldots
}
$$
De plus, ces applications définissent un morphisme de triangles distingués
$$
\xymatrix{
(F^n \to \ldots) \ar[r] \ar[d] &
(F^{n-1} \to \ldots) \ar[r] \ar[d] &
F^{n-1} \ar[r] \ar[d]_\beta &
(F^n \to \ldots)[1] \ar[d] \\
(F^n \to \ldots) \ar[r] &
M^\bullet \ar[r] &
C^\bullet \ar[r] &
(F^n \to \ldots)[1]
}
$$
Le choix de $\beta$ implique donc que le morphisme de complexes
$(F^{-1} \to \ldots) \to M^\bullet$ induit un isomorphisme sur la
cohomologie localisée en $f_j$ en degrés $\geq n$ et une surjection
en degré $n - 1$. Ceci achève la démonstration du lemme.
\end{proof}

\noindent
Les deux lemmes suivants devraient probablement être placés ailleurs.

\begin{lemma}
\label{perfect-lemma-vanishing-ext}
Soit $X$ un schéma quasi-compact et quasi-séparé.
Soit $E \in D^b_\QCoh(\mathcal{O}_X)$.
Il existe un entier $n_0 > 0$ tel que
$\Ext^n_{D(\mathcal{O}_X)}(\mathcal{E}, E) = 0$
pour tout
$\mathcal{O}_X$-module localement libre de rang fini $\mathcal{E}$ et tout $n \geq n_0$.
\end{lemma}

\begin{proof}
Rappelons que $\Ext^n_{D(\mathcal{O}_X)}(\mathcal{E}, E) =
\Hom_{D(\mathcal{O}_X)}(\mathcal{E}, E[n])$. Nous disposons de
Mayer-Vietoris pour les morphismes dans la catégorie dérivée, voir
Cohomologie, lemme \ref{cohomology-lemma-mayer-vietoris-hom}.
Ainsi, si $X = U \cup V$ et si la conclusion du lemme vaut
pour $E|_U$, $E|_V$ et $E|_{U \cap V}$ avec une borne $n_0$,
elle vaut pour $E$ avec la borne $n_0 + 1$.
Il suffit donc de démontrer le lemme lorsque $X$ est affine, voir
Cohomologie des schémas, lemme \ref{coherent-lemma-induction-principle}.

\medskip\noindent
Supposons que $X = \Spec(A)$ soit affine. Choisissons un complexe de $A$-modules
$M^\bullet$ dont le complexe de modules quasi-cohérents associé
représente $E$, voir le lemme \ref{perfect-lemma-affine-compare-bounded}.
Écrivons $\mathcal{E} = \widetilde{P}$ pour un $A$-module $P$.
Puisque $\mathcal{E}$ est localement libre de rang fini, $P$
est un $A$-module projectif de type fini. On a
\begin{align*}
\Hom_{D(\mathcal{O}_X)}(\mathcal{E}, E[n])
& = 
\Hom_{D(A)}(P, M^\bullet[n]) \\
& =
\Hom_{K(A)}(P, M^\bullet[n]) \\
& =
\Hom_A(P, H^n(M^\bullet))
\end{align*}
La première égalité résulte du lemme \ref{perfect-lemma-affine-compare-bounded},
la deuxième de
Catégories dérivées, lemme
\ref{derived-lemma-morphisms-from-projective-complex}, et
la dernière de l'exactitude du foncteur $\Hom_A(P, -)$.
Comme $E$, et donc $M^\bullet$, est borné,
on obtient zéro pour tout $n$ assez grand.
\end{proof}

\noindent
On peut renforcer le lemme suivant (l'annulation est uniforme
pour tous les $L$ dont les faisceaux de cohomologie non nuls
sont concentrés dans un intervalle fixé).

\begin{lemma}
\label{perfect-lemma-ext-from-perfect-into-bounded-QCoh}
Soit $X$ un schéma quasi-compact et quasi-séparé.
Soit $K$ un objet parfait de $D(\mathcal{O}_X)$. Alors
\begin{enumerate}
\item il existe des entiers $a \leq b$ tels que, pour tout
$L \in D_\QCoh(\mathcal{O}_X)$ avec $H^i(L) = 0$ pour $i \in [a, b]$,
on ait $\Hom_{D(\mathcal{O}_X)}(K, L) = 0$, et
\item si $L$ est borné, alors $\Ext^n_{D(\mathcal{O}_X)}(K, L)$
est nul sauf pour un nombre fini de valeurs de $n$.
\end{enumerate}
\end{lemma}

\begin{proof}
L'assertion (2) résulte de (1), car $\Ext^n_{D(\mathcal{O}_X)}(K, L) =
\Hom_{D(\mathcal{O}_X)}(K, L[n])$. Démontrons (1).
Puisque $K$ est parfait, on a
$$
\Hom_{D(\mathcal{O}_X)}(K, L) =
H^0(X, K^\vee \otimes_{\mathcal{O}_X}^\mathbf{L} L)
$$
où $K^\vee$ est le complexe parfait « dual » de $K$, voir
Cohomologie, lemme \ref{cohomology-lemma-dual-perfect-complex}.
Notons que $K^\vee \otimes_{\mathcal{O}_X}^\mathbf{L} L$
appartient à $D_\QCoh(X)$ d'après les
lemmes \ref{perfect-lemma-quasi-coherence-tensor-product} et
\ref{perfect-lemma-pseudo-coherent} (ce dernier montre qu'un complexe parfait
a des faisceaux de cohomologie quasi-cohérents). Supposons que $K^\vee$ ait
une amplitude de Tor contenue dans $[a, b]$. Alors la suite spectrale
$$
E_1^{p, q} = H^p(K^\vee \otimes_{\mathcal{O}_X}^\mathbf{L} H^q(L))
\Rightarrow
H^{p + q}(K^\vee \otimes_{\mathcal{O}_X}^\mathbf{L} L)
$$
montre que $H^j(K^\vee \otimes_{\mathcal{O}_X}^\mathbf{L} L)$
est nul si $H^q(L) = 0$ pour $q \in [j - b, j - a]$.
Soit $N$ l'entier $d$ de Cohomologie des schémas,
lemme \ref{coherent-lemma-vanishing-nr-affines-quasi-separated}.
Alors $H^0(X, K^\vee \otimes_{\mathcal{O}_X}^\mathbf{L} L)$
est nul si les faisceaux de cohomologie
$$
H^{-N}(K^\vee \otimes_{\mathcal{O}_X}^\mathbf{L} L),
\ H^{-N + 1}(K^\vee \otimes_{\mathcal{O}_X}^\mathbf{L} L),
\ \ldots,
\ H^0(K^\vee \otimes_{\mathcal{O}_X}^\mathbf{L} L)
$$
sont nuls. En effet, d'après le lemme cité et le
lemme \ref{perfect-lemma-application-nice-K-injective}, on a
$$
H^0(X, K^\vee \otimes_{\mathcal{O}_X}^\mathbf{L} L) =
H^0(X, \tau_{\geq -N}(K^\vee \otimes_{\mathcal{O}_X}^\mathbf{L} L))
$$
et, par annulation des faisceaux de cohomologie, ceci est égal à
$H^0(X, \tau_{\geq 1}(K^\vee \otimes_{\mathcal{O}_X}^\mathbf{L} L))$,
qui est nul d'après Catégories dérivées, lemme
\ref{derived-lemma-negative-vanishing}.
Il en résulte que $\Hom_{D(\mathcal{O}_X)}(K, L)$ est nul si
$H^i(L) = 0$ pour $i \in [-b - N, -a]$.
\end{proof}

\begin{lemma}
\label{perfect-lemma-lift-perfect-complex-plus-locally-free}
Soit $X$ un schéma affine. Soit $U \subset X$ un ouvert quasi-compact.
Pour tout objet parfait $E$ de $D(\mathcal{O}_U)$, il existe un entier
$r$ et un faisceau localement libre de rang fini $\mathcal{F}$ sur $U$ tels que
$\mathcal{F}[-r] \oplus E$ soit la restriction d'un objet parfait de
$D(\mathcal{O}_X)$.
\end{lemma}

\begin{proof}
Écrivons $X = \Spec(A)$. Rappelons qu'un complexe parfait est
pseudo-cohérent, voir
Cohomologie, lemme \ref{cohomology-lemma-perfect}.
D'après le lemme \ref{perfect-lemma-lift-pseudo-coherent}, on peut trouver un complexe borné supérieurement
$\mathcal{F}^\bullet$ de $A$-modules libres de rang fini tel que $E$ soit
isomorphe à $\mathcal{F}^\bullet|_U$ dans $D(\mathcal{O}_U)$.
D'après Cohomologie, lemme \ref{cohomology-lemma-perfect}, et puisque
$U$ est quasi-compact, $E$ est de dimension de Tor finie ; supposons que
$E$ ait une amplitude de Tor contenue dans $[a, b]$. Choisissons $r < a$ et posons
$$
\mathcal{K} = \Ker(\mathcal{F}^{r} \to \mathcal{F}^{r + 1})
= \Im(\mathcal{F}^{r - 1} \to \mathcal{F}^r).
$$
Puisque $E$ a une amplitude de Tor contenue dans $[a, b]$,
$\mathcal{F} = \mathcal{K}|_U$ est
plat (Cohomologie, lemme \ref{cohomology-lemma-last-one-flat}).
Ainsi $\mathcal{F}$ est plat et de présentation finie, donc localement
libre de rang fini (Propriétés, lemme \ref{properties-lemma-finite-locally-free}).
Il en résulte que
$$
\mathcal{F} \to \mathcal{F}^r|_U \to \mathcal{F}^{r + 1}|_U \to \ldots
$$
est un complexe strictement parfait sur $U$ représentant $E$. D'autre part,
le complexe $P = (\mathcal{F}^r \to \mathcal{F}^{r + 1} \to \ldots )$
est un complexe parfait sur $X$. Les
troncatures bêtes donnent un triangle distingué
$$
P|_U \to E \to \mathcal{F}[-r - 1] \to (P|_U)[1]
$$
Si le morphisme $E \to \mathcal{F}[-r - 1]$ est nul dans $D(\mathcal{O}_U)$,
alors $P|_U = \mathcal{F}[-r - 2] \oplus E$, voir
Catégories dérivées, lemme \ref{derived-lemma-split}.
C'est le cas pour $r \ll 0$, par exemple d'après le
lemme \ref{perfect-lemma-ext-from-perfect-into-bounded-QCoh}.
\end{proof}

\begin{lemma}
\label{perfect-lemma-lift-map}
Soit $X$ un schéma affine. Soit $U \subset X$ un ouvert quasi-compact.
Soient $E, E'$ des objets de $D_\QCoh(\mathcal{O}_X)$, avec $E$ parfait.
Pour tout morphisme $\alpha : E|_U \to E'|_U$, il existe des morphismes
$$
E \xleftarrow{\beta} E_1 \xrightarrow{\gamma} E'
$$
de complexes sur $X$, avec $E_1$ parfait, tels que $\beta : E_1 \to E$
se restreigne en un isomorphisme sur $U$ et que
$\alpha = \gamma|_U \circ \beta|_U^{-1}$.
De plus, on peut supposer que $E_1 = E \otimes_{\mathcal{O}_X}^\mathbf{L} I$
pour un complexe parfait $I$ sur $X$.
\end{lemma}

\begin{proof}
Écrivons $X = \Spec(A)$. Écrivons $U = D(f_1) \cup \ldots \cup D(f_r)$. Choisissons
un complexe borné de $A$-modules projectifs de type fini $M^\bullet$ représentant
$E$ (lemme \ref{perfect-lemma-perfect-affine}). Choisissons un complexe de $A$-modules
$(M')^\bullet$ représentant $E'$ (lemme \ref{perfect-lemma-affine-compare-bounded}).
Alors le complexe $H^\bullet = \Hom_A(M^\bullet, (M')^\bullet)$
est un complexe de $A$-modules dont le complexe de
$\mathcal{O}_X$-modules quasi-cohérents associé représente $R\SheafHom(E, E')$, voir
Cohomologie, lemme \ref{cohomology-lemma-Rhom-strictly-perfect}.
Le morphisme $\alpha$ détermine un élément $s$ de $H^0(U, R\SheafHom(E, E'))$, voir
Cohomologie, lemme \ref{cohomology-lemma-section-RHom-over-U}.
Il existe un $e$ et un morphisme
$$
\xi : I^\bullet(f_1^e, \ldots, f_r^e) \to \Hom_A(M^\bullet, (M')^\bullet)
$$
correspondant à $s$, voir la
proposition \ref{perfect-proposition-represent-cohomology-class-on-open}.
En prenant pour $E_1$ l'objet correspondant au
complexe de $\mathcal{O}_X$-modules quasi-cohérents
associé à
$$
\text{Tot}(I^\bullet(f_1^e, \ldots, f_r^e) \otimes_A M^\bullet)
$$
on obtient $E_1 \to E$ à l'aide du morphisme canonique
$I^\bullet(f_1^e, \ldots, f_r^e) \to A$, et $E_1 \to E'$
à l'aide de $\xi$ et de
Cohomologie, lemme \ref{cohomology-lemma-section-RHom-over-U}.
\end{proof}

\begin{lemma}
\label{perfect-lemma-lift-perfect-complex-plus-shift}
Soit $X$ un schéma affine. Soit $U \subset X$ un ouvert quasi-compact.
Pour tout objet parfait $F$ de $D(\mathcal{O}_U)$,
l'objet $F \oplus F[1]$ est la restriction
d'un objet parfait de $D(\mathcal{O}_X)$.
\end{lemma}

\begin{proof}
D'après le lemme \ref{perfect-lemma-lift-perfect-complex-plus-locally-free},
on peut trouver un objet parfait $E$ de $D(\mathcal{O}_X)$
tel que $E|_U = \mathcal{F}[r] \oplus F$ pour un
$\mathcal{O}_U$-module localement libre de rang fini $\mathcal{F}$.
D'après le lemme \ref{perfect-lemma-lift-map}, on peut trouver un morphisme de
complexes parfaits $\alpha : E_1 \to E$ tel que $(E_1)|_U \cong E|_U$
et que $\alpha|_U$ soit le morphisme
$$
\left(
\begin{matrix}
\text{id}_{\mathcal{F}[r]} & 0 \\
0 & 0
\end{matrix}
\right)
:
\mathcal{F}[r] \oplus F \to \mathcal{F}[r] \oplus F
$$
Le cône de $\alpha$ convient alors.
\end{proof}

\begin{lemma}
\label{perfect-lemma-perfect-into-support-on-T}
Soit $X$ un schéma quasi-compact et quasi-séparé.
Soit $f \in \Gamma(X, \mathcal{O}_X)$.
Pour tout morphisme $\alpha : E \to E'$ dans
$D_\QCoh(\mathcal{O}_X)$ tel que
\begin{enumerate}
\item $E$ soit parfait, et
\item $E'$ soit à support dans $T = V(f)$,
\end{enumerate}
il existe un $n \geq 0$ tel que $f^n \alpha  = 0$.
\end{lemma}

\begin{proof}
Nous disposons de Mayer-Vietoris pour les morphismes dans la catégorie dérivée, voir
Cohomologie, lemme \ref{cohomology-lemma-mayer-vietoris-hom}.
Ainsi, si $X = U \cup V$ et si la conclusion du lemme vaut
pour $f|_U$, $f|_V$ et $f|_{U \cap V}$, elle vaut pour $f$.
Il suffit donc de démontrer le lemme lorsque $X$ est affine, voir
Cohomologie des schémas, lemme \ref{coherent-lemma-induction-principle}.

\medskip\noindent
Soit $X = \Spec(A)$. Alors $f \in A$. Nous utiliserons
l'équivalence $D(A) = D_\QCoh(X)$ du
lemme \ref{perfect-lemma-affine-compare-bounded}
sans la mentionner davantage.
Représentons $E$ par un complexe borné de $A$-modules projectifs de type fini
$P^\bullet$. C'est possible d'après le lemme \ref{perfect-lemma-perfect-affine}.
Soit $t$ le plus grand entier tel que $P^t$ soit non nul.
Le triangle distingué
$$
P^t[-t] \to P^\bullet \to \sigma_{\leq t - 1}P^\bullet \to P^t[-t + 1]
$$
montre que, par récurrence sur la longueur du complexe $P^\bullet$,
on peut se ramener au cas où $P^\bullet$ n'a qu'un seul terme non nul.
Ceci, joint au foncteur de décalage, nous ramène au cas où $P^\bullet$
est constitué d'un seul $A$-module projectif de type fini $P$ en degré $0$.
Représentons $E'$ par un complexe $M^\bullet$ de $A$-modules.
Alors $\alpha$ correspond à un morphisme $P \to H^0(M^\bullet)$.
Puisque le module $H^0(M^\bullet)$ est à support dans $V(f)$ par l'hypothèse (2),
tout élément de $H^0(M^\bullet)$ est annulé par une puissance
de $f$. Comme $P$ est un $A$-module de type fini, le morphisme
$f^n\alpha : P \to H^0(M^\bullet)$ est nul pour un certain $n$, comme voulu.
\end{proof}

\begin{lemma}
\label{perfect-lemma-lift-perfect-complex-plus-shift-support}
Soit $X$ un schéma affine. Soit $T \subset X$ un sous-ensemble fermé
tel que $X \setminus T$ soit quasi-compact. Soit $U \subset X$ un
ouvert quasi-compact. Pour tout objet parfait $F$ de $D(\mathcal{O}_U)$
à support dans $T \cap U$, l'objet $F \oplus F[1]$ est la restriction
d'un objet parfait $E$ de $D(\mathcal{O}_X)$ à support dans $T$.
\end{lemma}

\begin{proof}
Écrivons $T = V(g_1, \ldots, g_s)$. En remplaçant $g_j$ par une puissance,
on peut supposer que la multiplication par $g_j$ est nulle sur $F$, voir le
lemme \ref{perfect-lemma-perfect-into-support-on-T}. Choisissons $E$ comme dans le
lemme \ref{perfect-lemma-lift-perfect-complex-plus-shift}.
Notons que $g_j : E \to E$ se restreint à zéro sur $U$.
Choisissons un triangle distingué
$$
E \xrightarrow{g_1} E \to C_1 \to E[1]
$$
D'après Catégories dérivées, lemme \ref{derived-lemma-split},
l'objet $C_1$ se restreint à
$F \oplus F[1] \oplus F[1] \oplus F[2]$ sur $U$.
De plus, $g_1 : C_1 \to C_1$ est de carré nul d'après
Catégories dérivées, lemme \ref{derived-lemma-third-map-square-zero}.
En effet, le diagramme
$$
\xymatrix{
E \ar[r] \ar[d]_0 & C_1 \ar[d]_{g_1} \ar[r] & E[1] \ar[d]_0 \\
E \ar[r] & C_1 \ar[r] & E[1]
}
$$
est commutatif, car les composés $E \xrightarrow{g_1} E \to C_1$ et
$C_1 \to E[1] \xrightarrow{g_1} E[1]$ sont nuls. En poursuivant, et en prenant
pour $C_{i + 1}$ le cône du morphisme $g_i : C_i \to C_i$, on obtient
un complexe parfait $C_s$ sur $X$ à support dans $T$
dont la restriction à $U$ est
$$
F \oplus F[1]^{\oplus s} \oplus F[2]^{\oplus {s \choose 2}}
\oplus \ldots \oplus F[s]
$$
Choisissons des morphismes de complexes parfaits $\beta : C' \to C_s$
et $\gamma : C' \to C_s$ comme dans le lemme \ref{perfect-lemma-lift-map},
tels que $\beta|_U$ soit un isomorphisme et que
$\gamma|_U \circ \beta|_U^{-1}$ soit le morphisme
$$
F \oplus F[1]^{\oplus s} \oplus F[2]^{\oplus {s \choose 2}}
\oplus \ldots \oplus F[s]
\to
F \oplus F[1]^{\oplus s} \oplus F[2]^{\oplus {s \choose 2}}
\oplus \ldots \oplus F[s]
$$
qui est l'identité sur tous les facteurs directs sauf sur $F$, où il est nul.
D'après le lemme \ref{perfect-lemma-lift-map}, on a aussi
$C' = C_s \otimes^\mathbf{L} I$ pour un complexe parfait
$I$ sur $X$. La nullité de $g_j^2\text{id}_{C_s}$ entraîne donc la
même propriété pour $C'$. Ainsi $C'$ est également à support dans $T$.
Alors $\text{Cone}(\gamma)$ convient.
\end{proof}

\noindent
On trouvera un cas particulier du lemme suivant dans
\cite{Neeman-Grothendieck}.

\begin{lemma}
\label{perfect-lemma-lift-map-from-perfect-complex-with-support}
Soit $X$ un schéma quasi-compact et quasi-séparé.
Soit $U \subset X$ un ouvert quasi-compact. Soit $T \subset X$
un sous-ensemble fermé tel que $X \setminus T$ soit rétrocompact dans $X$.
Soit $E$ un objet de $D_\QCoh(\mathcal{O}_X)$.
Soit $\alpha : P \to E|_U$ un morphisme, où $P$ est un objet parfait de
$D(\mathcal{O}_U)$ à support dans $T \cap U$. Il existe alors un morphisme
$\beta : R \to E$, où $R$ est un objet parfait de $D(\mathcal{O}_X)$
à support dans $T$, tel que $P$ soit facteur direct de $R|_U$ dans
$D(\mathcal{O}_U)$, de façon compatible avec $\alpha$ et $\beta|_U$.
\end{lemma}

\begin{proof}
Puisque $X$ est quasi-compact, il existe un entier $m$ tel que
$X = U \cup V_1 \cup \ldots \cup V_m$ pour des ouverts affines $V_j$ de $X$.
Par récurrence sur $m$, on voit que l'on peut supposer $m = 1$. Autrement
dit, on peut supposer que $X = U \cup V$, avec $V$ affine. D'après le
lemme \ref{perfect-lemma-lift-perfect-complex-plus-shift-support},
on peut choisir un objet parfait $Q$ de $D(\mathcal{O}_V)$
à support dans $T \cap V$ et un isomorphisme
$Q|_{U \cap V} \to (P \oplus P[1])|_{U \cap V}$.
D'après le lemme \ref{perfect-lemma-lift-map}, on peut remplacer $Q$ par
$Q \otimes^\mathbf{L} I$ (toujours à support dans $T \cap V$)
et supposer que le morphisme
$$
Q|_{U \cap V} \to (P \oplus P[1])|_{U \cap V}
\longrightarrow P|_{U \cap V}
\longrightarrow
E|_{U \cap V}
$$
se relève en $Q \to E|_V$. D'après
Cohomologie, lemme \ref{cohomology-lemma-glue},
on trouve un morphisme $a : R \to E$ de $D(\mathcal{O}_X)$
tel que $a|_U$ soit isomorphe à $P \oplus P[1] \to E|_U$
et $a|_V$ soit isomorphe à $Q \to E|_V$.
Ainsi $R$ est parfait et à support dans $T$, comme voulu.
\end{proof}

\begin{remark}
\label{perfect-remark-addendum}
La démonstration du lemme \ref{perfect-lemma-lift-map-from-perfect-complex-with-support}
montre que
$$
R|_U = P \oplus P^{\oplus n_1}[1] \oplus \ldots \oplus P^{\oplus n_m}[m]
$$
pour des entiers $m \geq 0$ et $n_j \geq 0$. Ainsi le faisceau de cohomologie
de plus haut degré de $R|_U$ est celui de $P$. En répétant la construction pour le morphisme
$P^{\oplus n_1}[1] \oplus \ldots \oplus P^{\oplus n_m}[m] \to R|_U$, en prenant
des cônes et en procédant par récurrence, on peut obtenir l'égalité des faisceaux de cohomologie
de $R|_U$ et de $P$ en tout degré supérieur à une borne donnée.
\end{remark}



\section{Approximation par des complexes parfaits}
\label{perfect-section-approximation}

\noindent
Dans cette section, nous étudions l'observation, due à Neeman et Lipman,
selon laquelle un complexe pseudo-cohérent peut être « approché » par des complexes parfaits.

\begin{definition}
\label{perfect-definition-approximation-holds}
Soit $X$ un schéma. Considérons les triplets $(T, E, m)$ où
\begin{enumerate}
\item $T \subset X$ est un sous-ensemble fermé,
\item $E$ est un objet de $D_\QCoh(\mathcal{O}_X)$, et
\item $m \in \mathbf{Z}$.
\end{enumerate}
On dit que {\it l'approximation est possible pour le triplet} $(T, E, m)$ s'il
existe un objet parfait $P$ de $D(\mathcal{O}_X)$ à support dans $T$
et un morphisme $\alpha : P \to E$ qui induit des isomorphismes $H^i(P) \to H^i(E)$
pour $i > m$ et une surjection $H^m(P) \to H^m(E)$.
\end{definition}

\noindent
L'approximation ne peut pas être possible pour tout triplet. En effet, il est clair que,
si l'approximation est possible pour le triplet $(T, E, m)$, alors
\begin{enumerate}
\item $E$ est $m$-pseudo-cohérent, voir
Cohomologie, définition \ref{cohomology-definition-pseudo-coherent}, et
\item les faisceaux de cohomologie $H^i(E)$ sont à support dans $T$ pour $i \geq m$.
\end{enumerate}
De plus, le « support » d'un complexe parfait est un sous-schéma fermé
dont le complémentaire est rétrocompact dans $X$ (détails omis). On ne peut donc
espérer que l'approximation soit possible sans cette hypothèse sur $T$.
Ceci explique en partie les conditions de la définition suivante.

\begin{definition}
\label{perfect-definition-approximation}
Soit $X$ un schéma. On dit que {\it l'approximation par des complexes parfaits est possible}
sur $X$ si, pour tout sous-ensemble fermé $T \subset X$ tel que $X \setminus T$
soit rétrocompact dans $X$, il existe un entier $r$ tel que,
pour tout triplet $(T, E, m)$ comme dans la
définition \ref{perfect-definition-approximation-holds} vérifiant
\begin{enumerate}
\item $E$ est $(m - r)$-pseudo-cohérent, et
\item $H^i(E)$ est à support dans $T$ pour $i \geq m - r$,
\end{enumerate}
l'approximation soit possible.
\end{definition}

\noindent
Nous démontrerons que l'approximation par des complexes parfaits est possible pour les
schémas quasi-compacts et quasi-séparés. Il semble que la deuxième
condition soit nécessaire à notre méthode de démonstration. Il est possible que la
première condition puisse être affaiblie en « $E$ est $m$-pseudo-cohérent »
par une analyse soigneuse des arguments ci-dessous.

\begin{lemma}
\label{perfect-lemma-open}
Soit $X$ un schéma. Soit $U \subset X$ un sous-schéma ouvert.
Soit $(T, E, m)$ un triplet comme dans la
définition \ref{perfect-definition-approximation-holds}.
Si
\begin{enumerate}
\item $T \subset U$,
\item l'approximation est possible pour $(T, E|_U, m)$, et
\item les faisceaux $H^i(E)$ pour $i \geq m$ sont à support dans $T$,
\end{enumerate}
alors l'approximation est possible pour $(T, E, m)$.
\end{lemma}

\begin{proof}
Soit $j : U \to X$ le morphisme d'inclusion.
Si $P \to E|_U$ est une approximation du triplet $(T, E|_U, m)$
sur $U$, alors $j_!P = Rj_*P \to j_!(E|_U) \to E$ est une approximation
de $(T, E, m)$ sur $X$.
Voir Cohomologie, lemmes \ref{cohomology-lemma-pushforward-restriction} et
\ref{cohomology-lemma-pushforward-perfect}.
\end{proof}

\begin{lemma}
\label{perfect-lemma-approximation-affine}
Soit $X$ un schéma affine. Alors l'approximation est possible pour tout
triplet $(T, E, m)$ comme dans la définition \ref{perfect-definition-approximation-holds}
tel qu'il existe un entier $r \geq 0$ vérifiant
\begin{enumerate}
\item $E$ est $m$-pseudo-cohérent,
\item $H^i(E)$ est à support dans $T$ pour $i \geq m - r + 1$,
\item $X \setminus T$ est la réunion de $r$ ouverts affines.
\end{enumerate}
En particulier, l'approximation par des complexes parfaits est possible pour les schémas affines.
\end{lemma}

\begin{proof}
Écrivons $X = \Spec(A)$. Écrivons $T = V(f_1, \ldots, f_r)$.
(Le cas $r = 0$, c'est-à-dire $T = X$, résulte immédiatement du
lemme \ref{perfect-lemma-pseudo-coherent-affine} et des définitions.)
Soit $(T, E, m)$ un triplet comme dans le lemme.
Soit $t$ le plus grand entier tel que $H^t(E)$ soit non nul.
Nous procéderons par récurrence sur $t$. Le cas initial est $t < m$ ;
le résultat est alors trivial. Supposons maintenant que $t \geq m$. D'après
Cohomologie, lemme \ref{cohomology-lemma-finite-cohomology},
le faisceau $H^t(E)$ est de type fini. Comme il est quasi-cohérent,
il est engendré par un nombre fini de sections
(Propriétés, lemme \ref{properties-lemma-finite-type-module}).
Pour tout $s \in \Gamma(X, H^t(E)) = H^t(X, E)$
(voir la démonstration du lemme \ref{perfect-lemma-affine-compare-bounded}),
on peut trouver un $e > 0$ et un morphisme $K_e[-t] \to E$
tels que $s$ appartienne à l'image de
$H^0(K_e) = H^t(K_e[-t]) \to H^t(E)$, voir le
lemme \ref{perfect-lemma-represent-cohomology-class-on-closed}.
Une somme directe finie de ces morphismes donne un morphisme
$P \to E$, où $P$ est un complexe parfait à support dans $T$,
où $H^i(P) = 0$ pour $i > t$, et où $H^t(P) \to E$ est
surjectif. Choisissons un triangle distingué
$$
P \to E \to E' \to P[1]
$$
Alors $E'$ est $m$-pseudo-cohérent
(Cohomologie, lemme \ref{cohomology-lemma-cone-pseudo-coherent}),
$H^i(E') = 0$ pour $i \geq t$, et
$H^i(E')$ est à support dans $T$ pour $i \geq m - r + 1$.
Par récurrence, on trouve une approximation $P' \to E'$
de $(T, E', m)$. Insérons le composé $P' \to E' \to P[1]$
dans un triangle distingué $P \to P'' \to P' \to P[1]$
et prolongeons les morphismes $P' \to E'$ et $P[1] \to P[1]$ en
un morphisme de triangles distingués
$$
\xymatrix{
P \ar[r] \ar[d] & P'' \ar[d] \ar[r] & P' \ar[d] \ar[r] & P[1] \ar[d] \\
P \ar[r] &  E \ar[r] & E' \ar[r] & P[1]
}
$$
à l'aide de TR3. Alors $P''$ est un complexe parfait
(Cohomologie, lemme \ref{cohomology-lemma-two-out-of-three-perfect})
à support dans $T$.
Une chasse au diagramme immédiate montre que $P'' \to E$ est l'approximation
cherchée.
\end{proof}

\begin{lemma}
\label{perfect-lemma-induction-step}
Soit $X$ un schéma. Soit $X = U \cup V$ un recouvrement ouvert,
avec $U$ quasi-compact, $V$ affine et $U \cap V$ quasi-compact.
Si l'approximation par des complexes parfaits est possible sur $U$,
elle est possible sur $X$.
\end{lemma}

\begin{proof}
Soit $T \subset X$ un sous-ensemble fermé tel que $X \setminus T$ soit rétrocompact
dans $X$. Soit $r_U$ l'entier de la définition \ref{perfect-definition-approximation}
adapté au couple $(U, T \cap U)$.
Posons $T' = T \setminus U$. Notons que
$T' \subset V$ et que $V \setminus T' = (X \setminus T) \cap U \cap V$
est quasi-compact par notre hypothèse sur $T$.
Soit $r'$ le nombre d'ouverts affines nécessaires pour recouvrir $V \setminus T'$.
Nous affirmons que $r = \max(r_U, r')$ convient pour le couple $(X, T)$.

\medskip\noindent
Pour le voir, choisissons un triplet $(T, E, m)$ tel que $E$ soit
$(m - r)$-pseudo-cohérent et que $H^i(E)$ soit à support dans $T$ pour
$i \geq m - r$. Soit $t$ le plus grand entier tel que
$H^t(E)|_U$ soit non nul. (Un tel entier existe, car $U$ est quasi-compact
et $E|_U$ est $(m - r)$-pseudo-cohérent.)
Nous démontrerons que $E$ peut être approché en procédant par récurrence sur $t$.

\medskip\noindent
Cas initial : $t \leq m - r'$. Cela signifie que $H^i(E)$ est à support
dans $T'$ pour $i \geq m - r'$. Le
lemme \ref{perfect-lemma-approximation-affine}
assure donc l'existence d'une approximation
$P \to E|_V$ de $(T', E|_V, m)$ sur $V$.
En appliquant le lemme \ref{perfect-lemma-open}, on voit que
$(T', E, m)$ peut être approché. Une telle approximation
est aussi une approximation de $(T, E, m)$.

\medskip\noindent
Étape de récurrence. Choisissons une approximation $P \to E|_U$
de $(T \cap U, E|_U, m)$. Elle donne en particulier une surjection
$H^t(P) \to H^t(E|_U)$. D'après le
lemme \ref{perfect-lemma-lift-perfect-complex-plus-shift-support},
on peut choisir un objet parfait $Q$ de $D(\mathcal{O}_V)$
à support dans $T \cap V$ et un isomorphisme
$Q|_{U \cap V} \to (P \oplus P[1])|_{U \cap V}$.
D'après le lemme \ref{perfect-lemma-lift-map}, on peut remplacer $Q$ par
$Q \otimes^\mathbf{L} I$
et supposer que le morphisme
$$
Q|_{U \cap V} \to (P \oplus P[1])|_{U \cap V}
\longrightarrow P|_{U \cap V}
\longrightarrow
E|_{U \cap V}
$$
se relève en $Q \to E|_V$. D'après
Cohomologie, lemme \ref{cohomology-lemma-glue},
on trouve un morphisme $a : R \to E$ de $D(\mathcal{O}_X)$
tel que $a|_U$ soit isomorphe à $P \oplus P[1] \to E|_U$
et $a|_V$ soit isomorphe à $Q \to E|_V$.
Ainsi $R$ est parfait et à support dans $T$,
et le morphisme $H^t(R) \to H^t(E)$ est surjectif après restriction à $U$.
Choisissons un triangle distingué
$$
R \to E \to E' \to R[1]
$$
Alors $E'$ est $(m - r)$-pseudo-cohérent
(Cohomologie, lemme \ref{cohomology-lemma-cone-pseudo-coherent}),
$H^i(E')|_U = 0$ pour $i \geq t$, et
$H^i(E')$ est à support dans $T$ pour $i \geq m - r$.
Par récurrence, on trouve une approximation $R' \to E'$
de $(T, E', m)$. Insérons le composé $R' \to E' \to R[1]$
dans un triangle distingué $R \to R'' \to R' \to R[1]$
et prolongeons les morphismes $R' \to E'$ et $R[1] \to R[1]$ en
un morphisme de triangles distingués
$$
\xymatrix{
R \ar[r] \ar[d] & R'' \ar[d] \ar[r] & R' \ar[d] \ar[r] & R[1] \ar[d] \\
R \ar[r] &  E \ar[r] & E' \ar[r] & R[1]
}
$$
à l'aide de TR3. Alors $R''$ est un complexe parfait
(Cohomologie, lemme \ref{cohomology-lemma-two-out-of-three-perfect})
à support dans $T$.
Une chasse au diagramme immédiate montre que $R'' \to E$ est l'approximation
cherchée.
\end{proof}

\begin{theorem}
\label{perfect-theorem-approximation}
Soit $X$ un schéma quasi-compact et quasi-séparé.
Alors l'approximation par des complexes parfaits est possible sur $X$.
\end{theorem}

\begin{proof}
Ceci résulte du principe de récurrence de
Cohomologie des schémas, lemme \ref{coherent-lemma-induction-principle},
et des lemmes \ref{perfect-lemma-induction-step} et \ref{perfect-lemma-approximation-affine}.
\end{proof}





\section{Générateurs des catégories dérivées}
\label{perfect-section-generating}

\noindent
Dans cette section, nous démontrons que la catégorie dérivée
$D_\QCoh(\mathcal{O}_X)$ d'un schéma quasi-compact
et quasi-séparé peut être engendrée par un seul objet parfait.
Nous invitons vivement le lecteur à lire la démonstration de ce résultat dans le remarquable article de
Bondal et van den Bergh, voir \cite{BvdB}.

\begin{lemma}
\label{perfect-lemma-direct-summand-of-a-restriction}
Soit $X$ un schéma quasi-compact et quasi-séparé.
Soit $U$ un sous-schéma ouvert quasi-compact.
Soit $P$ un objet parfait de $D(\mathcal{O}_U)$.
Alors $P$ est facteur direct de la restriction d'un objet parfait
de $D(\mathcal{O}_X)$.
\end{lemma}

\begin{proof}
C'est un cas particulier du lemme \ref{perfect-lemma-lift-map-from-perfect-complex-with-support}.
\end{proof}

\begin{lemma}
\label{perfect-lemma-orthogonal-koszul-complex}
\begin{reference}
\cite[proposition 6.1]{Bokstedt-Neeman}
\end{reference}
Dans la situation \ref{perfect-situation-complex}, notons $j : U \to X$ l'immersion
ouverte et soit $K$ l'objet parfait de $D(\mathcal{O}_X)$
correspondant au complexe de Koszul de $f_1, \ldots, f_r$ sur $A$.
Pour $E \in D_\QCoh(\mathcal{O}_X)$, les conditions suivantes sont équivalentes :
\begin{enumerate}
\item $E = Rj_*(E|_U)$, et
\item $\Hom_{D(\mathcal{O}_X)}(K[n], E) = 0$ pour tout $n \in \mathbf{Z}$.
\end{enumerate}
\end{lemma}

\begin{proof}
Choisissons un triangle distingué $E \to Rj_*(E|_U) \to N \to E[1]$.
Observons que
$$
\Hom_{D(\mathcal{O}_X)}(K[n], Rj_*(E|_U)) =
\Hom_{D(\mathcal{O}_U)}(K|_U[n], E) = 0
$$
pour tout $n$, car $K|_U = 0$. Il suffit donc de démontrer le résultat pour
$N$. Autrement dit, on peut supposer que $E$ se restreint à zéro sur $U$.
Observons qu'il existe des triangles distingués
$$
K^\bullet(f_1^{e_1}, \ldots, f_i^{e'_i}, \ldots, f_r^{e_r}) \to
K^\bullet(f_1^{e_1}, \ldots, f_i^{e'_i + e''_i}, \ldots, f_r^{e_r}) \to
K^\bullet(f_1^{e_1}, \ldots, f_i^{e''_i}, \ldots, f_r^{e_r}) \to \ldots
$$
de complexes de Koszul, voir
Compléments sur l'algèbre, lemme \ref{more-algebra-lemma-koszul-mult}.
Ainsi, si $\Hom_{D(\mathcal{O}_X)}(K[n], E) = 0$ pour tout $n \in \mathbf{Z}$,
la même propriété vaut en remplaçant $K$ par
$K_e$ comme dans le lemme \ref{perfect-lemma-represent-cohomology-class-on-closed}.
Notre lemme résulte donc immédiatement de ce dernier et du fait que $E$
est déterminé par le complexe de $A$-modules $R\Gamma(X, E)$, voir le
lemme \ref{perfect-lemma-affine-compare-bounded}.
\end{proof}

\begin{theorem}
\label{perfect-theorem-bondal-van-den-Bergh}
Soit $X$ un schéma quasi-compact et quasi-séparé. La catégorie
$D_\QCoh(\mathcal{O}_X)$ peut être engendrée par un seul
objet parfait. Plus précisément, il existe un objet parfait
$P$ de $D(\mathcal{O}_X)$ tel que, pour
$E \in D_\QCoh(\mathcal{O}_X)$, les conditions suivantes soient équivalentes :
\begin{enumerate}
\item $E = 0$, et
\item $\Hom_{D(\mathcal{O}_X)}(P[n], E) = 0$ pour tout $n \in \mathbf{Z}$.
\end{enumerate}
\end{theorem}

\begin{proof}
Nous utiliserons le principe de récurrence de
Cohomologie des schémas, lemme \ref{coherent-lemma-induction-principle}.

\medskip\noindent
Si $X$ est affine, alors $\mathcal{O}_X$ est un générateur parfait.
Ceci résulte du lemme \ref{perfect-lemma-affine-compare-bounded}.

\medskip\noindent
Supposons que $X = U \cup V$ soit un recouvrement ouvert, avec $U$ quasi-compact,
tel que le théorème soit vrai pour $U$, et que $V$ soit un ouvert affine.
Soit $P$ un objet parfait de $D(\mathcal{O}_U)$ qui engendre
$D_\QCoh(\mathcal{O}_U)$. À l'aide du
lemme \ref{perfect-lemma-direct-summand-of-a-restriction}, on peut
choisir un objet parfait
$Q$ de $D(\mathcal{O}_X)$ dont la restriction à $U$ est une somme directe dont
l'un des facteurs est $P$. Écrivons $V = \Spec(A)$. Soit $Z = X \setminus U$.
C'est un sous-ensemble fermé de $V$ tel que $V \setminus Z$ soit quasi-compact.
Choisissons $f_1, \ldots, f_r \in A$ tels que
$Z = V(f_1, \ldots, f_r)$. Soit $K \in D(\mathcal{O}_V)$ l'objet parfait
correspondant au complexe de Koszul de $f_1, \ldots, f_r$ sur $A$.
Comme $K$ est à support dans le fermé $Z \subset V$, l'image directe
$K' = R(V \to X)_*K$ est un objet parfait de $D(\mathcal{O}_X)$ dont la
restriction à $V$ est $K$ (voir
Cohomologie, lemme \ref{cohomology-lemma-pushforward-perfect}).
Nous affirmons que $Q \oplus K'$ engendre
$D_\QCoh(\mathcal{O}_X)$.

\medskip\noindent
Soit $E$ un objet de $D_\QCoh(\mathcal{O}_X)$ tel qu'il
n'existe aucun morphisme non nul d'un décalé quelconque de $Q \oplus K'$ vers $E$.
D'après Cohomologie, lemme \ref{cohomology-lemma-pushforward-restriction},
on a $K' =  R(V \to X)_! K$, et donc
$$
\Hom_{D(\mathcal{O}_X)}(K'[n], E) = \Hom_{D(\mathcal{O}_V)}(K[n], E|_V)
$$
D'après le lemme \ref{perfect-lemma-orthogonal-koszul-complex}, l'annulation de
ces groupes implique que $E|_V$ est isomorphe à
$R(U \cap V \to V)_*E|_{U \cap V}$. Il en résulte que $E = R(U \to X)_*E|_U$
(un petit détail est omis). Dans ce cas,
$$
\Hom_{D(\mathcal{O}_X)}(Q[n], E) = \Hom_{D(\mathcal{O}_U)}(Q|_U[n], E|_U)
$$
contient $\Hom_{D(\mathcal{O}_U)}(P[n], E|_U)$ comme facteur direct.
Par le choix de $P$, l'annulation de ces groupes implique donc que $E|_U$
est nul. Par suite, $E$ est nul.
\end{proof}

\noindent
Le résultat suivant renforce le
théorème \ref{perfect-theorem-bondal-van-den-Bergh}
et se démontre par exactement les mêmes méthodes.
Rappelons que, pour un sous-ensemble fermé $T$ d'un schéma $X$, on note
$D_T(\mathcal{O}_X)$ la sous-catégorie strictement pleine, saturée
et triangulée de $D(\mathcal{O}_X)$ formée des objets
à support dans $T$ (définition \ref{perfect-definition-supported-on}).
De même, on note $D_{\QCoh, T}(\mathcal{O}_X)$ la sous-catégorie strictement pleine, saturée
et triangulée de $D(\mathcal{O}_X)$ formée des
complexes dont les faisceaux de cohomologie sont quasi-cohérents et à support
dans $T$.

\begin{lemma}
\label{perfect-lemma-generator-with-support}
\begin{reference}
\cite[théorème 6.8]{Rouquier-dimensions}
\end{reference}
Soit $X$ un schéma quasi-compact et quasi-séparé. Soit $T \subset X$ un
sous-ensemble fermé tel que $X \setminus T$ soit quasi-compact. Avec les notations
ci-dessus, la catégorie $D_{\QCoh, T}(\mathcal{O}_X)$ est engendrée par un
seul objet parfait.
\end{lemma}

\begin{proof}
Nous utiliserons le principe de récurrence de
Cohomologie des schémas, lemme \ref{coherent-lemma-induction-principle}.

\medskip\noindent
Supposons que $X = \Spec(A)$ soit affine. Il existe alors
$f_1, \ldots, f_r \in A$ tels que $T = V(f_1, \ldots, f_r)$.
Soit $K$ le complexe de Koszul de $f_1, \ldots, f_r$ comme dans le
lemme \ref{perfect-lemma-orthogonal-koszul-complex}.
Alors $K$ est un objet parfait dont la cohomologie est à support dans
$T$, donc un objet parfait de $D_{\QCoh, T}(\mathcal{O}_X)$.
D'autre part, si $E \in D_{\QCoh, T}(\mathcal{O}_X)$ et si
$\Hom(K, E[n]) = 0$ pour tout $n$, le
lemme \ref{perfect-lemma-orthogonal-koszul-complex}
donne $E = Rj_*(E|_{X \setminus T}) = 0$.
Ainsi $K$ engendre $D_{\QCoh, T}(\mathcal{O}_X)$
(selon notre définition des générateurs de catégories triangulées dans
Catégories dérivées, définition \ref{derived-definition-generators}).

\medskip\noindent
Supposons que $X = U \cup V$ soit un recouvrement ouvert, avec $V$ affine et
$U$ quasi-compact, tel que le lemme soit vrai pour $U$.
Soit $P$ un objet parfait de $D(\mathcal{O}_U)$ à support dans $T \cap U$
qui engendre $D_{\QCoh, T \cap U}(\mathcal{O}_U)$. À l'aide du
lemme \ref{perfect-lemma-lift-map-from-perfect-complex-with-support},
on peut choisir un objet parfait $Q$ de $D(\mathcal{O}_X)$ à support dans $T$
dont la restriction à $U$ est une somme directe dont l'un des facteurs est $P$.
Écrivons $V = \Spec(B)$. Soit $Z = X \setminus U$. Alors $Z$ est un sous-ensemble fermé
de $V$ tel que $V \setminus Z$ soit quasi-compact. Comme $X$ est quasi-séparé,
il en résulte que $Z \cap T$ est un sous-ensemble fermé de $V$ tel que
$W = V \setminus (Z \cap T)$ soit quasi-compact. On peut donc choisir
$g_1, \ldots, g_s \in B$ tels que $Z \cap T = V(g_1, \ldots, g_r)$.
Soit $K \in D(\mathcal{O}_V)$ l'objet parfait correspondant au
complexe de Koszul de $g_1, \ldots, g_s$ sur $B$. Comme $K$ est
à support dans le fermé $(Z \cap T) \subset V$, l'image directe
$K' = R(V \to X)_*K$ est un objet parfait de $D(\mathcal{O}_X)$ dont la
restriction à $V$ est $K$ (voir
Cohomologie, lemme \ref{cohomology-lemma-pushforward-perfect}).
Nous affirmons que $Q \oplus K'$ engendre
$D_{\QCoh, T}(\mathcal{O}_X)$.

\medskip\noindent
Soit $E$ un objet de $D_{\QCoh, T}(\mathcal{O}_X)$ tel qu'il
n'existe aucun morphisme non nul d'un décalé quelconque de $Q \oplus K'$ vers $E$.
D'après Cohomologie, lemme \ref{cohomology-lemma-pushforward-restriction},
on a $K' =  R(V \to X)_! K$, et donc
$$
\Hom_{D(\mathcal{O}_X)}(K'[n], E) = \Hom_{D(\mathcal{O}_V)}(K[n], E|_V)
$$
D'après le lemme \ref{perfect-lemma-orthogonal-koszul-complex}, on a donc
$E|_V = Rj_*E|_W$, où $j : W \to V$ est l'inclusion. On a le diagramme
$$
\xymatrix{
W \ar[r]_j & V & Z \cap T \ar[l] \ar[d] \\
U \cap V \ar[u]^{j'} \ar[ru]_{j''} & & Z \ar[lu]
}
$$
Puisque $E$ est à support dans $T$, $E|_W$ est à support dans
$T \cap W = T \cap U \cap V$, qui est fermé dans $W$.
On en déduit que
$$
E|_V = Rj_*(E|_W) = Rj_*(Rj'_*(E|_{U \cap V})) = Rj''_*(E|_{U \cap V})
$$
où la deuxième égalité résulte de l'assertion (1) de
Cohomologie, lemme \ref{cohomology-lemma-pushforward-restriction}.
Il en résulte que $E = R(U \to X)_*E|_U$ (un petit détail est omis).
Dans ce cas,
$$
\Hom_{D(\mathcal{O}_X)}(Q[n], E) = \Hom_{D(\mathcal{O}_U)}(Q|_U[n], E|_U)
$$
contient $\Hom_{D(\mathcal{O}_U)}(P[n], E|_U)$ comme facteur direct.
Par le choix de $P$, l'annulation de ces groupes implique donc que $E|_U$
est nul. Par suite, $E$ est nul.
\end{proof}

\section{Un exemple de générateur}
\label{perfect-section-example-generator}

\noindent
Dans cette section, nous démontrons que la catégorie dérivée de l'espace projectif
sur un anneau est engendrée par un fibré vectoriel, et même par une somme
directe de tordus du faisceau structural.

\medskip\noindent
Le lemme suivant affirme que $\bigoplus_{n \geq 0} \mathcal{L}^{\otimes -n}$
est un générateur si $\mathcal{L}$ est ample.

\begin{lemma}
\label{perfect-lemma-nonzero-some-cohomology}
Soient $X$ un schéma et $\mathcal{L}$ un
$\mathcal{O}_X$-module inversible ample. Si $K$ est un objet non nul de
$D_\QCoh(\mathcal{O}_X)$, alors, pour un certain $n \geq 0$ et un certain $p \in \mathbf{Z}$,
le groupe de cohomologie
$H^p(X, K \otimes_{\mathcal{O}_X}^\mathbf{L} \mathcal{L}^{\otimes n})$
est non nul.
\end{lemma}

\begin{proof}
Rappelons que, puisque $X$ possède un faisceau inversible ample, il est quasi-compact
et séparé (Propriétés, définition \ref{properties-definition-ample} et
lemme \ref{properties-lemma-affine-s-opens-cover-quasi-separated}).
On peut donc appliquer la
proposition \ref{perfect-proposition-quasi-compact-affine-diagonal}
et représenter $K$ par un complexe $\mathcal{F}^\bullet$ de
modules quasi-cohérents. Choisissons un $p$ tel que
$\mathcal{H}^p = \Ker(\mathcal{F}^p \to \mathcal{F}^{p + 1})/
\Im(\mathcal{F}^{p - 1} \to \mathcal{F}^p)$ soit non nul.
Choisissons un point $x \in X$ tel que la fibre $\mathcal{H}^p_x$ soit
non nulle. Choisissons un $n \geq 0$ et un $s \in \Gamma(X, \mathcal{L}^{\otimes n})$
tels que $X_s$ soit un voisinage ouvert affine de $x$.
Choisissons $\tau \in \mathcal{H}^p(X_s)$ dont l'image soit un élément non nul
de la fibre $\mathcal{H}^p_x$ ; c'est possible,
car $\mathcal{H}^p$ est quasi-cohérent et $X_s$ est affine.
Puisque le foncteur des sections sur $X_s$ est exact sur les
modules quasi-cohérents, on peut trouver une section $\tau' \in \mathcal{F}^p(X_s)$
d'image nulle dans $\mathcal{F}^{p + 1}(X_s)$ et d'image
$\tau$ dans $\mathcal{H}^p(X_s)$. D'après
Propriétés, lemme \ref{properties-lemma-invert-s-sections},
il existe un $m$ tel que $\tau' \otimes s^{\otimes m}$
soit l'image d'une section
$\tau'' \in \Gamma(X, \mathcal{F}^p \otimes \mathcal{L}^{\otimes mn})$.
En appliquant encore une fois le même lemme, on trouve $l \geq 0$ tel que
$\tau'' \otimes s^{\otimes l}$ ait une image nulle dans
$\mathcal{F}^{p + 1} \otimes \mathcal{L}^{\otimes (m + l)n}$.
Alors $\tau''$ donne une classe non nulle dans
$H^p(X, K \otimes^\mathbf{L}_{\mathcal{O}_X} \mathcal{L}^{(m + l)n})$,
comme voulu.
\end{proof}

\begin{lemma}
\label{perfect-lemma-construct-the-next-one}
Soit $A$ un anneau. Soit $X = \mathbf{P}^n_A$. Pour tout $a \in \mathbf{Z}$,
il existe un complexe exact
$$
0 \to \mathcal{O}_X(a) \to \ldots
\to \mathcal{O}_X(a + i)^{\oplus {n + 1 \choose i}} \to
\ldots \to \mathcal{O}_X(a + n + 1) \to 0
$$
de fibrés vectoriels sur $X$.
\end{lemma}

\begin{proof}
Rappelons que $\mathbf{P}^n_A$ est $\text{Proj}(A[X_0, \ldots, X_n])$, voir
Constructions, définition \ref{constructions-definition-projective-space}.
Considérons le complexe de Koszul
$$
K_\bullet = K_\bullet(A[X_0, \ldots, X_n], X_0, \ldots, X_n)
$$
sur $S = A[X_0, \ldots, X_n]$ associé à $X_0, \ldots, X_n$.
Puisque $X_0, \ldots, X_n$ est manifestement une suite régulière dans
l'anneau de polynômes $S$, on voit
(Compléments sur l'algèbre, lemme \ref{more-algebra-lemma-regular-koszul-regular})
que le complexe de Koszul $K_\bullet$ est exact, sauf en degré $0$,
où la cohomologie est $S/(X_0, \ldots, X_n)$.
Notons que $K_\bullet$ devient un complexe de modules gradués si l'on
place les générateurs de $K_i$ en degré $+i$. On obtient donc un
complexe exact
$$
0 \to S(-n - 1) \to \ldots \to S(-n - 1 + i)^{\oplus {n \choose i}} \to \ldots
\to S \to S/(X_0, \ldots, X_n) \to 0
$$
En appliquant le foncteur exact $\tilde{\ }$ de Constructions,
lemme \ref{constructions-lemma-proj-sheaves}, et en utilisant le fait que
le dernier terme appartient au noyau de ce foncteur,
on obtient le complexe exact
$$
0 \to \mathcal{O}_X(-n - 1) \to \ldots
\to \mathcal{O}_X(-n - 1 + i)^{\oplus {n + 1 \choose i}} \to
\ldots \to \mathcal{O}_X \to 0
$$
En tensorisant par les faisceaux inversibles $\mathcal{O}_X(n + a + 1)$,
on obtient les complexes exacts du lemme.
\end{proof}

\begin{lemma}
\label{perfect-lemma-generator-P1}
Soit $A$ un anneau. Soit $X = \mathbf{P}^n_A$. Alors
$$
E =
\mathcal{O}_X \oplus \mathcal{O}_X(-1) \oplus \ldots \oplus \mathcal{O}_X(-n)
$$
est un générateur
(Catégories dérivées, définition \ref{derived-definition-generators})
de $D_\QCoh(X)$.
\end{lemma}

\begin{proof}
Soit $K \in D_\QCoh(\mathcal{O}_X)$. Supposons que
$\Hom(E, K[p]) = 0$ pour tout $p \in \mathbf{Z}$.
Nous devons montrer que $K = 0$.
D'après Catégories dérivées, lemme
\ref{derived-lemma-right-orthogonal},
$\Hom(E', K[p])$ est nul pour tout $E' \in \langle E \rangle$
et tout $p \in \mathbf{Z}$.
D'après le lemme \ref{perfect-lemma-construct-the-next-one}
appliqué avec $a = -n - 1$,
on a $\mathcal{O}_X(-n - 1) \in \langle E \rangle$,
car ce faisceau est quasi-isomorphe à un complexe borné
dont les termes sont des sommes directes finies de facteurs directs de $E$.
En répétant l'argument avec $a = -n - 2$, on obtient
$\mathcal{O}_X(-n - 2) \in \langle E \rangle$.
Par récurrence, on trouve $\mathcal{O}_X(-m) \in \langle E \rangle$
pour tout $m \geq 0$.
Puisque
$$
\Hom(\mathcal{O}_X(-m), K[p]) =
H^p(X, K \otimes_{\mathcal{O}_X}^\mathbf{L} \mathcal{O}_X(m)) =
H^p(X, K \otimes_{\mathcal{O}_X}^\mathbf{L} \mathcal{O}_X(1)^{\otimes m})
$$
on en déduit que $K = 0$ d'après le lemme \ref{perfect-lemma-nonzero-some-cohomology}.
(On utilise aussi le fait que $\mathcal{O}_X(1)$ est un
faisceau inversible ample sur $X$, ce qui résulte de
Propriétés, lemme \ref{properties-lemma-open-in-proj-ample}.)
\end{proof}

\begin{remark}
\label{perfect-remark-pullback-generator}
Soit $f : X \to Y$ un morphisme de schémas quasi-compacts et quasi-séparés.
Soit $E \in D_\QCoh(\mathcal{O}_Y)$ un générateur
(voir le théorème \ref{perfect-theorem-bondal-van-den-Bergh}).
Alors les conditions suivantes sont équivalentes :
\begin{enumerate}
\item pour $K \in D_\QCoh(\mathcal{O}_X)$, on a
$Rf_*K = 0$ si et seulement si $K = 0$,
\item $Rf_* : D_\QCoh(\mathcal{O}_X) \to D_\QCoh(\mathcal{O}_Y)$
est conservatif, et
\item $Lf^*E$ engendre $D_\QCoh(\mathcal{O}_X)$.
\end{enumerate}
L'équivalence de (1) et (2) est une conséquence formelle du fait que
$Rf_* : D_\QCoh(\mathcal{O}_X) \to D_\QCoh(\mathcal{O}_Y)$ est un
foncteur exact de catégories triangulées. De même, l'équivalence
de (1) et (3) résulte formellement du fait que $Lf^*$
est adjoint à gauche de $Rf_*$.
Ces conditions sont vérifiées si $f$ est affine (lemme \ref{perfect-lemma-affine-morphism}),
ou si $f$ est une immersion ouverte, ou si $f$ est un composé de tels morphismes.
On en déduit que
\begin{enumerate}
\item si $X$ est un schéma quasi-affine, alors $\mathcal{O}_X$ engendre
$D_\QCoh(\mathcal{O}_X)$,
\item si $X \subset \mathbf{P}^n_A$ est un
sous-schéma localement fermé quasi-compact, alors
$\mathcal{O}_X \oplus \mathcal{O}_X(-1) \oplus \ldots \oplus \mathcal{O}_X(-n)$
engendre $D_\QCoh(\mathcal{O}_X)$ d'après le
lemme \ref{perfect-lemma-generator-P1}.
\end{enumerate}
\end{remark}




\section{Objets compacts et parfaits}
\label{perfect-section-compact}

\noindent
Soit $X$ un schéma noethérien de dimension finie. D'après
Cohomologie, proposition \ref{cohomology-proposition-vanishing-Noetherian},
et
Cohomologie sur les sites, lemme \ref{sites-cohomology-lemma-when-jshriek-compact},
les faisceaux de modules $j_!\mathcal{O}_U$ sont des objets compacts
de $D(\mathcal{O}_X)$ pour tout ouvert $U \subset X$.
Ces faisceaux ne sont en général pas quasi-cohérents et ne donnent donc pas
des objets parfaits de la catégorie dérivée $D(\mathcal{O}_X)$.
Cependant, si l'on se restreint aux complexes dont les faisceaux de cohomologie
sont quasi-cohérents, cette situation ne se présente pas.
Voici l'énoncé précis.

\begin{proposition}
\label{perfect-proposition-compact-is-perfect}
Soit $X$ un schéma quasi-compact et quasi-séparé.
Un objet de $D_\QCoh(\mathcal{O}_X)$ est compact
si et seulement s'il est parfait.
\end{proposition}

\begin{proof}
Si $K$ est un objet parfait de $D(\mathcal{O}_X)$ de dual
$K^\vee$ (Cohomologie, lemme \ref{cohomology-lemma-dual-perfect-complex}),
on a
$$
\Hom_{D(\mathcal{O}_X)}(K, M) =
H^0(X, K^\vee \otimes_{\mathcal{O}_X}^\mathbf{L} M)
$$
fonctoriellement en $M$. Puisque $K^\vee \otimes_{\mathcal{O}_X}^\mathbf{L} -$
commute aux sommes directes et que $H^0(X, -)$ commute aux sommes
directes sur $D_\QCoh(\mathcal{O}_X)$ d'après le
lemme \ref{perfect-lemma-quasi-coherence-pushforward-direct-sums},
on en déduit que $K$ est compact dans $D_\QCoh(\mathcal{O}_X)$.

\medskip\noindent
Réciproquement, soit $K$ un objet compact de $D_\QCoh(\mathcal{O}_X)$.
Pour montrer que $K$ est parfait, il suffit de montrer que
$K|_U$ est parfait pour tout ouvert affine $U \subset X$, voir
Cohomologie, lemme \ref{cohomology-lemma-perfect-independent-representative}.
Observons que $j : U \to X$ est un morphisme quasi-compact et séparé.
Par conséquent,
$Rj_* : D_\QCoh(\mathcal{O}_U) \to D_\QCoh(\mathcal{O}_X)$
commute aux sommes directes, voir le
lemme \ref{perfect-lemma-quasi-coherence-pushforward-direct-sums}.
L'adjonction entre la restriction à $U$ et $Rj_*$ implique donc que
$K|_U$ est un objet compact de $D_\QCoh(\mathcal{O}_U)$.
On se ramène ainsi au cas où $X$ est affine.

\medskip\noindent
Supposons que $X = \Spec(A)$ soit affine. D'après le lemme \ref{perfect-lemma-affine-compare-bounded},
le problème se traduit en la même question pour $D(A)$.
Pour $D(A)$, le résultat est
Compléments sur l'algèbre, proposition \ref{more-algebra-proposition-perfect-is-compact}.
\end{proof}

\begin{remark}
\label{perfect-remark-classical-generator}
Soit $X$ un schéma quasi-compact et quasi-séparé. Soit $G$ un
objet parfait de $D(\mathcal{O}_X)$ qui engendre
$D_\QCoh(\mathcal{O}_X)$. D'après le théorème \ref{perfect-theorem-bondal-van-den-Bergh},
il en existe au moins un. En combinant le
lemme \ref{perfect-lemma-quasi-coherence-direct-sums} avec la
proposition \ref{perfect-proposition-compact-is-perfect} et avec
Catégories dérivées, proposition
\ref{derived-proposition-generator-versus-classical-generator},
on voit que $G$ est un générateur classique de $D_{perf}(\mathcal{O}_X)$.
\end{remark}

\noindent
Le résultat suivant renforce la
proposition \ref{perfect-proposition-compact-is-perfect}.
Soit $T \subset X$ un sous-ensemble fermé d'un schéma $X$. Comme précédemment,
$D_T(\mathcal{O}_X)$ désigne la sous-catégorie strictement pleine, saturée
et triangulée de $D(\mathcal{O}_X)$ formée des objets
à support dans $T$ (définition \ref{perfect-definition-supported-on}).
Puisque les sommes directes commutent aux faisceaux de cohomologie, il en résulte
que $D_T(\mathcal{O}_X)$ possède des sommes directes, qui coïncident
avec celles de $D(\mathcal{O}_X)$.

\begin{lemma}
\label{perfect-lemma-compact-is-perfect-with-support}
Soit $X$ un schéma quasi-compact et quasi-séparé.
Soit $T \subset X$ un sous-ensemble fermé tel que $X \setminus T$
soit quasi-compact. Un objet de $D_{\QCoh, T}(\mathcal{O}_X)$ est compact
si et seulement s'il est parfait en tant qu'objet de $D(\mathcal{O}_X)$.
\end{lemma}

\begin{proof}
Observons que $D_{\QCoh, T}(\mathcal{O}_X)$ est une catégorie triangulée
possédant des sommes directes, d'après la remarque précédant le lemme.
D'après la proposition \ref{perfect-proposition-compact-is-perfect},
les objets parfaits définissent des objets compacts de $D(\mathcal{O}_X)$,
donc a fortiori de toute sous-catégorie stable par sommes
directes. Pour la réciproque, nous utiliserons l'existence d'un générateur
$E \in D_{\QCoh, T}(\mathcal{O}_X)$ qui est un complexe parfait
de $\mathcal{O}_X$-modules, voir le
lemme \ref{perfect-lemma-generator-with-support}.
D'après ce qui précède, $E$ est donc compact. Il résulte alors de
Catégories dérivées, proposition
\ref{derived-proposition-generator-versus-classical-generator},
que $E$ est un générateur classique de la sous-catégorie pleine
des objets compacts de $D_{\QCoh, T}(\mathcal{O}_X)$.
Ainsi, tout objet compact se construit à partir de $E$ par
une suite finie d'opérations consistant à prendre
(a) des décalés, (b) des sommes directes finies, (c) des cônes, et
(d) des facteurs directs. Chacune de ces opérations conserve
la propriété d'être parfait, d'où le résultat.
\end{proof}

\begin{remark}
\label{perfect-remark-classical-generator-with-support}
Soit $X$ un schéma quasi-compact et quasi-séparé.
Soit $T \subset X$ un sous-ensemble fermé tel que $X \setminus T$
soit quasi-compact. Soit $G$ un
objet parfait de $D_{\QCoh, T}(\mathcal{O}_X)$ qui engendre
$D_{\QCoh, T}(\mathcal{O}_X)$. D'après le lemme \ref{perfect-lemma-generator-with-support},
il en existe au moins un. En combinant l'existence de sommes directes dans
$D_{\QCoh, T}(\mathcal{O}_X)$ avec le
lemme \ref{perfect-lemma-compact-is-perfect-with-support} et avec
Catégories dérivées, proposition
\ref{derived-proposition-generator-versus-classical-generator},
on voit que $G$ est un générateur classique de $D_{perf, T}(\mathcal{O}_X)$.
\end{remark}

\noindent
Le lemme suivant est une application des idées intervenant dans
la démonstration du lemme précédent.

\begin{lemma}
\label{perfect-lemma-map-from-pseudo-coherent-to-complex-with-support}
Soit $X$ un schéma quasi-compact et quasi-séparé. Soit $T \subset X$
un sous-ensemble fermé tel que $U = X \setminus T$ soit quasi-compact.
Soit $\alpha : P \to E$ un morphisme de $D_\QCoh(\mathcal{O}_X)$ vérifiant
l'une des conditions suivantes :
\begin{enumerate}
\item $P$ est parfait et $E$ est à support dans $T$, ou
\item $P$ est pseudo-cohérent, $E$ est à support dans $T$, et $E$ est borné inférieurement.
\end{enumerate}
Alors il existe un complexe parfait de $\mathcal{O}_X$-modules $I$
et un morphisme $I \to \mathcal{O}_X[0]$ tels que
$I \otimes^\mathbf{L} P \to E$ soit nul et que
$I|_U \to \mathcal{O}_U[0]$ soit un
isomorphisme.
\end{lemma}

\begin{proof}
Posons $\mathcal{D} = D_{\QCoh, T}(\mathcal{O}_X)$. Dans les deux cas, le complexe
$K = R\SheafHom(P, E)$ est un objet de $\mathcal{D}$. Voir le
lemme \ref{perfect-lemma-quasi-coherence-internal-hom} pour la quasi-cohérence.
Il est clair que $K$ est à support dans $T$, car la formation de $R\SheafHom$
commute à la restriction aux ouverts.
Le morphisme $\alpha$ définit un élément de
$H^0(K) = \Hom_{D(\mathcal{O}_X)}(\mathcal{O}_X[0], K)$.
Il suffit alors de démontrer le résultat pour le morphisme
$\alpha : \mathcal{O}_X[0] \to K$.

\medskip\noindent
Soit $E \in \mathcal{D}$ un générateur parfait, voir le
lemme \ref{perfect-lemma-generator-with-support}. Écrivons
$$
K = \text{hocolim} K_n
$$
comme dans Catégories dérivées, lemme \ref{derived-lemma-write-as-colimit},
en utilisant le générateur $E$. Puisque le foncteur $\mathcal{D} \to D(\mathcal{O}_X)$
commute aux sommes directes, on voit que $K = \text{hocolim} K_n$
vaut dans $D(\mathcal{O}_X)$. Puisque $\mathcal{O}_X$ est un objet compact
de $D(\mathcal{O}_X)$, on trouve un $n$ et un morphisme
$\alpha_n : \mathcal{O}_X \to K_n$ qui induit $\alpha$, voir
Catégories dérivées, lemme \ref{derived-lemma-commutes-with-countable-sums}.
D'après Catégories dérivées, lemme \ref{derived-lemma-factor-through},
appliqué au morphisme $\mathcal{O}_X[0] \to K_n$ dans la catégorie
ambiante $D(\mathcal{O}_X)$, on voit que $\alpha_n$ se factorise en
$\mathcal{O}_X[0] \to Q \to K_n$, où $Q$ est un objet
de $\langle E \rangle$. On en déduit que $Q$ est un complexe parfait
à support dans $T$.

\medskip\noindent
Choisissons un triangle distingué
$$
I \to \mathcal{O}_X[0] \to Q \to I[1]
$$
Par construction, $I$ est parfait, le morphisme $I \to \mathcal{O}_X[0]$
se restreint en un isomorphisme sur $U$, et le composé
$I \to K$ est nul, car $\alpha$ se factorise par $Q$.
Ceci démontre le lemme.
\end{proof}






\section{Catégories dérivées comme catégories de modules}
\label{perfect-section-derived-is-dga}

\noindent
Dans cette section, nous tirons quelques conséquences de ce qui précède.
Nous avons d'abord besoin de quelques lemmes supplémentaires.

\begin{lemma}
\label{perfect-lemma-tensor-with-QCoh-complex}
Soit $X$ un schéma. Soit $K^\bullet$ un complexe de $\mathcal{O}_X$-modules
dont les faisceaux de cohomologie sont quasi-cohérents.\newline Soit
$(E, d) = \Hom_{\text{Comp}^{dg}(\mathcal{O}_X)}(K^\bullet, K^\bullet)$
l'algèbre différentielle graduée des endomorphismes. Alors le foncteur
$$
- \otimes_E^\mathbf{L} K^\bullet :
D(E, \text{d}) \longrightarrow D(\mathcal{O}_X)
$$
de
Algèbre différentielle graduée, lemme
\ref{dga-lemma-tensor-with-complex-derived},
a son image contenue dans $D_\QCoh(\mathcal{O}_X)$.
\end{lemma}

\begin{proof}
Soit $P$ un $E$-module différentiel gradué ayant la propriété (P),
et soit $F_\bullet$ une filtration de $P$ comme dans
Algèbre différentielle graduée, section \ref{dga-section-P-resolutions}.
On a alors
$$
P \otimes_E K^\bullet = \text{hocolim}\ F_iP \otimes_E K^\bullet
$$
Chacun des $F_iP$ possède une filtration finie dont les quotients successifs
sont des sommes directes de $E[k]$. Le résultat s'ensuit aisément.
\end{proof}

\noindent
Le résultat suivant est tiré de \cite{BvdB}.

\begin{theorem}
\label{perfect-theorem-DQCoh-is-Ddga}
Soit $X$ un schéma quasi-compact et quasi-séparé.
Il existe alors une algèbre différentielle graduée $(E, \text{d})$
dont seul un nombre fini de groupes de cohomologie $H^i(E)$ sont non nuls,
telle que $D_\QCoh(\mathcal{O}_X)$ soit équivalente
à $D(E, \text{d})$.
\end{theorem}

\begin{proof}
Soit $K^\bullet$ un complexe K-injectif de $\mathcal{O}$-modules qui
soit parfait et engendre $D_\QCoh(\mathcal{O}_X)$. Un tel complexe
existe d'après le théorème \ref{perfect-theorem-bondal-van-den-Bergh}
et l'existence des résolutions K-injectives. Nous démontrerons que le
théorème est vérifié avec
$$
(E, \text{d}) = \Hom_{\text{Comp}^{dg}(\mathcal{O}_X)}(K^\bullet, K^\bullet)
$$
où $\text{Comp}^{dg}(\mathcal{O}_X)$ est la catégorie différentielle graduée
des complexes de $\mathcal{O}$-modules. Voir
Algèbre différentielle graduée, section \ref{dga-section-variant-base-change}.
Puisque $K^\bullet$ est K-injectif, on
a
\begin{equation}
\label{perfect-equation-E-is-OK}
H^n(E) = \Ext^n_{D(\mathcal{O}_X)}(K^\bullet, K^\bullet)
\end{equation}
pour tout $n \in \mathbf{Z}$. Seul un nombre fini de ces groupes Ext
sont non nuls, d'après le lemme \ref{perfect-lemma-ext-from-perfect-into-bounded-QCoh}.
Considérons le foncteur
$$
- \otimes_E^\mathbf{L} K^\bullet :
D(E, \text{d}) \longrightarrow D(\mathcal{O}_X)
$$
de
Algèbre différentielle graduée, lemme
\ref{dga-lemma-tensor-with-complex-derived}.
Puisque $K^\bullet$ est parfait, il définit un objet compact de
$D(\mathcal{O}_X)$, voir la proposition \ref{perfect-proposition-compact-is-perfect}.
En combinant ceci avec (\ref{perfect-equation-E-is-OK}), on voit que le foncteur ci-dessus est pleinement
fidèle, d'après
Algèbre différentielle graduée, lemmes
\ref{dga-lemma-fully-faithful-in-compact-case}. Il possède un adjoint à droite
$$
R\Hom(K^\bullet, - ) : D(\mathcal{O}_X) \longrightarrow D(E, \text{d})
$$
d'après Algèbre différentielle graduée, lemmes
\ref{dga-lemma-tensor-with-complex-hom-adjoint},
qui est un foncteur quasi-inverse à gauche par les propriétés générales des foncteurs
adjoints. D'autre part, il résulte du
lemme \ref{perfect-lemma-tensor-with-QCoh-complex} que l'on obtient
$$
- \otimes_E^\mathbf{L} K^\bullet :
D(E, \text{d}) \longrightarrow D_\QCoh(\mathcal{O}_X)
$$
et, par le choix de $K^\bullet$ comme générateur de
$D_\QCoh(\mathcal{O}_X)$, le noyau de l'adjoint
restreint à $D_\QCoh(\mathcal{O}_X)$ est nul.
Un argument formel montre que l'on obtient l'équivalence voulue, voir
Catégories dérivées, lemme
\ref{derived-lemma-fully-faithful-adjoint-kernel-zero}.
\end{proof}

\begin{remark}[Variante avec support]
\label{perfect-remark-DQCoh-is-Ddga-with-support}
Soit $X$ un schéma quasi-compact et quasi-séparé. Soit $T \subset X$
un sous-ensemble fermé tel que $X \setminus T$ soit quasi-compact.
L'analogue du théorème \ref{perfect-theorem-DQCoh-is-Ddga} est vrai
pour $D_{\QCoh, T}(\mathcal{O}_X)$.
Ceci résulte exactement du même argument que dans la démonstration
du théorème, en utilisant les
lemmes \ref{perfect-lemma-generator-with-support} et
\ref{perfect-lemma-compact-is-perfect-with-support}
et une variante avec supports du lemme \ref{perfect-lemma-tensor-with-QCoh-complex}.

Si nous en avons besoin, nous énoncerons ici précisément le
résultat et en donnerons une démonstration détaillée.
\end{remark}

\begin{remark}[Unicité de l'algèbre différentielle graduée]
\label{perfect-remark-independence-choice}
Soit $X$ un schéma quasi-compact et quasi-séparé sur un anneau $R$.
Par la construction de la démonstration du
théorème \ref{perfect-theorem-DQCoh-is-Ddga},
il existe une algèbre différentielle graduée $(A, \text{d})$ sur $R$
telle que $D_\QCoh(X)$ soit $R$-linéairement équivalente à
$D(A, \text{d})$ comme catégorie triangulée.
On peut se demander : dans quelle mesure $(A, \text{d})$ est-elle unique ?
La réponse n'est que légèrement plus forte que l'affirmation selon laquelle
$(A, \text{d})$ est bien définie à équivalence dérivée près.
Supposons en effet que $(B, \text{d})$ soit un second tel couple.
On a alors
$$
(A, \text{d}) = \Hom_{\text{Comp}^{dg}(\mathcal{O}_X)}(K^\bullet, K^\bullet)
$$
et
$$
(B, \text{d}) = \Hom_{\text{Comp}^{dg}(\mathcal{O}_X)}(L^\bullet, L^\bullet)
$$
pour des complexes K-injectifs $K^\bullet$ et $L^\bullet$
de $\mathcal{O}_X$-modules correspondant à des générateurs parfaits
de $D_\QCoh(\mathcal{O}_X)$. Posons
$$
\Omega = \Hom_{\text{Comp}^{dg}(\mathcal{O}_X)}(K^\bullet, L^\bullet)
\quad
\Omega' = \Hom_{\text{Comp}^{dg}(\mathcal{O}_X)}(L^\bullet, K^\bullet)
$$
Alors $\Omega$ est un $B^{opp} \otimes_R A$-module différentiel gradué,
et $\Omega'$ est un $A^{opp} \otimes_R B$-module différentiel gradué.
De plus, l'équivalence
$$
D(A, \text{d}) \to D_\QCoh(\mathcal{O}_X) \to
D(B, \text{d})
$$
est donnée par le foncteur $- \otimes_A^\mathbf{L} \Omega'$, et il en va
de même pour le quasi-inverse. Nous sommes donc dans la situation
d'Algèbre différentielle graduée, remarque \ref{dga-remark-hochschild-cohomology}.
Si nous avons besoin de cette remarque, nous donnerons ici un énoncé précis
et une démonstration détaillée.
\end{remark}




\section{Caractérisation des complexes pseudo-cohérents, I}
\label{perfect-section-pseudo-coherent-hocolim}

\noindent
Les méthodes précédentes permettent de caractériser les objets pseudo-cohérents
comme limites inductives homotopiques d'approximations par des objets parfaits.

\begin{lemma}
\label{perfect-lemma-pseudo-coherent-hocolim}
Soit $X$ un schéma quasi-compact et quasi-séparé.
Soit $K \in D(\mathcal{O}_X)$. Les conditions suivantes sont équivalentes :
\begin{enumerate}
\item $K$ est pseudo-cohérent, et
\item $K = \text{hocolim} K_n$, où
$K_n$ est parfait et $\tau_{\geq -n}K_n \to \tau_{\geq -n}K$
est un isomorphisme pour tout $n$.
\end{enumerate}
\end{lemma}

\begin{proof}
L'implication (2) $\Rightarrow$ (1) est vraie sur tout espace annelé.
Supposons en effet (2) vérifiée. Rappelons qu'un objet parfait de la catégorie
dérivée est pseudo-cohérent, voir
Cohomologie, lemme \ref{cohomology-lemma-perfect}.
Il résulte alors des définitions que
$\tau_{\geq -n}K_n$ est $(-n + 1)$-pseudo-cohérent,
donc que $\tau_{\geq -n}K$ est $(-n + 1)$-pseudo-cohérent,
et par suite que $K$ est $(-n + 1)$-pseudo-cohérent. Ceci est vrai pour
tout $n$, donc $K$ est pseudo-cohérent, voir
Cohomologie, définition \ref{cohomology-definition-pseudo-coherent}.

\medskip\noindent
Supposons (1) vérifiée. Commençons par choisir une approximation
$K_1 \to K$ de $(X, K, -2)$ par un complexe parfait $K_1$, voir les
définitions \ref{perfect-definition-approximation-holds} et
\ref{perfect-definition-approximation} et le
théorème \ref{perfect-theorem-approximation}.
Supposons avoir construit par récurrence
$$
K_1 \to K_2 \to \ldots \to K_n \to K
$$
avec $K_i$ parfait et tel que
$\tau_{\geq -i}K_i \to \tau_{\geq -i}K$ soit un isomorphisme
pour tout $1 \leq i \leq n$. Choisissons alors $a \leq b$ comme dans le
lemme \ref{perfect-lemma-ext-from-perfect-into-bounded-QCoh}
pour l'objet parfait $K_n$. Choisissons une approximation
$K_{n + 1} \to K$ de $(X, K, \min(a - 1, -n - 1))$.
Choisissons un triangle distingué
$$
K_{n + 1} \to K \to C \to K_{n + 1}[1]
$$
On voit alors que $C \in D_\QCoh(\mathcal{O}_X)$ vérifie
$H^i(C) = 0$ pour $i \geq a$. Par le choix de $a, b$,
on a donc $\Hom_{D(\mathcal{O}_X)}(K_n, C) = 0$.
Le composé $K_n \to K \to C$ est donc nul. D'après
Catégories dérivées, lemme \ref{derived-lemma-representable-homological},
on peut donc factoriser $K_n \to K$ par $K_{n + 1}$,
ce qui démontre l'étape de récurrence.

\medskip\noindent
Il reste à démontrer que $K = \text{hocolim} K_n$.
Ceci résulte de l'application de
Catégories dérivées, lemme \ref{derived-lemma-cohomology-of-hocolim},
aux foncteurs
$H^i( - ) : D(\mathcal{O}_X) \to \textit{Mod}(\mathcal{O}_X)$
et du choix des $K_n$.
\end{proof}

\begin{lemma}
\label{perfect-lemma-pseudo-coherent-hocolim-with-support}
Soit $X$ un schéma quasi-compact et quasi-séparé.
Soit $T \subset X$ un sous-ensemble fermé tel que $X \setminus T$
soit quasi-compact. Soit $K \in D(\mathcal{O}_X)$ à support dans $T$.
Les conditions suivantes sont équivalentes :
\begin{enumerate}
\item $K$ est pseudo-cohérent, et
\item $K = \text{hocolim} K_n$, où
$K_n$ est parfait, à support dans $T$, et
$\tau_{\geq -n}K_n \to \tau_{\geq -n}K$ est un isomorphisme pour tout $n$.
\end{enumerate}
\end{lemma}

\begin{proof}
La démonstration de ce lemme est exactement celle du
lemme \ref{perfect-lemma-pseudo-coherent-hocolim},
à ceci près que, pour choisir les approximations, on utilise
les triplets $(T, K, m)$.
\end{proof}









\section{Un exemple d'équivalence}
\label{perfect-section-example-equivalence}

\noindent
Dans la section \ref{perfect-section-example-generator},
nous avons démontré que la catégorie dérivée de l'espace projectif $\mathbf{P}^n_A$
sur un anneau $A$ est engendrée par un fibré vectoriel, et même par une somme
directe de tordus du faisceau structural. Dans cette section, nous démontrons que ceci
détermine une équivalence entre $D_\QCoh(\mathcal{O}_{\mathbf{P}^n_A})$
et la catégorie dérivée d'une $A$-algèbre.

\medskip\noindent
Avant d'énoncer le résultat, nous avons besoin de quelques notations.
Soit $A$ un anneau. Soit $X = \mathbf{P}^n_A = \text{Proj}(S)$,
où $S = A[X_0, \ldots, X_n]$. D'après le lemme \ref{perfect-lemma-generator-P1},
nous savons que
\begin{equation}
\label{perfect-equation-generator-Pn}
P =
\mathcal{O}_X \oplus \mathcal{O}_X(-1) \oplus \ldots \oplus \mathcal{O}_X(-n)
\end{equation}
est un générateur parfait de $D_\QCoh(\mathcal{O}_X)$.
Considérons la $A$-algèbre (non commutative)
\begin{equation}
\label{perfect-equation-algebra-for-Pn}
R = \Hom_{\mathcal{O}_X}(P, P) =
\left(
\begin{matrix}
S_0 & S_1 & S_2 & \ldots & \ldots \\
0 & S_0 & S_1 & \ldots & \ldots\\
0 & 0 & S_0 & \ldots & \ldots \\
\ldots & \ldots & \ldots & \ldots & \ldots \\
0 & \ldots & \ldots & \ldots & S_0
\end{matrix}
\right)
\end{equation}
munie de la multiplication et de l'addition évidentes. Si l'on considère $P$ comme un complexe
de $\mathcal{O}_X$-modules de la manière habituelle (c'est-à-dire avec $P$ en degré $0$
et zéro dans tous les autres degrés), on a
$$
R = \Hom_{\text{Comp}^{dg}(\mathcal{O}_X)}(P, P)
$$
où, dans le membre de droite, on considère $R$ comme une algèbre différentielle graduée
sur $A$
à différentielle nulle (c'est-à-dire avec $R$ en degré $0$ et zéro dans tous
les autres degrés). D'après la discussion de
Algèbre différentielle graduée, section \ref{dga-section-variant-base-change},
on obtient un foncteur dérivé
$$
- \otimes_R^\mathbf{L} P : D(R) \longrightarrow D(\mathcal{O}_X),
$$
voir en particulier Algèbre différentielle graduée, lemme
\ref{dga-lemma-tensor-with-complex-derived}.
D'après le lemme \ref{perfect-lemma-tensor-with-QCoh-complex},
l'image essentielle de ce foncteur
est contenue dans $D_\QCoh(\mathcal{O}_X)$.

\begin{lemma}
\label{perfect-lemma-Pn-module-category}
\begin{reference}
\cite{Beilinson}
\end{reference}
Soit $A$ un anneau. Soit $X = \mathbf{P}^n_A = \text{Proj}(S)$,
où $S = A[X_0, \ldots, X_n]$. Avec
$P$ comme dans (\ref{perfect-equation-generator-Pn}) et
$R$ comme dans (\ref{perfect-equation-algebra-for-Pn}),
le foncteur
$$
- \otimes_R^\mathbf{L} P : D(R) \longrightarrow D_\QCoh(\mathcal{O}_X)
$$
est une équivalence $A$-linéaire de catégories triangulées envoyant $R$
sur $P$.
\end{lemma}

\noindent
Autrement dit, la catégorie dérivée des modules quasi-cohérents sur
l'espace projectif est équivalente à la catégorie dérivée des modules
sur une algèbre (non commutative).
Cette propriété de l'espace projectif semble assez exceptionnelle parmi
tous les schémas projectifs sur $A$.

\begin{proof}
Pour démontrer que notre foncteur est pleinement fidèle, il suffit de démontrer
que $\Ext^i_X(P, P)$ est nul pour $i \not = 0$ et égal
à $R$ pour $i = 0$, voir
Algèbre différentielle graduée, lemme
\ref{dga-lemma-fully-faithful-in-compact-case}.
Comme dans la démonstration du
lemme \ref{perfect-lemma-ext-from-perfect-into-bounded-QCoh},
on voit que
$$
\Ext^i_X(P, P) = H^i(X, P^\wedge \otimes P) =
\bigoplus\nolimits_{0 \leq a, b \leq n} H^i(X, \mathcal{O}_X(a - b))
$$
Le calcul de la cohomologie de l'espace projectif
(Cohomologie des schémas, lemme
\ref{coherent-lemma-cohomology-projective-space-over-ring})
montre que
ces groupes $\Ext$ sont nuls sauf si $i = 0$.
Pour $i = 0$, on retrouve $R$, puisque c'est ainsi que nous avons défini $R$
dans (\ref{perfect-equation-algebra-for-Pn}).
D'après Algèbre différentielle graduée, lemme
\ref{dga-lemma-tensor-with-complex-hom-adjoint},
notre foncteur possède un adjoint à droite, à savoir
$R\Hom(P, -) : D_\QCoh(\mathcal{O}_X) \to D(R)$.
Puisque $P$ engendre $D_\QCoh(\mathcal{O}_X)$ d'après le
lemme \ref{perfect-lemma-generator-P1},
on voit que le noyau de $R\Hom(P, -)$ est nul.
Notre foncteur est donc une équivalence de catégories
triangulées, d'après Catégories dérivées, lemme
\ref{derived-lemma-fully-faithful-adjoint-kernel-zero}.
\end{proof}






\section{Nouvelle étude du cohérateur}
\label{perfect-section-better-coherator}

\noindent
Dans la section \ref{perfect-section-coherator}, nous avons construit et étudié
l'adjoint à droite $RQ_X$ du foncteur canonique
$D(\QCoh(\mathcal{O}_X)) \to D(\mathcal{O}_X)$.
Il a été construit comme le prolongement dérivé à droite du cohérateur
$Q_X : \textit{Mod}(\mathcal{O}_X) \to \QCoh(\mathcal{O}_X)$.
Dans cette section, nous étudions quand le foncteur d'inclusion
$$
D_\QCoh(\mathcal{O}_X) \longrightarrow D(\mathcal{O}_X)
$$
possède un adjoint à droite. Si cet adjoint à droite existe, nous le
noterons\footnote{Cette notation est probablement inhabituelle. Cependant, nous avons déjà
employé $Q_X$ pour le cohérateur et $RQ_X$ pour son prolongement dérivé.}
$$
DQ_X :
D(\mathcal{O}_X) \longrightarrow D_\QCoh(\mathcal{O}_X)
$$
Il se trouve que les schémas quasi-compacts et quasi-séparés
possèdent un tel adjoint à droite.

\begin{lemma}
\label{perfect-lemma-better-coherator}
Soit $X$ un schéma quasi-compact et quasi-séparé.
Le foncteur d'inclusion $D_\QCoh(\mathcal{O}_X) \to D(\mathcal{O}_X)$
possède un adjoint à droite $DQ_X$.
\end{lemma}

\begin{proof}[Première démonstration]
Nous utiliserons le principe de récurrence de
Cohomologie des schémas, lemme \ref{coherent-lemma-induction-principle},
pour démontrer ceci. Si $D(\QCoh(\mathcal{O}_X)) \to D_\QCoh(\mathcal{O}_X)$
est une équivalence, le lemme est vrai, car le foncteur
$RQ_X$ de la section \ref{perfect-section-coherator} est adjoint à droite du foncteur
$D(\QCoh(\mathcal{O}_X)) \to D(\mathcal{O}_X)$.
En particulier, notre lemme est vrai pour les schémas affines, voir le
lemme \ref{perfect-lemma-affine-coherator}.
Il suffit donc de montrer que, si $X = U \cup V$
est une réunion de deux ouverts quasi-compacts et si le
lemme est vrai pour $U$, $V$ et $U \cap V$, alors il est vrai pour $X$.

\medskip\noindent
L'adjoint existe si et seulement si, pour tout objet $K$ de
$D(\mathcal{O}_X)$, on peut trouver un triangle distingué
$$
E' \to E \to K \to E'[1]
$$
dans $D(\mathcal{O}_X)$
tel que $E'$ appartienne à $D_\QCoh(\mathcal{O}_X)$ et que
$\Hom(M, K) = 0$ pour tout $M$ de $D_\QCoh(\mathcal{O}_X)$. Voir
Catégories dérivées, lemme \ref{derived-lemma-right-adjoint}.
Considérons le triangle distingué
$$
E \to Rj_{U, *}E|_U \oplus Rj_{V, *}E|_V \to
Rj_{U \cap V, *}E|_{U \cap V} \to E[1]
$$
dans $D(\mathcal{O}_X)$ de
Cohomologie, lemme \ref{cohomology-lemma-exact-sequence-j-star}.
D'après Catégories dérivées, lemme \ref{derived-lemma-prepare-adjoint},
il suffit de construire les triangles distingués voulus
pour $Rj_{U, *}E|_U$, $Rj_{V, *}E|_V$ et
$Rj_{U \cap V, *}E|_{U \cap V}$. Ceci nous ramène à l'assertion
étudiée dans le paragraphe suivant.

\medskip\noindent
Soit $j : U \to X$ l'immersion ouverte correspondant à un ouvert quasi-compact $U$
pour lequel le lemme est vrai. Soit $L$ un objet de $D(\mathcal{O}_U)$.
Il existe alors un triangle distingué
$$
E' \to Rj_*L \to K \to E'[1]
$$
dans $D(\mathcal{O}_X)$
tel que $E'$ appartienne à $D_\QCoh(\mathcal{O}_X)$ et que
$\Hom(M, K) = 0$ pour tout $M$ de $D_\QCoh(\mathcal{O}_X)$.
Pour le voir, choisissons un triangle distingué
$$
L' \to L \to Q \to L'[1]
$$
dans $D(\mathcal{O}_U)$ tel que $L'$ appartienne à $D_\QCoh(\mathcal{O}_U)$
et que $\Hom(N, Q) = 0$ pour tout $N$ de $D_\QCoh(\mathcal{O}_U)$.
C'est possible, car l'énoncé de
Catégories dérivées, lemme \ref{derived-lemma-right-adjoint},
est une équivalence.
On obtient un triangle distingué
$$
Rj_*L' \to Rj_*L \to Rj_*Q \to Rj_*L'[1]
$$
dans $D(\mathcal{O}_X)$. Observons que $Rj_*L'$ appartient à $D_\QCoh(\mathcal{O}_X)$
d'après le lemme \ref{perfect-lemma-quasi-coherence-direct-image}.
D'autre part, si $M$ appartient à $D_\QCoh(\mathcal{O}_X)$, alors
$$
\Hom(M, Rj_*Q) = \Hom(Lj^*M, Q) = 0
$$
car $Lj^*M$ appartient à $D_\QCoh(\mathcal{O}_U)$ d'après le
lemme \ref{perfect-lemma-quasi-coherence-pullback}.
Ceci achève la démonstration.
\end{proof}

\begin{proof}[Deuxième démonstration]
L'adjoint existe d'après Catégories dérivées, proposition
\ref{derived-proposition-brown}. Les hypothèses sont vérifiées :
tout d'abord, $D_\QCoh(\mathcal{O}_X)$ possède des sommes directes,
et celles-ci commutent au foncteur d'inclusion
(lemme \ref{perfect-lemma-quasi-coherence-direct-sums}).
D'autre part, $D_\QCoh(\mathcal{O}_X)$
est engendrée par des objets compacts, car elle possède un générateur
parfait d'après le théorème \ref{perfect-theorem-bondal-van-den-Bergh},
et les objets parfaits sont compacts d'après la
proposition \ref{perfect-proposition-compact-is-perfect}.
\end{proof}

\begin{lemma}
\label{perfect-lemma-pushforward-better-coherator}
Soit $f : X \to Y$ un morphisme quasi-compact et quasi-séparé
de schémas. Si les adjoints à droite $DQ_X$ et $DQ_Y$
des foncteurs d'inclusion $D_\QCoh \to D$ existent pour $X$ et $Y$, alors
$$
Rf_* \circ DQ_X = DQ_Y \circ Rf_*
$$
\end{lemma}

\begin{proof}
L'énoncé a un sens, car $Rf_*$ envoie
$D_\QCoh(\mathcal{O}_X)$ dans $D_\QCoh(\mathcal{O}_Y)$ d'après le
lemme \ref{perfect-lemma-quasi-coherence-direct-image}.
Il est vrai, car $Lf^*$ envoie de même
$D_\QCoh(\mathcal{O}_Y)$ dans $D_\QCoh(\mathcal{O}_X)$
(lemme \ref{perfect-lemma-quasi-coherence-pullback}),
de sorte que $Rf_* \circ DQ_X$ et $DQ_Y \circ Rf_*$
sont tous deux adjoints à droite de $Lf^* : D_\QCoh(\mathcal{O}_Y) \to D(\mathcal{O}_X)$.
\end{proof}

\begin{remark}
\label{perfect-remark-explain-consequence}
Soit $X$ un schéma quasi-compact et quasi-séparé.
Soit $X = U \cup V$, avec $U$ et $V$ ouverts quasi-compacts.
D'après le lemme \ref{perfect-lemma-better-coherator}, les foncteurs
$DQ_X$, $DQ_U$, $DQ_V$, $DQ_{U \cap V}$ existent. De plus, il existe un
triangle distingué canonique
$$
DQ_X(K) \to Rj_{U, *}DQ_U(K|_U) \oplus Rj_{V, *}DQ_V(K|_V)
\to Rj_{U \cap V, *}DQ_{U \cap V}(K|_{U \cap V}) \to
$$
pour tout $K \in D(\mathcal{O}_X)$. Ceci résulte de l'application du
foncteur exact $DQ_X$ au triangle distingué de
Cohomologie, lemme \ref{cohomology-lemma-exact-sequence-j-star},
et de trois applications du lemme \ref{perfect-lemma-pushforward-better-coherator}.
\end{remark}

\begin{lemma}
\label{perfect-lemma-boundedness-better-coherator}
Soit $X$ un schéma quasi-compact et quasi-séparé. Le foncteur
$DQ_X$ du lemme \ref{perfect-lemma-better-coherator}
possède la propriété de bornitude suivante :
il existe un entier $N = N(X)$ tel que, si
$K$ appartient à $D(\mathcal{O}_X)$ et vérifie
$H^i(U, K) = 0$ pour $U$ ouvert affine de $X$ et $i \not \in [a, b]$, alors
les faisceaux de cohomologie $H^i(DQ_X(K))$ sont nuls pour
$i \not \in [a, b + N]$.
\end{lemma}

\begin{proof}
Nous utiliserons le principe de récurrence de
Cohomologie des schémas, lemme \ref{coherent-lemma-induction-principle}.

\medskip\noindent
Si $X$ est affine, le lemme est vrai avec $N = 0$, car alors
$RQ_X = DQ_X$ s'obtient en prenant le complexe de faisceaux
quasi-cohérents associé à $R\Gamma(X, K)$.
Voir les lemmes \ref{perfect-lemma-affine-compare-bounded} et \ref{perfect-lemma-affine-coherator}.

\medskip\noindent
Supposons que $U, V$ soient des ouverts quasi-compacts de $X$
et que le lemme soit vrai pour $U$, $V$ et $U \cap V$,
avec des entiers $N(U)$, $N(V)$ et $N(U \cap V)$.
Supposons maintenant que $K$ appartienne à $D(\mathcal{O}_X)$ et vérifie
$H^i(W, K) = 0$ pour tout ouvert affine $W \subset X$ et tout $i \not \in [a, b]$.
Alors $K|_U$, $K|_V$, $K|_{U \cap V}$ ont la même propriété.
On voit donc que $RQ_U(K|_U)$, $RQ_V(K|_V)$ et
$RQ_{U \cap V}(K|_{U \cap V})$ ont leurs faisceaux de cohomologie
nuls hors de l'intervalle $[a, b + \max(N(U), N(V), N(U \cap V))$.
Puisque les foncteurs $Rj_{U, *}$, $Rj_{V, *}$, $Rj_{U \cap V, *}$
sont de dimension cohomologique finie sur $D_\QCoh$ d'après le
lemme \ref{perfect-lemma-quasi-coherence-direct-image},
il existe un $N$ tel que
$Rj_{U, *}DQ_U(K|_U)$, $Rj_{V, *}DQ_V(K|_V)$ et
$Rj_{U \cap V, *}DQ_{U \cap V}(K|_{U \cap V})$ aient leurs faisceaux
de cohomologie nuls hors de l'intervalle $[a, b + N]$.
On conclut alors à l'aide du triangle distingué
de la remarque \ref{perfect-remark-explain-consequence}.
\end{proof}

\begin{example}
\label{perfect-example-inverse-limit-quasi-coherent}
Soit $X$ un schéma quasi-compact et quasi-séparé. Soit $(\mathcal{F}_n)$
un système projectif de faisceaux quasi-cohérents. Puisque $DQ_X$ est un adjoint à
droite, il commute aux produits, donc aux limites dérivées.
On a donc
$$
DQ_X(R\lim \mathcal{F}_n) =
(R\lim\text{ dans }D_\QCoh(\mathcal{O}_X))(\mathcal{F}_n)
$$
où la première limite $R\lim$ est prise dans $D(\mathcal{O}_X)$.
Posons $K = R\lim \mathcal{F}_n$ pour celle-ci.
Pour tout ouvert affine $U \subset X$, on a
$$
H^i(U, K) =
H^i(R\Gamma(U, R\lim \mathcal{F}_n)) =
H^i(R\lim R\Gamma(U, \mathcal{F}_n)) =
H^i(R\lim \Gamma(U, \mathcal{F}_n))
$$
car la cohomologie commute aux limites dérivées et les
faisceaux quasi-cohérents
$\mathcal{F}_n$ n'ont pas de cohomologie supérieure sur les affines.
Le calcul de $R\lim$ dans la catégorie des groupes
abéliens montre que $H^i(U, K) = 0$
sauf si $i \in [0, 1]$. On en déduit finalement que
la limite $R\lim$ dans $D_\QCoh(\mathcal{O}_X)$, qui est
$DQ_X(K)$ d'après ce qui précède, appartient à $D^b_\QCoh(\mathcal{O}_X)$
d'après le lemme \ref{perfect-lemma-boundedness-better-coherator}.
\end{example}













\section{Cohomologie et changement de base, IV}
\label{perfect-section-cohomology-and-base-change-perfect}

\noindent
Cette section poursuit la discussion de
Cohomologie des schémas, section
\ref{coherent-section-cohomology-and-base-change-perfect}.
Nous disposons d'abord d'une version très générale de la formule
de projection pour les morphismes quasi-compacts et quasi-séparés de schémas
et les complexes dont les faisceaux de cohomologie sont quasi-cohérents.

\begin{lemma}
\label{perfect-lemma-cohomology-base-change}
Soit $f : X \to Y$ un morphisme quasi-compact et quasi-séparé
de schémas. Pour $E$ dans $D_\QCoh(\mathcal{O}_X)$ et
$K$ dans $D_\QCoh(\mathcal{O}_Y)$, le morphisme
$$
Rf_*(E) \otimes_{\mathcal{O}_Y}^\mathbf{L} K
\longrightarrow
Rf_*(E \otimes_{\mathcal{O}_X}^\mathbf{L} Lf^*K)
$$
défini dans
Cohomologie, équation (\ref{cohomology-equation-projection-formula-map}),
est un isomorphisme.
\end{lemma}

\begin{proof}
Pour vérifier que ce morphisme est un isomorphisme, on peut travailler localement sur $Y$.
On se ramène donc au cas où $Y$ est affine.

\medskip\noindent
Supposons que $K = \bigoplus K_i$ soit une somme
directe de complexes $K_i \in D_\QCoh(\mathcal{O}_Y)$.
Si l'énoncé est vrai pour chaque $K_i$, il est vrai pour $K$.
En effet, les foncteurs $Lf^*$ et $\otimes^\mathbf{L}$ conservent
les sommes directes par construction, et $Rf_*$ commute aux sommes directes
(pour les complexes dont les faisceaux de cohomologie sont quasi-cohérents) d'après le
lemme \ref{perfect-lemma-quasi-coherence-pushforward-direct-sums}.
De plus, supposons que $K \to L \to M \to K[1]$ soit un triangle
distingué dans $D_\QCoh(Y)$. Si l'énoncé du
lemme est vrai pour deux des objets $K, L, M$, il est vrai pour le troisième
(car les foncteurs considérés sont des foncteurs exacts de catégories triangulées).

\medskip\noindent
Supposons $Y$ affine, et écrivons $Y = \Spec(A)$. Le foncteur
$\widetilde{\ } : D(A) \to D_\QCoh(\mathcal{O}_Y)$ est une équivalence
(lemme \ref{perfect-lemma-affine-compare-bounded}).
Soit $T$ la propriété, pour $K \in D(A)$, selon laquelle
l'énoncé du lemme est vrai pour $\widetilde{K}$.
La discussion ci-dessus et
Compléments sur l'algèbre, remarque \ref{more-algebra-remark-P-resolution},
montrent qu'il suffit de démontrer que $T$ est vraie pour $A[k]$.
Ceci achève la démonstration, car l'énoncé du lemme
est clair pour les décalés du faisceau structural.
\end{proof}

\begin{definition}
\label{perfect-definition-tor-independent}
Soit $S$ un schéma. Soient $X$, $Y$ des schémas sur $S$. On dit que $X$ et
$Y$ sont {\it Tor-indépendants sur $S$} si, pour tous $x \in X$ et
$y \in Y$ d'image commune $s \in S$, les anneaux
$\mathcal{O}_{X, x}$ et $\mathcal{O}_{Y, y}$ sont Tor-indépendants
sur $\mathcal{O}_{S, s}$ (voir
Compléments sur l'algèbre, définition \ref{more-algebra-definition-tor-independent}).
\end{definition}

\begin{lemma}
\label{perfect-lemma-tor-independent}
Soient $f : X \to S$ et $g : Y \to S$ des morphismes de schémas.
Les conditions suivantes sont équivalentes :
\begin{enumerate}
\item $X$ et $Y$ sont Tor-indépendants sur $S$, et
\item pour tous ouverts affines $U \subset X$, $V \subset Y$, $W \subset S$
avec $f(U) \subset W$ et $g(V) \subset W$, les anneaux
$\mathcal{O}_X(U)$ et $\mathcal{O}_Y(V)$ sont Tor-indépendants sur
$\mathcal{O}_S(W)$ ;
\item il existe un recouvrement ouvert affine $S = \bigcup W_i$ et,
pour chaque $i$, des recouvrements ouverts affines $f^{-1}(W_i) = \bigcup U_{ij}$
et $g^{-1}(W_i) = \bigcup V_{ik}$ tels que les anneaux
$\mathcal{O}_X(U_{ij})$ et $\mathcal{O}_Y(V_{ik})$ soient Tor-indépendants sur
$\mathcal{O}_S(W_i)$ pour tous $i, j, k$.
\end{enumerate}
\end{lemma}

\begin{proof}
Omise. Indication : utiliser Compléments sur l'algèbre, lemme
\ref{more-algebra-lemma-tor-independent}.
\end{proof}

\begin{lemma}
\label{perfect-lemma-flat-base-change-tor-independent}
Soient $X \to S$ et $Y \to S$ des morphismes de schémas. Soit $S' \to S$ un
morphisme de schémas et notons $X' = X \times_S S'$
et $Y' = Y \times_S S'$.
Si $X$ et $Y$ sont Tor-indépendants sur $S$ et si $S' \to S$ est plat,
alors $X'$ et $Y'$ sont Tor-indépendants sur $S'$.
\end{lemma}

\begin{proof}
Omise. Indication : utiliser le lemme \ref{perfect-lemma-tor-independent} et,
sur les ouverts affines, Compléments sur l'algèbre, lemme
\ref{more-algebra-lemma-flat-base-change-tor-independent}.
\end{proof}

\begin{lemma}
\label{perfect-lemma-compare-base-change}
Soit $g : S' \to S$ un morphisme de schémas.
Soit $f : X \to S$ quasi-compact et quasi-séparé.
Considérons le diagramme de changement de base
$$
\xymatrix{
X' \ar[r]_{g'} \ar[d]_{f'} &
X \ar[d]^f \\
S' \ar[r]^g &
S
}
$$
Si $X$ et $S'$ sont Tor-indépendants sur $S$, alors, pour tout
$E \in D_\QCoh(\mathcal{O}_X)$, la flèche canonique
$Lg^*Rf_*E \to Rf'_*L(g')^*E$ est un isomorphisme.
\end{lemma}

\begin{proof}
Pour tout objet $E$ de $D(\mathcal{O}_X)$, on peut utiliser
Cohomologie, remarque \ref{cohomology-remark-base-change}, pour obtenir un morphisme
canonique de changement de base $Lg^*Rf_*E \to Rf'_*L(g')^*E$. Pour vérifier qu'il
est un isomorphisme, on peut travailler localement sur $S'$. On peut donc supposer que
$g : S' \to S$ est un morphisme de schémas affines. En particulier, $g$
est affine, et il suffit de montrer que
$$
Rg_*Lg^*Rf_*E \to Rg_*Rf'_*L(g')^*E = Rf_*(Rg'_* L(g')^* E)
$$
est un isomorphisme, voir le lemme \ref{perfect-lemma-affine-morphism}
(et utiliser les lemmes \ref{perfect-lemma-quasi-coherence-pullback},
\ref{perfect-lemma-quasi-coherence-tensor-product} et
\ref{perfect-lemma-quasi-coherence-direct-image}
pour voir que les objets $Rf'_*L(g')^*E$ et $Lg^*Rf_*E$
ont des faisceaux de cohomologie quasi-cohérents). Notons que $g'$ est
également affine (Morphismes, lemme \ref{morphisms-lemma-base-change-affine}).
D'après le lemme \ref{perfect-lemma-affine-morphism-pull-push}, le morphisme s'identifie à
$$
Rf_*E \otimes_{\mathcal{O}_S}^\mathbf{L} g_*\mathcal{O}_{S'}
\longrightarrow
Rf_*(E \otimes_{\mathcal{O}_X}^\mathbf{L} g'_*\mathcal{O}_{X'})
$$
Observons que $g'_*\mathcal{O}_{X'} = f^*g_*\mathcal{O}_{S'}$ (par changement
de base affine, voir Cohomologie des schémas, lemme
\ref{coherent-lemma-affine-base-change}).
Ainsi, d'après le
lemme \ref{perfect-lemma-cohomology-base-change}, il suffit de démontrer que
$Lf^*g_*\mathcal{O}_{S'} = f^*g_*\mathcal{O}_{S'}$. Ceci résulte de notre
hypothèse que $X$ et $S'$ sont Tor-indépendants sur $S$. En effet, pour
le vérifier, on peut travailler localement sur $X$, donc aussi supposer $X$ affine.
Écrivons $X = \Spec(A)$, $S = \Spec(R)$ et $S' = \Spec(R')$. Notre hypothèse
implique que $A$ et $R'$ sont Tor-indépendants sur $R$
(Compléments sur l'algèbre, lemme \ref{more-algebra-lemma-tor-independent}), c'est-à-dire que
$\text{Tor}_i^R(A, R') = 0$ pour $i > 0$. Autrement dit,
$A \otimes_R^\mathbf{L} R' = A \otimes_R R'$, ce qui signifie exactement
que $Lf^*g_*\mathcal{O}_{S'} = f^*g_*\mathcal{O}_{S'}$
(utiliser le lemme \ref{perfect-lemma-quasi-coherence-pullback}).
\end{proof}

\noindent
Le lemme suivant sera utilisé dans le chapitre sur les complexes dualisants.

\begin{lemma}
\label{perfect-lemma-affine-morphism-and-hom-out-of-perfect}
Considérons un carré cartésien
$$
\xymatrix{
X' \ar[r]_{g'} \ar[d]_{f'} & X \ar[d]^f \\
S' \ar[r]^g & S
}
$$
de schémas quasi-compacts et quasi-séparés. Supposons que $g$ et $f$ soient
Tor-indépendants et que $S = \Spec(R)$, $S' = \Spec(R')$ soient affines. Pour
$M, K \in D(\mathcal{O}_X)$, le morphisme canonique
$$
R\Hom_X(M, K) \otimes^\mathbf{L}_R R'
\longrightarrow
R\Hom_{X'}(L(g')^*M, L(g')^*K)
$$
dans $D(R')$ est un isomorphisme dans les deux cas suivants :
\begin{enumerate}
\item $M \in D(\mathcal{O}_X)$ est parfait et $K \in D_\QCoh(X)$, ou
\item $M \in D(\mathcal{O}_X)$ est pseudo-cohérent,
$K \in D_\QCoh^+(X)$, et $R'$ est de dimension de Tor finie sur $R$.
\end{enumerate}
\end{lemma}

\begin{proof}
Il existe un morphisme canonique
$R\Hom_X(M, K) \to R\Hom_{X'}(L(g')^*M, L(g')^*K)$
dans $D(\Gamma(X, \mathcal{O}_X))$ entre complexes Hom globaux, voir
Cohomologie, section \ref{cohomology-section-global-RHom}.
Par restriction des scalaires, on peut le considérer comme un morphisme dans $D(R)$.
On utilise alors l'adjonction entre la restriction et
$- \otimes_R^\mathbf{L} R'$ pour obtenir le morphisme affiché dans le lemme.
Ce morphisme étant défini, il suffit de démontrer que c'est un isomorphisme
dans la catégorie dérivée des groupes abéliens.

\medskip\noindent
Le membre de droite est égal à
$$
R\Hom_X(M, R(g')_*L(g')^*K) =
R\Hom_X(M, K \otimes_{\mathcal{O}_X}^\mathbf{L} g'_*\mathcal{O}_{X'})
$$
d'après le lemme \ref{perfect-lemma-affine-morphism-pull-push}. Dans les deux cas, le complexe
$R\SheafHom(M, K)$ est un objet de $D_\QCoh(\mathcal{O}_X)$ d'après le
lemme \ref{perfect-lemma-quasi-coherence-internal-hom}. Il existe un morphisme naturel
$$
R\SheafHom(M, K) \otimes_{\mathcal{O}_X}^\mathbf{L} g'_*\mathcal{O}_{X'}
\longrightarrow
R\SheafHom(M, K \otimes_{\mathcal{O}_X}^\mathbf{L} g'_*\mathcal{O}_{X'})
$$
qui est un isomorphisme dans les deux cas d'après le
lemme \ref{perfect-lemma-internal-hom-evaluate-tensor-isomorphism}.
Pour voir que ce lemme s'applique dans le cas (2), notons que
$g'_*\mathcal{O}_{X'} = Rg'_*\mathcal{O}_{X'} =
Lf^*g_*\mathcal{O}_{S'}$, la deuxième égalité résultant du
lemme \ref{perfect-lemma-compare-base-change}.
En utilisant le lemme \ref{perfect-lemma-tor-dimension-affine} et
Cohomologie, lemme \ref{cohomology-lemma-tor-amplitude-pullback},
on en déduit que $g'_*\mathcal{O}_{X'}$ est de dimension de Tor finie.
Ainsi, dans les deux cas, en remplaçant $K$ par $R\SheafHom(M, K)$, on se ramène
à démontrer que
$$
R\Gamma(X, K) \otimes^\mathbf{L}_A A' \longrightarrow
R\Gamma(X, K \otimes^\mathbf{L}_{\mathcal{O}_X} g'_*\mathcal{O}_{X'})
$$
est un isomorphisme.
Notons que le membre de gauche est égal à $R\Gamma(X', L(g')^*K)$
d'après le lemme \ref{perfect-lemma-affine-morphism-pull-push}.
Le résultat découle donc du
lemme \ref{perfect-lemma-compare-base-change}.
\end{proof}

\begin{remark}
\label{perfect-remark-multiplication-map}
Reprenons les notations du lemme \ref{perfect-lemma-affine-morphism-and-hom-out-of-perfect}.
Le diagramme
$$
\xymatrix{
R\Hom_X(M, Rg'_*L) \otimes_R^\mathbf{L} R' \ar[r] \ar[d]_\mu &
R\Hom_{X'}(L(g')^*M, L(g')^*Rg'_*L) \ar[d]^a \\
R\Hom_X(M, R(g')_*L) \ar@{=}[r] &
R\Hom_{X'}(L(g')^*M, L)
}
$$
est commutatif ; la flèche horizontale supérieure est le morphisme du lemme,
$\mu$ est le morphisme de multiplication, et $a$ provient du morphisme d'adjonction
$L(g')^*Rg'_*L \to L$. Le morphisme de multiplication est le morphisme d'adjonction
$K' \otimes_R^\mathbf{L} R' \to K'$ pour tout $K' \in D(R')$.
\end{remark}

\begin{lemma}
\label{perfect-lemma-tor-independence-and-tor-amplitude}
Considérons un carré cartésien de schémas
$$
\xymatrix{
X' \ar[r]_{g'} \ar[d]_{f'} & X \ar[d]^f \\
S' \ar[r]^g & S
}
$$
Supposons que $g$ et $f$ soient Tor-indépendants.
\begin{enumerate}
\item Si $E \in D(\mathcal{O}_X)$ a une amplitude de Tor
contenue dans $[a, b]$ en tant que complexe de $f^{-1}\mathcal{O}_S$-modules,
alors $L(g')^*E$ a une amplitude de Tor
contenue dans $[a, b]$ en tant que complexe de $f^{-1}\mathcal{O}_{S'}$-modules.
\item Si $\mathcal{G}$ est un $\mathcal{O}_X$-module plat
sur $S$, alors $L(g')^*\mathcal{G} = (g')^*\mathcal{G}$.
\end{enumerate}
\end{lemma}

\begin{proof}
On peut calculer la dimension de Tor sur les fibres, voir
Cohomologie, lemme \ref{cohomology-lemma-tor-amplitude-stalk}.
Si $x' \in X'$ a pour image $x \in X$, alors
$$
(L(g')^*E)_{x'} =
E_x \otimes_{\mathcal{O}_{X, x}}^\mathbf{L} \mathcal{O}_{X', x'}
$$
Soient $s' \in S'$ et $s \in S$ les images de $x'$ et $x$.
Puisque $X$ et $S'$ sont Tor-indépendants sur $S$, on peut appliquer
Compléments sur l'algèbre, lemme \ref{more-algebra-lemma-base-change-comparison},
pour voir que le membre de droite de la formule affichée est égal à
$E_x \otimes_{\mathcal{O}_{S, s}}^\mathbf{L} \mathcal{O}_{S', s'}$
dans $D(\mathcal{O}_{S', s'})$.
L'assertion (1) résulte donc de
Compléments sur l'algèbre, lemme \ref{more-algebra-lemma-pull-tor-amplitude}.
Pour voir (2), observons que la platitude de $\mathcal{G}$ équivaut à
la condition que $\mathcal{G}[0]$ ait une amplitude de Tor contenue dans $[0, 0]$.
On conclut en appliquant (1).
\end{proof}

\begin{lemma}
\label{perfect-lemma-compare-base-change-closed-immersion}
Considérons un diagramme cartésien de schémas
$$
\xymatrix{
Z' \ar[r]_{i'} \ar[d]_g & X' \ar[d]^f \\
Z \ar[r]^i & X
}
$$
où $i$ est une immersion fermée. Si $Z$ et $X'$ sont
Tor-indépendants sur $X$, alors $Ri'_* \circ Lg^* = Lf^* \circ Ri_*$
comme foncteurs $D(\mathcal{O}_Z) \to D(\mathcal{O}_{X'})$.
\end{lemma}

\begin{proof}
Notons que le lemme doit être vrai pour tout $K \in D(\mathcal{O}_Z)$.
Observons que $i_*$ et $i'_*$ sont des foncteurs exacts ; par conséquent,
$Ri_*$ et $Ri'_*$ se calculent en appliquant $i_*$ et $i'_*$
à des représentants quelconques. Ainsi le morphisme de changement de base
$$
Lf^*(Ri_*(K)) \longrightarrow Ri'_*(Lg^*(K))
$$
sur les fibres en un point $z' \in Z'$ d'image $z \in Z$ est donné par
$$
K_z \otimes_{\mathcal{O}_{X, z}}^\mathbf{L} \mathcal{O}_{X', z'}
\longrightarrow
K_z \otimes_{\mathcal{O}_{Z, z}}^\mathbf{L} \mathcal{O}_{Z', z'}
$$
Ce morphisme est un isomorphisme d'après
Compléments sur l'algèbre, lemme \ref{more-algebra-lemma-base-change-comparison},
et la Tor-indépendance supposée.
\end{proof}








\section{Formule de K\"unneth, II}
\label{perfect-section-kunneth}

\noindent
Pour le cas où la base est un corps, voir
Variétés, section \ref{varieties-section-kunneth}.
Considérons un diagramme cartésien de schémas
$$
\xymatrix{
& X \times_S Y \ar[ld]^p \ar[rd]_q \ar[dd]^f \\
X \ar[rd]_a & & Y \ar[ld]^b \\
& S
}
$$
Soient $K \in D(\mathcal{O}_X)$ et $M \in D(\mathcal{O}_Y)$.
Il existe un morphisme canonique
\begin{equation}
\label{perfect-equation-kunneth}
Ra_*K \otimes_{\mathcal{O}_S}^\mathbf{L} Rb_*M
\longrightarrow
Rf_*(Lp^*K \otimes_{\mathcal{O}_{X \times_S Y}}^\mathbf{L} Lq^*M)
\end{equation}
En effet, on peut utiliser les morphismes
$Ra_*K \to Ra_*Rp_* Lp^*K = Rf_*Lp^*K$ et
$Rb_*M \to Rb_*Rq_* Lq^*M = Rf_*Lq^*M$, puis le cup-produit
relatif (Cohomologie, remarque \ref{cohomology-remark-cup-product}).

\medskip\noindent
Posons $A = \Gamma(S, \mathcal{O}_S)$. Il existe un morphisme de K\"unneth global
\begin{equation}
\label{perfect-equation-kunneth-global}
R\Gamma(X, K) \otimes_A^\mathbf{L} R\Gamma(Y, M)
\longrightarrow
R\Gamma(X \times_S Y,
Lp^*K \otimes_{\mathcal{O}_{X \times_S Y}}^\mathbf{L} Lq^*M)
\end{equation}
dans $D(A)$. Ce morphisme se construit à l'aide des morphismes de rappel
$R\Gamma(X, K) \to R\Gamma(X \times_S Y, Lp^*K)$ et
$R\Gamma(Y, M) \to R\Gamma(X \times_S Y, Lq^*M)$ et
du cup-produit construit dans
Cohomologie, section \ref{cohomology-section-cup-product}.

\begin{lemma}
\label{perfect-lemma-kunneth}
Dans la situation ci-dessus, si $a$ et $b$ sont quasi-compacts et quasi-séparés
et si $X$ et $Y$ sont Tor-indépendants sur $S$, alors (\ref{perfect-equation-kunneth})
est un isomorphisme pour $K \in D_\QCoh(\mathcal{O}_X)$ et
$M \in D_\QCoh(\mathcal{O}_Y)$. Si de plus $S = \Spec(A)$ est affine,
alors le morphisme (\ref{perfect-equation-kunneth-global}) est un isomorphisme.
\end{lemma}

\begin{proof}[Première démonstration]
Cela résulte de la suite d'isomorphismes suivante
\begin{align*}
Rf_*(Lp^*K \otimes_{\mathcal{O}_{X \times_S Y}}^\mathbf{L} Lq^*M)
& =
Ra_*Rp_*(Lp^*K \otimes_{\mathcal{O}_{X \times_S Y}}^\mathbf{L} Lq^*M) \\
& =
Ra_*(K \otimes_{\mathcal{O}_X}^\mathbf{L} Rp_*Lq^*M) \\
& =
Ra_*(K \otimes_{\mathcal{O}_X}^\mathbf{L} La^*Rb_*M) \\
& =
Ra_*K \otimes_{\mathcal{O}_S}^\mathbf{L} Rb_*M
\end{align*}
La première égalité vaut parce que $f = a \circ p$. La deuxième égalité
résulte du lemme \ref{perfect-lemma-cohomology-base-change}. La troisième, du
lemme \ref{perfect-lemma-compare-base-change}. La quatrième, du
lemme \ref{perfect-lemma-cohomology-base-change}.
Nous omettons la vérification que la composée de ces isomorphismes
est le morphisme (\ref{perfect-equation-kunneth}).
Si $S$ est affine, alors la source et le but de la flèche
(\ref{perfect-equation-kunneth-global}) s'obtiennent en appliquant
$R\Gamma(S, -)$ à la source et au but de (\ref{perfect-equation-kunneth}),
ce qui donne la dernière assertion ; nous omettons les détails.
\end{proof}

\begin{proof}[Deuxième démonstration]
La construction de la flèche (\ref{perfect-equation-kunneth}) est compatible
avec la restriction aux sous-schémas ouverts de $S$, comme il ressort
immédiatement de la construction du cup-produit relatif. Il suffit donc de montrer
que (\ref{perfect-equation-kunneth}) est un isomorphisme lorsque $S$ est affine.

\medskip\noindent
Supposons $S = \Spec(A)$ affine. Par Leray, on a
$R\Gamma(S, Rf_*K) = R\Gamma(X, K)$, et de même dans les autres cas.
D'après Cohomologie, lemme \ref{cohomology-lemma-compose-cup-product},
le morphisme (\ref{perfect-equation-kunneth}) induit le morphisme
(\ref{perfect-equation-kunneth-global}) après application de $R\Gamma(S, -)$.
D'autre part, d'après les lemmes \ref{perfect-lemma-quasi-coherence-direct-image} et
\ref{perfect-lemma-quasi-coherence-tensor-product},
la source et le but du morphisme
(\ref{perfect-equation-kunneth}) appartiennent à $D_\QCoh(\mathcal{O}_S)$.
Ainsi, d'après le lemme \ref{perfect-lemma-affine-compare-bounded}, il suffit de montrer que
(\ref{perfect-equation-kunneth-global}) est un isomorphisme.

\medskip\noindent
Supposons $S = \Spec(A)$, $X = \Spec(B)$ et $Y = \Spec(C)$ tous affines.
Nous utiliserons le lemme \ref{perfect-lemma-affine-compare-bounded} sans autre mention.
Dans ce cas, on peut choisir un complexe K-plat $K^\bullet$ de $B$-modules
dont les termes sont plats et tel que $K$ soit représenté par $\widetilde{K}^\bullet$.
De même, on peut choisir un complexe K-plat $M^\bullet$ de $C$-modules
dont les termes sont plats et tel que $M$ soit représenté par $\widetilde{M}^\bullet$.
Voir Compléments sur l'algèbre, lemme \ref{more-algebra-lemma-K-flat-resolution}.
Alors $\widetilde{K}^\bullet$ est un complexe K-plat de $\mathcal{O}_X$-modules,
et de même pour $\widetilde{M}^\bullet$ ; voir
le lemme \ref{perfect-lemma-affine-K-flat}. Ainsi,
$La^*K$ est représenté par
$$
a^*\widetilde{K}^\bullet = \widetilde{K^\bullet \otimes_A C}
$$
et de même pour $Lb^*M$. À son tour, c'est un complexe K-plat
de $\mathcal{O}_{X \times_S Y}$-modules d'après le lemme cité ci-dessus
et Compléments sur l'algèbre, lemme \ref{more-algebra-lemma-base-change-K-flat}.
On voit donc finalement que le complexe de
$\mathcal{O}_{X \times_S Y}$-modules associé à
$$
\text{Tot}((K^\bullet \otimes_A C) \otimes_{B \otimes_A C}
(B \otimes_A M^\bullet)) =
\text{Tot}(K^\bullet \otimes_A M^\bullet)
$$
représente $La^*K \otimes_{\mathcal{O}_{X \times_S Y}}^\mathbf{L} Lb^*M$
dans la catégorie dérivée de $X \times_S Y$. En prenant les sections globales,
on obtient $\text{Tot}(K^\bullet \otimes_A M^\bullet)$, qui est bien entendu
également le complexe représentant
$R\Gamma(X, K) \otimes_A^\mathbf{L} R\Gamma(Y, M)$.
Le fait que l'isomorphisme soit donné par le cup-produit résulte de la
relation entre le véritable cup-produit et le cup-produit naïf
dans Cohomologie, section \ref{cohomology-section-cup-product}
(et en particulier de
Cohomologie, lemme \ref{cohomology-lemma-cup-compatible-with-naive},
ainsi que de la discussion qui le suit).

\medskip\noindent
Supposons $S = \Spec(A)$ et $Y$ affines. Nous allons utiliser le principe
d'induction de
Cohomologie des schémas, lemme \ref{coherent-lemma-induction-principle},
pour démontrer l'assertion. Pour cela, il suffit de montrer ceci :
si $X = U \cup V$ est un recouvrement ouvert avec $U$ et $V$ quasi-compacts
et si le morphisme (\ref{perfect-equation-kunneth-global})
$$
R\Gamma(U, K) \otimes_A^\mathbf{L} R\Gamma(Y, M)
\longrightarrow
R\Gamma(U \times_S Y,
Lp^*K \otimes_{\mathcal{O}_{X \times_S Y}}^\mathbf{L} Lq^*M)
$$
pour $U$ et $Y$ sur $S$, le morphisme (\ref{perfect-equation-kunneth-global})
$$
R\Gamma(V, K) \otimes_A^\mathbf{L} R\Gamma(Y, M)
\longrightarrow
R\Gamma(V \times_S Y,
Lp^*K \otimes_{\mathcal{O}_{X \times_S Y}}^\mathbf{L} Lq^*M)
$$
pour $V$ et $Y$ sur $S$, et le morphisme (\ref{perfect-equation-kunneth-global})
$$
R\Gamma(U \cap V, K) \otimes_A^\mathbf{L} R\Gamma(Y, M)
\longrightarrow
R\Gamma((U \cap V) \times_S Y,
Lp^*K \otimes_{\mathcal{O}_{X \times_S Y}}^\mathbf{L} Lq^*M)
$$
pour $U \cap V$ et $Y$ sur $S$ sont des isomorphismes, alors le morphisme
(\ref{perfect-equation-kunneth-global}) pour $X$ et $Y$ sur $S$ en est aussi un.
Cependant, d'après Cohomologie, lemme \ref{cohomology-lemma-mayer-vietoris-cup},
ces morphismes s'insèrent dans un morphisme de triangles distingués dont
(\ref{perfect-equation-kunneth-global}) est la dernière branche ; on conclut donc
par Catégories dérivées, lemme \ref{derived-lemma-third-isomorphism-triangle}.

\medskip\noindent
Supposons $S = \Spec(A)$ affine. Pour achever la démonstration, on peut appliquer
le principe d'induction de
Cohomologie des schémas, lemme \ref{coherent-lemma-induction-principle},
à $Y$. En effet, ce qui précède montre déjà que notre morphisme est
un isomorphisme lorsque $Y$ est affine. La suite de l'argument est
exactement la même que dans le paragraphe précédent, en échangeant les rôles de
$X$ et $Y$.
\end{proof}

\begin{lemma}
\label{perfect-lemma-cohomology-de-rham-base-change}
Soit $a : X \to S$ un morphisme quasi-compact et quasi-séparé
de schémas. Soit $\mathcal{F}^\bullet$ un complexe localement borné
de $a^{-1}\mathcal{O}_S$-modules. Supposons que, pour tout $n \in \mathbf{Z}$,
le faisceau $\mathcal{F}^n$ soit un $a^{-1}\mathcal{O}_S$-module plat et que
$\mathcal{F}^n$ possède une structure de $\mathcal{O}_X$-module quasi-cohérent
compatible avec la structure de $a^{-1}\mathcal{O}_S$-module donnée (mais les
différentielles du complexe $\mathcal{F}^\bullet$ ne sont pas nécessairement
$\mathcal{O}_X$-linéaires). Alors les assertions suivantes sont vraies :
\begin{enumerate}
\item $Ra_*\mathcal{F}^\bullet$ est localement borné,
\item $Ra_*\mathcal{F}^\bullet$ appartient à $D_\QCoh(\mathcal{O}_S)$,
\item $Ra_*\mathcal{F}^\bullet$ est localement de dimension de Tor finie,
\item $\mathcal{G} \otimes_{\mathcal{O}_S}^\mathbf{L} Ra_*\mathcal{F}^\bullet =
Ra_*(a^{-1}\mathcal{G} \otimes_{a^{-1}\mathcal{O}_S} \mathcal{F}^\bullet)$
pour $\mathcal{G} \in \QCoh(\mathcal{O}_S)$, et
\item $K \otimes_{\mathcal{O}_S}^\mathbf{L} Ra_*\mathcal{F}^\bullet =
Ra_*(a^{-1}K \otimes_{a^{-1}\mathcal{O}_S}^\mathbf{L} \mathcal{F}^\bullet)$
pour $K \in D_\QCoh(\mathcal{O}_S)$.
\end{enumerate}
\end{lemma}

\begin{proof}
Les parties (1), (2) et (3) sont locales sur $S$ ; on peut donc supposer, et l'on suppose,
que $S$ est affine. Puisque $a$ est quasi-compact, on en déduit que $X$ est quasi-compact.
Comme $\mathcal{F}^\bullet$ est localement borné, on en déduit que
$\mathcal{F}^\bullet$ est borné.

\medskip\noindent
Pour (1) et (2), on peut utiliser la première suite spectrale
$R^pa_*\mathcal{F}^q \Rightarrow R^{p + q}a_*\mathcal{F}^\bullet$ de
Catégories dérivées, lemme \ref{derived-lemma-two-ss-complex-functor}.
En combinant Cohomologie des schémas, lemme
\ref{coherent-lemma-quasi-coherence-higher-direct-images},
et Homologie, lemme \ref{homology-lemma-biregular-ss-converges},
on conclut.

\medskip\noindent
Démontrons (3) au moyen du principe d'induction de
Cohomologie des schémas, lemme \ref{coherent-lemma-induction-principle}.
Plus précisément, pour un ouvert quasi-compact $U$ de $X$, considérons la
condition que $R(a|_U)_*(\mathcal{F}^\bullet|_U)$ soit de
dimension de Tor finie. Si $U, V$ sont des ouverts quasi-compacts de
$X$, alors on dispose d'un triangle distingué de Mayer--Vietoris relatif
$$
R(a|_{U \cup V})_*\mathcal{F}^\bullet|_{U \cup V} \to
R(a|_U)_*\mathcal{F}^\bullet|_U \oplus
R(a|_V)_*\mathcal{F}^\bullet|_V \to
R(a|_{U \cap V})_*\mathcal{F}^\bullet|_{U \cap V} \to
$$
d'après Cohomologie, lemme \ref{cohomology-lemma-unbounded-relative-mayer-vietoris}.
Le comportement de l'amplitude de Tor dans les triangles distingués
(voir Cohomologie, lemme \ref{cohomology-lemma-cone-tor-amplitude})
montre que, si le résultat est connu pour $U$, $V$ et $U \cap V$, alors
il vaut pour $U \cup V$. Cela nous ramène au cas où
$X$ est affine. Dans ce cas, on a
$$
Ra_*\mathcal{F}^\bullet = a_*\mathcal{F}^\bullet
$$
d'après le lemme d'acyclicité de Leray
(Catégories dérivées, lemme \ref{derived-lemma-leray-acyclicity})
et l'annulation des images directes supérieures des modules quasi-cohérents
par un morphisme affine
(Cohomologie des schémas, lemme \ref{coherent-lemma-relative-affine-vanishing}).
Puisque $\mathcal{F}^n$ est plat sur $S$ par hypothèse et que $X$ est affine, les modules
$a_*\mathcal{F}^n$ sont plats pour tout $n$. Par conséquent, $a_*\mathcal{F}^\bullet$
est un complexe borné de $\mathcal{O}_S$-modules plats, donc il est de
dimension de Tor finie.

\medskip\noindent
Démonstration de la partie (5). Notons
$a' : (X, a^{-1}\mathcal{O}_S) \to (S, \mathcal{O}_S)$
le morphisme plat évident d'espaces annelés. La partie (5) affirme que
$$
K \otimes_{\mathcal{O}_S}^\mathbf{L} Ra'_*\mathcal{F}^\bullet =
Ra'_*(L(a')^*K \otimes_{a^{-1}\mathcal{O}_S}^\mathbf{L}
\mathcal{F}^\bullet)
$$
Ainsi,
Cohomologie, équation (\ref{cohomology-equation-projection-formula-map}),
fournit un morphisme fonctoriel du membre de gauche vers celui de droite, et nous voulons
montrer que ce morphisme est un isomorphisme.
Cette question est locale sur $S$ ; on peut donc supposer, et l'on suppose, que $S$
est affine. Le reste de la démonstration est {\it exactement} identique à la
démonstration du lemme \ref{perfect-lemma-cohomology-base-change}, sauf qu'il faut
montrer que le foncteur
$K \mapsto Ra'_*(L(a')^*K \otimes_{a^{-1}\mathcal{O}_S}^\mathbf{L}
\mathcal{F}^\bullet)$ commute aux sommes directes.
C'est ici que nous utilisons le fait que $\mathcal{F}^n$ possède une structure
de $\mathcal{O}_X$-module quasi-cohérent. En effet, observons que
$K \mapsto L(a')^*K
\otimes_{a^{-1}\mathcal{O}_S}^\mathbf{L} \mathcal{F}^\bullet$
commute aux sommes directes arbitraires. Ensuite, si
$\mathcal{F}^\bullet$ consiste en un seul $\mathcal{O}_X$-module
quasi-cohérent $\mathcal{F}^\bullet = \mathcal{F}^n[-n]$,
alors $L(a')^*G
\otimes_{a^{-1}\mathcal{O}_S}^\mathbf{L} \mathcal{F}^\bullet =
La^*K \otimes_{\mathcal{O}_X}^\mathbf{L} \mathcal{F}^n[-n]$ ; voir
Cohomologie, lemme \ref{cohomology-lemma-variant-derived-pullback}.
Dans ce cas, la commutation aux sommes directes résulte donc du
lemme \ref{perfect-lemma-quasi-coherence-pushforward-direct-sums}.
Dans le cas général, puisque $S$ est affine (donc $X$ quasi-compact)
et que $\mathcal{F}^\bullet$ est localement borné, on voit que
$$
\mathcal{F}^\bullet = (\mathcal{F}^a \to \ldots \to \mathcal{F}^b)
$$
est borné. En raisonnant par récurrence sur $b - a$ et en considérant le
triangle distingué
$$
\mathcal{F}^b[-b] \to (\mathcal{F}^a \to \ldots \to \mathcal{F}^b)
\to (\mathcal{F}^a \to \ldots \to \mathcal{F}^{b - 1}) \to
\mathcal{F}^b[-b + 1]
$$
on achève la démonstration de cette partie. Certains détails sont omis.

\medskip\noindent
Démonstration de la partie (4). Soit $a' : (X, a^{-1}\mathcal{O}_S) \to (S, \mathcal{O}_S)$
comme ci-dessus. Puisque $\mathcal{F}^\bullet$ est un complexe localement borné
de $a^{-1}\mathcal{O}_S$-modules plats, on voit que le complexe
$a^{-1}\mathcal{G} \otimes_{a^{-1}\mathcal{O}_S} \mathcal{F}^\bullet$
représente $L(a')^*\mathcal{G}
\otimes_{a^{-1}\mathcal{O}_S}^\mathbf{L}
\mathcal{F}^\bullet$ dans $D(a^{-1}\mathcal{O}_S)$. L'assertion (4)
résulte donc de (5).
\end{proof}

\begin{lemma}
\label{perfect-lemma-K-flat}
Soit $f : X \to Y$ un morphisme de schémas, avec $Y = \Spec(A)$ affine.
Soit $\mathcal{U} : X = \bigcup_{i \in I} U_i$ un recouvrement ouvert affine fini
tel que toutes les intersections finies
$U_{i_0 \ldots i_p} = U_{i_0} \cap \ldots \cap U_{i_p}$
soient affines. Soit $\mathcal{F}^\bullet$ un complexe borné de
$f^{-1}\mathcal{O}_Y$-modules. Supposons que, pour tout $n \in \mathbf{Z}$,
le faisceau $\mathcal{F}^n$ soit un $f^{-1}\mathcal{O}_Y$-module plat et que
$\mathcal{F}^n$ possède une structure de $\mathcal{O}_X$-module quasi-cohérent
compatible avec la structure de $p^{-1}\mathcal{O}_Y$-module donnée (mais les
différentielles du complexe $\mathcal{F}^\bullet$ ne sont pas nécessairement
$\mathcal{O}_X$-linéaires). Alors le complexe
$\text{Tot}(\check{\mathcal{C}}^\bullet(\mathcal{U}, \mathcal{F}^\bullet))$
est K-plat comme complexe de $A$-modules.
\end{lemma}

\begin{proof}
On peut écrire
$$
\mathcal{F}^\bullet = (\mathcal{F}^a \to \ldots \to \mathcal{F}^b)
$$
En raisonnant par récurrence sur $b - a$ et en considérant le triangle distingué
$$
\mathcal{F}^b[-b] \to (\mathcal{F}^a \to \ldots \to \mathcal{F}^b)
\to (\mathcal{F}^a \to \ldots \to \mathcal{F}^{b - 1}) \to
\mathcal{F}^b[-b + 1]
$$
et en utilisant
Compléments sur l'algèbre, lemme \ref{more-algebra-lemma-K-flat-two-out-of-three},
on se ramène au cas où $\mathcal{F}^\bullet$ consiste en un seul
$\mathcal{O}_X$-module quasi-cohérent $\mathcal{F}$
placé en degré $0$. Dans ce cas, le complexe de {\v C}ech
de $\mathcal{F}$ relativement à $\mathcal{U}$ est homotopiquement équivalent au
complexe de {\v C}ech alterné ; voir
Cohomologie, lemme \ref{cohomology-lemma-alternating-usual}.
Comme $U_{i_0 \ldots i_p}$ est toujours affine, on voit que
$\mathcal{F}(U_{i_0 \ldots i_p})$ est plat sur $A$.
Par conséquent,
$\check{\mathcal{C}}_{alt}^\bullet(\mathcal{U}, \mathcal{F})$
est un complexe borné de $A$-modules plats, donc K-plat
d'après Compléments sur l'algèbre, lemme
\ref{more-algebra-lemma-derived-tor-quasi-isomorphism}.
\end{proof}

\noindent
Soient $X, Y, S, a, b, p, q, f$ comme dans l'introduction de cette section.
Soit $\mathcal{F}$ un $\mathcal{O}_X$-module.
Soit $\mathcal{G}$ un $\mathcal{O}_Y$-module.
Posons $A = \Gamma(S, \mathcal{O}_S)$.
Considérons le morphisme
\begin{equation}
\label{perfect-equation-kunneth-single-sheaves}
R\Gamma(X, \mathcal{F})
\otimes_A^\mathbf{L}
R\Gamma(Y, \mathcal{G})
\longrightarrow
R\Gamma(X \times_S Y,
p^*\mathcal{F}
\otimes_{\mathcal{O}_{X \times_S Y}}
q^*\mathcal{G})
\end{equation}
dans $D(A)$. Ce morphisme se construit à l'aide des morphismes de rappel
$R\Gamma(X, \mathcal{F}) \to R\Gamma(X \times_S Y, p^*\mathcal{F})$ et
$R\Gamma(Y, \mathcal{G}) \to R\Gamma(X \times_S Y, q^*\mathcal{G})$,
du cup-produit construit dans
Cohomologie, section \ref{cohomology-section-cup-product}, et
du morphisme canonique
$p^*\mathcal{F}
\otimes_{\mathcal{O}_{X \times_S Y}}^\mathbf{L}
q^*\mathcal{G}
\to
p^*\mathcal{F}
\otimes_{\mathcal{O}_{X \times_S Y}}
q^*\mathcal{G}$.

\begin{lemma}
\label{perfect-lemma-kunneth-single-sheaf}
Dans la situation ci-dessus, le morphisme (\ref{perfect-equation-kunneth-single-sheaves}) est un
isomorphisme si $S$ est affine, si $\mathcal{F}$ et $\mathcal{G}$ sont plats sur $S$ et
quasi-cohérents, et si $X$ et $Y$ sont quasi-compacts à diagonale affine.
\end{lemma}

\begin{proof}
Nous recommandons vivement au lecteur de lire d'abord la démonstration de
Variétés, lemme \ref{varieties-lemma-kunneth}.
Choisissons des recouvrements ouverts affines finis
$\mathcal{U} : X = \bigcup_{i \in I} U_i$ et
$\mathcal{V} : Y = \bigcup_{j \in J} V_j$.
Ils déterminent un recouvrement ouvert affine
$\mathcal{W} : X \times_S Y = \bigcup_{(i, j) \in I \times J} U_i \times_S V_j$.
Notons que $\mathcal{W}$ raffine
$\text{pr}_1^{-1}\mathcal{U}$ et $\text{pr}_2^{-1}\mathcal{V}$.
Ainsi, d'après la discussion de Cohomologie, section
\ref{cohomology-section-cech-cohomology-of-complexes},
on obtient des morphismes
$$
\check{\mathcal{C}}^\bullet(\mathcal{U}, \mathcal{F})
\to
\check{\mathcal{C}}^\bullet(\mathcal{W}, p^*\mathcal{F})
\quad\text{et}\quad
\check{\mathcal{C}}^\bullet(\mathcal{V}, \mathcal{G})
\to
\check{\mathcal{C}}^\bullet(\mathcal{W}, q^*\mathcal{G})
$$
bien définis à homotopie près et compatibles avec les morphismes de rappel en cohomologie.
Dans Cohomologie, équation (\ref{cohomology-equation-needs-signs}),
nous avons construit un morphisme de complexes
$$
\text{Tot}(
\check{\mathcal{C}}^\bullet(\mathcal{W}, p^*\mathcal{F})
\otimes_A
\check{\mathcal{C}}^\bullet(\mathcal{W}, q^*\mathcal{G}))
\longrightarrow
\check{\mathcal{C}}^\bullet(\mathcal{W},
p^*\mathcal{F} \otimes_{\mathcal{O}_{X \times_S Y}}
q^*\mathcal{G})
$$
qui est compatible avec le cup-produit en cohomologie d'après
Cohomologie, lemme \ref{cohomology-lemma-diagrams-commute}.
En combinant ce qui précède, on obtient un morphisme de complexes
\begin{equation}
\label{perfect-equation-kunneth-on-cech}
\text{Tot}(
\check{\mathcal{C}}^\bullet(\mathcal{U}, \mathcal{F})
\otimes_A
\check{\mathcal{C}}^\bullet(\mathcal{V}, \mathcal{G}))
\to
\check{\mathcal{C}}^\bullet(\mathcal{W},
p^*\mathcal{F}
\otimes_{\mathcal{O}_{X \times_S Y}}
q^*\mathcal{G})
\end{equation}
Nous affirmons qu'il s'agit du morphisme de l'énoncé du lemme, c'est-à-dire que
la source et le but de cette flèche s'identifient à la source
et au but de (\ref{perfect-equation-kunneth-single-sheaves}). En effet, d'après
Cohomologie des schémas, lemme
\ref{coherent-lemma-quasi-coherent-affine-cohomology-zero},
et
Cohomologie, lemme \ref{cohomology-lemma-cech-complex-complex-computes},
les morphismes canoniques
$$
\check{\mathcal{C}}^\bullet(\mathcal{U}, \mathcal{F})
\to
R\Gamma(X, \mathcal{F}),
\quad
\check{\mathcal{C}}^\bullet(\mathcal{V}, \mathcal{G})
\to
R\Gamma(Y, \mathcal{G})
$$
et
$$
\check{\mathcal{C}}^\bullet(\mathcal{W},
p^*\mathcal{F} \otimes_{\mathcal{O}_{X \times_S Y}} q^*\mathcal{G})
\to
R\Gamma(X \times_S Y,
p^*\mathcal{F} \otimes_{\mathcal{O}_{X \times_S Y}}
q^*\mathcal{G})
$$
sont des isomorphismes. D'autre part, le complexe
$\check{\mathcal{C}}^\bullet(\mathcal{U}, \mathcal{F})$
est K-plat d'après le lemme \ref{perfect-lemma-K-flat}, et on en déduit que
$\text{Tot}(
\check{\mathcal{C}}^\bullet(\mathcal{U}, \mathcal{F})
\otimes_A
\check{\mathcal{C}}^\bullet(\mathcal{V}, \mathcal{G}))$
représente le produit tensoriel dérivé
$R\Gamma(X, \mathcal{F}) \otimes_A^\mathbf{L} R\Gamma(Y, \mathcal{G})$,
comme annoncé.

\medskip\noindent
Il reste à montrer que (\ref{perfect-equation-kunneth-on-cech})
est un quasi-isomorphisme. Nous allons le faire par décalage de dimension.
Posons $d(\mathcal{F}) = \max \{d \mid H^d(X, \mathcal{F}) \not = 0\}$.
Supposons $d(\mathcal{F}) > 0$. Posons $U = \coprod\nolimits_{i \in I} U_i$.
C'est un schéma affine puisque $I$ est fini. Notons
$j : U \to X$ le morphisme qui est l'inclusion $U_i \to X$
sur chaque $U_i$. Comme la diagonale de $X$ est affine, le morphisme
$j$ est affine ; voir
Morphismes, lemme \ref{morphisms-lemma-affine-permanence}.
Il s'ensuit que $\mathcal{F}' = j_*j^*\mathcal{F}$ est plat sur $S$ ; voir
Morphismes, lemme \ref{morphisms-lemma-pushforward-flat-affine}.
Il s'ensuit aussi que $d(\mathcal{F}') = 0$ en combinant
Cohomologie des schémas, lemmes
\ref{coherent-lemma-relative-affine-cohomology} et
\ref{coherent-lemma-quasi-coherent-affine-cohomology-zero}.
Pour tout $x \in X$, le morphisme $\mathcal{F}_x  \to \mathcal{F}'_x$
est l'inclusion d'un facteur direct : si $x \in U_i$,
alors $\mathcal{F}' \to (U_i \to X)_*\mathcal{F}|_{U_i}$
fournit une rétraction. On en déduit que
$\mathcal{F} \to \mathcal{F}'$ est injectif et que
$\mathcal{F}'' = \mathcal{F}'/\mathcal{F}$
est lui aussi plat sur $S$. La suite exacte courte
$0 \to \mathcal{F} \to \mathcal{F}' \to \mathcal{F}'' \to 0$
de $\mathcal{O}_X$-modules quasi-cohérents plats
donne une suite exacte courte de complexes
$$
0 \to
\text{Tot}(
\check{\mathcal{C}}^\bullet(\mathcal{U}, \mathcal{F})
\otimes_A
\check{\mathcal{C}}^\bullet(\mathcal{V}, \mathcal{G})) \to
\text{Tot}(
\check{\mathcal{C}}^\bullet(\mathcal{U}, \mathcal{F}')
\otimes_A
\check{\mathcal{C}}^\bullet(\mathcal{V}, \mathcal{G})) \to
\text{Tot}(
\check{\mathcal{C}}^\bullet(\mathcal{U}, \mathcal{F}'')
\otimes_A
\check{\mathcal{C}}^\bullet(\mathcal{V}, \mathcal{G})) \to 0
$$
et une suite exacte courte de complexes
$$
0 \to
\check{\mathcal{C}}^\bullet(\mathcal{W},
p^*\mathcal{F}
\otimes_{\mathcal{O}_{X \times_S Y}}
q^*\mathcal{G}) \to
\check{\mathcal{C}}^\bullet(\mathcal{W},
p^*\mathcal{F}'
\otimes_{\mathcal{O}_{X \times_S Y}}
q^*\mathcal{G}) \to
\check{\mathcal{C}}^\bullet(\mathcal{W},
p^*\mathcal{F}''
\otimes_{\mathcal{O}_{X \times_S Y}}
q^*\mathcal{G}) \to 0
$$
De plus, les morphismes (\ref{perfect-equation-kunneth-on-cech}) entre celles-ci
sont compatibles avec ces suites exactes courtes. Il suffit donc de montrer que
(\ref{perfect-equation-kunneth-on-cech})
est un isomorphisme pour $\mathcal{F}'$ et $\mathcal{F}''$. Enfin,
on a $d(\mathcal{F}'') < d(\mathcal{F})$.
On se ramène ainsi au cas $d(\mathcal{F}) = 0$.

\medskip\noindent
En raisonnant de la même manière pour $\mathcal{G}$, on voit que l'on
peut supposer que $\mathcal{F}$ et $\mathcal{G}$
n'ont de cohomologie non nulle qu'en degré $0$.
Observons que cela signifie que $\Gamma(X, \mathcal{F})$
est quasi-isomorphe au complexe $K$-plat
$\check{\mathcal{C}}^\bullet(\mathcal{U}, \mathcal{F})$
de $A$-modules situé en degrés $\geq 0$.
Il s'ensuit que $\Gamma(X, \mathcal{F})$ est un $A$-module plat
(car on peut calculer les Tor supérieurs contre ce module
en tensorisant par le complexe de \v{C}ech).
Soit $V \subset Y$ un ouvert affine. Considérons le recouvrement ouvert affine
$\mathcal{U}_V : X \times_S V = \bigcup_{i \in I} U_i \times_S V$.
Il est immédiat que
$$
\check{\mathcal{C}}^\bullet(\mathcal{U}, \mathcal{F})
\otimes_A \mathcal{G}(V) =
\check{\mathcal{C}}^\bullet(\mathcal{U}_V,
p^*\mathcal{F} \otimes_{\mathcal{O}_{X \times Y}}
q^*\mathcal{G})
$$
(égalité de complexes). Comme $\mathcal{G}(V)$ est plat
sur $A$, on voit que
$\Gamma(X, \mathcal{F}) \otimes_A \mathcal{G}(V) \to
\check{\mathcal{C}}^\bullet(\mathcal{U}, \mathcal{F})
\otimes_A \mathcal{G}(V)$ est un quasi-isomorphisme.
Puisque le faisceau associé à
$V \mapsto \check{\mathcal{C}}^\bullet(\mathcal{U}_V,
p^*\mathcal{F} \otimes_{\mathcal{O}_{X \times Y}}
q^*\mathcal{G})$ représente
$Rq_*(p^*\mathcal{F} \otimes_{\mathcal{O}_{X \times Y}} q^*\mathcal{G})$
d'après Cohomologie des schémas, lemme
\ref{coherent-lemma-separated-case-relative-cech},
on en déduit que
$$
Rq_*(p^*\mathcal{F} \otimes_{\mathcal{O}_{X \times Y}} q^*\mathcal{G})
\cong
\Gamma(X, \mathcal{F}) \otimes_A \mathcal{G}
$$
sur $Y$, où la notation du membre de droite désigne le module
$$
b^*\widetilde{\Gamma(X, \mathcal{F})} \otimes_{\mathcal{O}_Y} \mathcal{G}
$$
En utilisant la suite spectrale de Leray pour $q$, on obtient
$$
H^n(X \times_S Y, p^*\mathcal{F} \otimes_{\mathcal{O}_{X \times Y}}
q^*\mathcal{G}) =
H^n(Y,
b^*\widetilde{\Gamma(X, \mathcal{F})} \otimes_{\mathcal{O}_Y} \mathcal{G})
$$
En appliquant le lemme \ref{perfect-lemma-cohomology-base-change} au morphisme
$b : Y \to S = \Spec(A)$ et en utilisant le fait que $\Gamma(X, \mathcal{F})$
est plat sur $A$, on conclut que
$H^n(X \times_S Y, p^*\mathcal{F} \otimes_{\mathcal{O}_{X \times Y}}
q^*\mathcal{G})$ est nul pour $n > 0$ et isomorphe à
$H^0(X, \mathcal{F}) \otimes_A H^0(Y, \mathcal{G})$ pour $n = 0$.
Bien entendu, nous utilisons ici aussi le fait que $\mathcal{G}$ n'a de
cohomologie qu'en degré $0$.
Cela achève la démonstration (à ceci près qu'il faudrait vérifier que
l'isomorphisme est bien donné par le cup-produit en degré $0$ ; nous omettons
cette vérification).
\end{proof}

\begin{remark}
\label{perfect-remark-annoying-compatibility}
Soit $S = \Spec(A)$ un schéma affine. Soient $a : X \to S$
et $b : Y \to S$ des morphismes de schémas. Soient $\mathcal{F}$ et $\mathcal{G}$
des $\mathcal{O}_X$-modules quasi-cohérents, et soit $\mathcal{E}$
un $\mathcal{O}_Y$-module quasi-cohérent. Soit
$\xi \in H^i(X, \mathcal{G})$, dont le rappel est
$p^*\xi \in H^i(X \times_S Y, p^*\mathcal{G})$.
Alors le diagramme suivant est commutatif
$$
\xymatrix{
R\Gamma(X, \mathcal{F})[-i] \otimes_A^\mathbf{L} R\Gamma(Y, \mathcal{E})
\ar[d] \ar[rr]_-{\xi \otimes \text{id}} & &
R\Gamma(X, \mathcal{G} \otimes_{\mathcal{O}_X} \mathcal{F})
\otimes_A^\mathbf{L} R\Gamma(Y, \mathcal{E}) \ar[d] \\
R\Gamma(X \times_S Y, p^*\mathcal{F} \otimes q^*\mathcal{E})[-i]
\ar[rr]^-{p^*\xi} & &
R\Gamma(X \times_S Y,
p^*(\mathcal{G} \otimes_{\mathcal{O}_X} \mathcal{F}) \otimes q^*\mathcal{E})
}
$$
où les produits tensoriels non décorés sont pris sur $\mathcal{O}_{X \times_S Y}$.
Les flèches horizontales proviennent de Cohomologie, remarque
\ref{cohomology-remark-cup-with-element-map-total-cohomology},
et les flèches verticales sont (\ref{perfect-equation-kunneth-global}) ; elles sont donc
données par le rappel suivi du cup-produit sur $X \times_S Y$.
Le diagramme commute parce que le cup-produit global (sur $X \times_S Y$
avec les faisceaux $p^*\mathcal{G}$, $p^*\mathcal{F}$ et $q^*\mathcal{E}$)
est associatif ; voir
Cohomologie, lemme \ref{cohomology-lemma-cup-product-associative}.
\end{remark}


\section{Formule de K\"unneth, III}
\label{perfect-section-kunneth-complexes}

\noindent
Soient $X, Y, S, a, b, p, q, f$ comme dans l'introduction de la
section \ref{perfect-section-kunneth}. Dans cette section, étant donnés un
$\mathcal{O}_X$-module $\mathcal{F}$ et un $\mathcal{O}_Y$-module
$\mathcal{G}$, posons
$$
\mathcal{F} \boxtimes \mathcal{G} =
p^*\mathcal{F} \otimes_{\mathcal{O}_{X \times_S Y}} q^*\mathcal{G}
$$
Notons que, contrairement à ce qui se produira dans une section ultérieure, nous prenons
ici le produit tensoriel non dérivé.

\medskip\noindent
Sur $X$, soit $\mathcal{F}^\bullet$ un complexe de faisceaux de groupes abéliens
dont les termes sont des $\mathcal{O}_X$-modules quasi-cohérents
et tel que les différentielles
$d^i_\mathcal{F} : \mathcal{F}^i \to \mathcal{F}^{i + 1}$
soient des opérateurs différentiels sur $X/S$ d'ordre fini ; voir
Morphismes, section \ref{morphisms-section-differential-operators}.
De même, sur $Y$, soit $\mathcal{G}^\bullet$ un complexe de faisceaux
de groupes abéliens dont les termes sont des $\mathcal{O}_Y$-modules quasi-cohérents
et tel que les différentielles
$d^j_\mathcal{G} : \mathcal{G}^j \to \mathcal{G}^{j + 1}$
soient des opérateurs différentiels sur $Y/S$ d'ordre fini.
En appliquant la construction de
Morphismes, lemme \ref{morphisms-lemma-base-change-differential-operators},
on obtient un complexe double
$$
\xymatrix{
\ldots &
\ldots &
\ldots &
\ldots \\
\ldots \ar[r] &
\mathcal{F}^i \boxtimes \mathcal{G}^{j + 1}
\ar[r]^{d_1^{i, j + 1}} \ar[u] &
\mathcal{F}^{i + 1} \boxtimes \mathcal{G}^{j + 1} \ar[r] \ar[u] &
\ldots \\
\ldots \ar[r] &
\mathcal{F}^i \boxtimes \mathcal{G}^j
\ar[r]^{d_1^{i, j}} \ar[u]^{d_2^{i, j}} &
\mathcal{F}^{i + 1} \boxtimes \mathcal{G}^j \ar[r] \ar[u]_{d_2^{i + 1, j}} &
\ldots \\
\ldots &
\ldots \ar[u] &
\ldots \ar[u] &
\ldots
}
$$
de modules quasi-cohérents dont les morphismes sont des opérateurs différentiels
d'ordre fini sur $X \times_S Y / S$. Voir la discussion de
Morphismes, remarque \ref{morphisms-remark-base-change-differential-operators},
et
Homologie, exemple \ref{homology-example-double-complex-as-tensor-product-of}.
Explicitement, on pose
$$
d_1^{i, j} = d^i_\mathcal{F} \boxtimes 1
\quad\text{et}\quad
d_2^{i, j} = 1 \boxtimes d^j_\mathcal{G}
$$
Dans la discussion ci-dessous, la notation
$$
\text{Tot}(\mathcal{F}^\bullet \boxtimes \mathcal{G}^\bullet)
$$
désigne le complexe total associé à ce complexe double.
Les termes de ce complexe sont des
$\mathcal{O}_{X \times_S Y}$-modules quasi-cohérents, et ses différentielles
sont des opérateurs différentiels d'ordre fini sur $X \times_S Y / S$.

\medskip\noindent
Dans la situation ci-dessus, il existe un morphisme de « cup-produit relatif »
\begin{equation}
\label{perfect-equation-relative-de-rham-kunneth}
Ra_*(\mathcal{F}^\bullet)
\otimes_{\mathcal{O}_S}^\mathbf{L}
Rb_*(\mathcal{G}^\bullet)
\longrightarrow
Rf_*\left(\text{Tot}(\mathcal{F}^\bullet \boxtimes \mathcal{G}^\bullet)\right)
\end{equation}
On peut en effet construire ce morphisme en combinant
\begin{enumerate}
\item $Ra_*(\mathcal{F}^\bullet) \to Rf_*(p^{-1}\mathcal{F}^\bullet)$,
\item $Rb_*(\mathcal{G}^\bullet) \to Rf_*(q^{-1}\mathcal{G}^\bullet)$,
\item
$Rf_*(p^{-1}\mathcal{F}^\bullet)
\otimes_{\mathcal{O}_S}^\mathbf{L}
Rf_*(q^{-1}\mathcal{G}^\bullet)
\to
Rf_*(p^{-1}\mathcal{F}^\bullet \otimes_{f^{-1}\mathcal{O}_S}^\mathbf{L}
q^{-1}\mathcal{G}^\bullet)$,
\item
$p^{-1}\mathcal{F}^\bullet \otimes_{f^{-1}\mathcal{O}_S}^\mathbf{L}
q^{-1}\mathcal{G}^\bullet \to
\text{Tot}(p^{-1}\mathcal{F}^\bullet \otimes_{f^{-1}\mathcal{O}_S}
q^{-1}\mathcal{G}^\bullet)$
\item
$\text{Tot}(p^{-1}\mathcal{F}^\bullet \otimes_{f^{-1}\mathcal{O}_S}
q^{-1}\mathcal{G}^\bullet) \to \text{Tot}(\mathcal{F}^\bullet \boxtimes
\mathcal{G}^\bullet)$.
\end{enumerate}
Les morphismes (1) et (2) sont des morphismes de rappel, le morphisme (3) est
le cup-produit relatif ; voir
Cohomologie, remarque \ref{cohomology-remark-cup-product} ;
le morphisme (4) compare les produits tensoriels dérivé et non dérivé, et
le morphisme (5) est donné par les morphismes évidents
$p^{-1}\mathcal{F}^i \otimes_{f^{-1}\mathcal{O}_S} q^{-1}\mathcal{G}^j
\to \mathcal{F}^i \boxtimes \mathcal{G}^j$ sur les complexes doubles sous-jacents.

\medskip\noindent
Posons $A = \Gamma(S, \mathcal{O}_S)$. Il existe un morphisme de « cup-produit global »
\begin{equation}
\label{perfect-equation-de-rham-kunneth}
R\Gamma(X, \mathcal{F}^\bullet)
\otimes_A^\mathbf{L}
R\Gamma(Y, \mathcal{G}^\bullet)
\longrightarrow
R\Gamma(X \times_S Y,
\text{Tot}(\mathcal{F}^\bullet \boxtimes \mathcal{G}^\bullet))
\end{equation}
dans $D(A)$. Ce morphisme se construit comme le cup-produit relatif ci-dessus,
en utilisant
\begin{enumerate}
\item $R\Gamma(X, \mathcal{F}^\bullet) \to
R\Gamma(X \times_S Y, p^{-1}\mathcal{F}^\bullet)$
\item $R\Gamma(Y, \mathcal{G}^\bullet) \to
R\Gamma(X \times_S Y, q^{-1}\mathcal{G}^\bullet)$,
\item $R\Gamma(X \times_S Y, p^{-1}\mathcal{F}^\bullet)
\otimes_A^\mathbf{L} R\Gamma(X \times_S Y, q^{-1}\mathcal{G}^\bullet) \to
R\Gamma(X \times_S Y,
p^{-1}\mathcal{F}^\bullet \otimes_{f^{-1}\mathcal{O}_S}^\mathbf{L}
q^{-1}\mathcal{G}^\bullet)$,
\item
$p^{-1}\mathcal{F}^\bullet \otimes_{f^{-1}\mathcal{O}_S}^\mathbf{L}
q^{-1}\mathcal{G}^\bullet \to
\text{Tot}(p^{-1}\mathcal{F}^\bullet \otimes_{f^{-1}\mathcal{O}_S}
q^{-1}\mathcal{G}^\bullet)$
\item
$\text{Tot}(p^{-1}\mathcal{F}^\bullet \otimes_{f^{-1}\mathcal{O}_S}
q^{-1}\mathcal{G}^\bullet) \to \text{Tot}(\mathcal{F}^\bullet \boxtimes
\mathcal{G}^\bullet)$.
\end{enumerate}
Ici, les morphismes (1) et (2) sont les morphismes de rappel, et le morphisme (3)
est le cup-produit construit dans Cohomologie, section \ref{cohomology-section-cup-product}.
Les morphismes (4) et (5) sont ceux indiqués au paragraphe précédent.

\begin{lemma}
\label{perfect-lemma-kunneth-special}
Dans la situation ci-dessus, le cup-produit (\ref{perfect-equation-de-rham-kunneth})
est un isomorphisme dans $D(A)$ si les hypothèses suivantes sont satisfaites :
\begin{enumerate}
\item $S = \Spec(A)$ est affine,
\item $X$ et $Y$ sont quasi-compacts à diagonale affine,
\item $\mathcal{F}^\bullet$ est borné,
\item $\mathcal{G}^\bullet$ est borné inférieurement,
\item $\mathcal{F}^n$ est plat sur $S$, et
\item $\mathcal{G}^m$ est plat sur $S$.
\end{enumerate}
\end{lemma}

\begin{proof}
Nous utiliserons les notations $\mathcal{A}_{X/S}$ et $\mathcal{A}_{Y/S}$
introduites dans Morphismes, remarque
\ref{morphisms-remark-base-change-differential-operators}.
Supposons que l'on ait des morphismes de complexes
$$
\mathcal{F}_1^\bullet \to
\mathcal{F}_2^\bullet \to
\mathcal{F}_3^\bullet \to
\mathcal{F}_1^\bullet[1]
$$
dans la catégorie $\mathcal{A}_{X/S}$.
La fonctorialité du cup-produit fournit alors un diagramme commutatif
$$
\xymatrix{
R\Gamma(X, \mathcal{F}_1^\bullet)
\otimes_A^\mathbf{L}
R\Gamma(Y, \mathcal{G}^\bullet)
\ar[r] \ar[d] &
R\Gamma(X \times_S Y,
\text{Tot}(\mathcal{F}_1^\bullet \boxtimes \mathcal{G}^\bullet)) \ar[d] \\
R\Gamma(X, \mathcal{F}_2^\bullet)
\otimes_A^\mathbf{L}
R\Gamma(Y, \mathcal{G}^\bullet)
\ar[r] \ar[d] &
R\Gamma(X \times_S Y,
\text{Tot}(\mathcal{F}_2^\bullet \boxtimes \mathcal{G}^\bullet)) \ar[d] \\
R\Gamma(X, \mathcal{F}_3^\bullet)
\otimes_A^\mathbf{L}
R\Gamma(Y, \mathcal{G}^\bullet)
\ar[r] \ar[d] &
R\Gamma(X \times_S Y,
\text{Tot}(\mathcal{F}_3^\bullet \boxtimes \mathcal{G}^\bullet)) \ar[d] \\
R\Gamma(X, \mathcal{F}_1^\bullet[1])
\otimes_A^\mathbf{L}
R\Gamma(Y, \mathcal{G}^\bullet)
\ar[r] &
R\Gamma(X \times_S Y,
\text{Tot}(\mathcal{F}_1^\bullet[1] \boxtimes \mathcal{G}^\bullet))
}
$$
Si les morphismes initiaux forment un triangle distingué dans la catégorie homotopique
de $\mathcal{A}_{X/S}$, alors
les colonnes de ce diagramme forment des triangles distingués dans $D(A)$.

\medskip\noindent
Dans la situation du lemme,
supposons $\mathcal{F}^n = 0$ pour $n < i$. On peut alors considérer la
suite exacte courte de complexes, scindée terme à terme,
$$
0 \to \sigma_{\geq i + 1}\mathcal{F}^\bullet \to
\mathcal{F}^\bullet \to \mathcal{F}^i[-i] \to 0
$$
où la troncature est celle de
Homologie, section \ref{homology-section-truncations}.
Elle donne le triangle distingué
$$
\sigma_{\geq i + 1}\mathcal{F}^\bullet \to
\mathcal{F}^\bullet \to
\mathcal{F}^i[-i] \to
(\sigma_{\geq i + 1}\mathcal{F}^\bullet)[1]
$$
dans la catégorie homotopique de $\mathcal{A}_{X/S}$,
où la dernière flèche est donnée par le morphisme de bord
$\mathcal{F}^i \to \mathcal{F}^{i + 1}$.
La discussion qui précède montre qu'il suffit de démontrer le lemme pour
$\mathcal{F}^i[-i]$ et $\sigma_{\geq i + 1}\mathcal{F}^\bullet$.
Puisque $\sigma_{\geq i + 1}\mathcal{F}^\bullet$ a moins de termes non nuls,
une récurrence montre que, si l'on peut démontrer le lemme lorsque $\mathcal{F}^\bullet$
n'est non nul qu'en un seul degré, alors le lemme en résulte.
On peut donc supposer que $\mathcal{F}^\bullet$ n'est non nul qu'en un seul degré.

\medskip\noindent
Supposons que $\mathcal{F}^\bullet$ soit le complexe ayant pour terme, plat sur $S$,
le $\mathcal{O}_X$-module quasi-cohérent $\mathcal{F}$ placé en degré $0$,
et nul dans les autres degrés. Observons que $R\Gamma(X, \mathcal{F})$ est de
dimension de Tor finie, par exemple d'après le lemme \ref{perfect-lemma-cohomology-de-rham-base-change}.
Disons que son amplitude de Tor est contenue dans $[i, j]$.
Choisissons $N \gg 0$ et considérons le triangle distingué
$$
\sigma_{\geq N + 1}\mathcal{G}^\bullet \to
\mathcal{G}^\bullet \to
\sigma_{\leq N}\mathcal{G}^\bullet \to
(\sigma_{\geq N + 1}\mathcal{G}^\bullet)[1]
$$
dans la catégorie homotopique de $\mathcal{A}_{Y/S}$. Observons maintenant que les deux complexes
$$
R\Gamma(X, \mathcal{F})
\otimes_A^\mathbf{L}
R\Gamma(Y, \sigma_{\geq N + 1}\mathcal{G}^\bullet)
\quad\text{et}\quad
R\Gamma(X \times_S Y,
\text{Tot}(\mathcal{F} \boxtimes \sigma_{\geq N + 1}\mathcal{G}^\bullet))
$$
ont une cohomologie nulle en degrés $\leq N + i$. Ainsi, en utilisant les arguments
donnés ci-dessus, si l'on veut démontrer notre assertion en un degré donné, on peut
supposer que $\mathcal{G}^\bullet$ est borné.
En répétant encore une fois les arguments précédents, on peut aussi supposer
que $\mathcal{G}^\bullet$ n'est non nul qu'en un seul degré.
Ce cas est traité par le lemme \ref{perfect-lemma-kunneth-single-sheaf}.
\end{proof}






\section{Formule de K\"unneth pour Ext}
\label{perfect-section-kunneth-Ext}

\noindent
Considérons un diagramme cartésien de schémas
$$
\xymatrix{
& X \times_S Y \ar[ld]^p \ar[rd]_q \ar[dd]^f \\
X \ar[rd]_a & & Y \ar[ld]^b \\
& S
}
$$
Pour $K \in D(\mathcal{O}_X)$ et $M \in D(\mathcal{O}_Y)$,
posons dans cette section
$$
K \boxtimes M =
Lp^*K \otimes_{\mathcal{O}_{X \times_S Y}}^\mathbf{L} Lq^*M
$$
Nous affirmons qu'il existe un morphisme canonique
\begin{equation}
\label{perfect-equation-kunneth-ext}
Ra_*R\SheafHom(K, K')
\otimes_{\mathcal{O}_S}^\mathbf{L}
Rb_*R\SheafHom(M, M')
\longrightarrow
Rf_*(R\SheafHom(K \boxtimes M, K' \boxtimes M'))
\end{equation}
pour $K, K' \in D(\mathcal{O}_X)$ et $M, M' \in D(\mathcal{O}_Y)$.
On peut en effet prendre le morphisme adjoint au morphisme
$$
\begin{matrix}
Lf^*\left(Ra_*R\SheafHom(K, K')
\otimes_{\mathcal{O}_S}^\mathbf{L}
Rb_* R\SheafHom(M, M')\right) = \\
Lf^* Ra_* R\SheafHom(K, K')
\otimes_{\mathcal{O}_{X \times_S Y}}^\mathbf{L}
Lf^* Rb_* R\SheafHom(M, M') = \\
Lp^* La^* Ra_* R\SheafHom(K, K')
\otimes_{\mathcal{O}_{X \times_S Y}}^\mathbf{L}
Lq^* Lb^* Rb_* R\SheafHom(M, M') \to \\
Lp^* R\SheafHom(K, K')
\otimes_{\mathcal{O}_{X \times_S Y}}^\mathbf{L}
Lq^* R\SheafHom(M, M') \to \\
R\SheafHom(Lp^*K, Lp^*K')
\otimes_{\mathcal{O}_{X \times_S Y}}^\mathbf{L}
R\SheafHom(Lq^*M, Lq^*M') \to \\
R\SheafHom(K \boxtimes M, K' \boxtimes M')
\end{matrix}
$$
Ici, la première égalité exprime la compatibilité des rappels avec les produits tensoriels,
Cohomologie, lemme \ref{cohomology-lemma-pullback-tensor-product}.
La deuxième égalité utilise $f = a \circ p = b \circ q$ et
la composition des rappels,
Cohomologie, lemme \ref{cohomology-lemma-derived-pullback-composition}.
La première flèche est donnée par les morphismes d'adjonction
$La^* Ra_* \to \text{id}$ et
$Lb^* Rb_* \to \text{id}$, car l'image directe et le rappel sont adjoints,
Cohomologie, lemme \ref{cohomology-lemma-adjoint}.
La deuxième flèche est donnée par
Cohomologie, remarque \ref{cohomology-remark-prepare-fancy-base-change}.
La troisième et dernière flèche est celle de
Cohomologie, remarque \ref{cohomology-remark-tensor-internal-hom}.
Le résultat suivant en est un cas particulier simple.

\begin{lemma}
\label{perfect-lemma-kunneth-Ext}
Dans la situation ci-dessus, supposons que $a$ et $b$ soient quasi-compacts et
quasi-séparés et que $X$ et $Y$ soient Tor-indépendants sur $S$.
Si $K$ est parfait, $K' \in D_\QCoh(\mathcal{O}_X)$, $M$ est parfait, et si
$M' \in D_\QCoh(\mathcal{O}_Y)$, alors (\ref{perfect-equation-kunneth-ext})
est un isomorphisme.
\end{lemma}

\begin{proof}
Dans ce cas, on a $R\SheafHom(K, K') = K' \otimes^\mathbf{L} K^\vee$,
$R\SheafHom(M, M') = M' \otimes^\mathbf{L} M^\vee$, et
$$
R\SheafHom(K \boxtimes M, K' \boxtimes M') =
(K' \otimes^\mathbf{L} K^\vee) \boxtimes
(M' \otimes^\mathbf{L} M^\vee)
$$
Voir Cohomologie, lemme \ref{cohomology-lemma-dual-perfect-complex} ;
on utilise également le fait que la perfection est préservée par rappel
et par produit tensoriel.
Ce cas résulte donc du lemme \ref{perfect-lemma-kunneth}.
(Nous omettons de vérifier que ces identifications
donnent le même morphisme.)
\end{proof}







\section{Cohomologie et changement de base, V}
\label{perfect-section-cohomology-base-change}

\noindent
Dans la section \ref{perfect-section-cohomology-and-base-change-perfect},
nous avons vu qu'un théorème de changement de base vaut lorsque les morphismes sont Tor-indépendants.
Même dans le cas affine, il ne peut exister de théorème de changement de base sans une telle
condition ; voir
Compléments sur l'algèbre, section \ref{more-algebra-section-tor-independence}.
Dans cette section, nous analysons quand on peut obtenir un résultat de changement de base
« un complexe à la fois ».

\medskip\noindent
Pour que cela fonctionne, supposons donné un diagramme commutatif
$$
\xymatrix{
X' \ar[r]_{g'} \ar[d]_{f'} &
X \ar[d]^f \\
S' \ar[r]^g &
S
}
$$
de schémas (en général, nous supposerons qu'il est cartésien).
Soit $K \in D_\QCoh(\mathcal{O}_X)$,
et soit $L(g')^*K \to K'$ un morphisme dans $D_\QCoh(\mathcal{O}_{X'})$.
Pour un point $x' \in X'$, posons $x = g'(x') \in X$,
$s' = f'(x') \in S'$ et $s = f(x) = g(s')$.
On peut alors considérer les morphismes
$$
K_x \otimes_{\mathcal{O}_{S, s}}^\mathbf{L} \mathcal{O}_{S', s'} \to
K_x \otimes_{\mathcal{O}_{X, x}}^\mathbf{L} \mathcal{O}_{X', x'} \to
K'_{x'}
$$
où la première flèche est celle de Compléments sur l'algèbre,
équation (\ref{more-algebra-equation-comparison-map}),
et la seconde provient de
$(L(g')^*K)_{x'} =
K_x \otimes_{\mathcal{O}_{X, x}}^\mathbf{L} \mathcal{O}_{X', x'}$
et du morphisme donné $L(g')^*K \to K'$.
Pour chaque $i \in \mathbf{Z}$, on obtient une structure de
$\mathcal{O}_{X, x} \otimes_{\mathcal{O}_{S, s}} \mathcal{O}_{S', s'}$-module
sur
$H^i(K_x \otimes_{\mathcal{O}_{S, s}}^\mathbf{L} \mathcal{O}_{S', s'})$.
En rassemblant le tout, on obtient des morphismes canoniques
\begin{equation}
\label{perfect-equation-bc}
H^i(K_x \otimes_{\mathcal{O}_{S, s}}^\mathbf{L} \mathcal{O}_{S', s'})
\otimes_{(\mathcal{O}_{X, x} \otimes_{\mathcal{O}_{S, s}} \mathcal{O}_{S', s'})}
\mathcal{O}_{X', x'}
\longrightarrow
H^i(K'_{x'})
\end{equation}
de $\mathcal{O}_{X', x'}$-modules.

\begin{lemma}
\label{perfect-lemma-single-complex-base-change-condition}
Soit
$$
\xymatrix{
X' \ar[r]_{g'} \ar[d]_{f'} &
X \ar[d]^f \\
S' \ar[r]^g &
S
}
$$
un diagramme cartésien de schémas. Soit $K \in D_\QCoh(\mathcal{O}_X)$,
et soit $L(g')^*K \to K'$ un morphisme dans $D_\QCoh(\mathcal{O}_{X'})$.
Les assertions suivantes sont équivalentes :
\begin{enumerate}
\item pour tous $x' \in X'$ et $i \in \mathbf{Z}$, le morphisme (\ref{perfect-equation-bc})
est un isomorphisme,
\item si $U \subset X$ et $V' \subset S'$ sont des ouverts affines dont les images sont
contenues dans l'ouvert affine $V \subset S$, et si $U' = V' \times_V U$,
la composée
$$
R\Gamma(U, K) \otimes_{\mathcal{O}_S(U)}^\mathbf{L} \mathcal{O}_{S'}(V')
\to
R\Gamma(U, K) \otimes_{\mathcal{O}_X(U)}^\mathbf{L} \mathcal{O}_{X'}(U')
\to
R\Gamma(U', K')
$$
est un isomorphisme dans $D(\mathcal{O}_{S'}(V'))$, et
\item il existe un ensemble $I$ de quadruplets $U_i, V_i', V_i, U_i'$, $i \in I$,
comme en (2), tels que $X' = \bigcup U'_i$.
\end{enumerate}
\end{lemma}

\begin{proof}
La deuxième flèche de (2) provient de l'égalité
$$
R\Gamma(U, K) \otimes_{\mathcal{O}_X(U)}^\mathbf{L} \mathcal{O}_{X'}(U') =
R\Gamma(U', L(g')^*K)
$$
du lemme \ref{perfect-lemma-quasi-coherence-pullback} et du morphisme donné
$L(g')^*K \to K'$. La première flèche de (2) est celle de
Compléments sur l'algèbre, équation (\ref{more-algebra-equation-comparison-map}).
Il est clair que (2) implique (3). Observons que (1) est local sur $X'$.
Il suffit donc de montrer que, si $X$, $S$, $S'$ et $X'$ sont affines, alors
(1) équivaut à la condition que
$$
R\Gamma(X, K) \otimes_{\mathcal{O}_S(S)}^\mathbf{L} \mathcal{O}_{S'}(S')
\to
R\Gamma(X, K) \otimes_{\mathcal{O}_X(X)}^\mathbf{L} \mathcal{O}_{X'}(X')
\to
R\Gamma(X', K')
$$
soit un isomorphisme dans $D(\mathcal{O}_{S'}(S'))$. Écrivons
$S = \Spec(R)$, $X = \Spec(A)$, $S' = \Spec(R')$, $X' = \Spec(A')$,
et supposons que $K$ corresponde au complexe $M^\bullet$ de $A$-modules, et que
$K'$ corresponde au complexe $N^\bullet$ de $A'$-modules.
Notons que $A' = A \otimes_R R'$. La condition ci-dessus signifie que la composée
$$
M^\bullet \otimes_R^\mathbf{L} R' \to
M^\bullet \otimes_A^\mathbf{L} A' \to
N^\bullet
$$
est un isomorphisme dans $D(R')$. De manière équivalente, pour tout
$i \in \mathbf{Z}$, le morphisme
$$
H^i(M^\bullet \otimes_R^\mathbf{L} R') \to
H^i(M^\bullet \otimes_A^\mathbf{L} A') \to
H^i(N^\bullet)
$$
est un isomorphisme. Observons qu'il s'agit d'un morphisme de $A \otimes_R R'$-modules,
c'est-à-dire de $A'$-modules. D'autre part, (1) exige
que, pour des idéaux premiers compatibles
$\mathfrak q' \subset A'$, $\mathfrak q \subset A$,
$\mathfrak p' \subset R'$, $\mathfrak p \subset R$,
la composée
$$
H^i(M^\bullet_\mathfrak q \otimes_{R_\mathfrak p}^\mathbf{L} R'_{\mathfrak p'})
\otimes_{(A_\mathfrak q \otimes_{R_\mathfrak p} R'_{\mathfrak p'})}
A'_{\mathfrak q'} \to
H^i(M^\bullet_{\mathfrak q}
\otimes_{A_\mathfrak q}^\mathbf{L} A'_{\mathfrak q'})
\to H^i(N^\bullet_{\mathfrak q'})
$$
soit un isomorphisme. Comme
$$
H^i(M^\bullet_\mathfrak q \otimes_{R_\mathfrak p}^\mathbf{L} R'_{\mathfrak p'})
\otimes_{(A_\mathfrak q \otimes_{R_\mathfrak p} R'_{\mathfrak p'})}
A'_{\mathfrak q'} =
H^i(M^\bullet \otimes_R^\mathbf{L} R') \otimes_{A'} A'_{\mathfrak q'}
$$
est la localisation en $\mathfrak q'$,
on voit que ces deux conditions sont équivalentes d'après
Algèbre, lemme \ref{algebra-lemma-characterize-zero-local}.
\end{proof}

\begin{lemma}
\label{perfect-lemma-single-complex-base-change}
Soit
$$
\xymatrix{
X' \ar[r]_{g'} \ar[d]_{f'} &
X \ar[d]^f \\
S' \ar[r]^g &
S
}
$$
un diagramme cartésien de schémas. Soit $K \in D_\QCoh(\mathcal{O}_X)$,
et soit $L(g')^*K \to K'$ un morphisme dans $D_\QCoh(\mathcal{O}_{X'})$.
Si
\begin{enumerate}
\item les conditions équivalentes du
lemme \ref{perfect-lemma-single-complex-base-change-condition} sont satisfaites, et
\item $f$ est quasi-compact et quasi-séparé,
\end{enumerate}
alors la composée $Lg^*Rf_*K \to Rf'_*L(g')^*K \to Rf'_*K'$
est un isomorphisme.
\end{lemma}

\begin{proof}
On pourrait le démontrer par la même méthode que dans la démonstration du
lemme \ref{perfect-lemma-compare-base-change}, mais nous allons plutôt utiliser
le principe d'induction et Mayer--Vietoris relatif.

\medskip\noindent
Pour vérifier que le morphisme est un isomorphisme, on peut travailler localement sur $S'$.
On peut donc supposer que $g : S' \to S$ est un morphisme de schémas affines.
En particulier, $X$ est un schéma quasi-compact et quasi-séparé.
Nous utiliserons le principe d'induction de
Cohomologie des schémas, lemme \ref{coherent-lemma-induction-principle},
pour démontrer que, pour tout ouvert quasi-compact $U \subset X$, le morphisme construit
de la même manière $Lg^*R(U \to S)_*K|_U \to R(U' \to S')_*K'|_{U'}$
est un isomorphisme. Ici $U' = (g')^{-1}(U)$.

\medskip\noindent
Si $U \subset X$ est un ouvert affine, alors le résultat est
vrai par hypothèse ; voir
le lemme \ref{perfect-lemma-single-complex-base-change-condition}, partie (2),
et la traduction algébrique fournie par les
lemmes \ref{perfect-lemma-affine-compare-bounded} et
\ref{perfect-lemma-quasi-coherence-pullback}.

\medskip\noindent
L'étape d'induction. Supposons que $X = U \cup V$ soit un recouvrement ouvert
avec $U$, $V$ et $U \cap V$
quasi-compacts, tel que le résultat soit vrai pour $U$, $V$ et $U \cap V$.
Notons $a = f|_U$, $b = f|_V$ et $c = f|_{U \cap V}$.
Soient $a' : U' \to S'$, $b' : V' \to S'$ et $c' : U' \cap V' \to S'$
les changements de base de $a$, $b$ et $c$.
En utilisant les triangles distingués de Mayer--Vietoris relatif
(Cohomologie, lemme \ref{cohomology-lemma-unbounded-relative-mayer-vietoris}),
on obtient un diagramme commutatif
$$
\xymatrix{
Lg^*Rf_*K \ar[r] \ar[d] &
Rf'_* K' \ar[d] \\
Lg^*Ra_* K|_U \oplus
Lg^*Rb_* K|_V \ar[r] \ar[d] &
Ra'_* K'|_{U'} \oplus
Rb'_* K'|_{V'} \ar[d] \\
Lg^*Rc_* K|_{U \cap V} \ar[r] \ar[d] &
Rc'_* K'|_{U' \cap V'} \ar[d] \\
Lg^*Rf_* K[1] \ar[r] &
Rf'_* K'[1]
}
$$
Puisque les deuxième et troisième flèches horizontales sont des isomorphismes, la première
l'est aussi (Catégories dérivées, lemme \ref{derived-lemma-third-isomorphism-triangle}),
ce qui achève la démonstration du lemme.
\end{proof}

\begin{lemma}
\label{perfect-lemma-single-complex-base-change-condition-inherited}
Soit
$$
\xymatrix{
X' \ar[r]_{g'} \ar[d]_{f'} &
X \ar[d]^f \\
S' \ar[r]^g &
S
}
$$
un diagramme cartésien de schémas. Soit $K \in D_\QCoh(\mathcal{O}_X)$,
et soit $L(g')^*K \to K'$ un morphisme dans $D_\QCoh(\mathcal{O}_{X'})$.
Si les conditions équivalentes du
lemme \ref{perfect-lemma-single-complex-base-change-condition} sont satisfaites, alors
\begin{enumerate}
\item pour $E \in D_\QCoh(\mathcal{O}_X)$, les conditions
équivalentes du lemme \ref{perfect-lemma-single-complex-base-change-condition} sont satisfaites
pour $L(g')^*(E \otimes^\mathbf{L} K) \to L(g')^*E \otimes^\mathbf{L} K'$,
\item si $E$ dans $D(\mathcal{O}_X)$ est parfait, les conditions équivalentes du
lemme \ref{perfect-lemma-single-complex-base-change-condition} sont satisfaites pour
$L(g')^*R\SheafHom(E, K) \to R\SheafHom(L(g')^*E, K')$, et
\item si $K$ est borné inférieurement et si $E$ dans $D(\mathcal{O}_X)$ est
pseudo-cohérent, les conditions équivalentes du
lemme \ref{perfect-lemma-single-complex-base-change-condition} sont satisfaites pour
$L(g')^*R\SheafHom(E, K) \to R\SheafHom(L(g')^*E, K')$.
\end{enumerate}
\end{lemma}

\begin{proof}
L'énoncé a un sens, car les complexes en jeu ont des faisceaux de cohomologie
quasi-cohérents d'après les lemmes
\ref{perfect-lemma-quasi-coherence-pullback},
\ref{perfect-lemma-quasi-coherence-tensor-product} et
\ref{perfect-lemma-quasi-coherence-internal-hom}, ainsi que les
lemmes de Cohomologie \ref{cohomology-lemma-pseudo-coherent-pullback} et
\ref{cohomology-lemma-perfect-pullback}.
Ceci étant dit, on peut vérifier que les morphismes (\ref{perfect-equation-bc})
sont des isomorphismes dans le cas (1) en calculant la source et le but
de (\ref{perfect-equation-bc}) au moyen de la transitivité du produit tensoriel ; voir
Compléments sur l'algèbre, lemme \ref{more-algebra-lemma-triple-tensor-product}.
Le morphisme dans (2) et (3) est la composée
$$
L(g')^*R\SheafHom(E, K) \to R\SheafHom(L(g')^*E, L(g')^*K)
\to R\SheafHom(L(g')^*E, K')
$$
où la première flèche est celle de la
remarque de Cohomologie \ref{cohomology-remark-prepare-fancy-base-change},
et la seconde provient du morphisme donné $L(g')^*K \to K'$.
Pour démontrer que les morphismes (\ref{perfect-equation-bc}) sont des isomorphismes, on représente
$E_x$ par un complexe borné de $\mathcal{O}_{X. x}$-modules projectifs de type fini
dans le cas (2), ou par un complexe borné supérieurement de modules libres de type fini dans le cas (3),
et l'on calcule la source et le but de la flèche.
Certains détails sont omis.
\end{proof}

\begin{lemma}
\label{perfect-lemma-base-change-tensor}
Soit $f : X \to S$ un morphisme quasi-compact et quasi-séparé de
schémas. Soit $E \in D_\QCoh(\mathcal{O}_X)$. Soit $\mathcal{G}^\bullet$
un complexe borné supérieurement de $\mathcal{O}_X$-modules
quasi-cohérents et plats sur $S$. Alors la formation de
$$
Rf_*(E \otimes^\mathbf{L}_{\mathcal{O}_X} \mathcal{G}^\bullet)
$$
commute aux changements de base arbitraires (voir la démonstration pour un énoncé précis).
\end{lemma}

\begin{proof}
L'énoncé signifie ce qui suit. Soit $g : S' \to S$ un morphisme de
schémas et considérons le diagramme de changement de base
$$
\xymatrix{
X' \ar[r]_{g'} \ar[d]_{f'} &
X \ar[d]^f \\
S' \ar[r]^g &
S
}
$$
autrement dit $X' = S' \times_S X$. Le lemme affirme que
$$
Lg^*Rf_*(E \otimes^\mathbf{L}_{\mathcal{O}_X} \mathcal{G}^\bullet)
\longrightarrow
Rf'_*\left(
L(g')^*E \otimes^\mathbf{L}_{\mathcal{O}_{X'}} (g')^*\mathcal{G}^\bullet
\right)
$$
est un isomorphisme. Observons que, dans le membre de droite, nous n'utilisons
{\bf pas} l'image inverse dérivée pour $\mathcal{G}^\bullet$.
Pour le démontrer, appliquons les lemmes \ref{perfect-lemma-single-complex-base-change} et
\ref{perfect-lemma-single-complex-base-change-condition-inherited} : il suffit alors de démontrer que le morphisme canonique
$$
L(g')^*\mathcal{G}^\bullet \to (g')^*\mathcal{G}^\bullet
$$
satisfait les conditions équivalentes du
lemme \ref{perfect-lemma-single-complex-base-change-condition}.
Cela résulte de la vérification de la condition sur les fibres, où l'assertion
découle immédiatement du fait que
$\mathcal{G}^\bullet_x \otimes_{\mathcal{O}_{S, s}} \mathcal{O}_{S', s'}$
calcule le produit tensoriel dérivé, par les hypothèses faites sur le complexe
$\mathcal{G}^\bullet$.
\end{proof}

\begin{lemma}
\label{perfect-lemma-base-change-RHom}
Soit $f : X \to S$ un morphisme quasi-compact et quasi-séparé de schémas.
Soit $E$ un objet de $D(\mathcal{O}_X)$.
Soit $\mathcal{G}^\bullet$ un complexe de
$\mathcal{O}_X$-modules quasi-cohérents. Si
\begin{enumerate}
\item $E$ est parfait, $\mathcal{G}^\bullet$ est borné supérieurement,
et $\mathcal{G}^n$ est plat sur $S$, ou bien
\item $E$ est pseudo-cohérent, $\mathcal{G}^\bullet$ est borné,
et $\mathcal{G}^n$ est plat sur $S$,
\end{enumerate}
alors la formation de
$$
Rf_*R\SheafHom(E, \mathcal{G}^\bullet)
$$
commute aux changements de base arbitraires (voir la démonstration pour un énoncé précis).
\end{lemma}

\begin{proof}
L'énoncé signifie ce qui suit. Soit $g : S' \to S$ un morphisme de
schémas et considérons le diagramme de changement de base
$$
\xymatrix{
X' \ar[r]_{g'} \ar[d]_{f'} &
X \ar[d]^f \\
S' \ar[r]^g &
S
}
$$
autrement dit $X' = S' \times_S X$. Le lemme affirme que
$$
Lg^*Rf_*R\SheafHom(E, \mathcal{G}^\bullet)
\longrightarrow
R(f')_*R\SheafHom(L(g')^*E, (g')^*\mathcal{G}^\bullet)
$$
est un isomorphisme. Observons que, dans le membre de droite, nous n'utilisons
{\bf pas} l'image inverse dérivée pour $\mathcal{G}^\bullet$. Pour le démontrer, appliquons
les lemmes \ref{perfect-lemma-single-complex-base-change} et
\ref{perfect-lemma-single-complex-base-change-condition-inherited} : il suffit alors de démontrer que le morphisme canonique
$$
L(g')^*\mathcal{G}^\bullet \to (g')^*\mathcal{G}^\bullet
$$
satisfait les conditions équivalentes du
lemme \ref{perfect-lemma-single-complex-base-change-condition}.
Cela a été démontré dans la démonstration du lemme \ref{perfect-lemma-base-change-tensor}.
\end{proof}




\section{Production de complexes parfaits}
\label{perfect-section-producing-perfect}

\noindent
Le lemme suivant est notre principal outil technique pour produire des
complexes parfaits. Des versions ultérieures de ce résultat s'y ramèneront
par approximation noethérienne ; voir la
section \ref{perfect-section-cohomology-and-base-change-final}.

\begin{lemma}
\label{perfect-lemma-perfect-direct-image}
Soit $S$ un schéma noethérien. Soit $f : X \to S$ un morphisme de schémas
localement de type fini. Soit $E \in D(\mathcal{O}_X)$ tel que
\begin{enumerate}
\item $E \in D^b_{\textit{Coh}}(\mathcal{O}_X)$,
\item le support de $H^i(E)$ est propre sur $S$ pour tout $i$, et
\item $E$ est de dimension Tor finie comme objet de $D(f^{-1}\mathcal{O}_S)$.
\end{enumerate}
Alors $Rf_*E$ est un objet parfait de $D(\mathcal{O}_S)$.
\end{lemma}

\begin{proof}
Le lemme \ref{perfect-lemma-direct-image-coherent} montre que $Rf_*E$ est un objet de
$D^b_{\textit{Coh}}(\mathcal{O}_S)$. Ainsi $Rf_*E$ est pseudo-cohérent
(lemme \ref{perfect-lemma-identify-pseudo-coherent-noetherian}).
Il suffit donc de montrer que $Rf_*E$ est de dimension Tor finie ; voir
Cohomologie, lemme \ref{cohomology-lemma-perfect}.
D'après le lemme \ref{perfect-lemma-tor-qc-qs}, il suffit de vérifier que
$Rf_*(E) \otimes_{\mathcal{O}_S}^\mathbf{L} \mathcal{F}$
est à cohomologie uniformément bornée pour tout faisceau quasi-cohérent
$\mathcal{F}$ sur $S$. Le fait d'être borné supérieurement est clair puisque $Rf_*(E)$
est borné supérieurement. Soit $T \subset X$ la réunion des supports
des $H^i(E)$ pour tous les $i$. Alors $T$ est propre sur $S$ en vertu des hypothèses
(1) et (2) ; voir Cohomologie des schémas, lemme
\ref{coherent-lemma-union-closed-proper-over-base}.
En particulier, il existe un ouvert quasi-compact
$X' \subset X$ contenant $T$. En posant $f' = f|_{X'}$, on a
$Rf_*(E) = Rf'_*(E|_{X'})$, car la restriction de $E$ à $X \setminus T$ est nulle.
On peut donc remplacer $X$ par $X'$ et supposer $f$ quasi-compact.
De plus, $f$ est quasi-séparé d'après Morphismes, lemme
\ref{morphisms-lemma-finite-type-Noetherian-quasi-separated}. Or
$$
Rf_*(E) \otimes_{\mathcal{O}_S}^\mathbf{L} \mathcal{F} =
Rf_*\left(E \otimes_{\mathcal{O}_X}^\mathbf{L} Lf^*\mathcal{F}\right) =
Rf_*\left(E \otimes_{f^{-1}\mathcal{O}_S}^\mathbf{L} f^{-1}\mathcal{F}\right)
$$
d'après
le lemme \ref{perfect-lemma-cohomology-base-change}
et le
lemme de Cohomologie \ref{cohomology-lemma-variant-derived-pullback}.
Par l'hypothèse (3), le complexe
$E \otimes_{f^{-1}\mathcal{O}_S}^\mathbf{L} f^{-1}\mathcal{F}$
a ses faisceaux de cohomologie dans un
intervalle fini fixé, disons $[a, b]$. Son image par $Rf_*$
a donc sa cohomologie dans l'intervalle $[a, \infty)$, ce qui conclut.
\end{proof}

\begin{lemma}
\label{perfect-lemma-tensor-perfect}
Soit $S$ un schéma noethérien. Soit $f : X \to S$ un morphisme de schémas
localement de type fini. Soit $E \in D(\mathcal{O}_X)$ parfait.
Soit $\mathcal{G}^\bullet$ un complexe borné de
$\mathcal{O}_X$-modules cohérents, plats sur $S$ et à support propre sur $S$.
Alors $K = Rf_*(E \otimes_{\mathcal{O}_X}^\mathbf{L} \mathcal{G}^\bullet)$
est un objet parfait de $D(\mathcal{O}_S)$.
\end{lemma}

\begin{proof}
L'objet $K$ est parfait d'après le lemme \ref{perfect-lemma-perfect-direct-image}.
Vérifions que ce lemme s'applique. Localement, $E$ est isomorphe à un complexe fini
de $\mathcal{O}_X$-modules libres de type fini. Ainsi, localement,
$E \otimes^\mathbf{L}_{\mathcal{O}_X} \mathcal{G}^\bullet$ est isomorphe
à un complexe fini dont les termes sont de la forme
$$
\bigoplus\nolimits_{i = a, \ldots, b} (\mathcal{G}^i)^{\oplus r_i}
$$
pour certains entiers $a, b, r_a, \ldots, r_b$. Il en résulte immédiatement que les
faisceaux de cohomologie $H^i(E \otimes^\mathbf{L}_{\mathcal{O}_X} \mathcal{G})$
sont cohérents. L'hypothèse sur la dimension Tor est également satisfaite puisque
$\mathcal{G}^i$ est plat sur $f^{-1}\mathcal{O}_S$.
\end{proof}

\begin{lemma}
\label{perfect-lemma-ext-perfect}
Soit $S$ un schéma noethérien. Soit $f : X \to S$ un morphisme de schémas
localement de type fini. Soit $E \in D(\mathcal{O}_X)$ parfait.
Soit $\mathcal{G}^\bullet$ un complexe borné de
$\mathcal{O}_X$-modules cohérents, plats sur $S$ et à support propre sur $S$.
Alors $K = Rf_*R\SheafHom(E, \mathcal{G}^\bullet)$ est un objet parfait de
$D(\mathcal{O}_S)$.
\end{lemma}

\begin{proof}
Puisque $E$ est un complexe parfait, il existe un complexe parfait dual
$E^\vee$ ; voir Cohomologie, lemme \ref{cohomology-lemma-dual-perfect-complex}.
Observons que $R\SheafHom(E, \mathcal{G}^\bullet) =
E^\vee \otimes^\mathbf{L}_{\mathcal{O}_X} \mathcal{G}^\bullet$.
Le caractère parfait de $K$ résulte donc du lemme \ref{perfect-lemma-tensor-perfect}.
\end{proof}

\noindent
Nous généraliserons le lemme suivant aux morphismes plats et propres
sur des bases quelconques dans le
lemme \ref{perfect-lemma-flat-proper-perfect-direct-image-general},
et aux morphismes parfaits propres dans
Compléments sur les morphismes, lemme
\ref{more-morphisms-lemma-perfect-proper-perfect-direct-image}.

\begin{lemma}
\label{perfect-lemma-flat-proper-perfect-direct-image}
Soit $S$ un schéma noethérien. Soit $f : X \to S$ un morphisme plat et propre
de schémas. Soit $E \in D(\mathcal{O}_X)$ parfait. Alors
$Rf_*E$ est un objet parfait de $D(\mathcal{O}_S)$.
\end{lemma}

\begin{proof}
Affirmons que le lemme \ref{perfect-lemma-perfect-direct-image} s'applique.
Les conditions (1) et (2) sont immédiates. La condition (3) est locale
sur $X$. On peut donc supposer $X$ et $S$ affines et $E$
représenté par un complexe strictement parfait de $\mathcal{O}_X$-modules.
Puisque $\mathcal{O}_X$ est plat comme faisceau de $f^{-1}\mathcal{O}_S$-modules,
la condition (3) est satisfaite.
\end{proof}






\section{Une formule de projection pour Ext}
\label{perfect-section-ext}

\noindent
Le lemme \ref{perfect-lemma-compute-ext} (ou des résultats analogues dans la littérature)
est parfois utilisé pour vérifier l'un des critères d'Artin pour les
foncteurs de Quot, les schémas de Hilbert et d'autres problèmes de modules.
Supposons que $f : X \to S$ soit un morphisme propre, plat et de présentation finie
de schémas et que $E \in D(\mathcal{O}_X)$ soit parfait.
Le lemme affirme alors que
$$
\Ext^i_X(E, f^*\mathcal{F}) =
\Ext^i_S((Rf_*E^\vee)^\vee, \mathcal{F})
$$
pour $\mathcal{F}$ quasi-cohérent sur $S$.
Sous cette forme, l'énoncé ressemble à une formule de projection
pour Ext ; de fait, le résultat découle assez
facilement du lemme \ref{perfect-lemma-cohomology-base-change}.

\begin{lemma}
\label{perfect-lemma-compute-tensor-perfect}
Sous les hypothèses et avec les notations du lemme \ref{perfect-lemma-tensor-perfect},
il existe des isomorphismes fonctoriels
$$
H^i(S, K \otimes^\mathbf{L}_{\mathcal{O}_S} \mathcal{F})
\longrightarrow
H^i(X, E \otimes_{\mathcal{O}_X}^\mathbf{L}
(\mathcal{G}^\bullet \otimes_{\mathcal{O}_X} f^*\mathcal{F}))
$$
pour $\mathcal{F}$ quasi-cohérent sur $S$,
compatibles aux morphismes de connexion (voir la démonstration).
\end{lemma}

\begin{proof}
On a
$$
\mathcal{G}^\bullet \otimes_{\mathcal{O}_X}^\mathbf{L} Lf^*\mathcal{F} =
\mathcal{G}^\bullet \otimes_{f^{-1}\mathcal{O}_S}^\mathbf{L} f^{-1}\mathcal{F} =
\mathcal{G}^\bullet \otimes_{f^{-1}\mathcal{O}_S} f^{-1}\mathcal{F} =
\mathcal{G}^\bullet \otimes_{\mathcal{O}_X} f^*\mathcal{F}
$$
où la première égalité résulte du
lemme de Cohomologie \ref{cohomology-lemma-variant-derived-pullback},
la deuxième de ce que $\mathcal{G}^n$ est un $f^{-1}\mathcal{O}_S$-module plat, et
la troisième de la définition des images inverses. On obtient donc
\begin{align*}
H^i(X, E \otimes^\mathbf{L}_{\mathcal{O}_X}
(\mathcal{G}^\bullet \otimes_{\mathcal{O}_X} f^*\mathcal{F}))
& =
H^i(X, E \otimes^\mathbf{L}_{\mathcal{O}_X} \mathcal{G}^\bullet
\otimes_{\mathcal{O}_X}^\mathbf{L} Lf^*\mathcal{F}) \\
& =
H^i(S,
Rf_*(E \otimes^\mathbf{L}_{\mathcal{O}_X} \mathcal{G}^\bullet
\otimes^\mathbf{L}_{\mathcal{O}_X} Lf^*\mathcal{F})) \\
& =
H^i(S, Rf_*(E \otimes^\mathbf{L}_{\mathcal{O}_X} \mathcal{G}^\bullet)
\otimes^\mathbf{L}_{\mathcal{O}_S} \mathcal{F}) \\
& =
H^i(S, K \otimes^\mathbf{L}_{\mathcal{O}_S} \mathcal{F}) 
\end{align*}
La première égalité résulte de ce qui précède, la deuxième de Leray
(Cohomologie, lemme \ref{cohomology-lemma-before-Leray}), et
la troisième du lemme \ref{perfect-lemma-cohomology-base-change}.
L'assertion relative aux morphismes de connexion signifie ce qui suit. Étant donnée une suite
exacte courte $0 \to \mathcal{F}_1 \to \mathcal{F}_2 \to \mathcal{F}_3 \to 0$
de $\mathcal{O}_S$-modules quasi-cohérents, les isomorphismes s'insèrent dans des
diagrammes commutatifs
$$
\xymatrix{
H^i(S, K \otimes^\mathbf{L}_{\mathcal{O}_S} \mathcal{F}_3)
\ar[r] \ar[d]_\delta &
H^i(X, E \otimes^\mathbf{L}_{\mathcal{O}_X}
(\mathcal{G}^\bullet \otimes_{\mathcal{O}_X} f^*\mathcal{F}_3))
\ar[d]^\delta \\
H^{i + 1}(S, K \otimes^\mathbf{L}_{\mathcal{O}_S} \mathcal{F}_1)
\ar[r] &
H^{i + 1}(X, E \otimes^\mathbf{L}_{\mathcal{O}_X}
(\mathcal{G}^\bullet \otimes_{\mathcal{O}_X} f^*\mathcal{F}_1))
}
$$
où les morphismes de connexion proviennent du triangle distingué
$$
K \otimes^\mathbf{L}_{\mathcal{O}_S} \mathcal{F}_1 \to
K \otimes^\mathbf{L}_{\mathcal{O}_S} \mathcal{F}_2 \to
K \otimes^\mathbf{L}_{\mathcal{O}_S} \mathcal{F}_3 \to
K \otimes^\mathbf{L}_{\mathcal{O}_S} \mathcal{F}_1[1]
$$
et du triangle distingué de $D(\mathcal{O}_X)$ associé à
la suite exacte courte
$$
0 \to
\mathcal{G}^\bullet \otimes_{\mathcal{O}_X} f^*\mathcal{F}_1 \to
\mathcal{G}^\bullet \otimes_{\mathcal{O}_X} f^*\mathcal{F}_2 \to
\mathcal{G}^\bullet \otimes_{\mathcal{O}_X} f^*\mathcal{F}_3 \to 0
$$
de complexes de $\mathcal{O}_X$-modules.
Cette suite est exacte parce que $\mathcal{G}^n$ est plat sur $S$.
Nous omettons la vérification de la commutativité du diagramme affiché.
\end{proof}

\begin{lemma}
\label{perfect-lemma-compute-ext-perfect}
Sous les hypothèses et avec les notations du lemme \ref{perfect-lemma-ext-perfect},
il existe des isomorphismes fonctoriels
$$
H^i(S, K \otimes^\mathbf{L}_{\mathcal{O}_S} \mathcal{F})
\longrightarrow
\Ext^i_{\mathcal{O}_X}(E,
\mathcal{G}^\bullet \otimes_{\mathcal{O}_X} f^*\mathcal{F})
$$
pour $\mathcal{F}$ quasi-cohérent sur $S$,
compatibles aux morphismes de connexion (voir la démonstration).
\end{lemma}

\begin{proof}
Comme dans la démonstration du lemme \ref{perfect-lemma-ext-perfect}, soit $E^\vee$ le
complexe parfait dual et rappelons que
$K = Rf_*(E^\vee \otimes_{\mathcal{O}_X}^\mathbf{L} \mathcal{G}^\bullet)$.
Puisque l'on a aussi
$$
\Ext^i_{\mathcal{O}_X}(E,
\mathcal{G}^\bullet \otimes_{\mathcal{O}_X} f^*\mathcal{F})
=
H^i(X, E^\vee \otimes^\mathbf{L}_{\mathcal{O}_X}
(\mathcal{G}^\bullet \otimes_{\mathcal{O}_X} f^*\mathcal{F}))
$$
par construction de $E^\vee$, l'existence des isomorphismes résulte
du lemme \ref{perfect-lemma-compute-tensor-perfect} appliqué à $E^\vee$
et $\mathcal{G}^\bullet$.
L'assertion relative aux morphismes de connexion signifie ce qui suit. Étant donnée une suite
exacte courte $0 \to \mathcal{F}_1 \to \mathcal{F}_2 \to \mathcal{F}_3 \to 0$,
les isomorphismes s'insèrent dans des diagrammes commutatifs
$$
\xymatrix{
H^i(S, K \otimes^\mathbf{L}_{\mathcal{O}_S} \mathcal{F}_3)
\ar[r] \ar[d]_\delta &
\Ext^i_{\mathcal{O}_X}(E,
\mathcal{G}^\bullet \otimes_{\mathcal{O}_X} f^*\mathcal{F}_3) \ar[d]^\delta \\
H^{i + 1}(S, K \otimes^\mathbf{L}_{\mathcal{O}_S} \mathcal{F}_1)
\ar[r] &
\Ext^{i + 1}_{\mathcal{O}_X}(E,
\mathcal{G}^\bullet \otimes_{\mathcal{O}_X} f^*\mathcal{F}_1)
}
$$
où les morphismes de connexion proviennent du triangle distingué
$$
K \otimes^\mathbf{L}_{\mathcal{O}_S} \mathcal{F}_1 \to
K \otimes^\mathbf{L}_{\mathcal{O}_S} \mathcal{F}_2 \to
K \otimes^\mathbf{L}_{\mathcal{O}_S} \mathcal{F}_3 \to
K \otimes^\mathbf{L}_{\mathcal{O}_S} \mathcal{F}_1[1]
$$
et du triangle distingué de $D(\mathcal{O}_X)$ associé à
la suite exacte courte
$$
0 \to
\mathcal{G}^\bullet \otimes_{\mathcal{O}_X} f^*\mathcal{F}_1 \to
\mathcal{G}^\bullet \otimes_{\mathcal{O}_X} f^*\mathcal{F}_2 \to
\mathcal{G}^\bullet \otimes_{\mathcal{O}_X} f^*\mathcal{F}_3 \to 0
$$
de complexes.
Cette suite est exacte parce que $\mathcal{G}$ est plat sur $S$.
Nous omettons la vérification de la commutativité du diagramme affiché.
\end{proof}

\begin{lemma}
\label{perfect-lemma-compute-ext}
Soient $f : X \to S$ un morphisme de schémas, $E \in D(\mathcal{O}_X)$
et $\mathcal{G}^\bullet$ un complexe de $\mathcal{O}_X$-modules.
Supposons que
\begin{enumerate}
\item $S$ soit noethérien,
\item $f$ soit localement de type fini,
\item $E \in D^-_{\textit{Coh}}(\mathcal{O}_X)$,
\item $\mathcal{G}^\bullet$ soit un complexe borné de
$\mathcal{O}_X$-modules cohérents, plats sur $S$ et à support propre sur $S$.
\end{enumerate}
Alors les deux assertions suivantes sont vraies :
\begin{enumerate}
\item[(A)] pour tout $m \in \mathbf{Z}$, il existe un objet parfait $K$
de $D(\mathcal{O}_S)$ et des morphismes fonctoriels
$$
\alpha^i_\mathcal{F} :
\Ext^i_{\mathcal{O}_X}(E,
\mathcal{G}^\bullet \otimes_{\mathcal{O}_X} f^*\mathcal{F})
\longrightarrow
H^i(S, K \otimes^\mathbf{L}_{\mathcal{O}_S} \mathcal{F})
$$
pour $\mathcal{F}$ quasi-cohérent sur $S$, compatibles aux morphismes de connexion
(voir la démonstration), tels que $\alpha^i_\mathcal{F}$ soit un isomorphisme pour $i \leq m$ ;
\item[(B)] il existe un $L \in D(\mathcal{O}_S)$ pseudo-cohérent
et des isomorphismes fonctoriels
$$
\Ext^i_{\mathcal{O}_S}(L, \mathcal{F}) \longrightarrow
\Ext^i_{\mathcal{O}_X}(E,
\mathcal{G}^\bullet \otimes_{\mathcal{O}_X} f^*\mathcal{F})
$$
pour $\mathcal{F}$ quasi-cohérent sur $S$, compatibles aux morphismes de connexion.
\end{enumerate}
\end{lemma}

\begin{proof}
Démonstration de (A).
Supposons $\mathcal{G}^i$ non nul seulement pour $i \in [a, b]$.
On peut remplacer $X$ par un voisinage ouvert quasi-compact de
la réunion des supports des $\mathcal{G}^i$.
On peut donc supposer $X$ noethérien.
Dans ce cas, $X$ et $f$ sont quasi-compacts et quasi-séparés.
Choisissons une approximation $P \to E$ par un complexe parfait $P$ de
$(X, E, -m - 1 + a)$
(ce qui est possible d'après le théorème \ref{perfect-theorem-approximation}).
Alors le morphisme induit
$$
\Ext^i_{\mathcal{O}_X}(E,
\mathcal{G}^\bullet \otimes_{\mathcal{O}_X} f^*\mathcal{F})
\longrightarrow
\Ext^i_{\mathcal{O}_X}(P,
\mathcal{G}^\bullet \otimes_{\mathcal{O}_X} f^*\mathcal{F})
$$
est un isomorphisme pour $i \leq m$. En effet, le noyau, resp.\ le conoyau, de ce
morphisme est un quotient, resp.\ un sous-module, de
$$
\Ext^i_{\mathcal{O}_X}(C,
\mathcal{G}^\bullet \otimes_{\mathcal{O}_X} f^*\mathcal{F})
\quad\text{resp.}\quad
\Ext^{i + 1}_{\mathcal{O}_X}(C,
\mathcal{G}^\bullet \otimes_{\mathcal{O}_X} f^*\mathcal{F})
$$
où $C$ est le cône de $P \to E$. Puisque les faisceaux de cohomologie de $C$
sont nuls en degrés $\geq -m - 1 + a$, ces groupes $\Ext$ sont nuls
pour $i \leq m + 1$ d'après
Catégories dérivées, lemme \ref{derived-lemma-negative-exts}.
On est ainsi ramené au cas où
$E$ est un complexe parfait, qui est le lemme \ref{perfect-lemma-compute-ext-perfect}.
L'assertion relative aux morphismes de connexion est expliquée dans la démonstration du
lemme \ref{perfect-lemma-compute-ext-perfect}.

\medskip\noindent
Démonstration de (B). Comme dans la démonstration de (A), on peut supposer $X$ noethérien.
Observons que $E$ est pseudo-cohérent d'après le
lemme \ref{perfect-lemma-identify-pseudo-coherent-noetherian}.
D'après le lemme \ref{perfect-lemma-pseudo-coherent-hocolim}, on peut écrire
$E = \text{hocolim} E_n$, où $E_n$ est parfait et où $E_n \to E$ induit
un isomorphisme sur les troncatures $\tau_{\geq -n}$. Soit $E_n^\vee$
le complexe parfait dual
(Cohomologie, lemme \ref{cohomology-lemma-dual-perfect-complex}).
On obtient un système projectif $\ldots \to E_3^\vee \to E_2^\vee \to E_1^\vee$
d'objets parfaits. Celui-ci donne à son tour un système projectif
$$
\ldots \to K_3 \to K_2 \to K_1\quad\text{avec}\quad
K_n = Rf_*(E_n^\vee \otimes_{\mathcal{O}_X}^\mathbf{L} \mathcal{G}^\bullet)
$$
d'objets parfaits sur $S$ ; voir le lemme \ref{perfect-lemma-tensor-perfect}.
D'après le lemme \ref{perfect-lemma-compute-ext-perfect} et sa démonstration, ainsi que
les arguments du paragraphe précédent (avec $P = E_n$),
pour tout $\mathcal{F}$ quasi-cohérent sur $S$, on dispose de
morphismes canoniques fonctoriels
$$
\xymatrix{
& \Ext^i_{\mathcal{O}_X}(E,
\mathcal{G}^\bullet \otimes_{\mathcal{O}_X} f^*\mathcal{F})
\ar[ld] \ar[rd] \\
H^i(S, K_{n + 1} \otimes_{\mathcal{O}_S}^\mathbf{L} \mathcal{F})
\ar[rr] & &
H^i(S, K_n \otimes_{\mathcal{O}_S}^\mathbf{L} \mathcal{F})
}
$$
qui sont des isomorphismes pour $i \leq n + a$.
Soit $L_n = K_n^\vee$ le complexe parfait dual.
On voit alors que $L_1 \to L_2 \to L_3 \to \ldots$
est un système d'objets parfaits de $D(\mathcal{O}_S)$
tel que, pour tout $\mathcal{F}$ quasi-cohérent sur $S$,
les morphismes
$$
\Ext^i_{\mathcal{O}_S}(L_{n + 1}, \mathcal{F})
\longrightarrow
\Ext^i_{\mathcal{O}_S}(L_n, \mathcal{F})
$$
sont des isomorphismes pour $i \leq n + a - 1$. Il en résulte que
$L_n \to L_{n + 1}$ induit un isomorphisme sur les troncatures
$\tau_{\geq -n - a + 2}$ (indication : prendre le cône de $L_n \to L_{n + 1}$
et considérer son dernier faisceau de cohomologie non nul).
Ainsi $L = \text{hocolim} L_n$ est pseudo-cohérent ; voir le
lemme \ref{perfect-lemma-pseudo-coherent-hocolim}. La propriété universelle
des colimites homotopiques donne
$\Ext^i_{\mathcal{O}_S}(L, \mathcal{F}) =
\Ext^i_{\mathcal{O}_S}(L_n, \mathcal{F})$
pour $i \leq n + a - 3$, ce qui achève la démonstration.
\end{proof}

\begin{remark}
\label{perfect-remark-base-change-of-L}
Le complexe pseudo-cohérent $L$ de la partie (B) du lemme \ref{perfect-lemma-compute-ext}
est canoniquement associé à la situation. Par exemple,
la formation de $L$ comme en (B) est compatible au changement de base.
Autrement dit, étant donné un diagramme cartésien
$$
\xymatrix{
X' \ar[r]_{g'} \ar[d]_{f'} &
X \ar[d]^f \\
S' \ar[r]^g &
S
}
$$
de schémas, on dispose d'isomorphismes canoniques fonctoriels
$$
\Ext^i_{\mathcal{O}_{S'}}(Lg^*L, \mathcal{F}') \longrightarrow
\Ext^i_{\mathcal{O}_X}(L(g')^*E,
(g')^*\mathcal{G}^\bullet \otimes_{\mathcal{O}_{X'}} (f')^*\mathcal{F}')
$$
pour $\mathcal{F}'$ quasi-cohérent sur $S'$. Observons que nous n'utilisons {\bf pas}
l'image inverse dérivée de $\mathcal{G}^\bullet$ dans le membre de droite.
Si nous avons un jour besoin de ce fait, nous
formulerons ici un résultat précis et en donnerons une démonstration détaillée.
\end{remark}





\section{Limites et catégories dérivées}
\label{perfect-section-limits}

\noindent
Dans cette section, nous rassemblons quelques résultats sur la catégorie dérivée
d'un schéma qui est la limite d'un système projectif de schémas.
Plus précisément, nous travaillerons dans la situation suivante.

\begin{situation}
\label{perfect-situation-descent}
Soit $S = \lim_{i \in I} S_i$ la limite d'un système filtrant de schémas
dont les morphismes de transition $f_{i'i} : S_{i'} \to S_i$ sont affines.
Nous supposons que $S_i$ est quasi-compact et quasi-séparé pour tout $i \in I$.
Nous notons $f_i : S \to S_i$ la projection. Fixons également un élément $0 \in I$.
\end{situation}

\begin{lemma}
\label{perfect-lemma-descend-homomorphisms}
Dans la situation \ref{perfect-situation-descent},
soient $E_0$ et $K_0$ des objets de
$D(\mathcal{O}_{S_0})$.
Posons $E_i = Lf_{i0}^*E_0$ et $K_i = Lf_{i0}^*K_0$ pour $i \geq 0$,
et posons $E = Lf_0^*E_0$ et $K = Lf_0^*K_0$. Alors le morphisme
$$
\colim_{i \geq 0} \Hom_{D(\mathcal{O}_{S_i})}(E_i, K_i)
\longrightarrow
\Hom_{D(\mathcal{O}_S)}(E, K)
$$
est un isomorphisme dans chacun des deux cas suivants :
\begin{enumerate}
\item $E_0$ est parfait et $K_0 \in D_\QCoh(\mathcal{O}_{S_0})$ ;
\item $E_0$ est pseudo-cohérent et
$K_0 \in D_\QCoh(\mathcal{O}_{S_0})$ est de dimension Tor finie.
\end{enumerate}
\end{lemma}

\begin{proof}
Pour tout ouvert $U_0 \subset S_0$, considérons la condition $P$ selon laquelle le morphisme canonique
$$
\colim_{i \geq 0} \Hom_{D(\mathcal{O}_{U_i})}(E_i|_{U_i}, K_i|_{U_i})
\longrightarrow
\Hom_{D(\mathcal{O}_U)}(E|_U, K|_U)
$$
est un isomorphisme, où $U = f_0^{-1}(U_0)$ et $U_i = f_{i0}^{-1}(U_0)$.
Nous allons démontrer que $P$ est satisfaite pour tout ouvert quasi-compact $U_0$
par le principe d'induction de
Cohomologie des schémas, lemme \ref{coherent-lemma-induction-principle}.
La condition (2) de ce lemme résulte immédiatement de Mayer--Vietoris
pour les Hom de la catégorie dérivée ; voir
Cohomologie, lemme \ref{cohomology-lemma-mayer-vietoris-hom}.
Il suffit donc de démontrer le lemme lorsque $S_0$ est affine.

\medskip\noindent
Supposons $S_0$ affine. Écrivons $S_0 = \Spec(A_0)$, $S_i = \Spec(A_i)$ et
$S = \Spec(A)$. Nous utiliserons le lemme \ref{perfect-lemma-affine-compare-bounded}
sans le mentionner davantage.

\medskip\noindent
Dans le cas (1), l'objet $E_0^\bullet$ correspond à un complexe fini
de $A_0$-modules projectifs de type fini ; voir le lemme \ref{perfect-lemma-perfect-affine}.
On peut représenter l'objet $K_0$ par un complexe K-plat $K_0^\bullet$
de $A_0$-modules. Dans cette situation, il s'agit de démontrer que
$$
\colim_{i \geq 0} \Hom_{D(A_i)}(E_0^\bullet \otimes_{A_0} A_i,
K_0^\bullet \otimes_{A_0} A_i)
\longrightarrow
\Hom_{D(A)}(E_0^\bullet \otimes_{A_0} A, K_0^\bullet \otimes_{A_0} A)
$$
est un isomorphisme. Comme $E_0^\bullet$ est un complexe borné supérieurement de modules projectifs,
on peut réécrire ce morphisme sous la forme
$$
\colim_{i \geq 0} \Hom_{K(A_0)}(E_0^\bullet,
K_0^\bullet \otimes_{A_0} A_i)
\longrightarrow
\Hom_{K(A_0)}(E_0^\bullet, K_0^\bullet \otimes_{A_0} A)
$$
Puisqu'il n'y a qu'un nombre fini de modules non nuls
$E_0^n$ et que ceux-ci sont tous de présentation finie, ce
morphisme est un isomorphisme.

\medskip\noindent
Dans le cas (2), l'objet $E_0$ correspond à un
complexe borné supérieurement $E_0^\bullet$ de $A_0$-modules libres de type fini ;
voir le lemme \ref{perfect-lemma-pseudo-coherent-affine}.
On peut représenter $K_0$ par un complexe fini $K_0^\bullet$
de $A_0$-modules plats ; voir le lemme \ref{perfect-lemma-tor-dimension-affine}
et
Compléments sur l'algèbre, lemme \ref{more-algebra-lemma-tor-amplitude}.
En particulier, $K_0^\bullet$ est K-plat, et l'on peut raisonner comme précédemment
pour obtenir le morphisme
$$
\colim_{i \geq 0} \Hom_{K(A_0)}(E_0^\bullet,
K_0^\bullet \otimes_{A_0} A_i)
\longrightarrow
\Hom_{K(A_0)}(E_0^\bullet, K_0^\bullet \otimes_{A_0} A)
$$
Il est clair que ce morphisme est un isomorphisme (seul un nombre fini de
termes intervient puisque $K_0^\bullet$ est borné).
\end{proof}

\begin{lemma}
\label{perfect-lemma-descend-perfect}
Dans la situation \ref{perfect-situation-descent}, la catégorie des objets parfaits
de $D(\mathcal{O}_S)$ est la colimite des catégories
d'objets parfaits de $D(\mathcal{O}_{S_i})$.
\end{lemma}

\begin{proof}
Pour tout ouvert $U_0 \subset S_0$, considérons la condition $P$ selon laquelle
le foncteur
$$
\colim_{i \geq 0} D_{perf}(\mathcal{O}_{U_i})
\longrightarrow
D_{perf}(\mathcal{O}_U)
$$
est une équivalence, où ${}_{perf}$ désigne la sous-catégorie pleine des
objets parfaits et où $U = f_0^{-1}(U_0)$ et $U_i = f_{i0}^{-1}(U_0)$.
Nous allons démontrer que $P$ est satisfaite pour tout ouvert quasi-compact $U_0$
par le principe d'induction de
Cohomologie des schémas, lemme \ref{coherent-lemma-induction-principle}.
D'abord, observons que le foncteur est déjà pleinement fidèle
d'après le lemme \ref{perfect-lemma-descend-homomorphisms}. Il suffit donc de démontrer
qu'il est essentiellement surjectif.

\medskip\noindent
Vérifions d'abord la condition (2) du principe d'induction. Supposons donc
que $S_0 = U_0 \cup V_0$ et que $P$ soit satisfaite pour
$U_0$, $V_0$ et $U_0 \cap V_0$. Soit $E$ un objet parfait
de $D(\mathcal{O}_S)$. On peut trouver $i \geq 0$, un objet $E_{U, i}$ parfait sur $U_i$
et un objet $E_{V, i}$ parfait sur $V_i$, dont les images inverses sur $U$ et $V$ sont isomorphes
à $E|_U$ et $E|_V$. Notons
$$
a : E_{U, i} \to (Rf_{i, *}E)|_{U_i}
\quad\text{et}\quad
b : E_{V, i} \to (Rf_{i, *}E)|_{V_i}
$$
les morphismes adjoints aux isomorphismes $Lf_i^*E_{U, i} \to E|_U$
et $Lf_i^*E_{V, i} \to E|_V$.
Par pleine fidélité, quitte à augmenter $i$,
on peut trouver un isomorphisme
$c : E_{U, i}|_{U_i \cap V_i} \to E_{V, i}|_{U_i \cap V_i}$
qui, par image inverse, donne les identifications
$$
Lf_i^*E_{U, i}|_{U \cap V} \to E|_{U \cap V} \to Lf_i^*E_{V, i}|_{U \cap V}.
$$
Appliquons Cohomologie, lemme \ref{cohomology-lemma-glue},
pour obtenir un objet $E_i$ sur $S_i$ et un morphisme $d : E_i \to Rf_{i, *}E$
dont les restrictions sont les morphismes $a$ et $b$ sur $U_i$ et $V_i$.
Il est alors clair que $E_i$ est parfait et que
$d$ est adjoint à un isomorphisme $Lf_i^*E_i \to E$.

\medskip\noindent
Enfin, vérifions la condition (1) du principe d'induction, autrement
dit vérifions le lemme lorsque $S_0$ est affine.
Écrivons $S_0 = \Spec(A_0)$, $S_i = \Spec(A_i)$ et
$S = \Spec(A)$. Les lemmes \ref{perfect-lemma-affine-compare-bounded}
et \ref{perfect-lemma-perfect-affine} montrent qu'il faut établir que
$$
D_{perf}(A) = \colim D_{perf}(A_i)
$$
Cela résulte du fait que les complexes parfaits sur des anneaux sont
donnés par des complexes finis de modules projectifs de type fini (donc de présentation finie).
Voir Compléments sur l'algèbre, lemme
\ref{more-algebra-lemma-colimit-perfect-complexes}, pour les détails.
\end{proof}





\section{Cohomologie et changement de base, VI}
\label{perfect-section-cohomology-and-base-change-final}

\noindent
Une dernière section sur la cohomologie et le changement de base, qui poursuit
la discussion des sections
\ref{perfect-section-cohomology-and-base-change-perfect},
\ref{perfect-section-cohomology-base-change} et
\ref{perfect-section-producing-perfect}.
Un cas particulier facile à saisir est donné dans la
remarque \ref{perfect-remark-explain-perfect-direct-image}.

\begin{lemma}
\label{perfect-lemma-base-change-tensor-perfect}
Soit $f : X \to S$ un morphisme de présentation finie.
Soit $E \in D(\mathcal{O}_X)$ un objet parfait. Soit $\mathcal{G}^\bullet$
un complexe borné de $\mathcal{O}_X$-modules de présentation finie,
plats sur $S$ et à support propre sur $S$. Alors
$$
K = Rf_*(E \otimes_{\mathcal{O}_X}^\mathbf{L} \mathcal{G}^\bullet)
$$
est un objet parfait de $D(\mathcal{O}_S)$ et sa formation
commute à tout changement de base.
\end{lemma}

\begin{proof}
L'assertion relative au changement de base est le lemme \ref{perfect-lemma-base-change-tensor}.
Il suffit donc de montrer que $K$ est un objet parfait. Si $S$ est
noethérien, cela résulte du
lemme \ref{perfect-lemma-tensor-perfect}.
Nous allons nous ramener à ce cas par approximation noethérienne.
Nous conseillons au lecteur de sauter le reste de cette démonstration.

\medskip\noindent
La question est locale sur $S$ ; on peut donc supposer $S$ affine.
Écrivons $S = \Spec(R)$. Écrivons $R = \colim R_i$ comme colimite filtrante
d'anneaux noethériens $R_i$. D'après Limites, lemme
\ref{limits-lemma-descend-finite-presentation},
il existe un $i$ et un schéma $X_i$ de présentation finie sur $R_i$
dont le changement de base à $R$ est $X$. D'après
Limites, lemme \ref{limits-lemma-descend-modules-finite-presentation},
quitte à augmenter $i$, on peut supposer qu'il existe un complexe borné
de $\mathcal{O}_{X_i}$-modules de présentation finie
$\mathcal{G}_i^\bullet$ dont l'image inverse sur $X$ est $\mathcal{G}^\bullet$.
Quitte à augmenter $i$, on peut supposer $\mathcal{G}_i^n$ plat sur $R_i$ ; voir
Limites, lemme \ref{limits-lemma-descend-module-flat-finite-presentation}.
Quitte à augmenter $i$, on peut supposer que le support de $\mathcal{G}_i^n$
est propre sur $R_i$ ; voir
Limites, lemme \ref{limits-lemma-eventually-proper-support},
et Cohomologie des schémas, lemme
\ref{coherent-lemma-module-support-proper-over-base}.
Enfin, d'après le lemme \ref{perfect-lemma-descend-perfect},
quitte à augmenter $i$, on peut supposer qu'il existe un objet parfait
$E_i$ de $D(\mathcal{O}_{X_i})$ dont l'image inverse sur
$X$ est $E$. En appliquant le lemme \ref{perfect-lemma-tensor-perfect}
à $X_i \to \Spec(R_i)$, $E_i$, $\mathcal{G}_i^\bullet$, et en utilisant la
propriété de changement de base déjà démontrée, on obtient le résultat.
\end{proof}

\begin{remark}
\label{perfect-remark-explain-perfect-direct-image}
Soit $R$ un anneau. Soit $X$ un schéma de présentation finie sur
$R$. Soit $\mathcal{G}$ un $\mathcal{O}_X$-module de présentation finie,
plat sur $R$ et à support propre sur $R$. D'après le
lemme \ref{perfect-lemma-base-change-tensor-perfect},
il existe un complexe fini de $R$-modules projectifs de type fini
$M^\bullet$ tel que l'on ait
$$
R\Gamma(X_{R'}, \mathcal{G}_{R'}) = M^\bullet \otimes_R R'
$$
fonctoriellement en la $R$-algèbre $R'$.
\end{remark}

\begin{lemma}
\label{perfect-lemma-base-change-tensor-pseudo-coherent}
Soit $f : X \to S$ un morphisme de présentation finie.
Soit $E \in D(\mathcal{O}_X)$ un objet pseudo-cohérent.
Soit $\mathcal{G}^\bullet$ un complexe borné supérieurement de
$\mathcal{O}_X$-modules de présentation finie, plats sur $S$
et à support propre sur $S$. Alors
$$
K = Rf_*(E \otimes_{\mathcal{O}_X}^\mathbf{L} \mathcal{G}^\bullet)
$$
est un objet pseudo-cohérent de $D(\mathcal{O}_S)$ et sa formation
commute à tout changement de base.
\end{lemma}

\begin{proof}
L'assertion relative au changement de base est le lemme \ref{perfect-lemma-base-change-tensor}.
Il suffit donc de montrer que $K$ est un objet pseudo-cohérent.
Cela résultera du lemme \ref{perfect-lemma-base-change-tensor-perfect}
par approximation par des complexes parfaits. Nous conseillons au lecteur de
sauter le reste de la démonstration.

\medskip\noindent
La question est locale sur $S$ ; on peut donc supposer $S$ affine.
Alors $X$ est quasi-compact et quasi-séparé. De plus, il
existe un entier $N$ tel que l'image directe totale
$Rf_* : D_\QCoh(\mathcal{O}_X) \to D_\QCoh(\mathcal{O}_S)$
soit de dimension cohomologique $N$, comme expliqué dans le
lemme \ref{perfect-lemma-quasi-coherence-direct-image}.
Choisissons un entier $b$ tel que $\mathcal{G}^i = 0$ pour $i > b$.
Il suffit de montrer que $K$ est $m$-pseudo-cohérent pour
tout $m$. Choisissons une approximation $P \to E$ par un complexe parfait $P$
de $(X, E, m - N - 1 - b)$. Cela est possible d'après le
théorème \ref{perfect-theorem-approximation}.
Choisissons un triangle distingué
$$
P \to E \to C \to P[1]
$$
dans $D_\QCoh(\mathcal{O}_X)$. Les faisceaux de cohomologie de $C$ sont nuls
en degrés $\geq m - N - 1 - b$. Les faisceaux de cohomologie de
$C \otimes^\mathbf{L} \mathcal{G}^\bullet$ sont donc nuls en degrés
$\geq m - N - 1$. Ainsi, les faisceaux de cohomologie de
$Rf_*(C \otimes^\mathbf{L} \mathcal{G}^\bullet)$
sont nuls en degrés $\geq m - 1$.
Par conséquent,
$$
Rf_*(P \otimes^\mathbf{L} \mathcal{G}^\bullet) \to
Rf_*(E \otimes^\mathbf{L} \mathcal{G}^\bullet)
$$
est un isomorphisme sur les faisceaux de cohomologie en degrés $\geq m$.
Supposons ensuite que $H^i(P) = 0$ pour $i > a$. Alors
$
P \otimes^\mathbf{L} \sigma_{\geq m - N - 1 - a}\mathcal{G}^\bullet
\longrightarrow
P \otimes^\mathbf{L} \mathcal{G}^\bullet
$
est un isomorphisme sur les faisceaux de cohomologie en degrés $\geq m - N - 1$.
On trouve donc, de nouveau, que
$$
Rf_*(P \otimes^\mathbf{L} \sigma_{\geq m - N - 1 - a}\mathcal{G}^\bullet) \to
Rf_*(P \otimes^\mathbf{L} \mathcal{G}^\bullet)
$$
est un isomorphisme sur les faisceaux de cohomologie en degrés $\geq m$.
D'après le lemme \ref{perfect-lemma-base-change-tensor-perfect}, la source
est un complexe parfait.
On conclut que $K$ est $m$-pseudo-cohérent, comme voulu.
\end{proof}

\begin{lemma}
\label{perfect-lemma-flat-proper-perfect-direct-image-general}
Soit $S$ un schéma. Soit $f : X \to S$ un morphisme propre
de présentation finie.
\begin{enumerate}
\item Soit $E \in D(\mathcal{O}_X)$ parfait et supposons $f$ plat. Alors
$Rf_*E$ est un objet parfait de $D(\mathcal{O}_S)$ et sa formation
commute à tout changement de base.
\item Soit $\mathcal{G}$ un $\mathcal{O}_X$-module de présentation finie,
plat sur $S$. Alors $Rf_*\mathcal{G}$ est un objet parfait de
$D(\mathcal{O}_S)$ et sa formation commute à tout changement de base.
\end{enumerate}
\end{lemma}

\begin{proof}
Ce sont des cas particuliers du
lemme \ref{perfect-lemma-base-change-tensor-perfect}, appliqué avec
(1) $\mathcal{G}^\bullet$ égal à $\mathcal{O}_X$ placé en degré $0$
et (2) $E = \mathcal{O}_X$ et $\mathcal{G}^\bullet$ constitué de
$\mathcal{G}$ placé en degré $0$.
\end{proof}

\begin{lemma}
\label{perfect-lemma-flat-proper-pseudo-coherent-direct-image-general}
Soit $S$ un schéma. Soit $f : X \to S$ un morphisme plat et propre
de présentation finie. Soit $E \in D(\mathcal{O}_X)$
pseudo-cohérent. Alors $Rf_*E$ est un objet pseudo-cohérent de
$D(\mathcal{O}_S)$ et sa formation commute à tout changement de base.
\end{lemma}

\noindent
Plus généralement, si $f : X \to S$ est propre et si $E$ sur $X$ est pseudo-cohérent
relativement à $S$ (Compléments sur les morphismes, définition
\ref{more-morphisms-definition-relative-pseudo-coherence}),
alors $Rf_*E$ est pseudo-cohérent
(mais sa formation ne commute pas au changement de base dans cette généralité).
Voir \cite{Kiehl}.

\begin{proof}
C'est un cas particulier du
lemme \ref{perfect-lemma-base-change-tensor-pseudo-coherent}, appliqué avec
$\mathcal{G}^\bullet$ égal à $\mathcal{O}_X$ placé en degré $0$.
\end{proof}

\begin{lemma}
\label{perfect-lemma-pullback-and-limits}
Soit $R$ un anneau. Soit $X$ un schéma, et soit
$f : X \to \Spec(R)$ un morphisme propre, plat et
de présentation finie. Soit $(M_n)$ un système projectif
de $R$-modules dont les morphismes de transition sont surjectifs.
Alors le morphisme canonique
$$
\mathcal{O}_X \otimes_R (\lim M_n)
\longrightarrow
\lim \mathcal{O}_X \otimes_R M_n
$$
induit un isomorphisme de la source vers l'image du but par $DQ_X$.
\end{lemma}

\begin{proof}
L'énoncé signifie que, pour tout objet $E$ de
$D_\QCoh(\mathcal{O}_X)$, le morphisme induit
$$
\Hom(E, \mathcal{O}_X \otimes_R (\lim M_n))
\longrightarrow
\Hom(E, \lim \mathcal{O}_X \otimes_R M_n)
$$
est un isomorphisme. Puisque $D_\QCoh(\mathcal{O}_X)$ possède
un générateur parfait (théorème \ref{perfect-theorem-bondal-van-den-Bergh}),
il suffit de le vérifier pour $E$ parfait.
D'après le lemme \ref{perfect-lemma-Rlim-quasi-coherent}, on a
$\lim \mathcal{O}_X \otimes_R M_n = R\lim \mathcal{O}_X \otimes_R M_n$.
Le foncteur exact
$R\Hom_X(E, -) : D_\QCoh(\mathcal{O}_X) \to D(R)$
de Cohomologie, section \ref{cohomology-section-global-RHom},
commute aux produits et donc aux limites dérivées, d'où
$$
R\Hom_X(E, \lim \mathcal{O}_X \otimes_R M_n) =
R\lim R\Hom_X(E, \mathcal{O}_X \otimes_R M_n)
$$
Soit $E^\vee$ le complexe parfait dual ; voir
Cohomologie, lemme \ref{cohomology-lemma-dual-perfect-complex}.
On a
$$
R\Hom_X(E, \mathcal{O}_X \otimes_R M_n) =
R\Gamma(X, E^\vee \otimes_{\mathcal{O}_X}^\mathbf{L} Lf^*M_n) =
R\Gamma(X, E^\vee) \otimes_R^\mathbf{L} M_n
$$
d'après le lemme \ref{perfect-lemma-cohomology-base-change}.
Le lemme \ref{perfect-lemma-flat-proper-perfect-direct-image-general}
montre que $R\Gamma(X, E^\vee)$ est un complexe parfait de $R$-modules.
En particulier, c'est un complexe pseudo-cohérent, et le
lemme de Compléments sur l'algèbre \ref{more-algebra-lemma-pseudo-coherent-tensor-limit}
donne
$$
R\lim R\Gamma(X, E^\vee) \otimes_R^\mathbf{L} M_n =
R\Gamma(X, E^\vee) \otimes_R^\mathbf{L} \lim M_n
$$
comme voulu.
\end{proof}

\begin{lemma}
\label{perfect-lemma-base-change-RHom-perfect}
Soit $f : X \to S$ un morphisme de présentation finie.
Soit $E \in D(\mathcal{O}_X)$ un objet parfait. Soit $\mathcal{G}^\bullet$
un complexe borné de $\mathcal{O}_X$-modules de présentation finie,
plats sur $S$ et à support propre sur $S$. Alors
$$
K = Rf_*R\SheafHom(E, \mathcal{G}^\bullet)
$$
est un objet parfait de $D(\mathcal{O}_S)$ et sa formation
commute à tout changement de base.
\end{lemma}

\begin{proof}
L'assertion relative au changement de base est le lemme \ref{perfect-lemma-base-change-RHom}.
Il suffit donc de montrer que $K$ est un objet parfait. Si $S$ est
noethérien, cela résulte du
lemme \ref{perfect-lemma-ext-perfect}.
Nous allons nous ramener à ce cas par approximation noethérienne.
Nous conseillons au lecteur de sauter le reste de cette démonstration.

\medskip\noindent
La question est locale sur $S$ ; on peut donc supposer $S$ affine.
Écrivons $S = \Spec(R)$. Écrivons $R = \colim R_i$ comme colimite filtrante
d'anneaux noethériens $R_i$. D'après Limites, lemme
\ref{limits-lemma-descend-finite-presentation},
il existe un $i$ et un schéma $X_i$ de présentation finie sur $R_i$
dont le changement de base à $R$ est $X$. D'après
Limites, lemme \ref{limits-lemma-descend-modules-finite-presentation},
quitte à augmenter $i$, on peut supposer qu'il existe un complexe borné
de $\mathcal{O}_{X_i}$-modules de présentation finie $\mathcal{G}_i^\bullet$
dont l'image inverse sur $X$ est $\mathcal{G}^\bullet$. Quitte à augmenter $i$,
on peut supposer $\mathcal{G}_i^n$ plat sur $R_i$ ; voir
Limites, lemme \ref{limits-lemma-descend-module-flat-finite-presentation}.
Quitte à augmenter $i$, on peut supposer que le support de $\mathcal{G}_i^n$
est propre sur $R_i$ ; voir
Limites, lemme \ref{limits-lemma-eventually-proper-support}, et
Cohomologie des schémas, lemme
\ref{coherent-lemma-module-support-proper-over-base}.
Enfin, d'après le lemme \ref{perfect-lemma-descend-perfect},
quitte à augmenter $i$, on peut supposer qu'il existe un objet parfait
$E_i$ de $D(\mathcal{O}_{X_i})$ dont l'image inverse sur
$X$ est $E$. En appliquant le lemme \ref{perfect-lemma-ext-perfect}
à $X_i \to \Spec(R_i)$, $E_i$, $\mathcal{G}_i^\bullet$, et en utilisant la
propriété de changement de base déjà démontrée, on obtient le résultat.
\end{proof}









\section{Complexes parfaits}
\label{perfect-section-perfect-complexes}

\noindent
Nous parlons d'abord des lieux de saut des nombres de Betti des complexes parfaits.
Étant donnés un complexe $E$ sur un schéma $X$ et un point $x$ de $X$, nous écrivons souvent
$E \otimes_{\mathcal{O}_X}^\mathbf{L} \kappa(x)$ au lieu de l'expression plus correcte
$Li_x^*E$, où $i_x : x \to X$ est le morphisme canonique.

\begin{lemma}
\label{perfect-lemma-jump-loci}
Soit $X$ un schéma. Soit $E \in D(\mathcal{O}_X)$ pseudo-cohérent
(par exemple parfait). Pour tout $i \in \mathbf{Z}$, considérons la fonction
$$
\beta_i : X \longrightarrow \{0, 1, 2, \ldots\},\quad
x \longmapsto
\dim_{\kappa(x)}
H^i(E \otimes_{\mathcal{O}_X}^\mathbf{L} \kappa(x))
$$
Alors :
\begin{enumerate}
\item la formation de $\beta_i$ commute à tout changement de base ;
\item les fonctions $\beta_i$ sont semi-continues supérieurement ;
\item les ensembles de niveau de $\beta_i$ sont localement constructibles dans $X$.
\end{enumerate}
\end{lemma}

\begin{proof}
Considérons un morphisme de schémas $f : Y \to X$ et un point $y \in Y$.
Soit $x$ l'image de $y$, et considérons le diagramme commutatif
$$
\xymatrix{
y \ar[r]_j \ar[d]_g & Y \ar[d]^f \\
x \ar[r]^i & X
}
$$
On voit que $Lg^* \circ Li^* = Lj^* \circ Lf^*$. Il en résulte que
la fonction $\beta'_i$ associée au complexe pseudo-cohérent $Lf^*E$
est l'image inverse de la fonction $\beta_i$, c'est-à-dire
$\beta'_i = \beta_i \circ f$. Tel est le sens de (1).

\medskip\noindent
Fixons $i$ et soit $x \in X$. Il suffit de démontrer (2) et (3)
dans un voisinage ouvert de $x$ ; on peut donc supposer $X$ affine.
On peut alors représenter $E$ par un complexe borné supérieurement $\mathcal{F}^\bullet$
de modules libres de type fini (lemme \ref{perfect-lemma-lift-pseudo-coherent}).
Alors $P = \sigma_{\geq i - 1}\mathcal{F}^\bullet$ est un objet parfait,
et $P \to E$ induit un isomorphisme
$$
H^i(P \otimes_{\mathcal{O}_X}^\mathbf{L} \kappa(x')) \to
H^i(E \otimes_{\mathcal{O}_X}^\mathbf{L} \kappa(x'))
$$
pour tout $x' \in X$. On peut donc supposer $E$ parfait. Dans ce cas,
d'après Compléments sur l'algèbre, lemme
\ref{more-algebra-lemma-lift-perfect-from-residue-field},
il existe un voisinage ouvert affine $U$ de $x$ et
$a \leq b$ tels que $E|_U$ soit représenté par un complexe
$$
\ldots \to 0 \to \mathcal{O}_U^{\oplus \beta_a(x)}
\to \mathcal{O}_U^{\oplus \beta_{a + 1}(x)} \to
\ldots \to
\mathcal{O}_U^{\oplus \beta_{b - 1}(x)} \to
\mathcal{O}_U^{\oplus \beta_b(x)} \to 0
\to \ldots
$$
(On utilise également ici des résultats antérieurs pour ramener le problème à l'algèbre, par exemple les
lemmes \ref{perfect-lemma-affine-compare-bounded} et
\ref{perfect-lemma-perfect-affine}.)
Il en résulte immédiatement que $\beta_i(x') \leq \beta_i(x)$
pour tout $x' \in U$. Cela démontre que $\beta_i$ est semi-continue
supérieurement.

\medskip\noindent
Pour démontrer (3), on peut supposer que $X$ est affine et que
$E$ est donné par un complexe de
$\mathcal{O}_X$-modules libres de type fini (par exemple en raisonnant comme au paragraphe
précédent, ou en utilisant Cohomologie, lemme
\ref{cohomology-lemma-perfect-on-locally-ringed}).
Il faut donc montrer qu'étant donné un complexe
$$
\mathcal{O}_X^{\oplus a} \to
\mathcal{O}_X^{\oplus b} \to
\mathcal{O}_X^{\oplus c}
$$
la fonction qui associe à un point $x \in X$ la dimension de la cohomologie
au milieu de $\kappa_x^{\oplus a} \to \kappa_x^{\oplus b} \to \kappa_x^{\oplus c}$
a des ensembles de niveau constructibles. Soit
$A \in \text{Mat}(a \times b, \Gamma(X, \mathcal{O}_X))$ la matrice
de la première flèche. Le rang de l'image de $A$ dans
$\text{Mat}(a \times b, \kappa(x))$ est égal à $r$ si tous les
mineurs $(r + 1) \times (r + 1)$ de $A$ s'annulent en $x$ et s'il existe un
mineur $r \times r$ de $A$ qui ne s'annule pas en $x$. L'ensemble
des points où le rang vaut $r$ est donc un ensemble localement fermé constructible.
En raisonnant de même pour la seconde flèche et en réunissant le tout,
on obtient le résultat voulu.
\end{proof}

\begin{lemma}
\label{perfect-lemma-chi-locally-constant}
Soit $X$ un schéma. Soit $E \in D(\mathcal{O}_X)$ parfait.
La fonction
$$
\chi_E : X \longrightarrow \mathbf{Z},\quad
x \longmapsto \sum (-1)^i
\dim_{\kappa(x)} H^i(E \otimes_{\mathcal{O}_X}^\mathbf{L} \kappa(x))
$$
est localement constante sur $X$.
\end{lemma}

\begin{proof}
D'après Cohomologie, lemme
\ref{cohomology-lemma-perfect-on-locally-ringed},
on peut, localement sur $X$, représenter $E$ par un complexe fini
$\mathcal{E}^\bullet$ de $\mathcal{O}_X$-modules libres de type fini.
Sur un tel ouvert, la fonction $\chi_E$ est constante de valeur
$\sum (-1)^i \text{rank}(\mathcal{E}^i)$.
\end{proof}

\begin{lemma}
\label{perfect-lemma-open-where-cohomology-in-degree-i-rank-r}
Soit $X$ un schéma. Soit $E \in D(\mathcal{O}_X)$ parfait.
Étant donnés $i, r \in \mathbf{Z}$, il existe un
sous-schéma ouvert $U \subset X$ caractérisé par les propriétés suivantes :
\begin{enumerate}
\item $E|_U \cong H^i(E|_U)[-i]$ et $H^i(E|_U)$ est un
$\mathcal{O}_U$-module localement libre de rang $r$ ;
\item un morphisme $f : Y \to X$ se factorise par $U$ si et seulement si
$Lf^*E$ est isomorphe à un module localement libre de rang $r$
placé en degré $i$.
\end{enumerate}
\end{lemma}

\begin{proof}
Soient $\beta_j : X \to \{0, 1, 2, \ldots\}$, pour $j \in \mathbf{Z}$,
les fonctions du lemme \ref{perfect-lemma-jump-loci}. Alors l'ensemble
$$
W = \{x \in X \mid \beta_j(x) \leq 0\text{ pour tout }j \not = i\}
$$
est ouvert dans $X$, et sa formation commute à l'image inverse sur tout
$Y$ au-dessus de $X$. Cela résulte du lemme et du fait que,
a priori, dans un voisinage de tout point, seul un nombre fini
des $\beta_j$ est non nul. On peut donc remplacer $X$ par $W$
et supposer que $\beta_j(x) = 0$ pour tout $x \in X$ et tout $j \not = i$.
Dans ce cas, $H^i(E)$ est un module localement libre de type fini et
$E \cong H^i(E)[-i]$ ; voir par exemple
Compléments sur l'algèbre, lemme
\ref{more-algebra-lemma-lift-perfect-from-residue-field}.
Ainsi, $X$ est la réunion disjointe des sous-schémas ouverts sur lesquels le
rang de $H^i(E)$ est fixé, ce qui conclut.
\end{proof}

\begin{lemma}
\label{perfect-lemma-locally-closed-where-H0-locally-free}
Soit $X$ un schéma. Soit $E \in D(\mathcal{O}_X)$ parfait
d'amplitude Tor contenue dans $[a, b]$ pour certains $a, b \in \mathbf{Z}$.
Soit $r \geq 0$.
Alors il existe un sous-schéma localement fermé $j : Z \to X$
caractérisé par les propriétés suivantes :
\begin{enumerate}
\item $H^a(Lj^*E)$ est un $\mathcal{O}_Z$-module localement libre de rang $r$ ;
\item un morphisme $f : Y \to X$ se factorise par $Z$
si et seulement si, pour tout morphisme $g : Y' \to Y$, le
$\mathcal{O}_{Y'}$-module $H^a(L(f \circ g)^*E)$ est localement libre
de rang $r$.
\end{enumerate}
De plus, $j : Z \to X$ est de présentation finie, et l'on a :
\begin{enumerate}
\item[(3)] si $f : Y \to X$ se factorise sous la forme $Y \xrightarrow{g} Z \to X$, alors
$H^a(Lf^*E) = g^*H^a(Lj^*E)$ ;
\item[(4)] si $\beta_a(x) \leq r$ pour tout $x \in X$, alors
$j$ est une immersion fermée et, pour tout $f : Y \to X$, les assertions suivantes
sont équivalentes :
\begin{enumerate}
\item $f : Y \to X$ se factorise par $Z$ ;
\item $H^a(Lf^*E)$ est un $\mathcal{O}_Y$-module localement libre de rang $r$ ;
\end{enumerate}
et, si $r = 1$, elles sont encore équivalentes à
\begin{enumerate}
\item[(c)] $\mathcal{O}_Y \to \SheafHom_{\mathcal{O}_Y}(H^a(Lf^*E), H^a(Lf^*E))$
est injectif.
\end{enumerate}
\end{enumerate}
\end{lemma}

\begin{proof}
Tout d'abord, soit $U \subset X$ le sous-schéma ouvert
localement constructible où la fonction
$\beta_a$ du lemme \ref{perfect-lemma-jump-loci} prend des valeurs $\leq r$.
Soit $f : Y \to X$ comme en (2). Alors, pour tout $y \in Y$,
on a $\beta_a(Lf^*E) = r$ ; le point $y$ est donc envoyé dans $U$ d'après le
lemme \ref{perfect-lemma-jump-loci}. Par conséquent, $f$ comme en (2) se factorise par $U$.
On peut ainsi remplacer $X$ par $U$ et supposer que
$\beta_a(x) \in \{0, 1, \ldots, r\}$ pour tout $x \in X$.
Nous allons montrer que, dans ce cas,
il existe un sous-schéma fermé $Z \subset X$ défini par un idéal
quasi-cohérent de type fini, caractérisé
par l'équivalence de (4)(a), (b), et de (4)(c) si $r = 1$,
et que (3) est satisfaite.
Cela achèvera la démonstration, car il en résultera a fortiori
que les morphismes comme en (2) se factorisent par $Z$.

\medskip\noindent
Si $x \in X$ et $\beta_a(x) < r$, il existe un
voisinage ouvert de $x$ sur lequel $\beta_a < r$
(lemme \ref{perfect-lemma-jump-loci}). On voit ainsi que,
au moins ensemblistement, $Z$ est un sous-ensemble fermé.

\medskip\noindent
Pour obtenir une structure schématique, considérons un point $x \in X$
tel que $\beta_a(x) = r$. Posons $\beta = \beta_{a + 1}(x)$.
D'après Compléments sur l'algèbre, lemme
\ref{more-algebra-lemma-lift-perfect-from-residue-field},
il existe un voisinage ouvert affine $U$ de $x$
tel que $K|_U$ soit représenté par un complexe
$$
\ldots \to 0 \to \mathcal{O}_U^{\oplus r}
\xrightarrow{(f_{ij})} \mathcal{O}_U^{\oplus \beta} \to
\ldots \to
\mathcal{O}_U^{\oplus \beta_{b - 1}(x)} \to
\mathcal{O}_U^{\oplus \beta_b(x)} \to 0
\to \ldots
$$
(On utilise également ici des résultats antérieurs pour ramener le problème à l'algèbre, par exemple les
lemmes \ref{perfect-lemma-affine-compare-bounded} et
\ref{perfect-lemma-perfect-affine}.) Or, si $g : Y \to U$ est un morphisme quelconque
de schémas tel que $g^\sharp(f_{ij})$ soit non nul pour un certain couple $i, j$, alors
$H^a(Lg^*E)$ n'est pas un $\mathcal{O}_Y$-module localement libre de rang $r$.
Voir Compléments sur l'algèbre, lemme \ref{more-algebra-lemma-coker-injective-free}.
Il est immédiat que $H^a(Lg^*E)$ est un $\mathcal{O}_Y$-module localement libre si
$g^\sharp(f_{ij}) = 0$ pour tous $i, j$. On voit donc que, sur $U$, le
sous-schéma fermé défini par tous les $f_{ij}$ satisfait
(3), et que l'on a l'équivalence de (4)(a) et (b).
La caractérisation de $Z$ montre que les morceaux construits localement
se recollent (détails omis). Enfin, si $r = 1$, alors
(4)(c) équivaut à (4)(b), car dans ce cas, localement,
$H^a(Lg^*E) \subset \mathcal{O}_Y$ est l'annulateur de l'idéal engendré
par les éléments $g^\sharp(f_{ij})$.
\end{proof}





\section{Applications}
\label{perfect-section-applications}

\noindent
Principalement des applications de la cohomologie et du changement de base. À l'avenir, nous pourrons
généraliser ces résultats à la situation examinée dans le
lemme \ref{perfect-lemma-base-change-tensor-perfect}.

\begin{lemma}
\label{perfect-lemma-jump-loci-geometric}
Soit $f : X \to S$ un morphisme propre de présentation finie.
Soit $\mathcal{F}$ un $\mathcal{O}_X$-module de présentation finie,
plat sur $S$. Pour $i \in \mathbf{Z}$ fixé, considérons la fonction
$$
\beta_i : S \to \{0, 1, 2, \ldots\},\quad
s \longmapsto \dim_{\kappa(s)} H^i(X_s, \mathcal{F}_s)
$$
Alors :
\begin{enumerate}
\item la formation de $\beta_i$ commute à tout changement de base ;
\item les fonctions $\beta_i$ sont semi-continues supérieurement ;
\item les ensembles de niveau de $\beta_i$ sont localement constructibles dans $S$.
\end{enumerate}
\end{lemma}

\begin{proof}
D'après la cohomologie et le changement de base (plus précisément d'après le
lemme \ref{perfect-lemma-flat-proper-perfect-direct-image-general}),
l'objet $K = Rf_*\mathcal{F}$ est un objet parfait de la catégorie dérivée
de $S$, dont la formation commute à tout changement de base.
En particulier, on a
$$
H^i(X_s, \mathcal{F}_s) = H^i(K \otimes_{\mathcal{O}_S}^\mathbf{L} \kappa(s))
$$
Le lemme résulte donc du
lemme \ref{perfect-lemma-jump-loci}.
\end{proof}

\begin{lemma}
\label{perfect-lemma-chi-locally-constant-geometric}
Soit $f : X \to S$ un morphisme propre de présentation finie.
Soit $\mathcal{F}$ un $\mathcal{O}_X$-module de présentation finie,
plat sur $S$. La fonction
$$
s \longmapsto \chi(X_s, \mathcal{F}_s)
$$
est localement constante sur $S$. La formation de cette fonction commute au
changement de base.
\end{lemma}

\begin{proof}
D'après la cohomologie et le changement de base (plus précisément d'après le
lemme \ref{perfect-lemma-flat-proper-perfect-direct-image-general}),
l'objet $K = Rf_*\mathcal{F}$ est un objet parfait de la catégorie dérivée
de $S$, dont la formation commute à tout changement de base.
Il faut donc montrer que l'application
$$
s \longmapsto \sum (-1)^i \dim_{\kappa(s)}
H^i(K \otimes^\mathbf{L}_{\mathcal{O}_S} \kappa(s))
$$
est localement constante sur $S$. C'est le lemme \ref{perfect-lemma-chi-locally-constant}.
\end{proof}

\begin{lemma}
\label{perfect-lemma-open-where-cohomology-in-degree-i-rank-r-geometric}
Soit $f : X \to S$ un morphisme propre de présentation finie.
Soit $\mathcal{F}$ un $\mathcal{O}_X$-module de présentation finie,
plat sur $S$. Fixons $i, r \in \mathbf{Z}$.
Alors il existe un sous-schéma ouvert
$U \subset S$ ayant la propriété suivante :
un morphisme $T \to S$ se factorise par $U$ si et seulement si
$Rf_{T, *}\mathcal{F}_T$ est isomorphe à un
module localement libre de type fini de rang $r$ placé en degré $i$.
\end{lemma}

\begin{proof}
D'après la cohomologie et le changement de base (plus précisément d'après le
lemme \ref{perfect-lemma-flat-proper-perfect-direct-image-general}),
l'objet $K = Rf_*\mathcal{F}$ est un objet parfait de la catégorie dérivée
de $S$, dont la formation commute à tout changement de base.
Ce lemme résulte donc immédiatement du
lemme \ref{perfect-lemma-open-where-cohomology-in-degree-i-rank-r}.
\end{proof}

\begin{lemma}
\label{perfect-lemma-vanishing-implies-locally-free}
Soit $f : X \to S$ un morphisme de présentation finie.
Soit $\mathcal{F}$ un $\mathcal{O}_X$-module de présentation finie,
plat sur $S$, dont le support est propre sur $S$. Si $R^if_*\mathcal{F} = 0$
pour $i > 0$, alors $f_*\mathcal{F}$ est localement libre et sa formation
commute à tout changement de base (voir la démonstration pour une explication).
\end{lemma}

\begin{proof}
D'après le lemme \ref{perfect-lemma-base-change-tensor-perfect},
l'objet $E = Rf_*\mathcal{F}$ de $D(\mathcal{O}_S)$
est parfait et sa formation commute à tout changement de base,
au sens où $Rf'_*(g')^*\mathcal{F} = Lg^*E$
pour tout diagramme cartésien
$$
\xymatrix{
X' \ar[r]_{g'} \ar[d]_{f'} &
X \ar[d]^f \\
S' \ar[r]^g &
S
}
$$
de schémas.
Puisqu'il n'y a jamais de cohomologie en degrés $< 0$, on voit que
$E$ a (localement) une amplitude de Tor contenue dans $[0, b]$ pour un certain $b$.
Si $H^i(E) = R^if_*\mathcal{F} = 0$ pour $i > 0$,
alors $E$ a une amplitude de Tor contenue dans $[0, 0]$. Par conséquent,
$E = H^0(E)[0]$. On en déduit que $H^0(E) = f_*\mathcal{F}$
est localement libre de type fini d'après
Compléments sur l'algèbre, lemme \ref{more-algebra-lemma-perfect}
(et la caractérisation des modules projectifs de type fini donnée dans
Algèbre, lemme \ref{algebra-lemma-finite-projective}).
La commutation au changement de base signifie que
$g^*f_*\mathcal{F} = f'_*(g')^*\mathcal{F}$ pour
un diagramme comme ci-dessus, et elle résulte de la commutation au changement
de base déjà établie pour $E$.
\end{proof}

\begin{lemma}
\label{perfect-lemma-proper-flat-h0}
Soit $f : X \to S$ un morphisme de schémas. Supposons que
\begin{enumerate}
\item $f$ soit propre, plat et de présentation finie ;
\item pour tout $s \in S$, on ait $\kappa(s) = H^0(X_s, \mathcal{O}_{X_s})$.
\end{enumerate}
Alors :
\begin{enumerate}
\item[(a)] $f_*\mathcal{O}_X = \mathcal{O}_S$, et
cela reste vrai après tout changement de base ;
\item[(b)] localement sur $S$, on a
$$
Rf_*\mathcal{O}_X = \mathcal{O}_S \oplus P
$$
dans $D(\mathcal{O}_S)$,
où $P$ est parfait d'amplitude de Tor contenue dans $[1, \infty)$.
\end{enumerate}
\end{lemma}

\begin{proof}
D'après la cohomologie et le changement de base
(lemme \ref{perfect-lemma-flat-proper-perfect-direct-image-general}),
le complexe $E = Rf_*\mathcal{O}_X$
est parfait et sa formation commute à tout changement de base.
Cela implique d'abord que $E$ a une amplitude de Tor contenue dans $[0, \infty)$.
Ensuite, cela implique que, pour $s \in S$, on a
$H^0(E \otimes^\mathbf{L} \kappa(s)) =
H^0(X_s, \mathcal{O}_{X_s}) = \kappa(s)$.
Il en résulte que l'application $\mathcal{O}_S \to Rf_*\mathcal{O}_X = E$
induit un isomorphisme
$\mathcal{O}_S \otimes \kappa(s) \to H^0(E \otimes^\mathbf{L} \kappa(s))$.
Ainsi, $H^0(E) \otimes \kappa(s) \to H^0(E \otimes^\mathbf{L} \kappa(s))$
est surjective, et l'on peut appliquer
Compléments sur l'algèbre, lemme \ref{more-algebra-lemma-better-cut-complex-in-two},
pour voir qu'après avoir remplacé $S$ par un voisinage ouvert affine de $s$,
on a une décomposition $E = H^0(E) \oplus \tau_{\geq 1}E$
où $\tau_{\geq 1}E$ est parfait d'amplitude de Tor contenue dans $[1, \infty)$.
Puisque $E$ a une amplitude de Tor contenue dans $[0, \infty)$, on trouve que
$H^0(E)$ est un $\mathcal{O}_S$-module plat.
Il en résulte que $H^0(E)$ est un $\mathcal{O}_S$-module plat et parfait,
donc localement libre de type fini ; voir
Compléments sur l'algèbre, lemme \ref{more-algebra-lemma-perfect}
(ainsi que le fait que les modules projectifs de type fini sont localement libres de type fini d'après
Algèbre, lemme \ref{algebra-lemma-finite-projective}).
Il en résulte que l'application $\mathcal{O}_S \to H^0(E)$ est
un isomorphisme, puisque cela se vérifie après tensorisation par les
corps résiduels (Algèbre, lemme \ref{algebra-lemma-cokernel-flat}).
\end{proof}

\begin{lemma}
\label{perfect-lemma-proper-flat-geom-red-connected}
Soit $f : X \to S$ un morphisme de schémas. Supposons que
\begin{enumerate}
\item $f$ soit propre, plat et de présentation finie ;
\item les fibres géométriques de $f$ soient réduites et connexes.
\end{enumerate}
Alors $f_*\mathcal{O}_X = \mathcal{O}_S$, et cela reste vrai
après tout changement de base.
\end{lemma}

\begin{proof}
D'après le lemme \ref{perfect-lemma-proper-flat-h0},
il suffit de montrer que $\kappa(s) = H^0(X_s, \mathcal{O}_{X_s})$
pour tout $s \in S$. Cela résulte de
Variétés, lemme
\ref{varieties-lemma-proper-geometrically-reduced-global-sections},
et du fait que $X_s$ est géométriquement connexe et géométriquement réduit.
\end{proof}

\begin{lemma}
\label{perfect-lemma-proper-idempotent-on-fibre}
Soit $f : X \to S$ un morphisme propre de schémas. Soit $s \in S$,
et soit $e \in H^0(X_s, \mathcal{O}_{X_s})$ un idempotent.
Alors $e$ appartient à l'image de l'application
$(f_*\mathcal{O}_X)_s \to H^0(X_s, \mathcal{O}_{X_s})$.
\end{lemma}

\begin{proof}
Soit $X_s = T_1 \amalg T_2$ la décomposition en union disjointe
où $T_1$ et $T_2$ sont non vides, ouverts et fermés dans $X_s$,
correspondant à $e$, c'est-à-dire telle que $e$ soit identiquement égal à $1$
sur $T_1$ et identiquement égal à $0$ sur $T_2$.

\medskip\noindent
Supposons $S$ noethérien. Nous utiliserons le théorème des fonctions formelles
sous la forme de Cohomologie des schémas, lemme
\ref{coherent-lemma-formal-functions-stalk}.
Il nous dit que
$$
(f_*\mathcal{O}_X)_s^\wedge = \lim_n H^0(X_n, \mathcal{O}_{X_n})
$$
où $X_n$ est le $n$-ième voisinage infinitésimal de $X_s$.
Puisque l'espace topologique sous-jacent à $X_n$ est égal à celui
de $X_s$, on obtient pour tout $n $ une décomposition de schémas en union disjointe
$X_n = T_{1, n} \amalg T_{2, n}$ où l'espace topologique sous-jacent
à $T_{i, n}$ est $T_i$ pour $i = 1, 2$. Cela signifie que
$H^0(X_n, \mathcal{O}_{X_n})$ contient un idempotent non trivial $e_n$,
à savoir la fonction identiquement égale à $1$ sur $T_{1, n}$ et
identiquement égale à $0$ sur $T_{2, n}$. Il est clair que $e_{n + 1}$
se restreint en $e_n$ sur $X_n$. Ainsi, $e_\infty = \lim e_n$
est un idempotent non trivial de la limite. Par conséquent, $e_\infty$
est un élément du complété de $(f_*\mathcal{O}_X)_s$
qui s'envoie sur $e$ dans $H^0(X_s, \mathcal{O}_{X_s})$.
Puisque l'application $(f_*\mathcal{O}_X)_s^\wedge \to H^0(X_s, \mathcal{O}_{X_s})$
se factorise par
$(f_*\mathcal{O}_X)^\wedge_s / \mathfrak m_s (f_*\mathcal{O}_X)_s^\wedge =
(f_*\mathcal{O}_X)_s / \mathfrak m_s (f_*\mathcal{O}_X)_s$
(Algèbre, lemme \ref{algebra-lemma-hathat-finitely-generated}),
on en déduit que $e$ appartient à l'image de
l'application $(f_*\mathcal{O}_X)_s \to H^0(X_s, \mathcal{O}_{X_s})$,
comme désiré.

\medskip\noindent
Cas général : on ramène le cas général au cas noethérien par des
arguments de limites. Nous invitons le lecteur à omettre la démonstration.
On peut remplacer $S$ par un voisinage ouvert affine de $s$.
On peut donc supposer, et l'on suppose, que $S$ est affine. D'après Limites, lemme
\ref{limits-lemma-proper-limit-of-proper-finite-presentation-noetherian},
on peut écrire $(f : X \to S) = \lim (f_i : X_i \to S_i)$, avec $f_i$
propre et $S_i$ noethérien. Notons $s_i \in S_i$ l'image de $s$.
Alors $s = \lim s_i$ ; voir
Limites, lemme \ref{limits-lemma-inverse-limit-irreducibles}.
Alors $X_s = X \times_S s = \lim X_i \times_{S_i} s_i = \lim X_{i, s_i}$,
car les limites commutent aux limites
(Catégories, lemme \ref{categories-lemma-colimits-commute}).
Ainsi, $e$ est l'image d'un certain idempotent
$e_i \in H^0(X_{i, s_i}, \mathcal{O}_{X_{i, s_i}})$
d'après Limites, lemme \ref{limits-lemma-descend-section}.
D'après le cas noethérien, il existe un élément $\tilde e_i$
de la fibre $(f_{i, *}\mathcal{O}_{X_i})_{s_i}$ qui s'envoie
sur $e_i$. En prenant l'image inverse de $\tilde e_i$, on obtient un
élément $\tilde e$ de $(f_*\mathcal{O}_X)_s$ qui s'envoie
sur $e$, ce qui achève la démonstration.
\end{proof}

\begin{lemma}
\label{perfect-lemma-proper-flat-geom-red}
Soit $f : X \to S$ un morphisme de schémas. Soit $s \in S$. Supposons que
\begin{enumerate}
\item $f$ soit propre, plat et de présentation finie ;
\item la fibre $X_s$ soit géométriquement réduite.
\end{enumerate}
Alors, après avoir remplacé $S$ par un voisinage ouvert de $s$, il
existe une décomposition en somme directe
$Rf_*\mathcal{O}_X = f_*\mathcal{O}_X \oplus P$
dans $D(\mathcal{O}_S)$, où $f_*\mathcal{O}_X$ est une algèbre finie \'etale
sur $\mathcal{O}_S$, et
$P$ est parfait d'amplitude de Tor contenue dans $[1, \infty)$.
\end{lemma}

\begin{proof}
La démonstration de ce lemme est analogue à celle du
lemme \ref{perfect-lemma-proper-flat-h0},
que nous suggérons au lecteur de lire d'abord.
D'après la cohomologie et le changement de base
(lemme \ref{perfect-lemma-flat-proper-perfect-direct-image-general}),
le complexe $E = Rf_*\mathcal{O}_X$
est parfait et sa formation commute à tout changement de base.
Cela implique d'abord que $E$ a une amplitude de Tor contenue dans $[0, \infty)$.

\medskip\noindent
Nous affirmons qu'après avoir remplacé $S$ par un voisinage ouvert de
$s$, on peut trouver une décomposition en somme directe
$E = H^0(E) \oplus \tau_{\geq 1}E$ dans $D(\mathcal{O}_S)$
où $\tau_{\geq 1}E$ a une amplitude de Tor contenue dans $[1, \infty)$.
Admettons cette assertion pour le moment et supposons le remplacement effectué,
de sorte que l'on dispose de la décomposition en somme directe.
Puisque $E$ a une amplitude de Tor contenue dans $[0, \infty)$, on trouve que
$H^0(E)$ est un $\mathcal{O}_S$-module plat.
Ainsi, $H^0(E)$ est un $\mathcal{O}_S$-module plat et parfait,
donc localement libre de type fini ; voir
Compléments sur l'algèbre, lemme \ref{more-algebra-lemma-perfect}
(ainsi que le fait que les modules projectifs de type fini sont localement libres de type fini d'après
Algèbre, lemme \ref{algebra-lemma-finite-projective}).
Bien entendu, $H^0(E) = f_*\mathcal{O}_X$ est une $\mathcal{O}_S$-algèbre.
D'après la cohomologie et le changement de base, on obtient
$H^0(E) \otimes \kappa(s) = H^0(X_s, \mathcal{O}_{X_s})$.
D'après Variétés, lemme
\ref{varieties-lemma-proper-geometrically-reduced-global-sections},
et l'hypothèse que $X_s$ est géométriquement réduit, on
voit que $\kappa(s) \to H^0(E) \otimes \kappa(s)$
est finie \'etale. D'après Morphismes, lemme
\ref{morphisms-lemma-set-points-where-fibres-etale},
appliqué au morphisme fini localement libre
$\underline{\Spec}_S(H^0(E)) \to S$,
on en déduit qu'après avoir rétréci $S$, la $\mathcal{O}_S$-algèbre
$H^0(E)$ est finie \'etale.

\medskip\noindent
Il reste à démontrer l'assertion. Pour cela, il suffit de montrer que l'application
$$
(f_*\mathcal{O}_X)_s
\longrightarrow
H^0(X_s, \mathcal{O}_{X_s}) = H^0(E \otimes^\mathbf{L} \kappa(s))
$$
est surjective ; voir
Compléments sur l'algèbre, lemme \ref{more-algebra-lemma-better-cut-complex-in-two}.
Choisissons un homomorphisme local plat d'anneaux $\mathcal{O}_{S, s} \to A$
tel que le corps résiduel $k$ de $A$ soit algébriquement clos ; voir
Algèbre, lemme \ref{algebra-lemma-flat-local-given-residue-field}.
D'après le changement de base plat (Cohomologie des schémas, lemme
\ref{coherent-lemma-flat-base-change-cohomology}),
on obtient $H^0(X_A, \mathcal{O}_{X_A}) =
(f_*\mathcal{O}_X)_s \otimes_{\mathcal{O}_{S, s}} A$
et $H^0(X_k, \mathcal{O}_{X_k}) =
H^0(X_s, \mathcal{O}_{X_s}) \otimes_{\kappa(s)} k$.
Il suffit donc de montrer que
$H^0(X_A, \mathcal{O}_{X_A}) \to H^0(X_k, \mathcal{O}_{X_k})$
est surjective. Puisque $X_k$ est un schéma propre réduit sur $k$
et que $k$ est algébriquement clos, on voit que
$H^0(X_k, \mathcal{O}_{X_k})$ est un produit fini de copies
de $k$ d'après Variétés, lemme déjà utilisé,
\ref{varieties-lemma-proper-geometrically-reduced-global-sections}.
Puisque, d'après le lemme \ref{perfect-lemma-proper-idempotent-on-fibre},
les idempotents de cette $k$-algèbre sont dans l'image
de $H^0(X_A, \mathcal{O}_{X_A}) \to H^0(X_k, \mathcal{O}_{X_k})$, on conclut.
\end{proof}





\section{Autres applications}
\label{perfect-section-other-applications}

\noindent
Dans cette section, nous énonçons et démontrons quelques résultats qui se déduisent
de la théorie développée ci-dessus.

\begin{lemma}
\label{perfect-lemma-cohomology-over-coherent-ring}
Soit $R$ un anneau cohérent. Soit $X$ un schéma de présentation finie sur $R$.
Soit $\mathcal{G}$ un $\mathcal{O}_X$-module de présentation finie,
plat sur $R$, dont le support est propre sur $R$. Alors
$H^i(X, \mathcal{G})$ est un $R$-module cohérent.
\end{lemma}

\begin{proof}
On combine le lemme \ref{perfect-lemma-base-change-tensor-perfect} avec les
lemmes de Compléments sur l'algèbre \ref{more-algebra-lemma-coherent-pseudo-coherent} et
\ref{more-algebra-lemma-perfect}.
\end{proof}

\begin{lemma}
\label{perfect-lemma-countable-cohomology}
Soit $X$ un schéma quasi-compact et quasi-séparé.
Soit $K$ un objet de $D_\QCoh(\mathcal{O}_X)$
tel que les faisceaux de cohomologie $H^i(K)$ aient des ensembles
dénombrables de sections sur les ouverts affines. Alors, pour tout ouvert quasi-compact
$U \subset X$ et tout objet parfait $E$ de $D(\mathcal{O}_X)$,
les ensembles
$$
H^i(U, K \otimes^\mathbf{L} E),\quad \Ext^i(E|_U, K|_U)
$$
sont dénombrables.
\end{lemma}

\begin{proof}
En utilisant Cohomologie, lemme \ref{cohomology-lemma-dual-perfect-complex},
on voit qu'il suffit de démontrer le résultat
pour les groupes $H^i(U, K \otimes^\mathbf{L} E)$.
Nous utiliserons le principe d'induction pour démontrer le lemme ; voir
Cohomologie des schémas, lemme \ref{coherent-lemma-induction-principle}.

\medskip\noindent
Montrons d'abord que l'assertion vaut lorsque $U = \Spec(A)$ est affine. En effet, on peut
représenter $K$ par un complexe de $A$-modules $K^\bullet$, et $E$ par un
complexe fini de $A$-modules projectifs de type fini $P^\bullet$.
Voir les lemmes \ref{perfect-lemma-affine-compare-bounded} et
\ref{perfect-lemma-perfect-affine},
ainsi que notre définition des complexes parfaits de $A$-modules
(Compléments sur l'algèbre, définition \ref{more-algebra-definition-perfect}).
Alors $(E \otimes^\mathbf{L} K)|_U$ est représenté par
le complexe total associé au complexe double
$P^\bullet \otimes_A K^\bullet$
(lemme \ref{perfect-lemma-quasi-coherence-tensor-product}).
En raisonnant par récurrence sur la longueur du complexe
$P^\bullet$ (ou en utilisant une suite spectrale appropriée),
on voit qu'il suffit de montrer que
$H^i(P^a \otimes_A K^\bullet)$ est dénombrable pour tout $a$.
Puisque $P^a$ est un facteur direct de $A^{\oplus n}$ pour
un certain $n$, cela résulte de l'hypothèse que
le groupe de cohomologie $H^i(K^\bullet)$ est dénombrable.

\medskip\noindent
Pour achever la démonstration, il suffit de montrer ceci : si $U = V \cup W$
et si le résultat vaut pour $V$, $W$ et $V \cap W$, alors
il vaut pour $U$. C'est une conséquence immédiate
de la suite de Mayer-Vietoris ; voir
Cohomologie, lemme \ref{cohomology-lemma-unbounded-mayer-vietoris}.
\end{proof}

\begin{lemma}
\label{perfect-lemma-countable}
Soit $X$ un schéma quasi-compact et quasi-séparé tel que
les ensembles de sections de $\mathcal{O}_X$ sur les ouverts affines soient dénombrables.
Soit $K$ un objet de $D_\QCoh(\mathcal{O}_X)$. Les
conditions suivantes sont équivalentes :
\begin{enumerate}
\item $K = \text{hocolim} E_n$, où $E_n$ est un objet parfait de
$D(\mathcal{O}_X)$ ;
\item les faisceaux de cohomologie $H^i(K)$ ont des ensembles
dénombrables de sections sur les ouverts affines.
\end{enumerate}
\end{lemma}

\begin{proof}
Si (1) est vraie, alors (2) l'est, car les colimites homotopiques commutent
à la prise des faisceaux de cohomologie
(d'après Catégories dérivées, lemme \ref{derived-lemma-cohomology-of-hocolim})
et parce qu'un complexe parfait est
localement isomorphe à un complexe fini de $\mathcal{O}_X$-modules libres de type fini
et satisfait donc (2) d'après l'hypothèse sur $X$.

\medskip\noindent
Supposons (2).
Choisissons un complexe K-injectif $\mathcal{K}^\bullet$ représentant $K$.
Choisissons un générateur parfait $E$ de $D_\QCoh(\mathcal{O}_X)$ et
représentons-le par un complexe K-injectif $\mathcal{I}^\bullet$.
D'après le théorème \ref{perfect-theorem-DQCoh-is-Ddga}
et sa démonstration, il existe une équivalence
de catégories triangulées $F : D_\QCoh(\mathcal{O}_X) \to D(A, \text{d})$,
où $(A, \text{d})$ est l'algèbre différentielle graduée
$$
(A, \text{d}) =
\Hom_{\text{Comp}^{dg}(\mathcal{O}_X)}
(\mathcal{I}^\bullet, \mathcal{I}^\bullet)
$$
qui envoie $K$ sur le module différentiel gradué
$$
M = \Hom_{\text{Comp}^{dg}(\mathcal{O}_X)}
(\mathcal{I}^\bullet, \mathcal{K}^\bullet)
$$
Notons que $H^i(A) = \Ext^i(E, E)$ et
$H^i(M) = \Ext^i(E, K)$.
De plus, puisque $F$ est une équivalence, elle et son quasi-inverse commutent
aux colimites homotopiques.
Il suffit donc d'écrire $M$ comme une colimite homotopique
d'objets compacts de $D(A, \text{d})$.
D'après Algèbre différentielle graduée, lemme \ref{dga-lemma-countable},
il suffit de montrer que $\Ext^i(E, E)$ et
$\Ext^i(E, K)$ sont dénombrables pour tout $i$.
Cela résulte du lemme \ref{perfect-lemma-countable-cohomology}.
\end{proof}

\begin{lemma}
\label{perfect-lemma-computing-sections-as-colim}
Soit $A$ un anneau. Soit $X$ un schéma de présentation finie sur $A$.
Soit $f : U \to X$ un morphisme plat de présentation finie. Alors
\begin{enumerate}
\item il existe un système projectif d'objets parfaits $L_n$ de
$D(\mathcal{O}_X)$ tel que
$$
R\Gamma(U, Lf^*K) = \text{hocolim}\ R\Hom_X(L_n, K)
$$
dans $D(A)$, fonctoriellement en $K$ dans $D_\QCoh(\mathcal{O}_X)$ ;
\item il existe un système d'objets parfaits $E_n$ de
$D(\mathcal{O}_X)$ tel que
$$
R\Gamma(U, Lf^*K) = \text{hocolim}\ R\Gamma(X, E_n \otimes^\mathbf{L} K)
$$
dans $D(A)$, fonctoriellement en $K$ dans $D_\QCoh(\mathcal{O}_X)$.
\end{enumerate}
\end{lemma}

\begin{proof}
D'après le lemme \ref{perfect-lemma-cohomology-base-change}, on a
$$
R\Gamma(U, Lf^*K) = R\Gamma(X, Rf_*\mathcal{O}_U \otimes^\mathbf{L} K)
$$
fonctoriellement en $K$. Observons que $R\Gamma(X, -)$ commute aux
colimites homotopiques, car il commute aux sommes directes d'après le
lemme \ref{perfect-lemma-quasi-coherence-pushforward-direct-sums}.
De même, $- \otimes^\mathbf{L} K$ commute aux colimites dérivées,
car $- \otimes^\mathbf{L} K$ commute aux sommes directes
(puisque les sommes directes dans $D(\mathcal{O}_X)$
sont données par les sommes directes des complexes qui les représentent).
Ainsi, pour démontrer (2), il suffit d'écrire
$Rf_*\mathcal{O}_U = \text{hocolim} E_n$ pour un système
d'objets parfaits $E_n$ de $D(\mathcal{O}_X)$. Cela fait,
on obtient (1) en posant $L_n = E_n^\vee$ ; voir
Cohomologie, lemme \ref{cohomology-lemma-dual-perfect-complex}.

\medskip\noindent
Écrivons $A = \colim A_i$, où $A_i$ est de type fini sur $\mathbf{Z}$.
D'après Limites, lemme \ref{limits-lemma-descend-finite-presentation},
on peut trouver un $i$ et des morphismes $U_i \to X_i \to \Spec(A_i)$
de présentation finie dont le changement de base à $\Spec(A)$ redonne
$U \to X \to \Spec(A)$.
Après avoir augmenté $i$, on peut supposer que $f_i : U_i \to X_i$ est
plat ; voir Limites, lemme \ref{limits-lemma-descend-flat-finite-presentation}.
D'après le lemme \ref{perfect-lemma-compare-base-change},
l'image inverse dérivée de $Rf_{i, *}\mathcal{O}_{U_i}$
par $g : X \to X_i$ est égale à $Rf_*\mathcal{O}_U$.
Puisque $Lg^*$ commute aux colimites dérivées, il suffit
de démontrer ce que l'on veut pour $f_i$. On peut donc supposer que
$U$ et $X$ sont de type fini sur $\mathbf{Z}$.

\medskip\noindent
Supposons que $f : U \to X$ soit un morphisme entre des schémas de type fini sur $\mathbf{Z}$.
Pour achever la démonstration, nous allons montrer que $Rf_*\mathcal{O}_U$ est une
colimite homotopique de complexes parfaits. Pour le voir, appliquons le lemme \ref{perfect-lemma-countable}.
Il suffit donc de montrer que $R^if_*\mathcal{O}_U$
a des ensembles dénombrables de sections sur les ouverts affines.
Cela résulte du lemme \ref{perfect-lemma-countable-cohomology}
appliqué au faisceau structural.
\end{proof}




\section{Caractérisation des complexes pseudo-cohérents, II}
\label{perfect-section-pseudo-coherent}

\noindent
Cette section prolonge la
section \ref{perfect-section-pseudo-coherent-hocolim}.
Nous y examinons des caractérisations des complexes pseudo-cohérents
en termes de cohomologie. On trouvera d'autres résultats de cette nature dans
Compléments sur les morphismes, section
\ref{more-morphisms-section-characterize-pseudo-coherent}.

\begin{lemma}
\label{perfect-lemma-pseudo-coherent-over-algebra}
Soit $A$ un anneau. Soit $R$ une $A$-algèbre (éventuellement non commutative)
qui est libre de type fini comme $A$-module. Alors tout objet $M$ de $D(R)$
qui est pseudo-cohérent dans $D(A)$ peut être représenté par un
complexe borné supérieurement de $R$-modules (à droite) libres de type fini.
\end{lemma}

\begin{proof}
Choisissons un complexe $M^\bullet$ de $R$-modules à droite représentant $M$.
Puisque $M$ est pseudo-cohérent, on a $H^i(M) = 0$ pour $i$ assez grand.
Soit $m$ le plus petit indice tel que $H^m(M)$ soit non nul.
Alors $H^m(M)$ est un $A$-module de type fini d'après
Compléments sur l'algèbre, lemme \ref{more-algebra-lemma-finite-cohomology}.
On peut donc choisir un $R$-module libre de type fini $F^m$ et une application
$F^m \to M^m$ tels que $F^m \to M^m \to M^{m + 1}$ soit nulle
et que $F^m \to H^m(M)$ soit surjective.
Diagramme :
$$
\xymatrix{
& F^m \ar[d]^\alpha \ar[r] & 0 \ar[d] \ar[r] & \ldots \\
M^{m - 1} \ar[r] & M^m \ar[r] & M^{m + 1} \ar[r] & \ldots
}
$$
Par récurrence descendante sur $n \leq m$, nous allons construire
des $R$-modules libres de type fini $F^i$ pour $i \geq n$, des différentielles
$d^i : F^i \to F^{i + 1}$ pour $i \geq n$, et des applications $\alpha : F^i \to K^i$
compatibles aux différentielles, telles que
(1) $H^i(\alpha)$ soit un isomorphisme pour $i > n$ et une surjection pour $i = n$, et
(2) $F^i = 0$ pour $i > m$. Diagramme :
$$
\xymatrix{
& F^n \ar[r] \ar[d]^\alpha & F^{n + 1} \ar[d]^\alpha \ar[r] & \ldots \ar[r]
& F^i \ar[d]^\alpha \ar[r] & 0 \ar[d] \ar[r] & \ldots \\
M^{n - 1} \ar[r] & M^n \ar[r] & M^{n + 1} \ar[r] & \ldots \ar[r] &
M^i \ar[r] & M^{i + 1} \ar[r] & \ldots
}
$$
Le cas initial $n = m$ a été traité ci-dessus.
Étape de récurrence. Soit $C^\bullet$ le cône de $\alpha$
(Catégories dérivées, définition \ref{derived-definition-cone}).
La suite exacte longue
de cohomologie montre que $H^i(C^\bullet) = 0$ pour $i \geq n$.
Observons que $F^\bullet$ est pseudo-cohérent comme complexe de $A$-modules,
car $R$ est libre de type fini comme $A$-module. Par conséquent,
d'après Compléments sur l'algèbre, lemme \ref{more-algebra-lemma-cone-pseudo-coherent},
on voit que $C^\bullet$ est $(n - 1)$-pseudo-cohérent comme complexe de
$A$-modules. D'après
Compléments sur l'algèbre, lemme \ref{more-algebra-lemma-finite-cohomology},
on voit que $H^{n - 1}(C^\bullet)$ est un $A$-module de type fini.
Choisissons un $R$-module libre de type fini $F^{n - 1}$ et une application
$\beta : F^{n - 1} \to C^{n - 1}$ tels que la composée
$F^{n - 1} \to C^{n - 1} \to C^n$ soit nulle et que $F^{n - 1}$
se surjecte sur $H^{n - 1}(C^\bullet)$. Puisque
$C^{n - 1} = M^{n - 1} \oplus F^n$,
on peut écrire $\beta = (\alpha^{n - 1}, -d^{n - 1})$. L'annulation de la
composée $F^{n - 1} \to C^{n - 1} \to C^n$ implique que
ces applications s'insèrent dans un morphisme de complexes
$$
\xymatrix{
& F^{n - 1} \ar[d]^{\alpha^{n - 1}} \ar[r]_{d^{n - 1}} &
F^n \ar[r] \ar[d]^\alpha &
F^{n + 1} \ar[d]^\alpha \ar[r] & \ldots \\
\ldots \ar[r] &
M^{n - 1} \ar[r] & M^n \ar[r] & M^{n + 1} \ar[r] & \ldots
}
$$
De plus, ces applications définissent un morphisme de triangles distingués
$$
\xymatrix{
(F^n \to \ldots) \ar[r] \ar[d] &
(F^{n - 1} \to \ldots) \ar[r] \ar[d] &
F^{n - 1} \ar[r] \ar[d]_\beta &
(F^n \to \ldots)[1] \ar[d] \\
(F^n \to \ldots) \ar[r] &
M^\bullet \ar[r] &
C^\bullet \ar[r] &
(F^n \to \ldots)[1]
}
$$
Notre choix de $\beta$ implique donc que le morphisme de complexes
$(F^{n - 1} \to \ldots) \to M^\bullet$ induit un isomorphisme en
cohomologie dans les degrés $\geq n$ et une surjection en degré $n - 1$.
Cela achève la démonstration du lemme.
\end{proof}

\begin{lemma}
\label{perfect-lemma-pseudo-coherent-on-projective-space}
Soit $A$ un anneau. Soit $n \geq 0$. Soit
$K \in D_\QCoh(\mathcal{O}_{\mathbf{P}^n_A})$.
Les conditions suivantes sont équivalentes :
\begin{enumerate}
\item $K$ est pseudo-cohérent ;
\item $R\Gamma(\mathbf{P}^n_A, E \otimes^\mathbf{L} K)$ est un objet pseudo-cohérent
de $D(A)$ pour tout objet pseudo-cohérent $E$ de
$D(\mathcal{O}_{\mathbf{P}^n_A})$ ;
\item $R\Gamma(\mathbf{P}^n_A, E \otimes^\mathbf{L} K)$ est un objet pseudo-cohérent
de $D(A)$ pour tout objet parfait $E$ de
$D(\mathcal{O}_{\mathbf{P}^n_A})$ ;
\item $R\Hom_{\mathbf{P}^n_A}(E, K)$ est un objet pseudo-cohérent
de $D(A)$ pour tout objet parfait $E$ de
$D(\mathcal{O}_{\mathbf{P}^n_A})$ ;
\item $R\Gamma(\mathbf{P}^n_A,
K \otimes^\mathbf{L} \mathcal{O}_{\mathbf{P}^n_A}(d))$ est un objet pseudo-cohérent
de $D(A)$ pour $d = 0, 1, \ldots, n$.
\end{enumerate}
\end{lemma}

\begin{proof}
Rappelons que
$$
R\Hom_{\mathbf{P}^n_A}(E, K) =
R\Gamma(\mathbf{P}^n_A, R\SheafHom_{\mathcal{O}_{\mathbf{P}^n_A}}(E, K))
$$
par définition ; voir Cohomologie, section \ref{cohomology-section-global-RHom}.
Ainsi, les assertions (4) et (3) sont équivalentes d'après
Cohomologie, lemme \ref{cohomology-lemma-dual-perfect-complex}.

\medskip\noindent
Puisque tout complexe parfait est pseudo-cohérent, il est clair que
(2) implique (3).

\medskip\noindent
Supposons (1). Alors $E \otimes^\mathbf{L} K$ est pseudo-cohérent
pour tout $E$ pseudo-cohérent ; voir
Cohomologie, lemme \ref{cohomology-lemma-tensor-pseudo-coherent}.
D'après le lemme \ref{perfect-lemma-flat-proper-pseudo-coherent-direct-image-general},
l'image directe d'un tel complexe pseudo-cohérent est pseudo-cohérente,
et l'on voit que (2) est vraie.

\medskip\noindent
L'assertion (3) implique (5), car on peut prendre $E = \mathcal{O}_{\mathbf{P}^n_A}(d)$
pour $d = 0, 1, \ldots, n$.

\medskip\noindent
Pour achever la démonstration, il faut montrer que (5) implique (1).
Soit $P$ comme dans (\ref{perfect-equation-generator-Pn}), et
$R$ comme dans (\ref{perfect-equation-algebra-for-Pn}).
D'après le lemme \ref{perfect-lemma-Pn-module-category}, on a une équivalence
$$
- \otimes^\mathbf{L}_R P :
D(R) \longrightarrow D_\QCoh(\mathcal{O}_{\mathbf{P}^n_A})
$$
Soit $M \in D(R)$ un objet tel que $M \otimes^\mathbf{L} P = K$.
D'après Algèbre différentielle graduée, lemme
\ref{dga-lemma-upgrade-tensor-with-complex-derived},
il existe un isomorphisme
$$
R\Hom(R, M) = R\Hom_{\mathbf{P}^n_A}(P, K)
$$
dans $D(A)$. En raisonnant comme ci-dessus, on obtient
$$
R\Hom_{\mathbf{P}^n_A}(P, K) = R\Gamma(\mathbf{P}^n_A,
R\SheafHom_{\mathcal{O}_{\mathbf{P}^n_A}}(E, K)) =
R\Gamma(\mathbf{P}^n_A,
P^\vee \otimes^\mathbf{L}_{\mathcal{O}_{\mathbf{P}^n_A}} K).
$$
Puisque $P^\vee$ est la somme directe des
$\mathcal{O}_{\mathbf{P}^n_A}(d)$ pour $d = 0, 1, \ldots, n$,
il résulte de (5) que $R\Hom(R, M)$ est pseudo-cohérent comme complexe
de $A$-modules. Bien entendu, $M = R\Hom(R, M)$ dans $D(A)$.
Ainsi, $M$ est pseudo-cohérent comme complexe de $A$-modules.
D'après le lemme \ref{perfect-lemma-pseudo-coherent-over-algebra},
on peut représenter $M$ par un complexe borné supérieurement $F^\bullet$
de $R$-modules libres de type fini. Alors
$F^\bullet = \bigcup_{p \geq 0} \sigma_{\geq p}F^\bullet$
est une filtration qui montre que $F^\bullet$ est un module différentiel
gradué sur $R$ ayant la propriété (P) ; voir
Algèbre différentielle graduée, section \ref{dga-section-P-resolutions}.
Par conséquent, $K = M \otimes^\mathbf{L}_R P$ est représenté
par $F^\bullet \otimes_R P$ (cela résulte de la construction du
foncteur de produit tensoriel dérivé ; voir, par exemple, la démonstration de
Algèbre différentielle graduée, lemme \ref{dga-lemma-tensor-with-complex-derived}).
Puisque $F^\bullet \otimes_R P$
est un complexe borné supérieurement dont les termes sont des sommes directes de
copies de $P$, on en déduit que le lemme est vrai.
\end{proof}

\begin{lemma}
\label{perfect-lemma-perfect-enough}
Soit $A$ un anneau. Soit $X$ un schéma quasi-compact
et quasi-séparé sur $A$. Soit $K \in D^-_\QCoh(\mathcal{O}_X)$.
Si $R\Gamma(X, E \otimes^\mathbf{L} K)$ est pseudo-cohérent
dans $D(A)$ pour tout $E$ parfait dans $D(\mathcal{O}_X)$,
alors $R\Gamma(X, E \otimes^\mathbf{L} K)$ est pseudo-cohérent
dans $D(A)$ pour tout $E$ pseudo-cohérent dans $D(\mathcal{O}_X)$.
\end{lemma}

\begin{proof}
Il existe un entier $N$ tel que
$R\Gamma(X, -) : D_\QCoh(\mathcal{O}_X) \to D(A)$
soit de dimension cohomologique $N$, comme expliqué dans le
lemme \ref{perfect-lemma-quasi-coherence-direct-image}.
Soit $b \in \mathbf{Z}$ tel que $H^i(K) = 0$ pour $i > b$.
Soit $E$ pseudo-cohérent sur $X$.
Il suffit de montrer que $R\Gamma(X, E \otimes^\mathbf{L} K)$
est $m$-pseudo-cohérent pour tout $m$.
Choisissons une approximation $P \to E$ par un complexe parfait $P$
de $(X, E, m - N - 1 - b)$. Cela est possible d'après le
théorème \ref{perfect-theorem-approximation}.
Choisissons un triangle distingué
$$
P \to E \to C \to P[1]
$$
dans $D_\QCoh(\mathcal{O}_X)$. Les faisceaux de cohomologie de $C$ sont nuls
dans les degrés $\geq m - N - 1 - b$. Par conséquent, les faisceaux de cohomologie
de $C \otimes^\mathbf{L} K$ sont nuls dans les degrés $\geq m - N - 1$.
Ainsi, les groupes de cohomologie de $R\Gamma(X, C \otimes^\mathbf{L} K)$
sont nuls dans les degrés $\geq m - 1$. Par conséquent,
$$
R\Gamma(X, P \otimes^\mathbf{L} K) \to R\Gamma(X, E \otimes^\mathbf{L} K)
$$
est un isomorphisme sur la cohomologie dans les degrés $\geq m$.
Par hypothèse, la source est pseudo-cohérente.
On en déduit que $R\Gamma(X, E \otimes^\mathbf{L} K)$
est $m$-pseudo-cohérent, comme désiré.
\end{proof}









\section{Objets relativement parfaits}
\label{perfect-section-relatively-perfect}

\noindent
Dans cette section, nous introduisons une notion tirée de
\cite{lieblich-complexes}.

\begin{definition}
\label{perfect-definition-relatively-perfect}
Soit $f : X \to S$ un morphisme de schémas plat et
localement de présentation finie. Un objet $E$ de $D(\mathcal{O}_X)$ est
{\it parfait relativement à $S$} ou
{\it $S$-parfait} si $E$ est pseudo-cohérent
(Cohomologie, définition \ref{cohomology-definition-pseudo-coherent}) et si
$E$ a localement une dimension de Tor finie comme objet de
$D(f^{-1}\mathcal{O}_S)$
(Cohomologie, définition \ref{cohomology-definition-tor-amplitude}).
\end{definition}

\noindent
Voir la remarque \ref{perfect-remark-discuss-rel-perfect} pour une discussion.

\begin{example}
\label{perfect-example-relatively-perfect-field}
Soit $k$ un corps. Soit $X$ un schéma de présentation finie sur $k$
(en particulier, $X$ est quasi-compact). Alors un objet $E$ de $D(\mathcal{O}_X)$
est $k$-parfait si et seulement s'il est borné et pseudo-cohérent
(par définition), c'est-à-dire si et seulement s'il appartient à $D^b_{\textit{Coh}}(X)$
(d'après le lemme \ref{perfect-lemma-identify-pseudo-coherent-noetherian}).
Ainsi, être relativement parfait ne signifie {\bf pas} « parfait sur les fibres ».
\end{example}

\noindent
La notion algébrique correspondante est étudiée dans
Compléments sur l'algèbre, section \ref{more-algebra-section-relatively-perfect}.
On peut relier la notion pour les schémas à la notion algébrique
comme suit.

\begin{lemma}
\label{perfect-lemma-affine-locally-rel-perfect}
Soit $f : X \to S$ un morphisme de schémas plat et
localement de présentation finie. Soit $E$ un objet de
$D_\QCoh(\mathcal{O}_X)$. Les conditions suivantes sont équivalentes :
\begin{enumerate}
\item $E$ est $S$-parfait ;
\item pour tout ouvert affine $U \subset X$ qui s'envoie dans un ouvert affine
$V \subset S$, le complexe $R\Gamma(U, E)$ est $\mathcal{O}_S(V)$-parfait ;
\item il existe un recouvrement ouvert affine $S = \bigcup V_i$
et, pour tout $i$, un recouvrement ouvert affine $f^{-1}(V_i) = \bigcup U_{ij}$
tels que le complexe $R\Gamma(U_{ij}, E)$ soit $\mathcal{O}_S(V_i)$-parfait.
\end{enumerate}
\end{lemma}

\begin{proof}
La pseudo-cohérence est une propriété locale, et
« avoir localement une dimension de Tor finie » est une propriété locale.
Le lemme se ramène donc immédiatement à l'énoncé suivant :
si $X$ et $S$ sont affines, alors $E$ est $S$-parfait
si et seulement si $K = R\Gamma(X, E)$ est $\mathcal{O}_S(S)$-parfait.
Écrivons $X = \Spec(A)$, $S = \Spec(R)$, et supposons que $E$ corresponde à
$K \in D(A)$, c'est-à-dire que $K = R\Gamma(X, E)$ ; voir le
lemme \ref{perfect-lemma-affine-compare-bounded}.

\medskip\noindent
Observons que $K$ est $R$-parfait si et seulement si $K$ est
pseudo-cohérent et a une dimension de Tor finie comme complexe
de $R$-modules (Compléments sur l'algèbre, définition
\ref{more-algebra-definition-relatively-perfect}).
D'après le lemme \ref{perfect-lemma-pseudo-coherent-affine},
on voit que $E$ est pseudo-cohérent si et seulement si
$K$ est pseudo-cohérent.
D'après le lemme \ref{perfect-lemma-tor-dimension-rel-affine}, on voit que
$E$ a une dimension de Tor finie sur $f^{-1}\mathcal{O}_S$
si et seulement si $K$ a une dimension de Tor finie comme complexe de
$R$-modules.
\end{proof}

\begin{lemma}
\label{perfect-lemma-triangulated}
Soit $f : X \to S$ un morphisme de schémas
plat et localement de présentation finie. La sous-catégorie pleine
de $D(\mathcal{O}_X)$ formée des objets $S$-parfaits est
une sous-catégorie triangulée saturée\footnote{Catégories dérivées, définition
\ref{derived-definition-saturated}.}.
\end{lemma}

\begin{proof}
Cela résulte de Cohomologie, lemmes
\ref{cohomology-lemma-cone-pseudo-coherent},
\ref{cohomology-lemma-summands-pseudo-coherent},
\ref{cohomology-lemma-cone-tor-amplitude} et
\ref{cohomology-lemma-summands-tor-amplitude}.
\end{proof}

\begin{lemma}
\label{perfect-lemma-perfect-relatively-perfect}
Soit $f : X \to S$ un morphisme de schémas plat et localement
de présentation finie. Tout objet parfait de $D(\mathcal{O}_X)$ est $S$-parfait.
Si $K, M \in D(\mathcal{O}_X)$, alors $K \otimes_{\mathcal{O}_X}^\mathbf{L} M$
est $S$-parfait si $K$ est parfait et $M$ est $S$-parfait.
\end{lemma}

\begin{proof}
Première démonstration : on se ramène au cas affine à l'aide du
lemme \ref{perfect-lemma-affine-locally-rel-perfect},
puis on applique Compléments sur l'algèbre, lemme
\ref{more-algebra-lemma-perfect-relatively-perfect}.
\end{proof}

\begin{lemma}
\label{perfect-lemma-base-change-relatively-perfect}
Soit $f : X \to S$ un morphisme de schémas plat et
localement de présentation finie.
Soit $g : S' \to S$ un morphisme de schémas. Posons $X' = S' \times_S X$
et notons $g' : X' \to X$ la projection.
Si $K \in D(\mathcal{O}_X)$ est $S$-parfait, alors $L(g')^*K$
est $S'$-parfait.
\end{lemma}

\begin{proof}
Première démonstration : on se ramène au cas affine à l'aide du
lemme \ref{perfect-lemma-affine-locally-rel-perfect},
puis on applique Compléments sur l'algèbre, lemme
\ref{more-algebra-lemma-base-change-relatively-perfect}.

\medskip\noindent
Deuxième démonstration : $L(g')^*K$ est pseudo-cohérent d'après
Cohomologie, lemme \ref{cohomology-lemma-pseudo-coherent-pullback},
et la propriété de dimension de Tor bornée résulte du
lemme \ref{perfect-lemma-tor-independence-and-tor-amplitude}.
\end{proof}

\begin{situation}
\label{perfect-situation-relative-descent}
Soit $S = \lim_{i \in I} S_i$ la limite d'un système filtrant de schémas
dont les morphismes de transition affines sont $g_{i'i} : S_{i'} \to S_i$.
On suppose que $S_i$ est quasi-compact et quasi-séparé pour tout $i \in I$.
On note $g_i : S \to S_i$ la projection. On fixe un élément $0 \in I$
et un morphisme plat de présentation finie $X_0 \to S_0$.
On pose $X_i = S_i \times_{S_0} X_0$ et $X = S \times_{S_0} X_0$,
et l'on note les morphismes de transition $f_{i'i} : X_{i'} \to X_i$
et les projections $f_i : X \to X_i$.
\end{situation}

\begin{lemma}
\label{perfect-lemma-relative-descend-homomorphisms}
Dans la situation \ref{perfect-situation-relative-descent},
soient $K_0$ et $L_0$ des objets de $D(\mathcal{O}_{X_0})$.
Posons $K_i = Lf_{i0}^*K_0$ et $L_i = Lf_{i0}^*L_0$ pour $i \geq 0$,
et posons $K = Lf_0^*K_0$ et $L = Lf_0^*L_0$. Alors l'application
$$
\colim_{i \geq 0} \Hom_{D(\mathcal{O}_{X_i})}(K_i, L_i)
\longrightarrow
\Hom_{D(\mathcal{O}_X)}(K, L)
$$
est un isomorphisme si $K_0$ est pseudo-cohérent et si
$L_0 \in D_\QCoh(\mathcal{O}_{X_0})$ a (localement)
une dimension de Tor finie comme objet de
$D((X_0 \to S_0)^{-1}\mathcal{O}_{S_0})$.
\end{lemma}

\begin{proof}
Pour tout ouvert quasi-compact $U_0 \subset X_0$, considérons la
condition $P$ selon laquelle
$$
\colim_{i \geq 0} \Hom_{D(\mathcal{O}_{U_i})}(K_i|_{U_i}, L_i|_{U_i})
\longrightarrow
\Hom_{D(\mathcal{O}_U)}(K|_U, L|_U)
$$
est un isomorphisme, où $U = f_0^{-1}(U_0)$ et $U_i = f_{i0}^{-1}(U_0)$.
Si $P$ vaut pour $U_0$, $V_0$ et $U_0 \cap V_0$, alors elle vaut
pour $U_0 \cup V_0$ par Mayer-Vietoris
pour les Hom dans la catégorie dérivée ; voir
Cohomologie, lemme \ref{cohomology-lemma-mayer-vietoris-hom}.

\medskip\noindent
Notons $\pi_0 : X_0 \to S_0$ le morphisme donné.
On peut d'abord considérer $U_0 = \pi_0^{-1}(W_0)$ avec
$W_0 \subset S_0$ ouvert quasi-compact. En appliquant le principe d'induction de
Cohomologie des schémas, lemme \ref{coherent-lemma-induction-principle},
aux ouverts quasi-compacts de $S_0$
et en utilisant la remarque ci-dessus, on constate qu'il suffit de démontrer
$P$ pour $U_0 = \pi_0^{-1}(W_0)$ avec $W_0$ affine.
Autrement dit, on s'est ramené au cas où $S_0$ est affine.
Ensuite, on applique de nouveau le principe d'induction, cette fois à tous les
ouverts quasi-compacts et quasi-séparés de $X_0$, pour se ramener également au
cas où $X_0$ est affine.

\medskip\noindent
Si $X_0$ et $S_0$ sont affines, le résultat découle de
Compléments sur l'algèbre, lemme \ref{more-algebra-lemma-colimit-relatively-perfect}.
En effet, d'après les lemmes \ref{perfect-lemma-pseudo-coherent} et
\ref{perfect-lemma-affine-compare-bounded},
l'énoncé se traduit en calculs de morphismes dans les
catégories dérivées de modules. Puis le
lemme \ref{perfect-lemma-pseudo-coherent-affine}
montre que le complexe de modules correspondant à $K_0$
est pseudo-cohérent. Et le
lemme \ref{perfect-lemma-tor-dimension-rel-affine}
montre que le complexe de modules correspondant à $L_0$
a une dimension de Tor finie sur $\mathcal{O}_{S_0}(S_0)$.
Les hypothèses de
Compléments sur l'algèbre, lemme \ref{more-algebra-lemma-colimit-relatively-perfect},
sont donc satisfaites, et l'on conclut.
\end{proof}

\begin{lemma}
\label{perfect-lemma-descend-relatively-perfect}
Dans la situation \ref{perfect-situation-relative-descent}, la catégorie des
objets $S$-parfaits de $D(\mathcal{O}_X)$ est la colimite des catégories
d'objets $S_i$-parfaits de $D(\mathcal{O}_{X_i})$.
\end{lemma}

\begin{proof}
Pour tout ouvert quasi-compact $U_0 \subset X_0$, considérons la condition $P$
selon laquelle le foncteur
$$
\colim_{i \geq 0} D_{S_i\text{-parfait}}(\mathcal{O}_{U_i})
\longrightarrow
D_{S\text{-parfait}}(\mathcal{O}_U)
$$
est une équivalence, où $U = f_0^{-1}(U_0)$ et $U_i = f_{i0}^{-1}(U_0)$.
Observons que nous savons déjà que ce foncteur est pleinement fidèle
d'après le lemme \ref{perfect-lemma-relative-descend-homomorphisms}. Il suffit donc de démontrer
la surjectivité essentielle.

\medskip\noindent
Supposons que $P$ vaille pour des ouverts quasi-compacts $U_0$, $V_0$ de $X_0$.
Nous affirmons que $P$ vaut pour $U_0 \cup V_0$. Nous utiliserons les notations
$U_i = f_{i0}^{-1}U_0$, $U = f_0^{-1}U_0$, $V_i = f_{i0}^{-1}V_0$,
et $V = f_0^{-1}V_0$, et nous emploierons abusivement le symbole
$f_i$ pour tous les morphismes $U \to U_i$, $V \to V_i$,
$U \cap V \to U_i \cap V_i$ et $U \cup V \to U_i \cup V_i$.
Soit $E$ un objet $S$-parfait de $D(\mathcal{O}_{U \cup V})$.
But : montrer que $E$ appartient à l'image essentielle du foncteur.
Par hypothèse,
on peut trouver $i \geq 0$, un objet $S_i$-parfait $E_{U, i}$ sur $U_i$,
un objet $S_i$-parfait $E_{V, i}$ sur $V_i$, et
des isomorphismes $Lf_i^*E_{U, i} \to E|_U$ et $Lf_i^*E_{V, i} \to E|_V$.
Soient
$$
a : E_{U, i} \to (Rf_{i, *}E)|_{U_i}
\quad\text{et}\quad
b : E_{V, i} \to (Rf_{i, *}E)|_{V_i}
$$
les morphismes adjoints aux isomorphismes $Lf_i^*E_{U, i} \to E|_U$
et $Lf_i^*E_{V, i} \to E|_V$.
Par pleine fidélité, après avoir augmenté $i$,
on peut trouver un isomorphisme
$c : E_{U, i}|_{U_i \cap V_i} \to E_{V, i}|_{U_i \cap V_i}$
dont l'image inverse donne les identifications
$$
Lf_i^*E_{U, i}|_{U \cap V} \to E|_{U \cap V} \to Lf_i^*E_{V, i}|_{U \cap V}.
$$
Appliquons Cohomologie, lemme \ref{cohomology-lemma-glue},
pour obtenir un objet $E_i$ sur $U_i \cup V_i$ et un morphisme $d : E_i \to Rf_{i, *}E$
qui se restreint aux morphismes $a$ et $b$ sur $U_i$ et $V_i$.
Il est alors clair que $E_i$ est $S_i$-parfait (car être
relativement parfait est une propriété locale) et que
$d$ est adjoint à un isomorphisme $Lf_i^*E_i \to E$.

\medskip\noindent
Par exactement le même argument que dans
la démonstration du lemme \ref{perfect-lemma-relative-descend-homomorphisms},
en utilisant le principe d'induction
(Cohomologie des schémas, lemme \ref{coherent-lemma-induction-principle}),
on se ramène au cas où $X_0$ et $S_0$ sont tous deux
affines. (On travaille d'abord avec les ouverts de $S_0$ pour se ramener à
$S_0$ affine, puis avec les ouverts de $X_0$ pour se ramener à
$X_0$ affine.) Dans le cas affine, le résultat découle de
Compléments sur l'algèbre, lemme \ref{more-algebra-lemma-colimit-relatively-perfect}.
La traduction en algèbre est donnée par le
lemme \ref{perfect-lemma-affine-locally-rel-perfect}.
\end{proof}

\begin{lemma}
\label{perfect-lemma-derived-pushforward-rel-perfect}
Soit $f : X \to S$ un morphisme de schémas plat, propre et
de présentation finie. Soit $E \in D(\mathcal{O}_X)$ un objet $S$-parfait.
Alors $Rf_*E$ est un objet parfait de $D(\mathcal{O}_S)$,
et sa formation commute à tout changement de base.
\end{lemma}

\begin{proof}
L'énoncé sur le changement de base est le lemme \ref{perfect-lemma-compare-base-change}.
Il suffit donc de montrer que $Rf_*E$ est un objet parfait. Nous allons nous ramener
au cas où $S$ est noethérien affine par un argument de limites.

\medskip\noindent
La question est locale sur $S$ ; on peut donc supposer $S$ affine.
Écrivons $S = \Spec(R)$. Écrivons $R = \colim R_i$ comme colimite filtrante
d'anneaux noethériens $R_i$. D'après Limites, lemme
\ref{limits-lemma-descend-finite-presentation},
il existe un $i$ et un schéma $X_i$ de présentation finie sur $R_i$
dont le changement de base à $R$ est $X$. D'après
Limites, lemmes \ref{limits-lemma-eventually-proper} et
\ref{limits-lemma-descend-flat-finite-presentation},
on peut supposer $X_i$ propre et plat sur $R_i$.
D'après le lemme \ref{perfect-lemma-descend-relatively-perfect},
on peut supposer qu'il existe un objet $R_i$-parfait $E_i$ de
$D(\mathcal{O}_{X_i})$ dont l'image inverse sur $X$ est $E$.
En appliquant le lemme \ref{perfect-lemma-perfect-direct-image}
à $X_i \to \Spec(R_i)$ et à $E_i$, puis en utilisant la
propriété de changement de base déjà démontrée, on obtient le résultat.
\end{proof}

\begin{lemma}
\label{perfect-lemma-compute-ext-rel-perfect}
Soit $f : X \to S$ un morphisme de schémas. Soient $E, K \in D(\mathcal{O}_X)$.
Supposons que
\begin{enumerate}
\item $S$ soit quasi-compact et quasi-séparé ;
\item $f$ soit propre, plat et de présentation finie ;
\item $E$ soit $S$-parfait ;
\item $K$ soit pseudo-cohérent.
\end{enumerate}
Alors il existe un $L \in D(\mathcal{O}_S)$ pseudo-cohérent tel que
$$
Rf_*R\SheafHom(K, E) = R\SheafHom(L, \mathcal{O}_S)
$$
et il en va de même après tout changement de base : étant donné
$$
\vcenter{
\xymatrix{
X' \ar[r]_{g'} \ar[d]_{f'} &
X \ar[d]^f \\
S' \ar[r]^g &
S
}
}
\quad\quad
\begin{matrix}
\text{cartésien, alors on a } \\
Rf'_*R\SheafHom(L(g')^*K, L(g')^*E) \\
= R\SheafHom(Lg^*L, \mathcal{O}_{S'})
\end{matrix}
$$
\end{lemma}

\begin{proof}
Puisque $S$ est quasi-compact et quasi-séparé, il en va de même de $X$.
D'après le lemme \ref{perfect-lemma-pseudo-coherent-hocolim}, on peut écrire
$K = \text{hocolim} K_n$, où $K_n$ est parfait et où $K_n \to K$ induit
un isomorphisme sur les troncatures $\tau_{\geq -n}$. Soit $K_n^\vee$
le complexe parfait dual
(Cohomologie, lemme \ref{cohomology-lemma-dual-perfect-complex}).
On obtient un système projectif $\ldots \to K_3^\vee \to K_2^\vee \to K_1^\vee$
d'objets parfaits. D'après le lemme \ref{perfect-lemma-perfect-relatively-perfect},
on voit que $K_n^\vee \otimes_{\mathcal{O}_X} E$ est $S$-parfait.
On peut donc appliquer le lemme \ref{perfect-lemma-derived-pushforward-rel-perfect}
à $K_n^\vee \otimes_{\mathcal{O}_X} E$, et l'on obtient un système projectif
$$
\ldots \to M_3 \to M_2 \to M_1
$$
de complexes parfaits sur $S$, avec
$$
M_n = Rf_*(K_n^\vee \otimes_{\mathcal{O}_X}^\mathbf{L} E) =
Rf_*R\SheafHom(K_n, E)
$$
De plus, la formation de ces complexes commute à tout
changement de base, à savoir $Lg^*M_n =
Rf'_*((L(g')^*K_n)^\vee \otimes_{\mathcal{O}_{X'}}^\mathbf{L} L(g')^*E) =
Rf'_*R\SheafHom(L(g')^*K_n, L(g')^*E)$.

\medskip\noindent
Puisque $K_n \to K$ induit un isomorphisme sur $\tau_{\geq -n}$, on voit que
$K_n \to K_{n + 1}$ induit un isomorphisme sur $\tau_{\geq -n}$.
Il en résulte que $K_{n + 1}^\vee \to K_n^\vee$
induit un isomorphisme sur $\tau_{\leq n}$, car
$K_n^\vee = R\SheafHom(K_n, \mathcal{O}_X)$.
Supposons que $E$ ait une amplitude de Tor contenue dans $[a, b]$ comme complexe
de $f^{-1}\mathcal{O}_Y$-modules. Il en va alors de même après
tout changement de base ; voir le lemme \ref{perfect-lemma-tor-independence-and-tor-amplitude}.
On constate que
$K_{n + 1}^\vee \otimes_{\mathcal{O}_X} E \to
K_n^\vee \otimes_{\mathcal{O}_X} E$
induit un isomorphisme sur $\tau_{\leq n + a}$,
et il en va de même après tout changement de base.
En appliquant le foncteur dérivé droit $Rf_*$,
on en déduit que les morphismes $M_{n + 1} \to M_n$
induisent des isomorphismes sur $\tau_{\leq n + a}$,
et il en va de même après tout changement de base.
Choisissons un triangle distingué
$$
M_{n + 1} \to M_n \to C_n \to M_{n + 1}[1]
$$
Prenons pour $S'$ le spectre du corps résiduel en un point
$s \in S$ et effectuons l'image inverse pour voir que
$C_n \otimes_{\mathcal{O}_S}^\mathbf{L} \kappa(s)$
n'a de cohomologie non nulle que dans les degrés $\geq n + a$. D'après
Compléments sur l'algèbre, lemme
\ref{more-algebra-lemma-lift-perfect-from-residue-field},
on voit que le complexe parfait $C_n$ a une amplitude de Tor contenue dans
$[n + a, m_n]$ pour un certain entier $m_n$.
En particulier, le complexe parfait dual $C_n^\vee$ a une amplitude de Tor contenue dans
$[-m_n, -n - a]$.

\medskip\noindent
Soit $L_n = M_n^\vee$ le complexe parfait dual. La
conclusion de la discussion du paragraphe précédent est que
$L_n \to L_{n + 1}$ induit des isomorphismes sur $\tau_{\geq -n - a}$.
Ainsi, $L = \text{hocolim} L_n$ est pseudo-cohérent ; voir le
lemme \ref{perfect-lemma-pseudo-coherent-hocolim}.
Puisque l'on a
$$
R\SheafHom(K, E) = R\SheafHom(\text{hocolim} K_n, E) =
R\lim R\SheafHom(K_n, E) = R\lim K_n^\vee \otimes_{\mathcal{O}_X} E
$$
(Cohomologie, lemme \ref{cohomology-lemma-colim-and-lim-of-duals})
et puisque $R\lim$ commute à $Rf_*$, on trouve que
$$
Rf_*R\SheafHom(K, E) = R\lim M_n = R\lim R\SheafHom(L_n, \mathcal{O}_S) =
R\SheafHom(L, \mathcal{O}_S)
$$
Cela démontre la formule sur $S$. Puisque la construction de $M_n$ est
compatible au changement de base, la formule reste vraie après
tout changement de base.
\end{proof}

\begin{remark}
\label{perfect-remark-compare-L}
Le lecteur aura peut-être remarqué la similitude entre le
lemme \ref{perfect-lemma-compute-ext-rel-perfect} et le
lemme \ref{perfect-lemma-compute-ext}.
En effet, le complexe pseudo-cohérent $L$ du
lemme \ref{perfect-lemma-compute-ext-rel-perfect}
peut être caractérisé comme l'unique complexe pseudo-cohérent
sur $S$ tel qu'il existe des isomorphismes fonctoriels
$$
\Ext^i_{\mathcal{O}_S}(L, \mathcal{F}) \longrightarrow
\Ext^i_{\mathcal{O}_X}(K,
E \otimes_{\mathcal{O}_X}^\mathbf{L} Lf^*\mathcal{F})
$$
compatibles aux morphismes de bord, lorsque $\mathcal{F}$ parcourt
$\QCoh(\mathcal{O}_S)$. Si cela nous est un jour nécessaire, nous
formulerons ici un résultat précis et en donnerons une démonstration détaillée.
\end{remark}

\begin{lemma}
\label{perfect-lemma-bounded-on-fibres}
Soit $f : X \to S$ un morphisme de schémas plat et
localement de présentation finie. Soit $E$ un objet pseudo-cohérent
de $D(\mathcal{O}_X)$. Les conditions suivantes sont équivalentes :
\begin{enumerate}
\item $E$ est $S$-parfait ;
\item $E$ est localement borné inférieurement et, pour tout point $s \in S$,
l'objet $L(X_s \to X)^*E$ de $D(\mathcal{O}_{X_s})$
est localement borné inférieurement.
\end{enumerate}
\end{lemma}

\begin{proof}
Puisque tout est local, on se ramène immédiatement au
cas où $X$ et $S$ sont affines ; voir le lemme
\ref{perfect-lemma-affine-locally-rel-perfect}.
Supposons que $X \to S$ corresponde à $\Spec(A) \to \Spec(R)$ et que
$E$ corresponde à $K$ dans $D(A)$. Si $s$ correspond à
l'idéal premier $\mathfrak p \subset R$, alors $L(X_s \to X)^*E$
correspond à $K \otimes_R^\mathbf{L} \kappa(\mathfrak p)$,
car $R \to A$ est plat ; voir, par exemple, le
lemme \ref{perfect-lemma-compare-base-change}.
On voit donc que notre lemme résulte du résultat algébrique
correspondant ; voir Compléments sur l'algèbre, lemme
\ref{more-algebra-lemma-bounded-on-fibres-relatively-perfect}.
\end{proof}

\begin{remark}
\label{perfect-remark-discuss-rel-perfect}
Notre définition \ref{perfect-definition-relatively-perfect} d'un
complexe relativement parfait est équivalente à celle donnée
dans \cite{lieblich-complexes} chaque fois que notre définition s'applique\footnote{Pour
le voir, utiliser le lemme \ref{perfect-lemma-affine-locally-rel-perfect} et
Compléments sur l'algèbre, lemme \ref{more-algebra-lemma-structure-relatively-perfect}.}.
Supposons maintenant que $f : X \to S$ soit seulement localement
de type fini (sans être nécessairement plat ni localement de
présentation finie). La définition de l'article cité ci-dessus est que
$E \in D(\mathcal{O}_X)$ est relativement parfait si
\begin{enumerate}
\item[(A)] localement sur $X$, l'objet $E$ est
quasi-isomorphe à un complexe fini de
modules plats sur $S$ et de présentation finie comme $\mathcal{O}_X$-modules.
\end{enumerate}
D'autre part, la généralisation naturelle de notre
définition \ref{perfect-definition-relatively-perfect} est la suivante :
\begin{enumerate}
\item[(B)] $E$ est pseudo-cohérent relativement à $S$
(Compléments sur les morphismes, définition
\ref{more-morphisms-definition-relative-pseudo-coherence})
et $E$ a localement une dimension de Tor finie comme objet de
$D(f^{-1}\mathcal{O}_S)$
(Cohomologie, définition \ref{cohomology-definition-tor-amplitude}).
\end{enumerate}
L'avantage de la condition (B) est qu'elle définit manifestement une sous-catégorie
triangulée de $D(\mathcal{O}_X)$, tandis que nous soupçonnons que ce n'est pas
le cas de la condition (A). L'avantage de la condition (A)
est qu'elle est plus facile à manier, en particulier en ce qui concerne les limites.
\end{remark}








\section{La propriété de résolution}
\label{perfect-section-resolution-property}

\noindent
Cette notion est étudiée dans l'article \cite{totaro_resolution} ;
la discussion se poursuit dans \cite{Gross-thesis},
\cite{Gross-surface} et \cite{Gross-stack}.
On ignore actuellement si tout schéma propre sur un corps
possède toujours la propriété de résolution ou si cela est faux.
Si vous connaissez la réponse à cette question, veuillez écrire à
\href{mailto:stacks.project@gmail.com}{stacks.project@gmail.com}.

\medskip\noindent
On peut donner la définition suivante, bien qu'il ne soit guère pertinent
de la considérer pour des schémas généraux.

\begin{definition}
\label{perfect-definition-resolution-property}
Soit $X$ un schéma. On dit que $X$ possède la {\it propriété de résolution}
si tout $\mathcal{O}_X$-module quasi-cohérent de type fini
est quotient d'un $\mathcal{O}_X$-module localement libre de type fini.
\end{definition}

\noindent
Si $X$ est un schéma quasi-compact et quasi-séparé, il suffit de vérifier que
tout $\mathcal{O}_X$-module de présentation finie (automatiquement
quasi-cohérent) est quotient d'un module localement libre de type fini
sur $\mathcal{O}_X$ ; voir Propriétés, lemme
\ref{properties-lemma-finite-directed-colimit-surjective-maps}.
Si $X$ est un schéma noethérien, les modules quasi-cohérents de type fini
sont exactement les $\mathcal{O}_X$-modules cohérents ; voir
Cohomologie des schémas, lemme \ref{coherent-lemma-coherent-Noetherian}.

\begin{lemma}
\label{perfect-lemma-resolution-property-ample}
Soit $X$ un schéma. Si $X$ possède un $\mathcal{O}_X$-module inversible ample,
alors $X$ possède la propriété de résolution.
\end{lemma}

\begin{proof}
C'est une conséquence immédiate de Propriétés, proposition
\ref{properties-proposition-characterize-ample}.
\end{proof}

\begin{lemma}
\label{perfect-lemma-resolution-property-ample-relative}
Soit $f : X \to Y$ un morphisme de schémas. Supposons que
\begin{enumerate}
\item $Y$ soit quasi-compact et quasi-séparé et possède la propriété de résolution ;
\item il existe un module inversible $f$-ample sur $X$.
\end{enumerate}
Alors $X$ possède la propriété de résolution.
\end{lemma}

\begin{proof}
Soit $\mathcal{F}$ un $\mathcal{O}_X$-module quasi-cohérent de type fini.
Soit $\mathcal{L}$ un module inversible $f$-ample.
Choisissons un recouvrement ouvert affine $Y = V_1 \cup \ldots \cup V_m$.
Posons $U_j = f^{-1}(V_j)$. D'après Propriétés, proposition
\ref{properties-proposition-characterize-ample},
pour tout $j$, on sait qu'il existe un nombre fini de morphismes
$s_{j, i} : \mathcal{L}^{\otimes n_{j, i}}|_{U_j} \to \mathcal{F}|_{U_j}$
qui sont conjointement surjectifs. Considérons les
$\mathcal{O}_Y$-modules quasi-cohérents
$$
\mathcal{H}_n =
f_*(\mathcal{F} \otimes_{\mathcal{O}_X} \mathcal{L}^{\otimes n})
$$
On peut considérer $s_{j, i}$ comme une section sur $V_j$ du faisceau
$\mathcal{H}_{-n_{j, i}}$. Supposons que l'on puisse trouver des
$\mathcal{O}_Y$-modules localement libres de type fini $\mathcal{E}_{i, j}$ et des morphismes
$\mathcal{E}_{i, j} \to \mathcal{H}_{-n_{j, i}}$ tels que
$s_{j, i}$ appartienne à leur image. Alors les morphismes correspondants
$$
f^*\mathcal{E}_{i, j}
\otimes_{\mathcal{O}_X}
\mathcal{L}^{\otimes n_{i, j}} \longrightarrow \mathcal{F}
$$
seront conjointement surjectifs, ce qui démontrera le lemme. D'après
Propriétés, lemme \ref{properties-lemma-quasi-coherent-colimit-finite-type},
pour tout $i, j$, on peut trouver un sous-module quasi-cohérent
de type fini
$\mathcal{H}'_{i, j} \subset  \mathcal{H}_{-n_{j, i}}$
qui contient la section $s_{i, j}$ sur $V_j$.
La propriété de résolution de $Y$ fournit donc des surjections
$\mathcal{E}_{i, j} \to \mathcal{H}'_{j, i}$, et l'on conclut.
\end{proof}

\begin{lemma}
\label{perfect-lemma-resolution-property-goes-up-affine}
Soit $f : X \to Y$ un morphisme affine ou quasi-affine de schémas, où
$Y$ est quasi-compact et quasi-séparé.
Si $Y$ possède la propriété de résolution, alors $X$ la possède aussi.
\end{lemma}

\begin{proof}
D'après Morphismes, lemme \ref{morphisms-lemma-quasi-affine-O-ample},
c'est un cas particulier du
lemme \ref{perfect-lemma-resolution-property-ample-relative}.
\end{proof}

\noindent
Voici un cas où l'on peut démontrer que la propriété de résolution descend.

\begin{lemma}
\label{perfect-lemma-resolution-property-goes-down-finite-flat}
Soit $f : X \to Y$ un morphisme de schémas surjectif et fini localement libre.
Si $X$ possède la propriété de résolution, alors $Y$ la possède aussi.
\end{lemma}

\begin{proof}
La condition signifie que $f$ est affine et que $f_*\mathcal{O}_X$
est un $\mathcal{O}_Y$-module localement libre de type fini et de rang positif.
Soit $\mathcal{G}$ un $\mathcal{O}_Y$-module quasi-cohérent
de type fini. Par hypothèse, il existe une surjection
$\mathcal{E} \to f^*\mathcal{G}$ pour un certain $\mathcal{O}_X$-module
$\mathcal{E}$ localement libre de type fini. Puisque $f_*$ est exact sur les
modules quasi-cohérents
(Cohomologie des schémas, lemme \ref{coherent-lemma-relative-affine-vanishing}),
on obtient une surjection
$$
f_*\mathcal{E}
\longrightarrow
f_*f^*\mathcal{G} = \mathcal{G} \otimes_{\mathcal{O}_Y} f_*\mathcal{O}_X
$$
En prenant les duaux, on obtient une surjection
$$
f_*\mathcal{E}
\otimes_{\mathcal{O}_Y}
\SheafHom_{\mathcal{O}_Y}(f_*\mathcal{O}_X, \mathcal{O}_Y)
\longrightarrow
\mathcal{G}
$$
Puisque $f_*\mathcal{E}$ est localement libre de type fini\footnote{En effet,
si $A \to B$ est un homomorphisme d'anneaux fini localement libre
et si $N$ est un $B$-module localement libre de type fini, alors $N$
est un $A$-module localement libre de type fini. Pour le voir, notons d'abord
que le fait que $N$ soit localement libre de type fini sur $B$ implique que $N$ est
plat et de présentation finie comme $B$-module ; voir
Algèbre, lemme \ref{algebra-lemma-finite-projective}. Ensuite,
$N$ est un $A$-module de présentation finie d'après
Algèbre, lemme \ref{algebra-lemma-finite-finitely-presented-extension},
et un $A$-module plat d'après Algèbre, lemme \ref{algebra-lemma-composition-flat}.
On conclut alors en appliquant
Algèbre, lemme \ref{algebra-lemma-finite-projective}, sur $A$.},
on conclut.
\end{proof}

\begin{lemma}
\label{perfect-lemma-resolution-property-finite-number}
Soit $X$ un schéma. Supposons donnés
\begin{enumerate}
\item un recouvrement ouvert affine fini $X = U_1 \cup \ldots \cup U_m$ ;
\item des idéaux quasi-cohérents de type fini $\mathcal{I}_j$
tels que $V(\mathcal{I}_j) = X \setminus U_j$.
\end{enumerate}
Alors $X$ possède la propriété de résolution si et seulement si $\mathcal{I}_j$
est quotient d'un $\mathcal{O}_X$-module localement libre de type fini
pour $j = 1, \ldots, m$.
\end{lemma}

\begin{proof}
Une implication du lemme est triviale. Pour l'autre,
supposons que $\mathcal{E}_j \to \mathcal{I}_j$ soit une surjection, avec
$\mathcal{E}_j$ localement libre de type fini. Dans le paragraphe suivant, nous
nous ramenons au cas noethérien ; nous suggérons au lecteur de l'omettre.

\medskip\noindent
La première observation est que $U_j \cap U_{j'}$ est quasi-compact,
car c'est le complémentaire du schéma des zéros de l'idéal quasi-cohérent de type fini
$\mathcal{I}_{j'}|{U_j}$ sur le schéma affine $U_j$ ; voir
Propriétés, lemme \ref{properties-lemma-quasi-coherent-finite-type-ideals}.
Ainsi, $X$ est quasi-compact et quasi-séparé ; voir
Schémas, lemme \ref{schemes-lemma-characterize-quasi-separated}.
D'après Limites, proposition \ref{limits-proposition-approximate},
on peut écrire $X = \lim X_i$ comme limite d'un système inductif de
schémas noethériens à morphismes de transition affines. Pour tout $j$,
on peut trouver un $i$ et un $\mathcal{O}_{X_i}$-module localement libre de type fini
$\mathcal{E}_{i, j}$ dont l'image inverse est $\mathcal{E}_j$ ; voir
Limites, lemme \ref{limits-lemma-descend-invertible-modules}.
Après avoir augmenté $i$, on peut supposer que la composée
$\mathcal{E}_j \to \mathcal{I}_j \to \mathcal{O}_X$ est l'image
inverse d'un morphisme $\mathcal{E}_{i, j} \to \mathcal{O}_{X_i}$ ; voir
Limites, lemme \ref{limits-lemma-descend-modules-finite-presentation}.
Notons $\mathcal{I}_{i, j} \subset \mathcal{O}_{X_i}$ l'image
de ce morphisme ; c'est un faisceau d'idéaux quasi-cohérent sur le schéma noethérien
$X_i$, dont l'image inverse sur $X$ est $\mathcal{I}_j$.
En notant $U_{i, j} \subset X_i$ les ouverts complémentaires, on
peut supposer qu'ils sont affines pour tous $i, j$ ; voir
Limites, lemme \ref{limits-lemma-limit-affine}.
Si l'on peut démontrer le lemme pour les ouverts $U_{i, j}$ et les
faisceaux d'idéaux $\mathcal{I}_{i, j}$ sur $X_i$, alors $X$, étant affine
sur $X_i$, possédera la propriété de résolution d'après le
lemme \ref{perfect-lemma-resolution-property-goes-up-affine}.
De cette manière, on se ramène au cas d'un schéma noethérien.

\medskip\noindent
Supposons $X$ noethérien. Pour tout module cohérent $\mathcal{F}$,
on peut choisir une liste finie de sections $s_{jk} \in \mathcal{F}(U_j)$,
$k = 1, \ldots, e_j$, qui engendrent la restriction de $\mathcal{F}$ à $U_j$.
D'après Cohomologie des schémas, lemme \ref{coherent-lemma-homs-over-open},
on peut prolonger $s_{jk}$ en un morphisme
$s'_{jk} : \mathcal{I}_i^{n_{jk}} \to \mathcal{F}$ pour un certain $n_{jk} \geq 1$.
On peut alors considérer les composées
$$
\mathcal{E}_j^{\otimes n_{jk}} \to \mathcal{I}_j^{n_{jk}} \to \mathcal{F}
$$
pour conclure.
\end{proof}

\begin{lemma}
\label{perfect-lemma-resolution-property-ample-family}
Soit $X$ un schéma. Si $X$ possède une famille ample de modules inversibles
(Morphismes, définition
\ref{morphisms-definition-family-ample-invertible-modules}),
alors $X$ possède la propriété de résolution.
\end{lemma}

\begin{proof}
Puisque $X$ est quasi-compact, il existe un $n$ et des couples
$(\mathcal{L}_i, s_i)$, $i = 1, \ldots, n$, où
$\mathcal{L}_i$ est un $\mathcal{O}_X$-module inversible et
$s_i \in \Gamma(X, \mathcal{L}_i)$ est une section
tels que l'ensemble des points $U_i \subset X$ où $s_i$ ne s'annule pas
soit affine et que $X = U_1 \cup \ldots \cup U_n$.
Soit $\mathcal{I}_i \subset \mathcal{O}_X$ l'image
de $s_i : \mathcal{L}_i^{\otimes -1} \to \mathcal{O}_X$.
En appliquant le lemme \ref{perfect-lemma-resolution-property-finite-number},
on constate que $X$ possède la propriété de résolution.
\end{proof}

\begin{lemma}
\label{perfect-lemma-regular-resolution-property}
Soit $X$ un schéma quasi-compact régulier à diagonale affine.
Alors $X$ possède la propriété de résolution.
\end{lemma}

\begin{proof}
On combine Diviseurs, lemme \ref{divisors-lemma-regular-ample-family}, et le
lemme \ref{perfect-lemma-resolution-property-ample-family} ci-dessus.
\end{proof}

\begin{lemma}
\label{perfect-lemma-resolution-property-descends}
Soit $X = \lim X_i$ la limite d'un système inductif de schémas quasi-compacts
et quasi-séparés à morphismes de transition affines.
Alors $X$ possède la propriété de résolution si et seulement si $X_i$ possède
la propriété de résolution pour un certain $i$.
\end{lemma}

\begin{proof}
Si $X_i$ possède la propriété de résolution, alors $X$ la possède d'après le
lemme \ref{perfect-lemma-resolution-property-goes-up-affine}.
Supposons que $X$ possède la propriété de résolution.
Choisissons $i \in I$. Choisissons un recouvrement ouvert affine fini
$X_i = U_{i, 1} \cup \ldots \cup U_{i, m}$.
Pour tout $j$, choisissons un faisceau quasi-cohérent de type fini
d'idéaux $\mathcal{I}_{i, j} \subset \mathcal{O}_{X_i}$
tel que $X_i \setminus V(\mathcal{I}_{i, j}) = U_{i, j}$ ; voir
Propriétés, lemme \ref{properties-lemma-quasi-coherent-finite-type-ideals}.
Notons $U_j \subset X$ l'image inverse de $U_{i, j}$,
et notons $\mathcal{I}_j \subset \mathcal{O}_X$
l'image inverse de $\mathcal{I}_{i, j}$.
Puisque $X$ possède la propriété de résolution, on peut
choisir des $\mathcal{O}_X$-modules localement libres de type fini $\mathcal{E}_j$
et des surjections $\mathcal{E}_j \to \mathcal{I}_j$.
D'après Limites, lemmes \ref{limits-lemma-descend-invertible-modules} et
\ref{limits-lemma-descend-modules-finite-presentation},
après avoir augmenté $i$, on peut trouver des
$\mathcal{O}_{X_i}$-modules localement libres de type fini $\mathcal{E}_{i, j}$
et des morphismes $\mathcal{E}_{i, j} \to \mathcal{O}_{X_i}$
dont les changements de base à $X$ redonnent les composées
$\mathcal{E}_j \to \mathcal{I}_j \to \mathcal{O}_X$, $j = 1, \ldots, m$.
Les images inverses des $\mathcal{O}_{X_i}$-modules de présentation finie
$\Coker(\mathcal{E}_{i, j} \to \mathcal{O}_{X_i})$
et $\mathcal{O}_{X_i}/\mathcal{I}_{i, j}$
sur $X$ coïncident comme quotients de $\mathcal{O}_X$.
Par conséquent, d'après Limites, lemme
\ref{limits-lemma-descend-modules-finite-presentation},
on peut supposer qu'ils coïncident, autrement dit que
l'image de $\mathcal{E}_{i, j} \to \mathcal{O}_{X_i}$
est égale à $\mathcal{I}_{i, j}$.
On en déduit que $X_i$ possède la propriété de résolution
d'après le lemme \ref{perfect-lemma-resolution-property-finite-number}.
\end{proof}

\begin{lemma}
\label{perfect-lemma-resolution-property-affine-diagonal}
\begin{reference}
Cas particulier de \cite[Proposition 1.3]{totaro_resolution}.
\end{reference}
Soit $X$ un schéma quasi-compact et quasi-séparé possédant
la propriété de résolution. Alors $X$ a une diagonale affine.
\end{lemma}

\begin{proof}
En combinant Limites, proposition \ref{limits-proposition-approximate},
et le lemme \ref{perfect-lemma-resolution-property-descends},
on se ramène au cas où $X$ est noethérien (un petit détail est omis).
Supposons $X$ noethérien.
Rappelons que $X \times X$ est recouvert par les ouverts affines
$U \times V$, où $U$, $V$ sont des ouverts affines de $X$ ; voir
Schémas, section \ref{schemes-section-fibre-products}.
Ainsi, pour montrer que la diagonale $\Delta : X \to X \times X$
est affine, il suffit de montrer que $U \cap V = \Delta^{-1}(U \times V)$
est affine pour tous ouverts affines $U$, $V$ de $X$ ; voir
Morphismes, lemme \ref{morphisms-lemma-characterize-affine}.
En particulier, il suffit de montrer que le morphisme d'inclusion
$j : U \to X$ est affine si $U$ est un ouvert affine de $X$.
D'après Cohomologie des schémas, lemme \ref{coherent-lemma-criterion-affine-morphism},
il suffit de montrer que $R^1j_*\mathcal{G} = 0$ pour tout
$\mathcal{O}_U$-module quasi-cohérent $\mathcal{G}$.
D'après la proposition \ref{perfect-proposition-Noetherian} (c'est ici que l'on utilise
la réduction au cas noethérien),
on peut représenter $Rj_*\mathcal{G}$ par un complexe
$\mathcal{H}^\bullet$ de $\mathcal{O}_X$-modules quasi-cohérents.
Supposons que
$$
H^1(\mathcal{H}^\bullet) =
\Ker(\mathcal{H}^1 \to \mathcal{H}^2)/\Im(\mathcal{H}^0 \to \mathcal{H}^1)
$$
soit non nul, afin d'obtenir une contradiction. On peut alors trouver un
$\mathcal{O}_X$-module cohérent $\mathcal{F}$ et un morphisme
$$
\mathcal{F} \longrightarrow
\Ker(\mathcal{H}^1 \to \mathcal{H}^2)
$$
tel que sa composée avec la projection sur $H^1(\mathcal{H}^\bullet)$
soit non nulle. En effet, on peut écrire $\Ker(\mathcal{H}^1 \to \mathcal{H}^2)$
comme réunion filtrante de ses sous-modules cohérents d'après
Propriétés, lemme \ref{properties-lemma-quasi-coherent-colimit-finite-type},
et l'un d'eux conviendra.
Ensuite, on choisit un $\mathcal{O}_X$-module localement libre de type fini
$\mathcal{E}$ et une surjection $\mathcal{E} \to \mathcal{F}$ en utilisant
la propriété de résolution de $X$.
Cela fournit un morphisme dans la catégorie dérivée
$$
\mathcal{E}[-1] \longrightarrow Rj_*\mathcal{G}
$$
qui est non nul sur les faisceaux de cohomologie et donc non nul dans $D(\mathcal{O}_X)$.
Par adjonction, cela revient à un morphisme
$$
j^*\mathcal{E}[-1] \to \mathcal{G}
$$
non nul dans $D(\mathcal{O}_U)$. Puisque $\mathcal{E}$ est localement libre
de type fini, cela revient à un élément non nul de
$$
H^1(U, j^*\mathcal{E}^\vee \otimes_{\mathcal{O}_U} \mathcal{G})
$$
où
$\mathcal{E}^\vee = \SheafHom_{\mathcal{O}_X}(\mathcal{E}, \mathcal{O}_X)$
est le module localement libre de type fini dual. Cependant, ce groupe est nul d'après
Cohomologie des schémas, lemme
\ref{coherent-lemma-quasi-coherent-affine-cohomology-zero},
ce qui est la contradiction recherchée.
(En cas de doute sur l'étape utilisant les duaux, voir le résultat plus général
Cohomologie, lemme \ref{cohomology-lemma-dual-perfect-complex}.)
\end{proof}





\section{La propriété de résolution et les complexes parfaits}
\label{perfect-section-resolution-property-perfect}

\noindent
Dans cette section, nous examinons la relation entre les complexes parfaits
et les complexes strictement parfaits sur les schémas qui possèdent la propriété de résolution.

\begin{lemma}
\label{perfect-lemma-construct-strictly-perfect}
Soit $X$ un schéma quasi-compact et quasi-séparé possédant la
propriété de résolution.
Soit $\mathcal{F}^\bullet$ un complexe borné inférieurement de
$\mathcal{O}_X$-modules quasi-cohérents représentant un objet parfait de
$D(\mathcal{O}_X)$. Alors il existe un complexe borné
$\mathcal{E}^\bullet$ de $\mathcal{O}_X$-modules localement libres de type fini
et un quasi-isomorphisme $\mathcal{E}^\bullet \to \mathcal{F}^\bullet$.
\end{lemma}

\begin{proof}
Soient $a, b \in \mathbf{Z}$ des entiers tels que $\mathcal{F}^\bullet$
ait une amplitude de Tor contenue dans $[a, b]$ et tels que $\mathcal{F}^n = 0$ pour
$n < a$. L'existence d'un tel couple d'entiers
résulte de Cohomologie, lemme \ref{cohomology-lemma-perfect},
et du fait que $X$ est quasi-compact.
Si $b < a$, alors $\mathcal{F}^\bullet$ est nul dans la
catégorie dérivée et le lemme en résulte.
Nous allons démontrer par récurrence sur
$b - a \geq 0$ qu'il existe un complexe
$\mathcal{E}^a \to \ldots \to \mathcal{E}^b$
où chaque $\mathcal{E}^i$ est localement libre de type fini,
et un quasi-isomorphisme $\mathcal{E}^\bullet \to \mathcal{F}^\bullet$.

\medskip\noindent
Le cas initial est celui où $b - a = 0$. Dans ce cas,
$H^b(\mathcal{F}^\bullet) = H^a(\mathcal{F}^\bullet) =
\Ker(\mathcal{F}^a \to \mathcal{F}^{a + 1})$
est localement libre de type fini. En effet, c'est un
$\mathcal{O}_X$-module de présentation finie, ayant une dimension de Tor égale à $0$, donc
localement libre de type fini. Voir Cohomologie, lemmes \ref{cohomology-lemma-perfect} et
\ref{cohomology-lemma-finite-cohomology}, ainsi que
Propriétés, lemme \ref{properties-lemma-finite-locally-free}.
Ainsi, on peut prendre pour $\mathcal{E}^\bullet$ le complexe
$H^b(\mathcal{F}^\bullet)$ placé en degré $b$.
Le reste de la démonstration est consacré à l'étape de récurrence.

\medskip\noindent
Supposons $b > a$. Observons que
$$
H^b(\mathcal{F}^\bullet) =
\Ker(\mathcal{F}^b \to \mathcal{F}^{b + 1})/
\Im(\mathcal{F}^{b - 1} \to \mathcal{F}^b)
$$
est un $\mathcal{O}_X$-module quasi-cohérent de type fini ; voir
Cohomologie, lemmes \ref{cohomology-lemma-perfect} et
\ref{cohomology-lemma-finite-cohomology}. On peut alors trouver un
$\mathcal{O}_X$-module quasi-cohérent de type fini $\mathcal{F}$ et un morphisme
$$
\mathcal{F} \longrightarrow
\Ker(\mathcal{F}^b \to \mathcal{F}^{b + 1})
$$
tel que sa composée avec la projection sur
$H^b(\mathcal{F}^\bullet)$ soit surjective.
En effet, on peut écrire $\Ker(\mathcal{F}^b \to \mathcal{F}^{b + 1})$
comme la réunion filtrante de ses sous-modules quasi-cohérents de type fini d'après
Propriétés, lemme \ref{properties-lemma-quasi-coherent-colimit-finite-type},
et l'un d'eux conviendra.
Ensuite, on choisit un $\mathcal{O}_X$-module localement libre de type fini
$\mathcal{E}^b$ et une surjection $\mathcal{E}^b \to \mathcal{F}$ en utilisant
la propriété de résolution de $X$. Considérons le morphisme de complexes
$$
\alpha : \mathcal{E}^b[-b] \to \mathcal{F}^\bullet
$$
et son cône $C(\alpha)^\bullet$ ; voir
Catégories dérivées, définition \ref{derived-definition-cone}.
Observons que $C(\alpha)^\bullet$ n'est non nul qu'en degrés
$\geq a$, a une amplitude de Tor contenue dans $[a, b]$ d'après
Cohomologie, lemme \ref{cohomology-lemma-cone-tor-amplitude},
et vérifie $H^b(C(\alpha)^\bullet) = 0$ par construction.
Ainsi, on trouve en fait que $C(\alpha)^\bullet$ a une amplitude de Tor
contenue dans $[a, b - 1]$. L'hypothèse de récurrence s'applique donc à
$C(\alpha)^\bullet$ et l'on obtient un morphisme de complexes
$$
(\mathcal{E}^a \to \ldots \to \mathcal{E}^{b - 1})
\longrightarrow
C(\alpha)^\bullet
$$
ayant les propriétés énoncées dans l'hypothèse de récurrence. En explicitant
la définition du cône, on obtient le diagramme commutatif
$$
\xymatrix{
\ldots \ar[r] &
\mathcal{E}^{b - 2} \ar[r] \ar[d] &
\mathcal{E}^{b - 1} \ar[r] \ar[d] &
0 \ar[r] \ar[d] &
\ldots \\
\ldots \ar[r] &
\mathcal{F}^{b - 2} \ar[r] &
\mathcal{F}^{b - 1} \oplus \mathcal{E}^b \ar[r] &
\mathcal{F}^b \ar[r] &
\ldots
}
$$
Il est clair que l'on obtient un morphisme de complexes
$(\mathcal{E}^a \to \ldots \to \mathcal{E}^b) \to \mathcal{F}^\bullet$.
Nous omettons de vérifier que ce morphisme est un quasi-isomorphisme.
\end{proof}

\begin{lemma}
\label{perfect-lemma-resolution-property-perfect-complex}
Soit $X$ un schéma quasi-compact et quasi-séparé possédant la
propriété de résolution. Alors tout objet parfait de $D(\mathcal{O}_X)$
peut être représenté par un complexe borné de
$\mathcal{O}_X$-modules localement libres de type fini.
\end{lemma}

\begin{proof}
Soit $E$ un objet parfait de $D(\mathcal{O}_X)$.
D'après le lemme \ref{perfect-lemma-resolution-property-affine-diagonal},
on voit que $X$ a une diagonale affine. Par conséquent, d'après la
proposition \ref{perfect-proposition-quasi-compact-affine-diagonal},
on peut représenter $E$ par un complexe $\mathcal{F}^\bullet$
de $\mathcal{O}_X$-modules quasi-cohérents.
Observons que $E$ appartient à $D^b(\mathcal{O}_X)$ parce que
$X$ est quasi-compact. Par conséquent, $\tau_{\geq n}\mathcal{F}^\bullet$
est un complexe borné inférieurement de $\mathcal{O}_X$-modules quasi-cohérents
qui représente $E$ si $n \ll 0$. On peut donc appliquer le
lemme \ref{perfect-lemma-construct-strictly-perfect} pour conclure.
\end{proof}

\begin{lemma}
\label{perfect-lemma-resolution-property-map-perfect-complex}
Soit $X$ un schéma quasi-compact et quasi-séparé possédant la
propriété de résolution. Soient $\mathcal{E}^\bullet$ et $\mathcal{F}^\bullet$
des complexes finis de $\mathcal{O}_X$-modules localement libres de type fini.
Alors tout
$\alpha \in \Hom_{D(\mathcal{O}_X)}(\mathcal{E}^\bullet, \mathcal{F}^\bullet)$
peut être représenté par un diagramme
$$
\mathcal{E}^\bullet \leftarrow \mathcal{G}^\bullet \to \mathcal{F}^\bullet
$$
où $\mathcal{G}^\bullet$ est un complexe borné dont chaque terme est localement libre de type fini en tant que
$\mathcal{O}_X$-module et où $\mathcal{G}^\bullet \to \mathcal{E}^\bullet$
est un quasi-isomorphisme.
\end{lemma}

\begin{proof}
D'après le lemme \ref{perfect-lemma-resolution-property-affine-diagonal},
on voit que $X$ a une diagonale affine. Par conséquent, d'après la
proposition \ref{perfect-proposition-quasi-compact-affine-diagonal},
on peut représenter $\alpha$ par un diagramme
$$
\mathcal{E}^\bullet \leftarrow \mathcal{H}^\bullet \to \mathcal{F}^\bullet
$$
où $\mathcal{H}^\bullet$ est un complexe de
$\mathcal{O}_X$-modules quasi-cohérents et où
$\mathcal{H}^\bullet \to \mathcal{E}^\bullet$
est un quasi-isomorphisme. Pour $n \ll 0$, les morphismes
$\mathcal{H}^\bullet \to \mathcal{E}^\bullet$ et
$\mathcal{H}^\bullet \to \mathcal{F}^\bullet$
se factorisent par le quasi-isomorphisme
$\mathcal{H}^\bullet \to \tau_{\geq n}\mathcal{H}^\bullet$
simplement parce que $\mathcal{E}^\bullet$ et $\mathcal{F}^\bullet$
sont des complexes bornés. On peut donc remplacer $\mathcal{H}^\bullet$
par $\tau_{\geq n}\mathcal{H}^\bullet$ et supposer que $\mathcal{H}^\bullet$
est borné inférieurement.
On peut alors appliquer le
lemme \ref{perfect-lemma-construct-strictly-perfect} pour conclure.
\end{proof}

\begin{lemma}
\label{perfect-lemma-resolution-property-homotopy-map-perfect-complex}
Soit $X$ un schéma quasi-compact et quasi-séparé possédant la
propriété de résolution. Soient $\mathcal{E}^\bullet$ et $\mathcal{F}^\bullet$
des complexes finis de $\mathcal{O}_X$-modules localement libres de type fini.
Soient $\alpha^\bullet, \beta^\bullet :\mathcal{E}^\bullet \to \mathcal{F}^\bullet$
deux morphismes de complexes définissant le même morphisme dans $D(\mathcal{O}_X)$.
Alors il existe un quasi-isomorphisme
$\gamma^\bullet : \mathcal{G}^\bullet \to \mathcal{E}^\bullet$
où $\mathcal{G}^\bullet$ est un complexe borné dont chaque terme est localement libre de type fini en tant que
$\mathcal{O}_X$-module,
tel que $\alpha^\bullet \circ \gamma^\bullet$ et
$\beta^\bullet \circ \gamma^\bullet$ soient des morphismes de complexes homotopes.
\end{lemma}

\begin{proof}
D'après le lemme \ref{perfect-lemma-resolution-property-affine-diagonal},
on voit que $X$ a une diagonale affine. Par conséquent, d'après la
proposition \ref{perfect-proposition-quasi-compact-affine-diagonal}
(et la définition de la catégorie dérivée),
il existe un quasi-isomorphisme
$\gamma^\bullet : \mathcal{G}^\bullet \to \mathcal{E}^\bullet$
où $\mathcal{G}^\bullet$ est un complexe de
$\mathcal{O}_X$-modules quasi-cohérents,
tel que $\alpha^\bullet \circ \gamma^\bullet$ et
$\beta^\bullet \circ \gamma^\bullet$ soient des morphismes de complexes homotopes.
Choisissons une homotopie $h^i : \mathcal{G}^i \to \mathcal{F}^{i - 1}$
attestant ce fait. Choisissons $n \ll 0$. Le morphisme
$\gamma^\bullet$ se factorise alors canoniquement par le morphisme quotient
$\mathcal{G}^\bullet \to \tau_{\geq n}\mathcal{G}^\bullet$
puisque $\mathcal{E}^\bullet$ est borné inférieurement. Pour exactement la même
raison, les morphismes $h^i$ se factorisent par les surjections
$\mathcal{G}^i \to (\tau_{\geq n}\mathcal{G})^i$.
On voit donc que l'on peut remplacer $\mathcal{G}^\bullet$
par $\tau_{\geq n}\mathcal{G}^\bullet$.
On peut alors appliquer le lemme \ref{perfect-lemma-construct-strictly-perfect} pour conclure.
\end{proof}

\begin{proposition}
\label{perfect-proposition-perfect-resolution-property}
Soit $X$ un schéma quasi-compact et quasi-séparé possédant la
propriété de résolution. Notons
\begin{enumerate}
\item $\mathcal{A}$ la catégorie additive des
$\mathcal{O}_X$-modules localement libres de type fini,
\item $K^b(\mathcal{A})$ la catégorie homotopique des complexes bornés
de $\mathcal{A}$ ; voir
Catégories dérivées, section \ref{derived-section-homotopy}, et
\item $D_{perf}(\mathcal{O}_X)$ la sous-catégorie strictement pleine, saturée
et triangulée de $D(\mathcal{O}_X)$ formée des
objets parfaits.
\end{enumerate}
Avec ces notations, le foncteur évident
$$
K^b(\mathcal{A}) \longrightarrow D_{perf}(\mathcal{O}_X)
$$
est un foncteur exact de catégories triangulées qui se factorise par une
équivalence $S^{-1}K^b(\mathcal{A}) \to D_{perf}(\mathcal{O}_X)$
de catégories triangulées,
où $S$ est le système multiplicatif saturé des quasi-isomorphismes
de $K^b(\mathcal{A})$.
\end{proposition}

\begin{proof}
Si vous parvenez à analyser l'énoncé de la proposition, sautez donc ce
premier paragraphe. Pour certaines des définitions employées, voir
Catégories dérivées, définition
\ref{derived-definition-triangulated-subcategory}
(sous-catégorie triangulée),
Catégories dérivées, définition \ref{derived-definition-saturated}
(sous-catégorie triangulée saturée),
Catégories dérivées, définition \ref{derived-definition-localization}
(système multiplicatif compatible avec la structure triangulée), et
Catégories,
définition \ref{categories-definition-saturated-multiplicative-system}
(système multiplicatif saturé).
Observons que $D_{perf}(\mathcal{O}_X)$ est une sous-catégorie triangulée saturée
de $D(\mathcal{O}_X)$ d'après
Cohomologie, lemmes \ref{cohomology-lemma-two-out-of-three-perfect} et
\ref{cohomology-lemma-summands-perfect}. Notons également que
$K^b(\mathcal{A})$ est une catégorie triangulée ; voir
Catégories dérivées, lemme
\ref{derived-lemma-bounded-triangulated-subcategories}.

\medskip\noindent
Il est clair que le foncteur transforme les triangles distingués en
triangles distingués, c'est-à-dire qu'il est exact. Alors $S$ est un système
multiplicatif saturé compatible avec la structure triangulée
de $K^b(\mathcal{A})$ d'après
Catégories dérivées, lemme \ref{derived-lemma-triangle-functor-localize}.
Par conséquent, la localisation $S^{-1}K^b(\mathcal{A})$ existe et est
une catégorie triangulée d'après
Catégories dérivées, proposition
\ref{derived-proposition-construct-localization}.
On obtient une factorisation par le foncteur exact
$S^{-1}K^b(\mathcal{A}) \to D_{perf}(\mathcal{O}_X)$ d'après
Catégories dérivées, lemme
\ref{derived-lemma-universal-property-localization}.
D'après les lemmes \ref{perfect-lemma-resolution-property-perfect-complex},
\ref{perfect-lemma-resolution-property-map-perfect-complex} et
\ref{perfect-lemma-resolution-property-homotopy-map-perfect-complex},
ce foncteur est une équivalence. Enfin, le foncteur
$S^{-1}K^b(\mathcal{A}) \to D_{perf}(\mathcal{O}_X)$
est une équivalence de catégories triangulées (au sens où
les triangles distingués se correspondent) d'après
Catégories dérivées, lemme \ref{derived-lemma-exact-equivalence}.
\end{proof}



\section{Groupes K}
\label{perfect-section-K-groups}

\noindent
Quelques mots au sujet de $K_0$ des diverses catégories associées aux schémas.
On trouvera des résultats antérieurs dans
Algèbre, section \ref{algebra-section-K-groups},
Homologie, section \ref{homology-section-K-groups},
Catégories dérivées, section \ref{derived-section-K-groups}, et
Compléments sur l'algèbre, lemme \ref{more-algebra-lemma-perfect-to-K-group-universal}.

\medskip\noindent
Par analogie avec Algèbre, section \ref{algebra-section-K-groups},
nous définirons deux groupes de $K$-théorie, $K'_0(X)$ et $K_0(X)$, pour tout schéma
noethérien $X$. Le premier fera intervenir les $\mathcal{O}_X$-modules cohérents
et le second les $\mathcal{O}_X$-modules localement libres de type fini.

\begin{lemma}
\label{perfect-lemma-Noetherian-Kprime}
Soit $X$ un schéma noethérien. Alors
$$
K_0(\textit{Coh}(\mathcal{O}_X)) =
K_0(D^b(\textit{Coh}(\mathcal{O}_X)) =
K_0(D^b_{\textit{Coh}}(\mathcal{O}_X))
$$
\end{lemma}

\begin{proof}
La première égalité est donnée par
Catégories dérivées, lemme \ref{derived-lemma-K-bounded-derived}.
On a $K_0(\textit{Coh}(\mathcal{O}_X)) =
K_0(D^b_{\textit{Coh}}(\mathcal{O}_X))$ d'après
Catégories dérivées, lemme \ref{derived-lemma-DBA-map-K}.
Ceci démontre le lemme.
(On peut aussi utiliser l'égalité
$D^b(\textit{Coh}(\mathcal{O}_X)) = D^b_{\textit{Coh}}(\mathcal{O}_X)$
de la proposition \ref{perfect-proposition-DCoh} pour obtenir la seconde
égalité.)
\end{proof}

\noindent
Voici la définition.

\begin{definition}
\label{perfect-definition-K-group}
Soit $X$ un schéma.
\begin{enumerate}
\item On note $K_0(X)$ le {\it groupe de Grothendieck de $X$}. C'est le
groupe de K-théorie d'indice zéro de la sous-catégorie triangulée
strictement pleine et saturée $D_{perf}(\mathcal{O}_X)$ de $D(\mathcal{O}_X)$
formée des objets parfaits. Autrement dit,
$$
K_0(X) = K_0(D_{perf}(\mathcal{O}_X))
$$
\item Si $X$ est localement noethérien, on note $K'_0(X)$ le
{\it groupe de Grothendieck des faisceaux cohérents sur $X$}. C'est le
groupe $K$ d'indice zéro de la catégorie abélienne
des $\mathcal{O}_X$-modules cohérents. Autrement dit,
$$
K'_0(X) = K_0(\textit{Coh}(\mathcal{O}_X))
$$
\end{enumerate}
\end{definition}

\noindent
Nous montrerons que notre définition de $K_0(X)$ coïncide avec la définition
couramment utilisée en termes de modules localement libres de type fini si $X$
possède la propriété de résolution (par exemple si $X$ possède un module
inversible ample). Voir le lemme \ref{perfect-lemma-K-is-old-K}.

\begin{lemma}
\label{perfect-lemma-K-agrees-affine}
Soit $X = \Spec(R)$ un schéma affine. Alors $K_0(X) = K_0(R)$
et, si $R$ est noethérien, $K'_0(X) = K'_0(R)$.
\end{lemma}

\begin{proof}
Rappelons que $K'_0(R)$ et $K_0(R)$ ont été définis dans
Algèbre, section \ref{algebra-section-K-groups}.

\medskip\noindent
D'après Compléments sur l'algèbre, lemme \ref{more-algebra-lemma-perfect-to-K-group-universal},
on a $K_0(R) = K_0(D_{perf}(R))$.
D'après les lemmes \ref{perfect-lemma-perfect-affine} et \ref{perfect-lemma-affine-compare-bounded},
on a $D_{perf}(R) = D_{perf}(\mathcal{O}_X)$.
Ceci démontre l'égalité $K_0(R) = K_0(X)$.

\medskip\noindent
L'égalité $K'_0(R) = K'_0(X)$ résulte du fait que
$\textit{Coh}(\mathcal{O}_X)$ est équivalente à la catégorie
des $R$-modules finis d'après Cohomologie des schémas, lemme
\ref{coherent-lemma-coherent-Noetherian}. De plus, les définitions
montrent immédiatement que $K'_0(R)$ est le groupe de K-théorie d'indice zéro de la catégorie
des $R$-modules finis.
\end{proof}

\noindent
Soit $X$ un schéma noethérien. Alors $K'_0(X)$ et $K_0(X)$
sont tous deux définis. Dans ce cas, il existe une application canonique
$$
K_0(X) = K_0(D_{perf}(\mathcal{O}_X))
\longrightarrow
K_0(D^b_{\textit{Coh}}(\mathcal{O}_X)) = K'_0(X)
$$
En effet, les complexes parfaits appartiennent à $D^b_{\textit{Coh}}(\mathcal{O}_X)$
(d'après le lemme \ref{perfect-lemma-identify-pseudo-coherent-noetherian}), le foncteur
d'inclusion
$D_{perf}(\mathcal{O}_X) \to D^b_{\textit{Coh}}(\mathcal{O}_X)$
induit un homomorphisme entre les groupes $K$ d'indice zéro
(Catégories dérivées, lemme \ref{derived-lemma-map-K}),
et l'égalité de droite est donnée par le
lemme \ref{perfect-lemma-Noetherian-Kprime}.

\begin{lemma}
\label{perfect-lemma-Kprime-K}
Soit $X$ un schéma noethérien régulier. Alors
l'application $K_0(X) \to K'_0(X)$ est un isomorphisme.
\end{lemma}

\begin{proof}
Cela résulte immédiatement du lemme \ref{perfect-lemma-perfect-on-regular}
et de la construction ci-dessus de l'application $K_0(X) \to K'_0(X)$.
\end{proof}

\noindent
Soit $X$ un schéma. Notons $\textit{Vect}(X)$ la catégorie
des $\mathcal{O}_X$-modules localement libres de type fini. Bien que
$\textit{Vect}(X)$ ne soit pas en général une catégorie abélienne, la notion
de suite exacte courte dans $\textit{Vect}(X)$ est claire.
Notons $K_0(\textit{Vect}(X))$ l'unique groupe abélien ayant les
propriétés suivantes\footnote{Le cadre général correct serait ici
de définir $K_0$ pour toute catégorie exacte ; voir
Injectifs, remarque \ref{injectives-remark-embed-exact-category}.} :
\begin{enumerate}
\item Pour tout $\mathcal{O}_X$-module localement libre de type fini $\mathcal{E}$,
un élément $[\mathcal{E}]$ de $K_0(\textit{Vect}(X))$ est donné,
\item pour toute suite exacte courte
$0 \to \mathcal{E}' \to \mathcal{E} \to \mathcal{E}'' \to 0$
de $\mathcal{O}_X$-modules localement libres de type fini, on a la relation
$[\mathcal{E}] = [\mathcal{E}'] + [\mathcal{E}'']$ dans $K_0(\textit{Vect}(X))$,
\item le groupe $K_0(\textit{Vect}(X))$ est engendré par les éléments
$[\mathcal{E}]$, et
\item toutes les relations dans $K_0(\textit{Vect}(X))$ entre les générateurs
$[\mathcal{E}]$ sont des combinaisons $\mathbf{Z}$-linéaires
des relations provenant des suites exactes comme ci-dessus.
\end{enumerate}
Nous omettons la construction détaillée de $K_0(\textit{Vect}(X))$.
Il existe une application naturelle
$$
K_0(\textit{Vect}(X)) \longrightarrow K_0(X)
$$
En effet, étant donné un $\mathcal{O}_X$-module localement libre de type fini $\mathcal{E}$,
notons $\mathcal{E}[0]$ le complexe parfait sur $X$ dans lequel
$\mathcal{E}$ est placé en degré $0$ et dont les autres termes sont nuls.
Étant donnée une suite exacte courte
$0 \to \mathcal{E} \to \mathcal{E}' \to \mathcal{E}'' \to 0$ de
$\mathcal{O}_X$-modules localement libres de type fini, on obtient un triangle distingué
$\mathcal{E}[0] \to \mathcal{E}'[0] \to \mathcal{E}''[0] \to \mathcal{E}[1]$ ;
voir Catégories dérivées, section \ref{derived-section-canonical-delta-functor}.
Cela montre que l'on obtient une application
$K_0(\textit{Vect}(X)) \to K_0(D_{perf}(\mathcal{O}_X)) = K_0(X)$
en envoyant $[\mathcal{E}]$ sur $[\mathcal{E}[0]]$,
en priant le lecteur d'excuser cette notation épouvantable.

\begin{lemma}
\label{perfect-lemma-K-is-old-K}
Soit $X$ un schéma quasi-compact et quasi-séparé possédant la
propriété de résolution. Alors l'application $K_0(\textit{Vect}(X)) \to K_0(X)$
est un isomorphisme.
\end{lemma}

\begin{proof}
Ce lemme résulte directement des
lemmes \ref{perfect-lemma-resolution-property-perfect-complex},
\ref{perfect-lemma-resolution-property-map-perfect-complex} et
\ref{perfect-lemma-resolution-property-homotopy-map-perfect-complex},
dont nous utiliserons les résultats sans autre mention.
Construisons une application réciproque
$$
c : K_0(X) = K_0(D_{perf}(\mathcal{O}_X)) \longrightarrow
K_0(\textit{Vect}(X))
$$
Tout objet de $D_{perf}(\mathcal{O}_X)$ peut être représenté
par un complexe borné $\mathcal{E}^\bullet$ de
$\mathcal{O}_X$-modules localement libres de type fini. On pose alors
$$
c([\mathcal{E}^\bullet]) = \sum (-1)^i[\mathcal{E}^i]
$$
Bien entendu, il faut montrer que cette définition est indépendante du choix.
Pour l'instant, nous considérons $c$ comme une application définie sur les complexes
bornés de $\mathcal{O}_X$-modules localement libres de type fini.

\medskip\noindent
Supposons que $\mathcal{E}^\bullet \to \mathcal{F}^\bullet$ soit une application surjective
de complexes bornés de $\mathcal{O}_X$-modules localement libres de type fini.
Soit $\mathcal{K}^\bullet$ son noyau. On obtient alors des suites exactes
courtes de $\mathcal{O}_X$-modules
$$
0 \to \mathcal{K}^n \to \mathcal{E}^n \to \mathcal{F}^n \to 0
$$
qui sont scindées localement parce que $\mathcal{F}^n$ est localement libre de type fini.
Ainsi, $\mathcal{K}^\bullet$ est lui aussi un complexe borné de
$\mathcal{O}_X$-modules localement libres de type fini, et l'on a
$c(\mathcal{E}^\bullet) = c(\mathcal{K}^\bullet) + c(\mathcal{F}^\bullet)$
dans $K_0(\textit{Vect}(X))$.

\medskip\noindent
Supposons donné un complexe borné $\mathcal{E}^\bullet$
de $\mathcal{O}_X$-modules localement libres de type fini qui soit acyclique.
Supposons que $\mathcal{E}^n = 0$ pour $n \not \in [a, b]$. On peut alors
décomposer $\mathcal{E}^\bullet$ en suites exactes courtes
$$
\begin{matrix}
0 \to \mathcal{E}^a \to \mathcal{E}^{a + 1} \to \mathcal{F}^{a + 1} \to 0,\\
0 \to \mathcal{F}^{a + 1} \to \mathcal{E}^{a + 2} \to
\mathcal{F}^{a + 3} \to 0, \\
\ldots \\
0 \to \mathcal{F}^{b - 3} \to \mathcal{E}^{b - 2} \to
\mathcal{F}^{b - 2} \to 0, \\
0 \to \mathcal{F}^{b - 2} \to \mathcal{E}^{b - 1} \to \mathcal{E}^b \to 0
\end{matrix}
$$
En raisonnant par récurrence descendante, on voit que
$\mathcal{F}^{b - 2}, \ldots, \mathcal{F}^{a + 1}$
sont des $\mathcal{O}_X$-modules localement libres de type fini et que
$$
c(\mathcal{E}^\bullet) = \sum (-1)[\mathcal{E}^n] =
\sum (-1)^n([\mathcal{F}^{n - 1}] + [\mathcal{F}^n]) = 0
$$
Ainsi, notre construction donne zéro sur les complexes acycliques.

\medskip\noindent
Il résulte des deux paragraphes précédents que $c$
est bien définie. En effet, supposons que les complexes bornés
$\mathcal{E}^\bullet$ et $\mathcal{F}^\bullet$ de $\mathcal{O}_X$-modules localement libres
de type fini représentent le même objet de $D(\mathcal{O}_X)$.
On peut alors trouver des quasi-isomorphismes
$a : \mathcal{G}^\bullet \to \mathcal{E}^\bullet$ et
$b : \mathcal{G}^\bullet \to \mathcal{F}^\bullet$,
où $\mathcal{G}^\bullet$ est un complexe borné de $\mathcal{O}_X$-modules localement libres
de type fini.
On obtient une suite exacte courte de complexes
$$
0 \to \mathcal{E}^\bullet \to C(a)^\bullet \to \mathcal{G}^\bullet[1] \to 0
$$
voir Catégories dérivées, définition \ref{derived-definition-cone}.
Puisque $a$ est un quasi-isomorphisme, le cône $C(a)^\bullet$ est
acyclique (cela résulte par exemple de la discussion dans
Catégories dérivées, section \ref{derived-section-canonical-delta-functor}).
Par conséquent,
$$
0 = c(C(f)^\bullet) = c(\mathcal{E}^\bullet) + c(\mathcal{G}^\bullet[1]) =
c(\mathcal{E}^\bullet) - c(\mathcal{G}^\bullet)
$$
comme voulu. Le même raisonnement appliqué à $b$ montre que
$0 = c(\mathcal{F}^\bullet) - c(\mathcal{G}^\bullet)$.
On obtient donc $c(\mathcal{E}^\bullet) = c(\mathcal{F}^\bullet)$,
et $c$ est bien définie.

\medskip\noindent
Un raisonnement analogue utilisant le cône d'une application
$\mathcal{E}^\bullet \to \mathcal{F}^\bullet$
de complexes bornés de $\mathcal{O}_X$-modules localement libres de type fini
montre que $c(Y) = c(X) + c(Z)$ si $X \to Y \to Z$ est un triangle distingué
dans $D_{perf}(\mathcal{O}_X)$. Les détails sont omis.
On obtient ainsi l'homomorphisme souhaité
de groupes abéliens $c : K_0(X) \to K_0(\textit{Vect}(X))$.

\medskip\noindent
Il est clair que la composée
$K_0(\textit{Vect}(X)) \to K_0(X) \to K_0(\textit{Vect}(X))$
est l'identité. D'autre part, soit $\mathcal{E}^\bullet$
un complexe borné de $\mathcal{O}_X$-modules localement libres de type fini.
L'existence des triangles distingués associés aux « troncatures bêtes »
(voir Homologie, section \ref{homology-section-truncations})
$$
\sigma_{\geq n}\mathcal{E}^\bullet \to
\sigma_{\geq n - 1}\mathcal{E}^\bullet \to
\mathcal{E}^{n - 1}[-n + 1] \to
(\sigma_{\geq n}\mathcal{E}^\bullet)[1]
$$
et une récurrence montrent que
$$
[\mathcal{E}^\bullet] = \sum (-1)^i[\mathcal{E}^i[0]]
$$
dans $K_0(X) = K_0(D_{perf}(\mathcal{O}_X))$, en priant le lecteur d'excuser la notation.
Ainsi, l'application $K_0(\textit{Vect}(X)) \to K_0(D_{perf}(\mathcal{O}_X)) = K_0(X)$
est surjective, ce qui achève la preuve.
\end{proof}

\begin{remark}
\label{perfect-remark-K-ring}
Soit $X$ un schéma. Le groupe de K-théorie $K_0(X)$ est canoniquement un anneau commutatif.
En effet, le produit tensoriel dérivé
$$
\otimes = \otimes^\mathbf{L}_{\mathcal{O}_X} :
D_{perf}(\mathcal{O}_X) \times D_{perf}(\mathcal{O}_X)
\longrightarrow
D_{perf}(\mathcal{O}_X)
$$
et Catégories dérivées, lemme \ref{derived-lemma-bilinear-map-K},
donnent une multiplication bilinéaire. Comme $K \otimes L \cong L \otimes K$,
ce produit est commutatif. Comme
$(K \otimes L) \otimes M = K \otimes (L \otimes M)$,
ce produit est associatif.
Enfin, l'unité de $K_0(X)$ est l'élément $1 = [\mathcal{O}_X]$.

\medskip\noindent
Si $\textit{Vect}(X)$ et $K_0(\textit{Vect}(X))$ sont comme ci-dessus, alors
il est clair que $K_0(\textit{Vect}(X))$ possède lui aussi une
structure d'anneau : si $\mathcal{E}$ et $\mathcal{F}$ sont des
$\mathcal{O}_X$-modules localement libres de type fini, on pose
$$
[\mathcal{E}] \cdot [\mathcal{F}] =
[\mathcal{E} \otimes_{\mathcal{O}_X} \mathcal{F}]
$$
Le lecteur vérifie aisément que cela définit bien un produit bilinéaire,
commutatif et associatif. Les détails sont omis. L'application
$$
K_0(\textit{Vect}(X)) \longrightarrow K_0(X)
$$
construite ci-dessus est un homomorphisme d'anneaux pour ces définitions.

\medskip\noindent
Supposons maintenant $X$ noethérien. Le produit tensoriel dérivé fournit aussi
une application
$$
\otimes = \otimes^\mathbf{L}_{\mathcal{O}_X} :
D_{perf}(\mathcal{O}_X) \times D^b_{\textit{Coh}}(\mathcal{O}_X)
\longrightarrow
D^b_{\textit{Coh}}(\mathcal{O}_X)
$$
En utilisant de nouveau Catégories dérivées, lemme \ref{derived-lemma-bilinear-map-K},
on obtient une multiplication bilinéaire $K_0(X) \times K'_0(X) \to K'_0(X)$,
car $K'_0(X) = K_0(D^b_{\textit{Coh}}(\mathcal{O}_X))$ d'après le
lemme \ref{perfect-lemma-Noetherian-Kprime}.
Le lecteur montre aisément que cela munit $K'_0(X)$ d'une structure
de module sur l'anneau $K_0(X)$.
\end{remark}

\begin{remark}
\label{perfect-remark-pushforward-K}
Soit $f : X \to Y$ un morphisme propre de schémas localement noethériens.
Il existe une application
$$
f_* : K'_0(X) \longrightarrow K'_0(Y)
$$
qui envoie $[\mathcal{F}]$ sur
$$
[\bigoplus\nolimits_{i \geq 0} R^{2i}f_*\mathcal{F}] -
[\bigoplus\nolimits_{i \geq 0} R^{2i + 1}f_*\mathcal{F}]
$$
Cette application est bien définie parce que les faisceaux $R^if_*\mathcal{F}$
sont cohérents (Cohomologie des schémas, proposition
\ref{coherent-proposition-proper-pushforward-coherent}), parce que localement
seul un nombre fini d'entre eux sont non nuls, et parce qu'une
suite exacte courte de faisceaux cohérents sur $X$ donne une suite exacte
longue de $R^if_*$ sur $Y$. Si $Y$ est quasi-compact (le seul
cas généralement utilisé en pratique), on peut réécrire ce qui précède sous la forme
$$
f_*[\mathcal{F}] = \sum (-1)^i[R^if_*\mathcal{F}] = [Rf_*\mathcal{F}]
$$
où l'on a utilisé l'égalité $K'_0(Y) = K_0(D^b_{\textit{Coh}}(Y))$ du
lemme \ref{perfect-lemma-Noetherian-Kprime}.
\end{remark}

\begin{lemma}
\label{perfect-lemma-projection-formula}
Soit $f : X \to Y$ un morphisme propre de schémas localement noethériens.
Alors $f_*(\alpha \cdot f^*\beta) = f_*\alpha \cdot \beta$
pour $\alpha \in K'_0(X)$ et $\beta \in K_0(Y)$.
\end{lemma}

\begin{proof}
Cela résulte du lemme \ref{perfect-lemma-cohomology-base-change}, de la discussion de la
remarque \ref{perfect-remark-pushforward-K} et de la définition du produit
$K'_0(X) \times K_0(X) \to K'_0(X)$ dans la remarque \ref{perfect-remark-K-ring}.
\end{proof}

\begin{remark}
\label{perfect-remark-perf-Z}
Soit $X$ un schéma. Soit $Z \subset X$ un sous-schéma fermé. Considérons la
sous-catégorie triangulée strictement pleine et saturée
$$
D_{Z, perf}(\mathcal{O}_X) \subset D(\mathcal{O}_X)
$$
formée des complexes parfaits de $\mathcal{O}_X$-modules
dont les faisceaux de cohomologie sont à support ensembliste dans $Z$.
Le groupe $K$ d'indice zéro $K_0(D_{Z, perf}(\mathcal{O}_X))$
de cette catégorie triangulée est parfois noté
$K_Z(X)$ ou $K_{0, Z}(X)$. En utilisant le produit tensoriel dérivé exactement
comme dans la remarque \ref{perfect-remark-K-ring}, on voit que $K_0(D_{Z, perf}(\mathcal{O}_X))$
possède une multiplication associative et commutative,
mais qu'en général $K_0(D_{Z, perf}(\mathcal{O}_X))$ n'a pas d'unité.
\end{remark}










\section{Déterminants des complexes}
\label{perfect-section-det-two-terms}

\noindent
Cette section fait suite à
Compléments sur l'algèbre, section \ref{more-algebra-section-determinants-complexes}.
Pour tout espace annelé $(X, \mathcal{O}_X)$, il existe un foncteur
$$
\det :
\left\{
\begin{matrix}
\text{catégorie des complexes parfaits} \\
\text{les morphismes sont les isomorphismes}
\end{matrix}
\right\}
\longrightarrow
\left\{
\begin{matrix}
\text{catégorie des modules inversibles} \\
\text{les morphismes sont les isomorphismes}
\end{matrix}
\right\}
$$
De plus, si $(L, F)$ est un objet de la catégorie dérivée filtrée
$DF(\mathcal{O}_X)$ dont la filtration est finie et dont les gradués
sont des complexes parfaits, il existe un isomorphisme canonique
$\det(\text{gr}L) \to \det(L)$. Voir \cite{determinant} pour l'exposé
original. Nous ajouterons ces développements ultérieurement (insérer une référence future).

\medskip\noindent
Pour le moment, nous présentons une construction ad hoc dans le cas
où $X$ est un schéma et où nous considérons des objets parfaits $L$ de
$D(\mathcal{O}_X)$ d'amplitude de Tor contenue dans $[-1, 0]$.

\begin{lemma}
\label{perfect-lemma-determinant-two-term-complexes}
Soit $X$ un schéma. Il existe un foncteur
$$
\det :
\left\{
\begin{matrix}
\text{catégorie des complexes parfaits} \\
\text{d'amplitude de Tor contenue dans }[-1, 0] \\
\text{les morphismes sont les isomorphismes}
\end{matrix}
\right\}
\longrightarrow
\left\{
\begin{matrix}
\text{catégorie des modules inversibles} \\
\text{les morphismes sont les isomorphismes}
\end{matrix}
\right\}
$$
De plus, si l'on se donne un objet parfait de rang $0$, noté $L$, de $D(\mathcal{O}_X)$,
d'amplitude de Tor contenue dans $[-1, 0]$, il existe un élément canonique
$\delta(L) \in \Gamma(X, \det(L))$ tel que, pour tout isomorphisme
$a : L \to K$ dans $D(\mathcal{O}_X)$, on ait $\det(a)(\delta(L)) = \delta(K)$.
De plus, la construction est donnée localement sur les ouverts affines par celle
de Compléments sur l'algèbre, section \ref{more-algebra-section-determinants-complexes}.
\end{lemma}

\begin{proof}
Soit $L$ un objet de la catégorie figurant à gauche. Si $\Spec(A) = U \subset X$
est un ouvert affine, alors $L|_U$ correspond à un complexe parfait $L^\bullet$
de $A$-modules d'amplitude de Tor contenue dans $[-1, 0]$, voir les
lemmes \ref{perfect-lemma-affine-compare-bounded},
\ref{perfect-lemma-tor-dimension-affine} et
\ref{perfect-lemma-perfect-affine}.
On peut alors considérer le $A$-module inversible $\det(L^\bullet)$ construit dans
Compléments sur l'algèbre, lemme \ref{more-algebra-lemma-determinant-two-term-complexes}.
Si $\Spec(B) = V \subset U$ est un autre ouvert affine contenu dans $U$,
alors $\det(L^\bullet) \otimes_A B = \det(L^\bullet \otimes_A B)$,
et cette construction est donc compatible avec les applications de restriction
(voir le lemme \ref{perfect-lemma-quasi-coherence-pullback} et noter que $A \to B$ est plat).
Ainsi, nous pouvons recoller ces modules inversibles pour obtenir un module inversible
$\det(L)$ sur $X$. La fonctorialité et les sections canoniques
se construisent exactement de la même manière. Détails omis.
\end{proof}

\begin{remark}
\label{perfect-remark-functorial-det}
La construction du lemme \ref{perfect-lemma-determinant-two-term-complexes}
est compatible avec les images inverses. Plus précisément, soient un morphisme
$f : X \to Y$ de schémas et un objet parfait $K$ de $D(\mathcal{O}_Y)$
d'amplitude de Tor contenue dans $[-1, 0]$. Alors $Lf^*K$ est un
objet parfait $K$ de $D(\mathcal{O}_X)$
d'amplitude de Tor contenue dans $[-1, 0]$ et nous avons une identification canonique
$$
f^*\det(K) \longrightarrow \det(Lf^*K)
$$
De plus, si $K$ est de rang $0$, alors l'image réciproque de $\delta(K)$ s'identifie à
$\delta(Lf^*K)$ par ce morphisme. Cela résulte immédiatement de la construction
du déterminant sur les ouverts affines.
\end{remark}




\section{Détection de la bornitude}
\label{perfect-section-detecting-boundedness}

\noindent
Dans cette section, nous montrons que les générateurs compacts de $D_\QCoh$ d'un
schéma quasi-compact et quasi-séparé, tels qu'ils sont construits dans la
section \ref{perfect-section-generating}, possèdent une propriété particulière.
Nous recommandons de lire d'abord cette dernière section, car elle est très semblable à celle-ci.

\begin{lemma}
\label{perfect-lemma-orthogonal-koszul-first-variant}
Dans la situation \ref{perfect-situation-complex}, notons $j : U \to X$ l'immersion
ouverte et soit $K$ l'objet parfait de $D(\mathcal{O}_X)$
correspondant au complexe de Koszul de $f_1, \ldots, f_r$ sur $A$.
Soient $E \in D_\QCoh(\mathcal{O}_X)$ et $a \in \mathbf{Z}$.
Considérons les conditions suivantes :
\begin{enumerate}
\item Le morphisme canonique $\tau_{\geq a}E \to \tau_{\geq a} Rj_*(E|_U)$
est un isomorphisme.
\item On a $\Hom_{D(\mathcal{O}_X)}(K[-n], E) = 0$ pour tout $n \geq a$.
\end{enumerate}
Alors (2) implique (1), et (1) implique (2) lorsque $a$ est remplacé par $a + 1$.
\end{lemma}

\begin{proof}
Choisissons un triangle distingué $N \to E \to Rj_*(E|_U) \to N[1]$.
Alors (1) implique $\tau_{\geq a + 1} N = 0$, tandis que
$\tau_{\geq a}N = 0$ implique (1). Observons que
$$
\Hom_{D(\mathcal{O}_X)}(K[-n], Rj_*(E|_U)) =
\Hom_{D(\mathcal{O}_U)}(K|_U[-n], E) = 0
$$
pour tout $n$, puisque $K|_U = 0$. Ainsi (2) équivaut à
$\Hom_{D(\mathcal{O}_X)}(K[-n], N) = 0$ pour tout $n \geq a$. 
Observons qu'il existe des triangles distingués
$$
K^\bullet(f_1^{e_1}, \ldots, f_i^{e'_i}, \ldots, f_r^{e_r}) \to
K^\bullet(f_1^{e_1}, \ldots, f_i^{e'_i + e''_i}, \ldots, f_r^{e_r}) \to
K^\bullet(f_1^{e_1}, \ldots, f_i^{e''_i}, \ldots, f_r^{e_r}) \to \ldots
$$
de complexes de Koszul, voir
Compléments sur l'algèbre, lemme \ref{more-algebra-lemma-koszul-mult}. Par conséquent,
$\Hom_{D(\mathcal{O}_X)}(K[-n], N) = 0$ pour tout $n \geq a$
équivaut à
$\Hom_{D(\mathcal{O}_X)}(K_e[-n], N) = 0$ pour tout $n \geq a$ et
tout $e \geq 1$, avec $K_e$ comme dans le
lemme \ref{perfect-lemma-represent-cohomology-class-on-closed}.
Puisque $N|_U = 0$, ce lemme implique que cette condition équivaut à son tour à
$H^n(X, N) = 0$ pour $n \geq a$. Nous en concluons que (2) équivaut
à $\tau_{\geq a}N = 0$, puisque $N$ est déterminé par le complexe de
$A$-modules $R\Gamma(X, N)$, voir le lemme \ref{perfect-lemma-affine-compare-bounded}.
Cela démontre le lemme.
\end{proof}

\begin{lemma}
\label{perfect-lemma-orthogonal-koszul-second-variant}
Dans la situation \ref{perfect-situation-complex}, notons $j : U \to X$ l'immersion
ouverte et soit $K$ l'objet parfait de $D(\mathcal{O}_X)$
correspondant au complexe de Koszul de $f_1, \ldots, f_r$ sur $A$.
Soient $E \in D_\QCoh(\mathcal{O}_X)$ et $a \in \mathbf{Z}$. Considérons
les conditions suivantes :
\begin{enumerate}
\item Le morphisme canonique $\tau_{\leq a}E \to \tau_{\leq a} Rj_*(E|_U)$
est un isomorphisme, et
\item $\Hom_{D(\mathcal{O}_X)}(K[-n], E) = 0$ pour tout $n \leq a$.
\end{enumerate}
Alors (2) implique (1), et (1) implique (2) lorsque $a$ est remplacé par $a - 1$.
\end{lemma}

\begin{proof}
Choisissons un triangle distingué $E \to Rj_*(E|_U) \to N \to E[1]$. Alors (1)
implique $\tau_{\leq a - 1}N = 0$, tandis que $\tau_{\leq a}N = 0$ implique (1).
Observons que
$$
\Hom_{D(\mathcal{O}_X)}(K[-n], Rj_*(E|_U)) =
\Hom_{D(\mathcal{O}_U)}(K|_U[-n], E) = 0
$$
pour tout $n$, puisque $K|_U = 0$. Ainsi (2) équivaut à
$\Hom_{D(\mathcal{O}_X)}(K[-n], N) = 0$ pour tout $n \leq a$. 
Observons qu'il existe des triangles distingués
$$
K^\bullet(f_1^{e_1}, \ldots, f_i^{e'_i}, \ldots, f_r^{e_r}) \to
K^\bullet(f_1^{e_1}, \ldots, f_i^{e'_i + e''_i}, \ldots, f_r^{e_r}) \to
K^\bullet(f_1^{e_1}, \ldots, f_i^{e''_i}, \ldots, f_r^{e_r}) \to \ldots
$$
de complexes de Koszul, voir
Compléments sur l'algèbre, lemme \ref{more-algebra-lemma-koszul-mult}. Par conséquent,
$\Hom_{D(\mathcal{O}_X)}(K[-n], N) = 0$ pour tout $n \leq a$
équivaut à
$\Hom_{D(\mathcal{O}_X)}(K_e[-n], N) = 0$ pour tout $n \leq a$ et tout $e \geq 1$,
avec $K_e$ comme dans le lemme \ref{perfect-lemma-represent-cohomology-class-on-closed}.
Puisque $N|_U = 0$, ce lemme implique que cette condition équivaut à son tour à
$H^n(X, N) = 0$ pour $n \leq a$. Nous en concluons que (2) équivaut à
$\tau_{\leq a}N = 0$, puisque $N$ est déterminé par le complexe de
$A$-modules $R\Gamma(X, N)$, voir le lemme \ref{perfect-lemma-affine-compare-bounded}.
Cela démontre le lemme.
\end{proof}

\begin{lemma}
\label{perfect-lemma-bounded-truncation}
Soit $X$ un schéma quasi-compact et quasi-séparé.
Soient $P \in D_{perf}(\mathcal{O}_X)$ et $E \in D_{\QCoh}(\mathcal{O}_X)$. 
Soit $a \in \mathbf{Z}$. Les conditions suivantes sont équivalentes :
\begin{enumerate}
\item $\Hom_{D(\mathcal{O}_X)}(P[-i], E) = 0$ pour $i \gg 0$, et
\item $\Hom_{D(\mathcal{O}_X)}(P[-i], \tau_{\geq a} E) = 0$ pour $i \gg 0$.
\end{enumerate}
\end{lemma}

\begin{proof}
En utilisant le triangle $ \tau_{< a} E \to E \to \tau_{\geq a} E \to$,
on voit qu'il suffit, pour obtenir l'équivalence, de montrer que
$$
\Hom_{D(\mathcal{O}_X)}(P[-i], \tau_{< a} E) =
\Hom_{D(\mathcal{O}_X)}(P, (\tau_{< a} E)[i]) = 0 
$$
pour $i \gg 0$. Comme $P$ est parfait, cela résulte du
lemme \ref{perfect-lemma-ext-from-perfect-into-bounded-QCoh}.
\end{proof}

\begin{lemma}
\label{perfect-lemma-bounded-below-truncation}
Soit $X$ un schéma quasi-compact et quasi-séparé. Soient
$P \in D_{perf}(\mathcal{O}_X)$ et $E \in D_{\QCoh}(\mathcal{O}_X)$.
Soit $a \in \mathbf{Z}$. Les conditions suivantes sont équivalentes :
\begin{enumerate}
\item $\Hom_{D(\mathcal{O}_X)}(P[-i], E) = 0$ pour $i \ll 0$, et
\item $\Hom_{D(\mathcal{O}_X)}(P[-i], \tau_{\leq a} E) = 0$ pour $i \ll 0$.
\end{enumerate}
\end{lemma}

\begin{proof}
En utilisant le triangle $ \tau_{\leq a} E \to E \to \tau_{> a} E \to$,
on voit qu'il suffit, pour obtenir l'équivalence, de montrer que
$$
\Hom_{D(\mathcal{O}_X)}(P[-i], \tau_{> a} E) =
\Hom_{D(\mathcal{O}_X)}(P, (\tau_{> a} E)[i]) = 0
$$
pour $i \ll 0$. Comme $P$ est parfait, cela résulte du
lemme \ref{perfect-lemma-ext-from-perfect-into-bounded-QCoh}.
\end{proof}

\begin{proposition}
\label{perfect-proposition-detecting-bounded-above}
Soit $X$ un schéma quasi-compact et quasi-séparé. Soit
$G \in D_{perf}(\mathcal{O}_X)$ un complexe parfait qui engendre 
$D_\QCoh (\mathcal{O}_X)$. Soit $E \in D_\QCoh (\mathcal{O}_X)$.
Les conditions suivantes sont équivalentes :
\begin{enumerate}
\item $E \in D^-_\QCoh (\mathcal{O}_X)$,
\item $\Hom_{D(\mathcal{O}_X)}(G[-i], E) = 0$ pour $i \gg 0$,
\item $\Ext^i_X(G, E) = 0$ pour $i \gg 0$,
\item $R\Hom_X(G, E)$ appartient à $D^-(\mathbf{Z})$,
\item $H^i(X, G^\vee \otimes_{\mathcal{O}_X}^\mathbf{L} E) = 0$
pour $i \gg 0$,
\item $R\Gamma(X, G^\vee \otimes_{\mathcal{O}_X}^\mathbf{L} E)$
appartient à $D^-(\mathbf{Z})$,
\item pour tout objet parfait $P$ de $D(\mathcal{O}_X)$,
\begin{enumerate}
\item les assertions (2), (3), (4) sont vraies lorsque $G$ est remplacé par $P$, et
\item $H^i(X, P \otimes_{\mathcal{O}_X}^\mathbf{L} E) = 0$ pour $i \gg 0$,
\item $R\Gamma(X, P \otimes_{\mathcal{O}_X}^\mathbf{L} E)$
appartient à $D^-(\mathbf{Z})$.
\end{enumerate}
\end{enumerate}
\end{proposition}

\begin{proof}
Supposons (1).\newline Puisque
$\Hom_{D(\mathcal{O}_X)}(G[-i], E) = \Hom_{D(\mathcal{O}_X)}(G, E[i])$,
ce groupe est nul pour $i \gg 0$ d'après le
lemme \ref{perfect-lemma-ext-from-perfect-into-bounded-QCoh}. Ceci démontre
que (1) implique (2).

\medskip\noindent
Les parties (2), (3), (4) sont équivalentes d'après la discussion de
Cohomologie, section \ref{cohomology-section-global-RHom}.
Les parties (5) et (6) sont équivalentes, puisque $H^i(X, -) = H^i(R\Gamma(X, -))$
par définition. Les conditions équivalentes (2), (3), (4) sont
équivalentes aux conditions équivalentes (5), (6) d'après
Cohomologie, lemme \ref{cohomology-lemma-dual-perfect-complex},
et le fait que $(G[-i])^\vee = G^\vee[i]$.

\medskip\noindent
Il est clair que (7) implique (2). Réciproquement, 
démontrons que les conditions équivalentes (2) -- (6) impliquent (7).
Rappelons que $G$ est un générateur classique de $D_{perf}(\mathcal{O}_X)$ d'après la
remarque \ref{perfect-remark-classical-generator}.
Pour $P \in D_{perf}(\mathcal{O}_X)$, notons $T(P)$ l'assertion selon laquelle
$R\Hom_X(P, E)$ appartient à $D^-(\mathbf{Z})$.
Il est clair que $T$ est stable par sommes directes,
vérifie la propriété des deux sur trois pour les triangles
distingués, et est stable par facteurs directs et par décalages.
Par conséquent, d'après Catégories dérivées, remarque \ref{derived-remark-check-on-generator},
on voit que (4) implique que $T$ vaut sur tout $D_{perf}(\mathcal{O}_X)$.
Le même argument vaut pour toutes les autres propriétés, si ce n'est que, pour les propriétés
(7)(b) et (7)(c), on utilise aussi que $P \mapsto P^\vee$ est une
autoéquivalence de $D_{perf}(\mathcal{O}_X)$. Un petit détail est omis.

\medskip\noindent
Nous allons démontrer que les conditions équivalentes (2) -- (7) impliquent (1)
en utilisant le principe de récurrence de
Cohomologie des schémas, lemme \ref{coherent-lemma-induction-principle}.

\medskip\noindent
D'abord, démontrons que (2) -- (7) $\Rightarrow$ (1) si $X$ est affine.
Posons $P = \mathcal{O}_X[0]$. La condition (7) donne $H^i (X, E) = 0$ pour $i \gg 0$.
Ainsi (1) en résulte, puisque $E$ est déterminé par $R\Gamma (X, E)$,
voir le lemme \ref{perfect-lemma-affine-compare-bounded}.

\medskip\noindent
Supposons maintenant $X = U \cup V$, où $U$ est un ouvert quasi-compact de $X$ et
$V$ un ouvert affine, et supposons que l'implication (2) -- (7) $\Rightarrow$ (1)
soit connue pour les schémas $U$, $V$ et $U \cap V$.
Supposons que $E \in D_\QCoh(\mathcal{O}_X)$ satisfasse aux conditions (2) -- (7).
D'après le lemme \ref{perfect-lemma-direct-summand-of-a-restriction} et le
théorème \ref{perfect-theorem-bondal-van-den-Bergh}, il existe un complexe parfait
$Q$ sur $X$ tel que $Q|_U$ engendre $D_\QCoh (\mathcal{O}_U)$. 
Soient $f_1, \dots , f_r \in \Gamma (V, \mathcal{O}_V)$
tels que $V \setminus U = V(f_1, \dots , f_r)$ comme sous-ensembles de
$V$. Soit $K \in D_{perf}(\mathcal{O}_V)$ l'objet
correspondant au complexe de Koszul de $f_1, \dots , f_r$.
Soit $K' \in D_{perf}(\mathcal{O}_X)$ défini par
\begin{equation}
\label{perfect-equation-detecting-bounded-above}
K' = R (V \to X)_* K = R (V \to X)_! K,
\end{equation}
voir Cohomologie, lemmes \ref{cohomology-lemma-pushforward-restriction} et
\ref{cohomology-lemma-pushforward-perfect}. C'est un complexe parfait
sur $X$ à support dans le fermé $X \setminus U \subset V$
et isomorphe à $K$ sur $V$. Par hypothèse, nous savons que
$R\Hom_{\mathcal{O}_X}(Q, E)$ et
$R\Hom_{\mathcal{O}_X}(K', E)$ sont bornés supérieurement.

\medskip\noindent
D'après la seconde description de $K'$ dans (\ref{perfect-equation-detecting-bounded-above}),
on a
$$
\Hom_{D(\mathcal{O}_V)}(K[-i], E|_V) = \Hom_{D(\mathcal{O}_X)}(K'[-i], E) = 0
$$
pour $i \gg 0$. Nous pouvons donc appliquer le
lemme \ref{perfect-lemma-orthogonal-koszul-first-variant} à $E|_V$ pour
obtenir un entier $a$ tel que
$\tau_{\geq a}(E|_V) = \tau_{\geq a} R (U \cap V \to V)_* (E|_{U \cap V})$.
Alors $\tau_{\geq a} E = \tau_{\geq a} R (U \to X)_* (E |_U)$
(vérifier que le morphisme canonique est un isomorphisme après restriction à
$U$ et à $V$). Par conséquent, en appliquant deux fois le lemme \ref{perfect-lemma-bounded-truncation},
on voit que
$$
\Hom_{D(\mathcal{O}_U)}(Q|_U [-i], E|_U) =
\Hom_{D(\mathcal{O}_X)}(Q[-i], R (U \to X)_* (E|_U)) = 0
$$
pour $i \gg 0$. Puisque la proposition est vraie pour $U$ et le générateur
$Q|_U$, on a $E|_U \in D^-_\QCoh(\mathcal{O}_U)$. Mais alors, puisque
le foncteur $R (U \to X)_*$ préserve $D^-_\QCoh$ 
(d'après le lemme \ref{perfect-lemma-quasi-coherence-direct-image}), on obtient
$\tau_{\geq a}E \in D^-_\QCoh(\mathcal{O}_X)$. Ainsi 
$E \in D^-_\QCoh (\mathcal{O}_X)$. 
\end{proof}

\begin{proposition}
\label{perfect-proposition-detecting-bounded-below}
Soit $X$ un schéma quasi-compact et quasi-séparé.
Soit $G \in D_{perf}(\mathcal{O}_X)$
un complexe parfait qui engendre $D_\QCoh (\mathcal{O}_X)$. Soit
$E \in D_\QCoh (\mathcal{O}_X)$. Les conditions suivantes sont équivalentes :
\begin{enumerate}
\item $E \in D^+_\QCoh (\mathcal{O}_X)$,
\item $\Hom_{D(\mathcal{O}_X)}(G[-i], E) = 0$ pour $i \ll 0$,
\item $\Ext^i_X(G, E) = 0$ pour $i \ll 0$,
\item $R\Hom_X(G, E)$ appartient à $D^+(\mathbf{Z})$,
\item $H^i(X, G^\vee \otimes_{\mathcal{O}_X}^\mathbf{L} E) = 0$
pour $i \ll 0$,
\item $R\Gamma(X, G^\vee \otimes_{\mathcal{O}_X}^\mathbf{L} E)$
appartient à $D^+(\mathbf{Z})$,
\item pour tout objet parfait $P$ de $D(\mathcal{O}_X)$,
\begin{enumerate}
\item les assertions (2), (3), (4) sont vraies lorsque $G$ est remplacé par $P$, et
\item $H^i(X, P \otimes_{\mathcal{O}_X}^\mathbf{L} E) = 0$ pour $i \ll 0$,
\item $R\Gamma(X, P \otimes_{\mathcal{O}_X}^\mathbf{L} E)$
appartient à $D^+(\mathbf{Z})$.
\end{enumerate}
\end{enumerate}
\end{proposition}

\begin{proof}
Supposons (1).\newline Puisque
$\Hom_{D(\mathcal{O}_X)}(G[-i], E) = \Hom_{D(\mathcal{O}_X)}(G, E[i])$,
ce groupe est nul pour $i \ll 0$ d'après le
lemme \ref{perfect-lemma-ext-from-perfect-into-bounded-QCoh}. Ceci démontre
que (1) implique (2).

\medskip\noindent
Les parties (2), (3), (4) sont équivalentes d'après la discussion de
Cohomologie, section \ref{cohomology-section-global-RHom}.
Les parties (5) et (6) sont équivalentes, puisque $H^i(X, -) = H^i(R\Gamma(X, -))$
par définition. Les conditions équivalentes (2), (3), (4) sont
équivalentes aux conditions équivalentes (5), (6) d'après
Cohomologie, lemme \ref{cohomology-lemma-dual-perfect-complex},
et le fait que $(G[-i])^\vee = G^\vee[i]$.

\medskip\noindent
Il est clair que (7) implique (2). Réciproquement, 
démontrons que les conditions équivalentes (2) -- (6) impliquent (7).
Rappelons que $G$ est un générateur classique de $D_{perf}(\mathcal{O}_X)$ d'après la
remarque \ref{perfect-remark-classical-generator}.
Pour $P \in D_{perf}(\mathcal{O}_X)$, notons $T(P)$ l'assertion selon laquelle
$R\Hom_X(P, E)$ appartient à $D^+(\mathbf{Z})$.
Il est clair que $T$ est stable par sommes directes,
vérifie la propriété des deux sur trois pour les triangles
distingués, et est stable par facteurs directs et par décalages.
Par conséquent, d'après Catégories dérivées, remarque \ref{derived-remark-check-on-generator},
on voit que (4) implique que $T$ vaut sur tout $D_{perf}(\mathcal{O}_X)$.
Le même argument vaut pour toutes les autres propriétés, si ce n'est que, pour les propriétés
(7)(b) et (7)(c), on utilise aussi que $P \mapsto P^\vee$ est une
autoéquivalence de $D_{perf}(\mathcal{O}_X)$. Un petit détail est omis.

\medskip\noindent
Nous allons démontrer que les conditions équivalentes (2) -- (7) impliquent (1)
en utilisant le principe de récurrence de
Cohomologie des schémas, lemme \ref{coherent-lemma-induction-principle}.

\medskip\noindent
D'abord, démontrons que (2) -- (7) $\Rightarrow$ (1) si $X$ est affine.
Posons $P = \mathcal{O}_X[0]$. La condition (7) donne $H^i (X, E) = 0$
pour $i \ll 0$. Ainsi (1) en résulte, puisque $E$ est
déterminé par $R\Gamma (X, E)$, voir le lemme \ref{perfect-lemma-affine-compare-bounded}.

\medskip\noindent
Supposons maintenant $X = U \cup V$, où $U$ est un ouvert quasi-compact de $X$ et
$V$ un ouvert affine, et supposons que l'implication (2) -- (7) $\Rightarrow$ (1)
soit connue pour les schémas $U$, $V$ et $U \cap V$.
Supposons que $E \in D_\QCoh(\mathcal{O}_X)$ satisfasse aux conditions (2) -- (7).
D'après le lemme \ref{perfect-lemma-direct-summand-of-a-restriction} et le
théorème \ref{perfect-theorem-bondal-van-den-Bergh}, il existe un complexe parfait
$Q$ sur $X$ tel que $Q|_U$ engendre $D_\QCoh (\mathcal{O}_U)$. 
Soient $f_1, \dots , f_r \in \Gamma (V, \mathcal{O}_V)$
tels que $V \setminus U = V(f_1, \dots , f_r)$ comme sous-ensembles de
$V$. Soit $K \in D_{perf}(\mathcal{O}_V)$ l'objet
correspondant au complexe de Koszul de $f_1, \dots , f_r$.
Soit $K' \in D_{perf}(\mathcal{O}_X)$ défini par
\begin{equation}
\label{perfect-equation-detecting-bounded-below}
K' = R (V \to X)_* K = R (V \to X)_! K,
\end{equation}
voir Cohomologie, lemmes \ref{cohomology-lemma-pushforward-restriction} et
\ref{cohomology-lemma-pushforward-perfect}. C'est un complexe parfait
sur $X$ à support dans le fermé $X \setminus U \subset V$
et isomorphe à $K$ sur $V$. Par hypothèse, nous savons que
$R\Hom_{\mathcal{O}_X}(Q, E)$ et
$R\Hom_{\mathcal{O}_X}(K', E)$ sont bornés inférieurement.

\medskip\noindent
D'après la seconde description de $K'$ dans (\ref{perfect-equation-detecting-bounded-below}),
on a
$$
\Hom_{D(\mathcal{O}_V)}(K[-i], E|_V) = \Hom_{D(\mathcal{O}_X)}(K'[-i], E) = 0
$$
pour $i \ll 0$. Nous pouvons donc appliquer le
lemme \ref{perfect-lemma-orthogonal-koszul-second-variant} à $E|_V$ pour
obtenir un entier $a$ tel que
$\tau_{\leq a}(E|_V) = \tau_{\leq a} R(U \cap V \to V)_*(E|_{U \cap V})$.
Alors $\tau_{\leq a} E = \tau_{\leq a} R(U \to X)_*(E|_U)$
(vérifier que le morphisme canonique est un isomorphisme après restriction à
$U$ et à $V$). Par conséquent, en appliquant deux fois le lemme \ref{perfect-lemma-bounded-below-truncation},
on voit que
$$
\Hom_{D(\mathcal{O}_U)}(Q|_U [-i], E|_U) =
\Hom_{D(\mathcal{O}_X)}(Q[-i], R (U \to X)_* (E|_U)) = 0
$$
pour $i \ll 0$. Puisque la proposition est vraie pour $U$ et le générateur
$Q|_U$, on a $E|_U \in D^+_\QCoh (\mathcal{O}_U)$. Mais alors, puisque
le foncteur $R(U \to X)_*$ préserve les objets bornés inférieurement
(voir Cohomologie, section \ref{cohomology-section-derived-functors}), on obtient
$\tau_{\leq a} E \in D^+_\QCoh(\mathcal{O}_X)$. Ainsi 
$E \in D^+_\QCoh (\mathcal{O}_X)$.
\end{proof}







\section{Objets quasi-cohérents dans la catégorie dérivée}
\label{perfect-section-QC}

\noindent
Soit $X$ un schéma. Rappelons que $X_{affine, Zar}$
désigne la catégorie des ouverts affines de $X$ munie de la topologie donnée
par les recouvrements de Zariski standards ; voir
Topologies, définition \ref{topologies-definition-big-small-Zariski}.
Rappelons au lecteur que le topos de $X_{affine, Zar}$
est le petit topos de Zariski de $X$ ; voir Topologies, lemme
\ref{topologies-lemma-alternative-zariski}. Le site $X_{affine, Zar}$
est muni d'un faisceau structural $\mathcal{O}$ et il existe une
équivalence de topos annelés
$$
(\Sh(X_{affine, Zar}), \mathcal{O})
\longrightarrow
(\Sh(X_{Zar}), \mathcal{O})
$$
Voir Descente, équation (\ref{descent-equation-alternative-small-ringed}),
ainsi que la discussion qui l'entoure dans
Descente, section \ref{descent-section-alternative-quasi-coherent},
où une notation légèrement différente est employée.

\medskip\noindent
Dans cette section, nous notons $X_{affine}$ la catégorie sous-jacente à
$X_{affine, Zar}$, munie de la topologie chaotique, de sorte que les faisceaux
coïncident avec les préfaisceaux. En particulier, le faisceau structural $\mathcal{O}$
devient également un faisceau sur $X_{affine}$.
Nous obtenons un morphisme de sites annelés
$$
\epsilon :
(X_{affine, Zar}, \mathcal{O})
\longrightarrow
(X_{affine}, \mathcal{O})
$$
comme dans Cohomologie sur les sites, section \ref{sites-cohomology-section-compare}.
Dans cette section, nous identifierons $D_\QCoh(\mathcal{O}_X)$ à
la catégorie $\mathit{QC}(X_{affine}, \mathcal{O})$ introduite dans
Cohomologie sur les sites, section \ref{sites-cohomology-section-modules-cohomology}.

\begin{lemma}
\label{perfect-lemma-DQCoh-alternative-small}
Dans la situation ci-dessus, il existe des équivalences exactes canoniques entre
les catégories triangulées suivantes :
\begin{enumerate}
\item $D_\QCoh(\mathcal{O}_X)$,
\item $D_\QCoh(X_{Zar}, \mathcal{O})$,
\item $D_\QCoh(X_{affine, Zar}, \mathcal{O})$,
\item $D_\QCoh(X_{affine}, \mathcal{O}_X)$, et
\item $\mathit{QC}(X_{affine}, \mathcal{O})$.
\end{enumerate}
\end{lemma}

\begin{proof}
Si $U \subset V \subset X$ sont des ouverts affines, alors le morphisme d'anneaux
$\mathcal{O}(V) \to \mathcal{O}(U)$ est plat. Par conséquent,
l'équivalence entre (4) et (5) est un cas particulier de
Cohomologie sur les sites, lemme
\ref{sites-cohomology-lemma-cartesian-flat-quasi-coherent}
(la démonstration éclaire également l'énoncé).

\medskip\noindent
Le site annelé $(X_{Zar}, \mathcal{O})$ et l'espace annelé
$(X, \mathcal{O}_X)$ ont la même catégorie de modules d'après
Descente, remarque \ref{descent-remark-Zariski-site-space}.
Par cette équivalence, les modules quasi-cohérents se correspondent d'après
Descente, proposition \ref{descent-proposition-equivalence-quasi-coherent}.
On obtient donc une équivalence exacte canonique entre les catégories triangulées
de (1) et (2).

\medskip\noindent
La discussion précédant le lemme montre que nous avons une équivalence de
topos annelés
$(\Sh(X_{affine, Zar}), \mathcal{O}) \to (\Sh(X_{Zar}), \mathcal{O})$
et donc une équivalence entre les catégories abéliennes de modules.
Comme la notion de module quasi-cohérent est intrinsèque
(Modules sur les sites, lemme \ref{sites-modules-lemma-special-locally-free}),
nous voyons que cette équivalence préserve les sous-catégories
de modules quasi-cohérents. Ainsi, nous obtenons une équivalence exacte canonique
entre les catégories triangulées de (2) et (3).

\medskip\noindent
Pour obtenir une équivalence exacte entre les catégories triangulées
de (3) et (4), nous allons appliquer Cohomologie sur les sites, lemme
\ref{sites-cohomology-lemma-compare-topologies-derived-adequate-modules},
au morphisme $\epsilon : (X_{affine, Zar}, \mathcal{O}) \to
(X_{affine}, \mathcal{O})$ ci-dessus. Nous prenons
$\mathcal{B} = \Ob(X_{affine})$ et nous prenons pour
$\mathcal{A} \subset \textit{PMod}(X_{affine}, \mathcal{O})$
la sous-catégorie pleine des préfaisceaux $\mathcal{F}$
tels que $\mathcal{F}(V) \otimes_{\mathcal{O}(V)} \mathcal{O}(U)
\to \mathcal{F}(U)$ soit un isomorphisme. Observons que, d'après
Descente, lemme \ref{descent-lemma-quasi-coherent-alternative-small},
les objets de $\mathcal{A}$ sont exactement les faisceaux pour la topologie de Zariski
qui sont des modules quasi-cohérents sur $(X_{affine, Zar}, \mathcal{O})$.
D'autre part, d'après Modules sur les sites, lemme
\ref{sites-modules-lemma-chaotic-quasi-coherent},
les objets de $\mathcal{A}$ sont exactement les modules quasi-cohérents
sur $(X_{affine}, \mathcal{O})$, c'est-à-dire pour la topologie chaotique.
Ainsi, si nous montrons que Cohomologie sur les sites, lemme
\ref{sites-cohomology-lemma-compare-topologies-derived-adequate-modules},
s'applique, alors nous obtenons bien l'équivalence canonique entre les
catégories de (3) et (4) au moyen de $\epsilon^*$ et $R\epsilon_*$.

\medskip\noindent
Nous devons vérifier quatre conditions :
\begin{enumerate}
\item Tout objet de $\mathcal{A}$ est un faisceau pour la topologie de Zariski.
Nous l'avons vu ci-dessus.
\item $\mathcal{A}$ est une sous-catégorie de Serre faible de
$\textit{Mod}(X_{affine, Zar}, \mathcal{O})$. Nous avons vu
ci-dessus que $\mathcal{A} = \QCoh(X_{affine, Zar}, \mathcal{O})$ et
nous avons également vu que ses objets,
par l'équivalence $\textit{Mod}(X_{affine, Zar}, \mathcal{O}) =
\textit{Mod}(X, \mathcal{O}_X)$, correspondent aux modules quasi-cohérents
sur $X$. Le résultat découle donc de la discussion de
Schémas, section \ref{schemes-section-quasi-coherent}.
\item Tout objet de $X_{affine}$ possède un recouvrement pour la topologie
chaotique dont les membres sont des éléments de $\mathcal{B}$. C'est le cas
car $\mathcal{B}$ contient tous les objets.
\item Pour tout objet $U$ de $X_{affine}$ et tout $\mathcal{F}$ de $\mathcal{A}$,
nous avons $H^p_{Zar}(U, \mathcal{F}) = 0$ pour $p > 0$.
Cela résulte de l'annulation de la cohomologie des modules quasi-cohérents
sur les schémas affines ; voir
Cohomologie des schémas, lemme
\ref{coherent-lemma-quasi-coherent-affine-cohomology-zero}.
\end{enumerate}
Ceci achève la démonstration.
\end{proof}

\begin{remark}
\label{perfect-remark-QC-compare}
Soit $S$ un schéma.\newline Nous montrerons plus loin que
$\mathit{QC}((\textit{Aff}/S), \mathcal{O})$
est également canoniquement équivalente à $D_\QCoh(\mathcal{O}_S)$.
Voir Faisceaux sur les champs, proposition \ref{stacks-sheaves-proposition-QC-compare}.
\end{remark}












% OK, start here.
%
\chapter{Compléments sur les morphismes}



\phantomsection
\label{more-morphisms-section-phantom}


\section{Introduction}
\label{more-morphisms-section-introduction}

\noindent
Dans ce chapitre, nous poursuivons l'étude des propriétés des morphismes de schémas.
Une référence fondamentale est \cite{EGA}.






\section{Épaississements}
\label{more-morphisms-section-thickenings}

\noindent
La terminologie suivante n'est peut-être pas entièrement usuelle, mais elle est commode.

\begin{definition}
\label{more-morphisms-definition-thickening}
Épaississements.
\begin{enumerate}
\item On dit qu'un schéma $X'$ est un {\it épaississement} d'un schéma $X$ si
$X$ est un sous-schéma fermé de $X'$ et si les espaces topologiques sous-jacents
sont égaux.
\item On dit qu'un schéma $X'$ est un {\it épaississement du premier ordre} d'un schéma $X$ si
$X$ est un sous-schéma fermé de $X'$ et si le faisceau quasi-cohérent d'idéaux
$\mathcal{I} \subset \mathcal{O}_{X'}$ définissant $X$ est de carré nul.
\item On dit qu'un schéma $X'$ est un {\it épaississement d'ordre fini} d'un schéma $X$
si $X$ est un sous-schéma fermé de $X'$ et si le faisceau quasi-cohérent d'idéaux
$\mathcal{I} \subset \mathcal{O}_{X'}$ définissant $X$ est nilpotent, c'est-à-dire
s'il existe un entier $n \geq 0$ tel que $\mathcal{I}^{n + 1} = 0$.
\item On dit qu'un schéma $X'$ est un {\it épaississement d'ordre $n$} d'un schéma $X$
si $X$ est un sous-schéma fermé de $X$ et si $\mathcal{I}^{n + 1} = 0$,
où $\mathcal{I} \subset \mathcal{O}_{X'}$ est le
faisceau quasi-cohérent d'idéaux définissant $X$.
\item Étant donnés deux épaississements $X \subset X'$ et $Y \subset Y'$, un
{\it morphisme d'épaississements} est un morphisme $f' : X' \to Y'$ tel que
$f'(X) \subset Y$, c'est-à-dire tel que $f'|_X$ se factorise par le sous-schéma
fermé $Y$. Dans cette situation, on pose $f = f'|_X : X \to Y$ et on dit
que $(f, f') : (X \subset X') \to (Y \subset Y')$ est un morphisme
d'épaississements.
\item Soit $S$ un schéma. On définit de même les {\it épaississements sur $S$} et
les {\it morphismes d'épaississements sur $S$}. Cela signifie que les schémas
$X, X', Y, Y'$ ci-dessus sont des schémas sur $S$ et que les morphismes
$X \to X'$, $Y \to Y'$ et $f' : X' \to Y'$ sont des morphismes sur $S$.
\end{enumerate}
\end{definition}

\noindent
Soit $i_X : X \to X'$ un épaississement. Toute section locale du noyau
$\mathcal{I} = \Ker(i_X^\sharp)$ est localement nilpotente. Cela donne
un certain contrôle sur le faisceau structural $\mathcal{O}_{X'}$, mais,
techniquement, la classe des épaississements d'ordre fini $X \subset X'$
est beaucoup plus facile à manier. En effet, si $X'$ est un épaississement d'ordre $n$,
on dispose d'une filtration
$$
0 = \mathcal{I}^{n + 1} \subset
\mathcal{I}^n \subset
\mathcal{I}^{n - 1} \subset \ldots \subset
\mathcal{I} \subset \mathcal{O}_{X'}
$$
et l'on voit que $X'$ est filtré par des sous-espaces fermés
$$
X = X_1 \subset X_2 \subset \ldots \subset X_n \subset X_{n + 1} = X'
$$
tels que chaque paire $X_i \subset X_{i + 1}$ soit un épaississement du premier ordre
sur $S$. Par de simples raisonnements par récurrence, de nombreux résultats établis pour
les épaississements du premier ordre se reformulent en résultats sur les épaississements d'ordre fini.

\medskip\noindent
Les épaississements du premier ordre se décrivent comme suit (voir
Modules, lemme \ref{modules-lemma-double-structure-gives-derivation}).

\begin{lemma}
\label{more-morphisms-lemma-first-order-thickening}
Soit $X$ un schéma sur une base $S$. Considérons une suite exacte courte
$$
0 \to \mathcal{I} \to \mathcal{A} \to \mathcal{O}_X \to 0
$$
de faisceaux sur $X$, où $\mathcal{A}$ est un faisceau de
$f^{-1}\mathcal{O}_S$-algèbres,
$\mathcal{A} \to \mathcal{O}_X$ est une surjection
de faisceaux de $f^{-1}\mathcal{O}_S$-algèbres et $\mathcal{I}$ en est le noyau.
Si
\begin{enumerate}
\item $\mathcal{I}$ est un idéal de carré nul dans $\mathcal{A}$ et
\item $\mathcal{I}$ est quasi-cohérent en tant que $\mathcal{O}_X$-module,
\end{enumerate}
alors $X' = (X, \mathcal{A})$ est un schéma et $X \to X'$ est un épaississement
du premier ordre sur $S$. De plus, tout épaississement du premier ordre sur
$S$ est de cette forme.
\end{lemma}

\begin{proof}
Il est clair que $X'$ est un espace localement annelé. Soit $U = \Spec(B)$
un ouvert affine de $X$. Posons $A = \Gamma(U, \mathcal{A})$. Remarquons que,
puisque $H^1(U, \mathcal{I}) = 0$ (voir Cohomologie des schémas, lemme
\ref{coherent-lemma-quasi-coherent-affine-cohomology-zero}),
l'application $A \to B$ est surjective. Par hypothèse, le noyau
$I = \mathcal{I}(U)$ est un idéal de carré nul de l'anneau $A$.
D'après
Schémas, lemme \ref{schemes-lemma-morphism-into-affine},
il existe un morphisme canonique d'espaces localement annelés
$$
(U, \mathcal{A}|_U) \longrightarrow \Spec(A)
$$
provenant de l'application $B \to \Gamma(U, \mathcal{A})$. Ce morphisme
s'inscrit dans le diagramme commutatif
$$
\xymatrix{
(U, \mathcal{O}_X|_U) \ar[d] \ar[r] & \Spec(B) \ar[d] \\
(U, \mathcal{A}|_U) \ar[r] & \Spec(A)
}
$$
et l'on voit donc qu'il est un homéomorphisme sur les espaces topologiques sous-jacents.
Pour montrer que c'est un isomorphisme, il suffit donc de vérifier qu'il induit
un isomorphisme sur les anneaux locaux.
Pour $u \in U$ correspondant à l'idéal premier $\mathfrak p \subset A$,
on obtient un diagramme commutatif de suites exactes courtes
$$
\xymatrix{
0 \ar[r] &
I_{\mathfrak p} \ar[r] \ar[d] &
A_{\mathfrak p} \ar[r] \ar[d] &
B_{\mathfrak p} \ar[r] \ar[d] & 0 \\
0 \ar[r] &
\mathcal{I}_u \ar[r] &
\mathcal{A}_u \ar[r] &
\mathcal{O}_{X, u} \ar[r] & 0.
}
$$
Les flèches verticales de gauche et de droite sont des isomorphismes, car
$\mathcal{I}$ et $\mathcal{O}_X$ sont des faisceaux quasi-cohérents.
La flèche du milieu est donc elle aussi un isomorphisme. Ainsi, tout point
de $X' = (X, \mathcal{A})$ possède un voisinage affine et $X'$ est un
schéma, comme voulu.
\end{proof}

\begin{lemma}
\label{more-morphisms-lemma-thickening-affine-scheme}
\begin{slogan}
Le caractère affine est invariant par épaississement
\end{slogan}
\begin{reference}
Le cas d'un épaississement d'ordre fini est \cite[Proposition 5.1.9]{EGA1}.
\end{reference}
Tout épaississement d'un schéma affine est affine.
\end{lemma}

\begin{proof}
C'est un cas particulier de
Limites, proposition \ref{limits-proposition-affine}.
\end{proof}

\begin{proof}[Démonstration pour un épaississement d'ordre fini]
Supposons que $X \subset X'$ soit un épaississement d'ordre fini avec $X$ affine.
On peut alors appliquer le critère de Serre pour prouver que $X'$ est affine. Plus précisément,
on utilisera Cohomologie des schémas, lemme
\ref{coherent-lemma-quasi-compact-h1-zero-covering}. Soit $\mathcal{F}$ un
$\mathcal{O}_{X'}$-module quasi-cohérent. Il suffit de montrer que
$H^1(X', \mathcal{F}) = 0$. Notons $i : X \to X'$ l'immersion fermée donnée
et posons
$\mathcal{I} = \Ker(i^\sharp : \mathcal{O}_{X'} \to i_*\mathcal{O}_X)$.
D'après la discussion des épaississements d'ordre fini (qui suit
la définition \ref{more-morphisms-definition-thickening}), il existe un entier $n \geq 0$
et une filtration
$$
0 = \mathcal{F}_{n + 1} \subset \mathcal{F}_n \subset
\mathcal{F}_{n - 1} \subset \ldots \subset
\mathcal{F}_0 = \mathcal{F}
$$
par des sous-modules quasi-cohérents tels que $\mathcal{F}_a/\mathcal{F}_{a + 1}$ soit
annulé par $\mathcal{I}$. On peut en effet prendre
$\mathcal{F}_a = \mathcal{I}^a\mathcal{F}$. Alors
$\mathcal{F}_a/\mathcal{F}_{a + 1} = i_*\mathcal{G}_a$ pour un certain
$\mathcal{O}_X$-module quasi-cohérent $\mathcal{G}_a$, voir Morphismes, lemme
\ref{morphisms-lemma-i-star-equivalence}. On obtient
$$
H^1(X', \mathcal{F}_a/\mathcal{F}_{a + 1}) =
H^1(X', i_*\mathcal{G}_a) = H^1(X, \mathcal{G}_a) = 0
$$
La deuxième égalité résulte de Cohomologie des schémas, lemme
\ref{coherent-lemma-relative-affine-cohomology},
et la dernière de Cohomologie des schémas, lemme
\ref{coherent-lemma-quasi-coherent-affine-cohomology-zero}.
Ainsi, $\mathcal{F}$ possède une filtration finie dont les quotients successifs
ont leur premier groupe de cohomologie nul ; un simple
raisonnement par récurrence montre donc que $H^1(X', \mathcal{F}) = 0$.
\end{proof}

\begin{lemma}
\label{more-morphisms-lemma-base-change-thickening}
Soit $S \subset S'$ un épaississement de schémas. Soit $X' \to S'$ un morphisme
et posons $X = S \times_{S'} X'$. Alors $(X \subset X') \to (S \subset S')$
est un morphisme d'épaississements. Si $S \subset S'$ est un épaississement du premier
ordre (resp.\ d'ordre fini), alors $X \subset X'$ est un épaississement du premier
ordre (resp.\ d'ordre fini).
\end{lemma}

\begin{proof}
Omise.
\end{proof}

\begin{lemma}
\label{more-morphisms-lemma-composition-thickening}
\begin{slogan}
Les composés d'épaississements sont des épaississements
\end{slogan}
Si $S \subset S'$ et $S' \subset S''$ sont des épaississements, alors
$S \subset S''$ en est un également.
\end{lemma}

\begin{proof}
Omise.
\end{proof}

\begin{lemma}
\label{more-morphisms-lemma-descending-property-thickening}
La propriété d'être un épaississement est locale pour la topologie fpqc.
Il en va de même pour les épaississements du premier ordre.
\end{lemma}

\begin{proof}
L'énoncé signifie ceci : soit $X \to X'$ un morphisme
de schémas et soit $\{g_i : X'_i \to X'\}$
un recouvrement fpqc tel que le changement de base $X_i \to X'_i$
soit un épaississement pour tout $i$. Alors $X \to X'$ est un épaississement.
Puisque les morphismes $g_i$ sont conjointement surjectifs, on en déduit
que $X \to X'$ est surjectif. D'après
Descente, lemme \ref{descent-lemma-descending-property-closed-immersion},
on en déduit que $X \to X'$ est une immersion fermée.
Ainsi, $X \to X'$ est un épaississement. On omet la démonstration dans le
cas des épaississements du premier ordre.
\end{proof}









\section{Morphismes d'épaississements}
\label{more-morphisms-section-morphisms-thickenings}

\noindent
Si $(f, f') : (X \subset X') \to (Y \subset Y')$ est un morphisme
d'épaississements de schémas, alors les propriétés du morphisme
$f$ se transmettent souvent à $f'$. Il existe plusieurs variantes.

\begin{lemma}
\label{more-morphisms-lemma-thicken-property-morphisms}
Soit $(f, f') : (X \subset X') \to (S \subset S')$ un morphisme
d'épaississements. Alors
\begin{enumerate}
\item $f$ est un morphisme affine si et seulement si $f'$ est un morphisme affine,
\item $f$ est un morphisme surjectif si et seulement si $f'$ est un morphisme surjectif,
\item $f$ est quasi-compact si et seulement si $f'$ est quasi-compact,
\item $f$ est universellement fermé si et seulement si $f'$ est universellement fermé,
\item $f$ est entier si et seulement si $f'$ est entier,
\item $f$ est (quasi-)séparé si et seulement si $f'$ est (quasi-)séparé,
\item $f$ est universellement injectif si et seulement si $f'$ est universellement injectif,
\item $f$ est universellement ouvert si et seulement si $f'$ est universellement ouvert,
\item $f$ est quasi-affine si et seulement si $f'$ est quasi-affine, et
\item ajouter d'autres propriétés ici.
\end{enumerate}
\end{lemma}

\begin{proof}
Remarquons que $S \to S'$ et $X \to X'$ sont des homéomorphismes universels
(voir par exemple
Morphismes, lemme \ref{morphisms-lemma-reduction-universal-homeomorphism}).
Cela entraîne immédiatement (2), (3), (4), (7) et (8).
Le point (1) résulte du lemme \ref{more-morphisms-lemma-thickening-affine-scheme},
qui établit une correspondance biunivoque entre
les ouverts affines de $S$ et de $S'$, ainsi qu'entre ceux de $X$ et de $X'$.
Le point (9) résulte de
Limites, lemme \ref{limits-lemma-thickening-quasi-affine},
et de la remarque précédente sur les ouverts affines de $S$ et de $S'$.
Le point (5) résulte de (1) et (4), d'après
Morphismes, lemme \ref{morphisms-lemma-integral-universally-closed}.
Enfin, remarquons que

$$
S \times_X S = S \times_{X'} S \to S \times_{X'} S' \to S' \times_{X'} S'
$$
est un épaississement (les deux flèches sont des épaississements d'après
le lemme \ref{more-morphisms-lemma-base-change-thickening}).
En appliquant (3) et (4) au morphisme
$(S \subset S') \to (S \times_X S \to S' \times_{X'} S')$,
on obtient donc (6).
\end{proof}

\begin{lemma}
\label{more-morphisms-lemma-thicken-property-relatively-ample}
Soit $(f, f') : (X \subset X') \to (S \subset S')$ un morphisme
d'épaississements. Soit $\mathcal{L}'$ un faisceau inversible sur $X'$
et notons $\mathcal{L}$ sa restriction à $X$.
Alors $\mathcal{L}'$ est $f'$-ample si et seulement si
$\mathcal{L}$ est $f$-ample.
\end{lemma}

\begin{proof}
Rappelons que l'amplitude relative se vérifie sur chaque
ouvert affine de la base, voir
Morphismes, définition \ref{morphisms-definition-relatively-ample}.
D'après le lemme \ref{more-morphisms-lemma-thickening-affine-scheme},
il existe une correspondance biunivoque entre
les ouverts affines de $S$ et de $S'$.
On peut donc supposer que $S$ et $S'$ sont affines,
et l'on est ramené à prouver que
$\mathcal{L}'$ est ample si et seulement si
$\mathcal{L}$ est ample.
C'est Limites, lemme \ref{limits-lemma-ample-on-reduction}.
\end{proof}

\begin{lemma}
\label{more-morphisms-lemma-thicken-property-morphisms-cartesian}
Soit $(f, f') : (X \subset X') \to (S \subset S')$ un morphisme
d'épaississements tel que $X = S \times_{S'} X'$. Si $S \subset S'$
est un épaississement d'ordre fini, alors
\begin{enumerate}
\item $f$ est une immersion fermée si et seulement si $f'$ est une immersion fermée,
\item $f$ est localement de type fini si et seulement si $f'$ est
localement de type fini,
\item $f$ est localement quasi-fini si et seulement si $f'$ est localement
quasi-fini,
\item $f$ est localement de type fini et de dimension relative $d$ si et
seulement si $f'$ est localement de type fini et de dimension relative $d$,
\item $\Omega_{X/S} = 0$ si et seulement si $\Omega_{X'/S'} = 0$,
\item $f$ est non ramifié si et seulement si $f'$ est non ramifié,
\item $f$ est propre si et seulement si $f'$ est propre,
\item $f$ est fini si et seulement si $f'$ est fini,
\item $f$ est un monomorphisme si et seulement si $f'$ est un monomorphisme,
\item $f$ est une immersion si et seulement si $f'$ est une immersion, et
\item ajouter d'autres propriétés ici.
\end{enumerate}
\end{lemma}

\begin{proof}
Les propriétés $\mathcal{P}$ énumérées dans le lemme sont toutes stables
par changement de base ; si $f'$ possède la propriété $\mathcal{P}$,
il en va donc de même de $f$. Voir
Schémas, lemmes \ref{schemes-lemma-base-change-immersion} et
\ref{schemes-lemma-base-change-monomorphism},
et
Morphismes, lemmes
\ref{morphisms-lemma-base-change-finite-type},
\ref{morphisms-lemma-base-change-quasi-finite},
\ref{morphisms-lemma-base-change-relative-dimension-d},
\ref{morphisms-lemma-base-change-differentials},
\ref{morphisms-lemma-base-change-unramified},
\ref{morphisms-lemma-base-change-proper} et
\ref{morphisms-lemma-base-change-finite}.

\medskip\noindent
Dans chaque cas, le sens intéressant consiste donc à supposer
que $f$ possède la propriété et à en déduire que $f'$ la possède aussi.
Par récurrence sur l'ordre de l'épaississement, on peut
supposer que $S \subset S'$ est un épaississement du premier ordre, voir
la discussion qui suit immédiatement
la définition \ref{more-morphisms-definition-thickening}.

\medskip\noindent
La plupart des démonstrations utiliseront une réduction au cas affine. Soient
$U' \subset S'$ un ouvert affine et $V' \subset X'$ un ouvert affine
au-dessus de $U'$. Écrivons $U' = \Spec(A')$ et notons $I \subset A'$ l'idéal
définissant le sous-schéma fermé $U' \cap S$. Écrivons $V' = \Spec(B')$.
Alors $V' \cap X = \Spec(B'/IB')$. En posant $A = A'/I$ et
$B = B'/IB'$, on obtient un diagramme commutatif
$$
\xymatrix{
0 \ar[r] &
IB' \ar[r] &
B' \ar[r] &
B \ar[r] & 0 \\
0 \ar[r] &
IA' \ar[r] \ar[u] &
A' \ar[r] \ar[u] &
A \ar[r] \ar[u] & 0
}
$$
dont les lignes sont exactes et où $I^2 = 0$.

\medskip\noindent
Traduction algébrique de (1) : si $A \to B$ est surjectif,
alors $A' \to B'$ est surjectif. Cela résulte du
lemme de Nakayama (Algèbre, lemme \ref{algebra-lemma-NAK}).

\medskip\noindent
Traduction algébrique de (2) : si $A \to B$ est un homomorphisme d'anneaux
de type fini, alors $A' \to B'$ est de type fini. Cela résulte du
lemme de Nakayama (Algèbre, lemme \ref{algebra-lemma-NAK})
appliqué à un homomorphisme $A'[x_1, \ldots, x_n] \to B'$ tel que
$A[x_1, \ldots, x_n] \to B$ soit surjectif.

\medskip\noindent
Démonstration de (3). Elle résulte de (2) et du fait que la quasi-finitude d'un morphisme
localement de type fini se vérifie sur les fibres, voir
Morphismes, lemme \ref{morphisms-lemma-quasi-finite-at-point-characterize}.

\medskip\noindent
Démonstration de (4). Elle résulte de (2) et du fait que la propriété supplémentaire d'« être de
dimension relative $d$ » se vérifie sur les fibres (par définition, voir
Morphismes, définition \ref{morphisms-definition-relative-dimension-d}).

\medskip\noindent
Traduction algébrique de (5) : si $\Omega_{B/A} = 0$, alors
$\Omega_{B'/A'} = 0$. D'après
Algèbre, lemme \ref{algebra-lemma-differentials-base-change},
on a $0 = \Omega_{B/A} = \Omega_{B'/A'}/I\Omega_{B'/A'}$.
Ainsi, $\Omega_{B'/A'} = 0$ par le
lemme de Nakayama (Algèbre, lemme \ref{algebra-lemma-NAK}).

\medskip\noindent
Traduction algébrique de (6) : si $A \to B$ est un homomorphisme non ramifié,
alors $A' \to B'$ est non ramifié. Puisque $A \to B$ est de type
fini, on voit que $A' \to B'$ est de type fini d'après (2).
Comme $A \to B$ est non ramifié, on a $\Omega_{B/A} = 0$. D'après
(5), on a $\Omega_{B'/A'} = 0$. Ainsi, $A' \to B'$ est non ramifié.

\medskip\noindent
Démonstration de (7). Elle résulte de la combinaison de (2) avec
les résultats du lemme \ref{more-morphisms-lemma-thicken-property-morphisms}
et du fait que propre signifie quasi-compact $+$
séparé $+$ localement de type fini $+$ universellement fermé.

\medskip\noindent
Démonstration de (8). Elle résulte de la combinaison de (2) avec
les résultats du lemme \ref{more-morphisms-lemma-thicken-property-morphisms}
et du fait que fini signifie entier $+$ localement
de type fini (Morphismes, lemme \ref{morphisms-lemma-finite-integral}).

\medskip\noindent
Démonstration de (9). Puisque $f$ est un monomorphisme, on a $X = X \times_S X$.
On peut appliquer les résultats déjà démontrés au morphisme d'épaississements
$(X \subset X') \to (X \times_S X \subset X' \times_{S'} X')$.
On en déduit que $X' \to X' \times_{S'} X'$ est une immersion fermée d'après (1).
En fait, c'est un épaississement du premier ordre, car l'idéal définissant
l'immersion fermée
$X' \to X' \times_{S'} X'$ est contenu dans l'image inverse de l'idéal
$\mathcal{I} \subset \mathcal{O}_{S'}$ définissant $S$ dans $S'$.
En effet, $X = X \times_S X = (X' \times_{S'} X') \times_{S'} S$
est contenu dans $X'$. D'après
Morphismes, lemme \ref{morphisms-lemma-differentials-diagonal},
il suffit donc de montrer que
$\Omega_{X'/S'} = 0$, ce qui résulte de (5)
et de l'énoncé correspondant pour $X/S$.

\medskip\noindent
Démonstration de (10). Si $f : X \to S$ est une immersion, il se factorise en
$X \to U \to S$, où $U \to S$ est une immersion ouverte et $X \to U$ une
immersion fermée. Soit $U' \subset S'$ le sous-schéma ouvert dont
l'espace topologique sous-jacent est celui de $U$. Alors $X' \to S'$
se factorise par $U'$ et l'on en déduit que $X' \to U'$ est une immersion
fermée d'après (1). Cela achève la démonstration.
\end{proof}

\noindent
Le lemme suivant est une variante du précédent. Au lieu de supposer
que les épaississements considérés sont d'ordre fini (ce qui permet de transmettre
la propriété d'être localement de type fini de $f$ à $f'$),
on suppose directement que $f$ et $f'$ sont tous deux localement de
type fini.

\begin{lemma}
\label{more-morphisms-lemma-properties-that-extend-over-thickenings}
Soit $(f, f') : (X \subset X') \to (Y \to Y')$ un morphisme
d'épaississements. Supposons que $f$ et $f'$ soient localement de type fini
et que $X = Y \times_{Y'} X'$. Alors
\begin{enumerate}
\item $f$ est localement quasi-fini si et seulement si $f'$ est localement quasi-fini,
\item $f$ est fini si et seulement si $f'$ est fini,
\item $f$ est une immersion fermée si et seulement si $f'$ est une immersion fermée,
\item $\Omega_{X/Y} = 0$ si et seulement si $\Omega_{X'/Y'} = 0$,
\item $f$ est non ramifié si et seulement si $f'$ est non ramifié,
\item $f$ est un monomorphisme si et seulement si $f'$ est un monomorphisme,
\item $f$ est une immersion si et seulement si $f'$ est une immersion,
\item $f$ est propre si et seulement si $f'$ est propre, et
\item ajouter d'autres propriétés ici.
\end{enumerate}
\end{lemma}

\begin{proof}
Les propriétés $\mathcal{P}$ énumérées dans le lemme sont toutes stables
par changement de base ; si $f'$ possède la propriété $\mathcal{P}$,
il en va donc de même de $f$. Voir
Schémas, lemmes \ref{schemes-lemma-base-change-immersion} et
\ref{schemes-lemma-base-change-monomorphism},
et
Morphismes, lemmes
\ref{morphisms-lemma-base-change-quasi-finite},
\ref{morphisms-lemma-base-change-relative-dimension-d},
\ref{morphisms-lemma-base-change-differentials},
\ref{morphisms-lemma-base-change-unramified},
\ref{morphisms-lemma-base-change-proper} et
\ref{morphisms-lemma-base-change-finite}.
Dans chaque cas, il suffit donc de prouver que, si $f$ possède
la propriété voulue, $f'$ la possède aussi.

\medskip\noindent
Un morphisme est localement quasi-fini si et seulement s'il est localement
de type fini et si ses fibres schématiques sont des espaces discrets, voir
Morphismes, lemme \ref{morphisms-lemma-locally-quasi-finite-fibres}.
Puisque le passage à un épaississement ne modifie pas
l'espace topologique sous-jacent, on voit que $f'$ est localement quasi-fini si
(et seulement si) $f$ l'est. Cela prouve (1).

\medskip\noindent
Le cas (2) résulte du cas (5) du lemme \ref{more-morphisms-lemma-thicken-property-morphisms}
et du fait que les morphismes finis sont précisément
les morphismes entiers localement de type fini
(Morphismes, lemme \ref{morphisms-lemma-finite-integral}).

\medskip\noindent
Cas (3). Il résulte immédiatement de
Morphismes, lemme \ref{morphisms-lemma-check-closed-infinitesimally}.

\medskip\noindent
Le cas (4) résulte de l'énoncé algébrique suivant : soient $A$ un anneau et
$I \subset A$ un idéal localement nilpotent. Soit $B$ une
$A$-algèbre de type fini. Si $\Omega_{(B/IB)/(A/I)} = 0$, alors $\Omega_{B/A} = 0$.
En effet, l'hypothèse signifie que $I\Omega_{B/A} = 0$, voir
Algèbre, lemme \ref{algebra-lemma-differentials-base-change}.
D'autre part, $\Omega_{B/A}$ est un $B$-module de type fini, voir
Algèbre, lemme \ref{algebra-lemma-differentials-finitely-generated}.
L'annulation de $\Omega_{B/A}$ résulte donc du lemme de
Nakayama (Algèbre, lemme \ref{algebra-lemma-NAK}) et du fait
que $IB$ est contenu dans le radical de Jacobson de $B$.

\medskip\noindent
Le cas (5) résulte immédiatement de (4) et de
Morphismes, lemme \ref{morphisms-lemma-unramified-omega-zero}.

\medskip\noindent
Démonstration de (6). Puisque $f$ est un monomorphisme, on a $X = X \times_Y X$.
On peut appliquer les résultats déjà démontrés au morphisme d'épaississements
$(X \subset X') \to (X \times_Y X \subset X' \times_{Y'} X')$.
On en déduit que $\Delta_{X'/Y'} : X' \to X' \times_{Y'} X'$
est une immersion fermée d'après (3). En fait, $\Delta_{X'/Y'}$ est une bijection sur
les ensembles sous-jacents ; $\Delta_{X'/Y'}$ est donc un épaississement. D'autre part,
$\Delta_{X'/Y'}$ est localement de présentation finie d'après
Morphismes, lemme \ref{morphisms-lemma-diagonal-morphism-finite-type}.
Autrement dit, $\Delta_{X'/Y'}(X')$ est défini par
un faisceau quasi-cohérent d'idéaux
$\mathcal{J} \subset \mathcal{O}_{X' \times_{Y'} X'}$ de type fini.
Puisque $\Omega_{X'/Y'} = 0$ d'après (5), on voit que
le faisceau conormal de $X' \to X' \times_{Y'} X'$ est nul d'après
Morphismes, lemme \ref{morphisms-lemma-differentials-diagonal}.
Autrement dit, $\mathcal{J}/\mathcal{J}^2 = 0$.
Cela entraîne que $\Delta_{X'/Y'}$ est un isomorphisme, par exemple
d'après Algèbre, lemme \ref{algebra-lemma-ideal-is-squared-union-connected}.

\medskip\noindent
Démonstration de (7). Si $f : X \to Y$ est une immersion, il se factorise en
$X \to V \to Y$, où $V \to Y$ est une immersion ouverte et $X \to V$ une
immersion fermée. Soit $V' \subset Y'$ le sous-schéma ouvert dont
l'espace topologique sous-jacent est celui de $V$. Alors $X' \to V'$
se factorise par $V'$ et l'on en déduit que $X' \to V'$ est une immersion
fermée d'après (3).

\medskip\noindent
Le cas (8) résulte du lemme \ref{more-morphisms-lemma-thicken-property-morphisms}
et de la définition des morphismes propres comme morphismes quasi-compacts,
universellement fermés, séparés et localement de type fini.
\end{proof}








\section{Groupes de Picard des épaississements}
\label{more-morphisms-section-picard-group-thickening}

\noindent
Quelques résultats sur les groupes de Picard des épaississements.

\begin{lemma}
\label{more-morphisms-lemma-picard-group-first-order-thickening}
Soit $X \subset X'$ un épaississement du premier ordre
de faisceau d'idéaux $\mathcal{I}$. Il existe alors une suite
exacte canonique
$$
\xymatrix{
0 \ar[r] &
H^0(X, \mathcal{I}) \ar[r] &
H^0(X', \mathcal{O}_{X'}^*) \ar[r] &
H^0(X, \mathcal{O}^*_X) \ar `r[d] `d[l] `l[llld] `d[dll] [dll] \\
& H^1(X, \mathcal{I}) \ar[r] &
\Pic(X') \ar[r] &
\Pic(X) \ar `r[d] `d[l] `l[llld] `d[dll] [dll] \\
& H^2(X, \mathcal{I}) \ar[r] & \ldots \ar[r] & \ldots
}
$$
de groupes abéliens.
\end{lemma}

\begin{proof}
C'est la suite exacte longue de cohomologie associée à la
suite exacte courte de faisceaux de groupes abéliens
$$
0 \to \mathcal{I} \to \mathcal{O}_{X'}^* \to \mathcal{O}_X^* \to 0
$$
dans laquelle la première application envoie une section locale $f$ de $\mathcal{I}$
sur la section inversible $1 + f$ de $\mathcal{O}_{X'}$.
On utilise aussi l'identification du groupe de Picard d'un
espace annelé au premier groupe de cohomologie du faisceau
des fonctions inversibles, voir
Cohomologie, lemme \ref{cohomology-lemma-h1-invertible}.
\end{proof}

\begin{lemma}
\label{more-morphisms-lemma-torsion-pic-thickening}
Soit $X \subset X'$ un épaississement. Soit $n$ un entier
inversible dans $\mathcal{O}_X$. Alors l'application
$\Pic(X')[n] \to \Pic(X)[n]$ est bijective.
\end{lemma}

\begin{proof}[Démonstration pour un épaississement d'ordre fini]
D'après le principe général expliqué à la suite de
la définition \ref{more-morphisms-definition-thickening},
on est ramené au cas d'un épaississement du premier ordre.
On peut alors utiliser le lemme \ref{more-morphisms-lemma-picard-group-first-order-thickening}
pour voir qu'il suffit de montrer que
$H^1(X, \mathcal{I})[n]$, $H^1(X, \mathcal{I})/n$ et
$H^2(X, \mathcal{I})[n]$ sont nuls.
Cela résulte du fait que la multiplication par $n$ sur $\mathcal{I}$
est un isomorphisme, puisque c'est un $\mathcal{O}_X$-module.
\end{proof}

\begin{proof}[Démonstration dans le cas général]
Soit $\mathcal{I} \subset \mathcal{O}_{X'}$ le faisceau quasi-cohérent d'idéaux
définissant $X$. On dispose alors d'une suite exacte courte de
groupes abéliens
$$
0 \to (1 + \mathcal{I})^* \to \mathcal{O}_{X'}^* \to \mathcal{O}_X^* \to 0
$$
On obtient une suite exacte longue de cohomologie comme dans l'énoncé du
lemme \ref{more-morphisms-lemma-picard-group-first-order-thickening},
où $H^i(X, \mathcal{I})$ est remplacé par $H^i(X, (1 + \mathcal{I})^*)$.
Il suffit donc de montrer que l'élévation à la puissance $n$ est un
isomorphisme $(1 + \mathcal{I})^* \to (1 + \mathcal{I})^*$.
En prenant les sections sur les ouverts affines, cela résulte de
Algèbre, lemme \ref{algebra-lemma-lift-nth-roots}.
\end{proof}





\section{Voisinages infinitésimaux}
\label{more-morphisms-section-first-order-infinitesimal-neighbourhood}

\noindent
Voici une construction naturelle d'épaississements d'ordre fini.
Soit $i : Z \to X$ une immersion de schémas. Choisissons un
sous-schéma ouvert $U \subset X$ tel que $i$ identifie $Z$ à un sous-schéma
fermé $Z \subset U$. Soit $\mathcal{I} \subset \mathcal{O}_U$ le
faisceau quasi-cohérent d'idéaux définissant $Z$ dans $U$. Pour $n \geq 1$,
on peut considérer le sous-schéma fermé $Z_n \subset U$
défini par le faisceau quasi-cohérent d'idéaux $\mathcal{I}^{n + 1}$.

\begin{definition}
\label{more-morphisms-definition-first-order-infinitesimal-neighbourhood}
Soit $i : Z \to X$ une immersion de schémas.
\begin{enumerate}
\item Le {\it voisinage infinitésimal du premier ordre} de $Z$ dans $X$ est
l'épaississement du premier ordre $Z \subset Z_1$ sur $X$ décrit ci-dessus.
\item Le {\it voisinage infinitésimal d'ordre $n$} de $Z$ dans $X$ est
l'épaississement d'ordre $n$ $Z \subset Z_n$ sur $X$ décrit ci-dessus.
\end{enumerate}
\end{definition}

\noindent
Ces épaississements possèdent la propriété universelle suivante (qui dissipe
toute crainte que la construction ci-dessus ne dépende du choix de l'ouvert
$U$).

\begin{lemma}
\label{more-morphisms-lemma-first-order-infinitesimal-neighbourhood}
Soit $i : Z \to X$ une immersion de schémas.
\begin{enumerate}
\item Le voisinage infinitésimal du premier ordre $Z'$ de $Z$ dans $X$
possède la propriété universelle suivante : pour tout diagramme commutatif
$$
\xymatrix{
Z \ar[d]_i & T \ar[l]^a \ar[d] \\
X & T' \ar[l]_b
}
$$
où $T \subset T'$ est un épaississement du premier ordre sur $X$, il existe
un unique morphisme $(a', a) : (T \subset T') \to (Z \subset Z')$
d'épaississements sur $X$.
\item Pour $n \geq 1$, le voisinage infinitésimal d'ordre $n$, $Z_n$,
de $Z$ dans $X$ possède la propriété universelle suivante : pour tout diagramme
commutatif
$$
\xymatrix{
Z \ar[d]_i & T \ar[l]^a \ar[d] \\
X & T' \ar[l]_b
}
$$
où $T \subset T'$ est un épaississement d'ordre $n$ sur $X$, il existe
un unique morphisme $(a', a) : (T \subset T') \to (Z \subset Z_n)$
d'épaississements sur $X$.
\end{enumerate}
\end{lemma}

\begin{proof}
On démontrera seulement (1).
Soit $U \subset X$ l'ouvert utilisé dans la construction de $Z'$, c'est-à-dire un
ouvert tel que $Z$ soit identifié à un sous-schéma fermé de $U$ défini par
le faisceau quasi-cohérent d'idéaux $\mathcal{I}$.
Puisque $|T| = |T'|$, on voit que $b(T') \subset U$. On peut donc
considérer $b$ comme un morphisme à valeurs dans $U$. Soit $\mathcal{J} \subset \mathcal{O}_{T'}$
l'idéal définissant $T$. Puisque $b(T) \subset Z$ d'après le diagramme ci-dessus,
on voit que $b^\sharp(b^{-1}\mathcal{I}) \subset \mathcal{J}$. Comme
$T'$ est un épaississement du premier ordre de $T$, on a $\mathcal{J}^2 = 0$,
d'où $b^\sharp(b^{-1}(\mathcal{I}^2)) = 0$. D'après
Schémas, lemme \ref{schemes-lemma-characterize-closed-subspace},
cela entraîne que $b$ se factorise par $Z'$. Notons $a' : T' \to Z'$
cette factorisation ; tout est alors clair.
\end{proof}

\begin{lemma}
\label{more-morphisms-lemma-infinitesimal-neighbourhood-conormal}
Soit $i : Z \to X$ une immersion de schémas. Soit $Z \subset Z'$
le voisinage infinitésimal du premier ordre de $Z$ dans $X$.
Alors le diagramme
$$
\xymatrix{
Z \ar[r] \ar[d] & Z' \ar[d] \\
Z \ar[r] & X
}
$$
induit un homomorphisme de faisceaux conormaux $\mathcal{C}_{Z/X} \to \mathcal{C}_{Z/Z'}$ d'après
Morphismes, lemme \ref{morphisms-lemma-conormal-functorial}.
Cet homomorphisme est un isomorphisme.
\end{lemma}

\begin{proof}
Cela résulte immédiatement de la construction de $Z'$ ci-dessus.
\end{proof}

\section{Morphismes formellement non ramifiés}
\label{more-morphisms-section-formally-unramified}

\noindent
Rappelons qu'un homomorphisme d'anneaux $R \to A$ est dit {\it formellement non ramifié}
(voir Algèbre, définition \ref{algebra-definition-formally-unramified})
si, pour tout diagramme commutatif de flèches pleines
$$
\xymatrix{
A \ar[r] \ar@{-->}[rd] & B/I \\
R \ar[r] \ar[u] & B \ar[u]
}
$$
où $I \subset B$ est un idéal de carré nul, il existe au plus une flèche
pointillée qui rende le diagramme commutatif. Cela conduit
à l'analogue suivant pour les morphismes de schémas.

\begin{definition}
\label{more-morphisms-definition-formally-unramified}
Soit $f : X \to S$ un morphisme de schémas.
On dit que $f$ est {\it formellement non ramifié} si, pour tout diagramme commutatif de flèches pleines
$$
\xymatrix{
X \ar[d]_f & T \ar[d]^i \ar[l] \\
S & T' \ar[l] \ar@{-->}[lu]
}
$$
où $T \subset T'$ est un épaississement du premier ordre de schémas affines sur $S$,
il existe au plus une flèche pointillée rendant le diagramme commutatif.
\end{definition}

\noindent
On prouve d'abord quelques lemmes formels, c'est-à-dire des lemmes dont la démonstration
consiste à tracer les diagrammes correspondants.

\begin{lemma}
\label{more-morphisms-lemma-formally-unramified-not-affine}
Si $f : X \to S$ est un morphisme formellement non ramifié, alors, pour
tout diagramme commutatif de flèches pleines
$$
\xymatrix{
X \ar[d]_f & T \ar[d]^i \ar[l] \\
S & T' \ar[l] \ar@{-->}[lu]
}
$$
où $T \subset T'$ est un épaississement du premier ordre de schémas sur $S$,
il existe au plus une flèche pointillée rendant le diagramme commutatif.
Autrement dit, dans
la définition \ref{more-morphisms-definition-formally-unramified},
on peut omettre la condition que $T$ soit affine.
\end{lemma}

\begin{proof}
Cela vient de ce qu'un morphisme est déterminé par ses restrictions
aux ouverts affines.
\end{proof}

\begin{lemma}
\label{more-morphisms-lemma-composition-formally-unramified}
Un composé de morphismes formellement non ramifiés est formellement non ramifié.
\end{lemma}

\begin{proof}
C'est formel.
\end{proof}

\begin{lemma}
\label{more-morphisms-lemma-base-change-formally-unramified}
Tout changement de base d'un morphisme formellement non ramifié est formellement non ramifié.
\end{lemma}

\begin{proof}
C'est formel.
\end{proof}

\begin{lemma}
\label{more-morphisms-lemma-formally-unramified-on-opens}
Soit $f : X \to S$ un morphisme de schémas.
Soient $U \subset X$ et $V \subset S$ des ouverts tels que
$f(U) \subset V$. Si $f$ est formellement non ramifié, il en va de même de $f|_U : U \to V$.
\end{lemma}

\begin{proof}
Considérons un diagramme de flèches pleines
$$
\xymatrix{
U \ar[d]_{f|_U} & T \ar[d]^i \ar[l]^a \\
V & T' \ar[l] \ar@{-->}[lu]
}
$$
comme dans la définition \ref{more-morphisms-definition-formally-unramified}. Si $f$ est formellement
non ramifié, il existe au plus un
$S$-morphisme $a' : T' \to X$ tel que $a'|_T = a$.
Il existe donc, a fortiori, au plus un tel morphisme à valeurs dans $U$.
\end{proof}

\begin{lemma}
\label{more-morphisms-lemma-affine-formally-unramified}
Soit $f : X \to S$ un morphisme de schémas.
Supposons que $X$ et $S$ soient affines.
Alors $f$ est formellement non ramifié si et seulement si
$\mathcal{O}_S(S) \to \mathcal{O}_X(X)$ est un homomorphisme d'anneaux
formellement non ramifié.
\end{lemma}

\begin{proof}
Cela résulte immédiatement des définitions
(définition \ref{more-morphisms-definition-formally-unramified} et
Algèbre, définition \ref{algebra-definition-formally-unramified}),
par l'équivalence entre les catégories des anneaux et des schémas affines,
voir
Schémas, lemme \ref{schemes-lemma-category-affine-schemes}.
\end{proof}

\noindent
Voici une caractérisation en termes du faisceau des différentielles.

\begin{lemma}
\label{more-morphisms-lemma-formally-unramified-differentials}
Soit $f : X \to S$ un morphisme de schémas.
Alors $f$ est formellement non ramifié si et seulement si $\Omega_{X/S} = 0$.
\end{lemma}

\begin{proof}
Rappelons quelques arguments de la démonstration de
Morphismes, lemme \ref{morphisms-lemma-differentials-affine}.
Soit $W \subset X \times_S X$ un ouvert tel que
$\Delta : X \to X \times_S X$ induise une immersion fermée dans $W$.
Soit $\mathcal{J} \subset \mathcal{O}_W$ le faisceau d'idéaux de cette
immersion fermée. Soit $X' \subset W$ le sous-schéma fermé
défini par le faisceau quasi-cohérent d'idéaux $\mathcal{J}^2$.
Considérons les deux morphismes $p_1, p_2 : X' \to X$ induits par
les deux projections $X \times_S X \to X$.
Remarquons que $p_1$ et $p_2$ coïncident après composition avec $\Delta : X \to X'$
et que $X \to X'$ est une immersion fermée définie par un idéal
de carré nul. De plus, on dispose d'une suite exacte courte
$$
0 \to \mathcal{J}/\mathcal{J}^2 \to \mathcal{O}_{X'} \to \mathcal{O}_X \to 0
$$
et $\Omega_{X/S} = \mathcal{J}/\mathcal{J}^2$. De plus,
$\mathcal{J}/\mathcal{J}^2$ est engendré par les sections
locales $p_1^\sharp(f) - p_2^\sharp(f)$, où $f$ est une section locale de
$\mathcal{O}_X$.

\medskip\noindent
Supposons que $f : X \to S$ soit formellement non ramifié.
Par hypothèse, cela signifie que $p_1 = p_2$ après restriction à tout
ouvert affine $T' \subset X'$. Donc $p_1 = p_2$. D'après ce qui précède,
on en déduit que $\Omega_{X/S} = \mathcal{J}/\mathcal{J}^2 = 0$.

\medskip\noindent
Réciproquement, supposons que $\Omega_{X/S} = 0$. Alors $X' = X$. Prenons deux
morphismes $f'_1, f'_2 : T' \to X$ pouvant servir de flèches pointillées dans
le diagramme de
la définition \ref{more-morphisms-definition-formally-unramified}.
On obtient ainsi un morphisme $(f'_1, f'_2) : T' \to X \times_S X$.
Puisque $f'_1|_T = f'_2|_T$ et $|T| =|T'|$, on voit que l'image de $T'$
par $(f'_1, f'_2)$ est contenue dans l'ouvert $W$ choisi ci-dessus. Comme
$(f'_1, f'_2)(T) \subset \Delta(X)$ et que $T$ est défini par un idéal
de carré nul dans $T'$, on voit que $(f'_1, f'_2)$ se factorise par $X'$.
Puisque $X' = X$, on en déduit que $f_1' = f'_2$, comme voulu.
\end{proof}

\begin{lemma}
\label{more-morphisms-lemma-formally-unramified}
Soit $f : X \to Y$ un morphisme de schémas. Les conditions suivantes sont équivalentes :
\begin{enumerate}
\item $f$ est formellement non ramifié,
\item pour tout $x \in X$, il existe des ouverts $x \in U \subset X$ et
$f(x) \in V \subset Y$ avec $f(U) \subset V$ tels que
$f|_U : U \to V$ soit formellement non ramifié,
\item pour tout couple d'ouverts affines $U \subset X$ et $V \subset Y$
avec $f(U) \subset V$, l'homomorphisme d'anneaux $\mathcal{O}_Y(V) \to \mathcal{O}_X(U)$
est formellement non ramifié, et
\item il existe un recouvrement ouvert affine $Y = \bigcup V_j$ et,
pour chaque $j$, un recouvrement ouvert affine $f^{-1}(V_j) = \bigcup U_{ji}$
tels que $\mathcal{O}_Y(V) \to \mathcal{O}_X(U)$ soit un homomorphisme d'anneaux
formellement non ramifié pour tous $j$ et $i$.
\end{enumerate}
\end{lemma}

\begin{proof}
En combinant les lemmes \ref{more-morphisms-lemma-formally-unramified-on-opens} et
\ref{more-morphisms-lemma-affine-formally-unramified}, on voit que (1) $\Rightarrow$ (3).
L'implication (3) $\Rightarrow$ (4) est immédiate.
On a (4) $\Rightarrow$ (2) d'après le lemme \ref{more-morphisms-lemma-affine-formally-unramified}.
Si (2) est vérifiée, alors $\Omega_{X/S}$ s'annule au voisinage
de chaque point d'après le lemme \ref{more-morphisms-lemma-formally-unramified-differentials}
(on utilise aussi Morphismes, lemme
\ref{morphisms-lemma-differentials-restrict-open}) ; le
lemme \ref{more-morphisms-lemma-formally-unramified-differentials} montre alors que $f$
est formellement non ramifié.
\end{proof}

\begin{lemma}
\label{more-morphisms-lemma-unramified-formally-unramified}
\begin{slogan}
Les morphismes non ramifiés sont exactement les morphismes formellement non ramifiés
qui sont localement de type fini.
\end{slogan}
Soit $f : X \to S$ un morphisme de schémas.
Les conditions suivantes sont équivalentes :
\begin{enumerate}
\item Le morphisme $f$ est non ramifié (resp.\ G-non ramifié), et
\item le morphisme $f$ est localement de type fini (resp.\ localement de
présentation finie) et formellement non ramifié.
\end{enumerate}
\end{lemma}

\begin{proof}
Utiliser le lemme \ref{more-morphisms-lemma-formally-unramified-differentials} et
Morphismes, lemme \ref{morphisms-lemma-unramified-omega-zero}.
\end{proof}










\section{Épaississements universels du premier ordre}
\label{more-morphisms-section-universal-thickening}

\noindent
Soit $h : Z \to X$ un morphisme de schémas. Un {\it épaississement universel du premier
ordre} de $Z$ sur $X$ est un épaississement du premier ordre $Z \subset Z'$
sur $X$ tel que, pour tout épaississement du premier ordre $T \subset T'$
sur $X$ et tout diagramme commutatif de flèches pleines
$$
\xymatrix{
& Z \ar[ld] & & T \ar[rd] \ar[ll]^a \\
Z' \ar[rrd] & & & & T' \ar@{..>}[llll]_{a'} \ar[lld]^b \\
 & & X
}
$$
il existe une unique flèche pointillée rendant le diagramme commutatif.
Remarquons que, dans cette situation, $(a, a') : (T \subset T') \to (Z \subset Z')$
est un morphisme d'épaississements sur $X$. Ainsi, s'il existe un épaississement universel
du premier ordre, il est unique à isomorphisme unique près.
En général, un épaississement universel du premier ordre
n'existe pas, mais il en existe un si $h$ est formellement non ramifié.

\begin{lemma}
\label{more-morphisms-lemma-universal-thickening}
Soit $h : Z \to X$ un morphisme de schémas formellement non ramifié.
Il existe un épaississement universel du premier ordre $Z \subset Z'$ de
$Z$ sur $X$.
\end{lemma}

\begin{proof}
Au cours de cette démonstration, on dira que $Z \subset Z'$ est un épaississement universel du premier
ordre de $Z$ sur $X$ s'il satisfait à la condition du lemme.
On construira l'épaississement universel du premier ordre $Z \subset Z'$ sur $X$
par recollement, en partant du cas affine, donné par
Algèbre, lemme \ref{algebra-lemma-universal-thickening}.
Commençons par quelques remarques générales.

\medskip\noindent
S'il existe un épaississement universel du premier ordre de $Z$ sur $X$, il est unique
à isomorphisme unique près. De plus, supposons que $V \subset Z$ et
$U \subset X$ soient des sous-schémas ouverts tels que $h(V) \subset U$. Soit
$Z \subset Z'$ un épaississement universel du premier ordre de $Z$ sur $X$.
Soit $V' \subset Z'$ le sous-schéma ouvert tel que $V = Z \cap V'$.
On affirme alors que $V \subset V'$ est l'épaississement universel du premier ordre de
$V$ sur $U$. En effet, supposons donné un diagramme
$$
\xymatrix{
V \ar[d]_h & T \ar[l]^a \ar[d] \\
U & T' \ar[l]_b
}
$$
où $T \subset T'$ est un épaississement du premier ordre sur $U$. Par la propriété
universelle de $Z'$, on obtient $(a, a') : (T \subset T') \to (Z \subset Z')$.
Mais, puisqu'on a l'égalité $|T| = |T'|$ des espaces topologiques sous-jacents,
on voit que $a'(T') \subset V'$. On peut donc considérer $(a, a')$
comme un morphisme d'épaississements $(a, a') : (T \subset T') \to (V \subset V')$
sur $U$. L'unicité est également claire. De façon tout à fait analogue, on prouve
que, si $h(Z) \subset U$ et si $Z \subset Z'$ est un épaississement universel du premier
ordre sur $U$, alors $Z \subset Z'$ est un épaississement universel du premier ordre
sur $X$.

\medskip\noindent
Avant de recoller les morceaux affines, montrons que le lemme est vrai si
$Z$ et $X$ sont affines. Écrivons $X = \Spec(R)$ et $Z = \Spec(S)$. D'après
Algèbre, lemme \ref{algebra-lemma-universal-thickening},
il existe un épaississement du premier ordre $Z \subset Z'$ sur $X$
qui possède la propriété universelle du lemme pour les diagrammes
$$
\xymatrix{
Z \ar[d]_h & T \ar[l]^a \ar[d] \\
X & T' \ar[l]_b
}
$$
où $T, T'$ sont affines. Pour un diagramme quelconque, on peut choisir un recouvrement
ouvert affine $T' = \bigcup T'_i$ et l'on obtient des morphismes
$a'_i : T'_i \to Z'$ sur $X$ tels que $a'_i|_{T_i} = a|_{T_i}$.
Par unicité, on voit que $a'_i$ et $a'_j$ coïncident sur tout ouvert affine
de $T'_i \cap T'_j$. Les morphismes $a'_i$ se recollent donc en un morphisme global
$a' : T' \to Z'$ sur $X$, comme voulu. Le lemme est ainsi prouvé si $X$ et $Z$
sont affines.

\medskip\noindent
Choisissons un recouvrement ouvert affine $Z = \bigcup Z_i$ tel que chaque $Z_i$
s'envoie dans un ouvert affine $U_i$ de $X$. D'après
le lemme \ref{more-morphisms-lemma-formally-unramified-on-opens},
les morphismes $Z_i \to U_i$ sont formellement non ramifiés.
Le cas affine fournit donc des épaississements universels du premier ordre
$Z_i \subset Z_i'$ sur $U_i$. D'après les remarques générales ci-dessus,
$Z_i \subset Z_i'$ est aussi un épaississement universel du premier ordre de
$Z_i$ sur $X$. Soit $Z'_{i, j} \subset Z'_i$ le sous-schéma ouvert
tel que $Z_i \cap Z_j = Z'_{i, j} \cap Z_i$. D'après les remarques générales,
on voit que $Z'_{i, j}$ et $Z'_{j, i}$ sont tous deux des épaississements universels du premier
ordre de $Z_i \cap Z_j$ sur $X$. Ainsi, d'après
la première de ces remarques, il existe un isomorphisme canonique
$\varphi_{ij} : Z'_{i, j} \to Z'_{j, i}$ induisant l'identité sur
$Z_i \cap Z_j$. On affirme que ces morphismes vérifient la condition de cocycle de
Schémas, section \ref{schemes-section-glueing-schemes}.
(Vérification omise. Indication : utiliser le fait que $Z'_{i, j} \cap Z'_{i, k}$ est
l'épaississement universel du premier ordre de $Z_i \cap Z_j \cap Z_k$, ce qui le détermine
à isomorphisme unique près d'après ce qui précède.)
On peut donc utiliser les résultats de
Schémas, section \ref{schemes-section-glueing-schemes},
pour obtenir un épaississement du premier ordre $Z \subset Z'$ sur $X$ tel que
le sous-schéma ouvert $Z'_i \subset Z'$ vérifiant $Z_i = Z'_i \cap Z$
soit un épaississement universel du premier ordre de $Z_i$ sur $X$.

\medskip\noindent
Cela entraîne formellement que $Z'$ est un épaississement universel du premier
ordre de $Z$ sur $X$. En effet, on dispose de la propriété universelle pour tout
diagramme
$$
\xymatrix{
Z \ar[d]_h & T \ar[l]^a \ar[d] \\
X & T' \ar[l]_b
}
$$
où $a(T)$ est contenu dans l'un des $Z_i$. Pour un diagramme quelconque, on peut
choisir un recouvrement ouvert $T' = \bigcup T'_i$ tel que $a(T_i) \subset Z_i$.
On obtient des morphismes $a'_i : T'_i \to Z'$ sur $X$ tels que
$a'_i|_{T_i} = a|_{T_i}$. On voit que $a'_i$ et $a'_j$ coïncident nécessairement
sur $T'_i \cap T'_j$, car $a'_i|_{T'_i \cap T'_j}$ et
$a'_j|_{T'_i \cap T'_j}$ résolvent tous deux le problème de morphisme à valeurs dans
l'épaississement universel du premier ordre $Z'_i \cap Z'_j$ de $Z_i \cap Z_j$ sur $X$.
Les morphismes $a'_i$ se recollent donc en un morphisme global
$a' : T' \to Z'$ sur $X$, comme voulu. Cela achève la démonstration.
\end{proof}

\begin{definition}
\label{more-morphisms-definition-universal-thickening}
Soit $h : Z \to X$ un morphisme de schémas formellement non ramifié.
\begin{enumerate}
\item L'{\it épaississement universel du premier ordre} de $Z$ sur $X$
est l'épaississement $Z \subset Z'$ construit dans
le lemme \ref{more-morphisms-lemma-universal-thickening}.
\item Le {\it faisceau conormal de $Z$ sur $X$} est le faisceau conormal
de $Z$ dans son épaississement universel du premier ordre $Z'$ sur $X$.
\end{enumerate}
Dans cette situation, on note souvent le faisceau conormal $\mathcal{C}_{Z/X}$.
\end{definition}

\noindent
On voit ainsi qu'il existe une suite exacte courte de faisceaux
$$
0 \to \mathcal{C}_{Z/X} \to \mathcal{O}_{Z'} \to \mathcal{O}_Z \to 0
$$
sur $Z$.
Le lemme suivant prouve que cette définition est compatible
avec celle donnée dans le cas où $Z \to X$ est une immersion.

\begin{lemma}
\label{more-morphisms-lemma-immersion-universal-thickening}
Soit $i : Z \to X$ une immersion de schémas. Alors
\begin{enumerate}
\item $i$ est formellement non ramifié,
\item l'épaississement universel du premier ordre de $Z$ sur $X$ est le voisinage
infinitésimal du premier ordre de $Z$ dans $X$ de
la définition \ref{more-morphisms-definition-first-order-infinitesimal-neighbourhood}, et
\item le faisceau conormal de $i$ au sens de
Morphismes, définition \ref{morphisms-definition-conormal-sheaf},
coïncide avec le faisceau conormal de $i$ au sens de
la définition \ref{more-morphisms-definition-universal-thickening}.
\end{enumerate}
\end{lemma}

\begin{proof}
D'après
Morphismes, lemmes \ref{morphisms-lemma-open-immersion-unramified} et
\ref{morphisms-lemma-closed-immersion-unramified},
une immersion est non ramifiée, donc formellement non ramifiée d'après
le lemme \ref{more-morphisms-lemma-unramified-formally-unramified}.
Les autres assertions résultent de la combinaison
des lemmes \ref{more-morphisms-lemma-first-order-infinitesimal-neighbourhood} et
\ref{more-morphisms-lemma-infinitesimal-neighbourhood-conormal}
avec les définitions.
\end{proof}

\begin{lemma}
\label{more-morphisms-lemma-universal-thickening-unramified}
Soit $Z \to X$ un morphisme de schémas formellement non ramifié.
Alors l'épaississement universel du premier ordre $Z'$ est formellement
non ramifié sur $X$.
\end{lemma}

\begin{proof}
Il y a deux démonstrations. La première consiste à montrer que $\Omega_{Z'/X} = 0$
en travaillant localement sur des ouverts affines et en appliquant
Algèbre, lemme \ref{algebra-lemma-differentials-universal-thickening}.
Alors
le lemme \ref{more-morphisms-lemma-formally-unramified-differentials}
donne le résultat voulu.
La seconde est l'argument direct suivant.

\medskip\noindent
Soit $T \subset T'$ un épaississement du premier ordre. Soit
$$
\xymatrix{
Z' \ar[d] & T \ar[l]^c \ar[d] \\
X & T' \ar[l] \ar[lu]^{a, b}
}
$$
un diagramme commutatif. Considérons deux morphismes $a, b : T' \to Z'$
s'inscrivant dans ce diagramme. Posons $T_0 = c^{-1}(Z) \subset T$ et
$T'_a = a^{-1}(Z)$ (au sens schématique).
Puisque $Z'$ est un épaississement du premier ordre de $Z$, on voit que $T'$
est un épaississement du premier ordre de $T'_a$. De plus, puisque $c = a|_T$, on voit que
$T_0 = T \cap T'_a$ (au sens schématique). Comme $T'$ est un épaississement du premier
ordre de $T$, il en résulte que $T'_a$
est un épaississement du premier ordre de $T_0$. Or $a|_{T'_a}$ et $b|_{T'_a}$
sont des morphismes de $T'_a$ dans $Z'$ sur $X$ qui coïncident sur $T_0$ en tant que
morphismes à valeurs dans $Z$. Par la propriété universelle de $Z'$, on en déduit donc que
$a|_{T'_a} = b|_{T'_a}$. Ainsi, $a$ et $b$ sont des morphismes définis sur
l'épaississement du premier ordre $T'$ de $T'_a$ dont les restrictions à
$T'_a$ coïncident comme morphismes à valeurs dans $Z$. En utilisant encore la propriété universelle de
$Z'$, on en déduit donc que $a = b$. Autrement dit, la propriété définissant
un morphisme formellement non ramifié est vérifiée par $Z' \to X$, comme voulu.
\end{proof}

\begin{lemma}
\label{more-morphisms-lemma-universal-thickening-functorial}
Considérons un diagramme commutatif de schémas
$$
\xymatrix{
Z \ar[r]_h \ar[d]_f & X \ar[d]^g \\
W \ar[r]^{h'} & Y
}
$$
où $h$ et $h'$ sont formellement non ramifiés. Soit $Z \subset Z'$ l'épaississement universel
du premier ordre de $Z$ sur $X$. Soit $W \subset W'$ l'épaississement universel
du premier ordre de $W$ sur $Y$. Il existe un morphisme canonique
$(f, f') : (Z, Z') \to (W, W')$ d'épaississements sur $Y$ qui s'inscrit dans
le diagramme commutatif suivant
$$
\xymatrix{
& & & Z' \ar[ld] \ar[d]^{f'} \\
Z \ar[rr] \ar[d]_f \ar[rrru] & & X \ar[d] & W' \ar[ld] \\
W \ar[rrru]|!{[rr];[rruu]}\hole \ar[rr] & & Y
}
$$
En particulier, le morphisme $(f, f')$ d'épaississements induit un homomorphisme
de faisceaux conormaux $f^*\mathcal{C}_{W/Y} \to \mathcal{C}_{Z/X}$.
\end{lemma}

\begin{proof}
La première assertion résulte immédiatement de la propriété universelle de $W'$.
L'homomorphisme induit sur les faisceaux conormaux est celui de
Morphismes, lemme \ref{morphisms-lemma-conormal-functorial},
appliqué à $(Z \subset Z') \to (W \subset W')$.
\end{proof}

\begin{lemma}
\label{more-morphisms-lemma-universal-thickening-fibre-product}
Soit
$$
\xymatrix{
Z \ar[r]_h \ar[d]_f & X \ar[d]^g \\
W \ar[r]^{h'} & Y
}
$$
un diagramme cartésien dans la catégorie des schémas, où
$h'$ est formellement non ramifié. Alors $h$ est formellement non ramifié et, si
$W \subset W'$ est l'épaississement universel du premier ordre de $W$ sur $Y$,
alors $Z = X \times_Y W \subset X \times_Y W'$ est l'épaississement universel
du premier ordre de $Z$ sur $X$. En particulier, l'homomorphisme canonique
$f^*\mathcal{C}_{W/Y} \to \mathcal{C}_{Z/X}$ du
lemme \ref{more-morphisms-lemma-universal-thickening-functorial}
est surjectif.
\end{lemma}

\begin{proof}
Le morphisme $h$ est formellement non ramifié d'après
le lemme \ref{more-morphisms-lemma-base-change-formally-unramified}.
Il est clair que $X \times_Y W'$ est un épaississement du premier ordre.
On vérifie directement qu'il possède la propriété universelle,
car $W'$ la possède (par la propriété universelle des
produits fibrés). Voir
Morphismes, lemme \ref{morphisms-lemma-conormal-functorial-flat},
pour en déduire la surjectivité de l'homomorphisme de faisceaux conormaux.
\end{proof}

\begin{lemma}
\label{more-morphisms-lemma-universal-thickening-fibre-product-flat}
Soit
$$
\xymatrix{
Z \ar[r]_h \ar[d]_f & X \ar[d]^g \\
W \ar[r]^{h'} & Y
}
$$
un diagramme cartésien dans la catégorie des schémas, où
$h'$ est formellement non ramifié et $g$ plat. Dans ce cas, le morphisme
correspondant $Z' \to W'$ des épaississements universels du premier ordre est plat, et
$f^*\mathcal{C}_{W/Y} \to \mathcal{C}_{Z/X}$ est un isomorphisme.
\end{lemma}

\begin{proof}
La platitude est préservée par changement de base, voir
Morphismes, lemme \ref{morphisms-lemma-base-change-flat}.
La première assertion résulte donc de la description de
$W'$ dans le lemme \ref{more-morphisms-lemma-universal-thickening-fibre-product}.
Il est clair que $X \times_Y W'$ est un épaississement du premier ordre.
On vérifie directement qu'il possède la propriété universelle,
car $W'$ la possède (par la propriété universelle des
produits fibrés). Voir
Morphismes, lemme \ref{morphisms-lemma-conormal-functorial-flat},
pour en déduire que l'homomorphisme de faisceaux conormaux est un isomorphisme.
\end{proof}

\begin{lemma}
\label{more-morphisms-lemma-universal-thickening-localize}
La formation des épaississements universels du premier ordre commute à la restriction aux ouverts.
Plus précisément, soit $h : Z \to X$ un morphisme de schémas formellement non ramifié.
Soient $V \subset Z$, $U \subset X$ des ouverts tels que $h(V) \subset U$.
Soit $Z'$ l'épaississement universel du premier ordre de $Z$ sur $X$.
Alors $h|_V : V \to U$ est formellement non ramifié et l'épaississement universel du premier
ordre de $V$ sur $U$ est le sous-schéma ouvert $V' \subset Z'$
tel que $V = Z \cap V'$. En particulier,
$\mathcal{C}_{Z/X}|_V = \mathcal{C}_{V/U}$.
\end{lemma}

\begin{proof}
La première assertion est
le lemme \ref{more-morphisms-lemma-formally-unramified-on-opens}.
La compatibilité des épaississements universels peut se déduire de la démonstration du
lemme \ref{more-morphisms-lemma-universal-thickening},
ou d'Algèbre,
lemme \ref{algebra-lemma-universal-thickening-localize},
ou encore du
lemme \ref{more-morphisms-lemma-universal-thickening-fibre-product-flat}.
\end{proof}

\begin{lemma}
\label{more-morphisms-lemma-differentials-universally-unramified}
Soit $h : Z \to X$ un morphisme formellement non ramifié de schémas sur $S$.
Soit $Z \subset Z'$ l'épaississement universel du premier ordre de $Z$
sur $X$, de morphisme structural $h' : Z' \to X$. L'homomorphisme canonique
$$
c_{h'} : (h')^*\Omega_{X/S} \longrightarrow \Omega_{Z'/S}
$$
induit un isomorphisme
$h^*\Omega_{X/S} \to \Omega_{Z'/S} \otimes \mathcal{O}_Z$.
\end{lemma}

\begin{proof}
L'homomorphisme $c_{h'}$ est celui défini dans
Morphismes, lemme \ref{morphisms-lemma-functoriality-differentials}.
Si $i : Z \to Z'$ est l'immersion fermée donnée, alors
$i^*c_{h'}$ est un homomorphisme
$h^*\Omega_{X/S} \to \Omega_{Z'/S} \otimes \mathcal{O}_Z$.
Pour vérifier que c'est un isomorphisme, on se ramène au cas affine
par localisation, voir
le lemme \ref{more-morphisms-lemma-universal-thickening-localize}
et
Morphismes, lemme \ref{morphisms-lemma-differentials-restrict-open}.
Dans ce cas, le résultat est
Algèbre, lemme \ref{algebra-lemma-differentials-universal-thickening}.
\end{proof}

\begin{lemma}
\label{more-morphisms-lemma-universally-unramified-differentials-sequence}
Soit $h : Z \to X$ un morphisme formellement non ramifié de schémas sur $S$.
Il existe une suite exacte canonique
$$
\mathcal{C}_{Z/X} \to h^*\Omega_{X/S} \to \Omega_{Z/S} \to 0.
$$
La première flèche est induite par $\text{d}_{Z'/S}$, où
$Z'$ est le voisinage universel du premier ordre de $Z$ sur $X$.
\end{lemma}

\begin{proof}
On sait qu'il existe une suite exacte canonique
$$
\mathcal{C}_{Z/Z'} \to
\Omega_{Z'/S} \otimes \mathcal{O}_Z \to
\Omega_{Z/S} \to 0.
$$
voir
Morphismes, lemme \ref{morphisms-lemma-differentials-relative-immersion}.
Le résultat s'en déduit donc en appliquant
le lemme \ref{more-morphisms-lemma-differentials-universally-unramified}.
\end{proof}

\begin{lemma}
\label{more-morphisms-lemma-two-unramified-morphisms}
Soit
$$
\xymatrix{
Z \ar[r]_i \ar[rd]_j & X \ar[d] \\
& Y
}
$$
un diagramme commutatif de schémas où $i$ et $j$ sont formellement
non ramifiés. Il existe alors une suite exacte canonique
$$
\mathcal{C}_{Z/Y} \to
\mathcal{C}_{Z/X} \to
i^*\Omega_{X/Y} \to 0
$$
dont la première flèche provient du
lemme \ref{more-morphisms-lemma-universal-thickening-functorial}
et la seconde du
lemme \ref{more-morphisms-lemma-universally-unramified-differentials-sequence}.
\end{lemma}

\begin{proof}
Notons $Z \to Z'$ l'épaississement universel du premier ordre de $Z$ sur $X$.
Notons $Z \to Z''$ l'épaississement universel du premier ordre de $Z$ sur $Y$.
D'après
le lemme \ref{more-morphisms-lemma-universally-unramified-differentials-sequence},
il existe un morphisme canonique $Z' \to Z''$ donnant un diagramme
commutatif
$$
\xymatrix{
Z \ar[r]_{i'} \ar[rd]_{j'} & Z' \ar[r] \ar[d] & X \ar[d] \\
& Z'' \ar[r] & Y
}
$$
Appliquons
Morphismes, lemme \ref{morphisms-lemma-two-immersions},
au triangle de gauche pour obtenir une suite exacte
$$
\mathcal{C}_{Z/Z''} \to
\mathcal{C}_{Z/Z'} \to
(i')^*\Omega_{Z'/Z''} \to 0
$$
Comme $Z''$ est formellement non ramifié sur $Y$ (voir
le lemme \ref{more-morphisms-lemma-universal-thickening-unramified}),
on a
$\Omega_{Z'/Z''} = \Omega_{Z/Y}$ (en combinant
le lemme \ref{more-morphisms-lemma-formally-unramified-differentials}
et
Morphismes, lemme \ref{morphisms-lemma-triangle-differentials}).
On a alors $(i')^*\Omega_{Z'/Y} = i^*\Omega_{X/Y}$ d'après
le lemme \ref{more-morphisms-lemma-differentials-universally-unramified}.
\end{proof}

\begin{lemma}
\label{more-morphisms-lemma-transitivity-conormal}
Soient $Z \to Y \to X$ des morphismes de schémas formellement non ramifiés.
\begin{enumerate}
\item Si $Z \subset Z'$ est l'épaississement universel du premier ordre de $Z$
sur $X$ et si $Y \subset Y'$ est l'épaississement universel du premier ordre de $Y$
sur $X$, alors il existe un morphisme $Z' \to Y'$ et $Y \times_{Y'} Z'$ est
l'épaississement universel du premier ordre de $Z$ sur $Y$.
\item Il existe une suite exacte canonique
$$
i^*\mathcal{C}_{Y/X} \to
\mathcal{C}_{Z/X} \to
\mathcal{C}_{Z/Y} \to 0
$$
dont les homomorphismes proviennent du
lemme \ref{more-morphisms-lemma-universal-thickening-functorial}
et où $i : Z \to Y$ est le premier morphisme.
\end{enumerate}
\end{lemma}

\begin{proof}
Le morphisme $h : Z' \to Y'$ de (1) provient du
lemme \ref{more-morphisms-lemma-universal-thickening-functorial}.
Le fait que $Y \times_{Y'} Z'$ soit l'épaississement universel du premier
ordre de $Z$ sur $Y$ résulte immédiatement des propriétés universelles
de $Z'$ et de $Y'$. D'après
Morphismes, lemme \ref{morphisms-lemma-transitivity-conormal},
on dispose d'une suite exacte
$$
(i')^*\mathcal{C}_{Y \times_{Y'} Z'/Z'} \to
\mathcal{C}_{Z/Z'} \to
\mathcal{C}_{Z/Y \times_{Y'} Z'} \to 0
$$
où $i' : Z \to Y \times_{Y'} Z'$ est le morphisme donné. D'après
Morphismes, lemme \ref{morphisms-lemma-conormal-functorial-flat},
il existe une surjection
$h^*\mathcal{C}_{Y/Y'} \to \mathcal{C}_{Y \times_{Y'} Z'/Z'}$.
En combinant cela avec les égalités
$\mathcal{C}_{Y/Y'} = \mathcal{C}_{Y/X}$,
$\mathcal{C}_{Z/Z'} = \mathcal{C}_{Z/X}$ et
$\mathcal{C}_{Z/Y \times_{Y'} Z'} = \mathcal{C}_{Z/Y}$,
on obtient le lemme.
\end{proof}












\section{Morphismes formellement étales}
\label{more-morphisms-section-formally-etale}

\noindent
Rappelons qu'un homomorphisme d'anneaux $R \to A$ est dit {\it formellement étale}
(voir Algèbre, définition \ref{algebra-definition-formally-etale})
si, pour tout diagramme commutatif de flèches pleines
$$
\xymatrix{
A \ar[r] \ar@{-->}[rd] & B/I \\
R \ar[r] \ar[u] & B \ar[u]
}
$$
où $I \subset B$ est un idéal de carré nul, il existe
exactement une flèche pointillée qui rende le diagramme commutatif. Cela conduit
à l'analogue suivant pour les morphismes de schémas.

\begin{definition}
\label{more-morphisms-definition-formally-etale}
Soit $f : X \to S$ un morphisme de schémas.
On dit que $f$ est {\it formellement étale} si, pour tout diagramme commutatif de flèches pleines
$$
\xymatrix{
X \ar[d]_f & T \ar[d]^i \ar[l] \\
S & T' \ar[l] \ar@{-->}[lu]
}
$$
où $T \subset T'$ est un épaississement du premier ordre de schémas affines sur $S$,
il existe exactement une flèche pointillée rendant le diagramme commutatif.
\end{definition}

\noindent
Il est clair qu'un morphisme formellement étale est formellement non ramifié.
Ainsi, si $f : X \to S$ est formellement étale, alors $\Omega_{X/S}$ est nul, voir
le lemme \ref{more-morphisms-lemma-formally-unramified-differentials}.

\begin{lemma}
\label{more-morphisms-lemma-formally-etale-not-affine}
Si $f : X \to S$ est un morphisme formellement étale, alors, pour
tout diagramme commutatif de flèches pleines
$$
\xymatrix{
X \ar[d]_f & T \ar[d]^i \ar[l] \\
S & T' \ar[l] \ar@{-->}[lu]
}
$$
où $T \subset T'$ est un épaississement du premier ordre de schémas sur $S$,
il existe exactement une flèche pointillée rendant le diagramme commutatif.
Autrement dit, dans
la définition \ref{more-morphisms-definition-formally-etale},
on peut omettre la condition que $T$ soit affine.
\end{lemma}

\begin{proof}
Soit $T' = \bigcup T'_i$ un recouvrement ouvert affine et posons
$T_i = T \cap T'_i$. On obtient alors des morphismes $a'_i : T'_i \to X$ s'inscrivant
dans le diagramme. Par unicité, on voit que $a'_i$ et $a'_j$ coïncident sur
tout sous-schéma ouvert affine de $T'_i \cap T'_j$. Ainsi, $a'_i$ et
$a'_j$ coïncident sur $T'_i \cap T'_j$. On voit donc que les morphismes $a'_i$
se recollent en un morphisme global $a' : T' \to X$. L'unicité de
$a'$ a été établie dans
le lemme \ref{more-morphisms-lemma-formally-unramified-not-affine}.
\end{proof}

\begin{lemma}
\label{more-morphisms-lemma-composition-formally-etale}
Un composé de morphismes formellement étales est formellement étale.
\end{lemma}

\begin{proof}
C'est formel.
\end{proof}

\begin{lemma}
\label{more-morphisms-lemma-base-change-formally-etale}
Tout changement de base d'un morphisme formellement étale est formellement étale.
\end{lemma}

\begin{proof}
C'est formel.
\end{proof}

\begin{lemma}
\label{more-morphisms-lemma-formally-etale-on-opens}
Soit $f : X \to S$ un morphisme de schémas.
Soient $U \subset X$ et $V \subset S$ des sous-schémas ouverts tels que
$f(U) \subset V$. Si $f$ est formellement étale, il en va de même de $f|_U : U \to V$.
\end{lemma}

\begin{proof}
Considérons un diagramme de flèches pleines
$$
\xymatrix{
U \ar[d]_{f|_U} & T \ar[d]^i \ar[l]^a \\
V & T' \ar[l] \ar@{-->}[lu]
}
$$
comme dans la définition \ref{more-morphisms-definition-formally-etale}. Si $f$ est formellement
étale, il existe exactement un $S$-morphisme $a' : T' \to X$
tel que $a'|_T = a$. Puisque $|T'| = |T|$, on en déduit que $a'(T') \subset U$,
ce qui donne l'unique morphisme de $T'$ dans $U$ voulu.
\end{proof}

\begin{lemma}
\label{more-morphisms-lemma-characterize-formally-etale}
Soit $f : X \to S$ un morphisme de schémas.
Les conditions suivantes sont équivalentes :
\begin{enumerate}
\item $f$ est formellement étale,
\item $f$ est formellement non ramifié et l'épaississement universel du premier ordre
de $X$ sur $S$ est égal à $X$,
\item $f$ est formellement non ramifié et $\mathcal{C}_{X/S} = 0$, et
\item $\Omega_{X/S} = 0$ et $\mathcal{C}_{X/S} = 0$.
\end{enumerate}
\end{lemma}

\begin{proof}
En fait, la dernière assertion n'a de sens que parce que $\Omega_{X/S} = 0$
entraîne que $\mathcal{C}_{X/S}$ est défini en vertu du
lemme \ref{more-morphisms-lemma-formally-unramified-differentials}
et de
la définition \ref{more-morphisms-definition-universal-thickening}.
Cela montre aussi clairement que (3) et (4) sont équivalentes.

\medskip\noindent
Chacune des hypothèses (1), (2) et (3) entraîne que $f$ est formellement
non ramifié. On peut donc supposer $f$ formellement non ramifié. L'équivalence
de (1), (2) et (3) résulte de la propriété universelle de l'épaississement universel
du premier ordre $X'$ de $X$ sur $S$ et du fait que
$X = X' \Leftrightarrow \mathcal{C}_{X/S} = 0$, puisque
$\mathcal{C}_{X/S} = \mathcal{C}_{X/X'}$ est, par définition,
le faisceau d'idéaux de $X$ dans $X'$.
\end{proof}

\begin{lemma}
\label{more-morphisms-lemma-unramified-flat-formally-etale}
Un morphisme non ramifié et plat est formellement étale.
\end{lemma}

\begin{proof}
Supposons $X \to S$ non ramifié et plat. Alors $\Delta : X \to X \times_S X$
est une immersion ouverte, voir
Morphismes, lemme \ref{morphisms-lemma-diagonal-unramified-morphism}.
Il faut montrer que $\mathcal{C}_{X/S}$ est nul.
Considérons les deux projections $p, q : X \times_S X \to X$.
Comme $f$ est formellement non ramifié (voir
le lemme \ref{more-morphisms-lemma-unramified-formally-unramified}),
$q$ est formellement non ramifié (voir
le lemme \ref{more-morphisms-lemma-base-change-formally-unramified}).
Comme $f$ est plat, $p$ est plat, voir
Morphismes, lemme \ref{morphisms-lemma-base-change-flat}.
Ainsi, $p^*\mathcal{C}_{X/S} = \mathcal{C}_q$ d'après
le lemme \ref{more-morphisms-lemma-universal-thickening-fibre-product-flat},
où $\mathcal{C}_q$ désigne le faisceau conormal du morphisme formellement
non ramifié $q : X \times_S X \to X$.
Mais $\Delta(X) \subset X \times_S X$ est un sous-schéma ouvert
qui s'envoie isomorphiquement sur $X$ par $q$. D'après
le lemme \ref{more-morphisms-lemma-universal-thickening-localize},
on voit donc que $\mathcal{C}_q|_{\Delta(X)} = \mathcal{C}_{X/X} = 0$.
Autrement dit, l'image inverse de $\mathcal{C}_{X/S}$ sur $X$ par
le morphisme identité est nulle, c'est-à-dire $\mathcal{C}_{X/S} = 0$.
\end{proof}

\begin{lemma}
\label{more-morphisms-lemma-affine-formally-etale}
Soit $f : X \to S$ un morphisme de schémas.
Supposons que $X$ et $S$ soient affines.
Alors $f$ est formellement étale si et seulement si
$\mathcal{O}_S(S) \to \mathcal{O}_X(X)$ est un homomorphisme d'anneaux
formellement étale.
\end{lemma}

\begin{proof}
Cela résulte immédiatement des définitions
(définition \ref{more-morphisms-definition-formally-etale} et
Algèbre, définition \ref{algebra-definition-formally-etale}),
par l'équivalence entre les catégories des anneaux et des schémas affines,
voir
Schémas, lemme \ref{schemes-lemma-category-affine-schemes}.
\end{proof}

\begin{lemma}
\label{more-morphisms-lemma-formally-etale}
Soit $f : X \to Y$ un morphisme de schémas. Les conditions suivantes sont équivalentes :
\begin{enumerate}
\item $f$ est formellement étale,
\item pour tout $x \in X$, il existe des ouverts $x \in U \subset X$ et
$f(x) \in V \subset Y$ avec $f(U) \subset V$ tels que
$f|_U : U \to V$ soit formellement étale,
\item pour tout couple d'ouverts affines $U \subset X$ et $V \subset Y$
avec $f(U) \subset V$, l'homomorphisme d'anneaux $\mathcal{O}_Y(V) \to \mathcal{O}_X(U)$
est formellement étale, et
\item il existe un recouvrement ouvert affine $Y = \bigcup V_j$ et,
pour chaque $j$, un recouvrement ouvert affine $f^{-1}(V_j) = \bigcup U_{ji}$
tels que $\mathcal{O}_Y(V) \to \mathcal{O}_X(U)$ soit un homomorphisme d'anneaux
formellement étale pour tous $j$ et $i$.
\end{enumerate}
\end{lemma}

\begin{proof}
En combinant les lemmes \ref{more-morphisms-lemma-formally-etale-on-opens} et
\ref{more-morphisms-lemma-affine-formally-etale}, on voit que (1) $\Rightarrow$ (3).
L'implication (3) $\Rightarrow$ (4) est immédiate.
On a (4) $\Rightarrow$ (2) d'après le lemme \ref{more-morphisms-lemma-affine-formally-etale}.
Si (2) est vérifiée, alors $\Omega_{X/S}$ et $\mathcal{C}_{X/S}$ s'annulent
au voisinage de chaque point d'après
le lemme \ref{more-morphisms-lemma-characterize-formally-etale}
(on utilise aussi Morphismes, lemme
\ref{morphisms-lemma-differentials-restrict-open}, et
le lemme \ref{more-morphisms-lemma-universal-thickening-localize}) ; le
lemme \ref{more-morphisms-lemma-characterize-formally-etale} montre alors que $f$
est formellement étale.
\end{proof}

\begin{lemma}
\label{more-morphisms-lemma-etale-formally-etale}
Soit $f : X \to S$ un morphisme de schémas.
Les conditions suivantes sont équivalentes :
\begin{enumerate}
\item Le morphisme $f$ est étale, et
\item le morphisme $f$ est localement de présentation finie et
formellement étale.
\end{enumerate}
\end{lemma}

\begin{proof}
Supposons $f$ étale.
Un morphisme étale est localement de présentation finie, plat et non ramifié, voir
Morphismes, section \ref{morphisms-section-etale}.
Ainsi, $f$ est localement de présentation finie et formellement étale, voir
le lemme \ref{more-morphisms-lemma-unramified-flat-formally-etale}.

\medskip\noindent
Réciproquement, supposons que $f$ soit localement de présentation finie et
formellement étale. La propriété d'être étale est locale pour la topologie de Zariski sur
$X$ et $S$, voir
Morphismes, lemme \ref{morphisms-lemma-etale-characterize}.
D'après
le lemme \ref{more-morphisms-lemma-formally-etale-on-opens},
on peut recouvrir $X$ par des ouverts affines $U$ s'envoyant dans des ouverts affines
$V$ tels que $U \to V$ soit formellement étale (et de présentation finie, voir
Morphismes,
lemme \ref{morphisms-lemma-locally-finite-presentation-characterize}).
D'après
le lemme \ref{more-morphisms-lemma-affine-formally-etale},
on voit que les homomorphismes d'anneaux $\mathcal{O}(V) \to \mathcal{O}(U)$ sont
formellement étales (et de présentation finie).
On conclut par
Algèbre, lemme \ref{algebra-lemma-formally-etale-etale}.
(On donnera une autre démonstration de cette implication lors de l'étude
des morphismes formellement lisses.)
\end{proof}














\section{Déformations infinitésimales des morphismes}
\label{more-morphisms-section-action-by-derivations}

\noindent
Dans cette section, on explique comment utiliser une dérivation pour
déformer infinitésimalement un morphisme. Dans toute cette section, on utilise le fait qu'un
faisceau sur un épaississement $X'$ de $X$ peut être considéré comme un faisceau sur $X$.

\begin{lemma}
\label{more-morphisms-lemma-difference-derivation}
Soit $S$ un schéma.
Soient $X \subset X'$ et $Y \subset Y'$ deux épaississements du premier ordre
sur $S$. Soient $(a, a'), (b, b') : (X \subset X') \to (Y \subset Y')$
deux morphismes d'épaississements sur $S$. Supposons que
\begin{enumerate}
\item $a = b$, et
\item les deux homomorphismes $a^*\mathcal{C}_{Y/Y'} \to \mathcal{C}_{X/X'}$
(Morphismes, lemme \ref{morphisms-lemma-conormal-functorial})
soient égaux.
\end{enumerate}
Alors l'application $(a')^\sharp - (b')^\sharp$ se factorise en
$$
\mathcal{O}_{Y'} \to \mathcal{O}_Y \xrightarrow{D}
a_*\mathcal{C}_{X/X'} \to a_*\mathcal{O}_{X'}
$$
où $D$ est une $\mathcal{O}_S$-dérivation.
\end{lemma}

\begin{proof}
Au lieu de travailler sur $Y$, on travaille sur $X$. L'avantage est que le foncteur image inverse
$a^{-1}$ est exact. Grâce à (1) et (2), on obtient un diagramme commutatif
à lignes exactes
$$
\xymatrix{
0 \ar[r] &
\mathcal{C}_{X/X'} \ar[r] &
\mathcal{O}_{X'} \ar[r] &
\mathcal{O}_X \ar[r] & 0 \\
0 \ar[r] &
a^{-1}\mathcal{C}_{Y/Y'} \ar[r] \ar[u] &
a^{-1}\mathcal{O}_{Y'}
\ar[r] \ar@<1ex>[u]^{(a')^\sharp} \ar@<-1ex>[u]_{(b')^\sharp} &
a^{-1}\mathcal{O}_Y \ar[r] \ar[u] & 0
}
$$
Or, dans une telle situation, c'est un fait général que la différence des homomorphismes de
$\mathcal{O}_S$-algèbres $(a')^\sharp$ et $(b')^\sharp$ est une
$\mathcal{O}_S$-dérivation de $a^{-1}\mathcal{O}_Y$ dans $\mathcal{C}_{X/X'}$.
Par adjonction des foncteurs $a^{-1}$ et $a_*$, cela revient
à se donner une $\mathcal{O}_S$-dérivation de
$\mathcal{O}_Y$ dans $a_*\mathcal{C}_{X/X'}$. Certains détails sont omis.
\end{proof}

\noindent
Remarquons que, dans la situation du lemme précédent, on peut écrire
$D$ sous la forme
\begin{equation}
\label{more-morphisms-equation-D}
D = \text{d}_{Y/S} \circ \theta
\end{equation}
où $\theta$ est un homomorphisme $\mathcal{O}_Y$-linéaire
$\theta : \Omega_{Y/S} \to a_*\mathcal{C}_{X/X'}$.
Bien entendu, par une nouvelle adjonction, on peut alors considérer $\theta$ comme un
homomorphisme $\mathcal{O}_X$-linéaire
$\theta : a^*\Omega_{Y/S} \to \mathcal{C}_{X/X'}$.

\begin{lemma}
\label{more-morphisms-lemma-action-by-derivations}
Soit $S$ un schéma.
Soit $(a, a') : (X \subset X') \to (Y \subset Y')$
un morphisme d'épaississements du premier ordre sur $S$.
Soit
$$
\theta : a^*\Omega_{Y/S} \to \mathcal{C}_{X/X'}
$$
un homomorphisme $\mathcal{O}_X$-linéaire. Il existe alors un unique morphisme de couples
$(b, b') : (X \subset X') \to (Y \subset Y')$ tel que
(1) et (2) du
lemme \ref{more-morphisms-lemma-difference-derivation}
soient vérifiées et que la dérivation $D$ et $\theta$ soient liées par
l'équation (\ref{more-morphisms-equation-D}).
\end{lemma}

\begin{proof}
On pose simplement $b = a$ et on définit $(b')^\sharp$ comme l'application
$$
(a')^\sharp + D : a^{-1}\mathcal{O}_{Y'} \to \mathcal{O}_{X'}
$$
où $D$ est donné par l'équation (\ref{more-morphisms-equation-D}). On omet de vérifier
que $(b')^\sharp$ est un homomorphisme de faisceaux de $\mathcal{O}_S$-algèbres et
que (1) et (2) du
lemme \ref{more-morphisms-lemma-difference-derivation}
sont vérifiées. L'équation (\ref{more-morphisms-equation-D}) est satisfaite par construction.
\end{proof}

\begin{remark}
\label{more-morphisms-remark-action-by-derivations}
Les hypothèses et notations sont celles du lemme \ref{more-morphisms-lemma-action-by-derivations}.
L'action d'une section locale $\theta$ sur $a'$ est parfois notée
$\theta \cdot a'$. Cela signifie simplement
que $(a')^\sharp$ et $(\theta \cdot a')^\sharp$ diffèrent d'une dérivation
$D$ liée à $\theta$ par l'équation (\ref{more-morphisms-equation-D}).
\end{remark}

\begin{lemma}
\label{more-morphisms-lemma-sheaf}
Soit $S$ un schéma.
Soient $X \subset X'$ et $Y \subset Y'$ des épaississements du premier ordre
sur $S$. Supposons donnés un morphisme $a : X \to Y$ et un homomorphisme
$A : a^*\mathcal{C}_{Y/Y'} \to \mathcal{C}_{X/X'}$ de
$\mathcal{O}_X$-modules. Pour un sous-schéma ouvert $U' \subset X'$,
considérons les morphismes $a' : U' \to Y'$ tels que
\begin{enumerate}
\item $a'$ soit un morphisme sur $S$,
\item $a'|_U = a|_U$, et
\item l'homomorphisme induit
$a^*\mathcal{C}_{Y/Y'}|_U \to \mathcal{C}_{X/X'}|_U$
soit la restriction de $A$ à $U$.
\end{enumerate}
Ici, $U = X \cap U'$. Alors la règle
\begin{equation}
\label{more-morphisms-equation-sheaf}
U' \mapsto
\{a' : U' \to Y'\text{ vérifiant (1), (2), (3).}\}
\end{equation}
définit un faisceau d'ensembles sur $X'$.
\end{lemma}

\begin{proof}
Notons $\mathcal{F}$ la règle du lemme.
L'application de restriction $\mathcal{F}(U') \to \mathcal{F}(V')$, pour
$V' \subset U' \subset X'$,
de $\mathcal{F}$ est simplement l'application de restriction $a' \mapsto a'|_{V'}$.
Avec cette définition, il est clair que $\mathcal{F}$ est un
faisceau, puisque les morphismes se définissent localement.
\end{proof}

\noindent
Dans le lemme suivant, on identifie les faisceaux sur $X$ à ceux sur un épaississement
quelconque de $X$.

\begin{lemma}
\label{more-morphisms-lemma-action-sheaf}
Les notations et hypothèses sont celles du lemme \ref{more-morphisms-lemma-sheaf}.
Il existe une action du faisceau
$$
\SheafHom_{\mathcal{O}_X}(a^*\Omega_{Y/S}, \mathcal{C}_{X/X'})
$$
sur le faisceau (\ref{more-morphisms-equation-sheaf}). De plus, cette action
est simplement transitive pour tout ouvert $U' \subset X'$ sur lequel le faisceau
(\ref{more-morphisms-equation-sheaf}) possède une section.
\end{lemma}

\begin{proof}
Cela résulte de la combinaison des
lemmes \ref{more-morphisms-lemma-difference-derivation},
\ref{more-morphisms-lemma-action-by-derivations}
et \ref{more-morphisms-lemma-sheaf}.
\end{proof}

\begin{remark}
\label{more-morphisms-remark-special-case}
Un cas particulier des
lemmes \ref{more-morphisms-lemma-difference-derivation},
\ref{more-morphisms-lemma-action-by-derivations},
\ref{more-morphisms-lemma-sheaf} et
\ref{more-morphisms-lemma-action-sheaf}
est celui où $Y = Y'$. Dans ce cas, l'homomorphisme $A$ est toujours nul.
Le faisceau du
lemme \ref{more-morphisms-lemma-sheaf}
est simplement donné par la règle
$$
U' \mapsto
\{a' : U' \to Y\text{ sur }S\text{ avec } a'|_U = a|_U\}
$$
et l'on fait agir sur lui le faisceau
$\SheafHom_{\mathcal{O}_X}(a^*\Omega_{Y/S}, \mathcal{C}_{X/X'})$.
\end{remark}

\begin{remark}
\label{more-morphisms-remark-another-special-case}
Un autre cas particulier des
lemmes \ref{more-morphisms-lemma-difference-derivation},
\ref{more-morphisms-lemma-action-by-derivations},
\ref{more-morphisms-lemma-sheaf} et
\ref{more-morphisms-lemma-action-sheaf}
est celui où $S$ est lui-même un épaississement $Z \subset Z' = S$
et où $Y = Z \times_{Z'} Y'$. Voici le diagramme :
$$
\xymatrix{
(X \subset X') \ar@{..>}[rr]_{(a, ?)} \ar[rd]_{(g, g')} & &
(Y \subset Y') \ar[ld]^{(h, h')} \\
& (Z \subset Z')
}
$$
Dans ce cas, l'homomorphisme $A : a^*\mathcal{C}_{Y/Y'} \to \mathcal{C}_{X/X'}$
est déterminé par $a$ : l'homomorphisme
$h^*\mathcal{C}_{Z/Z'} \to \mathcal{C}_{Y/Y'}$ est surjectif (car
on a supposé $Y = Z \times_{Z'} Y'$),
donc son image inverse $g^*\mathcal{C}_{Z/Z'} = a^*h^*\mathcal{C}_{Z/Z'} \to
a^*\mathcal{C}_{Y/Y'}$ est surjective, et le composé
$g^*\mathcal{C}_{Z/Z'} \to a^*\mathcal{C}_{Y/Y'} \to \mathcal{C}_{X/X'}$
doit être l'homomorphisme canonique induit par $g'$. Ainsi, le faisceau du
lemme \ref{more-morphisms-lemma-sheaf}
est simplement donné par la règle
$$
U' \mapsto
\{a' : U' \to Y'\text{ sur }Z'\text{ avec } a'|_U = a|_U\}
$$
et l'on fait agir sur lui le faisceau
$\SheafHom_{\mathcal{O}_X}(a^*\Omega_{Y/Z}, \mathcal{C}_{X/X'})$.
\end{remark}

\begin{lemma}
\label{more-morphisms-lemma-omega-deformation}
Soit $S$ un schéma. Soit $X \subset X'$ un épaississement du premier ordre sur
$S$. Soit $Y$ un schéma sur $S$. Soient
$a', b' : X' \to Y$ deux morphismes sur $S$ tels que
$a = a'|_X = b'|_X$. On obtient ainsi un diagramme commutatif
$$
\xymatrix{
X \ar[r] \ar[d]_a & X' \ar[d]^{(b', a')} \\
Y \ar[r]^-{\Delta_{Y/S}} & Y \times_S Y
}
$$
Puisque les flèches horizontales sont des immersions de faisceaux conormaux
$\mathcal{C}_{X/X'}$ et $\Omega_{Y/S}$, d'après
Morphismes, lemme \ref{morphisms-lemma-conormal-functorial},
on obtient un homomorphisme $\theta : a^*\Omega_{Y/S} \to \mathcal{C}_{X/X'}$.
Alors cet homomorphisme $\theta$ et la dérivation $D$ du
lemme \ref{more-morphisms-lemma-difference-derivation}
sont liés par l'équation (\ref{more-morphisms-equation-D}).
\end{lemma}

\begin{proof}
Omise. Indication : l'égalité se vérifie sur les ouverts affines, où elle
résulte du calcul suivant. Si $f$ est une section locale de
$\mathcal{O}_Y$, alors $1 \otimes f - f \otimes 1$ est une section locale
de $\mathcal{C}_{Y/(Y \times_S Y)}$ correspondant à $\text{d}_{Y/S}(f)$.
Elle s'envoie sur la section locale $(a')^\sharp(f) - (b')^\sharp(f) = D(f)$
de $\mathcal{C}_{X/X'}$. Autrement dit, $\theta(\text{d}_{Y/S}(f)) = D(f)$.
\end{proof}

\noindent
Pour la suite, il nous faut un résultat affirmant, en substance, que la
construction du
lemme \ref{more-morphisms-lemma-action-by-derivations}
est compatible avec la localisation étale.

\begin{lemma}
\label{more-morphisms-lemma-sheaf-differentials-etale-localization}
Soit
$$
\xymatrix{
X_1 \ar[d] & X_2 \ar[l]^f \ar[d] \\
S_1 & S_2 \ar[l]
}
$$
un diagramme commutatif de schémas dans lequel $X_2 \to X_1$ et $S_2 \to S_1$
sont étales. Alors l'homomorphisme $c_f : f^*\Omega_{X_1/S_1} \to \Omega_{X_2/S_2}$ de
Morphismes, lemme \ref{morphisms-lemma-functoriality-differentials},
est un isomorphisme.
\end{lemma}

\begin{proof}
Rappelons qu'un morphisme étale $U \to V$ est un morphisme lisse
tel que $\Omega_{U/V} = 0$. En utilisant cela, on voit que
Morphismes, lemme \ref{morphisms-lemma-triangle-differentials},
entraîne $\Omega_{X_2/S_2} = \Omega_{X_2/S_1}$, et que
Morphismes, lemme \ref{morphisms-lemma-triangle-differentials-smooth},
entraîne que l'homomorphisme $f^*\Omega_{X_1/S_1} \to \Omega_{X_2/S_1}$
(pour le morphisme $f$ considéré comme morphisme sur $S_1$)
est un isomorphisme. Le lemme s'en déduit.
\end{proof}

\begin{lemma}
\label{more-morphisms-lemma-action-by-derivations-etale-localization}
Considérons un diagramme commutatif d'épaississements du premier ordre
$$
\vcenter{
\xymatrix{
(T_2 \subset T_2') \ar[d]_{(h, h')} \ar[rr]_{(a_2, a_2')} & &
(X_2 \subset X_2') \ar[d]^{(f, f')} \\
(T_1 \subset T_1') \ar[rr]^{(a_1, a_1')} & &
(X_1 \subset X_1')
}
}
\quad
\begin{matrix}
\text{et un diagramme commutatif} \\
\text{de schémas}
\end{matrix}
\quad
\vcenter{
\xymatrix{
X_2' \ar[r] \ar[d] & S_2 \ar[d] \\
X_1' \ar[r] & S_1
}
}
$$
où $X_2 \to X_1$ et $S_2 \to S_1$ sont étales.
Pour tout homomorphisme $\mathcal{O}_{T_1}$-linéaire
$\theta_1 : a_1^*\Omega_{X_1/S_1} \to \mathcal{C}_{T_1/T'_1}$, notons
$\theta_2$ le composé
$$
\xymatrix{
a_2^*\Omega_{X_2/S_2} \ar@{=}[r] &
h^*a_1^*\Omega_{X_1/S_1} \ar[r]^-{h^*\theta_1} &
h^*\mathcal{C}_{T_1/T'_1} \ar[r] &
\mathcal{C}_{T_2/T'_2}
}
$$
(le signe d'égalité est expliqué dans la démonstration). Alors le diagramme
$$
\xymatrix{
T_2' \ar[rr]_{\theta_2 \cdot a_2'} \ar[d] & & X'_2 \ar[d] \\
T_1' \ar[rr]^{\theta_1 \cdot a_1'} & & X'_1
}
$$
est commutatif, où les actions $\theta_2 \cdot a_2'$ et $\theta_1 \cdot a_1'$
sont celles de la remarque \ref{more-morphisms-remark-action-by-derivations}.
\end{lemma}

\begin{proof}
Le signe d'égalité provient de l'identification
$f^*\Omega_{X_1/S_1} = \Omega_{X_2/S_2}$ du
lemme \ref{more-morphisms-lemma-sheaf-differentials-etale-localization}.
En effet, grâce à celle-ci, on a
$a_2^*\Omega_{X_2/S_2} = a_2^*f^*\Omega_{X_1/S_1} =
h^*a_1^*\Omega_{X_1/S_1}$, car $f \circ a_2 = a_1 \circ h$.
Cela étant, la commutativité du diagramme peut se vérifier
sur les ouverts affines. On peut donc supposer que les schémas du grand
diagramme initial sont affines. On obtient ainsi des diagrammes commutatifs
$$
\vcenter{
\xymatrix{
(B'_2, I_2) & & (A'_2, J_2) \ar[ll]^{a_2'} \\
(B'_1, I_1) \ar[u]^{h'} & & (A'_1, J_1) \ar[ll]_{a_1'} \ar[u]_{f'}
}
}
\quad\text{et}\quad
\vcenter{
\xymatrix{
A'_2 & & R_2 \ar[ll] \\
A'_1 \ar[u] & & R_1 \ar[ll] \ar[u]
}
}
$$
Cette notation signifie que $I_1, I_2, J_1, J_2$ sont des idéaux de carré
nul et que les morphismes de couples sont des homomorphismes d'anneaux envoyant les idéaux dans les idéaux.
Posons $A_1 = A'_1/J_1$, $A_2 = A'_2/J_2$, $B_1 = B'_1/I_1$ et
$B_2 = B'_2/I_2$. On sait que
$$
A_2 \otimes_{A_1} \Omega_{A_1/R_1} \longrightarrow \Omega_{A_2/R_2}
$$
est un isomorphisme. Alors
$\theta_1 : B_1 \otimes_{A_1} \Omega_{A_1/R_1} \to I_1$
est $B_1$-linéaire. Cela donne une $R_1$-dérivation
$D_1 = \theta_1 \circ \text{d}_{A_1/R_1} : A_1 \to I_1$.
De même, on voit que
$\theta_2 : B_2 \otimes_{A_2} \Omega_{A_2/R_2} \to I_2$
donne une $R_2$-dérivation
$D_2 = \theta_2 \circ \text{d}_{A_2/R_2} : A_2 \to I_2$.
La construction de $\theta_2$ entraîne la compatibilité suivante entre
$\theta_1$ et $\theta_2$ : pour tout $x \in A_1$, on a
$$
h'(D_1(x)) = D_2(f'(x))
$$
dans $I_2$. On peut considérer $D_1$ comme une application $A'_1 \to B'_1$
à l'aide de $A'_1 \to A_1 \xrightarrow{D_1} I_1 \to B_1$ ; de même,
on peut considérer $D_2$ comme une application $A'_2 \to B'_2$. L'égalité
affichée est alors valable pour $x \in A'_1$.
D'après la construction de l'action dans
le lemme \ref{more-morphisms-lemma-action-by-derivations} et
la remarque \ref{more-morphisms-remark-action-by-derivations},
on sait que $\theta_1 \cdot a_1'$ correspond à l'homomorphisme d'anneaux
$a_1' + D_1 : A'_1 \to B'_1$ et que $\theta_2 \cdot a_2'$ correspond
à l'homomorphisme d'anneaux $a_2' + D_2 : A'_2 \to B'_2$. L'égalité affichée
donne donc
$h' \circ (a_1' + D_1) = (a_2' + D_2) \circ f'$,
comme voulu.
\end{proof}

\begin{remark}
\label{more-morphisms-remark-tiny-improvement}
Le lemme \ref{more-morphisms-lemma-action-by-derivations-etale-localization}
peut être amélioré de la manière suivante.
Supposons donnés des diagrammes commutatifs comme dans
le lemme \ref{more-morphisms-lemma-action-by-derivations-etale-localization},
sans supposer que $X_2 \to X_1$
et $S_2 \to S_1$ soient étales. Supposons ensuite que l'on dispose de
$\theta_1 : a_1^*\Omega_{X_1/S_1} \to \mathcal{C}_{T_1/T'_1}$
et de
$\theta_2 : a_2^*\Omega_{X_2/S_2} \to \mathcal{C}_{T_2/T'_2}$
tels que
$$
\xymatrix{
f_*\mathcal{O}_{X_2} \ar[rr]_{f_*D_2} & &
f_*a_{2, *}\mathcal{C}_{T_2/T_2'} \\
\mathcal{O}_{X_1} \ar[rr]^{D_1} \ar[u]^{f^\sharp} & &
a_{1, *}\mathcal{C}_{T_1/T_1'} \ar[u]_{\text{induit par }(h')^\sharp}
}
$$
soit commutatif, où $D_i$ correspond à $\theta_i$ comme dans
l'équation (\ref{more-morphisms-equation-D}). On obtient alors la conclusion du
lemme \ref{more-morphisms-lemma-action-by-derivations-etale-localization}.
L'intérêt de l'hypothèse que $X_2 \to X_1$ et
$S_2 \to S_1$ soient tous deux étales est qu'elle permet de construire un $\theta_2$
à partir de $\theta_1$.
\end{remark}










\section{Déformations infinitésimales des schémas}
\label{more-morphisms-section-deform}

\noindent
Le lemme élémentaire suivant est souvent commode pour vérifier
qu'une déformation infinitésimale d'un morphisme est plate.

\begin{lemma}
\label{more-morphisms-lemma-deform}
Soit $(f, f') : (X \subset X') \to (S \subset S')$ un morphisme
d'épaississements du premier ordre. Supposons $f$ plat.
Alors les conditions suivantes sont équivalentes :
\begin{enumerate}
\item $f'$ est plat et $X = S \times_{S'} X'$, et
\item l'homomorphisme canonique $f^*\mathcal{C}_{S/S'} \to \mathcal{C}_{X/X'}$
est un isomorphisme.
\end{enumerate}
\end{lemma}

\begin{proof}
Le problème étant local sur $X'$, on peut supposer que $X, X', S, S'$
sont des schémas affines. Écrivons $S' = \Spec(A')$, $X' = \Spec(B')$,
$S = \Spec(A)$, $X = \Spec(B)$, avec $A = A'/I$ et $B = B'/J$
pour certains idéaux de carré nul. On obtient alors le diagramme commutatif
suivant :
$$
\xymatrix{
0 \ar[r] &
J \ar[r] &
B' \ar[r] &
B \ar[r] & 0 \\
0 \ar[r] &
I \ar[r] \ar[u] &
A' \ar[r] \ar[u] &
A \ar[r] \ar[u] & 0
}
$$
dont les lignes sont exactes. L'homomorphisme canonique du lemme est l'homomorphisme
$$
I \otimes_A B = I \otimes_{A'} B' \longrightarrow J.
$$
L'hypothèse que $f$ est plat signifie que $A \to B$ est plat.

\medskip\noindent
Supposons (1). Alors $A' \to B'$ est plat et $J = IB'$. La platitude entraîne
$\text{Tor}_1^{A'}(B', A) = 0$ (voir
Algèbre, lemme \ref{algebra-lemma-characterize-flat}).
Cela signifie que $I \otimes_{A'} B' \to B'$ est injectif (voir
Algèbre, remarque \ref{algebra-remark-Tor-ring-mod-ideal}).
On voit donc que $I \otimes_A B \to J$ est un isomorphisme.

\medskip\noindent
Supposons (2). Il en résulte que $J = IB'$, donc que $X = S \times_{S'} X'$.
De plus, on obtient $\text{Tor}_1^{A'}(B', A'/I) = 0$ en renversant
les implications du paragraphe précédent. Ainsi, $B'$ est plat sur $A'$ d'après
Algèbre, lemme \ref{algebra-lemma-what-does-it-mean}.
\end{proof}

\noindent
Le lemme suivant est la version « nilpotente » du
« critère de platitude par fibres », voir
la section \ref{more-morphisms-section-criterion-flat-fibres}.

\begin{lemma}
\label{more-morphisms-lemma-flatness-morphism-thickenings}
Considérons un diagramme commutatif
$$
\xymatrix{
(X \subset X') \ar[rr]_{(f, f')} \ar[rd] & & (Y \subset Y') \ar[ld] \\
& (S \subset S')
}
$$
d'épaississements. Supposons que
\begin{enumerate}
\item $X'$ soit plat sur $S'$,
\item $f$ soit plat,
\item $S \subset S'$ soit un épaississement d'ordre fini, et
\item $X = S \times_{S'} X'$ et $Y = S \times_{S'} Y'$.
\end{enumerate}
Alors $f'$ est plat et $Y'$ est plat sur $S'$ en tout point de
l'image de $f'$.
\end{lemma}

\begin{proof}
C'est une conséquence immédiate du résultat d'Algèbre,
lemme \ref{algebra-lemma-criterion-flatness-fibre-nilpotent}.
\end{proof}

\noindent
De nombreuses propriétés des morphismes de schémas sont préservées par déformation
plate.

\begin{lemma}
\label{more-morphisms-lemma-deform-property}
Considérons un diagramme commutatif
$$
\xymatrix{
(X \subset X') \ar[rr]_{(f, f')} \ar[rd] & & (Y \subset Y') \ar[ld] \\
& (S \subset S')
}
$$
d'épaississements. Supposons que $S \subset S'$ soit un épaississement d'ordre fini,
que $X'$ soit plat sur $S'$, que $X = S \times_{S'} X'$ et que
$Y = S \times_{S'} Y'$. Alors
\begin{enumerate}
\item $f$ est plat si et seulement si $f'$ est plat,
\label{more-morphisms-item-flat}
\item $f$ est un isomorphisme si et seulement si $f'$ est un isomorphisme,
\label{more-morphisms-item-isomorphism}
\item $f$ est une immersion ouverte si et seulement si $f'$ est une immersion ouverte,
\label{more-morphisms-item-open-immersion}
\item $f$ est quasi-compact si et seulement si $f'$ est quasi-compact,
\label{more-morphisms-item-quasi-compact}
\item $f$ est universellement fermé si et seulement si $f'$ est universellement fermé,
\label{more-morphisms-item-universally-closed}
\item $f$ est (quasi-)séparé si et seulement si $f'$ est (quasi-)séparé,
\label{more-morphisms-item-separated}
\item $f$ est un monomorphisme si et seulement si $f'$ est un monomorphisme,
\label{more-morphisms-item-monomorphism}
\item $f$ est surjectif si et seulement si $f'$ est surjectif,
\label{more-morphisms-item-surjective}
\item $f$ est universellement injectif si et seulement si $f'$ est universellement injectif,
\label{more-morphisms-item-universally-injective}
\item $f$ est affine si et seulement si $f'$ est affine,
\label{more-morphisms-item-affine}
\item
\label{more-morphisms-item-finite-type}
$f$ est localement de type fini si et seulement si $f'$ est localement de type fini,
\item $f$ est localement quasi-fini si et seulement si $f'$ est localement quasi-fini,
\label{more-morphisms-item-quasi-finite}
\item
\label{more-morphisms-item-finite-presentation}
$f$ est localement de présentation finie si et seulement si $f'$ est localement de
présentation finie,
\item
\label{more-morphisms-item-relative-dimension-d}
$f$ est localement de type fini et de dimension relative $d$ si et seulement si
$f'$ est localement de type fini et de dimension relative $d$,
\item $f$ est universellement ouvert si et seulement si $f'$ est universellement ouvert,
\label{more-morphisms-item-universally-open}
\item $f$ est syntomique si et seulement si $f'$ est syntomique,
\label{more-morphisms-item-syntomic}
\item $f$ est lisse si et seulement si $f'$ est lisse,
\label{more-morphisms-item-smooth}
\item $f$ est non ramifié si et seulement si $f'$ est non ramifié,
\label{more-morphisms-item-unramified}
\item $f$ est étale si et seulement si $f'$ est étale,
\label{more-morphisms-item-etale}
\item $f$ est propre si et seulement si $f'$ est propre,
\label{more-morphisms-item-proper}
\item $f$ est entier si et seulement si $f'$ est entier,
\label{more-morphisms-item-integral}
\item $f$ est fini si et seulement si $f'$ est fini,
\label{more-morphisms-item-finite}
\item
\label{more-morphisms-item-finite-locally-free}
$f$ est fini localement libre (de rang $d$) si et seulement si $f'$
est fini localement libre (de rang $d$), et
\item ajouter d'autres propriétés ici.
\end{enumerate}
\end{lemma}

\begin{proof}
Les hypothèses sur $X$ et $Y$ signifient que $f$ est le changement de base de
$f'$ par $X \to X'$.
Les propriétés $\mathcal{P}$ énumérées en (1) -- (23) ci-dessus sont toutes stables
par changement de base ; si $f'$ possède la propriété $\mathcal{P}$,
il en va donc de même de $f$. Voir
Schémas, lemmes \ref{schemes-lemma-base-change-immersion},
\ref{schemes-lemma-quasi-compact-preserved-base-change},
\ref{schemes-lemma-separated-permanence} et
\ref{schemes-lemma-base-change-monomorphism},
et
Morphismes, lemmes
\ref{morphisms-lemma-base-change-surjective},
\ref{morphisms-lemma-base-change-universally-injective},
\ref{morphisms-lemma-base-change-affine},
\ref{morphisms-lemma-base-change-finite-type},
\ref{morphisms-lemma-base-change-quasi-finite},
\ref{morphisms-lemma-base-change-finite-presentation},
\ref{morphisms-lemma-base-change-relative-dimension-d},
\ref{morphisms-lemma-base-change-syntomic},
\ref{morphisms-lemma-base-change-smooth},
\ref{morphisms-lemma-base-change-unramified},
\ref{morphisms-lemma-base-change-etale},
\ref{morphisms-lemma-base-change-proper},
\ref{morphisms-lemma-base-change-finite} et
\ref{morphisms-lemma-base-change-finite-locally-free}.

\medskip\noindent
Dans chaque cas, le sens intéressant consiste donc à supposer
que $f$ possède la propriété et à en déduire que $f'$ la possède aussi.
Par récurrence sur l'ordre de l'épaississement, on peut
supposer que $S \subset S'$ est un épaississement du premier ordre, voir
la discussion qui suit immédiatement
la définition \ref{more-morphisms-definition-thickening}.
Faisons quelques remarques générales que l'on utilisera sans autre
mention dans les arguments ci-dessous.
(I) Soient $W' \subset S'$ un ouvert affine et $U' \subset X'$
et $V' \subset Y'$ des ouverts affines au-dessus de $W'$ tels que $f'(U') \subset V'$.
Écrivons $W' = \Spec(R')$ et notons $I \subset R'$ l'idéal
définissant le sous-schéma fermé $W' \cap S$. Écrivons $U' = \Spec(B')$
et $V' = \Spec(A')$. On obtient alors un diagramme commutatif
$$
\xymatrix{
0 \ar[r] &
IB' \ar[r] &
B' \ar[r] &
B \ar[r] & 0 \\
0 \ar[r] &
IA' \ar[r] \ar[u] &
A' \ar[r] \ar[u] &
A \ar[r] \ar[u] & 0
}
$$
à lignes exactes. De plus, $IB' \cong I \otimes_R B$, voir la démonstration du
lemme \ref{more-morphisms-lemma-deform}.
(II) Les morphismes $X \to X'$ et $Y \to Y'$ sont des homéomorphismes universels.
Les applications $f$ et $f'$ ont donc la même topologie (après tout changement de base).
(III) Si $f$ est plat, alors $f'$ est plat et
$Y' \to S'$ est plat en tout point de l'image de $f'$, voir
le lemme \ref{more-morphisms-lemma-flatness-morphism-thickenings}.

\medskip\noindent
Pour (\ref{more-morphisms-item-flat}). C'est la remarque générale (III).

\medskip\noindent
Pour (\ref{more-morphisms-item-isomorphism}). Supposons que $f$ soit un isomorphisme.
D'après (III), on voit que $Y' \to S'$ est plat. Choisissons un
ouvert affine $V' \subset Y'$ et posons $U' = (f')^{-1}(V')$. Alors
$V = Y \cap V'$ est affine, ce qui entraîne que
$V \cong f^{-1}(V) = U = Y \times_{Y'} U'$ est affine. D'après
le lemme \ref{more-morphisms-lemma-thickening-affine-scheme},
on voit que $U'$ est affine. On dispose donc d'un diagramme comme dans
la remarque générale (I), et de plus $IA \cong I \otimes_R A$, car
$R' \to A'$ est plat. Alors
$IB' \cong I \otimes_R B \cong I \otimes_R A \cong IA'$
et $A \cong B$. L'exactitude des lignes
du diagramme ci-dessus donne $A' \cong B'$,
c'est-à-dire $U' \cong V'$. Ainsi, $f'$ est un isomorphisme.

\medskip\noindent
Pour (\ref{more-morphisms-item-open-immersion}). Supposons que $f$ soit une immersion ouverte.
Alors $f$ est un isomorphisme de $X$ sur un sous-schéma ouvert $V \subset Y$.
Soit $V' \subset Y'$ le sous-schéma ouvert dont l'espace topologique
sous-jacent est $V$. Alors $f'$ est un morphisme de $X'$ dans $V'$, qui est un isomorphisme d'après
(\ref{more-morphisms-item-isomorphism}). Ainsi, $f'$ est une immersion ouverte.

\medskip\noindent
Pour (\ref{more-morphisms-item-quasi-compact}). C'est immédiat d'après la remarque (II). Voir aussi
le lemme \ref{more-morphisms-lemma-thicken-property-morphisms} pour un énoncé plus général.

\medskip\noindent
Pour (\ref{more-morphisms-item-universally-closed}). C'est immédiat d'après la remarque (II). Voir aussi
le lemme \ref{more-morphisms-lemma-thicken-property-morphisms} pour un énoncé plus général.

\medskip\noindent
Pour (\ref{more-morphisms-item-separated}). Remarquons que
$X \times_Y X = Y \times_{Y'} (X' \times_{Y'} X')$, de sorte que
$X' \times_{Y'} X'$ est un épaississement de $X \times_Y X$.
Les applications $\Delta_{X/Y}$ et $\Delta_{X'/Y'}$ ont donc la même
topologie, ce qui permet de conclure. Voir aussi
le lemme \ref{more-morphisms-lemma-thicken-property-morphisms} pour un énoncé plus général.

\medskip\noindent
Pour (\ref{more-morphisms-item-monomorphism}). Supposons que $f$ soit un monomorphisme.
Considérons le morphisme diagonal $\Delta_{X'/Y'} : X' \to X' \times_{Y'} X'$.
Le changement de base de $\Delta_{X'/Y'}$ par $S \to S'$ est $\Delta_{X/Y}$,
qui est un isomorphisme par hypothèse. D'après (\ref{more-morphisms-item-isomorphism}),
on en déduit que $\Delta_{X'/Y'}$ est un isomorphisme.

\medskip\noindent
Pour (\ref{more-morphisms-item-surjective}). C'est clair. Voir aussi
le lemme \ref{more-morphisms-lemma-thicken-property-morphisms} pour un énoncé plus général.

\medskip\noindent
Pour (\ref{more-morphisms-item-universally-injective}). C'est immédiat d'après la remarque (II). Voir aussi
le lemme \ref{more-morphisms-lemma-thicken-property-morphisms} pour un énoncé plus général.

\medskip\noindent
Pour (\ref{more-morphisms-item-affine}). Supposons que $f$ soit affine. Choisissons un
ouvert affine $V' \subset Y'$ et posons $U' = (f')^{-1}(V')$.
Alors $V = Y \cap V'$ est affine, ce qui entraîne que
$U = Y \times_{Y'} U'$ est affine. D'après
le lemme \ref{more-morphisms-lemma-thickening-affine-scheme},
on voit que $U'$ est affine. Ainsi, $f'$ est affine. Voir aussi
le lemme \ref{more-morphisms-lemma-thicken-property-morphisms} pour un énoncé plus général.

\medskip\noindent
Pour (\ref{more-morphisms-item-finite-type}). D'après la remarque (I), il s'agit de prouver que $A' \to B'$
est de type fini si $A \to B$ est de type fini. Supposons que
$x_1, \ldots, x_n \in B'$ soient des éléments dont les images dans $B$ engendrent $B$
comme $A$-algèbre. Alors $A'[x_1, \ldots, x_n] \to B$ est surjectif,
car $A'[x_1, \ldots, x_n] \to B$ est surjectif et
$I \otimes_R A[x_1, \ldots, x_n] \to I \otimes_R B$ est surjectif. Voir aussi
le lemme \ref{more-morphisms-lemma-thicken-property-morphisms-cartesian}
pour un énoncé plus général.

\medskip\noindent
Pour (\ref{more-morphisms-item-quasi-finite}). Cela résulte de (\ref{more-morphisms-item-finite-type}) et du fait que
la quasi-finitude d'un morphisme de type fini se vérifie sur les fibres, voir
Morphismes, lemme \ref{morphisms-lemma-quasi-finite-at-point-characterize}.
Voir aussi le lemme \ref{more-morphisms-lemma-thicken-property-morphisms-cartesian}
pour un énoncé plus général.

\medskip\noindent
Pour (\ref{more-morphisms-item-finite-presentation}). D'après la remarque (I), il s'agit de prouver que
$A' \to B'$ est de présentation finie si $A \to B$ est de présentation finie.
On peut supposer que $B' = A'[x_1, \ldots, x_n]/K'$ pour un certain idéal $K'$ d'après
(\ref{more-morphisms-item-finite-type}). On obtient une suite exacte courte
$$
0 \to K' \to A'[x_1, \ldots, x_n] \to B' \to 0
$$
Comme $B'$ est plat sur $R'$, on voit que $K' \otimes_{R'} R$ est le noyau de
la surjection $A[x_1, \ldots, x_n] \to B$. Par hypothèse sur $A \to B$,
il existe un nombre fini d'éléments $f'_1, \ldots, f'_m \in K'$ dont les images dans
$A[x_1, \ldots, x_n]$ engendrent ce noyau. Puisque $I$ est nilpotent, on voit
que $f'_1, \ldots, f'_m$ engendrent $K'$ par le lemme de Nakayama, voir
Algèbre, lemme \ref{algebra-lemma-NAK}.

\medskip\noindent
Pour (\ref{more-morphisms-item-relative-dimension-d}). Cela résulte de (\ref{more-morphisms-item-finite-type})
et de la remarque générale (II). Voir aussi
le lemme \ref{more-morphisms-lemma-thicken-property-morphisms-cartesian}
pour un énoncé plus général.

\medskip\noindent
Pour (\ref{more-morphisms-item-universally-open}). C'est immédiat d'après la remarque générale (II). Voir aussi
le lemme \ref{more-morphisms-lemma-thicken-property-morphisms} pour un énoncé plus général.

\medskip\noindent
Pour (\ref{more-morphisms-item-syntomic}). Supposons $f$ syntomique. D'après
(\ref{more-morphisms-item-finite-presentation}), $f'$ est localement de présentation finie ;
d'après la remarque générale (III), $f'$ est plat, et les fibres de $f'$ sont celles
de $f$. Ainsi, $f'$ est syntomique d'après
Morphismes, lemme \ref{morphisms-lemma-syntomic-flat-fibres}.

\medskip\noindent
Pour (\ref{more-morphisms-item-smooth}). Supposons $f$ lisse. D'après
(\ref{more-morphisms-item-finite-presentation}), $f'$ est localement de présentation finie ;
d'après la remarque générale (III), $f'$ est plat, et les fibres de $f'$ sont les
fibres de $f$. Ainsi, $f'$ est lisse d'après
Morphismes, lemme \ref{morphisms-lemma-smooth-flat-smooth-fibres}.

\medskip\noindent
Pour (\ref{more-morphisms-item-unramified}). Supposons $f$ non ramifié. D'après
(\ref{more-morphisms-item-finite-type}), $f'$ est localement de type fini,
et les fibres de $f'$ sont celles de $f$.
Ainsi, $f'$ est non ramifié d'après
Morphismes, lemme \ref{morphisms-lemma-unramified-etale-fibres}. Voir aussi
le lemme \ref{more-morphisms-lemma-thicken-property-morphisms-cartesian}
pour un énoncé plus général.

\medskip\noindent
Pour (\ref{more-morphisms-item-etale}). Supposons $f$ étale. D'après
(\ref{more-morphisms-item-finite-presentation}), $f'$ est localement de présentation finie ;
d'après la remarque générale (III), $f'$ est plat, et les fibres de $f'$ sont celles
de $f$. Ainsi, $f'$ est étale d'après
Morphismes, lemme \ref{morphisms-lemma-etale-flat-etale-fibres}.

\medskip\noindent
Pour (\ref{more-morphisms-item-proper}). Cela résulte de la combinaison de
(\ref{more-morphisms-item-separated}), (\ref{more-morphisms-item-finite-type}), (\ref{more-morphisms-item-quasi-compact})
et (\ref{more-morphisms-item-universally-closed}). Voir aussi
le lemme \ref{more-morphisms-lemma-thicken-property-morphisms-cartesian}
pour un énoncé plus général.

\medskip\noindent
Pour (\ref{more-morphisms-item-integral}). Combiner (\ref{more-morphisms-item-universally-closed}) et
(\ref{more-morphisms-item-affine}) avec
Morphismes, lemme \ref{morphisms-lemma-integral-universally-closed}. Voir aussi
le lemme \ref{more-morphisms-lemma-thicken-property-morphisms} pour un énoncé plus général.

\medskip\noindent
Pour (\ref{more-morphisms-item-finite}). Combiner (\ref{more-morphisms-item-integral})
et (\ref{more-morphisms-item-finite-type}) avec
Morphismes, lemme \ref{morphisms-lemma-finite-integral}. Voir aussi
le lemme \ref{more-morphisms-lemma-thicken-property-morphisms-cartesian}
pour un énoncé plus général.

\medskip\noindent
Pour (\ref{more-morphisms-item-finite-locally-free}). Supposons $f$ fini localement libre. D'après
(\ref{more-morphisms-item-finite}), on voit que $f'$ est fini ; d'après la remarque générale (III),
$f'$ est plat, et d'après (\ref{more-morphisms-item-finite-presentation}), $f'$ est localement de
présentation finie. Ainsi, $f'$ est fini localement libre d'après
Morphismes, lemme \ref{morphisms-lemma-finite-flat}.
\end{proof}


\noindent
Le lemme suivant est la version « localement nilpotente » du
« critère de platitude par fibres », voir
la section \ref{more-morphisms-section-criterion-flat-fibres}.

\begin{lemma}
\label{more-morphisms-lemma-flatness-morphism-thickenings-fp-over-ft}
Considérons un diagramme commutatif
$$
\xymatrix{
(X \subset X') \ar[rr]_{(f, f')} \ar[rd] & & (Y \subset Y') \ar[ld] \\
& (S \subset S')
}
$$
d'épaississements. Supposons que
\begin{enumerate}
\item $Y' \to S'$ soit localement de type fini,
\item $X' \to S'$ soit plat et localement de présentation finie,
\item $f$ soit plat, et
\item $X = S \times_{S'} X'$ et $Y = S \times_{S'} Y'$.
\end{enumerate}
Alors $f'$ est plat et, pour tout $y' \in Y'$ dans l'image de $f'$,
l'anneau local $\mathcal{O}_{Y', y'}$ est
plat et essentiellement de présentation finie sur $\mathcal{O}_{S', s'}$.
\end{lemma}

\begin{proof}
C'est une conséquence immédiate du résultat d'Algèbre,
lemme \ref{algebra-lemma-criterion-flatness-fibre-locally-nilpotent}.
\end{proof}

\noindent
De nombreuses propriétés des morphismes de schémas sont préservées par les déformations
plates considérées dans le lemme précédent.

\begin{lemma}
\label{more-morphisms-lemma-deform-property-fp-over-ft}
Considérons un diagramme commutatif
$$
\xymatrix{
(X \subset X') \ar[rr]_{(f, f')} \ar[rd] & & (Y \subset Y') \ar[ld] \\
& (S \subset S')
}
$$
d'épaississements. Supposons $Y' \to S'$ localement de type fini,
$X' \to S'$ plat et localement de présentation finie,
$X = S \times_{S'} X'$ et $Y = S \times_{S'} Y'$. Alors
\begin{enumerate}
\item $f$ est plat si et seulement si $f'$ est plat,
\label{more-morphisms-item-flat-fp-over-ft}
\item $f$ est un isomorphisme si et seulement si $f'$ est un isomorphisme,
\label{more-morphisms-item-isomorphism-fp-over-ft}
\item $f$ est une immersion ouverte si et seulement si $f'$ est une immersion ouverte,
\label{more-morphisms-item-open-immersion-fp-over-ft}
\item $f$ est quasi-compact si et seulement si $f'$ est quasi-compact,
\label{more-morphisms-item-quasi-compact-fp-over-ft}
\item $f$ est universellement fermé si et seulement si $f'$ est universellement fermé,
\label{more-morphisms-item-universally-closed-fp-over-ft}
\item $f$ est (quasi-)séparé si et seulement si $f'$ est (quasi-)séparé,
\label{more-morphisms-item-separated-fp-over-ft}
\item $f$ est un monomorphisme si et seulement si $f'$ est un monomorphisme,
\label{more-morphisms-item-monomorphism-fp-over-ft}
\item $f$ est surjectif si et seulement si $f'$ est surjectif,
\label{more-morphisms-item-surjective-fp-over-ft}
\item $f$ est universellement injectif si et seulement si $f'$ est universellement injectif,
\label{more-morphisms-item-universally-injective-fp-over-ft}
\item $f$ est affine si et seulement si $f'$ est affine,
\label{more-morphisms-item-affine-fp-over-ft}
\item $f$ est localement quasi-fini si et seulement si $f'$ est localement quasi-fini,
\label{more-morphisms-item-quasi-finite-fp-over-ft}
\item
\label{more-morphisms-item-relative-dimension-d-fp-over-ft}
$f$ est localement de type fini et de dimension relative $d$ si et seulement si
$f'$ est localement de type fini et de dimension relative $d$,
\item $f$ est universellement ouvert si et seulement si $f'$ est universellement ouvert,
\label{more-morphisms-item-universally-open-fp-over-ft}
\item $f$ est syntomique si et seulement si $f'$ est syntomique,
\label{more-morphisms-item-syntomic-fp-over-ft}
\item $f$ est lisse si et seulement si $f'$ est lisse,
\label{more-morphisms-item-smooth-fp-over-ft}
\item $f$ est non ramifié si et seulement si $f'$ est non ramifié,
\label{more-morphisms-item-unramified-fp-over-ft}
\item $f$ est étale si et seulement si $f'$ est étale,
\label{more-morphisms-item-etale-fp-over-ft}
\item $f$ est propre si et seulement si $f'$ est propre,
\label{more-morphisms-item-proper-fp-over-ft}
\item $f$ est fini si et seulement si $f'$ est fini,
\label{more-morphisms-item-finite-fp-over-ft}
\item
\label{more-morphisms-item-finite-locally-free-fp-over-ft}
$f$ est fini localement libre (de rang $d$) si et seulement si $f'$
est fini localement libre (de rang $d$), et
\item ajouter d'autres propriétés ici.
\end{enumerate}
\end{lemma}

\begin{proof}
Les hypothèses sur $X$ et $Y$ signifient que $f$ est le changement de base de
$f'$ par $X \to X'$.
Les propriétés $\mathcal{P}$ énumérées en (1) -- (20) ci-dessus sont toutes stables
par changement de base ; si $f'$ possède la propriété $\mathcal{P}$,
il en va donc de même de $f$. Voir
Schémas, lemmes \ref{schemes-lemma-base-change-immersion},
\ref{schemes-lemma-quasi-compact-preserved-base-change},
\ref{schemes-lemma-separated-permanence} et
\ref{schemes-lemma-base-change-monomorphism},
et
Morphismes, lemmes
\ref{morphisms-lemma-base-change-surjective},
\ref{morphisms-lemma-base-change-universally-injective},
\ref{morphisms-lemma-base-change-affine},
\ref{morphisms-lemma-base-change-quasi-finite},
\ref{morphisms-lemma-base-change-relative-dimension-d},
\ref{morphisms-lemma-base-change-syntomic},
\ref{morphisms-lemma-base-change-smooth},
\ref{morphisms-lemma-base-change-unramified},
\ref{morphisms-lemma-base-change-etale},
\ref{morphisms-lemma-base-change-proper},
\ref{morphisms-lemma-base-change-finite} et
\ref{morphisms-lemma-base-change-finite-locally-free}.

\medskip\noindent
Dans chaque cas, le sens intéressant consiste donc à supposer
que $f$ possède la propriété et à en déduire que $f'$ la possède aussi.
Faisons quelques remarques générales que l'on utilisera sans autre
mention dans les arguments ci-dessous.
(I) Soient $W' \subset S'$ un ouvert affine et $U' \subset X'$
et $V' \subset Y'$ des ouverts affines au-dessus de $W'$ tels que $f'(U') \subset V'$.
Écrivons $W' = \Spec(R')$ et notons $I \subset R'$ l'idéal
définissant le sous-schéma fermé $W' \cap S$. Écrivons $U' = \Spec(B')$
et $V' = \Spec(A')$. On obtient alors un diagramme commutatif
$$
\xymatrix{
0 \ar[r] &
IB' \ar[r] &
B' \ar[r] &
B \ar[r] & 0 \\
0 \ar[r] &
IA' \ar[r] \ar[u] &
A' \ar[r] \ar[u] &
A \ar[r] \ar[u] & 0
}
$$
à lignes exactes.
(II) Les morphismes $X \to X'$ et $Y \to Y'$ sont des homéomorphismes universels.
Les applications $f$ et $f'$ ont donc la même topologie (après tout changement de base).
(III) Si $f$ est plat, alors $f'$ est plat et $Y' \to S'$ est plat en tout
point de l'image de $f'$, voir
le lemme \ref{more-morphisms-lemma-flatness-morphism-thickenings}.

\medskip\noindent
Pour (\ref{more-morphisms-item-flat-fp-over-ft}). C'est la remarque générale (III).

\medskip\noindent
Pour (\ref{more-morphisms-item-isomorphism-fp-over-ft}). Supposons que $f$ soit un isomorphisme.
Choisissons un ouvert affine $V' \subset Y'$ et posons $U' = (f')^{-1}(V')$.
Alors $V = Y \cap V'$ est affine, ce qui entraîne que
$V \cong f^{-1}(V) = U = Y \times_{Y'} U'$ est affine. D'après
le lemme \ref{more-morphisms-lemma-thickening-affine-scheme},
on voit que $U'$ est affine. On dispose donc d'un diagramme comme dans
la remarque générale (I). D'après Algèbre, lemme
\ref{algebra-lemma-isomorphism-modulo-locally-nilpotent},
on voit que $A' \to B'$ est un isomorphisme, c'est-à-dire $U' \cong V'$.
Ainsi, $f'$ est un isomorphisme.

\medskip\noindent
Pour (\ref{more-morphisms-item-open-immersion-fp-over-ft}). Supposons que $f$ soit une immersion ouverte.
Alors $f$ est un isomorphisme de $X$ sur un sous-schéma ouvert $V \subset Y$.
Soit $V' \subset Y'$ le sous-schéma ouvert dont l'espace topologique
sous-jacent est $V$. Alors $f'$ est un morphisme de $X'$ dans $V'$, qui est un isomorphisme d'après
(\ref{more-morphisms-item-isomorphism-fp-over-ft}). Ainsi, $f'$ est une immersion ouverte.

\medskip\noindent
Pour (\ref{more-morphisms-item-quasi-compact-fp-over-ft}). C'est immédiat d'après la remarque (II). Voir aussi
le lemme \ref{more-morphisms-lemma-thicken-property-morphisms} pour un énoncé plus général.

\medskip\noindent
Pour (\ref{more-morphisms-item-universally-closed-fp-over-ft}). C'est immédiat d'après la remarque (II). Voir
aussi le lemme \ref{more-morphisms-lemma-thicken-property-morphisms} pour un énoncé plus général.

\medskip\noindent
Pour (\ref{more-morphisms-item-separated-fp-over-ft}). Remarquons que
$X \times_Y X = Y \times_{Y'} (X' \times_{Y'} X')$, de sorte que
$X' \times_{Y'} X'$ est un épaississement de $X \times_Y X$.
Les applications $\Delta_{X/Y}$ et $\Delta_{X'/Y'}$ ont donc la même
topologie, ce qui permet de conclure. Voir aussi
le lemme \ref{more-morphisms-lemma-thicken-property-morphisms} pour un énoncé plus général.

\medskip\noindent
Pour (\ref{more-morphisms-item-monomorphism-fp-over-ft}). Supposons que $f$ soit un monomorphisme.
Considérons le morphisme diagonal $\Delta_{X'/Y'} : X' \to X' \times_{Y'} X'$.
Remarquons que $X' \times_{Y'} X' \to S'$ est localement de type fini.
Le changement de base de $\Delta_{X'/Y'}$ par $S \to S'$ est $\Delta_{X/Y}$,
qui est un isomorphisme par hypothèse. D'après (\ref{more-morphisms-item-isomorphism-fp-over-ft}),
on en déduit que $\Delta_{X'/Y'}$ est un isomorphisme.

\medskip\noindent
Pour (\ref{more-morphisms-item-surjective-fp-over-ft}). C'est clair. Voir aussi
le lemme \ref{more-morphisms-lemma-thicken-property-morphisms} pour un énoncé plus général.

\medskip\noindent
Pour (\ref{more-morphisms-item-universally-injective-fp-over-ft}).
C'est immédiat d'après la remarque (II). Voir aussi
le lemme \ref{more-morphisms-lemma-thicken-property-morphisms} pour un énoncé plus général.

\medskip\noindent
Pour (\ref{more-morphisms-item-affine-fp-over-ft}). Supposons que $f$ soit affine. Choisissons un
ouvert affine $V' \subset Y'$ et posons $U' = (f')^{-1}(V')$.
Alors $V = Y \cap V'$ est affine, ce qui entraîne que
$U = Y \times_{Y'} U'$ est affine. D'après
le lemme \ref{more-morphisms-lemma-thickening-affine-scheme},
on voit que $U'$ est affine. Ainsi, $f'$ est affine. Voir aussi
le lemme \ref{more-morphisms-lemma-thicken-property-morphisms} pour un énoncé plus général.

\medskip\noindent
Pour (\ref{more-morphisms-item-quasi-finite-fp-over-ft}). Cela résulte du fait que $f'$
est localement de type fini
(d'après Morphismes, lemme \ref{morphisms-lemma-permanence-finite-type}) et du fait que
la quasi-finitude d'un morphisme de type fini se vérifie sur les fibres, voir
Morphismes, lemme \ref{morphisms-lemma-quasi-finite-at-point-characterize}.

\medskip\noindent
Pour (\ref{more-morphisms-item-relative-dimension-d-fp-over-ft}).
Cela résulte de la remarque générale (II) et du fait que $f'$
est localement de type fini
(Morphismes, lemme \ref{morphisms-lemma-permanence-finite-type}).

\medskip\noindent
Pour (\ref{more-morphisms-item-universally-open-fp-over-ft}).
C'est immédiat d'après la remarque générale (II). Voir aussi
le lemme \ref{more-morphisms-lemma-thicken-property-morphisms} pour un énoncé plus général.

\medskip\noindent
Pour (\ref{more-morphisms-item-syntomic-fp-over-ft}). Supposons $f$ syntomique. D'après
Morphismes, lemme \ref{morphisms-lemma-finite-presentation-permanence},
$f'$ est localement de présentation finie.
D'après la remarque générale (III), $f'$ est plat. Les fibres de $f'$ sont celles
de $f$. Ainsi, $f'$ est syntomique d'après
Morphismes, lemme \ref{morphisms-lemma-syntomic-flat-fibres}.

\medskip\noindent
Pour (\ref{more-morphisms-item-smooth-fp-over-ft}). Supposons $f$ lisse. D'après
Morphismes, lemme \ref{morphisms-lemma-finite-presentation-permanence},
$f'$ est localement de présentation finie.
D'après la remarque générale (III), $f'$ est plat. Les fibres de $f'$ sont les
fibres de $f$. Ainsi, $f'$ est lisse d'après
Morphismes, lemme \ref{morphisms-lemma-smooth-flat-smooth-fibres}.

\medskip\noindent
Pour (\ref{more-morphisms-item-unramified-fp-over-ft}). Supposons $f$ non ramifié. D'après
Morphismes, lemme \ref{morphisms-lemma-permanence-finite-type},
$f'$ est localement de type fini. Les fibres de $f'$ sont celles de $f$.
Ainsi, $f'$ est non ramifié d'après
Morphismes, lemme \ref{morphisms-lemma-unramified-etale-fibres}.

\medskip\noindent
Pour (\ref{more-morphisms-item-etale-fp-over-ft}). Supposons $f$ étale. D'après
Morphismes, lemme \ref{morphisms-lemma-finite-presentation-permanence},
$f'$ est localement de présentation finie.
D'après la remarque générale (III), $f'$ est plat.
Les fibres de $f'$ sont celles de $f$. Ainsi, $f'$ est étale d'après
Morphismes, lemme \ref{morphisms-lemma-etale-flat-etale-fibres}.

\medskip\noindent
Pour (\ref{more-morphisms-item-proper-fp-over-ft}). Cela résulte de la combinaison de
(\ref{more-morphisms-item-separated-fp-over-ft}), du fait que $f$ est localement
de type fini (Morphismes, lemme \ref{morphisms-lemma-permanence-finite-type}),
de (\ref{more-morphisms-item-quasi-compact-fp-over-ft})
et de (\ref{more-morphisms-item-universally-closed-fp-over-ft}).

\medskip\noindent
Pour (\ref{more-morphisms-item-finite-fp-over-ft}).
Combiner (\ref{more-morphisms-item-universally-closed-fp-over-ft}),
(\ref{more-morphisms-item-affine-fp-over-ft}),
Morphismes, lemme \ref{morphisms-lemma-integral-universally-closed},
le fait que $f$ est localement de type fini
(Morphismes, lemme \ref{morphisms-lemma-permanence-finite-type}) et
Morphismes, lemme \ref{morphisms-lemma-finite-integral}.

\medskip\noindent
Pour (\ref{more-morphisms-item-finite-locally-free-fp-over-ft}).
Supposons $f$ fini localement libre. D'après
(\ref{more-morphisms-item-finite-fp-over-ft}), on voit que $f'$ est fini.
D'après la remarque générale (III), $f'$ est plat.
D'après Morphismes, lemme \ref{morphisms-lemma-finite-presentation-permanence},
$f'$ est localement de présentation finie. Ainsi, $f'$ est fini localement libre d'après
Morphismes, lemme \ref{morphisms-lemma-finite-flat}.
\end{proof}


\begin{lemma}[Déformations des schémas projectifs]
\label{more-morphisms-lemma-deform-projective}
Soit $f : X \to S$ un morphisme de schémas propre, plat et
de présentation finie. Soit $\mathcal{L}$ un faisceau $f$-ample. Supposons
$S$ quasi-compact. Il existe un entier $d_0 \geq 0$ tel que,
pour tout diagramme cartésien
$$
\vcenter{
\xymatrix{
X \ar[r]_{i'} \ar[d]_f & X' \ar[d]^{f'} \\
S \ar[r]^i & S'
}
}
\quad\text{et}\quad
\begin{matrix}
\text{un }\mathcal{O}_{X'}\text{-module inversible}\\
\mathcal{L}'\text{ avec }\mathcal{L} \cong (i')^*\mathcal{L}'
\end{matrix}
$$
où $S \subset S'$ est un épaississement et où $f'$ est
propre, plat et de présentation finie, on ait
\begin{enumerate}
\item $R^p(f')_*(\mathcal{L}')^{\otimes d} = 0$
pour tous $p > 0$ et $d \geq d_0$,
\item $\mathcal{A}'_d = (f')_*(\mathcal{L}')^{\otimes d}$
est localement libre de rang fini pour $d \geq d_0$,
\item $\mathcal{A}' =
\mathcal{O}_{S'} \oplus \bigoplus_{d \geq d_0} \mathcal{A}'_d$
est une $\mathcal{O}_{S'}$-algèbre quasi-cohérente de présentation finie,
\item il existe un isomorphisme canonique
$r' : X' \to \underline{\text{Proj}}_{S'}(\mathcal{A}')$, et
\item il existe un isomorphisme canonique
$\theta' : (r')^*\mathcal{O}_{\underline{\text{Proj}}_{S'}(\mathcal{A}')}(1)
\to \mathcal{L}'$.
\end{enumerate}
La construction de $\mathcal{A}'$, $r'$, $\theta'$
est fonctorielle en les données $(X', S', i, i', f', \mathcal{L}')$.
\end{lemma}

\begin{proof}
Décrivons d'abord les morphismes $r'$ et $\theta'$.
Remarquons que $\mathcal{L}'$ est $f'$-ample, voir
le lemme \ref{more-morphisms-lemma-thicken-property-relatively-ample}.
Il existe un homomorphisme canonique de
$\mathcal{O}_{S'}$-algèbres graduées quasi-cohérentes
$\mathcal{A}' \to \bigoplus_{d \geq 0} (f')_*(\mathcal{L}')^{\otimes d}$
qui est un isomorphisme en degrés $\geq d_0$.
Il induit donc un isomorphisme des Proj relatifs,
compatible avec les faisceaux structuraux tordus de Serre, voir
Constructions, lemme
\ref{constructions-lemma-eventual-iso-graded-rings-map-relative-proj}.
On obtient ainsi le morphisme $r'$ par
Morphismes, lemme \ref{morphisms-lemma-characterize-relatively-ample}
(qui utilise à son tour la construction donnée dans
Constructions, lemme
\ref{constructions-lemma-invertible-map-into-relative-proj}),
et c'est un isomorphisme d'après
Morphismes, lemme \ref{morphisms-lemma-proper-ample-is-proj}.
L'homomorphisme $\theta'$ est fourni par Constructions, lemme
\ref{constructions-lemma-invertible-map-into-relative-proj}.
D'après Propriétés, lemme \ref{properties-lemma-ample-gcd-is-one},
on voit que $\theta'$ est un isomorphisme
(on utilise aussi le fait que le morphisme $r'$, au-dessus des ouverts
affines de $S'$, est le même que celui de
Propriétés, lemme \ref{properties-lemma-map-into-proj},
comme l'explique la démonstration
de Morphismes, lemme \ref{morphisms-lemma-proper-ample-is-proj}).

\medskip\noindent
En admettant l'annulation et la liberté locale énoncées en
(1) et (2), on peut voir la fonctorialité de la construction comme suit.
Soit $h : T \to S'$ un morphisme de schémas ; notons
$f_T : X'_T \to T$ le changement de base de $f'$ et
$\mathcal{L}_T$ l'image inverse de $\mathcal{L}$ sur $X'_T$.
Par les théorèmes de cohomologie et changement de base
(tels qu'ils sont formulés, par exemple, dans Catégories dérivées des schémas,
lemme \ref{perfect-lemma-compare-base-change}),
on a l'annulation correspondante sur $T$ et, de plus,
$h^*\mathcal{A}'_d = f_{T, *}\mathcal{L}_T^{\otimes d}$
(donc aussi la liberté locale des images directes,
ainsi que le caractère de type fini de la
$\mathcal{O}_T$-algèbre graduée correspondante $\mathcal{A}_T$).
Ainsi, le morphisme
$r_T : X_T \to
\underline{\text{Proj}}_T(\bigoplus f_{T, *}\mathcal{L}_T^{\otimes d})$
est simplement le changement de base de $r'$ à $T$, et l'image inverse de
$\theta'$ est l'homomorphisme $\theta_T$.

\medskip\noindent
Après ces observations, on voit qu'il suffit de prouver
(1), (2) et (3). Choisissons $d_0$ tel que
$R^pf_*\mathcal{L}^{\otimes d} = 0$ pour tous $d \geq d_0$ et $p > 0$, voir
Cohomologie des schémas, lemme \ref{coherent-lemma-coherent-proper-ample}.
On affirme que $d_0$ convient.

\medskip\noindent
Par les théorèmes de cohomologie et changement de base
(Catégories dérivées des schémas,
lemme \ref{perfect-lemma-flat-proper-perfect-direct-image-general}),
on voit que $E'_d = Rf'_*(\mathcal{L}')^{\otimes d}$
est un objet parfait de $D(\mathcal{O}_{S'})$
et que sa formation commute à tout changement de base.
En particulier, $E_d = Li^*E'_d = Rf_*\mathcal{L}^{\otimes d}$.
D'après Catégories dérivées des schémas, lemme
\ref{perfect-lemma-vanishing-implies-locally-free},
on voit que, pour $d \geq d_0$, le complexe $E_d$ est isomorphe au
$\mathcal{O}_S$-module localement libre de rang fini
$f_*\mathcal{L}^{\otimes d}$ placé
en degré cohomologique $0$. D'après
Catégories dérivées des schémas, lemme
\ref{perfect-lemma-open-where-cohomology-in-degree-i-rank-r},
on en déduit que $E'_d$ est isomorphe à un module localement libre de rang fini
placé en degré cohomologique $0$.
Bien entendu, cela signifie que $E'_d = \mathcal{A}'_d[0]$,
que $R^pf'_*(\mathcal{L}')^{\otimes d} = 0$ pour $p > 0$
et que $\mathcal{A}'_d$ est localement libre de rang fini.
Cela prouve (1) et (2).

\medskip\noindent
Il reste à montrer que
$\mathcal{A}'$ est de présentation finie comme faisceau de $\mathcal{O}_{S'}$-algèbres
(cette notion a été introduite dans Propriétés, section
\ref{properties-section-extending-quasi-coherent-sheaves}).
Soit $U' = \Spec(R') \subset S'$ un ouvert affine.
Alors $A' = \mathcal{A}'(U')$ est une $R'$-algèbre graduée
dont les composantes homogènes sont des $R'$-modules projectifs de type fini.
Il faut montrer que $A'$ est une $R'$-algèbre de présentation finie.
On le prouvera en se ramenant au cas noethérien.
En effet, on peut trouver une sous-$\mathbf{Z}$-algèbre de type fini
$R'_0 \subset R'$ et un couple\footnote{Possédant les mêmes propriétés
que $X' \to S'$ et $\mathcal{L}'$, c'est-à-dire
que $X'_0 \to \Spec(R'_0)$ est plat et propre et que $\mathcal{L}'_0$
est ample.} $(X'_0, \mathcal{L}'_0)$ sur $R'_0$
dont le changement de base est $(X'_{U'}, \mathcal{L}'|_{X'_{U'}})$, voir
Limites, lemmes
\ref{limits-lemma-descend-modules-finite-presentation},
\ref{limits-lemma-descend-invertible-modules},
\ref{limits-lemma-eventually-proper},
\ref{limits-lemma-descend-flat-finite-presentation} et
\ref{limits-lemma-limit-ample}.
Cohomologie des schémas, lemme \ref{coherent-lemma-coherent-proper-ample},
entraîne que
$A'_0 = \bigoplus_{d \geq 0} H^0(X'_0, (\mathcal{L}'_0)^{\otimes d})$
est une $R'_0$-algèbre graduée de type fini et
qu'il existe un entier $d'_0$ tel que
$H^p(X'_0, (\mathcal{L}'_0)^{\otimes d}) = 0$, $p > 0$, pour $d \geq d'_0$.
En appliquant les arguments ci-dessus à $X'_0 \to \Spec(R'_0)$ et
$\mathcal{L}'_0$, on voit que $(A'_0)_d$ est un $R'_0$-module projectif de type fini
et que
$$
A'_d = \mathcal{A}'_d(U') =
H^0(X'_{U'}, (\mathcal{L}')^{\otimes d}|_{X'_{U'}}) =
H^0(X'_0, (\mathcal{L}'_0)^{\otimes d}) \otimes_{R'_0} R' =
(A'_0)_d \otimes_{R'_0} R'
$$
pour $d \geq d'_0$. Une petite difficulté est que l'on ne sait pas
si l'on peut choisir $d'_0$ égal à $d_0$\footnote{En fait,
on peut se ramener à ce cas en utilisant davantage d'arguments de passage à la limite.}. Pour
la contourner, on achève la démonstration par la succession d'arguments
suivante :
\begin{enumerate}
\item[(a)] L'algèbre
$B = R'_0 \oplus \bigoplus_{d \geq \max(d_0, d'_0)} (A'_0)_d$
est une $R'_0$-algèbre de type fini : appliquer
le lemme d'Artin--Tate à $B \subset A'_0$, voir
Algèbre, lemme \ref{algebra-lemma-Artin-Tate}.
\item[(b)] Puisque $R'_0$ est noethérien, on voit que
$B$ est une $R'_0$-algèbre de présentation finie.
\item[(c)] Par exactitude à droite du produit tensoriel, on voit que
$B \otimes_{R'_0} R'$ est une $R'$-algèbre de présentation finie.
\item[(d)] D'après les égalités affichées, cela signifie exactement que
$C = R' \oplus \bigoplus_{d \geq \max(d_0, d'_0)} A'_d$
est une $R'$-algèbre de présentation finie.
\item[(e)] Le quotient $A'/C$ est la somme directe des
$R'$-modules projectifs de type fini $A'_d$, $d_0 \leq d \leq \max(d_0, d'_0)$ ;
il est donc de présentation finie comme $R'$-module.
\item[(f)] Le quotient $A'/C$ est de présentation finie
comme $C$-module d'après Algèbre, lemme
\ref{algebra-lemma-finitely-presented-over-subring}.
\item[(g)] Ainsi, $A'$ est de présentation finie comme $C$-module d'après
Algèbre, lemme \ref{algebra-lemma-extension}.
\item[(h)] D'après Algèbre, lemme \ref{algebra-lemma-finite-finite-type},
cela entraîne que $A'$ est de présentation finie comme $C$-algèbre.
\item[(i)] Enfin, d'après
Algèbre, lemme \ref{algebra-lemma-compose-finite-type},
appliqué à $R' \to C \to A'$,
cela entraîne que $A'$ est de présentation finie comme $R'$-algèbre.
\end{enumerate}
Cela achève la démonstration.
\end{proof}












\section{Morphismes formellement lisses}
\label{more-morphisms-section-formally-smooth}

\noindent
Selon Michael Artin, les critères différentiels de lissité (par exemple,
Morphismes, lemme \ref{morphisms-lemma-smooth-at-point}) sont
essentiellement inutiles (en pratique). Dans cette section, on introduit la
notion de morphisme formellement lisse $X \to S$. Un tel morphisme est
caractérisé par la propriété que les points à valeurs dans $T$ de $X$ se relèvent
aux épaississements infinitésimaux de $T$, pourvu que $T$ soit affine.
Le résultat principal est qu'un morphisme formellement lisse et
localement de présentation finie est lisse, voir
le lemme \ref{more-morphisms-lemma-smooth-formally-smooth}.
Il se trouve que ce critère est souvent plus facile à utiliser que les
critères différentiels mentionnés ci-dessus.

\medskip\noindent
Rappelons qu'un homomorphisme d'anneaux $R \to A$ est dit {\it formellement lisse}
(voir Algèbre, définition \ref{algebra-definition-formally-smooth})
si, pour tout diagramme commutatif de flèches pleines
$$
\xymatrix{
A \ar[r] \ar@{-->}[rd] & B/I \\
R \ar[r] \ar[u] & B \ar[u]
}
$$
où $I \subset B$ est un idéal de carré nul, il existe une flèche
pointillée rendant le diagramme commutatif. Cela conduit
à l'analogue suivant pour les morphismes de schémas.

\begin{definition}
\label{more-morphisms-definition-formally-smooth}
Soit $f : X \to S$ un morphisme de schémas.
On dit que $f$ est {\it formellement lisse} si, pour tout diagramme commutatif de flèches pleines
$$
\xymatrix{
X \ar[d]_f & T \ar[d]^i \ar[l] \\
S & T' \ar[l] \ar@{-->}[lu]
}
$$
où $T \subset T'$ est un épaississement du premier ordre de schémas affines sur $S$,
il existe une flèche pointillée rendant le diagramme commutatif.
\end{definition}

\noindent
Dans le cas des morphismes formellement non ramifiés et formellement étales,
on pouvait omettre la condition que $T'$ soit affine, voir
les lemmes \ref{more-morphisms-lemma-formally-unramified-not-affine} et
\ref{more-morphisms-lemma-formally-etale-not-affine}.
Ce n'est plus vrai dans le cas des morphismes formellement lisses.
Une condition légèrement plus naturelle serait, en fait, que l'on puisse
compléter le diagramme par la flèche pointillée localement pour la topologie de Zariski sur $T'$.
L'analyse de la démonstration du
lemme \ref{more-morphisms-lemma-formally-smooth}
montre que cette condition serait équivalente à la définition
actuelle. En particulier, la propriété d'être formellement lisse est
locale sur la source pour la topologie de Zariski (et même pour la topologie lisse ;
insérer ici une référence ultérieure).

\begin{lemma}
\label{more-morphisms-lemma-composition-formally-smooth}
Un composé de morphismes formellement lisses est formellement lisse.
\end{lemma}

\begin{proof}
Omise.
\end{proof}

\begin{lemma}
\label{more-morphisms-lemma-base-change-formally-smooth}
Tout changement de base d'un morphisme formellement lisse est formellement lisse.
\end{lemma}

\begin{proof}
Omise ; voir cependant Algèbre, lemme \ref{algebra-lemma-base-change-fs},
pour la version algébrique.
\end{proof}

\begin{lemma}
\label{more-morphisms-lemma-formally-etale-unramified-smooth}
Soit $f : X \to S$ un morphisme de schémas.
Alors $f$ est formellement étale si et seulement si
$f$ est formellement lisse et formellement non ramifié.
\end{lemma}

\begin{proof}
Omise.
\end{proof}

\begin{lemma}
\label{more-morphisms-lemma-formally-smooth-on-opens}
Soit $f : X \to S$ un morphisme de schémas.
Soient $U \subset X$ et $V \subset S$ des sous-schémas ouverts tels que
$f(U) \subset V$. Si $f$ est formellement lisse, il en va de même de $f|_U : U \to V$.
\end{lemma}

\begin{proof}
Considérons un diagramme de flèches pleines
$$
\xymatrix{
U \ar[d]_{f|_U} & T \ar[d]^i \ar[l]^a \\
V & T' \ar[l] \ar@{-->}[lu]
}
$$
comme dans la définition \ref{more-morphisms-definition-formally-smooth}. Si $f$ est formellement
lisse, il existe un $S$-morphisme $a' : T' \to X$ tel que
$a'|_T = a$. Puisque $T$ et $T'$ ont le même ensemble sous-jacent,
on voit que $a'$ est un morphisme à valeurs dans $U$ (voir Schémas, section
\ref{schemes-section-open-immersion}). C'est manifestement aussi un $V$-morphisme.
On obtient donc la flèche pointillée voulue.
\end{proof}

\begin{lemma}
\label{more-morphisms-lemma-affine-formally-smooth}
Soit $f : X \to S$ un morphisme de schémas.
Supposons que $X$ et $S$ soient affines.
Alors $f$ est formellement lisse si et seulement si
$\mathcal{O}_S(S) \to \mathcal{O}_X(X)$ est un homomorphisme d'anneaux
formellement lisse.
\end{lemma}

\begin{proof}
Cela résulte immédiatement des définitions
(définition \ref{more-morphisms-definition-formally-smooth} et
Algèbre, définition \ref{algebra-definition-formally-smooth}),
par l'équivalence entre les catégories des anneaux et des schémas affines,
voir
Schémas, lemme \ref{schemes-lemma-category-affine-schemes}.
\end{proof}

\noindent
Le lemme suivant est le résultat principal de cette section. C'est une victoire du
point de vue fonctoriel : combiné avec
Limites,
proposition \ref{limits-proposition-characterize-locally-finite-presentation},
il permet de reconnaître si un morphisme $f : X \to S$ est lisse à l'aide de
propriétés « simples » du foncteur $h_X : \Sch/S \to \textit{Ensembles}$.

\begin{lemma}[Critère de relèvement infinitésimal]
\label{more-morphisms-lemma-smooth-formally-smooth}
Soit $f : X \to S$ un morphisme de schémas.
Les conditions suivantes sont équivalentes :
\begin{enumerate}
\item Le morphisme $f$ est lisse, et
\item le morphisme $f$ est localement de présentation finie et
formellement lisse.
\end{enumerate}
\end{lemma}

\begin{proof}
Supposons $f : X \to S$ localement de présentation finie et formellement lisse.
Considérons deux ouverts affines $\Spec(A) = U \subset X$ et
$\Spec(R) = V \subset S$
tels que $f(U) \subset V$. D'après le lemme \ref{more-morphisms-lemma-formally-smooth-on-opens},
on voit que $U \to V$ est formellement lisse. D'après le lemme
\ref{more-morphisms-lemma-affine-formally-smooth}, on voit que $R \to A$ est formellement
lisse. D'après
Morphismes, lemme \ref{morphisms-lemma-locally-finite-presentation-characterize},
on voit que $R \to A$ est de présentation finie.
D'après Algèbre, proposition \ref{algebra-proposition-smooth-formally-smooth},
on voit que $R \to A$ est lisse.
La définition d'un morphisme lisse montre donc que $X \to S$ est lisse.

\medskip\noindent
Réciproquement, supposons que $f : X \to S$ soit lisse. Considérons un diagramme commutatif
de flèches pleines
$$
\xymatrix{
X \ar[d]_f & T \ar[d]^i \ar[l]^a \\
S & T' \ar[l] \ar@{-->}[lu]
}
$$
comme dans la définition \ref{more-morphisms-definition-formally-smooth}.
On montrera que la flèche pointillée existe, ce qui
prouvera que $f$ est formellement lisse.

\medskip\noindent
Soit $\mathcal{F}$ le faisceau d'ensembles sur $T'$ du lemme \ref{more-morphisms-lemma-sheaf},
dans le cas particulier étudié à la remarque \ref{more-morphisms-remark-special-case}.
Soit
$$
\mathcal{H} =
\SheafHom_{\mathcal{O}_T}(a^*\Omega_{X/S}, \mathcal{C}_{T/T'})
$$
le faisceau de $\mathcal{O}_T$-modules muni de l'action
$\mathcal{H} \times \mathcal{F} \to \mathcal{F}$ donnée par
le lemme \ref{more-morphisms-lemma-action-sheaf}. Il s'agit simplement
de montrer que $\mathcal{F}(T) \not = \emptyset$. Autrement dit, on
cherche à montrer que $\mathcal{F}$ est un $\mathcal{H}$-torseur trivial
sur $T$ (voir Cohomologie, section \ref{cohomology-section-h1-torsors}).
Il y a deux étapes : (I) Pour montrer que $\mathcal{F}$ est un torseur,
il faut montrer que $\mathcal{F}_t \not = \emptyset$ pour tout $t \in T$ (voir
Cohomologie, définition \ref{cohomology-definition-torsor}).
(II) Pour montrer que $\mathcal{F}$ est le torseur trivial, il suffit
de montrer que $H^1(T, \mathcal{H}) = 0$ (voir
Cohomologie, lemme \ref{cohomology-lemma-torsors-h1} --
on peut utiliser la cohomologie
de $\mathcal{H}$ comme faisceau abélien ou comme $\mathcal{O}_T$-module,
voir Cohomologie, lemme \ref{cohomology-lemma-modules-abelian}).

\medskip\noindent
Prouvons d'abord (I). Pour tout $t \in T$, on peut
choisir un voisinage ouvert affine $U \subset T$ de $t$
tel que $a(U)$ soit contenu
dans un ouvert affine $\Spec(A) = W \subset X$
qui s'envoie dans un ouvert affine $\Spec(R) = V \subset S$.
D'après Morphismes, lemme \ref{morphisms-lemma-smooth-characterize},
l'homomorphisme d'anneaux $R \to A$ est lisse.
D'après Algèbre, proposition \ref{algebra-proposition-smooth-formally-smooth},
l'homomorphisme d'anneaux $R \to A$ est donc formellement lisse.
Le lemme \ref{more-morphisms-lemma-affine-formally-smooth}
entraîne à son tour que $W \to V$ est formellement lisse.
On peut donc relever $a|_U : U \to W$ en un $V$-morphisme
$a' : U' \to W \subset X$, ce qui montre que $\mathcal{F}(U) \not = \emptyset$.

\medskip\noindent
Prouvons enfin (II).
D'après Morphismes, lemme \ref{morphisms-lemma-finite-presentation-differentials},
on voit que $\Omega_{X/S}$ est de présentation finie
(il est même localement libre de rang fini d'après
Morphismes, lemme \ref{morphisms-lemma-smooth-omega-finite-locally-free}).
Ainsi, $a^*\Omega_{X/S}$ est de présentation finie (voir
Modules, lemme \ref{modules-lemma-pullback-finite-presentation}).
Le faisceau
$\mathcal{H} =
\SheafHom_{\mathcal{O}_T}(a^*\Omega_{X/S}, \mathcal{C}_{T/T'})$
est donc quasi-cohérent d'après la discussion de
Schémas, section \ref{schemes-section-quasi-coherent}.
D'après Cohomologie des schémas, lemme
\ref{coherent-lemma-quasi-coherent-affine-cohomology-zero},
on a donc $H^1(T, \mathcal{H}) = 0$, comme voulu.
\end{proof}

\noindent
Les modules quasi-cohérents localement projectifs sont définis dans
Propriétés, section \ref{properties-section-locally-projective}.

\begin{lemma}
\label{more-morphisms-lemma-formally-smooth-sheaf-differentials}
Soit $f : X \to Y$ un morphisme de schémas formellement lisse.
Alors $\Omega_{X/Y}$ est localement projectif sur $X$.
\end{lemma}

\begin{proof}
Choisissons des ouverts affines $U \subset X$ et $V \subset Y$ tels que
$f(U) \subset V$. D'après
le lemme \ref{more-morphisms-lemma-formally-smooth-on-opens},
$f|_U : U \to V$ est formellement lisse. Ainsi,
$\Gamma(V, \mathcal{O}_V) \to \Gamma(U, \mathcal{O}_U)$ est
un homomorphisme d'anneaux formellement lisse, voir
le lemme \ref{more-morphisms-lemma-affine-formally-smooth}.
D'après
Algèbre, lemme \ref{algebra-lemma-characterize-formally-smooth-again},
le $\Gamma(U, \mathcal{O}_U)$-module
$\Omega_{\Gamma(U, \mathcal{O}_U)/\Gamma(V, \mathcal{O}_V)}$
est donc projectif. Ainsi, $\Omega_{U/V}$ est localement projectif, voir
Propriétés, section \ref{properties-section-locally-projective}.
\end{proof}

\begin{lemma}
\label{more-morphisms-lemma-h1-is-zero}
Soit $T$ un schéma affine. Soient $\mathcal{F}$, $\mathcal{G}$ des
$\mathcal{O}_T$-modules quasi-cohérents. Considérons
$\mathcal{H} = \SheafHom_{\mathcal{O}_T}(\mathcal{F}, \mathcal{G})$.
Si $\mathcal{F}$ est localement projectif, alors $H^1(T, \mathcal{H}) = 0$.
\end{lemma}

\begin{proof}
D'après la définition d'un faisceau localement projectif sur un schéma (voir
Propriétés, définition \ref{properties-definition-locally-projective}),
on voit que $\mathcal{F}$ est un facteur direct d'un
$\mathcal{O}_T$-module libre. On peut donc supposer que
$\mathcal{F} = \bigoplus_{i \in I} \mathcal{O}_T$ est un module libre.
Dans ce cas, $\mathcal{H} = \prod_{i \in I} \mathcal{G}$ est
un produit de modules quasi-cohérents. D'après
Cohomologie, lemme \ref{cohomology-lemma-cohomology-products},
on en déduit que $H^1 = 0$, car la cohomologie d'un faisceau quasi-cohérent
sur un schéma affine est nulle, voir Cohomologie des schémas, lemme
\ref{coherent-lemma-quasi-coherent-affine-cohomology-zero}.
\end{proof}

\begin{lemma}
\label{more-morphisms-lemma-formally-smooth}
Soit $f : X \to Y$ un morphisme de schémas. Les conditions suivantes sont équivalentes :
\begin{enumerate}
\item $f$ est formellement lisse,
\item pour tout $x \in X$, il existe des ouverts $x \in U \subset X$ et
$f(x) \in V \subset Y$ avec $f(U) \subset V$ tels que
$f|_U : U \to V$ soit formellement lisse,
\item pour tout couple d'ouverts affines $U \subset X$ et $V \subset Y$
avec $f(U) \subset V$, l'homomorphisme d'anneaux $\mathcal{O}_Y(V) \to \mathcal{O}_X(U)$
est formellement lisse, et
\item il existe un recouvrement ouvert affine $Y = \bigcup V_j$ et,
pour chaque $j$, un recouvrement ouvert affine $f^{-1}(V_j) = \bigcup U_{ji}$
tels que $\mathcal{O}_Y(V) \to \mathcal{O}_X(U)$ soit un homomorphisme d'anneaux
formellement lisse pour tous $j$ et $i$.
\end{enumerate}
\end{lemma}

\begin{proof}
Les implications (1) $\Rightarrow$ (2),
(1) $\Rightarrow$ (3) et (2) $\Rightarrow$ (4) résultent du
lemme \ref{more-morphisms-lemma-formally-smooth-on-opens}.
L'implication (3) $\Rightarrow$ (4) est immédiate.

\medskip\noindent
Supposons (4). La démonstration du fait que $f$ est formellement lisse est la même
que la seconde partie de la démonstration du lemme \ref{more-morphisms-lemma-smooth-formally-smooth}.
Considérons un diagramme commutatif de flèches pleines
$$
\xymatrix{
X \ar[d]_f & T \ar[d]^i \ar[l]^a \\
Y & T' \ar[l] \ar@{-->}[lu]
}
$$
comme dans la définition \ref{more-morphisms-definition-formally-smooth}.
On montrera que la flèche pointillée existe, ce qui
prouvera que $f$ est formellement lisse.
Soit $\mathcal{F}$ le faisceau d'ensembles sur $T'$ du
lemme \ref{more-morphisms-lemma-sheaf}, dans le cas particulier étudié à
la remarque \ref{more-morphisms-remark-special-case}.
Soit
$$
\mathcal{H} =
\SheafHom_{\mathcal{O}_T}(a^*\Omega_{X/Y}, \mathcal{C}_{T/T'})
$$
le faisceau de $\mathcal{O}_T$-modules sur $T$
muni de l'action $\mathcal{H} \times \mathcal{F} \to \mathcal{F}$ du
lemme \ref{more-morphisms-lemma-action-sheaf}.
L'action $\mathcal{H} \times \mathcal{F} \to \mathcal{F}$
fait de $\mathcal{F}$ un pseudo-$\mathcal{H}$-torseur, voir
Cohomologie, définition \ref{cohomology-definition-torsor}.
Il s'agit de montrer que $\mathcal{F}$ est un $\mathcal{H}$-torseur trivial.
Il y a deux étapes : (I) Pour montrer que $\mathcal{F}$ est un torseur,
il faut montrer que $\mathcal{F}$ possède localement une
section. (II) Pour montrer que $\mathcal{F}$ est le torseur trivial,
il suffit de montrer que $H^1(T, \mathcal{H}) = 0$, voir
Cohomologie, lemme \ref{cohomology-lemma-torsors-h1}.

\medskip\noindent
Prouvons d'abord (I). Pour tout $t \in T$, on peut
choisir un voisinage ouvert affine $W \subset T$ de $t$
tel que $a(W)$ soit contenu dans $U_{ji}$ pour certains $i, j$.
Soit $W' \subset T'$ le sous-schéma ouvert correspondant.
Par l'hypothèse (4), on peut relever $a|_W : W \to U_{ji}$
en un $V_j$-morphisme $a' : W' \to U_{ji}$, ce qui montre que
$\mathcal{F}(W')$ est non vide.

\medskip\noindent
Prouvons enfin (II). D'après
le lemme \ref{more-morphisms-lemma-formally-smooth-sheaf-differentials},
on voit que $\Omega_{U_{ji}/V_j}$ est localement projectif.
Ainsi, $\Omega_{X/Y}$ est localement projectif, voir
Propriétés, lemme \ref{properties-lemma-locally-projective}.
Ainsi, $a^*\Omega_{X/Y}$ est localement projectif, voir
Propriétés, lemme \ref{properties-lemma-locally-projective-pullback}.
On a donc
$$
H^1(T, \mathcal{H}) =
H^1(T, \SheafHom_{\mathcal{O}_T}(a^*\Omega_{X/Y}, \mathcal{C}_{T/T'})) = 0
$$
d'après
le lemme \ref{more-morphisms-lemma-h1-is-zero},
comme voulu.
\end{proof}

\begin{lemma}
\label{more-morphisms-lemma-triangle-differentials-formally-smooth}
Soient $f : X \to Y$, $g : Y \to S$ des morphismes de schémas.
Supposons $f$ formellement lisse. Alors
$$
0 \to f^*\Omega_{Y/S} \to \Omega_{X/S} \to \Omega_{X/Y} \to 0
$$
(voir
Morphismes, lemme \ref{morphisms-lemma-triangle-differentials})
est une suite exacte courte.
\end{lemma}

\begin{proof}
La version algébrique de ce lemme est la suivante :
pour des homomorphismes d'anneaux $A \to B \to C$ avec $B \to C$ formellement lisse,
la suite
$$
0 \to C \otimes_B \Omega_{B/A} \to \Omega_{C/A} \to \Omega_{C/B} \to 0
$$
donnée dans
Algèbre, lemme \ref{algebra-lemma-exact-sequence-differentials},
est exacte. C'est
Algèbre, lemme \ref{algebra-lemma-ses-formally-smooth}.
\end{proof}

\begin{lemma}
\label{more-morphisms-lemma-differentials-formally-unramified-formally-smooth}
Soit $h : Z \to X$ un morphisme formellement non ramifié de schémas sur $S$.
Supposons que $Z$ soit formellement lisse sur $S$. Alors la
suite exacte canonique
$$
0 \to \mathcal{C}_{Z/X} \to h^*\Omega_{X/S} \to \Omega_{Z/S} \to 0
$$
du
lemme \ref{more-morphisms-lemma-universally-unramified-differentials-sequence}
est une suite exacte courte.
\end{lemma}

\begin{proof}
Soit $Z \to Z'$ l'épaississement universel du premier ordre de $Z$ sur $X$.
La démonstration du
lemme \ref{more-morphisms-lemma-universally-unramified-differentials-sequence}
montre que notre suite s'identifie à la suite
$$
\mathcal{C}_{Z/Z'} \to \Omega_{Z'/S} \otimes \mathcal{O}_Z \to
\Omega_{Z/S} \to 0.
$$
Puisque $Z \to S$ est formellement lisse, on peut trouver, localement sur $Z'$,
un inverse à gauche $Z' \to Z$ sur $S$ de l'inclusion $Z \to Z'$.
La suite est donc localement scindée, voir
Morphismes, lemme \ref{morphisms-lemma-differentials-relative-immersion-section}.
\end{proof}

\begin{lemma}
\label{more-morphisms-lemma-two-unramified-morphisms-formally-smooth}
Soit
$$
\xymatrix{
Z \ar[r]_i \ar[rd]_j & X \ar[d]^f \\
& Y
}
$$
un diagramme commutatif de schémas où $i$ et $j$ sont formellement
non ramifiés et où $f$ est formellement lisse. Alors la suite exacte canonique
$$
0 \to
\mathcal{C}_{Z/Y} \to
\mathcal{C}_{Z/X} \to
i^*\Omega_{X/Y} \to 0
$$
du
lemme \ref{more-morphisms-lemma-two-unramified-morphisms}
est exacte et localement scindée.
\end{lemma}

\begin{proof}
Notons $Z \to Z'$ l'épaississement universel du premier ordre de $Z$ sur $X$.
Notons $Z \to Z''$ l'épaississement universel du premier ordre de $Z$ sur $Y$.
D'après
le lemme \ref{more-morphisms-lemma-universally-unramified-differentials-sequence},
il existe un morphisme canonique $Z' \to Z''$ donnant un diagramme
commutatif
$$
\xymatrix{
Z \ar[r]_{i'} \ar[rd]_{j'} & Z' \ar[r]_a \ar[d]^k & X \ar[d]^f \\
& Z'' \ar[r]^b & Y
}
$$
Dans la démonstration du
lemme \ref{more-morphisms-lemma-two-unramified-morphisms},
on a identifié la suite ci-dessus à la suite
$$
\mathcal{C}_{Z/Z''} \to
\mathcal{C}_{Z/Z'} \to
(i')^*\Omega_{Z'/Z''} \to 0
$$
Soit $U'' \subset Z''$ un ouvert affine. Notons $U \subset Z$ et
$U' \subset Z'$ les sous-schémas ouverts affines correspondants.
Comme $f$ est formellement lisse, il existe un morphisme $h : U'' \to X$
qui coïncide avec $i$ sur $U$ et tel que $f \circ h$ soit égal à $b|_{U''}$.
Puisque $Z'$ est l'épaississement universel du premier ordre, on obtient un unique
morphisme $g : U'' \to Z'$ tel que $g = a \circ h$. La propriété
universelle de $Z''$ entraîne que $k \circ g$ est l'inclusion
$U'' \to Z''$. Ainsi, $g$ est un inverse à droite de $k$. Voici le diagramme :
$$
\xymatrix{
U \ar[d] \ar[r] & Z' \ar[d]^k \\
U'' \ar[r] \ar[ru]^g & Z''
}
$$
Ainsi, $g$ induit un homomorphisme $\mathcal{C}_{Z/Z'}|_U \to \mathcal{C}_{Z/Z''}|_U$
qui est un inverse à gauche de l'homomorphisme
$\mathcal{C}_{Z/Z''} \to \mathcal{C}_{Z/Z'}$ sur $U$.
\end{proof}




































\section{Lissité sur une base noethérienne}
\label{more-morphisms-section-smooth-Noetherian}

\noindent
Lorsque la base est noethérienne, on peut alléger
la formulation de la lissité formelle. En un certain sens, les lemmes suivants
constituent les premiers éléments de la théorie des déformations.

\begin{lemma}
\label{more-morphisms-lemma-lifting-along-artinian-at-point}
Soit $f : X \to S$ un morphisme de schémas.
Soit $x \in X$.
Supposons $S$ localement noethérien et $f$ localement de type fini.
Les conditions suivantes sont équivalentes :
\begin{enumerate}
\item $f$ est lisse en $x$,
\item pour tout diagramme commutatif de flèches pleines
$$
\xymatrix{
X \ar[d]_f & \Spec(B) \ar[d]^i \ar[l]^-\alpha \\
S & \Spec(B') \ar[l]_-{\beta} \ar@{-->}[lu]
}
$$
où $B' \to B$ est une surjection d'anneaux locaux telle que
$\Ker(B' \to B)$ soit de carré nul et où $\alpha$ envoie le
point fermé de $\Spec(B)$ sur $x$, il existe
une flèche pointillée rendant le diagramme commutatif,
\item même condition qu'en (2), mais en faisant parcourir à $B' \to B$ les petites
extensions (voir Algèbre, définition \ref{algebra-definition-small-extension}),
et
\item même condition qu'en (2), mais en faisant parcourir à $B' \to B$ les petites
extensions telles que $\alpha$ induise un isomorphisme
$\kappa(x) \to \kappa(\mathfrak m)$, où $\mathfrak m \subset B$
est l'idéal maximal.
\end{enumerate}
\end{lemma}

\begin{proof}
Choisissons un voisinage affine $V \subset S$ de $f(x)$ et un
voisinage affine $U \subset X$ de $x$ tel que $f(U) \subset V$.
Pour tout diagramme « test » comme en (2), le morphisme $\alpha$ enverra
$\Spec(B)$ dans $U$ et le morphisme $\beta$ enverra $\Spec(B')$
dans $V$ (voir Schémas, section \ref{schemes-section-points}).
Le lemme se ramène donc au morphisme $f|_U : U \to V$ de schémas affines.
(En effet, $V$ est noethérien et $f|_U$ est de type fini, voir
Propriétés, lemme \ref{properties-lemma-locally-Noetherian}, et
Morphismes, lemme \ref{morphisms-lemma-locally-finite-type-characterize}.)
Dans ce cas affine, le lemme est identique à
Algèbre, lemme \ref{algebra-lemma-smooth-test-artinian}.
\end{proof}

\noindent
Il est parfois utile de savoir qu'il suffit de vérifier le
critère de relèvement pour les petites extensions « centrées » en des points
de type fini (voir
Morphismes, section \ref{morphisms-section-points-finite-type}).

\begin{lemma}
\label{more-morphisms-lemma-lifting-along-artinian}
Soit $f : X \to S$ un morphisme de schémas.
Supposons $S$ localement noethérien et $f$ localement de type fini.
Les conditions suivantes sont équivalentes :
\begin{enumerate}
\item $f$ est lisse,
\item pour tout diagramme commutatif de flèches pleines
$$
\xymatrix{
X \ar[d]_f & \Spec(B) \ar[d]^i \ar[l]^-\alpha \\
S & \Spec(B') \ar[l]_-{\beta} \ar@{-->}[lu]
}
$$
où $B' \to B$ est une petite extension d'anneaux locaux artiniens
et où $\beta$ est de type fini (!), il existe une flèche pointillée rendant
le diagramme commutatif.
\end{enumerate}
\end{lemma}

\begin{proof}
Si $f$ est lisse, le critère de relèvement infinitésimal
(lemme \ref{more-morphisms-lemma-smooth-formally-smooth}) affirme que
$f$ est formellement lisse, et (2) est vérifiée.

\medskip\noindent
Supposons (2). L'ensemble des points $x \in X$ où $f$ n'est pas lisse
forme un fermé $T$ de $X$. D'après la discussion de Morphismes,
section \ref{morphisms-section-points-finite-type}, si $T \not = \emptyset$,
il existe un point $x \in T \subset X$ tel que le morphisme
$$
\Spec(\kappa(x)) \to X \to S
$$
soit de type fini (il suffit de choisir un point $x$ de $T$ qui soit fermé
dans un ouvert affine de $X$). D'après
Morphismes, lemme \ref{morphisms-lemma-artinian-finite-type}, pour tout
anneau local artinien $B'$ de corps résiduel $\kappa(x)$, tout
morphisme $\beta : \Spec(B') \to S$ est de type fini. Ainsi,
on voit que tous les diagrammes utilisés dans
le lemme \ref{more-morphisms-lemma-lifting-along-artinian-at-point} (4) correspondent
à des diagrammes du présent lemme (2). Donc $X \to S$ est lisse
en $x$, contradiction.
\end{proof}

\noindent
Voici une application utile.

\begin{lemma}
\label{more-morphisms-lemma-check-smoothness-on-infinitesimal-nbhds}
Soit $f : X \to S$ un morphisme de type fini entre schémas localement noethériens.
Soit $Z \subset S$ un sous-schéma fermé de voisinage infinitésimal d'ordre $n$
$Z_n \subset S$. Posons $X_n = Z_n \times_S X$.
\begin{enumerate}
\item Si $X_n \to Z_n$ est lisse pour tout $n$, alors $f$
est lisse en tout point de $f^{-1}(Z)$.
\item Si $X_n \to Z_n$ est étale pour tout $n$, alors $f$
est étale en tout point de $f^{-1}(Z)$.
\end{enumerate}
\end{lemma}

\begin{proof}
Supposons $X_n \to Z_n$ lisse pour tout $n$.
Soit $x \in X$ un point au-dessus d'un point de $Z$.
Étant donnés une petite extension $B' \to B$ et des morphismes $\alpha$, $\beta$ comme dans
le lemme \ref{more-morphisms-lemma-lifting-along-artinian-at-point} (3),
l'idéal maximal de $B'$ est nilpotent (puisque $B'$ est artinien) ;
le morphisme $\beta$ se factorise donc par $Z_n$ et $\alpha$
par $X_n$ pour un entier $n$ convenable. La propriété de relèvement de
$X_n \to Z_n$ fournit alors la flèche pointillée voulue.
Cela prouve (1). Le point (2) résulte de (1) et du fait qu'un morphisme
est étale si et seulement s'il est lisse de dimension relative $0$.
\end{proof}

\begin{lemma}
\label{more-morphisms-lemma-check-flatness-on-infinitesimal-nbhds}
Soit $f : X \to S$ un morphisme de schémas localement noethériens.
Soit $Z \subset S$ un sous-schéma fermé de voisinage infinitésimal d'ordre $n$
$Z_n \subset S$. Posons $X_n = Z_n \times_S X$.
Si $X_n \to Z_n$ est plat pour tout $n$, alors $f$
est plat en tout point de $f^{-1}(Z)$.
\end{lemma}

\begin{proof}
C'est la traduction d'Algèbre, lemme \ref{algebra-lemma-flat-module-powers},
dans le langage des schémas.
\end{proof}










\section{Le complexe cotangent naïf}
\label{more-morphisms-section-netherlander}

\noindent
Cette section poursuit
Modules, section \ref{modules-section-netherlander},
qui prolonge elle-même la discussion d'Algèbre,
section \ref{algebra-section-netherlander}.

\begin{definition}
\label{more-morphisms-definition-netherlander}
Soit $f : X \to Y$ un morphisme de schémas.
Le {\it complexe cotangent naïf de $f$}
est le complexe défini dans Modules, définition
\ref{modules-definition-cotangent-complex-morphism-ringed-topoi}.
Notation : $\NL_f$ ou $\NL_{X/Y}$.
\end{definition}

\begin{lemma}
\label{more-morphisms-lemma-NL-affine}
Soit $f : X \to Y$ un morphisme de schémas. Soient
$\Spec(A) = U \subset X$ et $\Spec(R) = V \subset S$
des ouverts affines tels que $f(U) \subset V$.
Il existe un morphisme canonique
$$
\widetilde{\NL_{A/R}} \longrightarrow \NL_{X/Y}|_U
$$
de complexes, qui est un isomorphisme dans $D(\mathcal{O}_U)$.
\end{lemma}

\begin{proof}
La construction de $\NL_{X/Y}$ dans
Modules, section \ref{modules-section-netherlander},
fournit un morphisme canonique de complexes
$\NL_{\mathcal{O}_X(U)/f^{-1}\mathcal{O}_Y(U)} \to \NL_{X/Y}(U)$
de $A = \mathcal{O}_X(U)$-modules, compatible
aux restrictions ultérieures. À l'aide de l'homomorphisme canonique
$R \to f^{-1}\mathcal{O}_Y(U)$, on obtient un morphisme canonique
$\NL_{A/R} \to \NL_{\mathcal{O}_X(U)/f^{-1}\mathcal{O}_Y(U)}$
de complexes de $A$-modules.
Par la propriété universelle du foncteur $\widetilde{\ }$
(voir Schémas, lemme \ref{schemes-lemma-compare-constructions}),
on obtient un morphisme comme dans l'énoncé du lemme.
Pour vérifier qu'il induit des isomorphismes sur les faisceaux de cohomologie,
on peut vérifier qu'il induit des isomorphismes sur leurs fibres.
Cela résulte des résultats d'Algèbre,
lemmes \ref{algebra-lemma-NL-localize-bottom} et
\ref{algebra-lemma-localize-NL},
et de
Modules, lemme \ref{modules-lemma-stalk-NL}
(ainsi que de la description des fibres de
$\mathcal{O}_X$ et $f^{-1}\mathcal{O}_Y$
en un point $\mathfrak p \in \Spec(A)$ comme $A_\mathfrak p$ et
$R_\mathfrak q$, où $\mathfrak q = R \cap \mathfrak p$ ; les références
utilisées sont Schémas, lemme \ref{schemes-lemma-spec-sheaves},
et
Faisceaux, lemme \ref{sheaves-lemma-stalk-pullback}).
\end{proof}

\begin{lemma}
\label{more-morphisms-lemma-netherlander-quasi-coherent}
Soit $f : X \to Y$ un morphisme de schémas. Les faisceaux de cohomologie
du complexe $\NL_{X/Y}$ sont quasi-cohérents, nuls en dehors des
degrés $-1$, $0$, et égaux à $\Omega_{X/Y}$ en degré $0$.
\end{lemma}

\begin{proof}
Par construction du complexe cotangent naïf dans
Modules, section \ref{modules-section-netherlander},
$\NL_{X/Y}$ est un complexe placé en degrés $-1$, $0$,
et sa cohomologie en degré $0$ est $\Omega_{X/Y}$.
Le faisceau des différentielles est quasi-cohérent (d'après
Morphismes, lemme \ref{morphisms-lemma-differentials-diagonal}).
Pour achever la démonstration, il suffit de montrer que $H^{-1}(\NL_{X/Y})$
est quasi-cohérent. Cela se vérifie sur les ouverts affines
à l'aide du lemme \ref{more-morphisms-lemma-NL-affine}.
\end{proof}

\begin{lemma}
\label{more-morphisms-lemma-netherlander-fp}
Soit $f : X \to Y$ un morphisme de schémas. Si $f$ est localement de
présentation finie, alors $\NL_{X/Y}$ est, localement sur $X$, quasi-isomorphe à
un complexe
$$
\ldots \to 0 \to \mathcal{F}^{-1} \to \mathcal{F}^0 \to 0 \to \ldots
$$
de $\mathcal{O}_X$-modules quasi-cohérents,
où $\mathcal{F}^0$ est de présentation finie
et $\mathcal{F}^{-1}$ de type fini.
\end{lemma}

\begin{proof}
D'après le lemme \ref{more-morphisms-lemma-NL-affine}, il suffit de montrer que $\NL_{A/R}$
a cette forme si $R \to A$ est un homomorphisme d'anneaux de présentation finie.
Écrivons $A = R[x_1, \ldots, x_n]/I$ avec $I$ de type fini.
Alors $I/I^2$ est un
$A$-module de type fini et $\NL_{A/R}$ est quasi-isomorphe à
$$
\ldots \to 0 \to I/I^2 \to
\bigoplus\nolimits_{i = 1, \ldots, n} A\text{d}x_i \to 0 \to \ldots
$$
d'après Algèbre, section \ref{algebra-section-netherlander},
et en particulier
Algèbre, lemme \ref{algebra-lemma-NL-homotopy}.
\end{proof}

\begin{lemma}
\label{more-morphisms-lemma-NL-formally-smooth}
Soit $f : X \to Y$ un morphisme de schémas. Les conditions suivantes sont équivalentes :
\begin{enumerate}
\item $f$ est formellement lisse,
\item $H^{-1}(\NL_{X/Y}) = 0$ et $H^0(\NL_{X/Y}) = \Omega_{X/Y}$
est localement projectif.
\end{enumerate}
\end{lemma}

\begin{proof}
Cela résulte d'Algèbre, proposition
\ref{algebra-proposition-characterize-formally-smooth},
et du lemme \ref{more-morphisms-lemma-formally-smooth}.
\end{proof}

\begin{lemma}
\label{more-morphisms-lemma-NL-formally-etale}
Soit $f : X \to Y$ un morphisme de schémas. Les conditions suivantes sont équivalentes :
\begin{enumerate}
\item $f$ est formellement étale,
\item $H^{-1}(\NL_{X/Y}) = H^0(\NL_{X/Y}) = 0$.
\end{enumerate}
\end{lemma}

\begin{proof}
Un morphisme formellement étale est formellement lisse ; on a donc
$H^{-1}(\NL_{X/Y}) = 0$ d'après le lemme \ref{more-morphisms-lemma-NL-formally-smooth}.
D'autre part, on a $\Omega_{X/Y} = 0$ d'après
le lemme \ref{more-morphisms-lemma-characterize-formally-etale}.
Réciproquement, si (2) est vérifiée, alors $f$ est formellement lisse d'après
le lemme \ref{more-morphisms-lemma-NL-formally-smooth}
et formellement non ramifié d'après
le lemme \ref{more-morphisms-lemma-formally-unramified-differentials},
donc formellement étale d'après
le lemme \ref{more-morphisms-lemma-formally-etale-unramified-smooth}.
\end{proof}

\begin{lemma}
\label{more-morphisms-lemma-NL-smooth}
Soit $f : X \to Y$ un morphisme de schémas. Les conditions suivantes sont équivalentes :
\begin{enumerate}
\item $f$ est lisse, et
\item $f$ est localement de présentation finie,
$H^{-1}(\NL_{X/Y}) = 0$, et $H^0(\NL_{X/Y}) = \Omega_{X/Y}$
est localement libre de rang fini.
\end{enumerate}
\end{lemma}

\begin{proof}
Cela résulte de la définition d'un homomorphisme d'anneaux lisse
(Algèbre, définition \ref{algebra-definition-smooth}),
du lemme \ref{more-morphisms-lemma-NL-affine} et
de la définition d'un morphisme de schémas lisse
(Morphismes, définition \ref{morphisms-definition-smooth}).
On utilise aussi le fait que, pour les modules sur un anneau,
être localement libre de rang fini équivaut à être projectif de type fini
(Algèbre, lemme \ref{algebra-lemma-finite-projective}).
\end{proof}

\begin{lemma}
\label{more-morphisms-lemma-NL-etale}
Soit $f : X \to Y$ un morphisme de schémas. Les conditions suivantes sont équivalentes :
\begin{enumerate}
\item $f$ est étale, et
\item $f$ est localement de présentation finie et
$H^{-1}(\NL_{X/Y}) = H^0(\NL_{X/Y}) = 0$.
\end{enumerate}
\end{lemma}

\begin{proof}
Cela résulte de la définition d'un homomorphisme d'anneaux étale
(Algèbre, définition \ref{algebra-definition-etale}),
du lemme \ref{more-morphisms-lemma-NL-affine} et
de la définition d'un morphisme de schémas étale
(Morphismes, définition \ref{morphisms-definition-etale}).
\end{proof}


\begin{lemma}
\label{more-morphisms-lemma-NL-immersion}
Soit $i : Z \to X$ une immersion de schémas. Alors $\NL_{Z/X}$
est isomorphe à $\mathcal{C}_{Z/X}[1]$ dans $D(\mathcal{O}_Z)$,
où $\mathcal{C}_{Z/X}$ est le faisceau conormal de $Z$ dans $X$.
\end{lemma}

\begin{proof}
Cela résulte d'Algèbre, lemme \ref{algebra-lemma-NL-surjection},
de Morphismes, lemme \ref{morphisms-lemma-affine-conormal}, et
du lemme \ref{more-morphisms-lemma-NL-affine}.
\end{proof}

\begin{lemma}
\label{more-morphisms-lemma-exact-sequence-NL}
Soient $f : X \to Y$ et $g : Y \to Z$ des morphismes de schémas.
Il existe une suite exacte canonique à six termes
$$
H^{-1}(f^*\NL_{Y/Z}) \to
H^{-1}(\NL_{X/Z}) \to
H^{-1}(\NL_{X/Y}) \to
f^*\Omega_{Y/Z} \to \Omega_{X/Z} \to \Omega_{X/Y} \to 0
$$
de faisceaux de cohomologie.
\end{lemma}

\begin{proof}
C'est un cas particulier de
Modules, lemme \ref{modules-lemma-exact-sequence-NL-ringed-topoi}.
\end{proof}

\begin{lemma}
\label{more-morphisms-lemma-get-triangle-NL}
Soient $f : X \to Y$ et $Y \to Z$ des morphismes de schémas. Supposons que
$X \to Y$ soit un morphisme d'intersection complète. Il existe alors
un triangle distingué canonique
$$
f^*\NL_{Y/Z} \to \NL_{X/Z} \to \NL_{X/Y} \to f^*\NL_{Y/Z}[1]
$$
dans $D(\mathcal{O}_X)$, qui redonne la suite exacte à $6$ termes du
lemme \ref{more-morphisms-lemma-exact-sequence-NL}.
\end{lemma}

\begin{proof}
Il suffit de montrer que le morphisme canonique
$$
f^*\NL_{Y/Z} \to \text{Cone}(\NL_{X/Z} \to \NL_{X/Y})[-1]
$$
de Modules, lemme \ref{modules-lemma-exact-sequence-NL-ringed-topoi},
est un isomorphisme dans $D(\mathcal{O}_X)$. Pour cela,
il suffit de montrer que la suite à $6$ termes commence par
un zéro à gauche, c'est-à-dire que $H^{-1}(f^*\NL_{Y/Z}) \to H^{-1}(\NL_{X/Z})$
est injectif. Localement sur les ouverts affines, cela résulte du
résultat algébrique correspondant dans Compléments d'algèbre, lemme
\ref{more-algebra-lemma-transitive-lci-at-end}.
Pour effectuer cette traduction algébrique, utiliser le lemme \ref{more-morphisms-lemma-NL-affine}.
\end{proof}

\begin{lemma}
\label{more-morphisms-lemma-smooth-etale-permanence}
Soient $X \to Y \to Z$ des morphismes de schémas. Supposons $X \to Z$ lisse
et $Y \to Z$ étale. Alors $X \to Y$ est lisse.
\end{lemma}

\begin{proof}
Le morphisme $X \to Y$ est localement de présentation finie d'après
Morphismes, lemme \ref{morphisms-lemma-finite-presentation-permanence}.
D'après le lemme \ref{more-morphisms-lemma-NL-smooth}, on a $H^{-1}(\NL_{X/Z}) = 0$
et le module $\Omega_{X/Z}$ est localement libre de rang fini.
D'après le lemme \ref{more-morphisms-lemma-NL-etale}, on a
$H^{-1}(\NL_{Y/Z}) = H^0(\NL_{Y/Z}) = 0$.
D'après le lemme \ref{more-morphisms-lemma-exact-sequence-NL}, on obtient
$H^{-1}(\NL_{X/Y}) = 0$ et $\Omega_{X/Y} \cong \Omega_{X/Z}$
est localement libre de rang fini.
D'après le lemme \ref{more-morphisms-lemma-NL-smooth}, le morphisme $X \to Y$ est lisse.
\end{proof}

\begin{lemma}
\label{more-morphisms-lemma-get-NL}
Soit $f : X \to Y$ un morphisme de schémas qui se factorise
en $f = g \circ i$, où $i$ est une immersion et $g : P \to Y$
est formellement lisse (par exemple lisse). Il existe alors un isomorphisme canonique
$$
\NL_{X/Y} \cong \left(\mathcal{C}_{X/P} \to i^*\Omega_{P/Y}\right)
$$
dans $D(\mathcal{O}_X)$, où le faisceau conormal $\mathcal{C}_{X/P}$
est placé en degré $-1$.
\end{lemma}

\begin{proof}
(Pour l'assertion entre parenthèses, voir le lemme \ref{more-morphisms-lemma-smooth-formally-smooth}.)
D'après les lemmes \ref{more-morphisms-lemma-NL-immersion} et \ref{more-morphisms-lemma-NL-formally-smooth}, on a
$\NL_{X/P} = \mathcal{C}_{X/P}[1]$ et $\NL_{P/Y} = \Omega_{P/Y}$, avec
$\Omega_{P/Y}$ localement projectif. Cela entraîne que
$i^*\NL_{P/Y} \to i^*\Omega_{P/Y}$ est aussi un quasi-isomorphisme
(un petit détail est omis ; la raison est que $i^*\NL_{P/Y}$ est
la même chose que $\tau_{\geq -1}Li^*\NL_{P/Y}$, voir Compléments d'algèbre, lemme
\ref{more-algebra-lemma-tensor-NL}).
Ainsi, le morphisme canonique
$$
i^*\NL_{P/Y} \to \text{Cone}(\NL_{X/Y} \to \NL_{X/P})[-1]
$$
de Modules, lemme \ref{modules-lemma-exact-sequence-NL-ringed-topoi},
est un isomorphisme dans $D(\mathcal{O}_X)$, car le groupe de cohomologie
$H^{-1}(i^*\NL_{P/Y})$ est nul d'après ce qui précède.
Autrement dit, on dispose d'un triangle distingué
$$
i^*\NL_{P/Y} \to \NL_{X/Y} \to \NL_{X/P} \to i^*\NL_{P/Y}[1]
$$
Cela signifie manifestement que $\NL_{X/Y}$ est le cône du morphisme
$\NL_{X/P}[-1] \to i^*\NL_{P/Y}$, ce qui équivaut
à l'énoncé du lemme, d'après le calcul effectué ci-dessus des faisceaux
de cohomologie de ces objets de la catégorie dérivée.
\end{proof}

\begin{lemma}
\label{more-morphisms-lemma-base-change-NL}
Considérons un diagramme cartésien de schémas
$$
\xymatrix{
X' \ar[r]_{g'} \ar[d] & X \ar[d] \\
Y' \ar[r] & Y
}
$$
Le morphisme canonique $(g')^*\NL_{X/Y} \to \NL_{X'/Y'}$ induit
un isomorphisme sur $H^0$ et une surjection sur $H^{-1}$.
\end{lemma}

\begin{proof}
La traduction algébrique de cet énoncé est
Compléments d'algèbre, lemme \ref{more-algebra-lemma-base-change-NL}.
Pour effectuer la traduction, utiliser le lemme \ref{more-morphisms-lemma-NL-affine}.
\end{proof}

\begin{lemma}
\label{more-morphisms-lemma-flat-base-change-NL}
Considérons un diagramme cartésien de schémas
$$
\xymatrix{
X' \ar[d] \ar[r]_{g'} & X \ar[d] \\
Y' \ar[r] & Y
}
$$
Si $Y' \to Y$ est plat, alors le morphisme canonique
$(g')^*\NL_{X/Y} \to \NL_{X'/Y'}$ est un quasi-isomorphisme.
\end{lemma}

\begin{proof}
D'après le lemme \ref{more-morphisms-lemma-NL-affine}, cela résulte d'Algèbre,
lemme \ref{algebra-lemma-change-base-NL}.
\end{proof}

\begin{lemma}
\label{more-morphisms-lemma-base-change-NL-flat}
Considérons un diagramme cartésien de schémas
$$
\xymatrix{
X' \ar[r]_{g'} \ar[d] & X \ar[d] \\
Y' \ar[r] & Y
}
$$
Si $X \to Y$ est plat, alors le morphisme canonique
$(g')^*\NL_{X/Y} \to \NL_{X'/Y'}$ est un quasi-isomorphisme.
Si, de plus, $\NL_{X/Y}$ a une amplitude de Tor contenue dans $[-1, 0]$,
alors $L(g')^*\NL_{X/Y} \to \NL_{X'/Y'}$ est aussi un quasi-isomorphisme.
\end{lemma}

\begin{proof}
La traduction algébrique de cet énoncé est
Compléments d'algèbre, lemme \ref{more-algebra-lemma-base-change-NL-flat}.
Pour effectuer la traduction, utiliser le lemme \ref{more-morphisms-lemma-NL-affine}
et Catégories dérivées des schémas, lemmes
\ref{perfect-lemma-affine-compare-bounded} et
\ref{perfect-lemma-tor-dimension-affine}.
\end{proof}














\section{Sommes amalgamées dans la catégorie des schémas, I}
\label{more-morphisms-section-pushouts}

\noindent
Dans cette section, on construit les sommes amalgamées de $Y \leftarrow X \rightarrow X'$
lorsque $X \to Y$ est affine et que $X \to X'$ est un épaississement. Ce cas
jouera un rôle important par la suite, ce qui justifie une discussion détaillée.
Dans la section \ref{more-morphisms-section-pushouts-II}, on traite un cas plus intéressant
et plus difficile. Voir Catégories, section
\ref{categories-section-pushouts}, pour une discussion générale
des sommes amalgamées dans une catégorie quelconque.

\begin{lemma}
\label{more-morphisms-lemma-basic-example-pushout}
Soient $A' \to A$ une surjection d'anneaux et $B \to A$ un homomorphisme d'anneaux.
Soit $B' = B \times_A A'$ le produit fibré d'anneaux. Posons
$S = \Spec(A)$, $S' = \Spec(A')$, $T = \Spec(B)$ et $T' = \Spec(B')$.
Alors
$$
\vcenter{
\xymatrix{
S \ar[r]_i \ar[d]_f & S' \ar[d]^{f'} \\
T \ar[r]^{i'} & T'
}
}
\quad\text{correspondant à}\quad
\vcenter{
\xymatrix{
A & A' \ar[l] \\
B \ar[u] & B' \ar[l] \ar[u]
}
}
$$
est une somme amalgamée de schémas.
\end{lemma}

\begin{proof}
D'après Compléments d'algèbre, lemme \ref{more-algebra-lemma-points-of-fibre-product},
on a $T' = T \amalg_S S'$ comme espaces topologiques, c'est-à-dire que le diagramme
est une somme amalgamée dans la catégorie des espaces topologiques. Considérons ensuite
l'homomorphisme
$$
((i')^\sharp, (f')^\sharp) :
\mathcal{O}_{T'}
\longrightarrow
i'_*\mathcal{O}_T \times_{g_*\mathcal{O}_S} f'_*\mathcal{O}_{S'}
$$
où $g = i' \circ f = f' \circ i$. On affirme que cet homomorphisme est un isomorphisme de
faisceaux d'anneaux. En effet, on peut considérer les deux membres comme des
$\mathcal{O}_{T'}$-modules quasi-cohérents (utiliser
Schémas, lemme \ref{schemes-lemma-push-forward-quasi-coherent},
pour le membre de droite), et l'homomorphisme est $\mathcal{O}_{T'}$-linéaire.
Il suffit donc de montrer qu'il est un isomorphisme au niveau
des sections globales
(Schémas, lemme \ref{schemes-lemma-equivalence-quasi-coherent}).
Sur les sections globales, on retrouve l'identification
$B' \to B \times_A A'$ de l'énoncé du lemme (c'est ainsi
que l'on a choisi $B'$).

\medskip\noindent
Soit $X$ un schéma. Supposons donnés des morphismes de schémas
$m' : S' \to X$ et $n : T \to X$ tels que $m' \circ i = n \circ f$
(notons ce morphisme $m$). On obtient une unique application d'espaces topologiques
$n' : T' \to X$ compatible avec $m'$ et $n$, puisque
$T' = T \amalg_S S'$ (voir ci-dessus).
D'après la description de $\mathcal{O}_{T'}$ au paragraphe
précédent, on obtient un unique homomorphisme de faisceaux d'anneaux
$$
(n')^\sharp :
\mathcal{O}_X
\longrightarrow
(n')_*\mathcal{O}_{T'} =
m'_*\mathcal{O}_T \times_{m_*\mathcal{O}_T} n_*\mathcal{O}_S
$$
donné par $(m')^\sharp$ et $n^\sharp$.
Ainsi, $(n', (n')^\sharp)$ est l'unique morphisme
d'espaces annelés $T' \to X$ compatible avec $m'$ et $n$.
Pour achever la démonstration, il suffit de montrer que $n'$
est un morphisme de schémas, c'est-à-dire un morphisme d'espaces localement annelés.

\medskip\noindent
Soit $t' \in T'$ d'image $x \in X$.
Il faut montrer que $\mathcal{O}_{X, x} \to \mathcal{O}_{T', t'}$
est local. Si $t' \not \in T$, alors $t'$ est l'image d'un unique
point $s' \in S'$ et $\mathcal{O}_{T', t'} = \mathcal{O}_{S', s'}$.
En effet, $S' \setminus S \to T' \setminus T$ est un isomorphisme de
schémas, car $B' \to A'$ induit un isomorphisme
$\Ker(B' \to B) = \Ker(A' \to A)$.
Si $t'$ est l'image de $t \in T$,
on sait que le composé
$\mathcal{O}_{X, x} \to \mathcal{O}_{T', t'} \to \mathcal{O}_{T, t}$
est local, ce qui permet aussi de conclure.
\end{proof}

\begin{lemma}
\label{more-morphisms-lemma-pushout-fpqc-local}
Soit $\mathcal{I} \to (\Sch/S)_{fppf}$, $i \mapsto X_i$, un diagramme de schémas.
Soit $(W, X_i \to W)$ un cocône sur ce diagramme dans la catégorie des schémas
(Catégories, remarque \ref{categories-remark-cones-and-cocones}).
S'il existe un recouvrement fpqc $\{W_a \to W\}_{a \in A}$ de schémas tel que
\begin{enumerate}
\item pour tout $a \in A$, on ait
$W_a = \colim X_i \times_W W_a$
dans la catégorie des schémas, et
\item pour tous $a, b \in A$, on ait
$W_a \times_W W_b = \colim X_i \times_W W_a \times_W W_b$
dans la catégorie des schémas,
\end{enumerate}
alors $W = \colim X_i$ dans la catégorie des schémas.
\end{lemma}

\begin{proof}
En effet, pour un schéma $T$, un morphisme $W \to T$ revient à se donner
une famille de morphismes $W_a \to T$, $a \in A$, qui coïncident sur les
intersections $W_a \times_W W_b$, voir
Descente, lemme \ref{descent-lemma-fpqc-universal-effective-epimorphisms}.
\end{proof}

\begin{lemma}
\label{more-morphisms-lemma-pushout-along-thickening}
Soit $X \to X'$ un épaississement de schémas et soit $X \to Y$ un morphisme
affine de schémas. Il existe alors une somme amalgamée
$$
\xymatrix{
X \ar[r] \ar[d]_f
&
X' \ar[d]^{f'}
\\
Y \ar[r]
&
Y'
}
$$
dans la catégorie des schémas. De plus, $Y \subset Y'$ est un
épaississement, $X = Y \times_{Y'} X'$, et
$$
\mathcal{O}_{Y'} = \mathcal{O}_Y \times_{f_*\mathcal{O}_X} f'_*\mathcal{O}_{X'}
$$
comme faisceaux sur $|Y| = |Y'|$.
\end{lemma}

\begin{proof}
Construisons d'abord $Y'$ comme espace annelé. Comme espace
topologique, on prend $Y' = Y$. Notons $f' : X' \to Y'$ l'application d'espaces topologiques
égale à $f$. Comme faisceau structural $\mathcal{O}_{Y'}$, on prend
le membre de droite de l'égalité du lemme. Pour montrer que
$Y'$ est un schéma, il faut montrer que tout point possède un voisinage
affine. Puisque la formation du produit fibré de faisceaux
commute à la restriction aux ouverts, on peut supposer $Y$ affine.
Alors $X$ est affine (puisque $f$ est affine) et $X'$ l'est aussi
(voir le lemme \ref{more-morphisms-lemma-thickening-affine-scheme}).
Supposons que $Y \leftarrow X \rightarrow X'$ corresponde
à $B \rightarrow A \leftarrow A'$. Posons $B' = B \times_A A'$ ;
c'est l'anneau des sections globales de $\mathcal{O}_{Y'}$. Comme $A' \to A$ est surjectif
à noyau localement nilpotent, on voit que $B' \to B$ est surjectif
à noyau localement nilpotent. Ainsi, $\Spec(B') = \Spec(B)$ (comme
espaces topologiques). On affirme que $Y' = \Spec(B')$. Pour cela,
on montrera que, pour $g' \in B'$ d'image $g \in B$,
on a $\mathcal{O}_{Y'}(D(g)) = B'_{g'}$. En effet, d'après
Compléments d'algèbre, lemme \ref{more-algebra-lemma-diagram-localize}, on a
$$
(B')_{g'} = B_g \times_{A_h} A'_{h'}
$$
où $h \in A$, $h' \in A'$ sont les images de $g'$. Puisque
$B_g$ (resp.\ $A_h$, resp.\ $A'_{h'}$) est égal à $\mathcal{O}_Y(D(g))$
(resp.\ $f_*\mathcal{O}_X(D(g))$, resp.\ $f'_*\mathcal{O}_{X'}(D(g))$),
l'assertion s'ensuit.

\medskip\noindent
Il reste à montrer que $Y'$ est la somme amalgamée.
La discussion ci-dessus montre que le schéma $Y'$
possède un recouvrement ouvert affine $Y' = \bigcup W'_i$
tel que les ouverts correspondants
$U'_i \subset X'$, $W_i \subset Y$ et
$U_i \subset X$ soient affines.
De plus, si $A'_i$, $B_i$, $A_i$ sont les anneaux correspondant à
$U'_i$, $W_i$, $U_i$, alors
$W'_i$ correspond à $B_i \times_{A_i} A'_i$.
On peut donc appliquer les lemmes \ref{more-morphisms-lemma-basic-example-pushout} et
\ref{more-morphisms-lemma-pushout-fpqc-local} pour conclure que notre construction est une somme amalgamée
dans la catégorie des schémas.
\end{proof}

\noindent
Dans le lemme suivant, on utilise le produit fibré de catégories
défini dans
Catégories, exemple \ref{categories-example-2-fibre-product-categories}.

\begin{lemma}
\label{more-morphisms-lemma-equivalence-categories-schemes-over-pushout}
Soit $X \to X'$ un épaississement de schémas et soit $X \to Y$ un
morphisme affine de schémas. Soit $Y' = Y \amalg_X X'$ la somme amalgamée
(voir le lemme \ref{more-morphisms-lemma-pushout-along-thickening}). Le changement de base donne
un foncteur
$$
F : (\Sch/Y') \longrightarrow (\Sch/Y) \times_{(\Sch/X)} (\Sch/X')
$$
défini par $V' \longmapsto (V' \times_{Y'} Y, V' \times_{Y'} X', 1)$,
qui admet un adjoint à gauche
$$
G : (\Sch/Y) \times_{(\Sch/X)} (\Sch/X') \longrightarrow (\Sch/Y')
$$
envoyant le triplet $(V, U', \varphi)$ sur la somme amalgamée
$V \amalg_{(V \times_Y X)} U'$. Enfin, $F \circ G$ est isomorphe au
foncteur identité.
\end{lemma}

\begin{proof}
Soit $(V, U', \varphi)$ un objet de la catégorie produit fibré.
Posons $U = U' \times_{X'} X$. Remarquons que $U \to U'$ est un épaississement.
Puisque $\varphi : V \times_Y X \to U' \times_{X'} X = U$ est un isomorphisme,
on dispose d'un morphisme $U \to V$ au-dessus de $X \to Y$ qui identifie $U$ au
produit fibré $X \times_Y V$. En particulier, $U \to V$ est affine, voir
Morphismes, lemme \ref{morphisms-lemma-base-change-affine}.
On peut donc appliquer le lemme \ref{more-morphisms-lemma-pushout-along-thickening}
pour obtenir une somme amalgamée $V' = V \amalg_U U'$. Notons $V' \to Y'$ le morphisme
fourni par la propriété de somme amalgamée de $V'$ et par
les morphismes donnés $V \to Y$ et $U' \to X'$, qui coïncident sur $U$
comme morphismes à valeurs dans $Y'$. En posant $G(V, U', \varphi) = V'$,
on obtient le foncteur $G$.

\medskip\noindent
Prouvons que $G$ est adjoint à gauche de $F$. Soit $Z$ un schéma
sur $Y'$. Il faut montrer que
$$
\Mor(V', Z) = \Mor((V, U', \varphi), F(Z))
$$
où les ensembles de morphismes sont pris dans leurs catégories respectives.
Soit $g' : V' \to Z$ un morphisme. Notons $\tilde g$ (resp.\ $\tilde f'$)
le composé de $g'$ avec le morphisme $V \to V'$ (resp.\ $U' \to V'$).
Effectuons le changement de base de $\tilde g$ (resp.\ $\tilde f'$) par $Y \to Y'$ (resp.\ $X' \to Y'$)
pour obtenir un morphisme $g : V \to Z \times_{Y'} Y$
(resp.\ $f' : U' \to Z \times_{Y'} X'$). Alors $(g, f')$ est un élément
du membre de droite de l'égalité ci-dessus (détails omis).
Réciproquement, supposons que $(g, f') : (V, U', \varphi) \to F(Z)$ soit un
élément du membre de droite.
On peut considérer le composé $\tilde g : V \to Z$
(resp.\ $\tilde f' : U' \to Z$) de $g$ (resp.\ $f$) avec
$Z \times_{Y'} X' \to Z$ (resp.\ $Z \times_{Y'} Y \to Z$).
Alors $\tilde g$ et $\tilde f'$ coïncident comme morphismes de $U$ dans $Z$.
Par la propriété universelle de la somme amalgamée, on obtient un morphisme
$g' : V' \to Z$, c'est-à-dire un élément du membre de gauche.
On omet de vérifier que ces constructions sont inverses l'une de l'autre.

\medskip\noindent
Pour prouver que $F \circ G$ est isomorphe à l'identité, il faut
montrer que le morphisme d'adjonction
$(V, U', \varphi) \to F(G(V, U', \varphi))$ est un isomorphisme.
On peut pour cela travailler localement sur des ouverts affines. Écrivons $X = \Spec(A)$, $X' = \Spec(A')$
et $Y = \Spec(B)$. Alors $A' \to A$ et $B \to A$ sont des homomorphismes d'anneaux comme dans
Compléments d'algèbre, lemme \ref{more-algebra-lemma-module-over-fibre-product},
et $Y' = \Spec(B')$, avec $B' = B \times_A A'$. Supposons ensuite que
$V = \Spec(D)$, $U' = \Spec(C')$ et que $\varphi$ soit donné par un
isomorphisme de $A$-algèbres $D \otimes_B A \to C' \otimes_{A'} A = C'/IC'$.
Posons $D' = D \times_{C'/IC'} C'$. Dans ce cas, il faut
prouver que $D' \otimes_{B'} B \cong D$ et $D' \otimes_{B'} A' \cong C'$.
C'est un cas particulier de Compléments d'algèbre, lemme
\ref{more-algebra-lemma-module-over-fibre-product}.
\end{proof}

\begin{lemma}
\label{more-morphisms-lemma-scheme-over-pushout-flat-modules}
Soit $X \to X'$ un épaississement de schémas et soit $X \to Y$ un
morphisme affine de schémas. Soit $Y' = Y \amalg_X X'$ la somme amalgamée
(voir le lemme \ref{more-morphisms-lemma-pushout-along-thickening}). Soit $V' \to Y'$
un morphisme de schémas. Posons
$V = Y \times_{Y'} V'$, $U' = X' \times_{Y'} V'$ et $U = X \times_{Y'} V'$.
Il existe une équivalence de catégories entre
\begin{enumerate}
\item les $\mathcal{O}_{V'}$-modules quasi-cohérents plats sur $Y'$, et
\item la catégorie des triplets $(\mathcal{G}, \mathcal{F}', \varphi)$, où
\begin{enumerate}
\item $\mathcal{G}$ est un $\mathcal{O}_V$-module quasi-cohérent plat sur $Y$,
\item $\mathcal{F}'$ est un $\mathcal{O}_{U'}$-module quasi-cohérent plat
sur $X'$, et
\item $\varphi : (U \to V)^*\mathcal{G} \to (U \to U')^*\mathcal{F}'$
est un isomorphisme de $\mathcal{O}_U$-modules.
\end{enumerate}
\end{enumerate}
L'équivalence envoie $\mathcal{G}'$ sur
$((V \to V')^*\mathcal{G}', (U' \to V')^*\mathcal{G}', can)$.
Supposons que $\mathcal{G}'$ corresponde au triplet
$(\mathcal{G}, \mathcal{F}', \varphi)$. Alors
\begin{enumerate}
\item[(a)] $\mathcal{G}'$ est un $\mathcal{O}_{V'}$-module de type fini si et
seulement si $\mathcal{G}$ et $\mathcal{F}'$ sont des
$\mathcal{O}_Y$- et $\mathcal{O}_{U'}$-modules de type fini.
\item[(b)] si $V' \to Y'$ est localement de présentation finie, alors
$\mathcal{G}'$ est un $\mathcal{O}_{V'}$-module de
présentation finie si et seulement si $\mathcal{G}$ et $\mathcal{F}'$ sont des
$\mathcal{O}_Y$- et $\mathcal{O}_{U'}$-modules de présentation finie.
\end{enumerate}
\end{lemma}

\begin{proof}
Un foncteur quasi-inverse associe au triplet
$(\mathcal{G}, \mathcal{F}', \varphi)$ le produit fibré
$$
(V \to V')_*\mathcal{G}
\times_{(U \to V')_*\mathcal{F}}
(U' \to V')_*\mathcal{F}'
$$
où $\mathcal{F} = (U \to U')^*\mathcal{F}'$. Cette construction convient, car,
sur les ouverts affines, on retrouve l'équivalence de Compléments d'algèbre, lemme
\ref{more-algebra-lemma-relative-flat-module-over-fibre-product}.
Certains détails sont omis.

\medskip\noindent
Les points (a) et (b) résultent de
Compléments d'algèbre, lemmes
\ref{more-algebra-lemma-relative-finite-module-over-fibre-product} et
\ref{more-algebra-lemma-relative-finitely-presented-module-over-fibre-product}.
\end{proof}

\begin{lemma}
\label{more-morphisms-lemma-equivalence-categories-schemes-over-pushout-flat}
Dans la situation du
lemme \ref{more-morphisms-lemma-equivalence-categories-schemes-over-pushout},
si $V' = G(V, U', \varphi)$ pour un certain triplet $(V, U', \varphi)$, alors
\begin{enumerate}
\item $V' \to Y'$ est localement de type fini si et seulement si $V \to Y$ et
$U' \to X'$ sont localement de type fini,
\item $V' \to Y'$ est plat si et seulement si $V \to Y$ et $U' \to X'$ sont plats,
\item $V' \to Y'$ est plat et localement de présentation finie si et seulement si
$V \to Y$ et $U' \to X'$ sont plats et localement de présentation finie,
\item $V' \to Y'$ est lisse si et seulement si $V \to Y$ et $U' \to X'$ sont lisses,
\item $V' \to Y'$ est étale si et seulement si $V \to Y$ et $U' \to X'$
sont étales, et
\item ajouter ici d'autres propriétés selon les besoins.
\end{enumerate}
Si $W'$ est plat sur $Y'$, alors le morphisme d'adjonction
$G(F(W')) \to W'$ est un isomorphisme. Ainsi, $F$ et $G$ définissent des foncteurs
quasi-inverses l'un de l'autre entre la catégorie des schémas plats sur $Y'$
et la catégorie des triplets $(V, U', \varphi)$ tels que $V \to Y$
et $U' \to X'$ soient plats.
\end{lemma}

\begin{proof}
En se plaçant sur des ouverts affines, les assertions de ce lemme
sont équivalentes aux assertions correspondantes de
Compléments d'algèbre, lemme
\ref{more-algebra-lemma-properties-algebras-over-fibre-product}.
\end{proof}


\section{Ouverture du lieu de platitude}
\label{more-morphisms-section-open-flat}

\noindent
La démonstration de ce résultat demande un certain travail et mérite
(peut-être) une section à elle seule. La voici.

\begin{theorem}
\label{more-morphisms-theorem-openness-flatness}
\begin{reference}
\cite[IV, théorème 11.3.1]{EGA}
\end{reference}
Soit $S$ un schéma.
Soit $f : X \to S$ un morphisme localement de présentation finie.
Soit $\mathcal{F}$ un $\mathcal{O}_X$-module quasi-cohérent
localement de présentation finie. Alors
$$
U = \{x \in X \mid \mathcal{F}\text{ est plat sur }S\text{ en }x\}
$$
est ouvert dans $X$.
\end{theorem}

\begin{proof}
L'ouverture se vérifie localement sur $X$ ; on peut donc supposer
que $f$ est un morphisme de schémas affines. Dans ce cas, le
théorème est exactement le
théorème \ref{algebra-theorem-openness-flatness} d'Algèbre.
\end{proof}

\begin{lemma}
\label{more-morphisms-lemma-flat-locus-base-change}
Soit $S$ un schéma.
Soit
$$
\xymatrix{
X' \ar[r]_{g'} \ar[d]_{f'} & X \ar[d]^f \\
S' \ar[r]^g & S
}
$$
un diagramme cartésien de schémas.
Soit $\mathcal{F}$ un $\mathcal{O}_X$-module quasi-cohérent.
Soit $x' \in X'$ d'images
$x = g'(x')$ et $s' = f'(x')$.
\begin{enumerate}
\item Si $\mathcal{F}$ est plat sur $S$ en $x$, alors
$(g')^*\mathcal{F}$ est plat sur $S'$ en $x'$.
\item Si $g$ est plat en $s'$ et si $(g')^*\mathcal{F}$ est plat sur $S'$ en
$x'$, alors $\mathcal{F}$ est plat sur $S$ en $x$.
\end{enumerate}
En particulier, si $g$ est plat, si $f$ est localement de présentation finie
et si $\mathcal{F}$ est localement de présentation finie,
alors la formation de l'ouvert du
théorème \ref{more-morphisms-theorem-openness-flatness}
commute au changement de base.
\end{lemma}

\begin{proof}
Considérons le diagramme commutatif d'anneaux locaux
$$
\xymatrix{
\mathcal{O}_{X', x'} & \mathcal{O}_{X, x} \ar[l] \\
\mathcal{O}_{S', s'} \ar[u] & \mathcal{O}_{S, s} \ar[l] \ar[u]
}
$$
Notons que $\mathcal{O}_{X', x'}$
est un localisé de
$\mathcal{O}_{X, x} \otimes_{\mathcal{O}_{S, s}} \mathcal{O}_{S', s'}$
et que $((g')^*\mathcal{F})_{x'}$ est égal à
$\mathcal{F}_x \otimes_{\mathcal{O}_{X, x}} \mathcal{O}_{X', x'}$.
Le lemme résulte donc du
lemme \ref{algebra-lemma-base-change-flat-up-down} d'Algèbre.
\end{proof}






\section{Critère de platitude par fibres}
\label{more-morphisms-section-criterion-flat-fibres}

\noindent
Considérons un diagramme commutatif de schémas (diagramme de gauche)
$$
\xymatrix{
X \ar[rr]_f \ar[dr] & & Y \ar[dl] \\
& S
}
\quad
\xymatrix{
X_s \ar[rr]_{f_s} \ar[rd] & & Y_s \ar[dl] \\
& \Spec(\kappa(s))
}
$$
et un $\mathcal{O}_X$-module quasi-cohérent $\mathcal{F}$.
Étant donné un point $x \in X$ au-dessus de $s \in S$, d'image $y = f(x)$,
on se pose la question : $\mathcal{F}$ est-il plat
sur $Y$ en $x$ ? Si $\mathcal{F}$ est plat sur $S$ en $x$, le
théorème affirme que cette question est étroitement liée à celle
de savoir si la restriction de $\mathcal{F}$ à la fibre
$$
\mathcal{F}_s = (X_s \to X)^*\mathcal{F}
$$
est plate sur $Y_s$ en $x$. On trouvera ci-dessous une version
« noethérienne » et une version « de présentation finie » ; une version
« nilpotente » a déjà été traitée, voir le
lemme \ref{more-morphisms-lemma-flatness-morphism-thickenings}.

\begin{theorem}
\label{more-morphisms-theorem-criterion-flatness-fibre-Noetherian}
Soit $S$ un schéma.
Soit $f : X \to Y$ un morphisme de schémas sur $S$.
Soit $\mathcal{F}$ un $\mathcal{O}_X$-module quasi-cohérent.
Soit $x \in X$. Posons $y = f(x)$ et notons $s \in S$ l'image de $x$ dans $S$.
Supposons $S$, $X$, $Y$ localement noethériens,
$\mathcal{F}$ cohérent et $\mathcal{F}_x \not = 0$.
Alors les conditions suivantes sont équivalentes :
\begin{enumerate}
\item $\mathcal{F}$ est plat sur $S$ en $x$ et
$\mathcal{F}_s$ est plat sur $Y_s$ en $x$, et
\item $Y$ est plat sur $S$ en $y$ et $\mathcal{F}$ est
plat sur $Y$ en $x$.
\end{enumerate}
\end{theorem}

\begin{proof}
Considérons les homomorphismes d'anneaux
$$
\mathcal{O}_{S, s} \longrightarrow
\mathcal{O}_{Y, y} \longrightarrow \mathcal{O}_{X, x}
$$
et le module $\mathcal{F}_x$. La fibre de $\mathcal{F}_s$ en $x$
est le module $\mathcal{F}_x/\mathfrak m_s \mathcal{F}_x$ et
l'anneau local de $Y_s$ en $y$ est
$\mathcal{O}_{Y, y}/\mathfrak m_s \mathcal{O}_{Y, y}$.
L'implication (1) $\Rightarrow$ (2) est donc le
lemme \ref{algebra-lemma-criterion-flatness-fibre-Noetherian} d'Algèbre.
Si (2) est vérifiée, le premier homomorphisme d'anneaux est fidèlement plat
et $\mathcal{F}_x$ est plat sur $\mathcal{O}_{Y, y}$ ; le
lemme \ref{algebra-lemma-composition-flat} d'Algèbre
montre donc que $\mathcal{F}_x$ est plat sur $\mathcal{O}_{S, s}$.
De plus, $\mathcal{F}_x/\mathfrak m_s \mathcal{F}_x$ s'obtient
à partir du module plat $\mathcal{F}_x$ par le changement de base
$\mathcal{O}_{Y, y} \to \mathcal{O}_{Y, y}/\mathfrak m_s \mathcal{O}_{Y, y}$ ;
il est donc plat d'après le
lemme \ref{algebra-lemma-flat-base-change} d'Algèbre.
\end{proof}

\noindent
Voici la version non noethérienne.

\begin{theorem}
\label{more-morphisms-theorem-criterion-flatness-fibre}
Soit $S$ un schéma.
Soit $f : X \to Y$ un morphisme de schémas sur $S$.
Soit $\mathcal{F}$ un $\mathcal{O}_X$-module quasi-cohérent.
Supposons que
\begin{enumerate}
\item $X$ soit localement de présentation finie sur $S$,
\item $\mathcal{F}$ soit un $\mathcal{O}_X$-module de présentation finie, et
\item $Y$ soit localement de type fini sur $S$.
\end{enumerate}
Soit $x \in X$. Posons $y = f(x)$ et soit $s \in S$ l'image de $x$ dans $S$.
Si $\mathcal{F}_x \not = 0$, les conditions suivantes sont équivalentes :
\begin{enumerate}
\item $\mathcal{F}$ est plat sur $S$ en $x$ et
$\mathcal{F}_s$ est plat sur $Y_s$ en $x$, et
\item $Y$ est plat sur $S$ en $y$ et $\mathcal{F}$ est
plat sur $Y$ en $x$.
\end{enumerate}
De plus, l'ensemble des points $x$ où (1) et (2) sont vérifiées est ouvert dans
$\text{Supp}(\mathcal{F})$.
\end{theorem}

\begin{proof}
Considérons les homomorphismes d'anneaux
$$
\mathcal{O}_{S, s} \longrightarrow
\mathcal{O}_{Y, y} \longrightarrow \mathcal{O}_{X, x}
$$
et le module $\mathcal{F}_x$. La fibre de $\mathcal{F}_s$ en $x$
est le module $\mathcal{F}_x/\mathfrak m_s \mathcal{F}_x$ et
l'anneau local de $Y_s$ en $y$ est
$\mathcal{O}_{Y, y}/\mathfrak m_s \mathcal{O}_{Y, y}$.
L'implication (1) $\Rightarrow$ (2) est donc le
lemme \ref{algebra-lemma-criterion-flatness-fibre-fp-over-ft} d'Algèbre.
Si (2) est vérifiée, le premier homomorphisme d'anneaux est fidèlement plat
et $\mathcal{F}_x$ est plat sur $\mathcal{O}_{Y, y}$ ; le
lemme \ref{algebra-lemma-composition-flat} d'Algèbre
montre donc que $\mathcal{F}_x$ est plat sur $\mathcal{O}_{S, s}$.
De plus, $\mathcal{F}_x/\mathfrak m_s \mathcal{F}_x$ s'obtient
à partir du module plat $\mathcal{F}_x$ par le changement de base
$\mathcal{O}_{Y, y} \to \mathcal{O}_{Y, y}/\mathfrak m_s \mathcal{O}_{Y, y}$ ;
il est donc plat d'après le
lemme \ref{algebra-lemma-flat-base-change} d'Algèbre.

\medskip\noindent
D'après
Morphismes, lemme \ref{morphisms-lemma-finite-presentation-permanence},
le morphisme $f$ est localement de présentation finie.
Considérons l'ensemble
\begin{equation}
\label{more-morphisms-equation-open}
U = \{x \in X \mid \mathcal{F} \text{ plat en }x
\text{ à la fois sur }Y\text{ et }S\}.
\end{equation}
Cet ensemble est ouvert dans $X$ d'après le
théorème \ref{more-morphisms-theorem-openness-flatness}.
Notons que si $x \in U$, alors $\mathcal{F}_s$ est plat en
$x$ sur $Y_s$, comme changement de base d'un module plat par le
morphisme $Y_s \to Y$, voir
Morphismes, lemme \ref{morphisms-lemma-base-change-module-flat}.
Ainsi, en tout point de $U \cap \text{Supp}(\mathcal{F})$,
la condition (1) est vérifiée. D'autre part, il est
clair que si $x \in \text{Supp}(\mathcal{F})$ vérifie
(1) et (2), alors $x \in U$. L'ouvert recherché est donc
$U \cap \text{Supp}(\mathcal{F})$.
\end{proof}

\noindent
On utilise souvent ces théorèmes sous les formes simplifiées suivantes.
Nous ne donnons que les énoncés globaux ; il existe bien entendu aussi des
versions ponctuelles.

\begin{lemma}
\label{more-morphisms-lemma-morphism-between-flat-Noetherian}
Soit $S$ un schéma.
Soit $f : X \to Y$ un morphisme de schémas sur $S$.
Supposons que
\begin{enumerate}
\item $S$, $X$, $Y$ soient localement noethériens,
\item $X$ soit plat sur $S$,
\item pour tout $s \in S$, le morphisme
$f_s : X_s \to Y_s$ soit plat.
\end{enumerate}
Alors $f$ est plat. Si $f$ est en outre surjectif, $Y$ est plat sur $S$.
\end{lemma}

\begin{proof}
C'est un cas particulier du
théorème \ref{more-morphisms-theorem-criterion-flatness-fibre-Noetherian}.
\end{proof}

\begin{lemma}
\label{more-morphisms-lemma-morphism-between-flat}
Soit $S$ un schéma.
Soit $f : X \to Y$ un morphisme de schémas sur $S$.
Supposons que
\begin{enumerate}
\item $X$ soit localement de présentation finie sur $S$,
\item $X$ soit plat sur $S$,
\item pour tout $s \in S$, le morphisme
$f_s : X_s \to Y_s$ soit plat, et
\item $Y$ soit localement de type fini sur $S$.
\end{enumerate}
Alors $f$ est plat. Si $f$ est en outre surjectif, $Y$ est plat sur $S$.
\end{lemma}

\begin{proof}
C'est un cas particulier du
théorème \ref{more-morphisms-theorem-criterion-flatness-fibre}.
\end{proof}

\begin{lemma}
\label{more-morphisms-lemma-base-change-criterion-flatness-fibre}
Soit $S$ un schéma. Soit $f : X \to Y$ un morphisme de schémas sur $S$.
Soit $\mathcal{F}$ un $\mathcal{O}_X$-module quasi-cohérent.
Supposons que
\begin{enumerate}
\item $X$ soit localement de présentation finie sur $S$,
\item $\mathcal{F}$ soit un $\mathcal{O}_X$-module de présentation finie,
\item $\mathcal{F}$ soit plat sur $S$, et
\item $Y$ soit localement de type fini sur $S$.
\end{enumerate}
Alors l'ensemble
$$
U = \{x \in X \mid \mathcal{F} \text{ plat en }x \text{ sur }Y\}.
$$
est ouvert dans $X$ et sa formation commute à tout changement de base :
si $S' \to S$ est un morphisme de schémas et si $U'$ est l'ensemble des points
de $X' = X \times_S S'$ où $\mathcal{F}' = \mathcal{F} \times_S S'$
est plat sur $Y' = Y \times_S S'$, alors $U' = U \times_S S'$.
\end{lemma}

\begin{proof}
D'après
Morphismes, lemme \ref{morphisms-lemma-finite-presentation-permanence},
le morphisme $f$ est localement de présentation finie.
Ainsi, $U$ est ouvert d'après le
théorème \ref{more-morphisms-theorem-openness-flatness}.
Puisque l'on a supposé $\mathcal{F}$ plat sur $S$, le
théorème \ref{more-morphisms-theorem-criterion-flatness-fibre}
donne
$$
U = \{x \in X \mid \mathcal{F}_s \text{ plat en }x \text{ sur }Y_s\}.
$$
où $s$ désigne toujours l'image de $x$ dans $S$. (Cette description
reste trivialement valable lorsque $\mathcal{F}_x = 0$.) De plus, les hypothèses
du lemme restent vérifiées pour le morphisme $f' : X' \to Y'$
et le faisceau $\mathcal{F}'$. Ainsi, $U'$ admet une description analogue.
Autrement dit, il suffit de prouver que, pour
$s' \in S'$ d'image $s \in S$, l'ensemble
$$
\{x' \in X'_{s'} \mid
\mathcal{F}'_{s'} \text{ plat en }x' \text{ sur }Y'_{s'}\}
$$
est l'image réciproque du lieu correspondant dans $X_s$.
C'est vrai d'après le
lemme \ref{more-morphisms-lemma-flat-locus-base-change},
car, dans le diagramme cartésien
$$
\xymatrix{
X'_{s'} \ar[d] \ar[r] & X_s \ar[d] \\
Y'_{s'} \ar[r] & Y_s
}
$$
les morphismes horizontaux sont plats, puisqu'ils s'obtiennent par changement de
base par le morphisme plat $\Spec(\kappa(s')) \to \Spec(\kappa(s))$.
\end{proof}

\begin{lemma}
\label{more-morphisms-lemma-base-change-flatness-fibres}
Soit $S$ un schéma. Soit $f : X \to Y$ un morphisme de schémas sur $S$.
Supposons que
\begin{enumerate}
\item $X$ soit localement de présentation finie sur $S$,
\item $X$ soit plat sur $S$, et
\item $Y$ soit localement de type fini sur $S$.
\end{enumerate}
Alors l'ensemble
$$
U = \{x \in X \mid X\text{ plat en }x \text{ sur }Y\}.
$$
est ouvert dans $X$ et sa formation commute à tout changement de base.
\end{lemma}

\begin{proof}
C'est un cas particulier du
lemme \ref{more-morphisms-lemma-base-change-criterion-flatness-fibre}.
\end{proof}

\noindent
Le lemme suivant est une variante du
lemme \ref{algebra-lemma-free-fibre-flat-free} d'Algèbre.
Notons que l'hypothèse
que $(\mathcal{F}_s)_x$ est un $\mathcal{O}_{X_s, x}$-module plat
signifie que $(\mathcal{F}_s)_x$ est un $\mathcal{O}_{X_s, x}$-module libre,
ce qui est toujours le cas si $x \in X_s$ est un point générique d'une
composante irréductible de $X_s$ et si $X_s$ est réduit (en effet,
dans ce cas $\mathcal{O}_{X_s, x}$ est un corps, voir
Algèbre, lemme \ref{algebra-lemma-minimal-prime-reduced-ring}).

\begin{lemma}
\label{more-morphisms-lemma-flat-and-free-at-point-fibre}
Soit $f : X \to S$ un morphisme de schémas de présentation finie.
Soit $\mathcal{F}$ un $\mathcal{O}_X$-module de présentation finie.
Soit $x \in X$ d'image $s \in S$.
Si $\mathcal{F}$ est plat en $x$ sur $S$ et si $(\mathcal{F}_s)_x$ est un
$\mathcal{O}_{X_s, x}$-module plat, alors $\mathcal{F}$
est libre de rang fini dans un voisinage de $x$.
\end{lemma}

\begin{proof}
Si $\mathcal{F}_x \otimes \kappa(x)$ est nul, alors $\mathcal{F}_x = 0$
d'après le lemme de Nakayama (Algèbre, lemme \ref{algebra-lemma-NAK}), donc
$\mathcal{F}$ est nul dans un voisinage de $x$
(Modules, lemme \ref{modules-lemma-finite-type-stalk-zero})
et le lemme est établi. On peut donc supposer $\mathcal{F}_x \otimes \kappa(x)$
non nul, et l'on voit que le
théorème \ref{more-morphisms-theorem-criterion-flatness-fibre}
s'applique avec $f = \text{id} : X \to X$. On en conclut que $\mathcal{F}_x$
est plat sur $\mathcal{O}_{X, x}$. Ainsi, $\mathcal{F}_x$ est libre, voir
par exemple Algèbre, lemme \ref{algebra-lemma-finite-flat-local}.
Choisissons un voisinage ouvert $x \in U \subset X$ et des sections
$s_1, \ldots, s_r \in \mathcal{F}(U)$ dont les images forment une base de
$\mathcal{F}_x$. Le morphisme correspondant
$\psi : \mathcal{O}_U^{\oplus r} \to \mathcal{F}|_U$ est surjectif après
restriction de $U$ (Modules, lemme \ref{modules-lemma-finite-type-stalk-zero}).
Alors $\Ker(\psi)$ est de type fini (voir Modules, lemme
\ref{modules-lemma-kernel-surjection-finite-free-onto-finite-presentation})
et $\Ker(\psi)_x = 0$. En restreignant encore $U$,
$\psi$ devient donc un isomorphisme.
\end{proof}

\begin{lemma}
\label{more-morphisms-lemma-finite-free-open}
Soit $f : X \to S$ un morphisme de schémas
localement de présentation finie.
Soit $\mathcal{F}$ un $\mathcal{O}_X$-module de présentation finie
plat sur $S$. Alors l'ensemble
$$
\{x \in X : \mathcal{F}\text{ libre dans un voisinage de }x\}
$$
est ouvert dans $X$ et sa formation commute à tout changement de base
$S' \to S$.
\end{lemma}

\begin{proof}
L'ouverture est immédiate. Soit $x \in X$ d'image $s \in S$.
D'après le lemme \ref{more-morphisms-lemma-flat-and-free-at-point-fibre},
$x$ appartient à notre ensemble si et seulement si
$\mathcal{F}|_{X_s}$ est plat en $x$ sur $X_s$.
Ceci équivaut aussi à la platitude de $\mathcal{F}$
en $x$ sur $X$ (car cette propriété est
impliquée par la liberté de $\mathcal{F}_x$ et implique
la platitude de $\mathcal{F}|_{X_s}$ en $x$ sur $X_s$).
L'assertion sur le changement de base résulte donc du
lemme \ref{more-morphisms-lemma-base-change-criterion-flatness-fibre}
appliqué à $\text{id} : X \to X$ sur $S$.
\end{proof}
\section{Immersions fermées entre schémas lisses}
\label{more-morphisms-section-closed-immersions-smooth}

\noindent
Quelques résultats qui ne trouvent guère leur place ailleurs.

\begin{lemma}
\label{more-morphisms-lemma-etale-local-structure}
Soit $S$ un schéma. Soit $Y \to X$ une immersion fermée de schémas
lisses sur $S$. Pour tout $y \in Y$, il existe des entiers
$0 \leq m, n$ et un diagramme commutatif
$$
\xymatrix{
Y \ar[d] &
V \ar[l] \ar[d] \ar[r] &
\mathbf{A}^m_S
\ar[d]^{(a_1, \ldots, a_m) \mapsto (a_1, \ldots, a_m, 0 \ldots, 0)} \\
X &
U \ar[l] \ar[r]^-\pi &
\mathbf{A}^{m + n}_S
}
$$
où $U \subset X$ est ouvert, $V = Y \cap U$,
$\pi$ est étale, $V = \pi^{-1}(\mathbf{A}^m_S)$ et $y \in V$.
\end{lemma}

\begin{proof}
La question est locale sur $X$ ; on peut donc remplacer $X$ par
un voisinage ouvert de $y$. Puisque $Y \to X$ est une immersion régulière
d'après Diviseurs, lemme
\ref{divisors-lemma-immersion-smooth-into-smooth-regular-immersion},
on peut supposer $X = \Spec(A)$ affine et qu'il existe une suite régulière
$f_1, \ldots, f_n \in A$ telle que $Y = V(f_1, \ldots, f_n)$.
En restreignant encore $X$ (et donc $Y$),
on peut supposer qu'il existe un morphisme étale $Y \to \mathbf{A}^m_S$, voir
Morphismes, lemme \ref{morphisms-lemma-smooth-etale-over-affine-space}.
Soient $\overline{g}_1, \ldots, \overline{g}_m$ dans $\mathcal{O}_Y(Y)$
les fonctions coordonnées de ce morphisme étale.
Choisissons des relèvements $g_1, \ldots, g_m \in A$ de ces fonctions
et considérons le morphisme
$$
(g_1, \ldots, g_m, f_1, \ldots, f_n) :
X
\longrightarrow
\mathbf{A}^{m + n}_S
$$
sur $S$. C'est un morphisme entre schémas localement de présentation finie
sur $S$ ; il est donc localement de présentation finie
(Morphismes, lemme \ref{morphisms-lemma-finite-presentation-permanence}).
La restriction de ce morphisme à
$\mathbf{A}^m_S \subset \mathbf{A}^{m + n}_S$
est étale par construction. Par conséquent, pour montrer que
$X \to \mathbf{A}^{m + n}_S$ est étale en $y$,
il suffit de montrer que $X \to \mathbf{A}^{m + n}_S$ est plat en $y$,
voir Morphismes, lemme \ref{morphisms-lemma-etale-at-point}.
Soit $s \in S$ l'image de $y$. Il suffit de vérifier que
$X_s \to \mathbf{A}^{m + n}_s$ est plat en $y$, voir le
théorème \ref{more-morphisms-theorem-criterion-flatness-fibre}.
Soit $z \in \mathbf{A}^{m + n}_s$ l'image de $y$.
L'homomorphisme local d'anneaux
$$
\mathcal{O}_{\mathbf{A}^{m + n}_s, z}
\longrightarrow
\mathcal{O}_{X_s, y}
$$
est plat d'après Algèbre, lemme \ref{algebra-lemma-CM-over-regular-flat}.
En effet, les schémas lisses sur un corps sont réguliers et les anneaux réguliers
sont de Cohen-Macaulay, voir Variétés, lemme \ref{varieties-lemma-smooth-regular}
et Algèbre, lemme \ref{algebra-lemma-regular-ring-CM}.
La source et le but sont donc des anneaux locaux réguliers (et donc de Cohen-Macaulay).
Ils ont la même dimension : en effet, on a
$\dim(\mathcal{O}_{Y_s, y}) = \dim(\mathcal{O}_{\mathbf{A}^m_s, z})$
d'après Compléments d'algèbre, lemme \ref{more-algebra-lemma-dimension-etale-extension},
on a $\dim(\mathcal{O}_{\mathbf{A}^{m + n}_s, z}) = n +
\dim(\mathcal{O}_{\mathbf{A}^m_s, z})$, et l'on a
$\dim(\mathcal{O}_{X_s, y}) = n + \dim(\mathcal{O}_{Y_s, y})$,
car $\mathcal{O}_{Y_s, y}$ est le quotient de
$\mathcal{O}_{X_s, y}$ par la suite régulière $f_1, \ldots, f_n$
de longueur $n$ (voir
Diviseurs, remarque \ref{divisors-remark-relative-regular-immersion-elements}).
Enfin, l'anneau de la fibre de la flèche ci-dessus est fini sur $\kappa(z)$,
puisque $Y_s \to \mathbf{A}^m_s$ est étale en $y$.
Ceci achève la démonstration.
\end{proof}

\begin{remark}
\label{more-morphisms-remark-relative-blowup}
Fixons un anneau $R$ et posons $S = \Spec(R)$. Fixons des entiers $0 \leq m$ et
$1 \leq n$. Considérons l'immersion fermée
$$
Z = \mathbf{A}^m_S \longrightarrow \mathbf{A}^{m + n}_S = X,\quad
(a_1, \ldots, a_m) \mapsto (a_1, \ldots, a_m, 0, \ldots 0).
$$
Nous allons considérer l'éclatement $X'$ de $X$
le long du sous-schéma fermé $Z$. Écrivons
$$
X = \Spec(A)
\quad\text{avec}\quad
A = R[x_1, \ldots, x_m, y_1, \ldots, y_n]
$$
Alors $X'$ est le Proj de l'algèbre de Rees de $A$ relativement à
l'idéal $(y_1, \ldots, y_n)$. Cette algèbre de Rees est égale à
$B = A[T_1, \ldots, T_n]/(y_iT_j - y_jT_i)$ ; les détails sont omis.
Ainsi, $X' = \text{Proj}(B)$ est lisse sur $S$, car il est
recouvert par les ouverts affines
\begin{align*}
D_+(T_i)
& =
\Spec(B_{(T_i)}) \\
& =
\Spec(A[t_1, \ldots, \hat t_i, \ldots t_n]/(y_j - y_i t_j)) \\
& =
\Spec(R[x_1, \ldots, x_m, y_i, t_1, \ldots, \hat t_i, \ldots, t_n])
\end{align*}
qui sont isomorphes à $\mathbf{A}^{n + m}_S$.
Dans cette carte, le diviseur exceptionnel est défini par
$y_i = 0$ ; il est donc lui aussi lisse
sur $S$.
\end{remark}

\begin{lemma}
\label{more-morphisms-lemma-blowup}
Soit $S$ un schéma. Soit $Z \to X$ une immersion fermée de schémas
lisses sur $S$. Soit $b : X' \to X$ l'éclatement de $Z$,
de diviseur exceptionnel $E \subset X'$. Alors $X'$ et $E$ sont lisses
sur $S$. Le morphisme $p : E \to Z$ est canoniquement isomorphe
au fibré projectif
$$
\mathbf{P}(\mathcal{I}/\mathcal{I}^2) \longrightarrow Z
$$
où $\mathcal{I} \subset \mathcal{O}_X$ est le faisceau d'idéaux
de $Z$. Le faisceau relatif $\mathcal{O}_E(1)$ provenant de la structure
de fibré projectif est isomorphe à la restriction de
$\mathcal{O}_{X'}(-E)$ à $E$.
\end{lemma}

\begin{proof}
D'après Diviseurs, lemme
\ref{divisors-lemma-immersion-smooth-into-smooth-regular-immersion},
l'immersion $Z \to X$ est régulière ; le faisceau d'idéaux
$\mathcal{I}$ est donc de type fini, et $b$ est un morphisme projectif
admettant pour faisceau inversible relativement ample
$\mathcal{O}_{X'}(1) = \mathcal{O}_{X'}(-E)$, voir
Diviseurs, lemmes
\ref{divisors-lemma-blowing-up-gives-effective-Cartier-divisor} et
\ref{divisors-lemma-blowing-up-projective}.
Le morphisme canonique $\mathcal{I} \to b_*\mathcal{O}_{X'}(1)$
donne une immersion fermée
$$
X' \longrightarrow
\mathbf{P}\left(\bigoplus\nolimits_{n \geq 0}
\text{Sym}^n_{\mathcal{O}_X}(\mathcal{I})\right)
$$
par la construction même de l'éclatement. La restriction de ce morphisme
à $E$ donne un morphisme canonique
$$
E \longrightarrow
\mathbf{P}\left(\bigoplus\nolimits_{n \geq 0}
\text{Sym}^n_{\mathcal{O}_Z}(\mathcal{I}/\mathcal{I}^2)\right)
$$
sur $Z$. Puisque $\mathcal{I}/\mathcal{I}^2$ est localement libre de rang fini,
si ce morphisme canonique est un isomorphisme, la dernière assertion du
lemme est vérifiée. Après ces observations, la question est locale
sur $X$ pour la topologie étale. En effet, l'éclatement commute au changement
de base plat d'après Diviseurs, lemme \ref{divisors-lemma-flat-base-change-blowing-up},
et la lissité peut se vérifier après précomposition par un morphisme
étale surjectif. Ainsi, la structure locale étale d'une
immersion fermée de schémas sur $S$ donnée au lemme
\ref{more-morphisms-lemma-etale-local-structure} nous ramène au
cas traité dans la remarque \ref{more-morphisms-remark-relative-blowup}.
\end{proof}
\section{Modules plats et assassins relatifs}
\label{more-morphisms-section-flat-relative-assassin}

\noindent
Dans cette section, nous montrerons que, dans les situations géométriques,
le support d'un module plat est, en un certain sens, équidimensionnel sur la base.
Pour le cas noethérien, nous renvoyons à
\cite[IV Proposition 12.1.1.5]{EGA}.
Commençons par deux lemmes auxiliaires.

\begin{lemma}
\label{more-morphisms-lemma-mod-injective-valuation-ring}
Soit $A$ un anneau de valuation. Soit $A \to B$ un homomorphisme local
d'anneaux locaux essentiellement de type fini.
Soit $u : N \to M$ un homomorphisme de $B$-modules de type fini.
Supposons $M$ plat sur $A$ et
$\overline{u} : N/\mathfrak m_A N \to M/\mathfrak m_A M$ injectif.
Alors $u$ est injectif et $M/u(N)$ est plat sur $A$.
\end{lemma}

\begin{proof}
Nous déduirons ce lemme du
lemme \ref{algebra-lemma-mod-injective-general} d'Algèbre
(notons que les rôles de $M$ et $N$ ont été intervertis).
Pour cela, nous utiliserons Compléments d'algèbre, lemme
\ref{more-algebra-lemma-valuation-ring-flat-essentially-finite-type},
afin de nous ramener au cas de présentation finie.

\medskip\noindent
Par hypothèse, on peut écrire $B$ comme quotient du localisé d'une
algèbre de polynômes $P = A[x_1, \ldots, x_n]$ en un idéal premier
$\mathfrak q$. On peut alors considérer $u : N \to M$ comme un homomorphisme
de $P_\mathfrak q$-modules de type fini. On peut donc supposer, et l'on suppose,
que $B$ est essentiellement de présentation finie sur $A$.

\medskip\noindent
Le $B$-module $N$ est de type fini, mais pas nécessairement de présentation finie.
Écrivons $N = \colim N_\lambda$ comme limite inductive filtrante de
$B$-modules de présentation finie, avec des morphismes de transition surjectifs.
Par exemple, choisissons une présentation
$0 \to K \to B^{\oplus r} \to N \to 0$,
écrivons $K$ comme réunion de ses sous-modules de type fini $K_\lambda$ et
posons $N_\lambda = \Coker(K_\lambda \to B^{\oplus r})$.
Le module $N/\mathfrak m_A N$ est de présentation finie
sur l'anneau noethérien $B/\mathfrak m_A B$.
Par conséquent, pour $\lambda$ assez grand, on a
$N_\lambda/\mathfrak m_A N_\lambda = N/\mathfrak m_A N$.
Si l'on peut démontrer le lemme pour le composé
$u_\lambda : N_\lambda \to M$, on en déduit que
$N_\lambda = N$ et que le résultat vaut pour $u$.
On peut donc supposer, et l'on suppose, que $N$ est de présentation finie sur $B$.

\medskip\noindent
D'après Compléments d'algèbre, lemme
\ref{more-algebra-lemma-valuation-ring-flat-essentially-finite-type},
le module $M$ est de présentation finie sur $B$.
Toutes les hypothèses du
lemme \ref{algebra-lemma-mod-injective-general} d'Algèbre
sont donc vérifiées, ce qui permet de conclure.
\end{proof}

\begin{lemma}
\label{more-morphisms-lemma-make-base-change}
\begin{reference}
Ce résultat figure dans la démonstration de
\cite[IV Proposition 12.1.1.5]{EGA}
\end{reference}
Soit $f : X \to S$ un morphisme de schémas. Soit $y \in X$ un point
d'image $t \in S$. Notons $Y \subset X$ l'adhérence de $\{y\}$,
munie de sa structure de sous-schéma fermé intègre de $X$. Soit $s \in S$ et soit
$x \in Y_s$ un point générique d'une composante irréductible de $Y_s$.
Il existe un diagramme cartésien
$$
\xymatrix{
X' \ar[r]_{g'} \ar[d]_{f'} & X \ar[d]^f \\
S' \ar[r]^g & S
}
$$
ayant les propriétés suivantes :
\begin{enumerate}
\item $S'$ est le spectre d'un anneau de valuation
de point générique $t'$ et de point fermé $s'$,
\item $g(t') = t$ et $g(s') = s$,
\item il existe un point $y' \in X'_{t'}$ qui est
un point générique d'une composante irréductible de
$(S' \times_S Y)_{t'} = Y_t \times_t t'$
et vérifie $g'(y') = y$,
\item si l'on note $Y' \subset X'$ l'adhérence de $\{y'\}$,
munie de sa structure de sous-schéma fermé intègre de $X'$,
il existe un point $x' \in Y'_{s'}$ qui est un point générique
d'une composante irréductible de $Y'_{s'}$
et vérifie $g'(x') = x$.
\end{enumerate}
\end{lemma}

\begin{proof}
Choisissons un anneau de valuation $R$, posons $S' = \Spec(R)$, de
point générique $t'$ et de point fermé $s'$, et choisissons un morphisme
$h : S' \to X$ tel que $h(t') = y$ et $h(s') = x$.
Voir Schémas, lemme \ref{schemes-lemma-points-specialize}.
Posons $g = f \circ h$, de sorte que $g(t') = t$ et $g(s') = s$.
Considérons le changement de base
$$
\xymatrix{
X' \ar[r]_{g'} \ar[d] & X \ar[d] \\
S' \ar@/^1em/[u]^\sigma \ar[r]^-g & S
}
$$
On obtient une section $\sigma$ du morphisme changé de base telle que
$h = g' \circ \sigma$.

\medskip\noindent
Bien entendu, $\sigma$ se factorise par le changement de base $S' \times_S Y$ de $Y$,
puisque $h$ se factorise par $Y$. Soit $y' \in X'_{t'} \subset X'$ le
point générique d'une composante irréductible de la fibre
$$
(S' \times_S Y)_{t'} = Y_t \times_t t'
$$
contenant le point $\sigma(t')$, c'est-à-dire telle que
$y' \leadsto \sigma(t')$.
Puisque $g'(y') \in Y_t$ et $g(y') \leadsto g(\sigma(t')) = y$,
on obtient $g'(y') = y$, car $y$ est le point générique de
la fibre $Y_t$.
Notons $Y' \subset X'$ l'adhérence de $\{y'\}$ dans $X'$,
munie de sa structure de sous-schéma fermé intègre.
Alors $\sigma$ se factorise par $Y'$, puisque $\sigma(t') \in Y'$.
Choisissons un point générique $x' \in Y'_{s'}$ d'une composante
irréductible de $Y'_{s'}$ contenant $\sigma(s')$ ; on
a donc $x' \leadsto \sigma(s')$, puis $g'(x') \leadsto g'(\sigma(s')) = x$.
À nouveau, puisque $x$ est par hypothèse un point générique d'une
composante irréductible de $Y_s$
et puisque $g'(Y') \subset Y$, on en conclut que $g'(x') = x$.
\end{proof}

\begin{lemma}
\label{more-morphisms-lemma-associated-point-specializes}
Soit $f : X \to S$ un morphisme de schémas localement de type fini.
Soit $\mathcal{F}$ un $\mathcal{O}_X$-module quasi-cohérent de type fini.
Soit $y \in \text{Ass}_{X/S}(\mathcal{F})$ d'image $t \in S$.
Notons $Y \subset X$ l'adhérence de $\{y\}$ dans $X$,
munie de sa structure de sous-schéma fermé intègre. Soit $s \in S$
et soit $x \in Y_s$ un point générique d'une composante irréductible
de $Y_s$. Si $\mathcal{F}$ est plat sur $S$ en $x$, alors
$x \in \text{Ass}_{X/S}(\mathcal{F})$ et
$\dim_x(Y_s) = \dim(Y_t)$.
\end{lemma}

\begin{proof}
Choisissons un diagramme comme au lemme \ref{more-morphisms-lemma-make-base-change}.
Posons $\mathcal{F}' = (g')^*\mathcal{F}$.
Diviseurs, lemme \ref{divisors-lemma-base-change-relative-assassin},
implique que $y' \in \text{Ass}_{X'/S'}(\mathcal{F}')$.
Le choix de $y'$ donne aussi $\dim(Y'_{t'}) = \dim(Y_t)$,
voir par exemple Algèbre, lemme
\ref{algebra-lemma-inequalities-under-field-extension}.
D'après Algèbre, lemme
\ref{algebra-lemma-finite-type-domain-over-valuation-ring-dim-fibres},
$Y'_{s'}$ est équidimensionnel de dimension égale à $\dim(Y_t)$.
Puisque $\mathcal{F}$ est plat en $x$ sur $S$,
$\mathcal{F}'$ est plat en $x'$ sur $S'$, voir
Morphismes, lemme \ref{morphisms-lemma-base-change-module-flat}.

\medskip\noindent
Supposons que l'on puisse montrer $x' \in \text{Ass}_{X'/S}(\mathcal{F}')$. Alors
Diviseurs, lemme \ref{divisors-lemma-base-change-relative-assassin},
implique que $x \in \text{Ass}_{X/S}(\mathcal{F})$ et que
la composante irréductible $C'$ de $Y'_{s'}$ contenant $x'$
est une composante irréductible de $C \times_s s'$, où $C \subset Y_s$
est la composante irréductible contenant $x$.
D'où $\dim(C) = \dim(C') = \dim(Y_t)$ (voir ci-dessus), ce qui achève
la démonstration.
On est ainsi ramené au cas traité dans le paragraphe suivant.

\medskip\noindent
Supposons $S = \Spec(A)$, où $A$ est un anneau de valuation,
et où $t$ et $s$ sont les points générique et fermé de $S$.
Nous allons supposer $x \not \in \text{Ass}_{X/S}(\mathcal{F})$
afin d'aboutir à une contradiction. Autrement dit, on suppose
$x \not \in \text{Ass}_{X_s}(\mathcal{F}_s)$, où
$\mathcal{F}_s$ est l'image réciproque de $\mathcal{F}$ sur $X_s$.
Considérons l'homomorphisme d'anneaux
$$
A \longrightarrow \mathcal{O}_{X, x} = B
$$
et le module $N = \mathcal{F}_x$ sur $B = \mathcal{O}_{X, x}$.
Alors $B/\mathfrak m_A B = \mathcal{O}_{X_s, x}$ et
$N/\mathfrak m_A N$ est la fibre de $\mathcal{F}_s$ au point $x$.
Notons $\mathfrak q \subset B$ l'idéal premier correspondant au
point $y$, voir Schémas, lemme \ref{schemes-lemma-specialize-points}.
Puisque $x$ est un point générique de $Y_s$, le radical de
$\mathfrak q + \mathfrak m_A B$ est $\mathfrak m_B$.
Alors $\text{Ass}_{B/\mathfrak m_A B}(N/\mathfrak m_A N)$
est un ensemble fini d'idéaux premiers (Algèbre, lemme \ref{algebra-lemma-finite-ass})
qui ne contient pas l'idéal maximal de $B/\mathfrak m_A B$,
puisque $x \not \in \text{Ass}_{X/S}(\mathcal{F})$.
L'image de $\mathfrak q$ dans $B/\mathfrak m_A B$ n'est donc contenue
dans aucun de ces idéaux premiers. Par conséquent,
le lemme d'évitement des idéaux premiers (Algèbre, lemme \ref{algebra-lemma-silly})
permet de trouver un élément $g \in \mathfrak q$
dont l'image dans $B/\mathfrak m_A B$ est non diviseur de zéro sur
$N/\mathfrak m_A N$ (on utilise ici la description des diviseurs de zéro
d'Algèbre, lemme \ref{algebra-lemma-ass-zero-divisors}). Puisque
$N = \mathcal{F}_x$ est $A$-plat, le
lemme \ref{more-morphisms-lemma-mod-injective-valuation-ring} montre que
$$
g : N \longrightarrow N
$$
est injectif. En particulier, si $K = \text{Frac}(A)$ est le
corps des fractions de $A$, alors
$$
g : N \otimes_A K \longrightarrow N \otimes_A K
$$
est injectif. Observons que $\mathfrak q$ correspond à un
idéal premier de $B \otimes_A K$. Notons $\mathcal{F}_t$ la
restriction de $\mathcal{F}$ à la fibre générique $X_t$. On a
$(B \otimes_A K)_{\mathfrak q} = \mathcal{O}_{X_t, y}$
et $(N \otimes_A K)_\mathfrak q$
est la fibre en $y$ de $\mathcal{F}_t$. On obtient donc
que $g \in \mathfrak m_y \subset \mathcal{O}_{X_t, y}$
est non diviseur de zéro sur la fibre $(\mathcal{F}_t)_y$,
ce qui contredit l'hypothèse $y \in \text{Ass}_{X/S}(\mathcal{F})$.
\end{proof}

\begin{lemma}
\label{more-morphisms-lemma-relative-dimension-support-flat}
Soit $f : X \to S$ un morphisme de schémas localement de type fini.
Soit $\mathcal{F}$ un $\mathcal{O}_X$-module quasi-cohérent de type fini,
plat sur $S$. Supposons $S$ irréductible, de point générique $\eta$.
Si $\dim(\text{Supp}(\mathcal{F}_\eta)) \leq r$, alors, pour tout $s \in S$,
on a $\dim(\text{Supp}(\mathcal{F}_s)) \leq r$.
\end{lemma}

\begin{proof}
Soit $x \in \text{Supp}(\mathcal{F}_s)$ un point générique d'une composante
irréductible de $\text{Supp}(\mathcal{F}_s)$. D'après
Algèbre, lemme \ref{algebra-lemma-going-down-flat-module},
on peut trouver une spécialisation $y \leadsto x$ dans $\text{Supp}(\mathcal{F})$
avec $f(y) = \eta$. On peut bien entendu supposer que $y$ est un point générique
d'une composante irréductible de $\text{Supp}(\mathcal{F}_\eta)$.
Le lemme \ref{more-morphisms-lemma-associated-point-specializes}
permet de conclure que la dimension de $\overline{\{x\}}$ est au plus $r$.
\end{proof}

\begin{lemma}
\label{more-morphisms-lemma-flat-associated-equidimensional}
Soit $f : X \to S$ un morphisme de schémas localement de type fini.
Soit $\mathcal{F}$ un $\mathcal{O}_X$-module quasi-cohérent de type fini.
Soit $y \in \text{Ass}_{X/S}(\mathcal{F})$.
Notons $Y \subset X$ l'adhérence de $\{y\}$ dans $X$,
munie de sa structure de sous-schéma fermé intègre. Notons $T \subset S$ l'adhérence
de $\{f(y)\}$, munie de sa structure de sous-schéma fermé intègre. On obtient
un diagramme commutatif
$$
\xymatrix{
Y \ar[r] \ar[d] & X \ar[d] \\
T \ar[r] & S
}
$$
où $Y \to T$ est dominant. Supposons $\mathcal{F}$ plat sur $S$
en tous les points génériques des composantes irréductibles des fibres de $Y \to T$
(par exemple si $\mathcal{F}$ est plat sur $S$). Alors
\begin{enumerate}
\item si $s \in S$ et si $x \in Y_s$ est le point générique d'une
composante irréductible de $Y_s$, alors $x \in \text{Ass}_{X/S}(\mathcal{F})$, et
\item il existe un entier $d \geq 0$ tel que
$Y \to T$ soit de dimension relative $d$, voir
Morphismes, définition \ref{morphisms-definition-relative-dimension-d}.
\end{enumerate}
\end{lemma}

\begin{proof}
Cela résulte immédiatement de la version ponctuelle,
le lemme \ref{more-morphisms-lemma-associated-point-specializes}.
Notons que, pour calculer la dimension des schémas
localement algébriques $Y_s$, il suffit de se placer au voisinage
des points génériques, voir Variétés, section
\ref{varieties-section-algebraic-schemes}.
\end{proof}

\begin{remark}
\label{more-morphisms-remark-flat-equidimensional}
Voici quelques cas où les résultats ci-dessus, en particulier le
lemme \ref{more-morphisms-lemma-flat-associated-equidimensional},
permettent de conclure qu'un morphisme de schémas $f : X \to S$
est de dimension relative $d$ au sens de
Morphismes, définition \ref{morphisms-definition-relative-dimension-d}.
C'est par exemple le cas si
\begin{enumerate}
\item $X$ est intègre, de point générique $\xi$,
\item le degré de transcendance de $\kappa(\xi)$ sur $\kappa(f(\xi))$ est $d$,
\item $f$ est localement de type fini, et
\item il existe un $\mathcal{O}_X$-module quasi-cohérent $\mathcal{F}$
de type fini, plat sur $S$ et tel que $\text{Supp}(\mathcal{F}) = X$.
\end{enumerate}
Voici un autre ensemble d'hypothèses qui convient :
\begin{enumerate}
\item $S$ est irréductible, de point générique $\eta$,
\item $X_\eta$ est dense dans $X$,
\item toute composante irréductible de $X_\eta$ est de dimension $d$,
\item $f$ est localement de type fini, et
\item il existe un $\mathcal{O}_X$-module quasi-cohérent $\mathcal{F}$
de type fini, plat sur $S$ et tel que $\text{Supp}(\mathcal{F}) = X$.
\end{enumerate}
On peut bien entendu affaiblir la condition de platitude sur $\mathcal{F}$
et demander seulement que $\mathcal{F}$ soit plat sur $S$ en codimension $0$,
c'est-à-dire que $\mathcal{F}$ soit plat sur $S$ en tout point générique de
toute fibre. Si nous avons un jour besoin de ces résultats, nous les énoncerons
et les démontrerons ici avec soin.
\end{remark}
\section{Retour sur la normalisation}
\label{more-morphisms-section-normalization}

\noindent
La normalisation commute au changement de base lisse.

\begin{lemma}
\label{more-morphisms-lemma-integral-closure-smooth-pullback}
Soit $f : Y \to X$ un morphisme lisse de schémas. Soit $\mathcal{A}$ un
faisceau quasi-cohérent de $\mathcal{O}_X$-algèbres. La clôture intégrale
de $\mathcal{O}_Y$ dans $f^*\mathcal{A}$ est égale à $f^*\mathcal{A}'$,
où $\mathcal{A}' \subset \mathcal{A}$ est la clôture intégrale de
$\mathcal{O}_X$ dans $\mathcal{A}$.
\end{lemma}

\begin{proof}
C'est la traduction du
lemme \ref{algebra-lemma-integral-closure-commutes-smooth} d'Algèbre
dans le langage des schémas. Les détails sont omis.
\end{proof}

\begin{lemma}[La normalisation commute au changement de base lisse]
\label{more-morphisms-lemma-normalization-smooth-localization}
Soit
$$
\xymatrix{
Y_2 \ar[r] \ar[d]_{f_2} & Y_1 \ar[d]^{f_1} \\
X_2 \ar[r]^\varphi & X_1
}
$$
un carré cartésien dans la catégorie des schémas. Supposons $f_1$ quasi-compact
et quasi-séparé, et $\varphi$ lisse.
Soit $Y_i \to X_i' \to X_i$ la normalisation de $X_i$ dans $Y_i$.
Alors $X_2' \cong X_2 \times_{X_1} X_1'$.
\end{lemma}

\begin{proof}
Le changement de base de la factorisation $Y_1 \to X_1' \to X_1$ à $X_2$
est une factorisation $Y_2 \to X_2 \times_{X_1} X_1' \to X_2$, et
$X_2 \times_{X_1} X_1' \to X_2$ est entier
(Morphismes, lemme \ref{morphisms-lemma-base-change-finite}).
On obtient donc un morphisme
$h : X_2' \to X_2 \times_{X_1} X_1'$ par la propriété universelle de
Morphismes, lemme \ref{morphisms-lemma-characterize-normalization}.
Observons que $X_2'$ est le spectre relatif de la clôture intégrale
de $\mathcal{O}_{X_2}$ dans $f_{2, *}\mathcal{O}_{Y_2}$.
Si $\mathcal{A}' \subset f_{1, *}\mathcal{O}_{Y_1}$ désigne la clôture
intégrale de $\mathcal{O}_{X_1}$, alors $X_2 \times_{X_1} X_1'$ est le
spectre relatif de $\varphi^*\mathcal{A}'$, voir
Constructions, lemme \ref{constructions-lemma-spec-properties}.
D'après
Cohomologie des schémas, lemme \ref{coherent-lemma-flat-base-change-cohomology},
on sait que $f_{2, *}\mathcal{O}_{Y_2} = \varphi^*f_{1, *}\mathcal{O}_{Y_1}$.
Le résultat découle donc du
lemme \ref{more-morphisms-lemma-integral-closure-smooth-pullback}.
\end{proof}

\begin{lemma}[Normalisation et morphismes lisses]
\label{more-morphisms-lemma-normalization-and-smooth}
Soit $X \to Y$ un morphisme lisse de schémas. Supposons que tout ouvert
quasi-compact de $Y$ ait un nombre fini de composantes irréductibles. Il en est alors
de même de $X$, et il existe un unique isomorphisme $X^\nu = X \times_Y Y^\nu$
sur $X$, où $X^\nu$, $Y^\nu$ sont les normalisations de $X$, $Y$.
\end{lemma}

\begin{proof}
D'après Descente, lemme \ref{descent-lemma-locally-finite-nr-irred-local-fppf},
tout ouvert quasi-compact de $X$ a un nombre fini de composantes irréductibles.
Notons que $X_{red} = X \times_Y Y_{red}$, car un schéma lisse sur un schéma
réduit est réduit, voir
Descente, lemme \ref{descent-lemma-reduced-local-smooth}.
On peut donc supposer que $X$ et $Y$ sont réduits (car la normalisation
d'un schéma est, par définition, égale à celle de son réduit).
Notons ensuite que $X' = X \times_Y Y^\nu$ est un schéma normal
d'après Descente, lemme \ref{descent-lemma-normal-local-smooth}.
Le morphisme $X' \to Y^\nu$ est lisse (donc plat) ; les points génériques
des composantes irréductibles de $X'$ sont donc au-dessus des points génériques
des composantes irréductibles de $Y^\nu$. Puisque $Y^\nu \to Y$ est birationnel,
on en conclut que $X' \to X$ est aussi birationnel (car $X' \to Y^\nu$
induit un isomorphisme sur les fibres au-dessus des points génériques de $Y$).
On en déduit qu'il existe une factorisation
$X^\nu \to X' \to X$, voir
Morphismes, lemme \ref{morphisms-lemma-normalization-normal},
dont le premier morphisme est un isomorphisme, puisque $X'$ est normal et entier sur $X$.
\end{proof}

\begin{lemma}[Normalisation et hensélisation]
\label{more-morphisms-lemma-normalization-and-henselization}
Soit $X$ un schéma localement noethérien. Soit $\nu : X^\nu \to X$
le morphisme de normalisation. Alors, pour tout point $x \in X$,
le changement de base
$$
X^\nu \times_X \Spec(\mathcal{O}_{X, x}^h) \to \Spec(\mathcal{O}_{X, x}^h),
\quad\text{resp.}\quad
X^\nu \times_X \Spec(\mathcal{O}_{X, x}^{sh}) \to \Spec(\mathcal{O}_{X, x}^{sh})
$$
est la normalisation de $\Spec(\mathcal{O}_{X, x}^h)$,
resp.\ $\Spec(\mathcal{O}_{X, x}^{sh})$.
\end{lemma}

\begin{proof}
Soient $\eta_1, \ldots, \eta_r$ les points génériques des composantes
irréductibles de $X$ passant par $x$. Le changement de base de la normalisation à
$\Spec(\mathcal{O}_{X, x})$ est le spectre de la
clôture intégrale de $\mathcal{O}_{X, x}$ dans $\prod \kappa(\eta_i)$.
Cela résulte de notre construction de la normalisation de $X$ dans
Morphismes, définition \ref{morphisms-definition-normalization},
et de Morphismes, lemme \ref{morphisms-lemma-integral-closure} ;
on peut aussi utiliser la description de la normalisation
dans Morphismes, lemme \ref{morphisms-lemma-description-normalization}.
On est ainsi ramené au problème d'algèbre suivant.
Soit $A$ un anneau local noethérien ; rappelons que son
hensélisé $A^h$ et son hensélisé strict $A^{sh}$
sont alors eux aussi noethériens (Compléments d'algèbre, lemme
\ref{more-algebra-lemma-henselization-noetherian}).
Soient $\mathfrak p_1, \ldots, \mathfrak p_r$
ses idéaux premiers minimaux. Soit $A'$ la clôture intégrale
de $A$ dans $\prod \kappa(\mathfrak p_i)$.
Le problème est de montrer que
$A' \otimes_A A^h$, resp.\ $A' \otimes_A A^{sh}$, s'obtient
à partir de l'anneau local noethérien $A^h$, resp.\ $A^{sh}$,
de la même manière.

\medskip\noindent
Puisque $A^h$, resp.\ $A^{sh}$, sont des limites inductives de $A$-algèbres étales,
les idéaux premiers minimaux de $A$ et $A^{sh}$ sont exactement
les idéaux premiers de $A^h$, resp.\ $A^{sh}$, au-dessus des idéaux premiers minimaux
de $A$ (par descente des idéaux premiers, voir
Algèbre, lemmes \ref{algebra-lemma-flat-going-down} et
\ref{algebra-lemma-minimal-prime-image-minimal-prime}).
Ainsi, Compléments d'algèbre, lemme \ref{more-algebra-lemma-fibres-henselization},
montre que
$A^h \otimes_A \prod \kappa(\mathfrak p_i)$,
resp.
$A^{sh} \otimes_A \prod \kappa(\mathfrak p_i)$,
est le produit des corps résiduels aux idéaux premiers minimaux
de $A^h$, resp.\ $A^{sh}$. On sait que la formation de la
clôture intégrale dans un suranneau commute au changement de base
étale, voir
Algèbre, lemme \ref{algebra-lemma-integral-closure-commutes-etale}.
En écrivant $A^h$ et $A^{sh}$ comme limites inductives de $A$-algèbres étales,
on voit qu'il en est de même pour le changement de base à
$A^h$ et $A^{sh}$ (on peut aussi utiliser le résultat plus général
d'Algèbre, lemme \ref{algebra-lemma-integral-closure-commutes-colim-smooth}).
\end{proof}








\section{Morphismes normaux}
\label{more-morphisms-section-normal}

\noindent
L'article \cite{DM} de Deligne et Mumford mentionne la notion de morphisme
normal. Il s'agit d'un type parmi toute une série\footnote{
Les autres types sont coprof $\leq k$, Cohen-Macaulay, $(S_k)$,
régulier, $(R_k)$ et réduit. Voir \cite[IV, définition 6.8.1.]{EGA}.
Les morphismes de Gorenstein seront définis dans
Dualité pour les schémas, section \ref{duality-section-gorenstein}.}
de morphismes qui peuvent tous être définis de manière analogue. Nous les
introduirons au besoin, chacun dans sa propre section.

\begin{definition}
\label{more-morphisms-definition-normal}
Soit $f : X \to Y$ un morphisme de schémas.
Supposons que toutes les fibres $X_y$ soient des schémas localement noethériens.
\begin{enumerate}
\item Soient $x \in X$ et $y = f(x)$. On dit que $f$ est {\it normal en $x$}
si $f$ est plat en $x$ et si le schéma $X_y$ est géométriquement
normal en $x$ sur $\kappa(y)$ (voir
Variétés, définition \ref{varieties-definition-geometrically-normal}).
\item On dit que $f$ est un {\it morphisme normal} si $f$ est normal
en tout point de $X$.
\end{enumerate}
\end{definition}

\noindent
La normalité du morphisme $X \to Y$ est donc une condition
plus forte que celle demandant seulement que toutes les
fibres soient des schémas normaux localement noethériens.

\begin{lemma}
\label{more-morphisms-lemma-normal}
Soit $f : X \to Y$ un morphisme de schémas.
Supposons toutes les fibres de $f$ localement noethériennes.
Les conditions suivantes sont équivalentes :
\begin{enumerate}
\item $f$ est normal, et
\item $f$ est plat et ses fibres sont des schémas géométriquement normaux.
\end{enumerate}
\end{lemma}

\begin{proof}
Cela résulte directement des définitions.
\end{proof}

\begin{lemma}
\label{more-morphisms-lemma-smooth-normal}
Un morphisme lisse est normal.
\end{lemma}

\begin{proof}
Soit $f : X \to Y$ un morphisme lisse.
Puisque $f$ est localement de présentation finie, voir
Morphismes, lemme \ref{morphisms-lemma-smooth-locally-finite-presentation},
les fibres $X_y$ sont localement de type fini sur un corps, donc
localement noethériennes. De plus, $f$ est plat, voir
Morphismes, lemme \ref{morphisms-lemma-smooth-flat}.
Enfin, les fibres $X_y$ sont lisses sur un corps (d'après
Morphismes, lemme \ref{morphisms-lemma-base-change-smooth}),
donc géométriquement normales d'après
Variétés, lemme \ref{varieties-lemma-smooth-geometrically-normal}.
Ainsi, $f$ est normal d'après le
lemme \ref{more-morphisms-lemma-normal}.
\end{proof}

\noindent
Nous voulons montrer que cette notion est locale à la source et au but
pour la topologie lisse. Traitons d'abord la propriété d'avoir des
fibres localement noethériennes.

\begin{lemma}
\label{more-morphisms-lemma-locally-Noetherian-fibres-fppf-local-source-and-target}
La propriété $\mathcal{P}(f)=$« les fibres de $f$ sont localement noethériennes »
est locale pour la topologie fppf à la source et au but.
\end{lemma}

\begin{proof}
Soit $f : X \to Y$ un morphisme de schémas.
Soit $\{\varphi_i : Y_i \to Y\}_{i \in I}$ un recouvrement fppf de $Y$.
Notons $f_i : X_i \to Y_i$ le changement de base de $f$ par $\varphi_i$.
Soit $i \in I$ et soit $y_i \in Y_i$ un point.
Posons $y = \varphi_i(y_i)$. Notons que
$$
X_{i, y_i} = \Spec(\kappa(y_i)) \times_{\Spec(\kappa(y))} X_y.
$$
De plus, puisque $\varphi_i$ est de présentation finie, l'extension de corps
$\kappa(y_i)/\kappa(y)$ est de type fini.
Dans cette situation, $X_y$ est donc localement noethérien si et
seulement si $X_{i, y_i}$ est localement noethérien, voir
Variétés, lemme \ref{varieties-lemma-locally-Noetherian-base-change}.
Ce fait donne le caractère local au but.

\medskip\noindent
Soit $\{X_i \to X\}$ un recouvrement fppf de $X$.
Soit $y \in Y$. Dans ce cas, $\{X_{i, y} \to X_y\}$ est un
recouvrement fppf de la fibre. Le caractère local à la source
résulte donc de Descente, lemme \ref{descent-lemma-Noetherian-local-fppf}.
\end{proof}

\begin{lemma}
\label{more-morphisms-lemma-normal-fppf-local-source-and-target}
La propriété
$\mathcal{P}(f)=$« les fibres de $f$ sont localement noethériennes et $f$ est normal »
est locale pour la topologie fppf au but et
locale pour la topologie lisse à la source.
\end{lemma}

\begin{proof}
On a
$\mathcal{P}(f) =
\mathcal{P}_1(f) \wedge \mathcal{P}_2(f) \wedge \mathcal{P}_3(f)$,
où
$\mathcal{P}_1(f)=$« les fibres de $f$ sont localement noethériennes »,
$\mathcal{P}_2(f)=$« $f$ est plat », et
$\mathcal{P}_3(f)=$« les fibres de $f$ sont géométriquement normales ».
On a déjà vu que $\mathcal{P}_1$ et $\mathcal{P}_2$
sont locales pour la topologie fppf à la source et au but, voir le
lemme \ref{more-morphisms-lemma-locally-Noetherian-fibres-fppf-local-source-and-target},
et Descente, lemmes \ref{descent-lemma-descending-property-flat} et
\ref{descent-lemma-flat-fpqc-local-source}. Il reste donc à traiter
$\mathcal{P}_3$.

\medskip\noindent
Soit $f : X \to Y$ un morphisme de schémas.
Soit $\{\varphi_i : Y_i \to Y\}_{i \in I}$ un recouvrement fpqc de $Y$.
Notons $f_i : X_i \to Y_i$ le changement de base de $f$ par $\varphi_i$.
Soit $i \in I$ et soit $y_i \in Y_i$ un point.
Posons $y = \varphi_i(y_i)$. Notons que
$$
X_{i, y_i} = \Spec(\kappa(y_i)) \times_{\Spec(\kappa(y))} X_y.
$$
Dans cette situation, $X_y$ est donc géométriquement normal si et
seulement si $X_{i, y_i}$ est géométriquement normal, voir
Variétés, lemme \ref{varieties-lemma-geometrically-normal-upstairs}.
Ce fait montre que $\mathcal{P}_3$ est locale pour la topologie fpqc au but.

\medskip\noindent
Soit $\{X_i \to X\}$ un recouvrement lisse de $X$.
Soit $y \in Y$. Dans ce cas, $\{X_{i, y} \to X_y\}$ est un
recouvrement lisse de la fibre. Le caractère local de $\mathcal{P}_3$
pour la topologie lisse à la source résulte donc de
Descente, lemme \ref{descent-lemma-normal-local-smooth}.
En combinant ces observations, on obtient le lemme.
\end{proof}
\section{Morphismes réguliers}
\label{more-morphisms-section-regular}

\noindent
Comparer avec la section \ref{more-morphisms-section-normal}. La version algébrique de
cette notion est étudiée dans
Compléments d'algèbre, section \ref{more-algebra-section-regular}.

\begin{definition}
\label{more-morphisms-definition-regular}
Soit $f : X \to Y$ un morphisme de schémas.
Supposons que toutes les fibres $X_y$ soient des schémas localement noethériens.
\begin{enumerate}
\item Soient $x \in X$ et $y = f(x)$. On dit que $f$ est {\it régulier en $x$}
si $f$ est plat en $x$ et si le schéma $X_y$ est géométriquement
régulier en $x$ sur $\kappa(y)$ (voir
Variétés, définition \ref{varieties-definition-geometrically-regular}).
\item On dit que $f$ est un {\it morphisme régulier} si $f$ est régulier
en tout point de $X$.
\end{enumerate}
\end{definition}

\noindent
La régularité du morphisme $X \to Y$ est une condition
plus forte que celle demandant seulement que toutes les
fibres soient des schémas réguliers localement noethériens.

\begin{lemma}
\label{more-morphisms-lemma-regular}
Soit $f : X \to Y$ un morphisme de schémas.
Supposons toutes les fibres de $f$ localement noethériennes.
Les conditions suivantes sont équivalentes :
\begin{enumerate}
\item $f$ est régulier,
\item $f$ est plat et ses fibres sont des schémas géométriquement réguliers,
\item pour toute paire d'ouverts affines $U \subset X$, $V \subset Y$
telle que $f(U) \subset V$, l'homomorphisme d'anneaux $\mathcal{O}_Y(V) \to \mathcal{O}_X(U)$
est régulier,
\item il existe un recouvrement ouvert $Y = \bigcup_{j \in J} V_j$
et des recouvrements ouverts $f^{-1}(V_j) = \bigcup_{i \in I_j} U_i$ tels
que chacun des morphismes $U_i \to V_j$ soit régulier, et
\item il existe un recouvrement ouvert affine $Y = \bigcup_{j \in J} V_j$
et des recouvrements ouverts affines $f^{-1}(V_j) = \bigcup_{i \in I_j} U_i$ tels
que les homomorphismes d'anneaux $\mathcal{O}_Y(V_j) \to \mathcal{O}_X(U_i)$ soient réguliers.
\end{enumerate}
\end{lemma}

\begin{proof}
L'équivalence de (1) et (2) résulte immédiatement des définitions.
Soit $x \in X$ et posons $y = f(x)$. Par définition, $f$ est plat en $x$
si et seulement si $\mathcal{O}_{Y, y} \to \mathcal{O}_{X, x}$ est un
homomorphisme d'anneaux plat, et $X_y$ est géométriquement régulier en $x$ sur
$\kappa(y)$ si et seulement si
$\mathcal{O}_{X_y, x} = \mathcal{O}_{X, x}/\mathfrak m_y\mathcal{O}_{X, x}$
est une algèbre géométriquement régulière sur $\kappa(y)$.
La régularité de $f$ en $x$
ne dépend donc que de l'homomorphisme local d'anneaux locaux
$\mathcal{O}_{Y, y} \to \mathcal{O}_{X, x}$.
L'équivalence de (1) et (4) est donc claire.

\medskip\noindent
Rappelons (Compléments d'algèbre, définition \ref{more-algebra-definition-regular})
qu'un homomorphisme d'anneaux $A \to B$ est régulier si et seulement s'il est plat
et si les anneaux des fibres $B \otimes_A \kappa(\mathfrak p)$ sont noethériens
et géométriquement réguliers pour tout idéal premier $\mathfrak p \subset A$.
D'après Variétés, lemme \ref{varieties-lemma-geometrically-regular},
cela équivaut à ce que $\Spec(B \otimes_A \kappa(\mathfrak p))$
soit un schéma géométriquement régulier sur $\kappa(\mathfrak p)$.
Ainsi, (2) implique (3). Il est clair que (3) implique
(5). Enfin, supposons (5). Alors $f$ est plat
(voir Morphismes, lemme \ref{morphisms-lemma-flat-characterize}).
De plus, si $y \in Y$, alors $y \in V_j$ pour un certain $j$, et
$X_y = \bigcup_{i \in I_j} U_{i, y}$, chaque $U_{i, y}$
étant géométriquement régulier sur $\kappa(y)$ d'après
Variétés, lemme \ref{varieties-lemma-geometrically-regular}.
Une nouvelle application de
Variétés, lemme \ref{varieties-lemma-geometrically-regular},
montre que $X_y$ est géométriquement régulier. Ainsi, (2) est vérifiée,
ce qui achève la démonstration du lemme.
\end{proof}

\begin{lemma}
\label{more-morphisms-lemma-smooth-regular}
Un morphisme lisse est régulier.
\end{lemma}

\begin{proof}
Soit $f : X \to Y$ un morphisme lisse.
Puisque $f$ est localement de présentation finie, voir
Morphismes, lemme \ref{morphisms-lemma-smooth-locally-finite-presentation},
les fibres $X_y$ sont localement de type fini sur un corps, donc
localement noethériennes. De plus, $f$ est plat, voir
Morphismes, lemme \ref{morphisms-lemma-smooth-flat}.
Enfin, les fibres $X_y$ sont lisses sur un corps (d'après
Morphismes, lemme \ref{morphisms-lemma-base-change-smooth}),
donc géométriquement régulières d'après
Variétés, lemme \ref{varieties-lemma-smooth-geometrically-normal}.
Ainsi, $f$ est régulier d'après le
lemme \ref{more-morphisms-lemma-regular}.
\end{proof}

\begin{lemma}
\label{more-morphisms-lemma-regular-fppf-local-source-and-target}
La propriété $\mathcal{P}(f)=$« les fibres de $f$ sont
localement noethériennes et $f$ est régulier »
est locale pour la topologie fppf au but et
locale pour la topologie lisse à la source.
\end{lemma}

\begin{proof}
On a
$\mathcal{P}(f) =
\mathcal{P}_1(f) \wedge \mathcal{P}_2(f) \wedge \mathcal{P}_3(f)$,
où
$\mathcal{P}_1(f)=$« les fibres de $f$ sont localement noethériennes »,
$\mathcal{P}_2(f)=$« $f$ est plat », et
$\mathcal{P}_3(f)=$« les fibres de $f$ sont géométriquement régulières ».
On a déjà vu que $\mathcal{P}_1$ et $\mathcal{P}_2$
sont locales pour la topologie fppf à la source et au but, voir le
lemme \ref{more-morphisms-lemma-locally-Noetherian-fibres-fppf-local-source-and-target},
et Descente, lemmes \ref{descent-lemma-descending-property-flat} et
\ref{descent-lemma-flat-fpqc-local-source}. Il reste donc à traiter
$\mathcal{P}_3$.

\medskip\noindent
Soit $f : X \to Y$ un morphisme de schémas.
Soit $\{\varphi_i : Y_i \to Y\}_{i \in I}$ un recouvrement fpqc de $Y$.
Notons $f_i : X_i \to Y_i$ le changement de base de $f$ par $\varphi_i$.
Soit $i \in I$ et soit $y_i \in Y_i$ un point.
Posons $y = \varphi_i(y_i)$. Notons que
$$
X_{i, y_i} = \Spec(\kappa(y_i)) \times_{\Spec(\kappa(y))} X_y.
$$
Dans cette situation, $X_y$ est donc géométriquement régulier si et
seulement si $X_{i, y_i}$ est géométriquement régulier, voir
Variétés, lemme \ref{varieties-lemma-geometrically-regular-upstairs}.
Ce fait montre que $\mathcal{P}_3$ est locale pour la topologie fpqc au but.

\medskip\noindent
Soit $\{X_i \to X\}$ un recouvrement lisse de $X$.
Soit $y \in Y$. Dans ce cas, $\{X_{i, y} \to X_y\}$ est un
recouvrement lisse de la fibre. Le caractère local de $\mathcal{P}_3$
pour la topologie lisse à la source résulte donc de
Descente, lemme \ref{descent-lemma-regular-local-smooth}.
En combinant ces observations, on obtient le lemme.
\end{proof}
\section{Morphismes de Cohen-Macaulay}
\label{more-morphisms-section-CM}

\noindent
Comparer avec la section \ref{more-morphisms-section-normal}.
Notons que, comme il est signalé dans
Algèbre, section \ref{algebra-section-geometrically-CM},
et dans
Variétés, section \ref{varieties-section-CM},
« géométriquement de Cohen-Macaulay » équivaut simplement à « de Cohen-Macaulay ».

\begin{definition}
\label{more-morphisms-definition-CM}
Soit $f : X \to Y$ un morphisme de schémas.
Supposons que toutes les fibres $X_y$ soient des schémas localement noethériens.
\begin{enumerate}
\item Soient $x \in X$ et $y = f(x)$. On dit que $f$ est
{\it de Cohen-Macaulay en $x$} si $f$ est plat en $x$ et si
l'anneau local du schéma $X_y$ en $x$ est de Cohen-Macaulay.
\item On dit que $f$ est un {\it morphisme de Cohen-Macaulay} si $f$ est
de Cohen-Macaulay en tout point de $X$.
\end{enumerate}
\end{definition}

\noindent
Voici une reformulation.

\begin{lemma}
\label{more-morphisms-lemma-CM}
Soit $f : X \to Y$ un morphisme de schémas.
Supposons toutes les fibres de $f$ localement noethériennes.
Les conditions suivantes sont équivalentes :
\begin{enumerate}
\item $f$ est de Cohen-Macaulay, et
\item $f$ est plat et ses fibres sont des schémas de Cohen-Macaulay.
\end{enumerate}
\end{lemma}

\begin{proof}
Cela résulte directement des définitions.
\end{proof}

\begin{lemma}
\label{more-morphisms-lemma-CM-dimension}
Soit $f : X \to Y$ un morphisme de schémas localement noethériens,
localement de type fini et de Cohen-Macaulay.
Pour tout point $x$ de $X$ d'image $y$ dans $Y$,
$$
\dim_x(X) = \dim_y(Y) + \dim_x(X_y),
$$
où $X_y$ désigne la fibre au-dessus de $y$.
\end{lemma}

\begin{proof}
Après remplacement de $X$ par un voisinage ouvert de $x$,
il existe un entier naturel $d$ tel que toutes les fibres
de $X \to Y$ soient de dimension $d$ en tout point, voir
Morphismes, lemme
\ref{morphisms-lemma-flat-finite-presentation-CM-fibres-relative-dimension}.
Alors $f$ est plat, localement de type fini
et de dimension relative $d$. Le résultat découle donc de
Morphismes, lemme \ref{morphisms-lemma-rel-dimension-dimension}.
\end{proof}

\begin{lemma}
\label{more-morphisms-lemma-composition-CM}
Soient $f : X \to Y$ et $g : Y \to Z$ des morphismes de schémas. Supposons que les
fibres de $f$, $g$ et $g \circ f$ soient localement noethériennes.
Soit $x \in X$ d'images $y \in Y$ et $z \in Z$.
\begin{enumerate}
\item Si $f$ est de Cohen-Macaulay en $x$ et si $g$ est de Cohen-Macaulay
en $f(x)$, alors $g \circ f$ est de Cohen-Macaulay en $x$.
\item Si $f$ et $g$ sont de Cohen-Macaulay, alors $g \circ f$ est de Cohen-Macaulay.
\item Si $g \circ f$ est de Cohen-Macaulay en $x$ et si $f$ est plat en $x$,
alors $f$ est de Cohen-Macaulay en $x$ et $g$ est de Cohen-Macaulay en $f(x)$.
\item Si $g \circ f$ est de Cohen-Macaulay et si $f$ est plat, alors
$f$ est de Cohen-Macaulay et $g$ est de Cohen-Macaulay en tout point
de l'image de $f$.
\end{enumerate}
\end{lemma}

\begin{proof}
Considérons l'homomorphisme d'anneaux locaux noethériens
$$
\mathcal{O}_{Y_z, y} \to \mathcal{O}_{X_z, x}
$$
et observons que sa fibre est
$$
\mathcal{O}_{X_z, x}/\mathfrak m_{Y_z, y}\mathcal{O}_{X_z, x} =
\mathcal{O}_{X_y, x}
$$
Le lemme résulte donc du
lemme \ref{algebra-lemma-CM-goes-up} d'Algèbre.
\end{proof}

\begin{lemma}
\label{more-morphisms-lemma-flat-morphism-from-CM-scheme}
Soit $f : X \to Y$ un morphisme plat de schémas localement noethériens.
Si $X$ est de Cohen-Macaulay, alors $f$ est de Cohen-Macaulay et
$\mathcal{O}_{Y, f(x)}$ est de Cohen-Macaulay pour tout $x \in X$.
\end{lemma}

\begin{proof}
Après traduction en algèbre, cela résulte du
lemme \ref{algebra-lemma-CM-goes-up} d'Algèbre.
\end{proof}

\begin{lemma}
\label{more-morphisms-lemma-base-change-CM}
Soit $f : X \to Y$ un morphisme de schémas.
Supposons que toutes les fibres $X_y$ soient des schémas localement noethériens.
Soit $Y' \to Y$ un morphisme localement de type fini. Soit $f' : X' \to Y'$
le changement de base de $f$.
Soit $x' \in X'$ un point d'image $x \in X$.
\begin{enumerate}
\item Si $f$ est de Cohen-Macaulay en $x$, alors
$f' : X' \to Y'$ est de Cohen-Macaulay en $x'$.
\item Si $f$ est plat en $x$ et si $f'$ est de Cohen-Macaulay en $x'$, alors $f$
est de Cohen-Macaulay en $x$.
\item Si $Y' \to Y$ est plat en $f'(x')$ et si $f'$ est de Cohen-Macaulay en
$x'$, alors $f$ est de Cohen-Macaulay en $x$.
\end{enumerate}
\end{lemma}

\begin{proof}
Notons que l'hypothèse sur $Y' \to Y$ implique que, pour $y' \in Y'$
d'image $y \in Y$, l'extension de corps $\kappa(y')/\kappa(y)$
est de type fini. Toutes les fibres
$X'_{y'} = (X_y)_{\kappa(y')}$ sont donc elles aussi localement noethériennes, voir
Variétés, lemme \ref{varieties-lemma-locally-Noetherian-base-change}.
Le lemme a donc un sens. Posons $y' = f'(x')$ et $y = f(x)$.
On obtient le diagramme commutatif d'anneaux locaux suivant :
$$
\xymatrix{
\mathcal{O}_{X', x'} & \mathcal{O}_{X, x} \ar[l] \\
\mathcal{O}_{Y', y'} \ar[u] & \mathcal{O}_{Y, y} \ar[l] \ar[u]
}
$$
où l'anneau en haut à gauche est un localisé du produit tensoriel des anneaux
en haut à droite et en bas à gauche sur l'anneau en bas à droite.

\medskip\noindent
Supposons $f$ de Cohen-Macaulay en $x$.
La platitude de $\mathcal{O}_{Y, y} \to \mathcal{O}_{X, x}$
implique celle de $\mathcal{O}_{Y', y'} \to \mathcal{O}_{X', x'}$, voir
Algèbre, lemme \ref{algebra-lemma-base-change-flat-up-down}.
Le fait que $\mathcal{O}_{X, x}/\mathfrak m_y\mathcal{O}_{X, x}$
soit de Cohen-Macaulay implique que
$\mathcal{O}_{X', x'}/\mathfrak m_{y'}\mathcal{O}_{X', x'}$
est de Cohen-Macaulay, voir
Variétés, lemme \ref{varieties-lemma-CM-base-change}. Ainsi, $f'$
est de Cohen-Macaulay en $x'$.

\medskip\noindent
Supposons $f$ plat en $x$ et $f'$ de Cohen-Macaulay en $x'$.
Le fait que $\mathcal{O}_{X', x'}/\mathfrak m_{y'}\mathcal{O}_{X', x'}$
soit de Cohen-Macaulay implique que
$\mathcal{O}_{X, x}/\mathfrak m_y\mathcal{O}_{X, x}$
est de Cohen-Macaulay, voir
Variétés, lemme \ref{varieties-lemma-CM-base-change}.
Ainsi, $f$ est de Cohen-Macaulay en $x$.

\medskip\noindent
Supposons $Y' \to Y$ plat en $y'$ et $f'$ de Cohen-Macaulay en
$x'$. La platitude de $\mathcal{O}_{Y', y'} \to \mathcal{O}_{X', x'}$
et de $\mathcal{O}_{Y, y} \to \mathcal{O}_{Y', y'}$ implique celle
de $\mathcal{O}_{Y, y} \to \mathcal{O}_{X, x}$, voir
Algèbre, lemme \ref{algebra-lemma-base-change-flat-up-down}.
Le fait que $\mathcal{O}_{X', x'}/\mathfrak m_{y'}\mathcal{O}_{X', x'}$
soit de Cohen-Macaulay implique que
$\mathcal{O}_{X, x}/\mathfrak m_y\mathcal{O}_{X, x}$
est de Cohen-Macaulay, voir
Variétés, lemme \ref{varieties-lemma-CM-base-change}. Ainsi, $f$
est de Cohen-Macaulay en $x$.
\end{proof}

\begin{lemma}
\label{more-morphisms-lemma-flat-finite-presentation-CM-open}
\begin{reference}
\cite[IV, corollaire 12.1.7(iii)]{EGA}
\end{reference}
Soit $f : X \to S$ un morphisme de schémas plat et localement
de présentation finie. Posons
$$
W = \{x \in X \mid f\text{ est de Cohen-Macaulay en }x\}
$$
Alors
\begin{enumerate}
\item $W = \{x \in X \mid \mathcal{O}_{X_{f(x)}, x}\text{ est de Cohen-Macaulay}\}$,
\item $W$ est ouvert dans $X$,
\item $W$ est dense dans toute fibre de $X \to S$,
\item la formation de $W$ commute à tout changement de base de $f$ :
pour tout morphisme $g : S' \to S$, considérons
le changement de base $f' : X' \to S'$ de $f$ et la
projection $g' : X' \to X$. Alors l'ensemble
$W'$ correspondant au morphisme $f'$ est égal à $W' = (g')^{-1}(W)$.
\end{enumerate}
\end{lemma}

\begin{proof}
Puisque $f$ est plat à fibres localement noethériennes, l'égalité de (1)
vaut par définition. Les assertions (2) et (3) résultent du
lemme \ref{algebra-lemma-generic-CM-flat-finite-presentation} d'Algèbre.
L'assertion (4) résulte soit du
lemme \ref{algebra-lemma-CM-locus-commutes-base-change} d'Algèbre,
soit de
Variétés, lemme \ref{varieties-lemma-CM-base-change}.
\end{proof}

\begin{lemma}
\label{more-morphisms-lemma-flat-finite-presentation-characterize-CM}
Soit $f : X \to S$ un morphisme de schémas plat et localement
de présentation finie. Soit $x \in X$ d'image $s \in S$.
Posons $d = \dim_x(X_s)$.
Les conditions suivantes sont équivalentes :
\begin{enumerate}
\item $f$ est de Cohen-Macaulay en $x$,
\item il existe un voisinage ouvert $U \subset X$ de $x$
et un morphisme localement quasi-fini $U \to \mathbf{A}^d_S$ sur $S$
qui est plat en $x$,
\item il existe un voisinage ouvert $U \subset X$ de $x$
et un morphisme plat localement quasi-fini $U \to \mathbf{A}^d_S$ sur $S$,
\item pour tout $S$-morphisme $g : U \to \mathbf{A}^d_S$
d'un voisinage ouvert $U \subset X$ de $x$, on a :
$g$ est quasi-fini en $x$ $\Rightarrow$ $g$ est plat en $x$.
\end{enumerate}
\end{lemma}

\begin{proof}
L'ouverture du lieu de platitude montre que (2) et (3)
sont équivalentes, voir le théorème \ref{more-morphisms-theorem-openness-flatness}.

\medskip\noindent
Choisissons des ouverts affines $U = \Spec(A) \subset X$ avec $x \in U$ et
$V = \Spec(R) \subset S$ avec $f(U) \subset V$. Alors $R \to A$
est un homomorphisme d'anneaux plat de présentation finie. Soit $\mathfrak p \subset A$
l'idéal premier correspondant à $x$. Après remplacement de $A$ par un
localisé principal, on peut supposer qu'il existe un homomorphisme quasi-fini
$R[x_1, \ldots, x_d]  \to A$, voir
Algèbre, lemme \ref{algebra-lemma-quasi-finite-over-polynomial-algebra}.
Il existe donc au moins un couple $(U, g)$ formé d'un
voisinage ouvert $U \subset X$ de $x$ et d'un morphisme
localement\footnote{Si $S$ est quasi-séparé, alors $g$ sera quasi-fini.} quasi-fini
$g : U \to \mathbf{A}^d_S$.

\medskip\noindent
Affirmation : étant donnés $R \to A$ plat et de présentation finie, un idéal premier
$\mathfrak p \subset A$ et $\varphi : R[x_1, \ldots, x_d] \to A$
quasi-fini en $\mathfrak p$, on a : $\Spec(\varphi)$
est plat en $\mathfrak p$ si et seulement si $\Spec(A) \to \Spec(R)$
est de Cohen-Macaulay en $\mathfrak p$. En effet, d'après le
théorème \ref{more-morphisms-theorem-criterion-flatness-fibre},
la platitude peut se vérifier sur les fibres. Il en est de même
de la propriété d'être de Cohen-Macaulay (puisque $A$ est déjà supposé plat sur $R$).
L'affirmation résulte donc du
lemme \ref{algebra-lemma-where-CM} d'Algèbre.

\medskip\noindent
Cette affirmation montre que (1) équivaut à (4) ; compte tenu
de la construction d'un couple convenable $(U, g)$ au deuxième
paragraphe, elle montre aussi que (1) équivaut à (2).
\end{proof}

\begin{lemma}
\label{more-morphisms-lemma-flat-finite-presentation-CM-pieces}
Soit $f : X \to S$ un morphisme de schémas plat et localement
de présentation finie. Pour $d \geq 0$, il existe des ouverts $U_d \subset X$
ayant les propriétés suivantes :
\begin{enumerate}
\item $W = \bigcup_{d \geq 0} U_d$ est dense dans toute fibre de $f$, et
\item $U_d \to S$ est de dimension relative $d$ (voir
Morphismes, définition \ref{morphisms-definition-relative-dimension-d}).
\end{enumerate}
\end{lemma}

\begin{proof}
Cela résulte de la combinaison du
lemme \ref{more-morphisms-lemma-flat-finite-presentation-CM-open}
avec
Morphismes, lemme
\ref{morphisms-lemma-flat-finite-presentation-CM-fibres-relative-dimension}.
\end{proof}

\begin{lemma}
\label{more-morphisms-lemma-flat-finite-presentation-specialization-dimension}
Soit $f : X \to S$ un morphisme de schémas plat et localement
de présentation finie.
Supposons que $x' \leadsto x$ soit une spécialisation de points de $X$,
d'image $s' \leadsto s$ dans $S$. Si $x$ est un point générique d'une
composante irréductible de $X_s$, alors $\dim_{x'}(X_{s'}) = \dim_x(X_s)$.
\end{lemma}

\begin{proof}
Le point $x$ appartient à $U_d$ pour un certain $d$, où $U_d$ est comme au
lemme \ref{more-morphisms-lemma-flat-finite-presentation-CM-pieces}.
\end{proof}

\begin{lemma}
\label{more-morphisms-lemma-CM-local-source-and-target}
La propriété
$\mathcal{P}(f)=$« les fibres de $f$ sont localement noethériennes et $f$ est
de Cohen-Macaulay » est locale pour la topologie fppf au but et
locale pour la topologie syntomique à la source.
\end{lemma}

\begin{proof}
On a
$\mathcal{P}(f) =
\mathcal{P}_1(f) \wedge \mathcal{P}_2(f)$,
où
$\mathcal{P}_1(f)=$« $f$ est plat », et
$\mathcal{P}_2(f)=$« les fibres de $f$ sont localement noethériennes
et de Cohen-Macaulay ».
On sait que $\mathcal{P}_1$ est
locale pour la topologie fppf à la source et au but, voir
Descente, lemmes \ref{descent-lemma-descending-property-flat} et
\ref{descent-lemma-flat-fpqc-local-source}. Il reste donc à traiter
$\mathcal{P}_2$.

\medskip\noindent
Soit $f : X \to Y$ un morphisme de schémas.
Soit $\{\varphi_i : Y_i \to Y\}_{i \in I}$ un recouvrement fppf de $Y$.
Notons $f_i : X_i \to Y_i$ le changement de base de $f$ par $\varphi_i$.
Soit $i \in I$ et soit $y_i \in Y_i$ un point.
Posons $y = \varphi_i(y_i)$. Notons que
$$
X_{i, y_i} = \Spec(\kappa(y_i)) \times_{\Spec(\kappa(y))} X_y.
$$
et que $\kappa(y_i)/\kappa(y)$ est une extension de corps
de type fini. Ainsi, si $X_y$ est localement noethérien,
$X_{i, y_i}$ est localement noethérien, voir
Variétés, lemme \ref{varieties-lemma-locally-Noetherian-base-change}.
Si, de plus, $X_y$ est de Cohen-Macaulay,
alors $X_{i, y_i}$ est de Cohen-Macaulay, voir
Variétés, lemme \ref{varieties-lemma-CM-base-change}.
Ainsi, $\mathcal{P}_2$ est locale pour la topologie fppf au but.

\medskip\noindent
Soit $\{X_i \to X\}$ un recouvrement syntomique de $X$.
Soit $y \in Y$. Dans ce cas, $\{X_{i, y} \to X_y\}$ est un
recouvrement syntomique de la fibre. Le caractère local de $\mathcal{P}_2$
pour la topologie syntomique à la source résulte donc de
Descente, lemme \ref{descent-lemma-CM-local-syntomic}.
En combinant ces observations, on obtient le lemme.
\end{proof}

\section{sections des morphismes de Cohen-Macaulay}
\label{more-morphisms-section-slicing-CM}

\noindent
Les résultats de cette section conduisent finalement à l'assertion selon laquelle
la topologie fppf est la même que la topologie
« de présentation finie, plate et quasi-finie ».
Le lemme suivant est très proche de
Diviseurs, lemme \ref{divisors-lemma-fibre-Cartier}.

\begin{lemma}
\label{more-morphisms-lemma-slice-given-element}
Soit $f : X \to S$ un morphisme de schémas.
Soit $x \in X$ un point d'image $s \in S$.
Soit $h \in \mathfrak m_x \subset \mathcal{O}_{X, x}$.
Supposons que
\begin{enumerate}
\item $f$ soit localement de présentation finie,
\item $f$ soit plat en $x$, et
\item l'image $\overline{h}$ de $h$ dans
$\mathcal{O}_{X_s, x} = \mathcal{O}_{X, x}/\mathfrak m_s\mathcal{O}_{X, x}$
soit non diviseur de zéro.
\end{enumerate}
Alors il existe un voisinage ouvert affine $U \subset X$ de $x$
tel que $h$ provienne d'un élément $h \in \Gamma(U, \mathcal{O}_U)$ et
tel que $D = V(h)$ soit un diviseur de Cartier effectif de $U$, avec $x \in D$ et
$D \to S$ plat et localement de présentation finie.
\end{lemma}

\begin{proof}
Nous allons le démontrer en nous ramenant au cas noethérien.
Par ouverture du lieu de platitude (voir le
théorème \ref{more-morphisms-theorem-openness-flatness}),
on peut supposer, après remplacement de $X$ par un
voisinage ouvert de $x$, que $X \to S$ est plat.
On peut aussi supposer $X$ et $S$ affines.
Après restriction éventuelle de $X$, on peut supposer qu'il existe
un élément $h \in \Gamma(X, \mathcal{O}_X)$ dont l'image est l'élément $h$ donné.

\medskip\noindent
On peut écrire $S = \Spec(A)$ et $A = \colim_i A_i$
comme limite inductive filtrante de $\mathbf{Z}$-algèbres de type fini.
Alors, d'après
Algèbre, lemme \ref{algebra-lemma-flat-finite-presentation-limit-flat},
ou
Limites, lemmes \ref{limits-lemma-descend-finite-presentation},
\ref{limits-lemma-descend-affine-finite-presentation} et
\ref{limits-lemma-descend-finite-presentation},
on peut trouver un diagramme cartésien
$$
\xymatrix{
X \ar[r] \ar[d]_f & X_0 \ar[d]^{f_0} \\
S \ar[r] & S_0
}
$$
avec $f_0$ plat et de présentation finie, $X_0$ affine et
$S_0$ affine et noethérien. Soit $x_0 \in X_0$, resp.\ $s_0 \in S_0$,
l'image de $x$, resp.\ $s$. On peut aussi supposer qu'il existe un élément
$h_0 \in \Gamma(X_0, \mathcal{O}_{X_0})$ qui se restreint à $h$ sur $X$.
(Si l'on utilise la référence algébrique ci-dessus, c'est clair ; si l'on utilise
les références au chapitre sur les limites, cela résulte de
Limites, lemme \ref{limits-lemma-descend-finite-presentation},
en considérant $h$ comme un morphisme $X \to \mathbf{A}^1_S$.)
Notons que $\mathcal{O}_{X_s, x}$ est un localisé de
$\mathcal{O}_{(X_0)_{s_0}, x_0} \otimes_{\kappa(s_0)} \kappa(s)$, de sorte que
$\mathcal{O}_{(X_0)_{s_0}, x_0} \to \mathcal{O}_{X_s, x}$ est un homomorphisme
local plat, donc fidèlement plat. Par conséquent, l'image
$\overline{h}_0 \in \mathcal{O}_{(X_0)_{s_0}, x_0}$
appartient à $\mathfrak m_{(X_0)_{s_0}, x_0}$ et est non diviseur de zéro.
Nous affirmons qu'après remplacement de $X_0$ par un voisinage ouvert principal de
$x_0$, l'élément $h_0$ est non diviseur de zéro dans
$B_0 = \Gamma(X_0, \mathcal{O}_{X_0})$ et $B_0/h_0B_0$ est plat
sur $A_0 = \Gamma(S_0, \mathcal{O}_{S_0})$.
Dans ce cas,
$$
0 \to B_0 \xrightarrow{h_0} B_0 \to B_0/h_0B_0 \to 0
$$
est une suite exacte courte de $A_0$-modules plats. Elle reste donc exacte
après tensorisation par $A$ (d'après
Algèbre, lemme \ref{algebra-lemma-flat-tor-zero}),
et le lemme s'ensuit.

\medskip\noindent
Il reste à démontrer l'affirmation ci-dessus. L'énoncé algébrique correspondant
est le suivant (on omet ici l'indice ${}_0$) :
soit $A \to B$ un homomorphisme d'anneaux noethériens, plat et de type fini.
Soit $\mathfrak q \subset B$ un idéal premier au-dessus de $\mathfrak p \subset A$.
Supposons que $h \in \mathfrak q$ ait pour image un non-diviseur de zéro dans
$B_{\mathfrak q}/\mathfrak p B_{\mathfrak q}$.
Il s'agit de montrer qu'après remplacement éventuel de $B$ par $B_g$ pour un
$g \in B$, $g \not \in \mathfrak q$, l'élément $h$ devient non diviseur de zéro
et $B/hB$ devient plat sur $A$. D'après
Algèbre, lemme \ref{algebra-lemma-grothendieck},
$h$ est non diviseur de zéro dans $B_{\mathfrak q}$ et
$B_{\mathfrak q}/hB_{\mathfrak q}$ est plat sur $A$.
Par ouverture du lieu de platitude, voir
Algèbre, théorème \ref{algebra-theorem-openness-flatness},
ou le
théorème \ref{more-morphisms-theorem-openness-flatness},
$B/hB$ est plat sur $A$ après remplacement de $B$ par $B_g$ pour un
$g \in B$, $g \not \in \mathfrak q$. Enfin, soit
$I = \{b \in B \mid hb = 0\}$ l'annulateur de $h$. Alors
$IB_{\mathfrak q} = 0$, puisque $h$ est non diviseur de zéro dans $B_{\mathfrak q}$.
De plus, $I$ est de type fini, car $B$ est noethérien.
Il existe donc un $g \in B$, $g \not \in \mathfrak q$, tel que $IB_g = 0$.
Après remplacement de $B$ par $B_g$, l'élément $h$ est non diviseur de zéro.
\end{proof}

\begin{lemma}
\label{more-morphisms-lemma-slice-given-elements}
Soit $f : X \to S$ un morphisme de schémas.
Soit $x \in X$ un point d'image $s \in S$.
Soient $h_1, \ldots, h_r \in \mathcal{O}_{X, x}$.
Supposons que
\begin{enumerate}
\item $f$ soit localement de présentation finie,
\item $f$ soit plat en $x$, et
\item les images de $h_1, \ldots, h_r$ dans
$\mathcal{O}_{X_s, x} = \mathcal{O}_{X, x}/\mathfrak m_s\mathcal{O}_{X, x}$
forment une suite régulière.
\end{enumerate}
Alors il existe un voisinage ouvert affine $U \subset X$ de $x$
tel que $h_1, \ldots, h_r$ proviennent d'éléments
$h_1, \ldots, h_r \in \Gamma(U, \mathcal{O}_U)$ et
tel que $Z = V(h_1, \ldots, h_r) \to U$ soit une immersion régulière, avec
$x \in Z$ et $Z \to S$ plat et localement de présentation finie.
De plus, le changement de base $Z_{S'} \to U_{S'}$ est une immersion régulière
pour tout schéma $S'$ sur $S$.
\end{lemma}

\begin{proof}
(Nos conventions sur les suites régulières impliquent que $h_i \in \mathfrak m_x$
pour tout $i$.) Le cas $r = 1$ résulte du
lemme \ref{more-morphisms-lemma-slice-given-element}
combiné avec
Diviseurs, lemme \ref{divisors-lemma-relative-Cartier},
qui montre que $V(h_1)$ reste un diviseur de Cartier effectif après changement de base.
Le cas $r > 1$ résulte d'une récurrence immédiate sur $r$ (en appliquant
le résultat pour $r = 1$ exactement $r$ fois ; les détails sont omis).

\medskip\noindent
Une autre démonstration du lemme utilise les résultats de
Diviseurs, section \ref{divisors-section-relative-regular-immersion}.
En effet, d'abord, par ouverture du lieu de platitude (voir le
théorème \ref{more-morphisms-theorem-openness-flatness}),
on peut supposer, après remplacement de $X$ par un
voisinage ouvert de $x$, que $X \to S$ est plat.
On peut aussi supposer $X$ et $S$ affines.
Après restriction éventuelle de $X$, on peut supposer que l'on a
$h_1, \ldots, h_r \in \Gamma(X, \mathcal{O}_X)$. Posons
$Z = V(h_1, \ldots, h_r)$. Notons que $X_s$ est un schéma noethérien
(car c'est un $\kappa(s)$-schéma algébrique, voir
Variétés, section \ref{varieties-section-algebraic-schemes})
et que la topologie de $X_s$ est induite par celle de $X$
(voir
Schémas, lemme \ref{schemes-lemma-fibre-topological}).
Ainsi, en restreignant encore $X$,
on peut supposer que $Z_s \subset X_s$ est une immersion régulière
définie par les $r$ éléments $h_i|_{X_s}$, voir
Diviseurs, lemme \ref{divisors-lemma-Noetherian-scheme-regular-ideal}
et sa démonstration. Il est aussi clair que $r = \dim_x(X_s) - \dim_x(Z_s)$, car
\begin{align*}
\dim_x(X_s) & = \dim(\mathcal{O}_{X_s, x}) +
\text{trdeg}_{\kappa(s)}(\kappa(x)), \\
\dim_x(Z_s) & = \dim(\mathcal{O}_{Z_s, x}) +
\text{trdeg}_{\kappa(s)}(\kappa(x)), \\
\dim(\mathcal{O}_{X_s, x}) & = \dim(\mathcal{O}_{Z_s, x}) + r
\end{align*}
les deux premières égalités résultant du
lemme \ref{algebra-lemma-dimension-at-a-point-finite-type-field} d'Algèbre,
et la troisième de $r$ applications du
lemme \ref{algebra-lemma-one-equation} d'Algèbre.
Ainsi,
Diviseurs, lemme \ref{divisors-lemma-fibre-quasi-regular}, assertion (3),
s'applique et montre qu'après restriction de $X$ pour la topologie de Zariski, le morphisme
$Z \to X$ est une immersion régulière à laquelle s'applique
Diviseurs, lemme
\ref{divisors-lemma-relative-regular-immersion-flat-in-neighbourhood}
(ce qui donne la platitude et l'assertion sur le changement de base).
\end{proof}

\begin{lemma}
\label{more-morphisms-lemma-slice-once}
Soit $f : X \to S$ un morphisme de schémas.
Soit $x \in X$ un point d'image $s \in S$.
Supposons que
\begin{enumerate}
\item $f$ soit localement de présentation finie,
\item $f$ soit plat en $x$, et
\item $\mathcal{O}_{X_s, x}$ vérifie $\text{depth} \geq 1$.
\end{enumerate}
Alors il existe un voisinage ouvert affine $U \subset X$ de $x$
et un diviseur de Cartier effectif $D \subset U$ contenant $x$ tels que
$D \to S$ soit plat et de présentation finie.
\end{lemma}

\begin{proof}
Choisir un élément $h \in \mathfrak m_x \subset \mathcal{O}_{X, x}$ dont
l'image soit non diviseur de zéro dans $\mathcal{O}_{X_s, x}$ et appliquer le
lemme \ref{more-morphisms-lemma-slice-given-element}.
\end{proof}

\begin{lemma}
\label{more-morphisms-lemma-slice-CM}
\begin{reference}
\cite[IV Proposition 17.16.1]{EGA}
\end{reference}
Soit $f : X \to S$ un morphisme de schémas.
Soit $x \in X$ un point d'image $s \in S$.
Supposons que
\begin{enumerate}
\item $f$ soit localement de présentation finie,
\item $f$ soit de Cohen-Macaulay en $x$, et
\item $x$ soit un point fermé de $X_s$.
\end{enumerate}
Alors il existe une immersion régulière $Z \to X$ contenant $x$ telle que
\begin{enumerate}
\item[(a)] $Z \to S$ soit plat et localement de présentation finie,
\item[(b)] $Z \to S$ soit localement quasi-fini, et
\item[(c)] $Z_s = \{x\}$ ensemblistement.
\end{enumerate}
\end{lemma}

\begin{proof}
On peut remplacer, et l'on remplace, $S$ par un voisinage ouvert affine de $s$.
Nous démontrerons le lemme pour $S$ affine par récurrence sur $d = \dim_x(X_s)$.

\medskip\noindent
Cas $d = 0$. Dans ce cas, nous montrons que l'on peut prendre pour $Z$
un voisinage ouvert de $x$. (Notons qu'une immersion ouverte est
une immersion régulière.) En effet, si $d = 0$, alors $X \to S$
est quasi-fini en $x$, voir
Morphismes, lemme \ref{morphisms-lemma-locally-quasi-finite-rel-dimension-0}.
Il existe donc un voisinage ouvert affine $U \subset X$ tel
que $U \to S$ soit quasi-fini, voir
Morphismes, lemme \ref{morphisms-lemma-quasi-finite-points-open}.
Ainsi, après remplacement de $X$ par $U$, la fibre $X_s$ est un ensemble
fini discret. En remplaçant $X$ par un nouveau voisinage ouvert affine
de $X$, on obtient donc $f^{-1}(\{s\}) = \{x\}$ (car la topologie
de $X_s$ est induite par celle de $X$, voir
Schémas, lemme \ref{schemes-lemma-fibre-topological}).
Cela démontre le lemme dans ce cas.

\medskip\noindent
Supposons maintenant $d > 0$. Puisque $x$ est un point fermé de sa
fibre, l'extension $\kappa(x)/\kappa(s)$ est finie (d'après le
théorème des zéros de Hilbert, voir
Morphismes, lemme \ref{morphisms-lemma-closed-point-fibre-locally-finite-type}).
On a donc
$$
\text{depth}(\mathcal{O}_{X_s, x}) = \dim(\mathcal{O}_{X_s, x}) = d > 0
$$
la première égalité provenant de ce que $\mathcal{O}_{X_s, x}$ est de Cohen-Macaulay,
et la seconde de
Morphismes, lemme \ref{morphisms-lemma-dimension-fibre-at-a-point}.
On peut donc appliquer le
lemme \ref{more-morphisms-lemma-slice-once}
pour trouver un diagramme
$$
\xymatrix{
D \ar[r] \ar[rrd] & U \ar[r] \ar[rd] & X \ar[d] \\
& & S
}
$$
avec $x \in D$. Notons que
$\mathcal{O}_{D_s, x} = \mathcal{O}_{X_s, x}/(\overline{h})$ pour un certain
non-diviseur de zéro $\overline{h}$, voir
Diviseurs, lemme \ref{divisors-lemma-relative-Cartier}.
Ainsi, $\mathcal{O}_{D_s, x}$ est de Cohen-Macaulay, de dimension
inférieure d'une unité à celle de $\mathcal{O}_{X_s, x}$, voir
Algèbre, lemme \ref{algebra-lemma-reformulate-CM},
par exemple. Le morphisme $D \to S$ est donc plat,
localement de présentation finie et de Cohen-Macaulay en $x$, avec
$\dim_x(D_s) = \dim_x(X_s) - 1 = d - 1$. Par l'hypothèse de récurrence,
on peut trouver une immersion régulière $Z \to D$ ayant les propriétés (a), (b), (c).
Puisque les deux morphismes $Z \to D \to U$ sont des immersions régulières,
$Z \to U$ est aussi une immersion régulière d'après
Diviseurs, lemme \ref{divisors-lemma-composition-regular-immersion}.
Ceci achève la démonstration.
\end{proof}

\begin{lemma}
\label{more-morphisms-lemma-qf-fp-flat-neighbourhood-dominates-fppf}
Soit $f : X \to S$ un morphisme plat de schémas
localement de présentation finie. Soit $s \in S$ un point de l'image de $f$.
Alors il existe un diagramme commutatif
$$
\xymatrix{
S' \ar[rr] \ar[rd]_g & & X \ar[ld]^f \\
& S
}
$$
où $g : S' \to S$ est plat, localement de présentation finie,
localement quasi-fini, et où $s \in g(S')$.
\end{lemma}

\begin{proof}
La fibre $X_s$ est non vide par hypothèse. Il existe donc un point
fermé $x \in X_s$ où $f$ est de Cohen-Macaulay, voir le
lemme \ref{more-morphisms-lemma-flat-finite-presentation-CM-open}.
Appliquer le
lemme \ref{more-morphisms-lemma-slice-CM}
et poser $S' = Z$.
\end{proof}

\noindent
Le lemme suivant montre que les faisceaux pour la topologie fppf
sont exactement les faisceaux pour la topologie
« quasi-finie, plate et de présentation finie ».

\begin{lemma}
\label{more-morphisms-lemma-qf-fp-flat-dominates-fppf}
Soit $S$ un schéma. Soit $\mathcal{U} = \{S_i \to S\}_{i \in I}$ un recouvrement
fppf de $S$, voir
Topologies, définition \ref{topologies-definition-fppf-covering}.
Alors il existe un recouvrement fppf $\mathcal{V} = \{T_j \to S\}_{j \in J}$
qui raffine (voir
Sites, définition \ref{sites-definition-morphism-coverings})
$\mathcal{U}$ et tel que chaque morphisme $T_j \to S$ soit localement quasi-fini.
\end{lemma}

\begin{proof}
Pour tout $s \in S$, il existe un $i \in I$ tel que $s$ appartienne à
l'image de $S_i \to S$. D'après le
lemme \ref{more-morphisms-lemma-qf-fp-flat-neighbourhood-dominates-fppf},
on peut trouver un morphisme $g_s : T_s \to S$ tel que $s \in g_s(T_s)$,
plat, localement de présentation finie et localement quasi-fini,
et tel que $g_s$ se factorise par $S_i \to S$. Ainsi,
$\{T_s \to S\}$ est le recouvrement recherché de $S$ raffinant $\mathcal{U}$.
\end{proof}

\section{Fibres génériques}
\label{more-morphisms-section-generic}

\noindent
Quelques résultats sur les rapports entre les fibres génériques et les
fibres voisines.

\begin{lemma}
\label{more-morphisms-lemma-empty-generic-fibre}
Soit $f : X \to Y$ un morphisme de type fini de schémas. Supposons que
$Y$ soit irréductible, de point générique $\eta$. Si $X_\eta = \emptyset$,
alors il existe un ouvert non vide $V \subset Y$ tel que
$X_V = V \times_Y X = \emptyset$.
\end{lemma}

\begin{proof}
Cela résulte immédiatement du résultat plus général
Morphismes,
lemme \ref{morphisms-lemma-quasi-compact-generic-point-not-in-image}.
\end{proof}

\begin{lemma}
\label{more-morphisms-lemma-nonempty-generic-fibre}
Soit $f : X \to Y$ un morphisme de type fini de schémas. Supposons que
$Y$ soit irréductible, de point générique $\eta$. Si $X_\eta \not = \emptyset$,
alors il existe un ouvert non vide $V \subset Y$ tel que
$X_V = V \times_Y X \to V$ soit surjectif.
\end{lemma}

\begin{proof}
Cela résulte, en prenant des ouverts affines, de
Algèbre, lemme \ref{algebra-lemma-characterize-image-finite-type}.
(Bien entendu, cela résulte aussi de la platitude générique.)
\end{proof}

\begin{lemma}
\label{more-morphisms-lemma-nowhere-dense-generic-fibre}
Soit $f : X \to Y$ un morphisme de type fini de schémas. Supposons que
$Y$ soit irréductible, de point générique $\eta$.
Si $Z \subset X$ est un sous-ensemble fermé tel que $Z_\eta$ soit nulle part dense
dans $X_\eta$, alors il existe un ouvert non vide $V \subset Y$ tel
que $Z_y$ soit nulle part dense dans $X_y$ pour tout $y \in V$.
\end{lemma}

\begin{proof}
Soit $Y' \subset Y$ le réduit de $Y$.
Posons $X' = Y' \times_Y X$ et $Z' = Y' \times_Y Z$.
Comme $Y' \to Y$ est un homéomorphisme universel d'après
Morphismes, lemme \ref{morphisms-lemma-reduction-universal-homeomorphism},
il suffit de démontrer le lemme pour $Z' \subset X' \to Y'$.
On peut donc supposer que $Y$ est intègre, voir
Propriétés, lemme \ref{properties-lemma-characterize-integral}.
D'après
Morphismes, proposition \ref{morphisms-proposition-generic-flatness},
il existe un ouvert affine non vide $V \subset Y$ tel que
$X_V \to V$ et $Z_V \to V$ soient plats et de présentation finie.
Nous affirmons que $V$ convient.
Choisissons $y \in V$. Si $Z_y$ est d'intérieur non vide, alors $Z_y$ contient
un point générique $\xi$ d'une composante irréductible de $X_y$.
Notons que $\eta \leadsto f(\xi)$. Comme $Z_V \to V$ est plat, on peut
choisir une spécialisation $\xi' \leadsto \xi$, $\xi' \in Z$,
avec $f(\xi') = \eta$, voir
Morphismes, lemme \ref{morphisms-lemma-generalizations-lift-flat}.
D'après le
lemme \ref{more-morphisms-lemma-flat-finite-presentation-specialization-dimension},
on a
$$
\dim_{\xi'}(Z_\eta) = \dim_{\xi}(Z_y) = \dim_{\xi}(X_y) = \dim_{\xi'}(X_\eta).
$$
Ainsi, une composante irréductible de $Z_\eta$ passant par $\xi'$ est
de dimension $\dim_{\xi'}(X_\eta)$, ce qui contredit l'hypothèse selon laquelle
$Z_\eta$ est nulle part dense dans $X_\eta$, et le résultat est démontré.
\end{proof}

\begin{lemma}
\label{more-morphisms-lemma-scheme-theoretically-dense-generic-fibre}
Soit $f : X \to Y$ un morphisme de type fini de schémas.
Supposons que $Y$ soit irréductible, de point générique $\eta$.
Soit $U \subset X$ un sous-schéma ouvert tel que $U_\eta$ soit
schématiquement dense dans $X_\eta$.
Alors il existe un ouvert non vide $V \subset Y$ tel
que $U_y$ soit schématiquement dense dans $X_y$ pour tout $y \in V$.
\end{lemma}

\begin{proof}
Soit $Y' \subset Y$ le réduit de $Y$.
Posons $X' = Y' \times_Y X$ et $U' = Y' \times_Y U$.
Comme $Y' \to Y$ induit une bijection sur les points, et comme
$U' \to U$ et $X' \to X$ induisent des isomorphismes sur les fibres schématiques,
on peut remplacer $Y$ par $Y'$ et $X$ par $X'$.
On peut donc supposer que $Y$ est intègre, voir
Propriétés, lemme \ref{properties-lemma-characterize-integral}.
On peut aussi remplacer $Y$ par un ouvert affine non vide. Autrement dit, on
peut supposer que $Y = \Spec(A)$, où $A$ est un anneau intègre de corps des
fractions $K$.

\medskip\noindent
Puisque $f$ est de type fini, $X$ est quasi-compact.
Écrivons $X = X_1 \cup \ldots \cup X_n$, où les $X_i$ sont des ouverts affines. D'après
Morphismes, définition \ref{morphisms-definition-scheme-theoretically-dense},
$U_i = X_i \cap U$ est un sous-schéma ouvert de $X_i$ tel que
$U_{i, \eta}$ soit schématiquement dense dans $X_{i, \eta}$.
Il suffit donc de démontrer le résultat pour les couples $(X_i, U_i)$ ;
autrement dit, on peut supposer que $X$ est affine.

\medskip\noindent
Écrivons $X = \Spec(B)$. Notons que $B_K$ est noethérien, puisque c'est une
$K$-algèbre de type fini. Ainsi, $U_\eta$ est quasi-compact. On peut donc
trouver un nombre fini d'éléments $g_1, \ldots, g_m \in B$ tels que $D(g_j) \subset U$
et tels que $U_\eta = D(g_1)_\eta \cup \ldots \cup D(g_m)_\eta$.
Le fait que $U_\eta$ soit schématiquement dense dans
$X_\eta$ signifie que $B_K \to \bigoplus_j (B_K)_{g_j}$
est injectif, voir
Morphismes, exemple \ref{morphisms-example-scheme-theoretic-closure}.
D'après
Algèbre, lemme \ref{algebra-lemma-when-injective-covering},
cela équivaut à l'injectivité de
$B_K \to \bigoplus\nolimits_{j = 1, \ldots, m} B_K$,
$b \mapsto (g_1b, \ldots, g_mb)$. Soit $M$ le conoyau de cette
application sur $A$, c'est-à-dire que l'on a une suite exacte
$$
0 \to I \to B \xrightarrow{(g_1, \ldots, g_m)}
\bigoplus\nolimits_{j = 1, \ldots, m} B \to M \to 0
$$
Quitte à remplacer $A$ par $A_h$ pour un certain $h$ non nul, on peut supposer que $B$
est une $A$-algèbre plate et de présentation finie, et que $M$
est plat sur $A$, voir
Algèbre, lemme \ref{algebra-lemma-generic-flatness}.
La platitude de $B$ sur $A$ implique que $B$ est sans torsion comme
$A$-module, voir
Compléments sur l'algèbre, lemme \ref{more-algebra-lemma-flat-torsion-free}.
Ainsi $B \subset B_K$. Par hypothèse, $I_K = 0$, ce qui implique que $I = 0$
(car $I \subset B \subset B_K$ est inclus dans $I_K$). On a donc maintenant
une suite exacte courte
$$
0 \to B \xrightarrow{(g_1, \ldots, g_m)}
\bigoplus\nolimits_{j = 1, \ldots, m} B \to M \to 0
$$
où $M$ est plat sur $A$. Pour tout homomorphisme $A \to \kappa$, où
$\kappa$ est un corps, on obtient donc une suite exacte courte
$$
0 \to B \otimes_A \kappa \xrightarrow{(g_1 \otimes 1, \ldots, g_m \otimes 1)}
\bigoplus\nolimits_{j = 1, \ldots, m} B \otimes_A \kappa \to
M \otimes_A \kappa \to 0
$$
voir
Algèbre, lemme \ref{algebra-lemma-flat-tor-zero}.
En reprenant en sens inverse les arguments ci-dessus,
cela signifie que $\bigcup D(g_j \otimes 1)$ est schématiquement
dense dans $\Spec(B \otimes_A \kappa)$.
Comme $\bigcup D(g_j \otimes 1) = \bigcup D(g_j)_\kappa \subset U_\kappa$,
on en déduit que $U_\kappa$ est schématiquement dense dans $X_\kappa$,
ce qu'il fallait démontrer.
\end{proof}

\noindent
Soient un morphisme de schémas $f : X \to Y$ et
un point $y \in Y$. Rappelons que la fibre $X_y$ est homéomorphe
au sous-ensemble $f^{-1}(\{y\})$ de $X$ muni de la topologie induite, voir
Schémas, lemme \ref{schemes-lemma-fibre-topological}.
Soit un sous-ensemble fermé $T(y) \subset X_y$.
Soit $T$ l'adhérence de $T(y)$ dans $X$.
Munissons $T$ de la structure de schéma réduit induite.
Alors $T$ est un sous-schéma fermé de $X$ tel
que $T_y = T(y)$ ensemblistement. En fait, $T$ est le plus petit
sous-schéma fermé de $X$ ayant cette propriété. Il est donc « sans inconvénient »
de désigner un sous-ensemble fermé de $X_y$ par $T_y$ si on le souhaite.
Dans le lemme suivant, on applique ceci à la fibre générique de $f$.

\begin{lemma}
\label{more-morphisms-lemma-cover-generic-fibre-neighbourhood}
Soit $f : X \to Y$ un morphisme de type fini de schémas. Supposons que
$Y$ soit irréductible, de point générique $\eta$. Soit
$X_\eta = Z_{1, \eta} \cup \ldots \cup Z_{n, \eta}$ un recouvrement de
la fibre générique par des sous-ensembles fermés de $X_\eta$.
Soit $Z_i$ l'adhérence de $Z_{i, \eta}$ dans $X$ (voir la discussion ci-dessus).
Alors il existe un ouvert non vide $V \subset Y$ tel
que $X_y = Z_{1, y} \cup \ldots \cup Z_{n, y}$ pour tout $y \in V$.
\end{lemma}

\begin{proof}
Si $Y$ est noethérien, alors $U = X \setminus (Z_1 \cup \ldots \cup Z_n)$
est de type fini sur $Y$, et on peut appliquer directement le
lemme \ref{more-morphisms-lemma-empty-generic-fibre}
pour obtenir $U_V = \emptyset$ pour un ouvert non vide $V \subset Y$.
Dans le cas général, raisonnons comme suit. Puisque la question est topologique,
on peut remplacer $Y$ par son réduit. Ainsi $Y$ est intègre, voir
Propriétés, lemme \ref{properties-lemma-characterize-integral}.
Quitte à rétrécir $Y$, on peut supposer que $X \to Y$ est plat, voir
Morphismes, proposition \ref{morphisms-proposition-generic-flatness}.
Dans ce cas, tout point $x$ de $X_y$ est une spécialisation d'un point
$x' \in X_\eta$, d'après
Morphismes, lemme \ref{morphisms-lemma-generalizations-lift-flat}.
Comme les $Z_i$ sont fermés dans $X$ et recouvrent la fibre générique, ceci
implique que $X_y = \bigcup Z_{i, y}$ pour $y \in Y$, comme désiré.
\end{proof}

\noindent
Le lemme suivant affirme que les fibres génériques de morphismes dont la source est
réduite sont réduites.

\begin{lemma}
\label{more-morphisms-lemma-reduction-generic-fibre}
Soit $f : X \to Y$ un morphisme de schémas. Soit $\eta \in Y$ un point générique
d'une composante irréductible de $Y$. Alors
$(X_\eta)_{red} = (X_{red})_\eta$.
\end{lemma}

\begin{proof}
Choisissons un voisinage affine $\Spec(A) \subset Y$ de $\eta$.
Choisissons un ouvert affine $\Spec(B) \subset X$ envoyé dans
$\Spec(A)$ par le morphisme $f$. Soit $\mathfrak p \subset A$ l'idéal premier minimal
correspondant à $\eta$. Soit $B_{red}$ le quotient de $B$ par
le nilradical $\sqrt{(0)}$. Le contenu algébrique du lemme est que
$C = B_{red} \otimes_A \kappa(\mathfrak p)$ est réduit.
Notons $I \subset A$ le nilradical, de sorte que $A_{red} = A/I$.
Notons $\mathfrak p_{red} = \mathfrak p/I$ ;
c'est un idéal premier minimal de $A_{red}$ et
$\kappa(\mathfrak p) = \kappa(\mathfrak p_{red})$.
Comme $A \to B_{red}$ et $A \to \kappa(\mathfrak p)$
se factorisent tous deux par $A \to A_{red}$, on a
$C = B_{red} \otimes_{A_{red}} \kappa(\mathfrak p_{red})$.
Or $\kappa(\mathfrak p_{red}) = (A_{red})_{\mathfrak p_{red}}$
est une localisation d'après
Algèbre, lemme \ref{algebra-lemma-minimal-prime-reduced-ring}.
Ainsi $C$ est une localisation de $B_{red}$
(Algèbre, lemme \ref{algebra-lemma-tensor-localization})
et est donc réduit.
\end{proof}

\begin{lemma}
\label{more-morphisms-lemma-make-generic-fibre-geometrically-reduced}
Soit $f : X \to Y$ un morphisme de schémas.
Supposons que $Y$ soit irréductible et que $f$ soit de type fini.
Il existe un diagramme
$$
\xymatrix{
X' \ar[d]_{f'} \ar[r]_{g'} & X_V \ar[r] \ar[d] & X \ar[d]^f \\
Y' \ar[r]^g & V \ar[r] & Y
}
$$
où
\begin{enumerate}
\item $V$ est un ouvert non vide de $Y$,
\item $X_V = V \times_Y X$,
\item $g : Y' \to V$ est un homéomorphisme universel fini,
\item $X' = (Y' \times_Y X)_{red} = (Y' \times_V X_V)_{red}$,
\item $g'$ est un homéomorphisme universel fini,
\item $Y'$ est un schéma affine intègre,
\item $f'$ est plat et de présentation finie, et
\item la fibre générique de $f'$ est géométriquement réduite.
\end{enumerate}
\end{lemma}

\begin{proof}
Soit $V = \Spec(A)$ un ouvert affine non vide de $Y$.
Par hypothèse, le nilradical $\mathfrak p = \sqrt{(0)}$ de $A$
est un idéal premier.
Soit $K = \kappa(\mathfrak p)$.
Soit $p$ la caractéristique de $K$ si elle est positive, et $1$
si elle est nulle. D'après
Variétés, lemme \ref{varieties-lemma-finite-extension-geometrically-reduced},
il existe une extension de corps finie purement inséparable
$K'/K$ telle que $(X_{K'})_{red}$ soit géométriquement réduit sur $K'$.
Choisissons des éléments $x_1, \ldots, x_n \in K'$ qui engendrent $K'$ sur
$K$ et tels qu'une puissance d'exposant une puissance de $p$ de $x_i$ appartienne à $A/\mathfrak p$.
Soit $A' \subset K'$ la $A$-sous-algèbre finie de $K'$ engendrée par
$x_1, \ldots, x_n$. Notons que $A'$ est un anneau intègre de corps des fractions $K'$. D'après
Algèbre, lemme \ref{algebra-lemma-p-ring-map},
l'application $A \to A'$ induit un homéomorphisme universel sur les spectres.
Posons $Y' = \Spec(A')$. Posons $X' = (Y' \times_Y X)_{red}$.
La fibre générique de $X' \to Y'$ est $(X_{K'})_{red}$ d'après le
lemme \ref{more-morphisms-lemma-reduction-generic-fibre},
et elle est géométriquement réduite par construction.
Notons que $X' \to X_V$ est un homéomorphisme universel fini, car c'est la
composée du morphisme de réduction $X' \to Y' \times_Y X$ (voir
Morphismes, lemme \ref{morphisms-lemma-reduction-universal-homeomorphism})
et du changement de base de $g$.
À ce stade, toutes les propriétés du lemme sont satisfaites, à l'exception
possible de (7). On peut obtenir cette dernière en rétrécissant $Y'$, et donc $V$, voir
Morphismes, proposition \ref{morphisms-proposition-generic-flatness}.
\end{proof}

\begin{lemma}
\label{more-morphisms-lemma-make-components-generic-fibre-geometrically-irreducible}
Soit $f : X \to Y$ un morphisme de schémas.
Supposons que $Y$ soit irréductible et que $f$ soit de type fini.
Il existe un diagramme
$$
\xymatrix{
X' \ar[d]_{f'} \ar[r]_{g'} & X_V \ar[r] \ar[d] & X \ar[d]^f \\
Y' \ar[r]^g & V \ar[r] & Y
}
$$
où
\begin{enumerate}
\item $V$ est un ouvert non vide de $Y$,
\item $X_V = V \times_Y X$,
\item $g : Y' \to V$ est fini étale surjectif,
\item $X' = Y' \times_Y X = Y' \times_V X_V$,
\item $g'$ est fini étale surjectif,
\item $Y'$ est un schéma affine irréductible, et
\item toutes les composantes irréductibles de la fibre générique de $f'$
sont géométriquement irréductibles.
\end{enumerate}
\end{lemma}

\begin{proof}
Soit $V = \Spec(A)$ un ouvert affine non vide de $Y$.
Par hypothèse, le nilradical $\mathfrak p = \sqrt{(0)}$ de $A$ est un idéal premier.
Soit $K = \kappa(\mathfrak p)$. D'après
Variétés, lemme
\ref{varieties-lemma-finite-extension-geometrically-irreducible-components},
il existe une extension de corps finie séparable
$K'/K$ telle que toutes les composantes irréductibles de $X_{K'}$ soient
géométriquement irréductibles sur $K'$.
Choisissons un élément $\alpha \in K'$ qui engendre $K'$ sur
$K$, voir
Corps, lemme \ref{fields-lemma-primitive-element}.
Soit $P(T) \in K[T]$ le polynôme minimal de $\alpha$ sur $K$.
Quitte à remplacer $\alpha$ par $f \alpha$ pour un certain
$f \in A$, $f \not \in \mathfrak p$,
on peut supposer qu'il existe un polynôme unitaire
$T^d + a_1T^{d - 1} + \ldots + a_d \in A[T]$ dont l'image est
$P(T) \in K[T]$ par l'application $A[T] \to K[T]$.
Posons $A' = A[T]/(P)$. Alors $A \to A'$ est un homomorphisme d'anneaux fini libre
tel qu'il existe un unique idéal premier $\mathfrak q$ au-dessus de
$\mathfrak p$, que
$K = \kappa(\mathfrak p) \subset \kappa(\mathfrak q) = K'$
est une extension finie séparable, et que $\mathfrak pA'_{\mathfrak q}$
est l'idéal maximal de $A'_{\mathfrak q}$.
Ainsi $g : Y' = \Spec(A') \to V = \Spec(A)$
est étale en $\mathfrak q$, voir
Algèbre, lemme \ref{algebra-lemma-characterize-etale}.
Cela signifie qu'il existe un ouvert $W \subset \Spec(A')$ tel
que $g|_W : W \to \Spec(A)$ soit étale.
Comme $g$ est fini et comme $\mathfrak q$ est le seul point au-dessus de
$\mathfrak p$, on voit que $Z = g(Y' \setminus W)$ est un sous-ensemble fermé de $V$
qui ne contient pas $\mathfrak p$. Ainsi, quitte à remplacer $V$ par un ouvert affine
principal de $V$ qui ne rencontre pas $Z$, on obtient que $g$ est fini
étale.
\end{proof}

\section{Assassins relatifs}
\label{more-morphisms-section-assassin}

\begin{lemma}
\label{more-morphisms-lemma-relative-assassin-in-neighbourhood}
Soit $f : X \to S$ un morphisme de schémas.
Soit $\mathcal{F}$ un $\mathcal{O}_X$-module quasi-cohérent.
Soit $\xi \in \text{Ass}_{X/S}(\mathcal{F})$ et posons
$Z = \overline{\{\xi\}} \subset X$.
Si $f$ est localement de type fini et si $\mathcal{F}$ est un
$\mathcal{O}_X$-module de type fini, alors il existe un ouvert non vide
$V \subset Z$ tel que, pour tout $s \in f(V)$, les points génériques
de $V_s$ appartiennent à $\text{Ass}_{X/S}(\mathcal{F})$.
\end{lemma}

\begin{proof}
On peut remplacer $S$ par un voisinage ouvert affine de $f(\xi)$
et $X$ par un voisinage ouvert affine de $\xi$. On peut donc supposer
$S = \Spec(A)$, $X = \Spec(B)$ et que $f$ est donné par
l'homomorphisme d'anneaux de type fini $A \to B$, voir
Morphismes, lemme \ref{morphisms-lemma-locally-finite-type-characterize}.
De plus, on peut écrire $\mathcal{F} = \widetilde{M}$ pour un certain
$B$-module de type fini $M$, voir
Propriétés, lemme \ref{properties-lemma-finite-type-module}.
Soit $\mathfrak q \subset B$ l'idéal premier correspondant à $\xi$, et
soit $\mathfrak p \subset A$ l'idéal premier correspondant de $A$.
Par hypothèse, $\mathfrak q \in \text{Ass}_B(M \otimes_A \kappa(\mathfrak p))$,
voir
Algèbre, remarque \ref{algebra-remark-bourbaki}
et
Diviseurs, lemme \ref{divisors-lemma-associated-affine-open}.
Avec ces notations, $Z = V(\mathfrak q) \subset \Spec(B)$.
En particulier, $f(Z) \subset V(\mathfrak p)$. Il suffit donc clairement de
démontrer le lemme après avoir remplacé $A$, $B$ et $M$ par
$A/\mathfrak pA$, $B/\mathfrak pB$ et $M/\mathfrak pM$.
Autrement dit, on peut supposer que $A$ est un anneau intègre de corps des fractions
$K$ et que $\mathfrak q \subset B$ est un idéal premier associé de $M \otimes_A K$.

\medskip\noindent
On peut maintenant utiliser la platitude générique. En effet, d'après
Algèbre, lemme \ref{algebra-lemma-generic-flatness},
il existe un élément non nul $g \in A$ tel que $M_g$ soit plat
comme $A_g$-module. Après avoir remplacé $A$ par $A_g$, on peut supposer que
$M$ est plat comme $A$-module.

\medskip\noindent
Dans ce cas, d'après
Algèbre, lemme \ref{algebra-lemma-post-bourbaki},
$\mathfrak q$ est aussi un idéal premier associé de $M$.
On obtient donc une application injective de $B$-modules $B/\mathfrak q \to M$.
Soit $Q$ son conoyau ; on obtient ainsi une suite exacte courte
$$
0 \to B/\mathfrak q \to M \to Q \to 0
$$
de $B$-modules de type fini. En appliquant une nouvelle fois la platitude générique,
Algèbre, lemme \ref{algebra-lemma-generic-flatness},
cette fois au $B$-module $Q$, on peut supposer que $Q$
est un $A$-module plat. En particulier, on peut supposer que la suite exacte
courte ci-dessus est universellement injective, voir
Algèbre, lemme \ref{algebra-lemma-flat-tor-zero}.
Dans cette situation,
$(B/\mathfrak q) \otimes_A \kappa(\mathfrak p')
\subset M \otimes_A \kappa(\mathfrak p')$
pour tout idéal premier $\mathfrak p'$ de $A$. Le lemme en résulte, car tout idéal
premier minimal $\mathfrak q'$ du support de
$(B/\mathfrak q) \otimes_A \kappa(\mathfrak p')$
est un idéal premier associé de $(B/\mathfrak q) \otimes_A \kappa(\mathfrak p')$ d'après
Diviseurs, lemme \ref{divisors-lemma-minimal-support-in-ass}.
\end{proof}

\begin{lemma}
\label{more-morphisms-lemma-bad-case}
Soit $f : X \to Y$ un morphisme de schémas. Soit $\mathcal{F}$ un
$\mathcal{O}_X$-module quasi-cohérent. Soit $U \subset X$ un sous-schéma
ouvert. Supposons que
\begin{enumerate}
\item $f$ soit de type fini,
\item $\mathcal{F}$ soit de type fini,
\item $Y$ soit irréductible, de point générique $\eta$, et
\item $\text{Ass}_{X_\eta}(\mathcal{F}_\eta)$ ne soit pas contenu dans $U_\eta$.
\end{enumerate}
Alors il existe un sous-schéma ouvert non vide $V \subset Y$ tel que,
pour tout $y \in V$, l'ensemble $\text{Ass}_{X_y}(\mathcal{F}_y)$ ne soit pas
contenu dans $U_y$.
\end{lemma}

\begin{proof}
Soit $Z \subset X$ le support schématique de $\mathcal{F}$, voir
Morphismes, définition \ref{morphisms-definition-scheme-theoretic-support}.
Alors $Z_\eta$ est le support schématique de $\mathcal{F}_\eta$
(Morphismes, lemme \ref{morphisms-lemma-flat-pullback-support}).
Ainsi, les points génériques des composantes irréductibles de $Z_\eta$
appartiennent à $\text{Ass}_{X_\eta}(\mathcal{F}_\eta)$ d'après
Diviseurs, lemme \ref{divisors-lemma-minimal-support-in-ass}.
On voit donc que $Z_\eta \cap U_\eta = \emptyset$.
Ainsi $T = Z \setminus U$ est un sous-ensemble fermé de $Z$ tel que $T_\eta = \emptyset$.
Si l'on munit $T$ de la structure de schéma réduit induite, alors
$T \to Y$ est un morphisme de type fini. D'après le
lemme \ref{more-morphisms-lemma-empty-generic-fibre},
il existe un ouvert non vide $V \subset Y$ tel que $T_V = \emptyset$.
Alors $V$ convient.
\end{proof}

\begin{lemma}
\label{more-morphisms-lemma-good-case}
Soit $f : X \to Y$ un morphisme de schémas. Soit $\mathcal{F}$ un
$\mathcal{O}_X$-module quasi-cohérent. Soit $U \subset X$ un sous-schéma
ouvert. Supposons que
\begin{enumerate}
\item $f$ soit de type fini,
\item $\mathcal{F}$ soit de type fini,
\item $Y$ soit irréductible, de point générique $\eta$, et
\item $\text{Ass}_{X_\eta}(\mathcal{F}_\eta) \subset U_\eta$.
\end{enumerate}
Alors il existe un sous-schéma ouvert non vide $V \subset Y$ tel que,
pour tout $y \in V$, on ait $\text{Ass}_{X_y}(\mathcal{F}_y) \subset U_y$.
\end{lemma}

\begin{proof}
(Cette démonstration est identique à celle du
lemme \ref{more-morphisms-lemma-scheme-theoretically-dense-generic-fibre}.
Nous invitons le lecteur à lire d'abord cette dernière.)
Puisque l'énoncé porte sur les fibres, on peut clairement remplacer
$Y$ par son réduit. On peut donc supposer que $Y$ est intègre, voir
Propriétés, lemme \ref{properties-lemma-characterize-integral}.
On peut aussi supposer que $Y = \Spec(A)$ est affine. Alors $A$
est un anneau intègre de corps des fractions $K$.

\medskip\noindent
Puisque $f$ est de type fini, $X$ est quasi-compact.
Écrivons $X = X_1 \cup \ldots \cup X_n$, où les $X_i$ sont des ouverts affines,
et posons $\mathcal{F}_i = \mathcal{F}|_{X_i}$. Par
hypothèse, la fibre générique de $U_i = X_i \cap U$ contient
$\text{Ass}_{X_{i, \eta}}(\mathcal{F}_{i, \eta})$.
Il suffit donc de démontrer le résultat pour les triplets
$(X_i, \mathcal{F}_i, U_i)$ ;
autrement dit, on peut supposer que $X$ est affine.

\medskip\noindent
Écrivons $X = \Spec(B)$. Soit $N$ un $B$-module de type fini tel que
$\mathcal{F} = \widetilde{N}$.
Notons que $B_K$ est noethérien, puisque c'est une
$K$-algèbre de type fini. Ainsi $U_\eta$ est quasi-compact. On peut donc
trouver un nombre fini d'éléments $g_1, \ldots, g_m \in B$ tels que $D(g_j) \subset U$
et tels que $U_\eta = D(g_1)_\eta \cup \ldots \cup D(g_m)_\eta$.
Puisque $\text{Ass}_{X_\eta}(\mathcal{F}_\eta) \subset U_\eta$,
l'application $N_K \to \bigoplus_j (N_K)_{g_j}$ est injective. D'après
Algèbre, lemme \ref{algebra-lemma-when-injective-covering},
cela équivaut à l'injectivité de
$N_K \to \bigoplus\nolimits_{j = 1, \ldots, m} N_K$,
$n \mapsto (g_1n, \ldots, g_mn)$. Soient $I$ et $M$ le noyau et le
conoyau de cette application sur $A$, c'est-à-dire que l'on a une suite exacte
$$
0 \to I \to N \xrightarrow{(g_1, \ldots, g_m)}
\bigoplus\nolimits_{j = 1, \ldots, m} N \to M \to 0
$$
Quitte à remplacer $A$ par $A_h$ pour un certain $h$ non nul, on peut supposer que
$B$ est une $A$-algèbre plate et de présentation finie et que
$M$ et $N$ sont tous deux plats sur $A$, voir
Algèbre, lemme \ref{algebra-lemma-generic-flatness}.
La platitude de $N$ sur $A$ implique que $N$ est sans torsion comme
$A$-module, voir
Compléments sur l'algèbre, lemme \ref{more-algebra-lemma-flat-torsion-free}.
Ainsi $N \subset N_K$. Par construction, $I_K = 0$, ce qui implique que $I = 0$
(car $I \subset N \subset N_K$ est inclus dans $I_K$). On a donc maintenant
une suite exacte courte
$$
0 \to N \xrightarrow{(g_1, \ldots, g_m)}
\bigoplus\nolimits_{j = 1, \ldots, m} N \to M \to 0
$$
où $M$ est plat sur $A$. Pour tout homomorphisme $A \to \kappa$, où
$\kappa$ est un corps, on obtient donc une suite exacte courte
$$
0 \to N \otimes_A \kappa \xrightarrow{(g_1 \otimes 1, \ldots, g_m \otimes 1)}
\bigoplus\nolimits_{j = 1, \ldots, m} N \otimes_A \kappa \to
M \otimes_A \kappa \to 0
$$
voir
Algèbre, lemme \ref{algebra-lemma-flat-tor-zero}.
En reprenant en sens inverse les arguments ci-dessus,
cela signifie que $\bigcup D(g_j \otimes 1)$ contient
$\text{Ass}_{B \otimes_A \kappa}(N \otimes_A \kappa)$.
Comme $\bigcup D(g_j \otimes 1) = \bigcup D(g_j)_\kappa \subset U_\kappa$,
on en déduit que $U_\kappa$ contient
$\text{Ass}_{X \otimes \kappa}(\mathcal{F} \otimes \kappa)$,
ce qu'il fallait démontrer.
\end{proof}

\begin{lemma}
\label{more-morphisms-lemma-base-change-assassin-in-U}
Soit $f : X \to S$ un morphisme localement de type fini.
Soit $\mathcal{F}$ un $\mathcal{O}_X$-module quasi-cohérent
de type fini. Soit $U \subset X$ un sous-schéma ouvert.
Soit $g : S' \to S$ un morphisme de schémas, soit
$f' : X' = X_{S'} \to S'$ le changement de base de $f$,
soit $g' : X' \to X$ la projection, posons
$\mathcal{F}' = (g')^*\mathcal{F}$, et posons
$U' = (g')^{-1}(U)$. Enfin, soit $s' \in S'$ d'image $s = g(s')$.
Dans ce cas,
$$
\text{Ass}_{X_s}(\mathcal{F}_s) \subset U_s
\Leftrightarrow
\text{Ass}_{X'_{s'}}(\mathcal{F}'_{s'}) \subset U'_{s'}.
$$
\end{lemma}

\begin{proof}
Cela résulte immédiatement de
Diviseurs, lemme \ref{divisors-lemma-base-change-relative-assassin}.
Voir aussi
Diviseurs, remarque \ref{divisors-remark-base-change-relative-assassin}.
\end{proof}

\begin{lemma}
\label{more-morphisms-lemma-relative-assassin-constructible}
Soit $f : X \to Y$ un morphisme de présentation finie.
Soit $\mathcal{F}$ un $\mathcal{O}_X$-module quasi-cohérent
de présentation finie. Soit $U \subset X$ un sous-schéma ouvert
tel que $U \to Y$ soit quasi-compact. Alors l'ensemble
$$
E = \{y \in Y \mid \text{Ass}_{X_y}(\mathcal{F}_y) \subset U_y\}
$$
est localement constructible dans $Y$.
\end{lemma}

\begin{proof}
Soit $y \in Y$. Il faut montrer qu'il existe un voisinage ouvert
$V$ de $y$ dans $Y$ tel que $E \cap V$ soit constructible dans $V$. On peut donc
supposer que $Y$ est affine. Écrivons $Y = \Spec(A)$ et
$A = \colim A_i$ comme limite inductive filtrante de
$\mathbf{Z}$-algèbres de type fini. D'après
Limites, lemme \ref{limits-lemma-descend-finite-presentation},
on peut trouver un $i$ et un morphisme $f_i : X_i \to \Spec(A_i)$ de
présentation finie dont le changement de base à $Y$ redonne $f$.
Quitte à augmenter $i$, on peut supposer qu'il existe un
$\mathcal{O}_{X_i}$-module quasi-cohérent $\mathcal{F}_i$ de présentation
finie dont l'image inverse sur $X$ est isomorphe à $\mathcal{F}$, voir
Limites, lemme \ref{limits-lemma-descend-modules-finite-presentation}.
Quitte à augmenter encore une fois $i$, on peut supposer qu'il existe
un sous-schéma ouvert $U_i \subset X_i$ dont l'image inverse dans $X$
est $U$, voir
Limites, lemme \ref{limits-lemma-descend-opens}.
D'après le
lemme \ref{more-morphisms-lemma-base-change-assassin-in-U},
il suffit de démontrer le lemme pour $f_i$. On se ramène donc au
cas où $Y$ est le spectre d'un anneau noethérien.

\medskip\noindent
Nous allons utiliser le critère de
Topologie, lemme \ref{topology-lemma-characterize-constructible-Noetherian},
pour démontrer que $E$ est constructible lorsque $Y$ est un schéma noethérien.
Pour le voir, soit $Z \subset Y$ un sous-schéma fermé irréductible.
Il faut montrer que $E \cap Z$ contient un ouvert non vide
ou n'est pas dense dans $Z$. Cela résulte des
lemmes \ref{more-morphisms-lemma-bad-case} et
\ref{more-morphisms-lemma-good-case},
appliqués au changement de base $(X, \mathcal{F}, U) \times_Y Z$ sur $Z$.
\end{proof}

\section{Fibres réduites}
\label{more-morphisms-section-reduced}

\begin{lemma}
\label{more-morphisms-lemma-nonreduced-in-neighbourhood}
Soit $f : X \to Y$ un morphisme de schémas. Supposons que $Y$ soit irréductible, de
point générique $\eta$, et que $f$ soit de type fini. Si $X_\eta$ n'est pas réduit,
alors il existe un ouvert non vide $V \subset Y$
tel que, pour tout $y \in V$, la fibre $X_y$ ne soit pas réduite.
\end{lemma}

\begin{proof}
Soit $Y' \subset Y$ le réduit de $Y$. Soit $X' \to Y'$
le changement de base de $f$. Notons que $Y' \to Y$
induit une bijection sur les points et que $X' \to X$ identifie les fibres.
On peut donc supposer que $Y'$ est réduit, c'est-à-dire intègre, voir
Propriétés, lemme \ref{properties-lemma-characterize-integral}.
On peut aussi remplacer $Y$ par un ouvert affine. On peut donc supposer que
$Y = \Spec(A)$, où $A$ est un anneau intègre. Notons $K$ le
corps des fractions de $A$.
Choisissons un ouvert affine $\Spec(B) = U \subset X$ et une section
$h_\eta \in \Gamma(U_\eta, \mathcal{O}_{U_\eta}) = B_K$
qui soit non nulle et nilpotente.
Quitte à rétrécir $Y$, on peut supposer que $h$ provient de
$h \in \Gamma(U, \mathcal{O}_U) = B$. Quitte à rétrécir encore un peu $Y$,
on peut supposer que $h$ est nilpotent. Soit
$I = \{b \in B \mid hb = 0\}$ l'annulateur de $h$.
Alors $C = B/I$ est une $A$-algèbre de type fini dont la fibre générique
$(B/I)_K$ est non nulle (car $h_\eta \not = 0$). Appliquons
la platitude générique à $A \to C$ et $A \to B/hB$, voir
Algèbre, lemme \ref{algebra-lemma-generic-flatness} ;
on obtient un $g \in A$, $g \not = 0$ tel que $C_g$ soit libre comme
$A_g$-module et que $(B/hB)_g$ soit plat comme $A_g$-module.
Remplaçons $Y$ par $D(g) \subset Y$. On a maintenant la suite exacte courte
$$
0 \to C \to B \to B/hB \to 0.
$$
où $B/hB$ est plat sur $A$ et où $C$ est un $A$-module libre non nul.
Il en résulte que, pour tout homomorphisme $A \to \kappa$ vers un corps,
l'anneau $C \otimes_A \kappa$ est non nul et la suite
$$
0 \to C \otimes_A \kappa \to B \otimes_A \kappa \to
B/hB \otimes_A \kappa \to 0
$$
est exacte, voir
Algèbre, lemme \ref{algebra-lemma-flat-tor-zero}.
Notons que
$B/hB \otimes_A \kappa = (B \otimes_A \kappa) / h(B \otimes_A \kappa)$
par exacte à droite du produit tensoriel. On en conclut que
la multiplication par $h$ n'est pas nulle sur $B \otimes_A \kappa$.
Cela signifie clairement que, pour tout point $y \in Y$, l'élément $h$
se restreint en un élément non nul de $U_y$, d'où $X_y$ n'est pas réduit.
\end{proof}

\begin{lemma}
\label{more-morphisms-lemma-base-change-fibres-geometrically-reduced}
Soit $f : X \to Y$ un morphisme de schémas.
Soit $g : Y' \to Y$ un morphisme quelconque, et notons
$f' : X' \to Y'$ le changement de base de $f$.
Alors
\begin{align*}
\{y' \in Y' \mid X'_{y'}\text{ est géométriquement réduit}\} \\
= g^{-1}(\{y \in Y \mid X_y\text{ est géométriquement réduit}\}).
\end{align*}
\end{lemma}

\begin{proof}
Cela revient à dire que, pour $y' \in Y'$ d'image
$y \in Y$, la fibre $X'_{y'} = X_y \times_y y'$ est géométriquement
réduite sur $\kappa(y')$ si et seulement si $X_y$ est géométriquement
réduit sur $\kappa(y)$. Cela résulte de
Variétés, lemme \ref{varieties-lemma-geometrically-reduced-upstairs}.
\end{proof}

\begin{lemma}
\label{more-morphisms-lemma-not-geometrically-reduced-in-neighbourhood}
Soit $f : X \to Y$ un morphisme de schémas. Supposons que $Y$ soit irréductible, de
point générique $\eta$, et que $f$ soit de type fini. Si $X_\eta$ n'est pas
géométriquement réduit, alors il existe un ouvert non vide $V \subset Y$
tel que, pour tout $y \in V$, la fibre $X_y$ ne soit pas géométriquement réduite.
\end{lemma}

\begin{proof}
Appliquons le
lemme \ref{more-morphisms-lemma-make-generic-fibre-geometrically-reduced}
pour obtenir
$$
\xymatrix{
X' \ar[d]_{f'} \ar[r]_{g'} & X_V \ar[r] \ar[d] & X \ar[d]^f \\
Y' \ar[r]^g & V \ar[r] & Y
}
$$
avec toutes les propriétés mentionnées dans ce lemme.
Soit $\eta'$ le point générique de $Y'$.
Considérons le morphisme $X' \to X_{Y'}$ (qui est le morphisme de
réduction) et le morphisme de fibres génériques qui en résulte,
$X'_{\eta'} \to X_{\eta'}$.
Puisque $X'_{\eta'}$ est géométriquement réduit, tandis que $X_\eta$
ne l'est pas, ce morphisme ne peut être un isomorphisme, voir
Variétés, lemme \ref{varieties-lemma-geometrically-reduced-upstairs}.
Ainsi $X_{\eta'}$ n'est pas réduit. D'après le
lemme \ref{more-morphisms-lemma-nonreduced-in-neighbourhood},
les fibres de $X_{Y'} \to Y'$ ne sont donc pas réduites en tout point $y' \in V'$
d'un ouvert non vide $V' \subset Y'$. Puisque $g : Y' \to V$ est un homéomorphisme, le
lemme \ref{more-morphisms-lemma-base-change-fibres-geometrically-reduced}
montre que $g(V')$ est l'ouvert recherché.
\end{proof}

\begin{lemma}
\label{more-morphisms-lemma-geometrically-reduced-generic-fibre}
Soit $f : X \to Y$ un morphisme de schémas.
Supposons que
\begin{enumerate}
\item $Y$ soit irréductible, de point générique $\eta$,
\item $X_\eta$ soit géométriquement réduit, et
\item $f$ soit de type fini.
\end{enumerate}
Alors il existe un sous-schéma ouvert non vide $V \subset Y$
tel que $X_V \to V$ ait des fibres géométriquement réduites.
\end{lemma}

\begin{proof}
Soit $Y' \subset Y$ le réduit de $Y$. Soit $X' \to Y'$
le changement de base de $f$. Notons que $Y' \to Y$
induit une bijection sur les points et que $X' \to X$ identifie les fibres.
On peut donc supposer que $Y'$ est réduit, c'est-à-dire intègre, voir
Propriétés, lemme \ref{properties-lemma-characterize-integral}.
On peut aussi remplacer $Y$ par un ouvert affine. On peut donc supposer que
$Y = \Spec(A)$, où $A$ est un anneau intègre. Notons $K$ le
corps des fractions de $A$. Quitte à rétrécir un peu $Y$, on peut également supposer que
$X \to Y$ est plat et de présentation finie, voir
Morphismes, proposition \ref{morphisms-proposition-generic-flatness}.

\medskip\noindent
Puisque $X_\eta$ est géométriquement réduit, il existe un sous-ensemble ouvert dense
$V \subset X_\eta$ tel que $V \to \Spec(K)$ soit lisse, voir
Variétés, lemme \ref{varieties-lemma-geometrically-reduced-dense-smooth-open}.
Soit $U \subset X$ l'ensemble des points où $f$ est lisse. D'après
Morphismes, lemme \ref{morphisms-lemma-set-points-where-fibres-smooth},
on a $V \subset U_\eta$. Ainsi, la fibre générique de $U$ est dense
dans la fibre générique de $X$. Puisque $X_\eta$ est réduit, il en résulte
que $U_\eta$ est schématiquement dense dans $X_\eta$, voir
Morphismes, lemme \ref{morphisms-lemma-reduced-scheme-theoretically-dense}.
Notons que, puisque $U \to Y$ est lisse, toutes les fibres de $U \to Y$
sont géométriquement réduites. Il suffit donc de montrer que, après avoir
rétréci $Y$, pour tout $y \in Y$, le schéma $U_y$ est schématiquement
dense dans $X_y$, voir
Morphismes, lemme \ref{morphisms-lemma-reduced-subscheme-closure}.
Cela résulte du
lemme \ref{more-morphisms-lemma-scheme-theoretically-dense-generic-fibre}.
\end{proof}

\begin{lemma}
\label{more-morphisms-lemma-geometrically-reduced-constructible}
Soit $f : X \to Y$ un morphisme quasi-compact et
localement de présentation finie. Alors l'ensemble
$$
E = \{y \in Y \mid X_y\text{ est géométriquement réduit}\}
$$
est localement constructible dans $Y$.
\end{lemma}

\begin{proof}
Soit $y \in Y$. Nous devons montrer qu'il existe un voisinage ouvert
$V$ de $y$ dans $Y$ tel que $E \cap V$ soit constructible dans $V$. On peut donc
supposer que $Y$ est affine. Alors $X$ est quasi-compact.
Choisissons un recouvrement ouvert affine fini $X = U_1 \cup \ldots \cup U_n$.
Les fibres des $U_i \to Y$ en $y$ forment alors un recouvrement ouvert affine
de la fibre de $X \to Y$ en $y$. On peut donc aussi supposer que $X$ est affine.
Écrivons $Y = \Spec(A)$.
Écrivons $A = \colim A_i$ comme limite inductive d'algèbres
de type fini sur $\mathbf{Z}$. D'après
Limites, lemme \ref{limits-lemma-descend-finite-presentation},
on peut trouver un $i$ et un morphisme $f_i : X_i \to \Spec(A_i)$ de
présentation finie dont le changement de base à $Y$ redonne $f$. D'après le
lemme \ref{more-morphisms-lemma-base-change-fibres-geometrically-reduced},
il suffit de démontrer le lemme pour $f_i$. On se ramène ainsi au
cas où $Y$ est le spectre d'un anneau noethérien.

\medskip\noindent
Nous utiliserons le critère du
Topologie, lemme \ref{topology-lemma-characterize-constructible-Noetherian},
pour montrer que $E$ est constructible lorsque $Y$ est un schéma noethérien.
Pour cela, soit $Z \subset Y$ un sous-schéma fermé irréductible
de point générique $\xi$.
Nous devons montrer que $E \cap Z$ contient un ouvert non vide
ou n'est pas dense dans $Z$. Si $X_\xi$ est géométriquement réduit, alors le
lemme \ref{more-morphisms-lemma-geometrically-reduced-generic-fibre}
(appliqué au morphisme $X_Z \to Z$)
implique que toutes les fibres $X_y$ sont géométriquement réduites
pour un ouvert non vide $V \subset Z$.
Si $X_\xi$ n'est pas géométriquement réduit, alors le
lemme \ref{more-morphisms-lemma-not-geometrically-reduced-in-neighbourhood}
(appliqué au morphisme $X_Z \to Z$)
implique que toutes les fibres $X_y$ ne sont pas géométriquement réduites
pour un ouvert non vide $V \subset Z$. Cela conclut.
\end{proof}

\begin{lemma}
\label{more-morphisms-lemma-proper-flat-over-dvr-reduced-fibre}
Soit $f : X \to \Spec(R)$ un morphisme propre, où $R$ est un
anneau de valuation discrète. Supposons que toute composante irréductible
de $X$ domine $\Spec(R)$ (par exemple si $f$ est plat) et
que la fibre spéciale soit réduite. Alors
$X$ et la fibre générique $X_\eta$ sont réduits.
\end{lemma}

\begin{proof}
(Le fait qu'un morphisme plat satisfasse l'hypothèse sur les composantes
irréductibles résulte par exemple du
Diviseurs, lemme \ref{divisors-lemma-bourbaki},
ou de la descente des idéaux premiers pour les homomorphismes d'anneaux plats.)
Supposons que la fibre spéciale $X_s$ soit réduite.
Soit $x \in X$ un point quelconque, et montrons que $\mathcal{O}_{X, x}$
est réduit ; cela montrera que $X$ et $X_\eta$ sont réduits.
Soit $x \leadsto x'$ une spécialisation telle que $x'$
appartienne à la fibre spéciale ; une telle spécialisation existe
puisqu'un morphisme propre est fermé. Considérons l'anneau local
$A = \mathcal{O}_{X, x'}$. Alors $\mathcal{O}_{X, x}$
est un localisé de $A$, de sorte qu'il suffit de montrer que
$A$ est réduit. Soit $\pi \in R$ une uniformisante.
Si $a \in A$ est non nul, il existe un $n \geq 0$ et un élément
$a' \in A$ tels que $a = \pi^n a'$ et $a' \not \in \pi A$.
Cela résulte du théorème d'intersection de Krull
(Algèbre, lemme \ref{algebra-lemma-intersect-powers-ideal-module-zero}).
Si $a$ est nilpotent, alors $a$ appartient à tous les idéaux premiers minimaux de $A$,
et il en va donc de même de $a'$, car $\pi$ n'appartient à aucun
des idéaux premiers minimaux de $A$ en vertu de notre hypothèse sur les composantes
irréductibles de $X$. Ainsi $a'$ est nilpotent.
Mais $a'$ s'envoie sur un élément non nul de l'anneau réduit
$A/\pi A = \mathcal{O}_{X_s, x'}$.
C'est une contradiction à moins que $A$ ne soit réduit, ce qui
est ce que nous voulions montrer.
\end{proof}

\begin{lemma}
\label{more-morphisms-lemma-geometrically-reduced-open}
Soit $f : X \to Y$ un morphisme plat et propre de présentation finie.
Alors l'ensemble $\{y \in Y \mid X_y\text{ est géométriquement réduit}\}$
est ouvert dans $Y$.
\end{lemma}

\begin{proof}
On peut supposer que $Y$ est affine. Alors $Y$ est une limite projective filtrante
de schémas affines de type fini sur $\mathbf{Z}$.
On peut donc supposer que $X \to Y$ est le
changement de base de $X_0 \to Y_0$, où $Y_0$ est le spectre d'une algèbre de type
fini sur $\mathbf{Z}$ et $X_0 \to Y_0$ est plat et propre.
Voir Limites, lemme \ref{limits-lemma-descend-finite-presentation},
\ref{limits-lemma-descend-flat-finite-presentation}, et
\ref{limits-lemma-eventually-proper}. Puisque la formation de
l'ensemble des points où les fibres sont géométriquement
réduites commute au changement de base
(lemme \ref{more-morphisms-lemma-base-change-fibres-geometrically-reduced}),
on peut supposer que la base est noethérienne.

\medskip\noindent
Supposons que $Y$ soit noethérien. L'ensemble est constructible d'après le
lemme \ref{more-morphisms-lemma-geometrically-reduced-constructible}.
Il suffit donc de montrer que l'ensemble est stable par générisation
(Topologie, lemme \ref{topology-lemma-characterize-closed-Noetherian}). D'après
Propriétés, lemme \ref{properties-lemma-locally-Noetherian-specialization-dvr},
on se ramène au cas où $Y = \Spec(R)$, où $R$ est un anneau de
valuation discrète, et où la fibre fermée $X_y$ est géométriquement
réduite. Il faut montrer que la fibre générique $X_\eta$ est géométriquement réduite.

\medskip\noindent
Sinon, il existe une extension finie $L$ du corps des fractions
de $R$ telle que $X_L$ ne soit pas réduit, voir
Variétés, lemme \ref{varieties-lemma-geometrically-reduced}.
Il existe un anneau de valuation discrète
$R' \subset L$ de corps des fractions $L$ et dominant $R$, voir
Algèbre, lemme \ref{algebra-lemma-integral-closure-Dedekind}.
Après avoir remplacé $R$ par $R'$, on se ramène au
lemme \ref{more-morphisms-lemma-proper-flat-over-dvr-reduced-fibre}.
\end{proof}

\section{Composantes irréductibles des fibres}
\label{more-morphisms-section-irreducible}

\begin{lemma}
\label{more-morphisms-lemma-irreducible-components-in-neighbourhood}
Soit $f : X \to Y$ un morphisme de schémas. Supposons $Y$ irréductible de
point générique $\eta$ et $f$ de type fini. Si $X_\eta$ a $n$
composantes irréductibles, alors il existe un ouvert non vide $V \subset Y$
tel que, pour tout $y \in V$, la fibre $X_y$ ait au moins $n$
composantes irréductibles.
\end{lemma}

\begin{proof}
La question étant purement topologique, on peut remplacer $X$ et $Y$ par
leurs réduits. En particulier, cela implique que $Y$ est intègre, voir
Propriétés, lemme \ref{properties-lemma-characterize-integral}.
Soit $X_\eta = X_{1, \eta} \cup \ldots \cup X_{n, \eta}$
la décomposition de $X_\eta$ en composantes irréductibles.
Soit $X_i \subset X$ le sous-schéma fermé réduit dont la fibre
générique est $X_{i, \eta}$. Notons que $Z_{i, j} = X_i \cap X_j$
est un sous-ensemble fermé de $X_i$ dont la fibre générique $Z_{i, j, \eta}$
est nulle part dense dans $X_{i, \eta}$. Après avoir rétréci $Y$, on peut donc
supposer que $Z_{i, j, y}$
est nulle part dense dans $X_{i, y}$ pour tout $y \in Y$, voir
lemme \ref{more-morphisms-lemma-nowhere-dense-generic-fibre}.
Après avoir encore rétréci $Y$, on peut supposer que
$X_y = \bigcup X_{i, y}$ pour $y \in Y$, voir
lemme \ref{more-morphisms-lemma-cover-generic-fibre-neighbourhood}.
De plus, après avoir rétréci $Y$, on peut supposer que chaque $X_i \to Y$
est plat et de présentation finie, voir
Morphismes, proposition \ref{morphisms-proposition-generic-flatness}.
Les morphismes $X_i \to Y$ sont ouverts, voir
Morphismes, lemme \ref{morphisms-lemma-fppf-open}.
Il existe donc un voisinage ouvert $V$ de $\eta$ qui est contenu
dans $f(X_i)$ pour chaque $i$.
Pour tout $y \in V$, les schémas $X_{i, y}$ sont des
sous-ensembles fermés non vides de $X_y$, on a $X_y = \bigcup X_{i, y}$,
et les intersections $Z_{i, j, y} = X_{i, y} \cap X_{j, y}$
ne sont pas denses dans $X_{i, y}$. Cela implique clairement que
$X_y$ a au moins $n$ composantes irréductibles.
\end{proof}

\begin{lemma}
\label{more-morphisms-lemma-base-change-fibres-geometrically-irreducible}
Soit $f : X \to Y$ un morphisme de schémas.
Soit $g : Y' \to Y$ un morphisme quelconque, et notons
$f' : X' \to Y'$ le changement de base de $f$.
Alors
\begin{align*}
\{y' \in Y' \mid X'_{y'}\text{ est géométriquement irréductible}\} \\
= g^{-1}(\{y \in Y \mid X_y\text{ est géométriquement irréductible}\}).
\end{align*}
\end{lemma}

\begin{proof}
Cela se ramène à l'énoncé suivant : pour $y' \in Y'$ d'image
$y \in Y$, la fibre $X'_{y'} = X_y \times_y y'$ est géométriquement
irréductible sur $\kappa(y')$ si et seulement si $X_y$ est géométriquement
irréductible sur $\kappa(y)$. Cela résulte de
Variétés,
lemme \ref{varieties-lemma-geometrically-irreducible-check-after-extension}.
\end{proof}

\begin{lemma}
\label{more-morphisms-lemma-base-change-fibres-nr-geometrically-irreducible-components}
Soit $f : X \to Y$ un morphisme de schémas. Soit
$$
n_{X/Y} : Y \to \{0, 1, 2, 3, \ldots, \infty\}
$$
la fonction qui associe à $y \in Y$ le nombre de composantes irréductibles
de $(X_y)_K$, où $K$ est une extension séparablement close
de $\kappa(y)$. Cette fonction est bien définie et, si $g : Y' \to Y$ est un morphisme,
alors
$$
n_{X'/Y'} = n_{X/Y} \circ g
$$
où $X' \to Y'$ est le changement de base de $f$.
\end{lemma}

\begin{proof}
Supposons que $y' \in Y'$ ait pour image $y \in Y$.
Supposons que $K \supset \kappa(y)$ et $K' \supset \kappa(y')$ soient des extensions
séparablement closes. On peut alors choisir un diagramme commutatif
$$
\xymatrix{
K \ar[r] & K'' & K' \ar[l] \\
\kappa(y) \ar[u] \ar[rr] & & \kappa(y') \ar[u]
}
$$
de corps. Le résultat s'ensuit, car les morphismes de schémas
$$
\xymatrix{
(X'_{y'})_{K'} &
(X'_{y'})_{K''} = (X_y)_{K''} \ar[l] \ar[r] &
(X_y)_K
}
$$
induisent des bijections entre les composantes irréductibles, voir
Variétés,
lemme \ref{varieties-lemma-separably-closed-field-irreducible-components}.
\end{proof}

\begin{lemma}
\label{more-morphisms-lemma-irreducible-polynomial-over-domain}
Soit $A$ un anneau intègre de corps des fractions $K$.
Soit $P \in A[x_1, \ldots, x_n]$.
Notons $\overline{K}$ la clôture algébrique de $K$.
Supposons que $P$ soit irréductible dans $\overline{K}[x_1, \ldots, x_n]$.
Alors il existe un $f \in A$ tel que
$P^\varphi \in \kappa[x_1, \ldots, x_n]$ soit irréductible pour tout
homomorphisme $\varphi : A_f \to \kappa$ vers un corps.
\end{lemma}

\begin{proof}
Il existe un automorphisme $\Psi$ de $A[x_1, \ldots, x_n]$ sur $A$
tel que $\Psi(P) = ax_n^d +$ des termes de degré inférieur en $x_n$, avec
$a \not = 0$, voir
Algèbre, lemme \ref{algebra-lemma-helper-polynomial}.
On peut remplacer $P$ par $\Psi(P)$ et $A$ par $A_a$.
On peut ainsi supposer que $P$ est unitaire en $x_n$ de degré $d > 0$.
Pour $i = 1, \ldots, n - 1$, soit $d_i$ le degré de $P$ en $x_i$.
Notons que cela implique que $P^\varphi$ est unitaire de degré $d$ en $x_n$
et de degré $\leq d_i$ en $x_i$ pour tout homomorphisme
$\varphi : A \to \kappa$, où $\kappa$ est un corps.
Ainsi, si $P^\varphi$ est réductible, on peut écrire
$$
P^\varphi = Q_1 Q_2
$$
avec $Q_1, Q_2$ unitaires de degrés $e_1, e_2 \geq 0$ en $x_n$, où
$e_1 + e_2 = d$, et de degré $\leq d_i$ en $x_i$ pour
$i = 1, \ldots, n - 1$. Autrement dit, on peut écrire
\begin{equation}
\label{more-morphisms-equation-factors}
Q_j = x_n^{e_j} + \sum\nolimits_{0 \leq l < e_j}
\left( \sum\nolimits_{L \in \mathcal{L}} a_{j, l, L} x^L \right) x_n^l
\end{equation}
où la somme porte sur l'ensemble $\mathcal{L}$ des multi-indices $L$
de la forme $L = (l_1, \ldots, l_{n - 1})$ avec $0 \leq l_i \leq d_i$.
Pour tous $e_1, e_2 \geq 0$ tels que $e_1 + e_2 = d$, considérons la $A$-algèbre
$$
B_{e_1, e_2} =
A[\{a_{1, l, L}\}_{0 \leq l < e_1, L \in \mathcal{L}},
\{a_{2, l, L}\}_{0 \leq l < e_2, L \in \mathcal{L}}]/(\text{relations})
$$
où $(\text{relations})$ est l'idéal engendré par les coefficients
du polynôme
$$
P - Q_1Q_2 \in
A[\{a_{1, l, L}\}_{0 \leq l < e_1, L \in \mathcal{L}},
\{a_{2, l, L}\}_{0 \leq l < e_2, L \in \mathcal{L}}][x_1, \ldots, x_n]
$$
où $Q_1$ et $Q_2$ sont définis comme dans (\ref{more-morphisms-equation-factors}).
L'hypothèse que $P$ est irréductible sur $\overline{K}$ implique qu'il
n'existe aucun homomorphisme de $A$-algèbres
$B_{e_1, e_2} \to \overline{K}$. D'après le Nullstellensatz de Hilbert, voir
Algèbre, théorème \ref{algebra-theorem-nullstellensatz},
cela signifie que $B_{e_1, e_2} \otimes_A K = 0$.
Comme $B_{e_1, e_2}$ est une $A$-algèbre de type fini, cela signifie que
l'on peut trouver un $f_{e_1, e_2} \in A$ tel que
$(B_{e_1, e_2})_{f_{e_1, e_2}} = 0$. Par construction, cela signifie que,
si $\varphi : A_{f_{e_1, e_2}} \to \kappa$ est un homomorphisme vers un corps,
alors $P^\varphi$ n'a pas de factorisation $P^\varphi = Q_1 Q_2$
avec $Q_1$ de degré $e_1$ en $x_n$ et $Q_2$ de degré $e_2$ en $x_n$.
Ainsi, en prenant
$f = \prod_{e_1, e_2 \geq 0, e_1 + e_2 = d} f_{e_1, e_2}$, on conclut.
\end{proof}

\begin{lemma}
\label{more-morphisms-lemma-geom-irreducible-generic-fibre}
Soit $f : X \to Y$ un morphisme de schémas.
Supposons que
\begin{enumerate}
\item $Y$ soit irréductible, de point générique $\eta$,
\item $X_\eta$ soit géométriquement irréductible, et
\item $f$ soit de type fini.
\end{enumerate}
Alors il existe un sous-schéma ouvert non vide $V \subset Y$
tel que $X_V \to V$ ait des fibres géométriquement irréductibles.
\end{lemma}

\begin{proof}[Première démonstration du lemme \ref{more-morphisms-lemma-geom-irreducible-generic-fibre}]
Nous donnons deux démonstrations du lemme. Elles sont essentiellement équivalentes ;
la seconde est plus autonome, mais un peu plus longue.
Choisissons un diagramme
$$
\xymatrix{
X' \ar[d]_{f'} \ar[r]_{g'} & X_V \ar[r] \ar[d] & X \ar[d]^f \\
Y' \ar[r]^g & V \ar[r] & Y
}
$$
comme dans le
lemme \ref{more-morphisms-lemma-make-generic-fibre-geometrically-reduced}.
Notons que la fibre générique de $f'$ est le réduit de la
fibre générique de $f$ (voir le
lemme \ref{more-morphisms-lemma-reduction-generic-fibre})
et est donc géométriquement irréductible.
Supposons que le lemme soit vrai pour le morphisme $f'$. Après avoir rétréci
$V$, toutes les fibres de $f'$ sont alors géométriquement irréductibles.
Comme $X' = (Y' \times_V X_V)_{red}$, cela implique que toutes les fibres
de $Y' \times_V X_V$ sont géométriquement irréductibles. D'après le
lemme \ref{more-morphisms-lemma-base-change-fibres-geometrically-irreducible},
toutes les fibres de $X_V \to V$ sont donc géométriquement irréductibles et
on conclut. De cette manière, on voit que l'on peut supposer que la fibre générique
est à la fois géométriquement réduite et géométriquement irréductible,
et l'on peut supposer $Y = \Spec(A)$ avec $A$ un anneau intègre.

\medskip\noindent
Soit $x \in X_\eta$ le point générique. Puisque $X_\eta$ est géométriquement
irréductible et réduit, on voit que $L = \kappa(x)$ est une extension de type fini
de $K = \kappa(\eta)$ qui est géométriquement réduite et
géométriquement irréductible, voir
Variétés, lemmes \ref{varieties-lemma-geometrically-reduced-at-point} et
\ref{varieties-lemma-geometrically-irreducible-function-field}.
En particulier, l'extension de corps $L/K$ est séparable, voir
Algèbre, lemme \ref{algebra-lemma-characterize-separable-field-extensions}.
On peut donc trouver $x_1, \ldots, x_{r + 1} \in L$ qui engendrent $L$
sur $K$ et tels que $x_1, \ldots, x_r$ soient une base de transcendance de
$L$ sur $K$, voir
Algèbre, lemme
\ref{algebra-lemma-generating-finitely-generated-separable-field-extensions}.
Soit $P \in K(x_1, \ldots, x_r)[T]$ le polynôme minimal de
$x_{r + 1}$. En chassant les dénominateurs, on peut supposer que
$P$ a ses coefficients dans $A[x_1, \ldots, x_r]$.
Notons que, puisque $L$ est géométriquement réduit et géométriquement irréductible
sur $K$, le polynôme $P$ est irréductible dans
$\overline{K}[x_1, \ldots, x_r, T]$, où $\overline{K}$ est la
clôture algébrique de $K$. Notons
$$
B' = A[x_1, \ldots, x_{r + 1}]/(P(x_{r + 1}))
$$
et posons $X' = \Spec(B')$. Par construction, le corps des fractions de $B'$
est isomorphe à $L = \kappa(x)$ comme extension de $K$. Il existe donc un
ouvert $U \subset X$, un ouvert $U' \subset X'$ et un $Y$-isomorphisme
$U \to U'$, voir
Morphismes, lemme \ref{morphisms-lemma-common-open}.
Voici un diagramme :
$$
\xymatrix{
X \ar[rd] &
U \ar[l] \ar@{=}[r] \ar[d] &
U' \ar[r] \ar[d] &
X' \ar[ld] \ar@{=}[r] & \Spec(B') \\
& Y \ar@{=}[r] & Y &
}
$$
Notons que $U_\eta \subset X_\eta$ et $U'_\eta \subset X'_\eta$ sont des
ouverts denses. Ainsi, après avoir rétréci $Y$ en appliquant le
lemme \ref{more-morphisms-lemma-nowhere-dense-generic-fibre},
on obtient que $U_y$ est dense dans $X_y$ et que $U'_y$ est dense dans $X'_y$
pour tout $y \in Y$. Il suffit donc de démontrer le lemme pour
$X' \to Y$, ce qui est le contenu du
lemme \ref{more-morphisms-lemma-irreducible-polynomial-over-domain}.
\end{proof}

\begin{proof}[Seconde démonstration du lemme \ref{more-morphisms-lemma-geom-irreducible-generic-fibre}]
Soit $Y' \subset Y$ le réduit de $Y$. Soit $X' \to X$ le morphisme de réduction
de $X$. Notons que $X' \to X  \to Y$ se factorise par $Y'$, voir
Schémas, lemme \ref{schemes-lemma-map-into-reduction}.
Comme $Y' \to Y$ et $X' \to X$ sont des homéomorphismes
universels d'après
Morphismes, lemme \ref{morphisms-lemma-reduction-universal-homeomorphism},
on voit qu'il suffit de démontrer le lemme pour $X' \to Y'$. On peut donc
supposer que $X$ et $Y$ sont réduits. En particulier, $Y$ est intègre, voir
Propriétés, lemme \ref{properties-lemma-characterize-integral}.
Ainsi, d'après
Morphismes, proposition \ref{morphisms-proposition-generic-flatness},
il existe un ouvert affine non vide $V \subset Y$ tel que $X_V \to V$ soit
plat et de présentation finie. Après avoir remplacé $Y$ par $V$, on peut
supposer, outre (1), (2), (3), que $Y$ est affine intègre, que $X$
est réduit et que $f$ est plat et de présentation finie. En particulier,
$f$ est universellement ouvert, voir
Morphismes, lemme \ref{morphisms-lemma-fppf-open}.

\medskip\noindent
Choisissons un ouvert affine non vide $U \subset X$. Alors $U \to Y$ est plat et de
présentation finie, à fibre générique géométriquement irréductible.
Le complémentaire $X_\eta \setminus U_\eta$ est nulle part dense. Ainsi, après
avoir rétréci $Y$, on peut supposer $U_y \subset X_y$ ouvert dense pour tout
$y \in Y$, voir le
lemme \ref{more-morphisms-lemma-nowhere-dense-generic-fibre}.
On peut donc remplacer $X$ par $U$ et on se ramène au
cas où $Y$ est affine intègre et où $X$ est affine réduit, plat et de présentation
finie sur $Y$, à fibre générique $X_\eta$ géométriquement irréductible.

\medskip\noindent
Écrivons $X = \Spec(B)$ et $Y = \Spec(A)$. Alors $A$ est un anneau intègre,
$B$ est réduit, $A \to B$ est plat et de présentation finie, et $B_K$ est
géométriquement irréductible sur le corps des fractions $K$ de $A$.
En particulier, on voit que $B_K$ est un anneau intègre. Soit $L$ le corps des fractions
de $B_K$. Notons que
$L$ est une extension de corps de type fini de $K$, puisque $B$ est une $A$-algèbre
de présentation finie. Soit $K'/K$ une extension finie purement inséparable
telle que $(L \otimes_K K')_{red}$ soit une extension de corps séparablement engendrée,
voir
Algèbre, lemme \ref{algebra-lemma-make-separable}.
Choisissons $x_1, \ldots, x_n \in K'$ qui engendrent l'extension de corps
$K'$ sur $K$, et tels que $x_i^{q_i} \in A$ pour une certaine puissance d'un nombre premier
$q_i$ (démonstration de l'existence des $x_i$ omise). Soit $A'$ la sous-$A$-algèbre
de $K'$ engendrée par $x_1, \ldots, x_n$. Alors $A'$ est une
sous-$A$-algèbre finie $A' \subset K'$ dont le corps des fractions est $K'$. Notons que
$\Spec(A') \to \Spec(A)$ est un homéomorphisme universel, voir
Algèbre, lemme \ref{algebra-lemma-p-ring-map}.
Il suffit donc de démontrer le résultat après changement de base à $\Spec(A')$.
Nous allons remplacer $A$ par $A'$ et $B$ par $(B \otimes_A A')_{red}$
pour arriver à la situation où $L$ est une extension de corps séparablement engendrée
de $K$. Bien sûr, il peut arriver que $(B \otimes_A A')_{red}$ ne soit plus
plat ou de présentation finie sur $A'$, mais on peut y remédier en
remplaçant $A'$ par $A'_f$ pour un $f \in A'$ convenable, voir
Algèbre, lemme \ref{algebra-lemma-generic-flatness}.

\medskip\noindent
À ce stade, on sait que $A$ est un anneau intègre, que $B$ est réduit, que $A \to B$
est plat et de présentation finie, et que $B_K$ est un anneau intègre dont le corps des fractions
$L$ est une extension de corps séparablement engendrée du corps des fractions $K$
de $A$. D'après Algèbre, lemme
\ref{algebra-lemma-generating-finitely-generated-separable-field-extensions},
on peut écrire $L = K(x_1, \ldots, x_{r + 1})$
où $x_1, \ldots, x_r$ sont algébriquement indépendants sur $K$, et où
$x_{r + 1}$ est séparable sur $K(x_1, \ldots, x_r)$. Après avoir chassé
les dénominateurs, on peut supposer que le polynôme minimal
$P \in K(x_1, \ldots, x_r)[T]$ de $x_{r + 1}$ sur $K(x_1, \ldots, x_r)$
a ses coefficients dans $A[x_1, \ldots, x_r]$. Notons que, puisque
$L/K$ est séparable et que $L$ est géométriquement irréductible sur
$K$, le polynôme $P$ est irréductible sur la clôture algébrique
$\overline{K}$ de $K$. Notons
$$
B' = A[x_1, \ldots, x_{r + 1}]/(P(x_{r + 1})).
$$
Par construction, les corps des fractions de $B$ et $B'$ sont isomorphes comme
extensions de $K$. Il existe donc un isomorphisme de $A$-algèbres
$B_h \cong B'_{h'}$ pour des $h \in B$ et $h' \in B'$ convenables, voir
Morphismes, lemme \ref{morphisms-lemma-common-open}.
Autrement dit, $X$ et $X' = \Spec(B')$ ont un ouvert affine commun $U$.
Voici un diagramme :
$$
\xymatrix{
X = \Spec(B) \ar[rd] &
U \ar[l] \ar[r] \ar[d] &
\Spec(B') = X' \ar[ld] \\
& Y = \Spec(A) &
}
$$
Après avoir encore rétréci $Y$ (en appliquant le
lemme \ref{more-morphisms-lemma-nowhere-dense-generic-fibre}
à $Z = X \setminus U$ dans $X$ et à $Z' = X' \setminus U$ dans $X'$),
on voit que $U_y$ est dense dans $X_y$ et que $U_y$ est dense dans $X'_y$
pour tout $y \in Y$. Il suffit donc de démontrer le lemme pour
$X' \to Y$, ce qui est le contenu du
lemme \ref{more-morphisms-lemma-irreducible-polynomial-over-domain}.
\end{proof}

\begin{lemma}
\label{more-morphisms-lemma-nr-geom-irreducible-components-good}
Soit $f : X \to Y$ un morphisme de schémas. Soit
$n_{X/Y}$ la fonction sur $Y$ qui compte le nombre de composantes
géométriquement irréductibles des fibres de $f$, introduite dans le
lemme \ref{more-morphisms-lemma-base-change-fibres-nr-geometrically-irreducible-components}.
Supposons $f$ de type fini.
Soit $y \in Y$ un point. Alors il existe un ouvert non vide
$V \subset \overline{\{y\}}$ tel que $n_{X/Y}|_V$ soit constant.
\end{lemma}

\begin{proof}
Soit $Z$ le schéma muni de la structure induite réduite sur $\overline{\{y\}}$.
Soit $f_Z : X_Z \to Z$ le changement de base de $f$. Il suffit clairement de démontrer
le lemme pour $f_Z$ et le point générique de $Z$. On peut donc supposer que
$Y$ est un schéma intègre, voir
Propriétés, lemme \ref{properties-lemma-characterize-integral}.
Notre but, dans ce cas, est de produire un ouvert non vide $V \subset Y$ tel que
$n_{X/Y}|_V$ soit constant.

\medskip\noindent
Appliquons le
lemme \ref{more-morphisms-lemma-make-components-generic-fibre-geometrically-irreducible}
à $f : X \to Y$ ; on obtient $g : Y' \to V \subset Y$. Comme $g : Y' \to V$ est
fini étale surjectif, donc en particulier ouvert (voir
Morphismes, lemme \ref{morphisms-lemma-etale-open}),
il suffit de démontrer qu'il existe un ouvert $V' \subset Y'$
tel que $n_{X'/Y'}|_{V'}$ soit constant, voir
lemme \ref{more-morphisms-lemma-base-change-fibres-nr-geometrically-irreducible-components}.
On voit ainsi que l'on peut supposer que toutes les composantes irréductibles de
la fibre générique $X_\eta$ sont géométriquement irréductibles sur $\kappa(\eta)$.

\medskip\noindent
À ce stade, supposons que
$X_\eta = X_{1, \eta} \bigcup \ldots \bigcup X_{n, \eta}$
soit la décomposition de la fibre générique en composantes
(géométriquement) irréductibles.
En particulier, $n_{X/Y}(\eta) = n$.
Soit $X_i$ l'adhérence de
$X_{i, \eta}$ dans $X$. Après avoir rétréci $Y$, on peut supposer que
$X = \bigcup X_i$, voir
lemme \ref{more-morphisms-lemma-cover-generic-fibre-neighbourhood}.
Après avoir encore rétréci $Y$, on voit que chaque fibre de
$f$ a au moins $n$ composantes irréductibles, voir
lemme \ref{more-morphisms-lemma-irreducible-components-in-neighbourhood}.
Ainsi $n_{X/Y}(y) \geq n$ pour tout $y \in Y$.
Après avoir encore rétréci $Y$, on obtient que $X_{i, y}$
est géométriquement irréductible pour chaque $i$ et tout $y \in Y$, voir
lemme \ref{more-morphisms-lemma-geom-irreducible-generic-fibre}.
Comme $X_y = \bigcup X_{i, y}$,
cela montre que $n_{X/Y}(y) \leq n$ et achève la démonstration.
\end{proof}

\begin{lemma}
\label{more-morphisms-lemma-nr-geom-irreducible-components-constructible}
Soit $f : X \to Y$ un morphisme de schémas. Soit
$n_{X/Y}$ la fonction sur $Y$ qui compte le nombre de composantes
géométriquement irréductibles des fibres de $f$, introduite dans le
lemme \ref{more-morphisms-lemma-base-change-fibres-nr-geometrically-irreducible-components}.
Supposons $f$ de présentation finie. Alors les ensembles de niveau
$$
E_n = \{y \in Y \mid n_{X/Y}(y) = n\}
$$
de $n_{X/Y}$ sont localement constructibles dans $Y$.
\end{lemma}

\begin{proof}
Fixons $n$. Soit $y \in Y$. Nous devons montrer qu'il existe un voisinage ouvert
$V$ de $y$ dans $Y$ tel que $E_n \cap V$ soit constructible dans $V$. On peut donc
supposer que $Y$ est affine. Écrivons $Y = \Spec(A)$ et
$A = \colim A_i$ comme limite inductive d'algèbres
de type fini sur $\mathbf{Z}$. D'après
Limites, lemme \ref{limits-lemma-descend-finite-presentation},
on peut trouver un $i$ et un morphisme $f_i : X_i \to \Spec(A_i)$ de
présentation finie dont le changement de base à $Y$ redonne $f$. D'après le
lemme \ref{more-morphisms-lemma-base-change-fibres-nr-geometrically-irreducible-components},
il suffit de démontrer le lemme pour $f_i$. On se ramène ainsi au
cas où $Y$ est le spectre d'un anneau noethérien.

\medskip\noindent
Nous utiliserons le critère du
Topologie, lemme \ref{topology-lemma-characterize-constructible-Noetherian},
pour montrer que $E_n$ est constructible lorsque $Y$ est un schéma noethérien.
Pour cela, soit $Z \subset Y$ un sous-schéma fermé irréductible.
Nous devons montrer que $E_n \cap Z$ contient un ouvert non vide
ou n'est pas dense dans $Z$. Soit $\xi \in Z$ le point générique. Alors le
lemme \ref{more-morphisms-lemma-nr-geom-irreducible-components-good}
montre que $n_{X/Y}$ est constant dans un voisinage de $\xi$ dans $Z$.
Cela implique clairement le résultat voulu.
\end{proof}

\section{Composantes connexes des fibres}
\label{more-morphisms-section-connected}

\begin{lemma}
\label{more-morphisms-lemma-connected-components-in-neighbourhood}
Soit $f : X \to Y$ un morphisme de schémas. Supposons $Y$ irréductible de
point générique $\eta$ et $f$ de type fini. Si $X_\eta$ a $n$
composantes connexes, alors il existe un ouvert non vide $V \subset Y$
tel que, pour tout $y \in V$, la fibre $X_y$ ait au moins $n$
composantes connexes.
\end{lemma}

\begin{proof}
La question étant purement topologique, on peut remplacer $X$ et $Y$ par
leurs réduits. En particulier, cela implique que $Y$ est intègre, voir
Propriétés, lemme \ref{properties-lemma-characterize-integral}.
Soit $X_\eta = X_{1, \eta} \cup \ldots \cup X_{n, \eta}$
la décomposition de $X_\eta$ en composantes connexes.
Soit $X_i \subset X$ le sous-schéma fermé réduit dont la fibre
générique est $X_{i, \eta}$. Notons que $Z_{i, j} = X_i \cap X_j$
est un sous-ensemble fermé de $X$ dont la fibre générique $Z_{i, j, \eta}$ est vide.
Après avoir rétréci $Y$, on peut donc supposer que $Z_{i, j} = \emptyset$, voir
lemme \ref{more-morphisms-lemma-empty-generic-fibre}.
Après avoir encore rétréci $Y$, on peut supposer que
$X_y = \bigcup X_{i, y}$ pour $y \in Y$, voir
lemme \ref{more-morphisms-lemma-cover-generic-fibre-neighbourhood}.
De plus, après avoir rétréci $Y$, on peut supposer que chaque $X_i \to Y$
est plat et de présentation finie, voir
Morphismes, proposition \ref{morphisms-proposition-generic-flatness}.
Les morphismes $X_i \to Y$ sont ouverts, voir
Morphismes, lemme \ref{morphisms-lemma-fppf-open}.
Il existe donc un voisinage ouvert $V$ de $\eta$ qui est contenu
dans $f(X_i)$ pour chaque $i$.
Pour tout $y \in V$, les schémas $X_{i, y}$ sont des
sous-ensembles fermés non vides de $X_y$, on a $X_y = \bigcup X_{i, y}$,
et les intersections $Z_{i, j, y} = X_{i, y} \cap X_{j, y}$
sont vides ! Cela implique clairement que
$X_y$ a au moins $n$ composantes connexes.
\end{proof}

\begin{lemma}
\label{more-morphisms-lemma-base-change-fibres-geometrically-connected}
Soit $f : X \to Y$ un morphisme de schémas.
Soit $g : Y' \to Y$ un morphisme quelconque, et notons
$f' : X' \to Y'$ le changement de base de $f$.
Alors
\begin{align*}
\{y' \in Y' \mid X'_{y'}\text{ est géométriquement connexe}\} \\
= g^{-1}(\{y \in Y \mid X_y\text{ est géométriquement connexe}\}).
\end{align*}
\end{lemma}

\begin{proof}
Cela se ramène à l'énoncé suivant : pour $y' \in Y'$ d'image
$y \in Y$, la fibre $X'_{y'} = X_y \times_y y'$ est géométriquement
connexe sur $\kappa(y')$ si et seulement si $X_y$ est géométriquement connexe
sur $\kappa(y)$. Cela résulte de
Variétés,
lemme \ref{varieties-lemma-geometrically-connected-check-after-extension}.
\end{proof}

\begin{lemma}
\label{more-morphisms-lemma-base-change-fibres-nr-geometrically-connected-components}
Soit $f : X \to Y$ un morphisme de schémas. Soit
$$
n_{X/Y} : Y \to \{0, 1, 2, 3, \ldots, \infty\}
$$
la fonction qui associe à $y \in Y$ le nombre de composantes connexes
de $(X_y)_K$, où $K$ est une extension séparablement close
de $\kappa(y)$. Cette fonction est bien définie et, si $g : Y' \to Y$ est un morphisme,
alors
$$
n_{X'/Y'} = n_{X/Y} \circ g
$$
où $X' \to Y'$ est le changement de base de $f$.
\end{lemma}

\begin{proof}
Supposons que $y' \in Y'$ ait pour image $y \in Y$.
Supposons que $K \supset \kappa(y)$ et $K' \supset \kappa(y')$ soient des extensions
séparablement closes. On peut alors choisir un diagramme commutatif
$$
\xymatrix{
K \ar[r] & K'' & K' \ar[l] \\
\kappa(y) \ar[u] \ar[rr] & & \kappa(y') \ar[u]
}
$$
de corps. Le résultat s'ensuit, car les morphismes de schémas
$$
\xymatrix{
(X'_{y'})_{K'} &
(X'_{y'})_{K''} = (X_y)_{K''} \ar[l] \ar[r] &
(X_y)_K
}
$$
induisent des bijections entre les composantes connexes, voir
Variétés,
lemme \ref{varieties-lemma-separably-closed-field-connected-components}.
\end{proof}

\begin{lemma}
\label{more-morphisms-lemma-geometrically-connected-generic-fibre}
Soit $f : X \to Y$ un morphisme de schémas.
Supposons que
\begin{enumerate}
\item $Y$ soit irréductible, de point générique $\eta$,
\item $X_\eta$ soit géométriquement connexe, et
\item $f$ soit de type fini.
\end{enumerate}
Alors il existe un sous-schéma ouvert non vide $V \subset Y$
tel que $X_V \to V$ ait des fibres géométriquement connexes.
\end{lemma}

\begin{proof}
Choisissons un diagramme
$$
\xymatrix{
X' \ar[d]_{f'} \ar[r]_{g'} & X_V \ar[r] \ar[d] & X \ar[d]^f \\
Y' \ar[r]^g & V \ar[r] & Y
}
$$
comme dans le
lemme \ref{more-morphisms-lemma-make-components-generic-fibre-geometrically-irreducible}.
Notons que la fibre générique de $f'$ est géométriquement connexe
(par exemple d'après le
lemme \ref{more-morphisms-lemma-base-change-fibres-nr-geometrically-connected-components}).
Supposons que le lemme soit vrai pour le morphisme $f'$. Cela signifie qu'il
existe un ouvert non vide $W \subset Y'$ tel que toute fibre de
$X' \to Y'$ au-dessus de $W$ soit géométriquement connexe.
Alors, comme $g$ est un morphisme ouvert d'après
Morphismes, lemme \ref{morphisms-lemma-etale-open},
toutes les fibres de $f$ aux points de l'ouvert non vide $V = g(W)$ sont
géométriquement connexes, voir
lemme \ref{more-morphisms-lemma-base-change-fibres-nr-geometrically-connected-components}.
On voit ainsi que l'on peut supposer que les composantes irréductibles
de la fibre générique $X_\eta$ sont géométriquement irréductibles.

\medskip\noindent
Soit $Y'$ le réduit de $Y$, et posons $X' = Y' \times_Y X$.
Il suffit alors de démontrer le lemme pour le morphisme $X' \to Y'$
(par exemple d'après le
lemme \ref{more-morphisms-lemma-base-change-fibres-nr-geometrically-connected-components}
une fois encore). Comme la fibre générique de $X' \to Y'$ est la même que la
fibre générique de $X \to Y$, on voit que l'on peut supposer que $Y$ est
irréductible et réduit (c'est-à-dire intègre, voir
Propriétés, lemme \ref{properties-lemma-characterize-integral})
et que les composantes irréductibles de la fibre générique $X_\eta$
sont géométriquement irréductibles.

\medskip\noindent
À ce stade, supposons que
$X_\eta = X_{1, \eta} \bigcup \ldots \bigcup X_{n, \eta}$
soit la décomposition de la fibre générique en composantes
(géométriquement) irréductibles.
Soit $X_i$ l'adhérence de $X_{i, \eta}$ dans $X$.
Après avoir rétréci $Y$, on peut supposer que
$X = \bigcup X_i$, voir
lemme \ref{more-morphisms-lemma-cover-generic-fibre-neighbourhood}.
Soit $Z_{i, j} = X_i \cap X_j$.
Soit
$$
\{1, \ldots, n\} \times \{1, \ldots, n\} = I \amalg J
$$
où $(i, j) \in I$ si $Z_{i, j, \eta} = \emptyset$ et
$(i, j) \in J$ si $Z_{i, j, \eta} \not = \emptyset$.
Après avoir rétréci $Y$, on peut supposer que $Z_{i, j} = \emptyset$
pour tout $(i, j) \in I$, voir
lemme \ref{more-morphisms-lemma-empty-generic-fibre}.
Après avoir rétréci $Y$, on obtient que $X_{i, y}$
est géométriquement irréductible pour chaque $i$ et tout $y \in Y$, voir
lemme \ref{more-morphisms-lemma-geom-irreducible-generic-fibre}.
Après avoir encore rétréci $Y$, on arrive à la situation où
chaque $Z_{i, j} \to Y$ est plat et de présentation finie pour
tout $(i, j) \in J$, voir
Morphismes, proposition \ref{morphisms-proposition-generic-flatness}.
Cela signifie que $f(Z_{i, j}) \subset Y$ est ouvert, voir
Morphismes, lemme \ref{morphisms-lemma-fppf-open}.
Nous affirmons que
$$
V  = \bigcap\nolimits_{(i, j) \in J} f(Z_{i, j})
$$
convient, c'est-à-dire que $X_y$ est géométriquement connexe pour tout
$y \in V$. En effet, le fait que $X_\eta$ soit connexe implique que
la relation d'équivalence engendrée par les paires de $J$ n'a qu'une
seule classe d'équivalence. Or, si $y \in V$ et si $K \supset \kappa(y)$
est une extension séparablement close, alors les composantes irréductibles
de $(X_y)_K$ sont les fibres $(X_{i, y})_K$. De plus, la construction et
$y \in V$ montrent que $(X_{i, y})_K$ rencontre $(X_{j, y})_K$
si et seulement si $(i, j) \in J$. La remarque sur les classes d'équivalence
montre donc que $(X_y)_K$ est connexe, et on conclut.
\end{proof}

\begin{lemma}
\label{more-morphisms-lemma-nr-geom-connected-components-good}
Soit $f : X \to Y$ un morphisme de schémas. Soit
$n_{X/Y}$ la fonction sur $Y$ qui compte le nombre de composantes
géométriquement connexes des fibres de $f$, introduite dans le
lemme \ref{more-morphisms-lemma-base-change-fibres-nr-geometrically-connected-components}.
Supposons $f$ de type fini.
Soit $y \in Y$ un point. Alors il existe un ouvert non vide
$V \subset \overline{\{y\}}$ tel que $n_{X/Y}|_V$ soit constant.
\end{lemma}

\begin{proof}
Soit $Z$ le schéma muni de la structure induite réduite sur $\overline{\{y\}}$.
Soit $f_Z : X_Z \to Z$ le changement de base de $f$. Il suffit clairement de démontrer
le lemme pour $f_Z$ et le point générique de $Z$. On peut donc supposer que
$Y$ est un schéma intègre, voir
Propriétés, lemme \ref{properties-lemma-characterize-integral}.
Notre but, dans ce cas, est de produire un ouvert non vide $V \subset Y$ tel que
$n_{X/Y}|_V$ soit constant.

\medskip\noindent
Appliquons le
lemme \ref{more-morphisms-lemma-make-components-generic-fibre-geometrically-irreducible}
à $f : X \to Y$ ; on obtient $g : Y' \to V \subset Y$. Comme $g : Y' \to V$ est
fini étale surjectif, donc en particulier ouvert (voir
Morphismes, lemme \ref{morphisms-lemma-etale-open}),
il suffit de démontrer qu'il existe un ouvert $V' \subset Y'$
tel que $n_{X'/Y'}|_{V'}$ soit constant, voir
lemme \ref{more-morphisms-lemma-base-change-fibres-nr-geometrically-irreducible-components}.
On voit ainsi que l'on peut supposer que toutes les composantes irréductibles de
la fibre générique $X_\eta$ sont géométriquement irréductibles sur $\kappa(\eta)$.
D'après
Variétés, lemme
\ref{varieties-lemma-irreducible-components-geometrically-irreducible},
cela implique que les composantes connexes de $X_\eta$ sont elles aussi
géométriquement connexes.

\medskip\noindent
À ce stade, supposons que
$X_\eta = X_{1, \eta} \bigcup \ldots \bigcup X_{n, \eta}$
soit la décomposition de la fibre générique en composantes
(géométriquement) connexes.
En particulier, $n_{X/Y}(\eta) = n$.
Soit $X_i$ l'adhérence de
$X_{i, \eta}$ dans $X$. Après avoir rétréci $Y$, on peut supposer que
$X = \bigcup X_i$, voir
lemme \ref{more-morphisms-lemma-cover-generic-fibre-neighbourhood}.
Après avoir encore rétréci $Y$, on voit que chaque fibre de
$f$ a au moins $n$ composantes connexes, voir
lemme \ref{more-morphisms-lemma-connected-components-in-neighbourhood}.
Ainsi $n_{X/Y}(y) \geq n$ pour tout $y \in Y$.
Après avoir encore rétréci $Y$, on obtient que $X_{i, y}$
est géométriquement connexe pour chaque $i$ et tout $y \in Y$, voir
lemme \ref{more-morphisms-lemma-geometrically-connected-generic-fibre}.
Comme $X_y = \bigcup X_{i, y}$,
cela montre que $n_{X/Y}(y) \leq n$ et achève la démonstration.
\end{proof}

\begin{lemma}
\label{more-morphisms-lemma-nr-geom-connected-components-constructible}
Soit $f : X \to Y$ un morphisme de schémas. Soit
$n_{X/Y}$ la fonction sur $Y$ qui compte le nombre de composantes
géométriquement connexes des fibres de $f$, introduite dans le
lemme \ref{more-morphisms-lemma-base-change-fibres-nr-geometrically-connected-components}.
Supposons $f$ de présentation finie. Alors les ensembles de niveau
$$
E_n = \{y \in Y \mid n_{X/Y}(y) = n\}
$$
de $n_{X/Y}$ sont localement constructibles dans $Y$.
\end{lemma}

\begin{proof}
Fixons $n$. Soit $y \in Y$. Nous devons montrer qu'il existe un voisinage ouvert
$V$ de $y$ dans $Y$ tel que $E_n \cap V$ soit constructible dans $V$. On peut donc
supposer que $Y$ est affine. Écrivons $Y = \Spec(A)$ et
$A = \colim A_i$ comme limite inductive d'algèbres
de type fini sur $\mathbf{Z}$. D'après
Limites, lemme \ref{limits-lemma-descend-finite-presentation},
on peut trouver un $i$ et un morphisme $f_i : X_i \to \Spec(A_i)$ de
présentation finie dont le changement de base à $Y$ redonne $f$. D'après le
lemme \ref{more-morphisms-lemma-base-change-fibres-nr-geometrically-connected-components},
il suffit de démontrer le lemme pour $f_i$. On se ramène ainsi au
cas où $Y$ est le spectre d'un anneau noethérien.

\medskip\noindent
Nous utiliserons le critère du
Topologie, lemme \ref{topology-lemma-characterize-constructible-Noetherian},
pour montrer que $E_n$ est constructible lorsque $Y$ est un schéma noethérien.
Pour cela, soit $Z \subset Y$ un sous-schéma fermé irréductible.
Nous devons montrer que $E_n \cap Z$ contient un ouvert non vide
ou n'est pas dense dans $Z$. Soit $\xi \in Z$ le point générique. Alors le
lemme \ref{more-morphisms-lemma-nr-geom-connected-components-good}
montre que $n_{X/Y}$ est constant dans un voisinage de $\xi$ dans $Z$.
Cela implique clairement le résultat voulu.
\end{proof}

\begin{lemma}
\label{more-morphisms-lemma-connected-flat-over-dvr}
\begin{slogan}
Une dégénérescence plate d'un schéma non connexe est soit non connexe,
soit non réduite.
\end{slogan}
Soit $f : X \to S$ un morphisme de schémas.
Supposons que
\begin{enumerate}
\item $S$ soit le spectre d'un anneau de valuation discrète,
\item $f$ soit plat,
\item $X$ soit connexe,
\item la fibre fermée $X_s$ soit réduite.
\end{enumerate}
Alors la fibre générique $X_\eta$ est connexe.
\end{lemma}

\begin{proof}
Écrivons $S = \Spec(R)$ et soit $\pi \in R$ une uniformisante.
Pour obtenir une contradiction, supposons que $X_\eta$ ne soit pas connexe.
Cela signifie qu'il existe un idempotent non trivial
$e \in \Gamma(X_\eta, \mathcal{O}_{X_\eta})$.
Soit $U = \Spec(A)$ un ouvert affine quelconque de $X$.
Notons que $\pi$ est un non-diviseur de zéro dans $A$, puisque $A$ est plat sur $R$, voir
Plus sur l'algèbre, lemme \ref{more-algebra-lemma-flat-torsion-free},
par exemple. Alors $e|_{U_\eta}$ correspond à un élément $e \in A[1/\pi]$.
Soit $z \in A$ un élément tel que $e = z/\pi^n$, avec $n \geq 0$ minimal.
Notons que $z^2 = \pi^nz$. Cela signifie que $z \bmod \pi A$ est nilpotent
si $n > 0$. Par hypothèse, $A/\pi A$ est réduit, et la minimalité de
$n$ implique donc $n = 0$. On conclut ainsi que $e \in A$ ! Autrement dit,
$e \in \Gamma(X, \mathcal{O}_X)$. Comme $X$ est connexe, il s'ensuit
que $e$ est un idempotent trivial, ce qui est une contradiction.
\end{proof}

\section{Composantes connexes rencontrant une section}
\label{more-morphisms-section-connected-components}

\noindent
Les résultats de cette section s'appliquent en particulier à un schéma en groupes
$G \to S$ et à sa section neutre $e : S \to G$.

\begin{situation}
\label{more-morphisms-situation-connected-along-section}
Ici, $f : X \to Y$ est un morphisme de schémas, et
$s : Y \to X$ est une section de $f$.
Pour tout $y \in Y$, on note $X^0_y$ la composante connexe de $X_y$
contenant $s(y)$. Enfin, on pose $X^0 = \bigcup_{y \in Y} X^0_y$.
\end{situation}

\begin{lemma}
\label{more-morphisms-lemma-base-change-connected-along-section}
Soient $f : X \to Y$, $s : Y \to X$ comme dans la
situation \ref{more-morphisms-situation-connected-along-section}.
Si $g : Y' \to Y$ est un morphisme quelconque, considérons le diagramme de changement de base
$$
\xymatrix{
X' \ar[r]_{g'} \ar[d]^{f'} & X \ar[d]_f \\
Y' \ar@/^1pc/[u]^{s'} \ar[r]^g & Y \ar@/_1pc/[u]_s
}
$$
de sorte que l'on obtient $(X')^0 \subset X'$.
Alors $(X')^0 = (g')^{-1}(X^0)$.
\end{lemma}

\begin{proof}
Soit $y' \in Y'$ d'image $y \in Y$. On peut considérer
$X^0_y$ comme un sous-schéma fermé de $X_y$, voir par exemple
Morphismes,
définition \ref{morphisms-definition-scheme-structure-connected-component}.
Comme $s(y) \in X^0_y$, il résulte de
Variétés, lemme
\ref{varieties-lemma-geometrically-connected-if-connected-and-point}
que $X_y^0$ est un schéma géométriquement connexe sur $\kappa(y)$.
Ainsi $X_y^0 \times_y y' \to X'_{y'}$ est un sous-schéma fermé connexe
qui contient $s'(y')$. Donc $X_y^0 \times_y y' \subset (X'_{y'})^0$.
L'autre inclusion $X_y^0 \times_y y' \supset (X'_{y'})^0$ est claire,
car l'image de $(X'_{y'})^0$ dans $X_y$ est un sous-ensemble connexe de $X_y$ qui
contient $s(y)$.
\end{proof}

\begin{lemma}
\label{more-morphisms-lemma-connected-along-section-good}
Soient $f : X \to Y$, $s : Y \to X$ comme dans la
situation \ref{more-morphisms-situation-connected-along-section}.
Supposons $f$ de type fini. Soit $y \in Y$ un point.
Alors il existe un ouvert non vide $V \subset \overline{\{y\}}$ tel que
l'image réciproque de $X^0$ dans le changement de base $X_V$ soit ouverte et fermée dans
$X_V$.
\end{lemma}

\begin{proof}
Soit $Z \subset Y$ le sous-schéma fermé muni de la structure induite réduite
sur $\overline{\{y\}}$. Soient $f_Z : X_Z \to Z$ et $s_Z : Z \to X_Z$
les changements de base de $f$ et $s$. D'après le
lemme \ref{more-morphisms-lemma-base-change-connected-along-section},
on a $(X_Z)^0 = (X^0)_Z$. Il suffit donc de démontrer le lemme pour
le morphisme $X_Z \to Z$ et le point $x \in X_Z$ qui s'envoie sur le point générique
de $Z$. Autrement dit, on s'est ramené au cas
où $Y$ est un schéma intègre (voir
Propriétés, lemme \ref{properties-lemma-characterize-integral})
de point générique $\eta$. Notre but est de montrer qu'après avoir rétréci
$Y$, le sous-ensemble $X^0$ devient un sous-ensemble ouvert et fermé de $X$.

\medskip\noindent
Notons que le schéma $X_\eta$ est de type fini sur un corps, donc noethérien.
Ses composantes connexes sont ainsi à la fois ouvertes et fermées. On peut donc écrire
$X_\eta = X_\eta^0 \amalg T_\eta$ pour un certain sous-ensemble ouvert et fermé
$T_\eta$ de $X_\eta$. Soit ensuite $T \subset X$ l'adhérence de $T_\eta$,
et soit $X^{00} \subset X$ l'adhérence de $X_\eta^0$. Notons que
$T_\eta$, respectivement $X^0_\eta$, est la fibre générique de $T$, respectivement de $X^{00}$,
voir la discussion précédant le
lemme \ref{more-morphisms-lemma-cover-generic-fibre-neighbourhood}.
De plus, ce lemme implique qu'après avoir rétréci $Y$, on peut supposer que
$X = X^{00} \cup T$ (au sens des ensembles).
Notons que $(T \cap X^{00})_\eta = T_\eta \cap X^0_\eta = \emptyset$.
Après avoir rétréci $Y$, on peut donc supposer que $T \cap X^{00} = \emptyset$, voir
lemme \ref{more-morphisms-lemma-empty-generic-fibre}.
En particulier, $X^{00}$ est ouvert dans $X$. Notons que $X^0_\eta$ est connexe
et possède un point rationnel, à savoir $s(\eta)$ ; il est donc géométriquement
connexe, voir
Variétés,
lemme \ref{varieties-lemma-geometrically-connected-if-connected-and-point}.
Ainsi, après avoir rétréci $Y$, on peut supposer que toutes les fibres de $X^{00} \to Y$
sont géométriquement connexes, voir
lemme \ref{more-morphisms-lemma-geometrically-connected-generic-fibre}.
Il s'ensuit alors que les fibres $X^{00}_y$ sont des
sous-ensembles ouverts, fermés et connexes de $X_y$ contenant $s(y)$.
On en déduit que $X^0 = X^{00}$, et on conclut.
\end{proof}

\begin{lemma}
\label{more-morphisms-lemma-connected-along-section-locally-constructible}
Soient $f : X \to Y$, $s : Y \to X$ comme dans la
situation \ref{more-morphisms-situation-connected-along-section}.
Si $f$ est de présentation finie, alors $X^0$ est localement constructible
dans $X$.
\end{lemma}

\begin{proof}
Soit $x \in X$. Nous devons montrer qu'il existe un voisinage ouvert
$U$ de $x$ tel que $X^0 \cap U$ soit constructible dans $U$.
Cela nous ramène au cas où $Y$ est affine.
Écrivons $Y = \Spec(A)$ et $A = \colim A_i$ comme limite
inductive d'algèbres de type fini sur $\mathbf{Z}$. D'après
Limites, lemme \ref{limits-lemma-descend-finite-presentation},
on peut trouver un $i$ et un morphisme $f_i : X_i \to \Spec(A_i)$ de
présentation finie, muni d'une section $s_i : \Spec(A_i) \to X_i$
dont le changement de base à $Y$ redonne $f$ et la section $s$. D'après le
lemme \ref{more-morphisms-lemma-base-change-connected-along-section},
il suffit de démontrer le lemme pour $f_i, s_i$. On se ramène ainsi au
cas où $Y$ est le spectre d'un anneau noethérien.

\medskip\noindent
Supposons que $Y$ soit un schéma affine noethérien. Puisque $f$ est de présentation finie,
c'est-à-dire de type fini, on voit que $X$ est lui aussi un schéma noethérien, voir
Morphismes, lemme \ref{morphisms-lemma-finite-type-noetherian}.
Pour démontrer le lemme dans
ce cas, il suffit de montrer que, pour tout sous-ensemble fermé irréductible
$Z \subset X$, l'intersection $Z \cap X^0$ contient un ouvert non vide
de $Z$ ou n'est pas dense dans $Z$, voir
Topologie, lemme \ref{topology-lemma-characterize-constructible-Noetherian}.
Soit $x \in Z$ le point générique, et soit $y = f(x)$. D'après le
lemme \ref{more-morphisms-lemma-connected-along-section-good},
il existe un sous-ensemble ouvert non vide $V \subset \overline{\{y\}}$ tel
que $X^0 \cap X_V$ soit ouvert et fermé dans $X_V$. Comme
$f(Z) \subset \overline{\{y\}}$ et $f(x) = y \in V$, on voit que
$W = f^{-1}(V) \cap Z$ est un sous-ensemble ouvert non vide de $Z$. Il s'ensuit que
$X^0 \cap W$ est ouvert et fermé dans $W$. Puisque $W$ est irréductible,
on voit que $X^0 \cap W$ est soit vide, soit égal à $W$.
Cela démontre le lemme.
\end{proof}

\begin{lemma}
\label{more-morphisms-lemma-connected-along-section-open-neighbourhood}
Soient $f : X \to Y$, $s : Y \to X$ comme dans la
situation \ref{more-morphisms-situation-connected-along-section}.
Soit $y \in Y$ un point.
Supposons que
\begin{enumerate}
\item $f$ soit de présentation finie et plat, et
\item la fibre $X_y$ soit géométriquement réduite.
\end{enumerate}
Alors $X^0$ est un voisinage de $X^0_y$ dans $X$.
\end{lemma}

\begin{proof}
On peut remplacer $Y$ par un voisinage ouvert affine de $y$.
Écrivons $Y = \Spec(A)$ et $A = \colim A_i$ comme limite
inductive d'algèbres de type fini sur $\mathbf{Z}$. D'après
Limites, lemme \ref{limits-lemma-descend-finite-presentation},
on peut trouver un $i$ et un morphisme $f_i : X_i \to \Spec(A_i)$ de
présentation finie, muni d'une section $s_i : \Spec(A_i) \to X_i$
dont le changement de base à $Y$ redonne $f$ et la section $s$.
Après avoir éventuellement augmenté $i$, on peut aussi supposer que $f_i$ est plat, voir
Limites, lemme \ref{limits-lemma-descend-flat-finite-presentation}.
Soit $y_i$ l'image de $y$ dans $Y_i$. Notons que
$X_y = (X_{i, y_i}) \times_{y_i} y$. Ainsi $X_{i, y_i}$ est géométriquement
réduit, voir
Variétés, lemme \ref{varieties-lemma-geometrically-reduced-upstairs}.
D'après le
lemme \ref{more-morphisms-lemma-base-change-connected-along-section},
il suffit de démontrer le lemme pour le système $f_i, s_i, y_i \in Y_i$.
On se ramène ainsi au cas où $Y$ est le spectre d'un anneau noethérien.

\medskip\noindent
Supposons que $Y$ soit le spectre d'un anneau noethérien.
Puisque $f$ est de présentation finie,
c'est-à-dire de type fini, on voit que $X$ est lui aussi un schéma noethérien, voir
Morphismes, lemme \ref{morphisms-lemma-finite-type-noetherian}.
Soit $x \in X^0$ un point au-dessus de $y$. D'après
Topologie, lemme \ref{topology-lemma-constructible-neighbourhood-Noetherian},
il suffit de démontrer que, pour tout fermé irréductible $Z \subset X$
passant par $x$, l'intersection $X^0 \cap Z$ est dense dans $Z$.
En particulier, il suffit de démontrer que le point générique $x' \in Z$
appartient à $X^0$. D'après
Propriétés, lemme \ref{properties-lemma-locally-Noetherian-specialization-dvr},
on peut trouver un anneau de valuation discrète $R$ et un morphisme
$\Spec(R) \to X$ qui envoie le point spécial sur $x$ et
le point générique sur $x'$. Nous considérerons $\Spec(R)$
comme un schéma sur $Y$ via la composée $\Spec(R) \to X \to Y$. D'après le
lemme \ref{more-morphisms-lemma-base-change-connected-along-section},
on sait que $(X_R)^0$ est l'image réciproque de $X^0$.
Par construction, on dispose d'une seconde section $t : \Spec(R) \to X_R$
(outre le changement de base $s_R$ de $s$)
du morphisme structural $X_R \to \Spec(R)$ telle que
$t(\eta_R)$ soit un point de $X_R$ qui s'envoie sur $x'$ et
$t(0_R)$ un point de $X_R$ qui s'envoie sur $x$. Notons que
$t(0_R)$ appartient à $(X_R)^0$ et que $t(\eta_R) \leadsto t(0_R)$.
Il suffit donc de démontrer que cela implique $t(\eta_R) \in (X_R)^0$.
Il suffit par conséquent de démontrer le lemme lorsque $Y$
est le spectre d'un anneau de valuation discrète et $y$ son point fermé.

\medskip\noindent
Supposons que $Y$ soit le spectre d'un anneau de valuation discrète et que $y$ soit son point
fermé. Notre but est de démontrer que $X^0$ est un voisinage de $X^0_y$.
Notons que $X^0_y$ est ouvert et fermé dans $X_y$, puisque $X_y$ possède un nombre fini
de composantes irréductibles. Le complémentaire $C = X_y \setminus X_y^0$
est donc fermé dans $X$. Ainsi $U = X \setminus C$ est un voisinage ouvert de
$X^0_y$, et $U^0 = X^0$. Il suffit donc de démontrer le résultat pour le
morphisme $U \to Y$. Autrement dit, on peut supposer que
$X_y$ est connexe. Supposons que $X$ ne soit pas connexe, et écrivons
$X = X_1 \amalg \ldots \amalg X_n$ comme une décomposition en composantes
connexes. Alors $s(Y)$ est entièrement contenu dans l'un des $X_i$.
Disons $s(Y) \subset X_1$. Alors $X^0 \subset X_1$. On peut donc remplacer
$X$ par $X_1$ et supposer que $X$ est connexe. À ce stade, le
lemme \ref{more-morphisms-lemma-connected-flat-over-dvr}
implique que $X_\eta$ est connexe, c'est-à-dire $X^0 = X$, et on conclut.
\end{proof}

\begin{lemma}
\label{more-morphisms-lemma-connected-along-section-open}
Soient $f : X \to Y$, $s : Y \to X$ comme dans la
situation \ref{more-morphisms-situation-connected-along-section}.
Supposons que
\begin{enumerate}
\item $f$ soit de présentation finie et plat, et
\item toutes les fibres de $f$ soient géométriquement réduites.
\end{enumerate}
Alors $X^0$ est ouvert dans $X$.
\end{lemma}

\begin{proof}
C'est une conséquence immédiate du
lemme \ref{more-morphisms-lemma-connected-along-section-open-neighbourhood}.
\end{proof}

\section{Dimension des fibres}
\label{more-morphisms-section-dimension}

\begin{lemma}
\label{more-morphisms-lemma-dimension-in-neighbourhood}
Soit $f : X \to Y$ un morphisme de schémas. Supposons $Y$ irréductible de
point générique $\eta$ et $f$ de type fini. Si $X_\eta$ est de dimension $n$,
alors il existe un ouvert non vide $V \subset Y$
tel que, pour tout $y \in V$, la fibre $X_y$ soit de dimension $n$.
\end{lemma}

\begin{proof}
Soit $Z = \{x \in X \mid \dim_x(X_{f(x)}) > n \}$. D'après
Morphismes, lemme \ref{morphisms-lemma-openness-bounded-dimension-fibres},
c'est un sous-ensemble fermé de $X$. Par hypothèse, $Z_\eta = \emptyset$.
Ainsi, d'après le
lemme \ref{more-morphisms-lemma-empty-generic-fibre},
on peut rétrécir $Y$ et supposer que $Z = \emptyset$. Soit
$Z' = \{x \in X \mid \dim_x(X_{f(x)}) > n - 1 \} =
\{x \in X \mid \dim_x(X_{f(x)}) = n\}$. Comme précédemment, c'est un sous-ensemble fermé
de $X$. Par hypothèse, on a $Z'_\eta \not = \emptyset$. Ainsi, après
avoir rétréci $Y$, on peut supposer que $Z' \to Y$ est surjectif, voir
lemme \ref{more-morphisms-lemma-nonempty-generic-fibre}.
On conclut.
\end{proof}

\begin{lemma}
\label{more-morphisms-lemma-base-change-dimension-fibres}
Soit $f : X \to Y$ un morphisme de type fini. Soit
$$
n_{X/Y} : Y \to \{0, 1, 2, 3, \ldots, \infty\}
$$
la fonction qui associe à $y \in Y$ la dimension de $X_y$.
Si $g : Y' \to Y$ est un morphisme, alors
$$
n_{X'/Y'} = n_{X/Y} \circ g
$$
où $X' \to Y'$ est le changement de base de $f$.
\end{lemma}

\begin{proof}
Cela résulte de
Morphismes, lemme \ref{morphisms-lemma-dimension-fibre-after-base-change}.
\end{proof}

\begin{lemma}
\label{more-morphisms-lemma-dimension-fibres-constructible}
Soit $f : X \to Y$ un morphisme de schémas. Soit
$n_{X/Y}$ la fonction sur $Y$ donnant la dimension des fibres de $f$,
introduite dans le
lemme \ref{more-morphisms-lemma-base-change-dimension-fibres}.
Supposons $f$ de présentation finie. Alors les ensembles de niveau
$$
E_n = \{y \in Y \mid n_{X/Y}(y) = n\}
$$
de $n_{X/Y}$ sont localement constructibles dans $Y$.
\end{lemma}

\begin{proof}
Fixons $n$. Soit $y \in Y$. Nous devons montrer qu'il existe un voisinage ouvert
$V$ de $y$ dans $Y$ tel que $E_n \cap V$ soit constructible dans $V$. On peut donc
supposer que $Y$ est affine. Écrivons $Y = \Spec(A)$ et
$A = \colim A_i$ comme limite inductive d'algèbres
de type fini sur $\mathbf{Z}$. D'après
Limites, lemme \ref{limits-lemma-descend-finite-presentation},
on peut trouver un $i$ et un morphisme $f_i : X_i \to \Spec(A_i)$ de
présentation finie dont le changement de base à $Y$ redonne $f$. D'après le
lemme \ref{more-morphisms-lemma-base-change-dimension-fibres},
il suffit de démontrer le lemme pour $f_i$. On se ramène ainsi au
cas où $Y$ est le spectre d'un anneau noethérien.

\medskip\noindent
Nous utiliserons le critère de
Topologie, lemme \ref{topology-lemma-characterize-constructible-Noetherian},
pour démontrer que $E_n$ est constructible lorsque $Y$ est un schéma noethérien.
Pour cela, soit $Z \subset Y$ un sous-schéma fermé irréductible.
Nous devons montrer que $E_n \cap Z$ contient un sous-ensemble ouvert non vide
ou n'est pas dense dans $Z$. Soit $\xi \in Z$ le point générique. Alors le
lemme \ref{more-morphisms-lemma-dimension-in-neighbourhood}
montre que $n_{X/Y}$ est constante dans un voisinage de $\xi$ dans $Z$.
Cela donne le résultat voulu.
\end{proof}

\begin{lemma}
\label{more-morphisms-lemma-dimension-fibres-flat}
Soit $f : X \to Y$ un morphisme plat de schémas, de présentation finie. Soit
$n_{X/Y}$ la fonction sur $Y$ donnant la dimension des fibres de $f$,
introduite dans le lemme \ref{more-morphisms-lemma-base-change-dimension-fibres}.
Alors $n_{X/Y}$ est semi-continue inférieurement.
\end{lemma}

\begin{proof}
Soit $W \subset X$, $W = \coprod_{d \geq 0} U_d$ l'ouvert construit dans les
lemmes \ref{more-morphisms-lemma-flat-finite-presentation-CM-open} et
\ref{more-morphisms-lemma-flat-finite-presentation-CM-pieces}.
Soit $y \in Y$ un point. Si $n_{X/Y}(y) = \dim(X_y) = n$, alors
$y$ appartient à l'image de $U_n \to Y$.
D'après Morphismes, lemme \ref{morphisms-lemma-fppf-open},
on voit que $f(U_n)$ est ouvert dans $Y$.
Il existe donc un voisinage ouvert de $y$ où
$n_{X/Y}$ est $\geq n$.
\end{proof}

\begin{lemma}
\label{more-morphisms-lemma-dimension-fibres-proper}
Soit $f : X \to Y$ un morphisme propre de schémas. Soit
$n_{X/Y}$ la fonction sur $Y$ donnant la dimension des fibres de $f$,
introduite dans le lemme \ref{more-morphisms-lemma-base-change-dimension-fibres}.
Alors $n_{X/Y}$ est semi-continue supérieurement.
\end{lemma}

\begin{proof}
Soit $Z_d = \{x \in X \mid \dim_x(X_{f(x)}) > d\}$.
Alors $Z_d$ est un sous-ensemble fermé de $X$ d'après
Morphismes, lemme \ref{morphisms-lemma-openness-bounded-dimension-fibres}.
Puisque $f$ est propre, $f(Z_d)$ est fermé.
Comme $y \in f(Z_d) \Leftrightarrow n_{X/Y}(y) > d$,
on voit que le lemme est vrai.
\end{proof}

\begin{lemma}
\label{more-morphisms-lemma-dimension-fibres-proper-flat}
Soit $f : X \to Y$ un morphisme propre et plat de schémas, de présentation finie.
Soit $n_{X/Y}$ la fonction sur $Y$ donnant la dimension des fibres de $f$,
introduite dans le lemme \ref{more-morphisms-lemma-base-change-dimension-fibres}.
Alors $n_{X/Y}$ est localement constante.
\end{lemma}

\begin{proof}
C'est une conséquence immédiate des
lemmes \ref{more-morphisms-lemma-dimension-fibres-flat} et
\ref{more-morphisms-lemma-dimension-fibres-proper}.
\end{proof}

\section{Normalisation de Noether relative faible}
\label{more-morphisms-section-relative-noether-normalization}

\noindent
Le but de cette section est de démontrer le
lemme \ref{more-morphisms-lemma-weak-relative-noether-normalization}.

\begin{lemma}
\label{more-morphisms-lemma-silly}
Soit $R$ un anneau. Soient $\mathfrak p_1, \ldots, \mathfrak p_r$
des idéaux premiers de $R$ tels que $\mathfrak p_i \not \subset \mathfrak p_j$
si $i \not = j$. Soient $k_i \subset \kappa(\mathfrak p_i)$ des
sous-corps tels que les extensions $\kappa(\mathfrak p_i)/k_i$
ne soient pas algébriques. Soit $J \subset R$ un idéal qui n'est contenu
dans aucun des $\mathfrak p_i$. Alors il existe un élément $x \in J$
tel que l'image de $x$ dans $\kappa(\mathfrak p_i)$
est transcendante sur $k_i$ pour $i = 1, \ldots, r$.
\end{lemma}

\begin{proof}
L'idéal
$J_i = J \mathfrak p_1 \ldots \hat{\mathfrak p}_i \ldots \mathfrak p_r$ n'est pas
contenu dans $\mathfrak p_i$, voir
Algèbre, lemme \ref{algebra-lemma-product-ideals-in-prime}.
Il s'ensuit que tout élément $\xi$ de
$\kappa(\mathfrak p_i) = \text{Frac}(B/\mathfrak p_i)$
est de la forme $\xi = a/b$, avec $a, b \in J_i$
et $b \not \in \mathfrak p_i$. En choisissant $\xi$ transcendant
sur $k_i$, on voit que soit $a$, soit $b$ s'envoie sur un élément
de $\kappa(\mathfrak p_i)$ transcendant sur $k_i$.
On conclut que, pour tout $i = 1, \ldots, r$,
on peut trouver un élément
$x_i \in J_i = J \mathfrak p_1 \ldots \hat{\mathfrak p}_i \ldots \mathfrak p_r$
qui s'envoie sur un élément de $\kappa(\mathfrak p_i)$
transcendant sur $k_i$. Alors $x = x_1 + \ldots + x_r$ convient.
\end{proof}

\begin{lemma}
\label{more-morphisms-lemma-sum-is-ok}
Soit $R \to S$ un homomorphisme d'anneaux de type fini. Soient $d \geq 0$ et $a, b \in S$.
Supposons que les fibres de
$$
f_a : \Spec(S) \longrightarrow \mathbf{A}^1_R
$$
donné par l'homomorphisme de $R$-algèbres $R[x] \to S$ qui envoie $x$ sur $a$
soient de dimension $\leq d$. Alors il existe un $n_0$ tel que, pour
$n \geq n_0$, les fibres de
$$
f_{a^n + b} : \Spec(S) \longrightarrow \mathbf{A}^1_R
$$
donné par l'homomorphisme de $R$-algèbres $R[x] \to S$ qui envoie $x$ sur $a^n + b$
soient de dimension $\leq d$.
\end{lemma}

\begin{proof}
Dans ce paragraphe, on se ramène au cas où $R \to S$ est de présentation
finie. En effet, écrivons $S = R[A, B, x_1, \ldots, x_n]/J$ pour un idéal
$J \subset R[x_1, \ldots, x_n]$, où $A$ et $B$ s'envoient sur $a$ et $b$ dans $S$.
Alors $J$ est la réunion de ses idéaux de type fini $J_\lambda \subset J$.
Posons $S_\lambda = R[A, B, x_1, \ldots, x_n]/J_\lambda$ et
notons $a_\lambda, b_\lambda \in S_\lambda$ les images de $A$ et $B$.
Alors, pour un certain $\lambda$, les fibres de
$$
f_{a_\lambda} : \Spec(S_\lambda) \longrightarrow \mathbf{A}^1_R
$$
sont de dimension $\leq d$, voir Limites, lemme \ref{limits-lemma-limit-dimension}.
Fixons un tel $\lambda$. Si l'on peut trouver un $n_0$ qui convient pour $R \to S_\lambda$,
$a_\lambda$, $b_\lambda$, alors $n_0$ convient pour $R \to S$. En effet, les fibres
de $f_{a_\lambda^n + b_\lambda} : \Spec(S_\lambda) \to \mathbf{A}^1_R$
contiennent les fibres de $f_{a^n + b} : \Spec(S) \to \mathbf{A}^1_R$.
On est ainsi ramené au cas traité dans le paragraphe suivant.

\medskip\noindent
Supposons que $R \to S$ soit de présentation finie. Dans ce paragraphe, on se ramène
au cas où $R$ est de type fini sur $\mathbf{Z}$. D'après
Algèbre, lemme \ref{algebra-lemma-limit-module-finite-presentation},
on peut trouver un ensemble ordonné filtrant $\Lambda$ et un système d'homomorphismes d'anneaux
$R_\lambda \to S_\lambda$ indexé par $\Lambda$, dont la colimite est
$R \to S$, tels que $S_\mu = S_\lambda \otimes_{R_\lambda} R_\mu$
pour $\mu \geq \lambda$ et que chacun des $R_\lambda$ et $S_\lambda$
soit de type fini sur $\mathbf{Z}$. Choisissons $\lambda_0 \in \Lambda$ et des
éléments $a_{\lambda_0}, b_{\lambda_0} \in S_{\lambda_0}$ qui s'envoient sur
$a, b \in S$. Pour $\lambda \geq \lambda_0$, notons
$a_\lambda, b_\lambda \in S_\lambda$
l'image de $a_{\lambda_0}, b_{\lambda_0}$.
Alors, pour $\lambda \geq \lambda_0$ assez grand, les fibres de
$$
f_{a_\lambda} : 
\Spec(S_\lambda) \longrightarrow \mathbf{A}^1_{R_\lambda}
$$
sont de dimension $\leq d$, voir
Limites, lemme \ref{limits-lemma-descend-dimension-d}.
Fixons un tel $\lambda$. Si l'on peut trouver un $n_0$ qui convient pour
$R_\lambda \to S_\lambda$, $a_\lambda$, $b_\lambda$, alors $n_0$
convient pour $R \to S$. En effet, toute fibre de
$f_{a^n + b} : \Spec(S) \to \mathbf{A}^1_R$
a la même dimension qu'une fibre de
$f_{a_\lambda^n + b_\lambda} : \Spec(S_\lambda) \to \mathbf{A}^1_{R_\lambda}$
d'après Morphismes, lemme \ref{morphisms-lemma-dimension-fibre-after-base-change}.
On est ainsi ramené au cas traité dans le paragraphe suivant.

\medskip\noindent
Supposons que $R$ et $S$ soient de type fini sur $\mathbf{Z}$. En particulier,
la dimension de $R$ est finie, et on peut raisonner par récurrence sur $\dim(R)$.
On peut donc supposer le résultat vrai pour toutes les situations où
$R' \to S'$, $a$, $b$ sont comme dans le lemme, avec $R'$ et $S'$ de type fini
sur $\mathbf{Z}$, mais où $\dim(R') < \dim(R)$.

\medskip\noindent
Puisque l'énoncé ne concerne que la topologie du spectre de $S$,
on peut supposer $S$ réduit. Soit $S^\nu$ la normalisation de $S$.
Alors $S \subset S^\nu$ est une extension finie, puisque $S$ est excellent, voir
Algèbre, proposition \ref{algebra-proposition-ubiquity-nagata},
et Morphismes, lemme \ref{morphisms-lemma-nagata-normalization}.
Ainsi $\Spec(S^\nu) \to \Spec(S)$ est surjectif et fini
(Algèbre, lemme \ref{algebra-lemma-integral-overring-surjective}).
Il s'ensuit que, si le résultat est vrai pour $R \to S^\nu$ et les
images de $a$, $b$ dans $S^\nu$, alors il est vrai pour $R \to S$, $a$, $b$.
(Un petit détail est omis.) On est ainsi ramené au cas traité dans le
paragraphe suivant.

\medskip\noindent
Supposons que $R$ et $S$ soient de type fini sur $\mathbf{Z}$ et que
$S$ soit normal. Alors $S = S_1 \times \ldots \times S_r$ pour certains
anneaux intègres normaux $S_i$. Si le résultat est vrai pour chaque $R \to S_i$
et pour les images de $a$, $b$ dans $S_i$, alors il est vrai pour
$R \to S$, $a$, $b$. (Un petit détail est omis.)
On est ainsi ramené au cas traité dans le paragraphe suivant.

\medskip\noindent
Supposons que $R$ et $S$ soient de type fini sur $\mathbf{Z}$ et que
$S$ soit un anneau intègre normal. On peut remplacer $R$ par l'image de $R$ dans $S$
(cela n'augmente pas la dimension de $R$).
On est ainsi ramené au cas traité dans le paragraphe suivant.

\medskip\noindent
Supposons que $R \subset S$ soient de type fini sur $\mathbf{Z}$ et que
$S$ soit un anneau intègre normal. Considérons le morphisme
$$
f_a : \Spec(S) \to \mathbf{A}^1_R
$$
L'hypothèse nous dit que $f_a$ a des fibres de dimension $\leq d$.
Les fibres de $f : \Spec(S) \to \Spec(R)$ sont donc de dimension $\leq d + 1$
(Morphismes, lemme \ref{morphisms-lemma-dimension-fibre-at-a-point-additive}).
Considérons le morphisme de schémas intègres
$$
\phi : \Spec(S) \to \mathbf{A}^2_R = \Spec(R[x, y])
$$
correspondant à l'homomorphisme de $R$-algèbres $R[x, y] \to S$ qui envoie $x$ sur $a$
et $y$ sur $b$. Il y a deux cas à considérer :
\begin{enumerate}
\item $\phi$ est dominant, et
\item $\phi$ n'est pas dominant.
\end{enumerate}
Nous affirmons que, dans les deux cas, il existe un entier $n_0$
et un ouvert non vide $V \subset \Spec(R)$ tels que, pour $n \geq n_0$,
les fibres de $f_{a^n + b}$ aux points $q \in \mathbf{A}^1_V$
soient de dimension $\leq d$.

\medskip\noindent
Démontrons l'affirmation dans le cas (1). On a $f_{a^n + b} = \pi_n \circ \phi$,
où
$$
\pi_n :  \mathbf{A}^2_R \to \mathbf{A}^1_R
$$
est le morphisme plat correspondant à l'homomorphisme de $R$-algèbres $R[x] \to R[x, y]$
qui envoie $x$ sur $x^n + y$. Puisque $\phi$ est dominant, il existe un ouvert dense
$U \subset \Spec(S)$ tel que $\phi|_U : U \to \mathbf{A}^2_R$ soit plat
(cela résulte par exemple de la platitude générique, voir
Morphismes, proposition \ref{morphisms-proposition-generic-flatness}).
Alors la composée
$$
f_{a^n + b}|_U :
U \xrightarrow{\phi|_U}
\mathbf{A}^2_R
\xrightarrow{\pi_n}
\mathbf{A}^1_R
$$
est elle aussi plate. Les fibres de ce morphisme sont donc de codimension au moins
$1$ dans les fibres de $f|_U : U \to \Spec(R)$ d'après
Morphismes, lemme \ref{morphisms-lemma-dimension-fibre-at-a-point-additive}.
Autrement dit, les fibres de $f_{a^n + b}|_U$ sont de dimension $\leq d$.
D'autre part, puisque $U$ est dense dans $\Spec(S)$, on peut trouver un ouvert non vide
$V \subset \Spec(R)$ tel que $U \cap f^{-1}(p) \subset f^{-1}(p)$
soit dense pour tout $p \in V$ (voir par exemple le lemme
\ref{more-morphisms-lemma-nowhere-dense-generic-fibre}).
Ainsi $\dim(f^{-1}(p) \setminus U \cap f^{-1}(p)) \leq d$, et
on conclut que notre affirmation est vraie (puisque toute fibre de
$f_{a^n + b} : \Spec(S) \to \mathbf{A}^1_R$
est contenue dans une fibre de $f$).

\medskip\noindent
Cas (2). Dans ce cas, on peut trouver un élément non nul $g = \sum c_{ij} x^i y^j$
de $R[x, y]$ tel que $\Im(\phi) \subset V(g)$. On peut en fait supposer
$g$ irréductible sur $\text{Frac}(R)$. Si $g \in R[x]$, disons de coefficient
dominant $c$, alors, sur $V = D(c) \subset \Spec(R)$, les fibres de
$f$ sont déjà de dimension $\leq d$ (car l'image de $f_a$ est contenue
dans $V(g) \subset \mathbf{A}^1_R$, qui a des fibres finies sur $V$). On peut donc
supposer que $g$ n'appartient pas à $R[x]$.
Soit $s \geq 1$ le degré de $g$ comme polynôme en $y$, et
soit $t$ le degré de $\sum c_{is} x^i$ comme polynôme en $x$.
Alors $c_{ts}$ est non nul et
$$
g(x, -x^n) = (-1)^s c_{ts} x^{t + sn} + l.o.t.
$$
pourvu que $n$ soit supérieur au degré de $g$ comme polynôme en $x$
(un petit détail est omis). Pour un tel $n$, le polynôme $g(x, -x^n)$ 
est un polynôme non nul en $x$ et s'envoie sur un polynôme non nul de
$\kappa(\mathfrak p)[x]$ pour tout $\mathfrak p \subset R$,
$c_{st} \not \in \mathfrak p$. On conclut que notre affirmation est
vraie pour $V$ égal à l'ouvert principal $D(c_{ts})$ de $\Spec(R)$.

\medskip\noindent
Nous pouvons maintenant raisonner par récurrence sur $\dim(R)$. En effet, soit
$I \subset R$ un idéal tel que $V(I) = \Spec(R) \setminus V$.
Observons que $\dim(R/I) < \dim(R)$, puisque $R$ est un anneau intègre.
Soit $n'_0$ l'entier fourni par l'hypothèse de récurrence sur $\dim(R)$
pour $R/I \to S/IS$ et les images de $a$ et $b$ dans $S/IS$.
Alors $\max(n_0, n'_0)$ convient.
\end{proof}

\begin{lemma}
\label{more-morphisms-lemma-weak-relative-noether-normalization}
Soit $R \to S$ un homomorphisme d'anneaux de type fini. Soit $d$ le maximum
des dimensions des fibres de $\Spec(S) \to \Spec(R)$.
Alors il existe un homomorphisme d'anneaux quasi-fini
$R[t_1, \ldots, t_d] \to S$.
\end{lemma}

\begin{proof}
Dans ce paragraphe, on se ramène au cas où $R \to S$ est de présentation
finie. En effet, écrivons $S = R[x_1, \ldots, x_n]/J$ pour un idéal
$J \subset R[x_1, \ldots, x_n]$.
Alors $J$ est la réunion de ses idéaux de type fini $J_\lambda \subset J$.
Posons $S_\lambda = R[x_1, \ldots, x_n]/J_\lambda$.
Alors, pour un certain $\lambda$, les fibres de
$\Spec(S_\lambda) \to \Spec(R)$
sont de dimension $\leq d$, voir Limites, lemme \ref{limits-lemma-limit-dimension}.
Fixons un tel $\lambda$. Si l'on peut trouver un homomorphisme quasi-fini
$R[t_1, \ldots, t_d] \to S_\lambda$, alors la composée
$R[t_1, \ldots, t_d] \to S$ est bien sûr quasi-finie.
On est ainsi ramené au cas traité dans le paragraphe suivant.

\medskip\noindent
Supposons que $R \to S$ soit de présentation finie. Dans ce paragraphe, on se ramène
au cas où $R$ est de type fini sur $\mathbf{Z}$. D'après
Algèbre, lemme \ref{algebra-lemma-limit-module-finite-presentation},
on peut trouver un ensemble ordonné filtrant $\Lambda$ et un système d'homomorphismes d'anneaux
$R_\lambda \to S_\lambda$ indexé par $\Lambda$, dont la colimite est
$R \to S$, tels que $S_\mu = S_\lambda \otimes_{R_\lambda} R_\mu$
pour $\mu \geq \lambda$ et
que chacun des $R_\lambda$ et $S_\lambda$ soit de type
fini sur $\mathbf{Z}$.
Alors, pour $\lambda$ assez grand, les fibres de
$\Spec(S_\lambda) \to \Spec(R_\lambda)$
sont de dimension $\leq d$, voir
Limites, lemme \ref{limits-lemma-descend-dimension-d}.
Fixons un tel $\lambda$. Si l'on peut trouver un homomorphisme d'anneaux quasi-fini
$R_\lambda[t_1, \ldots, t_d] \to S_\lambda$, alors le changement de base
$R[t_1, \ldots, t_d] \to S$ est lui aussi quasi-fini
(Algèbre, lemme \ref{algebra-lemma-quasi-finite-base-change}).
On est ainsi ramené au cas traité dans le paragraphe suivant.

\medskip\noindent
Supposons que $R$ et $S$ soient de type fini sur $\mathbf{Z}$. Si $d = 0$, alors
l'homomorphisme d'anneaux est quasi-fini et la démonstration est achevée. Supposons $d > 0$.
Nous allons trouver un élément $a \in S$ tel que
les fibres de l'homomorphisme de $R$-algèbres $R[x] \to S$, $x \mapsto a$,
soient de dimension $< d$. Cela achèvera la démonstration par récurrence.

\medskip\noindent
Nous allons démontrer l'existence de $a$ par récurrence sur $\dim(R)$.

\medskip\noindent
Soient $\mathfrak q_1, \ldots, \mathfrak q_r \subset S$ ceux
des idéaux premiers minimaux de $S$ tels que $\dim_{\mathfrak q_i}(S/R) = d$.
Pour la notation, voir
Algèbre, définition \ref{algebra-definition-relative-dimension}.
Disons que $\mathfrak q_i$ est au-dessus de l'idéal premier $\mathfrak p_i \subset R$.
On a $\text{trdeg}_{\kappa(\mathfrak p_i)}(\kappa(\mathfrak q_i)) = d$,
car $\mathfrak q_i$ est un point générique de sa fibre ; on peut par exemple
appliquer Algèbre, lemme \ref{algebra-lemma-dimension-at-a-point-finite-type-field},
à $S \otimes_R \kappa(\mathfrak p_i)$.
Ainsi, d'après le lemme \ref{more-morphisms-lemma-silly}, on peut trouver un élément $a \in S$ tel que
l'image de $a$ dans $\kappa(\mathfrak q_i)$ soit transcendante sur
$\kappa(\mathfrak p_i)$ pour $i = 1, \ldots, r$.
Considérons le morphisme
$$
f_a : \Spec(S) \longrightarrow \mathbf{A}^1_R
$$
correspondant à l'homomorphisme de $R$-algèbres $R[x] \to S$
qui envoie $x$ sur $a$. Soit $U \subset \Spec(S)$ l'ouvert
où les fibres sont de dimension $\leq d - 1$, voir
Morphismes, lemme \ref{morphisms-lemma-openness-bounded-dimension-fibres}.
Par construction, $U$ contient tous les points génériques de $\Spec(S)$.
En particulier, on voit que $U$ contient tous les points génériques de
toutes les fibres génériques de $\Spec(S) \to \Spec(R)$, puisque de tels points sont
nécessairement des points génériques de $\Spec(S)$. Posons $T = \Spec(S) \setminus U$,
considéré comme un sous-schéma fermé réduit de $\Spec(S)$.
Il résulte de ce que l'on vient de dire et de l'hypothèse
$\dim(S/R) \leq d$ que
les fibres génériques de $T \to \Spec(R)$ sont de dimension $\leq d - 1$.
Ainsi, d'après le lemme \ref{more-morphisms-lemma-dimension-in-neighbourhood},
appliqué plusieurs fois pour produire des voisinages ouverts des points
génériques de $\Spec(R)$, on peut trouver un ouvert dense $V \subset \Spec(R)$ tel que
$T_V \to V$ ait des fibres de dimension $\leq d - 1$.
On conclut que, pour $q \in \mathbf{A}^1_V$, la fibre de
$f_a$ au-dessus de $q$ est de dimension $\leq d - 1$ (car on a borné la dimension
de la fibre de $f_a|_U$ et de celle de $f_a|_T$).

\medskip\noindent
Par évitement des idéaux premiers, on peut supposer que $V = D(f)$ pour un certain $f \in R$.
On voit alors que l'homomorphisme d'anneaux $R_f[x] \to S_f$, $x \mapsto a$,
a des fibres de dimension $\leq d - 1$.
On peut remplacer $a$ par $fa$ et supposer $a \in (f)$.
Par récurrence sur $\dim(R)$, on peut trouver un élément
$\overline{b} \in S/fS$
tel que les fibres de $\Spec(S/fS) \to \Spec(R/fR[x])$,
$x \mapsto \overline{b}$, soient de dimension $\leq d - 1$.
Soit $b \in S$ un relèvement de $\overline{b}$.
D'après le lemme \ref{more-morphisms-lemma-sum-is-ok}, il existe un $n > 0$ tel
que $a^n + b$ convienne encore pour $R_f \to S_f$.
D'autre part, l'image de $a^n + b$ dans $S/fS$ est $\overline{b}$,
et la démonstration est achevée.
\end{proof}

\section{Théorèmes de Bertini}
\label{more-morphisms-section-bertini}

\noindent
Nous poursuivons la discussion commencée dans
Variétés, section \ref{varieties-section-bertini}.
Dans cette section, nous démontrons que les sections hyperplanes générales
des variétés géométriquement irréductibles sont géométriquement irréductibles,
en suivant l'argument remarquable donné dans \cite{Jou}.

\begin{lemma}
\label{more-morphisms-lemma-amazing-bertini-lemma}
\begin{reference}
Voir les pages 71 et 72 de \cite{Jou}
\end{reference}
Soit $K/k$ une extension de corps géométriquement irréductible et de type fini.
Soit $n \geq 1$.
Soient $g_1, \ldots, g_n \in K$ des éléments tels qu'il
existe $c_1, \ldots, c_n \in k$ pour lesquels les éléments
$$
x_1, \ldots, x_n, \sum g_ix_i, \sum c_ig_i \in K(x_1, \ldots, x_n)
$$
soient algébriquement indépendants sur $k$. Alors
$K(x_1, \ldots, x_n)$ est géométriquement irréductible sur
$k(x_1, \ldots, x_n, \sum g_ix_i)$.
\end{lemma}

\begin{proof}
Soient $c_1, \ldots, c_n \in k$ comme dans l'énoncé du lemme.
Écrivons $\xi = \sum g_ix_i$ et $\delta = \sum c_ig_i$.
Pour $a \in k$, considérons l'automorphisme $\sigma_a$
de $K(x_1, \ldots, x_n)$ qui est l'identité sur $K$ et qui est donné par les règles
$$
\sigma_a(x_i) = x_i + a c_i
$$
Observons que $\sigma_a(\xi) = \xi + a \delta$ et $\sigma_a(\delta) = \delta$.
Considérons la tour d'extensions de corps
$$
K_0 = k(x_1, \ldots, x_n) \subset
K_1 = K_0(\xi) \subset
K_2 = K_0(\xi, \delta) \subset K(x_1, \ldots, x_n) = \Omega
$$
Observons que $\sigma_a(K_0) = K_0$ et $\sigma_a(K_2) = K_2$.
Soit $\theta \in \Omega$ un élément algébrique séparable sur $K_1$.
Nous devons montrer que $\theta \in K_1$, voir Algèbre, lemme
\ref{algebra-lemma-geometrically-irreducible-separable-elements}.

\medskip\noindent
Notons $K'_2$ la fermeture algébrique séparable de $K_2$ dans $\Omega$.
Comme $K'_2/K_2$ est finie (Algèbre, lemme
\ref{algebra-lemma-make-geometrically-irreducible}) et
séparable, il n'y a qu'un nombre fini de corps intermédiaires entre
$K'_2$ et $K_2$ (Corps, lemme \ref{fields-lemma-primitive-element}).
Si $k$ est infini\footnote{Nous traiterons le cas d'un corps fini dans le
dernier paragraphe de la démonstration.}, on peut alors trouver des éléments distincts
$a_1, a_2$ de $k$ tels que
$$
K_2(\sigma_{a_1}(\theta)) = K_2(\sigma_{a_2}(\theta))
$$
comme sous-corps de $\Omega$. Écrivons $\theta_i = \sigma_{a_i}(\theta)$
et $\xi_i = \sigma_{a_i}(\xi) = \xi + a_i \delta$. Observons que
$$
K_2 = K_0(\xi_1, \xi_2)
$$
car on a $\xi_i = \xi + a_i \delta$,
$\xi = (a_2 \xi_1 - a_1 \xi_2)/(a_2 - a_1)$, et
$\delta = (\xi_1 - \xi_2)/(a_1 - a_2)$.
Puisque $K_2/K_0$ est purement transcendante de degré $2$, on conclut
que $\xi_1$ et $\xi_2$ sont algébriquement indépendants sur $K_0$.
Comme $\theta_1$ est algébrique sur $K_0(\xi_1)$, on conclut que
$\xi_2$ est transcendant sur $K_0(\xi_1, \theta_1)$.

\medskip\noindent
Par hypothèse, $K/k$ est géométriquement irréductible. Cela implique
que $K(x_1, \ldots, x_n)/K_0$ est géométriquement irréductible
(Algèbre, lemme
\ref{algebra-lemma-geometrically-irreducible-base-change-transcendental}).
Cela implique à son tour que $K_0(\xi_1, \theta_1)/K_0$
est géométriquement irréductible en tant que sous-extension
(Algèbre, lemme \ref{algebra-lemma-subalgebra-geometrically-irreducible}).
Puisque $\xi_2$ est transcendant sur $K_0(\xi_1, \theta_1)$,
on conclut que $K_0(\xi_1, \xi_2, \theta_1)/K_0(\xi_2)$
est géométriquement irréductible (Algèbre, lemme
\ref{algebra-lemma-geometrically-irreducible-add-transcendental}).
Par notre choix de $a_1, a_2$ ci-dessus, on a
$$
K_0(\xi_1, \xi_2, \theta_1) =
K_2(\sigma_{a_1}(\theta)) =
K_2(\sigma_{a_2}(\theta)) =
K_0(\xi_1, \xi_2, \theta_2)
$$
Comme $\theta_2$ est algébrique séparable sur $K_0(\xi_2)$,
on conclut à nouveau, d'après Algèbre, lemme
\ref{algebra-lemma-geometrically-irreducible-separable-elements}, que
$\theta_2 \in K_0(\xi_2)$. En appliquant $\sigma_{a_2}^{-1}$
à cette relation, on obtient $\theta \in K_0(\xi) = K_1$, comme souhaité.

\medskip\noindent
Cela achève la démonstration lorsque $k$ est infini. Si $k$ est fini,
on peut choisir une variable $t$ et considérer l'extension
$K(t)/k(t)$, qui est géométriquement irréductible d'après
Algèbre, lemme
\ref{algebra-lemma-geometrically-irreducible-base-change-transcendental}.
Comme il reste vrai que
$x_1, \ldots, x_n, \sum g_ix_i, \sum c_ig_i$
dans $K(t, x_1, \ldots, x_n)$ sont algébriquement indépendants sur $k(t)$,
on conclut que $K(t, x_1, \ldots, x_n)$
est géométriquement irréductible sur
$k(t, x_1, \ldots, x_n, \sum g_ix_i)$
par l'argument déjà donné.
Il suffit alors d'utiliser encore une fois Algèbre, lemme
\ref{algebra-lemma-geometrically-irreducible-base-change-transcendental},
pour achever la démonstration.
\end{proof}

\begin{lemma}
\label{more-morphisms-lemma-algebraically-independent}
Soit $A$ un anneau intègre de type fini sur un corps $k$. Soit $n \geq 2$.
Soient $g_1, \ldots, g_n \in A$ des éléments tels que $V(g_1, g_2)$
ait une composante irréductible de dimension $\dim(A) - 2$.
Alors il existe $c_1, \ldots, c_n \in k$ tels que les éléments
$$
x_1, \ldots, x_n, \sum g_ix_i, \sum c_ig_i \in
\text{Frac}(A)(x_1, \ldots, x_n)
$$
soient algébriquement indépendants sur $k$.
\end{lemma}

\begin{proof}
L'indépendance algébrique sur $k$ signifie que le morphisme
$$
T = \Spec(A[x_1, \ldots, x_n])
\longrightarrow
\Spec(k[x_1, \ldots, x_n, y, z]) = S
$$
donné par $y = \sum g_ix_i$ et $z = \sum c_ig_i$ est dominant.
Posons $d = \dim(A)$. Si $T \to S$ n'est pas dominant, alors son image
est de dimension $< n + 2$, et toute composante irréductible
de toute fibre est donc de dimension $> d + n - (n + 2) = d - 2$, voir
Variétés, lemme
\ref{varieties-lemma-dimension-fibres-locally-algebraic}.
Choisissons un point fermé $u \in V(g_1, g_2)$ contenu dans une composante irréductible
de dimension $d - 2$ et dans aucune autre composante de $V(g_1, g_2)$.
Considérons le point fermé $t = (u, 1, 0, \ldots 0)$
de $T$ situé au-dessus de $u$. Posons $(c_1, \ldots, c_n) = (0, 1, 0, \ldots, 0)$.
Alors $t$ s'envoie sur le point $s = (1, 0, \ldots, 0)$ de $S$.
La fibre de $T \to S$ au-dessus de $s$ est définie par
$$
x_1 - 1, x_2, \ldots, x_n, \sum x_ig_i, g_2
$$
et est donc, de manière équivalente, définie par
$$
x_1 - 1, x_2, \ldots, x_n, g_1, g_2
$$
D'après notre condition sur $g_1, g_2$, ce sous-schéma possède une composante irréductible
de dimension $d - 2$.
\end{proof}

\begin{lemma}
\label{more-morphisms-lemma-bertini-irreducible}
\begin{reference}
\cite[théorème 6.3, partie 4)]{Jou}
\end{reference}
Dans Variétés, situation \ref{varieties-situation-family-divisors}, supposons que
\begin{enumerate}
\item $X$ soit de type fini sur $k$,
\item $X$ soit géométriquement irréductible sur $k$,
\item il existe $v_1, v_2, v_3 \in V$ et une composante irréductible
$Z$ de $H_{v_2} \cap H_{v_3}$ telles que $Z \not \subset H_{v_1}$ et
$\text{codim}(Z, X) = 2$, et
\item toute composante irréductible $Y$ de $\bigcap_{v \in V} H_v$
ait $\text{codim}(Y, X) \geq 2$.
\end{enumerate}
Alors, pour $v \in V \otimes_k k'$ général,
le schéma $H_v$ est géométriquement irréductible sur $k'$.
\end{lemma}

\begin{proof}
Pour que l'hypothèse (3) soit satisfaite, les éléments $v_1, v_2, v_3$
doivent être linéairement indépendants sur $k$ dans $V$ (un petit détail est omis).
On peut donc choisir une base $v_1, \ldots, v_r$ de $V$ contenant
ces éléments comme ses $3$ premiers éléments. Rappelons que
$H_{univ} \subset \mathbf{A}^r_k \times_k X$ est le « diviseur universel ».
Considérons la projection $q : H_{univ} \to \mathbf{A}^r_k$
dont les fibres schématiques sont les diviseurs $H_v$. D'après le
lemme \ref{more-morphisms-lemma-geom-irreducible-generic-fibre},
il suffit de montrer que la fibre générique de $q$ est géométriquement
irréductible.
Pour cela, on peut remplacer $X$ par son réduit ; on peut donc
supposer que $X$ est un schéma intègre de type fini sur $k$.

\medskip\noindent
Soit $U \subset X$ un ouvert affine non vide tel que
$\mathcal{L}|_U \cong \mathcal{O}_U$. Écrivons $U = \Spec(A)$.
Notons $f_i \in A$ l'élément correspondant à la section $\psi(v_i)|_U$
via l'isomorphisme $\mathcal{L}|_U \cong \mathcal{O}_U$.
Alors $H_{univ} \cap (\mathbf{A}^r_k \times_k U)$ est donné par
$$
H_U = \Spec(A[x_1, \ldots, x_r]/(x_1f_1 + \ldots + x_rf_r))
$$
Par notre choix de la base, on voit que $f_1$ ne peut pas être nul, car cela
signifierait que $v_1 = 0$ et donc que $H_{v_1} = X$, ce qui contredit
l'hypothèse (3). Ainsi $\sum x_if_i$ est un non-diviseur de zéro dans
$A[x_1, \ldots, x_r]$.
Il s'ensuit que toute composante irréductible de $H_U$ est de dimension
$d + r - 1$, où $d = \dim(X) = \dim(A)$. Si $U' = U \cap D(f_1)$,
alors on voit que
$$
H_{U'} =
\Spec(A_{f_1}[x_1, \ldots, x_r]/(x_1f_1 + \ldots + x_rf_r)) \cong
\Spec(A_{f_1}[x_2, \ldots x_r]) =
\mathbf{A}^{r - 1}_k \times_k U'
$$
est irréductible. D'autre part, on a
$$
H_U \setminus H_{U'} =
\Spec(A/(f_1)[x_1, \ldots, x_r]/(x_2f_2 + \ldots + x_rf_r))
$$
qui est de dimension au plus $d + r - 2$. En effet, pour $i \not = 1$,
le schéma $(H_U \setminus H_{U'}) \times_U D(f_i)$ est soit vide
(si $f_i = 0$), soit, par le même argument que ci-dessus, isomorphe à un espace affine de dimension $r - 1$
sur un ouvert de $\Spec(A/(f_1))$, et est donc de
dimension au plus $d + r - 2$.
D'autre part, $(H_U \setminus H_{U'}) \times_U V(f_2, \ldots, f_r)$
est un espace affine de dimension $r$ sur $\Spec(A/(f_1, \ldots, f_r))$ ;
l'hypothèse (4) nous dit donc qu'il est de dimension au plus
$d + r - 2$. On conclut que $H_U$ est irréductible pour tout $U$ comme ci-dessus.
Il s'ensuit que $H_{univ}$ est irréductible.

\medskip\noindent
Il suffit donc de montrer que le point générique de $H_{univ}$
est géométriquement irréductible sur le point générique de $\mathbf{A}^r_k$, voir
Variétés, lemme
\ref{varieties-lemma-geometrically-irreducible-function-field}.
Choisissons un ouvert affine non vide $U = \Spec(A)$ de $X$,
contenu dans $X \setminus H_{v_1}$ et qui rencontre la composante
irréductible $Z$ de $H_{v_2} \cap H_{v_3}$ dont l'existence est affirmée dans
l'hypothèse (3). Avec les notations ci-dessus, il faut démontrer que
l'extension de corps
$$
\text{Frac}(A[x_1, \ldots, x_r]/(x_1f_1 + \ldots + x_rf_r))/
k(x_1, \ldots , x_r)
$$
est géométriquement irréductible. Observons que $f_1$ est inversible dans $A$
par notre choix de $U$.
Posons $K = \text{Frac}(A)$, le corps des fractions de $A$.
En éliminant la variable $x_1$ comme ci-dessus,
on trouve qu'il faut montrer que l'extension de corps
$$
K(x_2, \ldots, x_r)/
k(x_2, \ldots, x_r, -\sum\nolimits_{i = 2, \ldots, r} f_1^{-1}f_i x_i)
$$
est géométriquement irréductible. D'après le lemme \ref{more-morphisms-lemma-amazing-bertini-lemma},
il suffit de montrer que, pour certains $c_2, \ldots, c_r \in k$, les éléments
$$
x_2, \ldots, x_r, \sum\nolimits_{i = 2, \ldots, r} f_1^{-1}f_i x_i,
\sum\nolimits_{i = 2, \ldots, r} c_if_1^{-1}f_i
$$
sont algébriquement indépendants sur $k$ dans le corps des fractions de
$A[x_2, \ldots, x_r]$. Cela résulte du
lemme \ref{more-morphisms-lemma-algebraically-independent}
et du fait que $Z \cap U$ est une composante irréductible
de $V(f_1^{-1}f_2, f_1^{-1}f_3) \subset U$.
\end{proof}

\begin{remark}
\label{more-morphisms-remark-geometric-proof-bertini-irreducible}
Esquissons une démonstration « géométrique » d'un cas particulier du
lemme \ref{more-morphisms-lemma-bertini-irreducible}. Supposons que $k$
soit un corps algébriquement clos et que $X \subset \mathbf{P}^n_k$ soit lisse et
irréductible de dimension $\geq 2$. Nous affirmons qu'il existe un
hyperplan $H \subset \mathbf{P}^n_k$ tel que $X \cap H$ soit
lisse et irréductible. En effet, d'après
Variétés, lemme \ref{varieties-lemma-bertini},
pour $v \in V = kT_0 \oplus \ldots \oplus kT_n$ général,
la section hyperplane correspondante $X \cap H_v$ est lisse.
D'autre part, d'après Enriques--Severi--Zariski, le schéma
$X \cap H_v$ est connexe, voir
Variétés, lemme \ref{varieties-lemma-connectedness-ample-divisor}.
Ainsi $X \cap H_v$ est lisse et irréductible.
\end{remark}

\section{Théorème du cube}
\label{more-morphisms-section-theorem-cube}

\noindent
Le lemme suivant nous dit que la diagonale du foncteur de Picard
est représentable par des immersions localement fermées sous
les hypothèses faites dans le lemme.

\begin{lemma}
\label{more-morphisms-lemma-diagonal-picard-flat-proper}
Soit $f : X \to S$ un morphisme plat et propre de présentation finie.
Soit $\mathcal{E}$ un $\mathcal{O}_X$-module localement libre de rang fini.
Pour un morphisme $g : T \to S$, considérons le diagramme de changement de base
$$
\xymatrix{
X_T \ar[d]_p \ar[r]_q & X \ar[d]^f \\
T \ar[r]^g & S
}
$$
Supposons que $\mathcal{O}_T \to p_*\mathcal{O}_{X_T}$ soit un
isomorphisme pour tout $g : T \to S$. Alors il existe une
immersion $j : Z \to S$ de présentation finie telle qu'un
morphisme $g : T \to S$ se factorise par $Z$ si et seulement s'il
existe un $\mathcal{O}_T$-module localement libre de rang fini $\mathcal{N}$
tel que $p^*\mathcal{N} \cong q^*\mathcal{E}$.
\end{lemma}

\begin{proof}
Observons que les fibres $X_s$ de $f$ sont connexes par notre hypothèse
$H^0(X_s, \mathcal{O}_{X_s}) = \kappa(s)$. Le rang de
$\mathcal{E}$ est donc constant sur les fibres. Puisque $f$ est ouvert
(Morphismes, lemme \ref{morphisms-lemma-fppf-open}) et fermé,
on conclut qu'il existe une décomposition $S = \coprod S_r$
de $S$ en sous-schémas ouverts et fermés telle que $\mathcal{E}$
soit de rang constant $r$ sur l'image réciproque de $S_r$.
On peut donc supposer que $\mathcal{E}$ est de rang constant $r$.
On notera $\mathcal{E}^\vee = \SheafHom(\mathcal{E}, \mathcal{O}_X)$
le module dual, de rang $r$.

\medskip\noindent
D'après la cohomologie et le changement de base (plus précisément, d'après
Catégories dérivées de schémas, lemme
\ref{perfect-lemma-flat-proper-perfect-direct-image-general}),
on voit que $E = Rf_*\mathcal{E}$ est un objet parfait de la
catégorie dérivée de $S$ et que sa formation commute au
changement de base arbitraire. Il en va de même pour $E' = Rf_*\mathcal{E}^\vee$.
Comme il n'y a jamais de cohomologie en degrés $< 0$, on voit que
$E$ et $E'$ ont (localement) une amplitude de Tor dans $[0, b]$ pour un certain $b$.
Observons que, pour tout $g : T \to S$, on a
$p_*(q^*\mathcal{E}) = H^0(Lg^*E)$ et
$p_*(q^*\mathcal{E}^\vee) = H^0(Lg^*E')$.
Soient $j : Z \to S$ et $j' : Z' \to S$ les immersions
de présentation finie construites dans Catégories dérivées de schémas, lemme
\ref{perfect-lemma-locally-closed-where-H0-locally-free},
pour $E$ et $E'$, avec $a = 0$ et $r = r$ ; elles sont, en gros,
caractérisées par la propriété que $H^0(Lj^*E)$ et $H^0((j')^*E')$
sont des modules localement libres de rang fini et que la formation de ces modules commute au changement de base.

\medskip\noindent
Soit $g : T \to S$ un morphisme. S'il existe un $\mathcal{N}$
comme dans le lemme, alors, en utilisant la formule de projection
Cohomologie, lemme \ref{cohomology-lemma-projection-formula},
on voit que les modules
$$
p_*(q^*\mathcal{E}) \cong
p_*(p^*\mathcal{N}) \cong
\mathcal{N} \otimes_{\mathcal{O}_T} p_*\mathcal{O}_{X_T} \cong
\mathcal{N}\quad\text{et de même }\quad
p_*(q^*\mathcal{E}^\vee) \cong \mathcal{N}^\vee
$$
sont des modules localement libres de rang $r$
et le restent après tout changement de base ultérieur $T' \to T$.
Dans ce cas, $T \to S$ se factorise donc par $j$ et par $j'$.
On peut ainsi remplacer $S$ par $Z \times_S Z'$ et supposer que
$f_*\mathcal{E}$ et $f_*\mathcal{E}^\vee$ sont des
$\mathcal{O}_S$-modules localement libres de rang $r$
dont la formation commute au changement de base arbitraire
(un petit détail est omis).

\medskip\noindent
Dans cette situation, si $g : T \to S$ est un morphisme et s'il existe un
$\mathcal{N}$ comme dans le lemme, alors l'application (le cup-produit en degré $0$)
$$
p_*(q^*\mathcal{E})
\otimes_{\mathcal{O}_T}
p_*(q^*\mathcal{E}^\vee)
\longrightarrow \mathcal{O}_T
$$
est un accouplement parfait. Réciproquement, si cette application de cup-produit est un
accouplement parfait, on voit que, localement sur $T$, on peut choisir une
base de sections
$\sigma_1, \ldots, \sigma_r$ dans $p_*(q^*\mathcal{E})$
et $\tau_1, \ldots, \tau_r$ dans $p_*(q^*\mathcal{E}^\vee)$
dont les produits vérifient $\sigma_i \tau_j = \delta_{ij}$.
En considérant $\sigma_i$ comme une section de $q^*\mathcal{E}$ sur $X_T$
et $\tau_j$ comme une section de $q^*\mathcal{E}^\vee$ sur $X_T$,
on conclut que
$$
\sigma_1, \ldots, \sigma_r :
\mathcal{O}_{X_T}^{\oplus r}
\longrightarrow
q^*\mathcal{E}
$$
est un isomorphisme dont l'inverse est donné par
$$
\tau_1, \ldots, \tau_r :
q^*\mathcal{E}
\longrightarrow
\mathcal{O}_{X_T}^{\oplus r}
$$
Autrement dit, on voit que $p^*p_*q^*\mathcal{E} \cong q^*\mathcal{E}$.
Or, la condition que le cup-produit soit non dégénéré définit
un sous-schéma ouvert rétrocompact (à savoir, le lieu où un déterminant
convenable est non nul), et la démonstration est achevée.
\end{proof}

\noindent
Le lemme précédent nous dit en particulier que, si un fibré vectoriel est
trivial sur les fibres d'une famille propre et plate d'espaces propres, alors
il est l'image réciproque d'un fibré vectoriel. Détaillons un peu cela.

\begin{lemma}
\label{more-morphisms-lemma-trivial-on-fibres}
Soit $f : X \to S$ un morphisme plat et propre de présentation finie
tel que $f_*\mathcal{O}_X = \mathcal{O}_S$, et que cela reste
vrai après tout changement de base. Soit $\mathcal{E}$ un
$\mathcal{O}_X$-module localement libre de rang fini. Supposons que
\begin{enumerate}
\item $\mathcal{E}|_{X_s}$ soit isomorphe à
$\mathcal{O}_{X_s}^{\oplus r_s}$ pour tout $s \in S$, et
\item $S$ soit réduit.
\end{enumerate}
Alors $\mathcal{E} = f^*\mathcal{N}$ pour un certain
$\mathcal{O}_S$-module localement libre de rang fini $\mathcal{N}$.
\end{lemma}

\begin{proof}
En effet, dans ce cas, l'immersion localement fermée $j : Z \to S$ du
lemme \ref{more-morphisms-lemma-diagonal-picard-flat-proper}
est bijective, et est donc une
immersion fermée. Mais puisque $S$ est réduit, $j$ est un isomorphisme.
\end{proof}

\begin{lemma}
\label{more-morphisms-lemma-triviality-generic-fibre-valuation-ring}
Soit $f : X \to S$ un morphisme propre et plat de présentation finie.
Soit $\mathcal{L}$ un $\mathcal{O}_X$-module inversible.
Supposons que
\begin{enumerate}
\item $S$ soit le spectre d'un anneau de valuation,
\item $\mathcal{L}$ soit trivial sur la fibre générique $X_\eta$ de $f$,
\item la fibre fermée $X_0$ de $f$ soit intègre,
\item $H^0(X_\eta, \mathcal{O}_{X_\eta})$ soit égal au corps des fonctions de $S$.
\end{enumerate}
Alors $\mathcal{L}$ est trivial.
\end{lemma}

\begin{proof}
Écrivons $S = \Spec(A)$. Nous allons d'abord démontrer le lemme lorsque $A$ est un
anneau de valuation discrète (car c'est le cas le plus souvent utilisé en pratique).
Soit $\pi \in A$ une uniformisante.
Prenons une section trivialisante $s \in \Gamma(X_\eta, \mathcal{L}_\eta)$.
Après avoir remplacé $s$ par $\pi^n s$ si nécessaire,
on peut supposer que $s \in \Gamma(X, \mathcal{L})$.
Si $s|_{X_0} = 0$, on voit alors
que $s$ est divisible par $\pi$ (parce que $X_0$ est la fibre
schématique et que $X$ est plat sur $A$). On peut donc supposer que $s|_{X_0}$
est non nul. Le lieu des zéros $Z(s)$ de $s$ est alors contenu dans $X_0$,
mais ne contient pas le point générique de $X_0$ (parce que $X_0$ est intègre).
Cela signifie que $Z(s)$ est de codimension $\geq 2$ dans $X$, ce qui contredit
Diviseurs, lemme \ref{divisors-lemma-effective-Cartier-codimension-1},
à moins que $Z(s) = \emptyset$, comme souhaité.

\medskip\noindent
Démonstration dans le cas général. Puisque l'anneau de valuation $A$ est
cohérent (Algèbre, exemple \ref{algebra-example-valuation-ring-coherent}),
on voit que $H^0(X, \mathcal{L})$ est un $A$-module cohérent, voir
Catégories dérivées de schémas, lemme
\ref{perfect-lemma-cohomology-over-coherent-ring}.
De manière équivalente, $H^0(X, \mathcal{L})$ est un
$A$-module de présentation finie (Algèbre, lemme \ref{algebra-lemma-coherent-ring}).
Puisque $H^0(X, \mathcal{L})$ est sans torsion (par platitude de $X$ sur $A$),
on déduit de Plus sur l'algèbre, lemme
\ref{more-algebra-lemma-generalized-valuation-ring-modules},
que $H^0(X, \mathcal{L}) = A^{\oplus n}$ pour un certain $n$.
Par changement de base plat (Cohomologie des schémas, lemme
\ref{coherent-lemma-flat-base-change-cohomology}),
on a
$$
K = H^0(X_\eta, \mathcal{O}_{X_\eta}) \cong
H^0(X_\eta, \mathcal{L}_\eta) =
H^0(X, \mathcal{L}) \otimes_A K
$$
où $K$ est le corps des fractions de $A$. Ainsi $n = 1$.
Choisissons un générateur $s \in H^0(X, \mathcal{L})$.
Soit $\mathfrak m \subset A$ l'idéal maximal.
Alors $\kappa = A/\mathfrak m = \colim A/\pi$, où
il s'agit de la colimite filtrante sur les $\pi \in \mathfrak m$ non nuls
(on utilise ici le fait que $A$ est un anneau de valuation).
Ainsi $X_0 = \lim X \times_S \Spec(A/\pi)$.
Si $s|_{X_0}$ est nul, alors, pour un certain $\pi$,
on voit que la restriction de $s$ à $X \times_S \Spec(A/\pi)$ est nulle, voir
Limites, lemme \ref{limits-lemma-descend-section}.
Mais si cela se produit, alors $\pi^{-1} s$ est
une section globale de $\mathcal{L}$, ce qui contredit
le fait que $s$ soit un générateur de $H^0(X, \mathcal{L})$.
Ainsi $s|_{X_0}$ n'est pas nul. Soit $Z(s) \subset X$ le schéma des zéros de $s$.
Puisque $s|_{X_0}$ n'est pas nul et que $X_0$ est intègre,
on voit que $Z(s)_0 \subset X_0$ est un diviseur de Cartier effectif.
Puisque $f$ est propre et que $S$ est local, tout point de $Z(s)$
se spécialise en un point de $Z(s)_0$. Ainsi, d'après
Diviseurs, lemme \ref{divisors-lemma-fibre-Cartier}, partie (3),
on voit que $Z(s)$ est un diviseur de Cartier effectif relatif ;
en particulier, $Z(s) \to S$ est plat.
Par conséquent, si $Z(s)$ était non vide, alors $Z(s)_\eta$ serait non vide,
ce qui contredit le fait que $s|_{X_\eta}$ soit une trivialisation
de $\mathcal{L}_\eta$. Ainsi $Z(s) = \emptyset$, comme souhaité.
\end{proof}

\begin{lemma}
\label{more-morphisms-lemma-get-a-closed}
Soient $f : X \to S$ et $\mathcal{E}$ comme dans le
lemme \ref{more-morphisms-lemma-diagonal-picard-flat-proper},
et supposons en outre que $\mathcal{E}$ soit un $\mathcal{O}_X$-module inversible.
Si, de plus, les fibres géométriques de $f$ sont
intègres, alors $Z$ est fermé dans $S$.
\end{lemma}

\begin{proof}
Puisque $j : Z \to S$ est de présentation finie, il suffit
de montrer ceci : pour tout morphisme $g : \Spec(A) \to S$, où $A$ est un
anneau de valuation de corps des fractions $K$, tel que $g(\Spec(K)) \in j(Z)$,
on a $g(\Spec(A)) \subset j(Z)$. Voir
Morphismes, lemme \ref{morphisms-lemma-reach-points-scheme-theoretic-image}.
Cela résulte du lemme \ref{more-morphisms-lemma-triviality-generic-fibre-valuation-ring}
et de la caractérisation de $j : Z \to S$ dans le
lemme \ref{more-morphisms-lemma-diagonal-picard-flat-proper}.
\end{proof}

\begin{lemma}
\label{more-morphisms-lemma-H1-O-picard-flat-proper}
Considérons un diagramme commutatif de schémas
$$
\xymatrix{
X' \ar[rr] \ar[dr]_{f'} & & X \ar[dl]^f \\
& S
}
$$
où $f' : X' \to S$ et $f : X \to S$ satisfont les hypothèses du
lemme \ref{more-morphisms-lemma-diagonal-picard-flat-proper}.
Soit $\mathcal{L}$ un $\mathcal{O}_X$-module inversible,
et soit $\mathcal{L}'$ son image réciproque sur $X'$. Soit $Z \subset S$,
resp.\ $Z' \subset S$, le sous-schéma localement fermé construit
dans le lemme \ref{more-morphisms-lemma-diagonal-picard-flat-proper}
pour $(f, \mathcal{L})$, resp.\ $(f', \mathcal{L}')$,
de sorte que $Z \subset Z'$. Si $s \in Z$ et si
$$
H^1(X_s, \mathcal{O}) \longrightarrow H^1(X'_s, \mathcal{O})
$$
est injectif, alors $Z \cap U = Z' \cap U$ pour un certain voisinage ouvert
$U$ de $s$.
\end{lemma}

\begin{proof}
On peut remplacer $S$ par $Z'$. Après avoir remplacé $S$ par un voisinage ouvert affine
de $s$, on peut supposer que $\mathcal{L}' = \mathcal{O}_{X'}$.
Soient $E = Rf_*\mathcal{L}$ et $E' = Rf'_*\mathcal{L}' = Rf'_*\mathcal{O}_{X'}$.
Ce sont des complexes parfaits dont la formation commute à tout
changement de base (Catégories dérivées de schémas, lemme
\ref{perfect-lemma-flat-proper-perfect-direct-image-general}).
En particulier, on voit que
$$
E \otimes_{\mathcal{O}_S}^\mathbf{L} \kappa(s) =
R\Gamma(X_s, \mathcal{L}_s) = R\Gamma(X_s, \mathcal{O}_{X_s})
$$
La seconde égalité vient de ce que $s \in Z$. Posons
$h_i = \dim_{\kappa(s)} H^i(X_s, \mathcal{O}_{X_s})$.
Après avoir rétréci $S$, on peut représenter $E$ par un complexe
$$
\mathcal{O}_S \to \mathcal{O}_S^{\oplus h_1} \to
\mathcal{O}_S^{\oplus h_2} \to \ldots
$$
voir Plus sur l'algèbre, lemme
\ref{more-algebra-lemma-lift-perfect-from-residue-field}
(à strictement parler, cela utilise aussi
Catégories dérivées de schémas, lemmes
\ref{perfect-lemma-affine-compare-bounded} et
\ref{perfect-lemma-perfect-affine}). De même, on peut supposer que $E'$
est représenté par un complexe
$$
\mathcal{O}_S \to \mathcal{O}_S^{\oplus h'_1} \to
\mathcal{O}_S^{\oplus h'_2} \to \ldots
$$
où $h'_i = \dim_{\kappa(s)} H^i(X'_s, \mathcal{O}_{X'_s})$.
Par fonctorialité de la cohomologie, on dispose d'un morphisme
$$
E \longrightarrow E'
$$
dans $D(\mathcal{O}_S)$ dont la formation commute à tout changement de base.
Puisque le complexe représentant $E$ est un complexe fini de modules libres
de type fini et que $S$ est affine, on peut choisir un morphisme de complexes
$$
\xymatrix{
\mathcal{O}_S \ar[r]_d \ar[d]_a &
\mathcal{O}_S^{\oplus h_1} \ar[r] \ar[d]_b &
\mathcal{O}_S^{\oplus h_2} \ar[r] \ar[d]_c & \ldots \\
\mathcal{O}_S \ar[r]^{d'} &
\mathcal{O}_S^{\oplus h'_1} \ar[r] &
\mathcal{O}_S^{\oplus h'_2} \ar[r] & \ldots
}
$$
représentant le morphisme donné $E \to E'$. Puisque $s \in Z$, on voit que
la section trivialisante de $\mathcal{L}_s$ donne par image réciproque une section trivialisante
de $\mathcal{L}'_s = \mathcal{O}_{X'_s}$. Ainsi
$a \otimes \kappa(s)$ est un isomorphisme ; après avoir rétréci $S$,
on voit donc que $a$ est un isomorphisme. Enfin, on utilise l'hypothèse
que $H^1(X_s, \mathcal{O}) \to H^1(X'_s, \mathcal{O})$
est injectif pour voir qu'il existe un mineur $h_1 \times h_1$ de la
matrice définissant $b$ dont l'image est un élément non nul
de $\kappa(s)$. Après avoir rétréci $S$, on peut donc supposer
que $b$ est injectif. Or, puisque $\mathcal{L}' = \mathcal{O}_{X'}$,
on voit que $d' = 0$. Il s'ensuit que $d = 0$. On voit ainsi
que la section trivialisante de $\mathcal{L}_s$ se relève en une section
de $\mathcal{L}$ sur $X$. Un argument topologique immédiat (omis)
montre que cela signifie que $\mathcal{L}$ est trivial après avoir éventuellement
encore un peu rétréci $S$.
\end{proof}

\begin{lemma}
\label{more-morphisms-lemma-H1-O-multiple-picard-flat-proper}
Considérons $n$ diagrammes commutatifs de schémas
$$
\xymatrix{
X_i \ar[rr] \ar[dr]_{f_i} & & X \ar[dl]^f \\
& S
}
$$
où $f_i : X_i \to S$ et $f : X \to S$ satisfont les hypothèses du
lemme \ref{more-morphisms-lemma-diagonal-picard-flat-proper}.
Soit $\mathcal{L}$ un $\mathcal{O}_X$-module inversible,
et soit $\mathcal{L}_i$ son image réciproque sur $X_i$. Soit $Z \subset S$,
resp.\ $Z_i \subset S$, le sous-schéma localement fermé construit
dans le lemme \ref{more-morphisms-lemma-diagonal-picard-flat-proper}
pour $(f, \mathcal{L})$, resp.\ $(f_i, \mathcal{L}_i)$,
de sorte que $Z \subset \bigcap_{i = 1, \ldots, n} Z_i$. Si $s \in Z$ et si
$$
H^1(X_s, \mathcal{O}) \longrightarrow
\bigoplus\nolimits_{i = 1, \ldots, n} H^1(X_{i, s}, \mathcal{O})
$$
est injectif, alors $Z \cap U = (\bigcap_{i = 1, \ldots, n} Z_i) \cap U$
(intersection schématique) pour un certain voisinage ouvert $U$ de $s$.
\end{lemma}

\begin{proof}
Ce lemme est une variante du lemme \ref{more-morphisms-lemma-H1-O-picard-flat-proper},
et nous recommandons vivement au lecteur de lire d'abord cette démonstration ;
la présente démonstration en est essentiellement une copie, avec de légères modifications.
Il résulte de la description des points (à valeurs dans des schémas) de $Z$ et des $Z_i$
que $Z \subset \bigcap_{i = 1, \ldots, n} Z_i$, où l'intersection
est schématique. On peut donc remplacer $S$ par l'intersection
schématique $\bigcap_{i = 1, \ldots, n} Z_i$. Après avoir remplacé
$S$ par un voisinage ouvert affine de $s$, on peut supposer que
$\mathcal{L}_i = \mathcal{O}_{X_i}$ pour $i = 1, \ldots, n$.
Soient $E = Rf_*\mathcal{L}$ et
$E_i = Rf_{i, *}\mathcal{L}_i = Rf_{i, *}\mathcal{O}_{X_i}$.
Ce sont des complexes parfaits dont la formation commute à tout
changement de base (Catégories dérivées de schémas, lemme
\ref{perfect-lemma-flat-proper-perfect-direct-image-general}).
En particulier, on voit que
$$
E \otimes_{\mathcal{O}_S}^\mathbf{L} \kappa(s) =
R\Gamma(X_s, \mathcal{L}_s) = R\Gamma(X_s, \mathcal{O}_{X_s})
$$
La seconde égalité vient de ce que $s \in Z$. Posons
$h_j = \dim_{\kappa(s)} H^j(X_s, \mathcal{O}_{X_s})$.
Après avoir rétréci $S$, on peut représenter $E$ par un complexe
$$
\mathcal{O}_S \to \mathcal{O}_S^{\oplus h_1} \to
\mathcal{O}_S^{\oplus h_2} \to \ldots
$$
voir Plus sur l'algèbre, lemme
\ref{more-algebra-lemma-lift-perfect-from-residue-field}
(à strictement parler, cela utilise aussi
Catégories dérivées de schémas, lemmes
\ref{perfect-lemma-affine-compare-bounded} et
\ref{perfect-lemma-perfect-affine}). De même, on peut supposer que $E_i$
est représenté par un complexe
$$
\mathcal{O}_S \to \mathcal{O}_S^{\oplus h_{i, 1}} \to
\mathcal{O}_S^{\oplus h_{i, 2}} \to \ldots
$$
où $h_{i, j} = \dim_{\kappa(s)} H^j(X_{i, s}, \mathcal{O}_{X_{i, s}})$.
Par fonctorialité de la cohomologie, on dispose d'un morphisme
$$
E \longrightarrow E_i
$$
dans $D(\mathcal{O}_S)$ dont la formation commute à tout changement de base.
Puisque le complexe représentant $E$ est un complexe fini de modules libres
de type fini et que $S$ est affine, on peut choisir un morphisme de complexes
$$
\xymatrix{
\mathcal{O}_S \ar[r]_d \ar[d]_{a_i} &
\mathcal{O}_S^{\oplus h_1} \ar[r] \ar[d]_{b_i} &
\mathcal{O}_S^{\oplus h_2} \ar[r] \ar[d]_{c_i} & \ldots \\
\mathcal{O}_S \ar[r]^{d_i} &
\mathcal{O}_S^{\oplus h_{i, 1}} \ar[r] &
\mathcal{O}_S^{\oplus h_{i, 2}} \ar[r] & \ldots
}
$$
représentant le morphisme donné $E \to E_i$. Puisque $s \in Z$, on voit que
la section trivialisante de $\mathcal{L}_s$ donne par image réciproque une section trivialisante
de $\mathcal{L}_{i, s} = \mathcal{O}_{X_{i, s}}$. Ainsi
$a_i \otimes \kappa(s)$ est un isomorphisme ; après avoir rétréci $S$,
on voit donc que $a_i$ est un isomorphisme. Enfin, on utilise l'hypothèse
que $H^1(X_s, \mathcal{O}) \to
\bigoplus_{i = 1, \ldots, n} H^1(X_{i, s}, \mathcal{O})$
est injectif pour voir qu'il existe un mineur $h_1 \times h_1$ de la
matrice définissant $\oplus b_i$ dont l'image est un élément non nul
de $\kappa(s)$. Après avoir rétréci $S$, on peut donc supposer
que $(b_1, \ldots, b_n) : \mathcal{O}_S^{h_1}
\to \bigoplus_{i = 1, \ldots, n} \mathcal{O}_S^{h_{i, 1}}$
est injectif. Or, puisque $\mathcal{L}_i = \mathcal{O}_{X_i}$,
on voit que $d_i = 0$ pour $i = 1, \ldots n$. Il s'ensuit que $d = 0$,
car $(b_1, \ldots, b_n) \circ d = (\oplus d_i) \circ (a_1, \ldots, a_n)$.
On voit ainsi
que la section trivialisante de $\mathcal{L}_s$ se relève en une section
de $\mathcal{L}$ sur $X$. Un argument topologique immédiat (omis)
montre que cela signifie que $\mathcal{L}$ est trivial après avoir éventuellement
encore un peu rétréci $S$.
\end{proof}

\begin{lemma}
\label{more-morphisms-lemma-pic-of-product}
Soient $f : X \to S$ et $g : Y \to S$ des morphismes de schémas
satisfaisant les hypothèses du lemme \ref{more-morphisms-lemma-diagonal-picard-flat-proper}.
Soient $\sigma : S \to X$ et $\tau : S \to Y$ des sections de
$f$ et de $g$. Soit $s \in S$.
Soit $\mathcal{L}$ un faisceau inversible sur $X \times_S Y$.
Si $(1 \times \tau)^*\mathcal{L}$ sur $X$, $(\sigma \times 1)^*\mathcal{L}$
sur $Y$ et $\mathcal{L}|_{(X \times_S Y)_s}$ sont triviales, alors
il existe un voisinage ouvert $U$ de $s$ tel que
$\mathcal{L}$ soit trivial sur $(X \times_S Y)_U$.
\end{lemma}

\begin{proof}
Par la formule de K\"unneth (Variétés, lemme \ref{varieties-lemma-kunneth}),
le morphisme
$$
H^1(X_s \times_{\Spec(\kappa(s))} Y_s, \mathcal{O}) \to
H^1(X_s, \mathcal{O}) \oplus H^1(Y_s, \mathcal{O})
$$
est injectif. On peut donc
appliquer le lemme \ref{more-morphisms-lemma-H1-O-multiple-picard-flat-proper}
aux deux morphismes
$$
1 \times \tau : X \to X \times_S Y
\quad\text{et}\quad
\sigma \times 1 : Y \to X \times_S Y
$$
pour conclure.
\end{proof}

\begin{theorem}[Théorème du cube]
\label{more-morphisms-theorem-of-the-cube}
Soit $S$ un schéma. Soient $X$, $Y$ et $Z$ des schémas sur $S$.
Soient $x : S \to X$ et $y : S \to Y$ des sections des morphismes structuraux.
Soit $\mathcal{L}$ un module inversible sur $X \times_S Y \times_S Z$. Si
\begin{enumerate}
\item $X \to S$ et $Y \to S$ sont des morphismes plats, propres
et de présentation finie, à fibres géométriquement intègres,
\item les images réciproques de $\mathcal{L}$ par
$x \times \text{id}_Y \times \text{id}_Z$ et
$\text{id}_X \times y \times \text{id}_Z$
sont triviales sur $Y \times_S Z$ et $X \times_S Z$,
\item il existe un point $z \in Z$ tel que la restriction de $\mathcal{L}$
à $X \times_S Y \times_S z$ soit triviale, et
\item $Z$ est connexe,
\end{enumerate}
alors $\mathcal{L}$ est trivial.
\end{theorem}

\noindent
Un cas particulier fréquemment utilisé est le suivant.
Soit $k$ un corps. Soient $X, Y, Z$ des variétés munies respectivement de
points $k$-rationnels $x, y, z$. Soit $\mathcal{L}$ un module inversible
sur $X \times Y \times Z$. Si
\begin{enumerate}
\item $\mathcal{L}$ est trivial sur
$x \times Y \times Z$, $X \times y \times Z$ et $X \times Y \times z$, et
\item $X$ et $Y$ sont géométriquement intègres et propres sur $k$,
\end{enumerate}
alors $\mathcal{L}$ est trivial.

\begin{proof}
Remarquons que le morphisme $X \times_S Y \to S$ est plat, propre
et de présentation finie, à fibres géométriquement intègres
(voir Variétés, lemmes \ref{varieties-lemma-geometrically-integral},
\ref{varieties-lemma-bijection-irreducible-components}
et \ref{varieties-lemma-geometrically-reduced-any-base-change} pour
l'assertion concernant les fibres). Par Catégories dérivées de schémas, lemme
\ref{perfect-lemma-proper-flat-geom-red-connected},
on voit que l'image directe du faisceau structural par $X \to S$, $Y \to S$ ou
$X \times_S Y \to S$ est le faisceau structural de $S$, et que cela reste vrai
après tout changement de base. On peut donc appliquer
le lemme \ref{more-morphisms-lemma-diagonal-picard-flat-proper} au morphisme
$$
p : X \times_S Y \times_S Z \longrightarrow Z
$$
et au module inversible $\mathcal{L}$ pour obtenir un sous-schéma localement fermé
« universel » $Z' \subset Z$ tel que $\mathcal{L}|_{X \times_S Y \times_S Z'}$
soit l'image réciproque d'un module inversible $\mathcal{N}$ sur $Z'$.
L'existence de $z$ montre que $Z'$ est non vide. D'après le
lemme \ref{more-morphisms-lemma-get-a-closed},
on voit que $Z' \subset Z$ est un sous-schéma fermé.
Soit $z' \in Z'$ un point.
Remarquons que l'on peut écrire $p$ comme le morphisme produit
$$
(X \times_S Z) \times_Z (Y \times_S Z) \longrightarrow Z
$$
On peut donc appliquer le lemme \ref{more-morphisms-lemma-pic-of-product}
au morphisme $p$, au point $z'$ et aux sections
$\sigma : Z \to X \times_S Z$ et $\tau : Z \to Y \times_S Z$
données par $x$ et $y$. On en déduit que $Z'$ est ouvert.
Ainsi $Z' = Z$ et $\mathcal{L} = p^*\mathcal{N}$
pour un certain module inversible $\mathcal{N}$ sur $Z$.
En prenant l'image réciproque par
$x \times y \times \text{id}_Z : Z  \to X \times_S Y \times_S Z$,
on obtient d'une part $\mathcal{N}$ et, d'autre part,
le module inversible trivial d'après l'hypothèse (2).
Ainsi $\mathcal{N} = \mathcal{O}_Z$, ce qui achève la démonstration.
\end{proof}









\section{Arguments de passage à la limite}
\label{more-morphisms-section-limits}

\noindent
Quelques lemmes faisant intervenir des limites de schémas, ainsi que l'approximation
noethérienne. Nous nous en tenons principalement au cas affine. Certains de ces lemmes sont des
cas particuliers de lemmes du chapitre consacré aux limites.

\begin{lemma}
\label{more-morphisms-lemma-Noetherian-approximation}
Soit $f : X \to S$ un morphisme de schémas affines, de
présentation finie. Alors il existe un diagramme cartésien
$$
\xymatrix{
X_0 \ar[d]_{f_0} & X \ar[l]^g \ar[d]^f \\
S_0 & S \ar[l]
}
$$
tel que
\begin{enumerate}
\item $X_0$, $S_0$ soient des schémas affines,
\item $S_0$ soit de type fini sur $\mathbf{Z}$,
\item $f_0$ soit de type fini.
\end{enumerate}
\end{lemma}

\begin{proof}
Écrivons $S = \Spec(A)$ et $X = \Spec(B)$.
Comme $f$ est de présentation finie, on voit que
$B$ est une $A$-algèbre de présentation finie, voir
Morphismes,
Lemme \ref{morphisms-lemma-locally-finite-presentation-characterize}.
Le lemme résulte donc de
Algèbre, Lemme \ref{algebra-lemma-limit-module-finite-presentation}.
\end{proof}

\begin{lemma}
\label{more-morphisms-lemma-Noetherian-approximation-module}
Soit $f : X \to S$ un morphisme de schémas affines, de
présentation finie. Soit $\mathcal{F}$ un $\mathcal{O}_X$-module quasi-cohérent
de présentation finie. Alors il existe un diagramme comme dans le
Lemme \ref{more-morphisms-lemma-Noetherian-approximation}
tel qu'il existe un $\mathcal{O}_{X_0}$-module cohérent $\mathcal{F}_0$
vérifiant $g^*\mathcal{F}_0 = \mathcal{F}$.
\end{lemma}

\begin{proof}
Écrivons $S = \Spec(A)$, $X = \Spec(B)$ et
$\mathcal{F} = \widetilde{M}$. Comme $f$ est de présentation finie, on voit que
$B$ est une $A$-algèbre de présentation finie, voir
Morphismes,
Lemme \ref{morphisms-lemma-locally-finite-presentation-characterize}.
Comme $\mathcal{F}$ est de présentation finie sur $\mathcal{O}_X$, on voit que
$M$ est un $B$-module de présentation finie, voir
Propriétés, Lemme \ref{properties-lemma-finite-presentation-module}.
Le lemme résulte donc de
Algèbre, Lemme \ref{algebra-lemma-limit-module-finite-presentation}.
\end{proof}

\begin{lemma}
\label{more-morphisms-lemma-Noetherian-approximation-flat-module}
Soit $f : X \to S$ un morphisme de schémas affines, de
présentation finie. Soit $\mathcal{F}$ un $\mathcal{O}_X$-module quasi-cohérent
de présentation finie et plat sur $S$. On peut alors choisir un diagramme comme dans le
Lemme \ref{more-morphisms-lemma-Noetherian-approximation-module}
et un faisceau $\mathcal{F}_0$ tels qu'en outre $\mathcal{F}_0$
soit plat sur $S_0$.
\end{lemma}

\begin{proof}
Écrivons $S = \Spec(A)$, $X = \Spec(B)$ et
$\mathcal{F} = \widetilde{M}$. Comme $f$ est de présentation finie, on voit que
$B$ est une $A$-algèbre de présentation finie, voir
Morphismes,
Lemme \ref{morphisms-lemma-locally-finite-presentation-characterize}.
Comme $\mathcal{F}$ est de présentation finie sur $\mathcal{O}_X$, on voit que
$M$ est un $B$-module de présentation finie, voir
Propriétés, Lemme \ref{properties-lemma-finite-presentation-module}.
Comme $\mathcal{F}$ est plat sur $S$, on voit que $M$ est plat sur $A$, voir
Morphismes, Lemme \ref{morphisms-lemma-flat-module-characterize}.
Le lemme résulte donc de
Algèbre, Lemme \ref{algebra-lemma-flat-finite-presentation-limit-flat}.
\end{proof}

\begin{lemma}
\label{more-morphisms-lemma-Noetherian-approximation-flat}
Soit $f : X \to S$ un morphisme de schémas affines, de
présentation finie et plat. Alors il existe un diagramme comme dans le
Lemme \ref{more-morphisms-lemma-Noetherian-approximation}
tel qu'en outre $f_0$ soit plat.
\end{lemma}

\begin{proof}
C'est un cas particulier du
Lemme \ref{more-morphisms-lemma-Noetherian-approximation-flat-module}.
\end{proof}

\begin{lemma}
\label{more-morphisms-lemma-Noetherian-approximation-smooth}
Soit $f : X \to S$ un morphisme lisse de schémas affines.
Alors il existe un diagramme comme dans le
Lemme \ref{more-morphisms-lemma-Noetherian-approximation}
tel qu'en outre $f_0$ soit lisse.
\end{lemma}

\begin{proof}
Écrivons $S = \Spec(A)$, $X = \Spec(B)$ ;
puisque $f$ est lisse, on voit que $B$ est une $A$-algèbre lisse, voir
Morphismes,
Lemme \ref{morphisms-lemma-smooth-characterize}.
Le lemme résulte donc de
Algèbre, Lemme \ref{algebra-lemma-finite-presentation-fs-Noetherian}.
\end{proof}

\begin{lemma}
\label{more-morphisms-lemma-Noetherian-approximation-geometrically-reduced}
Soit $f : X \to S$ un morphisme de schémas affines, de
présentation finie et dont les fibres sont géométriquement réduites.
Alors il existe un diagramme comme dans le
Lemme \ref{more-morphisms-lemma-Noetherian-approximation}
tel qu'en outre les fibres de $f_0$ soient géométriquement réduites.
\end{lemma}

\begin{proof}
Appliquons le
Lemme \ref{more-morphisms-lemma-Noetherian-approximation}
pour obtenir un diagramme cartésien
$$
\xymatrix{
X_0 \ar[d]_{f_0} & X \ar[l]^g \ar[d]^f \\
S_0 & S \ar[l]_h
}
$$
de schémas affines, où $X_0 \to S_0$ est un morphisme de type fini entre
schémas de type fini sur $\mathbf{Z}$. D'après le
Lemme \ref{more-morphisms-lemma-geometrically-reduced-constructible},
l'ensemble $E \subset S_0$ des points en lesquels la fibre de
$f_0$ est géométriquement réduite est une partie constructible. D'après le
Lemme \ref{more-morphisms-lemma-base-change-fibres-geometrically-reduced},
on a $h(S) \subset E$. Écrivons $S_0 = \Spec(A_0)$ et
$S = \Spec(A)$. Écrivons $A = \colim_i A_i$ comme une
colimite filtrante de $A_0$-algèbres de type fini. D'après
Limites, Lemme \ref{limits-lemma-limit-contained-in-constructible},
on voit que l'image de $\Spec(A_i) \to S_0$ est contenue dans $E$
pour un certain $i$. En remplaçant $S_0$ par $\Spec(A_i)$ et
$X_0$ par $X_0 \times_{S_0} \Spec(A_i)$, on voit que
toutes les fibres de $f_0$ sont géométriquement réduites.
\end{proof}

\begin{lemma}
\label{more-morphisms-lemma-Noetherian-approximation-geometrically-irreducible}
Soit $f : X \to S$ un morphisme de schémas affines, de
présentation finie et dont les fibres sont géométriquement irréductibles.
Alors il existe un diagramme comme dans le
Lemme \ref{more-morphisms-lemma-Noetherian-approximation}
tel qu'en outre les fibres de $f_0$ soient géométriquement irréductibles.
\end{lemma}

\begin{proof}
Appliquons le
Lemme \ref{more-morphisms-lemma-Noetherian-approximation}
pour obtenir un diagramme cartésien
$$
\xymatrix{
X_0 \ar[d]_{f_0} & X \ar[l]^g \ar[d]^f \\
S_0 & S \ar[l]_h
}
$$
de schémas affines, où $X_0 \to S_0$ est un morphisme de type fini entre
schémas de type fini sur $\mathbf{Z}$. D'après le
Lemme \ref{more-morphisms-lemma-nr-geom-irreducible-components-constructible},
l'ensemble $E \subset S_0$ des points en lesquels la fibre de
$f_0$ est géométriquement irréductible est une partie constructible. D'après le
Lemme \ref{more-morphisms-lemma-base-change-fibres-geometrically-irreducible},
on a $h(S) \subset E$. Écrivons $S_0 = \Spec(A_0)$ et
$S = \Spec(A)$. Écrivons $A = \colim_i A_i$ comme une
colimite filtrante de $A_0$-algèbres de type fini. D'après
Limites, Lemme \ref{limits-lemma-limit-contained-in-constructible},
on voit que l'image de $\Spec(A_i) \to S_0$ est contenue dans $E$
pour un certain $i$. En remplaçant $S_0$ par $\Spec(A_i)$ et
$X_0$ par $X_0 \times_{S_0} \Spec(A_i)$, on voit que
toutes les fibres de $f_0$ sont géométriquement irréductibles.
\end{proof}

\begin{lemma}
\label{more-morphisms-lemma-Noetherian-approximation-geometrically-connected}
Soit $f : X \to S$ un morphisme de schémas affines, de
présentation finie et dont les fibres sont géométriquement connexes.
Alors il existe un diagramme comme dans le
Lemme \ref{more-morphisms-lemma-Noetherian-approximation}
tel qu'en outre les fibres de $f_0$ soient géométriquement connexes.
\end{lemma}

\begin{proof}
Appliquons le
Lemme \ref{more-morphisms-lemma-Noetherian-approximation}
pour obtenir un diagramme cartésien
$$
\xymatrix{
X_0 \ar[d]_{f_0} & X \ar[l]^g \ar[d]^f \\
S_0 & S \ar[l]_h
}
$$
de schémas affines, où $X_0 \to S_0$ est un morphisme de type fini entre
schémas de type fini sur $\mathbf{Z}$. D'après le
Lemme \ref{more-morphisms-lemma-nr-geom-connected-components-constructible},
l'ensemble $E \subset S_0$ des points en lesquels la fibre de
$f_0$ est géométriquement connexe est une partie constructible. D'après le
Lemme \ref{more-morphisms-lemma-base-change-fibres-geometrically-connected},
on a $h(S) \subset E$. Écrivons $S_0 = \Spec(A_0)$ et
$S = \Spec(A)$. Écrivons $A = \colim_i A_i$ comme une
colimite filtrante de $A_0$-algèbres de type fini. D'après
Limites, Lemme \ref{limits-lemma-limit-contained-in-constructible},
on voit que l'image de $\Spec(A_i) \to S_0$ est contenue dans $E$
pour un certain $i$. En remplaçant $S_0$ par $\Spec(A_i)$ et
$X_0$ par $X_0 \times_{S_0} \Spec(A_i)$, on voit que
toutes les fibres de $f_0$ sont géométriquement connexes.
\end{proof}

\begin{lemma}
\label{more-morphisms-lemma-Noetherian-approximation-dimension-d}
Soit $d \geq 0$ un entier.
Soit $f : X \to S$ un morphisme de schémas affines, de
présentation finie, dont toutes les fibres sont de dimension $d$.
Alors il existe un diagramme comme dans le
Lemme \ref{more-morphisms-lemma-Noetherian-approximation}
tel qu'en outre toutes les fibres de $f_0$ soient de dimension $d$.
\end{lemma}

\begin{proof}
Appliquons le
Lemme \ref{more-morphisms-lemma-Noetherian-approximation}
pour obtenir un diagramme cartésien
$$
\xymatrix{
X_0 \ar[d]_{f_0} & X \ar[l]^g \ar[d]^f \\
S_0 & S \ar[l]_h
}
$$
de schémas affines, où $X_0 \to S_0$ est un morphisme de type fini entre
schémas de type fini sur $\mathbf{Z}$. D'après le
Lemme \ref{more-morphisms-lemma-dimension-fibres-constructible},
l'ensemble $E \subset S_0$ des points en lesquels la fibre de
$f_0$ est de dimension $d$ est une partie constructible. D'après le
Lemme \ref{more-morphisms-lemma-base-change-dimension-fibres},
on a $h(S) \subset E$. Écrivons $S_0 = \Spec(A_0)$ et
$S = \Spec(A)$. Écrivons $A = \colim_i A_i$ comme une
colimite filtrante de $A_0$-algèbres de type fini. D'après
Limites, Lemme \ref{limits-lemma-limit-contained-in-constructible},
on voit que l'image de $\Spec(A_i) \to S_0$ est contenue dans $E$
pour un certain $i$. En remplaçant $S_0$ par $\Spec(A_i)$ et
$X_0$ par $X_0 \times_{S_0} \Spec(A_i)$, on voit que
toutes les fibres de $f_0$ sont de dimension $d$.
\end{proof}

\begin{lemma}
\label{more-morphisms-lemma-Noetherian-approximation-standard-syntomic}
Soit $f : X \to S$ un morphisme de schémas affines qui est
syntomique standard (voir
Morphismes, Définition \ref{morphisms-definition-syntomic}).
Alors il existe un diagramme comme dans le
Lemme \ref{more-morphisms-lemma-Noetherian-approximation}
tel qu'en outre $f_0$ soit syntomique standard.
\end{lemma}

\begin{proof}
Ce lemme est une copie de
Algèbre,
Lemme \ref{algebra-lemma-relative-global-complete-intersection-Noetherian}.
\end{proof}

\begin{lemma}
\label{more-morphisms-lemma-Noetherian-approximation-combine}
(Approximation noethérienne et combinaison de propriétés.)
Soient $P$, $Q$ des propriétés de morphismes de schémas qui sont stables
par changement de base. Soit $f : X \to S$ un morphisme de présentation finie
entre schémas affines. Supposons que l'on puisse trouver des diagrammes cartésiens
$$
\vcenter{
\xymatrix{
X_1 \ar[d]_{f_1} & X \ar[l] \ar[d]^f \\
S_1 & S \ar[l]
}
}
\quad\text{et}\quad
\vcenter{
\xymatrix{
X_2 \ar[d]_{f_2} & X \ar[l] \ar[d]^f \\
S_2 & S \ar[l]
}
}
$$
de schémas affines, où $S_1$, $S_2$ sont de type fini sur $\mathbf{Z}$
et $f_1$, $f_2$ sont de type fini, tels que $f_1$ ait la propriété $P$
et $f_2$ ait la propriété $Q$. Alors on peut trouver un diagramme cartésien
$$
\xymatrix{
X_0 \ar[d]_{f_0} & X \ar[l] \ar[d]^f \\
S_0 & S \ar[l]
}
$$
de schémas affines, où $S_0$ est de type fini sur $\mathbf{Z}$
et $f_0$ est de type fini, tel que $f_0$ ait à la fois la propriété $P$ et
la propriété $Q$.
\end{lemma}

\begin{proof}
Les deux diagrammes donnés correspondent à des diagrammes cocartésiens d'anneaux
$$
\vcenter{
\xymatrix{
B_1 \ar[r] & B \\
A_1 \ar[u] \ar[r] & A \ar[u]
}
}
\quad\text{et}\quad
\vcenter{
\xymatrix{
B_2 \ar[r] & B \\
A_2 \ar[u] \ar[r] & A \ar[u]
}
}
$$
Soit $A_0 \subset A$ une sous-$\mathbf{Z}$-algèbre de type fini de $A$
contenant les images de $A_1 \to A$ et de $A_2 \to A$. Une telle sous-algèbre
existe car, par hypothèse, $A_1$ et $A_2$ sont toutes deux de type fini sur
$\mathbf{Z}$. Remarquons que les anneaux $B_{0, 1} = B_1 \otimes_{A_1} A_0$
et $B_{0, 2} = B_2 \otimes_{A_2} A_0$ sont des $A_0$-algèbres de type fini
qui vérifient
$B_{0, 1} \otimes_{A_0} A \cong B \cong B_{0, 2} \otimes_{A_0} A$
en tant que $A$-algèbres. Comme $A$ est la colimite filtrante de ses
sous-$A_0$-algèbres de type fini, d'après
Limites, Lemme \ref{limits-lemma-descend-finite-presentation},
on peut supposer, quitte à agrandir $A_0$, qu'il existe un isomorphisme
$B_{0, 1} \cong B_{0, 2}$ de $A_0$-algèbres. Puisque les propriétés $P$ et $Q$
sont supposées stables par changement de base, on conclut qu'en posant
$S_0 = \Spec(A_0)$ et
$$
X_0 = X_1 \times_{S_1} S_0 =
\Spec(B_{0, 1}) \cong \Spec(B_{0, 2}) = X_2 \times_{S_2} S_0
$$
on obtient ce que l'on veut.
\end{proof}












\section{Voisinages \'etales}
\label{more-morphisms-section-etale-neighbourhoods}

\noindent
Il s'avère que certaines propriétés des morphismes sont plus faciles à étudier
après un changement de base \'etale. Il est commode d'introduire la
terminologie suivante.

\begin{definition}
\label{more-morphisms-definition-etale-neighbourhood}
Soit $S$ un schéma. Soit $s \in S$ un point.
\begin{enumerate}
\item Un {\it voisinage \'etale de $(S, s)$} est un
couple $(U, u)$ muni d'un morphisme \'etale
de schémas $\varphi : U \to S$ tel que $\varphi(u) = s$.
\item Un {\it morphisme de voisinages \'etales} $f : (V, v) \to (U, u)$
de $(S, s)$ est simplement un morphisme de $S$-schémas $f : V \to U$ tel
que $f(v) = u$.
\item Un {\it voisinage \'etale élémentaire} est un voisinage \'etale
$\varphi : (U, u) \to (S, s)$ tel que $\kappa(s) = \kappa(u)$.
\end{enumerate}
\end{definition}

\noindent
La notion de voisinage \'etale élémentaire porte de nombreux noms
dans la littérature ; ces objets sont par exemple parfois appelés
``voisinages \'etales'' (\cite[Page 36]{Milne} ou
``fortement \'etales'' (\cite[Page 108]{KPR}).
Nous suivons ici la convention de l'article
\cite{GruRay} en les appelant voisinages \'etales élémentaires.

\medskip\noindent
Si $f : (V, v) \to (U, u)$ est un morphisme de voisinages \'etales,
alors $f$ est automatiquement \'etale, voir
Morphismes, Lemme \ref{morphisms-lemma-etale-permanence}.
Ainsi, $(V, v)$ devient un voisinage \'etale de
$(U, u)$. Bien entendu, puisque la composée de morphismes \'etales
est \'etale (Morphismes, Lemme \ref{morphisms-lemma-composition-etale}),
on voit réciproquement que tout voisinage \'etale $(V, v)$ de
$(U, u)$ est aussi un voisinage \'etale de $(S, s)$.
Remarquons également que, si $U \subset S$ est un voisinage ouvert
de $s$, alors $(U, s) \to (S, s)$ est un voisinage \'etale.
Cela résulte du fait qu'une immersion ouverte est
\'etale (Morphismes, Lemme \ref{morphisms-lemma-open-immersion-etale}).
Nous utiliserons ces remarques sans autre mention dans toute cette
section.

\medskip\noindent
Remarquons que $\kappa(u)/\kappa(s)$ est une extension finie séparable
si $(U, u) \to (S, s)$ est un voisinage \'etale,
voir Morphismes, Lemme \ref{morphisms-lemma-etale-at-point}.

\begin{lemma}
\label{more-morphisms-lemma-realize-prescribed-residue-field-extension-etale}
Soit $S$ un schéma.
Soit $s \in S$.
Soit $k/\kappa(s)$ une extension finie séparable.
Alors il existe un voisinage \'etale $(U, u) \to (S, s)$
tel que l'extension de corps $\kappa(u)/\kappa(s)$ soit
isomorphe à $k/\kappa(s)$.
\end{lemma}

\begin{proof}
On peut supposer $S$ affine.
Dans ce cas, le lemme résulte de
Algèbre, Lemme \ref{algebra-lemma-make-etale-map-prescribed-residue-field}.
\end{proof}

\begin{lemma}
\label{more-morphisms-lemma-etale-neighbourhoods-not-quite-filtered}
Soient $S$ un schéma et $s$ un point de $S$.
La catégorie des voisinages \'etales possède les propriétés suivantes :
\begin{enumerate}
\item Soient $(U_i, u_i)_{i=1, 2}$ deux voisinages \'etales de
$s$ dans $S$. Alors il existe un troisième voisinage \'etale
$(U, u)$ et des morphismes
$(U, u) \to (U_i, u_i)$, $i = 1, 2$.
\item Soient $h_1, h_2: (U, u) \to (U', u')$ deux
morphismes entre voisinages \'etales de $s$.
Supposons que $h_1$, $h_2$ induisent la même application $\kappa(u') \to \kappa(u)$ entre corps
résiduels. Alors il existe un voisinage \'etale $(U'', u'')$ et un morphisme
$h : (U'', u'') \to (U, u)$
qui égalise $h_1$ et $h_2$, c'est-à-dire tel que
$h_1 \circ h = h_2 \circ h$.
\end{enumerate}
\end{lemma}

\begin{proof}
Pour (1), considérons le produit fibré $U = U_1 \times_S U_2$.
Il est \'etale sur $U_1$ et sur $U_2$, car les morphismes \'etales sont
préservés par changement de base, voir
Morphismes, Lemme \ref{morphisms-lemma-base-change-etale}.
Il existe un point de $U$ qui s'envoie sur $u_1$ et sur $u_2$, par exemple
d'après la description des points d'un produit fibré dans
Schémas, Lemme \ref{schemes-lemma-points-fibre-product}.
Pour (2), définissons $U''$ comme le produit fibré
$$
\xymatrix{
U'' \ar[r] \ar[d] & U \ar[d]^{(h_1, h_2)} \\
U' \ar[r]^-\Delta & U' \times_S U'.
}
$$
Puisque $h_1$ et $h_2$ induisent la même application de corps résiduels
$\kappa(u') \to \kappa(u)$, il existe un point $u'' \in U''$
au-dessus de $u'$ tel que $\kappa(u'') = \kappa(u')$.
En particulier, $U'' \not = \emptyset$.
De plus, puisque $U'$ est \'etale sur $S$, le produit fibré
$U'\times_S U'$ l'est aussi (voir
Morphismes, Lemmes \ref{morphisms-lemma-base-change-etale} et
\ref{morphisms-lemma-composition-etale}).
La flèche verticale $(h_1, h_2)$ est donc \'etale d'après
Morphismes, Lemme \ref{morphisms-lemma-etale-permanence}.
Par conséquent, $U''$ est \'etale sur $U'$ par changement de base, et donc aussi
\'etale sur $S$ (car les composées de morphismes \'etales sont \'etales).
Ainsi, $(U'', u'')$ est une solution au problème.
\end{proof}

\begin{lemma}
\label{more-morphisms-lemma-elementary-etale-neighbourhoods}
Soient $S$ un schéma et $s$ un point de $S$.
La catégorie des voisinages \'etales élémentaires de $(S, s)$
est cofiltrante (voir
Catégories, Définition \ref{categories-definition-codirected}).
\end{lemma}

\begin{proof}
Cela résulte immédiatement des définitions et du
Lemme \ref{more-morphisms-lemma-etale-neighbourhoods-not-quite-filtered}.
\end{proof}

\begin{lemma}
\label{more-morphisms-lemma-describe-henselization}
Soit $S$ un schéma. Soit $s \in S$. Alors on a
$$
\mathcal{O}_{S, s}^h =
\colim_{(U, u)} \mathcal{O}(U)
$$
où la colimite est prise sur la catégorie filtrante qui est l'opposée de la
catégorie des voisinages \'etales élémentaires $(U, u)$ de $(S, s)$.
\end{lemma}

\begin{proof}
Soit $\Spec(A) \subset S$ un voisinage affine de $s$.
Soit $\mathfrak p \subset A$ l'idéal premier correspondant à $s$.
Avec ces choix, on dispose d'isomorphismes canoniques
$\mathcal{O}_{S, s} = A_{\mathfrak p}$ et $\kappa(s) = \kappa(\mathfrak p)$.
Un système cofinal de voisinages \'etales élémentaires est donné par les
voisinages \'etales élémentaires $(U, u)$ tels que $U$ soit affine et que
$U \to S$ se factorise par $\Spec(A)$. Autrement dit, on voit que
le membre de droite est égal à $\colim_{(B, \mathfrak q)} B$,
où la colimite porte sur les $A$-algèbres \'etales $B$ munies d'un idéal premier
$\mathfrak q$ au-dessus de $\mathfrak p$ et telles que
$\kappa(\mathfrak p) = \kappa(\mathfrak q)$.
Le lemme résulte donc de
Algèbre, Lemme \ref{algebra-lemma-henselization-different}.
\end{proof}

\noindent
On peut relever à l'espace total les voisinages \'etales de points
des fibres.

\begin{lemma}
\label{more-morphisms-lemma-lift-etale-neighbourhood-fibre}
\begin{slogan}
Relever à l'espace total un voisinage \'etale d'un point de la fibre.
\end{slogan}
Soit $X \to S$ un morphisme de schémas. Soit $x \in X$ d'image $s \in S$.
Soit $(V, v) \to (X_s, x)$ un voisinage \'etale.
Alors il existe un voisinage \'etale $(U, u) \to (X, x)$
tel qu'il existe un morphisme $(U_s, u) \to (V, v)$
de voisinages \'etales de $(X_s, x)$ qui soit une immersion ouverte.
\end{lemma}

\begin{proof}
On peut supposer $X$, $V$ et $S$ affines. Supposons que le morphisme
$X \to S$ soit donné par $A \to B$, le point $x$ par un idéal premier
$\mathfrak q \subset B$, le point $s$ par $\mathfrak p = A \cap \mathfrak q$,
et le morphisme $V \to X_s$ par $B \otimes_A \kappa(\mathfrak p) \to C$.
Puisque $\kappa(\mathfrak p)$ est un localisé de
$A/\mathfrak p$, il existe $f \in A$, $f \not \in \mathfrak p$
et un morphisme d'anneaux \'etale $B \otimes_A (A/\mathfrak p)_f \to D$
tel que
$$
C = (B \otimes_A \kappa(\mathfrak p))
\otimes_{B \otimes_A (A/\mathfrak p)_f} D
$$
Voir Algèbre, Lemme \ref{algebra-lemma-etale}, partie (9).
Après avoir remplacé $A$ par $A_f$ et $B$ par $B_f$, on peut supposer que
$D$ est \'etale sur $B \otimes_A A/\mathfrak p = B/\mathfrak p B$.
On peut alors appliquer Algèbre, Lemme \ref{algebra-lemma-lift-etale}.
Cela démontre le lemme.
\end{proof}














\section{Voisinages \'etales et branches}
\label{more-morphisms-section-etale-nbhds-branches}

\noindent
Le nombre de branches (géométriques) d'un schéma en un point a été
défini dans Propriétés, section \ref{properties-section-unibranch}.
Dans Variétés, section \ref{varieties-section-number-of-branches},
nous l'avons relié aux fibres du morphisme de normalisation.
Dans cette section, nous donnons une caractérisation de ce nombre au moyen de
voisinages \'etales.

\begin{lemma}
\label{more-morphisms-lemma-nr-minimal-primes}
Soit $R = \colim R_i$ la colimite d'un système inductif d'anneaux
dont les morphismes de transition sont fidèlement plats.
Alors le nombre d'idéaux premiers minimaux de $R$,
considéré comme un élément de $\{0, 1, 2, \ldots, \infty\}$,
est la borne supérieure des nombres d'idéaux premiers minimaux des $R_i$.
\end{lemma}

\begin{proof}
Si $A \to B$ est un morphisme plat d'anneaux, alors $\Spec(B) \to \Spec(A)$
envoie les idéaux premiers minimaux sur des idéaux premiers minimaux en vertu du théorème de descente
(Algèbre, lemme \ref{algebra-lemma-flat-going-down}).
Si $A \to B$ est fidèlement plat, alors tout idéal premier
minimal est l'image d'un idéal premier minimal (d'après
Algèbre, lemmes \ref{algebra-lemma-ff-rings} et
\ref{algebra-lemma-minimal-prime-image-minimal-prime}).
Par conséquent, le nombre d'idéaux premiers minimaux de $R_i$ est
$\geq$ au nombre d'idéaux premiers minimaux de $R_{i'}$ si $i \leq i'$.
D'après Algèbre, lemme \ref{algebra-lemma-colimit-faithfully-flat},
chacun des morphismes $R_i \to R$ est
fidèlement plat, et l'on voit aussi que
le nombre d'idéaux premiers minimaux de $R$ est
$\geq$ au nombre d'idéaux premiers minimaux de $R_i$.
Enfin, supposons que $\mathfrak q_1, \ldots, \mathfrak q_n$
soient des idéaux premiers minimaux de $R$ deux à deux distincts. On peut alors
trouver un $i$ tel que $R_i \cap \mathfrak q_1, \ldots, R_i \cap \mathfrak q_n$
soient deux à deux distincts (comme ensembles, et donc comme idéaux premiers).
Ceci démontre le lemme.
\end{proof}

\begin{lemma}
\label{more-morphisms-lemma-nr-branches}
Soient $X$ un schéma et $x \in X$ un point. Alors
\begin{enumerate}
\item le nombre de branches de $X$ en $x$ est égal à
la borne supérieure du nombre de composantes irréductibles de $U$
passant par $u$, pour les voisinages \'etales élémentaires
$(U, u) \to (X, x)$,
\item le nombre de branches géométriques de $X$ en $x$ est égal à
la borne supérieure du nombre de composantes irréductibles de $U$
passant par $u$, pour les voisinages \'etales
$(U, u) \to (X, x)$,
\item $X$ est unibranche en $x$ si et seulement si, pour tout
voisinage \'etale élémentaire $(U, u) \to (X, x)$, il existe
exactement une composante irréductible de $U$ passant par $u$, et
\item $X$ est géométriquement unibranche en $x$ si et seulement si, pour tout
voisinage \'etale $(U, u) \to (X, x)$, il existe
exactement une composante irréductible de $U$ passant par $u$.
\end{enumerate}
\end{lemma}

\begin{proof}
Les assertions (3) et (4) se déduisent des assertions (1) et (2) au moyen de
Propriétés, lemme \ref{properties-lemma-number-of-branches-1}.

\medskip\noindent
Démonstration de (1). Soit $\Spec(A)$ un voisinage ouvert affine
de $x$, et soit $\mathfrak p \subset A$ l'idéal premier
correspondant à $x$. On peut remplacer $X$ par $\Spec(A)$, et
il suffit, dans la borne supérieure, de considérer les voisinages \'etales élémentaires affines
$(U, u)$, puisqu'ils forment un sous-système cofinal.
Rappelons que l'hensélisé $A_\mathfrak p^h$
est la colimite des anneaux $B_\mathfrak q$ sur la catégorie
des couples $(B, \mathfrak q)$, où $B$ est une $A$-algèbre \'etale
et $\mathfrak q$ un idéal premier au-dessus de $\mathfrak p$ tel que
$\kappa(\mathfrak q) = \kappa(\mathfrak p)$ ; voir
Algèbre, lemme \ref{algebra-lemma-henselization-different}.
Ces couples $(B, \mathfrak q)$ correspondent exactement aux
voisinages \'etales élémentaires affines $(U, u)$
par la correspondance entre anneaux et schémas affines.
Observons que les composantes irréductibles de $\Spec(B)$
passant par $\mathfrak q$ correspondent exactement aux idéaux
premiers minimaux de $B_\mathfrak q$. Le nombre d'idéaux premiers
minimaux de $A_\mathfrak p^h$ est le nombre de branches
de $X$ en $x$, d'après Propriétés, définition
\ref{properties-definition-number-of-branches}.
Observons que les morphismes de transition
$B_\mathfrak q \to B'_{\mathfrak q'}$ du système
sont tous plats. Puisqu'un morphisme local plat d'anneaux est fidèlement plat
(Algèbre, lemme \ref{algebra-lemma-local-flat-ff}),
on voit que l'assertion découle du
lemme \ref{more-morphisms-lemma-nr-minimal-primes}.

\medskip\noindent
Démonstration de (2). La démonstration est la même que celle de (1), si ce n'est que l'on utilise
Algèbre, lemme \ref{algebra-lemma-strict-henselization-different}.
Il existe une petite différence : étant donnée une clôture algébrique séparable
$\kappa^{sep}$ de $\kappa(x)$, pour tout voisinage \'etale
$(U, u)$ on peut choisir un plongement de $\kappa(x)$-algèbres
$\phi : \kappa(u) \to \kappa^{sep}$,
car l'extension $\kappa(u)/\kappa(x)$ est finie et séparable
(Morphismes, lemme \ref{morphisms-lemma-etale-at-point}).
On peut donc prendre la borne supérieure sur tous les triplets
$(U, u, \phi)$, où $(U, u) \to (X, x)$ est un voisinage
\'etale affine et $\phi : \kappa(u) \to \kappa^{sep}$
est un plongement de $\kappa(x)$-algèbres. Ces triplets correspondent
exactement à ceux de
Algèbre, lemme \ref{algebra-lemma-strict-henselization-different},
et le reste de la démonstration est exactement le même.
\end{proof}

\noindent
Nous aurons besoin d'une variante relative du lemme précédent.

\begin{lemma}
\label{more-morphisms-lemma-nr-branches-fibre}
Soit $X \to S$ un morphisme de schémas et soit $x \in X$ un point d'image $s$.
Alors
\begin{enumerate}
\item le nombre de branches de la fibre $X_s$ en $x$ est égal à
la borne supérieure du nombre de composantes irréductibles de la fibre $U_s$
passant par $u$, pour les voisinages \'etales élémentaires
$(U, u) \to (X, x)$,
\item le nombre de branches géométriques de la fibre $X_s$ en $x$ est égal à
la borne supérieure du nombre de composantes irréductibles de la fibre $U_s$
passant par $u$, pour les voisinages \'etales
$(U, u) \to (X, x)$,
\item la fibre $X_s$ est unibranche en $x$ si et seulement si, pour tout
voisinage \'etale élémentaire $(U, u) \to (X, x)$, il existe
exactement une composante irréductible de la fibre $U_s$ passant par $u$, et
\item $X$ est géométriquement unibranche en $x$ si et seulement si, pour tout
voisinage \'etale $(U, u) \to (X, x)$, il existe
exactement une composante irréductible de $U_s$ passant par $u$.
\end{enumerate}
\end{lemma}

\begin{proof}
Combiner les lemmes \ref{more-morphisms-lemma-nr-branches} et
\ref{more-morphisms-lemma-lift-etale-neighbourhood-fibre}.
\end{proof}

\begin{lemma}
\label{more-morphisms-lemma-number-of-branches-and-smooth}
Soit $X \to S$ un morphisme lisse de schémas.
Soit $x \in X$ d'image $s \in S$.
Alors
\begin{enumerate}
\item Le nombre de branches géométriques de $X$ en $x$
est égal au nombre de branches géométriques de $S$ en $s$.
\item Si $\kappa(x)/\kappa(s)$ est une extension de corps purement inséparable\footnote{
En fait, il suffirait que $\kappa(x)$ soit géométriquement
irréductible sur $\kappa(s)$. Si nous en avons un jour besoin, nous ajouterons une démonstration détaillée.},
alors le nombre de branches de $X$ en $x$
est égal au nombre de branches de $S$ en $s$.
\end{enumerate}
\end{lemma}

\begin{proof}
Ceci résulte immédiatement de Plus sur l'algèbre, lemme
\ref{more-algebra-lemma-invariance-number-branches-smooth},
et des définitions.
\end{proof}




\section{Morphismes non ramifiés et \'etales}
\label{more-morphisms-section-unramified-is-etale}

\noindent
Les morphismes non ramifiés sont parfois automatiquement \'etales.

\begin{lemma}
\label{more-morphisms-lemma-unramified-dominant-unibranch-is-etale}
Soit $f : X \to Y$ un morphisme de schémas. Soit $x \in X$
d'image $y \in Y$. Supposons que
\begin{enumerate}
\item $Y$ soit intègre et géométriquement unibranche en $y$,
\item $f$ soit localement de type fini,
\item il existe une spécialisation $x' \leadsto x$ telle que $f(x')$
soit le point générique de $Y$,
\item $f$ soit non ramifié en $x$.
\end{enumerate}
Alors $f$ est \'etale en $x$.
\end{lemma}

\begin{proof}
On peut remplacer $X$ et $Y$ par des voisinages ouverts affines convenables
de $x$ et $y$. Alors $Y$ est le spectre d'un anneau intègre $A$ et
$X$ est le spectre d'une $A$-algèbre de type fini $B$.
Soient $\mathfrak q \subset B$ l'idéal premier correspondant
à $x$ et $\mathfrak p \subset A$ l'idéal premier correspondant à $y$.
L'anneau local $A_\mathfrak p = \mathcal{O}_{Y, y}$ est géométriquement unibranche.
Le morphisme d'anneaux $A \to B$ est non ramifié en $\mathfrak q$.
En outre, le point $x'$ de (3) correspond à un idéal premier
$\mathfrak q' \subset \mathfrak q$ tel que $A \cap \mathfrak q' = (0)$.
Il s'ensuit que $A_\mathfrak p \to B_\mathfrak q$ est injectif.
On conclut par
Plus sur l'algèbre, lemme
\ref{more-algebra-lemma-local-unramified-extension-unibranch-domain-is-etale}.
\end{proof}

\begin{lemma}
\label{more-morphisms-lemma-global-unramified-dominant-unibranch-is-etale}
\begin{reference}
\cite[Exposé I, Corollaire 9.11]{SGA1}
\end{reference}
Soit $f : X \to Y$ un morphisme de schémas. Supposons que
\begin{enumerate}
\item $Y$ soit intègre et géométriquement unibranche,
\item au moins une composante irréductible de $X$ domine $Y$,
\item $f$ soit non ramifié, et
\item $X$ soit connexe.
\end{enumerate}
Alors $f$ est \'etale et $X$ est irréductible.
\end{lemma}

\begin{proof}
Soit $X' \subset X$ la composante irréductible qui domine $Y$.
Cela signifie que le point générique de $X'$ s'envoie sur le point générique de $Y$
(voir par exemple Morphismes, lemme
\ref{morphisms-lemma-dominant-finite-number-irreducible-components}).
D'après le lemme \ref{more-morphisms-lemma-unramified-dominant-unibranch-is-etale},
$f$ est \'etale en tout point de $X'$. En particulier,
le sous-schéma ouvert $U \subset X$ où $f$ est \'etale contient $X'$.
Notons que tout ouvert quasi-compact de $U$ possède un nombre fini de composantes
irréductibles ; voir
Descente, lemme \ref{descent-lemma-locally-finite-nr-irred-local-fppf}.
D'autre part, puisque $Y$ est géométriquement unibranche et que $U$ est \'etale
sur $Y$, le schéma $U$ est géométriquement unibranche. En particulier, par
chaque point de $U$ passe au plus une composante irréductible.
Un argument topologique simple montre alors que $X' \subset U$
est à la fois ouvert et fermé. Ainsi, $X'$ est ouvert et fermé dans $X$ et, par
connexité, on obtient $X = U = X'$, comme souhaité.
\end{proof}

\begin{lemma}
\label{more-morphisms-lemma-weird-permanence-etale}
Soient $f : X \to Y$ et $g : Y \to Z$ des morphismes de schémas.
Soit $x \in X$ d'image $y \in Y$. Supposons que
\begin{enumerate}
\item $Y$ soit intègre et géométriquement unibranche en $y$,
\item $g \circ f$ soit \'etale en $x$,
\item il existe une spécialisation $x' \leadsto x$ telle que $f(x')$
soit le point générique de $Y$.
\end{enumerate}
Alors $f$ est \'etale en $x$ et $g$ est \'etale en $y$.
\end{lemma}

\begin{proof}
Après avoir remplacé $X$ par un voisinage ouvert de $x$,
on peut supposer que $g \circ f$ est \'etale.
On trouve alors que $f$ est non ramifié par Morphismes, lemmes
\ref{morphisms-lemma-unramified-permanence}
et \ref{morphisms-lemma-etale-smooth-unramified}.
Ainsi, $f$ est \'etale en $x$ d'après le
lemme \ref{more-morphisms-lemma-unramified-dominant-unibranch-is-etale}.
Par descente \'etale, on voit alors que $g$ est \'etale en $y$ ; voir
par exemple Descente, lemme
\ref{descent-lemma-syntomic-smooth-etale-permanence}.
\end{proof}

\begin{lemma}
\label{more-morphisms-lemma-global-weird-permanence-etale}
Soient $f : X \to Y$ et $g : Y \to Z$ des morphismes de schémas.
Supposons que
\begin{enumerate}
\item $Y$ soit intègre et géométriquement unibranche,
\item $g \circ f$ soit \'etale,
\item toute composante irréductible de $X$ domine $Y$.
\end{enumerate}
Alors $f$ est \'etale et $g$ est \'etale en tout point de
l'image de $f$.
\end{lemma}

\begin{proof}
Ceci résulte immédiatement de la version ponctuelle,
lemme \ref{more-morphisms-lemma-weird-permanence-etale}.
\end{proof}










\section{Découpage des morphismes lisses}
\label{more-morphisms-section-etale-over-smooth}

\noindent
Dans cette section, nous expliquons un résultat qui affirme, en gros, que
les recouvrements lisses d'un schéma $S$ peuvent être raffinés par des recouvrements \'etales.
La technique de démonstration repose sur un argument de découpage.

\begin{lemma}
\label{more-morphisms-lemma-slice-smooth-given-element}
Soit $f : X \to S$ un morphisme de schémas.
Soit $x \in X$ un point d'image $s \in S$.
Soit $h \in \mathfrak m_x \subset \mathcal{O}_{X, x}$.
Supposons que
\begin{enumerate}
\item $f$ soit lisse en $x$, et
\item l'image $\text{d}\overline{h}$ de $\text{d}h$ dans
$$
\Omega_{X_s/s, x} \otimes_{\mathcal{O}_{X_s, x}} \kappa(x) =
\Omega_{X/S, x} \otimes_{\mathcal{O}_{X, x}} \kappa(x)
$$
soit non nulle.
\end{enumerate}
Alors il existe un voisinage ouvert affine $U \subset X$ de $x$
tel que $h$ provienne d'un $h \in \Gamma(U, \mathcal{O}_U)$ et
que $D = V(h)$ soit un diviseur de Cartier effectif dans $U$, avec $x \in D$, et que
$D \to S$ soit lisse.
\end{lemma}

\begin{proof}
Puisque $f$ est lisse en $x$, on peut supposer, après avoir remplacé $X$ par un voisinage
ouvert de $x$, que $f$ est lisse. En particulier, on voit que
$f$ est plat et localement de présentation finie. D'après le
lemme \ref{more-morphisms-lemma-slice-given-element},
on sait déjà qu'il existe un voisinage ouvert $U \subset X$ de $x$
tel que $h$ provienne d'un $h \in \Gamma(U, \mathcal{O}_U)$ et
que $D = V(h)$ soit un diviseur de Cartier effectif dans $U$, avec $x \in D$, et que
$D \to S$ soit plat et de présentation finie. D'après
Morphismes, lemme \ref{morphisms-lemma-differentials-relative-immersion},
on dispose d'une suite exacte à droite
$$
\mathcal{C}_{D/U} \to i^*\Omega_{U/S} \to \Omega_{D/S} \to 0
$$
où $i : D \to U$ est l'immersion fermée et $\mathcal{C}_{D/U}$
est le faisceau conormal de $D$ dans $U$. Puisque $D$ est un diviseur de Cartier
effectif défini par $h \in \Gamma(U, \mathcal{O}_U)$, on voit que
$\mathcal{C}_{D/U} = h \cdot \mathcal{O}_S$. Comme $U \to S$ est lisse,
le faisceau $\Omega_{U/S}$ est localement libre de type fini ; son image réciproque
$i^*\Omega_{U/S}$ est donc elle aussi localement libre de type fini. La première flèche de
la suite envoie le générateur libre $h$ sur la section $\text{d}h|_D$
de $i^*\Omega_{U/S}$, dont la valeur dans la fibre
$\Omega_{U/S, x} \otimes \kappa(x)$ est non nulle par hypothèse. Par exactitude à droite
de $\otimes \kappa(x)$, on conclut que
$$
\dim_{\kappa(x)} \left( \Omega_{D/S, x} \otimes \kappa(x) \right)
=
\dim_{\kappa(x)} \left( \Omega_{U/S, x} \otimes \kappa(x) \right) - 1.
$$
D'après
Morphismes, lemme \ref{morphisms-lemma-smooth-at-point},
on voit que $\Omega_{U/S, x} \otimes \kappa(x)$ peut être engendré par
au plus $\dim_x(U_s)$ éléments. D'après la formule affichée, on voit que
$\Omega_{D/S, x} \otimes \kappa(x)$ peut être engendré par au plus
$\dim_x(U_s) - 1$ éléments. Notons que
$\dim_x(D_s) = \dim_x(U_s) - 1$, par exemple parce que
$\dim(\mathcal{O}_{D_s, x}) = \dim(\mathcal{O}_{U_s, x}) - 1$ d'après
Algèbre, lemme \ref{algebra-lemma-one-equation}
(de plus, $D_s \subset U_s$ est un diviseur de Cartier effectif ; voir
Diviseurs, lemme \ref{divisors-lemma-relative-Cartier}),
puis en utilisant
Morphismes, lemme \ref{morphisms-lemma-dimension-fibre-at-a-point}.
On conclut ainsi que $\Omega_{D/S, x} \otimes \kappa(x)$ peut être engendré
par au plus $\dim_x(D_s)$ éléments, et l'on conclut que $D \to S$
est lisse en $x$ en appliquant de nouveau
Morphismes, lemme \ref{morphisms-lemma-smooth-at-point}.
Après avoir rétréci $U$, on obtient que $D \to S$ est lisse, ce qui conclut.
\end{proof}

\begin{lemma}
\label{more-morphisms-lemma-slice-smooth-once}
Soit $f : X \to S$ un morphisme de schémas.
Soit $x \in X$ un point d'image $s \in S$.
Supposons que
\begin{enumerate}
\item $f$ soit lisse en $x$, et
\item le morphisme
$$
\Omega_{X_s/s, x} \otimes_{\mathcal{O}_{X_s, x}} \kappa(x)
\longrightarrow
\Omega_{\kappa(x)/\kappa(s)}
$$
ait un noyau non nul.
\end{enumerate}
Alors il existe un voisinage ouvert affine $U \subset X$ de $x$
et un diviseur de Cartier effectif $D \subset U$ contenant $x$, tels que
$D \to S$ soit lisse.
\end{lemma}

\begin{proof}
Posons $k = \kappa(s)$ et $R = \mathcal{O}_{X_s, x}$.
Notons $\mathfrak m$ l'idéal maximal de $R$ et posons
$\kappa = R/\mathfrak m$, de sorte que $\kappa = \kappa(x)$.
Comme la formation des modules de différentielles commute à la localisation (voir
Algèbre, lemme \ref{algebra-lemma-differentials-localize}),
on a $\Omega_{X_s/s, x} = \Omega_{R/k}$. D'après
Algèbre, lemme \ref{algebra-lemma-differential-seq},
il existe une suite exacte
$$
\mathfrak m/\mathfrak m^2 \xrightarrow{\text{d}}
\Omega_{R/k} \otimes_R \kappa \to
\Omega_{\kappa/k} \to 0.
$$
Par conséquent, si (2) est satisfaite, il existe un élément $\overline{h} \in \mathfrak m$
tel que $\text{d}\overline{h}$ soit non nul. Choisissons un relèvement
$h \in \mathcal{O}_{X, x}$ de $\overline{h}$ et appliquons le
lemme \ref{more-morphisms-lemma-slice-smooth-given-element}.
\end{proof}

\begin{remark}
\label{more-morphisms-remark-necessary-condition-slice-smooth}
La seconde condition du
lemme \ref{more-morphisms-lemma-slice-smooth-once}
est nécessaire, même lorsque $x$ est un point fermé d'une fibre de dimension
strictement positive. En voici un exemple. Soit $k$ un corps
imparfait de caractéristique $p > 0$. Soit $a \in k$ un
élément qui n'est pas une puissance $p$-ième. Posons
$\mathfrak m = (x, y^p - a) \subset k[x, y]$. Ceci correspond à un point fermé
$w$ de $X = \mathbf{A}^2_k$. Posons $S = \mathbf{A}^1_k$ et
soit $f : X \to S$ le morphisme correspondant à $k[x] \to k[x, y]$.
Il n'existe alors aucun diagramme commutatif
$$
\xymatrix{
S' \ar[rr]_h \ar[rd]_g & & X \ar[ld]^f \\
& S
}
$$
avec $g$ \'etale et $w$ dans l'image de $h$. En effet, l'extension de corps
résiduels $\kappa(w)/\kappa(f(w))$ est purement inséparable,
tandis que, pour tout $s' \in S'$ tel que $g(s') = f(w)$, l'extension
$\kappa(s')/\kappa(f(w))$ serait séparable.
\end{remark}

\noindent
Si l'on suppose que l'extension de corps résiduels est séparable, le
phénomène de la
remarque \ref{more-morphisms-remark-necessary-condition-slice-smooth}
ne se produit pas. Voici l'énoncé précis.

\begin{lemma}
\label{more-morphisms-lemma-slice-smooth-once-separable-residue-field-extension}
Soit $f : X \to S$ un morphisme de schémas.
Soit $x \in X$ un point d'image $s \in S$.
Supposons que
\begin{enumerate}
\item $f$ soit lisse en $x$,
\item l'extension de corps résiduels $\kappa(x)/\kappa(s)$
soit séparable, et
\item $x$ ne soit pas un point générique de $X_s$.
\end{enumerate}
Alors il existe un voisinage ouvert affine $U \subset X$ de $x$
et un diviseur de Cartier effectif $D \subset U$ contenant $x$, tels que
$D \to S$ soit lisse.
\end{lemma}

\begin{proof}
Posons $k = \kappa(s)$ et $R = \mathcal{O}_{X_s, x}$.
Notons $\mathfrak m$ l'idéal maximal de $R$ et posons
$\kappa = R/\mathfrak m$, de sorte que $\kappa = \kappa(x)$.
Comme la formation des modules de différentielles commute à la localisation (voir
Algèbre, lemme \ref{algebra-lemma-differentials-localize}),
on a $\Omega_{X_s/s, x} = \Omega_{R/k}$. D'après l'hypothèse (2) et
Algèbre, lemme \ref{algebra-lemma-computation-differential},
le morphisme
$$
\text{d} :
\mathfrak m/\mathfrak m^2
\longrightarrow
\Omega_{R/k} \otimes_R \kappa(\mathfrak m)
$$
est injectif. L'hypothèse (3) implique que
$\mathfrak m/\mathfrak m^2 \not = 0$.
Il existe donc un élément $\overline{h} \in \mathfrak m$
tel que $\text{d}\overline{h}$ soit non nul. Choisissons un relèvement
$h \in \mathcal{O}_{X, x}$ de $\overline{h}$ et appliquons le
lemme \ref{more-morphisms-lemma-slice-smooth-given-element}.
\end{proof}

\noindent
Le sous-schéma $Z$ construit dans le lemme suivant est en fait une intersection
complète dans un voisinage ouvert affine de $x$. Si nous en avons un jour besoin,
nous formulerons explicitement un lemme séparé énonçant ce fait.

\begin{lemma}
\label{more-morphisms-lemma-slice-smooth}
Soit $f : X \to S$ un morphisme de schémas.
Soit $x \in X$ un point d'image $s \in S$.
Supposons que
\begin{enumerate}
\item $f$ soit lisse en $x$, et
\item $x$ soit un point fermé de $X_s$ et que $\kappa(s) \subset \kappa(x)$
soit séparable.
\end{enumerate}
Alors il existe une immersion $Z \to X$ dont l'image contient $x$ telle que
\begin{enumerate}
\item $Z \to S$ soit \'etale, et
\item $Z_s = \{x\}$ ensemblistement.
\end{enumerate}
\end{lemma}

\begin{proof}
On peut, et l'on va, remplacer $S$ par un voisinage ouvert affine de $s$.
On peut, et l'on va, remplacer $X$ par un voisinage ouvert affine de $x$
tel que $X \to S$ soit lisse.
Nous allons démontrer le lemme pour les morphismes lisses entre schémas affines
par récurrence sur $d = \dim_x(X_s)$.

\medskip\noindent
Le cas $d = 0$. Dans ce cas, nous montrons que l'on peut prendre pour $Z$
un voisinage ouvert de $x$. En effet, si $d = 0$, alors $X \to S$
est quasi-fini en $x$ ; voir
Morphismes, lemme \ref{morphisms-lemma-locally-quasi-finite-rel-dimension-0}.
Il existe donc un voisinage ouvert affine $U \subset X$
tel que $U \to S$ soit quasi-fini ; voir
Morphismes, lemme \ref{morphisms-lemma-quasi-finite-points-open}.
Ainsi, après avoir remplacé $X$ par $U$, on voit que
$X$ est quasi-fini et lisse sur $S$, donc
lisse de dimension relative $0$ sur $S$, et donc
\'etale sur $S$. De plus, la fibre $X_s$ est un ensemble fini
discret. Après avoir remplacé $X$ par un voisinage ouvert affine supplémentaire
de $X$, on voit donc que $f^{-1}(\{s\}) = \{x\}$ (car la topologie
de $X_s$ est induite par celle de $X$ ; voir
Schémas, lemme \ref{schemes-lemma-fibre-topological}).
Ceci démontre le lemme dans ce cas.

\medskip\noindent
Supposons maintenant que $d > 0$. Notons que, puisque $x$ est un point fermé de sa
fibre, l'extension $\kappa(x)/\kappa(s)$ est finie (d'après le
théorème des zéros de Hilbert ; voir
Morphismes, lemme \ref{morphisms-lemma-closed-point-fibre-locally-finite-type}).
On a donc $\Omega_{\kappa(x)/\kappa(s)} = 0$, puisque ceci vaut pour les
extensions de corps algébriques séparables.
On peut donc appliquer le
lemme \ref{more-morphisms-lemma-slice-smooth-once}
pour obtenir un diagramme
$$
\xymatrix{
D \ar[r] \ar[rrd] & U \ar[r] \ar[rd] & X \ar[d] \\
& & S
}
$$
avec $x \in D$. Notons que
$\dim_x(D_s) = \dim_x(X_s) - 1$, par exemple parce que
$\dim(\mathcal{O}_{D_s, x}) = \dim(\mathcal{O}_{X_s, x}) - 1$ d'après
Algèbre, lemme \ref{algebra-lemma-one-equation}
(de plus, $D_s \subset X_s$ est un diviseur de Cartier effectif ; voir
Diviseurs, lemme \ref{divisors-lemma-relative-Cartier}),
puis en utilisant
Morphismes, lemme \ref{morphisms-lemma-dimension-fibre-at-a-point}.
Ainsi, le morphisme $D \to S$ est lisse et
$\dim_x(D_s) = \dim_x(X_s) - 1 = d - 1$. Par l'hypothèse de récurrence,
on peut trouver une immersion $Z \to D$ comme souhaité, ce qui achève la démonstration.
\end{proof}

\begin{lemma}
\label{more-morphisms-lemma-etale-nbhd-dominates-smooth}
\begin{slogan}
Les morphismes lisses admettent des sections locales pour la topologie \'etale.
\end{slogan}
\begin{reference}
Voir \cite[Corollaire 17.16.3 (ii)]{EGA4}.
\end{reference}
Soit $f : X \to S$ un morphisme lisse de schémas.
Soit $s \in S$ un point de l'image de $f$.
Alors il existe un voisinage \'etale $(S', s') \to (S, s)$
et un $S$-morphisme $S' \to X$.
\end{lemma}

\begin{proof}[Première démonstration du lemme \ref{more-morphisms-lemma-etale-nbhd-dominates-smooth}]
Par hypothèse, $X_s \not = \emptyset$. D'après
Variétés, lemme \ref{varieties-lemma-smooth-separable-closed-points-dense},
il existe un point fermé $x \in X_s$ tel que $\kappa(x)$
soit une extension finie séparable de $\kappa(s)$.
Par conséquent, d'après le
lemme \ref{more-morphisms-lemma-slice-smooth},
il existe une immersion $Z \to X$ telle que $Z \to S$ soit \'etale et
que $x \in Z$. Prenons $(S' , s') = (Z, x)$.
\end{proof}

\begin{proof}[Seconde démonstration du lemme \ref{more-morphisms-lemma-etale-nbhd-dominates-smooth}]
Choisissons un point $x \in X$ tel que $f(x) = s$.
Choisissons un diagramme
$$
\xymatrix{
X \ar[d] & U \ar[l] \ar[d] \ar[r]_-\pi & \mathbf{A}^d_V \ar[ld] \\
S & V \ar[l]
}
$$
avec $\pi$ \'etale, $x \in U$ et $V = \Spec(R)$ affine ; voir
Morphismes, lemme \ref{morphisms-lemma-smooth-etale-over-affine-space}.
En particulier, $s \in V$. Le morphisme
$\pi : U \to \mathbf{A}^d_V$ est ouvert ; voir
Morphismes, lemme \ref{morphisms-lemma-etale-open}.
Ainsi, $W = \pi(U) \cap \mathbf{A}^d_s$ est un ouvert non vide de
$\mathbf{A}^d_s$. Soit $w \in W$ un point tel que l'extension $\kappa(s) \subset \kappa(w)$
soit finie et séparable ; voir
Variétés, lemme \ref{varieties-lemma-affine-space-over-field}.
D'après
Algèbre, lemme \ref{algebra-lemma-dim-affine-space},
il existe $d$ éléments
$\overline{f}_1, \ldots, \overline{f}_d \in \kappa(s)[x_1, \ldots, x_d]$
qui engendrent l'idéal maximal correspondant à $w$ dans
$\kappa(s)[x_1, \ldots, x_d]$.
Après avoir remplacé $R$ par une localisation principale,
on peut supposer qu'il existe des $f_1, \ldots, f_d \in R[x_1, \ldots, x_d]$
dont les images sont
$\overline{f}_1, \ldots, \overline{f}_d \in \kappa(s)[x_1, \ldots, x_d]$.
Considérons la $R$-algèbre
$$
R' = R[x_1, \ldots, x_d]/(f_1, \ldots, f_d)
$$
et posons $S' = \Spec(R')$. Par construction, nous avons une
immersion fermée $j : S' \to \mathbf{A}^d_V$ au-dessus de $V$.
Par construction, la fibre de $S' \to V$ au-dessus de $s$ est constituée d'un unique
point $s'$ dont le corps résiduel est une extension finie séparable de $\kappa(s)$.
Soit $\mathfrak q' \subset R'$ l'idéal premier correspondant. D'après
Algèbre, lemme \ref{algebra-lemma-localize-relative-complete-intersection},
on voit que $(R')_g$ est une intersection complète globale relative sur $R$
pour un certain $g \in R'$, $g \not \in \mathfrak q$.
Ainsi, $S' \to V$ est plat et de présentation finie dans un
voisinage de $s'$ ; voir
Algèbre, lemme \ref{algebra-lemma-relative-global-complete-intersection}.
Par construction, la fibre schématique
de $S' \to V$ au-dessus de $s$ est $\Spec(\kappa(s'))$. Il résulte donc de
Morphismes, lemme \ref{morphisms-lemma-etale-at-point},
que $S' \to S$ est \'etale en $s'$. Posons
$$
S'' = U \times_{\pi, \mathbf{A}^d_V, j} S'.
$$
Par construction, il existe un point $s'' \in S''$ qui s'envoie sur
$s'$ par la projection $p : S'' \to S'$. Notons que $p$ est \'etale
comme changement de base du morphisme \'etale $\pi$ ; voir
Morphismes, lemme \ref{morphisms-lemma-base-change-etale}.
Choisissons un petit voisinage ouvert affine $S''' \subset S''$ de $s''$
qui s'envoie dans le voisinage ouvert non vide de $s' \in S'$
où le morphisme $S' \to S$ est \'etale. Alors le voisinage \'etale
$(S''', s'') \to (S, s)$ est une solution au problème posé par le lemme.
\end{proof}

\noindent
Le lemme suivant montre que les faisceaux pour la topologie lisse sont
la même chose que les faisceaux pour la topologie \'etale.

\begin{lemma}
\label{more-morphisms-lemma-etale-dominates-smooth}
Soit $S$ un schéma. Soit $\mathcal{U} = \{S_i \to S\}_{i \in I}$ un recouvrement
lisse de $S$ ; voir
Topologies, définition \ref{topologies-definition-smooth-covering}.
Alors il existe un recouvrement \'etale $\mathcal{V} = \{T_j \to S\}_{j \in J}$
(voir
Topologies, définition \ref{topologies-definition-etale-covering})
qui raffine (voir
Sites, définition \ref{sites-definition-morphism-coverings})
$\mathcal{U}$.
\end{lemma}

\begin{proof}
Pour tout $s \in S$, il existe un $i \in I$ tel que $s$ appartienne à
l'image de $S_i \to S$. D'après le
lemme \ref{more-morphisms-lemma-etale-nbhd-dominates-smooth},
on peut trouver un morphisme \'etale $g_s : T_s \to S$ tel que $s \in g_s(T_s)$
et tel que $g_s$ se factorise par $S_i \to S$. Ainsi,
$\{T_s \to S\}$ est un recouvrement \'etale de $S$ qui raffine $\mathcal{U}$.
\end{proof}

\begin{lemma}
\label{more-morphisms-lemma-cover-smooth-by-special}
Soit $f : X \to S$ un morphisme lisse de schémas. Alors il existe un
recouvrement \'etale $\{U_i \to X\}_{i \in I}$ tel que $U_i \to S$
se factorise sous la forme $U_i \to V_i \to S$, où $V_i \to S$ est \'etale et
$U_i \to V_i$ est un morphisme lisse de schémas affines, qui
admet une section et dont les fibres sont géométriquement connexes.
\end{lemma}

\begin{proof}
Soit $s \in S$. D'après
Variétés, lemme \ref{varieties-lemma-smooth-separable-closed-points-dense},
l'ensemble des points fermés $x \in X_s$ tels que $\kappa(x)/\kappa(s)$
soit séparable est dense dans $X_s$. Il suffit donc de construire un
morphisme \'etale $U \to X$ dont l'image contient $x$
et tel que $U \to S$ se factorise de la manière décrite dans le lemme.
Pour ce faire, choisissons une immersion $Z \to X$ passant par $x$
telle que $Z \to S$ soit \'etale (lemme \ref{more-morphisms-lemma-slice-smooth}).
Après avoir remplacé $S$ par $Z$ et $X$ par $Z \times_S X$,
on voit que l'on peut supposer que $X \to S$ admet une section $\sigma : S \to X$
avec $\sigma(s) = x$. On peut alors d'abord remplacer $S$ par un voisinage
ouvert affine de $s$, puis remplacer $X$ par un voisinage ouvert affine
de $x$. Enfin, considérons le sous-ensemble
$X^0 \subset X$ de la section \ref{more-morphisms-section-connected-components}.
D'après les lemmes \ref{more-morphisms-lemma-connected-along-section-open} et
\ref{more-morphisms-lemma-connected-along-section-locally-constructible},
c'est un sous-schéma ouvert rétrocompact contenant $\sigma$
tel que les fibres de $X^0 \to S$ soient géométriquement connexes.
Si $X^0$ n'est pas affine, choisissons un ouvert affine $U \subset X^0$
contenant $x$. Comme $X^0 \to S$ est lisse, l'image de $U$
est ouverte. Choisissons un voisinage ouvert affine $V \subset S$ de $s$
contenu dans $\sigma^{-1}(U)$ et dans l'image de $U \to S$.
Enfin, le lecteur vérifiera que $U \cap f^{-1}(V) \to V$
possède toutes les propriétés souhaitées. Par exemple, $U \cap f^{-1}(V)$,
qui est égal à $U \times_S V$, est affine comme produit fibré de schémas
affines. En outre, les fibres géométriques de $U \cap f^{-1}(V) \to V$ sont des
sous-schémas ouverts non vides des fibres irréductibles de $X^0 \to S$
et sont donc connexes. Certains détails sont omis.
\end{proof}











\section{Voisinages \'etales et approximation d'Artin}
\label{more-morphisms-section-etale-nbhds-artin}

\noindent
Dans cette section, nous démontrons des résultats du type suivant : si deux
schémas pointés ont des anneaux locaux complets isomorphes, alors
ils ont des voisinages \'etales isomorphes. Nous nous appuierons
sur le théorème de Popescu ; voir
Lissage des morphismes d'anneaux, théorème \ref{smoothing-theorem-popescu}.

\begin{lemma}
\label{more-morphisms-lemma-map-approximation}
Soit $S$ un schéma localement noethérien. Soient $X$, $Y$ des
schémas localement de type
fini sur $S$. Soient $x \in X$ et $y \in Y$ des points situés au-dessus du
même point $s \in S$. Supposons que $\mathcal{O}_{S, s}$ soit un G-anneau.
Supposons en outre donné un homomorphisme local de $\mathcal{O}_{S, s}$-algèbres
$$
\varphi : \mathcal{O}_{Y, y} \longrightarrow \mathcal{O}_{X, x}^\wedge
$$
Pour tout $N \geq 1$,
il existe un voisinage \'etale élémentaire
$(U, u) \to (X, x)$ et un $S$-morphisme
$f : U \to Y$ envoyant $u$ sur $y$ tels que le diagramme
$$
\xymatrix{
\mathcal{O}_{X, x}^\wedge \ar[r] &
\mathcal{O}_{U, u}^\wedge \\
\mathcal{O}_{Y, y} \ar[r]^{f^\sharp_u} \ar[u]^\varphi &
\mathcal{O}_{U, u} \ar[u]
}
$$
soit commutatif modulo $\mathfrak m_u^N$.
\end{lemma}

\begin{proof}
La question est locale sur $X$ ; on peut donc supposer que $X$, $Y$, $S$ sont affines.
Écrivons $S = \Spec(R)$, $X = \Spec(A)$, $Y = \Spec(B)$.
Écrivons $B = R[x_1, \ldots, x_n]/(f_1, \ldots, f_m)$.
Soit $\mathfrak p \subset A$ l'idéal premier correspondant à $x$.
L'anneau local $\mathcal{O}_{X, x} = A_\mathfrak p$ est un G-anneau d'après
Compléments d'algèbre, proposition
\ref{more-algebra-proposition-finite-type-over-G-ring}.
L'homomorphisme $\varphi$ est un homomorphisme
$$
B_\mathfrak q^\wedge \longrightarrow A_\mathfrak p^\wedge
$$
où $\mathfrak q \subset B$ est l'idéal premier correspondant à $y$.
Soient $a_1, \ldots, a_n \in A_\mathfrak p^\wedge$ les images
de $x_1, \ldots, x_n$ par
$R[x_1, \ldots, x_n] \to B \to B_\mathfrak q^\wedge \to A_\mathfrak p^\wedge$.
On peut alors appliquer Lissage des morphismes d'anneaux, lemme
\ref{smoothing-lemma-approximation-property-variant},
pour obtenir un homomorphisme \'etale d'anneaux $A \to A'$, un idéal premier
$\mathfrak p' \subset A'$ et des éléments $b_1, \ldots, b_n \in A'$ tels que
$\kappa(\mathfrak p) = \kappa(\mathfrak p')$,
$a_i - b_i \in (\mathfrak p')^N(A'_{\mathfrak p'})^\wedge$, et
$f_j(b_1, \ldots, b_n) = 0$ pour $j = 1, \ldots, n$.
Ceci détermine un homomorphisme de $R$-algèbres $B \to A'$ envoyant la
classe de $x_i$ sur $b_i \in A'$. Ceci achève la démonstration
en prenant $U = \Spec(A') \to \Spec(B)$ pour le morphisme $f$
et en prenant $u = \mathfrak p'$.
\end{proof}

\begin{lemma}
\label{more-morphisms-lemma-isomorphism-approximation}
Soit $S$ un schéma localement noethérien. Soient $X$, $Y$ des
schémas localement de type
fini sur $S$. Soient $x \in X$ et $y \in Y$ des points situés au-dessus du
même point $s \in S$. Supposons que $\mathcal{O}_{S, s}$ soit un G-anneau.
Supposons donné un isomorphisme de $\mathcal{O}_{S, s}$-algèbres
$$
\varphi : \mathcal{O}_{Y, y}^\wedge \longrightarrow \mathcal{O}_{X, x}^\wedge
$$
entre les anneaux locaux complets. Alors, pour tout $N \geq 1$,
il existe des morphismes
$$
(X, x) \leftarrow (U, u) \rightarrow (Y, y)
$$
de schémas pointés au-dessus de $S$ tels que les deux flèches définissent des
voisinages \'etales élémentaires et tels que le diagramme
$$
\xymatrix{
& \mathcal{O}_{U, u}^\wedge \\
\mathcal{O}_{Y, y}^\wedge \ar[rr]^\varphi \ar[ru] & &
\mathcal{O}_{X, x}^\wedge \ar[lu]
}
$$
soit commutatif modulo $\mathfrak m_u^N$.
\end{lemma}

\begin{proof}
On peut supposer $N \geq 2$. Appliquons le lemme \ref{more-morphisms-lemma-map-approximation} pour obtenir
$(U, u) \to (X, x)$ et $f : (U, u) \to (Y, y)$.
Nous affirmons que $f$ est \'etale en $u$, ce qui achèvera la démonstration.
En fait, nous allons montrer que l'homomorphisme induit
$\mathcal{O}_{Y, y}^\wedge \to \mathcal{O}_{U, u}^\wedge$
est un isomorphisme. Une fois ceci démontré, le
lemme \ref{more-morphisms-lemma-lifting-along-artinian-at-point}
montrera que $f$ est lisse en $u$ et, bien entendu,
$f$ est également non ramifié en $u$ ; ainsi,
Morphismes, lemme \ref{morphisms-lemma-etale-smooth-unramified},
nous dit que $f$ est \'etale en $u$.
Pour un anneau local $(R, \mathfrak m)$, posons
$\text{Gr}_\mathfrak m(R) =
\bigoplus_{n \geq 0} \mathfrak m^n/\mathfrak m^{n + 1}$.
Pour démontrer l'affirmation, considérons le diagramme induit
d'anneaux gradués
$$
\xymatrix{
& \text{Gr}_{\mathfrak m_u}(\mathcal{O}_{U, u}) \\
\text{Gr}_{\mathfrak m_y}(\mathcal{O}_{Y, y}) \ar[rr]^\varphi \ar[ru] & &
\text{Gr}_{\mathfrak m_x}(\mathcal{O}_{X, x}) \ar[lu]
}
$$
Puisque $N \geq 2$, ce diagramme est effectivement commutatif, car les
algèbres graduées affichées sont engendrées en degré $1$ !
Par hypothèse, la flèche inférieure est un isomorphisme.
D'après Compléments d'algèbre, lemme \ref{more-algebra-lemma-flat-unramified}
(par exemple), l'homomorphisme
$\mathcal{O}_{X, x}^\wedge \to \mathcal{O}_{U, u}^\wedge$
est un isomorphisme ; la flèche nord-ouest
du diagramme est donc un isomorphisme. On en conclut que
$f$ induit un isomorphisme
$\text{Gr}_{\mathfrak m_x}(\mathcal{O}_{X, x}) \to
\text{Gr}_{\mathfrak m_y}(\mathcal{O}_{U, u})$.
Par récurrence, en utilisant les suites exactes courtes
$$
0 \to \text{Gr}^n_{\mathfrak m}(R) \to R/\mathfrak m^{n + 1} \to
R/\mathfrak m^n \to 0
$$
pour les deux anneaux locaux, on conclut (par le lemme du serpent)
que $f$ induit des isomorphismes
$\mathcal{O}_{Y, y}/\mathfrak m_y^n \to \mathcal{O}_{U, u}/\mathfrak m_u^n$
pour tout $n$, ce que nous voulions montrer.
\end{proof}

\begin{lemma}
\label{more-morphisms-lemma-relative-map-approximation-pre}
Donnons-nous $X \to S$, $Y \to T$, $x$, $s$, $y$, $t$, $\sigma$, $y_\sigma$ et
$\varphi$ comme suit : nous avons des morphismes de schémas
$$
\vcenter{
\xymatrix{
X \ar[d] & Y \ar[d] \\
S & T
}
}
\quad\text{avec les points}\quad
\vcenter{
\xymatrix{
x \ar[d] & y \ar[d] \\
s & t
}
}
$$
Ici, $S$ est localement noethérien et $T$ est de type fini sur $\mathbf{Z}$.
Les morphismes $X \to S$ et $Y \to T$ sont localement de type fini.
L'anneau local $\mathcal{O}_{S, s}$ est un G-anneau. L'homomorphisme
$$
\sigma : \mathcal{O}_{T, t} \longrightarrow \mathcal{O}_{S, s}^\wedge
$$
est un homomorphisme local. Posons
$Y_\sigma = Y \times_{T, \sigma} \Spec(\mathcal{O}_{S, s}^\wedge)$.
Ensuite, $y_\sigma$ est un point de $Y_\sigma$ qui s'envoie sur $y$ et sur
le point fermé de $\Spec(\mathcal{O}_{S, s}^\wedge)$. Enfin,
$$
\varphi :
\mathcal{O}_{X, x}^\wedge
\longrightarrow
\mathcal{O}_{Y_\sigma, y_\sigma}^\wedge
$$
est un isomorphisme de $\mathcal{O}_{S, s}^\wedge$-algèbres.
Dans cette situation, il existe un diagramme commutatif
$$
\xymatrix{
X \ar[d] &
W \ar[l] \ar[rd] \ar[rr] & &
Y \times_{T, \tau} V \ar[r] \ar[ld] & Y \ar[d] \\
S & &
V \ar[ll] \ar[rr]^\tau & &
T
}
$$
de schémas, ainsi que des points $w \in W$, $v \in V$, tels que
\begin{enumerate}
\item $(V, v) \to (S, s)$ est un voisinage \'etale élémentaire,
\item $(W, w) \to (X, x)$ est un voisinage \'etale élémentaire, et
\item $\tau(v) = t$.
\end{enumerate}
Soit $y_\tau \in Y \times_T V$ le point qui correspond à $y_\sigma$
par l'identification $(Y_\sigma)_s = (Y \times_T V)_v$.
Alors
\begin{enumerate}
\item[(4)] $(W, w) \to (Y \times_{T, \tau} V, y_\tau)$ est un voisinage
\'etale élémentaire.
\end{enumerate}
\end{lemma}

\begin{proof}
Notons $X_\sigma = X \times_S \Spec(\mathcal{O}_{S, s}^\wedge)$
et $x_\sigma \in X_\sigma$ l'unique point situé au-dessus de $x$.
Observons que $\mathcal{O}_{S, s}^\wedge$ est un G-anneau d'après
Compléments d'algèbre, proposition
\ref{more-algebra-proposition-Noetherian-complete-G-ring}.
D'après le lemme \ref{more-morphisms-lemma-isomorphism-approximation},
on peut choisir
$$
(X_\sigma, x_\sigma) \leftarrow (U, u) \rightarrow (Y_\sigma, y_\sigma)
$$
où les deux flèches sont des voisinages \'etales élémentaires.

\medskip\noindent
Après avoir remplacé $S$ par un voisinage ouvert de $s$, on peut
supposer que $S = \Spec(R)$ est affine. Puisque $\mathcal{O}_{S, s}$
est un G-anneau, le théorème \ref{smoothing-theorem-popescu} de Lissage des morphismes d'anneaux montre que
l'anneau $\mathcal{O}_{S, s}^\wedge$ est une colimite filtrante de
$R$-algèbres lisses. On peut donc écrire
$$
\Spec(\mathcal{O}_{S, s}^\wedge) = \lim S_i
$$
comme limite filtrante de schémas affines $S_i$ lisses sur $S$.
Notons $s_i \in S_i$ l'image du point fermé de
$\Spec(\mathcal{O}_{S, s}^\wedge)$. Observons que $\kappa(s) = \kappa(s_i)$.
Posons $X_i = X \times_S S_i$ et notons $x_i \in X_i$ l'unique
point s'envoyant sur $x$. Notons que $\kappa(x) = \kappa(x_i)$.
Comme $T$ est de type fini sur $\mathbf{Z}$, d'après Limites, proposition
\ref{limits-proposition-characterize-locally-finite-presentation},
on peut choisir un $i$ et un morphisme $\sigma_i : (S_i, s_i) \to (T, t)$
de schémas pointés dont la composée avec
$\Spec(\mathcal{O}_{S, s}^\wedge) \to S_i$ est égale à $\sigma$.
Posons $Y_i = Y \times_T S_i$ et notons $y_i$ l'image de
$y_\sigma$. Notons que $\kappa(y_i) = \kappa(y_\sigma)$.
D'après Limites, lemme \ref{limits-lemma-descend-finite-presentation},
on peut choisir un $i$ et un diagramme
$$
\xymatrix{
X_i \ar[rd] &
U_i \ar[l] \ar[d] \ar[r] &
Y_i \ar[ld] \\
& S_i
}
$$
dont le changement de base à $\Spec(\mathcal{O}_{S, s}^\wedge)$
redonne $X_\sigma \leftarrow U \rightarrow Y_\sigma$.
D'après Limites, lemme \ref{limits-lemma-descend-etale},
quitte à augmenter $i$, on peut supposer que les morphismes
$X_i \leftarrow U_i \rightarrow Y_i$ sont \'etales.
Soit $u_i \in U_i$ l'image de $u$. Alors $u_i \mapsto x_i$,
d'où
$\kappa(x) = \kappa(x_\sigma) = \kappa(u) \supset \kappa(u_i) \supset
\kappa(x_i) = \kappa(x)$, et l'on voit que $\kappa(u_i) = \kappa(x_i)$.
Ainsi, $(X_i, x_i) \leftarrow (U_i, u_i)$ est un voisinage
\'etale élémentaire. Comme, de plus, $\kappa(y_i) = \kappa(y_\sigma) = \kappa(u)$,
on voit que $(U_i, u_i) \to (Y_i, y_i)$ est également un voisinage
\'etale élémentaire.

\medskip\noindent
À ce stade, nous avons construit un diagramme
$$
\xymatrix{
X \ar[d] &
X \times_S S_i \ar[l] \ar[rd] &
U_i \ar[l] \ar[r] \ar[d] &
Y \times_T S_i \ar[r] \ar[ld] &
Y \ar[d] \\
S & &
S_i \ar[ll] \ar[rr] & &
T
}
$$
comme dans l'énoncé du lemme, à ceci près que $S_i \to S$ est lisse.
D'après le lemme \ref{more-morphisms-lemma-slice-smooth}, et après avoir rétréci $S_i$,
on peut supposer qu'il existe un sous-schéma fermé $V \subset S_i$ passant
par $s_i$ tel que $V \to S$ soit \'etale.
En prenant pour $W$ l'image réciproque schématique de $V$
dans $U_i$, on conclut.
\end{proof}

\noindent
Nous recommandons vivement au lecteur de ne pas lire le reste de cette section.

\begin{lemma}
\label{more-morphisms-lemma-relative-map-approximation}
Considérons un diagramme
$$
\vcenter{
\xymatrix{
X \ar[d] & Y \ar[d] \\
S & T \ar[l]
}
}
\quad\text{avec les points}\quad
\vcenter{
\xymatrix{
x \ar[d] & y \ar[d] \\
s & t \ar[l]
}
}
$$
où $S$ est un schéma localement noethérien et les morphismes sont
localement de type fini. Supposons que $\mathcal{O}_{S, s}$ soit un G-anneau.
Supposons en outre donné un homomorphisme local de $\mathcal{O}_{S, s}$-algèbres
$$
\sigma : \mathcal{O}_{T, t} \longrightarrow \mathcal{O}_{S, s}^\wedge
$$
et un homomorphisme local de $\mathcal{O}_{S, s}$-algèbres
$$
\varphi :
\mathcal{O}_{X, x}
\longrightarrow
\mathcal{O}_{Y_\sigma, y_\sigma}^\wedge
$$
où $Y_\sigma = Y \times_{T, \sigma} \Spec(\mathcal{O}_{S, s}^\wedge)$
et $y_\sigma$ est l'unique point de $Y_\sigma$ situé au-dessus de $y$.
Pour tout $N \geq 1$, il existe un diagramme commutatif
$$
\xymatrix{
X \ar[d] & X \times_S V \ar[l] \ar[rd] &
W \ar[l]^-f \ar[r] \ar[d] &
Y \times_{T, \tau} V \ar[r] \ar[ld] & Y \ar[d] \\
S & & V \ar[ll] \ar[rr]^\tau & & T
}
$$
de schémas au-dessus de $S$, ainsi que des points $w \in W$, $v \in V$, tels que
\begin{enumerate}
\item $v \mapsto s$, $\tau(v) = t$, $f(w) = (x, v)$, et $w \mapsto (y, v)$,
\item $(V, v) \to (S, s)$ est un voisinage \'etale élémentaire,
\item le diagramme
$$
\xymatrix{
\mathcal{O}_{S, s}^\wedge \ar[r] & \mathcal{O}_{V, v}^\wedge \\
\mathcal{O}_{T, t} \ar[r]^{\tau^\sharp_v} \ar[u]_\sigma &
\mathcal{O}_{V, v} \ar[u]
}
$$
est commutatif modulo $\mathfrak m_v^N$,
\item $(W, w) \to (Y \times_{T, \tau} V, (y, v))$ est un
voisinage \'etale élémentaire,
\item le diagramme
$$
\xymatrix{
\mathcal{O}_{X, x} \ar[r]_\varphi &
\mathcal{O}_{Y_\sigma, y_\sigma}^\wedge \ar[r] &
\mathcal{O}_{Y_\sigma, y_\sigma}/\mathfrak m_{y_\sigma}^N \ar@{=}[r] &
\mathcal{O}_{Y \times_{T, \tau} V, (y, v)}/\mathfrak m_{(y, v)}^N
\ar[d]_{\cong} \\
\mathcal{O}_{X, x} \ar[r] \ar@{=}[u] &
\mathcal{O}_{X \times_S V, (x, v)} \ar[r]^{f^\sharp_w} &
\mathcal{O}_{W, w} \ar[r] &
\mathcal{O}_{W, w}/\mathfrak m_w^N
}
$$
est commutatif. L'égalité vient du fait que
$Y_\sigma$ et $Y \times_{T, \tau} V$ sont canoniquement isomorphes au-dessus de
$\mathcal{O}_{V, v}/\mathfrak m_v^N = \mathcal{O}_{S, s}/\mathfrak m_s^N$
d'après les parties (2) et (3).
\end{enumerate}
\end{lemma}

\begin{proof}
Après avoir remplacé $X$, $S$, $T$, $Y$ par des sous-schémas ouverts affines, on
peut supposer que le diagramme de l'énoncé du lemme provient de
l'application de $\Spec$ à un diagramme
$$
\vcenter{
\xymatrix{
A & B \\
R \ar[u] \ar[r] & C \ar[u]
}
}
\quad\text{avec les idéaux premiers}\quad
\vcenter{
\xymatrix{
\mathfrak p_A & \mathfrak p_B \\
\mathfrak p_R \ar@{-}[u] \ar@{-}[r] & \mathfrak p_C \ar@{-}[u]
}
}
$$
d'anneaux noethériens et d'homomorphismes d'anneaux de type fini.
Dans cette démonstration, tout anneau $E$ sera une $R$-algèbre noethérienne munie
d'un idéal premier $\mathfrak p_E$ au-dessus de $\mathfrak p_R$, et tous les homomorphismes d'anneaux
seront des homomorphismes de $R$-algèbres compatibles avec les idéaux premiers donnés. De plus,
si nous écrivons $E^\wedge$, nous entendons le complété du localisé
de $E$ en $\mathfrak p_E$. Nous utiliserons aussi sans autre mention
qu'un homomorphisme \'etale d'anneaux $E_1 \to E_2$ tel que
$\kappa(\mathfrak p_{E_1}) = \kappa(\mathfrak p_{E_2})$ induit
un isomorphisme $E_1^\wedge = E_2^\wedge$ d'après
Compléments d'algèbre, lemme \ref{more-algebra-lemma-flat-unramified}.

\medskip\noindent
Avec ces notations, $\sigma$ et $\varphi$ correspondent aux homomorphismes d'anneaux
$$
\sigma : C \to R^\wedge
\quad\text{et}\quad
\varphi : A \longrightarrow (B \otimes_{C, \sigma} R^\wedge)^\wedge
$$
Voici un diagramme :
$$
\xymatrix{
A \ar@/^1em/[rrr]^\varphi &
B \ar[r] &
B \otimes_{C, \sigma} R^\wedge \ar[r] &
(B \otimes_{C, \sigma} R^\wedge)^\wedge \\
R \ar[r] \ar[u] &
C \ar[r]^\sigma \ar[u] & R^\wedge \ar[u] \ar[ru]
}
$$
Observons que $R^\wedge$ est un G-anneau d'après
Compléments d'algèbre, proposition
\ref{more-algebra-proposition-Noetherian-complete-G-ring}.
Ainsi, $B \otimes_{C, \sigma} R^\wedge$ est un G-anneau d'après
Compléments d'algèbre, proposition
\ref{more-algebra-proposition-finite-type-over-G-ring}.
D'après le lemme \ref{more-morphisms-lemma-map-approximation} (sous sa forme algébrique),
il existe un homomorphisme \'etale d'anneaux $B \otimes_{C, \sigma} R^\wedge \to B'$,
qui induit un isomorphisme
$\kappa(\mathfrak p_{B \otimes_{C, \sigma} R^\wedge})
\to \kappa(\mathfrak p_{B'})$,
et un homomorphisme de $R$-algèbres $A \to B'$ tel que la composée
$$
A \to B' \to (B')^\wedge = (B \otimes_{C, \sigma} R^\wedge)^\wedge
$$
soit égale à $\varphi$ modulo
$(\mathfrak p_{(B \otimes_{C, \sigma} R^\wedge)^\wedge})^N$.
On peut donc remplacer $\varphi$ par cette composée, car la seule
façon dont $\varphi$ intervient dans la conclusion est la condition de
commutativité de la partie (5) de l'énoncé du lemme.
Diagramme :
$$
\xymatrix{
& & B' \ar[r] & (B')^\wedge \ar@{=}[d] \\
A \ar[rru] &
B \ar[r] &
B \otimes_{C, \sigma} R^\wedge \ar[r] \ar[u] &
(B \otimes_{C, \sigma} R^\wedge)^\wedge \\
R \ar[r] \ar[u] &
C \ar[r]^\sigma \ar[u] & R^\wedge \ar[u] \ar[ru]
}
$$
Ensuite, on utilise le fait que $R^\wedge$ est une colimite filtrante de
$R$-algèbres lisses (Lissage des morphismes d'anneaux, théorème \ref{smoothing-theorem-popescu}),
car $R_{\mathfrak p_R}$ est un G-anneau par hypothèse. Comme $C$ est de
présentation finie sur $R$, on obtient une factorisation
$$
C \to R' \to R^\wedge
$$
pour un certain homomorphisme lisse $R \to R'$ ; voir
Algèbre, lemme \ref{algebra-lemma-characterize-finite-presentation}.
Quitte à remplacer $R'$ par un terme ultérieur, on peut supposer qu'il existe
une algèbre \'etale sur $B \otimes_C R'$, notée $B''$, dont le changement de base
à $B \otimes_{C, \sigma} R^\wedge$ est $B'$ ; voir
Algèbre, lemme \ref{algebra-lemma-etale}.
Alors $B'$ est la colimite filtrante de ces $B''$, et l'on
conclut que, quitte à remplacer $R'$ par un terme ultérieur, on peut supposer qu'il existe
un homomorphisme de $R$-algèbres $A \to B''$ tel que $A \to B'' \to B'$ soit
l'homomorphisme construit précédemment (même référence que ci-dessus). Diagramme :
$$
\xymatrix{
& & B'' \ar[r] & B' \ar[r] & (B')^\wedge \ar@{=}[d] \\
A \ar[rru] &
B \ar[r] &
B \otimes_C R' \ar[r] \ar[u] &
B \otimes_{C, \sigma} R^\wedge \ar[r] \ar[u] &
(B \otimes_{C, \sigma} R^\wedge)^\wedge \\
R \ar[r] \ar[u] &
C \ar[r] \ar[u] &
R' \ar[r] \ar[u] &
R^\wedge \ar[u] \ar[ru]
}
$$
et
$$
B' = B'' \otimes_{(B \otimes_C R')} (B \otimes_{C, \sigma} R^\wedge)
$$
Ceci signifie que l'on peut remplacer $C$ par $R'$, $\sigma : C \to R^\wedge$ par
$R' \to R^\wedge$ et $B$ par $B''$, de sorte que l'on se ramène au diagramme
$$
\xymatrix{
A \ar[r] &
B \ar[r] &
B \otimes_{C, \sigma} R^\wedge \\
R \ar[r] \ar[u] &
C \ar[r]^\sigma \ar[u] & R^\wedge \ar[u]
}
$$
où $\varphi$ est égal à la composée des flèches horizontales
suivie de l'homomorphisme canonique de $B \otimes_{C, \sigma} R^\wedge$
vers son complété.
La dernière étape de la démonstration consiste à appliquer le lemme \ref{more-morphisms-lemma-map-approximation}
(ou sa démonstration)
une fois de plus à $\Spec(C)$ et $\Spec(R)$ au-dessus de $\Spec(R)$ et à l'homomorphisme
$C \to R^\wedge$. Le lemme fournit un homomorphisme d'anneaux $C \to D$
tel que $R \to D$ soit \'etale, que
$\kappa(\mathfrak p_R) = \kappa(\mathfrak p_D)$, et que
$$
C \to D \to D^\wedge = R^\wedge
$$
coïncide avec $\sigma : C \to R^\wedge$ modulo $(\mathfrak p_{R^\wedge})^N$.
Nous pouvons alors prendre
$$
V = \Spec(D)
\quad\text{et}\quad
W = \Spec(B \otimes_C D)
$$
comme solution au problème posé par le lemme. En effet, le diagramme
$$
\xymatrix{
A \ar[r] &
B \otimes_{C, \sigma} R^\wedge \ar[r] &
B \otimes_{C, \sigma} R^\wedge/(\mathfrak p_{R^\wedge})^N \ar@{=}[r] &
B \otimes_C D/(\mathfrak p_D)^N \\
A \ar@{=}[u] \ar[r] &
A \otimes_R D \ar[r] &
B \otimes_R D \ar[r] &
B \otimes_C D/(\mathfrak p_D)^N \ar@{=}[u]
}
$$
est commutatif parce que la composée $C \to D \to D^\wedge = R^\wedge$ coïncide avec
$\sigma$ modulo $(\mathfrak p_{R^\wedge})^N$. Ceci démontre la partie (5),
et les autres propriétés résultent immédiatement de la construction.
\end{proof}

\begin{lemma}
\label{more-morphisms-lemma-control-agreement}
Soit $T \to S$ un morphisme de type fini entre schémas noethériens.
Soit $t \in T$ d'image $s \in S$, et soit
$\sigma : \mathcal{O}_{T, t} \to \mathcal{O}_{S, s}^\wedge$
un morphisme local de $\mathcal{O}_{S, s}$-algèbres. Pour tout $N \geq 1$,
il existe un morphisme de type fini $(T', t') \to (T, t)$
tel que $\sigma$ se factorise par
$\mathcal{O}_{T, t} \to \mathcal{O}_{T', t'}$
et tel que, pour tout morphisme local de $\mathcal{O}_{S, s}$-algèbres
$\sigma' : \mathcal{O}_{T, t} \to \mathcal{O}_{S, s}^\wedge$
qui se factorise par $\mathcal{O}_{T, t} \to \mathcal{O}_{T', t'}$,
les morphismes $\sigma$ et $\sigma'$ coïncident modulo $\mathfrak m_s^N$.
\end{lemma}

\begin{proof}
Nous pouvons supposer $S$ et $T$ affines. Écrivons $S = \Spec(R)$ et
$T = \Spec(C)$. Soient $c_1, \ldots, c_n \in C$ des générateurs
de $C$ comme $R$-algèbre. Soit $\mathfrak p \subset R$ l'idéal
premier correspondant à $s$. Écrivons $\mathfrak p = (f_1, \ldots, f_m)$.
Après avoir remplacé $R$ par un localisé principal
(afin de chasser les dénominateurs dans $R_\mathfrak p$),
nous pouvons supposer qu'il existe
$r_1, \ldots, r_n \in R$ et $a_{i, I} \in \mathcal{O}_{S, s}^\wedge$,
où $I = (i_1, \ldots, i_m)$ avec $\sum i_j = N$, tels que
$$
\sigma(c_i) = r_i + \sum\nolimits_I a_{i, I} f_1^{i_1} \ldots f_m^{i_m}
$$
dans $\mathcal{O}_{S, s}^\wedge$. Considérons alors
$$
C' = C[t_{i, I}]/
\left(c_i - r_i - \sum\nolimits_I t_{i, I} f_1^{i_1} \ldots f_m^{i_m}\right)
$$
avec $\mathfrak p' = \mathfrak pC' + (t_{i, I})$, et la factorisation
de $\sigma : C \to \mathcal{O}_{S, s}^\wedge$ par $C'$
obtenue en envoyant $t_{i, I}$ sur $a_{i, I}$.
Le choix $T' = \Spec(C')$ convient, car tout $\sigma'$ comme dans l'énoncé
du lemme envoie $c_i$ sur $r_i$ modulo la
puissance $N$-ième de l'idéal maximal.
\end{proof}

\begin{lemma}
\label{more-morphisms-lemma-control-graded}
Soient $Y \to T \to S$ des morphismes de type fini entre schémas noethériens.
Soit $t \in T$ d'image $s \in S$, et soit
$\sigma : \mathcal{O}_{T, t} \to \mathcal{O}_{S, s}^\wedge$
un morphisme local de $\mathcal{O}_{S, s}$-algèbres.
Il existe un morphisme de type fini $(T', t') \to (T, t)$
tel que $\sigma$ se factorise par
$\mathcal{O}_{T, t} \to \mathcal{O}_{T', t'}$
et tel que, pour tout morphisme local de $\mathcal{O}_{S, s}$-algèbres
$\sigma' : \mathcal{O}_{T, t} \to \mathcal{O}_{S, s}^\wedge$
qui se factorise par $\mathcal{O}_{T, t} \to \mathcal{O}_{T', t'}$,
les immersions fermées
$$
Y \times_{T, \sigma} \Spec(\mathcal{O}_{S, s}^\wedge) = Y_\sigma
\longleftarrow Y_t \longrightarrow
Y_{\sigma'} =
Y \times_{T, \sigma'} \Spec(\mathcal{O}_{S, s}^\wedge)
$$
ont des algèbres conormales isomorphes.
\end{lemma}

\begin{proof}
Une observation utile est que $\kappa(s) = \kappa(t)$ par l'existence de
$\sigma$. Remarquons que l'énoncé a un sens, puisque les fibres de $Y_\sigma$
et de $Y_{\sigma'}$ au-dessus de $s \in \Spec(\mathcal{O}_{S, s}^\wedge)$
sont toutes deux canoniquement isomorphes à $Y_t$. Nous considérerons la propriété
«~$\sigma'$ se factorise par $\mathcal{O}_{T, t} \to \mathcal{O}_{T', t'}$~»
comme une contrainte sur $\sigma'$. Si nous avons plusieurs telles contraintes,
disons données par $(T'_i, t'_i) \to (T, t)$, $i = 1, \ldots, n$,
nous pouvons les combiner en considérant
$(T'_1 \times_T \ldots \times_T T'_n, (t'_1, \ldots, t'_n)) \to
(T, t)$. Nous utiliserons ceci sans autre mention dans la suite.

\medskip\noindent
D'après le lemme \ref{more-morphisms-lemma-control-agreement}, nous pouvons supposer que tout $\sigma'$
comme dans l'énoncé du lemme est égal à $\sigma$ modulo
$\mathfrak m_s^2$. Remarquons que l'algèbre conormale de $Y_t$ dans
$Y_\sigma$ n'est autre que la $\mathcal{O}_{Y_t}$-algèbre graduée quasi-cohérente
$$
\bigoplus\nolimits_{n \geq 0}
\mathfrak m_s^n\mathcal{O}_{Y_\sigma}/
\mathfrak m_s^{n + 1}\mathcal{O}_{Y_\sigma}
$$
et de même pour $Y_{\sigma'}$. Puisque $\sigma$ et $\sigma'$
coïncident modulo $\mathfrak m_s^2$, nous voyons que ces deux algèbres
sont les mêmes en degrés $0$ et $1$. D'autre part, ces
algèbres conormales sont engendrées en degré $1$ sur leur partie de degré $0$.
Par conséquent, s'il existe un isomorphisme prolongeant l'isomorphisme
qui vient d'être construit en degrés $0$ et $1$, alors il est unique.

\medskip\noindent
Nous pouvons supposer $S$ et $T$ affines. Soit $Y = Y_1 \cup \ldots \cup Y_n$
un recouvrement ouvert affine. Si nous pouvons construire $(T_i', t'_i) \to (T, t)$
comme dans le lemme de telle sorte que l'isomorphisme souhaité (voir le paragraphe précédent)
existe pour $Y_i \to T \to S$ et $\sigma$, alors ces isomorphismes se recollent par
unicité et donnent le résultat pour $Y \to T$. Nous pouvons donc supposer $Y$ affine.

\medskip\noindent
Écrivons $S = \Spec(R)$, $T = \Spec(C)$ et $Y = \Spec(B)$.
Choisissons une présentation $B = C[x_1, \ldots, x_n]/(f_1, \ldots, f_m)$.
Notons $R^\wedge = \mathcal{O}_{S, s}^\wedge$.
Soient $a_{kj} \in R^\wedge[x_1, \ldots, x_n]$ des polynômes
tels que
$$
\sum\nolimits_{j = 1, \ldots, m} a_{kj}\sigma(f_j) = 0,\quad
\text{pour }k = 1, \ldots, K
$$
forment un système de générateurs du module des relations entre
les $\sigma(f_j) \in R^\wedge[x_1, \ldots, x_n]$.
Nous avons donc une suite exacte
\begin{equation}
\label{more-morphisms-equation-resolution}
R^\wedge[x_1, \ldots, x_n]^{\oplus K} \to
R^\wedge[x_1, \ldots, x_n]^{\oplus m} \to
R^\wedge[x_1, \ldots, x_n] \to B \otimes_{C, \sigma} R^\wedge \to 0
\end{equation}
Soit $c$ un entier pour lequel le lemme d'Artin--Rees s'applique
à la fois au premier et au second morphisme de cette suite, ainsi qu'à l'idéal
$\mathfrak m_{R^\wedge}R^\wedge[x_1, \ldots, x_n]$, comme dans
Compléments d'algèbre, section \ref{more-algebra-section-artin-rees}.
Écrivons
$$
a_{kj} = \sum\nolimits_{I \in \Omega} a_{kj, I} x^I
\quad\text{et}\quad
f_j = \sum\nolimits_{I \in \Omega} f_{j, I} x^I
$$
en notation multi-indicielle, où $a_{kj, I} \in R^\wedge$, $f_{j, I} \in C$,
et $\Omega$ est un ensemble fini de multi-indices. Nous voyons alors que
$$
\sum\nolimits_{j = 1, \ldots, m,\ I, I' \in \Omega,\ I + I' = I''}
a_{kj, I} \sigma(f_{j, I'}) = 0,\quad
I''\text{ un multi-indice}
$$
dans $R^\wedge$. Prenons donc
$$
C' = C[t_{jk, I}]/
\left(
\sum\nolimits_{j = 1, \ldots, m,\ I, I' \in \Omega,\ I + I' = I''}
t_{kj, I} f_{j, I'},\ I''\text{ un multi-indice}\right)
$$
Alors $\sigma$ se factorise par un morphisme $\tilde\sigma : C' \to R^\wedge$
envoyant $t_{kj, I}$ sur $a_{jk, I}$.
Ainsi, $T' = \Spec(C')$ possède un point $t' \in T'$ tel que
$\sigma$ se factorise par $\mathcal{O}_{T, t} \to \mathcal{O}_{T', t'}$.
Soit $t_{kj} = \sum t_{kj, I} x^I$ dans $C'[x_1, \ldots, x_n]$.
Nous voyons alors que nous avons un complexe
\begin{equation}
\label{more-morphisms-equation-resolution-new}
C'[x_1, \ldots, x_n]^{\oplus K} \to
C'[x_1, \ldots, x_n]^{\oplus m} \to
C'[x_1, \ldots, x_n] \to B \otimes_C C' \to 0
\end{equation}
qui est exact en $C'[x_1, \ldots, x_n]$ et dont le changement de base
par $\tilde\sigma$ donne (\ref{more-morphisms-equation-resolution}).

\medskip\noindent
D'après le lemme \ref{more-morphisms-lemma-control-agreement},
nous pouvons trouver un morphisme supplémentaire $(T'', t'') \to (T', t')$
tel que $\tilde\sigma$
se factorise par $\mathcal{O}_{T', t'} \to \mathcal{O}_{T'', t''}$
et tel que, si $\sigma' : C \to R^\wedge$
se factorise par $\mathcal{O}_{T'', t''}$, alors le morphisme induit
$\tilde \sigma' : C' \to R^\wedge$
coïncide modulo $\mathfrak m_s^{c + 1}$ avec $\tilde \sigma$.
Ainsi, si $\sigma'$ est un tel morphisme, nous obtenons un complexe
$$
R^\wedge[x_1, \ldots, x_n]^{\oplus K} \to
R^\wedge[x_1, \ldots, x_n]^{\oplus m} \to
R^\wedge[x_1, \ldots, x_n] \to B \otimes_{C, \sigma'} R^\wedge \to 0
$$
sur $R^\wedge[x_1, \ldots, x_n]$ en appliquant $\tilde\sigma'$
aux polynômes $t_{kj}$ et $f_j$. Autrement dit, il
s'agit du changement de base du complexe (\ref{more-morphisms-equation-resolution-new})
par $\tilde\sigma'$. Les matrices définissant ce complexe
sont congruentes modulo $\mathfrak m_s^{c + 1}$ aux matrices
définissant le complexe (\ref{more-morphisms-equation-resolution}), car
$\tilde \sigma$ et $\tilde \sigma'$ sont congrus modulo
$\mathfrak m_s^{c + 1}$. Puisque (\ref{more-morphisms-equation-resolution}) est exacte,
nous pouvons appliquer
Compléments d'algèbre, lemme \ref{more-algebra-lemma-approximate-complex-graded}
pour conclure que
$$
\text{Gr}_{\mathfrak m_s}(B \otimes_{C, \sigma'} R^\wedge)
\cong
\text{Gr}_{\mathfrak m_s}(B \otimes_{C, \sigma} R^\wedge)
$$
comme souhaité.
\end{proof}

\begin{lemma}
\label{more-morphisms-lemma-relative-isomorphism-approximation}
Avec les notations et les hypothèses du lemme \ref{more-morphisms-lemma-relative-map-approximation},
supposons que $\varphi$ induise un isomorphisme sur les complétés.
Nous pouvons alors choisir notre diagramme de telle sorte que $f$ soit \'etale.
\end{lemma}

\begin{proof}
Nous pouvons supposer $N \geq 2$ et remplacer $(T, t)$ par $(T', t')$ comme dans le
lemme \ref{more-morphisms-lemma-control-graded}. Puisque $(V, v) \to (S, s)$ est un voisinage
\'etale élémentaire, il en va de même de $(X \times_S V, (x, v)) \to (X, x)$.
Ainsi, $\mathcal{O}_{X, x} \to \mathcal{O}_{X \times_S V, (x, v)}$
induit un isomorphisme sur les complétés d'après
Compléments d'algèbre, lemme \ref{more-algebra-lemma-flat-unramified}.
Nous affirmons que $\mathcal{O}_{X, x} \to \mathcal{O}_{W, w}$ induit
un isomorphisme sur les complétés. Une fois ceci démontré, le
lemme \ref{more-morphisms-lemma-lifting-along-artinian-at-point}
montrera que $f$ est lisse en $w$ et, bien entendu,
$f$ est aussi non ramifié en $u$; par conséquent,
Morphismes, lemme \ref{morphisms-lemma-etale-smooth-unramified}
nous dit que $f$ est \'etale en $w$.

\medskip\noindent
Nous utilisons d'abord la commutativité de la partie (5) du
lemme \ref{more-morphisms-lemma-relative-map-approximation}
pour voir que, pour $i \leq N$, il existe un diagramme commutatif
$$
\xymatrix{
\text{Gr}^i_{\mathfrak m_x}(\mathcal{O}_{X, x})
\ar[r]_-\varphi &
\text{Gr}^i_{\mathfrak m_{y_\sigma}}(\mathcal{O}_{Y_\sigma, y_\sigma}^\wedge)
\ar@{=}[r] &
\text{Gr}^i_{\mathfrak m_{(y, v)}}(\mathcal{O}_{Y \times_{T, \tau} V, (y, v)})
\ar[d]_{\cong} \\
\text{Gr}^i_{\mathfrak m_x}(\mathcal{O}_{X, x})
\ar[r]^-{\cong} \ar@{=}[u] &
\text{Gr}^i_{\mathfrak m_{(x, v)}}(\mathcal{O}_{X \times_S V, (x, v)})
\ar[r]^{f^\sharp_w} &
\text{Gr}^i_{\mathfrak m_w}(\mathcal{O}_{W, w})
}
$$
Ceci implique que $f^\sharp_w$ définit un isomorphisme
$\kappa(x) \to \kappa(w)$ sur les corps résiduels et un isomorphisme
$\mathfrak m_x/\mathfrak m_x^2 \to \mathfrak m_w/\mathfrak m_w^2$
sur les espaces cotangents. Par conséquent, $f^\sharp_w$ définit une surjection
$\mathcal{O}_{X, x}^\wedge \to \mathcal{O}_{W, w}^\wedge$
entre les anneaux locaux complets.

\medskip\noindent
D'après le lemme \ref{more-morphisms-lemma-control-graded}, il existe un isomorphisme de
$\text{Gr}_{\mathfrak m_s}(\mathcal{O}_{(Y \times_{T, \tau} V, (y, v)})$
sur
$\text{Gr}_{\mathfrak m_s}(\mathcal{O}_{Y_\sigma, y_\sigma})$.
Ceci s'obtient en prenant les germes de l'isomorphisme des faisceaux conormaux
au point $y$. Puisque nos anneaux locaux sont noethériens,
la formation du gradué associé relativement à $\mathfrak m_s$
commute à la complétion, car la complétion relativement à un idéal
est un foncteur exact sur les modules finis sur les anneaux noethériens.
Ceci fournit l'isomorphisme vertical de droite dans le diagramme d'anneaux gradués
$$
\xymatrix{
\text{Gr}_{\mathfrak m_s}(\mathcal{O}_{W, w}^\wedge) &
\text{Gr}_{\mathfrak m_s}
(\mathcal{O}_{(Y \times_{T, \tau} V, (y, v)}^\wedge) \ar[l] \\
\text{Gr}_{\mathfrak m_s}(\mathcal{O}_{X, x}^\wedge)
\ar[r]^\varphi \ar[u] &
\text{Gr}_{\mathfrak m_s}(\mathcal{O}_{Y_\sigma, y_\sigma}^\wedge)
\ar[u]_{\cong}
}
$$
Nous n'affirmons pas que le diagramme soit commutatif. D'après le résultat du paragraphe
précédent, la flèche de gauche est surjective. Les trois autres flèches
sont des isomorphismes. Il s'ensuit que la flèche de gauche est un morphisme surjectif
entre des anneaux noethériens isomorphes. C'est donc un isomorphisme
d'après Algèbre, lemme \ref{algebra-lemma-surjective-endo-noetherian-ring-is-iso}
(on peut également le démontrer directement au moyen des fonctions de Hilbert).
En particulier, $\mathcal{O}_{X, x}^\wedge \to \mathcal{O}_{W, w}^\wedge$
doit être injectif aussi bien que surjectif, ce qui achève la démonstration.
\end{proof}















\section{Les morphismes finis localement libres dominent localement les morphismes \'etales}
\label{more-morphisms-section-finite-free-over-etale}

\noindent
Dans cette section, nous expliquons un résultat qui énonce, en substance, que
les recouvrements \'etales d'un schéma $S$ peuvent être raffinés par des recouvrements
de Zariski de schémas finis localement libres sur $S$.

\begin{lemma}
\label{more-morphisms-lemma-dominate-etale-neighbourhood-finite-flat}
Soit $S$ un schéma. Soit $s \in S$.
Soit $f : (U, u) \to (S, s)$ un voisinage \'etale.
Il existe un voisinage ouvert affine $s \in V \subset S$
et un morphisme fini localement libre et surjectif $\pi : T \to V$
tels que, pour tout $t \in \pi^{-1}(s)$, il existe un
voisinage ouvert $t \in W_t \subset T$ et un
diagramme commutatif
$$
\xymatrix{
T \ar[d]^\pi & W_t \ar[l] \ar[rr]_{h_t} \ar[rd] & & U \ar[dl] \\
V \ar[rr] & & S
}
$$
avec $h_t(t) = u$.
\end{lemma}

\begin{proof}
Le problème est local sur $S$; nous pouvons donc remplacer $S$ par un
voisinage ouvert quelconque de $s$.
Nous pouvons aussi remplacer $U$ par un voisinage ouvert de $u$.
Ainsi, d'après Morphismes, lemme \ref{morphisms-lemma-etale-locally-standard-etale},
nous pouvons supposer que
$U \to S$ est un morphisme \'etale standard de schémas affines.
Dans ce cas, le lemme (avec $V = S$) résulte de
Algèbre, lemme \ref{algebra-lemma-standard-etale-finite-flat-Zariski}.
\end{proof}

\begin{lemma}
\label{more-morphisms-lemma-dominate-etale-affine-finite-flat}
Soit $f : U \to S$ un morphisme \'etale surjectif de schémas affines.
Il existe un morphisme fini localement libre et surjectif
$\pi : T \to S$ et un recouvrement ouvert fini
$T = T_1 \cup \ldots \cup T_n$ tels que chacun des morphismes
$T_i \to S$ se factorise par $U \to S$. Diagramme :
$$
\xymatrix{
& \coprod T_i  \ar[rd] \ar[ld] & \\
T \ar[rd]^\pi & & U \ar[ld]_f \\
& S &
}
$$
où la flèche sud-ouest est un recouvrement de Zariski.
\end{lemma}

\begin{proof}
C'est une reformulation de
Algèbre, lemme \ref{algebra-lemma-etale-finite-flat-zariski}.
\end{proof}

\begin{remark}
\label{more-morphisms-remark-topologies}
En termes de topologies, les
lemmes \ref{more-morphisms-lemma-dominate-etale-neighbourhood-finite-flat} et
\ref{more-morphisms-lemma-dominate-etale-affine-finite-flat} signifient ce qui suit.
Soit $S$ un schéma quelconque. Soit $\{f_i : U_i \to S\}$ un recouvrement \'etale
de $S$. Il existe un recouvrement ouvert de Zariski $S = \bigcup V_j$,
pour tout $j$ un morphisme fini localement libre et surjectif
$W_j \to V_j$, et pour tout $j$ un recouvrement ouvert de Zariski
$\{W_{j, k} \to W_j\}$ tel que la famille
$\{W_{j, k} \to S\}$ raffine le recouvrement \'etale donné
$\{f_i : U_i \to S\}$. Que signifie ceci en pratique ?
Par exemple, supposons que nous ayons un problème de descente que nous
sachions résoudre pour les recouvrements de Zariski et pour les recouvrements fppf
de la forme $\{\pi : T \to S\}$, où $\pi$ est fini localement libre
et surjectif. Alors ce problème de descente admet également une réponse
affirmative pour les recouvrements \'etales. Cette astuce a été utilisée par
Gabber dans sa démonstration de $\text{Br}(X) = \text{Br}'(X)$
pour un schéma affine $X$; voir \cite{Hoobler}.
\end{remark}




\section{Localisation \'etale des morphismes quasi-finis}
\label{more-morphisms-section-etale-localization}

\noindent
Nous abordons maintenant une série de lemmes autour du thème
«~les morphismes quasi-finis deviennent finis après localisation \'etale~».
L'idée générale est la suivante. Soient un morphisme
de schémas $f : X \to S$ et un point $s \in S$. Soit
$\varphi : (U, u) \to (S, s)$ un voisinage \'etale de $s$ dans $S$.
Considérons le produit fibré $X_U = U \times_S X$ et le
diagramme fondamental
\begin{equation}
\label{more-morphisms-equation-basic-diagram}
\vcenter{
\xymatrix{
V \ar[r] \ar[dr] & X_U \ar[d] \ar[r] & X \ar[d]^f \\
& U \ar[r]^\varphi & S
}
}
\end{equation}
où $V \subset X_U$ est ouvert.
Existe-t-il un modèle standard pour le morphisme $f_U : X_U \to U$, ou pour
le morphisme $V \to U$ lorsque $V$ est un ouvert convenable ?
Bien entendu, la réponse est non en général. Mais, pour les morphismes quasi-finis,
nous pouvons dire quelque chose.

\begin{lemma}
\label{more-morphisms-lemma-etale-makes-quasi-finite-finite-at-point}
Soit $f : X \to S$ un morphisme de schémas.
Soit $x \in X$. Posons $s = f(x)$.
Supposons que
\begin{enumerate}
\item $f$ soit localement de type fini, et
\item $x \in X_s$ soit isolé\footnote{En présence de (1),
cela signifie que $f$ est
quasi-fini en $x$; voir
Morphismes, lemme \ref{morphisms-lemma-quasi-finite-at-point-characterize}.}.
\end{enumerate}
Alors il existe
\begin{enumerate}
\item[(a)] un voisinage \'etale élémentaire $(U, u) \to (S, s)$,
\item[(b)] un sous-schéma ouvert $V \subset X_U$
(voir \ref{more-morphisms-equation-basic-diagram})
\end{enumerate}
tels que
\begin{enumerate}
\item[(\romannumeral1)] $V \to U$ soit un morphisme fini,
\item[(\romannumeral2)] il existe un unique point $v$ de $V$
d'image $u$ dans $U$, et
\item[(\romannumeral3)] le point $v$ s'envoie sur $x$
par le morphisme $X_U \to X$, qui induit $\kappa(x) = \kappa(v)$.
\end{enumerate}
De plus, pour tout voisinage \'etale élémentaire $(U', u') \to (U, u)$,
si l'on pose $V' = U' \times_U V \subset X_{U'}$, le triplet $(U', u', V')$
satisfait encore les propriétés
(\romannumeral1), (\romannumeral2) et (\romannumeral3).
\end{lemma}

\begin{proof}
Soient $Y \subset X$, $W \subset S$ des ouverts affines tels que
$f(Y) \subset W$ et $x \in Y$. Remarquons que $x$ est
aussi un point isolé de la fibre du morphisme $f|_Y : Y \to W$.
Si nous pouvons démontrer le résultat pour $f|_Y : Y \to W$, alors le
résultat en découle pour $f$. Nous nous ramenons donc au cas où
$f$ est un morphisme de schémas affines. Ce cas est traité dans
Algèbre, lemme \ref{algebra-lemma-etale-makes-quasi-finite-finite-one-prime}.
\end{proof}

\noindent
Dans le lemme précédent et dans le suivant, nous ne supposons pas que le morphisme
$f$ soit séparé. Cela signifie que les ouverts $V$, $V_i$ qui y sont construits
ne sont pas nécessairement fermés dans $X_U$. De plus, si nous choisissons
le voisinage $U$ affine, alors chaque $V_i$ est affine, mais
les intersections $V_i \cap V_j$ ne sont pas nécessairement affines (dans le cas non
séparé).

\begin{lemma}
\label{more-morphisms-lemma-etale-makes-quasi-finite-finite-multiple-points}
Soit $f : X \to S$ un morphisme de schémas.
Soient $x_1, \ldots, x_n \in X$ des points de même image $s$ dans $S$.
Supposons que
\begin{enumerate}
\item $f$ soit localement de type fini, et
\item $x_i \in X_s$ soit isolé pour $i = 1, \ldots, n$.
\end{enumerate}
Alors il existe
\begin{enumerate}
\item[(a)] un voisinage \'etale élémentaire $(U, u) \to (S, s)$,
\item[(b)] pour chaque $i$, un sous-schéma ouvert $V_i \subset X_U$,
\end{enumerate}
tels que, pour chaque $i$, les propriétés suivantes soient vérifiées :
\begin{enumerate}
\item[(\romannumeral1)] $V_i \to U$ est un morphisme fini,
\item[(\romannumeral2)] il existe un unique point $v_i$ de $V_i$
d'image $u$ dans $U$, et
\item[(\romannumeral3)] le point $v_i$ s'envoie sur $x_i$ dans $X$ et
$\kappa(x_i) = \kappa(v_i)$.
\end{enumerate}
\end{lemma}

\begin{proof}
Nous raisonnons par récurrence sur $n$.
Supposons en effet que $(U, u) \to (S, s)$ et $V_i \subset X_U$,
$i = 1, \ldots, n - 1$, conviennent pour $x_1, \ldots, x_{n - 1}$. Puisque
$\kappa(s) = \kappa(u)$, la fibre $(X_U)_u = X_s$. Il existe donc
un unique point $x'_n \in X_u \subset X_U$ correspondant à
$x_n \in X_s$. De plus, $x'_n$ est isolé dans $X_u$. Ainsi, d'après le
lemme \ref{more-morphisms-lemma-etale-makes-quasi-finite-finite-at-point}, il
existe un voisinage \'etale élémentaire $(U', u') \to (U, u)$
et un ouvert $V_n \subset X_{U'}$ qui convient pour $x'_n$, et donc
pour $x_n$.
D'après la dernière assertion du
lemme \ref{more-morphisms-lemma-etale-makes-quasi-finite-finite-at-point},
les sous-schémas ouverts $V'_i = U'\times_U V_i$, pour $i = 1, \ldots, n - 1$, conviennent encore
relativement à $x_1, \ldots, x_{n - 1}$. D'où le résultat.
\end{proof}

\noindent
Si nous autorisons une extension de corps non triviale $\kappa(u)/\kappa(s)$, c'est-à-dire
des voisinages \'etales généraux, nous pouvons scinder les points comme suit.

\begin{lemma}
\label{more-morphisms-lemma-etale-makes-quasi-finite-finite-multiple-points-var}
Soit $f : X \to S$ un morphisme de schémas.
Soient $x_1, \ldots, x_n \in X$ des points de même image $s$ dans $S$.
Supposons que
\begin{enumerate}
\item $f$ soit localement de type fini, et
\item $x_i \in X_s$ soit isolé pour $i = 1, \ldots, n$.
\end{enumerate}
Alors il existe
\begin{enumerate}
\item[(a)] un voisinage \'etale $(U, u) \to (S, s)$,
\item[(b)] pour chaque $i$, un entier $m_i$ et des
sous-schémas ouverts $V_{i, j} \subset X_U$, $j = 1, \ldots, m_i$,
\end{enumerate}
tels que l'on ait
\begin{enumerate}
\item[(\romannumeral1)] chaque $V_{i, j} \to U$ soit un morphisme fini,
\item[(\romannumeral2)] il existe un unique point $v_{i, j}$ de $V_{i, j}$
d'image $u$ dans $U$, tel que $\kappa(u) \subset \kappa(v_{i, j})$ soit
une extension finie purement inséparable,
\item[(\romannumeral4)] si $v_{i, j} = v_{i', j'}$, alors $i = i'$ et
$j = j'$, et
\item[(\romannumeral3)] les points $v_{i, j}$ s'envoient sur $x_i$ dans $X$ et
aucun autre point de $(X_U)_u$ ne s'envoie sur $x_i$.
\end{enumerate}
\end{lemma}

\begin{proof}
Cette démonstration est une variante de la démonstration de
Algèbre, lemme \ref{algebra-lemma-etale-makes-quasi-finite-finite-variant},
formulée dans le langage des schémas.
D'après Morphismes, lemme \ref{morphisms-lemma-quasi-finite-at-point-characterize},
le morphisme $f$ est quasi-fini en chacun des points $x_i$.
Ainsi, l'extension $\kappa(s) \subset \kappa(x_i)$ est finie pour tout $i$
(Morphismes, lemme \ref{morphisms-lemma-residue-field-quasi-finite}).
Pour chaque $i$, soit $\kappa(s) \subset L_i \subset \kappa(x_i)$
le sous-corps tel que $L_i/\kappa(s)$ soit séparable et que
$\kappa(x_i)/L_i$ soit purement inséparable. Choisissons une extension finie
galoisienne $L/\kappa(s)$ telle qu'il existe des
$\kappa(s)$-plongements $L_i \to L$ pour $i = 1, \ldots, n$.
Choisissons un voisinage \'etale $(U, u) \to (S, s)$ tel que
$L \cong \kappa(u)$ comme extensions de $\kappa(s)$
(lemme \ref{more-morphisms-lemma-realize-prescribed-residue-field-extension-etale}).

\medskip\noindent
Soient $y_{i, j}$, $j = 1, \ldots, m_i$, les points de $X_U$
au-dessus de $x_i \in X$ et de $u \in U$. D'après
Schémas, lemme \ref{schemes-lemma-points-fibre-product},
ces points $y_{i, j}$ correspondent exactement aux idéaux premiers des anneaux
$\kappa(u) \otimes_{\kappa(s)} \kappa(x_i)$. Ceci explique également
pourquoi ils sont en nombre fini; en fait,
$m_i = [L_i : \kappa(s)]$, mais nous n'en avons pas besoin.
Par notre choix de
$L$ (et par la théorie élémentaire des corps),
nous voyons que l'extension $\kappa(u) \subset \kappa(y_{i, j})$ est
finie purement inséparable pour chaque paire $i, j$.
De plus, d'après Morphismes, lemme \ref{morphisms-lemma-base-change-quasi-finite},
par exemple, le morphisme
$X_U \to U$ est quasi-fini aux points $y_{i, j}$ pour
tous $i, j$.

\medskip\noindent
Appliquons le lemme \ref{more-morphisms-lemma-etale-makes-quasi-finite-finite-multiple-points}
au morphisme $X_U \to U$, au point $u \in U$
et aux points $y_{i, j} \in (X_U)_u$. Nous obtenons un voisinage \'etale
$(U', u') \to (U, u)$ avec $\kappa(u) = \kappa(u')$ et des
ouverts $V_{i, j} \subset X_{U'}$ possédant les propriétés
(\romannumeral1), (\romannumeral2) et (\romannumeral3)
de ce lemme. Nous affirmons que le voisinage \'etale
$(U', u') \to (S, s)$ et les ouverts $V_{i, j} \subset X_{U'}$
constituent une solution au problème posé par le lemme.
Nous omettons les vérifications.
\end{proof}

\begin{lemma}
\label{more-morphisms-lemma-etale-splits-off-quasi-finite-part-technical}
Soit $f : X \to S$ un morphisme de schémas.
Soit $s \in S$. Soient $x_1, \ldots, x_n \in X_s$. Supposons que
\begin{enumerate}
\item $f$ soit localement de type fini,
\item $f$ soit séparé, et
\item $x_1, \ldots, x_n$ soient des points de $X_s$, isolés et deux à deux distincts.
\end{enumerate}
Il existe alors un voisinage \'etale élémentaire $(U, u) \to (S, s)$
et une décomposition
$$
U \times_S X = W \amalg V_1 \amalg \ldots \amalg V_n
$$
en sous-schémas ouverts et fermés telle que les morphismes
$V_i \to U$ soient finis, que les fibres de $V_i \to U$ au-dessus de $u$ soient
les singletons $\{v_i\}$, que chaque $v_i$ s'envoie sur $x_i$ avec
$\kappa(x_i) = \kappa(v_i)$, et que la fibre de $W \to U$
au-dessus de $u$ ne contienne aucun point s'envoyant sur l'un des $x_i$.
\end{lemma}

\begin{proof}
Choisissons $(U, u) \to (S, s)$ et $V_i \subset X_U$ comme dans le
lemme \ref{more-morphisms-lemma-etale-makes-quasi-finite-finite-multiple-points}.
Puisque $X_U \to U$ est séparé
(Schémas, lemme \ref{schemes-lemma-separated-permanence})
et que $V_i \to U$ est fini, donc propre
(Morphismes, lemme \ref{morphisms-lemma-finite-proper}),
on voit que $V_i \subset X_U$ est fermé d'après
Morphismes, lemme \ref{morphisms-lemma-image-proper-scheme-closed}.
Ainsi, $V_i \cap V_j$ est un sous-ensemble fermé de $V_i$ qui
ne contient pas $v_i$. Par conséquent, l'image de $V_i \cap V_j$
dans $U$ est un fermé (car $V_i \to U$ est propre) qui ne
contient pas $u$. Après avoir rétréci $U$, on peut donc supposer
que $V_i \cap V_j = \emptyset$ pour tous $i, j$. Cela donne la
décomposition de l'énoncé.
\end{proof}

\noindent
Voici la variante où l'on se ramène à des extensions de corps
purement inséparables.

\begin{lemma}
\label{more-morphisms-lemma-etale-splits-off-quasi-finite-part-technical-variant}
Soit $f : X \to S$ un morphisme de schémas.
Soient $s \in S$ et $x_1, \ldots, x_n \in X_s$. Supposons que
\begin{enumerate}
\item $f$ soit localement de type fini,
\item $f$ soit séparé, et
\item $x_1, \ldots, x_n$ soient des points isolés deux à deux distincts de $X_s$.
\end{enumerate}
Il existe alors un voisinage \'etale $(U, u) \to (S, s)$
et une décomposition
$$
U \times_S X =
W \amalg
\ \coprod\nolimits_{i = 1, \ldots, n}
\ \coprod\nolimits_{j = 1, \ldots, m_i}
V_{i, j}
$$
en sous-schémas ouverts et fermés telle que les morphismes
$V_{i, j} \to U$ soient finis, que les fibres de $V_{i, j} \to U$ au-dessus de $u$ soient
les singletons $\{v_{i, j}\}$, que chaque $v_{i, j}$ s'envoie sur $x_i$,
que $\kappa(u) \subset \kappa(v_{i, j})$ soit purement inséparable,
et que la fibre de $W \to U$ au-dessus de $u$ ne contienne aucun point s'envoyant
sur l'un des $x_i$.
\end{lemma}

\begin{proof}
On procède exactement comme dans la démonstration du
lemme \ref{more-morphisms-lemma-etale-splits-off-quasi-finite-part-technical}, si ce n'est qu'on
utilise le lemme \ref{more-morphisms-lemma-etale-makes-quasi-finite-finite-multiple-points-var}
au lieu du lemme \ref{more-morphisms-lemma-etale-makes-quasi-finite-finite-multiple-points}.
\end{proof}

\noindent
La version suivante est peut-être un peu plus facile à lire.

\begin{lemma}
\label{more-morphisms-lemma-etale-splits-off-quasi-finite-part}
Soit $f : X \to S$ un morphisme de schémas.
Soit $s \in S$. Supposons que
\begin{enumerate}
\item $f$ soit localement de type fini,
\item $f$ soit séparé, et
\item $X_s$ ait au plus un nombre fini de points isolés.
\end{enumerate}
Il existe alors un voisinage \'etale élémentaire $(U, u) \to (S, s)$
et une décomposition
$$
U \times_S X = W \amalg V
$$
en sous-schémas ouverts et fermés telle que le morphisme
$V \to U$ soit fini et que la fibre $W_u$ du
morphisme $W \to U$ ne contienne aucun point isolé.
En particulier, si $f^{-1}(s)$ est un ensemble fini, alors $W_u = \emptyset$.
\end{lemma}

\begin{proof}
Cela résulte immédiatement du
lemme \ref{more-morphisms-lemma-etale-splits-off-quasi-finite-part-technical},
en choisissant pour $x_1, \ldots, x_n$ l'ensemble de tous les
points isolés de $X_s$ et en posant $V = \bigcup V_i$.
\end{proof}






\section{Localisation \'etale des morphismes entiers}
\label{more-morphisms-section-etale-localization-integral}

\noindent
Voici quelques variantes des résultats de la section \ref{more-morphisms-section-etale-localization}
dans le cas des morphismes entiers.

\begin{lemma}
\label{more-morphisms-lemma-etale-makes-integral-split}
Soit $R \to S$ un homomorphisme d'anneaux entier. Soit $\mathfrak p \subset R$ un idéal
premier. Supposons
\begin{enumerate}
\item qu'il n'y ait qu'un nombre fini d'idéaux premiers $\mathfrak q_1, \ldots, \mathfrak q_n$
au-dessus de $\mathfrak p$, et
\item que, pour tout $i$, la sous-extension séparable maximale
$\kappa(\mathfrak q)/\kappa(\mathfrak q_i)_{sep}/\kappa(\mathfrak p)$
(Corps, lemme \ref{fields-lemma-separable-first})
soit finie sur $\kappa(\mathfrak p)$.
\end{enumerate}
Il existe alors un homomorphisme d'anneaux \'etale $R \to R'$ et un idéal premier
$\mathfrak p'$ au-dessus de $\mathfrak p$ tels que
$$
S \otimes_R R' = A_1 \times \ldots \times A_m
$$
où $R' \to A_j$ est entier et possède un unique idéal premier $\mathfrak r_j$
au-dessus de $\mathfrak p'$ tel que $\kappa(\mathfrak r_j)/\kappa(\mathfrak p')$
soit purement inséparable.
\end{lemma}

\begin{proof}[Première démonstration]
Cette démonstration utilise
Algèbre, lemme \ref{algebra-lemma-etale-makes-quasi-finite-finite-variant}.
Choisissons en effet un générateur $\theta_i \in \kappa(\mathfrak q_i)_{sep}$
de ce corps sur $\kappa(\mathfrak p)$
(Corps, lemme \ref{fields-lemma-primitive-element}).
Le spectre de l'anneau fibre $S \otimes_R \kappa(\mathfrak p)$
est fini et discret, ses points correspondant à
$\mathfrak q_1, \ldots, \mathfrak q_n$.
D'après le théorème chinois des restes
(Algèbre, lemme \ref{algebra-lemma-chinese-remainder}),
on voit que $S \otimes_R \kappa(\mathfrak p) \to \prod \kappa(\mathfrak q_i)$
est surjectif. Par conséquent, après avoir remplacé $R$ par $R_g$ pour un certain
$g \in R$, $g \not \in \mathfrak p$, on peut supposer que
$(0, \ldots, 0, \theta_i, 0, \ldots, 0) \in \prod \kappa(\mathfrak q_i)$
est l'image d'un élément $x_i \in S$.
Soit $S' \subset S$ la $R$-sous-algèbre engendrée par les éléments $x_i$.
Puisque $\Spec(S) \to \Spec(S')$ est surjectif
(Algèbre, lemme \ref{algebra-lemma-integral-overring-surjective}),
on conclut que
$\mathfrak q_i' = S' \cap \mathfrak q_i$ sont les idéaux premiers de
$S'$ au-dessus de $\mathfrak p$. Le choix des $x_i$ montre que
ces idéaux premiers sont distincts et que
$\kappa(\mathfrak q'_i)_{sep} = \kappa(\mathfrak q_i)_{sep}$.
En particulier, les extensions de corps
$\kappa(\mathfrak q_i)/\kappa(\mathfrak q'_i)$ sont purement
inséparables.
Puisque $R \to S'$ est fini, on peut appliquer
Algèbre, lemme \ref{algebra-lemma-etale-makes-quasi-finite-finite-variant},
et l'on obtient $R \to R'$ et $\mathfrak p'$, ainsi qu'une décomposition
$$
S' \otimes_R R' = A'_1 \times \ldots \times A'_m \times B'
$$
où $R' \to A'_j$ est entier et possède un unique idéal premier
$\mathfrak r'_j$ au-dessus de $\mathfrak p'$ tel que
$\kappa(\mathfrak r'_j)/\kappa(\mathfrak p')$
soit purement inséparable, et où $B'$ ne possède aucun idéal premier
au-dessus de $\mathfrak p'$. Comme $R' \to B'$ est fini
(car $R \to S'$ est fini), on peut,
après avoir localisé $R'$ en un certain $g' \in R'$, $g' \not \in \mathfrak p'$,
supposer que $B' = 0$. Le morphisme $S' \otimes_R R' \to S \otimes_R R'$
nous fournit la décomposition correspondante pour $S$.
\end{proof}

\begin{proof}[Seconde démonstration]
Cette démonstration utilise l'hensélisation stricte. Supposons d'abord que $R$ soit strictement
hensélien, d'idéal maximal $\mathfrak p$. Alors $S/\mathfrak p S$
possède un nombre fini d'idéaux premiers correspondant à
$\mathfrak q_1, \ldots, \mathfrak q_n$, tous maximaux,
et dont les corps résiduels sont purement inséparables sur $\kappa(\mathfrak p)$.
Ainsi, $S/\mathfrak p S$ est égal à $\prod (S/\mathfrak p S)_{\mathfrak p_i}$.
D'après Plus sur l'algèbre, lemme \ref{more-algebra-lemma-characterize-henselian-pair},
on peut relever cette décomposition en produit en une décomposition en produit de $S$
comme dans l'énoncé.

\medskip\noindent
Dans le cas général, soit $R^{sh}$ l'hensélisé strict de
$R_\mathfrak p$. On peut alors appliquer le résultat du premier paragraphe
à $R^{sh} \to S \otimes_R R^{sh}$. Considérons les $m$ idempotents deux à deux orthogonaux
de $S \otimes_R R^{sh}$ correspondant à la décomposition en produit.
Puisque $R^{sh}$ est une colimite filtrante de morphismes d'anneaux \'etales
$(R, \mathfrak p) \to (R', \mathfrak p')$ d'après
Algèbre, lemme \ref{algebra-lemma-strict-henselization-different},
on voit que ces idempotents sont définis sur un certain $R'$, comme voulu.
\end{proof}







\section{Théorème principal de Zariski}
\label{more-morphisms-section-application-etale-neighbourhoods}

\noindent
Dans cette section, on démontre le théorème principal de Zariski sous la forme donnée par Grothendieck.
Souvent, lorsque l'on dit « théorème principal de Zariski » dans ce contexte, on entend l'un des
lemmes \ref{more-morphisms-lemma-finite-type-separated},
\ref{more-morphisms-lemma-quasi-finite-separated-quasi-affine} ou
\ref{more-morphisms-lemma-quasi-finite-separated-pass-through-finite}.
Dans la plupart des textes, c'est au dernier de ces énoncés que l'on donne le nom de
théorème principal de Zariski.

\medskip\noindent
On a déjà démontré la version algébrique dans
Algèbre, théorème \ref{algebra-theorem-main-theorem},
et on a déjà reformulé cette version algébrique
dans le langage des schémas ; voir
Morphismes, théorème \ref{morphisms-theorem-main-theorem}.
La version de cette section est plus subtile ; pour obtenir le résultat
complet, on utilise les techniques de localisation \'etale
de la section \ref{more-morphisms-section-etale-localization} afin de se ramener au
cas algébrique.

\begin{lemma}
\label{more-morphisms-lemma-finite-type-separated}
Soit $f : X \to S$ un morphisme de schémas.
Supposons que $f$ soit de type fini et séparé.
Soit $S'$ la normalisation de $S$ dans $X$ ; voir
Morphismes, définition \ref{morphisms-definition-normalization-X-in-Y}.
Diagramme :
$$
\xymatrix{
X \ar[rd]_f \ar[rr]_{f'} & & S' \ar[ld]^\nu \\
& S &
}
$$
Il existe alors un sous-schéma ouvert $U' \subset S'$ tel que
\begin{enumerate}
\item $(f')^{-1}(U') \to U'$ soit un isomorphisme, et
\item $(f')^{-1}(U') \subset X$ soit l'ensemble des points où
$f$ est quasi-fini.
\end{enumerate}
\end{lemma}

\begin{proof}
D'après Morphismes, lemme \ref{morphisms-lemma-quasi-finite-points-open},
le sous-ensemble $U \subset X$ des points où $f$ est quasi-fini est ouvert.
Le lemme équivaut aux trois assertions suivantes :
\begin{enumerate}
\item[(a)] $U' = f'(U) \subset S'$ est ouvert,
\item[(b)] $U = (f')^{-1}(U')$, et
\item[(c)] $U \to U'$ est un isomorphisme.
\end{enumerate}
Soit $x \in U$ arbitraire. On affirme qu'il existe un voisinage ouvert
$f'(x) \in V \subset S'$ tel que $(f')^{-1}V \to V$ soit un
isomorphisme. Montrons d'abord que cette assertion implique le lemme.
En effet, $(f')^{-1}V \cong V$ est alors à la fois localement de type
fini sur $S$ (comme sous-schéma ouvert de $X$), et, pour $v \in V$, l'extension de corps
résiduels $\kappa(v)/\kappa(\nu(v))$ est algébrique (car
$V \subset S'$ et $S'$ est entier sur $S$). Les fibres
de $V \to S$ sont donc discrètes (Morphismes, lemme
\ref{morphisms-lemma-algebraic-residue-field-extension-closed-point-fibre})
et $(f')^{-1}V \to S$ est localement quasi-fini
(Morphismes, lemme \ref{morphisms-lemma-locally-quasi-finite-fibres}).
Il s'ensuit que $(f')^{-1}V \subset U$ et $V \subset U'$. Puisque $x$ était
arbitraire, les assertions (a), (b) et (c) sont vraies.

\medskip\noindent
Posons $s = f(x)$. Soit $(T, t) \to (S, s)$ un voisinage \'etale
élémentaire. On note par un indice ${}_T$ le changement de base à $T$.
Soit $y = (x, t) \in X_T$ l'unique point de
la fibre $X_t$ situé au-dessus de $x$. Remarquons que $U_T \subset X_T$
est l'ensemble des points où $f_T$ est quasi-fini ; voir
Morphismes, lemme \ref{morphisms-lemma-base-change-quasi-finite}.
Remarquons que
$$
X_T \xrightarrow{f'_T} S'_T \xrightarrow{\nu_T} T
$$
est la normalisation de $T$ dans $X_T$ ; voir
lemme \ref{more-morphisms-lemma-normalization-smooth-localization}.
Supposons que l'assertion soit vraie pour $y \in U_T \subset X_T \to S'_T \to T$, c'est-à-dire
que l'on puisse trouver un voisinage ouvert
$f'_T(y) \in V' \subset S'_T$ tel que $(f'_T)^{-1}V' \to V'$ soit un
isomorphisme. Le morphisme $S'_T \to S'$ est \'etale ; l'image
$V \subset S'$ de $V'$ est donc ouverte. On a $f'(x) \in V$ puisque $f'_T(y) \in V'$.
Observons que
$$
\xymatrix{
(f'_T)^{-1}V' \ar[r] \ar[d] & (f')^{-1}(V) \ar[d] \\
V' \ar[r] & V
}
$$
est un carré cartésien (car $S'_T \times_{S'} X = X_T$).
Puisque la flèche verticale de gauche est un isomorphisme
et que $\{V' \to V\}$ est un recouvrement \'etale, on conclut que la flèche verticale de droite
est un isomorphisme d'après
Descente, lemme \ref{descent-lemma-descending-property-isomorphism}.
Autrement dit, l'assertion est vraie pour $x \in U \subset X \to S' \to S$.

\medskip\noindent
D'après le paragraphe précédent, on peut remplacer $S$ par un
voisinage \'etale élémentaire de $s = f(x)$ pour démontrer l'assertion.
On peut donc supposer qu'il existe une décomposition
$$
X = V \amalg W
$$
en sous-schémas ouverts et fermés, où $V \to S$ est fini et $x \in V$ ;
voir le lemme \ref{more-morphisms-lemma-etale-splits-off-quasi-finite-part-technical}.
Puisque $X$ est la réunion disjointe de $V$ et $W$ au-dessus de $S$ et que
$V \to S$ est fini, on voit que la
normalisation de $S$ dans $X$ est le morphisme
$$
X = V \amalg W \longrightarrow V \amalg W' \longrightarrow S
$$
où $W'$ est la normalisation de $S$ dans $W$ ; voir
Morphismes, lemmes \ref{morphisms-lemma-normalization-in-disjoint-union},
\ref{morphisms-lemma-finite-integral} et
\ref{morphisms-lemma-normalization-in-integral}.
L'assertion s'ensuit, ce qui achève la démonstration.
\end{proof}

\begin{lemma}
\label{more-morphisms-lemma-quasi-finite-separated-quasi-affine}
\begin{slogan}
Les morphismes quasi-finis et séparés sont quasi-affines
\end{slogan}
Soit $f : X \to S$ un morphisme de schémas.
Supposons que $f$ soit quasi-fini et séparé.
Soit $S'$ la normalisation de $S$ dans $X$ ; voir
Morphismes, définition \ref{morphisms-definition-normalization-X-in-Y}.
Diagramme :
$$
\xymatrix{
X \ar[rd]_f \ar[rr]_{f'} & & S' \ar[ld]^\nu \\
& S &
}
$$
Alors $f'$ est une immersion ouverte quasi-compacte et $\nu$ est entier.
En particulier, $f$ est quasi-affine.
\end{lemma}

\begin{proof}
Cela résulte du lemme \ref{more-morphisms-lemma-finite-type-separated}. En effet, d'après
ce lemme, il existe un sous-schéma ouvert $U' \subset S'$ tel que
$(f')^{-1}(U') = X$ et que $X \to U'$ soit un isomorphisme. Autrement
dit, $f'$ est une immersion ouverte. Remarquons que $f'$ est quasi-compact puisque
$f$ est quasi-compact et que $\nu : S' \to S$ est séparé
(Schémas, lemme \ref{schemes-lemma-quasi-compact-permanence}).
Il s'ensuit que $f$ est quasi-affine d'après
Morphismes, lemme \ref{morphisms-lemma-characterize-quasi-affine}.
\end{proof}

\begin{lemma}[Théorème principal de Zariski]
\label{more-morphisms-lemma-quasi-finite-separated-pass-through-finite}
\begin{reference}
\cite[IV Corollaire 18.12.13]{EGA}
\end{reference}
Soit $f : X \to S$ un morphisme de schémas.
Supposons que $f$ soit quasi-fini et séparé, et que
$S$ soit quasi-compact et quasi-séparé. Il existe alors
une factorisation
$$
\xymatrix{
X \ar[rd]_f \ar[rr]_j & & T \ar[ld]^\pi \\
& S &
}
$$
où $j$ est une immersion ouverte quasi-compacte et $\pi$ est fini.
\end{lemma}

\begin{proof}
Soit $X \to S' \to S$ la factorisation fournie par
le lemme \ref{more-morphisms-lemma-quasi-finite-separated-quasi-affine}.
D'après
Propriétés, lemme
\ref{properties-lemma-integral-algebra-directed-colimit-finite},
on peut écrire
$\nu_*\mathcal{O}_{S'} = \colim_{i \in I} \mathcal{A}_i$ comme une
colimite filtrante de $\mathcal{O}_S$-algèbres quasi-cohérentes finies
$\mathcal{A}_i \subset \nu_*\mathcal{O}_{S'}$. Alors
$\pi_i : T_i = \underline{\Spec}_S(\mathcal{A}_i) \to S$
est un morphisme fini pour tout $i$.
Remarquons que les morphismes de transition $T_{i'} \to T_i$ sont affines
et que $S' = \lim T_i$.

\medskip\noindent
D'après Limites, lemme \ref{limits-lemma-descend-opens},
il existe un $i$ et un ouvert quasi-compact
$U_i \subset T_i$ dont l'image réciproque dans $S'$ est égale à
$f'(X)$. Pour $i' \geq i$, soit $U_{i'}$ l'image réciproque
de $U_i$ dans $T_{i'}$. Alors $X \cong f'(X) = \lim_{i' \geq i} U_{i'}$ ; voir
Limites, lemme \ref{limits-lemma-directed-inverse-system-has-limit}.
D'après Limites, lemme \ref{limits-lemma-finite-type-eventually-closed},
on voit que $X \to U_{i'}$ est une immersion fermée pour un certain $i' \geq i$.
(En fait, $X \cong U_{i'}$ pour tout indice suffisamment
grand $i'$, mais on n'en a pas besoin.) Ainsi, $X \to T_{i'}$ est une immersion. D'après
Morphismes, lemme \ref{morphisms-lemma-factor-quasi-compact-immersion},
on peut la factoriser en $X \to T \to T_{i'}$, où la première flèche
est une immersion ouverte et la seconde une immersion fermée. Cela achève la démonstration.
\end{proof}

\begin{lemma}
\label{more-morphisms-lemma-quasi-finite-separated-pass-through-finite-addendum}
Avec les notations et sous les hypothèses du
lemme \ref{more-morphisms-lemma-quasi-finite-separated-pass-through-finite},
supposons de plus que $f$ soit localement de présentation finie. On peut alors
choisir la factorisation de sorte que $T$ soit fini et de
présentation finie sur $S$.
\end{lemma}

\begin{proof}
D'après Limites, lemme
\ref{limits-lemma-finite-in-finite-and-finite-presentation}, on peut écrire
$T = \lim T_i$, où tous les $T_i$ sont finis et de présentation finie
sur $S$ et où les morphismes de transition $T_{i'} \to T_i$ sont des immersions
fermées. D'après
Limites, lemme \ref{limits-lemma-descend-opens},
il existe un $i$ et un sous-schéma ouvert $U_i \subset T_i$ dont l'image
réciproque dans $T$ est $X$. D'après
Limites, lemme
\ref{limits-lemma-finite-type-eventually-closed},
on voit que $X \cong U_i$ pour $i$ suffisamment grand.
Remplacer $T$ par $T_i$ achève la démonstration.
\end{proof}








\section{Applications du théorème principal de Zariski, I}
\label{more-morphisms-section-applications-zmt}

\noindent
Une première application est la caractérisation des morphismes finis
comme les morphismes propres à fibres finies.

\begin{lemma}
\label{more-morphisms-lemma-characterize-finite}
Soit $f : X \to S$ un morphisme de schémas.
Les assertions suivantes sont équivalentes :
\begin{enumerate}
\item $f$ est fini,
\item $f$ est propre à fibres finies,
\item $f$ est propre et localement quasi-fini,
\item $f$ est universellement fermé, séparé, localement de type fini
et à fibres finies.
\end{enumerate}
\end{lemma}

\begin{proof}
(1) implique (2) d'après
Morphismes, lemmes \ref{morphisms-lemma-finite-proper},
\ref{morphisms-lemma-quasi-finite}
et \ref{morphisms-lemma-finite-quasi-finite}.
(2) implique (3) d'après Morphismes, lemme \ref{morphisms-lemma-finite-fibre}.
(3) implique (4) d'après la définition des morphismes propres et
Morphismes, lemmes \ref{morphisms-lemma-quasi-finite-locally-quasi-compact} et
\ref{morphisms-lemma-quasi-finite}.

\medskip\noindent
Supposons (4). Choisissons $s \in S$. D'après
Morphismes, lemme \ref{morphisms-lemma-finite-fibre}, on
voit que les points, en nombre fini, de $X_s$ sont tous isolés dans $X_s$.
Choisissons un voisinage \'etale élémentaire $(U, u) \to (S, s)$
et une décomposition $X_U = V \amalg W$ comme dans le
lemme \ref{more-morphisms-lemma-etale-splits-off-quasi-finite-part}.
Remarquons que $W_u = \emptyset$ car tous les points de $X_s$ sont isolés.
Puisque $f$ est universellement fermé, on voit que
l'image de $W$ dans $U$ est un fermé qui ne contient pas $u$.
Après avoir rétréci $U$, on peut supposer que $W = \emptyset$.
Autrement dit, on voit que $X_U = V$ est fini sur $U$.
Comme $s \in S$ était arbitraire,
il existe une famille $\{U_i \to S\}$
de morphismes \'etales dont les images recouvrent $S$ et telle que
les changements de base $X_{U_i} \to U_i$ soient finis.
Remarquons que $\{U_i \to S\}$ est un recouvrement \'etale ;
voir Topologies, définition \ref{topologies-definition-etale-covering}.
C'est donc un recouvrement fpqc ; voir
Topologies,
lemme \ref{topologies-lemma-zariski-etale-smooth-syntomic-fppf-fpqc}.
On conclut alors que $f$ est fini d'après
Descente, lemme \ref{descent-lemma-descending-property-finite}.
\end{proof}

\noindent
On en déduit les résultats utiles suivants.

\begin{lemma}
\label{more-morphisms-lemma-proper-finite-fibre-finite-in-neighbourhood}
Soit $f : X \to S$ un morphisme de schémas.
Soit $s \in S$.
Supposons que $f$ soit propre et que $f^{-1}(\{s\})$ soit un ensemble fini.
Il existe alors un voisinage ouvert $V \subset S$ de $s$
tel que $f|_{f^{-1}(V)} : f^{-1}(V) \to V$ soit fini.
\end{lemma}

\begin{proof}
Le morphisme $f$ est quasi-fini en tous les points de $f^{-1}(\{s\})$
d'après Morphismes, lemme \ref{morphisms-lemma-finite-fibre}.
D'après Morphismes, lemme \ref{morphisms-lemma-quasi-finite-points-open}, l'ensemble
des points où $f$ est quasi-fini est un ouvert $U \subset X$.
Posons $Z = X \setminus U$. Alors $s \not \in f(Z)$. Puisque $f$ est propre,
l'ensemble $f(Z) \subset S$ est fermé. Choisissons un voisinage ouvert quelconque
$V \subset S$ de $s$ tel que $f(Z) \cap V = \emptyset$. Alors
$f^{-1}(V) \to V$ est localement quasi-fini et propre.
Il est donc quasi-fini
(Morphismes, lemme \ref{morphisms-lemma-quasi-finite-locally-quasi-compact}),
donc à fibres finies
(Morphismes, lemme \ref{morphisms-lemma-quasi-finite}), et il
est ainsi fini d'après le lemme \ref{more-morphisms-lemma-characterize-finite}.
\end{proof}

\begin{lemma}
\label{more-morphisms-lemma-flat-proper-family-cannot-collapse-fibre}
Considérons un diagramme commutatif de schémas
$$
\xymatrix{
X \ar[rr]_h \ar[rd]_f & & Y \ar[ld]^g \\
& S
}
$$
Soit $s \in S$. Supposons
\begin{enumerate}
\item que $X \to S$ soit un morphisme propre,
\item que $Y \to S$ soit séparé et localement de type fini, et
\item que l'image de $X_s \to Y_s$ soit finie.
\end{enumerate}
Il existe alors un ouvert
$U \subset S$ contenant $s$ tel que $X_U \to Y_U$
se factorise par un sous-schéma fermé $Z \subset Y_U$ fini sur $U$.
\end{lemma}

\begin{proof}
Soit $Z \subset Y$ l'image schématique de $h$ ; voir
Morphismes, section \ref{morphisms-section-scheme-theoretic-image}.
D'après Morphismes, lemme \ref{morphisms-lemma-scheme-theoretic-image-is-proper},
le morphisme $X \to Z$ est surjectif et $Z \to S$ est propre.
Ainsi, $X_s \to Z_s$ est surjectif. On voit alors que
(3) implique que $Z_s$ est fini.
Par conséquent, $Z \to S$ est fini dans un voisinage ouvert de $s$ d'après le
lemme \ref{more-morphisms-lemma-proper-finite-fibre-finite-in-neighbourhood}.
\end{proof}





\section{Applications du théorème principal de Zariski, II}
\label{more-morphisms-section-applications-zmt-II}

\noindent
Dans cette section, on donne quelques conséquences supplémentaires du théorème principal de Zariski
concernant la structure des morphismes quasi-finis.

\begin{lemma}
\label{more-morphisms-lemma-separated-locally-quasi-finite-over-affine}
Soit $f : X \to Y$ un morphisme séparé et localement quasi-fini,
où $Y$ est affine. Alors tout ensemble fini de points de $X$ est contenu
dans un ouvert affine de $X$.
\end{lemma}

\begin{proof}
Soient $x_1, \ldots, x_n \in X$. Choisissons un ouvert quasi-compact
$U \subset X$ tel que $x_i \in U$. Alors $U \to Y$ est quasi-affine d'après le
lemme \ref{more-morphisms-lemma-quasi-finite-separated-quasi-affine}.
Il existe donc un ouvert affine $V \subset U$ contenant
$x_1, \ldots, x_n$ d'après
Propriétés, lemme \ref{properties-lemma-ample-finite-set-in-affine}.
\end{proof}

\begin{lemma}
\label{more-morphisms-lemma-quasi-finite-finite-over-dense-open}
Soit $f : Y \to X$ un morphisme quasi-fini.
Il existe un ouvert dense $U \subset X$ tel que
$f|_{f^{-1}(U)} : f^{-1}(U) \to U$ soit fini.
\end{lemma}

\begin{proof}
Si $U_i \subset X$, $i \in I$, est une famille d'ouverts telle que les
restrictions $f|_{f^{-1}(U_i)} : f^{-1}(U_i) \to U_i$ soient finies,
alors, pour $U = \bigcup U_i$, la restriction $f|_{f^{-1}(U)} : f^{-1}(U) \to U$
est finie ; voir
Morphismes, lemme \ref{morphisms-lemma-finite-local}.
Le problème est donc local sur $X$, et l'on peut supposer que $X$ est affine.

\medskip\noindent
Supposons $X$ affine.
Écrivons $Y = \bigcup_{j = 1, \ldots, m} V_j$, avec $V_j$ affine.
C'est possible puisque $f$ est quasi-fini et donc, en
particulier, quasi-compact. Chaque $V_j \to X$ est quasi-fini
et séparé. Soit $\eta \in X$ un point générique d'une composante irréductible
de $X$. On déduit de
Morphismes, lemmes
\ref{morphisms-lemma-quasi-finite} et \ref{morphisms-lemma-generically-finite}
qu'il existe un voisinage ouvert $\eta \in U_\eta$ tel que
$f^{-1}(U_\eta) \cap V_j \to U_\eta$ soit fini. On peut choisir $U_\eta$ de sorte
qu'il convienne pour chaque $j = 1, \ldots, m$.
Remarquons que l'ensemble des points génériques de $X$ est dense dans $X$.
On voit donc qu'il existe un ouvert dense $W = \bigcup_\eta U_\eta$
tel que chaque morphisme $f^{-1}(W) \cap V_j \to W$ soit fini.
Il suffit de montrer qu'il existe un ouvert dense $U \subset W$
tel que $f|_{f^{-1}(U)} : f^{-1}(U) \to U$ soit fini.
Ainsi, on peut remplacer $X$ par un sous-schéma ouvert affine de $W$ et
supposer que chaque $V_j \to X$ est fini.

\medskip\noindent
Supposons $X$ affine, $Y = \bigcup_{j = 1, \ldots, m} V_j$ avec les $V_j$ affines,
et les restrictions $f|_{V_j} : V_j \to X$ finies.
Posons
$$
\Delta_{ij} =
\Big(\overline{V_i \cap V_j} \setminus V_i \cap V_j\Big) \cap V_j.
$$
C'est un fermé nulle part dense de $V_j$, car c'est la frontière
du sous-ensemble ouvert $V_i \cap V_j$ dans $V_j$. D'après
Morphismes, lemme \ref{morphisms-lemma-image-nowhere-dense-finite},
l'image $f(\Delta_{ij})$ est un fermé nulle part dense de $X$. D'après
Topologie, lemme \ref{topology-lemma-nowhere-dense},
la réunion $T = \bigcup f(\Delta_{ij})$ est un fermé nulle part dense
de $X$. Ainsi, $U = X \setminus T$ est un ouvert dense de $X$.
Nous affirmons que $f|_{f^{-1}(U)} : f^{-1}(U) \to U$ est fini.
Pour le voir, soit $U' \subset U$ un ouvert affine.
Posons $Y' = f^{-1}(U') = U' \times_X Y$,
$V_j' = Y' \cap V_j = U' \times_X V_j$. Considérons la restriction
$$
f' = f|_{Y'} : Y' \longrightarrow U'
$$
de $f$. Ce morphisme possède alors la propriété suivante :
$Y' = \bigcup_{j = 1, \ldots, m} V'_j$ est un recouvrement ouvert affine,
chaque $V'_j \to U'$ est fini, et $V_i' \cap V_j'$ est (ouvert et) fermé
à la fois dans $V'_i$ et dans $V'_j$. Ainsi, $V_i' \cap V_j'$ est affine, et l'application
$$
\mathcal{O}(V'_i) \otimes_{\mathbf{Z}} \mathcal{O}(V'_j)
\longrightarrow
\mathcal{O}(V'_i \cap V'_j)
$$
est surjective. Cela entraîne que $Y'$ est séparé, voir
Schémas, lemme \ref{schemes-lemma-characterize-separated}.
Enfin, considérons le diagramme commutatif
$$
\xymatrix{
\coprod_{j = 1, \ldots, m} V'_j \ar[rd] \ar[rr] & & Y' \ar[ld] \\
& U' &
}
$$
La flèche sud-est est finie, donc propre, la flèche horizontale est
surjective, et la flèche sud-ouest est séparée. Par conséquent, d'après
Morphismes, lemme \ref{morphisms-lemma-image-proper-is-proper},
on conclut que $Y' \to U'$ est propre. Comme il est aussi quasi-fini,
on voit qu'il est fini d'après le lemme \ref{more-morphisms-lemma-characterize-finite},
ce qui conclut.
\end{proof}

\begin{lemma}
\label{more-morphisms-lemma-stratify-flat-fp-lqf-universally-bounded}
Soit $f : X \to S$ un morphisme plat, localement de présentation finie, séparé,
localement quasi-fini, à fibres universellement bornées. Il existe alors
des sous-ensembles fermés
$$
\emptyset = Z_{-1} \subset Z_0 \subset Z_1 \subset Z_2 \subset
\ldots \subset Z_n = S
$$
tels que, si $S_r = Z_r \setminus Z_{r - 1}$, la stratification
$S = \coprod_{r = 0, \ldots, n} S_r$ soit caractérisée par la propriété
universelle suivante : étant donné $g : T \to S$, la projection
$X \times_S T \to T$ est finie localement libre
de degré $r$ si et seulement si $g(T) \subset S_r$ (ensemblistement).
\end{lemma}

\begin{proof}
Soit $n$ un entier qui majore le degré des fibres de $X \to S$.
D'après Morphismes, lemme \ref{morphisms-lemma-base-change-universally-bounded},
on voit que les degrés des fibres de tout changement de base sont encore majorés par $n$.
En particulier, tous les entiers $r$ qui interviennent dans l'énoncé du lemme
seront $\leq n$. Nous allons démontrer le lemme par récurrence sur $n$. Le cas
initial $n = 0$ est évident.

\medskip\noindent
Nous affirmons que l'ensemble des points $s \in S$
tels que $\deg_{\kappa(s)}(X_s) = n$ est un sous-ensemble ouvert $S_n \subset S$
et que $X \times_S S_n \to S_n$ est fini localement libre de degré $n$.
En effet, soit $s \in S$ un tel point. Choisissons un morphisme étale
\'elementaire $(U, u) \to (S, s)$ et une décomposition
$U \times_S X = W \amalg V$ comme dans le
lemme \ref{more-morphisms-lemma-etale-splits-off-quasi-finite-part}.
Puisque $V \to U$ est fini, plat et localement de présentation finie,
on voit que $V \to U$ est fini localement libre, voir
Morphismes, lemme \ref{morphisms-lemma-finite-flat}.
Après avoir rétréci $U$ à un voisinage plus petit de $u$,
on peut supposer que $V \to U$ est fini localement libre d'un certain degré $d$, voir
Morphismes, lemme \ref{morphisms-lemma-finite-locally-free}.
Comme $u \mapsto s$ et $W_u = \emptyset$, on voit que $d = n$. Puisque $n$
est le degré maximal d'une fibre, on voit que $W = \emptyset$ !
Ainsi, $U \times_S X \to U$ est fini localement libre de degré $n$.
D'après Descente, lemme \ref{descent-lemma-descending-property-finite-locally-free},
on conclut que $X \to S$ est fini localement libre de degré $n$
au-dessus de $\Im(U \to S)$, qui est un voisinage ouvert de $s$
(Morphismes, lemme \ref{morphisms-lemma-etale-open}).
Ceci démontre l'affirmation.

\medskip\noindent
Soit $S' = S \setminus S_n$, muni de la structure de schéma réduit induite,
et posons $X' = X \times_S S'$. Remarquons que les degrés des fibres
de $X' \to S'$ sont universellement majorés par $n - 1$. Par récurrence, on trouve une
stratification $S' = S_0 \amalg \ldots \amalg S_{n - 1}$ adaptée
au morphisme $X' \to S'$. Nous affirmons que $S = \coprod_{r = 0, \ldots, n} S_r$
convient au morphisme $X \to S$. Soit $g : T \to S$ un morphisme de schémas
et supposons que $X \times_S T \to T$ soit fini localement libre de degré $r$.
Comme on l'a remarqué plus haut, cela implique que $r \leq n$. Si $r = n$, il est
clair que $T \to S$ se factorise par $S_n$. Si $r < n$, alors
$g(T) \subset S' = S \setminus S_d$ (ensemblistement), donc
$T_{red} \to S$ se factorise par $S'$, voir
Schémas, lemme \ref{schemes-lemma-map-into-reduction}.
Remarquons que $X \times_S T_{red} \to T_{red}$ est lui aussi
fini localement libre de degré $r$, comme changement de base.
Par la propriété universelle de la stratification
$S' = \coprod_{r = 0, \ldots, n - 1} S_r$, on voit que $g(T) = g(T_{red})$
est contenu dans $S_r$.
Réciproquement, supposons donné $g : T \to S$ tel que
$g(T) \subset S_r$ (ensemblistement).
Si $r = n$, alors $g$ se factorise par $S_n$ et
il est clair que $X \times_S T \to T$
est fini localement libre de degré $n$, comme changement de base.
Si $r < n$, alors $X \times_S T \to T$ est un morphisme
séparé, plat et localement de présentation finie, tel que
sa restriction à $T_{red}$ soit finie localement libre de degré $r$.
Puisque $T_{red} \to T$ est un homéomorphisme universel, on conclut
que $X \times_S T_{red} \to X \times_S T$ est aussi un homéomorphisme universel
et donc que $X \times_S T \to T$ est universellement fermé (puisque ceci
est vrai pour le morphisme fini $X \times_S T_{red} \to T_{red}$).
Il s'ensuit que $X \times_S T \to T$ est fini, par exemple d'après le
lemme \ref{more-morphisms-lemma-characterize-finite}. On peut alors utiliser
Morphismes, lemme \ref{morphisms-lemma-finite-flat},
pour voir que $X \times_S T \to T$ est fini localement libre.
Enfin, son degré est $r$, puisque toutes ses fibres sont de degré $r$.
\end{proof}

\begin{lemma}
\label{more-morphisms-lemma-stratify-flat-fp-qf}
Soit $f : X \to S$ un morphisme de schémas plat, localement de
présentation finie, séparé et quasi-fini. Il existe alors
des sous-ensembles fermés
$$
\emptyset = Z_{-1} \subset Z_0 \subset Z_1 \subset Z_2 \subset
\ldots \subset S
$$
tels que, si $S_r = Z_r \setminus Z_{r - 1}$, la stratification
$S = \coprod S_r$ soit caractérisée par la propriété universelle suivante :
pour tout morphisme $g : T \to S$, la projection $X \times_S T \to T$ est
finie localement libre de degré $r$ si et seulement si $g(T) \subset S_r$
(ensemblistement). De plus, les morphismes d'inclusion $S_r \to S$ sont
quasi-compacts.
\end{lemma}

\begin{proof}
La question est locale sur $S$ ; on peut donc supposer $S$ affine.
D'après Morphismes, lemme
\ref{morphisms-lemma-locally-quasi-finite-qc-source-universally-bounded},
les fibres de $f$ sont alors universellement bornées.
L'existence de la stratification résulte donc du
lemme \ref{more-morphisms-lemma-stratify-flat-fp-lqf-universally-bounded}.

\medskip\noindent
Nous allons montrer que $U_r = S \setminus Z_r \to S$ est quasi-compact pour
chaque $r \geq 0$. Cela démontrera la dernière assertion par topologie élémentaire.
Comme une composée d'applications quasi-compactes est quasi-compacte,
il suffit de démontrer que $U_r \to U_{r - 1}$ est quasi-compact.
Choisissons un ouvert affine $W \subset U_{r - 1}$. Écrivons $W = \Spec(A)$.
Alors $Z_r \cap W = V(I)$ pour un certain idéal $I \subset A$,
et $X \times_S \Spec(A/I) \to \Spec(A/I)$ est fini localement
libre de degré $r$. Remarquons que $A/I = \colim A/I_i$, où $I_i \subset I$
parcourt les idéaux de type fini. D'après
Limites, lemme \ref{limits-lemma-descend-finite-locally-free},
on voit que $X \times_S \Spec(A/I_i) \to \Spec(A/I_i)$
est fini localement libre de degré $r$ pour un certain $i$. (On utilise ici
que $X \to S$ est de présentation finie, puisqu'il est localement de présentation
finie, séparé et quasi-compact.)
Ainsi, $\Spec(A/I_i) \to \Spec(A) = W$ se factorise (ensemblistement)
par $Z_r \cap W$. Il s'ensuit que $Z_r \cap W = V(I_i)$ est l'ensemble des zéros
d'une famille finie d'éléments de $A$. Cela signifie que
$W \setminus Z_r$ est une réunion finie d'ouverts principaux, donc est quasi-compact,
comme souhaité.
\end{proof}

\begin{lemma}
\label{more-morphisms-lemma-stratify-flat-fp-lqf}
Soit $f : X \to S$ un morphisme de schémas plat, localement de présentation finie,
séparé et localement quasi-fini. Il existe alors
des sous-schémas ouverts
$$
S = U_0 \supset U_1 \supset U_2 \supset \ldots
$$
tels qu'un morphisme $\Spec(k) \to S$, où $k$ est un corps,
se factorise par $U_d$ si et
seulement si $X \times_S \Spec(k)$ est de degré $\geq d$ sur $k$.
\end{lemma}

\begin{proof}
L'énoncé signifie simplement que l'ensemble des points où le degré
de la fibre est $\geq d$ est ouvert. On peut donc travailler localement sur $S$ et
supposer $S$ affine. Dans ce cas, pour tout ouvert quasi-compact $W \subset X$,
l'ensemble $U_d(W)$ des points où les fibres de $W \to S$ sont de
degré $\geq d$ est ouvert d'après le lemme \ref{more-morphisms-lemma-stratify-flat-fp-qf}.
Puisque $U_d = \bigcup_W U_d(W)$, le résultat s'ensuit.
\end{proof}

\begin{lemma}
\label{more-morphisms-lemma-go-down-with-annihilators}
Soit $f : X \to S$ un morphisme de schémas plat, localement de
présentation finie et localement quasi-fini. Soit
$g \in \Gamma(X, \mathcal{O}_X)$ non nul. Il existe alors
un ouvert $V \subset X$ tel que $g|_V \not = 0$, un ouvert
$U \subset S$ s'insérant dans un diagramme commutatif
$$
\xymatrix{
V \ar[r] \ar[d]_\pi & X \ar[d]^f \\
U \ar[r] & S,
}
$$
un sous-faisceau quasi-cohérent $\mathcal{F} \subset \mathcal{O}_U$, un entier
$r > 0$, et un homomorphisme injectif de $\mathcal{O}_U$-modules
$\mathcal{F}^{\oplus r} \to \pi_*\mathcal{O}_V$
dont l'image contient $g|_V$.
\end{lemma}

\begin{proof}
On peut supposer $X$ et $S$ affines. On obtient une filtration
$\emptyset = Z_{-1} \subset Z_0 \subset Z_1 \subset Z_2 \subset \ldots
\subset Z_n = S$ comme dans les
lemmes \ref{more-morphisms-lemma-stratify-flat-fp-lqf-universally-bounded} et
\ref{more-morphisms-lemma-stratify-flat-fp-qf}.
Soit $T \subset X$ le support schématique du module fini
sur $\mathcal{O}_X$ donné par $\Im(g : \mathcal{O}_X \to \mathcal{O}_X)$.
Remarquons que $T$ est le support de $g$ comme section de $\mathcal{O}_X$
(Modules, définition \ref{modules-definition-support}) et que,
pour tout ouvert $V \subset X$, on a $g|_V \not = 0$ si et seulement si
$V \cap T \not = \emptyset$.
Soit $r$ le plus petit entier tel que $f(T) \subset Z_r$
ensemblistement. Soit $\xi \in T$ un point générique d'une composante irréductible
de $T$ tel que $f(\xi) \not \in Z_{r - 1}$ (et donc
$f(\xi) \in Z_r$). On peut remplacer $S$ par un voisinage affine de
$f(\xi)$ contenu dans $S \setminus Z_{r - 1}$. Écrivons $S = \Spec(A)$
et soit $I = (a_1, \ldots, a_m) \subset A$ un idéal de type fini
tel que $V(I) = Z_r$ (ensemblistement, voir
Algèbre, lemme \ref{algebra-lemma-qc-open}).
Puisque le support de $g$ est contenu dans $f^{-1}V(I)$ par notre choix de $r$,
on voit qu'il existe un entier $N$ tel que
$a_j^N g = 0$ pour $j = 1, \ldots, m$. En remplaçant $a_j$ par $a_j^r$,
on peut supposer que $Ig = 0$. Pour tout $A$-module $M$, notons
$M[I]$ le sous-module des éléments annulés par $I$ dans $M$, c'est-à-dire $M[I] = \{m \in M \mid Im = 0\}$.
Écrivons $X = \Spec(B)$, de sorte que $g \in B[I]$. Comme $A \to B$ est plat, on
voit que
$$
B[I] = A[I] \otimes_A B \cong A[I] \otimes_{A/I} B/IB
$$
Par notre choix de $Z_r$, le $A/I$-module $B/IB$ est
fini localement libre de rang $r$. Ainsi, après avoir remplacé $S$ par
un voisinage ouvert affine plus petit de $f(\xi)$, on peut supposer
que $B/IB \cong (A/IA)^{\oplus r}$ comme $A/I$-modules.
Choisissons une application $\psi : A^{\oplus r} \to B$ dont la réduction modulo $I$ est
l'isomorphisme de la phrase précédente. On voit alors que
l'application induite
$$
A[I]^{\oplus r} \longrightarrow B[I]
$$
est un isomorphisme. Le lemme s'ensuit en prenant pour $\mathcal{F}$ le
faisceau quasi-cohérent associé au $A$-module $A[I]$, et pour
l'application $\mathcal{F}^{\oplus r} \to \pi_*\mathcal{O}_V$ celle
qui correspond à $A[I]^{\oplus r} \subset A^{\oplus r} \to B$.
\end{proof}

\begin{lemma}
\label{more-morphisms-lemma-there-is-a-scheme-integral-over}
Soit $U \to X$ un morphisme \'etale surjectif de schémas. Supposons $X$
quasi-compact et quasi-séparé. Il existe alors un morphisme entier surjectif
$Y \to X$ tel que, pour
tout $y \in Y$, il existe un voisinage ouvert $V \subset Y$
tel que $V \to X$ se factorise par $U$. En fait, on peut supposer
$Y \to X$ fini et de présentation finie.
\end{lemma}

\begin{proof}
Puisque $X$ est quasi-compact, il existe un nombre fini d'ouverts affines
$U_i \subset U$ tels que $U' = \coprod U_i \to X$ soit surjectif.
Après avoir remplacé $U$ par $U'$, on voit que l'on peut supposer $U$ affine.
En particulier, $U \to X$ est séparé
(Schémas, lemme \ref{schemes-lemma-affine-separated}).
Il existe alors un entier $d$ majorant le degré des fibres géométriques
de $U \to X$ (voir Morphismes, lemme
\ref{morphisms-lemma-locally-quasi-finite-qc-source-universally-bounded}).
Nous allons démontrer le lemme par récurrence sur $d$ pour tous les schémas
quasi-compacts et séparés $U$ se projetant de manière \'etale et surjective sur $X$.
Si $d = 1$, alors $U = X$ et le résultat vaut avec $Y = U$.
Supposons $d > 1$.

\medskip\noindent
On applique le lemme \ref{more-morphisms-lemma-quasi-finite-separated-quasi-affine}
et on obtient une factorisation
$$
\xymatrix{
U \ar[rr]_j \ar[rd] & & Y \ar[ld]^\pi \\
& X
}
$$
où $\pi$ est entier et $j$ une immersion ouverte quasi-compacte. On peut, et l'on va,
supposer que $j(U)$ est schématiquement dense dans $Y$. Remarquons que
$$
U \times_X Y = U \amalg W
$$
où le premier terme est l'image de $U \to U \times_X Y$
(qui est fermée d'après
Schémas, lemme \ref{schemes-lemma-semi-diagonal},
et ouverte parce qu'elle est \'etale comme morphisme entre schémas \'etales sur $Y$),
et le second facteur est le complémentaire (ouvert et fermé).
L'image $V \subset Y$ de $W$ est un sous-schéma ouvert contenant
$Y \setminus U$.

\medskip\noindent
Le morphisme \'etale $W \to Y$ a des fibres géométriques de cardinal $< d$.
En effet, cela est clair par inspection pour les points géométriques de $U \subset Y$.
Puisque $U \subset Y$ est dense, cela vaut pour tous les points géométriques de $Y$,
par exemple d'après le lemme
\ref{more-morphisms-lemma-stratify-flat-fp-lqf-universally-bounded}
(le degré des fibres d'un morphisme \'etale séparé quasi-compact
ne croît pas par spécialisation). On peut donc appliquer l'hypothèse de récurrence
à $W \to V$ et trouver un morphisme entier surjectif
$Z \to V$, avec $Z$ un schéma, qui se factorise localement pour la topologie de Zariski par $W$.
Choisissons une factorisation $Z \to Z' \to Y$ où $Z' \to Y$ est entier et
$Z \to Z'$ une immersion ouverte
(lemme \ref{more-morphisms-lemma-quasi-finite-separated-quasi-affine}).
Après avoir remplacé $Z'$ par l'adhérence schématique de $Z$ dans $Z'$,
on peut supposer $Z$ schématiquement dense dans $Z'$.
Après cette opération, on a $Z' \times_Y V = Z$. Enfin,
soit $T \subset Y$ le sous-schéma fermé réduit induit
sur $Y \setminus V$. Considérons le morphisme
$$
Z' \amalg T \longrightarrow X
$$
C'est, par construction, un morphisme entier surjectif.
Puisque $T \subset U$, il est clair que le morphisme $T \to X$
se factorise par $U$. D'autre part, soit $z \in Z'$
un point. Si $z \not \in Z$, alors $z$ s'envoie sur un point de
$Y \setminus V \subset U$, et l'on trouve un voisinage de $z$
sur lequel le morphisme se factorise par $U$.
Si $z \in Z$, alors il existe un voisinage $\Omega \subset Z$
qui se factorise par $W \subset U \times_X Y$, et donc par $U$.
Ceci démontre l'existence.

\medskip\noindent
Supposons avoir trouvé un morphisme entier surjectif $Y \to X$ qui se factorise
localement pour la topologie de Zariski par $U$. Choisissons un recouvrement ouvert affine fini
$Y = \bigcup V_j$ tel que $V_j \to X$ se factorise par $U$. On peut
écrire $Y = \lim Y_i$, où $Y_i \to X$ est fini et de présentation
finie, voir Limites, lemme
\ref{limits-lemma-integral-limit-finite-and-finite-presentation}.
Pour $i$ assez grand, on peut trouver des ouverts affines $V_{i, j} \subset Y_i$
dont les images réciproques dans $Y$ redonnent les $V_j$, voir
Limites, lemme \ref{limits-lemma-descend-opens}.
Pour $i$ encore plus grand, les morphismes $V_j \to U$ au-dessus de $X$
proviennent de morphismes $V_{i, j} \to U$ au-dessus de $X$, voir
Limites, proposition
\ref{limits-proposition-characterize-locally-finite-presentation}.
Ceci achève la démonstration.
\end{proof}







\section{Application aux morphismes à fibres connexes}
\label{more-morphisms-section-application-geometrically-connected}

\noindent
Dans cette section, nous démontrons quelques lemmes qui produisent des morphismes
dont toutes les fibres sont géométriquement connexes ou géométriquement intègres.
Ils nous seront utiles plus loin dans notre étude de la structure locale des morphismes
de type fini.

\begin{lemma}
\label{more-morphisms-lemma-descent-connected-fibres}
Considérons un diagramme de morphismes de schémas
$$
\xymatrix{
Z \ar[r]_{\sigma} \ar[rd] & X \ar[d] \\
& Y
}
$$
et un point $y \in Y$. Supposons que
\begin{enumerate}
\item $X \to Y$ soit de présentation finie et plat,
\item $Z \to Y$ soit fini localement libre,
\item $Z_y \not = \emptyset$,
\item toutes les fibres de $X \to Y$ soient géométriquement réduites, et
\item $X_y$ soit géométriquement connexe sur $\kappa(y)$.
\end{enumerate}
Il existe alors un ouvert quasi-compact $X^0 \subset X$ tel que $X^0_y = X_y$
et tel que toutes les fibres non vides de $X^0 \to Y$ soient géométriquement connexes.
\end{lemma}

\begin{proof}
Dans cette démonstration, nous utiliserons que les propriétés d'être plat, d'être de présentation
finie et d'être fini localement libre sont préservées par changement de base et par composition.
Nous utiliserons aussi qu'un morphisme fini localement libre est à la fois ouvert et
fermé. On trouvera ces faits dans
Morphismes, lemmes
\ref{morphisms-lemma-base-change-flat},
\ref{morphisms-lemma-base-change-finite-presentation},
\ref{morphisms-lemma-base-change-finite-locally-free},
\ref{morphisms-lemma-composition-flat},
\ref{morphisms-lemma-composition-finite-presentation},
\ref{morphisms-lemma-composition-finite-locally-free},
\ref{morphisms-lemma-fppf-open}, et
\ref{morphisms-lemma-finite-proper}.

\medskip\noindent
Remarquons que $X_Z \to Z$ est un morphisme plat de présentation finie
qui possède une section $s$ provenant de $\sigma$. Notons $X_Z^0$
le sous-ensemble de $X_Z$ défini dans la
situation \ref{more-morphisms-situation-connected-along-section}.
D'après le
lemme \ref{more-morphisms-lemma-connected-along-section-open},
c'est un sous-ensemble ouvert de $X_Z$.

\medskip\noindent
L'image réciproque $X_{Z \times_Y Z}$ de $X$ sur $Z \times_Y Z$ est munie
de deux sections $s_0, s_1$, à savoir les changements de base de $s$ par
$\text{pr}_0, \text{pr}_1 : Z \times_Y Z \to Z$. La construction de la
situation \ref{more-morphisms-situation-connected-along-section}
donne deux sous-ensembles $(X_{Z \times_Y Z})_{s_0}^0$ et
$(X_{Z \times_Y Z})_{s_1}^0$. D'après le
lemme \ref{more-morphisms-lemma-base-change-connected-along-section},
ce sont les images réciproques de $X_Z^0$ par les morphismes
$1_X \times \text{pr}_0, 1_X \times \text{pr}_1 : X_{Z \times_Y Z} \to X_Z$.
En particulier, ces sous-ensembles sont ouverts.

\medskip\noindent
Soit $(Z \times_Y Z)_y = \{z_1, \ldots, z_n\}$.
Comme $X_y$ est géométriquement connexe, on voit que les fibres de
$(X_{Z \times_Y Z})_{s_0}^0$ et de $(X_{Z \times_Y Z})_{s_1}^0$
au-dessus de chaque $z_i$ coïncident (elles sont égales à la fibre entière). Une autre
manière d'exprimer ceci est
$$
s_0(z_i) \in (X_{Z \times_Y Z})_{s_1}^0
\quad\text{et}\quad
s_1(z_i) \in (X_{Z \times_Y Z})_{s_0}^0.
$$
Puisque les ensembles $(X_{Z \times_Y Z})_{s_0}^0$ et $(X_{Z \times_Y Z})_{s_1}^0$
sont ouverts dans $X_{Z \times_Y Z}$, il existe un voisinage ouvert
$W \subset Z \times_Y Z$ de $(Z \times_Y Z)_y$ tel que
$$
s_0(W) \subset (X_{Z \times_Y Z})_{s_1}^0
\quad\text{et}\quad
s_1(W) \subset (X_{Z \times_Y Z})_{s_0}^0.
$$
Il résulte alors directement de la construction de la
situation \ref{more-morphisms-situation-connected-along-section}
que
$$
p^{-1}(W) \cap (X_{Z \times_Y Z})_{s_0}^0
=
p^{-1}(W) \cap (X_{Z \times_Y Z})_{s_1}^0
$$
où $p : X_{Z \times_Y Z} \to Z \times_W Z$ est la projection.
Puisque $Z \times_Y Z \to Y$ est fini localement libre, donc ouvert et fermé,
il existe un voisinage ouvert affine $V \subset Y$ de $y$ tel que
$q^{-1}(V) \subset W$, où $q : Z \times_Y Z \to Y$ est le
morphisme structural. Pour démontrer le lemme, on peut remplacer $Y$ par $V$.
Après ce remplacement, on voit que $X_Z^0 \subset Y_Z$ est un ouvert tel que
$$
(1_X \times \text{pr}_0)^{-1}(X_Z^0) =
(1_X \times \text{pr}_1)^{-1}(X_Z^0).
$$
Cela signifie que l'image $X^0 \subset X$ de $X_Z^0$ est un ouvert tel
que $(X_Z \to X)^{-1}(X^0) = X_Z^0$, voir
Descente, lemme \ref{descent-lemma-open-fpqc-covering}.
Enfin, $X^0$ est quasi-compact parce que $X_Z^0$ est quasi-compact
d'après le lemme \ref{more-morphisms-lemma-connected-along-section-locally-constructible}
(on utilise qu'à ce stade $Y$ est affine, donc $X$ est quasi-compact et
quasi-séparé, de sorte que localement constructible équivaut à constructible
et entraîne en particulier quasi-compact ; détails omis).
On voit ainsi que $X^0$ possède toutes les propriétés voulues.
\end{proof}

\begin{lemma}
\label{more-morphisms-lemma-fibre-geometrically-connected-reduced}
Soit $h : Y \to S$ un morphisme de schémas.
Soit $s \in S$ un point.
Soit $T \subset Y_s$ un sous-schéma ouvert.
Supposons que
\begin{enumerate}
\item $h$ soit plat et de présentation finie,
\item toutes les fibres de $h$ soient géométriquement réduites, et
\item $T$ soit géométriquement connexe sur $\kappa(s)$.
\end{enumerate}
On peut alors trouver un voisinage \'etale élémentaire affine
$(S', s') \to (S, s)$
et un ouvert quasi-compact $V \subset Y_{S'}$ tels que
\begin{enumerate}
\item[(a)] toutes les fibres de $V \to S'$ soient géométriquement connexes,
\item[(b)] $V_{s'} = T \times_s s'$.
\end{enumerate}
\end{lemma}

\begin{proof}
Le problème est manifestement local sur $S$ ; on peut donc remplacer $S$ par un
voisinage ouvert affine de $s$.
La topologie de $Y_s$ est induite par celle de $Y$, voir
Schémas, lemme \ref{schemes-lemma-fibre-topological}.
On peut donc trouver un ouvert quasi-compact $V \subset Y$ tel que $V_s = T$.
La restriction de $h$ à $V$ est quasi-compacte (car $S$ est affine et $V$
quasi-compact), quasi-séparée, localement de présentation finie et
plate, donc plate de présentation finie.
Ainsi, après avoir remplacé $Y$ par $V$, on peut supposer, en plus
de (1) et (2), que $Y_s = T$ et que $S$ est affine.

\medskip\noindent
Choisissons un point fermé $y \in Y_s$ tel que $h$ soit de Cohen--Macaulay en $y$, voir
le lemme \ref{more-morphisms-lemma-flat-finite-presentation-CM-open}.
D'après le
lemme \ref{more-morphisms-lemma-slice-CM},
il existe un diagramme
$$
\xymatrix{
Z \ar[r] \ar[rd] & Y \ar[d] \\
& S
}
$$
tel que $Z \to S$ soit plat, localement de présentation finie, localement
quasi-fini, avec $Z_s = \{y\}$. Appliquons le
lemme \ref{more-morphisms-lemma-etale-makes-quasi-finite-finite-at-point}
pour trouver un voisinage élémentaire $(S', s') \to (S, s)$ et un ouvert
$Z' \subset Z_{S'} = S' \times_S Z$ tel que $Z' \to S'$ soit fini et possède un unique
point $z' \in Z'$ au-dessus de $s$. Remarquons que $Z' \to S'$ est aussi
localement de présentation finie et plat (comme ouvert du changement de base
de $Z \to S$), donc $Z' \to S'$ est fini localement libre, voir
Morphismes, lemme \ref{morphisms-lemma-finite-flat}.
Remarquons que $Y_{S'} \to S'$ est plat et de présentation finie,
à fibres géométriquement réduites, comme changement de base de $h$.
De plus, $Y_{s'} = Y_s$ est géométriquement connexe.
Appliquons le lemme \ref{more-morphisms-lemma-descent-connected-fibres}
à $Z' \to Y_{S'}$ au-dessus de $S'$ pour obtenir un ouvert quasi-compact $V \subset Y_{S'}$
satisfaisant (2), dont les fibres sur $S'$ sont soit vides, soit
géométriquement connexes. Comme $V \to S'$ est ouvert
(Morphismes, lemme \ref{morphisms-lemma-fppf-open}), après avoir remplacé
$S'$ par un voisinage ouvert affine de $s'$,
on peut supposer $V \to S'$ surjectif, d'où (1).
\end{proof}

\begin{lemma}
\label{more-morphisms-lemma-cover-by-geometrically-connected}
Soit $f : X \to S$ un morphisme de schémas
localement de présentation finie et plat, à fibres géométriquement
réduites. Il existe alors
un recouvrement \'etale $\{X_i \to X\}_{i \in I}$
tel que $X_i \to S$ se factorise en $X_i \to S_i \to S$,
où $S_i \to S$ est \'etale et $X_i \to S_i$ est
plat de présentation finie, à fibres géométriquement connexes
et géométriquement réduites.
\end{lemma}

\begin{proof}
Choisissons un point $x \in X$ d'image $s \in S$. Nous allons construire
un diagramme
$$
\xymatrix{
X' \ar[r] \ar[rd] & S' \times_S X \ar[r] \ar[d] & X \ar[d] \\
& S' \ar[r] & S
}
$$
et des points $s' \in S'$, $x' \in X'$, $y \in S' \times_S X$
tels que $x'$ s'envoie sur $x$, que $(S', s') \to (S, s)$
soit un voisinage \'etale, que $(X', x') \to (S' \times_S X, y)$
soit un voisinage \'etale\footnote{La démonstration donne en fait
un ouvert $X' \subset S' \times_S X$.}, et que
$X' \to S'$ ait des fibres géométriquement
connexes. Si l'on peut faire cela pour tout $x \in X$, alors
le lemme s'ensuit (les membres du recouvrement étant donnés par la
famille des morphismes \'etales $X' \to X$ ainsi obtenus).
La première étape consiste à remplacer $X$ et $S$ par des voisinages ouverts affines
de $x$ et de $s$, ce qui nous ramène au cas où $X$ et $S$ sont affines
(et donc où $f$ est de présentation finie).

\medskip\noindent
Choisissons une extension algébrique séparable $\overline{k}$ de $\kappa(s)$.
Notons $X_{\overline{k}}$ le changement de base de $X_s$.
Choisissons un point $\overline{x}$ de $X_{\overline{k}}$ s'envoyant sur $x \in X_s$.
Choisissons un voisinage ouvert connexe quasi-compact
$\overline{V} \subset X_{\overline{k}}$
de $\overline{x}$. (Cela est possible parce que tout schéma
localement de type fini sur un corps est localement connexe
comme espace topologique localement noethérien.)
D'après Variétés, lemme \ref{varieties-lemma-Galois-action-quasi-compact-open},
on peut trouver une extension finie séparable $k'/\kappa(s)$
et un ouvert quasi-compact $V' \subset X_{k'}$ dont le
changement de base est $\overline{V}$. En particulier, $V'$ est
géométriquement connexe sur $k'$, voir
Variétés, lemme \ref{varieties-lemma-characterize-geometrically-connected}. D'après
le lemme \ref{more-morphisms-lemma-realize-prescribed-residue-field-extension-etale},
on peut trouver un voisinage \'etale $(S', s') \to (S, s)$
tel que $\kappa(s')$ soit isomorphe à $k'$ comme extension
de $\kappa(s)$.
Notons $x' \in (S' \times_S X)_{s'} = X_{k'}$ l'image de $\overline{x}$.
Ainsi, après avoir remplacé $(S, s)$ par $(S', s')$ et $(X, x)$ par
$(S' \times_S X, x')$, on se ramène au cas traité au
paragraphe suivant.

\medskip\noindent
Supposons qu'il existe un ouvert quasi-compact $V \subset X_s$
qui contienne $x$ et soit géométriquement connexe.
On peut alors appliquer le lemme \ref{more-morphisms-lemma-fibre-geometrically-connected-reduced}
pour trouver un voisinage \'etale affine $(S', s') \to (S, s)$
et un ouvert quasi-compact $X' \subset S' \times_S X$ tels que
$X' \to S'$ ait des fibres géométriquement connexes
et que $X'$ contienne un point au-dessus de $x$.
Ceci achève la démonstration.
\end{proof}

\begin{lemma}
\label{more-morphisms-lemma-normal-morphism-irreducible}
Soit $h : Y \to S$ un morphisme de schémas.
Soit $s \in S$ un point.
Soit $T \subset Y_s$ un sous-schéma ouvert.
Supposons que
\begin{enumerate}
\item $h$ soit de présentation finie,
\item $h$ soit normal, et
\item $T$ soit géométriquement irréductible sur $\kappa(s)$.
\end{enumerate}
Alors on peut trouver un voisinage \'etale élémentaire affine
$(S', s') \to (S, s)$ et un ouvert quasi-compact $V \subset Y_{S'}$ tels que
\begin{enumerate}
\item[(a)] toutes les fibres de $V \to S'$ soient géométriquement intègres,
\item[(b)] $V_{s'} = T \times_s s'$.
\end{enumerate}
\end{lemma}

\begin{proof}
Appliquons le
lemme \ref{more-morphisms-lemma-fibre-geometrically-connected-reduced}
pour trouver un voisinage \'etale élémentaire affine $(S', s') \to (S, s)$ et
un ouvert quasi-compact $V \subset Y_{S'}$ tels que toutes les fibres de
$V \to S'$ soient géométriquement connexes et $V_{s'} = T \times_s s'$.
Comme $V$ est un ouvert du changement de base de $h$, toutes les fibres de $V \to S'$
sont géométriquement normales, voir le lemme \ref{more-morphisms-lemma-normal}.
En particulier, elles sont géométriquement réduites. Pour achever la démonstration,
il reste à montrer qu'elles sont géométriquement irréductibles. Or, si $t \in S'$,
alors $V_t$ est de type fini sur $\kappa(t)$ et donc
$V_t \times_{\kappa(t)} \overline{\kappa(t)}$ est de type fini
sur $\overline{\kappa(t)}$, donc noethérien. Par le choix de $S' \to S$,
le schéma $V_t \times_{\kappa(t)} \overline{\kappa(t)}$ est connexe.
Ainsi $V_t \times_{\kappa(t)} \overline{\kappa(t)}$ est irréductible d'après
Propriétés, lemme \ref{properties-lemma-normal-Noetherian},
ce qui achève la démonstration.
\end{proof}








\section{Application à la structure des morphismes de type fini}
\label{more-morphisms-section-application-finite-type}

\noindent
Le résultat de cette section se trouve dans \cite{GruRay}.
En termes informels, il affirme qu'un morphisme de type fini est, localement pour la topologie \'etale
sur la source et le but, la composée d'un
morphisme fini suivi d'un morphisme lisse à fibres géométriquement connexes
de dimension relative égale à la dimension de la fibre du morphisme
initial.

\begin{lemma}
\label{more-morphisms-lemma-local-structure-finite-type}
Soit $f : X \to S$ un morphisme. Soit $x \in X$ et posons $s = f(x)$.
Supposons que $f$ soit localement de type fini et que $n = \dim_x(X_s)$.
Alors il existe un diagramme commutatif
$$
\xymatrix{
X \ar[dd] & X' \ar[l]^g \ar[d]^\pi & x \ar@{|->}[dd] &
x' \ar@{|->}[l]  \ar@{|->}[d] \\
& Y \ar[d]^h & & y \ar@{|->}[d] \\
S \ar@{=}[r] & S & s & s \ar@{=}[l]
}
$$
et un point $x' \in X'$ avec $g(x') = x$ tels que, si $y = \pi(x')$,
on ait
\begin{enumerate}
\item $h : Y \to S$ soit lisse de dimension relative $n$,
\item $g : (X', x') \to (X, x)$ soit un voisinage \'etale élémentaire,
\item $\pi$ soit fini, et $\pi^{-1}(\{y\}) = \{x'\}$, et
\item $\kappa(y)$ soit une extension purement transcendante de $\kappa(s)$.
\end{enumerate}
De plus, si $f$ est localement de présentation finie, alors $\pi$ est
de présentation finie.
\end{lemma}

\begin{proof}
Le problème est local sur $X$ et $S$ ; on peut donc supposer $X$ et
$S$ affines. D'après
Algèbre, lemme \ref{algebra-lemma-refined-quasi-finite-over-polynomial-algebra},
après avoir remplacé $X$ par un voisinage ouvert standard de $x$ dans $X$,
on peut supposer qu'il existe une factorisation
$$
\xymatrix{
X \ar[r]^\pi & \mathbf{A}^n_S \ar[r] & S
}
$$
telle que $\pi$ soit quasi-fini et que $\kappa(\pi(x))$
soit une extension purement transcendante de $\kappa(s)$. D'après le
lemme \ref{more-morphisms-lemma-etale-makes-quasi-finite-finite-at-point},
il existe un voisinage \'etale élémentaire
$$
(Y, y) \to (\mathbf{A}^n_S, \pi(x))
$$
et un ouvert $X' \subset X \times_{\mathbf{A}^n_S} Y$ qui contient un
unique point $x'$ au-dessus de $y$, tel que $X' \to Y$ soit fini.
Ceci démontre les assertions (1) -- (4). Pour la dernière assertion, utiliser
Morphismes, lemme \ref{morphisms-lemma-finite-presentation-permanence}.
\end{proof}

\begin{lemma}
\label{more-morphisms-lemma-local-local-structure-finite-type}
\begin{slogan}
Un morphisme de type fini est, dans des voisinages \'etales, fini au-dessus d'un
morphisme lisse.
\end{slogan}
Soit $f : X \to S$ un morphisme. Soit $x \in X$ et posons $s = f(x)$.
Supposons que $f$ soit localement de type fini et que $n = \dim_x(X_s)$.
Alors il existe un diagramme commutatif
$$
\xymatrix{
X \ar[dd] & X' \ar[l]^g \ar[d]^\pi & x \ar@{|->}[dd] &
x' \ar@{|->}[l] \ar@{|->}[d] \\
& Y' \ar[d]^h & & y' \ar@{|->}[d] \\
S & S' \ar[l]_e & s & s' \ar@{|->}[l]
}
$$
et un point $x' \in X'$ avec $g(x') = x$ tels que, si $y' = \pi(x')$,
$s' = h(y')$, on ait
\begin{enumerate}
\item $h : Y' \to S'$ soit lisse de dimension relative $n$,
\item toutes les fibres de $Y' \to S'$ soient géométriquement intègres,
\item $g : (X', x') \to (X, x)$ soit un voisinage \'etale élémentaire,
\item $\pi$ soit fini, et $\pi^{-1}(\{y'\}) = \{x'\}$,
\item $\kappa(y')$ soit une extension purement transcendante de $\kappa(s')$, et
\item $e : (S', s') \to (S, s)$ soit un voisinage \'etale élémentaire.
\end{enumerate}
De plus, si $f$ est localement de présentation finie, alors $\pi$ est
de présentation finie.
\end{lemma}

\begin{proof}
La question est locale sur $S$ ; on peut donc remplacer $S$ par un voisinage
ouvert affine de $s$. Appliquons ensuite le
lemme \ref{more-morphisms-lemma-local-structure-finite-type}
pour obtenir un diagramme commutatif
$$
\xymatrix{
X \ar[dd] & X' \ar[l]^g \ar[d]^\pi & x \ar@{|->}[dd] &
x' \ar@{|->}[l] \ar@{|->}[d] \\
& Y \ar[d]^h & & y \ar@{|->}[d] \\
S \ar@{=}[r] & S & s & s \ar@{=}[l]
}
$$
où $h$ est lisse de dimension relative $n$ et $\kappa(y)$ est
une extension purement transcendante de $\kappa(s)$. Puisque la question est
aussi locale sur $X$, on peut remplacer $Y$ par un voisinage affine de $y$
(et $X'$ par son image réciproque par $\pi$). Comme $S$ est affine,
ceci assure que $Y \to S$ est quasi-compact, séparé et lisse,
en particulier de présentation finie.
Soit $T$ la composante connexe de $Y_s$ contenant $y$.
Comme $Y_s$ est noethérien, on voit que $T$ est ouvert.
On voit également que $T$ est géométriquement connexe sur $\kappa(s)$ d'après
Variétés,
lemme \ref{varieties-lemma-geometrically-connected-if-connected-and-point}.
Puisque $T$ est aussi lisse sur $\kappa(s)$, il est géométriquement normal, voir
Variétés, lemme \ref{varieties-lemma-smooth-geometrically-normal}.
On en conclut que $T$ est géométriquement irréductible sur $\kappa(s)$ (car un
schéma noethérien normal connexe est irréductible, voir
Propriétés, lemme \ref{properties-lemma-normal-Noetherian}).
Enfin, observons que le morphisme lisse $h$ est normal d'après le
lemme \ref{more-morphisms-lemma-smooth-normal}.
Nous avons alors vérifié que toutes les hypothèses du
lemme \ref{more-morphisms-lemma-normal-morphism-irreducible}
sont satisfaites pour le morphisme $h : Y \to S$ et l'ouvert $T \subset Y_s$.
En appliquant le
lemme \ref{more-morphisms-lemma-normal-morphism-irreducible},
on obtient $e : S' \to S$, $s' \in S'$ et $Y'$ comme
dans le diagramme commutatif
$$
\xymatrix{
X \ar[dd] & X' \ar[l]^g \ar[d]^\pi & X' \times_Y Y' \ar[l] \ar[d] &
x \ar@{|->}[dd] & x' \ar@{|->}[l] \ar@{|->}[d] &
(x', s') \ar@{|->}[l] \ar@{|->}[d] \\
& Y \ar[d]^h & Y' \ar[d] \ar[l] & & y \ar@{|->}[d] &
(y, s') \ar@{|->}[l] \ar@{|->}[d] \\
S \ar@{=}[r] & S & S' \ar[l]_e & s & s \ar@{=}[l] & s' \ar@{|->}[l]
}
$$
où $e : (S', s') \to (S, s)$ est un voisinage \'etale élémentaire,
et où $Y' \subset Y_{S'}$ est un voisinage ouvert dont toutes les fibres
sur $S'$ sont géométriquement irréductibles, tel que $Y'_{s'} = T$ via
l'identification $Y_s = Y_{S', s'}$. Soit $(y, s') \in Y'$ le point
correspondant à $y \in T$ ; c'est aussi l'unique point de $Y \times_S S'$
au-dessus de $y$, de corps résiduel égal à $\kappa(y)$ et dont l'image est $s'$
dans $S'$. De même, soit $(x', s') \in X' \times_Y Y' \subset X' \times_S S'$
l'unique point au-dessus de $x'$, de corps résiduel égal à $\kappa(x')$
et situé au-dessus de $s'$. La partie extérieure de ce diagramme fournit alors une solution au
problème posé dans le lemme. Quelques détails mineurs sont omis.
\end{proof}

\begin{lemma}
\label{more-morphisms-lemma-local-local-structure-finite-type-affine}
Hypothèses et notations comme dans le
lemme \ref{more-morphisms-lemma-local-local-structure-finite-type}.
Outre les propriétés (1) -- (6), on peut imposer en plus que
\begin{enumerate}
\item[(7)] $S'$, $Y'$, $X'$ soient affines.
\end{enumerate}
\end{lemma}

\begin{proof}
Observons que si $Y'$ est affine, alors $X'$ est affine puisque $\pi$ est fini.
Choisissons un voisinage ouvert affine $U' \subset S'$ de $s'$.
Choisissons un voisinage ouvert affine $V' \subset h^{-1}(U')$ de $y'$.
Posons $W' = h(V')$. C'est un voisinage ouvert de $s'$ dans $S'$, voir
Morphismes, lemme \ref{morphisms-lemma-smooth-open},
contenu dans $U'$. Choisissons un voisinage ouvert affine $U'' \subset W'$
de $s'$. Alors $h^{-1}(U'') \cap V'$ est affine, car il est égal à
$U'' \times_{U'} V'$. Par construction,
$h^{-1}(U'') \cap V' \to U''$ est un morphisme lisse et surjectif dont les
fibres sont des sous-schémas ouverts non vides de fibres géométriquement intègres
de $Y' \to S'$, et sont donc géométriquement intègres. On peut ainsi remplacer
$S'$ par $U''$ et $Y'$ par $h^{-1}(U'') \cap V'$.
\end{proof}

\noindent
La portée de la propriété $\pi^{-1}(\{y'\}) = \{x'\}$ est en partie
expliquée par le lemme suivant.

\begin{lemma}
\label{more-morphisms-lemma-finite-morphism-single-point-in-fibre}
Soit $\pi : X \to Y$ un morphisme fini.
Soit $x \in X$, avec $y = \pi(x)$, tel que $\pi^{-1}(\{y\}) = \{x\}$.
Alors
\begin{enumerate}
\item Pour tout voisinage $U \subset X$ de $x$ dans $X$, il
existe un voisinage $V \subset Y$ de $y$ tel que
$\pi^{-1}(V) \subset U$.
\item Le morphisme d'anneaux $\mathcal{O}_{Y, y} \to \mathcal{O}_{X, x}$
est fini.
\item Si $\pi$ est de présentation finie, alors
$\mathcal{O}_{Y, y} \to \mathcal{O}_{X, x}$ est de présentation finie.
\item Pour tout $\mathcal{O}_X$-module quasi-cohérent $\mathcal{F}$,
on a $\mathcal{F}_x = \pi_*\mathcal{F}_y$ comme
$\mathcal{O}_{Y, y}$-modules.
\end{enumerate}
\end{lemma}

\begin{proof}
La première assertion est purement topologique ; utiliser le fait que
$\pi$ est une application continue et fermée telle que $\pi^{-1}(\{y\}) = \{x\}$.
Pour démontrer les deuxième et troisième assertions, on peut supposer que
$X = \Spec(B)$ et $Y = \Spec(A)$. Alors
$A \to B$ est un morphisme fini d'anneaux et $y$ correspond à un idéal premier
$\mathfrak p$ de $A$ tel qu'il existe un unique idéal premier $\mathfrak q$ de
$B$ au-dessus de $\mathfrak p$. Alors
$B_{\mathfrak q} = B_{\mathfrak p}$, voir
Algèbre, lemme \ref{algebra-lemma-unique-prime-over-localize-below}.
Autrement dit, le morphisme $A_{\mathfrak p} \to B_{\mathfrak q}$
est égal au morphisme $A_{\mathfrak p} \to B_{\mathfrak p}$ que l'on obtient
en localisant $A \to B$ en $\mathfrak p$.
Ainsi, (2) et (3) résultent de propriétés élémentaires de la localisation
(certains détails sont omis). Pour la dernière assertion, supposons que
$\mathcal{F} = \widetilde M$ pour un $B$-module $M$.
Alors $\mathcal{F} = M_{\mathfrak q}$ et
$\pi_*\mathcal{F}_y = M_{\mathfrak p}$. D'après ce qui précède, ces
localisations coïncident. On peut aussi utiliser la partie (1) et
la définition des germes pour voir que $\mathcal{F}_x = \pi_*\mathcal{F}_y$
directement.
\end{proof}






\section{Application à la topologie fppf}
\label{more-morphisms-section-application-fppf}

\noindent
On peut utiliser les techniques de localisation \'etale ci-dessus pour démontrer le
résultat suivant, qui décrit la topologie fppf comme la topologie
« engendrée par » les recouvrements de Zariski et les recouvrements de la
forme $\{f : T \to S\}$, où $f$ est fini localement libre et surjectif.

\begin{lemma}
\label{more-morphisms-lemma-dominate-fppf-etale-locally}
Soit $S$ un schéma. Soit $\{S_i \to S\}_{i \in I}$ un recouvrement fppf.
Alors il existe
\begin{enumerate}
\item un recouvrement \'etale $\{S'_a \to S\}$,
\item des morphismes finis localement libres et surjectifs $V_a \to S'_a$,
\end{enumerate}
tels que le recouvrement fppf $\{V_a \to S\}$ raffine le
recouvrement donné $\{S_i \to S\}$.
\end{lemma}

\begin{proof}
On peut supposer que chaque $S_i \to S$ est localement quasi-fini, voir le
lemme \ref{more-morphisms-lemma-qf-fp-flat-dominates-fppf}.

\medskip\noindent
Fixons un point $s \in S$. Choisissons $i \in I$ et un point
$s_i \in S_i$ au-dessus de $s$. Choisissons un voisinage \'etale élémentaire
$(S', s) \to (S, s)$ tel qu'il existe un ouvert
$$
S_i \times_S S'  \supset V
$$
qui contienne un unique point $v \in V$ au-dessus de $s \in S'$
et tel que $V \to S'$ soit fini, voir le
lemme \ref{more-morphisms-lemma-etale-makes-quasi-finite-finite-at-point}.
Alors $V \to S'$ est fini localement libre, puisqu'il est fini
et que $S_i \times_S S' \to S'$ est plat et localement de présentation finie
comme changement de base du morphisme $S_i \to S$, voir
Morphismes, lemmes \ref{morphisms-lemma-base-change-finite-presentation},
\ref{morphisms-lemma-base-change-flat} et
\ref{morphisms-lemma-finite-flat}.
Ainsi $V \to S'$ est ouvert et, après avoir rétréci $S'$,
on peut supposer que $V \to S'$ est fini localement libre et surjectif.
Comme on peut procéder ainsi pour tout point de $S$, on conclut que
$\{S_i \to S\}$ peut être raffiné par un recouvrement de la forme
$\{V_a \to S\}_{a \in A}$ où chaque $V_a \to S$ se factorise en
$V_a \to S'_a \to S$, avec $S'_a \to S$ \'etale et $V_a \to S'_a$ fini
localement libre et surjectif.
\end{proof}

\begin{lemma}
\label{more-morphisms-lemma-dominate-fppf}
Soit $S$ un schéma. Soit $\{S_i \to S\}_{i \in I}$ un recouvrement fppf.
Alors il existe
\begin{enumerate}
\item un recouvrement ouvert de Zariski $S = \bigcup U_j$,
\item des morphismes finis localement libres et surjectifs $W_j \to U_j$,
\item des recouvrements ouverts de Zariski $W_j = \bigcup_k W_{j, k}$,
\item des morphismes finis localement libres et surjectifs $T_{j, k} \to W_{j, k}$
\end{enumerate}
tels que le recouvrement fppf $\{T_{j, k} \to S\}$ raffine le
recouvrement donné $\{S_i \to S\}$.
\end{lemma}

\begin{proof}
Soit $\{V_a \to S\}_{a \in A}$ le recouvrement fppf obtenu dans le
lemme \ref{more-morphisms-lemma-dominate-fppf-etale-locally}.
Autrement dit, ce recouvrement raffine
$\{S_i \to S\}$
et chaque $V_a \to S$ se factorise en
$V_a \to S'_a \to S$, avec $S'_a \to S$ \'etale et $V_a \to S'_a$
fini localement libre et surjectif.

\medskip\noindent
D'après la
remarque \ref{more-morphisms-remark-topologies},
il existe un recouvrement ouvert de Zariski $S = \bigcup U_j$,
pour chaque $j$ un morphisme fini localement libre et surjectif
$W_j \to U_j$, et pour chaque $j$ un recouvrement ouvert de Zariski
$\{W_{j, k} \to W_j\}$ tel que la famille
$\{W_{j, k} \to S\}$ raffine le recouvrement \'etale
$\{S'_a \to S\}$, c'est-à-dire que, pour chaque paire $j, k$, il existe
un $a(j, k)$ et une factorisation $W_{j, k} \to S'_a \to S$
du morphisme $W_{j, k} \to S$. Posons
$T_{j, k} = W_{j, k} \times_{S'_a} V_a$ ; tout est alors clair.
\end{proof}

\begin{lemma}
\label{more-morphisms-lemma-extend-integral-surjective-morphisms}
Soit $S$ un schéma. Si $U \subset S$ est ouvert et $V \to U$ est un morphisme
intégral surjectif, alors il existe un morphisme intégral surjectif
$\overline{V} \to S$ tel que $\overline{V} \times_S U$ soit
isomorphe à $V$ comme schéma sur $U$.
\end{lemma}

\begin{proof}
Soit $V' \to S$ la normalisation de $S$ dans $U$, voir
Morphismes, section \ref{morphisms-section-normalization-X-in-Y}.
Par construction, $V' \to S$ est intégral.
D'après Morphismes, lemmes \ref{morphisms-lemma-normalization-localization} et
\ref{morphisms-lemma-normalization-in-integral}, on voit que
l'image réciproque de $U$ dans $V'$ est $V$. Soit $Z$ la structure de sous-schéma réduit
induite sur $S \setminus U$. Alors
$\overline{V} = V' \amalg Z$ convient.
\end{proof}

\begin{lemma}
\label{more-morphisms-lemma-extend-finite-surjective-morphisms}
Soit $S$ un schéma quasi-compact et quasi-séparé.
Si $U \subset S$ est un ouvert quasi-compact
et $V \to U$ un morphisme fini surjectif, alors il existe un
morphisme fini surjectif $\overline{V} \to S$ tel que $\overline{V} \times_S U$
soit isomorphe à $V$ comme schéma sur $U$.
\end{lemma}

\begin{proof}
D'après le théorème principal de Zariski
(lemme \ref{more-morphisms-lemma-quasi-finite-separated-pass-through-finite}),
on peut supposer que $V$ est un ouvert quasi-compact d'un schéma $V'$
fini sur $S$. Après avoir remplacé $V'$ par l'image schématique
de $V$, on peut supposer que $V$ est dense dans $V'$.
Il s'ensuit que $V' \times_S U = V$, car $V \to V' \times_S U$
est fermé puisque $V$ est fini sur $U$. Soit $Z$ la structure de sous-schéma réduit
induite sur $S \setminus U$. Alors
$\overline{V} = V' \amalg Z$ convient.
\end{proof}

\begin{lemma}
\label{more-morphisms-lemma-dominate-fppf-integral}
Soit $S$ un schéma. Soit $\{S_i \to S\}_{i \in I}$ un recouvrement fppf.
Alors il existe un morphisme intégral surjectif $S' \to S$ et un
recouvrement ouvert $S' = \bigcup U'_\alpha$ tels que, pour tout $\alpha$, le
morphisme $U'_\alpha \to S$ se factorise par $S_i \to S$ pour un certain $i$.
\end{lemma}

\begin{proof}
Choisissons $S = \bigcup U_j$, $W_j \to U_j$, $W_j = \bigcup W_{j, k}$ et
$T_{j, k} \to W_{j, k}$ comme dans le lemme \ref{more-morphisms-lemma-dominate-fppf}.
D'après le lemme \ref{more-morphisms-lemma-extend-integral-surjective-morphisms},
on peut prolonger $W_j \to U_j$ en un morphisme intégral surjectif
$\overline{W}_j \to S$.
On peut ensuite prolonger $T_{j, k} \to W_{j, k}$ en un morphisme
intégral surjectif $\overline{T}_{j, k} \to \overline{W}_j$.
On définit $\overline{T}_j$ comme le produit de tous les schémas
$\overline{T}_{j, k}$ sur $\overline{W}_j$
(Limites, lemme \ref{limits-lemma-infinite-product}).
Puis on définit $S'$ comme le produit de tous les schémas
$\overline{T}_j$ sur $S$.
Si $x \in S'$, il existe un $j$ tel que l'image de $x$ dans
$S$ appartienne à $U_j$. Il existe donc un $k$ tel que l'image de
$x$ par la projection $S' \to \overline{W}_j$ appartienne à $W_{j, k}$.
Ainsi, par la projection $S' \to \overline{T}_j \to \overline{T}_{j, k}$,
l'image du point $x$ appartient à $T_{j, k}$. Or $T_{j, k} \to S$
se factorise par $S_i$ pour un certain $i$.
Enfin, le morphisme $S' \to S$ est intégral et surjectif
d'après Limites, lemmes \ref{limits-lemma-infinite-product-integral} et
\ref{limits-lemma-infinite-product-surjective}.
\end{proof}

\begin{lemma}
\label{more-morphisms-lemma-dominate-fppf-finite}
Soit $S$ un schéma quasi-compact et quasi-séparé.
Soit $\{S_i \to S\}_{i \in I}$ un recouvrement fppf.
Alors il existe un morphisme fini surjectif $S' \to S$
de présentation finie et un
recouvrement ouvert $S' = \bigcup U'_\alpha$ tel que, pour tout $\alpha$, le
morphisme $U'_\alpha \to S$ se factorise par $S_i \to S$ pour un certain $i$.
\end{lemma}

\begin{proof}
Soit $Y \to X$ le morphisme intégral surjectif obtenu dans le
lemme \ref{more-morphisms-lemma-dominate-fppf-integral}.
Choisissons un recouvrement ouvert affine fini $Y = \bigcup V_j$
tel que $V_j \to X$ se factorise par $S_{i(j)}$.
On peut écrire $Y = \lim Y_\lambda$, avec
$Y_\lambda \to X$ fini et de présentation finie, voir
Limites, lemme \ref{limits-lemma-integral-limit-finite-and-finite-presentation}.
Pour $\lambda$ assez grand, on peut trouver des ouverts affines
$V_{\lambda, j} \subset Y_\lambda$
dont les images réciproques dans $Y$ redonnent les $V_j$, voir
Limites, lemme \ref{limits-lemma-descend-opens}.
Pour $\lambda$ encore plus grand, les morphismes $V_j \to S_{i(j)}$
sur $X$ proviennent de morphismes $V_{\lambda, j} \to S_{i(j)}$ sur
$X$, voir
Limites, proposition
\ref{limits-proposition-characterize-locally-finite-presentation}.
Poser $S' = Y_\lambda$ pour ce $\lambda$ achève la démonstration.
\end{proof}

\begin{lemma}
\label{more-morphisms-lemma-fppf-ph}
Un recouvrement fppf de schémas est un recouvrement ph.
\end{lemma}

\begin{proof}
Soit $\{T_i \to T\}$ un recouvrement fppf de schémas, voir
Topologies, définition \ref{topologies-definition-fppf-covering}.
Observons que $T_i \to T$ est localement de type fini.
Soit $U \subset T$ un ouvert affine.
Il suffit de montrer que $\{T_i \times_T U \to U\}$
peut être raffiné par un recouvrement ph standard, voir
Topologies, définition \ref{topologies-definition-ph-covering}.
Ceci résulte immédiatement du lemme \ref{more-morphisms-lemma-dominate-fppf-finite}
et du fait qu'un morphisme fini est propre
(Morphismes, lemme \ref{morphisms-lemma-finite-proper}).
\end{proof}

\begin{remark}
\label{more-morphisms-remark-change-topologies}
Comme conséquence du lemme \ref{more-morphisms-lemma-fppf-ph}, on obtient un morphisme de comparaison
$$
\epsilon : (\Sch/S)_{ph} \longrightarrow (\Sch/S)_{fppf}
$$
C'est le morphisme de sites donné par le foncteur identité
sur les catégories sous-jacentes (avec des choix convenables de sites
comme dans Topologies, remarque \ref{topologies-remark-choice-sites}).
Le foncteur $\epsilon_*$ est l'identité sur les préfaisceaux sous-jacents
et le foncteur $\epsilon^{-1}$ associe à un faisceau fppf
son faisceau ph associé.
Par composition, on peut en outre comparer la topologie ph
aux topologies syntomique, lisse, \'etale et de Zariski.
\end{remark}




\section{Schémas quasi-projectifs}
\label{more-morphisms-section-quasi-projective}

\noindent
Le terme « schéma quasi-projectif » n'a pas encore été défini.
Une définition possible serait : un schéma qui possède un faisceau inversible
ample. Toutefois, si $X$ est un schéma sur un schéma de base $S$, alors
on dit que {\it $X$ est quasi-projectif sur $S$} si le morphisme
$X \to S$ est quasi-projectif
(Morphismes, définition \ref{morphisms-definition-quasi-projective}).
Comme le morphisme identité de tout schéma est quasi-projectif, on
voit qu'un schéma quasi-projectif sur $S$ ne possède pas nécessairement
un faisceau inversible ample. Pour cette raison, il semble préférable de
ne pas définir le terme « schéma quasi-projectif ».

\begin{lemma}
\label{more-morphisms-lemma-quasi-projective}
Soit $S$ un schéma qui possède un faisceau inversible ample.
Soit $f : X \to S$ un morphisme de schémas. Les conditions suivantes sont
équivalentes :
\begin{enumerate}
\item $X \to S$ est quasi-projectif,
\item $X \to S$ est H-quasi-projectif,
\item il existe une immersion ouverte quasi-compacte $X \to X'$ de schémas
sur $S$ telle que $X' \to S$ soit projectif,
\item $X \to S$ est de type fini et $X$ possède un faisceau inversible
ample, et
\item $X \to S$ est de type fini et il existe un
faisceau inversible $f$-très ample.
\end{enumerate}
\end{lemma}

\begin{proof}
L'implication (2) $\Rightarrow$ (1) est donnée par
Morphismes, lemme \ref{morphisms-lemma-H-quasi-projective-quasi-projective}.
L'implication (1) $\Rightarrow$ (2) est donnée par
Morphismes, lemme
\ref{morphisms-lemma-projective-over-quasi-projective-is-H-projective}.
L'implication (2) $\Rightarrow$ (3) est donnée par
Morphismes, lemme \ref{morphisms-lemma-H-quasi-projective-open-H-projective}.

\medskip\noindent
Supposons que $X \subset X'$ soit comme en (3). En particulier, $X \to S$ est
de type fini. D'après
Morphismes, lemme \ref{morphisms-lemma-H-quasi-projective-open-H-projective},
le morphisme $X \to S$ est H-projectif.
Il existe donc une immersion quasi-compacte $i : X \to \mathbf{P}^n_S$.
Ainsi $\mathcal{L} = i^*\mathcal{O}_{\mathbf{P}^n_S}(1)$
est $f$-très ample. Comme $X \to S$ est quasi-compact, on déduit de
Morphismes, lemme \ref{morphisms-lemma-ample-very-ample},
que $\mathcal{L}$ est $f$-ample. Ainsi $X \to S$ est quasi-projectif
par définition.

\medskip\noindent
L'implication (4) $\Rightarrow$ (2) est donnée par
Morphismes, lemme \ref{morphisms-lemma-quasi-projective-finite-type-over-S}.

\medskip\noindent
L'implication (5) $\Rightarrow$ (4) résulte de
Morphismes, lemme \ref{morphisms-lemma-pullback-ample-tensor-relatively-ample}.

\medskip\noindent
L'équivalence de (1) et (5) résulte de
Morphismes, lemmes \ref{morphisms-lemma-ample-very-ample} et
\ref{morphisms-lemma-finite-type-ample-very-ample}.
\end{proof}

\begin{lemma}
\label{more-morphisms-lemma-category-quasi-projective}
Soit $S$ un schéma qui possède un faisceau inversible ample.
Soit $\text{QP}_S$ la sous-catégorie pleine de la
catégorie des schémas sur $S$ qui satisfont aux conditions
équivalentes du lemme \ref{more-morphisms-lemma-quasi-projective}.
\begin{enumerate}
\item si $S' \to S$ est un morphisme de schémas et que $S'$ possède
un faisceau inversible ample, alors le changement de base définit
un foncteur $\text{QP}_S \to \text{QP}_{S'}$,
\item si $X \in \text{QP}_S$ et $Y \in \text{QP}_X$, alors $Y \in \text{QP}_S$,
\item la catégorie $\text{QP}_S$ est stable par produits fibrés,
\item la catégorie $\text{QP}_S$ est stable par
réunions disjointes finies,
\item si $X \to S$ est projectif, alors $X \in \text{QP}_S$,
\item si $X \to S$ est quasi-affine de type fini, alors
$X$ appartient à $\text{QP}_S$,
\item si $X \to S$ est quasi-fini et séparé, alors
$X \in \text{QP}_S$,
\item si $X \to S$ est une immersion quasi-compacte, alors
$X \in \text{QP}_S$,
\item ajouter d'autres assertions ici.
\end{enumerate}
\end{lemma}

\begin{proof}
L'assertion (1) résulte de Morphismes, lemme
\ref{morphisms-lemma-base-change-quasi-projective}.

\medskip\noindent
L'assertion (2) résulte de la quatrième caractérisation du
lemme \ref{more-morphisms-lemma-quasi-projective}.

\medskip\noindent
Si $X \to S$ et $Y \to S$ sont quasi-projectifs, alors
$X \times_S Y \to Y$ est quasi-projectif d'après
Morphismes, lemme \ref{morphisms-lemma-base-change-quasi-projective}.
Ainsi, (3) résulte de (2).

\medskip\noindent
Si $X = Y \amalg Z$ est une réunion disjointe de schémas
et si $\mathcal{L}$ est un $\mathcal{O}_X$-module inversible
tel que $\mathcal{L}|_Y$ et $\mathcal{L}|_Z$ soient amples, alors
$\mathcal{L}$ est ample (détails omis). Ainsi,
l'assertion (4) résulte de la quatrième caractérisation du
lemme \ref{more-morphisms-lemma-quasi-projective}.

\medskip\noindent
L'assertion (5) résulte de
Morphismes, lemme \ref{morphisms-lemma-projective-quasi-projective}.

\medskip\noindent
L'assertion (6) résulte de
Morphismes, lemme
\ref{morphisms-lemma-quasi-affine-finite-type-quasi-projective}.

\medskip\noindent
L'assertion (7) résulte de l'assertion (6) et du
lemme \ref{more-morphisms-lemma-quasi-finite-separated-quasi-affine}.

\medskip\noindent
L'assertion (8) résulte de l'assertion (7) et de
Morphismes, lemme \ref{morphisms-lemma-immersion-locally-quasi-finite}.
\end{proof}

\noindent
Le lemme suivant n'a pas vraiment sa place dans cette section, mais
il ne semble exister aucun meilleur endroit où le placer.

\begin{lemma}
\label{more-morphisms-lemma-integral-over-quasi-affine}
Soit $X$ un schéma quasi-affine. Soit $f : U \to X$ un morphisme
entier. Alors $U$ est quasi-affine et le diagramme
$$
\xymatrix{
U \ar[r] \ar[d] & \Spec(\Gamma(U, \mathcal{O}_U)) \ar[d] \\
X \ar[r] & \Spec(\Gamma(X, \mathcal{O}_X))
}
$$
est cartésien.
\end{lemma}

\begin{proof}
Le schéma $U$ est quasi-affine parce que les morphismes entiers sont affines,
les morphismes affines sont quasi-affines, un schéma est quasi-affine si et seulement si
le morphisme structural vers $\Spec(\mathbf{Z})$ est quasi-affine, et que
les composés de morphismes quasi-affines sont quasi-affines.
Les deux premières assertions résultent immédiatement de la définition,
et la troisième résulte de
Morphismes, lemme \ref{morphisms-lemma-composition-quasi-affine}.
Posons $U' =
X \times_{\Spec(\Gamma(X, \mathcal{O}_X))} \Spec(\Gamma(U, \mathcal{O}_U))$
et considérons le diagramme prolongé
$$
\xymatrix{
U \ar[r]_j \ar[rd] & U' \ar[d] \ar[r] &
\Spec(\Gamma(U, \mathcal{O}_U)) \ar[d] \\
& X \ar[r] & \Spec(\Gamma(X, \mathcal{O}_X))
}
$$
Le morphisme $j$ est fermé en vertu de
Morphismes, lemme \ref{morphisms-lemma-image-proper-scheme-closed},
ainsi que du fait qu'un morphisme entier est universellement fermé
(Morphismes, lemme \ref{morphisms-lemma-integral-universally-closed}),
et du fait que les flèches verticales du diagramme sont séparées.
D'autre part, $j$ est ouvert parce que les flèches
horizontales du diagramme du lemme sont ouvertes d'après
Propriétés, lemme \ref{properties-lemma-quasi-affine}.
Ainsi, $j$ identifie $U$ à un sous-schéma ouvert et fermé de $U'$.
Si $U \not = U'$, alors $U$ n'est pas dense dans $U'$ et, a fortiori,
n'est pas dense dans le spectre de $\Gamma(U, \mathcal{O}_U)$.
Cependant, l'image schématique de
$U$ dans $\Spec(\Gamma(U, \mathcal{O}_U))$ est $\Spec(\Gamma(U, \mathcal{O}_U))$,
car tout idéal de $\Gamma(U, \mathcal{O}_U)$
définissant un sous-schéma fermé par lequel $U$
se factoriserait devrait être nul.
Par conséquent, $U$ est dense dans $\Spec(\Gamma(U, \mathcal{O}_U))$, par exemple d'après
Morphismes, lemme \ref{morphisms-lemma-quasi-compact-scheme-theoretic-image}.
Ainsi, $U = U'$, ce qui conclut.
\end{proof}









\section{Schémas projectifs}
\label{more-morphisms-section-projective}

\noindent
Cette section est l'analogue de la section \ref{more-morphisms-section-quasi-projective}
pour les morphismes projectifs.

\begin{lemma}
\label{more-morphisms-lemma-projective}
Soit $S$ un schéma qui possède un faisceau inversible ample.
Soit $f : X \to S$ un morphisme de schémas. Les conditions suivantes sont
équivalentes :
\begin{enumerate}
\item $X \to S$ est projectif,
\item $X \to S$ est H-projectif,
\item $X \to S$ est quasi-projectif et propre,
\item $X \to S$ est H-quasi-projectif et propre,
\item $X \to S$ est propre et $X$ possède un faisceau inversible ample,
\item $X \to S$ est propre et il existe un faisceau inversible $f$-ample,
\item $X \to S$ est propre et il existe un faisceau inversible $f$-très ample,
\item il existe une $\mathcal{O}_S$-algèbre graduée quasi-cohérente $\mathcal{A}$
engendrée par $\mathcal{A}_1$ sur $\mathcal{A}_0$, où $\mathcal{A}_1$ est un
$\mathcal{O}_S$-module de type fini, telle que
$X = \underline{\text{Proj}}_S(\mathcal{A})$.
\end{enumerate}
\end{lemma}

\begin{proof}
Observons d'abord que, dans chacun des cas, le morphisme $f$ est propre, voir
Morphismes, lemmes \ref{morphisms-lemma-H-projective} et
\ref{morphisms-lemma-locally-projective-proper}.
Il suffit donc de démontrer l'équivalence des notions dans le
cas où $f$ est un morphisme propre. Nous l'utiliserons dans la suite sans
autre mention.

\medskip\noindent
Les équivalences (1) $\Leftrightarrow$ (3) et
(2) $\Leftrightarrow$ (4) sont données par
Morphismes, lemme \ref{morphisms-lemma-projective-is-quasi-projective-proper}.

\medskip\noindent
L'implication (2) $\Rightarrow$ (1) est donnée par
Morphismes, lemme \ref{morphisms-lemma-H-projective}.

\medskip\noindent
Les implications (1) $\Rightarrow$ (2) et (3) $\Rightarrow$ (4) sont données par
Morphismes, lemme
\ref{morphisms-lemma-projective-over-quasi-projective-is-H-projective}.

\medskip\noindent
L'implication (1) $\Rightarrow$ (7) résulte immédiatement de
Morphismes, définitions \ref{morphisms-definition-projective} et
\ref{morphisms-definition-very-ample}.

\medskip\noindent
Les conditions (3) et (6) sont équivalentes d'après
Morphismes, définition \ref{morphisms-definition-quasi-projective}.

\medskip\noindent
Ainsi, (1) -- (4) et (6) sont équivalentes et entraînent (7). D'après le
lemme \ref{more-morphisms-lemma-quasi-projective},
les conditions (3), (5) et (7) sont équivalentes.
On voit donc que (1) -- (7) sont équivalentes.

\medskip\noindent
D'après Diviseurs, lemme \ref{divisors-lemma-relative-proj-projective},
on voit que (8) entraîne (1). Réciproquement, si (2) est satisfaite, alors
on peut choisir une immersion fermée
$$
i :
X
\longrightarrow
\mathbf{P}^n_S = \underline{\text{Proj}}_S(\mathcal{O}_S[T_0, \ldots, T_n]).
$$
Voir Constructions, lemme \ref{constructions-lemma-projective-space-bundle}
pour cette égalité. D'après
Diviseurs, lemme \ref{divisors-lemma-closed-subscheme-proj},
on voit que $X$ est le Proj relatif d'une algèbre quotient graduée
quasi-cohérente $\mathcal{A}$ de $\mathcal{O}_S[T_0, \ldots, T_n]$.
Alors $\mathcal{A}$ satisfait aux conditions de (8).
\end{proof}

\begin{lemma}
\label{more-morphisms-lemma-category-projective}
Soit $S$ un schéma qui possède un faisceau inversible ample.
Soit $\text{P}_S$ la sous-catégorie pleine de la
catégorie des schémas sur $S$ qui satisfont aux conditions
équivalentes du lemme \ref{more-morphisms-lemma-projective}.
\begin{enumerate}
\item si $S' \to S$ est un morphisme de schémas et que $S'$ possède
un faisceau inversible ample, alors le changement de base définit
un foncteur $\text{P}_S \to \text{P}_{S'}$,
\item si $X \in \text{P}_S$ et $Y \in \text{P}_X$, alors $Y \in \text{P}_S$,
\item la catégorie $\text{P}_S$ est stable par produits fibrés,
\item la catégorie $\text{P}_S$ est stable par
réunions disjointes finies,
\item si $X \to S$ est fini, alors $X$ appartient à $\text{P}_S$,
\item ajouter d'autres assertions ici.
\end{enumerate}
\end{lemma}

\begin{proof}
L'assertion (1) résulte de Morphismes, lemme
\ref{morphisms-lemma-base-change-projective}.

\medskip\noindent
L'assertion (2) résulte de la cinquième caractérisation du
lemme \ref{more-morphisms-lemma-projective} et du fait que les composés
de morphismes propres sont propres
(Morphismes, lemme \ref{morphisms-lemma-composition-proper}).

\medskip\noindent
Si $X \to S$ et $Y \to S$ sont projectifs, alors
$X \times_S Y \to Y$ est projectif d'après
Morphismes, lemme \ref{morphisms-lemma-base-change-projective}.
Ainsi, (3) résulte de (2).

\medskip\noindent
Si $X = Y \amalg Z$ est une réunion disjointe de schémas
et si $\mathcal{L}$ est un $\mathcal{O}_X$-module inversible
tel que $\mathcal{L}|_Y$ et $\mathcal{L}|_Z$ soient amples, alors
$\mathcal{L}$ est ample (détails omis). Ainsi,
l'assertion (4) résulte de la cinquième caractérisation du
lemme \ref{more-morphisms-lemma-projective}.

\medskip\noindent
L'assertion (5) résulte de
Morphismes, lemme \ref{morphisms-lemma-finite-projective}.
\end{proof}

\noindent
Voici un résultat d'un type légèrement différent.

\begin{lemma}
\label{more-morphisms-lemma-ample-in-neighbourhood}
\begin{reference}
\cite[IV Corollaire 9.6.4]{EGA}
\end{reference}
Soit $f : X \to Y$ un morphisme propre de schémas.
Soit $\mathcal{L}$ un $\mathcal{O}_X$-module inversible.
Soit $y \in Y$ un point tel que $\mathcal{L}_y$ soit ample
sur $X_y$. Alors il existe un voisinage ouvert $V \subset Y$
de $y$ tel que $\mathcal{L}|_{f^{-1}(V)}$ soit ample sur $f^{-1}(V)/V$.
\end{lemma}

\begin{proof}
On peut supposer $Y$ affine. On trouve alors un ensemble ordonné filtrant $I$
et un système projectif de morphismes $X_i \to Y_i$ de schémas,
où $Y_i$ est de type fini sur $\mathbf{Z}$, dont les
morphismes de transition $X_i \to X_{i'}$ et $Y_i \to Y_{i'}$ sont affines,
où $X_i \to Y_i$ est propre, tels que $X \to Y = \lim (X_i \to Y_i)$.
Voir Limites, lemme
\ref{limits-lemma-proper-limit-of-proper-finite-presentation-noetherian}.
Après avoir restreint $I$, on peut supposer disposer d'un système compatible de
$\mathcal{O}_{X_i}$-modules inversibles $\mathcal{L}_i$
dont les images réciproques donnent $\mathcal{L}$, voir
Limites, lemme \ref{limits-lemma-descend-invertible-modules}.
Soit $y_i \in Y_i$ l'image de $y$.
Alors $\kappa(y) = \colim \kappa(y_i)$.
Par conséquent, pour un certain $i$, $\mathcal{L}_{i, y_i}$
est ample sur $X_{i, y_i}$ d'après
Limites, lemme \ref{limits-lemma-limit-ample}.
D'après Cohomologie des schémas, lemme \ref{coherent-lemma-ample-in-neighbourhood},
on trouve un voisinage ouvert
$V_i \subset Y_i$ de $y_i$ tel que
la restriction de $\mathcal{L}_i$ à $f_i^{-1}(V_i)$
soit ample relativement à $V_i$.
En prenant pour $V \subset Y$ l'image réciproque de
$V_i$, on achève la démonstration (indications : utiliser
Morphismes, lemme \ref{morphisms-lemma-ample-base-change},
le fait que $X \to Y \times_{Y_i} X_i$ est affine,
et le fait que l'image réciproque d'un
faisceau inversible ample par un morphisme affine est ample d'après
Morphismes, lemme \ref{morphisms-lemma-pullback-ample-tensor-relatively-ample}).
\end{proof}







\section{Proj et Spec}
\label{more-morphisms-section-proj-spec}

\noindent
Dans cette section, nous précisons la relation entre le Proj et le spectre
d'un anneau gradué.

\medskip\noindent
Soit $R$ un anneau. Soit $A$ une $R$-algèbre graduée, voir
Algèbre, section \ref{algebra-section-graded}.
Pour $m \geq 0$, nous notons $A_{\geq m} = \bigoplus_{d \geq m} A_d$.
Considérons l'anneau gradué
$$
B = \bigoplus\nolimits_{d \geq 0} A_{\geq d}
$$
Pour $d' \geq d$ et $a \in A_{d'}$, notons $a^{(d)} \in B$
l'élément de $B_d$ correspondant à $a$.
Notons $\sigma : A \to B$ et $\psi : A \to B$
les deux homomorphismes évidents d'anneaux : si $a \in A_d$, alors
$\sigma(a) = a^{(0)}$ et $\psi(a) = a^{(d)}$.
Alors $\psi$ est un homomorphisme d'anneaux gradués et $\sigma$ munit $B$
d'une structure d'algèbre graduée sur $A$. Il existe aussi un
homomorphisme surjectif d'anneaux gradués $\tau : B \to A$
qui, pour $d' \geq d$ et $a \in A_{d'}$, envoie
$a^{(d)}$ sur $0$ si $d' > d$, et sur $a$ si $d' = d$.

\medskip\noindent
Schémas affines et spectres.
Posons $X = \Spec(A)$. L'idéal non pertinent $A_+$ définit un sous-schéma fermé
$Z = V(A_+) = \Spec(A/A_+) = \Spec(A_0)$. Posons $U = X \setminus Z$.
$$
U \longrightarrow X \longrightarrow Z
$$
Schémas projectifs et Proj. Posons $P = \text{Proj}(A)$. Nous pouvons
et allons considérer $P$ comme un schéma sur $\Spec(A_0) = Z$.
Posons $L = \text{Proj}(B)$. Nous pouvons et allons considérer $L$ comme un schéma
sur $\Spec(B_0) = \Spec(A) = X$ ; observons que l'identification
de $B_0$ avec $A$ est donnée par $\sigma$.
La surjection $\tau$ définit une immersion fermée $0 : P \to L$.
Comme $A \xrightarrow{\sigma} B \to A$ est égal à l'homomorphisme $A \to A_0 \to A$,
on conclut que
$$
\xymatrix{
P \ar[d] \ar[r]_0 & L \ar[d] \\
Z \ar[r] & X
}
$$
est commutatif.

\medskip\noindent
Nous affirmons que $\psi$ définit un morphisme $L \to P$.
Pour le voir, d'après Constructions, lemme \ref{constructions-lemma-morphism-proj},
il suffit de vérifier que $\psi(A_+) \not \subset \mathfrak p$ pour
tout idéal premier homogène
$\mathfrak p \subset B$ tel que $B_+ \not \subset \mathfrak p$.
Choisissons d'abord $g \in B_+$ homogène tel que $g \not \in \mathfrak p$.
On peut alors écrire $g$ comme une somme finie $g = \sum a_i^{(d)}$
avec $a_i \in A_{d_i}$ pour certains $d_i \geq d$.
On conclut qu'il existe $d' \geq d$ et $a \in A_{d'}$
tels que $a^{(d)} \not \in \mathfrak p$.
Alors
$$
(a^{(d)})^{d'} =
(a^{d'})^{(d'd)} =
a^{(d)} (a^{d' - 1})^{(d(d' - 1))} =
\psi(a) (a^{d' - 1})^{(d(d' - 1))}
$$
(la notation laisse quelque peu à désirer) n'appartient pas à $\mathfrak p$.
Ainsi $\psi(a) \not \in \mathfrak p$, ce qui démontre l'assertion.
On peut donc prolonger notre diagramme ci-dessus en un diagramme commutatif
$$
\xymatrix{
P \ar[d] \ar[r]_0 & L \ar[d] \ar[r]_\pi & P \ar[d] \\
Z \ar[r] & X \ar[r] & Z
}
$$
où $X \to Z$ est donné par $A_0 \to A$.
Comme $\tau \circ \psi = \text{id}_A$, on voit que $\pi \circ 0 = \text{id}_P$.

\medskip\noindent
Observons que $\pi$ est un morphisme affine. Cela ressort clairement de la construction
donnée dans Constructions, lemme \ref{constructions-lemma-morphism-proj}. En effet, si
$f \in A_d$ pour un certain $d > 0$, alors, en
posant $g = \psi(f)$, on a $\pi^{-1}(D_+(f)) = D_+(g)$.
Dans ce cas, on a l'égalité suivante entre composantes homogènes :
$$
(B[1/g])_{m'} = \bigoplus\nolimits_{m \geq m'} (A[1/f])_m
$$
Cet isomorphisme est compatible avec les localisations ultérieures.
En prenant $m' = 0$, on voit que $\pi_*\mathcal{O}_L$ est la
somme directe des $\mathcal{O}_P(m)$ pour $m \geq 0$\footnote{On obtient de même
que $\pi_*\mathcal{O}_L(i) = \bigoplus_{m \geq -i} \mathcal{O}_P(m)$.}.
On en conclut que $L$ s'identifie au spectre relatif :
$$
L = \underline{\Spec}_P
\left(
\bigoplus\nolimits_{m \geq 0} \mathcal{O}_P(m)
\right)
$$
En particulier, $L \to P$ est un cône\footnote{Souvent, $L$ est un fibré en droites
sur $P$, voir ci-dessous.}, voir
Constructions, section \ref{constructions-section-cone}.
De plus, il est clair que $0 : P \to L$ est le sommet du cône.

\medskip\noindent
Soient $f \in A_d$ pour un certain $d > 0$ et $g = \psi(f) \in B_d$,
comme au paragraphe précédent. En examinant la structure des homomorphismes d'anneaux
$$
\xymatrix{
A_0 \ar[r] \ar[d] & A \ar[d]^\sigma \ar[r] & A_0 \ar[d] \\
(A[1/f])_0 \ar[r]^-\psi &
(B[1/g])_0 = \bigoplus\nolimits_{m \geq 0} (A[1/f])_m
\ar[r]^-\tau &
(A[1/f])_0
}
$$
quelques calculs\footnote{Les assertions (1) et (2) sont claires.
Pour voir (3), notons que si $a \in A_d$, alors
$\sigma(a) = \sigma(f) \psi(a/f)$. Pour (4), notons que
$b/g^m$ appartient au noyau de $\tau$ si et seulement si
$b \in A_{\geq md}$ s'envoie sur zéro dans $A_{md}$.
Il suffit donc de montrer que, si $m' > md$ et $a \in A_{m'}$,
alors une certaine puissance de $a^{(md)}/g^m$ appartient à l'idéal engendré par
$\sigma(f)$. Prenons $e$ tel que $em' - emd \geq d$. Alors
$$
(a^{(md)}/g^m)^e = (a^e)^{(emd)}/g^{em} =
(fa^e)^{(emd + d)}/g^{em + 1} = \sigma(f) \cdot (a^e)^{(emd + d)}/g^{em + 1}
$$
comme voulu (toutes nos excuses pour cette notation épouvantable). Pour voir (5), raisonnons comme
précédemment et notons que $a^{(md)}/g^m = \sigma(f) \cdot a^{(md + 1)}/g^{m + 1}$
si $d = 1$.}
dans les anneaux gradués montrent que
\begin{enumerate}
\item $\sigma(A_+)(B[1/g])_0 \subset \Ker(\tau : (B[1/g])_0 \to (A[1/f])_0)$,
\item $\sigma(f) \in (B[1/g])_0$ est un élément non diviseur de zéro,
\item $\sigma(f) (B[1/g])_0 = \sigma(A_d) (B[1/g])_0$ en tant qu'idéaux,
\item $\sigma(f) (B[1/g])_0$ et
$\Ker(\tau : (B[1/g])_0 \to (A[1/f])_0)$ ont le même radical,
\item si $d = 1$, alors
$\sigma(f) (B[1/g])_0 = \Ker(\tau : (B[1/g])_0 \to (A[1/f])_0)$.
\end{enumerate}
On voit en particulier que
$$
0(D_+(f)) = V(\sigma(f)) \subset D_+(g) = \Spec((B[1/g])_0)
$$
ensemblistement. En d'autres termes, l'idéal engendré par
$\sigma(A_d)$ définit un diviseur de Cartier effectif sur
$D_+(g)$, dont l'ensemble sous-jacent est
égal à l'image de l'immersion fermée $0 : P \to L$.

\medskip\noindent
Nous affirmons que $L \to X$ est un isomorphisme au-dessus de $U$.
En effet, si $f \in A_d$ pour un certain $d > 0$, alors
$$
\Spec(A_f) \times_X L = \text{Proj}(A_f \otimes_A B) =
\text{Proj}(B_{\sigma(f)})
$$
Pour tout $e$, on a
$(B_{\sigma(f)})_e = A_f \otimes_B B_e = A_f \otimes_A A_{\geq e} = A_f$,
la dernière égalité étant induite par l'injection $A_{\geq e} \subset A$.
Par conséquent, $B_{\sigma(f)} \cong A_f[T]$, où $T$ est de degré $1$.
Cela démontre l'assertion, puisque $\text{Proj}(A_f[T]) \to \Spec(A_f)$
est un isomorphisme. Dorénavant, nous identifions $U$ à l'ouvert correspondant
de $L$.

\medskip\noindent
L'identification effectuée au paragraphe précédent nous permet de
considérer la restriction $\pi|_U : U \to P$. Choisissons
$f \in A_d$ pour un certain $d > 0$ et $g = \psi(f) \in B_d$,
comme nous l'avons fait plusieurs fois ci-dessus. Alors
$$
U \cap \pi^{-1}(D_+(f)) = U \cap D_+(g)
$$
est le complémentaire du lieu des zéros de $\sigma(f) \in (B[1/g])_0$
via l'identification de $D_+(g)$ avec le spectre de $(B[1/g])_0$.
C'est l'assertion (4) ci-dessus. Par conséquent,
$U \cap D_+(g)$ est affine et
$$
\mathcal{O}_L(U \cap D_+(g)) = (B[1/g])_0[1/\sigma(f)] =
\bigoplus\nolimits_{m \in \mathbf{Z}} (A[1/f])_m
$$
où le dernier signe d'égalité est le prolongement naturel de
l'identification $(B[1/g])_0 = \bigoplus_{m \geq 0} (A[1/f])_m$
effectuée ci-dessus. Exactement comme précédemment pour $\pi : L \to P$,
on conclut que $\pi|_U : U \to P$ est affine et que
$$
U = \underline{\Spec}_P
\left(
\bigoplus\nolimits_{m \in \mathbf{Z}} \mathcal{O}_P(m)
\right)
$$
comme schémas sur $P$.

\medskip\noindent
En résumé, nos constructions donnent un diagramme commutatif
\begin{equation}
\label{more-morphisms-equation-proj-and-spec}
\vcenter{
\xymatrix{
\underline{\Spec}_P \left(
\bigoplus\nolimits_{m \in \mathbf{Z}} \mathcal{O}_P(m)
\right)
\ar[r] \ar@{=}[d] &
L =
\underline{\Spec}_P \left(
\bigoplus\nolimits_{m \geq 0} \mathcal{O}_P(m)
\right) \ar[d]^\sigma \ar[r]_-\pi & P \ar[d] \\
U \ar[r] & X \ar[r] & Z
}
}
\end{equation}
de schémas, où $\pi$ est un cône dont la section nulle $0 : P \to L$
a pour image ensembliste l'image réciproque de $Z$ dans $L$.

\medskip\noindent
Soit $W \subset P$ le plus grand ouvert tel que $\mathcal{O}_P(1)|_W$
soit inversible et que les morphismes naturels induisent des isomorphismes
$\mathcal{O}_P(m)|_W \cong \mathcal{O}_P(1)^{\otimes m}|_W$
pour tout $m \in \mathbf{Z}$, c'est-à-dire l'ouvert de
Constructions, lemme \ref{constructions-lemma-where-invertible} pour $d = 1$.
On voit alors que $L|_W = \pi^{-1}(W) \to W$ est un fibré vectoriel
(Constructions, section \ref{constructions-section-vector-bundle})
de rang $1$, à savoir
$$
L|_W = \mathbf{V}(\mathcal{O}_P(1)|_W)
$$
selon la notation grothendieckienne. Cela résulte immédiatement de ce qui précède, qui montre
que $L|_W$ est égal au spectre relatif de l'algèbre symétrique du
$\mathcal{O}_W$-module $\mathcal{O}_P(1)|_W$. Il est alors clair que le
morphisme $0|_W : W \to L|_W$ est la section nulle de ce
fibré vectoriel. En particulier, $0(W)$ est un diviseur de Cartier effectif
sur $L|_W$. De plus, l'ouvert $U|_W = (\pi|_U)^{-1}(W)$
est le complémentaire de la section nulle.

\medskip\noindent
Si $A$ est engendrée par $f_1, \ldots, f_r \in A_1$ sur $A_0$, alors
$(f_1, \ldots, f_r)^m = A_{\geq m}$ pour tout $m \geq 0$ et, par conséquent, notre
$B$ ci-dessus est l'algèbre de Rees de $A_+ = (f_1, \ldots, f_r)$. Ainsi, dans ce
cas, $L \to X$ est le morphisme d'éclatement de centre $Z$ et $W = P$, où $W$ désigne l'ouvert défini au
paragraphe précédent.

\medskip\noindent
Si $P$ est quasi-compact, alors, pour $d$ suffisamment divisible, le
sous-schéma fermé $D \subset L$ défini par $\sigma(A_d)\mathcal{O}_L$
est un diviseur de Cartier effectif, $0 : P \to L$ se factorise par $D$,
et $0(P) = D$ ensemblistement. Cela résulte de
Constructions, lemme \ref{constructions-lemma-proj-quasi-compact},
et des assertions (1), (2), (3) et (4) démontrées ci-dessus. (On peut prendre pour $d$ tout multiple du
ppcm des degrés des éléments trouvés dans le lemme.)

\medskip\noindent
Nous continuons de supposer $P$ quasi-compact.
Soit $\mathcal{F}$ un $\mathcal{O}_P$-module quasi-cohérent.
Posons $\mathcal{F}_U = \pi^*\mathcal{F}|_U$. On a alors
\begin{equation}
\label{more-morphisms-equation-cohomology-torsor}
R\Gamma(U, \mathcal{F}_U) =
\bigoplus\nolimits_{m \in \mathbf{Z}}
R\Gamma(P, \mathcal{F} \otimes_{\mathcal{O}_P} \mathcal{O}_P(m))
\end{equation}
De plus, cette décomposition en somme directe est fonctorielle en $\mathcal{F}$
et la structure de $A$-module induite sur le membre de droite est la même que la
structure de $A$-module du membre de gauche provenant de $U \subset X$.
Pour démontrer la formule, puisque $\pi|_U$ est affine et que
$(\pi|_U)_*\mathcal{O}_U = \bigoplus_{m \in \mathbf{Z}} \mathcal{O}_P(m)$,
on obtient
\begin{align*}
R(\pi|_U)_*\mathcal{F}_U
& =
(\pi|_U)_*\mathcal{F}_U \\
& =
(\pi|_U)_*(\pi|_U)^*\mathcal{F} \\
& =
\mathcal{F} \otimes_{\mathcal{O}_P} 
\bigoplus\nolimits_{m \in \mathbf{Z}} \mathcal{O}_P(m) \\
& =
\bigoplus\nolimits_{m \in \mathbf{Z}}
\mathcal{F} \otimes_{\mathcal{O}_P} \mathcal{O}_P(m)
\end{align*}
Par Leray, on trouve que
$R\Gamma(U, \mathcal{F}_U) =
R\Gamma(P, R(\pi|_U)_*\mathcal{F}_U)$, voir
Cohomologie, lemme \ref{cohomology-lemma-apply-Leray}.
La démonstration est achevée, car la prise de cohomologie commute aux
sommes directes dans ce cas, voir Catégories dérivées des schémas, lemme
\ref{perfect-lemma-quasi-coherence-pushforward-direct-sums}.
C'est ici que l'on utilise que $P$ est quasi-compact ; $P$ est séparé
d'après Constructions, lemme \ref{constructions-lemma-proj-separated}.

\begin{lemma}
\label{more-morphisms-lemma-apply-proj-spec}
Soit $R$ un anneau. Soit $P$ un schéma propre sur $R$ et soit
$\mathcal{L}$ un $\mathcal{O}_P$-module inversible ample.
Posons $A = \bigoplus_{m \geq 0} \Gamma(P, \mathcal{L}^{\otimes m})$.
Alors $P = \text{Proj}(A)$ et le diagramme (\ref{more-morphisms-equation-proj-and-spec})
devient le diagramme
$$
\xymatrix{
\underline{\Spec}_P \left(
\bigoplus\nolimits_{m \in \mathbf{Z}} \mathcal{L}^{\otimes m}
\right)
\ar[r] \ar@{=}[d] &
L =
\underline{\Spec}_P \left(
\bigoplus\nolimits_{m \geq 0} \mathcal{L}^{\otimes m}
\right) \ar[d]^\sigma \ar[r]_-\pi & P \ar[d] \\
U \ar[r] & X \ar[r] & Z
}
$$
ayant les propriétés expliquées ci-dessus.
\end{lemma}

\begin{proof}
On a $P = \text{Proj}(A)$ d'après
Morphismes, lemme \ref{morphisms-lemma-proper-ample-is-proj}.
De plus, d'après Propriétés, lemme \ref{properties-lemma-ample-gcd-is-one},
via cette identification, on a $\mathcal{O}_P(m) = \mathcal{L}^{\otimes m}$
pour tout $m \in \mathbf{Z}$.
\end{proof}









\section{Points fermés dans les fibres}
\label{more-morphisms-section-closed-points-fibres}

\noindent
Une partie du contenu de cette section est tirée de la prépublication
\cite{Osserman-Payne}.

\begin{lemma}
\label{more-morphisms-lemma-locally-principal-vertical}
Soit $f : X \to S$ un morphisme de schémas.
Soit $Z \subset X$ un sous-schéma fermé.
Soit $s \in S$.
Supposons que
\begin{enumerate}
\item $S$ est irréductible, de point générique $\eta$,
\item $X$ est irréductible,
\item $f$ est dominant,
\item $f$ est localement de type fini,
\item $\dim(X_s) \leq \dim(X_\eta)$,
\item $Z$ est localement principal dans $X$, et
\item $Z_\eta = \emptyset$.
\end{enumerate}
Alors la fibre $Z_s$ est (ensemblistement) une réunion de
composantes irréductibles de $X_s$.
\end{lemma}

\begin{proof}
Notons $X_{red}$ la réduction de $X$. Alors $Z \cap X_{red}$ est
un sous-schéma fermé localement principal de $X_{red}$, voir
Diviseurs, lemme \ref{divisors-lemma-pullback-locally-principal}.
On peut donc supposer que $X$ est réduit. Autrement dit, $X$ est intègre, voir
Propriétés, lemme \ref{properties-lemma-characterize-integral}.
Dans ce cas, le morphisme $X \to S$ se factorise par $S_{red}$, voir
Schémas, lemme \ref{schemes-lemma-map-into-reduction}.
Ainsi, on peut remplacer $S$ par $S_{red}$ et supposer que $S$ est également intègre.

\medskip\noindent
Dire que $f$ est dominant signifie que le point générique de $X$
s'envoie sur $\eta$, voir
Morphismes,
lemme \ref{morphisms-lemma-dominant-finite-number-irreducible-components}.
De plus, le schéma $X_\eta$ est un schéma intègre qui est localement de
type fini sur le corps $\kappa(\eta)$. Ainsi,
$d = \dim(X_\eta) \geq 0$ est égal à $\dim_\xi(X_\eta)$ pour
tout point $\xi$ de $X_\eta$, voir
Algèbre, lemmes \ref{algebra-lemma-dimension-spell-it-out} et
\ref{algebra-lemma-dimension-at-a-point-finite-type-over-field}.
Au vu de
Morphismes, lemme \ref{morphisms-lemma-openness-bounded-dimension-fibres}
et de la condition (5), on conclut que $\dim_x(X_s) = d$
pour tout $x \in X_s$.

\medskip\noindent
Dans le cas noethérien, l'assertion se démontre comme suit.
Si le lemme est faux, il existe $x \in Z_s$ qui est un point générique
d'une composante irréductible de $Z_s$, mais pas un point générique
d'une composante irréductible de $X_s$. On voit alors que
$\dim_x(Z_s) \leq d - 1$, car $\dim_x(X_s) = d$ et, dans un voisinage
de $x$ dans $X_s$, le sous-schéma fermé $Z_s$ ne contient aucune des
composantes irréductibles de $X_s$. Ainsi, après avoir remplacé $X$ par un
voisinage ouvert de $x$, on peut supposer que
$\dim_z(Z_{f(z)}) \leq d - 1$ pour tout $z \in Z$, voir
Morphismes, lemme \ref{morphisms-lemma-openness-bounded-dimension-fibres}.
Soit $\xi' \in Z$ un point générique d'une composante irréductible de $Z$
et posons $s' = f(\xi)$. Comme $Z \not = X$ et que ce sous-schéma est localement principal, on voit que
$\dim(\mathcal{O}_{X, \xi}) = 1$, voir
Algèbre, lemme \ref{algebra-lemma-minimal-over-1}
(c'est ici que l'on utilise que $X$ est noethérien).
Soit $\xi \in X$ le point générique de $X$ et
soit $\xi_1$ un point générique d'une composante irréductible quelconque
de $X_{s'}$ qui contient $\xi'$. On voit alors qu'on a
les spécialisations
$$
\xi \leadsto \xi_1 \leadsto \xi'.
$$
Comme $\dim(\mathcal{O}_{X, \xi}) = 1$, l'une des deux spécialisations
doit être une égalité.
Par hypothèse, $s' \not = \eta$ ; la première spécialisation
n'est donc pas une égalité.
Ainsi, $\xi' = \xi_1$ est un point générique d'une composante irréductible de
$X_{s'}$. En appliquant
Morphismes, lemme \ref{morphisms-lemma-openness-bounded-dimension-fibres}
une fois encore, ceci entraîne
$\dim_{\xi'}(Z_{s'}) = \dim_{\xi'}(X_{s'}) \geq \dim(X_\eta) = d$
ce qui donne la contradiction recherchée.

\medskip\noindent
Dans le cas général, on se ramène au cas noethérien comme suit.
Si le lemme est faux, alors il existe un point
$x \in X$ au-dessus de $s$ tel que $x$ soit un point générique d'une
composante irréductible de $Z_s$, mais
ne soit le point générique d'aucune des composantes irréductibles de $X_s$.
Soit $U \subset S$ un voisinage affine de $s$ et soit
$V \subset X$ un voisinage affine de $x$ tel que $f(V) \subset U$.
Écrivons $U = \Spec(A)$ et $V = \Spec(B)$, de sorte que $f|_V$
est donné par un homomorphisme d'anneaux $A \to B$. Soient $\mathfrak q \subset B$,
resp.\ $\mathfrak p \subset A$, les idéaux premiers correspondant à $x$, resp.\ $s$.
Quitte à rétrécir $V$, on peut supposer que $Z \cap V$ est défini par
un élément $g \in B$. Notons $K$ le corps des fractions de $A$.
Ce que nous savons à ce stade est ce qui suit :
\begin{enumerate}
\item $A \subset B$ est une extension de type fini d'anneaux intègres,
\item l'élément $g \otimes 1$ est inversible dans $B \otimes_A K$,
\item $d = \dim(B \otimes_A K) = \dim(B \otimes_A \kappa(\mathfrak p))$,
\item $g \otimes 1$ n'est pas une unité de $B \otimes_A \kappa(\mathfrak p)$, et
\item $g \otimes 1$ n'appartient à aucun des idéaux premiers minimaux de
$B \otimes_A \kappa(\mathfrak p)$.
\end{enumerate}
Nous cherchons une contradiction.

\medskip\noindent
Choisissons des éléments $x_1, \ldots, x_n \in B$ qui engendrent $B$ sur $A$.
Pour une $\mathbf{Z}$-algèbre de type fini $A_0 \subset A$,
soit $B_0 \subset B$ la sous-$A_0$-algèbre engendrée par
$x_1, \ldots, x_n$, notons $K_0$ le corps des fractions de $A_0$, et posons
$\mathfrak p_0 = A_0 \cap \mathfrak p$.
Nous affirmons que, lorsque $A_0$ est assez grande, (1) -- (5) sont encore satisfaites pour
le système $(A_0 \subset B_0, g, \mathfrak p_0)$.

\medskip\noindent
Démontrons successivement chacune des conditions. La condition (1) est satisfaite par construction.
Pour la partie (2), écrivons $(g \otimes 1) h = 1$ pour un certain
$h \otimes 1/a \in B \otimes_A K$. Écrivons
$g = \sum a_I x^I$, $h = \sum a'_I x^I$ (notation multi-indicielle)
pour certains coefficients $a_I, a'_I \in A$. Dès que $A_0$ contient
$a$ et les $a_I, a'_I$, alors (2) est vraie parce que
$B_0 \otimes_{A_0} K_0 \subset B \otimes_A K$ (comme localisations de
l'homomorphisme injectif $B_0 \to B$).
Pour obtenir (3), considérons la suite exacte
$$
0 \to I \to A[X_1, \ldots, X_n] \to B \to 0
$$
qui définit $I$, où le second homomorphisme envoie $X_i$ sur $x_i$. Comme $\otimes$
est exact à droite, on voit que $I \otimes_A K$ et, respectivement,
$I \otimes_A \kappa(\mathfrak p)$ sont les noyaux des surjections
$K[X_1, \ldots, X_n] \to B \otimes_A K$ et, respectivement,
$\kappa(\mathfrak p)[X_1, \ldots, X_n] \to B \otimes_A \kappa(\mathfrak p)$.
Comme un anneau de polynômes sur un corps est noethérien,
il existe un nombre fini d'éléments $h_j \in I$, $j = 1, \ldots, m$,
qui engendrent $I \otimes_A K$ et $I \otimes_A \kappa(\mathfrak p)$.
Écrivons $h_j = \sum a_{j, I}X^I$. Dès que
$A_0$ contient tous les $a_{j, I}$, on se trouve dans la situation où
$$
B_0 \otimes_{A_0} K_0 \otimes_{K_0} K = B \otimes_A K
\quad\text{et}\quad
B_0 \otimes_{A_0} \kappa(\mathfrak p_0)
\otimes_{\kappa(\mathfrak p_0)} \kappa(\mathfrak p)
=
B \otimes_A \kappa(\mathfrak p).
$$
En appliquant au choix
Morphismes, lemme \ref{morphisms-lemma-dimension-fibre-after-base-change}
ou
Algèbre, lemme \ref{algebra-lemma-dimension-preserved-field-extension},
on voit que les égalités de dimension de (3) sont satisfaites.
La partie (4) est immédiate. Comme
$B_0 \otimes_{A_0} \kappa(\mathfrak p_0) \subset
B \otimes_A \kappa(\mathfrak p)$, tout idéal premier minimal de
$B_0 \otimes_{A_0} \kappa(\mathfrak p_0)$ est la contraction d'un idéal
premier minimal de $B \otimes_A \kappa(\mathfrak p)$ d'après
Algèbre, lemme \ref{algebra-lemma-image-dense-generic-points}.
Cela montre que (5) est satisfaite.
On réduit ainsi le problème au cas noethérien, qui a été
traité ci-dessus.
\end{proof}

\noindent
Voici une application algébrique du lemme précédent.
La quatrième hypothèse du lemme est satisfaite si $A \to B$ est plat, voir
le lemme \ref{more-morphisms-lemma-equality-dimensions}.

\begin{lemma}
\label{more-morphisms-lemma-horizontal}
Soit $A \to B$ un homomorphisme local d'anneaux locaux, et soit
$g \in \mathfrak m_B$. Supposons que
\begin{enumerate}
\item $A$ et $B$ sont des anneaux intègres et $A \subset B$,
\item $B$ est essentiellement de type fini sur $A$,
\item $g$ n'appartient à aucun idéal premier minimal contenant $\mathfrak m_AB$, et
\item $\dim(B/\mathfrak m_AB) +
\text{trdeg}_{\kappa(\mathfrak m_A)}(\kappa(\mathfrak m_B)) =
\text{trdeg}_A(B)$.
\end{enumerate}
Alors $A \subset B/gB$, c'est-à-dire que le point générique de $\Spec(A)$
appartient à l'image du morphisme $\Spec(B/gB) \to \Spec(A)$.
\end{lemma}

\begin{proof}
Notons que les deux assertions sont équivalentes d'après
Algèbre, lemme \ref{algebra-lemma-image-dense-generic-points}.
Pour commencer la démonstration, soit $C$ une $A$-algèbre de type fini
et soit $\mathfrak q$ un idéal premier de $C$ tel que $B = C_{\mathfrak q}$.
On peut bien sûr supposer que $C$ est un anneau intègre et que $g \in C$.
Après avoir remplacé $C$ par un localisé, on voit que
$\dim(C/\mathfrak m_AC) = \dim(B/\mathfrak m_AB) +
\text{trdeg}_{\kappa(\mathfrak m_A)}(\kappa(\mathfrak m_B))$, voir
Morphismes, lemme \ref{morphisms-lemma-dimension-fibre-at-a-point}.
En posant $K$ égal au corps des fractions de $A$,
on voit par la même référence que
$\dim(C \otimes_A K) = \text{trdeg}_A(B)$. Ainsi, l'hypothèse
(4) signifie que les fibres générique et fermée du morphisme
$\Spec(C) \to \Spec(A)$ ont même dimension.

\medskip\noindent
Supposons le lemme faux. Alors $(B/gB) \otimes_A K = 0$.
Cela signifie que $g \otimes 1$ est inversible dans $B \otimes_A K
= C_{\mathfrak q} \otimes_A K$. Comme $C_{\mathfrak q}$ est une limite
de localisés principaux, on conclut que $g \otimes 1$
est inversible dans $C_h \otimes_A K$ pour un certain
$h \in C$, $h \not \in \mathfrak q$. Ainsi, après avoir remplacé $C$
par $C_h$, on peut supposer que $(C/gC) \otimes_A K = 0$.
On remplace encore une fois $C$ afin que les idéaux premiers minimaux
de $C/\mathfrak m_AC$ correspondent bijectivement aux idéaux premiers minimaux
de $B/\mathfrak m_AB$. À ce stade, on applique
le lemme \ref{more-morphisms-lemma-locally-principal-vertical}
à $X = \Spec(C) \to \Spec(A) = S$ et au sous-schéma localement fermé
$Z = \Spec(C/gC)$. Comme $Z_K = \emptyset$, on voit que
$Z \otimes \kappa(\mathfrak m_A)$ doit contenir une composante irréductible
de
$X \otimes \kappa(\mathfrak m_A) = \Spec(C/\mathfrak m_AC)$.
Mais ceci contredit l'hypothèse que $g$ n'appartient
à aucun idéal premier minimal contenant $\mathfrak m_AB$. Le lemme en résulte.
\end{proof}

\begin{lemma}
\label{more-morphisms-lemma-equality-dimensions}
Soit $A \to B$ un homomorphisme local d'anneaux locaux. Supposons que
\begin{enumerate}
\item $A$ et $B$ sont des anneaux intègres et $A \subset B$,
\item $B$ est essentiellement de type fini sur $A$, et
\item $B$ est plat sur $A$.
\end{enumerate}
Alors on a
$$
\dim(B/\mathfrak m_AB) +
\text{trdeg}_{\kappa(\mathfrak m_A)}(\kappa(\mathfrak m_B)) =
\text{trdeg}_A(B).
$$
\end{lemma}

\begin{proof}
Soit $C$ une $A$-algèbre de type fini et soit $\mathfrak q$ un idéal premier de $C$
tel que $B = C_{\mathfrak q}$. On peut supposer que $C$ est un anneau intègre.
On a
$\dim_{\mathfrak q}(C/\mathfrak m_AC) = \dim(B/\mathfrak m_AB) +
\text{trdeg}_{\kappa(\mathfrak m_A)}(\kappa(\mathfrak m_B))$, voir
Morphismes, lemme \ref{morphisms-lemma-dimension-fibre-at-a-point}.
En posant $K$ égal au corps des fractions de $A$,
on voit par la même référence que
$\dim(C \otimes_A K) = \text{trdeg}_A(B)$.
Ainsi, il s'agit en fait de démontrer que
$\dim_{\mathfrak q}(C/\mathfrak m_AC) = \dim(C \otimes_A K)$.
Choisissons un anneau de valuation $A'$ de $K$ dominant $A$, voir
Algèbre, lemme \ref{algebra-lemma-dominate}.
Posons $C' = C \otimes_A A'$.
Choisissons un idéal premier $\mathfrak q'$ de $C'$ au-dessus de $\mathfrak q$ ; un tel
idéal premier existe car
$$
C'/\mathfrak m_{A'}C' =
C/\mathfrak m_AC \otimes_{\kappa(\mathfrak m_A)} \kappa(\mathfrak m_{A'})
$$
ce qui démontre que $C/\mathfrak m_AC \to C'/\mathfrak m_{A'}C'$ est fidèlement
plat. Ceci démontre également que
$\dim_{\mathfrak q}(C/\mathfrak m_AC) =
\dim_{\mathfrak q'}(C'/\mathfrak m_{A'}C')$, voir
Algèbre,
lemme \ref{algebra-lemma-dimension-at-a-point-preserved-field-extension}.
Notons que $B' = C'_{\mathfrak q'}$ est un localisé de $B \otimes_A A'$.
Ainsi, $B'$ est plat sur $A'$. La fibre générique $B' \otimes_{A'} K$
est un localisé de $B \otimes_A K$. Ainsi, $B'$ est un anneau intègre.
Si l'on démontre le lemme pour $A' \subset B'$, on obtient alors l'égalité
$\dim_{\mathfrak q'}(C'/\mathfrak m_{A'}C') = \dim(C' \otimes_{A'} K)$
ce qui implique l'égalité recherchée
$\dim_{\mathfrak q}(C/\mathfrak m_AC) = \dim(C \otimes_A K)$
d'après ce qui précède. Cela ramène le
lemme au cas où $A$ est un anneau de valuation.

\medskip\noindent
Soit $A \subset B$ comme dans le lemme, avec $A$ un anneau de valuation.
Comme précédemment, écrivons $B = C_{\mathfrak q}$ pour un anneau intègre $C$ de type
fini sur $A$. D'après
Algèbre,
lemme \ref{algebra-lemma-finite-type-domain-over-valuation-ring-dim-fibres},
on obtient $\dim(C/\mathfrak m_AC) = \dim(C \otimes_A K)$, ce qui achève la démonstration.
\end{proof}

\begin{lemma}
\label{more-morphisms-lemma-closed-point-nearby-fibre}
Soit $f : X \to S$ un morphisme de schémas.
Soit $x \leadsto x'$ une spécialisation de points de $X$.
Posons $s = f(x)$ et $s' = f(x')$.
Supposons que
\begin{enumerate}
\item $x'$ est un point fermé de $X_{s'}$, et
\item $f$ est localement de type fini.
\end{enumerate}
Alors l'ensemble
$$
\{x_1 \in X
\text{ tel que }
f(x_1) = s
\text{ et }
x_1\text{ est fermé dans }X_s
\text{ et }
x \leadsto x_1 \leadsto x'
\}
$$
est dense dans l'adhérence de $x$ dans $X_s$.
\end{lemma}

\begin{proof}
Appliquons
Schémas, lemme \ref{schemes-lemma-points-specialize},
à la spécialisation $x \leadsto x'$.
Cela fournit un morphisme $\varphi : \Spec(B) \to X$,
où $B$ est un anneau de valuation tel que $\varphi$ envoie le
point générique sur $x$ et le point fermé sur $x'$. On peut aussi
supposer que $\kappa(x)$ est le corps des fractions de $B$.
Soit $A = B \cap \kappa(s)$. Notons que c'est un anneau de valuation (voir
Algèbre, lemme \ref{algebra-lemma-valuation-ring-cap-field})
qui domine l'image de $\mathcal{O}_{S, s'} \to \kappa(s)$.
Considérons le diagramme commutatif
$$
\xymatrix{
\Spec(B) \ar[rd] \ar[r] &
X_A \ar[d] \ar[r] & X \ar[d] \\
& \Spec(A) \ar[r] & S
}
$$
Le point générique (resp.\ fermé) de $B$ s'envoie sur un point $x_A$
(resp.\ $x'_A$) de $X_A$ situé au-dessus du point générique (resp.\ fermé)
de $\Spec(A)$. Notons que $x'_A$ est un point fermé
de la fibre spéciale de $X_A$ d'après
Morphismes,
lemme \ref{morphisms-lemma-base-change-closed-point-fibre-locally-finite-type}.
Notons que la fibre générique de $X_A \to \Spec(A)$ est isomorphe
à $X_s$. Ainsi, nous sommes ramenés au cas où $S$ est
le spectre d'un anneau de valuation, $s = \eta \in S$ est le point générique, et
$s' \in S$ est le point fermé.

\medskip\noindent
Nous allons démontrer le lemme par récurrence sur $\dim_x(X_\eta)$.
Si $\dim_x(X_\eta) = 0$, alors il n'existe aucun autre point de $X_\eta$
se spécialisant en $x$, et $x$ est fermé dans sa fibre, voir
Morphismes, lemme \ref{morphisms-lemma-quasi-finite-at-point-characterize},
et le résultat en découle. Supposons que $\dim_x(X_\eta) > 0$.

\medskip\noindent
Soit $X' \subset X$ le sous-schéma fermé irréductible
$\overline{\{x\}}$ de $X$, muni de sa structure de schéma induite réduite, voir
Schémas, définition \ref{schemes-definition-reduced-induced-scheme}.
Pour démontrer le lemme, on peut remplacer $X$ par $X'$, puisque ceci ne fait que diminuer
$\dim_x(X_\eta)$. On peut donc aussi supposer que $X$ est un schéma intègre
et que $x$ est son point générique. En outre, on peut remplacer $X$ par un
voisinage affine de $x'$. Ainsi, on a $X = \Spec(B)$, où
$A \subset B$ est une extension de type fini d'anneaux intègres. Notons que, dans
ce cas, $\dim_x(X_\eta) = \dim(X_\eta) = \dim(X_{s'})$, et qu'en fait
$X_{s'}$ est équidimensionnel, voir
Algèbre,
lemme \ref{algebra-lemma-finite-type-domain-over-valuation-ring-dim-fibres}.

\medskip\noindent
Soit $W \subset X_\eta$ un fermé propre (c'est le
sous-ensemble que l'on veut « éviter »). Comme $X_s$ est de type fini sur un corps,
on voit que $W$ possède un nombre fini de composantes irréductibles
$W = W_1 \cup \ldots \cup W_n$. Soient
$\mathfrak q_j \subset B$, $j = 1, \ldots, r$
les idéaux premiers correspondants. Soit $\mathfrak q \subset B$
l'idéal maximal correspondant au point $x'$.
Soient $\mathfrak p_1, \ldots, \mathfrak p_s \subset B$ les
idéaux premiers minimaux contenant $\mathfrak m_AB$. Il y en a un nombre fini,
puisqu'ils correspondent aux composantes irréductibles du
schéma noethérien $X_{s'}$. De plus, chacune de ces composantes irréductibles
est de dimension $> 0$ (voir ci-dessus) ; on voit donc que
$\mathfrak p_i \not = \mathfrak q$ pour tout $i$.
Choisissons maintenant un élément $g \in \mathfrak q$ tel que
$g \not \in \mathfrak q_j$ pour tout $j$ et $g \not \in \mathfrak p_i$
pour tout $i$, voir
Algèbre, lemme \ref{algebra-lemma-silly}.
Notons $Z \subset X$ le sous-schéma fermé localement principal défini par $g$.
Écrivons $Z_\eta = Z_{1, \eta} \cup \ldots \cup Z_{n, \eta}$, $n \geq 0$,
pour la décomposition de la fibre générique de $Z$ en composantes irréductibles
(en nombre fini puisque la fibre générique est noethérienne).
Notons $Z_i \subset X$ l'adhérence de $Z_{i, \eta}$.
Après avoir remplacé $X$ par un voisinage affine plus petit,
on peut supposer que $x' \in Z_i$ pour tout $i = 1, \ldots, n$.
Par construction, $Z \cap X_{s'}$ ne contient aucune composante irréductible
de $X_{s'}$. Ainsi, d'après
le lemme \ref{more-morphisms-lemma-locally-principal-vertical},
on conclut que $Z_\eta \not = \emptyset$ ! Autrement dit,
$n \geq 1$. Soit $x_1 \in Z_1$ le point générique ; on voit
que $x_1 \leadsto x'$ et $f(x_1) = \eta$.
De plus, par construction, $Z_{1, \eta} \cap W_j \subset W_j$
est un fermé propre. Ainsi, toute composante irréductible de
$Z_{1, \eta} \cap W_j$ est de codimension $\geq 2$ dans $X_\eta$,
tandis que $\text{codim}(Z_{1, \eta}, X_\eta) = 1$ d'après
Algèbre, lemme \ref{algebra-lemma-minimal-over-1}.
Ainsi, $W \cap Z_{1, \eta}$ est un fermé propre.
À ce stade, on voit que l'hypothèse de récurrence s'applique à
$Z_1 \to S$ et à la spécialisation $x_1 \leadsto x'$.
Ceci fournit un point fermé $x_2$ de $Z_{1, \eta}$ qui n'appartient pas
à $W$ et qui se spécialise en $x'$. Ainsi, on obtient
$x \leadsto x_2 \leadsto x'$ ; le point $x_2$ est fermé dans $X_\eta$,
et $x_2 \not \in W$, comme souhaité.
\end{proof}

\begin{remark}
\label{more-morphisms-remark-full-specialization-sequence}
La démonstration du
lemme \ref{more-morphisms-lemma-closed-point-nearby-fibre}
montre en fait qu'il existe une suite de spécialisations
$$
x \leadsto x_1 \leadsto x_2 \leadsto \ldots \leadsto x_d \leadsto x'
$$
où tous les $x_i$ appartiennent à la fibre $X_s$, chaque spécialisation est
immédiate, et $x_d$ est un point fermé de $X_s$. L'entier
$d = \text{trdeg}_{\kappa(s)}(\kappa(x)) = \dim(\overline{\{x\}})$
où l'adhérence est prise dans $X_s$. De plus, les points
$x_i$ peuvent être choisis de manière à éviter tout fermé de $X_s$ qui
ne contient pas le point $x$.
\end{remark}

\noindent
Exemples, section \ref{examples-section-topology-finite-type},
montre que le lemme suivant est faux si $A$ n'est pas supposé
noethérien.

\begin{lemma}
\label{more-morphisms-lemma-quasi-finite-quasi-section-meeting-nearby-open}
Soit $\varphi : A \to B$ un homomorphisme local d'anneaux locaux.
Soit $V \subset \Spec(B)$ un sous-schéma ouvert
qui contient au moins un idéal premier ne se trouvant pas au-dessus de $\mathfrak m_A$.
Supposons que $A$ soit noethérien, que $\varphi$ soit essentiellement de type fini, et que
l'extension $A/\mathfrak m_A \subset B/\mathfrak m_B$ soit finie.
Alors il existe un $\mathfrak q \in V$,
$\mathfrak m_A \not = \mathfrak q \cap A$, tel que
$A \to B/\mathfrak q$ soit le localisé d'un homomorphisme d'anneaux quasi-fini.
\end{lemma}

\begin{proof}
Comme $A$ est noethérien et que $A \to B$ est essentiellement de type fini,
on sait que $B$ est également noethérien. D'après
Propriétés, lemme \ref{properties-lemma-complement-closed-point-Jacobson},
l'espace topologique $\Spec(B) \setminus \{\mathfrak m_B\}$
est jacobsonien. On peut donc choisir un point fermé $\mathfrak q$
qui appartient à l'ouvert non vide
$$
V \setminus \{\mathfrak q \subset B \mid \mathfrak m_A = \mathfrak q \cap A\}.
$$
(Non vide par hypothèse, ouvert parce que $\{\mathfrak m_A\}$ est un fermé
de $\Spec(A)$.)
Alors $\Spec(B/\mathfrak q)$ possède deux points, à savoir $\mathfrak m_B$
et $\mathfrak q$, et $\mathfrak q$ ne se trouve pas au-dessus de $\mathfrak m_A$.
Écrivons $B/\mathfrak q = C_{\mathfrak m}$ pour une $A$-algèbre de type fini
$C$ et un idéal premier $\mathfrak m$. Alors $A \to C$ est quasi-fini en
$\mathfrak m$ d'après
Algèbre, lemme \ref{algebra-lemma-isolated-point-fibre} (2).
Par conséquent, d'après
Algèbre, lemme \ref{algebra-lemma-quasi-finite-open},
on voit qu'après avoir remplacé $C$ par un localisé principal, l'homomorphisme
$A \to C$ est quasi-fini.
\end{proof}

\begin{lemma}
\label{more-morphisms-lemma-quasi-finite-quasi-section-meeting-nearby-open-X}
Soit $f : X \to S$ un morphisme de schémas.
Soit $x \in X$ d'image $s \in S$.
Soit $U \subset X$ un sous-schéma ouvert.
Supposons que $f$ soit localement de type fini, que $S$ soit localement noethérien, et que $x$ soit un point fermé
de $X_s$, et qu'il existe un point $x' \in U$ tel que
$x' \leadsto x$ et $f(x') \not = s$. Alors il existe un sous-schéma fermé
$Z \subset X$ tel que (a) $x \in Z$, (b) $f|_Z : Z \to S$ soit
quasi-fini en $x$, et (c) il existe un $z \in Z$, $z \in U$,
tel que $z \leadsto x$ et $f(z) \not = s$.
\end{lemma}

\begin{proof}
C'est une reformulation du
lemme \ref{more-morphisms-lemma-quasi-finite-quasi-section-meeting-nearby-open}.
En effet, posons $A = \mathcal{O}_{S, s}$ et $B = \mathcal{O}_{X, x}$.
Notons $V \subset \Spec(B)$ l'image réciproque de $U$.
L'homomorphisme d'anneaux $f^\sharp : A \to B$ est essentiellement de type fini.
Par hypothèse, il existe au moins un point de $V$ qui ne
s'envoie pas sur le point fermé de $\Spec(A)$. Toutes les hypothèses du
lemme \ref{more-morphisms-lemma-quasi-finite-quasi-section-meeting-nearby-open}
sont donc satisfaites, et l'on obtient un idéal premier $\mathfrak q \subset B$ qui ne
se trouve pas au-dessus de $\mathfrak m_A$ et tel que $A \to B/\mathfrak q$ soit
le localisé d'un homomorphisme d'anneaux quasi-fini. Soit $z \in X$
l'image du point $\mathfrak q$ par le morphisme canonique
$\Spec(B) \to X$. Posons $Z = \overline{\{z\}}$
muni de la structure de schéma induite réduite. Comme $z \leadsto x$,
on voit que $x \in Z$ et $\mathcal{O}_{Z, x} = B/\mathfrak q$.
Par construction, $Z \to S$ est quasi-fini en $x$.
\end{proof}

\begin{remark}
\label{more-morphisms-remark-alternative-closed-point-nearby-fibre}
On peut utiliser le
lemme \ref{more-morphisms-lemma-quasi-finite-quasi-section-meeting-nearby-open}
ou sa variante, le
lemme \ref{more-morphisms-lemma-quasi-finite-quasi-section-meeting-nearby-open-X},
pour donner une autre démonstration du
lemme \ref{more-morphisms-lemma-closed-point-nearby-fibre}
lorsque $S$ est localement noethérien.
En voici une esquisse.
On remplace d'abord $S$ par
le spectre de l'anneau local en $s'$. Puis on peut raisonner par récurrence
sur $\dim(S)$. Le cas $\dim(S) = 0$ est trivial, car alors $s' = s$.
Remplaçons $X$ par $\overline{\{x\}}$, muni de sa structure de schéma induite réduite.
Appliquons le
lemme \ref{more-morphisms-lemma-quasi-finite-quasi-section-meeting-nearby-open-X}
à $X \to S$ et à $x' \mapsto s'$, ainsi qu'à tout
ouvert non vide $U \subset X$ contenant $x$. On obtient un sous-schéma fermé
$x' \in Z \subset X$ et un point $z \in Z$
tels que $Z \to S$ soit quasi-fini en $x'$ et que $f(z) \not = s'$.
Alors $z$ est un point fermé de $X_{f(z)}$, et $z \leadsto x'$.
Comme $f(z) \not = s'$, on a $\dim(\mathcal{O}_{S, f(z)}) < \dim(S)$.
Puisque $x$ est le point générique de $X$, on a $x \leadsto z$, donc
$s = f(x) \leadsto f(z)$.
Appliquons l'hypothèse de récurrence à $s \leadsto f(z)$ et à $z \mapsto f(z)$
pour conclure.
\end{remark}

\begin{lemma}
\label{more-morphisms-lemma-change-hypotheses}
Supposons que $f : X \to S$ soit localement de type fini et que $S$ soit localement noethérien,
que $x \in X$ soit un point fermé de sa fibre $X_s$ et que $U \subset X$ soit un sous-schéma
ouvert tel que $U \cap X_s = \emptyset$ et $x \in \overline{U}$ ; alors
les conclusions du
lemme \ref{more-morphisms-lemma-quasi-finite-quasi-section-meeting-nearby-open-X}
sont valables.
\end{lemma}

\begin{proof}
En effet, on peut se ramener au lemme cité comme suit : d'abord, on
remplace $X$ et $S$ par des voisinages affines de $x$ et $s$. Alors $X$ est
noethérien ; en particulier, $U$ est quasi-compact (voir
Morphismes, lemme \ref{morphisms-lemma-finite-type-noetherian}
et
Topologie, lemmes \ref{topology-lemma-Noetherian} et
\ref{topology-lemma-Noetherian-quasi-compact}).
Il existe donc une spécialisation $x' \leadsto x$ avec $x' \in U$ (voir
Morphismes, lemme \ref{morphisms-lemma-reach-points-scheme-theoretic-image}).
Notons que $f(x') \not = s$. Ainsi, toutes les hypothèses du lemme
sont satisfaites, ce qui conclut.
\end{proof}



\section{Factorisation de Stein}
\label{more-morphisms-section-stein-factorization}

\noindent
La factorisation de Stein affirme qu'un morphisme propre $f : X \to S$
tel que $f_*\mathcal{O}_X = \mathcal{O}_S$ a des fibres connexes.

\begin{lemma}
\label{more-morphisms-lemma-stein-universally-closed}
Soit $S$ un schéma. Soit $f : X \to S$ un morphisme universellement fermé et
quasi-séparé. Il existe une factorisation
$$
\xymatrix{
X \ar[rr]_{f'} \ar[rd]_f & & S' \ar[dl]^\pi \\
& S &
}
$$
ayant les propriétés suivantes :
\begin{enumerate}
\item le morphisme $f'$ est universellement fermé, quasi-compact, quasi-séparé,
et surjectif,
\item le morphisme $\pi : S' \to S$ est intégral,
\item on a $f'_*\mathcal{O}_X = \mathcal{O}_{S'}$,
\item on a $S' = \underline{\Spec}_S(f_*\mathcal{O}_X)$, et
\item $S'$ est la normalisation de $S$ dans $X$, voir
Morphismes, définition \ref{morphisms-definition-normalization-X-in-Y}.
\end{enumerate}
La formation de la factorisation $f = \pi \circ f'$ commute
au changement de base plat.
\end{lemma}

\begin{proof}
D'après Morphismes, lemme \ref{morphisms-lemma-universally-closed-quasi-compact},
le morphisme $f$ est quasi-compact. Par conséquent, la normalisation $S'$ de $S$ dans
$X$ est définie (Morphismes, définition
\ref{morphisms-definition-normalization-X-in-Y})
et l'on a la factorisation $X \to S' \to S$. D'après
Morphismes, lemme \ref{morphisms-lemma-normalization-in-universally-closed},
on a (2), (4) et (5). Le morphisme $f'$ est universellement fermé d'après
Morphismes, lemme \ref{morphisms-lemma-image-proper-scheme-closed}.
Il est quasi-compact d'après
Schémas, lemme \ref{schemes-lemma-quasi-compact-permanence},
et quasi-séparé d'après
Schémas, lemme \ref{schemes-lemma-compose-after-separated}.

\medskip\noindent
Pour démontrer les assertions restantes, on peut supposer que le schéma de base $S$ est affine,
disons $S = \Spec(R)$. Alors $S' = \Spec(A)$ avec
$A = \Gamma(X, \mathcal{O}_X)$ une $R$-algèbre entière.
Il est donc clair que $f'_*\mathcal{O}_X$
est $\mathcal{O}_{S'}$ (car $f'_*\mathcal{O}_X$ est quasi-cohérent,
d'après
Schémas, lemme
\ref{schemes-lemma-push-forward-quasi-coherent},
et donc égal à $\widetilde{A}$). Cela démontre (3).

\medskip\noindent
Montrons que $f'$ est surjectif. Comme $f'$ est universellement fermé (voir ci-dessus),
l'image de $f'$ est un sous-ensemble fermé
$V(I) \subset S' = \Spec(A)$. Prenons $h \in I$. Alors
$h|_X = f^\sharp(h)$ est une section globale du faisceau structural de
$X$ qui s'annule en tout point. Comme $X$ est quasi-compact, cela signifie
que $h|_X$ est une section nilpotente, c'est-à-dire que $h^n|X = 0$ pour un certain $n > 0$.
Mais $A = \Gamma(X, \mathcal{O}_X)$, donc $h^n = 0$.
Tout élément de $I$ étant nilpotent, on conclut que $V(I) = S'$, comme voulu.

\medskip\noindent
D'après Cohomologie des schémas, lemme \ref{coherent-lemma-flat-base-change-cohomology},
on voit que la formation de $f_*\mathcal{O}_X$ commute au changement de base plat.
La formation du spectre relatif commute à tout changement de base d'après
Constructions, lemme \ref{constructions-lemma-spec-properties}.
Ainsi, la formation de la factorisation commute au changement de base plat.
\end{proof}

\begin{lemma}
\label{more-morphisms-lemma-stein-universally-closed-residue-fields}
Dans le lemme \ref{more-morphisms-lemma-stein-universally-closed}, supposons en outre que
$f$ soit localement de type fini. Alors, pour $s \in S$, la fibre
$\pi^{-1}(\{s\}) = \{s_1, \ldots, s_n\}$ est finie et les extensions de corps
$\kappa(s_i)/\kappa(s)$ sont finies.
\end{lemma}

\begin{proof}
Rappelons qu'il n'y a pas de spécialisations entre les points de $\pi^{-1}(\{s\})$ ;
voir Algèbre, lemme \ref{algebra-lemma-integral-no-inclusion}.
Comme $f'$ est surjectif, on obtient que $|X_s| \to \pi^{-1}(\{s\})$ est surjectif.
Observons que $X_s$ est un schéma quasi-séparé de type fini
sur un corps (la quasi-compacité a été démontrée dans la preuve du
lemme cité). Ainsi, $X_s$ est noethérien
(Morphismes, lemme \ref{morphisms-lemma-finite-type-noetherian}).
Un argument topologique (omis) montre alors que $\pi^{-1}(\{s\})$ est fini.
Pour chaque $i$, on peut choisir un point de type fini $x_i \in X_s$ s'envoyant sur $s_i$
(Morphismes, lemme \ref{morphisms-lemma-enough-finite-type-points}).
On conclut que $\kappa(s_i)/\kappa(s)$ est finie :
$x_i$ peut être représenté par un morphisme $\Spec(k_i) \to X_s$
de type fini (d'après notre définition des points de type fini)
et, par conséquent, $\Spec(k_i) \to s = \Spec(\kappa(s))$ est de type fini
(comme composé de morphismes de type fini) ;
ainsi, $k_i/\kappa(s)$ est finie (Morphismes, lemme
\ref{morphisms-lemma-point-finite-type}).
\end{proof}

\begin{lemma}
\label{more-morphisms-lemma-characterize-geometrically-connected-fibres}
Soit $f : X \to S$ un morphisme de schémas.
Soit $s \in S$. Alors $X_s$ est géométriquement connexe si et
seulement si, pour tout voisinage \'etale $(U, u) \to (S, s)$,
le changement de base $X_U \to U$ a une fibre connexe $X_u$.
\end{lemma}

\begin{proof}
Si $X_s$ est géométriquement connexe, tout changement de base de celui-ci est connexe.
Réciproquement, supposons que $X_s$ ne soit pas géométriquement connexe.
Alors, d'après
Variétés, lemme
\ref{varieties-lemma-characterize-geometrically-disconnected},
on voit que $X_s \times_{\Spec(\kappa(s))} \Spec(k)$ est
non connexe pour une certaine
extension de corps séparable finie $k/\kappa(s)$. D'après
le lemme \ref{more-morphisms-lemma-realize-prescribed-residue-field-extension-etale},
il existe un voisinage \'etale affine $(U, u) \to (S, s)$ tel que
$\kappa(u)/\kappa(s)$ soit identifié à $k/\kappa(s)$.
Dans ce cas, $X_u$ est non connexe.
\end{proof}

\begin{theorem}[Factorisation de Stein ; cas noethérien]
\label{more-morphisms-theorem-stein-factorization-Noetherian}
Soit $S$ un schéma localement noethérien.
Soit $f : X \to S$ un morphisme propre.
Il existe une factorisation
$$
\xymatrix{
X \ar[rr]_{f'} \ar[rd]_f & & S' \ar[dl]^\pi \\
& S &
}
$$
ayant les propriétés suivantes :
\begin{enumerate}
\item le morphisme $f'$ est propre à fibres géométriquement connexes,
\item le morphisme $\pi : S' \to S$ est fini,
\item on a $f'_*\mathcal{O}_X = \mathcal{O}_{S'}$,
\item on a $S' = \underline{\Spec}_S(f_*\mathcal{O}_X)$, et
\item $S'$ est la normalisation de $S$ dans $X$, voir
Morphismes, définition \ref{morphisms-definition-normalization-X-in-Y}.
\end{enumerate}
\end{theorem}

\begin{proof}
Soit $f = \pi \circ f'$ la factorisation du
lemme \ref{more-morphisms-lemma-stein-universally-closed}. Notons qu'outre les
conclusions du lemme \ref{more-morphisms-lemma-stein-universally-closed}, on
a aussi que $f'$ est séparé
(Schémas, lemme \ref{schemes-lemma-compose-after-separated})
et de type fini
(Morphismes, lemme \ref{morphisms-lemma-permanence-finite-type}).
Ainsi, $f'$ est propre. D'après
Cohomologie des schémas, proposition
\ref{coherent-proposition-proper-pushforward-coherent},
on voit que $f_*\mathcal{O}_X$ est un $\mathcal{O}_S$-module cohérent.
On voit donc que $\pi$ est fini, c'est-à-dire que (2) est satisfaite.

\medskip\noindent
Cela démontre toutes les assertions sauf la plus intéressante, à savoir que
toutes les fibres de $f'$ sont géométriquement connexes.
Il résulte clairement de la discussion précédente que l'on peut remplacer $S$ par $S'$ ;
on peut donc supposer que $S$ est noethérien, affine,
que $f : X \to S$ est propre et que $f_*\mathcal{O}_X = \mathcal{O}_S$.
Soit $s \in S$ un point de $S$. Il faut montrer que $X_s$ est
géométriquement connexe. D'après le lemme
\ref{more-morphisms-lemma-characterize-geometrically-connected-fibres},
on voit qu'il suffit de montrer que $X_u$ est connexe
pour tout voisinage \'etale $(U, u) \to (S, s)$.
On peut supposer que $U$ est affine. Alors $U$ est noethérien
(Morphismes, lemme \ref{morphisms-lemma-finite-type-noetherian}),
le changement de base $f_U : X_U \to U$ est propre
(Morphismes, lemme \ref{morphisms-lemma-base-change-proper}),
et l'on a aussi $(f_U)_*\mathcal{O}_{X_U} = \mathcal{O}_U$
(Cohomologie des schémas, lemme \ref{coherent-lemma-flat-base-change-cohomology}).
Ainsi, après avoir remplacé
$(f : X \to S, s)$ par le changement de base $(f_U : X_U \to U, u)$,
il suffit de démontrer que la fibre $X_s$ est connexe lorsque
$f_*\mathcal{O}_X = \mathcal{O}_S$. On peut le déduire de
Catégories dérivées de schémas, lemme
\ref{perfect-lemma-proper-idempotent-on-fibre}
(en considérant les idempotents du faisceau structural de $X_s$),
mais nous donnerons aussi ci-dessous un argument direct.

\medskip\noindent
Plus précisément, appliquons le théorème des fonctions formelles,
sous la forme de Cohomologie des schémas, lemme
\ref{coherent-lemma-formal-functions-stalk}.
Il nous dit que
$$
\mathcal{O}^\wedge_{S, s} = (f_*\mathcal{O}_X)_s^\wedge =
\lim_n H^0(X_n, \mathcal{O}_{X_n})
$$
où $X_n$ est le $n$-ième voisinage infinitésimal de $X_s$.
Puisque l'espace topologique sous-jacent à $X_n$ est égal à celui
de $X_s$, on voit que, si $X_s = T_1 \amalg T_2$ est une réunion disjointe
de sous-schémas ouverts et fermés non vides, alors, de même,
$X_n = T_{1, n} \amalg T_{2, n}$ pour tout $n$. Cela signifie à son tour que
$H^0(X_n, \mathcal{O}_{X_n})$ contient un idempotent non trivial $e_{1, n}$,
à savoir la fonction identiquement égale à $1$ sur $T_{1, n}$ et
identiquement égale à $0$ sur $T_{2, n}$. Il est clair que $e_{1, n + 1}$
se restreint à $e_{1, n}$ sur $X_n$. Ainsi, $e_1 = \lim e_{1, n}$
est un idempotent non trivial de la limite. Cela contredit le fait
que $\mathcal{O}^\wedge_{S, s}$ soit un anneau local. Ainsi,
l'hypothèse était fausse, c'est-à-dire que $X_s$ est connexe, ce qui conclut.
\end{proof}

\begin{theorem}[Factorisation de Stein ; cas général]
\label{more-morphisms-theorem-stein-factorization-general}
Soit $S$ un schéma.
Soit $f : X \to S$ un morphisme propre.
Il existe une factorisation
$$
\xymatrix{
X \ar[rr]_{f'} \ar[rd]_f & & S' \ar[dl]^\pi \\
& S &
}
$$
ayant les propriétés suivantes :
\begin{enumerate}
\item le morphisme $f'$ est propre à fibres géométriquement connexes,
\item le morphisme $\pi : S' \to S$ est entier,
\item on a $f'_*\mathcal{O}_X = \mathcal{O}_{S'}$,
\item on a $S' = \underline{\Spec}_S(f_*\mathcal{O}_X)$, et
\item $S'$ est la normalisation de $S$ dans $X$, voir
Morphismes, définition \ref{morphisms-definition-normalization-X-in-Y}.
\end{enumerate}
\end{theorem}

\begin{proof}
On peut appliquer le lemme \ref{more-morphisms-lemma-stein-universally-closed} pour obtenir le
morphisme $f' : X \to S'$.
Notons qu'outre les
conclusions du lemme \ref{more-morphisms-lemma-stein-universally-closed}, on
a aussi que $f'$ est séparé
(Schémas, lemme \ref{schemes-lemma-compose-after-separated})
et de type fini
(Morphismes, lemme \ref{morphisms-lemma-permanence-finite-type}).
Ainsi, $f'$ est propre. À ce stade, on a démontré toutes les
assertions, sauf celle selon laquelle
$f'$ a des fibres géométriquement connexes.

\medskip\noindent
On peut supposer que $S = \Spec(R)$ est affine.
Posons $R' = \Gamma(X, \mathcal{O}_X)$. Alors $S' = \Spec(R')$.
On peut donc remplacer $S$ par $S'$ et supposer que
$S = \Spec(R)$ est affine et que $R = \Gamma(X, \mathcal{O}_X)$.
Ensuite, soit $s \in S$ un point. Soit $U \to S$ un morphisme \'etale
de schémas affines et soit $u \in U$ un point s'envoyant sur $s$.
Soit $X_U \to U$ le changement de base de $X$. D'après
le lemme \ref{more-morphisms-lemma-characterize-geometrically-connected-fibres},
il suffit de montrer que la fibre de $X_U \to U$ au-dessus de $u$ est
connexe. D'après
Cohomologie des schémas, lemme \ref{coherent-lemma-flat-base-change-cohomology},
on voit que
$\Gamma(X_U, \mathcal{O}_{X_U}) = \Gamma(U, \mathcal{O}_U)$.
Il faut donc montrer ceci : étant donnés
$S = \Spec(R)$ affine, $X \to S$ propre avec $\Gamma(X, \mathcal{O}_X) = R$
et $s \in S$ un point, la fibre $X_s$ est connexe.

\medskip\noindent
Pour cela, il suffit de montrer que les seuls idempotents
$e \in H^0(X_s, \mathcal{O}_{X_s})$ sont $0$ et $1$ (on sait déjà
que $X_s$ est non vide d'après le lemme \ref{more-morphisms-lemma-stein-universally-closed}).
D'après Catégories dérivées de schémas, lemme
\ref{perfect-lemma-proper-idempotent-on-fibre},
après avoir remplacé $R$ par une localisation principale,
on peut supposer que $e$ est l'image d'un élément de $R$.
Puisque $R \to H^0(X_s, \mathcal{O}_{X_s})$ se factorise par
$\kappa(s)$, on conclut.
\end{proof}

\noindent
Voici une application.

\begin{lemma}
\label{more-morphisms-lemma-geometrically-connected-fibres-towards-normal}
Soit $f : X \to S$ un morphisme de schémas. Supposons que
\begin{enumerate}
\item $f$ soit propre,
\item $S$ soit intègre de point générique $\xi$,
\item $S$ soit normal,
\item $X$ soit réduit,
\item tout point générique d'une composante irréductible de $X$ s'envoie sur $\xi$,
\item on ait $H^0(X_\xi, \mathcal{O}) = \kappa(\xi)$.
\end{enumerate}
Alors $f_*\mathcal{O}_X = \mathcal{O}_S$ et $f$
a des fibres géométriquement connexes.
\end{lemma}

\begin{proof}
Appliquons le théorème \ref{more-morphisms-theorem-stein-factorization-general} pour obtenir une
factorisation $X \to S' \to S$. Il suffit de montrer que $S' = S$.
Cela résultera de Morphismes, lemme
\ref{morphisms-lemma-finite-birational-over-normal}.
En effet, $S'$ est réduit parce que $X$ est réduit
(Morphismes, lemme \ref{morphisms-lemma-normalization-in-reduced}).
Le morphisme $S' \to S$ est entier d'après le théorème cité ci-dessus.
Tout point générique de $S'$ est au-dessus de $\xi$ d'après
Morphismes, lemme \ref{morphisms-lemma-normalization-generic},
et l'hypothèse (5). D'autre part, puisque $S'$ est le spectre relatif
de $f_*\mathcal{O}_X$, on voit que la fibre schématique
$S'_\xi$ est le spectre de $H^0(X_\xi, \mathcal{O})$, qui est
égal à $\kappa(\xi)$ par hypothèse. Ainsi, $S'$ est un schéma
intègre dont le corps des fonctions est égal à celui de $S$.
Cela conclut la preuve.
\end{proof}

\noindent
Voici une autre application.

\begin{lemma}
\label{more-morphisms-lemma-proper-flat-nr-geom-conn-comps-lower-semicontinuous}
Soit $X \to S$ un morphisme plat propre de présentation finie. Soit
$n_{X/S}$ la fonction sur $S$ qui compte le nombre de composantes
connexes géométriques des fibres de $f$ introduite dans le
lemme \ref{more-morphisms-lemma-base-change-fibres-nr-geometrically-connected-components}.
Alors $n_{X/S}$ est semi-continue inférieurement.
\end{lemma}

\begin{proof}
Soit $s \in S$. Posons $n = n_{X/S}(s)$. Notons que $n < \infty$, car la fibre
géométrique de $X \to S$ en $s$ est un schéma propre sur un corps, donc noethérien,
et possède donc un nombre fini de composantes connexes. Il faut trouver un
voisinage ouvert $V$ de $s$ tel que $n_{X/S}|_V \geq n$.
Soit $X \to S' \to S$ la factorisation de Stein du
théorème \ref{more-morphisms-theorem-stein-factorization-general}.
D'après le lemme \ref{more-morphisms-lemma-stein-universally-closed-residue-fields},
il y a un nombre fini de points $s'_1, \ldots, s'_m \in S'$ au-dessus de $s$,
et les extensions $\kappa(s'_i)/\kappa(s)$ sont finies.
Le lemme \ref{more-morphisms-lemma-etale-makes-integral-split}
nous dit alors qu'après avoir remplacé $S$ par un voisinage \'etale
de $s$, on peut supposer que $S' = V_1 \amalg \ldots \amalg V_m$ comme schéma,
avec $s'_i \in V_i$ et $\kappa(s'_i)/\kappa(s)$ purement inséparable.
Alors les schémas $X_{s_i'}$ sont géométriquement connexes sur
$\kappa(s)$, donc $m = n$. Les schémas
$X_i = (f')^{-1}(V_i)$, $i = 1, \ldots, n$,
sont plats et de présentation finie sur $S$. Par conséquent, l'image de $X_i \to S$
est ouverte (Morphismes, lemme \ref{morphisms-lemma-fppf-open}).
Ainsi, dans un voisinage de $s$, on voit que $n_{X/S}$ est
au moins égal à $n$.
\end{proof}

\begin{lemma}
\label{more-morphisms-lemma-proper-flat-geom-red}
Soit $f : X \to S$ un morphisme de schémas. Supposons que
\begin{enumerate}
\item $f$ soit propre, plat et de présentation finie, et
\item les fibres géométriques de $f$ soient réduites.
\end{enumerate}
Alors la fonction $n_{X/S} : S \to \mathbf{Z}$,
qui compte le nombre de composantes connexes géométriques
des fibres de $f$, est localement constante.
\end{lemma}

\begin{proof}
D'après le lemme \ref{more-morphisms-lemma-proper-flat-nr-geom-conn-comps-lower-semicontinuous},
la fonction $n_{X/S}$ est semi-continue inférieurement.
Pour $s \in S$, considérons la $\kappa(s)$-algèbre
$$
A_s = H^0(X_s, \mathcal{O}_{X_s})
$$
Définissons une fonction $\beta_0 : X \to \mathbf{Z}$ qui envoie
$s$ sur $\dim_{\kappa(s)} A_s$.
D'après Variétés, lemme
\ref{varieties-lemma-proper-geometrically-reduced-global-sections},
et le fait que $X_s$ est géométriquement réduit,
$A_s$ est un produit fini d'extensions séparables finies de $\kappa(s)$.
Ainsi, $A_s \otimes_{\kappa(s)} \kappa(\overline{s})$ est un produit
de $\beta_0(s)$ copies de $\kappa(\overline{s})$. Par conséquent,
$X_{\overline{s}}$ a $\beta_0(s)$ composantes connexes.
Autrement dit, on a $n_{X/S} = \beta_0$
comme fonctions sur $S$. Ainsi, $n_{X/S}$ est semi-continue supérieurement d'après
Catégories dérivées de schémas, lemme \ref{perfect-lemma-jump-loci-geometric}.
Cela conclut la preuve.
\end{proof}

\noindent
Une dernière application.

\begin{lemma}
\label{more-morphisms-lemma-split-off-proper-part-henselian}
\begin{reference}
Une référence pour le cas d'une base noethérienne adique est
\cite[III, Proposition 5.5.1]{EGA}
\end{reference}
Soit $(A, I)$ un couple hensélien. Soit $X \to \Spec(A)$
séparé et de type fini. Posons $X_0 = X \times_{\Spec(A)} \Spec(A/I)$.
Soit $Y \subset X_0$ un sous-schéma ouvert et fermé tel que
$Y \to \Spec(A/I)$ soit propre. Alors il existe un sous-schéma ouvert et fermé
$W \subset X$, propre sur $A$, tel que
$W \times_{\Spec(A)} \Spec(A/I) = Y$.
\end{lemma}

\begin{proof}
Nous noterons $T \mapsto T_0$ le changement de base par $\Spec(A/I) \to \Spec(A)$.
D'après le lemme de Chow (sous la forme du
lemme de Limites \ref{limits-lemma-chow-finite-type}),
il existe un morphisme propre surjectif $\varphi : X' \to X$ tel
que $X'$ admette une immersion dans $\mathbf{P}^n_A$.
Posons $Y' = \varphi^{-1}(Y)$. C'est un sous-schéma ouvert et fermé
de $X'_0$. Supposons le lemme vrai pour $(X', Y')$. Soit $W' \subset X'$
le sous-schéma ouvert et fermé propre sur $A$ tel que $Y' = W'_0$.
D'après Morphismes, lemme \ref{morphisms-lemma-image-proper-scheme-closed},
$W = \varphi(W') \subset X$ et
$Q = \varphi(X' \setminus W') \subset X$ sont des sous-ensembles fermés et, d'après
Morphismes, lemme \ref{morphisms-lemma-image-proper-is-proper},
$W$ est propre sur $A$. L'image de $W \cap Q$ dans $\Spec(A)$ est fermée.
Puisque $(A, I)$ est hensélien, si $W \cap Q$ est non vide, alors
il existe dans $W \cap Q$ un point au-dessus de $\Spec(A/I)$.
C'est impossible puisque $W'_0 = Y' = \varphi^{-1}(Y)$.
On conclut que $W$ est un sous-schéma ouvert et fermé
de $X$, propre sur $A$, avec $W_0 = Y$.
On se ramène ainsi au cas décrit au paragraphe suivant.

\medskip\noindent
Supposons qu'il existe une immersion $j : X \to \mathbf{P}^n_A$ sur $A$.
Soit $\overline{X}$ l'image schématique de $j$.
Puisque $j$ est un morphisme quasi-compact
(Schémas, lemme \ref{schemes-lemma-quasi-compact-permanence}),
on voit que $j : X \to \overline{X}$ est une immersion ouverte
(Morphismes, lemme \ref{morphisms-lemma-quasi-compact-immersion}).
Ainsi, le changement de base $j_0 : X_0 \to \overline{X}_0$
est également une immersion ouverte.
Par conséquent, $j_0(Y) \subset \overline{X}_0$ est ouvert.
Il est aussi fermé d'après Morphismes, lemme
\ref{morphisms-lemma-image-proper-scheme-closed}.
Supposons le lemme vrai pour $(\overline{X}, j_0(Y))$.
Soit $\overline{W} \subset \overline{X}$ le
sous-schéma ouvert et fermé correspondant, propre sur $A$,
tel que $j_0(Y) = \overline{W}_0$.
Alors $T = \overline{W} \setminus j(X)$ est fermé dans $\overline{W}$,
et a donc une image fermée dans $\Spec(A)$, par propreté de $\overline{W}$
sur $A$. Puisque $(A, I)$ est hensélien, on voit que si $T$
est non vide, il existe un point de $T$ s'envoyant dans $\Spec(A/I)$.
C'est impossible parce que $j_0(Y) = \overline{W}_0$ est contenu dans $j(X)$.
Ainsi, $\overline{W}$ est contenu dans $j(X)$ et l'on peut
prendre pour $W \subset X$ l'unique sous-schéma ouvert et fermé
s'envoyant isomorphiquement sur $\overline{W}$ par $j$.
On se ramène ainsi au cas décrit au paragraphe suivant.

\medskip\noindent
Supposons que $X \subset \mathbf{P}^n_A$ soit un sous-schéma fermé.
Alors $X \to \Spec(A)$ est un morphisme propre.
Soit $Z = X_0 \setminus Y$. C'est un sous-schéma ouvert et fermé
de $X_0$, et $X_0 = Y \amalg Z$.
Soit $X \to X' \to \Spec(A)$ la factorisation de Stein du
théorème \ref{more-morphisms-theorem-stein-factorization-general}.
Soient $Y' \subset X'_0$ et $Z' \subset X'_0$ les images de
$Y$ et de $Z$.
Puisque les fibres de $X \to Z$ sont géométriquement connexes,
on voit que $Y' \cap Z' = \emptyset$.
Ainsi, $X'_0 = Y' \amalg Z'$, car $X \to X'$ est surjectif.
Puisque $X' \to \Spec(A)$ est entier, on voit que
$X'$ est le spectre d'une $A$-algèbre entière sur $A$.
Rappelons que les sous-ensembles ouverts et fermés des spectres correspondent
$1$ à $1$ aux idempotents de l'anneau correspondant, voir
Algèbre, lemme \ref{algebra-lemma-disjoint-decomposition}.
Par conséquent, d'après
Compléments d'algèbre, lemme \ref{more-algebra-lemma-characterize-henselian-pair},
on voit que l'on peut écrire $X' = W' \amalg V'$
avec $W'$ et $V'$ ouverts et fermés, et
avec $Y' = W'_0$ et $Z' = V'_0$.
Soit $W$ l'image réciproque du premier de ces sous-schémas dans $X$,
ce qui conclut la preuve.
\end{proof}









\section{Stratification par platitude générique}
\label{more-morphisms-section-generic-flatness-stratification}

\noindent
On peut utiliser la platitude générique pour construire une
stratification de la base telle qu'un module donné
devienne plat sur les strates.

\begin{lemma}[Stratification par platitude générique]
\label{more-morphisms-lemma-generic-flatness-stratification}
Soit $f : X \to S$ un morphisme de présentation finie entre des schémas quasi-compacts
et quasi-séparés. Soit $\mathcal{F}$ un $\mathcal{O}_X$-module
de présentation finie. Alors il existe un $t \geq 0$ et des sous-schémas
fermés
$$
S \supset S_0 \supset S_1 \supset \ldots \supset S_t = \emptyset
$$
tels que $S_i \to S$ soit défini par un faisceau d'idéaux de type fini,
que $S_0 \subset S$ soit un épaississement et que le tiré en arrière de $\mathcal{F}$ sur
$X \times_S (S_i \setminus S_{i + 1})$ soit plat sur $S_i \setminus S_{i + 1}$.
\end{lemma}

\begin{proof}
On peut trouver un diagramme cartésien
$$
\xymatrix{
X \ar[d] \ar[r] & X_0 \ar[d] \\
S \ar[r] & S_0
}
$$
et un $\mathcal{O}_{X_0}$-module de présentation finie $\mathcal{F}_0$
dont le tiré en arrière est $\mathcal{F}$, où $X_0$ et $S_0$ sont de
type fini sur $\mathbf{Z}$. Voir
Limites, proposition \ref{limits-proposition-approximate}, et
les lemmes \ref{limits-lemma-descend-finite-presentation} et
\ref{limits-lemma-descend-modules-finite-presentation}.
Ainsi, on peut supposer que $X$ et $S$ sont de type fini sur $\mathbf{Z}$
et que $\mathcal{F}$ est un $\mathcal{O}_X$-module cohérent.

\medskip\noindent
Supposons que $X$ et $S$ soient de type fini sur $\mathbf{Z}$
et que $\mathcal{F}$ soit un $\mathcal{O}_X$-module cohérent.
Dans ce cas, tout idéal quasi-cohérent est de type fini ; ainsi,
on n'a pas à vérifier que $S_i$ est défini
par un idéal de type fini. Posons $S_0 = S_{red}$, le réduit de $S$.
D'après la platitude générique telle qu'énoncée dans Morphismes, proposition
\ref{morphisms-proposition-generic-flatness-reduced},
il existe un ouvert dense $U_0 \subset S_0$ tel que le tiré en arrière de $\mathcal{F}$
sur $X \times_S U_0$ soit plat sur $U_0$.
Soit $S_1 \subset S_0$ le sous-schéma fermé réduit dont
le sous-ensemble fermé sous-jacent est $S \setminus U_0$. On continue de cette
manière, pourvu que $S_1 \not = \emptyset$, afin d'obtenir
$S_0 \supset S_1 \supset \ldots$. Comme $S$
est noethérien, toute chaîne descendante de sous-ensembles fermés se stabilise ;
on voit donc que $S_t = \emptyset$ pour un certain $t \geq 0$.
\end{proof}

\begin{lemma}
\label{more-morphisms-lemma-generic-flatness-stratification-scheme}
Soit $f : X \to S$ un morphisme de présentation finie entre des schémas quasi-compacts
et quasi-séparés. Alors il existe un $t \geq 0$ et des sous-schémas
fermés
$$
S \supset S_0 \supset S_1 \supset \ldots \supset S_t = \emptyset
$$
tels que $S_i \to S$ soit défini par un faisceau d'idéaux de type fini,
que $S_0 \subset S$ soit un épaississement et que
$X \times_S (S_i \setminus S_{i + 1})$ soit plat sur $S_i \setminus S_{i + 1}$.
\end{lemma}

\begin{proof}
Appliquons le lemme \ref{more-morphisms-lemma-generic-flatness-stratification}
avec $\mathcal{F} = \mathcal{O}_X$.
\end{proof}

\begin{lemma}[Longke Tang]
\label{more-morphisms-lemma-flatness-criterion}
Soit $f : X \to S$ un morphisme surjectif de présentation finie.
Soit $M$ un objet de $D_\QCoh(\mathcal{O}_S)$ tel que
$H^i(M) = 0$ pour $i > 0$. Les conditions suivantes sont équivalentes :
\begin{enumerate}
\item $M$ est isomorphe à un $\mathcal{O}_S$-module plat
placé en degré $0$,
\item $Lf^*M$ est isomorphe à un $\mathcal{O}_X$-module plat
placé en degré $0$.
\end{enumerate}
\end{lemma}

\begin{proof}
On omet la preuve que (1) implique (2). Supposons (2). Démontrer (1)
est une question locale ; on peut donc supposer que $S$ est affine.
Choisissons une stratification
$S \supset S_0 \supset S_1 \supset \ldots \supset S_t = \emptyset$
comme dans le lemme \ref{more-morphisms-lemma-generic-flatness-stratification-scheme}.
Posons $T_i = S_i \setminus S_{i - 1}$ et $Y_i = X \times_S T_i$.
Notons $g_i : T_i \to S$ et $h_i : Y_i \to X$ les morphismes d'inclusion.
Comme $f$ est surjectif, les morphismes $f_i : X_i \to T_i$
sont surjectifs et plats. Puisque
$Lh_i^* \circ Lf^* = Lf_i^* \circ Lg_i^*$, on déduit de (2) que
$Lf_i^* Lg_i^*M$ est isomorphe à un $\mathcal{O}_{Y_i}$-module plat
placé en degré $0$. Comme $f_i$ est plat et surjectif, on voit que
$Lg_i^*M$ est isomorphe à un $\mathcal{O}_{T_i}$-module plat
placé en degré $0$.

\medskip\noindent
On démontrera que cela implique le résultat par récurrence sur $t$.
On utilisera les notations suivantes : $S = \Spec(A)$ et $M$
est l'objet de $D_\QCoh(\mathcal{O}_S)$ correspondant à
$N$ dans $D(A)$, voir Catégories dérivées de schémas, lemme
\ref{perfect-lemma-affine-compare-bounded}.

\medskip\noindent
Cas de base : $t = 1$. Écrivons $S_0 = \Spec(A/I)$, où
$I$ est un idéal de type fini. Alors $I$ est nilpotent, puisque $S_0 = S$
ensemblistement. L'hypothèse est que $N \otimes_A^\mathbf{L} A/I$
a une amplitude de Tor dans $[0, 0]$. D'après Compléments d'algèbre, lemme
\ref{more-algebra-lemma-check-tor-dimension-modulo-nilpotent},
il en va de même pour $N$.

\medskip\noindent
Étape de récurrence. Supposons que $t > 1$. Écrivons $S_{t - 1} = \Spec(A/I)$,
où $I = (f_1, \ldots, f_r)$ est un idéal de type fini.
On raisonne par récurrence sur $r$. Observons que
$N_{f_r} = N \otimes_A^\mathbf{L} A_{f_r}$ dans $D(A_{f_r})$
a une amplitude de Tor dans $[0, 0]$ par l'hypothèse de récurrence
(car la stratification sur l'ouvert principal $D(f_r)$ est de longueur
au plus $t - 1$). D'autre part, considérons
$N' = N \otimes_A^\mathbf{L} A/f_rA$ dans $D(A/f_rA)$.
En posant $S' = \Spec(A/f_rA)$, on obtient une stratification
$$
S' \supset S' \cap S_0 \supset S' \cap S_1 \supset \ldots \supset
S' \cap S_{t - 1} \supset
S' \cap S_t = \emptyset
$$
où $S' \cap S_{t - 1}$ est découpé par $r - 1$ équations dans $S'$.
Comme précédemment, les tirés en arrière dérivés de $M$ sur les parties
$S' \cap S_i \setminus S' \cap S_{i + 1}$ ont une amplitude de Tor
dans $[0, 0]$. Par récurrence sur $r$, on conclut donc que
$N'$ a une amplitude de Tor dans $[0, 0]$.
On conclut finalement que $N$ a une amplitude de Tor dans $[0, 0]$
d'après Compléments d'algèbre, lemme \ref{more-algebra-lemma-f-flat}.
\end{proof}

\begin{lemma}
\label{more-morphisms-lemma-cokernel-flat}
Soit $R$ un anneau intègre noethérien. Soient $R \to A \to B$ des morphismes d'anneaux de type fini.
Soit $M$ un $A$-module de type fini et soit $N$ un $B$-module de type fini.
Soit $M \to N$ un morphisme $A$-linéaire. Il existe un élément non nul $f \in R$
tel que le conoyau de $M_f \to N_f$ soit un $R_f$-module plat.
\end{lemma}

\begin{proof}
En remplaçant $M$ par l'image de $M \to N$, on peut supposer que $M \subset N$.
Choisissons une filtration $0 = N_0 \subset N_1 \subset \ldots \subset N_t = N$
telle que $N_i/N_{i - 1} = B/\mathfrak q_i$ pour un idéal premier
$\mathfrak q_i \subset B$, voir
Algèbre, lemme \ref{algebra-lemma-filter-Noetherian-module}.
Posons $M_i = M \cap N_i$. Alors $Q = N/M$ possède une filtration par les sous-modules
$Q_i = N_i/M_i$. Il suffit de montrer que $Q_i/Q_{i - 1}$ devient plat
après localisation en un élément non nul $f$ (puisque les extensions de
modules plats sont plates d'après Algèbre, lemme \ref{algebra-lemma-flat-ses}).
Puisque $Q_i/Q_{i - 1}$ est isomorphe au conoyau
du morphisme $M_i/M_{i - 1} \to N_i/N_{i - 1}$, on se ramène au
cas traité au paragraphe suivant.

\medskip\noindent
Supposons que $B$ soit un anneau intègre et que $M \subset N = B$. Après avoir remplacé $A$ par
l'image de $A$ dans $B$, on peut supposer que $A \subset B$. Par platitude générique,
on peut supposer que $A$ et $B$ sont plats sur $R$
(Algèbre, lemme \ref{algebra-lemma-generic-flatness-Noetherian}).
Il suffit maintenant de montrer que $M \to B$ devient $R$-universellement
injectif après remplacement de $R$ par une localisation
principale (Algèbre, lemme \ref{algebra-lemma-ui-flat-domain}).
Par liberté générique, on peut trouver un élément non nul $g \in A$ tel que
$B_g$ soit un $A_g$-module libre
(Algèbre, lemme \ref{algebra-lemma-generic-flatness-Noetherian}).
On peut donc choisir un facteur direct $M' \subset B_g$ comme $A_g$-module,
qui soit un $A_g$-module libre de type fini, et
tel que $M \to B \to B_g$ se factorise par $M'$.
Il suffit manifestement de montrer que $M \to M'$
devient $R$-universellement injectif après remplacement de
$R$ par une localisation principale.

\medskip\noindent
Écrivons $M' = A_g^{\oplus n}$. Puisque $M \subset M'$ est un $A$-module de type fini,
on voit que $M$ est contenu dans $(1/g^m)A^{\oplus n}$ pour un certain $m \geq 0$.
Après changement de la base choisie de $M'$, on peut supposer que $M \subset A^{\oplus n}$.
Il suffit alors de montrer que $A^{\oplus n}/M$ et $A_g/A$ deviennent
$R$-plats après remplacement de $R$ par une localisation principale. En effet, alors
$M \to A^{\oplus n}$ et $A^{\oplus n} \to A_g^{\oplus n}$ sont
universellement injectifs d'après Algèbre, lemme \ref{algebra-lemma-flat-tor-zero},
et il en va donc de même du composé $M \to M' = A_g^{\oplus n}$.

\medskip\noindent
Par platitude générique (voir la référence ci-dessus), on peut supposer que le
module $A^{\oplus n}/M$ est $R$-plat. Pour le quotient $A_g/A$,
on utilise le fait que
$$
A_g/A = \colim (1/g^m)A/A \cong \colim A/g^mA
$$
et que le module $A/g^mA$ possède une filtration de longueur $m$ dont les
quotients successifs sont isomorphes à $A/gA$. De nouveau par platitude
générique, on peut supposer que $A/gA$ est $R$-plat, donc chaque $A/g^mA$
est $R$-plat, et par suite $A_g/A$ l'est aussi.
\end{proof}

\noindent
Soit $f : X \to Y$ un morphisme de schémas sur un schéma de base $S$.
Soit $Z \subset Y$ l'image schématique de $f$, voir
Morphismes, section \ref{morphisms-section-scheme-theoretic-image}.
Soit $g : S' \to S$ un morphisme de schémas et soit
$f' : X \times_S S' \to Y \times_S S'$ le changement de base de $f$ par $g$.
Il n'est pas toujours vrai que $Z \times_S S' \subset Y \times_S S'$
soit l'image schématique de $f'$.
Disons que {\it la formation de l'image schématique de $f/S$
commute à tout changement de base} si, pour tout $g$ comme ci-dessus,
l'image schématique de $f'$ est égale à $Z \times_S S'$.

\begin{lemma}
\label{more-morphisms-lemma-base-change-scheme-theoretic-image}
Soit $S$ un schéma quasi-compact et quasi-séparé.
Soit $f : X \to Y$ un morphisme de schémas sur $S$, les schémas
$X$ et $Y$ étant tous deux de présentation finie sur $S$.
Il existe alors un entier $t \geq 0$ et des sous-schémas fermés
$$
S \supset S_0 \supset S_1 \supset \ldots \supset S_t = \emptyset
$$
possédant les propriétés suivantes :
\begin{enumerate}
\item $S_i \to S$ est défini par un faisceau d'idéaux de type fini,
\item $S_0 \subset S$ est un épaississement, et
\item si $T_i = S_i \setminus S_{i + 1}$ et si $f_i$ est le changement de
base de $f$ à $T_i$, alors :
la formation de l'image schématique de $f_i/T_i$
commute à tout changement de base (voir la discussion précédant le lemme).
\end{enumerate}
\end{lemma}

\begin{proof}
On peut trouver un diagramme commutatif
$$
\xymatrix{
X \ar[d] \ar[r] & Y \ar[d] \ar[r] & S \ar[d] \\
U \ar[r] & V \ar[r] & W
}
$$
à carrés cartésiens tel que $U$, $V$, $W$
soient de type fini sur $\mathbf{Z}$. En effet, on écrit d'abord $S$
comme limite cofiltrée de schémas de type fini sur $\mathbf{Z}$
dont les morphismes de transition sont affines,
en utilisant Limites, proposition \ref{limits-proposition-approximate},
puis on fait descendre le morphisme $X \to Y$ à l'aide de
Limites, lemme \ref{limits-lemma-descend-finite-presentation}.
On se ramène ainsi au cas traité au paragraphe suivant.

\medskip\noindent
Supposons que $S$ soit noethérien.  Dans ce cas, tout idéal quasi-cohérent
est de type fini ; il n'est donc pas nécessaire de vérifier que
$S_i$ est défini par un idéal de type fini. Posons $S_0 = S_{red}$, le
réduit de $S$. Soit $\eta \in S_0$ un point générique
d'une composante irréductible de $S_0$. Par récurrence noethérienne sur
l'espace topologique sous-jacent de $S_0$, on peut supposer que le résultat
vaut pour tout sous-schéma fermé de $S_0$ ne contenant pas $\eta$.
Il suffit donc de montrer qu'il existe un voisinage ouvert
$U_0 \subset S_0$ tel que le changement de base $f_0$ de $f$ à $U_0$
possède la propriété (3).

\medskip\noindent
Soit $R$ un anneau intègre noethérien. Soit $f : X \to Y$ un morphisme
de schémas de type fini sur $R$. D'après la discussion du paragraphe précédent,
il suffit de montrer qu'après remplacement de $R$ par $R_g$ pour un certain $g \in R$
non nul, et de $X$, $Y$ par leurs changements de base à $R_g$,
la formation de l'image schématique de $f/R$ commute
à tout changement de base.

\medskip\noindent
Soit $Y = V_1 \cup \ldots V_n$ un recouvrement ouvert affine.
Posons $U_i = f^{-1}(V_i)$. Si l'assertion est vraie pour chacun des
morphismes $U_i \to V_i$ sur $R$, alors elle est vraie pour $f$.
En effet, l'image schématique de $U_i \to V_i$
est l'intersection de $V_i$ avec l'image schématique de
$f : X \to Y$, d'après Morphismes, lemme
\ref{morphisms-lemma-quasi-compact-scheme-theoretic-image}.
On peut donc supposer que $Y$ est affine.

\medskip\noindent
Soit $X = U_1 \cup \ldots U_n$ un recouvrement ouvert affine.
Alors l'image schématique de $X \to Y$ est la même que
l'image schématique de $\coprod U_i \to Y$.
On peut donc supposer que $X$ est affine.

\medskip\noindent
Écrivons $X = \Spec(A)$ et $Y = \Spec(B)$, et supposons que $f$ corresponde
au morphisme de $R$-algèbres $\varphi : A \to B$.
Alors l'image schématique de $f$ est $\Spec(A/\Ker(\varphi))$,
et il en va de même après changement de base (par un morphisme affine, mais il
suffit de le vérifier pour ceux-ci).
Ainsi, la formation de l'image schématique commute
au changement de base si $\Ker(\varphi \otimes_R R') = \Ker(\varphi) \otimes_R R'$
pour tout morphisme d'anneaux $R \to R'$.

\medskip\noindent
Après remplacement de $R$, $A$, $B$ par $R_g$, $A_g$, $B_g$ pour un
élément non nul convenable $g$ de $R$, on peut supposer que $A$ et $B$ sont plats sur $R$.
D'après le lemme \ref{more-morphisms-lemma-cokernel-flat},
on peut également supposer que $B/A$ est un $R$-module plat.
Alors $0 \to \Ker(\varphi) \to A \to B \to B/A \to 0$
est une suite exacte de $R$-modules plats, ce qui implique
l'assertion voulue sur le changement de base.
\end{proof}






\section{Stratification d'un morphisme}
\label{more-morphisms-section-stratifying-morphisms}

\noindent
Soit $f : X \to S$ un morphisme de présentation finie entre schémas quasi-compacts
et quasi-séparés.
Dans la section \ref{more-morphisms-section-generic-flatness-stratification},
nous avons vu que l'on peut stratifier $S$ de façon que $X$ soit plat
sur les strates. Dans cette section, nous cherchons des stratifications de $S$
et de $X$ telles que les strates obtenues soient lisses ; cela ne fonctionne pas tout à fait,
et nous aurons également besoin d'un changement de base par des morphismes finis localement libres.

\begin{lemma}
\label{more-morphisms-lemma-smoothness-stratification-at-generic-point}
Soit $f : X \to S$ un morphisme de schémas de présentation finie.
Soit $\eta \in S$ un point générique d'une composante irréductible de $S$.
Supposons que $S$ soit réduit. Il existe alors
\begin{enumerate}
\item un sous-schéma ouvert $U \subset S$ contenant $\eta$,
\item un morphisme fini localement libre, surjectif et universellement injectif
$V \to U$,
\item un entier $t \geq 0$ et des sous-schémas fermés
$$
X \times_S V \supset Z_0 \supset Z_1 \supset \ldots \supset Z_t = \emptyset
$$
tels que $Z_i \to X \times_S V$ soit défini par un faisceau d'idéaux de type fini,
que $Z_0 \subset X \times_S V$ soit un épaississement et que le morphisme
$Z_i \setminus Z_{i + 1} \to V$ soit lisse.
\end{enumerate}
\end{lemma}

\begin{proof}
Il est clair que l'on peut remplacer $S$ par un voisinage ouvert de $\eta$
et $X$ par sa restriction à cet ouvert.
On peut donc supposer que $S = \Spec(A)$, où $A$ est un anneau réduit,
et que $\eta$ correspond à un idéal premier minimal $\mathfrak p$.
Rappelons que, dans ce cas, l'anneau local $\mathcal{O}_{S, \eta} = A_\mathfrak p$
est égal à $\kappa(\mathfrak p)$, voir
Algèbre, lemme \ref{algebra-lemma-minimal-prime-reduced-ring}.

\medskip\noindent
Appliquons Variétés, lemme \ref{varieties-lemma-smooth-stratification},
au schéma $X_\eta$ sur $k = \kappa(\eta)$.
Notons $k'/k$ l'extension de corps purement inséparable ainsi obtenue.
Au paragraphe suivant, on se ramène au cas $k' = k$.
(Cette étape correspond à la construction du morphisme $V \to U$ de
l'énoncé du lemme ; en particulier, on peut prendre $V = U$
si le corps $\kappa(\mathfrak p)$ est de caractéristique zéro.)

\medskip\noindent
Si le corps $k = \kappa(\mathfrak p)$ est de caractéristique zéro, alors
$k' = k$. Si le corps $k = \kappa(\mathfrak p)$
est de caractéristique $p > 0$, alors $p$ s'envoie sur zéro dans $A_\mathfrak p = \kappa(\mathfrak p)$.
Ainsi, après remplacement de $A$ par une localisation principale (c'est-à-dire après
restriction de $S$), on peut supposer que $p = 0$ dans $A$. Si $k' \not = k$, alors
il existe un élément $\beta \in k'$, $\beta \not \in k$,
tel que $\beta^p \in k$. Après remplacement de $A$ par une localisation
principale, on peut supposer qu'il existe un élément $a \in A$ tel
que $\beta^p = a$. Posons $A' = A[x]/(x^p - a)$.
Alors $S' = \Spec(A') \to \Spec(A) = S$ est fini localement libre,
surjectif et universellement injectif. De plus, si $\mathfrak p' \subset A'$
désigne l'unique idéal premier au-dessus de $\mathfrak p$,
alors $A'_{\mathfrak p'} = k(\beta)$ et $k'/k(\beta)$
est de degré plus petit. Ainsi, après remplacement de $S$ par $S'$
et de $\eta$ par le point $\eta'$ correspondant à $\mathfrak p'$,
on voit que le degré de $k'$ sur le corps résiduel de $\eta$
a diminué. En poursuivant ainsi, une récurrence nous ramène
au cas $k' = \kappa(\mathfrak p) = \kappa(\eta)$.

\medskip\noindent
On peut donc supposer que $S$ est affine et réduit, et qu'il existe un entier $t \geq 0$
et des sous-schémas fermés
$$
X_\eta \supset Z_{\eta, 0} \supset Z_{\eta, 1}
\supset \ldots \supset Z_{\eta, t} = \emptyset
$$
tels que $Z_{\eta, 0} = (X_\eta)_{red}$ et que
$Z_{\eta, i} \setminus Z_{\eta, i + 1}$
soit lisse sur $\eta$ pour tout $i$.
Rappelons que $\kappa(\eta) = \kappa(\mathfrak p) = A_\mathfrak p$ est
la limite inductive filtrante des $A_a$ pour $a \in A$, $a \not \in \mathfrak p$.
Voir Algèbre, lemme \ref{algebra-lemma-localization-colimit}.
On peut donc faire descendre le diagramme ci-dessus en un diagramme correspondant
sur $\Spec(A_a)$ pour un certain $a \in A$, $a \not \in \mathfrak p$.
Plus précisément, après remplacement de $S$ par $\Spec(A_a)$,
on peut supposer qu'il existe un entier $t \geq 0$ et des sous-schémas fermés
$$
X \supset Z_0 \supset Z_1 \supset \ldots \supset Z_t = \emptyset
$$
tels que $Z_i \to X$ soit une immersion fermée de présentation finie,
que $Z_0 \to X$ soit un épaississement et que $Z_i \setminus Z_{i + 1}$
soit lisse sur $S$. Autrement dit, le lemme est démontré.
Plus précisément, on utilise d'abord
Limites, lemme \ref{limits-lemma-descend-finite-presentation},
pour obtenir des morphismes
$$
Z_t \to Z_{t - 1} \to \ldots \to Z_0 \to X
$$
sur $S$, chacun de présentation finie, et dont le
changement de base à $\eta$ redonne les inclusions entre les
sous-schémas fermés donnés ci-dessus. Quitte à rétrécir encore $S$, on peut supposer
que chacun de ces morphismes est une immersion fermée, voir
Limites, lemme
\ref{limits-lemma-descend-closed-immersion-finite-presentation}.
Quitte à rétrécir $S$, on peut supposer que $Z_0 \to X$ est surjectif,
donc est un épaississement, voir
Limites, lemme \ref{limits-lemma-descend-surjective}.
Quitte à rétrécir encore une fois $S$, on peut supposer que
$Z_i \setminus Z_{i + 1} \to S$ est lisse, voir
Limites, lemme \ref{limits-lemma-descend-smooth}.
Cela achève la démonstration.
\end{proof}

\begin{lemma}
\label{more-morphisms-lemma-smoothness-stratification}
Soit $f : X \to S$ un morphisme de présentation finie entre schémas quasi-compacts
et quasi-séparés. Il existe alors un entier $t \geq 0$ et des
sous-schémas fermés
$$
S \supset S_0 \supset S_1 \supset \ldots \supset S_t = \emptyset
$$
tels que
\begin{enumerate}
\item $S_i \to S$ soit défini par un faisceau d'idéaux de type fini,
\item $S_0 \subset S$ soit un épaississement,
\item pour chaque $i$, il existe un morphisme fini localement libre et surjectif
$T_i \to S_i \setminus S_{i + 1}$,
\item pour chaque $i$, il existe un entier $t_i \geq 0$ et des sous-schémas fermés
$$
X_i = X \times_S T_i \supset Z_{i, 0}
\supset Z_{i, 1} \supset \ldots \supset Z_{i, t_i} = \emptyset
$$
tels que $Z_{i, j} \to X_i$ soit défini par un faisceau d'idéaux de type fini,
que $Z_{i, 0} \subset X_i$ soit un épaississement et que le morphisme
$Z_{i, j} \setminus Z_{i, j + 1} \to T_i$ soit lisse.
\end{enumerate}
\end{lemma}

\begin{proof}
On peut trouver un diagramme cartésien
$$
\xymatrix{
X \ar[d] \ar[r] & X_0 \ar[d] \\
S \ar[r] & S_0
}
$$
tel que $X_0$ et $S_0$ soient de type fini sur $\mathbf{Z}$. Voir
Limites, proposition \ref{limits-proposition-approximate} et
lemme \ref{limits-lemma-descend-finite-presentation}.
On peut donc supposer que $X$ et $S$ sont de type fini sur $\mathbf{Z}$.
En effet, une solution du problème posé par le lemme pour $X_0 \to S_0$
donnera par changement de base une solution sur $S$ ; détails omis.

\medskip\noindent
Supposons que $X$ et $S$ soient de type fini sur $\mathbf{Z}$.
Dans ce cas, tout idéal quasi-cohérent est de type fini ; il n'est donc
pas nécessaire de vérifier que $S_i$ est défini
par un idéal de type fini. Posons $S_0 = S_{red}$, le réduit de $S$.
Soit $\eta \in S_0$ un point générique d'une composante irréductible.
D'après le lemme \ref{more-morphisms-lemma-smoothness-stratification-at-generic-point},
on peut trouver un sous-schéma ouvert $U \subset S_0$,
un morphisme fini localement libre, surjectif et universellement injectif
$V \to U$, un entier $t_0 \geq 0$ et des sous-schémas fermés
$$
X \times_S V \supset Z_{0, 0} \supset Z_{0, 1} \supset \ldots \supset
Z_{0, t_0} = \emptyset
$$
tels que $Z_{0, i} \to X \times_S V$ soit défini par un faisceau d'idéaux de type fini,
que $Z_{0, 0} \subset X \times_S V$ soit un épaississement et que le morphisme
$Z_{0, i} \setminus Z_{0, i + 1} \to V$ soit lisse.
On munit alors $S_1 \subset S_0$ de la structure de sous-schéma réduit induit
sur $S_0 \setminus U$. Par récurrence noethérienne sur l'espace
topologique sous-jacent de $S$, on peut supposer que le lemme vaut pour
$X \times_S S_1 \to S_1$. Cela fournit un entier $t \geq 1$ et
$$
S_1 = S_1 \supset S_2 \supset \ldots \supset S_t = \emptyset
$$
ainsi que des $t_i$ et des $Z_{i, j}$ comme dans l'énoncé du lemme.
Cela démontre le lemme.
\end{proof}


















\section{Amélioration des morphismes de dimension relative égale à un}
\sectionmark{Morphismes de dimension relative un}
\label{more-morphisms-section-make-good-curves}

\noindent
On peut rendre toute courbe lisse et projective après extension du
corps de base, compactification et normalisation. Cela implique aussi des
résultats sur les morphismes de type fini dont les fibres génériques sont de dimension $1$.

\begin{lemma}
\label{more-morphisms-lemma-make-good-curves}
Soit $f : X \to S$ un morphisme de schémas. Soit $\eta \in S$ un
point générique d'une composante irréductible de $S$. Supposons que $f$ soit
séparé, de présentation finie, et que $\dim(X_\eta) \leq 1$.
Il existe alors un diagramme commutatif
$$
\xymatrix{
\overline{Y}_1 \amalg \ldots \amalg \overline{Y}_n \ar[rd] &
Y_1 \amalg \ldots \amalg Y_n \ar[r]_-\nu \ar[d] \ar[l]^j &
X_V \ar[r] \ar[d] &
X_U \ar[r] \ar[d] &
X \ar[d]^f \\
& T_1 \amalg \ldots \amalg T_n \ar[r] &
V \ar[r] &
U \ar[r] &
S
}
$$
de schémas possédant les propriétés suivantes :
\begin{enumerate}
\item $U \subset S$ est un voisinage ouvert de $\eta$,
\item $V \to U$ est un morphisme fini, surjectif et universellement injectif,
\item $X_U = U \times_S X$ et $X_V = V \times_S X$ sont les changements de base,
\item $\nu$ est fini, surjectif, et il existe un ouvert $W \subset X_V$
tel que
\begin{enumerate}
\item $W$ soit dense dans toutes les fibres de $X_V \to V$,
\item $\nu^{-1}(W) \cap Y_i$ soit dense dans toutes les fibres de $Y_i \to T_i$, et
\item $\nu^{-1}(W) \to W$ soit un épaississement,
\end{enumerate}
\item $j$ est une immersion ouverte,
\item $T_i \to V$ est fini \'etale,
\item $Y_i \to T_i$ est surjectif et lisse,
\item $\overline{Y}_i \to T_i$ est lisse, propre, à fibres géométriquement
connexes de dimension $\leq 1$.
\end{enumerate}
\end{lemma}

\begin{proof}
Il est clair que l'on peut remplacer $S$ par un voisinage ouvert de $\eta$
et $X$ par sa restriction à cet ouvert.
On peut de plus remplacer $S$ par son réduit et $X$ par
le changement de base à ce réduit.
On peut donc supposer que $S = \Spec(A)$, où $A$ est un anneau réduit,
et que $\eta$ correspond à un idéal premier minimal $\mathfrak p$.
Rappelons que, dans ce cas, l'anneau local $\mathcal{O}_{S, \eta} = A_\mathfrak p$
est égal à $\kappa(\mathfrak p)$, voir
Algèbre, lemme \ref{algebra-lemma-minimal-prime-reduced-ring}.

\medskip\noindent
Appliquons Variétés, lemme
\ref{varieties-lemma-dim-1-projective-completion-after-insep},
au schéma $X_\eta$ sur $k = \kappa(\eta)$.
Notons $k'/k$ l'extension de corps purement inséparable ainsi obtenue.
Au paragraphe suivant, on se ramène au cas $k' = k$.
(Cette étape correspond à la construction du morphisme $V \to U$ de
l'énoncé du lemme ; en particulier, on peut prendre $V = U$
si le corps $\kappa(\mathfrak p)$ est de caractéristique zéro.)

\medskip\noindent
Si le corps $k = \kappa(\mathfrak p)$ est de caractéristique zéro, alors
$k' = k$. Si le corps $k = \kappa(\mathfrak p)$
est de caractéristique $p > 0$, alors $p$ s'envoie sur zéro dans $A_\mathfrak p = \kappa(\mathfrak p)$.
Ainsi, après remplacement de $A$ par une localisation principale (c'est-à-dire après
restriction de $S$), on peut supposer que $p = 0$ dans $A$. Si $k' \not = k$, alors
il existe un élément $\beta \in k'$, $\beta \not \in k$,
tel que $\beta^p \in k$. Après remplacement de $A$ par une localisation
principale, on peut supposer qu'il existe un élément $a \in A$ tel
que $\beta^p = a$. Posons $A' = A[x]/(x^p - a)$.
Alors $S' = \Spec(A') \to \Spec(A) = S$ est fini, surjectif et
universellement injectif. De plus, si $\mathfrak p' \subset A'$
désigne l'unique idéal premier au-dessus de $\mathfrak p$,
alors $A'_{\mathfrak p'} = k(\beta)$ et $k'/k(\beta)$
est de degré plus petit. Ainsi, après remplacement de $S$ par $S'$
et de $\eta$ par le point $\eta'$ correspondant à $\mathfrak p'$,
on voit que le degré de $k'$ sur le corps résiduel de $\eta$
a diminué. En poursuivant ainsi, une récurrence nous ramène
au cas $k' = \kappa(\mathfrak p) = \kappa(\eta)$.

\medskip\noindent
On peut donc supposer que $S$ est affine et réduit, et que l'on dispose d'un diagramme
$$
\xymatrix{
\overline{Y}_{1, \eta} \amalg \ldots \amalg \overline{Y}_{n, \eta} \ar[rd] &
Y_{1, \eta} \amalg \ldots \amalg Y_{n, \eta} \ar[r]_-\nu \ar[d] \ar[l]^j &
X_\eta \ar[d] \\
& \Spec(k_1) \amalg \ldots \amalg \Spec(k_n) \ar[r] &
\eta
}
$$
de schémas possédant les propriétés suivantes :
\begin{enumerate}
\item $\nu$ est la normalisation de $X_\eta$,
\item $j$ est une immersion ouverte d'image dense,
\item $k_i/\kappa(\eta)$
est une extension finie séparable pour $i = 1, \ldots, n$,
\item $\overline{Y}_{i, \eta}$ est lisse, projectif et
géométriquement irréductible de dimension $\leq 1$ sur $k_i$.
\end{enumerate}
Rappelons que $\kappa(\eta) = \kappa(\mathfrak p) = A_\mathfrak p$ est
la limite inductive filtrante des $A_a$ pour $a \in A$, $a \not \in \mathfrak p$.
Voir Algèbre, lemme \ref{algebra-lemma-localization-colimit}.
On peut donc faire descendre le diagramme ci-dessus en un diagramme correspondant
sur $\Spec(A_a)$ pour un certain $a \in A$, $a \not \in \mathfrak p$.
Plus précisément, après remplacement de $S$ par $\Spec(A_a)$,
on peut supposer que l'on dispose d'un diagramme commutatif
$$
\xymatrix{
\overline{Y}_1 \amalg \ldots \amalg \overline{Y}_n \ar[rd] &
Y_1 \amalg \ldots \amalg Y_n \ar[r]_-\nu \ar[d] \ar[l]^j &
X \ar[d] \\
& T_1 \amalg \ldots \amalg T_n \ar[r] & S
}
$$
de schémas dont le changement de base à $\eta$ est le diagramme ci-dessus
et qui possèdent les propriétés suivantes :
\begin{enumerate}
\item $\nu$ est un morphisme fini et surjectif,
\item $j$ est une immersion ouverte,
\item $T_i \to S$ est fini \'etale pour $i = 1, \ldots, n$,
\item $Y_i \to T_i$ est lisse et surjectif,
\item $\overline{Y}_i \to T_i$ est lisse et propre, à fibres
géométriquement connexes de dimension $\leq 1$.
\end{enumerate}
Pour cela, on utilise d'abord
Limites, lemme \ref{limits-lemma-descend-finite-presentation},
afin d'obtenir un diagramme dont le changement de base redonne le diagramme précédent.
On utilise ensuite Limites, lemmes \ref{limits-lemma-descend-etale},
\ref{limits-lemma-descend-smooth},
\ref{limits-lemma-descend-finite-finite-presentation},
\ref{limits-lemma-limit-affine},
\ref{limits-lemma-descend-open-immersion},
\ref{limits-lemma-eventually-proper}, et
\ref{limits-lemma-descend-surjective},
pour obtenir que $\nu$ est fini et surjectif, que $j$ est une immersion ouverte, que $T_i \to S$ est fini
\'etale, que $Y_i \to T$ est lisse, et que $\overline{Y}_i \to T_i$ est propre et
lisse. Puisque $Y_i$ ne peut être vide, puisque les morphismes lisses
sont ouverts, et puisque $T_i \to S$ est fini \'etale, quitte à rétrécir $S$,
on peut supposer que $Y_i \to T_i$ est surjectif. Enfin, la fibre de
$\overline{Y}_i \to T_i$ au-dessus de l'unique point $\eta_i = \Spec(k_i)$ 
de $T_i$ au-dessus de $\eta$ est géométriquement connexe.
Par un nouveau rétrécissement, on peut donc supposer qu'il en va de même
pour toutes les fibres, voir
lemme \ref{more-morphisms-lemma-proper-flat-geom-red}.

\medskip\noindent
Il reste à démontrer l'existence d'un ouvert $W \subset X$
satisfaisant (a), (b) et (c). Puisque $\nu_\eta : \coprod Y_{i, \eta} \to X_\eta$
est le morphisme de normalisation, on sait d'après
Variétés, lemme \ref{varieties-lemma-normalization-locally-algebraic},
qu'il existe un ouvert dense $W_\eta \subset X_\eta$
tel que $\nu^{-1}(W_\eta) \to W_\eta$ soit égal à
l'inclusion du réduit de $W_\eta$ dans $W_\eta$.
Soit $W \subset X$ un ouvert quasi-compact dont la fibre
au-dessus de $\eta$ est l'ouvert $W_\eta$ que l'on vient de trouver.
Après remplacement de $A = \Gamma(S, \mathcal{O}_S)$
par une autre localisation, on peut supposer que $\nu^{-1}(W) \to W$
est une immersion fermée, voir Limites, lemme
\ref{limits-lemma-descend-closed-immersion-finite-presentation}.
Comme $\nu$ est aussi surjectif, on en conclut que
$\nu^{-1}(W) \to W$ est un épaississement.
Posons $W_i = \nu^{-1}(W) \cap Y_i$.
Quitte à rétrécir encore $S$, on peut supposer que $W_i \to T_i$ est
surjectif pour tout $i$ (par le même argument que ci-dessus).
On en déduit que $W_i \subset Y_i$
est dense dans toutes les fibres de $Y_i \to T_i$, puisque $Y_i \to T_i$ a des
fibres géométriquement irréductibles.
Comme $\nu$ est fini et surjectif, il en résulte alors
que $W = \nu(\nu^{-1}(W))$ est également dense dans toutes les fibres
de $X \to S$.
\end{proof}











\section{Descente des morphismes séparés localement quasi-finis}
\label{more-morphisms-section-separated-locally-quasi-finite}

\noindent
Dans cette section, nous montrons que « les morphismes séparés localement quasi-finis
satisfont à la descente pour les recouvrements fppf ». Voir Descente, définition
\ref{descent-definition-descending-types-morphisms} pour la terminologie.
Ce résultat se trouve dans le remarquable article de Raynaud et Gruson
(à bien des égards), dissimulé dans la démonstration
de \cite[Lemme 5.7.1]{GruRay}.
On le trouve également dans \cite{Murre-representation}, et dans
\cite[Expos\'e X, Lemme 5.4]{SGA3},
sous l'hypothèse supplémentaire
que le morphisme est localement de présentation finie.
Voici l'énoncé formel.

\begin{lemma}
\label{more-morphisms-lemma-separated-locally-quasi-finite-morphisms-fppf-descend}
Soit $S$ un schéma.
Soit $\{X_i \to S\}_{i\in I}$ un recouvrement fppf, voir
Topologies, définition \ref{topologies-definition-fppf-covering}.
Soit $(V_i/X_i, \varphi_{ij})$ une donnée de descente
relative à $\{X_i \to S\}$. Si chaque morphisme
$V_i \to X_i$ est séparé et localement quasi-fini,
alors la donnée de descente est effective.
\end{lemma}

\begin{proof}
Être séparé et être localement quasi-fini
sont des propriétés des morphismes de schémas
qui sont préservées par tout changement de base, voir
Schémas, lemme \ref{schemes-lemma-separated-permanence} et
Morphismes, lemme \ref{morphisms-lemma-base-change-quasi-finite}.
Par conséquent, Descente, lemme \ref{descent-lemma-descending-types-morphisms},
s'applique, et il suffit de démontrer l'assertion du lemme
lorsque le recouvrement fppf est donné par un unique
morphisme $\{X \to S\}$ plat, surjectif et de présentation finie entre schémas affines.
Écrivons $X = \Spec(A)$ et $S = \Spec(R)$, de sorte
que $R \to A$ soit un morphisme d'anneaux fidèlement plat.
Soit $(V, \varphi)$ une donnée de descente relativement à $X$ sur $S$
et supposons que $\pi : V \to X$ soit séparé et
localement quasi-fini.

\medskip\noindent
Soit $W^1 \subset V$ un ouvert affine quelconque.
Considérons $W = \text{pr}_1(\varphi(W^1 \times_S X)) \subset V$.
Voici un diagramme
$$
\xymatrix{
W^1 \times_S X \ar[rrrrr] \ar[ddd] \ar[rd]
& & & & &
\varphi(W^1 \times_S X) \ar[ddd] \ar[ld] \\
& V \times_S X \ar[rrr]^\varphi \ar[rd] \ar[dd]
& & &
X \times_S V \ar[ld] \ar[dd] & \\
& &
X \times_S X \ar[r]^1 \ar[d]_{\text{pr}_0}
&
X \times_S X \ar[d]^{\text{pr}_1}
& & \\
W^1 \ar[r] &
V \ar[r] &
X &
X &
V \ar[l] &
W \ar[l]
}
$$
Or, puisque $X \to S$ est plat et de présentation finie, il
est universellement ouvert (Morphismes, lemme \ref{morphisms-lemma-fppf-open}).
Nous en déduisons que $W$ est ouvert. De plus, il est
également clair que $W$ est quasi-compact, et que
$W^1 \subset W$. En outre, nous remarquons que
$\varphi(W \times_S X) = X \times_S W$ par la condition de cocycle
pour $\varphi$. Nous obtenons donc une nouvelle donnée de descente
$(W, \varphi')$ en restreignant $\varphi$ à $W \times_S X$.
Remarquons que le morphisme $W \to X$ est quasi-compact, séparé
et localement quasi-fini. Cela implique qu'il est
séparé et quasi-fini par définition. Il est donc quasi-affine d'après le
lemme \ref{more-morphisms-lemma-quasi-finite-separated-quasi-affine}.
Ainsi, d'après
Descente, lemme \ref{descent-lemma-quasi-affine},
la donnée de descente
$(W, \varphi')$ est effective.

\medskip\noindent
Autrement dit, on trouve un recouvrement ouvert
$V = \bigcup W_i$ par des ouverts quasi-compacts $W_i$ qui sont
stables par le morphisme de descente $\varphi$.
De plus, pour chacun de ces ouverts quasi-compacts $W \subset V$,
la donnée de descente correspondante $(W, \varphi')$ est effective.
Cela signifie que la donnée de descente initiale est effective, par recollement des
schémas obtenus en descendant les ouverts $W_i$ ; voir
Descente, lemme \ref{descent-lemma-effective-for-fpqc-is-local-upstairs}.
\end{proof}






\section{Présentation finie relative}
\label{more-morphisms-section-finite-type-finite-presentation}

\noindent
Soit $R \to A$ un homomorphisme d'anneaux de type fini. Soit $M$ un $A$-module. Dans
Compléments d'algèbre,
section \ref{more-algebra-section-relative-finite-presentation},
nous avons défini ce que signifie pour $M$ d'être de présentation finie relativement à $R$.
Nous avons aussi démontré que cette notion se localise bien et se recolle.
Nous pouvons donc définir la notion globale correspondante comme suit.

\begin{definition}
\label{more-morphisms-definition-relatively-finitely-presented-sheaf}
Soit $f : X \to S$ un morphisme de schémas localement de type fini.
Soit $\mathcal{F}$ un $\mathcal{O}_X$-module quasi-cohérent. On dit que
$\mathcal{F}$ est {\it finiment présenté relativement à $S$} ou
{\it de présentation finie relativement à $S$}
s'il existe un recouvrement par des ouverts affines $S = \bigcup V_i$ et,
pour tout $i$, un recouvrement par des ouverts affines
$f^{-1}(V_i) = \bigcup_j U_{ij}$ tels que $\mathcal{F}(U_{ij})$
soit un $\mathcal{O}_X(U_{ij})$-module de présentation finie relativement
à $\mathcal{O}_S(V_i)$.
\end{definition}

\noindent
Remarquons que cela implique que $\mathcal{F}$ est un
$\mathcal{O}_X$-module de type fini. Si $X \to S$ est seulement localement de type
fini, alors $\mathcal{F}$ peut être de présentation finie relativement
à $S$, sans que $X \to S$ soit localement de présentation finie.
Nous verrons que $X \to S$ est localement de présentation finie si
et seulement si $\mathcal{O}_X$ est de présentation finie relativement à $S$.

\begin{lemma}
\label{more-morphisms-lemma-relative-finite-presentation-characterize}
Soit $f : X \to S$ un morphisme de schémas localement de type fini.
Soit $\mathcal{F}$ un $\mathcal{O}_X$-module quasi-cohérent. Les conditions suivantes
sont équivalentes :
\begin{enumerate}
\item $\mathcal{F}$ est de présentation finie relativement à $S$,
\item pour tous ouverts affines $U \subset X$, $V \subset S$
tels que $f(U) \subset V$, le $\mathcal{O}_X(U)$-module $\mathcal{F}(U)$
est de présentation finie relativement à $\mathcal{O}_S(V)$.
\end{enumerate}
De plus, si ces conditions sont satisfaites, alors, pour tous sous-schémas ouverts
$U \subset X$ et $V \subset S$ tels que $f(U) \subset V$,
la restriction $\mathcal{F}|_U$ est de présentation finie relativement à $V$.
\end{lemma}

\begin{proof}
La dernière assertion résulte clairement de l'équivalence de (1) et (2).
Il est également clair que (2) implique (1). Supposons (1) satisfaite.
Soient $S = \bigcup V_i$ et $f^{-1}(V_i) = \bigcup U_{ij}$ des
recouvrements par des ouverts affines comme dans la
définition \ref{more-morphisms-definition-relatively-finitely-presented-sheaf}.
Soient $U \subset X$ et $V \subset S$ comme dans (2).
D'après Compléments d'algèbre, lemme
\ref{more-algebra-lemma-glue-relative-finite-presentation},
il suffit de trouver un recouvrement $U = \bigcup U_k$ de $U$ par des ouverts principaux
tel que $\mathcal{F}(U_k)$ soit de présentation finie relativement à
$\mathcal{O}_S(V)$. Autrement dit, pour tout $u \in U$, il suffit
de trouver un ouvert affine principal $u \in U' \subset U$ tel que
$\mathcal{F}(U')$ soit de présentation finie relativement à $\mathcal{O}_S(V)$.
Choisissons $i$ tel que $f(u) \in V_i$, puis $j$ tel que
$u \in U_{ij}$. D'après
Schémas, lemme \ref{schemes-lemma-standard-open-two-affines},
on peut trouver $v \in V' \subset V \cap V_i$ qui est un ouvert affine principal
dans $V'$ et dans $V_i$. Alors $f^{-1}V'  \cap U$ et, respectivement,\ $f^{-1}V' \cap U_{ij}$
sont des ouverts affines principaux de $U$ et, respectivement,\ de $U_{ij}$.
En appliquant de nouveau le lemme, on peut trouver
$u \in U' \subset f^{-1}V' \cap U \cap U_{ij}$ qui est un ouvert affine principal
à la fois dans $f^{-1}V'  \cap U$ et dans $f^{-1}V' \cap U_{ij}$.
Ainsi, $U'$ est aussi un ouvert affine principal de $U$ et de $U_{ij}$.
D'après Compléments d'algèbre, lemme
\ref{more-algebra-lemma-localize-relative-finite-presentation},
l'hypothèse selon laquelle $\mathcal{F}(U_{ij})$ est de présentation finie
relativement à $\mathcal{O}_S(V_i)$ implique que
$\mathcal{F}(U')$ est de présentation finie relativement à $\mathcal{O}_S(V_i)$.
Puisque $\mathcal{O}_X(U') =
\mathcal{O}_X(U') \otimes_{\mathcal{O}_S(V_i)} \mathcal{O}_S(V')$,
Compléments d'algèbre, lemme
\ref{more-algebra-lemma-base-change-relative-finite-presentation},
montre que $\mathcal{F}(U')$ est de présentation finie relativement à $\mathcal{O}_S(V')$.
En appliquant encore une fois Compléments d'algèbre, lemme
\ref{more-algebra-lemma-localize-relative-finite-presentation},
nous concluons que
$\mathcal{F}(U')$ est de présentation finie relativement à $\mathcal{O}_S(V)$.
Ceci achève la démonstration.
\end{proof}

\begin{lemma}
\label{more-morphisms-lemma-relative-finite-presentation}
Soit $f : X \to S$ un morphisme de schémas localement de type
fini. Soit $\mathcal{F}$ un $\mathcal{O}_X$-module quasi-cohérent.
\begin{enumerate}
\item Si $f$ est localement de présentation finie, alors $\mathcal{F}$
est de présentation finie relativement à $S$ si et seulement si $\mathcal{F}$
est de présentation finie.
\item Le morphisme $f$ est localement de présentation finie si et seulement
si $\mathcal{O}_X$ est de présentation finie relativement à $S$.
\end{enumerate}
\end{lemma}

\begin{proof}
Cela résulte immédiatement des définitions ; voir la
discussion qui suit la définition dans
Compléments d'algèbre, définition
\ref{more-algebra-definition-relatively-finitely-presented}.
\end{proof}

\begin{lemma}
\label{more-morphisms-lemma-finite-morphism-relative-finite-presentation}
Soit $\pi : X \to Y$ un morphisme fini entre des schémas localement de type
fini sur un schéma de base $S$. Soit $\mathcal{F}$ un
$\mathcal{O}_X$-module quasi-cohérent. Alors $\mathcal{F}$ est de présentation finie
relativement à $S$ si et seulement si $\pi_*\mathcal{F}$ est de présentation finie
relativement à $S$.
\end{lemma}

\begin{proof}
C'est la traduction du résultat de
Compléments d'algèbre, lemme \ref{more-algebra-lemma-finite-extension},
dans le langage des schémas.
\end{proof}

\begin{lemma}
\label{more-morphisms-lemma-base-change-relative-finite-presentation}
Soit $f : X \to S$ un morphisme de schémas localement de type
fini. Soit $\mathcal{F}$ un $\mathcal{O}_X$-module quasi-cohérent.
Soit $S' \to S$ un morphisme de schémas, posons $X' = X \times_S S'$
et notons $\mathcal{F}'$ l'image inverse de $\mathcal{F}$ sur $X'$.
Si $\mathcal{F}$ est de présentation finie relativement à $S$, alors
$\mathcal{F}'$ est de présentation finie relativement à $S'$.
\end{lemma}

\begin{proof}
C'est la traduction du résultat de
Compléments d'algèbre, lemme
\ref{more-algebra-lemma-base-change-relative-finite-presentation},
dans le langage des schémas.
\end{proof}

\begin{lemma}
\label{more-morphisms-lemma-pull-relative-finite-presentation}
Soient $X \to Y \to S$ des morphismes de schémas localement de type
fini. Soit $\mathcal{G}$ un $\mathcal{O}_Y$-module quasi-cohérent.
Si $f : X \to Y$ est localement de présentation finie et si
$\mathcal{G}$ est de présentation finie relativement à $S$, alors
$f^*\mathcal{G}$ est de présentation finie relativement à $S$.
\end{lemma}

\begin{proof}
C'est la traduction du résultat de
Compléments d'algèbre, lemme
\ref{more-algebra-lemma-pull-relative-finite-presentation},
dans le langage des schémas.
\end{proof}

\begin{lemma}
\label{more-morphisms-lemma-composition-relative-finite-presentation}
Soient $X \to Y \to S$ des morphismes de schémas localement de type
fini. Soit $\mathcal{F}$ un $\mathcal{O}_X$-module quasi-cohérent.
Si $Y \to S$ est localement de présentation finie et si $\mathcal{F}$
est de présentation finie relativement à $Y$, alors $\mathcal{F}$
est de présentation finie relativement à $S$.
\end{lemma}

\begin{proof}
C'est la traduction du résultat de
Compléments d'algèbre, lemme
\ref{more-algebra-lemma-composition-relative-finite-presentation},
dans le langage des schémas.
\end{proof}

\begin{lemma}
\label{more-morphisms-lemma-ses-relatively-finite-presentation}
Soit $X \to S$ un morphisme de schémas localement de type fini.
Soit $0 \to \mathcal{F}' \to \mathcal{F} \to \mathcal{F}'' \to 0$
une suite exacte courte de $\mathcal{O}_X$-modules quasi-cohérents.
\begin{enumerate}
\item Si $\mathcal{F}', \mathcal{F}''$ sont de présentation finie relativement à
$S$, alors il en est de même de $\mathcal{F}$.
\item Si $\mathcal{F}'$ est un $\mathcal{O}_X$-module de type fini
et si $\mathcal{F}$ est de présentation finie relativement à $S$, alors
$\mathcal{F}''$ est de présentation finie relativement à $S$.
\end{enumerate}
\end{lemma}

\begin{proof}
C'est la traduction du résultat de
Compléments d'algèbre, lemme
\ref{more-algebra-lemma-ses-relatively-finite-presentation},
dans le langage des schémas.
\end{proof}

\begin{lemma}
\label{more-morphisms-lemma-sum-relatively-finite-presentation}
\begin{slogan}
Un facteur direct d'un module hérite de la propriété d'être de présentation
finie relativement à une base.
\end{slogan}
Soit $X \to S$ un morphisme de schémas localement de type fini.
Soient $\mathcal{F}, \mathcal{F}'$ des $\mathcal{O}_X$-modules quasi-cohérents.
Si $\mathcal{F} \oplus \mathcal{F}'$ est de présentation finie relativement à $S$,
alors il en est de même de $\mathcal{F}$ et de $\mathcal{F}'$.
\end{lemma}

\begin{proof}
C'est la traduction du résultat de
Compléments d'algèbre, lemme
\ref{more-algebra-lemma-sum-relatively-finite-presentation},
dans le langage des schémas.
\end{proof}










\section{Pseudo-cohérence relative}
\label{more-morphisms-section-relative-pseudo-coherence}

\noindent
Cette section est, pour les schémas, l'analogue de
Compléments d'algèbre, section \ref{more-algebra-section-relative-pseudo-coherent}.
Nous recommandons vivement au lecteur de consulter d'abord cette
section. Bien que nous ayons développé les résultats de cette section
ainsi que ceux sur les complexes pseudo-cohérents dans
Cohomologie, sections \ref{cohomology-section-strictly-perfect},
\ref{cohomology-section-pseudo-coherent},
\ref{cohomology-section-tor} et
\ref{cohomology-section-perfect},
pour des complexes arbitraires de $\mathcal{O}_X$-modules, si
$X$ est un schéma, travailler exclusivement avec des objets de
$D_\QCoh(\mathcal{O}_X)$
simplifie considérablement de nombreux lemmes et arguments, en
ramenant souvent d'emblée le problème à son analogue algébrique.
De plus, l'une des premières choses que nous faisons est de montrer que
la pseudo-cohérence relative implique que les faisceaux de cohomologie
sont quasi-cohérents ; voir le lemme \ref{more-morphisms-lemma-relative-pseudo-coherence}.
Pour une première lecture, nous suggérons donc au lecteur de travailler exclusivement
avec des objets de $D_\QCoh(\mathcal{O}_X)$.

\begin{lemma}
\label{more-morphisms-lemma-relatively-pseudo-coherent}
Soit $X \to S$ un morphisme de type fini de schémas affines.
Soit $E$ un objet de $D(\mathcal{O}_X)$.
Soit $m \in \mathbf{Z}$.
Les conditions suivantes sont équivalentes :
\begin{enumerate}
\item pour une immersion fermée $i : X \to \mathbf{A}^n_S$,
l'objet $Ri_*E$ de $D(\mathcal{O}_{\mathbf{A}^n_S})$
est $m$-pseudo-cohérent, et
\item pour toute immersion fermée $i : X \to \mathbf{A}^n_S$,
l'objet $Ri_*E$ de $D(\mathcal{O}_{\mathbf{A}^n_S})$
est $m$-pseudo-cohérent.
\end{enumerate}
\end{lemma}

\begin{proof}
Écrivons $S = \Spec(R)$ et $X = \Spec(A)$. Soit $i$ l'immersion correspondant à la surjection
$\alpha : R[x_1, \ldots, x_n] \to A$ et soit $X \to \mathbf{A}^m_S$
correspondant à $\beta : R[y_1, \ldots, y_m] \to A$.
Choisissons $f_j \in R[x_1, \ldots, x_n]$ tels que $\alpha(f_j) = \beta(y_j)$
et $g_i \in R[y_1, \ldots, y_m]$ tels que $\beta(g_i) = \alpha(x_i)$.
On obtient alors un diagramme commutatif
$$
\xymatrix{
R[x_1, \ldots, x_n, y_1, \ldots, y_m]
\ar[d]^{x_i \mapsto g_i} \ar[rr]_-{y_j \mapsto f_j} & &
R[x_1, \ldots, x_n] \ar[d] \\
R[y_1, \ldots, y_m] \ar[rr] & & A
}
$$
qui correspond au diagramme commutatif d'immersions fermées
$$
\xymatrix{
\mathbf{A}^{n + m}_S & \mathbf{A}^n_S \ar[l] \\
\mathbf{A}^m_S \ar[u] & X \ar[u] \ar[l]
}
$$
Il suffit donc de montrer que, pour une immersion fermée
$$
f : \mathbf{A}^m_S \to \mathbf{A}^{n + m}_S
$$
un objet $E$ de $D(\mathcal{O}_{\mathbf{A}^m_S})$ est
$m$-pseudo-cohérent si et seulement si $Rf_*E$ est $m$-pseudo-cohérent.
Cela résulte de
Catégories dérivées de schémas, lemme
\ref{perfect-lemma-closed-push-pseudo-coherent},
et du fait que $f_*\mathcal{O}_{\mathbf{A}^m_S}$ est
un $\mathcal{O}_{\mathbf{A}^{n + m}_S}$-module pseudo-cohérent.
La pseudo-cohérence de $f_*\mathcal{O}_{\mathbf{A}^m_S}$
se démontre directement sans difficulté, mais résulte aussi de
Catégories dérivées de schémas, lemme \ref{perfect-lemma-pseudo-coherent-affine},
et de
Compléments d'algèbre, lemme \ref{more-algebra-lemma-relatively-pseudo-coherent}.
\end{proof}

\noindent
Rappelons que, si $f : X \to S$ est un morphisme de schémas localement de
type fini, alors, pour toute paire d'ouverts affines $U \subset X$ et
$V \subset S$ tels que $f(U) \subset V$, l'homomorphisme d'anneaux
$\mathcal{O}_S(V) \to \mathcal{O}_X(U)$ est de type fini
(Morphismes, lemme \ref{morphisms-lemma-locally-finite-type-characterize}).
Il existe donc toujours des
immersions fermées $U \to \mathbf{A}^n_V$, et la définition suivante a un sens.

\begin{definition}
\label{more-morphisms-definition-relative-pseudo-coherence}
Soit $f : X \to S$ un morphisme de schémas localement de type fini.
Soit $E$ un objet de $D(\mathcal{O}_X)$. Soit $\mathcal{F}$ un
$\mathcal{O}_X$-module. Fixons $m \in \mathbf{Z}$.
\begin{enumerate}
\item On dit que $E$ est {\it $m$-pseudo-cohérent relativement à $S$}
s'il existe un recouvrement par des ouverts affines $S = \bigcup V_i$ et,
pour tout $i$, un recouvrement par des ouverts affines $f^{-1}(V_i) = \bigcup U_{ij}$
tel que les conditions équivalentes du
lemme \ref{more-morphisms-lemma-relatively-pseudo-coherent}
soient satisfaites pour chacune des paires $(U_{ij} \to V_i, E|_{U_{ij}})$.
\item On dit que $E$ est {\it pseudo-cohérent relativement à $S$}
si $E$ est $m$-pseudo-cohérent relativement à $S$ pour tout $m \in \mathbf{Z}$.
\item On dit que $\mathcal{F}$ est {\it $m$-pseudo-cohérent relativement à $S$} si
$\mathcal{F}$, considéré comme un objet de $D(\mathcal{O}_X)$, est
$m$-pseudo-cohérent relativement à $S$.
\item On dit que $\mathcal{F}$ est {\it pseudo-cohérent relativement à $S$} si
$\mathcal{F}$, considéré comme un objet de $D(\mathcal{O}_X)$, est
pseudo-cohérent relativement à $S$.
\end{enumerate}
\end{definition}

\noindent
Si $X$ est quasi-compact et si $E$ est $m$-pseudo-cohérent relativement à $S$
pour un certain $m$, alors $E$ est borné supérieurement. Si $E$ est
pseudo-cohérent relativement à $S$, alors $E$ a des faisceaux de cohomologie
quasi-cohérents.

\begin{lemma}
\label{more-morphisms-lemma-relative-pseudo-coherence}
Soit $f : X \to S$ un morphisme de schémas localement de type fini.
Si $E$ dans $D(\mathcal{O}_X)$ est $m$-pseudo-cohérent relativement à $S$,
alors $H^i(E)$ est un $\mathcal{O}_X$-module quasi-cohérent pour $i > m$.
Si $E$ est pseudo-cohérent relativement à $S$, alors $E$ est un objet de
$D_\QCoh(\mathcal{O}_X)$.
\end{lemma}

\begin{proof}
Choisissons un recouvrement par des ouverts affines $S = \bigcup V_i$ et,
pour tout $i$, un recouvrement par des ouverts affines $f^{-1}(V_i) = \bigcup U_{ij}$
tel que les conditions équivalentes du
lemme \ref{more-morphisms-lemma-relatively-pseudo-coherent}
soient satisfaites pour chacune des paires $(U_{ij} \to V_i, E|_{U_{ij}})$.
La quasi-cohérence étant une propriété locale sur $X$, nous pouvons supposer
qu'il existe une immersion fermée $i : X \to \mathbf{A}^n_S$
telle que $Ri_*E$ soit $m$-pseudo-cohérent sur $\mathbf{A}^n_S$.
D'après Catégories dérivées de schémas, lemme \ref{perfect-lemma-pseudo-coherent},
cela signifie que $H^q(Ri_*E)$ est quasi-cohérent pour $q > m$.
Puisque $i_*$ est un foncteur exact, on a $i_*H^q(E) = H^q(Ri_*E)$,
qui est quasi-cohérent sur $\mathbf{A}^n_S$.
D'après Morphismes, lemme \ref{morphisms-lemma-i-star-equivalence},
cela implique que $H^q(E)$ est quasi-cohérent, comme voulu
(strictement parlant, cela implique qu'il existe un
$\mathcal{O}_X$-module quasi-cohérent $\mathcal{F}$ tel que
$i_*\mathcal{F} = i_*H^q(E)$, puis, d'après
Modules, lemme \ref{modules-lemma-i-star-equivalence},
on a $\mathcal{F} \cong H^q(E)$, d'où le résultat).
\end{proof}

\noindent
Montrons maintenant que la condition de pseudo-cohérence relative se localise bien.

\begin{lemma}
\label{more-morphisms-lemma-localize-relative-pseudo-coherent}
Soit $S$ un schéma affine. Soit $V \subset S$ un ouvert principal.
Soit $X \to V$ un morphisme de type fini de schémas affines.
Soit $U \subset X$ un ouvert affine. Soit $E$ un objet de
$D(\mathcal{O}_X)$. Si les conditions équivalentes du
lemme \ref{more-morphisms-lemma-relatively-pseudo-coherent}
sont satisfaites pour la paire $(X \to V, E)$, alors
les conditions équivalentes du
lemme \ref{more-morphisms-lemma-relatively-pseudo-coherent}
sont satisfaites pour la paire $(U \to S, E|_U)$.
\end{lemma}

\begin{proof}
Écrivons $S = \Spec(R)$, $V = D(f)$, $X = \Spec(A)$ et $U = D(g)$.
Supposons que les conditions équivalentes du
lemme \ref{more-morphisms-lemma-relatively-pseudo-coherent}
soient satisfaites pour la paire $(X \to V, E)$.

\medskip\noindent
Choisissons une surjection $R_f[x_1, \ldots, x_n] \to A$. Écrivons
$R_f = R[x_0]/(fx_0 - 1)$. Alors $R[x_0, x_1, \ldots, x_n] \to A$
est surjectif, et $R_f[x_1, \ldots, x_n]$ est pseudo-cohérent comme
$R[x_0, \ldots, x_n]$-module. Nous avons donc
$$
X \to \mathbf{A}^n_V \to \mathbf{A}^{n + 1}_S
$$
et nous pouvons appliquer
Catégories dérivées de schémas,
lemme \ref{perfect-lemma-closed-push-pseudo-coherent},
pour conclure que l'image directe $E'$ de $E$ sur $\mathbf{A}^{n + 1}_S$
est $m$-pseudo-cohérente.

\medskip\noindent
Choisissons un élément $g' \in R[x_0, x_1, \ldots, x_n]$ dont l'image est
$g \in A$. Considérons la surjection
$R[x_0, \ldots, x_{n + 1}] \to R[x_0, \ldots, x_n, 1/g']$.
On obtient
$$
\xymatrix{
X \ar[d] & U \ar[d] \ar[l] \ar[dr] \\
\mathbf{A}^{n + 1}_S & D(g')\ar[l] \ar[r] & \mathbf{A}^{n + 2}_S
}
$$
où la flèche inférieure gauche est une immersion ouverte et la flèche inférieure droite
une immersion fermée. On conclut comme précédemment que l'image directe de
$E'|_{D(g')}$ sur $\mathbf{A}^{n + 2}_S$ est $m$-pseudo-cohérente.
Comme c'est également l'image directe de $E|_U$ sur $\mathbf{A}^{n + 2}_S$,
ce qui démontre le lemme.
\end{proof}

\begin{lemma}
\label{more-morphisms-lemma-glue-relative-pseudo-coherent}
Soit $X \to S$ un morphisme de type fini de schémas affines. Soit $E$ un
objet de $D(\mathcal{O}_X)$. Soit $m \in \mathbf{Z}$.
Soit $X = \bigcup U_i$ un recouvrement par des ouverts affines principaux.
Les conditions suivantes sont équivalentes :
\begin{enumerate}
\item les conditions équivalentes du
lemme \ref{more-morphisms-lemma-relatively-pseudo-coherent}
sont satisfaites pour les paires $(U_i \to S, E|_{U_i})$,
\item les conditions équivalentes du
lemme \ref{more-morphisms-lemma-relatively-pseudo-coherent}
sont satisfaites pour la paire $(X \to S, E)$.
\end{enumerate}
\end{lemma}

\begin{proof}
L'implication (2) $\Rightarrow$ (1) est le
lemme \ref{more-morphisms-lemma-localize-relative-pseudo-coherent}.
Supposons (1). Écrivons $S = \Spec(R)$ et $X = \Spec(A)$, et
$U_i = D(f_i)$. Écrivons $1 = \sum f_ig_i$ dans $A$.
Considérons les surjections
$$
R[x_i, y_i, z_i] \to R[x_i, y_i, z_i]/(\sum y_iz_i - 1) \to A.
$$
La composée envoie $y_i$ sur $f_i$ et $z_i$ sur $g_i$. Remarquons que
$R[x_i, y_i, z_i]/(\sum y_iz_i - 1)$ est pseudo-cohérent comme
$R[x_i, y_i, z_i]$-module. Il suffit donc de démontrer que
l'image directe de $E$ sur $T = \Spec(R[x_i, y_i, z_i]/(\sum y_iz_i - 1))$
est $m$-pseudo-cohérente ; voir
Catégories dérivées de schémas,
lemme \ref{perfect-lemma-closed-push-pseudo-coherent}.
Pour chaque $i_0$, il suffit de démontrer que la restriction de cette
image directe à
$W_{i_0} = \Spec(R[x_i, y_i, z_i, 1/y_{i_0}]/(\sum y_iz_i - 1))$
est $m$-pseudo-cohérente. Remarquons qu'il existe un diagramme commutatif
$$
\xymatrix{
X \ar[d] & U_{i_0} \ar[l] \ar[d] \\
T & W_{i_0} \ar[l]
}
$$
qui implique que l'image directe de $E$ sur $T$, restreinte à $W_{i_0}$,
est l'image directe de $E|_{U_{i_0}}$ sur $W_{i_0}$. Puisque
$R[x_i, y_i, z_i, 1/y_{i_0}]/(\sum y_iz_i - 1)$ est isomorphe
à un anneau de polynômes sur $R$, cela démontre ce que nous voulons.
\end{proof}

\begin{lemma}
\label{more-morphisms-lemma-relative-pseudo-coherence-characterize}
Soit $f : X \to S$ un morphisme de schémas localement de type fini.
Soit $E$ un objet de $D(\mathcal{O}_X)$.
Fixons $m \in \mathbf{Z}$. Les conditions suivantes sont équivalentes :
\begin{enumerate}
\item $E$ est $m$-pseudo-cohérent relativement à $S$,
\item pour tous ouverts affines $U \subset X$ et $V \subset S$
tels que $f(U) \subset V$, les conditions équivalentes du
lemme \ref{more-morphisms-lemma-relatively-pseudo-coherent}
sont satisfaites pour la paire $(U \to V, E|_U)$.
\end{enumerate}
De plus, si ces conditions sont satisfaites, alors, pour tous sous-schémas ouverts
$U \subset X$ et $V \subset S$ tels que $f(U) \subset V$,
la restriction $E|_U$ est $m$-pseudo-cohérente relativement à $V$.
\end{lemma}

\begin{proof}
La dernière assertion résulte clairement de l'équivalence de (1) et (2).
Il est également clair que (2) implique (1). Supposons (1) satisfaite.
Soient $S = \bigcup V_i$ et $f^{-1}(V_i) = \bigcup U_{ij}$ des
recouvrements par des ouverts affines comme dans la
définition \ref{more-morphisms-definition-relative-pseudo-coherence}.
Soient $U \subset X$ et $V \subset S$ comme dans (2).
D'après le lemme \ref{more-morphisms-lemma-glue-relative-pseudo-coherent},
il suffit de trouver un recouvrement $U = \bigcup U_k$ de $U$ par des ouverts principaux
tel que les conditions équivalentes du
lemme \ref{more-morphisms-lemma-relatively-pseudo-coherent}
soient satisfaites pour les paires $(U_k \to V, E|_{U_k})$.
Autrement dit, pour tout $u \in U$, il suffit
de trouver un ouvert affine principal $u \in U' \subset U$ tel que
les conditions équivalentes du
lemme \ref{more-morphisms-lemma-relatively-pseudo-coherent}
soient satisfaites pour la paire $(U' \to V, E|_{U'})$.
Choisissons $i$ tel que $f(u) \in V_i$, puis $j$ tel que
$u \in U_{ij}$. D'après
Schémas, lemme \ref{schemes-lemma-standard-open-two-affines},
on peut trouver $v \in V' \subset V \cap V_i$ qui est un ouvert affine principal
dans $V'$ et dans $V_i$. Alors $f^{-1}V'  \cap U$ et, respectivement,\ $f^{-1}V' \cap U_{ij}$
sont des ouverts affines principaux de $U$ et, respectivement,\ de $U_{ij}$.
En appliquant de nouveau le lemme, on peut trouver
$u \in U' \subset f^{-1}V' \cap U \cap U_{ij}$ qui est un ouvert affine principal
à la fois dans $f^{-1}V'  \cap U$ et dans $f^{-1}V' \cap U_{ij}$.
Ainsi, $U'$ est aussi un ouvert affine principal de $U$ et de $U_{ij}$.
D'après le lemme \ref{more-morphisms-lemma-localize-relative-pseudo-coherent},
l'hypothèse selon laquelle les conditions équivalentes du
lemme \ref{more-morphisms-lemma-relatively-pseudo-coherent}
sont satisfaites pour la paire $(U_{ij} \to V_i, E|_{U_{ij}})$
implique que les conditions équivalentes du
lemme \ref{more-morphisms-lemma-relatively-pseudo-coherent}
sont satisfaites pour la paire $(U' \to V, E|_{U'})$.
\end{proof}

\noindent
Pour les objets de la catégorie dérivée dont les faisceaux de cohomologie sont
quasi-cohérents, on peut relier la $m$-pseudo-cohérence relative
à la notion définie dans Compléments d'algèbre, définition
\ref{more-algebra-definition-relatively-pseudo-coherent}.
Nous utiliserons le fait que, pour un schéma affine
$U = \Spec(A)$, le foncteur $R\Gamma(U, -)$ induit une équivalence
entre $D_\QCoh(\mathcal{O}_U)$ et $D(A)$ ; voir
Catégories dérivées de schémas, lemme
\ref{perfect-lemma-affine-compare-bounded}.
Ce foncteur commute aux images inverses :
si $E$ est un objet de $D_\QCoh(\mathcal{O}_U)$
et $A \to B$ est un homomorphisme d'anneaux correspondant à un morphisme de
schémas affines $g : V = \Spec(B) \to \Spec(A) = U$, alors
$R\Gamma(V, Lg^*E) = R\Gamma(U, E) \otimes_A^\mathbf{L} B$.
Voir Catégories dérivées de schémas, lemme
\ref{perfect-lemma-quasi-coherence-pullback}.

\begin{lemma}
\label{more-morphisms-lemma-qcoh-relative-pseudo-coherence-characterize}
Soit $f : X \to S$ un morphisme de schémas localement de type fini.
Soit $E$ un objet de $D_\QCoh(\mathcal{O}_X)$.
Fixons $m \in \mathbf{Z}$. Les assertions suivantes sont équivalentes :
\begin{enumerate}
\item $E$ est $m$-pseudo-cohérent relativement à $S$ ;
\item il existe un recouvrement ouvert affine $S = \bigcup V_i$ et,
pour tout $i$, un recouvrement ouvert affine $f^{-1}(V_i) = \bigcup U_{ij}$
tels que le complexe de $\mathcal{O}_X(U_{ij})$-modules
$R\Gamma(U_{ij}, E)$ soit $m$-pseudo-cohérent relativement à
$\mathcal{O}_S(V_i)$ ;
\item pour tous ouverts affines $U \subset X$ et $V \subset S$
tels que $f(U) \subset V$, le complexe de $\mathcal{O}_X(U)$-modules
$R\Gamma(U, E)$ est $m$-pseudo-cohérent relativement à $\mathcal{O}_S(V)$.
\end{enumerate}
\end{lemma}

\begin{proof}
Soient $U$ et $V$ comme dans (2), et choisissons une immersion fermée
$i : U \to \mathbf{A}^n_V$. Un argument formel, utilisant le
lemme \ref{more-morphisms-lemma-relative-pseudo-coherence-characterize}, montre qu'il
suffit de prouver que
$Ri_*(E|_U)$ est $m$-pseudo-cohérent si et seulement si $R\Gamma(U, E)$
est $m$-pseudo-cohérent relativement à $\mathcal{O}_S(V)$.
Écrivons $U = \Spec(A)$, $V = \Spec(R)$ et
$\mathbf{A}^n_V = \Spec(R[x_1, \ldots, x_n]$. D'après les remarques
précédant le lemme, $E|_U$ est quasi-isomorphe au
complexe de faisceaux quasi-cohérents sur $U$ associé à l'objet
$R\Gamma(U, E)$ de $D(A)$. Remarquons que
$R\Gamma(U, E) = R\Gamma(\mathbf{A}^n_V, Ri_*(E|_U))$, car $i$ est une
immersion fermée (et donc $i_*$ est exact). Ainsi, $Ri_*E$
est associé à $R\Gamma(U, E)$ considéré comme un objet de
$D(R[x_1, \ldots, x_n])$. On conclut, puisque la $m$-pseudo-cohérence
de $Ri_*(E|_U)$ équivaut à la $m$-pseudo-cohérence de
$R\Gamma(U, E)$ dans $D(R[x_1, \ldots, x_n])$ d'après
Catégories dérivées de schémas, lemme \ref{perfect-lemma-pseudo-coherent-affine},
ce qui équivaut à dire que $R\Gamma(U, E)$ est $m$-pseudo-cohérent
relativement à $R = \mathcal{O}_S(V)$ par définition.
\end{proof}

\begin{lemma}
\label{more-morphisms-lemma-closed-morphism-relative-pseudo-coherence}
Soit $i : X \to Y$ un morphisme de schémas localement de type fini sur un
schéma de base $S$. Supposons que $i$ induise un homéomorphisme de $X$ sur une partie
fermée de $Y$. Soit $E$ un objet de $D(\mathcal{O}_X)$.
Alors $E$ est $m$-pseudo-cohérent relativement à $S$ si et seulement si
$Ri_*E$ est $m$-pseudo-cohérent relativement à $S$.
\end{lemma}

\begin{proof}
D'après Morphismes, lemme \ref{morphisms-lemma-homeomorphism-affine},
le morphisme $i$ est affine. On peut donc supposer $S$, $Y$ et $X$ affines.
Écrivons $S = \Spec(R)$, $Y = \Spec(A)$ et $X = \Spec(B)$.
La condition signifie que $A/\text{rad}(A) \to B/\text{rad}(B)$ est
surjectif ; ici $\text{rad}(A)$ et $\text{rad}(B)$ désignent
les radicaux de Jacobson de $A$ et de $B$.
Comme $B$ est de type fini sur $A$, on peut trouver
$b_1, \ldots, b_m \in \text{rad}(B)$ qui engendrent $B$ comme
$A$-algèbre. Écrivons $b_j^N = 0$ pour tout $j$. Considérons le diagramme
d'anneaux
$$
\xymatrix{
B & R[x_i, y_j]/(y_j^N) \ar[l] & R[x_i, y_j] \ar[l] \\
A \ar[u] & R[x_i] \ar[l] \ar[u] \ar[ru]
}
$$
qui se traduit par un diagramme
$$
\xymatrix{
X \ar[d] \ar[r] & T \ar[d] \ar[r] & \mathbf{A}^{n + m}_S \ar[ld] \\
Y \ar[r] & \mathbf{A}^n_S
}
$$
de schémas affines. D'après le lemme \ref{more-morphisms-lemma-relative-pseudo-coherence-characterize},
on voit que $E$ est $m$-pseudo-cohérent relativement à $S$ si et seulement si son
image directe dans $\mathbf{A}^{n + m}_S$ est $m$-pseudo-cohérente.
D'après Catégories dérivées de schémas, lemme
\ref{perfect-lemma-closed-push-pseudo-coherent},
on voit que cela a lieu si et seulement si son image directe dans $T$ est
$m$-pseudo-cohérente. Le même lemme montre que cela a lieu si et seulement si
l'image directe dans $\mathbf{A}^n_S$ est $m$-pseudo-cohérente.
De nouveau, d'après le
lemme \ref{more-morphisms-lemma-relative-pseudo-coherence-characterize},
cela a lieu si et seulement si $Ri_*E$ est $m$-pseudo-cohérent relativement à $S$.
\end{proof}

\begin{lemma}
\label{more-morphisms-lemma-finite-morphism-relative-pseudo-coherence}
Soit $\pi : X \to Y$ un morphisme fini de schémas localement de type
fini sur un schéma de base $S$. Soit $E$ un objet de
$D_\QCoh(\mathcal{O}_X)$. Alors $E$ est $m$-pseudo-cohérent
relativement à $S$ si et seulement si $R\pi_*E$ est $m$-pseudo-cohérent
relativement à $S$.
\end{lemma}

\begin{proof}
C'est la traduction du résultat de
Plus sur l'algèbre, lemme
\ref{more-algebra-lemma-finite-extension-pseudo-coherent}
dans le langage des schémas. Observons que $R\pi_*$ envoie bien
$D_\QCoh(\mathcal{O}_X)$ dans $D_\QCoh(\mathcal{O}_Y)$
d'après Catégories dérivées de schémas, lemme
\ref{perfect-lemma-quasi-coherence-direct-image}.
Pour effectuer la traduction, utiliser le
lemme \ref{more-morphisms-lemma-relative-pseudo-coherence-characterize}.
\end{proof}

\begin{lemma}
\label{more-morphisms-lemma-cone-relatively-pseudo-coherent}
Soit $f : X \to S$ un morphisme de schémas localement de type fini.
Soit $(E, E', E'')$ un triangle distingué de
$D(\mathcal{O}_X)$. Soit $m \in \mathbf{Z}$.
\begin{enumerate}
\item Si $E$ est $(m + 1)$-pseudo-cohérent relativement à $S$ et
$E'$ est $m$-pseudo-cohérent relativement à $S$, alors $E''$ est
$m$-pseudo-cohérent relativement à $S$.
\item Si $E, E''$ sont $m$-pseudo-cohérents relativement à $S$,
alors $E'$ est $m$-pseudo-cohérent relativement à $S$.
\item Si $E'$ est $(m + 1)$-pseudo-cohérent relativement à $S$
et $E''$ est $m$-pseudo-cohérent relativement à $S$, alors
$E$ est $(m + 1)$-pseudo-cohérent relativement à $S$.
\end{enumerate}
De plus, si deux des trois objets $E, E', E''$ sont pseudo-cohérents
relativement à $S$, alors le troisième l'est aussi.
\end{lemma}

\begin{proof}
Cela résulte immédiatement du lemme \ref{more-morphisms-lemma-relative-pseudo-coherence-characterize} et de
Cohomologie, lemme \ref{cohomology-lemma-cone-pseudo-coherent}.
\end{proof}

\begin{lemma}
\label{more-morphisms-lemma-rel-n-pseudo-module}
Soit $X \to S$ un morphisme de schémas localement de type fini.
Soit $\mathcal{F}$ un $\mathcal{O}_X$-module. Alors
\begin{enumerate}
\item $\mathcal{F}$ est $m$-pseudo-cohérent relativement à $S$ pour tout $m > 0$ ;
\item $\mathcal{F}$ est $0$-pseudo-cohérent relativement à $S$ si et seulement si
$\mathcal{F}$ est un $\mathcal{O}_X$-module de type fini ;
\item $\mathcal{F}$ est $(-1)$-pseudo-cohérent relativement à $S$ si et seulement si
$\mathcal{F}$ est quasi-cohérent et de présentation finie relativement à $S$.
\end{enumerate}
\end{lemma}

\begin{proof}
L'assertion (1) résulte immédiatement de la définition. Pour voir l'assertion (3),
on peut travailler localement sur $X$ (les deux propriétés sont locales). Ainsi, on
peut supposer $X$ et $S$ affines. Choisissons une immersion fermée
$i : X \to \mathbf{A}^n_S$. On voit alors que $\mathcal{F}$ est
$(-1)$-pseudo-cohérent relativement à $S$ si et seulement si $i_*\mathcal{F}$
est $(-1)$-pseudo-cohérent, ce qui a lieu si et seulement si $i_*\mathcal{F}$
est un $\mathcal{O}_{\mathbf{A}^n_S}$-module de présentation finie, voir
Cohomologie, lemme \ref{cohomology-lemma-finite-cohomology}.
Un module de présentation finie est quasi-cohérent, voir
Modules, lemme \ref{modules-lemma-finite-presentation-quasi-coherent}.
D'après Morphismes, lemme \ref{morphisms-lemma-i-star-equivalence},
on voit que $\mathcal{F}$ est quasi-cohérent si et seulement si $i_*\mathcal{F}$
est quasi-cohérent. Cela étant dit, l'assertion (3) en découle. La démonstration de (2)
est analogue, mais plus simple.
\end{proof}

\begin{lemma}
\label{more-morphisms-lemma-summands-relative-pseudo-coherent}
Soit $X \to S$ un morphisme de schémas localement de type fini.
Soit $m \in \mathbf{Z}$. Soient $E, K$ des objets de $D(\mathcal{O}_X)$.
Si $E \oplus K$ est $m$-pseudo-cohérent relativement à $S$, alors $E$ et $K$ le sont aussi.
\end{lemma}

\begin{proof}
Cela résulte de
Cohomologie, lemme \ref{cohomology-lemma-summands-pseudo-coherent},
et des définitions.
\end{proof}

\begin{lemma}
\label{more-morphisms-lemma-complex-relative-pseudo-coherent-modules}
Soit $X \to S$ un morphisme de schémas localement de type fini.
Soit $m \in \mathbf{Z}$. Soit $\mathcal{F}^\bullet$ un complexe (localement) borné
supérieurement de $\mathcal{O}_X$-modules tel que $\mathcal{F}^i$ soit
$(m - i)$-pseudo-cohérent relativement à $S$ pour tout $i$. Alors
$\mathcal{F}^\bullet$ est $m$-pseudo-cohérent relativement à $S$.
\end{lemma}

\begin{proof}
Cela résulte de
Cohomologie, lemme \ref{cohomology-lemma-complex-pseudo-coherent-modules},
et des définitions.
\end{proof}

\begin{lemma}
\label{more-morphisms-lemma-cohomology-relative-pseudo-coherent}
Soit $X \to S$ un morphisme de schémas localement de type fini.
Soit $m \in \mathbf{Z}$. Soit $E$ un objet de $D(\mathcal{O}_X)$.
Si $E$ est (localement) borné supérieurement et si $H^i(E)$ est $(m - i)$-pseudo-cohérent
relativement à $S$ pour tout $i$, alors $E$ est $m$-pseudo-cohérent relativement à $S$.
\end{lemma}

\begin{proof}
Cela résulte de
Cohomologie, lemme \ref{cohomology-lemma-cohomology-pseudo-coherent},
et des définitions.
\end{proof}

\begin{lemma}
\label{more-morphisms-lemma-base-change-relative-pseudo-coherent}
Soit $X \to S$ un morphisme de schémas localement de type fini.
Soit $m \in \mathbf{Z}$. Soit $E$ un objet de $D(\mathcal{O}_X)$
qui est $m$-pseudo-cohérent relativement à $S$. Soit $S' \to S$ un
morphisme de schémas. Posons $X' = X \times_S S'$ et notons $E'$
l'image inverse dérivée de $E$ sur $X'$. Si $S'$ et $X$ sont
Tor-indépendants sur $S$, alors $E'$
est $m$-pseudo-cohérent relativement à $S'$.
\end{lemma}

\begin{proof}
Le problème est local sur $X$ et $X'$ ; on peut donc supposer $X$, $S$, $S'$ et
$X'$ affines. Choisissons une immersion fermée $i : X \to \mathbf{A}^n_S$
et notons $i' : X' \to \mathbf{A}^n_{S'}$ son changement de base à $S'$.
Notons $g : X' \to X$ et $g' : \mathbf{A}^n_{S'} \to \mathbf{A}^n_S$
les projections, de sorte que $E' = Lg^*E$. Puisque $X$ et $S'$ sont tor-indépendants
sur $S$, le morphisme de changement de base
(Cohomologie, remarque \ref{cohomology-remark-base-change})
induit un isomorphisme
$$
Ri'_*(Lg^*E) = L(g')^*Ri_*E
$$
En effet, pour un point $x' \in X'$ au-dessus de $x \in X$, le morphisme de changement de base
sur les fibres en $x'$ est le morphisme
$$
E_x \otimes_{\mathcal{O}_{\mathbf{A}^n_S, x}}^\mathbf{L}
\mathcal{O}_{\mathbf{A}^n_{S'}, x'}
\longrightarrow
E_x \otimes_{\mathcal{O}_{X, x}}^\mathbf{L}
\mathcal{O}_{X', x'}
$$
provenant des immersions fermées $i$ et $i'$. Remarquons que la source
est quasi-isomorphe à une localisation de
$E_x \otimes_{\mathcal{O}_{S, s}}^\mathbf{L} \mathcal{O}_{S', s'}$,
qui est isomorphe au but, puisque
$\mathcal{O}_{X', x'}$ est isomorphe à la (même) localisation de
$\mathcal{O}_{X, x} \otimes_{\mathcal{O}_{S, s}}^\mathbf{L}
\mathcal{O}_{S', s'}$ par hypothèse. On conclut le lemme
en appliquant
Cohomologie, lemme \ref{cohomology-lemma-pseudo-coherent-pullback}.
\end{proof}

\begin{lemma}
\label{more-morphisms-lemma-pull-relative-pseudo-coherent}
Soit $f : X \to Y$ un morphisme de schémas localement de type fini
sur une base $S$. Soit $m \in \mathbf{Z}$. Soit $E$ un objet de
$D(\mathcal{O}_Y)$. Supposons que
\begin{enumerate}
\item $\mathcal{O}_X$ soit pseudo-cohérent relativement à $Y$\footnote{Cela
signifie que $f$ est pseudo-cohérent, voir la
définition \ref{more-morphisms-definition-pseudo-coherent}.} ;
\item $E$ soit $m$-pseudo-cohérent relativement à $S$.
\end{enumerate}
Alors $Lf^*E$ est $m$-pseudo-cohérent relativement à $S$.
\end{lemma}

\begin{proof}
Le problème est local sur $X$. On peut donc supposer $X$, $Y$ et $S$ affines.
En raisonnant comme dans la démonstration de Plus sur l'algèbre, lemme
\ref{more-algebra-lemma-pull-relative-pseudo-coherent},
on peut trouver un diagramme commutatif
$$
\xymatrix{
X \ar[r]_i \ar[d]_f &
\mathbf{A}^d_Y \ar[r]_j \ar[ld]^p &
\mathbf{A}^{n + d}_S \ar[ld] \\
Y \ar[r] &
\mathbf{A}^n_S
}
$$
Observons que
$$
Ri_* Lf^*E = Ri_* Li^* Lp^*E =
Lp^*E \otimes_{\mathcal{O}_{\mathbf{A}_Y^n}}^\mathbf{L} Ri_*\mathcal{O}_X
$$
d'après Cohomologie, lemme
\ref{cohomology-lemma-projection-formula-closed-immersion}.
Par hypothèse et puisque $Y$ est affine, on peut représenter
$Ri_*\mathcal{O}_X = i_*\mathcal{O}_X$ par un complexe de modules libres de type fini
sur $\mathcal{O}_{\mathbf{A}_Y^n}$, noté $\mathcal{F}^\bullet$, avec
$\mathcal{F}^q = 0$ pour $q > 0$
(détails omis ; utiliser Catégories dérivées de schémas, lemme
\ref{perfect-lemma-pseudo-coherent-affine}
et
Plus sur l'algèbre, lemme
\ref{more-algebra-lemma-rel-n-pseudo-module}).
Par hypothèse, $E$ est borné supérieurement ; écrivons $H^q(E) = 0$ pour $q > a$.
Représentons $E$ par un complexe $\mathcal{E}^\bullet$ de $\mathcal{O}_Y$-modules
tel que $\mathcal{E}^q = 0$ pour $q > a$. Alors le produit tensoriel dérivé ci-dessus
est représenté par $\text{Tot}(p^*\mathcal{E}^\bullet
\otimes_{\mathcal{O}_{\mathbf{A}_Y^n}} \mathcal{F}^\bullet)$.

\medskip\noindent
Puisque $j$ est une immersion fermée, le foncteur $j_*$ est exact et
$Rj_*$ se calcule en appliquant $j_*$ à n'importe quel complexe représentatif
de faisceaux. Il faut donc montrer que
$j_*\text{Tot}(p^*\mathcal{E}^\bullet \otimes_{\mathcal{O}_{\mathbf{A}_Y^n}}
\mathcal{F}^\bullet)$ est $m$-pseudo-cohérent
comme complexe de $\mathcal{O}_{\mathbf{A}^{n + m}_S}$-modules.
Remarquons que
$\text{Tot}(p^*\mathcal{E}^\bullet \otimes_{\mathcal{O}_{\mathbf{A}_Y^n}}
\mathcal{F}^\bullet)$ possède une filtration par
sous-complexes dont les quotients successifs sont les complexes
$p^*\mathcal{E}^\bullet
\otimes_{\mathcal{O}_{\mathbf{A}_Y^n}} \mathcal{F}^q[-q]$.
Remarquons que, pour $q \ll 0$, les complexes
$p^*\mathcal{E}^\bullet \otimes_{\mathcal{O}_{\mathbf{A}_Y^n}}
\mathcal{F}^q[-q]$
ont une cohomologie nulle en degrés $\leq m$ et sont donc $m$-pseudo-cohérents.
Ainsi, en appliquant le
lemme \ref{more-morphisms-lemma-cone-relatively-pseudo-coherent}
et en raisonnant par récurrence, il suffit de montrer que
$p^*\mathcal{E}^\bullet \otimes_{\mathcal{O}_{\mathbf{A}_Y^n}}
\mathcal{F}^q[-q]$
est pseudo-cohérent relativement à $S$ pour tout $q$. Remarquons que $\mathcal{F}^q = 0$
pour $q > 0$. Puisque $\mathcal{F}^q$ est aussi libre de type fini, cela
se ramène à prouver que $p^*\mathcal{E}^\bullet$ est
$m$-pseudo-cohérent relativement à $S$, ce qui résulte du
lemme \ref{more-morphisms-lemma-base-change-relative-pseudo-coherent}
par exemple.
\end{proof}

\begin{lemma}
\label{more-morphisms-lemma-composition-relative-pseudo-coherent}
Soit $f : X \to Y$ un morphisme de schémas localement de type fini
sur une base $S$. Soit $m \in \mathbf{Z}$. Soit $E$ un objet de
$D(\mathcal{O}_X)$. Supposons que $\mathcal{O}_Y$ soit pseudo-cohérent relativement
à $S$\footnote{Cela signifie que $Y \to S$ est pseudo-cohérent, voir la
définition \ref{more-morphisms-definition-pseudo-coherent}.}.
Alors les assertions suivantes sont équivalentes :
\begin{enumerate}
\item $E$ est $m$-pseudo-cohérent relativement à $Y$ ;
\item $E$ est $m$-pseudo-cohérent relativement à $S$.
\end{enumerate}
\end{lemma}

\begin{proof}
La question est locale sur $X$ ; on peut donc supposer $X$, $Y$ et $S$ affines.
En raisonnant comme dans la démonstration de Plus sur l'algèbre, lemme
\ref{more-algebra-lemma-pull-relative-pseudo-coherent},
on peut trouver un diagramme commutatif
$$
\xymatrix{
X \ar[r]_i \ar[d]_f &
\mathbf{A}^m_Y \ar[r]_j \ar[ld]^p &
\mathbf{A}^{n + m}_S \ar[ld] \\
Y \ar[r] &
\mathbf{A}^n_S
}
$$
L'hypothèse que $\mathcal{O}_Y$ est pseudo-cohérent relativement à $S$
entraîne que $\mathcal{O}_{\mathbf{A}^m_Y}$ est pseudo-cohérent relativement
à $\mathbf{A}^m_S$ (par changement de base plat ; on peut le voir en utilisant,
par exemple, le lemme \ref{more-morphisms-lemma-base-change-relative-pseudo-coherent}).
Cela entraîne à son tour que $j_*\mathcal{O}_{\mathbf{A}^n_Y}$ est
pseudo-cohérent comme
$\mathcal{O}_{\mathbf{A}^{n + m}_S}$-module. L'équivalence du
lemme résulte alors de
Catégories dérivées de schémas, lemme
\ref{perfect-lemma-closed-push-pseudo-coherent}.
\end{proof}

\begin{lemma}
\label{more-morphisms-lemma-check-relative-pseudo-coherence-on-charts}
Soit
$$
\xymatrix{
X \ar[rd] \ar[rr]_i & & P \ar[ld] \\
& S
}
$$
un diagramme commutatif de schémas. Supposons que $i$ soit une immersion fermée
et que $P \to S$ soit plat et localement de présentation finie. Soit $E$
un objet de $D(\mathcal{O}_X)$. Alors les assertions suivantes
sont équivalentes :
\begin{enumerate}
\item $E$ est $m$-pseudo-cohérent relativement à $S$ ;
\item $Ri_*E$ est $m$-pseudo-cohérent relativement à $S$ ;
\item $Ri_*E$ est $m$-pseudo-cohérent sur $P$.
\end{enumerate}
\end{lemma}

\begin{proof}
L'équivalence de (1) et (2) est le
lemme \ref{more-morphisms-lemma-finite-morphism-relative-pseudo-coherence}.
L'équivalence de (2) et (3) résulte du
lemme
\ref{more-morphisms-lemma-composition-relative-pseudo-coherent}
appliqué à $\text{id} : P \to P$,
pourvu que l'on montre que $\mathcal{O}_P$ est
pseudo-cohérent relativement à $S$. Cela résulte de
Plus sur l'algèbre, lemme
\ref{more-algebra-lemma-flat-finite-presentation-perfect},
et des définitions.
\end{proof}









\section{Morphismes pseudo-cohérents}
\label{more-morphisms-section-pseudo-coherent}

\noindent
Évitez à tout prix de lire cette section.
Si vous avez besoin d'une partie de ce matériau, consultez d'abord les
sections algébriques correspondantes, voir
Plus sur l'algèbre, sections \ref{more-algebra-section-pseudo-coherent},
\ref{more-algebra-section-relative-pseudo-coherent} et
\ref{more-algebra-section-pseudo-coherent-perfect-ring-map}.
Pour l'instant, la seule chose à savoir est qu'un homomorphisme d'anneaux
$A \to B$ est pseudo-cohérent si et seulement si $B = A[x_1, \ldots, x_n]/I$
et si $B$, comme $A[x_1, \ldots, x_n]$-module, possède une résolution par des
$A[x_1, \ldots, x_n]$-modules libres de type fini.

\begin{lemma}
\label{more-morphisms-lemma-pseudo-coherent}
Soit $f : X \to S$ un morphisme de schémas. Les assertions suivantes sont équivalentes :
\begin{enumerate}
\item il existe un recouvrement ouvert affine $S = \bigcup V_j$ et, pour chaque $j$,
un recouvrement ouvert affine $f^{-1}(V_j) = \bigcup U_{ji}$ tel que
$\mathcal{O}_S(V_j) \to \mathcal{O}_X(U_{ij})$ soit un homomorphisme
d'anneaux pseudo-cohérent ;
\item pour toute paire d'ouverts affines $U \subset X$, $V \subset S$
telle que $f(U) \subset V$, l'homomorphisme d'anneaux
$\mathcal{O}_S(V) \to \mathcal{O}_X(U)$ est pseudo-cohérent ;
\item $f$ est localement de type fini et $\mathcal{O}_X$
est pseudo-cohérent relativement à $S$.
\end{enumerate}
\end{lemma}

\begin{proof}
Pour voir l'équivalence de (1) et (2), il suffit de vérifier les conditions
(1)(a), (b), (c) de
Morphismes, définition \ref{morphisms-definition-property-local},
pour la propriété d'être un homomorphisme d'anneaux pseudo-cohérent.
Ces propriétés résultent (en utilisant que la localisation est plate) de
Plus sur l'algèbre, lemmes
\ref{more-algebra-lemma-base-change-relative-pseudo-coherent},
\ref{more-algebra-lemma-localize-relative-pseudo-coherent} et
\ref{more-algebra-lemma-glue-relative-pseudo-coherent}.

\medskip\noindent
Si (1) est vérifiée, alors $f$ est localement de type fini, car un homomorphisme
d'anneaux pseudo-cohérent est de type fini par définition. De plus, (1) entraîne,
via le lemme \ref{more-morphisms-lemma-qcoh-relative-pseudo-coherence-characterize}
et les définitions, que $\mathcal{O}_X$ est pseudo-cohérent
relativement à $S$. Réciproquement, si (3) est vérifiée, on voit
que, pour tous $U$ et $V$ comme dans (2), l'anneau
$\mathcal{O}_X(U)$ est de type fini sur $\mathcal{O}_S(V)$
et que $\mathcal{O}_X(U)$ est, comme module,
pseudo-cohérent relativement à $\mathcal{O}_S(V)$, voir les
lemmes \ref{more-morphisms-lemma-relative-pseudo-coherence-characterize} et
\ref{more-morphisms-lemma-qcoh-relative-pseudo-coherence-characterize}.
C'est la définition d'un homomorphisme d'anneaux pseudo-cohérent ; ainsi,
(2) et (1) sont vérifiées.
\end{proof}

\begin{definition}
\label{more-morphisms-definition-pseudo-coherent}
Un morphisme de schémas $f : X \to S$ est dit {\it pseudo-cohérent}
si les conditions équivalentes du
lemme \ref{more-morphisms-lemma-pseudo-coherent}
sont satisfaites. Dans ce cas, on dit aussi que $X$ est pseudo-cohérent
sur $S$.
\end{definition}

\noindent
Attention : un changement de base d'un morphisme pseudo-cohérent n'est pas
pseudo-cohérent en général.

\begin{lemma}
\label{more-morphisms-lemma-flat-base-change-pseudo-coherent}
Un changement de base plat d'un morphisme pseudo-cohérent est pseudo-cohérent.
\end{lemma}

\begin{proof}
Cela se traduit par le résultat algébrique suivant :
soit $A \to B$ un homomorphisme d'anneaux pseudo-cohérent.
Soit $A \to A'$ plat. Alors $A' \to B \otimes_A A'$ est
pseudo-cohérent. Cela résulte du résultat plus général
Plus sur l'algèbre,
lemme \ref{more-algebra-lemma-base-change-relative-pseudo-coherent}.
\end{proof}

\begin{lemma}
\label{more-morphisms-lemma-composition-pseudo-coherent}
Un composé de morphismes de schémas pseudo-cohérents est
pseudo-cohérent.
\end{lemma}

\begin{proof}
Cela se traduit par le résultat algébrique suivant :
si $A \to B \to C$ sont des homomorphismes d'anneaux pseudo-cohérents composables,
alors $A \to C$ est pseudo-cohérent. Cela résulte soit de
Plus sur l'algèbre,
lemme \ref{more-algebra-lemma-pull-relative-pseudo-coherent},
soit de
Plus sur l'algèbre,
lemme \ref{more-algebra-lemma-composition-relative-pseudo-coherent}.
\end{proof}

\begin{lemma}
\label{more-morphisms-lemma-pseudo-coherent-finite-presentation}
Un morphisme pseudo-cohérent est localement de présentation finie.
\end{lemma}

\begin{proof}
Cela résulte immédiatement des définitions.
\end{proof}

\begin{lemma}
\label{more-morphisms-lemma-flat-finite-presentation-pseudo-coherent}
Un morphisme plat et localement de présentation finie est pseudo-cohérent.
\end{lemma}

\begin{proof}
Cela résulte du fait qu'un homomorphisme d'anneaux plat de présentation finie est
pseudo-cohérent (et même parfait), voir
Plus sur l'algèbre,
lemme \ref{more-algebra-lemma-flat-finite-presentation-perfect}.
\end{proof}

\begin{lemma}
\label{more-morphisms-lemma-permanence-pseudo-coherent}
Soit $f : X \to Y$ un morphisme de schémas pseudo-cohérents
sur un schéma de base $S$. Alors $f$ est pseudo-cohérent.
\end{lemma}

\begin{proof}
Cela se traduit par le résultat algébrique suivant :
si $R \to A \to B$ sont des homomorphismes d'anneaux composables
et si $R \to A$, $R \to B$ sont pseudo-cohérents, alors
$R \to B$ est pseudo-cohérent. Cela résulte de
Plus sur l'algèbre,
lemme \ref{more-algebra-lemma-composition-relative-pseudo-coherent}.
\end{proof}

\begin{lemma}
\label{more-morphisms-lemma-finite-pseudo-coherent}
Soit $f : X \to S$ un morphisme fini de schémas.
Alors $f$ est pseudo-cohérent si et seulement si $f_*\mathcal{O}_X$
est pseudo-cohérent comme $\mathcal{O}_S$-module.
\end{lemma}

\begin{proof}
Traduit en algèbre, ce lemme affirme ceci : si $R \to A$
est un homomorphisme d'anneaux fini, alors $R \to A$ est pseudo-cohérent comme homomorphisme d'anneaux
(ce qui signifie par définition que $A$, comme $A$-module, est
pseudo-cohérent relativement à $R$) si et seulement si
$A$ est pseudo-cohérent comme $R$-module. Cela résulte du résultat
plus général de
Plus sur l'algèbre, lemme
\ref{more-algebra-lemma-finite-extension-pseudo-coherent}.
\end{proof}

\begin{lemma}
\label{more-morphisms-lemma-Noetherian-pseudo-coherent}
Soit $f : X \to S$ un morphisme de schémas.
Si $S$ est localement noethérien, alors $f$ est pseudo-cohérent si
et seulement si $f$ est localement de type fini.
\end{lemma}

\begin{proof}
Cela se traduit par le résultat algébrique suivant :
si $R \to A$ est un homomorphisme d'anneaux de type fini avec $R$ noethérien, alors
$R \to A$ est pseudo-cohérent si et seulement si $R \to A$ est de type fini.
Pour le voir, remarquons qu'un homomorphisme d'anneaux pseudo-cohérent est de type fini par
définition. Réciproquement, si $R \to A$ est de type fini, alors
on peut écrire $A = R[x_1, \ldots, x_n]/I$ et il résulte de
Plus sur l'algèbre,
lemme \ref{more-algebra-lemma-Noetherian-pseudo-coherent},
que $A$ est pseudo-cohérent comme $R[x_1, \ldots, x_n]$-module, c'est-à-dire que
$R \to A$ est un homomorphisme d'anneaux pseudo-cohérent.
\end{proof}

\begin{lemma}
\label{more-morphisms-lemma-descending-property-pseudo-coherent}
La propriété $\mathcal{P}(f) =$``$f$ est pseudo-cohérent''
est fpqc locale sur la base.
\end{lemma}

\begin{proof}
Nous allons utiliser le critère de
Descente, lemme \ref{descent-lemma-descending-properties-morphisms}
pour le démontrer. Par
la définition \ref{more-morphisms-definition-pseudo-coherent},
la propriété d'être pseudo-cohérent est locale pour la topologie de Zariski sur la base. Par
le lemme \ref{more-morphisms-lemma-flat-base-change-pseudo-coherent},
la propriété d'être pseudo-cohérent est préservée par changement de base plat.
La dernière hypothèse (3) du
lemme \ref{descent-lemma-descending-properties-morphisms} de Descente
se traduit par l'énoncé algébrique suivant :
Soit $A \to B$ un homomorphisme fidèlement plat d'anneaux.
Soit $C = A[x_1, \ldots, x_n]/I$ une $A$-algèbre.
Si $C \otimes_A B$ est pseudo-cohérent comme $B[x_1, \ldots, x_n]$-module,
alors $C$ est pseudo-cohérent comme $A[x_1, \ldots, x_n]$-module.
C'est le
lemme \ref{more-algebra-lemma-flat-descent-pseudo-coherent} de Compléments d'algèbre.
\end{proof}

\begin{lemma}
\label{more-morphisms-lemma-quotient-of-flat-finitely-presented}
Soit $A \to B$ un homomorphisme plat d'anneaux de présentation finie.
Soit $I \subset B$ un idéal. Alors $A \to B/I$ est pseudo-cohérent
si et seulement si $I$ est pseudo-cohérent comme $B$-module.
\end{lemma}

\begin{proof}
Choisissons une présentation $B = A[x_1, \ldots, x_n]/J$.
Remarquons que $B$ est pseudo-cohérent comme $A[x_1, \ldots, x_n]$-module
car $A \to B$ est un homomorphisme pseudo-cohérent d'anneaux d'après le
lemme \ref{more-morphisms-lemma-flat-finite-presentation-pseudo-coherent}.
Remarquons que $A \to B/I$ est pseudo-cohérent si et seulement si
$B/I$ est pseudo-cohérent comme $A[x_1, \ldots, x_n]$-module. D'après le
lemme \ref{more-algebra-lemma-finite-push-pseudo-coherent} de Compléments d'algèbre,
nous voyons que cela équivaut à demander que $B/I$ soit
pseudo-cohérent comme $B$-module. Cela démontre le lemme, car la
suite exacte courte $0 \to I \to B \to B/I \to 0$
montre que $I$ est pseudo-cohérent si et seulement si $B/I$ l'est (voir
Compléments d'algèbre,
lemme \ref{more-algebra-lemma-two-out-of-three-pseudo-coherent}).
\end{proof}

\noindent
Le lemme suivant sera supplanté par le
lemme \ref{more-morphisms-lemma-pseudo-coherent-fppf-local-source}, plus fort.

\begin{lemma}
\label{more-morphisms-lemma-pseudo-coherent-syntomic-local-source}
La propriété $\mathcal{P}(f) =$``$f$ est pseudo-cohérent''
est locale pour la topologie syntomique sur la source.
\end{lemma}

\begin{proof}
Nous allons utiliser le critère du
lemme \ref{descent-lemma-properties-morphisms-local-source} de Descente
pour le démontrer. Il résulte des
lemmes \ref{more-morphisms-lemma-flat-finite-presentation-pseudo-coherent} et
\ref{more-morphisms-lemma-composition-pseudo-coherent}
que la propriété d'être pseudo-cohérent est préservée par précomposition avec des
morphismes plats localement de présentation finie, en particulier par
précomposition avec des morphismes syntomiques (voir
Morphismes, lemmes \ref{morphisms-lemma-syntomic-flat} et
\ref{morphisms-lemma-syntomic-locally-finite-presentation}).
Il ressort clairement de la
définition \ref{more-morphisms-definition-pseudo-coherent}
que la propriété d'être pseudo-cohérent est
locale pour la topologie de Zariski sur la source et le but.
Ainsi, d'après le
lemme \ref{descent-lemma-properties-morphisms-local-source} de Descente déjà cité,
il suffit de démontrer ce qui suit : supposons que $X' \to X \to Y$ soient des
morphismes de schémas affines, avec $X' \to X$ syntomique et $X' \to Y$
pseudo-cohérent. Alors $X \to Y$ est pseudo-cohérent.
Pour le voir, remarquons que, dans tous les cas, $X \to Y$ est de présentation finie d'après
Descente, lemme \ref{descent-lemma-flat-finitely-presented-permanence-algebra}.
Choisissons une immersion fermée $X \to \mathbf{A}^n_Y$. D'après
Algèbre, lemme \ref{algebra-lemma-lift-syntomic},
nous pouvons trouver un recouvrement d'ouverts affines $X' = \bigcup_{i = 1, \ldots, n} X'_i$
et des morphismes syntomiques $W_i \to \mathbf{A}^n_Y$ relevant les morphismes
$X'_i \to X$, c'est-à-dire tels qu'il existe des diagrammes cartésiens
$$
\xymatrix{
X'_i \ar[d] \ar[r] & W_i \ar[d] \\
X \ar[r] & \mathbf{A}^n_Y
}
$$
Après avoir remplacé $X'$ par $\coprod X'_i$ et posé $W = \coprod W_i$,
nous obtenons un diagramme cartésien
$$
\xymatrix{
X' \ar[d] \ar[r] & W \ar[d]^h \\
X \ar[r] & \mathbf{A}^n_Y
}
$$
où $W \to \mathbf{A}^n_Y$ est plat et de présentation finie, et
$X' \to Y$ est encore pseudo-cohérent. Puisque
$W \to \mathbf{A}^n_Y$ est ouvert (voir
Morphismes, lemme \ref{morphisms-lemma-fppf-open})
et que $X' \to X$ est surjectif, nous pouvons trouver
$f \in \Gamma(\mathbf{A}^n_Y, \mathcal{O})$ tel que
$X \subset D(f) \subset \Im(h)$. Écrivons
$Y = \Spec(R)$, $X = \Spec(A)$, $X' = \Spec(A')$
et $W = \Spec(B)$, $A = R[x_1, \ldots, x_n]/I$ et
$A' = B/IB$. Alors $R \to A'$ est pseudo-cohérent. Considérons le diagramme
$$
\xymatrix{
A' = B/IB & B \ar[l] \\
A = R[x_1, \ldots, x_n]/I \ar[u] & R[x_1, \ldots, x_n] \ar[l] \ar[u]
}
$$
D'après le
lemme \ref{more-morphisms-lemma-quotient-of-flat-finitely-presented},
nous voyons que $IB$ est pseudo-cohérent comme $B$-module.
L'homomorphisme d'anneaux $R[x_1, \ldots, x_n]_f \to B_f$ est fidèlement plat par
notre choix de $f$ ci-dessus. Cela entraîne que
$I_f \subset R[x_1, \ldots, x_n]_f$
est pseudo-cohérent, voir
Compléments d'algèbre, lemme \ref{more-algebra-lemma-flat-descent-pseudo-coherent}.
En appliquant encore une fois le
lemme \ref{more-morphisms-lemma-quotient-of-flat-finitely-presented},
nous voyons que $R \to A$ est pseudo-cohérent.
\end{proof}

\begin{lemma}
\label{more-morphisms-lemma-pseudo-coherent-fppf-local-source}
La propriété $\mathcal{P}(f) =$``$f$ est pseudo-cohérent''
est locale pour la topologie fppf sur la source.
\end{lemma}

\begin{proof}
Soit $f : X \to S$ un morphisme de schémas.
Soit $\{g_i : X_i \to X\}$ un recouvrement fppf tel que chaque composée
$f \circ g_i$ soit pseudo-cohérente. D'après le
lemme \ref{more-morphisms-lemma-dominate-fppf},
il existe
\begin{enumerate}
\item un recouvrement d'ouverts de Zariski $X = \bigcup U_j$,
\item des morphismes surjectifs finis localement libres $W_j \to U_j$,
\item des recouvrements d'ouverts de Zariski $W_j = \bigcup_k W_{j, k}$,
\item des morphismes surjectifs finis localement libres $T_{j, k} \to W_{j, k}$
\end{enumerate}
tels que le recouvrement fppf $\{h_{j, k} : T_{j, k} \to X\}$ raffine le
recouvrement donné $\{X_i \to X\}$. Notons $\psi_{j, k} : T_{j, k} \to X_{\alpha(j, k)}$
les morphismes qui témoignent du fait que $\{T_{j, k} \to X\}$ raffine
le recouvrement donné $\{X_i \to X\}$. Remarquons que $T_{j, k} \to X$ est un morphisme plat
localement de présentation finie ; ainsi, aussi bien $X_i$ que $T_{j, k}$ sont
pseudo-cohérents sur $X$ d'après le
lemme \ref{more-morphisms-lemma-flat-finite-presentation-pseudo-coherent}.
Ainsi, $\psi_{j, k} : T_{j, k} \to X_i$ est pseudo-cohérent, voir le
lemme \ref{more-morphisms-lemma-permanence-pseudo-coherent}.
Ainsi, $T_{j, k} \to S$ est pseudo-cohérent comme composée
de $\psi_{j, k}$ et de $f \circ g_{\alpha(j, k)}$, voir le
lemme \ref{more-morphisms-lemma-composition-pseudo-coherent}.
Nous voyons donc que le lemme est ramené au cas d'un recouvrement
d'ouverts de Zariski (qui convient) et au cas d'un recouvrement
donné par un seul morphisme fini localement libre et surjectif, que nous
traitons dans le paragraphe suivant.

\medskip\noindent
Supposons que $X' \to X \to S$ soit une suite de morphismes de schémas,
où $X' \to X$ est fini localement libre et surjectif et $X' \to Y$ pseudo-cohérent.
Notre but est de montrer que $X \to S$ est pseudo-cohérent.
Remarquons que, d'après
Descente, lemme \ref{descent-lemma-flat-finitely-presented-permanence},
le morphisme $X \to S$ est localement de présentation finie.
Il est clair que le problème se ramène au cas où $X'$, $X$ et $S$
sont affines et où $X' \to X$ est libre de rang $r > 0$. Le problème
algébrique correspondant est le suivant : supposons que $R \to A \to A'$ soient des homomorphismes d'anneaux
tels que $R \to A'$ soit pseudo-cohérent, que $R \to A$ soit de présentation finie,
et que $A' \cong A^{\oplus r}$ comme $A$-module. But : montrer que $R \to A$ est
pseudo-cohérent. L'hypothèse que $R \to A'$ est pseudo-cohérent
signifie que $A'$, comme $A'$-module, est pseudo-cohérent relativement à $R$. D'après
Compléments d'algèbre,
lemme \ref{more-algebra-lemma-finite-extension-pseudo-coherent},
cela entraîne que $A'$, comme $A$-module, est pseudo-cohérent relativement à $R$.
Puisque $A' \cong A^{\oplus r}$ comme $A$-module, nous voyons que
$A$, comme $A$-module, est pseudo-cohérent relativement à $R$, voir
Compléments d'algèbre,
lemme \ref{more-algebra-lemma-summands-relative-pseudo-coherent}.
Cela signifie, par définition, que $R \to A$ est pseudo-cohérent, ce qui conclut.
\end{proof}








\section{Morphismes parfaits}
\label{more-morphisms-section-perfect}

\noindent
Pour comprendre les résultats de cette
section, il suffit de connaître un peu ceux de la section
sur les morphismes pseudo-cohérents.
Pour le moment, la seule chose à savoir est qu'un homomorphisme d'anneaux
$A \to B$ est parfait si et seulement s'il est pseudo-cohérent
et si $B$ est de dimension de Tor finie comme $A$-module.

\begin{lemma}
\label{more-morphisms-lemma-perfect}
Soit $f : X \to S$ un morphisme de schémas localement de type fini.
Les assertions suivantes sont équivalentes :
\begin{enumerate}
\item il existe un recouvrement d'ouverts affines $S = \bigcup V_j$ et, pour chaque $j$,
un recouvrement d'ouverts affines $f^{-1}(V_j) = \bigcup U_{ji}$ tel que
$\mathcal{O}_S(V_j) \to \mathcal{O}_X(U_{ij})$ soit un homomorphisme
d'anneaux parfait, et
\item pour toute paire d'ouverts affines $U \subset X$, $V \subset S$
tels que $f(U) \subset V$, l'homomorphisme d'anneaux
$\mathcal{O}_S(V) \to \mathcal{O}_X(U)$ est parfait.
\end{enumerate}
\end{lemma}

\begin{proof}
Supposons (1) et soient $U, V$ comme dans (2).
Il résulte du
lemme \ref{more-morphisms-lemma-pseudo-coherent}
que $\mathcal{O}_S(V) \to \mathcal{O}_X(U)$ est pseudo-cohérent.
Il suffit donc de démontrer que la propriété, pour un homomorphisme d'anneaux,
d'être « de dimension de Tor finie » satisfait aux
conditions (1)(a), (b), (c) de
Morphismes, définition \ref{morphisms-definition-property-local}.
Ces propriétés résultent de
Compléments d'algèbre,
lemmes \ref{more-algebra-lemma-flat-push-tor-amplitude},
\ref{more-algebra-lemma-flat-base-change-finite-tor-dimension} et
\ref{more-algebra-lemma-glue-tor-amplitude}.
Certains détails sont omis.
\end{proof}

\begin{definition}
\label{more-morphisms-definition-perfect}
Un morphisme de schémas $f : X \to S$ est dit {\it parfait}
si les conditions équivalentes du
lemme \ref{more-morphisms-lemma-perfect}
sont satisfaites. Dans ce cas, nous disons aussi que $X$ est parfait
sur $S$.
\end{definition}

\noindent
Remarquons qu'un morphisme parfait est en particulier pseudo-cohérent, donc
localement de présentation finie. Attention : un changement de base d'un
morphisme parfait n'est pas parfait en général.

\begin{lemma}
\label{more-morphisms-lemma-flat-base-change-perfect}
Un changement de base plat d'un morphisme parfait est parfait.
\end{lemma}

\begin{proof}
Cela se traduit par le résultat algébrique suivant :
soit $A \to B$ un homomorphisme parfait d'anneaux.
Soit $A \to A'$ plat. Alors $A' \to B \otimes_A A'$ est
parfait. Nous avons vu ce résultat pour les homomorphismes pseudo-cohérents d'anneaux dans le
lemme \ref{more-morphisms-lemma-flat-base-change-pseudo-coherent}.
Le fait correspondant pour la dimension de Tor finie résulte de
Compléments d'algèbre,
lemme \ref{more-algebra-lemma-flat-base-change-finite-tor-dimension}.
\end{proof}

\begin{lemma}
\label{more-morphisms-lemma-composition-perfect}
Une composée de morphismes parfaits de schémas est parfaite.
\end{lemma}

\begin{proof}
Cela se traduit par le résultat algébrique suivant :
si $A \to B \to C$ sont des homomorphismes parfaits d'anneaux composables,
alors $A \to C$ est parfait. Nous avons vu que c'est le cas pour la
pseudo-cohérence dans le
lemme \ref{more-morphisms-lemma-composition-pseudo-coherent}
et sa démonstration. Par hypothèse, il existe des entiers $n$, $m$ tels
que $B$ soit de dimension de Tor $\leq n$ sur $A$ et $C$ de dimension de Tor
$\leq m$ sur $B$. Alors, pour tout $A$-module $M$, nous avons
$$
M \otimes_A^{\mathbf{L}} C =
(M \otimes_A^{\mathbf{L}} B) \otimes_B^{\mathbf{L}} C
$$
et la suite spectrale de
Compléments d'algèbre, exemple \ref{more-algebra-example-tor},
montre que $\text{Tor}^A_p(M, C) = 0$ pour $p > n + m$, comme souhaité.
\end{proof}

\begin{lemma}
\label{more-morphisms-lemma-flat-finite-presentation-perfect}
Soit $f : X \to S$ un morphisme de schémas.
Les assertions suivantes sont équivalentes :
\begin{enumerate}
\item $f$ est plat et parfait, et
\item $f$ est plat et localement de présentation finie.
\end{enumerate}
\end{lemma}

\begin{proof}
L'implication (2) $\Rightarrow$ (1) est donnée par
Compléments d'algèbre,
lemme \ref{more-algebra-lemma-flat-finite-presentation-perfect}.
La réciproque résulte du fait qu'un morphisme pseudo-cohérent
est localement de présentation finie, voir le
lemme \ref{more-morphisms-lemma-pseudo-coherent-finite-presentation}.
\end{proof}

\begin{lemma}
\label{more-morphisms-lemma-regular-target-perfect}
Soit $f : X \to S$ un morphisme de schémas.
Supposons que $S$ soit régulier et que $f$ soit localement de type fini.
Alors $f$ est parfait.
\end{lemma}

\begin{proof}
Voir
Compléments d'algèbre, lemme \ref{more-algebra-lemma-regular-perfect-ring-map}.
\end{proof}

\begin{lemma}
\label{more-morphisms-lemma-regular-immersion-perfect}
Une immersion régulière de schémas est parfaite.
Une immersion Koszul-régulière de schémas est parfaite.
\end{lemma}

\begin{proof}
Puisqu'une immersion régulière est une immersion Koszul-régulière, voir
Diviseurs, lemme \ref{divisors-lemma-regular-quasi-regular-immersion},
il suffit de démontrer la seconde assertion. Cela se traduit par l'énoncé
algébrique suivant : supposons que $I \subset A$ soit un
idéal engendré par une suite Koszul-régulière $f_1, \ldots, f_r$ de $A$.
Alors $A \to A/I$ est un homomorphisme parfait d'anneaux. Puisque $A \to A/I$ est surjectif,
ceci est une présentation de $A/I$ par une algèbre polynomiale sur $A$.
Il suffit donc de voir que $A/I$ est pseudo-cohérent comme $A$-module
et est de dimension de Tor finie. Par définition d'une suite de Koszul,
le complexe de Koszul $K(A, f_1, \ldots, f_r)$ est une résolution libre finie
de $A/I$. Ainsi, $A/I$ est un complexe parfait de $A$-modules, ce qui conclut.
\end{proof}

\begin{lemma}
\label{more-morphisms-lemma-perfect-permanence}
Soit
$$
\xymatrix{
X \ar[rr]_f \ar[rd] & & Y \ar[ld] \\
& S
}
$$
un diagramme commutatif de morphismes de schémas. Supposons que $Y \to S$ soit
lisse et que $X \to S$ soit parfait. Alors $f : X \to Y$ est parfait.
\end{lemma}

\begin{proof}
Nous pouvons factoriser $f$ comme la composée
$$
X \longrightarrow X \times_S Y \longrightarrow Y
$$
où le premier morphisme est l'application $i = (1, f)$ et le second
morphisme est la projection. Puisque $Y \to S$ est plat, voir
Morphismes, lemme \ref{morphisms-lemma-smooth-flat},
nous voyons que $X \times_S Y \to Y$ est parfait d'après le
lemme \ref{more-morphisms-lemma-flat-base-change-perfect}.
Puisque $Y \to S$ est lisse, $X \times_S Y \to X$ est également lisse, voir
Morphismes, lemme \ref{morphisms-lemma-base-change-smooth}.
Ainsi, $i$ est une section d'un morphisme lisse ; par conséquent, $i$ est
une immersion régulière, voir
Diviseurs, lemme \ref{divisors-lemma-section-smooth-regular-immersion}.
Cela entraîne que $i$ est parfait, voir le
lemme \ref{more-morphisms-lemma-regular-immersion-perfect}.
Nous concluons que $f$ est parfait, car la composée de morphismes
parfaits est parfaite, voir le
lemme \ref{more-morphisms-lemma-composition-perfect}.
\end{proof}

\begin{remark}
\label{more-morphisms-remark-perfect-permanence}
Il n'est pas vrai qu'un morphisme entre des schémas $X, Y$ parfaits sur une base $S$
soit parfait. Un exemple est $S = \Spec(k)$, $X = \Spec(k)$,
$Y = \Spec(k[x]/(x^2)$ et $X \to Y$ l'unique $S$-morphisme.
\end{remark}

\begin{lemma}
\label{more-morphisms-lemma-descending-property-perfect}
La propriété $\mathcal{P}(f) =$``$f$ est parfait''
est locale pour la topologie fpqc sur la base.
\end{lemma}

\begin{proof}
Nous allons utiliser le critère du
lemme \ref{descent-lemma-descending-properties-morphisms} de Descente
pour le démontrer. D'après la
définition \ref{more-morphisms-definition-perfect},
la propriété d'être parfait est locale pour la topologie de Zariski sur la base. D'après le
lemme \ref{more-morphisms-lemma-flat-base-change-perfect},
la propriété d'être parfait est préservée par changement de base plat.
La dernière hypothèse (3) du
lemme \ref{descent-lemma-descending-properties-morphisms} de Descente
se traduit par l'énoncé algébrique suivant :
soit $A \to B$ un homomorphisme fidèlement plat d'anneaux.
Soit $C = A[x_1, \ldots, x_n]/I$ une $A$-algèbre.
Si $C \otimes_A B$ est parfait comme $B[x_1, \ldots, x_n]$-module,
alors $C$ est parfait comme $A[x_1, \ldots, x_n]$-module.
C'est le
lemme \ref{more-algebra-lemma-flat-descent-perfect} de Compléments d'algèbre.
\end{proof}

\begin{lemma}
\label{more-morphisms-lemma-check-perfect-stalks}
Soit $f : X \to S$ un morphisme pseudo-cohérent de schémas.
Les assertions suivantes sont équivalentes :
\begin{enumerate}
\item $f$ est parfait,
\item $\mathcal{O}_X$ est localement de dimension de Tor finie comme
faisceau de $f^{-1}\mathcal{O}_S$-modules, et
\item pour tout $x \in X$, l'anneau $\mathcal{O}_{X, x}$ est de dimension de Tor
finie comme $\mathcal{O}_{S, f(x)}$-module.
\end{enumerate}
\end{lemma}

\begin{proof}
Le problème est local sur $X$ et $S$. Nous pouvons donc supposer que
$X = \Spec(B)$, $S = \Spec(A)$ et que $f$ correspond à un homomorphisme pseudo-cohérent
d'anneaux $A \to B$.

\medskip\noindent
Si (1) est satisfaite, alors $B$ est de dimension de Tor finie, égale à $d$, comme $A$-module. Alors
$B_\mathfrak q$ est de dimension de Tor $d$ comme $A_\mathfrak p$-module pour tout
idéal premier $\mathfrak q \subset B$ avec $\mathfrak p = A \cap \mathfrak q$, voir
Compléments d'algèbre, lemme \ref{more-algebra-lemma-tor-amplitude-localization}.
Alors $\mathcal{O}_X$ est de dimension de Tor $d$ comme
faisceau de $f^{-1}\mathcal{O}_S$-modules d'après
Cohomologie, lemme \ref{cohomology-lemma-tor-amplitude-stalk}.
Ainsi, (1) implique (2).

\medskip\noindent
D'après Cohomologie, lemme \ref{cohomology-lemma-tor-amplitude-stalk}, (2) implique
(3).

\medskip\noindent
Supposons (3). Nous ne pouvons pas utiliser
Compléments d'algèbre, lemme \ref{more-algebra-lemma-tor-amplitude-localization},
pour conclure, car nous ne savons pas que la dimension de Tor de
$B_\mathfrak q$ sur $A_\mathfrak p$ est bornée indépendamment de $\mathfrak q$.
Choisissons une présentation $A[x_1, \ldots, x_n] \to B$. Alors $B$ est
pseudo-cohérent comme $A[x_1, \ldots, x_n]$-module. Soit
$\mathfrak q \subset A[x_1, \ldots, x_n]$ un idéal premier
au-dessus de $\mathfrak p \subset A$. Alors, soit $B_\mathfrak q$ est nul,
soit, par hypothèse, il est de dimension de Tor finie comme
$A_\mathfrak p$-module. Puisque les fibres de $A \to A[x_1, \ldots, x_n]$
sont de dimension globale finie, nous pouvons appliquer le
lemme \ref{more-algebra-lemma-perfect-over-polynomial-ring} de Compléments d'algèbre
à $A_\mathfrak p \to A[x_1, \ldots, x_n]_\mathfrak q$
pour voir que $B_\mathfrak q$ est un
$A[x_1, \ldots, x_n]_\mathfrak q$-module parfait. Par conséquent,
$B$ est un $A[x_1, \ldots, x_n]$-module parfait d'après
Compléments d'algèbre, lemme \ref{more-algebra-lemma-check-perfect-stalks}.
Ainsi, $A \to B$ est un homomorphisme parfait d'anneaux par définition.
\end{proof}

\begin{lemma}
\label{more-morphisms-lemma-perfect-closed-immersion-perfect-direct-image}
Soit $i : Z \to X$ une immersion fermée parfaite de schémas.
Alors $i_*\mathcal{O}_Z$ est un $\mathcal{O}_X$-module parfait, c'est-à-dire
un objet parfait de $D(\mathcal{O}_X)$.
\end{lemma}

\begin{proof}
Cela résulte presque immédiatement de la définition. En effet, soit
$U = \Spec(A)$ un ouvert affine de $X$. Alors $i^{-1}(U) = \Spec(A/I)$
pour un idéal $I \subset A$, et $A/I$ admet une résolution finie par des
$A$-modules projectifs de type fini d'après
Compléments d'algèbre, lemme \ref{more-algebra-lemma-perfect-ring-map}.
Ainsi, $i_*\mathcal{O}_Z|_U$ peut être représenté par un complexe de longueur
finie de $\mathcal{O}_U$-modules localement libres de type fini.
C'est ce qu'il fallait montrer, voir
Cohomologie, section \ref{cohomology-section-perfect}.
\end{proof}

\begin{lemma}
\label{more-morphisms-lemma-perfect-proper-perfect-direct-image}
Soit $S$ un schéma noethérien. Soit $f : X \to S$ un morphisme parfait et propre
de schémas. Soit $E \in D(\mathcal{O}_X)$ parfait. Alors
$Rf_*E$ est un objet parfait de $D(\mathcal{O}_S)$.
\end{lemma}

\begin{proof}
Nous affirmons que le lemme
\ref{perfect-lemma-perfect-direct-image} de Catégories dérivées de schémas s'applique.
Les conditions (1) et (2) sont immédiates. La condition (3) est locale
sur $X$. Nous pouvons donc supposer $X$ et $S$ affines et $E$
représenté par un complexe strictement parfait de $\mathcal{O}_X$-modules.
Il suffit donc de montrer que $\mathcal{O}_X$ est de dimension de Tor finie
comme faisceau de $f^{-1}\mathcal{O}_S$-modules.
Cela équivaut à être parfait d'après le
lemme \ref{more-morphisms-lemma-check-perfect-stalks}.
\end{proof}

\begin{lemma}
\label{more-morphisms-lemma-perfect-fppf-local-source}
La propriété $\mathcal{P}(f) =$``$f$ est parfait''
est locale pour la topologie fppf sur la source.
\end{lemma}

\begin{proof}
Soit $\{g_i : X_i \to X\}_{i \in I}$ un recouvrement fppf de schémas et soit
$f : X \to S$ un morphisme tel que chaque $f \circ g_i$ soit
parfait. D'après le
lemme \ref{more-morphisms-lemma-pseudo-coherent-fppf-local-source},
nous concluons que $f$ est pseudo-cohérent.
Ainsi, d'après le
lemme \ref{more-morphisms-lemma-check-perfect-stalks},
il suffit de vérifier que $\mathcal{O}_{X, x}$ est un
$\mathcal{O}_{S, f(x)}$-module de dimension de Tor finie pour tout $x \in X$.
Choisissons $i \in I$ et $x_i \in X_i$ s'envoyant sur $x$. Nous voyons alors que
$\mathcal{O}_{X_i, x_i}$ est de dimension de Tor finie sur
$\mathcal{O}_{S, f(x)}$ et que
$\mathcal{O}_{X, x} \to \mathcal{O}_{X_i, x_i}$ est fidèlement plat.
La conclusion voulue résulte du
lemme \ref{more-algebra-lemma-flat-descent-tor-amplitude} de Compléments d'algèbre.
\end{proof}

\begin{lemma}
\label{more-morphisms-lemma-factor-regular-immersion}
Soient $i : Z \to Y$ et $j : Y \to X$ des immersions de schémas.
Supposons que
\begin{enumerate}
\item $X$ soit localement noethérien,
\item $j \circ i$ soit une immersion régulière, et
\item $i$ soit parfait.
\end{enumerate}
Alors $i$ et $j$ sont des immersions régulières.
\end{lemma}

\begin{proof}
Puisque $X$ (et donc $Y$) est localement noethérien, les quatre notions d'immersion
régulière coïncident ; de plus, nous pouvons vérifier si un morphisme est une
immersion régulière au niveau des anneaux locaux, voir
Diviseurs, lemme \ref{divisors-lemma-Noetherian-scheme-regular-ideal}.
Le résultat découle donc de
Algèbre à puissances divisées, lemme \ref{dpa-lemma-regular-sequence}.
\end{proof}









\section{Morphismes d'intersection complète locale}
\label{more-morphisms-section-lci}

\noindent
Dans
Diviseurs, section \ref{divisors-section-regular-immersions},
nous avons défini quatre types différents d'immersions régulières : régulières,
Koszul-régulières, $H_1$-régulières et quasi-régulières. Dans cette section,
nous considérons les morphismes $f : X \to S$ qui, localement sur $X$, se factorisent comme
$$
\xymatrix{
X \ar[rr]_i \ar[rd] & & \mathbf{A}^n_S \ar[ld] \\
& S
}
$$
où $i$ est une immersion $*$-régulière pour
$* \in \{\emptyset, Koszul, H_1, quasi\}$.
Cependant, nous ne savons pas démontrer que cette condition ne dépend pas
de la factorisation si $* = \emptyset$, c'est-à-dire lorsque nous demandons que $i$
soit une immersion régulière. D'autre part, nous voulons qu'un
morphisme d'intersection complète locale soit parfait, ce qui ne sera vrai
que si $* = Koszul$ ou $* = \emptyset$. Nous définirons donc un
{\it morphisme d'intersection complète locale} ou un
{\it morphisme de Koszul} comme un morphisme de schémas $f : X \to S$
qui, localement sur $X$, admet une factorisation comme ci-dessus où $i$ est une immersion
Koszul-régulière. Pour voir que cette définition convient, nous démontrons d'abord qu'elle ne dépend pas
des factorisations choisies.

\begin{lemma}
\label{more-morphisms-lemma-koszul-independence-factorization}
Soit $S$ un schéma. Soient $U$, $P$, $P'$ des schémas sur $S$.
Soit $u \in U$. Soient $i : U \to P$, $i' : U \to P'$ des immersions sur $S$.
Supposons $P$ et $P'$ lisses sur $S$. Les assertions suivantes sont alors équivalentes :
\begin{enumerate}
\item $i$ est une immersion Koszul-régulière dans un voisinage de $u$, et
\item $i'$ est une immersion Koszul-régulière dans un voisinage de $u$.
\end{enumerate}
\end{lemma}

\begin{proof}
Supposons que $i$ soit une immersion Koszul-régulière dans un voisinage de $u$.
Considérons le morphisme $j = (i, i') : U \to P \times_S P' = P''$.
Puisque $P'' = P \times_S P' \to P$ est lisse, il résulte de
Diviseurs, lemme \ref{divisors-lemma-lift-regular-immersion-to-smooth},
que $j$ est une immersion Koszul-régulière ; il résulte alors de
Diviseurs, lemme \ref{divisors-lemma-push-regular-immersion-thru-smooth},
que $i'$ est une immersion Koszul-régulière.
\end{proof}

\noindent
Avant d'énoncer la définition, faisons la remarque élémentaire
suivante. Soit $f : X \to S$ un morphisme de schémas localement
de type fini. Soit $x \in X$. Il existe alors un voisinage ouvert
$U \subset X$ et une factorisation de $f|_U$ comme composée d'une
immersion $i : U \to \mathbf{A}^n_S$ et de la projection
$\mathbf{A}^n_S \to S$, qui est lisse. Considérons le diagramme
$$
\xymatrix{
X \ar[rd] & U \ar[l] \ar[d] \ar[r]_-i & \mathbf{A}^n_S = P \ar[ld]^\pi \\
& S
}
$$
En fait, on peut procéder ainsi avec tout voisinage ouvert affine
$U$ de $x$ dans $X$ ; voir
Morphismes, lemme \ref{morphisms-lemma-quasi-affine-finite-type-over-S}.

\begin{definition}
\label{more-morphisms-definition-lci}
Soit $f : X \to S$ un morphisme de schémas.
\begin{enumerate}
\item Soit $x \in X$. On dit que $f$ est {\it de Koszul en $x$} si $f$
est de type fini en $x$ et s'il existe un voisinage ouvert
et une factorisation de $f|_U$ sous la forme $\pi \circ i$, où $i : U \to P$
est une immersion Koszul-régulière et $\pi : P \to S$ est lisse.
\item On dit que $f$ est un {\it morphisme de Koszul}, ou que
$f$ est un {\it morphisme d'intersection complète locale},
si $f$ est de Koszul en tout point.
\end{enumerate}
\end{definition}

\noindent
Nous avons vu plus haut que le choix de la factorisation
$f|_U = \pi \circ i$ est sans importance : étant donnée une factorisation
de $f|_U$ comme une immersion $i$ suivie d'un morphisme lisse $\pi$, le fait que
$i$ soit ou non Koszul-régulière dans un voisinage de $x$ est une propriété
intrinsèque de $f$ en $x$. Consignons-le ici explicitement sous forme de lemme
afin de pouvoir nous y référer.

\begin{lemma}
\label{more-morphisms-lemma-lci}
Soit $f : X \to S$ un morphisme d'intersection complète locale.
Soit $P$ un schéma lisse sur $S$. Soit $U \subset X$ un sous-schéma ouvert
et soit $i : U \to P$ une immersion de schémas sur $S$.
Alors $i$ est une immersion Koszul-régulière.
\end{lemma}

\begin{proof}
C'est la propriété définissant un morphisme d'intersection complète
locale. Voir la discussion ci-dessus.
\end{proof}

\noindent
Il paraît opportun de rassembler ici quelques propriétés communes
à tous les morphismes de Koszul.

\begin{lemma}
\label{more-morphisms-lemma-lci-properties}
Soit $f : X \to S$ un morphisme d'intersection complète locale.
Alors
\begin{enumerate}
\item $f$ est localement de présentation finie,
\item $f$ est pseudo-cohérent, et
\item $f$ est parfait.
\end{enumerate}
\end{lemma}

\begin{proof}
Puisqu'un morphisme parfait est pseudo-cohérent
(car un homomorphisme d'anneaux parfait est pseudo-cohérent)
et qu'un morphisme pseudo-cohérent est localement de présentation finie
(car un homomorphisme d'anneaux pseudo-cohérent est de présentation finie),
il suffit de démontrer la dernière assertion. La perfection étant une propriété
locale, on peut supposer que $f$ se factorise sous la forme $\pi \circ i$, où
$\pi$ est lisse et $i$ est une immersion Koszul-régulière.
Une immersion Koszul-régulière est parfaite ; voir
le lemme \ref{more-morphisms-lemma-regular-immersion-perfect}.
Un morphisme lisse est parfait puisqu'il est plat et localement de
présentation finie ; voir
le lemme \ref{more-morphisms-lemma-flat-finite-presentation-perfect}.
Enfin, une composée de morphismes parfaits est parfaite ; voir
le lemme \ref{more-morphisms-lemma-composition-perfect}.
\end{proof}

\begin{lemma}
\label{more-morphisms-lemma-affine-lci}
Soit $f : X = \Spec(B) \to S = \Spec(A)$ un morphisme de schémas affines.
Alors $f$ est un morphisme d'intersection complète locale si et seulement si
$A \to B$ est un homomorphisme d'intersection complète locale ; voir
Compléments d'algèbre, définition
\ref{more-algebra-definition-local-complete-intersection}.
\end{lemma}

\begin{proof}
Cela résulte immédiatement des définitions.
\end{proof}

\noindent
Attention : un changement de base d'un morphisme de Koszul n'est pas
en général de Koszul.

\begin{lemma}
\label{more-morphisms-lemma-flat-base-change-lci}
Un changement de base plat d'un morphisme d'intersection complète locale est
un morphisme d'intersection complète locale.
\end{lemma}

\begin{proof}
Omis. Indication : cela est vrai parce qu'un changement de base d'un morphisme lisse
est lisse et qu'un changement de base plat d'une immersion Koszul-régulière est une
immersion Koszul-régulière ; voir
Diviseurs, lemme \ref{divisors-lemma-regular-immersion-noetherian}.
\end{proof}

\begin{lemma}
\label{more-morphisms-lemma-composition-lci}
Une composée de morphismes d'intersection complète locale
est un morphisme d'intersection complète locale.
\end{lemma}

\begin{proof}
Soient $g : Y \to S$ et $f : X \to Y$ des morphismes d'intersection complète
locale. Soit $x \in X$ et posons $y = f(x)$. Choisissons un voisinage ouvert
$V \subset Y$ de $y$ et une factorisation $g|_V = \pi \circ i$, où
$i : V \to P$ est une immersion Koszul-régulière et $\pi : P \to S$ un morphisme lisse.
Choisissons ensuite un voisinage ouvert $U$ de $x \in X$ et une factorisation
$f|_U = \pi' \circ i'$, où $i' : U \to P'$ est une immersion Koszul-régulière
et $\pi' : P' \to Y$ un morphisme lisse. On peut en fait supposer que
$P' = \mathbf{A}^n_V$ ; voir les discussions qui précèdent et qui suivent la
définition \ref{more-morphisms-definition-lci}. Diagramme :
$$
\xymatrix{
X \ar[d] & U \ar[l] \ar[r]_-{i'} & P' = \mathbf{A}^n_V \ar[d] \\
Y \ar[d] &  & V \ar[ll] \ar[r]_i & P \ar[d] \\
S & & & S \ar[lll]
}
$$
Posons $P'' = \mathbf{A}^n_P$. Alors $U \to P' \to P''$ est une
immersion Koszul-régulière en tant que composée
d'immersions Koszul-régulières, à savoir $i'$ et le changement de base plat de
$i$ par $P'' \to P$ ; voir
Diviseurs,
lemme \ref{divisors-lemma-regular-immersion-noetherian}
et
Diviseurs, lemme \ref{divisors-lemma-composition-regular-immersion}.
De plus, $P'' \to P \to S$ est lisse en tant que composée de morphismes lisses ;
voir
Morphismes, lemme \ref{morphisms-lemma-composition-smooth}.
On en conclut que $X \to S$ est de Koszul en $x$, comme voulu.
\end{proof}

\begin{lemma}
\label{more-morphisms-lemma-flat-lci}
\begin{slogan}
Un morphisme est plat et d'intersection complète locale si et seulement s'il est syntomique.
\end{slogan}
Soit $f : X \to S$ un morphisme de schémas.
Les assertions suivantes sont équivalentes :
\begin{enumerate}
\item $f$ est plat et est un morphisme d'intersection complète locale, et
\item $f$ est syntomique.
\end{enumerate}
\end{lemma}

\begin{proof}
En travaillant localement dans le cas affine, c'est le contenu de
Compléments d'algèbre, lemme \ref{more-algebra-lemma-syntomic-lci}.
Donnons aussi une démonstration plus géométrique.

\medskip\noindent
Supposons (2). D'après
Morphismes, lemme \ref{morphisms-lemma-syntomic-locally-standard-syntomic},
pour tout point $x$ de $X$, il existe des voisinages ouverts affines
$U$ de $x$ et $V$ de $f(x)$ tels que $f|_U : U \to V$ soit syntomique standard.
Cela signifie que $U = \Spec(R[x_1, \ldots, x_n]/(f_1, \ldots, f_c))
\to V = \Spec(R)$, où $R[x_1, \ldots, x_n]/(f_1, \ldots, f_c)$ est une
intersection complète globale relative sur $R$. D'après
Algèbre,
lemme \ref{algebra-lemma-relative-global-complete-intersection-conormal},
la suite $f_1, \ldots, f_c$ est régulière dans tout anneau local
$R[x_1, \ldots, x_n]_{\mathfrak q}$, pour tout idéal premier
$\mathfrak q \supset (f_1, \ldots, f_c)$. Considérons le complexe de Koszul
$K_\bullet = K_\bullet(R[x_1, \ldots, x_n], f_1, \ldots, f_c)$
dont les groupes d'homologie sont $H_i = H_i(K_\bullet)$. D'après
Compléments d'algèbre, lemme \ref{more-algebra-lemma-regular-koszul-regular},
on voit que $(H_i)_{\mathfrak q} = 0$, $i > 0$, pour tout $\mathfrak q$
comme ci-dessus. D'autre part, d'après
Compléments d'algèbre, lemme \ref{more-algebra-lemma-homotopy-koszul},
on voit que $H_i$ est annulé par $(f_1, \ldots, f_c)$. Par conséquent,
$H_i = 0$, $i > 0$, et $f_1, \ldots, f_c$ est une suite Koszul-régulière.
Cela prouve que $U \to V$ se factorise en une immersion Koszul-régulière
$U \to \mathbf{A}^n_V$ suivie d'un morphisme lisse, comme voulu.

\medskip\noindent
Supposons (1). Alors $f$ est plat et localement de présentation finie
(lemme \ref{more-morphisms-lemma-lci-properties}).
Ainsi, d'après
Morphismes, lemme \ref{morphisms-lemma-syntomic-locally-standard-syntomic},
il suffit de montrer que les anneaux locaux $\mathcal{O}_{X_s, x}$
sont des anneaux d'intersection complète locale. Choisissons, localement sur $X$, une factorisation
$f = \pi \circ i$, où $i : X \to P$ est une immersion Koszul-régulière
et $\pi : P \to S$ un morphisme lisse. Remarquons que $X \to P$ est
une immersion quasi-régulière relative sur $S$ ; voir
Diviseurs, définition \ref{divisors-definition-relative-H1-regular-immersion}.
Ainsi, d'après
Diviseurs,
lemme \ref{divisors-lemma-relative-regular-immersion-flat-in-neighbourhood},
on voit que $X \to P$ est une immersion régulière et qu'il en reste de même
après tout changement de base. Chaque fibre est donc une immersion régulière, si bien que
tous les anneaux locaux de toutes les fibres de $X$ sont des intersections complètes locales.
\end{proof}

\begin{lemma}
\label{more-morphisms-lemma-regular-immersion-lci}
Une immersion régulière de schémas est un morphisme d'intersection complète locale.
Une immersion Koszul-régulière de schémas est un morphisme d'intersection complète
locale.
\end{lemma}

\begin{proof}
Puisqu'une immersion régulière est une immersion Koszul-régulière ; voir
Diviseurs, lemme \ref{divisors-lemma-regular-quasi-regular-immersion},
il suffit de démontrer la seconde assertion. Celle-ci
résulte immédiatement de la définition.
\end{proof}

\begin{lemma}
\label{more-morphisms-lemma-lci-permanence}
Soit
$$
\xymatrix{
X \ar[rr]_f \ar[rd] & & Y \ar[ld] \\
& S
}
$$
un diagramme commutatif de morphismes de schémas. Supposons que $Y \to S$
soit lisse et que $X \to S$ soit un morphisme d'intersection complète locale.
Alors $f : X \to Y$ est un morphisme d'intersection complète locale.
\end{lemma}

\begin{proof}
Cela résulte immédiatement des définitions.
\end{proof}

\begin{lemma}
\label{more-morphisms-lemma-morphism-regular-schemes-is-lci}
Soit $f : X \to Y$ un morphisme de schémas. Si $f$ est localement
de type fini et si $X$ et $Y$ sont réguliers, alors
$f$ est un morphisme d'intersection complète locale.
\end{lemma}

\begin{proof}
On peut supposer qu'il existe une factorisation $X \to \mathbf{A}^n_Y \to Y$
où la première flèche est une immersion.
Puisque $Y$ est régulier, $\mathbf{A}^n_Y$ l'est aussi d'après
Algèbre, lemme \ref{algebra-lemma-regular-goes-up}.
Ainsi, $X \to \mathbf{A}^n_Y$ est une immersion régulière d'après
Diviseurs, lemme \ref{divisors-lemma-immersion-regular-regular-immersion}.
\end{proof}

\noindent
Le lemme suivant est d'une autre nature.

\begin{lemma}
\label{more-morphisms-lemma-lci-avramov}
Soit
$$
\xymatrix{
X \ar[rr]_f \ar[rd] & & Y \ar[ld] \\
& S
}
$$
un diagramme commutatif de morphismes de schémas. Supposons que
\begin{enumerate}
\item $S$ soit localement noethérien,
\item $Y \to S$ soit localement de type fini,
\item $f : X \to Y$ soit parfait,
\item $X \to S$ soit un morphisme d'intersection complète locale.
\end{enumerate}
Alors $X \to Y$ est un morphisme d'intersection complète locale
et $Y \to S$ est de Koszul en $f(x)$ pour tout $x \in X$.
\end{lemma}

\begin{proof}
Au cours de cette démonstration, tous les schémas seront localement noethériens
et tous les anneaux seront noethériens. Nous utiliserons sans autre mention
le fait que les suites régulières et les suites Koszul-régulières coïncident dans ce
cadre ; voir Compléments d'algèbre, lemme
\ref{more-algebra-lemma-noetherian-finite-all-equivalent}.
De plus, la régularité d'un idéal (resp.\ d'un faisceau d'idéaux)
peut se vérifier sur les anneaux locaux (resp.\ sur les germes) ; voir
Algèbre, lemme \ref{algebra-lemma-regular-sequence-in-neighbourhood}
(resp.\ Diviseurs, lemme \ref{divisors-lemma-Noetherian-scheme-regular-ideal}).

\medskip\noindent
La question est locale. On peut donc supposer que $S$, $X$ et $Y$ sont
affines. Dans cette situation, on peut choisir un diagramme commutatif
$$
\xymatrix{
\mathbf{A}^{n + m}_S \ar[d] & X \ar[l] \ar[d] \\
\mathbf{A}^n_S \ar[d] & Y \ar[l] \ar[ld] \\
S
}
$$
dont les flèches horizontales sont des immersions fermées. Soit $x \in X$ un
point, et considérons le diagramme commutatif correspondant d'anneaux
locaux
$$
\xymatrix{
J \ar[r] &
\mathcal{O}_{\mathbf{A}^{n + m}_S, x} \ar[r] &
\mathcal{O}_{X, x} \\
I \ar[r] \ar[u] &
\mathcal{O}_{\mathbf{A}^n_S, f(x)} \ar[r] \ar[u] &
\mathcal{O}_{Y, f(x)} \ar[u]
}
$$
où $J$ et $I$ sont les noyaux des flèches horizontales.
Puisque $X \to S$ est un morphisme d'intersection complète locale, l'idéal
$J$ est engendré par une suite régulière. Puisque $X \to Y$ est
parfait, l'anneau $\mathcal{O}_{X, x}$ est de dimension de Tor finie sur
$\mathcal{O}_{Y, f(x)}$. On peut donc appliquer
Algèbre à puissances divisées, lemme \ref{dpa-lemma-perfect-map-ci},
pour conclure que $I$ et $J/I$ sont engendrés par des suites régulières.
D'après nos remarques initiales, cela achève la démonstration.
\end{proof}

\begin{lemma}
\label{more-morphisms-lemma-lci-to-regular}
Soit
$$
\xymatrix{
X \ar[rr]_f \ar[rd] & & Y \ar[ld] \\
& S
}
$$
un diagramme commutatif de morphismes de schémas. Supposons que
$S$ soit localement noethérien, que $Y \to S$ soit localement de type fini,
que $Y$ soit régulier et que $X \to S$ soit un morphisme d'intersection complète locale.
Alors $f : X \to Y$ est un morphisme d'intersection complète locale
et $Y \to S$ est de Koszul en $f(x)$ pour tout $x \in X$.
\end{lemma}

\begin{proof}
C'est un cas particulier du lemme \ref{more-morphisms-lemma-lci-avramov}
compte tenu du lemme \ref{more-morphisms-lemma-regular-target-perfect}
(et de Morphismes, lemme \ref{morphisms-lemma-permanence-finite-type}).
\end{proof}

\begin{lemma}
\label{more-morphisms-lemma-perfect-conormal-free-lci}
Soit $i : X \to Y$ une immersion. Si
\begin{enumerate}
\item $i$ est parfait,
\item $Y$ est localement noethérien, et
\item le faisceau conormal $\mathcal{C}_{X/Y}$ est localement libre de rang fini,
\end{enumerate}
alors $i$ est une immersion régulière.
\end{lemma}

\begin{proof}
Traduit en algèbre, c'est
Algèbre à puissances divisées, proposition \ref{dpa-proposition-regular-ideal}.
\end{proof}

\begin{lemma}
\label{more-morphisms-lemma-lci-NL}
Soit $f : X \to Y$ un morphisme d'intersection complète locale.
Alors le complexe cotangent naïf $\NL_{X/Y}$ est un objet parfait
de $D(\mathcal{O}_X)$ d'amplitude de Tor contenue dans $[-1, 0]$.
\end{lemma}

\begin{proof}
Traduit en algèbre, c'est
Compléments d'algèbre, lemme \ref{more-algebra-lemma-lci-NL}.
Pour effectuer la traduction, utiliser
les lemmes \ref{more-morphisms-lemma-affine-lci} et
\ref{more-morphisms-lemma-NL-affine}, ainsi que
Catégories dérivées des schémas, lemmes
\ref{perfect-lemma-affine-compare-bounded},
\ref{perfect-lemma-tor-dimension-affine} et
\ref{perfect-lemma-perfect-affine}.
\end{proof}

\begin{lemma}
\label{more-morphisms-lemma-perfect-NL-lci}
Soit $f : X \to Y$ un morphisme parfait de schémas localement noethériens.
Les assertions suivantes sont équivalentes :
\begin{enumerate}
\item $f$ est un morphisme d'intersection complète locale,
\item $\NL_{X/Y}$ a une amplitude de Tor contenue dans $[-1, 0]$, et
\item $\NL_{X/Y}$ est parfait d'amplitude de Tor contenue dans $[-1, 0]$.
\end{enumerate}
\end{lemma}

\begin{proof}
Traduit en algèbre, c'est
Algèbre à puissances divisées, lemme \ref{dpa-lemma-perfect-NL-lci}.
Pour effectuer la traduction, utiliser
les lemmes \ref{more-morphisms-lemma-affine-lci} et
\ref{more-morphisms-lemma-NL-affine}, ainsi que
Catégories dérivées des schémas, lemmes
\ref{perfect-lemma-affine-compare-bounded},
\ref{perfect-lemma-tor-dimension-affine} et
\ref{perfect-lemma-perfect-affine}.
\end{proof}

\begin{lemma}
\label{more-morphisms-lemma-flat-fp-NL-lci}
Soit $f : X \to Y$ un morphisme plat de présentation finie.
Les assertions suivantes sont équivalentes :
\begin{enumerate}
\item $f$ est un morphisme d'intersection complète locale,
\item $f$ est syntomique,
\item $\NL_{X/Y}$ a une amplitude de Tor contenue dans $[-1, 0]$, et
\item $\NL_{X/Y}$ est parfait d'amplitude de Tor contenue dans $[-1, 0]$.
\end{enumerate}
\end{lemma}

\begin{proof}
Traduit en algèbre, c'est
Algèbre à puissances divisées, lemme \ref{dpa-lemma-flat-fp-NL-lci}.
Pour effectuer la traduction, utiliser
les lemmes \ref{more-morphisms-lemma-affine-lci} et
\ref{more-morphisms-lemma-NL-affine}, ainsi que
Catégories dérivées des schémas, lemmes
\ref{perfect-lemma-affine-compare-bounded},
\ref{perfect-lemma-tor-dimension-affine} et
\ref{perfect-lemma-perfect-affine}.
\end{proof}

\noindent
Le lemme suivant donne une caractérisation des morphismes lisses comme morphismes
plats dont la diagonale est parfaite.

\begin{lemma}
\label{more-morphisms-lemma-smooth-diagonal-perfect}
Soit $f : X \to Y$ un morphisme de type fini entre schémas localement noethériens.
Notons $\Delta : X \to X \times_Y X$ le morphisme diagonal.
Les assertions suivantes sont équivalentes :
\begin{enumerate}
\item $f$ est lisse,
\item $f$ est plat et $\Delta : X \to X \times_Y X$ est une immersion régulière,
\item $f$ est plat et $\Delta : X \to X \times_Y X$ est un
morphisme d'intersection complète locale,
\item $f$ est plat et $\Delta : X \to X \times_Y X$ est parfait.
\end{enumerate}
\end{lemma}

\begin{proof}
Supposons (1). Alors $f$ est plat d'après
Morphismes, lemme \ref{morphisms-lemma-smooth-flat}.
Les projections $X \times_Y X \to X$ sont lisses d'après
Morphismes, lemme \ref{morphisms-lemma-base-change-smooth}. La diagonale
est donc une section d'un morphisme lisse, donc une immersion régulière ; voir
Diviseurs, lemme \ref{divisors-lemma-section-smooth-regular-immersion}.
Ainsi, (1) $\Rightarrow$ (2).
L'implication (2) $\Rightarrow$ (3) est le
lemme \ref{more-morphisms-lemma-regular-immersion-lci}.
L'implication (3) $\Rightarrow$ (4) est le
lemme \ref{more-morphisms-lemma-lci-properties}.
L'implication intéressante (4) $\Rightarrow$ (1) résulte immédiatement
d'Algèbre à puissances divisées, lemme \ref{dpa-lemma-perfect-diagonal}.
\end{proof}

\begin{lemma}
\label{more-morphisms-lemma-descending-property-lci}
La propriété $\mathcal{P}(f) =$``$f$ est un morphisme d'intersection complète
locale'' est locale pour la topologie fpqc sur la base.
\end{lemma}

\begin{proof}
Soit $f : X \to S$ un morphisme de schémas.
Soit $\{S_i \to S\}$ un recouvrement fpqc de $S$.
Supposons que chaque changement de base $f_i : X_i \to S_i$ de $f$ soit
un morphisme d'intersection complète locale.
Remarquons que cela implique notamment que $f$ est localement de type
fini ; voir
le lemme \ref{more-morphisms-lemma-lci-properties}
et
Descente, lemme \ref{descent-lemma-descending-property-locally-finite-type}.
Soit $x \in X$. Choisissons un voisinage ouvert $U$ de $x$ et
une immersion $j : U \to \mathbf{A}^n_S$ sur $S$ (voir la
discussion qui précède la
définition \ref{more-morphisms-definition-lci}).
Il faut montrer que $j$ est une immersion Koszul-régulière.
Puisque $f_i$ est un morphisme d'intersection complète locale, on voit
que le changement de base $j_i : U \times_S S_i \to \mathbf{A}^n_{S_i}$
est une immersion Koszul-régulière ; voir
le lemme \ref{more-morphisms-lemma-lci}.
Comme $\{\mathbf{A}^n_{S_i} \to \mathbf{A}^n_S\}$ est un
recouvrement fpqc, il résulte de
Descente, lemme \ref{descent-lemma-descending-property-regular-immersion},
que $j$ est une immersion Koszul-régulière, comme voulu.
\end{proof}

\begin{lemma}
\label{more-morphisms-lemma-lci-syntomic-local-source}
La propriété $\mathcal{P}(f) =$``$f$ est un morphisme d'intersection complète
locale'' est locale pour la topologie syntomique sur la source.
\end{lemma}

\begin{proof}
Nous utiliserons le critère de
Descente, lemme \ref{descent-lemma-properties-morphisms-local-source},
pour le démontrer. Il résulte des
lemmes \ref{more-morphisms-lemma-flat-lci} et
\ref{more-morphisms-lemma-composition-lci}
que la propriété d'être un morphisme d'intersection complète locale est préservée par
précomposition par des morphismes syntomiques. Il résulte clairement de la
définition \ref{more-morphisms-definition-lci}
que la propriété d'être un morphisme d'intersection complète locale est locale pour la topologie de Zariski sur la
source et le but. Ainsi, d'après le lemme susmentionné de
Descente, lemme \ref{descent-lemma-properties-morphisms-local-source},
il suffit de prouver ce qui suit : supposons que $X' \to X \to Y$ soient des
morphismes de schémas affines, avec $X' \to X$ syntomique et $X' \to Y$
un morphisme d'intersection complète locale. Alors $X \to Y$ est un morphisme d'intersection complète
locale. Pour le voir, remarquons que, dans tous les cas, $X \to Y$ est de
présentation finie d'après
Descente, lemme \ref{descent-lemma-flat-finitely-presented-permanence-algebra}.
Choisissons une immersion fermée $X \to \mathbf{A}^n_Y$. D'après
Algèbre, lemme \ref{algebra-lemma-lift-syntomic},
on peut trouver un recouvrement ouvert affine $X' = \bigcup_{i = 1, \ldots, n} X'_i$
et des morphismes syntomiques $W_i \to \mathbf{A}^n_Y$ relevant les morphismes
$X'_i \to X$, c'est-à-dire tels que l'on ait des diagrammes cartésiens
$$
\xymatrix{
X'_i \ar[d] \ar[r] & W_i \ar[d] \\
X \ar[r] & \mathbf{A}^n_Y
}
$$
Après avoir remplacé $X'$ par $\coprod X'_i$ et posé $W = \coprod W_i$,
on obtient un diagramme cartésien de schémas affines
$$
\xymatrix{
X' \ar[d] \ar[r] & W \ar[d]^h \\
X \ar[r] & \mathbf{A}^n_Y
}
$$
où $h : W \to \mathbf{A}^n_Y$ est syntomique et où $X' \to Y$ reste un morphisme d'intersection
complète locale. Puisque $W \to \mathbf{A}^n_Y$ est ouvert (voir
Morphismes, lemme \ref{morphisms-lemma-fppf-open})
et que $X' \to X$ est surjectif, on voit que $X$ est contenu dans l'image
de $W \to \mathbf{A}^n_Y$. Choisissons une immersion fermée
$W \to \mathbf{A}^{n + m}_Y$ au-dessus de $\mathbf{A}^n_Y$. Le diagramme prend alors la forme
$$
\xymatrix{
X' \ar[d] \ar[r] & W \ar[d]^h \ar[r] & \mathbf{A}^{n + m}_Y \ar[ld] \\
X \ar[r] & \mathbf{A}^n_Y
}
$$
Comme $h$ est syntomique, et donc un morphisme d'intersection complète locale (voir
ci-dessus), le morphisme $W \to \mathbf{A}^{n + m}_Y$ est une immersion Koszul-régulière.
Comme $X' \to Y$ est un morphisme d'intersection complète locale, le morphisme
$X' \to \mathbf{A}^{n + m}_Y$ est une immersion Koszul-régulière. On déduit de
Diviseurs, lemme \ref{divisors-lemma-permanence-regular-immersion},
que $X' \to W$ est une immersion Koszul-régulière. Ainsi, puisque la propriété d'être
une immersion Koszul-régulière est locale pour la topologie fpqc sur le but (voir
Descente, lemme \ref{descent-lemma-descending-property-regular-immersion}),
on conclut que $X \to \mathbf{A}^n_Y$ est une immersion Koszul-régulière,
ce qu'il fallait démontrer.
\end{proof}

\begin{lemma}
\label{more-morphisms-lemma-base-change-lci-fibres}
Soit $S$ un schéma. Soit $f : X \to Y$ un morphisme de schémas sur $S$.
Supposons que $X$ et $Y$ soient tous deux plats et localement de présentation finie sur $S$.
Alors l'ensemble
$$
\{x \in X \mid f\text{ est de Koszul en }x\}.
$$
est ouvert dans $X$, et sa formation commute à tout changement de base
$S' \to S$.
\end{lemma}

\begin{proof}
L'ensemble est ouvert par définition (voir la
définition \ref{more-morphisms-definition-lci}).
Soit $S' \to S$ un morphisme de schémas. Posons $X' = S' \times_S X$,
$Y' = S' \times_S Y$, et notons $f' : X' \to Y'$ le changement de base de $f$.
Soit $x' \in X'$ un point tel que $f'$ soit de Koszul en $x'$. Notons
$s' \in S'$, $x \in X$, $y' \in Y'$, $y \in Y$, $s \in S$ les images
de $x'$. Remarquons que $f$ est localement de présentation finie ; voir
Morphismes, lemme \ref{morphisms-lemma-finite-presentation-permanence}.
On peut donc choisir un voisinage affine $U \subset X$ de
$x$ et une immersion $i : U \to \mathbf{A}^n_Y$. Notons $U' = S' \times_S U$
et $i' : U' \to \mathbf{A}^n_{Y'}$ le changement de base de $i$.
L'hypothèse que $f'$ est de Koszul en $x'$ implique que $i'$ est une
immersion Koszul-régulière dans un voisinage de $x'$ ; voir
le lemme \ref{more-morphisms-lemma-lci}.
Le schéma $X'$ est plat et localement de
présentation finie sur $S'$ en tant que changement de base de $X$ (voir
Morphismes, lemmes \ref{morphisms-lemma-base-change-flat} et
\ref{morphisms-lemma-base-change-finite-presentation}).
Ainsi, $i'$ est une immersion $H_1$-régulière relative sur $S'$
dans un voisinage de $x'$ (voir
Diviseurs, définition \ref{divisors-definition-relative-H1-regular-immersion}).
Le changement de base $i'_{s'} : U'_{s'} \to \mathbf{A}^n_{Y'_{s'}}$ est donc
une immersion $H_1$-régulière dans un voisinage ouvert de $x'$ ; voir
Diviseurs, lemme \ref{divisors-lemma-relative-regular-immersion}
et la discussion qui suit
Diviseurs, définition \ref{divisors-definition-relative-H1-regular-immersion}.
Puisque $s' = \Spec(\kappa(s')) \to \Spec(\kappa(s)) = s$
est un morphisme surjectif, plat et universellement ouvert (voir
Morphismes, lemme \ref{morphisms-lemma-scheme-over-field-universally-open}),
on conclut que le changement de base $i_s : U_s \to \mathbf{A}^n_{Y_s}$ est une
immersion $H_1$-régulière dans un voisinage de $x$ ; voir
Descente, lemme \ref{descent-lemma-descending-property-regular-immersion}.
Enfin, remarquons que $\mathbf{A}^n_Y$ est plat et localement de
présentation finie sur $S$ ; par conséquent,
Diviseurs, lemme \ref{divisors-lemma-fibre-quasi-regular}
implique que $i$ est une immersion (Koszul-)régulière dans un voisinage de $x$
comme souhaité.
\end{proof}

\begin{lemma}
\label{more-morphisms-lemma-unramified-lci}
Soit $f : X \to Y$ un morphisme d'intersection complète locale de schémas.
Alors $f$ est non ramifié si et seulement si $f$ est formellement non ramifié et,
dans ce cas, le faisceau conormal $\mathcal{C}_{X/Y}$ est localement libre de rang fini
sur $X$.
\end{lemma}

\begin{proof}
La première assertion résulte immédiatement du
lemme \ref{more-morphisms-lemma-unramified-formally-unramified}
et du fait qu'un morphisme d'intersection complète locale est localement
de type fini. Pour calculer le faisceau conormal de $f$, choisissons, localement
sur $X$, une factorisation de $f$ sous la forme $f = p \circ i$, où $i : X \to V$
est une immersion Koszul-régulière et $V \to Y$ est lisse. D'après le
lemme \ref{more-morphisms-lemma-two-unramified-morphisms-formally-smooth},
on voit que $\mathcal{C}_{X/Y}$ est localement un facteur direct de
$\mathcal{C}_{X/V}$, qui est localement libre de rang fini puisque $i$ est une immersion
Koszul-régulière (donc quasi-régulière) ; voir
Diviseurs, lemme \ref{divisors-lemma-quasi-regular-immersion}.
\end{proof}

\begin{lemma}
\label{more-morphisms-lemma-transitivity-conormal-lci}
Soient $Z \to Y \to X$ des morphismes de schémas formellement non ramifiés.
Supposons que $Z \to Y$ soit un morphisme d'intersection complète locale.
La suite exacte
$$
0 \to i^*\mathcal{C}_{Y/X} \to
\mathcal{C}_{Z/X} \to
\mathcal{C}_{Z/Y} \to 0
$$
du
lemme \ref{more-morphisms-lemma-transitivity-conormal}
est une suite exacte courte.
\end{lemma}

\begin{proof}
La question est locale sur $Z$ ; on peut donc supposer qu'il existe une factorisation
$Z \to \mathbf{A}^n_Y \to Y$ du morphisme $Z \to Y$. On obtient alors un
diagramme commutatif
$$
\xymatrix{
Z \ar[r]_{i'} \ar@{=}[d] &
\mathbf{A}^n_Y \ar[r] \ar[d] &
\mathbf{A}^n_X \ar[d] \\
Z \ar[r]^i & Y \ar[r] & X
}
$$
Comme $Z \to Y$ est un morphisme d'intersection complète locale, on voit que
$Z \to \mathbf{A}^n_Y$ est une immersion Koszul-régulière. Par conséquent, d'après
Diviseurs, lemme \ref{divisors-lemma-transitivity-conormal-quasi-regular},
la suite
$$
0 \to (i')^*\mathcal{C}_{\mathbf{A}^n_Y/\mathbf{A}^n_X} \to
\mathcal{C}_{Z/\mathbf{A}^n_X} \to
\mathcal{C}_{Z/\mathbf{A}^n_Y} \to 0
$$
est exacte et localement scindée. Remarquons que
$i^*\mathcal{C}_{Y/X} = (i')^*\mathcal{C}_{\mathbf{A}^n_Y/\mathbf{A}^n_X}$
d'après le
lemme \ref{more-morphisms-lemma-universal-thickening-fibre-product-flat},
et remarquons que le diagramme
$$
\xymatrix{
(i')^*\mathcal{C}_{\mathbf{A}^n_Y/\mathbf{A}^n_X} \ar[r] &
\mathcal{C}_{Z/\mathbf{A}^n_X} \\
i^*\mathcal{C}_{Y/X} \ar[u]^{\cong} \ar[r] & \mathcal{C}_{Z/X} \ar[u]
}
$$
est commutatif. La flèche horizontale inférieure est donc une injection
localement scindée. Ceci démontre le lemme.
\end{proof}













\section{Suites exactes de différentielles et de faisceaux conormaux}
\label{more-morphisms-section-exact}

\noindent
Dans cette section, nous rassemblons quelques résultats sur les suites exactes de faisceaux
conormaux et de faisceaux de différentielles. En un certain sens, ce sont toutes des
réalisations du triangle des complexes cotangents associé à une paire de
morphismes de schémas composables.

\medskip\noindent
Soient $g : Z \to Y$ et $f : Y \to X$ des morphismes de schémas.
\begin{enumerate}
\item Il existe une suite exacte canonique
$$
g^*\Omega_{Y/X} \to \Omega_{Z/X} \to \Omega_{Z/Y} \to 0,
$$
voir
Morphismes, lemme \ref{morphisms-lemma-triangle-differentials}.
Si $g : Z \to Y$ est lisse ou, plus généralement, formellement lisse, alors cette
suite est une suite exacte courte ; voir
Morphismes, lemme \ref{morphisms-lemma-triangle-differentials-smooth},
ou le lemme \ref{more-morphisms-lemma-triangle-differentials-formally-smooth}.
\item Si $g$ est une immersion ou, plus généralement, est formellement non ramifié,
alors il existe une suite exacte canonique
$$
\mathcal{C}_{Z/Y} \to g^*\Omega_{Y/X} \to \Omega_{Z/X} \to 0,
$$
voir Morphismes, lemme \ref{morphisms-lemma-differentials-relative-immersion},
ou le lemme \ref{more-morphisms-lemma-universally-unramified-differentials-sequence}.
Si $f \circ g : Z \to X$ est lisse ou, plus généralement, formellement lisse, alors cette
suite est une suite exacte courte ; voir
Morphismes, lemme \ref{morphisms-lemma-differentials-relative-immersion-smooth},
ou le lemme \ref{more-morphisms-lemma-differentials-formally-unramified-formally-smooth}.
\item Si $g$ et $f \circ g$ sont des immersions ou, plus généralement,
sont formellement non ramifiés, alors il existe une suite exacte canonique
$$
\mathcal{C}_{Z/X} \to \mathcal{C}_{Z/Y} \to g^*\Omega_{Y/X} \to 0,
$$
voir
Morphismes, lemme \ref{morphisms-lemma-two-immersions}, ou le
lemme \ref{more-morphisms-lemma-two-unramified-morphisms}. Si $f : Y \to X$ est lisse ou,
plus généralement, formellement lisse, alors cette suite est une suite exacte courte ;
voir Morphismes, lemme \ref{morphisms-lemma-two-immersions-smooth}, ou le
lemme \ref{more-morphisms-lemma-two-unramified-morphisms-formally-smooth}.
\item Si $g$ et $f$ sont des immersions ou, plus généralement, sont formellement non ramifiés, alors
il existe une suite exacte canonique
$$
g^*\mathcal{C}_{Y/X} \to \mathcal{C}_{Z/X} \to \mathcal{C}_{Z/Y} \to 0.
$$
Voir
Morphismes, lemme \ref{morphisms-lemma-transitivity-conormal}, ou le
lemme \ref{more-morphisms-lemma-transitivity-conormal}.
Si $g : Z \to Y$ est une immersion régulière\footnote{Il suffit que $g$ soit
une immersion $H_1$-régulière. Remarquons qu'une immersion qui est un
morphisme d'intersection complète locale est Koszul-régulière.}, ou
plus généralement un morphisme d'intersection complète locale, alors cette
suite est une suite exacte courte ; voir
Diviseurs, lemme \ref{divisors-lemma-transitivity-conormal-quasi-regular},
ou le lemme \ref{more-morphisms-lemma-transitivity-conormal-lci}.
\end{enumerate}










\section{Morphismes faiblement \'etales}
\label{more-morphisms-section-weakly-etale}

\noindent
Un homomorphisme d'anneaux $A \to B$ est faiblement \'etale si $A \to B$ et
$B \otimes_A B \to B$ sont tous deux plats ; voir
Compléments d'algèbre, définition \ref{more-algebra-definition-weakly-etale}.
La notion analogue pour les morphismes de schémas est la suivante.

\begin{definition}
\label{more-morphisms-definition-weakly-etale}
Un morphisme de schémas $X \to Y$ est dit {\it faiblement \'etale} ou
{\it absolument plat} si $X \to Y$ ainsi que le morphisme diagonal
$X \to X \times_Y X$ sont plats.
\end{definition}

\noindent
Un morphisme \'etale est faiblement \'etale et, réciproquement,
il se trouve qu'un morphisme faiblement \'etale ressemble effectivement, dans une certaine mesure,
à un morphisme \'etale. Par exemple, si $X \to Y$ est faiblement
\'etale, alors $L_{X/Y} = 0$, comme il résulte de
Complexe cotangent, lemme \ref{cotangent-lemma-when-zero}.
Nous démontrerons plus loin un résultat très précis reliant les morphismes faiblement \'etales
aux morphismes \'etales (voir
Cohomologie pro-\'etale, section \ref{proetale-section-weakly-etale}).
Dans cette section, nous nous en tenons aux faits élémentaires.

\begin{lemma}
\label{more-morphisms-lemma-check-weakly-etale-stalks}
Soit $f : X \to Y$ un morphisme de schémas. Les assertions suivantes sont équivalentes :
\begin{enumerate}
\item $X \to Y$ est faiblement \'etale, et
\item pour tout $x \in X$, l'homomorphisme d'anneaux
$\mathcal{O}_{Y, f(x)} \to \mathcal{O}_{X, x}$ est faiblement \'etale.
\end{enumerate}
\end{lemma}

\begin{proof}
Remarquons que, sous chacune des hypothèses (1) et (2), le morphisme $f$ est plat.
On peut donc supposer que $f$ est plat. Soit $x \in X$, d'image $y = f(x)$ dans $Y$.
Il existe des homomorphismes canoniques d'anneaux
$$
\mathcal{O}_{X, x} \otimes_{\mathcal{O}_{Y, y}} \mathcal{O}_{X, x}
\longrightarrow
\mathcal{O}_{X \times_Y X, \Delta_{X/Y}(x)}
\longrightarrow
\mathcal{O}_{X, x}
$$
où le premier homomorphisme est une localisation (donc plat) et le second est une
surjection. La condition (1) signifie que la seconde flèche est plate pour tout $x$.
La condition (2) signifie que la composée est plate pour tout $x$. L'équivalence résulte donc
d'Algèbre, lemme \ref{algebra-lemma-flat-localization}, partie (2).
\end{proof}

\begin{lemma}
\label{more-morphisms-lemma-key}
Soit $X \to Y$ un morphisme de schémas tel que
$X \to X \times_Y X$ soit plat. Soit $\mathcal{F}$ un $\mathcal{O}_X$-module.
Si $\mathcal{F}$ est plat sur $Y$, alors $\mathcal{F}$ est plat sur $X$.
\end{lemma}

\begin{proof}
Soit $x \in X$, d'image $y = f(x)$ dans $Y$.
Puisque $X \to X \times_Y X$ est plat, on voit que
$\mathcal{O}_{X, x} \otimes_{\mathcal{O}_{Y, y}} \mathcal{O}_{X, x} \to
\mathcal{O}_{X, x}$ est plat. Le résultat découle donc du
lemme \ref{more-algebra-lemma-key} de Compléments d'algèbre
et des définitions.
\end{proof}

\begin{lemma}
\label{more-morphisms-lemma-weakly-etale-characterize}
Soit $f : X \to S$ un morphisme de schémas. Les assertions suivantes sont équivalentes :
\begin{enumerate}
\item Le morphisme $f$ est faiblement \'etale.
\item Pour tous ouverts affines $U \subset X$, $V \subset S$
tels que $f(U) \subset V$, l'homomorphisme d'anneaux
$\mathcal{O}_S(V) \to \mathcal{O}_X(U)$ est faiblement \'etale.
\item Il existe un recouvrement ouvert $S = \bigcup_{j \in J} V_j$
et des recouvrements ouverts $f^{-1}(V_j) = \bigcup_{i \in I_j} U_i$ tels
que chacun des morphismes $U_i \to V_j$, $j\in J, i\in I_j$,
soit faiblement \'etale.
\item Il existe un recouvrement ouvert affine $S = \bigcup_{j \in J} V_j$
et des recouvrements ouverts affines $f^{-1}(V_j) = \bigcup_{i \in I_j} U_i$ tels
que l'homomorphisme d'anneaux $\mathcal{O}_S(V_j) \to \mathcal{O}_X(U_i)$ soit
faiblement \'etale, pour tous $j\in J, i\in I_j$.
\end{enumerate}
De plus, si $f$ est faiblement \'etale, alors, pour
tous sous-schémas ouverts $U \subset X$, $V \subset S$ tels que $f(U) \subset V$,
la restriction $f|_U : U \to V$ est faiblement \'etale.
\end{lemma}

\begin{proof}
Supposons donnés des sous-schémas ouverts $U \subset X$, $V \subset S$ tels que
$f(U) \subset V$. Alors $U \times_V U \subset X \times_Y X$ est ouvert
(Schémas, lemme \ref{schemes-lemma-open-fibre-product})
et la diagonale $\Delta_{U/V}$ de $f|_U : U \to V$ est
la restriction $\Delta_{X/Y}|_U : U \to U \times_V U$.
Puisque la platitude est une propriété locale des morphismes de schémas
(Morphismes, lemme \ref{morphisms-lemma-flat-characterize}),
l'assertion finale du lemme en résulte,
ainsi que l'équivalence de (1) et (3).
Si $X$ et $Y$ sont affines, alors $X \to Y$ est faiblement \'etale
si et seulement si $\mathcal{O}_Y(Y) \to \mathcal{O}_X(X)$ est
faiblement \'etale (utiliser de nouveau
Morphismes, lemme \ref{morphisms-lemma-flat-characterize}).
Ainsi, (1) et (3) sont également équivalentes à (2) et (4).
\end{proof}

\begin{lemma}
\label{more-morphisms-lemma-composition-weakly-etale}
Soient $X \to Y \to Z$ des morphismes de schémas.
\begin{enumerate}
\item Si $X \to X \times_Y X$ et $Y \to Y \times_Z Y$ sont plats,
alors $X \to X \times_Z X$ est plat.
\item Si $X \to Y$ et $Y \to Z$ sont faiblement \'etales, alors
$X \to Z$ est faiblement \'etale.
\end{enumerate}
\end{lemma}

\begin{proof}
La partie (1) résulte de la factorisation
$$
X \to X \times_Y X \to X \times_Z X
$$
de la diagonale de $X$ sur $Z$, de l'égalité
$$
X \times_Y X = (X \times_Z X) \times_{(Y \times_Z Y)} Y,
$$
du fait qu'un changement de base d'un morphisme plat est plat, et
du fait que la composée de morphismes plats est plate
(Morphismes, lemmes \ref{morphisms-lemma-base-change-flat} et
\ref{morphisms-lemma-composition-flat}).
La partie (2) résulte de la partie (1) et du fait (que l'on vient d'utiliser)
que la composée de morphismes plats est plate.
\end{proof}

\begin{lemma}
\label{more-morphisms-lemma-base-change-weakly-etale}
Soient $X \to Y$ et $Y' \to Y$ des morphismes de schémas, et soit
$X' = Y' \times_Y X$ le changement de base de $X$.
\begin{enumerate}
\item Si $X \to X \times_Y X$ est plat, alors $X' \to X' \times_{Y'} X'$
est plat.
\item Si $X \to Y$ est faiblement \'etale, alors $X' \to Y'$ est faiblement \'etale.
\end{enumerate}
\end{lemma}

\begin{proof}
Supposons que $X \to X \times_Y X$ soit plat. Le morphisme $X' \to X' \times_{Y'} X'$
est le changement de base de $X \to X \times_Y X$ par $Y' \to Y$. Il
est donc plat d'après Morphismes, lemme \ref{morphisms-lemma-base-change-flat}.
Ceci démontre (1). La partie (2) résulte de (1) et du fait (que l'on vient d'utiliser)
qu'un changement de base d'un morphisme plat est plat.
\end{proof}

\begin{lemma}
\label{more-morphisms-lemma-go-down}
Soient $X \to Y \to Z$ des morphismes de schémas. Supposons que $X \to Y$ soit
plat et surjectif, et que $X \to X \times_Z X$ soit plat.
Alors $Y \to Y \times_Z Y$ est plat.
\end{lemma}

\begin{proof}
Considérons le diagramme commutatif
$$
\xymatrix{
X \ar[r] \ar[d] & X \times_Z X \ar[d] \\
Y \ar[r] & Y \times_Z Y
}
$$
La flèche horizontale supérieure est plate et les flèches verticales sont plates.
Ainsi, $X$ est plat sur $Y \times_Z Y$. D'après
Morphismes, lemme \ref{morphisms-lemma-flat-permanence},
on voit que $Y$ est plat sur $Y \times_Z Y$.
\end{proof}

\begin{lemma}
\label{more-morphisms-lemma-weakly-etale-formally-unramified}
Soit $f : X \to Y$ un morphisme de schémas faiblement \'etale.
Alors $f$ est formellement non ramifié, c'est-à-dire que $\Omega_{X/Y} = 0$.
\end{lemma}

\begin{proof}
Rappelons que $f$ est formellement non ramifié si et seulement si $\Omega_{X/Y} = 0$, d'après le
lemme \ref{more-morphisms-lemma-formally-unramified-differentials}.
Grâce au lemme \ref{more-morphisms-lemma-weakly-etale-characterize} et à
Morphismes, lemme \ref{morphisms-lemma-differentials-affine},
ceci résulte du cas des anneaux, qui est le
lemme \ref{more-algebra-lemma-formally-unramified} de Compléments d'algèbre.
\end{proof}

\begin{lemma}
\label{more-morphisms-lemma-when-weakly-etale}
Soit $f : X \to Y$ un morphisme de schémas. Alors $X \to Y$ est faiblement \'etale
dans chacun des cas suivants :
\begin{enumerate}
\item $X \to Y$ est un monomorphisme plat,
\item $X \to Y$ est une immersion ouverte,
\item $X \to Y$ est plat et non ramifié,
\item $X \to Y$ est \'etale.
\end{enumerate}
\end{lemma}

\begin{proof}
Si (1) est satisfaite, alors $\Delta_{X/Y}$ est un isomorphisme
(Schémas, lemme \ref{schemes-lemma-monomorphism}), donc, assurément,
$f$ est faiblement \'etale. Le cas (2) est un cas particulier de (1).
La diagonale d'un morphisme non ramifié est une immersion ouverte
(Morphismes, lemme \ref{morphisms-lemma-diagonal-unramified-morphism}),
donc elle est plate. Ainsi, un morphisme plat et non ramifié est faiblement \'etale.
Un morphisme \'etale est plat et non ramifié
(Morphismes, lemme \ref{morphisms-lemma-etale-smooth-unramified}) ;
ainsi, (4) résulte de (3).
\end{proof}

\begin{lemma}
\label{more-morphisms-lemma-reduced-goes-up}
Soit $f : X \to Y$ un morphisme de schémas.
Si $Y$ est réduit et si $f$ est faiblement \'etale, alors $X$ est réduit.
\end{lemma}

\begin{proof}
Grâce au lemme \ref{more-morphisms-lemma-weakly-etale-characterize},
ceci résulte du cas des anneaux, qui est le
lemme de Compléments d'algèbre
\ref{more-algebra-lemma-absolutely-flat-over-absolutely-flat}.
\end{proof}

\noindent
Le lemme suivant utilise un résultat non trivial sur les homomorphismes d'anneaux
faiblement \'etales.

\begin{lemma}
\label{more-morphisms-lemma-weakly-etale-strictly-henselian-local-rings}
Soit $f : X \to Y$ un morphisme de schémas.
Les assertions suivantes sont équivalentes :
\begin{enumerate}
\item $f$ est faiblement \'etale, et
\item pour $x \in X$, l'homomorphisme d'anneaux locaux
$\mathcal{O}_{Y, f(x)} \to \mathcal{O}_{X, x}$ induit un isomorphisme
sur les hensélisés stricts.
\end{enumerate}
\end{lemma}

\begin{proof}
Soit $x \in X$ un point d'image $y = f(x)$ dans $Y$.
Choisissons une clôture algébrique séparable $\kappa^{sep}$ de $\kappa(x)$.
Soit $\mathcal{O}_{X, x}^{sh}$ le hensélisé strict
correspondant à $\kappa^{sep}$, et $\mathcal{O}_{Y, y}^{sh}$
le hensélisé strict relatif à la clôture algébrique séparable
de $\kappa(y)$ dans $\kappa^{sep}$.
Considérons le diagramme commutatif
$$
\xymatrix{
\mathcal{O}_{X, x} \ar[r] & \mathcal{O}_{X, x}^{sh} \\
\mathcal{O}_{Y, y} \ar[u] \ar[r] & \mathcal{O}_{Y, y}^{sh} \ar[u]
}
$$
d'homomorphismes locaux d'anneaux locaux ; voir
Algèbre, lemme \ref{algebra-lemma-strictly-henselian-functorial}.
Puisque le hensélisé strict est une colimite filtrante d'homomorphismes d'anneaux \'etales,
Compléments d'algèbre, lemme \ref{more-algebra-lemma-when-weakly-etale},
montre que les homomorphismes horizontaux sont faiblement \'etales.
De plus, les homomorphismes horizontaux sont fidèlement plats d'après
Compléments d'algèbre, lemme \ref{more-algebra-lemma-dumb-properties-henselization}.

\medskip\noindent
Supposons $f$ faiblement \'etale. D'après le lemme \ref{more-morphisms-lemma-check-weakly-etale-stalks},
la flèche verticale de gauche est faiblement \'etale. D'après
Compléments d'algèbre, lemmes \ref{more-algebra-lemma-composition-weakly-etale} et
\ref{more-algebra-lemma-weakly-etale-permanence},
la flèche verticale de droite est faiblement \'etale. D'après
Compléments d'algèbre, théorème \ref{more-algebra-theorem-olivier},
on conclut que la flèche verticale de droite est un isomorphisme.

\medskip\noindent
Supposons que $\mathcal{O}_{Y, y}^{sh} \to \mathcal{O}_{X, x}^{sh}$ soit un isomorphisme.
Alors $\mathcal{O}_{Y, y} \to \mathcal{O}_{X, x}^{sh}$ est faiblement \'etale.
Comme $\mathcal{O}_{X, x} \to \mathcal{O}_{X, x}^{sh}$ est fidèlement
plat, on conclut que $\mathcal{O}_{Y, y} \to \mathcal{O}_{X, x}$
est faiblement \'etale d'après
Compléments d'algèbre, lemme \ref{more-algebra-lemma-go-down}.
Ainsi, (2) implique (1) d'après le lemme \ref{more-morphisms-lemma-check-weakly-etale-stalks}.
\end{proof}

\begin{lemma}
\label{more-morphisms-lemma-normal-goes-up}
Soit $f : X \to Y$ un morphisme de schémas. Si $Y$ est un schéma normal
et si $f$ est faiblement \'etale, alors $X$ est un schéma normal.
\end{lemma}

\begin{proof}
D'après Compléments d'algèbre, lemme \ref{more-algebra-lemma-henselization-normal},
un schéma $S$ est normal si et seulement si, pour tout $s \in S$,
le hensélisé strict de $\mathcal{O}_{S, s}$ est un domaine normal.
Le lemme résulte donc du
lemme \ref{more-morphisms-lemma-weakly-etale-strictly-henselian-local-rings}.
\end{proof}

\begin{lemma}
\label{more-morphisms-lemma-weakly-etale-permanence}
Soit $S$ un schéma. Soit $f : X \to Y$ un morphisme de schémas sur $S$.
Si $X$ et $Y$ sont faiblement \'etales sur $S$, alors $f$ est faiblement \'etale.
\end{lemma}

\begin{proof}
Nous utiliserons sans autre mention Morphismes, lemmes \ref{morphisms-lemma-base-change-flat} et
\ref{morphisms-lemma-composition-flat}.
Écrivons $X \to Y$ comme la composée $X \to X \times_S Y \to Y$.
Le second morphisme est plat, car c'est le changement de base du morphisme plat
$X \to S$. Le premier est le changement de base du morphisme plat
$Y \to Y \times_S Y$ par le morphisme $X \times_S Y \to Y \times_S Y$ ;
il est donc plat. Ainsi, $X \to Y$ est plat. Le morphisme
$X \times_Y X \to X \times_S X$ est une immersion.
Ainsi, le lemme \ref{more-morphisms-lemma-key} implique que, puisque
$X$ est plat sur $X \times_S X$, il s'ensuit que $X$ est
plat sur $X \times_Y X$.
\end{proof}

\noindent
Le lemme suivant est une généralisation schématique de l'observation
selon laquelle une extension de corps à la fois séparable et purement inséparable
est nécessairement un isomorphisme.

\begin{lemma}
\label{more-morphisms-lemma-weakly-etale-universal-homeomorphism}
Soit $f : X \to Y$ un morphisme de schémas. Si $f$ est faiblement \'etale et
est un homéomorphisme universel, alors c'est un isomorphisme.
\end{lemma}

\begin{proof}
Puisque $f$ est un homéomorphisme universel, la diagonale
$\Delta : X \to X \times_Y X$ est une immersion fermée surjective d'après
Morphismes, lemmes \ref{morphisms-lemma-homeomorphism-affine} et
\ref{morphisms-lemma-universally-injective}. Comme $\Delta$ est aussi
plate, on voit que $\Delta$ est nécessairement un isomorphisme d'après
Morphismes, lemme \ref{morphisms-lemma-characterize-flat-closed-immersions}.
Autrement dit, $f$ est un monomorphisme
(Schémas, lemme \ref{schemes-lemma-monomorphism}).
Puisque $f$ est un homéomorphisme universel, il est certainement
quasi-compact. Ainsi, d'après Descente, lemme
\ref{descent-lemma-flat-surjective-quasi-compact-monomorphism-isomorphism},
on trouve que $f$ est un isomorphisme.
\end{proof}

\noindent
Le résultat suivant est une généralisation faiblement \'etale de
Morphismes \'etales, lemme \ref{etale-lemma-relative-frobenius-etale}.

\begin{lemma}
\label{more-morphisms-lemma-relative-frobenius-weakly-etale}
Soit $U \to X$ un morphisme de schémas faiblement \'etale, où $X$ est un schéma
de caractéristique $p$. Alors le Frobenius relatif
$F_{U/X} : U \to U \times_{X, F_X} X$ est un isomorphisme.
\end{lemma}

\begin{proof}
Le morphisme $F_{U/X}$ est un homéomorphisme universel d'après
Variétés, lemme \ref{varieties-lemma-relative-frobenius}.
Le morphisme $F_{U/X}$ est faiblement \'etale en tant que morphisme
entre des schémas faiblement \'etales sur $X$, d'après le
lemme \ref{more-morphisms-lemma-weakly-etale-permanence}. Ainsi, $F_{U/X}$
est un isomorphisme d'après le lemme \ref{more-morphisms-lemma-weakly-etale-universal-homeomorphism}.
\end{proof}






\section{Théorème de la fibre réduite}
\label{more-morphisms-section-reduced-fibre-theorem}

\noindent
Dans cette section, nous examinons le type le plus simple de théorème parmi ceux
annoncés par le titre. Bien que la démonstration du résultat soit quelque peu
laborieuse, elle découle essentiellement de manière directe du
résultat d'Epp sur l'élimination de la ramification ; voir
Compléments d'algèbre, théorème \ref{more-algebra-theorem-epp}.

\medskip\noindent
Soit $A$ un anneau de Dedekind de corps des fractions $K$.
Soit $X$ un schéma plat et de type fini sur $A$.
Soit $L$ une extension finie de $K$. Soit $B$ la clôture intégrale
de $A$ dans $L$. Alors $B$ est un anneau de Dedekind
(Algèbre, lemme \ref{algebra-lemma-integral-closure-Dedekind}).
Soit $X_B = X \times_{\Spec(A)} \Spec(B)$ le changement de base.
Alors $X_B \to \Spec(B)$ est de type fini
(Morphismes, lemme \ref{morphisms-lemma-base-change-finite-type}).
Par conséquent, $X_B$ est noethérien
(Morphismes, lemme \ref{morphisms-lemma-finite-type-noetherian}).
La normalisation $\nu : Y \to X_B$ existe donc (voir
Morphismes, définition \ref{morphisms-definition-normalization},
ainsi que la discussion qui suit). On a le diagramme
\begin{equation}
\label{more-morphisms-equation-normalized-base-change}
\xymatrix{
Y \ar[rd] \ar[r]_\nu & X_B \ar[r] \ar[d] & X \ar[d] \\
 & \Spec(B) \ar[r] & \Spec(A)
}
\end{equation}
On appelle parfois $Y$ le {\it changement de base normalisé} de $X$.
En général, le morphisme $\nu$ n'est pas nécessairement fini. Mais si $A$ est
un anneau de Nagata (condition presque toujours satisfaite en
pratique), alors $\nu$ est fini et $Y$ est de type fini sur $B$ ; voir
Morphismes, lemmes \ref{morphisms-lemma-nagata-normalization} et
\ref{morphisms-lemma-finite-type-nagata}.

\medskip\noindent
La formation du changement de base normalisé commute à la composition.
Plus précisément, si $M/L/K$ sont des extensions finies
de corps, de clôtures intégrales $A \subset B \subset C$,
alors le changement de base normalisé $Z$ de $Y \to \Spec(B)$
relativement à $M/L$ est égal au changement de base normalisé
de $X \to \Spec(A)$ relativement à $M/K$.

\begin{theorem}
\label{more-morphisms-theorem-normalized-base-change-with-reduced-fibre}
Soit $A$ un anneau de Dedekind de corps des fractions $K$.
Soit $X$ un schéma plat et de type fini sur $A$.
Supposons que $A$ soit un anneau de Nagata.
Il existe une extension finie $L/K$ telle que
le changement de base normalisé $Y$ soit lisse sur $\Spec(B)$
en tous les points génériques de toutes les fibres.
\end{theorem}

\begin{proof}
Au cours de la démonstration, nous utiliserons à plusieurs reprises que la formation de l'ensemble des points
où un morphisme (plat et de présentation finie), tel que $X \to \Spec(A)$, est
lisse commute au changement de base ; voir
Morphismes, lemme \ref{morphisms-lemma-set-points-where-fibres-smooth}.

\medskip\noindent
Choisissons d'abord une extension finie $L/K$ telle que
$(X_L)_{red}$ soit géométriquement réduit sur $L$ ; voir
Variétés, lemme \ref{varieties-lemma-finite-extension-geometrically-reduced}.
Puisque $Y \to (X_B)_{red}$ est birationnel, on voit, en appliquant
Variétés, lemme \ref{varieties-lemma-generic-points-geometrically-reduced},
que $Y_L$ est lui aussi géométriquement réduit sur $L$.
Par conséquent, $Y_L \to \Spec(L)$ est lisse sur un ouvert dense $V \subset Y_L$ d'après
Variétés, lemme \ref{varieties-lemma-geometrically-reduced-dense-smooth-open}.
Ainsi, le lieu de lissité $U \subset Y$ du morphisme $Y \to \Spec(B)$
est ouvert (d'après Morphismes, définition \ref{morphisms-definition-smooth})
et est dense dans la fibre générique. En remplaçant $A$ par $B$ et $X$ par $Y$,
on se ramène au cas traité dans le paragraphe suivant.

\medskip\noindent
Supposons que $X$ soit normal et que le lieu de lissité $U \subset X$ de $X \to \Spec(A)$
soit dense dans la fibre générique. Cela implique que $U$ est dense dans toutes les fibres sauf
un nombre fini ; voir le lemme \ref{more-morphisms-lemma-nowhere-dense-generic-fibre}.
Soient $x_1, \ldots, x_r \in X \setminus U$ les points génériques, en nombre fini,
des composantes irréductibles de $X \setminus U$ qui sont en outre
des points génériques de composantes irréductibles de fibres de $X \to \Spec(A)$.
Posons $\mathcal{O}_i = \mathcal{O}_{X, x_i}$. Soit $A_i$ la localisation de
$A$ en l'idéal maximal correspondant à l'image de $x_i$ dans $\Spec(A)$.
D'après
Compléments d'algèbre, proposition
\ref{more-algebra-proposition-epp-essentially-finite-type},
il existe des extensions finies
$K_i/K$ qui sont des solutions pour l'extension d'anneaux de valuation
discrète $A_i \to \mathcal{O}_i$. Soit $L/K$ une extension finie
dominant toutes les extensions $K_i/K$. Alors $L/K$
est encore une solution pour $A_i \to \mathcal{O}_i$ d'après
Compléments d'algèbre, lemme \ref{more-algebra-lemma-solution-goes-up}.

\medskip\noindent
Considérons le diagramme (\ref{more-morphisms-equation-normalized-base-change})
avec l'extension $L/K$ que nous venons de construire. Notons que $U_B \subset X_B$
est lisse sur $B$, donc normal (on peut par exemple utiliser
Algèbre, lemme \ref{algebra-lemma-normal-goes-up}).
Ainsi, $Y \to X_B$ est un isomorphisme au-dessus de $U_B$.
Soit $y \in Y$ un point générique d'une composante irréductible
d'une fibre de $Y \to \Spec(B)$ située au-dessus de l'idéal maximal
$\mathfrak m \subset B$. Supposons que $y \not \in U_B$.
Alors $y$ s'envoie sur l'un des points $x_i$. Il s'ensuit que
$\mathcal{O}_{Y, y}$ est un anneau local de la clôture intégrale
de $\mathcal{O}_i$ dans $R(X) \otimes_K L$ (détails omis).
Par conséquent, puisque $L/K$ est une solution pour
$A_i \to \mathcal{O}_i$, on voit que
$B_\mathfrak m \to \mathcal{O}_{Y, y}$ est formellement lisse
pour la topologie $\mathfrak m_y$-adique
(c'est la définition du fait d'être une « solution »).
Autrement dit, $\mathfrak m\mathcal{O}_{Y, y} = \mathfrak m_y$
et l'extension des corps résiduels est séparable ; voir
Compléments d'algèbre, lemme \ref{more-algebra-lemma-extension-dvrs-formally-smooth}.
L'anneau local
de la fibre en $y$ est donc $\kappa(y)$.
Cela implique que la fibre est lisse sur $\kappa(\mathfrak m)$
en $y$, par exemple d'après Algèbre, lemme \ref{algebra-lemma-separable-smooth}.
Ceci achève la démonstration.
\end{proof}

\begin{lemma}[Variante sur les courbes]
\label{more-morphisms-lemma-normalized-base-change-with-reduced-fibre-over-curve}
Soit $f : X \to S$ un morphisme de schémas plat et de type fini.
Supposons que $S$ soit un schéma de Nagata, intègre, de corps des fonctions $K$, et
régulier de dimension $1$. Il existe alors une extension finie $L/K$
telle que, dans le diagramme
$$
\xymatrix{
Y \ar[rd]_g \ar[r]_-\nu & X \times_S T \ar[d] \ar[r] & X \ar[d]_f \\
& T \ar[r] & S
}
$$
le morphisme $g$ soit lisse en tous les points génériques des fibres. Ici,
$T$ est la normalisation de $S$ dans $\Spec(L)$ et $\nu : Y \to X \times_S T$
est la normalisation.
\end{lemma}

\begin{proof}
Choisissons un recouvrement ouvert affine fini $S = \bigcup \Spec(A_i)$.
Alors $K$ est égal au corps des fractions de $A_i$ pour tout $i$.
Posons $X_i = X \times_S \Spec(A_i)$.
Choisissons $L_i/K$ comme dans le
théorème \ref{more-morphisms-theorem-normalized-base-change-with-reduced-fibre}
pour le morphisme $X_i \to \Spec(A_i)$.
Soit $B_i \subset L_i$ la clôture intégrale de $A_i$ et
soit $Y_i$ le changement de base normalisé de $X$ par $B_i$.
Soit $L/K$ une extension finie dominant chacun des $L_i$.
Soit $T_i \subset T$ l'image réciproque de $\Spec(A_i)$.
Pour tout $i$, on obtient un diagramme commutatif
$$
\xymatrix{
g^{-1}(T_i) \ar[r] \ar[d] & Y_i \ar[r] \ar[d] & X \times_S \Spec(A_i) \ar[d] \\
T_i \ar[r] & \Spec(B_i) \ar[r] & \Spec(A_i)
}
$$
et, en fait, le carré de gauche est un changement de base normalisé,
comme on l'a expliqué au début de la section. Dans la démonstration
du théorème \ref{more-morphisms-theorem-normalized-base-change-with-reduced-fibre},
nous avons vu que le lieu de lissité de $Y \to T$ contient
l'image réciproque dans $g^{-1}(T_i)$ de l'ensemble des points
où $Y_i$ est lisse sur $B_i$. Cela démontre le lemme.
\end{proof}

\begin{lemma}[Variante avec une extension séparable]
\label{more-morphisms-lemma-normalized-base-change-with-reduced-fibre-separable}
Soit $A$ un anneau de Dedekind de corps des fractions $K$.
Soit $X$ un schéma plat et de type fini sur $A$.
Supposons que $A$ soit un anneau de Nagata et que, pour tout point générique
$\eta$ d'une composante irréductible de $X$, l'extension de corps
$\kappa(\eta)/K$ soit séparable.
Il existe alors une extension finie séparable $L/K$ telle que
le changement de base normalisé $Y$ soit lisse sur $\Spec(B)$
en tous les points génériques de toutes les fibres.
\end{lemma}

\begin{proof}
La démonstration est exactement la même que celle du
théorème \ref{more-morphisms-theorem-normalized-base-change-with-reduced-fibre},
à quelques modifications mineures près. Le changement le plus important
consiste à utiliser Compléments d'algèbre, lemme
\ref{more-algebra-lemma-epp-essentially-finite-type-separable},
au lieu de Compléments d'algèbre, proposition
\ref{more-algebra-proposition-epp-essentially-finite-type}.
Au cours de la démonstration, nous utiliserons à plusieurs reprises que la formation de l'ensemble des points
où un morphisme (plat et de présentation finie), tel que $X \to \Spec(A)$, est
lisse commute au changement de base ; voir
Morphismes, lemme \ref{morphisms-lemma-set-points-where-fibres-smooth}.

\medskip\noindent
Puisque $X$ est plat sur $A$, tout point générique $\eta$ de $X$ s'envoie sur le
point générique de $\Spec(A)$.
Après avoir remplacé $X$ par sa réduction, on peut supposer que $X$ est réduit.
Dans ce cas, $X_K$ est géométriquement réduit sur $K$
d'après Variétés, lemme \ref{varieties-lemma-generic-points-geometrically-reduced}.
Par conséquent, $X_K \to \Spec(K)$ est lisse sur un ouvert dense d'après
Variétés, lemme \ref{varieties-lemma-geometrically-reduced-dense-smooth-open}.
Ainsi, le lieu de lissité $U \subset X$ du morphisme $X \to \Spec(A)$
est ouvert (d'après Morphismes, définition \ref{morphisms-definition-smooth})
et est dense dans la fibre générique. On est ainsi ramené à la situation
du paragraphe suivant.

\medskip\noindent
Supposons que $X$ soit normal et que le lieu de lissité $U \subset X$ de $X \to \Spec(A)$
soit dense dans la fibre générique. Cela implique que $U$ est dense dans toutes les fibres sauf
un nombre fini ; voir le lemme \ref{more-morphisms-lemma-nowhere-dense-generic-fibre}.
Soient $x_1, \ldots, x_r \in X \setminus U$ les points génériques, en nombre fini,
des composantes irréductibles de $X \setminus U$ qui sont en outre
des points génériques de composantes irréductibles de fibres de $X \to \Spec(A)$.
Posons $\mathcal{O}_i = \mathcal{O}_{X, x_i}$. Observons que le corps des fractions
de $\mathcal{O}_i$ est le corps résiduel d'un point générique de $X$.
Soit $A_i$ la localisation de $A$ en l'idéal maximal correspondant
à l'image de $x_i$ dans $\Spec(A)$. On peut appliquer Compléments d'algèbre, lemme
\ref{more-algebra-lemma-epp-essentially-finite-type-separable},
et l'on trouve des extensions finies séparables $K_i/K$ qui sont des
solutions pour $A_i \to \mathcal{O}_i$. Soit $L/K$ une extension finie
séparable dominant toutes les extensions $K_i/K$.
Alors $L/K$ est encore une solution pour $A_i \to \mathcal{O}_i$ d'après
Compléments d'algèbre, lemme \ref{more-algebra-lemma-solution-goes-up}.

\medskip\noindent
Considérons le diagramme (\ref{more-morphisms-equation-normalized-base-change})
avec l'extension $L/K$ que nous venons de construire. Notons que $U_B \subset X_B$
est lisse sur $B$, donc normal (on peut par exemple utiliser
Algèbre, lemme \ref{algebra-lemma-normal-goes-up}).
Ainsi, $Y \to X_B$ est un isomorphisme au-dessus de $U_B$.
Soit $y \in Y$ un point générique d'une composante irréductible
d'une fibre de $Y \to \Spec(B)$ située au-dessus de l'idéal maximal
$\mathfrak m \subset B$. Supposons que $y \not \in U_B$.
Alors $y$ s'envoie sur l'un des points $x_i$. Il s'ensuit que
$\mathcal{O}_{Y, y}$ est un anneau local de la clôture intégrale
de $\mathcal{O}_i$ dans $R(X) \otimes_K L$ (détails omis).
Par conséquent, puisque $L/K$ est une solution pour
$A_i \to \mathcal{O}_i$, on voit que
$B_\mathfrak m \to \mathcal{O}_{Y, y}$ est formellement lisse
(c'est la définition du fait d'être une « solution »).
Autrement dit, $\mathfrak m\mathcal{O}_{Y, y} = \mathfrak m_y$
et l'extension des corps résiduels est séparable. L'anneau local
de la fibre en $y$ est donc $\kappa(y)$.
Cela implique que la fibre est lisse sur $\kappa(\mathfrak m)$
en $y$, par exemple d'après Algèbre, lemme \ref{algebra-lemma-separable-smooth}.
Ceci achève la démonstration.
\end{proof}

\begin{lemma}[Variante avec des extensions séparables sur les courbes]
\label{more-morphisms-lemma-normalized-base-change-with-reduced-fibre-over-curve-separable}
Soit $f : X \to S$ un morphisme de schémas plat et de type fini.
Supposons que $S$ soit un schéma de Nagata, intègre, de corps des fonctions $K$, et
régulier de dimension $1$. Supposons que les extensions de corps $\kappa(\eta)/K$
soient séparables pour tout point générique $\eta$ d'une composante
irréductible de $X$. Il existe alors une extension finie séparable $L/K$
telle que, dans le diagramme
$$
\xymatrix{
Y \ar[rd]_g \ar[r]_-\nu & X \times_S T \ar[d] \ar[r] & X \ar[d]_f \\
& T \ar[r] & S
}
$$
le morphisme $g$ soit lisse en tous les points génériques des fibres. Ici,
$T$ est la normalisation de $S$ dans $\Spec(L)$ et $\nu : Y \to X \times_S T$
est la normalisation.
\end{lemma}

\begin{proof}
Cela résulte du
lemme \ref{more-morphisms-lemma-normalized-base-change-with-reduced-fibre-separable}
exactement de la même manière que le
lemme \ref{more-morphisms-lemma-normalized-base-change-with-reduced-fibre-over-curve}
résulte du
théorème \ref{more-morphisms-theorem-normalized-base-change-with-reduced-fibre}.
\end{proof}








\section{Morphismes ind-quasi-affines}
\label{more-morphisms-section-ind-quasi-affine}

\noindent
Un peu de théorie qui sera utilisée plus loin.

\begin{definition}
\label{more-morphisms-definition-ind-quasi-affine}
Un schéma $X$ est dit {\it ind-quasi-affine} si tout ouvert quasi-compact de
$X$ est quasi-affine. De même, un morphisme de schémas $X \to Y$
est dit {\it ind-quasi-affine} si $f^{-1}(V)$ est ind-quasi-affine
pour tout ouvert affine $V$ de $Y$.
\end{definition}

\noindent
Un ouvert d'un schéma affine est un exemple de schéma ind-quasi-affine.
Si $X = \bigcup_{i \in I} U_i$ est une réunion d'ouverts quasi-affines telle que
deux quelconques des $U_i$ soient contenus dans un troisième, alors $X$ est ind-quasi-affine.
Un schéma ind-quasi-affine $X$ est séparé, car deux ouverts affines
$U, V$ quelconques sont contenus dans un sous-schéma ouvert séparé de $X$, à savoir
$U \cup V$. De même, un morphisme ind-quasi-affine est séparé.

\begin{lemma}
\label{more-morphisms-lemma-ind-quasi-affine-alternative-definition}
Pour un morphisme de schémas $f : X \to Y$, les assertions suivantes sont équivalentes :
\begin{enumerate}
\item $f$ est ind-quasi-affine,
\item pour tout sous-schéma ouvert affine $V \subset Y$ et
tout sous-schéma ouvert quasi-compact $U \subset f^{-1}(V)$,
le morphisme induit $U \to V$ est quasi-affine.
\item
pour un recouvrement $\{ V_j \}_{j \in J}$ de $Y$ par des sous-schémas
ouverts quasi-compacts et quasi-séparés
$V_j \subset Y$, pour tout $j \in J$ et tout sous-schéma
ouvert quasi-compact $U \subset f^{-1}(V_j)$, le morphisme induit
$U \to V_j$ est quasi-affine.
\item pour tout sous-schéma ouvert quasi-compact et quasi-séparé
$V \subset Y$ et tout sous-schéma ouvert quasi-compact
$U \subset f^{-1}(V)$, le morphisme induit $U \to V$ est quasi-affine.
\end{enumerate}
En particulier, la propriété d'être un morphisme ind-quasi-affine
est locale pour la topologie de Zariski sur la base.
\end{lemma}

\begin{proof}
L'équivalence (1) $\Leftrightarrow$ (2)
résulte des définitions et de
Morphismes, lemme \ref{morphisms-lemma-characterize-quasi-affine}.
Pour (2) $\Rightarrow$ (4), soient $U$ et $V$ comme dans (4). D'après
Schémas, lemme \ref{schemes-lemma-quasi-compact-permanence}, le
morphisme induit $U \to V$ est quasi-compact. Ainsi, pour tout ouvert
affine $V' \subset V$, le produit fibré $V' \times_V U$ est quasi-compact,
donc, d'après (2), le morphisme induit $V' \times_V U \to V'$ est quasi-affine.
Ainsi, $U \to V$ est lui aussi quasi-affine d'après
Morphismes, lemme \ref{morphisms-lemma-characterize-quasi-affine}.
Cet argument donne également (3) $\Rightarrow$ (4) : en effet, en conservant les
mêmes notations, les ouverts affines $V' \subset V$ qui sont contenus dans l'un
des $V_j$ recouvrent $V$ ; il faut donc montrer que le morphisme
quasi-compact $V' \times_V U \to V'$ est quasi-affine.
Or, d'après (3), le composé $V' \times_V U \to V' \to V_j$
est quasi-affine et, d'après
Schémas, lemme \ref{schemes-lemma-compose-after-separated}, le morphisme
$V' \to V_j$ est quasi-séparé. Ainsi, $V' \times_V U \to V'$
est quasi-affine d'après
Morphismes, lemme \ref{morphisms-lemma-quasi-affine-permanence}.
Les dernières implications (4) $\Rightarrow$ (2) et (4) $\Rightarrow$ (3)
sont évidentes.
\end{proof}

\begin{lemma}
\label{more-morphisms-lemma-ind-quasi-affine-composition}
La propriété d'être un morphisme ind-quasi-affine est stable par composition.
\end{lemma}

\begin{proof}
Soient $f : X \to Y$ et $g : Y \to Z$ des morphismes ind-quasi-affines.
Soient $V \subset Z$ et $U \subset f^{-1}(g^{-1}(V))$ des ouverts quasi-compacts,
avec $V$ de plus quasi-séparé. L'image $f(U)$ est une partie
quasi-compacte de $g^{-1}(V)$ ; elle est donc contenue dans un
ouvert quasi-compact $W \subset g^{-1}(V)$ (réunion d'un nombre fini d'ouverts affines).
On obtient une factorisation $U \to W \to V$. Le morphisme $W \to V$ est quasi-affine
d'après le lemme \ref{more-morphisms-lemma-ind-quasi-affine-alternative-definition} ; en
particulier, $W$ est quasi-séparé. Alors, encore d'après le
lemme \ref{more-morphisms-lemma-ind-quasi-affine-alternative-definition}, $U \to W$
est lui aussi quasi-affine. Par conséquent, d'après Morphismes,
lemme \ref{morphisms-lemma-composition-quasi-affine}, le composé
$U \to V$ est lui aussi quasi-affine, et il reste à appliquer
une fois de plus le lemme \ref{more-morphisms-lemma-ind-quasi-affine-alternative-definition}.
\end{proof}

\begin{lemma}
\label{more-morphisms-lemma-ind-quasi-affine-examples}
Tout morphisme quasi-affine est ind-quasi-affine.
Toute immersion est ind-quasi-affine.
\end{lemma}

\begin{proof}
La première assertion résulte immédiatement des définitions.
En particulier, les morphismes affines, comme les immersions fermées,
sont ind-quasi-affines. Ainsi, d'après le
lemme \ref{more-morphisms-lemma-ind-quasi-affine-composition}, il reste
à montrer qu'une immersion ouverte est ind-quasi-affine.
Mais cela résulte immédiatement des définitions.
\end{proof}

\begin{lemma}
\label{more-morphisms-lemma-ind-quasi-affine-permanence}
Si $f : X \to Y$ et $g : Y \to Z$ sont des morphismes de schémas
tels que $g \circ f$ soit ind-quasi-affine, alors $f$ est ind-quasi-affine.
\end{lemma}

\begin{proof}
D'après le lemme \ref{more-morphisms-lemma-ind-quasi-affine-alternative-definition}, on peut
travailler localement pour la topologie de Zariski sur $Z$, puis sur $Y$ ; on ne perd donc
rien à supposer que $Z$, puis également $Y$, est affine. Tout ouvert quasi-compact
de $X$ est alors quasi-affine, et le
lemme \ref{more-morphisms-lemma-ind-quasi-affine-alternative-definition} donne l'assertion.
\end{proof}

\begin{lemma}
\label{more-morphisms-lemma-base-change-ind-quasi-affine}
La propriété d'être ind-quasi-affine est stable par changement de base.
\end{lemma}

\begin{proof}
Soit $f : X \to Y$ un morphisme ind-quasi-affine.
Pour vérifier que tout changement de base de $f$ est ind-quasi-affine, d'après le
lemme \ref{more-morphisms-lemma-ind-quasi-affine-alternative-definition}, on peut travailler
localement pour la topologie de Zariski sur $Y$ ; on suppose donc que $Y$ est affine.
En outre, on peut aussi supposer que, dans le morphisme de changement de base
$Z \to Y$, le schéma $Z$ est affine. Le changement de base
$X \times_Y Z \to X$ est un morphisme affine ; d'après les
lemmes \ref{more-morphisms-lemma-ind-quasi-affine-composition} et
\ref{more-morphisms-lemma-ind-quasi-affine-examples},
le morphisme $X \times_Y Z \to Y$ est donc ind-quasi-affine. Alors, d'après le
lemme \ref{more-morphisms-lemma-ind-quasi-affine-permanence}, le changement
de base $X \times_Y Z \to Z$ est ind-quasi-affine, comme souhaité.
\end{proof}

\begin{lemma}
\label{more-morphisms-lemma-descending-property-ind-quasi-affine}
La propriété d'être ind-quasi-affine est locale pour la topologie fpqc sur la base.
\end{lemma}

\begin{proof}
La stabilité de l'ind-quasi-affinité par changement de base,
fournie par le lemme \ref{more-morphisms-lemma-base-change-ind-quasi-affine},
donne une implication. Pour l'autre, soit $f : X \to Y$
un morphisme de schémas et soit $\{g_i : Y_i \to Y\}$
un recouvrement fpqc tel que le changement de base $f_i : X_i \to Y_i$
soit ind-quasi-affine pour tout $i$. Il faut montrer que $f$ est ind-quasi-affine.

\medskip\noindent
D'après le lemme \ref{more-morphisms-lemma-ind-quasi-affine-alternative-definition}, on peut travailler
localement pour la topologie de Zariski sur $Y$ ; on suppose donc que $Y$ est affine.
On utilise alors la stabilité par changement de base assurée par le
lemme \ref{more-morphisms-lemma-base-change-ind-quasi-affine} pour raffiner le recouvrement
et supposer qu'il est donné par un unique morphisme affine et fidèlement plat
$g : Y' \to Y$. Pour tout ouvert quasi-compact $U \subset X$, son
changement de base par $Y'$, $U \times_Y Y' \subset X \times_Y Y'$, est lui aussi quasi-compact.
Il reste à observer que, d'après
Descente, lemme \ref{descent-lemma-descending-property-quasi-affine},
le morphisme $U \to Y$ est quasi-affine si et seulement si $U \times_Y Y' \to Y'$ l'est.
\end{proof}

\begin{lemma}
\label{more-morphisms-lemma-etale-separated-ind-quasi-affine}
Un morphisme de schémas séparé et localement quasi-fini est ind-quasi-affine.
\end{lemma}

\begin{proof}
Soit $f : X \to Y$ un morphisme de schémas séparé et localement quasi-fini.
Soient $V \subset Y$ un ouvert affine et $U \subset f^{-1}(V)$ un ouvert
quasi-compact. Il faut montrer que $U$ est quasi-affine. Comme $U \to V$ est un
morphisme de schémas séparé et quasi-fini, cela résulte du
théorème principal de Zariski. Voir le
lemme \ref{more-morphisms-lemma-quasi-finite-separated-quasi-affine}.
\end{proof}




\section{Sommes amalgamées dans la catégorie des schémas, II}
\label{more-morphisms-section-pushouts-II}

\noindent
Cette section prolonge la section \ref{more-morphisms-section-pushouts}.
Dans cette section, nous construisons des sommes amalgamées de $Y \leftarrow Z \rightarrow X$
lorsque $Z \to X$ est une immersion fermée, que $Z \to Y$ est entier et
qu'une condition supplémentaire est satisfaite.
On trouvera une discussion détaillée dans \cite{Ferrand-Conducteur}.

\begin{situation}
\label{more-morphisms-situation-pushout-along-closed-immersion-and-integral}
Ici, $S$ est un schéma, et $i : Z \to X$ et $j : Z \to Y$
sont des morphismes de schémas sur $S$. On suppose que
\begin{enumerate}
\item $i$ est une immersion fermée,
\item $j$ est un morphisme entier de schémas,
\item pour tout $y \in Y$, il existe un ouvert affine $U \subset X$
tel que $j^{-1}(\{y\}) \subset i^{-1}(U)$.
\end{enumerate}
\end{situation}

\begin{lemma}
\label{more-morphisms-lemma-prepare-pushout-along-closed-immersion-and-integral}
Dans la situation \ref{more-morphisms-situation-pushout-along-closed-immersion-and-integral},
pour tout $y \in Y$, il existe des ouverts affines $U \subset X$ et
$V \subset Y$ tels que $i^{-1}(U) = j^{-1}(V)$ et $y \in V$.
\end{lemma}

\begin{proof}
Soit $y \in Y$. Choisissons un ouvert affine $U \subset X$
tel que $j^{-1}(\{y\}) \subset i^{-1}(U)$ (ce qui est possible par hypothèse).
Choisissons un voisinage ouvert affine $V \subset Y$ de $y$
tel que $j^{-1}(V) \subset i^{-1}(U)$.
C'est possible parce que $j : Z \to Y$ est un morphisme fermé
(Morphismes, lemme \ref{morphisms-lemma-integral-universally-closed}) et que
$i^{-1}(U)$ contient la fibre au-dessus de $y$.
Comme $j$ est entier, la fibre schématique $Z_y$
est le spectre d'une algèbre entière sur un corps.
D'après Limites, lemme \ref{limits-lemma-ample-profinite-set-in-principal-affine},
on peut trouver un $\overline{f} \in \Gamma(i^{-1}(U), \mathcal{O}_{i^{-1}(U)})$
tel que $Z_y \subset D(\overline{f}) \subset j^{-1}(V)$.
Puisque $i|_{i^{-1}(U)} : i^{-1}(U) \to U$ est une immersion fermée
de schémas affines, on peut choisir un $f \in \Gamma(U, \mathcal{O}_U)$
dont la restriction à $i^{-1}(U)$ est $\overline{f}$.
Après avoir remplacé $U$ par l'ouvert principal $D(f) \subset U$,
on trouve des ouverts affines $y \in V \subset Y$ et $U \subset X$ tels que
$$
j^{-1}(\{y\}) \subset i^{-1}(U) \subset j^{-1}(V)
$$
Nous répétons maintenant (en un certain sens) l'argument. Plus précisément, nous choisissons
$g \in \Gamma(V, \mathcal{O}_V)$ tel que $y \in D(g)$ et
$j^{-1}(D(g)) \subset i^{-1}(U)$ (ce qui est possible par le même argument
que ci-dessus). On peut alors prendre $f \in \Gamma(U, \mathcal{O}_U)$
dont la restriction à $i^{-1}(U)$ est l'image réciproque de $g$
par $i^{-1}(U) \to V$ (à nouveau pour la même raison que ci-dessus).
On dispose finalement d'ouverts affines $y \in V' = D(g) \subset V \subset Y$
et $U' = D(f) \subset U \subset X$ tels que $j^{-1}(V') = i^{-1}(V')$.
\end{proof}

\begin{proposition}
\label{more-morphisms-proposition-pushout-along-closed-immersion-and-integral}
\begin{reference}
\cite[Théorème 7.1, partie iii]{Ferrand-Conducteur}
\end{reference}
Dans la situation \ref{more-morphisms-situation-pushout-along-closed-immersion-and-integral},
la somme amalgamée $Y \amalg_Z X$ existe dans la catégorie des schémas. On a le diagramme
$$
\xymatrix{
Z \ar[r]_i \ar[d]_j & X \ar[d]^a \\
Y \ar[r]^-b & Y \amalg_Z X
}
$$
Le diagramme est un carré cartésien, le morphisme $a$ est entier,
le morphisme $b$ est une immersion fermée, et
$$
\mathcal{O}_{Y \amalg_Z X} =
b_*\mathcal{O}_Y \times_{c_*\mathcal{O}_Z} a_*\mathcal{O}_X
$$
comme faisceaux d'anneaux, où $c = a \circ i = b \circ j$.
\end{proposition}

\begin{proof}
Comme espace topologique, on définit $Y \amalg_Z X$ comme la somme amalgamée du
diagramme dans la catégorie des espaces topologiques (Topologie, section
\ref{topology-section-colimits}). Il s'agit simplement de la somme amalgamée
des ensembles sous-jacents (Topologie, lemme \ref{topology-lemma-colimits}),
munie de la topologie quotient.
Sur $Y \amalg_Z X$, on dispose des morphismes de faisceaux d'anneaux
$$
b_*\mathcal{O}_Y \longrightarrow c_*\mathcal{O}_Z
\longleftarrow a_*\mathcal{O}_X
$$
et l'on peut définir
$$
\mathcal{O}_{Y \amalg_Z X} =
b_*\mathcal{O}_Y \times_{c_*\mathcal{O}_Z} a_*\mathcal{O}_X
$$
comme le produit fibré dans la catégorie des faisceaux d'anneaux. Pour montrer que l'on
obtient un schéma, il faut montrer que tout point possède un
voisinage ouvert affine. C'est clair pour les points qui ne sont pas dans l'image de $c$,
car l'image de $c$ est une partie fermée dont le complément est
isomorphe, comme espace annelé, à $(Y \setminus j(Z)) \amalg (X \setminus i(Z))$.

\medskip\noindent
Un point de l'image de $c$ correspond à un unique $y \in Y$
dans l'image de $j$. D'après le
lemme \ref{more-morphisms-lemma-prepare-pushout-along-closed-immersion-and-integral},
on trouve des ouverts affines $U \subset X$ et $V \subset Y$ tels que
$y \in V$ et $i^{-1}(U) = j^{-1}(V)$.
Comme la construction du premier paragraphe est manifestement compatible
avec la restriction à des sous-schémas ouverts compatibles, pour montrer qu'elle
produit un schéma, on peut supposer que $X$, $Y$ et $Z$ sont affines.

\medskip\noindent
Si $X = \Spec(A)$, $Y = \Spec(B)$ et $Z = \Spec(C)$ sont affines, alors
Compléments d'algèbre, lemme \ref{more-algebra-lemma-points-of-fibre-product},
montre que $Y \amalg_Z X = \Spec(B \times_C A)$ comme espaces topologiques.
Pour achever la démonstration du fait que $Y \times_Z X$ est un schéma, il suffit de montrer
que, sur $\Spec(B \times_C A)$, le faisceau structural est le produit fibré
des images directes. Cela résulte de l'application de
Compléments d'algèbre, lemme \ref{more-algebra-lemma-diagram-localize},
aux ouverts affines principaux de $\Spec(B \times_C A)$.

\medskip\noindent
La discussion ci-dessus montre que le schéma $Y \amalg_Z X$
possède un recouvrement ouvert affine $Y \amalg_Z X = \bigcup W_i$
tel que $U_i = a^{-1}(W_i)$, $V_i = b^{-1}(W_i)$ et
$\Omega_i = c^{-1}(W_i)$ soient des ouverts affines de $X$, $Y$ et $Z$.
Ainsi, $a$ et $b$ sont affines.
De plus, si $A_i$, $B_i$, $C_i$ sont les anneaux correspondant à
$U_i$, $V_i$, $\Omega_i$, alors $A_i \to C_i$ est surjectif et
$W_i$ correspond à $A_i \times_{C_i} B_i$, qui s'envoie surjectivement sur
$B_i$. Ainsi, $b$ est une immersion fermée.
Le morphisme d'anneaux $A_i \times_{C_i} B_i \to A_i$ est entier d'après
Compléments d'algèbre, lemme \ref{more-algebra-lemma-fibre-product-integral} ;
donc $a$ est entier. Le diagramme est cartésien parce que
$$
C_i \cong B_i \otimes_{B_i \times_{C_i} A_i} A_i
$$
Cela résulte du fait que $B_i \times_{C_i} A_i \to B_i$
et $A_i \to C_i$ sont des morphismes surjectifs de même noyau.

\medskip\noindent
Enfin, on peut appliquer les lemmes \ref{more-morphisms-lemma-basic-example-pushout} et
\ref{more-morphisms-lemma-pushout-fpqc-local} pour conclure que notre construction est une somme amalgamée
dans la catégorie des schémas.
\end{proof}

\begin{lemma}
\label{more-morphisms-lemma-pushout-separated}
Dans la situation \ref{more-morphisms-situation-pushout-along-closed-immersion-and-integral},
si $X$ et $Y$ sont séparés, alors la somme amalgamée $Y \amalg_Z X$
(proposition \ref{more-morphisms-proposition-pushout-along-closed-immersion-and-integral})
est séparée. Il en va de même pour « séparé sur $S$ », « quasi-séparé » et
« quasi-séparé sur $S$ ».
\end{lemma}

\begin{proof}
Le morphisme $Y \amalg X \to Y \amalg_Z X$ est surjectif
et universellement fermé. On peut donc appliquer
Morphismes, lemme \ref{morphisms-lemma-image-universally-closed-separated}.
\end{proof}

\begin{lemma}
\label{more-morphisms-lemma-pushout-finite-type}
Dans la situation \ref{more-morphisms-situation-pushout-along-closed-immersion-and-integral},
supposons que $S$ soit un schéma localement noethérien et que $X$, $Y$ et $Z$
soient localement de type fini sur $S$. Alors la somme amalgamée $Y \amalg_Z X$
(proposition \ref{more-morphisms-proposition-pushout-along-closed-immersion-and-integral})
est localement de type fini sur $S$.
\end{lemma}

\begin{proof}
En se plaçant sur des ouverts affines, on retrouve le résultat de
Compléments d'algèbre, lemme \ref{more-algebra-lemma-fibre-product-finite-type}.
\end{proof}

\begin{lemma}
\label{more-morphisms-lemma-pushout-functor}
Dans la situation \ref{more-morphisms-situation-pushout-along-closed-immersion-and-integral},
supposons donné un diagramme commutatif
$$
\xymatrix{
Y' \ar[d]^g & Z' \ar[l]^{j'} \ar[r]_{i'} \ar[d]^h & X' \ar[d]^f \\
Y & Z \ar[l] \ar[r] & X
}
$$
dont les carrés sont cartésiens et où $f, g, h$ sont séparés et localement quasi-finis. Alors
\begin{enumerate}
\item les sommes amalgamées $Y \amalg_Z X$ et $Y' \amalg_{Z'} X'$ existent,
\item $Y' \amalg_{Z'} X' \to Y \amalg_Z X$ est
séparé et localement quasi-fini, et
\item les carrés
$$
\xymatrix{
Y' \ar[r] \ar[d] & Y' \amalg_{Z'} X' \ar[d] & X' \ar[l] \ar[d] \\
Y \ar[r] & Y \amalg_Z X & X \ar[l]
}
$$
sont cartésiens.
\end{enumerate}
\end{lemma}

\begin{proof}
La somme amalgamée $Y \amalg_Z X$ existe d'après la
proposition \ref{more-morphisms-proposition-pushout-along-closed-immersion-and-integral}.
Pour voir que la somme amalgamée $Y' \amalg_{Z'} X'$ existe, vérifions que la
condition (3) de la
situation \ref{more-morphisms-situation-pushout-along-closed-immersion-and-integral}
est satisfaite pour $(X', Y', Z', i', j')$.
En effet, soit $y' \in Y'$ et notons $y \in Y$ son image.
Choisissons un ouvert affine $U \subset X$ tel que $i(j^{-1}(y)) \subset U$.
Choisissons un ouvert quasi-compact $U' \subset X'$ contenu dans
$f^{-1}(U)$ et contenant le sous-ensemble quasi-compact $i'((j')^{-1}(\{y'\}))$.
D'après le lemme \ref{more-morphisms-lemma-etale-separated-ind-quasi-affine},
on voit que $U'$ est quasi-affine. Puisque $Z'_{y'}$ est le spectre
d'une algèbre entière sur un corps, on peut appliquer
Limites, lemme \ref{limits-lemma-ample-profinite-set-in-principal-affine},
et l'on trouve qu'il existe un sous-schéma ouvert affine de $U'$ contenant
$i'((j')^{-1}(\{y'\}))$, comme souhaité.

\medskip\noindent
L'existence étant vérifiée, établissons les autres assertions.
Localement pour la topologie affine, nous sommes exactement dans la situation de Compléments d'algèbre, lemme
\ref{more-algebra-lemma-properties-algebras-over-fibre-product},
avec $B \to D$ et $A' \to C'$ localement quasi-finis\footnote{Plus précisément,
$X, Y, Z, Y \amalg_Z X, X', Y', Z', Y' \amalg_{Z'} X'$
correspondent à $A', B, A, B', C', D, C, D'$.}.
En particulier, le morphisme $Y' \amalg_{Z'} X' \to Y \amalg_Z X$ est localement
de type fini. Les carrés du diagramme sont cartésiens d'après
Compléments d'algèbre, lemme \ref{more-algebra-lemma-module-over-fibre-product}.
Puisque la propriété d'être localement quasi-fini se vérifie sur les fibres
(Morphismes, lemme \ref{morphisms-lemma-quasi-finite-at-point-characterize}),
on conclut que $Y' \amalg_{Z'} X' \to Y \amalg_Z X$
est localement quasi-fini.

\medskip\noindent
Il reste à vérifier que $Y' \amalg_{Z'} X' \to Y \amalg_Z X$ est séparé.
Remarquons que $Y' \amalg X' \to Y' \amalg_{Z'} X'$ est universellement fermé
et surjectif d'après la
proposition \ref{more-morphisms-proposition-pushout-along-closed-immersion-and-integral}.
Comme le morphisme $Y' \amalg X' \to Y \amalg_Z X$ est aussi séparé
(puisqu'il se factorise sous la forme $Y' \amalg X' \to Y \amalg X \to Y \amalg_Z X$),
on conclut par
Morphismes, lemme \ref{morphisms-lemma-image-universally-closed-separated}.
\end{proof}

\begin{lemma}
\label{more-morphisms-lemma-pushout-functor-equivalence-flat}
Dans la situation \ref{more-morphisms-situation-pushout-along-closed-immersion-and-integral},
la catégorie des schémas plats, séparés et localement quasi-finis
sur la somme amalgamée $Y \amalg_Z X$ est équivalente à la catégorie des
$(X', Y', Z', i', j', f, g, h)$ comme dans le lemme \ref{more-morphisms-lemma-pushout-functor},
où $f, g, h$ sont plats. Même énoncé en remplaçant « plat » par
« \'etale ».
\end{lemma}

\begin{proof}
Si l'on part de $(X', Y', Z', i', j', f, g, h)$
comme dans le lemme \ref{more-morphisms-lemma-pushout-functor}, avec $f, g, h$ plats
ou \'etales, alors $Y' \amalg_{Z'} X' \to Y \amalg_Z X$ est plat
ou \'etale d'après Compléments d'algèbre, lemme
\ref{more-algebra-lemma-properties-algebras-over-fibre-product}.

\medskip\noindent
Réciproquement, soit
$W \to Y \amalg_Z X$ un morphisme séparé et localement quasi-fini.
Posons $X' = W \times_{Y \amalg_Z X} X$, $Y' = W \times_{Y \amalg_Z X} Y$ et
$Z' = W \times_{Y \amalg_Z X} Z$, avec les morphismes évidents $i', j', f, g, h$.
Formons la somme amalgamée $Y' \amalg_{Z'} X'$. On obtient un morphisme
$$
Y' \amalg_{Z'} X' \longrightarrow W
$$
de schémas au-dessus de $Y \amalg_X Z$ par la propriété universelle de la
somme amalgamée. Si l'on ne suppose pas que $W \to Y \amalg_Z X$ est plat,
alors ce morphisme n'est en général pas un isomorphisme.
(En fait, Compléments d'algèbre, lemme
\ref{more-algebra-lemma-module-over-fibre-product-bis}
montre que la flèche affichée est une immersion fermée, mais pas un isomorphisme
en général.) Cependant, si $W \to Y \times_Z X$ est plat, alors
c'est un isomorphisme d'après Compléments d'algèbre, lemme
\ref{more-algebra-lemma-properties-algebras-over-fibre-product}.
\end{proof}

\noindent
Nous examinons ensuite l'existence dans le cas où
les deux morphismes sont des immersions fermées.

\begin{lemma}
\label{more-morphisms-lemma-pushout-along-closed-immersions}
Soient $i : Z \to X$ et $j : Z \to Y$ des immersions fermées de schémas.
Alors la somme amalgamée $Y \amalg_Z X$ existe dans la catégorie des schémas. Diagramme :
$$
\xymatrix{
Z \ar[r]_i \ar[d]_j & X \ar[d]^a \\
Y \ar[r]^-b & Y \amalg_Z X
}
$$
Le diagramme est un carré cartésien, les morphismes $a$ et $b$
sont des immersions fermées, et l'on a une suite exacte courte
$$
0 \to \mathcal{O}_{Y \amalg_Z X} \to
a_*\mathcal{O}_X \oplus b_*\mathcal{O}_Y \to
c_*\mathcal{O}_Z \to 0
$$
où $c = a \circ i = b \circ j$.
\end{lemma}

\begin{proof}
C'est un cas particulier de la
proposition \ref{more-morphisms-proposition-pushout-along-closed-immersion-and-integral}.
Remarquons que l'hypothèse (3) de la
situation \ref{more-morphisms-situation-pushout-along-closed-immersion-and-integral}
est immédiate, car
les fibres de $j$ ont au plus un point. Enfin, inversons les rôles des flèches
pour conclure que $a$ et $b$ sont tous deux des immersions fermées.
\end{proof}

\begin{lemma}
\label{more-morphisms-lemma-pushout-along-closed-immersions-properties-above}
Soient $i : Z \to X$ et $j : Z \to Y$ des immersions fermées de schémas.
Soient $f : X' \to X$ et $g : Y' \to Y$ des morphismes de schémas, et soit
$\varphi : X' \times_{X, i} Z \to Y' \times_{Y, j} Z$
un isomorphisme de schémas au-dessus de $Z$. Considérons le morphisme
$$
h :
X' \amalg_{X' \times_{X, i} Z, \varphi} Y'
\longrightarrow
X \amalg_Z Y
$$
Alors :
\begin{enumerate}
\item $h$ est localement de type fini si et seulement si $f$ et $g$ sont
localement de type fini ;
\item $h$ est plat si et seulement si $f$ et $g$ sont plats ;
\item $h$ est plat et localement de présentation finie si et seulement si
$f$ et $g$ sont plats et localement de présentation finie ;
\item $h$ est lisse si et seulement si $f$ et $g$ sont lisses ;
\item $h$ est \'etale si et seulement si $f$ et $g$ sont \'etales ; et
\item ajouter d'autres propriétés ici si nécessaire.
\end{enumerate}
\end{lemma}

\begin{proof}
On sait que les sommes amalgamées existent d'après le
lemme \ref{more-morphisms-lemma-pushout-along-closed-immersions}.
En particulier, on obtient le morphisme $h$.
On peut donc remplacer tous les schémas considérés par des
schémas affines. Dans ce cas, les assertions du lemme
sont équivalentes aux assertions correspondantes de
Compléments d'algèbre, lemme
\ref{more-algebra-lemma-properties-algebras-over-fibre-product}.
\end{proof}





\section{Morphismes relatifs}
\label{more-morphisms-section-relative-morphisms}

\noindent
Dans cette section, nous démontrons un résultat de représentabilité que nous utiliserons dans
Groupes fondamentaux, section \ref{pione-section-finite-etale},
pour démontrer un résultat sur la catégorie des revêtements \'etales finis d'un schéma.
Le contenu de cette section est traité dans la généralité appropriée dans
Critères de représentabilité, section
\ref{criteria-section-relative-morphisms}.

\medskip\noindent
Soit $S$ un schéma. Soient $Z$ et $X$ des schémas sur $S$.
Étant donné un schéma $T$ sur $S$, on peut considérer les morphismes
$b : T \times_S Z \to T \times_S X$ au-dessus de $S$. Diagramme :
\begin{equation}
\label{more-morphisms-equation-hom}
\vcenter{
\xymatrix{
T \times_S Z \ar[rd] \ar[rr]_b & &
T \times_S X \ar[ld] & Z \ar[rd] & & X \ar[ld] \\
& T \ar[rrr] & & & S
}
}
\end{equation}
Bien entendu, on peut aussi voir $b$ comme un morphisme
$b : T \times_S Z \to X$ tel que le diagramme
$$
\xymatrix{
T \times_S Z \ar[r] \ar[d] \ar@/^1pc/[rrr]_-b &
Z \ar[rd] & & X \ar[ld] \\
T \ar[rr] & & S
}
$$
soit commutatif. Dans cette situation, on peut définir un foncteur
\begin{equation}
\label{more-morphisms-equation-hom-functor}
\mathit{Mor}_S(Z, X) : (\Sch/S)^{opp} \longrightarrow \textit{Ensembles},
\quad
T \longmapsto \{b\text{ comme ci-dessus}\}
\end{equation}
Voici un résultat élémentaire de représentabilité.

\begin{lemma}
\label{more-morphisms-lemma-hom-from-finite-free-into-affine}
Soient $Z \to S$ et $X \to S$ des morphismes de schémas affines.
Supposons que $\Gamma(Z, \mathcal{O}_Z)$ soit un
$\Gamma(S, \mathcal{O}_S)$-module libre de rang fini. Alors $\mathit{Mor}_S(Z, X)$
est représentable par un schéma affine sur $S$.
\end{lemma}

\begin{proof}
Écrivons $S = \Spec(R)$. Choisissons une base $\{e_1, \ldots, e_m\}$
de $\Gamma(Z, \mathcal{O}_Z)$ sur $R$. Choisissons une présentation
$$
\Gamma(X, \mathcal{O}_X) = R[\{x_i\}_{i \in I}]/(\{f_k\}_{k \in K}).
$$
Nous noterons $\overline{x}_i$ l'image de $x_i$ dans ce quotient.
Écrivons
$$
P = R[\{a_{ij}\}_{i \in I, 1 \leq j \leq m}].
$$
Considérons l'homomorphisme de $R$-algèbres
$$
\Psi :
R[\{x_i\}_{i \in I}]
\longrightarrow
P \otimes_R \Gamma(Z, \mathcal{O}_Z), \quad
x_i \longmapsto \sum\nolimits_j a_{ij} \otimes e_j.
$$
Écrivons $\Psi(f_k) = \sum c_{kj} \otimes e_j$ avec $c_{kj} \in P$.
Enfin, notons $J \subset P$ l'idéal engendré par les éléments
$c_{kj}$, $k \in K$, $1 \leq j \leq m$. Nous affirmons que
$W = \Spec(P/J)$ représente le foncteur $\mathit{Mor}_S(Z, X)$.

\medskip\noindent
Remarquons d'abord que, par construction, $P/J$ est une $R$-algèbre, donc définit
un morphisme $W \to S$. Ensuite, par construction, l'homomorphisme
$\Psi$ se factorise par $\Gamma(X, \mathcal{O}_X)$ ; on obtient donc
un homomorphisme de $P/J$-algèbres
$$
P/J \otimes_R \Gamma(X, \mathcal{O}_X)
\longrightarrow
P/J \otimes_R \Gamma(Z, \mathcal{O}_Z)
$$
qui détermine un morphisme
$b_{univ} : W \times_S Z \to W \times_S X$.
Par le lemme de Yoneda, $b_{univ}$ détermine une
transformation de foncteurs $W \to \mathit{Mor}_S(Z, X)$ dont nous
affirmons que c'est un isomorphisme. Pour le montrer, il suffit
de montrer qu'elle induit une bijection d'ensembles
$W(T) \to \mathit{Mor}_S(Z, X)(T)$ pour tout schéma affine
$T$.

\medskip\noindent
Supposons que $T = \Spec(R')$ soit un schéma affine sur $S$
et que $b \in \mathit{Mor}_S(Z, X)(T)$. Le morphisme structural $T \to S$
définit une structure de $R$-algèbre sur $R'$, et $b$ définit un homomorphisme de $R'$-algèbres
$$
b^\sharp :
R' \otimes_R \Gamma(X, \mathcal{O}_X)
\longrightarrow
R' \otimes_R \Gamma(Z, \mathcal{O}_Z).
$$
En particulier, on peut écrire
$b^\sharp(1 \otimes \overline{x}_i) = \sum \alpha_{ij} \otimes e_j$
pour certains $\alpha_{ij} \in R'$. Cela correspond à un homomorphisme de $R$-algèbres
$P \to R'$ déterminé par la règle $a_{ij} \mapsto \alpha_{ij}$. Cet
homomorphisme se factorise par le quotient $P/J$ par construction de l'idéal
$J$, ce qui donne un homomorphisme $P/J \to R'$. Celui-ci correspond à son tour à un morphisme
$T \to W$ tel que $b$ soit le changement de base de $b_{univ}$.
Certains détails sont omis.
\end{proof}

\begin{lemma}
\label{more-morphisms-lemma-hom-from-finite-locally-free-into-affine}
Soient $Z \to S$ et $X \to S$ des morphismes de schémas.
Si $Z \to S$ est fini localement libre et $X \to S$ est affine,
alors $\mathit{Mor}_S(Z, X)$ est représentable par un schéma
affine sur $S$.
\end{lemma}

\begin{proof}
Choisissons un recouvrement ouvert affine $S = \bigcup U_i$ tel que
$\Gamma(Z \times_S U_i, \mathcal{O}_{Z \times_S U_i})$ soit
libre de rang fini sur $\mathcal{O}_S(U_i)$. Soit $F_i \subset \mathit{Mor}_S(Z, X)$
le sous-foncteur qui associe à $T/S$ l'ensemble vide si
$T \to S$ ne se factorise pas par $U_i$, et $\mathit{Mor}_S(Z, X)(T)$
sinon. Alors la famille de ces sous-foncteurs satisfait les conditions
(2)(a), (2)(b), (2)(c) de
Schémas, lemme \ref{schemes-lemma-glue-functors}, ce qui démontre le lemme.
La condition (2)(a) résulte du
lemme \ref{more-morphisms-lemma-hom-from-finite-free-into-affine},
et les deux autres résultent d'arguments immédiats.
\end{proof}

\noindent
La condition imposée au morphisme $f : X \to S$ dans le lemme ci-dessous est très
utile pour démontrer des énoncés de ce type. Elle est satisfaite si l'une des conditions suivantes
l'est : $X$ est quasi-affine, $f$ est quasi-affine, $f$ est quasi-projectif,
$f$ est localement projectif, il existe un faisceau inversible ample sur $X$,
il existe un faisceau inversible $f$-ample sur $X$, ou
il existe un faisceau inversible $f$-très ample sur $X$.

\begin{lemma}
\label{more-morphisms-lemma-hom-from-finite-locally-free-representable}
Soient $Z \to S$ et $X \to S$ des morphismes de schémas.
Supposons que
\begin{enumerate}
\item $Z \to S$ soit fini localement libre ; et
\item pour tout $(s, x_1, \ldots, x_d)$ où $s \in S$ et
$x_1, \ldots, x_d \in X_s$, il existe un ouvert affine $U \subset X$
tel que $x_1, \ldots, x_d \in U$.
\end{enumerate}
Alors $\mathit{Mor}_S(Z, X)$ est représentable par un schéma.
\end{lemma}

\begin{proof}
Considérons l'ensemble $I$ des couples $(U, V)$ où $U \subset X$ et $V \subset S$
sont des ouverts affines et où $U \to S$ se factorise par $V$. Pour $i \in I$, notons
$(U_i, V_i)$ le couple correspondant. Posons
$F_i = \mathit{Mor}_{V_i}(Z_{V_i}, U_i)$.
Il est immédiat que $F_i$ est un sous-foncteur de $\mathit{Mor}_S(Z, X)$.
Nous affirmons alors que les conditions
(2)(a), (2)(b), (2)(c) de
Schémas, lemme \ref{schemes-lemma-glue-functors}, sont satisfaites, ce qui démontre le lemme.

\medskip\noindent
La condition (2)(a) résulte du
lemme \ref{more-morphisms-lemma-hom-from-finite-locally-free-into-affine}.

\medskip\noindent
Pour vérifier la condition (2)(b), considérons $T/S$ et $b \in \mathit{Mor}_S(Z, X)(T)$.
En considérant $b$ comme un morphisme $T \times_S Z \to X$, on obtient un ouvert
$b^{-1}(U_i) \subset T \times_S Z$. Il est clair que $b \in F_i(T)$
si et seulement si $b^{-1}(U_i) = T \times_S Z$. Puisque la projection
$p : T \times_S Z \to T$ est finie, donc fermée, l'ensemble
$U_{i, b} \subset T$ des points $t \in T$ tels que
$p^{-1}(\{t\}) \subset b^{-1}(U_i)$ est ouvert.
Alors $f : T' \to T$ se factorise par $U_{i, b}$ si et seulement
si $b \circ f \in F_i(T')$, et la vérification de (2)(b) est achevée.

\medskip\noindent
Enfin, vérifions la condition (2)(c) ; c'est ici qu'intervient notre condition
sur $X \to S$. En effet, considérons
$T/S$ et $b \in \mathit{Mor}_S(Z, X)(T)$.
Il suffit de démontrer que tout $t \in T$
appartient à l'un des ouverts $U_{i, b}$ définis
au paragraphe précédent.
Ceci équivaut à la condition que
$b(p^{-1}(\{t\})) \subset U_i$ pour un certain $i$
où $p : T \times_S Z \to T$ est la projection et
$b : T \times_S Z \to X$ est le morphisme donné.
Puisque $p$ est fini, l'ensemble $b(p^{-1}(\{t\})) \subset X$
est fini et contenu dans la fibre de $X \to S$ au-dessus de
l'image $s$ de $t$ dans $S$.
Ainsi, notre condition sur $X \to S$ montre exactement qu'un
couple convenable existe.
\end{proof}

\begin{lemma}
\label{more-morphisms-lemma-hom-from-finite-locally-free-separated-lqf}
Soient $Z \to S$ et $X \to S$ des morphismes de schémas.
Supposons que $Z \to S$ soit fini localement libre et que $X \to S$
soit séparé et localement quasi-fini.
Alors $\mathit{Mor}_S(Z, X)$ est représentable par un schéma.
\end{lemma}

\begin{proof}
Ceci résulte des
lemmes \ref{more-morphisms-lemma-hom-from-finite-locally-free-representable} et
\ref{more-morphisms-lemma-separated-locally-quasi-finite-over-affine}.
\end{proof}









\section{Caractérisation des complexes pseudo-cohérents, III}
\label{more-morphisms-section-characterize-pseudo-coherent}

\noindent
Dans cette section, nous étudions des caractérisations des complexes pseudo-cohérents
en termes de cohomologie. Ceci fait suite à
Catégories dérivées de schémas, section
\ref{perfect-section-pseudo-coherent}.
Un outil fondamental consistera à nous ramener au cas de l'espace projectif
au moyen d'une version dérivée du lemme de Chow ; voir le lemme \ref{more-morphisms-lemma-derived-chow}.

\begin{lemma}
\label{more-morphisms-lemma-case-of-tor-independence}
Considérons un diagramme commutatif de schémas
$$
\xymatrix{
Z' \ar[d] \ar[r] & Y' \ar[d] \\
X' \ar[r] & S'
}
$$
Soit $S \to S'$ un morphisme. Notons $X$ et $Y$ les changements de
base de $X'$ et $Y'$ à $S$.
Supposons que $Y' \to S'$ et $Z' \to X'$ soient plats.
Alors $X \times_S Y$ et $Z'$ sont Tor-indépendants sur $X' \times_{S'} Y'$.
\end{lemma}

\begin{proof}
La question est locale ; on peut donc supposer tous les schémas affines
(certains détails sont omis). Remarquons que
$$
\xymatrix{
X \times_S Y \ar[r] \ar[d] & X' \times_{S'} Y' \ar[d] \\
X \ar[r] & X'
}
$$
est cartésien, avec des flèches verticales plates.
Écrivons $X = \Spec(A)$, $X' = \Spec(A')$ et
$X' \times_{S'} Y' = \Spec(B')$. Alors
$X \times_S Y = \Spec(A \otimes_{A'} B')$.
Écrivons $Z' = \Spec(C')$. Il faut montrer que
$$
\text{Tor}_p^{B'}(A \otimes_{A'} B', C') = 0,
\quad\text{pour } p > 0
$$
Puisque $A' \to B'$ est plat,
on a $A \otimes_{A'} B' = A \otimes_{A'}^\mathbf{L} B'$.
Par conséquent,
$$
(A \otimes_{A'} B') \otimes_{B'}^\mathbf{L} C' =
(A \otimes_{A'}^\mathbf{L} B') \otimes_{B'}^\mathbf{L} C' =
A \otimes_{A'}^\mathbf{L} C' =
A \otimes_{A'} C'
$$
La deuxième égalité résulte de Compléments d'algèbre, lemme
\ref{more-algebra-lemma-double-base-change}.
La dernière égalité vient de ce que $A' \to C'$ est plat. Ceci démontre le lemme.
\end{proof}

\begin{lemma}[Lemme de Chow dérivé]
\label{more-morphisms-lemma-derived-chow}
Soit $A$ un anneau. Soit $X$ un schéma séparé de présentation finie
sur $A$. Soit $x \in X$. Alors il existe
un voisinage ouvert $U \subset X$ de $x$,
un entier $n \geq 0$,
un ouvert $V \subset \mathbf{P}^n_A$,
un sous-schéma fermé $Z \subset X \times_A \mathbf{P}^n_A$,
un point $z \in Z$, et
un objet $E$ de $D(\mathcal{O}_{X \times_A \mathbf{P}^n_A})$ tels que
\begin{enumerate}
\item $Z \to X \times_A \mathbf{P}^n_A$ est de présentation finie ;
\item $b : Z \to X$ est un isomorphisme au-dessus de $U$ et $b(z) = x$ ;
\item $c : Z \to \mathbf{P}^n_A$ est une immersion fermée au-dessus de $V$ ;
\item $b^{-1}(U) = c^{-1}(V)$, en particulier $c(z) \in V$ ;
\item $E|_{X \times_A V} \cong
(b, c)_*\mathcal{O}_Z|_{X \times_A V}$ ;
\item $E$ est pseudo-cohérent et à support dans $Z$.
\end{enumerate}
\end{lemma}

\begin{proof}
On peut trouver une sous-$\mathbf{Z}$-algèbre de type fini $A' \subset A$
et un schéma $X'$ séparé et de présentation finie sur $A'$
dont le changement de base à $A$ est $X$. Voir
Limites, lemmes \ref{limits-lemma-descend-finite-presentation} et
\ref{limits-lemma-descend-separated-finite-presentation}.
Soit $x' \in X'$ l'image de $x$.
Si l'on peut démontrer le lemme pour $x' \in X'/A'$, alors
le lemme en résulte pour $x \in X/A$.
En effet, si $U', n', V', Z', z', E'$ fournissent la solution
pour $x' \in X'/A'$, alors on peut prendre pour
$U \subset X$ l'image réciproque de $U'$,
prendre $n = n'$,
prendre pour $V \subset \mathbf{P}^n_A$ l'image réciproque de $V'$,
prendre pour $Z \subset X \times \mathbf{P}^n$
l'image réciproque schématique de $Z'$,
prendre pour $z \in Z$ l'unique point ayant pour image $x$, et
prendre pour $E$ l'image réciproque dérivée de $E'$.
Remarquons que $E$ est pseudo-cohérent d'après
Cohomologie, lemme \ref{cohomology-lemma-pseudo-coherent-pullback}.
Il reste seulement à vérifier (5). Pour cela,
posons $W = b^{-1}(U) = c^{-1}(V)$ et $W' = (b')^{-1}(U) = (c')^{-1}(V')$ 
et considérons le carré cartésien
$$
\xymatrix{
W \ar[d]_{(b, c)} \ar[r] & W' \ar[d]^{(b', c')} \\
X \times_A V \ar[r] & X' \times_{A'} V'
}
$$
D'après le lemme \ref{more-morphisms-lemma-case-of-tor-independence}, les schémas
$X \times_A V$ et $W'$ sont Tor-indépendants sur $X' \times_{A'} V'$.
Par conséquent, l'image réciproque dérivée de
$(b', c')_*\mathcal{O}_{W'}$ sur $X \times_A V$
est $(b, c)_*\mathcal{O}_W$ d'après
Catégories dérivées de schémas,
lemme \ref{perfect-lemma-compare-base-change}.
On utilise aussi le fait que $R(b', c')_*\mathcal{O}_{Z'} = (b', c')_*\mathcal{O}_{Z'}$
car $(b', c')$ est une immersion fermée, et de même pour
$(b, c)_*\mathcal{O}_Z$.
Puisque $E'|_{U' \times_{A'} V'} =
(b', c')_*\mathcal{O}_{W'}$, on obtient
$E|_{U \times_A V} = (b, c)_*\mathcal{O}_W$,
et (5) est vérifiée.
Ceci nous ramène à la situation décrite au
paragraphe suivant.

\medskip\noindent
Supposons que $A$ soit de type fini sur $\mathbf{Z}$.
Choisissons un voisinage ouvert affine $U \subset X$ de $x$.
Alors $U$ est de type fini sur $A$.
Choisissons une immersion fermée $U \to \mathbf{A}^n_A$ et notons
$j : U \to \mathbf{P}^n_A$ l'immersion obtenue en la composant
avec l'immersion ouverte $\mathbf{A}^n_A \to \mathbf{P}^n_A$.
Soit $Z$ l'adhérence schématique de
$$
(\text{id}_U, j) : U \longrightarrow X \times_A \mathbf{P}^n_A
$$
Puisque la projection $X \times \mathbf{P}^n \to X$ est séparée,
on conclut de Morphismes, lemme
\ref{morphisms-lemma-scheme-theoretic-image-of-partial-section},
que $b : Z \to X$ est un isomorphisme au-dessus de $U$.
Soit $z \in Z$ l'unique point situé au-dessus de $x$.

\medskip\noindent
Soit $Y \subset \mathbf{P}^n_A$ l'adhérence schématique
de l'image de $j$. Il est alors clair que $Z \subset X \times_A Y$
est l'adhérence schématique de
$(\text{id}_U, j) : U \to X \times_A Y$.
Comme $X$ est séparé, le morphisme
$X \times_A Y \to Y$ est également séparé.
On voit donc que $Z \to Y$ est un isomorphisme au-dessus du
sous-schéma ouvert $j(U) \subset Y$, d'après le même lemme que ci-dessus.
Choisissons un ouvert $V \subset \mathbf{P}^n_A$ tel que $V \cap Y = j(U)$.
On voit alors que (3) et (4) sont satisfaites.

\medskip\noindent
Puisque $A$ est noethérien, on voit que $X$ et $X \times_A \mathbf{P}^n_A$
sont des schémas noethériens. On peut donc prendre $E = (b, c)_*\mathcal{O}_Z$
dans ce cas ; voir Catégories dérivées de schémas, lemme
\ref{perfect-lemma-identify-pseudo-coherent-noetherian}.
Ceci achève la démonstration.
\end{proof}

\begin{lemma}
\label{more-morphisms-lemma-compute-Fourier-Mukai-for-derived-chow}
Soient $A$, $x \in X$, ainsi que
$U, n, V, Z, z, E$ comme dans le lemme \ref{more-morphisms-lemma-derived-chow}.
Pour tout $K \in D_\QCoh(\mathcal{O}_X)$, on a
$$
Rq_*(Lp^*K \otimes^\mathbf{L} E)|_V = R(U \to V)_*K|_U
$$
où $p : X \times_A \mathbf{P}^n_A \to X$ et
$q : X \times_A \mathbf{P}^n_A \to \mathbf{P}^n_A$ sont
les projections et où le morphisme $U \to V$ est
l'immersion fermée de présentation finie $c \circ (b|_U)^{-1}$.
\end{lemma}

\begin{proof}
Puisque $b^{-1}(U) = c^{-1}(V)$ et que $c$ est une immersion fermée
au-dessus de $V$, on voit que $c \circ (b|_U)^{-1}$ est une immersion fermée.
Elle est de présentation finie parce que $U$ et $V$ sont de présentation
finie sur $A$ ; voir
Morphismes, lemme \ref{morphisms-lemma-finite-presentation-permanence}.
On a d'abord
$$
Rq_*(Lp^*K \otimes^\mathbf{L} E)|_V =
Rq'_*\left((Lp^*K \otimes^\mathbf{L} E)|_{X \times_A V}\right)
$$
où $q' : X \times_A V \to V$ est la projection, car
la formation de l'image directe totale commute à la localisation.
Posons $W = b^{-1}(U) = c^{-1}(V)$ et notons $i : W \to X \times_A V$
l'immersion fermée $i = (b, c)|_W$. Alors
$$
Rq'_*\left((Lp^*K \otimes^\mathbf{L} E)|_{X \times_A V}\right) =
Rq'_*(Lp^*K|_{X \times_A V} \otimes^\mathbf{L} i_*\mathcal{O}_W)
$$
d'après la propriété (5). Puisque $i$ est une immersion fermée, on a
$i_*\mathcal{O}_W = Ri_*\mathcal{O}_W$.
En utilisant
Catégories dérivées de schémas,
lemme \ref{perfect-lemma-cohomology-base-change},
on peut réécrire ceci sous la forme
$$
Rq'_* Ri_* Li^* Lp^*K|_{X \times_A V} =
R(q' \circ i)_* Lb^*K|_W =
R(U \to V)_* K|_U
$$
ce qui est le résultat voulu.
\end{proof}

\begin{lemma}
\label{more-morphisms-lemma-characterize-pseudo-coherent}
Soit $A$ un anneau. Soit $X$ un schéma séparé et
de présentation finie sur $A$. Soit $K \in D_\QCoh(\mathcal{O}_X)$.
Si $R\Gamma(X, E \otimes^\mathbf{L} K)$ est pseudo-cohérent
dans $D(A)$ pour tout objet pseudo-cohérent $E$ de $D(\mathcal{O}_X)$,
alors $K$ est pseudo-cohérent relativement à $A$.
\end{lemma}

\begin{proof}
Supposons que $K \in D_\QCoh(\mathcal{O}_X)$ et que
$R\Gamma(X, E \otimes^\mathbf{L} K)$ soit pseudo-cohérent
dans $D(A)$ pour tout objet pseudo-cohérent $E$ de $D(\mathcal{O}_X)$.
Soit $x \in X$. Nous allons montrer que $K$ est pseudo-cohérent relativement à $A$
dans un voisinage de $x$, ce qui démontrera le lemme.

\medskip\noindent
Choisissons $U, n, V, Z, z, E$ comme dans le lemme \ref{more-morphisms-lemma-derived-chow}.
Notons $p : X \times \mathbf{P}^n \to X$ et
$q : X \times \mathbf{P}^n \to \mathbf{P}^n_A$
les projections.
Alors, pour tout $i \in \mathbf{Z}$, on a
\begin{align*}
& R\Gamma(\mathbf{P}^n_A,
Rq_*(Lp^*K \otimes^\mathbf{L} E)
\otimes^\mathbf{L}
\mathcal{O}_{\mathbf{P}^n_A}(i)) \\
& =
R\Gamma(X \times \mathbf{P}^n,
Lp^*K \otimes^\mathbf{L} E \otimes^\mathbf{L}
Lq^*\mathcal{O}_{\mathbf{P}^n_A}(i)) \\
& =
R\Gamma(X,
K \otimes^\mathbf{L} Rp_*(E \otimes^\mathbf{L}
Lq^*\mathcal{O}_{\mathbf{P}^n_A}(i)))
\end{align*}
d'après
Catégories dérivées de schémas,
lemme \ref{perfect-lemma-cohomology-base-change}.
D'après
Catégories dérivées de schémas,
lemme \ref{perfect-lemma-flat-proper-pseudo-coherent-direct-image-general},
le complexe $Rp_*(E \otimes^\mathbf{L} Lq^*\mathcal{O}_{\mathbf{P}^n_A}(i))$
est pseudo-cohérent sur $X$. L'hypothèse nous dit donc que l'expression
de la formule affichée est un objet pseudo-cohérent de $D(A)$.
D'après
Catégories dérivées de schémas,
lemme \ref{perfect-lemma-pseudo-coherent-on-projective-space},
on conclut que $Rq_*(Lp^*K \otimes^\mathbf{L} E)$ est
pseudo-cohérent sur $\mathbf{P}^n_A$.
D'après le lemme \ref{more-morphisms-lemma-compute-Fourier-Mukai-for-derived-chow},
on a
$$
Rq_*(Lp^*K \otimes^\mathbf{L} E)|_{X \times_A V} =
R(U \to V)_*K|_U
$$
Comme $U \to V$ est une immersion fermée dans un sous-schéma ouvert de
$\mathbf{P}^n_A$, cela signifie que $K|_U$ est pseudo-cohérent relativement à $A$
d'après le lemme \ref{more-morphisms-lemma-check-relative-pseudo-coherence-on-charts}.
\end{proof}

\begin{lemma}
\label{more-morphisms-lemma-characterize-pseudo-coh-improved}
Soit $A$ un anneau. Soit $X$ un schéma séparé et
de présentation finie sur $A$. Soit
$K \in D_\QCoh(\mathcal{O}_X).$ Si
$R \Gamma (X, E \otimes ^{\mathbf{L}} K)$ est
pseudo-cohérent dans $D(A)$ pour tout objet parfait
$E \in D(\mathcal{O}_X)$, alors $K$ est pseudo-cohérent
relativement à $A$.
\end{lemma}

\begin{proof}
Compte tenu du lemme \ref{more-morphisms-lemma-characterize-pseudo-coherent}, il suffit
de montrer que $R \Gamma (X, E \otimes ^{\mathbf{L}} K)$ est
pseudo-cohérent dans $D(A)$ pour tout objet pseudo-cohérent
$E \in D(\mathcal{O}_X)$. D'après Catégories dérivées de schémas,
proposition \ref{perfect-proposition-detecting-bounded-above},
il s'ensuit que $K \in D^-_\QCoh (\mathcal{O}_X)$. Le
résultat résulte alors de Catégories dérivées de schémas, lemme
\ref{perfect-lemma-perfect-enough}.
\end{proof}

\begin{lemma}
\label{more-morphisms-lemma-characterize-relatively-perfect}
Soit $A$ un anneau. Soit $X$ un schéma séparé, de
présentation finie et plat sur $A$. Soit
$K \in D_\QCoh(\mathcal{O}_X).$ Si
$R \Gamma (X, E \otimes^\mathbf{L} K)$ est parfait dans
$D(A)$ pour tout objet parfait $E \in D(\mathcal{O}_X)$, alors $K$ est
parfait relativement à $\Spec(A)$.
\end{lemma}

\begin{proof}
D'après le lemme \ref{more-morphisms-lemma-characterize-pseudo-coh-improved},
$K$ est pseudo-cohérent relativement à $A$. D'après le lemme
\ref{more-morphisms-lemma-check-relative-pseudo-coherence-on-charts},
$K$ est pseudo-cohérent dans $D( \mathcal{O}_X)$.
D'après Catégories dérivées de schémas, proposition
\ref{perfect-proposition-detecting-bounded-below},
on voit que $K$ appartient à $D^-(\mathcal{O}_X)$.
Soit $\mathfrak{p}$ un idéal premier de $A$ et notons
$i : Y \to X$ l'inclusion de la fibre schématique
au-dessus de $\mathfrak{p}$, c'est-à-dire que $Y$ est un schéma sur $\kappa(\mathfrak p)$.
D'après Catégories dérivées de schémas, lemme \ref{perfect-lemma-bounded-on-fibres},
il suffira de montrer que $Li^*(K)$ est borné inférieurement.
Soit $G \in D_{perf} (\mathcal{O}_X)$ un complexe parfait
qui engendre $D_\QCoh (\mathcal{O}_X)$ ;
voir Catégories dérivées de schémas, théorème
\ref{perfect-theorem-bondal-van-den-Bergh}.
On a
\begin{align*}
R\Hom _{\mathcal{O}_Y}(Li^*(G), Li^*(K))
& =
R\Gamma(Y, Li^*(G ^\vee \otimes ^\mathbf{L} K)) \\
& =
R\Gamma(X, G^\vee \otimes ^{\mathbf{L}} K)
\otimes^\mathbf{L}_A \kappa(\mathfrak{p})
\end{align*}
La première égalité utilise le fait que $Li^*$ préserve les objets parfaits et les duaux,
ainsi que Cohomologie, lemme \ref{cohomology-lemma-dual-perfect-complex} ; nous omettons
quelques détails. La seconde égalité résulte de
Catégories dérivées de schémas, lemme
\ref{perfect-lemma-compare-base-change},
car $X$ est plat sur $A$. Il résulte de notre hypothèse que ceci est un
objet parfait de $D(\kappa(\mathfrak{p}))$. L'objet
$Li^*(G) \in D_{perf}(\mathcal{O}_Y)$ engendre $D_\QCoh(\mathcal{O}_Y)$ d'après
Catégories dérivées de schémas, remarque \ref{perfect-remark-pullback-generator}.
Par conséquent, Catégories dérivées de schémas, proposition
\ref{perfect-proposition-detecting-bounded-below},
implique maintenant que $Li^*(K)$ est borné inférieurement, ce qui conclut.
\end{proof}




\section{Descente des propriétés de finitude des complexes}
\label{more-morphisms-section-descent-finiteness}

\noindent
Cette section fait suite à Catégories dérivées de schémas,
section \ref{perfect-section-descent-finiteness}.

\begin{lemma}
\label{more-morphisms-lemma-relative-pseudo-coherent-descends-fppf}
Soit $X \to S$ localement de type fini.
Soit $\{f_i : X_i \to X\}$ un recouvrement fppf de schémas.
Soient $E \in D_\QCoh(\mathcal{O}_X)$ et $m \in \mathbf{Z}$.
Alors $E$ est $m$-pseudo-cohérent relativement à $S$
si et seulement si chaque $Lf_i^*E$ est $m$-pseudo-cohérent relativement à $S$.
\end{lemma}

\begin{proof}
Supposons que $E$ soit $m$-pseudo-cohérent relativement à $S$.
Les morphismes $f_i$ sont pseudo-cohérents d'après le
lemme \ref{more-morphisms-lemma-flat-finite-presentation-pseudo-coherent}.
Ainsi, $Lf_i^*E$ est $m$-pseudo-cohérent relativement à $S$
d'après le lemme \ref{more-morphisms-lemma-pull-relative-pseudo-coherent}.

\medskip\noindent
Réciproquement, supposons que $Lf_i^*E$ soit $m$-pseudo-cohérent relativement à $S$
pour tout $i$. Choisissons $S = \bigcup U_j$, $W_j \to U_j$,
$W_j = \bigcup W_{j, k}$, $T_{j, k} \to W_{j, k}$, et des
morphismes $\alpha_{j, k} : T_{j, k} \to X_{i(j, k)}$ au-dessus de $S$ comme dans le
lemme \ref{more-morphisms-lemma-dominate-fppf}.
Comme le morphisme $T_{j, K} \to S$ est plat et de présentation finie,
on voit que $\alpha_{j, k}$ est pseudo-cohérent d'après le
lemme \ref{more-morphisms-lemma-permanence-pseudo-coherent}.
Ainsi,
$$
L\alpha_{j, k}^*Lf_{i(j, k)}^*E = L(T_{i, k} \to S)^*E
$$
est $m$-pseudo-cohérent relativement à $S$ d'après le
lemme \ref{more-morphisms-lemma-pull-relative-pseudo-coherent}.
Nous voulons maintenant faire descendre cette propriété
le long des recouvrements $\{T_{j, k} \to W_{j, k}\}$,
$W_j = \bigcup W_{j, k}$, $\{W_j \to U_j\}$ et $S = \bigcup U_j$.
Comme le résultat est vrai pour les recouvrements de Zariski
(par définition de la $m$-pseudo-cohérence relativement à $S$),
cela signifie que nous pouvons supposer qu'il existe un morphisme
fini localement libre et surjectif $\pi : Y \to X$
tel que $L\pi^*E$ soit pseudo-cohérent relativement à $S$.
Dans ce cas, $R\pi_*L\pi^*E$ est pseudo-cohérent relativement à $S$
d'après le lemme \ref{more-morphisms-lemma-finite-morphism-relative-pseudo-coherence}
(c'est la première fois que nous utilisons le fait que les faisceaux de cohomologie de $E$ sont quasi-cohérents).
On a
$R\pi_*L\pi^*E = E \otimes^\mathbf{L}_{\mathcal{O}_X} \pi_*\mathcal{O}_Y$,
par exemple d'après
Catégories dérivées de schémas, lemme
\ref{perfect-lemma-cohomology-base-change},
et, localement sur $X$, l'application $\mathcal{O}_X \to \pi_*\mathcal{O}_Y$
est l'inclusion d'un facteur direct. On conclut donc d'après le
lemme \ref{more-morphisms-lemma-summands-relative-pseudo-coherent}.
\end{proof}

\begin{lemma}
\label{more-morphisms-lemma-relative-pseudo-coherent-post-compose}
Soient $X \to T \to S$ des morphismes de schémas. Supposons que $T \to S$
soit plat et localement de présentation finie et que $X \to T$ soit
localement de type fini. Soient $E \in D(\mathcal{O}_X)$ et $m \in \mathbf{Z}$.
Alors $E$ est $m$-pseudo-cohérent relativement à $S$
si et seulement si $E$ est $m$-pseudo-cohérent relativement à $T$.
\end{lemma}

\begin{proof}
Localement sur $X$, on peut choisir une immersion fermée $i : X \to \mathbf{A}^n_T$.
Alors $\mathbf{A}^n_T \to S$ est plat et localement de présentation finie.
Nous pouvons donc
appliquer le lemme \ref{more-morphisms-lemma-composition-relative-pseudo-coherent}
pour obtenir l'équivalence.
\end{proof}

\begin{lemma}
\label{more-morphisms-lemma-relative-pseudo-coherent-descends-fppf-base}
Soit $f : X \to S$ localement de type fini.
Soit $\{S_i \to S\}$ un recouvrement fppf de schémas.
Notons $f_i : X_i \to S_i$ le changement de base de $f$
et $g_i : X_i \to X$ la projection.
Soient $E \in D_\QCoh(\mathcal{O}_X)$ et $m \in \mathbf{Z}$.
Alors $E$ est $m$-pseudo-cohérent relativement à $S$
si et seulement si chaque $Lg_i^*E$ est $m$-pseudo-cohérent relativement à $S_i$.
\end{lemma}

\begin{proof}
Ceci résulte formellement des
lemmes \ref{more-morphisms-lemma-relative-pseudo-coherent-descends-fppf} et
\ref{more-morphisms-lemma-relative-pseudo-coherent-post-compose}.
En effet, si $E$ est $m$-pseudo-cohérent relativement à $S$,
alors $Lg_i^*E$ est $m$-pseudo-cohérent relativement à $S$ (d'après le premier lemme),
donc $Lg_i^*E$ est $m$-pseudo-cohérent relativement à $S_i$ (d'après le second).
Réciproquement, si
$Lg_i^*E$ est $m$-pseudo-cohérent relativement à $S_i$, alors
$Lg_i^*E$ est $m$-pseudo-cohérent relativement à $S$ (d'après le second lemme),
donc $E$ est $m$-pseudo-cohérent relativement à $S$ (d'après le premier lemme).
\end{proof}






\section{Objets relativement parfaits}
\label{more-morphisms-section-relatively-perfect}

\noindent
Cette section fait suite à la discussion de
Catégories dérivées de schémas, section
\ref{perfect-section-relatively-perfect}.

\begin{lemma}
\label{more-morphisms-lemma-thickening-pseudo-coherent}
Soit $i : X \to X'$ un épaississement d'ordre fini de schémas. Soit
$K' \in D(\mathcal{O}_{X'})$ un objet tel que
$K = Li^*K'$ soit pseudo-cohérent. Alors $K'$ est pseudo-cohérent.
\end{lemma}

\begin{proof}
Montrons d'abord que les faisceaux de cohomologie de $K'$ sont quasi-cohérents.
Pour ce faire, on peut se ramener au cas d'un épaississement du premier ordre, voir
section \ref{more-morphisms-section-thickenings}. Soit $\mathcal{I} \subset \mathcal{O}_{X'}$
le faisceau quasi-cohérent d'idéaux définissant $X$.
En tensorisant la suite exacte courte
$$
0 \to \mathcal{I} \to \mathcal{O}_{X'} \to i_*\mathcal{O}_X \to 0
$$
par $K'$, on obtient un triangle distingué
$$
K' \otimes_{\mathcal{O}_{X'}}^\mathbf{L} \mathcal{I}
\to K' \to
K' \otimes_{\mathcal{O}_{X'}}^\mathbf{L} i_*\mathcal{O}_X
\to
(K' \otimes_{\mathcal{O}_{X'}}^\mathbf{L} \mathcal{I})[1]
$$
Comme $i_* = Ri_*$ et que l'on peut considérer $\mathcal{I}$
comme un $\mathcal{O}_X$-module quasi-cohérent (puisqu'il s'agit d'un épaississement
du premier ordre), on peut réécrire ceci sous la forme
$$
i_*(K \otimes_{\mathcal{O}_X}^\mathbf{L} \mathcal{I})
\to K' \to
i_*K \to
i_*(K \otimes_{\mathcal{O}_X}^\mathbf{L} \mathcal{I})[1]
$$
On utilise Cohomologie, lemme
\ref{cohomology-lemma-projection-formula-closed-immersion},
pour identifier les termes. Comme $K$ appartient à
$D_\QCoh(\mathcal{O}_X)$, on conclut que
$K'$ appartient à $D_\QCoh(\mathcal{O}_{X'})$ ; ceci utilise
Catégories dérivées de schémas, lemmes
\ref{perfect-lemma-pseudo-coherent},
\ref{perfect-lemma-quasi-coherence-tensor-product} et
\ref{perfect-lemma-quasi-coherence-direct-image}.

\medskip\noindent
Supposons que $K'$ appartienne à $D_\QCoh(\mathcal{O}_{X'})$.
La question est locale sur $X'$ ; on peut donc supposer que $X'$ est affine.
Écrivons $X' = \Spec(A')$ et $X = \Spec(A)$ avec $A = A'/I$ et $I$ nilpotent.
Alors $K'$ provient d'un objet $M' \in D(A')$, voir
Catégories dérivées de schémas, lemme
\ref{perfect-lemma-affine-compare-bounded}.
Ainsi, $M = M' \otimes_{A'}^\mathbf{L} A$ est un objet pseudo-cohérent
de $D(A)$ d'après Catégories dérivées de schémas,
lemme \ref{perfect-lemma-pseudo-coherent-affine}, et notre hypothèse sur $K$.
Il s'ensuit que $M'$ est pseudo-cohérent d'après
Compléments d'algèbre, lemme
\ref{more-algebra-lemma-check-pseudo-coherent-modulo-nilpotent}.
Ceci achève la démonstration.
\end{proof}

\begin{lemma}
\label{more-morphisms-lemma-thickening-relatively-perfect}
Considérons un diagramme cartésien
$$
\xymatrix{
X \ar[r]_i \ar[d]_f & X' \ar[d]^{f'} \\
Y \ar[r]^j & Y'
}
$$
de schémas. Supposons que $X' \to Y'$ soit plat et localement
de présentation finie et que $Y \to Y'$ soit un épaississement d'ordre fini.
Soit $E' \in D(\mathcal{O}_{X'})$. Si $E = Li^*(E')$ est parfait relativement à $Y$,
alors $E'$ est parfait relativement à $Y'$.
\end{lemma}

\begin{proof}
Rappelons que la perfection relative à $Y$ de $E$ signifie que $E$ est
pseudo-cohérent et possède localement une dimension de Tor finie comme complexe
de $f^{-1}\mathcal{O}_Y$-modules
(Catégories dérivées de schémas,
définition \ref{perfect-definition-relatively-perfect}).
D'après le lemme \ref{more-morphisms-lemma-thickening-pseudo-coherent},
on trouve que $E'$ est pseudo-cohérent.
En particulier, $E'$ appartient à $D_\QCoh(\mathcal{O}_{X'})$, voir
Catégories dérivées de schémas, lemme \ref{perfect-lemma-pseudo-coherent}.
Pour montrer que $E'$ possède localement une dimension de Tor finie,
on peut travailler localement sur $X'$. On peut donc supposer que
$X'$, $S'$, $X$, $S$ sont affines, disons associés aux
anneaux $A'$, $R'$, $A$, $R$.
On se ramène alors à la version d'algèbre commutative grâce à
Catégories dérivées de schémas,
lemme \ref{perfect-lemma-affine-locally-rel-perfect}.
La version d'algèbre commutative est donnée dans Compléments d'algèbre, lemme
\ref{more-algebra-lemma-thickening-relatively-perfect}.
\end{proof}

\begin{lemma}
\label{more-morphisms-lemma-henselian-relatively-perfect}
Soit $(R, I)$ une paire formée d'un anneau et d'un idéal $I$
contenu dans le radical de Jacobson. Posons $S = \Spec(R)$ et $S_0 = \Spec(R/I)$.
Soit $f : X \to S$ propre, plat et de présentation finie.
Notons $X_0 = S_0 \times_S X$. Soit $E \in D(\mathcal{O}_X)$
pseudo-cohérent. Si la restriction dérivée $E_0$ de $E$
à $X_0$ est parfaite relativement à $S_0$, alors $E$ est parfait relativement à $S$.
\end{lemma}

\begin{proof}
Choisissons un recouvrement fini par des ouverts affines $X = U_1 \cup \ldots \cup U_n$.
Pour tout $i$, on peut choisir une immersion fermée $U_i \to \mathbf{A}^{d_i}_S$.
Posons $U_{i, 0} = S_0 \times_S U_i$.
Pour tout $i$, le complexe $E_0|_{U_{i, 0}}$ a une amplitude de Tor
contenue dans $[a_i, b_i]$ pour certains $a_i, b_i \in \mathbf{Z}$.
Soit $x \in X$ un point.
Nous allons montrer que l'amplitude de Tor de
$E_x$ sur $R$ est contenue dans $[a_i - d_i, b_i]$ pour un certain $i$.
Ceci achèvera la démonstration, car l'amplitude de Tor peut être
lue sur les germes d'après
Cohomologie, lemme \ref{cohomology-lemma-tor-amplitude-stalk}.

\medskip\noindent
Comme $f$ est propre, $f(\overline{\{x\}})$ est un sous-ensemble fermé de $S$.
Comme $I$ est contenu dans le radical de Jacobson, on voit que
$f(\overline{\{x\}})$ rencontre le sous-ensemble fermé $S_0 \subset S$.
Il existe donc une spécialisation $x \leadsto x_0$ avec $x_0 \in X_0$.
Choisissons $i$ tel que $x_0 \in U_i$ ; ainsi $x_0 \in U_{i, 0}$.
Nous fixerons $i$ pour le reste de la démonstration.
Écrivons $U_i = \Spec(A)$. Alors $A$ est une algèbre plate de présentation
finie sur $R$, qui est un quotient d'une $R$-algèbre polynomiale en
$d_i$ variables. La restriction $E|_{U_i}$ correspond
(d'après Catégories dérivées de schémas, lemmes
\ref{perfect-lemma-affine-compare-bounded} et
\ref{perfect-lemma-pseudo-coherent-affine})
à un objet pseudo-cohérent $K$ de $D(A)$.
Remarquons que $E_0$ correspond à $K \otimes_A^\mathbf{L} A/IA$.
Soient $\mathfrak q \subset \mathfrak q_0 \subset A$ les idéaux
premiers correspondant à $x \leadsto x_0$.
Alors $E_x = K_{\mathfrak q}$ et $K_{\mathfrak q}$
est une localisation de $K_{\mathfrak q_0}$. Il suffit donc
de montrer que $K_{\mathfrak q_0}$ a une amplitude de Tor contenue dans
$[a_i - d_i, b_i]$ comme complexe de $R$-modules.
Soit $I \subset \mathfrak p_0 \subset R$ l'idéal
premier correspondant à $f(x_0)$.
Alors on a
\begin{align*}
K \otimes_R^\mathbf{L} \kappa(\mathfrak p_0)
& =
(K \otimes_R^\mathbf{L} R/I) \otimes_{R/I}^\mathbf{L}
\kappa(\mathfrak p_0) \\
& =
(K \otimes_A^\mathbf{L} A/IA) \otimes_{R/I}^\mathbf{L} \kappa(\mathfrak p_0)
\end{align*}
la seconde égalité valant parce que $R \to A$ est plat.
Par notre choix de $a_i, b_i$, ce complexe n'a de cohomologie
qu'aux degrés de l'intervalle $[a_i, b_i]$.
On peut donc finalement appliquer
Compléments d'algèbre, lemme
\ref{more-algebra-lemma-lift-from-fibre-relatively-perfect},
à $R \to A$, $\mathfrak q_0$, $\mathfrak p_0$ et $K$
pour conclure.
\end{proof}










\section{Contraction de courbes rationnelles}
\label{more-morphisms-section-contracting}

\noindent
Dans cette section, nous étudions les morphismes propres $f : X \to Y$ dont les fibres
sont de dimension $\leq 1$ et qui vérifient $R^1f_*\mathcal{O}_X = 0$.
Pour comprendre le titre de cette section, voir
Courbes algébriques, sections
\ref{curves-section-contracting-rational-tails},
\ref{curves-section-contracting-rational-bridges} et
\ref{curves-section-contracting-to-stable}.

\begin{lemma}
\label{more-morphisms-lemma-check-h1-fibre-zero}
Soit $f : X \to Y$ un morphisme propre de schémas. Soit $y \in Y$
un point tel que $\dim(X_y) \leq 1$. Si
\begin{enumerate}
\item $R^1f_*\mathcal{O}_X = 0$, ou, plus généralement,
\item il existe un morphisme $g : Y' \to Y$ tel que $y$ appartienne à l'image
de $g$ et tel que $R^1f'_*\mathcal{O}_{X'} = 0$, où $f' : X' \to Y'$
est le changement de base de $f$ par $g$,
\end{enumerate}
alors $H^1(X_y, \mathcal{O}_{X_y}) = 0$.
\end{lemma}

\begin{proof}
Pour démontrer le lemme, on peut remplacer $Y$ par un voisinage ouvert de $y$.
On peut donc supposer que $Y$ est affine
et que toutes les fibres de $f$ sont de dimension $\leq 1$, voir
Morphismes, lemme \ref{morphisms-lemma-openness-bounded-dimension-fibres}.
Dans ce cas, $R^1f_*\mathcal{O}_X$ est un $\mathcal{O}_Y$-module quasi-cohérent
de type fini et sa formation commute à tout changement de base, voir
Limites, lemmes \ref{limits-lemma-proper-top-cohomology-finite-type} et
\ref{limits-lemma-higher-direct-images-zero-above-dimension-fibre}.
Le lemme en résulte immédiatement.
\end{proof}

\begin{lemma}
\label{more-morphisms-lemma-h1-fibre-zero}
Soit $f : X \to Y$ un morphisme propre de schémas. Soit $y \in Y$
un point tel que $\dim(X_y) \leq 1$ et $H^1(X_y, \mathcal{O}_{X_y}) = 0$.
Alors il existe un voisinage ouvert $V \subset Y$ de $y$ tel que
$R^1f_*\mathcal{O}_X|_V = 0$
et que la même propriété soit vraie après tout changement de base $Y' \to V$.
\end{lemma}

\begin{proof}
Pour démontrer le lemme, on peut remplacer $Y$ par un voisinage ouvert de $y$.
On peut donc supposer que $Y$ est affine et que toutes les fibres de $f$ sont de
dimension $\leq 1$, voir
Morphismes, lemme \ref{morphisms-lemma-openness-bounded-dimension-fibres}.
Dans ce cas, $R^1f_*\mathcal{O}_X$ est un $\mathcal{O}_Y$-module quasi-cohérent
de type fini et sa formation commute à tout changement de base, voir
Limites, lemmes \ref{limits-lemma-proper-top-cohomology-finite-type} et
\ref{limits-lemma-higher-direct-images-zero-above-dimension-fibre}.
Écrivons $Y = \Spec(A)$ ; supposons que $y$ corresponde à l'idéal premier $\mathfrak p \subset A$ et que
$R^1f_*\mathcal{O}_X$ corresponde au $A$-module fini $M$.
Alors $H^1(X_y, \mathcal{O}_{X_y}) = 0$ signifie que
$\mathfrak pM_\mathfrak p = M_\mathfrak p$ d'après l'énoncé
sur le changement de base. Par le lemme de Nakayama,
on conclut que $M_\mathfrak p = 0$. Comme $M$ est fini, on trouve
un $f \in A$, $f \not \in \mathfrak p$, tel que $M_f = 0$.
En prenant pour $V$ l'ouvert principal $D(f)$, on obtient donc le résultat voulu.
\end{proof}

\begin{lemma}
\label{more-morphisms-lemma-globally-generated-vanishing}
Soit $f : X \to Y$ un morphisme propre de schémas tel
que $\dim(X_y) \leq 1$ et $H^1(X_y, \mathcal{O}_{X_y}) = 0$
pour tout $y \in Y$. Soit $\mathcal{F}$ quasi-cohérent sur $X$. Alors
\begin{enumerate}
\item $R^pf_*\mathcal{F} = 0$ pour $p > 1$, et
\item $R^1f_*\mathcal{F} = 0$ s'il existe une surjection
$f^*\mathcal{G} \to \mathcal{F}$ avec $\mathcal{G}$ quasi-cohérent
sur $Y$.
\end{enumerate}
Si $Y$ est affine, on a aussi
\begin{enumerate}
\item[(3)] $H^p(X, \mathcal{F}) = 0$ pour $p \not \in \{0, 1\}$, et
\item[(4)] $H^1(X, \mathcal{F}) = 0$ si $\mathcal{F}$ est engendré par ses sections globales.
\end{enumerate}
\end{lemma}

\begin{proof}
L'annulation de (1) est donnée par Limites, lemme
\ref{limits-lemma-higher-direct-images-zero-above-dimension-fibre}.
Pour démontrer (2), on peut travailler localement sur $Y$ et supposer que $Y$ est affine.
Alors $R^1f_*\mathcal{F}$ est le module quasi-cohérent sur $Y$
associé au module $H^1(X, \mathcal{F})$.
On utilise ici le fait que $Y$ est affine, la quasi-cohérence des images directes
supérieures (Cohomologie des schémas, lemme
\ref{coherent-lemma-quasi-coherence-higher-direct-images}) et
Cohomologie des schémas, lemme
\ref{coherent-lemma-quasi-coherence-higher-direct-images-application}.
Comme $Y$ est affine, le module quasi-cohérent $\mathcal{G}$
est engendré par ses sections globales ; il en est donc de même de $f^*\mathcal{G}$
et de $\mathcal{F}$.
On voit ainsi que (4) implique (2).
L'assertion (3) résulte de (1), ainsi que des remarques qui viennent d'être faites
sur la quasi-cohérence des images directes. Ainsi,
il ne reste qu'à démontrer (4).
Si $\mathcal{F}$ est engendré par ses sections globales, il existe une surjection
$\bigoplus_{i \in I} \mathcal{O}_X \to \mathcal{F}$. D'après l'assertion (1)
et la suite exacte longue de cohomologie, celle-ci
induit une surjection sur $H^1$. Comme $H^1(X, \mathcal{O}_X) = 0$
parce que $R^1f_*\mathcal{O}_X = 0$ d'après le
lemme \ref{more-morphisms-lemma-h1-fibre-zero}, et
comme $H^1(X, -)$ commute aux sommes directes
(Cohomologie, lemme \ref{cohomology-lemma-quasi-separated-cohomology-colimit}),
on conclut.
\end{proof}

\begin{lemma}
\label{more-morphisms-lemma-h1-fibre-zero-check-h0-kappa}
Soit $f : X \to Y$ un morphisme propre de schémas. Supposons que
\begin{enumerate}
\item pour tout $y \in Y$, on ait $\dim(X_y) \leq 1$ et
$H^1(X_y, \mathcal{O}_{X_y}) = 0$, et
\item $\mathcal{O}_Y \to f_*\mathcal{O}_X$ soit surjective.
\end{enumerate}
Alors $\mathcal{O}_{Y'} \to f'_*\mathcal{O}_{X'}$ est surjective
pour tout changement de base $f' : X' \to Y'$ de $f$.
\end{lemma}

\begin{proof}
On peut supposer que $Y$ et $Y'$ sont affines. On peut alors choisir une immersion
fermée $Y' \to Y''$ avec $Y'' \to Y$ un morphisme plat de schémas affines.
Par changement de base plat
(Cohomologie des schémas, lemme \ref{coherent-lemma-flat-base-change-cohomology}),
on voit que le résultat vaut pour $X'' \to Y''$.
On peut donc supposer que $Y'$ est un sous-schéma fermé de $Y$.
Soit $\mathcal{I} \subset \mathcal{O}_Y$ l'idéal définissant $Y'$.
Il existe alors une suite exacte courte
$$
0 \to \mathcal{I}\mathcal{O}_X \to
\mathcal{O}_X \to \mathcal{O}_{X'} \to 0
$$
où l'on considère $\mathcal{O}_{X'}$ comme un module quasi-cohérent sur $X$.
D'après le lemme \ref{more-morphisms-lemma-globally-generated-vanishing},
on a $H^1(X, \mathcal{I}\mathcal{O}_X) = 0$.
Il s'ensuit que
$$
H^0(Y, \mathcal{O}_Y) \to
H^0(Y, f_*\mathcal{O}_X) = H^0(X, \mathcal{O}_X) \to
H^0(X, \mathcal{O}_{X'})
$$
est surjective, comme voulu. La première flèche est surjective puisque $Y$
est affine et que nous avons supposé $\mathcal{O}_Y \to f_*\mathcal{O}_X$
surjective, et la seconde l'est par la suite exacte longue de
cohomologie associée à la suite exacte courte ci-dessus et par
l'annulation qui vient d'être démontrée.
\end{proof}

\begin{lemma}
\label{more-morphisms-lemma-h1-fibre-zero-h0-kappa}
Considérons un diagramme commutatif
$$
\xymatrix{
X \ar[rr]_f \ar[rd] & & Y \ar[ld] \\
& S
}
$$
de morphismes de schémas. Soit $s \in S$ un point. Supposons que
\begin{enumerate}
\item $X \to S$ soit localement de présentation finie et plat aux
points de $X_s$,
\item $f$ soit propre,
\item les fibres de $f_s : X_s \to Y_s$ soient de dimension $\leq 1$
et que $R^1f_{s, *}\mathcal{O}_{X_s} = 0$,
\item $\mathcal{O}_{Y_s} \to f_{s, *}\mathcal{O}_{X_s}$ soit surjective.
\end{enumerate}
Alors il existe un ouvert $Y_s \subset V \subset Y$ tel que
(a) $f^{-1}(V)$ soit plat sur $S$,
(b) $\dim(X_y) \leq 1$ pour $y \in V$,
(c) $R^1f_*\mathcal{O}_X|_V = 0$,
(d) $\mathcal{O}_V \to f_*\mathcal{O}_X|_V$
soit surjective,
et tel que (b), (c) et (d) restent vrais après tout changement de base $Y' \to V$.
\end{lemma}

\begin{proof}
Soit $y \in Y$ un point au-dessus de $s$.
Il suffit de trouver un voisinage ouvert de $y$
possédant les propriétés voulues. Dans un premier temps, on remplace $Y$
par l'ouvert $V$ trouvé dans le lemme \ref{more-morphisms-lemma-h1-fibre-zero}, de sorte que
$R^1f_*\mathcal{O}_X$ soit universellement nul (l'hypothèse du
lemme est satisfaite d'après le lemme \ref{more-morphisms-lemma-check-h1-fibre-zero}).
On rétrécit aussi $Y$ afin que toutes les fibres de $f$ soient de dimension $\leq 1$
(on utilise Morphismes, lemme
\ref{morphisms-lemma-openness-bounded-dimension-fibres},
et la propreté de $f$). On peut donc supposer que (b) et (c) sont satisfaites
avec $V = Y$ et après tout changement de base $Y' \to Y$.
Ainsi, d'après le lemme \ref{more-morphisms-lemma-h1-fibre-zero-check-h0-kappa},
il suffit maintenant de montrer (d) sur $Y$.
On peut encore rétrécir $Y$ ; par exemple, on peut supposer, et on suppose,
que $Y$ et $S$ sont affines.

\medskip\noindent
D'après le théorème \ref{more-morphisms-theorem-openness-flatness},
il existe un ouvert $U \subset X$ sur lequel $X \to S$
est plat et qui contient $X_s$ par hypothèse.
Alors $f(X \setminus U)$ est un fermé
qui ne contient pas $y$. Ainsi, après avoir rétréci $Y$,
on peut supposer que $X$ est plat sur $S$.

\medskip\noindent
Écrivons $S = \Spec(R)$. Choisissons une immersion fermée $Y \to Y'$, où $Y'$
est le spectre d'un anneau de polynômes $R[x_e; e \in E]$ sur un ensemble $E$.
Notons $f' : X \to Y'$ la composée de $f$ et de $Y \to Y'$.
Alors les hypothèses (1) -- (4), ainsi que (b) et (c),
sont satisfaites pour $f'$ et $s$. Si nous montrons que $\mathcal{O}_{Y'} \to f'_*\mathcal{O}_X$
est surjectif dans un voisinage ouvert de $y$, il en va de même pour
$\mathcal{O}_Y \to f_*\mathcal{O}_X$. Ainsi, on peut supposer que $Y$ est
le spectre de $R[x_e; e \in E]$.

\medskip\noindent
À ce stade, $X$ et $Y$ sont plats sur $S$. Alors $Y_s$ et $X$
sont Tor-indépendants au-dessus de $Y$. Nous invitons le lecteur à en trouver sa propre
démonstration, mais cela résulte aussi du
lemme \ref{more-morphisms-lemma-case-of-tor-independence} appliqué au
carré de sommets $X, Y, S, S$ et à son changement de base par $s \to S$.
Par conséquent,
$$
Rf_{s, *}\mathcal{O}_{X_s} = L(Y_s \to Y)^*Rf_*\mathcal{O}_X
$$
d'après le lemme
\ref{perfect-lemma-compare-base-change} de Catégories dérivées de schémas.
Compte tenu de l'annulation déjà établie, ceci implique
$f_{s, *}\mathcal{O}_{X_s} = (Y_s \to Y)^*f_*\mathcal{O}_X$.
On en conclut que $\mathcal{O}_Y \to f_*\mathcal{O}_X$ est un morphisme de
$\mathcal{O}_Y$-modules quasi-cohérents dont l'image inverse
sur $Y_s$ est surjective. Nous affirmons que
$f_*\mathcal{O}_X$ est un $\mathcal{O}_Y$-module de type fini.
Si cela est vrai, alors le conoyau $\mathcal{F}$ de
$\mathcal{O}_Y \to f_*\mathcal{O}_X$
est un $\mathcal{O}_Y$-module quasi-cohérent de type fini
tel que $\mathcal{F}_y \otimes \kappa(y) = 0$.
D'après le lemme de Nakayama (Algèbre, lemme \ref{algebra-lemma-NAK}),
on a $\mathcal{F}_y = 0$. Ainsi, $\mathcal{F}$ est nul
dans un voisinage ouvert de $y$
(Modules, lemme \ref{modules-lemma-finite-type-stalk-zero}),
ce qui achève la démonstration.

\medskip\noindent
Démonstration de l'affirmation.
Pour une partie finie $E' \subset E$, posons $Y' = \Spec(R[x_e; e \in E'])$.
Pour $E'$ assez grand, le morphisme $f' : X \to Y \to Y'$ est propre, voir
Limites, lemme \ref{limits-lemma-finite-type-eventually-proper}.
Dans la suite, on fixe $E'$ et $Y'$.
Écrivons $R = \colim R_i$ comme la colimite de ses sous-algèbres de type fini
sur $\mathbf{Z}$. Posons $S_i = \Spec(R_i)$
et $Y'_i = \Spec(R_i[x_e; e \in E'])$.
Pour $i$ assez grand, on peut trouver un diagramme
$$
\xymatrix{
X \ar[d] \ar[r]_{f'} & Y' \ar[d] \ar[r] & S \ar[d] \\
X_i \ar[r]^{f'_i} & Y'_i \ar[r] & S_i
}
$$
à carrés cartésiens, tel que $X_i$ soit plat sur $S_i$ et que
$X_i \to Y'_i$ soit propre. Voir
Limites, lemmes \ref{limits-lemma-descend-finite-presentation},
\ref{limits-lemma-descend-flat-finite-presentation} et
\ref{limits-lemma-eventually-proper}.
Le même argument que ci-dessus montre que $Y'$ et $X_i$
sont Tor-indépendants au-dessus de $Y'_i$ et donc que
$$
R\Gamma(X, \mathcal{O}_X) =
R\Gamma(X_i, \mathcal{O}_{X_i})
\otimes^\mathbf{L}_{R_i[x_e; e \in E']}
R[x_e; e \in E']
$$
d'après la même référence que ci-dessus. D'après Cohomologie des schémas, lemme
\ref{coherent-lemma-proper-over-affine-cohomology-finite},
le complexe $R\Gamma(X_i, \mathcal{O}_{X_i})$ est pseudo-cohérent
dans la catégorie dérivée de l'anneau noethérien $R_i[x_e; e \in E']$ (voir
Compléments d'algèbre, lemme \ref{more-algebra-lemma-Noetherian-pseudo-coherent}).
Ainsi, $R\Gamma(X, \mathcal{O}_X)$ est pseudo-cohérent dans la catégorie dérivée
de $R[x_e; e \in E']$, voir
Compléments d'algèbre, lemme \ref{more-algebra-lemma-pull-pseudo-coherent}.
Puisque son seul module de cohomologie non nul
est $H^0(X, \mathcal{O}_X)$, on en conclut qu'il s'agit d'un
$R[x_e; e \in E']$-module de type fini, voir
Compléments d'algèbre, lemme \ref{more-algebra-lemma-n-pseudo-module}.
Ceci achève la démonstration.
\end{proof}

\begin{lemma}
\label{more-morphisms-lemma-h1-fibre-zero-h0-flat}
Considérons un diagramme commutatif
$$
\xymatrix{
X \ar[rr]_f \ar[rd] & & Y \ar[ld] \\
& S
}
$$
de morphismes de schémas. Supposons que $X \to S$
soit plat, que $f$ soit propre, que $\dim(X_y) \leq 1$ pour $y \in Y$, et que
$R^1f_*\mathcal{O}_X = 0$. Alors $f_*\mathcal{O}_X$
est $S$-plat et la formation de $f_*\mathcal{O}_X$ commute
à tout changement de base $S' \to S$.
\end{lemma}

\begin{proof}
On peut supposer que $Y$ et $S$ sont affines, disons $S = \Spec(A)$.
Pour montrer que le $\mathcal{O}_Y$-module quasi-cohérent
$f_*\mathcal{O}_X$ est plat relativement à $S$, il suffit
de montrer que $H^0(X, \mathcal{O}_X)$ est plat sur $A$
(certains détails sont omis).
D'après le lemme \ref{more-morphisms-lemma-globally-generated-vanishing}, on a
$H^1(X, \mathcal{O}_X \otimes_A M) = 0$ pour tout $A$-module $M$.
Puisque $\mathcal{O}_X$ est aussi plat sur $A$, on en déduit que le foncteur
$M \mapsto H^0(X, \mathcal{O}_X \otimes_A M)$ est exact.
De plus, ce foncteur commute aux sommes directes d'après
Cohomologie, lemme \ref{cohomology-lemma-quasi-separated-cohomology-colimit}.
Il est alors facile de vérifier que
$H^0(X, \mathcal{O}_X \otimes_A M) = M \otimes_A H^0(X, \mathcal{O}_X)$
fonctoriellement en $M$, ce qui donne la platitude voulue.
Enfin, si $S' \to S$ est un morphisme de schémas affines donné par
l'homomorphisme d'anneaux $A \to A'$, alors, dans le cas affine qui vient d'être traité,
on voit que
$$
H^0(X \times_S S', \mathcal{O}_{X \times_S S'}) =
H^0(X, \mathcal{O}_X \otimes_A A') = H^0(X, \mathcal{O}_X) \otimes_A A'
$$
Cela montre que la formation de $f_*\mathcal{O}_X$ commute
à tout changement de base $S' \to S$. Certains détails sont omis.
\end{proof}

\begin{lemma}
\label{more-morphisms-lemma-h1-fibre-zero-isom}
Considérons un diagramme commutatif
$$
\xymatrix{
X \ar[rr]_f \ar[rd] & & Y \ar[ld] \\
& S
}
$$
de morphismes de schémas. Soit $s \in S$ un point. Supposons que
\begin{enumerate}
\item $X \to S$ soit localement de présentation finie et plat aux
points de $X_s$,
\item $Y \to S$ soit localement de présentation finie,
\item $f$ soit propre,
\item les fibres de $f_s : X_s \to Y_s$ soient de dimension $\leq 1$
et que $R^1f_{s, *}\mathcal{O}_{X_s} = 0$,
\item $\mathcal{O}_{Y_s} \to f_{s, *}\mathcal{O}_{X_s}$ soit un isomorphisme.
\end{enumerate}
Alors il existe un ouvert $Y_s \subset V \subset Y$ tel que
(a) $V$ soit plat sur $S$,
(b) $f^{-1}(V)$ soit plat sur $S$,
(c) $\dim(X_y) \leq 1$ pour $y \in V$,
(d) $R^1f_*\mathcal{O}_X|_V = 0$,
(e) $\mathcal{O}_V \to f_*\mathcal{O}_X|_V$
soit un isomorphisme, et tel que (a) -- (e)
restent vraies après le changement de base de $f^{-1}(V) \to V$ par tout $S' \to S$.
\end{lemma}

\begin{proof}
Soit $y \in Y_s$. On peut toujours remplacer $Y$ par un voisinage ouvert de $y$.
On peut donc supposer que $Y$ et $S$ sont affines. On peut également supposer que
$X$ est plat sur $S$, que $\dim(X_y) \leq 1$ pour $y \in Y$,
que $R^1f_*\mathcal{O}_X = 0$ universellement et que
$\mathcal{O}_Y \to f_*\mathcal{O}_X$ est surjectif, voir
le lemme \ref{more-morphisms-lemma-h1-fibre-zero-h0-kappa}. (Nous n'utiliserons pas tout ceci.)

\medskip\noindent
Supposons $S$ et $Y$ affines.
Écrivons $S = \lim S_i$ comme limite cofiltrante de schémas affines noethériens $S_i$.
D'après Limites, lemme \ref{limits-lemma-descend-finite-presentation},
il existe un élément $0 \in I$ et un diagramme
$$
\xymatrix{
X_0 \ar[rr]_{f_0} \ar[rd] & & Y_0 \ar[ld] \\
& S_0
}
$$
de morphismes de type fini entre schémas dont le changement de base à $S$
est le diagramme du lemme. Après avoir augmenté $0$, on peut supposer que $Y_0$
est affine et que $X_0 \to S_0$ est propre, voir
Limites, lemmes \ref{limits-lemma-eventually-proper} et
\ref{limits-lemma-limit-affine}.
Soit $s_0 \in S_0$ l'image de $s$.
Puisque $Y_s$ est affine, on voit que
$R^1f_{s, *}\mathcal{O}_{X_s} = 0$ équivaut
à $H^1(X_s, \mathcal{O}_{X_s}) = 0$.
Puisque $X_s$ est le changement de base de $X_{0, s_0}$ par
l'homomorphisme fidèlement plat $\kappa(s_0) \to \kappa(s)$,
on voit que $H^1(X_{0, s_0}, \mathcal{O}_{X_{0, s_0}}) = 0$
et donc que $R^1f_{0, *}\mathcal{O}_{X_{0, s_0}} = 0$.
De même, puisque $\mathcal{O}_{Y_s} \to f_{s, *}\mathcal{O}_{X_s}$
est un isomorphisme, il en va de même de
$\mathcal{O}_{Y_{0, s_0}} \to f_{0, *}\mathcal{O}_{X_{0, s_0}}$.
Puisque les dimensions des fibres de $X_s \to Y_s$ sont au plus $1$,
il en va de même pour le morphisme $X_{0, s_0} \to Y_{0, s_0}$.
Enfin, puisque $X \to S$ est plat, après avoir augmenté $0$,
on peut supposer que $X_0$ est plat sur $S_0$, voir
Limites, lemme \ref{limits-lemma-descend-flat-finite-presentation}.
Il suffit donc de démontrer le lemme pour
$X_0 \to Y_0 \to S_0$ et le point $s_0$.

\medskip\noindent
En combinant les réductions ci-dessus, on se ramène au cas où
$S$ et $Y$ sont affines, $S$ est noethérien, les fibres de $f$ sont de dimension
$\leq 1$, et $R^1f_*\mathcal{O}_X = 0$ universellement.
Soit $y \in Y_s$ un point. Affirmation :
$$
\mathcal{O}_{Y, y} \longrightarrow (f_*\mathcal{O}_X)_y
$$
est un isomorphisme. L'affirmation implique le lemme. En effet, puisque
$f_*\mathcal{O}_X$ est cohérent (Cohomologie des schémas, proposition
\ref{coherent-proposition-proper-pushforward-coherent}),
l'affirmation signifie que l'on peut remplacer $Y$ par un voisinage ouvert de $y$
et obtenir un isomorphisme $\mathcal{O}_Y \to f_*\mathcal{O}_X$.
On en conclut alors que $Y$ est plat sur $S$
d'après le lemme \ref{more-morphisms-lemma-h1-fibre-zero-h0-flat}.
Enfin, l'isomorphisme $\mathcal{O}_Y \to f_*\mathcal{O}_X$
reste un isomorphisme après tout changement de base $S' \to S$
d'après la dernière assertion du lemme \ref{more-morphisms-lemma-h1-fibre-zero-h0-flat}.

\medskip\noindent
Démonstration de l'affirmation. On sait déjà que
$\mathcal{O}_{Y, y} \longrightarrow (f_*\mathcal{O}_X)_y$
est surjectif (lemme \ref{more-morphisms-lemma-h1-fibre-zero-h0-kappa}),
que $(f_*\mathcal{O}_X)_y$ est
plat sur $\mathcal{O}_{S, s}$ (lemme \ref{more-morphisms-lemma-h1-fibre-zero-h0-flat})
et que le morphisme induit
$$
\mathcal{O}_{Y_s, y} =  \mathcal{O}_{Y, y}/\mathfrak m_s\mathcal{O}_{Y, y}
\longrightarrow
(f_*\mathcal{O}_X)_y/\mathfrak m_s (f_*\mathcal{O}_X)_y
\to
(f_{s, *}\mathcal{O}_{X_s})_y
$$
est injectif d'après l'hypothèse du lemme. Il résulte alors du
lemme \ref{algebra-lemma-mod-injective} d'Algèbre
que $\mathcal{O}_{Y, y} \longrightarrow (f_*\mathcal{O}_X)_y$
est injectif, comme souhaité.
\end{proof}

\begin{lemma}
\label{more-morphisms-lemma-bijection-on-Pic}
Soit $f : X \to Y$ un morphisme propre de schémas noethériens
tel que $f_*\mathcal{O}_X = \mathcal{O}_Y$, dont
les fibres de $f$ sont de dimension $\leq 1$, et tel que
$H^1(X_y, \mathcal{O}_{X_y}) = 0$ pour $y \in Y$.
Alors $f^* : \Pic(Y) \to \Pic(X)$ est une bijection sur
le sous-groupe des $\mathcal{L} \in \Pic(X)$ tels que
$\mathcal{L}|_{X_y} \cong \mathcal{O}_{X_y}$
pour tout $y \in Y$.
\end{lemma}

\begin{proof}
D'après la formule de projection
(Cohomologie, lemme \ref{cohomology-lemma-projection-formula}),
on voit que $f_*f^*\mathcal{N} \cong \mathcal{N}$ pour
$\mathcal{N} \in \Pic(Y)$.
Nous affirmons que, pour $\mathcal{L} \in \Pic(X)$ tel que
$\mathcal{L}|_{X_y} \cong \mathcal{O}_{X_y}$
pour tout $y \in Y$, le module $\mathcal{N} = f_*\mathcal{L}$
est inversible et que $\mathcal{L} \cong f^*\mathcal{N}$.
Cela achèvera la démonstration.

\medskip\noindent
Le $\mathcal{O}_Y$-module $\mathcal{N} = f_*\mathcal{L}$
est cohérent d'après Cohomologie des schémas, proposition
\ref{coherent-proposition-proper-pushforward-coherent}.
Ainsi, pour voir qu'il s'agit d'un $\mathcal{O}_Y$-module inversible,
il suffit de le vérifier sur les tiges (Algèbre, lemme
\ref{algebra-lemma-finite-projective}).
Puisque l'homomorphisme d'un anneau local noethérien vers son complété
est fidèlement plat, il suffit de vérifier que le complété
$(f_*\mathcal{L})_y^\wedge$ est libre (voir
Algèbre, section \ref{algebra-section-completion-noetherian} et
lemme \ref{algebra-lemma-finite-projective-descends}).
Pour cela, nous utiliserons le théorème des fonctions formelles sous la forme donnée dans
Cohomologie des schémas, lemme \ref{coherent-lemma-formal-functions-stalk}.
Puisque $f_*\mathcal{O}_X = \mathcal{O}_Y$ et donc
$(f_*\mathcal{O}_X)_y^\wedge \cong \mathcal{O}_{Y, y}^\wedge$,
il suffit de montrer que $\mathcal{L}|_{X_n} \cong \mathcal{O}_{X_n}$
pour tout $n$ (de manière compatible lorsque $n$ varie.
D'après le lemme \ref{more-morphisms-lemma-picard-group-first-order-thickening},
on a une suite exacte
$$
H^1(X_y, \mathfrak m_y^n\mathcal{O}_X/\mathfrak m_y^{n + 1}\mathcal{O}_X)
\to \Pic(X_{n + 1}) \to \Pic(X_n)
$$
avec les notations du théorème des fonctions formelles.
Observons que l'on a une surjection
$$
\mathcal{O}_{X_y}^{\oplus r_n} \cong
\mathfrak m_y^n/\mathfrak m_y^{n + 1} \otimes_{\kappa(y)} \mathcal{O}_{X_y}
\longrightarrow
\mathfrak m_y^n\mathcal{O}_X/\mathfrak m_y^{n + 1}\mathcal{O}_X
$$
pour certains entiers $r_n \geq 0$.
Puisque $\dim(X_y) \leq 1$, cette surjection induit une surjection
sur les groupes de cohomologie de degré un (par l'annulation de la cohomologie en degrés $\geq 2$
qui résulte de Cohomologie, proposition
\ref{cohomology-proposition-vanishing-Noetherian}).
Le $H^1$ de la suite est donc nul et les morphismes de transition
$\Pic(X_{n + 1}) \to \Pic(X_n)$ sont injectifs, comme souhaité.

\medskip\noindent
Il reste à montrer que $f^*\mathcal{N} \cong \mathcal{L}$.
Cela se démontre par la même méthode et nous omettons les détails.
\end{proof}









\section{Stratifications affines}
\label{more-morphisms-section-affine-stratifications}

\noindent
Ce contenu est tiré de \cite{RV}. Nous invitons le lecteur à se familiariser quelque peu
avec les stratifications dans
Topologie, section \ref{topology-section-stratifications},
avant de lire cette section.

\medskip\noindent
Si $X$ est un schéma, une stratification de $X$ désigne en général une
stratification de l'espace topologique sous-jacent à $X$.
Les strates sont des parties localement fermées. Nous considérerons ces
strates comme des sous-schémas localement fermés réduits de $X$ au moyen de
Schémas, remarque \ref{schemes-remark-reduced-induced-locally-closed}.

\begin{definition}
\label{more-morphisms-definition-affine-stratification}
Soit $X$ un schéma. Une {\it stratification affine} est une
stratification localement finie $X = \coprod_{i \in I} X_i$
dont les strates $X_i$ sont affines et telle que
les morphismes d'inclusion $X_i \to X$ soient affines.
\end{definition}

\noindent
La condition qu'une stratification $X = \coprod X_i$ soit localement finie est,
en présence de la condition que les morphismes d'inclusion $X_i \to X$
soient quasi-compacts, équivalente à la condition que les strates soient des parties localement
constructibles de $X$, voir Propriétés, lemme
\ref{properties-lemma-stratification-locally-finite-constructible}.

\medskip\noindent
La condition que $X_i \to X$ soit un morphisme affine ne dépend pas
de la structure de schéma que l'on met sur la partie localement fermée $X_i$, voir
le lemme \ref{more-morphisms-lemma-thicken-property-morphisms}. De plus,
si $X$ est séparé (ou, plus généralement, est à diagonale affine) et si
$X = \coprod X_i$ est une stratification localement finie à strates affines,
alors les morphismes $X_i \to X$ sont affines. Voir
Morphismes, lemme \ref{morphisms-lemma-affine-permanence}.
Cela permet, dans la plupart des cas qui nous intéressent, de faire abstraction de la condition d'affinité des
morphismes d'inclusion $X_i \to X$.

\medskip\noindent
Nous nous intéressons souvent au cas où l'ensemble partiellement ordonné d'indices $I$
de la stratification est fini. Rappelons que la {\it longueur} d'un
ensemble partiellement ordonné $I$ est le supremum des longueurs $p$ des
chaînes $i_0 < i_1 < \ldots < i_p$ d'éléments de $I$.

\begin{lemma}
\label{more-morphisms-lemma-improve-stratification}
Soit $X$ un schéma. Soit $X = \coprod_{i \in I} X_i$ une stratification
affine finie. Il existe une stratification affine
dont l'ensemble d'indices est $\{0, \ldots, n\}$, où $n$ est la longueur de $I$.
\end{lemma}

\begin{proof}
Rappelons que $I$ est muni d'un ordre partiel tel que
l'adhérence de $X_i$ soit contenue dans $\bigcup_{j \leq i} X_j$
pour tout $i \in I$.
Soit $I' \subset I$ l'ensemble des indices maximaux de $I$.
Si $i \in I'$, alors $X_i$ est ouvert dans $X$, car la réunion des
adhérences des autres strates est le complémentaire de $X_i$.
Soit $U = \bigcup_{i \in I'} X_i$, considéré comme un sous-schéma ouvert de $X$,
de sorte que $U_{red} = \coprod_{i \in I'} X_i$ en tant que schémas.
Alors $U$ est un schéma affine d'après
Schémas, lemme \ref{schemes-lemma-disjoint-union-affines},
et le lemme \ref{more-morphisms-lemma-thickening-affine-scheme}. Le morphisme $U \to X$
est affine, puisque chacun des $X_i \to X$, $i \in I'$, est affine, par le même raisonnement
utilisant le lemme \ref{more-morphisms-lemma-thicken-property-morphisms}.
Le complémentaire $Z = X \setminus U$, muni de la
structure de schéma réduite induite, possède la stratification
affine $Z = \bigcup_{i \in I \setminus I'} X_i$.
On utilise ici le fait qu'un morphisme de schémas $T \to Z$ est affine
si et seulement si la composée $T \to X$ est affine ; cela résulte
de Morphismes, lemmes \ref{morphisms-lemma-closed-immersion-affine},
\ref{morphisms-lemma-composition-affine} et
\ref{morphisms-lemma-affine-permanence}.
Observons que l'ensemble partiellement ordonné $I \setminus I'$
a une longueur exactement égale à celle de $I$ moins un.
Par récurrence, on trouve donc que $Z$ possède une stratification
affine $Z = Z_0 \amalg \ldots \amalg Z_{n - 1}$
dont l'ensemble d'indices est $\{1, \ldots, n\}$. En posant $Z_n = U$,
on obtient la stratification voulue de $X$.
\end{proof}

\noindent
Si un schéma $X$ possède une stratification affine finie, alors
$X$ est bien entendu quasi-compact. Il est un peu moins évident que
cela force également $X$ à être quasi-séparé.

\begin{lemma}
\label{more-morphisms-lemma-qc-affine-stratification}
Soit $X$ un schéma. Les assertions suivantes sont équivalentes :
\begin{enumerate}
\item $X$ possède une stratification affine finie, et
\item $X$ est quasi-compact et quasi-séparé.
\end{enumerate}
\end{lemma}

\begin{proof}
Soit $X = \bigcup X_i$ une stratification affine finie. Puisque chaque $X_i$
est affine, donc quasi-compact, on en conclut que $X$ est quasi-compact.
Soient $U, V \subset X$ des ouverts affines. Alors $U \cap X_i$ et $V \cap X_i$
sont des ouverts affines de $X_i$, puisque $X_i \to X$ est un morphisme affine. Ainsi,
$U \cap V \cap X_i$ est un ouvert affine du schéma affine $X_i$ (voir, par exemple,
Schémas, lemme \ref{schemes-lemma-characterize-separated}).
Par conséquent, $U \cap V = \coprod U \cap V \cap X_i$ est quasi-compact,
étant une réunion finie de strates affines. On en conclut que $X$ est
quasi-séparé d'après Schémas, lemme
\ref{schemes-lemma-characterize-quasi-separated}.

\medskip\noindent
Supposons que $X$ soit quasi-compact et quasi-séparé. On peut utiliser le
principe de récurrence de
Cohomologie des schémas, lemme \ref{coherent-lemma-induction-principle},
pour démontrer que $X$ possède une stratification affine finie.
Si $X$ est vide, il possède une stratification affine vide.
Si $X$ est affine non vide, il possède une stratification affine
à une strate. Supposons ensuite que $X = U \cup V$, où $U$ est un ouvert quasi-compact,
$V$ est un ouvert affine, et que l'on dispose de stratifications affines finies
$U = \bigcup_{i \in I} U_i$ et $U \cap V = \coprod_{j \in J} W_j$.
Notons $Z = X \setminus V$ et $Z' = X \setminus U$.
Remarquons que $Z$ est fermé dans $U$ et que $Z'$ est fermé dans $V$.
Observons que $U_i \cap Z$ et $U_i \cap W_j = U_i \times_U W_j$ sont des schémas
affines et affines au-dessus de $U$. (Indications : utiliser le fait que $U_i \times_U W_j \to W_j$ est
affine comme changement de base de $U_i \to U$, donc que $U_i \cap W_j$ est affine,
donc que $U_i \cap W_j \to U_i$ est affine, et donc que $U_i \cap W_j \to U$ est affine.)
Il s'ensuit que
$$
U = \coprod\nolimits_{i \in I} (U_i \cap Z) \amalg
\coprod\nolimits_{(i, j) \in I \times J} (U_i \cap W_j)
$$
est une stratification affine finie munie de l'ordre partiel sur
$I \amalg I \times J$ donné par $i' \leq (i, j) \Leftrightarrow i' \leq i$
et $(i', j') \leq (i, j) \Leftrightarrow i' \leq i$ et $j' \leq j$.
Observons que $(U_i \cap Z) \times_X V = \emptyset$ et que
$(U_i \cap W_j) \times_X V = U_i \cap W_j$ sont affines.
Ainsi, les morphismes $U_i \cap Z \to X$ et $U_i \cap W_j \to X$
sont affines, car on peut vérifier l'affinité d'un morphisme
localement sur le but
(Morphismes, lemme \ref{morphisms-lemma-characterize-affine})
et l'on a l'affinité au-dessus de $U$ et de $V$.
Pour achever la démonstration, on prend la stratification ci-dessus
et on ajoute une strate supplémentaire, à savoir $Z'$,
dont on ajoute l'indice comme élément minimal de l'ensemble partiellement ordonné.
\end{proof}

\begin{definition}
\label{more-morphisms-definition-affine-stratification-number}
Soit $X$ un schéma non vide, quasi-compact et quasi-séparé. Le
{\it nombre de stratification affine} est le plus petit entier $n \geq 0$
tel que les conditions équivalentes suivantes soient satisfaites :
\begin{enumerate}
\item il existe une stratification affine finie
$X = \coprod_{i \in I} X_i$ où $I$ est de longueur $n$,
\item il existe une stratification affine
$X = X_0 \amalg X_1 \amalg \ldots \amalg X_n$ dont
l'ensemble d'indices est $\{0, \ldots, n\}$.
\end{enumerate}
\end{definition}

\noindent
L'équivalence des conditions résulte du
lemme \ref{more-morphisms-lemma-improve-stratification}.
L'existence d'une stratification affine finie
est démontrée dans le lemme \ref{more-morphisms-lemma-qc-affine-stratification}.

\begin{lemma}
\label{more-morphisms-lemma-affine-stratification-number-bound}
Soit $X$ un schéma séparé qui possède un recouvrement ouvert par
$n + 1$ schémas affines. Alors le nombre de stratification affine de $X$
est au plus $n$.
\end{lemma}

\begin{proof}
Soit $X = U_0 \cup \ldots \cup U_n$ un recouvrement ouvert affine.
Posons
$$
X_i = (U_i \cup \ldots \cup U_n) \setminus
(U_{i + 1} \cup \ldots \cup U_n)
$$
Alors $X_i$ est affine en tant que sous-schéma fermé de $U_i$. Le morphisme $X_i \to X$
est affine d'après Morphismes, lemme \ref{morphisms-lemma-affine-permanence}.
Enfin, on a $\overline{X_i} \subset X_i \cup X_{i - 1} \cup \ldots X_0$.
\end{proof}

\begin{lemma}
\label{more-morphisms-lemma-affine-stratification-number-bound-Noetherian}
Soit $X$ un schéma noethérien de dimension $\infty > d \geq 0$.
Alors le nombre de stratification affine de $X$ est au plus $d$.
\end{lemma}

\begin{proof}
On raisonne par récurrence sur $d$. Si $d = 0$, alors $X$ est affine, voir
Propriétés, lemme \ref{properties-lemma-locally-Noetherian-dimension-0}.
Supposons $d > 0$. Soient $\eta_1, \ldots, \eta_n$ les points génériques
des composantes irréductibles de $X$ (Propriétés, lemme
\ref{properties-lemma-Noetherian-irreducible-components}).
On peut recouvrir $X$ par des ouverts affines contenant $\eta_1, \ldots, \eta_n$, voir
Propriétés, lemme \ref{properties-lemma-point-and-maximal-points-affine}.
Puisque $X$ est quasi-compact, on peut trouver un recouvrement ouvert affine fini
$X = \bigcup_{j = 1, \ldots, m} U_j$ tel que
$\eta_1, \ldots, \eta_n \in U_j$ pour tout $j = 1, \ldots, m$.
Choisissons un ouvert affine $U \subset U_1 \cap \ldots \cap U_m$
contenant $\eta_1, \ldots, \eta_n$ (ce qui est possible d'après le lemme
déjà cité). Alors le morphisme $U \to X$ est affine,
car $U \to U_j$ est affine pour tout $j$, voir
Morphismes, lemme \ref{morphisms-lemma-characterize-affine}.
Soit $Z = X \setminus U$.
Par construction, $\dim(Z) < \dim(X)$.
Par hypothèse de récurrence, on peut trouver une
stratification affine $Z = \bigcup_{i \in \{0, \ldots, n\}} Z_i$
de $Z$ avec $n \leq \dim(Z)$.
En posant $U = X_{n + 1}$ et $X_i = Z_i$ pour $i \leq n$,
on conclut.
\end{proof}

\begin{proposition}
\label{more-morphisms-proposition-vanishing-affine-stratification-number}
Soit $X$ un schéma non vide, quasi-compact et quasi-séparé, de
nombre de stratification affine $n$. Alors $H^p(X, \mathcal{F}) = 0$, $p > n$,
pour tout $\mathcal{O}_X$-module quasi-cohérent $\mathcal{F}$.
\end{proposition}

\begin{proof}
Nous allons le démontrer par récurrence sur le nombre de stratification affine $n$.
Si $n = 0$, alors $X$ est affine et le résultat est le
lemme
\ref{coherent-lemma-quasi-coherent-affine-cohomology-zero} de Cohomologie des schémas.
Supposons $n > 0$. D'après la définition \ref{more-morphisms-definition-affine-stratification-number},
il existe un schéma affine $U$ et une immersion ouverte affine $j : U \to X$
tels que le complémentaire $Z$ ait pour nombre de stratification affine $n - 1$.
Puisque $U$ et $j$ sont affines, on a $H^p(X, j_*(\mathcal{F}|_U)) = 0$
pour $p > 0$, voir Cohomologie des schémas, lemmes
\ref{coherent-lemma-relative-affine-cohomology} et
\ref{coherent-lemma-relative-affine-vanishing}.
Notons $\mathcal{K}$ et $\mathcal{Q}$ le noyau et le conoyau
du morphisme $\mathcal{F} \to j_*(\mathcal{F}|_U)$. On obtient ainsi
une suite exacte
$$
0 \to \mathcal{K} \to \mathcal{F} \to j_*(\mathcal{F}|_U)
\to \mathcal{Q} \to 0
$$
de $\mathcal{O}_X$-modules quasi-cohérents (voir
Schémas, section \ref{schemes-section-quasi-coherent}).
Un argument standard, consistant à décomposer notre suite exacte en suites exactes
courtes et à utiliser la suite exacte longue de cohomologie,
montre qu'il suffit de démontrer $H^p(X, \mathcal{K}) = 0$
et $H^p(X, \mathcal{Q}) = 0$ pour $p \geq n$.
Puisque $\mathcal{F} \to j_*(\mathcal{F}|_U)$ se restreint
à un isomorphisme sur $U$, on voit que $\mathcal{K}$
et $\mathcal{Q}$ sont supportés sur $Z$.
D'après Propriétés, lemme
\ref{properties-lemma-quasi-coherent-colimit-finite-type},
on peut écrire ces modules comme les colimites filtrantes de
leurs sous-modules quasi-cohérents de type fini.
En utilisant le fait que la cohomologie des faisceaux sur $X$ commute
aux colimites filtrantes, voir
Cohomologie, lemme \ref{cohomology-lemma-quasi-separated-cohomology-colimit},
on conclut qu'il suffit de montrer que, si $\mathcal{G}$
est un module quasi-cohérent de type fini dont le support
est contenu dans $Z$, alors $H^p(X, \mathcal{G}) = 0$ pour $p \geq n$.
Soit $Z' \subset X$ le support schématique de
$\mathcal{G} \oplus \mathcal{O}_Z$ ; on peut considérer, et l'on considère,
$\mathcal{G}$ comme un module quasi-cohérent sur $Z'$, voir
Morphismes, section \ref{morphisms-section-support}.
Alors $Z'$ et $Z$ ont le même espace topologique sous-jacent
et donc le même nombre de stratification affine, à savoir $n - 1$.
Ainsi, $H^p(X, \mathcal{G}) = H^p(Z', \mathcal{G})$
(égalité d'après Cohomologie des schémas, lemme
\ref{coherent-lemma-relative-affine-cohomology})
s'annule pour $p \geq n$ par hypothèse de récurrence.
\end{proof}

\begin{example}
\label{more-morphisms-example-affine-stratification-number}
Soit $k$ un corps et soit $X = \mathbf{P}^n_k$ l'espace projectif de dimension $n$
sur $k$. Le lemme \ref{more-morphisms-lemma-affine-stratification-number-bound}
s'y applique d'après Constructions, lemme
\ref{constructions-lemma-standard-covering-projective-space}.
Par conséquent, le nombre de stratification affine de $\mathbf{P}^n_k$ est
au plus $n$. D'autre part, on a une cohomologie non nulle
en degré $n$ pour certains modules quasi-cohérents sur
$\mathbf{P}^n_k$, voir
Cohomologie des schémas, lemme
\ref{coherent-lemma-cohomology-projective-space-over-ring}.
En appliquant la proposition \ref{more-morphisms-proposition-vanishing-affine-stratification-number},
on conclut que le nombre de stratification affine de
$\mathbf{P}^n_k$ est égal à $n$.
\end{example}








\section{Morphismes universellement ouverts}
\label{more-morphisms-section-universally-open}

\noindent
Quelques résultats sur les morphismes universellement ouverts.

\begin{lemma}
\label{more-morphisms-lemma-test-universally-open}
Soit $f : X \to S$ un morphisme de schémas.
Les conditions suivantes sont équivalentes :
\begin{enumerate}
\item $f$ est universellement ouvert,
\item pour tout morphisme $S' \to S$ localement de présentation finie,
le changement de base $X_{S'} \to S'$ est ouvert, et
\item pour tout $n$, le morphisme
$\mathbf{A}^n \times X \to \mathbf{A}^n \times S$
est ouvert.
\end{enumerate}
\end{lemma}

\begin{proof}
Il est clair que (1) implique (2) et que (2) implique (3).
Montrons que (3) implique (1).
Supposons que le changement de base $X_T \to T$ ne soit pas ouvert
pour un certain morphisme de schémas $g : T \to S$.
On peut alors trouver des ouverts affines
$V \subset S$, $U \subset X$, $W \subset T$
tels que $f(U) \subset V$ et $g(W) \subset V$
et que $U \times_V W \to W$ ne soit pas ouvert.
Si l'on peut montrer que cela implique que
$\mathbf{A}^n \times U \to \mathbf{A}^n \times V$
n'est pas ouvert, alors $\mathbf{A}^n \times X \to \mathbf{A}^n \times S$
n'est pas ouvert et la démonstration est achevée. On est ainsi ramené
au résultat démontré dans le paragraphe suivant.

\medskip\noindent
Soit $A \to B$ un homomorphisme d'anneaux tel que $A' \to B' = A' \otimes_A B$
n'induise pas une application ouverte entre spectres pour une certaine $A$-algèbre $A'$.
Comme les ouverts principaux forment une base de la topologie de $\Spec(B')$,
on conclut que l'image de $D(g)$ dans $\Spec(A')$
n'est pas ouverte pour un certain $g \in B'$. Écrivons
$g = \sum_{i = 1, \ldots, n} a'_i \otimes b_i$
pour un certain $n$, avec $a'_i \in A'$ et $b_i \in B$.
Considérons l'élément $h = \sum_{i = 1, \ldots, n} x_i b_i$
de $B[x_1, \ldots, x_n]$. Supposons, afin d'obtenir une contradiction, que $D(h)$ ait une image
ouverte par le morphisme
$$
\Spec(B[x_1, \ldots, x_n]) \longrightarrow \Spec(A[x_1, \ldots, x_n])
$$
Alors $D(h)$ s'enverrait surjectivement sur un ouvert quasi-compact
$U \subset \Spec(A[x_1, \ldots, x_n])$.
Soit $A[x_1, \ldots, x_n] \to A'$ l'homomorphisme de $A$-algèbres
qui envoie $x_i$ sur $a'_i$. Il induit également un homomorphisme de $B$-algèbres
$B[x_1, \ldots, x_n] \to B'$ qui envoie $h$ sur $g$.
Puisque
$$
\xymatrix{
\Spec(B[x_1, \ldots, x_n]) \ar[d] &
\Spec(B') \ar[l] \ar[d] \\
\Spec(A[x_1, \ldots, x_n]) &
\Spec(A') \ar[l]
}
$$
est cartésien, l'image de $D(g)$ dans $\Spec(A')$ est égale à
l'image réciproque de $U$ dans $\Spec(A')$, et est donc ouverte, ce qui est
la contradiction recherchée.
\end{proof}

\begin{lemma}
\label{more-morphisms-lemma-quasi-finite-Noetherian-universally-open}
Soit $f : X \to Y$ un morphisme de schémas. Si
\begin{enumerate}
\item $f$ est localement quasi-fini,
\item $Y$ est géométriquement unibranche et localement noethérien, et
\item toute composante irréductible de $X$ domine
une composante irréductible de $Y$,
\end{enumerate}
alors $f$ est universellement ouvert.
\end{lemma}

\begin{proof}
Pour tout $n$, le schéma $\mathbf{A}^n \times Y$ est géométriquement unibranche
d'après le lemme \ref{more-morphisms-lemma-number-of-branches-and-smooth} et
Propriétés, lemme \ref{properties-lemma-number-of-branches-1}. Par conséquent,
les hypothèses du lemme sont satisfaites par les morphismes
$\mathbf{A}^n \times X \to \mathbf{A}^n \times Y$ pour tout $n$.
D'après le lemme \ref{more-morphisms-lemma-test-universally-open},
il suffit de montrer que $f$ est ouvert. D'après Morphismes, lemme
\ref{morphisms-lemma-locally-finite-presentation-universally-open},
il suffit de montrer que les générisations se relèvent le long de $f$.
Supposons que $y' \leadsto y$ soit une spécialisation de points de
$Y$ et que $x \in X$ soit un point s'envoyant sur $y$.
Comme dans le lemme \ref{more-morphisms-lemma-etale-makes-quasi-finite-finite-at-point},
choisissons un diagramme
$$
\xymatrix{
u \ar[d] & U \ar[d] \ar[r] & X \ar[d] \\
v & V \ar[r] & Y
}
$$
où $(V, v) \to (Y, y)$ est un voisinage \'etale élémentaire,
$U \to V$ est fini, $u$ est l'unique point de $U$ s'envoyant sur $v$,
$U \subset V \times_Y X$ est ouvert, et $v \mapsto y$ et $u \mapsto x$.
Soit $E$ une composante irréductible de $U$ passant par $u$
(il en existe au moins une). Comme $U \to X$ est \'etale, $E$
s'envoie dans une composante irréductible de $X$,
qui à son tour domine une composante irréductible de $Y$ (par hypothèse).
Comme $U \to V$ est fini, donc fermé, on conclut que
l'image $E' \subset V$ de $E$ est une partie fermée irréductible
passant par $v$ et dominant une composante irréductible de $Y$.
Comme $V \to Y$ est \'etale, $E'$ est nécessairement une composante irréductible
de $V$ passant par $v$.
Puisque $Y$ est géométriquement unibranche, on voit que $E'$ est l'unique composante irréductible
de $V$ passant par $v$ (lemme \ref{more-morphisms-lemma-nr-branches}).
Comme $V$ est localement noethérien, on peut, après avoir rétréci $V$,
supposer que $E' = V$ (égalité d'ensembles).

\medskip\noindent
Comme $V \to Y$ est \'etale, on peut trouver une spécialisation
$v' \leadsto v$ dont l'image est $y' \leadsto y$.
D'après ce qui précède, on peut trouver $u' \in U$ s'envoyant sur $v'$.
Alors $u' \leadsto u$, car $u$ est l'unique point de
$U$ s'envoyant sur $v$ et $U \to V$ est fermé.
Enfin, l'image $x' \in X$ de $u'$ est un point
qui se spécialise en $x$ et s'envoie sur $y'$, ce qui achève la démonstration.
\end{proof}

\begin{lemma}
\label{more-morphisms-lemma-characterize-universally-open-finite}
Soit $A \to B$ un homomorphisme d'anneaux. Supposons que $B$ soit engendré comme $A$-module par
$b_1, \ldots, b_d \in B$. Posons $h = \sum x_ib_i \in B[x_1, \ldots, x_d]$.
Alors $\Spec(B) \to \Spec(A)$ est universellement ouvert si et seulement si l'image de
$D(h)$ dans $\Spec(A[x_1, \ldots, x_d])$ est ouverte.
\end{lemma}

\begin{proof}
Si $\Spec(B) \to \Spec(A)$ est universellement ouvert, alors, bien entendu,
l'image de $D(h)$ est ouverte. Réciproquement, supposons que l'image $U$ de
$D(h)$ soit ouverte. Soit $A \to A'$ un homomorphisme d'anneaux. Il suffit de montrer
que l'image de tout ouvert principal $D(g) \subset \Spec(A' \otimes_A B)$
dans $\Spec(A')$ est ouverte. On peut écrire
$g = \sum_{i = 1, \ldots, d} a'_i \otimes b_i$ avec certains $a'_i \in A'$.
Soit $A[x_1, \ldots, x_n] \to A'$ l'homomorphisme de $A$-algèbres
qui envoie $x_i$ sur $a'_i$. Il induit également un homomorphisme de $B$-algèbres
$B[x_1, \ldots, x_n] \to A' \otimes_A B$ qui envoie $h$ sur $g$.
Puisque
$$
\xymatrix{
\Spec(B[x_1, \ldots, x_n]) \ar[d] &
\Spec(B') \ar[l] \ar[d] \\
\Spec(A[x_1, \ldots, x_n]) &
\Spec(A') \ar[l]
}
$$
est cartésien, l'image de $D(g)$ dans $\Spec(A')$ est égale à
l'image réciproque de $U$ dans $\Spec(A')$, et est donc ouverte.
\end{proof}

\begin{lemma}
\label{more-morphisms-lemma-descend-quasi-finite-universally-open}
Soit $S = \lim S_i$ la limite d'un système cofiltrant de schémas
dont les morphismes de transition sont affines.
Soit $0 \in I$ et soit $f_0 : X_0 \to Y_0$ un morphisme de schémas sur $S_0$.
Supposons $S_0$, $X_0$, $Y_0$ quasi-compacts et quasi-séparés.
Soit $f_i : X_i \to Y_i$ le changement de base de $f_0$ à $S_i$ et
soit $f : X \to Y$ le changement de base de $f_0$ à $S$.
Si
\begin{enumerate}
\item $f$ est localement quasi-fini et universellement ouvert, et
\item $f_0$ est localement de présentation finie,
\end{enumerate}
alors il existe $i \geq 0$ tel que $f_i$ soit localement quasi-fini
et universellement ouvert.
\end{lemma}

\begin{proof}
D'après Limites, lemme \ref{limits-lemma-descend-quasi-finite}, après avoir augmenté
$0$, on peut supposer que $f_0$ est localement quasi-fini. Soit $x \in X$.
Par localisation \'etale des morphismes quasi-finis,
on peut trouver un diagramme
$$
\xymatrix{
X \ar[d] & U \ar[l] \ar[d] \\
Y & V \ar[l]
}
$$
où $V \to Y$ est \'etale, $U \subset X_V$ est ouvert, $U \to V$ est fini,
et $x$ appartient à l'image de $U \to X$, voir
le lemme \ref{more-morphisms-lemma-etale-makes-quasi-finite-finite-at-point}.
Après avoir rétréci $V$, on peut supposer que $V$ et $U$ sont affines.
Comme $X$ est quasi-compact, il s'ensuit, en prenant une somme disjointe finie
de tels $V$ et $U$, que l'on peut construire un diagramme comme ci-dessus
tel que $U \to X$ soit surjectif. D'après
Limites, lemmes \ref{limits-lemma-descend-finite-presentation},
\ref{limits-lemma-descend-opens},
\ref{limits-lemma-descend-surjective},
\ref{limits-lemma-descend-finite-finite-presentation},
\ref{limits-lemma-descend-etale}, et
\ref{limits-lemma-limit-affine}, après avoir éventuellement augmenté $0$,
on peut supposer que l'on dispose d'un diagramme
$$
\xymatrix{
X_0 \ar[d] & U_0 \ar[l] \ar[d] \\
Y_0 & V_0 \ar[l]
}
$$
où $V_0$ est affine, $V_0 \to Y_0$ est \'etale,
$U_0 \subset (X_0)_{V_0}$ est ouvert, $U_0 \to V_0$ est fini,
et $U_0 \to X_0$ est surjectif. Comme $V_i \to Y_i$ est \'etale
et donc universellement ouvert, il s'ensuit qu'il suffit
de démontrer que $U_i \to V_i$ est universellement ouvert pour
$i$ assez grand. On est ainsi ramené au cas traité dans le paragraphe
suivant.

\medskip\noindent
Soit $A = \colim A_i$ une colimite filtrante d'anneaux.
Soit $A_0 \to B_0$ un homomorphisme d'anneaux. Posons $B = A \otimes_{A_0} B_0$ et
$B_i = A_i \otimes_{A_0} B_0$. Supposons que $A_0 \to B_0$ soit fini,
de présentation finie, et que $A \to B$ soit universellement ouvert.
Il faut montrer que $A_i \to B_i$ est universellement ouvert pour $i$ assez grand.
Choisissons $b_{0, 1}, \ldots, b_{0, d} \in B_0$ qui engendrent $B_0$
comme $A_0$-module. Posons $h_0 = \sum_{j = 1, \ldots, d} x_jb_{0, j}$
dans $B_0[x_1, \ldots, x_d]$. Notons $h$, resp.\ $h_i$, l'image de $h_0$
dans $B[x_1, \ldots, x_d]$, resp.\ $B_i[x_1, \ldots, x_d]$.
L'image $U$ de $D(h)$ dans $\Spec(A[x_1, \ldots, x_d])$ est ouverte,
car $A \to B$ est universellement ouvert. Bien entendu, $U$ est quasi-compact
puisqu'il est l'image d'un schéma affine. Pour $i$ assez grand, il
existe un ouvert quasi-compact $U_i \subset \Spec(A_i[x_1, \ldots, x_d])$
dont l'image réciproque dans $\Spec(A[x_1, \ldots, x_d])$ est $U$, voir
Limites, lemme \ref{limits-lemma-descend-opens}.
Après avoir augmenté $i$, on peut supposer que $D(h_i)$ s'envoie
dans $U_i$ ; cela résulte du même lemme appliqué à
l'image réciproque de $U_i$ dans $D(h_i)$. Enfin, pour $i$ encore plus grand,
le morphisme de schémas $D(h_i) \to U_i$ sera surjectif
par application du lemme déjà utilisé,
Limites, lemme \ref{limits-lemma-descend-surjective}.
On conclut que $A_i \to B_i$ est universellement ouvert d'après le
lemme \ref{more-morphisms-lemma-characterize-universally-open-finite}.
\end{proof}

\begin{lemma}
\label{more-morphisms-lemma-count-geometric-fibres}
Soit $f : X \to Y$ un morphisme localement quasi-fini. Alors
\begin{enumerate}
\item les fonctions $n_{X/Y}$ des
lemmes \ref{more-morphisms-lemma-base-change-fibres-nr-geometrically-irreducible-components}
et \ref{more-morphisms-lemma-base-change-fibres-nr-geometrically-connected-components}
coïncident,
\item si $X$ est quasi-compact, alors $n_{X/Y}$ atteint un maximum $d < \infty$.
\end{enumerate}
\end{lemma}

\begin{proof}
L'égalité des fonctions résulte immédiatement du fait que les
fibres (géométriques) d'un morphisme localement quasi-fini sont discrètes, voir
Morphismes, lemme \ref{morphisms-lemma-locally-quasi-finite-fibres}.
Le caractère borné résulte de
Morphismes, lemmes \ref{morphisms-lemma-characterize-universally-bounded} et
\ref{morphisms-lemma-locally-quasi-finite-qc-source-universally-bounded}.
\end{proof}

\begin{lemma}
\label{more-morphisms-lemma-count-geometric-fibres-universally-open}
Soit $f : X \to Y$ un morphisme de schémas séparé, localement quasi-fini et universellement ouvert.
Soit $n_{X/Y}$ comme dans le
lemme \ref{more-morphisms-lemma-count-geometric-fibres}.
Si $n_{X/Y}(y) \geq d$ pour un certain $y \in Y$ et $d \geq 0$,
alors $n_{X/Y} \geq d$ dans un voisinage ouvert de $y$.
\end{lemma}

\begin{proof}
La question est locale sur $Y$ ; on peut donc supposer $Y$ affine.
Soit $K$ une clôture algébrique du corps résiduel $\kappa(y)$.
Notre hypothèse est que $(X_y)_K$ possède $\geq d$ composantes connexes.
Alors, pour un ouvert quasi-compact convenable $X' \subset X$,
le schéma $(X'_y)_K$ possède $\geq d$ composantes connexes ; détails omis.
Après avoir remplacé $X$ par $X'$, on peut supposer $X$ quasi-compact.
Alors $f$ est quasi-fini. Soient $x_1, \ldots, x_n$ les points de $X$
au-dessus de $y$. Appliquons le
lemme \ref{more-morphisms-lemma-etale-splits-off-quasi-finite-part-technical-variant}
pour obtenir un voisinage \'etale $(U, u) \to (Y, y)$ et une décomposition
$$
U \times_Y X =
W \amalg
\ \coprod\nolimits_{i = 1, \ldots, n}
\ \coprod\nolimits_{j = 1, \ldots, m_i}
V_{i, j}
$$
comme au lemme cité. Observons que $n_{X/Y}(y) = \sum_i m_i$
dans cette situation ; certains détails sont omis.
Comme $f$ est universellement ouvert, on voit
que $V_{i, j} \to U$ est ouvert pour tous $i, j$. Après avoir rétréci
$U$, on peut donc supposer $V_{i, j} \to U$ surjectif
pour tous $i, j$. Cela démontre que
$n_{U \times_Y X/U} \geq \sum_i m_i = n_{X/Y}(y) \geq d$.
Comme la construction de $n_{X/Y}$ est compatible avec le
changement de base, la démonstration est achevée.
\end{proof}

\begin{lemma}
\label{more-morphisms-lemma-large-open}
Soit $f : X \to Y$ un morphisme de schémas séparé, localement quasi-fini et universellement ouvert.
Soit $n_{X/Y}$ comme dans le
lemme \ref{more-morphisms-lemma-count-geometric-fibres}.
Si $n_{X/Y}$ atteint un maximum $d < \infty$, alors l'ensemble
$$
Y_d = \{y \in Y \mid n_{X/Y}(y) = d\}
$$
est ouvert dans $Y$ et le morphisme $f^{-1}(Y_d) \to Y_d$ est fini.
\end{lemma}

\begin{proof}
L'ouverture de $Y_d$ résulte immédiatement du
lemme \ref{more-morphisms-lemma-count-geometric-fibres-universally-open}.
Pour démontrer la finitude au-dessus de $Y_d$, reprenons l'argument de la
démonstration de ce lemme. Soit donc $y \in Y_d$. Il y a alors
au plus $d$ points de $X$ au-dessus de $y$. Notons
$x_1, \ldots, x_n$ les points de $X$
au-dessus de $y$. Appliquons le
lemme \ref{more-morphisms-lemma-etale-splits-off-quasi-finite-part-technical-variant}
pour obtenir un voisinage \'etale $(U, u) \to (Y, y)$ et une décomposition
$$
U \times_Y X =
W \amalg
\ \coprod\nolimits_{i = 1, \ldots, n}
\ \coprod\nolimits_{j = 1, \ldots, m_i}
V_{i, j}
$$
comme au lemme cité. Observons que $d = n_{X/Y}(y) = \sum_i m_i$
dans cette situation ; certains détails sont omis.
Comme $f$ est universellement ouvert, on voit
que $V_{i, j} \to U$ est ouvert pour tous $i, j$. Après avoir rétréci
$U$, on peut donc supposer $V_{i, j} \to U$ surjectif
pour tous $i, j$, et l'on peut supposer que $U$ s'envoie dans $W$.
Cela démontre que $n_{U \times_Y X/U} \geq \sum_i m_i = d$.
Comme la construction de $n_{X/Y}$ est compatible avec le
changement de base, on sait que $n_{U \times_Y X/U} = d$. Cela signifie que
$W$ doit être vide, et on conclut que $U \times_Y X \to U$ est fini.
D'après Descente, lemme \ref{descent-lemma-descending-property-finite},
cela implique que $X \to Y$ est fini au-dessus de
l'image du morphisme ouvert $U \to Y$. Autrement dit,
on voit que $f$ est fini au-dessus d'un voisinage ouvert de $y$,
comme souhaité.
\end{proof}






\section{Pondérations}
\label{more-morphisms-section-weightings}

\noindent
Le contenu de cette section est tiré de \cite[Exposé XVII, 6.2.4]{SGA4}.

\medskip\noindent
Soit $\pi : U \to V$ un morphisme de schémas localement quasi-fini à
fibres finies. Étant donnée une fonction $w : U \to \mathbf{Z}$, on définit une fonction
$$
\textstyle{\int}_\pi w : V \longrightarrow \mathbf{Z},\quad
v \longmapsto
\sum\nolimits_{u \in U,\ \pi(u) = v} w(u) [\kappa(u) : \kappa(v)]_s
$$
Remarquons que les extensions de corps sont finies
(Morphismes, lemme \ref{morphisms-lemma-residue-field-quasi-finite}),
que $[\kappa' : \kappa]_s$ est le degré séparable
(Corps, définition \ref{fields-definition-insep-degree}),
et que la somme est finie puisque les fibres de $\pi$ sont supposées finies.
Une autre manière de calculer la valeur de $\int_\pi w$ en un point $v \in V$
est la suivante.
Choisissons un corps algébriquement clos $k$ et un morphisme
$\overline{v} : \Spec(k) \to V$ dont l'image est $v$. On a alors
$$
(\textstyle{\int}_\pi w)(v) =
\sum\nolimits_{\overline{u} \in U_{\overline{v}}} w(\overline{u})
$$
où, bien entendu, $w(\overline{u})$ désigne la valeur de $w$ en
l'image $u$ du point $\overline{u}$ par le morphisme
$U_{\overline{v}} \to U$.
Remarquons que l'on peut considérer tout $\overline{u} \in U_{\overline{v}}$ comme un morphisme
$\overline{u} : \Spec(k) \to U$ tel que
$\pi \circ \overline{u} = \overline{v}$. En effet, puisque
$U \to V$ est localement quasi-fini à fibres finies,
le schéma $U_{\overline{v}}$ est le spectre
d'une algèbre de dimension finie sur $k$ dont tous les idéaux premiers
sont des idéaux maximaux de corps résiduel $k$. Pour voir que l'égalité
est satisfaite, remarquons que le nombre de morphismes $\overline{u}$ au-dessus
d'un $u$ donné est égal à $[\kappa(u) : \kappa(v)]_s$ d'après
Corps, lemme \ref{fields-lemma-separable-degree}.

\begin{lemma}
\label{more-morphisms-lemma-weighting-check-after-etale-base-change}
Étant donnés un carré cartésien
$$
\xymatrix{
U \ar[d]_\pi & U' \ar[l]^h \ar[d]^{\pi'} \\
V & V' \ar[l]_g
}
$$
où $\pi$ est localement quasi-fini à fibres finies,
et une fonction $w : U \to \mathbf{Z}$,
on a $(\int_\pi w) \circ g = \int_{\pi'} (w \circ h)$.
\end{lemma}

\begin{proof}
Cela résulte immédiatement de la seconde description de $\int_\pi w$
donnée ci-dessus. Pour le démontrer à partir de la définition, on utilise le fait que, si
$E/F$ est une extension finie de corps et $F'/F$ une autre extension de corps,
alors, en écrivant $(E \otimes_F F')_{red} = \prod E'_i$ comme un produit de corps
finis sur $F'$, on a
$$
[E : F]_s = \sum [E'_i : F']_s
$$
Pour démontrer cette égalité, choisissons une extension algébriquement close
$\Omega/F'$ et observons que
\begin{align*}
[E : F]_s
& =
|\Mor_F(E, \Omega)| \\
& =
|\Mor_{F'}(E \otimes_F F', \Omega)| \\
& =
|\Mor_{F'}((E \otimes_F F')_{red}, \Omega)| \\
& =
\sum |\Mor_{F'}(E'_i, \Omega)| \\
& =
\sum [E'_i : F']_s
\end{align*}
où l'on a utilisé Corps, lemme \ref{fields-lemma-separable-degree}.
\end{proof}

\begin{definition}
\label{more-morphisms-definition-weighting}
Soit $f : X \to Y$ un morphisme localement quasi-fini. Une
{\it fonction de poids} ou une {\it pond\'eration} de $f$ est une application
$w : X \to \mathbf{Z}$ telle que, pour tout diagramme
$$
\xymatrix{
X \ar[d]_f & U \ar[l]^h \ar[d]^\pi \\
Y & V \ar[l]_g
}
$$
où $V \to Y$ est \'etale, $U \subset X_V$ est ouvert et $U \to V$ fini,
la fonction $\int_\pi (w \circ h)$ soit localement constante.
\end{definition}

\noindent
Bien entendu, en prenant $w = 0$, on obtient une pondération de tout morphisme
$f$ localement quasi-fini, bien qu'elle ne soit guère intéressante. Il s'avérera que les
pondérations {\it positives}, c'est-à-dire $w : X \to \mathbf{Z}_{> 0}$, sont
les plus intéressantes à diverses fins.

\begin{lemma}
\label{more-morphisms-lemma-weighting-base-change}
Soit $f : X \to Y$ un morphisme localement quasi-fini.
Soit $w : X \to \mathbf{Z}$ une pondération. Soit $f' : X' \to Y'$
le changement de base de $f$ par un morphisme $Y' \to Y$. Alors la
composée $w' : X' \to \mathbf{Z}$ de $w$ et de la projection $X' \to X$
est une pondération de $f'$.
\end{lemma}

\begin{proof}
Considérons un diagramme
$$
\xymatrix{
X' \ar[d]_{f'} & U' \ar[l]^{h'} \ar[d]^{\pi'} \\
Y' & V' \ar[l]_{g'}
}
$$
comme dans la définition \ref{more-morphisms-definition-weighting} pour le morphisme $f'$.
Pour tout $v' \in V'$, il faut montrer que $\int_{\pi'} (w' \circ h')$ est
constante dans un voisinage ouvert de $v'$. D'après le
lemme \ref{more-morphisms-lemma-weighting-check-after-etale-base-change}
(et le fait que les morphismes \'etales sont ouverts),
on peut remplacer $V'$ par n'importe quel voisinage \'etale de $v'$.
Après avoir remplacé $V'$ par un voisinage \'etale de $v'$,
on peut supposer que $U' = U'_1 \amalg \ldots \amalg U'_n$,
où chaque $U'_i$ possède un unique point $u'_i$ au-dessus de $v'$,
tel que $\kappa(u'_i)/\kappa(v')$ soit purement inséparable,
voir le lemme \ref{more-morphisms-lemma-etale-splits-off-quasi-finite-part-technical-variant}.
Il suffit clairement de démontrer que $\int_{U'_i \to V'} w'|_{U'_i}$
est constante dans un voisinage de $v'$.
On est ainsi ramené au cas traité dans le paragraphe suivant.

\medskip\noindent
On a $v' \in V'$ et il existe un unique point $u'$ de $U'$
au-dessus de $v'$ tel que $\kappa(u')/\kappa(v')$ soit purement inséparable.
Notons $x \in X$ et $y \in Y$ les images de $u'$ et de $v'$.
On peut trouver un voisinage \'etale $(V, v) \to (Y, y)$
et un ouvert $U \subset X_V$ tels que $\pi : U \to V$ soit fini
et qu'il existe un unique point $u \in U$ au-dessus de $v$
qui s'envoie sur $x \in X$ par la projection $h : U \to X$,
et tel que, de plus, $\kappa(u)/\kappa(v)$ soit
purement inséparable. Cela est possible d'après le lemme utilisé ci-dessus.
Considérons le morphisme
$$
U'' = U \times_X U' \longrightarrow V \times_Y V' = V''
$$
Comme $u$ et $u'$ s'envoient tous deux sur $x \in X$, il existe un point
$u'' \in U''$ s'envoyant sur $(u, u')$. Notons $v'' \in V''$
l'image de $u''$. Après avoir remplacé $V', v'$ par $V'', v''$,
on peut supposer que la composée $V' \to Y' \to Y$ se factorise
par un morphisme de voisinages \'etales $(V', v') \to (V, v)$
tel que le morphisme induit $X'_{V'} = X_{V'} \to X_V$ envoie
$u'$ sur $u$. Dans le changement de base $X'_{V'} = X_{V'}$, on a
deux sous-schémas ouverts, à savoir $U'$ et l'image réciproque $U_{V'}$ de
$U \subset X_V$. Par construction, tous deux contiennent un unique point
au-dessus de $v'$, à savoir $u'$ dans les deux cas.
Ainsi, après avoir rétréci $V'$, on peut supposer que ces ouverts
sont égaux ; en effet, $U' \setminus (U' \cap U_{V'})$ et
$U_{V'} \setminus (U' \cap U_{V'})$ ont
une image fermée dans $V'$, et ces images ne contiennent pas $v'$.
Ainsi $U' = U_{V'}$, et l'on obtient un diagramme cartésien comme dans le
lemme \ref{more-morphisms-lemma-weighting-check-after-etale-base-change}.
Comme $\int_\pi (w \circ h)$ est localement constante
par hypothèse, on conclut.
\end{proof}

\begin{lemma}
\label{more-morphisms-lemma-weighting-open}
Soit $f : X \to Y$ un morphisme localement quasi-fini. Soit
$w : X \to \mathbf{Z}$ une pondération de $f$. Si $X' \subset X$ est ouvert,
alors $w|_{X'}$ est une pondération de $f|_{X'} : X' \to Y$.
\end{lemma}

\begin{proof}
Cela résulte immédiatement de la définition.
\end{proof}

\begin{lemma}
\label{more-morphisms-lemma-weighting-composition}
Soient $f : X \to Y$ et $g : Y \to Z$ des morphismes localement quasi-finis.
Soit $w_f : X \to \mathbf{Z}$ une pondération de $f$ et soit
$w_g : Y \to \mathbf{Z}$ une pondération de $g$. Alors la fonction
$$
X \longrightarrow \mathbf{Z},\quad
x \longmapsto w_f(x) w_g(f(x))
$$
est une pondération de $g \circ f$.
\end{lemma}

\begin{proof}
Posons $w_{g \circ f}(x) = w_f(x) w_g(f(x))$ pour $x \in X$.
Considérons un diagramme
$$
\xymatrix{
X \ar[d]_{g \circ f} & U \ar[l] \ar[d]^\pi \\
Z & W \ar[l]
}
$$
où $W \to Z$ est \'etale, $U \subset X_W$ est ouvert et $U \to W$ fini.
Il faut montrer que $\int_\pi w_{g \circ f}|_U$ est localement constante.
Choisissons un point $w \in W$. D'après le
lemme \ref{more-morphisms-lemma-weighting-check-after-etale-base-change}
(et le fait que les morphismes \'etales sont ouverts), il suffit de montrer que
$\int_\pi w_{g \circ f}|_U$ est constante après avoir remplacé
$(W, w)$ par un voisinage \'etale.
Après avoir remplacé $(W, w)$ par un voisinage \'etale, on peut
supposer $U = U_1 \amalg \ldots \amalg U_n$, où chaque $U_i$
possède un unique point $u_i$ au-dessus de $w$, tel que
$\kappa(u_i)/\kappa(w)$ soit purement inséparable, voir le
lemme \ref{more-morphisms-lemma-etale-splits-off-quasi-finite-part-technical-variant}.
Il suffit clairement de montrer que $\int_{U_i \to W} w_{g \circ f}|_{U_i}$
est constante dans un voisinage \'etale de $w$.
On est ainsi ramené au cas traité dans le paragraphe suivant.

\medskip\noindent
On a $w \in W$ et il existe un unique point $u \in U$
au-dessus de $w$ tel que $\kappa(u)/\kappa(w)$ soit purement inséparable.
Considérons le point $v = f(u) \in Y$. Après avoir remplacé
$(W, w)$ par un voisinage \'etale élémentaire, on peut
supposer qu'il existe un voisinage ouvert $V \subset Y_W$
de $v$ tel que $V \to W$ soit fini, voir le
lemme \ref{more-morphisms-lemma-etale-makes-quasi-finite-finite-at-point}.
Alors $f_W^{-1}(V) \cap U$ est un voisinage ouvert de $u$,
où $f_W : X_W \to Y_W$ est le changement de base de $f$ à $W$.
Après avoir rétréci $W$ au sens de Zariski, on peut donc supposer $f_W(U) \subset V$.
On obtient ainsi des morphismes
$$
U \xrightarrow{a} V \xrightarrow{b} W
$$
et $U \to V$ est fini puisque $V \to W$ est séparé (car fini).
Puisque $w_f$ et $w_g$ sont des pondérations de $f$ et $g$,
nous voyons que $\int_a w_f|_U$ est localement constante sur $V$ et que
$\int_b w_g|_V$ est localement constante sur $W$. Ainsi, quitte à rétrécir
$W$ une fois encore, nous pouvons supposer que ces fonctions sont constantes,
disons de valeurs $n$ et $m$. Il s'ensuit immédiatement que
$\int_\pi w_{g \circ f}|_U = \int_{b \circ a} w_{g \circ f}|_U$
est constante de valeur $nm$, comme souhaité.
\end{proof}

\begin{lemma}
\label{more-morphisms-lemma-weighting-universally-open}
Soit $f : X \to Y$ un morphisme localement quasi-fini.
Soit $w : X \to \mathbf{Z}$ une pondération. Si $w(x) > 0$
pour tout $x \in X$, alors $f$ est universellement ouvert.
\end{lemma}

\begin{proof}
Puisque la propriété est préservée par changement de base, voir le
lemme \ref{more-morphisms-lemma-weighting-base-change}, il suffit
de montrer que $f$ est ouvert. Puisque nous pouvons aussi remplacer
$X$ par n'importe quel ouvert de $X$, il suffit de montrer que $f(X)$
est ouvert. Soit $y \in f(X)$. Choisissons $x \in X$ tel que $f(x) = y$.
Il suffit de montrer que $f(X)$ contient un voisinage ouvert
de $y$, et il suffit de le faire après avoir remplacé $Y$ par un
voisinage \'etale de $y$. Par localisation \'etale des
morphismes quasi-finis, voir la section \ref{more-morphisms-section-etale-localization},
nous pouvons supposer
qu'il existe un voisinage ouvert $U \subset X$ de $x$
tel que $\pi = f|_U : U \to Y$ soit fini. Alors
$\int_\pi w|_U$ est localement constante et prend une valeur positive en $y$.
Par conséquent, $\pi(U)$ contient un voisinage ouvert de $y$,
ce qui achève la démonstration.
\end{proof}

\begin{lemma}
\label{more-morphisms-lemma-weighting-flat-quasi-finite}
Soit $f : X \to Y$ un morphisme de schémas. Supposons $f$
localement quasi-fini, localement de présentation finie et plat.
Alors il existe une pondération positive $w : X \to \mathbf{Z}_{> 0}$ de $f$
donnée par la règle qui associe à $x \in X$, situé au-dessus de $y \in Y$,
$$
w(x) =
\text{longueur}_{\mathcal{O}_{X, x}}
(\mathcal{O}_{X, x}/\mathfrak m_y \mathcal{O}_{X, x})
[\kappa(x) : \kappa(y)]_i
$$
où $[\kappa' : \kappa]_i$ est le degré d'inséparabilité
(Corps, définition \ref{fields-definition-insep-degree}).
\end{lemma}

\begin{proof}
Considérons un diagramme comme dans la définition \ref{more-morphisms-definition-weighting}.
Soit $u \in U$, d'images $x, y, v$ dans $X, Y, V$. Nous affirmons alors que
$$
\text{longueur}_{\mathcal{O}_{X, x}}
(\mathcal{O}_{X, x}/\mathfrak m_y \mathcal{O}_{X, x}) =
\text{longueur}_{\mathcal{O}_{U, u}}
(\mathcal{O}_{U, u}/\mathfrak m_v \mathcal{O}_{U, u})
$$
et
$$
[\kappa(x) : \kappa(y)]_i =
[\kappa(u) : \kappa(v)]_i
$$
La première égalité découle du fait que $\mathcal{O}_{X, x} \to \mathcal{O}_{U, u}$
est un homomorphisme local plat tel que
$\mathfrak m_y \mathcal{O}_{U, u} = \mathfrak m_v \mathcal{O}_{U, u}$
et $\mathfrak m_x \mathcal{O}_{U, u} = \mathfrak m_u$
(car $\mathcal{O}_{Y, y} \to \mathcal{O}_{V, v}$ et
$\mathcal{O}_{X, x} \to \mathcal{O}_{U, u}$ sont non ramifiés),
d'où l'égalité par Algèbre, lemme \ref{algebra-lemma-pullback-module}.
La seconde égalité découle du fait que $\kappa(v)/\kappa(y)$ est une extension
séparable finie et que $\kappa(u)$ est un facteur de
$\kappa(x) \otimes_{\kappa(y)} \kappa(v)$, si bien que le degré
d'inséparabilité est inchangé. Cela étant dit, nous voyons que la formation
de la fonction du lemme commute aux changements de base \'etales.
Ceci ramène le problème à la discussion du paragraphe suivant.

\medskip\noindent
Supposons que $f$ soit un morphisme fini, plat et de présentation finie.
Nous devons montrer que $\int_f w$ est localement constante sur $Y$.
En fait, $f$ est fini localement libre
(Morphismes, lemme \ref{morphisms-lemma-finite-flat}),
et nous allons montrer que $\int_f w$ est égale au degré de $f$
(qui est une fonction localement constante sur $Y$). En effet,
pour $y \in Y$, nous voyons que
\begin{align*}
(\textstyle{\int}_f w)(y)
& =
\sum\nolimits_{f(x) = y}
\text{longueur}_{\mathcal{O}_{X, x}}
(\mathcal{O}_{X, x}/\mathfrak m_y \mathcal{O}_{X, x})
[\kappa(x) : \kappa(y)]_i
[\kappa(x) : \kappa(y)]_s \\
& =
\sum\nolimits_{f(x) = y}
\text{longueur}_{\mathcal{O}_{X, x}}
(\mathcal{O}_{X, x}/\mathfrak m_y \mathcal{O}_{X, x})
[\kappa(x) : \kappa(y)] \\
& =
\text{longueur}_{\mathcal{O}_{Y, y}}((f_*\mathcal{O}_X)_y/
\mathfrak m_y (f_*\mathcal{O}_X)_y)
\end{align*}
La dernière égalité résulte d'Algèbre, lemme \ref{algebra-lemma-pushdown-module}.
Le dernier nombre est le rang de $f_*\mathcal{O}_X$ en $y$, comme souhaité.
\end{proof}

\begin{lemma}
\label{more-morphisms-lemma-weighting-quasi-finite-Noetherian}
Soit $f : X \to Y$ un morphisme de schémas. Supposons que
\begin{enumerate}
\item $f$ soit localement quasi-fini, et
\item $Y$ soit géométriquement unibranche et localement noethérien.
\end{enumerate}
Alors il existe une pondération $w : X \to \mathbf{Z}_{\geq 0}$ donnée par
la règle qui associe à $x \in X$, situé au-dessus de $y \in Y$, le
``degré séparable générique''
de $\mathcal{O}_{X, x}^{sh}$ sur $\mathcal{O}_{Y, y}^{sh}$.
\end{lemma}

\begin{proof}
Il résulte d'Algèbre, lemme
\ref{algebra-lemma-quasi-finite-strict-henselization}
que $\mathcal{O}_{Y, y}^{sh} \to \mathcal{O}_{X, x}^{sh}$
est fini. Puisque $Y$ est géométriquement unibranche,
il existe un unique idéal premier minimal $\mathfrak p$ de
$\mathcal{O}_{Y, y}^{sh}$, voir
Compléments d'algèbre, lemme \ref{more-algebra-lemma-geometrically-unibranch}.
Écrivons
$$
(\kappa(\mathfrak p) \otimes_{\mathcal{O}_{Y, y}^{sh}}
\mathcal{O}_{X, x}^{sh})_{red} =
\prod K_i
$$
comme un produit fini de corps.
Nous posons $w(x) = \sum [K_i : \kappa(\mathfrak p)]_s$.

\medskip\noindent
Puisque cette définition est manifestement insensible à la localisation \'etale,
pour montrer que $w$ est une pondération, nous nous ramenons à montrer que, si
$f$ est un morphisme fini, alors $\int_f w$ est localement constante.
Observons que la valeur de $\int_f w$ en un point générique $\eta$
de $Y$ est simplement le nombre de points de la fibre géométrique
$X_{\overline{\eta}}$ de $X \to Y$ au-dessus de $\eta$. De plus, puisque
$Y$ est unibranche, un point $y$ de $Y$ est la spécialisation d'un unique
point générique $\eta$. Il suffit donc de montrer que $(\int_f w)(y)$
est égal au nombre de points de $X_{\overline{\eta}}$.
Après passage à un voisinage affine de $y$, nous pouvons supposer que
$X \to Y$ est donné par un homomorphisme fini d'anneaux $A \to B$. Supposons que
$\mathcal{O}_{Y, y}^{sh}$ soit construit au moyen d'un homomorphisme
$\kappa(y) \to k$ vers un corps algébriquement clos $k$.
Alors
$$
\mathcal{O}_{Y, y}^{sh} \otimes_A B =
\prod\nolimits_{f(x) = y}
\prod\nolimits_{\varphi \in \Mor_{\kappa(y)}(\kappa(x), k)}
\mathcal{O}_{X, x}^{sh}
$$
par Algèbre, lemme \ref{algebra-lemma-finite-over-henselian}
et le lemme utilisé ci-dessus.
Observons que l'idéal premier minimal $\mathfrak p$ de $\mathcal{O}_{Y, y}^{sh}$
s'envoie sur l'idéal premier de $A$ correspondant à $\eta$. Nous voyons donc que
l'égalité voulue est satisfaite, car le nombre de points d'une fibre
géométrique est inchangé par extension du corps de base.
\end{proof}

\begin{lemma}
\label{more-morphisms-lemma-weighting-specialization}
Soit $f : X \to Y$ un morphisme localement quasi-fini de schémas.
Soit $w : X \to \mathbf{Z}$ une pondération de $f$.
Soit $y \leadsto y'$ une spécialisation de points de $Y$.
Supposons $w(x) = 0$ pour tout $x \in X$ tel que $f(x) = y$. Alors
$w(x') = 0$ pour tout $x' \in X$ tel que $f(x') = y'$.
\end{lemma}

\begin{proof}
Soit $x' \in X$ un point s'envoyant sur $y'$. D'après le
lemme \ref{more-morphisms-lemma-etale-makes-quasi-finite-finite-at-point},
nous pouvons choisir un diagramme
$$
\xymatrix{
U \ar[r] \ar[dr]_\pi & X_V \ar[d] \ar[r] & X \ar[d]^f \\
& V \ar[r]^\varphi & Y
}
$$
où $V \to Y$ est \'etale, avec un point $v' \in V$ s'envoyant sur $y'$,
avec $U \subset Y_V$ ouvert, tel que $U \to V$ soit fini, et
tel qu'il existe un unique point $u' \in U$ s'envoyant sur $v'$,
qui s'envoie en outre sur $x'$ dans $X$. Notons que, puisque $V \to Y$ est ouvert,
il existe un point $v \in V$ se spécialisant en $v'$ et s'envoyant sur $y$.
Par la définition \ref{more-morphisms-definition-weighting}, la fonction
$\int_\pi (w \circ h)$ est localement constante, où $h : U \to X$ est
la composée des flèches horizontales supérieures.
Puisque $\int_\pi (w \circ h)$ est nulle en $v$ d'après
notre hypothèse sur $w$, nous en concluons qu'elle est nulle
en $v'$. Cependant, sa valeur en $v'$ est $w(x')[\kappa(u'):\kappa(v')]_s$,
et nous concluons.
\end{proof}




\section{Compléments sur les pondérations}
\label{more-morphisms-section-more-weightings}

\noindent
Nous démontrons quelques propriétés élémentaires supplémentaires des pondérations. Bien
qu'il semble à première vue que les pondérations puissent être très irrégulières, il
s'avère en fait que la condition imposée dans la
définition \ref{more-morphisms-definition-weighting} est assez forte.

\begin{lemma}
\label{more-morphisms-lemma-jumps-w}
Soit $f : X \to Y$ un morphisme localement quasi-fini.
Soit $w : X \to \mathbf{Z}$ une pondération de $f$. Alors
les ensembles de niveau de la fonction $w$ sont localement constructibles dans $X$.
\end{lemma}

\begin{proof}
Dans la démonstration ci-dessous, nous utiliserons les lemmes \ref{more-morphisms-lemma-weighting-open} et
\ref{more-morphisms-lemma-weighting-base-change} sans autre mention.
Nous utiliserons également les propriétés élémentaires des parties constructibles
des schémas et des espaces topologiques, voir
Topologie, section \ref{topology-section-constructible} et
Propriétés, section \ref{properties-section-constructible}.
À l'aide de celles-ci, le lecteur voit que la question est locale sur $X$ et $Y$ ;
les détails sont omis. Nous pouvons donc supposer $X$ et $Y$ affines.
Si nous pouvons trouver un morphisme surjectif $Y' \to Y$
de présentation finie tel que les ensembles de niveau de $w$
aient pour images réciproques des parties localement constructibles de $X' = Y' \times_Y X$,
alors nous concluons par Morphismes, théorème \ref{morphisms-theorem-chevalley}.

\medskip\noindent
Supposons $X$ et $Y$ affines. Nous pouvons choisir une immersion $X \to T$
où $T \to Y$ est fini, voir le
lemme \ref{more-morphisms-lemma-quasi-finite-separated-pass-through-finite}.
Par Morphismes, lemme \ref{morphisms-lemma-massage-finite},
après avoir remplacé $Y$ par $Y'$, fini localement libre et surjectif sur $Y$,
remplacé $X$ par $Y' \times_Y X$ et $T$ par un schéma fini localement
libre sur $Y'$ contenant $Y' \times_Y T$ comme sous-schéma fermé,
nous pouvons supposer que $T$ est fini localement libre sur $Y$,
et qu'il contient des sous-schémas fermés $T_i$ s'envoyant isomorphiquement sur $Y$
tels que $T = \bigcup_{i = 1, \ldots, n} T_i$ (ensemblistement).
Puisque $T_i \subset T$ est une partie fermée constructible (étant l'image
d'un morphisme de présentation finie $Y \to T$), nous voyons
que, pour $I \subset \{1, \ldots, n\}$, l'intersection
$\bigcap_{i \in I} T_i$ est une partie fermée constructible de $T$
et son image est donc une partie fermée constructible de $Y$.

\medskip\noindent
Pour une décomposition en union disjointe
$\{1, \ldots, n\} = I_1 \amalg \ldots \amalg I_r$ à parties non vides,
considérons la partie $Y_{I_1, \ldots, I_r} \subset Y$
formée des points $y \in Y$ tels que $T_y = \{x_1, \ldots, x_r\}$
soit formé d'exactement $r$ points, avec $x_j \in T_i \Leftrightarrow i \in I_j$.
D'après les remarques précédentes, ceci donne une partition constructible de $Y$.
Il existe un schéma affine $Y'$ de présentation finie
sur $Y$ tel que l'image de $Y' \to Y$ soit exactement
$Y_{I_1, \ldots, I_r}$, voir Algèbre, lemme
\ref{algebra-lemma-constructible-is-image}.
Nous pouvons donc supposer que $Y = Y_{I_1, \ldots, I_r}$ pour une
décomposition en union disjointe
$\{1, \ldots, n\} = I_1 \amalg \ldots \amalg I_r$.
Dans ce cas, l'égalité $T = T(1) \amalg \ldots \amalg T(r)$, où
$T(j) = \bigcap_{i \in I_j} T_i$,
donne une décomposition de $T$ en parties fermées disjointes (et donc ouvertes).
En intersectant avec le sous-schéma localement fermé $X$, nous obtenons une
décomposition analogue $X = X(1) \amalg \ldots \amalg X(r)$ en parties ouvertes et fermées.
Le morphisme $X(j) \to Y$ est une immersion.
Puisque $w$ est une pondération, il s'ensuit que $w|_{X(j)}$
est localement constante\footnote{En fait, si $f : X \to Y$ est une
immersion et si $w$ est une pondération de $f$, alors la restriction de $f$ est une
application ouverte sur le lieu où $w$ est non nulle.}, et nous concluons.
\end{proof}

\begin{lemma}
\label{more-morphisms-lemma-weighting-on-dense-set}
Soit $f : X \to Y$ un morphisme localement quasi-fini de schémas.
Soit $w : X \to \mathbf{Z}$ une pondération de $f$.
Soit $E \subset X$ une partie dense.
Supposons $w(x) = 0$ pour tout $x \in E$. Alors $w = 0$.
\end{lemma}

\begin{proof}
Nous conseillons de lire d'abord le lemme \ref{more-morphisms-lemma-weighting-specialization}.
Soit $x \in X$, d'image $y \in Y$. Nous allons montrer que $w(x) = 0$. D'après le
lemme \ref{more-morphisms-lemma-etale-makes-quasi-finite-finite-at-point},
nous pouvons choisir un diagramme
$$
\xymatrix{
U \ar[r] \ar[dr]_\pi & X_V \ar[d] \ar[r] & X \ar[d]^f \\
& V \ar[r]^\varphi & Y
}
$$
où $V \to Y$ est \'etale, avec un point $v \in V$ s'envoyant sur $y$,
avec $U \subset Y_V$ ouvert, tel que $U \to V$ soit fini, et
tel qu'il existe un unique point $u \in U$ s'envoyant sur $v$,
qui s'envoie en outre sur $x$ dans $X$. Nous pouvons supposer que $V$
(et donc $U$) est affine. Par la définition \ref{more-morphisms-definition-weighting},
la fonction $\int_\pi (w \circ h)$ est localement constante, où
$h : U \to X$ est la composée des flèches horizontales supérieures.
L'ensemble $U_0$ des points où $w \circ h$ vaut $0$ est une
partie constructible de $U$ par le lemme \ref{more-morphisms-lemma-jumps-w},
et elle est dense dans $U$ puisqu'elle contient $h^{-1}(E)$ (utiliser le fait que $h$ est ouvert).
Par Topologie, lemme \ref{topology-lemma-dense-in-constructible},
nous voyons que $U_0$ contient un ouvert dense,
c'est-à-dire que l'ensemble des points où $w \circ h$ est non nul est contenu
dans une partie fermée nulle part dense $Z \subset U$. Alors $\pi(Z)$
est une partie fermée nulle part dense de $V$ par
Morphismes, lemme \ref{morphisms-lemma-image-nowhere-dense-finite}.
Ainsi, $\int_\pi (w \circ h)$ est nulle sur un ouvert dense, donc
s'annule partout.
Cependant, sa valeur en $v$ est $w(x)[\kappa(u):\kappa(v)]_s$,
et nous concluons.
\end{proof}

\begin{lemma}
\label{more-morphisms-lemma-jumps-int-w}
Soit $f : X \to Y$ un morphisme localement quasi-fini de présentation
finie. Soit $w : X \to \mathbf{Z}$ une pondération de $f$. Alors
les ensembles de niveau de la fonction $\int_f w$ sont localement constructibles dans $Y$.
\end{lemma}

\begin{proof}
Par le lemme \ref{more-morphisms-lemma-weighting-check-after-etale-base-change},
la formation de la fonction $\int_f w$ commute à tout
changement de base et, par le lemme \ref{more-morphisms-lemma-weighting-base-change},
après changement de base, nous avons encore une pondération.
Cela signifie que, si nous pouvons trouver $Y' \to Y$
surjectif et de présentation finie, alors il
suffit de démontrer le résultat après changement de base à $Y'$, voir
Morphismes, théorème \ref{morphisms-theorem-chevalley}.

\medskip\noindent
La question est locale sur $Y$ ; nous pouvons donc supposer $Y$ affine.
Alors $X$ est quasi-compact et quasi-séparé
(puisque $f$ est de présentation finie). Supposons que $X = U \cup V$
soit une réunion d'ouverts quasi-compacts. Alors nous avons
$$
\textstyle{\int}_f w =
\textstyle{\int}_{f|_U} w|_U +
\textstyle{\int}_{f|_V} w|_V -
\textstyle{\int}_{f|_{U \cap V}} w|_{U \cap V}
$$
Ainsi, si nous connaissons le résultat pour $w|_U$, $w|_V$, $w|_{U \cap V}$,
alors nous le connaissons pour $w$. Par le principe d'induction
(Cohomologie des schémas, lemme \ref{coherent-lemma-induction-principle}),
il suffit de démontrer le lemme lorsque $X$ est affine.

\medskip\noindent
Supposons $X$ et $Y$ affines. Nous pouvons choisir une immersion ouverte
$X \to T$ où $T \to Y$ est fini, voir le
lemme \ref{more-morphisms-lemma-quasi-finite-separated-pass-through-finite}.
Puisque nous pouvons encore effectuer un changement de base par un $Y' \to Y$ convenable,
nous pouvons utiliser Morphismes, lemme \ref{morphisms-lemma-massage-finite},
pour nous ramener au cas où
toutes les extensions de corps résiduels induites par le morphisme $T \to Y$
(et a fortiori par $X \to Y$) sont triviales.
Dans cette situation, $\int_f w$ consiste simplement à prendre les sommes
des valeurs de $w$ dans les fibres. Les ensembles de niveau de $w$
sont localement constructibles dans $X$ (lemme \ref{more-morphisms-lemma-jumps-w}).
La fonction $w$ ne prend qu'un nombre fini de valeurs d'après
Propriétés, lemme
\ref{properties-lemma-stratification-locally-finite-constructible}.
Nous concluons donc par
Morphismes, théorème \ref{morphisms-theorem-chevalley},
et quelques arguments élémentaires sur les sommes d'entiers.
\end{proof}

\begin{lemma}
\label{more-morphisms-lemma-semicontinuous-w}
Soit $f : X \to Y$ un morphisme localement quasi-fini.
Soit $w : X \to \mathbf{Z}_{> 0}$ une pondération positive de $f$.
Alors $w$ est semi-continue supérieurement.
\end{lemma}

\begin{proof}
Soit $x \in X$, d'image $y \in Y$. Choisissons un voisinage \'etale
$(V, v) \to (Y, y)$ et un ouvert $U \subset X_V$ tels que
$\pi : U \to V$ soit fini et qu'il existe un unique point $u \in U$
s'envoyant sur $v$, avec $\kappa(u)/\kappa(v)$ purement inséparable.
Voir le lemme \ref{more-morphisms-lemma-etale-makes-quasi-finite-finite-multiple-points-var}.
Alors $(\int_\pi w|_U)(v) = w(u)$.
Il résulte de la définition \ref{more-morphisms-definition-weighting}
que, quitte à remplacer $V$ par un voisinage de $v$, nous
avons $w|_U(u') \leq w|_U(u) = w(x)$ pour tout $u' \in U$.
En effet, $w|_U(u')$ intervient comme un terme de la somme dans l'expression
de $(\int_\pi w|_U)(\pi(u'))$.
Ceci démontre le lemme, car le morphisme \'etale
$U \to X$ est ouvert.
\end{proof}

\begin{lemma}
\label{more-morphisms-lemma-semicontinuous-int-w}
Soit $f : X \to Y$ un morphisme séparé et localement quasi-fini
à fibres finies.
Soit $w : X \to \mathbf{Z}_{> 0}$ une pondération positive de $f$.
Alors $\int_f w$ est semi-continue inférieurement.
\end{lemma}

\begin{proof}
Soit $y \in Y$. Soient $x_1, \ldots, x_r \in X$ les points situés au-dessus de $y$.
Appliquons le
lemme \ref{more-morphisms-lemma-etale-splits-off-quasi-finite-part-technical-variant}
pour obtenir un voisinage \'etale $(U, u) \to (Y, y)$ et une décomposition
$$
U \times_Y X =
W \amalg
\ \coprod\nolimits_{i = 1, \ldots, n}
\ \coprod\nolimits_{j = 1, \ldots, m_i}
V_{i, j}
$$
comme dans le passage cité. Observons que $(\int_f w)(y) = \sum w(v_{i, j})$,
où $w(v_{i, j}) = w(x_i)$. Puisque $\int_{V_{i, j} \to U} w|_{V_{i, j}}$
est localement constante par définition, nous pouvons, quitte à rétrécir $U$,
supposer que ces fonctions sont constantes de valeur $w(v_{i, j})$.
Nous en concluons que
$$
\textstyle{\int}_{U \times_Y X \to U} w|_{U \times_Y X} =
\textstyle{\int}_{W \to U} w|_W +
\sum \textstyle{\int}_{V_{i, j} \to U} w|_{V_{i, j}} =
\textstyle{\int}_{W \to U} w|_W + (\int_f w)(y)
$$
Ceci est $\geq (\int_f w)(y)$, et nous concluons puisque $U \to Y$
est ouvert et que la formation de l'intégrale commute aux changements de base
(lemme \ref{more-morphisms-lemma-weighting-check-after-etale-base-change}).
\end{proof}

\begin{lemma}
\label{more-morphisms-lemma-max-int-w}
Soit $f : X \to Y$ un morphisme localement quasi-fini,
avec $X$ quasi-compact. Soit $w : X \to \mathbf{Z}$ une pondération de $f$.
Alors $\int_f w$ atteint son maximum.
\end{lemma}

\begin{proof}
Il résulte du lemme \ref{more-morphisms-lemma-jumps-w} et de
Propriétés, lemme
\ref{properties-lemma-stratification-locally-finite-constructible}
que $w$ ne prend qu'un nombre fini de valeurs sur $X$.
Il résulte de Morphismes, lemme
\ref{morphisms-lemma-locally-quasi-finite-qc-source-universally-bounded}
que $X \to Y$ a des fibres géométriques de cardinal borné.
Ceci montre que $\int_f w$ est bornée.
\end{proof}

\begin{lemma}
\label{more-morphisms-lemma-max-int-finite}
Soit $f : X \to Y$ un morphisme séparé et localement quasi-fini.
Soit $w : X \to \mathbf{Z}_{> 0}$ une pondération positive de $f$.
Supposons que $\int_w f$ atteigne son maximum $d$, et soit $Y_d \subset Y$
l'ouvert des points $y$ tels que $(\int_f w)(y) = d$. Alors
le morphisme $f^{-1}(Y_d) \to Y_d$ est fini.
\end{lemma}

\begin{proof}
Observons que $Y_d$ est ouvert par le lemme \ref{more-morphisms-lemma-semicontinuous-int-w}.
Soit $y \in Y_d$. Soient $x_1, \ldots, x_n$ les points de $X$
situés au-dessus de $y$. Appliquons le
lemme \ref{more-morphisms-lemma-etale-splits-off-quasi-finite-part-technical-variant}
pour obtenir un voisinage \'etale $(U, u) \to (Y, y)$ et une décomposition
$$
U \times_Y X =
W \amalg
\ \coprod\nolimits_{i = 1, \ldots, n}
\ \coprod\nolimits_{j = 1, \ldots, m_i}
V_{i, j}
$$
comme dans le passage cité. Observons que $d = \sum w(v_{i, j})$, où
$w(v_{i, j}) = w(x_i)$. Puisque $\int_{V_{i, j} \to U} w|_{V_{i, j}}$
est localement constante par définition, nous pouvons, quitte à rétrécir $U$,
supposer que ces fonctions sont constantes de valeur $w(v_{i, j})$.
Nous en concluons que
$$
\textstyle{\int}_{U \times_Y X \to U} w|_{U \times_Y X} =
\textstyle{\int}_{W \to U} w|_W +
\sum \textstyle{\int}_{V_{i, j} \to U} w|_{V_{i, j}} =
\textstyle{\int}_{W \to U} w|_W + (\int_f w)(y)
$$
Ceci est $\geq (\int_f w)(y) = d$, et nous concluons que $W$
doit être vide. Ainsi, $U \times_Y X \to U$ est fini.
Par Descente, lemme \ref{descent-lemma-descending-property-finite},
ceci implique que $X \to Y$ est fini au-dessus de
l'image du morphisme ouvert $U \to Y$. Autrement dit,
nous voyons que $f$ est fini au-dessus d'un voisinage ouvert de $y$,
comme souhaité.
\end{proof}

\begin{lemma}
\label{more-morphisms-lemma-open-and-closed-in-finite}
Soit $A \to B$ un homomorphisme d'anneaux fini et de présentation finie.
Il existe un homomorphisme d'anneaux de présentation finie $A \to A_{univ}$
et un idempotent $e_{univ} \in B \otimes_A A_{univ}$
tels que, pour tout homomorphisme d'anneaux $A \to A'$ et tout idempotent $e \in B \otimes_A A'$,
il existe un homomorphisme d'anneaux $A_{univ} \to A'$ envoyant $e_{univ}$ sur $e$.
\end{lemma}

\begin{proof}
Choisissons $b_1, \ldots, b_n \in B$ engendrant $B$ comme $A$-module.
Pour chaque $i$, choisissons un polynôme unitaire $P_i \in A[x]$ tel que $P_i(b_i) = 0$
dans $B$, voir Algèbre, lemme \ref{algebra-lemma-finite-is-integral}.
Ainsi, $B$ est un quotient de la $A$-algèbre libre de rang fini
$B' = A[x_1, \ldots, x_n]/(P_1(x_1), \ldots, P_n(x_n))$.
Soit $J \subset B'$ le noyau de la surjection $B' \to B$.
Alors $J =(f_1, \ldots, f_m)$ est engendré par un nombre fini d'éléments, puisque $B$
est une $A$-algèbre de type fini, voir
Algèbre, lemme \ref{algebra-lemma-compose-finite-type}.
Choisissons, sur $A$, une base $b'_1, \ldots, b'_N$ de $B'$.
Considérons l'algèbre
$$
A_{univ} = A[z_1, \ldots, z_N, y_1, \ldots, y_m]/I
$$
où $I$ est l'idéal engendré par les coefficients dans
$A[z_1, \ldots, z_n, y_1, \ldots, y_m]$
de l'expression suivante relativement à la base $b'_1, \ldots, b'_N$ :
$$
(\sum z_j b'_j)^2 - \sum z_j b'_j + \sum y_k f_k
$$
dans $B'[z_1, \ldots, z_N, y_1, \ldots, y_m]$. Par construction,
l'élément $\sum z_j b'_j$ s'envoie sur un idempotent $e_{univ}$ de
l'algèbre $B \otimes_A A_{univ}$. De plus, si $e \in B \otimes_A A'$
est un idempotent, nous pouvons relever $e$ en un élément de la forme
$\sum b'_j \otimes a'_j$ dans $B' \otimes_A A'$, et nous pouvons trouver
des $a''_k \in A'$ tels que
$$
(\sum b'_j \otimes a'_j)^2 - \sum b'_j \otimes a'_j + \sum f_k \otimes a''_k
$$
soit nul dans $B' \otimes_A A'$. Nous obtenons donc un homomorphisme de $A$-algèbres
$A_{univ} \to A$ envoyant $z_j$ sur $a'_j$ et $y_k$ sur $a''_k$,
et $e_{univ}$ sur $e$. Ceci achève la démonstration.
\end{proof}

\begin{lemma}
\label{more-morphisms-lemma-open-and-closed-in-quasi-finite}
Soit $X \to Y$ un morphisme de schémas affines quasi-fini et
de présentation finie. Il existe un morphisme $Y_{univ} \to Y$
de présentation finie et un sous-schéma ouvert
$U_{univ} \subset Y_{univ} \times_Y X$ tel que
$U_{univ} \to Y_{univ}$ soit fini et vérifie la propriété suivante :
pour tout morphisme $Y' \to Y$ de schémas affines
et tout sous-schéma ouvert $U' \subset Y' \times_Y X$
tel que $U' \to Y'$ soit fini, il existe un morphisme
$Y' \to Y_{univ}$ tel que l'image réciproque de $U_{univ}$ soit $U'$.
\end{lemma}

\begin{proof}
Rappelons qu'un morphisme de type fini est quasi-fini si et seulement
s'il est de dimension relative $0$, voir
Morphismes, lemme \ref{morphisms-lemma-locally-quasi-finite-rel-dimension-0}.
Par le lemme \ref{more-morphisms-lemma-Noetherian-approximation-dimension-d},
appliqué avec $d = 0$, nous nous ramenons au cas où
$X$ et $Y$ sont noethériens. Nous pouvons choisir une immersion ouverte
$X \to X'$ telle que $X' \to Y$ soit fini, voir
Algèbre, lemme \ref{algebra-lemma-quasi-finite-open-integral-closure}.
Notons que, si nous avons $Y' \to Y$ et $U'$ comme en (2), alors
$$
U' \to Y' \times_Y X \to Y' \times_Y X'
$$
est une immersion ouverte entre des schémas finis sur $Y'$, et est donc
aussi fermée. Nous en concluons que $U'$ correspond à un
idempotent dans
$$
\Gamma(Y', \mathcal{O}_{Y'})
\otimes_{\Gamma(Y, \mathcal{O}_Y)}
\Gamma(X', \mathcal{O}_{X'})
$$
dont la partie ouverte et fermée correspondante est contenue dans
l'ouvert $Y' \times_Y X$. Soient $Y'_{univ} \to Y$ et l'idempotent
$$
e'_{univ} \in
\Gamma(Y_{univ}, \mathcal{O}_{Y_{univ}})
\otimes_{\Gamma(Y, \mathcal{O}_Y)}
\Gamma(X', \mathcal{O}_{X'})
$$
le couple construit dans le lemme \ref{more-morphisms-lemma-open-and-closed-in-finite}
pour l'homomorphisme d'anneaux $\Gamma(Y, \mathcal{O}_Y) \to \Gamma(X', \mathcal{O}_{X'})$
(ici, nous utilisons le fait que $Y$ est noethérien pour voir que $X'$ est de présentation finie
sur $Y$). Soit $U'_{univ} \subset Y'_{univ} \times_Y X'$ le sous-schéma
ouvert et fermé correspondant. Nous voyons alors que
$$
U'_{univ} \setminus Y'_{univ} \times_Y X
$$
est une partie fermée de $U'_{univ}$, et a donc une image fermée
$T \subset Y'_{univ}$. Si nous posons $Y_{univ} = Y'_{univ} \setminus T$
et si $U_{univ}$ est la restriction de $U'_{univ}$ à
$Y_{univ} \times_Y X$, alors nous voyons que le lemme est vrai.
\end{proof}

\begin{lemma}
\label{more-morphisms-lemma-descend-weighting}
Soit $Y = \lim Y_i$ une limite cofiltrante de schémas affines. Soit $0 \in I$,
et soit $f_0 : X_0 \to Y_0$ un morphisme de schémas affines quasi-fini
et de présentation finie. Soient $f : X  \to Y$
et $f_i : X_i \to Y_i$, pour $i \geq 0$, les changements de base de $f_0$.
Si $w : X \to \mathbf{Z}$ est une pondération de $f$, alors, pour tout indice
$i$ assez grand, il existe une pondération $w_i : X_i \to \mathbf{Z}$
de $f_i$ dont le tiré en arrière sur $X$ est $w$.
\end{lemma}

\begin{proof}
D'après le lemme \ref{more-morphisms-lemma-jumps-w}, les ensembles de niveau de $w$ sont les sous-ensembles constructibles
$E_k$ de $X$. Il en résulte que la fonction $w$ ne prend qu'un nombre fini
de valeurs, d'après Propriétés, lemme
\ref{properties-lemma-stratification-locally-finite-constructible}.
Il existe donc un $i$ tel que $E_k$ descende en un sous-ensemble
constructible $E_{i, k}$ de $X_i$ pour tout $k$ ; de plus, nous pouvons supposer que
$X_i = \coprod E_{i, k}$. Cela résulte du fait que l'espace topologique
de $X$ est la limite, dans la catégorie des espaces topologiques,
des espaces spectraux $X_i$ le long d'un système dirigé dont les
applications de transition sont spectrales. Voir
Limites, section \ref{limits-section-descent}
et
Topologie, section \ref{topology-section-limits-spectral}.
Nous définissons $w_i : X_i \to \mathbf{Z}$ en demandant que ses ensembles
de niveau soient les ensembles constructibles $E_{i, k}$.

\medskip\noindent
Choisissons $Y_{i, univ} \to Y_i$ et
$U_{i, univ} \subset Y_{i, univ} \times_{Y_i} X_i$
comme dans le lemme \ref{more-morphisms-lemma-open-and-closed-in-quasi-finite}.
D'après la propriété universelle de la construction, pour
montrer que $w_i$ est une pondération, il suffirait de montrer
que
$$
\tau_i = \textstyle{\int}_{U_{i, univ} \to Y_{i, univ}} w_i|_{U_{i, univ}}
$$
est localement constante sur $Y_{i, univ}$. D'après le lemme \ref{more-morphisms-lemma-jumps-int-w},
cette fonction a des ensembles de niveau constructibles, mais elle peut
ne pas être (encore) localement constante. Posons
$Y_{univ} = Y_{i, univ} \times_{Y_i} Y$
et soit $U_{univ} \subset Y_{univ} \times_Y X$
l'image réciproque de $U_{i, univ}$.
Alors, puisque le tiré en arrière de $w$ sur $Y_{univ} \times_Y X$
est une pondération de $Y_{univ} \times_Y X \to Y_{univ}$
(lemme \ref{more-morphisms-lemma-weighting-base-change}),
nous avons bien que
$$
\tau = \textstyle{\int}_{U_{univ} \to Y_{univ}} w_i|_{U_{univ}}
$$
est localement constante sur $Y_{univ}$. Les ensembles de niveau de
$\tau$ sont donc ouverts et fermés. Enfin, nous avons
$Y_{univ} = \lim_{i' \geq i} Y_{i', univ}$,
et les ensembles de niveau de $\tau$ sont les limites projectives des
ensembles de niveau de $\tau_{i'}$ (définis de manière analogue).
Les références ci-dessus impliquent donc que, pour un indice
$i'$ assez grand, les ensembles de niveau de $\tau_{i'}$ sont eux aussi ouverts.
Pour un tel indice $i'$, nous concluons que $w_{i'}$
est une pondération de $f_{i'}$, comme voulu.
\end{proof}







\section{Pondérations et nombres de stratification affine}
\label{more-morphisms-section-bounds-asn}

\noindent
Dans cette section, nous donnons une borne du nombre de stratification affine d'un
schéma qui possède un certain type de revêtement par un schéma affine.

\begin{lemma}
\label{more-morphisms-lemma-affineness-of-large-open}
Soit $f : X \to Y$ un morphisme de schémas affines qui est
quasi-fini et de présentation finie.
Soit $w : X \to \mathbf{Z}_{> 0}$ une pondération positive de $f$.
Soit $d < \infty$ la valeur maximale de $\int_f w$. L'ouvert
$$
Y_d = \{y \in Y \mid (\textstyle{\int}_f w)(y) = d \}
$$
de $Y$ est affine.
\end{lemma}

\begin{proof}
Observons que $\int_f w$ atteint son maximum d'après le lemme \ref{more-morphisms-lemma-max-int-w}.
L'ensemble $Y_d$ est ouvert d'après le lemme \ref{more-morphisms-lemma-semicontinuous-int-w}.
L'énoncé du lemme a donc bien un sens.

\medskip\noindent
Réduction au cas noethérien ; veuillez omettre ce paragraphe.
Rappelons qu'un morphisme de type fini est quasi-fini si et seulement
s'il est de dimension relative $0$, voir
Morphismes, lemme \ref{morphisms-lemma-locally-quasi-finite-rel-dimension-0}.
En appliquant le lemme \ref{more-morphisms-lemma-Noetherian-approximation-dimension-d}
avec $d = 0$, nous pouvons trouver un morphisme quasi-fini $f_0 : X_0 \to Y_0$
de schémas affines noethériens et un morphisme $Y \to Y_0$ tels que $f$
soit le changement de base de $f_0$. Nous pouvons alors écrire $Y = \lim Y_i$ comme limite
d'un système dirigé de schémas affines de type fini sur $Y_0$, voir
Algèbre, lemme \ref{algebra-lemma-ring-colimit-fp}.
D'après le lemme \ref{more-morphisms-lemma-descend-weighting},
nous pouvons trouver un $i$ tel que notre pondération $w$
descende en une pondération $w_i$ du changement de base $f_i : X_i \to Y_i$
de $f_0$. Si le lemme vaut maintenant pour $f_i, w_i$, il implique
alors le lemme pour $f$, puisque la formation de $\int_f w$ commute au
changement de base, voir le lemme \ref{more-morphisms-lemma-weighting-check-after-etale-base-change}.

\medskip\noindent
Supposons $X$ et $Y$ noethériens. Soit $X' \to Y'$ le changement de base de $f$
par un morphisme $g : Y' \to Y$. La formation de $\int_f w$, et par suite
celle de l'ouvert $Y_d$, commute au changement de base. Si $g$ est fini et surjectif, alors
$Y'_d \to Y_d$ est fini et surjectif. Dans ce cas, montrer que
$Y_d$ est affine équivaut à montrer que $Y'_d$ est affine, voir
Cohomologie des schémas, lemme
\ref{coherent-lemma-image-affine-finite-morphism-affine-Noetherian}.

\medskip\noindent
Nous pouvons choisir une immersion $X \to T$ avec $T$ fini sur $Y$, voir le lemme
\ref{more-morphisms-lemma-quasi-finite-separated-pass-through-finite}.
Nous allons appliquer Morphismes, lemme \ref{morphisms-lemma-massage-finite}
au morphisme fini $T \to Y$. Ce lemme nous dit qu'il
existe un morphisme fini surjectif $Y' \to Y$ tel que
$Y' \times_Y T$ soit un sous-schéma fermé d'un schéma $T'$ fini sur $Y'$
qui possède une forme particulière.
D'après la discussion du premier paragraphe, nous pouvons remplacer $Y$ par $Y'$,
$T$ par $T'$, et $X$ par $Y' \times_Y X$.
Nous pouvons donc supposer qu'il existe une immersion $X \to T$ (qui n'est pas nécessairement
ouverte ou fermée) et des sous-schémas fermés
$T_i \subset T$, $i = 1, \ldots, n$ tels que
\begin{enumerate}
\item $T \to Y$ est fini (et localement libre),
\item $T_i \to Y$ est un isomorphisme, et
\item $T = \bigcup_{i = 1, \ldots, n} T_i$ ensemblistement.
\end{enumerate}
Soit $Y' = \coprod Y_k$ la réunion disjointe des composantes irréductibles
de $Y$ (considérées comme des sous-schémas fermés intègres de $Y$).
Nous pouvons alors effectuer encore une fois le changement de base par $Y' \to Y$ ; nous
utilisons ici le fait que $Y$ est noethérien. Nous pouvons ainsi supposer en
outre que $Y$ est intègre et noethérien.

\medskip\noindent
Nous pouvons et voulons aussi supposer que $T_i \not = T_j$ si $i \not = j$, en
supprimant les répétitions. Puisque $Y$, et donc tous les $T_i$, sont intègres, cela
signifie que, si $T_i$ et $T_j$ se rencontrent, leur intersection est un
sous-ensemble fermé dont l'image est un sous-ensemble fermé propre de $Y$.

\medskip\noindent
Observons que $V_i = X \cap T_i$ est un sous-ensemble localement fermé
qui est en outre un sous-schéma fermé de $X$, donc affine.
Soient $\eta \in Y$ et $\eta_i \in T_i$ les points génériques.
Si $\eta \not \in Y_d$, alors $Y_d = \emptyset$ et nous avons terminé.
Supposons $\eta \in Y_d$. Notons $I \in \{1, \ldots, n\}$
le sous-ensemble des indices $i$ tels que $\eta_i \in V_i$.
Pour $i \in I$, le sous-ensemble localement fermé $V_i \subset T_i$
contient le point générique de l'espace irréductible $T_i$
et est donc ouvert. D'autre part, puisque $f$ est ouvert
(lemme \ref{more-morphisms-lemma-weighting-universally-open}),
pour tout $x \in X$, nous pouvons trouver un $i \in I$ et une spécialisation
$\eta_i \leadsto x$. Il s'ensuit que $x \in T_i$, et donc
$x \in V_i$. Autrement dit, nous voyons que $X = \bigcup_{i \in I} V_i$
ensemblistement.
Nous affirmons que $Y_d = \bigcap_{i \in I} \Im(V_i \to Y)$ ; cela
achèvera la preuve, puisque l'intersection des ouverts affines
$\Im(V_i \to Y)$ de $Y$ est affine.

\medskip\noindent
Pour $y \in Y$, écrivons $f^{-1}(\{y\}) = \{x_1, \ldots, x_r\}$ dans $X$.
Pour chaque $i \in I$, il existe au plus un $j(i) \in \{1, \ldots, x_r\}$
tel que $\eta_i \leadsto x_{j(i)}$. En fait, $j(i)$ existe et est
égal à $j$ si et seulement si $x_j \in V_i$. Si $i \in I$ est tel que
$j = j(i)$ soit défini, alors $V_i \to Y$ est un isomorphisme au voisinage
de $x_j \mapsto y$. Par suite, $\bigcup_{i \in I,\ j(i) = j} V_i \to Y$
est fini après remplacement de la source et du but par des voisinages respectifs de
$x_j \mapsto y$. La définition d'une pondération nous donne donc
$w(x_j) = \sum_{i \in I,\ j(i) = j} w(\eta_i)$.
Nous voyons ainsi que
$$
(\textstyle{\int}_f w)(\eta) =
\sum\nolimits_{i \in I} w(\eta_i) \geq
\sum\nolimits_{j(i)\text{ existe}} w(\eta_i) =
\sum\nolimits_j w(x_j) = (\textstyle{\int}_f w)(y)
$$
Il y a donc égalité si et seulement si $y$ appartient à
$\bigcap_{i \in I} \Im(V_i \to Y)$, ce qu'il fallait montrer.
\end{proof}

\begin{proposition}
\label{more-morphisms-proposition-asn-weighting}
Soit $f : X \to Y$ un morphisme surjectif quasi-fini de schémas.
Soit $w : X \to \mathbf{Z}_{> 0}$ une pondération positive de $f$.
Supposons $X$ affine et $Y$ séparé\footnote{Il suffit que la
diagonale de $Y$ soit affine.} et non vide. Alors le nombre de stratification
affine de $Y$ est au plus égal au nombre de valeurs distinctes de $\int_f w$
moins $1$.
\end{proposition}

\begin{proof}
Notons que, puisque $Y$ est séparé, le morphisme $X \to Y$ est affine
(Morphismes, lemme \ref{morphisms-lemma-affine-permanence}).
La fonction $\int_f w$ atteint sa valeur maximale $d$ d'après le
lemme \ref{more-morphisms-lemma-max-int-w}. Nous allons raisonner par récurrence sur $d$.
Considérons le sous-schéma ouvert $Y_d = \{y \in Y \mid (\int_f w)(y) = d\}$
de $Y$ et rappelons que $f^{-1}(Y_d) \to Y_d$ est fini, voir le
lemme \ref{more-morphisms-lemma-max-int-finite}.
D'après le lemme \ref{more-morphisms-lemma-affineness-of-large-open},
pour tout ouvert affine $W \subset Y$, l'intersection $Y_d \cap W$ est affine
(on utilise ici que $W \times_Y X$ est affine, puisqu'il est affine sur $X$).
Ainsi $Y_d \to Y$ est un morphisme affine de schémas. Nous
en concluons que $f^{-1}(Y_d) = Y_d \times_Y X$ est
un schéma affine, puisqu'il est affine sur $X$.
Alors $f^{-1}(Y_d) \to Y_d$ est surjectif, et
$Y_d$ est donc affine d'après Limites, lemme \ref{limits-lemma-affine}.
Si $Y_d = Y$, alors le nombre de stratification affine de $Y$ vaut $0$
et nous concluons que le lemme est vrai. Sinon, posons
$Y' = Y \setminus Y_d$, considéré comme un sous-schéma fermé réduit de $Y$.
Posons $X' = Y' \times_Y X$.
Puisque $X'$ est fermé dans $X$, il est affine. Puisque
$Y'$ est fermé dans $Y$, il est séparé.
Le morphisme $f' : X' \to Y'$ est surjectif, et
$w$ induit une pondération $w'$ de $f'$ (voir
le lemme \ref{more-morphisms-lemma-weighting-base-change}) tandis que
$\int_{f'} w'$ est la restriction de $\int_f w$ à $Y'$.
Par récurrence, $Y'$ possède une stratification affine dont la
longueur est au plus le nombre ($\leq$) de valeurs distinctes de
$\int_{f'} w'$ moins $1$, ce qui achève la preuve.
\end{proof}









\section{Morphismes complètement décomposés}
\label{more-morphisms-section-completely-decomposed}

\noindent
Nishnevich a étudié la notion de famille complètement décomposée
de morphismes \'etales, afin de définir ce que l'on appelle aujourd'hui la
topologie de Nisnevich ; voir par exemple \cite{Nishnevich}.

\begin{definition}
\label{more-morphisms-definition-cd-morphism}
Un morphisme $f : X \to Y$ de schémas est dit
{\it complètement décomposé}\footnote{Cette terminologie n'est peut-être pas standard.}
si, pour tout point $y \in Y$, il
existe un point $x \in X$ tel que $f(x) = y$ et que l'extension de corps
$\kappa(x)/\kappa(y)$ soit triviale.
Une famille de morphismes $\{f_i : X_i \to Y\}_{i \in I}$ de
schémas à but fixé est dite {\it complètement décomposée}
si $\coprod f_i : \coprod Y_i \to X$ est complètement décomposé.
\end{definition}

\noindent
Nous commençons par quelques lemmes élémentaires.

\begin{lemma}
\label{more-morphisms-lemma-composition-cd}
Le composé de deux morphismes complètement décomposés de schémas
est complètement décomposé.
Si $\{f_i : X_i \to Y\}_{i \in I}$ est complètement décomposée
et si, pour chaque $i$, nous avons une famille $\{X_{ij} \to X_i\}_{j \in J_i}$
qui est complètement décomposée, alors la famille
$\{X_{ij} \to Y\}_{i \in I, j \in J_i}$ est complètement décomposée.
\end{lemma}

\begin{proof}
Omis.
\end{proof}

\begin{lemma}
\label{more-morphisms-lemma-base-change-cd}
Tout changement de base d'un morphisme complètement décomposé de schémas
est complètement décomposé.
Si $\{f_i : X_i \to Y\}_{i \in I}$ est complètement décomposée
et si $Y' \to Y$ est un morphisme de schémas, alors
$\{X_i \times_Y Y' \to Y'\}_{i \in I}$ est complètement
décomposée.
\end{lemma}

\begin{proof}
Soient $f : X \to Y$ et $g : Y' \to Y$ des morphismes de schémas.
Soit $y' \in Y'$ un point d'image $y = g(y')$ dans $Y$.
Si $x \in X$ est un point tel que $f(x) = y$ et $\kappa(x) = \kappa(y)$,
alors il existe un unique point $x' \in X' = X \times_Y Y'$
qui s'envoie sur $y'$ dans $Y'$ et sur $x$ dans $X$, et de plus
$\kappa(x') = \kappa(y')$, voir
Schémas, lemme \ref{schemes-lemma-points-fibre-product}.
Le lemme découle aisément de ce fait ; nous omettons les détails.
\end{proof}

\begin{lemma}
\label{more-morphisms-lemma-decompose}
\begin{reference}
\cite[Lemme 2.1.2]{EHIK}
\end{reference}
Soit $f : X \to Y$ un morphisme de schémas. Supposons que
$f$ soit complètement décomposé,
$f$ soit localement de présentation finie, et que
$Y$ soit quasi-compact et quasi-séparé.
Alors il existe $n \geq 0$ et des morphismes
$Z_i \to Y$, $i = 1, \ldots, n$, ayant les propriétés suivantes :
\begin{enumerate}
\item $\coprod Z_i \to Y$ est surjectif,
\item $Z_i \to Y$ est une immersion pour tout $i$,
\item $Z_i \to Y$ est de présentation finie pour tout $i$, et
\item le changement de base $X \times_Y Z_i \to Z_i$ possède une section
pour tout $i$.
\end{enumerate}
\end{lemma}

\begin{proof}
Soit $y \in Y$. Par hypothèse, il existe un morphisme
$\sigma : \Spec(\kappa(y)) \to X$ au-dessus de $Y$. Nous pouvons écrire $\Spec(\kappa(y))$
comme la limite d'un système dirigé de schémas affines $Z$ au-dessus de $Y$ tels que
$Z \to Y$ soit une immersion de présentation finie.
En effet, choisissons un ouvert affine $y \in \Spec(A) \subset Y$
et disons que $y$ correspond à l'idéal premier $\mathfrak p$ de $A$.
Alors $\kappa(\mathfrak p)$ est la colimite filtrante des
anneaux $(A/I)_f$, où $I \subset \mathfrak p$ est un idéal de type fini
et $f \in A$, $f \not \in \mathfrak p$.
Les morphismes $Z = \Spec((A/I)_f) \to Y$ sont des immersions de
présentation finie ; la quasi-compacité de $Z \to Y$
résulte du fait que $Y$ est quasi-séparé, voir
Schémas, lemme \ref{schemes-lemma-quasi-compact-permanence}.
D'après Limites, proposition
\ref{limits-proposition-characterize-locally-finite-presentation},
pour l'un de ces $Z$, il existe un morphisme $\sigma' : Z \to X$ au-dessus de $Y$
qui coïncide avec $\sigma$ sur le spectre de $\kappa(\mathfrak p)$.
Comme $\sigma'$ est un morphisme au-dessus de $Y$, nous obtenons une section
de la projection $X \times_Y Z \to Z$.

\medskip\noindent
Nous concluons que $Y$ est la réunion des images d'immersions
$Z \to Y$ de présentation finie telles que $X \times_Y Z \to Z$
possède une section. Puisque l'image de $Z \to Y$ est constructible
(Morphismes, lemme \ref{morphisms-lemma-chevalley})
et puisque $Y$ est compact pour la topologie constructible
(Propriétés, lemme
\ref{properties-lemma-quasi-compact-quasi-separated-spectral} et
Topologie, lemme \ref{topology-lemma-constructible-hausdorff-quasi-compact}),
nous voyons qu'un nombre fini de ces immersions suffit.
\end{proof}

\begin{lemma}
\label{more-morphisms-lemma-descend-cd}
Soit $S = \lim_{\lambda \in \Lambda} S_\lambda$
la limite d'un système dirigé de schémas dont les morphismes de transition sont affines.
Soit $0 \in \Lambda$ et soit $f_0 : X_0 \to Y_0$
un morphisme de schémas au-dessus de $S_0$.
Pour $\lambda \geq 0$, soit $f_\lambda : X_\lambda \to Y_\lambda$
le changement de base de $f_0$ à $S_\lambda$, et
soit $f : X \to Y$ le changement de base de $f_0$ à $S$. Si
\begin{enumerate}
\item $f$ est complètement décomposé,
\item $Y_0$ est quasi-compact et quasi-séparé, et
\item $f_0$ est localement de présentation finie,
\end{enumerate}
alors il existe un $\lambda \geq 0$ tel que $f_\lambda$
soit complètement décomposé.
\end{lemma}

\begin{proof}
Puisque $Y_0$ est quasi-compact et quasi-séparé, le schéma $Y$,
qui est affine sur $Y_0$, est quasi-compact et quasi-séparé.
Choisissons $n \geq 0$ et $Z_i \to Y$, $i = 1, \ldots, n$, comme dans le
lemme \ref{more-morphisms-lemma-decompose}. Notons $\sigma_i : Z_i \to X$
des morphismes au-dessus de $Y$ qui existent par notre choix des $Z_i$.
Quitte à remplacer $0 \in \Lambda$ par un indice plus grand, nous pouvons supposer qu'il existe
des morphismes $Z_{i, 0} \to Y_0$ de présentation finie
dont les changements de base à $S$ sont les morphismes $Z_i \to Y$, voir
Limites, lemme \ref{limits-lemma-descend-finite-presentation}.
D'après Limites, lemme \ref{limits-lemma-descend-immersion},
nous pouvons supposer, quitte à remplacer $0$ par un indice plus grand, que $Z_{i, 0} \to Y_0$
est une immersion. Puisque $\coprod Z_i \to Y$ est surjectif, nous pouvons supposer,
quitte à remplacer $0$ par un indice plus grand, que
$\coprod Z_{i, 0} \to Y_0$ est surjectif, voir
Limites, lemme \ref{limits-lemma-descend-surjective}. Observons que
$Z_i = \lim_{\lambda \geq 0} Z_{i, \lambda}$,
où $Z_{i, \lambda} = Y_\lambda \times_{Y_0} Z_{i, 0}$.
Considérons les composés
$$
Z_i \xrightarrow{\sigma_i} X \to X_0
$$
comme des morphismes au-dessus de $Y_0$. Puisque $f_0$ est localement de présentation
finie, Limites, proposition
\ref{limits-proposition-characterize-locally-finite-presentation}
nous permet de trouver un $\lambda \geq 0$ tel qu'il existe des
morphismes $\sigma'_{i, \lambda} : Z_{i, \lambda} \to X_0$
au-dessus de $Y_0$ dont les précomposés avec $Z_i \to Z_{i, \lambda}$
sont les flèches affichées. Bien entendu, $\sigma'_{i, \lambda}$
détermine alors un morphisme $\sigma_{i, \lambda} : Z_{i, \lambda} \to
X_\lambda = X_0 \times_{Y_0} Y_\lambda$ au-dessus de $Y_\lambda$.
Puisque $\coprod Z_{i, \lambda} \to Y_\lambda$ est surjectif,
nous concluons que $X_\lambda \to Y_\lambda$ est complètement décomposé.
\end{proof}




\section{Familles de modules inversibles amples}
\label{more-morphisms-section-families-ample-invertible-modules}

\noindent
Nous poursuivons la discussion de
Morphismes, section \ref{morphisms-section-families-ample-invertible-modules}.

\begin{lemma}
\label{more-morphisms-lemma-ample-family-ample-relative}
Soit $f : X \to Y$ un morphisme de schémas. Supposons que
\begin{enumerate}
\item $Y$ possède une famille ample de modules inversibles,
\item il existe un module inversible $f$-ample sur $X$.
\end{enumerate}
Alors $X$ possède une famille ample de modules inversibles.
\end{lemma}

\begin{proof}
Soit $\mathcal{L}$ un module inversible $f$-ample sur $X$.
Cela implique en particulier que $f$ est quasi-compact, voir
Morphismes, définition \ref{morphisms-definition-relatively-ample}.
Puisque $Y$ est quasi-compact d'après Morphismes, définition
\ref{morphisms-definition-family-ample-invertible-modules},
nous voyons que $X$ est quasi-compact (et que $X$ lui-même satisfait donc la
première condition de Morphismes, définition
\ref{morphisms-definition-family-ample-invertible-modules}).
Soit $x \in X$ d'image $y \in Y$. Par l'hypothèse (2), nous pouvons
trouver un $\mathcal{O}_Y$-module inversible $\mathcal{N}$ et une section
$t \in \Gamma(Y, \mathcal{N})$ tels que le lieu $Y_t$ où $t$
ne s'annule pas soit affine. Alors $\mathcal{L}$ est ample sur
$f^{-1}(Y_t) = X_{f^*t}$, et nous pouvons donc trouver une section
$s \in \Gamma(X_{f^*t}, \mathcal{L})$ telle que $(X_{f^*t})_s$ soit affine
et contienne $x$. D'après
Propriétés, lemme \ref{properties-lemma-invert-s-sections},
pour un certain $n \geq 0$, le produit $(f^*t)^n s$ se prolonge en une section
$s' \in \Gamma(X, f^*\mathcal{N}^{\otimes n} \otimes \mathcal{L})$.
Enfin, la section $s'' = f^* ts'$ de
$f^*\mathcal{N}^{\otimes n + 1} \otimes \mathcal{L}$
s'annule en tout point de $X \setminus X_{f^*t}$ ; ainsi,
nous voyons que $X_{s''} = (X_{f^*t})_s$ est affine, comme voulu.
\end{proof}

\begin{lemma}
\label{more-morphisms-lemma-resolution-property-goes-up-affine}
Soit $f : X \to Y$ un morphisme affine ou quasi-affine de schémas.
Si $Y$ possède une famille ample de modules inversibles, il en va de même pour $X$.
\end{lemma}

\begin{proof}
D'après Morphismes, lemme \ref{morphisms-lemma-quasi-affine-O-ample},
c'est un cas particulier du
lemme \ref{more-morphisms-lemma-ample-family-ample-relative}.
\end{proof}






\section{Éclatements et familles amples de modules inversibles}
\label{more-morphisms-section-ample-family-by-blowing-up}

\noindent
Nous démontrons un résultat de \cite{Gross-thesis}.

\begin{lemma}
\label{more-morphisms-lemma-get-ample-family}
Soit $X$ un schéma. Supposons donnés des diviseurs de Cartier effectifs
$D_1, \ldots, D_m$ sur $X$ et des modules inversibles
$\mathcal{L}_1, \ldots, \mathcal{L}_m$ tels que
$\bigcap D_i = \emptyset$ et que $\mathcal{L}_i|_{X \setminus D_i}$
soit ample. Alors $X$ possède une famille ample de modules inversibles.
\end{lemma}

\begin{proof}
Soit $x \in X$. Choisissons un indice $i \in \{1, \ldots, m\}$
tel que $x \not \in D_i$. Posons $U_i = X \setminus D_i$.
Puisque $\mathcal{L}_i|_{U_i}$ est ample, nous pouvons trouver un $n \geq 1$ et
une section $s \in \Gamma(U_i, \mathcal{L}_i^{\otimes n})$
telle que le lieu $(U_i)_s$ où $s$ ne s'annule pas soit affine
(Propriétés, définition \ref{properties-definition-ample}).
Puisque $U_i$ est le lieu où la section canonique
$1 \in \mathcal{O}_X(D_i)$ ne s'annule pas, nous déduisons de
Propriétés, lemme \ref{properties-lemma-invert-s-sections}
qu'il existe un $N \geq 0$ tel que $s$ se prolonge en une section
$$
s' \in \Gamma(X, \mathcal{L}_i^{\otimes n} \otimes_{\mathcal{O}_X}
\mathcal{O}_X(N D_i))
$$
Après avoir remplacé $N$ par $N + 1$, nous voyons que $s'$ s'annule en tout
point de $D_i$, et donc que $X_{s'} = (U_i)_s$ est affine.
Cela prouve que $X$ possède une famille ample de modules inversibles, voir
Morphismes, définition
\ref{morphisms-definition-family-ample-invertible-modules}.
\end{proof}

\begin{lemma}
\label{more-morphisms-lemma-blow-up-ample-family}
\begin{reference}
\cite[Proposition 1.3.1]{Gross-thesis}
\end{reference}
Soit $X$ un schéma quasi-compact et quasi-séparé n'ayant qu'un nombre fini de
composantes irréductibles. Il existe un ouvert dense quasi-compact
$U \subset X$ et un éclatement $U$-admissible $X' \to X$ tels que le
schéma $X'$ possède une famille ample de modules inversibles.
\end{lemma}

\begin{proof}
Soient $\eta_1, \ldots, \eta_n \in X$ les points génériques des
composantes irréductibles de $X$. D'après
Propriétés, lemme \ref{properties-lemma-point-and-maximal-points-affine}
et la quasi-compacité de $X$, nous pouvons trouver un recouvrement ouvert affine
fini $X = U_1 \cup \ldots \cup U_m$ tel que
chaque $U_i$ contienne $\eta_1, \ldots, \eta_n$. En particulier,
le sous-ensemble ouvert quasi-compact $U = U_1 \cap \ldots \cap U_m$ est dense dans $X$.
Soit $\mathcal{I}_i \subset \mathcal{O}_X$ un faisceau quasi-cohérent d'idéaux de type
fini tel que $U_i = X \setminus Z_i$, où $Z_i = V(\mathcal{I}_i)$,
voir Propriétés, lemme \ref{properties-lemma-quasi-coherent-finite-type-ideals}.
Soit
$$
f : X' \longrightarrow X
$$
l'éclatement de $X$ le long du faisceau d'idéaux
$\mathcal{I} = \mathcal{I}_1 \cdots \mathcal{I}_m$.
Notons que $f$ est un éclatement $U$-admissible, puisque $V(\mathcal{I})$
est (ensemblistement) la réunion des $Z_i$, qui sont disjoints de $U$.
De plus, $f$ est un morphisme projectif et
$\mathcal{O}_{X'}(1)$ est relativement ample pour $f$, voir
Diviseurs, lemme \ref{divisors-lemma-blowing-up-projective}.
D'après Diviseurs, lemme \ref{divisors-lemma-blowing-up-two-ideals},
pour chaque $i$, le morphisme $f'$ se factorise en $X' \to X'_i \to X$,
où $X'_i \to X$ est l'éclatement le long de $\mathcal{I}_i$
et $X' \to X'_i$ est un autre éclatement (à savoir, le long du tiré en arrière
du produit des idéaux $\mathcal{I}_j$, en omettant $\mathcal{I}_i$).
Il en résulte que $D_i = f^{-1}(Z_i) \subset X'$ est un diviseur de Cartier
effectif, voir
Diviseurs, lemmes \ref{divisors-lemma-blow-up-pullback-effective-Cartier} et
\ref{divisors-lemma-blowing-up-gives-effective-Cartier-divisor}.
Nous avons $X' \setminus D_i = f^{-1}(U_i)$. Comme $\mathcal{O}_{X'}(1)$ est
$f$-ample, la restriction de
$\mathcal{O}_{X'}(1)$ à $X' \setminus D_i$ est ample. Il résulte du
lemme \ref{more-morphisms-lemma-get-ample-family}
que $X'$ possède une famille ample de modules inversibles.
\end{proof}

\begin{proposition}
\label{more-morphisms-proposition-envelope-with-resolution-property}
Soit $X$ un schéma quasi-compact et quasi-séparé. Il existe un
morphisme $f : Y \to X$ qui est de présentation finie, propre et
complètement décomposé (définition \ref{more-morphisms-definition-cd-morphism}),
tel que le schéma $Y$ possède une famille ample de modules inversibles.
\end{proposition}

\begin{proof}
D'après Limites, proposition \ref{limits-proposition-approximate},
il existe un morphisme affine $X \to X_0$, où $X_0$
est un schéma de type fini sur $\mathbf{Z}$. Nous construisons ci-dessous
un morphisme $Y_0 \to X_0$ ayant toutes les propriétés voulues.
En posant alors $Y = X \times_{X_0} Y_0$ et en prenant pour $f$
la projection $f : Y \to X$, nous concluons.
En effet, $f$ est de présentation finie
(Morphismes, lemme \ref{morphisms-lemma-base-change-finite-presentation}),
$f$ est propre
(Morphismes, lemme \ref{morphisms-lemma-base-change-proper}),
$f$ est complètement décomposé
(lemme \ref{more-morphisms-lemma-base-change-cd}). Enfin, puisque $Y \to Y_0$ est
affine (étant le changement de base de $X \to X_0$), nous voyons que $Y$ possède
une famille ample de modules inversibles d'après le
lemme \ref{more-morphisms-lemma-resolution-property-goes-up-affine}.
Nous sommes ainsi ramenés au cas traité dans le paragraphe suivant.

\medskip\noindent
Supposons $X$ de type fini sur $\mathbf{Z}$. En particulier,
$\dim(X) < \infty$. Nous allons raisonner par récurrence sur $\dim(X)$.
Si $\dim(X) = 0$, alors $X$ est affine et possède la propriété de résolution.
En général, il existe un ouvert dense $U \subset X$ et un
éclatement $U$-admissible $X' \to X$ tel que $X'$ possède
une famille ample de modules inversibles, voir
le lemme \ref{more-morphisms-lemma-blow-up-ample-family}.
Puisque $f : X' \to X$ est un isomorphisme au-dessus de $U$,
nous voyons que tout point de $U$ se relève en un point de $X'$
ayant le même corps résiduel.
Soit $Z = X \setminus U$, muni de la structure de sous-schéma réduit induite.
Alors $\dim(Z) < \dim(X)$, puisque $U$ est dense dans $X$ (voir ci-dessus).
Par récurrence, nous trouvons un morphisme propre, complètement décomposé
$W \to Z$ tel que $W$ possède une famille ample de modules inversibles.
Il s'ensuit alors que $Y = X' \amalg W \to X$ est le morphisme
voulu.
\end{proof}




\section{Le critère extensif pour les immersions fermées}
\label{more-morphisms-section-extensive-criterion-closed-immersions}

\noindent
Dans cette section, nous donnons un critère pour qu'un morphisme de schémas 
soit une immersion fermée.

\begin{lemma}
\label{more-morphisms-lemma-affine-extensive-criterion}
Un morphisme $f : X \to Y$ de schémas affines est une immersion fermée 
si et seulement si, pour tout homomorphisme injectif d'anneaux $A \to B$ et tout 
carré commutatif
$$
\xymatrix{
\Spec(B) \ar[d] \ar[r] & X \ar[d]^f \\
\Spec(A) \ar[r] \ar@{..>}[ur] & Y
}
$$
il existe un relèvement $\Spec(A) \to X$ rendant les deux triangles commutatifs.
\end{lemma}

\begin{proof}
Soit le morphisme $f$ donné par l'homomorphisme d'anneaux $\phi : R \to S$.
Alors $f$ est une immersion fermée si et seulement si $\phi$ est surjectif.

\medskip\noindent
Supposons tout d'abord que $\phi$ soit surjectif.
Soit $\psi : A \to B$ un homomorphisme injectif d'anneaux, et supposons donné
un diagramme commutatif
$$
\xymatrix{
R \ar[r]^\alpha \ar[d]^\phi & A \ar[d]^\psi \\
S \ar[r]^\beta \ar@{..>}[ur] & B
}
$$
Nous définissons alors un relèvement $S \to A$ par $s \mapsto \alpha(r)$, où
$r \in R$ est tel que $\phi(r) = s$.
Cette définition est bien posée, car $\psi$ est injectif et le carré est commutatif.
Comme le passage au spectre établit une anti-équivalence entre les anneaux
commutatifs et les schémas affines, la propriété de relèvement voulue pour
$f$ est vérifiée.

\medskip\noindent
Supposons ensuite que $\phi$ admette des relèvements par rapport à tous les
homomorphismes injectifs d'anneaux $\psi: A \to B$.
Remarquons que $\phi(R)$ est un sous-anneau de $S$ ; on obtient donc un carré
commutatif
$$
\xymatrix{
R \ar[r] \ar[d]^\phi  & \phi(R) \ar[d] \\
S \ar@{=}[r] \ar@{..>}[ur] & S
}
$$
dans lequel il existe un relèvement $S \to \phi(R)$.
Par suite, l'inclusion $\phi(R) \to S$ est nécessairement un isomorphisme, ce qui
montre que $\phi$ est surjectif, et le résultat s'ensuit.
\end{proof}

\begin{lemma}
\label{more-morphisms-lemma-scheme-affine-if-aff-is-split-mono}
Soit $X$ un schéma.
Si le morphisme canonique $X \to \Spec(\Gamma(X, \mathcal{O}_X))$
considéré dans Schémas, lemme \ref{schemes-lemma-morphism-into-affine},
admet une rétraction, alors $X$ est un schéma affine.
\end{lemma}

\begin{proof}
Posons $S = \Spec(\Gamma(X, \mathcal{O}_X))$ et notons $f : X \to S$ le morphisme
donné dans le lemme. Soit $s : S \to X$ une rétraction ; ainsi $\text{id}_X = sf$.
Alors $f s f = \text{id}_S f$. Comme $f$ induit un isomorphisme
$\Gamma(S, \mathcal{O}_S) \to \Gamma(X, \mathcal{O}_X)$,
cela signifie que $fs$ et $\text{id}_S$ induisent la même application
sur $\Gamma(S, \mathcal{O}_S)$. D'où $f s = \text{id}_S$, puisque $S$
est affine. Par suite, $f$ est un isomorphisme et $X$ est un schéma affine,
comme il fallait le démontrer.
\end{proof}

\begin{lemma}
\label{more-morphisms-lemma-aff-is-injective-on-sheaves}
Soit $X$ un schéma. Soit $f : X \to S = \Spec(\Gamma(X, \mathcal{O}_X))$
le morphisme canonique considéré dans
Schémas, lemme \ref{schemes-lemma-morphism-into-affine}.
Le plus grand $\mathcal{O}_S$-module quasi-cohérent contenu
dans le noyau de $f^\sharp : \mathcal{O}_S \to f_*\mathcal{O}_X$
est nul. Si $X$ est quasi-compact, alors $f^\sharp$ est injectif.
En particulier, si $X$ est quasi-compact, alors $f$ est un
morphisme dominant.
\end{lemma}

\begin{proof}
Soit $M \subset \Gamma(S, \mathcal{O}_S)$ le sous-module correspondant
au plus grand $\mathcal{O}_S$-module quasi-cohérent contenu
dans le noyau de $f^\sharp$. Alors tout élément $a \in M$ est envoyé
sur zéro par $f^\sharp$. Or $f^\sharp(a)$ est l'élément de
$$
\Gamma(S, f_*\mathcal{O}_X) = \Gamma(X, \mathcal{O}_X) =
\Gamma(S, \mathcal{O}_S)
$$
qui correspond à $a$ lui-même ! Ainsi $a = 0$. Par conséquent $M = 0$, ce qui prouve
la première assertion. Remarquons que cela équivaut à dire que le morphisme
$f : X \to S$ est schématiquement surjectif.

\medskip\noindent
Si $X$ est quasi-compact, alors $\Ker(f^\sharp)$ est quasi-cohérent d'après
Morphismes, lemme \ref{morphisms-lemma-quasi-compact-scheme-theoretic-image}.
Ainsi $\Ker(f^\sharp) = 0$ et $f^\sharp$ est injectif.
Dans ce cas, $f$ est un morphisme dominant en vertu de la partie (4) de
Morphismes, lemme \ref{morphisms-lemma-quasi-compact-scheme-theoretic-image}.
\end{proof}

\begin{lemma}
\label{more-morphisms-lemma-extensive-criterion}
Soit $f: X \to Y$ un morphisme quasi-compact de schémas.
Alors $f$ est une immersion fermée si et seulement si, pour tout homomorphisme
injectif d'anneaux $A \to B$ et tout carré commutatif
$$
\xymatrix{
\Spec(B) \ar[d] \ar[r] & X \ar[d]^f \\
\Spec(A) \ar[r] \ar@{..>}[ur] & Y
}
$$
il existe un relèvement $\Spec A \to X$ rendant le diagramme commutatif.
\end{lemma}

\begin{proof}
Supposons que $f$ soit une immersion fermée. Soit $A \to B$ un homomorphisme injectif
d'anneaux et considérons un carré commutatif
$$
\xymatrix{
\Spec(B) \ar[d] \ar[r]  & X \ar[d]^f \\
\Spec(A) \ar[r] \ar@{..>}[ur] & Y
}
$$
Alors $\Spec(A) \times_Y X \to \Spec(A)$ est une immersion fermée et
il existe donc un idéal $I \subset A$ et un diagramme commutatif
$$
\xymatrix{
\Spec(B) \ar[d] \ar[r] & \Spec(A/I) \ar[r] \ar[d] & X \ar[d]^f \\
\Spec(A) \ar[r] \ar@{..>}[ur] & \Spec(A) \ar[r] & Y
}
$$
Le lemme \ref{more-morphisms-lemma-affine-extensive-criterion} fournit un relèvement.

\medskip\noindent
Supposons que $f$ satisfasse à la propriété de relèvement énoncée dans le lemme.
La propriété pour $f$ d'être une immersion fermée est locale sur $Y$ ;
nous pouvons donc supposer, ce que nous ferons, que $Y$ est affine.
En particulier, $Y$ est quasi-compact, et $X$ est donc quasi-compact.
Il existe donc un recouvrement fini par des ouverts affines $X = U_1 \cup \ldots \cup U_n$.
La source du morphisme
$$
\pi : U = \coprod U_i \longrightarrow X
$$
est affine et l'homomorphisme d'anneaux induit
$\Gamma(X, \mathcal{O}_X) \to \Gamma(U, \mathcal{O}_U)$
est injectif. Par hypothèse, il existe un relèvement dans le diagramme
$$
\xymatrix{
U \ar[r]^\pi \ar[d] & X \ar[d]^f \\
\Spec(\Gamma(X, \mathcal{O}_X)) \ar[r]^-{f'} \ar@{..>}[ur]^h & Y
}
$$
où $f'$ est le morphisme de schémas affines correspondant à
l'homomorphisme d'anneaux $\Gamma(Y, \mathcal{O}_Y) \to \Gamma(X, \mathcal{O}_X)$.
Il résulte du fait que $\pi$ est un épimorphisme
que le morphisme $h$ est une rétraction du morphisme canonique
$X \to \Spec(\Gamma(X, \mathcal{O}_X))$ ; nous omettons les détails. Ainsi $X$ est affine d'après le
lemme \ref{more-morphisms-lemma-scheme-affine-if-aff-is-split-mono}.
Le lemme \ref{more-morphisms-lemma-affine-extensive-criterion} permet alors de
conclure que $f$ est une immersion fermée.
\end{proof}









% OK, start here.
%

\chapter{Compléments sur la platitude}


\phantomsection
\label{flat-section-phantom}




\section{Introduction}
\label{flat-section-introduction}

\noindent
Dans ce chapitre, nous exposons des résultats avancés sur les modules plats et les
morphismes plats de schémas, ainsi que des applications. La plupart des résultats de platitude
figurent dans l'article \cite{GruRay} de Raynaud et Gruson.

\medskip\noindent
Avant de lire ce chapitre, nous conseillons au lecteur de prendre connaissance
des résultats suivants (cette liste sert aussi à retrouver
les résultats précédents) :
\begin{enumerate}
\item Généralités sur les modules plats dans
Algèbre, section \ref{algebra-section-flat}.
\item Relations entre les groupes $\text{Tor}$ et la platitude, voir
Algèbre, section \ref{algebra-section-tor}.
\item Critères de platitude, voir
Algèbre, section \ref{algebra-section-criteria-flatness}
(cas noethérien),
Algèbre, section \ref{algebra-section-flatness-artinian}
(cas artinien),
Algèbre, section \ref{algebra-section-more-flatness-criteria}
(cas non noethérien), et enfin
Compléments sur les morphismes, section \ref{more-morphisms-section-criterion-flat-fibres}.
\item Platitude générique, voir
Algèbre, section \ref{algebra-section-generic-flatness}
et
Morphismes, section \ref{morphisms-section-generic-flatness}.
\item Ouverture du lieu de platitude, voir
Algèbre, section \ref{algebra-section-open-flat}
et
Compléments sur les morphismes, section \ref{more-morphisms-section-open-flat}.
\item Platification, voir
Compléments d'algèbre, sections
\ref{more-algebra-section-flattening},
\ref{more-algebra-section-flattening-artinian},
\ref{more-algebra-section-flattening-local-base},
\ref{more-algebra-section-flattening-local-source-base}, et
\ref{more-algebra-section-flattening-Noetherian-complete-local}.
\item Résultats supplémentaires dans
Compléments d'algèbre, sections \ref{more-algebra-section-descent-flatness-integral},
\ref{more-algebra-section-torsion-flat},
\ref{more-algebra-section-flat-finite}, et
\ref{more-algebra-section-blowup-flat}.
\end{enumerate}
Comme applications des résultats de platitude, nous traitons les sujets
suivants : une version non noethérienne du théorème d'existence de Grothendieck,
les éclatements et la platitude, le théorème de compactification de Nagata,
la topologie h, les carrés d'éclatement et la descente, la normalisation faible,
la descente des fibrés vectoriels en caractéristique positive et
la structure locale des complexes parfaits pour la topologie h.






\section{Lemmes sur la localisation \'etale}
\label{flat-section-etale-localization}

\noindent
Nous réunissons dans cette section quelques lemmes de localisation \'etale qui seront
utiles dans la suite de ce chapitre. On peut omettre cette section en première lecture.

\begin{lemma}
\label{flat-lemma-lift-etale}
Soit $i : Z \to X$ une immersion fermée de schémas affines.
Soit $Z' \to Z$ un morphisme \'etale, avec $Z'$ affine.
Il existe alors un morphisme \'etale $X' \to X$, avec $X'$
affine, tel que $Z' \cong Z \times_X X'$ comme schémas sur $Z$.
\end{lemma}

\begin{proof}
Voir
Algèbre, Lemme \ref{algebra-lemma-lift-etale}.
\end{proof}

\begin{lemma}
\label{flat-lemma-etale-at-point}
Soit
$$
\xymatrix{
X \ar[d] & X' \ar[l] \ar[d] \\
S & S' \ar[l]
}
$$
un diagramme commutatif de schémas où $X' \to X$ et $S' \to S$ sont \'etales.
Soit $s' \in S'$ un point. Alors
$$
X' \times_{S'} \Spec(\mathcal{O}_{S', s'})
\longrightarrow
X \times_S \Spec(\mathcal{O}_{S', s'})
$$
est \'etale.
\end{lemma}

\begin{proof}
Cela résulte du fait que $X' \to X_{S'}$ est \'etale, en tant que morphisme de
schémas \'etales sur $X$, voir
Morphismes, Lemme \ref{morphisms-lemma-etale-permanence},
et du fait que le changement de base d'un morphisme \'etale est \'etale, voir
Morphismes, Lemme \ref{morphisms-lemma-base-change-etale}.
\end{proof}

\begin{lemma}
\label{flat-lemma-etale-flat-up-down}
Soient $X \to T \to S$ des morphismes de schémas, avec $T \to S$ \'etale.
Soit $\mathcal{F}$ un $\mathcal{O}_X$-module quasi-cohérent.
Soit $x \in X$ un point. Alors
$$
\mathcal{F}\text{ plat sur }S\text{ en }x
\Leftrightarrow
\mathcal{F}\text{ plat sur }T\text{ en }x
$$
En particulier, $\mathcal{F}$ est plat sur $S$ si et seulement si $\mathcal{F}$
est plat sur $T$.
\end{lemma}

\begin{proof}
Puisqu'un morphisme \'etale est un morphisme plat (voir
Morphismes, Lemme \ref{morphisms-lemma-etale-flat}),
l'implication ``$\Leftarrow$'' résulte de
Algèbre, Lemme \ref{algebra-lemma-composition-flat}.
Réciproquement, supposons que $\mathcal{F}$ soit plat en $x$ sur $S$.
Notons $\tilde x \in X \times_S T$ le point au-dessus de $x$ dans $X$
et de l'image de $x$ dans $T$, prise comme point de $T$.
Alors $(X \times_S T \to X)^*\mathcal{F}$ est plat en $\tilde x$ sur $T$
pour $\text{pr}_2 : X \times_S T \to T$, voir
Morphismes, Lemme \ref{morphisms-lemma-base-change-module-flat}.
La diagonale $\Delta_{T/S} : T \to T \times_S T$ est une immersion ouverte ;
il suffit de combiner
Morphismes, Lemmes \ref{morphisms-lemma-diagonal-unramified-morphism} et
\ref{morphisms-lemma-etale-smooth-unramified}.
Ainsi $X$ s'identifie à un sous-schéma ouvert de $X \times_S T$,
la restriction de $\text{pr}_2$ à cet ouvert est le morphisme donné $X \to T$,
le point $\tilde x$ correspond au point $x$ de cet ouvert et
la restriction de $(X \times_S T \to X)^*\mathcal{F}$ à cet ouvert est $\mathcal{F}$.
On en déduit que $\mathcal{F}$ est plat en $x$ sur $T$.
\end{proof}

\begin{lemma}
\label{flat-lemma-etale-flat-up-down-local-ring}
Soit $T \to S$ un morphisme \'etale. Soit $t \in T$ d'image $s \in S$.
Soit $M$ un $\mathcal{O}_{T, t}$-module. Alors
$$
M\text{ plat sur }\mathcal{O}_{S, s}
\Leftrightarrow
M\text{ plat sur }\mathcal{O}_{T, t}.
$$
\end{lemma}

\begin{proof}
On peut remplacer $S$ par un voisinage affine de $s$, puis
$T$ par un voisinage affine de $t$.
Posons $\mathcal{F} = (\Spec(\mathcal{O}_{T, t}) \to T)_*\widetilde M$.
C'est un faisceau quasi-cohérent (voir
Schémas, Lemme \ref{schemes-lemma-push-forward-quasi-coherent},
ou raisonner directement)
sur $T$, dont la fibre en $t$ est $M$ (détails omis).
On applique
le Lemme \ref{flat-lemma-etale-flat-up-down}.
\end{proof}

\begin{lemma}
\label{flat-lemma-flat-up-down-henselization}
Soient $S$ un schéma et $s \in S$ un point. Notons $\mathcal{O}_{S, s}^h$
(resp.\ $\mathcal{O}_{S, s}^{sh}$) le hensélisé (resp.\ le hensélisé
strict), voir
Algèbre, Définition \ref{algebra-definition-henselization}.
Soit $M^{sh}$ un $\mathcal{O}_{S, s}^{sh}$-module.
Les conditions suivantes sont équivalentes :
\begin{enumerate}
\item $M^{sh}$ est plat sur $\mathcal{O}_{S, s}$,
\item $M^{sh}$ est plat sur $\mathcal{O}_{S, s}^h$, et
\item $M^{sh}$ est plat sur $\mathcal{O}_{S, s}^{sh}$.
\end{enumerate}
Si $M^{sh} = M^h \otimes_{\mathcal{O}_{S, s}^h} \mathcal{O}_{S, s}^{sh}$,
elles sont aussi équivalentes à
\begin{enumerate}
\item[(4)] $M^h$ est plat sur $\mathcal{O}_{S, s}$, et
\item[(5)] $M^h$ est plat sur $\mathcal{O}_{S, s}^h$.
\end{enumerate}
Si $M^h = M \otimes_{\mathcal{O}_{S, s}} \mathcal{O}_{S, s}^h$,
elles sont aussi équivalentes à
\begin{enumerate}
\item[(6)] $M$ est plat sur $\mathcal{O}_{S, s}$.
\end{enumerate}
\end{lemma}

\begin{proof}
D'après Compléments d'algèbre, Lemme
\ref{more-algebra-lemma-dumb-properties-henselization},
les homomorphismes locaux
$\mathcal{O}_{S, s} \to \mathcal{O}_{S, s}^h \to \mathcal{O}_{S, s}^{sh}$
sont fidèlement plats.
Ainsi (3) $\Rightarrow$ (2) $\Rightarrow$ (1) et
(5) $\Rightarrow$ (4) résultent de
Algèbre, Lemme \ref{algebra-lemma-composition-flat}.
Par fidèle platitude, les équivalences (6) $\Leftrightarrow$ (5) et
(5) $\Leftrightarrow$ (3) résultent de
Algèbre, Lemme \ref{algebra-lemma-flatness-descends}.
Il suffit donc de montrer que
(1) $\Rightarrow$ (2) $\Rightarrow$ (3) et
(4) $\Rightarrow$ (5).
Pour cela, on peut supposer que $S$ est un schéma affine.

\medskip\noindent
Supposons (1). D'après
le Lemme \ref{flat-lemma-etale-flat-up-down-local-ring},
on voit que $M^{sh}$ est plat sur $\mathcal{O}_{T, t}$ pour
tout voisinage \'etale $(T, t) \to (S, s)$. Puisque $\mathcal{O}_{S, s}^h$
et $\mathcal{O}_{S, s}^{sh}$ sont des limites inductives filtrantes d'anneaux locaux
de la forme $\mathcal{O}_{T, t}$ (voir
Algèbre, Lemmes \ref{algebra-lemma-henselization-different}
et \ref{algebra-lemma-strict-henselization-different}),
on en conclut que $M^{sh}$ est plat sur $\mathcal{O}_{S, s}^h$
et $\mathcal{O}_{S, s}^{sh}$ d'après
Algèbre, Lemme \ref{algebra-lemma-colimit-rings-flat}.
Ainsi (1) implique (2) et (3). Cela entraîne aussi
(2) $\Rightarrow$ (3) en remplaçant $\mathcal{O}_{S, s}$ par
$\mathcal{O}_{S, s}^h$. Le même raisonnement démontre
(4) $\Rightarrow$ (5).
\end{proof}

\begin{lemma}
\label{flat-lemma-tor-amplitude-up-down-henselization}
Soient $S$ un schéma et $s \in S$ un point. Notons $\mathcal{O}_{S, s}^h$
(resp.\ $\mathcal{O}_{S, s}^{sh}$) le hensélisé (resp.\ le hensélisé
strict), voir
Algèbre, Définition \ref{algebra-definition-henselization}.
Soit $M^{sh}$ un objet de $D(\mathcal{O}_{S, s}^{sh})$.
Soient $a, b \in \mathbf{Z}$.
Les conditions suivantes sont équivalentes :
\begin{enumerate}
\item $M^{sh}$ est de Tor-amplitude contenue dans $[a, b]$ sur $\mathcal{O}_{S, s}$,
\item $M^{sh}$ est de Tor-amplitude contenue dans $[a, b]$ sur $\mathcal{O}_{S, s}^h$, et
\item $M^{sh}$ est de Tor-amplitude contenue dans $[a, b]$ sur $\mathcal{O}_{S, s}^{sh}$.
\end{enumerate}
Si $M^{sh} =
M^h \otimes_{\mathcal{O}_{S, s}^h}^\mathbf{L} \mathcal{O}_{S, s}^{sh}$
pour $M^h \in D(\mathcal{O}_{S, s}^h)$, elles sont aussi équivalentes à
\begin{enumerate}
\item[(4)] $M^h$ est de Tor-amplitude contenue dans $[a, b]$ sur $\mathcal{O}_{S, s}$, et
\item[(5)] $M^h$ est de Tor-amplitude contenue dans $[a, b]$ sur $\mathcal{O}_{S, s}^h$.
\end{enumerate}
Si $M^h = M \otimes_{\mathcal{O}_{S, s}}^\mathbf{L} \mathcal{O}_{S, s}^h$
pour $M \in D(\mathcal{O}_{S, s})$,
elles sont aussi équivalentes à
\begin{enumerate}
\item[(6)] $M$ est de Tor-amplitude contenue dans $[a, b]$ sur $\mathcal{O}_{S, s}$.
\end{enumerate}
\end{lemma}

\begin{proof}
D'après Compléments d'algèbre, Lemme
\ref{more-algebra-lemma-dumb-properties-henselization},
les homomorphismes locaux
$\mathcal{O}_{S, s} \to \mathcal{O}_{S, s}^h \to \mathcal{O}_{S, s}^{sh}$
sont fidèlement plats.
Ainsi (3) $\Rightarrow$ (2) $\Rightarrow$ (1) et
(5) $\Rightarrow$ (4) résultent de
Compléments d'algèbre, Lemme \ref{more-algebra-lemma-flat-push-tor-amplitude}.
Par fidèle platitude, les équivalences (6) $\Leftrightarrow$ (5) et
(5) $\Leftrightarrow$ (3) résultent de
Compléments d'algèbre, Lemme \ref{more-algebra-lemma-flat-descent-tor-amplitude}.
Il suffit donc de montrer que
(1) $\Rightarrow$ (3), (2) $\Rightarrow$ (3), et (4) $\Rightarrow$ (5).

\medskip\noindent
Supposons (1). En particulier, la cohomologie de $M^{sh}$ est nulle
en degrés $< a$ et $> b$. On peut donc représenter $M^{sh}$ par un complexe
$P^\bullet$ de $\mathcal{O}_{X, x}^{sh}$-modules libres tel que
$P^i = 0$ pour $i > b$
(voir, par exemple, le résultat très général
Catégories dérivées, Lemme \ref{derived-lemma-subcategory-left-resolution}).
Notons que $P^n$ est plat sur $\mathcal{O}_{S, s}$ pour tout $n$.
Considérons $\Coker(d_P^{a - 1})$. D'après Compléments d'algèbre, Lemme
\ref{more-algebra-lemma-last-one-flat},
c'est un $\mathcal{O}_{S, s}$-module plat.
Le Lemme \ref{flat-lemma-flat-up-down-henselization} montre donc que
c'est un $\mathcal{O}_{S, s}^{sh}$-module plat.
Ainsi $\tau_{\geq a}P^\bullet$ est un complexe de
$\mathcal{O}_{S, s}^{sh}$-modules plats représentant $M^{sh}$
dans $D(\mathcal{O}_{S, s}^{sh}$ et l'on en déduit que
$M^{sh}$ est de Tor-amplitude contenue dans $[a, b]$, voir
Compléments d'algèbre, Lemme \ref{more-algebra-lemma-tor-amplitude}.
Ainsi (1) implique (3). Cela entraîne aussi
(2) $\Rightarrow$ (3) en remplaçant $\mathcal{O}_{S, s}$ par
$\mathcal{O}_{S, s}^h$. Le même raisonnement démontre
(4) $\Rightarrow$ (5).
\end{proof}

\begin{lemma}
\label{flat-lemma-finite-flat-weak-assassin-up-down}
Soit $g : T \to S$ un morphisme fini et plat de schémas.
Soit $\mathcal{G}$ un $\mathcal{O}_S$-module quasi-cohérent.
Soit $t \in T$ un point d'image $s \in S$. Alors
$$
t \in \text{WeakAss}(g^*\mathcal{G})
\Leftrightarrow
s \in \text{WeakAss}(\mathcal{G})
$$
\end{lemma}

\begin{proof}
L'implication ``$\Leftarrow$'' résulte immédiatement de
Diviseurs, Lemme \ref{divisors-lemma-weakly-ass-pullback}.
Supposons $t \in \text{WeakAss}(g^*\mathcal{G})$.
Soit $\Spec(A) \subset S$ un voisinage ouvert affine de $s$.
Soit $\mathcal{G}$ le faisceau quasi-cohérent associé au $A$-module $M$.
Soit $\mathfrak p \subset A$ l'idéal premier correspondant à $s$.
Puisque $g$ est fini et plat, on a $g^{-1}(\Spec(A)) = \Spec(B)$
pour une certaine $A$-algèbre finie et plate $B$. Notons que
$g^*\mathcal{G}$ est le $\mathcal{O}_{\Spec(B)}$-module quasi-cohérent
associé au $B$-module $M \otimes_A B$ et que $g_*g^*\mathcal{G}$ est le
$\mathcal{O}_{\Spec(A)}$-module quasi-cohérent associé au
$A$-module $M \otimes_A B$. D'après
Algèbre, Lemme \ref{algebra-lemma-finite-flat-local},
on a $B_{\mathfrak p} \cong A_{\mathfrak p}^{\oplus n}$
pour un entier $n \geq 0$. Notons que $n \geq 1$, puisqu'on a supposé qu'il
existe au moins un point de $T$ au-dessus de $s$. En considérant
les fibres, on voit donc que
$$
s \in \text{WeakAss}(\mathcal{G})
\Leftrightarrow
s \in \text{WeakAss}(g_*g^*\mathcal{G})
$$
Or l'hypothèse $t \in \text{WeakAss}(g^*\mathcal{G})$
entraîne $s \in \text{WeakAss}(g_*g^*\mathcal{G})$ d'après
Diviseurs, Lemme \ref{divisors-lemma-weakly-associated-finite},
et donc, d'après ce qui précède, $s \in  \text{WeakAss}(\mathcal{G})$.
\end{proof}

\begin{lemma}
\label{flat-lemma-etale-weak-assassin-up-down}
Soit $h : U \to S$ un morphisme \'etale de schémas.
Soit $\mathcal{G}$ un $\mathcal{O}_S$-module quasi-cohérent.
Soit $u \in U$ un point d'image $s \in S$. Alors
$$
u \in \text{WeakAss}(h^*\mathcal{G})
\Leftrightarrow
s \in \text{WeakAss}(\mathcal{G})
$$
\end{lemma}

\begin{proof}
En remplaçant $S$ et $U$ par des voisinages affines de $s$ et $u$,
on peut supposer que $g$ est un morphisme \'etale standard de schémas affines, voir
Morphismes, Lemme \ref{morphisms-lemma-etale-locally-standard-etale}.
On peut donc supposer $S = \Spec(A)$ et
$X = \Spec(A[x, 1/g]/(f))$, où $f$ est unitaire et $f'$
est inversible dans $A[x, 1/g]$.
Notons que $A[x, 1/g]/(f) = (A[x]/(f))_g$ est aussi un localisé
de la $A$-algèbre finie et libre $A[x]/(f)$. On peut donc considérer
$U$ comme un sous-schéma ouvert du schéma $T = \Spec(A[x]/(f))$,
qui est fini localement libre sur $S$. On est ainsi ramené au
Lemme \ref{flat-lemma-finite-flat-weak-assassin-up-down}
ci-dessus.
\end{proof}

\begin{lemma}
\label{flat-lemma-weakly-associated-henselization}
Soient $S$ un schéma et $s \in S$ un point. Notons $\mathcal{O}_{S, s}^h$
(resp.\ $\mathcal{O}_{S, s}^{sh}$) le hensélisé (resp.\ le hensélisé
strict), voir
Algèbre, Définition \ref{algebra-definition-henselization}.
Soit $\mathcal{F}$ un $\mathcal{O}_S$-module quasi-cohérent.
Les conditions suivantes sont équivalentes :
\begin{enumerate}
\item $s$ est un point faiblement associé à $\mathcal{F}$,
\item $\mathfrak m_s$ est un idéal premier faiblement associé à $\mathcal{F}_s$,
\item $\mathfrak m_s^h$ est un idéal premier faiblement associé à
$\mathcal{F}_s \otimes_{\mathcal{O}_{S, s}} \mathcal{O}_{S, s}^h$, et
\item $\mathfrak m_s^{sh}$ est un idéal premier faiblement associé à
$\mathcal{F}_s \otimes_{\mathcal{O}_{S, s}} \mathcal{O}_{S, s}^{sh}$.
\end{enumerate}
\end{lemma}

\begin{proof}
L'équivalence de (1) et (2) est la définition, voir
Diviseurs, Définition \ref{divisors-definition-weakly-associated}.
Les implications (2) $\Rightarrow$ (3) $\Rightarrow$ (4)
résultent de Diviseurs, Lemme \ref{divisors-lemma-weakly-ass-pullback},
appliqué aux morphismes plats (Compléments d'algèbre, Lemme
\ref{more-algebra-lemma-dumb-properties-henselization})
suivants :
$$
\Spec(\mathcal{O}_{S, s}^{sh}) \to
\Spec(\mathcal{O}_{S, s}^h) \to
\Spec(\mathcal{O}_{S, s})
$$
et aux points fermés. Pour montrer (4) $\Rightarrow$ (2), on peut remplacer
$S$ par un voisinage affine. Supposons que
$x \in \mathcal{F}_s \otimes_{\mathcal{O}_{S, s}} \mathcal{O}_{S, s}^{sh}$
soit un élément dont l'annulateur a pour radical $\mathfrak m_s^{sh}$.
(Voir Algèbre, Lemme \ref{algebra-lemma-weakly-ass-local}.)
Puisque $\mathcal{O}_{S, s}^{sh}$ est la limite
des $\mathcal{O}_{U, u}$ sur les voisinages \'etales
$f : (U, u) \to (S, s)$ d'après Algèbre, Lemme
\ref{algebra-lemma-strict-henselization-different},
on peut supposer que $x$ est l'image d'un élément
$x' \in \mathcal{F}_s \otimes_{\mathcal{O}_{S, s}} \mathcal{O}_{U, u}$.
L'homomorphisme local $\mathcal{O}_{U, u} \to \mathcal{O}_{S, s}^{sh}$
est fidèlement plat (puisqu'il s'agit du hensélisé strict), donc
universellement injectif
(Algèbre, Lemme \ref{algebra-lemma-faithfully-flat-universally-injective}).
Il s'ensuit que l'annulateur de $x'$ est l'image réciproque de
l'annulateur de $x$. Le radical de cet annulateur
est donc égal à $\mathfrak m_u$.
Ainsi $u$ est un point faiblement associé à $f^*\mathcal{F}$.
D'après le Lemme \ref{flat-lemma-etale-weak-assassin-up-down},
on voit que $s$ est un point faiblement associé à $\mathcal{F}$.
\end{proof}





\section{Structure locale d'un module de type fini}
\label{flat-section-local-structure-module}

\noindent
Le lemme technique essentiel sur lequel reposent de nombreux raisonnements de ce
chapitre est le lemme géométrique
\ref{flat-lemma-elementary-devissage}.

\begin{lemma}
\label{flat-lemma-sheaf-lives-on-subscheme}
Soit $f : X \to S$ un morphisme de type fini de schémas affines.
Soit $\mathcal{F}$ un $\mathcal{O}_X$-module quasi-cohérent de type fini.
Soit $x \in X$ d'image $s = f(x)$ dans $S$.
Posons $\mathcal{F}_s = \mathcal{F}|_{X_s}$.
Il existe alors une immersion fermée $i : Z \to X$ de présentation finie
et un $\mathcal{O}_Z$-module quasi-cohérent de type fini $\mathcal{G}$
tels que $i_*\mathcal{G} = \mathcal{F}$ et
$Z_s = \text{Supp}(\mathcal{F}_s)$.
\end{lemma}

\begin{proof}
Supposons que le morphisme $f : X \to S$ soit donné par l'homomorphisme d'anneaux
$A \to B$ et que $\mathcal{F}$ soit le faisceau quasi-cohérent
associé au $B$-module $M$. D'après
Morphismes, Lemme \ref{morphisms-lemma-locally-finite-type-characterize},
on sait que $A \to B$ est un homomorphisme d'anneaux de type fini, et d'après
Propriétés, Lemme \ref{properties-lemma-finite-type-module},
on sait que $M$ est un $B$-module de type fini. En particulier, le
support de $\mathcal{F}$ est le sous-schéma fermé de $\Spec(B)$
défini par l'annulateur
$I = \{x \in B \mid xm = 0\ \forall m \in M\}$ de $M$, voir
Algèbre, Lemme \ref{algebra-lemma-support-closed}.
Soit $\mathfrak q \subset B$ l'idéal premier correspondant à $x$
et soit $\mathfrak p \subset A$ l'idéal premier correspondant à $s$.
Notons que $X_s = \Spec(B \otimes_A \kappa(\mathfrak p))$ et
que $\mathcal{F}_s$ est le faisceau quasi-cohérent associé au
$B \otimes_A \kappa(\mathfrak p)$-module $M \otimes_A \kappa(\mathfrak p)$. D'après
Morphismes, Lemme \ref{morphisms-lemma-support-finite-type},
le support de $\mathcal{F}_s$ est égal à
$V(I(B \otimes_A \kappa(\mathfrak p)))$. Puisque
$B \otimes_A \kappa(\mathfrak p)$ est de type fini sur $\kappa(\mathfrak p)$,
il existe un nombre fini d'éléments $f_1, \ldots, f_m \in I$
tels que
$$
I(B \otimes_A \kappa(\mathfrak p)) =
(f_1, \ldots, f_n)(B \otimes_A \kappa(\mathfrak p)).
$$
Notons $i : Z \to X$ le sous-schéma fermé défini par
$(f_1, \ldots, f_m)$, autrement dit $Z = \Spec(B/(f_1, \ldots, f_m))$.
Puisque $M$ est annulé par $I$, on peut considérer $M$ comme un
$B/(f_1, \ldots, f_m)$-module. En d'autres termes, $\mathcal{F}$ est
l'image directe d'un module de type fini sur $Z$.
Comme $Z_s = \text{Supp}(\mathcal{F}_s)$ par construction, cela
démontre le lemme.
\end{proof}

\begin{lemma}
\label{flat-lemma-elementary-devissage}
Soit $f : X \to S$ un morphisme de schémas localement de type fini.
Soit $\mathcal{F}$ un $\mathcal{O}_X$-module quasi-cohérent de type fini.
Soit $x \in X$ d'image $s = f(x)$ dans $S$.
Posons $\mathcal{F}_s = \mathcal{F}|_{X_s}$ et
$n = \dim_x(\text{Supp}(\mathcal{F}_s))$.
On peut alors construire
\begin{enumerate}
\item des voisinages \'etales élémentaires $g : (X', x') \to (X, x)$,
$e : (S', s') \to (S, s)$,
\item un diagramme commutatif
$$
\xymatrix{
X \ar[dd]_f & X' \ar[dd] \ar[l]^g & Z' \ar[l]^i \ar[d]^\pi \\
& & Y' \ar[d]^h \\
S & S' \ar[l]_e & S' \ar@{=}[l]
}
$$
\item un point $z' \in Z'$ tel que $i(z') = x'$, $y' = \pi(z')$, $h(y') = s'$,
\item un $\mathcal{O}_{Z'}$-module quasi-cohérent de type fini $\mathcal{G}$,
\end{enumerate}
de sorte que les propriétés suivantes soient vérifiées :
\begin{enumerate}
\item $X'$, $Z'$, $Y'$, $S'$ sont des schémas affines,
\item $i$ est une immersion fermée de présentation finie,
\item $i_*(\mathcal{G}) \cong g^*\mathcal{F}$,
\item $\pi$ est fini et $\pi^{-1}(\{y'\}) = \{z'\}$,
\item l'extension $\kappa(y')/\kappa(s')$ est transcendante pure,
\item $h$ est lisse de dimension relative $n$
et à fibres géométriquement intègres.
\end{enumerate}
\end{lemma}

\begin{proof}
Soit $V \subset S$ un voisinage affine de $s$.
Soit $U \subset f^{-1}(V)$ un voisinage affine de $x$.
Il suffit alors de démontrer le lemme pour $f|_U : U \to V$ et
$\mathcal{F}|_U$. Dans la suite de la démonstration, on suppose donc que
$X$ et $S$ sont affines.

\medskip\noindent
Supposons d'abord que $X_s = \text{Supp}(\mathcal{F}_s)$ ; en particulier,
$n = \dim_x(X_s)$. Appliquons
Compléments sur les morphismes,
Lemmes \ref{more-morphisms-lemma-local-local-structure-finite-type} et
\ref{more-morphisms-lemma-local-local-structure-finite-type-affine}.
On obtient un diagramme commutatif
$$
\xymatrix{
X \ar[dd] & X' \ar[l]^g \ar[d]^\pi \\
& Y' \ar[d]^h  \\
S & S' \ar[l]_e
}
$$
et un point $x' \in X'$. On pose $Z' = X'$, $i = \text{id}$ et
$\mathcal{G} = g^*\mathcal{F}$, ce qui donne une solution dans ce cas.

\medskip\noindent
Dans le cas général, choisissons une immersion fermée $Z \to X$ et un faisceau
$\mathcal{G}$ sur $Z$ comme dans
le Lemme \ref{flat-lemma-sheaf-lives-on-subscheme}.
En appliquant le résultat du paragraphe précédent à $Z \to S$ et
$\mathcal{G}$, on obtient un diagramme
$$
\xymatrix{
X \ar[dd]_f & Z \ar[l] \ar[dd]_{f|_Z} & Z' \ar[l]^g \ar[d]^\pi \\
& & Y' \ar[d]^h \\
S & S \ar@{=}[l] & S' \ar[l]_e
}
$$
et un point $z' \in Z'$ vérifiant toutes les propriétés requises.
Nous allons utiliser
le Lemme \ref{flat-lemma-lift-etale}
pour plonger $Z'$ dans un schéma \'etale sur $X$. On ne peut appliquer directement le lemme,
car on veut que $X'$ soit un schéma sur $S'$. On
considère donc les morphismes
$$
\xymatrix{
Z' \ar[r] & Z \times_S S' \ar[r] & X \times_S S'
}
$$
Le premier morphisme est \'etale d'après
Morphismes, Lemme \ref{morphisms-lemma-etale-permanence}.
Le second est une immersion fermée, puisqu'il s'obtient par changement de base d'une immersion fermée.
Enfin, puisque $X$, $S$, $S'$, $Z$, $Z'$ sont tous affines, on peut appliquer
le Lemme \ref{flat-lemma-lift-etale}
pour obtenir un morphisme \'etale de schémas affines $X' \to X \times_S S'$ tel que
$$
Z' = (Z \times_S S') \times_{(X \times_S S')} X' = Z \times_X X'.
$$
Puisque $Z \to X$ est une immersion fermée de présentation finie, il en est de même de $Z' \to X'$.
Soit $x' \in X'$ le point correspondant à $z' \in Z'$.
Le diagramme complété
$$
\xymatrix{
X \ar[dd] & X' \ar[dd] \ar[l] & Z' \ar[l]^i \ar[d]^\pi \\
& & Y' \ar[d]^h \\
S & S' \ar[l]_e & S' \ar@{=}[l]
}
$$
fournit alors une solution du problème initial.
\end{proof}

\begin{lemma}
\label{flat-lemma-devissage-finite-presentation}
Sous les hypothèses et avec les notations du
Lemme \ref{flat-lemma-elementary-devissage},
si $f$ est localement de présentation finie,
alors $\pi$ est de présentation finie.
Dans ce cas, les conditions suivantes sont équivalentes :
\begin{enumerate}
\item $\mathcal{F}$ est un $\mathcal{O}_X$-module de présentation finie
au voisinage de $x$,
\item $\mathcal{G}$ est un $\mathcal{O}_{Z'}$-module de présentation finie
au voisinage de $z'$, et
\item $\pi_*\mathcal{G}$ est un $\mathcal{O}_{Y'}$-module
de présentation finie au voisinage de $y'$.
\end{enumerate}
Toujours sous l'hypothèse que $f$ est localement de présentation finie, les conditions
suivantes sont équivalentes :
\begin{enumerate}
\item[(a)] $\mathcal{F}_x$ est un $\mathcal{O}_{X, x}$-module de présentation
finie,
\item[(b)] $\mathcal{G}_{z'}$ est un $\mathcal{O}_{Z', z'}$-module
de présentation finie, et
\item[(c)] $(\pi_*\mathcal{G})_{y'}$ est un $\mathcal{O}_{Y', y'}$-module
de présentation finie.
\end{enumerate}
\end{lemma}

\begin{proof}
Supposons $f$ localement de présentation finie. Alors $Z' \to S$ est localement
de présentation finie, comme composé de tels morphismes, voir
Morphismes, Lemme \ref{morphisms-lemma-composition-finite-presentation}.
Notons que $Y' \to S$ est aussi localement de présentation finie, comme composé
d'un morphisme lisse et d'un morphisme \'etale. Ainsi
Morphismes, Lemme \ref{morphisms-lemma-finite-presentation-permanence},
entraîne que $\pi$ est localement de présentation finie.
Puisque $\pi$ est fini, il est aussi séparé et
quasi-compact ; par suite, $\pi$ est bien de présentation finie.

\medskip\noindent
Pour démontrer l'équivalence de (1), (2) et (3), considérons aussi :
(4) $g^*\mathcal{F}$ est un $\mathcal{O}_{X'}$-module de présentation finie
au voisinage de $x'$. L'image réciproque d'un module de présentation finie
est de présentation finie, voir
Modules, Lemme \ref{modules-lemma-pullback-finite-presentation}.
Ainsi (1) $\Rightarrow$ (4).
Le morphisme \'etale $g$ est ouvert, voir
Morphismes, Lemme \ref{morphisms-lemma-etale-open}.
Pour tout voisinage ouvert $U' \subset X'$ de $x'$, l'image
$g(U')$ est donc un voisinage ouvert de $x$ et
$\{U' \to g(U')\}$ est un recouvrement \'etale. Ainsi (4) $\Rightarrow$ (1) d'après
Descente, Lemme \ref{descent-lemma-finite-presentation-descends}.
En utilisant
Descente, Lemme \ref{descent-lemma-finite-finitely-presented-module},
et quelques arguments topologiques élémentaires (voir
Compléments sur les morphismes,
Lemme \ref{more-morphisms-lemma-finite-morphism-single-point-in-fibre}),
on voit que
(4) $\Leftrightarrow$ (2) $\Leftrightarrow$ (3).

\medskip\noindent
Pour démontrer l'équivalence de (a), (b), (c), considérons les homomorphismes d'anneaux
$$
\mathcal{O}_{X, x} \to
\mathcal{O}_{X', x'} \to
\mathcal{O}_{Z', z'} \leftarrow
\mathcal{O}_{Y', y'}
$$
Le premier homomorphisme est fidèlement plat. Par suite,
$\mathcal{F}_x$ est de présentation finie sur $\mathcal{O}_{X, x}$
si et seulement si $g^*\mathcal{F}_{x'}$ est de présentation finie sur
$\mathcal{O}_{X', x'}$, voir
Algèbre, Lemme \ref{algebra-lemma-descend-properties-modules}.
Le deuxième homomorphisme est surjectif (donc fini) et
de présentation finie par hypothèse ; par suite,
$g^*\mathcal{F}_{x'}$ est de présentation finie sur $\mathcal{O}_{X', x'}$
si et seulement si $\mathcal{G}_{z'}$ est de présentation finie sur
$\mathcal{O}_{Z', z'}$, voir
Algèbre, Lemme \ref{algebra-lemma-finite-finitely-presented-extension}.
Puisque $\pi$ est fini, de présentation finie et que
$\pi^{-1}(\{y'\}) = \{x'\}$, l'homomorphisme d'anneaux
$\mathcal{O}_{Y', y'} \leftarrow \mathcal{O}_{Z', z'}$ est fini
et de présentation finie, voir
Compléments sur les morphismes,
Lemme \ref{more-morphisms-lemma-finite-morphism-single-point-in-fibre}.
Par suite, $\mathcal{G}_{z'}$ est de présentation finie sur $\mathcal{O}_{Z', z'}$
si et seulement si $\pi_*\mathcal{G}_{y'}$ est de présentation finie sur
$\mathcal{O}_{Y', y'}$, voir
Algèbre, Lemme \ref{algebra-lemma-finite-finitely-presented-extension}.
\end{proof}

\begin{lemma}
\label{flat-lemma-devissage-flat}
Sous les hypothèses et avec les notations du
Lemme \ref{flat-lemma-elementary-devissage},
les conditions suivantes sont équivalentes :
\begin{enumerate}
\item $\mathcal{F}$ est plat sur $S$ au voisinage de $x$,
\item $\mathcal{G}$ est plat sur $S'$ au voisinage de $z'$, et
\item $\pi_*\mathcal{G}$ est plat sur $S'$ au voisinage de $y'$.
\end{enumerate}
Les conditions suivantes sont aussi équivalentes :
\begin{enumerate}
\item[(a)] $\mathcal{F}_x$ est plat sur $\mathcal{O}_{S, s}$,
\item[(b)] $\mathcal{G}_{z'}$ est plat sur $\mathcal{O}_{S', s'}$, et
\item[(c)] $(\pi_*\mathcal{G})_{y'}$ est plat sur $\mathcal{O}_{S', s'}$.
\end{enumerate}
\end{lemma}

\begin{proof}
Pour démontrer l'équivalence de (1), (2) et (3), considérons aussi :
(4) $g^*\mathcal{F}$ est plat sur $S$ au voisinage de $x'$.
Nous utiliserons
le Lemme \ref{flat-lemma-etale-flat-up-down}
pour identifier la platitude sur $S$ et sur $S'$ sans le mentionner de nouveau.
Le morphisme \'etale $g$ est plat et ouvert, voir
Morphismes, Lemme \ref{morphisms-lemma-etale-open}.
Pour tout voisinage ouvert $U' \subset X'$ de $x'$, l'image
$g(U')$ est donc un voisinage ouvert de $x$ et le morphisme
$U' \to g(U')$ est surjectif et plat.
Ainsi (4) $\Leftrightarrow$ (1) d'après
Morphismes, Lemme \ref{morphisms-lemma-flat-permanence}.
Notons que
$$
\Gamma(X', g^*\mathcal{F}) =
\Gamma(Z', \mathcal{G}) =
\Gamma(Y', \pi_*\mathcal{G})
$$
Les conditions de platitude de $g^*\mathcal{F}$, de $\mathcal{G}$ et de $\pi_*\mathcal{G}$
sur $S'$ sont donc équivalentes (on utilise ici le fait que $X'$, $Z'$, $Y'$ et
$S'$ sont tous affines). Des arguments topologiques laissés au lecteur (comparer
Compléments sur les morphismes,
Lemme \ref{more-morphisms-lemma-finite-morphism-single-point-in-fibre})
portant sur les voisinages affines montrent alors que
(4) $\Leftrightarrow$ (2) $\Leftrightarrow$ (3).

\medskip\noindent
Pour démontrer l'équivalence de (a), (b), (c), considérons le diagramme commutatif
d'homomorphismes locaux
$$
\xymatrix{
\mathcal{O}_{X', x'} \ar[r]_\iota &
\mathcal{O}_{Z', z'} &
\mathcal{O}_{Y', y'} \ar[l]^\alpha &
\mathcal{O}_{S', s'} \ar[l]^\beta \\
\mathcal{O}_{X, x} \ar[u]^\gamma & & &
\mathcal{O}_{S, s} \ar[lll]_\varphi \ar[u]_\epsilon
}
$$
Nous utiliserons
le Lemme \ref{flat-lemma-etale-flat-up-down-local-ring}
pour identifier la platitude sur $\mathcal{O}_{S, s}$ et sur $\mathcal{O}_{S', s'}$
sans le mentionner de nouveau.
L'homomorphisme $\gamma$ est fidèlement plat. Par suite,
$\mathcal{F}_x$ est plat sur $\mathcal{O}_{S, s}$
si et seulement si $g^*\mathcal{F}_{x'}$ est plat sur
$\mathcal{O}_{S', s'}$, voir
Algèbre, Lemme \ref{algebra-lemma-flatness-descends-more-general}.
Comme $\mathcal{O}_{S', s'}$-modules, les modules
$g^*\mathcal{F}_{x'}$, $\mathcal{G}_{z'}$ et
$\pi_*\mathcal{G}_{y'}$ sont tous isomorphes, voir
Compléments sur les morphismes,
Lemme \ref{more-morphisms-lemma-finite-morphism-single-point-in-fibre}.
Cela achève la démonstration.
\end{proof}









\section{D\'evissage en une étape}
\label{flat-section-one-step-devissage}

\noindent
Dans cette section, nous expliquons ce qu'est un d\'evissage en une étape d'un
module. Un d\'evissage en une étape existe localement pour la topologie \'etale sur la base et le but.
Nous étudions le changement de base, la restriction à des ouverts de Zariski et la localisation \'etale
d'un d\'evissage en une étape.

\begin{definition}
\label{flat-definition-one-step-devissage}
Soit $S$ un schéma.
Soit $X$ localement de type fini sur $S$.
Soit $\mathcal{F}$ un $\mathcal{O}_X$-module quasi-cohérent de type fini.
Soit $s \in S$ un point.
Un {\it d\'evissage en une étape de $\mathcal{F}/X/S$ au-dessus de $s$}
est donné par des morphismes de schémas sur $S$
$$
\xymatrix{
X & Z \ar[l]_i \ar[r]^\pi & Y
}
$$
et un $\mathcal{O}_Z$-module quasi-cohérent $\mathcal{G}$ de type fini
tels que
\begin{enumerate}
\item $X$, $S$, $Z$ et $Y$ soient affines,
\item $i$ soit une immersion fermée de présentation finie,
\item $\mathcal{F} \cong i_*\mathcal{G}$,
\item $\pi$ soit fini, et
\item le morphisme structural $Y \to S$ soit lisse à
fibres géométriquement irréductibles de
dimension $\dim(\text{Supp}(\mathcal{F}_s))$.
\end{enumerate}
On dit alors que $(Z, Y, i, \pi, \mathcal{G})$ est un
d\'evissage en une étape de $\mathcal{F}/X/S$ au-dessus de $s$.
\end{definition}

\noindent
Un tel d\'evissage en une étape ne peut exister que si $X$ et $S$
sont affines. Dans la définition ci-dessus, on exige seulement que $X$ soit
(localement) de type fini sur $S$, et l'on conserve ce cadre
dans la suite. Dans \cite{GruRay}, les auteurs se placent systématiquement dans le cas
où $X \to S$ est localement de présentation finie et où $\mathcal{F}$
est un $\mathcal{O}_X$-module quasi-cohérent de type fini. Ce choix présente
l'avantage de donner un sens à la condition que $\mathcal{F}$ soit
de présentation finie comme $\mathcal{O}_X$-module, tandis que, dans notre
cadre, cette condition n'a pas de sens intrinsèque. Voir
Compléments sur les morphismes, section
\ref{more-morphisms-section-finite-type-finite-presentation},
pour une discussion ; les observations qui y sont faites montrent que, dans notre cadre,
on peut considérer la condition que $\mathcal{F}$ soit « localement de présentation
finie relativement à $S$ », et l'on pourrait employer systématiquement cette
notion. Nous préférons cependant nous appuyer sur les résultats du
Lemme \ref{flat-lemma-devissage-finite-presentation}
et les observations de la
Remarque \ref{flat-remark-finite-presentation}
pour traiter cette question au cas par cas lorsqu'elle se présente.

\begin{definition}
\label{flat-definition-one-step-devissage-at-x}
Soit $S$ un schéma.
Soit $X$ localement de type fini sur $S$.
Soit $\mathcal{F}$ un $\mathcal{O}_X$-module quasi-cohérent de type fini.
Soit $x \in X$ un point d'image $s$ dans $S$.
Un {\it d\'evissage en une étape de $\mathcal{F}/X/S$ en $x$}
est un système $(Z, Y, i, \pi, \mathcal{G}, z, y)$, où
$(Z, Y, i, \pi, \mathcal{G})$ est un d\'evissage en une étape de
$\mathcal{F}/X/S$ au-dessus de $s$ et où
\begin{enumerate}
\item $\dim_x(\text{Supp}(\mathcal{F}_s)) = \dim(\text{Supp}(\mathcal{F}_s))$,
\item $z \in Z$ est un point tel que $i(z) = x$ et $\pi(z) = y$,
\item on a $\pi^{-1}(\{y\}) = \{z\}$,
\item l'extension $\kappa(y)/\kappa(s)$ est transcendante
pure.
\end{enumerate}
\end{definition}

\noindent
Un d\'evissage en une étape de $\mathcal{F}/X/S$ en $x$ ne peut exister que si
$X$ et $S$ sont affines. La condition (1) assure que $Y \to S$ est de
dimension relative égale à $\dim_x(\text{Supp}(\mathcal{F}_s))$,
en vertu de la condition (5) de la
Définition \ref{flat-definition-one-step-devissage}.

\begin{lemma}
\label{flat-lemma-elementary-devissage-variant}
Soit $f : X \to S$ un morphisme de schémas localement de type fini.
Soit $\mathcal{F}$ un $\mathcal{O}_X$-module quasi-cohérent de type fini.
Soit $x \in X$ d'image $s = f(x)$ dans $S$.
Il existe alors un diagramme commutatif de schémas pointés
$$
\xymatrix{
(X, x) \ar[d]_f & (X', x') \ar[l]^g \ar[d] \\
(S, s) & (S', s') \ar[l] \\
}
$$
tel que $(S', s') \to (S, s)$ et $(X', x') \to (X, x)$
soient des voisinages \'etales élémentaires et tel que
$g^*\mathcal{F}/X'/S'$ admette un d\'evissage en une étape en $x'$.
\end{lemma}

\begin{proof}
Cela résulte immédiatement de la
Définition \ref{flat-definition-one-step-devissage-at-x}
et du
Lemme \ref{flat-lemma-elementary-devissage}.
\end{proof}

\begin{lemma}
\label{flat-lemma-base-change-one-step}
Soient $S$, $X$, $\mathcal{F}$, $s$ comme dans la
Définition \ref{flat-definition-one-step-devissage}.
Soit $(Z, Y, i, \pi, \mathcal{G})$ un d\'evissage en une étape
de $\mathcal{F}/X/S$ au-dessus de $s$.
Soit $(S', s') \to (S, s)$ un morphisme quelconque de schémas pointés.
Ces données étant fixées, soient $X', Z', Y', i', \pi'$ les objets obtenus par changement
de base à partir de $X, Z, Y, i, \pi$ suivant $S' \to S$.
Soit $\mathcal{F}'$ l'image réciproque de $\mathcal{F}$ sur $X'$
et soit $\mathcal{G}'$ l'image réciproque de $\mathcal{G}$ sur $Z'$.
Si $S'$ est affine, alors $(Z', Y', i', \pi', \mathcal{G}')$
est un d\'evissage en une étape de $\mathcal{F}'/X'/S'$ au-dessus de $s'$.
\end{lemma}

\begin{proof}
Les produits fibrés de schémas affines sont affines, voir
Schémas, Lemme \ref{schemes-lemma-fibre-product-affines}.
Le changement de base préserve
les immersions fermées,
les morphismes de présentation finie,
les morphismes finis,
les morphismes lisses,
les morphismes à fibres géométriquement irréductibles et
les morphismes de dimension relative $n$, voir
Morphismes, Lemmes \ref{morphisms-lemma-base-change-closed-immersion},
\ref{morphisms-lemma-base-change-finite-presentation},
\ref{morphisms-lemma-base-change-finite},
\ref{morphisms-lemma-base-change-smooth},
\ref{morphisms-lemma-base-change-relative-dimension-d}, et
Compléments sur les morphismes, Lemme
\ref{more-morphisms-lemma-base-change-fibres-geometrically-irreducible}.
On a $i'_*\mathcal{G}' \cong \mathcal{F}'$, car l'image directe
par le morphisme fini $i$ commute au changement de base, voir
Cohomologie des schémas, Lemme \ref{coherent-lemma-affine-base-change}.
On a
$\dim(\text{Supp}(\mathcal{F}_s)) = \dim(\text{Supp}(\mathcal{F}'_{s'}))$
d'après
Morphismes, Lemme \ref{morphisms-lemma-dimension-fibre-after-base-change},
puisque
$$
\text{Supp}(\mathcal{F}_s) \times_s s' = \text{Supp}(\mathcal{F}'_{s'}).
$$
Cela démontre le lemme.
\end{proof}

\begin{lemma}
\label{flat-lemma-base-change-one-step-at-x}
Soient $S$, $X$, $\mathcal{F}$, $x$, $s$ comme dans la
Définition \ref{flat-definition-one-step-devissage-at-x}.
Soit $(Z, Y, i, \pi, \mathcal{G}, z, y)$ un d\'evissage en une étape
de $\mathcal{F}/X/S$ en $x$.
Soit $(S', s') \to (S, s)$ un morphisme de schémas pointés
induisant un isomorphisme $\kappa(s) = \kappa(s')$.
Soit $(Z', Y', i', \pi', \mathcal{G}')$ le système construit dans le
Lemme \ref{flat-lemma-base-change-one-step},
et soit $x' \in X'$ (resp.\ $z' \in Z'$, $y' \in Y'$) l'unique
point s'envoyant à la fois sur $x \in X$ (resp.\ $z \in Z$, $y \in Y$)
et sur $s' \in S'$.
Si $S'$ est affine, alors $(Z', Y', i', \pi', \mathcal{G}', z', y')$
est un d\'evissage en une étape de $\mathcal{F}'/X'/S'$ en $x'$.
\end{lemma}

\begin{proof}
D'après le
Lemme \ref{flat-lemma-base-change-one-step},
$(Z', Y', i', \pi', \mathcal{G}')$ est un d\'evissage en une étape de
$\mathcal{F}'/X'/S'$ au-dessus de $s'$. Les propriétés (1) -- (4) de la
Définition \ref{flat-definition-one-step-devissage-at-x}
sont vérifiées pour $(Z', Y', i', \pi', \mathcal{G}', z', y')$,
car l'hypothèse $\kappa(s) = \kappa(s')$ assure que les fibres
$X'_{s'}$, $Z'_{s'}$ et $Y'_{s'}$ sont respectivement isomorphes à
$X_s$, $Z_s$ et $Y_s$.
\end{proof}

\begin{definition}
\label{flat-definition-shrink}
Soient $S$, $X$, $\mathcal{F}$, $x$, $s$ comme dans la
Définition \ref{flat-definition-one-step-devissage-at-x}.
Soit $(Z, Y, i, \pi, \mathcal{G}, z, y)$ un d\'evissage en une étape
de $\mathcal{F}/X/S$ en $x$. Appelons
{\it restriction standard} de cette situation la donnée
d'ouverts standards $S' \subset S$, $X' \subset X$, $Z' \subset Z$
et $Y' \subset Y$ tels que $s \in S'$, $x \in X'$, $z \in Z'$ et
$y \in Y'$, et tels que
$$
(Z', Y', i|_{Z'}, \pi|_{Z'}, \mathcal{G}|_{Z'}, z, y)
$$
soit un d\'evissage en une étape de $\mathcal{F}|_{X'}/X'/S'$ en $x$.
\end{definition}

\begin{lemma}
\label{flat-lemma-shrink}
Sous les hypothèses et avec les notations de la
Définition \ref{flat-definition-shrink},
on a :
\begin{enumerate}
\item
\label{flat-item-shrink-base}
Si $S' \subset S$ est un voisinage ouvert standard de $s$, alors,
en posant $X' = X_{S'}$, $Z' = Z_{S'}$ et $Y' = Y_{S'}$, on obtient une
restriction standard.
\item
\label{flat-item-shrink-on-Y}
Soit $W \subset Y$ un voisinage ouvert standard de $y$.
Il existe alors une restriction standard telle que $Y' = W \times_S S'$.
\item
\label{flat-item-shrink-on-X}
Soit $U \subset X$ un voisinage ouvert de $x$.
Il existe alors une restriction standard telle que $X' \subset U$.
\end{enumerate}
\end{lemma}

\begin{proof}
L'assertion (1) résulte immédiatement du
Lemme \ref{flat-lemma-base-change-one-step-at-x}
et du fait que l'image réciproque d'un ouvert standard par un morphisme
de schémas affines est un ouvert standard, voir
Algèbre, Lemme \ref{algebra-lemma-spec-functorial}.

\medskip\noindent
Soit $W \subset Y$ comme dans (2). Puisque $Y \to S$ est lisse, il est ouvert, voir
Morphismes, Lemme \ref{morphisms-lemma-smooth-open}.
On peut donc trouver un voisinage ouvert standard $S'$ de $s$
contenu dans l'image de $W$. Les fibres de $W_{S'} \to S'$
sont alors des sous-schémas ouverts non vides des fibres de $Y \to S$ au-dessus de $S'$
et sont donc elles aussi géométriquement irréductibles. En posant $Y' = W_{S'}$
et $Z' = \pi^{-1}(Y')$, on voit que $Z' \subset Z$ est un voisinage ouvert standard
de $z$. Soit $\overline{h} \in \Gamma(Z, \mathcal{O}_Z)$
une fonction telle que $Z' = D(\overline{h})$. Puisque $i : Z \to X$
est une immersion fermée, on peut trouver une fonction $h \in \Gamma(X, \mathcal{O}_X)$
telle que $i^\sharp(h) = \overline{h}$. Prenons $X' = D(h) \subset X$.
On obtient ainsi une restriction standard comme dans (2).

\medskip\noindent
Soit $U \subset X$ comme dans (3). Quitte à restreindre $U$, on peut supposer que
$U$ est un ouvert standard. D'après
Compléments sur les morphismes,
Lemme \ref{more-morphisms-lemma-finite-morphism-single-point-in-fibre},
il existe un ouvert standard $W \subset Y$, voisinage de $y$, tel
que $\pi^{-1}(W) \subset i^{-1}(U)$. Appliquons (2) pour obtenir une restriction
standard $X', S', Z', Y'$ avec $Y' = W_{S'}$. Puisque
$Z' \subset \pi^{-1}(W) \subset i^{-1}(U)$, on peut remplacer $X'$ par
$X' \cap U$ (qui est encore un ouvert standard, puisque $U$ l'est aussi)
sans modifier aucune des conditions définissant une restriction standard.
On obtient ainsi le résultat voulu.
\end{proof}

\begin{lemma}
\label{flat-lemma-elementary-etale-neighbourhood}
Soient $S$, $X$, $\mathcal{F}$, $x$, $s$ comme dans la
Définition \ref{flat-definition-one-step-devissage-at-x}.
Soit $(Z, Y, i, \pi, \mathcal{G}, z, y)$ un d\'evissage en une étape
de $\mathcal{F}/X/S$ en $x$. Soit
$$
\xymatrix{
(Y, y) \ar[d] & (Y', y') \ar[l] \ar[d] \\
(S, s) & (S', s') \ar[l]
}
$$
un diagramme commutatif de schémas pointés dont les flèches horizontales
sont des voisinages \'etales élémentaires. Il existe alors
un diagramme commutatif
$$
\xymatrix{
& & (X'', x'') \ar[lld] \ar[d] & (Z'', z'') \ar[l] \ar[lld] \ar[d] \\
(X, x) \ar[d] & (Z, z) \ar[l] \ar[d] &
(S'', s'') \ar[lld] & (Y'', y'') \ar[lld] \ar[l] \\
(S, s) & (Y, y) \ar[l]
}
$$
de schémas pointés vérifiant les propriétés suivantes :
\begin{enumerate}
\item $(S'', s'') \to (S', s')$ est un voisinage \'etale élémentaire et
le morphisme $S'' \to S$ est le composé $S'' \to S' \to S$,
\item $Y''$ est un sous-schéma ouvert de $Y' \times_{S'} S''$,
\item $Z'' = Z \times_Y Y''$,
\item $(X'', x'') \to (X, x)$ est un voisinage \'etale élémentaire, et
\item $(Z'', Y'', i'', \pi'', \mathcal{G}'', z'', y'')$ est un
d\'evissage en une étape en $x''$ du faisceau $\mathcal{F}''$.
\end{enumerate}
Ici $\mathcal{F}''$ (resp.\ $\mathcal{G}''$) est l'image réciproque de
$\mathcal{F}$ (resp.\ $\mathcal{G}$) par le morphisme $X'' \to X$
(resp.\ $Z'' \to Z$), et $i'' : Z'' \to X''$ et $\pi'' : Z'' \to Y''$
sont les morphismes du diagramme.
\end{lemma}

\begin{proof}
Soit $(S'', s'') \to (S', s')$ un voisinage \'etale élémentaire quelconque
avec $S''$ affine. Soit $Y'' \subset Y' \times_{S'} S''$ un voisinage
ouvert affine quelconque du point $y'' = (y', s'')$. On
obtient alors un schéma affine $(Z'', z'')$ par (3). De plus, $Z_{S''} \to X_{S''}$
est une immersion fermée et $Z'' \to Z_{S''}$ est un morphisme \'etale.
Le
Lemme \ref{flat-lemma-lift-etale}
s'applique donc, et l'on peut trouver un morphisme \'etale $X'' \to X_{S'}$ de schémas affines
tel que $Z'' \cong X'' \times_{X_{S'}} Z_{S'}$. Notons $i'' : Z'' \to X''$
l'immersion fermée correspondante. En posant $x'' = i''(z'')$, on obtient un
diagramme commutatif comme dans le lemme.
Les propriétés (1), (2), (3) et (4) sont vérifiées par construction.
Il suffit donc de montrer que (5) est vérifiée pour un choix convenable de
$(S'', s'') \to (S', s')$ et de $Y''$.

\medskip\noindent
Énumérons d'abord les propriétés vérifiées pour tout choix de
$(S'', s'') \to (S', s')$ et de $Y''$ comme dans le premier paragraphe.
Puisqu'on a $Z'' = X'' \times_X Z$ par construction, on voit que
$i''_*\mathcal{G}'' = \mathcal{F}''$ (avec les notations de
l'énoncé du lemme), voir
Cohomologie des schémas, Lemme \ref{coherent-lemma-affine-base-change}.
Posons $n = \dim(\text{Supp}(\mathcal{F}_s)) = \dim_x(\text{Supp}(\mathcal{F}_s))$.
Le morphisme $Y'' \to S''$ est lisse de dimension relative $n$
(car $Y' \to S'$ est lisse de dimension relative $n$,
comme composé $Y' \to Y_{S'} \to S'$ d'un morphisme \'etale et
d'un morphisme lisse de dimension relative $n$, et parce que le changement de base
préserve les morphismes lisses de dimension relative $n$).
On a $\kappa(y'') = \kappa(y)$ et $\kappa(s) = \kappa(s'')$ ;
ainsi $\kappa(y'')$ est une extension transcendante pure de $\kappa(s'')$.
Le morphisme de fibres $X''_{s''} \to X_s$ est un morphisme \'etale de schémas
affines sur $\kappa(s) = \kappa(s'')$, envoyant le point $x''$ sur le
point $x$ et dont l'image réciproque de $\mathcal{F}_s$ est $\mathcal{F}''_{s''}$.
Par suite,
$$
\dim(\text{Supp}(\mathcal{F}''_{s''})) =
\dim(\text{Supp}(\mathcal{F}_s)) = n =
\dim_x(\text{Supp}(\mathcal{F}_s)) =
\dim_{x''}(\text{Supp}(\mathcal{F}''_{s''}))
$$
car la dimension est invariante par localisation \'etale, voir
Descente, Lemme \ref{descent-lemma-dimension-at-point-local}.
Puisque $\pi'' : Z'' \to Y''$ est obtenu par changement de base à partir de $\pi$, on voit que
$\pi''$ est fini et, puisque $\kappa(y) = \kappa(y'')$, on voit que
$\pi^{-1}(\{y''\}) = \{z''\}$.

\medskip\noindent
À ce stade, on a vérifié toutes les conditions de la
Définition \ref{flat-definition-one-step-devissage},
sauf que l'on n'a pas vérifié que $Y'' \to S''$ est à fibres géométriquement
irréductibles. Cela n'est évidemment pas vrai en général,
et c'est ici que l'on utilise le fait que
$\kappa(s) \subset \kappa(y)$ est une extension transcendante pure. En effet,
soit $T \subset Y'_{s'}$ la composante irréductible de
$Y'_{s'}$ contenant $y' = (y, s')$. Notons que $T$ est un sous-schéma ouvert
de $Y'_{s'}$, car ce dernier est un schéma lisse sur $\kappa(s')$. D'après
Variétés,
Lemme \ref{varieties-lemma-geometrically-connected-if-connected-and-point},
on voit que $T$ est géométriquement connexe, puisque $\kappa(s') = \kappa(s)$
est algébriquement fermé dans $\kappa(y') = \kappa(y)$.
Puisque $T$ est lisse, on voit que $T$ est géométriquement irréductible. Le
Lemme de Compléments sur les morphismes,
\ref{more-morphisms-lemma-normal-morphism-irreducible},
s'applique donc, et l'on peut trouver un morphisme \'etale élémentaire
$(S'', s'') \to (S', s')$ et un ouvert affine $Y'' \subset Y'_{S''}$
tels que toutes les fibres de $Y'' \to S''$ soient géométriquement irréductibles
et que $T = Y''_{s''}$. Quitte à restreindre d'abord $Y''$, puis $S''$,
on peut supposer que $Y''$ et $S''$ sont tous deux affines.
Cela achève la démonstration du lemme.
\end{proof}

\begin{lemma}
\label{flat-lemma-existence-alpha}
Soient $S$, $X$, $\mathcal{F}$, $s$ comme dans la
Définition \ref{flat-definition-one-step-devissage}.
Soit $(Z, Y, i, \pi, \mathcal{G})$ un d\'evissage en une étape
de $\mathcal{F}/X/S$ au-dessus de $s$.
Soit $\xi \in Y_s$ l'unique point générique.
Il existe alors un entier $r > 0$ et un homomorphisme de $\mathcal{O}_Y$-modules
$$
\alpha : \mathcal{O}_Y^{\oplus r} \longrightarrow \pi_*\mathcal{G}
$$
tels que
$$
\alpha :
\kappa(\xi)^{\oplus r}
\longrightarrow
(\pi_*\mathcal{G})_\xi \otimes_{\mathcal{O}_{Y, \xi}} \kappa(\xi)
$$
soit un isomorphisme. De plus, on a alors
$$
\dim(\text{Supp}(\Coker(\alpha)_s)) < \dim(\text{Supp}(\mathcal{F}_s)).
$$
\end{lemma}

\begin{proof}
Par hypothèse, les schémas $S$ et $Y$ sont affines.
Écrivons $S = \Spec(A)$ et $Y = \Spec(B)$.
Puisque $\pi$ est fini, le $\mathcal{O}_Y$-module $\pi_*\mathcal{G}$
est un $\mathcal{O}_Y$-module quasi-cohérent de type fini.
On a donc $\pi_*\mathcal{G} = \widetilde{N}$ pour un certain $B$-module de type fini $N$.
Soit $\mathfrak p \subset B$ l'idéal premier correspondant à $\xi$.
Pour obtenir $\alpha$, posons
$r = \dim_{\kappa(\mathfrak p)} N \otimes_B \kappa(\mathfrak p)$
et choisissons $x_1, \ldots, x_r \in N$ dont les images forment une base de
$N \otimes_B \kappa(\mathfrak p)$. Prenons pour $\alpha : B^{\oplus r} \to N$
l'homomorphisme donné par la formule $\alpha(b_1, \ldots, b_r) = \sum b_ix_i$.
Il est clair que
$\alpha : \kappa(\mathfrak p)^{\oplus r} \to N \otimes_B \kappa(\mathfrak p)$
est un isomorphisme, comme souhaité. Enfin, soit $\alpha$ un homomorphisme quelconque ayant cette
propriété. Alors $N' = \Coker(\alpha)$ est un $B$-module de type fini
tel que $N' \otimes \kappa(\mathfrak p) = 0$. Par le lemme de Nakayama
(Algèbre, Lemme \ref{algebra-lemma-NAK}),
on voit que $N'_{\mathfrak p} = 0$. Puisque la fibre $Y_s$ est
géométriquement irréductible de dimension $n$, de point générique $\xi$,
et que l'on vient de voir que $\xi$ n'appartient pas au support de
$\Coker(\alpha)$, la dernière assertion du lemme est vérifiée.
\end{proof}


\section{D\'evissage complet}
\label{flat-section-complete-devissage}

\noindent
Dans cette section, nous expliquons ce qu'est un d\'evissage complet d'un
module et démontrons l'existence de tels dévissages. Le contenu de cette
section consiste surtout à organiser les constructions précédentes.

\begin{definition}
\label{flat-definition-complete-devissage}
Soit $S$ un schéma.
Soit $X$ localement de type fini sur $S$.
Soit $\mathcal{F}$ un $\mathcal{O}_X$-module quasi-cohérent de type fini.
Soit $s \in S$ un point.
Un {\it d\'evissage complet de $\mathcal{F}/X/S$ au-dessus de $s$} est donné par un
diagramme
$$
\xymatrix{
X & Z_1 \ar[l]^{i_1} \ar[d]^{\pi_1} \\
& Y_1 & Z_2 \ar[l]^{i_2} \ar[d]^{\pi_2} \\
& & Y_2 & Z_3 \ar[l] \ar[d] \\
& & & ... & ... \ar[l] \ar[d] \\
& & & & Y_n
}
$$
de schémas sur $S$, des $\mathcal{O}_{Z_k}$-modules quasi-cohérents de type fini
$\mathcal{G}_k$ et des homomorphismes de $\mathcal{O}_{Y_k}$-modules
$$
\alpha_k :
\mathcal{O}_{Y_k}^{\oplus r_k}
\longrightarrow
\pi_{k, *}\mathcal{G}_k,
\quad
k = 1, \ldots, n
$$
vérifiant les propriétés suivantes :
\begin{enumerate}
\item $(Z_1, Y_1, i_1, \pi_1, \mathcal{G}_1)$ est un
d\'evissage en une étape de $\mathcal{F}/X/S$ au-dessus de $s$,
\item l'homomorphisme $\alpha_k$ induit un isomorphisme
$$
\kappa(\xi_k)^{\oplus r_k} \longrightarrow
(\pi_{k, *}\mathcal{G}_k)_{\xi_k}
\otimes_{\mathcal{O}_{Y_k, \xi_k}} \kappa(\xi_k)
$$
où $\xi_k \in (Y_k)_s$ est l'unique point générique,
\item pour $k = 2, \ldots, n$, le système
$(Z_k, Y_k, i_k, \pi_k, \mathcal{G}_k)$
est un d\'evissage en une étape de $\Coker(\alpha_{k - 1})/Y_{k - 1}/S$
au-dessus de $s$,
\item $\Coker(\alpha_n) = 0$.
\end{enumerate}
On dit alors que
$(Z_k, Y_k, i_k, \pi_k, \mathcal{G}_k, \alpha_k)_{k = 1, \ldots, n}$
est un d\'evissage complet de $\mathcal{F}/X/S$ au-dessus de $s$.
\end{definition}

\begin{definition}
\label{flat-definition-complete-devissage-at-x}
Soit $S$ un schéma.
Soit $X$ localement de type fini sur $S$.
Soit $\mathcal{F}$ un $\mathcal{O}_X$-module quasi-cohérent de type fini.
Soit $x \in X$ un point d'image $s \in S$.
Un {\it d\'evissage complet de $\mathcal{F}/X/S$ en $x$} est donné par un
système
$$
(Z_k, Y_k, i_k, \pi_k, \mathcal{G}_k, \alpha_k, z_k, y_k)_{k = 1, \ldots, n}
$$
tel que $(Z_k, Y_k, i_k, \pi_k, \mathcal{G}_k, \alpha_k)$ soit un
d\'evissage complet de $\mathcal{F}/X/S$ au-dessus de $s$, et tel que
\begin{enumerate}
\item $(Z_1, Y_1, i_1, \pi_1, \mathcal{G}_1, z_1, y_1)$ soit un
d\'evissage en une étape de $\mathcal{F}/X/S$ en $x$,
\item pour $k = 2, \ldots, n$, le système
$(Z_k, Y_k, i_k, \pi_k, \mathcal{G}_k, z_k, y_k)$
soit un d\'evissage en une étape de $\Coker(\alpha_{k - 1})/Y_{k - 1}/S$
en $y_{k - 1}$.
\end{enumerate}
\end{definition}

\noindent
Remarquons encore qu'un d\'evissage complet ne peut exister que si $X$ et
$S$ sont affines.

\begin{lemma}
\label{flat-lemma-base-change-complete}
Soient $S$, $X$, $\mathcal{F}$, $s$ comme dans la
Définition \ref{flat-definition-complete-devissage}.
Soit $(S', s') \to (S, s)$ un morphisme quelconque de schémas pointés.
Soit $(Z_k, Y_k, i_k, \pi_k, \mathcal{G}_k, \alpha_k)_{k = 1, \ldots, n}$
un d\'evissage complet de $\mathcal{F}/X/S$ au-dessus de $s$.
Ces données étant fixées, soient $X', Z'_k, Y'_k, i'_k, \pi'_k$ les objets obtenus par changement
de base à partir de $X, Z_k, Y_k, i_k, \pi_k$ suivant $S' \to S$.
Soit $\mathcal{F}'$ l'image réciproque de $\mathcal{F}$ sur $X'$
et soit $\mathcal{G}'_k$ l'image réciproque de $\mathcal{G}_k$ sur $Z'_k$.
Soit $\alpha'_k$ l'image réciproque de $\alpha_k$ sur $Y'_k$.
Si $S'$ est affine, alors
$(Z'_k, Y'_k, i'_k, \pi'_k, \mathcal{G}'_k, \alpha'_k)_{k = 1, \ldots, n}$
est un d\'evissage complet de $\mathcal{F}'/X'/S'$ au-dessus de $s'$.
\end{lemma}

\begin{proof}
D'après le
Lemme \ref{flat-lemma-base-change-one-step},
on sait que le changement de base d'un d\'evissage en une étape est un
d\'evissage en une étape. Il suffit donc de montrer que la formation de
$\Coker(\alpha_k)$ commute au changement de base et que la
condition (2) de la
Définition \ref{flat-definition-complete-devissage}
est préservée par changement de base. La première assertion résulte du fait que
$\pi'_{k, *}\mathcal{G}'_k$ est l'image réciproque de
$\pi_{k, *}\mathcal{G}_k$ (d'après
Cohomologie des schémas, Lemme \ref{coherent-lemma-affine-base-change})
et de l'exactitude à droite de $\otimes$. Pour la seconde,
le même raisonnement donne
$$
(\pi_{k, *}\mathcal{G}_k)_{\xi_k}
\otimes_{\mathcal{O}_{Y_k, \xi_k}} \kappa(\xi_k)
\otimes_{\kappa(\xi_k)} \kappa(\xi'_k)
\cong
(\pi'_{k, *}\mathcal{G}'_k)_{\xi'_k}
\otimes_{\mathcal{O}_{Y'_k, \xi'_k}} \kappa(\xi'_k)
$$
avec des notations évidentes.
\end{proof}

\begin{lemma}
\label{flat-lemma-base-change-complete-at-x}
Soient $S$, $X$, $\mathcal{F}$, $x$, $s$ comme dans la
Définition \ref{flat-definition-complete-devissage-at-x}.
Soit $(S', s') \to (S, s)$ un morphisme de schémas pointés
induisant un isomorphisme $\kappa(s) = \kappa(s')$. Soit
$(Z_k, Y_k, i_k, \pi_k, \mathcal{G}_k, \alpha_k, z_k, y_k)_{k = 1, \ldots, n}$
un d\'evissage complet de $\mathcal{F}/X/S$ en $x$.
Soit
$(Z'_k, Y'_k, i'_k, \pi'_k, \mathcal{G}'_k, \alpha'_k)_{k = 1, \ldots, n}$
le système construit dans le
Lemme \ref{flat-lemma-base-change-complete},
et soit $x' \in X'$ (resp.\ $z'_k \in Z'$, $y'_k \in Y'$) l'unique
point s'envoyant à la fois sur $x \in X$ (resp.\ $z_k \in Z_k$, $y_k \in Y_k$)
et sur $s' \in S'$.
Si $S'$ est affine, alors
$(Z'_k, Y'_k, i'_k, \pi'_k, \mathcal{G}'_k, \alpha'_k,
z'_k, y'_k)_{k = 1, \ldots, n}$
est un d\'evissage complet de $\mathcal{F}'/X'/S'$ en $x'$.
\end{lemma}

\begin{proof}
On combine
le Lemme \ref{flat-lemma-base-change-complete}
et
le Lemme \ref{flat-lemma-base-change-one-step-at-x}.
\end{proof}

\begin{definition}
\label{flat-definition-shrink-complete}
Soient $S$, $X$, $\mathcal{F}$, $x$, $s$ comme dans la
Définition \ref{flat-definition-complete-devissage-at-x}.
Considérons un d\'evissage complet
$(Z_k, Y_k, i_k, \pi_k, \mathcal{G}_k, \alpha_k, z_k, y_k)_{k = 1, \ldots, n}$
de $\mathcal{F}/X/S$ en $x$. Appelons
{\it restriction standard} de cette situation la donnée
d'ouverts standards $S' \subset S$, $X' \subset X$,
$Z'_k \subset Z_k$ et $Y'_k \subset Y_k$ tels que $s_k \in S'$,
$x_k \in X'$, $z_k \in Z'$ et $y_k \in Y'$, et tels que
$$
(Z'_k, Y'_k, i'_k, \pi'_k,
\mathcal{G}'_k, \alpha'_k, z_k, y_k)_{k = 1, \ldots, n}
$$
soit un d\'evissage en une étape de $\mathcal{F}'/X'/S'$ en $x$, où
$\mathcal{G}'_k = \mathcal{G}_k|_{Z'_k}$ et
$\mathcal{F}' = \mathcal{F}|_{X'}$.
\end{definition}

\begin{lemma}
\label{flat-lemma-shrink-complete}
Sous les hypothèses et avec les notations de la
Définition \ref{flat-definition-shrink-complete},
on a :
\begin{enumerate}
\item
\label{flat-item-shrink-base-complete}
Si $S' \subset S$ est un voisinage ouvert standard de $s$, alors,
en posant $X' = X_{S'}$, $Z'_k = Z_{S'}$ et $Y'_k = Y_{S'}$, on obtient une
restriction standard.
\item
\label{flat-item-shrink-on-Y-complete}
Soit $W \subset Y_n$ un voisinage ouvert standard de $y$.
Il existe alors une restriction standard telle que $Y'_n = W \times_S S'$.
\item
\label{flat-item-shrink-on-X-complete}
Soit $U \subset X$ un voisinage ouvert de $x$.
Il existe alors une restriction standard telle que $X' \subset U$.
\end{enumerate}
\end{lemma}

\begin{proof}
L'assertion (1) résulte immédiatement des
Lemmes \ref{flat-lemma-base-change-complete-at-x} et
\ref{flat-lemma-shrink}.

\medskip\noindent
Démonstration de (2). Pour simplifier les notations, posons $X = Y_0$. On applique
le Lemme \ref{flat-lemma-shrink} (\ref{flat-item-shrink-on-Y})
pour trouver une restriction standard
$S', Y'_{n - 1}, Z'_n, Y'_n$
du d\'evissage en une étape de $\Coker(\alpha_{n - 1})/Y_{n - 1}/S$
en $y_{n - 1}$ avec $Y'_n = W \times_S S'$. On peut répéter cette opération
et trouver une restriction standard
$S'', Y''_{n - 2}, Z''_{n - 1}, Y''_{n - 1}$
du d\'evissage en une étape de $\Coker(\alpha_{n - 2})/Y_{n - 2}/S$
en $y_{n - 2}$ avec $Y''_{n - 1} = Y'_{n - 1} \times_S S''$.
On continue ainsi jusqu'à obtenir
$S^{(n)}, Y^{(n)}_0, Z^{(n)}_1, Y^{(n)}_1$.
Il est alors clair que l'on obtient la restriction standard cherchée
en prenant $S^{(n)}$, $X^{(n)}$, $Z_k^{(n - k)} \times_S S^{(n)}$ et
$Y_k^{(n - k)} \times_S S^{(n)}$, avec la propriété voulue.

\medskip\noindent
Démonstration de (3). On raisonne par récurrence sur la longueur du
d\'evissage complet. On applique d'abord
le Lemme \ref{flat-lemma-shrink} (\ref{flat-item-shrink-on-X})
pour trouver une restriction standard
$S', X', Z'_1, Y'_1$
du d\'evissage en une étape de $\mathcal{F}/X/S$ en $x$
avec $X' \subset U$. Si $n = 1$, la démonstration est terminée.
Si $n > 1$, l'hypothèse de récurrence fournit une restriction standard
$S''$, $Y''_1$, $Z''_k$ et $Y''_k$ du d\'evissage complet
$(Z_k, Y_k, i_k, \pi_k, \mathcal{G}_k, \alpha_k, z_k, y_k)_{k = 2, \ldots, n}$
de $\Coker(\alpha_1)/Y_1/S$ en $x$, telle que
$Y''_1 \subset Y'_1$. En utilisant
le Lemme \ref{flat-lemma-shrink} (\ref{flat-item-shrink-on-Y}),
on peut trouver $S''' \subset S'$, $X''' \subset X'$, $Z'''_1$ et
$Y'''_1 = Y''_1 \times_S S'''$ qui constituent une restriction standard.
La solution du problème consiste à prendre
$$
S''', X''', Z'''_1, Y'''_1, Z''_2 \times_S S''',
Y''_2 \times_S S''', \ldots, Z''_n \times_S S''', Y''_n \times_S S'''
$$
Cela achève la démonstration du lemme.
\end{proof}

\begin{proposition}
\label{flat-proposition-existence-complete-at-x}
Soit $S$ un schéma.
Soit $X$ localement de type fini sur $S$.
Soit $x \in X$ un point d'image $s \in S$.
Il existe un diagramme commutatif
$$
\xymatrix{
(X, x) \ar[d] & (X', x') \ar[l]^g \ar[d] \\
(S, s) & (S', s') \ar[l]
}
$$
de schémas pointés dont les flèches
horizontales sont des voisinages \'etales élémentaires
et tel que $g^*\mathcal{F}/X'/S'$ admette un
d\'evissage complet en $x$.
\end{proposition}

\begin{proof}
On démontre cette assertion par récurrence sur l'entier
$d = \dim_x(\text{Supp}(\mathcal{F}_s))$.
D'après le
Lemme \ref{flat-lemma-elementary-devissage-variant},
il existe un diagramme
$$
\xymatrix{
(X, x) \ar[d] & (X', x') \ar[l]^g \ar[d] \\
(S, s) & (S', s') \ar[l]
}
$$
de schémas pointés dont les flèches
horizontales sont des voisinages \'etales élémentaires
et tel que $g^*\mathcal{F}/X'/S'$ admette un d\'evissage en une étape en $x'$.
Le caractère local du problème permet de remplacer
$(X, x) \to (S, s)$ par $(X', x') \to (S', s')$. Après ce remplacement,
on peut donc supposer qu'il existe un d\'evissage en une étape
$(Z_1, Y_1, i_1, \pi_1, \mathcal{G}_1)$ de $\mathcal{F}/X/S$ en $x$.

\medskip\noindent
On applique
le Lemme \ref{flat-lemma-existence-alpha}
pour trouver un homomorphisme
$$
\alpha_1 :
\mathcal{O}_{Y_1}^{\oplus r_1}
\longrightarrow
\pi_{1, *}\mathcal{G}_1
$$
qui induit un isomorphisme d'espaces vectoriels sur $\kappa(\xi_1)$,
où $\xi_1 \in Y_1$ est l'unique point générique de la fibre de
$Y_1$ au-dessus de $s$. De plus,
$\dim_{y_1}(\text{Supp}(\Coker(\alpha_1)_s)) < d$.
Il peut arriver que la fibre du faisceau $\Coker(\alpha_1)_s$
en $y_1$ soit nulle. On peut alors restreindre $Y_1$ en utilisant
le Lemme \ref{flat-lemma-shrink} (\ref{flat-item-shrink-on-Y})
et supposer que $\Coker(\alpha_1) = 0$, obtenant ainsi un
d\'evissage complet de longueur nulle.

\medskip\noindent
Supposons maintenant que la fibre du faisceau $\Coker(\alpha_1)_s$
en $y_1$ ne soit pas nulle. L'hypothèse de récurrence fournit alors un
diagramme commutatif
\begin{equation}
\label{flat-equation-overcome-this}
\vcenter{
\xymatrix{
(Y_1, y_1) \ar[d] & (Y'_1, y'_1) \ar[l]^h \ar[d] \\
(S, s) & (S', s') \ar[l]
}
}
\end{equation}
de schémas pointés dont les flèches
horizontales sont des voisinages \'etales élémentaires
et tel que $h^*\Coker(\alpha_1)/Y'_1/S'$ admette un
d\'evissage complet
$$
(Z_k, Y_k, i_k, \pi_k, \mathcal{G}_k, \alpha_k, z_k, y_k)_{k = 2, \ldots, n}
$$
en $y'_1$. (En particulier, $i_2 : Z_2 \to Y'_1$ est une immersion fermée dans
$Y'_2$.) On applique alors
le Lemme \ref{flat-lemma-elementary-etale-neighbourhood}
à $S, X, \mathcal{F}, x, s$, au système
$(Z_1, Y_1, i_1, \pi_1, \mathcal{G}_1)$ et au
diagramme (\ref{flat-equation-overcome-this}). On obtient un diagramme
$$
\xymatrix{
& & (X'', x'') \ar[lld] \ar[d] & (Z''_1, z''_1) \ar[l] \ar[lld] \ar[d] \\
(X, x) \ar[d] & (Z_1, z_1) \ar[l] \ar[d] &
(S'', s'') \ar[lld] & (Y''_1, y''_1) \ar[lld] \ar[l] \\
(S, s) & (Y_1, y_1) \ar[l]
}
$$
possédant toutes les propriétés énumérées dans le lemme cité.
En particulier, $Y''_1 \subset Y'_1 \times_{S'} S''$. Posons
$X_1 = Y'_1 \times_{S'} S''$ et notons $\mathcal{F}_1$
l'image réciproque de $\Coker(\alpha_1)$. D'après le
Lemme \ref{flat-lemma-base-change-complete-at-x},
le système
\begin{equation}
\label{flat-equation-shrink-this}
(Z_k \times_{S'} S'',
Y_k \times_{S'} S'', i''_k, \pi''_k, \mathcal{G}''_k,
\alpha''_k, z''_k, y''_k)_{k = 2, \ldots, n}
\end{equation}
est un d\'evissage complet de $\mathcal{F}_1$
sur $X_1$. De nouveau, le caractère local du problème
permet de remplacer $(X, x) \to (S, s)$ par $(X'', x'') \to (S'', s'')$.
On voit ainsi que l'on peut supposer :
\begin{enumerate}
\item[(a)] Il existe un d\'evissage en une étape
$(Z_1, Y_1, i_1, \pi_1, \mathcal{G}_1)$ de $\mathcal{F}/X/S$ en $x$,
\item[(b)] il existe un $\alpha_1 : \mathcal{O}_{Y_1}^{\oplus r_1}
\to \pi_{1, *}\mathcal{G}_1$ tel que $\alpha \otimes \kappa(\xi_1)$
soit un isomorphisme,
\item[(c)] $Y_1 \subset X_1$ est ouvert, $y_1 = x_1$, et
$\mathcal{F}_1|_{Y_1} \cong \Coker(\alpha_1)$, et
\item[(d)] il existe un d\'evissage complet
$(Z_k, Y_k, i_k, \pi_k, \mathcal{G}_k, \alpha_k, z_k, y_k)_{k = 2, \ldots, n}$
de $\mathcal{F}_1/X_1/S$ en $x_1$.
\end{enumerate}
Pour achever la démonstration, il suffit de restreindre le d\'evissage en une étape
et le d\'evissage complet de façon à les assembler en un
d\'evissage complet. (Nous suggérons au lecteur de le faire lui-même en utilisant
les Lemmes \ref{flat-lemma-shrink} et
\ref{flat-lemma-shrink-complete},
au lieu de lire la démonstration qui suit.) Puisque $Y_1 \subset X_1$
est un voisinage ouvert de $x_1$, on peut appliquer
le Lemme \ref{flat-lemma-shrink-complete} (\ref{flat-item-shrink-on-X-complete})
pour trouver une restriction standard $S', X'_1, Z'_2, Y'_2, \ldots, Y'_n$
de la donnée (d), telle que $X'_1 \subset Y_1$. Notons que $X'_1$ est aussi
un ouvert standard du schéma affine $Y_1$. On restreint ensuite la donnée
(a) comme suit : on restreint d'abord la base $S$ à $S'$, voir
le Lemme \ref{flat-lemma-shrink} (\ref{flat-item-shrink-base}), puis
on restreint le résultat à $S''$, $X''$, $Z''_1$, $Y''_1$ en utilisant
le Lemme \ref{flat-lemma-shrink} (\ref{flat-item-shrink-on-Y}),
de sorte que finalement $Y''_1 = X'_1 \times_S S''$ et $S'' \subset S'$.
On voit alors que
$$
Z''_1, Y''_1, Z'_2 \times_{S'} S'', Y'_2 \times_{S'} S'', \ldots,
Y'_n \times_{S'} S''
$$
fournit le d\'evissage complet cherché.
\end{proof}

\noindent
En organisant encore les constructions, on obtient la conséquence suivante.

\begin{lemma}
\label{flat-lemma-existence-complete}
Soit $X \to S$ un morphisme de type fini de schémas.
Soit $\mathcal{F}$ un $\mathcal{O}_X$-module quasi-cohérent de type fini.
Soit $s \in S$ un point.
Il existe un voisinage \'etale élémentaire
$(S', s') \to (S, s)$ et des morphismes \'etales
$h_i : Y_i \to X_{S'}$, $i = 1, \ldots, n$, tels que, pour chaque
$i$, il existe un d\'evissage complet de $\mathcal{F}_i/Y_i/S'$ au-dessus de $s'$,
où $\mathcal{F}_i$ est l'image réciproque de $\mathcal{F}$ sur $Y_i$,
et tels que $X_s = (X_{S'})_{s'} \subset \bigcup h_i(Y_i)$.
\end{lemma}

\begin{proof}
Pour tout point $x \in X_s$, on peut trouver un diagramme
$$
\xymatrix{
(X, x) \ar[d] & (X', x') \ar[l]^g \ar[d] \\
(S, s) & (S', s') \ar[l]
}
$$
de schémas pointés dont les flèches horizontales sont des voisinages
\'etales élémentaires et tel que $g^*\mathcal{F}/X'/S'$ admette un
d\'evissage complet en $x'$. Puisque $X \to S$ est de type fini, la
fibre $X_s$ est quasi-compacte et, puisque chaque $g : X' \to X$ comme ci-dessus
est ouvert, on peut recouvrir $X_s$ par une réunion finie de $g(X'_{s'})$.
On peut donc trouver une famille finie de tels diagrammes
$$
\vcenter{
\xymatrix{
(X, x) \ar[d] & (X'_i, x'_i) \ar[l]^{g_i} \ar[d] \\
(S, s) & (S'_i, s'_i) \ar[l]
}
}
\quad i = 1, \ldots, n
$$
telle que $X_s = \bigcup g_i(X'_i)$. Posons
$S' = S'_1 \times_S \ldots \times_S S'_n$
et soit $Y_i = X_i \times_{S'_i} S'$ le changement de base de $X'_i$ à $S'$. D'après
le Lemme \ref{flat-lemma-base-change-complete},
on voit que l'image réciproque de $\mathcal{F}$ sur $Y_i$ admet un d\'evissage complet
au-dessus de $s$, ce qui donne le résultat voulu.
\end{proof}



\section{Traduction en termes d'algèbre}
\label{flat-section-translation}

\noindent
Il peut être utile d'expliciter en termes algébriques ce que signifie l'existence d'un
d\'evissage complet. Nous introduisons la notion suivante (qui n'est pas
très utile, d'où le nom démesurément long que nous lui donnons).

\begin{definition}
\label{flat-definition-elementary-etale-neighbourhood}
Soit $R \to S$ un homomorphisme d'anneaux. Soit $\mathfrak q$ un idéal premier de $S$ au-dessus
de l'idéal premier $\mathfrak p$ de $R$. Une {\it localisation \'etale élémentaire de
l'homomorphisme d'anneaux $R \to S$ en $\mathfrak q$} est donnée par un diagramme commutatif
d'anneaux et par les idéaux premiers correspondants
$$
\xymatrix{
S \ar[r] & S' \\
R \ar[u] \ar[r] & R' \ar[u]
}
\quad\quad
\xymatrix{
\mathfrak q \ar@{-}[r] & \mathfrak q' \\
\mathfrak p \ar@{-}[u] \ar@{-}[r] & \mathfrak p' \ar@{-}[u]
}
$$
tels que $R \to R'$ et $S \to S'$ soient des homomorphismes d'anneaux \'etales et que
$\kappa(\mathfrak p) = \kappa(\mathfrak p')$ et
$\kappa(\mathfrak q) = \kappa(\mathfrak q')$.
\end{definition}

\begin{definition}
\label{flat-definition-complete-devissage-algebra}
Soit $R \to S$ un homomorphisme d'anneaux de type fini.
Soit $\mathfrak r$ un idéal premier de $R$.
Soit $N$ un $S$-module de type fini.
Un {\it d\'evissage complet de $N/S/R$ au-dessus de $\mathfrak r$}
est donné par des homomorphismes de $R$-algèbres
$$
\xymatrix{
& A_1 & & A_2 & & ... & & A_n \\
S \ar[ru] & & B_1 \ar[lu] \ar[ru] & & ... \ar[lu] \ar[ru] & &
... \ar[lu] \ar[ru] & & B_n \ar[lu]
}
$$
des $A_i$-modules de type fini $M_i$ et des homomorphismes de $B_i$-modules
$\alpha_i : B_i^{\oplus r_i} \to M_i$ tels que
\begin{enumerate}
\item $S \to A_1$ soit surjectif et de présentation finie,
\item $B_i \to A_{i + 1}$ soit surjectif et de présentation finie,
\item $B_i \to A_i$ soit fini,
\item $R \to B_i$ soit lisse à fibres géométriquement irréductibles,
\item $N \cong M_1$ comme $S$-modules,
\item $\Coker(\alpha_i) \cong M_{i + 1}$ comme $B_i$-modules,
\item $\alpha_i : \kappa(\mathfrak p_i)^{\oplus r_i}
\to M_i \otimes_{B_i} \kappa(\mathfrak p_i)$ soit un isomorphisme,
où $\mathfrak p_i = \mathfrak rB_i$, et
\item $\Coker(\alpha_n) = 0$.
\end{enumerate}
On dit alors que
$(A_i, B_i, M_i, \alpha_i)_{i = 1, \ldots, n}$
est un d\'evissage complet de $N/S/R$ au-dessus de $\mathfrak r$.
\end{definition}

\begin{remark}
\label{flat-remark-finite-presentation}
Notons que les $R$-algèbres $B_i$, pour tout $i$, et $A_i$, pour $i \geq 2$,
sont de présentation finie sur $R$. Si $S$ est de présentation finie sur
$R$, alors $A_1$ est également de présentation finie sur
$R$. Dans ce cas, tous les homomorphismes d'anneaux du d\'evissage complet sont de
présentation finie. Voir
Algèbre, Lemme \ref{algebra-lemma-compose-finite-type}.
Toujours sous l'hypothèse que $S$ est de présentation finie sur $R$,
les conditions suivantes sont équivalentes :
\begin{enumerate}
\item $M$ est de présentation finie sur $S$,
\item $M_1$ est de présentation finie sur $A_1$,
\item $M_1$ est de présentation finie sur $B_1$,
\item chaque $M_i$ est de présentation finie à la fois comme $A_i$-module
et comme $B_i$-module.
\end{enumerate}
Les équivalences (1) $\Leftrightarrow$ (2) et (2) $\Leftrightarrow$ (3)
résultent de
Algèbre, Lemme \ref{algebra-lemma-finite-finitely-presented-extension}.
Si $M_1$ est de présentation finie, il en est de même de $\Coker(\alpha_1)$ (voir
Algèbre, Lemme \ref{algebra-lemma-extension}),
donc de $M_2$, et ainsi de suite.
\end{remark}

\begin{definition}
\label{flat-definition-complete-devissage-at-x-algebra}
Soit $R \to S$ un homomorphisme d'anneaux de type fini.
Soit $\mathfrak q$ un idéal premier de $S$ au-dessus de l'idéal premier $\mathfrak r$ de $R$.
Soit $N$ un $S$-module de type fini.
Un {\it d\'evissage complet de $N/S/R$ en $\mathfrak q$} est donné par un
d\'evissage complet $(A_i, B_i, M_i, \alpha_i)_{i = 1, \ldots, n}$
de $N/S/R$ au-dessus de $\mathfrak r$ et par des idéaux premiers $\mathfrak q_i \subset B_i$
au-dessus de $\mathfrak r$ tels que
\begin{enumerate}
\item $\kappa(\mathfrak r) \subset \kappa(\mathfrak q_i)$ soit purement
transcendante,
\item il existe un unique idéal premier $\mathfrak q'_i \subset A_i$
au-dessus de $\mathfrak q_i \subset B_i$,
\item $\mathfrak q = \mathfrak q'_1 \cap S$ et
$\mathfrak q_i = \mathfrak q'_{i + 1} \cap A_i$,
\item $R \to B_i$ soit de dimension relative
$\dim_{\mathfrak q_i}(\text{Supp}(M_i \otimes_R \kappa(\mathfrak r)))$.
\end{enumerate}
\end{definition}

\begin{remark}
\label{flat-remark-same-notion}
Soit $A \to B$ un homomorphisme d'anneaux de type fini et soit $N$ un
$B$-module de type fini. Soit $\mathfrak q$ un idéal premier de $B$ au-dessus de l'idéal premier
$\mathfrak r$ de $A$. Posons $X = \Spec(B)$, $S = \Spec(A)$ et
$\mathcal{F} = \widetilde{N}$ sur $X$. Soit $x$ le point correspondant
à $\mathfrak q$ et soit $s \in S$ le point correspondant à
$\mathfrak p$. Alors
\begin{enumerate}
\item s'il existe un d\'evissage complet de $\mathcal{F}/X/S$
au-dessus de $s$, il existe un d\'evissage complet de
$N/B/A$ au-dessus de $\mathfrak p$, et
\item il existe un d\'evissage complet de $\mathcal{F}/X/S$
en $x$ si et seulement s'il existe un d\'evissage complet de
$N/B/A$ en $\mathfrak q$.
\end{enumerate}
La seule nuance est que nous avons omis la condition sur
la dimension relative dans la formulation « un d\'evissage complet de
$N/B/A$ au-dessus de $\mathfrak p$ », ce qui explique que l'implication (1)
ne soit donnée que dans un sens.
La notion de d\'evissage complet en
$\mathfrak q$ inclut bien cette condition. Quoi qu'il en soit, nous utiliserons
seulement le fait que l'existence pour $\mathcal{F}/X/S$
entraîne l'existence pour $N/B/A$.
\end{remark}

\begin{lemma}
\label{flat-lemma-existence-algebra}
Soit $R \to S$ un homomorphisme d'anneaux de type fini.
Soit $M$ un $S$-module de type fini.
Soit $\mathfrak q$ un idéal premier de $S$.
Il existe une localisation \'etale élémentaire
$R' \to S', \mathfrak q', \mathfrak p'$ de
l'homomorphisme d'anneaux $R \to S$ en $\mathfrak q$ telle qu'il
existe un d\'evissage complet de
$(M \otimes_S S')/S'/R'$ en $\mathfrak q'$.
\end{lemma}

\begin{proof}
C'est une reformulation de la
Proposition \ref{flat-proposition-existence-complete-at-x}
au moyen de la
Remarque \ref{flat-remark-same-notion}
\end{proof}




\section{Localisation et homomorphismes universellement injectifs}
\label{flat-section-localize-universally-injective}


\begin{lemma}
\label{flat-lemma-homothety-spectrum}
Soit $R \to S$ un homomorphisme d'anneaux.
Soit $N$ un $S$-module.
Supposons que
\begin{enumerate}
\item $R$ soit un anneau local d'idéal maximal $\mathfrak m$,
\item $\overline{S} = S/\mathfrak m S$ soit noethérien, et
\item $\overline{N} = N/\mathfrak m_R N$ soit un $\overline{S}$-module de type fini.
\end{enumerate}
Soit $\Sigma \subset S$ la partie multiplicative des éléments qui ne sont pas
diviseurs de zéro sur $\overline{N}$. Alors $\Sigma^{-1}S$ est un anneau semi-local
dont le spectre est formé des idéaux premiers $\mathfrak q \subset S$ contenus dans un
élément de $\text{Ass}_S(\overline{N})$. De plus, tout idéal
maximal de $\Sigma^{-1}S$ correspond à un idéal premier associé à
$\overline{N}$ sur $\overline{S}$.
\end{lemma}

\begin{proof}
Notons que
$\text{Ass}_S(\overline{N}) = \text{Ass}_{\overline{S}}(\overline{N})$, voir
Algèbre, Lemme \ref{algebra-lemma-ass-quotient-ring}.
Cet ensemble est fini d'après
Algèbre, Lemme \ref{algebra-lemma-finite-ass}.
Écrivons $\{\mathfrak q_1, \ldots, \mathfrak q_r\} = \text{Ass}_S(\overline{N})$.
On a $\Sigma = S \setminus (\bigcup \mathfrak q_i)$ d'après
Algèbre, Lemme \ref{algebra-lemma-ass-zero-divisors}.
D'après la description de $\Spec(\Sigma^{-1}S)$ donnée dans
Algèbre, Lemme \ref{algebra-lemma-spec-localization},
et d'après
Algèbre, Lemme \ref{algebra-lemma-silly},
on voit que les idéaux premiers de $\Sigma^{-1}S$ correspondent aux idéaux premiers de
$S$ contenus dans l'un des $\mathfrak q_i$.
Les idéaux maximaux de $\Sigma^{-1}S$ correspondent donc bijectivement aux
éléments maximaux (pour\ l'inclusion) de l'ensemble
$\{\mathfrak q_1, \ldots, \mathfrak q_r\}$. Cela démontre le lemme.
\end{proof}

\begin{lemma}
\label{flat-lemma-homothety-universally-injective}
Sous les hypothèses et avec les notations du
Lemme \ref{flat-lemma-homothety-spectrum},
supposons de plus que
\begin{enumerate}
\item $S$ soit local et que $R \to S$ soit un homomorphisme local,
\item $S$ soit essentiellement de présentation finie sur $R$,
\item $N$ soit de présentation finie sur $S$, et
\item $N$ soit plat sur $R$.
\end{enumerate}
Alors tout $s \in \Sigma$ définit un
homomorphisme universellement injectif de $R$-modules $s : N \to N$, et
l'homomorphisme $N \to \Sigma^{-1}N$ est $R$-universellement injectif.
\end{lemma}

\begin{proof}
D'après
Algèbre, Lemme \ref{algebra-lemma-mod-injective-general},
la suite $0 \to N \to N \to N/sN \to 0$ est exacte et
$N/sN$ est plat sur $R$. Il en résulte que $s : N \to N$
est universellement injectif, voir
Algèbre, Lemme \ref{algebra-lemma-flat-tor-zero}.
L'homomorphisme $N \to \Sigma^{-1}N$ est universellement injectif comme limite
inductive filtrante des homomorphismes $s : N \to N$.
\end{proof}

\begin{lemma}
\label{flat-lemma-base-change-universally-flat-local}
Soit $R \to S$ un homomorphisme d'anneaux.
Soit $N$ un $S$-module.
Soit $S \to S'$ un homomorphisme d'anneaux.
Supposons que
\begin{enumerate}
\item $R \to S$ soit un homomorphisme local d'anneaux locaux,
\item $S$ soit essentiellement de présentation finie sur $R$,
\item $N$ soit de présentation finie sur $S$,
\item $N$ soit plat sur $R$,
\item $S \to S'$ soit plat, et
\item l'image de $\Spec(S') \to \Spec(S)$ contienne
tous les idéaux premiers $\mathfrak q$ de $S$ au-dessus de $\mathfrak m_R$
tels que $\mathfrak q$ soit un idéal premier associé à $N/\mathfrak m_R N$.
\end{enumerate}
Alors $N \to N \otimes_S S'$ est $R$-universellement injectif.
\end{lemma}

\begin{proof}
Posons $N' = N \otimes_R S'$. Considérons le diagramme commutatif
$$
\xymatrix{
N \ar[d] \ar[r] & N' \ar[d] \\
\Sigma^{-1}N \ar[r] & \Sigma^{-1}N'
}
$$
où $\Sigma \subset S$ est l'ensemble des éléments qui ne sont pas
diviseurs de zéro sur $N/\mathfrak m_R N$. S'il est établi que l'homomorphisme
$N \to \Sigma^{-1}N'$ est universellement injectif, alors $N \to N'$
l'est aussi (voir
Algèbre, Lemme \ref{algebra-lemma-universally-injective-permanence}).

\medskip\noindent
D'après le
Lemme \ref{flat-lemma-homothety-spectrum},
l'anneau $\Sigma^{-1}S$ est un anneau semi-local dont les idéaux maximaux
correspondent aux idéaux premiers associés à $N/\mathfrak m_R N$.
L'image de
$\Spec(\Sigma^{-1}S') \to \Spec(\Sigma^{-1}S)$
contient donc tous ces idéaux maximaux par hypothèse. D'après
Algèbre, Lemme \ref{algebra-lemma-ff-rings},
l'homomorphisme d'anneaux $\Sigma^{-1}S \to \Sigma^{-1}S'$ est fidèlement plat.
Ainsi $\Sigma^{-1}N \to \Sigma^{-1}N'$, qui est l'homomorphisme
$$
N \otimes_S \Sigma^{-1}S \longrightarrow N \otimes_S \Sigma^{-1}S'
$$
est universellement injectif, voir
Algèbre, Lemmes \ref{algebra-lemma-faithfully-flat-universally-injective} et
\ref{algebra-lemma-universally-injective-tensor}.
Enfin, on applique
le Lemme \ref{flat-lemma-homothety-universally-injective}
pour voir que $N \to \Sigma^{-1}N$ est universellement injectif.
Puisque le composé d'homomorphismes de modules universellement injectifs est universellement
injectif (voir
Algèbre, Lemme \ref{algebra-lemma-composition-universally-injective}),
on en conclut que $N \to \Sigma^{-1}N'$ est universellement injectif, ce qui donne le résultat voulu.
\end{proof}

\begin{lemma}
\label{flat-lemma-base-change-universally-flat}
Soit $R \to S$ un homomorphisme d'anneaux.
Soit $N$ un $S$-module.
Soit $S \to S'$ un homomorphisme d'anneaux.
Supposons que
\begin{enumerate}
\item $R \to S$ soit de présentation finie et que $N$ soit de présentation finie
sur $S$,
\item $N$ soit plat sur $R$,
\item $S \to S'$ soit plat, et
\item l'image de $\Spec(S') \to \Spec(S)$ contienne
tous les idéaux premiers $\mathfrak q$ tels que $\mathfrak q$ soit un idéal premier associé
à $N \otimes_R \kappa(\mathfrak p)$, où $\mathfrak p$ est l'image
réciproque de $\mathfrak q$ dans $R$.
\end{enumerate}
Alors $N \to N \otimes_S S'$ est $R$-universellement injectif.
\end{lemma}

\begin{proof}
D'après
Algèbre, Lemme \ref{algebra-lemma-universally-injective-check-stalks},
il suffit de montrer que $N_{\mathfrak q} \to (N \otimes_R S')_{\mathfrak q}$
est $R_{\mathfrak p}$-universellement injectif pour tout idéal premier $\mathfrak q$
de $S$ au-dessus de $\mathfrak p$ dans $R$. On peut donc appliquer
le Lemme \ref{flat-lemma-base-change-universally-flat-local}
aux homomorphismes d'anneaux
$R_{\mathfrak p} \to S_{\mathfrak q} \to S'_{\mathfrak q}$
et au module $N_{\mathfrak q}$.
\end{proof}

\noindent
Le lecteur pourra comparer le lemme suivant à
Algèbre, Lemmes \ref{algebra-lemma-mod-injective} et
\ref{algebra-lemma-mod-injective-general}, ainsi qu'aux résultats de la
section \ref{flat-section-variants-mod-injective}.
Dans chaque cas, la conclusion est que l'homomorphisme $u : M \to N$ est
universellement injectif, de conoyau plat.

\begin{lemma}
\label{flat-lemma-universally-injective-local}
Soit $(R, \mathfrak m)$ un anneau local. Soit $u : M \to N$ un homomorphisme de $R$-modules.
Si $M$ est un $R$-module projectif, si $N$ est un $R$-module plat et si
$\overline{u} : M/\mathfrak mM \to N/\mathfrak mN$ est injectif,
alors $u$ est universellement injectif.
\end{lemma}

\begin{proof}
D'après
Algèbre, Théorème \ref{algebra-theorem-projective-free-over-local-ring},
le module $M$ est libre. Si l'on établit le résultat pour tout facteur direct de type
fini de $M$, le lemme en résultera. On peut donc
supposer que $M$ est libre de type fini. Écrivons $N = \colim_i N_i$ comme
limite inductive filtrante de modules libres de type fini, voir
Algèbre, Théorème \ref{algebra-theorem-lazard}.
Notons que $u : M \to N$ se factorise par $N_i$ pour un certain $i$ (puisque $M$ est libre
de type fini). Notons $u_i : M \to N_i$ l'homomorphisme de $R$-modules correspondant.
Puisque $\overline{u}$ est injectif, on voit que
$\overline{u_i} : M/\mathfrak mM \to N_i/\mathfrak mN_i$ est
injectif et le reste après composition avec les homomorphismes
$N_i/\mathfrak mN_i \to N_{i'}/\mathfrak mN_{i'}$ pour tout $i' \geq i$.
Puisque $M$ et $N_{i'}$ sont libres de type fini sur l'anneau local $R$, cela entraîne
que $M \to N_{i'}$ est une injection scindée pour tout $i' \geq i$. Par suite,
pour tout $R$-module $Q$, on voit que $M \otimes_R Q \to N_{i'} \otimes_R Q$
est injectif pour tout $i' \geq i$. Puisque $- \otimes_R Q$ commute aux
limites inductives, on en conclut que $M \otimes_R Q \to N_{i'} \otimes_R Q$
est injectif, comme souhaité.
\end{proof}

\begin{lemma}
\label{flat-lemma-invert-universally-injective}
Sous les hypothèses et avec les notations du
Lemme \ref{flat-lemma-homothety-spectrum},
supposons de plus que $N$ soit projectif comme $R$-module.
Alors tout $s \in \Sigma$ définit un
homomorphisme universellement injectif de $R$-modules $s : N \to N$, et
l'homomorphisme $N \to \Sigma^{-1}N$ est $R$-universellement injectif.
\end{lemma}

\begin{proof}
Choisissons $s \in \Sigma$. D'après le
Lemme \ref{flat-lemma-universally-injective-local},
l'homomorphisme $s : N \to N$ est universellement injectif.
L'homomorphisme $N \to \Sigma^{-1}N$ est universellement injectif comme limite
inductive filtrante des homomorphismes $s : N \to N$.
\end{proof}






\section{Complétion et modules de Mittag-Leffler}
\label{flat-section-completion-ML}


\begin{lemma}
\label{flat-lemma-completed-direct-sum-ML}
Soit $R$ un anneau. Soit $I \subset R$ un idéal. Soit $A$ un ensemble.
Supposons que $R$ soit noethérien et complet pour la topologie définie par $I$. Le complété
$(\bigoplus\nolimits_{\alpha \in A} R)^\wedge$
est plat et de Mittag-Leffler.
\end{lemma}

\begin{proof}
D'après
Compléments d'algèbre, Lemme
\ref{more-algebra-lemma-ui-completion-direct-sum-into-product},
l'homomorphisme $(\bigoplus\nolimits_{\alpha \in A} R)^\wedge
\to \prod_{\alpha \in A} R$ est universellement injectif.
Ainsi, d'après
Algèbre, Lemmes \ref{algebra-lemma-ui-flat-domain} et
\ref{algebra-lemma-pure-submodule-ML},
il suffit de montrer que $\prod_{\alpha \in A} R$ est plat et de Mittag-Leffler.
D'après
Algèbre, Proposition \ref{algebra-proposition-characterize-coherent}
(et
Algèbre, Lemme \ref{algebra-lemma-Noetherian-coherent}),
on voit que $\prod_{\alpha \in A} R$ est plat.
On conclut donc, car un produit de copies de $R$ est de Mittag-Leffler, voir
Algèbre, Lemme \ref{algebra-lemma-product-over-Noetherian-ring}.
\end{proof}

\begin{lemma}
\label{flat-lemma-lift-ML}
Soit $R$ un anneau. Soit $I \subset R$ un idéal.
Soit $M$ un $R$-module.
Supposons que
\begin{enumerate}
\item $R$ soit noethérien et complet pour la topologie $I$-adique,
\item $M$ soit plat sur $R$, et
\item $M/IM$ soit un $R/I$-module projectif.
\end{enumerate}
Alors le complété $I$-adique $M^\wedge$ est un
$R$-module plat de Mittag-Leffler.
\end{lemma}

\begin{proof}
Choisissons une surjection $F \to M$, où $F$ est un $R$-module libre. D'après
Algèbre, Lemme \ref{algebra-lemma-split-completed-sequence},
le module $M^\wedge$ est un facteur direct du module $F^\wedge$.
Il suffit donc de démontrer le lemme pour $F$.
Dans ce cas, le lemme résulte du
Lemme \ref{flat-lemma-completed-direct-sum-ML}.
\end{proof}

\noindent
Dans les
Lemmes \ref{flat-lemma-universally-injective-to-completion} et
\ref{flat-lemma-universally-injective-to-completion-flat},
l'hypothèse que $S$ est noethérien est satisfaite si $R \to S$ est de type fini, voir
Algèbre, Lemme \ref{algebra-lemma-Noetherian-permanence}.

\begin{lemma}
\label{flat-lemma-universally-injective-to-completion}
Soit $R$ un anneau.
Soit $I \subset R$ un idéal.
Soient $R \to S$ un homomorphisme d'anneaux et $N$ un $S$-module.
Supposons que
\begin{enumerate}
\item $R$ soit un anneau noethérien,
\item $S$ soit un anneau noethérien,
\item $N$ soit un $S$-module de type fini, et
\item pour tout $R$-module de type fini $Q$, tout
$\mathfrak q \in \text{Ass}_S(Q \otimes_R N)$
vérifie $IS + \mathfrak q \not = S$.
\end{enumerate}
Alors l'homomorphisme $N \to N^\wedge$ de $N$ dans le complété $I$-adique de $N$
est universellement injectif comme homomorphisme de $R$-modules.
\end{lemma}

\begin{proof}
Il faut montrer que, pour tout $R$-module de type fini $Q$, l'homomorphisme
$Q \otimes_R N \to Q \otimes_R N^\wedge$ est injectif, voir
Algèbre, Théorème \ref{algebra-theorem-universally-exact-criteria}.
Puisqu'il existe un homomorphisme canonique $Q \otimes_R N^\wedge \to (Q \otimes_R N)^\wedge$,
il suffit de démontrer que l'homomorphisme canonique
$Q \otimes_R N \to (Q \otimes_R N)^\wedge$ est injectif.
On peut donc remplacer $N$ par $Q \otimes_R N$, et il suffit de démontrer
l'injectivité de l'homomorphisme $N \to N^\wedge$.

\medskip\noindent
Soit $K = \Ker(N \to N^\wedge)$. Il suffit de montrer que
$K_{\mathfrak q} = 0$ pour $\mathfrak q \in \text{Ass}(N)$, puisque $N$ est un
sous-module de $\prod_{\mathfrak q \in \text{Ass}(N)} N_{\mathfrak q}$, voir
Algèbre, Lemme \ref{algebra-lemma-zero-at-ass-zero}.
Choisissons $\mathfrak q \in \text{Ass}(N)$. La dernière hypothèse montre
qu'il existe un idéal premier $\mathfrak q' \supset IS + \mathfrak q$.
Puisque $K_{\mathfrak q}$ est un localisé de $K_{\mathfrak q'}$,
il suffit de démontrer la nullité de $K_{\mathfrak q'}$.
Notons que $K = \bigcap I^nN$ ; par suite,
$K_{\mathfrak q'} \subset \bigcap I^nN_{\mathfrak q'}$.
Ainsi $K_{\mathfrak q'} = 0$ d'après
Algèbre, Lemme \ref{algebra-lemma-intersect-powers-ideal-module-zero}.
\end{proof}

\begin{lemma}
\label{flat-lemma-universally-injective-to-completion-flat}
Soit $R$ un anneau.
Soit $I \subset R$ un idéal.
Soient $R \to S$ un homomorphisme d'anneaux et $N$ un $S$-module.
Supposons que
\begin{enumerate}
\item $R$ soit un anneau noethérien,
\item $S$ soit un anneau noethérien,
\item $N$ soit un $S$-module de type fini,
\item $N$ soit plat sur $R$, et
\item pour tout idéal premier $\mathfrak q \subset S$ associé à
$N \otimes_R \kappa(\mathfrak p)$, où $\mathfrak p = R \cap \mathfrak q$,
on ait $IS + \mathfrak q \not = S$.
\end{enumerate}
Alors l'homomorphisme $N \to N^\wedge$ de $N$ dans le complété $I$-adique de $N$
est universellement injectif comme homomorphisme de $R$-modules.
\end{lemma}

\begin{proof}
Cela résulte du
Lemme \ref{flat-lemma-universally-injective-to-completion},
car
Algèbre, Lemme
\ref{algebra-lemma-bourbaki-fibres} et
Remarque \ref{algebra-remark-bourbaki}
assurent que les ensembles d'idéaux premiers associés aux produits tensoriels
$N \otimes_R Q$ sont contenus dans l'ensemble des idéaux premiers associés aux
modules $N \otimes_R \kappa(\mathfrak p)$.
\end{proof}




\section{Modules projectifs}
\label{flat-section-projective}

\noindent
Le lemme suivant permet de démontrer la projectivité par
récurrence noethérienne sur la base, voir
le Lemme \ref{flat-lemma-fibres-irreducible-flat-projective}.

\begin{lemma}
\label{flat-lemma-flat-pure-over-complete-projective}
Soit $R$ un anneau.
Soit $I \subset R$ un idéal.
Soient $R \to S$ un homomorphisme d'anneaux et $N$ un $S$-module.
Supposons que
\begin{enumerate}
\item $R$ soit noethérien et complet pour la topologie $I$-adique,
\item $R \to S$ soit de type fini,
\item $N$ soit un $S$-module de type fini,
\item $N$ soit plat sur $R$,
\item $N/IN$ soit projectif comme $R/I$-module, et
\item pour tout idéal premier $\mathfrak q \subset S$ associé à
$N \otimes_R \kappa(\mathfrak p)$, où $\mathfrak p = R \cap \mathfrak q$,
on ait $IS + \mathfrak q \not = S$.
\end{enumerate}
Alors $N$ est projectif comme $R$-module.
\end{lemma}

\begin{proof}
D'après le
Lemme \ref{flat-lemma-universally-injective-to-completion-flat},
l'homomorphisme $N \to N^\wedge$ est universellement injectif.
D'après le
Lemme \ref{flat-lemma-lift-ML},
le module $N^\wedge$ est de Mittag-Leffler.
D'après
Algèbre, Lemme \ref{algebra-lemma-pure-submodule-ML},
on en conclut que $N$ est de Mittag-Leffler.
Ainsi $N$ est engendré par une famille dénombrable, plat et de Mittag-Leffler comme $R$-module,
donc projectif d'après
Algèbre, Lemme \ref{algebra-lemma-countgen-projective}.
\end{proof}

\begin{lemma}
\label{flat-lemma-fibres-irreducible-flat-projective}
Soit $R$ un anneau.
Soit $R \to S$ un homomorphisme d'anneaux.
Supposons que
\begin{enumerate}
\item $R$ soit noethérien,
\item $R \to S$ soit de type fini et plat, et
\item tout anneau de fibre $S \otimes_R \kappa(\mathfrak p)$ soit
géométriquement intègre sur $\kappa(\mathfrak p)$.
\end{enumerate}
Alors $S$ est projectif comme $R$-module.
\end{lemma}

\begin{proof}
Considérons l'ensemble
$$
\{I \subset R \mid S/IS\text{ non projectif comme }R/I\text{-module}\}
$$
Il faut montrer que cet ensemble est vide. Raisonnons par l'absurde en le supposant
non vide. Il contient alors un élément maximal $I$.
Soit $J = \sqrt{I}$ son radical. Si $I \not = J$, alors
$S/JS$ est projectif comme $R/J$-module, $S/IS$ est plat sur $R/I$
et $J/I$ est un idéal nilpotent de $R/I$. En appliquant
Algèbre, Lemme \ref{algebra-lemma-lift-projective},
on voit que $S/IS$ est un $R/I$-module projectif, ce qui est contradictoire.
On peut donc supposer que $I$ est un idéal radical. Autrement dit, on
est ramené à démontrer le lemme lorsque $R$ est un anneau réduit et que
$S/IS$ est un $R/I$-module projectif pour tout idéal non nul $I$
de $R$.

\medskip\noindent
Supposons que $R$ soit un anneau réduit et que $S/IS$ soit un $R/I$-module projectif
pour tout idéal non nul $I$ de $R$. Par platitude générique,
Algèbre, Lemme \ref{algebra-lemma-generic-flatness-Noetherian}
(appliqué à un localisé $R_g$ qui est intègre), ou par le résultat plus général
Algèbre, Lemme \ref{algebra-lemma-generic-flatness-reduced},
il existe un élément non nul $f \in R$ tel que $S_f$ soit libre comme
$R_f$-module. Notons $R^\wedge = \lim R/(f^n)$ le complété $(f)$-adique
de $R$. Notons que l'homomorphisme d'anneaux
$$
R \longrightarrow R_f \times R^\wedge
$$
est fidèlement plat, voir
Algèbre, Lemme \ref{algebra-lemma-completion-flat}.
Par descente fidèlement plate de la projectivité, voir
Algèbre, Théorème \ref{algebra-theorem-ffdescent-projectivity},
il suffit donc de démontrer que $S \otimes_R R^\wedge$ est un
$R^\wedge$-module projectif. Pour cela, on utilisera le critère du
Lemme \ref{flat-lemma-flat-pure-over-complete-projective}.
Tout d'abord, notons que $S/fS = (S \otimes_R R^\wedge)/f(S \otimes_R R^\wedge)$
est un $R/(f)$-module projectif et que $S \otimes_R R^\wedge$ est plat
et de type fini sur $R^\wedge$, par changement de base.
Supposons ensuite que $\mathfrak p^\wedge$ soit un idéal premier
de $R^\wedge$. Soit $\mathfrak p \subset R$ l'idéal premier correspondant
de $R$. Puisque $R \to S$ est à anneaux de fibres géométriquement intègres,
il en est de même après tout changement de base. Par suite,
$\mathfrak q^\wedge = \mathfrak p^\wedge(S \otimes_R R^\wedge)$,
est un idéal premier au-dessus de $\mathfrak p^\wedge$
et c'est l'unique idéal premier associé à
$S \otimes_R \kappa(\mathfrak p^\wedge)$. Il suffit donc de montrer que
$f(S \otimes_R R^\wedge) + \mathfrak q^\wedge \not = S \otimes_R R^\wedge$.
Cela résulte de $\mathfrak p^\wedge + fR^\wedge \not = R^\wedge$, puisque
$f$ appartient au radical de Jacobson de l'anneau $f$-adiquement complet $R^\wedge$,
et du fait que $R^\wedge \to S \otimes_R R^\wedge$ est surjectif sur les spectres,
ses fibres étant non vides (les espaces irréductibles sont non vides).
\end{proof}

\begin{lemma}
\label{flat-lemma-fibres-irreducible-flat-projective-nonnoetherian}
Soit $R$ un anneau. Soit $R \to S$ un homomorphisme d'anneaux.
Supposons que
\begin{enumerate}
\item $R \to S$ soit de présentation finie et plat, et
\item tout anneau de fibre $S \otimes_R \kappa(\mathfrak p)$ soit
géométriquement intègre sur $\kappa(\mathfrak p)$.
\end{enumerate}
Alors $S$ est projectif comme $R$-module.
\end{lemma}

\begin{proof}
On peut trouver un diagramme cocartésien d'anneaux
$$
\xymatrix{
S_0 \ar[r] & S \\
R_0 \ar[u] \ar[r] & R \ar[u]
}
$$
tel que $R_0$ soit de type fini sur $\mathbf{Z}$ et que l'homomorphisme
$R_0 \to S_0$ soit de type fini et plat, à fibres géométriquement
intègres, voir
Compléments sur les morphismes,
Lemmes \ref{more-morphisms-lemma-Noetherian-approximation-flat},
\ref{more-morphisms-lemma-Noetherian-approximation-geometrically-reduced},
\ref{more-morphisms-lemma-Noetherian-approximation-geometrically-irreducible},
et \ref{more-morphisms-lemma-Noetherian-approximation-combine}.
D'après le
Lemme \ref{flat-lemma-fibres-irreducible-flat-projective},
on voit que $S_0$ est un $R_0$-module projectif. Ainsi $S = S_0 \otimes_{R_0} R$
est un $R$-module projectif, voir
Algèbre, Lemme \ref{algebra-lemma-ascend-properties-modules}.
\end{proof}

\begin{remark}
\label{flat-remark-how-in-RG}
Le Lemme \ref{flat-lemma-fibres-irreducible-flat-projective-nonnoetherian}
est une étape essentielle dans l'établissement des résultats de ce chapitre. Son analogue
dans \cite{GruRay} est \cite[I Proposition 3.3.1]{GruRay} :
si $R \to S$ est lisse à fibres géométriquement intègres, alors $S$
est projectif comme $R$-module. C'est un cas particulier du
Lemme \ref{flat-lemma-fibres-irreducible-flat-projective-nonnoetherian},
mais, puisque nous améliorerons encore ce lemme dans la suite, il est peu utile
de disposer dès à présent d'un résultat plus fort.
Esquissons brièvement la démonstration donnée dans \cite{GruRay}.
\begin{enumerate}
\item On se ramène d'abord au cas où $R$ est noethérien, comme ci-dessus.
\item Puisque la projectivité descend par les homomorphismes d'anneaux fidèlement plats, voir
Algèbre, Théorème \ref{algebra-theorem-ffdescent-projectivity},
on peut travailler localement pour la topologie fppf sur $R$, donc supposer
que $R \to S$ admet une section $\sigma : S \to R$. (Il suffit, comme d'habitude,
d'effectuer le changement de base à $S$.) Posons $I = \Ker(S \to R)$.
\item En localisant encore sur $R$, on peut supposer que $I/I^2$ est un
$R$-module libre et que le complété $S^\wedge$ de $S$ pour la topologie définie par $I$
est isomorphe à $R[[t_1, \ldots, t_n]]$, voir
Morphismes, Lemme \ref{morphisms-lemma-section-smooth-morphism}.
On utilise ici le fait que $R \to S$ est lisse.
\item Pour démontrer que $S$ est projectif comme $R$-module, il suffit de
montrer que $S$ est plat, engendré par une famille dénombrable et de Mittag-Leffler comme
$R$-module, voir
Algèbre, Lemme \ref{algebra-lemma-countgen-projective}.
Les deux premières propriétés sont évidentes. Il suffit donc de montrer que $S$
est de Mittag-Leffler comme $R$-module. D'après
Algèbre, Lemme \ref{algebra-lemma-power-series-ML},
le module $R[[t_1, \ldots, t_n]]$ est de Mittag-Leffler sur $R$. Ainsi
Algèbre, Lemme \ref{algebra-lemma-pure-submodule-ML},
montre qu'il suffit d'établir que
$S \to S^\wedge$ est universellement injectif comme homomorphisme de $R$-modules.
\item On applique
le Lemme \ref{flat-lemma-base-change-universally-flat}
pour voir que $S \to S^\wedge$ est $R$-universellement injectif.
En effet, puisque $R \to S$ est à fibres géométriquement intègres, tout point associé
de tout anneau de fibre est précisément le point générique de cet anneau de fibre,
qui appartient à l'image de $\Spec(S^\wedge) \to \Spec(S)$.
\end{enumerate}
Il existe une analogie entre la démonstration ainsi esquissée et
les arguments qui conduisent à la démonstration du
Lemme \ref{flat-lemma-fibres-irreducible-flat-projective-nonnoetherian}.
Dans les deux cas, une complétion joue un rôle essentiel, et l'hypothèse
que les fibres sont géométriquement intègres assure chaque fois que
l'homomorphisme de $S$ dans le complété de $S$ est $R$-universellement injectif.
\end{remark}













\section{Modules plats de type fini, I}
\label{flat-section-finite-type-flat-I}

\noindent
Dans certains cas, étant donnés un homomorphisme d'anneaux $R \to S$ de présentation finie et
un $S$-module de type fini $N$, la platitude de $N$ sur $R$ entraîne que $N$
est de présentation finie. Nous démontrons dans cette section que cette assertion est vraie
« ponctuellement ». Remarquons que la première démonstration de la
Proposition \ref{flat-proposition-finite-type-flat-at-point}
utilise les résultats géométriques de la
section \ref{flat-section-local-structure-module},
mais non l'existence d'un d\'evissage complet.

\begin{lemma}
\label{flat-lemma-induction-step}
Soit $(R, \mathfrak m)$ un anneau local. Soit $R \to S$ un homomorphisme d'anneaux de présentation finie,
plat et à fibres géométriquement intègres. Posons
$\mathfrak p = \mathfrak mS$. Soit $\mathfrak q \subset S$ un idéal premier
au-dessus de $\mathfrak m$. Soit $N$ un $S$-module de type fini.
Il existe $r \geq 0$ et un homomorphisme de $S$-modules
$$
\alpha : S^{\oplus r} \longrightarrow N
$$
tels que
$\alpha : \kappa(\mathfrak p)^{\oplus r} \to N \otimes_S \kappa(\mathfrak p)$
soit un isomorphisme. Pour tout tel $\alpha$, les conditions suivantes sont équivalentes :
\begin{enumerate}
\item $N_{\mathfrak q}$ est $R$-plat,
\item $\alpha$ est $R$-universellement injectif et
$\Coker(\alpha)_{\mathfrak q}$ est $R$-plat,
\item $\alpha$ est injectif et
$\Coker(\alpha)_{\mathfrak q}$ est $R$-plat,
\item $\alpha_{\mathfrak p}$ est un isomorphisme et
$\Coker(\alpha)_{\mathfrak q}$ est $R$-plat, et
\item $\alpha_{\mathfrak q}$ est injectif et
$\Coker(\alpha)_{\mathfrak q}$ est $R$-plat.
\end{enumerate}
\end{lemma}

\begin{proof}
Pour obtenir $\alpha$, posons
$r = \dim_{\kappa(\mathfrak p)} N \otimes_S \kappa(\mathfrak p)$ et choisissons
$x_1, \ldots, x_r \in N$ dont les images forment une base de
$N \otimes_S \kappa(\mathfrak p)$. Définissons
$\alpha(s_1, \ldots, s_r) = \sum s_i x_i$. Cela démontre l'existence.

\medskip\noindent
Fixons un tel $\alpha$. L'implication la plus intéressante est
(1) $\Rightarrow$ (2), que nous démontrons d'abord. Supposons (1).
Puisque $S/\mathfrak mS$ est un anneau intègre dont le corps des fractions est $\kappa(\mathfrak p)$,
on voit que
$(S/\mathfrak mS)^{\oplus r} \to
N_{\mathfrak p}/\mathfrak mN_{\mathfrak p} = N \otimes_S \kappa(\mathfrak p)$
est injectif. Par suite, d'après les
Lemmes \ref{flat-lemma-universally-injective-local} et
\ref{flat-lemma-fibres-irreducible-flat-projective-nonnoetherian},
l'homomorphisme $S^{\oplus r} \to N_{\mathfrak p}$ est $R$-universellement injectif.
Il en résulte que $S^{\oplus r} \to N$ est $R$-universellement injectif, voir
Algèbre, Lemme \ref{algebra-lemma-universally-injective-permanence}.
Le localisé $\alpha_{\mathfrak q}$ est alors lui aussi $R$-universellement
injectif, voir
Algèbre, Lemme \ref{algebra-lemma-universally-injective-localize}.
On en conclut que $\Coker(\alpha)_{\mathfrak q}$ est $R$-plat d'après
Algèbre, Lemme \ref{algebra-lemma-ui-flat-domain}.

\medskip\noindent
L'implication (2) $\Rightarrow$ (3) est immédiate. Si (3) est vérifiée, alors
$\alpha_{\mathfrak p}$ est injectif, comme localisé d'un homomorphisme injectif
de modules. Par le lemme de Nakayama
(Algèbre, Lemme \ref{algebra-lemma-NAK}),
$\alpha_{\mathfrak p}$ est aussi surjectif. Ainsi (3) $\Rightarrow$ (4).
Si (4) est vérifiée, alors $\alpha_{\mathfrak p}$ est un isomorphisme ; par suite,
$\alpha$ est injectif, puisque $S_{\mathfrak q} \to S_{\mathfrak p}$ est injectif.
En effet, les éléments de $S \setminus \mathfrak p$ ne sont pas diviseurs de zéro sur $S$,
comme on le voit en combinant les
Lemmes \ref{flat-lemma-invert-universally-injective} et
\ref{flat-lemma-fibres-irreducible-flat-projective-nonnoetherian}.
Ainsi (4) $\Rightarrow$ (5). Enfin, si (5) est vérifiée, alors
$N_{\mathfrak q}$ est $R$-plat, comme extension de modules plats, voir
Algèbre, Lemme \ref{algebra-lemma-flat-ses}.
Ainsi (5) $\Rightarrow$ (1), ce qui achève la démonstration.
\end{proof}

\begin{lemma}
\label{flat-lemma-complete-devissage-flat-finite-type-module}
Soit $(R, \mathfrak m)$ un anneau local.
Soit $R \to S$ un homomorphisme d'anneaux de présentation finie.
Soit $N$ un $S$-module de type fini.
Soit $\mathfrak q$ un idéal premier de $S$ au-dessus de $\mathfrak m$.
Supposons que $N_{\mathfrak q}$ soit plat sur $R$ et
qu'il existe un d\'evissage complet de $N/S/R$ en $\mathfrak q$.
Alors $N$ est un $S$-module de présentation finie, libre comme $R$-module,
et il existe un isomorphisme
$$
N \cong B_1^{\oplus r_1} \oplus \ldots \oplus B_n^{\oplus r_n}
$$
de $R$-modules, où chaque $B_i$ est une $R$-algèbre lisse à fibres géométriquement
irréductibles.
\end{lemma}

\begin{proof}
Soit $(A_i, B_i, M_i, \alpha_i, \mathfrak q_i)_{i = 1, \ldots, n}$
le d\'evissage complet donné. On démontre le lemme par récurrence sur $n$.
Notons que $N$ est de présentation finie comme $S$-module si et seulement si
$M_1$ est de présentation finie comme $B_1$-module, voir
la Remarque \ref{flat-remark-finite-presentation}.
Notons que $N_{\mathfrak q} \cong (M_1)_{\mathfrak q_1}$ comme $R$-modules,
car (a) $N_{\mathfrak q} \cong (M_1)_{\mathfrak q'_1}$, où
$\mathfrak q'_1$ est l'unique idéal premier de $A_1$ au-dessus de $\mathfrak q_1$,
et (b) $(A_1)_{\mathfrak q'_1} = (A_1)_{\mathfrak q_1}$ d'après
Algèbre, Lemme \ref{algebra-lemma-unique-prime-over-localize-below},
d'où (c) $(M_1)_{\mathfrak q'_1} \cong (M_1)_{\mathfrak q_1}$.
Ainsi $(M_1)_{\mathfrak q_1}$ est un $R$-module plat. On peut donc remplacer
$(S, N)$ par $(B_1, M_1)$ pour démontrer le lemme. D'après le
Lemme \ref{flat-lemma-induction-step},
l'homomorphisme $\alpha_1 : B_1^{\oplus r_1} \to M_1$ est $R$-universellement injectif
et $\Coker(\alpha_1)_{\mathfrak q}$ est $R$-plat.
Notons que $(A_i, B_i, M_i, \alpha_i, \mathfrak q_i)_{i = 2, \ldots, n}$
est un d\'evissage complet de $\Coker(\alpha_1)/B_1/R$ en
$\mathfrak q_1$. L'hypothèse de récurrence
entraîne donc que $\Coker(\alpha_1)$ est de présentation finie comme
$B_1$-module, libre comme $R$-module et possède une décomposition comme dans le lemme.
Il en résulte que $M_1$ est de présentation finie comme $B_1$-module, voir
Algèbre, Lemme \ref{algebra-lemma-extension}.
Il en résulte aussi que
$M_1 \cong B_1^{\oplus r_1} \oplus \Coker(\alpha_1)$
comme $R$-modules, d'où une décomposition comme dans le lemme.
Enfin, $B_1$ est projectif comme $R$-module d'après le
Lemme \ref{flat-lemma-fibres-irreducible-flat-projective-nonnoetherian},
donc libre comme $R$-module d'après
Algèbre, Théorème \ref{algebra-theorem-projective-free-over-local-ring}.
Cela achève la démonstration.
\end{proof}

\begin{proposition}
\label{flat-proposition-finite-type-flat-at-point}
Soit $f : X \to S$ un morphisme de schémas.
Soit $\mathcal{F}$ un faisceau quasi-cohérent sur $X$.
Soit $x \in X$ d'image $s \in S$.
Supposons que
\begin{enumerate}
\item $f$ soit localement de présentation finie,
\item $\mathcal{F}$ soit de type fini, et
\item $\mathcal{F}$ soit plat en $x$ sur $S$.
\end{enumerate}
Il existe alors un voisinage \'etale élémentaire $(S', s') \to (S, s)$
et un sous-schéma ouvert
$$
V \subset X \times_S \Spec(\mathcal{O}_{S', s'})
$$
qui contient l'unique point de
$X \times_S \Spec(\mathcal{O}_{S', s'})$ s'envoyant sur $x$,
tels que l'image réciproque de $\mathcal{F}$ sur $V$ soit un $\mathcal{O}_V$-module
de présentation finie et plat sur $\mathcal{O}_{S', s'}$.
\end{proposition}

\begin{proof}[Première démonstration]
Cette démonstration est plus longue, mais n'utilise pas l'existence d'un d\'evissage complet.
Le problème est local au voisinage de $x$ et de $s$ ; on peut donc supposer que $X$
et $S$ sont affines. Au cours de la démonstration, on remplacera un nombre fini de fois
$S$ par un voisinage \'etale élémentaire de $(S, s)$. Le but sera alors de trouver,
après un tel remplacement, un ouvert
$V \subset X \times_S \Spec(\mathcal{O}_{S, s})$ contenant $x$
tel que $\mathcal{F}|_V$ soit plat sur $S$ et de présentation finie.
On peut évidemment aussi remplacer $S$ par $\Spec(\mathcal{O}_{S, s})$
à tout moment de la démonstration, c'est-à-dire supposer que $S$ est un schéma local.
On démontrera la proposition par récurrence sur l'entier
$n = \dim_x(\text{Supp}(\mathcal{F}_s))$.

\medskip\noindent
On peut choisir
\begin{enumerate}
\item des voisinages \'etales élémentaires $g : (X', x') \to (X, x)$,
$e : (S', s') \to (S, s)$,
\item un diagramme commutatif
$$
\xymatrix{
X \ar[dd]_f & X' \ar[dd] \ar[l]^g & Z' \ar[l]^i \ar[d]^\pi \\
& & Y' \ar[d]^h \\
S & S' \ar[l]_e & S' \ar@{=}[l]
}
$$
\item un point $z' \in Z'$ tel que $i(z') = x'$, $y' = \pi(z')$, $h(y') = s'$,
\item un $\mathcal{O}_{Z'}$-module quasi-cohérent de type fini $\mathcal{G}$,
\end{enumerate}
comme dans
le Lemme \ref{flat-lemma-elementary-devissage}.
Nous allons remplacer $S$ par $\Spec(\mathcal{O}_{S', s'})$, voir
les remarques du premier paragraphe de la démonstration. Considérons le diagramme
$$
\xymatrix{
X_{\mathcal{O}_{S', s'}} \ar[ddr]_f &
X'_{\mathcal{O}_{S', s'}} \ar[dd] \ar[l]^g &
Z'_{\mathcal{O}_{S', s'}} \ar[l]^i \ar[d]^\pi \\
& & Y'_{\mathcal{O}_{S', s'}} \ar[dl]^h \\
& \Spec(\mathcal{O}_{S', s'})
}
$$
On a ici effectué le changement de base des schémas $X', Z', Y'$ sur $S'$ suivant
$\Spec(\mathcal{O}_{S', s'}) \to S'$ et celui du schéma $X$ sur $S$ suivant
$\Spec(\mathcal{O}_{S', s'}) \to S$. Le morphisme
$g$ est encore \'etale, voir
le Lemme \ref{flat-lemma-etale-at-point}.
Après avoir remplacé $X$ par $X_{\mathcal{O}_{S', s'}}$,
$X'$ par $X'_{\mathcal{O}_{S', s'}}$,
$Z'$ par $Z'_{\mathcal{O}_{S', s'}}$ et
$Y'$ par $Y'_{\mathcal{O}_{S', s'}}$,
on peut supposer que l'on dispose d'un diagramme comme dans le
Lemme \ref{flat-lemma-elementary-devissage},
où, de plus, $S = S'$ est un schéma local de point fermé $s$. D'après les
Lemmes \ref{flat-lemma-devissage-finite-presentation} et
\ref{flat-lemma-devissage-flat},
le résultat pour $Y' \to S$, le faisceau $\pi_*\mathcal{G}$ et le
point $y'$ entraîne le résultat pour $X \to S$, $\mathcal{F}$ et $x$.
On peut donc supposer que $S$ est local et que $X \to S$ est un morphisme lisse
de schémas affines, à fibres géométriquement irréductibles de dimension $n$.

\medskip\noindent
Cas initial de la récurrence : $n = 0$.
Puisque $X \to S$ est lisse à fibres géométriquement
irréductibles de dimension $0$, on voit que $X \to S$ est une immersion
ouverte, voir
Descente, Lemme \ref{descent-lemma-universally-injective-etale-open-immersion}.
Puisque $S$ est local et que son point fermé appartient à l'image de $X \to S$,
on en conclut que $X = S$. Ainsi $\mathcal{F}$ correspond
à un $\mathcal{O}_{S, s}$-module plat de type fini. Dans ce cas, le résultat
découle de
Algèbre, Lemme \ref{algebra-lemma-finite-flat-local},
qui affirme que $\mathcal{F}$ est en fait libre de type fini.

\medskip\noindent
Étape de récurrence. Supposons le résultat établi lorsque la dimension
du support dans la fibre fermée est $< n$. Écrivons $S = \Spec(A)$,
$X = \Spec(B)$ et $\mathcal{F} = \widetilde{N}$ pour un certain $B$-module
$N$. Notons que $A$ est un anneau local ; désignons par $\mathfrak m$ son idéal maximal.
Alors $\mathfrak p = \mathfrak mB$ est l'unique idéal premier minimal au-dessus de
$\mathfrak m$, puisque $X \to S$ est à fibres géométriquement irréductibles. Enfin,
soit $\mathfrak q \subset B$ l'idéal premier correspondant à $x$. D'après le
Lemme \ref{flat-lemma-induction-step},
on peut choisir un homomorphisme
$$
\alpha : B^{\oplus r} \to N
$$
tel que $\kappa(\mathfrak p)^{\oplus r} \to N \otimes_B \kappa(\mathfrak p)$
soit un isomorphisme. De plus, puisque $N_{\mathfrak q}$ est $A$-plat, le lemme
montre aussi que $\alpha$ est injectif et que
$\Coker(\alpha)_{\mathfrak q}$ est $A$-plat.
Posons $Q = \Coker(\alpha)$. Notons que le support de $Q/\mathfrak mQ$
ne contient pas $\mathfrak p$. On a donc certainement
$\dim_{\mathfrak q}(\text{Supp}(Q/\mathfrak mQ)) < n$.
En réunissant les propriétés connues de $Q$, on voit
que l'hypothèse de récurrence s'applique à $Q$. Il existe donc
un morphisme \'etale élémentaire $(S', s) \to (S, s)$ tel que la conclusion
soit vérifiée pour $Q \otimes_A A'$ sur $B \otimes_A A'$, où
$A' = \mathcal{O}_{S', s'}$. Après avoir remplacé $A$ par $A'$, on dispose d'une
suite exacte
$$
0 \to B^{\oplus r} \to N \to Q \to 0
$$
(on utilise ici l'injectivité de $\alpha$ indiquée plus haut)
de $B$-modules de type fini, et l'on obtient aussi un élément
$g \in B$, $g \not \in \mathfrak q$, tel que
$Q_g$ soit de présentation finie sur $B_g$ et plat sur $A$. La localisation
étant exacte, on voit que
$$
0 \to B_g^{\oplus r} \to N_g \to Q_g \to 0
$$
est encore exacte. Puisque $B_g$ et $Q_g$ sont plats sur $A$, on en conclut que
$N_g$ est plat sur $A$, voir
Algèbre, Lemme \ref{algebra-lemma-flat-ses},
et puisque $B_g$ et $Q_g$ sont de présentation finie sur $B_g$, il en est de même
de $N_g$, voir
Algèbre, Lemme \ref{algebra-lemma-extension}.
\end{proof}

\begin{proof}[Deuxième démonstration]
On applique la
Proposition \ref{flat-proposition-existence-complete-at-x}
pour trouver un diagramme commutatif
$$
\xymatrix{
(X, x) \ar[d] & (X', x') \ar[l]^g \ar[d] \\
(S, s) & (S', s') \ar[l]
}
$$
de schémas pointés dont les flèches
horizontales sont des voisinages \'etales élémentaires
et tel que $g^*\mathcal{F}/X'/S'$ admette un
d\'evissage complet en $x$.
(En particulier, $S'$ et $X'$ sont affines.) D'après
Morphismes, Lemme \ref{morphisms-lemma-flat-permanence},
on voit que $g^*\mathcal{F}$ est plat en $x'$ sur $S$ et, d'après le
Lemme \ref{flat-lemma-etale-flat-up-down},
on voit qu'il est plat en $x'$ sur $S'$. Au moyen de la
Remarque \ref{flat-remark-same-notion},
on en déduit que
$$
\Gamma(X', g^*\mathcal{F})/
\Gamma(X', \mathcal{O}_{X'})/
\Gamma(S', \mathcal{O}_{S'})
$$
admet un d\'evissage complet en l'idéal premier de $\Gamma(X', \mathcal{O}_{X'})$
correspondant à $x'$. On peut effectuer le changement de base de ce
d\'evissage complet à l'anneau local $\mathcal{O}_{S', s'}$
de $\Gamma(S', \mathcal{O}_{S'})$ en l'idéal premier correspondant à
$s'$. Le
Lemme \ref{flat-lemma-complete-devissage-flat-finite-type-module}
entraîne donc que
$$
\Gamma(X', \mathcal{F}')
\otimes_{\Gamma(S', \mathcal{O}_{S'})}
\mathcal{O}_{S', s'}
$$
est plat sur $\mathcal{O}_{S', s'}$ et de présentation finie sur
$\Gamma(X', \mathcal{O}_{X'})
\otimes_{\Gamma(S', \mathcal{O}_{S'})} \mathcal{O}_{S', s'}$.
Autrement dit, la restriction de $\mathcal{F}$ à
$X' \times_{S'} \Spec(\mathcal{O}_{S', s'})$
est de présentation finie et plate sur $\mathcal{O}_{S', s'}$.
Puisque le morphisme
$X' \times_{S'} \Spec(\mathcal{O}_{S', s'})
\to X \times_S \Spec(\mathcal{O}_{S', s'})$
est \'etale
(Lemme \ref{flat-lemma-etale-at-point}),
son image $V \subset X \times_S \Spec(\mathcal{O}_{S', s'})$
est un sous-schéma ouvert et, par descente \'etale, la restriction
de $\mathcal{F}$ à $V$ est de présentation finie et plate sur
$\mathcal{O}_{S', s'}$. (Résultats utilisés :
Morphismes, Lemme \ref{morphisms-lemma-etale-open},
Descente, Lemme \ref{descent-lemma-finite-presentation-descends}, et
Morphismes, Lemme \ref{morphisms-lemma-flat-permanence}.)
\end{proof}

\begin{lemma}
\label{flat-lemma-open-in-fibre-where-flat}
Soit $f : X \to S$ un morphisme de schémas localement de type fini.
Soit $\mathcal{F}$ un $\mathcal{O}_X$-module quasi-cohérent de type fini.
Soit $s \in S$. Alors l'ensemble
$$
\{x \in X_s \mid \mathcal{F} \text{ plat sur }S\text{ en }x\}
$$
est ouvert dans la fibre $X_s$.
\end{lemma}

\begin{proof}
Supposons $x \in U$. Choisissons un voisinage \'etale élémentaire
$(S', s') \to (S, s)$ et un ouvert
$V \subset X \times_S \Spec(\mathcal{O}_{S', s'})$ comme dans la
Proposition \ref{flat-proposition-finite-type-flat-at-point}.
Notons que $X_{s'} = X_s$, puisque $\kappa(s) = \kappa(s')$.
Si $x' \in V \cap X_{s'}$, alors l'image réciproque de $\mathcal{F}$ sur
$X \times_S S'$ est plate sur $S'$ en $x'$. Ainsi $\mathcal{F}$ est
plat en $x'$ sur $S$, voir
Morphismes, Lemme \ref{morphisms-lemma-flat-permanence}.
Autrement dit, $X_s \cap V \subset U$ est un voisinage ouvert
de $x$ dans $U$.
\end{proof}

\begin{lemma}
\label{flat-lemma-finite-type-flat-at-point}
Soit $f : X \to S$ un morphisme de schémas.
Soit $\mathcal{F}$ un faisceau quasi-cohérent sur $X$.
Soit $x \in X$ d'image $s \in S$.
Supposons que
\begin{enumerate}
\item $f$ soit localement de type fini,
\item $\mathcal{F}$ soit de type fini, et
\item $\mathcal{F}$ soit plat en $x$ sur $S$.
\end{enumerate}
Il existe alors un voisinage \'etale élémentaire $(S', s') \to (S, s)$
et un sous-schéma ouvert
$$
V \subset X \times_S \Spec(\mathcal{O}_{S', s'})
$$
qui contient l'unique point de
$X \times_S \Spec(\mathcal{O}_{S', s'})$ s'envoyant sur $x$,
tels que l'image réciproque de $\mathcal{F}$ sur $V$ soit plate sur
$\mathcal{O}_{S', s'}$.
\end{lemma}

\begin{proof}
(La seule différence avec la
Proposition \ref{flat-proposition-finite-type-flat-at-point}
est que l'on ne suppose pas $f$ de présentation finie.)
La question est locale sur $X$ et sur $S$ ; on peut donc supposer que $X$ et $S$
sont affines. Écrivons $X = \Spec(B)$, $S = \Spec(A)$ et
$B = A[x_1, \ldots, x_n]/I$. Autrement dit, on obtient une immersion fermée
$i : X \to \mathbf{A}^n_S$. Notons $t = i(x) \in \mathbf{A}^n_S$.
On peut appliquer la
Proposition \ref{flat-proposition-finite-type-flat-at-point}
à $\mathbf{A}^n_S \to S$, au faisceau $i_*\mathcal{F}$
et au point $t$. On obtient un voisinage
\'etale élémentaire $(S', s') \to (S, s)$ et un sous-schéma ouvert
$$
W \subset \mathbf{A}^n_{\mathcal{O}_{S', s'}}
$$
tels que l'image réciproque de $i_*\mathcal{F}$ sur $W$ soit plate sur
$\mathcal{O}_{S', s'}$. Cela signifie que
$V := W \cap \big(X \times_S \Spec(\mathcal{O}_{S', s'})\big)$
est le sous-schéma ouvert cherché.
\end{proof}

\begin{lemma}
\label{flat-lemma-finite-type-flat-along-fibre}
Soit $f : X \to S$ un morphisme de schémas.
Soit $\mathcal{F}$ un faisceau quasi-cohérent sur $X$.
Soit $s \in S$.
Supposons que
\begin{enumerate}
\item $f$ soit de présentation finie,
\item $\mathcal{F}$ soit de type fini, et
\item $\mathcal{F}$ soit plat sur $S$ en tout point de la fibre $X_s$.
\end{enumerate}
Il existe alors un voisinage \'etale élémentaire $(S', s') \to (S, s)$
et un sous-schéma ouvert
$$
V \subset X \times_S \Spec(\mathcal{O}_{S', s'})
$$
qui contient la fibre $X_s = X \times_S s'$, tels que l'image réciproque
de $\mathcal{F}$ sur $V$ soit un $\mathcal{O}_V$-module
de présentation finie et plat sur $\mathcal{O}_{S', s'}$.
\end{lemma}

\begin{proof}
Pour tout point $x \in X_s$, on peut utiliser la
Proposition \ref{flat-proposition-finite-type-flat-at-point}
pour trouver un voisinage \'etale élémentaire $(S_x, s_x) \to (S, s)$
et un ouvert $V_x \subset X \times_S \Spec(\mathcal{O}_{S_x, s_x})$
tels que $x \in X_s = X \times_S s_x$ appartienne à $V_x$ et que
l'image réciproque de $\mathcal{F}$ sur $V_x$ soit un
$\mathcal{O}_{V_x}$-module de présentation finie et plat sur
$\mathcal{O}_{S_x, s_x}$. En particulier, on peut considérer la fibre
$(V_x)_{s_x}$ comme un voisinage ouvert de $x$ dans $X_s$.
Puisque $X_s$ est quasi-compact, on peut trouver un nombre fini de points
$x_1, \ldots, x_n \in X_s$ tels que $X_s$ soit la réunion
des $(V_{x_i})_{s_{x_i}}$. Choisissons un voisinage \'etale élémentaire
$(S' , s') \to (S, s)$ qui domine chacun des voisinages
$(S_{x_i}, s_{x_i})$, voir
Compléments sur les morphismes,
Lemme \ref{more-morphisms-lemma-elementary-etale-neighbourhoods}.
Posons $V = \bigcup V_i$, où $V_i$ est l'image réciproque de l'ouvert
$V_{x_i}$ par le morphisme
$$
X \times_S \Spec(\mathcal{O}_{S', s'})
\longrightarrow
X \times_S \Spec(\mathcal{O}_{S_{x_i}, s_{x_i}})
$$
Par construction, $V$ contient $X_s$ et l'image réciproque
de $\mathcal{F}$ sur $V$ est un $\mathcal{O}_V$-module
de présentation finie et plat sur $\mathcal{O}_{S', s'}$.
\end{proof}

\begin{lemma}
\label{flat-lemma-finite-type-flat-along-fibre-variant}
Soit $f : X \to S$ un morphisme de schémas.
Soit $\mathcal{F}$ un faisceau quasi-cohérent sur $X$.
Soit $s \in S$.
Supposons que
\begin{enumerate}
\item $f$ soit de type fini,
\item $\mathcal{F}$ soit de type fini, et
\item $\mathcal{F}$ soit plat sur $S$ en tout point de la fibre $X_s$.
\end{enumerate}
Il existe alors un voisinage \'etale élémentaire $(S', s') \to (S, s)$
et un sous-schéma ouvert
$$
V \subset X \times_S \Spec(\mathcal{O}_{S', s'})
$$
qui contient la fibre $X_s = X \times_S s'$, tels que l'image réciproque
de $\mathcal{F}$ sur $V$ soit plate sur $\mathcal{O}_{S', s'}$.
\end{lemma}

\begin{proof}
(La seule différence avec le
Lemme \ref{flat-lemma-finite-type-flat-along-fibre}
est que l'on ne suppose pas $f$ de présentation finie.)
Pour tout point $x \in X_s$, on peut utiliser le
Lemme \ref{flat-lemma-finite-type-flat-at-point}
pour trouver un voisinage \'etale élémentaire $(S_x, s_x) \to (S, s)$
et un ouvert $V_x \subset X \times_S \Spec(\mathcal{O}_{S_x, s_x})$
tels que $x \in X_s = X \times_S s_x$ appartienne à $V_x$ et que
l'image réciproque de $\mathcal{F}$ sur $V_x$ soit plate sur
$\mathcal{O}_{S_x, s_x}$. En particulier, on peut considérer la fibre
$(V_x)_{s_x}$ comme un voisinage ouvert de $x$ dans $X_s$.
Puisque $X_s$ est quasi-compact, on peut trouver un nombre fini de points
$x_1, \ldots, x_n \in X_s$ tels que $X_s$ soit la réunion
des $(V_{x_i})_{s_{x_i}}$. Choisissons un voisinage \'etale élémentaire
$(S' , s') \to (S, s)$ qui domine chacun des voisinages
$(S_{x_i}, s_{x_i})$, voir
Compléments sur les morphismes,
Lemme \ref{more-morphisms-lemma-elementary-etale-neighbourhoods}.
Posons $V = \bigcup V_i$, où $V_i$ est l'image réciproque de l'ouvert
$V_{x_i}$ par le morphisme
$$
X \times_S \Spec(\mathcal{O}_{S', s'})
\longrightarrow
X \times_S \Spec(\mathcal{O}_{S_{x_i}, s_{x_i}})
$$
Par construction, $V$ contient $X_s$ et l'image réciproque
de $\mathcal{F}$ sur $V$ est plate sur $\mathcal{O}_{S', s'}$.
\end{proof}

\begin{lemma}
\label{flat-lemma-finite-type-flat-at-point-X}
Soit $S$ un schéma. Soit $X$ localement de type fini sur $S$.
Soit $x \in X$ d'image $s \in S$.
Si $X$ est plat en $x$ sur $S$, il existe alors un voisinage
\'etale élémentaire $(S', s') \to (S, s)$ et un sous-schéma ouvert
$$
V \subset X \times_S \Spec(\mathcal{O}_{S', s'})
$$
qui contient l'unique point de
$X \times_S \Spec(\mathcal{O}_{S', s'})$ s'envoyant sur $x$,
tels que $V \to \Spec(\mathcal{O}_{S', s'})$
soit plat et de présentation finie.
\end{lemma}

\begin{proof}
La question est locale sur $X$ et sur $S$ ; on peut donc supposer que $X$ et $S$
sont affines. Écrivons $X = \Spec(B)$, $S = \Spec(A)$ et
$B = A[x_1, \ldots, x_n]/I$. Autrement dit, on obtient une immersion fermée
$i : X \to \mathbf{A}^n_S$. Notons $t = i(x) \in \mathbf{A}^n_S$.
On peut appliquer la
Proposition \ref{flat-proposition-finite-type-flat-at-point}
à $\mathbf{A}^n_S \to S$, au faisceau $\mathcal{F} = i_*\mathcal{O}_X$
et au point $t$. On obtient un voisinage
\'etale élémentaire $(S', s') \to (S, s)$ et un sous-schéma ouvert
$$
W \subset \mathbf{A}^n_{\mathcal{O}_{S', s'}}
$$
tels que l'image réciproque de $i_*\mathcal{O}_X$ soit plate et de présentation
finie. Cela signifie que
$V := W \cap \big(X \times_S \Spec(\mathcal{O}_{S', s'})\big)$
est le sous-schéma ouvert cherché.
\end{proof}

\begin{lemma}
\label{flat-lemma-finite-type-flat-at-point-local}
Soit $f : X \to S$ un morphisme localement de présentation finie.
Soit $\mathcal{F}$ un $\mathcal{O}_X$-module quasi-cohérent de type fini.
Si $x \in X$ et si $\mathcal{F}$ est plat en $x$ sur $S$, alors
$\mathcal{F}_x$ est un $\mathcal{O}_{X, x}$-module de présentation finie.
\end{lemma}

\begin{proof}
Soit $s = f(x)$. D'après la
Proposition \ref{flat-proposition-finite-type-flat-at-point},
il existe un voisinage \'etale élémentaire $(S', s') \to (S, s)$
tel que l'image réciproque de $\mathcal{F}$ sur
$X \times_S \Spec(\mathcal{O}_{S', s'})$ soit de
présentation finie au voisinage du point $x' \in X_{s'} = X_s$
correspondant à $x$. L'homomorphisme d'anneaux
$$
\mathcal{O}_{X, x} \longrightarrow
\mathcal{O}_{X \times_S \Spec(\mathcal{O}_{S', s'}), x'}
=
\mathcal{O}_{X \times_S S', x'}
$$
est plat et local, comme localisé d'un homomorphisme d'anneaux \'etale. Par suite,
$\mathcal{F}_x$ est de présentation finie sur $\mathcal{O}_{X, x}$
par descente, voir
Algèbre, Lemme \ref{algebra-lemma-descend-properties-modules}
(et le fait qu'un homomorphisme local plat est fidèlement plat, voir
Algèbre, Lemme \ref{algebra-lemma-local-flat-ff}).
\end{proof}

\begin{lemma}
\label{flat-lemma-finite-type-flat-at-point-local-X}
Soit $f : X \to S$ un morphisme localement de type fini.
Soit $x \in X$ d'image $s \in S$. Si $f$ est plat en $x$ sur $S$, alors
$\mathcal{O}_{X, x}$ est essentiellement de présentation finie sur
$\mathcal{O}_{S, s}$.
\end{lemma}

\begin{proof}
On peut supposer $X$ et $S$ affines. Écrivons $X = \Spec(B)$,
$S = \Spec(A)$ et $B = A[x_1, \ldots, x_n]/I$.
Autrement dit, on obtient une immersion fermée $i : X \to \mathbf{A}^n_S$.
Notons $t = i(x) \in \mathbf{A}^n_S$. On peut appliquer
le Lemme \ref{flat-lemma-finite-type-flat-at-point-local}
à $\mathbf{A}^n_S \to S$, au faisceau $\mathcal{F} = i_*\mathcal{O}_X$
et au point $t$. On en conclut que $\mathcal{O}_{X, x}$ est
de présentation finie sur $\mathcal{O}_{\mathbf{A}^n_S, t}$,
ce qui donne l'assertion cherchée.
\end{proof}







\section{Extension de propriétés à partir d'un ouvert}
\label{flat-section-extending-properties}

\noindent
Nous réunissons dans cette section plusieurs résultats de la forme suivante : si $f : X \to S$
est un morphisme plat de schémas et si $f$ possède une certaine propriété au-dessus d'un
ouvert dense de $S$, alors $f$ possède cette même propriété au-dessus de tout $S$.

\begin{lemma}
\label{flat-lemma-flat-finite-type-finitely-presented-over-dense-open}
\begin{slogan}
Les prolongements $S$-plats et de type fini de modules de présentation finie
sur un (bon) ouvert sont aussi de présentation finie sur $X$.
\end{slogan}
Soit $f : X \to S$ un morphisme de schémas. Soit $\mathcal{F}$ un
$\mathcal{O}_X$-module quasi-cohérent. Soit $U \subset S$ un ouvert.
Supposons que
\begin{enumerate}
\item $f$ soit localement de présentation finie,
\item $\mathcal{F}$ soit de type fini et plat sur $S$,
\item $U \subset S$ soit rétrocompact et schématiquement dense,
\item $\mathcal{F}|_{f^{-1}U}$ soit de présentation finie.
\end{enumerate}
Alors $\mathcal{F}$ est de présentation finie.
\end{lemma}

\begin{proof}
Le problème est local sur $X$ et sur $S$ ; on peut donc supposer $X$ et $S$ affines.
Écrivons $S = \Spec(A)$ et $X = \Spec(B)$. Soit $N$ un $B$-module de type fini tel
que $\mathcal{F}$ soit le faisceau quasi-cohérent associé à $N$.
On a $U = D(f_1) \cup \ldots \cup D(f_n)$ pour certains $f_i \in A$, voir
Algèbre, Lemme \ref{algebra-lemma-qc-open}.
Puisque $U$ est schématiquement dense, l'homomorphisme
$A \to A_{f_1} \times \ldots \times A_{f_n}$ est injectif.
Choisissons un idéal premier $\mathfrak q \subset B$ au-dessus de $\mathfrak p \subset A$,
correspondant à $x \in X$ s'envoyant sur $s \in S$.
D'après le Lemme \ref{flat-lemma-finite-type-flat-at-point-local},
le module $N_\mathfrak q$ est de présentation finie sur $B_\mathfrak q$.
Choisissons une surjection $\varphi : B^{\oplus m} \to N$ de $B$-modules.
Choisissons $k_1, \ldots, k_t \in \Ker(\varphi)$ et
posons $N' = B^{\oplus m}/\sum Bk_j$. Il existe une surjection canonique
$N' \to N$, et $N$ est la limite inductive filtrante des $B$-modules $N'$
ainsi construits. On voit donc que l'on peut choisir
$k_1, \ldots, k_t$ de sorte que (a) $N'_{f_i} \cong N_{f_i}$, $i = 1, \ldots, n$,
et (b) $N'_\mathfrak q \cong N_\mathfrak q$.
Cela entraîne en particulier que $N'_\mathfrak q$ est plat sur $A$.
Par ouverture du lieu de platitude, voir
Algèbre, Théorème \ref{algebra-theorem-openness-flatness},
on en conclut qu'il existe $g \in B$, $g \not \in \mathfrak q$,
tel que $N'_g$ soit plat sur $A$. Considérons le diagramme commutatif
$$
\xymatrix{
N'_g \ar[r] \ar[d] & N_g \ar[d] \\
\prod N'_{gf_i} \ar[r] & \prod N_{gf_i}
}
$$
La flèche inférieure est un isomorphisme par le choix de $k_1, \ldots, k_t$.
La flèche verticale de gauche est injective, puisque
$A \to \prod A_{f_i}$ est injectif et que $N'_g$ est plat sur $A$.
La flèche horizontale supérieure est donc injective, donc un isomorphisme.
Cela démontre que $N_g$ est de présentation finie sur $B_g$.
On conclut en appliquant
Algèbre, Lemme \ref{algebra-lemma-cover}.
\end{proof}

\begin{lemma}
\label{flat-lemma-flat-finite-type-finitely-presented-over-dense-open-X}
Soit $f : X \to S$ un morphisme de schémas. Soit $U \subset S$ un ouvert.
Supposons que
\begin{enumerate}
\item $f$ soit localement de type fini et plat,
\item $U \subset S$ soit rétrocompact et schématiquement dense,
\item $f|_{f^{-1}U} : f^{-1}U \to U$ soit localement de présentation finie.
\end{enumerate}
Alors $f$ est localement de présentation finie.
\end{lemma}

\begin{proof}
La question est locale sur $X$ et sur $S$ ; on peut donc supposer $X$ et
$S$ affines. Choisissons une immersion fermée $i : X \to \mathbf{A}^n_S$
et appliquons
le Lemme \ref{flat-lemma-flat-finite-type-finitely-presented-over-dense-open}
à $i_*\mathcal{O}_X$. Nous omettons certains détails.
\end{proof}

\begin{lemma}
\label{flat-lemma-flat-finite-presentation-dimension-over-dense-open}
Soit $f : X \to S$ un morphisme de schémas plat et localement
de type fini. Soit $U \subset S$ un ouvert dense tel que
$X_U \to U$ soit de dimension relative $\leq e$, voir
Morphismes, Définition \ref{morphisms-definition-relative-dimension-d}.
Si, de plus, l'une des conditions suivantes est vérifiée :
\begin{enumerate}
\item $f$ est localement de présentation finie, ou
\item $U \subset S$ est rétrocompact,
\end{enumerate}
alors $f$ est de dimension relative $\leq e$.
\end{lemma}

\begin{proof}
Démonstration dans le cas (1). Soit $W \subset X$ le sous-schéma ouvert construit
et étudié dans Compléments sur les morphismes, Lemmes
\ref{more-morphisms-lemma-flat-finite-presentation-CM-open} et
\ref{more-morphisms-lemma-flat-finite-presentation-CM-pieces}.
Notons que tout point générique de toute fibre appartient à $W$ ;
il suffit donc de démontrer le résultat pour $W$. Puisque
$W = \bigcup_{d \geq 0} U_d$, il suffit de montrer que $U_d = \emptyset$
pour $d > e$. Puisque $f$ est plat et localement de présentation finie, il est ouvert ;
ainsi $f(U_d)$ est ouvert (Morphismes, Lemme \ref{morphisms-lemma-fppf-open}).
Si $U_d$ est non vide, on a donc $f(U_d) \cap U \not = \emptyset$,
comme souhaité.

\medskip\noindent
Démonstration dans le cas (2). On peut remplacer $S$ par son réduit. Alors $U$ est
schématiquement dense. Par suite, $f$ est localement de présentation finie
d'après le Lemme \ref{flat-lemma-flat-finite-type-finitely-presented-over-dense-open-X}.
On est ainsi ramené au cas (1).
\end{proof}

\begin{lemma}
\label{flat-lemma-proper-flat-finite-over-dense-open}
Soit $f : X \to S$ un morphisme de schémas plat et propre.
Soit $U \subset S$ un ouvert dense tel que $X_U \to U$ soit fini.
Si, de plus, $f$ est localement de présentation finie ou si
$U \subset S$ est rétrocompact, alors $f$ est fini.
\end{lemma}

\begin{proof}
D'après le Lemme \ref{flat-lemma-flat-finite-presentation-dimension-over-dense-open},
les fibres de $f$ sont de dimension nulle.
Ainsi $f$ est quasi-fini
(Morphismes, Lemme \ref{morphisms-lemma-locally-quasi-finite-rel-dimension-0}),
donc à fibres finies
(Morphismes, Lemme \ref{morphisms-lemma-quasi-finite}).
Par suite, $f$ est fini d'après
Compléments sur les morphismes, Lemme \ref{more-morphisms-lemma-characterize-finite}.
\end{proof}

\begin{lemma}
\label{flat-lemma-zariski}
Soient $f : X \to S$ un morphisme de schémas et $U \subset S$ un ouvert. Si
\begin{enumerate}
\item $f$ est séparé, localement de type fini et plat,
\item $f^{-1}(U) \to U$ est un isomorphisme, et
\item $U \subset S$ est rétrocompact et schématiquement dense,
\end{enumerate}
alors $f$ est une immersion ouverte.
\end{lemma}

\begin{proof}
D'après le Lemme \ref{flat-lemma-flat-finite-type-finitely-presented-over-dense-open-X},
le morphisme $f$ est localement de présentation finie.
L'image $f(X) \subset S$ est ouverte
(Morphismes, Lemme \ref{morphisms-lemma-fppf-open}) ;
on peut donc remplacer $S$ par $f(X)$. Il faut ainsi démontrer que
$f$ est un isomorphisme. On peut supposer $S$ affine. On peut se ramener
au cas où $X$ est quasi-compact, car il suffit de montrer
que tout ouvert quasi-compact $X' \subset X$ dont l'image est $S$
s'envoie isomorphiquement sur $S$. On peut donc supposer $f$ quasi-compact.
Toutes les fibres de $f$ sont de dimension $0$, voir
le Lemme \ref{flat-lemma-flat-finite-presentation-dimension-over-dense-open}.
Ainsi $f$ est quasi-fini, voir
Morphismes, Lemme \ref{morphisms-lemma-locally-quasi-finite-rel-dimension-0}.
Soit $s \in S$. Choisissons un voisinage \'etale
élémentaire $g : (T, t) \to (S, s)$ tel que $X \times_S T = V \amalg W$
avec $V \to T$ fini et $W_t = \emptyset$, voir
Compléments sur les morphismes, Lemme
\ref{more-morphisms-lemma-etale-splits-off-quasi-finite-part}.
Notons $\pi : V \amalg W \to T$ le morphisme donné. Puisque $\pi$ est
plat et localement de présentation finie, on voit que $\pi(V)$ est
ouvert dans $T$ (Morphismes, Lemme \ref{morphisms-lemma-fppf-open}).
Quitte à restreindre $T$, on peut supposer que $T = \pi(V)$.
Puisque $f$ est un isomorphisme au-dessus de $U$, on voit que $\pi$ est un
isomorphisme au-dessus de $g^{-1}U$. Puisque $\pi(V) = T$, cela entraîne
que $\pi^{-1}g^{-1}U$ est contenu dans $V$.
D'après Morphismes, Lemme
\ref{morphisms-lemma-flat-morphism-scheme-theoretically-dense-open},
on voit que $\pi^{-1}g^{-1}U \subset V \amalg W$ est schématiquement
dense. On en déduit donc que $W = \emptyset$.
Ainsi $X \times_S T = V$ est fini sur $T$.
Il en résulte que $f$ est fini (après avoir remplacé $S$ par un
voisinage ouvert de $s$), par exemple d'après
Descente, Lemme \ref{descent-lemma-descending-property-finite}.
Alors $f$ est fini localement libre
(Morphismes, Lemme \ref{morphisms-lemma-finite-flat})
et, après avoir restreint $S$ à un voisinage ouvert plus petit de $s$,
on voit que $f$ est fini localement libre d'un certain degré $d$
(Morphismes, Lemme \ref{morphisms-lemma-finite-locally-free}).
Mais $d = 1$, puisque le degré vaut $1$ au-dessus
de l'ouvert dense $U$. Ainsi $f$ est un isomorphisme.
\end{proof}





\section{Modules plats de présentation finie}
\label{flat-section-finitely-presented-flat}

\noindent
Dans certains cas, étant donnés un homomorphisme d'anneaux $R \to S$ de présentation finie et
un $S$-module de présentation finie $N$, la platitude de $N$ sur $R$ entraîne
que $N$ est projectif comme $R$-module, au moins après avoir remplacé $S$
par une extension \'etale. Nous réunissons dans cette section quelques résultats
de cette nature.

\begin{lemma}
\label{flat-lemma-induction-step-fp}
Soit $R$ un anneau. Soit $R \to S$ un homomorphisme d'anneaux de présentation finie,
plat et à fibres géométriquement intègres. Soit
$\mathfrak q \subset S$ un idéal premier au-dessus de l'idéal premier
$\mathfrak r \subset R$. Posons $\mathfrak p = \mathfrak r S$.
Soit $N$ un $S$-module de présentation finie.
Il existe $r \geq 0$ et un homomorphisme de $S$-modules
$$
\alpha : S^{\oplus r} \longrightarrow N
$$
tels que
$\alpha : \kappa(\mathfrak p)^{\oplus r} \to N \otimes_S \kappa(\mathfrak p)$
soit un isomorphisme. Pour tout tel $\alpha$, les conditions suivantes sont équivalentes :
\begin{enumerate}
\item $N_{\mathfrak q}$ est $R$-plat,
\item il existe $f \in R$, $f \not \in \mathfrak r$, tel que
$\alpha_f : S_f^{\oplus r} \to N_f$ soit $R_f$-universellement injectif et
un $g \in S$, $g \not \in \mathfrak q$, tel que $\Coker(\alpha)_g$
soit $R$-plat,
\item $\alpha_{\mathfrak r}$ est $R_{\mathfrak r}$-universellement injectif et
$\Coker(\alpha)_{\mathfrak q}$ est $R$-plat,
\item $\alpha_{\mathfrak r}$ est injectif et
$\Coker(\alpha)_{\mathfrak q}$ est $R$-plat,
\item $\alpha_{\mathfrak p}$ est un isomorphisme et
$\Coker(\alpha)_{\mathfrak q}$ est $R$-plat, et
\item $\alpha_{\mathfrak q}$ est injectif et
$\Coker(\alpha)_{\mathfrak q}$ est $R$-plat.
\end{enumerate}
\end{lemma}

\begin{proof}
Pour obtenir $\alpha$, posons
$r = \dim_{\kappa(\mathfrak p)} N \otimes_S \kappa(\mathfrak p)$ et choisissons
$x_1, \ldots, x_r \in N$ dont les images forment une base de
$N \otimes_S \kappa(\mathfrak p)$. Définissons
$\alpha(s_1, \ldots, s_r) = \sum s_i x_i$. Cela démontre l'existence.

\medskip\noindent
Fixons un choix de $\alpha$.
On peut appliquer
le Lemme \ref{flat-lemma-induction-step}
à l'homomorphisme
$\alpha_{\mathfrak r} : S_{\mathfrak r}^{\oplus r} \to N_{\mathfrak r}$.
On voit donc que (1), (3), (4), (5) et (6) sont équivalentes.
Puisqu'il est aussi clair que (2) implique (3), il reste seulement
à montrer que (1) implique (2).

\medskip\noindent
Supposons (1). Par ouverture du lieu de platitude, voir
Algèbre, Théorème \ref{algebra-theorem-openness-flatness},
l'ensemble
$$
U_1 = \{\mathfrak q' \subset S \mid N_{\mathfrak q'}\text{ est plat sur }R\}
$$
est ouvert dans $\Spec(S)$. Il contient $\mathfrak q$ par hypothèse,
donc $\mathfrak p$. Puisque $S^{\oplus r}$ et $N$ sont des
$S$-modules de présentation finie, l'ensemble
$$
U_2 = \{\mathfrak q' \subset S \mid
\alpha_{\mathfrak q'}\text{ est un isomorphisme}\}
$$
est ouvert dans $\Spec(S)$, voir
Algèbre, Lemme \ref{algebra-lemma-map-between-finitely-presented}.
Il contient $\mathfrak p$ d'après (5). Puisque $R \to S$
est de présentation finie et plat, l'application
$\Phi : \Spec(S) \to \Spec(R)$ est ouverte, voir
Algèbre, Proposition \ref{algebra-proposition-fppf-open}.
Pour tout idéal premier $\mathfrak r' \in \Phi(U_1 \cap U_2)$, on voit
qu'il existe un idéal premier $\mathfrak q'$ au-dessus de $\mathfrak r'$ tel que
$N_{\mathfrak q'}$ soit plat et que $\alpha_{\mathfrak q'}$ soit
un isomorphisme, ce qui entraîne que $\alpha \otimes \kappa(\mathfrak p')$
est un isomorphisme, où $\mathfrak p' = \mathfrak r' S$. Ainsi
$\alpha_{\mathfrak r'}$ est $R_{\mathfrak r'}$-universellement injectif
par l'implication (1) $\Rightarrow$ (3).
Par suite, si l'on choisit $f \in R$, $f \not \in \mathfrak r$, tel que
$D(f) \subset \Phi(U_1 \cap U_2)$, on en conclut que
$\alpha_f$ est $R_f$-universellement injectif, voir
Algèbre, Lemme \ref{algebra-lemma-universally-injective-check-stalks}.
Le même raisonnement montre aussi que, pour tout
$\mathfrak q' \in U_1 \cap \Phi^{-1}(\Phi(U_1 \cap U_2))$,
le module $\Coker(\alpha)_{\mathfrak q'}$ est $R$-plat.
Notons que $\mathfrak q \in U_1 \cap \Phi^{-1}(\Phi(U_1 \cap U_2))$.
On peut donc trouver $g \in S$, $g \not \in \mathfrak q$, tel
que $D(g) \subset U_1 \cap \Phi^{-1}(\Phi(U_1 \cap U_2))$,
ce qui donne le résultat voulu.
\end{proof}

\begin{lemma}
\label{flat-lemma-complete-devissage-flat-finitely-presented-module}
Soit $R \to S$ un homomorphisme d'anneaux de présentation finie.
Soit $N$ un $S$-module de présentation finie, plat sur $R$.
Soit $\mathfrak r \subset R$ un idéal premier.
Supposons qu'il existe un d\'evissage complet de $N/S/R$ au-dessus de $\mathfrak r$.
Il existe alors $f \in R$, $f \not \in \mathfrak r$,
tel que
$$
N_f \cong B_1^{\oplus r_1} \oplus \ldots \oplus B_n^{\oplus r_n}
$$
comme $R$-modules, où chaque $B_i$ est une $R_f$-algèbre lisse à fibres géométriquement
irréductibles. De plus, $N_f$ est projectif comme $R_f$-module.
\end{lemma}

\begin{proof}
Soit $(A_i, B_i, M_i, \alpha_i)_{i = 1, \ldots, n}$ le
d\'evissage complet donné. On démontre le lemme par récurrence sur $n$.
Notons que les assertions du lemme portent uniquement sur la structure
de $N$ comme $R$-module. On peut donc remplacer $N$ par $M_1$ et
considérer $M_1$ comme un $B_1$-module. Voir
la Remarque \ref{flat-remark-finite-presentation}
pour vérifier que $M_1$ est de présentation finie comme $B_1$-module. D'après le
Lemme \ref{flat-lemma-induction-step-fp},
on peut, après avoir remplacé $R$ par $R_f$ pour un certain
$f \in R$, $f \not \in \mathfrak r$, supposer que
l'homomorphisme $\alpha_1 : B_1^{\oplus r_1} \to M_1$ est $R$-universellement injectif.
Puisque $M_1$ et $B_1^{\oplus r_1}$ sont $R$-plats et de présentation finie comme
$B_1$-modules, on voit que $\Coker(\alpha_1)$ est $R$-plat
(Algèbre, Lemme \ref{algebra-lemma-ui-flat-domain})
et de présentation finie comme $B_1$-module. Notons que
$(A_i, B_i, M_i, \alpha_i)_{i = 2, \ldots, n}$ est un
d\'evissage complet de $\Coker(\alpha_1)$. L'hypothèse de récurrence
entraîne donc qu'après avoir remplacé
$R$ par $R_f$ pour un certain $f \in R$, $f \not \in \mathfrak r$,
on peut supposer que $\Coker(\alpha_1)$ possède une décomposition
comme dans le lemme et est projectif. En particulier,
$M_1 = B_1^{\oplus r_1} \oplus \Coker(\alpha_1)$.
Cela démontre l'assertion relative à la décomposition.
L'assertion de projectivité en résulte, puisque $B_1$ est projectif comme
$R$-module d'après le
Lemme \ref{flat-lemma-fibres-irreducible-flat-projective-nonnoetherian}.
\end{proof}

\begin{remark}
\label{flat-remark-complete-devissage-flat-finitely-presented-module}
Il existe une variante du
Lemme \ref{flat-lemma-complete-devissage-flat-finitely-presented-module}
où l'on affaiblit la condition de platitude en supposant seulement que $N$
est plat en un idéal premier donné $\mathfrak q$ au-dessus de $\mathfrak r$,
mais où l'on renforce la condition de d\'evissage en supposant
l'existence d'un d\'evissage complet {\it en $\mathfrak q$}. Comparer avec le
Lemme \ref{flat-lemma-complete-devissage-flat-finite-type-module}.
\end{remark}

\noindent
Voici le résultat principal de cette section.

\begin{proposition}
\label{flat-proposition-finite-presentation-flat-at-point}
Soit $f : X \to S$ un morphisme de schémas.
Soit $\mathcal{F}$ un faisceau quasi-cohérent sur $X$.
Soit $x \in X$ d'image $s \in S$.
Supposons que
\begin{enumerate}
\item $f$ soit localement de présentation finie,
\item $\mathcal{F}$ soit de présentation finie, et
\item $\mathcal{F}$ soit plat en $x$ sur $S$.
\end{enumerate}
Il existe alors un diagramme commutatif de schémas pointés
$$
\xymatrix{
(X, x) \ar[d] & (X', x') \ar[l]^g \ar[d] \\
(S, s) & (S', s') \ar[l]
}
$$
dont les flèches horizontales sont des voisinages \'etales élémentaires,
tel que $X'$, $S'$ soient affines et que
$\Gamma(X', g^*\mathcal{F})$ soit un
$\Gamma(S', \mathcal{O}_{S'})$-module projectif.
\end{proposition}

\begin{proof}
Par ouverture du lieu de platitude, voir
Compléments sur les morphismes, Théorème \ref{more-morphisms-theorem-openness-flatness},
on peut remplacer $X$ par un voisinage ouvert de $x$ et supposer que
$\mathcal{F}$ est plat sur $S$. On applique ensuite la
Proposition \ref{flat-proposition-existence-complete-at-x}
pour trouver un diagramme comme dans l'énoncé de la proposition,
tel que $g^*\mathcal{F}/X'/S'$ admette un d\'evissage complet au-dessus de $s'$.
(En particulier, $S'$ et $X'$ sont affines.) D'après
Morphismes, Lemme \ref{morphisms-lemma-flat-permanence},
on voit que $g^*\mathcal{F}$ est plat sur $S$ et, d'après le
Lemme \ref{flat-lemma-etale-flat-up-down},
on voit qu'il est plat sur $S'$. Au moyen de la
Remarque \ref{flat-remark-same-notion},
on en déduit que
$$
\Gamma(X', g^*\mathcal{F})/
\Gamma(X', \mathcal{O}_{X'})/
\Gamma(S', \mathcal{O}_{S'})
$$
admet un d\'evissage complet au-dessus de l'idéal premier de $\Gamma(S', \mathcal{O}_{S'})$
correspondant à $s'$. Le
Lemme \ref{flat-lemma-complete-devissage-flat-finitely-presented-module}
entraîne donc que le résultat de la proposition est vérifié après avoir remplacé
$S'$ par un voisinage ouvert standard de $s'$.
\end{proof}

\noindent
Dans la suite de cette section, nous démontrons plusieurs variantes
de ce résultat. La première est une version « globale ».

\begin{lemma}
\label{flat-lemma-finite-presentation-flat-along-fibre}
Soit $f : X \to S$ un morphisme de schémas.
Soit $\mathcal{F}$ un faisceau quasi-cohérent sur $X$.
Soit $s \in S$.
Supposons que
\begin{enumerate}
\item $f$ soit de présentation finie,
\item $\mathcal{F}$ soit de présentation finie, et
\item $\mathcal{F}$ soit plat sur $S$ en tout point de la fibre $X_s$.
\end{enumerate}
Il existe alors un voisinage \'etale élémentaire
$(S', s') \to (S, s)$ et un diagramme commutatif de schémas
$$
\xymatrix{
X \ar[d] & X' \ar[l]^g \ar[d] \\
S & S' \ar[l]
}
$$
tels que $g$ soit \'etale, $X_s \subset g(X')$, que les schémas
$X'$, $S'$ soient affines et que
$\Gamma(X', g^*\mathcal{F})$ soit un
$\Gamma(S', \mathcal{O}_{S'})$-module projectif.
\end{lemma}

\begin{proof}
Pour tout point $x \in X_s$, on peut utiliser la
Proposition \ref{flat-proposition-finite-presentation-flat-at-point}
pour trouver un diagramme commutatif
$$
\xymatrix{
(X, x) \ar[d] & (Y_x, y_x) \ar[l]^{g_x} \ar[d] \\
(S, s) & (S_x, s_x) \ar[l]
}
$$
dont les flèches horizontales sont des voisinages \'etales élémentaires,
tel que $Y_x$, $S_x$ soient affines et que
$\Gamma(Y_x, g_x^*\mathcal{F})$ soit un
$\Gamma(S_x, \mathcal{O}_{S_x})$-module projectif. En particulier,
$g_x(Y_x) \cap X_s$ est un voisinage ouvert de $x$ dans $X_s$.
Puisque $X_s$ est quasi-compact, on peut trouver un nombre fini de points
$x_1, \ldots, x_n \in X_s$ tels que $X_s$ soit la réunion
des $g_{x_i}(Y_{x_i}) \cap X_s$. Choisissons un voisinage \'etale élémentaire
$(S' , s') \to (S, s)$ qui domine chacun des voisinages
$(S_{x_i}, s_{x_i})$, voir
Compléments sur les morphismes,
Lemme \ref{more-morphisms-lemma-elementary-etale-neighbourhoods}.
On peut aussi supposer $S'$ affine.
Posons $X' = \coprod Y_{x_i} \times_{S_{x_i}} S'$ et munissons-le du
morphisme évident $g : X' \to X$.
Par construction, $g(X')$ contient $X_s$ et
$$
\Gamma(X', g^*\mathcal{F})
=
\bigoplus \Gamma(Y_{x_i}, g_{x_i}^*\mathcal{F})
\otimes_{\Gamma(S_{x_i}, \mathcal{O}_{S_{x_i}})}
\Gamma(S', \mathcal{O}_{S'}).
$$
C'est un $\Gamma(S', \mathcal{O}_{S'})$-module projectif, voir
Algèbre, Lemme \ref{algebra-lemma-ascend-properties-modules}.
\end{proof}

\noindent
Les deux lemmes suivants reformulent les résultats
précédents lorsque $\mathcal{F} = \mathcal{O}_X$.

\begin{lemma}
\label{flat-lemma-finite-presentation-flat-at-point-X}
Soit $f : X \to S$ localement de présentation finie.
Soit $x \in X$ d'image $s \in S$.
Si $f$ est plat en $x$ sur $S$, il existe alors un diagramme commutatif
de schémas pointés
$$
\xymatrix{
(X, x) \ar[d] & (X', x') \ar[l]^g \ar[d] \\
(S, s) & (S', s') \ar[l]
}
$$
dont les flèches horizontales sont des voisinages \'etales élémentaires,
tel que $X'$, $S'$ soient affines et que
$\Gamma(X', \mathcal{O}_{X'})$ soit un
$\Gamma(S', \mathcal{O}_{S'})$-module projectif.
\end{lemma}

\begin{proof}
C'est un cas particulier de la
Proposition \ref{flat-proposition-finite-presentation-flat-at-point}.
\end{proof}

\begin{lemma}
\label{flat-lemma-finite-presentation-flat-along-fibre-X}
Soit $f : X \to S$ de présentation finie.
Soit $s \in S$.
Si $X$ est plat sur $S$ en tout point de $X_s$,
il existe alors un voisinage \'etale élémentaire
$(S', s') \to (S, s)$ et un diagramme commutatif de schémas
$$
\xymatrix{
X \ar[d] & X' \ar[l]^g \ar[d] \\
S & S' \ar[l]
}
$$
avec $g$ \'etale, $X_s \subset g(X')$, tels que $X'$, $S'$
soient affines et que
$\Gamma(X', \mathcal{O}_{X'})$ soit un
$\Gamma(S', \mathcal{O}_{S'})$-module projectif.
\end{lemma}

\begin{proof}
C'est un cas particulier du
Lemme \ref{flat-lemma-finite-presentation-flat-along-fibre}.
\end{proof}

\noindent
Les lemmes suivants donnent des conséquences de la
Proposition \ref{flat-proposition-finite-presentation-flat-at-point}
lorsque l'on suppose seulement que le morphisme et le faisceau sont de type fini
(et pas nécessairement de présentation finie).

\begin{lemma}
\label{flat-lemma-finite-type-flat-at-point-free}
Soit $f : X \to S$ un morphisme de schémas.
Soit $\mathcal{F}$ un faisceau quasi-cohérent sur $X$.
Soit $x \in X$ d'image $s \in S$.
Supposons que
\begin{enumerate}
\item $f$ soit localement de présentation finie,
\item $\mathcal{F}$ soit de type fini, et
\item $\mathcal{F}$ soit plat en $x$ sur $S$.
\end{enumerate}
Il existe alors un voisinage \'etale élémentaire $(S', s') \to (S, s)$
et un diagramme commutatif de schémas pointés
$$
\xymatrix{
(X, x) \ar[d] & (X', x') \ar[l]^g \ar[d] \\
(S, s) & (\Spec(\mathcal{O}_{S', s'}), s') \ar[l]
}
$$
tels que $X' \to X \times_S \Spec(\mathcal{O}_{S', s'})$
soit \'etale, $\kappa(x) = \kappa(x')$, que le schéma $X'$ soit
affine de présentation finie sur $\mathcal{O}_{S', s'}$,
que le faisceau $g^*\mathcal{F}$ soit de présentation finie sur $\mathcal{O}_{X'}$
et que $\Gamma(X', g^*\mathcal{F})$ soit un
$\mathcal{O}_{S', s'}$-module libre.
\end{lemma}

\begin{proof}
Pour démontrer le lemme, on peut remplacer $(S, s)$ par un voisinage \'etale
élémentaire quelconque, et l'on peut aussi remplacer $S$ par
$\Spec(\mathcal{O}_{S, s})$. Ainsi, d'après la
Proposition \ref{flat-proposition-finite-type-flat-at-point},
on peut supposer que $\mathcal{F}$ est de présentation finie et plat sur
$S$ au voisinage de $x$. Dans ce cas, le résultat découle de la
Proposition \ref{flat-proposition-finite-presentation-flat-at-point},
car
Algèbre, Théorème \ref{algebra-theorem-projective-free-over-local-ring},
assure que projectif $=$ libre sur un anneau local.
\end{proof}

\begin{lemma}
\label{flat-lemma-finite-type-flat-at-point-free-variant}
Soit $f : X \to S$ un morphisme de schémas.
Soit $\mathcal{F}$ un faisceau quasi-cohérent sur $X$.
Soit $x \in X$ d'image $s \in S$.
Supposons que
\begin{enumerate}
\item $f$ soit localement de type fini,
\item $\mathcal{F}$ soit de type fini, et
\item $\mathcal{F}$ soit plat en $x$ sur $S$.
\end{enumerate}
Il existe alors un voisinage \'etale élémentaire $(S', s') \to (S, s)$
et un diagramme commutatif de schémas pointés
$$
\xymatrix{
(X, x) \ar[d] & (X', x') \ar[l]^g \ar[d] \\
(S, s) & (\Spec(\mathcal{O}_{S', s'}), s') \ar[l]
}
$$
tels que $X' \to X \times_S \Spec(\mathcal{O}_{S', s'})$
soit \'etale, $\kappa(x) = \kappa(x')$, que le schéma $X'$ soit
affine et que $\Gamma(X', g^*\mathcal{F})$ soit un
$\mathcal{O}_{S', s'}$-module libre.
\end{lemma}

\begin{proof}
(La seule différence avec le
Lemme \ref{flat-lemma-finite-type-flat-at-point-free}
est que l'on ne suppose pas $f$ de présentation finie.)
Le problème est local sur $X$ et sur $S$. On peut donc supposer $X$ et
$S$ affines, et écrire $X = \Spec(B)$ et $S = \Spec(A)$.
Puisque $B$ est une $A$-algèbre de type fini, on peut trouver une surjection
$A[x_1, \ldots, x_n] \to B$. Autrement dit, on peut choisir une immersion
fermée $i : X \to \mathbf{A}^n_S$. Posons $t = i(x)$ et
$\mathcal{G} = i_*\mathcal{F}$. Notons que $\mathcal{G}_t \cong \mathcal{F}_x$
comme $\mathcal{O}_{S, s}$-modules. Ainsi $\mathcal{G}$ est plat sur $S$ en $t$.
On applique
le Lemme \ref{flat-lemma-finite-type-flat-at-point-free}
au morphisme $\mathbf{A}^n_S \to S$, au point $t$ et au
faisceau $\mathcal{G}$. On peut donc trouver un
voisinage \'etale élémentaire $(S', s') \to (S, s)$
et un diagramme commutatif de schémas pointés
$$
\xymatrix{
(\mathbf{A}^n_S, t) \ar[d] & (Y, y) \ar[l]^h \ar[d] \\
(S, s) & (\Spec(\mathcal{O}_{S', s'}), s') \ar[l]
}
$$
tels que $Y \to \mathbf{A}^n_{\mathcal{O}_{S', s'}}$
soit \'etale, $\kappa(t) = \kappa(y)$, que le schéma $Y$ soit
affine et que $\Gamma(Y, h^*\mathcal{G})$ soit un
$\mathcal{O}_{S', s'}$-module projectif. Une solution du problème
initial est alors donnée par le sous-schéma fermé
$X' = Y \times_{\mathbf{A}^n_S} X$ de $Y$.
\end{proof}

\begin{lemma}
\label{flat-lemma-finite-type-flat-along-fibre-free}
Soit $f : X \to S$ un morphisme de schémas.
Soit $\mathcal{F}$ un faisceau quasi-cohérent sur $X$.
Soit $s \in S$.
Supposons que
\begin{enumerate}
\item $f$ soit de présentation finie,
\item $\mathcal{F}$ soit de type fini, et
\item $\mathcal{F}$ soit plat sur $S$ en tout point de $X_s$.
\end{enumerate}
Il existe alors un voisinage \'etale élémentaire $(S', s') \to (S, s)$
et un diagramme commutatif de schémas
$$
\xymatrix{
X \ar[d] & X' \ar[l]^g \ar[d] \\
S & \Spec(\mathcal{O}_{S', s'}) \ar[l]
}
$$
tels que $X' \to X \times_S \Spec(\mathcal{O}_{S', s'})$
soit \'etale, $X_s = g((X')_{s'})$, que le schéma $X'$ soit
affine de présentation finie sur $\mathcal{O}_{S', s'}$,
que le faisceau $g^*\mathcal{F}$ soit de présentation finie sur $\mathcal{O}_{X'}$
et que $\Gamma(X', g^*\mathcal{F})$ soit un
$\mathcal{O}_{S', s'}$-module libre.
\end{lemma}

\begin{proof}
Pour tout point $x \in X_s$, on peut utiliser le
Lemme \ref{flat-lemma-finite-type-flat-at-point-free}
pour trouver un voisinage \'etale élémentaire $(S_x , s_x) \to (S, s)$
et un diagramme commutatif
$$
\xymatrix{
(X, x) \ar[d] & (Y_x, y_x) \ar[l]^{g_x} \ar[d] \\
(S, s) & (\Spec(\mathcal{O}_{S_x, s_x}), s_x) \ar[l]
}
$$
tels que $Y_x \to X \times_S \Spec(\mathcal{O}_{S_x, s_x})$
soit \'etale, $\kappa(x) = \kappa(y_x)$, que le schéma $Y_x$ soit affine
de présentation finie sur $\mathcal{O}_{S_x, s_x}$, que le faisceau
$g_x^*\mathcal{F}$ soit de présentation finie sur $\mathcal{O}_{Y_x}$ et
que $\Gamma(Y_x, g_x^*\mathcal{F})$ soit un
$\mathcal{O}_{S_x, s_x}$-module libre. En particulier,
$g_x((Y_x)_{s_x})$ est un voisinage ouvert de $x$ dans $X_s$.
Puisque $X_s$ est quasi-compact, on peut trouver un nombre fini de points
$x_1, \ldots, x_n \in X_s$ tels que $X_s$ soit la réunion
des $g_{x_i}((Y_{x_i})_{s_{x_i}})$. Choisissons un voisinage \'etale élémentaire
$(S' , s') \to (S, s)$ qui domine chacun des voisinages
$(S_{x_i}, s_{x_i})$, voir
Compléments sur les morphismes,
Lemme \ref{more-morphisms-lemma-elementary-etale-neighbourhoods}.
Posons
$$
X' = \coprod Y_{x_i} \times_{\Spec(\mathcal{O}_{S_{x_i}, s_{x_i}})}
\Spec(\mathcal{O}_{S', s'})
$$
et munissons-le du morphisme évident $g : X' \to X$.
Par construction, $X_s = g(X'_{s'})$ et
$$
\Gamma(X', g^*\mathcal{F})
=
\bigoplus \Gamma(Y_{x_i}, g_{x_i}^*\mathcal{F})
\otimes_{\mathcal{O}_{S_{x_i}, s_{x_i}}}
\mathcal{O}_{S', s'}.
$$
C'est un $\mathcal{O}_{S', s'}$-module libre, comme somme directe
de changements de base de modules libres. Nous omettons quelques détails mineurs.
\end{proof}

\begin{lemma}
\label{flat-lemma-finite-type-flat-along-fibre-free-variant}
Soit $f : X \to S$ un morphisme de schémas.
Soit $\mathcal{F}$ un faisceau quasi-cohérent sur $X$.
Soit $s \in S$.
Supposons que
\begin{enumerate}
\item $f$ soit de type fini,
\item $\mathcal{F}$ soit de type fini, et
\item $\mathcal{F}$ soit plat sur $S$ en tout point de $X_s$.
\end{enumerate}
Il existe alors un voisinage \'etale élémentaire $(S', s') \to (S, s)$
et un diagramme commutatif de schémas
$$
\xymatrix{
X \ar[d] & X' \ar[l]^g \ar[d] \\
S & \Spec(\mathcal{O}_{S', s'}) \ar[l]
}
$$
tels que $X' \to X \times_S \Spec(\mathcal{O}_{S', s'})$
soit \'etale, $X_s = g((X')_{s'})$, que le schéma $X'$ soit affine
et que $\Gamma(X', g^*\mathcal{F})$ soit un
$\mathcal{O}_{S', s'}$-module libre.
\end{lemma}

\begin{proof}
(La seule différence avec le
Lemme \ref{flat-lemma-finite-type-flat-along-fibre-free}
est que l'on ne suppose pas $f$ de présentation finie.)
Pour tout point $x \in X_s$, on peut utiliser le
Lemme \ref{flat-lemma-finite-type-flat-at-point-free-variant}
pour trouver un voisinage \'etale élémentaire $(S_x , s_x) \to (S, s)$
et un diagramme commutatif
$$
\xymatrix{
(X, x) \ar[d] & (Y_x, y_x) \ar[l]^{g_x} \ar[d] \\
(S, s) & (\Spec(\mathcal{O}_{S_x, s_x}), s_x) \ar[l]
}
$$
tels que $Y_x \to X \times_S \Spec(\mathcal{O}_{S_x, s_x})$
soit \'etale, $\kappa(x) = \kappa(y_x)$, que le schéma $Y_x$ soit affine et
que $\Gamma(Y_x, g_x^*\mathcal{F})$ soit un
$\mathcal{O}_{S_x, s_x}$-module libre. En particulier,
$g_x((Y_x)_{s_x})$ est un voisinage ouvert de $x$ dans $X_s$.
Puisque $X_s$ est quasi-compact, on peut trouver un nombre fini de points
$x_1, \ldots, x_n \in X_s$ tels que $X_s$ soit la réunion
des $g_{x_i}((Y_{x_i})_{s_{x_i}})$. Choisissons un voisinage \'etale élémentaire
$(S' , s') \to (S, s)$ qui domine chacun des voisinages
$(S_{x_i}, s_{x_i})$, voir
Compléments sur les morphismes,
Lemme \ref{more-morphisms-lemma-elementary-etale-neighbourhoods}.
Posons
$$
X' = \coprod Y_{x_i} \times_{\Spec(\mathcal{O}_{S_{x_i}, s_{x_i}})}
\Spec(\mathcal{O}_{S', s'})
$$
et munissons-le du morphisme évident $g : X' \to X$.
Par construction, $X_s = g(X'_{s'})$ et
$$
\Gamma(X', g^*\mathcal{F})
=
\bigoplus \Gamma(Y_{x_i}, g_{x_i}^*\mathcal{F})
\otimes_{\mathcal{O}_{S_{x_i}, s_{x_i}}}
\mathcal{O}_{S', s'}.
$$
C'est un $\mathcal{O}_{S', s'}$-module libre, comme somme directe
de changements de base de modules libres.
\end{proof}




\section{Modules plats de type fini, II}
\label{flat-section-finite-type-flat-II}

\noindent
Nous aurons besoin du lemme suivant.

\begin{lemma}
\label{flat-lemma-weak-bourbaki-pre-pre}
Soit $R \to S$ un homomorphisme d'anneaux de présentation finie. Soit $N$ un
$S$-module de présentation finie. Soit $\mathfrak q \subset S$ un idéal premier
au-dessus de $\mathfrak p \subset R$. Posons
$\overline{S} = S \otimes_R \kappa(\mathfrak p)$,
$\overline{\mathfrak q} = \mathfrak q \overline{S}$ et
$\overline{N} = N \otimes_R \kappa(\mathfrak p)$. Alors
il existe un $g \in S$ tel que
$g \not \in \mathfrak q$ et que
$\overline{g} \in \mathfrak r$ pour tout
$\mathfrak r \in \text{Ass}_{\overline{S}}(\overline{N})$
tel que $\mathfrak r \not \subset \overline{\mathfrak q}$.
\end{lemma}

\begin{proof}
En effet, si $\text{Ass}_{\overline{S}}(\overline{N}) =
\{\mathfrak r_1, \ldots, \mathfrak r_n\}$
(la finitude résulte d'Algèbre, lemme \ref{algebra-lemma-finite-ass}),
alors, quitte à renuméroter, nous pouvons supposer que
$$
\mathfrak r_1 \subset \overline{\mathfrak q},
\ldots,
\mathfrak r_r \subset \overline{\mathfrak q}, \quad
\mathfrak r_{r + 1} \not \subset \overline{\mathfrak q},
\ldots,
\mathfrak r_n \not \subset \overline{\mathfrak q}
$$
Comme $\overline{\mathfrak q}$ est un idéal premier, le produit
$\mathfrak r_{r + 1} \ldots \mathfrak r_n$ n'est pas contenu dans
$\overline{\mathfrak q}$ ; nous pouvons donc choisir un élément
$a$ de $\overline{S}$ appartenant à
$\mathfrak r_{r + 1}, \ldots, \mathfrak r_n$ mais non à
$\overline{\mathfrak q}$.
S'il existe $g \in S$ d'image $a$, alors $g$
convient. En général, nous pouvons trouver un élément non nul
$\lambda \in \kappa(\mathfrak p)$
tel que $\lambda a$ soit l'image d'un $g \in S$.
\end{proof}

\noindent
Le lemme suivant admet une variante légèrement plus forte,
le lemme \ref{flat-lemma-weak-bourbaki}
ci-dessous.

\begin{lemma}
\label{flat-lemma-weak-bourbaki-pre}
Soit $R \to S$ un homomorphisme d'anneaux de présentation finie.
Soit $N$ un $S$-module de présentation finie
qui est plat comme $R$-module. Soit $M$ un $R$-module.
Soit $\mathfrak q$ un idéal premier de $S$ au-dessus de $\mathfrak p \subset R$.
Alors
$$
\mathfrak q \in \text{WeakAss}_S(M \otimes_R N)
\Leftrightarrow
\Big(
\mathfrak p \in \text{WeakAss}_R(M)
\text{ et }
\overline{\mathfrak q} \in \text{Ass}_{\overline{S}}(\overline{N})
\Big)
$$
Ici, $\overline{S} = S \otimes_R \kappa(\mathfrak p)$,
$\overline{\mathfrak q} = \mathfrak q \overline{S}$ et
$\overline{N} = N \otimes_R \kappa(\mathfrak p)$.
\end{lemma}

\begin{proof}
Choisissons $g \in S$ comme au lemme \ref{flat-lemma-weak-bourbaki-pre-pre}.
Appliquons la proposition \ref{flat-proposition-finite-presentation-flat-at-point}
au morphisme de schémas $\Spec(S_g) \to \Spec(R)$, au module quasi-cohérent
associé à $N_g$ et aux points
correspondant aux idéaux premiers $\mathfrak qS_g$ et $\mathfrak p$. La traduction
algébrique donne un diagramme commutatif d'anneaux
$$
\xymatrix{
S \ar[r] & S_g \ar[r] & S' \\
& R \ar[lu] \ar[u] \ar[r] & R' \ar[u]
}
\quad\quad
\xymatrix{
\mathfrak q \ar@{-}[r] \ar@{-}[rd] &
\mathfrak qS_g \ar@{-}[d] \ar@{-}[r] & \mathfrak q' \ar@{-}[d] \\
& \mathfrak p \ar@{-}[r] & \mathfrak p'
}
$$
muni des idéaux premiers indiqués ; les flèches horizontales sont \'etales,
et $N \otimes_S S'$ est projectif comme $R'$-module. Posons
$N' = N \otimes_S S'$, $M' = M \otimes_R R'$,
$\overline{S}' = S' \otimes_{R'} \kappa(\mathfrak q')$,
$\overline{\mathfrak q}' = \mathfrak q' \overline{S}'$,
et
$$
\overline{N}' = N' \otimes_{R'} \kappa(\mathfrak p') =
\overline{N} \otimes_{\overline{S}} \overline{S}'
$$
D'après le lemme \ref{flat-lemma-etale-weak-assassin-up-down}, nous avons
\begin{align*}
\text{WeakAss}_{S'}(M' \otimes_{R'} N') & =
(\Spec(S') \to \Spec(S))^{-1}\text{WeakAss}_S(M \otimes_R N) \\
\text{WeakAss}_{R'}(M') & =
(\Spec(R') \to \Spec(R))^{-1}\text{WeakAss}_R(M) \\
\text{Ass}_{\overline{S}'}(\overline{N}') & =
(\Spec(\overline{S}') \to \Spec(\overline{S}))^{-1}
\text{Ass}_{\overline{S}}(\overline{N})
\end{align*}
Utilisons Algèbre, lemme \ref{algebra-lemma-ass-weakly-ass},
pour $\overline{N}$ et $\overline{N}'$. En particulier, nous avons
\begin{align*}
\mathfrak q \in \text{WeakAss}_S(M \otimes_R N)
& \Leftrightarrow
\mathfrak q' \in \text{WeakAss}_{S'}(M' \otimes_{R'} N') \\
\mathfrak p \in \text{WeakAss}_R(M)
& \Leftrightarrow
\mathfrak p' \in \text{WeakAss}_{R'}(M') \\
\overline{\mathfrak q} \in \text{Ass}_{\overline{S}}(\overline{N})
& \Leftrightarrow
\overline{\mathfrak q}' \in \text{WeakAss}_{\overline{S}'}(\overline{N}')
\end{align*}
Notre choix précis de $g$ et la formule donnant
$\text{Ass}_{\overline{S}'}(\overline{N}')$ ci-dessus montrent que
\begin{equation}
\label{flat-equation-key-observation}
\text{si }\mathfrak r' \in \text{Ass}_{\overline{S}'}(\overline{N}')
\text{ est au-dessus de }\mathfrak r \subset \overline{S}\text{, alors }
\mathfrak r \subset \overline{\mathfrak q}
\end{equation}
Cette observation sera essentielle dans la suite de la démonstration. Nous utiliserons
la caractérisation des idéaux premiers faiblement associés donnée dans
Algèbre, lemme \ref{algebra-lemma-weakly-ass-local}, sans autre mention.

\medskip\noindent
Supposons que
$\overline{\mathfrak q} \not \in \text{Ass}_{\overline{S}}(\overline{N})$.
Alors
$\overline{\mathfrak q}' \not \in \text{Ass}_{\overline{S}'}(\overline{N}')$.
D'après
Algèbre, lemmes \ref{algebra-lemma-ass-zero-divisors},
\ref{algebra-lemma-finite-ass} et
\ref{algebra-lemma-silly},
il existe un élément $\overline{a}' \in \overline{\mathfrak q}'$
qui n'est pas diviseur de zéro sur $\overline{N}'$.
Quitte à remplacer $\overline{a}'$ par $\lambda \overline{a}'$ pour un élément non nul
$\lambda \in \kappa(\mathfrak p)$, nous pouvons trouver
$a' \in \mathfrak q'$ d'image $\overline{a}'$. D'après le
lemme \ref{flat-lemma-invert-universally-injective},
l'application $a' : N'_{\mathfrak p'} \to N'_{\mathfrak p'}$ est
$R'_{\mathfrak p'}$-universellement injective. En particulier,
$a' : M' \otimes_{R'} N' \to M' \otimes_{R'} N'$ est
injective après localisation en $\mathfrak p'$, donc après
localisation en $\mathfrak q'$. Il en résulte clairement que
$\mathfrak q' \not \in \text{WeakAss}_{S'}(M' \otimes_{R'} N')$.
Nous concluons que $\mathfrak q \in \text{WeakAss}_S(M \otimes_R N)$ entraîne
$\overline{\mathfrak q} \in \text{Ass}_{\overline{S}}(\overline{N})$.

\medskip\noindent
Supposons $\mathfrak q \in \text{WeakAss}_S(M \otimes_R N)$. Nous voulons
montrer que $\mathfrak p \in \text{WeakAss}_S(M)$.
Soit $z \in M \otimes_R N$ un élément tel que $\mathfrak q$
soit minimal au-dessus de $J = \text{Ann}_S(z)$.
Soient $f_i \in \mathfrak p$, $i \in I$, des générateurs de
l'idéal $\mathfrak p$. Comme $\mathfrak q$ est au-dessus de $\mathfrak p$, pour tout $i$,
nous pouvons choisir $n_i \geq 1$ et $g_i \in S$, $g_i \not \in \mathfrak q$,
tels que $g_i f_i^{n_i} \in J$, c'est-à-dire $g_i f_i^{n_i} z = 0$.
Soit $z' \in (M' \otimes_{R'} N')_{\mathfrak p'}$ l'image de $z$.
Observons que $z'$ est non nul car $z$ a une image non nulle dans
$(M \otimes_R N)_\mathfrak q$ et car $S_\mathfrak q \to S'_{\mathfrak q'}$
est fidèlement plat. Nous affirmons que $f_i^{n_i} z' = 0$.

\medskip\noindent
Démontrons cette assertion : soit $g'_i \in S'$ l'image de $g_i$.
D'après l'observation essentielle (\ref{flat-equation-key-observation}),
l'image $\overline{g}'_i \in \overline{S}'$
n'appartient à $\mathfrak r'$ pour aucun
$\mathfrak r' \in \text{Ass}_{\overline{S}'}(\overline{N})$.
Ainsi, d'après le lemme \ref{flat-lemma-invert-universally-injective},
l'application $g'_i : N'_{\mathfrak p'} \to N'_{\mathfrak p'}$ est
$R'_{\mathfrak p'}$-universellement injective. En particulier,
$g'_i : M' \otimes_{R'} N' \to M' \otimes_{R'} N'$ est
injective après localisation en $\mathfrak p'$. L'assertion
en résulte puisque $g_i f_i^{n_i} z' = 0$.

\medskip\noindent
Cette assertion montre que l'annulateur de $z'$ dans $R'_{\mathfrak p'}$
contient les éléments $f_i^{n_i}$. Comme $R \to R'$ est \'etale, nous avons
$\mathfrak p'R'_{\mathfrak p'} = \mathfrak pR'_{\mathfrak p'}$
d'après Algèbre, lemme \ref{algebra-lemma-etale-at-prime}.
L'annulateur de $z'$ dans $R'_{\mathfrak p'}$ a donc pour radical
$\mathfrak p' R_{\mathfrak p'}$ (nous utilisons ici que $z'$ est non nul).
D'autre part,
$$
z' \in (M' \otimes_{R'} N')_{\mathfrak p'} =
M'_{\mathfrak p'} \otimes_{R'_{\mathfrak p'}} N'_{\mathfrak p'}
$$
Le module $N'_{\mathfrak p'}$ est projectif sur l'anneau local
$R'_{\mathfrak p'}$, donc libre
(Algèbre, théorème \ref{algebra-theorem-projective-free-over-local-ring}).
Nous pouvons donc trouver un facteur direct libre de type fini $F' \subset N'_{\mathfrak p'}$
tel que $z' \in M'_{\mathfrak p'} \otimes_{R'_{\mathfrak p'}} F'$.
Si $F'$ est de rang $n$, nous en déduisons que
$\mathfrak p' R'_{\mathfrak p'} \in
\text{WeakAss}_{R'_{\mathfrak p'}}({M'_{\mathfrak p'}}^{\oplus n})$.
Cela entraîne
$\mathfrak p'R'_{\mathfrak p'} \in \text{WeakAss}(M'_{\mathfrak p'})$,
par exemple d'après Algèbre, lemme \ref{algebra-lemma-weakly-ass}.
Alors $\mathfrak p' \in \text{WeakAss}_{R'}(M')$,
ce qui donne à son tour $\mathfrak p \in \text{WeakAss}_R(M)$.
Cela achève la démonstration de l'implication
``$\Rightarrow$'' de l'équivalence du lemme.

\medskip\noindent
Supposons que $\mathfrak p \in \text{WeakAss}_R(M)$ et
$\overline{\mathfrak q} \in \text{Ass}_{\overline{S}}(\overline{N})$.
Nous voulons montrer que $\mathfrak q$ est faiblement associé à $M \otimes_R N$.
Notons que $\overline{\mathfrak q}'$ est un élément maximal de
$\text{Ass}_{\overline{S}'}(\overline{N}')$.
Cela résulte de (\ref{flat-equation-key-observation})
et du fait qu'il n'y a pas d'inclusions entre les idéaux premiers
de $\overline{S}'$ au-dessus de $\overline{\mathfrak q}$
(les fibres des morphismes \'etales étant discrètes,
Morphismes, lemme \ref{morphisms-lemma-etale-over-field}).
Ainsi, en remplaçant $R, S, \mathfrak p, \mathfrak q, M, N$ par
$R', S', \mathfrak p', \mathfrak q', M', N'$,
nous pouvons supposer, outre les hypothèses du lemme, que
\begin{enumerate}
\item $\mathfrak p \in \text{WeakAss}_R(M)$,
\item $\overline{\mathfrak q} \in \text{Ass}_{\overline{S}}(\overline{N})$,
\item $N$ est projectif comme $R$-module, et
\item $\overline{\mathfrak q}$ est maximal dans
$\text{Ass}_{\overline{S}}(\overline{N})$.
\end{enumerate}
Une réduction supplémentaire consiste à remplacer
$R, S, M, N$ par leurs localisés en $\mathfrak p$.
Cela fournit une condition supplémentaire, à savoir :
\begin{enumerate}
\item[(5)] $R$ est un anneau local d'idéal maximal $\mathfrak p$.
\end{enumerate}
Nous terminerons en montrant que (1) -- (5) entraînent
$\mathfrak q \in \text{WeakAss}(M \otimes_R N)$.

\medskip\noindent
Comme $R$ est local et $\mathfrak p \in \text{WeakAss}_R(M)$,
nous pouvons choisir un $y \in M$ dont l'annulateur $I$ a pour radical
$\mathfrak p$.
Écrivons $\overline{\mathfrak q} = (\overline{g}_1, \ldots, \overline{g}_n)$
pour des $\overline{g}_i \in \overline{S}$. Choisissons $g_i \in S$
d'image $\overline{g}_i$.
Alors $\mathfrak q = \mathfrak pS + g_1S + \ldots + g_nS$.
Considérons l'application
$$
\Psi : N/IN \longrightarrow (N/IN)^{\oplus n}, \quad
z \longmapsto (g_1z, \ldots, g_nz).
$$
C'est un homomorphisme de $R/I$-modules projectifs.
L'anneau local $R/I$ est auto-associé
(Compléments d'algèbre, définition \ref{more-algebra-definition-auto-ass})
car $\mathfrak p/I$ est localement nilpotent.
L'application $\Psi \otimes \kappa(\mathfrak p)$ n'est pas injective, car
$\overline{\mathfrak q} \in \text{Ass}_{\overline{S}}(\overline{N})$.
Ainsi, Compléments d'algèbre, lemme \ref{more-algebra-lemma-P-fPD-zero},
entraîne que $\Psi$ n'est pas injective. Choisissons $z \in N/IN$
non nul dans le noyau de $\Psi$. L'annulateur $J = \text{Ann}_S(z)$
contient $IS$ et $g_i$ par construction. Ainsi,
$\sqrt{J} \subset S$ contient $\mathfrak q$.
Soit $\mathfrak s \subset S$ un idéal premier minimal au-dessus de $J$.
Alors $\mathfrak q \subset \mathfrak s$,
$\mathfrak s$ est au-dessus de $\mathfrak p$, et
$\mathfrak s \in \text{WeakAss}_S(N/IN)$.
Ce dernier fait résulte de la définition des idéaux premiers faiblement associés.
Appliquons le sens ``$\Rightarrow$'' du lemme (déjà démontré)
à l'homomorphisme d'anneaux $R \to S$ et aux modules $R/I$ et $N$
pour en déduire que
$\overline{\mathfrak s} \in \text{Ass}_{\overline{S}}(\overline{N})$.
Comme $\overline{\mathfrak q} \subset \overline{\mathfrak s}$,
la maximalité de $\overline{\mathfrak q}$, voir la condition (4) ci-dessus,
entraîne que $\overline{\mathfrak q} = \overline{\mathfrak s}$.
Cela montre que $\mathfrak q = \mathfrak s$ et donne
la conclusion voulue.
\end{proof}

\begin{lemma}
\label{flat-lemma-bourbaki-finite-type-general-base-at-point}
Soit $S$ un schéma.
Soit $f : X \to S$ localement de type fini.
Soit $x \in X$ d'image $s \in S$.
Soit $\mathcal{F}$ un faisceau quasi-cohérent de type fini sur $X$.
Soit $\mathcal{G}$ un faisceau quasi-cohérent sur $S$.
Si $\mathcal{F}$ est plat en $x$ sur $S$, alors
$$
x \in \text{WeakAss}_X(\mathcal{F} \otimes_{\mathcal{O}_X} f^*\mathcal{G})
\Leftrightarrow
s \in \text{WeakAss}_S(\mathcal{G})
\text{ et }
x \in \text{Ass}_{X_s}(\mathcal{F}_s).
$$
\end{lemma}

\begin{proof}
Dans ce paragraphe, nous nous ramenons au cas où $f$ est de présentation finie.
La question est locale sur $X$ et $S$ ; nous pouvons donc supposer $X$ et $S$
affines. Écrivons $X = \Spec(B)$, $S = \Spec(A)$ et
$B = A[x_1, \ldots, x_n]/I$. Autrement dit, nous obtenons une immersion fermée
$i : X \to \mathbf{A}^n_S$ sur $S$. Notons $t = i(x) \in \mathbf{A}^n_S$.
Notons que $i_*\mathcal{F}$ est un faisceau quasi-cohérent de type fini sur
$\mathbf{A}^n_S$, plat en $t$ sur $S$, et que
$$
i_*(\mathcal{F} \otimes_{\mathcal{O}_X} f^*\mathcal{G}) =
i_*\mathcal{F} \otimes_{\mathcal{O}_{\mathbf{A}^n_S}} p^*\mathcal{G}
$$
où $p : \mathbf{A}^n_S \to S$ est la projection. Notons que
$t$ est un point faiblement associé de
$i_*(\mathcal{F} \otimes_{\mathcal{O}_X} f^*\mathcal{G})$
si et seulement si $x$ est un point faiblement associé de
$\mathcal{F} \otimes_{\mathcal{O}_X} f^*\mathcal{G}$, voir
Diviseurs, lemme \ref{divisors-lemma-weakly-associated-finite}.
De même, $x \in \text{Ass}_{X_s}(\mathcal{F}_s)$ si et seulement
si $t \in \text{Ass}_{\mathbf{A}^n_s}((i_*\mathcal{F})_s)$ (voir
Algèbre, lemme \ref{algebra-lemma-ass-quotient-ring}).
Il suffit donc de démontrer le lemme lorsque $X = \mathbf{A}^n_S$.
Nous pouvons ainsi supposer que $X \to S$ est de présentation finie.

\medskip\noindent
Dans ce paragraphe, nous nous ramenons au cas où $\mathcal{F}$ est de présentation finie
et plat sur $S$.
Choisissons un voisinage \'etale élémentaire $e : (S', s') \to (S, s)$
et un ouvert $V \subset X \times_S \Spec(\mathcal{O}_{S', s'})$
comme dans la proposition \ref{flat-proposition-finite-type-flat-at-point}.
Soit $x' \in X' = X \times_S S'$ l'unique point d'images $x$
et $s'$. Il suffit alors de démontrer l'énoncé pour
$X' \to S'$, $x'$, $s'$, $(X' \to X)^*\mathcal{F}$ et $e^*\mathcal{G}$, voir
le lemme \ref{flat-lemma-etale-weak-assassin-up-down}.
Soit $v \in V$ l'unique point d'image $x'$
et soit $s' \in \Spec(\mathcal{O}_{S', s'})$ le point fermé.
Alors $\mathcal{O}_{V, v} = \mathcal{O}_{X', x'}$
et $\mathcal{O}_{\Spec(\mathcal{O}_{S', s'}), s'} =
\mathcal{O}_{S', s'}$, et de même pour les fibres des images réciproques de
$\mathcal{F}$ et $\mathcal{G}$.
De plus, $V_{s'} \subset X'_{s'}$ est un sous-schéma ouvert.
Comme la propriété d'être un point faiblement associé
ne dépend que de la fibre du faisceau, nous pouvons
remplacer
$X' \to S'$, $x'$, $s'$, $(X' \to X)^*\mathcal{F}$ et $e^*\mathcal{G}$
par
$V \to \Spec(\mathcal{O}_{S', s'})$, $v$, $s'$, $(V \to X)^*\mathcal{F}$
et $(\Spec(\mathcal{O}_{S', s'}) \to S)^*\mathcal{G}$.
Nous pouvons donc supposer que $f$ est de présentation finie et que
$\mathcal{F}$ est de présentation finie et plat sur $S$.

\medskip\noindent
Supposons $f$ de présentation finie et
$\mathcal{F}$ de présentation finie et plat sur $S$.
Après restriction de $X$ et $S$ à des voisinages affines
de $x$ et $s$, ce cas est traité par le
lemme \ref{flat-lemma-weak-bourbaki-pre}.
\end{proof}

\begin{lemma}
\label{flat-lemma-weak-bourbaki}
Soit $R \to S$ un homomorphisme d'anneaux essentiellement de type fini.
Soit $N$ un localisé d'un $S$-module de type fini, plat sur $R$.
Soit $M$ un $R$-module. Alors
$$
\text{WeakAss}_S(M \otimes_R N)
=
\bigcup\nolimits_{\mathfrak p \in \text{WeakAss}_R(M)}
\text{Ass}_{S \otimes_R \kappa(\mathfrak p)}(N \otimes_R \kappa(\mathfrak p))
$$
\end{lemma}

\begin{proof}
Ce lemme est la traduction
du lemme \ref{flat-lemma-bourbaki-finite-type-general-base-at-point}
en termes algébriques. Les détails de cette traduction sont omis.
\end{proof}

\begin{lemma}
\label{flat-lemma-bourbaki-finite-type-general-base}
Soit $f : X \to S$ un morphisme localement de type fini.
Soit $\mathcal{F}$ un faisceau quasi-cohérent de type fini sur $X$
qui est plat sur $S$. Soit $\mathcal{G}$ un faisceau quasi-cohérent sur $S$.
Alors nous avons
$$
\text{WeakAss}_X(\mathcal{F} \otimes_{\mathcal{O}_X} f^*\mathcal{G}) =
\bigcup\nolimits_{s \in \text{WeakAss}_S(\mathcal{G})}
\text{Ass}_{X_s}(\mathcal{F}_s)
$$
\end{lemma}

\begin{proof}
Conséquence immédiate du
lemme \ref{flat-lemma-bourbaki-finite-type-general-base-at-point}.
\end{proof}

\begin{theorem}
\label{flat-theorem-finite-type-flat}
\begin{slogan}
Le lieu de platitude est ouvert (version non noethérienne).
\end{slogan}
Soit $f : X \to S$ un morphisme de schémas.
Soit $\mathcal{F}$ un $\mathcal{O}_X$-module quasi-cohérent.
Supposons que
\begin{enumerate}
\item $X \to S$ soit localement de présentation finie,
\item $\mathcal{F}$ soit un $\mathcal{O}_X$-module de type fini, et
\item l'ensemble des points faiblement associés de $S$ soit localement fini dans $S$.
\end{enumerate}
Alors $U = \{x \in X \mid \mathcal{F}\text{ plat en }x\text{ sur }S\}$
est ouvert dans $X$ et $\mathcal{F}|_U$ est un $\mathcal{O}_U$-module
de présentation finie et plat sur $S$.
\end{theorem}

\begin{proof}
Soit $x \in X$ tel que $\mathcal{F}$ soit plat en $x$ sur $S$.
Il faut trouver un voisinage ouvert de $x$ sur lequel $\mathcal{F}$ se restreigne
en un module de présentation finie plat sur $S$.
Le problème est local sur $X$ et $S$ ; nous pouvons donc supposer $X$ et $S$
affines. Comme $\mathcal{F}_x$ est un
$\mathcal{O}_{X, x}$-module de présentation finie d'après le
lemme \ref{flat-lemma-finite-type-flat-at-point-local},
nous déduisons d'
Algèbre, lemme \ref{algebra-lemma-construct-fp-module-from-stalk},
qu'il existe un $\mathcal{O}_X$-module de présentation finie $\mathcal{F}'$
et une application $\varphi : \mathcal{F}' \to \mathcal{F}$ induisant
un isomorphisme $\varphi_x : \mathcal{F}'_x \to \mathcal{F}_x$. En particulier,
$\mathcal{F}'$ est plat sur $S$ en $x$ ; par ouverture
du lieu de platitude,
Compléments sur les morphismes, théorème \ref{more-morphisms-theorem-openness-flatness},
nous pouvons donc, après restriction de $X$, supposer que
$\mathcal{F}'$ est plat sur $S$. Comme $\mathcal{F}$ est de type fini,
nous pouvons, après restriction de $X$, supposer $\varphi$ surjective, voir
Modules, lemme \ref{modules-lemma-finite-type-surjective-on-stalk},
ou utiliser le lemme de Nakayama
(Algèbre, lemme \ref{algebra-lemma-NAK}).
D'après le
lemme \ref{flat-lemma-bourbaki-finite-type-general-base},
nous avons
$$
\text{WeakAss}_X(\mathcal{F}') \subset
\bigcup\nolimits_{s \in \text{WeakAss}(S)} \text{Ass}_{X_s}(\mathcal{F}'_s)
$$
Comme $\text{WeakAss}(S)$ est fini par hypothèse et puisque
$\text{Ass}_{X_s}(\mathcal{F}'_s)$ est fini d'après
Diviseurs, lemme \ref{divisors-lemma-finite-ass},
nous concluons que $\text{WeakAss}_X(\mathcal{F}')$ est fini. À l'aide d'
Algèbre, lemme \ref{algebra-lemma-silly},
nous pouvons, après une nouvelle restriction de $X$, supposer que
$\text{WeakAss}_X(\mathcal{F}')$ est contenu dans l'ensemble des généralisations
de $x$. Considérons maintenant $\mathcal{K} = \Ker(\varphi)$. Nous avons
$\text{WeakAss}_X(\mathcal{K}) \subset \text{WeakAss}_X(\mathcal{F}')$
(d'après
Diviseurs, lemme \ref{divisors-lemma-ses-weakly-ass}),
mais, d'autre part, $\varphi_x$ est un isomorphisme, donc $\varphi_{x'}$
est aussi un isomorphisme pour tout $x' \leadsto x$. Nous concluons que
$\text{WeakAss}_X(\mathcal{K}) = \emptyset$, d'où
$\mathcal{K} = 0$ d'après
Diviseurs, lemme \ref{divisors-lemma-weakly-ass-zero}.
\end{proof}

\begin{lemma}
\label{flat-lemma-finite-type-flat-algebra}
Soit $R \to S$ un homomorphisme d'anneaux de présentation finie.
Soit $M$ un $S$-module de type fini. Supposons $\text{WeakAss}_R(R)$ fini.
Alors
$$
U = \{\mathfrak q \subset S \mid M_{\mathfrak q}\text{ plat sur }R\}
$$
est ouvert dans $\Spec(S)$ et, pour tout $g \in S$ tel que
$D(g) \subset U$, le localisé $M_g$ est un
$S_g$-module de présentation finie, plat sur $R$.
\end{lemma}

\begin{proof}
Cela résulte immédiatement du
théorème \ref{flat-theorem-finite-type-flat}.
\end{proof}

\begin{lemma}
\label{flat-lemma-finite-type-flat-X}
Soit $f : X \to S$ un morphisme de schémas localement de type
fini. Supposons l'ensemble des points faiblement associés de $S$ localement fini
dans $S$. Alors l'ensemble des points $x \in X$ où $f$ est plat est un sous-schéma
ouvert $U \subset X$ et $U \to S$ est plat et localement de présentation
finie.
\end{lemma}

\begin{proof}
Le problème est local sur $X$ et $S$ ; nous pouvons donc supposer
$X$ et $S$ affines. Alors $X \to S$ correspond à un homomorphisme d'anneaux
de type fini $A \to B$. Choisissons une surjection $A[x_1, \ldots, x_n] \to B$
et considérons $B$ comme un $A[x_1, \ldots, x_n]$-module. Une application du
lemme \ref{flat-lemma-finite-type-flat-algebra}
achève la démonstration.
\end{proof}

\begin{lemma}
\label{flat-lemma-finite-type-flat-over-integral}
Soit $f : X \to S$ un morphisme de schémas
localement de type fini et plat. Si $S$ est intègre, alors $f$
est localement de présentation finie.
\end{lemma}

\begin{proof}
Cas particulier du
lemme \ref{flat-lemma-finite-type-flat-X}.
\end{proof}

\begin{proposition}
\label{flat-proposition-flat-finite-type-finite-presentation-domain}
Soit $R$ un anneau intègre. Soit $R \to S$ un homomorphisme d'anneaux de type fini.
Soit $M$ un $S$-module de type fini.
\begin{enumerate}
\item Si $S$ est plat sur $R$, alors $S$ est une $R$-algèbre de présentation finie.
\item Si $M$ est plat comme $R$-module, alors $M$ est de présentation finie
comme $S$-module.
\end{enumerate}
\end{proposition}

\begin{proof}
La partie (1) est un cas particulier du
lemme \ref{flat-lemma-finite-type-flat-over-integral}.
Pour la partie (2), choisissons une surjection $R[x_1, \ldots, x_n] \to S$.
D'après le lemme \ref{flat-lemma-finite-type-flat-algebra}, $M$
est de présentation finie comme $R[x_1, \ldots, x_n]$-module.
Nous concluons par Algèbre, lemme
\ref{algebra-lemma-finitely-presented-over-subring}.
\end{proof}

\begin{lemma}[Version de type fini du théorème \ref{flat-theorem-finite-type-flat}]
\label{flat-lemma-finite-type-flat}
Soit $f : X \to S$ un morphisme de schémas.
Soit $\mathcal{F}$ un $\mathcal{O}_X$-module quasi-cohérent.
Supposons que
\begin{enumerate}
\item $X \to S$ soit localement de type fini,
\item $\mathcal{F}$ soit un $\mathcal{O}_X$-module de type fini, et
\item l'ensemble des points faiblement associés de $S$ soit localement fini dans $S$.
\end{enumerate}
Alors $U = \{x \in X \mid \mathcal{F}\text{ plat en }x\text{ sur }S\}$
est ouvert dans $X$ et $\mathcal{F}|_U$ est plat sur $S$ et localement
de présentation finie relativement à $S$ (voir
Compléments sur les morphismes, définition
\ref{more-morphisms-definition-relatively-finitely-presented-sheaf}).
\end{lemma}

\begin{proof}
La question est locale sur $X$ et $S$. Nous pouvons donc supposer $X$ et $S$ affines.
Nous pouvons alors choisir une immersion fermée $i : X \to \mathbf{A}^n_S$.
Appliquons le théorème \ref{flat-theorem-finite-type-flat} à $X' = \mathbf{A}^n_S \to S$
et au module quasi-cohérent $\mathcal{F}' = i_*\mathcal{F}$ de type fini ;
nous trouvons que
$$
U' = \{x' \in X' \mid \mathcal{F}'\text{ plat en }x'\text{ sur }S\}
$$
est ouvert dans $X'$ et que $\mathcal{F}'|_{U'}$ est de présentation finie.
Comme $\mathcal{F}'$ se restreint à zéro sur $X' \setminus i(X)$ et
comme $\mathcal{F}'_{i(x)} \cong \mathcal{F}_x$ pour tout $x \in X$, nous voyons que
$$
U' = i(U) \amalg (X' \setminus i(X))
$$
Ainsi, $U = i^{-1}(U')$ est ouvert. De plus, il est clair que
$\mathcal{F}'|_{U'} = (i|_U)_*(\mathcal{F}|_U)$.
Nous en déduisons que $\mathcal{F}|_U$ est de présentation finie
relativement à $S$ d'après Compléments sur les morphismes, lemmes
\ref{more-morphisms-lemma-relative-finite-presentation} et
\ref{more-morphisms-lemma-finite-morphism-relative-finite-presentation}.
\end{proof}

\begin{lemma}
\label{flat-lemma-finite-type-flat-algebra-bis}
Soit $R \to S$ un homomorphisme d'anneaux de type fini.
Soit $M$ un $S$-module de type fini.
Supposons $\text{WeakAss}_R(R)$ fini.
Alors
$$
U = \{\mathfrak q \subset S \mid M_{\mathfrak q}\text{ plat sur }R\}
$$
est ouvert dans $\Spec(S)$ et, pour tout $g \in S$ tel que
$D(g) \subset U$, le localisé $M_g$ est plat sur $R$ et
est un $S_g$-module de présentation finie relativement à $R$ (voir
Compléments d'algèbre, définition
\ref{more-algebra-definition-relatively-finitely-presented}).
\end{lemma}

\begin{proof}
C'est le lemme \ref{flat-lemma-finite-type-flat} traduit en termes algébriques.
\end{proof}








\section{Exemples de modules relativement purs}
\label{flat-section-examples-pure-modules}

\noindent
Dans cette courte section, nous discutons quelques exemples de résultats qui motiveront
la notion de {\it module relativement pur} et le
concept d'{\it impureté} que nous introduirons plus loin. Chaque
exemple est énoncé sous forme de lemme. Notons l'analogie entre la condition sur
les idéaux premiers associés et les conditions figurant dans les
lemmes \ref{flat-lemma-base-change-universally-flat},
\ref{flat-lemma-universally-injective-to-completion},
\ref{flat-lemma-universally-injective-to-completion-flat} et
\ref{flat-lemma-flat-pure-over-complete-projective}.
Voir aussi
Algèbre, lemme \ref{algebra-lemma-compare-relative-assassins},
pour une discussion.

\begin{lemma}
\label{flat-lemma-explain-why-pure}
Soit $R$ un anneau local d'idéal maximal $\mathfrak m$.
Soit $R \to S$ un homomorphisme d'anneaux. Soit $N$ un $S$-module.
Supposons que
\begin{enumerate}
\item $N$ soit projectif comme $R$-module, et
\item $S/\mathfrak mS$ soit noethérien et $N/\mathfrak mN$ soit un
$S/\mathfrak mS$-module de type fini.
\end{enumerate}
Alors, pour tout idéal premier $\mathfrak q \subset S$ associé à
$N \otimes_R \kappa(\mathfrak p)$, où $\mathfrak p = R \cap \mathfrak q$,
nous avons $\mathfrak q + \mathfrak m S \not = S$.
\end{lemma}

\begin{proof}
Notons que les hypothèses des
lemmes \ref{flat-lemma-homothety-spectrum} et
\ref{flat-lemma-invert-universally-injective}
sont satisfaites. Nous utiliserons leurs conclusions sans autre
mention. Soit $\Sigma \subset S$ l'ensemble multiplicatif des éléments
qui ne sont pas diviseurs de zéro sur $N/\mathfrak mN$. L'application
$N \to \Sigma^{-1}N$ est $R$-universellement injective. Ainsi,
tout $\mathfrak q \subset S$ qui est un idéal premier associé à
$N \otimes_R \kappa(\mathfrak p)$ est aussi un idéal premier associé à
$\Sigma^{-1}N \otimes_R \kappa(\mathfrak p)$. Il en résulte clairement que
$\mathfrak q$ correspond à un idéal premier de $\Sigma^{-1}S$.
Ainsi, $\mathfrak q \subset \mathfrak q'$, où $\mathfrak q'$
correspond à un idéal premier associé à $N/\mathfrak mN$, ce qui conclut.
\end{proof}

\noindent
Le lemme suivant donne un autre exemple, un peu trivial, de ce phénomène.

\begin{lemma}
\label{flat-lemma-explain-why-pure-complete}
Soit $R$ un anneau. Soit $I \subset R$ un idéal.
Soit $R \to S$ un homomorphisme d'anneaux. Soit $N$ un $S$-module.
Si $N$ est complet pour la topologie $I$-adique, alors, pour tout $R$-module $M$ et
pour tout idéal premier $\mathfrak q \subset S$ associé à
$N \otimes_R M$, nous avons $\mathfrak q + I S \not = S$.
\end{lemma}

\begin{proof}
Notons $S^\wedge$ le complété $I$-adique de $S$.
Notons que $N$ est un $S^\wedge$-module, donc
$N \otimes_R M$ est aussi un $S^\wedge$-module.
Soit $z \in N \otimes_R M$ un élément tel que
$\mathfrak q = \text{Ann}_S(z)$. Comme $z \not = 0$, nous voyons
que $\text{Ann}_{S^\wedge}(z) \not = S^\wedge$. Ainsi,
$\mathfrak q S^\wedge \not = S^\wedge$. Il existe donc un
idéal maximal $\mathfrak m \subset S^\wedge$ tel que
$\mathfrak q S^\wedge \subset \mathfrak m$. Puisque
$IS^\wedge \subset \mathfrak m$ d'après
Algèbre, lemme \ref{algebra-lemma-radical-completion},
nous obtenons la conclusion.
\end{proof}

\noindent
Notons que le lemme suivant donne une autre démonstration du
lemme \ref{flat-lemma-explain-why-pure},
car un module projectif sur un anneau local est libre, voir
Algèbre, théorème \ref{algebra-theorem-projective-free-over-local-ring}.

\begin{lemma}
\label{flat-lemma-explain-why-pure-direct-sum-finite-modules}
Soit $R$ un anneau local d'idéal maximal $\mathfrak m$.
Soit $R \to S$ un homomorphisme d'anneaux. Soit $N$ un $S$-module.
Supposons $N$ isomorphe, comme $R$-module, à une somme
directe de $R$-modules de type fini. Alors, pour tout $R$-module $M$ et
pour tout idéal premier $\mathfrak q \subset S$ associé à
$N \otimes_R M$, nous avons $\mathfrak q + \mathfrak m S \not = S$.
\end{lemma}

\begin{proof}
Écrivons $N = \bigoplus_{i \in I} M_i$, chaque $M_i$ étant un $R$-module de type fini.
Soient $M$ un $R$-module et $\mathfrak q \subset S$ un idéal premier
associé à $N \otimes_R M$ tel que $\mathfrak q + \mathfrak m S = S$. Soit
$z \in N \otimes_R M$ un élément tel que $\mathfrak q = \text{Ann}_S(z)$.
En modifiant légèrement la décomposition en somme directe, nous pouvons supposer
$z \in M_1 \otimes_R M$ pour un élément $1 \in I$. Écrivons
$1 = f + \sum x_j g_j$ pour des $f \in \mathfrak q$, $x_j \in \mathfrak m$
et $g_j \in S$. Pour tout $g \in S$, notons $g'$ l'application $R$-linéaire
$$
M_1 \to N \xrightarrow{g} N \to M_1
$$
dont la première flèche est l'inclusion, la deuxième la multiplication
par $g$ et la troisième la projection. Comme chaque $x_j \in R$,
nous obtenons l'égalité
$$
f' + \sum x_j g'_j = \text{id}_{M_1} \in \text{End}_R(M_1)
$$
D'après le lemme de Nakayama
(Algèbre, lemme \ref{algebra-lemma-NAK}),
$f'$ est surjective ; ainsi, d'après
Algèbre, lemme \ref{algebra-lemma-fun},
$f'$ est un isomorphisme. En particulier, l'application
$$
M_1 \otimes_R M \to N \otimes_R M \xrightarrow{f} N \otimes_R M
\to M_1 \otimes_R M
$$
est un isomorphisme. Cela contredit l'hypothèse $fz = 0$.
\end{proof}

\begin{lemma}
\label{flat-lemma-explain-why-pure-ML}
Soit $R$ un anneau local hensélien d'idéal maximal $\mathfrak m$.
Soit $R \to S$ un homomorphisme d'anneaux. Soit $N$ un $S$-module.
Supposons $N$ engendré par un ensemble dénombrable et de Mittag-Leffler comme $R$-module.
Alors, pour tout $R$-module $M$ et pour tout idéal premier $\mathfrak q \subset S$
associé à $N \otimes_R M$, nous avons
$\mathfrak q + \mathfrak m S \not = S$.
\end{lemma}

\begin{proof}
Ce lemme se ramène au
lemme \ref{flat-lemma-explain-why-pure-direct-sum-finite-modules}
d'après
Algèbre, lemme \ref{algebra-lemma-split-ML-henselian}.
\end{proof}

\noindent
Supposons que $f : X \to S$ soit un morphisme de schémas et que
$\mathcal{F}$ soit un module quasi-cohérent sur $X$.
Soient $\xi \in \text{Ass}_{X/S}(\mathcal{F})$ et $Z = \overline{\{\xi\}}$.
Considérons le diagramme
$$
\xymatrix{
\xi \ar@{|->}[d] & Z \ar[r] \ar[rd] & X \ar[d]^f \\
f(\xi) & & S
}
$$
Notons que $f(Z) \subset \overline{\{f(\xi)\}}$ et que $f(Z)$ est fermé
si et seulement s'il y a égalité, c'est-à-dire $f(Z) = \overline{\{f(\xi)\}}$.
Il résulte du
lemme \ref{flat-lemma-explain-why-pure}
que, si $S$, $X$ sont affines, si les fibres $X_s$ sont noethériennes,
si $\mathcal{F}$ est de type fini et si $\Gamma(X, \mathcal{F})$
est un $\Gamma(S, \mathcal{O}_S)$-module projectif, alors
$f(Z) = \overline{\{f(\xi)\}}$ est une partie fermée.
Des énoncés analogues légèrement différents valent dans les cas décrits aux
lemmes \ref{flat-lemma-explain-why-pure-complete},
\ref{flat-lemma-explain-why-pure-direct-sum-finite-modules} et
\ref{flat-lemma-explain-why-pure-ML}.




\section{Impuretés}
\label{flat-section-impure}

\noindent
Nous voulons formaliser le phénomène dont nous avons donné des exemples dans la
section \ref{flat-section-examples-pure-modules}
en termes de spécialisations des points de $\text{Ass}_{X/S}(\mathcal{F})$.
Nous voulons aussi travailler localement au voisinage d'un point $s \in S$. À cette fin,
nous introduisons les définitions suivantes.

\begin{situation}
\label{flat-situation-pre-pure}
Ici, $S$, $X$ sont des schémas et $f : X \to S$ est un morphisme de type fini.
De plus, $\mathcal{F}$ est un $\mathcal{O}_X$-module quasi-cohérent de type fini.
Enfin, $s$ est un point de $S$.
\end{situation}

\noindent
Dans cette situation, considérons un morphisme $g : T \to S$, un point $t \in T$
tel que $g(t) = s$, une spécialisation $t' \leadsto t$ et un point
$\xi \in X_T$ du changement de base de $X$, au-dessus de $t'$. Diagramme :
\begin{equation}
\label{flat-equation-impurity}
\vcenter{
\xymatrix{
\xi \ar@{|->}[d] & \\
t' \ar@{~>}[r] & t \ar@{|->}[r] & s
}
}
\quad\quad
\vcenter{
\xymatrix{
X_T \ar[d] \ar[r] & X \ar[d] \\
T \ar[r]^g & S
}
}
\end{equation}
Notons en outre $\mathcal{F}_T$ l'image réciproque de $\mathcal{F}$ sur $X_T$.

\begin{definition}
\label{flat-definition-impurity}
Dans la
situation \ref{flat-situation-pre-pure},
nous disons qu'un diagramme (\ref{flat-equation-impurity}) définit une
{\it impureté de $\mathcal{F}$ au-dessus de $s$}
si $\xi \in \text{Ass}_{X_T/T}(\mathcal{F}_T)$ et
$\overline{\{\xi\}} \cap X_t = \emptyset$. Nous exprimerons
cela en disant : « soit $(g : T \to S, t' \leadsto t, \xi)$ une
impureté de $\mathcal{F}$ au-dessus de $s$ ».
\end{definition}

\begin{lemma}
\label{flat-lemma-impure-finite-presentation}
Plaçons-nous dans la situation \ref{flat-situation-pre-pure}.
S'il existe une impureté de $\mathcal{F}$ au-dessus de $s$, alors
il existe une impureté $(g : T \to S, t' \leadsto t, \xi)$
de $\mathcal{F}$ au-dessus de $s$ telle que $g$ soit localement de présentation
finie et $t$ soit un point fermé de la fibre de $g$ au-dessus de $s$.
\end{lemma}

\begin{proof}
Soit $(g : T \to S, t' \leadsto t, \xi)$ une impureté quelconque de
$\mathcal{F}$ au-dessus de $s$. Appliquons
Limites, lemme \ref{limits-lemma-separate},
à $t \in T$ et $Z = \overline{\{\xi\}}$ pour obtenir un voisinage ouvert
$V \subset T$ de $t$, un diagramme commutatif
$$
\xymatrix{
V \ar[d] \ar[r]_a & T' \ar[d]^b \\
T \ar[r]^g & S,
}
$$
et un sous-schéma fermé $Z' \subset X_{T'}$ tels que
\begin{enumerate}
\item le morphisme $b : T' \to S$ soit localement de présentation finie,
\item on ait $Z' \cap X_{a(t)} = \emptyset$, et
\item $Z \cap X_V$ soit envoyé dans $Z'$ par le morphisme $X_V \to X_{T'}$.
\end{enumerate}
Comme $t'$ se spécialise en $t$, nous pouvons remplacer $T$ par le voisinage ouvert
$V$ de $t$. Nous avons donc un diagramme commutatif
$$
\xymatrix{
X_T \ar[d] \ar[r] &
X_{T'} \ar[d] \ar[r] &
X \ar[d] \\
T \ar[r]^a & T' \ar[r]^b & S
}
$$
où $b \circ a = g$. Notons $\xi' \in X_{T'}$
l'image de $\xi$. D'après
Diviseurs, lemme \ref{divisors-lemma-base-change-relative-assassin},
nous avons $\xi' \in \text{Ass}_{X_{T'}/T'}(\mathcal{F}_{T'})$.
De plus, par construction, l'adhérence de $\overline{\{\xi'\}}$
est contenue dans la partie fermée $Z'$, qui ne rencontre pas la fibre
$X_{a(t)}$. Ainsi, $(T' \to S, a(t') \leadsto a(t), \xi')$
est une impureté de $\mathcal{F}$ au-dessus de $s$.

\medskip\noindent
Nous pouvons donc supposer $g : T \to S$ localement de présentation finie.
Soit $Z = \overline{\{\xi\}}$. Par hypothèse, $Z_t = \emptyset$. D'après
Compléments sur les morphismes, lemme \ref{more-morphisms-lemma-empty-generic-fibre},
cela signifie que $Z_{t''} = \emptyset$ pour $t''$ dans une partie ouverte
de $\overline{\{t\}}$. Comme la fibre de
$T \to S$ au-dessus de $s$ est un schéma de Jacobson, voir
Morphismes, lemme \ref{morphisms-lemma-ubiquity-Jacobson-schemes},
il existe un point fermé $t'' \in \overline{\{t\}}$ tel que
$Z_{t''} = \emptyset$. Alors $(g : T \to S, t' \leadsto t'', \xi)$ est
l'impureté cherchée.
\end{proof}

\begin{lemma}
\label{flat-lemma-impure-limit}
Plaçons-nous dans la situation \ref{flat-situation-pre-pure}.
Soit $(g : T \to S, t' \leadsto t, \xi)$ une impureté de
$\mathcal{F}$ au-dessus de $s$. Supposons que $T = \lim_{i \in I} T_i$
soit une limite dirigée de schémas affines sur $S$. Alors, pour
un certain $i$, le triplet $(T_i \to S, t'_i \leadsto t_i, \xi_i)$
est une impureté de $\mathcal{F}$ au-dessus de $s$.
\end{lemma}

\begin{proof}
Les notations de l'énoncé ont le sens suivant : soient $p_i : T \to T_i$
les projections, $t_i = p_i(t)$ et $t'_i = p_i(t')$.
Enfin, $\xi_i \in X_{T_i}$ est l'image de $\xi$. D'après
Diviseurs, lemme \ref{divisors-lemma-base-change-relative-assassin},
$\xi_i$ est bien un point de l'assassin
relatif de $\mathcal{F}_{T_i}$ sur $T_i$. Il reste donc seulement à
montrer que $\overline{\{\xi_i\}} \cap X_{t_i} = \emptyset$ pour un certain $i$.

\medskip\noindent
Première démonstration. Soient $Z_i = \overline{\{\xi_i\}} \subset X_{T_i}$
et $Z = \overline{\{\xi\}} \subset X_T$
munis des structures de sous-schémas réduits induites.
Alors $Z = \lim Z_i$ d'après
Limites, lemme \ref{limits-lemma-inverse-limit-irreducibles}.
Choisissons un corps $k$ et un morphisme $\Spec(k) \to T$ d'image $t$.
Alors
$$
\emptyset =
Z \times_T \Spec(k) = (\lim Z_i) \times_{(\lim T_i)} \Spec(k)
= \lim Z_i \times_{T_i} \Spec(k)
$$
car les limites commutent aux produits fibrés (les limites commutent aux limites).
Chaque $Z_i \times_{T_i} \Spec(k)$ est quasi-compact, car $X_{T_i} \to T_i$
est de type fini, donc $Z_i \to T_i$ est de type fini.
Ainsi, $Z_i \times_{T_i} \Spec(k)$ est vide pour un certain $i$ d'après
Limites, lemme \ref{limits-lemma-limit-nonempty}.
Comme l'image du composé $\Spec(k) \to T \to T_i$ est $t_i$,
nous obtenons la conclusion voulue.

\medskip\noindent
Deuxième démonstration. Posons $Z = \overline{\{\xi\}}$. Appliquons
Limites, lemme \ref{limits-lemma-separate},
à cette situation pour obtenir un voisinage ouvert
$V \subset T$ de $t$, un diagramme commutatif
$$
\xymatrix{
V \ar[d] \ar[r]_a & T' \ar[d]^b \\
T \ar[r]^g & S,
}
$$
et un sous-schéma fermé $Z' \subset X_{T'}$ tels que
\begin{enumerate}
\item le morphisme $b : T' \to S$ soit localement de présentation finie,
\item on ait $Z' \cap X_{a(t)} = \emptyset$, et
\item $Z \cap X_V$ soit envoyé dans $Z'$ par le morphisme $X_V \to X_{T'}$.
\end{enumerate}
Nous pouvons supposer $V$ ouvert affine de $T$ ; d'après
Limites, lemmes \ref{limits-lemma-descend-opens} et
\ref{limits-lemma-limit-affine},
nous pouvons donc trouver un $i$ et un ouvert affine $V_i \subset T_i$ tels que
$V = f_i^{-1}(V_i)$. D'après
Limites,
proposition \ref{limits-proposition-characterize-locally-finite-presentation},
nous pouvons, quitte à augmenter $i$, trouver un morphisme
$a_i : V_i \to T'$ tel que $a = a_i \circ f_i|_V$.
Le morphisme induit $X_{V_i} \to X_{T'}$ envoie $\xi_i$ dans
$Z'$. Comme $Z' \cap X_{a(t)} = \emptyset$, nous concluons que
$(T_i \to S, t'_i \leadsto t_i, \xi_i)$ est une impureté de
$\mathcal{F}$ au-dessus de $s$.
\end{proof}

\begin{lemma}
\label{flat-lemma-quasi-finite-impurity-elementary}
Plaçons-nous dans la situation \ref{flat-situation-pre-pure}.
S'il existe une impureté $(g : T \to S, t' \leadsto t, \xi)$
de $\mathcal{F}$ au-dessus de $s$ avec $g$ quasi-fini en $t$, alors il
existe une impureté $(g : T \to S, t' \leadsto t, \xi)$ telle que
$(T, t) \to (S, s)$ soit un voisinage \'etale élémentaire.
\end{lemma}

\begin{proof}
Soit $(g : T \to S, t' \leadsto t, \xi)$ une impureté de
$\mathcal{F}$ au-dessus de $s$ telle que $g$ soit quasi-fini en $t$.
Après restriction de $T$, nous pouvons supposer $g$ localement de type fini.
Appliquons
Compléments sur les morphismes,
lemme \ref{more-morphisms-lemma-etale-makes-quasi-finite-finite-at-point},
à $T \to S$ et $t \mapsto s$. Nous obtenons un diagramme
$$
\xymatrix{
T \ar[d] & T \times_S U \ar[l] \ar[d] & V \ar[l] \ar[ld] \\
S & U \ar[l]
}
$$
où $(U, u) \to (S, s)$ est un voisinage \'etale élémentaire
et $V \subset T \times_S U$ est un voisinage ouvert de $v = (t, u)$
tel que $V \to U$ soit fini et que $v$ soit l'unique point de $V$
au-dessus de $u$. Comme le morphisme $V \to T$ est \'etale,
donc plat, il existe une spécialisation $v' \leadsto v$ telle
que $v' \mapsto t'$. Notons que $\kappa(t') \subset \kappa(v')$
est une extension finie séparable. Choisissons un point quelconque $\zeta \in X_{v'}$ d'image
$\xi \in X_{t'}$. D'après
Diviseurs, lemme \ref{divisors-lemma-base-change-relative-assassin},
nous avons $\zeta \in \text{Ass}_{X_V/V}(\mathcal{F}_V)$.
De plus, l'adhérence $\overline{\{\zeta\}}$ ne rencontre pas
la fibre $X_v$, puisque, par hypothèse, l'adhérence $\overline{\{\xi\}}$
ne rencontre pas $X_t$. Autrement dit, $(V \to S, v' \leadsto v, \zeta)$
est une impureté de $\mathcal{F}$ au-dessus de $S$.

\medskip\noindent
Ensuite, soient $u' \in U'$ l'image de $v'$ et
$\theta \in X_U$ l'image de $\zeta$.
Alors $\theta \mapsto u'$ et $u' \leadsto u$.
D'après
Diviseurs, lemme \ref{divisors-lemma-base-change-relative-assassin},
nous avons $\theta \in \text{Ass}_{X_U/U}(\mathcal{F})$.
De plus, comme $\pi : X_V \to X_U$ est fini, nous avons
$\pi\big(\overline{\{\zeta\}}\big) = \overline{\{\pi(\zeta)\}}$. Puisque
$v$ est l'unique point de $V$ au-dessus de $u$, nous avons
$X_u \cap \overline{\{\pi(\zeta)\}} = \emptyset$, car
$X_v \cap \overline{\{\zeta\}} = \emptyset$. Nous concluons ainsi que
$(U \to S, u' \leadsto u, \theta)$ est une impureté de
$\mathcal{F}$ au-dessus de $s$, ce qui achève la démonstration.
\end{proof}

\begin{lemma}
\label{flat-lemma-Noetherian-impurity-quasi-finite}
Plaçons-nous dans la situation \ref{flat-situation-pre-pure}.
Supposons $S$ localement noethérien.
S'il existe une impureté de $\mathcal{F}$ au-dessus de $s$, alors
il existe une impureté $(g : T \to S, t' \leadsto t, \xi)$
de $\mathcal{F}$ au-dessus de $s$ telle que $g$ soit quasi-fini en $t$.
\end{lemma}

\begin{proof}
Nous pouvons remplacer $S$ par un voisinage affine de $s$. D'après le
lemme \ref{flat-lemma-impure-finite-presentation},
nous pouvons supposer donnée une impureté $(g : T \to S, t' \leadsto t, \xi)$
telle que $g$ soit localement de type fini et $t$ un point fermé de la
fibre de $g$ au-dessus de $s$. Nous pouvons remplacer $T$ par le sous-schéma réduit
induit sur $\overline{\{t'\}}$. Soit
$Z = \overline{\{\xi\}} \subset X_T$. Par hypothèse, $Z_t = \emptyset$
et l'image de $Z \to T$ contient $t'$. D'après
Compléments sur les morphismes,
lemme \ref{more-morphisms-lemma-relative-assassin-in-neighbourhood},
il existe un ouvert non vide $V \subset Z$ tel que, pour tout
$w \in f(V)$, tout point générique $\xi'$ de $V_w$ appartienne à
$\text{Ass}_{X_T/T}(\mathcal{F}_T)$. D'après
Compléments sur les morphismes, lemme \ref{more-morphisms-lemma-nonempty-generic-fibre},
il existe un ouvert non vide $W \subset T$ tel que $W \subset f(V)$. D'après
Compléments sur les morphismes, lemme
\ref{more-morphisms-lemma-quasi-finite-quasi-section-meeting-nearby-open-X},
il existe un sous-schéma fermé $T' \subset T$ tel que
$t \in T'$, que $T' \to S$ soit quasi-fini en $t$ et qu'il existe un point
$z \in T' \cap W$, $z \leadsto t$, qui ne s'envoie pas sur $s$.
Choisissons un point générique quelconque $\xi'$ du schéma non vide $V_z$.
Alors $(T' \to S, z \leadsto t, \xi')$ est l'impureté cherchée.
\end{proof}

\noindent
Dans la suite, nous utiliserons l'hensélisé
$S^h = \Spec(\mathcal{O}_{S, s}^h)$
de $S$ en $s$, voir
Cohomologie \'etale,
définition \ref{etale-cohomology-definition-etale-local-rings}.
Comme $S^h \to S$ envoie le point fermé de $S^h$ sur $s$ et
induit un isomorphisme des corps résiduels, nous noterons
aussi $s \in S^h$ ce point fermé. Ainsi, $(S^h, s) \to (S, s)$ est
un morphisme de schémas pointés.

\begin{lemma}
\label{flat-lemma-impurity-on-henselization}
Plaçons-nous dans la situation \ref{flat-situation-pre-pure}.
S'il existe une impureté $(S^h \to S, s' \leadsto s, \xi)$
de $\mathcal{F}$ au-dessus de $s$, alors il existe une impureté
$(T \to S, t' \leadsto t, \xi)$ de $\mathcal{F}$ au-dessus de $s$
où $(T, t) \to (S, s)$ est un voisinage \'etale élémentaire.
\end{lemma}

\begin{proof}
Nous pouvons remplacer $S$ par un voisinage affine de $s$.
Disons que $S = \Spec(A)$ et que $s$ correspond à l'idéal premier
$\mathfrak p \subset A$. Alors
$\mathcal{O}_{S, s}^h = \colim_{(T, t)} \Gamma(T, \mathcal{O}_T)$,
la limite étant prise sur la catégorie opposée à la
catégorie cofiltrante des voisinages \'etales élémentaires affines
$(T, t)$ de $(S, s)$, voir
Compléments sur les morphismes,
lemme \ref{more-morphisms-lemma-describe-henselization}
et sa démonstration. Ainsi, $S^h = \lim_i T_i$ et nous concluons par le
lemme \ref{flat-lemma-impure-limit}.
\end{proof}

\begin{lemma}
\label{flat-lemma-pure-along-X-s}
Dans la situation \ref{flat-situation-pre-pure}, les conditions suivantes
sont équivalentes :
\begin{enumerate}
\item il existe une impureté $(S^h \to S, s' \leadsto s, \xi)$
de $\mathcal{F}$ au-dessus de $s$, où $S^h$ est l'hensélisé de $S$ en $s$,
\item il existe une impureté $(T \to S, t' \leadsto t, \xi)$
de $\mathcal{F}$ au-dessus de $s$ telle que $(T, t) \to (S, s)$ soit un
voisinage \'etale élémentaire, et
\item il existe une impureté $(T \to S, t' \leadsto t, \xi)$
de $\mathcal{F}$ au-dessus de $s$ telle que $T \to S$ soit quasi-fini en $t$.
\end{enumerate}
\end{lemma}

\begin{proof}
Comme un morphisme \'etale est localement quasi-fini, il est clair que
(2) entraîne (3). Nous avons vu que (3) entraîne (2) dans le
lemme \ref{flat-lemma-quasi-finite-impurity-elementary}.
Nous avons vu que (1) entraîne (2) dans le
lemme \ref{flat-lemma-impurity-on-henselization}.
Enfin, si $(T \to S, t' \leadsto t, \xi)$ est une impureté
de $\mathcal{F}$ au-dessus de $s$ telle que $(T, t) \to (S, s)$ soit un
voisinage \'etale élémentaire, nous pouvons choisir une factorisation
$S^h \to T \to S$ du morphisme structural $S^h \to S$.
Choisissons un point quelconque $s' \in S^h$ d'image $t'$ et un point quelconque
$\xi' \in X_{s'}$ d'image $\xi \in X_{t'}$. Alors
$(S^h \to S, s' \leadsto s, \xi')$ est une impureté de
$\mathcal{F}$ au-dessus de $s$. Nous omettons les détails.
\end{proof}





\section{Modules relativement purs}
\label{flat-section-pure}

\noindent
La notion de module pur relativement à une base a été introduite dans \cite{GruRay}.

\begin{definition}
\label{flat-definition-pure}
Soit $f : X \to S$ un morphisme de schémas de type fini.
Soit $\mathcal{F}$ un $\mathcal{O}_X$-module quasi-cohérent de type fini.
\begin{enumerate}
\item Soit $s \in S$. Nous disons que $\mathcal{F}$ est {\it pur le long de $X_s$}
s'il n'existe aucune impureté $(g : T \to S, t' \leadsto t, \xi)$
de $\mathcal{F}$ au-dessus de $s$ avec $(T, t) \to (S, s)$ un
voisinage \'etale élémentaire.
\item Nous disons que $\mathcal{F}$ est {\it universellement pur le long de $X_s$}
s'il n'existe aucune impureté de $\mathcal{F}$ au-dessus de $s$.
\item Nous disons que $X$ est {\it pur le long de $X_s$} si $\mathcal{O}_X$
est pur le long de $X_s$.
\item Nous disons que $\mathcal{F}$ est {\it universellement $S$-pur}, ou
{\it universellement pur relativement à $S$}, si $\mathcal{F}$ est universellement
pur le long de $X_s$ pour tout $s \in S$.
\item Nous disons que $\mathcal{F}$ est {\it $S$-pur}, ou
{\it pur relativement à $S$}, si $\mathcal{F}$ est pur le long de $X_s$
pour tout $s \in S$.
\item Nous disons que $X$ est {\it $S$-pur} ou {\it pur relativement à $S$}
si $\mathcal{O}_X$ est pur relativement à $S$.
\end{enumerate}
\end{definition}

\noindent
Nous nous restreignons ici délibérément aux morphismes
de type fini, et non aux morphismes seulement localement de
type fini ; voir la
remarque \ref{flat-remark-discuss-finite-type}
pour une discussion. Dans la situation de la définition, le
lemme \ref{flat-lemma-pure-along-X-s}
nous dit que les conditions suivantes sont équivalentes :
\begin{enumerate}
\item $\mathcal{F}$ est pur le long de $X_s$,
\item il n'existe aucune impureté $(g : T \to S, t' \leadsto t, \xi)$ avec $g$
quasi-fini en $t$,
\item il n'existe aucune impureté de la forme
$(S^h \to S, s' \leadsto s, \xi)$, où $S^h$ est l'hensélisé
de $S$ en $s$.
\end{enumerate}
Si nous notons $X^h = X \times_S S^h$ et $\mathcal{F}^h$ l'image réciproque
de $\mathcal{F}$ sur $X^h$, nous pouvons formuler cette dernière condition
sous la forme positive suivante :
\begin{enumerate}
\item[(4)] Tous les points de $\text{Ass}_{X^h/S^h}(\mathcal{F}^h)$ se spécialisent
en des points de $X_s$.
\end{enumerate}
En particulier, il est clair que $\mathcal{F}$ est pur le long de $X_s$
si et seulement si l'image réciproque de $\mathcal{F}$ sur
$X \times_S \Spec(\mathcal{O}_{S, s})$ est pure le long de $X_s$.

\begin{remark}
\label{flat-remark-discuss-finite-type}
Soient $f : X \to S$ un morphisme localement de type fini
et $\mathcal{F}$ un $\mathcal{O}_X$-module quasi-cohérent de type fini.
Dans ce cas, il est encore vrai que (1) et (2) ci-dessus sont équivalentes,
car la démonstration du
lemme \ref{flat-lemma-quasi-finite-impurity-elementary}
n'utilise pas la quasi-compacité de $f$. Il est également clair que
(3) et (4) sont équivalentes. Cependant, nous ignorons si (1) et (3) sont
équivalentes. Dans ce cas, il peut parfois être plus commode de définir
la pureté par les conditions équivalentes (3) et (4), comme dans \cite{GruRay}.
D'autre part, pour de nombreuses applications, il semble que la notion appropriée
soit bien celle de pureté universelle.
\end{remark}

\noindent
Il est naturel de se demander si la propriété de pureté relativement à
la base est conservée par changement de base, autrement dit si être pur revient
à être universellement pur. Il se trouve que c'est vrai
sur une base noethérienne (voir le
lemme \ref{flat-lemma-Noetherian-base-change}),
ou si le faisceau est plat (voir les
lemmes \ref{flat-lemma-finite-type-flat-pure-along-fibre-is-universal} et
\ref{flat-lemma-finite-type-flat-pure-is-universal}).
Ce n'est pas vrai en général, même si le morphisme et le faisceau sont de
présentation finie ; voir
Exemples, section \ref{examples-section-pure-not-universally},
pour un contre-exemple. Vérifions d'abord que notre emploi de « universellement »
correspond à la notion usuelle.

\begin{lemma}
\label{flat-lemma-base-change-universally}
Soit $f : X \to S$ un morphisme de schémas de type fini.
Soit $\mathcal{F}$ un $\mathcal{O}_X$-module quasi-cohérent de type fini.
Soit $s \in S$. Les conditions suivantes sont équivalentes :
\begin{enumerate}
\item $\mathcal{F}$ est universellement pur le long de $X_s$, et
\item pour tout morphisme de schémas pointés $(S', s') \to (S, s)$,
l'image réciproque $\mathcal{F}_{S'}$ est pure le long de $X_{s'}$.
\end{enumerate}
En particulier, $\mathcal{F}$ est universellement pur relativement à $S$ si et
seulement si tout changement de base $\mathcal{F}_{S'}$ de $\mathcal{F}$ est
pur relativement à $S'$.
\end{lemma}

\begin{proof}
C'est formel.
\end{proof}

\begin{lemma}
\label{flat-lemma-quasi-finite-base-change}
Soit $f : X \to S$ un morphisme de schémas de type fini.
Soit $\mathcal{F}$ un $\mathcal{O}_X$-module quasi-cohérent de type fini.
Soit $s \in S$. Soit $(S', s') \to (S, s)$ un morphisme de schémas pointés.
Si $S' \to S$ est quasi-fini en $s'$ et $\mathcal{F}$ est pur le long de $X_s$,
alors $\mathcal{F}_{S'}$ est pur le long de $X_{s'}$.
\end{lemma}

\begin{proof}
Si $(T \to S', t' \leadsto t, \xi)$ est une impureté de
$\mathcal{F}_{S'}$ au-dessus de $s'$ avec $T \to S'$ quasi-fini en $t$,
alors $(T \to S, t' \to t, \xi)$ est une impureté de $\mathcal{F}$
au-dessus de $s$ avec $T \to S$ quasi-fini en $t$, voir
Morphismes, lemme \ref{morphisms-lemma-composition-quasi-finite}.
Le lemme résulte donc immédiatement de la caractérisation (2)
de la pureté donnée après la
définition \ref{flat-definition-pure}.
\end{proof}

\begin{lemma}
\label{flat-lemma-Noetherian-base-change}
Soit $f : X \to S$ un morphisme de schémas de type fini.
Soit $\mathcal{F}$ un $\mathcal{O}_X$-module quasi-cohérent de type fini.
Soit $s \in S$. Si $\mathcal{O}_{S, s}$ est noethérien, alors
$\mathcal{F}$ est pur le long de $X_s$ si et seulement si $\mathcal{F}$
est universellement pur le long de $X_s$.
\end{lemma}

\begin{proof}
Nous pouvons d'abord remplacer $S$ par $\Spec(\mathcal{O}_{S, s})$, c'est-à-dire
supposer $S$ noethérien. Utilisons ensuite le
lemme \ref{flat-lemma-Noetherian-impurity-quasi-finite}
et la caractérisation (2) de la pureté donnée dans la discussion suivant la
définition \ref{flat-definition-pure}
pour conclure.
\end{proof}

\noindent
La pureté satisfait à la descente plate.

\begin{lemma}
\label{flat-lemma-flat-descend-pure}
Soit $f : X \to S$ un morphisme de schémas de type fini.
Soit $\mathcal{F}$ un $\mathcal{O}_X$-module quasi-cohérent de type fini.
Soit $s \in S$. Soit $(S', s') \to (S, s)$ un morphisme de schémas pointés.
Supposons $S' \to S$ plat en $s'$.
\begin{enumerate}
\item Si $\mathcal{F}_{S'}$ est pur le long de $X_{s'}$,
alors $\mathcal{F}$ est pur le long de $X_s$.
\item Si $\mathcal{F}_{S'}$ est universellement pur le long de $X_{s'}$,
alors $\mathcal{F}$ est universellement pur le long de $X_s$.
\end{enumerate}
\end{lemma}

\begin{proof}
Soit $(T \to S, t' \leadsto t, \xi)$ une impureté de
$\mathcal{F}$ au-dessus de $s$. Posons $T_1 = T \times_S S'$ et soit $t_1$
l'unique point de $T_1$ d'images $t$ et $s'$. Comme
$T_1 \to T$ est plat en $t_1$, voir
Morphismes, lemme \ref{morphisms-lemma-base-change-flat},
il existe une spécialisation $t'_1 \leadsto t_1$ au-dessus de
$t' \leadsto t$, voir
Algèbre, section \ref{algebra-section-going-up}.
Choisissons un point $\xi_1 \in X_{t'_1}$ correspondant à un point
générique de $\Spec(\kappa(t'_1) \otimes_{\kappa(t')} \kappa(\xi))$, voir
Schémas, lemme \ref{schemes-lemma-points-fibre-product}.
D'après
Diviseurs, lemme \ref{divisors-lemma-base-change-relative-assassin},
nous avons $\xi_1 \in \text{Ass}_{X_{T_1}/T_1}(\mathcal{F}_{T_1})$.
Comme l'adhérence de Zariski de $\{\xi_1\}$ dans $X_{T_1}$ s'envoie dans
l'adhérence de Zariski de $\{\xi\}$ dans $X_T$, nous concluons que
cette adhérence est disjointe de $X_{t_1}$. Ainsi,
$(T_1 \to S', t'_1 \leadsto t_1, \xi_1)$
est une impureté de $\mathcal{F}_{S'}$ au-dessus de $s'$.
Autrement dit, nous avons démontré la contraposée de la partie (2) du
lemme. Enfin, si $(T, t) \to (S, s)$ est un voisinage
\'etale élémentaire, alors $(T_1, t_1) \to (S', s')$ est aussi un
voisinage \'etale élémentaire ; nous voyons ainsi que (1) est vraie.
\end{proof}

\begin{lemma}
\label{flat-lemma-supported-on-closed}
Soit $i : Z \to X$ une immersion fermée de schémas de type fini sur
un schéma $S$. Soit $s \in S$. Soit $\mathcal{F}$ un
faisceau quasi-cohérent de type fini sur $Z$. Alors $\mathcal{F}$ est
(universellement) pur le long de $Z_s$ si et seulement si $i_*\mathcal{F}$
est (universellement) pur le long de $X_s$.
\end{lemma}

\begin{proof}
Cela résulte de
Diviseurs, lemme \ref{divisors-lemma-relative-weak-assassin-finite}.
\end{proof}








\section{Exemples de faisceaux relativement purs}
\label{flat-section-examples-pure-sheaves}

\noindent
Voici quelques exemples où l'on peut expliciter la signification de la pureté.

\begin{lemma}
\label{flat-lemma-proper-pure}
Soit $f : X \to S$ un morphisme de schémas de type fini.
Soit $\mathcal{F}$ un $\mathcal{O}_X$-module quasi-cohérent de type fini.
\begin{enumerate}
\item Si le support de $\mathcal{F}$ est propre sur $S$, alors
$\mathcal{F}$ est universellement pur relativement à $S$.
\item Si $f$ est propre, alors
$\mathcal{F}$ est universellement pur relativement à $S$.
\item Si $f$ est propre, alors $X$ est universellement pur relativement à $S$.
\end{enumerate}
\end{lemma}

\begin{proof}
Ramenons d'abord (1) à (2). Soit $Z \subset X$ le
support schématique de $\mathcal{F}$. Soit $i : Z \to X$
l'immersion fermée correspondante et écrivons
$\mathcal{F} = i_*\mathcal{G}$ pour un
$\mathcal{O}_Z$-module quasi-cohérent de type fini $\mathcal{G}$, voir
Morphismes, section \ref{morphisms-section-support}.
Dans le cas (1), $Z \to S$ est propre par hypothèse.
Ainsi, d'après le lemme \ref{flat-lemma-supported-on-closed}, le cas (1) se ramène au cas (2).

\medskip\noindent
Supposons $f$ propre.
Soit $(g : T \to S, t' \leadsto t, \xi)$ une impureté de $\mathcal{F}$
au-dessus de $s \in S$. Comme $f$ est propre, il est universellement fermé. Ainsi,
$f_T : X_T \to T$ est fermé. Comme $f_T(\xi) = t'$, cela entraîne que
$t \in f(\overline{\{\xi\}})$, d'où une contradiction.
\end{proof}

\begin{lemma}
\label{flat-lemma-quasi-finite-pure}
Soit $f : X \to S$ un morphisme de schémas séparé et de type fini.
Soit $\mathcal{F}$ un $\mathcal{O}_X$-module quasi-cohérent de type fini.
Supposons $\text{Supp}(\mathcal{F}_s)$ fini pour tout $s \in S$.
Alors les conditions suivantes sont équivalentes :
\begin{enumerate}
\item $\mathcal{F}$ est pur relativement à $S$,
\item le support schématique de $\mathcal{F}$ est fini sur $S$, et
\item $\mathcal{F}$ est universellement pur relativement à $S$.
\end{enumerate}
En particulier, pour un morphisme quasi-fini séparé $X \to S$,
$X$ est pur relativement à $S$ si et seulement si $X \to S$ est fini.
\end{lemma}

\begin{proof}
Soit $Z \subset X$ le support schématique de $\mathcal{F}$, voir
Morphismes, définition \ref{morphisms-definition-scheme-theoretic-support}.
Alors $Z \to S$ est un morphisme de schémas séparé de type fini à
fibres finies. Il est donc séparé et quasi-fini, voir
Morphismes, lemme \ref{morphisms-lemma-quasi-finite}.
D'après le
lemme \ref{flat-lemma-supported-on-closed},
il suffit de démontrer le lemme pour $Z \to S$ et le faisceau $\mathcal{F}$
considéré comme module quasi-cohérent de type fini sur $Z$. Nous pouvons donc
supposer que $X \to S$ est séparé et quasi-fini et que
$\text{Supp}(\mathcal{F}) = X$.

\medskip\noindent
Il résulte du
lemme \ref{flat-lemma-proper-pure}
et de
Morphismes, lemme \ref{morphisms-lemma-finite-proper},
que (2) entraîne (3). Il est trivial que (3) entraîne (1). Supposons (1) vérifiée.
Montrons (2). Il est clair que nous pouvons supposer $S$ affine. D'après
Compléments sur les morphismes,
lemme \ref{more-morphisms-lemma-quasi-finite-separated-pass-through-finite},
nous pouvons trouver un diagramme
$$
\xymatrix{
X \ar[rd]_f \ar[rr]_j & & T \ar[ld]^\pi \\
& S &
}
$$
où $\pi$ est fini et $j$ est une immersion ouverte quasi-compacte.
Si nous montrons que $j$ est fermé, alors $j$ est une immersion fermée,
et nous concluons que $f = \pi \circ j$ est fini.
Pour montrer que $j$ est fermé, il suffit de montrer que les spécialisations
se relèvent le long de $j$, voir
Schémas, lemme \ref{schemes-lemma-quasi-compact-closed}.
Soit $x \in X$, posons $t' = j(x)$ et soit $t' \leadsto t$ une spécialisation.
Il faut montrer $t \in j(X)$. Posons $s' = f(x)$ et $s = \pi(t)$, de sorte que
$s' \leadsto s$. D'après
Compléments sur les morphismes, lemme
\ref{more-morphisms-lemma-etale-splits-off-quasi-finite-part-technical},
nous pouvons trouver un voisinage \'etale élémentaire
$(U, u) \to (S, s)$ et une décomposition
$$
T_U = T \times_S U = V \amalg W
$$
en sous-schémas ouverts et fermés, telle que $V \to U$ soit fini,
qu'il existe un unique point $v$ de $V$ d'image $u$, et que
$v$ ait pour image $t$ dans $T$. Comme $V \to T$ est \'etale, nous pouvons relever
les généralisations, voir
Morphismes, lemmes \ref{morphisms-lemma-generalizations-lift-flat} et
\ref{morphisms-lemma-etale-flat}.
Il existe donc une spécialisation $v' \leadsto v$ telle que $v'$
s'envoie sur $t' \in T$. En particulier, $v' \in X_U \subset T_U$.
Notons $u' \in U$ l'image de $t'$. Notons que
$v' \in \text{Ass}_{X_U/U}(\mathcal{F})$, car $X_{u'}$ est un ensemble
fini discret et $X_{u'} = \text{Supp}(\mathcal{F}_{u'})$.
Comme $\mathcal{F}$ est pur relativement à $S$, le point $v'$ doit
se spécialiser en un point de $X_u$. Puisque $v$ est l'unique point de
$V$ au-dessus de $u$ (et qu'aucun point de $W$ ne peut être une spécialisation
de $v'$), nous voyons que $v \in X_u$. Ainsi, $t \in X$.
\end{proof}

\begin{lemma}
\label{flat-lemma-flat-geometrically-integral-fibres-pure}
Soit $f : X \to S$ un morphisme de schémas plat de type fini
à fibres géométriquement intègres. Alors $X$ est universellement pur
sur $S$.
\end{lemma}

\begin{proof}
Soient $\xi \in X$ avec $s' = f(\xi)$ et $s' \leadsto s$ une spécialisation
de $S$. Si $\xi$ est un point associé de $X_{s'}$, alors $\xi$ est
l'unique point générique, car $X_{s'}$ est un schéma intègre. Soit
$\xi_0$ l'unique point générique de $X_s$. Comme $X \to S$ est plat,
nous pouvons relever $s' \leadsto s$ en une spécialisation
$\xi' \leadsto \xi_0$ dans $X$, voir
Morphismes, lemme \ref{morphisms-lemma-generalizations-lift-flat}.
Alors $\xi \leadsto \xi'$, car $\xi$ est le point générique de $X_{s'}$,
donc $\xi \leadsto \xi_0$. Cela signifie que $(\text{id}_S, s' \to s, \xi)$
n'est pas une impureté de $\mathcal{O}_X$ au-dessus de $s$. Comme l'hypothèse
que $f$ est de type fini, plat et à fibres géométriquement intègres
est conservée par changement de base, il n'existe aucune
impureté après un changement de base quelconque. Nous voyons ainsi que $X$ est
universellement $S$-pur.
\end{proof}

\begin{lemma}
\label{flat-lemma-affine-locally-projective-pure}
Soit $f : X \to S$ un morphisme de schémas affine et de type fini.
Soit $\mathcal{F}$ un $\mathcal{O}_X$-module quasi-cohérent de type fini
tel que $f_*\mathcal{F}$ soit localement projectif sur $S$, voir
Propriétés, définition \ref{properties-definition-locally-projective}.
Alors $\mathcal{F}$ est universellement pur sur $S$.
\end{lemma}

\begin{proof}
Après réduction au cas où $S$ est le spectre d'un anneau local
hensélien, cela résulte du
lemme \ref{flat-lemma-explain-why-pure}.
\end{proof}




\section{Un critère de pureté}
\label{flat-section-criterion-purity}

\noindent
Nous démontrons d'abord que, pour une famille plate de faisceaux
quasi-cohérents de type fini, les points de l'assassin relatif
se spécialisent en des points des assassins relatifs des fibres voisines
(lorsqu'ils admettent une spécialisation dans ces fibres).

\begin{lemma}
\label{flat-lemma-associated-point-specializes}
Soit $f : X \to S$ un morphisme de schémas de type fini.
Soit $\mathcal{F}$ un $\mathcal{O}_X$-module quasi-cohérent de type fini.
Soit $s \in S$.
Supposons $\mathcal{F}$ plat sur $S$ en tout point de $X_s$.
Soit $x' \in \text{Ass}_{X/S}(\mathcal{F})$ avec $f(x') = s'$
tel que $s' \leadsto s$ soit une spécialisation dans $S$. Si
$x'$ se spécialise en un point de $X_s$, alors $x' \leadsto x$
pour un $x \in \text{Ass}_{X_s}(\mathcal{F}_s)$.
\end{lemma}

\begin{proof}
Disons que $x' \leadsto t$ avec $t \in X_s$. Nous pouvons alors trouver des spécialisations
$x' \leadsto x \leadsto t$, où $x$ correspond au point générique
d'une composante irréductible de
$\overline{\{x'\}} \cap f^{-1}(\{s\})$. Par hypothèse, $\mathcal{F}$ est plat
sur $S$ en $x$. D'après Compléments sur les morphismes, lemme
\ref{more-morphisms-lemma-associated-point-specializes},
nous avons $x \in \text{Ass}_{X/S}(\mathcal{F})$, comme voulu.
\end{proof}

\begin{lemma}
\label{flat-lemma-criterion}
Soit $f : X \to S$ un morphisme de schémas de type fini.
Soit $\mathcal{F}$ un $\mathcal{O}_X$-module quasi-cohérent
de type fini. Soit $s \in S$. Soit $(S', s') \to (S, s)$ un
voisinage \'etale élémentaire et soit
$$
\xymatrix{
X \ar[d] & X' \ar[l]^g \ar[d] \\
S & S' \ar[l]
}
$$
un diagramme commutatif de morphismes de schémas. Supposons que
\begin{enumerate}
\item $\mathcal{F}$ soit plat sur $S$ en tout point de $X_s$,
\item $X' \to S'$ soit de type fini,
\item $g^*\mathcal{F}$ soit pur le long de $X'_{s'}$,
\item $g : X' \to X$ soit \'etale, et
\item $g(X')$ contienne $\text{Ass}_{X_s}(\mathcal{F}_s)$.
\end{enumerate}
Dans cette situation, $\mathcal{F}$ est pur le long de $X_s$ si et seulement
si l'image de $X' \to X \times_S S'$ contient les points de
$\text{Ass}_{X \times_S S'/S'}(\mathcal{F} \times_S S')$
au-dessus des points de $S'$ qui se spécialisent en $s'$.
\end{lemma}

\begin{proof}
Comme le morphisme $S' \to S$ est \'etale, si $\mathcal{F}$
est pur le long de $X_s$, alors $\mathcal{F} \times_S S'$ est pur le long de
$X_s$, voir le
lemme \ref{flat-lemma-quasi-finite-base-change}.
Comme la pureté satisfait à la descente plate, voir le
lemme \ref{flat-lemma-flat-descend-pure},
si $\mathcal{F} \times_S S'$ est pur le long de $X_{s'}$, alors
$\mathcal{F}$ est pur le long de $X_s$. Nous pouvons donc remplacer $S$ par $S'$
et supposer $S = S'$, de sorte que $g : X' \to X$ est un morphisme \'etale
entre schémas de type fini sur $S$. De plus, nous pouvons remplacer
$S$ par $\Spec(\mathcal{O}_{S, s})$ et supposer $S$ local.

\medskip\noindent
Supposons d'abord $\mathcal{F}$ pur le long de $X_s$.
Dans ce cas, tout point de $\text{Ass}_{X/S}(\mathcal{F})$
se spécialise en un point de $X_s$, par pureté. Ainsi, d'après le
lemme \ref{flat-lemma-associated-point-specializes},
tout point de $\text{Ass}_{X/S}(\mathcal{F})$
se spécialise en un point de $\text{Ass}_{X_s}(\mathcal{F}_s)$.
Tout point de $\text{Ass}_{X/S}(\mathcal{F})$ appartient donc à
l'image de $g$ (car cette image est ouverte et contient
$\text{Ass}_{X_s}(\mathcal{F}_s)$).

\medskip\noindent
Réciproquement, supposons que $g(X')$ contienne $\text{Ass}_{X/S}(\mathcal{F})$.
Soit $S^h = \Spec(\mathcal{O}_{S, s}^h)$ l'hensélisé
de $S$ en $s$. Notons $g^h : (X')^h \to X^h$ le changement de base de $g$
par $S^h \to S$, et $\mathcal{F}^h$ l'image réciproque de $\mathcal{F}$
sur $X^h$. D'après
Diviseurs, lemme \ref{divisors-lemma-base-change-relative-assassin} et
remarque \ref{divisors-remark-base-change-relative-assassin},
l'assassin relatif $\text{Ass}_{X^h/S^h}(\mathcal{F}^h)$
est l'image réciproque de $\text{Ass}_{X/S}(\mathcal{F})$ par la projection
$X^h \to X$. Comme nous avons supposé que $g(X')$ contient
$\text{Ass}_{X/S}(\mathcal{F})$, nous concluons que le changement de base
$g^h((X')^h) = g(X') \times_S S^h$ contient
$\text{Ass}_{X^h/S^h}(\mathcal{F}^h)$. Nous nous ramenons ainsi
au cas où $S$ est le spectre d'un anneau local hensélien.
Soit $x \in \text{Ass}_{X/S}(\mathcal{F})$. Pour achever la démonstration du
lemme, il faut montrer que $x$ se spécialise en un point de $X_s$, voir
le critère (4) de pureté dans la discussion suivant la
définition \ref{flat-definition-pure}.
Par hypothèse, il existe un $x' \in X'$ tel que $g(x') = x$.
Comme $g : X' \to X$ est \'etale, nous avons
$x' \in \text{Ass}_{X'/S}(g^*\mathcal{F})$, voir le
lemme \ref{flat-lemma-etale-weak-assassin-up-down} (appliqué au
morphisme de fibres $X'_w \to X_w$, où $w \in S$ est l'image de $x'$).
Comme $g^*\mathcal{F}$ est pur le long de $X'_s$, nous avons $x' \leadsto y$
pour un $y \in X'_s$. Ainsi, $x = g(x') \leadsto g(y)$ et
$g(y) \in X_s$, comme voulu.
\end{proof}

\begin{lemma}
\label{flat-lemma-finite-type-flat-pure-along-fibre-is-universal}
Soit $f : X \to S$ un morphisme de schémas.
Soit $\mathcal{F}$ un $\mathcal{O}_X$-module quasi-cohérent.
Soit $s \in S$.
Supposons que
\begin{enumerate}
\item $f$ soit de type fini,
\item $\mathcal{F}$ soit de type fini,
\item $\mathcal{F}$ soit plat sur $S$ en tout point de $X_s$, et
\item $\mathcal{F}$ soit pur le long de $X_s$.
\end{enumerate}
Alors $\mathcal{F}$ est universellement pur le long de $X_s$.
\end{lemma}

\begin{proof}
Faisons d'abord une remarque préliminaire. Supposons que $(S', s') \to (S, s)$
soit un voisinage \'etale élémentaire. Notons $\mathcal{F}'$
l'image réciproque de $\mathcal{F}$ sur $X' = X \times_S S'$.
D'après la discussion suivant la
définition \ref{flat-definition-pure},
$\mathcal{F}'$ est pur le long de $X'_{s'}$. De plus, $\mathcal{F}'$
est plat sur $S'$ le long de $X'_{s'}$. Il suffit alors de démontrer
que $\mathcal{F}'$ est universellement pur le long de $X'_{s'}$. En effet, pour
tout morphisme $(T, t) \to (S, s)$ de schémas pointés,
le produit fibré $(T', t') = (T \times_S S', (t, s'))$ est plat sur $(T, t)$ ;
ainsi, si $\mathcal{F}_{T'}$ est pur le long de $X_{t'}$, alors
$\mathcal{F}_T$ est pur le long de $X_t$ d'après le
lemme \ref{flat-lemma-flat-descend-pure}.
Ainsi, au cours de la démonstration, nous pouvons toujours remplacer $(s, S)$ par un voisinage
\'etale élémentaire.
Nous pouvons aussi remplacer $S$ par $\Spec(\mathcal{O}_{S, s})$
en raison du caractère local du problème.

\medskip\noindent
Choisissons un voisinage \'etale élémentaire $(S', s') \to (S, s)$ et
un diagramme commutatif
$$
\xymatrix{
X \ar[d] & X' \ar[l]^g \ar[d] \\
S & \Spec(\mathcal{O}_{S', s'}) \ar[l]
}
$$
tels que $X' \to X \times_S \Spec(\mathcal{O}_{S', s'})$
soit \'etale, que $X_s = g((X')_{s'})$, que le schéma $X'$ soit affine
et que $\Gamma(X', g^*\mathcal{F})$ soit un
$\mathcal{O}_{S', s'}$-module libre, voir le
lemme \ref{flat-lemma-finite-type-flat-along-fibre-free-variant}.
Notons que $X' \to \Spec(\mathcal{O}_{S', s'})$ est de type fini
(c'est un morphisme quasi-compact composé d'un morphisme \'etale
et du changement de base d'un morphisme de type fini).
D'après les remarques préliminaires du premier paragraphe de la démonstration,
nous pouvons remplacer $S$ par $\Spec(\mathcal{O}_{S', s'})$. Nous pouvons donc
supposer qu'il existe un diagramme commutatif
$$
\xymatrix{
X \ar[dr] & & X' \ar[ll]^g \ar[ld] \\
& S &
}
$$
de schémas de type fini sur $S$, où $g$ est \'etale, $X_s \subset g(X')$,
où $S$ est local de point fermé $s$, où $X'$ est affine et où
$\Gamma(X', g^*\mathcal{F})$ est un $\Gamma(S, \mathcal{O}_S)$-module libre.
Notons que, dans ce cas, $g^*\mathcal{F}$ est universellement pur sur $S$, voir le
lemme \ref{flat-lemma-affine-locally-projective-pure}.

\medskip\noindent
Dans cette situation, nous appliquons le
lemme \ref{flat-lemma-criterion}
pour déduire que $\text{Ass}_{X/S}(\mathcal{F}) \subset g(X')$
de l'hypothèse que $\mathcal{F}$ est pur le long de $X_s$
et plat sur $S$ le long de $X_s$. D'après
Diviseurs, lemme \ref{divisors-lemma-base-change-relative-assassin} et
remarque \ref{divisors-remark-base-change-relative-assassin},
pour tout morphisme de schémas pointés
$(T, t) \to (S, s)$, nous avons
$$
\text{Ass}_{X_T/T}(\mathcal{F}_T) \subset
(X_T \to X)^{-1}(\text{Ass}_{X/S}(\mathcal{F})) \subset
g(X') \times_S T = g_T(X'_T).
$$
Ainsi, par le
lemme \ref{flat-lemma-criterion}
appliqué au changement de base à $(T, t)$ du diagramme ci-dessus,
nous concluons que $\mathcal{F}_T$ est pur le long de $X_t$, comme voulu.
\end{proof}

\begin{lemma}
\label{flat-lemma-finite-type-flat-pure-is-universal}
Soit $f : X \to S$ un morphisme de schémas de type fini.
Soit $\mathcal{F}$ un $\mathcal{O}_X$-module quasi-cohérent de type fini.
Supposons $\mathcal{F}$ plat sur $S$. Dans ce cas,
$\mathcal{F}$ est pur relativement à $S$ si et seulement si $\mathcal{F}$
est universellement pur relativement à $S$.
\end{lemma}

\begin{proof}
Conséquence immédiate du
lemme \ref{flat-lemma-finite-type-flat-pure-along-fibre-is-universal}
et des définitions.
\end{proof}

\begin{lemma}
\label{flat-lemma-limit-purity}
Soit $I$ un ensemble ordonné filtrant.
Soit $(S_i, g_{ii'})$ un système projectif de schémas affines indexé par $I$.
Posons $S = \lim_i S_i$ et soit $s \in S$.
Notons $g_i : S \to S_i$ les projections et posons $s_i = g_i(s)$.
Supposons que $f : X \to S$ soit un morphisme de présentation finie,
et $\mathcal{F}$ un $\mathcal{O}_X$-module quasi-cohérent de présentation finie,
pur le long de $X_s$ et plat sur $S$ en tout point de $X_s$.
Alors il existe un $i \in I$, un morphisme de présentation finie
$X_i \to S_i$, un $\mathcal{O}_{X_i}$-module quasi-cohérent $\mathcal{F}_i$
de présentation finie, pur le long de $(X_i)_{s_i}$ et plat sur $S_i$
en tout point de $(X_i)_{s_i}$, tels que $X \cong X_i \times_{S_i} S$
et que l'image réciproque de $\mathcal{F}_i$ sur $X$ soit isomorphe
à $\mathcal{F}$.
\end{lemma}

\begin{proof}
Soit $U \subset X$ l'ensemble des points où $\mathcal{F}$ est
plat sur $S$. D'après
Compléments sur les morphismes, théorème \ref{more-morphisms-theorem-openness-flatness},
c'est un sous-schéma ouvert de $X$. Par hypothèse, $X_s \subset U$.
Comme $X_s$ est quasi-compact, nous pouvons trouver un ouvert quasi-compact
$U' \subset U$ tel que $X_s \subset U'$. D'après
Limites, lemme \ref{limits-lemma-descend-finite-presentation},
nous pouvons trouver un $i \in I$ et un morphisme de présentation finie
$f_i : X_i \to S_i$ dont le changement de base à $S$ est isomorphe à $f_i$.
Fixons un tel choix et posons $X_{i'} = X_i \times_{S_i} S_{i'}$.
Alors $X = \lim_{i'} X_{i'}$, avec des morphismes de transition affines. D'après
Limites, lemme \ref{limits-lemma-descend-modules-finite-presentation},
nous pouvons, quitte à augmenter $i$, supposer qu'il existe un
$\mathcal{O}_{X_i}$-module quasi-cohérent $\mathcal{F}_i$
de présentation finie dont le changement de base à $S$ est isomorphe à
$\mathcal{F}$. D'après
Limites, lemme \ref{limits-lemma-descend-opens},
nous pouvons, quitte à augmenter $i$, supposer qu'il existe un
ouvert $U'_i \subset X_i$ dont l'image réciproque dans $X$ est $U'$.
Notons qu'en particulier $(X_i)_{s_i} \subset U'_i$. D'après
Limites, lemme \ref{limits-lemma-descend-module-flat-finite-presentation},
(en augmentant encore $i$),
nous pouvons supposer $\mathcal{F}_i$ plat sur $U'_i$.
En particulier, $\mathcal{F}_i$ est plat le long de $(X_i)_{s_i}$.

\medskip\noindent
Utilisons ensuite le
lemme \ref{flat-lemma-finite-presentation-flat-along-fibre}
pour choisir un voisinage \'etale élémentaire
$(S_i', s_i') \to (S_i, s_i)$ et un diagramme commutatif de schémas
$$
\xymatrix{
X_i \ar[d] & X_i' \ar[l]^{g_i} \ar[d] \\
S_i & S_i' \ar[l]
}
$$
tels que $g_i$ soit \'etale, que $(X_i)_{s_i} \subset g_i(X_i')$, que les schémas
$X_i'$, $S_i'$ soient affines et que
$\Gamma(X_i', g_i^*\mathcal{F}_i)$ soit un
$\Gamma(S_i', \mathcal{O}_{S_i'})$-module projectif.
Notons que $g_i^*\mathcal{F}_i$ est universellement pur sur $S'_i$, voir le
lemme \ref{flat-lemma-affine-locally-projective-pure}.
Par changement de base, le diagramme ci-dessus donne un diagramme comportant les morphismes
$(S'_{i'}, s'_{i'}) \to (S_{i'}, s_{i'})$ et
$g_{i'} : X'_{i'} \to X_{i'}$ sur $S_{i'}$
pour tout $i' \geq i$, ainsi qu'un diagramme
comportant les morphismes $(S', s') \to (S, s)$ et $g : X' \to X$ sur $S$.

\medskip\noindent
Nous pouvons maintenant utiliser notre critère de pureté.
Soit $W'_i \subset X_i \times_{S_i} S'_i$ l'image du morphisme
\'etale $X'_i \to X_i \times_{S_i} S'_i$. Pour tout $i' \geq i$,
nous avons de même l'image $W'_{i'} \subset X_{i'} \times_{S_{i'}} S'_{i'}$,
ainsi que l'image $W' \subset X \times_S S'$. La formation des images commute
au changement de base ; ainsi, $W'_{i'} = W'_i \times_{S'_i} S'_{i'}$ et
$W' = W_i \times_{S'_i} S'$. Comme
$\mathcal{F}$ est pur le long de $X_s$, le
lemme \ref{flat-lemma-criterion}
entraîne que
\begin{equation}
\label{flat-equation-inclusion}
f^{-1}(\Spec(\mathcal{O}_{S', s'})) \cap
\text{Ass}_{X \times_S S'/S'}(\mathcal{F} \times_S S') \subset W'
\end{equation}
D'après
Compléments sur les morphismes,
lemme \ref{more-morphisms-lemma-relative-assassin-constructible},
les ensembles
$$
E = \{t \in S' \mid \text{Ass}_{X_t}(\mathcal{F}_t) \subset W' \}
\quad\text{et}\quad
E_{i'} = \{t \in S'_{i'} \mid
\text{Ass}_{X_t}(\mathcal{F}_{i', t}) \subset W'_{i'} \}
$$
sont des parties localement constructibles de $S'$ et $S'_{i'}$. D'après
Compléments sur les morphismes,
lemme \ref{more-morphisms-lemma-base-change-assassin-in-U},
$E_{i'}$ est l'image réciproque de $E_i$ par le morphisme
$S'_{i'} \to S'_i$, et $E$ est l'image réciproque de $E_i$ par
le morphisme $S' \to S'_i$. Ainsi,
l'équation (\ref{flat-equation-inclusion})
équivaut à dire que
$\Spec(\mathcal{O}_{S', s'})$ s'envoie dans $E_i$. Comme
$\mathcal{O}_{S', s'} =
\colim_{i' \geq i} \mathcal{O}_{S'_{i'}, s'_{i'}}$,
nous voyons que $\Spec(\mathcal{O}_{S'_{i'}, s'_{i'}})$
s'envoie dans $E_i$ pour un certain $i' \geq i$, voir
Limites, lemme \ref{limits-lemma-limit-contained-in-constructible}.
En appliquant alors le
lemme \ref{flat-lemma-criterion}
à la situation sur $S_{i'}$,
nous concluons que $\mathcal{F}_{i'}$ est pur le long de $(X_{i'})_{s_{i'}}$.
\end{proof}

\begin{lemma}
\label{flat-lemma-flat-finite-presentation-purity-open}
Soit $f : X \to S$ un morphisme de présentation finie.
Soit $\mathcal{F}$ un $\mathcal{O}_X$-module quasi-cohérent
de présentation finie, plat sur $S$. Alors l'ensemble
$$
U = \{s \in S \mid \mathcal{F}\text{ est pur le long de }X_s\}
$$
est ouvert dans $S$.
\end{lemma}

\begin{proof}
Soit $s \in U$. À l'aide du
lemme \ref{flat-lemma-finite-presentation-flat-along-fibre},
nous pouvons trouver un voisinage \'etale élémentaire
$(S', s') \to (S, s)$ et un diagramme commutatif
$$
\xymatrix{
X \ar[d] & X' \ar[l]^g \ar[d] \\
S & S' \ar[l]
}
$$
tels que $g$ soit \'etale, que $X_s \subset g(X')$, que les schémas
$X'$, $S'$ soient affines et que $\Gamma(X', g^*\mathcal{F})$
soit un $\Gamma(S', \mathcal{O}_{S'})$-module projectif.
Notons que $g^*\mathcal{F}$ est universellement pur sur $S'$, voir le
lemme \ref{flat-lemma-affine-locally-projective-pure}.
Soit $W' \subset X \times_S S'$ l'image du morphisme \'etale
$X' \to X \times_S S'$. Notons que $W$ est ouvert et quasi-compact sur
$S'$. Posons
$$
E = \{t \in S' \mid \text{Ass}_{X_t}(\mathcal{F}_t) \subset W' \}.
$$
D'après
Compléments sur les morphismes,
lemme \ref{more-morphisms-lemma-relative-assassin-constructible},
$E$ est une partie constructible de $S'$. D'après le
lemme \ref{flat-lemma-criterion},
nous avons $\Spec(\mathcal{O}_{S', s'}) \subset E$.
D'après
Morphismes, lemme \ref{morphisms-lemma-constructible-containing-open},
$E$ contient un voisinage ouvert $V'$ de $s'$.
En appliquant le
lemme \ref{flat-lemma-criterion}
une nouvelle fois, nous voyons que, pour tout point $s_1$ de l'image de $V'$ dans $S$,
le faisceau $\mathcal{F}$ est pur le long de $X_{s_1}$. Comme
$S' \to S$ est \'etale, l'image de $V'$ dans $S$ est ouverte, ce qui conclut.
\end{proof}


\section{Applications de la pureté}
\label{flat-section-applications-purity}

\noindent
Voici quelques exemples d'applications de la pureté. Le premier lemme
utilise en fait une forme légèrement plus faible de pureté.

\begin{lemma}
\label{flat-lemma-injectivity-map-source-flat-pure}
Soit $f : X \to S$ un morphisme de type fini.
Soit $\mathcal{F}$ un faisceau quasi-cohérent de type fini sur $X$.
Supposons $S$ local de point fermé $s$.
Supposons $\mathcal{F}$ pur le long de $X_s$ et
$\mathcal{F}$ plat sur $S$.
Soit $\varphi : \mathcal{F} \to \mathcal{G}$ un morphisme de
$\mathcal{O}_X$-modules quasi-cohérents. Alors les conditions suivantes sont équivalentes :
\begin{enumerate}
\item l'application sur les fibres $\varphi_x$ est injective pour tout
$x \in \text{Ass}_{X_s}(\mathcal{F}_s)$, et
\item $\varphi$ est injectif.
\end{enumerate}
\end{lemma}

\begin{proof}
Soit $\mathcal{K} = \Ker(\varphi)$. Notre but est de montrer que
$\mathcal{K} = 0$. Pour cela, il suffit de montrer que
$\text{WeakAss}_X(\mathcal{K}) = \emptyset$, voir
Diviseurs, lemme \ref{divisors-lemma-weakly-ass-zero}.
Nous avons
$\text{WeakAss}_X(\mathcal{K}) \subset \text{WeakAss}_X(\mathcal{F})$, voir
Diviseurs, lemme \ref{divisors-lemma-ses-weakly-ass}.
Comme $\mathcal{F}$ est plat, le
lemme \ref{flat-lemma-bourbaki-finite-type-general-base}
montre que $\text{WeakAss}_X(\mathcal{F}) \subset \text{Ass}_{X/S}(\mathcal{F})$.
Par pureté, tout point $x'$ de $\text{Ass}_{X/S}(\mathcal{F})$
est une généralisation d'un point de $X_s$, et est donc la
spécialisation d'un point $x \in \text{Ass}_{X_s}(\mathcal{F}_s)$, d'après le
lemme \ref{flat-lemma-associated-point-specializes}.
L'injectivité de $\varphi_x$ entraîne donc celle de
$\varphi_{x'}$, d'où $\mathcal{K}_{x'} = 0$.
\end{proof}

\begin{proposition}
\label{flat-proposition-finite-presentation-flat-pure-is-projective}
Soit $f : X \to S$ un morphisme de schémas affine de présentation finie.
Soit $\mathcal{F}$ un $\mathcal{O}_X$-module quasi-cohérent de
présentation finie, plat sur $S$. Alors les conditions suivantes
sont équivalentes :
\begin{enumerate}
\item $f_*\mathcal{F}$ est localement projectif sur $S$, et
\item $\mathcal{F}$ est pur relativement à $S$.
\end{enumerate}
En particulier, pour un homomorphisme d'anneaux $A \to B$ de présentation finie et
un $B$-module de présentation finie $N$, plat sur $A$, nous avons :
$N$ est projectif comme $A$-module si et seulement si $\widetilde{N}$
sur $\Spec(B)$ est pur relativement à $\Spec(A)$.
\end{proposition}

\begin{proof}
L'implication (1) $\Rightarrow$ (2) est le
lemme \ref{flat-lemma-affine-locally-projective-pure}.
Supposons $\mathcal{F}$ pur relativement à $S$.
Notons que, d'après le
lemme \ref{flat-lemma-finite-type-flat-pure-along-fibre-is-universal},
cela entraîne que $\mathcal{F}$ reste pur après tout changement de base. D'après
Descente, lemme \ref{descent-lemma-locally-projective-descends},
il suffit de montrer que $f_*\mathcal{F}$ est localement projectif sur $S$ pour la topologie fpqc.
Choisissons $s \in S$. Nous allons montrer que la restriction de
$f_*\mathcal{F}$ à un voisinage \'etale de $s$ est localement projective.
En effet, d'après le
lemme \ref{flat-lemma-finite-presentation-flat-along-fibre},
après remplacement de $S$ par un voisinage \'etale élémentaire affine
de $s$, nous pouvons supposer qu'il existe un diagramme
$$
\xymatrix{
X \ar[dr] & & X' \ar[ll]^g \ar[ld] \\
& S &
}
$$
de schémas affines et de présentation finie sur $S$,
où $g$ est \'etale, $X_s \subset g(X')$ et
$\Gamma(X', g^*\mathcal{F})$ est un $\Gamma(S, \mathcal{O}_S)$-module projectif.
Notons que, dans ce cas, $g^*\mathcal{F}$ est universellement pur sur $S$, voir le
lemme \ref{flat-lemma-affine-locally-projective-pure}.
Ainsi, d'après le
lemme \ref{flat-lemma-criterion},
l'ouvert $g(X')$ contient les points de
$\text{Ass}_{X/S}(\mathcal{F})$ au-dessus de $\Spec(\mathcal{O}_{S, s})$.
Posons
$$
E = \{t \in S \mid \text{Ass}_{X_t}(\mathcal{F}_t) \subset g(X') \}.
$$
D'après
Compléments sur les morphismes,
lemme \ref{more-morphisms-lemma-relative-assassin-constructible},
$E$ est une partie constructible de $S$. Nous avons vu que
$\Spec(\mathcal{O}_{S, s}) \subset E$. D'après
Morphismes, lemme \ref{morphisms-lemma-constructible-containing-open},
$E$ contient un voisinage ouvert de $s$. Ainsi, après
remplacement de $S$ par un voisinage affine de $s$, nous pouvons supposer que
$\text{Ass}_{X/S}(\mathcal{F}) \subset g(X')$.
D'après le
lemme \ref{flat-lemma-base-change-universally-flat},
cela signifie que
$$
\Gamma(X, \mathcal{F}) \longrightarrow \Gamma(X', g^*\mathcal{F})
$$
est $\Gamma(S, \mathcal{O}_S)$-universellement injective.
D'après Algèbre, lemme \ref{algebra-lemma-pure-submodule-ML},
nous concluons que $\Gamma(X, \mathcal{F})$ est de Mittag-Leffler comme
$\Gamma(S, \mathcal{O}_S)$-module. Comme
$\Gamma(X, \mathcal{F})$ est engendré par un ensemble dénombrable et plat comme
$\Gamma(S, \mathcal{O}_S)$-module, nous concluons
qu'il est projectif d'après
Algèbre, lemme \ref{algebra-lemma-countgen-projective}.
\end{proof}

\noindent
La proposition permet d'améliorer certains résultats précédents.
Le lemme suivant améliore la
proposition \ref{flat-proposition-finite-presentation-flat-at-point}.

\begin{lemma}
\label{flat-lemma-flat-finite-presentation-affine-neighbourhood-projective}
Soit $f : X \to S$ un morphisme localement de présentation finie.
Soit $\mathcal{F}$ un $\mathcal{O}_X$-module quasi-cohérent
de présentation finie. Soit $x \in X$ avec $s = f(x) \in S$.
Si $\mathcal{F}$ est plat en $x$ sur $S$, il existe un voisinage
\'etale élémentaire affine $(S', s') \to (S, s)$ et
un ouvert affine $U' \subset X \times_S S'$ contenant $x' = (x, s')$
tels que $\Gamma(U', \mathcal{F}|_{U'})$ soit un
$\Gamma(S', \mathcal{O}_{S'})$-module projectif.
\end{lemma}

\begin{proof}
Au cours de la démonstration, nous pouvons remplacer $X$ par un voisinage ouvert de $x$
et $S$ par un voisinage \'etale élémentaire de $s$.
Ainsi, par ouverture du lieu de platitude (voir
Compléments sur les morphismes, théorème \ref{more-morphisms-theorem-openness-flatness}),
nous pouvons supposer $\mathcal{F}$ plat sur $S$.
Nous pouvons supposer $S$ et $X$ affines.
Après une nouvelle restriction de $X$, nous pouvons supposer que tout
point de $\text{Ass}_{X_s}(\mathcal{F}_s)$ est une généralisation de $x$.
Cette propriété est conservée lorsqu'on remplace $(S, s)$ par un voisinage
\'etale élémentaire. Nous pouvons donc appliquer le
lemme \ref{flat-lemma-finite-presentation-flat-along-fibre}
pour nous ramener à la situation où il existe un diagramme
$$
\xymatrix{
X \ar[dr] & & X' \ar[ll]^g \ar[ld] \\
& S &
}
$$
de schémas affines et de présentation finie sur $S$,
où $g$ est \'etale, $X_s \subset g(X')$ et
$\Gamma(X', g^*\mathcal{F})$ est un $\Gamma(S, \mathcal{O}_S)$-module projectif.
Notons que, dans ce cas, $g^*\mathcal{F}$ est universellement pur sur $S$, voir le
lemme \ref{flat-lemma-affine-locally-projective-pure}.

\medskip\noindent
Soit $U \subset g(X')$ un voisinage ouvert affine de $x$.
Nous affirmons que $\mathcal{F}|_U$ est pur le long de $U_s$. Si cela est démontré,
le lemme en résulte, car $\mathcal{F}|_U$ sera pur relativement à $S$
après restriction de $S$, voir le
lemme \ref{flat-lemma-flat-finite-presentation-purity-open},
et la projectivité résultera alors de la
proposition \ref{flat-proposition-finite-presentation-flat-pure-is-projective}.
Pour démontrer l'assertion, il faut montrer, après remplacement de $(S, s)$
par un voisinage \'etale élémentaire quelconque, que tout point $\xi$ de
$\text{Ass}_{U/S}(\mathcal{F}|_U)$ au-dessus d'un
$s' \in S$, $s' \leadsto s$, se spécialise en un point de $U_s$.
Comme $U \subset g(X')$, nous pouvons trouver un $\xi' \in X'$ tel que
$g(\xi') = \xi$. Puisque $g^*\mathcal{F}$ est pur sur $S$, le
lemme \ref{flat-lemma-associated-point-specializes}
montre qu'il existe une spécialisation $\xi' \leadsto x'$ avec
$x' \in \text{Ass}_{X'_s}(g^*\mathcal{F}_s)$. Alors
$g(x') \in \text{Ass}_{X_s}(\mathcal{F}_s)$ (voir par exemple le
lemme \ref{flat-lemma-etale-weak-assassin-up-down}
appliqué au morphisme \'etale $X'_s \to X_s$ de schémas noethériens),
donc $g(x') \leadsto x$ par le choix de $X$ ci-dessus ! Puisque
$x \in U$, nous concluons que $g(x') \in U$. Ainsi,
$\xi = g(\xi') \leadsto g(x') \in U_s$, comme voulu.
\end{proof}

\noindent
Le lemme suivant améliore le
lemme \ref{flat-lemma-finite-type-flat-at-point-free-variant}.

\begin{lemma}
\label{flat-lemma-flat-finite-type-affine-neighbourhood-projective}
Soit $f : X \to S$ un morphisme localement de type fini.
Soit $\mathcal{F}$ un $\mathcal{O}_X$-module quasi-cohérent
de type fini. Soit $x \in X$ avec $s = f(x) \in S$.
Si $\mathcal{F}$ est plat en $x$ sur $S$, il existe un voisinage
\'etale élémentaire affine $(S', s') \to (S, s)$ et
un ouvert affine $U' \subset X \times_S \Spec(\mathcal{O}_{S', s'})$
contenant $x' = (x, s')$ tels que
$\Gamma(U', \mathcal{F}|_{U'})$ soit un
$\mathcal{O}_{S', s'}$-module libre.
\end{lemma}

\begin{proof}
La question est locale pour la topologie de Zariski sur $X$ et $S$. Nous pouvons donc supposer
$X$ et $S$ affines. Nous pouvons alors trouver une immersion fermée
$i : X \to \mathbf{A}^n_S$ sur $S$. Il est clair qu'il suffit de
démontrer le lemme pour le faisceau $i_*\mathcal{F}$ sur $\mathbf{A}^n_S$
et le point $i(x)$. Nous nous ramenons ainsi au cas où $X \to S$ est
de présentation finie. Après remplacement de $S$ par
$\Spec(\mathcal{O}_{S', s'})$ et de $X$ par un ouvert de
$X \times_S \Spec(\mathcal{O}_{S', s'})$, nous pouvons supposer que
$\mathcal{F}$ est de présentation finie, voir la
proposition \ref{flat-proposition-finite-type-flat-at-point}.
Dans ce cas, nous pouvons invoquer le
lemme \ref{flat-lemma-flat-finite-presentation-affine-neighbourhood-projective}
et
Algèbre, théorème \ref{algebra-theorem-projective-free-over-local-ring},
pour conclure.
\end{proof}

\begin{lemma}
\label{flat-lemma-flat-finite-type-local-colimit-free}
Soit $A \to B$ un homomorphisme local d'anneaux locaux essentiellement de
type fini. Soit $N$ un $B$-module de type fini, plat comme $A$-module.
Si $A$ est hensélien, alors $N$ est une limite inductive filtrante
$$
N = \colim_i F_i
$$
de $A$-modules libres $F_i$ telle que toutes les applications de transition
$u_i : F_i \to F_{i'}$ du système induisent des applications injectives
$\overline{u}_i : F_i/\mathfrak m_AF_i \to F_{i'}/\mathfrak m_AF_{i'}$.
De plus, $N$ est un $A$-module de Mittag-Leffler.
\end{lemma}

\begin{proof}
Nous pouvons trouver un morphisme de type fini $X \to S = \Spec(A)$,
un point $x \in X$ au-dessus du point fermé $s$ de $S$ et un
$\mathcal{O}_X$-module quasi-cohérent de type fini $\mathcal{F}$ tels que
$\mathcal{F}_x \cong N$ comme $A$-module. Après restriction de $X$,
nous pouvons supposer que tout point de $\text{Ass}_{X_s}(\mathcal{F}_s)$ se spécialise
en $x$. D'après le
lemme \ref{flat-lemma-flat-finite-type-affine-neighbourhood-projective},
il existe un système fondamental de voisinages ouverts affines
$U_i \subset X$ de $x$ tels que $\Gamma(U_i, \mathcal{F})$ soit
un $A$-module libre $F_i$. Notons que, si $U_{i'} \subset U_i$, alors
$$
F_i/\mathfrak m_AF_i = \Gamma(U_{i, s}, \mathcal{F}_s)
\longrightarrow
\Gamma(U_{i', s}, \mathcal{F}_s) = F_{i'}/\mathfrak m_AF_{i'}
$$
est injective, car une section du noyau aurait son support dans
une partie fermée de $X_s$ ne rencontrant pas $x$, ce qui contredit
le choix de $X$ ci-dessus. Comme les applications $F_i \to F_{i'}$ sont
$A$-universellement injectives (lemme \ref{flat-lemma-universally-injective-local}),
il en résulte que $N$ est
de Mittag-Leffler d'après
Algèbre, lemme \ref{algebra-lemma-colimit-universally-injective-ML}.
\end{proof}

\noindent
On pourra omettre le lemme suivant lors d'une première lecture.

\begin{lemma}
\label{flat-lemma-flat-finite-type-local-valuation-ring-has-content}
Soit $A \to B$ un homomorphisme local d'anneaux locaux essentiellement de
type fini. Soit $N$ un $B$-module de type fini, plat comme $A$-module.
Si $A$ est un anneau de valuation, alors tout élément de $N$ possède un idéal de contenu
$I \subset A$ (Compléments d'algèbre, définition
\ref{more-algebra-definition-content-ideal}). De plus, $I$ est un
idéal principal.
\end{lemma}

\begin{proof}
La dernière assertion résulte du fait que $I$ est un idéal de type
fini d'après
Compléments d'algèbre, lemme \ref{more-algebra-lemma-content-finitely-generated},
et Algèbre, lemme \ref{algebra-lemma-characterize-valuation-ring}.

\medskip\noindent
Démontrons l'existence de $I$.
Soit $A \subset A^h$ l'hensélisé. Soit $B'$ le localisé
de $B \otimes_A A^h$ en l'idéal maximal
$\mathfrak m_B \otimes A^h + B \otimes \mathfrak m_{A^h}$.
Alors $B \to B'$ est plat, donc fidèlement plat.
Soit $N' = N \otimes_B B'$.
Soit $x \in N$ et soit $x' \in N'$ son image.
Nous affirmons que, pour un idéal $I \subset A$, nous avons
$x \in IN \Leftrightarrow x' \in IN'$.
En effet, $N/IN \to N'/IN'$ est le produit tensoriel de $B \to B'$
avec $N/IN$, et $B \to B'$ est universellement injectif d'après
Algèbre, lemme \ref{algebra-lemma-faithfully-flat-universally-injective}.
D'après Compléments d'algèbre, lemme \ref{more-algebra-lemma-henselization-valuation-ring},
et Algèbre, lemme \ref{algebra-lemma-ideals-valuation-ring},
l'homomorphisme $A \to A^h$ définit une bijection préservant
l'inclusion $I \mapsto IA^h$ entre les ensembles d'idéaux. Nous concluons que
$x$ possède un idéal de contenu dans $A$ si et seulement si $x'$ possède un idéal de
contenu dans $A^h$. L'assertion pour $x' \in N'$ résulte du
lemme \ref{flat-lemma-flat-finite-type-local-colimit-free} et d'
Algèbre, lemme \ref{algebra-lemma-minimal-contains}.
\end{proof}

\noindent
Voici une application.

\begin{lemma}
\label{flat-lemma-proper-flat-over-dvr-reduced-fibre}
Soit $X \to \Spec(R)$ un morphisme propre et plat, où $R$ est un anneau de valuation.
Si la fibre spéciale est réduite, alors $X$ et toutes les fibres de $X \to \Spec(R)$
sont réduits.
\end{lemma}

\begin{proof}
Supposons la fibre spéciale $X_s$ réduite.
Soit $x \in X$ un point quelconque et montrons que $\mathcal{O}_{X, x}$
est réduit ; cela prouvera que $X$ est réduit.
Soit $x \leadsto x'$ une spécialisation avec $x'$
dans la fibre spéciale ; une telle spécialisation existe,
car un morphisme propre est fermé. Considérons l'anneau
local $A = \mathcal{O}_{X, x'}$. Alors $\mathcal{O}_{X, x}$
est un localisé de $A$ ; il suffit donc de montrer que
$A$ est réduit. Soit $a \in A$ non nul et soit $I = (\pi) \subset R$ son
idéal de contenu, voir le
lemme \ref{flat-lemma-flat-finite-type-local-valuation-ring-has-content}.
Alors $a = \pi a'$ et $a'$ s'envoie sur un élément non nul de
$A/\mathfrak mA$, où $\mathfrak m \subset R$ est l'idéal maximal.
Si $a$ est nilpotent, il en est de même de $a'$, car $\pi$ n'est pas diviseur de zéro
par platitude de $A$ sur $R$.
Mais $a'$ s'envoie sur un élément non nul de l'anneau réduit
$A/\mathfrak m A = \mathcal{O}_{X_s, x'}$.
Cela donne une contradiction à moins que $A$ ne soit réduit, ce
que nous voulions montrer.

\medskip\noindent
Bien entendu, si $X$ est réduit, il en est de même de la fibre générique de $X$ sur $R$.
Si $\mathfrak p \subset R$ est un
idéal premier, alors $R/\mathfrak p$ est un anneau de valuation d'après
Algèbre, lemme \ref{algebra-lemma-make-valuation-rings}.
En reprenant l'argument pour le changement de base de $X$
à $R/\mathfrak p$, on montre donc que la fibre au-dessus de $\mathfrak p$
est réduite.
\end{proof}






\section{Foncteurs de platification}
\label{flat-section-flattening-functors}

\noindent
Soit $S$ un schéma. Rappelons qu'un foncteur
$F : (\Sch/S)^{opp} \to \textit{Ensembles}$ commute aux limites projectives filtrantes
si, pour tout système projectif filtrant
$\{T_i\}_{i \in I}$ de schémas affines de limite $T$, nous avons
$F(T) = \colim_i F(T_i)$.

\begin{situation}
\label{flat-situation-iso}
Soit $f : X \to S$ un morphisme de schémas.
Soit $u : \mathcal{F} \to \mathcal{G}$ un homomorphisme de
$\mathcal{O}_X$-modules quasi-cohérents. Pour tout schéma $T$ sur
$S$, nous noterons $u_T : \mathcal{F}_T \to \mathcal{G}_T$ le
changement de base de $u$ à $T$ ; autrement dit, $u_T$ est l'image réciproque
de $u$ par la projection $X_T = X \times_S T \to X$.
Dans cette situation, nous pouvons considérer le foncteur
\begin{equation}
\label{flat-equation-iso}
F_{iso} : (\Sch/S)^{opp} \longrightarrow \textit{Ensembles}, \quad
T \longrightarrow \left\{
\begin{matrix}
\{*\} & \text{si} & u_T \text{ est un isomorphisme}, \\
\emptyset & \text{sinon.} &
\end{matrix}
\right.
\end{equation}
Il existe des variantes $F_{inj}$, $F_{surj}$, $F_{zero}$ où l'on demande que
$u_T$ soit injectif, surjectif ou nul.
\end{situation}

\begin{lemma}
\label{flat-lemma-iso-sheaf}
Plaçons-nous dans la situation \ref{flat-situation-iso}.
\begin{enumerate}
\item Chacun des foncteurs $F_{iso}$, $F_{inj}$, $F_{surj}$, $F_{zero}$
satisfait à la condition de faisceau pour la topologie fpqc.
\item Si $f$ est quasi-compact et $\mathcal{G}$ de type fini,
alors $F_{surj}$ commute aux limites projectives filtrantes.
\item Si $f$ est quasi-compact et $\mathcal{F}$ de type fini, alors
$F_{zero}$ commute aux limites projectives filtrantes.
\item Si $f$ est quasi-compact, $\mathcal{F}$ de type fini et
$\mathcal{G}$ de présentation finie, alors $F_{iso}$ commute aux limites projectives filtrantes.
\end{enumerate}
\end{lemma}

\begin{proof}
Soit $\{T_i \to T\}_{i \in I}$ un recouvrement fpqc de schémas sur $S$.
Posons $X_i = X_{T_i} = X \times_S T_i$ et $u_i = u_{T_i}$.
Notons que $\{X_i \to X_T\}_{i \in I}$ est un recouvrement fpqc de $X_T$, voir
Topologies, lemme \ref{topologies-lemma-fpqc}.
En particulier, pour tout $x \in X_T$, il existe un $i \in I$
et un $x_i \in X_i$ d'image $x$. Comme
$\mathcal{O}_{X_T, x} \to \mathcal{O}_{X_i, x_i}$ est plat, donc
fidèlement plat (voir
Algèbre, lemme \ref{algebra-lemma-local-flat-ff}),
nous concluons que $(u_i)_{x_i}$ est injectif, surjectif, bijectif ou nul
si et seulement si $(u_T)_x$ est injectif, surjectif, bijectif ou nul.
D'où la partie (1) du lemme.

\medskip\noindent
Démonstration de (2). Supposons $f$ quasi-compact et $\mathcal{G}$ de type fini.
Soit $T = \lim_{i \in I} T_i$ une limite dirigée de $S$-schémas affines
et supposons $u_T$ surjectif.
Posons $X_i = X_{T_i} = X \times_S T_i$ et
$u_i = u_{T_i} : \mathcal{F}_i = \mathcal{F}_{T_i}
\to \mathcal{G}_i = \mathcal{G}_{T_i}$.
Pour démontrer (2), il faut montrer que $u_i$ est surjectif pour un certain $i$.
Choisissons $i_0 \in I$ et remplaçons $I$ par $\{i \mid i \geq i_0\}$.
Comme $f$ est quasi-compact, le schéma $X_{i_0}$ est quasi-compact.
Nous pouvons donc choisir des ouverts affines $W_1, \ldots, W_m \subset X$
et un recouvrement ouvert affine
$X_{i_0} = U_{1, i_0} \cup \ldots \cup U_{m, i_0}$ tels que
$U_{j, i_0}$ s'envoie dans $W_j$ par la projection $X_{i_0} \to X$.
Pour tout $i \in I$, soit $U_{j, i}$ l'image réciproque de $U_{j, i_0}$.
En posant $U_j = \lim_i U_{j, i}$, nous voyons que $X_T = U_1 \cup \ldots \cup U_m$
est un recouvrement ouvert affine de $X_T$. Il suffit alors de montrer, pour un
$j \in \{1, \ldots, m\}$ donné, que $u_i|_{U_{j, i}}$ est surjectif pour un certain
$i = i(j) \in I$. À l'aide de
Propriétés, lemme \ref{properties-lemma-finite-type-module},
cela se traduit par le problème algébrique suivant :
soient $A$ un anneau et $u : M \to N$ une application de $A$-modules.
Supposons que $R = \colim_{i \in I} R_i$ soit une limite inductive filtrante
de $A$-algèbres. Si $N$ est un $A$-module de type fini et si
$u \otimes 1 : M \otimes_A R \to N \otimes_A R$ est surjective, alors,
pour un certain $i$, l'application
$u \otimes 1 : M \otimes_A R_i \to N \otimes_A R_i$ est surjective.
C'est
Algèbre, lemme \ref{algebra-lemma-module-map-property-in-colimit}, partie (2).

\medskip\noindent
Démonstration de (3). Les mêmes arguments que dans la démonstration de (2)
ramènent au problème algébrique suivant :
soient $A$ un anneau et $u : M \to N$ une application de $A$-modules.
Supposons que $R = \colim_{i \in I} R_i$ soit une limite inductive filtrante
de $A$-algèbres. Si $M$ est un $A$-module de type fini et si
$u \otimes 1 : M \otimes_A R \to N \otimes_A R$ est nulle, alors,
pour un certain $i$, l'application
$u \otimes 1 : M \otimes_A R_i \to N \otimes_A R_i$ est nulle.
C'est
Algèbre, lemme \ref{algebra-lemma-module-map-property-in-colimit}, partie (1).

\medskip\noindent
Démonstration de (4).
Supposons $f$ quasi-compact et $\mathcal{F}, \mathcal{G}$ de présentation finie.
En raisonnant exactement comme dans le paragraphe précédent
(et en utilisant en outre
Propriétés, lemme \ref{properties-lemma-finite-presentation-module}),
la partie (3) se traduit par l'énoncé algébrique suivant :
soient $A$ un anneau et $u : M \to N$ une application de $A$-modules.
Supposons que $R = \colim_{i \in I} R_i$ soit une limite inductive filtrante
de $A$-algèbres. Supposons $M$ de type fini comme $A$-module, $N$ de présentation
finie comme $A$-module, et
$u \otimes 1 : M \otimes_A R \to N \otimes_A R$ un isomorphisme.
Alors, pour un certain $i$, l'application
$u \otimes 1 : M \otimes_A R_i \to N \otimes_A R_i$ est un isomorphisme.
C'est
Algèbre, lemme \ref{algebra-lemma-module-map-property-in-colimit}, partie (3).
\end{proof}

\begin{situation}
\label{flat-situation-flat-at-point}
Soit $(A, \mathfrak m_A)$ un anneau local.
Notons $\mathcal{C}$ la catégorie dont les objets sont les $A$-algèbres
$A'$ qui sont des anneaux locaux et dont la structure d'algèbre
$A \to A'$ est un homomorphisme local d'anneaux locaux.
Un morphisme entre objets $A', A''$ de $\mathcal{C}$ est un
homomorphisme local $A' \to A''$ de $A$-algèbres.
Soient $A \to B$ un homomorphisme local d'anneaux locaux et $M$ un $B$-module.
Si $A'$ est un objet de $\mathcal{C}$, nous posons $B' = B \otimes_A A'$
et $M' = M \otimes_A A'$, considéré comme $B'$-module.
Pour $A' \in \Ob(\mathcal{C})$, considérons la condition
\begin{equation}
\label{flat-equation-flat-at-primes}
\forall \mathfrak q \in V(\mathfrak m_{A'}B' + \mathfrak m_B B')
\subset \Spec(B') :
M'_{\mathfrak q}\text{ est plat sur }A'.
\end{equation}
Notons l'analogie avec
Compléments d'algèbre, équation (\ref{more-algebra-equation-flat-at-primes}).
En particulier, si $A' \to A''$ est un morphisme de $\mathcal{C}$ et si
(\ref{flat-equation-flat-at-primes}) est satisfaite pour $A'$, elle l'est pour $A''$, voir
Compléments d'algèbre,
lemme \ref{more-algebra-lemma-base-change-flat-at-primes}.
Nous obtenons donc un foncteur
\begin{equation}
\label{flat-equation-flat-at-point}
F_{lf} : \mathcal{C} \longrightarrow \textit{Ensembles}, \quad
A' \longrightarrow \left\{
\begin{matrix}
\{*\} & \text{si }(\ref{flat-equation-flat-at-primes})\text{ est satisfaite}, \\
\emptyset & \text{sinon.} &
\end{matrix}
\right.
\end{equation}
\end{situation}

\begin{lemma}
\label{flat-lemma-flat-at-point}
Plaçons-nous dans la situation \ref{flat-situation-flat-at-point}.
\begin{enumerate}
\item Si $A' \to A''$ est un morphisme plat dans $\mathcal{C}$,
alors $F_{lf}(A') = F_{lf}(A'')$.
\item Si $A \to B$ est essentiellement de présentation finie et
$M$ est un $B$-module de présentation finie, alors $F_{lf}$ est compatible aux
limites : si $\{A_i\}_{i \in I}$ est un
système inductif filtrant d'objets de $\mathcal{C}$, alors
$F_{lf}(\colim_i A_i) = \colim_i F_{lf}(A_i)$.
\end{enumerate}
\end{lemma}

\begin{proof}
La partie (1) est un cas particulier de
Compléments d'algèbre,
lemme \ref{more-algebra-lemma-flat-descent-flat-at-primes}.
La partie (2) est un cas particulier de
Compléments d'algèbre,
lemme \ref{more-algebra-lemma-limit-preserving-flat-at-primes}.
\end{proof}

\begin{lemma}
\label{flat-lemma-flat-at-point-finite}
Dans la situation \ref{flat-situation-flat-at-point}, soient $B \to C$ un homomorphisme local de
$A$-algèbres locales et $N$ un $C$-module. Notons
$F'_{lf} : \mathcal{C} \to \textit{Ensembles}$ le foncteur associé au couple
$(C, N)$. Si $M \cong N$ comme $B$-modules et si $B \to C$ est fini, alors
$F_{lf} = F'_{lf}$.
\end{lemma}

\begin{proof}
Soit $A'$ un objet de $\mathcal{C}$. Posons $C' = C \otimes_A A'$
et $N' = N \otimes_A A'$, comme dans les définitions de $B'$, $M'$ de la
situation \ref{flat-situation-flat-at-point}.
Notons que $M' \cong N'$ comme $B'$-modules.
L'hypothèse que $B \to C$ est fini a deux conséquences :
(a) $\mathfrak m_C = \sqrt{\mathfrak m_B C}$ et (b)
$B' \to C'$ est fini. La conséquence (a) entraîne que
$$
V(\mathfrak m_{A'}C' + \mathfrak m_C C')
=
\left(
\Spec(C') \to \Spec(B')
\right)^{-1}V(\mathfrak m_{A'}B' + \mathfrak m_B B').
$$
Supposons $\mathfrak q \subset V(\mathfrak m_{A'}B' + \mathfrak m_B B')$.
Alors $M'_{\mathfrak q}$ est plat sur $A'$ si et seulement si
le $C'_{\mathfrak q}$-module $N'_{\mathfrak q}$ est plat sur $A'$
(car ils sont isomorphes comme $A'$-modules), si et seulement si,
pour tout idéal maximal $\mathfrak r$ de $C'_{\mathfrak q}$,
le module $N'_{\mathfrak r}$ est plat sur $A'$ (voir
Algèbre, lemme \ref{algebra-lemma-flat-localization}).
Comme $B'_{\mathfrak q} \to C'_{\mathfrak q}$ est fini d'après (b),
les idéaux maximaux de $C'_{\mathfrak q}$ correspondent exactement
aux idéaux premiers de $C'$ au-dessus de $\mathfrak q$ (voir
Algèbre, lemme \ref{algebra-lemma-integral-going-up}),
et ces idéaux premiers appartiennent tous à
$V(\mathfrak m_{A'}C' + \mathfrak m_C C')$ d'après
l'égalité ci-dessus. Le résultat du lemme en découle.
\end{proof}

\begin{lemma}
\label{flat-lemma-flat-at-point-go-up}
Dans la
situation \ref{flat-situation-flat-at-point},
supposons que $B \to C$ soit un homomorphisme local plat d'anneaux
locaux. Posons $N = M \otimes_B C$. Notons
$F'_{lf} : \mathcal{C} \to \textit{Ensembles}$ le foncteur associé
au couple $(C, N)$. Alors $F_{lf} = F'_{lf}$.
\end{lemma}

\begin{proof}
Soit $A'$ un objet de $\mathcal{C}$. Posons $C' = C \otimes_A A'$
et $N' = N \otimes_A A' = M' \otimes_{B'} C'$, comme dans les définitions
de $B'$, $M'$ de la
situation \ref{flat-situation-flat-at-point}. Notons que
$$
V(\mathfrak m_{A'}B' + \mathfrak m_B B')
=
\Spec( \kappa(\mathfrak m_B) \otimes_A \kappa(\mathfrak m_{A'}) )
$$
et de même pour $V(\mathfrak m_{A'}C' + \mathfrak m_C C')$.
L'homomorphisme d'anneaux
$$
\kappa(\mathfrak m_B) \otimes_A \kappa(\mathfrak m_{A'})
\longrightarrow
\kappa(\mathfrak m_C) \otimes_A \kappa(\mathfrak m_{A'})
$$
est fidèlement plat, donc
$V(\mathfrak m_{A'}C' + \mathfrak m_C C') \to
V(\mathfrak m_{A'}B' + \mathfrak m_B B')$ est surjectif.
Enfin, si $\mathfrak r \in V(\mathfrak m_{A'}C' + \mathfrak m_C C')$
s'envoie sur $\mathfrak q \in V(\mathfrak m_{A'}B' + \mathfrak m_B B')$, alors
$M'_{\mathfrak q}$ est plat sur $A'$ si et seulement si
$N'_{\mathfrak r}$ est plat sur $A'$, car $B' \to C'$ est plat, voir
Algèbre, lemme \ref{algebra-lemma-flatness-descends-more-general}.
Le lemme résulte formellement de ces remarques.
\end{proof}

\begin{situation}
\label{flat-situation-free-at-generic-points}
Soit $f : X \to S$ un morphisme lisse à fibres géométriquement
irréductibles. Soit $\mathcal{F}$ un $\mathcal{O}_X$-module quasi-cohérent de
type fini. Pour tout schéma $T$ sur $S$, nous noterons
$\mathcal{F}_T$ le changement de base de $\mathcal{F}$ à $T$ ; autrement
dit, $\mathcal{F}_T$ est l'image réciproque de $\mathcal{F}$ par la
projection $X_T = X \times_S T \to X$. Notons que $X_T \to T$
est lisse à fibres géométriquement irréductibles, voir
Morphismes, lemme \ref{morphisms-lemma-base-change-smooth}, et
Compléments sur les morphismes,
lemme \ref{more-morphisms-lemma-base-change-fibres-geometrically-irreducible}.
Soit $p \geq 0$ un entier. Pour un point $t \in T$, considérons la
condition
\begin{equation}
\label{flat-equation-free-at-generic-point-fibre}
\mathcal{F}_T \text{ est libre de rang }p\text{ au voisinage de }\xi_t
\end{equation}
où $\xi_t$ est le point générique de la fibre $X_t$. Cette condition
pour tout $t \in T$ est stable par changement de base ; nous obtenons donc un foncteur
\begin{equation}
\label{flat-equation-free-at-generic-points}
H_p : (\Sch/S)^{opp} \longrightarrow \textit{Ensembles}, \quad
T \longrightarrow \left\{
\begin{matrix}
\{*\} & \text{si }\mathcal{F}_T\text{ vérifie
(\ref{flat-equation-free-at-generic-point-fibre}) }\forall t\in T, \\
\emptyset & \text{sinon.}
\end{matrix}
\right.
\end{equation}
\end{situation}

\begin{lemma}
\label{flat-lemma-free-at-generic-points}
Plaçons-nous dans la situation \ref{flat-situation-free-at-generic-points}.
\begin{enumerate}
\item Le foncteur $H_p$ satisfait à la condition de faisceau pour la topologie fpqc.
\item Si $\mathcal{F}$ est de présentation finie, alors le foncteur $H_p$ est
commute aux limites projectives filtrantes.
\end{enumerate}
\end{lemma}

\begin{proof}
Soit $\{T_i \to T\}_{i \in I}$ un recouvrement fpqc\footnote{Il est assez facile de
montrer que $H_p$
est un faisceau pour la topologie fppf en utilisant le fait que les morphismes plats de présentation
finie sont ouverts. C'est tout ce dont nous aurons besoin par la suite. Mais il est assez
plaisant de démontrer directement qu'il satisfait aussi à la condition de faisceau
pour la topologie fpqc.} de schémas sur $S$.
Posons $X_i = X_{T_i} = X \times_S T_i$ et notons $\mathcal{F}_i$
l'image réciproque de $\mathcal{F}$ sur $X_i$. Supposons que $\mathcal{F}_i$ satisfasse à
(\ref{flat-equation-free-at-generic-point-fibre}) pour tout $i$.
Choisissons $t \in T$ et notons $\xi_t \in X_T$ le point générique de $X_t$.
Il faut montrer que $\mathcal{F}$ est libre dans un voisinage de $\xi_t$.
Pour un certain $i \in I$, nous pouvons trouver un $t_i \in T_i$ d'image $t$.
Notons $\xi_i \in X_i$ le point générique de $X_{t_i}$, de sorte que
$\xi_i$ s'envoie sur $\xi_t$. Le fait que $\mathcal{F}_i$ soit libre de rang
$p$ dans un voisinage de $\xi_i$ entraîne que
$(\mathcal{F}_i)_{x_i} \cong \mathcal{O}_{X_i, x_i}^{\oplus p}$,
ce qui entraîne que
$\mathcal{F}_{T, \xi_t} \cong \mathcal{O}_{X_T, \xi_t}^{\oplus p}$,
puisque $\mathcal{O}_{X_T, \xi_t} \to \mathcal{O}_{X_i, x_i}$ est plat, voir
par exemple
Algèbre, lemme \ref{algebra-lemma-finite-projective-descends}.
Il existe donc un voisinage affine $U$ de $\xi_t$ dans $X_T$
et une surjection
$\mathcal{O}_U^{\oplus p} \to \mathcal{F}_U = \mathcal{F}_T|_U$, voir
Modules, lemme \ref{modules-lemma-finite-type-surjective-on-stalk}.
Après restriction de $T$, nous pouvons supposer $U \to T$ surjectif. Ainsi,
$U \to T$ est un morphisme lisse de schémas affines à fibres géométriquement
irréductibles. De plus, pour tout $t' \in T$, l'application induite
$$
\alpha :
\mathcal{O}_{U, \xi_{t'}}^{\oplus p}
\longrightarrow
\mathcal{F}_{U, \xi_{t'}}
$$
est un isomorphisme (car, par le même argument que précédemment, le module de droite
est libre de rang $p$). Il résulte du
lemme \ref{flat-lemma-induction-step}
que
$$
\Gamma(U, \mathcal{O}_U^{\oplus p})
\otimes_{\Gamma(T, \mathcal{O}_T)} \mathcal{O}_{T, t'}
\longrightarrow
\Gamma(U, \mathcal{F}_U)
\otimes_{\Gamma(T, \mathcal{O}_T)} \mathcal{O}_{T, t'}
$$
est injective pour tout $t' \in T$. La surjection $\alpha$ est donc un
isomorphisme. Cela achève la démonstration de (1).

\medskip\noindent
Supposons $\mathcal{F}$ de présentation finie.
Soit $T = \lim_{i \in I} T_i$ une limite dirigée de $S$-schémas affines
et supposons que $\mathcal{F}_T$ satisfasse à
(\ref{flat-equation-free-at-generic-point-fibre}).
Posons $X_i = X_{T_i} = X \times_S T_i$ et notons $\mathcal{F}_i$
l'image réciproque de $\mathcal{F}$ sur $X_i$.
Notons $U \subset X_T$ le sous-schéma ouvert des points où
$\mathcal{F}_T$ est plat sur $T$, voir
Compléments sur les morphismes, théorème \ref{more-morphisms-theorem-openness-flatness}.
Par hypothèse, tout point générique de toute fibre appartient à $U$, c'est-à-dire que
$U \to T$ est un morphisme lisse surjectif à fibres géométriquement irréductibles.
Nous pouvons restreindre $U$ et supposer $U$ quasi-compact.
À l'aide de
Limites, lemme \ref{limits-lemma-descend-opens},
nous pouvons trouver un $i \in I$ et un ouvert quasi-compact $U_i \subset X_i$
dont l'image réciproque dans $X_T$ est $U$. En augmentant $i$, nous pouvons
supposer $\mathcal{F}_i|_{U_i}$ plat sur $T_i$, voir
Limites, lemme \ref{limits-lemma-descend-module-flat-finite-presentation}.
En particulier, $\mathcal{F}_i|_{U_i}$ est localement libre de type fini,
donc définit une fonction rang
localement constante $\rho : U_i \to \{0, 1, 2, \ldots \}$.
Notons $(U_i)_p \subset U_i$ la partie ouverte et fermée
où $\rho$ a pour valeur $p$. Soit $V_i \subset T_i$
l'image de $(U_i)_p$ ; notons que $V_i$ est ouvert et quasi-compact.
Par hypothèse, l'image de $T \to T_i$ est contenue dans $V_i$.
Il existe donc un $i' \geq i$ tel que $T_{i'} \to T_i$ se factorise
par $V_i$, d'après
Limites, lemme \ref{limits-lemma-descend-opens}.
Alors $\mathcal{F}_{i'}$ satisfait à (\ref{flat-equation-free-at-generic-point-fibre}),
comme voulu. Certains détails sont omis.
\end{proof}

\begin{lemma}
\label{flat-lemma-pre-flat-dimension-n}
Soit $f : X \to S$ un morphisme de schémas localement de type fini.
Soit $\mathcal{F}$ un $\mathcal{O}_X$-module quasi-cohérent de type
fini. Soit $n \geq 0$. Les conditions suivantes sont équivalentes :
\begin{enumerate}
\item pour $s \in S$, la partie fermée $Z \subset X_s$ des points
où $\mathcal{F}$ n'est pas plat sur $S$ (voir le
lemme \ref{flat-lemma-open-in-fibre-where-flat})
vérifie $\dim(Z) < n$, et
\item pour $x \in X$ tel que $\mathcal{F}$ ne soit pas plat en $x$
sur $S$, nous avons $\text{trdeg}_{\kappa(f(x))}(\kappa(x)) < n$.
\end{enumerate}
Si ces conditions sont satisfaites, elles le restent après tout changement de base.
\end{lemma}

\begin{proof}
Soit $x \in X$ un point au-dessus de $s \in S$.
Alors la dimension de l'adhérence de $\{x\}$ dans $X_s$
est $\text{trdeg}_{\kappa(s)}(\kappa(x))$, d'après
Variétés, lemme \ref{varieties-lemma-dimension-locally-algebraic}.
Réciproquement, si $Z \subset X_s$ est une partie fermée
de dimension $d$, il existe un point $x \in Z$
tel que $\text{trdeg}_{\kappa(s)}(\kappa(x)) = d$ (même référence).
Les conditions (1) et (2) sont donc équivalentes (même fibre par fibre).
L'assertion sur le changement de base résulte de
Morphismes, lemmes \ref{morphisms-lemma-base-change-module-flat} et
\ref{morphisms-lemma-dimension-fibre-after-base-change}.
\end{proof}

\begin{definition}
\label{flat-definition-flat-dimension-n}
Soit $f : X \to S$ un morphisme de schémas localement de type fini.
Soit $\mathcal{F}$ un $\mathcal{O}_X$-module quasi-cohérent de type fini.
Soit $n \geq 0$.
Nous disons que {\it $\mathcal{F}$ est plat sur $S$ en dimensions $\geq n$}
si les conditions équivalentes du lemme \ref{flat-lemma-pre-flat-dimension-n}
sont satisfaites.
\end{definition}

\begin{situation}
\label{flat-situation-flat-dimension-n}
Soit $f : X \to S$ un morphisme de schémas localement de type fini.
Soit $\mathcal{F}$ un $\mathcal{O}_X$-module quasi-cohérent de type
fini. Pour tout schéma $T$ sur $S$, nous noterons $\mathcal{F}_T$ le
changement de base de $\mathcal{F}$ à $T$ ; autrement dit, $\mathcal{F}_T$
est l'image réciproque de $\mathcal{F}$ par la projection
$X_T = X \times_S T \to X$. Notons que $X_T \to T$ est de type fini
et que $\mathcal{F}_T$ est un $\mathcal{O}_{X_T}$-module
de type fini (Morphismes, lemme
\ref{morphisms-lemma-base-change-finite-type}, et
Modules, lemme \ref{modules-lemma-pullback-finite-type}).
Soit $n \geq 0$. La définition \ref{flat-definition-flat-dimension-n} et le
lemme \ref{flat-lemma-pre-flat-dimension-n} donnent un foncteur
\begin{equation}
\label{flat-equation-flat-dimension-n}
F_n : (\Sch/S)^{opp} \longrightarrow \textit{Ensembles}, \quad
T \longrightarrow \left\{
\begin{matrix}
\{*\} & \text{si }\mathcal{F}_T\text{ est plat sur }T\text{ en dimension }\dim \geq n, \\
\emptyset & \text{sinon.}
\end{matrix}
\right.
\end{equation}
\end{situation}

\begin{lemma}
\label{flat-lemma-flat-dimension-n}
Plaçons-nous dans la situation \ref{flat-situation-flat-dimension-n}.
\begin{enumerate}
\item Le foncteur $F_n$ satisfait à la condition de faisceau pour la topologie fpqc.
\item Si $f$ est quasi-compact et localement de présentation finie
et si $\mathcal{F}$ est de présentation finie, alors le foncteur $F_n$ est
commute aux limites projectives filtrantes.
\end{enumerate}
\end{lemma}

\begin{proof}
Soit $\{T_i \to T\}_{i \in I}$ un recouvrement fpqc de schémas sur $S$.
Posons $X_i = X_{T_i} = X \times_S T_i$ et notons $\mathcal{F}_i$
l'image réciproque de $\mathcal{F}$ sur $X_i$. Supposons $\mathcal{F}_i$
plat sur $T_i$ en dimensions $\geq n$ pour tout $i$. Soit $t \in T$.
Choisissons un indice $i$ et un point $t_i \in T_i$ d'image $t$.
Considérons le diagramme cartésien
$$
\xymatrix{
X_{\Spec(\mathcal{O}_{T, t})} \ar[d] &
X_{\Spec(\mathcal{O}_{T_i, t_i})} \ar[d] \ar[l] \\
\Spec(\mathcal{O}_{T, t}) &
\Spec(\mathcal{O}_{T_i, t_i}) \ar[l]
}
$$
Comme le morphisme horizontal inférieur est plat, il résulte de
Compléments sur les morphismes, lemme \ref{more-morphisms-lemma-flat-locus-base-change},
que l'ensemble $Z_i \subset X_{t_i}$ où $\mathcal{F}_i$ n'est pas plat
sur $T_i$ et l'ensemble $Z \subset X_t$ où $\mathcal{F}_T$ n'est pas plat
sur $T$ sont liés par la relation $Z_i = Z_{\kappa(t_i)}$. Nous en déduisons
que $\mathcal{F}_T$ est plat sur $T$ en dimensions $\geq n$, d'après
Morphismes, lemme \ref{morphisms-lemma-dimension-fibre-after-base-change}.

\medskip\noindent
Supposons que $f$ soit quasi-compact et localement de présentation finie et
que $\mathcal{F}$ soit de présentation finie.
Dans ce paragraphe, nous ramenons d'abord la démonstration de (2) au cas où
$f$ est de présentation finie.
Soit $T = \lim_{i \in I} T_i$ une limite dirigée de
$S$-schémas affines et supposons $\mathcal{F}_T$ plat en dimensions
$\geq n$. Posons $X_i = X_{T_i} = X \times_S T_i$ et notons $\mathcal{F}_i$
l'image réciproque de $\mathcal{F}$ sur $X_i$. Il faut montrer que
$\mathcal{F}_i$ est plat en dimensions $\geq n$ pour un certain $i$.
Choisissons $i_0 \in I$ et remplaçons $I$ par $\{i \mid i \geq i_0\}$.
Comme $T_{i_0}$ est affine (donc quasi-compact), il existe un nombre fini
d'ouverts affines $W_j \subset S$, $j = 1, \ldots, m$, et un recouvrement ouvert affine
$T_{i_0} = \bigcup_{j = 1, \ldots, m} V_{j, i_0}$
tels que $T_{i_0} \to S$ envoie $V_{j, i_0}$ dans $W_j$.
Pour $i \geq i_0$, notons $V_{j, i}$ l'image réciproque de $V_{j, i_0}$
dans $T_i$. Si nous pouvons montrer, pour chaque $j$, qu'il existe un $i$ tel que
$\mathcal{F}_{V_{j, i_0}}$ soit plat en dimensions $\geq n$,
nous aurons conclu. Nous nous ramenons ainsi au cas où $S$ est affine.
Dans ce cas, $X$ est quasi-compact et admet un recouvrement ouvert
affine fini $X = W_1 \cup \ldots \cup W_m$. Dans ce cas,
le résultat pour $(X \to S, \mathcal{F})$ équivaut au résultat pour
$(\coprod W_j, \coprod \mathcal{F}|_{W_j})$. Nous pouvons donc supposer
$f$ de présentation finie.

\medskip\noindent
Supposons $f$ de présentation finie et $\mathcal{F}$ de présentation
finie. Notons $U \subset X_T$ le sous-schéma ouvert des points où
$\mathcal{F}_T$ est plat sur $T$, voir
Compléments sur les morphismes, théorème \ref{more-morphisms-theorem-openness-flatness}.
Par hypothèse, toute fibre de $Z = X_T \setminus U$ sur
$T$ est de dimension $< n$. D'après
Limites, lemme \ref{limits-lemma-approximate-given-relative-dimension},
nous pouvons trouver un sous-schéma fermé $Z \subset Z' \subset X_T$
tel que $\dim(Z'_t) < n$ pour tout $t \in T$ et que
$Z' \to X_T$ soit de présentation finie. D'après
Limites, lemmes \ref{limits-lemma-descend-finite-presentation} et
\ref{limits-lemma-descend-closed-immersion-finite-presentation},
il existe un $i \in I$ et un sous-schéma fermé $Z'_i \subset X_i$
de présentation finie dont le changement de base à $T$ est $Z'$. D'après
Limites, lemme \ref{limits-lemma-limit-dimension},
nous pouvons supposer toutes les fibres de $Z'_i \to T_i$ de dimension $< n$. D'après
Limites, lemme \ref{limits-lemma-descend-module-flat-finite-presentation},
nous pouvons supposer $\mathcal{F}_i|_{X_i \setminus T'_i}$
plat sur $T_i$. Cela entraîne que $\mathcal{F}_i$ est
plat en dimensions $\geq n$ ; nous utilisons ici que $Z' \to X_T$ est de présentation
finie, donc que le complémentaire $X_T \setminus Z'$ est quasi-compact !
La partie (2) est ainsi démontrée et la démonstration du lemme est achevée.
\end{proof}

\begin{situation}
\label{flat-situation-flat}
Soit $f : X \to S$ un morphisme de schémas.
Soit $\mathcal{F}$ un $\mathcal{O}_X$-module quasi-cohérent.
Pour tout schéma $T$ sur $S$, nous noterons $\mathcal{F}_T$ le
changement de base de $\mathcal{F}$ à $T$ ; autrement dit, $\mathcal{F}_T$
est l'image réciproque de $\mathcal{F}$ par la projection
$X_T = X \times_S T \to X$. Comme le changement de base d'un module plat
est plat, nous obtenons un foncteur
\begin{equation}
\label{flat-equation-flat}
F_{flat} : (\Sch/S)^{opp} \longrightarrow \textit{Ensembles}, \quad
T \longrightarrow \left\{
\begin{matrix}
\{*\} & \text{si } \mathcal{F}_T \text{ est plat sur }T, \\
\emptyset & \text{sinon.}
\end{matrix}
\right.
\end{equation}
\end{situation}

\begin{lemma}
\label{flat-lemma-flat}
Plaçons-nous dans la situation \ref{flat-situation-flat}.
\begin{enumerate}
\item Le foncteur $F_{flat}$ satisfait à la condition de faisceau pour la topologie fpqc.
\item Si $f$ est quasi-compact et localement de présentation finie
et si $\mathcal{F}$ est de présentation finie, alors le foncteur
$F_{flat}$ commute aux limites projectives filtrantes.
\end{enumerate}
\end{lemma}

\begin{proof}
La partie (1) résulte de l'énoncé suivant : si $T' \to T$ est un morphisme plat
surjectif de schémas sur $S$, alors $\mathcal{F}_{T'}$ est plat sur $T'$
si et seulement si $\mathcal{F}_T$ est plat sur $T$, voir
Compléments sur les morphismes, lemme \ref{more-morphisms-lemma-flat-locus-base-change}.
La partie (2) résulte de
Limites, lemme \ref{limits-lemma-descend-module-flat-finite-presentation},
après réduction au cas où $X$ et $S$ sont affines (comparer avec
la démonstration du
lemme \ref{flat-lemma-flat-dimension-n}).
\end{proof}




\section{Stratifications de platification}
\label{flat-section-flattening}

\noindent
Seulement les définitions. Le lecteur qui cherche une
``stratification de platitude générique'' consultera
Compléments sur les morphismes, section
\ref{more-morphisms-section-generic-flatness-stratification}.

\begin{definition}
\label{flat-definition-flattening}
Soit $X \to S$ un morphisme de schémas.
Soit $\mathcal{F}$ un $\mathcal{O}_X$-module quasi-cohérent.
Nous disons que {\it le platificateur universel de $\mathcal{F}$ existe}
si le foncteur $F_{flat}$ défini dans la situation \ref{flat-situation-flat}
est représentable par un schéma $S'$ sur $S$.
Nous disons que {\it le platificateur universel de $X$ existe}
si le platificateur universel de $\mathcal{O}_X$ existe.
\end{definition}

\noindent
Notons que si le platificateur universel $S'$\footnote{Le schéma $S'$ est parfois
appelé {\it flatificateur universel}. Dans \cite{GruRay}, il est appelé
le {\it platificateur universel}. L'existence du platificateur universel
ne doit pas être confondue avec le type de résultats étudiés dans
Compléments d'algèbre, section \ref{more-algebra-section-blowup-flat}.} de
$\mathcal{F}$ existe, alors le morphisme $S' \to S$ est un monomorphisme
surjectif de schémas tel que $\mathcal{F}_{S'}$ soit plat sur $S'$
et tel qu'un morphisme $T \to S$ se factorise par $S'$ si et seulement si
$\mathcal{F}_T$ est plat sur $T$.

\begin{example}
\label{flat-example-no-universal-flattening}
Soit $X = S = \Spec(k[x, y])$, où $k$ est un corps. Soit
$\mathcal{F} = \widetilde{M}$, où $M = k[x, x^{-1}, y]/(y)$.
Pour une $k[x, y]$-algèbre $A$, posons $F_{flat}(A) = F_{flat}(\Spec(A))$.
Alors $F_{flat}(k[x, y]/(x, y)^n) = \{*\}$ pour tout $n$, tandis que
$F_{flat}(k[[x, y]]) = \emptyset$. Cela signifie que $F_{flat}$ n'est pas
représentable (même par un espace algébrique, voir
Espaces formels, Lemme
\ref{formal-spaces-lemma-map-into-algebraic-space}).
Ainsi, le platificateur universel n'existe pas dans ce cas.
\end{example}

\noindent
Nous définissons (comparer à
Topologie, Remarque \ref{topology-remark-locally-finite-stratification})
une {\it stratification} (localement finie, au sens schématique) d'un schéma $S$
comme la donnée de sous-schémas fermés $Z_i \subset S$ indexés par un
ensemble partiellement ordonné $I$ tels que
$S = \bigcup Z_i$ (ensemblistement), que tout point de $S$ possède
un voisinage ne rencontrant qu'un nombre fini de $Z_i$, et que
$$
Z_i \cap Z_j = \bigcup\nolimits_{k \leq i, j} Z_k.
$$
En posant $S_i = Z_i \setminus \bigcup_{j < i} Z_j$, la stratification
proprement dite est la décomposition $S = \coprod S_i$ en
sous-schémas localement fermés. Souvent, nous indiquons seulement les strates
$S_i$ et laissons au lecteur la construction des sous-schémas fermés $Z_i$.
Étant donnée une stratification, nous obtenons un monomorphisme
$$
S' = \coprod\nolimits_{i \in I} S_i \longrightarrow S.
$$
Nous l'appellerons le {\it monomorphisme associé à la stratification}.
Avec cette terminologie, nous pouvons définir ce que signifie posséder une
stratification de platification.

\begin{definition}
\label{flat-definition-flattening-stratification}
Soit $X \to S$ un morphisme de schémas.
Soit $\mathcal{F}$ un $\mathcal{O}_X$-module quasi-cohérent.
Nous disons que $\mathcal{F}$ admet une {\it stratification de platification}
si le foncteur $F_{flat}$ défini dans la situation \ref{flat-situation-flat}
est représentable par un monomorphisme $S' \to S$ associé
à une stratification de $S$ par des sous-schémas localement fermés.
Nous disons que $X$ admet une {\it stratification de platification}
si $\mathcal{O}_X$ admet une stratification de platification.
\end{definition}

\noindent
Lorsqu'une stratification de platification existe, il importe souvent
de comprendre l'ensemble d'indices qui étiquette les strates et son ordre partiel.
Cela est souvent lié aux rangs de modules. Par exemple, si
$X = S$ et si $\mathcal{F}$ est un $\mathcal{O}_S$-module de présentation finie,
alors la stratification de platification existe et est donnée par les idéaux de Fitting
de $\mathcal{F}$, voir
Diviseurs, Lemme \ref{divisors-lemma-finite-presentation-module}.




\section{Stratification de platification sur un anneau artinien}
\label{flat-section-flattening-artinian}

\noindent
Une stratification de platification existe lorsque le schéma de base est le spectre
d'un anneau artinien.

\begin{lemma}
\label{flat-lemma-flattening-stratification-artinian}
Soit $S$ le spectre d'un anneau artinien.
Pour tout schéma $X$ sur $S$ et tout $\mathcal{O}_X$-module quasi-cohérent,
il existe un platificateur universel. En fait, le platificateur universel
est donné par une immersion fermée $S' \to S$ et constitue donc aussi une
stratification de platification pour $\mathcal{F}$.
\end{lemma}

\begin{proof}
Choisissons un recouvrement ouvert affine $X = \bigcup U_i$.
Alors $F_{flat}$ est le produit des foncteurs associés à
chacun des couples $(U_i, \mathcal{F}|_{U_i})$.
Il suffit donc de démontrer le résultat pour chaque
$(U_i, \mathcal{F}|_{U_i})$.
Dans le cas affine, le lemme résulte immédiatement de
Compléments d'algèbre,
Lemme \ref{more-algebra-lemma-flattening-artinian-universal-property}.
\end{proof}






\section{Platification d'une application}
\label{flat-section-flattening-map}

\noindent
Le théorème \ref{flat-theorem-flattening-map}
est la clef des énoncés de platification qui suivent.

\begin{lemma}
\label{flat-lemma-universally-separating}
Soit $S$ un schéma.
Soit $g : X' \to X$ un morphisme plat de schémas sur $S$,
où $X$ est localement de type fini sur $S$.
Soit $\mathcal{F}$ un $\mathcal{O}_X$-module quasi-cohérent de type fini
qui est plat sur $S$. Si $\text{Ass}_{X/S}(\mathcal{F}) \subset g(X')$,
alors l'application canonique
$$
\mathcal{F} \longrightarrow g_*g^*\mathcal{F}
$$
est injective et le reste après tout changement de base.
\end{lemma}

\begin{proof}
La dernière assertion signifie que $\mathcal{F}_T \to (g_T)_*g_T^*\mathcal{F}_T$
est injective pour tout morphisme $T \to S$. L'hypothèse
$\text{Ass}_{X/S}(\mathcal{F}) \subset g(X')$ est préservée par changement de base, voir
Diviseurs, Lemme \ref{divisors-lemma-base-change-relative-assassin} et
Remarque \ref{divisors-remark-base-change-relative-assassin}.
Il en va de même de l'hypothèse de platitude et de type fini.
Il suffit donc de démontrer l'injectivité de la flèche affichée.
Soit $\mathcal{K} = \Ker(\mathcal{F} \to g_*g^*\mathcal{F})$.
Notre but est de montrer que $\mathcal{K} = 0$.
Pour ce faire, il suffit de montrer que
$\text{WeakAss}_X(\mathcal{K}) = \emptyset$, voir
Diviseurs, Lemme \ref{divisors-lemma-weakly-ass-zero}.
Nous avons
$\text{WeakAss}_X(\mathcal{K}) \subset \text{WeakAss}_X(\mathcal{F})$, voir
Diviseurs, Lemme \ref{divisors-lemma-ses-weakly-ass}.
Puisque $\mathcal{F}$ est plat, il résulte du
lemme \ref{flat-lemma-bourbaki-finite-type-general-base}
que $\text{WeakAss}_X(\mathcal{F}) \subset \text{Ass}_{X/S}(\mathcal{F})$.
Par hypothèse, tout point $x$ de $\text{Ass}_{X/S}(\mathcal{F})$
est l'image d'un certain $x' \in X'$. Comme $g$ est plat, l'homomorphisme
d'anneaux locaux $\mathcal{O}_{X, x} \to \mathcal{O}_{X', x'}$
est fidèlement plat ; par conséquent, l'application
$$
\mathcal{F}_x
\longrightarrow
g^*\mathcal{F}_{x'} =
\mathcal{F}_x \otimes_{\mathcal{O}_{X, x}} \mathcal{O}_{X', x'}
$$
est injective (voir
Algèbre, Lemme \ref{algebra-lemma-faithfully-flat-universally-injective}).
Cela implique que $\mathcal{K}_x = 0$, comme souhaité.
\end{proof}

\begin{lemma}
\label{flat-lemma-flattening-module-map}
Soit $A$ un anneau. Soit $u : M \to N$ une application surjective de $A$-modules.
Si $M$ est projectif comme $A$-module, alors il existe un idéal
$I \subset A$ tel que, pour tout homomorphisme d'anneaux $\varphi : A \to B$,
les conditions suivantes soient équivalentes :
\begin{enumerate}
\item $u \otimes 1 : M \otimes_A B \to N \otimes_A B$ est un
isomorphisme ;
\item $\varphi(I) = 0$.
\end{enumerate}
\end{lemma}

\begin{proof}
Comme $M$ est projectif, nous pouvons trouver un $A$-module projectif $C$
tel que $F = M \oplus C$ soit un $A$-module libre.
En remplaçant $u$ par $u \oplus 1 : F = M \oplus C \to N \oplus C$,
nous voyons que nous pouvons supposer $M$ libre. Dans ce cas, soit $I$
l'idéal de $A$ engendré par les coefficients de tous les éléments de
$\Ker(u)$ relativement à une base (fixée) de $M$.
Ce choix convient : puisque $u$ est surjective et que
$\otimes_A B$ exact à droite, $\Ker(u \otimes 1)$ est
l'image de $\Ker(u) \otimes_A B$ dans $M \otimes_A B$.
\end{proof}

\begin{theorem}
\label{flat-theorem-flattening-map}
Dans la
situation \ref{flat-situation-iso},
supposons que
\begin{enumerate}
\item $f$ soit de présentation finie ;
\item $\mathcal{F}$ soit de présentation finie, plat sur $S$ et
pur relativement à $S$ ;
\item $u$ soit surjective.
\end{enumerate}
Alors $F_{iso}$ est représentable par une immersion fermée $Z \to S$.
De plus, $Z \to S$ est de présentation finie si $\mathcal{G}$ est
de présentation finie.
\end{theorem}

\begin{proof}
Nous utiliserons sans autre mention le fait que $\mathcal{F}$ est universellement pur
sur $S$, voir
le lemme \ref{flat-lemma-finite-type-flat-pure-along-fibre-is-universal}.
D'après
le lemme \ref{flat-lemma-iso-sheaf}
et
Descente, Lemmes \ref{descent-lemma-closed-immersion} et
\ref{descent-lemma-descent-data-sheaves},
la question est locale pour la topologie étale sur $S$.
Il suffit donc de montrer, étant donné $s \in S$, qu'il existe
un voisinage étale de $(S, s)$ sur lequel le théorème soit vrai.

\medskip\noindent
En utilisant
le lemme \ref{flat-lemma-finite-presentation-flat-along-fibre}
et après avoir remplacé $S$ par un voisinage étale élémentaire de $s$,
nous pouvons supposer qu'il existe un diagramme commutatif
$$
\xymatrix{
X \ar[dr] & & X' \ar[ll]^g \ar[ld] \\
& S &
}
$$
de schémas de présentation finie sur $S$,
où $g$ est étale, $X_s \subset g(X')$, les schémas $X'$ et $S$ sont affines,
et $\Gamma(X', g^*\mathcal{F})$ est un $\Gamma(S, \mathcal{O}_S)$-module projectif.
Notons que $g^*\mathcal{F}$ est universellement pur sur $S$, voir
le lemme \ref{flat-lemma-affine-locally-projective-pure}.
Ainsi, d'après
le lemme \ref{flat-lemma-criterion},
nous voyons que l'ouvert $g(X')$ contient les points de
$\text{Ass}_{X/S}(\mathcal{F})$ situés au-dessus de $\Spec(\mathcal{O}_{S, s})$.
Posons
$$
E = \{t \in S \mid \text{Ass}_{X_t}(\mathcal{F}_t) \subset g(X') \}.
$$
D'après
Compléments sur les morphismes,
Lemme \ref{more-morphisms-lemma-relative-assassin-constructible},
$E$ est une partie constructible de $S$. Nous avons vu que
$\Spec(\mathcal{O}_{S, s}) \subset E$. D'après
Morphismes, Lemme \ref{morphisms-lemma-constructible-containing-open},
nous voyons que $E$ contient un voisinage ouvert de $s$. Ainsi, après avoir
remplacé $S$ par un voisinage affine plus petit de $s$, nous pouvons supposer que
$\text{Ass}_{X/S}(\mathcal{F}) \subset g(X')$.

\medskip\noindent
Puisque nous avons supposé $u$ surjective, nous avons $F_{iso} = F_{inj}$. D'après
le lemme \ref{flat-lemma-universally-separating},
il s'ensuit que $u : \mathcal{F} \to \mathcal{G}$ est injective si et seulement si
$g^*u : g^*\mathcal{F} \to g^*\mathcal{G}$ est injective, et cela reste vrai
après tout changement de base. Nous nous sommes donc ramenés au cas où,
outre les hypothèses du théorème, $X \to S$ est un morphisme de
schémas affines et $\Gamma(X, \mathcal{F})$ est un
$\Gamma(S, \mathcal{O}_S)$-module projectif. Ce cas résulte immédiatement du
lemme \ref{flat-lemma-flattening-module-map}.

\medskip\noindent
Pour voir que $Z$ est de présentation finie si $\mathcal{G}$ est de présentation
finie, combinons la partie (4) du
lemme \ref{flat-lemma-iso-sheaf}
avec
Limites, Remarque \ref{limits-remark-limit-preserving}.
\end{proof}

\begin{lemma}
\label{flat-lemma-Weil-restriction-closed-subschemes}
Soit $f:X\to S$ un morphisme de schémas de présentation finie,
plat et pur. Soit $Y$ un sous-schéma fermé de $X$. Soit $F=f_*Y$ le
foncteur de restriction de Weil de $Y$ le long de $f$, défini par
$$
F : (\Sch/S)^{opp} \to \textit{Ensembles}, \quad
T \mapsto
\left\{
\begin{matrix}
\{*\} & \text{si} & Y_T\to X_T \text{ est un isomorphisme, }\\
\emptyset & \text{sinon.} &
\end{matrix}
\right.
$$
Alors $F$ est représentable par une immersion fermée $Z\to S$. De plus,
$Z\to S$ est de présentation finie si $Y\to S$ l'est.
\end{lemma}

\begin{proof}
Soit $\mathcal{I}$ le faisceau d'idéaux définissant $Y$ dans $X$ et soit
$u:\mathcal{O}_X\to\mathcal{O}_X/\mathcal{I}$ la surjection.
Alors, pour un $S$-schéma $T$, l'immersion fermée $Y_T\to X_T$
est un isomorphisme si et seulement si $u_T$ est un isomorphisme. Le résultat
découle donc du
théorème \ref{flat-theorem-flattening-map}.
\end{proof}



\section{Platification dans le cas local}
\label{flat-section-flattening-local}

\noindent
Dans cette section, nous commençons à appliquer les résultats précédents afin d'obtenir une
ébauche de la stratification de platification.

\begin{theorem}
\label{flat-theorem-flattening-local}
Dans la
situation \ref{flat-situation-flat-at-point},
supposons $A$ hensélien, $B$ essentiellement de type fini sur $A$, et
$M$ un $B$-module fini. Il existe alors un idéal
$I \subset A$ tel que $A/I$ coreprésente le foncteur $F_{lf}$ sur la catégorie
$\mathcal{C}$. Autrement dit, étant donné un homomorphisme local d'anneaux locaux
$\varphi : A \to A'$, avec $B' = B \otimes_A A'$ et $M' = M \otimes_A A'$,
les conditions suivantes sont équivalentes :
\begin{enumerate}
\item $\forall \mathfrak q \in V(\mathfrak m_{A'}B' + \mathfrak m_B B')
\subset \Spec(B') :
M'_{\mathfrak q}\text{ est plat sur }A'$ ;
\item $\varphi(I) = 0$.
\end{enumerate}
Si $B$ est essentiellement de présentation finie sur $A$ et $M$
de présentation finie sur $B$, alors $I$ est un idéal de type fini.
\end{theorem}

\begin{proof}
Choisissons un homomorphisme d'anneaux de type fini $A \to C$, un $C$-module fini $N$
et un idéal premier $\mathfrak q$ de $C$ tels que $B = C_{\mathfrak q}$
et $M = N_{\mathfrak q}$. Dans la suite, lorsque nous disons
``le théorème vaut pour $(N/C/A, \mathfrak q)$'', nous entendons qu'il
vaut pour $(A \to B, M)$, où $B = C_{\mathfrak q}$ et
$M = N_{\mathfrak q}$. D'après le
lemme \ref{flat-lemma-flat-at-point-go-up},
le foncteur $F_{lf}$ ne change pas si nous remplaçons $B$ par un anneau local
plat sur $B$. Ainsi, puisque $A$ est hensélien, nous pouvons appliquer le
lemme \ref{flat-lemma-existence-algebra}
et supposer qu'il existe un dévissage complet de
$N/C/A$ en $\mathfrak q$.

\medskip\noindent
Soit $(A_i, B_i, M_i, \alpha_i, \mathfrak q_i)_{i = 1, \ldots, n}$
un tel dévissage complet de $N/C/A$ en $\mathfrak q$. Soit
$\mathfrak q'_i \subset A_i$ l'unique idéal premier situé au-dessus de
$\mathfrak q_i \subset B_i$, comme dans la
définition \ref{flat-definition-complete-devissage-at-x-algebra}.
Puisque $C \to A_1$ est surjectif et $N \cong M_1$ comme $C$-modules,
nous voyons d'après le
lemme \ref{flat-lemma-flat-at-point-finite}
qu'il suffit de montrer que le théorème vaut pour $(M_1/A_1/A, \mathfrak q'_1)$.
Puisque $B_1 \to A_1$ est fini et que $\mathfrak q_1$ est l'unique idéal premier
de $B_1$ au-dessus de $\mathfrak q'_1$, nous voyons que
$(A_1)_{\mathfrak q'_1} \to (B_1)_{\mathfrak q_1}$ est fini (voir
Algèbre, Lemme \ref{algebra-lemma-unique-prime-over-localize-below}, ou
Compléments sur les morphismes,
Lemme \ref{more-morphisms-lemma-finite-morphism-single-point-in-fibre}).
Ainsi, d'après le
lemme \ref{flat-lemma-flat-at-point-finite},
il suffit de montrer que le théorème vaut pour $(M_1/B_1/A, \mathfrak q_1)$.

\medskip\noindent
À ce stade, nous pouvons supposer, par récurrence sur la longueur $n$ du
dévissage, que le théorème vaut pour $(M_2/B_2/A, \mathfrak q_2)$.
(Si $n = 1$, alors $M_2 = 0$, qui est plat sur $A$.)
En inversant les deux dernières étapes du paragraphe précédent et en utilisant
$M_2 \cong \Coker(\alpha_2)$ comme $B_1$-modules, nous voyons
que le théorème vaut pour $(\Coker(\alpha_1)/B_1/A, \mathfrak q_1)$.

\medskip\noindent
Soit $A'$ un objet de $\mathcal{C}$. À ce stade, nous utilisons le
lemme \ref{flat-lemma-induction-step}
pour voir que si $(M_1 \otimes_A A')_{\mathfrak q'}$ est plat
sur $A'$ pour un idéal premier $\mathfrak q'$ de $B_1 \otimes_A A'$
situé au-dessus de $\mathfrak m_{A'}$, alors
$(\Coker(\alpha_1) \otimes_A A')_{\mathfrak q'}$ est plat sur $A'$.
Nous en concluons que $F_{lf}$ est un sous-foncteur du
foncteur $F'_{lf}$ associé au module
$\Coker(\alpha_1)_{\mathfrak q_1}$ sur $(B_1)_{\mathfrak q_1}$.
D'après le paragraphe précédent, $F'_{lf}$ est coreprésenté par
$A/J$ pour un certain idéal $J \subset A$. Nous pouvons donc remplacer $A$ par
$A/J$ et supposer que $\Coker(\alpha_1)_{\mathfrak q_1}$ est
plat sur $A$.

\medskip\noindent
Puisque $\Coker(\alpha_1)$ est un $B_1$-module pour lequel
il existe un dévissage complet de $N_1/B_1/A$ en $\mathfrak q_1$
et puisque $\Coker(\alpha_1)_{\mathfrak q_1}$ est
plat sur $A$, d'après le
lemme \ref{flat-lemma-complete-devissage-flat-finite-type-module},
nous voyons que $\Coker(\alpha_1)$ est libre comme $A$-module, donc en particulier
plat comme $A$-module. Par conséquent, le
lemme \ref{flat-lemma-induction-step}
implique que $F_{lf}(A')$ est non vide si et seulement si $\alpha \otimes 1_{A'}$
est injective. Soit $N_1 = \Im(\alpha_1) \subset M_1$, de sorte que
nous avons les suites exactes
$$
0 \to N_1 \to M_1 \to \Coker(\alpha_1) \to 0
\quad\text{et}\quad
B_1^{\oplus r_1} \to N_1 \to 0
$$
La platitude de $\Coker(\alpha_1)$ implique que la première suite
est universellement exacte (voir
Algèbre, Lemme \ref{algebra-lemma-flat-universally-injective}).
Ainsi, $\alpha \otimes 1_{A'}$ est injective si et seulement si
$B_1^{\oplus r_1} \otimes_A A' \to N_1 \otimes_A A'$
est un isomorphisme. Enfin, le
théorème \ref{flat-theorem-flattening-map}
montre que ce foncteur est coreprésentable par $A/I$ pour un certain idéal $I$,
et nous concluons que $F_{lf}$ est lui aussi coreprésentable par $A/I$.

\medskip\noindent
Pour démontrer la dernière assertion, supposons que $A \to B$ soit essentiellement de
présentation finie et $M$ de présentation finie sur $B$. Soit $I \subset A$
l'idéal tel que $F_{lf}$ soit coreprésenté par $A/I$.
Écrivons $I = \bigcup I_\lambda$, où $I_\lambda$ parcourt les idéaux de type fini
contenus dans $I$. Comme $F_{lf}(A/I) = \{*\}$,
nous voyons alors que $F_{lf}(A/I_\lambda) = \{*\}$ pour un certain $\lambda$, voir
la partie (2) du lemme \ref{flat-lemma-flat-at-point}.
Cela implique clairement que $I = I_\lambda$.
\end{proof}

\begin{remark}
\label{flat-remark-flattening-local-scheme-theoretic}
Voici une reformulation schématique du
théorème \ref{flat-theorem-flattening-local}.
Soit $(X, x) \to (S, s)$ un morphisme de schémas pointés
qui est localement de type fini. Soit $\mathcal{F}$ un
$\mathcal{O}_X$-module quasi-cohérent de type fini.
Supposons $S$ local hensélien, de point fermé $s$.
Il existe un sous-schéma fermé $Z \subset S$ ayant la propriété suivante :
pour tout morphisme de schémas pointés $(T, t) \to (S, s)$, les conditions
suivantes sont équivalentes :
\begin{enumerate}
\item $\mathcal{F}_T$ est plat sur $T$ en tous les points de la fibre
$X_t$ qui s'envoient sur $x \in X_s$ ;
\item $\Spec(\mathcal{O}_{T, t}) \to S$ se factorise par $Z$.
\end{enumerate}
De plus, si $X \to S$ est de présentation finie en $x$ et $\mathcal{F}_x$
de présentation finie sur $\mathcal{O}_{X, x}$, alors $Z \to S$
est de présentation finie.
\end{remark}

\noindent
À ce stade, nous pouvons déduire gratuitement du résultat précédent des énoncés
très généraux. Notons que le cas le plus intéressant est peut-être
celui où $E = X_s$ !

\begin{lemma}
\label{flat-lemma-freebie}
Soit $S$ le spectre d'un anneau local hensélien, de point fermé $s$.
Soit $X \to S$ un morphisme de schémas localement de type fini.
Soit $\mathcal{F}$ un $\mathcal{O}_X$-module quasi-cohérent de type fini.
Soit $E \subset X_s$ une partie. Il existe un sous-schéma fermé
$Z \subset S$ ayant la propriété suivante : pour tout morphisme de schémas pointés
$(T, t) \to (S, s)$, les conditions suivantes sont équivalentes :
\begin{enumerate}
\item $\mathcal{F}_T$ est plat sur $T$ en tous les points de la fibre
$X_t$ qui s'envoient sur un point de $E \subset X_s$ ;
\item $\Spec(\mathcal{O}_{T, t}) \to S$ se factorise par $Z$.
\end{enumerate}
De plus, si $X \to S$ est localement de présentation finie,
si $\mathcal{F}$ est de présentation finie et si $E \subset X_s$ est
fermé et quasi-compact, alors $Z \to S$ est de présentation finie.
\end{lemma}

\begin{proof}
Pour $x \in X_s$, notons $Z_x \subset S$ le sous-schéma fermé obtenu dans la
remarque \ref{flat-remark-flattening-local-scheme-theoretic}.
Il est alors clair que $Z = \bigcap_{x \in E} Z_x$ convient !

\medskip\noindent
Pour démontrer la dernière assertion, supposons $X$ localement de présentation finie,
$\mathcal{F}$ de présentation finie et $Z$ fermé et quasi-compact.
Choisissons d'abord un nombre fini d'ouverts affines $W_j \subset X$ tels que
$E \subset \bigcup W_j$. Il suffit clairement de démontrer le
résultat pour chaque morphisme $W_j \to S$, muni du faisceau $\mathcal{F}|_{X_j}$
et de la partie fermée $E \cap W_j$. Nous pouvons donc supposer $X$ affine.
Dans ce cas,
Compléments d'algèbre, Lemme \ref{more-algebra-lemma-limit-preserving-flat-at-primes},
montre que le foncteur défini par (1) ``commute aux limites projectives filtrantes''.
Nous pouvons donc montrer que $Z \to S$ est de présentation finie exactement
comme dans la dernière partie de la démonstration du
théorème \ref{flat-theorem-flattening-local}.
\end{proof}

\begin{remark}
\label{flat-remark-flattening-complete-noetherian}
En remontant la démonstration du
lemme \ref{flat-lemma-freebie}
jusqu'à ses origines, nous découvrons un chemin long et sinueux. Mais si nous supposons que
\begin{enumerate}
\item $f$ est de type fini ;
\item $\mathcal{F}$ est un $\mathcal{O}_X$-module de type fini ;
\item $E = X_s$ ;
\item $S$ est le spectre d'un anneau local noethérien complet,
\end{enumerate}
alors il existe une démonstration reposant entièrement sur une algèbre plus élémentaire,
comme suit : nous nous ramenons d'abord au cas où $X$ est affine en choisissant
un recouvrement ouvert affine fini. Dans ce cas, $Z$ existe d'après
Compléments d'algèbre,
Lemme \ref{more-algebra-lemma-flattening-complete-local-universal-property}.
L'étape essentielle de cette démonstration consiste à construire le sous-schéma fermé $Z$
pas à pas à l'intérieur des troncations
$\Spec(\mathcal{O}_{S, s}/\mathfrak m_s^n)$.
Cela repose sur le fait que les stratifications de platification existent toujours
lorsque la base est artinienne, et sur le fait que
$\mathcal{O}_{S, s} = \lim \mathcal{O}_{S, s}/\mathfrak m_s^n$.
\end{remark}




\section{Variantes d'un lemme}
\label{flat-section-variants-mod-injective}

\noindent
Dans cette section, nous étudions des variantes des
lemmes d'Algèbre \ref{algebra-lemma-mod-injective-general} et
\ref{algebra-lemma-mod-injective}.
La version la plus générale est la
proposition \ref{flat-proposition-finite-type-injective-into-flat-mod-m} ;
elle est énoncée comme \cite[Lemme 4.2.2]{GruRay}, mais la démonstration de
loc.cit.\ ne donne que le résultat plus faible énoncé dans le
lemme \ref{flat-lemma-upstairs-finite-type-injective-into-flat-mod-m}.
La démonstration délicate de la
proposition \ref{flat-proposition-finite-type-injective-into-flat-mod-m}
est due à Ofer Gabber. Comme nous n'avons actuellement aucune application de
cette proposition, nous encourageons le lecteur à passer à la section suivante
après avoir lu la démonstration du
lemme \ref{flat-lemma-upstairs-finite-type-injective-into-flat-mod-m} ;
ce lemme sera utilisé dans la section suivante pour démontrer le
théorème \ref{flat-theorem-check-flatness-at-associated-points}.

\begin{situation}
\label{flat-situation-mod-injective}
Soit $\varphi : A \to B$ un homomorphisme local d'anneaux locaux
qui est essentiellement de type fini. Soient $M$ un $A$-module plat,
$N$ un $B$-module fini et $u : N \to M$ une application de $A$-modules telle que
$\overline{u} : N/\mathfrak m_AN \to M/\mathfrak m_AM$ soit injective.
\end{situation}

\noindent
Dans cette situation, notre but est de montrer que $u$ est $A$-universellement injective,
que $N$ est de présentation finie sur $B$ et que $N$ est plat comme $A$-module.
Si cela est vrai, nous dirons que {\it le lemme vaut} dans la situation donnée.

\begin{lemma}
\label{flat-lemma-Noetherian-finite-type-injective-into-flat-mod-m}
Si, dans la situation \ref{flat-situation-mod-injective}, l'anneau $A$ est noethérien,
alors le lemme vaut.
\end{lemma}

\begin{proof}
En appliquant Algèbre, Lemme \ref{algebra-lemma-mod-injective}, nous voyons que
$u$ est injective et que $N/u(M)$ est plat sur $A$. Alors $u$ est
$A$-universellement injective
(Algèbre, Lemme \ref{algebra-lemma-flat-tor-zero}) et $N$ est $A$-plat
(Algèbre, Lemme \ref{algebra-lemma-flat-ses}). Puisque $B$ est noethérien
dans ce cas, nous voyons que $N$ est de présentation finie.
\end{proof}

\begin{lemma}
\label{flat-lemma-reduce-finite-type-injective-into-flat-mod-m}
Soit $A_0$ un anneau local. Si le lemme vaut pour toute
situation \ref{flat-situation-mod-injective} avec $A = A_0$, où $B$ est une
localisation d'une algèbre de polynômes sur $A$ et $N$ est de présentation finie
sur $B$, alors le lemme vaut pour toute
situation \ref{flat-situation-mod-injective} avec $A = A_0$.
\end{lemma}

\begin{proof}
Soient $A \to B$ et $u : N \to M$ comme dans la situation \ref{flat-situation-mod-injective}.
Écrivons $B = C/I$, où $C$ est la localisation d'une algèbre de polynômes
sur $A$ en un idéal premier. Si nous pouvons montrer que $N$ est de présentation finie comme
$C$-module, cela montre a fortiori que $N$ est de présentation finie
comme $B$-module (voir
Algèbre, Lemme \ref{algebra-lemma-finitely-presented-over-subring}).
Nous pouvons donc supposer que $B$ est la localisation d'une algèbre de polynômes.
Écrivons ensuite $N = B^{\oplus n}/K$ pour un sous-module $K \subset B^{\oplus n}$.
Comme $B/\mathfrak m_AB$ est noethérien (étant essentiellement de type fini
sur un corps), il existe un nombre fini d'éléments $k_1, \ldots, k_s \in K$
tels que, pour $K' = \sum Bk_i$ et $N' = B^{\oplus n}/K'$, la
surjection canonique $N' \to N$ induise un isomorphisme
$N'/\mathfrak m_AN' \cong N/\mathfrak m_AN$.
Maintenant, si le lemme vaut pour la composée $u' : N' \to M$,
alors $u'$ est injective, donc $N' = N$ et $u' = u$. Le lemme vaut donc pour
la situation initiale.
\end{proof}

\begin{lemma}
\label{flat-lemma-henselian-finite-type-injective-into-flat-mod-m}
Si, dans la situation \ref{flat-situation-mod-injective}, l'anneau $A$ est hensélien,
alors le lemme vaut.
\end{lemma}

\begin{proof}
Il suffit de le démontrer lorsque $B$ est essentiellement de présentation finie
sur $A$ et $N$ est de présentation finie sur $B$, voir le
lemme \ref{flat-lemma-reduce-finite-type-injective-into-flat-mod-m}.
Faisons temporairement l'hypothèse supplémentaire que $N$ est plat sur $A$.
Alors $N$ est une colimite filtrante $N = \colim_i F_i$
de $A$-modules libres $F_i$ telle que les applications de transition
$u_{ii'} : F_i \to F_{i'}$ soient injectives modulo $\mathfrak m_A$, voir le
lemme \ref{flat-lemma-flat-finite-type-local-colimit-free}.
Chacune des composées $u_i : F_i \to M$ est $A$-universellement
injective d'après le
lemme \ref{flat-lemma-universally-injective-local} ;
par conséquent, $u = \colim u_i$ est $A$-universellement injective, comme souhaité.

\medskip\noindent
Supposons $A$ local hensélien, $B$ essentiellement
de présentation finie sur $A$ et $N$ de présentation finie sur $B$. D'après le
théorème \ref{flat-theorem-flattening-local},
il existe un idéal de type fini $I \subset A$ tel que
$N/IN$ soit plat sur $A/I$ et que $N/I^2N$ ne soit pas plat sur
$A/I^2$, sauf si $I = 0$. Le résultat du paragraphe précédent montre que
le lemme vaut pour $u \bmod I : N/IN \to M/IM$ sur $A/I$.
Considérons le diagramme commutatif
$$
\xymatrix{
0 \ar[r] &
M \otimes_A I/I^2 \ar[r] &
M/I^2M \ar[r] &
M/IM \ar[r] & 0 \\
&
N \otimes_A I/I^2 \ar[r] \ar[u]^u &
N/I^2N \ar[r] \ar[u]^u &
N/IN \ar[r] \ar[u]^u & 0
}
$$
dont les lignes sont exactes par exactitude à droite de $\otimes$ et parce que
$M$ est plat sur $A$. Notons que la flèche verticale de gauche est l'application
$N/IN \otimes_{A/I} I/I^2 \to M/IM \otimes_{A/I} I/I^2$ et qu'elle est donc
injective. Une chasse au diagramme montre que la flèche inférieure gauche est injective,
c'est-à-dire $\text{Tor}^1_{A/I^2}(I/I^2, M/I^2) = 0$, voir
Algèbre, Remarque \ref{algebra-remark-Tor-ring-mod-ideal}.
Ainsi, $N/I^2N$ est plat sur $A/I^2$ d'après
Algèbre, Lemme \ref{algebra-lemma-what-does-it-mean},
ce qui est une contradiction, sauf si $I = 0$.
\end{proof}

\noindent
Le lemme suivant traite le cas particulier de la
situation \ref{flat-situation-mod-injective} où $M$ est muni
d'une structure de $B$-module et $u$ est $B$-linéaire. C'est le cas le plus
souvent utilisé en pratique, et sa démonstration est nettement plus facile
que celle du cas général.

\begin{lemma}
\label{flat-lemma-upstairs-finite-type-injective-into-flat-mod-m}
Soit $A \to B$ un homomorphisme local d'anneaux locaux qui est
essentiellement de type fini. Soit $u : N \to M$ une application de $B$-modules.
Si $N$ est un $B$-module fini, $M$ est plat sur $A$ et
$\overline{u} : N/\mathfrak m_A N \to M/\mathfrak m_A M$ est injective,
alors $u$ est $A$-universellement injective, $N$ est de présentation finie sur
$B$ et $N$ est plat sur $A$.
\end{lemma}

\begin{proof}
Soit $A \to A^h$ l'hensélisé de $A$. Soit $B'$ la localisation
de $B \otimes_A A^h$ en l'idéal maximal
$\mathfrak m_B \otimes A^h + B \otimes \mathfrak m_{A^h}$.
Comme $B \to B'$ est plat (donc fidèlement plat, voir
Algèbre, Lemme \ref{algebra-lemma-local-flat-ff}),
nous pouvons remplacer $A \to B$ par $A^h \to B'$,
le module $M$ par $M \otimes_B B'$, le module $N$ par $N \otimes_B B'$
et $u$ par $u \otimes \text{id}_{B'}$, voir
Algèbre, Lemmes \ref{algebra-lemma-descend-properties-modules} et
\ref{algebra-lemma-flatness-descends-more-general}.
Nous pouvons ainsi supposer que $A$ est un anneau local hensélien.
Dans ce cas, notre lemme résulte du lemme plus général
\ref{flat-lemma-henselian-finite-type-injective-into-flat-mod-m}.
\end{proof}

\begin{lemma}
\label{flat-lemma-valuation-ring-finite-type-injective-into-flat-mod-m}
Si, dans la situation \ref{flat-situation-mod-injective}, l'anneau $A$ est un
anneau de valuation, alors le lemme vaut.
\end{lemma}

\begin{proof}
Rappelons qu'un $A$-module est plat si et seulement s'il est sans torsion, voir
Compléments d'algèbre, Lemme
\ref{more-algebra-lemma-valuation-ring-torsion-free-flat}.
Soit $T \subset N$ la $A$-torsion. Alors $u(T) = 0$ et
$N/T$ est $A$-plat. Ainsi, $N/T$ est de présentation finie sur $B$, voir
Compléments d'algèbre, Lemme
\ref{more-algebra-lemma-flat-finite-type-valuation-ring-finite-presentation}.
Par conséquent, $T$ est un $B$-module fini, voir
Algèbre, Lemme \ref{algebra-lemma-extension}.
Comme $N/T$ est $A$-plat, nous voyons que
$T/\mathfrak m_A T \subset N/\mathfrak m_A N$, voir
Algèbre, Lemme \ref{algebra-lemma-flat-tor-zero}.
Puisque $\overline{u}$ est injective mais $u(T) = 0$, nous concluons que
$T/\mathfrak m_A T = 0$. Ainsi, $T = 0$ par le lemme de Nakayama, voir
Algèbre, Lemme \ref{algebra-lemma-NAK}. À ce stade, nous avons
démontré deux des trois assertions ($N$ est $A$-plat et
de présentation finie sur $B$) ; il reste à montrer que
$u$ est universellement injective.

\medskip\noindent
D'après Algèbre, Théorème \ref{algebra-theorem-universally-exact-criteria},
il suffit de montrer que $N \otimes_A Q \to M \otimes_A Q$ est injective
pour tout $A$-module de présentation finie $Q$. D'après
Compléments d'algèbre, Lemme
\ref{more-algebra-lemma-generalized-valuation-ring-modules},
nous pouvons supposer $Q = A/(a)$, avec $a \in \mathfrak m_A$ non nul.
Il suffit donc de montrer que $N/aN \to M/aM$ est injective.
Soit $x \in N$ tel que $u(x) \in aM$. D'après le
lemme \ref{flat-lemma-flat-finite-type-local-valuation-ring-has-content},
nous savons que $x$ possède un idéal de contenu $I \subset A$. Comme
$I$ est de type fini
(Compléments d'algèbre, Lemme \ref{more-algebra-lemma-content-finitely-generated})
et $A$ est un anneau de valuation, nous avons $I = (b)$ pour un certain $b$
(par Algèbre, Lemme \ref{algebra-lemma-characterize-valuation-ring}).
D'après Compléments d'algèbre, Lemme \ref{more-algebra-lemma-equal-content},
l'élément $u(x)$ a lui aussi pour idéal de contenu $I$.
Puisque $u(x) \in aM$, nous voyons que $(b) \subset (a)$
d'après Compléments d'algèbre, Définition \ref{more-algebra-definition-content-ideal}.
Comme $x \in bN$, nous concluons $x \in aN$, comme souhaité.
\end{proof}

\noindent
Considérons la situation suivante :
\begin{equation}
\label{flat-equation-star}
\begin{matrix}
A \to B\text{ de présentation finie, }S \subset B
\text{ une partie multiplicative, et }\\
N\text{ un }S^{-1}B\text{-module de présentation finie}
\end{matrix}
\end{equation}
Dans cette situation, un {\it prolongement pur} est la donnée d'un ouvert affine
$U \subset \Spec(B)$ tel que $\Spec(S^{-1}B) \subset U$ et d'un
$\mathcal{O}(U)$-module de présentation finie $N'$ prolongeant $N$
tel que $N'$ soit $A$-projectif et que $N' \to N = S^{-1}N'$
soit $A$-universellement injective.

\medskip\noindent
Dans (\ref{flat-equation-star}), si $A \to A_1$ est un homomorphisme d'anneaux, nous
pouvons effectuer un changement de base : prenons $B_1 = B \otimes_A A_1$, soit
$S_1 \subset B_1$ l'image de $S$, et soit $N_1 = N \otimes_A A_1$. Cette construction convient
car $S_1^{-1}B_1 = S^{-1}B \otimes_A A_1$. Nous l'utiliserons
sans autre mention dans la suite.

\begin{lemma}
\label{flat-lemma-properties-pure-spreadout}
Dans (\ref{flat-equation-star}), s'il existe un prolongement pur, alors
\begin{enumerate}
\item les éléments de $N$ possèdent des idéaux de contenu dans $A$ ;
\item si $u : N \to M$ est un morphisme vers un $A$-module plat $M$
tel que $N/\mathfrak m N \to M/\mathfrak m M$ soit injectif
pour tout idéal maximal $\mathfrak m$ de $A$, alors $u$ est
$A$-universellement injectif.
\end{enumerate}
\end{lemma}

\begin{proof}
Choisissons $U$, $N'$ comme dans la définition d'un prolongement pur.
Tout élément $x' \in N'$ possède un idéal de contenu dans $A$ parce que
$N'$ est $A$-projectif (cela se voit facilement directement, mais
résulte aussi de Compléments d'algèbre, Lemme
\ref{more-algebra-lemma-content-exists-flat-Mittag-Leffler}, et
Algèbre, Exemple \ref{algebra-example-ML}).
Puisque $N' \to N$ est $A$-universellement injective, nous voyons que
l'image $x \in N$ de tout $x' \in N'$ possède un idéal de contenu dans $A$
(c'est le même que l'idéal de contenu de $x'$).
Pour un élément général $x \in N$, nous choisissons $s \in S$ tel que
$s x$ appartienne à l'image de $N' \to N$ et utilisons le fait que $x$ et $sx$
ont le même idéal de contenu.

\medskip\noindent
Soit $u : N \to M$ comme en (2). Pour montrer que $u$ est $A$-universellement
injective, nous pouvons remplacer $A$ par une localisation en un idéal maximal
(nous omettons ce petit détail). Supposons $A$ local, d'idéal maximal $\mathfrak m$.
Choisissons $s \in S$ et considérons la composée
$$
N' \to N \xrightarrow{1/s} N \xrightarrow{u} M
$$
Chacune de ces applications est injective modulo $\mathfrak m$ ; la composée
est donc $A$-universellement injective d'après le
lemme \ref{flat-lemma-universally-injective-local}.
Puisque $N = \colim_{s \in S} (1/s)N'$, nous concluons que $u$
est $A$-universellement injective comme colimite d'applications universellement injectives.
\end{proof}

\begin{lemma}
\label{flat-lemma-find-pure-spreadout}
Dans (\ref{flat-equation-star}), pour tout $\mathfrak p \in \Spec(A)$,
il existe un idéal de type fini $I \subset \mathfrak pA_\mathfrak p$
tel que nous ayons un prolongement pur sur $A_\mathfrak p/I$.
\end{lemma}

\begin{proof}
Nous pouvons remplacer $A$ par $A_\mathfrak p$. Ainsi, nous pouvons supposer $A$
local et $\mathfrak p$ égal à l'idéal maximal $\mathfrak m$ de $A$.
Nous pouvons écrire $N = S^{-1}N'$ pour un $B$-module de présentation finie $N'$,
en chassant les dénominateurs dans une présentation de $N$ sur $S^{-1}B$.
Puisque $B/\mathfrak m B$ est noethérien, le noyau $K$ de
$N'/\mathfrak m N' \to N/\mathfrak m N$ est de type fini.
Nous pouvons donc choisir $s \in S$ tel que $K$ soit annulé par $s$.
Après avoir remplacé $B$ par $B_s$, ce qui est permis puisque cela revient simplement à passer
à un sous-schéma ouvert affine de $\Spec(B)$, nous constatons que les éléments de $S$
agissent injectivement sur $N'/\mathfrak m N'$. Nous choisissons alors
un sous-anneau local $A_0 \subset A$ essentiellement de type fini sur $\mathbf{Z}$,
un homomorphisme d'anneaux de type fini $A_0 \to B_0$ tel que $B = A \otimes_{A_0} B_0$,
et un $B_0$-module fini $N'_0$ tel que
$N' = B \otimes_{B_0} N'_0 = A \otimes_{A_0} N'_0$.
Nous affirmons que $I = \mathfrak m_{A_0} A$ convient.
En effet, nous avons
$$
N'/IN' = N'_0/\mathfrak m_{A_0} N'_0 \otimes_{\kappa_{A_0}} A/I
$$
qui est libre sur $A/I$. La multiplication par les éléments de $S$
est injective après passage au quotient par l'idéal maximal ; par conséquent,
$N'/IN' \to N/IN$ est universellement injective, par exemple d'après le
lemme \ref{flat-lemma-invert-universally-injective}.
\end{proof}

\begin{lemma}
\label{flat-lemma-universally-injective-if-flat}
Dans (\ref{flat-equation-star}), supposons que $N$ soit $A$-plat, que $M$ soit un $A$-module plat,
et que $u : N \to M$ soit un homomorphisme de $A$-modules tel que
$u \otimes \text{id}_{\kappa(\mathfrak p)}$ soit injectif pour tout
$\mathfrak p \in \Spec(A)$. Alors $u$ est $A$-universellement injectif.
\end{lemma}

\begin{proof}
D'après Algèbre, Lemme \ref{algebra-lemma-check-universally-injective-into-flat},
il suffit de vérifier que $N/IN \to M/IM$ est injectif pour tout
idéal $I \subset A$. Après avoir remplacé $A$ par $A/I$, nous voyons qu'il suffit
de démontrer que $u$ est injectif.

\medskip\noindent
Démonstration de l'injectivité de $u$. Soit $x \in N$ un élément non nul du
noyau de $u$. Il existe alors un idéal premier faiblement associé $\mathfrak p$
au module $Ax$ ; voir Algèbre, Lemme \ref{algebra-lemma-weakly-ass-zero}.
En remplaçant $A$ par $A_\mathfrak p$, nous pouvons supposer que $A$ est local et
nous trouvons un élément non nul $y \in Ax$ dont l'annulateur a pour radical
$\mathfrak m_A$ ; voir
Algèbre, Lemme \ref{algebra-lemma-weakly-ass-local}.
Ainsi, $\text{Supp}(y) \subset \Spec(S^{-1}B)$ est non vide et
contenu dans la fibre fermée de $\Spec(S^{-1}B) \to \Spec(A)$.
Soit $I \subset \mathfrak m_A$ un
idéal de type fini tel que nous ayons un prolongement pur sur $A/I$ ; voir le
lemme \ref{flat-lemma-find-pure-spreadout}. Alors $I^n y = 0$ pour un certain $n$. Or,
$y \in \text{Ann}_M(I^n) = \text{Ann}_A(I^n) \otimes_R N$ par platitude.
Ainsi, pour obtenir la contradiction souhaitée, il suffit de montrer que
$$
\text{Ann}_A(I^n) \otimes_R N
\longrightarrow
\text{Ann}_A(I^n) \otimes_R M
$$
est injectif. Puisque $N$ et $M$ sont plats et que $\text{Ann}_A(I^n)$
est annulé par $I^n$, il suffit de montrer que
$Q \otimes_A N \to Q \otimes_A M$ est injectif pour tout $A$-module
$Q$ annulé par $I$. Cela découle de notre choix de $I$ et du
lemme \ref{flat-lemma-properties-pure-spreadout}, partie (2).
\end{proof}

\begin{lemma}
\label{flat-lemma-big-intersection-is-zero}
Soit $A$ un anneau intègre local qui n'est pas un corps.
Soit $S$ un ensemble d'idéaux de type fini de $A$.
Supposons que $S$ soit stable par produits et que
$\bigcup_{I \in S} V(I)$ soit le complémentaire du point générique de $\Spec(A)$.
Alors $\bigcap_{I \in S} I = (0)$.
\end{lemma}

\begin{proof}
Puisque $\mathfrak m_A \subset A$ n'est pas le point générique de $\Spec(A)$,
nous voyons que $I \subset \mathfrak m_A$ pour au moins un $I \in S$.
Par conséquent, $\bigcap_{I \in S} I \subset \mathfrak m_A$.
Soit $f \in \mathfrak m_A$ non nul. Alors
$V(f) \subset \bigcup_{I \in S} V(I)$.
Comme la topologie constructible sur $V(f)$ est quasi-compacte
(Topologie, Lemme \ref{topology-lemma-constructible-hausdorff-quasi-compact}
et
Algèbre, Lemme \ref{algebra-lemma-spec-spectral}),
nous trouvons que $V(f) \subset V(I_1) \cup \ldots \cup V(I_n)$
pour certains $I_j \in S$. Comme $I_1 \ldots I_n \in S$, nous voyons que
$V(f) \subset V(I)$ pour un certain $I$. Puisque $I$ est de type fini,
cela implique que $I^m \subset (f)$ pour un certain $m$ et, puisque
$S$ est stable par produits, nous voyons que $I \subset (f^2)$ pour
un certain $I \in S$. Il est alors impossible que $f \in I$.
\end{proof}

\begin{lemma}
\label{flat-lemma-closed-points-complement}
Soit $A$ un anneau local. Soient $I, J \subset A$ des idéaux.
Si $J$ est de type fini et $I \subset J^n$ pour tout $n \geq 1$,
alors $V(I)$ contient les points fermés de $\Spec(A) \setminus V(J)$.
\end{lemma}

\begin{proof}
Soit $\mathfrak p \subset A$ un point fermé de $\Spec(A) \setminus V(J)$.
Nous voulons montrer que $I \subset \mathfrak p$. Sinon, un certain $f \in I$
s'envoie sur un élément non nul de $A/\mathfrak p$. Remarquons que
$V(J) \cap \Spec(A/\mathfrak p)$ est l'ensemble des points non génériques.
Le lemme \ref{flat-lemma-big-intersection-is-zero}, appliqué
à la famille d'idéaux $J^nA/\mathfrak p$, entraîne donc que
l'image de $f$ est nulle dans $A/\mathfrak p$.
\end{proof}

\begin{lemma}
\label{flat-lemma-make-smaller-flatness-ideal}
Soit $A$ un anneau local. Soit $I \subset A$ un idéal.
Soit $U \subset \Spec(A)$ un ouvert quasi-compact.
Soit $M$ un $A$-module. Supposons que
\begin{enumerate}
\item $M/IM$ soit plat sur $A/I$,
\item $M$ soit plat sur $U$.
\end{enumerate}
Alors $M/I_2M$ est plat sur $A/I_2$, où
$I_2 = \Ker(I \to \Gamma(U, I/I^2))$.
\end{lemma}

\begin{proof}
Il suffit de montrer que
$M \otimes_A I/I_2 \to IM/I_2M$ est injectif ; voir
Algèbre, Lemme \ref{algebra-lemma-what-does-it-mean-again}.
C'est vrai sur $U$ par l'hypothèse (2). Il suffit donc de montrer
que $M \otimes_A I/I_2$ s'injecte dans ses sections sur $U$.
Nous avons $M \otimes_A I/I_2 = M/IM \otimes_A I/I_2$ et
$M/IM$ est une colimite filtrante de $A/I$-modules libres finis
(Algèbre, Théorème \ref{algebra-theorem-lazard}).
Il suffit donc de montrer que $I/I_2$ s'injecte dans ses sections
sur $U$, ce qui résulte de la construction de $I_2$.
\end{proof}

\begin{proposition}
\label{flat-proposition-finite-type-injective-into-flat-mod-m}
Soit $A \to B$ un homomorphisme local d'anneaux locaux
qui est essentiellement de type fini. Soient $M$ un $A$-module plat,
$N$ un $B$-module fini et $u : N \to M$ un homomorphisme de $A$-modules tel que
$\overline{u} : N/\mathfrak m_AN \to M/\mathfrak m_AM$ soit injectif.
Alors $u$ est $A$-universellement injectif, $N$ est de présentation finie sur
$B$, et $N$ est plat sur $A$.
\end{proposition}

\begin{proof}
Nous pouvons supposer que $B$ est la localisation d'une
$A$-algèbre de présentation finie $B_0$ et que $N$ est la localisation d'un
$B_0$-module de présentation finie $M_0$ ; voir le
lemme \ref{flat-lemma-reduce-finite-type-injective-into-flat-mod-m}.
D'après Compléments sur les morphismes, Lemme
\ref{more-morphisms-lemma-generic-flatness-stratification},
il existe une « stratification de platitude générique »
pour $\widetilde{M_0}$ sur $\Spec(B_0)$ au-dessus de $\Spec(A)$.
En revenant à $N$, nous obtenons une suite de sous-schémas fermés
$$
S = \Spec(A) \supset S_0 \supset S_1 \supset \ldots \supset S_t = \emptyset
$$
où $S_i \subset S$ est défini par un idéal de type fini de $A$
et telle que l'image inverse de $\widetilde{N}$ sur
$\Spec(B) \times_S (S_i \setminus S_{i + 1})$ soit plate sur
$S_i \setminus S_{i + 1}$. Nous allons démontrer la proposition par
récurrence sur $t$ (le cas initial $t = 1$ sera démontré en même temps
que les autres étapes). Soit $\Spec(A/J_i)$ l'adhérence schématique
de $S_i \setminus S_{i + 1}$.

\medskip\noindent
{\bf Assertion 1.} $N/J_iN$ est plat sur $A/J_i$. C'est immédiat pour
$i = t - 1$ et cela résulte de l'hypothèse de récurrence pour $i > 0$.
Nous pouvons donc supposer $t > 1$, $S_{t - 1} \not = \emptyset$, et
$J_0 = 0$, et nous devons démontrer que $N$ est plat. Soit $J \subset A$
l'idéal définissant $S_1$. En raisonnant de nouveau par récurrence sur $t$, nous avons aussi
la platitude modulo les puissances de $J$. Soit $A^h$ l'hensélisé de $A$
et soit $B'$ la localisation de $B \otimes_A A^h$ en l'idéal maximal
$\mathfrak m_B \otimes A^h + B \otimes \mathfrak m_{A^h}$. Alors $B \to B'$
est fidèlement plat. Posons $N' = N \otimes_B B'$. Remarquons que $N'$
est $A^h$-plat si et seulement si $N$ est $A$-plat. D'après le
théorème \ref{flat-theorem-flattening-local}, il existe un plus petit idéal
$I \subset A^h$ tel que $N'/IN'$ soit plat sur $A^h/I$, et
$I$ est de type fini. D'après ce qui précède, $I \subset J^nA^h$ pour
tout $n \geq 1$. Soit $S_i^h \subset \Spec(A^h)$ l'image inverse
de $S_i \subset \Spec(A)$. D'après le lemme \ref{flat-lemma-closed-points-complement},
nous voyons que $V(I)$ contient les points fermés de $U = \Spec(A^h) - S_1^h$.
Par construction, $N'$ est $A^h$-plat sur $U$.
D'après le lemme \ref{flat-lemma-make-smaller-flatness-ideal}, nous voyons que $N'/I_2N'$
est plat sur $A/I_2$, où
$I_2 = \Ker(I \to \Gamma(U, I/I^2))$. Ainsi $I = I_2$ par minimalité
de $I$. Cela implique que $I = I^2$ localement sur $U$, c'est-à-dire
que nous avons $I\mathcal{O}_{U, u} = (0)$ ou $I\mathcal{O}_{U, u} = (1)$
pour tout $u \in U$. Comme $V(I)$ contient les points fermés de $U$,
nous voyons que $I = 0$ sur $U$. Puisque $U \subset \Spec(A^h)$ est schématiquement
dense (parce que nous avons remplacé $A$ par $A/J_0$ au début
de ce paragraphe), nous voyons que $I = 0$. Ainsi $N'$ est $A^h$-plat
et, par conséquent, l'assertion 1 est vraie.

\medskip\noindent
Revenons à la situation décrite avant l'assertion 1. Avec
$A^h$ l'hensélisé de $A$, avec $B'$ la localisation
de $B \otimes_A A^h$ en l'idéal maximal
$\mathfrak m_B \otimes A^h + B \otimes \mathfrak m_{A^h}$, et avec
$N' = N \otimes_B B'$, nous voyons maintenant que l'idéal de platification $I \subset A^h$
du théorème \ref{flat-theorem-flattening-local} est nilpotent.
Si $nil(A^h)$ désigne l'idéal des éléments nilpotents, alors
$nil(A^h) = nil(A) A^h$
(Compléments d'algèbre, Lemme \ref{more-algebra-lemma-henselization-nil}).
Il existe donc un idéal de type fini
nilpotent $I_0 \subset A$ tel que $N/I_0N$ soit plat sur $A/I_0$.

\medskip\noindent
{\bf Assertion 2.} Pour tout idéal premier $\mathfrak p \subset A$,
l'application
$\kappa(\mathfrak p) \otimes_A N \to \kappa(\mathfrak p) \otimes_A M$
est injective. Nous disons que $\mathfrak p$ est mauvais si cela est faux.
Supposons que $C$ soit une chaîne non vide de premiers mauvais et posons
$\mathfrak p^* = \bigcup_{\mathfrak p \in C} \mathfrak p$.
D'après le lemme \ref{flat-lemma-find-pure-spreadout},
il existe un idéal de type fini
$\mathfrak a \subset \mathfrak p^*A_{\mathfrak p^*}$
tel qu'il existe un prolongement pur sur $V(\mathfrak a)$.
Si $\mathfrak p^*$ était bon, il résulterait du
lemme \ref{flat-lemma-properties-pure-spreadout}
que les points de $V(\mathfrak a)$ sont bons.
Cependant, comme $\mathfrak a$ est de type fini et que
$\mathfrak p^*A_{\mathfrak p^*} = \bigcup_{\mathfrak p \in C}A_{\mathfrak p^*}$,
nous voyons que $V(\mathfrak a)$ contient un $\mathfrak p \in C$, contradiction.
Ainsi $\mathfrak p^*$ est mauvais. D'après le lemme de Zorn, s'il existe un
premier mauvais, il en existe un maximal, disons $\mathfrak p$.
Autrement dit, nous pouvons supposer que tout $\mathfrak p' \supset \mathfrak p$,
$\mathfrak p' \not = \mathfrak p$, est bon.
Dans ce cas, nous voyons que, pour tout $f \in A$, $f \not \in \mathfrak p$,
l'application $u \otimes \text{id}_{A/(\mathfrak p + f)}$ est universellement
injective ; voir le lemme \ref{flat-lemma-universally-injective-if-flat}.
Il suffit donc de montrer que $N/\mathfrak p N$ est séparé
pour la topologie définie par les sous-modules $f(N/\mathfrak pN)$.
Puisque $B \to B'$ est fidèlement plat, il suffit de démontrer le
même énoncé pour le module $N'/\mathfrak p N'$.
D'après le lemme \ref{flat-lemma-flat-finite-type-local-colimit-free} et
Compléments d'algèbre, Lemme
\ref{more-algebra-lemma-content-exists-flat-Mittag-Leffler},
les éléments de $N'/\mathfrak pN'$ possèdent des idéaux de contenu dans $A^h/\mathfrak pA^h$.
Il suffit donc de montrer que
$\bigcap_{f \in A, f \not \in \mathfrak p} f(A^h/\mathfrak p A^h) = 0$.
Il suffit ensuite de montrer la même chose pour
$A^h/\mathfrak q A^h$ pour tout idéal premier $\mathfrak q \subset A^h$
minimal au-dessus de $\mathfrak p A^h$. Comme $A \to A^h$ est l'hensélisation,
tout $\mathfrak q$ se contracte en $\mathfrak p$ et tout
$\mathfrak q' \supset \mathfrak q$, $\mathfrak q' \not = \mathfrak q$,
se contracte en un idéal premier $\mathfrak p'$ contenant strictement $\mathfrak p$.
L'annulation des intersections découle donc du
lemme \ref{flat-lemma-big-intersection-is-zero}.

\medskip\noindent
Nous pouvons maintenant tout rassembler. En effet, en utilisant
les assertions 1 et 2, nous voyons que $N/I_0 N \to M/I_0M$ est
$A/I_0$-universellement injectif d'après le
lemme \ref{flat-lemma-universally-injective-if-flat}.
Alors les diagrammes
$$
\xymatrix{
N \otimes_A (I_0^n/I_0^{n + 1}) \ar[r] \ar[d] &
M \otimes_A (I_0^n/I_0^{n + 1}) \ar@{=}[d] \\
I_0^n N /I_0^{n + 1} N \ar[r] &
I_0^n M /I_0^{n + 1} M
}
$$
montrent que les flèches verticales de gauche sont injectives. Ainsi, d'après
Algèbre, Lemme \ref{algebra-lemma-what-does-it-mean-again},
nous voyons que $N$ est plat. De manière analogue, l'injectivité
universelle de $u$ peut être ramenée (même sans
démontrer d'abord la platitude de $N$) à celle modulo $I_0$. Ceci achève
la démonstration.
\end{proof}


\section{Modules plats de type fini, III}
\label{flat-section-finite-type-flat-III}

\noindent
Le résultat suivant est l'un des résultats principaux de ce chapitre.

\begin{theorem}
\label{flat-theorem-check-flatness-at-associated-points}
Soit $f : X \to S$ localement de type fini.
Soit $\mathcal{F}$ un $\mathcal{O}_X$-module quasi-cohérent de type fini.
Soit $x \in X$ d'image $s \in S$.
Les conditions suivantes sont équivalentes :
\begin{enumerate}
\item $\mathcal{F}$ est plat en $x$ sur $S$ ;
\item pour tout $x' \in \text{Ass}_{X_s}(\mathcal{F}_s)$ qui
se spécialise en $x$, le module $\mathcal{F}$ est plat en $x'$ sur $S$.
\end{enumerate}
\end{theorem}

\begin{proof}
Il est clair que (1) implique (2), puisque $\mathcal{F}_{x'}$ est une localisation
de $\mathcal{F}_x$ pour tout point qui se spécialise en $x$.
Posons $A = \mathcal{O}_{S, s}$, $B = \mathcal{O}_{X, x}$ et
$N = \mathcal{F}_x$. Soit $\Sigma \subset B$ la partie multiplicative
de $B$ formée des éléments qui agissent comme non-diviseurs de zéro sur $N/\mathfrak m_AN$.
L'hypothèse (2) implique que $\Sigma^{-1}N$ est $A$-plat d'après la description
de $\Spec(\Sigma^{-1}N)$ donnée dans le
lemme \ref{flat-lemma-homothety-spectrum}.
D'autre part, l'application $N \to \Sigma^{-1}N$ est injective modulo
$\mathfrak m_A$ par construction. En appliquant donc le
lemme \ref{flat-lemma-upstairs-finite-type-injective-into-flat-mod-m},
nous concluons.
\end{proof}

\noindent
Nous appliquons maintenant directement ce résultat afin d'obtenir les résultats utiles suivants.

\begin{lemma}
\label{flat-lemma-check-along-closed-fibre}
Soit $S$ un schéma local de point fermé $s$.
Soit $f : X \to S$ localement de type fini.
Soit $\mathcal{F}$ un $\mathcal{O}_X$-module quasi-cohérent de type fini.
Supposons que
\begin{enumerate}
\item tout point de $\text{Ass}_{X/S}(\mathcal{F})$ se spécialise
en un point de la fibre fermée $X_s$\footnote{Cela vaut par exemple si
$f$ est de type fini et $\mathcal{F}$ est pur le long de $X_s$, ou
si $f$ est propre.},
\item $\mathcal{F}$ est plat sur $S$ en tout point de $X_s$.
\end{enumerate}
Alors $\mathcal{F}$ est plat sur $S$.
\end{lemma}

\begin{proof}
Cela résulte immédiatement du fait qu'il suffit de vérifier la
platitude aux points de l'assassin relatif de $\mathcal{F}$
sur $S$, d'après le
théorème \ref{flat-theorem-check-flatness-at-associated-points}.
\end{proof}











\section{Platificateur universel}
\label{flat-section-flattening-final}

\noindent
Si $f : X \to S$ est un morphisme propre de présentation finie
de schémas, alors on peut trouver un platificateur universel de $f$.
Dans cette section, nous étudions ce résultat ainsi que certaines de ses variantes.

\begin{lemma}
\label{flat-lemma-free-at-generic-points-representable}
Dans la situation \ref{flat-situation-free-at-generic-points}.
Pour chaque $p \geq 0$, le foncteur $H_p$
(\ref{flat-equation-free-at-generic-points}) est représentable
par une immersion localement fermée $S_p \to S$. Si $\mathcal{F}$
est de présentation finie, alors $S_p \to S$ est de présentation finie.
\end{lemma}

\begin{proof}
Pour chaque $S$, nous allons démontrer simultanément l'énoncé pour tout $p \geq 0$.
Le foncteur $H_p$ est un faisceau pour la topologie fppf d'après le
lemme \ref{flat-lemma-free-at-generic-points}.
En combinant donc
Descente, Lemme \ref{descent-lemma-descent-data-sheaves},
Compléments sur les morphismes, Lemme
\ref{more-morphisms-lemma-separated-locally-quasi-finite-morphisms-fppf-descend}
, et
Descente, Lemme \ref{descent-lemma-descending-fppf-property-immersion},
nous voyons que la question est locale pour la topologie \'etale sur $S$.
En particulier, la question est locale pour la topologie de Zariski sur $S$.

\medskip\noindent
Pour $s \in S$, notons $\xi_s$ l'unique point générique de la fibre $X_s$.
Remarquons que, pour tout $s \in S$, la restriction $\mathcal{F}_s$ de
$\mathcal{F}$ est localement libre d'un certain rang $p(s) \geq 0$ dans un
voisinage de $\xi_s$. (Comme $X_s$ est irréductible
et lisse, cela résulte de la platitude générique de $\mathcal{F}_s$ sur
$X_s$ ; voir
Algèbre, Lemme \ref{algebra-lemma-generic-flatness-Noetherian},
bien que ce soit excessif.) Pour référence ultérieure, remarquons que
$$
p(s) =
\dim_{\kappa(\xi_s)}(
\mathcal{F}_{\xi_s} \otimes_{\mathcal{O}_{X, \xi_s}} \kappa(\xi_s)
).
$$
En particulier, $H_{p(s)}(s)$ est non vide et $H_q(s)$ est vide
si $q \not = p(s)$.

\medskip\noindent
Soit $U \subset X$ un sous-schéma ouvert.
Comme $f : X \to S$ est lisse, il est ouvert.
Il résulte immédiatement de (\ref{flat-equation-free-at-generic-points})
que le foncteur $H_p$ pour la paire $(f|_U : U \to f(U), \mathcal{F}|_U)$
et le foncteur $H_p$ pour la paire
$(f|_{f^{-1}(f(U))}, \mathcal{F}|_{f^{-1}(f(U))})$
sont les mêmes. Ainsi, pour démontrer l'existence de $S_p$ sur $f(U)$, nous pouvons
toujours remplacer $X$ par $U$.

\medskip\noindent
Choisissons $s \in S$. Il existe un voisinage ouvert affine $U$
de $\xi_s$ tel que $\mathcal{F}|_U$ puisse être engendré par au plus
$p(s)$ éléments. D'après les arguments précédents, pour démontrer
l'énoncé concernant $H_{p(s)}$ dans un voisinage de $s$, nous pouvons supposer
que $\mathcal{F}$ est engendré par $p(s)$ éléments, c'est-à-dire qu'il existe
une surjection
$$
u : \mathcal{O}_X^{\oplus p(s)} \longrightarrow \mathcal{F}
$$
Dans ce cas, il est clair que $H_{p(s)}$ est égal à $F_{iso}$
(\ref{flat-equation-iso}) pour l'application $u$ (cela résulte immédiatement du
lemme \ref{flat-lemma-injectivity-map-source-flat-pure},
mais aussi du
lemme \ref{flat-lemma-induction-step-fp},
après avoir encore un peu rétréci de sorte que $S$ et $X$ soient tous deux affines.)
Nous pouvons donc appliquer le
théorème \ref{flat-theorem-flattening-map}
pour voir que $H_{p(s)}$ est représentable par une immersion fermée dans un
voisinage de $s$.

\medskip\noindent
Le résultat découle formellement de ce qui précède.
En effet, les arguments ci-dessus montrent que, localement sur $S$, la fonction
$s \mapsto p(s)$ est bornée. Nous pouvons donc raisonner par récurrence
sur $p = \max_{s \in S} p(s)$. Le foncteur $H_p$ est représentable
par une immersion fermée $S_p \to S$ d'après ce qui précède. Remplaçons $S$ par
$S \setminus S_p$, ce qui diminue le maximum d'au moins un, et
nous concluons par l'hypothèse de récurrence.

\medskip\noindent
Supposons que $\mathcal{F}$ soit de présentation finie.
Alors $S_p \to S$ est localement de présentation finie d'après le
lemme \ref{flat-lemma-free-at-generic-points}, partie (2), combiné à
Limites, Remarque \ref{limits-remark-limit-preserving}.
Nous reprenons alors l'argument de récurrence du paragraphe précédent pour
voir que chaque $S_p$ est quasi-compact lorsque $S$ est affine :
d'abord, si $p = \max_{s \in S} p(s)$, alors $S_p \subset S$
est fermé (voir ci-dessus), donc quasi-compact. Ensuite, $U = S \setminus S_p$
est un ouvert quasi-compact de $S$, car $S_p \to S$ est une immersion
fermée de présentation finie (voir par exemple la discussion dans Morphismes, section
\ref{morphisms-section-constructible}). Puis $S_{p - 1} \to U$
est une immersion fermée de présentation finie, de sorte que $S_{p - 1}$
est quasi-compact et $U' = S \setminus (S_p \cup S_{p - 1})$
est quasi-compact. Et ainsi de suite.
\end{proof}

\begin{lemma}
\label{flat-lemma-localize-flat-dimension-n}
Dans la situation \ref{flat-situation-flat-dimension-n}.
Soit $h : X' \to X$ un morphisme \'etale.
Posons $\mathcal{F}' = h^*\mathcal{F}$ et $f' = f \circ h$.
Soit $F_n'$ le foncteur (\ref{flat-equation-flat-dimension-n})
associé à $(f' : X' \to S, \mathcal{F}')$.
Alors $F_n$ est un sous-foncteur de $F_n'$ et, si
$h(X') \supset \text{Ass}_{X/S}(\mathcal{F})$, alors $F_n = F'_n$.
\end{lemma}

\begin{proof}
Soit $T \to S$ un morphisme quelconque. Alors $h_T : X'_T \to X_T$ est \'etale comme
changement de base du morphisme \'etale $g$. Pour $t \in T$, notons
$Z \subset X_t$ l'ensemble des points où $\mathcal{F}_T$ n'est pas
plat sur $T$, et notons de même $Z' \subset X'_t$ l'ensemble des
points où $\mathcal{F}'_T$ n'est pas plat sur $T$. Comme
$\mathcal{F}'_T = h_T^*\mathcal{F}_T$, nous voyons que
$Z' = h_t^{-1}(Z)$ ; voir
Morphismes, Lemme \ref{morphisms-lemma-flat-permanence}.
Ainsi, $Z' \to Z$ est un morphisme \'etale, donc $\dim(Z') \leq \dim(Z)$
(par exemple d'après
Descente, Lemme \ref{descent-lemma-dimension-at-point-local},
ou simplement parce qu'un morphisme \'etale est lisse de dimension relative $0$).
Cela implique que $F_n \subset F_n'$.

\medskip\noindent
Enfin, supposons que $h(X') \supset \text{Ass}_{X/S}(\mathcal{F})$
et que $T \to S$ soit un morphisme tel que $F_n'(T)$ soit non vide, c'est-à-dire
tel que $\mathcal{F}'_T$ soit plat en dimensions $\geq n$ sur $T$.
Choisissons un point $t \in T$ et soient $Z \subset X_t$ et $Z' \subset X'_t$
comme ci-dessus. Pour obtenir une contradiction, supposons que $\dim(Z) \geq n$.
Choisissons un point générique $\xi \in Z$ correspondant à une composante
de dimension $\geq n$. Soit $x \in \text{Ass}_{X_t}(\mathcal{F}_t)$
une généralisation de $\xi$. Alors $x$ s'envoie sur un point de
$\text{Ass}_{X/S}(\mathcal{F})$ d'après
Diviseurs, Lemme \ref{divisors-lemma-base-change-relative-assassin} et
Remarque \ref{divisors-remark-base-change-relative-assassin}.
Nous voyons donc que $x$ appartient à l'image de $h_T$, disons
$x = h_T(x')$ pour un certain $x' \in X'_T$. Mais $x' \not \in Z'$
puisque $x \leadsto \xi$ et $\dim(Z') < n$. Ainsi $\mathcal{F}'_T$
est plat sur $T$ en $x'$, ce qui implique que $\mathcal{F}_T$ est plat
en $x$ sur $T$ (d'après
Morphismes, Lemme \ref{morphisms-lemma-flat-permanence}).
Comme cela vaut pour tout $x$ de cette forme, nous concluons
que $\mathcal{F}_T$ est plat sur $T$ en $\xi$ d'après le
théorème \ref{flat-theorem-check-flatness-at-associated-points},
ce qui est la contradiction souhaitée.
\end{proof}

\begin{lemma}
\label{flat-lemma-compare-H-F}
Supposons que $X \to S$ soit un morphisme lisse de schémas affines
à fibres géométriquement irréductibles de dimension $d$ et que
$\mathcal{F}$ soit un $\mathcal{O}_X$-module quasi-cohérent de
présentation finie. Alors $F_d = \coprod_{p = 0, \ldots, c} H_p$
pour un certain $c \geq 0$, où $F_d$ est comme dans
(\ref{flat-equation-flat-dimension-n}) et $H_p$ comme dans
(\ref{flat-equation-free-at-generic-points}).
\end{lemma}

\begin{proof}
Comme $X$ est affine et $\mathcal{F}$ est quasi-cohérent de présentation finie,
nous savons que $\mathcal{F}$ peut être engendré par $c \geq 0$ éléments.
Alors $\dim_{\kappa(x)}(\mathcal{F}_x \otimes \kappa(x))$
en tout point $x \in X$ ne dépasse jamais $c$. En particulier, $H_p = \emptyset$
pour $p > c$. De plus, remarquons qu'il existe certainement une inclusion
$\coprod H_p \to F_d$. Cela étant dit, le contenu
du lemme est que, si un changement de base $\mathcal{F}_T$ est plat en
dimensions $\geq d$ sur $T$ et si $t \in T$, alors $\mathcal{F}_T$ est
libre d'un certain rang $r$ dans un voisinage ouvert $U \subset X_T$
de l'unique point générique $\xi$ de $X_t$. En effet, $H_r$
contient alors l'image de $U$, qui est un voisinage ouvert de $t$.
L'existence de $U$ résulte de
Compléments sur les morphismes, Lemme
\ref{more-morphisms-lemma-flat-and-free-at-point-fibre}.
\end{proof}

\begin{lemma}
\label{flat-lemma-flat-dimension-n-representable}
Dans la situation \ref{flat-situation-flat-dimension-n}.
Soient $s \in S$ et $d \geq 0$. Supposons que
\begin{enumerate}
\item il existe un d\'evissage complet
de $\mathcal{F}/X/S$ au-dessus d'un certain point $s \in S$,
\item $X$ soit de présentation finie sur $S$,
\item $\mathcal{F}$ soit un $\mathcal{O}_X$-module de présentation finie, et
\item $\mathcal{F}$ soit plat en dimensions $\geq d + 1$ sur $S$.
\end{enumerate}
Alors, après avoir éventuellement remplacé $S$ par un voisinage ouvert
de $s$, le foncteur $F_d$ (\ref{flat-equation-flat-dimension-n})
est représentable par un monomorphisme $Z_d \to S$ de présentation finie.
\end{lemma}

\begin{proof}
Une remarque préliminaire est que $X$ et $S$ sont des schémas affines et qu'il
suffit de démontrer que $F_d$ est représentable par un monomorphisme de présentation
finie $Z_d \to S$ dans la catégorie des schémas affines sur $S$.
(Bien entendu, nous n'exigeons pas que $Z_d$ soit affine.)
Ainsi, tout au long de la démonstration du
lemme, nous travaillons dans la catégorie des schémas affines sur $S$.

\medskip\noindent
Soit $(Z_k, Y_k, i_k, \pi_k, \mathcal{G}_k, \alpha_k)_{k = 1, \ldots, n}$
un d\'evissage complet de $\mathcal{F}/X/S$ au-dessus de $s$ ; voir la
définition \ref{flat-definition-complete-devissage}.
Nous allons raisonner par récurrence sur la longueur $n$ du d\'evissage.
Rappelons que $Y_k \to S$ est lisse à fibres géométriquement irréductibles ; voir la
définition \ref{flat-definition-one-step-devissage}.
Soit $d_k$ la dimension relative de $Y_k$ sur $S$.
Rappelons que $i_{k, *}\mathcal{G}_k = \Coker(\alpha_k)$ et
que $i_k$ est une immersion fermée.
D'après les définitions citées ci-dessus, nous avons
$d_1 = \dim(\text{Supp}(\mathcal{F}_s))$ et
$$
d_k = \dim(\text{Supp}(\Coker(\alpha_{k - 1})_s))
= \dim(\text{Supp}(\mathcal{G}_{k, s}))
$$
pour $k = 2, \ldots, n$. Il s'ensuit que $d_1 > d_2 > \ldots > d_n \geq 0$
parce que $\alpha_k$ est un isomorphisme au point générique de $(Y_k)_s$.

\medskip\noindent
Remarquons que $i_1$ est une immersion fermée et que
$\mathcal{F} = i_{1, *}\mathcal{G}_1$.
Ainsi, pour tout morphisme de schémas $T \to S$ avec $T$ affine,
nous avons $\mathcal{F}_T = i_{1, T, *}\mathcal{G}_{1, T}$ et
$i_{1, T}$ est encore une immersion fermée de schémas sur $T$.
Ainsi, $\mathcal{F}_T$ est plat en dimensions $\geq d$ sur $T$
si et seulement si $\mathcal{G}_{1, T}$ est plat en dimensions $\geq d$ sur $T$.
Comme $\pi_1 : Z_1 \to Y_1$ est fini, nous voyons de même que
$\mathcal{G}_{1, T}$ est plat en dimensions $\geq d$ sur $T$
si et seulement si $\pi_{1, T, *}\mathcal{G}_{1, T}$ est plat en dimensions
$\geq d$ sur $T$. Les mêmes arguments s'appliquent à
« plat en dimensions $\geq d + 1$ » et nous concluons en particulier que
$\pi_{1, *}\mathcal{G}_1$ est plat sur $S$ en dimensions $\geq d + 1$
d'après notre hypothèse sur $\mathcal{F}$.

\medskip\noindent
Supposons que $d_1 > d$. Il résulte de la discussion précédente qu'en
particulier $\pi_{1, *}\mathcal{G}_1$ est plat sur $S$ au
point générique de $(Y_1)_s$. D'après le
lemme \ref{flat-lemma-induction-step-fp},
nous pouvons remplacer $S$ par un voisinage affine de $s$ et supposer que
$\alpha_1$ est $S$-universellement injectif. Puisque $\alpha_1$ est
$S$-universellement injectif, pour tout morphisme $T \to S$ avec $T$ affine,
nous avons une suite exacte courte
$$
0 \to \mathcal{O}_{Y_{1, T}}^{\oplus r_1}
\to \pi_{1, T, *}\mathcal{G}_{1, T} \to \Coker(\alpha_1)_T \to 0
$$
et la première flèche reste $T$-universellement injective. Ainsi, l'ensemble
des points de $(Y_1)_T$ où $\pi_{1, T, *}\mathcal{G}_{1, T}$ est plat sur
$T$ est le même que l'ensemble des points de $(Y_1)_T$ où
$\Coker(\alpha_1)_T$ est plat sur $S$. De cette manière, la question
se ramène au faisceau $\Coker(\alpha_1)$, qui possède un d\'evissage
complet de longueur $n - 1$, et nous concluons par récurrence.

\medskip\noindent
Si $d_1 < d$, alors $F_d$ est représenté par $S$ et nous concluons.

\medskip\noindent
Le dernier cas est celui où $d_1 = d$. Ce cas résulte d'une combinaison du
lemme \ref{flat-lemma-compare-H-F}
et du
lemme \ref{flat-lemma-free-at-generic-points-representable}.
\end{proof}

\begin{theorem}
\label{flat-theorem-flat-dimension-n-representable}
Dans la situation \ref{flat-situation-flat-dimension-n}.
Supposons en outre que $f$ soit de présentation finie, que
$\mathcal{F}$ soit un $\mathcal{O}_X$-module de présentation finie,
et que $\mathcal{F}$ soit pur relativement à $S$.
Alors $F_n$ est représentable par un monomorphisme
$Z_n \to S$ de présentation finie.
\end{theorem}

\begin{proof}
Le foncteur $F_n$ est un faisceau pour la topologie fppf d'après le
lemme \ref{flat-lemma-flat-dimension-n}.
Observons qu'un monomorphisme de présentation finie est
séparé et quasi-fini (Morphismes, Lemme
\ref{morphisms-lemma-monomorphism-loc-finite-type-loc-quasi-finite}).
En combinant donc
Descente, Lemme \ref{descent-lemma-descent-data-sheaves},
Compléments sur les morphismes, Lemme
\ref{more-morphisms-lemma-separated-locally-quasi-finite-morphisms-fppf-descend}
, et
Descente, Lemmes \ref{descent-lemma-descending-property-monomorphism} et
\ref{descent-lemma-descending-property-finite-presentation},
nous voyons que la question est locale pour la topologie \'etale sur $S$.

\medskip\noindent
En particulier, la situation est locale pour la topologie de Zariski sur $S$
et nous pouvons supposer que $S$ est affine. Dans ce cas, la dimension des
fibres de $f$ est majorée ; nous voyons donc que $F_n$ est représentable
pour $n$ assez grand. Nous pouvons ainsi raisonner par récurrence descendante sur $n$.
Supposons que nous sachions que $F_{n + 1}$ est représentable par un monomorphisme
$Z_{n + 1} \to S$ de présentation finie. Considérons le changement de base
$X_{n + 1} = Z_{n + 1} \times_S X$ et l'image inverse $\mathcal{F}_{n + 1}$
de $\mathcal{F}$ sur $X_{n + 1}$. Le morphisme $Z_{n + 1} \to S$ est
quasi-fini, puisqu'il s'agit d'un monomorphisme de présentation finie ; le
lemme \ref{flat-lemma-quasi-finite-base-change}
implique donc que $\mathcal{F}_{n + 1}$ est pur relativement à $Z_{n + 1}$.
Puisque $F_n$ est un sous-foncteur de $F_{n + 1}$, nous concluons que, pour
démontrer le résultat pour $F_n$, il suffit de le démontrer pour le
foncteur correspondant à la situation
$\mathcal{F}_{n + 1}/X_{n + 1}/Z_{n + 1}$.
Nous nous ramenons ainsi à démontrer le résultat pour $F_n$ dans le cas où
$S_{n + 1} = S$, c'est-à-dire que nous pouvons supposer que $\mathcal{F}$ est plat
en dimensions $\geq n + 1$ sur $S$.

\medskip\noindent
Fixons $n$ et supposons que $\mathcal{F}$ soit plat en dimensions $\geq n + 1$
sur $S$. Pour achever la démonstration, nous devons montrer que $F_n$ est représentable
par un monomorphisme $Z_n \to S$ de présentation finie.
Puisque la question est locale pour la topologie \'etale sur $S$, il suffit de
montrer que, pour tout $s \in S$, il existe un voisinage \'etale élémentaire
$(S', s') \to (S, s)$ tel que le résultat soit vrai après changement de base à $S'$.
Ainsi, d'après le
lemme \ref{flat-lemma-existence-complete},
nous pouvons supposer qu'il existe des morphismes \'etales $h_j : Y_j \to X$,
$j = 1, \ldots, m$, tels que, pour chaque $j$, il existe un d\'evissage complet
de $\mathcal{F}_j/Y_j/S$ au-dessus de $s$,
où $\mathcal{F}_j$ est l'image inverse de $\mathcal{F}$ sur $Y_j$,
et tels que $X_s \subset \bigcup h_j(Y_j)$. Remarquons que, d'après le
lemme \ref{flat-lemma-localize-flat-dimension-n},
les faisceaux $\mathcal{F}_j$ sont encore plats
en dimensions $\geq n + 1$ sur $S$.
Posons $W = \bigcup h_j(Y_j)$, qui est un ouvert quasi-compact de $X$.
Comme $\mathcal{F}$ est pur le long de $X_s$, nous voyons que
$$
E = \{t \in S \mid \text{Ass}_{X_t}(\mathcal{F}_t) \subset W \}.
$$
contient toutes les généralisations de $s$. D'après
Compléments sur les morphismes,
Lemme \ref{more-morphisms-lemma-relative-assassin-constructible},
$E$ est une partie constructible de $S$. Nous avons vu que
$\Spec(\mathcal{O}_{S, s}) \subset E$. D'après
Morphismes, Lemme \ref{morphisms-lemma-constructible-containing-open},
nous voyons que $E$ contient un voisinage ouvert de $s$.
Ainsi, après avoir rétréci $S$, nous pouvons supposer que $E = S$.
Il résulte du
lemme \ref{flat-lemma-localize-flat-dimension-n}
qu'il suffit de démontrer le lemme pour le foncteur $F_n$ associé à
$X = \coprod Y_j$ et $\mathcal{F} = \coprod \mathcal{F}_j$.
Si $F_{j, n}$ désigne le foncteur pour $Y_j \to S$ et le faisceau
$\mathcal{F}_i$, nous voyons que $F_n = \prod F_{j, n}$. Il suffit donc
de démontrer que chaque $F_{j, n}$ est représentable par un monomorphisme
$Z_{j, n} \to S$ de présentation finie, puisque alors
$$
Z_n = Z_{1, n} \times_S \ldots \times_S Z_{m, n}
$$
Nous avons ainsi ramené le théorème au cas particulier traité dans le
lemme \ref{flat-lemma-flat-dimension-n-representable}.
\end{proof}

\noindent
Nous explicitons dans le lemme suivant la signification du théorème en termes
de platificateurs universels.

\begin{lemma}
\label{flat-lemma-when-universal-flattening}
Soit $f : X \to S$ un morphisme de schémas.
Soit $\mathcal{F}$ un $\mathcal{O}_X$-module quasi-cohérent.
\begin{enumerate}
\item Si $f$ est de présentation finie, si $\mathcal{F}$ est un
$\mathcal{O}_X$-module de présentation finie et si $\mathcal{F}$ est
pur relativement à $S$, alors il existe un platificateur universel
$S' \to S$ de $\mathcal{F}$. De plus, $S' \to S$ est un monomorphisme
de présentation finie.
\item Si $f$ est de présentation finie et $X$ est pur relativement à $S$,
alors il existe un platificateur universel $S' \to S$ de $X$.
De plus, $S' \to S$ est un monomorphisme de présentation finie.
\item Si $f$ est propre et de présentation finie et si $\mathcal{F}$ est un
$\mathcal{O}_X$-module de présentation finie, alors il existe un
platificateur universel $S' \to S$ de $\mathcal{F}$. De plus, $S' \to S$ est
un monomorphisme de présentation finie.
\item Si $f$ est propre et de présentation finie,
alors il existe un platificateur universel $S' \to S$ de $X$.
\end{enumerate}
\end{lemma}

\begin{proof}
Ces énoncés résultent immédiatement du
théorème \ref{flat-theorem-flat-dimension-n-representable}
appliqué à $F_0 = F_{flat}$,
et du fait que, si $f$ est propre, alors $\mathcal{F}$ est automatiquement
pur sur la base ; voir le
lemme \ref{flat-lemma-proper-pure}.
\end{proof}







\section{Théorème d'existence de Grothendieck, IV}
\label{flat-section-existence}

\noindent
Cette section poursuit l'étude entreprise dans
Cohomologie des schémas, sections \ref{coherent-section-existence},
\ref{coherent-section-existence-proper}, et
\ref{coherent-section-existence-proper-support}.
Nous nous plaçons dans la situation suivante.

\begin{situation}
\label{flat-situation-existence}
On dispose ici d'un système projectif d'anneaux $(A_n)$ dont les morphismes de transition
sont surjectifs et dont les noyaux sont localement nilpotents. Posons $A = \lim A_n$. On dispose d'un
schéma $X$ séparé et de présentation finie sur $A$.
Posons $X_n = X \times_{\Spec(A)} \Spec(A_n)$ et
considérons-le comme un sous-schéma fermé de $X$. On suppose en outre donné un système
$(\mathcal{F}_n, \varphi_n)$ où $\mathcal{F}_n$ est un
$\mathcal{O}_{X_n}$-module de présentation finie, plat sur $A_n$, à support propre sur $A_n$,
et
$$
\varphi_n :
\mathcal{F}_n \otimes_{\mathcal{O}_{X_n}} \mathcal{O}_{X_{n - 1}}
\longrightarrow
\mathcal{F}_{n - 1}
$$
est un isomorphisme (notation fondée sur l'équivalence de
Morphismes, lemme \ref{morphisms-lemma-i-star-equivalence}).
\end{situation}

\noindent
Notre but est de déterminer s'il existe un faisceau quasi-cohérent $\mathcal{F}$
sur $X$ tel que
$\mathcal{F}_n = \mathcal{F} \otimes_{\mathcal{O}_X} \mathcal{O}_{X_n}$
pour tout $n$.

\begin{lemma}
\label{flat-lemma-compute-what-it-should-be}
Dans la situation \ref{flat-situation-existence}, considérons
$$
K = R\lim_{D_\QCoh(\mathcal{O}_X)}(\mathcal{F}_n) =
DQ_X(R\lim_{D(\mathcal{O}_X)}\mathcal{F}_n)
$$
Alors $K$ appartient à $D^b_{\QCoh}(\mathcal{O}_X)$ et, en fait,
$K$ n'a de faisceaux de cohomologie non nuls qu'en degrés $\geq 0$.
\end{lemma}

\begin{proof}
C'est un cas particulier de
Catégories dérivées de schémas, exemple
\ref{perfect-example-inverse-limit-quasi-coherent}.
\end{proof}

\begin{lemma}
\label{flat-lemma-compute-against-perfect}
Dans la situation \ref{flat-situation-existence}, soit $K$ comme dans le
lemme \ref{flat-lemma-compute-what-it-should-be}. Pour tout
objet parfait $E$ de $D(\mathcal{O}_X)$, on a
\begin{enumerate}
\item $M = R\Gamma(X, K \otimes^\mathbf{L} E)$ est un objet parfait de $D(A)$
et il existe un isomorphisme canonique
$R\Gamma(X_n, \mathcal{F}_n \otimes^\mathbf{L} E|_{X_n}) =
M \otimes_A^\mathbf{L} A_n$
dans $D(A_n)$,
\item $N = R\Hom_X(E, K)$ est un objet parfait de $D(A)$
et il existe un isomorphisme canonique
$R\Hom_{X_n}(E|_{X_n}, \mathcal{F}_n) = N \otimes_A^\mathbf{L} A_n$
dans $D(A_n)$.
\end{enumerate}
Dans les deux assertions, $E|_{X_n}$ désigne l'image inverse dérivée
de $E$ sur $X_n$.
\end{lemma}

\begin{proof}
Démontrons (2). Posons $E_n = E|_{X_n}$ et
$N_n = R\Hom_{X_n}(E_n, \mathcal{F}_n)$.
Rappelons que $R\Hom_{X_n}(-, -)$ est égal à
$R\Gamma(X_n, R\SheafHom(-, -))$, voir
Cohomologie, section \ref{cohomology-section-global-RHom}.
Il résulte alors de Catégories dérivées de schémas, lemme
\ref{perfect-lemma-base-change-RHom-perfect}
que $N_n$ est un objet parfait de $D(A_n)$
dont la formation commute au changement de base. Ainsi, les morphismes
$N_n \otimes_{A_n}^\mathbf{L} A_{n - 1} \to N_{n - 1}$
induits par $\varphi_n$ sont des isomorphismes.
D'après Compléments d'algèbre, lemme \ref{more-algebra-lemma-Rlim-perfect-gives-perfect}, nous constatons que
$R\lim N_n$ est parfait et
que son changement de base à $A_n$ redonne $N_n$.
D'autre part, le foncteur exact
$R\Hom_X(E, -) : D_\QCoh(\mathcal{O}_X) \to D(A)$
entre catégories triangulées commute aux produits
et donc aux limites dérivées, d'où
$$
R\Hom_X(E, K) =
R\lim R\Hom_X(E, \mathcal{F}_n) =
R\lim R\Hom_X(E_n, \mathcal{F}_n) =
R\lim N_n
$$
Ceci prouve (2). Pour établir (1), on se ramène à (2)
à l'aide de Cohomologie, lemme \ref{cohomology-lemma-dual-perfect-complex}.
\end{proof}

\begin{lemma}
\label{flat-lemma-relative-pseudo-coherence}
Dans la situation \ref{flat-situation-existence}, soit $K$ comme dans le
lemme \ref{flat-lemma-compute-what-it-should-be}. Alors $K$
est pseudo-cohérent relativement à $A$.
\end{lemma}

\begin{proof}
En combinant le lemme \ref{flat-lemma-compute-against-perfect} et
Catégories dérivées de schémas, lemme \ref{perfect-lemma-perfect-enough},
on voit que $R\Gamma(X, K \otimes^\mathbf{L} E)$
est pseudo-cohérent dans $D(A)$ pour tout objet pseudo-cohérent
$E$ de $D(\mathcal{O}_X)$. Le lemme découle donc de
Compléments sur les morphismes, lemme
\ref{more-morphisms-lemma-characterize-pseudo-coherent}.
\end{proof}

\begin{lemma}
\label{flat-lemma-compute-over-affine}
Dans la situation \ref{flat-situation-existence}, soit $K$ comme dans le
lemme \ref{flat-lemma-compute-what-it-should-be}. Pour tout
ouvert quasi-compact $U \subset X$, on a
$$
R\Gamma(U, K) \otimes_A^\mathbf{L} A_n =
R\Gamma(U_n, \mathcal{F}_n)
$$
dans $D(A_n)$, où $U_n = U \cap X_n$.
\end{lemma}

\begin{proof}
Fixons $n$. D'après Catégories dérivées de schémas, lemme
\ref{perfect-lemma-computing-sections-as-colim}
il existe un système de complexes parfaits $E_m$
sur $X$ tel que
$R\Gamma(U, K) = \text{hocolim} R\Gamma(X, K \otimes^\mathbf{L} E_m)$.
En fait, cette formule vaut non seulement pour $K$, mais pour tout objet de
$D_\QCoh(\mathcal{O}_X)$.
En l'appliquant à $\mathcal{F}_n$
nous obtenons
\begin{align*}
R\Gamma(U_n, \mathcal{F}_n)
& =
R\Gamma(U, \mathcal{F}_n) \\
& =
\text{hocolim}_m R\Gamma(X, \mathcal{F}_n \otimes^\mathbf{L} E_m) \\
& =
\text{hocolim}_m R\Gamma(X_n, \mathcal{F}_n \otimes^\mathbf{L} E_m|_{X_n})
\end{align*}
En utilisant le lemme \ref{flat-lemma-compute-against-perfect}
et le fait que $- \otimes_A^\mathbf{L} A_n$
commute aux colimites homotopiques, nous obtenons le résultat.
\end{proof}

\begin{lemma}
\label{flat-lemma-finitely-presented}
Dans la situation \ref{flat-situation-existence}, soit $K$ comme dans le
lemme \ref{flat-lemma-compute-what-it-should-be}.
Notons $X_0 \subset X$ le sous-ensemble fermé
constitué des points situés au-dessus du sous-ensemble fermé
$\Spec(A_1) = \Spec(A_2) = \ldots$ de $\Spec(A)$.
Il existe un ouvert $W \subset X$ contenant $X_0$
tel que
\begin{enumerate}
\item $H^i(K)|_W$ est nul sauf si $i = 0$,
\item $\mathcal{F} = H^0(K)|_W$ est de présentation finie, et
\item $\mathcal{F}_n = \mathcal{F} \otimes_{\mathcal{O}_X} \mathcal{O}_{X_n}$.
\end{enumerate}
\end{lemma}

\begin{proof}
Fixons $n \geq 1$. Par construction, il existe un morphisme canonique
$K \to \mathcal{F}_n$ dans $D_\QCoh(\mathcal{O}_X)$
et donc un morphisme canonique $H^0(K) \to \mathcal{F}_n$
de faisceaux quasi-cohérents. Ceci explique le sens du point (3).

\medskip\noindent
Soit $x \in X_0$ un point.
Nous allons trouver un voisinage ouvert $W$
de $x$ tel que (1), (2) et (3) soient vrais. Puisque $X_0$ est quasi-compact,
ceci démontrera le lemme. Soit $U \subset X$ un voisinage ouvert affine
de $x$. Écrivons $U = \Spec(B)$.
Choisissons une surjection $P \to B$ avec $P$ lisse sur $A$.
D'après le lemme \ref{flat-lemma-relative-pseudo-coherence}
et la définition de la pseudo-cohérence relative,
il existe un complexe $F^\bullet$ borné supérieurement
de $P$-modules libres de rang fini représentant
$Ri_*K$, où $i : U \to \Spec(P)$ est l'immersion
fermée induite par la présentation.
Soit $M_n$ le $B$-module correspondant à $\mathcal{F}_n|_U$.
D'après le lemme \ref{flat-lemma-compute-over-affine},
$$
H^i(F^\bullet \otimes_A A_n) =
\left\{
\begin{matrix}
0 & \text{si} & i \not = 0 \\
M_n & \text{si} & i = 0
\end{matrix}
\right.
$$
Soit $i$ l'indice maximal tel que $F^i$ soit non nul.
Si $i \leq 0$, alors (1), (2) et (3) sont vrais.
Sinon, $i > 0$ et nous voyons que le rang du morphisme
$$
F^{i - 1} \to F^i
$$
au point $x$ est maximal. Il est donc maximal dans un voisinage ouvert
de $x$ dans $\Spec(P)$. Ainsi, après avoir remplacé
$P$ par une localisation principale, nous pouvons supposer que le morphisme
affiché est surjectif. Puisque $F^i$ est libre de rang fini, nous pouvons choisir
une décomposition $F^{i - 1} = F' \oplus F^i$. Nous pouvons alors
remplacer $F^\bullet$ par le complexe
$$
\ldots \to F^{i - 2} \to F' \to 0 \to \ldots
$$
et conclure par récurrence sur $i$.
\end{proof}

\begin{lemma}
\label{flat-lemma-proper-support}
Dans la situation \ref{flat-situation-existence}, soit $K$ comme dans le
lemme \ref{flat-lemma-compute-what-it-should-be}. Soit $W \subset X$
comme dans le lemme \ref{flat-lemma-finitely-presented}.
Posons $\mathcal{F} = H^0(K)|_W$. Alors, après avoir éventuellement restreint l'ouvert $W$,
le support de $\mathcal{F}$ est propre sur $A$.
\end{lemma}

\begin{proof}
Fixons $n \geq 1$. Soit $I_n = \Ker(A \to A_n)$.
D'après le lemme \ref{more-algebra-lemma-limit-henselian} de Compléments d'algèbre,
le couple $(A, I_n)$ est hensélien.
Soit $Z \subset W$ le support de $\mathcal{F}$.
C'est un sous-ensemble fermé puisque $\mathcal{F}$ est de présentation finie.
D'après le point (3) du lemme \ref{flat-lemma-finitely-presented},
nous voyons que $Z \times_{\Spec(A)} \Spec(A_n)$
est égal au support de $\mathcal{F}_n$ et est donc
propre sur $\Spec(A/I)$.
D'après le lemme de Compléments sur les morphismes
\ref{more-morphisms-lemma-split-off-proper-part-henselian}
nous pouvons écrire $Z = Z_1 \amalg Z_2$, où $Z_1, Z_2$ sont ouverts et
fermés dans $Z$, où $Z_1$ est propre
sur $A$ et où $Z_1 \times_{\Spec(A)} \Spec(A/I_n)$
est égal au support de $\mathcal{F}_n$.
Autrement dit, $Z_2$ ne rencontre pas $X_0$.
Ainsi, après avoir remplacé $W$ par $W \setminus Z_2$, nous obtenons le lemme.
\end{proof}

\begin{lemma}
\label{flat-lemma-monomorphism-isomorphism}
Soit $A = \lim A_n$ la limite d'un système d'anneaux
dont les morphismes de transition sont surjectifs et les noyaux
localement nilpotents. Soit $S = \Spec(A)$. Soit $T \to S$ un monomorphisme
localement de type fini. Si $\Spec(A_n) \to S$
se factorise par $T$ pour tout $n$, alors $T = S$.
\end{lemma}

\begin{proof}
Posons $S_n = \Spec(A_n)$. Soit $T_0 \subset T$ l'image
commune des factorisations $S_n \to T$. Alors $T_0$ est quasi-compact.
Soit $T' \subset T$ un ouvert quasi-compact contenant $T_0$.
Alors $S_n \to T$ se factorise par $T'$.
Si nous pouvons montrer que $T' = S$, alors $T' = T = S$.
Nous pouvons donc supposer que $T$ est quasi-compact.

\medskip\noindent
Supposons que $T$ soit quasi-compact.
Dans ce cas, $T \to S$ est séparé et quasi-fini
(lemme de Morphismes
\ref{morphisms-lemma-monomorphism-loc-finite-type-loc-quasi-finite}).
En utilisant le théorème principal de Zariski
(sous la forme du lemme de Compléments sur les morphismes
\ref{more-morphisms-lemma-quasi-finite-separated-pass-through-finite})
nous choisissons une factorisation $T \to W \to S$ avec $W \to S$ fini
et $T \to W$ une immersion ouverte. Écrivons $W = \Spec(B)$.
Les factorisations (uniques) $S_n \to T$ peuvent être vues
comme des morphismes vers $W$ et nous obtenons
$$
A \longrightarrow B \longrightarrow \lim A_n = A
$$
Considérons le morphisme $h : S = \Spec(A) \to \Spec(B) = W$ provenant
de la flèche de droite. Alors
$$
T \times_{W, h} S
$$
est un sous-schéma ouvert de $S$ contenant l'image de $S_n \to S$ pour tout $n$.
Pour achever la démonstration, il suffit de montrer que tout ouvert $U \subset S$
contenant l'image de $S_n \to S$ pour un certain $n \geq 1$ est égal à $S$.
C'est vrai parce que $(A, \Ker(A \to A_n))$ est un couple hensélien
(lemme \ref{more-algebra-lemma-limit-henselian} de Compléments d'algèbre)
et que, par conséquent, tout point fermé de $S$ appartient à l'image de $S_n \to S$.
\end{proof}

\begin{theorem}[Théorème d'existence de Grothendieck]
\label{flat-theorem-existence}
Dans la situation \ref{flat-situation-existence},
il existe un $\mathcal{O}_X$-module de présentation finie
$\mathcal{F}$, plat sur $A$, dont le support est propre sur $A$,
tel que
$\mathcal{F}_n = \mathcal{F} \otimes_{\mathcal{O}_X} \mathcal{O}_{X_n}$
pour tout $n$, de manière compatible avec les morphismes $\varphi_n$.
\end{theorem}

\begin{proof}
Appliquons les lemmes \ref{flat-lemma-compute-what-it-should-be},
\ref{flat-lemma-compute-against-perfect},
\ref{flat-lemma-relative-pseudo-coherence},
\ref{flat-lemma-compute-over-affine},
\ref{flat-lemma-finitely-presented}, et
\ref{flat-lemma-proper-support}
pour obtenir un sous-schéma ouvert $W \subset X$ contenant tous les points
situés au-dessus de $\Spec(A_n)$
et un $\mathcal{O}_W$-module de présentation finie $\mathcal{F}$
dont le support est propre sur $A$, avec
$\mathcal{F}_n = \mathcal{F} \otimes_{\mathcal{O}_W} \mathcal{O}_{X_n}$
pour tout $n \geq 1$. (Ceci a un sens puisque $X_n \subset W$.)
D'après le lemme \ref{flat-lemma-proper-pure}, nous voyons que $\mathcal{F}$
est universellement pur relativement à $\Spec(A)$.
D'après le théorème \ref{flat-theorem-flat-dimension-n-representable}
(pour une explication, voir le lemme \ref{flat-lemma-when-universal-flattening}),
il existe un platificateur universel $S' \to \Spec(A)$
de $\mathcal{F}$ et, de plus, le morphisme $S' \to \Spec(A)$
est un monomorphisme de présentation finie.
Puisque le changement de base de $\mathcal{F}$ à $\Spec(A_n)$
est $\mathcal{F}_n$, nous trouvons que $\Spec(A_n) \to \Spec(A)$
se factorise (de manière unique) par $S'$ pour tout $n$.
D'après le lemme \ref{flat-lemma-monomorphism-isomorphism},
nous voyons que $S' = \Spec(A)$.
Cela signifie que $\mathcal{F}$ est plat sur $A$.
Enfin, puisque le support schématique $Z$ de $\mathcal{F}$
est propre sur $\Spec(A)$, le morphisme $Z \to X$ est fermé.
Par conséquent, l'image directe $(W \to X)_*\mathcal{F}$ a son support
dans $W$ et possède toutes les propriétés voulues.
\end{proof}






\section{Théorème d'existence de Grothendieck, V}
\label{flat-section-existence-derived}

\noindent
Dans cette section, nous démontrons un analogue du théorème d'existence de Grothendieck
dans la catégorie dérivée, en suivant la méthode utilisée dans la
section \ref{flat-section-existence} pour les modules quasi-cohérents.
Le cas classique est étudié dans
Cohomologie des schémas, sections \ref{coherent-section-existence},
\ref{coherent-section-existence-proper}, et
\ref{coherent-section-existence-proper-support}.
Nous travaillerons dans la situation suivante.

\begin{situation}
\label{flat-situation-existence-derived}
Nous avons ici un système projectif d'anneaux $(A_n)$ dont les morphismes de transition
sont surjectifs et dont les noyaux sont localement nilpotents. Posons $A = \lim A_n$. Nous avons un
schéma $X$ propre, plat et de présentation finie sur $A$.
Nous posons $X_n = X \times_{\Spec(A)} \Spec(A_n)$ et le
considérons comme un sous-schéma fermé de $X$. Supposons en outre donné un système
$(K_n, \varphi_n)$ où $K_n$ est un objet pseudo-cohérent de
$D(\mathcal{O}_{X_n})$ et
$$
\varphi_n : K_n \longrightarrow K_{n - 1}
$$
est un morphisme dans $D(\mathcal{O}_{X_n})$ qui induit un isomorphisme
$K_n \otimes_{\mathcal{O}_{X_n}}^\mathbf{L} \mathcal{O}_{X_{n - 1}}
\to K_{n - 1}$ dans $D(\mathcal{O}_{X_{n - 1}})$.
\end{situation}

\noindent
Plus précisément, nous devrions écrire
$\varphi_n : K_n \to Ri_{n - 1, *}K_{n - 1}$
où $i_{n - 1} : X_{n - 1} \to X_n$ est le morphisme d'inclusion,
et, dans cette notation, la condition est que le morphisme
adjoint $Li_{n - 1}^*K_n \to K_{n - 1}$ soit un isomorphisme.
Notre but est de trouver un $K \in D(\mathcal{O}_X)$ pseudo-cohérent
tel que $K_n = K \otimes_{\mathcal{O}_X}^\mathbf{L} \mathcal{O}_{X_n}$
pour tout $n$ (avec le même abus de notation).

\begin{lemma}
\label{flat-lemma-compute-what-it-should-be-derived}
Dans la situation \ref{flat-situation-existence-derived}, considérons
$$
K = R\lim_{D_\QCoh(\mathcal{O}_X)}(K_n) =
DQ_X(R\lim_{D(\mathcal{O}_X)} K_n)
$$
Alors $K$ appartient à $D^-_{\QCoh}(\mathcal{O}_X)$.
\end{lemma}

\begin{proof}
Le foncteur $DQ_X$ existe parce que $X$ est quasi-compact et
quasi-séparé, voir le lemme de Catégories dérivées de schémas
\ref{perfect-lemma-better-coherator}.
Puisque $DQ_X$ est un adjoint à droite, il commute aux produits
et donc aux limites dérivées. D'où l'égalité
dans l'énoncé du lemme.

\medskip\noindent
D'après le lemme de Catégories dérivées de schémas
\ref{perfect-lemma-boundedness-better-coherator},
le foncteur $DQ_X$ est de dimension cohomologique bornée.
Il suffit donc de montrer que $R\lim K_n \in D^-(\mathcal{O}_X)$.
Pour le voir, soit $U \subset X$ un ouvert affine.
Il existe alors une suite exacte canonique
$$
0 \to
R^1\lim H^{m - 1}(U, K_n) \to H^m(U, R\lim K_n) \to
\lim H^m(U, K_n) \to 0
$$
d'après le lemme \ref{cohomology-lemma-RGamma-commutes-with-Rlim} de Cohomologie.
Puisque $U$ est affine et $K_n$ est pseudo-cohérent (et possède donc des
faisceaux de cohomologie quasi-cohérents d'après le lemme
\ref{perfect-lemma-pseudo-coherent} de Catégories dérivées de schémas),
nous voyons que $H^m(U, K_n) = H^m(K_n)(U)$ d'après le
lemme de Catégories dérivées de schémas
\ref{perfect-lemma-affine-compare-bounded}.
Nous concluons ainsi qu'il suffit de montrer que $K_n$
est borné supérieurement indépendamment de $n$.

\medskip\noindent
Puisque $K_n$ est pseudo-cohérent, nous avons $K_n \in D^-(\mathcal{O}_{X_n})$.
Supposons que $a_n$ soit maximal tel que $H^{a_n}(K_n)$ soit non nul.
Bien entendu, $a_1 \leq a_2 \leq a_3 \leq \ldots$.
Notons que $H^{a_n}(K_n)$ est un
$\mathcal{O}_{X_n}$-module de présentation finie
(lemme \ref{cohomology-lemma-finite-cohomology} de Cohomologie).
Nous avons $H^{a_n}(K_{n - 1}) =
H^{a_n}(K_n) \otimes_{\mathcal{O}_{X_n}} \mathcal{O}_{X_{n - 1}}$.
Puisque $X_{n - 1} \to X_n$ est un épaississement, il résulte du
lemme de Nakayama (lemme \ref{algebra-lemma-NAK} d'Algèbre) que si
$H^{a_n}(K_n) \otimes_{\mathcal{O}_{X_n}} \mathcal{O}_{X_{n - 1}}$
est nul, alors $H^{a_n}(K_n)$ est lui aussi nul. Ainsi $a_n = a_{n - 1}$
pour tout $n$, et nous concluons.
\end{proof}

\begin{lemma}
\label{flat-lemma-compute-against-perfect-derived}
Dans la situation \ref{flat-situation-existence-derived}, soit $K$ comme dans le
lemme \ref{flat-lemma-compute-what-it-should-be-derived}. Pour tout objet parfait
$E$ de $D(\mathcal{O}_X)$, la cohomologie
$$
M = R\Gamma(X, K \otimes^\mathbf{L} E)
$$
est un objet pseudo-cohérent de $D(A)$ et il existe un isomorphisme canonique
$$
R\Gamma(X_n, K_n \otimes^\mathbf{L} E|_{X_n}) = M \otimes_A^\mathbf{L} A_n
$$
dans $D(A_n)$. Ici, $E|_{X_n}$ désigne l'image inverse dérivée de $E$ sur $X_n$.
\end{lemma}

\begin{proof}
Écrivons $E_n = E|_{X_n}$ et
$M_n = R\Gamma(X_n, K_n \otimes^\mathbf{L} E|_{X_n})$.
D'après le lemme de Catégories dérivées de schémas
\ref{perfect-lemma-flat-proper-pseudo-coherent-direct-image-general}
nous voyons que $M_n$ est un objet pseudo-cohérent de $D(A_n)$
dont la formation commute au changement de base. Ainsi, les morphismes
$M_n \otimes_{A_n}^\mathbf{L} A_{n - 1} \to M_{n - 1}$
provenant de $\varphi_n$ sont des isomorphismes. D'après le
lemme de Compléments d'algèbre
\ref{more-algebra-lemma-Rlim-pseudo-coherent-gives-pseudo-coherent}
nous trouvons que $R\lim M_n$ est pseudo-cohérent et
que son changement de base à $A_n$ redonne $M_n$.
D'autre part, le foncteur exact
$R\Gamma(X, -) : D_\QCoh(\mathcal{O}_X) \to D(A)$
entre catégories triangulées commute aux produits
et donc aux limites dérivées, d'où
$$
R\Gamma(X, E \otimes^\mathbf{L} K) =
R\lim R\Gamma(X,  E \otimes^\mathbf{L} K_n) =
R\lim R\Gamma(X_n, E_n \otimes^\mathbf{L} K_n) =
R\lim M_n
$$
comme souhaité.
\end{proof}

\begin{lemma}
\label{flat-lemma-relative-pseudo-coherence-derived}
Dans la situation \ref{flat-situation-existence-derived}, soit $K$ comme dans le
lemme \ref{flat-lemma-compute-what-it-should-be-derived}. Alors $K$
est pseudo-cohérent sur $X$.
\end{lemma}

\begin{proof}
En combinant le lemme \ref{flat-lemma-compute-against-perfect-derived} et le
lemme de Catégories dérivées de schémas
\ref{perfect-lemma-perfect-enough}
nous voyons que $R\Gamma(X, K \otimes^\mathbf{L} E)$
est pseudo-cohérent dans $D(A)$ pour tout objet pseudo-cohérent
$E$ de $D(\mathcal{O}_X)$. Il résulte donc du
lemme de Compléments sur les morphismes
\ref{more-morphisms-lemma-characterize-pseudo-coherent}
que $K$ est pseudo-cohérent relativement à $A$.
Puisque $X$ est plat et de présentation finie
sur $A$, cela équivaut à être pseudo-cohérent sur $X$, voir le
lemme de Compléments sur les morphismes
\ref{more-morphisms-lemma-check-relative-pseudo-coherence-on-charts}.
\end{proof}

\begin{lemma}
\label{flat-lemma-compute-over-affine-derived}
Dans la situation \ref{flat-situation-existence-derived}, soit $K$ comme dans le
lemme \ref{flat-lemma-compute-what-it-should-be-derived}. Pour tout ouvert quasi-compact
$U \subset X$, nous avons
$$
R\Gamma(U, K) \otimes_A^\mathbf{L} A_n =
R\Gamma(U_n, K_n)
$$
dans $D(A_n)$, où $U_n = U \cap X_n$.
\end{lemma}

\begin{proof}
Fixons $n$. D'après le lemme de Catégories dérivées de schémas
\ref{perfect-lemma-computing-sections-as-colim}
il existe un système de complexes parfaits $E_m$
sur $X$ tel que
$R\Gamma(U, K) = \text{hocolim} R\Gamma(X, K \otimes^\mathbf{L} E_m)$.
En fait, cette formule vaut non seulement pour $K$, mais pour tout objet de
$D_\QCoh(\mathcal{O}_X)$.
En l'appliquant à $K_n$,
nous obtenons
\begin{align*}
R\Gamma(U_n, K_n)
& =
R\Gamma(U, K_n) \\
& =
\text{hocolim}_m R\Gamma(X, K_n \otimes^\mathbf{L} E_m) \\
& =
\text{hocolim}_m R\Gamma(X_n, K_n \otimes^\mathbf{L} E_m|_{X_n})
\end{align*}
En utilisant le lemme \ref{flat-lemma-compute-against-perfect-derived}
et le fait que $- \otimes_A^\mathbf{L} A_n$
commute aux colimites homotopiques, nous obtenons le résultat.
\end{proof}

\begin{theorem}[Théorème d'existence de Grothendieck dérivé]
\label{flat-theorem-existence-derived}
Dans la situation \ref{flat-situation-existence-derived},
il existe un $K$ pseudo-cohérent dans $D(\mathcal{O}_X)$
tel que $K_n = K \otimes_{\mathcal{O}_X}^\mathbf{L} \mathcal{O}_{X_n}$
pour tout $n$, de manière compatible avec les morphismes $\varphi_n$.
\end{theorem}

\begin{proof}
Appliquons les lemmes \ref{flat-lemma-compute-what-it-should-be-derived},
\ref{flat-lemma-compute-against-perfect-derived},
\ref{flat-lemma-relative-pseudo-coherence-derived}
pour obtenir un objet pseudo-cohérent $K$ de $D(\mathcal{O}_X)$.
En choisissant des ouverts affines dans le lemme
\ref{flat-lemma-compute-over-affine-derived}
il s'ensuit immédiatement que $K$ a pour restriction $K_n$ sur $X_n$.
\end{proof}

\begin{remark}
\label{flat-remark-correct-generality}
Le résultat de cette section peut être généralisé. Il est probablement valable
si nous supposons seulement que $X \to \Spec(A)$ est séparé et de présentation finie,
et que $K_n$ est pseudo-cohérent relativement à $A_n$ et à support dans un sous-ensemble
fermé de $X_n$ propre sur $A_n$. Le résultat sera un $K$ qui
est pseudo-cohérent relativement à $A$ et à support dans un sous-ensemble fermé
propre sur $A$. Si nous en avons un jour besoin, nous
formulerons un énoncé précis et le démontrerons ici.
\end{remark}






\section{Éclatements et platitude}
\label{flat-section-blowup-flat}

\noindent
Dans cette section, nous poursuivons l'étude de résultats de la forme : « Après un
éclatement, le transformé strict devient plat », voir la
section \ref{more-algebra-section-blowup-flat} de Compléments d'algèbre et la
section \ref{divisors-section-blowup-flat} de Diviseurs.
Nous utiliserons dans cette section la notation suivante
(plus ou moins standard). Si $X \to S$ est
un morphisme de schémas, si $\mathcal{F}$ est un module quasi-cohérent
sur $X$ et si $T \to S$ est un morphisme de schémas, alors nous notons
$\mathcal{F}_T$ l'image inverse de $\mathcal{F}$ sur le changement de base
$X_T = X \times_S T$.

\begin{remark}
\label{flat-remark-successive-blowups}
Soit $S$ un schéma quasi-compact et quasi-séparé. Soit $f : X \to S$
un morphisme de schémas. Soit $\mathcal{F}$ un module quasi-cohérent sur $X$.
Soit $U \subset S$ un sous-schéma ouvert quasi-compact. Étant donné un éclatement $U$-admissible
$S' \to S$, nous notons $X'$ le transformé strict de $X$ et $\mathcal{F}'$
le transformé strict de $\mathcal{F}$, considéré comme un module quasi-cohérent
sur $X'$ (au moyen du lemme \ref{divisors-lemma-strict-transform} de Diviseurs).
Soit $P$ une propriété de $\mathcal{F}/X/S$ qui est stable par passage au transformé
strict (comme ci-dessus) pour les éclatements $U$-admissibles. Le problème général de
cette section est de montrer (sous des conditions auxiliaires sur $\mathcal{F}/X/S$)
qu'il existe un éclatement $U$-admissible $S' \to S$
tel que le transformé strict $\mathcal{F}'/X'/S'$ possède $P$.

\medskip\noindent
La stratégie générale consistera à utiliser le fait qu'une composée
d'éclatements $U$-admissibles est un éclatement $U$-admissible, voir le
lemme \ref{divisors-lemma-composition-admissible-blowups} de Diviseurs.
En fait, nous utiliserons le lemme plus précis
\ref{divisors-lemma-composition-finite-type-blowups} de Diviseurs
et le combinerons avec le
lemme \ref{divisors-lemma-strict-transform-composition-blowups} de Diviseurs.
Il en résulte qu'il suffit de trouver une suite d'éclatements $U$-admissibles
successifs
$$
S = S_0 \leftarrow S_1 \leftarrow \ldots \leftarrow S_n
$$
telle que, en posant $\mathcal{F}_0 = \mathcal{F}$ et $X_0 = X$, puis en définissant
$\mathcal{F}_i/X_i$ comme le transformé strict de
$\mathcal{F}_{i - 1}/X_{i - 1}$, nous
aboutissions à $\mathcal{F}_n/X_n/S_n$ possédant la propriété $P$.

\medskip\noindent
En particulier, choisissons un faisceau quasi-cohérent d'idéaux de type fini
$\mathcal{I} \subset \mathcal{O}_S$ tel que $V(\mathcal{I}) = S \setminus U$,
voir le lemme \ref{properties-lemma-quasi-coherent-finite-type-ideals} de Propriétés.
Soit $S' \to S$ l'éclatement le long de $\mathcal{I}$ et soit $E \subset S'$
le diviseur exceptionnel (lemme de Diviseurs
\ref{divisors-lemma-blowing-up-gives-effective-Cartier-divisor}).
Nous voyons alors que nous avons ramené le
problème au cas où il existe un diviseur de Cartier effectif
$D \subset S$ dont le support est $X \setminus U$. En particulier, nous pouvons
supposer que $U$ est schématiquement dense dans $S$
(lemme \ref{divisors-lemma-complement-effective-Cartier-divisor} de Diviseurs).

\medskip\noindent
Supposons que $P$ soit locale sur $S$ : si $S = \bigcup S_i$ est un recouvrement ouvert
fini par des ouverts quasi-compacts et si $P$ vaut pour
$\mathcal{F}_{S_i}/X_{S_i}/S_i$, alors $P$ vaut pour $\mathcal{F}/X/S$.
Dans ce cas, le problème général ci-dessus est lui aussi local sur $S$, c'est-à-dire que,
si, pour tout $s \in S$, nous pouvons trouver un voisinage ouvert quasi-compact $W$ de $s$
tel que le problème pour $\mathcal{F}_W/X_W/W$ soit résoluble, alors le
problème est résoluble pour $\mathcal{F}/X/S$. Cela résulte des
lemmes \ref{divisors-lemma-extend-admissible-blowups} et
\ref{divisors-lemma-dominate-admissible-blowups}.
\end{remark}

\begin{lemma}
\label{flat-lemma-helper-blowup-affine-space}
Soit $R$ un anneau et soit $f \in R$. Soit $r\geq 0$ un entier.
Soit $R \to S$ un morphisme d'anneaux et soit $M$ un $S$-module. Supposons que
\begin{enumerate}
\item $R \to S$ est de présentation finie et plat,
\item tout anneau fibre $S \otimes_R \kappa(\mathfrak p)$ est
géométriquement intègre sur $R$,
\item $M$ est un $S$-module de type fini,
\item $M_f$ est un $S_f$-module de présentation finie,
\item pour tout $\mathfrak p \in R$, $f \not \in \mathfrak p$, en posant
$\mathfrak q = \mathfrak pS$, le module $M_{\mathfrak q}$ est libre
de rang $r$ sur $S_\mathfrak q$.
\end{enumerate}
Alors il existe un idéal de type fini $I \subset R$ avec
$V(f) = V(I)$ tel que, pour tout $a \in I$, en posant $R' = R[\frac{I}{a}]$,
le quotient
$$
M' = (M \otimes_R R')/a\text{-torsion}
$$
sur $S' = S \otimes_R R'$ possède la propriété suivante : pour tout idéal premier
$\mathfrak p' \subset R'$, il existe $g \in S'$,
$g \not \in \mathfrak p'S'$, tel que $M'_g$ soit un $S'_g$-module libre
de rang $r$.
\end{lemma}

\begin{proof}
Ce lemme est une généralisation du
lemme \ref{more-algebra-lemma-blowup-module} de Compléments d'algèbre ;
nous recommandons au lecteur de lire d'abord cette démonstration.
Choisissons une surjection $S^{\oplus n} \to M$, ce qui est possible par (1).
Choisissons un sous-module de type fini $K \subset \Ker(S^{\oplus n} \to M)$
tel que $S^{\oplus n}/K \to M$ devienne un isomorphisme après inversion de $f$.
Ceci est possible par (4). Posons $M_1 = S^{\oplus n}/K$ et supposons que nous puissions
démontrer le lemme pour $M_1$. Soit $I \subset R$ l'idéal correspondant.
Alors, pour $a \in I$, le morphisme
$$
M_1' = (M_1 \otimes_R R')/a\text{-torsion}
\longrightarrow
M' = (M \otimes_R R')/a\text{-torsion}
$$
est surjectif. C'est aussi un isomorphisme après inversion de $a$ dans $R'$,
car $R'_a = R_f$, voir le lemme \ref{algebra-lemma-blowup-in-principal} d'Algèbre.
Mais $a$ est un non-diviseur de zéro dans $M'_1$, de sorte que le morphisme affiché est un
isomorphisme. Il suffit donc de démontrer le lemme lorsque $M$ est un module
de présentation finie sur $S$.

\medskip\noindent
Supposons que $M$ soit un $S$-module de présentation finie vérifiant (3).
Alors $J = \text{Fit}_r(M) \subset S$ est un idéal de type fini.
D'après le lemme \ref{flat-lemma-fibres-irreducible-flat-projective-nonnoetherian},
nous pouvons écrire $S$ comme facteur direct d'un
$R$-module libre : $\bigoplus_{\alpha \in A} R = S \oplus C$.
Pour tout élément $h \in S$, en écrivant $h = \sum a_\alpha$ dans la
décomposition ci-dessus, nous disons que les $a_\alpha$ sont les coefficients de $h$.
Soit $I' \subset R$ l'idéal des coefficients
des éléments de $J$. La multiplication par un élément de $S$ définit
un morphisme $R$-linéaire $S \to S$ ; par conséquent, $I'$ est engendré par les coefficients
des générateurs de $J$, c'est-à-dire que $I'$ est un idéal de type fini.
Nous affirmons que $I = fI'$ convient.

\medskip\noindent
Vérifions d'abord que $V(f) = V(I)$. L'inclusion $V(f) \subset V(I)$ est
claire. Réciproquement, si $f \not \in \mathfrak p$, alors
$\mathfrak q =  \mathfrak p S$ n'est pas un élément de $V(J)$ d'après la
propriété (5) et le lemme de Compléments d'algèbre
\ref{more-algebra-lemma-fitting-ideal-generate-locally}.
Il existe donc un
élément de $J$ dont l'image dans $S \otimes_R \kappa(\mathfrak p)$ n'est pas nulle.
Ainsi, il existe un élément de $I'$ qui n'appartient pas à
$\mathfrak p$, donc $\mathfrak p \not \in V(fI') = V(I)$.

\medskip\noindent
Soit $a \in I$ et posons $R' = R[\frac{I}{a}]$. Nous pouvons écrire $a = fa'$
pour un certain $a' \in I'$. D'après les lemmes \ref{algebra-lemma-affine-blowup} et
\ref{algebra-lemma-blowup-add-principal} d'Algèbre, nous voyons que $I' R' = a'R'$
et que $a'$ est un non-diviseur de zéro dans $R'$. Posons $S' = S \otimes_S R'$.
Tout élément $g$ de $JS' = \text{Fit}_r(M \otimes_S S')$ peut
s'écrire $g = \sum_\alpha c_\alpha$ avec $c_\alpha \in I'R'$.
Puisque $I'R' = a'R'$, nous pouvons écrire $c_\alpha = a'c'_\alpha$ pour un certain
$c'_\alpha \in R'$ et $g = (\sum c'_\alpha)a' = g' a'$ dans $S'$.
De plus, il existe $g_0 \in J$ tel que $a' = c_\alpha$
pour un certain $\alpha$. Pour cet élément, nous avons $g_0 = g'_0 a'$ dans $S'$,
où $g'_0$ est une unité de $S'$.
Soit $\mathfrak p' \subset R'$
un idéal premier et $\mathfrak q' = \mathfrak p'S'$.
D'après ce qui précède, nous voyons que $JS'_{\mathfrak q'}$ est l'idéal
principal engendré par le non-diviseur de zéro $a'$.
Il résulte du lemme de Compléments d'algèbre
\ref{more-algebra-lemma-principal-fitting-ideal}
que $M'_{\mathfrak q'}$ peut être engendré par $r$ éléments.
Puisque $M'$ est de type fini, il existe $m_1, \ldots, m_r \in M'$ et
$g \in S'$, $g \not \in \mathfrak q'$, tels que le morphisme correspondant
$(S')^{\oplus r} \to M'$ devienne surjectif après inversion de $g$.

\medskip\noindent
Enfin, considérons l'idéal $J' = \text{Fit}_{r - 1}(M')$.
Notons que $J'S'_g$ est engendré par les coefficients des relations entre
$m_1, \ldots, m_r$ (compatibilité de l'idéal de Fitting avec le changement de base).
Il suffit donc de montrer que $J' = 0$, voir le
lemme de Compléments d'algèbre
\ref{more-algebra-lemma-fitting-ideal-finite-locally-free}.
Puisque $R'_a = R_f$ (lemme \ref{algebra-lemma-blowup-in-principal} d'Algèbre)
et $M'_a = M_f$, nous voyons d'après (5)
que l'image de $J'_a$ est nulle dans $S_{\mathfrak q''}$ pour tout idéal premier
$\mathfrak q'' \subset S'$ de la forme $\mathfrak q'' = \mathfrak p''S'$,
où $\mathfrak p'' \subset R'_a$. Puisque
$S'_a \subset \prod_{\mathfrak q''\text{ comme ci-dessus}} S'_{\mathfrak q''}$
(car $(S'_a)_{\mathfrak p''} \subset S'_{\mathfrak q''}$ d'après le
lemme \ref{flat-lemma-base-change-universally-flat}),
nous voyons que $J'R'_a = 0$. Puisque $a$ est un non-diviseur de zéro dans $R'$, nous
concluons que $J' = 0$, ce qui achève la démonstration.
\end{proof}

\begin{lemma}
\label{flat-lemma-flatten-module-pre}
Soit $S$ un schéma quasi-compact et quasi-séparé.
Soit $X \to S$ un morphisme de schémas.
Soit $\mathcal{F}$ un module quasi-cohérent sur $X$.
Soit $U \subset S$ un ouvert quasi-compact. Supposons que
\begin{enumerate}
\item $X \to S$ est affine, de présentation finie, plat, à
fibres géométriquement intègres,
\item $\mathcal{F}$ est un module de type fini,
\item $\mathcal{F}_U$ est de présentation finie,
\item $\mathcal{F}$ est plat sur $S$ en tous les points génériques des
fibres situées au-dessus de points de $U$.
\end{enumerate}
Alors il existe un éclatement $U$-admissible $S' \to S$
et un sous-schéma ouvert $V \subset X_{S'}$
tels que (a) le transformé strict $\mathcal{F}'$ de $\mathcal{F}$
se restreigne en un $\mathcal{O}_V$-module localement libre de rang fini et que
(b) $V \to S'$ soit surjectif.
\end{lemma}

\begin{proof}
Étant donnés $\mathcal{F}/X/S$ et $U \subset S$ sous les hypothèses du
lemme, notons $P$ la propriété « $\mathcal{F}$ est plat sur $S$ en
tous les points génériques des fibres ». Il est clair que $P$ est préservée par passage au
transformé strict, voir le
lemme \ref{divisors-lemma-strict-transform-flat} de Diviseurs
et le lemme \ref{morphisms-lemma-base-change-module-flat} de Morphismes.
Il est également clair que $P$ est locale sur $S$. Par conséquent, toutes les observations
de la remarque \ref{flat-remark-successive-blowups} s'appliquent au problème posé par
le lemme.

\medskip\noindent
Considérons la fonction $r : U \to \mathbf{Z}_{\geq 0}$ qui associe à
$u \in U$ l'entier
$$
r(u) = \dim_{\kappa(\xi_u)}(\mathcal{F}_{\xi_u} \otimes \kappa(\xi_u))
$$
où $\xi_u$ est le point générique de la fibre $X_u$.
D'après le lemme de Compléments sur les morphismes
\ref{more-morphisms-lemma-flat-and-free-at-point-fibre}
et le fait que l'image d'un ouvert de $X_S$ dans $S$ est ouverte,
nous voyons que $r(u)$ est localement constant. Par conséquent,
$U = U_0 \amalg U_1 \amalg \ldots \amalg U_c$ est une réunion disjointe finie
de sous-schémas ouverts et fermés sur lesquels $r$ est constante de valeur
$i$ sur $U_i$. D'après le
lemme \ref{divisors-lemma-separate-disjoint-opens-by-blowing-up} de Diviseurs,
nous pouvons trouver un éclatement $U$-admissible qui décompose $S$ en
la réunion disjointe de deux schémas, le premier contenant $U_0$ et
le second $U_1 \cup \ldots \cup U_c$. En répétant ceci encore $c - 1$
fois, nous pouvons supposer que $S$ est une réunion disjointe
$S = S_0 \amalg S_1 \amalg \ldots \amalg S_c$ avec $U_i \subset S_i$.
Nous pouvons donc supposer que la fonction $r$ définie ci-dessus est constante, disons
de valeur $r$.

\medskip\noindent
D'après la remarque \ref{flat-remark-successive-blowups}, nous voyons que nous pouvons supposer qu'il
existe un diviseur de Cartier effectif $D \subset S$ dont le support est
$S \setminus U$. Une autre application de la remarque \ref{flat-remark-successive-blowups},
combinée avec le
lemme \ref{divisors-lemma-characterize-effective-Cartier-divisor} de Diviseurs,
nous indique que nous pouvons supposer que
$S = \Spec(R)$ et $D = \Spec(R/(f))$ pour un certain non-diviseur de zéro
$f \in R$. Ce cas est traité par le
lemme \ref{flat-lemma-helper-blowup-affine-space}.
\end{proof}

\begin{lemma}
\label{flat-lemma-trick-fitting-ideal}
Soit $A \to C$ un morphisme d'anneaux fini et localement libre de rang $d$.
Soit $h \in C$ un élément tel que $C_h$ soit \'etale sur $A$.
Soit $J \subset C$ un idéal. Posons $I = \text{Fit}_0(C/J)$, où nous
considérons $C/J$ comme un $A$-module de type fini. Alors $IC_h = JJ'$ pour un certain idéal
$J' \subset C_h$. Si $J$ est de type fini, $I$ et $J'$ le sont aussi.
\end{lemma}

\begin{proof}
Nous utiliserons les propriétés élémentaires des idéaux de Fitting, voir le
lemme \ref{more-algebra-lemma-fitting-ideal-basics} de Compléments d'algèbre.
Alors $IC$ est l'idéal de Fitting de $C/J \otimes_A C$.
Notons que $C \to C \otimes_A C$, $c \mapsto 1 \otimes c$, possède une section
(le morphisme de multiplication). Par hypothèse, $C \to C \otimes_A C$ est \'etale
en tout idéal premier de l'image de $\Spec(C_h)$ par cette section.
Par conséquent, le morphisme de multiplication $C \otimes_A C_h \to C_h$ est \'etale
donc en particulier plat, voir le lemme \ref{algebra-lemma-map-between-etale} d'Algèbre.
Il existe donc une $C_h$-algèbre telle que
$C \otimes_A C_h \cong C_h \oplus C'$ comme $C_h$-algèbres, voir le
lemme \ref{algebra-lemma-surjective-flat-finitely-presented} d'Algèbre.
Ainsi $(C/J) \otimes_A C_h \cong (C_h/J_h) \oplus C'/I'$ comme
$C_h$-modules pour un certain idéal $I' \subset C'$.
Par conséquent, $IC_h = JJ'$ avec $J' = \text{Fit}_0(C'/I')$, où nous considérons
$C'/J'$ comme un $C_h$-module.
\end{proof}

\begin{lemma}
\label{flat-lemma-push-ideal}
Soit $A \to B$ un morphisme d'anneaux \'etale. Soit $a \in A$ un non-diviseur de zéro.
Soit $J \subset B$ un idéal de type fini avec $V(J) \subset V(aB)$.
Pour tout idéal premier $\mathfrak q \subset B$, il existe un idéal de type fini
$I \subset A$ avec $V(I) \subset V(a)$ et un élément
$g \in B$, $g \not \in \mathfrak q$, tels que
$IB_g = JJ'$ pour un certain idéal de type fini $J' \subset B_g$.
\end{lemma}

\begin{proof}
Nous pouvons remplacer $B$ par une localisation principale en
un élément $g \in B$, $g \not \in \mathfrak q$. Nous pouvons donc supposer
que $B$ est \'etale standard, voir la
proposition \ref{algebra-proposition-etale-locally-standard} d'Algèbre.
Nous pouvons donc supposer que $B$ est une localisation
de $C = A[x]/(f)$ pour un certain polynôme unitaire $f \in A[x]$ de degré $d$.
Écrivons $B = C_h$ pour un certain $h \in C$. Choisissons des éléments $h_1, \ldots, h_n \in C$
qui engendrent $J$
sur $B$. La condition $V(J) \subset V(aB)$ signifie que
$a^m = \sum b_i h_i$ dans $B$ pour un certain entier $m$ suffisamment grand. Posons $h_{n + 1} = a^m$.
Comme dans le lemme \ref{flat-lemma-trick-fitting-ideal}, prenons
$I = \text{Fit}_0(C/(h_1, \ldots, h_{r + 1}))$.
Puisque le module $C/(h_1, \ldots, h_{r + 1})$
est annulé par $a^m$, nous voyons que $a^{dm} \in I$, ce qui implique
que $V(I) \subset V(a)$.
\end{proof}

\begin{lemma}
\label{flat-lemma-flatten-module-etale-localize}
Soit $S$ un schéma quasi-compact et quasi-séparé.
Soit $X \to S$ un morphisme de schémas.
Soit $\mathcal{F}$ un module quasi-cohérent sur $X$.
Soit $U \subset S$ un ouvert quasi-compact.
Supposons qu'il existe un nombre fini de diagrammes commutatifs
$$
\xymatrix{
& X_i \ar[r]_{j_i} \ar[d] & X \ar[d] \\
S_i^* \ar[r] & S_i \ar[r]^{e_i} & S
}
$$
où
\begin{enumerate}
\item $e_i : S_i \to S$ sont des morphismes \'etales quasi-compacts et
$S = \bigcup e_i(S_i)$,
\item $j_i : X_i \to X$ sont des morphismes \'etales et
$X = \bigcup j_i(X_i)$,
\item $S^*_i \to S_i$ est un éclatement $e_i^{-1}(U)$-admissible
tel que le transformé strict $\mathcal{F}_i^*$ de $j_i^*\mathcal{F}$
soit plat sur $S^*_i$.
\end{enumerate}
Alors il existe un éclatement $U$-admissible $S' \to S$ tel que
le transformé strict de $\mathcal{F}$ soit plat sur $S'$.
\end{lemma}

\begin{proof}
Nous affirmons que les hypothèses du lemme sont préservées par les éclatements $U$-admissibles.
En effet, supposons que $b : S' \to S$ soit un éclatement $U$-admissible le long du
faisceau quasi-cohérent d'idéaux $\mathcal{I}$. De plus, soit $S^*_i \to S_i$
l'éclatement du faisceau quasi-cohérent d'idéaux $\mathcal{J}_i$.
Alors la famille de morphismes $e'_i : S'_i = S_i \times_S S' \to S'$
et $j'_i : X_i' = X_i \times_S S' \to X \times_S S'$ satisfait aux conditions
(1), (2), (3) pour le transformé strict $\mathcal{F}'$ de $\mathcal{F}$
relativement à l'éclatement $S' \to S$. Premièrement, observons que $S_i'$ est
l'éclatement de $S_i$ le long de l'image inverse de $\mathcal{I}$, voir le
lemme \ref{divisors-lemma-flat-base-change-blowing-up} de Diviseurs.
Deuxièmement, considérons l'éclatement
$S_i^{\prime *} \to S_i'$ de $S_i'$ le long de l'image inverse de l'idéal
$\mathcal{J}_i$.
D'après le lemme \ref{divisors-lemma-blowing-up-two-ideals} de Diviseurs,
nous obtenons un diagramme commutatif
$$
\xymatrix{
S_i^{\prime *} \ar[r] \ar[rd] \ar[d] & S'_i \ar[d] \\
S_i^* \ar[r] & S_i
}
$$
et tous les morphismes du diagramme ci-dessus sont des éclatements. Par conséquent, d'après les
lemmes \ref{divisors-lemma-strict-transform-flat} et
\ref{divisors-lemma-strict-transform-composition-blowups} de Diviseurs,
nous voyons
\begin{align*}
& \text{ le transformé strict de }(j'_i)^*\mathcal{F}'\text{ par }
S_i^{\prime *} \to S_i' \\
= &
\text{ le transformé strict de }j_i^*\mathcal{F}\text{ par }
S_i^{\prime *} \to S_i \\
= &
\text{ le transformé strict de }\mathcal{F}_i'\text{ par }
S_i^{\prime *} \to S_i' \\
= &
\text{ l'image inverse de }\mathcal{F}_i^*\text{ par }
X_i \times_{S_i} S_i^{\prime *} \to X_i
\end{align*}
qui est donc plat sur $S_i^{\prime *}$
(lemme \ref{morphisms-lemma-base-change-module-flat} de Morphismes).
Ceci étant dit, nous voyons que toutes les observations de la
remarque \ref{flat-remark-successive-blowups} s'appliquent au
problème consistant à trouver un éclatement $U$-admissible tel que le
transformé strict de $\mathcal{F}$ devienne plat sur la base
sous les hypothèses du lemme. En particulier, nous pouvons supposer
que $S \setminus U$ est le support d'un diviseur de Cartier effectif
$D \subset S$. Une autre application de la remarque \ref{flat-remark-successive-blowups},
combinée avec le
lemme \ref{divisors-lemma-characterize-effective-Cartier-divisor} de Diviseurs,
montre que nous pouvons supposer $S = \Spec(A)$ et $D = \Spec(A/(a))$
pour un certain non-diviseur de zéro $a \in A$.

\medskip\noindent
Choisissons un $i$ et $s \in S_i$. Le lemme \ref{flat-lemma-push-ideal}
implique que nous pouvons trouver un voisinage ouvert $s \in W_i \subset S_i$
et un idéal quasi-cohérent de type fini $\mathcal{I} \subset \mathcal{O}_S$
tel que $\mathcal{I} \cdot \mathcal{O}_{W_i} = \mathcal{J}_i \mathcal{J}'_i$
pour un certain idéal quasi-cohérent de type fini
$\mathcal{J}'_i \subset \mathcal{O}_{W_i}$
et tel que $V(\mathcal{I}) \subset V(a) = S \setminus U$.
Puisque $S_i$ est quasi-compact, nous pouvons remplacer $S_i$ par une famille finie
$W_1, \ldots, W_n$ de tels ouverts et supposer que, pour tout $i$, il existe
un faisceau quasi-cohérent d'idéaux $\mathcal{I}_i \subset \mathcal{O}_S$ tel
que $\mathcal{I}_i \cdot \mathcal{O}_{S_i} = \mathcal{J}_i \mathcal{J}'_i$
pour un certain idéal quasi-cohérent de type fini
$\mathcal{J}'_i \subset \mathcal{O}_{S_i}$.
Comme dans la discussion du premier paragraphe de la démonstration, considérons
l'éclatement $S'$ de $S$ le long du produit $\mathcal{I}_1 \ldots \mathcal{I}_n$
(cet éclatement est $U$-admissible par construction). Le changement de base de $S' \to S$
à $S_i$ est l'éclatement le long de
$$
\mathcal{J}_i \cdot
\mathcal{J}'_i \mathcal{I}_1 \ldots \hat{\mathcal{I}_i} \ldots \mathcal{I}_n
$$
qui se factorise par l'éclatement donné $S_i^* \to S_i$
(lemme \ref{divisors-lemma-blowing-up-two-ideals} de Diviseurs). Dans la notation
du diagramme ci-dessus, cela signifie que $S_i^{\prime *} = S_i'$. Par conséquent,
après avoir remplacé $S$ par $S'$, nous arrivons à la situation où
$j_i^*\mathcal{F}$ est plat sur $S_i$. Par conséquent, $j_i^*\mathcal{F}$ est plat
sur $S$, voir le
lemme \ref{flat-lemma-etale-flat-up-down}.
D'après le lemme \ref{morphisms-lemma-flat-permanence} de Morphismes,
nous voyons que $\mathcal{F}$ est plat sur $S$.
\end{proof}

\begin{theorem}
\label{flat-theorem-flatten-module}
Soit $S$ un schéma quasi-compact et quasi-séparé.
Soit $X$ un schéma sur $S$.
Soit $\mathcal{F}$ un module quasi-cohérent sur $X$.
Soit $U \subset S$ un ouvert quasi-compact. Supposons que
\begin{enumerate}
\item $X$ est quasi-compact,
\item $X$ est localement de présentation finie sur $S$,
\item $\mathcal{F}$ est un module de type fini,
\item $\mathcal{F}_U$ est de présentation finie, et
\item $\mathcal{F}_U$ est plat sur $U$.
\end{enumerate}
Alors il existe un éclatement $U$-admissible $S' \to S$ tel que le
transformé strict $\mathcal{F}'$ de $\mathcal{F}$ soit un
$\mathcal{O}_{X \times_S S'}$-module de présentation finie et
plat sur $S'$.
\end{theorem}

\begin{proof}
Montrons d'abord que nous pouvons trouver un éclatement $U$-admissible tel que le
transformé strict soit plat. La question est locale pour la topologie \'etale sur la source
et sur le but, voir le
lemme \ref{flat-lemma-flatten-module-etale-localize} pour un énoncé précis.
En particulier, nous pouvons supposer que $S = \Spec(R)$ et $X = \Spec(A)$
sont affines. Pour $s \in S$, écrivons $\mathcal{F}_s = \mathcal{F}|_{X_s}$
(image inverse de $\mathcal{F}$ sur la fibre). Comme $X \to S$ est de type fini,
$d = \max_{s \in S} \dim(\text{Supp}(\mathcal{F}_s))$
est un entier. Nous allons raisonner par récurrence sur $d$.

\medskip\noindent
Soit $x \in X$ un point de $X$ situé au-dessus de $s \in S$ avec
$\dim_x(\text{Supp}(\mathcal{F}_s)) = d$.
Appliquons le lemme \ref{flat-lemma-elementary-devissage} pour obtenir
$g : X' \to X$, $e : S' \to S$, $i : Z' \to X'$, et $\pi : Z' \to Y'$.
Observons que $Y' \to S'$ est
un morphisme lisse de schémas affines à fibres géométriquement irréductibles
de dimension $d$. Puisque le problème est local pour la topologie \'etale, il suffit
de démontrer le théorème pour $g^*\mathcal{F}/X'/S'$. Puisque $i : Z' \to X'$
est une immersion fermée de présentation finie (et puisque le transformé strict
commute à l'image directe affine, voir le
lemme \ref{divisors-lemma-strict-transform-affine} de Diviseurs),
il suffit de démontrer le résultat de platification pour $\mathcal{G}$.
Puisque $\pi$ est fini (donc également affine), il suffit de démontrer le
résultat de platification pour $\pi_*\mathcal{G}/Y'/S'$. Nous pouvons donc supposer
que $X \to S$ est un morphisme lisse de schémas affines à fibres
géométriquement irréductibles de dimension $d$.

\medskip\noindent
Ensuite, nous appliquons un éclatement comme dans le lemme \ref{flat-lemma-flatten-module-pre}.
Nous arrivons ainsi à la situation où il existe un
ouvert $V \subset X$ surjectif sur $S$ tel que $\mathcal{F}|_V$
soit localement libre de rang fini. Soit $\xi \in X$ le point générique de $X_s$. Soit
$r = \dim_{\kappa(\xi)} \mathcal{F}_\xi \otimes \kappa(\xi)$.
Choisissons un morphisme $\alpha : \mathcal{O}_X^{\oplus r} \to \mathcal{F}$
qui induit un isomorphisme
$\kappa(\xi)^{\oplus r} \to \mathcal{F}_\xi \otimes \kappa(\xi)$.
Puisque $\mathcal{F}$ est localement libre sur $V$, nous trouvons un voisinage ouvert
$W$ de $\xi$ où $\alpha$ est un isomorphisme. Restreignons $S$ à un
voisinage ouvert affine de $s$ tel que $W \to S$ soit surjectif. Disons que $\mathcal{F}$
est le module quasi-cohérent associé au $A$-module $N$. Puisque
$\mathcal{F}$ est plat sur $S$ en tous les points génériques des fibres
(en fait en tous les points de $W$), nous voyons que
$$
\alpha_\mathfrak p : A_\mathfrak p^{\oplus r} \to N_\mathfrak p
$$
est universellement injectif pour tout idéal premier $\mathfrak p$ de $R$, voir le
lemme \ref{flat-lemma-induction-step}. Par conséquent, $\alpha$ est universellement injectif,
voir le lemme \ref{algebra-lemma-universally-injective-check-stalks} d'Algèbre.
Posons $\mathcal{H} = \Coker(\alpha)$.
D'après le lemme \ref{divisors-lemma-strict-transform-universally-injective} de Diviseurs,
nous voyons que, pour tout éclatement $U$-admissible $S' \to S$,
les transformés stricts $\mathcal{F}'$ et $\mathcal{H}'$
s'insèrent dans une suite exacte
$$
0 \to \mathcal{O}_{X \times_S S'}^{\oplus r} \to \mathcal{F}'
\to \mathcal{H}' \to 0
$$
Le lemme \ref{flat-lemma-induction-step} montre donc aussi que $\mathcal{F}'$
est plat en un point $x'$ si et seulement si
$\mathcal{H}'$ est plat en ce point. En particulier, $\mathcal{H}_U$ est
plat sur $U$ et $\mathcal{H}_U$ est un module de présentation finie.
L'hypothèse de récurrence appliquée à $\mathcal{H}$ montre l'existence
d'un éclatement $U$-admissible tel que le transformé strict
$\mathcal{H}'$ soit plat, comme souhaité.

\medskip\noindent
Pour achever la démonstration du théorème, il nous reste à montrer que $\mathcal{F}'$
est un module de présentation finie (après éventuellement un autre
éclatement $U$-admissible). Cela résulte du
lemme \ref{flat-lemma-flat-finite-type-finitely-presented-over-dense-open},
car nous pouvons supposer que $U \subset S$ est schématiquement dense (voir le
troisième paragraphe de la remarque \ref{flat-remark-successive-blowups}).
Ceci achève la démonstration du théorème.
\end{proof}





\section{Applications}
\label{flat-section-applications}

\noindent
Dans cette section, nous appliquons certains des résultats ci-dessus.

\begin{lemma}
\label{flat-lemma-flat-after-blowing-up}
Soit $S$ un schéma quasi-compact et quasi-séparé.
Soit $X$ un schéma sur $S$.
Soit $U \subset S$ un ouvert quasi-compact.
Supposons que
\begin{enumerate}
\item $X \to S$ soit de type fini et quasi-séparé, et que
\item $X_U \to U$ soit plat et localement de présentation finie.
\end{enumerate}
Alors il existe un éclatement $U$-admissible $S' \to S$ tel que
le transformé strict de $X$ soit plat et de présentation finie
sur $S'$.
\end{lemma}

\begin{proof}
Puisque $X \to S$ est quasi-compact et quasi-séparé par hypothèse,
le transformé strict de $X$ relativement à un éclatement $S' \to S$
est lui aussi quasi-compact et quasi-séparé. Ainsi, pour démontrer le lemme,
il suffit de trouver un éclatement $U$-admissible tel que le transformé
strict soit plat et localement de présentation finie.
Soit $X = W_1 \cup \ldots \cup W_n$ un recouvrement ouvert affine fini.
Si nous pouvons trouver un éclatement $U$-admissible $S_i \to S$ tel que le
transformé strict de $W_i$ soit plat et localement de présentation finie,
alors il existe un éclatement $U$-admissible $S' \to S$ dominant
tous les $S_i \to S$ et qui convient (voir
Diviseurs, lemme \ref{divisors-lemma-dominate-admissible-blowups} ;
voir aussi la remarque \ref{flat-remark-successive-blowups}).
Nous pouvons donc supposer que $X$ est affine.

\medskip\noindent
Supposons que $X$ soit affine. D'après
Morphismes, lemme \ref{morphisms-lemma-quasi-affine-finite-type-over-S},
nous pouvons choisir une immersion $j : X \to \mathbf{A}^n_S$ sur $S$.
Soit $V \subset \mathbf{A}^n_S$ un sous-schéma ouvert quasi-compact
tel que $j$ induise une immersion fermée $i : X \to V$ sur $S$. Appliquons
le théorème \ref{flat-theorem-flatten-module}
à $V \to S$ et au module quasi-cohérent $i_*\mathcal{O}_X$
pour obtenir un éclatement $U$-admissible $S' \to S$ tel que le transformé
strict de $i_*\mathcal{O}_X$ soit plat sur $S'$ et de présentation finie
comme $\mathcal{O}_{V \times_S S'}$-module. Soit $X'$ le transformé strict
de $X$ relativement à $S' \to S$. Soit $i' : X' \to V \times_S S'$
le morphisme induit.
Comme la formation du transformé strict commute à l'image directe par les
morphismes affines (Diviseurs, lemme \ref{divisors-lemma-strict-transform-affine}),
nous voyons que $i'_*\mathcal{O}_{X'}$ est plat sur $S$ et de
présentation finie comme $\mathcal{O}_{V \times_S S'}$-module.
Cela implique le lemme.
\end{proof}

\begin{lemma}
\label{flat-lemma-finite-after-blowing-up}
Soit $S$ un schéma quasi-compact et quasi-séparé.
Soit $X$ un schéma sur $S$.
Soit $U \subset S$ un ouvert quasi-compact.
Supposons que
\begin{enumerate}
\item $X \to S$ soit propre, et que
\item $X_U \to U$ soit fini localement libre.
\end{enumerate}
Alors il existe un éclatement $U$-admissible $S' \to S$ tel que
le transformé strict de $X$ soit fini localement libre sur $S'$.
\end{lemma}

\begin{proof}
D'après le lemme \ref{flat-lemma-flat-after-blowing-up}, nous pouvons supposer que
$X \to S$ est plat et de présentation finie. Après avoir remplacé
$S$ par un éclatement $U$-admissible si nécessaire, nous pouvons supposer
que $U \subset S$ est schématiquement dense. Alors $f$ est
fini d'après le lemme \ref{flat-lemma-proper-flat-finite-over-dense-open}.
Ainsi $f$ est fini localement libre d'après
Morphismes, lemme \ref{morphisms-lemma-finite-flat}.
\end{proof}

\begin{lemma}
\label{flat-lemma-zariski-after-blowup}
Soit $\varphi : X \to S$ un morphisme séparé de type fini, avec
$S$ quasi-compact et quasi-séparé. Soit $U \subset S$ un
ouvert quasi-compact tel que $\varphi^{-1}U \to U$ soit un isomorphisme.
Alors il existe un éclatement $U$-admissible $S' \to S$ tel que
le transformé strict $X'$ de $X$ soit isomorphe à un sous-schéma ouvert
de $S'$.
\end{lemma}

\begin{proof}
La discussion de la remarque \ref{flat-remark-successive-blowups} s'applique.
Nous pouvons donc effectuer un premier éclatement $U$-admissible et supposer que le complémentaire
$S \setminus U$ est le support d'un diviseur de Cartier effectif $D$.
En particulier, $U$ est schématiquement dense dans $S$.
Ensuite, nous effectuons un autre éclatement $U$-admissible pour nous ramener au cas où
$X \to S$ est plat et de présentation finie, voir
le lemme \ref{flat-lemma-flat-after-blowing-up}.
Dans ce cas, le résultat découle du lemme \ref{flat-lemma-zariski}.
\end{proof}

\noindent
Le lemme suivant affirme qu'une modification propre peut être dominée
par un éclatement.

\begin{lemma}
\label{flat-lemma-dominate-modification-by-blowup}
Soit $\varphi : X \to S$ un morphisme propre, où
$S$ est quasi-compact et quasi-séparé. Soit $U \subset S$ un
ouvert quasi-compact tel que $\varphi^{-1}U \to U$ soit un isomorphisme.
Alors il existe un éclatement $U$-admissible $S' \to S$
qui domine $X$, c'est-à-dire tel qu'il existe une factorisation
$S' \to X \to S$ du morphisme d'éclatement.
\end{lemma}

\begin{proof}
La discussion de la remarque \ref{flat-remark-successive-blowups} s'applique.
Nous pouvons donc effectuer un premier éclatement $U$-admissible et supposer que le complémentaire
$S \setminus U$ est le support d'un diviseur de Cartier effectif $D$.
En particulier, $U$ est schématiquement dense dans $S$.
Choisissons un autre éclatement $U$-admissible $S' \to S$ tel que le transformé
strict $X'$ de $X$ soit un sous-schéma ouvert de $S'$, voir
le lemme \ref{flat-lemma-zariski-after-blowup}.
Puisque $X' \to S'$ est propre et que
$U \subset S'$ est dense, nous voyons que $X' = S'$. Quelques détails sont omis.
\end{proof}

\begin{lemma}
\label{flat-lemma-get-section-after-blowup}
Soit $S$ un schéma. Soient $U \subset W \subset S$ des sous-schémas ouverts.
Soit $f : X \to W$ un morphisme et soit $s : U \to X$ un
morphisme tel que $f \circ s = \text{id}_U$. Supposons que
\begin{enumerate}
\item $f$ soit propre,
\item $S$ soit quasi-compact et quasi-séparé, et que
\item $U$ et $W$ soient quasi-compacts.
\end{enumerate}
Alors il existe un éclatement $U$-admissible $b : S' \to S$ et un morphisme
$s' : b^{-1}(W) \to X$ prolongeant $s$, avec $f \circ s' = b|_{b^{-1}(W)}$.
\end{lemma}

\begin{proof}
Nous pouvons, et nous allons, remplacer $X$ par l'image schématique de $s$.
Alors $X \to W$ est un isomorphisme au-dessus de $U$, voir
Morphismes, lemme
\ref{morphisms-lemma-scheme-theoretic-image-of-partial-section}.
D'après le lemme \ref{flat-lemma-dominate-modification-by-blowup},
il existe un éclatement $U$-admissible $W' \to W$ et un
prolongement $W' \to X$ de $s$.
Nous achevons la démonstration en appliquant
Diviseurs, lemme \ref{divisors-lemma-extend-admissible-blowups},
afin de prolonger $W' \to W$ en un éclatement $U$-admissible de $S$.
\end{proof}









\section{Compactifications}
\label{flat-section-compactify}

\noindent
Soit $S$ un schéma quasi-compact et quasi-séparé. Nous dirons qu'un
schéma $X$ sur $S$ {\it admet une compactification sur $S$}
ou {\it est compactifiable sur $S$} s'il existe
une immersion ouverte quasi-compacte $X \to \overline{X}$ dans un schéma
$\overline{X}$ propre sur $S$. Si $X$ admet une compactification sur
$S$, alors $X \to S$ est séparé et de type fini. Un théorème de
Nagata, voir \cite{Lutkebohmert}, \cite{Conrad-Nagata}, \cite{Nagata-1},
\cite{Nagata-2}, \cite{Nagata-3} et \cite{Nagata-4},
affirme que la réciproque est également vraie. Nous démontrerons ce théorème
dans la section suivante, voir le théorème \ref{flat-theorem-nagata}.

\medskip\noindent
Soit $S$ un schéma quasi-compact et quasi-séparé.
Soit $X \to S$ un morphisme de schémas séparé et de type fini. La catégorie
des {\it compactifications de $X$ sur $S$} est la catégorie définie
comme suit :
\begin{enumerate}
\item Ses objets sont les immersions ouvertes $j : X \to \overline{X}$ sur $S$ telles que
$\overline{X} \to S$ soit propre.
\item Ses morphismes $(j' : X \to \overline{X}') \to (j : X \to \overline{X})$
sont les morphismes $f : \overline{X}' \to \overline{X}$ de schémas sur $S$
tels que $f \circ j' = j$.
\end{enumerate}
Si $j : X \to \overline{X}$ est une compactification, alors $j$ est une
immersion ouverte quasi-compacte, voir
Schémas, remarque \ref{schemes-remark-quasi-compact-and-quasi-separated}.

\medskip\noindent
{\bf Attention.} Nous ne supposons {\it pas} que les compactifications $j : X \to \overline{X}$
aient une image dense. Par conséquent, si $f : \overline{X}' \to \overline{X}$
est un morphisme de compactifications, il se peut que
$f^{-1}(j(X)) = j'(X)$ ne soit pas vérifiée.

\begin{lemma}
\label{flat-lemma-compactifications-cofiltered}
Soit $S$ un schéma quasi-compact et quasi-séparé.
Soit $X$ un schéma compactifiable sur $S$.
\begin{enumerate}
\item[(a)] La catégorie des compactifications de $X$ sur $S$ est
cofiltrante.
\item[(b)] La sous-catégorie pleine formée des compactifications
$j : X \to \overline{X}$ telles que $j(X)$ soit dense et
schématiquement dense dans $\overline{X}$ est initiale
(Catégories, définition \ref{categories-definition-initial}).
\item[(c)] Si $f : \overline{X}' \to \overline{X}$ est un morphisme
de compactifications de $X$ tel que $j'(X)$ soit dense dans $\overline{X}'$,
alors $f^{-1}(j(X)) = j'(X)$.
\end{enumerate}
\end{lemma}

\begin{proof}
Pour démontrer (a), nous devons vérifier les conditions (1), (2), (3) de
Catégories, définition \ref{categories-definition-codirected}.
La condition (1) est satisfaite précisément parce que nous avons supposé $X$
compactifiable.
Soient $j_i : X \to \overline{X}_i$, $i = 1, 2$, deux compactifications.
Nous pouvons alors considérer l'image schématique $\overline{X}$
de $(j_1, j_2) : X \to \overline{X}_1 \times_S \overline{X}_2$.
Cela détermine une troisième compactification $j : X \to \overline{X}$
qui domine les deux $j_i$ :
$$
\xymatrix{
(X, \overline{X}_1) & (X, \overline{X}) \ar[l] \ar[r] & (X, \overline{X}_2)
}
$$
Ainsi (2) est satisfaite. Soient $f_1, f_2 : \overline{X}_1 \to \overline{X}_2$
deux morphismes entre les compactifications
$j_i : X \to \overline{X}_i$, $i = 1, 2$.
Soit $\overline{X} \subset \overline{X}_1$ l'égalisateur de
$f_1$ et $f_2$. Comme $\overline{X}_2 \to S$ est séparé, nous voyons
que $\overline{X}$ est un sous-schéma fermé de $\overline{X}_1$
et est donc propre sur $S$. De plus, nous obtenons une
immersion ouverte $X \to \overline{X}$ puisque $f_1|_X = f_2|_X = \text{id}_X$.
Le morphisme $(X \to \overline{X}) \to (j_1 : X \to \overline{X}_1)$
donné par l'immersion fermée $\overline{X} \to \overline{X}_1$
égalise $f_1$ et $f_2$, ce qui démontre la condition (3).

\medskip\noindent
Démonstration de (b). Soit $j : X \to \overline{X}$ une compactification.
Si $\overline{X}'$ désigne l'adhérence schématique de $X$
dans $\overline{X}$, alors $X$ est dense et schématiquement dense
dans $\overline{X}'$ d'après
Morphismes, lemme \ref{morphisms-lemma-quasi-compact-immersion}.
Cela démontre la première condition de
Catégories, définition \ref{categories-definition-initial}.
Puisque nous avons déjà montré que la catégorie des compactifications
de $X$ est cofiltrante, la seconde condition de
Catégories, définition \ref{categories-definition-initial},
découle de la première (nous omettons la résolution de cet
exercice catégorique).

\medskip\noindent
Démonstration de (c). Après avoir remplacé $\overline{X}'$ par l'adhérence schématique
de $j'(X)$ (ce qui ne modifie pas l'espace topologique sous-jacent),
cela découle de Morphismes, lemme
\ref{morphisms-lemma-scheme-theoretic-image-of-partial-section}.
\end{proof}

\noindent
Nous pouvons également considérer la catégorie de toutes les compactifications (pour $X$ variable).
Il se trouve que cette catégorie, localisée par rapport à l'ensemble des morphismes
qui induisent un isomorphisme sur l'intérieur,
est équivalente à la catégorie des schémas compactifiables sur $S$.

\begin{lemma}
\label{flat-lemma-compactifyable}
Soit $S$ un schéma quasi-compact et quasi-séparé. Soit $f : X \to Y$
un morphisme de schémas sur $S$, avec $Y$ séparé et de type fini
sur $S$ et $X$ compactifiable sur $S$. Alors $X$ admet une compactification
sur $Y$.
\end{lemma}

\begin{proof}
Soit $j : X \to \overline{X}$ une compactification de $X$ sur $S$.
Nous définissons alors $\overline{X}'$ comme
l'image schématique de $(j, f) : X \to \overline{X} \times_S Y$.
Le morphisme $\overline{X}' \to Y$ est propre, car
$\overline{X} \times_S Y \to Y$ est propre comme changement de base de
$\overline{X} \to S$. D'autre part, puisque $Y$ est séparé
sur $S$, le morphisme $(1, f) : X \to X \times_S Y$ est une
immersion fermée (Schémas, lemme \ref{schemes-lemma-semi-diagonal})
et donc $X \to \overline{X}'$ est une immersion ouverte d'après Morphismes, lemme
\ref{morphisms-lemma-scheme-theoretic-image-of-partial-section}, appliqué
à la ``section partielle'' $s = (j, f)$ de la projection
$\overline{X} \times_S Y \to \overline{X}$.
\end{proof}

\noindent
Soit $S$ un schéma quasi-compact et quasi-séparé.
Nous définissons la {\it catégorie des compactifications} comme la catégorie
dont les objets sont les paires $(X, \overline{X})$ où $\overline{X}$
est un schéma propre sur $S$ et $X \subset \overline{X}$ est un
ouvert quasi-compact, et dont les morphismes
sont les diagrammes commutatifs
$$
\xymatrix{
X \ar[d] \ar[r]_f & Y \ar[d] \\
\overline{X} \ar[r]^{\overline{f}} & \overline{Y}
}
$$
de morphismes de schémas sur $S$.

\begin{lemma}
\label{flat-lemma-right-multiplicative-system}
Soit $S$ un schéma quasi-compact et quasi-séparé.
La famille des morphismes
$(u, \overline{u}) : (X', \overline{X}') \to (X, \overline{X})$
tels que $u$ soit un isomorphisme forme un système multiplicatif à droite
(Catégories, définition \ref{categories-definition-multiplicative-system})
de flèches dans la catégorie des compactifications.
\end{lemma}

\begin{proof}
L'axiome RMS1 se vérifie immédiatement. Vérifions que RMS2 est satisfait.
Supposons donné un diagramme
$$
\xymatrix{
& (X', \overline{X}') \ar[d]_{(u, \overline{u})} \\
(Y, \overline{Y}) \ar[r]^{(f, \overline{f})} & (X, \overline{X})
}
$$
où $u : X' \to X$ est un isomorphisme. Posons alors $Y' = Y \times_X X'$
et notons $v : Y' \to Y$ la projection (qui est un isomorphisme). Posons aussi
$\overline{Y}' = \overline{Y} \times_{\overline{X}} \overline{X}'$
et notons $\overline{v} : \overline{Y}' \to \overline{Y}$ la projection.
Il est clair que $Y' \to \overline{Y}'$ est une immersion ouverte.
Le diagramme
$$
\xymatrix{
(Y', \overline{Y}') \ar[r]_{(g, \overline{g})} \ar[d]_{(v, \overline{v})} &
(X', \overline{X}') \ar[d]_{(u, \overline{u})} \\
(Y, \overline{Y}) \ar[r]^{(f, \overline{f})} & (X, \overline{X})
}
$$
montre que l'axiome RMS2 est satisfait.

\medskip\noindent
Vérifions que RMS3 est satisfait. Supposons donnés deux morphismes
$(f, \overline{f}), (g, \overline{g}) :
(X, \overline{X}) \to (Y, \overline{Y})$
de compactifications et un morphisme
$(v, \overline{v}) : (Y, \overline{Y}) \to (Y', \overline{Y}')$
tel que $v$ soit un isomorphisme et que
$(v, \overline{v}) \circ (f, \overline{f}) =
(v, \overline{v}) \circ (g, \overline{g})$. Alors $f = g$.
Ainsi, si nous notons $\overline{X}' \subset \overline{X}$
l'égalisateur de $\overline{f}$ et $\overline{g}$,
alors $(u, \overline{u}) : (X, \overline{X}') \to (X, \overline{X})$
sera un morphisme de la catégorie des compactifications
tel que $(f, \overline{f}) \circ (u, \overline{u}) =
(g, \overline{g}) \circ (u, \overline{u})$, comme voulu.
\end{proof}

\begin{lemma}
\label{flat-lemma-invert-right-multiplicative-system}
Soit $S$ un schéma quasi-compact et quasi-séparé.
Le foncteur $(X, \overline{X}) \mapsto X$ définit une
équivalence de la catégorie des compactifications localisée
(Catégories, lemme \ref{categories-lemma-right-localization})
par rapport au système
multiplicatif à droite du lemme \ref{flat-lemma-right-multiplicative-system}
vers la catégorie des schémas compactifiables sur $S$.
\end{lemma}

\begin{proof}
Notons $\mathcal{C}$ la catégorie des compactifications et
notons $Q : \mathcal{C} \to \mathcal{C}'$ le foncteur de
localisation de Catégories, lemme
\ref{categories-lemma-properties-right-localization}.
Notons $\mathcal{D}$ la catégorie des schémas compactifiables
sur $S$. Il est clair, d'après le lemme que nous venons de citer et notre
choix du système multiplicatif, que nous
obtenons un foncteur $\mathcal{C}' \to \mathcal{D}$.
Ce foncteur est manifestement essentiellement surjectif.
Si $f : X \to Y$ est un morphisme de schémas
compactifiables, choisissons une immersion ouverte $Y \to \overline{Y}$
dans un schéma propre sur $S$, puis choisissons un plongement
$X \to \overline{X}$ dans un schéma $\overline{X}$ propre sur
$\overline{Y}$ (ce qui est possible d'après le lemme \ref{flat-lemma-compactifyable}
appliqué à $X \to \overline{Y}$). Nous obtenons ainsi un morphisme
$(X, \overline{X}) \to (Y, \overline{Y})$ de compactifications
qui induit le morphisme donné $X \to Y$.
Enfin, supposons donnés deux morphismes de la
catégorie localisée ayant même source et même but, disons
$$
a = ((f, \overline{f}) : (X', \overline{X}') \to (Y, \overline{Y}),
(u, \overline{u}) : (X', \overline{X}') \to (X, \overline{X}))
$$
et
$$
b = ((g, \overline{g}) : (X'', \overline{X}'') \to (Y, \overline{Y}),
(v, \overline{v}) : (X'', \overline{X}'') \to (X, \overline{X}))
$$
qui induisent le même morphisme $X \to Y$ sur $S$, autrement dit
$f \circ u^{-1} = g \circ v^{-1}$. D'après
Catégories, lemme \ref{categories-lemma-morphisms-right-localization},
nous pouvons supposer que $(X', \overline{X}') = (X'', \overline{X}'')$
et que $(u, \overline{u}) = (v, \overline{v})$. Dans ce cas, nous
pouvons considérer l'égalisateur $\overline{X}''' \subset \overline{X}'$
de $\overline{f}$ et $\overline{g}$. Le morphisme
$(w, \overline{w}) : (X', \overline{X}''') \to (X', \overline{X}')$ appartient
au système multiplicatif et nous voyons que $a = b$ dans la catégorie
localisée en précomposant avec $(w, \overline{w})$.
\end{proof}






\section{Compactification de Nagata}
\label{flat-section-compactifications}

\noindent
Dans cette section, nous démontrons le théorème annoncé à la
section \ref{flat-section-compactify}.

\begin{lemma}
\label{flat-lemma-check-separated}
Soit $X \to S$ un morphisme de schémas. Si $X = U \cup V$
est un recouvrement ouvert tel que $U \to S$ et $V \to S$ soient séparés
et que $U \cap V \to U \times_S V$ soit fermé, alors
$X \to S$ est séparé.
\end{lemma}

\begin{proof}
Omis. Indication : vérifier que $\Delta : X \to X \times_S X$ est
fermé en utilisant le recouvrement ouvert de $X \times_S X$ donné par
$U \times_S U$, $U \times_S V$, $V \times_S U$ et $V \times_S V$.
\end{proof}

\begin{lemma}
\label{flat-lemma-separate-disjoint-locally-closed-by-blowing-up}
Soit $X$ un schéma quasi-compact et quasi-séparé.
Soit $U \subset X$ un ouvert quasi-compact.
\begin{enumerate}
\item Si $Z_1, Z_2 \subset X$ sont des sous-schémas fermés de
présentation finie tels que $Z_1 \cap Z_2 \cap U = \emptyset$, alors
il existe un éclatement $U$-admissible $X' \to X$
tel que les transformés stricts de $Z_1$ et $Z_2$ soient disjoints.
\item Si $T_1, T_2 \subset U$ sont des parties disjointes fermées pour la topologie constructible, alors
il existe un éclatement $U$-admissible $X' \to X$ tel que les adhérences de
$T_1$ et $T_2$ soient disjointes.
\end{enumerate}
\end{lemma}

\begin{proof}
Démonstration de (1). L'hypothèse que $Z_i \to X$ est de présentation finie
signifie que le faisceau quasi-cohérent d'idéaux $\mathcal{I}_i$ de $Z_i$
est de type fini, voir
Morphismes, lemme \ref{morphisms-lemma-closed-immersion-finite-presentation}.
Notons $Z \subset X$ le sous-schéma fermé
découpé par le produit $\mathcal{I}_1 \mathcal{I}_2$.
Observons que $Z \cap U$ est la réunion disjointe
de $Z_1 \cap U$ et $Z_2 \cap U$. D'après Diviseurs, lemme
\ref{divisors-lemma-separate-disjoint-opens-by-blowing-up},
il existe un éclatement $U \cap Z$-admissible $Z' \to Z$ tel que
les transformés stricts de $Z_1$ et $Z_2$ soient disjoints.
Notons $Y \subset Z$ le centre de cet éclatement.
Alors $Y \to X$ est une immersion fermée de présentation finie comme composée
de $Y \to Z$ et $Z \to X$ (Diviseurs, définition
\ref{divisors-definition-admissible-blowup}
et Morphismes, lemme \ref{morphisms-lemma-composition-finite-presentation}).
Ainsi, l'éclatement $X' \to X$ de centre $Y$ est un éclatement $U$-admissible.
D'après les propriétés générales des transformés stricts, les
transformés stricts de $Z_1, Z_2$ relativement à $X' \to X$
sont les mêmes que les transformés stricts de $Z_1, Z_2$ relativement
à $Z' \to Z$, voir
Diviseurs, lemme \ref{divisors-lemma-strict-transform}.
Cela démontre (1).

\medskip\noindent
Démonstration de (2). D'après Propriétés, lemme
\ref{properties-lemma-quasi-coherent-finite-type-ideals},
il existe un faisceau quasi-cohérent d'idéaux de type fini
$\mathcal{J}_i \subset \mathcal{O}_U$ tel que
$T_i = V(\mathcal{J}_i)$ (ensemblistement).
D'après Propriétés, lemme \ref{properties-lemma-extend},
il existe un faisceau quasi-cohérent de type fini
d'idéaux $\mathcal{I}_i \subset \mathcal{O}_X$
dont la restriction à $U$ est $\mathcal{J}_i$.
Appliquons le résultat de la partie (1) aux
sous-schémas fermés $Z_i = V(\mathcal{I}_i)$ pour conclure.
\end{proof}

\begin{lemma}
\label{flat-lemma-blowup-iso-along}
Soit $f : X \to Y$ un morphisme propre de schémas quasi-compacts et
quasi-séparés. Soit $V \subset Y$ un ouvert quasi-compact,
et soit $U = f^{-1}(V)$. Soit $T \subset V$ une partie fermée telle que
$f|_U : U \to V$ soit un isomorphisme au-dessus d'un voisinage ouvert de $T$
dans $V$. Alors il existe un éclatement $V$-admissible $Y' \to Y$
tel que le transformé strict $f' : X' \to Y'$ de $f$
soit un isomorphisme au-dessus d'un voisinage ouvert de l'adhérence
de $T$ dans $Y'$.
\end{lemma}

\begin{proof}
Soit $T' \subset V$ le complémentaire du plus grand ouvert au-dessus duquel
$f|_U$ est un isomorphisme. Alors $T', T$ sont fermés dans $V$ et
$T \cap T' = \emptyset$. Puisque $V$ est un espace topologique
spectral, nous pouvons trouver des parties $T_c, T'_c$ fermées pour la topologie constructible
telles que $T \subset T_c$, $T' \subset T'_c$ et
$T_c \cap T'_c = \emptyset$ (choisir un ouvert quasi-compact
$W$ de $V$ contenant $T'$ et ne rencontrant pas $T$,
et poser $T_c = V \setminus W$, puis choisir un ouvert quasi-compact
$W'$ de $V$ contenant $T_c$ et ne rencontrant pas $T'$,
et poser $T'_c = V \setminus W'$).
D'après le lemme \ref{flat-lemma-separate-disjoint-locally-closed-by-blowing-up},
nous pouvons, après avoir remplacé $Y$ par un éclatement $V$-admissible,
supposer que $T_c$ et $T'_c$ ont des adhérences disjointes dans $Y$.
Posons $Y_0 = Y \setminus \overline{T}'_c$, $V_0 = V \setminus T'_c$,
$U_0 = U \times_V V_0$ et $X_0 = X \times_Y Y_0$.
Puisque $U_0 \to V_0$ est un isomorphisme, nous pouvons trouver un
éclatement $V_0$-admissible $Y'_0 \to Y_0$ tel que le
transformé strict $X'_0$ de $X_0$ s'envoie isomorphiquement sur $Y'_0$, voir
le lemme \ref{flat-lemma-zariski-after-blowup}.
D'après Diviseurs, lemme \ref{divisors-lemma-extend-admissible-blowups},
il existe un éclatement $V$-admissible $Y' \to Y$ dont la restriction
à $Y_0$ est $Y'_0 \to Y_0$. Si $f' : X' \to Y'$ désigne le
transformé strict de $f$, alors nous obtenons le résultat voulu, car
$f'$ se restreint en un isomorphisme au-dessus de $Y'_0$.
\end{proof}

\begin{lemma}
\label{flat-lemma-find-common-blowups}
Soit $S$ un schéma quasi-compact et quasi-séparé.
Soient $U \to X_1$ et $U \to X_2$ des immersions ouvertes
de schémas sur $S$, et supposons que $U$, $X_1$, $X_2$ soient de type
fini et séparés sur $S$. Alors il existe un diagramme commutatif
$$
\xymatrix{
X_1' \ar[d] \ar[r] & X & X_2' \ar[l] \ar[d] \\
X_1 & U \ar[l] \ar[lu] \ar[u] \ar[ru] \ar[r] & X_2
}
$$
de schémas sur $S$, où $X_i' \to X_i$ est un éclatement $U$-admissible,
$X_i' \to X$ est une immersion ouverte et $X$ est séparé et de
type fini sur $S$.
\end{lemma}

\begin{proof}
Tout au long de la démonstration, tous les schémas seront séparés et de type fini sur $S$.
Cela implique en particulier que ces schémas sont quasi-compacts et quasi-séparés,
et que les morphismes entre eux sont quasi-compacts et séparés.
Voir Schémas, sections \ref{schemes-section-quasi-compact} et
\ref{schemes-section-separation-axioms}.
Nous utiliserons le fait que, si $U \to W$ est une immersion de tels schémas sur $S$,
alors l'image schématique $Z$ de $U$ dans $W$ est un sous-schéma fermé
de $W$, que $U \to Z$ est une immersion ouverte, que $U \subset Z$ est
schématiquement dense et que $U \subset Z$ est topologiquement dense. Voir
Morphismes, lemme
\ref{morphisms-lemma-quasi-compact-immersion}.

\medskip\noindent
Soit $X_{12} \subset X_1 \times_S X_2$ l'image schématique
de $U \to X_1 \times_S X_2$. Les projections $p_i : X_{12} \to X_i$
induisent des isomorphismes $p_i^{-1}(U) \to U$ d'après
Morphismes, lemme
\ref{morphisms-lemma-scheme-theoretic-image-of-partial-section}.
Choisissons un éclatement $U$-admissible $X_i^i \to X_i$ tel que
le transformé strict $X_{12}^i$ de $X_{12}$ soit isomorphe à un
sous-schéma ouvert de $X_i^i$, voir
le lemme \ref{flat-lemma-zariski-after-blowup}.
Soit $\mathcal{I}_i \subset \mathcal{O}_{X_i}$ le faisceau quasi-cohérent
d'idéaux de type fini correspondant.
Rappelons que $X_{12}^i \to X_{12}$ est l'éclatement de
$p_i^{-1}\mathcal{I}_i \mathcal{O}_{X_{12}}$, voir
Diviseurs, lemme \ref{divisors-lemma-strict-transform}.
Soit $X_{12}'$ l'éclatement de $X_{12}$ de centre
$p_1^{-1}\mathcal{I}_1 p_2^{-1}\mathcal{I}_2 \mathcal{O}_{X_{12}}$, voir
Diviseurs, lemme \ref{divisors-lemma-blowing-up-two-ideals}
pour la signification de cette construction. Nous obtenons en particulier un diagramme commutatif
$$
\xymatrix{
X_{12}' \ar[d] \ar[r] & X_{12}^2 \ar[d] \\
X_{12}^1 \ar[r] & X_{12}
}
$$
où tous les morphismes sont des éclatements $U$-admissibles.
Puisque $X_{12}^i \subset X_i^i$ est un ouvert, nous pouvons choisir un éclatement $U$-admissible
$X_i' \to X_i^i$ dont la restriction est $X_{12}' \to X_{12}^i$, voir
Diviseurs, lemme \ref{divisors-lemma-extend-admissible-blowups}.
Alors $X_{12}' \subset X_i'$ est un sous-schéma ouvert et le diagramme
$$
\xymatrix{
X_{12}' \ar[d] \ar[r] & X_i' \ar[d] \\
X_{12}^i \ar[r] & X_i^i
}
$$
est commutatif, les flèches verticales étant des éclatements et les flèches horizontales
des immersions ouvertes. Remarquons que $X'_{12} \to X_1' \times_S X_2'$ est
une immersion propre (utiliser que $X'_{12} \to X_{12}$ est propre,
que $X_{12} \to X_1 \times_S X_2$ est fermé, que $X_1' \times_S X_2' \to
X_1 \times_S X_2$ est séparé et appliquer Morphismes, lemme
\ref{morphisms-lemma-image-proper-scheme-closed}).
Ainsi $X'_{12} \to  X_1' \times_S X_2'$ est une immersion fermée.
Il s'ensuit que, si nous définissons $X$ en recollant $X_1'$ et $X_2'$
le long du sous-schéma ouvert commun $X_{12}'$, alors $X \to S$ est de type fini
et séparé (lemme \ref{flat-lemma-check-separated}).
Comme les composées d'éclatements $U$-admissibles sont des éclatements $U$-admissibles
(Diviseurs, lemme \ref{divisors-lemma-composition-admissible-blowups}),
le lemme est démontré.
\end{proof}

\begin{lemma}
\label{flat-lemma-replaced-by-strict-transform}
Soient $X \to S$ et $Y \to S$ des morphismes de schémas.
Soit $U \subset X$ un sous-schéma ouvert.
Soit $V \to X \times_S Y$ un morphisme quasi-compact
dont la composée avec la première projection est à valeurs dans $U$.
Soit $Z \subset X \times_S Y$ l'image schématique de
$V \to X \times_S Y$. Soit $X' \to X$ un éclatement $U$-admissible.
Alors l'image schématique de $V \to X' \times_S Y$ est le
transformé strict de $Z$ relativement à l'éclatement.
\end{lemma}

\begin{proof}
Notons $Z' \to Z$ le transformé strict. Le morphisme $Z' \to X'$
induit un morphisme $Z' \to X' \times_S Y$ qui est une immersion fermée
(car $Z'$ est par définition un sous-schéma fermé de $X' \times_X Z$).
Ainsi, pour achever la démonstration, il suffit de montrer que l'image schématique
$Z''$ de $V \to Z'$ est $Z'$. Observons que $Z'' \subset Z'$
est un sous-schéma fermé tel que $V \to Z'$ se factorise par $Z''$.
Puisque $V \to X \times_S Y$ et $V \to X' \times_S Y$ sont tous deux
quasi-compacts (pour le second, cela découle de Schémas, lemme
\ref{schemes-lemma-quasi-compact-permanence},
et du fait que $X' \times_S Y \to X \times_S Y$ est séparé
comme changement de base d'un morphisme propre), d'après Morphismes, lemme
\ref{morphisms-lemma-quasi-compact-scheme-theoretic-image},
nous voyons que $Z \cap (U \times_S Y) = Z'' \cap (U \times_S Y)$.
Ainsi, le morphisme d'inclusion $Z'' \to Z'$ est un isomorphisme
hors du diviseur exceptionnel $E$ de $Z' \to Z$. Cependant, le
faisceau structural de $Z'$ n'a aucune section non nulle à support
dans $E$ (par définition des transformés stricts), et nous concluons que
la surjection $\mathcal{O}_{Z'} \to \mathcal{O}_{Z''}$
doit être un isomorphisme.
\end{proof}

\begin{lemma}
\label{flat-lemma-compactification-dominates}
Soit $S$ un schéma quasi-compact et quasi-séparé. Soit $U$ un
schéma de type fini et séparé sur $S$. Soit $V \subset U$ un
ouvert quasi-compact. Si $V$ admet une compactification $V \subset Y$
sur $S$, alors il existe un éclatement $V$-admissible $Y' \to Y$ et un
ouvert $V \subset V' \subset Y'$ tel que $V \to U$
se prolonge en un morphisme propre $V' \to U$.
\end{lemma}

\begin{proof}
Considérons l'image schématique $Z \subset Y \times_S U$
du morphisme ``diagonal'' $V \to Y \times_S U$. Si nous remplaçons
$Y$ par un éclatement $V$-admissible, alors $Z$ est remplacé par
le transformé strict relativement à cet éclatement, voir
le lemme \ref{flat-lemma-replaced-by-strict-transform}. Ainsi, d'après
le lemme \ref{flat-lemma-zariski-after-blowup}, nous pouvons supposer que $Z \to Y$
est une immersion ouverte. Si $V' \subset Y$ désigne son image, alors nous
voyons que le morphisme induit $V' \to U$ est propre, car la
projection $Y \times_S U \to U$ est propre et $V' \cong Z$
est un sous-schéma fermé de $Y \times_S U$.
\end{proof}

\noindent
Le lemme suivant n'est formulé que dans le cas noethérien.
La version pour les schémas quasi-compacts et quasi-séparés est
également vraie, mais elle découlera trivialement du théorème principal
de cette section.

\begin{lemma}
\label{flat-lemma-two-compactifications}
Soit $S$ un schéma noethérien. Soit $U$ un schéma de type fini
et séparé sur $S$. Soit $U = U_1 \cup U_2$ avec des ouverts tels que
$U_1$ et $U_2$ admettent des compactifications sur $S$ et que
$U_1 \cap U_2$ soit dense dans $U$. Alors $U$ admet une compactification sur $S$.
\end{lemma}

\begin{proof}
Choisissons une compactification $U_i \subset X_i$ pour $i = 1, 2$. Nous pouvons
supposer que $U_i$ est schématiquement dense dans $X_i$. Nous pouvons supposer qu'il
existe un ouvert $V_i \subset X_i$ et un morphisme propre
$\psi_i : V_i \to U$ prolongeant $\text{id} : U_i \to U_i$, voir
le lemme \ref{flat-lemma-compactification-dominates}. Diagramme :
$$
\xymatrix{
U_i \ar[r] \ar[d] & V_i \ar[r] \ar[dl]^{\psi_i} & X_i \\
U
}
$$
Pour $\{i, j\} = \{1, 2\}$, posons
$Z_i = U \setminus U_j = U_i \setminus (U_1 \cap U_2)$
et
$Z_j = U \setminus U_i = U_j \setminus (U_1 \cap U_2)$.
Ainsi, nous avons
$$
U = U_1 \amalg Z_2 = Z_1 \amalg U_2 = Z_1 \amalg (U_1 \cap U_2) \amalg Z_2
$$
Notons $Z_{i, i} \subset V_i$ l'image réciproque de $Z_i$ par $\psi_i$.
Observons que $\psi_i$ est un isomorphisme au-dessus d'un voisinage ouvert de $Z_i$.
Notons $Z_{i, j} \subset V_i$ l'image réciproque de $Z_j$ par $\psi_i$.
Observons que $\psi_i : Z_{i, j} \to Z_j$ est un morphisme propre.
Puisque $Z_i$ et $Z_j$ sont des parties fermées disjointes de
$U$, nous voyons que $Z_{i, i}$ et $Z_{i, j}$ sont des parties fermées disjointes
de $V_i$.

\medskip\noindent
Notons $\overline{Z}_{i, i}$ et $\overline{Z}_{i, j}$ les adhérences de
$Z_{i, i}$ et $Z_{i, j}$ dans $X_i$. Après avoir remplacé $X_i$ par un
éclatement $V_i$-admissible, nous pouvons supposer que
$\overline{Z}_{i, i}$ et $\overline{Z}_{i, j}$ sont disjointes, voir
le lemme \ref{flat-lemma-separate-disjoint-locally-closed-by-blowing-up}.
Nous supposons que cela vaut à la fois pour $X_1$ et $X_2$.
Observons que cette propriété est préservée si nous remplaçons $X_i$
par un nouvel éclatement $V_i$-admissible.

\medskip\noindent
Posons $V_{12} = V_1 \times_U V_2$. Nous avons une immersion
$V_{12} \to X_1 \times_S X_2$ qui est la composée de l'immersion fermée
$V_{12} = V_1 \times_U V_2 \to V_1 \times_S V_2$
(Schémas, lemme \ref{schemes-lemma-fibre-product-after-map})
et de l'immersion ouverte $V_1 \times_S V_2 \to X_1 \times_S X_2$.
Soit $X_{12} \subset X_1 \times_S X_2$ l'image schématique
de $V_{12} \to X_1 \times_S X_2$. Les morphismes de projection
$$
p_1 : X_{12} \to X_1
\quad\text{et}\quad
p_2 : X_{12} \to X_2
$$
sont propres, car $X_1$ et $X_2$ sont propres sur $S$. Si nous remplaçons $X_1$ par un
éclatement $V_1$-admissible, alors $X_{12}$ est remplacé par
le transformé strict relativement à cet éclatement, voir
le lemme \ref{flat-lemma-replaced-by-strict-transform}.

\medskip\noindent
Notons $\psi : V_{12} \to U$ le morphisme donné par les composées
$\psi = \psi_1 \circ p_1|_{V_{12}} = \psi_2 \circ p_2|_{V_{12}}$.
Considérons le sous-schéma fermé
$$
Z_{12, 2} =
(p_1|_{V_{12}})^{-1}(Z_{1, 2}) =
(p_2|_{V_{12}})^{-1}(Z_{2, 2}) =
\psi^{-1}(Z_2) \subset V_{12}
$$
Le morphisme $p_1|_{V_{12}} : V_{12} \to V_1$ est un isomorphisme
au-dessus d'un voisinage ouvert de $Z_{1, 2}$, car $\psi_2 : V_2 \to U$
est un isomorphisme au-dessus d'un voisinage ouvert de $Z_2$ et
$V_{12} = V_1 \times_U V_2$.
D'après le lemme \ref{flat-lemma-blowup-iso-along},
il existe un éclatement $V_1$-admissible $X_1' \to X_1$
tel que le transformé strict $p'_1 : X'_{12} \to X'_1$
de $p_1$ soit un isomorphisme au-dessus d'un voisinage ouvert de
l'adhérence de $Z_{1, 2}$ dans $X'_1$.
Après avoir remplacé $X_1$ par $X'_1$ et $X_{12}$ par $X'_{12}$,
nous pouvons supposer que $p_1$ est un isomorphisme au-dessus d'un voisinage
ouvert de $\overline{Z}_{1, 2}$.

\medskip\noindent
La réduction du paragraphe précédent nous dit que
$$
X_{12} \cap (\overline{Z}_{1, 2} \times_S \overline{Z}_{2, 1}) = \emptyset
$$
où l'intersection est prise dans $X_1 \times_S X_2$. En effet, l'image
réciproque $p_1^{-1}(\overline{Z}_{1, 2})$ dans $X_{12}$ est envoyée isomorphiquement
sur $\overline{Z}_{1, 2}$. En particulier, nous voyons que $Z_{12, 2}$
est dense dans $p_1^{-1}(\overline{Z}_{1, 2})$. Ainsi $p_2$ envoie
$p_1^{-1}(\overline{Z}_{1, 2})$ dans $\overline{Z}_{2, 2}$.
Puisque $\overline{Z}_{2, 2} \cap \overline{Z}_{2, 1} = \emptyset$,
nous concluons.

\medskip\noindent
Considérons les schémas
$$
W_i = U \coprod\nolimits_{U_i} (X_i \setminus \overline{Z}_{i, j}),
\quad i = 1, 2
$$
obtenus par recollement. Appliquons le lemme \ref{flat-lemma-check-separated}
pour voir que $W_i \to S$ est séparé. Tout d'abord,
$U \to S$ et $X_i \to S$ sont séparés. L'immersion
$U_i \to U \times_S (X_i \setminus \overline{Z}_{i, j})$
est fermée, car toute spécialisation $u_i \leadsto u$
avec $u_i \in U_i$ et $u \in U \setminus U_i$
se relève de manière unique en une spécialisation
$u_i \leadsto v_i$ dans $V_i$ le long du morphisme propre
$\psi_i : V_i \to U$, et alors $v_i$ appartient nécessairement à $Z_{i, j}$.
Ainsi, l'image de l'immersion est fermée, et l'immersion
est donc une immersion fermée.

\medskip\noindent
D'autre part, pour tout anneau de valuation $A$ sur $S$, ayant $K$ pour corps des fractions,
et tout morphisme $\gamma : \Spec(K) \to (U_1 \cap U_2)$ sur $S$, il
existe un $i$ et un prolongement de $\gamma$ en un morphisme $h_i : \Spec(A) \to W_i$.
En effet, pour chacun des $i = 1, 2$, il existe un morphisme
$g_i : \Spec(A) \to X_i$ prolongeant $\gamma$ d'après le
critère valuatif de propreté de $X_i$ sur $S$, voir
Morphismes, lemme \ref{morphisms-lemma-characterize-proper}.
Ainsi, nous ne rencontrons une difficulté
que si $g_i(\mathfrak m_A) \in \overline{Z}_{i, j}$ pour $i = 1, 2$. Cela est
impossible, car l'intersection de $X_{12}$ et
$\overline{Z}_{1, 2} \times_S \overline{Z}_{2, 1}$ est vide, comme démontré ci-dessus.

\medskip\noindent
Considérons un diagramme
$$
\xymatrix{
W_1' \ar[d] \ar[r] & W & W_2' \ar[l] \ar[d] \\
W_1 & U \ar[l] \ar[lu] \ar[u] \ar[ru] \ar[r] & W_2
}
$$
comme dans le lemme \ref{flat-lemma-find-common-blowups}. D'après le paragraphe précédent,
pour tout diagramme aux flèches pleines
$$
\xymatrix{
\Spec(K) \ar[r]_\gamma  \ar[d] & W \ar[d] \\
\Spec(A) \ar@{..>}[ru] \ar[r] & S
}
$$
où $\Im(\gamma) \subset U_1 \cap U_2$, il existe un $i$ et
un prolongement $h_i : \Spec(A) \to W_i$ de $\gamma$.
En utilisant le critère valuatif de propreté pour $W'_i \to W_i$,
nous pouvons alors relever $h_i$ en $h'_i : \Spec(A) \to W'_i$.
La flèche pointillée du diagramme existe donc. Comme $W$
est séparé sur $S$, nous voyons que cette flèche est également unique.
Cela implique que $W \to S$ est universellement fermé d'après
Morphismes, lemme
\ref{morphisms-lemma-refined-valuative-criterion-universally-closed}.
Comme $W \to S$ est déjà de type fini et séparé, on conclut.
\end{proof}

\begin{theorem}
\label{flat-theorem-nagata}
\begin{reference}
Voir \cite{Lutkebohmert}, \cite{Conrad-Nagata}, \cite{Nagata-1},
\cite{Nagata-2}, \cite{Nagata-3} et \cite{Nagata-4}
\end{reference}
Soit $S$ un schéma quasi-compact et quasi-séparé. Soit
$X \to S$ un morphisme séparé de type fini.
Alors $X$ admet une compactification sur $S$.
\end{theorem}

\begin{proof}
Nous nous ramenons d'abord au cas noethérien. Nous recommandons vivement au lecteur
de sauter ce paragraphe. Il existe une immersion fermée
$X \to X'$ telle que $X' \to S$ soit de présentation finie et séparé.
Voir Limites, proposition
\ref{limits-proposition-separated-closed-in-finite-presentation}.
Si nous trouvons une compactification de $X'$ sur $S$, alors
prendre l'image schématique de $X$ dans celle-ci donnera
une compactification de $X$ sur $S$. Nous pouvons donc supposer que
$X \to S$ est séparé et de présentation finie.
Nous pouvons écrire $S = \lim S_i$ comme limite d'un
système filtrant de schémas noethériens à morphismes de transition affines.
Voir Limites, proposition \ref{limits-proposition-approximate}.
Nous pouvons choisir un $i$ et un morphisme $X_i \to S_i$ de présentation
finie dont le changement de base à $S$ est $X \to S$, voir
Limites, lemme \ref{limits-lemma-descend-finite-presentation}.
Quitte à augmenter $i$, nous pouvons supposer que $X_i \to S_i$ est séparé, voir
Limites, lemme \ref{limits-lemma-descend-separated-finite-presentation}.
Si nous pouvons trouver une compactification de $X_i$ sur $S_i$, alors le
changement de base de celle-ci à $S$ sera une compactification de $X$ sur $S$.
Cela nous ramène au cas traité dans le paragraphe suivant.

\medskip\noindent
Supposons que $S$ soit noethérien. Nous pouvons choisir un recouvrement ouvert affine fini
$X = \bigcup_{i = 1, \ldots, n} U_i$ tel que $U_1 \cap \ldots \cap U_n$
soit dense dans $X$. Cela découle de
Propriétés, lemme \ref{properties-lemma-point-and-maximal-points-affine},
et du fait que $X$ est quasi-compact et possède un nombre fini de
composantes irréductibles. Pour chaque $i$, nous pouvons choisir un $n_i \geq 0$ et une
immersion $U_i \to \mathbf{A}^{n_i}_S$ d'après
Morphismes, lemme \ref{morphisms-lemma-quasi-affine-finite-type-over-S}.
Ainsi, $U_i$ admet une compactification sur $S$ pour $i = 1, \ldots, n$
en prenant l'image schématique dans $\mathbf{P}^{n_i}_S$.
En appliquant le lemme \ref{flat-lemma-two-compactifications}
$(n - 1)$ fois, nous concluons que le théorème est vrai.
\end{proof}






\section{La topologie h}
\label{flat-section-h}

\noindent
Pour nous, en termes informels, un faisceau h est un faisceau pour la topologie de Zariski
qui satisfait à la propriété de faisceau pour les morphismes propres surjectifs
de présentation finie ; voir le lemme \ref{flat-lemma-characterize-sheaf-h}.
Il peut toutefois être utile de souligner que la définition de la topologie h
sur la catégorie des schémas dépend de la référence choisie.

\medskip\noindent
Voevodsky a initialement défini un
recouvrement h comme une famille finie de morphismes de type fini qui sont
conjointement universellement submersifs (Morphismes, définition
\ref{morphisms-definition-submersive}). Voir \cite[définition 3.1.2]{Voevodsky}.
Cette définition est particulièrement adaptée lorsque la catégorie de schémas sous-jacente est
restreinte aux schémas de type fini sur un schéma de base noethérien fixé.
Dans ce cadre, Voevodsky relie les recouvrements h aux recouvrements ph.
La topologie ph est engendrée par les recouvrements de Zariski et les morphismes propres
surjectifs. Voir Topologies, section \ref{topologies-section-ph}
pour plus de détails.

\medskip\noindent
Dans Topologies, section \ref{topologies-section-V}, nous étudions la topologie V.
Un morphisme quasi-compact $X \to Y$ définit un recouvrement V si toute spécialisation
de points de $Y$ est l'image d'une spécialisation de points de $X$
et si cette propriété subsiste après tout changement de base
(Topologies, lemme \ref{topologies-lemma-refine-qcqs-V}).
Dans ce cas, $X \to Y$ est universellement submersif
(Topologies, lemme \ref{topologies-lemma-V-covering-universally-submersive}).
Il s'avère que la notion de recouvrement V remplace avantageusement les
``familles de morphismes de but fixé qui sont
conjointement universellement submersifs'' lorsqu'on travaille avec des schémas non noethériens.

\medskip\noindent
Notre démarche consistera d'abord à démontrer l'équivalence entre les recouvrements ph et
les recouvrements V pour les familles (éventuellement infinies) de morphismes qui sont
localement de présentation finie. Nous utiliserons ensuite ces familles
comme notion de recouvrement h dans le projet Champs.
Pour les schémas noethériens et les familles finies, ces recouvrements
coïncident avec ceux de la définition de Voevodsky ; voir le
lemme \ref{flat-lemma-Noetherian-h-covering}.
Sur la catégorie des schémas de présentation finie sur un schéma de base
fixé $S$, quasi-compact et quasi-séparé, ces recouvrements
définissent la même topologie que celle de \cite[définition 2.7]{Witt-Grass}.

\begin{lemma}
\label{flat-lemma-equivalence-h-v-locally-finite-presentation}
Soit $\{f_i : X_i \to X\}_{i \in I}$ une famille de morphismes
de schémas de but fixé, où $f_i$ est localement de
présentation finie pour tout $i$. Les assertions suivantes sont équivalentes :
\begin{enumerate}
\item $\{X_i \to X\}$ est un recouvrement ph, et
\item $\{X_i \to X\}$ est un recouvrement V.
\end{enumerate}
\end{lemma}

\begin{proof}
Soit $U \subset X$ un ouvert affine. D'après
Topologies, définitions \ref{topologies-definition-ph-covering}
et \ref{topologies-definition-V-covering}, il suffit de montrer que
la famille obtenue par changement de base $\{X_i \times_X U \to U\}$ admet un raffinement
par un recouvrement ph standard si et seulement si elle admet un raffinement par
un recouvrement V standard. On peut donc supposer $X$ affine et il faut montrer que
$\{X_i \to X\}$ admet un raffinement par un recouvrement ph standard
si et seulement si elle admet un raffinement par un recouvrement V standard.
Puisqu'un recouvrement ph standard est un recouvrement V standard, voir
Topologies, lemme \ref{topologies-lemma-standard-ph-standard-V},
il suffit de démontrer la réciproque.

\medskip\noindent
Supposons $X$ affine et supposons que $\{f_i : X_i \to X\}_{i \in I}$
admette un raffinement par un recouvrement V standard
$\{g_j : Y_j \to X\}_{j = 1, \ldots, m}$.
Pour chaque $j$, choisissons un indice $i_j$ et un morphisme
$h_j : Y_j \to X_{i_j}$ tels que $g_j = f_{i_j} \circ h_j$.
Puisque $Y_j$ est affine, donc quasi-compact,
nous pouvons, pour chaque $j$, trouver un nombre fini d'ouverts affines
$U_{j, k} \subset X_{i_j}$ tels que $\Im(h_j) \subset \bigcup U_{j, k}$.
Alors $\{U_{j, k} \to X\}_{j, k}$ raffine $\{X_i \to X\}$
et constitue un recouvrement V standard (car c'est une famille finie de morphismes
entre schémas affines et que sa propriété de relèvement pour les anneaux de valuation
découle de la propriété correspondante de $\{Y_j \to X\}$).
Nous sommes ainsi ramenés au cas traité au paragraphe suivant.

\medskip\noindent
Supposons que $\{f_i : X_i \to X\}_{i = 1, \ldots, n}$
soit un recouvrement V standard, les $f_i$ étant de présentation finie.
Il faut montrer que $\{X_i \to X\}$ admet un raffinement par un recouvrement ph standard.
Choisissons une stratification de platitude générique
$$
X = S \supset S_0 \supset S_1 \supset \ldots \supset S_t = \emptyset
$$
comme dans Compléments sur les morphismes, lemme
\ref{more-morphisms-lemma-generic-flatness-stratification-scheme}
pour le morphisme de présentation finie
$$
\coprod\nolimits_{i = 1, \ldots, n} f_i :
\coprod\nolimits_{i = 1, \ldots, n} X_i
\longrightarrow
X
$$
de schémas affines. Nous utiliserons toutes les propriétés de la stratification
sans plus le préciser. Par construction, le changement de base de chaque $f_i$ à
$U_k = S_k \setminus S_{k + 1}$ est plat.
Notons $Y_k$ l'adhérence schématique de $U_k$ dans $S_k$. Comme
$U_k \to S_k$ est une immersion ouverte quasi-compacte (voir
Propriétés, lemme \ref{properties-lemma-quasi-coherent-finite-type-ideals}),
on voit que $U_k \subset Y_k$ est une immersion ouverte quasi-compacte et dense
(et schématiquement dense) ; voir
Morphismes, lemme \ref{morphisms-lemma-quasi-compact-scheme-theoretic-image}.
Le morphisme $\coprod_{k = 0, \ldots, t - 1} Y_k \to X$
est fini et surjectif ; $\{Y_k \to X\}$ est donc un recouvrement ph standard
et, par conséquent, un recouvrement V standard (voir ci-dessus). Par la propriété de transitivité
des recouvrements V standards
(Topologies, lemme \ref{topologies-lemma-composition-standard-V}),
il suffit de montrer que le tiré en arrière
du recouvrement $\{X_i \to X\}$ sur chaque $Y_k$ admet un raffinement par un
recouvrement V standard. Cela nous ramène au cas décrit dans le
paragraphe suivant.

\medskip\noindent
Supposons que $\{f_i : X_i \to X\}_{i = 1, \ldots, n}$ soit un recouvrement V standard
où chaque $f_i$ est de présentation finie et qu'il existe un ouvert quasi-compact dense
$U \subset X$ tel que $X_i \times_X U \to U$ soit plat.
D'après le théorème \ref{flat-theorem-flatten-module}
il existe un éclatement $U$-admissible $X' \to X$ tel que
le transformé strict $f'_i : X'_i \to X'$ de $f_i$ soit plat.
Observons que le morphisme projectif (donc fermé) $X' \to X$
est surjectif puisque $U \subset X$ est dense et que $U$ s'identifie
à un ouvert de $X'$. Après avoir remplacé $X'$ par un nouvel
éclatement $U$-admissible si nécessaire,
nous pouvons également supposer que $U \subset X'$ est schématiquement dense
(voir la remarque \ref{flat-remark-successive-blowups}).
Ainsi, pour tout point $x \in X'$, il existe un anneau de valuation $V$
et un morphisme $g : \Spec(V) \to X'$ tels que le point générique
de $\Spec(V)$ s'envoie dans $U$ et que le point fermé de
$\Spec(V)$ s'envoie sur $x$, voir Morphismes, lemme
\ref{morphisms-lemma-reach-points-scheme-theoretic-image}.
Puisque $\{X_i \to X\}$ est un recouvrement V standard, nous pouvons choisir
une extension d'anneaux de valuation $V \subset W$, un indice $i$ et un morphisme
$\Spec(W) \to X_i$ tels que le diagramme
$$
\xymatrix{
\Spec(W) \ar[d] \ar[rr] & & X_i \ar[d] \\
\Spec(V) \ar[r] & X' \ar[r] & X
}
$$
soit commutatif. Puisque $X'_i \subset X' \times_X X_i$ est un sous-schéma fermé
contenant l'ouvert $U \times_X X_i$, puisque $\Spec(W)$ est un schéma intègre,
et puisque le morphisme induit $h : \Spec(W) \to X' \times_X X_i$ envoie
le point générique de $\Spec(W)$ dans $U \times_X X_i$, nous en déduisons
que $h$ se factorise par le sous-schéma fermé $X'_i \subset X' \times_X X_i$.
Nous en déduisons que $\{f'_i : X'_i \to X'\}$ est un recouvrement V.
En particulier, $\coprod f'_i$ est surjectif. En particulier,
$\{X'_i \to X'\}$ est un recouvrement fppf. Puisqu'un recouvrement fppf est un recouvrement ph
(Compléments sur les morphismes, lemme \ref{more-morphisms-lemma-fppf-ph}),
il existe un recouvrement ph standard $\{Y_j \to X'\}$ qui raffine
$\{X'_i \to X\}$. Disons que ce recouvrement est donné par un morphisme propre surjectif
$Y \to X'$ et un recouvrement fini par des ouverts affines
$Y = \bigcup Y_j$. Alors la composée $Y \to X$ est propre et surjective,
et nous en déduisons que $\{Y_j \to X\}$ est un recouvrement ph standard.
Cela achève la démonstration.
\end{proof}

\noindent
Voici notre définition.

\begin{definition}
\label{flat-definition-h-covering}
Soit $T$ un schéma. Un {\it recouvrement h de $T$} est une famille de morphismes
$\{f_i : T_i \to T\}_{i \in I}$ telle que chaque $f_i$ soit
localement de présentation finie et que l'une des conditions équivalentes du
lemme \ref{flat-lemma-equivalence-h-v-locally-finite-presentation} soit satisfaite.
\end{definition}

\noindent
Pour les schémas noethériens, cette notion coïncide avec celle de recouvrement ph
(nous le consignons dans le lemme \ref{flat-lemma-Noetherian-h-ph} ci-dessous), et
nous retrouvons la notion de Voevodsky.

\begin{lemma}
\label{flat-lemma-Noetherian-h-covering}
Soit $X$ un schéma noethérien. Soit $\{X_i \to X\}_{i \in I}$
une famille finie de morphismes de type fini. Les assertions suivantes sont équivalentes :
\begin{enumerate}
\item $\coprod_{i \in I} X_i \to X$ est universellement submersif
(Morphismes, définition
\ref{morphisms-definition-submersive}), et
\item $\{X_i \to X\}_{i \in I}$ est un recouvrement h.
\end{enumerate}
\end{lemma}

\begin{proof}
L'implication (2) $\Rightarrow$ (1) découle du résultat plus général
Topologies, lemme \ref{topologies-lemma-V-covering-universally-submersive}
et de notre définition des recouvrements h. Supposons que $\coprod X_i \to X$
soit universellement submersif. Nous allons montrer que $\{X_i \to X\}$
admet un raffinement qui est un recouvrement ph ; cela suffira d'après
Topologies, lemme \ref{topologies-lemma-refine-by-ph} et
notre définition des recouvrements h.
L'argument sera le même que celui employé dans la démonstration du
lemme \ref{flat-lemma-equivalence-h-v-locally-finite-presentation}.

\medskip\noindent
Choisissons une stratification de platitude générique
$$
X = S \supset S_0 \supset S_1 \supset \ldots \supset S_t = \emptyset
$$
comme dans Compléments sur les morphismes, lemme
\ref{more-morphisms-lemma-generic-flatness-stratification-scheme}
pour le morphisme de présentation finie
$$
\coprod\nolimits_{i = 1, \ldots, n} f_i :
\coprod\nolimits_{i = 1, \ldots, n} X_i
\longrightarrow
X
$$
Nous allons utiliser toutes les propriétés de la stratification
sans autre mention. Par construction, le changement de base de chaque
$f_i$ à $U_k = S_k \setminus S_{k + 1}$ est plat.
Notons $Y_k$ l'adhérence schématique de $U_k$ dans $S_k$. Puisque
$U_k \to S_k$ est une immersion ouverte quasi-compacte (tous les schémas de
ce paragraphe sont noethériens),
nous voyons que $U_k \subset Y_k$ est une immersion ouverte quasi-compacte et dense
(et schématiquement dense), voir
Morphismes, lemme \ref{morphisms-lemma-quasi-compact-scheme-theoretic-image}.
Le morphisme $\coprod_{k = 0, \ldots, t - 1} Y_k \to X$
est fini et surjectif, donc $\{Y_k \to X\}$ est un recouvrement ph.
Par la propriété de transitivité des recouvrements ph
(Topologies, lemme \ref{topologies-lemma-ph}),
il suffit de montrer que le recouvrement obtenu par changement de base
de $\{X_i \to X\}$ à chaque $Y_k$ admet un raffinement qui est un
recouvrement ph. Cela nous ramène au cas décrit dans le
paragraphe suivant.

\medskip\noindent
Supposons que $\coprod X_i \to X$ soit universellement submersif et
qu'il existe un ouvert dense $U \subset X$ tel que
$X_i \times_X U \to U$ soit plat pour tout $i$.
D'après le théorème \ref{flat-theorem-flatten-module}
il existe un éclatement $U$-admissible $X' \to X$ tel que
le transformé strict $f'_i : X'_i \to X'$ de $f_i$ soit plat pour tout $i$.
Observons que le morphisme projectif (donc fermé) $X' \to X$
est surjectif puisque $U \subset X$ est dense et que $U$ s'identifie
à un ouvert de $X'$. En remplaçant au besoin $X'$ par un nouvel
éclatement $U$-admissible, nous pouvons aussi supposer que $U \subset X'$ est dense
(voir la remarque \ref{flat-remark-successive-blowups}).
Ainsi, pour tout point $x \in X'$, il existe un anneau de valuation discrète $A$
et un morphisme $g : \Spec(A) \to X'$ tel que le point
générique de $\Spec(A)$ soit envoyé dans $U$ et que le point fermé de
$\Spec(A)$ soit envoyé sur $x$; voir Limites, lemme
\ref{limits-lemma-reach-point-closure-Noetherian}.
Posons
$$
W = \Spec(A) \times_X \coprod X_i = \coprod \Spec(A) \times_X X_i
$$
Puisque $\coprod X_i \to X$ est universellement submersif,
il existe une spécialisation $w' \leadsto w$ dans $W$
telle que $w'$ ait pour image le point générique de $\Spec(A)$
et que $w$ ait pour image le point fermé de $\Spec(A)$.
(Sinon, la fibre fermée de $W \to \Spec(A)$
serait stable par générisation, donc ouverte, ce qui
contredirait le fait que $W \to \Spec(A)$ est submersif.)
Disons que $w' \in \Spec(A) \times_X X_i$; alors, bien sûr,
$w \in \Spec(A) \times_X X_i$ également. Soit
$x'_i \leadsto x_i$ l'image de $w' \leadsto w$ dans
$X' \times_X X_i$. Puisque $x'_i \in X'_i$ et que
$X'_i \subset X' \times_X X_i$ est un sous-schéma fermé,
nous voyons que $x_i \in X'_i$. Puisque $x_i$ a pour image $x \in X'$,
nous concluons que
$\coprod X'_i \to X'$ est surjectif ! En particulier,
$\{X'_i \to X'\}$ est un recouvrement fppf. Or, tout recouvrement fppf est un recouvrement ph
(Compléments sur les morphismes, lemme \ref{more-morphisms-lemma-fppf-ph}).
Comme $X' \to X$ est propre et surjectif, nous concluons
que $\{X'_i \to X\}$ est un recouvrement ph, ce qui achève la démonstration.
\end{proof}

\begin{lemma}
\label{flat-lemma-Noetherian-h-ph}
Soit $X$ un schéma localement noethérien. Une famille de morphismes
$\{f_i : X_i \to X\}_{i \in I}$ de but $X$ est un recouvrement h
si et seulement si elle est un recouvrement ph.
\end{lemma}

\begin{proof}
D'après la définition \ref{flat-definition-h-covering}, tout recouvrement h est un recouvrement ph.
Réciproquement, si $\{f_i : X_i \to X\}$ est un recouvrement ph, alors les morphismes
$f_i$ sont localement de type fini (Topologies, définition
\ref{topologies-definition-ph-covering}). Puisque $X$ est localement noethérien,
chaque $f_i$ est localement de présentation finie, et nous voyons qu'il s'agit
d'un recouvrement h par définition.
\end{proof}

\noindent
Le lemme suivant et \cite[théorème 8.4]{rydh_descent}
montrent que notre définition coïncide avec (ou du moins est
étroitement liée à) celle de l'article \cite{rydh_descent}
de David Rydh. Pour simplifier, nous nous limitons à une base affine.

\begin{lemma}
\label{flat-lemma-approximate-h-cover}
Soit $X$ un schéma affine. Soit $\{X_i \to X\}_{i \in I}$
un recouvrement h. Alors il existe un morphisme propre surjectif
$$
Y \longrightarrow X
$$
de présentation finie (!) et un recouvrement ouvert affine fini
$Y = \bigcup_{j = 1, \ldots, m} Y_j$ tel que
$\{Y_j \to X\}_{j = 1, \ldots, m}$ raffine $\{X_i \to X\}_{i \in I}$.
\end{lemma}

\begin{proof}
Par hypothèse, il existe un morphisme propre surjectif
$Y \to X$ et un recouvrement ouvert affine fini
$Y = \bigcup_{j = 1, \ldots, m} Y_j$ tel que
$\{Y_j \to X\}_{j = 1, \ldots, m}$ raffine $\{X_i \to X\}_{i \in I}$.
Cela signifie que, pour tout $j$, il existe un indice $i_j \in I$
et un morphisme $h_j : Y_j \to X_{i_j}$ au-dessus de $X$.
Voir la définition \ref{flat-definition-h-covering} et
Topologies, définition \ref{topologies-definition-ph-covering}.
Le problème est que nous ne savons pas si $Y \to X$ est de
présentation finie.
D'après
Limites, lemme \ref{limits-lemma-proper-limit-of-proper-finite-presentation},
nous pouvons écrire
$$
Y = \lim Y_\lambda
$$
comme limite d'un système projectif filtrant de schémas $Y_\lambda$ propres et de présentation finie
sur $X$, tels que les morphismes $Y \to Y_\lambda$ et les
morphismes de transition soient des immersions fermées. Observons que
chaque $Y_\lambda \to X$ est surjectif.
D'après Limites, lemme \ref{limits-lemma-descend-opens},
nous pouvons trouver un $\lambda$ et des ouverts quasi-compacts
$Y_{\lambda, j} \subset Y_\lambda$, $j = 1, \ldots, m$,
qui recouvrent $Y_\lambda$ et induisent $Y_j$ sur $Y$.
Alors $Y_j = \lim Y_{\lambda, j}$.
Quitte à augmenter $\lambda$, nous pouvons supposer $Y_{\lambda, j}$
affine pour tout $j$; voir
Limites, lemme \ref{limits-lemma-limit-affine}.
Enfin, puisque $X_i \to X$ est localement de présentation finie,
nous pouvons utiliser la caractérisation fonctorielle des morphismes
localement de présentation finie
(Limites, proposition
\ref{limits-proposition-characterize-locally-finite-presentation})
pour trouver un $\lambda$ tel que, pour tout $j$, il existe
un morphisme $h_{\lambda, j} : Y_{\lambda, j} \to X_{i_j}$
dont la restriction à $Y_j$ est le morphisme $h_j$ choisi ci-dessus.
Ainsi $\{Y_{\lambda, j} \to X\}$ raffine
$\{X_i \to X\}$, ce qui achève la démonstration.
\end{proof}

\noindent
Nous revenons au développement de la théorie générale des recouvrements h.

\begin{lemma}
\label{flat-lemma-zariski-h}
Un recouvrement fppf est un recouvrement h. Par conséquent, les recouvrements syntomiques, lisses, \'etales,
et de Zariski sont également des recouvrements h.
\end{lemma}

\begin{proof}
Cela est vrai car, dans un recouvrement fppf, on exige que les morphismes soient
localement de présentation finie, et car les recouvrements fppf sont des
recouvrements ph, voir Compléments sur les morphismes,
lemme \ref{more-morphisms-lemma-fppf-ph}.
La seconde assertion découle de la première et de
Topologies, lemme \ref{topologies-lemma-zariski-etale-smooth-syntomic-fppf}.
\end{proof}

\begin{lemma}
\label{flat-lemma-surjective-proper-finite-presentation-h}
Soit $f : Y \to X$ un morphisme surjectif et propre de schémas
qui est de présentation finie. Alors $\{Y \to X\}$ est un recouvrement h.
\end{lemma}

\begin{proof}
Combiner les lemmes de Topologies
\ref{topologies-lemma-zariski-etale-smooth-syntomic-fppf-fpqc-ph-V} et
\ref{topologies-lemma-surjective-proper-ph}.
\end{proof}

\begin{lemma}
\label{flat-lemma-refine-by-h}
Soit $T$ un schéma. Soit $\{f_i : T_i \to T\}_{i \in I}$ une famille
de morphismes telle que $f_i$ soit localement de présentation finie pour tout $i$.
Les assertions suivantes sont équivalentes :
\begin{enumerate}
\item $\{T_i \to T\}_{i \in I}$ est un recouvrement h,
\item il existe un recouvrement h qui raffine $\{T_i \to T\}_{i \in I}$, et
\item $\{\coprod_{i \in I} T_i \to T\}$ est un recouvrement h.
\end{enumerate}
\end{lemma}

\begin{proof}
Cela découle de l'énoncé analogue pour les recouvrements ph
(Topologies, lemme \ref{topologies-lemma-refine-by-ph})
ou de l'énoncé analogue pour les recouvrements V
(Topologies, lemme \ref{topologies-lemma-refine-by-V}).
\end{proof}

\noindent
Ensuite, nous montrons que notre notion de recouvrement h satisfait aux conditions de
Sites, définition \ref{sites-definition-site}.

\begin{lemma}
\label{flat-lemma-h}
Soit $T$ un schéma.
\begin{enumerate}
\item Si $T' \to T$ est un isomorphisme, alors $\{T' \to T\}$
est un recouvrement h de $T$.
\item Si $\{T_i \to T\}_{i\in I}$ est un recouvrement h et si, pour chaque
$i$, nous avons un recouvrement h $\{T_{ij} \to T_i\}_{j\in J_i}$, alors
$\{T_{ij} \to T\}_{i \in I, j\in J_i}$ est un recouvrement h.
\item Si $\{T_i \to T\}_{i\in I}$ est un recouvrement h
et si $T' \to T$ est un morphisme de schémas, alors
$\{T' \times_T T_i \to T'\}_{i\in I}$ est un recouvrement h.
\end{enumerate}
\end{lemma}

\begin{proof}
Cela découle immédiatement de l'énoncé correspondant pour
les recouvrements ph ou les recouvrements V
(Topologies, lemme \ref{topologies-lemma-ph} ou
\ref{topologies-lemma-V})
et du fait que la classe des morphismes localement
de présentation finie est stable par changement de base et par
composition.
\end{proof}

\noindent
Ensuite, nous définissons les grands sites h avec lesquels nous travaillerons dans le
projet Champs. Il convient de lire la discussion générale
dans Topologies, section \ref{topologies-section-procedure}
avant de poursuivre.

\begin{definition}
\label{flat-definition-big-h-site}
Un {\it grand site h} est un site $\Sch_h$ quelconque comme dans
Sites, définition \ref{sites-definition-site}, construit comme suit :
\begin{enumerate}
\item Choisir un ensemble quelconque de schémas $S_0$ et un ensemble quelconque de recouvrements h
$\text{Cov}_0$ parmi ces schémas.
\item Prendre comme catégorie sous-jacente une catégorie quelconque $\Sch_\alpha$
construite comme dans Ensembles, lemme \ref{sets-lemma-construct-category}
à partir de l'ensemble $S_0$.
\item Choisir un ensemble quelconque de recouvrements comme dans
Ensembles, lemme \ref{sets-lemma-coverings-site}, à partir de la
catégorie $\Sch_\alpha$, de la classe des recouvrements h
et de l'ensemble $\text{Cov}_0$ choisi ci-dessus.
\end{enumerate}
\end{definition}

\noindent
Voir les remarques qui suivent
Topologies, définition \ref{topologies-definition-big-zariski-site},
pour la motivation et les explications concernant la définition des grands sites.

\begin{definition}
\label{flat-definition-standard-h}
Soit $T$ un schéma affine. Un {\it recouvrement h standard} de $T$
est une famille $\{f_i : T_i \to T\}_{i = 1, \ldots, n}$ dans laquelle chaque $T_i$ est
affine et chaque $f_i$ est de présentation finie, et qui satisfait à l'une des
conditions équivalentes suivantes : (1) $\{T_i \to T\}$ admet un raffinement par
un recouvrement ph standard ou (2) $\{T_i \to T\}$ est un recouvrement V.
\end{definition}

\noindent
L'équivalence de ces conditions résulte du
lemme \ref{flat-lemma-equivalence-h-v-locally-finite-presentation}, de
Topologies, définition \ref{topologies-definition-ph-covering}, et du
lemme \ref{topologies-lemma-refine-by-ph}.

\medskip\noindent
Avant de poursuivre l'introduction du grand site h d'un
schéma $S$, signalons que la topologie sur un grand site h
$\Sch_h$ est, en un certain sens, induite par la topologie h
sur la catégorie de tous les schémas.

\begin{lemma}
\label{flat-lemma-h-induced}
Soit $\Sch_h$ un grand site h comme dans la
définition \ref{flat-definition-big-h-site}.
Soit $T \in \Ob(\Sch_h)$.
Soit $\{T_i \to T\}_{i \in I}$ un recouvrement h arbitraire de $T$.
\begin{enumerate}
\item Il existe un recouvrement $\{U_j \to T\}_{j \in J}$ de $T$ dans le site
$\Sch_h$ qui raffine $\{T_i \to T\}_{i \in I}$.
\item Si $\{T_i \to T\}_{i \in I}$ est un recouvrement h standard, alors
il est tautologiquement équivalent à un recouvrement de $\Sch_h$.
\item Si $\{T_i \to T\}_{i \in I}$ est un recouvrement de Zariski, alors
il est tautologiquement équivalent à un recouvrement de $\Sch_h$.
\end{enumerate}
\end{lemma}

\begin{proof}
Omis. Indication : il s'agit exactement de la même démonstration que celle de
Topologies, lemme \ref{topologies-lemma-ph-induced}.
\end{proof}

\begin{definition}
\label{flat-definition-big-small-h}
Soit $S$ un schéma. Soit $\Sch_h$ un grand site h contenant $S$.
\begin{enumerate}
\item Le {\it grand site h de $S$}, noté
$(\Sch/S)_h$, est le site $\Sch_h/S$
introduit dans Sites, section \ref{sites-section-localize}.
\item Le {\it grand site h affine de $S$}, noté
$(\textit{Aff}/S)_h$, est la sous-catégorie pleine de
$(\Sch/S)_h$ dont les objets sont les $U/S$ affines.
Un recouvrement de $(\textit{Aff}/S)_h$ est tout recouvrement
$\{U_i \to U\}$ de $(\Sch/S)_h$ qui est un recouvrement h standard.
\end{enumerate}
\end{definition}

\noindent
Précisons explicitement que le grand site h affine est un site.

\begin{lemma}
\label{flat-lemma-verify-site-h}
Soit $S$ un schéma. Soit $\Sch_h$ un grand
site h contenant $S$. Alors $(\textit{Aff}/S)_h$ est un site.
\end{lemma}

\begin{proof}
En raisonnant comme dans la démonstration de
Topologies, lemme \ref{topologies-lemma-verify-site-etale},
il suffit de montrer que la collection des recouvrements h standards
satisfait aux propriétés (1), (2) et (3) de
Sites, définition \ref{sites-definition-site}.
Cela est clair puisque, par exemple, étant donné un recouvrement h standard
$\{T_i \to T\}_{i\in I}$ et, pour chaque
$i$, un recouvrement h standard $\{T_{ij} \to T_i\}_{j \in J_i}$, alors
$\{T_{ij} \to T\}_{i \in I, j\in J_i}$ est un recouvrement h
(lemme \ref{flat-lemma-h}), $\bigcup_{i\in I} J_i$ est fini et
chaque $T_{ij}$ est affine. Ainsi $\{T_{ij} \to T\}_{i \in I, j\in J_i}$
est un recouvrement h standard.
\end{proof}

\begin{lemma}
\label{flat-lemma-fibre-products-h}
Soit $S$ un schéma. Soit $\Sch_h$ un grand
site h contenant $S$. Les catégories sous-jacentes aux sites
$\Sch_h$, $(\Sch/S)_h$ et $(\textit{Aff}/S)_h$ possèdent des produits fibrés.
Dans chaque cas, le foncteur évident vers la catégorie $\Sch$ de
tous les schémas commute aux produits fibrés. La catégorie
$(\Sch/S)_h$ possède un objet final, à savoir $S/S$.
\end{lemma}

\begin{proof}
Pour $\Sch_h$, cela est vrai par construction, voir
Ensembles, lemme \ref{sets-lemma-what-is-in-it}.
Supposons donnés des morphismes de schémas $U \to S$, $V \to U$, $W \to U$
avec $U, V, W \in \Ob(\Sch_h)$.
Le produit fibré $V \times_U W$ dans $\Sch_h$
est un produit fibré dans $\Sch$ et
est le produit fibré de $V/S$ et de $W/S$ au-dessus de $U/S$ dans
la catégorie de tous les schémas au-dessus de $S$, et donc également un
produit fibré dans $(\Sch/S)_h$.
Cela démontre le résultat pour $(\Sch/S)_h$.
Si $U, V, W$ sont affines, alors $V \times_U W$ l'est aussi, d'où le
résultat pour $(\textit{Aff}/S)_h$.
\end{proof}

\noindent
Ensuite, vérifions que le grand site affine définit le même
topos que le grand site.

\begin{lemma}
\label{flat-lemma-affine-big-site-h}
Soit $S$ un schéma. Soit $\Sch_h$ un grand
site h contenant $S$.
Le foncteur $(\textit{Aff}/S)_h \to (\Sch/S)_h$
est cocontinu et induit une équivalence de topos de
$\Sh((\textit{Aff}/S)_h)$ vers
$\Sh((\Sch/S)_h)$.
\end{lemma}

\begin{proof}
La notion de foncteur cocontinu spécial est introduite dans
Sites, définition \ref{sites-definition-special-cocontinuous-functor}.
Nous devons donc vérifier les hypothèses (1) -- (5) de
Sites, lemme \ref{sites-lemma-equivalence}.
Notons le foncteur d'inclusion
$u : (\textit{Aff}/S)_h \to (\Sch/S)_h$.
La cocontinuité résulte du fait que tout recouvrement h de
$T/S$, avec $T$ affine, peut être raffiné par un recouvrement h standard,
par exemple d'après le lemme \ref{flat-lemma-approximate-h-cover}. Ainsi (1) est satisfaite.
Nous voyons que $u$ est continu simplement parce qu'un recouvrement h standard
est un recouvrement h.
Ainsi, (2) est vérifiée. Les assertions (3) et (4) découlent immédiatement du fait
que $u$ est pleinement fidèle. Enfin, la condition (5) découle du
fait que tout schéma possède un recouvrement ouvert affine (qui est
un recouvrement h).
\end{proof}

\begin{lemma}
\label{flat-lemma-characterize-sheaf-h}
Soit $\mathcal{F}$ un préfaisceau sur $(\Sch/S)_h$.
Alors $\mathcal{F}$ est un faisceau si et seulement si
\begin{enumerate}
\item $\mathcal{F}$ satisfait la condition de faisceau pour les
recouvrements de Zariski, et
\item si $f : V \to U$ est propre, surjectif et de présentation finie, alors
l'application de $\mathcal{F}(U)$ vers l'égalisateur
des deux applications $\mathcal{F}(V) \to \mathcal{F}(V \times_U V)$ est bijective.
\end{enumerate}
De plus, en présence de (1), la propriété (2) équivaut à la
propriété
\begin{enumerate}
\item[(2')] la propriété de faisceau pour $\{V \to U\}$ comme en (2), avec $U$ affine.
\end{enumerate}
\end{lemma}

\begin{proof}
Nous allons montrer que, si (1) et (2) sont vérifiées, alors $\mathcal{F}$ est un faisceau.
Soit $\{T_i \to T\}$ un recouvrement dans $(\Sch/S)_h$.
Vérifions la condition de faisceau pour ce recouvrement.
Soient $s_i \in \mathcal{F}(T_i)$ des sections dont les restrictions donnent la même
section sur $T_i \times_T T_{i'}$. Montrons qu'il existe une
unique section $s \in \mathcal{F}(T)$ se restreignant à $s_i$ sur $T_i$.
Soit $T = \bigcup U_j$ un recouvrement ouvert affine.
D'après la propriété (1), il suffit de construire des sections $s_j \in \mathcal{F}(U_j)$
qui coïncident sur $U_j \cap U_{j'}$ pour obtenir $s$.
Considérons les recouvrements $\{T_i \times_T U_j \to U_j\}$.
Alors les $s_{ji} = s_i|_{T_i \times_T U_j}$ sont des sections qui coïncident
sur $(T_i \times_T U_j) \times_{U_j} (T_{i'} \times_T U_j)$.
Choisissons un morphisme propre surjectif $V_j \to U_j$ de présentation finie
et un recouvrement ouvert affine fini $V_j = \bigcup V_{jk}$
tel que $\{V_{jk} \to U_j\}$ raffine $\{T_i \times_T U_j \to U_j\}$.
Voir le lemme \ref{flat-lemma-approximate-h-cover}.
Si $s_{jk} \in \mathcal{F}(V_{jk})$
désigne le tiré en arrière de $s_{ji}$ sur $V_{jk}$ par les
morphismes sous-entendus, alors les $s_{jk}$ se recollent en une section
$s'_j \in \mathcal{F}(V_j)$. En utilisant encore la coïncidence sur les intersections,
nous voyons que $s'_j$ appartient à l'égalisateur des deux
applications $\mathcal{F}(V_j) \to \mathcal{F}(V_j \times_{U_j} V_j)$.
Ainsi, d'après (2), $s'_j$ provient d'une unique section
$s_j \in \mathcal{F}(U_j)$. Nous omettons de vérifier que ces
sections $s_j$ possèdent toutes les propriétés voulues.

\medskip\noindent
Démonstration de l'équivalence de (2) et (2') en présence de (1).
Supposons que $V \to U$ soit un morphisme de $(\Sch/S)_h$,
propre, surjectif et de présentation finie. Choisissons un
recouvrement ouvert affine $U = \bigcup U_i$ et posons $V_i = V \times_U U_i$.
Alors $\mathcal{F}(U) \to \mathcal{F}(V)$
est injective, car $\mathcal{F}(U_i) \to \mathcal{F}(V_i)$
est injective d'après (2') et $\mathcal{F}(U) \to \prod \mathcal{F}(U_i)$
est injective d'après (1). Enfin, supposons que l'on se donne un élément
$t \in \mathcal{F}(V)$ appartenant à l'égalisateur des deux applications
$\mathcal{F}(V) \to \mathcal{F}(V \times_U V)$.
Alors $t|_{V_i}$ appartient à l'égalisateur des deux applications
$\mathcal{F}(V_i) \to \mathcal{F}(V_i \times_{U_i} V_i)$
pour tout $i$. Nous obtenons donc une unique section $s_i \in \mathcal{F}(U_i)$
qui est envoyée sur $t|_{V_i}$ pour tout $i$, d'après (2').
Nous omettons de vérifier que $s_i|_{U_i \cap U_j} = s_j|_{U_i \cap U_j}$
pour tous $i, j$; cela utilise la propriété d'unicité que nous venons d'établir.
La propriété de faisceau pour le recouvrement $U = \bigcup U_i$ donne
une section $s \in \mathcal{F}(U)$. Nous omettons de démontrer que $s$
a pour image $t$ dans $\mathcal{F}(V)$.
\end{proof}

\noindent
Établissons maintenant quelques relations entre les topos
associés à ces sites.

\begin{lemma}
\label{flat-lemma-morphism-big-h}
Soit $\Sch_h$ un grand site h.
Soit $f : T \to S$ un morphisme dans $\Sch_h$.
Le foncteur
$$
u : (\Sch/T)_h \longrightarrow (\Sch/S)_h,
\quad
V/T \longmapsto V/S
$$
est cocontinu et admet un adjoint à droite continu
$$
v : (\Sch/S)_h \longrightarrow (\Sch/T)_h,
\quad
(U \to S) \longmapsto (U \times_S T \to T).
$$
Ces deux foncteurs induisent le même morphisme de topos
$$
f_{big} :
\Sh((\Sch/T)_h)
\longrightarrow
\Sh((\Sch/S)_h)
$$
Nous avons $f_{big}^{-1}(\mathcal{G})(U/T) = \mathcal{G}(U/S)$.
Nous avons $f_{big, *}(\mathcal{F})(U/S) = \mathcal{F}(U \times_S T/T)$.
De plus, $f_{big}^{-1}$ admet un adjoint à gauche $f_{big!}$ qui commute aux
produits fibrés et aux égalisateurs.
\end{lemma}

\begin{proof}
Le foncteur $u$ est cocontinu, continu et commute aux produits fibrés
et aux égalisateurs. Par conséquent,
Sites, lemmes \ref{sites-lemma-when-shriek} et
\ref{sites-lemma-preserve-equalizers}
s'appliquent, et nous en déduisons la formule
donnant $f_{big}^{-1}$ ainsi que l'existence de $f_{big!}$. De plus,
le foncteur $v$ est adjoint à droite car, étant donnés $U/T$ et $V/S$,
nous avons $\Mor_S(u(U), V) = \Mor_T(U, V \times_S T)$,
comme voulu. Nous pouvons donc appliquer
Sites, lemmes \ref{sites-lemma-have-functor-other-way} et
\ref{sites-lemma-have-functor-other-way-morphism} pour obtenir la
formule pour $f_{big, *}$.
\end{proof}

\begin{lemma}
\label{flat-lemma-composition-h}
Étant donnés des schémas $X$, $Y$, $Y$ dans $(\Sch/S)_h$
et des morphismes $f : X \to Y$, $g : Y \to Z$, nous avons
$g_{big} \circ f_{big} = (g \circ f)_{big}$.
\end{lemma}

\begin{proof}
Cela résulte de la description élémentaire des images directe
et inverse pour les foncteurs entre les grands sites donnée dans le
lemme \ref{flat-lemma-morphism-big-h}.
\end{proof}










\section{Compléments sur la topologie h}
\label{flat-section-h-more}

\noindent
Dans cette section, nous démontrons quelques résultats supplémentaires sur la topologie h.
Commençons par quelques contre-exemples.

\begin{example}
\label{flat-example-structure-sheaf-not-h-sheaf}
Le ``faisceau structural'' $\mathcal{O}$ n'est pas un faisceau pour la topologie h.
Considérons par exemple une immersion fermée surjective de
présentation finie $X \to Y$. Alors $\{X \to Y\}$ est un recouvrement h,
par exemple d'après le lemme \ref{flat-lemma-surjective-proper-finite-presentation-h}.
De plus, remarquons que $X \times_Y X = X$.
Ainsi, si $\mathcal{O}$ était un faisceau pour la topologie h, alors
$\mathcal{O}_Y(Y) \to \mathcal{O}_X(X)$ serait bijective.
Ce n'est pas le cas dès que $X$ et $Y$ sont affines et que le morphisme
$X \to Y$ n'est pas un isomorphisme.
\end{example}

\begin{example}
\label{flat-example-representable-sheaf-not-h-sheaf}
Sur chacun des sites $(\Sch/S)_h$, la topologie n'est pas sous-canonique,
autrement dit, les préfaisceaux représentables ne sont pas tous des faisceaux.
En effet, le ``faisceau structural'' $\mathcal{O}$ est représentable
puisque $\mathcal{O}(X) = \Mor_S(X, \mathbf{A}^1_S)$
dans $(\Sch/S)_h$, et nous avons vu dans
l'exemple \ref{flat-example-structure-sheaf-not-h-sheaf}
que $\mathcal{O}$ n'est pas un faisceau.
\end{example}

\begin{lemma}
\label{flat-lemma-limit-h-topology}
Soit $T$ un schéma affine qui s'écrit comme une limite
$T = \lim_{i \in I} T_i$ d'un système projectif filtrant de schémas affines.
\begin{enumerate}
\item Soit $\mathcal{V} = \{V_j \to T\}_{j = 1, \ldots, m}$ un
recouvrement h standard de $T$, voir la définition
\ref{flat-definition-standard-h}.
Il existe alors un indice $i$ et un recouvrement h standard
$\mathcal{V}_i = \{V_{i, j} \to T_i\}_{j = 1, \ldots, m}$
dont le changement de base $T \times_{T_i} \mathcal{V}_i$ à $T$
est isomorphe à $\mathcal{V}$.
\item Soient $\mathcal{V}_i$, $\mathcal{V}'_i$ deux recouvrements h standards
de $T_i$. Si
$f : T \times_{T_i} \mathcal{V}_i \to T \times_{T_i} \mathcal{V}'_i$ est
un morphisme de recouvrements de $T$, alors il existe un indice
$i' \geq i$ et un morphisme
$f_{i'} : T_{i'} \times_{T_i} \mathcal{V} \to
T_{i'} \times_{T_i} \mathcal{V}'_i$
dont le changement de base à $T$ est $f$.
\item Si
$f, g : \mathcal{V} \to \mathcal{V}'_i$
sont des morphismes de recouvrements h standards de $T_i$ dont les
changements de base $f_T, g_T$ à $T$ sont égaux, alors il existe un
indice $i' \geq i$ tel que $f_{T_{i'}} = g_{T_{i'}}$.
\end{enumerate}
Autrement dit, la catégorie des recouvrements h standards de $T$ est
la colimite sur $I$ des catégories de recouvrements h standards de $T_i$.
\end{lemma}

\begin{proof}
D'après Limites, lemme \ref{limits-lemma-descend-finite-presentation},
la catégorie des schémas de présentation finie sur $T$ est la
colimite sur $I$ des catégories de schémas de présentation finie sur $T_i$. D'après
Limites, lemme \ref{limits-lemma-descend-affine-finite-presentation},
il en va de même pour la catégorie des schémas qui sont affines et
de présentation finie sur $T$.
Pour achever la démonstration du lemme, il suffit de montrer que si
$\{V_{j, i} \to T_i\}_{j = 1, \ldots, m}$ est une famille finie de
morphismes de présentation finie, les $V_{j, i}$ étant affines, et que la
famille obtenue par changement de base $\{T \times_{T_i} V_{j, i} \to T\}$ est
un recouvrement h, alors, pour un certain $i' \geq i$, la famille
$\{T_{i'} \times_{T_i} V_{j, i} \to T_{i'}\}$ est un recouvrement h.
Pour le voir, appliquons le lemme \ref{flat-lemma-approximate-h-cover} afin de
choisir un morphisme de présentation finie, propre et surjectif
$Y \to T$ ainsi qu'un recouvrement fini par des ouverts affines
$Y = \bigcup_{k = 1, \ldots, n} Y_k$ de sorte que
$\{Y_k \to T\}_{k = 1, \ldots, n}$ raffine
$\{T \times_{T_i} V_{j, i} \to T\}$.
En utilisant les arguments ci-dessus ainsi que
Limites, lemmes \ref{limits-lemma-eventually-proper},
\ref{limits-lemma-descend-surjective}, et
\ref{limits-lemma-descend-opens},
nous pouvons trouver un $i' \geq i$ et un morphisme de présentation finie, surjectif et propre
$Y_{i'} \to T_{i'}$ ainsi qu'un recouvrement par des ouverts affines
$Y_{i'} = \bigcup_{k = 1, \ldots, n} Y_{i', k}$
de sorte que, de plus, $\{Y_{i', k} \to Y_{i'}\}$ raffine
$\{T_{i'} \times_{T_i} V_{j, i} \to T_{i'}\}$.
Il s'ensuit que cette dernière famille est
un recouvrement h, ce qui achève la démonstration.
\end{proof}

\begin{lemma}
\label{flat-lemma-extend-sheaf-h}
Soit $S$ un schéma contenu dans un grand site $\Sch_h$.
Soit $F : (\Sch/S)_h^{opp} \to \textit{Ensembles}$ un faisceau h satisfaisant à
la propriété (b) de Topologies, lemme \ref{topologies-lemma-extend},
avec $\mathcal{C} = (\Sch/S)_h$.
Alors l'extension $F'$ de $F$ à la catégorie de tous les
schémas sur $S$ satisfait à la condition de faisceau pour tout recouvrement h
et commute aux limites projectives filtrantes (Limites, remarque \ref{limits-remark-limit-preserving}).
\end{lemma}

\begin{proof}
Cela résulte des arguments donnés dans les démonstrations de
Topologies, lemmes \ref{topologies-lemma-extend-sheaf-general} et
\ref{topologies-lemma-extend-sheaf}, en utilisant les
lemmes \ref{flat-lemma-limit-h-topology} et \ref{flat-lemma-h-induced}.
Détails omis.
\end{proof}








\section{Carrés d'éclatement et topologie ph}
\label{flat-section-blow-up-ph}

\noindent
Soit $X$ un schéma. Soit $Z \subset X$ un sous-schéma fermé
tel que le morphisme d'inclusion soit de présentation finie, c'est-à-dire que
le faisceau quasi-cohérent d'idéaux correspondant à $Z$ soit de
type fini. Soit $b : X' \to X$ l'éclatement de $X$ le long de $Z$
et soit $E = b^{-1}(Z)$ le diviseur exceptionnel. Voir
Diviseurs, section \ref{divisors-section-blowing-up}.
Dans cette situation et dans toute cette section, nous dirons que
\begin{equation}
\label{flat-equation-blow-up-square}
\vcenter{
\xymatrix{
E \ar[d] \ar[r] & X' \ar[d]^b \\
Z \ar[r] & X
}
}
\end{equation}
est un {\it carré d'éclatement}.

\begin{lemma}
\label{flat-lemma-blow-up-square-ph}
Soit $\mathcal{F}$ un faisceau sur un site $(\Sch/S)_{ph}$, voir
Topologies, définition \ref{topologies-definition-big-small-ph}.
Alors, pour tout carré d'éclatement (\ref{flat-equation-blow-up-square})
dans la catégorie $(\Sch/S)_{ph}$, le diagramme
$$
\xymatrix{
\mathcal{F}(E) & \mathcal{F}(X') \ar[l] \\
\mathcal{F}(Z) \ar[u] & \mathcal{F}(X) \ar[u] \ar[l]
}
$$
est cartésien dans la catégorie des ensembles.
\end{lemma}

\begin{proof}
Puisque $Z \amalg X' \to X$ est un morphisme propre surjectif,
on voit que $\{Z \amalg X' \to X\}$ est un recouvrement ph
(Topologies, lemme \ref{topologies-lemma-surjective-proper-ph}).
Nous avons
$$
(Z \amalg X') \times_X (Z \amalg X') =
Z \amalg E \amalg E \amalg X' \times_X X'
$$
Puisque $\mathcal{F}$ est un faisceau de Zariski, on voit que
$\mathcal{F}$ transforme les unions disjointes en produits.
La condition de faisceau pour le recouvrement
$\{Z \amalg X' \to X\}$ affirme donc que
$\mathcal{F}(X) \to \mathcal{F}(Z) \times \mathcal{F}(X')$
est injective et a pour image l'ensemble des couples $(t, s')$ tels que
(a) $t|_E = s'|_E$ et (b) $s'$ appartient à l'égalisateur des deux applications
$\mathcal{F}(X') \to \mathcal{F}(X' \times_X X')$.
Observons ensuite que le morphisme évident
$$
E \times_Z E \amalg X' \longrightarrow X' \times_X X'
$$
est un morphisme propre surjectif, puisque $b$ induit
un isomorphisme $X' \setminus E \to X \setminus Z$. Nous en déduisons que
$\mathcal{F}(X' \times_X X') \to
\mathcal{F}(E \times_Z E) \times \mathcal{F}(X')$ est injective.
Il s'ensuit que (a) $\Rightarrow$ (b), ce qui prouve le lemme.
\end{proof}

\begin{lemma}
\label{flat-lemma-thickening-ph}
Soit $\mathcal{F}$ un faisceau sur un site $(\Sch/S)_{ph}$
comme dans Topologies, définition \ref{topologies-definition-big-small-ph}.
Soit $X \to X'$ un morphisme de $(\Sch/S)_{ph}$ qui est
un épaississement. Alors
$\mathcal{F}(X') \to \mathcal{F}(X)$ est bijective.
\end{lemma}

\begin{proof}
Remarquons que $X \to X'$ est un morphisme propre surjectif et que
$X \times_{X'} X = X$.
La propriété de faisceau pour le recouvrement ph $\{X \to X'\}$
(Topologies, lemme
\ref{topologies-lemma-surjective-proper-ph})
permet de conclure.
\end{proof}








\section{Carrés de quasi-éclatement et topologie h}
\label{flat-section-blow-up-h}

\noindent
Considérons un carré d'éclatement (\ref{flat-equation-blow-up-square}).
Bien que le morphisme $b : X' \to X$
soit projectif (Diviseurs, lemme \ref{divisors-lemma-blowing-up-projective}),
il n'existe en général aucun moyen simple de garantir que
$b$ soit de présentation finie. Comme les recouvrements h sont construits
à l'aide de morphismes de présentation finie, il nous faut une variante.
Précisément, nous dirons qu'un diagramme commutatif
\begin{equation}
\label{flat-equation-almost-blow-up-square}
\vcenter{
\xymatrix{
E \ar[d] \ar[r] & X' \ar[d]^b \\
Z \ar[r] & X
}
}
\end{equation}
de schémas est un {\it carré de quasi-éclatement} si les conditions suivantes
sont satisfaites
\begin{enumerate}
\item $Z \to X$ est une immersion fermée de présentation finie,
\item $E = b^{-1}(Z)$ est un sous-schéma fermé localement principal de $X'$,
\item $b$ est propre et de présentation finie,
\item le sous-schéma fermé $X'' \subset X'$ défini par l'idéal quasi-cohérent
des sections de $\mathcal{O}_{X'}$ à support dans $E$
(Propriétés, lemme \ref{properties-lemma-sections-supported-on-closed-subset})
est l'éclatement de $X$ le long de $Z$.
\end{enumerate}
Il s'ensuit que le morphisme $b$ induit un isomorphisme
$X' \setminus E \to X \setminus Z$.
Pour quelques exemples très simples de carrés de quasi-éclatement, voir les
exemples \ref{flat-example-one-generator} et
\ref{flat-example-two-generators}.

\medskip\noindent
Le changement de base d'un éclatement n'est en général pas un éclatement,
mais les quasi-éclatements sont compatibles avec le changement de base.

\begin{lemma}
\label{flat-lemma-base-change-almost-blow-up}
Considérons un carré de quasi-éclatement (\ref{flat-equation-almost-blow-up-square}).
Soit $Y \to X$ un morphisme quelconque. Alors le changement de base
$$
\xymatrix{
Y \times_X E \ar[d] \ar[r] & Y \times_X X' \ar[d] \\
Y \times_X Z \ar[r] & Y
}
$$
est lui aussi un carré de quasi-éclatement.
\end{lemma}

\begin{proof}
Le morphisme $Y \times_X X' \to Y$ est propre et de présentation finie
d'après Morphismes, lemmes \ref{morphisms-lemma-base-change-proper} et
\ref{morphisms-lemma-base-change-finite-presentation}.
Le morphisme $Y \times_X Z \to Y$ est une immersion fermée
(Morphismes, lemme \ref{morphisms-lemma-base-change-closed-immersion}) de
présentation finie. L'image inverse de $Y \times_X Z$ dans $Y \times_X X'$
est égale à l'image inverse de $E$ dans $Y \times_X X'$ et est donc
localement principale
(Diviseurs, lemme \ref{divisors-lemma-pullback-locally-principal}).
Soient $X'' \subset X'$, resp.\ $Y'' \subset Y \times_X X'$, les sous-schémas fermés
correspondant aux idéaux quasi-cohérents de sections de
$\mathcal{O}_{X'}$, resp.\ $\mathcal{O}_{Y \times_Y X'}$
à support dans $E$, resp.\ $Y \times_X E$.
Manifestement, $Y'' \subset Y \times_X X''$ est le sous-schéma fermé
correspondant à l'idéal quasi-cohérent des sections de
$\mathcal{O}_{Y \times_Y X''}$ à support dans $Y \times_X (E \cap X'')$.
Ainsi, $Y''$ est le transformé strict de $Y$ relativement à l'éclatement
$X'' \to X$, voir
Diviseurs, définition \ref{divisors-definition-strict-transform}.
Ainsi, d'après Diviseurs, lemme \ref{divisors-lemma-strict-transform},
nous voyons que $Y''$ est l'éclatement de centre $Y \times_X Z$ sur $Y$.
\end{proof}

\noindent
On peut rétrécir les carrés de quasi-éclatement.

\begin{lemma}
\label{flat-lemma-shrink-almost-blow-up}
Considérons un carré de quasi-éclatement (\ref{flat-equation-almost-blow-up-square}).
Soit $W \to X'$ une immersion fermée de présentation finie.
Les assertions suivantes sont équivalentes
\begin{enumerate}
\item $X' \setminus E$ est schématiquement contenu dans $W$,
\item l'éclatement $X''$ de $X$ le long de $Z$ est schématiquement contenu dans $W$,
\item le diagramme
$$
\xymatrix{
E \cap W \ar[d] \ar[r] & W \ar[d] \\
Z \ar[r] & X
}
$$
est un carré de quasi-éclatement. Ici, $E \cap W$ est l'intersection
schématique.
\end{enumerate}
\end{lemma}

\begin{proof}
Supposons (1). Alors la surjection $\mathcal{O}_{X'} \to \mathcal{O}_W$
est un isomorphisme au-dessus de l'ouvert $X' \subset E$. Comme le faisceau d'idéaux
de $X'' \subset X'$ est constitué des sections de $\mathcal{O}_{X'}$ à support dans $E$
(par notre définition des carrés de quasi-éclatement),
nous concluons que (2) est vraie. Si (2) est vraie, alors (3) l'est.
Si (3) est vraie, alors (1) l'est parce que $X'' \cap (X' \setminus E)$
est isomorphe à $X \setminus Z$, qui est à son tour isomorphe
à $X' \setminus E$.
\end{proof}

\noindent
L'éclatement proprement dit est la limite des rétrécissements de tout quasi-éclatement donné.

\begin{lemma}
\label{flat-lemma-blow-up-limit-almost-blow-up}
Considérons un carré de quasi-éclatement (\ref{flat-equation-almost-blow-up-square})
où $X$ est quasi-compact et quasi-séparé. Alors l'éclatement $X''$ de $X$
le long de $Z$ peut s'écrire
$$
X'' = \lim X'_i
$$
où la limite porte sur le système filtrant des sous-schémas fermés
$X'_i \subset X'$ de présentation finie satisfaisant aux conditions équivalentes
du lemme \ref{flat-lemma-shrink-almost-blow-up}.
\end{lemma}

\begin{proof}
Soit $\mathcal{I} \subset \mathcal{O}_{X'}$ le faisceau quasi-cohérent
d'idéaux correspondant à $X''$. D'après
Propriétés, lemme \ref{properties-lemma-quasi-coherent-colimit-finite-type},
nous pouvons écrire $\mathcal{I}$ comme la colimite filtrante
$\mathcal{I} = \colim \mathcal{I}_i$ de ses sous-modules quasi-cohérents
de type fini. Comme ces modules correspondent
$1$-à-$1$ aux sous-schémas fermés $X'_i$, la démonstration est terminée.
\end{proof}

\noindent
Il existe des carrés de quasi-éclatement.

\begin{lemma}
\label{flat-lemma-almost-blow-up-square}
Soit $X$ un schéma quasi-compact et quasi-séparé. Soit $Z \subset X$
un sous-schéma fermé défini par un faisceau quasi-cohérent d'idéaux
de type fini. Alors il existe un carré de quasi-éclatement comme dans
(\ref{flat-equation-almost-blow-up-square}).
\end{lemma}

\begin{proof}
Nous pouvons écrire $X = \lim X_i$ comme la limite d'un système projectif filtrant
de schémas noethériens dont les morphismes de transition sont affines, voir
Limites, proposition \ref{limits-proposition-approximate}.
Nous pouvons trouver un indice $i$ et une immersion fermée $Z_i \to X_i$
dont le changement de base à $X$ est l'immersion fermée $Z \to X$.
Voir Limites, lemmes \ref{limits-lemma-descend-finite-presentation} et
\ref{limits-lemma-descend-closed-immersion-finite-presentation}.
Soit $b_i : X'_i \to X_i$ l'éclatement de centre $Z_i$.
Cela donne un carré d'éclatement
$$
\xymatrix{
E_i \ar[r] \ar[d] & X'_i \ar[d]^{b_i} \\
Z_i \ar[r] & X_i
}
$$
où tous les morphismes sont des morphismes de type fini entre schémas noethériens
et sont donc de présentation finie. Il s'agit donc d'un carré de quasi-éclatement.
D'après le lemme \ref{flat-lemma-base-change-almost-blow-up},
le changement de base de ce diagramme à $X$ fournit le
carré de quasi-éclatement recherché.
\end{proof}

\noindent
Les carrés de quasi-éclatement sont uniques à rétrécissement près, comme dans le
lemme \ref{flat-lemma-shrink-almost-blow-up}.

\begin{lemma}
\label{flat-lemma-almost-blow-up-unique}
Soit $X$ un schéma quasi-compact et quasi-séparé et soit $Z \subset X$
un sous-schéma fermé défini par un faisceau quasi-cohérent d'idéaux de type fini.
Supposons donnés des carrés de quasi-éclatement (\ref{flat-equation-almost-blow-up-square})
$$
\xymatrix{
E_k \ar[r] \ar[d] & X_k' \ar[d] \\
Z \ar[r] & X
}
$$
pour $k = 1, 2$ ; alors il existe un carré de quasi-éclatement
$$
\xymatrix{
E \ar[r] \ar[d] & X' \ar[d] \\
Z \ar[r] & X
}
$$
et des immersions fermées $i_k : X' \to X'_k$ au-dessus de $X$
telles que $E = i_k^{-1}(E_k)$.
\end{lemma}

\begin{proof}
Notons $X'' \to X$ l'éclatement de centre $Z$ dans $X$.
Nous considérons $X''$ comme un sous-schéma fermé de $X'_1$ et de $X'_2$.
Écrivons $X'' = \lim X'_{1, i}$ comme dans le
lemme \ref{flat-lemma-blow-up-limit-almost-blow-up}.
D'après Limites, proposition
\ref{limits-proposition-characterize-locally-finite-presentation},
il existe un $i$ et un morphisme $h : X'_{1, i} \to X'_2$ qui coïncide
avec les inclusions $X'' \subset X'_{1, i}$ et $X'' \subset X'_2$.
D'après Limites, le lemme \ref{limits-lemma-eventually-closed-immersion},
la restriction de $h$ à $X'_{1, i'}$ est une immersion fermée
pour un certain $i' \geq i$. Ceci achève la démonstration.
\end{proof}

\noindent
Nos techniques de platification par éclatement s'étendent aux
quasi-éclatements dans les situations favorables.

\begin{lemma}
\label{flat-lemma-flat-after-almost-blowing-up}
Soit $Y$ un schéma quasi-compact et quasi-séparé.
Soit $X$ un schéma de présentation finie sur $Y$.
Soit $V \subset Y$ un ouvert quasi-compact tel que
$X_V \to V$ soit plat. Alors il existe un diagramme commutatif
$$
\xymatrix{
E \ar[ddd] \ar[rd] & & & D \ar[lll] \ar[ddd] \ar[ld] \\
& Y' \ar[d] & X' \ar[l] \ar[d] \\
& Y & X \ar[l] \\
Z \ar[ru] & & & T \ar[lll] \ar[lu]
}
$$
dont les carrés de droite et de gauche sont des carrés de quasi-éclatement,
dont les carrés inférieur et supérieur sont cartésiens, et tel que
$Z \cap V = \emptyset$, et
tel que $X' \to Y'$ soit plat (et de présentation finie).
\end{lemma}

\begin{proof}
Si $Y$ est un schéma noethérien, ce lemme découle immédiatement
du lemme \ref{flat-lemma-flat-after-blowing-up},
car, dans ce cas, les carrés d'éclatement sont des carrés de quasi-éclatement
(nous utilisons aussi le fait que les transformés stricts sont des éclatements).
Le cas général se ramène au cas noethérien par
approximation noethérienne absolue.

\medskip\noindent
Nous pouvons écrire $Y = \lim Y_i$ comme la limite d'un système projectif filtrant
de schémas noethériens dont les morphismes de transition sont affines, voir
Limites, proposition \ref{limits-proposition-approximate}.
Nous pouvons trouver un indice $i$ et un morphisme $X_i \to Y_i$
de présentation finie dont le changement de base à $Y$ est $X \to Y$.
Voir Limites, lemme \ref{limits-lemma-descend-finite-presentation}.
Quitte à augmenter $i$, nous pouvons supposer que $V$ est l'image réciproque
d'un sous-schéma ouvert $V_i \subset Y_i$, voir
Limites, lemme \ref{limits-lemma-descend-opens}.
Enfin, quitte à augmenter $i$, nous pouvons supposer que $X_{i, V_i} \to V_i$
est plat, voir
Limites, lemme \ref{limits-lemma-descend-flat-finite-presentation}.
Dans le cas noethérien, nous pouvons construire un diagramme comme dans le
lemme pour $X_i \to Y_i \supset V_i$. Le changement de base de ce
diagramme par $Y \to Y_i$ donne la solution. Utilisons le fait que le changement de base
préserve les propriétés des morphismes, voir Morphismes, lemmes
\ref{morphisms-lemma-base-change-proper},
\ref{morphisms-lemma-base-change-finite-presentation},
\ref{morphisms-lemma-base-change-closed-immersion}, et
\ref{morphisms-lemma-base-change-flat},
ainsi que le fait que le changement de base d'un carré de
quasi-éclatement est un carré de quasi-éclatement, voir
le lemme \ref{flat-lemma-base-change-almost-blow-up}.
\end{proof}

\begin{lemma}
\label{flat-lemma-blow-up-square-h}
Soit $\mathcal{F}$ un faisceau sur l'un des sites $(\Sch/S)_h$
construits dans la définition \ref{flat-definition-big-small-h}.
Alors, pour tout carré de quasi-éclatement (\ref{flat-equation-almost-blow-up-square})
dans la catégorie $(\Sch/S)_h$, le diagramme
$$
\xymatrix{
\mathcal{F}(E) & \mathcal{F}(X') \ar[l] \\
\mathcal{F}(Z) \ar[u] & \mathcal{F}(X) \ar[u] \ar[l]
}
$$
est cartésien dans la catégorie des ensembles.
\end{lemma}

\begin{proof}
Puisque $Z \amalg X' \to X$ est un morphisme propre surjectif
de présentation finie, on voit que $\{Z \amalg X' \to X\}$ est un recouvrement h
(lemme \ref{flat-lemma-surjective-proper-finite-presentation-h}).
Nous avons
$$
(Z \amalg X') \times_X (Z \amalg X') =
Z \amalg E \amalg E \amalg X' \times_X X'
$$
Puisque $\mathcal{F}$ est un faisceau pour la topologie de Zariski, on voit que
$\mathcal{F}$ envoie les unions disjointes sur des produits.
Ainsi, la condition de faisceau pour le recouvrement
$\{Z \amalg X' \to X\}$ dit que
$\mathcal{F}(X) \to \mathcal{F}(Z) \times \mathcal{F}(X')$
est injective et a pour image l'ensemble des couples $(t, s')$ tels que
(a) $t|_E = s'|_E$ et (b) $s'$ appartient à l'égalisateur des deux applications
$\mathcal{F}(X') \to \mathcal{F}(X' \times_X X')$.
Observons ensuite que le morphisme évident
$$
E \times_Z E \amalg X' \longrightarrow X' \times_X X'
$$
est un morphisme propre surjectif de présentation finie puisque $b$ induit
un isomorphisme $X' \setminus E \to X \setminus Z$. Nous en déduisons que
$\mathcal{F}(X' \times_X X') \to
\mathcal{F}(E \times_Z E) \times \mathcal{F}(X')$ est injective.
Il s'ensuit que (a) $\Rightarrow$ (b), ce qui prouve le lemme.
\end{proof}

\begin{lemma}
\label{flat-lemma-thickening-h}
Soit $\mathcal{F}$ un faisceau sur l'un des sites $(\Sch/S)_h$
construits dans la définition \ref{flat-definition-big-small-h}.
Soit $X \to X'$ un morphisme de $(\Sch/S)_h$ qui est
un épaississement et qui est de présentation finie. Alors
$\mathcal{F}(X') \to \mathcal{F}(X)$ est bijective.
\end{lemma}

\begin{proof}
Première démonstration. Remarquons que $X \to X'$ est un morphisme propre surjectif
de présentation finie et que $X \times_{X'} X = X$.
Par la propriété de faisceau pour le recouvrement h $\{X \to X'\}$
(lemme \ref{flat-lemma-surjective-proper-finite-presentation-h}),
nous pouvons conclure.

\medskip\noindent
Deuxième démonstration (sotte). L'éclatement de $X'$ le long de $X$ est le schéma vide.
Cela tient au fait que l'algèbre affine d'éclatement $A[\frac{I}{a}]$
(Algèbre, section \ref{algebra-section-blow-up})
est nulle si $a$ est un élément nilpotent de $A$. Les détails sont omis.
On obtient donc un carré de quasi-éclatement de la forme
$$
\xymatrix{
\emptyset \ar[r] \ar[d] & \emptyset \ar[d] \\
X \ar[r] & X'
}
$$
Puisque $\mathcal{F}$ est un faisceau, $\mathcal{F}(\emptyset)$
est un singleton.
En appliquant le lemme \ref{flat-lemma-blow-up-square-h}, on obtient la conclusion.
\end{proof}

\begin{proposition}
\label{flat-proposition-check-h}
Soit $\mathcal{F}$ un préfaisceau sur l'un des sites $(\Sch/S)_h$
construits dans la définition \ref{flat-definition-big-small-h}.
Alors $\mathcal{F}$ est un faisceau si et seulement si les
conditions suivantes sont satisfaites
\begin{enumerate}
\item $\mathcal{F}$ est un faisceau pour la topologie de Zariski,
\item pour tout morphisme $f : X \to Y$ de $(\Sch/S)_h$ tel que $Y$ soit affine
et que $f$ soit surjectif, plat, propre et de présentation finie,
$\mathcal{F}(Y)$ est l'égalisateur des deux applications
$\mathcal{F}(X) \to \mathcal{F}(X \times_Y X)$,
\item pour tout carré de quasi-éclatement (\ref{flat-equation-almost-blow-up-square})
avec $X$ affine dans la catégorie $(\Sch/S)_h$, le diagramme
$$
\xymatrix{
\mathcal{F}(E) & \mathcal{F}(X') \ar[l] \\
\mathcal{F}(Z) \ar[u] & \mathcal{F}(X) \ar[u] \ar[l]
}
$$
est cartésien dans la catégorie des ensembles.
\end{enumerate}
\end{proposition}

\begin{proof}
Supposons que $\mathcal{F}$ soit un faisceau. La condition (1) est vérifiée car
un recouvrement de Zariski est un recouvrement h, voir
le lemme \ref{flat-lemma-zariski-h}.
La condition (2) est vérifiée car, pour $f$ comme en (2),
$\{X \to Y\}$ est un recouvrement fppf (c'est clair)
et donc un recouvrement h, voir le lemme \ref{flat-lemma-zariski-h}.
La condition (3) résulte du lemme \ref{flat-lemma-blow-up-square-h}.

\medskip\noindent
Réciproquement, supposons que $\mathcal{F}$ satisfasse (1), (2) et (3).
Nous allons montrer que $\mathcal{F}$ est un faisceau en appliquant
le lemme \ref{flat-lemma-characterize-sheaf-h}. Considérons
un morphisme surjectif, de présentation finie et propre
$f : X \to Y$ dans $(\Sch/S)_h$, avec $Y$ affine. Il suffit de montrer
que $\mathcal{F}(Y)$ est l'égalisateur des deux applications
$\mathcal{F}(X) \to \mathcal{F}(X \times_Y X)$.

\medskip\noindent
Supposons d'abord qu'en outre $f : X \to Y$ soit une immersion fermée
(autrement dit, $f$ est un épaississement). Alors l'éclatement de $Y$ le long de $X$
est le schéma vide, ce qui fournit un carré de quasi-éclatement
dont les sommets sont $\emptyset, \emptyset, X, Y$
(comparer avec la deuxième démonstration du lemme \ref{flat-lemma-thickening-h}).
La condition (3) nous dit donc que
$$
\xymatrix{
\mathcal{F}(\emptyset) & \mathcal{F}(\emptyset) \ar[l] \\
\mathcal{F}(X) \ar[u] & \mathcal{F}(Y) \ar[u] \ar[l]
}
$$
est cartésien dans la catégorie des ensembles. Puisque $\mathcal{F}$ est un faisceau
pour la topologie de Zariski, on voit que $\mathcal{F}(\emptyset)$
est un singleton. On en déduit que $\mathcal{F}(X) = \mathcal{F}(Y)$.

\medskip\noindent
Intermède A : soit $T \to T'$ un morphisme de $(\Sch/S)_h$ qui est
un épaississement et qui est de présentation finie. Alors
$\mathcal{F}(T') \to \mathcal{F}(T)$ est bijective.
En effet, choisissons un recouvrement ouvert affine $T' = \bigcup T'_i$
et posons $T_i = T \times_{T'} T'_i$ ; ce sont les ouverts affines correspondants
de $T$. Alors $\mathcal{F}(T'_i) \to \mathcal{F}(T_i)$
est bijective pour tout $i$ d'après le résultat du paragraphe précédent.
La propriété de faisceau pour la topologie de Zariski montre que
$\mathcal{F}(T') \to \mathcal{F}(T)$ est injective. En répétant le
raisonnement, on constate qu'elle est bijective. Nous omettons les détails mineurs.

\medskip\noindent
Intermède B : considérons un carré de quasi-éclatement
(\ref{flat-equation-almost-blow-up-square}) dans la catégorie $(\Sch/S)_h$.
Nous affirmons alors que le diagramme
$$
\xymatrix{
\mathcal{F}(E) & \mathcal{F}(X') \ar[l] \\
\mathcal{F}(Z) \ar[u] & \mathcal{F}(X) \ar[u] \ar[l]
}
$$
est cartésien dans la catégorie des ensembles. Cela résulte de la condition
(3) de la manière suivante : on choisit un recouvrement ouvert affine de $X$ et on raisonne
comme dans l'intermède A. Nous omettons les détails.

\medskip\noindent
Ensuite, soit $f : X \to Y$ un morphisme surjectif, propre et de présentation finie
dans $(\Sch/S)_h$, avec $Y$ affine. Choisissons une stratification par platitude générique
$$
Y \supset Y_0 \supset Y_1 \supset \ldots \supset Y_t = \emptyset
$$
comme dans Compléments sur les morphismes, lemme
\ref{more-morphisms-lemma-generic-flatness-stratification-scheme},
pour $f : X \to Y$.
Nous allons utiliser toutes les propriétés de cette stratification
sans les mentionner davantage. Posons $X_0 = X \times_Y Y_0$.
D'après l'intermède B, nous avons $\mathcal{F}(Y_0) = \mathcal{F}(Y)$,
$\mathcal{F}(X_0) = \mathcal{F}(X)$, et
$\mathcal{F}(X_0 \times_{Y_0} X_0) = \mathcal{F}(X \times_Y X)$.

\medskip\noindent
Nous allons démontrer le résultat par récurrence sur $t$. Si $t = 1$,
alors $X_0 \to Y_0$ est surjectif, propre, plat et de présentation finie,
et le résultat découle de la propriété (2).
Pour $t > 1$, nous pouvons remplacer $Y$ par $Y_0$ et $X$ par $X_0$
(voir ci-dessus) et supposer que $Y = Y_0$.

\medskip\noindent
Considérons le sous-schéma ouvert quasi-compact
$V = Y \setminus Y_1 = Y_0 \setminus Y_1$.
Choisissons un diagramme
$$
\xymatrix{
E \ar[ddd] \ar[rd] & & & D \ar[lll] \ar[ddd] \ar[ld] \\
& Y' \ar[d] & X' \ar[l] \ar[d] \\
& Y & X \ar[l] \\
Z \ar[ru] & & & T \ar[lll] \ar[lu]
}
$$
comme dans le lemme \ref{flat-lemma-flat-after-almost-blowing-up}
pour $f : X \to Y \supset V$. Alors $f' : X' \to Y'$ est plat et de
présentation finie. De plus, $f'$ est propre (utiliser
Morphismes, lemmes \ref{morphisms-lemma-composition-proper} et
\ref{morphisms-lemma-image-proper-scheme-closed} pour le voir).
Ainsi, l'image $W = f'(X') \subset Y'$ est un sous-schéma ouvert
(Morphismes, lemme \ref{morphisms-lemma-fppf-open}) et fermé
de $Y'$. Observons que $Y' \setminus E$ est contenu
dans $W$. D'après le lemme \ref{flat-lemma-shrink-almost-blow-up},
cela signifie que nous pouvons remplacer $Y'$ par $W$
dans le diagramme ci-dessus. Autrement dit, nous pouvons supposer, et supposons,
que $f'$ est surjectif. À ce stade, nous savons que
$$
\vcenter{
\xymatrix{
\mathcal{F}(E) & \mathcal{F}(Y') \ar[l] \\
\mathcal{F}(Z) \ar[u] & \mathcal{F}(Y) \ar[u] \ar[l]
}
}
\quad\text{et}\quad
\vcenter{
\xymatrix{
\mathcal{F}(D) & \mathcal{F}(X') \ar[l] \\
\mathcal{F}(T) \ar[u] & \mathcal{F}(X) \ar[u] \ar[l]
}
}
$$
sont cartésiens d'après l'intermède B. Remarquons que
$Z \cap Y_1 \to Z$ est un épaississement de
présentation finie (puisque $Z$ est ensemblistement contenu dans $Y_1$
comme sous-schéma fermé de $Y$ disjoint de $V$).
Nous obtenons donc une filtration
$$
Z \supset Z \cap Y_1 \supset Z \cap Y_2 \subset \ldots \subset
Z \cap Y_t = \emptyset
$$
comme ci-dessus pour la restriction $T = Z \times_Y X \to Z$ de $f$ à $T$.
Par l'hypothèse de récurrence, nous trouvons donc que
$\mathcal{F}(Z) \to \mathcal{F}(T)$
est une application injective d'ensembles dont l'image est l'égalisateur
des deux applications $\mathcal{F}(T) \to
\mathcal{F}(T \times_Z T)$.

\medskip\noindent
Soit $s \in \mathcal{F}(X)$ un élément de l'égalisateur des deux applications
$\mathcal{F}(X) \to \mathcal{F}(X \times_Y X)$.
D'après ce qui précède, la restriction $s|_T$
provient d'un unique élément $t \in \mathcal{F}(Z)$,
et, de même, la restriction $s|_{X'}$
provient d'un unique élément $t' \in \mathcal{F}(Y')$.
En chassant les sections au moyen des applications de restriction de $\mathcal{F}$
correspondant aux flèches du grand diagramme commutatif ci-dessus,
le lecteur constate que $t$ et $t'$ se restreignent au même élément de
$\mathcal{F}(E)$, car ils se restreignent au même élément de
$\mathcal{F}(D)$ et que (2) s'applique ; nous utilisons ici le fait que $D \to E$ est
surjectif, plat, propre et de présentation finie en tant que restriction
de $X' \to Y'$. Ainsi, le premier des deux
carrés cartésiens affichés ci-dessus nous donne une unique section
$u \in \mathcal{F}(Y)$
dont les restrictions valent respectivement $t$ et $t'$ sur $Z$ et $Y'$.
Pour voir que $u$ se restreint à $s$ sur $X$, utiliser le second diagramme.
\end{proof}

\begin{example}
\label{flat-example-one-generator}
Soit $A$ un anneau. Soit $f \in A$ un élément. Soit $J \subset A$
un idéal de type fini annulé par une puissance de $f$.
Alors
$$
\xymatrix{
E = \Spec(A/fA + J) \ar[r] \ar[d] & \Spec(A/J) = X' \ar[d] \\
Z = \Spec(A/fA) \ar[r] & \Spec(A) = X
}
$$
est un carré de quasi-éclatement.
\end{example}

\begin{example}
\label{flat-example-two-generators}
Soit $A$ un anneau. Soient $f_1, f_2 \in A$ des éléments.
$$
\xymatrix{
E = \text{Proj}(A/(f_1, f_2)[T_0, T_1]) \ar[r] \ar[d] &
\text{Proj}(A[T_0, T_1]/(f_2 T_0 - f_1 T_1) = X' \ar[d] \\
Z = \Spec(A/(f_1, f_2)) \ar[r] & \Spec(A) = X
}
$$
est un carré de quasi-éclatement.
\end{example}

\begin{lemma}
\label{flat-lemma-refine-check-h}
Soit $\mathcal{F}$ un préfaisceau sur l'un des sites $(\Sch/S)_h$
construits dans la définition \ref{flat-definition-big-small-h}.
Alors $\mathcal{F}$ est un faisceau si et seulement si les
conditions suivantes sont satisfaites
\begin{enumerate}
\item $\mathcal{F}$ est un faisceau pour la topologie de Zariski,
\item étant donné un morphisme $f : X \to Y$ de $(\Sch/S)_h$ avec $Y$ affine
et $f$ surjectif, plat, propre et de présentation finie, alors
$\mathcal{F}(Y)$ est l'égalisateur des deux applications
$\mathcal{F}(X) \to \mathcal{F}(X \times_Y X)$,
\item $\mathcal{F}$ transforme un carré de quasi-éclatement comme dans
l'exemple \ref{flat-example-one-generator} dans la catégorie $(\Sch/S)_h$
en un diagramme cartésien d'ensembles, et
\item $\mathcal{F}$ transforme un carré de quasi-éclatement comme dans
l'exemple \ref{flat-example-two-generators} dans la catégorie $(\Sch/S)_h$
en un diagramme cartésien d'ensembles.
\end{enumerate}
\end{lemma}

\begin{proof}
D'après la proposition \ref{flat-proposition-check-h}, il suffit de montrer qu'étant donné
un carré de quasi-éclatement (\ref{flat-equation-almost-blow-up-square})
avec $X$ affine dans la catégorie $(\Sch/S)_h$, le diagramme
$$
\xymatrix{
\mathcal{F}(E) & \mathcal{F}(X') \ar[l] \\
\mathcal{F}(Z) \ar[u] & \mathcal{F}(X) \ar[u] \ar[l]
}
$$
est cartésien dans la catégorie des ensembles. L'idée générale de la démonstration
consiste à dominer le morphisme par d'autres carrés de quasi-éclatement auxquels
nous pouvons appliquer localement les hypothèses (3) et (4).

\medskip\noindent
Supposons que nous ayons un carré de quasi-éclatement
(\ref{flat-equation-almost-blow-up-square}) dans la catégorie $(\Sch/S)_h$,
un recouvrement ouvert $X = \bigcup U_i$ et des recouvrements ouverts
$U_i \cap U_j = \bigcup U_{ijk}$ tels que les diagrammes
$$
\vcenter{
\xymatrix{
\mathcal{F}(E \cap b^{-1}(U_i)) & \mathcal{F}(b^{-1}(U_i)) \ar[l] \\
\mathcal{F}(Z \cap U_i) \ar[u] & \mathcal{F}(U_i) \ar[u] \ar[l]
}
}
\quad\text{et}\quad
\vcenter{
\xymatrix{
\mathcal{F}(E \cap b^{-1}(U_{ijk})) &
\mathcal{F}(b^{-1}(U_{ijk})) \ar[l] \\
\mathcal{F}(Z \cap U_{ijk}) \ar[u] &
\mathcal{F}(U_{ijk}) \ar[u] \ar[l]
}
}
$$
sont cartésiens ; il en va alors de même pour
$$
\xymatrix{
\mathcal{F}(E) & \mathcal{F}(X') \ar[l] \\
\mathcal{F}(Z) \ar[u] & \mathcal{F}(X) \ar[u] \ar[l]
}
$$
Cela découle du fait que $\mathcal{F}$ est un faisceau pour la topologie de Zariski.

\medskip\noindent
En particulier, si nous avons un carré d'éclatement
(\ref{flat-equation-almost-blow-up-square}) tel que $b : X' \to X$
soit une immersion fermée et que $Z$ soit un sous-schéma fermé
localement principal, alors nous voyons que
$\mathcal{F}(X) = \mathcal{F}(X') \times_{\mathcal{F}(E)} \mathcal{F}(Z)$.
En effet, sur un recouvrement de $X$ par des ouverts affines, nous obtenons un carré de quasi-éclatement comme en (3).

\medskip\noindent
Soient $Z \subset X$, $E_k \subset X'_k \to X$, $E \subset X' \to X$, et
$i_k : X' \to X'_k$ comme dans l'énoncé du
lemme \ref{flat-lemma-almost-blow-up-unique}. Alors
$$
\xymatrix{
E \ar[d] \ar[r] & X' \ar[d] \\
E_k \ar[r] & X'_k
}
$$
est un carré de quasi-éclatement du type considéré au paragraphe
précédent. Ainsi
$$
\mathcal{F}(X'_k) = \mathcal{F}(X') \times_{\mathcal{F}(E)} \mathcal{F}(E_k)
$$
pour $k = 1, 2$, d'après le résultat du paragraphe précédent. Il s'ensuit que
$$
\mathcal{F}(X) \longrightarrow
\mathcal{F}(X'_k) \times_{\mathcal{F}(E_k)} \mathcal{F}(Z)
$$
est bijective pour $k = 1$ si et seulement si elle est bijective pour $k = 2$.
Ainsi, étant donnée une immersion fermée $Z \to X$ de présentation finie,
avec $X$ quasi-compact et quasi-séparé, le fait que
$\mathcal{F}(X) = \mathcal{F}(X') \times_{\mathcal{F}(E)} \mathcal{F}(Z)$
soit vérifiée ou non est indépendant du choix du carré de quasi-éclatement
(\ref{flat-equation-almost-blow-up-square}) effectué.
(De plus, d'après le lemme \ref{flat-lemma-almost-blow-up-square},
il existe effectivement un carré de
quasi-éclatement pour $Z \subset X$.)

\medskip\noindent
Enfin, considérons un objet affine $X$ de $(\Sch/S)_h$
et une immersion fermée $Z \to X$ de présentation finie.
Nous allons démontrer la propriété voulue pour la paire $(X, Z)$
par récurrence sur le nombre $r$ de générateurs de l'idéal
définissant $Z$ dans $X$. Si le nombre de générateurs est $\leq 2$,
nous pouvons choisir notre carré de quasi-éclatement comme
dans l'exemple \ref{flat-example-two-generators}
et conclure par l'hypothèse (4).

\medskip\noindent
Étape de récurrence. Supposons $X = \Spec(A)$ et $Z = \Spec(A/(f_1, \ldots, f_r))$
avec $r > 2$. Choisissons un carré d'éclatement
(\ref{flat-equation-almost-blow-up-square})
pour la paire $(X, Z)$. Posons $Z_1 = \Spec(A/(f_1, f_2))$
et soit
$$
\xymatrix{
E_1 \ar[d] \ar[r] & Y \ar[d] \\
Z_1 \ar[r] & X
}
$$
le carré de quasi-éclatement construit dans
l'exemple \ref{flat-example-two-generators}.
D'après le lemme \ref{flat-lemma-base-change-almost-blow-up}, les changements de base
$$
(I)
\vcenter{
\xymatrix{
Y \times_X E \ar[r] \ar[d] & Y \times_X X' \ar[d] \\
Y \times_X Z \ar[r] & Y
}
}
\quad\text{et}\quad
(II)
\vcenter{
\xymatrix{
E \ar[r] \ar[d] & Z_1 \times_X X' \ar[d] \\
Z \ar[r] & Z_1
}
}
$$
sont des carrés de quasi-éclatement. L'idéal de $Z$ dans $Z_1$ est engendré
par $r - 2$ éléments. L'idéal de $Y \times_X Z$
est engendré par les images réciproques de $f_1, \ldots, f_r$ dans $Y$. Localement sur $Y$,
l'idéal engendré par $f_1, f_2$ peut être engendré par
un élément ; ainsi, $Y \times_X Z$ est, sur des ouverts affines de $Y$,
défini par au plus $r - 1$ éléments.
Par les hypothèses de récurrence et la discussion ci-dessus,
$$
\mathcal{F}(Y) =
\mathcal{F}(Y \times_X X') \times_{\mathcal{F}(Y \times_X E)}
\mathcal{F}(Y \times_X Z)
$$
et
$$
\mathcal{F}(Z_1) =
\mathcal{F}(Z_1 \times_X X') \times_{\mathcal{F}(E)}
\mathcal{F}(Z)
$$
Par l'hypothèse (4), nous avons
$$
\mathcal{F}(X) =
\mathcal{F}(Y) \times_{\mathcal{F}(E_1)} \mathcal{F}(Z_1)
$$
Supposons maintenant que nous ayons une paire $(s', t)$ avec $s' \in \mathcal{F}(X')$
et $t \in \mathcal{F}(Z)$ ayant la même restriction dans $\mathcal{F}(E)$.
Alors $(s'|{Z_1 \times_X X'}, t)$ est l'image d'un unique élément
$t_1 \in \mathcal{F}(Z_1)$. De même,
$(s'|_{Y \times_X X'}, t|_{Y \times_X Z})$ est l'image d'un
unique élément $s_Y \in \mathcal{F}(Y)$.
Nous affirmons que $s_Y$ et $t_1$ se restreignent au même élément
de $\mathcal{F}(E_1)$. Cela est vrai parce que le
carré de quasi-éclatement
$$
\xymatrix{
E_1 \times_X E \ar[r] \ar[d] & E_1 \times_X X' \ar[d] \\
E_1 \times_X Z \ar[r] & E_1
}
$$
est le changement de base du carré de quasi-éclatement (I) par $E_1 \to Y$ et
le changement de base du carré de quasi-éclatement (II) par $E_1 \to Z_1$, et
parce que les paires de sections utilisées pour construire $s_Y$ et $t_1$ coïncident.
Ainsi, par la troisième égalité de produits fibrés, nous voyons qu'il existe
un unique $s \in \mathcal{F}(X)$ dont l'image est $s_Y$ dans $\mathcal{F}(Y)$
et $t_1$ dans $\mathcal{F}(Z)$.
Nous omettons de vérifier que l'image de $s$ est $s'$ dans $\mathcal{F}(X')$
et $t$ dans $\mathcal{F}(Z)$ ; indication : utiliser l'unicité de $s$ que nous venons
de construire et raisonner localement sur des ouverts affines.
\end{proof}

\begin{lemma}
\label{flat-lemma-refine-check-h-stack}
Soit $p : \mathcal{S} \to (\Sch/S)_h$ une catégorie fibrée en groupoïdes.
Alors $\mathcal{S}$ est un champ en groupoïdes si et seulement si les conditions
suivantes sont satisfaites
\begin{enumerate}
\item $\mathcal{S}$ est un champ en groupoïdes pour la topologie de Zariski,
\item étant donné un morphisme $f : X \to Y$ de $(\Sch/S)_h$ avec $Y$ affine
et $f$ surjectif, plat, propre et de présentation finie, alors
$$
\mathcal{S}_Y \longrightarrow
\mathcal{S}_X \times_{\mathcal{S}_{X \times_Y X}} \mathcal{S}_X
$$
est une équivalence de catégories,
\item pour un carré de quasi-éclatement comme dans
l'exemple \ref{flat-example-one-generator} ou \ref{flat-example-two-generators}
dans la catégorie $(\Sch/S)_h$, le foncteur
$$
\mathcal{S}_X \longrightarrow
\mathcal{S}_Z \times_{\mathcal{S}_E} \mathcal{S}_{X'}
$$
est une équivalence de catégories.
\end{enumerate}
\end{lemma}

\begin{proof}
Ce lemme est une conséquence formelle du lemme \ref{flat-lemma-refine-check-h}
et de notre définition des champs en groupoïdes.
Par exemple, supposons (1), (2), (3). Pour montrer que
$\mathcal{S}$ est un champ, nous devons
démontrer la descente des morphismes et des objets, voir
Champs, définition \ref{stacks-definition-stack-in-groupoids}.

\medskip\noindent
Si $x, y$ sont des objets de $\mathcal{S}$ au-dessus d'un objet $U$ de
$(\Sch/S)_h$, alors nos hypothèses impliquent que $\mathit{Isom}(x, y)$
est un préfaisceau sur $(\Sch/U)_h$ qui satisfait
(1), (2), (3) et (4) du lemme \ref{flat-lemma-refine-check-h},
et est donc un faisceau. Nous omettons quelques détails.

\medskip\noindent
Soit $\{U_i \to U\}_{i \in I}$ un recouvrement de $(\Sch/S)_h$.
Soit $(x_i, \varphi_{ij})$ une donnée de descente dans $\mathcal{S}$
relative à la famille $\{U_i \to U\}_{i \in I}$, voir
Champs, définition \ref{stacks-definition-descent-data}.
Considérons la règle $F$ qui, à $V/U$ dans $(\Sch/U)_h$, associe
l'ensemble des paires $(y, \psi_i)$, où $y$ est un objet de $\mathcal{S}_V$
et $\psi_i : y|_{U_i \times_U V} \to x_i|_{U_i \times_U V}$
est un morphisme de $\mathcal{S}$ au-dessus de $U_i \times_U V$
tel que
$$
\varphi_{ij}|_{U_i \times_U U_j \times_U V}
\circ
\psi_i|_{U_i \times_U U_j \times_U V}
=
\psi_j|_{U_i \times_U U_j \times_U V}
$$
à isomorphisme près. Puisque nous disposons déjà de la descente des morphismes, il est
clair que $F(V/U)$ est soit vide, soit un singleton. D'autre part,
$F(U_{i_0}/U)$ est non vide, car il contient
$(x_{i_0}, \varphi_{i_0i})$. Puisque notre but est de montrer que $F(U/U)$
est non vide, il suffit de montrer que $F$ est un faisceau sur $(\Sch/U)_h$.
Pour ce faire, nous pouvons utiliser le critère du lemme \ref{flat-lemma-refine-check-h}.
Cependant, nos hypothèses (1), (2), (3) impliquent (en traçant certains
diagrammes commutatifs que nous omettons) que les propriétés (1), (2), (3) et (4)
du lemme \ref{flat-lemma-refine-check-h} sont satisfaites pour $F$.

\medskip\noindent
Nous omettons de vérifier que, si $\mathcal{S}$ est un champ en groupoïdes,
alors (1), (2) et (3) sont satisfaites.
\end{proof}





\section{Normalisation faible absolue et recouvrements h}
\label{flat-section-h-and-weak-normalization}

\noindent
Dans cette section, nous utilisons les critères obtenus dans la
section \ref{flat-section-blow-up-h}
pour exhiber quelques faisceaux h et nous relions la faisceautisation h
du faisceau structural à la normalisation faible absolue.
Nous aurons besoin pour cela du lemme élémentaire suivant.

\begin{lemma}
\label{flat-lemma-funny-blow-up}
Soient $Z, X, X', E$ les schémas d'un carré de quasi-éclatement comme dans
l'exemple \ref{flat-example-two-generators}.
Alors $H^p(X', \mathcal{O}_{X'}) = 0$ pour $p > 0$ et
$\Gamma(X, \mathcal{O}_X) \to \Gamma(X', \mathcal{O}_{X'})$
est un homomorphisme surjectif d'anneaux dont le noyau est un idéal de carré nul.
\end{lemma}

\begin{proof}
Supposons d'abord que $A = \mathbf{Z}[f_1, f_2]$ est l'anneau de polynômes.
Dans ce cas, notre carré de quasi-éclatement est l'éclatement de $X = \Spec(A)$
le long du sous-schéma fermé $Z$ et, en fait, $X' \subset \mathbf{P}^1_X$
est un diviseur de Cartier effectif découpé par la section globale
$f_2T_0 - f_1 T_1$ de $\mathcal{O}_{\mathbf{P}^1_X}(1)$.
Nous avons donc une résolution
$$
0 \to
\mathcal{O}_{\mathbf{P}^1_X}(-1) \to
\mathcal{O}_{\mathbf{P}^1_X} \to
\mathcal{O}_{X'} \to
0
$$
En utilisant la description de la cohomologie donnée dans
Cohomologie des schémas, section
\ref{coherent-section-cohomology-projective-space},
il s'ensuit que, dans ce cas,
$\Gamma(X, \mathcal{O}_X) \to \Gamma(X', \mathcal{O}_{X'})$
est un isomorphisme et $H^1(X', \mathcal{O}_{X'}) = 0$.

\medskip\noindent
Ensuite, nous observons que tout diagramme comme dans l'exemple \ref{flat-example-two-generators}
est le changement de base du diagramme du paragraphe précédent par
l'homomorphisme d'anneaux $\mathbf{Z}[f_1, f_2] \to A$. Par conséquent, par Compléments sur les morphismes, lemmes
\ref{more-morphisms-lemma-check-h1-fibre-zero},
\ref{more-morphisms-lemma-h1-fibre-zero}, et
\ref{more-morphisms-lemma-h1-fibre-zero-check-h0-kappa},
nous concluons que $H^1(X', \mathcal{O}_{X'})$ est nul en général
et que l'homomorphisme
$H^0(X, \mathcal{O}_X) \to H^0(X', \mathcal{O}_{X'})$ est surjectif en général.

\medskip\noindent
Ensuite, dans le cas général, étudions le noyau. Si
$a \in A$ s'envoie sur zéro, alors, en examinant les cartes affines,
nous voyons que
$$
a = (f_1x - f_2)(a_0 + a_1x + \ldots + a_rx^r)\text{ dans }A[x]
$$
pour certains $r \geq 0$ et $a_0, \ldots, a_r \in A$, et de même
$$
a = (f_1 - f_2y)(b_0 + b_1y + \ldots + b_s y^s)\text{ dans }A[y]
$$
pour certains $s \geq 0$ et $b_0, \ldots, b_s \in A$. Cela signifie que nous
avons
$$
a = f_2 a_0,\ f_1 a_0 = f_2 a_1,\ \ldots,\ f_1 a_r = 0,
\ a = f_1 b_0,\ f_2 b_0 = f_1 b_1,\ \ldots,\ f_2 b_s = 0
$$
Si $(a', r', a'_i, s', b'_j)$ est un second système de ce type, alors nous avons
$$
aa' = f_1f_2a_0b'_0 = f_1f_2a_1b'_1 = f_1f_2a_2b'_2 = \ldots = 0
$$
comme souhaité.
\end{proof}

\noindent
Pour une $\mathbf{F}_p$-algèbre $A$,
nous posons $\colim_F A$ égal à la colimite du système
$$
A \xrightarrow{F} A \xrightarrow{F} A \xrightarrow{F} \ldots
$$
où $F : A \to A$, $a \mapsto a^p$ est l'endomorphisme de Frobenius.

\begin{lemma}
\label{flat-lemma-h-sheaf-colim-F}
Soit $p$ un nombre premier. Soit $S$ un schéma sur $\mathbf{F}_p$.
Soit $(\Sch/S)_h$ un site comme dans la définition \ref{flat-definition-big-small-h}.
Il existe un unique faisceau $\mathcal{F}$ sur $(\Sch/S)_h$ tel que
$$
\mathcal{F}(X) = \colim_F \Gamma(X, \mathcal{O}_X)
$$
pour tout objet quasi-compact et quasi-séparé $X$ de $(\Sch/S)_h$.
\end{lemma}

\begin{proof}
Notons $\mathcal{F}$ la faisceautisation de Zariski
du foncteur
$$
X \longrightarrow \colim_F \Gamma(X, \mathcal{O}_X)
$$
Pour les schémas quasi-compacts et quasi-séparés $X$,
nous avons $\mathcal{F}(X) = \colim_F \Gamma(X, \mathcal{O}_X)$
par Faisceaux, lemme \ref{sheaves-lemma-directed-colimits-sections}
et le fait que $\mathcal{O}$ est un faisceau pour la topologie de Zariski.
Il suffit donc de montrer que $\mathcal{F}$ est un faisceau h.
Pour le démontrer, nous vérifions les conditions (1), (2), (3) et (4) du
lemme \ref{flat-lemma-refine-check-h}.
La condition (1) est satisfaite parce que nous avons effectué une faisceautisation
de Zariski (presque superflue). La condition (2) est satisfaite parce que
$\mathcal{O}$ est un faisceau fppf (Descente, lemme
\ref{descent-lemma-sheaf-condition-holds}) et que, si
$A$ est l'égalisateur de deux homomorphismes $B \to C$ de $\mathbf{F}_p$-algèbres,
alors $\colim_F A$ est l'égalisateur des deux homomorphismes
$\colim_F B \to \colim_F C$.

\medskip\noindent
Nous vérifions la condition (3). Soient $A, f, J$ comme dans
l'exemple \ref{flat-example-one-generator}.
Nous devons montrer que
$$
\colim_F A = \colim_F A/J \times_{\colim_F A/fA + J} \colim_F A/fA
$$
Cela se ramène à la question d'algèbre suivante : supposons que $a', a'' \in A$
soient tels que $F^n(a' - a'') \in fA + J$. Trouver $a \in A$ et $m \geq 0$
tels que $a - F^m(a') \in J$ et $a - F^m(a'') \in fA$, et montrer que
la paire $(a, m)$ est uniquement déterminée à un remplacement près de la
forme $(a, m) \mapsto (F(a), m + 1)$.
Pour cela, il suffit d'écrire $F^n(a' - a'') = f h + g$ avec $h \in A$ et $g \in J$,
de poser $a = F^n(a') - g = F^n(a'') + fh$ et de poser $m = n$.
Pour voir l'unicité, supposons que $(a_1, m_1)$ soit une seconde solution.
Par un remplacement de la forme donnée ci-dessus, nous pouvons supposer $m = m_1$.
Alors nous voyons que $a - a_1 \in J$ et $a - a_1 \in fA$.
Puisque $J$ est annulé par une puissance de $f$, nous voyons que
$a - a_1$ est un élément nilpotent. Par conséquent, $F^k(a - a_1)$ est nul
pour un certain $k$ assez grand. Ainsi, après d'autres remplacements, nous obtenons
$a = a_1$.

\medskip\noindent
Nous vérifions la condition (4). Soient $X, X', Z, E$ comme dans
l'exemple \ref{flat-example-two-generators}. Par le
lemme \ref{flat-lemma-funny-blow-up}, nous voyons que
$$
\mathcal{F}(X) = \colim_F \Gamma(X, \mathcal{O}_X)
\longrightarrow
\colim_F \Gamma(X', \mathcal{O}_{X'}) = \mathcal{F}(X')
$$
est bijectif. Puisque $E = \mathbf{P}^1_Z$ dans ce cas, nous
voyons aussi que $\mathcal{F}(Z) \to \mathcal{F}(E)$ est bijectif.
La conclusion vaut donc également dans ce cas.
\end{proof}

\noindent
Soit $p$ un nombre premier. Pour une $\mathbf{F}_p$-algèbre $A$,
nous posons $\lim_F A$ égal à la limite du système inverse
$$
\ldots \xrightarrow{F} A \xrightarrow{F} A \xrightarrow{F} A
$$
où $F : A \to A$, $a \mapsto a^p$ est l'endomorphisme de Frobenius.

\begin{lemma}
\label{flat-lemma-h-sheaf-lim-F}
Soit $p$ un nombre premier. Soit $S$ un schéma sur $\mathbf{F}_p$.
Soit $(\Sch/S)_h$ un site comme dans la définition \ref{flat-definition-big-small-h}.
La règle
$$
\mathcal{F}(X) = \lim_F \Gamma(X, \mathcal{O}_X)
$$
définit un faisceau sur $(\Sch/S)_h$.
\end{lemma}

\begin{proof}
Pour montrer que $\mathcal{F}$ est un faisceau, vérifions les conditions
(1), (2), (3) et (4) du lemme \ref{flat-lemma-refine-check-h}.
La condition (1) est satisfaite parce que les limites de faisceaux sont des faisceaux
et que $\mathcal{O}$ est un faisceau de Zariski. La condition (2) est satisfaite parce que
$\mathcal{O}$ est un faisceau fppf (Descente, lemme
\ref{descent-lemma-sheaf-condition-holds}) et que, si $A$ est l'égalisateur
de deux homomorphismes $B \to C$ de $\mathbf{F}_p$-algèbres, alors $\lim_F A$
est l'égalisateur des deux homomorphismes $\lim_F B \to \lim_F C$.

\medskip\noindent
Nous vérifions la condition (3). Soient $A, f, J$ comme dans
l'exemple \ref{flat-example-one-generator}.
Nous devons montrer que
\begin{align*}
\lim_F A
& \to
\lim_F A/J \times_{\lim_F A/fA + J} \lim_F A/fA \\
& =
\lim_F (A/J \times_{A/fA + J} A/fA) \\
& =
\lim_F A/(fA \cap J)
\end{align*}
est bijectif. Puisque $J$ est annulé par une puissance de $f$, nous voyons que
$\mathfrak a = fA \cap J$ est un idéal nilpotent, c'est-à-dire qu'il existe un
$n$ tel que $\mathfrak a^n = 0$. Il est immédiat
de vérifier que, dans ce cas, $\lim_F A \to \lim_F A/\mathfrak a$
est bijectif.

\medskip\noindent
Nous vérifions la condition (4). Soient $X, X', Z, E$ comme dans
l'exemple \ref{flat-example-two-generators}. Par le
lemme \ref{flat-lemma-funny-blow-up} et le même argument que ci-dessus,
nous voyons que
$$
\mathcal{F}(X) = \lim_F \Gamma(X, \mathcal{O}_X)
\longrightarrow
\lim_F \Gamma(X', \mathcal{O}_{X'}) = \mathcal{F}(X')
$$
est bijectif. Puisque $E = \mathbf{P}^1_Z$ dans ce cas, nous
voyons aussi que $\mathcal{F}(Z) \to \mathcal{F}(E)$ est bijectif.
La conclusion vaut donc également dans ce cas.
\end{proof}

\noindent
Dans le lemme suivant, nous utilisons la normalisation faible absolue $X^{awn}$
d'un schéma $X$, voir
Morphismes, section \ref{morphisms-section-seminormalization}.

\begin{lemma}
\label{flat-lemma-weak-normalization-ph-sheaf}
Soit $(\Sch/S)_{ph}$ un site comme dans
Topologies, définition \ref{topologies-definition-big-small-ph}.
La règle
$$
X \longmapsto \Gamma(X^{awn}, \mathcal{O}_{X^{awn}})
$$
est un faisceau sur $(\Sch/S)_{ph}$.
\end{lemma}

\begin{proof}
Pour montrer que $\mathcal{F}$ est un faisceau, vérifions les conditions
(1) et (2) de Topologies, lemme \ref{topologies-lemma-characterize-sheaf}.
La condition (1) est satisfaite parce que la formation de $X^{awn}$ commute aux
recouvrements ouverts, voir Morphismes, lemme \ref{morphisms-lemma-seminormalization}
et sa démonstration.

\medskip\noindent
Soit $\pi : Y \to X$ un morphisme propre surjectif. Nous devons montrer
que l'égalisateur des deux homomorphismes
$$
\Gamma(Y^{awn}, \mathcal{O}_{Y^{awn}}) \to
\Gamma((Y \times_X Y)^{awn}, \mathcal{O}_{(Y \times_X Y)^{awn}})
$$
est égal à $\Gamma(X^{awn}, \mathcal{O}_{X^{awn}})$. Soit $f$ un
élément de cet égalisateur. Considérons alors le morphisme
$$
f : Y^{awn} \longrightarrow \mathbf{A}^1_X
$$
Puisque $Y^{awn} \to X$ est universellement fermé, l'image schématique
$Z$ de $f$ est un sous-schéma fermé de $\mathbf{A}^1_X$, propre sur $X$,
et $f : Y^{awn} \to Z$ est surjectif.
Voir Morphismes, lemme \ref{morphisms-lemma-scheme-theoretic-image-is-proper}.
Ainsi, $Z \to X$ est fini
(Morphismes, lemme \ref{morphisms-lemma-finite-proper})
et surjectif.

\medskip\noindent
Soit $k$ un corps et soient $z_1, z_2 : \Spec(k) \to Z$ deux
morphismes égalisés par $Z \to X$. Nous affirmons que $z_1 = z_2$.
Il suffit de montrer que les images $\lambda_i = z_i^*f \in k$ coïncident
(puisque le faisceau structural de $Z$ est engendré par $f$
sur le faisceau structural de $X$). Pour le voir, nous
choisissons une extension de corps
$K/k$ et des morphismes $y_1, y_2 : \Spec(K) \to Y^{awn}$ tels que
$z_i \circ (\Spec(K) \to \Spec(k)) = f \circ y_i$. Cela est possible
par la surjectivité du morphisme $Y^{awn} \to Z$. Choisissons une extension
algébriquement close $\Omega/k$ de très grand cardinal.
Pour tous les homomorphismes de $k$-algèbres $\sigma_i : K \to \Omega$,
nous obtenons
$$
\Spec(\Omega)
\xrightarrow{\sigma_1, \sigma_2}
\Spec(K \otimes_k K)
\xrightarrow{y_1, y_2}
Y^{awn} \times_X Y^{awn}
$$
Puisque le morphisme canonique
$(Y \times_X Y)^{awn} \to Y^{awn} \times_X Y^{awn}$
est un homéomorphisme universel et puisque $\Omega$ est algébriquement clos,
nous pouvons relever de manière unique la composition ci-dessus en un morphisme
$\Spec(\Omega) \to (Y \times_X Y)^{awn}$. Puisque $f$ appartient à l'égalisateur
ci-dessus, cela montre que $\sigma_1(\lambda_1) = \sigma_2(\lambda_2)$.
Un lemme facile sur les extensions de corps montre que cela implique
$\lambda_1 = \lambda_2$ ; les détails sont omis.

\medskip\noindent
Nous concluons que $Z \to X$ est universellement injectif, c'est-à-dire que
$Z \to X$ est injectif sur les points et induit des
extensions de corps résiduels purement inséparables
(Morphismes, lemme \ref{morphisms-lemma-universally-injective}).
En définitive, nous concluons que $Z \to X$ est un homéomorphisme universel, voir
Morphismes, lemme \ref{morphisms-lemma-universal-homeomorphism}.

\medskip\noindent
Soit $g : X^{awn} \to Z$ le morphisme obtenu par la propriété universelle
de $X^{awn}$. Alors $Y^{awn} \to X^{awn} \to Z$ et $f : Y^{awn} \to Z$
sont deux morphismes au-dessus de $X$. Par la propriété universelle de
$Y^{awn} \to Y$, les deux morphismes correspondants
$Y^{awn} \to Y \times_X Z$ au-dessus de $Y$ doivent être égaux. Cela implique
que $g \circ \pi^{wan} = f$ comme morphismes vers $\mathbf{A}^1_X$,
et nous concluons que $g \in \Gamma(X^{awn}, \mathcal{O}_{X^{awn}})$
est l'élément que nous cherchions.
\end{proof}

\begin{lemma}
\label{flat-lemma-weak-normalization-h-sheaf}
Soit $S$ un schéma. Choisissons un site $(\Sch/S)_h$
comme dans la définition \ref{flat-definition-big-small-h}.
La règle
$$
X \longmapsto \Gamma(X^{awn}, \mathcal{O}_{X^{awn}})
$$
est la faisceautisation du « faisceau structural » $\mathcal{O}$
sur $(\Sch/S)_h$. Il en va de même pour la topologie ph.
\end{lemma}

\begin{proof}
Dans le lemme \ref{flat-lemma-weak-normalization-ph-sheaf},
nous avons vu que la règle $\mathcal{F}$ du lemme
définit un faisceau pour la topologie ph, et donc, a fortiori,
un faisceau pour la topologie h. Il existe clairement un morphisme canonique
de préfaisceaux d'anneaux $\mathcal{O} \to \mathcal{F}$.
Pour terminer la démonstration, il suffit de montrer
\begin{enumerate}
\item si $f \in \mathcal{O}(X)$ s'envoie sur zéro dans $\mathcal{F}(X)$,
alors il existe un recouvrement h $\{X_i \to X\}$ tel que $f|_{X_i} = 0$, et
\item étant donné $f \in \mathcal{F}(X)$, il existe un recouvrement h
$\{X_i \to X\}$ tel que $f|_{X_i}$ soit l'image de $f_i \in \mathcal{O}(X_i)$.
\end{enumerate}
Soit $f$ comme en (1). Alors $f|_{X^{awn}} = 0$. Cela signifie que $f$
est localement nilpotent. Ainsi, si $X' \subset X$ est le sous-schéma fermé
découpé par $f$, alors $X' \to X$ est une immersion fermée surjective de
présentation finie. Par conséquent, $\{X' \to X\}$ est le recouvrement h recherché.
Soit $f$ comme en (2). Après avoir remplacé $X$ par les membres d'un
recouvrement ouvert affine, nous pouvons supposer que $X = \Spec(A)$ est affine.
Alors $f \in A^{awn}$, voir
Morphismes, lemme \ref{morphisms-lemma-seminormalization-ring}.
Par Morphismes, lemme \ref{morphisms-lemma-universal-homeo-limit},
nous pouvons trouver un homomorphisme d'anneaux $A \to B$ de présentation finie
tel que $\Spec(B) \to \Spec(A)$ soit un homéomorphisme universel
et que $f$ soit l'image d'un élément $b \in B$ par
l'homomorphisme canonique $B \to A^{awn}$. Alors
$\{\Spec(B) \to \Spec(A)\}$ est un recouvrement h, et nous concluons.
L'assertion concernant la topologie ph se démontre de la même manière
(ou peut se déduire de l'assertion pour la topologie h).
\end{proof}

\noindent
Soit $p$ un nombre premier. Une $\mathbf{F}_p$-algèbre $A$ est
dite {\it parfaite} si l'homomorphisme $F : A \to A$, $x \mapsto x^p$,
est un automorphisme de $A$.

\begin{lemma}
\label{flat-lemma-perfect-weankly-normal}
Soit $p$ un nombre premier. Une $\mathbf{F}_p$-algèbre $A$ est
absolument faiblement normale si et seulement si elle est parfaite.
\end{lemma}

\begin{proof}
Il résulte immédiatement de la condition (2)(b) dans
Morphismes, définition \ref{morphisms-definition-seminormal-ring}
que, si $A$ est absolument faiblement normale, alors elle est parfaite.

\medskip\noindent
Supposons $A$ parfaite. Soient $x, y \in A$ tels que $x^3 = y^2$.
Si $p > 3$, alors nous pouvons écrire $p = 2n + 3m$ pour certains $n, m > 0$.
Choisissons $a, b \in A$ tels que $a^p = x$ et $b^p = y$.
En posant $c = a^n b^m$, nous avons
$$
c^{2p} = x^{2n} y^{2m} = x^{2n + 3m} = x^p
$$
et donc $c^2 = x$. De même, $c^3 = y$. Si $p = 2$, alors
écrivons $x = a^2$ pour obtenir $a^6 = y^2$, ce qui implique $a^3 = y$.
Si $p = 3$, alors écrivons $y = a^3$ pour obtenir $x^3 = a^6$, ce qui
implique $x = a^2$.

\medskip\noindent
Supposons $x, y \in A$ avec $\ell^\ell x = y^\ell$ pour un certain nombre
premier $\ell$. Si $\ell \not = p$, alors $a = y/\ell$ vérifie
$a^\ell = x$ et $\ell a = y$. Si $\ell = p$, alors
$y = 0$ et $x = a^p$ pour un certain $a$.
\end{proof}

\begin{lemma}
\label{flat-lemma-char-p}
Soit $p$ un nombre premier.
\begin{enumerate}
\item Si $A$ est une $\mathbf{F}_p$-algèbre, alors $\colim_F A = A^{awn}$.
\item Si $S$ est un schéma sur $\mathbf{F}_p$, alors la
faisceautisation h de $\mathcal{O}$ associe à tout schéma quasi-compact
et quasi-séparé $X$ l'anneau $\colim_F \Gamma(X, \mathcal{O}_X)$.
\end{enumerate}
\end{lemma}

\begin{proof}
Démonstration de (1). Observons que $A \to \colim_F A$ induit un
homéomorphisme universel sur les spectres par
Algèbre, lemme \ref{algebra-lemma-p-ring-map}.
Il suffit donc de montrer que $B = \colim_F A$ est
absolument faiblement normal, voir
Morphismes, lemme \ref{morphisms-lemma-seminormalization-ring}.
Notons que l'homomorphisme d'anneaux $F : B \to B$ est un automorphisme ;
autrement dit, $B$ est un anneau parfait. Par conséquent, le
lemme \ref{flat-lemma-perfect-weankly-normal} s'applique.

\medskip\noindent
Démonstration de (2). Cela résulte de (1) et des
lemmes \ref{flat-lemma-h-sheaf-colim-F} et
\ref{flat-lemma-weak-normalization-h-sheaf}
en raisonnant localement sur des ouverts affines.
\end{proof}






\section{Descente des fibrés vectoriels en caractéristique positive}
\label{flat-section-descent-vectorbundles-h}

\noindent
Une référence pour cette section est \cite{Witt-Grass}.

\medskip\noindent
Pour un schéma $S$, notons $\textit{Vect}(S)$ la catégorie
des $\mathcal{O}_S$-modules localement libres de rang fini.
Soit $p$ un nombre premier. Soit $S$ un schéma quasi-compact et quasi-séparé
sur $\mathbf{F}_p$. Dans cette section, nous travaillerons avec la
catégorie
$$
\colim_F \textit{Vect}(S) =
\colim
\left(
\textit{Vect}(S) \xrightarrow{F^*}
\textit{Vect}(S) \xrightarrow{F^*}
\textit{Vect}(S) \xrightarrow{F^*} \ldots
\right)
$$
où $F : S \to S$ est le morphisme de Frobenius absolu. Concrètement,
un objet de cette catégorie est une paire $(\mathcal{E}, n)$,
où $\mathcal{E}$ est un $\mathcal{O}_S$-module localement libre de rang fini
et $n \geq 0$ est un entier. Pour les morphismes, nous prenons
$$
\Hom_{\colim_F \textit{Vect}(S)}((\mathcal{E}, n), (\mathcal{G}, m)) =
\colim_N \Hom_S(F^{N - n, *}\mathcal{E}, F^{N - m, *}\mathcal{G})
$$
où $F : S \to S$ est le morphisme de Frobenius absolu de $S$.
Ainsi, l'objet $(\mathcal{E}, n)$ est isomorphe à l'objet
$(F^*\mathcal{E}, n + 1)$.

\begin{lemma}
\label{flat-lemma-colim-F-Vect}
Soit $p$ un nombre premier. Soit $S$ un schéma quasi-compact et quasi-séparé
sur $\mathbf{F}_p$. La catégorie $\colim_F \textit{Vect}(S)$
est équivalente à la catégorie des modules localement libres de rang fini
sur le faisceau d'anneaux $\colim_F \mathcal{O}_S$ sur $S$.
\end{lemma}

\begin{proof}
Omis.
\end{proof}

\begin{lemma}
\label{flat-lemma-vector-bundle-I}
Soit $p$ un nombre premier. Considérons un carré de quasi-éclatement $X, X', Z, E$
en caractéristique $p$ comme dans l'exemple \ref{flat-example-one-generator}.
Alors le foncteur
$$
\colim_F \textit{Vect}(X)
\longrightarrow
\colim_F \textit{Vect}(Z)
\times_{\colim_F \textit{Vect}(E)}
\colim_F \textit{Vect}(X')
$$
est une équivalence.
\end{lemma}

\begin{proof}
Soient $A, f, J$ comme dans l'exemple \ref{flat-example-one-generator}.
Puisque tous nos schémas sont affines et que nous disposons du
Hom interne dans la catégorie des fibrés vectoriels, la pleine fidélité
du foncteur résulte du fait que
$$
\colim P \otimes_{A, F^N} A =
\colim P \otimes_{A, F^N} A/J
\times_{\colim P \otimes_{A, F^N} A/fA + J}
\colim P \otimes_{A, F^N} A/fA
$$
pour un $A$-module projectif de type fini $P$. En écrivant $P$ comme facteur direct
d'un module libre fini, cela résulte du cas où $P$ est libre
fini. Ce cas se ramène immédiatement au cas $P = A$. Le cas
$P = A$ résulte du lemme \ref{flat-lemma-h-sheaf-colim-F}
(en fait, nous avons démontré ce cas directement dans la démonstration de ce lemme).

\medskip\noindent
Surjectivité essentielle. Nous obtenons ici le problème d'algèbre suivant.
Supposons que $P_1$ soit un $A/J$-module projectif de type fini,
que $P_2$ soit un $A/fA$-module projectif de type fini, et que
$$
\varphi :
P_1 \otimes_{A/J} A/fA + J
\longrightarrow
P_2 \otimes_{A/fA} A/fA + J
$$
soit un isomorphisme. Objectif : montrer qu'il existe un $N$, un
$A$-module projectif de type fini $P$, un isomorphisme
$\varphi_1 : P \otimes_A A/J \to P_1 \otimes_{A/J, F^N} A/J$,
et un isomorphisme
$\varphi_2 : P \otimes_A A/fA \to P_2 \otimes_{A/fA, F^N} A/fA$
compatible avec $\varphi$ de manière évidente.
On peut le voir comme suit. Observons d'abord que
$$
A/(J \cap fA) = A/J \times_{A/fA + J} A/fA
$$
Par conséquent, par Compléments sur l'algèbre, lemme
\ref{more-algebra-lemma-finitely-presented-module-over-fibre-product},
il existe un module projectif de type fini $P'$ sur
$A/(J \cap fA)$ muni d'isomorphismes
$\varphi'_1 : P' \otimes_A A/J \to P_1$ et
$\varphi_2 : P' \otimes_A A/fA \to P_2$
compatibles avec $\varphi$. Puisque $J$ est un idéal de type fini et
de torsion pour une puissance de $f$, nous voyons que $J \cap fA$ est un
idéal nilpotent. Ainsi, pour un certain $N$, il existe une factorisation
$$
A \xrightarrow{\alpha} A/(J \cap fA) \xrightarrow{\beta} A
$$
de $F^N$. En posant $P = P' \otimes_{A/(J \cap fA), \beta} A$,
nous concluons.
\end{proof}

\begin{lemma}
\label{flat-lemma-vector-bundle-II}
Soit $p$ un nombre premier. Considérons un carré de quasi-éclatement $X, X', Z, E$
en caractéristique $p$ comme dans l'exemple \ref{flat-example-two-generators}.
Alors le foncteur
$$
G :
\colim_F \textit{Vect}(X)
\longrightarrow
\colim_F \textit{Vect}(Z)
\times_{\colim_F \textit{Vect}(E)}
\colim_F \textit{Vect}(X')
$$
est une équivalence.
\end{lemma}

\begin{proof}
Pleine fidélité. Supposons que $(\mathcal{E}, n)$ et
$(\mathcal{F}, m)$ soient des objets de $\colim_F \textit{Vect}(X)$.
Soit $(a, b) : G(\mathcal{E}, n) \to G(\mathcal{F}, m)$
un morphisme dans le membre de droite. Nous pouvons choisir $N \gg 0$ et
considérer $a$ comme un homomorphisme
$a : F^{N - n, *}\mathcal{E}|_Z \to F^{N - m, *}\mathcal{F}|_Z$
et $b$ comme un homomorphisme
$b : F^{N - n, *}\mathcal{E}|_{X'} \to F^{N - m, *}\mathcal{F}|_{X'}$
qui coïncident sur $E$.
Choisissons un recouvrement fini par des ouverts affines
$X = X_1 \cup \ldots \cup X_n$ tel que $\mathcal{E}|_{X_i}$
et $\mathcal{F}|_{X_i}$ soient des $\mathcal{O}_{X_i}$-modules libres de rang fini.
Pour chaque $i$, le changement de base
$$
\xymatrix{
E_i \ar[r] \ar[d] & X'_i \ar[d] \\
Z_i \ar[r] & X_i
}
$$
est un autre carré de quasi-éclatement comme dans l'exemple \ref{flat-example-two-generators}.
Pour ces carrés, nous savons que
$$
\colim_F H^0(X_i, \mathcal{O}_{X_i}) =
\colim_F H^0(Z_i, \mathcal{O}_{Z_i})
\times_{\colim_F H^0(E_i, \mathcal{O}_{E_i})}
\colim_F H^0(X'_i, \mathcal{O}_{X'_i})
$$
par le lemme \ref{flat-lemma-h-sheaf-colim-F} (voir la démonstration du lemme).
Ainsi, après avoir augmenté $N$, nous pouvons supposer que
les homomorphismes $a|_{Z_i}$ et $b|_{X'_i}$ proviennent
d'homomorphismes $c_i : F^{N - n, *}\mathcal{E}|_{X_i} \to F^{N - m, *}\mathcal{F}|_{X_i}$.
Après avoir éventuellement augmenté $N$, nous pouvons supposer que $c_i$ et $c_j$
coïncident sur $X_i \cap X_j$. Ces homomorphismes se recollent donc pour donner le
morphisme recherché $(\mathcal{E}, n) \to (\mathcal{F}, m)$
dans le membre de gauche.

\medskip\noindent
Surjectivité essentielle. Soit $(\mathcal{F}, \mathcal{G}, \varphi)$ un
triplet constitué de
un $\mathcal{O}_Z$-module localement libre de rang fini $\mathcal{F}$,
un $\mathcal{O}_{X'}$-module localement libre de rang fini $\mathcal{G}$, et
un isomorphisme $\varphi : \mathcal{F}|_E \to \mathcal{G}|_E$.
Nous devons montrer qu'après avoir remplacé ce triplet par un tiré en arrière par une
puissance du Frobenius, il provient d'un $\mathcal{O}_X$-module localement libre de rang fini.

\medskip\noindent
Réduction au cas noethérien ; nous conseillons vivement au lecteur de sauter ce paragraphe.
Rappelons que $X = \Spec(A)$ et $Z = \Spec(A/(f_1, f_2))$,
$X' = \text{Proj}(A[T_0, T_1]/(f_2T_0 - f_1T_1))$, et
$E = \mathbf{P}^1_Z$. Par Limites, lemme
\ref{limits-lemma-descend-invertible-modules},
nous pouvons trouver une $\mathbf{F}_p$-sous-algèbre de type fini
$A_0 \subset A$ contenant $f_1$ et $f_2$ telle que le triplet
$(\mathcal{F}, \mathcal{G}, \varphi)$ provienne de données définies sur
$X_0 = \Spec(A_0)$ et $Z_0 = \Spec(A_0/(f_1, f_2))$,
$X_0' = \text{Proj}(A_0[T_0, T_1]/(f_2T_0 - f_1T_1))$, et
$E_0 = \mathbf{P}^1_{Z_0}$. Nous pouvons donc supposer que nos schémas
sont noethériens.

\medskip\noindent
Supposons $X$ noethérien. Nous pouvons choisir un recouvrement fini par des ouverts affines
$X = X_1 \cup \ldots \cup X_n$ tel que $\mathcal{F}|_{Z \cap X_i}$ soit libre.
Puisque nous pouvons recoller les objets de $\colim_F \textit{Vect}(X)$
pour la topologie de Zariski (lemme \ref{flat-lemma-colim-F-Vect}), et
puisque nous connaissons déjà la
pleine fidélité sur $X_i$ et $X_i \cap X_j$ (voir le premier paragraphe
de la démonstration), il suffit de démontrer l'existence sur chaque $X_i$.
Cela nous ramène au cas traité dans le paragraphe suivant.

\medskip\noindent
Supposons $X$ noethérien et $\mathcal{F} = \mathcal{O}_Z^{\oplus r}$.
En utilisant $\varphi$, nous obtenons un isomorphisme
$\mathcal{O}_E^{\oplus r} \to \mathcal{G}|_E$.
Soit $I = (f_1, f_2) \subset A$.
Soit $\mathcal{I} \subset \mathcal{O}_{X'}$
le faisceau d'idéaux de $E$ ; il est globalement engendré par $f_1$ et $f_2$.
Pour tout $n$, il existe une surjection
$$
(\mathcal{I}^n/\mathcal{I}^{n + 1})^{\oplus r} =
\mathcal{I}^n/\mathcal{I}^{n + 1} \otimes_{\mathcal{O}_E}
\mathcal{G}|_E \longrightarrow
\mathcal{I}^n\mathcal{G}/\mathcal{I}^{n + 1}\mathcal{G}
$$
Le premier groupe de cohomologie de ce module est donc nul.
Nous utilisons ici le fait que $E = \mathbf{P}^1_Z$ et que, par conséquent, son faisceau
structural, et même tout module quasi-cohérent globalement engendré,
a un $H^1$ nul. Comparer avec Compléments sur les morphismes, lemme
\ref{more-morphisms-lemma-globally-generated-vanishing}.
Alors, en utilisant les suites exactes courtes
$$
0 \to  \mathcal{I}^n\mathcal{G}/\mathcal{I}^{n + 1}\mathcal{G} \to
\mathcal{G}/\mathcal{I}^{n + 1}\mathcal{G} \to
\mathcal{G}/\mathcal{I}^n\mathcal{G} \to 0
$$
et en raisonnant par récurrence, nous voyons que
$$
\lim H^0(X', \mathcal{G}/\mathcal{I}^n\mathcal{G})
\to
H^0(E, \mathcal{G}|_E) = H^0(E, \mathcal{O}_E^{\oplus r}) =
A/I^{\oplus r}
$$
est surjectif. Par le théorème des fonctions formelles
(Cohomologie des schémas, théorème \ref{coherent-theorem-formal-functions}),
cela implique que
$$
H^0(X', \mathcal{G}) \to
H^0(E, \mathcal{G}|_E) = H^0(E, \mathcal{O}_E^{\oplus r}) =
A/I^{\oplus r}
$$
est surjectif. Nous pouvons donc choisir un homomorphisme
$\alpha : \mathcal{O}_{X'}^{\oplus r} \to \mathcal{G}$
compatible avec la trivialisation donnée
de $\mathcal{G}|_E$. Ainsi, $\alpha$ est un isomorphisme sur
un voisinage ouvert de $E$ dans $X'$. Tout point
de $Z$ possède donc un voisinage ouvert affine sur lequel
nous pouvons résoudre le problème. Puisque $X' \setminus E \to X \setminus Z$
est un isomorphisme, il en va de même pour les points de $X$ qui ne sont pas dans $Z$.
Un autre argument de recollement de Zariski achève donc la démonstration.
\end{proof}

\begin{proposition}
\label{flat-proposition-h-descent-vector-bundles-p}
Soit $p$ un nombre premier. Soit $S$ un schéma de caractéristique $p$.
Alors la catégorie fibrée en groupoïdes
$$
p : \mathcal{S} \longrightarrow (\Sch/S)_h
$$
dont la catégorie fibre au-dessus de $U$ est la catégorie
des $\colim_F \mathcal{O}_U$-modules localement libres de rang fini sur $U$
est un champ en groupoïdes. De plus, si $U$ est quasi-compact
et quasi-séparé, alors $\mathcal{S}_U$ est $\colim_F \textit{Vect}(U)$.
\end{proposition}

\begin{proof}
La dernière assertion est le contenu du lemme \ref{flat-lemma-colim-F-Vect}.
Pour démontrer la proposition, nous vérifierons les conditions (1), (2) et (3) du
lemme \ref{flat-lemma-refine-check-h-stack}.

\medskip\noindent
La condition (1) est satisfaite parce que, par définition, nous avons le recollement
pour la topologie de Zariski.

\medskip\noindent
Pour voir la condition (2), supposons que $f : X \to Y$ soit un morphisme
plat, propre, surjectif et de présentation finie sur $S$, avec $Y$ affine.
Puisque $Y, X, X \times_Y X$ sont quasi-compacts et quasi-séparés,
nous pouvons utiliser la description des catégories fibres donnée dans
l'énoncé de la proposition. Il suffit alors clairement de
montrer que
$$
\textit{Vect}(Y)
\longrightarrow
\textit{Vect}(X)
\times_{\textit{Vect}(X \times_Y X)}
\textit{Vect}(X)
$$
est une équivalence (car cela impliquera la même chose pour les colimites).
Cela résulte immédiatement de la descente fppf des modules localement libres de rang fini, voir
Descente, proposition \ref{descent-proposition-fpqc-descent-quasi-coherent} et
lemme \ref{descent-lemma-finite-locally-free-descends}.

\medskip\noindent
La condition (3) est le contenu des
lemmes \ref{flat-lemma-vector-bundle-I} et \ref{flat-lemma-vector-bundle-II}.
\end{proof}

\begin{lemma}
\label{flat-lemma-trivial-fibres-dvr}
Soit $f : X \to S$ un morphisme propre à fibres géométriquement connexes,
où $S$ est le spectre d'un anneau de valuation discrète. Notons $\eta \in S$
le point générique et notons $X_n \subset X$ le sous-schéma fermé
découpé par la $n$-ième puissance d'une uniformisante sur $S$.
Il existe alors
un entier $n$ tel que l'assertion suivante soit vraie : tout
$\mathcal{O}_X$-module localement libre de rang fini $\mathcal{E}$
tel que $\mathcal{E}|_{X_\eta}$ et $\mathcal{E}|_{X_n}$
soient libres, est libre.
\end{lemma}

\begin{proof}
Nous nous ramenons d'abord au cas où $X \to S$ possède une section. Écrivons $S = \Spec(A)$.
Choisissons un point fermé $\xi$ de $X_\eta$. Choisissons une extension
d'anneaux de valuation discrète $A \subset B$ telle que le corps des fractions
de $B$ soit $\kappa(\xi)$. Cela est possible par Krull--Akizuki
(Algèbre, lemme \ref{algebra-lemma-integral-closure-Dedekind})
et par le fait que $\kappa(\xi)$ est une extension finie du
corps des fractions de $A$.
Par le critère valuatif de propreté
(Morphismes, lemme \ref{morphisms-lemma-characterize-proper}),
nous obtenons un point à valeurs dans $B$, $\tau : \Spec(B) \to X$,
qui induit une section $\sigma : \Spec(B) \to X_B$.
Pour un $\mathcal{O}_X$-module localement libre de rang fini $\mathcal{E}$,
notons $\mathcal{E}_B$ son tiré en arrière sur $X_B$.
Par changement de base plat
(Cohomologie des schémas, lemme \ref{coherent-lemma-flat-base-change-cohomology}),
nous voyons que $H^0(X_B, \mathcal{E}_B) = H^0(X, \mathcal{E}) \otimes_A B$.
Ainsi, si $\mathcal{E}_B$ est libre de rang $r$, alors les sections de
$H^0(X, \mathcal{E})$ engendrent le $B$-module libre
$\tau^*\mathcal{E} = \sigma^*\mathcal{E}_B$.
En particulier, nous pouvons trouver $r$ sections globales $s_1, \ldots, s_r$
de $\mathcal{E}$ qui engendrent $\tau^*\mathcal{E}$. Alors
$$
s_1, \ldots, s_r :
\mathcal{O}_X^{\oplus r}
\longrightarrow
\mathcal{E}
$$
est un homomorphisme de $\mathcal{O}_X$-modules localement libres de rang $r$,
et son tiré en arrière sur $X_B$ est un homomorphisme de $\mathcal{O}_{X_B}$-modules libres
dont la restriction est un isomorphisme en un point de chaque fibre.
En prenant le déterminant, nous obtenons une fonction
$g \in \Gamma(X_\eta, \mathcal{O}_{X_B})$
qui est inversible en un point de chaque fibre.
Comme les fibres sont propres et connexes, nous voyons que $g$
doit être inversible (détails omis ; indication : utiliser Variétés, lemme
\ref{varieties-lemma-proper-geometrically-reduced-global-sections}).
Il suffit donc de démontrer le lemme pour le changement de base $X_B \to \Spec(B)$.

\medskip\noindent
Supposons que nous ayons une section $\sigma : S \to X$. Soit $\mathcal{E}$
un $\mathcal{O}_X$-module localement libre de rang fini que nous supposons
libre sur la fibre générique et sur $X_n$ (nous choisirons $n$ plus tard).
Choisissons un isomorphisme $\sigma^*\mathcal{E} = \mathcal{O}_S^{\oplus r}$.
Considérons l'homomorphisme
$$
K = R\Gamma(X, \mathcal{E}) \longrightarrow
R\Gamma(S, \sigma^*\mathcal{E}) = A^{\oplus r}
$$
dans $D(A)$. En raisonnant comme ci-dessus, nous voyons que $\mathcal{E}$ est libre si (et seulement si)
l'homomorphisme induit $H^0(K) = H^0(X, \mathcal{E}) \to A^{\oplus r}$ est surjectif.

\medskip\noindent
Posons $L = R\Gamma(X, \mathcal{O}_X^{\oplus r})$ et observons que
l'homomorphisme correspondant $L \to A^{\oplus r}$ possède la propriété voulue.
Observons que $K \otimes_A Q(A) \cong L \otimes_A Q(A)$
par changement de base plat et par l'hypothèse que $\mathcal{E}$
est libre sur la fibre générique. Soit $\pi \in A$ une uniformisante. Observons que
$$
K \otimes_A^\mathbf{L} A/\pi^m A =
R\Gamma(X, \mathcal{E} \xrightarrow{\pi^m} \mathcal{E})
$$
et de même pour $L$.
Notons $\mathcal{E}_{tors} \subset \mathcal{E}$ le sous-faisceau cohérent des
sections supportées par la fibre spéciale, et procédons de même pour les autres
$\mathcal{O}_X$-modules. Choisissons $k > 0$ tel que
$(\mathcal{O}_X)_{tors} \to \mathcal{O}_X/\pi^k \mathcal{O}_X$
soit injectif (Cohomologie des schémas, lemme \ref{coherent-lemma-Artin-Rees}).
Puisque $\mathcal{E}$ est localement libre, nous voyons
que $\mathcal{E}_{tors} \subset \mathcal{E}/\pi^k\mathcal{E}$.
Alors, pour $n \geq m + k$, nous avons des isomorphismes
\begin{align*}
(\mathcal{E} \xrightarrow{\pi^m} \mathcal{E})
& \cong
(\mathcal{E}/\pi^k\mathcal{E} \xrightarrow{\pi^m}
\mathcal{E}/\pi^{k + m}\mathcal{E}) \\
& \cong
(\mathcal{O}_X^{\oplus r}/\pi^k\mathcal{O}_X^{\oplus r} \xrightarrow{\pi^m}
\mathcal{O}_X^{\oplus r}/\pi^{k + m}\mathcal{O}_X^{\oplus r}) \\
& \cong
(\mathcal{O}_X^{\oplus r} \xrightarrow{\pi^m} \mathcal{O}_X^{\oplus r})
\end{align*}
dans $D(\mathcal{O}_X)$. Cela détermine un isomorphisme
$$
K \otimes_A^\mathbf{L} A/\pi^m A \cong L \otimes_A^\mathbf{L} A/\pi^m A
$$
dans $D(A)$ (valable lorsque $n \geq m + k$). Observons que ces isomorphismes
sont compatibles au passage à l'image inverse par $\sigma$ ; en particulier,
nous concluons que
$K \otimes_A^\mathbf{L} A/\pi^m A \to (A/\pi^m A)^{\oplus r}$
induit une surjection sur les modules de cohomologie en degré $0$ (puisque
c'est vrai pour $L$).
Puisque $A$ est un anneau de valuation discrète, nous avons
$$
K \cong \bigoplus H^i(K)[-i]
\quad\text{et}\quad
L \cong
\bigoplus H^i(L)[-i]
$$
dans $D(A)$. Voir Compléments sur l'algèbre, exemple
\ref{more-algebra-example-finite-injective-finite-global-dimension}.
Les groupes de cohomologie $H^i(K) = H^i(X, \mathcal{E})$ et
$H^i(L) = H^i(X, \mathcal{O}_X)^{\oplus r}$
sont des $A$-modules de type fini d'après Cohomologie des schémas, lemme
\ref{coherent-lemma-proper-over-affine-cohomology-finite}.
Par Compléments sur l'algèbre, lemme
\ref{more-algebra-lemma-generalized-valuation-ring-modules},
ces modules sont des sommes directes de modules cycliques.
Nous avons vu ci-dessus que le rang $\beta_i$ de la
partie libre de $H^i(K)$ et de $H^i(L)$ est le même.
Observons ensuite que
$$
H^i(L \otimes_A^\mathbf{L} A/\pi^m A) =
H^i(L)/\pi^m H^i(L) \oplus H^{i + 1}(L)[\pi^m]
$$
et de même pour $K$. Soit $e$ le plus grand entier
tel que $A/\pi^eA$ apparaisse comme facteur direct de $H^i(X, \mathcal{O}_X)$,
ou, de manière équivalente, de $H^i(L)$, pour un certain $i$. En prenant alors $m = e + 1$,
nous voyons que $H^i(L \otimes_A^\mathbf{L} A/\pi^m A)$ est une somme directe de
$\beta_i$ copies de $A/\pi^m A$ et d'autres modules cycliques,
chacun annulé par $\pi^e$. Par le même raisonnement pour $K$
et l'isomorphisme
$K \otimes_A^\mathbf{L} A/\pi^m A \cong L \otimes_A^\mathbf{L} A/\pi^m A$,
il s'ensuit que $H^i(K)$
ne peut avoir aucun facteur direct cyclique de la forme $A/\pi^l A$
avec $l > e$. (Il s'ensuit également que $K$ est isomorphe à $L$
comme objet de $D(A)$, mais nous n'en aurons pas besoin.)
La seule manière dont l'homomorphisme
$$
H^0(K \otimes^\mathbf{L}_A A/\pi^{e + 1} A) =
H^0(K)/\pi^{e + 1}H^0(K) \oplus H^1(K)[\pi^{e + 1}]
\longrightarrow
(A/\pi^{e + 1} A)^{\oplus r}
$$
puisse être surjectif est donc qu'il soit surjectif sur le
premier facteur direct. C'est ce que nous voulions montrer.
(Pour être précis, l'entier $n$ de l'énoncé du
lemme, s'il existe une section $\sigma$,
doit être égal à $k + e + 1$, où $k$ et $e$ sont comme ci-dessus
et ne dépendent que de $X$.)
\end{proof}

\begin{lemma}
\label{flat-lemma-trivial-over-dvrs}
Soit $f : X \to S$ un morphisme de schémas.
Soit $\mathcal{E}$ un $\mathcal{O}_X$-module localement libre de rang fini.
Supposons que
\begin{enumerate}
\item $f$ est plat et propre et $\mathcal{O}_S = f_*\mathcal{O}_X$,
\item $S$ est un schéma noethérien normal,
\item le tiré en arrière de $\mathcal{E}$ sur $X \times_S \Spec(\mathcal{O}_{S, s})$
est libre pour tout point de codimension $1$, noté $s \in S$.
\end{enumerate}
Alors $\mathcal{E}$ est isomorphe au tiré en arrière d'un
$\mathcal{O}_S$-module localement libre de rang fini.
\end{lemma}

\begin{proof}
Nous allons démontrer que le morphisme canonique
$$
\Phi : f^*f_*\mathcal{E} \longrightarrow \mathcal{E}
$$
est un isomorphisme. Par changement de base plat (Cohomologie des schémas, lemme
\ref{coherent-lemma-flat-base-change-cohomology})
et, d'après les hypothèses (1) et (3), nous voyons que
son tiré en arrière sur $X \times_S \Spec(\mathcal{O}_{S, s})$
est un isomorphisme pour tout point de codimension $1$, noté $s \in S$.
Par Diviseurs, lemme \ref{divisors-lemma-check-isomorphism-via-depth-and-ass},
il suffit de démontrer que $\text{profondeur}((f^*f_*\mathcal{E})_x) \geq 2$
pour tout point $x \in X$ s'envoyant sur un point $s \in S$ de codimension $\geq 2$.
Puisque $f$ est plat et que
$(f^*f_*\mathcal{E})_x = (f_*\mathcal{E})_s \otimes_{\mathcal{O}_{S, s}}
\mathcal{O}_{X, x}$, il suffit de démontrer que
$\text{profondeur}((f_*\mathcal{E})_s) \geq 2$, voir
Algèbre, lemme \ref{algebra-lemma-apply-grothendieck}.
Puisque $S$ est un schéma noethérien normal
et que $\dim(\mathcal{O}_{S, s}) \geq 2$,
nous avons $\text{profondeur}(\mathcal{O}_{S, s}) \geq 2$, voir
Propriétés, lemme \ref{properties-lemma-criterion-normal}.
Nous obtenons donc le résultat voulu d'après
Diviseurs, lemme \ref{divisors-lemma-depth-pushforward}.
\end{proof}

\noindent
Nous pouvons utiliser les résultats ci-dessus pour démontrer l'énoncé
miraculeux suivant.

\begin{theorem}
\label{flat-theorem-pullback-trivial-fibres}
Soit $p$ un nombre premier. Soit $Y$ un schéma quasi-compact et quasi-séparé
sur $\mathbf{F}_p$.
Soit $f : X \to Y$ un morphisme propre, surjectif et de présentation finie
à fibres géométriquement connexes.
Alors le foncteur
$$
\colim_F \textit{Vect}(Y) \longrightarrow \colim_F \textit{Vect}(X)
$$
est pleinement fidèle et son image essentielle se décrit comme suit.
Soit $\mathcal{E}$ un $\mathcal{O}_X$-module localement libre de rang fini.
Supposons que, pour tout $y \in Y$, il existe des entiers $n_y, r_y \geq 0$
tels que
$$
F^{n_y, *}\mathcal{E}|_{X_{y, red}}
\cong
\mathcal{O}_{X_{y, red}}^{\oplus r_y}
$$
Alors, pour un certain $n \geq 0$, le tiré en arrière par la $n$-ième puissance de Frobenius
$F^{n, *}\mathcal{E}$ est le tiré en arrière d'un
$\mathcal{O}_Y$-module localement libre de rang fini.
\end{theorem}

\begin{proof}
Démonstration de la pleine fidélité. Puisque les fibrés vectoriels sur $Y$ sont localement
triviaux, cela se ramène à l'assertion que
$$
\colim_F \Gamma(Y, \mathcal{O}_Y)
\longrightarrow
\colim_F \Gamma(X, \mathcal{O}_X)
$$
est bijectif. Puisque $\{X \to Y\}$ est un recouvrement h, cela
résultera du lemme \ref{flat-lemma-h-sheaf-colim-F} si nous pouvons montrer que les deux homomorphismes
$$
\colim_F \Gamma(X, \mathcal{O}_X)
\longrightarrow
\colim_F \Gamma(X \times_Y X, \mathcal{O}_{X \times_Y X})
$$
sont égaux. Soit $g \in \Gamma(X, \mathcal{O}_X)$,
et notons $g_1$ et $g_2$ les deux tirés en arrière de $g$ à $X \times_Y X$.
Puisque $X_{y, red}$ est géométriquement connexe, nous
voyons que $H^0(X_{y, red}, \mathcal{O}_{X_{y, red}})$ est
une extension purement inséparable de $\kappa(y)$, voir
Variétés, lemme
\ref{varieties-lemma-proper-geometrically-reduced-global-sections}.
Ainsi, $g^q|_{X_{y, red}}$ provient d'un élément de
$\kappa(y)$ pour une certaine puissance de $p$, $q$ (qui peut dépendre de $y$).
Il s'ensuit que $g_1^q$ et $g_2^q$ s'envoient sur le même
élément du corps résiduel en tout point de
$(X \times_Y X)_y = X_y \times_y X_y$.
Par conséquent, la restriction de $g_1 - g_2$ à $(X \times_Y X)_{red}$ est nulle.
Ainsi, $(g_1 - g_2)^n = 0$ pour un certain $n$ que nous pouvons choisir
égal à une puissance de $p$, comme souhaité.

\medskip\noindent
Description de l'image essentielle. Soit $\mathcal{E}$ comme dans l'énoncé
de la proposition. Nous nous ramenons d'abord au cas noethérien.

\medskip\noindent
Soit $y \in Y$ un point, considéré comme un morphisme
$y \to Y$ du spectre du corps résiduel vers $Y$.
Nous pouvons écrire $y \to Y$ comme une limite filtrante de morphismes $Y_i \to Y$ de
présentation finie avec $Y_i$ affine. (Il est préférable de le démontrer
soi-même, mais cela résulte aussi formellement de
Limites, lemme \ref{limits-lemma-relative-approximation} et
\ref{limits-lemma-limit-affine}.)
Pour chaque $i$, posons $Z_i = Y_i \times_Y X$. Alors $X_y = \lim Z_i$
et $X_{y, red} = \lim Z_{i, red}$.
Par Limites, lemme \ref{limits-lemma-descend-modules-finite-presentation},
nous pouvons trouver un $i$ tel que
$F^{n_y, *}\mathcal{E}|_{Z_{i, red}} \cong
\mathcal{O}_{Z_{i, red}}^{\oplus r_y}$.
Fixons $i$.
Nous avons $Z_{i, red} = \lim Z_{i, j}$, où $Z_{i, j} \to Z_i$
est un épaississement de présentation finie (Limites, lemme
\ref{limits-lemma-closed-is-limit-closed-and-finite-presentation}).
En utilisant le même lemme que précédemment, nous pouvons trouver un $j$ tel que
$F^{n_y, *}\mathcal{E}|_{Z_{i, j}} \cong \mathcal{O}_{Z_{i, j}}^{\oplus r_y}$.
Nous concluons que, pour tout $y \in Y$, il existe un morphisme
$Y_y \to Y$ de présentation finie dont l'image contient $y$
et un épaississement $Z_y \to Y_y \times_Y X$ tel que
$F^{n_y, *}\mathcal{E}|_{Z_y} \cong \mathcal{O}_{Z_y}^{\oplus r_y}$.
Observons que l'image de $Y_y \to Y$ est constructible
(Morphismes, lemme \ref{morphisms-lemma-chevalley}).
Puisque $Y$ est quasi-compact pour la topologie constructible
(Topologie, lemme \ref{topology-lemma-constructible-hausdorff-quasi-compact} et
Propriétés, lemme \ref{properties-lemma-quasi-compact-quasi-separated-spectral}),
nous concluons qu'il existe un nombre fini de morphismes
$$
Y_1 \to Y,\ Y_2 \to Y,\ \ldots,\ Y_N \to Y
$$
de présentation finie tels que $Y = \bigcup \Im(Y_a \to Y)$
ensemblistement et tels que, pour chaque $a \in \{1, \ldots, N\}$,
il existe des entiers $n_a, r_a \geq 0$ et un épaississement
$Z_a \subset Y_a \times_Y X$ de présentation finie tel que
$F^{n_a, *}\mathcal{E}|_{Z_a} \cong \mathcal{O}_{Z_a}^{\oplus r_a}$.

\medskip\noindent
Sous cette forme, la condition se descend à une approximation
noethérienne absolue. Nous conseillons vivement au lecteur de sauter
ce paragraphe. Écrivons d'abord $Y = \lim_{i \in I} Y_i$ comme une limite cofiltrante
de schémas de type fini sur $\mathbf{F}_p$ à morphismes de transition
affines (Limites, lemme \ref{limits-lemma-relative-approximation}).
Ensuite, nous pouvons supposer que nous avons des morphismes propres $f_i : X_i \to Y_i$
dont le changement de base à $Y$ redonne $f : X \to Y$, voir
Limites, lemme \ref{limits-lemma-descend-finite-presentation}.
Après avoir augmenté $i$, nous pouvons supposer qu'il existe un
$\mathcal{O}_{X_i}$-module localement libre de rang fini $\mathcal{E}_i$ dont le tiré en arrière
sur $X$ est isomorphe à $\mathcal{E}$, voir
Limites, lemme \ref{limits-lemma-descend-invertible-modules}.
Choisissons $0 \in I$ et notons $E \subset Y_0$ le sous-ensemble constructible
où les fibres géométriques de $f_0$ sont connexes, voir
Compléments sur les morphismes, lemme
\ref{more-morphisms-lemma-nr-geom-connected-components-constructible}.
Alors $Y \to Y_0$ s'envoie dans $E$, voir
Compléments sur les morphismes, lemme
\ref{more-morphisms-lemma-base-change-fibres-geometrically-connected}.
Ainsi, $Y_i \to Y_0$ s'envoie dans $E$ pour $i \gg 0$, voir
Limites, lemme \ref{limits-lemma-limit-contained-in-constructible}.
Nous voyons donc que les fibres de $f_i$ sont géométriquement connexes
pour $i \gg 0$. Par Limites, lemme \ref{limits-lemma-descend-finite-presentation},
pour $i$ assez grand, nous pouvons trouver des morphismes
$Y_{i, a} \to Y_i$ de type fini dont le changement de base à $Y$
redonne $Y_a \to Y$, $a  \in \{1, \ldots, N\}$.
Après avoir éventuellement augmenté $i$, nous pouvons trouver des épaississements
$Z_{i, a} \subset Y_{i, a} \times_{Y_i} X_i$ dont le changement de base
à $Y_a \times_Y X$ redonne $Z_a$ (même référence que précédemment,
combinée avec
Limites, lemmes
\ref{limits-lemma-descend-closed-immersion-finite-presentation} et
\ref{limits-lemma-descend-surjective}).
Puisque $Z_a = \lim Z_{i, a}$, nous constatons qu'après avoir augmenté $i$, nous pouvons supposer
$F^{n_a, *}\mathcal{E}_i|_{Z_{i, a}} \cong
\mathcal{O}_{Z_{i, a}}^{\oplus r_a}$, voir
Limites, lemme \ref{limits-lemma-descend-modules-finite-presentation}.
Enfin, après avoir augmenté $i$ une dernière fois, nous pouvons supposer que
$\coprod Y_{i, a} \to Y_i$ est surjectif d'après
Limites, lemme \ref{limits-lemma-descend-surjective}.
À ce stade, toutes les hypothèses sont satisfaites pour $X_i \to Y_i$
et $\mathcal{E}_i$, et nous voyons qu'il suffit de démontrer le résultat
pour $X_i \to Y_i$ et $\mathcal{E}_i$.

\medskip\noindent
Supposons $Y$ de type fini sur $\mathbf{F}_p$.
Pour démontrer le résultat, nous allons raisonner par récurrence sur $\dim(Y)$.
Nous cherchons un objet de
$\colim_F \textit{Vect}(Y)$ dont le tiré en arrière soit
l'objet de $\colim_F \textit{Vect}(X)$ déterminé par $\mathcal{E}$.
Par la pleine fidélité déjà démontrée et en vertu de la
proposition \ref{flat-proposition-h-descent-vector-bundles-p},
il suffit de construire une descente de $\mathcal{E}$
après avoir remplacé $Y$ par les membres d'un recouvrement h
et $X$ par le changement de base correspondant. Cela signifie
que nous pouvons remplacer $Y$ par un schéma muni d'un morphisme propre et surjectif
vers $Y$, pourvu que cela n'augmente pas la dimension de $Y$.
Si $T \subset T'$ est un épaississement de schémas de type fini
sur $\mathbf{F}_p$, alors
$\colim_F \textit{Vect}(T) = \colim_F \textit{Vect}(T')$
car $\{T \to T'\}$ est un recouvrement h tel que $T \times_{T'} T = T$.
Si $T' \to T$ est un homéomorphisme universel de schémas
de type fini sur $\mathbf{F}_p$, alors
$\colim_F \textit{Vect}(T) = \colim_F \textit{Vect}(T')$
car $\{T \to T'\}$ est un recouvrement h tel que la diagonale
$T \subset T \times_{T'} T$ est un épaississement.

\medskip\noindent
D'après les remarques générales faites ci-dessus, nous remplaçons
$X$ par sa réduction et pouvons supposer que $X$ est réduit.
Considérons la factorisation de Stein $X \to Y' \to Y$, voir
Compléments sur les morphismes, théorème
\ref{more-morphisms-theorem-stein-factorization-Noetherian}.
Alors $Y' \to Y$ est un homéomorphisme universel
de schémas de type fini sur $\mathbf{F}_p$.
D'après ce qui précède, nous pouvons remplacer $Y$ par $Y'$.
Nous pouvons donc supposer $f_*\mathcal{O}_X = \mathcal{O}_Y$
et que $Y$ est réduit. Cela nous ramène au cas traité
dans le paragraphe suivant.

\medskip\noindent
Supposons $Y$ réduit et $f_*\mathcal{O}_X = \mathcal{O}_Y$
sur un sous-schéma ouvert dense de $Y$.
Alors $X \to Y$ est plat sur un sous-schéma
ouvert dense $V \subset Y$, voir
Morphismes, proposition \ref{morphisms-proposition-generic-flatness-reduced}.
Par le lemme \ref{flat-lemma-flat-after-blowing-up},
il existe un éclatement $V$-admissible $Y' \to Y$ tel que
le transformé strict $X'$ de $X$ soit plat sur $Y'$.
Observons que $\dim(Y') = \dim(Y)$, puisque $Y$ et $Y'$ possèdent
un sous-schéma ouvert dense commun. Par Compléments sur les morphismes, lemme
\ref{more-morphisms-lemma-proper-flat-nr-geom-conn-comps-lower-semicontinuous}
et le fait que $V \subset Y'$ est dense,
toutes les fibres de $f' : X' \to Y'$ sont géométriquement connexes.
Nous avons encore $(f'_*\mathcal{O}_{X'})|_V = \mathcal{O}_V$.
Écrivons
$$
Y' \times_Y X = X' \cup E \times_Y X
$$
où $E \subset Y'$ est le diviseur exceptionnel de l'éclatement.
Par les remarques générales ci-dessus, il suffit de démontrer l'existence
pour $Y' \times_Y X \to Y'$ et la restriction de $\mathcal{E}$
à $Y' \times_Y X$.
Supposons que nous trouvions un objet $\xi'$ dans $\colim_F \textit{Vect}(Y')$
dont le tiré en arrière soit la restriction de $\mathcal{E}$ à $X'$ (considérée
comme un objet de la catégorie colimite).
Par récurrence sur $\dim(Y)$, nous pouvons trouver un objet $\xi''$ dans
$\colim_F \textit{Vect}(E)$ dont le tiré en arrière soit la restriction de
$\mathcal{E}$ à $E \times_Y X$. Alors la pleine fidélité
détermine un unique isomorphisme $\xi'|_E \to \xi''$
compatible avec les identifications données avec la restriction
de $\mathcal{E}$ à $E \times_{Y'} X'$. Puisque
$$
\{E \times_Y X \to Y' \times_Y X, X' \to Y' \times_Y X\}
$$
est un recouvrement h donné par une paire d'immersions fermées avec
$$
(E \times_Y X) \times_{(Y' \times_Y X)} X' = E \times_{Y'} X'
$$
nous concluons que le tiré en arrière de $\xi'$ est la restriction de
$\mathcal{E}$ à $Y' \times_Y X$. Il suffit donc de trouver
$\xi'$, et nous nous ramenons au cas traité dans le paragraphe suivant.

\medskip\noindent
Supposons $Y$ réduit, $f$ plat et $f_*\mathcal{O}_X = \mathcal{O}_Y$
sur un sous-schéma ouvert dense de $Y$. Dans ce cas, nous considérons la
normalisation $Y^\nu \to Y$ (Morphismes, section
\ref{morphisms-section-normalization}). C'est un
morphisme fini surjectif
(Morphismes, lemme \ref{morphisms-lemma-nagata-normalization} et
\ref{morphisms-lemma-ubiquity-nagata}) qui est un isomorphisme
sur un ouvert dense. Par conséquent, d'après nos remarques générales, nous pouvons
remplacer $Y$ par $Y^\nu$ et $X$ par $Y^\nu \times_Y X$.
Après ce remplacement, nous voyons que $\mathcal{O}_Y = f_*\mathcal{O}_X$
(car la factorisation de Stein doit être un isomorphisme
dans ce cas ; petit détail omis).

\medskip\noindent
Supposons que $Y$ soit noethérien normal, que $f$ soit plat et que
$f_*\mathcal{O}_X = \mathcal{O}_Y$. Après avoir remplacé $\mathcal{E}$
par un tiré en arrière par une puissance convenable de Frobenius, nous pouvons supposer que $\mathcal{E}$
est trivial sur les fibres schématiques de $f$ aux points génériques
des composantes irréductibles de $Y$ (car
$\colim_F \textit{Vect}(-)$ est invariant par les homéomorphismes
universels, voir ci-dessus). Comme dans les arguments précédents
(lors de la réduction au cas noethérien), nous concluons qu'il existe un
sous-schéma ouvert dense $V \subset Y$ tel que $\mathcal{E}|_{f^{-1}(V)}$ soit libre.
Soit $Z \subset Y$ un sous-schéma fermé tel que
$Y = V \amalg Z$ ensemblistement. Soient $z_1, \ldots, z_t \in Z$
les points génériques des composantes irréductibles de $Z$
de codimension $1$. Alors $A_i = \mathcal{O}_{Y, z_i}$ est
un anneau de valuation discrète. Soit $n_i$ l'entier obtenu dans le
lemme \ref{flat-lemma-trivial-fibres-dvr} pour le schéma $X_{A_i}$ sur $A_i$.
Après avoir remplacé $\mathcal{E}$ par un tiré en arrière par une puissance
convenable de Frobenius, nous pouvons supposer que $\mathcal{E}$ est libre sur
$X_{A_i/\mathfrak m_i^{n_i}}$ (là encore parce que la catégorie colimite
est invariante par homéomorphismes universels, voir ci-dessus).
Alors le lemme \ref{flat-lemma-trivial-fibres-dvr} nous dit que
$\mathcal{E}$ est libre sur $X_{A_i}$.
Nous concluons enfin en appliquant le lemme \ref{flat-lemma-trivial-over-dvrs}.
\end{proof}










\section{Éclatement de complexes}
\label{flat-section-blowup-complexes}

\noindent
Cette section établit des formes normales pour les objets parfaits de la catégorie
dérivée après éclatement.

\begin{lemma}
\label{flat-lemma-fitting-ideals-complex}
Soit $X$ un schéma. Soit $E \in D(\mathcal{O}_X)$ pseudo-cohérent.
Pour tous $p, k \in \mathbf{Z}$, il existe un faisceau quasi-cohérent
d'idéaux de type fini $\text{Fit}_{p, k}(E) \subset \mathcal{O}_X$
ayant la propriété suivante : pour tout ouvert $U \subset X$
tel que $E|_U$ soit isomorphe à
$$
\ldots \to
\mathcal{O}_U^{\oplus n_{b - 2}}
\xrightarrow{d_{b - 2}}
\mathcal{O}_U^{\oplus n_{b - 1}}
\xrightarrow{d_{b - 1}}
\mathcal{O}_U^{\oplus n_b} \to 0 \to \ldots
$$
la restriction $\text{Fit}_{p, k}(E)|_U$ est engendrée par les
mineurs de la matrice de $d_p$ de taille
$$
- k + n_{p + 1} - n_{p + 2} + \ldots + (-1)^{b - p + 1} n_b
$$
Convention : l'idéal engendré par les mineurs $r \times r$
est $\mathcal{O}_U$ si $r \leq 0$, et l'idéal engendré par les
mineurs $r \times r$, lorsque $r > \min(n_p, n_{p + 1})$, est nul.
\end{lemma}

\begin{proof}
Observons que $E$ a, localement sur $X$, la forme indiquée dans le lemme ; voir
la section \ref{more-algebra-section-pseudo-coherent} de Compléments d'algèbre,
la section \ref{cohomology-section-pseudo-coherent} de Cohomologie et
la section \ref{perfect-section-spell-out} de Catégories dérivées de schémas.
Il suffit donc de démontrer que l'idéal des mineurs est indépendant
du représentant choisi. Pour cela, il suffit de raisonner
sur les anneaux locaux. Sur un anneau local $(R, \mathfrak m, \kappa)$,
considérons un complexe borné supérieurement
$$
F^\bullet :
\ldots \to
R^{\oplus n_{b - 2}}
\xrightarrow{d_{b - 2}}
R^{\oplus n_{b - 1}}
\xrightarrow{d_{b - 1}}
R^{\oplus n_b} \to 0 \to \ldots
$$
Notons $\text{Fit}_{k, p}(F^\bullet) \subset R$ l'idéal engendré
par les mineurs d'ordre $k - n_{p + 1} + n_{p + 2} - \ldots + (-1)^{b - p} n_b$
de la matrice de $d_p$. Supposons qu'un coefficient matriciel d'une
différentielle de $F^\bullet$ soit inversible. Choisissons alors le plus grand
entier $i$ tel que $d_i$ ait un coefficient matriciel inversible.
D'après le lemme \ref{algebra-lemma-add-trivial-complex} d'Algèbre,
le complexe $F^\bullet$ est isomorphe à la somme directe d'un complexe trivial
$\ldots \to 0 \to R \to R \to 0 \to \ldots$, dont les termes non nuls
sont en degrés $i$ et $i + 1$, et d'un complexe $(F')^\bullet$.
Nous laissons au lecteur le soin de vérifier que
$\text{Fit}_{p, k}(F^\bullet) = \text{Fit}_{p, k}((F')^\bullet)$ ;
c'est ici qu'intervient la formule donnant la taille des mineurs.
Si $(F')^\bullet$ possède une autre différentielle ayant un coefficient
matriciel inversible, nous recommençons, et ainsi de suite.
En poursuivant de cette manière, nous obtenons finalement un complexe $(F^\infty)^\bullet$
dont toutes les différentielles ont des matrices à coefficients
dans $\mathfrak m$. Il peut être nécessaire d'effectuer une infinité
d'étapes, mais, pour toute borne inférieure fixée, seul un nombre fini de ces
étapes modifie le complexe dans les degrés $\geq $ cette borne.
Ainsi, la « limite » $(F^\infty)^\bullet$ est un complexe borné supérieurement
bien défini de modules libres de rang fini ; elle est munie
d'un quasi-isomorphisme $(F^\infty)^\bullet \to F^\bullet$
vers le complexe de départ, et
$\text{Fit}_{p, k}(F^\bullet) = \text{Fit}_{p, k}((F^\infty)^\bullet)$.
Puisque le complexe $(F^\infty)^\bullet$ est unique à isomorphisme près d'après
le lemme
\ref{more-algebra-lemma-lift-pseudo-coherent-from-residue-field} de Compléments d'algèbre,
la démonstration est achevée.
\end{proof}

\begin{lemma}
\label{flat-lemma-blowup-complex}
Soit $X$ un schéma. Soit $E \in D(\mathcal{O}_X)$ parfait.
Soit $U \subset X$ un sous-schéma ouvert schématiquement dense
tel que $H^i(E|_U)$ soit localement libre de rang constant $r_i$
pour tout $i \in \mathbf{Z}$.
Alors il existe un éclatement $U$-admissible $b : X' \to X$ tel que
$H^i(Lb^*E)$ soit un $\mathcal{O}_{X'}$-module parfait
de dimension de Tor $\leq 1$ pour tout $i \in \mathbf{Z}$.
\end{lemma}

\begin{proof}
Nous allons construire et étudier l'éclatement localement sur des ouverts affines. Supposons donc que
$V \subset X$ soit un sous-schéma ouvert affine tel que
$E|_V$ puisse être représenté par le complexe
$$
\mathcal{O}_V^{\oplus n_a} \xrightarrow{d_a}
\ldots \xrightarrow{d_{b - 1}} \mathcal{O}_V^{\oplus n_b}
$$
Posons $k_ i = r_{i + 1} - r_{i + 2} + \ldots + (-1)^{b - i + 1}r_b$.
Un calcul, que nous omettons, montre que, sur $U \cap V$, le rang de
$d_i$ est
$$
\rho_i = - k_i + n_{i + 1} - n_{i + 2} + \ldots + (-1)^{b - i + 1}n_b
$$
au sens où le conoyau de $d_i$ est localement
libre de rang $n_{i + 1} - \rho_i$. Soit
$\mathcal{I}_i \subset \mathcal{O}_V$ l'idéal engendré par les mineurs
$\rho_i \times \rho_i$ de la matrice de $d_i$.

\medskip\noindent
D'une part, en comparant avec le lemme \ref{flat-lemma-fitting-ideals-complex},
nous voyons que l'idéal $\mathcal{I}_i$ correspond à l'idéal global
$\text{Fit}_{i, k_i}(E)$, dont on a démontré qu'il était
indépendant du choix du complexe représentant $E|_V$.
D'autre part, $\mathcal{I}_i$ est le $(n_{i + 1} - \rho_i)$-ième
idéal de Fitting de $\Coker(d_i)$. Gardons ce fait à l'esprit.

\medskip\noindent
Soit $b : X' \to X$ l'éclatement le long du produit des idéaux
$\text{Fit}_{i, k_i}(E)$ ; cette définition a un sens puisque, localement sur $X$,
presque tous ces idéaux sont égaux à l'idéal unité (voir ci-dessus).
Cet éclatement domine les éclatements $b_i : X'_i \to X$ le long des
idéaux $\text{Fit}_{i, k_i}(E)$ ; voir
le lemme \ref{divisors-lemma-blowing-up-two-ideals} de Diviseurs.
D'après le lemme \ref{divisors-lemma-blowup-fitting-ideal} de Diviseurs,
chaque $b_i$ est un éclatement $U$-admissible. Il s'ensuit que
$b$ est un éclatement $U$-admissible (nous omettons un petit détail ; comparer avec
la démonstration du lemme \ref{divisors-lemma-dominate-admissible-blowups} de Diviseurs).
Enfin, $U$ demeure un sous-schéma ouvert schématiquement dense de $X'$.
Après avoir remplacé $X$ par $X'$, nous nous trouvons donc dans la situation
étudiée au paragraphe suivant.

\medskip\noindent
Supposons que $\text{Fit}_{i, k_i}(E)$ soit un idéal inversible pour tout $i$.
Choisissons un ouvert affine $V$ et un complexe de modules libres de rang fini
représentant $E|_V$ comme ci-dessus. Il résulte du
lemme \ref{divisors-lemma-blowup-fitting-ideal} de Diviseurs,
que $\Coker(d_i)$ est de dimension de Tor $\leq 1$. Ainsi,
$\Im(d_i)$ est localement libre de rang fini, en tant que noyau d'une application
d'un module localement libre de rang fini vers un module de présentation finie
et de dimension de Tor $\leq 1$. Par conséquent, $\Ker(d_i)$
est lui aussi localement libre de rang fini (par le même argument). La suite
exacte courte
$$
0 \to \Im(d_{i - 1}) \to \Ker(d_i) \to H^i(E)|_V \to 0
$$
donne alors le résultat voulu, ce qui achève la démonstration.
\end{proof}

\begin{lemma}
\label{flat-lemma-blowup-complex-integral}
Soit $X$ un schéma intègre. Soit $E \in D(\mathcal{O}_X)$ parfait.
Alors il existe un ouvert non vide $U \subset X$
tel que $H^i(E|_U)$ soit localement libre de rang constant $r_i$
pour tout $i \in \mathbf{Z}$, ainsi qu'un éclatement $U$-admissible
$b : X' \to X$ tel que $H^i(Lb^*E)$ soit un
$\mathcal{O}_{X'}$-module parfait de dimension de Tor $\leq 1$ pour tout $i \in \mathbf{Z}$.
\end{lemma}

\begin{proof}
Nous invitons vivement le lecteur à trouver sa propre démonstration de
l'existence de $U$. Soit $\eta \in X$ le point générique.
La restriction de $E$ à $\eta$ est isomorphe dans $D(\kappa(\eta))$ à un
complexe fini $V^\bullet$ d'espaces vectoriels de dimension finie dont les
différentielles sont nulles. Posons $r_i = \dim_{\kappa(\eta)} V^i$.
Alors l'objet parfait $E'$ de $D(\mathcal{O}_X)$ représenté par le complexe
dont les termes sont $\mathcal{O}_X^{\oplus r_i}$ et les différentielles nulles
devient isomorphe à $E$ après restriction à $\eta$.
Par conséquent, d'après le lemme
\ref{perfect-lemma-descend-relatively-perfect} de Catégories dérivées de schémas,
il existe un voisinage ouvert $U$ de $\eta$ tel que $E|_U$ et $E'|_U$
soient isomorphes. Cela démontre la première assertion. La seconde résulte de la
première et du lemme \ref{flat-lemma-blowup-complex}, puisque tout ouvert non vide est
schématiquement dense dans le schéma intègre $X$.
\end{proof}

\begin{remark}
\label{flat-remark-when-you-have-a-complex}
Soit $X$ un schéma. Soit $E \in D(\mathcal{O}_X)$ un objet parfait tel que
$H^i(E)$ soit un $\mathcal{O}_X$-module parfait de dimension de Tor $\leq 1$
pour tout $i \in \mathbf{Z}$. Cette propriété permet parfois de ramener
des questions sur $E$ à des questions sur $H^i(E)$. Par exemple, supposons que
$$
\mathcal{E}^a \xrightarrow{d^a} \ldots
\xrightarrow{d^{b - 2}} \mathcal{E}^{b - 1}
\xrightarrow{d^{b - 1}} \mathcal{E}^b
$$
soit un complexe borné de $\mathcal{O}_X$-modules localement libres
de rang fini représentant $E$. Alors $\Im(d^i)$ et $\Ker(d^i)$ sont des
$\mathcal{O}_X$-modules localement libres de rang fini pour tout $i$. En effet, supposons par
récurrence que cela soit vrai pour tous les indices strictement supérieurs à $i$.
Nous pouvons alors commencer par utiliser la suite exacte courte
$$
0 \to \Im(d^i) \to \Ker(d^{i + 1}) \to H^{i + 1}(E) \to 0
$$
et l'hypothèse selon laquelle $H^{i + 1}(E)$ est parfait de dimension de Tor $\leq 1$
pour conclure que $\Im(d^i)$ est localement libre de rang fini.
Le même argument, appliqué ensuite à la suite exacte courte
$$
0 \to \Ker(d^i) \to \mathcal{E}^i \to \Im(d^i) \to 0
$$
montre que $\Ker(d^i)$ est localement libre de rang fini.
Il en résulte que les triangles distingués
$$
\tau_{\leq k - 1}E \to \tau_{\leq k}E \to H^k(E)[-k] \to
(\tau_{\leq k - 1}E)[1]
$$
sont représentés par les suites exactes courtes suivantes de
complexes bornés de modules localement libres de rang fini
$$
\begin{matrix}
& &
& &
& &
0 \\
& &
& &
& &
\downarrow \\
\mathcal{E}^a & \to &
\ldots & \to &
\mathcal{E}^{k - 2} & \to &
\Ker(d^{k - 1}) \\
\downarrow & &
& &
\downarrow & &
\downarrow \\
\mathcal{E}^a & \to &
\ldots & \to &
\mathcal{E}^{k - 2} & \to &
\mathcal{E}^{k - 1} & \to &
\Ker(d^k) \\
& &
& &
& &
\downarrow & &
\downarrow \\
& &
& &
& &
\Im(d^{k - 1}) & \to &
\Ker(d^k) \\
& &
& &
& &
\downarrow \\
& &
& &
& &
0
\end{matrix}
$$
Ici, les complexes sont les lignes, et les zéros « évidents »
sont omis dans le diagramme.
\end{remark}





\section{Éclatement de modules parfaits}
\label{flat-section-blowup-perfect-modules}

\noindent
Cette section tente d'établir des formes normales pour les modules parfaits de dimension de Tor
$\leq 1$ après éclatement. Nous n'y parvenons que partiellement.

\begin{lemma}
\label{flat-lemma-blowup-pd1-derived}
Soit $X$ un schéma. Soit $\mathcal{F}$ un
$\mathcal{O}_X$-module parfait de dimension de Tor $\leq 1$. Pour tout
éclatement $b : X' \to X$, on a $Lb^*\mathcal{F} = b^*\mathcal{F}$,
et $b^*\mathcal{F}$ est un $\mathcal{O}_X$-module parfait
de dimension de Tor $\leq 1$.
\end{lemma}

\begin{proof}
Nous pouvons supposer que $X = \Spec(A)$ est affine et que le $A$-module $M$
correspondant à $\mathcal{F}$ admet une présentation
$$
0 \to A^{\oplus m} \to A^{\oplus n} \to M \to 0
$$
Supposons que $I \subset A$ soit un idéal et que $a \in I$. Rappelons que
l'algèbre affine d'éclatement $A[\frac{I}{a}]$ est un sous-anneau de $A_a$.
Comme la localisation est exacte, $A_a^{\oplus m} \to A_a^{\oplus n}$
est injective. Par suite, $A[\frac{I}{a}]^{\oplus m} \to A[\frac{I}{a}]^{\oplus n}$
est elle aussi injective. Cela démontre le lemme.
\end{proof}

\begin{lemma}
\label{flat-lemma-blowup-pd1}
Soit $X$ un schéma. Soit $\mathcal{F}$ un $\mathcal{O}_X$-module parfait
de dimension de Tor $\leq 1$. Soit $U \subset X$ un ouvert schématiquement
dense tel que $\mathcal{F}|_U$ soit localement libre de rang constant
$r$. Alors il existe un éclatement $U$-admissible $b : X' \to X$ tel qu'il
existe une suite exacte courte canonique
$$
0 \to \mathcal{K} \to b^*\mathcal{F} \to \mathcal{Q} \to 0
$$
où $\mathcal{Q}$ est localement libre de rang $r$ et
$\mathcal{K}$ est un $\mathcal{O}_X$-module parfait
de dimension de Tor $\leq 1$ dont la restriction à $U$ est nulle.
\end{lemma}

\begin{proof}
Combiner le lemme \ref{divisors-lemma-blowup-fitting-ideal} de Diviseurs et
le lemme \ref{flat-lemma-blowup-pd1-derived}.
\end{proof}

\begin{lemma}
\label{flat-lemma-canonical-blowup-torsion-pd1}
Soit $X$ un schéma. Soit $\mathcal{F}$ un $\mathcal{O}_X$-module parfait
de dimension de Tor $\leq 1$. Soit $U \subset X$ un ouvert tel que
$\mathcal{F}|_U = 0$. Alors il existe un éclatement $U$-admissible
$$
b : X' \to X
$$
tel que $\mathcal{F}' = b^*\mathcal{F}$ soit muni de deux filtrations
canoniques localement finies
$$
0 = F^0 \subset F^1 \subset F^2 \subset \ldots \subset \mathcal{F}'
\quad\text{et}\quad
\mathcal{F}' = F_1 \supset F_2 \supset F_3 \supset \ldots \supset 0
$$
telles que, pour tout $n \geq 1$, il existe un diviseur de Cartier effectif
$D_n \subset X'$ tel que
$$
F^i/F^{i - 1}
\quad\text{et}\quad
F_i/F_{i + 1}
$$
sont localement libres de rang $i$ sur $D_i$.
\end{lemma}

\begin{proof}
Choisissons un ouvert affine $V \subset X$ tel qu'il existe une présentation
$$
0 \to \mathcal{O}_V^{\oplus n} \xrightarrow{A} \mathcal{O}_V^{\oplus n} \to
\mathcal{F} \to 0
$$
pour un certain $n$ et une certaine matrice $A$. L'idéal le long duquel nous allons éclater
est le produit des idéaux de Fitting $\text{Fit}_k(\mathcal{F})$
pour $k \geq 0$. Cette définition a un sens puisque, dans la situation affine
ci-dessus, $\text{Fit}_k(\mathcal{F})|_V = \mathcal{O}_V$
pour $k > n$. Il est clair qu'il s'agit d'un éclatement $U$-admissible. D'après
le lemme \ref{divisors-lemma-blowing-up-two-ideals} de Diviseurs,
nous voyons que, sur $X'$, les idéaux $\text{Fit}_k(\mathcal{F})$ sont inversibles.
Nous sommes donc ramenés au cas étudié au
paragraphe suivant.

\medskip\noindent
Supposons que $\text{Fit}_k(\mathcal{F})$ soit un idéal inversible pour
$k \geq 0$. Si $E_k \subset X$ est le diviseur de Cartier
effectif défini par $\text{Fit}_k(\mathcal{F})$
pour $k \geq 0$, alors les diviseurs de Cartier effectifs
$D_k$ de l'énoncé du lemme vérifieront
$$
E_k = D_{k + 1} + 2 D_{k + 2} + 3 D_{k + 3} + \ldots
$$
Cette somme a un sens puisque la famille $D_k$ sera localement finie.
De plus, elle détermine de manière unique les diviseurs de Cartier effectifs $D_k$ ;
il suffit donc de construire $D_k$ localement.

\medskip\noindent
Choisissons un ouvert affine $V \subset X$ et une présentation de $\mathcal{F}|_V$
comme ci-dessus. Nous allons construire les diviseurs et les filtrations par récurrence
sur l'entier $n$ de la présentation. Posons $D_k|_V = \emptyset$
pour $k > n$ et $D_n|V = E_{n - 1}|_V$.
Après avoir restreint $V$, nous pouvons supposer que
$\text{Fit}_{n - 1}(\mathcal{F})|_V$ est engendré par un
seul non-diviseur de zéro $f \in \Gamma(V, \mathcal{O}_V)$.
Comme $\text{Fit}_{n - 1}(\mathcal{F})|_V$ est l'idéal engendré
par les coefficients de $A$, il existe une matrice
$A'$ à coefficients dans $\Gamma(V, \mathcal{O}_V)$ telle que $A = fA'$.
Définissons $\mathcal{F}'$ sur $V$ par la suite exacte courte
$$
0 \to \mathcal{O}_V^{\oplus n} \xrightarrow{A'} \mathcal{O}_V^{\oplus n} \to
\mathcal{F}' \to 0
$$
Puisque les coefficients de $A'$ engendrent l'idéal unité de
$\Gamma(V, \mathcal{O}_V)$, le module $\mathcal{F}'$ admet localement sur $V$
une présentation où $n$ est diminué de
$1$ ; voir le lemme \ref{algebra-lemma-add-trivial-complex} d'Algèbre.
Notons en outre que
$f^{n - k}\text{Fit}_k(\mathcal{F}') = \text{Fit}_k(\mathcal{F})|_V$
pour $k = 0, \ldots, n$. Ainsi, $\text{Fit}_k(\mathcal{F}')$
est un idéal inversible pour tout $k$. Par récurrence,
il existe donc des diviseurs de Cartier effectifs $D'_k \subset V$
tels que $\mathcal{F}'$ admette deux filtrations canoniques comme dans l'énoncé
du lemme. Posons alors $D_k|_V = D'_k$ pour $k = 1, \ldots, n - 1$.
Les égalités entre diviseurs de Cartier effectifs
affichées ci-dessus sont vérifiées avec ces choix. Il reste
à construire les filtrations. En effet,
nous disposons des suites exactes courtes
$$
0 \to
\mathcal{O}_{D_n \cap V}^{\oplus n} \to
\mathcal{F} \to \mathcal{F}' \to 0
\quad\text{et}\quad
0 \to
\mathcal{F}' \to \mathcal{F} \to
\mathcal{O}_{D_n \cap V}^{\oplus n} \to 0
$$
qui proviennent des deux factorisations $A = A'f = f A'$ de $A$.
Ces suites sont canoniques parce que, dans la première,
le sous-module est $\Ker(f : \mathcal{F} \to \mathcal{F})$
et, dans la seconde, le module quotient est
$\Coker(f : \mathcal{F} \to \mathcal{F})$.
\end{proof}

\begin{lemma}
\label{flat-lemma-blowup-map-pd1}
Soit $X$ un schéma. Soit $\varphi : \mathcal{F} \to \mathcal{G}$
un homomorphisme de $\mathcal{O}_X$-modules parfaits de dimension de Tor $\leq 1$.
Soit $U \subset X$ un ouvert schématiquement dense
tel que $\mathcal{F}|_U = 0$ et $\mathcal{G}|_U = 0$.
Alors il existe un éclatement $U$-admissible $b : X' \to X$ tel que
le noyau, l'image et le conoyau de $b^*\varphi$ soient des
$\mathcal{O}_{X'}$-modules parfaits de dimension de Tor $\leq 1$.
\end{lemma}

\begin{proof}
Les hypothèses montrent que l'objet $(\mathcal{F} \to \mathcal{G})$
de $D(\mathcal{O}_X)$ est parfait. Nous obtenons donc un éclatement $U$-admissible
qui convient au conoyau et au noyau grâce aux lemmes \ref{flat-lemma-blowup-complex}
et \ref{flat-lemma-blowup-pd1-derived} (le second permettant de déterminer la forme du complexe
après image inverse). L'image est le noyau du conoyau ;
elle est donc elle aussi parfaite de dimension de Tor $\leq 1$.
\end{proof}











\section{Un opérateur introduit par Berthelot et Ogus}
\label{flat-section-eta}

\noindent
Lire d'abord la section \ref{cohomology-section-eta} de Cohomologie.

\medskip\noindent
Soit $X$ un schéma. Soit $D \subset X$ un diviseur de Cartier effectif.
Soit $\mathcal{I} = \mathcal{I}_D \subset \mathcal{O}_X$
le faisceau d'idéaux de $D$ ; voir
la section
\ref{divisors-section-effective-Cartier-invertible} de Diviseurs.
On peut manifestement appliquer les considérations de la section
\ref{cohomology-section-eta} de Cohomologie à $X$ et $\mathcal{I}$.

\begin{lemma}
\label{flat-lemma-eta-stalks}
Soit $X$ un schéma. Soit $D \subset X$ un diviseur de Cartier effectif
de faisceau d'idéaux $\mathcal{I} \subset \mathcal{O}_X$.
Soit $\mathcal{F}^\bullet$ un complexe de
$\mathcal{O}_X$-modules quasi-cohérents tel que $\mathcal{F}^i$ soit
sans $\mathcal{I}$-torsion pour tout $i$. Alors
$\eta_\mathcal{I}\mathcal{F}^\bullet$ est un complexe de
$\mathcal{O}_X$-modules quasi-cohérents. De plus,
si $U = \Spec(A) \subset X$ est un ouvert affine et $D \cap U = V(f)$,
alors $\eta_f(\mathcal{F}^\bullet(U))$ est canoniquement isomorphe
à $(\eta_\mathcal{I}\mathcal{F}^\bullet)(U)$.
\end{lemma}

\begin{proof}
Omis.
\end{proof}

\begin{lemma}
\label{flat-lemma-Leta}
Soit $X$ un schéma. Soit $D \subset X$ un diviseur de Cartier effectif
de faisceau d'idéaux $\mathcal{I} \subset \mathcal{O}_X$.
Le foncteur $L\eta_\mathcal{I} : D(\mathcal{O}_X) \to D(\mathcal{O}_X)$ du
lemme \ref{cohomology-lemma-Leta} de Cohomologie envoie
$D_\QCoh(\mathcal{O}_X)$ dans lui-même.
De plus, si $X = \Spec(A)$ est affine et $D = V(f)$,
alors le foncteur $L\eta_f$ sur $D(A)$ défini dans le
lemme \ref{more-algebra-lemma-Leta} de Compléments d'algèbre
et le foncteur $L\eta_\mathcal{I}$ sur $D_\QCoh(\mathcal{O}_X)$
se correspondent par l'équivalence du lemme
\ref{perfect-lemma-affine-compare-bounded} de Catégories dérivées de schémas.
\end{lemma}

\begin{proof}
Omis.
\end{proof}





\section{Éclatement de complexes, II}
\label{flat-section-blowup-complexes-II}

\noindent
Le contenu de cette section sera utilisé pour construire une version de la
construction du graphe de Macpherson dans la section \ref{flat-section-blowup-complexes-III}.

\begin{situation}
\label{flat-situation-complex-and-divisor}
Ici $X$ est un schéma, $D \subset X$ est un diviseur de Cartier effectif
de faisceau d'idéaux $\mathcal{I} \subset \mathcal{O}_X$, et
$M$ est un objet parfait de $D(\mathcal{O}_X)$.
\end{situation}

\noindent
Soit $(X, D, M)$ un triplet comme dans la
situation \ref{flat-situation-complex-and-divisor}.
Considérons un ouvert affine $U = \Spec(A) \subset X$
tel que
\begin{enumerate}
\item $D \cap U = V(f)$ pour un non-diviseur de zéro $f \in A$, et
\item il existe un complexe borné $M^\bullet$ de
$A$-modules libres de rang fini représentant $M|_U$ (par l'équivalence de
Catégories dérivées des schémas, lemme
\ref{perfect-lemma-affine-compare-bounded}).
\end{enumerate}
Nous dirons que $(U, A, f, M^\bullet)$ est une
{\it carte affine} pour $(X, D, M)$.
Considérons les idéaux $I_i(M^\bullet, f) \subset A$ définis dans
Compléments d'algèbre, section \ref{more-algebra-section-perfect-eta}.
Disons que $(X, S, M)$ est un {\it bon triplet} si, pour tout $x \in D$,
il existe une carte affine $(U, A, f, M^\bullet)$
avec $x \in U$ et telle que les idéaux $I_i(M^\bullet, f)$ soient principaux pour tout
$i \in \mathbf{Z}$.

\begin{lemma}
\label{flat-lemma-pullback-triple-ideals-good}
Dans la situation \ref{flat-situation-complex-and-divisor}, soit $h : Y \to X$ un
morphisme de schémas tel que l'image inverse $E = h^{-1}D$ de $D$
soit définie (Diviseurs, définition
\ref{divisors-definition-pullback-effective-Cartier-divisor}).
Soit $(U, A, f, M^\bullet)$ une carte affine pour $(X, D, M)$.
Soit $V = \Spec(B) \subset Y$ un ouvert affine avec $h(V) \subset U$.
Notons $g \in B$ l'image de $f \in A$.
Alors
\begin{enumerate}
\item $(V, B, g, M^\bullet \otimes_A B)$ est une carte affine pour $(Y, E, Lh^*M)$,
\item $I_i(M^\bullet, f)B = I_i(M^\bullet \otimes_A B, g)$ dans $B$, et
\item si $(X, D, M)$ est un bon triplet, alors
$(Y, E, Lh^*M)$ est un bon triplet.
\end{enumerate}
\end{lemma}

\begin{proof}
La première assertion résulte des observations suivantes :
$g$ est un non-diviseur de zéro dans $B$ qui définit $E \cap V \subset V$
et $M^\bullet \otimes_A B$ représente $M^\bullet \otimes_A^\mathbf{L} B$
et représente donc l'image inverse de $M$ sur $V$ d'après
Catégories dérivées des schémas, lemme
\ref{perfect-lemma-quasi-coherence-pullback}.
La partie (2) résulte de la partie (1) et de
Compléments d'algèbre, lemme \ref{more-algebra-lemma-eta-base-change-pre}.
En combinant ceci avec
Compléments d'algèbre, lemme \ref{more-algebra-lemma-eta-base-change-pre},
nous concluons que la deuxième assertion du lemme est vraie.
\end{proof}

\begin{lemma}
\label{flat-lemma-good-triple}
Soient $X, D, \mathcal{I}, M$ comme dans la
situation \ref{flat-situation-complex-and-divisor}.
Si $(X, D, M)$ est un bon triplet, alors
$L\eta_\mathcal{I}M$ est un objet parfait
de $D(\mathcal{O}_X)$.
\end{lemma}

\begin{proof}
Traduction de
Compléments d'algèbre, lemme \ref{more-algebra-lemma-eta-locally-free}.
Pour effectuer la traduction, utiliser le lemme \ref{flat-lemma-Leta}.
\end{proof}

\begin{lemma}
\label{flat-lemma-good-triple-bdd-loc-free}
Soient $X, D, \mathcal{I}, M$ comme dans la
situation \ref{flat-situation-complex-and-divisor}.
Supposons que $(X, D, M)$ soit un bon triplet.
S'il existe un complexe localement borné $\mathcal{M}^\bullet$
de $\mathcal{O}_X$-modules localement libres de rang fini représentant $M$,
alors il existe un complexe localement borné $\mathcal{Q}^\bullet$
de $\mathcal{O}_{X'}$-modules localement libres de rang fini représentant
$L\eta_\mathcal{I}M$.
\end{lemma}

\begin{proof}
D'après Cohomologie, lemme \ref{cohomology-lemma-Leta},
le complexe $\mathcal{Q}^\bullet = \eta_\mathcal{I}\mathcal{M}^\bullet$
représente $L\eta_\mathcal{I}M$. Pour vérifier que ce complexe
est localement borné et formé de modules localement libres de rang fini, nous pouvons travailler
localement sur des ouverts affines. La bornitude est alors claire.
Choisissons une carte affine $(U, A, f, M^\bullet)$ pour $(X, D, M)$
telle que les idéaux $I_i(M^\bullet, f)$ soient principaux et telle que
$\mathcal{M}^i|_U$ soit libre de rang fini pour chaque $i$. Notre hypothèse
que $(X, D, M)$ est un bon triplet nous permet de le faire. En posant
$N^i = \Gamma(U, \mathcal{M}^i|_U)$, nous obtenons un complexe borné
$N^\bullet$ de $A$-modules libres de rang fini représentant le
même objet de $D(A)$ que le complexe $M^\bullet$ (d'après
Catégories dérivées des schémas, lemme
\ref{perfect-lemma-affine-compare-bounded}).
Alors $I_i(N^\bullet, f)$ est un idéal principal pour tout $i$
d'après Compléments d'algèbre, lemme \ref{more-algebra-lemma-ideal-well-defined}.
Par conséquent, le complexe $\eta_fN^\bullet$ est un complexe borné
de $A$-modules localement libres de rang fini. Comme
$\mathcal{Q}^i|_U$ est le $\mathcal{O}_U$-module quasi-cohérent
correspondant à $\eta_fN^i$ d'après le lemme \ref{flat-lemma-eta-stalks}, nous concluons.
\end{proof}

\begin{lemma}
\label{flat-lemma-complex-and-divisor-eta-pull}
Dans la situation \ref{flat-situation-complex-and-divisor}, soit $h : Y \to X$
un morphisme de schémas tel que l'image inverse $E = h^{-1}D$
soit définie. Si $(X, D, M)$ est un bon triplet, alors
$$
Lh^*(L\eta_\mathcal{I}M) = L\eta_\mathcal{J}(Lh^*M)
$$
dans $D(\mathcal{O}_Y)$, où $\mathcal{J}$ est le faisceau d'idéaux de $E$.
\end{lemma}

\begin{proof}
Traduction de
Compléments d'algèbre, lemme \ref{more-algebra-lemma-eta-base-change}.
Utiliser les lemmes \ref{flat-lemma-eta-stalks} et \ref{flat-lemma-Leta}
pour effectuer la traduction.
\end{proof}

\begin{lemma}
\label{flat-lemma-complex-and-divisor-blowup-pre}
Dans la situation \ref{flat-situation-complex-and-divisor},
il existe un unique morphisme $b : X' \to X$ tel que
\begin{enumerate}
\item l'image inverse $D' = b^{-1}D$ soit définie et que
$(X', D', M')$ soit un bon triplet, où $M' = Lb^*M$, et
\item pour tout morphisme de schémas $h : Y \to X$ tel que
l'image inverse $E = h^{-1}D$ soit définie et que $(Y, E, Lh^*M)$
soit un bon triplet, $h$ se factorise de manière unique par $b$.
\end{enumerate}
De plus, pour toute carte affine $(U, A, f, M^\bullet)$, la restriction
$b^{-1}(U) \to U$ est l'éclatement dans le produit des idéaux
$I_i(M^\bullet, f)$ et, pour tout ouvert quasi-compact $W \subset X$, la
restriction $b|_{b^{-1}(W)} : b^{-1}(W) \to W$ est, relativement à $W \setminus D$, un
éclatement admissible.
\end{lemma}

\begin{proof}
La preuve consiste simplement à éclater localement $X$ dans les idéaux produits
$I_i(M^\bullet, f)$ pour toute carte affine $(U, A, f, M^\bullet)$.
Les premiers lemmes de Compléments d'algèbre, section
\ref{more-algebra-section-perfect-eta}, montrent que ceci est bien défini.
La propriété universelle (2) résulte alors de la
propriété universelle de l'éclatement.
On trouvera les détails ci-dessous.

\medskip\noindent
Soit $U, A, f, M^\bullet$ une carte affine pour $(X, D, M)$.
Tous les idéaux $I_i(M^\bullet, f)$ sauf un nombre fini sont
égaux à $A$, il est donc légitime de considérer
$$
I = \prod\nolimits_i I_i(M^\bullet, f)
$$
et c'est un idéal de type fini de $A$. Notons
$$
b_U : U' \to U
$$
l'éclatement de $U$ dans $I$. Alors $b_U^{-1}(U \cap D)$ est défini d'après
Diviseurs, lemme
\ref{divisors-lemma-blow-up-pullback-effective-Cartier}.
Rappelons que $f^{r_i} \in I_i(M^\bullet, f)$ et que, par conséquent,
$b_U$ est un éclatement $(U \setminus D)$-admissible.
D'après Diviseurs, lemme \ref{divisors-lemma-blowing-up-two-ideals},
pour chaque $i$, le morphisme $b_U$ se factorise en $U' \to U'_i \to U$,
où $U'_i \to U$ est l'éclatement dans $I_i(M^\bullet, f)$
et $U' \to U'_i$ est un autre éclatement. Il s'ensuit que
l'image inverse $I_i(M^\bullet, f)\mathcal{O}_{U'}$
de $I_i(M^\bullet, f)$ sur $U'$ est un faisceau d'idéaux inversible, voir
Diviseurs, lemmes \ref{divisors-lemma-blow-up-pullback-effective-Cartier} et
\ref{divisors-lemma-blowing-up-gives-effective-Cartier-divisor}.
Il s'ensuit que $(U', b^{-1}D, Lb^*M|_U)$ est un bon triplet, voir
le lemme \ref{flat-lemma-pullback-triple-ideals-good}
pour le comportement des idéaux $I_i(-,-)$ par image inverse.
Enfin, nous affirmons que $b_U : U' \to U$ possède la propriété universelle
mentionnée dans la partie (2) de l'énoncé du lemme. En effet, supposons que
$h : Y \to U$ soit un morphisme de schémas tel que
l'image inverse $E = h^{-1}(D \cap U)$ soit définie et que $(Y, E, Lh^*M)$
soit un bon triplet. Alors $Y$ est recouvert par des cartes affines
$(V, B, g, N^\bullet)$ telles que $I_i(N^\bullet, g)$ soit
un idéal inversible pour chaque $i$. Alors $g$ et l'image de
$f$ dans $B$ diffèrent par une unité, puisqu'ils découpent tous deux le diviseur de Cartier
effectif $E \cap V$. Nous pouvons donc supposer que $g$ est l'image de $f$ d'après
Compléments d'algèbre, lemme \ref{more-algebra-lemma-eta-change-unit}.
Alors $I_i(N^\bullet, g)$ est isomorphe à $I_i(M^\bullet \otimes_A B, g)$
comme $B$-module d'après Compléments d'algèbre, lemme
\ref{more-algebra-lemma-ideal-well-defined}.
Ainsi $I_i(M^\bullet \otimes_A B, g) = I_i(M^\bullet, f)B$
(lemme \ref{flat-lemma-pullback-triple-ideals-good})
est un $B$-module inversible.
Par conséquent, l'idéal $IB$ est inversible. Il s'ensuit que $I\mathcal{O}_Y$
est inversible. Nous obtenons donc que $h$ se factorise de manière unique par
$b_U$, d'après Diviseurs, lemme \ref{divisors-lemma-universal-property-blowing-up}.

\medskip\noindent
Soit $\mathcal{B}$ l'ensemble des ouverts affines $U \subset X$
pour lesquels il existe une carte affine $(U, A, f, M^\bullet)$
pour $(X, D, M)$. Alors $\mathcal{B}$ est une base de la topologie
de $X$ ; détails omis. Pour $U \in \mathcal{B}$, nous avons
le morphisme $b_U : U' \to U$ construit ci-dessus, qui satisfait
la propriété universelle sur $U$. Si $U_1 \subset U_2 \subset X$
appartiennent tous deux à $\mathcal{B}$, alors $b_{U_1} : U'_1 \to U_1$
est canoniquement isomorphe à
$$
b_{U_2}|_{b_{U_2}^{-1}(U_1)} : b_{U_2}^{-1}(U_1) \longrightarrow U_1
$$
par la propriété universelle. Autrement dit, nous obtenons un isomorphisme
$U'_1 \to b_{U_2}^{-1}(U_1)$ sur $U_1$. Ces isomorphismes satisfont
la condition de cocycle (de nouveau par la propriété universelle) et donc, d'après
Constructions, lemme \ref{constructions-lemma-relative-glueing},
nous obtenons un morphisme $b : X' \to X$ dont la restriction à chaque
$U$ de $\mathcal{B}$ est isomorphe à $U' \to U$.
Alors le morphisme $b : X' \to X$ satisfait les propriétés (1) et (2)
de l'énoncé du lemme, car ces propriétés peuvent être vérifiées
localement (détails omis).

\medskip\noindent
Il nous reste à démontrer la dernière assertion du lemme.
Soit $W \subset X$ un ouvert quasi-compact.
Choisissons un recouvrement fini $W = U_1 \cup \ldots \cup U_T$
tel que, pour chaque $1 \leq t \leq T$, il existe une carte affine
$(U_t, A_t, f_t, M_t^\bullet)$. Nous utiliserons ci-dessous le fait que, pour tout
ouvert affine $V = \Spec(B) \subset U_t \cap U_{t'}$, nous avons
(a) les images de $f_t$ et $f_{t'}$ dans $B$ diffèrent par une unité, et
(b) les complexes $M_t^\bullet \otimes_A B$ et $M_{t'} \otimes_A B$
définissent des objets isomorphes de $D(B)$. Pour $i \in \mathbf{Z}$, posons
$$
N_i = \max\nolimits_{t = 1, \ldots, T}
\left(\sum\nolimits_{j \geq i} (-1)^{j - i} rk(M_t^j)\right)
$$
Alors $N_t - \sum\nolimits_{j \geq i} (-1)^{j - i} rk(M_t^j) \geq 0$
et nous pouvons considérer les idéaux
$$
I_{t, i} =
f_t^{N_i - \sum\nolimits_{j \geq i} (-1)^{j - i} rk(M_t^j)}
I_i(M_t^\bullet, f_t)
\subset A_t
$$
Il résulte de Compléments d'algèbre, lemmes
\ref{more-algebra-lemma-eta-change-unit} et
\ref{more-algebra-lemma-ideal-well-defined},
que les idéaux $I_{t, i}$ se recollent en un idéal quasi-cohérent de type fini
$\mathcal{I}_i \subset \mathcal{O}_W$. De plus, tous ces idéaux sauf
un nombre fini sont égaux à $\mathcal{O}_W$.
Clairement, le morphisme $X' \to X$ construit ci-dessus
se restreint à l'éclatement de $W$ dans le produit
des idéaux $\mathcal{I}_i$. Ceci achève la preuve.
\end{proof}

\begin{lemma}
\label{flat-lemma-complex-and-divisor-blowup}
Dans la situation \ref{flat-situation-complex-and-divisor}, soit $b : X' \to X$
le morphisme du lemme \ref{flat-lemma-complex-and-divisor-blowup-pre}.
Considérons le diviseur de Cartier effectif $D' = b^{-1}D$ de faisceau d'idéaux
$\mathcal{I}' \subset \mathcal{O}_{X'}$. Alors $Q = L\eta_{\mathcal{I}'}Lb^*M$
est un objet parfait de $D(\mathcal{O}_{X'})$.
\end{lemma}

\begin{proof}
Cela résulte des lemmes \ref{flat-lemma-complex-and-divisor-blowup-pre}
et \ref{flat-lemma-good-triple}.
\end{proof}

\begin{lemma}
\label{flat-lemma-complex-and-divisor-blowup-base-change}
Dans la situation \ref{flat-situation-complex-and-divisor}, soit $h : Y \to X$
un morphisme de schémas tel que l'image inverse $E = h^{-1}D$
soit définie. Soient $b : X' \to X$, resp.\ $c : Y' \to Y$, construits dans le
lemme \ref{flat-lemma-complex-and-divisor-blowup-pre} pour
$D \subset X$ et $M$, resp.\ $E \subset Y$ et $Lh^*M$.
Alors $Y'$ est le transformé strict de $Y$ relativement à $b : X' \to X$
(voir la preuve pour une formulation précise) et
$$
L\eta_{\mathcal{J}'}L(h \circ c)^*M = L(Y' \to X')^*Q
$$
où $Q = L\eta_{\mathcal{I}'}Lb^*M$ est comme dans le
lemme \ref{flat-lemma-complex-and-divisor-blowup}.
En particulier, si $(Y, E, Lh^*M)$ est un bon triplet et si
$k : Y \to X'$ est l'unique morphisme tel que
$h = b \circ k$, alors $L\eta_\mathcal{J}Lh^*M = Lk^*Q$.
\end{lemma}

\begin{proof}
Notons $E' = c^{-1}E$. Alors $(Y', E', L(h \circ c)^*M)$ est un bon
triplet. Par conséquent, d'après la propriété universelle du
lemme \ref{flat-lemma-complex-and-divisor-blowup-pre},
il existe un unique morphisme
$$
h' : Y' \longrightarrow X'
$$
tel que $b \circ h' = h \circ c$. En particulier, il existe un morphisme
$(h', c) : Y' \to X' \times_X Y$. Nous affirmons que, si $W \subset X$ est un ouvert
quasi-compact tel que $b^{-1}(W) \to W$ soit un éclatement,
ce morphisme identifie $Y'|_W$ au transformé strict de $Y_W$
relativement à $b^{-1}(W) \to W$. Pour vérifier ce fait,
il suffit de raisonner localement sur $W$ ; nous pouvons donc démontrer
l'assertion sur une carte affine. C'est ce que nous faisons au paragraphe suivant.

\medskip\noindent
Soit $(U, A, f, M^\bullet)$ une carte affine pour $(X, D, M)$.
Rappelons que, d'après la preuve du lemme \ref{flat-lemma-complex-and-divisor-blowup},
la restriction de $b : X' \to X$ à $U$ est l'éclatement
de $U = \Spec(A)$ dans le produit des idéaux $I_i(M^\bullet, f)$.
Maintenant, si $V = \Spec(B) \subset Y$ est un ouvert affine avec $h(V) \subset U$,
alors $(V, B, g, M^\bullet \otimes_A B)$ est une carte affine pour
$(Y, E, Lh^*M)$, où $g \in B$ est l'image de $f$, voir le
lemme \ref{flat-lemma-pullback-triple-ideals-good}.
Ainsi, la restriction de $c : Y' \to Y$ à $V$
est l'éclatement dans le produit des idéaux
$I_i(M^\bullet, f)B$, c'est-à-dire que le morphisme $c : Y' \to Y$ au-dessus de
$h^{-1}(U)$ est l'éclatement de $h^{-1}(U)$ dans
l'idéal $\prod I_i(M^\bullet, f) \mathcal{O}_{h^{-1}(U)}$.
Puisque cela vaut aussi pour le transformé strict, nous voyons que
notre assertion sur les transformés stricts est vraie.

\medskip\noindent
Cela étant dit, l'égalité
$L\eta_{\mathcal{J}'}L(h \circ c)^*M = L(Y' \to X')^*Q$
résulte du
lemme \ref{flat-lemma-complex-and-divisor-eta-pull}.
La dernière assertion en est un cas particulier
(à savoir, le cas où $c = \text{id}_Y$ et $k = h'$).
\end{proof}

\begin{lemma}
\label{flat-lemma-complex-and-divisor-blowup-good}
Dans la situation \ref{flat-situation-complex-and-divisor}, soit $W \subset X$
le sous-schéma ouvert maximal sur lequel les faisceaux de cohomologie
de $M$ sont localement libres. Alors le morphisme $b : X' \to X$
du lemme \ref{flat-lemma-complex-and-divisor-blowup-pre} est un isomorphisme
au-dessus de $W$.
\end{lemma}

\begin{proof}
Cela est vrai car, pour toute carte affine $(U, A, f, M^\bullet)$
avec $U \subset W$, les idéaux $I_i(M^\bullet, f)$ sont localement
engendrés par une puissance de $f$ d'après
Compléments d'algèbre, lemme \ref{more-algebra-lemma-eta-cohomology-free}.
Puisque $f$ est un non-diviseur de zéro,
l'éclatement $b^{-1}(U) \to U$ est un isomorphisme.
\end{proof}

\begin{lemma}
\label{flat-lemma-complex-and-divisor-blowup-good-T}
Soient $X, D, \mathcal{I}, M$ comme dans la
situation \ref{flat-situation-complex-and-divisor}.
Si $(X, D, M)$ est un bon triplet, alors il existe une
immersion fermée
$$
i : T \longrightarrow D
$$
de présentation finie ayant les propriétés suivantes :
\begin{enumerate}
\item $T$ contient schématiquement $D \cap W$,
où $W \subset X$ est l'ouvert maximal sur lequel les
faisceaux de cohomologie de $M$ sont localement libres,
\item les faisceaux de cohomologie de $Li^*L\eta_\mathcal{I}M$
sont localement libres, et
\item pour tout point $t \in T$ d'image $x = i(t) \in W$, le rang
de $H^i(M)_x$ sur $\mathcal{O}_{X, x}$ et le rang
de $H^i(Li^*L\eta_\mathcal{I}M)_t$ sur $\mathcal{O}_{T, t}$ sont égaux.
\end{enumerate}
\end{lemma}

\begin{proof}
Soit $(U, A, f, M^\bullet)$ une carte affine pour $(X, D, M)$
telle que $I_i(M^\bullet, f)$ soit un idéal principal pour tout $i \in \mathbf{Z}$.
Nous définissons alors $T \cap U \subset D \cap U$ comme le sous-schéma fermé
défini par l'idéal
$$
J(M^\bullet, f) = \sum J_i(M^\bullet, f) \subset A/fA
$$
étudié dans Compléments d'algèbre, lemmes
\ref{more-algebra-lemma-eta-vanishing-beta-plus-pre} et
\ref{more-algebra-lemma-eta-vanishing-beta-plus} ;
d'après le second lemme, nous voyons que $T \cap U \to D \cap U$
est donné par l'homomorphisme d'anneaux $A/fA \to C$ qui y est étudié.
Puisque $(X, D, M)$ est un bon triplet, nous pouvons recouvrir $X$ par
des cartes affines de cette forme et, d'après le premier des deux
lemmes, cette construction se recolle. Nous obtenons donc un sous-schéma
fermé $T \subset D$ qui, sur les bonnes cartes affines ci-dessus,
est donné par l'idéal $J(M^\bullet, f)$. Les propriétés
(1) et (2) résultent alors du second lemme. Détails omis.
Petite observation pour aider le lecteur :
puisque $\eta_fM^\bullet$ est un complexe de modules localement libres
d'après Compléments d'algèbre, lemme \ref{more-algebra-lemma-eta-locally-free},
nous voyons que $Li^*L\eta_\mathcal{I}M|_{T \cap U}$
est représenté par le complexe $\eta_fM^\bullet \otimes_A C$
de $C$-modules. L'assertion (3) sur les rangs résulte de
Cohomologie, lemme \ref{cohomology-lemma-eta-cohomology-locally-free}.
\end{proof}

\begin{lemma}
\label{flat-lemma-complex-and-divisor-blowup-T}
Dans la situation \ref{flat-situation-complex-and-divisor}. Soient $b : X' \to X$
et $D'$ comme dans le lemme \ref{flat-lemma-complex-and-divisor-blowup-pre}. Soit
$Q = L\eta_{\mathcal{I}'}Lb^*M$ comme dans le
lemme \ref{flat-lemma-complex-and-divisor-blowup}.
Soit $W \subset X$ l'ouvert maximal sur lequel $M$ a des
modules de cohomologie localement libres.
Alors il existe une immersion fermée $i : T \to D'$ de présentation finie
telle que
\begin{enumerate}
\item $D' \cap b^{-1}(W) \subset T$ schématiquement,
\item $Li^*Q$ ait des faisceaux de cohomologie localement libres, et
\item pour $t \in T$ s'envoyant sur $w \in W$, le rang
de $H^i(Li^*Q)_t$ sur $\mathcal{O}_{T, t}$ soit égal au
rang de $H^i(M)_x$ sur $\mathcal{O}_{X, x}$.
\end{enumerate}
\end{lemma}

\begin{proof}
Le lemme \ref{flat-lemma-complex-and-divisor-blowup-good} nous dit que
$b$ est un isomorphisme au-dessus de $W$. Ainsi $b^{-1}(W) \subset X'$
est contenu dans l'ouvert maximal $W' \subset X'$ sur lequel $Lb^*M$
a des faisceaux de cohomologie localement libres. Les assertions du lemme elles-mêmes
sont alors une conséquence immédiate du
lemme \ref{flat-lemma-complex-and-divisor-blowup-good-T} appliqué
à $(X', D', Lb^*M)$ et des autres lemmes mentionnés dans l'énoncé.
\end{proof}

\begin{lemma}
\label{flat-lemma-complex-and-divisor-blowup-T-ranks}
Dans la situation \ref{flat-situation-complex-and-divisor}. Soient $b : X' \to X$,
$D' \subset X'$ et $Q$ comme dans le
lemme \ref{flat-lemma-complex-and-divisor-blowup}.
Soit $\rho = (\rho_i)_{i \in \mathbf{Z}}$ une famille d'entiers.
Soit $W(\rho) \subset X$ le sous-schéma ouvert maximal sur lequel
$H^i(M)$ est localement libre de rang $\rho_i$ pour tout $i$.
Soit $i : T \to D'$ comme dans le lemme \ref{flat-lemma-complex-and-divisor-blowup-T}.
Alors il existe un sous-schéma ouvert et fermé $T(\rho) \subset T$
contenant schématiquement $D' \cap b^{-1}(W(\rho))$
tel que $H^i(Li^*Q|_{T(\rho)})$ soit localement libre de rang $\rho_i$
pour tout $i$.
\end{lemma}

\begin{proof}
Soit $T(\rho) \subset T$ le sous-schéma ouvert et fermé
sur lequel $H^i(Li^*Q)$ est de rang $\rho_i$ pour tout $i$. L'assertion
résulte immédiatement de l'assertion du
lemme \ref{flat-lemma-complex-and-divisor-blowup-T}
sur les rangs des modules de cohomologie.
\end{proof}

\begin{lemma}
\label{flat-lemma-complex-and-divisor-blowup-T-represented-bdd-loc-free}
Dans la situation \ref{flat-situation-complex-and-divisor}. Soient $b : X' \to X$,
$D' \subset X'$ et $Q$ comme dans le
lemme \ref{flat-lemma-complex-and-divisor-blowup}.
S'il existe un complexe localement borné $\mathcal{M}^\bullet$
de $\mathcal{O}_X$-modules localement libres de rang fini représentant $M$,
alors il existe un complexe localement borné $\mathcal{Q}^\bullet$
de $\mathcal{O}_{X'}$-modules localement libres de rang fini représentant $Q$.
\end{lemma}

\begin{proof}
Rappelons que $Q = L\eta_{\mathcal{I}'}Lb^*M$, où $\mathcal{I}'$ est le
faisceau d'idéaux du diviseur de Cartier effectif $D'$.
Le complexe localement borné
$(\mathcal{M}')^\bullet = b^*\mathcal{M}^\bullet$ de
$\mathcal{O}_{X'}$-modules localement libres de rang fini
représente $Lb^*M$. Ainsi, le lemme résulte du
lemme \ref{flat-lemma-good-triple-bdd-loc-free}.
\end{proof}

\begin{lemma}
\label{flat-lemma-complex-and-divisor-derived}
Soit $X$ un schéma et soit $D \subset X$ un diviseur de Cartier effectif. Soit
$M \in D(\mathcal{O}_X)$ un objet parfait. Soit $W \subset X$ l'ouvert maximal
sur lequel les faisceaux de cohomologie $H^i(M)$ sont localement libres.
Il existe un morphisme propre $b : X' \longrightarrow X$
et un objet $Q$ de $D(\mathcal{O}_{X'})$ ayant les propriétés suivantes :
\begin{enumerate}
\item $b : X' \to X$ est un isomorphisme au-dessus de $X \setminus D$,
\item $b : X' \to X$ est un isomorphisme au-dessus de $W$,
\item $D' = b^{-1}D$ est un diviseur de Cartier effectif,
\item $Q = L\eta_{\mathcal{I}'}Lb^*M$, où $\mathcal{I}'$
est le faisceau d'idéaux de $D'$,
\item $Q$ est un objet parfait de $D(\mathcal{O}_{X'})$,
\item il existe une immersion fermée $i : T \to D'$ de présentation finie
telle que
\begin{enumerate}
\item $D' \cap b^{-1}(W) \subset T$ schématiquement,
\item $Li^*Q$ ait des faisceaux de cohomologie localement libres de rang fini,
\item pour $t \in T$ d'image $w \in W$, le rang
de $H^i(Li^*Q)_t$ sur $\mathcal{O}_{T, t}$ soit égal au
rang de $H^i(M)_x$ sur $\mathcal{O}_{X, x}$,
\end{enumerate}
\item pour toute carte affine $(U, A, f, M^\bullet)$ pour $(X, D, M)$,
la restriction de $b$ à $U$ est l'éclatement de $U = \Spec(A)$
dans l'idéal $I = \prod I_i(M^\bullet, f)$, et
\item pour toute carte affine $(V, B, g, N^\bullet)$ pour $(X', D', Lb^*N)$
telle que $I_i(N^\bullet, g)$ soit principal, nous avons
\begin{enumerate}
\item $Q|_V$ correspond à $\eta_gN^\bullet$,
\item $T \subset V \cap D'$ correspond à l'idéal
$J(N^\bullet, g) = \sum J_i(N^\bullet, g) \subset B/gB$
étudié dans
Compléments d'algèbre, lemme \ref{more-algebra-lemma-eta-vanishing-beta-plus}.
\end{enumerate}
\item Si $M$ peut être représenté par un complexe localement borné
de $\mathcal{O}_X$-modules localement libres de rang fini, alors $Q$ peut
être représenté par un complexe borné de
$\mathcal{O}_{X'}$-modules localement libres de rang fini.
\end{enumerate}
\end{lemma}

\begin{proof}
Cet énoncé rassemble les informations obtenues dans les
lemmes \ref{flat-lemma-pullback-triple-ideals-good},
\ref{flat-lemma-good-triple},
\ref{flat-lemma-complex-and-divisor-eta-pull},
\ref{flat-lemma-complex-and-divisor-blowup-pre},
\ref{flat-lemma-complex-and-divisor-blowup},
\ref{flat-lemma-complex-and-divisor-blowup-base-change},
\ref{flat-lemma-complex-and-divisor-blowup-good},
\ref{flat-lemma-complex-and-divisor-blowup-good-T},
\ref{flat-lemma-complex-and-divisor-blowup-T}, et
\ref{flat-lemma-complex-and-divisor-blowup-T-represented-bdd-loc-free}.
\end{proof}










\section{Éclatement de complexes, III}
\label{flat-section-blowup-complexes-III}

\noindent
Dans cette section, nous donnons une ``version algébrique'' de la
construction du graphe de Macpherson exposée dans \cite[section 18.1]{F}.

\medskip\noindent
Soit $X$ un schéma. Soit $E$ un objet parfait de $D(\mathcal{O}_X)$.
Soit $U \subset X$ le sous-schéma ouvert maximal tel que
$E|_U$ ait des faisceaux de cohomologie localement libres.

\medskip\noindent
Considérons le diagramme commutatif
$$
\xymatrix{
\mathbf{A}^1_X \ar[r] \ar[rd] &
\mathbf{P}^1_X \ar[d]^p &
(\mathbf{P}^1_X)_\infty \ar[l] \ar[ld] \\
& X \ar@/_1em/[ur]_\infty
}
$$
Rappelons ici que $\mathbf{A}^1 = D_+(T_0)$ est le premier ouvert
affine standard de $\mathbf{P}^1$ et que $\infty = V_+(T_0)$ est le
diviseur de Cartier effectif complémentaire, et que le diagramme ci-dessus
est l'image inverse de ces schémas sur $X$. Observons que
$\infty : X \to (\mathbf{P}^1_X)_\infty$ est un isomorphisme.
Alors
$$
(\mathbf{P}^1_X, (\mathbf{P}^1_X)_\infty, Lp^*E)
$$
est un triplet comme dans la situation \ref{flat-situation-complex-and-divisor}
de la section \ref{flat-section-blowup-complexes-II}.
Soient
$$
b : W \longrightarrow \mathbf{P}^1_X\quad\text{et}\quad
W_\infty = b^{-1}((\mathbf{P}^1_X)_\infty)
$$
l'éclatement et le diviseur de Cartier effectif construits
à partir de ce triplet dans le
lemme \ref{flat-lemma-complex-and-divisor-blowup-pre}.
Nous notons aussi
$$
Q = L\eta_{\mathcal{I}} Lb^*M = L\eta_\mathcal{I} L(p \circ b)^*E
$$
l'objet parfait de $D(\mathcal{O}_W)$
considéré dans le lemme \ref{flat-lemma-complex-and-divisor-blowup}.
Ici $\mathcal{I} \subset \mathcal{O}_W$ est le faisceau d'idéaux
de $W_\infty$.

\begin{lemma}
\label{flat-lemma-graph-construction}
La construction ci-dessus a les propriétés suivantes :
\begin{enumerate}
\item $b$ est un isomorphisme au-dessus de $\mathbf{P}^1_U \cup \mathbf{A}^1_X$,
\item la restriction de $Q$ à $\mathbf{A}^1_X$
est égale à l'image inverse de $E$,
\item il existe une immersion fermée $i : T \to W_\infty$
de présentation finie telle que $(W_\infty \to X)^{-1}U \subset T$
schématiquement et telle que $Li^*Q$ ait des faisceaux de cohomologie
localement libres,
\item pour $t \in T$ d'image $u \in U$, le
rang de $H^i(Li^*Q)_t$ sur $\mathcal{O}_{T, t}$ est égal au rang
de $H^i(M)_u$ sur $\mathcal{O}_{U, u}$,
\item si $E$ peut être représenté par un complexe localement borné de
$\mathcal{O}_X$-modules localement libres de rang fini, alors $Q$ peut être représenté
par un complexe localement borné de $\mathcal{O}_W$-modules localement libres de rang fini.
\end{enumerate}
\end{lemma}

\begin{proof}
Cela résulte immédiatement des résultats de la
section \ref{flat-section-blowup-complexes-II} ; pour un énoncé
qui rassemble tout ce qui est nécessaire, voir le
lemme \ref{flat-lemma-complex-and-divisor-derived}.
\end{proof}















\bibliography{my}
\bibliographystyle{amsalpha}
\end{document}
